% Options for packages loaded elsewhere
\PassOptionsToPackage{unicode}{hyperref}
\PassOptionsToPackage{hyphens}{url}
\PassOptionsToPackage{dvipsnames,svgnames,x11names}{xcolor}
\documentclass[
  letterpaper,
  11pt]{article}
\usepackage{xcolor}
\usepackage[margin=1.05in]{geometry}
\usepackage{amsmath,amssymb}
\setcounter{secnumdepth}{5}
\usepackage{iftex}
\ifPDFTeX
  \usepackage[T1]{fontenc}
  \usepackage[utf8]{inputenc}
  \usepackage{textcomp} % provide euro and other symbols
\else % if luatex or xetex
  \usepackage{unicode-math} % this also loads fontspec
  \defaultfontfeatures{Scale=MatchLowercase}
  \defaultfontfeatures[\rmfamily]{Ligatures=TeX,Scale=1}
\fi
\usepackage{lmodern}
\ifPDFTeX\else
  % xetex/luatex font selection
  \ifLuaTeX
    \usepackage{luaotfload}
    \directlua{luaotfload.add_fallback("mainfontfallback",{
      "DejaVuSans:mode=node"
    })}
  \fi
  \setmainfont[,RawFeature={fallback=mainfontfallback}]{Latin Modern Roman}
  \setmonofont[Scale=0.85]{DejaVu Sans Mono}
  \setmathfont[]{Latin Modern Math}
\fi
% Use upquote if available, for straight quotes in verbatim environments
\IfFileExists{upquote.sty}{\usepackage{upquote}}{}
\IfFileExists{microtype.sty}{% use microtype if available
  \usepackage[]{microtype}
  \UseMicrotypeSet[protrusion]{basicmath} % disable protrusion for tt fonts
}{}
\makeatletter
\@ifundefined{KOMAClassName}{% if non-KOMA class
  \IfFileExists{parskip.sty}{%
    \usepackage{parskip}
  }{% else
    \setlength{\parindent}{0pt}
    \setlength{\parskip}{6pt plus 2pt minus 1pt}}
}{% if KOMA class
  \KOMAoptions{parskip=half}}
\makeatother
\usepackage{longtable,booktabs,array}
\usepackage{caption}
\captionsetup[table]{skip=6pt}
\newcounter{none} % for unnumbered tables
\usepackage{calc} % for calculating minipage widths
% Correct order of tables after \paragraph or \subparagraph
\usepackage{etoolbox}
\makeatletter
\patchcmd\longtable{\par}{\if@noskipsec\mbox{}\fi\par}{}{}
\makeatother
\setlength{\emergencystretch}{3em} % prevent overfull lines
\providecommand{\tightlist}{%
  \setlength{\itemsep}{0pt}\setlength{\parskip}{0pt}}
\usepackage{mathtools}
\usepackage{amssymb}
\providecommand{\lt}{<}
\providecommand{\gt}{>}
\providecommand{\Bbb}{\mathbb}
\providecommand{\bold}{\mathbf}
\providecommand{\llbracket}{\lBrack}
\providecommand{\rrbracket}{\rBrack}
\usepackage{fancyhdr}
\pagestyle{fancy}\fancyhf{}\fancyhead[L]{\small Four-dimensional quantum Yang--Mills theory: the proof, outlined}\fancyhead[R]{\small\thepage}
\setlength{\headheight}{14pt}
\setlength{\emergencystretch}{3em}
\sloppy
\AtBeginDocument{\renewcommand{\setminus}{\mathbin{\backslash}}}
\usepackage{bookmark}
\IfFileExists{xurl.sty}{\usepackage{xurl}}{} % add URL line breaks if available
\urlstyle{same}
% fallback for those not using the hyperref driver hyperxmp:
\makeatletter
\@ifundefined{xmpquote}{\newcommand{\xmpquote}[1]{#1}}{}
\makeatother
\hypersetup{
  pdftitle={Four-dimensional quantum Yang--Mills theory: the proof{,} outlined (detailed notes)},
  colorlinks=true,
  linkcolor={[RGB]{0,0,114}},
  filecolor={Maroon},
  citecolor={Blue},
  urlcolor={[RGB]{0,0,114}},
  pdfcreator={LaTeX via pandoc}}

\title{Four-dimensional quantum Yang--Mills theory: the proof, outlined
(detailed notes)}
\usepackage{etoolbox}
\makeatletter
\providecommand{\subtitle}[1]{% add subtitle to \maketitle
  \apptocmd{\@title}{\par {\large #1 \par}}{}{}
}
\makeatother
\subtitle{A natural-language account of the Lean proof
(OpenAI-YM-20261028, commit 3456e5e)}
\author{}
\date{11 November 2026}

\begin{document}
\maketitle

{
\setcounter{tocdepth}{2}
\tableofcontents
}
\section{The proof in one reading}\label{the-proof-in-one-reading}

This is a substantial overview, about 10,000 words. The contents links
locate individual topics; the
\href{https://github.com/scottnarmstrong/yang-mills-informalization/blob/main/statement.md}{exact
statement} fixes the scope of every claim.

\subsubsection{The theorem}\label{the-theorem}

Let \(G\subset U(m)\) be a closed subgroup whose Lie algebra
\(\mathfrak g\) is simple and nonabelian, for instance \(SU(N)\) with
\(N\ge2\); \(G\) need not be connected. The main theorem is
\texttt{YangMills.jointExistence}; its precise statement is in
\href{https://github.com/scottnarmstrong/yang-mills-informalization/blob/main/statement.md}{statement.md}.
It asserts that there is one family of Wilson lattice gauge theories for
\(G\) on four-dimensional tori, consisting of:

\begin{itemize}
\tightlist
\item
  spacings \(a_N=1/D_N\downarrow0\) and physical side lengths
  \(L_p\uparrow\infty\);
\item
  bare couplings \(\beta_{N,p}\) with
  \(\beta_{N,p}/\log(1/a_N)\to b_\ast>0\) uniformly in \(p\);
\item
  insertions \(I_{N,p}(v,y)\) indexed by tokens \(v\), each a continuous
  gauge-invariant function of the links in a stencil of fixed size in
  lattice units around the site \(y\).
\end{itemize}

The family has the following property. For every strictly increasing
sequence \(p(k)\) of volume indices there are a further subsequence
\(p(q(k))\) and a continuum hierarchy satisfying (1)--(5) below.

\textbf{(1) Lattice origin.} Let
\(\mathrm{Conf}_n\subset(\mathbb R^4)^n\) be the open set of
configurations of distinct points. For every \(n\), every list of tokens
\(v\) and every \(f\in C_c^\infty(\mathrm{Conf}_n)\),

\[\lim_{k\to\infty}\,\lim_{N\to\infty}\,a_N^{4n}\sum_{y}f(\hat y)\,\Big\langle\prod_{i=1}^nI_{N,p(q(k))}(v_i,y_i)\Big\rangle_{N,p(q(k))}=\int f(x)\,S^{\mathrm{mix}}(\pi v;x)\,dx .\]

Here the sum runs over \(n\)-tuples of sites, \(\hat y\) are their
physical positions, and \(\langle\cdot\rangle_{N,p}\) is the Wilson
measure. - The inner limit runs over the whole sequence of spacings, at
every volume whose fundamental box contains the support of \(f\); this
holds for all large \(k\). - The test functions are joint functions of
the \(n\) points and vanish near all coincidences. - One subsequence
serves all \(n\), \(v\) and \(f\).

\textbf{(2) Phases.} On \(\mathrm{Conf}_n\),
\(S^{\mathrm{mix}}=\int_XS^\xi\,d\nu(\xi)\). Here \(X\) is a compact
metric space, \(\nu\) is a probability measure of full support, and
\(\xi\mapsto S^\xi(a;x)\) is continuous. All correlations \(S^s\),
\(s\in X\sqcup\lbrace\mathrm{mix}\rbrace\), are smooth and permutation
symmetric on \(\mathrm{Conf}_n\).

\textbf{(3) Reconstruction.} A selected phase \(\xi_\ast\) carries a
Wightman theory: - a strongly continuous unitary representation of the
proper orthochronous Poincaré group; - spectrum in the closed forward
cone; - a vacuum that is unique up to scalars among
translation-invariant vectors; - fields that are operator-valued
tempered distributions on the invariant domain spanned by field
monomials applied to the vacuum; - local commutativity and faithful
field labels.

Its Wightman functions are boundary values, on product test functions,
of the continuation of \(S^{\xi_\ast}\) that is holomorphic in the
ordered complex times.

\textbf{(4) Operator product expansion.} One system of coefficients,
with finitely many local fields of each degree, satisfies finite-order
OPE bounds in the mixed state and in every phase.

\textbf{(5) Nonfreeness.} Let \(e_4\) be the time direction. A field
\(\Theta\) of degree 4 with vanishing vacuum expectation satisfies

\[\lim_{r\downarrow0}\,r^8\log^2(1/r)\,G_\Theta(r)=\frac{3\dim_{\mathbb R}\mathfrak g}{\pi^4},\qquad G_\Theta(r)=S^{\xi_\ast}\big(\Theta(-\tfrac r2e_4)\,\Theta(\tfrac r2e_4)\big).\]

Moreover, \(G_\Theta\) agrees on no interval \((0,\delta)\) with a
finite Wick profile, that is, a finite sum of products of free
mass-shell kernels
\(\int E^q h(\mathbf p)(2E)^{-1}e^{-rE}\,d^3\mathbf p\).

The lattice data are fixed before the volume sequence. The refinement
and all continuum objects may depend on that sequence.

\textbf{The stronger statement.} The theorem is a projection of a longer
statement, \texttt{YangMills.constructed}. That statement holds for
every value \(b>0\) of the penalty parameter of the block-spin kernel;
the public theorem takes \(b=1\). The construction supplies the
following facts; the slope value is established in the proof of
regularity, rather than stated as a separate conjunct of
\texttt{YangMills.constructed}: - The slope is
\(b_\ast=\kappa_G=11C_{\mathfrak g}/(24\pi^2)\). - A token is a pair
\((d,c)\), with \(c\) in the space \(V_d\) of gauge-invariant
polynomials of weight at most \(d\) in the curvature and its covariant
derivatives (curvature has weight 2, and each derivative adds 1). - The
field spaces are \(F_d=V_d/N_d\), where \(N_d\) consists of the labels
whose mixed correlations with all separated spectators vanish. The label
space of the relativistic theory is the directed union of the \(F_d\),
so there is no finite field multiplet. - \(\Theta\) is the centred trace
of a constructed conserved stress tensor. - Every phase carries an
irreducible Wightman theory, with low-energy bounds that are uniform in
the phase. - The correlations extend to distributions across coincident
points; uniqueness is not asserted. - After finite-dimensional changes
of frame, the rescaled correlations converge at short distance to those
of \(\dim\mathfrak g\) free Maxwell fields. - Nested and integrated
forms of the OPE hold, as do further nonfreeness tests.

This overview describes the proof of the stronger statement.

\textbf{What is not claimed.} - No mass gap is proved; the available
spectral bounds in every phase are polynomial. - The infinite-volume
limit is not shown to be unique or independent of the subsequence. - The
phases are not shown to agree, and \(X\) is not shown to be a point. -
The mixed state is not shown to cluster. - The limit is not shown to be
independent of the lattice data. - The two limits are taken in a fixed
order. No diagonal limit, exchange of limits or rate of convergence is
asserted. - No phase is shown to be a limit of lattice states; the
lattice limit is the mixture. - The public theorem asserts nothing at
coincident points and no Euclidean axioms. The construction proves those
axioms internally. - There is no continuum measure on connections, no
Wilson-loop observable or area law, and no particle or scattering
theory. - Only finite Wick profiles are excluded. Generalized free
fields with a continuous mass distribution, infinite sums and general
composites of Gaussian fields are not excluded. - The OPE is a
finite-order bound and is not compared with perturbation theory.

\textbf{Normalizations.} On \(\mathfrak g\), put
\(\langle X,Y\rangle_{\mathfrak g}=-\tfrac1m\mathrm{Re}\,\mathrm{tr}(XY)\),
and define \(C_{\mathfrak g}>0\) by
\(\mathrm{tr}(\mathrm{ad}X\,\mathrm{ad}Y)=-C_{\mathfrak g}\langle X,Y\rangle_{\mathfrak g}\).
Two facts fix the coupling convention: -
\(1-\tfrac1m\mathrm{Re}\,\mathrm{tr}\,e^X=\tfrac12|X|_{\mathfrak g}^2+O(|X|^4)\);
- each unoriented plaquette is counted once in the Wilson action.

Hence \(\beta=1/g_0^2\) for the continuum action
\(\frac1{2g_0^2}\int\sum_{\mu<\nu}|F_{\mu\nu}|^2_{\mathfrak g}\). The
constructed couplings satisfy, uniformly in \(p\),

\[\frac{\beta_{N,p}}{\log(1/a_N)}\longrightarrow\kappa_G=\frac{11C_{\mathfrak g}}{24\pi^2}=2b_0,\qquad b_0=\frac{11C_{\mathfrak g}}{48\pi^2}.\]

For \(SU(N)\) in the defining representation, \(C_{\mathfrak g}=2N^2\)
and \(\kappa_G=11N^2/(12\pi^2)\). In that case \(\beta\) is the usual
Wilson coupling \(2N/g^2\), and the law above is the one-loop running
\(1/g^2\simeq(11N/24\pi^2)\log(1/a)\) of Gross--Wilczek and Politzer.

\subsubsection{The strategy}\label{the-strategy}

\textbf{Order of limits.} One lattice family is fixed first.
Correlations of renormalized insertions are taken to the limit
\(a_N\to0\) at fixed volume, along the whole sequence of spacings. The
volume then tends to infinity along a subsequence of the prescribed one.
The renormalization group (RG) is used only for the first limit. Its
levels run from the Wilson lattice (level \(N\)) to the torus of side
\(L_p\) at unit spacing (level \(0\)). The bare coupling is tuned so
that the effective coupling at level \(0\) is one fixed number
\(\bar\beta\), the same for all \(N\) and \(p\). No RG step is taken
above the unit scale; infinite volume is reached by compactness, from
local bounds that are uniform in the volume.

\textbf{Ultraviolet construction (Parts I--III).} At each level the
density is written as two pieces: - an explicit gauge-invariant action,
whose normalized curvature-squared coefficient is the running coupling
\(\beta_{N,p,m}\); - polymer activities: regular ones, controlled where
the curvature is small, and forced ones, which carry the large-field
regions.

One block-spin step integrates the fine links against an exactly
normalized, gauge-covariant probability kernel: - In the small-field
region this is a Gaussian integral around a stationary background field,
organized by a cluster expansion. - Large-field regions are kept and
paid for by a lower bound on the quadratic energy of the step. - The
marginal term is absorbed into the coupling. The coupling changes by the
one-loop amount \(\kappa_G\log L_m\) up to an error, and block ratios
that grow slowly along the levels make the accumulated error
\(o(\log(1/a_m))\).

The same steps transport inserted observables. This second layer
produces what later parts use: - chessboard bounds on renormalized
insertions and on their difference quotients, uniform in the volume; -
convergence at fixed volume, obtained from fixed-level convergence and
uniform observable tails. The construction also supplies analytic
estimates in a tuning parameter.

\textbf{Euclidean limit (Part IV).} - Chessboard bounds make the
correlations locally bounded, uniformly in the volume. - A
transfer-matrix representation and complex interpolation make them
equicontinuous, and Arzelà--Ascoli gives the subsequence. - Difference
quotients of insertions are close to insertions of higher degree. This
yields derivative maps on the field spaces and kernels that are smooth
on \(\mathrm{Conf}_n\). - The limit inherits translations, the signed
permutations of the axes and reflection positivity. It does not inherit
continuous rotations or clustering.

\textbf{Symmetry (Part VI).} - Lattice integration by parts gives shell
identities for a conserved tensor. Its normalization is fixed by
short-distance comparison with free Maxwell theory. - Reflection
positivity controls the flux at infinity, and analyticity extends the
Ward identities to all separated configurations. - The infinitesimal
rotations integrate to an \(O(4)\) action, again by comparison with free
Maxwell theory.

\textbf{Phases and reconstruction (Part V).} - The zero-energy subspace
of the Osterwalder--Schrader (OS) space of the mixed state carries a
commutative \(C^\ast\)-algebra; its proof of commutativity uses the
\(O(4)\) action. The characters of this algebra are the phases. - Each
phase is reconstructed as a Wightman theory, from bounds on operators
inserted in thin time slabs. Lorentz covariance comes from the Ward
identities. - A low-energy spectral bound, uniform in the phase and
derived from locality, upgrades vacuum uniqueness from almost every
phase to every phase.

\textbf{Identification (Part VII).} - Along sequences of scales tending
to zero, the theory converges, after changes of frame, to free Maxwell
theory. - The one-loop coefficient produces an inverse-logarithmic law
for the centered stress-tensor trace two-point function, and that law
excludes the specified finite Wick profiles. - The RG also yields the
OPE.

The logical order is not the order of the parts. Part VI uses the OS
spaces of the mixed state from Part V and the free comparison from Part
VII. Part V uses the \(O(4)\) action from Part VI. Part VIII assembles
everything for one family and one limit table.

\subsubsection{Part I. The lattice
theory}\label{part-i.-the-lattice-theory}

\emph{Unit:
\href{https://github.com/scottnarmstrong/yang-mills-informalization/blob/main/outline/units/I-1.md}{I.1
--- Lattice, symmetries, gauge groups and the Wilson measure}.}

On \(\Lambda_\ell=(\mathbb Z/\ell\mathbb Z)^4\), a configuration assigns
\(U(y,\mu)\in G\) to each positively oriented link. With \(h\) the
normalized Haar measure, the Wilson measure is

\[d\mu_{\ell,\beta}(U)=Z^{-1}\exp\Big[-\beta\sum_y\sum_{\mu<\nu}\big(1-\tfrac1m\mathrm{Re}\,\mathrm{tr}\,U_{y;\mu\nu}\big)\Big]\prod_{(y,\mu)}dh\big(U(y,\mu)\big).\]

The energy lies in \([0,12\ell^4]\), so this is a probability measure
for every real \(\beta\). It is invariant under gauge transformations,
translations and the signed permutations of the axes.

\textbf{Lattice schedule.} - Block ratios are odd integers \(L_m\ge3\),
with \(D_{m+1}=D_mL_m\) and \(a_m=1/D_m\). - The level-\(m\) torus at
volume \(p\) has \(L_pD_m\) sites per side, with \(4\mid L_p\). - Odd
ratios make the coarse lattice compatible with the signed permutations.
- Supports are measured by cover numbers by balls of radius \(B\), which
obey a coarsening inequality uniform in the volume.

\textbf{Group facts.} On \(G\) the unit proves: - a real-analytic
logarithm near \(1\); - an equivariant nearest-point retraction onto
\(G\), whose derivative is the orthogonal projection onto
\(\mathfrak g\); - a curvature map \(F(u)\in\mathfrak g\) with
\(|F(u)|\le\lVert u-1\rVert/\sqrt m\) and
\(F(vuv^{-1})=\mathrm{Ad}(v)F(u)\).

These are qualitative statements about single group elements;
quantitative control comes in Part II.

\textbf{Role of simplicity.} Simplicity enters through one fact: every
\(\mathrm{Ad}\)-invariant symmetric bilinear form on \(\mathfrak g\) is
a multiple of \(\langle\cdot,\cdot\rangle_{\mathfrak g}\). Together with
lattice symmetry, this leaves a single marginal coupling. The same fact
is used in the one-loop computation and in the uniqueness of the stress
tensor.

\textbf{Reflection positivity.} The Wilson measure is reflection
positive for every \(\beta\ge0\), for reflections through planes of
sites and through planes between sites (the argument of Osterwalder and
Seiler). This is proved in Part IV for the original lattice measure, not
for an effective action, and it is used only through the chessboard
inequality.

\subsubsection{Part II. Fields, insertions and polymer
activities}\label{part-ii.-fields-insertions-and-polymer-activities}

\emph{Units:
\href{https://github.com/scottnarmstrong/yang-mills-informalization/blob/main/outline/units/II-1.md}{II.1
--- Field labels and their lattice representatives};
\href{https://github.com/scottnarmstrong/yang-mills-informalization/blob/main/outline/units/II-2.md}{II.2
--- Polymer activities, norms, the Wilson family, the construction
record, reflection-positive cells}.}

\textbf{Labels.} In radial gauge, a connection near a point is described
by its Taylor coefficients, and the coefficient of order \(j\) has
weight \(j+2\). \(V_d\) is the space of \(\mathrm{Ad}(G)\)-invariant
complex polynomials of weight at most \(d\) in these coefficients.
Equivalently, it consists of the gauge-invariant polynomials of weight
at most \(d\) in the curvature and its covariant derivatives. It is
finite-dimensional, and \(1\in V_d\subset V_{d+1}\).

\textbf{Lattice representatives.} The basic lattice variables are
``slots'': plaquette curvatures transported to a base point along walks.
- For each \(d\), finitely many slots have linearizations with no common
kernel on the jets of weight at most \(d\). - The slot map then has an
analytic \(\mathrm{Ad}\)-equivariant left inverse near \(0\); it is
chosen non-constructively. - Compose a basis of \(V_d\) with the
order-\(d\) Taylor polynomial of this inverse, and average over the
signed permutations. For each label \(c\) this gives a gauge-invariant
polynomial \(\Phi_c\) in the link entries of a fixed box stencil.

Two properties follow: - On links obtained by parallel transport of
polynomial radial-gauge potentials at spacing \(h\), \(\Phi_c\)
reproduces the classical label to order \(h^d\). - Suppose the
plaquettes near the stencil have chord at most \(t\). Then a component
with \(f\ge1\) curvature factors, and its derivatives, are \(O(t^f)\).
Conversely, a function with such bounds determines a label with
uniformly bounded coefficients.

\textbf{Renormalized insertions.} The insertion of \(v\in V_d\) at a
site \(o\) is

\[O_{N,p}(v,o)=a_N^{-d}\,\Phi^{(d)}_{S_{1/\theta_N}\tau^{-1}v}(o).\]

\begin{itemize}
\tightlist
\item
  \(S_t\) multiplies the coefficient of a component with \(f\) curvature
  factors by \(t^f\).
\item
  \(\theta_N\) is the small-field threshold of the finest level.
\item
  \(\tau\) is a linear automorphism of \(V_d\). It commutes with
  conjugation and with the signed permutations, and it fixes the unit.
\end{itemize}

The RG defines \(\tau\). The running label coordinates of an insertion
obey a transfer relation from level \(N\) down to level \(m\), and
\(\tau\) is their value at a reference level \(r(d)\) fixed in advance.
So each insertion is normalized to have running coordinates exactly
\(v\) at level \(r(d)\). All labels of weight at most \(d\), including
the unit, may mix; no triangular structure is imposed. The large-field
part of the inserted density goes into reserved forced polymers, so the
construction never divides by the vacuum density.

\textbf{Equation relations.} Haar invariance gives the lattice
Schwinger--Dyson identity
\(\langle\langle A,\nabla_eS_W\rangle\,S\rangle_\beta=0\) for \(S\)
independent of the link \(e\). Products of such forces with slot
polynomials span the linearized free-Maxwell equation ideal, up to terms
of higher weight. They give lattice observables whose correlations with
distant spectators vanish exactly. Two facts are established later,
through the comparison with free theory: these relations become null in
the limit, and they exhaust \(N_d\).

\textbf{Polymers, budgets and marks.} - A polymer activity has a finite
support of sites. - Regular activities are anchored at a site and
measured by suprema of derivatives over small-field configurations. -
Forced activities carry the large-field regions, together with an
inventory of flagged sites. - Activities are controlled by budgets
summed through each site,
\(\sum_{X\ni x}e^{a\,\mathrm{cov}_B(X)}\lVert K_X\rVert\le\text{budget}\).
- Insertions carry marks: commuting variables \(\lambda\) with
\(\lambda^2=0\), one for each label or site that may occur. Exponentials
of marked quantities are finite sums, and logarithms are finite
recursions in the number of marks. - The weighted norm
\(\lVert a\rVert_w=\sum_x\prod_{\lambda\in x}w_\lambda\,|a(x)|\) is
submultiplicative, so bounds do not grow with the number of marks. - A
tree-graph bound without factorials controls connected clusters of
overlapping supports.

\subsubsection{Part III. The ultraviolet
construction}\label{part-iii.-the-ultraviolet-construction}

\paragraph{The block-spin step and the background
field}\label{the-block-spin-step-and-the-background-field}

\emph{Units:
\href{https://github.com/scottnarmstrong/yang-mills-informalization/blob/main/outline/units/III-1.md}{III.1
--- The block-spin step and the symmetric gauge fixing};
\href{https://github.com/scottnarmstrong/yang-mills-informalization/blob/main/outline/units/III-6.md}{III.6
--- The stationary branch: background field of the effective action};
\href{https://github.com/scottnarmstrong/yang-mills-informalization/blob/main/outline/units/III-3.md}{III.3
--- Free lattice kernels and covariances}.}

At level \(m\), relative to product Haar measure, the density is

\[\rho_m(U)=\exp\big(v|\Lambda_m|-\beta\,\mathcal K(U)-\mathcal L(U)-\mathcal E(U)\big)\,\mathcal G(U).\]

\begin{itemize}
\tightlist
\item
  \(\mathcal K\) is built from lattice paths: a plaquette term, smeared
  curvature squares, and cubic and quartic polynomials in transported
  curvatures. It is averaged over the signed permutations and normalized
  so that its quadratic symbol at zero momentum is the identity; this
  fixes the meaning of \(\beta\).
\item
  \(\mathcal L\) is a loop term independent of \(\beta\).
\item
  \(\mathcal E\) is the sum of the regular activities times small-field
  indicators.
\item
  \(\mathcal G\) is the gas of forced activities.
\item
  At the finest level, \(\rho_N=e^{-\beta_{N,p,N}S_W}\).
\end{itemize}

\textbf{The kernel.} One step is
\(\rho_m(V)=\int k_\beta(V,U)\,\rho_{m+1}(U)\,dU\). Given the fine
configuration \(U\), a finite tag \(t\) is drawn with probability
\(\lambda(U,t)\). Each coarse link is then drawn independently, with
Haar density

\[Z_\beta^{-1}\exp\Big(-\beta\kappa\big(1-\tfrac1m\mathrm{Re}\,\mathrm{tr}\,V_bQ_b(U,t)^{-1}\big)\Big),\]

where \(\kappa>0\) is the penalty. The predictor \(Q_b\) is built in two
steps: 1. On a patch of radius \((J+3)L\) around each coarse site, a
gauge frame is chosen. - If all route holonomies on the patch are close
to \(1\), the frame is the unique minimizer of an averaged holonomy
discrepancy. - Otherwise it is a uniformly drawn eligible axial gauge;
such gauges exist for every configuration. 2. Fine transports of length
\(L\), read in these frames, are combined by a nonlinear (minimizing)
average with a \(J\)-fold convolution weight, with a random fallback.

Properties of the kernel: - Haar invariance gives
\(\int k_\beta(V,U)\,dV=1\) exactly. - The kernel is exactly covariant
under fine gauge transformations and the lattice symmetries. - Hence
\(\int\rho_m\) equals the Wilson partition function at every level, and
the step is defined without any small-field assumption.

\textbf{Background field.} For the integration, the links of a spanning
forest are set to \(1\), and each remaining link is \(\exp(a)\) times a
fixed lift of the coarse field \(x\). The background is the critical
branch \(a=\rho(x)\) of the classical Taylor model of the action
(quadratic, cubic and quartic terms) plus the blocking penalty, at fixed
\(x\). The steps are: - The Hessian at the flat field is a positive
``bed'' form plus the penalty. - Each non-tree link is bounded by the
curvature in its block plus the response of the predictor. With this,
the Hessian is uniformly positive in the vertical directions, with a
constant independent of the volume and the forest. - The analytic
implicit function theorem then gives \(\rho\) as a locally unique
analytic germ with explicit linear part. Statements with a uniform
radius and localized jets give control independent of the volume. - The
branch is not shown to be a global minimizer, and this is not used. -
Changing the forest multiplies the Hessian determinant by the square of
a Jacobian and the Haar measure by the Jacobian itself. So the critical
value and the one-loop term do not depend on the forest.

\textbf{Free kernels.} The free covariances are written in Fourier space
through the actual averaging filters, which are entire and
multiplicative under composition of blockings. Aliased momenta carry
extra powers of \(1/L\). - All constants (decay rates, strip widths,
lower bounds) are fixed before the block ratio \(L\). - The quadratic
forms decay exponentially and are bounded below uniformly in the volume.
- An exact identity on each pair of tori splits a quadratic form into a
coarse part and a fluctuation part (a Schur complement).

\paragraph{Gaussian integration and the cluster
expansion}\label{gaussian-integration-and-the-cluster-expansion}

\emph{Units:
\href{https://github.com/scottnarmstrong/yang-mills-informalization/blob/main/outline/units/III-4.md}{III.4
--- Gaussian integration: change of Gaussian sector, determinants,
window law};
\href{https://github.com/scottnarmstrong/yang-mills-informalization/blob/main/outline/units/III-5.md}{III.5
--- Cluster expansion of the polymer gas};
\href{https://github.com/scottnarmstrong/yang-mills-informalization/blob/main/outline/units/III-8.md}{III.8
--- Connected activities and backward dependence cones}.}

\textbf{Change of variables.} Active links are written in exponential
coordinates around the background. Inactive links stay in Haar measure.
- The change of variables to a Gaussian integral with covariance \(C\)
and mean \(m(Y)\) is an exact identity. It holds for any
positive-definite \(C\), keeps all Jacobians, and uses no smallness
assumption. - The covariance actually used is the inverse of the
gauge-fixed Gram form on the active links, not the inverse of the Wilson
quadratic form. Its precision is bounded below uniformly in the volume.
- Window law: the global inverse is approximated by inverses on windows
of radius \(r\), with error \(Ce^{-vr}\) uniform in the volume. -
\(\log\det\) separates into local terms and far-field terms with
summable errors.

\textbf{Cluster expansion.} The gas is
\(\Xi(w)=\sum_s\prod_{i\in s}w_i\), summed over families of polymers
with pairwise disjoint supports and real weights. With
\(B=\lceil(\log\hat\beta)^2\rceil\), the smallness condition is

\[\sum_{i:\,x\in X_i}|w_i|\,e^{(a+\delta)\,\mathrm{cov}_B(X_i)}\le\gamma\ \ \text{for every site }x,\qquad 4(2B+1)^4\gamma/\delta\le q<1 .\]

Under this condition: - \(\log\Xi\) is an absolutely convergent sum of
connected clusters, bounded through spanning trees (Penrose's partition
scheme); - removing every polymer that meets a set \(Y\) changes
\(\log\Xi\) by at most \(|Y|\gamma/(1-q)\), uniformly in the volume.

So the vacuum logarithm is extensive, but a local change costs only its
size. A second, graph-free argument based on a deletion recursion gives
bounds of the same type.

\textbf{Connected activities and dependence cones.} - Fixed-order
connected coefficients come from the heat equation
\(\partial_tW=\tfrac12(\Gamma:\nabla^2W+\Gamma:\nabla W\,\nabla W)\) for
\(W_t=\log\mathbb E_{t\Gamma}e^V\). A rooted \(\ell^1\) norm with weight
\(e^{\sigma\,\mathrm{diam}}\) makes contraction and joining lossless, so
these coefficients are polynomially bounded uniformly in the volume.
There is no resummation. - Pulling a coarse observable back through the
block kernel gives a function of the fine links in a window of fixed
radius around its support. The backward cone therefore grows by a
bounded physical width per step. - For disjoint supports, the pull-back
factorizes exactly.

\paragraph{Large fields and one nonlinear
step}\label{large-fields-and-one-nonlinear-step}

\emph{Unit:
\href{https://github.com/scottnarmstrong/yang-mills-informalization/blob/main/outline/units/III-10.md}{III.10
--- One full nonlinear renormalization step: bounds on irrelevant
terms}.}

At reference coupling \(\hat\beta\) the small-field threshold is

\[t=\hat\beta^{-1/2}P(\hat\beta),\qquad P(\hat\beta)=\exp\big((\log\log\hat\beta)^2\big),\]

well above the Gaussian fluctuation scale \(\hat\beta^{-1/2}\).

\textbf{Seeds.} A seed is one of: - a fine plaquette with chord at least
\(t/4\); - a coarse link on which the coarse field departs from the
predictor by more than a fixed multiple of \(t\); - a flagged coarse
site.

\textbf{Energy lower bound.} For every store \(A\) of seeds,

\[\kappa\,t^2\,\sharp A\ \le\ \text{quadratic energy of the step integrand, after subtracting the stationary predictor}.\]

The energy is a sum over five channels, one of which is the blocking
penalty. The proof assigns charges to fine cells with bounded
multiplicity, and controls the linearization defects on cells without
seeds.

\textbf{Why this suffices.} Since \(\beta t^2\simeq P^2\), each seed
gains a factor \(e^{-cP^2}\). \(P^2\) grows faster than any power of
\(\log\hat\beta\), so this gain beats: - the polylogarithmic number of
connected seed sets through a site; - any per-seed loss of a fixed power
of \(\hat\beta\), or of order \(P^z\) with \(z<2\).

The decomposition into large-field sectors is exact: large-field sectors
are integrated, not discarded.

\textbf{One step.} The step takes inputs at level \(m+1\): regular norms
at most \(\hat\beta^{-11/10}\), forced norms at most \(e^{-P^{1/5}}\),
and the coupling in the band \([\hat\beta/2,2\hat\beta]\). For \(L\)
large, and then \(\hat\beta\) large, it produces: - site-summable
budgets for the new regular activities; - the bound
\(\exp\big((C\,\mathrm{bill}^5-\beta\mu t^2)\sharp A\big)\) for each
large-field sector, where \(\mathrm{bill}=1+B+(\log\hat\beta)^2\); -
bounds on the probabilities of the remaining catastrophic events.

The marginal and vacuum terms are extracted, and the level-\(m\) bounds
restored, by the reset below.

\paragraph{The marginal coupling and the
schedule}\label{the-marginal-coupling-and-the-schedule}

\emph{Units:
\href{https://github.com/scottnarmstrong/yang-mills-informalization/blob/main/outline/units/III-7.md}{III.7
--- Matching the Wilson energy to fourth order};
\href{https://github.com/scottnarmstrong/yang-mills-informalization/blob/main/outline/units/III-9.md}{III.9
--- The marginal flow: the running coupling and its reset};
\href{https://github.com/scottnarmstrong/yang-mills-informalization/blob/main/outline/units/III-2.md}{III.2
--- Choosing the constants: the order of choice and the RG schedule}.}

\textbf{Matching and the one-loop coefficient.} - At the finest level, a
path representation is constructed whose jets of order at most 4 agree
with those of the Wilson energy in exponential coordinates. - At later
levels, the cubic and quartic parts are represented in curvature
variables by quasi-local lifts (a cost-controlled Poincaré lemma). - In
each step, the coefficient \(\alpha_{N,m}\) of curvature squared
produced by the one-loop term satisfies
\(|\alpha_{N,m}+\kappa_G\log L_m|\le D\), with \(D\) fixed before \(L\).

The value of \(\kappa_G\) comes from an explicit momentum integral. The
comparison symbol is

\[h(k)=4C_{\mathfrak g}\Big(\frac{k_1^2}{|k|^6}-\frac{4k_1^2k_0^2}{|k|^8}-\frac1{|k|^4}\Big).\]

Its angular average is
\(4C_{\mathfrak g}(\tfrac14-\tfrac16-1)|k|^{-4}=-\tfrac{11}3C_{\mathfrak g}|k|^{-4}\),
and \(\int_{1/L<|k|\le1}|k|^{-4}\,d^4k=2\pi^2\log L\). Hence
\((2\pi)^{-4}\int_{1/L<|k|\le1}h=-\kappa_G\log L\). Comparing the step's
Hessian with \(h\) up to \(O(|k|^{-3})\) gives the bound on
\(\alpha_{N,m}\).

\textbf{The reset.} After the step, the blocked density is rewritten
exactly, at every configuration including large fields: - a new coupling
\(\beta_{N,p,m}=\beta_{N,p,m+1}+\alpha_{N,m}+\delta\); - a new vacuum
constant; - reorganized activities: flat values and isotropic quadratic
parts are subtracted, compensating path terms are added, and mismatches
of the small-field masks go into the forced gas.

The correction \(\delta\) is the isotropic Hessian coefficient of the
regular activities: - By simplicity and hyperoctahedral symmetry, their
quadratic part at a point is \(\delta\) times the identity on plaquette
orientations. - \(|\delta|\le1\), and \(\delta\) is fixed by a scalar
equation with a unique root.

Altogether, with \(\omega_j\) a drift allowance attached to regularity
index \(j\),

\[\big|\beta_{N,p,m+1}-\beta_{N,p,m}-\kappa_G\log L_m\big|\le\omega_{J_{m+1}} .\]

\textbf{Endpoint.} For real \(|b|<r_\ast\) there is a trajectory with
level-\(0\) coupling \(\bar\beta+b\). Its bare coupling is unique within
the construction, and the trajectory continues analytically to complex
\(|b|<r_\ast\), with \(r_\ast\) independent of \(N\) and \(p\). The
family uses \(b=0\), so \(\beta_{N,p,0}=\bar\beta\) for all \(N\) and
\(p\). An intermediate-value shooting argument also produces such
trajectories; the analytic version is the one used for convergence at
fixed volume below.

\textbf{Running.} Put \(s_m=\log(1/a_m)\) and
\(\hat\beta_m=\bar\beta+\kappa_Gs_m\). Induction from level \(0\) gives
\(|\beta_{N,p,m}-\hat\beta_m|\le\sum_{k\le m}\omega_{J_k}\), and the
schedule makes this \(o(s_m)\):

\[\beta_{N,p,m}=\bar\beta+\kappa_G\log(1/a_m)+o\big(\log(1/a_m)\big)\qquad\text{uniformly in }N\ge m\text{ and }p.\]

At \(m=N\) this is the slope law with \(b_\ast=\kappa_G\). Since
\(\bar\beta\) is chosen large, every level stays in the band
\([\hat\beta_m/2,2\hat\beta_m]\). The statement is leading-logarithmic
only: there is no two-loop term.

\textbf{Why the block ratios grow.} The one-loop error per step is
bounded but does not decrease with \(L\). For the accumulated error to
be \(o(s_m)\), the allowance per unit of logarithmic scale must tend to
zero. The schedule therefore imposes two conditions: -
\(\omega_{J_{m+1}}\le J_{m+1}^{-2}\log L_m\), which forces
\(\log L_m\to\infty\); - \(\log L_m/s_m\to0\), which keeps each step
small compared with the total scale.

The ratios are constant while the regularity index is fixed,
\(L_m=\ell_{J_{m+1}}\). The regularity index \(J\) is nondecreasing,
rises by at most one per level, tends to infinity toward the fine end,
and stays constant on long windows of levels. One sequence \((L_m)\)
serves every cutoff \(N\).

\textbf{Order of choices.} The constants are chosen in a fixed order: 1.
the kernel with its penalty, and for every index \(j\) its regularity
profile (derivative order, spatial power, support exponent, widths, path
counts, drift allowance), together with the one-loop error constant
\(D\); 2. all block ratios \(\ell_j\), so that the allowances and the
logarithms of the combinatorial costs are at most \(j^{-2}\log\ell_j\);
3. the smallness exponents; 4. the entry couplings and the volume
conditions; 5. \(\bar\beta\); 6. the assignment of regularity profiles
to levels.

Each degree \(d\) has a reference level \(r(d)\) with regularity index
at least \(d+3\). Beyond \(r(d)\) the derivative order exceeds \(d\) and
the spatial power is at least \(d\). All choices precede any observable
or extraction order, so every degree is served by the same family.

\paragraph{Renormalized observables, uniformly in the
volume}\label{renormalized-observables-uniformly-in-the-volume}

\emph{Units:
\href{https://github.com/scottnarmstrong/yang-mills-informalization/blob/main/outline/units/III-11.md}{III.11
--- Comparison with the free Gaussian measure: interacting rank and
determinant budgets};
\href{https://github.com/scottnarmstrong/yang-mills-informalization/blob/main/outline/units/III-12.md}{III.12
--- Bounds uniform in the volume and convergence as the spacing tends to
zero};
\href{https://github.com/scottnarmstrong/yang-mills-informalization/blob/main/outline/units/III-15.md}{III.15
--- Cost accounting for observables and the remainder comparison};
\href{https://github.com/scottnarmstrong/yang-mills-informalization/blob/main/outline/units/III-16.md}{III.16
--- Analytic toolkit}.}

\textbf{The chessboard root.} The quantity propagated for observables is
the chessboard root. For a function \(F\) of the links in one cell of a
tiling of the torus by cells of side \(a_m\),

\[\lVert F\rVert_T=\Big|\,\mathbb E_{\beta_{N,p}}\Big[\prod_{\text{cells}}\text{reflected copies of }F\Big]\Big|^{1/\sharp \text{cells}} .\]

The main bound reads as follows. For every degree \(d\), radius \(R\)
and \(\varepsilon>0\) there is \(A\) such that, for every \(m\ge r(d)\),
there is \(N_0\) independent of \(p\) with

\[\lVert O_{N,p}(v,o)\rVert_{T_m(o)}\le A\,a_m^{-d-\varepsilon}\qquad(N\ge N_0,\ \text{all }p,\ \lVert v\rVert\le R).\]

It is proved by running the insertion through the same RG steps;
\(\varepsilon\) absorbs losses smaller than any power. Reflection
positivity enters only afterwards, through the chessboard inequality.
That inequality makes the root a seminorm and bounds correlations of
separated observables by products of roots.

\textbf{No loss of \(\log\beta\) per site.} A root over \(c\) copies
tolerates a fixed cost per copy, but not a factor \(\log\beta\) per
site. The determinant budget avoids this by a gauge count on a terminal
torus with \(c=L^4\) sites: - A tree gauge leaves \(4c-(c-1)=3c+1\) free
links. - In the numerator, \(3c-3\) of them are integrated one at a
time, each being the only new link of a chosen plaquette. Each integral
gives at most \(b\,E^{\dim G}\), with \(E\propto g=\beta^{-1/2}\). The
four links that wrap around the torus, one per axis, are left over. -
The normalization is bounded below by confining all \(3c+1\) free links
to Haar balls of radius \(g/4\). - The ratio is therefore at most
\(\exp(c\,\pi-4\dim G\,\log g)\), where the cost \(\pi\) per site is
independent of \(g\). - Since the number of copies is \(c\), the root
carries only \(g^{-4\dim G/c}\).

\textbf{Field content.} Free Gaussian detectors give uniform rank
detection:
\(\lVert\pi x\rVert\le Z\,\lVert\mathrm{probe}(\pi x)\rVert\), where
\(\pi\) is the projection onto quotient labels. The limit field spaces
are isomorphic, degree by degree, to those of free Maxwell theory,
compatibly with conjugation and reflections.

\textbf{Difference quotients.} - A normalized lattice difference of an
insertion of degree \(d\) equals an insertion \(O(x)\) of degree
\(d+|l|\), up to a term whose root tends to zero. - Write
\(x=\pi x+(x-\pi x)\). The part \(x-\pi x\) lies in the Maxwell equation
ideal and has small root by the lattice equations of motion. The part
\(\pi x\) is bounded by the rank detector. - So difference quotients
satisfy root bounds of the same form. By the chessboard inequality,
lattice correlations with any number of differences are bounded at
\(r\)-separated points, uniformly in \(p\).

These bounds are propagated through the RG; they do not come from
differentiating a sup norm.

\textbf{A rank-free trace bound.} Let \(X\ge0\) be a matrix and let the
insertions \(b_a\) be separated by gaps of more than \(2r\) factors of
\(X\). Moving \(m\le r\) factors from one gap to another changes
\(\mathrm{tr}\prod_ab_aX^{Y_a}\) by at most \((\mathrm{cost}/r)\,m\).
The cost is a product of trace norms with no dimension factor. The proof
is a Hölder inequality for complex exponents, obtained from Hadamard's
three-lines theorem, followed by a Cauchy estimate. By quadrature it
extends to continuous positive kernels. Part IV uses it for
equicontinuity.

\textbf{Bookkeeping.} The numerical conditions of all steps can be met
in one fixed order: constants before \(L\), then \(L\), then the later
constants, then \(\beta\) large. This holds uniformly across regularity
profiles.

\paragraph{Convergence at fixed volume and the construction
run}\label{convergence-at-fixed-volume-and-the-construction-run}

\emph{Units:
\href{https://github.com/scottnarmstrong/yang-mills-informalization/blob/main/outline/units/III-12.md}{III.12};
\href{https://github.com/scottnarmstrong/yang-mills-informalization/blob/main/outline/units/III-13.md}{III.13
--- Extension to all labels: one common Wilson family};
\href{https://github.com/scottnarmstrong/yang-mills-informalization/blob/main/outline/units/III-14.md}{III.14
--- The construction run: flat representations and the ultraviolet
rate}.}

\textbf{The obstacle.} At fixed volume, two cutoffs \(N\) and \(N'\)
have different bare couplings, each fixed by the condition
\(\beta_{N,p,0}=\bar\beta\). Continuity in a fixed bare coupling
therefore compares nothing.

\textbf{Fixed-level convergence and tails.} The construction has a
one-parameter family of trajectories: - The level-\(0\) coupling is
\(\bar\beta+ws\), with a real tuning parameter \(s\in[-1,1]\). For each
\(s\) the bare coupling is unique within the recursion. - The bare
coupling and effective data have analytic dependence on the tuning
parameter. The estimates on that dependence are separate from the Cauchy
argument below. - For fixed \(j\), convergence of the trajectory and
stability of the observable transfers give convergence of the effective
insertions and their polymer parts. Uniform convergence passes through
the Haar integrals, and a positive lower bound on the normalized
partition function gives convergent correlation approximations. - The
omitted part is bounded, uniformly in the cutoff, by
\(2F_0e^{-B_j/64}\), where
\(B_j=\lceil(\log\hat\beta_j)^2\rceil\to\infty\). - Choose \(j\) to make
this tail small, then take the cutoffs large enough that their
level-\(j\) approximations are close. This proves the Cauchy property
first at lattice-fitted configurations; equicontinuity extends it to all
configurations distinct modulo the period. See
\href{https://github.com/scottnarmstrong/yang-mills-informalization/blob/main/outline/units/III-12.md\#45-convergence-in-the-mesh}{III.12,
§§4.5--4.6}.

Combined with the difference bounds, this gives at each fixed volume: -
limit kernels smooth away from periodic coincidences; - convergence of
difference quotients to derivatives; - convergence of smeared moments
along the whole sequence of spacings.

\textbf{The family used by the public theorem.}
\href{https://github.com/scottnarmstrong/yang-mills-informalization/blob/main/outline/units/VIII-1.md}{VIII.1}
assembles the construction for one family, using the bounds and
convergence results described in III.9--III.12.
\href{https://github.com/scottnarmstrong/yang-mills-informalization/blob/main/outline/units/III-13.md}{III.13}
describes an alternative coordinated family with root bounds, an
analytic shooting family and Maxwell rank comparisons. Its existence
statements do not identify that family with the one used by
\texttt{YangMills.constructed}; properties from the two constructions
should not be combined without such an identification.

\textbf{The construction run (III.14).} III.14 packages the construction
run. It also proves quantitative convergence of the order-zero quadratic
part of the effective action (its position-space Hessian kernel) at each
fixed level \(i\), at rate \((a_N/a_i)^c\). - The difference between two
cutoffs is written in curvature variables as \(d^{\mathsf T}Ad\), with
\(A\) defined modulo kernels annihilated by this map. - So the curvature
precision, which has zero modes, is never inverted. - This rate concerns
the quadratic germ only, not the full effective action or the final
limits.

\subsubsection{Part IV. Infinite volume}\label{part-iv.-infinite-volume}

\emph{Unit:
\href{https://github.com/scottnarmstrong/yang-mills-informalization/blob/main/outline/units/IV-1.md}{IV.1
--- Infinite volume: extraction of a subsequence and the Euclidean
axioms of the limit}.}

\textbf{Local bounds.} - Reflection positivity gives the chessboard
inequality \(|\langle\prod_cF_c\rangle|\le\prod_c\lVert F_c\rVert_T\)
for every \(\beta\ge0\). Log-convexity of roots under translation
handles cells not centred at the insertion. - Put each insertion in its
own cell of side \(a_m<r\). Then \(|C_{N,p}(v;x)|\le A^n\) whenever each
pair of points has some coordinate at distance at least \(r\), uniformly
in \(p\) and for large \(N\). - Near collisions, the same argument gives
\(|S(v;x)|\le Cr^{-4n(1+\max d_i)}\) when all points are at mutual
distance at least \(r\).

\textbf{Equicontinuity.} Let \(\mathrm{Gen}_n\) be the set of
configurations whose points have pairwise distinct coordinates on every
axis. 1. Integrating out the links along one axis gives a positive
transfer kernel \(K\). A correlation becomes
\(\mathrm{Tr}\prod_kB_kK^{Y_k}/\mathrm{Tr}\,K^M\), with slab kernels
\(B_k\) and gaps \(Y_k\). 2. Moving one point shifts units between two
gaps. The trace is entire in the size of the shift. 3. The rank-free
Hölder bound controls the trace on a disc, and a Cauchy estimate gives a
Lipschitz bound. 4. The trace norms of the slab kernels are themselves
chessboard roots of \(2q\) alternating copies, so the bound is uniform
in the volume.

The result is \(|C(x)-C(y)|\le16nA^n|x-y|/R\) on compact subsets of
\(\mathrm{Gen}_n\) whose coordinate gaps exceed \(R\). No derivative
bound is used.

\textbf{Extraction.} Arzelà--Ascoli is applied twice, over countably
many degrees, basis labels and compact sets: first to extract a spacing
subsequence that works for all volumes, then a volume subsequence. This
gives limit kernels on \(\mathrm{Gen}_n\). The complement
\(\mathrm{Conf}_n\setminus\mathrm{Gen}_n\) is a null set, and the
kernels are locally bounded on \(\mathrm{Conf}_n\), so dominated
convergence gives the smeared limits.

\textbf{Derivatives and smoothness.} - Exact summation by parts in the
lattice positions moves a difference quotient onto the test function. -
The ultraviolet estimates show that the lattice difference quotient of
an insertion is chessboard-close to an insertion of one degree higher. -
A dimension count, from the rank comparison, identifies the kernel of
the limiting projections with the null labels. - Finite-dimensional
compactness then gives fixed derivative maps
\(\partial_a:F_d\to F_{d+1}\) with

\[\partial_{x_i^a}S(\dots,v_i,\dots;x)=S(\dots,\partial_av_i,\dots;x)\qquad\text{in }\mathcal D'(\mathrm{Conf}_n).\]

Iteration and a Fourier argument give
\(S\in C^\infty(\mathrm{Conf}_n)\), with derivatives polynomially
bounded near collisions.

\textbf{What the limit inherits.} The limit inherits: - translations; -
the group of signed permutations of the axes, acting on points and
labels; - permutation symmetry, reality and the unit label; - reflection
positivity at coordinate planes, for point and smeared tests.

It does not inherit continuous rotations, clustering or uniqueness.

\textbf{Removing the spacing refinement.} The extraction refines both
the spacing and the volume sequences. At each admissible fixed volume,
the smeared limit along the whole spacing sequence exists (Part III) and
agrees with the subsequential one. This removes the spacing refinement
from the final statement.

\subsubsection{Part V. Phases and relativistic
reconstruction}\label{part-v.-phases-and-relativistic-reconstruction}

\emph{Units:
\href{https://github.com/scottnarmstrong/yang-mills-informalization/blob/main/outline/units/V-1.md}{V.1
--- Decomposition of the limit into phases};
\href{https://github.com/scottnarmstrong/yang-mills-informalization/blob/main/outline/units/V-2.md}{V.2
--- Osterwalder--Schrader reconstruction of a Wightman theory};
\href{https://github.com/scottnarmstrong/yang-mills-informalization/blob/main/outline/units/V-3.md}{V.3
--- Reconstruction for every phase and the reconstructed theory}.}

\textbf{The OS space of the mixture.} Positive-time field words, modulo
the null space of \(\langle\theta w,w'\rangle_S\) and completed, form a
Hilbert space. It carries: - a contraction semigroup \(e^{-tH}\) with
\(H\ge0\); - a unitary representation of the spatial translations.

The joint spectrum lies in a half-space. No positive spectral gap is
asserted. The zero-energy space may be large. Its projection \(\Pi\) is
the strong limit of \(e^{-tH}\) as \(t\to\infty\).

\textbf{Thin slabs.} A word \(w\) placed in a thin time slab defines an
operator \(B\) on the OS space. To bound it, consider the train
\(v_{n+1}=B^\ast Bv_n\), which contains \(n\) reflected copies of \(w\)
and an arbitrary exterior word. - The lattice bounds give
\(\lVert v_n\rVert^2\le B_vA^{4n}\). Here \(A\) is polynomial in the
word data and in the inverse buffer width, and \(B_v\) depends on the
exterior word but not on \(n\). - A power lemma removes \(B_v\): if
\(\lVert(BA)^nv\rVert\le K_vD^{2n}\) for all \(n\), then
\(\lVert Av\rVert\le D\lVert v\rVert\).

So the exterior word is paid once and disappears under the \(n\)-th
root.

\textbf{Phases.} - The compressions \(M(w)=\Pi\,B_w|_{\ker H}\) are
independent of the width and placement of the slab. -
\(\langle\Omega,M(w_1)\cdots M(w_k)\Omega\rangle\) is the limit of \(S\)
as the clusters are moved apart in time order.

The compressions commute: 1. Place four clusters on the rays
\(n(ue_4+ve_1)\), with \((u,v)=(0,0),(1,2),(2,1),(3,3)\). 2. Exchanging
the time axis with a spatial axis reverses the order of the middle pair.
3. \(O(4)\) covariance and permutation symmetry of \(S\) then give
\(\langle\Omega,M_xM_yM_wM_z\Omega\rangle=\langle\Omega,M_xM_wM_yM_z\Omega\rangle\).
4. By density, \([M_y,M_w]=0\).

The resulting phase decomposition: - The \(M(w)\) generate a commutative
unital \(C^\ast\)-algebra. Its Gelfand spectrum is \(X\), which is
compact and metrizable. - \(\nu\) represents the vacuum state, and it
has full support because the vacuum is separating. - The phase
correlations are \(S^\xi(w)=\xi(M(w))\). They are continuous in \(\xi\),
and \(\int S^\xi\,d\nu=S^{\mathrm{mix}}\). - Reflection positivity,
\(O(4)\) invariance, the derivative and stress identities and uniform
growth bounds hold for every \(\xi\). They are proved as operator
identities before any character is applied. - Every phase has the same
null spaces as the mixture. - A variance argument over a countable total
family of words gives a one-dimensional zero-energy space for
\(\nu\)-almost every phase. - The selected phase \(\xi_\ast\) is one
such phase, chosen by this existence argument; no further condition is
imposed. The trace law of Part VII is proved for the phase so chosen.

\textbf{Reconstruction.} For each phase, slab operators bounded by
\(C\,p^{-A}\) (with \(p\) the slab half-width) define

\[F_n(z;\mathbf x)=\big\langle\Omega,\,B(v_1)\,e^{i(z_1-z_2)H}U(\mathbf x_1-\mathbf x_2)\,B(v_2)\cdots\Omega\big\rangle .\]

This function is holomorphic for
\(\mathrm{Im}\,z_1<\dots<\mathrm{Im}\,z_n\), and it equals the Euclidean
correlation at imaginary times. The Wightman theory is obtained in five
steps: 1. \textbf{Boundary values.} Bounds of order \(k^{-m}\), when all
height gaps exceed \(k\), give tempered boundary values \(W_n\). They
converge against every joint Schwartz test as the vector of heights
tends to zero through all ordered vectors. 2. \textbf{Positivity and
fields.} This mode of convergence allows two independently regularized
vectors to be compared. So \(W\) is a limit of Gram matrices of OS
vectors: it is positive, and the Wightman space embeds isometrically
into the OS space. Fields act on finite words modulo the null space by
prepending a letter. This gives the invariant domain, the formal
adjoints and the cyclic vacuum. 3. \textbf{Covariance and spectrum.} The
rotation Ward identities of Part VI, in planes containing the time axis,
continue to infinitesimal boost identities. These integrate to Lorentz
covariance without any spectral assumption. The spectral measure of a
Wightman vector equals that of its OS image, so it is supported in
\(p^0\ge0\); covariance then moves the support into the forward cone. 4.
\textbf{Locality.} Only adjacent transpositions are needed. For a
spacelike adjacent pair, both orderings have holomorphic germs,
constructed with complex boosts. The germs agree to all orders on the
Euclidean slice, because the Euclidean kernel is symmetric and smooth
away from coincidences. 5. \textbf{Faithfulness.} A label whose Wightman
functions vanish has a vanishing tube function, hence a vanishing
Euclidean kernel.

\textbf{Every phase.} The low-energy bound uses only locality and the
spectrum condition. Take \(A\) a finite sum of field words with
compactly supported tests. - The commutator of time-translated copies of
\(A\) vanishes for times at most \(1\) and spatial separations beyond a
radius \(R\) independent of \(\xi\). - Average over unit time, pair with
a spatial-momentum bump of radius \(2u\) (whose Fourier transform is
\(O(u^3)\) in three space dimensions), and use \(|\mathbf p|\le E\) on
the forward cone. This bounds the spectral weight of each shell
\(u/2<E\le u\) by \(O(u^2)\), and dyadic summation gives the first bound
below.

\[\mu_{A\Omega_\xi}(0<E\le u)\le M_Au^2\ \ (0<u\le1),\qquad 0\le\big\langle A\Omega_\xi,(e^{-sH_\xi}-P_{0,\xi})A\Omega_\xi\big\rangle\le M_As^{-2}\ \ (s\ge1).\]

The constant \(M_A\) is independent of \(\xi\), by compactness of \(X\).

Vacuum uniqueness then extends from almost every phase to every phase: -
The function
\(\xi\mapsto\langle A\Omega_\xi,e^{-sH_\xi}A\Omega_\xi\rangle-|\langle\Omega_\xi,A\Omega_\xi\rangle|^2\)
is continuous. - For fixed \(A\) and \(s\), the set of phases where this
function lies in \([0,M_As^{-2}]\) is closed by continuity. It contains
the full-measure set of phases with unique vacuum; full support of
\(\nu\) makes this closed set all of \(X\). - Letting \(s\to\infty\)
gives a unique vacuum in every phase.

The comparison map from the Wightman space to the OS space is
surjective, proved independently by pushing time-ordered vectors apart.
This transfers vacuum uniqueness to the OS space. Each phase theory is
irreducible. Its Wightman distributions are continuous in \(\xi\), and
their \(\nu\)-average is the boundary value of the continuation of
\(S^{\mathrm{mix}}\). The bounds allow spectrum down to zero and give
only polynomial decay \(s^{-2}\) in Euclidean time, not exponential
clustering.

\subsubsection{Part VI. Ward identities, rotations and the stress
tensor}\label{part-vi.-ward-identities-rotations-and-the-stress-tensor}

\emph{Units:
\href{https://github.com/scottnarmstrong/yang-mills-informalization/blob/main/outline/units/VI-1.md}{VI.1
--- Ward identities, the rotation action and the conserved stress
tensor};
\href{https://github.com/scottnarmstrong/yang-mills-informalization/blob/main/outline/units/VI-2.md}{VI.2
--- The free Maxwell stress tensor and the local Ward charge}.}

The lattice has only the hypercubic symmetries. The proof obtains the
generators of rotations from stress fluxes, and the group action by
comparison with free Maxwell theory. It uses the OS spaces of the mixed
state in the coordinate directions and the free comparison of Part VII.
It uses no decay of correlations.

\textbf{Free Maxwell theory.} Consider \(\dim\mathfrak g\) free Maxwell
fields, with field spaces the \(G\)-invariant polynomials in on-shell
curvature jets. - The centred, real, symmetric tensors of degree 4 that
are covariant under signed permutations form a three-dimensional space.
- For simple \(\mathfrak g\), the divergence maps this space onto a
two-dimensional space. So the conserved tensors form a line, spanned by
the Maxwell stress tensor. - For every Killing field \(v\), the flux of
\(\sum_bv_bT_{ab}\) through a shell around an insertion equals the Lie
derivative of the insertion. For elementary fields this follows from
\(-\Delta G_0=\delta\); for polynomials, from the Leibniz rule and Wick
contraction. - An explicit three-point function fixes the sign and
normalization of the Ward constant.

\textbf{Conserved tensor.} There are scaling sequences
\(r_k\downarrow0\) with frames \(U_k\), isomorphisms for large \(k\).
Along them the rescaled interacting divergence converges to the free
one. - Surjectivity is an open condition, so the interacting divergence
is also onto, and its kernel is one-dimensional. - A nonzero conserved
\(J\) in this kernel converges, after rescaling, to a multiple of the
free stress tensor.

\textbf{Shell identities.} Lattice integration by parts in the link
variables, with ultraviolet control of the error terms, gives shell
identities. Take \(\chi\) equal to \(1\) near \(x_0\) and vanishing near
the spectators:

\[-\sum_a\int\partial_a\chi(y)\,S\big(J_{ab}(y)\,B(x_0)\cdots\big)\,dy=S\big((p_bB)(x_0)\cdots\big).\]

With weight \(K(y-x_0)\), for a rotation generator \(K\), the same
construction gives maps \(Y_K\). Here \(p_b\) and \(Y_K\) are linear
maps on labels, not yet identified with derivatives and rotations.
Moving the centre of the shell gives \(\partial p=p\,\partial\) and
\(\partial_bY_K+\sum_cK_{cb}\,p_c=Y_K\partial_b\).

\textbf{Flux at infinity.} - For rotations, cut off the rotation field
so that it is compactly supported and equal to the exact field on a box
far larger than the configuration. - Conservation and symmetry cancel
the cross terms. - The remaining exterior contribution is bounded
through spectral measures of the OS semigroups along and transverse to a
coordinate line, and it tends to zero. - This gives the angular identity
for configurations on one coordinate line, together with a line-sum
identity, and hence transversality of the label maps.

\textbf{Normalization.} In free theory, the finitely many relations
satisfied by \(p\) and \(\partial\) at a given degree have a
one-dimensional solution space. Injectivity is an open condition, so the
same holds along the scaling sequence. Hence \(p_b=c\,\partial_b\), with
\(c\) independent of the degree and \(c\ne0\). The tensor \(T=J/c\)
satisfies

\[-\sum_a\int\partial_a\chi(y)\,S^s\big(T_{ab}(y)\,B(x_0)\cdots\big)\,dy=\frac{\partial}{\partial x_0^b}\,S^s\big(B(x_0)\cdots\big)\]

in the mixed state. By the operator identities of Part V, it holds in
every phase. Since \(\partial\chi\) vanishes near every insertion, no
contact terms arise.

\textbf{All separated configurations.} Positivity of Gram functions,
together with the Ward values already obtained, gives analytic-vector
bounds \(\lVert f^{(m)}\rVert\le C(\rho/2)^{-m}m!\). The angular Ward
functional is therefore real-analytic along walks that move one point at
a time. Starting from configurations on a line, it vanishes on all
configurations whose points have pairwise distinct first coordinates,
and by continuity on all of \(\mathrm{Conf}_n\).

\textbf{Integration to \(O(4)\).} - The bracket defect
\([Y'_u,Y'_v]-Y'_{[u,v]}\) (with \(Y'=Y/c\)) satisfies a summed
identity. A rigidity statement, proved in free theory and transferred by
the comparison, forces the defect to vanish. - The infinitesimal action
converges to the Maxwell one:
\(U_k^{-1}Y'_KU_k\to\mathrm{Lie}^{\mathrm{Max}}_K\). - A perturbation of
spectral projections of the Casimir operators gives exact intertwiners
\(V_d\). The action on \(F_d\) is
\(R_d(k)=V_dR^{\mathrm{Max}}_d(k)V_d^{-1}\). - Since
\(t\mapsto S(R(e^{tK})v;e^{tK}x)\) has zero derivative, and reflections
act through the real structure, \(S(R(k)v;kx)=S(v;x)\) for all
\(k\in O(4)\). - Averaging \(J/c\) over \(O(4)\) gives a conserved,
real, symmetric, covariant tensor that satisfies both the translation
and the rotation Ward identities.

\subsubsection{Part VII. Identification as Yang--Mills
theory}\label{part-vii.-identification-as-yangmills-theory}

\paragraph{Free Maxwell fields at short
distance}\label{free-maxwell-fields-at-short-distance}

\emph{Units:
\href{https://github.com/scottnarmstrong/yang-mills-informalization/blob/main/outline/units/VII-1.md}{VII.1
--- Free Maxwell comparison theory};
\href{https://github.com/scottnarmstrong/yang-mills-informalization/blob/main/outline/units/VI-2.md}{VI.2
--- Interacting stress comparison};
\href{https://github.com/scottnarmstrong/yang-mills-informalization/blob/main/outline/units/VII-6.md}{VII.6
--- Scaling variants and fixed-field estimates}. The comparison is
assembled in
\href{https://github.com/scottnarmstrong/yang-mills-informalization/blob/main/outline/units/VIII-1.md}{VIII.1}.}

\textbf{The reference theory.} It consists of \(\dim\mathfrak g\) free
Maxwell fields built from \(\Gamma(x)=1/(4\pi^2|x|^2)\). - On the
lattice, corrected covariances of finite-difference jets converge to
\(\partial_a\partial_b\Gamma\) away from coincidences. - Wick moments
satisfy the same Gaussian recursion on the lattice and in the continuum,
so correlations of quotient labels converge. - A triangular basis built
from the Kelvin transform gives strict reflected Gram positivity for
every label outside the ideal \(I_d\) generated by the free Maxwell
equations and the Bianchi identities. - So the free null space is
exactly \(I_d\), and the free field space \(M_d=P_d/I_d\), where \(P_d\)
is the corresponding space of polynomials in free curvature jets, is
detected by finitely many correlations.

\textbf{The comparison.} There are scales \(r_k\downarrow0\) and linear
frames \(U_k\) on the field spaces with the following property. The
correlations of the mixed state, rescaled by \(r_k\) and transformed by
\(U_k\), converge to the free Maxwell correlations: - with all
derivatives; - uniformly on compact subsets of \(\mathrm{Conf}_n\); -
for all degrees along one sequence.

The stress tensor converges to the Maxwell stress tensor. An upper bound
on \(\dim F_d\) makes the frames isomorphisms for large \(k\), so
\(F_d\cong M_d\).

The frames are necessary: the claim is not that a fixed normalization of
each composite field has a limit along every sequence of radii. The
schedule contains a real parameter that moves the scale at which
observables are read, and the estimates are uniform in it. Tuning this
parameter, every sequence of radii tending to zero has a subsequence
along which the comparison holds.

\paragraph{Running coupling}\label{running-coupling}

\emph{Unit:
\href{https://github.com/scottnarmstrong/yang-mills-informalization/blob/main/outline/units/VII-5.md}{VII.5
--- Asymptotic freedom of the running coupling}.}

With \(g_{N,p}(m)=\beta_{N,p,m}^{-1/2}\), Part III gives

\[g_{N,p}(m)\,\sqrt{\log(1/a_m)}\longrightarrow\frac1{\sqrt{2b_0}}\qquad(m\to\infty),\]

uniformly in \(N\ge m\) and \(p\). This is the leading-logarithmic
running of the coupling of the construction's effective actions. No
two-loop term, no \(\Lambda\) parameter and no comparison with another
renormalization scheme is claimed. For \(SU(N)\), the conventionally
normalized coupling \(g_{\mathrm{std}}=\sqrt{2N}\,g\) satisfies
\(g_{\mathrm{std}}\sqrt{\log(1/a_m)}\to(11N/24\pi^2)^{-1/2}\).

\paragraph{The trace law and
nonfreeness}\label{the-trace-law-and-nonfreeness}

\emph{Units:
\href{https://github.com/scottnarmstrong/yang-mills-informalization/blob/main/outline/units/VII-2.md}{VII.2
--- The trace of the stress tensor and the trace anomaly};
\href{https://github.com/scottnarmstrong/yang-mills-informalization/blob/main/outline/units/VII-3.md}{VII.3
--- The theory is not free: Gram positivity and exclusion of free-field
profiles}.}

\textbf{Leading form of the trace.} The classical Maxwell and
Yang--Mills stress tensors are traceless. The quantum trace comes from
the scale dependence of the coefficient of curvature squared. - The
proof pulls the trace of the lattice stress observable back through the
RG steps. - Fourth-order matching identifies the leading part of the
trace as the derivative of the effective action with respect to scale. -
The one-loop coefficient enters through a momentum-shell estimate of the
kind used in Part III.

Along the scaling sequences, with \(\beta_k\) the running coupling at
the RG level that reads scale \(r_k\), and \(O_{F^2}\) the normalized
curvature-squared label,

\[\beta_k\,r_k^4\,U_k^{-1}\Theta\longrightarrow4d_2\,O_{F^2},\qquad d_2=-\frac{\kappa_G}4,\qquad\frac{\log(1/r_k)}{\beta_k}\longrightarrow\frac1{\kappa_G}.\]

In perturbative terms, this is the familiar
\(\Theta\simeq\frac{\beta(g)}{2g^3}F^2\), with
\(\beta(g)/(2g^3)=-b_0/2=d_2\) at one loop. The trace is a
curvature-squared field suppressed by
\(g^2(r)\simeq1/(\kappa_G\log(1/r))\).

\textbf{The two-point law.} Here
\(\Theta=\sum_\mu T_{\mu\mu}-S^{\xi_\ast}\big(\sum_\mu T_{\mu\mu}\big)\mathbf1\),
in the selected phase, which has a unique vacuum. 1. The free two-point
function of \(O_{F^2}\) is \(3\dim\mathfrak g/(\pi^4|x|^8)\). 2. So at
lattice-fitted configurations along the scaling sequence,
\(r_k^8\log^2(1/r_k)\langle\Theta\Theta\rangle\to(4d_2/\kappa_G)^2\cdot3\dim\mathfrak g/\pi^4=3\dim\mathfrak g/\pi^4\).
3. Rotation invariance moves the pair onto the time axis. 4. The
spectral representation in the phase gives
\(G_\Theta(r)=\int e^{-rE}\,d\mu(E)\) with a positive measure. So
\(G_\Theta\) is nonincreasing, which sandwiches \(G_\Theta(r)\) between
its values at fitted points. 5. If the limit failed along some sequence
of radii, the flexible scaling above would give a subsequence along
which it holds, a contradiction.

Hence the limit exists over all radii:

\[\lim_{r\downarrow0}\,r^8\log^2(1/r)\,G_\Theta(r)=\frac{3\dim_{\mathbb R}\mathfrak g}{\pi^4}.\]

The amplitude does not depend on the value of the one-loop coefficient.
\(\kappa_G\) enters once through the trace and once through the running,
and \((4d_2/\kappa_G)^2=1\) identically. So the test detects the
inverse-logarithmic suppression and the free normalization of \(F^2\),
not the value of \(b_0\); that value is established by the running law
of Part III.

\textbf{Exclusion of free profiles.} Consider a finite Wick profile
\(W(r)=\sum_\alpha c_\alpha\prod_\beta K_{\alpha\beta}(r)\). Each factor
\(K(r)=\int E^q\,h(\mathbf p)(2E)^{-1}e^{-rE}\,d^3\mathbf p\) has
\(E=\sqrt{\mathsf m^2+|\mathbf p|^2}\) and \(h\) homogeneous. - \(W\)
has an asymptotic expansion at \(r=0\) in the functions
\(r^a(\log r)^k\), with \(a\in\mathbb Z\) and \(k\ge0\). - Along
\(r=2^{-n}\), \(r^8\log^2(1/r)\,W(r)\) is then a finite sum of terms
\(c\,\rho^n n^{k+2}\), which can converge only to \(0\). - Since the
limit for \(G_\Theta\) is positive, \(G_\Theta\) differs from \(W\) at
points arbitrarily close to \(0\).

This excludes exactly the finite class of profiles in the statement, and
it concerns only \(G_\Theta\) on the time axis.

The stronger statement adds further tests: - a Wightman-side Wick test;
- Maxwell--Fock tests. Consider any isometric embedding of the selected
phase's Hilbert space into a Fock space of transverse photons that
preserves the vacuum and translations and has transverse symmetric
image. No such embedding makes \(\Theta\) local with respect to the Weyl
operators.

\paragraph{Operator product expansion and contact
extensions}\label{operator-product-expansion-and-contact-extensions}

\emph{Unit:
\href{https://github.com/scottnarmstrong/yang-mills-informalization/blob/main/outline/units/VII-4.md}{VII.4
--- Operator product expansion and contact extensions}.}

\textbf{The lattice statement.} The product of insertions in one cell
equals a single insertion of the same total degree, plus a remainder
with small chessboard root. The proof is an induction over the RG levels
between the cluster scale and the observation scale: the part of degree
at most \(D\) is kept, and the rest is paid into the remainder.

\textbf{The continuum coefficients.} They are read off rather than given
by a formula. - A prescription fixes, degree by degree, biorthogonal
systems of fields and detectors. The detectors are finite combinations
of mixed-state correlations of the cluster with fixed, separated
exterior words. - The prescription is averaged over \(O(4)\) and
symmetrized for reality. - The coefficient \(c_{j,\kappa}(a;x,o)\) is a
coordinate of the cluster's new component of degree \(j\) in a real
basis.

Consider an \((r,M,\eta)\)-cluster: points within \(Mr\) of \(o\) and at
mutual distance at least \(\eta r\), carrying fields of total degree
\(w\). Take spectators at distance at least \(c\) from \(o\) and from
each other. Then, for all sufficiently small \(r\),

\[|c_{j,\kappa}(a;x,o)|\le r^{\,j-w-\varepsilon},\qquad
\Big|S^s(a,x;b,y)-\sum_{j\le D}\sum_\kappa c_{j,\kappa}(a;x,o)\,S^s(O_{j,\kappa},o;b,y)\Big|\le r^{\,D+1-w-\varepsilon}.\]

\textbf{Uniformity and limits of the statement.} - The coefficients do
not depend on the spectators \(b\) and \(y\), on the state \(s\), or on
\(D\). - The threshold for \(r\) depends on the fields, on
\(M,\eta,c,\varepsilon,D\), and in the first bound on the channel, but
it is uniform in \(s\), \(x\), \(o\) and \(y\). - Uniformity in the
spectators comes from the train argument of Part V, in which the
exterior word is paid once. - Uniformity in the state comes from
bounding the error in the zero-energy algebra, whose norm dominates
every character and the mixture. - The loss \(\varepsilon\) absorbs
logarithms. There is no logarithmic refinement, so the bound is weaker
than the perturbative OPE. - Uniqueness of the coefficients is not
asserted. Neither convergence as \(D\to\infty\) nor agreement with
perturbation theory is claimed. - Nested versions over finite trees of
clusters, and integrated versions against measures of bounded variation,
also hold.

\textbf{Contact extensions.} These use Taylor subtractions over forests
of collapsing clusters on dyadic shells. The OPE remainder makes the
shell sums converge. For every state and every list of labels this gives
a distribution of finite order on \((\mathbb R^4)^n\) that agrees with
the correlation on \(\mathrm{Conf}_n\). Not established: uniqueness,
reflection positivity on overlapping supports, contact Ward identities,
real-time products, and convergence of the lattice correlations on the
diagonals.

\subsubsection{Part VIII. Assembly}\label{part-viii.-assembly}

\emph{Unit:
\href{https://github.com/scottnarmstrong/yang-mills-informalization/blob/main/outline/units/VIII-1.md}{VIII.1
--- Assembly of the main theorem}.}

The construction run depends only on \(G\) and the penalty \(b\). It is
chosen before any extraction order or observable. For each extraction
order (a spacing sequence and a volume sequence), the assembly composes
further refinements: - one refinement for the limits of the running
coupling; - one that produces the limit table, with the Euclidean
hierarchy, stress tensor, phases and reconstruction of Parts IV--VI.

To these it adds the reconstructions of all phases, the OPE
prescription, the nested and contact results and the Maxwell comparison.
Everything is transported to a common refinement, so all objects refer
to one table.

The public theorem is the case \(b=1\), with the full spacing sequence
and the given volume sequence; its \(q\) is the volume part of the
refinement.

No new estimate enters, but identifying all objects with one table is a
real constraint. Separate existence results for a lattice limit, a
Wightman theory, an OPE and a trace law would not show that they concern
the same correlation functions. The formal construction supplies the
hypotheses under which the limit theory is first proved. The outlines
have not explained every supplier in full; the
\href{https://github.com/scottnarmstrong/yang-mills-informalization/blob/main/outline/units/VIII-1.md\#7-what-remains-unexplained}{assembly
unit} records the limits of its source reading.

\subsubsection{Relation to earlier work, and where the difficulty
lies}\label{relation-to-earlier-work-and-where-the-difficulty-lies}

\textbf{Classical sources.} - Balaban's program is the model for the
ultraviolet analysis: block-spin transformations, gauge fixing around a
background field, fluctuation covariances built from averaging
operators, the small-field/large-field decomposition, and
renormalization of the marginal coupling. - The classical cluster
expansion (Glimm--Jaffe, Brydges, Kotecký--Preiss) organizes the polymer
activities. - Reflection positivity (Osterwalder--Seiler) and chessboard
estimates (Glimm--Jaffe--Spencer, Fröhlich--Israel--Lieb--Simon) connect
lattice estimates with bounds on observables. - Osterwalder--Schrader
and Glimm--Jaffe reconstruction turns Euclidean correlations into
Wightman fields. - The slope is the one-loop coefficient of
Gross--Wilczek and Politzer.

\textbf{Departures from Balaban's scheme.} - The block kernel is an
exactly normalized, gauge-covariant Markov kernel with random tags, not
a linear averaging operator in a fixed background gauge. - Integration
uses a spanning-forest gauge. - The block ratios grow along the levels.
- The threshold carries the factor \(e^{(\log\log\hat\beta)^2}\). - The
construction includes an analytic family in a tuning parameter. The unit
outlines do not establish a comprehensive comparison with Balaban's
papers.

\textbf{Beyond the ultraviolet analysis}, the proof does the following:
- One schedule and one family serve composite fields of every degree.
Each field has a fixed normalization, root bounds uniform in the volume,
and derivatives controlled by rank detection, with no coupling-dependent
loss per site. - Convergence at fixed volume is proved despite tuning
that depends on the cutoff. - Equicontinuity in infinite volume comes
from a rank-free trace inequality, without derivative bounds. -
Commutativity of the zero-energy algebra is proved, not assumed. -
Vacuum uniqueness in every phase follows from a uniform low-energy bound
that replaces a mass gap or a cluster argument. - Lorentz covariance
comes from Ward identities. - Rotation invariance is restored without
decay of correlations, and the group action is obtained by comparison
with free Maxwell theory. - The trace law holds as a limit over all
radii. - The OPE is derived from the RG, with one system of coefficients
for the mixture and every phase.

\textbf{Where the difficulty lies.} 1. \textbf{Gauge degeneracy and
large fields} must be controlled with a bounded cost per site and per
seed. This is done by the gauge count of the determinant budget, the
energy lower bound for seeds, and the exactness of the kernel. 2.
\textbf{Effective-action estimates must become estimates for
observables}: for every composite field, its derivatives and its
fixed-volume limit, inside one family. 3. \textbf{The passage from a
lattice mixture to relativistic phases} has no continuous symmetry and
no clustering to start from. It is handled by Ward identities with free
comparison, the commutative zero-energy algebra, and spectral bounds
based on locality.

The order of choice of constants and the transport of estimates through
the construction account for a substantial part of the formal
development. Its length alone does not measure mathematical novelty.

\subsubsection{Steps whose informal account is
incomplete}\label{steps-whose-informal-account-is-incomplete}

The theorem is proved in Lean, and the conditional statements of the
limit theory are discharged within its proof. The unit outlines,
however, describe the following steps only at the level of statements,
or with readings marked tentative. A written proof will have to supply
them.

\begin{itemize}
\tightlist
\item
  \textbf{One-loop identification (III.7, VII.2).} The slope bound is
  first proved for an auxiliary functional and then transferred to the
  completed Hessian of the actual step. The revised
  \href{https://github.com/scottnarmstrong/yang-mills-informalization/blob/main/outline/units/III-7.md\#7-what-remains-unexplained}{III.7}
  explains the determinant factorization and the transfer of the
  coefficient; several subsidiary real-quotient, complex-strip, tadpole,
  bubble and alias estimates remain inputs at statement level. The trace
  amplitude does not test this value.
\item
  \textbf{Covariance and Hessian (III.4).} The Gaussian covariance is
  the inverse of the gauge-fixed Gram form on the active links. Whether
  it equals the inverse Hessian of the effective action at the
  background is left open.
\item
  \textbf{Nonlinear step.}

  \begin{itemize}
  \tightlist
  \item
    The correspondence between the outputs of the one-step bounds
    (III.10) and the norms restored by the reset (III.9) is not written
    out.
  \item
    Several hypotheses of the step were read only as statements.
  \item
    The cluster expansion is proved for real weights (III.5), while the
    fixed-volume argument uses complex tuning.
  \end{itemize}
\item
  \textbf{Fixed-volume convergence (III.12, IV.1).}
  \href{https://github.com/scottnarmstrong/yang-mills-informalization/blob/main/outline/units/III-12.md\#7-what-remains-unexplained}{III.12}
  now explains fixed-level convergence, the uniform tail bound and the
  Cauchy argument; the underlying trajectory-convergence and multiscale
  estimates remain imported inputs.
  \href{https://github.com/scottnarmstrong/yang-mills-informalization/blob/main/outline/units/IV-1.md}{IV.1}
  locates the upgrade to the full spacing sequence, while leaving the
  upstream construction of its fixed-volume kernels outside that unit's
  review.
\item
  \textbf{Shell identities and the flux at infinity (VI.1, VI.2).} The
  lattice derivation of the shell identities is given only at statement
  level. The exterior-flux estimate was not read. The descent of the
  rotation action to exact intertwiners is marked tentative.
\item
  \textbf{Reconstruction (V.2).} The boundary values, the spectrum
  condition and the locality argument are described from statements;
  their mechanisms were confirmed in the Lean but not written out in the
  unit outlines.
\item
  \textbf{Contact extensions and nested OPE (VII.4).} The shell
  estimates and the nested rates are described from docstrings.
\item
  \textbf{Field relations (II.1).} It is not shown directly that the
  lattice equation relations exhaust the null space. The identification
  \(F_d\cong M_d\) comes from the comparison with free theory.
\end{itemize}

\subsubsection{Reading guide}\label{reading-guide}

The
\href{https://github.com/scottnarmstrong/yang-mills-informalization/blob/main/statement.md}{exact
statement} separates the public theorem from the stronger construction.
The
\href{https://github.com/scottnarmstrong/yang-mills-informalization/blob/main/outline/README.md}{architecture}
describes what each part hands to the next. For the stated inputs and
remaining gaps in this account, follow the relevant unit's ``What it
uses'' and ``What remains unexplained'' sections. The
\href{https://github.com/scottnarmstrong/yang-mills-informalization/blob/main/outline/units/VIII-1.md}{assembly
unit} describes how the construction is put together.

{\def\LTcaptype{none} % do not increment counter
\begin{longtable}[]{@{}
  >{\raggedright\arraybackslash}p{(\linewidth - 2\tabcolsep) * \real{0.5000}}
  >{\raggedright\arraybackslash}p{(\linewidth - 2\tabcolsep) * \real{0.5000}}@{}}
\toprule\noalign{}
\begin{minipage}[b]{\linewidth}\raggedright
Question
\end{minipage} & \begin{minipage}[b]{\linewidth}\raggedright
Units to read
\end{minipage} \\
\midrule\noalign{}
\endhead
\bottomrule\noalign{}
\endlastfoot
What are the fields, and how are their normalizations fixed? &
\href{https://github.com/scottnarmstrong/yang-mills-informalization/blob/main/outline/units/II-1.md}{II.1},
\href{https://github.com/scottnarmstrong/yang-mills-informalization/blob/main/outline/units/III-12.md}{III.12},
\href{https://github.com/scottnarmstrong/yang-mills-informalization/blob/main/outline/units/IV-1.md}{IV.1},
\href{https://github.com/scottnarmstrong/yang-mills-informalization/blob/main/outline/units/VII-1.md}{VII.1} \\
How does the ultraviolet analysis relate to Balaban's? &
\href{https://github.com/scottnarmstrong/yang-mills-informalization/blob/main/outline/units/III-1.md}{III.1},
\href{https://github.com/scottnarmstrong/yang-mills-informalization/blob/main/outline/units/III-3.md}{III.3},
\href{https://github.com/scottnarmstrong/yang-mills-informalization/blob/main/outline/units/III-4.md}{III.4},
\href{https://github.com/scottnarmstrong/yang-mills-informalization/blob/main/outline/units/III-6.md}{III.6},
\href{https://github.com/scottnarmstrong/yang-mills-informalization/blob/main/outline/units/III-5.md}{III.5},
\href{https://github.com/scottnarmstrong/yang-mills-informalization/blob/main/outline/units/III-10.md}{III.10} \\
Where does the one-loop coefficient enter? &
\href{https://github.com/scottnarmstrong/yang-mills-informalization/blob/main/outline/units/III-7.md}{III.7},
\href{https://github.com/scottnarmstrong/yang-mills-informalization/blob/main/outline/units/III-9.md}{III.9},
\href{https://github.com/scottnarmstrong/yang-mills-informalization/blob/main/outline/units/VII-5.md}{VII.5},
\href{https://github.com/scottnarmstrong/yang-mills-informalization/blob/main/outline/units/VII-2.md}{VII.2} \\
Why does one family serve every observable and every limit? &
\href{https://github.com/scottnarmstrong/yang-mills-informalization/blob/main/outline/units/III-2.md}{III.2},
\href{https://github.com/scottnarmstrong/yang-mills-informalization/blob/main/outline/units/II-1.md}{II.1},
\href{https://github.com/scottnarmstrong/yang-mills-informalization/blob/main/outline/units/III-12.md}{III.12},
\href{https://github.com/scottnarmstrong/yang-mills-informalization/blob/main/outline/units/VIII-1.md}{VIII.1};
\href{https://github.com/scottnarmstrong/yang-mills-informalization/blob/main/outline/units/III-13.md}{III.13}
treats an alternative family \\
How are bounds made uniform in the volume? &
\href{https://github.com/scottnarmstrong/yang-mills-informalization/blob/main/outline/units/III-11.md}{III.11},
\href{https://github.com/scottnarmstrong/yang-mills-informalization/blob/main/outline/units/III-12.md}{III.12},
\href{https://github.com/scottnarmstrong/yang-mills-informalization/blob/main/outline/units/III-16.md}{III.16} \\
How is infinite volume reached without an infrared RG? &
\href{https://github.com/scottnarmstrong/yang-mills-informalization/blob/main/outline/units/III-16.md}{III.16},
\href{https://github.com/scottnarmstrong/yang-mills-informalization/blob/main/outline/units/IV-1.md}{IV.1} \\
Why are there phases, and why does every phase reconstruct? &
\href{https://github.com/scottnarmstrong/yang-mills-informalization/blob/main/outline/units/V-1.md}{V.1},
\href{https://github.com/scottnarmstrong/yang-mills-informalization/blob/main/outline/units/V-2.md}{V.2},
\href{https://github.com/scottnarmstrong/yang-mills-informalization/blob/main/outline/units/V-3.md}{V.3};
\href{https://github.com/scottnarmstrong/yang-mills-informalization/blob/main/outline/units/VII-4.md}{VII.4}
for the train bounds \\
How are rotations and the stress tensor recovered? &
\href{https://github.com/scottnarmstrong/yang-mills-informalization/blob/main/outline/units/VI-2.md}{VI.2},
\href{https://github.com/scottnarmstrong/yang-mills-informalization/blob/main/outline/units/VI-1.md}{VI.1},
with
\href{https://github.com/scottnarmstrong/yang-mills-informalization/blob/main/outline/units/VII-1.md}{VII.1}
and
\href{https://github.com/scottnarmstrong/yang-mills-informalization/blob/main/outline/units/VII-6.md}{VII.6}
for the free comparison \\
What makes the theory nonfree? &
\href{https://github.com/scottnarmstrong/yang-mills-informalization/blob/main/outline/units/VII-2.md}{VII.2},
\href{https://github.com/scottnarmstrong/yang-mills-informalization/blob/main/outline/units/VII-3.md}{VII.3},
with
\href{https://github.com/scottnarmstrong/yang-mills-informalization/blob/main/outline/units/VII-6.md}{VII.6}
for arbitrary short-distance scales \\
What does the OPE say, and what is known at coincident points? &
\href{https://github.com/scottnarmstrong/yang-mills-informalization/blob/main/outline/units/VII-4.md}{VII.4},
then
\href{https://github.com/scottnarmstrong/yang-mills-informalization/blob/main/outline/units/VIII-1.md}{VIII.1}
for the assembly \\
\end{longtable}
}

\section{The architecture}\label{the-architecture}

This page describes how OpenAI's Lean proof of the construction of
four-dimensional Yang--Mills theory is organised: its eight parts, what
each part hands to the next, and the units that make up each part
(\href{https://github.com/scottnarmstrong/yang-mills-informalization/blob/main/outline/units.md}{detailed
map}).

It is read from the Lean statements and documentation and from the
dependency graph. The mathematical framing follows the Yang--Mills
group's
\href{https://github.com/yang-mills-group/YM/blob/main/docs/assessment-1.pdf}{assessment}.
Lean names link to the fixed source commit
\href{https://github.com/yang-mills-group/OpenAI-YM-20261028/tree/3456e5e75cf8d4cd558955a5c7638fba2d68f464}{\texttt{3456e5e}}.

\subsubsection{The order of limits}\label{the-order-of-limits}

The construction uses Wilson lattice theories on tori of side \(L_p\),
with lattice spacings \(a_N\) and couplings \(\beta_{N,p}\) satisfying
\(\beta_{N,p}/\log(1/a_N)\to 11C_G/(24\pi^2)\). The public theorem
leaves the positive limiting constant unspecified; see the
\href{https://github.com/scottnarmstrong/yang-mills-informalization/blob/main/statement.md}{exact
statement}.

Correlations of renormalized local observables are first taken to the
limit \textbf{\(N\to\infty\) at fixed volume}, along the whole sequence
of spacings. They are then taken to the limit \textbf{\(p\to\infty\)},
along a subsequence that may be extracted from any prescribed sequence
of volumes. The public form of this statement is
\href{https://github.com/yang-mills-group/OpenAI-YM-20261028/blob/3456e5e75cf8d4cd558955a5c7638fba2d68f464/PublicWilson.lean\#L85-L90}{\texttt{WilsonFamily.Converges}}.

The renormalization group is used only for the first limit. The second
is obtained by compactness.

\subsubsection{The eight parts}\label{the-eight-parts}

{\def\LTcaptype{none} % do not increment counter
\begin{longtable}[]{@{}
  >{\raggedright\arraybackslash}p{(\linewidth - 4\tabcolsep) * \real{0.3333}}
  >{\raggedright\arraybackslash}p{(\linewidth - 4\tabcolsep) * \real{0.3333}}
  >{\raggedright\arraybackslash}p{(\linewidth - 4\tabcolsep) * \real{0.3333}}@{}}
\toprule\noalign{}
\begin{minipage}[b]{\linewidth}\raggedright
part
\end{minipage} & \begin{minipage}[b]{\linewidth}\raggedright
name
\end{minipage} & \begin{minipage}[b]{\linewidth}\raggedright
units
\end{minipage} \\
\midrule\noalign{}
\endhead
\bottomrule\noalign{}
\endlastfoot
I & Setting &
\href{https://github.com/scottnarmstrong/yang-mills-informalization/blob/main/outline/units/I-1.md}{I.1} \\
II & Data of the construction &
\href{https://github.com/scottnarmstrong/yang-mills-informalization/blob/main/outline/units/II-1.md}{II.1},
\href{https://github.com/scottnarmstrong/yang-mills-informalization/blob/main/outline/units/II-2.md}{II.2} \\
III & Ultraviolet construction & III.1--III.16 \\
IV & Infinite volume &
\href{https://github.com/scottnarmstrong/yang-mills-informalization/blob/main/outline/units/IV-1.md}{IV.1} \\
V & Phases and reconstruction &
\href{https://github.com/scottnarmstrong/yang-mills-informalization/blob/main/outline/units/V-1.md}{V.1},
\href{https://github.com/scottnarmstrong/yang-mills-informalization/blob/main/outline/units/V-2.md}{V.2},
\href{https://github.com/scottnarmstrong/yang-mills-informalization/blob/main/outline/units/V-3.md}{V.3} \\
VI & Symmetry &
\href{https://github.com/scottnarmstrong/yang-mills-informalization/blob/main/outline/units/VI-1.md}{VI.1},
\href{https://github.com/scottnarmstrong/yang-mills-informalization/blob/main/outline/units/VI-2.md}{VI.2} \\
VII & Identification as Yang--Mills &
\href{https://github.com/scottnarmstrong/yang-mills-informalization/blob/main/outline/units/VII-1.md}{VII.1},
\href{https://github.com/scottnarmstrong/yang-mills-informalization/blob/main/outline/units/VII-2.md}{VII.2},
\href{https://github.com/scottnarmstrong/yang-mills-informalization/blob/main/outline/units/VII-3.md}{VII.3},
\href{https://github.com/scottnarmstrong/yang-mills-informalization/blob/main/outline/units/VII-4.md}{VII.4},
\href{https://github.com/scottnarmstrong/yang-mills-informalization/blob/main/outline/units/VII-5.md}{VII.5},
\href{https://github.com/scottnarmstrong/yang-mills-informalization/blob/main/outline/units/VII-6.md}{VII.6} \\
VIII & Assembly &
\href{https://github.com/scottnarmstrong/yang-mills-informalization/blob/main/outline/units/VIII-1.md}{VIII.1} \\
\end{longtable}
}

\paragraph{Ultraviolet units}\label{ultraviolet-units}

\begin{itemize}
\tightlist
\item
  Blocking, constants and Gaussian estimates:
  \href{https://github.com/scottnarmstrong/yang-mills-informalization/blob/main/outline/units/III-1.md}{III.1},
  \href{https://github.com/scottnarmstrong/yang-mills-informalization/blob/main/outline/units/III-2.md}{III.2},
  \href{https://github.com/scottnarmstrong/yang-mills-informalization/blob/main/outline/units/III-3.md}{III.3},
  \href{https://github.com/scottnarmstrong/yang-mills-informalization/blob/main/outline/units/III-4.md}{III.4}.
\item
  Cluster expansion, stationary field, matching and dependence:
  \href{https://github.com/scottnarmstrong/yang-mills-informalization/blob/main/outline/units/III-5.md}{III.5},
  \href{https://github.com/scottnarmstrong/yang-mills-informalization/blob/main/outline/units/III-6.md}{III.6},
  \href{https://github.com/scottnarmstrong/yang-mills-informalization/blob/main/outline/units/III-7.md}{III.7},
  \href{https://github.com/scottnarmstrong/yang-mills-informalization/blob/main/outline/units/III-8.md}{III.8}.
\item
  Coupling reset, nonlinear step and observable limits:
  \href{https://github.com/scottnarmstrong/yang-mills-informalization/blob/main/outline/units/III-9.md}{III.9},
  \href{https://github.com/scottnarmstrong/yang-mills-informalization/blob/main/outline/units/III-10.md}{III.10},
  \href{https://github.com/scottnarmstrong/yang-mills-informalization/blob/main/outline/units/III-11.md}{III.11},
  \href{https://github.com/scottnarmstrong/yang-mills-informalization/blob/main/outline/units/III-12.md}{III.12}.
\item
  Alternative coordinated family, construction data and analytic tools:
  \href{https://github.com/scottnarmstrong/yang-mills-informalization/blob/main/outline/units/III-13.md}{III.13},
  \href{https://github.com/scottnarmstrong/yang-mills-informalization/blob/main/outline/units/III-14.md}{III.14},
  \href{https://github.com/scottnarmstrong/yang-mills-informalization/blob/main/outline/units/III-15.md}{III.15},
  \href{https://github.com/scottnarmstrong/yang-mills-informalization/blob/main/outline/units/III-16.md}{III.16}.
\end{itemize}

The limit theory (Parts IV--VI) is written in the code around numbered
\href{https://github.com/scottnarmstrong/yang-mills-informalization/blob/main/outline/work-items.md}{work
items}; the units of those parts follow them.

\subsubsection{What each part hands to the
next}\label{what-each-part-hands-to-the-next}

\paragraph{Ultraviolet → infinite
volume}\label{ultraviolet-infinite-volume}

Let \(G\) be a compact matrix group with simple nonabelian Lie algebra,
and \(b>0\) a parameter of the block-spin kernel. The ultraviolet
construction provides a Wilson family with the following properties.

\begin{enumerate}
\def\labelenumi{\arabic{enumi}.}
\tightlist
\item
  \textbf{Specification.} The family satisfies its normalization,
  covariance and budget conditions, and its couplings obey the
  asymptotic-freedom law above.
\item
  \textbf{Chessboard bound, uniform in the volume.} For each degree
  \(d\) and \(\varepsilon>0\), every renormalized local observable at
  block scale \(m\) satisfies a chessboard bound of order
  \(a_m^{-d-\varepsilon}\), uniformly in the volume, once the lattice is
  fine enough. A chessboard bound is a bound on the root of a product of
  reflected copies.
\item
  \textbf{Differences.} The same holds for finite-difference quotients,
  so correlation kernels at \(r\)-separated points are bounded uniformly
  in the volume.
\item
  \textbf{Fixed volume.} At each fixed volume the correlation kernels
  converge along the whole sequence of spacings, to smooth limits.
\item
  \textbf{Corrected fields.} The same estimates hold for the corrected
  fields.
\end{enumerate}

For the subsequential extraction, the infinite-volume part uses 1--3 and
5, together with the reflection-positivity inequality on cells, which
holds for every \(\beta\ge0\). Its conclusion: for every increasing
sequence of volumes there is a subsequence along which all smeared
correlations converge to kernels that are smooth away from coincident
points. The fixed-volume convergence in item 4 then removes the spacing
refinement; see
\href{https://github.com/scottnarmstrong/yang-mills-informalization/blob/main/outline/units/IV-1.md}{IV.1}.
These two limits correspond to Theorems 4.1 and 5.1 of the group's
assessment.

A caution about the code. About three quarters of the code filed under
``full volume'' (folder
\href{https://github.com/yang-mills-group/OpenAI-YM-20261028/tree/3456e5e75cf8d4cd558955a5c7638fba2d68f464/YM/FullVolume}{\texttt{YM/FullVolume}})
belongs to the ultraviolet construction: it proves the volume-uniform
bounds. The genuinely infinite-volume argument is short, about 7,000
lines.

\paragraph{Infinite volume → phases, reconstruction and
symmetry}\label{infinite-volume-phases-reconstruction-and-symmetry}

The limit satisfies a Euclidean \emph{hierarchy}:

\begin{itemize}
\tightlist
\item
  translation homogeneity;
\item
  reflection positivity, including for smeared joint test functions;
\item
  smoothness on separated configurations, with controlled derivatives;
\item
  a compatible calculus of derivatives.
\end{itemize}

Reconstruction also needs a uniform bound on products of fields in thin
time slabs.

\paragraph{Reconstruction →
identification}\label{reconstruction-identification}

For every phase, a completed Wightman theory: unique vacuum, a
light-cone bound, heat-kernel decay, and a trivial weak commutant. The
heat bound is polynomial; no mass gap follows.

\paragraph{Symmetry → identification}\label{symmetry-identification}

A conserved stress tensor and its Ward identities. Both are compared,
along sequences of scales tending to zero, with the corresponding
objects for free Maxwell fields.

\paragraph{Into the assembly}\label{into-the-assembly}

A single Wilson family serves every part of the public theorem, from
extraction to the trace test;
\href{https://github.com/scottnarmstrong/yang-mills-informalization/blob/main/outline/units/VIII-1.md}{VIII.1}
describes the assembly. The construction run exists for every allowed
\(G\) and every \(b>0\).
\href{https://github.com/scottnarmstrong/yang-mills-informalization/blob/main/outline/units/III-13.md}{III.13}
describes an alternative coordinated family; its properties require a
separate identification before they can be used for this family.

\subsubsection{The work items of the limit
theory}\label{the-work-items-of-the-limit-theory}

Parts IV--VI are written in the code as about seventy numbered work
items. Each states one obligation, a ``proof request'' in the code's
vocabulary, and proves it in its own file.

{\def\LTcaptype{none} % do not increment counter
\begin{longtable}[]{@{}
  >{\raggedright\arraybackslash}p{(\linewidth - 6\tabcolsep) * \real{0.2500}}
  >{\raggedright\arraybackslash}p{(\linewidth - 6\tabcolsep) * \real{0.2500}}
  >{\raggedright\arraybackslash}p{(\linewidth - 6\tabcolsep) * \real{0.2500}}
  >{\raggedright\arraybackslash}p{(\linewidth - 6\tabcolsep) * \real{0.2500}}@{}}
\toprule\noalign{}
\begin{minipage}[b]{\linewidth}\raggedright
family
\end{minipage} & \begin{minipage}[b]{\linewidth}\raggedright
items
\end{minipage} & \begin{minipage}[b]{\linewidth}\raggedright
part
\end{minipage} & \begin{minipage}[b]{\linewidth}\raggedright
content
\end{minipage} \\
\midrule\noalign{}
\endhead
\bottomrule\noalign{}
\endlastfoot
D & 6 & IV & the limit passes to derivatives of the correlation kernels;
the Euclidean hierarchy \\
B & 18 & V & time-slab operators, decomposition into phases, reflection
positivity and smoothness for each phase \\
C & 29 & V & Osterwalder--Schrader → Wightman: continuation in time,
boost covariance, gluing of tubes, locality \\
A & 17 & VI & so(4) brackets, integration to the rotation action, the
averaged conserved stress tensor, boosts, comparison with free Maxwell
fields \\
\end{longtable}
}

These families interleave. D supplies the limit hierarchy; B begins with
its OS spaces and slab operators. A uses those spaces and the free
Maxwell comparison to obtain rotations and Ward identities. Those
symmetry results enter both the phase decomposition in B and the boost
and covariance arguments in C. The part numbers organize the subjects;
they do not give a linear dependency order. See
\href{https://github.com/scottnarmstrong/yang-mills-informalization/blob/main/summary.md\#the-strategy}{the
strategy} and the
\href{https://github.com/scottnarmstrong/yang-mills-informalization/blob/main/outline/work-items.md}{work-item
introductions}.

The same mechanism appears in the ultraviolet construction: families
E--I cover the energy chain of one renormalization step, its stationary
point and Hessian, label polynomials, and support bookkeeping.

See
\href{https://github.com/scottnarmstrong/yang-mills-informalization/blob/main/outline/work-items.md}{work-items.md}
for the complete list.

\subsubsection{What does not follow the order of the
parts}\label{what-does-not-follow-the-order-of-the-parts}

The Lean is not written in this order everywhere. Some definitions are
needed before the part that explains their properties:

\begin{itemize}
\tightlist
\item
  the reflection-positive cells and chessboard roots;
\item
  the record of the whole construction run;
\item
  the basic types of the renormalization group and of the labels.
\end{itemize}

Some sections are filed far from their mathematical home. Two examples:

\begin{itemize}
\tightlist
\item
  the rotation-invariance estimate, in folder \texttt{Next/ROTATIONS};
\item
  some renormalization-group pieces filed with the trace-anomaly code.
\end{itemize}

None of this involves new mathematics. A written account should simply
state these definitions early.

\subsubsection{How reliable is this}\label{how-reliable-is-this}

\begin{itemize}
\tightlist
\item
  Reachable declaration sets for seventeen named results were compared
  with OpenAI's compiled record. This comparison does not verify every
  edge or the mathematical interpretation of the interfaces
  (\href{https://github.com/scottnarmstrong/yang-mills-informalization/blob/main/method.md\#1-the-dependency-graph}{method}).
\item
  The plain-language interfaces summarize statements and documentation.
  The
  \href{https://github.com/scottnarmstrong/yang-mills-informalization/blob/main/outline/units.md}{unit
  index} gives each outline's review status; its local limitations
  describe which arguments remain at statement level. The
  \href{https://github.com/scottnarmstrong/yang-mills-informalization/blob/main/statement.md}{exact
  statement} distinguishes the public theorem from additional
  construction conclusions.
\item
  \href{https://github.com/scottnarmstrong/yang-mills-informalization/issues}{Report
  a correction} with the page, heading and source reference.
\end{itemize}

\section{The units}\label{the-units}

\subsection{I.1. Lattice, symmetries, gauge groups and the Wilson
measure}\label{i.1.-lattice-symmetries-gauge-groups-and-the-wilson-measure}

\textbf{Part:} I, Setting. \textbf{Lean:}
\href{https://github.com/yang-mills-group/OpenAI-YM-20261028/tree/3456e5e75cf8d4cd558955a5c7638fba2d68f464/YM/Lattice}{\texttt{YM/Lattice}}
(1,215 lines),
\href{https://github.com/yang-mills-group/OpenAI-YM-20261028/tree/3456e5e75cf8d4cd558955a5c7638fba2d68f464/YM/Gauge}{\texttt{YM/Gauge}}
(1,131),
\href{https://github.com/yang-mills-group/OpenAI-YM-20261028/tree/3456e5e75cf8d4cd558955a5c7638fba2d68f464/YM/Space}{\texttt{YM/Space}}
(430),
\href{https://github.com/yang-mills-group/OpenAI-YM-20261028/tree/3456e5e75cf8d4cd558955a5c7638fba2d68f464/YM/Corona}{\texttt{YM/Corona}}
(405),
\href{https://github.com/yang-mills-group/OpenAI-YM-20261028/tree/3456e5e75cf8d4cd558955a5c7638fba2d68f464/YM/Mink}{\texttt{YM/Mink}}
(363),
\href{https://github.com/yang-mills-group/OpenAI-YM-20261028/tree/3456e5e75cf8d4cd558955a5c7638fba2d68f464/YM/Wilson}{\texttt{YM/Wilson}}
(157). \textbf{Size:} 3,701 lines in 595 declarations used by the main
theorem.

\subsubsection{1. What this unit proves}\label{what-this-unit-proves}

The unit fixes the setting of the construction:

\begin{itemize}
\tightlist
\item
  the periodic lattices and their symmetry group;
\item
  parallel transport and gauge transformations;
\item
  the subdivision of a torus into blocks and the schedule of lattice
  spacings;
\item
  a combinatorial measure of the size of a lattice set;
\item
  the gauge group in its defining representation, with its Lie algebra,
  Casimir constant, local logarithm and nearest-point retraction;
\item
  the Wilson measure and its invariance;
\item
  the continuum symmetry groups of \(\mathbb R^4\).
\end{itemize}

The results require no input theorem from another unit. Simplicity of
\(\mathfrak g\) is assumed only in Proposition I.1.6(d)--(f); the other
hypotheses are stated locally.

\emph{Conventions.} \(G\) is a closed subgroup of \(U(N)\), with
\(N\ge1\) fixed. For an integer \(\ell\ge1\), write
\(\Lambda_\ell=(\mathbb Z/\ell\mathbb Z)^4\). Axes \(0,1,2\) are spatial
and axis \(3\) is time; \(e_i\) are the unit vectors. A link \((x,i)\)
joins \(x\) to \(x+e_i\). A configuration is a map
\(U:\Lambda_\ell\times\{0,1,2,3\}\to G\); the same transport and gauge
definitions apply on the cover \(\mathbb Z^4\). The norm \(\|\cdot\|_F\)
is the Frobenius norm on \(M_N(\mathbb C)\); norms on real coordinate
spaces are Euclidean unless marked otherwise.

\textbf{Proposition I.1.1 (torus geometry).} For \(y\in\mathbb Z^4\) let
\([y]_\ell\in\Lambda_\ell\) be its reduction mod \(\ell\). For
\(x\in\Lambda_\ell\) let \(r_\ell(x)\in\mathbb Z^4\) be the
representative with coordinates in \((-\ell/2,\ell/2]\) (the centred
lift; the tie \(x_i=\ell/2\) lifts to \(+\ell/2\)). Then:

\begin{itemize}
\tightlist
\item
  \([r_\ell(x)]_\ell=x\), and \(r_\ell(x)\) is the unique \(y\) with
  \([y]_\ell=x\) and \(-\ell<2y_i\le\ell\) for all \(i\);
\item
  reduction is injective on this centred box, and
  \(|\Lambda_\ell|=\ell^4\).
\end{itemize}

Put \(|x|_\ell=\max_i|r_\ell(x)_i|\). The separation \(|x-y|_\ell\) is a
metric on the torus:

\begin{itemize}
\tightlist
\item
  it is symmetric and is zero exactly when \(x=y\);
\item
  it satisfies the triangle inequality;
\item
  it is invariant under translations and under the hyperoctahedral group
  \(B_4\) of Proposition I.1.2;
\item
  reduction does not increase distances:
  \(|[y]_\ell|_\ell\le\max_i|y_i|\).
\end{itemize}

\emph{Status:} theorem. \emph{Lean:}
\href{https://github.com/yang-mills-group/OpenAI-YM-20261028/blob/3456e5e75cf8d4cd558955a5c7638fba2d68f464/YM/Lattice/Site.lean\#L19}{\texttt{Lattice.Site}},
\href{https://github.com/yang-mills-group/OpenAI-YM-20261028/blob/3456e5e75cf8d4cd558955a5c7638fba2d68f464/YM/Lattice/Site.lean\#L35-L38}{\texttt{wrap}},
\href{https://github.com/yang-mills-group/OpenAI-YM-20261028/blob/3456e5e75cf8d4cd558955a5c7638fba2d68f464/YM/Lattice/Site.lean\#L51}{\texttt{lift}},
\href{https://github.com/yang-mills-group/OpenAI-YM-20261028/blob/3456e5e75cf8d4cd558955a5c7638fba2d68f464/YM/Lattice/Site.lean\#L61-L68}{\texttt{lift\_eq\_iff}},
\href{https://github.com/yang-mills-group/OpenAI-YM-20261028/blob/3456e5e75cf8d4cd558955a5c7638fba2d68f464/YM/Lattice/Site.lean\#L70-L75}{\texttt{wrap\_injOn}},
\href{https://github.com/yang-mills-group/OpenAI-YM-20261028/blob/3456e5e75cf8d4cd558955a5c7638fba2d68f464/YM/Lattice/Site.lean\#L32-L33}{\texttt{site\_card}},
\href{https://github.com/yang-mills-group/OpenAI-YM-20261028/blob/3456e5e75cf8d4cd558955a5c7638fba2d68f464/YM/Lattice/Site.lean\#L81}{\texttt{separation}},
\href{https://github.com/yang-mills-group/OpenAI-YM-20261028/blob/3456e5e75cf8d4cd558955a5c7638fba2d68f464/YM/Lattice/Site.lean\#L93-L94}{\texttt{separation\_comm}},
\href{https://github.com/yang-mills-group/OpenAI-YM-20261028/blob/3456e5e75cf8d4cd558955a5c7638fba2d68f464/YM/Lattice/Site.lean\#L103-L105}{\texttt{separation\_triangle}},
\href{https://github.com/yang-mills-group/OpenAI-YM-20261028/blob/3456e5e75cf8d4cd558955a5c7638fba2d68f464/YM/Lattice/Site.lean\#L117-L118}{\texttt{separation\_translate}},
\href{https://github.com/yang-mills-group/OpenAI-YM-20261028/blob/3456e5e75cf8d4cd558955a5c7638fba2d68f464/YM/Lattice/Site.lean\#L119-L121}{\texttt{separation\_signed}},
\href{https://github.com/yang-mills-group/OpenAI-YM-20261028/blob/3456e5e75cf8d4cd558955a5c7638fba2d68f464/YM/Corona/Fill.lean\#L32-L39}{\texttt{Corona.lap}}.

\textbf{Proposition I.1.2 (lattice symmetries, transport, gauge
transformations).} Let \(B_4\) be the group of signed permutations of
the four axes; it has order \(2^4\cdot4!=384\). An element
\(\rho=(\pi,\epsilon)\), with \(\epsilon_i\in\{\pm1\}\), acts on
\(\mathbb Z^4\) and on \(\Lambda_\ell\) by
\((\rho x)_{\pi(i)}=\epsilon_ix_i\). On either lattice \(\Lambda\), an
affine map is \(A=(\rho,b)\), \(Ax=\rho x+b\), with \(b\in\Lambda\).

A walk \(w\) is a finite word of steps \(\pm e_i\). The transport
\(U(x;w)\in G\) is the ordered product along \(w\) from \(x\): a forward
step along the link \((y,i)\) contributes \(U(y,i)\), and a backward
step from \(y\) contributes \(U(y-e_i,i)^{-1}\).

The pullback \(A^*U\) is
\[(A^*U)(x,i)=\begin{cases}U(Ax,\pi(i)),&\epsilon_i=+1,\\ U(Ax-e_{\pi(i)},\pi(i))^{-1},&\epsilon_i=-1,\end{cases}\]
that is, the transport of \(U\) along the image of the link \((x,i)\).
Then:

\begin{itemize}
\item
  \begin{enumerate}
  \def\labelenumi{(\alph{enumi})}
  \tightlist
  \item
    \(1^*=\mathrm{id}\) and \(B^*A^*=(A\circ B)^*\),
    i.e.~\((\rho',b')^*(\rho,b)^*=(\rho\rho',\rho b'+b)^*\). For
    \(A=(1,b)\) the pullback is the translation
    \((\tau_bU)(x,i)=U(x+b,i)\).
  \end{enumerate}
\item
  (b) \((A^*U)(x;w)=U(Ax;\rho w)\), where \(\rho w\) is the image word.
\item
  (c) For \(g:\Lambda\to G\) put \(U^g(x,i)=g(x)U(x,i)g(x+e_i)^{-1}\).
  Then \(U^g(x;w)=g(x)\,U(x;w)\,g(x_{\rm end})^{-1}\); for a closed walk
  this is conjugation by \(g(x)\).
\item
  (d) (locality) If \(U\) and \(V\) agree on every link whose two
  endpoints are visited by \(w\) from \(x\), then \(U(x;w)=V(x;w)\).
\item
  (e) (plaquettes) A plaquette is \(p=(x;i<j)\) with holonomy
  \(U_p=U(x;e_i,e_j,-e_i,-e_j)\). Each \(A\) induces an injective map
  \$A\_\sharp \$ of plaquettes, bijective on \(\Lambda_\ell\). It sends
  the axis pair to \(\{\pi(i),\pi(j)\}\) in increasing order and the
  base point to \(Ax-\sum_{k\in\{i,j\}:\,\epsilon_k=-1}e_{\pi(k)}\). For
  every set \(M\) and every \(\phi:G\to M\) satisfying
  \(\phi(u^{-1})=\phi(u)\) and \(\phi(vuv^{-1})=\phi(u)\) for all
  \(u,v\in G\), one has \(\phi((A^*U)_p)=\phi(U_{A_\sharp p})\).
\item
  (f) (periodic extension) For \(b\in\Lambda_\ell\) let
  \((E_bU)(y,i)=U([y]_\ell+b,i)\) on \(\mathbb Z^4\). Then:

  \begin{itemize}
  \tightlist
  \item
    \((E_bU)(y;w)=U([y]_\ell+b;w)\);
  \item
    \(E_b((\rho,a)^*U)=(\rho,0)^*(E_{\rho b+a}U)\) for
    \(a\in\Lambda_\ell\);
  \item
    \(E_b(U^g)=(E_bU)^{g([\cdot]_\ell+b)}\);
  \item
    if \(F\) depends only on links inside \(X\subset\mathbb Z^4\), then
    \(F\circ E_b\) depends only on links inside \([X]_\ell+b\).
  \end{itemize}
\end{itemize}

A function \(F\) of configurations is \emph{gauge invariant} if
\(F(U^g)=F(U)\) for all \(g\). It is \emph{local on \(X\)} if
\(F(U)=F(V)\) whenever \(U,V\) agree on all links with both endpoints in
\(X\). These are the notions of (Reg3)--(Reg4) in the main statement.
\emph{Status:} theorem. \emph{Lean:}
\href{https://github.com/yang-mills-group/OpenAI-YM-20261028/blob/3456e5e75cf8d4cd558955a5c7638fba2d68f464/YM/Space/Index.lean\#L23-L26}{\texttt{Signed}}
with
\href{https://github.com/yang-mills-group/OpenAI-YM-20261028/blob/3456e5e75cf8d4cd558955a5c7638fba2d68f464/YM/Space/Index.lean\#L47-L51}{its
group law},
\href{https://github.com/yang-mills-group/OpenAI-YM-20261028/blob/3456e5e75cf8d4cd558955a5c7638fba2d68f464/YM/Space/Index.lean\#L62-L63}{\texttt{Signed.act}},
\href{https://github.com/yang-mills-group/OpenAI-YM-20261028/blob/3456e5e75cf8d4cd558955a5c7638fba2d68f464/YM/Lattice/Path.lean\#L105-L107}{\texttt{travel}},
\href{https://github.com/yang-mills-group/OpenAI-YM-20261028/blob/3456e5e75cf8d4cd558955a5c7638fba2d68f464/YM/Lattice/Path.lean\#L147-L148}{\texttt{pull}},
\href{https://github.com/yang-mills-group/OpenAI-YM-20261028/blob/3456e5e75cf8d4cd558955a5c7638fba2d68f464/YM/Lattice/Path.lean\#L173-L179}{\texttt{pull\_pull}},
\href{https://github.com/yang-mills-group/OpenAI-YM-20261028/blob/3456e5e75cf8d4cd558955a5c7638fba2d68f464/YM/Lattice/Path.lean\#L184-L186}{\texttt{pull\_one}},
\href{https://github.com/yang-mills-group/OpenAI-YM-20261028/blob/3456e5e75cf8d4cd558955a5c7638fba2d68f464/YM/Lattice/Path.lean\#L160-L167}{\texttt{travel\_pull}},
\href{https://github.com/yang-mills-group/OpenAI-YM-20261028/blob/3456e5e75cf8d4cd558955a5c7638fba2d68f464/YM/Lattice/Path.lean\#L132-L133}{\texttt{gauge}},
\href{https://github.com/yang-mills-group/OpenAI-YM-20261028/blob/3456e5e75cf8d4cd558955a5c7638fba2d68f464/YM/Lattice/Path.lean\#L141-L143}{\texttt{travel\_gauge}},
\href{https://github.com/yang-mills-group/OpenAI-YM-20261028/blob/3456e5e75cf8d4cd558955a5c7638fba2d68f464/YM/Lattice/Path.lean\#L135}{\texttt{GaugeInvariant}},
\href{https://github.com/yang-mills-group/OpenAI-YM-20261028/blob/3456e5e75cf8d4cd558955a5c7638fba2d68f464/YM/Lattice/Path.lean\#L192-L193}{\texttt{Agrees}},
\href{https://github.com/yang-mills-group/OpenAI-YM-20261028/blob/3456e5e75cf8d4cd558955a5c7638fba2d68f464/YM/Lattice/Path.lean\#L196-L197}{\texttt{LocalOn}},
\href{https://github.com/yang-mills-group/OpenAI-YM-20261028/blob/3456e5e75cf8d4cd558955a5c7638fba2d68f464/YM/Lattice/Path.lean\#L200-L220}{\texttt{travel\_local}},
\href{https://github.com/yang-mills-group/OpenAI-YM-20261028/blob/3456e5e75cf8d4cd558955a5c7638fba2d68f464/YM/Lattice/Faces.lean\#L115-L116}{\texttt{Face.send}},
\href{https://github.com/yang-mills-group/OpenAI-YM-20261028/blob/3456e5e75cf8d4cd558955a5c7638fba2d68f464/YM/Lattice/Faces.lean\#L119-L126}{\texttt{Face.injective\_send}},
\href{https://github.com/yang-mills-group/OpenAI-YM-20261028/blob/3456e5e75cf8d4cd558955a5c7638fba2d68f464/YM/Lattice/Faces.lean\#L128-L143}{\texttt{Face.send\_price}},
\href{https://github.com/yang-mills-group/OpenAI-YM-20261028/blob/3456e5e75cf8d4cd558955a5c7638fba2d68f464/YM/Lattice/Path.lean\#L224-L225}{\texttt{expandAt}},
\href{https://github.com/yang-mills-group/OpenAI-YM-20261028/blob/3456e5e75cf8d4cd558955a5c7638fba2d68f464/YM/Lattice/Path.lean\#L227-L235}{\texttt{expandAt\_pull}},
\href{https://github.com/yang-mills-group/OpenAI-YM-20261028/blob/3456e5e75cf8d4cd558955a5c7638fba2d68f464/YM/Lattice/Path.lean\#L240-L247}{\texttt{travel\_at}},
\href{https://github.com/yang-mills-group/OpenAI-YM-20261028/blob/3456e5e75cf8d4cd558955a5c7638fba2d68f464/YM/Lattice/Path.lean\#L249-L251}{\texttt{expandAt\_gauge}},
\href{https://github.com/yang-mills-group/OpenAI-YM-20261028/blob/3456e5e75cf8d4cd558955a5c7638fba2d68f464/YM/Lattice/Path.lean\#L253-L261}{\texttt{LocalOn.expandAt}}.

\textbf{Proposition I.1.3 (schedule of spacings; subdivision and
blocks).}

\begin{enumerate}
\def\labelenumi{(\alph{enumi})}
\tightlist
\item
  \emph{Schedule.} A schedule consists of odd integers \(L_m\ge3\) and a
  strictly increasing sequence of positive integers \(P_p\) with
  \(4\mid P_p\), indexed by \(m,p\in\mathbb N\). Put
  \[D_0=1,\quad D_{m+1}=D_mL_m,\quad a_m=1/D_m,\quad \ell(m,p)=P_pD_m,\quad s_m=\log(1/a_m).\]
  Then:
\end{enumerate}

\begin{itemize}
\tightlist
\item
  \(D_m\) is odd and strictly increasing;
\item
  \(a_0=1\), \(a_{m+1}=a_m/L_m\), and \(a_m\downarrow0\) strictly;
\item
  \(s_0=0\), \(s_{m+1}=s_m+\log L_m\), and \(s_m\to\infty\);
\item
  \(4\mid\ell(m,p)\) and \(a_m\ell(m,p)=P_p\);
\item
  the placement \(x\mapsto a_m\,r_{\ell(m,p)}(x)\) maps
  \(\Lambda_{\ell(m,p)}\) into \((-P_p/2,P_p/2]^4\).
\end{itemize}

For \(m\le m'\), \(\Lambda_{\ell(m,p)}\to\Lambda_{\ell(m',p)}\) is a
subdivision of odd ratio \(L_m\cdots L_{m'-1}\) that preserves
placements.

\begin{enumerate}
\def\labelenumi{(\alph{enumi})}
\setcounter{enumi}{1}
\item
  \emph{Mesh inequality.} If \(\log L_m/s_m\to0\) as \(m\to\infty\) (the
  schedule is \emph{gentle}), then there is \(c>0\), depending on the
  schedule, with \(a_{m+1}\ge c\,a_m^2\) for all \(m\in\mathbb N\). The
  ratio defining gentleness is considered for \(m\ge1\), when \(s_m>0\).
\item
  \emph{Subdivision.} Let \(L\ge1\) be an integer and \(\ell'=\ell L\).
  Put \(\iota(x)=[L\,r_\ell(x)]_{\ell'}\). Then:
\end{enumerate}

\begin{itemize}
\tightlist
\item
  \(r_{\ell'}(\iota x)=L\,r_\ell(x)\);
\item
  \(\iota:\Lambda_\ell\to\Lambda_{\ell'}\) is an injective group
  homomorphism commuting with \(B_4\).
\end{itemize}

The residue map \(\mathrm{res}:\Lambda_{\ell'}\to\Lambda_L\) (reduction
mod \(L\)) is additive with \(\mathrm{res}\circ\iota=0\). The map
\((x,y)\mapsto\iota(x)+[r_L(y)]_{\ell'}\) is a bijection
\(\Lambda_\ell\times\Lambda_L\to\Lambda_{\ell'}\), with inverse
\(z\mapsto(\mathrm{par}\,z,\mathrm{res}\,z)\); moreover
\(\mathrm{par}\circ\iota=\mathrm{id}\).

\begin{enumerate}
\def\labelenumi{(\alph{enumi})}
\setcounter{enumi}{3}
\tightlist
\item
  \emph{Odd ratio.} If \(L\) is odd, then:
\end{enumerate}

\begin{itemize}
\tightlist
\item
  \(r_L\) commutes with \(B_4\) and
  \(\mathrm{par}(\rho z)=\rho\,\mathrm{par}(z)\);
\item
  the block \(\mathrm{par}^{-1}(x)\) is
  \(\iota(x)+[\{-\tfrac{L-1}2,\dots,\tfrac{L-1}2\}^4]_{\ell'}\);
\item
  \(\mathrm{par}(z+e_i)=\mathrm{par}(z)+e_i\) if
  \(r_L(\mathrm{res}\,z)_i=\tfrac{L-1}2\), and
  \(\mathrm{par}(z+e_i)=\mathrm{par}(z)\) otherwise.
\end{itemize}

\begin{enumerate}
\def\labelenumi{(\alph{enumi})}
\setcounter{enumi}{4}
\tightlist
\item
  \emph{Metric comparison.} For every integer \(L\ge1\), all
  \(x,x'\in\Lambda_\ell\) and all \(z,z'\in\Lambda_{\ell'}\):
\end{enumerate}

\begin{itemize}
\tightlist
\item
  \(|z-\iota\,\mathrm{par}\,z|_{\ell'}\le L/2\);
\item
  \(L\,|x-x'|_\ell\le|\iota x-\iota x'|_{\ell'}\);
\item
  \(L\,|\mathrm{par}\,z-\mathrm{par}\,z'|_\ell\le|z-z'|_{\ell'}+L\)
  (foreshortening).
\end{itemize}

\emph{Status:} (a) the schedule is a definition (data with conditions);
the schedule of the construction is chosen in III.2. The consequences
are theorems. \emph{Lean:}
\href{https://github.com/yang-mills-group/OpenAI-YM-20261028/blob/3456e5e75cf8d4cd558955a5c7638fba2d68f464/YM/Lattice/Scale.lean\#L28-L34}{\texttt{Schedule}},
\href{https://github.com/yang-mills-group/OpenAI-YM-20261028/blob/3456e5e75cf8d4cd558955a5c7638fba2d68f464/YM/Lattice/Scale.lean\#L39-L41}{\texttt{denominator}},
\href{https://github.com/yang-mills-group/OpenAI-YM-20261028/blob/3456e5e75cf8d4cd558955a5c7638fba2d68f464/YM/Lattice/Scale.lean\#L67}{\texttt{side}},
\href{https://github.com/yang-mills-group/OpenAI-YM-20261028/blob/3456e5e75cf8d4cd558955a5c7638fba2d68f464/YM/Lattice/Scale.lean\#L76}{\texttt{side\_four}},
\href{https://github.com/yang-mills-group/OpenAI-YM-20261028/blob/3456e5e75cf8d4cd558955a5c7638fba2d68f464/YM/Lattice/Scale.lean\#L99}{\texttt{a}},
\href{https://github.com/yang-mills-group/OpenAI-YM-20261028/blob/3456e5e75cf8d4cd558955a5c7638fba2d68f464/YM/Lattice/Scale.lean\#L111-L113}{\texttt{strictAnti}},
\href{https://github.com/yang-mills-group/OpenAI-YM-20261028/blob/3456e5e75cf8d4cd558955a5c7638fba2d68f464/YM/Lattice/Scale.lean\#L108-L109}{\texttt{shrinks}},
\href{https://github.com/yang-mills-group/OpenAI-YM-20261028/blob/3456e5e75cf8d4cd558955a5c7638fba2d68f464/YM/Lattice/Scale.lean\#L115}{\texttt{scale}},
\href{https://github.com/yang-mills-group/OpenAI-YM-20261028/blob/3456e5e75cf8d4cd558955a5c7638fba2d68f464/YM/Lattice/Scale.lean\#L117-L120}{\texttt{divergent}},
\href{https://github.com/yang-mills-group/OpenAI-YM-20261028/blob/3456e5e75cf8d4cd558955a5c7638fba2d68f464/YM/Lattice/Scale.lean\#L169-L171}{\texttt{span\_side}},
\href{https://github.com/yang-mills-group/OpenAI-YM-20261028/blob/3456e5e75cf8d4cd558955a5c7638fba2d68f464/YM/Lattice/Scale.lean\#L167}{\texttt{place}},
\href{https://github.com/yang-mills-group/OpenAI-YM-20261028/blob/3456e5e75cf8d4cd558955a5c7638fba2d68f464/YM/Lattice/Scale.lean\#L173-L182}{\texttt{places\_period}},
\href{https://github.com/yang-mills-group/OpenAI-YM-20261028/blob/3456e5e75cf8d4cd558955a5c7638fba2d68f464/YM/Lattice/Scale.lean\#L83-L90}{\texttt{toDepth}},
\href{https://github.com/yang-mills-group/OpenAI-YM-20261028/blob/3456e5e75cf8d4cd558955a5c7638fba2d68f464/YM/Lattice/Scale.lean\#L92-L97}{\texttt{toDepth\_odd}},
\href{https://github.com/yang-mills-group/OpenAI-YM-20261028/blob/3456e5e75cf8d4cd558955a5c7638fba2d68f464/YM/Lattice/Scale.lean\#L191-L198}{\texttt{embeds\_place}},
\href{https://github.com/yang-mills-group/OpenAI-YM-20261028/blob/3456e5e75cf8d4cd558955a5c7638fba2d68f464/YM/Lattice/Scale.lean\#L129-L130}{\texttt{Gentle}},
\href{https://github.com/yang-mills-group/OpenAI-YM-20261028/blob/3456e5e75cf8d4cd558955a5c7638fba2d68f464/YM/Lattice/Scale.lean\#L133-L162}{\texttt{mesh\_inequality}},
\href{https://github.com/yang-mills-group/OpenAI-YM-20261028/blob/3456e5e75cf8d4cd558955a5c7638fba2d68f464/YM/Lattice/Subdivision.lean\#L22-L24}{\texttt{Ratio}},
\href{https://github.com/yang-mills-group/OpenAI-YM-20261028/blob/3456e5e75cf8d4cd558955a5c7638fba2d68f464/YM/Lattice/Subdivision.lean\#L29}{\texttt{Ratio.embed}},
\href{https://github.com/yang-mills-group/OpenAI-YM-20261028/blob/3456e5e75cf8d4cd558955a5c7638fba2d68f464/YM/Lattice/Subdivision.lean\#L67-L75}{\texttt{Ratio.root}},
\href{https://github.com/yang-mills-group/OpenAI-YM-20261028/blob/3456e5e75cf8d4cd558955a5c7638fba2d68f464/YM/Lattice/Subdivision.lean\#L78-L81}{\texttt{Ratio.signed}},
\href{https://github.com/yang-mills-group/OpenAI-YM-20261028/blob/3456e5e75cf8d4cd558955a5c7638fba2d68f464/YM/Lattice/Subdivision.lean\#L85-L88}{\texttt{Ratio.residue}},
\href{https://github.com/yang-mills-group/OpenAI-YM-20261028/blob/3456e5e75cf8d4cd558955a5c7638fba2d68f464/YM/Lattice/Subdivision.lean\#L103-L111}{\texttt{Ratio.assemble\_bijective}},
\href{https://github.com/yang-mills-group/OpenAI-YM-20261028/blob/3456e5e75cf8d4cd558955a5c7638fba2d68f464/YM/Lattice/Subdivision.lean\#L117}{\texttt{Ratio.parent}},
\href{https://github.com/yang-mills-group/OpenAI-YM-20261028/blob/3456e5e75cf8d4cd558955a5c7638fba2d68f464/YM/Lattice/Subdivision.lean\#L154-L157}{\texttt{Ratio.parent\_signed}},
\href{https://github.com/yang-mills-group/OpenAI-YM-20261028/blob/3456e5e75cf8d4cd558955a5c7638fba2d68f464/YM/Lattice/Crossing.lean\#L19-L25}{\texttt{Ratio.half\_box}},
\href{https://github.com/yang-mills-group/OpenAI-YM-20261028/blob/3456e5e75cf8d4cd558955a5c7638fba2d68f464/YM/Lattice/Crossing.lean\#L54-L84}{\texttt{Ratio.parents\_hop}},
\href{https://github.com/yang-mills-group/OpenAI-YM-20261028/blob/3456e5e75cf8d4cd558955a5c7638fba2d68f464/YM/Lattice/Crossing.lean\#L33-L36}{\texttt{Ratio.embed\_parent}},
\href{https://github.com/yang-mills-group/OpenAI-YM-20261028/blob/3456e5e75cf8d4cd558955a5c7638fba2d68f464/YM/Corona/Spend.lean\#L17-L22}{\texttt{Corona.near}},
\href{https://github.com/yang-mills-group/OpenAI-YM-20261028/blob/3456e5e75cf8d4cd558955a5c7638fba2d68f464/YM/Corona/Spend.lean\#L24-L30}{\texttt{Corona.below}},
\href{https://github.com/yang-mills-group/OpenAI-YM-20261028/blob/3456e5e75cf8d4cd558955a5c7638fba2d68f464/YM/Corona/Spend.lean\#L32-L36}{\texttt{Corona.foreshorten}}.

\textbf{Proposition I.1.4 (cover numbers of lattice sets).} For
\(X\subset\Lambda_\ell\) and \(M\in\mathbb N\) let
\(\mathrm{halo}_M(X)=\{y:\exists x\in X,\ |x-y|_\ell\le M\}\). Use the
graph joining sites at sup-distance at most \(1\); diagonal moves are
allowed. Connectivity means that every two vertices can be joined, so
the empty set is allowed. A set \(D\) \emph{covers \(X\) at scale \(M\)}
if \(X\subset\mathrm{halo}_M(D)\) and \(\mathrm{halo}_M(D)\) is
connected in this graph. Let \(\mathrm{cov}_M(X)\) be the least \(|D|\)
over such \(D\); it is finite since \(D=\Lambda_\ell\) covers. For all
\(X,Y,D\subset\Lambda_\ell\) and \(M\in\mathbb N\), the following bounds
hold independently of \(\ell\):

\begin{itemize}
\tightlist
\item
  \begin{enumerate}
  \def\labelenumi{(\alph{enumi})}
  \tightlist
  \item
    \(X\subset Y\) implies \(\mathrm{cov}_M(X)\le\mathrm{cov}_M(Y)\),
    and
    \(\mathrm{cov}_M(X\cup Y)\le\mathrm{cov}_M(X)+\mathrm{cov}_M(Y)\) if
    \(X\cap Y\neq\emptyset\).
  \end{enumerate}
\item
  (b) \(|\mathrm{halo}_M(D)|\le(2M+1)^4|D|\) and
  \(|X|\le(2M+1)^4\,\mathrm{cov}_M(X)\).
\item
  (c) \(\mathrm{halo}_M(D)\) is connected in this graph if and only if
  any two points of \(D\) are joined by a chain in \(D\) with steps of
  sup-distance at most \(2M+1\).
\item
  (d) (diameter) For \(x,y\in X\):
  \(|x-y|_\ell\le(4M+2)\,\mathrm{cov}_M(X)+2M\).
\item
  (e) (change of scale) For \(b,k,m\in\mathbb N\) with \(b\ge1\) and
  \(x\in\Lambda_\ell\),
  \(\mathrm{cov}_b(\mathrm{halo}_{kb}\{x\})\le(2k+1)^4\). If \(D\)
  covers \(X\) at scale \(mb\), then
  \(\mathrm{cov}_b(X)\le(2m+3)^4|D|\).
\item
  (f) (coarsening) Let \(\mathrm{par}:\Lambda_{\ell L}\to\Lambda_\ell\)
  be the parent map of a subdivision of integer ratio \(L\ge1\). Let
  \(X\subset\Lambda_{\ell L}\) and \(M,B,s\in\mathbb N\) with \(B\ge1\)
  and \(M\le2B\). Then
  \[L\cdot\mathrm{cov}_B\big(\mathrm{halo}_{sB}(\mathrm{par}\,X)\big)\le\big(L+10\,\mathrm{cov}_M(X)\big)\,(2s+13)^4,\]
  with the left side computed on the coarse torus. Thus a fine set of
  cover number \(k\) has coarse cover number \(\le(2s+13)^4(1+10k/L)\).
\end{itemize}

\emph{Status:} theorem. \emph{Lean:}
\href{https://github.com/yang-mills-group/OpenAI-YM-20261028/blob/3456e5e75cf8d4cd558955a5c7638fba2d68f464/YM/Lattice/Support.lean\#L43-L44}{\texttt{halo}},
\href{https://github.com/yang-mills-group/OpenAI-YM-20261028/blob/3456e5e75cf8d4cd558955a5c7638fba2d68f464/YM/Lattice/Support.lean\#L52-L55}{\texttt{covers}},
\href{https://github.com/yang-mills-group/OpenAI-YM-20261028/blob/3456e5e75cf8d4cd558955a5c7638fba2d68f464/YM/Lattice/Support.lean\#L69-L71}{\texttt{coverNumber}},
\href{https://github.com/yang-mills-group/OpenAI-YM-20261028/blob/3456e5e75cf8d4cd558955a5c7638fba2d68f464/YM/Corona/Halo.lean\#L10-L12}{\texttt{Corona.cover\_mono}},
\href{https://github.com/yang-mills-group/OpenAI-YM-20261028/blob/3456e5e75cf8d4cd558955a5c7638fba2d68f464/YM/Corona/Halo.lean\#L14-L30}{\texttt{Corona.meets\_sum}},
\href{https://github.com/yang-mills-group/OpenAI-YM-20261028/blob/3456e5e75cf8d4cd558955a5c7638fba2d68f464/YM/Corona/Fill.lean\#L120-L127}{\texttt{Corona.halo\_size}},
\href{https://github.com/yang-mills-group/OpenAI-YM-20261028/blob/3456e5e75cf8d4cd558955a5c7638fba2d68f464/YM/Corona/Fill.lean\#L129-L131}{\texttt{Corona.pays}},
\href{https://github.com/yang-mills-group/OpenAI-YM-20261028/blob/3456e5e75cf8d4cd558955a5c7638fba2d68f464/YM/Corona/Fill.lean\#L69-L93}{\texttt{Corona.ship\_halo}},
\href{https://github.com/yang-mills-group/OpenAI-YM-20261028/blob/3456e5e75cf8d4cd558955a5c7638fba2d68f464/YM/Corona/Fill.lean\#L95-L111}{\texttt{Corona.halo\_ship}},
\href{https://github.com/yang-mills-group/OpenAI-YM-20261028/blob/3456e5e75cf8d4cd558955a5c7638fba2d68f464/YM/Corona/Fit.lean\#L84-L95}{\texttt{Corona.extent}},
\href{https://github.com/yang-mills-group/OpenAI-YM-20261028/blob/3456e5e75cf8d4cd558955a5c7638fba2d68f464/YM/Corona/Halo.lean\#L66-L113}{\texttt{Corona.tile}},
\href{https://github.com/yang-mills-group/OpenAI-YM-20261028/blob/3456e5e75cf8d4cd558955a5c7638fba2d68f464/YM/Corona/Halo.lean\#L115-L125}{\texttt{Corona.tiling}},
\href{https://github.com/yang-mills-group/OpenAI-YM-20261028/blob/3456e5e75cf8d4cd558955a5c7638fba2d68f464/YM/Corona/Spend.lean\#L73-L105}{\texttt{Corona.collapse}}.

\textbf{Proposition I.1.5 (gauge group, plaquette cost, chord, Haar
measure).} \(G\) is compact and second countable, with its Borel
structure. For \(u\in G\) put
\[\sigma(u)=1-\tfrac1N\mathrm{Re}\,\mathrm{tr}\,u,\qquad \chi(u)=\|u-1\|_F/\sqrt N .\]
Then:

\begin{itemize}
\tightlist
\item
  \(\chi^2=2\sigma\), \(0\le\chi\le2\), \(0\le\sigma\le2\), and
  \(\sigma(1)=\chi(1)=0\);
\item
  \(\chi(uv)\le\chi(u)+\chi(v)\);
\item
  \(\sigma\) and \(\chi\) are continuous and invariant under
  \(u\mapsto u^{-1}\) and \(u\mapsto vuv^{-1}\).
\end{itemize}

The Frobenius norm is invariant under left and right multiplication by
unitaries, and \(\|u\|_F=\sqrt N\). The normalized Haar measure
\(\mu_G\) is invariant under left and right translations and under
inversion; hence \(u\mapsto vuw\) preserves \(\mu_G\). \emph{Status:}
theorem. \emph{Lean:}
\href{https://github.com/yang-mills-group/OpenAI-YM-20261028/blob/3456e5e75cf8d4cd558955a5c7638fba2d68f464/YM/Gauge/MatrixGroup.lean\#L23-L26}{\texttt{Gauge}},
\href{https://github.com/yang-mills-group/OpenAI-YM-20261028/blob/3456e5e75cf8d4cd558955a5c7638fba2d68f464/YM/Gauge/MatrixGroup.lean\#L61-L67}{\texttt{compactGaugeU}},
\href{https://github.com/yang-mills-group/OpenAI-YM-20261028/blob/3456e5e75cf8d4cd558955a5c7638fba2d68f464/YM/Gauge/MatrixGroup.lean\#L112}{\texttt{price}},
\href{https://github.com/yang-mills-group/OpenAI-YM-20261028/blob/3456e5e75cf8d4cd558955a5c7638fba2d68f464/YM/Gauge/MatrixGroup.lean\#L115}{\texttt{chord}},
\href{https://github.com/yang-mills-group/OpenAI-YM-20261028/blob/3456e5e75cf8d4cd558955a5c7638fba2d68f464/YM/Gauge/MatrixGroup.lean\#L117-L126}{\texttt{chord\_price}},
\href{https://github.com/yang-mills-group/OpenAI-YM-20261028/blob/3456e5e75cf8d4cd558955a5c7638fba2d68f464/YM/Gauge/MatrixGroup.lean\#L134-L139}{\texttt{chord\_le\_two}},
\href{https://github.com/yang-mills-group/OpenAI-YM-20261028/blob/3456e5e75cf8d4cd558955a5c7638fba2d68f464/YM/Gauge/MatrixGroup.lean\#L141-L144}{\texttt{price\_range}},
\href{https://github.com/yang-mills-group/OpenAI-YM-20261028/blob/3456e5e75cf8d4cd558955a5c7638fba2d68f464/YM/Gauge/MatrixGroup.lean\#L146-L155}{\texttt{chord\_triangle}},
\href{https://github.com/yang-mills-group/OpenAI-YM-20261028/blob/3456e5e75cf8d4cd558955a5c7638fba2d68f464/YM/Gauge/MatrixGroup.lean\#L157-L158}{\texttt{price\_inv}},
\href{https://github.com/yang-mills-group/OpenAI-YM-20261028/blob/3456e5e75cf8d4cd558955a5c7638fba2d68f464/YM/Gauge/MatrixGroup.lean\#L160-L164}{\texttt{price\_conj}},
\href{https://github.com/yang-mills-group/OpenAI-YM-20261028/blob/3456e5e75cf8d4cd558955a5c7638fba2d68f464/YM/Gauge/MatrixGroup.lean\#L166-L167}{\texttt{chord\_inv}},
\href{https://github.com/yang-mills-group/OpenAI-YM-20261028/blob/3456e5e75cf8d4cd558955a5c7638fba2d68f464/YM/Gauge/MatrixGroup.lean\#L169-L170}{\texttt{chord\_conj}},
\href{https://github.com/yang-mills-group/OpenAI-YM-20261028/blob/3456e5e75cf8d4cd558955a5c7638fba2d68f464/YM/Gauge/MatrixGroup.lean\#L45-L48}{\texttt{mul\_frob\_eq}},
\href{https://github.com/yang-mills-group/OpenAI-YM-20261028/blob/3456e5e75cf8d4cd558955a5c7638fba2d68f464/YM/Gauge/MatrixGroup.lean\#L187}{\texttt{probHaar}},
\href{https://github.com/yang-mills-group/OpenAI-YM-20261028/blob/3456e5e75cf8d4cd558955a5c7638fba2d68f464/YM/Gauge/MatrixGroup.lean\#L193-L200}{\texttt{probHaar\_rightInv}},
\href{https://github.com/yang-mills-group/OpenAI-YM-20261028/blob/3456e5e75cf8d4cd558955a5c7638fba2d68f464/YM/Gauge/MatrixGroup.lean\#L209-L210}{\texttt{preserve\_inv}},
\href{https://github.com/yang-mills-group/OpenAI-YM-20261028/blob/3456e5e75cf8d4cd558955a5c7638fba2d68f464/YM/Gauge/MatrixGroup.lean\#L212-L214}{\texttt{preserve\_two\_sided}}.

\textbf{Proposition I.1.6 (Lie algebra, colours, Casimir constant, slope
constant).}

\begin{enumerate}
\def\labelenumi{(\alph{enumi})}
\item
  \(\mathfrak g=\{X\in M_N(\mathbb C):e^{tX}\in G\ \forall t\in\mathbb R\}\)
  is a closed real subspace of skew-Hermitian matrices. It is closed
  under \([X,Y]=XY-YX\) and under \(X\mapsto uXu^{-1}\), \(u\in G\).
  \emph{Tangent criterion:} if \(f:\mathbb R\to M_N(\mathbb C)\) takes
  values in \(G\), \(f(0)=1\), and \(f\) is differentiable at \(0\) with
  \(f'(0)=X\), then \(X\in\mathfrak g\).
\item
  \emph{Colours.}
  \(\langle X,Y\rangle_{\mathfrak g}=-\frac1N\mathrm{tr}(XY)\) is real
  on \(\mathfrak g\) and is an \(\mathrm{Ad}(G)\)-invariant inner
  product. Fix an orthonormal basis and write
  \(\mathfrak g\cong\mathbb R^d\), \(c\mapsto X_c\),
  \(d=\dim\mathfrak g\). Then:
\end{enumerate}

\begin{itemize}
\tightlist
\item
  \(\|X_c\|_F=\sqrt N\,|c|\);
\item
  if \(P_{\mathfrak g}\) is the orthogonal projection onto
  \(\mathfrak g\) for the real Frobenius inner product
  \(\mathrm{Re}\,\mathrm{tr}(XY^*)\), define
  \(\Pi:M_N(\mathbb C)\to\mathbb R^d\) by
  \(X_{\Pi A}=P_{\mathfrak g}A\). Then \(\Pi X_c=c\) and
  \(|\Pi A|\le\|A\|_F/\sqrt N\);
\item
  \(\mathrm{Ad}(u)c=\Pi(uX_cu^{-1})\) is an orthogonal representation of
  \(G\);
\item
  \(\mathrm{ad}(c)c'=\Pi[X_c,X_{c'}]\) satisfies
  \(X_{\mathrm{ad}(c)c'}=[X_c,X_{c'}]\);
\item
  \(c\mapsto\exp c:=e^{X_c}\in G\) is continuous.
\end{itemize}

\begin{enumerate}
\def\labelenumi{(\alph{enumi})}
\setcounter{enumi}{2}
\tightlist
\item
  \(G\) is \emph{simple} if two conditions hold: there are
  \(X,Y\in\mathfrak g\) with \([X,Y]\neq0\), and every real subspace
  \(I\subset\mathfrak g\) with \([\mathfrak g,I]\subset I\) is \(0\) or
  \(\mathfrak g\). This agrees with the simplicity hypothesis of the
  main theorem: its Lie algebra is characterized by the same exponential
  condition
  (\href{https://github.com/yang-mills-group/OpenAI-YM-20261028/blob/3456e5e75cf8d4cd558955a5c7638fba2d68f464/WilsonBridge.lean\#L12-L17}{\texttt{GaugeGroup.coreSimple}},
  unit VIII.1).
\end{enumerate}

If \(G\) is simple, then:

\begin{itemize}
\tightlist
\item
  (d) for every symmetric real bilinear form \(B\) on \(\mathbb R^d\)
  satisfying \(B(\mathrm{ad}(x)y,z)=-B(y,\mathrm{ad}(x)z)\) for all
  \(x,y,z\), there is \(c\in\mathbb R\) such that
  \(B(x,y)=c\langle x,y\rangle\) for all \(x,y\);
\item
  (e)
  \(-\mathrm{tr}_{\mathbb R^d}(\mathrm{ad}(x)\mathrm{ad}(y))=C_{\mathfrak g}\langle x,y\rangle\)
  for all \(x,y\in\mathbb R^d\), where
  \(C_{\mathfrak g}:=-\frac1d\sum_{i=1}^d\mathrm{tr}_{\mathbb R^d}(\mathrm{ad}(b_i)^2)>0\)
  and \((b_i)_{i=1}^d\) is the standard orthonormal basis;
\item
  (f) the \emph{slope constant} \(\kappa:=11C_{\mathfrak g}/(24\pi^2)\)
  is positive.
\end{itemize}

\emph{Status:} theorem; \(C_{\mathfrak g}\) and \(\kappa\) are
definitions. \emph{Lean:}
\href{https://github.com/yang-mills-group/OpenAI-YM-20261028/blob/3456e5e75cf8d4cd558955a5c7638fba2d68f464/YM/Gauge/Lie.lean\#L126-L135}{\texttt{lie}},
\href{https://github.com/yang-mills-group/OpenAI-YM-20261028/blob/3456e5e75cf8d4cd558955a5c7638fba2d68f464/YM/Gauge/Lie.lean\#L141}{\texttt{closed\_lie}},
\href{https://github.com/yang-mills-group/OpenAI-YM-20261028/blob/3456e5e75cf8d4cd558955a5c7638fba2d68f464/YM/Gauge/Lie.lean\#L42-L44}{\texttt{generates\_skew}},
\href{https://github.com/yang-mills-group/OpenAI-YM-20261028/blob/3456e5e75cf8d4cd558955a5c7638fba2d68f464/YM/Gauge/Lie.lean\#L98-L111}{\texttt{tangent\_generates}},
\href{https://github.com/yang-mills-group/OpenAI-YM-20261028/blob/3456e5e75cf8d4cd558955a5c7638fba2d68f464/YM/Gauge/Lie.lean\#L159-L176}{\texttt{bracket\_mem}},
\href{https://github.com/yang-mills-group/OpenAI-YM-20261028/blob/3456e5e75cf8d4cd558955a5c7638fba2d68f464/YM/Gauge/Lie.lean\#L151-L157}{\texttt{conjugate\_mem}},
\href{https://github.com/yang-mills-group/OpenAI-YM-20261028/blob/3456e5e75cf8d4cd558955a5c7638fba2d68f464/YM/Gauge/Lie.lean\#L179-L182}{\texttt{Simple}},
\href{https://github.com/yang-mills-group/OpenAI-YM-20261028/blob/3456e5e75cf8d4cd558955a5c7638fba2d68f464/YM/Gauge/Colors.lean\#L52-L54}{\texttt{colorIso}},
\href{https://github.com/yang-mills-group/OpenAI-YM-20261028/blob/3456e5e75cf8d4cd558955a5c7638fba2d68f464/YM/Gauge/Colors.lean\#L64}{\texttt{colorMat}},
\href{https://github.com/yang-mills-group/OpenAI-YM-20261028/blob/3456e5e75cf8d4cd558955a5c7638fba2d68f464/YM/Gauge/Colors.lean\#L71-L75}{\texttt{trace\_color}},
\href{https://github.com/yang-mills-group/OpenAI-YM-20261028/blob/3456e5e75cf8d4cd558955a5c7638fba2d68f464/YM/Gauge/Colors.lean\#L77-L86}{\texttt{trace\_color\_real}},
\href{https://github.com/yang-mills-group/OpenAI-YM-20261028/blob/3456e5e75cf8d4cd558955a5c7638fba2d68f464/YM/Gauge/Colors.lean\#L91-L94}{\texttt{colorMat\_norm}},
\href{https://github.com/yang-mills-group/OpenAI-YM-20261028/blob/3456e5e75cf8d4cd558955a5c7638fba2d68f464/YM/Gauge/Colors.lean\#L96-L97}{\texttt{projectColor}},
\href{https://github.com/yang-mills-group/OpenAI-YM-20261028/blob/3456e5e75cf8d4cd558955a5c7638fba2d68f464/YM/Gauge/Colors.lean\#L99-L101}{\texttt{colorMat\_project}},
\href{https://github.com/yang-mills-group/OpenAI-YM-20261028/blob/3456e5e75cf8d4cd558955a5c7638fba2d68f464/YM/Gauge/Colors.lean\#L106-L107}{\texttt{project\_colorMat}},
\href{https://github.com/yang-mills-group/OpenAI-YM-20261028/blob/3456e5e75cf8d4cd558955a5c7638fba2d68f464/YM/Gauge/Colors.lean\#L109-L111}{\texttt{project\_norm}},
\href{https://github.com/yang-mills-group/OpenAI-YM-20261028/blob/3456e5e75cf8d4cd558955a5c7638fba2d68f464/YM/Gauge/Colors.lean\#L114-L117}{\texttt{adjoint}},
\href{https://github.com/yang-mills-group/OpenAI-YM-20261028/blob/3456e5e75cf8d4cd558955a5c7638fba2d68f464/YM/Gauge/Colors.lean\#L126-L128}{\texttt{adjoint\_mul}},
\href{https://github.com/yang-mills-group/OpenAI-YM-20261028/blob/3456e5e75cf8d4cd558955a5c7638fba2d68f464/YM/Gauge/Colors.lean\#L130-L134}{\texttt{norm\_adjoint}},
\href{https://github.com/yang-mills-group/OpenAI-YM-20261028/blob/3456e5e75cf8d4cd558955a5c7638fba2d68f464/YM/Gauge/Colors.lean\#L136-L150}{\texttt{ad}},
\href{https://github.com/yang-mills-group/OpenAI-YM-20261028/blob/3456e5e75cf8d4cd558955a5c7638fba2d68f464/YM/Gauge/Colors.lean\#L152-L154}{\texttt{colorMat\_ad}},
\href{https://github.com/yang-mills-group/OpenAI-YM-20261028/blob/3456e5e75cf8d4cd558955a5c7638fba2d68f464/YM/Gauge/Colors.lean\#L160}{\texttt{lieExp}},
\href{https://github.com/yang-mills-group/OpenAI-YM-20261028/blob/3456e5e75cf8d4cd558955a5c7638fba2d68f464/YM/Gauge/Colors.lean\#L165-L170}{\texttt{lieExp\_cont}},
\href{https://github.com/yang-mills-group/OpenAI-YM-20261028/blob/3456e5e75cf8d4cd558955a5c7638fba2d68f464/YM/Gauge/Colors.lean\#L180-L182}{\texttt{casimir}},
\href{https://github.com/yang-mills-group/OpenAI-YM-20261028/blob/3456e5e75cf8d4cd558955a5c7638fba2d68f464/YM/Gauge/Colors.lean\#L184}{\texttt{slopeConstant}},
\href{https://github.com/yang-mills-group/OpenAI-YM-20261028/blob/3456e5e75cf8d4cd558955a5c7638fba2d68f464/YM/Gauge/Casimir.lean\#L99-L125}{\texttt{invariant\_scalar}},
\href{https://github.com/yang-mills-group/OpenAI-YM-20261028/blob/3456e5e75cf8d4cd558955a5c7638fba2d68f464/YM/Gauge/Casimir.lean\#L140-L158}{\texttt{trace\_adjoint}},
\href{https://github.com/yang-mills-group/OpenAI-YM-20261028/blob/3456e5e75cf8d4cd558955a5c7638fba2d68f464/YM/Gauge/Casimir.lean\#L160-L172}{\texttt{casimir\_pos}},
\href{https://github.com/yang-mills-group/OpenAI-YM-20261028/blob/3456e5e75cf8d4cd558955a5c7638fba2d68f464/YM/Gauge/Casimir.lean\#L174-L176}{\texttt{slopeConstant\_pos}}.

\textbf{Proposition I.1.7 (Haar mass of small balls around \(1\)).} For
\(0<r\le1\), \(\mu_G\{u:\chi(u)\le r\}\ge(r/4)^{2N^2}\). The exponent is
the real dimension of \(M_N(\mathbb C)\), not \(\dim G\). \emph{Status:}
theorem (not sharp). \emph{Lean:}
\href{https://github.com/yang-mills-group/OpenAI-YM-20261028/blob/3456e5e75cf8d4cd558955a5c7638fba2d68f464/YM/Gauge/Size.lean\#L28}{\texttt{shell}},
\href{https://github.com/yang-mills-group/OpenAI-YM-20261028/blob/3456e5e75cf8d4cd558955a5c7638fba2d68f464/YM/Gauge/Size.lean\#L80-L130}{\texttt{shell\_ge}}.

\textbf{Proposition I.1.8 (local logarithm, nearest point, curvature of
a group element).} Let \(\mathfrak g^\perp\) be the orthogonal
complement of \(\mathfrak g\) in \(M_N(\mathbb C)\) for
\(\mathrm{Re}\,\mathrm{tr}(XY^*)\).

\begin{enumerate}
\def\labelenumi{(\alph{enumi})}
\tightlist
\item
  There are a map \(\log_G:M_N(\mathbb C)\to\mathbb R^d\), an open
  neighbourhood \(V\) of \(1\) in \(M_N(\mathbb C)\) and an open
  neighbourhood \(W\) of \(0\) in \(\mathbb R^d\) such that:
\end{enumerate}

\begin{itemize}
\tightlist
\item
  \(\log_G\) is real-analytic on \(V\) and \(\log_G(1)=0\);
\item
  \(\log_G(e^{X_c})=c\) for all \(c\in W\);
\item
  \(\exp(\log_G u)=u\) for all \(u\in G\cap V\).
\end{itemize}

It is the first component of the local inverse of
\((c,\nu)\mapsto e^{X_c}e^{\nu}\) on
\(\mathbb R^d\times\mathfrak g^\perp\). Connectedness of \(G\) is not
assumed, and nothing is claimed away from \(1\).

\begin{enumerate}
\def\labelenumi{(\alph{enumi})}
\setcounter{enumi}{1}
\tightlist
\item
  \(u\in G\) is a \emph{nearest point} to \(X\in M_N(\mathbb C)\) if
  \(\|X-u\|_F\le\|X-v\|_F\) for all \(v\in G\). Then:
\end{enumerate}

\begin{itemize}
\tightlist
\item
  nearest points exist;
\item
  for every \(v\in G\), \(u\) is nearest to \(X\) iff \(vu\) is nearest
  to \(vX\), iff \(uv\) is nearest to \(Xv\);
\item
  if \(u\) is nearest to \(X\), then \(u^*X-1\in\mathfrak g^\perp\);
\item
  there is \(\delta>0\), depending only on \(G\), such that \(X\) has
  exactly one nearest point whenever \(\|X-v\|_F<\delta\) for some
  \(v\in G\).
\end{itemize}

On a neighbourhood of \(1\) the nearest point is
\(\mathrm{foc}(X)=\exp(c)\), where \((c,\nu)\) is the local inverse of
\(\Psi(c,\nu)=e^{X_c}(1+\nu)\) at \(X\). As a matrix-valued map,
\(X\mapsto\mathrm{foc}(X)\) is real-analytic at \(1\), with strict
derivative the orthogonal projection \(P_{\mathfrak g}\) onto
\(\mathfrak g\). Consequently, for every \(\varepsilon>0\) there is
\(\delta>0\) such that, for all \(X,Y\in M_N(\mathbb C)\) and
\(u,v\in G\), the conditions \(\|X-1\|_F,\|Y-1\|_F<\delta\) and \(u,v\)
nearest to \(X,Y\) imply
\[\|u-v-P_{\mathfrak g}(X-Y)\|_F\le\varepsilon\|X-Y\|_F .\]

\begin{enumerate}
\def\labelenumi{(\alph{enumi})}
\setcounter{enumi}{2}
\tightlist
\item
  The curvature of \(u\in G\) is
  \(F(u)=\Pi\big(\tfrac12(u-u^{-1})\big)\in\mathbb R^d\), the projected
  anti-Hermitian part. It is continuous and satisfies:
\end{enumerate}

\begin{itemize}
\tightlist
\item
  \(F(1)=0\) and \(F(u^{-1})=-F(u)\);
\item
  \(F(vuv^{-1})=\mathrm{Ad}(v)F(u)\);
\item
  \(|F(u)|\le\chi(u)\);
\item
  \(\frac{d}{dt}F(\exp(tc))\big|_{t=0}=c\).
\end{itemize}

\emph{Status:} theorem. The neighbourhoods and \(\delta\) come from the
inverse function theorem and compactness and are not explicit.
\emph{Lean:}
\href{https://github.com/yang-mills-group/OpenAI-YM-20261028/blob/3456e5e75cf8d4cd558955a5c7638fba2d68f464/YM/Gauge/Chart.lean\#L138}{\texttt{plog}},
\href{https://github.com/yang-mills-group/OpenAI-YM-20261028/blob/3456e5e75cf8d4cd558955a5c7638fba2d68f464/YM/Gauge/Chart.lean\#L106}{\texttt{splitNormal}},
\href{https://github.com/yang-mills-group/OpenAI-YM-20261028/blob/3456e5e75cf8d4cd558955a5c7638fba2d68f464/YM/Gauge/Chart.lean\#L152-L155}{\texttt{plog\_diff}},
\href{https://github.com/yang-mills-group/OpenAI-YM-20261028/blob/3456e5e75cf8d4cd558955a5c7638fba2d68f464/YM/Gauge/Chart.lean\#L157-L162}{\texttt{plog\_left}},
\href{https://github.com/yang-mills-group/OpenAI-YM-20261028/blob/3456e5e75cf8d4cd558955a5c7638fba2d68f464/YM/Gauge/Chart.lean\#L164-L175}{\texttt{plog\_right}},
\href{https://github.com/yang-mills-group/OpenAI-YM-20261028/blob/3456e5e75cf8d4cd558955a5c7638fba2d68f464/YM/Gauge/Nearest.lean\#L26}{\texttt{Nearest}},
\href{https://github.com/yang-mills-group/OpenAI-YM-20261028/blob/3456e5e75cf8d4cd558955a5c7638fba2d68f464/YM/Gauge/Nearest.lean\#L28-L31}{\texttt{nearest\_exists}},
\href{https://github.com/yang-mills-group/OpenAI-YM-20261028/blob/3456e5e75cf8d4cd558955a5c7638fba2d68f464/YM/Gauge/Nearest.lean\#L33-L36}{\texttt{nearest\_left}},
\href{https://github.com/yang-mills-group/OpenAI-YM-20261028/blob/3456e5e75cf8d4cd558955a5c7638fba2d68f464/YM/Gauge/Nearest.lean\#L38-L41}{\texttt{nearest\_right}},
\href{https://github.com/yang-mills-group/OpenAI-YM-20261028/blob/3456e5e75cf8d4cd558955a5c7638fba2d68f464/YM/Gauge/Nearest.lean\#L65-L68}{\texttt{nearest\_normal}},
\href{https://github.com/yang-mills-group/OpenAI-YM-20261028/blob/3456e5e75cf8d4cd558955a5c7638fba2d68f464/YM/Gauge/Nearest.lean\#L91}{\texttt{focus}},
\href{https://github.com/yang-mills-group/OpenAI-YM-20261028/blob/3456e5e75cf8d4cd558955a5c7638fba2d68f464/YM/Gauge/Nearest.lean\#L128-L139}{\texttt{nearest\_is\_focus}},
\href{https://github.com/yang-mills-group/OpenAI-YM-20261028/blob/3456e5e75cf8d4cd558955a5c7638fba2d68f464/YM/Gauge/Nearest.lean\#L95-L98}{\texttt{focus\_diff}},
\href{https://github.com/yang-mills-group/OpenAI-YM-20261028/blob/3456e5e75cf8d4cd558955a5c7638fba2d68f464/YM/Gauge/Nearest.lean\#L141-L153}{\texttt{nearest\_tube}},
\href{https://github.com/yang-mills-group/OpenAI-YM-20261028/blob/3456e5e75cf8d4cd558955a5c7638fba2d68f464/YM/Gauge/Nearest.lean\#L155-L170}{\texttt{focus\_deriv}},
\href{https://github.com/yang-mills-group/OpenAI-YM-20261028/blob/3456e5e75cf8d4cd558955a5c7638fba2d68f464/YM/Gauge/Nearest.lean\#L172-L181}{\texttt{nearest\_approximate}},
\href{https://github.com/yang-mills-group/OpenAI-YM-20261028/blob/3456e5e75cf8d4cd558955a5c7638fba2d68f464/YM/Gauge/Colors.lean\#L173}{\texttt{curvature}},
\href{https://github.com/yang-mills-group/OpenAI-YM-20261028/blob/3456e5e75cf8d4cd558955a5c7638fba2d68f464/YM/Gauge/Colors.lean\#L175-L177}{\texttt{curvature\_inv}},
\href{https://github.com/yang-mills-group/OpenAI-YM-20261028/blob/3456e5e75cf8d4cd558955a5c7638fba2d68f464/YM/Gauge/Curvature.lean\#L55-L60}{\texttt{curvature\_conjugate}},
\href{https://github.com/yang-mills-group/OpenAI-YM-20261028/blob/3456e5e75cf8d4cd558955a5c7638fba2d68f464/YM/Gauge/Curvature.lean\#L62-L75}{\texttt{curvature\_norm}},
\href{https://github.com/yang-mills-group/OpenAI-YM-20261028/blob/3456e5e75cf8d4cd558955a5c7638fba2d68f464/YM/Gauge/Curvature.lean\#L78-L96}{\texttt{leading\_curvature}}.

\textbf{Proposition I.1.9 (the Wilson measure on the torus).} Let
\(\mu_\ell=\mu_G^{\otimes(\Lambda_\ell\times\{0,\dots,3\})}\) be the
product Haar measure on configurations. Let
\[E_\ell(U)=\sum_{x\in\Lambda_\ell}\sum_{i<j}\sigma(U_{(x;i<j)})\] be
the Wilson energy: \(6\ell^4\) terms, each unoriented plaquette counted
once. For every \(\beta\in\mathbb R\) put
\[Z_\ell(\beta)=\int e^{-\beta E_\ell}\,d\mu_\ell,\qquad d\mathbb P_{\ell,\beta}=Z_\ell(\beta)^{-1}e^{-\beta E_\ell}\,d\mu_\ell .\]
For integrable complex \(F\), write
\(\langle F\rangle_{\ell,\beta}=\int F\,d\mathbb P_{\ell,\beta}\). Then:

\begin{itemize}
\tightlist
\item
  \(E_\ell\) is continuous, with \(0\le E_\ell\le12\ell^4\);
\item
  \(Z_\ell(\beta)>0\) and \(\mathbb P_{\ell,\beta}\) is a probability
  measure;
\item
  every continuous complex function, and every bounded measurable
  complex function, is \(\mathbb P_{\ell,\beta}\)-integrable;
\item
  for integrable complex \(F\),
  \(\langle F\rangle_{\ell,\beta}=Z_\ell(\beta)^{-1}\int e^{-\beta E_\ell}F\,d\mu_\ell\).
\end{itemize}

These are the energy and measure of Definitions 1.3--1.4 of the main
statement, with \(h=\mu_G\); the identification is proved in VIII.1
(\href{https://github.com/yang-mills-group/OpenAI-YM-20261028/blob/3456e5e75cf8d4cd558955a5c7638fba2d68f464/WilsonBridge.lean\#L32-L41}{\texttt{energy\_eq}},
\href{https://github.com/yang-mills-group/OpenAI-YM-20261028/blob/3456e5e75cf8d4cd558955a5c7638fba2d68f464/WilsonBridge.lean\#L43-L46}{\texttt{law\_eq}}).
\emph{Status:} definition, with theorems. \emph{Lean:}
\href{https://github.com/yang-mills-group/OpenAI-YM-20261028/blob/3456e5e75cf8d4cd558955a5c7638fba2d68f464/YM/Wilson/Law.lean\#L24}{\texttt{Wilson.Config}},
\href{https://github.com/yang-mills-group/OpenAI-YM-20261028/blob/3456e5e75cf8d4cd558955a5c7638fba2d68f464/YM/Wilson/Law.lean\#L29}{\texttt{energy}},
\href{https://github.com/yang-mills-group/OpenAI-YM-20261028/blob/3456e5e75cf8d4cd558955a5c7638fba2d68f464/YM/Wilson/Law.lean\#L31-L32}{\texttt{energy\_nonneg}},
\href{https://github.com/yang-mills-group/OpenAI-YM-20261028/blob/3456e5e75cf8d4cd558955a5c7638fba2d68f464/YM/Wilson/Law.lean\#L42-L43}{\texttt{energy\_cont}},
\href{https://github.com/yang-mills-group/OpenAI-YM-20261028/blob/3456e5e75cf8d4cd558955a5c7638fba2d68f464/YM/Wilson/Law.lean\#L55}{\texttt{prior}},
\href{https://github.com/yang-mills-group/OpenAI-YM-20261028/blob/3456e5e75cf8d4cd558955a5c7638fba2d68f464/YM/Wilson/Law.lean\#L58}{\texttt{bare}},
\href{https://github.com/yang-mills-group/OpenAI-YM-20261028/blob/3456e5e75cf8d4cd558955a5c7638fba2d68f464/YM/Wilson/Law.lean\#L74}{\texttt{partition}},
\href{https://github.com/yang-mills-group/OpenAI-YM-20261028/blob/3456e5e75cf8d4cd558955a5c7638fba2d68f464/YM/Wilson/Law.lean\#L76-L80}{\texttt{partition\_pos}},
\href{https://github.com/yang-mills-group/OpenAI-YM-20261028/blob/3456e5e75cf8d4cd558955a5c7638fba2d68f464/YM/Wilson/Law.lean\#L87}{\texttt{law}},
\href{https://github.com/yang-mills-group/OpenAI-YM-20261028/blob/3456e5e75cf8d4cd558955a5c7638fba2d68f464/YM/Wilson/Law.lean\#L99}{\texttt{moment}},
\href{https://github.com/yang-mills-group/OpenAI-YM-20261028/blob/3456e5e75cf8d4cd558955a5c7638fba2d68f464/YM/Wilson/Law.lean\#L101-L103}{\texttt{continuous\_integrable}},
\href{https://github.com/yang-mills-group/OpenAI-YM-20261028/blob/3456e5e75cf8d4cd558955a5c7638fba2d68f464/YM/Wilson/Law.lean\#L105-L107}{\texttt{bounded\_integrable}},
\href{https://github.com/yang-mills-group/OpenAI-YM-20261028/blob/3456e5e75cf8d4cd558955a5c7638fba2d68f464/YM/Wilson/Law.lean\#L110-L121}{\texttt{moment\_formula}}.

\textbf{Proposition I.1.10 (invariance of the Wilson measure).} For
every \(\beta\in\mathbb R\), every lattice affine map \(A=(\rho,b)\) and
every \(g:\Lambda_\ell\to G\), the maps \(U\mapsto A^*U\) and
\(U\mapsto U^g\) preserve \(\mu_\ell\), \(E_\ell\) and
\(\mathbb P_{\ell,\beta}\).

This covers translations, axis permutations, and reflections
\(x_a\mapsto b_a-x_a\). When \(\ell\) is even, the reflection planes
pass through sites for even \(b_a\) and bisect links for odd \(b_a\).
Hence
\(\langle F\circ A^*\rangle_{\ell,\beta}=\langle F\rangle_{\ell,\beta}\)
for every integrable complex \(F\). More generally, any measurable map
preserving \(\mu_\ell\) and \(E_\ell\) preserves
\(\mathbb P_{\ell,\beta}\). Reflection positivity is a separate
assertion. \emph{Status:} theorem. \emph{Lean:}
\href{https://github.com/yang-mills-group/OpenAI-YM-20261028/blob/3456e5e75cf8d4cd558955a5c7638fba2d68f464/YM/Wilson/Law.lean\#L45-L50}{\texttt{energy\_pull}},
\href{https://github.com/yang-mills-group/OpenAI-YM-20261028/blob/3456e5e75cf8d4cd558955a5c7638fba2d68f464/YM/Wilson/Law.lean\#L52-L53}{\texttt{energy\_gauge}},
\href{https://github.com/yang-mills-group/OpenAI-YM-20261028/blob/3456e5e75cf8d4cd558955a5c7638fba2d68f464/YM/Wilson/Symmetry.lean\#L45-L67}{\texttt{prior\_pull}},
\href{https://github.com/yang-mills-group/OpenAI-YM-20261028/blob/3456e5e75cf8d4cd558955a5c7638fba2d68f464/YM/Wilson/Symmetry.lean\#L69-L72}{\texttt{prior\_gauge}},
\href{https://github.com/yang-mills-group/OpenAI-YM-20261028/blob/3456e5e75cf8d4cd558955a5c7638fba2d68f464/YM/Wilson/Symmetry.lean\#L75-L91}{\texttt{law\_preserved}},
\href{https://github.com/yang-mills-group/OpenAI-YM-20261028/blob/3456e5e75cf8d4cd558955a5c7638fba2d68f464/YM/Wilson/Symmetry.lean\#L93-L95}{\texttt{law\_pull}},
\href{https://github.com/yang-mills-group/OpenAI-YM-20261028/blob/3456e5e75cf8d4cd558955a5c7638fba2d68f464/YM/Wilson/Symmetry.lean\#L97-L99}{\texttt{law\_gauge}},
\href{https://github.com/yang-mills-group/OpenAI-YM-20261028/blob/3456e5e75cf8d4cd558955a5c7638fba2d68f464/YM/Wilson/Symmetry.lean\#L101-L104}{\texttt{law\_translate}},
\href{https://github.com/yang-mills-group/OpenAI-YM-20261028/blob/3456e5e75cf8d4cd558955a5c7638fba2d68f464/YM/Wilson/Symmetry.lean\#L106-L113}{\texttt{moment\_pull}}.

\textbf{Proposition I.1.11 (continuum symmetry groups).} \(\mathbb R^4\)
carries the Euclidean norm and Lebesgue measure. Write
\(x=(\vec x,x_t)\), with time the last coordinate. Let
\(\mathsf m(x)=x_t^2-|\vec x|^2\), \(\overline V_+=\{|\vec x|\le x_t\}\)
and \(V_+=\{|\vec x|<x_t\}\).

\begin{enumerate}
\def\labelenumi{(\alph{enumi})}
\item
  \emph{Euclidean.} There is a homomorphism \(B_4\to O(4)\) whose image
  acts on coordinates as in Proposition I.1.2. The mirror in axis \(a\)
  maps to the reflection \(x_a\mapsto-x_a\). The real skew matrices form
  \(\mathfrak{so}(4)\), with plane generators
  \(\Omega_{ij}=E_{ij}-E_{ji}\) for \(i<j\), and
  \(t\mapsto e^{t\Omega}\in O(4)\).
\item
  \emph{Lorentz group.}
  \(\mathcal L=\{g\in GL(\mathbb R^4):\det g=1,\ \mathsf m\circ g=\mathsf m,\ g\overline V_+=\overline V_+\}\);
  this is \(SO^+(1,3)\). It is a topological group of
  Lebesgue-preserving maps that preserves \(V_+\), and it contains every
  element of \(SO(4)\) fixing the time coordinate. For a unit spatial
  vector \(e\), the boost
  \[B_e(r):\ x_t\mapsto\cosh r\,x_t+\sinh r\,(\vec x\cdot e),\quad \vec x\mapsto\vec x+\big(\sinh r\,x_t+(\cosh r-1)(\vec x\cdot e)\big)e\]
  is a continuous one-parameter subgroup of \(\mathcal L\). Its
  generator at \(r=0\) is
  \(K_e:(\vec x,x_t)\mapsto\big(x_t\,e,\ \vec x\cdot e\big)\).
\item
  \emph{Dual action and cone.} With \(R(\vec x,x_t)=(-\vec x,x_t)\), put
  \(g^\vee=RgR\). Then \(g^\vee\in\mathcal L\), \((g^\vee)^\vee=g\),
  \(g^\vee=(g^{-1})^{T}\), and
  \(\langle gx,g^\vee p\rangle=\langle x,p\rangle\) for the Euclidean
  pairing. Moreover:
\end{enumerate}

\begin{itemize}
\tightlist
\item
  \(x\in\overline V_+\) iff \(\langle x,p\rangle\ge0\) for all
  \(p\in\overline V_+\);
\item
  if \(x\in V_+\) and \(p\in\overline V_+\setminus0\), then
  \(\langle x,p\rangle>0\);
\item
  \(x\in\overline V_+\) iff \((g^\vee x)_t\ge0\) for every
  \(g\in\mathcal L\).
\end{itemize}

\begin{enumerate}
\def\labelenumi{(\alph{enumi})}
\setcounter{enumi}{3}
\item
  \emph{Poincaré group.} \(\mathcal P=\mathcal L\ltimes\mathbb R^4\)
  consists of pairs \((\Lambda,b)\) and acts by \(x\mapsto\Lambda x+b\),
  with \((\Lambda,b)(\Lambda',b')=(\Lambda\Lambda',\Lambda b'+b)\). It
  is a topological group, determined by its action on points, and
  preserves Lebesgue measure on \(\mathbb R^4\) and, diagonally, on
  \((\mathbb R^4)^n\) for every \(n\in\mathbb N\).
\item
  \emph{Configurations.} For \(n\in\mathbb N\), \(r\in\mathbb R\) and
  \(x\in(\mathbb R^4)^n\), call \(x\):
\end{enumerate}

\begin{itemize}
\tightlist
\item
  \emph{separated} (\(x\in\mathrm{Conf}_n\)) if the \(x_i\) are pairwise
  distinct; \(\mathrm{Conf}_n\) is open;
\item
  \emph{generic} if for every axis the \(n\) coordinates on that axis
  are pairwise distinct (generic implies separated);
\item
  \emph{\(r\)-spaced} if \(|x_i-x_j|_\infty\ge r\) for \(i\ne j\)
  (\(r>0\) implies separated).
\end{itemize}

Lebesgue measure on \((\mathbb R^4)^n\) is the product of the
four-dimensional ones. \emph{Status:} theorem, apart from the
definitions. \emph{Lean:}
\href{https://github.com/yang-mills-group/OpenAI-YM-20261028/blob/3456e5e75cf8d4cd558955a5c7638fba2d68f464/YM/Space/Orthogonal.lean\#L93-L106}{\texttt{Signed.ortho}},
\href{https://github.com/yang-mills-group/OpenAI-YM-20261028/blob/3456e5e75cf8d4cd558955a5c7638fba2d68f464/YM/Space/Orthogonal.lean\#L108-L115}{\texttt{Signed.ortho\_move}},
\href{https://github.com/yang-mills-group/OpenAI-YM-20261028/blob/3456e5e75cf8d4cd558955a5c7638fba2d68f464/YM/Space/Orthogonal.lean\#L117-L120}{\texttt{Signed.inclusionO}},
\href{https://github.com/yang-mills-group/OpenAI-YM-20261028/blob/3456e5e75cf8d4cd558955a5c7638fba2d68f464/YM/Space/Orthogonal.lean\#L124-L126}{\texttt{Signed.mirrored}},
\href{https://github.com/yang-mills-group/OpenAI-YM-20261028/blob/3456e5e75cf8d4cd558955a5c7638fba2d68f464/YM/Space/Orthogonal.lean\#L132-L142}{\texttt{Ortho.skewSpace}},
\href{https://github.com/yang-mills-group/OpenAI-YM-20261028/blob/3456e5e75cf8d4cd558955a5c7638fba2d68f464/YM/Space/Orthogonal.lean\#L146-L148}{\texttt{Ortho.oneParameter}},
\href{https://github.com/yang-mills-group/OpenAI-YM-20261028/blob/3456e5e75cf8d4cd558955a5c7638fba2d68f464/YM/Space/Orthogonal.lean\#L152-L154}{\texttt{Ortho.plane}},
\href{https://github.com/yang-mills-group/OpenAI-YM-20261028/blob/3456e5e75cf8d4cd558955a5c7638fba2d68f464/YM/Mink/Lorentz.lean\#L32-L44}{\texttt{Mink.Proper}},
\href{https://github.com/yang-mills-group/OpenAI-YM-20261028/blob/3456e5e75cf8d4cd558955a5c7638fba2d68f464/YM/Mink/Lorentz.lean\#L75-L78}{\texttt{Lorentz.preserves}},
\href{https://github.com/yang-mills-group/OpenAI-YM-20261028/blob/3456e5e75cf8d4cd558955a5c7638fba2d68f464/YM/Mink/Lorentz.lean\#L54-L60}{\texttt{Lorentz.stays}},
\href{https://github.com/yang-mills-group/OpenAI-YM-20261028/blob/3456e5e75cf8d4cd558955a5c7638fba2d68f464/YM/Mink/Poincare.lean\#L18-L29}{\texttt{Lorentz.room}},
\href{https://github.com/yang-mills-group/OpenAI-YM-20261028/blob/3456e5e75cf8d4cd558955a5c7638fba2d68f464/YM/Mink/Boost.lean\#L92-L101}{\texttt{Ray.boost}},
\href{https://github.com/yang-mills-group/OpenAI-YM-20261028/blob/3456e5e75cf8d4cd558955a5c7638fba2d68f464/YM/Mink/Boost.lean\#L114-L117}{\texttt{Ray.generator}},
\href{https://github.com/yang-mills-group/OpenAI-YM-20261028/blob/3456e5e75cf8d4cd558955a5c7638fba2d68f464/YM/Mink/Lorentz.lean\#L95-L103}{\texttt{Lorentz.dual}},
\href{https://github.com/yang-mills-group/OpenAI-YM-20261028/blob/3456e5e75cf8d4cd558955a5c7638fba2d68f464/YM/Mink/Lorentz.lean\#L120-L122}{\texttt{Lorentz.dual\_sent}},
\href{https://github.com/yang-mills-group/OpenAI-YM-20261028/blob/3456e5e75cf8d4cd558955a5c7638fba2d68f464/YM/Mink/Lorentz.lean\#L132-L135}{\texttt{Lorentz.dual\_adjoint}},
\href{https://github.com/yang-mills-group/OpenAI-YM-20261028/blob/3456e5e75cf8d4cd558955a5c7638fba2d68f464/YM/Mink/Geometry.lean\#L113-L129}{\texttt{cone\_test}},
\href{https://github.com/yang-mills-group/OpenAI-YM-20261028/blob/3456e5e75cf8d4cd558955a5c7638fba2d68f464/YM/Mink/Geometry.lean\#L131-L141}{\texttt{strict\_pair}},
\href{https://github.com/yang-mills-group/OpenAI-YM-20261028/blob/3456e5e75cf8d4cd558955a5c7638fba2d68f464/YM/Mink/Boost.lean\#L135-L158}{\texttt{Lorentz.from\_half}},
\href{https://github.com/yang-mills-group/OpenAI-YM-20261028/blob/3456e5e75cf8d4cd558955a5c7638fba2d68f464/YM/Mink/Poincare.lean\#L32-L34}{\texttt{Poincare}},
\href{https://github.com/yang-mills-group/OpenAI-YM-20261028/blob/3456e5e75cf8d4cd558955a5c7638fba2d68f464/YM/Mink/Poincare.lean\#L44-L47}{\texttt{Poincare.by\_placement}},
\href{https://github.com/yang-mills-group/OpenAI-YM-20261028/blob/3456e5e75cf8d4cd558955a5c7638fba2d68f464/YM/Mink/Poincare.lean\#L117-L121}{\texttt{Poincare.together}},
\href{https://github.com/yang-mills-group/OpenAI-YM-20261028/blob/3456e5e75cf8d4cd558955a5c7638fba2d68f464/YM/Space/Configurations.lean\#L63}{\texttt{conf}},
\href{https://github.com/yang-mills-group/OpenAI-YM-20261028/blob/3456e5e75cf8d4cd558955a5c7638fba2d68f464/YM/Space/Configurations.lean\#L66}{\texttt{generic}},
\href{https://github.com/yang-mills-group/OpenAI-YM-20261028/blob/3456e5e75cf8d4cd558955a5c7638fba2d68f464/YM/Space/Configurations.lean\#L70}{\texttt{spaced}},
\href{https://github.com/yang-mills-group/OpenAI-YM-20261028/blob/3456e5e75cf8d4cd558955a5c7638fba2d68f464/YM/Space/Configurations.lean\#L83-L90}{\texttt{conf\_open}},
\href{https://github.com/yang-mills-group/OpenAI-YM-20261028/blob/3456e5e75cf8d4cd558955a5c7638fba2d68f464/YM/Space/Configurations.lean\#L153-L180}{\texttt{Points.volume\_product}}.

\textbf{Proposition I.1.12 (boxes, axis routes and slabs on the cover).}
For \(r\in\mathbb N\) let \(B_r=\{-r,\dots,r\}^4\subset\mathbb Z^4\);
then \(|B_r|=(2r+1)^4\) and \(\rho B_r=B_r\) for \(\rho\in B_4\). For
\(x,y\in\mathbb Z^4\) and an ordering \(o\) of the four axes, the
\emph{axis route} \(\mathrm{route}_o(x,y)\) goes from \(x\) in straight
legs along the axes in the order \(o\), adjusting one coordinate at a
time. It ends at \(y\), and:

\begin{itemize}
\tightlist
\item
  reversing the order gives the reverse of \(\mathrm{route}_o(y,x)\);
\item
  the image of \(\mathrm{route}_o(x,y)\) under \(A=(\rho,b)\) is
  \(\mathrm{route}_{\pi\circ o}(Ax,Ay)\).
\end{itemize}

For an anchor \(a\in\mathbb Z^4\) and a set \(J\) of axes, define
\(\mathrm{slab}(r,a,J)=\{y\in B_r:y_i=a_i\text{ for }i\notin J\}\). If
\(i\notin J\), then, for any function \(f\) with values in a commutative
additive monoid,

\[\sum_{y\in\mathrm{slab}(r,a,J\cup\{i\})}f(y)=\sum_{z=-r}^r\ \sum_{y\in\mathrm{slab}(r,a[i\leftarrow z],J)}f(y).\]

For \(z,w\in\{-r,\dots,r\}\), translation by \((w-z)e_i\) maps the layer
with pinned coordinate \(z\) bijectively onto the layer with pinned
coordinate \(w\). Write \(\mathrm{straight}_j(d)\) for the signed leg of
displacement \(de_j\), \(d\in\mathbb Z\). For every cover configuration
\(U\),

\[U(x;\mathrm{straight}_j(d+1))=U(x;\mathrm{straight}_j(d))\,U(x+de_j,j).\]
\emph{Status:} theorem. \emph{Lean:}
\href{https://github.com/yang-mills-group/OpenAI-YM-20261028/blob/3456e5e75cf8d4cd558955a5c7638fba2d68f464/YM/Lattice/Box.lean\#L13}{\texttt{box}},
\href{https://github.com/yang-mills-group/OpenAI-YM-20261028/blob/3456e5e75cf8d4cd558955a5c7638fba2d68f464/YM/Lattice/Box.lean\#L20-L31}{\texttt{box\_covariant}},
\href{https://github.com/yang-mills-group/OpenAI-YM-20261028/blob/3456e5e75cf8d4cd558955a5c7638fba2d68f464/YM/Corona/Fill.lean\#L117-L118}{\texttt{Corona.box\_size}},
\href{https://github.com/yang-mills-group/OpenAI-YM-20261028/blob/3456e5e75cf8d4cd558955a5c7638fba2d68f464/YM/Lattice/Box.lean\#L44}{\texttt{route}},
\href{https://github.com/yang-mills-group/OpenAI-YM-20261028/blob/3456e5e75cf8d4cd558955a5c7638fba2d68f464/YM/Lattice/Box.lean\#L74-L80}{\texttt{route\_end}},
\href{https://github.com/yang-mills-group/OpenAI-YM-20261028/blob/3456e5e75cf8d4cd558955a5c7638fba2d68f464/YM/Lattice/Route.lean\#L36-L37}{\texttt{back\_route}},
\href{https://github.com/yang-mills-group/OpenAI-YM-20261028/blob/3456e5e75cf8d4cd558955a5c7638fba2d68f464/YM/Lattice/Route.lean\#L55-L61}{\texttt{send\_route}},
\href{https://github.com/yang-mills-group/OpenAI-YM-20261028/blob/3456e5e75cf8d4cd558955a5c7638fba2d68f464/YM/Lattice/Foliation.lean\#L18-L19}{\texttt{slab}},
\href{https://github.com/yang-mills-group/OpenAI-YM-20261028/blob/3456e5e75cf8d4cd558955a5c7638fba2d68f464/YM/Lattice/Foliation.lean\#L62-L71}{\texttt{slab\_sum}},
\href{https://github.com/yang-mills-group/OpenAI-YM-20261028/blob/3456e5e75cf8d4cd558955a5c7638fba2d68f464/YM/Lattice/Foliation.lean\#L52-L60}{\texttt{slab\_update\_image}},
\href{https://github.com/yang-mills-group/OpenAI-YM-20261028/blob/3456e5e75cf8d4cd558955a5c7638fba2d68f464/YM/Lattice/Foliation.lean\#L91-L106}{\texttt{carry\_succ}}.

\subsubsection{2. What it uses}\label{what-it-uses}

No result from another unit is needed. The definitions and estimates are
used throughout the proof. Main consumers include:

\begin{itemize}
\tightlist
\item
  the nearest-point retraction (I.1.8(b)): gauge fixing and alignment in
  III.1;
\item
  the invariant-form lemma (I.1.6(d)): the one-loop computation in III.7
  and the stress-tensor forms in VI.1--VI.2;
\item
  the coarsening estimate (I.1.4(f)): III.5 and III.10;
\item
  the mesh inequality: V.1;
\item
  invariance of the law: IV.1, III.13 and VII.2;
\item
  positivity of \(C_{\mathfrak g}\): VIII.1 (the slope in (Reg1)) and
  VII.2;
\item
  the Haar-ball bound (I.1.7): a product lower bound for configurations
  whose links are near \(1\)
  (\href{https://github.com/yang-mills-group/OpenAI-YM-20261028/blob/3456e5e75cf8d4cd558955a5c7638fba2d68f464/YM/Vac/Shelter.lean\#L39-L45}{\texttt{YM.RG.space}}),
  a lower bound for the block-kernel noise normalizer
  (\href{https://github.com/yang-mills-group/OpenAI-YM-20261028/blob/3456e5e75cf8d4cd558955a5c7638fba2d68f464/YM/SectorNext/VACUUM_NORMALIZATION/Noise.lean\#L53-L73}{\texttt{YM.SectorNext.Vacuum.shell\_pZ}}),
  and positivity of the ball mass used in corrected-field tests
  (\href{https://github.com/yang-mills-group/OpenAI-YM-20261028/blob/3456e5e75cf8d4cd558955a5c7638fba2d68f464/YM/Push/CORRECTED_COLUMNS/Program.lean\#L179-L209}{\texttt{YM.Push.Corrected.Tests.Offer.stamp}}).
\end{itemize}

\textbf{Standard mathematics.} The underlying results from Mathlib,
together with references for the mathematical context, are:

\begin{itemize}
\tightlist
\item
  Existence and uniqueness of Haar measure on a compact group, and its
  right and inversion invariance (Halmos, \emph{Measure Theory},
  §§58--60; Folland, \emph{A Course in Abstract Harmonic Analysis}, Ch.
  2).
\item
  Product measures and invariance of product measures under coordinate
  permutations; change of variables for measure-preserving maps; Bochner
  integration.
\item
  Matrix exponential; the inverse function theorem in the real-analytic
  class.
\item
  The local structure of a closed matrix group near \(1\): the Lie
  algebra as the set of generators of one-parameter subgroups, the
  Lie--Trotter limit, and von Neumann's argument that a closed subgroup
  meets \(\exp(\mathfrak g^\perp)\) only at \(1\) near the identity (B.
  C. Hall, \emph{Lie Groups, Lie Algebras, and Representations}, Ch.
  2--3). The local structure is derived here from the exponential
  definition of the Lie algebra.
\item
  Spectral theorem for symmetric operators on a finite-dimensional
  Euclidean space; the Schur-type uniqueness of invariant forms on a
  simple real Lie algebra (Knapp, \emph{Lie Groups Beyond an
  Introduction}, Ch. I).
\item
  Scaling of Lebesgue measure of balls in a finite-dimensional space.
\item
  Lattice gauge theory: Wilson's action (K. Wilson, Phys. Rev.~D 10
  (1974) 2445); gauge-invariant lattice measures and their symmetries
  (Osterwalder--Seiler, Ann. Phys. 110 (1978) 440; Seiler, Lecture Notes
  in Physics 159).
\item
  Minkowski geometry and the restricted Lorentz and Poincaré groups
  (Streater--Wightman, \emph{PCT, Spin and Statistics, and All That},
  Ch. 1).
\end{itemize}

\textbf{Imported constants and assumptions.} The group \(G\) is an
arbitrary closed subgroup of \(U(N)\), possibly disconnected. Simplicity
is assumed only in I.1.6(d)--(f), and agrees with the main theorem's
hypothesis by the bridge cited in I.1.6(c). The local radii depend on
\(G\); the mesh constant also depends on the schedule.

\subsubsection{3. Strategy}\label{strategy}

The torus metric comes from the absolute values of centred
representatives. Odd subdivision ratios give symmetric residue windows,
so the parent map commutes with signed permutations even when the coarse
torus has even side. The schedule separates the physical period from the
mesh; gentle growth of the block factors gives the mesh inequality.
Cover numbers measure the size of lattice sets by connected unions of
balls. Sampling a path through a fine cover at intervals of \(LB\)
yields the coarsening bound \(O(1+k/L)\) for a set of fine cover number
\(k\), with constants independent of the period.

Affine pullback assigns to each link the transport along its image.
Reflections reverse links and shift their base points. The pullback
therefore permutes link variables and inverts some of them, preserving
product Haar measure. Image plaquette holonomies differ by conjugation
and possibly inversion; the cost \(\sigma\) is invariant under both.
Summing over the plaquette bijection proves invariance of the energy and
hence of the Wilson measure.

For the closed matrix group, the tangent criterion follows by
approximating \(e^X\) with \(f(1/n)^n\). It gives the linear and bracket
closure of \(\mathfrak g\). The inverse function theorem, together with
isolation of \(\exp(\mathfrak g^\perp)\) from \(G\) near \(1\), gives a
local logarithm. Applying the same theorem to \(e^{X_c}(1+\nu)\)
identifies the nearest-point retraction and its derivative. Finally, the
spectral theorem reduces an invariant symmetric form to an eigenspace;
real simplicity forces that eigenspace to be all of \(\mathfrak g\).
Applied to the negative Killing form, this gives \(C_{\mathfrak g}>0\).

\subsubsection{4. Steps}\label{steps}

\paragraph{4.1 Torus, subdivision,
schedule}\label{torus-subdivision-schedule}

\textbf{Lemma I.1.a (centred lift and separation).} Proposition I.1.1.
\emph{Proof sketch.} Coordinatewise, \(r_\ell\) is the representative of
minimal absolute value, with the tie sent to \(+\ell/2\). This gives
uniqueness, \(|r_\ell(a+b)_i|\le|r_\ell(a)_i|+|r_\ell(b)_i|\) (hence the
triangle inequality for the maximum), and
\(|r_\ell([y]_\ell)_i|\le|y_i|\). A zero lift reduces to the zero site,
so zero separation implies equality. Signed permutations permute
coordinates and change signs, and the centred absolute value is even,
including at the tie. Thus \(|\rho x|_\ell=|x|_\ell\). There are
\(\ell\) choices in each coordinate, giving \(\ell^4\) sites.
\emph{Lean:}
\href{https://github.com/yang-mills-group/OpenAI-YM-20261028/blob/3456e5e75cf8d4cd558955a5c7638fba2d68f464/YM/Lattice/Site.lean\#L61-L68}{\texttt{lift\_eq\_iff}},
\href{https://github.com/yang-mills-group/OpenAI-YM-20261028/blob/3456e5e75cf8d4cd558955a5c7638fba2d68f464/YM/Lattice/Site.lean\#L96-L101}{\texttt{pdist\_triangle}},
\href{https://github.com/yang-mills-group/OpenAI-YM-20261028/blob/3456e5e75cf8d4cd558955a5c7638fba2d68f464/YM/Lattice/Site.lean\#L107-L115}{\texttt{pdist\_invariant}},
\href{https://github.com/yang-mills-group/OpenAI-YM-20261028/blob/3456e5e75cf8d4cd558955a5c7638fba2d68f464/YM/Corona/Fill.lean\#L32-L39}{\texttt{Corona.lap}}.

\textbf{Lemma I.1.b (subdivision).} Proposition I.1.3(c). \emph{Proof
sketch.}

\begin{itemize}
\tightlist
\item
  \(L\,r_\ell(x)\) lies in the centred window of \(\Lambda_{\ell L}\),
  so \(r_{\ell L}(\iota x)=L\,r_\ell(x)\), and \(\iota\) is injective.
\item
  For every \(y\in\mathbb Z^4\), \(\iota([y]_\ell)=[Ly]_{\ell L}\),
  because changing \(y\) by \(\ell z\) changes \(Ly\) by \(\ell Lz\).
  This gives additivity and \(B_4\)-covariance.
\item
  Applying \(\mathrm{res}\) to an equality of assembled points gives
  equality of their residue coordinates. Injectivity of \(\iota\) then
  gives equality of the coarse coordinates. The count
  \(\ell^4L^4=(\ell L)^4\) makes the assembly map bijective. Since
  \(\iota x\) is assembled from \((x,0)\), its parent is \(x\).
\end{itemize}

\emph{Lean:}
\href{https://github.com/yang-mills-group/OpenAI-YM-20261028/blob/3456e5e75cf8d4cd558955a5c7638fba2d68f464/YM/Lattice/Subdivision.lean\#L31-L42}{\texttt{Ratio.lift\_embed}},
\href{https://github.com/yang-mills-group/OpenAI-YM-20261028/blob/3456e5e75cf8d4cd558955a5c7638fba2d68f464/YM/Lattice/Subdivision.lean\#L44-L52}{\texttt{Ratio.embed\_injective}},
\href{https://github.com/yang-mills-group/OpenAI-YM-20261028/blob/3456e5e75cf8d4cd558955a5c7638fba2d68f464/YM/Lattice/Subdivision.lean\#L67-L75}{\texttt{Ratio.root}},
\href{https://github.com/yang-mills-group/OpenAI-YM-20261028/blob/3456e5e75cf8d4cd558955a5c7638fba2d68f464/YM/Lattice/Subdivision.lean\#L78-L81}{\texttt{Ratio.signed}},
\href{https://github.com/yang-mills-group/OpenAI-YM-20261028/blob/3456e5e75cf8d4cd558955a5c7638fba2d68f464/YM/Lattice/Subdivision.lean\#L93-L95}{\texttt{Ratio.residue\_embed}},
\href{https://github.com/yang-mills-group/OpenAI-YM-20261028/blob/3456e5e75cf8d4cd558955a5c7638fba2d68f464/YM/Lattice/Subdivision.lean\#L103-L111}{\texttt{Ratio.assemble\_bijective}},
\href{https://github.com/yang-mills-group/OpenAI-YM-20261028/blob/3456e5e75cf8d4cd558955a5c7638fba2d68f464/YM/Lattice/Subdivision.lean\#L128-L130}{\texttt{Ratio.unpack\_at}}.

\textbf{Lemma I.1.c (odd blocks and their symmetry).} Proposition
I.1.3(d). \emph{Proof sketch.} For odd \(L\), the window
\((-L/2,L/2]\cap\mathbb Z=[-\frac{L-1}2,\frac{L-1}2]\cap\mathbb Z\) is
symmetric. Hence \(r_L(\rho y)=\rho\,r_L(y)\), the assembly map
intertwines \((\rho,\rho)\) with \(\rho\), and \(\mathrm{par}\) commutes
with \(\rho\). Every residue has exactly one representative in
\(B_{(L-1)/2}\), so the assembly bijection identifies each parent fibre
with the stated translated box. For the crossing rule, write
\(z=\iota(\mathrm{par}\,z)+[q]\) with
\(q=r_L(\mathrm{res}\,z)\in B_{(L-1)/2}\).

\begin{itemize}
\tightlist
\item
  If \(q_i<\frac{L-1}2\), then \(q+e_i\) stays in the block.
\item
  If \(q_i=\frac{L-1}2\), then
  \(z+e_i=\iota(\mathrm{par}\,z+e_i)+[q-(L-1)e_i]\), because
  \(\iota(x+e_i)=\iota(x)+Le_i\).
\end{itemize}

\emph{Lean:}
\href{https://github.com/yang-mills-group/OpenAI-YM-20261028/blob/3456e5e75cf8d4cd558955a5c7638fba2d68f464/YM/Lattice/Subdivision.lean\#L138-L148}{\texttt{Ratio.lift\_signed\_of\_odd}},
\href{https://github.com/yang-mills-group/OpenAI-YM-20261028/blob/3456e5e75cf8d4cd558955a5c7638fba2d68f464/YM/Lattice/Subdivision.lean\#L150-L152}{\texttt{Ratio.assemble\_signed}},
\href{https://github.com/yang-mills-group/OpenAI-YM-20261028/blob/3456e5e75cf8d4cd558955a5c7638fba2d68f464/YM/Lattice/Subdivision.lean\#L154-L157}{\texttt{Ratio.parent\_signed}},
\href{https://github.com/yang-mills-group/OpenAI-YM-20261028/blob/3456e5e75cf8d4cd558955a5c7638fba2d68f464/YM/Lattice/Crossing.lean\#L19-L25}{\texttt{Ratio.half\_box}},
\href{https://github.com/yang-mills-group/OpenAI-YM-20261028/blob/3456e5e75cf8d4cd558955a5c7638fba2d68f464/YM/Lattice/Crossing.lean\#L27-L31}{\texttt{Ratio.half\_lift}},
\href{https://github.com/yang-mills-group/OpenAI-YM-20261028/blob/3456e5e75cf8d4cd558955a5c7638fba2d68f464/YM/Lattice/Crossing.lean\#L54-L84}{\texttt{Ratio.parents\_hop}},
\href{https://github.com/yang-mills-group/OpenAI-YM-20261028/blob/3456e5e75cf8d4cd558955a5c7638fba2d68f464/YM/Lattice/Crossing.lean\#L86-L95}{\texttt{Ratio.internal}}.

\textbf{Lemma I.1.d (metric comparison across scales).} Proposition
I.1.3(e). \emph{Proof sketch.}

\begin{itemize}
\tightlist
\item
  The vector \(r_L(\mathrm{res}\,z)\) has coordinates bounded by
  \(L/2\). Its reduction is \(z-\iota(\mathrm{par}\,z)\), and reduction
  cannot increase the torus norm.
\item
  The lift of \(\iota(x-x')\) is \(L\,r_\ell(x-x')\), which gives the
  second inequality.
\item
  Foreshortening follows from the second inequality and two applications
  of the first.
\end{itemize}

\emph{Lean:}
\href{https://github.com/yang-mills-group/OpenAI-YM-20261028/blob/3456e5e75cf8d4cd558955a5c7638fba2d68f464/YM/Corona/Spend.lean\#L17-L22}{\texttt{Corona.near}},
\href{https://github.com/yang-mills-group/OpenAI-YM-20261028/blob/3456e5e75cf8d4cd558955a5c7638fba2d68f464/YM/Corona/Spend.lean\#L24-L30}{\texttt{Corona.below}},
\href{https://github.com/yang-mills-group/OpenAI-YM-20261028/blob/3456e5e75cf8d4cd558955a5c7638fba2d68f464/YM/Corona/Spend.lean\#L32-L36}{\texttt{Corona.foreshorten}}.

\textbf{Lemma I.1.e (schedule and mesh inequality).} Proposition
I.1.3(a),(b). \emph{Proof sketch.} The product \(D_m\) is odd, strictly
increasing and at least \(3^m\), which gives \(a_m\to0\) and
\(s_m\to\infty\). Placement follows by multiplying
\(-\ell<2r_\ell\le\ell\) by \(a_m\) and using \(a_m\ell(m,p)=P_p\). The
scaled-lift identity gives compatibility under subdivision. Gentleness
gives \(\log L_m<s_m\) for all \(m\ge m_0\), hence \(a_{m+1}>a_m^2\).
Take \(c\) to be the minimum of the finite nonempty set
\(\{1\}\cup\{a_{m+1}/a_m^2:m<m_0\}\); every member is positive.
\emph{Lean:}
\href{https://github.com/yang-mills-group/OpenAI-YM-20261028/blob/3456e5e75cf8d4cd558955a5c7638fba2d68f464/YM/Lattice/Scale.lean\#L60-L65}{\texttt{den\_strict}},
\href{https://github.com/yang-mills-group/OpenAI-YM-20261028/blob/3456e5e75cf8d4cd558955a5c7638fba2d68f464/YM/Lattice/Scale.lean\#L103-L104}{\texttt{a\_step}},
\href{https://github.com/yang-mills-group/OpenAI-YM-20261028/blob/3456e5e75cf8d4cd558955a5c7638fba2d68f464/YM/Lattice/Scale.lean\#L124-L127}{\texttt{scale\_succ}},
\href{https://github.com/yang-mills-group/OpenAI-YM-20261028/blob/3456e5e75cf8d4cd558955a5c7638fba2d68f464/YM/Lattice/Scale.lean\#L184-L189}{\texttt{embeds\_scale}},
\href{https://github.com/yang-mills-group/OpenAI-YM-20261028/blob/3456e5e75cf8d4cd558955a5c7638fba2d68f464/YM/Lattice/Scale.lean\#L133-L162}{\texttt{mesh\_inequality}}.

\paragraph{4.2 Cover numbers}\label{cover-numbers}

\textbf{Lemma I.1.f (elementary properties).} Proposition I.1.4(a),(b).
\emph{Proof sketch.} A cover of a set covers its subsets. For the union,
take optimal covers \(A\) of \(X\) and \(B\) of \(Y\). The halo of
\(A\cup B\) is the union of the two haloes, which is connected through a
common point of \(X\cap Y\). Counting: \(\mathrm{halo}_M(D)\) is the
image of \(D\times B_M\) under \((x,q)\mapsto x+[q]\), and
\(X\subset\mathrm{halo}_M(D)\) for an optimal \(D\). \emph{Lean:}
\href{https://github.com/yang-mills-group/OpenAI-YM-20261028/blob/3456e5e75cf8d4cd558955a5c7638fba2d68f464/YM/Corona/Halo.lean\#L10-L12}{\texttt{Corona.cover\_mono}},
\href{https://github.com/yang-mills-group/OpenAI-YM-20261028/blob/3456e5e75cf8d4cd558955a5c7638fba2d68f464/YM/Corona/Halo.lean\#L14-L30}{\texttt{Corona.meets\_sum}},
\href{https://github.com/yang-mills-group/OpenAI-YM-20261028/blob/3456e5e75cf8d4cd558955a5c7638fba2d68f464/YM/Corona/Fill.lean\#L120-L127}{\texttt{Corona.halo\_size}},
\href{https://github.com/yang-mills-group/OpenAI-YM-20261028/blob/3456e5e75cf8d4cd558955a5c7638fba2d68f464/YM/Corona/Fill.lean\#L129-L131}{\texttt{Corona.pays}}.

\textbf{Lemma I.1.g (connectivity of haloes).} Proposition I.1.4(c).
\emph{Proof sketch.} Use lattice segments
\(t\mapsto x+[\mathrm{clamp}_{[-t,t]}r_\ell(y-x)]\), which move by
\(\le1\) per unit time and reach \(y\) at \(t=|x-y|_\ell\). If
\(|x-y|_\ell\le2M+1\), the segment stays in \(\mathrm{halo}_M\{x,y\}\).
Conversely, a path with steps of sup-distance at most \(1\) in
\(\mathrm{halo}_M(D)\) yields a chain of centres with consecutive
distances \(\le2M+1\), by the triangle inequality. \emph{Lean:}
\href{https://github.com/yang-mills-group/OpenAI-YM-20261028/blob/3456e5e75cf8d4cd558955a5c7638fba2d68f464/YM/Corona/Fill.lean\#L69-L93}{\texttt{Corona.ship\_halo}},
\href{https://github.com/yang-mills-group/OpenAI-YM-20261028/blob/3456e5e75cf8d4cd558955a5c7638fba2d68f464/YM/Corona/Fill.lean\#L95-L111}{\texttt{Corona.halo\_ship}}.

\textbf{Lemma I.1.h (a path through a chained set; diameter).} Let
\(D\subset\Lambda_\ell\) and \(b\in\mathbb N\). Suppose any two points
of \(D\) can be joined in \(D\) by a chain with steps of sup-distance at
most \(b\). Then there are \(n\le2b|D|\) and sites \(R_0,\dots,R_n\)
visiting every point of \(D\), with \(|R_{k+1}-R_k|_\ell\le1\) for
\(k<n\). Consequently Proposition I.1.4(d) holds. \emph{Proof sketch.}
For nonempty \(D\), start at one point. The chain condition gives a
point of \(D\) outside the visited set within distance \(b\) of a
visited point. Insert a detour out and back along a
coordinatewise-clamped segment, costing at most \(2b\) steps. For
\(D=\emptyset\), take a constant path of length \(0\). A cover has chain
bound \(b=2M+1\), so the resulting path has length at most
\((4M+2)\,\mathrm{cov}_M(X)\). Distances along it are bounded by its
length, and the two halo radii add \(2M\). \emph{Lean:}
\href{https://github.com/yang-mills-group/OpenAI-YM-20261028/blob/3456e5e75cf8d4cd558955a5c7638fba2d68f464/YM/Corona/Fit.lean\#L38-L74}{\texttt{Corona.ring\_path}},
\href{https://github.com/yang-mills-group/OpenAI-YM-20261028/blob/3456e5e75cf8d4cd558955a5c7638fba2d68f464/YM/Corona/Fit.lean\#L76-L82}{\texttt{Corona.path\_span}},
\href{https://github.com/yang-mills-group/OpenAI-YM-20261028/blob/3456e5e75cf8d4cd558955a5c7638fba2d68f464/YM/Corona/Fit.lean\#L84-L95}{\texttt{Corona.extent}}.

\textbf{Lemma I.1.i (change of scale).} Proposition I.1.4(e).
\emph{Proof sketch.} The grid \(x+[b\,q]\), \(q\in B_k\), covers
\(\mathrm{halo}_{kb}\{x\}\) at scale \(b\):

\begin{itemize}
\tightlist
\item
  division with remainder puts every point within \(b\) of a grid point;
\item
  consecutive grid points along coordinatewise-clamped paths are at most
  \(b\) apart on the torus, so the grid is \(b\)-chained.
\end{itemize}

For a general \(D\) covering \(X\) at scale \(mb\), I.1.g gives chains
of centres with steps at most \(2mb+1\le2(m+1)b\). The balls
\(\mathrm{halo}_{(m+1)b}\{x\}\) along such a chain overlap. For each
centre, the grid above gives a cover of its ball whose \(b\)-halo is
connected and contains that ball. These connected haloes meet
successively, so their union covers \(X\) and is connected. The union of
their centres has at most \((2m+3)^4|D|\) points. \emph{Lean:}
\href{https://github.com/yang-mills-group/OpenAI-YM-20261028/blob/3456e5e75cf8d4cd558955a5c7638fba2d68f464/YM/Corona/Halo.lean\#L66-L113}{\texttt{Corona.tile}},
\href{https://github.com/yang-mills-group/OpenAI-YM-20261028/blob/3456e5e75cf8d4cd558955a5c7638fba2d68f464/YM/Corona/Halo.lean\#L115-L125}{\texttt{Corona.tiling}}.

\textbf{Lemma I.1.j (coarsening).} Proposition I.1.4(f). \emph{Proof
sketch.}

\begin{enumerate}
\def\labelenumi{\arabic{enumi}.}
\tightlist
\item
  Take an optimal cover \(A\) of \(X\) at scale \(M\). By I.1.h there is
  a fine path \(R_0,\dots,R_n\) visiting \(A\), with steps of
  sup-distance at most \(1\) and \(n\le(4M+2)|A|\).
\item
  Sample its parents every \(LB\) steps:
  \(D'=\{\mathrm{par}\,R_{jLB}:0\le j\le\lfloor n/(LB)\rfloor\}\). By
  foreshortening, parents of path points at most \(LB\) steps apart are
  within \(2B\) on the coarse torus. Thus \(D'\) is \(2B\)-chained,
  every \(\mathrm{par}\,R_k\) with \(0\le k\le n\) lies in
  \(\mathrm{halo}_{2B}(D')\), and \(LB|D'|\le n+LB\).
\item
  A point of \(\mathrm{halo}_{sB}(\mathrm{par}\,X)\) lies within
  \(sB+3B+2B\) of \(D'\). The middle term uses
  \(L|\mathrm{par}\,c-\mathrm{par}\,v|_\ell\le M+L\le3BL\) for
  \(c\in A\) within \(M\) of \(v\in X\). The \((s+5)B\)-halo of \(D'\)
  is connected by I.1.g, since \(D'\) is \(2B\)-chained, so \(D'\) is a
  cover at that scale.
\item
  Apply I.1.i with \(m=s+5\), giving a factor \((2s+13)^4\). Since
  \(M\le2B\) and \(B\ge1\), \(4M+2\le10B\). Thus
  \(LB|D'|\le LB+10B|A|\), and division by \(B\) gives
  \(L|D'|\le L+10|A|\).
\end{enumerate}

\emph{Lean:}
\href{https://github.com/yang-mills-group/OpenAI-YM-20261028/blob/3456e5e75cf8d4cd558955a5c7638fba2d68f464/YM/Corona/Spend.lean\#L38-L70}{\texttt{Corona.fuel}},
\href{https://github.com/yang-mills-group/OpenAI-YM-20261028/blob/3456e5e75cf8d4cd558955a5c7638fba2d68f464/YM/Corona/Spend.lean\#L73-L105}{\texttt{Corona.collapse}}.

\paragraph{4.3 Links, transport,
plaquettes}\label{links-transport-plaquettes}

\textbf{Lemma I.1.k (pullback is a right action).} Proposition I.1.2(a).
\emph{Proof sketch.} Evaluate at a link. The image step of \(\pm e_i\)
under \(\rho'\) and then \(\rho\) is its image under \(\rho\rho'\), and
the affine maps compose as
\((\rho,b)\circ(\rho',b')=(\rho\rho',\rho b'+b)\). \emph{Lean:}
\href{https://github.com/yang-mills-group/OpenAI-YM-20261028/blob/3456e5e75cf8d4cd558955a5c7638fba2d68f464/YM/Lattice/Path.lean\#L47-L50}{\texttt{Step.rotate\_mul}},
\href{https://github.com/yang-mills-group/OpenAI-YM-20261028/blob/3456e5e75cf8d4cd558955a5c7638fba2d68f464/YM/Lattice/Path.lean\#L169-L171}{\texttt{pull\_id}},
\href{https://github.com/yang-mills-group/OpenAI-YM-20261028/blob/3456e5e75cf8d4cd558955a5c7638fba2d68f464/YM/Lattice/Path.lean\#L173-L179}{\texttt{pull\_pull}},
\href{https://github.com/yang-mills-group/OpenAI-YM-20261028/blob/3456e5e75cf8d4cd558955a5c7638fba2d68f464/YM/Lattice/Path.lean\#L184-L186}{\texttt{pull\_one}}.

\textbf{Lemma I.1.l (transport under pullback and gauge; locality).}
Proposition I.1.2(b)--(d). \emph{Proof sketch.} Induction on the walk.
One step: the variable that \(A^*U\) assigns to a step equals the
variable that \(U\) assigns to the image step at the image point. For a
backward step, this is the inverse of the forward variable at the
shifted base point, which is why base points move under reflections.
Gauge: the product telescopes,
\(\prod_kg(x_k)U_kg(x_{k+1})^{-1}=g(x_0)\big(\prod_kU_k\big)g(x_n)^{-1}\).
Locality: each factor reads one link whose endpoints are visited.
\emph{Lean:}
\href{https://github.com/yang-mills-group/OpenAI-YM-20261028/blob/3456e5e75cf8d4cd558955a5c7638fba2d68f464/YM/Lattice/Path.lean\#L150-L158}{\texttt{Step.pull\_read}},
\href{https://github.com/yang-mills-group/OpenAI-YM-20261028/blob/3456e5e75cf8d4cd558955a5c7638fba2d68f464/YM/Lattice/Path.lean\#L160-L167}{\texttt{travel\_pull}},
\href{https://github.com/yang-mills-group/OpenAI-YM-20261028/blob/3456e5e75cf8d4cd558955a5c7638fba2d68f464/YM/Lattice/Path.lean\#L137-L139}{\texttt{Step.gauge\_read}},
\href{https://github.com/yang-mills-group/OpenAI-YM-20261028/blob/3456e5e75cf8d4cd558955a5c7638fba2d68f464/YM/Lattice/Path.lean\#L141-L143}{\texttt{travel\_gauge}},
\href{https://github.com/yang-mills-group/OpenAI-YM-20261028/blob/3456e5e75cf8d4cd558955a5c7638fba2d68f464/YM/Lattice/Path.lean\#L200-L220}{\texttt{travel\_local}}.

\textbf{Lemma I.1.m (plaquettes under symmetries).} Proposition
I.1.2(e). \emph{Proof sketch.}

\begin{itemize}
\tightlist
\item
  By I.1.l, \((A^*U)_p\) is the transport of \(U\) around the image
  square, read from the image of the base point along the image steps.
\item
  Reversing the orientation of a square inverts its holonomy, and
  starting it at the next corner conjugates its holonomy.
\item
  Hence \(\phi\) of the holonomy equals \(\phi\) of the holonomy of the
  positive-axis square based at
  \(Ax-\sum_{k\in\{i,j\}:\,\epsilon_k=-1}e_{\pi(k)}\), with axes in
  increasing order.
\item
  Injectivity of \$A\_\sharp \$: the axis pair is permuted bijectively,
  and the shifted base point and the axis pair determine \(Ax\). On the
  finite torus, injective implies bijective.
\end{itemize}

\emph{Lean:}
\href{https://github.com/yang-mills-group/OpenAI-YM-20261028/blob/3456e5e75cf8d4cd558955a5c7638fba2d68f464/YM/Lattice/Faces.lean\#L62-L66}{\texttt{oriented\_swap}},
\href{https://github.com/yang-mills-group/OpenAI-YM-20261028/blob/3456e5e75cf8d4cd558955a5c7638fba2d68f464/YM/Lattice/Faces.lean\#L69-L73}{\texttt{oriented\_next}},
\href{https://github.com/yang-mills-group/OpenAI-YM-20261028/blob/3456e5e75cf8d4cd558955a5c7638fba2d68f464/YM/Lattice/Faces.lean\#L80-L85}{\texttt{price\_flip}},
\href{https://github.com/yang-mills-group/OpenAI-YM-20261028/blob/3456e5e75cf8d4cd558955a5c7638fba2d68f464/YM/Lattice/Faces.lean\#L91-L107}{\texttt{price\_shift\_pos}},
\href{https://github.com/yang-mills-group/OpenAI-YM-20261028/blob/3456e5e75cf8d4cd558955a5c7638fba2d68f464/YM/Lattice/Faces.lean\#L112-L113}{\texttt{Face.lower}},
\href{https://github.com/yang-mills-group/OpenAI-YM-20261028/blob/3456e5e75cf8d4cd558955a5c7638fba2d68f464/YM/Lattice/Faces.lean\#L119-L126}{\texttt{Face.injective\_send}},
\href{https://github.com/yang-mills-group/OpenAI-YM-20261028/blob/3456e5e75cf8d4cd558955a5c7638fba2d68f464/YM/Lattice/Faces.lean\#L128-L143}{\texttt{Face.send\_price}}.

\textbf{Lemma I.1.n (periodic extension).} Proposition I.1.2(f).
\emph{Proof sketch.} \(E_bU\) evaluates \(U\) at the reduced, shifted
point. Reduction is additive, commutes with \(B_4\) and sends unit steps
to unit steps. So transport, pullback and gauge action commute with
\(E_b\) step by step, and an agreement set on the cover maps to its
image. \emph{Lean:}
\href{https://github.com/yang-mills-group/OpenAI-YM-20261028/blob/3456e5e75cf8d4cd558955a5c7638fba2d68f464/YM/Lattice/Path.lean\#L227-L235}{\texttt{expandAt\_pull}},
\href{https://github.com/yang-mills-group/OpenAI-YM-20261028/blob/3456e5e75cf8d4cd558955a5c7638fba2d68f464/YM/Lattice/Path.lean\#L240-L247}{\texttt{travel\_at}},
\href{https://github.com/yang-mills-group/OpenAI-YM-20261028/blob/3456e5e75cf8d4cd558955a5c7638fba2d68f464/YM/Lattice/Path.lean\#L249-L251}{\texttt{expandAt\_gauge}},
\href{https://github.com/yang-mills-group/OpenAI-YM-20261028/blob/3456e5e75cf8d4cd558955a5c7638fba2d68f464/YM/Lattice/Path.lean\#L253-L261}{\texttt{LocalOn.expandAt}}.

\textbf{Lemma I.1.o (boxes, routes, slabs).} Proposition I.1.12.
\emph{Proof sketch.} The endpoint of a route is computed axis by axis:
each leg changes one coordinate by \(y_a-x_a\), and the axes are
distinct. Reversal and covariance hold leg by leg, since the image of a
straight leg along \(a\) is a straight leg along \(\pi(a)\) of signed
length \(\epsilon_a(y_a-x_a)\). When \(i\notin J\), the slab sum is a
disjoint decomposition by \(y_i\); translating a pinned layer by
\((w-z)e_i\) changes only that coordinate. The transport identity
follows by appending a forward step for \(d\ge0\) and cancelling a final
backward step for \(d<0\). Box covariance follows coordinatewise from
the symmetric interval \([-r,r]\), and its cardinality is the product of
the four interval cardinalities. \emph{Lean:}
\href{https://github.com/yang-mills-group/OpenAI-YM-20261028/blob/3456e5e75cf8d4cd558955a5c7638fba2d68f464/YM/Lattice/Box.lean\#L74-L80}{\texttt{route\_end}},
\href{https://github.com/yang-mills-group/OpenAI-YM-20261028/blob/3456e5e75cf8d4cd558955a5c7638fba2d68f464/YM/Lattice/Route.lean\#L36-L37}{\texttt{back\_route}},
\href{https://github.com/yang-mills-group/OpenAI-YM-20261028/blob/3456e5e75cf8d4cd558955a5c7638fba2d68f464/YM/Lattice/Route.lean\#L55-L61}{\texttt{send\_route}},
\href{https://github.com/yang-mills-group/OpenAI-YM-20261028/blob/3456e5e75cf8d4cd558955a5c7638fba2d68f464/YM/Lattice/Foliation.lean\#L62-L71}{\texttt{slab\_sum}},
\href{https://github.com/yang-mills-group/OpenAI-YM-20261028/blob/3456e5e75cf8d4cd558955a5c7638fba2d68f464/YM/Lattice/Foliation.lean\#L91-L106}{\texttt{carry\_succ}}.

\paragraph{4.4 Gauge group}\label{gauge-group}

\textbf{Lemma I.1.p (cost, chord, Haar).} Proposition I.1.5. \emph{Proof
sketch.}

\begin{itemize}
\tightlist
\item
  \(\|u-1\|_F^2=\mathrm{tr}((u-1)(u-1)^*)=2N-2\,\mathrm{Re}\,\mathrm{tr}\,u\),
  using \(uu^*=1\).
\item
  \(\|u-1\|_F\le\|u\|_F+\|1\|_F=2\sqrt N\).
\item
  \(uv-1=(u-1)+u(v-1)\) and \(\|u(v-1)\|_F=\|v-1\|_F\) give the triangle
  inequality.
\item
  The trace is conjugation invariant, and
  \(\mathrm{tr}\,u^{-1}=\overline{\mathrm{tr}\,u}\).
\item
  A right translate of \(\mu_G\) is again a left Haar probability
  measure, hence equals \(\mu_G\) by uniqueness. Right invariance then
  makes the image under inversion a left Haar probability measure as
  well. Continuity of \(\sigma\) and \(\chi\) follows from continuity of
  the trace and norm.
\end{itemize}

\emph{Lean:}
\href{https://github.com/yang-mills-group/OpenAI-YM-20261028/blob/3456e5e75cf8d4cd558955a5c7638fba2d68f464/YM/Gauge/MatrixGroup.lean\#L34-L43}{\texttt{frob\_sq}},
\href{https://github.com/yang-mills-group/OpenAI-YM-20261028/blob/3456e5e75cf8d4cd558955a5c7638fba2d68f464/YM/Gauge/MatrixGroup.lean\#L55-L59}{\texttt{unitary\_frob}},
\href{https://github.com/yang-mills-group/OpenAI-YM-20261028/blob/3456e5e75cf8d4cd558955a5c7638fba2d68f464/YM/Gauge/MatrixGroup.lean\#L117-L126}{\texttt{chord\_price}},
\href{https://github.com/yang-mills-group/OpenAI-YM-20261028/blob/3456e5e75cf8d4cd558955a5c7638fba2d68f464/YM/Gauge/MatrixGroup.lean\#L146-L155}{\texttt{chord\_triangle}},
\href{https://github.com/yang-mills-group/OpenAI-YM-20261028/blob/3456e5e75cf8d4cd558955a5c7638fba2d68f464/YM/Gauge/MatrixGroup.lean\#L157-L158}{\texttt{price\_inv}},
\href{https://github.com/yang-mills-group/OpenAI-YM-20261028/blob/3456e5e75cf8d4cd558955a5c7638fba2d68f464/YM/Gauge/MatrixGroup.lean\#L193-L200}{\texttt{probHaar\_rightInv}},
\href{https://github.com/yang-mills-group/OpenAI-YM-20261028/blob/3456e5e75cf8d4cd558955a5c7638fba2d68f464/YM/Gauge/MatrixGroup.lean\#L202-L207}{\texttt{probHaar\_invInv}}.

\textbf{Lemma I.1.q (the Lie algebra).} Proposition I.1.6(a).
\emph{Proof sketch.}

\begin{itemize}
\item
  Skewness: differentiate \(f(t)f(t)^*=1\) at \(0\).
\item
  Tangent criterion: if \(f\) takes values in \(G\) with \(f(0)=1\) and
  \(f'(0)=X\), then \(X\) is skew and
  \[\|f(1/n)^n-e^X\|_F\le n\,\|f(1/n)-e^{X/n}\|_F=o(1),\] by the
  telescoping \(\|A^n-B^n\|_F\le n\|A-B\|_F\) for unitary \(A,B\).
  Unitary multiplication is isometric, so the bound has no exponential
  factor. Closedness of \(G\) gives \(e^X\in G\); apply this to
  \(f(t\,\cdot)\).
\item
  Additivity: apply the criterion to \(e^{tX}e^{tY}\).
\item
  Scalar multiplication follows by rescaling the one-parameter subgroup.
  The resulting real subspace is closed because it is finite
  dimensional. Conjugation preserves it since
  \(e^{s(uXu^{-1})}=ue^{sX}u^{-1}\) for \(u\in G\).
\item
  Brackets: \(e^{tX}Ye^{-tX}\in\mathfrak g\) and the closed subspace
  \(\mathfrak g\) contains the limit \([X,Y]\) of its difference
  quotients.
\end{itemize}

\emph{Lean:}
\href{https://github.com/yang-mills-group/OpenAI-YM-20261028/blob/3456e5e75cf8d4cd558955a5c7638fba2d68f464/YM/Gauge/Lie.lean\#L31-L40}{\texttt{tangent\_skew}},
\href{https://github.com/yang-mills-group/OpenAI-YM-20261028/blob/3456e5e75cf8d4cd558955a5c7638fba2d68f464/YM/Gauge/Lie.lean\#L47-L59}{\texttt{frob\_pow\_difference}},
\href{https://github.com/yang-mills-group/OpenAI-YM-20261028/blob/3456e5e75cf8d4cd558955a5c7638fba2d68f464/YM/Gauge/Lie.lean\#L61-L96}{\texttt{tangent\_approx}},
\href{https://github.com/yang-mills-group/OpenAI-YM-20261028/blob/3456e5e75cf8d4cd558955a5c7638fba2d68f464/YM/Gauge/Lie.lean\#L98-L111}{\texttt{tangent\_generates}},
\href{https://github.com/yang-mills-group/OpenAI-YM-20261028/blob/3456e5e75cf8d4cd558955a5c7638fba2d68f464/YM/Gauge/Lie.lean\#L114-L124}{\texttt{generates\_add}},
\href{https://github.com/yang-mills-group/OpenAI-YM-20261028/blob/3456e5e75cf8d4cd558955a5c7638fba2d68f464/YM/Gauge/Lie.lean\#L151-L157}{\texttt{conjugate\_mem}},
\href{https://github.com/yang-mills-group/OpenAI-YM-20261028/blob/3456e5e75cf8d4cd558955a5c7638fba2d68f464/YM/Gauge/Lie.lean\#L159-L176}{\texttt{bracket\_mem}}.

\textbf{Lemma I.1.r (colours, invariant forms, Casimir).} Proposition
I.1.6(b)--(f). \emph{Proof sketch.}

\begin{itemize}
\tightlist
\item
  The colour frame is a Frobenius-orthonormal basis of \(\mathfrak g\)
  scaled by \(\sqrt N\). For skew \(X,Y\),
  \(\mathrm{Re}\,\mathrm{tr}(XY^*)=-\mathrm{tr}(XY)\) and
  \(\mathrm{tr}(XY)\) is real. Orthogonal projection followed by the
  inverse colour frame gives \(\Pi\) and its norm bound. Unitary
  conjugation preserves \(\mathfrak g\) and the Frobenius inner product,
  so it induces the orthogonal adjoint representation. The exponential
  is continuous by the matrix exponential and the induced topology on
  \(G\).
\item
  (d): let \(A\) be the symmetric operator with
  \(\langle Ax,y\rangle=B(x,y)\). Invariance of \(B\) and skewness of
  \(\mathrm{ad}\) give \(A\,\mathrm{ad}_x=\mathrm{ad}_xA\).
  Nonabelianity gives \(d>0\). Choose a real eigenvalue \(c\) of \(A\);
  its nonzero eigenspace is an \(\mathrm{ad}\)-invariant subspace, hence
  an ideal. Simplicity makes it the whole space, so
  \(A=c\cdot\mathrm{id}\).
\item
  (e): \(K(x,y)=-\mathrm{tr}(\mathrm{ad}_x\mathrm{ad}_y)\) is symmetric
  and invariant, by
  \(\mathrm{ad}_{[x,y]}=[\mathrm{ad}_x,\mathrm{ad}_y]\) and cyclicity,
  so \(K=c\langle\cdot,\cdot\rangle\). Taking the trace over an
  orthonormal basis gives \(c=C_{\mathfrak g}\).
\item
  Positivity: \(K(x,x)=\sum_i|\mathrm{ad}_xb_i|^2\), which is \(>0\) for
  some \(x\) because \(\mathfrak g\) is nonabelian. This forces
  \(C_{\mathfrak g}>0\), and the positive factor \(11/(24\pi^2)\) gives
  \(\kappa>0\).
\end{itemize}

\emph{Lean:}
\href{https://github.com/yang-mills-group/OpenAI-YM-20261028/blob/3456e5e75cf8d4cd558955a5c7638fba2d68f464/YM/Gauge/Colors.lean\#L56-L62}{\texttt{inner\_frob\_color}},
\href{https://github.com/yang-mills-group/OpenAI-YM-20261028/blob/3456e5e75cf8d4cd558955a5c7638fba2d68f464/YM/Gauge/Casimir.lean\#L30-L50}{\texttt{Simple.ideals}},
\href{https://github.com/yang-mills-group/OpenAI-YM-20261028/blob/3456e5e75cf8d4cd558955a5c7638fba2d68f464/YM/Gauge/Casimir.lean\#L56-L62}{\texttt{ad\_jacobi}},
\href{https://github.com/yang-mills-group/OpenAI-YM-20261028/blob/3456e5e75cf8d4cd558955a5c7638fba2d68f464/YM/Gauge/Casimir.lean\#L64-L70}{\texttt{ad\_skew}},
\href{https://github.com/yang-mills-group/OpenAI-YM-20261028/blob/3456e5e75cf8d4cd558955a5c7638fba2d68f464/YM/Gauge/Casimir.lean\#L91-L96}{\texttt{minusKilling\_invariant}},
\href{https://github.com/yang-mills-group/OpenAI-YM-20261028/blob/3456e5e75cf8d4cd558955a5c7638fba2d68f464/YM/Gauge/Casimir.lean\#L99-L125}{\texttt{invariant\_scalar}},
\href{https://github.com/yang-mills-group/OpenAI-YM-20261028/blob/3456e5e75cf8d4cd558955a5c7638fba2d68f464/YM/Gauge/Casimir.lean\#L128-L137}{\texttt{minusKilling\_sq}},
\href{https://github.com/yang-mills-group/OpenAI-YM-20261028/blob/3456e5e75cf8d4cd558955a5c7638fba2d68f464/YM/Gauge/Casimir.lean\#L140-L158}{\texttt{trace\_adjoint}},
\href{https://github.com/yang-mills-group/OpenAI-YM-20261028/blob/3456e5e75cf8d4cd558955a5c7638fba2d68f464/YM/Gauge/Casimir.lean\#L160-L172}{\texttt{casimir\_pos}}.

\textbf{Lemma I.1.s (Haar mass of small balls).} Proposition I.1.7.
\emph{Proof sketch.} Points of \(G\) lie on the Frobenius sphere of
radius \(\sqrt N\) in \(M_N(\mathbb C)\cong\mathbb R^{2N^2}\). Take a
finite set \(p\subset G\) whose distinct points have Frobenius distance
strictly greater than \(r\sqrt N\):

\begin{itemize}
\tightlist
\item
  the balls of radius \(r\sqrt N/2\) around its points are disjoint and
  lie in the ball of radius \(2\sqrt N\), so \(|p|(r/4)^{2N^2}\le1\); in
  particular a maximal such set exists;
\item
  if \(p\) is maximal, then \(G=\bigcup_{i\in p}i\cdot\{\chi\le r\}\),
  and left invariance gives \(1\le|p|\,\mu_G\{\chi\le r\}\).
\end{itemize}

\emph{Lean:}
\href{https://github.com/yang-mills-group/OpenAI-YM-20261028/blob/3456e5e75cf8d4cd558955a5c7638fba2d68f464/YM/Gauge/Size.lean\#L80-L130}{\texttt{shell\_ge}}.

\textbf{Lemma I.1.t (local logarithm).} Proposition I.1.8(a).
\emph{Proof sketch.}

\begin{itemize}
\tightlist
\item
  \(\Phi(c,\nu)=e^{X_c}e^{\nu}\) on
  \(\mathbb R^d\times\mathfrak g^\perp\) is analytic, with
  \(D\Phi(0)(c,\nu)=X_c+\nu\) invertible, so it has a local analytic
  inverse.
\item
  \emph{Isolation:} near \(0\), \(e^\nu\in G\) with
  \(\nu\in\mathfrak g^\perp\) forces \(\nu=0\). Otherwise take
  \(\nu_k\to0\), \(\nu_k\neq0\), \(e^{\nu_k}\in G\), with
  \(\nu_k/\|\nu_k\|_F\to Y\) along a subsequence. For \(t\ge0\),
  \(e^{tY}=\lim(e^{\nu_k})^{\lfloor t/\|\nu_k\|_F\rfloor}\in G\) by
  closedness; inverses handle \(t<0\). So
  \(Y\in\mathfrak g\cap\mathfrak g^\perp=0\), contradicting
  \(\|Y\|_F=1\).
\item
  For \(u\in G\) near \(1\), write \(u=\Phi(c,\nu)\); then
  \(e^\nu=e^{-X_c}u\in G\) forces \(\nu=0\).
\end{itemize}

\emph{Lean:}
\href{https://github.com/yang-mills-group/OpenAI-YM-20261028/blob/3456e5e75cf8d4cd558955a5c7638fba2d68f464/YM/Gauge/Chart.lean\#L40-L69}{\texttt{direction\_in\_lie}},
\href{https://github.com/yang-mills-group/OpenAI-YM-20261028/blob/3456e5e75cf8d4cd558955a5c7638fba2d68f464/YM/Gauge/Chart.lean\#L71-L91}{\texttt{isolates}},
\href{https://github.com/yang-mills-group/OpenAI-YM-20261028/blob/3456e5e75cf8d4cd558955a5c7638fba2d68f464/YM/Gauge/Chart.lean\#L124-L133}{\texttt{splitDeriv}},
\href{https://github.com/yang-mills-group/OpenAI-YM-20261028/blob/3456e5e75cf8d4cd558955a5c7638fba2d68f464/YM/Gauge/Chart.lean\#L152-L155}{\texttt{plog\_diff}},
\href{https://github.com/yang-mills-group/OpenAI-YM-20261028/blob/3456e5e75cf8d4cd558955a5c7638fba2d68f464/YM/Gauge/Chart.lean\#L157-L162}{\texttt{plog\_left}},
\href{https://github.com/yang-mills-group/OpenAI-YM-20261028/blob/3456e5e75cf8d4cd558955a5c7638fba2d68f464/YM/Gauge/Chart.lean\#L164-L175}{\texttt{plog\_right}}.

\textbf{Lemma I.1.u (nearest point).} Proposition I.1.8(b). \emph{Proof
sketch.}

\begin{itemize}
\tightlist
\item
  Existence by compactness; equivariance because left and right unitary
  multiplication are Frobenius isometries.
\item
  First-order condition: if \(1\) is nearest to \(X\), then
  \(t\mapsto\|e^{tY}-X\|_F^2\) is minimal at \(0\), so
  \(\langle X-1,Y\rangle=0\) for \(Y\in\mathfrak g\). The general case
  follows by left translation.
\item
  Identification near \(1\): for any choice \(b(X)\) of nearest points,
  \(\|b(X)-1\|_F\le2\|X-1\|_F\). Hence
  \(\nu(X)=b(X)^*X-1\in\mathfrak g^\perp\) tends to \(0\), and
  \(X=\Psi(\log_Gb(X),\nu(X))\) with both arguments near \(0\).
  Injectivity of the local diffeomorphism \(\Psi\) gives
  \(b(X)=\mathrm{foc}(X)\). To obtain one neighbourhood valid for every
  nearest point, choose \(b(X)\) to differ from \(\mathrm{foc}(X)\)
  whenever such a nearest point exists. The same argument forces this
  selection to equal \(\mathrm{foc}\) on a neighbourhood, excluding all
  such exceptions there.
\item
  Uniqueness in a tube: translate by \(v^{-1}\).
\item
  The derivative of \(\mathrm{foc}\) at \(1\) is
  \(X_{(\cdot)}\circ\mathrm{pr}_1\circ D\Psi(0)^{-1}=P_{\mathfrak g}\),
  and the approximation bound is the strict differentiability of
  \(\mathrm{foc}\) at \(1\).
\end{itemize}

\emph{Lean:}
\href{https://github.com/yang-mills-group/OpenAI-YM-20261028/blob/3456e5e75cf8d4cd558955a5c7638fba2d68f464/YM/Gauge/Nearest.lean\#L28-L31}{\texttt{nearest\_exists}},
\href{https://github.com/yang-mills-group/OpenAI-YM-20261028/blob/3456e5e75cf8d4cd558955a5c7638fba2d68f464/YM/Gauge/Nearest.lean\#L33-L36}{\texttt{nearest\_left}},
\href{https://github.com/yang-mills-group/OpenAI-YM-20261028/blob/3456e5e75cf8d4cd558955a5c7638fba2d68f464/YM/Gauge/Nearest.lean\#L65-L68}{\texttt{nearest\_normal}},
\href{https://github.com/yang-mills-group/OpenAI-YM-20261028/blob/3456e5e75cf8d4cd558955a5c7638fba2d68f464/YM/Gauge/Nearest.lean\#L70}{\texttt{tubular}},
\href{https://github.com/yang-mills-group/OpenAI-YM-20261028/blob/3456e5e75cf8d4cd558955a5c7638fba2d68f464/YM/Gauge/Nearest.lean\#L78-L86}{\texttt{tubularDeriv}},
\href{https://github.com/yang-mills-group/OpenAI-YM-20261028/blob/3456e5e75cf8d4cd558955a5c7638fba2d68f464/YM/Gauge/Nearest.lean\#L128-L139}{\texttt{nearest\_is\_focus}},
\href{https://github.com/yang-mills-group/OpenAI-YM-20261028/blob/3456e5e75cf8d4cd558955a5c7638fba2d68f464/YM/Gauge/Nearest.lean\#L141-L153}{\texttt{nearest\_tube}},
\href{https://github.com/yang-mills-group/OpenAI-YM-20261028/blob/3456e5e75cf8d4cd558955a5c7638fba2d68f464/YM/Gauge/Nearest.lean\#L155-L170}{\texttt{focus\_deriv}},
\href{https://github.com/yang-mills-group/OpenAI-YM-20261028/blob/3456e5e75cf8d4cd558955a5c7638fba2d68f464/YM/Gauge/Nearest.lean\#L172-L181}{\texttt{nearest\_approximate}}.

\textbf{Lemma I.1.v (curvature of a group element).} Proposition
I.1.8(c). \emph{Proof sketch.} Conjugation by \(v\in G\) is a Frobenius
isometry preserving \(\mathfrak g\), hence also \(\mathfrak g^\perp\),
so \(\Pi(vAv^{-1})=\mathrm{Ad}(v)\Pi A\). For the bound:
\[|F(u)|\le\tfrac1{2\sqrt N}\|u-u^{-1}\|_F\le\tfrac12\big(\chi(u)+\chi(u^{-1})\big)=\chi(u).\]
The derivative of \(\frac12(e^{tX_c}-e^{-tX_c})\) at \(0\) is \(X_c\),
and \(\Pi X_c=c\). \emph{Lean:}
\href{https://github.com/yang-mills-group/OpenAI-YM-20261028/blob/3456e5e75cf8d4cd558955a5c7638fba2d68f464/YM/Gauge/Curvature.lean\#L35-L47}{\texttt{project\_conjugate}},
\href{https://github.com/yang-mills-group/OpenAI-YM-20261028/blob/3456e5e75cf8d4cd558955a5c7638fba2d68f464/YM/Gauge/Curvature.lean\#L55-L60}{\texttt{curvature\_conjugate}},
\href{https://github.com/yang-mills-group/OpenAI-YM-20261028/blob/3456e5e75cf8d4cd558955a5c7638fba2d68f464/YM/Gauge/Curvature.lean\#L62-L75}{\texttt{curvature\_norm}},
\href{https://github.com/yang-mills-group/OpenAI-YM-20261028/blob/3456e5e75cf8d4cd558955a5c7638fba2d68f464/YM/Gauge/Curvature.lean\#L78-L96}{\texttt{leading\_curvature}}.

\paragraph{4.5 The Wilson measure}\label{the-wilson-measure}

\textbf{Lemma I.1.w (the law).} Proposition I.1.9. \emph{Proof sketch.}
\(E_\ell\) is a finite sum of continuous functions on a compact space.
There are \(6\ell^4\) terms, each in \([0,2]\), giving the energy
bounds. The weight \(e^{-\beta E_\ell}\) is continuous and strictly
positive, so it is integrable and \(Z_\ell>0\). The density
\(e^{-\beta E_\ell}/Z_\ell\) has integral \(1\). Continuous complex
functions are bounded by compactness, and bounded measurable functions
are integrable against a probability measure. The moment formula follows
by integration against the density. \emph{Lean:}
\href{https://github.com/yang-mills-group/OpenAI-YM-20261028/blob/3456e5e75cf8d4cd558955a5c7638fba2d68f464/YM/Wilson/Law.lean\#L71-L72}{\texttt{bare\_integrable}},
\href{https://github.com/yang-mills-group/OpenAI-YM-20261028/blob/3456e5e75cf8d4cd558955a5c7638fba2d68f464/YM/Wilson/Law.lean\#L76-L80}{\texttt{partition\_pos}},
\href{https://github.com/yang-mills-group/OpenAI-YM-20261028/blob/3456e5e75cf8d4cd558955a5c7638fba2d68f464/YM/Wilson/Law.lean\#L110-L121}{\texttt{moment\_formula}}.

\textbf{Lemma I.1.x (invariance of Haar measure, energy and law).}
Proposition I.1.10. \emph{Proof sketch.}

\begin{itemize}
\tightlist
\item
  \emph{Haar.} \(A^*\) is the composite of two maps. The first is the
  relabelling \(U\mapsto U\circ e_A\), where \(e_A\) maps the link
  \((x,i)\) to the positively oriented link carrying its image; this is
  injective, hence a permutation of the finite link set. The second
  inverts the variables on links with \(\epsilon_i=-1\). Permutations of
  factors preserve a product measure, and inversion preserves \(\mu_G\).
  A gauge transformation acts on each factor by
  \(u\mapsto g(x)ug(x+e_i)^{-1}\), which preserves \(\mu_G\).
\item
  \emph{Energy.} Gauge: plaquette holonomies are conjugated. Pullback:
  by I.1.m the energy is re-summed over the bijection \$A\_\sharp \$.
\item
  \emph{Law.} A map preserving \(\mu_\ell\) and \(E_\ell\) preserves the
  density, hence the law.
\end{itemize}

\emph{Lean:}
\href{https://github.com/yang-mills-group/OpenAI-YM-20261028/blob/3456e5e75cf8d4cd558955a5c7638fba2d68f464/YM/Wilson/Symmetry.lean\#L25-L26}{\texttt{sendEdge}},
\href{https://github.com/yang-mills-group/OpenAI-YM-20261028/blob/3456e5e75cf8d4cd558955a5c7638fba2d68f464/YM/Wilson/Symmetry.lean\#L28-L37}{\texttt{sendEdge\_injective}},
\href{https://github.com/yang-mills-group/OpenAI-YM-20261028/blob/3456e5e75cf8d4cd558955a5c7638fba2d68f464/YM/Wilson/Symmetry.lean\#L45-L67}{\texttt{prior\_pull}},
\href{https://github.com/yang-mills-group/OpenAI-YM-20261028/blob/3456e5e75cf8d4cd558955a5c7638fba2d68f464/YM/Wilson/Symmetry.lean\#L69-L72}{\texttt{prior\_gauge}},
\href{https://github.com/yang-mills-group/OpenAI-YM-20261028/blob/3456e5e75cf8d4cd558955a5c7638fba2d68f464/YM/Lattice/Faces.lean\#L49-L51}{\texttt{face\_gauge}},
\href{https://github.com/yang-mills-group/OpenAI-YM-20261028/blob/3456e5e75cf8d4cd558955a5c7638fba2d68f464/YM/Wilson/Law.lean\#L45-L50}{\texttt{energy\_pull}},
\href{https://github.com/yang-mills-group/OpenAI-YM-20261028/blob/3456e5e75cf8d4cd558955a5c7638fba2d68f464/YM/Wilson/Symmetry.lean\#L75-L91}{\texttt{law\_preserved}},
\href{https://github.com/yang-mills-group/OpenAI-YM-20261028/blob/3456e5e75cf8d4cd558955a5c7638fba2d68f464/YM/Wilson/Symmetry.lean\#L106-L113}{\texttt{moment\_pull}}.

\paragraph{4.6 Continuum symmetry
groups}\label{continuum-symmetry-groups}

\textbf{Lemma I.1.y (Euclidean group and configurations).} Proposition
I.1.11(a),(e). \emph{Proof sketch.} The matrix of \(\rho\) has entry
\(\epsilon_a\) at \((\pi(a),a)\); it is orthogonal and multiplies like
\(B_4\). The exponential of a skew matrix is orthogonal. The set
\(\mathrm{Conf}_n\) is a finite intersection of open sets
\(\{x_i\neq x_j\}\). Distinct time coordinates imply distinct points, as
does a positive lower bound on pairwise sup-distances. Splitting off one
Euclidean four-dimensional factor at a time identifies Lebesgue measure
on \((\mathbb R^4)^n\) with the product measure. \emph{Lean:}
\href{https://github.com/yang-mills-group/OpenAI-YM-20261028/blob/3456e5e75cf8d4cd558955a5c7638fba2d68f464/YM/Space/Orthogonal.lean\#L108-L115}{\texttt{Signed.ortho\_move}},
\href{https://github.com/yang-mills-group/OpenAI-YM-20261028/blob/3456e5e75cf8d4cd558955a5c7638fba2d68f464/YM/Space/Orthogonal.lean\#L117-L120}{\texttt{Signed.inclusionO}},
\href{https://github.com/yang-mills-group/OpenAI-YM-20261028/blob/3456e5e75cf8d4cd558955a5c7638fba2d68f464/YM/Space/Configurations.lean\#L72-L74}{\texttt{generic\_conf}},
\href{https://github.com/yang-mills-group/OpenAI-YM-20261028/blob/3456e5e75cf8d4cd558955a5c7638fba2d68f464/YM/Space/Configurations.lean\#L76-L81}{\texttt{spaced\_conf}},
\href{https://github.com/yang-mills-group/OpenAI-YM-20261028/blob/3456e5e75cf8d4cd558955a5c7638fba2d68f464/YM/Space/Configurations.lean\#L83-L90}{\texttt{conf\_open}}.

\textbf{Lemma I.1.z (Lorentz group, boosts, cone, Poincaré group).}
Proposition I.1.11(b)--(d). \emph{Proof sketch.}

\begin{itemize}
\tightlist
\item
  The three defining conditions of \(\mathcal L\) are stable under
  products and inverses, and \(\det=1\) gives Lebesgue invariance.
\item
  Boosts: the addition law \(B_e(r)B_e(r')=B_e(r+r')\) follows from the
  addition formulas for \(\cosh,\sinh\), and \(\mathsf m\) is preserved
  by \(\cosh^2-\sinh^2=1\). For the determinant: \(RB_e(r)R=B_e(-r)\)
  gives \(\det B_e(r)=\det B_e(-r)=\det B_e(r)^{-1}\), and
  \(\det B_e(r)=\det B_e(r/2)^2>0\), so \(\det=1\). The cone is
  preserved by \(|\sinh r|\le\cosh r\). Differentiating the boost
  formula at \(r=0\) gives \(K_e\); continuity follows from \(\cosh\)
  and \(\sinh\).
\item
  Dual action: the identity
  \(\langle gx,g^\vee p\rangle=\langle x,p\rangle\) is polarization of
  \(\mathsf m(x)=\langle x,Rx\rangle\).
\item
  Cone duality: Cauchy--Schwarz gives
  \(\langle x,p\rangle\ge x_tp_t-|\vec x|\,|\vec p|\ge0\) for
  \(x,p\in\overline V_+\). Conversely, test against
  \(p=(-\vec x/|\vec x|,1)\) when \(\vec x\ne0\), and against \((0,1)\)
  otherwise. If \(x\in V_+\) and \(p\in\overline V_+\setminus\{0\}\),
  then \(p_t>0\) makes the bound strict.
\item
  Half-space characterization: for \(x\in\overline V_+\), every
  \(g^\vee x\) remains in that cone and has nonnegative time coordinate.
  If \(x_t<0\), take \(g=1\). Otherwise, if \(x\notin\overline V_+\),
  put \(e=\vec x/|\vec x|\), choose \(x_t/|\vec x|<s<1\), and set
  \(r=\mathrm{artanh}\,s>0\) and \(g=B_e(r)\). Then \(g^\vee=B_e(-r)\)
  and \((g^\vee x)_t=\cosh r\,(x_t-s|\vec x|)<0\).
\item
  The open cone is characterized by \(x_t>0\) and \(\mathsf m(x)>0\),
  and is preserved by the cone and quadratic-form identities. A
  Euclidean rotation fixing time preserves \(|\vec x|\) and hence both
  cones and \(\mathsf m\); determinant \(1\) places it in
  \(\mathcal L\). Affine composition gives the Poincaré group law.
  Evaluation at \(0\) recovers the translation, and subtracting it
  recovers the Lorentz map. Joint continuity follows from linear
  evaluation; the diagonal measure statement follows from preservation
  of each factor.
\end{itemize}

\emph{Lean:}
\href{https://github.com/yang-mills-group/OpenAI-YM-20261028/blob/3456e5e75cf8d4cd558955a5c7638fba2d68f464/YM/Mink/Boost.lean\#L48-L52}{\texttt{Ray.adds}},
\href{https://github.com/yang-mills-group/OpenAI-YM-20261028/blob/3456e5e75cf8d4cd558955a5c7638fba2d68f464/YM/Mink/Boost.lean\#L67-L84}{\texttt{Ray.positive}},
\href{https://github.com/yang-mills-group/OpenAI-YM-20261028/blob/3456e5e75cf8d4cd558955a5c7638fba2d68f464/YM/Mink/Boost.lean\#L86-L90}{\texttt{Ray.staysM}},
\href{https://github.com/yang-mills-group/OpenAI-YM-20261028/blob/3456e5e75cf8d4cd558955a5c7638fba2d68f464/YM/Mink/Lorentz.lean\#L120-L122}{\texttt{Lorentz.dual\_sent}},
\href{https://github.com/yang-mills-group/OpenAI-YM-20261028/blob/3456e5e75cf8d4cd558955a5c7638fba2d68f464/YM/Mink/Geometry.lean\#L113-L129}{\texttt{cone\_test}},
\href{https://github.com/yang-mills-group/OpenAI-YM-20261028/blob/3456e5e75cf8d4cd558955a5c7638fba2d68f464/YM/Mink/Boost.lean\#L135-L158}{\texttt{Lorentz.from\_half}},
\href{https://github.com/yang-mills-group/OpenAI-YM-20261028/blob/3456e5e75cf8d4cd558955a5c7638fba2d68f464/YM/Mink/Poincare.lean\#L101-L102}{\texttt{Poincare.preserving}},
\href{https://github.com/yang-mills-group/OpenAI-YM-20261028/blob/3456e5e75cf8d4cd558955a5c7638fba2d68f464/YM/Mink/Poincare.lean\#L117-L121}{\texttt{Poincare.together}}.

\subsubsection{5. Where the difficulty
lies}\label{where-the-difficulty-lies}

The local logarithm and nearest-point retraction require a
closed-subgroup argument that also handles disconnected groups. The
coarsening estimate must control arbitrary periodic sets uniformly in
the torus side. The invariant-form argument uses simplicity over
\(\mathbb R\); irreducibility of the defining representation is not
assumed. Invariance of the Wilson law reduces to Haar invariance, but
reflections require the correct shifts of link and plaquette base
points.

The representation \(G\subset U(N)\) fixes the cost and the colour inner
product, and hence \(C_{\mathfrak g}\) and \(\kappa\). Each
positive-axis plaquette is counted once. The later running-coupling
estimate says that the ratio of the level-\(m\) coupling to \(s_m\)
tends to \(\kappa\), uniformly in the ultraviolet cutoff \(N'\ge m\) and
the volume index \(p\), for a family satisfying the construction
conditions
(\href{https://github.com/yang-mills-group/OpenAI-YM-20261028/blob/3456e5e75cf8d4cd558955a5c7638fba2d68f464/YM/Pert/Asymptotics/RawAxis.lean\#L15-L30}{\texttt{zeroUp\_bound}}).
Its normalization is \(\kappa=2b_0\), where
\(b_0=11C_{\mathfrak g}/(48\pi^2)\)
(\href{https://github.com/yang-mills-group/OpenAI-YM-20261028/blob/3456e5e75cf8d4cd558955a5c7638fba2d68f464/YM/Pert/Asymptotics/RawAxis.lean\#L14}{\texttt{b0\_normalizes}}).

Even torus sides create tie points in the centred lift. Absolute values,
and hence separation, remain invariant under reflections. Placements
commute with signed permutations on the strict interior of the
fundamental box. Odd subdivision ratios remove this ambiguity from the
residue window, which gives equivariance of the parent map.

\subsubsection{6. Notation}\label{notation}

{\def\LTcaptype{none} % do not increment counter
\begin{longtable}[]{@{}
  >{\raggedright\arraybackslash}p{(\linewidth - 4\tabcolsep) * \real{0.3333}}
  >{\raggedright\arraybackslash}p{(\linewidth - 4\tabcolsep) * \real{0.3333}}
  >{\raggedright\arraybackslash}p{(\linewidth - 4\tabcolsep) * \real{0.3333}}@{}}
\toprule\noalign{}
\begin{minipage}[b]{\linewidth}\raggedright
object
\end{minipage} & \begin{minipage}[b]{\linewidth}\raggedright
code name
\end{minipage} & \begin{minipage}[b]{\linewidth}\raggedright
proposed notation
\end{minipage} \\
\midrule\noalign{}
\endhead
\bottomrule\noalign{}
\endlastfoot
torus, side, cover &
\href{https://github.com/yang-mills-group/OpenAI-YM-20261028/blob/3456e5e75cf8d4cd558955a5c7638fba2d68f464/YM/Lattice/Site.lean\#L19}{\texttt{Site}},
\href{https://github.com/yang-mills-group/OpenAI-YM-20261028/blob/3456e5e75cf8d4cd558955a5c7638fba2d68f464/YM/Lattice/Site.lean\#L17}{\texttt{Side}},
\href{https://github.com/yang-mills-group/OpenAI-YM-20261028/blob/3456e5e75cf8d4cd558955a5c7638fba2d68f464/YM/Lattice/Site.lean\#L20}{\texttt{Cover}}
& \(\Lambda_\ell=(\mathbb Z/\ell\mathbb Z)^4\), \(\ell\),
\(\mathbb Z^4\) \\
reduction, centred lift &
\href{https://github.com/yang-mills-group/OpenAI-YM-20261028/blob/3456e5e75cf8d4cd558955a5c7638fba2d68f464/YM/Lattice/Site.lean\#L35-L38}{\texttt{wrap}},
\href{https://github.com/yang-mills-group/OpenAI-YM-20261028/blob/3456e5e75cf8d4cd558955a5c7638fba2d68f464/YM/Lattice/Site.lean\#L51}{\texttt{lift}}
& \([y]_\ell\), \(r_\ell(x)\) \\
torus norm, separation &
\href{https://github.com/yang-mills-group/OpenAI-YM-20261028/blob/3456e5e75cf8d4cd558955a5c7638fba2d68f464/YM/Lattice/Site.lean\#L78}{\texttt{pdist}},
\href{https://github.com/yang-mills-group/OpenAI-YM-20261028/blob/3456e5e75cf8d4cd558955a5c7638fba2d68f464/YM/Lattice/Site.lean\#L81}{\texttt{separation}}
& \(\lvert x\rvert_\ell\), \(\lvert x-y\rvert_\ell\) \\
signed permutations, affine maps &
\href{https://github.com/yang-mills-group/OpenAI-YM-20261028/blob/3456e5e75cf8d4cd558955a5c7638fba2d68f464/YM/Space/Index.lean\#L23-L26}{\texttt{Signed}},
\href{https://github.com/yang-mills-group/OpenAI-YM-20261028/blob/3456e5e75cf8d4cd558955a5c7638fba2d68f464/YM/Lattice/Path.lean\#L145}{\texttt{affine}}
& \(\rho=(\pi,\epsilon)\in B_4\), \(A=(\rho,b)\) \\
step, walk, transport &
\href{https://github.com/yang-mills-group/OpenAI-YM-20261028/blob/3456e5e75cf8d4cd558955a5c7638fba2d68f464/YM/Lattice/Path.lean\#L23-L26}{\texttt{Step}},
\href{https://github.com/yang-mills-group/OpenAI-YM-20261028/blob/3456e5e75cf8d4cd558955a5c7638fba2d68f464/YM/Lattice/Path.lean\#L27}{\texttt{Walk}},
\href{https://github.com/yang-mills-group/OpenAI-YM-20261028/blob/3456e5e75cf8d4cd558955a5c7638fba2d68f464/YM/Lattice/Path.lean\#L105-L107}{\texttt{travel}}
& \(\pm e_i\), \(w\), \(U(x;w)\) \\
configuration &
\href{https://github.com/yang-mills-group/OpenAI-YM-20261028/blob/3456e5e75cf8d4cd558955a5c7638fba2d68f464/YM/Lattice/Path.lean\#L21}{\texttt{Links}},
\href{https://github.com/yang-mills-group/OpenAI-YM-20261028/blob/3456e5e75cf8d4cd558955a5c7638fba2d68f464/YM/Wilson/Law.lean\#L24}{\texttt{Wilson.Config}}
& \(U\) \\
pullback, translation &
\href{https://github.com/yang-mills-group/OpenAI-YM-20261028/blob/3456e5e75cf8d4cd558955a5c7638fba2d68f464/YM/Lattice/Path.lean\#L147-L148}{\texttt{pull}},
\href{https://github.com/yang-mills-group/OpenAI-YM-20261028/blob/3456e5e75cf8d4cd558955a5c7638fba2d68f464/YM/Lattice/Path.lean\#L182}{\texttt{translate}}
& \(A^*U\), \(\tau_bU\) \\
gauge transformation; invariance; locality &
\href{https://github.com/yang-mills-group/OpenAI-YM-20261028/blob/3456e5e75cf8d4cd558955a5c7638fba2d68f464/YM/Lattice/Path.lean\#L132-L133}{\texttt{gauge}},
\href{https://github.com/yang-mills-group/OpenAI-YM-20261028/blob/3456e5e75cf8d4cd558955a5c7638fba2d68f464/YM/Lattice/Path.lean\#L135}{\texttt{GaugeInvariant}},
\href{https://github.com/yang-mills-group/OpenAI-YM-20261028/blob/3456e5e75cf8d4cd558955a5c7638fba2d68f464/YM/Lattice/Path.lean\#L196-L197}{\texttt{LocalOn}}
& \(U^g\) \\
plaquette, holonomy, image plaquette &
\href{https://github.com/yang-mills-group/OpenAI-YM-20261028/blob/3456e5e75cf8d4cd558955a5c7638fba2d68f464/YM/Lattice/Faces.lean\#L14}{\texttt{Face}},
\href{https://github.com/yang-mills-group/OpenAI-YM-20261028/blob/3456e5e75cf8d4cd558955a5c7638fba2d68f464/YM/Lattice/Faces.lean\#L47}{\texttt{faceValue}},
\href{https://github.com/yang-mills-group/OpenAI-YM-20261028/blob/3456e5e75cf8d4cd558955a5c7638fba2d68f464/YM/Lattice/Faces.lean\#L115-L116}{\texttt{Face.send}}
& \(p=(x;i<j)\), \(U_p\), \(A_\sharp p\) \\
periodic extension &
\href{https://github.com/yang-mills-group/OpenAI-YM-20261028/blob/3456e5e75cf8d4cd558955a5c7638fba2d68f464/YM/Lattice/Path.lean\#L224-L225}{\texttt{expandAt}}
& \(E_bU\) \\
box, axis route, slab &
\href{https://github.com/yang-mills-group/OpenAI-YM-20261028/blob/3456e5e75cf8d4cd558955a5c7638fba2d68f464/YM/Lattice/Box.lean\#L13}{\texttt{box}},
\href{https://github.com/yang-mills-group/OpenAI-YM-20261028/blob/3456e5e75cf8d4cd558955a5c7638fba2d68f464/YM/Lattice/Box.lean\#L44}{\texttt{route}},
\href{https://github.com/yang-mills-group/OpenAI-YM-20261028/blob/3456e5e75cf8d4cd558955a5c7638fba2d68f464/YM/Lattice/Foliation.lean\#L18-L19}{\texttt{slab}}
& \(B_r\), \(\mathrm{route}_o(x,y)\), \(\mathrm{slab}(r,a,J)\) \\
schedule: block factors, periods &
\href{https://github.com/yang-mills-group/OpenAI-YM-20261028/blob/3456e5e75cf8d4cd558955a5c7638fba2d68f464/YM/Lattice/Scale.lean\#L28-L34}{\texttt{Schedule}}
& \((L_m)\), \((P_p)\); the statement writes \(L_p\) for \(P_p\) and
\(b_k\) for \(L_k\) \\
denominator, spacing, side, log-scale &
\href{https://github.com/yang-mills-group/OpenAI-YM-20261028/blob/3456e5e75cf8d4cd558955a5c7638fba2d68f464/YM/Lattice/Scale.lean\#L39-L41}{\texttt{denominator}},
\href{https://github.com/yang-mills-group/OpenAI-YM-20261028/blob/3456e5e75cf8d4cd558955a5c7638fba2d68f464/YM/Lattice/Scale.lean\#L99}{\texttt{a}},
\href{https://github.com/yang-mills-group/OpenAI-YM-20261028/blob/3456e5e75cf8d4cd558955a5c7638fba2d68f464/YM/Lattice/Scale.lean\#L67}{\texttt{side}},
\href{https://github.com/yang-mills-group/OpenAI-YM-20261028/blob/3456e5e75cf8d4cd558955a5c7638fba2d68f464/YM/Lattice/Scale.lean\#L115}{\texttt{scale}}
& \(D_m\), \(a_m\), \(\ell(m,p)\), \(s_m=\log(1/a_m)\) \\
placement &
\href{https://github.com/yang-mills-group/OpenAI-YM-20261028/blob/3456e5e75cf8d4cd558955a5c7638fba2d68f464/YM/Lattice/Scale.lean\#L167}{\texttt{place}}
& \(a_m\,r_\ell(x)\) \\
subdivision &
\href{https://github.com/yang-mills-group/OpenAI-YM-20261028/blob/3456e5e75cf8d4cd558955a5c7638fba2d68f464/YM/Lattice/Subdivision.lean\#L22-L24}{\texttt{Ratio}},
\href{https://github.com/yang-mills-group/OpenAI-YM-20261028/blob/3456e5e75cf8d4cd558955a5c7638fba2d68f464/YM/Lattice/Subdivision.lean\#L29}{\texttt{embed}},
\href{https://github.com/yang-mills-group/OpenAI-YM-20261028/blob/3456e5e75cf8d4cd558955a5c7638fba2d68f464/YM/Lattice/Subdivision.lean\#L85-L88}{\texttt{residue}},
\href{https://github.com/yang-mills-group/OpenAI-YM-20261028/blob/3456e5e75cf8d4cd558955a5c7638fba2d68f464/YM/Lattice/Subdivision.lean\#L117}{\texttt{parent}}
& ratio \(L\), \(\iota\), \(\mathrm{res}\), \(\mathrm{par}\) \\
halo, cover number &
\href{https://github.com/yang-mills-group/OpenAI-YM-20261028/blob/3456e5e75cf8d4cd558955a5c7638fba2d68f464/YM/Lattice/Support.lean\#L43-L44}{\texttt{halo}},
\href{https://github.com/yang-mills-group/OpenAI-YM-20261028/blob/3456e5e75cf8d4cd558955a5c7638fba2d68f464/YM/Lattice/Support.lean\#L69-L71}{\texttt{coverNumber}}
& \(\mathrm{halo}_M(X)\), \(\mathrm{cov}_M(X)\) (as in III.5) \\
gauge group &
\href{https://github.com/yang-mills-group/OpenAI-YM-20261028/blob/3456e5e75cf8d4cd558955a5c7638fba2d68f464/YM/Gauge/MatrixGroup.lean\#L23-L26}{\texttt{Gauge}},
\href{https://github.com/yang-mills-group/OpenAI-YM-20261028/blob/3456e5e75cf8d4cd558955a5c7638fba2d68f464/YM/Gauge/MatrixGroup.lean\#L78}{\texttt{El}},
\href{https://github.com/yang-mills-group/OpenAI-YM-20261028/blob/3456e5e75cf8d4cd558955a5c7638fba2d68f464/YM/Gauge/MatrixGroup.lean\#L94}{\texttt{mat}}
& \(G\subset U(N)\), \(N=n+1\) (the statement writes \(m\)) \\
plaquette cost, chord &
\href{https://github.com/yang-mills-group/OpenAI-YM-20261028/blob/3456e5e75cf8d4cd558955a5c7638fba2d68f464/YM/Gauge/MatrixGroup.lean\#L112}{\texttt{price}},
\href{https://github.com/yang-mills-group/OpenAI-YM-20261028/blob/3456e5e75cf8d4cd558955a5c7638fba2d68f464/YM/Gauge/MatrixGroup.lean\#L115}{\texttt{chord}}
& \(\sigma(u)=1-\frac1N\mathrm{Re}\,\mathrm{tr}\,u\), \(\chi(u)\) \\
Haar probability &
\href{https://github.com/yang-mills-group/OpenAI-YM-20261028/blob/3456e5e75cf8d4cd558955a5c7638fba2d68f464/YM/Gauge/MatrixGroup.lean\#L187}{\texttt{probHaar}}
& \(\mu_G\) \\
Lie algebra, colours, projection &
\href{https://github.com/yang-mills-group/OpenAI-YM-20261028/blob/3456e5e75cf8d4cd558955a5c7638fba2d68f464/YM/Gauge/Lie.lean\#L126-L135}{\texttt{lie}},
\href{https://github.com/yang-mills-group/OpenAI-YM-20261028/blob/3456e5e75cf8d4cd558955a5c7638fba2d68f464/YM/Gauge/Colors.lean\#L43}{\texttt{Color}},
\href{https://github.com/yang-mills-group/OpenAI-YM-20261028/blob/3456e5e75cf8d4cd558955a5c7638fba2d68f464/YM/Gauge/Colors.lean\#L64}{\texttt{colorMat}},
\href{https://github.com/yang-mills-group/OpenAI-YM-20261028/blob/3456e5e75cf8d4cd558955a5c7638fba2d68f464/YM/Gauge/Colors.lean\#L96-L97}{\texttt{projectColor}}
& \(\mathfrak g\cong\mathbb R^d\), \(c\mapsto X_c\), \(\Pi\);
\(X_{\Pi A}=P_{\mathfrak g}A\) \\
adjoint action, bracket, exponential &
\href{https://github.com/yang-mills-group/OpenAI-YM-20261028/blob/3456e5e75cf8d4cd558955a5c7638fba2d68f464/YM/Gauge/Colors.lean\#L114-L117}{\texttt{adjoint}},
\href{https://github.com/yang-mills-group/OpenAI-YM-20261028/blob/3456e5e75cf8d4cd558955a5c7638fba2d68f464/YM/Gauge/Colors.lean\#L136-L150}{\texttt{ad}},
\href{https://github.com/yang-mills-group/OpenAI-YM-20261028/blob/3456e5e75cf8d4cd558955a5c7638fba2d68f464/YM/Gauge/Colors.lean\#L160}{\texttt{lieExp}}
& \(\mathrm{Ad}(u)\), \(\mathrm{ad}(c)\), \(\exp c\) \\
simplicity &
\href{https://github.com/yang-mills-group/OpenAI-YM-20261028/blob/3456e5e75cf8d4cd558955a5c7638fba2d68f464/YM/Gauge/Lie.lean\#L179-L182}{\texttt{Simple}}
& \(G\) simple \\
Casimir, slope constant &
\href{https://github.com/yang-mills-group/OpenAI-YM-20261028/blob/3456e5e75cf8d4cd558955a5c7638fba2d68f464/YM/Gauge/Colors.lean\#L180-L182}{\texttt{casimir}},
\href{https://github.com/yang-mills-group/OpenAI-YM-20261028/blob/3456e5e75cf8d4cd558955a5c7638fba2d68f464/YM/Gauge/Colors.lean\#L184}{\texttt{slopeConstant}}
& \(C_{\mathfrak g}\), \(\kappa=11C_{\mathfrak g}/(24\pi^2)\)
(\(=b_*=2b_0\) of the statement) \\
local logarithm, nearest point &
\href{https://github.com/yang-mills-group/OpenAI-YM-20261028/blob/3456e5e75cf8d4cd558955a5c7638fba2d68f464/YM/Gauge/Chart.lean\#L138}{\texttt{plog}},
\href{https://github.com/yang-mills-group/OpenAI-YM-20261028/blob/3456e5e75cf8d4cd558955a5c7638fba2d68f464/YM/Gauge/Nearest.lean\#L26}{\texttt{Nearest}},
\href{https://github.com/yang-mills-group/OpenAI-YM-20261028/blob/3456e5e75cf8d4cd558955a5c7638fba2d68f464/YM/Gauge/Nearest.lean\#L91}{\texttt{focus}}
& \(\log_G\), nearest point, \(\mathrm{foc}(X)\) \\
curvature of a group element &
\href{https://github.com/yang-mills-group/OpenAI-YM-20261028/blob/3456e5e75cf8d4cd558955a5c7638fba2d68f464/YM/Gauge/Colors.lean\#L173}{\texttt{curvature}}
& \(F(u)\) \\
Wilson energy, weight, partition function &
\href{https://github.com/yang-mills-group/OpenAI-YM-20261028/blob/3456e5e75cf8d4cd558955a5c7638fba2d68f464/YM/Wilson/Law.lean\#L29}{\texttt{energy}},
\href{https://github.com/yang-mills-group/OpenAI-YM-20261028/blob/3456e5e75cf8d4cd558955a5c7638fba2d68f464/YM/Wilson/Law.lean\#L58}{\texttt{bare}},
\href{https://github.com/yang-mills-group/OpenAI-YM-20261028/blob/3456e5e75cf8d4cd558955a5c7638fba2d68f464/YM/Wilson/Law.lean\#L74}{\texttt{partition}}
& \(E_\ell(U)\), \(e^{-\beta E_\ell}\), \(Z_\ell(\beta)\) \\
product Haar, Wilson law, expectation &
\href{https://github.com/yang-mills-group/OpenAI-YM-20261028/blob/3456e5e75cf8d4cd558955a5c7638fba2d68f464/YM/Wilson/Law.lean\#L55}{\texttt{prior}},
\href{https://github.com/yang-mills-group/OpenAI-YM-20261028/blob/3456e5e75cf8d4cd558955a5c7638fba2d68f464/YM/Wilson/Law.lean\#L87}{\texttt{law}},
\href{https://github.com/yang-mills-group/OpenAI-YM-20261028/blob/3456e5e75cf8d4cd558955a5c7638fba2d68f464/YM/Wilson/Law.lean\#L99}{\texttt{moment}}
& \(\mu_\ell\), \(\mathbb P_{\ell,\beta}\),
\(\langle\cdot\rangle_{\ell,\beta}\) \\
spacetime, time, space, sup-distance &
\href{https://github.com/yang-mills-group/OpenAI-YM-20261028/blob/3456e5e75cf8d4cd558955a5c7638fba2d68f464/YM/Space/Configurations.lean\#L24}{\texttt{E4}},
\href{https://github.com/yang-mills-group/OpenAI-YM-20261028/blob/3456e5e75cf8d4cd558955a5c7638fba2d68f464/YM/Space/Configurations.lean\#L27}{\texttt{time}},
\href{https://github.com/yang-mills-group/OpenAI-YM-20261028/blob/3456e5e75cf8d4cd558955a5c7638fba2d68f464/YM/Space/Configurations.lean\#L28}{\texttt{spatial}},
\href{https://github.com/yang-mills-group/OpenAI-YM-20261028/blob/3456e5e75cf8d4cd558955a5c7638fba2d68f464/YM/Space/Configurations.lean\#L51}{\texttt{cubedist}}
& \(\mathbb R^4\), \(x_t\), \(\vec x\), \(\lvert x-y\rvert_\infty\) \\
\(n\)-point configurations &
\href{https://github.com/yang-mills-group/OpenAI-YM-20261028/blob/3456e5e75cf8d4cd558955a5c7638fba2d68f464/YM/Space/Configurations.lean\#L56}{\texttt{Points}},
\href{https://github.com/yang-mills-group/OpenAI-YM-20261028/blob/3456e5e75cf8d4cd558955a5c7638fba2d68f464/YM/Space/Configurations.lean\#L63}{\texttt{conf}},
\href{https://github.com/yang-mills-group/OpenAI-YM-20261028/blob/3456e5e75cf8d4cd558955a5c7638fba2d68f464/YM/Space/Configurations.lean\#L66}{\texttt{generic}},
\href{https://github.com/yang-mills-group/OpenAI-YM-20261028/blob/3456e5e75cf8d4cd558955a5c7638fba2d68f464/YM/Space/Configurations.lean\#L70}{\texttt{spaced}}
& \((\mathbb R^4)^n\), \(\mathrm{Conf}_n\), generic, \(r\)-spaced \\
orthogonal group, rotation generators &
\href{https://github.com/yang-mills-group/OpenAI-YM-20261028/blob/3456e5e75cf8d4cd558955a5c7638fba2d68f464/YM/Space/Orthogonal.lean\#L20}{\texttt{Ortho}},
\href{https://github.com/yang-mills-group/OpenAI-YM-20261028/blob/3456e5e75cf8d4cd558955a5c7638fba2d68f464/YM/Space/Orthogonal.lean\#L132-L142}{\texttt{skewSpace}},
\href{https://github.com/yang-mills-group/OpenAI-YM-20261028/blob/3456e5e75cf8d4cd558955a5c7638fba2d68f464/YM/Space/Orthogonal.lean\#L152-L154}{\texttt{plane}}
& \(O(4)\), \(\mathfrak{so}(4)\), \(\Omega_{ij}\) \\
Minkowski square, cones &
\href{https://github.com/yang-mills-group/OpenAI-YM-20261028/blob/3456e5e75cf8d4cd558955a5c7638fba2d68f464/YM/Mink/Geometry.lean\#L77}{\texttt{m}},
\href{https://github.com/yang-mills-group/OpenAI-YM-20261028/blob/3456e5e75cf8d4cd558955a5c7638fba2d68f464/YM/Mink/Geometry.lean\#L78}{\texttt{cone}},
\href{https://github.com/yang-mills-group/OpenAI-YM-20261028/blob/3456e5e75cf8d4cd558955a5c7638fba2d68f464/YM/Mink/Geometry.lean\#L79}{\texttt{inside}}
& \(\mathsf m(x)\), \(\overline V_+\), \(V_+\) \\
Lorentz group, dual, boost, Poincaré group &
\href{https://github.com/yang-mills-group/OpenAI-YM-20261028/blob/3456e5e75cf8d4cd558955a5c7638fba2d68f464/YM/Mink/Lorentz.lean\#L46}{\texttt{Mink.Lorentz}},
\href{https://github.com/yang-mills-group/OpenAI-YM-20261028/blob/3456e5e75cf8d4cd558955a5c7638fba2d68f464/YM/Mink/Lorentz.lean\#L95-L103}{\texttt{dual}},
\href{https://github.com/yang-mills-group/OpenAI-YM-20261028/blob/3456e5e75cf8d4cd558955a5c7638fba2d68f464/YM/Mink/Boost.lean\#L92-L101}{\texttt{Ray.boost}},
\href{https://github.com/yang-mills-group/OpenAI-YM-20261028/blob/3456e5e75cf8d4cd558955a5c7638fba2d68f464/YM/Mink/Poincare.lean\#L32-L34}{\texttt{Poincare}}
& \(\mathcal L=SO^+(1,3)\), \(g^\vee\), \(B_e(r)\), \(\mathcal P\) \\
\end{longtable}
}

\subsubsection{7. What remains
unexplained}\label{what-remains-unexplained}

\begin{enumerate}
\def\labelenumi{\arabic{enumi}.}
\tightlist
\item
  \textbf{Additional helper lemmas.} The assignment of declarations to
  I.1 includes helpers outside the six principal folders, including in
  \texttt{Cells}, \texttt{Extra}, \texttt{Further} and \texttt{Next}.
  Their statements and proofs are not covered here; their elementary
  character should not be assumed.
\item
  \textbf{Coarsening along the selected schedule.} Direct consumers
  retain \(M\le2B\) explicitly, for example
  \href{https://github.com/yang-mills-group/OpenAI-YM-20261028/blob/3456e5e75cf8d4cd558955a5c7638fba2d68f464/YM/Next/SCHEDULE/Support.lean\#L39-L65}{\texttt{YM.Next.Scheduling.Buy.footage}},
  or include it in a conjunction, as in
  \href{https://github.com/yang-mills-group/OpenAI-YM-20261028/blob/3456e5e75cf8d4cd558955a5c7638fba2d68f464/YM/Extra/ITERATION_NO_LOSS/Geometry.lean\#L75-L76}{\texttt{YM.Extra.Iteration.Shapes.raw}}.
  The derivation of this condition for the particular consecutive cover
  scales used in III.5 and III.10 remains to be traced through their
  schedule hypotheses.
\item
  \textbf{Bibliographic precision.} The references in §2 give context.
  Exact theorem locations and agreement of their conventions with the
  statements here remain to be checked.
\end{enumerate}

\subsection{II.1. Field labels and their lattice
representatives}\label{ii.1.-field-labels-and-their-lattice-representatives}

\textbf{Part:} II, Data of the construction. \textbf{Lean:} work items
G01--G07:
\href{https://github.com/yang-mills-group/OpenAI-YM-20261028/tree/3456e5e75cf8d4cd558955a5c7638fba2d68f464/YM/Labels}{\texttt{YM/Labels/}},
\href{https://github.com/yang-mills-group/OpenAI-YM-20261028/tree/3456e5e75cf8d4cd558955a5c7638fba2d68f464/YM/Source}{\texttt{YM/Source/}},
\href{https://github.com/yang-mills-group/OpenAI-YM-20261028/tree/3456e5e75cf8d4cd558955a5c7638fba2d68f464/YM/Work}{\texttt{YM/Work/}}
(G01--G07),
\href{https://github.com/yang-mills-group/OpenAI-YM-20261028/tree/3456e5e75cf8d4cd558955a5c7638fba2d68f464/YM/Next/FOUNDATIONS}{\texttt{YM/Next/FOUNDATIONS/}},
\href{https://github.com/yang-mills-group/OpenAI-YM-20261028/tree/3456e5e75cf8d4cd558955a5c7638fba2d68f464/YM/Push/RELAY_MARK}{\texttt{YM/Push/RELAY\_MARK/}},
\href{https://github.com/yang-mills-group/OpenAI-YM-20261028/tree/3456e5e75cf8d4cd558955a5c7638fba2d68f464/YM/Extra/SOURCE_HISTORIES}{\texttt{YM/Extra/SOURCE\_HISTORIES/}},
\href{https://github.com/yang-mills-group/OpenAI-YM-20261028/tree/3456e5e75cf8d4cd558955a5c7638fba2d68f464/YM/Push/LOCAL_SPLITTING}{\texttt{YM/Push/LOCAL\_SPLITTING/}},
\href{https://github.com/yang-mills-group/OpenAI-YM-20261028/tree/3456e5e75cf8d4cd558955a5c7638fba2d68f464/YM/Push/RELAY_DIFFERENCE}{\texttt{YM/Push/RELAY\_DIFFERENCE/}},
\href{https://github.com/yang-mills-group/OpenAI-YM-20261028/tree/3456e5e75cf8d4cd558955a5c7638fba2d68f464/YM/Planned/ATOM_GUARDED_LAYOUT}{\texttt{YM/Planned/ATOM\_GUARDED\_LAYOUT/}},
\href{https://github.com/yang-mills-group/OpenAI-YM-20261028/tree/3456e5e75cf8d4cd558955a5c7638fba2d68f464/YM/Further/LOCAL_LOG_JETS}{\texttt{YM/Further/LOCAL\_LOG\_JETS/}},
\href{https://github.com/yang-mills-group/OpenAI-YM-20261028/tree/3456e5e75cf8d4cd558955a5c7638fba2d68f464/YM/LocalH}{\texttt{YM/LocalH/}}.
\textbf{Size:} 10,805 lines in the main-theorem closure:

{\def\LTcaptype{none} % do not increment counter
\begin{longtable}[]{@{}ll@{}}
\toprule\noalign{}
package & lines \\
\midrule\noalign{}
\endhead
\bottomrule\noalign{}
\endlastfoot
equation sections (RELAY\_MARK) & 2,336 \\
source histories (SOURCE\_HISTORIES) & 1,801 \\
local splitting (LOCAL\_SPLITTING) & 1,485 \\
primary sections and readout (FOUNDATIONS) & 1,196 \\
comparison estimates (RELAY\_DIFFERENCE) & 858 \\
record, sections, transfers (Source) & 594 \\
work items G05--G07 & 934 \\
guarded layout, Haar chart, log jets & 749 \\
other files (labels, G01--G04) & 852 \\
\end{longtable}
}

\textbf{Conventions.}

\begin{itemize}
\tightlist
\item
  \textbf{Gauge group.} \(G\subset U(n)\) is a compact matrix group,
  with \(n\ge1\); nothing in this unit uses simplicity. Use the
  \(\mathrm{Ad}\)-invariant inner product
  \(\langle X,Y\rangle=-\mathrm{Re}\,\mathrm{tr}(XY)/n\) on its Lie
  algebra \(\mathfrak g\), and the Frobenius norm on matrices. For a
  plaquette \(P\), \(F_P(U)\in\mathfrak g\) is the orthogonal projection
  of \(\tfrac12(U_P-U_P^{-1})\), and
  \(\vartheta(U_P)=\lVert U_P-1\rVert/\sqrt n\) is its chord.
  \(S_W(U)=\sum_P(1-\mathrm{Re}\,\mathrm{tr}\,U_P/n)\) is the Wilson
  energy, summed once over each site and axis pair \(\mu<\nu\), and
  \(\langle\cdot\rangle_\beta\) is the Wilson expectation on a torus
  \(\Lambda_\ell=(\mathbb Z/\ell)^4\), where \(\ell\ge1\).
\item
  \textbf{Cover and placement.} Many objects live on the cover
  \(\mathbb Z^4\). A function \(f\) of finitely many cover links is
  \emph{placed} at a site \(o\in\Lambda_\ell\) by lifting the torus
  periodically around \(o\); the result is written \(f(o;U)\).
\item
  \textbf{Small fields.} \(\mathrm{Hy}_t(X)\) is the set of
  configurations in which no plaquette touching \(X\) has chord \(>t\).
\item
  \textbf{Levels.} \(m=0\) is the coarsest level and \(m=N\) the Wilson
  lattice. \(L_m\) is the odd block factor of the step \(m+1\to m\), the
  spacing satisfies \(a_m=L_ma_{m+1}\), and \(s_m=\log(1/a_m)\) is the
  scale.
\item
  \textbf{Couplings and level constants (III.9).}
  \(\hat\beta_m=\bar\beta+b_0s_m\) with \(b_0=11C_A/(24\pi^2)\). The
  small-field threshold is
  \(\theta_m=\hat\beta_m^{-1/2}\exp((\log\log\hat\beta_m)^2)\) and the
  \emph{forced allowance} is
  \(w_m=\exp(-\exp((\log\log\hat\beta_m)^2/5))\). \(J(m)\) is the
  nondecreasing stage index, \(\bar\kappa_j\ge1\) the stage-\(j\)
  constant, \(\kappa_m=\bar\kappa_{J(m)}\) the level-\(m\) constant, and
  \(b_m\) its budget point.
  \(\rho_m=\rho_{N,p,m}=e^{\mathcal V_m}\,\mathrm{gas}_m\) is the
  polymer representation of the level-\(m\) density, with vacuum
  exponent \(\mathcal V_m\) (II.2). \(T^{(m)}_\beta\) is the blocking
  push-forward from level \(m+1\) to level \(m\) (III.1).
\item
  \textbf{Budget norms (II.2).} A budget point \(b\) has order
  \(k_b\ge6\), power \(P_b\ge2\), threshold \(t_b\in(0,1]\), a support
  exponent and a cover size. \(f\) is \emph{\(b\)-regular on \(X\)} if
  it is local on links with both endpoints in \(X\), measurable and
  gauge invariant, and \(A\mapsto f(e^AU)\) is \(C^{k_b}\) at \(A=0\)
  for every \(U\in\mathrm{Hy}_{t_b}(X)\). All ordered left derivatives
  through order \(k_b\) must exist there. Write
  \(\lVert f\rVert_{b;o,X}\le c\) when \(c\ge0\) bounds every such
  derivative, including the value, in every list of directions
  satisfying \(\lvert A_e\rvert\le t_b(1+\mathrm{dist}(e_-,o))^{P_b}\).
  Here \(e_-\) is the initial vertex of the positive link \(e\), and
  \(\mathrm{dist}\) is the periodic \(\ell^\infty\) distance. All bounds
  are pointwise, including at the small-field boundary.
\end{itemize}

\textbf{What the unit does.} It defines the field labels and their
lattice representatives, and specifies the renormalized observables
\(O_{N,p}\).

\begin{enumerate}
\def\labelenumi{\arabic{enumi}.}
\tightlist
\item
  \textbf{Sections.} For each label it constructs gauge-invariant
  polynomials in the link matrices on a fixed stencil that reproduce the
  label to Taylor order \(d\) in the spacing, with symmetry and
  small-field bounds (§4.1--4.3).
\item
  \textbf{Readout.} It proves the converse: every regular lattice
  function has a label, and the label's size is bounded by the
  function's budget norm uniformly in the lattice side, support and norm
  parameters (§4.4).
\item
  \textbf{Equation sections.} From the lattice Schwinger--Dyson identity
  it builds a second kind of section on which every label of the
  free-Maxwell equation ideal becomes orthogonal to remote spectators
  (§4.5).
\item
  \textbf{Observables.} It specifies the record carrying a label through
  the RG history, proves the exact birth formula at the Wilson end, and
  collects the numerical conditions and comparison estimates used
  downstream (§4.6--4.8).
\end{enumerate}

\subsubsection{1. What this unit proves}\label{what-this-unit-proves-1}

\textbf{Proposition II.1.1 (label spaces and frames).} For \(j\ge0\) let
\(\mathrm{Tay}_j\) be the real space of \(4\)-tuples \(A=(A_\mu)\) of
symmetric \((j+1)\)-linear maps \((\mathbb R^4)^{j+1}\to\mathfrak g\)
satisfying the radial-gauge condition
\(\sum_\mu z_\mu A_\mu(z,\dots,z)=0\). Its elements are the Taylor
coefficients of weight \(j+2\) of a connection in radial gauge. A
\emph{jet} is a sequence \(J=(J_j)_{j\ge0}\). Fix a real basis of each
\(\mathrm{Tay}_j\), and give its coordinates weight \(j+2\). Let \(V_d\)
be the space of complex polynomials in these coordinates of total weight
\(\le d\) that are invariant under \(J\mapsto\mathrm{Ad}(g)J\) for all
\(g\in G\). Then:

\begin{enumerate}
\def\labelenumi{(\alph{enumi})}
\item
  \(V_d\) is finite-dimensional, \(1\in V_d\subset V_{d+1}\), and an
  element of \(V_d\) is determined by its values on real jets.
\item
  \(V_d=\bigoplus_{w,f}V_d^{(w,f)}\), where \((w,f)\) = (weight, number
  of factors). \(V_d^{(0,0)}=\mathbb C1\), and \(V_d^{(w,0)}=0\) for
  \(w>0\). Complex conjugation \(p\mapsto\bar p\) and the orthogonal
  group \(O(4)\), acting by rotating the jet, preserve the grades.
\item
  There is a coherent system of \emph{frames}. For each \(d\) there is a
  basis \(b_0=1,b_1,\dots,b_{K_d}\) of \(V_d\) with these properties:
\end{enumerate}

\begin{itemize}
\tightlist
\item
  each \(b_i\) is real (\(\bar b_i=b_i\)) and of pure grade
  \((w_i,f_i)\), with \((w_0,f_0)=(0,0)\), \(w_i\le d\) and \(f_i\ge1\)
  for \(i\neq0\);
\item
  the basis of \(V_d\) is contained in that of \(V_{d+1}\).
\end{itemize}

In coordinates \(c\in\mathbb C^{K_d+1}\) we write
\((S_tc)_i=t^{f_i}c_i\), \(\lVert c\rVert_1=\sum_i\lvert c_i\rvert\),
\(c^\circ=c-c_0e_0\), and \(s\cdot c\) for the \(O(4)\) action, defined
by \(p_{s\cdot c}(J)=p_c(s^{-1}J)\). \emph{Status:} definitions;
(a)--(c) are theorems. \emph{Lean:}
\href{https://github.com/yang-mills-group/OpenAI-YM-20261028/blob/3456e5e75cf8d4cd558955a5c7638fba2d68f464/YM/Labels/Jet.lean\#L39}{\texttt{Labels.Taylor}},
\href{https://github.com/yang-mills-group/OpenAI-YM-20261028/blob/3456e5e75cf8d4cd558955a5c7638fba2d68f464/YM/Labels/Jet.lean\#L22-L37}{\texttt{radial}},
\href{https://github.com/yang-mills-group/OpenAI-YM-20261028/blob/3456e5e75cf8d4cd558955a5c7638fba2d68f464/YM/Labels/Polynomial.lean\#L60-L61}{\texttt{limited}},
\href{https://github.com/yang-mills-group/OpenAI-YM-20261028/blob/3456e5e75cf8d4cd558955a5c7638fba2d68f464/YM/Labels/Polynomial.lean\#L103-L107}{\texttt{invariantSpace}},
\href{https://github.com/yang-mills-group/OpenAI-YM-20261028/blob/3456e5e75cf8d4cd558955a5c7638fba2d68f464/YM/Labels/Polynomial.lean\#L112}{\texttt{Label}},
\href{https://github.com/yang-mills-group/OpenAI-YM-20261028/blob/3456e5e75cf8d4cd558955a5c7638fba2d68f464/YM/Labels/Polynomial.lean\#L63-L94}{\texttt{cap\_finite}},
\href{https://github.com/yang-mills-group/OpenAI-YM-20261028/blob/3456e5e75cf8d4cd558955a5c7638fba2d68f464/YM/Labels/Polynomial.lean\#L38-L49}{\texttt{eval\_ext}},
\href{https://github.com/yang-mills-group/OpenAI-YM-20261028/blob/3456e5e75cf8d4cd558955a5c7638fba2d68f464/YM/Labels/Grading.lean\#L17-L18}{\texttt{Grade}},
\href{https://github.com/yang-mills-group/OpenAI-YM-20261028/blob/3456e5e75cf8d4cd558955a5c7638fba2d68f464/YM/Labels/Grading.lean\#L87-L105}{\texttt{decompose}},
\href{https://github.com/yang-mills-group/OpenAI-YM-20261028/blob/3456e5e75cf8d4cd558955a5c7638fba2d68f464/YM/Labels/Basis.lean\#L50-L57}{\texttt{Frame}},
\href{https://github.com/yang-mills-group/OpenAI-YM-20261028/blob/3456e5e75cf8d4cd558955a5c7638fba2d68f464/YM/Labels/Basis.lean\#L102-L104}{\texttt{Frames}},
\href{https://github.com/yang-mills-group/OpenAI-YM-20261028/blob/3456e5e75cf8d4cd558955a5c7638fba2d68f464/YM/Labels/Basis.lean\#L107-L111}{\texttt{frames}},
\href{https://github.com/yang-mills-group/OpenAI-YM-20261028/blob/3456e5e75cf8d4cd558955a5c7638fba2d68f464/YM/Labels/Reference.lean\#L58-L61}{\texttt{Frame.scaled}},
\href{https://github.com/yang-mills-group/OpenAI-YM-20261028/blob/3456e5e75cf8d4cd558955a5c7638fba2d68f464/YM/Labels/Reference.lean\#L83-L84}{\texttt{Frame.act}},
\href{https://github.com/yang-mills-group/OpenAI-YM-20261028/blob/3456e5e75cf8d4cd558955a5c7638fba2d68f464/YM/Labels/Reference.lean\#L123-L126}{\texttt{Frame.core}},
\href{https://github.com/yang-mills-group/OpenAI-YM-20261028/blob/3456e5e75cf8d4cd558955a5c7638fba2d68f464/YM/Labels/Reference.lean\#L128}{\texttt{Frame.coeffSize}}.
The field spaces \(F_d\) of the main theorem are quotients of \(V_d\) by
the null labels. That quotient is formed in V.1/V.2, not here.

\textbf{Proposition II.1.2 (test holonomies and sound sections).}

\begin{enumerate}
\def\labelenumi{(\alph{enumi})}
\item
  Fix \(d\), a jet \(J\), \(h\in\mathbb R\) and a cover edge
  \((x,\mu)\), and let
  \(A^{(d)}_\mu(z)=\sum_{j+2\le d}J_{j,\mu}(z,\ldots,z)\), with an empty
  sum when \(d<2\). Let \(u\colon\mathbb R\to G\) solve
  \(u'(t)=u(t)\,hA^{(d)}_\mu(h(x+te_\mu))\) with \(u(0)=1\). The
  solution exists and is unique, and we put
  \(U^{(d)}_{J,h}(x,\mu)=u(1)\). Moreover
  \(U^{(d)}_{J,h}=U^{(d)}_{\mathrm{sc}_hJ,1}\), where
  \(\mathrm{sc}_hJ=(h^{j+2}J_j)_j\), and \(U^{(d)}_{J,0}\equiv1\). A
  \emph{holonomy datum} is such a choice of solutions; by uniqueness its
  links do not depend on the choice.
\item
  A \emph{section} on a frame of \(V_d\) is a pair \((S,\Phi)\): a
  finite stencil \(S\subset\mathbb Z^4\) and a \(\mathbb C\)-linear map
  \(c\mapsto\Phi_c\) into polynomials in the entries, and their
  conjugates, of the link matrices on \(S\).
\item
  The section is \emph{sound} if:
\end{enumerate}

\begin{itemize}
\tightlist
\item
  \(0\in S\) and \(S\) is invariant under signed permutations;
\item
  each \(\Phi_c\) is gauge invariant, \(\Phi_{e_0}=1\), and
  \(\Phi_{\bar c}=\overline{\Phi_c}\);
\item
  \(\Phi_c(s^*U)=\Phi_{s\cdot c}(U)\) for every signed permutation
  \(s\), where \(s^*\) is the pull-back under \(y\mapsto sy\);
\item
  (\emph{projection}) for every \(c\) and every jet \(J\), the functions
  \(h\mapsto\Phi_c(U^{(d)}_{J,h})\) and
  \(h\mapsto p_c(\mathrm{sc}_hJ)\), with \(p_c=\sum_ic_ib_i\), are
  \(C^d\) at \(0\) and have equal derivatives of order \(\le d\) there.
\end{itemize}

\emph{Status:} definitions; the ODE facts are theorems. \emph{Lean:}
\href{https://github.com/yang-mills-group/OpenAI-YM-20261028/blob/3456e5e75cf8d4cd558955a5c7638fba2d68f464/YM/Source/Local.lean\#L55-L57}{\texttt{Source.Hol}},
\href{https://github.com/yang-mills-group/OpenAI-YM-20261028/blob/3456e5e75cf8d4cd558955a5c7638fba2d68f464/YM/Source/Local.lean\#L34-L35}{\texttt{partConnection}},
\href{https://github.com/yang-mills-group/OpenAI-YM-20261028/blob/3456e5e75cf8d4cd558955a5c7638fba2d68f464/YM/Source/Local.lean\#L42-L52}{\texttt{SolvesLink.unique}},
\href{https://github.com/yang-mills-group/OpenAI-YM-20261028/blob/3456e5e75cf8d4cd558955a5c7638fba2d68f464/YM/Source/ODE.lean\#L162-L168}{\texttt{Hol.exists\_ode}},
\href{https://github.com/yang-mills-group/OpenAI-YM-20261028/blob/3456e5e75cf8d4cd558955a5c7638fba2d68f464/YM/Source/Local.lean\#L82}{\texttt{scaleJet}},
\href{https://github.com/yang-mills-group/OpenAI-YM-20261028/blob/3456e5e75cf8d4cd558955a5c7638fba2d68f464/YM/Source/Local.lean\#L85-L89}{\texttt{Projects}},
\href{https://github.com/yang-mills-group/OpenAI-YM-20261028/blob/3456e5e75cf8d4cd558955a5c7638fba2d68f464/YM/Source/Local.lean\#L103-L105}{\texttt{SectionData}},
\href{https://github.com/yang-mills-group/OpenAI-YM-20261028/blob/3456e5e75cf8d4cd558955a5c7638fba2d68f464/YM/Source/Local.lean\#L140-L147}{\texttt{SectionData.Sound}}.

\textbf{Proposition II.1.3 (primary sections).} For every gauge group
\(G\), every holonomy datum and every coherent system of frames, there
are sections \(\Sigma_d=(S_d,\Phi^{(d)})\) on the frames of \(V_d\) such
that, for every \(d\):

\begin{enumerate}
\def\labelenumi{\arabic{enumi}.}
\tightlist
\item
  \(\Sigma_d\) is sound.
\item
  (\emph{slot structure}) Each column is
  \(\Phi^{(d)}_{e_i}(U)=P_i(\sigma_1(U),\dots,\sigma_m(U))\). Here the
  \(\sigma_l\) are slots with stencils in \(S_d\), \(P_i\) is a complex
  polynomial in their real components, and every monomial of \(P_i\) has
  degree \(\ge f_i\). A \emph{slot} \(\sigma=(\omega,P)\) is a lattice
  walk \(\omega\) from \(0\) followed by a plaquette \(P\) at its end,
  with \(\sigma(U)=\mathrm{Ad}(U_\omega)F_P(U)\in\mathfrak g\).
\item
  (\emph{cost}) For all \(k,P\in\mathbb N\) there is \(A\ge0\),
  depending on the chosen section and on \(k,P\), with the following
  property. For every budget point \(b\) of order \(k\) and power \(P\)
  (any threshold and support exponent), every side \(\ell\), site \(o\)
  and index \(i\), the placed column \(\Phi^{(d)}_{e_i}(o;\cdot)\) is
  \(b\)-regular on its footprint \(Y_o=o+S_d\) with sup bound
  \(A\,t_b^{f_i}\).
\end{enumerate}

For the fixed primary sections used in the final family, two further
facts hold:

\begin{enumerate}
\def\labelenumi{\arabic{enumi}.}
\setcounter{enumi}{3}
\tightlist
\item
  one nondecreasing bound \(\pi(D)\ge1\) serves in item 3 for all
  \(d,k,P\le D\);
\item
  there is \(c_d\ge1\) with \(\lvert\Phi^{(d)}_{e_i}\rvert\le c_d\) at
  every configuration and every placement.
\end{enumerate}

\emph{Status:} theorem, assembled from G01--G07; no simplicity
assumption on \(G\). \emph{Lean:}
\href{https://github.com/yang-mills-group/OpenAI-YM-20261028/blob/3456e5e75cf8d4cd558955a5c7638fba2d68f464/YM/Integrated.lean\#L216-L219}{\texttt{Integrated.source\_sections}},
\href{https://github.com/yang-mills-group/OpenAI-YM-20261028/blob/3456e5e75cf8d4cd558955a5c7638fba2d68f464/YM/Bid/Source.lean\#L76-L86}{\texttt{Bid.teaching}},
\href{https://github.com/yang-mills-group/OpenAI-YM-20261028/blob/3456e5e75cf8d4cd558955a5c7638fba2d68f464/YM/Bid/Source.lean\#L72-L74}{\texttt{Bid.Learns}},
\href{https://github.com/yang-mills-group/OpenAI-YM-20261028/blob/3456e5e75cf8d4cd558955a5c7638fba2d68f464/YM/Bid/Source.lean\#L55-L56}{\texttt{Bid.pop}},
\href{https://github.com/yang-mills-group/OpenAI-YM-20261028/blob/3456e5e75cf8d4cd558955a5c7638fba2d68f464/YM/Bid/Source.lean\#L63-L66}{\texttt{Bid.Credit}},
\href{https://github.com/yang-mills-group/OpenAI-YM-20261028/blob/3456e5e75cf8d4cd558955a5c7638fba2d68f464/YM/Next/FOUNDATIONS/Primary.lean\#L92}{\texttt{Foundations.hol}},
\href{https://github.com/yang-mills-group/OpenAI-YM-20261028/blob/3456e5e75cf8d4cd558955a5c7638fba2d68f464/YM/Next/FOUNDATIONS/Primary.lean\#L101}{\texttt{localSection}},
\href{https://github.com/yang-mills-group/OpenAI-YM-20261028/blob/3456e5e75cf8d4cd558955a5c7638fba2d68f464/YM/Next/FOUNDATIONS/Primary.lean\#L103-L104}{\texttt{sound}},
\href{https://github.com/yang-mills-group/OpenAI-YM-20261028/blob/3456e5e75cf8d4cd558955a5c7638fba2d68f464/YM/Next/FOUNDATIONS/Primary.lean\#L106-L107}{\texttt{valuation}},
\href{https://github.com/yang-mills-group/OpenAI-YM-20261028/blob/3456e5e75cf8d4cd558955a5c7638fba2d68f464/YM/Next/FOUNDATIONS/Primary.lean\#L126}{\texttt{price}},
\href{https://github.com/yang-mills-group/OpenAI-YM-20261028/blob/3456e5e75cf8d4cd558955a5c7638fba2d68f464/YM/Next/FOUNDATIONS/Primary.lean\#L133-L142}{\texttt{priced}},
\href{https://github.com/yang-mills-group/OpenAI-YM-20261028/blob/3456e5e75cf8d4cd558955a5c7638fba2d68f464/YM/Next/FOUNDATIONS/Primary.lean\#L147-L163}{\texttt{full\_value\_bound}}.
These are the sections of the final Wilson family (Prop. II.1.6).

\textbf{Proposition II.1.4 (readout of a lattice function).} Fix a
holonomy datum, \(d\), and a frame of \(V_d\).

\begin{enumerate}
\def\labelenumi{(\alph{enumi})}
\tightlist
\item
  Let \(b\) be a budget point of order \(\ge d\), let
  \(X\subset\Lambda_\ell\), \(o\in\Lambda_\ell\), and let \(f\) be
  \(b\)-regular on \(X\). There is a unique \(p\in V_d\) that
  \emph{reads} \(f\) at \(o\):
\end{enumerate}

\begin{itemize}
\tightlist
\item
  if \(X\) does not wrap around \(o\) (\(10\,\mathrm{dist}(x,o)<\ell\)
  for all \(x\in X\)), then for every jet \(J\) the function
  \(h\mapsto f(U^{(d),o}_{J,h})\) is \(C^d\) at \(0\) and has the same
  derivatives of order \(\le d\) there as
  \(h\mapsto p(\mathrm{sc}_hJ)\);
\item
  if \(X\) wraps, \(p=f(1)\cdot1\) by convention.
\end{itemize}

Here \(U^{(d),o}_{J,h}(x,\mu)=U^{(d)}_{J,h}(\widetilde{x-o},\mu)\),
where the tilde denotes the centred periodic lift. These are torus test
links, defined by lifting each initial vertex; they need not be periodic
samples of a connection across the seam.

\begin{enumerate}
\def\labelenumi{(\alph{enumi})}
\setcounter{enumi}{1}
\item
  There is \(C\ge1\), depending only on the holonomy datum and the
  frame, with the following property. For all \(\ell,o,X\), all budget
  points \(b\) with \(k_b\ge d\) and \(P_b\ge d\), all \(f\) that are
  \(b\)-regular on \(X\) with sup bound \(c\), and all \(p\) reading
  \(f\): \[\lVert S_{t_b}\,\mathrm{coord}(p)\rVert_1\le C\,c .\]
\item
  If \(\Sigma\) is sound and \(20(\lvert x_\mu\rvert+1)\le\ell\) for
  every \(x\in S\) and every axis \(\mu\), then the placed function
  \(\Phi_c(o;\cdot)\) reads \(p_c\) at \(o\).
\end{enumerate}

\emph{Status:} theorem. \emph{Lean:}
\href{https://github.com/yang-mills-group/OpenAI-YM-20261028/blob/3456e5e75cf8d4cd558955a5c7638fba2d68f464/YM/Next/FOUNDATIONS/Readout.lean\#L238-L247}{\texttt{exists\_read}},
\href{https://github.com/yang-mills-group/OpenAI-YM-20261028/blob/3456e5e75cf8d4cd558955a5c7638fba2d68f464/YM/Next/FOUNDATIONS/Readout.lean\#L218-L233}{\texttt{read\_projection}},
\href{https://github.com/yang-mills-group/OpenAI-YM-20261028/blob/3456e5e75cf8d4cd558955a5c7638fba2d68f464/YM/Next/FOUNDATIONS/JetReads.lean\#L186-L194}{\texttt{Reads}},
\href{https://github.com/yang-mills-group/OpenAI-YM-20261028/blob/3456e5e75cf8d4cd558955a5c7638fba2d68f464/YM/Next/FOUNDATIONS/TestBounds.lean\#L271-L290}{\texttt{price\_of\_readout}},
\href{https://github.com/yang-mills-group/OpenAI-YM-20261028/blob/3456e5e75cf8d4cd558955a5c7638fba2d68f464/YM/Next/FOUNDATIONS/JetReads.lean\#L253-L267}{\texttt{placed\_reads}}.
This bounds the local-splitting projection and the remainder comparisons
in III.12--III.13 and VIII.1.

\textbf{Proposition II.1.5 (equation sections).} Let
\(\mathrm{Sh}_j\subset\mathrm{Tay}_j\) be the radial potentials whose
curvature is a degree-\(j\) homogeneous Taylor coefficient of a
\(\mathfrak g\)-valued solution of the \emph{free} Maxwell equations.
That is: \(F\) is antisymmetric, satisfies the Bianchi identity, and
\(\sum_a\partial_aF_{ab}=0\). The \emph{equation ideal} is
\[\mathcal E_d=\{p\in V_d:\ p(J)=0\text{ for every jet with }J_j\in\mathrm{Sh}_j\text{ for all }j\}.\]
This is the weight-\(\le d\) part of the ideal of the linearized field
equations, imposed degree by degree. For every holonomy datum, every
\(d\) and every frame of \(V_d\), there are a section \(\Sigma\) and a
\(\mathbb C\)-linear automorphism \(R\) of \(V_d\) (the \emph{weave})
such that:

\begin{itemize}
\tightlist
\item
  \(\Sigma\) is sound and has the slot structure and cost of Prop.
  II.1.3;
\item
  \(R\) commutes with the signed-permutation action and with
  conjugation;
\item
  for every \(p\in\mathcal E_d\), the cover function
  \(\Phi_{\mathrm{coord}(Rp)}\) is \emph{micro}.
\end{itemize}

A function \(f\) is \emph{micro} if it is continuous and there are
\(r,M\in\mathbb N\) with the following property. For every side
\(\ell\ge M\), site \(o\), coupling \(\beta\neq0\) and bounded
measurable \(S\) depending only on links with both endpoints at distance
\(>r\) from \(o\), \[\big\langle f(o;\cdot)\,S\big\rangle_\beta=0 .\]
The same holds simultaneously for all \(d\) along a fixed system of
frames. The constructed \(R\) is unitriangular: if \(p\) has pure grade
\((w,f)\), every nonzero component of \(Rp-p\) has grade \((w',f')\)
with \((f',w')>(f,w)\) in lexicographic order. The existence statement
does not require this extra property. \emph{Status:} theorem,
unconditional. \emph{Lean:}
\href{https://github.com/yang-mills-group/OpenAI-YM-20261028/blob/3456e5e75cf8d4cd558955a5c7638fba2d68f464/YM/Maxwell/Equation.lean\#L72}{\texttt{Maxwell.equationIdeal}},
\href{https://github.com/yang-mills-group/OpenAI-YM-20261028/blob/3456e5e75cf8d4cd558955a5c7638fba2d68f464/YM/Maxwell/Equation.lean\#L74-L79}{\texttt{ideal\_meaning}},
\href{https://github.com/yang-mills-group/OpenAI-YM-20261028/blob/3456e5e75cf8d4cd558955a5c7638fba2d68f464/YM/Maxwell/Shell.lean\#L18-L42}{\texttt{Maxwell.onShell}},
\href{https://github.com/yang-mills-group/OpenAI-YM-20261028/blob/3456e5e75cf8d4cd558955a5c7638fba2d68f464/YM/Maxwell/Shell.lean\#L46-L47}{\texttt{gaugeJet}},
\href{https://github.com/yang-mills-group/OpenAI-YM-20261028/blob/3456e5e75cf8d4cd558955a5c7638fba2d68f464/YM/Push/RELAY_MARK/Micro.lean\#L33}{\texttt{Micro}},
\href{https://github.com/yang-mills-group/OpenAI-YM-20261028/blob/3456e5e75cf8d4cd558955a5c7638fba2d68f464/YM/Push/RELAY_MARK/Micro.lean\#L29-L31}{\texttt{MicroAt}},
\href{https://github.com/yang-mills-group/OpenAI-YM-20261028/blob/3456e5e75cf8d4cd558955a5c7638fba2d68f464/YM/Push/RELAY_MARK/Portfolio.lean\#L51-L58}{\texttt{Portfolio}},
\href{https://github.com/yang-mills-group/OpenAI-YM-20261028/blob/3456e5e75cf8d4cd558955a5c7638fba2d68f464/YM/Push/RELAY_MARK/Portfolio.lean\#L101-L105}{\texttt{Portfolio.build}},
\href{https://github.com/yang-mills-group/OpenAI-YM-20261028/blob/3456e5e75cf8d4cd558955a5c7638fba2d68f464/YM/Push/RELAY_MARK/Assembly.lean\#L30-L33}{\texttt{commissioned}},
\href{https://github.com/yang-mills-group/OpenAI-YM-20261028/blob/3456e5e75cf8d4cd558955a5c7638fba2d68f464/YM/Push/RELAY_MARK/Automorphism.lean\#L43-L50}{\texttt{regrade\_tri}}.
In the main chain these sections enter through the auxiliary-section
data in
\href{https://github.com/yang-mills-group/OpenAI-YM-20261028/blob/3456e5e75cf8d4cd558955a5c7638fba2d68f464/YM/FullVolume/Cross.lean\#L65-L71}{\texttt{Run.Ties}}
(III.12), as a second family of sections on the same history, together
with
\href{https://github.com/yang-mills-group/OpenAI-YM-20261028/blob/3456e5e75cf8d4cd558955a5c7638fba2d68f464/YM/FullVolume/Fire.lean\#L17-L20}{\texttt{Hearthside}}.
\href{https://github.com/yang-mills-group/OpenAI-YM-20261028/blob/3456e5e75cf8d4cd558955a5c7638fba2d68f464/YM/Push/RELAY_MARK/Assembly.lean\#L39-L41}{\texttt{Memento}}
and
\href{https://github.com/yang-mills-group/OpenAI-YM-20261028/blob/3456e5e75cf8d4cd558955a5c7638fba2d68f464/YM/Push/RELAY_MARK/Assembly.lean\#L57-L67}{\texttt{awarded}}
would make them the family's own sections, but they are outside the
main-theorem closure.

\textbf{Proposition II.1.6 (the construction record and the renormalized
observables).} A \emph{record} \(\mathsf C\) consists of:

\begin{itemize}
\tightlist
\item
  an RG history \(H\) (III.9) and a holonomy datum;
\item
  coherent frames, and sections \(\Sigma_d\) on them;
\item
  reference levels \(r(d)\);
\item
  for all \(d,N,p\), a transfer family
  \(\mathcal T_{d,N,p}=((E_m)_m,\tau)\) (Prop. II.1.7), with
  \(\tau=\tau_{d,N,p}\) a linear automorphism of \(\mathbb C^{K_d+1}\).
\end{itemize}

The renormalized observable of \(v\in\mathbb C^{K_d+1}\) at
\(o\in\Lambda_{N,p}\) is
\[O_{N,p}(v,o;U)=a_N^{-d}\,\Phi^{(d)}_{S_{1/\theta_N}\tau^{-1}v}(o;U).\]
The \emph{consistency predicate} \(\mathrm{Makes}(\mathsf C)\) states:

\begin{itemize}
\tightlist
\item
  \begin{enumerate}
  \def\labelenumi{(\roman{enumi})}
  \tightlist
  \item
    \(H\) is selected;
  \end{enumerate}
\item
  (ii) every \(\Sigma_d\) is sound;
\item
  (iii) \(d+3\le J(r(d))\), and for \(m\ge r(d)\) the level-\(m\) budget
  has order \(\ge d+1\) and power \(\ge d\);
\item
  (iv) if \(r(d)\le N\), then \(20(\lvert x_\mu\rvert+1)\le\ell_{N,p}\)
  for every \(x\in S_d\) and every axis \(\mu\);
\item
  (v) if \(r(d)\le N\), each placed column \(\Phi^{(d)}_{e_i}(o;\cdot)\)
  is \(b_N\)-regular on \(Y_o\) with sup bound
  \(\kappa_N\theta_N^{f_i}\);
\item
  (vi) if \(r(d)\le N\), then \(\mathcal T_{d,N,p}\) is based for
  \((\Sigma_d,r(d))\).
\end{itemize}

Then:

\begin{enumerate}
\def\labelenumi{(\alph{enumi})}
\item
  Without any hypothesis: \(v\mapsto O_{N,p}(v,o)\) is linear,
  \(O_{N,p}(v,o)\) is continuous and depends only on links in \(Y_o\),
  and \(O_{N,p}(v,o+b;U)=O_{N,p}(v,o;U(\cdot+b))\) (literal translation
  covariance).
\item
  If the sections are sound, \(O_{N,p}(v,o)\) is gauge invariant for
  every \(N\).
\item
  Under \(\mathrm{Makes}\) and \(r(d)\le N\):
\end{enumerate}

\begin{itemize}
\tightlist
\item
  reality: \(O_{N,p}(\bar v,o)=\overline{O_{N,p}(v,o)}\);
\item
  hypercubic covariance:
  \(O_{N,p}(s\cdot v,A_{s,b}o;U)=O_{N,p}(v,o;A_{s,b}^*U)\) for the
  lattice affine map \(A_{s,b}y=sy+b\);
\item
  normalization of the unit: \(O_{N,p}(a_{r(d)}^de_0,o)\equiv1\);
\item
  normalized readout: the level-\(r(d)\) coordinates of the transfer of
  \(\tau^{-1}v\) equal \(v\).
\end{itemize}

The renormalization map of the main theorem is
\(B_{N,p,d}=S_{1/\theta_N}\tau^{-1}\): writing \(P_d[c]=\Phi^{(d)}_c\),
its insertion formula is \(I_{N,p}=a_N^{-d}P_d[B_{N,p,d}c]\).
\emph{Status:} (a)--(c) are theorems. \(\mathrm{Makes}\) is proved for
the final family in
\href{https://github.com/yang-mills-group/OpenAI-YM-20261028/blob/3456e5e75cf8d4cd558955a5c7638fba2d68f464/YM/MergeProof/Marco.lean\#L30-L49}{\texttt{Bandage.squad}}
(VIII.1) and included in the construction record
\href{https://github.com/yang-mills-group/OpenAI-YM-20261028/blob/3456e5e75cf8d4cd558955a5c7638fba2d68f464/YM/FullVolume/Cross.lean\#L30-L48}{\texttt{Run}}
(III.12). That family uses the fixed coherent frames, holonomies and
primary sections of Prop. II.1.3. Clauses (ii)--(v) come from Prop.
II.1.9(c) of this unit and clause (vi) from G08 (III.10). \emph{Lean:}
\href{https://github.com/yang-mills-group/OpenAI-YM-20261028/blob/3456e5e75cf8d4cd558955a5c7638fba2d68f464/YM/Source/Construction.lean\#L24-L30}{\texttt{Family}},
\href{https://github.com/yang-mills-group/OpenAI-YM-20261028/blob/3456e5e75cf8d4cd558955a5c7638fba2d68f464/YM/Source/Construction.lean\#L42}{\texttt{Family.O}},
\href{https://github.com/yang-mills-group/OpenAI-YM-20261028/blob/3456e5e75cf8d4cd558955a5c7638fba2d68f464/YM/Source/Insertion.lean\#L77-L79}{\texttt{initial}},
\href{https://github.com/yang-mills-group/OpenAI-YM-20261028/blob/3456e5e75cf8d4cd558955a5c7638fba2d68f464/YM/Source/Construction.lean\#L44-L62}{\texttt{Family.Makes}},
\href{https://github.com/yang-mills-group/OpenAI-YM-20261028/blob/3456e5e75cf8d4cd558955a5c7638fba2d68f464/YM/Source/Construction.lean\#L117-L127}{\texttt{linearO}},
\href{https://github.com/yang-mills-group/OpenAI-YM-20261028/blob/3456e5e75cf8d4cd558955a5c7638fba2d68f464/YM/Source/Construction.lean\#L114-L115}{\texttt{continuousO}},
\href{https://github.com/yang-mills-group/OpenAI-YM-20261028/blob/3456e5e75cf8d4cd558955a5c7638fba2d68f464/YM/Source/Construction.lean\#L109-L112}{\texttt{localO}},
\href{https://github.com/yang-mills-group/OpenAI-YM-20261028/blob/3456e5e75cf8d4cd558955a5c7638fba2d68f464/YM/Source/Construction.lean\#L71-L72}{\texttt{translation}},
\href{https://github.com/yang-mills-group/OpenAI-YM-20261028/blob/3456e5e75cf8d4cd558955a5c7638fba2d68f464/YM/Source/Construction.lean\#L104-L107}{\texttt{invariant}},
\href{https://github.com/yang-mills-group/OpenAI-YM-20261028/blob/3456e5e75cf8d4cd558955a5c7638fba2d68f464/YM/Source/Construction.lean\#L74-L86}{\texttt{reality}},
\href{https://github.com/yang-mills-group/OpenAI-YM-20261028/blob/3456e5e75cf8d4cd558955a5c7638fba2d68f464/YM/Source/Construction.lean\#L88-L98}{\texttt{rotated}},
\href{https://github.com/yang-mills-group/OpenAI-YM-20261028/blob/3456e5e75cf8d4cd558955a5c7638fba2d68f464/YM/Source/Construction.lean\#L100-L102}{\texttt{constant}},
\href{https://github.com/yang-mills-group/OpenAI-YM-20261028/blob/3456e5e75cf8d4cd558955a5c7638fba2d68f464/YM/Source/Construction.lean\#L65-L68}{\texttt{normalized}}.

\textbf{Proposition II.1.7 (based transfer families).} \emph{Source
coefficients.} A source coefficient at level \(m\) is
\(E=(x,(X_i),(g_i),(h_i))_{i\in\mathbb N}\). Here
\(x\in\mathbb C^{K_d+1}\) are the running label coordinates, the \(X_i\)
are supports, the \(g_i\) are \emph{regular} remainders, and the \(h_i\)
are \emph{forced} functions attached to the forced-polymer addresses
\(i\) of \(\rho_m\), whose activities are \(K_i\).

\emph{Density.} Its density at \(o\) is the \(\lambda\)-coefficient, at
\(\lambda^2=0\), of the density with one inserted source:
\[\mathcal D_m(E)(U)=\rho_m(U)\Big[x_0+\mathbf 1_{\mathrm{Hy}_{\theta_m}(Y_o)}\Phi_{S_{1/\theta_m}x^\circ}(o;U)+\sum_i\mathbf 1_{\mathrm{Hy}_{\theta_m}(X_i)}g_i(U)\Big]+e^{\mathcal V_m(U)}\sum_{b\ \mathrm{adm}}\ \sum_{i\in b}h_i(U)\prod_{j\in b\setminus\{i\}}K_j(U).\]

\emph{Format at \(o\).} The \(g_i\) are \(b_m\)-regular on \(X_i\ni o\),
with \(g_i(1)=0\). When \(X_i\) does not wrap around \(o\),
\(h\mapsto g_i(U^{(d),o}_{J,h})\) is \(C^d\) at \(0\) and all its
derivatives of order \(\le d\) vanish there, for every \(J\). If
\(k_{b_m}\ge d\), this is exactly zero readout in Prop. II.1.4(a). Each
\(h_i\) is gauge invariant and measurable, vanishes unless polymer \(i\)
is eligible and carries \(o\), and depends only on the links and flags
of the support of polymer \(i\).

\emph{Remainder bound \((\kappa,c)\).} Require \(\kappa>0\) and
nonnegative \(a,b\) with \(a+b\le c\), together with nonnegative weights
\(e_i,e_i'\) with \(\sum_i\mathrm{mass}(X_i,e_i)\le a\),
\(\lVert g_i\rVert_{b_m;o,X_i}\le e_i\),
\(\sum_i\mathrm{mass}(\mathrm{supp}K_i,e'_i)\le\kappa b\) and
\(\lvert h_i\rvert\le e'_i\) at every configuration. Each weighted sum
is required to converge; \(\mathrm{supp}K_i\) denotes the prescribed
support of forced polymer \(i\).

\emph{Based family.} A family \((E_m(v))_m\) with terminal automorphism
\(\tau\) is \emph{based} for \((\Sigma,r)\) if \(r\le N\) and, for all
\(v\) and \(r\le m\le N\):

\begin{itemize}
\tightlist
\item
  \begin{enumerate}
  \def\labelenumi{(\arabic{enumi})}
  \tightlist
  \item
    \(E_m(v)\) has Format at \(0\) and remainder bound
    \((\sqrt{w_m},\ \kappa_m\lVert x_m(v)^\circ\rVert_1)\);
  \end{enumerate}
\item
  (2) \(x_m\), \(g\) and \(h\) are linear in \(v\), and the supports do
  not depend on \(v\);
\item
  (3) (\emph{bare end}) \(x_N(v)=v\), \(g\equiv0\), and
  \(\mathcal D_N(E_N(v))=\rho_N\,\Phi_{S_{1/\theta_N}v}(0;\cdot)\);
\item
  (4) (\emph{transfer})
  \(\mathcal D_m(E_m(v))=L_m^d\,T^{(m)}_{\beta_{N,p,m+1}}\mathcal D_{m+1}(E_{m+1}(v))\)
  for \(m<N\);
\item
  (5) \(\lVert x_{m+1}\rVert_1\le\kappa_{m+1}\lVert x_m\rVert_1\) and
  \(\lVert x_m\rVert_1\le\kappa_{m+1}L_m^d\lVert x_{m+1}\rVert_1\);
\item
  (6) \(x_r(v)=\tau v\), \(\tau\) commutes with conjugation and with the
  signed-permutation action, and \(x_m(a_N^de_0)=a_m^de_0\).
\end{itemize}

\emph{Status:} definition (specification). G08 (III.10) produces such
families from the per-step source estimates, together with the room and
bank inequalities of Prop. II.1.9. A coefficient with Format and a
remainder bound has bounded measurable density whenever the underlying
representation is lawful (Lemma 4.26). For a based family on a selected
history, and \(r\le m\le N\),
\[\int\mathcal D_m(E_m(\tau^{-1}v))\,dU=a_m^d\,Z_{N,p}\,\langle O_{N,p}(v,0)\rangle_{\beta_{N,p}} .\]
\emph{Lean:}
\href{https://github.com/yang-mills-group/OpenAI-YM-20261028/blob/3456e5e75cf8d4cd558955a5c7638fba2d68f464/YM/Source/Insertion.lean\#L26-L30}{\texttt{Coefficient}},
\href{https://github.com/yang-mills-group/OpenAI-YM-20261028/blob/3456e5e75cf8d4cd558955a5c7638fba2d68f464/YM/Source/Insertion.lean\#L48-L50}{\texttt{Coefficient.density}},
\href{https://github.com/yang-mills-group/OpenAI-YM-20261028/blob/3456e5e75cf8d4cd558955a5c7638fba2d68f464/YM/Source/Insertion.lean\#L37-L39}{\texttt{exponentialEntry}},
\href{https://github.com/yang-mills-group/OpenAI-YM-20261028/blob/3456e5e75cf8d4cd558955a5c7638fba2d68f464/YM/Source/Insertion.lean\#L52-L61}{\texttt{Format}},
\href{https://github.com/yang-mills-group/OpenAI-YM-20261028/blob/3456e5e75cf8d4cd558955a5c7638fba2d68f464/YM/Source/Insertion.lean\#L63-L68}{\texttt{remBound}},
\href{https://github.com/yang-mills-group/OpenAI-YM-20261028/blob/3456e5e75cf8d4cd558955a5c7638fba2d68f464/YM/Source/Insertion.lean\#L73-L75}{\texttt{TransferData}},
\href{https://github.com/yang-mills-group/OpenAI-YM-20261028/blob/3456e5e75cf8d4cd558955a5c7638fba2d68f464/YM/Source/Insertion.lean\#L89-L114}{\texttt{TransferData.Based}},
\href{https://github.com/yang-mills-group/OpenAI-YM-20261028/blob/3456e5e75cf8d4cd558955a5c7638fba2d68f464/YM/Source/Preservation.lean\#L100-L118}{\texttt{integral\_marked}}
(not used by the main theorem),
\href{https://github.com/yang-mills-group/OpenAI-YM-20261028/blob/3456e5e75cf8d4cd558955a5c7638fba2d68f464/YM/Bid/Loan.lean\#L101-L103}{\texttt{Bid.G08}}.

\textbf{Proposition II.1.8 (birth at the Wilson end).} Let \(\rho\) be a
lawful polymer representation on \(\Lambda_\ell\) with parameter \(w>0\)
(II.2) and threshold \(\theta\). Let \(o\in\Lambda_\ell\), and let
\(\Sigma\) be a sound section whose placed columns satisfy
\(\lvert\Phi_{e_i}(o;\cdot)\rvert\le B\) with \(B\ge1\). Put \(Y=Y_o\)
and \(f_{\max}=\max_if_i\), and assume the \emph{fee condition}
\[\mathrm{fee}:=B\,\theta^{-f_{\max}}\,\mathrm{mass}(Y,1)\,\big(e^{\lvert Y\rvert w}-1\big)\le\sqrt w .\]
Assume a \emph{shelf} of reserved forced addresses: for every pair of
finite supports \(\varnothing\ne I\subseteq X\) and every auxiliary
index in \(\mathbb N\), there is a distinct forced address with support
\(X\), flagged inventory \(I\), identically zero activity, and an
auxiliary eligibility test satisfied at every configuration. For every
fixed reservation index there is a coefficient \(E(v)\) with the
following properties:

\begin{itemize}
\tightlist
\item
  \(x=v\) and \(g\equiv0\);
\item
  the forced part is the large-field \emph{supplement}
  \(\Phi_{S_{1/\theta}v}(o;\cdot)-v_0-\mathbf 1_{\mathrm{Hy}_\theta(Y)}\Phi_{S_{1/\theta}v^\circ}(o;\cdot)\),
  spliced into the reserved addresses;
\item
  the components \(x,g,h\) depend linearly on \(v\), the supports are
  independent of \(v\), and \(E(v)\) has Format at \(o\) and remainder
  bound \((\sqrt w,\lVert v^\circ\rVert_1)\);
\item
  for every configuration \(U\),
  \[\mathcal D(E(v))(U)=\rho(U)\,\Phi_{S_{1/\theta}v}(o;U).\]
\end{itemize}

\emph{Status:} theorem. The main chain uses
\href{https://github.com/yang-mills-group/OpenAI-YM-20261028/blob/3456e5e75cf8d4cd558955a5c7638fba2d68f464/YM/Planned/ATOM_GUARDED_LAYOUT/Guards.lean\#L287-L326}{\texttt{Guard.summon}}
at level \(N\), with \(w=w_N\). It requires a selected history and
reserved addresses whose auxiliary tests need hold only when the support
has exactly the prescribed flagged inventory, together with a separate
zero-activity assumption. The section bound and fee condition are
unchanged. The address conditions are
\href{https://github.com/yang-mills-group/OpenAI-YM-20261028/blob/3456e5e75cf8d4cd558955a5c7638fba2d68f464/YM/Planned/ATOM_GUARDED_LAYOUT/Guards.lean\#L63-L68}{\texttt{Keeper}}
and
\href{https://github.com/yang-mills-group/OpenAI-YM-20261028/blob/3456e5e75cf8d4cd558955a5c7638fba2d68f464/YM/Planned/ATOM_GUARDED_LAYOUT/Guards.lean\#L89}{\texttt{Keeper.blank}}.
\emph{Lean:}
\href{https://github.com/yang-mills-group/OpenAI-YM-20261028/blob/3456e5e75cf8d4cd558955a5c7638fba2d68f464/YM/Extra/SOURCE_HISTORIES/Birth.lean\#L116-L161}{\texttt{launched}},
\href{https://github.com/yang-mills-group/OpenAI-YM-20261028/blob/3456e5e75cf8d4cd558955a5c7638fba2d68f464/YM/Extra/SOURCE_HISTORIES/Birth.lean\#L105-L106}{\texttt{born}},
\href{https://github.com/yang-mills-group/OpenAI-YM-20261028/blob/3456e5e75cf8d4cd558955a5c7638fba2d68f464/YM/Extra/SOURCE_HISTORIES/Birth.lean\#L24-L25}{\texttt{supplement}},
\href{https://github.com/yang-mills-group/OpenAI-YM-20261028/blob/3456e5e75cf8d4cd558955a5c7638fba2d68f464/YM/Extra/SOURCE_HISTORIES/Birth.lean\#L67}{\texttt{fullCost}},
\href{https://github.com/yang-mills-group/OpenAI-YM-20261028/blob/3456e5e75cf8d4cd558955a5c7638fba2d68f464/YM/Extra/SOURCE_HISTORIES/Birth.lean\#L108-L109}{\texttt{fee}}.

\textbf{Proposition II.1.9 (numerical conditions for the sections and
transfers).}

\begin{enumerate}
\def\labelenumi{(\alph{enumi})}
\tightlist
\item
  \emph{Plans.} Let \(P\) be a stage profile, \(C_j\ge1\), and let
  \((\mathsf K_j)\) be nondecreasing with \(\mathsf K_j\ge1\). For the
  fixed degree-\(d\) frame put
  \(\Omega_d=1+2\sum_{s,i}\lVert s\cdot e_i\rVert_1\), where \(s\)
  ranges over signed permutations and \(i\) over frame indices; the
  factor \(2\) accounts for conjugation. This constant bounds the finite
  symmetry average. Assume \(P\) satisfies the following, writing
  \(\bar\kappa_j\) for the constant of stage \(j\):
\end{enumerate}

\begin{itemize}
\tightlist
\item
  \(j\le\) order and \(j\le\) power of stage \(j\);
\item
  \(\pi(\max(\mathrm{order}_j,\mathrm{power}_j))\le\bar\kappa_j\);
\item
  \((C_j+1)\,\Omega_d\le\mathsf K_j\) for \(d\le j\);
\item
  \(64\mathsf K_{j+1}^2\le\bar\kappa_j\).
\end{itemize}

Then there are lower bounds \(n_j\) for the block factors and coupling
thresholds \(\beta_j^*\) such that every admissible level system with
profile \(P\), block factors at least \(n_j\) at stage \(j\), and
\(\hat\beta_m\ge\beta_{J(m+1)}^*\) satisfies the following numerical
conditions:

\begin{itemize}
\tightlist
\item
  \(J(m)\le\) order and power at level \(m\);
\item
  \(\pi(\max(\mathrm{order},\mathrm{power}))\le\kappa_m\);
\item
  the latch of (b) for all \(d\le J(m)\);
\item
  room: \(64\mathsf K^2_{J(m)+1}\le\kappa_m\);
\item
  bank: \(10^5\mathsf K_{j+1}^6(1+\kappa_{m+1})^2\le L_m^{1/8}\), with
  \(j=J(m+1)\);
\item
  for every \(d\le J(m)\), with \(j=J(m+1)\) and \(\eta=L_m^{-2/5}\),
  \[\eta(1+\kappa_{m+1})\Omega_d\le L_m^{-3/8},\qquad (C_j+\eta\kappa_{m+1})\Omega_d\le\mathsf K_j.\]
\end{itemize}

If a profile already satisfies
\(j\le\mathrm{order}_j,\mathrm{power}_j\), then, for any fixed
\(C_j\ge1\), one can choose \((\mathsf K_j)\) and enlarge only its
constants \(\bar\kappa_j\) to obtain the four stage conditions above.

\begin{enumerate}
\def\labelenumi{(\alph{enumi})}
\setcounter{enumi}{1}
\item
  \emph{Latch.} At level \(m\), \(\mathrm{Latch}(d,m)\) holds if, for
  every \(p\), \(\ell_{m,p}\) is at least the minimal side of
  \(\Sigma_d\), and the fee condition of Prop. II.1.8 holds with
  \(B=c_d\), \(\theta=\theta_m\), \(w=w_m\), at every placement. For
  each \(d\) and support exponent, \(\mathrm{Latch}(d,m)\) holds once
  \(\hat\beta_m\) is large enough.
\item
  \emph{Section clauses of Makes.} If the numerical conditions of (a)
  hold on a history and \(d+3\le J(r(d))\) for all \(d\), then the
  primary sections satisfy clauses (ii)--(v) of \(\mathrm{Makes}\).
\end{enumerate}

\emph{Status:} theorems. The two per-step inequalities are
\href{https://github.com/yang-mills-group/OpenAI-YM-20261028/blob/3456e5e75cf8d4cd558955a5c7638fba2d68f464/YM/Extra/SOURCE_HISTORIES/Replay.lean\#L47-L51}{\texttt{Buys}};
the finite symmetry constant is
\href{https://github.com/yang-mills-group/OpenAI-YM-20261028/blob/3456e5e75cf8d4cd558955a5c7638fba2d68f464/YM/Extra/SOURCE_HISTORIES/Motion.lean\#L121}{\texttt{orbitCost}}.
\emph{Lean:}
\href{https://github.com/yang-mills-group/OpenAI-YM-20261028/blob/3456e5e75cf8d4cd558955a5c7638fba2d68f464/YM/Extra/SOURCE_HISTORIES/Reserve.lean\#L57-L62}{\texttt{Doors.Plans}},
\href{https://github.com/yang-mills-group/OpenAI-YM-20261028/blob/3456e5e75cf8d4cd558955a5c7638fba2d68f464/YM/Extra/SOURCE_HISTORIES/Reserve.lean\#L65-L72}{\texttt{Doors.Budget}},
\href{https://github.com/yang-mills-group/OpenAI-YM-20261028/blob/3456e5e75cf8d4cd558955a5c7638fba2d68f464/YM/Extra/SOURCE_HISTORIES/Reserve.lean\#L76-L102}{\texttt{Doors.afford}},
\href{https://github.com/yang-mills-group/OpenAI-YM-20261028/blob/3456e5e75cf8d4cd558955a5c7638fba2d68f464/YM/Extra/SOURCE_HISTORIES/Reserve.lean\#L108-L133}{\texttt{alone}},
\href{https://github.com/yang-mills-group/OpenAI-YM-20261028/blob/3456e5e75cf8d4cd558955a5c7638fba2d68f464/YM/Extra/SOURCE_HISTORIES/Delays.lean\#L120-L123}{\texttt{Latch}},
\href{https://github.com/yang-mills-group/OpenAI-YM-20261028/blob/3456e5e75cf8d4cd558955a5c7638fba2d68f464/YM/Extra/SOURCE_HISTORIES/Delays.lean\#L126-L157}{\texttt{latch\_cost}},
\href{https://github.com/yang-mills-group/OpenAI-YM-20261028/blob/3456e5e75cf8d4cd558955a5c7638fba2d68f464/YM/Extra/SOURCE_HISTORIES/Delays.lean\#L162-L177}{\texttt{entrances}},
\href{https://github.com/yang-mills-group/OpenAI-YM-20261028/blob/3456e5e75cf8d4cd558955a5c7638fba2d68f464/YM/Extra/SOURCE_HISTORIES/Issue.lean\#L38-L45}{\texttt{Issue.refer}},
\href{https://github.com/yang-mills-group/OpenAI-YM-20261028/blob/3456e5e75cf8d4cd558955a5c7638fba2d68f464/YM/Next/FOUNDATIONS/Primary.lean\#L182-L193}{\texttt{Foundations.refLi}},
\href{https://github.com/yang-mills-group/OpenAI-YM-20261028/blob/3456e5e75cf8d4cd558955a5c7638fba2d68f464/YM/Bid/Loan.lean\#L141-L145}{\texttt{Bid.RefLi}}.

\textbf{Proposition II.1.10 (chessboard comparison at all offsets).} Let
\(\mathcal C_k(x)\) denote the tiling of \(\Lambda_{N,p}\) by link cells
of odd side \(D_{k,N}=a_k/a_N\), for \(k\le N\), centred at \(x\). Write
\(\mathrm{rt}_\beta\) for the chessboard root test of a tiling, and say
that \(F\) is \emph{interior} to a cell if it is measurable in that cell
with a uniform bound.

\begin{enumerate}
\def\labelenumi{(\alph{enumi})}
\item
  \emph{Centred symmetry.} Suppose \(F\) is interior to the central cell
  of \(\mathcal C_k(0)\). For every diagonal sign change \(s\) and every
  \(x\), the function \(F\circ A_{s,x}^*\) is interior to the central
  cell of \(\mathcal C_k(x)\) and has the same root test as \(F\) in
  \(\mathcal C_k(0)\). The same holds after complex conjugation.
  Diagonal sign changes do not permute the axes.
\item
  \emph{All offsets from one centred bound.} Let \(F_o=F_0(\cdot+o)\),
  and assume:
\end{enumerate}

\begin{itemize}
\tightlist
\item
  \(F_0\) is interior to the central cell of \(\mathcal C_k(0)\) with
  root test \(\le A\);
\item
  \(T\) is a tiling whose side is a multiple of the level-\(k\) cell
  side and which satisfies the reflection-positivity seam inequality at
  \(\beta\).
\end{itemize}

Then for every \(o\) in the inner margin of cell \(0\) of \(T\)
(distance \(\ge3\times\) the level-\(k\) side from its boundary),
\(F_o\) is interior to that cell and \(\mathrm{rt}_\beta^T(F_o)\le A\).

\begin{enumerate}
\def\labelenumi{(\alph{enumi})}
\setcounter{enumi}{2}
\tightlist
\item
  \emph{Combinations.} Assume \(\mathrm{Makes}\), fix \(\gamma>0\) and
  an extraction order \(f\), and let \(Q_i(N,p,o)\) be finitely many
  families. Each family is eventually a literal translate of its value
  at \(0\). For each \(i\) and every \(w\), assume there is \(k\ge w\)
  such that eventually \(k\le N\) and \(Q_i(N,p,0)\) is interior to the
  central cell of \(\mathcal C_k(0)\), with root test at most
  \(e^{-\gamma s_N}\). Suppose also that, for every \(\varepsilon>0\),
  eventually \(\lvert r_i(N,p)\rvert\le e^{\varepsilon s_N}\). The
  coefficients are independent of \(o\).
\end{enumerate}

``Eventually'' along \(f\) means: for all sufficiently large volume
indices, then for all sufficiently large mesh indices; this is the fused
order of IV.1. Put \(Q=\sum_i r_iQ_i\) and \(y(N,p)=e^{-\gamma s_N/2}\).
Then \(y\to0\) along this order and \(Q\) is \emph{close}: for every
\(j,m\in\mathbb N\) there is \(k\ge m\) such that eventually \(k\le N\)
and, for every link tiling \(T\) of level-\(k\) side and every
\(o\in T.\mathrm{margin}(0,\lfloor T.\mathrm{side}/(8(j+1))\rfloor)\),
\(Q(N,p,o)\) is interior to cell \(0\) and its root test is at most
\(y(N,p)\).

\emph{Status:} theorems. \emph{Lean:}
\href{https://github.com/yang-mills-group/OpenAI-YM-20261028/blob/3456e5e75cf8d4cd558955a5c7638fba2d68f464/YM/Push/RELAY_DIFFERENCE/Placement.lean\#L134-L160}{\texttt{center\_price}},
\href{https://github.com/yang-mills-group/OpenAI-YM-20261028/blob/3456e5e75cf8d4cd558955a5c7638fba2d68f464/YM/Push/RELAY_DIFFERENCE/Placement.lean\#L192-L226}{\texttt{corner\_root}},
\href{https://github.com/yang-mills-group/OpenAI-YM-20261028/blob/3456e5e75cf8d4cd558955a5c7638fba2d68f464/YM/Push/RELAY_DIFFERENCE/Placement.lean\#L232-L270}{\texttt{offset}},
\href{https://github.com/yang-mills-group/OpenAI-YM-20261028/blob/3456e5e75cf8d4cd558955a5c7638fba2d68f464/YM/Push/RELAY_DIFFERENCE/Combination.lean\#L65-L93}{\texttt{gather\_close}},
\href{https://github.com/yang-mills-group/OpenAI-YM-20261028/blob/3456e5e75cf8d4cd558955a5c7638fba2d68f464/YM/Bid/Drops.lean\#L25}{\texttt{Bid.close}},
\href{https://github.com/yang-mills-group/OpenAI-YM-20261028/blob/3456e5e75cf8d4cd558955a5c7638fba2d68f464/YM/Push/RELAY_DIFFERENCE/CalendarPrices.lean\#L84-L92}{\texttt{roof\_of\_tabs}}
(a bookkeeping corollary). Consumers: III.12, III.13, VII.4, VIII.1.

\subsubsection{2. What it uses}\label{what-it-uses-1}

\textbf{From other units.}

\begin{itemize}
\tightlist
\item
  I.1 (lattice and Wilson measure): tori and the cover \(\mathbb Z^4\),
  periodic placement, gauge transformations, signed permutations and
  their pull-backs, translations, Haar probability measure, Wilson
  energy and expectation
  (\href{https://github.com/yang-mills-group/OpenAI-YM-20261028/blob/3456e5e75cf8d4cd558955a5c7638fba2d68f464/YM/Wilson/Law.lean\#L99}{\texttt{Wilson.moment}}).
\item
  II.2 (polymer activities, norms, record, cells):

  \begin{itemize}
  \tightlist
  \item
    budget points, regularity and sup bounds
    (\href{https://github.com/yang-mills-group/OpenAI-YM-20261028/blob/3456e5e75cf8d4cd558955a5c7638fba2d68f464/YM/RG/Masks.lean\#L125-L136}{\texttt{BudgetPoint}},
    \href{https://github.com/yang-mills-group/OpenAI-YM-20261028/blob/3456e5e75cf8d4cd558955a5c7638fba2d68f464/YM/RG/Masks.lean\#L154-L161}{\texttt{Regular}},
    \href{https://github.com/yang-mills-group/OpenAI-YM-20261028/blob/3456e5e75cf8d4cd558955a5c7638fba2d68f464/YM/RG/Masks.lean\#L164-L166}{\texttt{supBounds}});
  \item
    healthy sets and masks
    (\href{https://github.com/yang-mills-group/OpenAI-YM-20261028/blob/3456e5e75cf8d4cd558955a5c7638fba2d68f464/YM/RG/Masks.lean\#L41}{\texttt{Healthy}},
    \href{https://github.com/yang-mills-group/OpenAI-YM-20261028/blob/3456e5e75cf8d4cd558955a5c7638fba2d68f464/YM/RG/Masks.lean\#L70-L72}{\texttt{masked}});
  \item
    slots
    (\href{https://github.com/yang-mills-group/OpenAI-YM-20261028/blob/3456e5e75cf8d4cd558955a5c7638fba2d68f464/YM/RG/PathActions.lean\#L38-L40}{\texttt{Slot}});
  \item
    polymer representations and lawfulness
    (\href{https://github.com/yang-mills-group/OpenAI-YM-20261028/blob/3456e5e75cf8d4cd558955a5c7638fba2d68f464/YM/RG/History.lean\#L165-L171}{\texttt{Rep}},
    \href{https://github.com/yang-mills-group/OpenAI-YM-20261028/blob/3456e5e75cf8d4cd558955a5c7638fba2d68f464/YM/RG/History.lean\#L181-L183}{\texttt{Rep.Lawful}}),
    together with mass, admissibility and the activity bound;
  \item
    tilings, root test, interior and seam inequality (used in Prop.
    II.1.10).
  \end{itemize}
\item
  III.1 (block-spin step): the blocking push-forward \(T^{(m)}_\beta\)
  and the fact that it preserves total mass (used in Lemma 4.26).
\item
  III.9 (marginal flow): the history, its selection, \(\hat\beta_m\),
  \(\theta_m\), \(w_m\), the stage index \(J(m)\) and the stage
  constants \(\kappa_m\).
\item
  III.2 (order of constants): sheets, profiles and the level systems
  quantified over in Prop. II.1.9.
\item
  III.10: G08 turns per-step estimates into based families; III.11
  supplies the per-step estimates used in VIII.1. Both use the formats
  fixed here.
\item
  III.16 (toolkit): block Taylor polynomials and jets of compositions,
  the Faà di Bruno formula, and the Maxwell exterior calculus (curl,
  recovery of a radial potential from its curvature). The radial-gauge
  identity identifies a homogeneous potential with its closed curvature
  jet.
\item
  IV.1: extraction orders and the fused order (used in Prop.
  II.1.10(c)).
\end{itemize}

\textbf{Standard mathematics.}

\begin{itemize}
\tightlist
\item
  Linear ODEs on matrix groups: existence by an exponential chart and
  left translation; uniqueness through \(\tfrac{d}{dt}(uv^*)=0\).
\item
  Analytic dependence on parameters through the Volterra equation and
  the Neumann series ({[}Dieudonné, \emph{Foundations of Modern
  Analysis}, ch.~X{]}).
\item
  The analytic inverse function theorem in finite dimensions.
\item
  Reynolds averaging over a finite group, and averaging over a compact
  group with Haar measure.
\item
  Polarization of symmetric multilinear forms, and norm equivalence on
  finite-dimensional spaces.
\item
  Radial (Fock--Schwinger) gauge and its Taylor expansion
  {[}Novikov--Shifman--Vainshtein--Zakharov, Fortschr. Phys. 32 (1984)
  585{]}.
\item
  Integration by parts for Haar measure, i.e.~the lattice
  Schwinger--Dyson identity {[}Makeenko--Migdal, Phys. Lett. B 88 (1979)
  135; Chatterjee, Comm. Math. Phys. 366 (2019) 203{]}.
\item
  Chessboard estimates from reflection positivity
  {[}Fröhlich--Israel--Lieb--Simon, Comm. Math. Phys. 62 (1978) 1{]}.
\end{itemize}

\textbf{Imported constants and assumptions.}

\begin{itemize}
\tightlist
\item
  In Prop. II.1.6 the record and \(\mathrm{Makes}\) are hypotheses of
  (c). They are discharged for the final family in VIII.1, using Prop.
  II.1.9(c) and G08.
\item
  In Prop. II.1.8 lawfulness and the fee condition are hypotheses. The
  fee condition at level \(N\) is the latch of Prop. II.1.9(b). Coupling
  thresholds fixed before the block sizes ensure it (III.2), since
  \(w_m\) decays faster than every inverse power of \(\hat\beta_m\).
\item
  The bounds \(\pi(D)\) and \(c_d\) depend on \(G\) and the fixed frames
  and sections, including the finite slot family chosen in G03. They are
  fixed before the block factors; the stage constants must satisfy
  \(\pi(\max(\mathrm{order}_j,\mathrm{power}_j))\le\bar\kappa_j\).
\end{itemize}

\subsubsection{3. Strategy}\label{strategy-1}

Labels are invariant polynomials in radial-gauge Taylor coefficients.
Test holonomies \(U_{J,h}\) compare them with functions of lattice links
by matching derivatives in \(h\) through order \(d\). Finitely many
transported plaquette curvatures give an analytic map of the truncated
jet whose derivative is injective. An equivariant left inverse, followed
by Taylor expansion, produces slot polynomials representing every label.
Averaging over signed permutations gives covariance. Every slot and each
of its admissible left derivatives costs one threshold factor, so a
column of factor degree \(f\) has norm \(O(\theta^f)\). Conversely,
Taylor expansion of a regular lattice function gives a unique read
label. Testing on jets of amplitude \(\theta\) bounds its scaled
coordinates uniformly in the lattice side and support.

The equation sections use Haar integration by parts in one link. For
\(\ell\ge2\), \(\beta\ne0\) and a bounded measurable multiplier
independent of that link, the expectation of the corresponding Wilson
force is zero. Transported forces are linear combinations of slots;
their leading continuum symbols are the free Maxwell divergence. Every
equation polynomial is a sum of products of a linear Maxwell equation
functional and other linear jet functionals. Represent the first factor
by a force and the remaining factors by slot polynomials avoiding its
link. Each product is orthogonal to remote spectators and reads the
prescribed generator modulo higher grade. A unitriangular change of
labels absorbs those higher terms. The resulting equation section
describes vanishing at separated lattice sites; its transfer to null
labels of the limit requires the later comparison arguments.

At the Wilson end the observable is inserted to first order in a
parameter \(\lambda\) with \(\lambda^2=0\). The density coefficient
consists of the vacuum density times a masked local section and regular
remainders, together with a gas in which one forced polymer is marked.
Splicing the section's large-field complement into reserved forced
addresses gives the exact bare-end identity. This uses no division by
the vacuum density. Along the RG history the running coordinates
describe mixing of local observables; their values at the reference
level define \(\tau\), hence the renormalization map
\(S_{1/\theta_N}\tau^{-1}\). This unit fixes the transfer specification
and proves the birth identity; III.10--III.11 supply the per-step
estimates.

\subsubsection{4. Steps}\label{steps-1}

\paragraph{4.1 Label spaces and frames}\label{label-spaces-and-frames}

\textbf{Lemma 4.1 (finiteness and evaluation).} \(V_d\) is
finite-dimensional, and a polynomial in the jet coordinates that
vanishes on all real jets is zero. \emph{Sketch:} every coordinate has
weight \(\ge2\). So a monomial of weight \(\le d\) has exponents
\(\le d\) and involves only \(\mathrm{Tay}_j\) with \(j\le d-2\), and
there are finitely many such monomials. A polynomial vanishing on
\(\mathbb R^n\subset\mathbb C^n\) is zero, because each variable ranges
over an infinite set. \emph{Lean:}
\href{https://github.com/yang-mills-group/OpenAI-YM-20261028/blob/3456e5e75cf8d4cd558955a5c7638fba2d68f464/YM/Labels/Polynomial.lean\#L63-L94}{\texttt{cap\_finite}},
\href{https://github.com/yang-mills-group/OpenAI-YM-20261028/blob/3456e5e75cf8d4cd558955a5c7638fba2d68f464/YM/Labels/Polynomial.lean\#L38-L49}{\texttt{eval\_ext}}.

\textbf{Lemma 4.2 (real graded frames).} Every linearly independent set
of real pure-grade labels in \(V_d\) containing \(1\) extends to a
frame. Consequently frames exist, and there is a coherent sequence of
them. \emph{Sketch:}

\begin{itemize}
\tightlist
\item
  \(V_d\) is spanned by the elements \(q+\bar q\) and \(i(q-\bar q)\),
  where \(q\) runs over the graded components; these elements are real
  and of pure grade.
\item
  Extend the given independent set inside this spanning set, and move
  \(1\) to index \(0\).
\item
  Recursion in \(d\) gives coherence, since \(V_d\subset V_{d+1}\)
  preserves reality and grades.
\item
  \(f_i>0\) for \(i\neq0\), since the only nonzero zero-factor grade is
  \(\mathbb C1\).
\end{itemize}

\emph{Lean:}
\href{https://github.com/yang-mills-group/OpenAI-YM-20261028/blob/3456e5e75cf8d4cd558955a5c7638fba2d68f464/YM/Labels/Basis.lean\#L16-L33}{\texttt{pure\_real\_span}},
\href{https://github.com/yang-mills-group/OpenAI-YM-20261028/blob/3456e5e75cf8d4cd558955a5c7638fba2d68f464/YM/Labels/Basis.lean\#L59-L81}{\texttt{extend\_subset}},
\href{https://github.com/yang-mills-group/OpenAI-YM-20261028/blob/3456e5e75cf8d4cd558955a5c7638fba2d68f464/YM/Labels/Basis.lean\#L83-L86}{\texttt{frame\_exists}},
\href{https://github.com/yang-mills-group/OpenAI-YM-20261028/blob/3456e5e75cf8d4cd558955a5c7638fba2d68f464/YM/Labels/Basis.lean\#L88-L100}{\texttt{extend\_frame}},
\href{https://github.com/yang-mills-group/OpenAI-YM-20261028/blob/3456e5e75cf8d4cd558955a5c7638fba2d68f464/YM/Labels/Basis.lean\#L107-L111}{\texttt{frames}}.

\paragraph{4.2 Test holonomies and the slot map
(G01--G04)}\label{test-holonomies-and-the-slot-map-g01g04}

\textbf{Lemma 4.3 (transport ODE; covariance, G01).}

\begin{itemize}
\tightlist
\item
  Existence and uniqueness hold as stated in Prop. II.1.2(a).
\item
  Scaling: \(U_{J,h}=U_{\mathrm{sc}_hJ,1}\), with only the components
  \(J_j\), \(j\le d-2\), entering.
\item
  Gauge covariance: the links of \(\mathrm{Ad}(g)J\) are the constant
  gauge transform by \(g\) of the links of \(J\).
\item
  Hypercubic covariance: \(s^*U^{(d)}_{J,h}=U^{(d)}_{s^{-1}J,h}\), where
  \(s^{-1}J\) is the jet rotated by \(s^{-1}\in O(4)\).
\item
  In particular each slot of \(\mathrm{Ad}(g)J\) is \(\mathrm{Ad}(g)\)
  applied to the slot of \(J\).
\end{itemize}

\emph{Sketch:}

\begin{itemize}
\tightlist
\item
  Uniqueness: if \(u,v\) solve the same equation, then \(uv^*\) has zero
  derivative.
\item
  Existence: take a local solution in an exponential chart,
  left-translate it, and glue the pieces.
\item
  Scaling and covariance: the transformed function solves the
  transformed equation, so uniqueness concludes; a reflected axis
  reverses the parameter of the ODE.
\end{itemize}

\emph{Lean:}
\href{https://github.com/yang-mills-group/OpenAI-YM-20261028/blob/3456e5e75cf8d4cd558955a5c7638fba2d68f464/YM/Source/Local.lean\#L42-L52}{\texttt{SolvesLink.unique}},
\href{https://github.com/yang-mills-group/OpenAI-YM-20261028/blob/3456e5e75cf8d4cd558955a5c7638fba2d68f464/YM/Source/ODE.lean\#L162-L168}{\texttt{Hol.exists\_ode}},
\href{https://github.com/yang-mills-group/OpenAI-YM-20261028/blob/3456e5e75cf8d4cd558955a5c7638fba2d68f464/YM/Bid/Source.lean\#L23}{\texttt{Bid.G01}},
\href{https://github.com/yang-mills-group/OpenAI-YM-20261028/blob/3456e5e75cf8d4cd558955a5c7638fba2d68f464/YM/Work/G01.lean\#L157-L164}{\texttt{Work.G01.answer}},
\href{https://github.com/yang-mills-group/OpenAI-YM-20261028/blob/3456e5e75cf8d4cd558955a5c7638fba2d68f464/YM/Work/G01.lean\#L138-L146}{\texttt{links\_colored}}.

\textbf{Lemma 4.4 (analytic germ, G02).}

\begin{enumerate}
\def\labelenumi{(\alph{enumi})}
\item
  Let \(x\) range over a finite-dimensional real vector space, let
  \(A_x(t)\) depend linearly on \(x\) and continuously on \(t\), and let
  \(u_x\) solve \(\dot u=uA_x\) with \(u_x(0)=1\). Then
  \(x\mapsto u_x(1)\) is analytic at \(0\), equals \(1\) there, and its
  derivative along every ray is \(\int_0^1A_x\).
\item
  Consequently each link, each walk and each slot is an analytic
  function of the truncated jet near \(0\). The derivative of the slot
  \((\omega,P)\) at \(0\) is the circulation of the edge-integrated
  potential around \(P\). The transport along \(\omega\) does not
  contribute at linear order. \emph{Sketch:} write the ODE as
  \(F=1+T(x)F\) on \(C([0,1])\), where \(T(x)\) is the Volterra
  operator, linear in \(x\), with \(\lVert T(x)\rVert<1\) near \(0\).
  Then \(F=(1-T(x))^{-1}1\), which is analytic since inversion is
  analytic near \(1\) in a Banach algebra. Differentiate the fixed-point
  equation at \(0\). The bounds are link by link; no analytic radius
  uniform over an infinite stencil is claimed. \emph{Lean:}
  \href{https://github.com/yang-mills-group/OpenAI-YM-20261028/blob/3456e5e75cf8d4cd558955a5c7638fba2d68f464/YM/Work/G02.lean\#L89-L130}{\texttt{endpointGerm}},
  \href{https://github.com/yang-mills-group/OpenAI-YM-20261028/blob/3456e5e75cf8d4cd558955a5c7638fba2d68f464/YM/Work/G02.lean\#L140-L177}{\texttt{linkGerm}},
  \href{https://github.com/yang-mills-group/OpenAI-YM-20261028/blob/3456e5e75cf8d4cd558955a5c7638fba2d68f464/YM/Work/G02.lean\#L191-L225}{\texttt{walkGerm}},
  \href{https://github.com/yang-mills-group/OpenAI-YM-20261028/blob/3456e5e75cf8d4cd558955a5c7638fba2d68f464/YM/Work/G02.lean\#L230-L273}{\texttt{slotGerm}},
  \href{https://github.com/yang-mills-group/OpenAI-YM-20261028/blob/3456e5e75cf8d4cd558955a5c7638fba2d68f464/YM/Source/Drops.lean\#L88-L91}{\texttt{Drops.Peg}}.
\end{enumerate}

\textbf{Lemma 4.5 (finitely many slots detect the jet, G03).} For each
\(d\) there are finitely many slots whose linearizations have no common
kernel on the truncated jet space
\(\bigoplus_{j\le d-2}\mathrm{Tay}_j\). \emph{Sketch:}

\begin{itemize}
\tightlist
\item
  \emph{All slots detect the jet.} Suppose all linearized slots vanish
  on \(A\). Write \(N_l(x,t)\) for the difference of the two iterated
  averages of partial derivatives over a square of side \(t\) based at
  \(x\). The unit-square circulation is \(\sum_lN_l(x,1)\). Homogeneity
  gives \(N_l(tx,t)=t^lN_l(x,1)\), hence
  \(\sum_{l\le k}t^{k-l}N_l(tx,t)=0\) when \(A_k\) is the top component
  and \(x\in\mathbb Z^4\).
\item
  \emph{Mesh to zero.} Letting \(t\to0\) with \(tx\to y\) gives
  \(\mathrm{curl}(A_k)(y)=0\) for the top component. In radial gauge the
  potential is recovered from its curvature, so \(A_k=0\). Descend in
  \(k\).
\item
  \emph{Finitely many suffice.} Among finite families choose one whose
  common kernel has minimal dimension. Adding any slot cannot lower the
  dimension, so this kernel lies in the common kernel of all slots,
  which is \(0\).
\end{itemize}

The choice is non-constructive. \emph{Lean:}
\href{https://github.com/yang-mills-group/OpenAI-YM-20261028/blob/3456e5e75cf8d4cd558955a5c7638fba2d68f464/YM/Bid/Source.lean\#L30}{\texttt{Bid.G03}},
\href{https://github.com/yang-mills-group/OpenAI-YM-20261028/blob/3456e5e75cf8d4cd558955a5c7638fba2d68f464/YM/Work/G03.lean\#L164-L185}{\texttt{top\_term}},
\href{https://github.com/yang-mills-group/OpenAI-YM-20261028/blob/3456e5e75cf8d4cd558955a5c7638fba2d68f464/YM/Work/G03.lean\#L136-L149}{\texttt{zero\_of\_zero\_partials}},
\href{https://github.com/yang-mills-group/OpenAI-YM-20261028/blob/3456e5e75cf8d4cd558955a5c7638fba2d68f464/YM/Work/G03.lean\#L212-L224}{\texttt{Work.G03.answer}}.

\textbf{Lemma 4.6 (equivariant left inverse, G04).} The slot map
\(\mathfrak S\) (finitely many slots) is analytic at \(0\). There is a
linear \(C\) with \(C\circ D\mathfrak S(0)=\mathrm{id}\) and
\(C\circ\mathrm{Ad}_g=\mathrm{Ad}_g\circ C\). \emph{Sketch:}
\(D\mathfrak S(0)\) is injective (Lemma 4.5) and equivariant (it is the
derivative of an equivariant map). Its range and orthogonal complement
are \(\mathrm{Ad}\)-invariant, so the inverse composed with the
orthogonal projection onto the range is equivariant. No Haar averaging
is needed. \emph{Lean:}
\href{https://github.com/yang-mills-group/OpenAI-YM-20261028/blob/3456e5e75cf8d4cd558955a5c7638fba2d68f464/YM/Bid/Source.lean\#L37}{\texttt{Bid.G04}},
\href{https://github.com/yang-mills-group/OpenAI-YM-20261028/blob/3456e5e75cf8d4cd558955a5c7638fba2d68f464/YM/Work/G04.lean\#L121-L139}{\texttt{Work.G04.answer}}.

\textbf{Lemma 4.7 (local analytic equivariant inverse).} There is \(f\),
analytic at \(0\), with \(f(0)=0\), \(f\circ\mathfrak S=\mathrm{id}\)
near \(0\), and \(f\circ\mathrm{Ad}_g=\mathrm{Ad}_g\circ f\) near \(0\).
\emph{Sketch:} apply the analytic inverse function theorem to
\(C\circ\mathfrak S\), whose derivative at \(0\) is the identity, and
put \(f=(C\circ\mathfrak S)^{-1}\circ C\). Equivariance follows from
that of \(C\) and \(\mathfrak S\), together with local injectivity.
\emph{Lean:}
\href{https://github.com/yang-mills-group/OpenAI-YM-20261028/blob/3456e5e75cf8d4cd558955a5c7638fba2d68f464/YM/Source/Drops.lean\#L145-L153}{\texttt{Drops.insight}},
\href{https://github.com/yang-mills-group/OpenAI-YM-20261028/blob/3456e5e75cf8d4cd558955a5c7638fba2d68f464/YM/Source/Drops.lean\#L104-L105}{\texttt{Vision}}.

\textbf{Lemma 4.8 (Haar measure in exponential coordinates).} There are
\(r,c>0\) and a measurable function \(J\) such that \(\exp\) is
injective on \(B_r(0)\subset\mathfrak g\), \(J\) is analytic and
positive there, \(J(0)=1\), and \(\exp_*(cJ(A)\mathbf1_{B_r(0)}\,dA)\)
is normalized Haar measure restricted to \(\exp(B_r(0))\), an open set.
\emph{Sketch:} take \(J\) to be the determinant of the derivative of
\(B\mapsto\log(e^{-A}e^B)\) at \(B=A\). It is analytic and equals \(1\)
at \(A=0\), so it is positive on a small ball. The chain rule gives its
transformation under local left translation. Change of variables then
proves the corresponding local measure identity; averaging over the
compact group identifies the scalar \(c\) with probability Haar
normalization. \emph{Lean:}
\href{https://github.com/yang-mills-group/OpenAI-YM-20261028/blob/3456e5e75cf8d4cd558955a5c7638fba2d68f464/YM/LocalH/Haar.lean\#L97-L146}{\texttt{LocalH.actual}},
\href{https://github.com/yang-mills-group/OpenAI-YM-20261028/blob/3456e5e75cf8d4cd558955a5c7638fba2d68f464/YM/Next/FOUNDATIONS/Primary.lean\#L84}{\texttt{Foundations.chart\_ok}}.
This supplies the measure normalization in the stationary-point analysis
(III.6) and the step estimates.

\paragraph{4.3 Primary sections
(G05--G07)}\label{primary-sections-g05g07}

\textbf{Lemma 4.9 (invariant slot polynomials, G05).} Let \(f\) be as in
Lemma 4.7, let \(T_df\) be its block Taylor polynomial of order \(d\) at
\(0\), and put \(P_i:=b_i\circ T_df\). Write \(\boldsymbol\sigma(U)\)
for the chosen finite tuple of slot values. Then:

\begin{itemize}
\tightlist
\item
  \(P_0=1\), each \(P_i\) is \(\mathrm{Ad}\)-invariant and real-valued
  on real slot values, and each monomial of \(P_i\) has degree
  \(\ge f_i\);
\item
  \(P_i(\boldsymbol\sigma(U_{J,h}))\) and \(b_i(\mathrm{sc}_hJ)\) have
  the same \(d\)-jet in \(h\).
\end{itemize}

\emph{Sketch:}

\begin{itemize}
\tightlist
\item
  \(T_df\) is equivariant because \(f\) is equivariant and
  \(\mathrm{Ad}\) is linear.
\item
  \(T_df(0)=f(0)=0\), so every coordinate of \(T_df\) lacks a constant
  term; since \(b_i\) is homogeneous of degree \(f_i\), every monomial
  of \(P_i\) has degree \(\ge f_i\).
\item
  Along \(X(h)=\mathrm{trunc}(\mathrm{sc}_hJ)\), which is smooth with
  \(X(0)=0\), the slots of \(U_{J,h}\) equal \(\mathfrak S(X(h))\).
\item
  Since \(f\circ\mathfrak S=\mathrm{id}\) near \(0\), and the \(d\)-jet
  of a composition depends only on the \(d\)-jets of its factors,
  \(T_df\circ\mathfrak S\circ X\) and \(X\) have the same \(d\)-jet.
\item
  Labels of weight \(\le d\) see only the components of weight
  \(\le d\).
\end{itemize}

All Taylor terms of order \(\le d\) are kept, so \(P_i\) may have degree
\(>d\) in the slots. \emph{Lean:}
\href{https://github.com/yang-mills-group/OpenAI-YM-20261028/blob/3456e5e75cf8d4cd558955a5c7638fba2d68f464/YM/Work/G05.lean\#L113-L139}{\texttt{Work.G05.answer}},
\href{https://github.com/yang-mills-group/OpenAI-YM-20261028/blob/3456e5e75cf8d4cd558955a5c7638fba2d68f464/YM/Work/G05.lean\#L47-L100}{\texttt{recovery\_projects}},
\href{https://github.com/yang-mills-group/OpenAI-YM-20261028/blob/3456e5e75cf8d4cd558955a5c7638fba2d68f464/YM/Work/G05/Taylor.lean\#L84-L102}{\texttt{blockTaylor\_jets}},
\href{https://github.com/yang-mills-group/OpenAI-YM-20261028/blob/3456e5e75cf8d4cd558955a5c7638fba2d68f464/YM/Bid/Source.lean\#L48-L50}{\texttt{Bid.Cumuli}}.

\textbf{Lemma 4.10 (covariant section, G06).} From the \(P_i\) and
hypercubic covariance of the links one obtains a sound section with the
slot structure, whose stencil is a box \([-r,r]^4\) containing all slot
stencils. \emph{Sketch:}

\begin{itemize}
\tightlist
\item
  Write \(\Psi_c(U)=\sum_i c_iP_i(\boldsymbol\sigma(U))\) before
  averaging, and let \(\mathcal H\) be the signed-permutation group. Set
  \(\Phi_c(U)=|\mathcal H|^{-1}\sum_{s\in\mathcal H}\Psi_{s^{-1}\cdot c}(s^*U)\).
  A transformed slot is a slot multiplied by a sign. The coefficient
  action preserves grades, so the average remains a slot polynomial with
  the same lower bound on degrees.
\item
  Gauge invariance: slots transform by \(\mathrm{Ad}\) of the gauge
  transformation at the base point, and the \(P_i\) are
  \(\mathrm{Ad}\)-invariant.
\item
  The projection property survives because jets and labels transform
  compatibly.
\item
  Each slot component is a polynomial in link entries and their
  conjugates, since \(U^{-1}=U^*\).
\end{itemize}

\emph{Lean:}
\href{https://github.com/yang-mills-group/OpenAI-YM-20261028/blob/3456e5e75cf8d4cd558955a5c7638fba2d68f464/YM/Work/G06.lean\#L43-L60}{\texttt{Work.G06.answer}},
\href{https://github.com/yang-mills-group/OpenAI-YM-20261028/blob/3456e5e75cf8d4cd558955a5c7638fba2d68f464/YM/Work/G06.lean\#L20-L28}{\texttt{enclose}},
\href{https://github.com/yang-mills-group/OpenAI-YM-20261028/blob/3456e5e75cf8d4cd558955a5c7638fba2d68f464/YM/Work/G06.lean\#L30-L41}{\texttt{realize}},
\href{https://github.com/yang-mills-group/OpenAI-YM-20261028/blob/3456e5e75cf8d4cd558955a5c7638fba2d68f464/YM/Work/G06/Formula.lean\#L137-L155}{\texttt{entryPoly\_slot}}.

\textbf{Lemma 4.11 (cost, G07).} For all \((k,P)\) there is \(A\) such
that the section has the cost property of Prop. II.1.3, item 3.
\emph{Sketch:}

\begin{itemize}
\tightlist
\item
  \emph{One slot costs \(t\).} On \(\mathrm{Hy}_t(S)\) a slot has norm
  \(O(t)\). Each left derivative in an admissible direction \(A\), with
  \(\lvert A_e\rvert\le t(1+\mathrm{dist})^P\), is also \(O(t)\), with
  constants depending on the stencil.
\item
  \emph{Algebra.} These bounds are stable under sums and products, with
  exponents of \(t\) adding, and \(t\le1\) lowers exponents.
\item
  \emph{Columns.} Every monomial of column \(i\) has \(\ge f_i\) slot
  factors, which gives \(c_it^{f_i}\); take \(A=\sum_ic_i\).
\end{itemize}

\(A\) is not uniform in \(d,k,P\). \emph{Lean:}
\href{https://github.com/yang-mills-group/OpenAI-YM-20261028/blob/3456e5e75cf8d4cd558955a5c7638fba2d68f464/YM/Work/G07.lean\#L79-L102}{\texttt{Work.G07.answer}},
\href{https://github.com/yang-mills-group/OpenAI-YM-20261028/blob/3456e5e75cf8d4cd558955a5c7638fba2d68f464/YM/Work/G07.lean\#L64-L77}{\texttt{paid\_lumin}},
\href{https://github.com/yang-mills-group/OpenAI-YM-20261028/blob/3456e5e75cf8d4cd558955a5c7638fba2d68f464/YM/Work/G07/Stencil.lean\#L180-L199}{\texttt{paid\_slot}},
\href{https://github.com/yang-mills-group/OpenAI-YM-20261028/blob/3456e5e75cf8d4cd558955a5c7638fba2d68f464/YM/Work/G07/Stencil.lean\#L154-L163}{\texttt{raw\_slot}}.

\textbf{Lemma 4.12 (assembly = Prop. II.1.3).} Lemmas 4.3--4.11 combine
into sections for all \(d\).

\begin{itemize}
\tightlist
\item
  The uniform price is \(\pi(D)=1+\sum_{d,k,P\le D}\max(0,A_{d,k,P})\),
  which is nondecreasing in \(D\).
\item
  The value bound \(c_d\) holds because each column is a continuous
  function of finitely many variables in the compact group \(G\).
\end{itemize}

\emph{Lean:}
\href{https://github.com/yang-mills-group/OpenAI-YM-20261028/blob/3456e5e75cf8d4cd558955a5c7638fba2d68f464/YM/Bid/Source.lean\#L76-L86}{\texttt{Bid.teaching}},
\href{https://github.com/yang-mills-group/OpenAI-YM-20261028/blob/3456e5e75cf8d4cd558955a5c7638fba2d68f464/YM/Integrated.lean\#L216-L219}{\texttt{Integrated.source\_sections}},
\href{https://github.com/yang-mills-group/OpenAI-YM-20261028/blob/3456e5e75cf8d4cd558955a5c7638fba2d68f464/YM/Next/FOUNDATIONS/Primary.lean\#L133-L142}{\texttt{priced}},
\href{https://github.com/yang-mills-group/OpenAI-YM-20261028/blob/3456e5e75cf8d4cd558955a5c7638fba2d68f464/YM/Next/FOUNDATIONS/Primary.lean\#L147-L163}{\texttt{full\_value\_bound}}.

\paragraph{4.4 Readout}\label{readout}

\textbf{Lemma 4.13 (logarithmic chart of test holonomies).} Fix \(d\), a
torus, a site \(o\) and a jet \(J\). Near \(h=0\), the torus test
configuration \(U^{(d),o}_{J,h}\) is the exponential of an analytic
function of \(h\) vanishing at \(0\). Thus
\(h\mapsto f(U^{(d),o}_{J,h})\) is \(C^d\) at \(0\) whenever
\(A\mapsto f(e^A)\) is \(C^d\) at \(A=0\). \emph{Sketch:} apply the
local group logarithm to each test link, use Lemma 4.4 and the
polynomial scaling of the jet, and intersect the finitely many
neighbourhoods on the torus. The logarithm inverts the exponential
there, so the assertion follows by composition. \emph{Lean:}
\href{https://github.com/yang-mills-group/OpenAI-YM-20261028/blob/3456e5e75cf8d4cd558955a5c7638fba2d68f464/YM/Next/FOUNDATIONS/Readout.lean\#L110-L114}{\texttt{testLog\_diff}},
\href{https://github.com/yang-mills-group/OpenAI-YM-20261028/blob/3456e5e75cf8d4cd558955a5c7638fba2d68f464/YM/Next/FOUNDATIONS/Readout.lean\#L119-L125}{\texttt{increment\_testLog}}.

\textbf{Lemma 4.14 (existence and uniqueness of the read label = Prop.
II.1.4(a)).} Fix \(d\), a torus and \(o\). If \(f\) is gauge invariant
and \(A\mapsto f(e^A)\) is \(C^d\) at \(A=0\), there is a unique
\(p\in V_d\) such that, for every \(J\), \(f(U^{(d),o}_{J,h})\) and
\(p(\mathrm{sc}_hJ)\) have equal \(d\)-jets at \(h=0\). \emph{Sketch:}

\begin{itemize}
\tightlist
\item
  \emph{Polynomial in the jet.} Take the block Taylor polynomial of
  order \(d\) at \(0\) of the map \(X\mapsto f(\text{test links of }X)\)
  on the truncated jet space. It is a polynomial in the jet coordinates;
  truncate it to weight \(\le d\).
\item
  \emph{Jets agree.} Since \(\mathrm{trunc}(\mathrm{sc}_hJ)\) has
  components \(h^{j+2}J_j\), truncation does not change
  \(h\)-derivatives of order \(\le d\).
\item
  \emph{Invariance.} Gauge invariance of \(f\) and covariance of the
  test links (Lemma 4.3) make each weight component of \(p\) invariant.
\item
  \emph{Uniqueness.} The \(h\)-derivatives at \(0\) determine the weight
  components, and then Lemma 4.1 applies.
\end{itemize}

For regular \(f\) with order \(\ge d\) the hypothesis holds at the
identity configuration, which is healthy; wrapping supports use the
constant convention. \emph{Lean:}
\href{https://github.com/yang-mills-group/OpenAI-YM-20261028/blob/3456e5e75cf8d4cd558955a5c7638fba2d68f464/YM/Next/FOUNDATIONS/Readout.lean\#L218-L233}{\texttt{read\_projection}},
\href{https://github.com/yang-mills-group/OpenAI-YM-20261028/blob/3456e5e75cf8d4cd558955a5c7638fba2d68f464/YM/Next/FOUNDATIONS/Readout.lean\#L238-L247}{\texttt{exists\_read}}.

\textbf{Lemma 4.15 (uniform bounds along scaled test jets).}

\begin{enumerate}
\def\labelenumi{(\alph{enumi})}
\item
  Fix \(d\). For each jet \(J\) and \(k\in\mathbb N\) there is \(C>0\)
  such that, for all torus sides, sites \(o\), positive links \(e\) and
  \(\theta\in[0,1]\),
  \[\left\lvert\left.\partial_h^k\log U^{(d),o}_{\theta J,h}(e)\right|_{h=0}\right\rvert\le C\,\theta\,(1+\mathrm{dist}(e_-,o))^{\max(k-1,0)}.\]
\item
  For \(k\le d\) there is \(C_{d,k,J}\ge0\) such that
  \(\lvert\partial_h^kf(U^{(d),o}_{t_bJ,h})\vert_{h=0}\rvert\le C_{d,k,J}c\)
  whenever \(d\le\min(k_b,P_b)\) and \(f\) is \(b\)-regular with sup
  bound \(c\). The bound is uniform in side, placement, support and
  \(b\). \emph{Sketch:}
\end{enumerate}

\begin{itemize}
\tightlist
\item
  (a): set \(r=(1+\mathrm{dist}(e_-,o))^{-1}\). The rescaled initial
  position lies in a fixed compact ball. Express the link at parameter
  \(hr\) by a polynomial-in-time ODE profile of amplitude \(\theta r\).
  Its logarithmic endpoint is analytic near the zero profile. The
  \(k\)-th \(h\)-derivative at \(0\) varies smoothly over the compact
  set of normalized positions, lengths and amplitudes, and vanishes at
  amplitude \(0\). A uniform Lipschitz bound in the amplitude gives
  \(C\theta r\). Undoing the rescaling gives \(C\theta r^{1-k}\) for
  \(k\ge1\); the order-zero logarithm is \(0\).
\item
  (b): apply Faà di Bruno to \((f\circ\exp)\circ\log U\) at the
  identity. Its terms contain symmetric Fréchet derivatives of
  \(A\mapsto f(e^A)\), evaluated on the logarithmic derivatives from
  (a). Pull each tensor back to the span of these finitely many
  directions, and bound it by its diagonal values through polarization.
  Repeated left derivatives in one direction equal derivatives along the
  corresponding exponential ray. This avoids a factor depending on the
  number of lattice links; it does not assume that ordered left
  derivatives commute.
\end{itemize}

\emph{Lean:}
\href{https://github.com/yang-mills-group/OpenAI-YM-20261028/blob/3456e5e75cf8d4cd558955a5c7638fba2d68f464/YM/Next/FOUNDATIONS/LinkBounds.lean\#L229-L273}{\texttt{testLog\_bound}},
\href{https://github.com/yang-mills-group/OpenAI-YM-20261028/blob/3456e5e75cf8d4cd558955a5c7638fba2d68f464/YM/Next/FOUNDATIONS/TestBounds.lean\#L194-L223}{\texttt{test\_bound}},
\href{https://github.com/yang-mills-group/OpenAI-YM-20261028/blob/3456e5e75cf8d4cd558955a5c7638fba2d68f464/YM/Next/FOUNDATIONS/TestBounds.lean\#L123-L184}{\texttt{mixed\_budget}}.

\textbf{Lemma 4.16 (readout estimate = Prop. II.1.4(b)).} \emph{Sketch:}

\begin{itemize}
\tightlist
\item
  \emph{Functionals.} For a jet \(J\) and \(k\le d\) let
  \(\lambda_{J,k}(c):=\partial_h^k\,p_c(\mathrm{sc}_hJ)\vert_{h=0}\).
\item
  \emph{Scaling.} Since \(p_{S_\theta c}(K)=p_c(\theta K)\), one has
  \(\lambda_{J,k}(S_\theta\,\mathrm{coord}(p))=\partial^k_hp(\mathrm{sc}_h(\theta J))\vert_{h=0}=\partial_h^kf(U^{(d),o}_{\theta J,h})\vert_{h=0}\).
\item
  \emph{Bound.} By Lemma 4.15(b) this is bounded by \(C_{J,k}c\).
\item
  \emph{Norm equivalence.} The \(\lambda_{J,k}\) separate points (Lemma
  4.14), and on a finite-dimensional space a separating family of
  functionals controls the norm.
\item
  \emph{Wrapping case.} \(p=f(1)\cdot1\), and \(\lvert f(1)\rvert\le c\)
  because the identity configuration is healthy.
\end{itemize}

The constant depends on the holonomy datum and the frame only.
\emph{Lean:}
\href{https://github.com/yang-mills-group/OpenAI-YM-20261028/blob/3456e5e75cf8d4cd558955a5c7638fba2d68f464/YM/Next/FOUNDATIONS/TestBounds.lean\#L234-L265}{\texttt{coeff\_bound}},
\href{https://github.com/yang-mills-group/OpenAI-YM-20261028/blob/3456e5e75cf8d4cd558955a5c7638fba2d68f464/YM/Next/FOUNDATIONS/TestBounds.lean\#L34-L67}{\texttt{budget\_probes}},
\href{https://github.com/yang-mills-group/OpenAI-YM-20261028/blob/3456e5e75cf8d4cd558955a5c7638fba2d68f464/YM/Next/FOUNDATIONS/JetReads.lean\#L158-L168}{\texttt{read\_frame\_scaled}},
\href{https://github.com/yang-mills-group/OpenAI-YM-20261028/blob/3456e5e75cf8d4cd558955a5c7638fba2d68f464/YM/Next/FOUNDATIONS/TestBounds.lean\#L271-L290}{\texttt{price\_of\_readout}}.

\textbf{Lemma 4.17 (placed sections read their labels = Prop.
II.1.4(c)).} Let \(\Sigma\) be sound. If
\(20(\lvert x_\mu\rvert+1)\le\ell\) on the stencil, the footprint does
not wrap around \(o\), the periodic lift of the stencil is injective,
and the placed function along the test holonomies is the cover function
along \(U_{J,h}\). Soundness (projection property) then concludes.
\emph{Lean:}
\href{https://github.com/yang-mills-group/OpenAI-YM-20261028/blob/3456e5e75cf8d4cd558955a5c7638fba2d68f464/YM/Next/FOUNDATIONS/JetReads.lean\#L240-L251}{\texttt{nonwrapping\_placement}},
\href{https://github.com/yang-mills-group/OpenAI-YM-20261028/blob/3456e5e75cf8d4cd558955a5c7638fba2d68f464/YM/Next/FOUNDATIONS/JetReads.lean\#L253-L267}{\texttt{placed\_reads}}.

\paragraph{4.5 Equation sections}\label{equation-sections}

\textbf{Lemma 4.18 (lattice Schwinger--Dyson identity).} Assume
\(\ell\ge2\). \emph{Gradient.} For a link \(e\) and \(A\in\mathfrak g\),
\(\tfrac{d}{dt}S_W(e^{tA}U_e)\vert_{0}=\langle A,\nabla_eS_W(U)\rangle\).
With \(F_{ab}(x)\) based at \(x\), put
\(\Delta_cF_{ab}(x)=F_{ab}(x)-\mathrm{Ad}(U_{x-e_c,c}^{-1})F_{ab}(x-e_c)\).
Then
\[\nabla_{(x,\nu)}S_W=\sum_{\nu<b}\Delta_bF_{\nu b}(x)-\sum_{a<\nu}\Delta_aF_{a\nu}(x).\]
This local formula also defines the force on cover configurations.
\emph{Integration by parts.} For continuous \(W,W'\) with \(W'\) the
derivative of \(W\) along \(e^{tA}U_e\), and bounded measurable \(S\)
invariant under that motion,
\[\big\langle(W'-\beta\langle A,\nabla_eS_W\rangle W)S\big\rangle_\beta=0 .\]
For \(\beta\neq0\) and \(W=1\) this gives
\(\langle\langle A,\nabla_eS_W\rangle S\rangle_\beta=0\). \emph{Sketch:}
the motion \(U_e\mapsto e^{tA}U_e\) preserves the product Haar measure,
and the Boltzmann factor is differentiated under the integral.
\emph{Lean:}
\href{https://github.com/yang-mills-group/OpenAI-YM-20261028/blob/3456e5e75cf8d4cd558955a5c7638fba2d68f464/YM/Push/RELAY_MARK/Wilson.lean\#L232-L237}{\texttt{field\_deriv}},
\href{https://github.com/yang-mills-group/OpenAI-YM-20261028/blob/3456e5e75cf8d4cd558955a5c7638fba2d68f464/YM/Push/RELAY_MARK/Wilson.lean\#L96-L101}{\texttt{turn\_preserving}},
\href{https://github.com/yang-mills-group/OpenAI-YM-20261028/blob/3456e5e75cf8d4cd558955a5c7638fba2d68f464/YM/Push/RELAY_MARK/Wilson.lean\#L321-L340}{\texttt{ibp}},
\href{https://github.com/yang-mills-group/OpenAI-YM-20261028/blob/3456e5e75cf8d4cd558955a5c7638fba2d68f464/YM/Push/RELAY_MARK/Wilson.lean\#L342-L348}{\texttt{equation\_silent}}.

\textbf{Lemma 4.19 (transported forces are micro).} Let \(b\) be a
positive cover link, let \(\omega\) be a walk from \(0\) avoiding \(b\),
and let \(g\) be continuous. Assume there is \(M\) such that, on every
torus of side \(\ell\ge M\) and at every placement \(o\), the placed
\(g\) is invariant under left increments of the placed link \(b\). Then
\(U\mapsto\langle X,\mathrm{Ad}(U_\omega)\nabla_bS_W(U)\rangle\,g(U)\)
is micro, for every \(X\in\mathfrak g\). \emph{Sketch:}

\begin{itemize}
\tightlist
\item
  Write
  \(\langle X,\mathrm{Ad}(U_\omega)Y\rangle=\sum_a\langle\mathrm{Ad}(U_\omega)^{-1}X,E_a\rangle\langle E_a,Y\rangle\)
  in an orthonormal basis.
\item
  Once the torus is large enough, the placed link \(b\) lies within a
  fixed radius of \(o\).
\item
  Then every factor other than \(\langle E_a,\nabla_bS_W\rangle\)
  (namely the transport, \(g\), and a spectator supported beyond that
  radius) is invariant under moving \(U_b\). Apply Lemma 4.18.
\end{itemize}

The periodic independence hypothesis holds, in particular, if \(g\)
depends on finitely many cover links, all distinct from \(b\). No
continuity or gauge invariance of the spectator is needed, and \(\beta\)
may be negative. \emph{Lean:}
\href{https://github.com/yang-mills-group/OpenAI-YM-20261028/blob/3456e5e75cf8d4cd558955a5c7638fba2d68f464/YM/Push/RELAY_MARK/Avoid.lean\#L170-L218}{\texttt{micro\_force}},
\href{https://github.com/yang-mills-group/OpenAI-YM-20261028/blob/3456e5e75cf8d4cd558955a5c7638fba2d68f464/YM/Push/RELAY_MARK/Avoid.lean\#L149-L150}{\texttt{force}}.

\textbf{Lemma 4.20 (forces span the equation functionals).}

\begin{enumerate}
\def\labelenumi{(\alph{enumi})}
\item
  Let \(J\) be truncated through \(J_j\), with \(j\ge1\), and write
  \(\mathfrak F_\nu(J;x)\) for the linearization at the zero jet of the
  transported force at \((x,\nu)\). It is independent of the transport
  path. For every \(y\in\mathbb R^4\) and \(x_t\in\mathbb Z^4\) with
  \(tx_t\to y\) as \(t\downarrow0\),
  \[\lim_{t\downarrow0}t^{j-1}\mathfrak F_\nu(J;x_t)=-\sum_a\partial_a(\mathrm{curl}\,J_j)_{a\nu}(y).\]
  For \(j=0\) the linearized force vanishes. At finite mesh the force is
  a transverse difference of plaquette averages.
\item
  For jets truncated through weight \(j+2\), if every linearized
  transported force on a walk ending at the initial vertex of its link
  and avoiding that link vanishes, then the component \(J_j\) lies in
  \(\mathrm{Sh}_j\).
\item
  Dually, for every real linear functional \(Q\) on \(\mathrm{Tay}_j\)
  vanishing on \(\mathrm{Sh}_j\), the functional
  \(Q\circ\mathrm{proj}_{\mathrm{top}}\) is a finite combination of
  linearized transported forces. \emph{Sketch:} express each transverse
  difference as an additional integral of a derivative of the plaquette
  average. Homogeneity and continuity of these averages give (a). If
  every force vanishes, every lattice point can be reached by a walk
  avoiding the varied link; (a) then makes the top curvature
  divergence-free. Its Bianchi identity follows because it is a curl,
  proving (b). In finite dimensions, the annihilator of the common force
  kernel is the span of the force functionals, giving (c). \emph{Lean:}
  \href{https://github.com/yang-mills-group/OpenAI-YM-20261028/blob/3456e5e75cf8d4cd558955a5c7638fba2d68f464/YM/Push/RELAY_MARK/Current.lean\#L208-L251}{\texttt{continuum\_current}},
  \href{https://github.com/yang-mills-group/OpenAI-YM-20261028/blob/3456e5e75cf8d4cd558955a5c7638fba2d68f464/YM/Push/RELAY_MARK/Current.lean\#L253-L280}{\texttt{top\_shell}},
  \href{https://github.com/yang-mills-group/OpenAI-YM-20261028/blob/3456e5e75cf8d4cd558955a5c7638fba2d68f464/YM/Push/RELAY_MARK/Current.lean\#L286-L297}{\texttt{force\_kernel\_last}},
  \href{https://github.com/yang-mills-group/OpenAI-YM-20261028/blob/3456e5e75cf8d4cd558955a5c7638fba2d68f464/YM/Push/RELAY_MARK/Span.lean\#L32-L41}{\texttt{force\_span}}.
\end{enumerate}

\textbf{Lemma 4.21 (generators of the equation ideal).}
\(p\in\mathcal E_d\) if and only if \(p\circ\Pi=0\), where \(\Pi\) is
the projection of each \(\mathrm{Tay}_j\) onto \(\mathrm{Sh}_j\). Every
\(p\in\mathcal E_d\) of pure grade \(m\) is a combination of products
\(\ell\cdot M\), of total grade \(m\), of a linear functional \(\ell\)
on \(\mathrm{Tay}_j\) vanishing on \(\mathrm{Sh}_j\) and a product \(M\)
of real linear jet functionals, with complex scalar coefficients. Grade
components of elements of \(\mathcal E_d\) lie in \(\mathcal E_d\), and
\(\mathcal E_d\) contains no nonzero constant. \emph{Sketch:} write
\(p=p-p\circ\Pi\). For a monomial \(M=\prod_r\ell_r\), the difference
\(M-M\circ\Pi\) telescopes into
\(\sum_r(\ell_r-\ell_r\circ\Pi)\cdot(\cdots)\), and each
\(\ell_r-\ell_r\circ\Pi\) vanishes on \(\mathrm{Sh}\). \emph{Lean:}
\href{https://github.com/yang-mills-group/OpenAI-YM-20261028/blob/3456e5e75cf8d4cd558955a5c7638fba2d68f464/YM/Push/RELAY_MARK/Ideal.lean\#L58-L68}{\texttt{erase\_shell}},
\href{https://github.com/yang-mills-group/OpenAI-YM-20261028/blob/3456e5e75cf8d4cd558955a5c7638fba2d68f464/YM/Push/RELAY_MARK/Ideal.lean\#L44-L45}{\texttt{multiplied}},
\href{https://github.com/yang-mills-group/OpenAI-YM-20261028/blob/3456e5e75cf8d4cd558955a5c7638fba2d68f464/YM/Push/RELAY_MARK/Ideal.lean\#L93-L112}{\texttt{generation\_ready}}.

\textbf{Lemma 4.22 (micro lifts at pure grade).} For
\(p\in\mathcal E_d\) of grade \(m\) (with \(m_1\le d\)) there are a
gauge-invariant lattice function \(f\) and a label \(q\) with:

\begin{itemize}
\tightlist
\item
  \(f\) represents \(q\) to order \(d\);
\item
  \(q-p\) has only components of strictly higher grade;
\item
  \(f\) is a slot polynomial with \(\ge m_2\) slot factors;
\item
  \(f\) is micro.
\end{itemize}

\emph{Sketch:}

\begin{itemize}
\tightlist
\item
  \emph{Generators.} For a generator \(\ell\cdot M\), write
  \(\ell\circ\mathrm{proj}_{\mathrm{top}}\) as a combination of
  linearized forces (Lemma 4.20).
\item
  \emph{Lift of a force.} A transported lattice force component is a
  degree-one slot polynomial. Its readout is the full linearized force
  plus terms with at least two jet factors.
\item
  \emph{Lift of the monomial.} Slots avoiding one specified link still
  detect every truncated jet: excluding that link removes only finitely
  many plaquette positions, which does not change the mesh-to-zero
  argument away from \(0\); continuity gives the value at \(0\). Thus
  every factor of \(M\) can be represented, through all retained
  weights, by a slot polynomial independent of that link after periodic
  placement.
\item
  \emph{Product.} The force combination matches \(\ell\) through its own
  weight \(j+2\); extending it to cutoff \(d\) adds only higher-weight
  one-factor terms. Nonlinear readout errors have more factors.
  Multiplication by the lifted \(M\) therefore gives \(\ell M\) modulo
  higher grade. The product is micro by Lemma 4.19. Combine linearly.
\item
  \emph{Gauge invariance.} Average over constant gauge transformations
  with Haar measure. The orbit spans a finite-dimensional space, so the
  average stays a slot polynomial and the relation above persists.
\end{itemize}

\emph{Lean:}
\href{https://github.com/yang-mills-group/OpenAI-YM-20261028/blob/3456e5e75cf8d4cd558955a5c7638fba2d68f464/YM/Push/RELAY_MARK/Rank.lean\#L100-L127}{\texttt{stanza}},
\href{https://github.com/yang-mills-group/OpenAI-YM-20261028/blob/3456e5e75cf8d4cd558955a5c7638fba2d68f464/YM/Push/RELAY_MARK/Rank.lean\#L129-L153}{\texttt{coherent\_verse}},
\href{https://github.com/yang-mills-group/OpenAI-YM-20261028/blob/3456e5e75cf8d4cd558955a5c7638fba2d68f464/YM/Push/RELAY_MARK/Rank.lean\#L157-L166}{\texttt{mounted\_read}},
\href{https://github.com/yang-mills-group/OpenAI-YM-20261028/blob/3456e5e75cf8d4cd558955a5c7638fba2d68f464/YM/Push/RELAY_MARK/Gauge.lean\#L110-L143}{\texttt{averaged}},
\href{https://github.com/yang-mills-group/OpenAI-YM-20261028/blob/3456e5e75cf8d4cd558955a5c7638fba2d68f464/YM/Push/RELAY_MARK/Gauge.lean\#L147-L163}{\texttt{positioned}},
\href{https://github.com/yang-mills-group/OpenAI-YM-20261028/blob/3456e5e75cf8d4cd558955a5c7638fba2d68f464/YM/Push/RELAY_MARK/Rank.lean\#L78-L83}{\texttt{Cone}}.

\textbf{Lemma 4.23 (a covariant lifting map).} There is a linear map
\(F:V_d\to V_d\times\{\text{lattice functions}\}\), with \(F(p)=(q,f)\),
\(F(1)=(1,1)\), such that for \(p\) of grade \(m\):

\begin{itemize}
\tightlist
\item
  \(f\) is gauge invariant and represents \(q\);
\item
  \(q-p\) has only components of higher grade;
\item
  \(f\) has \(\ge m_2\) slot factors;
\item
  \(f\) is micro whenever \(p\in\mathcal E_d\).
\end{itemize}

\(F\) commutes with signed permutations and conjugation. \emph{Sketch:}

\begin{itemize}
\tightlist
\item
  \emph{One grade.} Choose a linear projection
  \(P_m:V_d\to V_d^{(m)}\cap\mathcal E_d\). Apply the micro lift to
  \(P_mp\), and the primary-section lift
  \(q\mapsto(q,\Phi_{\mathrm{coord}(q)})\) to \(p-P_mp\); call the sum
  \(X_m(p)\). On pure grade \(m\) both summands have the required grade
  and slot bounds. On \(V_d^{(m)}\cap\mathcal E_d\) the second summand
  vanishes.
\item
  \emph{All grades.} Sum \(X_m\circ(\text{grade-}m\text{ part})\) over
  \(m\).
\item
  \emph{Symmetrization.} Average over the signed-permutation action, and
  with the conjugate. These operations preserve all properties above.
\end{itemize}

\emph{Lean:}
\href{https://github.com/yang-mills-group/OpenAI-YM-20261028/blob/3456e5e75cf8d4cd558955a5c7638fba2d68f464/YM/Push/RELAY_MARK/Linear.lean\#L99-L102}{\texttt{shelf}},
\href{https://github.com/yang-mills-group/OpenAI-YM-20261028/blob/3456e5e75cf8d4cd558955a5c7638fba2d68f464/YM/Push/RELAY_MARK/Linear.lean\#L33-L38}{\texttt{Offered}},
\href{https://github.com/yang-mills-group/OpenAI-YM-20261028/blob/3456e5e75cf8d4cd558955a5c7638fba2d68f464/YM/Push/RELAY_MARK/Linear.lean\#L65-L67}{\texttt{Payment}},
\href{https://github.com/yang-mills-group/OpenAI-YM-20261028/blob/3456e5e75cf8d4cd558955a5c7638fba2d68f464/YM/Push/RELAY_MARK/Linear.lean\#L69-L71}{\texttt{Campaign}},
\href{https://github.com/yang-mills-group/OpenAI-YM-20261028/blob/3456e5e75cf8d4cd558955a5c7638fba2d68f464/YM/Push/RELAY_MARK/Linear.lean\#L155-L168}{\texttt{unsym}},
\href{https://github.com/yang-mills-group/OpenAI-YM-20261028/blob/3456e5e75cf8d4cd558955a5c7638fba2d68f464/YM/Push/RELAY_MARK/HGroup.lean\#L106-L125}{\texttt{impression}}.

\textbf{Lemma 4.24 (weave and section = Prop. II.1.5).} \emph{Weave.}
Let \(R\) be the first component of \(F\). It is unitriangular for the
order of grades, hence an automorphism. \emph{Section.} The inverse of
\(R\) also does not lower grades, so \(R^{-1}b_i\) has no component
preceding the grade of \(b_i\). Hence the map
\(b_i\mapsto f(R^{-1}b_i)\) sends basis vectors to slot polynomials with
\(\ge f_i\) factors. Enclose all their stencils in one box and rewrite
them as entry polynomials; this gives a section \(\Sigma\) with
\(\Phi_{\mathrm{coord}(Rp)}=f(p)\). \emph{Properties.}

\begin{itemize}
\tightlist
\item
  Soundness: gauge invariance, unit, covariance and reality are
  inherited from \(F\); the projection property holds because \(f(p)\)
  represents \(Rp\).
\item
  Cost: G07 (Lemma 4.11).
\item
  Microness: for \(p\in\mathcal E_d\), \(f(p)\) is the sum of the micro
  grade components of \(F(p)\).
\end{itemize}

Choose these sections and automorphisms for every \(d\). \emph{Lean:}
\href{https://github.com/yang-mills-group/OpenAI-YM-20261028/blob/3456e5e75cf8d4cd558955a5c7638fba2d68f464/YM/Push/RELAY_MARK/Automorphism.lean\#L72-L73}{\texttt{twist}},
\href{https://github.com/yang-mills-group/OpenAI-YM-20261028/blob/3456e5e75cf8d4cd558955a5c7638fba2d68f464/YM/Push/RELAY_MARK/Automorphism.lean\#L43-L50}{\texttt{regrade\_tri}},
\href{https://github.com/yang-mills-group/OpenAI-YM-20261028/blob/3456e5e75cf8d4cd558955a5c7638fba2d68f464/YM/Push/RELAY_MARK/Automorphism.lean\#L100-L102}{\texttt{paid\_reel}},
\href{https://github.com/yang-mills-group/OpenAI-YM-20261028/blob/3456e5e75cf8d4cd558955a5c7638fba2d68f464/YM/Push/RELAY_MARK/Portfolio.lean\#L60-L97}{\texttt{weaveBuild}},
\href{https://github.com/yang-mills-group/OpenAI-YM-20261028/blob/3456e5e75cf8d4cd558955a5c7638fba2d68f464/YM/Push/RELAY_MARK/Portfolio.lean\#L101-L105}{\texttt{Portfolio.build}},
\href{https://github.com/yang-mills-group/OpenAI-YM-20261028/blob/3456e5e75cf8d4cd558955a5c7638fba2d68f464/YM/Push/RELAY_MARK/Assembly.lean\#L30-L33}{\texttt{commissioned}}.

\paragraph{4.6 Source histories}\label{source-histories}

\textbf{Lemma 4.25 (splicing a large-field multiplier).} Let \(f\) be
bounded and vanish on \(\mathrm{Hy}_\theta(Y)\), let \(\rho\) be lawful
with forced-activity parameter \(w\), and let a shelf as in Prop. II.1.8
be given. Then for every configuration \(U\),
\[\sum_b\sum_i\mathrm{line}_i(f)(U)\,\mathbf 1_{\{i\in b,\ b\ \mathrm{adm}\}}\prod_{j\in b\setminus\{i\}}K_j(U)=f(U)\,\mathrm{gas}(U),\]
as an identity on all configurations. The \emph{line} coefficients
satisfy
\(\sum_i\mathrm{mass}(X_i,p_i)\le\sup\lvert f\rvert\cdot\mathrm{mass}(Y,1)(e^{\lvert Y\rvert w}-1)\),
where \(p_i\ge\sup_U\lvert\mathrm{line}_i(f)(U)\rvert\). \emph{Sketch:}

\begin{itemize}
\tightlist
\item
  On a large-field configuration every admissible polymer set meets
  \(Y\), because the forced polymers cover the flagged sites.
\item
  For each nonempty compatible set \(c\) of old polymers meeting \(Y\),
  reserve one address with support \(Y\cup\bigcup_{i\in c}X_i\) and
  inventory \(\bigcup_{i\in c}I_i\). Its coefficient is
  \(f\prod_{i\in c}K_i\) when the flagged sites in that support equal
  the inventory, and \(0\) otherwise.
\item
  Replace all polymers meeting \(Y\) by that address and leave the
  remaining polymers unchanged. This is a bijection between old
  admissible sets meeting \(Y\) and admissible sets with a marked
  reserved address; its inverse expands the marked address. Paired
  summands agree. The forced-activity bounds give absolute convergence
  and the displayed weighted estimate, so the double sum can be
  reindexed.
\end{itemize}

\emph{Lean:}
\href{https://github.com/yang-mills-group/OpenAI-YM-20261028/blob/3456e5e75cf8d4cd558955a5c7638fba2d68f464/YM/Extra/SOURCE_HISTORIES/Purchase.lean\#L42-L78}{\texttt{Attach.paid}},
\href{https://github.com/yang-mills-group/OpenAI-YM-20261028/blob/3456e5e75cf8d4cd558955a5c7638fba2d68f464/YM/Extra/SOURCE_HISTORIES/Clusters.lean\#L198-L211}{\texttt{Attach.routes}},
\href{https://github.com/yang-mills-group/OpenAI-YM-20261028/blob/3456e5e75cf8d4cd558955a5c7638fba2d68f464/YM/Extra/SOURCE_HISTORIES/Stock.lean\#L120-L160}{\texttt{Attach.line\_budget}}.

\textbf{Lemma 4.26 (birth = Prop. II.1.8; summability; moment
identity).} \emph{Birth.} The supplement vanishes on
\(\mathrm{Hy}_\theta(Y)\), because \(\Phi_{e_0}=1\). It is gauge
invariant, measurable, determined by the links and flags of \(Y\), and
bounded by \(B\theta^{-f_{\max}}\lVert v^\circ\rVert_1\). Splice it by
Lemma 4.25. The local part of the density is then the masked healthy
part and the marked series is supplement × gas; their sum is
\(\rho\,\Phi_{S_{1/\theta}v}\). The forced budget is
\(\mathrm{fee}\cdot\lVert v^\circ\rVert_1\le\sqrt w\,\lVert v^\circ\rVert_1\).
\emph{Guarded birth.} Use the same reserved addresses and coefficients
after temporarily replacing the forced-polymer auxiliary tests by the
always-true test. The vacuum activities, supports and inventories stay
fixed, so the density and bounds are unchanged. Every nonzero inserted
coefficient enforces the correct inventory and therefore satisfies the
original auxiliary test. Restoring those tests gives the birth identity
on the original selected history.

\emph{Summability.} A remainder bound dominates the masked regular
series by a summable sequence. If the representation is lawful, combine
summable majorants for the activities and forced coefficients into
\(p_l\ge0\); the marked summand for a set \(b\) is bounded by
\(\prod_{l\in b}2p_l\), summable over finite sets. Format supplies
measurability of the coefficients. The lawful vacuum representation and
the finite section then give bounded measurable density. \emph{Moment
identity.} For a selected history, each transfer multiplies the total
integral by \(L_m^d\) (the push-forward preserves mass) and
\(a_m=L_ma_{m+1}\). This gives the identity
\(\int\mathcal D_m=a_m^dZ\langle O\rangle\) of Prop. II.1.7.
\emph{Lean:}
\href{https://github.com/yang-mills-group/OpenAI-YM-20261028/blob/3456e5e75cf8d4cd558955a5c7638fba2d68f464/YM/Extra/SOURCE_HISTORIES/Birth.lean\#L116-L161}{\texttt{launched}},
\href{https://github.com/yang-mills-group/OpenAI-YM-20261028/blob/3456e5e75cf8d4cd558955a5c7638fba2d68f464/YM/Extra/SOURCE_HISTORIES/Birth.lean\#L34-L39}{\texttt{healthy\_supplement}},
\href{https://github.com/yang-mills-group/OpenAI-YM-20261028/blob/3456e5e75cf8d4cd558955a5c7638fba2d68f464/YM/Extra/SOURCE_HISTORIES/Birth.lean\#L69-L101}{\texttt{supplement\_bound}},
\href{https://github.com/yang-mills-group/OpenAI-YM-20261028/blob/3456e5e75cf8d4cd558955a5c7638fba2d68f464/YM/Source/Preservation.lean\#L17-L23}{\texttt{regular\_series}},
\href{https://github.com/yang-mills-group/OpenAI-YM-20261028/blob/3456e5e75cf8d4cd558955a5c7638fba2d68f464/YM/Source/Preservation.lean\#L25-L51}{\texttt{marked\_series}},
\href{https://github.com/yang-mills-group/OpenAI-YM-20261028/blob/3456e5e75cf8d4cd558955a5c7638fba2d68f464/YM/Source/Preservation.lean\#L66-L84}{\texttt{envelope}},
\href{https://github.com/yang-mills-group/OpenAI-YM-20261028/blob/3456e5e75cf8d4cd558955a5c7638fba2d68f464/YM/Source/Preservation.lean\#L100-L118}{\texttt{integral\_marked}},
\href{https://github.com/yang-mills-group/OpenAI-YM-20261028/blob/3456e5e75cf8d4cd558955a5c7638fba2d68f464/YM/Extra/SOURCE_HISTORIES/Song.lean\#L48-L90}{\texttt{onTop}}.
\emph{(proofs read)}

\textbf{Lemma 4.27 (numerical conditions = Prop. II.1.9).} \emph{Latch.}
\(\theta_m^{-1}\le\hat\beta_m\). Once the cover size, which grows like
\((\log\hat\beta_m)^2\), exceeds the minimal side,
\(\mathrm{mass}(Y,1)\le e^{\text{support exponent}}\). Also
\(e^{\lvert Y\rvert w}-1\le\lvert Y\rvert we^{\lvert Y\rvert}\). Hence
\(\mathrm{fee}\le C\hat\beta_m^{f_{\max}}w_m\). Since
\(w_m=\exp(-\exp((\log\log\hat\beta_m)^2/5))\) decays faster than every
inverse power, \(\hat\beta^{n+1}w\le\sqrt w\) eventually. Thresholds for
all \(d\le J(m)\) are combined by summing finitely many thresholds.
\emph{Numerical conditions.}

\begin{itemize}
\tightlist
\item
  bank: \(L^{1/8}\to\infty\);
\item
  per-step prices: \(L^{-2/5}s\le1\) and \(L^{-2/5}(1+s)T\le L^{-3/8}\)
  for large \(L\);
\item
  block lower bounds: maximum of finitely many requirements per stage;
\item
  Enlarge each stage constant \(\bar\kappa_j\) to
  \(\max(\bar\kappa_j,\pi,64\mathsf K_{j+1}^2)\), with
  \(\mathsf K_k=1+\sum_{j\le k}\sum_{d\le j}(C_j+1)\,\Omega_d\).
\end{itemize}

\emph{Section clauses.} For \(m\ge r(d)\),
\(d+1\le d+3\le J(m)\le\mathrm{order}_m\) (and likewise for the power).
Latch gives the minimal side.
\(\pi(\max(\mathrm{order},\mathrm{power}))\le\kappa_N\) turns the bound
of Prop. II.1.3, item 4, into clause (v) of \(\mathrm{Makes}\).
\emph{Merger delay.} For each section and support-exponent ceiling \(u\)
there is a threshold \(q\) such that, whenever the level-\(m\) support
exponent is at most \(u\) and \(\hat\beta_m\ge q\),
\((2B_m+1)^4\sqrt{w_m}\le1\), where
\(B_m=\lceil(\log\hat\beta_m)^2\rceil\) is the cover size. The same
threshold ensures the bare-section bounds needed for local merging.
\emph{Lean:}
\href{https://github.com/yang-mills-group/OpenAI-YM-20261028/blob/3456e5e75cf8d4cd558955a5c7638fba2d68f464/YM/Extra/SOURCE_HISTORIES/Delays.lean\#L126-L157}{\texttt{latch\_cost}},
\href{https://github.com/yang-mills-group/OpenAI-YM-20261028/blob/3456e5e75cf8d4cd558955a5c7638fba2d68f464/YM/Extra/SOURCE_HISTORIES/Delays.lean\#L162-L177}{\texttt{entrances}},
\href{https://github.com/yang-mills-group/OpenAI-YM-20261028/blob/3456e5e75cf8d4cd558955a5c7638fba2d68f464/YM/Extra/SOURCE_HISTORIES/Reserve.lean\#L76-L102}{\texttt{afford}},
\href{https://github.com/yang-mills-group/OpenAI-YM-20261028/blob/3456e5e75cf8d4cd558955a5c7638fba2d68f464/YM/Extra/SOURCE_HISTORIES/Reserve.lean\#L108-L133}{\texttt{alone}},
\href{https://github.com/yang-mills-group/OpenAI-YM-20261028/blob/3456e5e75cf8d4cd558955a5c7638fba2d68f464/YM/Next/FOUNDATIONS/Primary.lean\#L182-L193}{\texttt{refLi}},
\href{https://github.com/yang-mills-group/OpenAI-YM-20261028/blob/3456e5e75cf8d4cd558955a5c7638fba2d68f464/YM/Extra/SOURCE_HISTORIES/Issue.lean\#L38-L45}{\texttt{Issue.refer}},
\href{https://github.com/yang-mills-group/OpenAI-YM-20261028/blob/3456e5e75cf8d4cd558955a5c7638fba2d68f464/YM/Push/LOCAL_SPLITTING/Delay.lean\#L56-L68}{\texttt{Push.Local.delay}}.

\textbf{Lemma 4.28 (local-merger identities and differentiation).}

\begin{enumerate}
\def\labelenumi{(\alph{enumi})}
\item
  Fix a valid forced-polymer family, a budget point, and a family
  \(f_i\) of measurable gauge-invariant functions. Each \(f_i\) is
  determined by links and flags on \(X_i\) and vanishes on
  \(\mathrm{Hy}_\theta(X_i)\). For every \(i\) and nonempty compatible
  bundle \(b\) of polymers meeting \(X_i\), assume a distinct reserved
  address with support \(X_i\cup\bigcup_{j\in b}X_j\), inventory
  \(\bigcup_{j\in b}I_j\), and auxiliary test satisfied everywhere. Pack
  \(f_i\prod_{j\in b}K_j\) into that address, multiplied by its
  eligibility indicator. If \(\sum_i|f_i(U)|<\infty\), the marked-gas
  series converges at \(U\) and equals
  \((\sum_i f_i(U))\mathrm{gas}(U)\). For a fixed admissible set, the
  marked summand is the sum over its possible marked addresses, with all
  other activities multiplied unchanged.
\item
  Suppose \(f\) is \(C^k\) in exponential coordinates at \(U\). Take a
  list of \(r\le k\) pairs \((A_j,B_j)\) of left and right directions.
  Assume that at each link either all \(A_j\) vanish or all \(B_j\)
  vanish. The ordered mixed derivative is the sum of \(2^r\) ordered
  left derivatives: for each pair choose its left direction or the
  direction \(\mathrm{Ad}(U)B_j\); put the chosen left directions first
  in their original order, followed by the chosen right directions in
  reverse order.
\item
  Let \(W\) be an open set of configurations. Suppose every \(f_i\) is
  \(C^k\) in exponential coordinates at every \(U\in W\). For every
  fixed ordered list \(l\) of at most \(k\) left directions, including
  the empty list, assume a summable majorant \(d_i(l)\) for
  \(|D_lf_i(U)|\), uniform in \(U\in W\). Then \(\sum_i f_i\) is \(C^k\)
  in those coordinates throughout \(W\), and
  \(D_l\sum_i f_i=\sum_iD_lf_i\) there.
\end{enumerate}

\emph{Sketch:} in (a), replace the polymers meeting \(X_i\) by their
reserved marked address. Expanding that address is the inverse
correspondence, and the activity bound and summability of \(f_i(U)\)
justify reindexing. In (b), move the right derivatives past the left
derivatives on disjoint link sets; replacing right increments by left
increments conjugates their directions by \(U\) and reverses their
order. For (c), work in one exponential chart contained in \(W\). Bounds
on finitely many basis derivatives give a summable bound on the Fréchet
derivative; the chart change is uniformly bounded on a smaller
neighbourhood. Differentiate the series there and induct on \(k\),
retaining the order of each left-derivative list.

\emph{Lean:}
\href{https://github.com/yang-mills-group/OpenAI-YM-20261028/blob/3456e5e75cf8d4cd558955a5c7638fba2d68f464/YM/Push/LOCAL_SPLITTING/Recovery.lean\#L29-L31}{\texttt{gasMark\_eq}},
\href{https://github.com/yang-mills-group/OpenAI-YM-20261028/blob/3456e5e75cf8d4cd558955a5c7638fba2d68f464/YM/Push/LOCAL_SPLITTING/Recovery.lean\#L184-L191}{\texttt{Slots.recover}},
\href{https://github.com/yang-mills-group/OpenAI-YM-20261028/blob/3456e5e75cf8d4cd558955a5c7638fba2d68f464/YM/Push/LOCAL_SPLITTING/Mixed.lean\#L179-L201}{\texttt{mix\_convert}},
\href{https://github.com/yang-mills-group/OpenAI-YM-20261028/blob/3456e5e75cf8d4cd558955a5c7638fba2d68f464/YM/Further/LOCAL_LOG_JETS/Series.lean\#L232-L255}{\texttt{smooth\_exchange}}.
These identities enter the conversion of local functions into source
coefficients and the differentiation of unmasked series in the per-step
estimates.

\paragraph{4.7 Properties of the
observables}\label{properties-of-the-observables}

\textbf{Lemma 4.29 (= Prop. II.1.6).} \emph{Sketch:}

\begin{itemize}
\tightlist
\item
  Translation is literal, since the placement commutes with
  translations.
\item
  Linearity, continuity and locality are properties of \(\Phi\) and of
  the linear maps \(S_{1/\theta_N}\) and \(\tau^{-1}\).
\item
  Gauge invariance follows from soundness of the section.
\item
  Reality and covariance use three facts: \(\tau\) commutes with
  conjugation and with the signed-permutation action (Prop. II.1.7(6));
  \(S_t\) commutes with the \(O(4)\) action because the grades are
  preserved; and the section is sound.
\item
  Unit: Prop. II.1.7(6) gives \(\tau(a_N^de_0)=x_r(a_N^de_0)=a_r^de_0\),
  hence \(\tau^{-1}(a_r^de_0)=a_N^de_0\). Since \(S_t e_0=e_0\) and
  \(\Phi_{e_0}=1\), the observable is \(a_N^{-d}a_N^d=1\).
\item
  Normalization is \(x_r=\tau\).
\end{itemize}

\emph{Lean:} see Prop. II.1.6 and
\href{https://github.com/yang-mills-group/OpenAI-YM-20261028/blob/3456e5e75cf8d4cd558955a5c7638fba2d68f464/YM/Source/Insertion.lean\#L121-L130}{\texttt{origin\_unit}}.

\paragraph{4.8 Comparison estimates}\label{comparison-estimates}

\textbf{Lemma 4.30 (= Prop. II.1.10).} \emph{Sketch:}

\begin{itemize}
\tightlist
\item
  \begin{enumerate}
  \def\labelenumi{(\alph{enumi})}
  \tightlist
  \item
    Translations and diagonal axis reflections carry the centred tiling
    and all its cell copies to the corresponding centred tiling at the
    new anchor. Reindex the cells and use invariance of the Wilson
    measure. Conjugating every copy conjugates the expectation, leaving
    its absolute value unchanged.
  \end{enumerate}
\item
  (b) At offsets \(o\) that are centres of the finer centred grid, apply
  the chessboard inequality for \(T\). Each cell copy of \(F_o\) is a
  reflected and translated copy of \(F_0\) planted in a distinct finer
  cell, so the product of roots is \(\le A^{\sharp \text{cells}}\) by
  (a).
\item
  (b), general offsets: bracket each coordinate between aligned corners,
  retaining the \(2\times\)-side margin for all intermediate positions.
  For literal translates \(F_u\) along an axis the seam inequality gives
  \(r_{u+1}^2\le r_ur_{u+2}\), where \(r_u\) is the root test. A
  nonnegative sequence satisfying this inequality cannot exceed both
  endpoint bounds. Iterate over the four coordinates.
\item
  (c) The root test is subadditive under finite sums. Slowly varying
  coefficients cost \(e^{\varepsilon s_N}\) with
  \(\varepsilon<\gamma/2\), which leaves the rate \(e^{-\gamma s_N/2}\).
\end{itemize}

\emph{Lean:}
\href{https://github.com/yang-mills-group/OpenAI-YM-20261028/blob/3456e5e75cf8d4cd558955a5c7638fba2d68f464/YM/Push/RELAY_DIFFERENCE/Placement.lean\#L134-L160}{\texttt{center\_price}},
\href{https://github.com/yang-mills-group/OpenAI-YM-20261028/blob/3456e5e75cf8d4cd558955a5c7638fba2d68f464/YM/Push/RELAY_DIFFERENCE/Placement.lean\#L192-L226}{\texttt{corner\_root}},
\href{https://github.com/yang-mills-group/OpenAI-YM-20261028/blob/3456e5e75cf8d4cd558955a5c7638fba2d68f464/YM/Push/RELAY_DIFFERENCE/Placement.lean\#L232-L270}{\texttt{offset}},
\href{https://github.com/yang-mills-group/OpenAI-YM-20261028/blob/3456e5e75cf8d4cd558955a5c7638fba2d68f464/YM/Push/RELAY_DIFFERENCE/Combination.lean\#L65-L93}{\texttt{gather\_close}}.

\subsubsection{5. Where the difficulty
lies}\label{where-the-difficulty-lies-1}

The main analytic estimate is the volume-independent readout bound.
Polarization in the full lattice direction space could introduce
constants depending on the volume. Pulling each tensor back to the
finitely many test directions avoids this, while rescaling each link
position gives the required distance weights. The analytic radii of
individual links need not be uniform.

The equation lift must serve two purposes at once: reproduce a
prescribed jet and permit integration by parts without a derivative of
its multiplier. Slot sampling away from the varied link supplies this
independence. The triangular argument orders grades first by factor
number, then by weight; higher-factor corrections need not have higher
engineering weight. The final family retains its primary sections, and
the equation sections form an auxiliary family whose comparison requires
later estimates.

The birth identity requires reserved addresses in the same
forced-polymer alphabet. Lawfulness alone does not supply them. The
guarded construction weakens the auxiliary-test requirement while
retaining exact eligibility whenever an inserted coefficient is nonzero.
Its numerical bound is possible because the forced allowance decays
faster than every inverse coupling power.

The finite slot family is chosen nonconstructively. Consequently the
stencils and constants are finite but not explicit. Their dependence on
degree and derivative parameters must be included before the block
factors are chosen.

\subsubsection{6. Notation}\label{notation-1}

{\def\LTcaptype{none} % do not increment counter
\begin{longtable}[]{@{}
  >{\raggedright\arraybackslash}p{(\linewidth - 4\tabcolsep) * \real{0.3333}}
  >{\raggedright\arraybackslash}p{(\linewidth - 4\tabcolsep) * \real{0.3333}}
  >{\raggedright\arraybackslash}p{(\linewidth - 4\tabcolsep) * \real{0.3333}}@{}}
\toprule\noalign{}
\begin{minipage}[b]{\linewidth}\raggedright
Written proof
\end{minipage} & \begin{minipage}[b]{\linewidth}\raggedright
Lean (glossary)
\end{minipage} & \begin{minipage}[b]{\linewidth}\raggedright
Meaning
\end{minipage} \\
\midrule\noalign{}
\endhead
\bottomrule\noalign{}
\endlastfoot
\(\mathrm{Tay}_j\), jet \(J\) &
\href{https://github.com/yang-mills-group/OpenAI-YM-20261028/blob/3456e5e75cf8d4cd558955a5c7638fba2d68f464/YM/Labels/Jet.lean\#L39}{\texttt{Labels.Taylor}},
\href{https://github.com/yang-mills-group/OpenAI-YM-20261028/blob/3456e5e75cf8d4cd558955a5c7638fba2d68f464/YM/Labels/Polynomial.lean\#L28}{\texttt{Jet}}
& radial-gauge Taylor coefficient of weight \(j+2\); sequence of them \\
\(V_d\) &
\href{https://github.com/yang-mills-group/OpenAI-YM-20261028/blob/3456e5e75cf8d4cd558955a5c7638fba2d68f464/YM/Labels/Polynomial.lean\#L112}{\texttt{Labels.Label}}
& invariant polynomials of weight \(\le d\) (labels) \\
\(p(J)\) &
\href{https://github.com/yang-mills-group/OpenAI-YM-20261028/blob/3456e5e75cf8d4cd558955a5c7638fba2d68f464/YM/Labels/Polynomial.lean\#L35}{\texttt{read}}
& value of a label on a jet \\
\((w,f)\), \(V_d^{(w,f)}\) &
\href{https://github.com/yang-mills-group/OpenAI-YM-20261028/blob/3456e5e75cf8d4cd558955a5c7638fba2d68f464/YM/Labels/Grading.lean\#L17-L18}{\texttt{Grade}}
& weight and number of factors \\
frame \((b_i)\), \(c\in\mathbb C^{K_d+1}\) &
\href{https://github.com/yang-mills-group/OpenAI-YM-20261028/blob/3456e5e75cf8d4cd558955a5c7638fba2d68f464/YM/Labels/Basis.lean\#L50-L57}{\texttt{Frame}},
\href{https://github.com/yang-mills-group/OpenAI-YM-20261028/blob/3456e5e75cf8d4cd558955a5c7638fba2d68f464/YM/Labels/Reference.lean\#L18}{\texttt{Frame.Coeff}}
& real graded basis starting with \(1\); coordinates \\
\(S_t\), \(c^\circ\), \(\lVert c\rVert_1\) &
\href{https://github.com/yang-mills-group/OpenAI-YM-20261028/blob/3456e5e75cf8d4cd558955a5c7638fba2d68f464/YM/Labels/Reference.lean\#L58-L61}{\texttt{Frame.scaled}},
\href{https://github.com/yang-mills-group/OpenAI-YM-20261028/blob/3456e5e75cf8d4cd558955a5c7638fba2d68f464/YM/Labels/Reference.lean\#L123-L126}{\texttt{core}},
\href{https://github.com/yang-mills-group/OpenAI-YM-20261028/blob/3456e5e75cf8d4cd558955a5c7638fba2d68f464/YM/Labels/Reference.lean\#L128}{\texttt{coeffSize}}
& \(c_i\mapsto t^{f_i}c_i\); drop the unit coordinate; \(\ell^1\)
norm \\
\(s\cdot c\), \(\bar c\) &
\href{https://github.com/yang-mills-group/OpenAI-YM-20261028/blob/3456e5e75cf8d4cd558955a5c7638fba2d68f464/YM/Labels/Reference.lean\#L83-L84}{\texttt{Frame.act}},
\href{https://github.com/yang-mills-group/OpenAI-YM-20261028/blob/3456e5e75cf8d4cd558955a5c7638fba2d68f464/YM/Labels/Action.lean\#L131-L140}{\texttt{realBar}}
& \(O(4)\) action; conjugation \\
\(U^{(d)}_{J,h}\) &
\href{https://github.com/yang-mills-group/OpenAI-YM-20261028/blob/3456e5e75cf8d4cd558955a5c7638fba2d68f464/YM/Source/Local.lean\#L55-L57}{\texttt{Hol}},
\href{https://github.com/yang-mills-group/OpenAI-YM-20261028/blob/3456e5e75cf8d4cd558955a5c7638fba2d68f464/YM/Source/Local.lean\#L62}{\texttt{Hol.links}}
& test holonomies \\
\(\mathrm{sc}_hJ\) &
\href{https://github.com/yang-mills-group/OpenAI-YM-20261028/blob/3456e5e75cf8d4cd558955a5c7638fba2d68f464/YM/Source/Local.lean\#L82}{\texttt{scaleJet}}
& \((h^{j+2}J_j)_j\) \\
slot \(\sigma=(\omega,P)\) &
\href{https://github.com/yang-mills-group/OpenAI-YM-20261028/blob/3456e5e75cf8d4cd558955a5c7638fba2d68f464/YM/RG/PathActions.lean\#L38-L40}{\texttt{RG.Slot}}
& \(\mathrm{Ad}(U_\omega)F_P(U)\) \\
slot map \(\mathfrak S\), truncated jets &
\href{https://github.com/yang-mills-group/OpenAI-YM-20261028/blob/3456e5e75cf8d4cd558955a5c7638fba2d68f464/YM/Source/Drops.lean\#L98}{\texttt{bring}},
\href{https://github.com/yang-mills-group/OpenAI-YM-20261028/blob/3456e5e75cf8d4cd558955a5c7638fba2d68f464/YM/Source/Drops.lean\#L26}{\texttt{Snap}}
& finitely many slots as a function of \((J_j)_{j\le d-2}\) \\
section \((S,\Phi)\); \(\Phi_c(o;U)\) &
\href{https://github.com/yang-mills-group/OpenAI-YM-20261028/blob/3456e5e75cf8d4cd558955a5c7638fba2d68f464/YM/Source/Local.lean\#L103-L105}{\texttt{SectionData}},
\href{https://github.com/yang-mills-group/OpenAI-YM-20261028/blob/3456e5e75cf8d4cd558955a5c7638fba2d68f464/YM/Source/Local.lean\#L112}{\texttt{places}}
& stencil and linear map to entry polynomials; placement \\
sound &
\href{https://github.com/yang-mills-group/OpenAI-YM-20261028/blob/3456e5e75cf8d4cd558955a5c7638fba2d68f464/YM/Source/Local.lean\#L140-L147}{\texttt{Sound}}
& symmetries, unit, projection property \\
slot structure; cost &
\href{https://github.com/yang-mills-group/OpenAI-YM-20261028/blob/3456e5e75cf8d4cd558955a5c7638fba2d68f464/YM/Bid/Source.lean\#L55-L56}{\texttt{Bid.pop}};
\href{https://github.com/yang-mills-group/OpenAI-YM-20261028/blob/3456e5e75cf8d4cd558955a5c7638fba2d68f464/YM/Bid/Source.lean\#L63-L66}{\texttt{Bid.Credit}}
& columns are slot polynomials with \(\ge f_i\) factors; \(t^{f_i}\)
budget bound \\
\(\pi(D)\), \(c_d\) &
\href{https://github.com/yang-mills-group/OpenAI-YM-20261028/blob/3456e5e75cf8d4cd558955a5c7638fba2d68f464/YM/Next/FOUNDATIONS/Primary.lean\#L126}{\texttt{price}},
\href{https://github.com/yang-mills-group/OpenAI-YM-20261028/blob/3456e5e75cf8d4cd558955a5c7638fba2d68f464/YM/Extra/SOURCE_HISTORIES/Delays.lean\#L20}{\texttt{cameras}}
& uniform price; value bound of the columns \\
\(p\) reads \(f\) at \(o\) &
\href{https://github.com/yang-mills-group/OpenAI-YM-20261028/blob/3456e5e75cf8d4cd558955a5c7638fba2d68f464/YM/Next/FOUNDATIONS/JetReads.lean\#L186-L194}{\texttt{Reads}}
& readout (constant convention when wrapping) \\
\(\mathrm{Sh}_j\), \(\mathcal E_d\) &
\href{https://github.com/yang-mills-group/OpenAI-YM-20261028/blob/3456e5e75cf8d4cd558955a5c7638fba2d68f464/YM/Maxwell/Shell.lean\#L18-L42}{\texttt{Maxwell.onShell}},
\href{https://github.com/yang-mills-group/OpenAI-YM-20261028/blob/3456e5e75cf8d4cd558955a5c7638fba2d68f464/YM/Maxwell/Shell.lean\#L46-L47}{\texttt{gaugeJet}},
\href{https://github.com/yang-mills-group/OpenAI-YM-20261028/blob/3456e5e75cf8d4cd558955a5c7638fba2d68f464/YM/Maxwell/Equation.lean\#L72}{\texttt{equationIdeal}}
& free-Maxwell components; free-Maxwell ideal \\
\(\nabla_eS_W\), transported force &
\href{https://github.com/yang-mills-group/OpenAI-YM-20261028/blob/3456e5e75cf8d4cd558955a5c7638fba2d68f464/YM/Push/RELAY_MARK/Wilson.lean\#L132}{\texttt{field}},
\href{https://github.com/yang-mills-group/OpenAI-YM-20261028/blob/3456e5e75cf8d4cd558955a5c7638fba2d68f464/YM/Push/RELAY_MARK/Avoid.lean\#L149-L150}{\texttt{force}}
& link gradient of the Wilson energy; its transport to \(0\) \\
micro &
\href{https://github.com/yang-mills-group/OpenAI-YM-20261028/blob/3456e5e75cf8d4cd558955a5c7638fba2d68f464/YM/Push/RELAY_MARK/Micro.lean\#L33}{\texttt{Micro}}
& orthogonal to remote bounded spectators \\
micro lift; lifting map \(F\) &
\href{https://github.com/yang-mills-group/OpenAI-YM-20261028/blob/3456e5e75cf8d4cd558955a5c7638fba2d68f464/YM/Push/RELAY_MARK/Rank.lean\#L78-L83}{\texttt{Cone}};
\href{https://github.com/yang-mills-group/OpenAI-YM-20261028/blob/3456e5e75cf8d4cd558955a5c7638fba2d68f464/YM/Push/RELAY_MARK/Linear.lean\#L69-L71}{\texttt{Campaign}}
& lift agreeing modulo higher grade; all grades at once \\
weave \(R\); equation section &
\href{https://github.com/yang-mills-group/OpenAI-YM-20261028/blob/3456e5e75cf8d4cd558955a5c7638fba2d68f464/YM/Push/RELAY_MARK/Automorphism.lean\#L72-L73}{\texttt{twist}};
\href{https://github.com/yang-mills-group/OpenAI-YM-20261028/blob/3456e5e75cf8d4cd558955a5c7638fba2d68f464/YM/Push/RELAY_MARK/Portfolio.lean\#L51-L58}{\texttt{Portfolio}}
& unitriangular regrading; section with micro equation part \\
record \(\mathsf C\); \(\mathrm{Makes}\) &
\href{https://github.com/yang-mills-group/OpenAI-YM-20261028/blob/3456e5e75cf8d4cd558955a5c7638fba2d68f464/YM/Source/Construction.lean\#L24-L30}{\texttt{Family}};
\href{https://github.com/yang-mills-group/OpenAI-YM-20261028/blob/3456e5e75cf8d4cd558955a5c7638fba2d68f464/YM/Source/Construction.lean\#L44-L62}{\texttt{Family.Makes}}
& Wilson family; consistency predicate \\
\(O_{N,p}(v,o)\) &
\href{https://github.com/yang-mills-group/OpenAI-YM-20261028/blob/3456e5e75cf8d4cd558955a5c7638fba2d68f464/YM/Source/Construction.lean\#L42}{\texttt{Family.O}}
& renormalized observable
\(a_N^{-d}\Phi_{S_{1/\theta_N}\tau^{-1}v}(o;\cdot)\) \\
\(r(d)\), \(a_r^de_0\) &
\href{https://github.com/yang-mills-group/OpenAI-YM-20261028/blob/3456e5e75cf8d4cd558955a5c7638fba2d68f464/YM/Source/Construction.lean\#L24-L30}{\texttt{Family.ref}},
\href{https://github.com/yang-mills-group/OpenAI-YM-20261028/blob/3456e5e75cf8d4cd558955a5c7638fba2d68f464/YM/Source/Construction.lean\#L37}{\texttt{Family.refUnit}}
& reference level; physical unit vector \\
\(E_m(v)\), \(\mathcal D_m\) &
\href{https://github.com/yang-mills-group/OpenAI-YM-20261028/blob/3456e5e75cf8d4cd558955a5c7638fba2d68f464/YM/Source/Insertion.lean\#L26-L30}{\texttt{Coefficient}},
\href{https://github.com/yang-mills-group/OpenAI-YM-20261028/blob/3456e5e75cf8d4cd558955a5c7638fba2d68f464/YM/Source/Insertion.lean\#L48-L50}{\texttt{density}}
& source coefficient; first-order inserted density \\
\(\tau\); based &
\href{https://github.com/yang-mills-group/OpenAI-YM-20261028/blob/3456e5e75cf8d4cd558955a5c7638fba2d68f464/YM/Source/Insertion.lean\#L73-L75}{\texttt{TransferData.terminal}};
\href{https://github.com/yang-mills-group/OpenAI-YM-20261028/blob/3456e5e75cf8d4cd558955a5c7638fba2d68f464/YM/Source/Insertion.lean\#L89-L114}{\texttt{Based}}
& terminal map; transfer specification \\
supplement; fee &
\href{https://github.com/yang-mills-group/OpenAI-YM-20261028/blob/3456e5e75cf8d4cd558955a5c7638fba2d68f464/YM/Extra/SOURCE_HISTORIES/Birth.lean\#L24-L25}{\texttt{supplement}};
\href{https://github.com/yang-mills-group/OpenAI-YM-20261028/blob/3456e5e75cf8d4cd558955a5c7638fba2d68f464/YM/Extra/SOURCE_HISTORIES/Birth.lean\#L108-L109}{\texttt{fee}}
& large-field part of the bare observable; birth condition \\
latch; Budget; Plans &
\href{https://github.com/yang-mills-group/OpenAI-YM-20261028/blob/3456e5e75cf8d4cd558955a5c7638fba2d68f464/YM/Extra/SOURCE_HISTORIES/Delays.lean\#L120-L123}{\texttt{Latch}};
\href{https://github.com/yang-mills-group/OpenAI-YM-20261028/blob/3456e5e75cf8d4cd558955a5c7638fba2d68f464/YM/Extra/SOURCE_HISTORIES/Reserve.lean\#L65-L72}{\texttt{Budget}};
\href{https://github.com/yang-mills-group/OpenAI-YM-20261028/blob/3456e5e75cf8d4cd558955a5c7638fba2d68f464/YM/Extra/SOURCE_HISTORIES/Reserve.lean\#L57-L62}{\texttt{Plans}}
& birth condition per level; level inequalities; stage-wise numerical
plan \\
price list \((\mathsf K_j)\) &
\href{https://github.com/yang-mills-group/OpenAI-YM-20261028/blob/3456e5e75cf8d4cd558955a5c7638fba2d68f464/YM/Bid/Loan.lean\#L29-L32}{\texttt{Bid.Worth}}
& nondecreasing constants \(\ge1\) \\
close &
\href{https://github.com/yang-mills-group/OpenAI-YM-20261028/blob/3456e5e75cf8d4cd558955a5c7638fba2d68f464/YM/Bid/Drops.lean\#L25}{\texttt{Bid.close}}
& root-test bounds tend to \(0\) in all prescribed link tilings and
margins \\
\end{longtable}
}

\subsubsection{7. What remains
unexplained}\label{what-remains-unexplained-1}

\begin{enumerate}
\def\labelenumi{\arabic{enumi}.}
\tightlist
\item
  \textbf{Identification in the limit (tentative).} The auxiliary
  equation sections enter the volume-uniform comparison through
  \href{https://github.com/yang-mills-group/OpenAI-YM-20261028/blob/3456e5e75cf8d4cd558955a5c7638fba2d68f464/YM/FullVolume/Cross.lean\#L65-L71}{\texttt{Run.Ties}}
  and
  \href{https://github.com/yang-mills-group/OpenAI-YM-20261028/blob/3456e5e75cf8d4cd558955a5c7638fba2d68f464/YM/FullVolume/Fire.lean\#L17-L20}{\texttt{Hearthside}}.
  This unit does not trace their regrading and transfer normalizations
  into the null-label quotient, or prove the converse identification
  with the free-Maxwell quotient. A literal equality
  \(N_d=\mathcal E_d\) in the original primary coordinates has not been
  verified.
\item
  \textbf{Later applications.} The section constants and reference-level
  conditions are explicit here. Their simultaneous compatibility with
  every later per-step estimate, and the full use of the local-merger
  identities, belong to III.2 and III.10--III.15 and are not verified by
  these local results.
\end{enumerate}

\subsection{II.2. Marker algebra, polymer combinatorics and Gaussian
estimates}\label{ii.2.-marker-algebra-polymer-combinatorics-and-gaussian-estimates}

\textbf{Part:} II, Data of the construction. \textbf{Lean:}
\href{https://github.com/yang-mills-group/OpenAI-YM-20261028/tree/3456e5e75cf8d4cd558955a5c7638fba2d68f464/YM/Marker}{\texttt{YM/Marker/}},
\href{https://github.com/yang-mills-group/OpenAI-YM-20261028/tree/3456e5e75cf8d4cd558955a5c7638fba2d68f464/YM/Ac}{\texttt{YM/Ac/}},
\href{https://github.com/yang-mills-group/OpenAI-YM-20261028/blob/3456e5e75cf8d4cd558955a5c7638fba2d68f464/YM/Ovens/Coffee.lean}{\texttt{YM/Ovens/Coffee.lean}},
\href{https://github.com/yang-mills-group/OpenAI-YM-20261028/blob/3456e5e75cf8d4cd558955a5c7638fba2d68f464/YM/Further/SHELL_BIRTH_49/Fusing/Coal.lean}{\texttt{YM/Further/SHELL\_BIRTH\_49/Fusing/Coal.lean}},
\href{https://github.com/yang-mills-group/OpenAI-YM-20261028/tree/3456e5e75cf8d4cd558955a5c7638fba2d68f464/YM/Extra/MULTISOURCE_TRAJECTORIES}{\texttt{YM/Extra/MULTISOURCE\_TRAJECTORIES/}},
\href{https://github.com/yang-mills-group/OpenAI-YM-20261028/blob/3456e5e75cf8d4cd558955a5c7638fba2d68f464/YM/Push/GRID_APPROXIMATION/Nil.lean}{\texttt{YM/Push/GRID\_APPROXIMATION/Nil.lean}},
\href{https://github.com/yang-mills-group/OpenAI-YM-20261028/tree/3456e5e75cf8d4cd558955a5c7638fba2d68f464/YM/MergeProof}{\texttt{YM/MergeProof/}},
\href{https://github.com/yang-mills-group/OpenAI-YM-20261028/tree/3456e5e75cf8d4cd558955a5c7638fba2d68f464/YM/Next/RANK}{\texttt{YM/Next/RANK/}},
\href{https://github.com/yang-mills-group/OpenAI-YM-20261028/blob/3456e5e75cf8d4cd558955a5c7638fba2d68f464/YM/Vac/Treasury.lean}{\texttt{YM/Vac/Treasury.lean}}.
\textbf{Size:} 3,691 lines in 793 declarations used by the main theorem.

This unit supplies the algebraic, combinatorial and numerical estimates
used to multiply and sum polymer activities in Part III, especially in
the Gaussian comparison of III.11. The activity format is defined in
III.1; the Wilson family and construction record in II.1;
reflection-positive cells in IV.1.

The results fall into five groups: marker algebra (II.2.1--II.2.4),
polymer combinatorics and forced-activity fusion (II.2.5--II.2.8),
switching (II.2.9), Gaussian polynomial calculus (II.2.10--II.2.11), and
numerical scale comparisons (II.2.12--II.2.13). They are used in
III.8--III.13, III.15, VII.2, VII.4 and VII.6; the vacuum-energy
estimate also enters V.2.

\subsubsection{1. What this unit proves}\label{what-this-unit-proves-2}

\textbf{Conventions.}

\begin{itemize}
\tightlist
\item
  \(\Lambda\) is a four-dimensional periodic lattice.
\item
  A \emph{support} is a finite subset of \(\Lambda\), and
  \(\mathrm{dist}\) is the lattice separation.
\item
  For \(M\in\mathbb N\) and a support \(X\), the \textbf{cover number}
  \(\mathrm{cov}_M(X)\) is the least cardinality of a set
  \(D\subseteq\Lambda\) such that the union \(B_M(D)\) of the balls of
  radius \(M\) around \(D\) contains \(X\) and is connected for
  nearest-neighbour steps
  (\href{https://github.com/yang-mills-group/OpenAI-YM-20261028/blob/3456e5e75cf8d4cd558955a5c7638fba2d68f464/YM/Lattice/Support.lean\#L69-L71}{\texttt{coverNumber}},
  \href{https://github.com/yang-mills-group/OpenAI-YM-20261028/blob/3456e5e75cf8d4cd558955a5c7638fba2d68f464/YM/Lattice/Support.lean\#L52-L55}{\texttt{covers}},
  I.1).
\item
  The parameters of the polymer norms are \(P=(r,q,e_P,B,t_P)\), with
  \(r,q,B\in\mathbb N\), \(r\ge6\), \(q\ge2\), \(B\ge1\), \(e_P>0\) and
  \(0<t_P\le1\): smoothness order, power, support exponent, cover size
  and threshold. The mass of an envelope \(c\) on support \(X\) is
  \(e^{e_P\mathrm{cov}_B(X)}c\).
\item
  Two facts from I.1 are used:

  \begin{itemize}
  \tightlist
  \item
    \(|X|\le(2M+1)^4\,\mathrm{cov}_M(X)\)
    (\href{https://github.com/yang-mills-group/OpenAI-YM-20261028/blob/3456e5e75cf8d4cd558955a5c7638fba2d68f464/YM/Corona/Fill.lean\#L129-L131}{\texttt{Corona.pays}});
  \item
    \(\mathrm{cov}_M(X\cup Y)\le\mathrm{cov}_M(X)+\mathrm{cov}_M(Y)\) if
    \(X\cap Y\neq\emptyset\)
    (\href{https://github.com/yang-mills-group/OpenAI-YM-20261028/blob/3456e5e75cf8d4cd558955a5c7638fba2d68f464/YM/Corona/Halo.lean\#L14-L30}{\texttt{Corona.meets\_sum}}).
  \end{itemize}
\end{itemize}

\paragraph{A. Marker algebra}\label{a.-marker-algebra}

\textbf{Proposition II.2.1 (marker algebra).} Let \(V\) be a finite set
and \(R\) a \emph{census} on \(V\): a finite family of finite subsets of
\(V\) that contains \(\emptyset\) and all singletons and is closed under
taking subsets. Let \(\mathcal A_R\) be the space of functions
\(a:R\to\mathbb C\), written \(a=\sum_{x\in R}a(x)\delta_x\), with
product
\[(ab)(z)=\sum_{x,y\in R:\ x\cap y=\emptyset,\ x\cup y=z}a(x)\,b(y)\qquad(z\in R).\]
Then:

\begin{enumerate}
\def\labelenumi{\arabic{enumi}.}
\tightlist
\item
  \(\mathcal A_R\) is a commutative unital \(\mathbb C\)-algebra with
  unit \(\delta_\emptyset\).
\item
  \(\delta_x\delta_y=\delta_{x\cup y}\) if \(x\cap y=\emptyset\) and
  \(x\cup y\in R\), and \(\delta_x\delta_y=0\) otherwise. So every
  marker \(\delta_{\{v\}}\) squares to zero, and products that leave
  \(R\) vanish.
\item
  \((ab)(\emptyset)=a(\emptyset)\,b(\emptyset)\).
\item
  If \(a(\emptyset)=0\), then \(a^k(x)=0\) whenever \(|x|<k\); in
  particular, \(a^k=0\) for \(k>|V|\).
\end{enumerate}

\emph{Status:} theorem. \emph{Lean:}
\href{https://github.com/yang-mills-group/OpenAI-YM-20261028/blob/3456e5e75cf8d4cd558955a5c7638fba2d68f464/YM/Marker/Console.lean\#L22-L26}{\texttt{Marker.Census}},
\href{https://github.com/yang-mills-group/OpenAI-YM-20261028/blob/3456e5e75cf8d4cd558955a5c7638fba2d68f464/YM/Marker/Console.lean\#L40}{\texttt{Poly}},
\href{https://github.com/yang-mills-group/OpenAI-YM-20261028/blob/3456e5e75cf8d4cd558955a5c7638fba2d68f464/YM/Marker/Console.lean\#L68}{\texttt{Poly.prod}},
\href{https://github.com/yang-mills-group/OpenAI-YM-20261028/blob/3456e5e75cf8d4cd558955a5c7638fba2d68f464/YM/Marker/Console.lean}{ring
structure},
\href{https://github.com/yang-mills-group/OpenAI-YM-20261028/blob/3456e5e75cf8d4cd558955a5c7638fba2d68f464/YM/Marker/Weight.lean\#L28}{\texttt{Poly.quiet\_mul}},
\href{https://github.com/yang-mills-group/OpenAI-YM-20261028/blob/3456e5e75cf8d4cd558955a5c7638fba2d68f464/YM/Marker/Truncated.lean\#L10-L22}{\texttt{Poly.high}},
\href{https://github.com/yang-mills-group/OpenAI-YM-20261028/blob/3456e5e75cf8d4cd558955a5c7638fba2d68f464/YM/Marker/Truncated.lean\#L26-L27}{\texttt{Poly.stops}}.

\textbf{Proposition II.2.2 (weighted \(\ell^1\) norm).} For
\(w:V\to(0,\infty)\) put \(w^x=\prod_{v\in x}w_v\) and
\(\|a\|_w=\sum_{x\in R}w^x|a(x)|\). Then \(\|1\|_w=1\),
\(\|c\,\delta_x\|_w=w^x|c|\), and
\[\|ab\|_w\le\|a\|_w\|b\|_w,\qquad \Big\|\prod_{i\in T}a_i\Big\|_w\le\prod_{i\in T}\|a_i\|_w,\qquad \|a^k\|_w\le\|a\|_w^k .\]
Set \(\mathrm{Exp}(a)=\sum_{k=0}^{|V|}a^k/k!\).

\begin{itemize}
\tightlist
\item
  \(\|\mathrm{Exp}(a)\|_w\le e^{\|a\|_w}\) for every \(a\).
\item
  If a countable family satisfies \(\|f_i\|_w\le b_i\) with
  \(\sum_ib_i<\infty\), then \(\sum_if_i\) converges in
  \(\mathcal A_R\), \(\|\sum_if_i\|_w\le\sum_ib_i\), and coefficients of
  the sum are sums of coefficients.
\item
  \(\max_x|a(x)|\le(\min_{x\in R}w^x)^{-1}\|a\|_w\).
\end{itemize}

The multiplicative and exponential bounds have constant \(1\),
independently of \(R\) and \(|V|\). The norm-comparison constant is used
to establish convergence.

\emph{Status:} theorem. \emph{Lean:}
\href{https://github.com/yang-mills-group/OpenAI-YM-20261028/blob/3456e5e75cf8d4cd558955a5c7638fba2d68f464/YM/Marker/Weight.lean\#L74}{\texttt{Poly.toll}},
\href{https://github.com/yang-mills-group/OpenAI-YM-20261028/blob/3456e5e75cf8d4cd558955a5c7638fba2d68f464/YM/Marker/Weight.lean\#L76}{\texttt{Poly.price}},
\href{https://github.com/yang-mills-group/OpenAI-YM-20261028/blob/3456e5e75cf8d4cd558955a5c7638fba2d68f464/YM/Marker/Weight.lean\#L139-L146}{\texttt{multiplication}},
\href{https://github.com/yang-mills-group/OpenAI-YM-20261028/blob/3456e5e75cf8d4cd558955a5c7638fba2d68f464/YM/Marker/Weight.lean\#L148-L152}{\texttt{multiple}},
\href{https://github.com/yang-mills-group/OpenAI-YM-20261028/blob/3456e5e75cf8d4cd558955a5c7638fba2d68f464/YM/Marker/Weight.lean\#L154-L157}{\texttt{powers}},
\href{https://github.com/yang-mills-group/OpenAI-YM-20261028/blob/3456e5e75cf8d4cd558955a5c7638fba2d68f464/YM/Marker/Truncated.lean\#L29}{\texttt{Poly.pass}},
\href{https://github.com/yang-mills-group/OpenAI-YM-20261028/blob/3456e5e75cf8d4cd558955a5c7638fba2d68f464/YM/Marker/Truncated.lean\#L39-L46}{\texttt{pass\_budget}},
\href{https://github.com/yang-mills-group/OpenAI-YM-20261028/blob/3456e5e75cf8d4cd558955a5c7638fba2d68f464/YM/Marker/Truncated.lean\#L51-L58}{\texttt{budgeted}},
\href{https://github.com/yang-mills-group/OpenAI-YM-20261028/blob/3456e5e75cf8d4cd558955a5c7638fba2d68f464/YM/Marker/Truncated.lean\#L60}{\texttt{share}},
\href{https://github.com/yang-mills-group/OpenAI-YM-20261028/blob/3456e5e75cf8d4cd558955a5c7638fba2d68f464/YM/Marker/Weight.lean\#L102-L112}{\texttt{equivalent\_norm}}.

\textbf{Proposition II.2.3 (exponential and connected polynomial).} Use
the marker algebra and exponential of II.2.1--II.2.2.

\begin{enumerate}
\def\labelenumi{(\alph{enumi})}
\item
  If \(a(\emptyset)=b(\emptyset)=0\), then
  \(\mathrm{Exp}(a+b)=\mathrm{Exp}(a)\,\mathrm{Exp}(b)\) and
  \(\mathrm{Exp}(-a)\,\mathrm{Exp}(a)=1\). More generally,
  \(\mathrm{Exp}(\sum_{i\in s}a_i)=\prod_{i\in s}\mathrm{Exp}(a_i)\) for
  a finite family with all \(a_i(\emptyset)=0\). The zero-constant-term
  hypotheses are necessary for these identities in general.
\item
  If \(x\neq\emptyset\), then
  \(\mathrm{Exp}(c\,\delta_x)=1+c\,\delta_x\). Hence, for a finite
  family of nonempty \(x_i\in R\) (not necessarily distinct),
  \[\mathrm{Exp}\Big(\sum_{i\in s}c_i\delta_{x_i}\Big)=\prod_{i\in s}(1+c_i\delta_{x_i}).\]
\item
  If \(F(\emptyset)=1\), there is exactly one \(C\in\mathcal A_R\) with
  \(C(\emptyset)=0\) and \(\mathrm{Exp}(C)=F\): the \textbf{connected
  polynomial} of \(F\). It is obtained by iteration:
\end{enumerate}

\begin{itemize}
\tightlist
\item
  \(C_0=0\) and \(C_{k+1}=C_k+F-\mathrm{Exp}(C_k)\);
\item
  \(F-\mathrm{Exp}(C_k)\) vanishes on all marker sets with at most \(k\)
  elements, so \(C=C_{|V|}\);
\item
  each \(C_k\) is a polynomial, hence continuous, function of \(F\).
\end{itemize}

\begin{enumerate}
\def\labelenumi{(\alph{enumi})}
\setcounter{enumi}{3}
\tightlist
\item
  If \(C(\emptyset)=0\) and \(\mathrm{Exp}(C)=F\), then
  \(F\,\mathrm{Exp}(-C)\,P=P\) for every \(P\).
\end{enumerate}

\emph{Status:} theorem. \emph{Lean:}
\href{https://github.com/yang-mills-group/OpenAI-YM-20261028/blob/3456e5e75cf8d4cd558955a5c7638fba2d68f464/YM/Next/RANK/Exponential.lean\#L56-L88}{\texttt{Exponential.pass\_add}},
\href{https://github.com/yang-mills-group/OpenAI-YM-20261028/blob/3456e5e75cf8d4cd558955a5c7638fba2d68f464/YM/Next/RANK/Exponential.lean\#L90}{\texttt{pass\_neg}},
\href{https://github.com/yang-mills-group/OpenAI-YM-20261028/blob/3456e5e75cf8d4cd558955a5c7638fba2d68f464/YM/Next/RANK/Exponential.lean\#L92-L100}{\texttt{pass\_sum}},
\href{https://github.com/yang-mills-group/OpenAI-YM-20261028/blob/3456e5e75cf8d4cd558955a5c7638fba2d68f464/YM/Extra/MULTISOURCE_TRAJECTORIES/Algebra.lean\#L113-L122}{\texttt{Trajectories.pass\_send}},
\href{https://github.com/yang-mills-group/OpenAI-YM-20261028/blob/3456e5e75cf8d4cd558955a5c7638fba2d68f464/YM/Extra/MULTISOURCE_TRAJECTORIES/Algebra.lean\#L124-L134}{\texttt{Trajectories.pass\_sum}},
\href{https://github.com/yang-mills-group/OpenAI-YM-20261028/blob/3456e5e75cf8d4cd558955a5c7638fba2d68f464/YM/Next/RANK/Exponential.lean\#L126}{\texttt{Expands}},
\href{https://github.com/yang-mills-group/OpenAI-YM-20261028/blob/3456e5e75cf8d4cd558955a5c7638fba2d68f464/YM/Next/RANK/Exponential.lean\#L130-L132}{\texttt{roll}},
\href{https://github.com/yang-mills-group/OpenAI-YM-20261028/blob/3456e5e75cf8d4cd558955a5c7638fba2d68f464/YM/Next/RANK/Exponential.lean\#L133-L149}{\texttt{rolls}},
\href{https://github.com/yang-mills-group/OpenAI-YM-20261028/blob/3456e5e75cf8d4cd558955a5c7638fba2d68f464/YM/Next/RANK/Exponential.lean\#L151-L156}{\texttt{roll\_cont}},
\href{https://github.com/yang-mills-group/OpenAI-YM-20261028/blob/3456e5e75cf8d4cd558955a5c7638fba2d68f464/YM/Next/RANK/Exponential.lean\#L161-L174}{\texttt{Expands.uniq}},
\href{https://github.com/yang-mills-group/OpenAI-YM-20261028/blob/3456e5e75cf8d4cd558955a5c7638fba2d68f464/YM/Next/RANK/Exponential.lean\#L179-L180}{\texttt{connected\_spec}},
\href{https://github.com/yang-mills-group/OpenAI-YM-20261028/blob/3456e5e75cf8d4cd558955a5c7638fba2d68f464/YM/Next/RANK/Exponential.lean\#L184-L185}{\texttt{correct}}.

\textbf{Proposition II.2.4 (tagged products; repacking).} Use positive
weights \(w\) as in II.2.2.

\begin{enumerate}
\def\labelenumi{(\alph{enumi})}
\tightlist
\item
  Let \(\tau:\iota\to R\) and let \(s\subseteq\iota\) be finite. Say
  that \textbf{\(s\) fits} if the sets \(\tau(i)\), \(i\in s\), are
  pairwise disjoint and their union lies in \(R\).
\end{enumerate}

\begin{itemize}
\tightlist
\item
  If \(s\) fits, then
  \(\prod_{i\in s}c_i\delta_{\tau(i)}=(\prod_{i\in s}c_i)\,\delta_{\bigcup_{i\in s}\tau(i)}\)
  and \(w^{\bigcup_{i\in s}\tau(i)}=\prod_{i\in s}w^{\tau(i)}\).
\item
  Conversely, if \(\prod_{i\in s}c_i\delta_{\tau(i)}\neq0\), then \(s\)
  fits.
\end{itemize}

\begin{enumerate}
\def\labelenumi{(\alph{enumi})}
\setcounter{enumi}{1}
\tightlist
\item
  Let:
\end{enumerate}

\begin{itemize}
\tightlist
\item
  \(\sigma:\iota\to\kappa\) assign a \emph{slot} to each entry;
\item
  ``admissible'' be any property of finite sets of slots;
\item
  \(F:\iota\to\mathcal A_R\) satisfy \(\sum_i\|F_i\|_w<\infty\).
\end{itemize}

For finite \(s\subseteq\iota\) put \(\mathrm{val}(s)=\prod_{i\in s}F_i\)
if \(\sigma\) is injective on \(s\) and \(\sigma(s)\) is admissible, and
\(\mathrm{val}(s)=0\) otherwise. Then \(\mathrm{val}\) is summable over
finite subsets of \(\iota\), and
\[\sum_{s}\mathrm{val}(s)=\sum_{S\subseteq\kappa\text{ finite, admissible}}\ \prod_{k\in S}\ \sum_{i:\,\sigma(i)=k}F_i .\]

\emph{Status:} theorem. \emph{Lean:}
\href{https://github.com/yang-mills-group/OpenAI-YM-20261028/blob/3456e5e75cf8d4cd558955a5c7638fba2d68f464/YM/Extra/MULTISOURCE_TRAJECTORIES/Algebra.lean\#L44-L45}{\texttt{Trajectories.Fits}},
\href{https://github.com/yang-mills-group/OpenAI-YM-20261028/blob/3456e5e75cf8d4cd558955a5c7638fba2d68f464/YM/Extra/MULTISOURCE_TRAJECTORIES/Algebra.lean\#L57-L69}{\texttt{tagged\_prod}},
\href{https://github.com/yang-mills-group/OpenAI-YM-20261028/blob/3456e5e75cf8d4cd558955a5c7638fba2d68f464/YM/Extra/MULTISOURCE_TRAJECTORIES/Algebra.lean\#L71-L87}{\texttt{fits\_of\_tagged\_prod}},
\href{https://github.com/yang-mills-group/OpenAI-YM-20261028/blob/3456e5e75cf8d4cd558955a5c7638fba2d68f464/YM/Extra/MULTISOURCE_TRAJECTORIES/Algebra.lean\#L89-L90}{\texttt{toll\_union}},
\href{https://github.com/yang-mills-group/OpenAI-YM-20261028/blob/3456e5e75cf8d4cd558955a5c7638fba2d68f464/YM/Extra/MULTISOURCE_TRAJECTORIES/Packets.lean\#L30-L38}{\texttt{Packets.summable\_value}},
\href{https://github.com/yang-mills-group/OpenAI-YM-20261028/blob/3456e5e75cf8d4cd558955a5c7638fba2d68f464/YM/Extra/MULTISOURCE_TRAJECTORIES/Packets.lean\#L114-L124}{\texttt{Packets.repack}}.

\paragraph{B. Polymer combinatorics}\label{b.-polymer-combinatorics}

\textbf{Proposition II.2.5 (connected components).} Let
\(X:\iota\to\{\text{supports}\}\) and let \(s\subseteq\iota\) be finite.

\begin{itemize}
\tightlist
\item
  Two labels of \(s\) are \emph{chained in \(s\)} if they are joined by
  a sequence of labels of \(s\) whose successive supports intersect.
\item
  A \textbf{spark} is a nonempty \(s\) in which all labels are chained.
\item
  \(\mathrm{Parts}(s)\) is the set of chain classes of \(s\), and
  \(X(p)=\bigcup_{i\in p}X_i\).
\end{itemize}

Then:

\begin{enumerate}
\def\labelenumi{\arabic{enumi}.}
\tightlist
\item
  The parts of \(s\) are sparks, pairwise disjoint, with union \(s\) and
  with pairwise disjoint supports \(X(p)\).
\item
  If \(S\) is a finite set of sparks with pairwise disjoint supports,
  then \(\mathrm{Parts}(\bigcup S)=S\).
\item
  If \(s\) is a spark and \(i\in s\), every label of \(s\setminus\{i\}\)
  is chained in \(s\setminus\{i\}\) to a label whose support meets
  \(X_i\).
\item
  If every label of \(s\) is chained in \(s\) to a label meeting a
  support \(Y\) (call \(s\) a \textbf{pack for \(Y\)}), then every part
  of \(s\) contains a label meeting \(Y\).
\end{enumerate}

The same component statements hold when the supports are finite subsets
of any set.

\emph{Status:} theorem. \emph{Lean:}
\href{https://github.com/yang-mills-group/OpenAI-YM-20261028/blob/3456e5e75cf8d4cd558955a5c7638fba2d68f464/YM/Ac/Nodes.lean\#L49}{\texttt{Ac.Spark}},
\href{https://github.com/yang-mills-group/OpenAI-YM-20261028/blob/3456e5e75cf8d4cd558955a5c7638fba2d68f464/YM/Ac/Nodes.lean\#L55}{\texttt{Parts}},
\href{https://github.com/yang-mills-group/OpenAI-YM-20261028/blob/3456e5e75cf8d4cd558955a5c7638fba2d68f464/YM/Ac/Nodes.lean\#L60-L63}{\texttt{reconstruct}},
\href{https://github.com/yang-mills-group/OpenAI-YM-20261028/blob/3456e5e75cf8d4cd558955a5c7638fba2d68f464/YM/Ac/Nodes.lean\#L72-L77}{\texttt{separate}},
\href{https://github.com/yang-mills-group/OpenAI-YM-20261028/blob/3456e5e75cf8d4cd558955a5c7638fba2d68f464/YM/Ac/Nodes.lean\#L79-L101}{\texttt{separates}},
\href{https://github.com/yang-mills-group/OpenAI-YM-20261028/blob/3456e5e75cf8d4cd558955a5c7638fba2d68f464/YM/Ac/Nodes.lean\#L103-L113}{\texttt{excise}},
\href{https://github.com/yang-mills-group/OpenAI-YM-20261028/blob/3456e5e75cf8d4cd558955a5c7638fba2d68f464/YM/Ac/Nodes.lean\#L115-L118}{\texttt{packHits}},
\href{https://github.com/yang-mills-group/OpenAI-YM-20261028/blob/3456e5e75cf8d4cd558955a5c7638fba2d68f464/YM/Further/SHELL_BIRTH_49/Fusing/Coal.lean\#L91-L111}{\texttt{Collide.uncluttered}}.

\textbf{Proposition II.2.6 (tree bound in cover-number weights).} Let:

\begin{itemize}
\tightlist
\item
  \(\iota\) be any set of labels and \(X:\iota\to\{\text{supports}\}\);
\item
  \(M\in\mathbb N\), \(f:\iota\to[0,\infty)\), \(a\ge0\) and
  \(\gamma\ge0\) with \((2M+1)^4\gamma\le1\).
\end{itemize}

Assume the \textbf{hit condition with exponent \(a+1\)}: for every site
\(x\) the family
\(\big(\mathbf 1[x\in X_i]\,f_i\,e^{(a+1)\mathrm{cov}_M(X_i)}\big)_i\)
is summable, with sum at most \(\gamma\). Then for every site \(x\),
\[\sum_{s\ \text{spark},\ x\in X(s)}e^{a\,\mathrm{cov}_M(X(s))}\prod_{i\in s}f_i\ \le\ e^{\gamma}-1 .\]
The sum runs over finite sets of labels, so each label is used at most
once, and the family is summable.

\emph{Status:} theorem. Used by III.9--III.13, VII.2 and VII.6.
\emph{Lean:}
\href{https://github.com/yang-mills-group/OpenAI-YM-20261028/blob/3456e5e75cf8d4cd558955a5c7638fba2d68f464/YM/Ac/Tree.lean\#L12-L13}{\texttt{Ac.Hit}},
\href{https://github.com/yang-mills-group/OpenAI-YM-20261028/blob/3456e5e75cf8d4cd558955a5c7638fba2d68f464/YM/Ac/Tree.lean\#L84-L85}{\texttt{Ac.roots}},
\href{https://github.com/yang-mills-group/OpenAI-YM-20261028/blob/3456e5e75cf8d4cd558955a5c7638fba2d68f464/YM/Ac/Tree.lean\#L89-L127}{\texttt{Ac.forest}}.

\textbf{Proposition II.2.7 (unique recovery for two kinds of overlap).}
A \textbf{kiln} on a set \(I\) of labels consists of:

\begin{itemize}
\tightlist
\item
  finite sets \(\mathrm{men}(i)\subseteq V\) (``mentions'');
\item
  finite sets \(\mathrm{wall}(i)\subseteq W\) (\(W\) another set);
\item
  a set of \emph{dirty} labels.
\end{itemize}

For finite \(s\subseteq I\):

\begin{itemize}
\tightlist
\item
  a label \(i\in s\) is \textbf{hot in \(s\)} if it is joined to a dirty
  label of \(s\) by a chain in \(s\) of overlapping mentions;
\item
  the \textbf{fuel} of \(i\) in \(s\) is the subset of \(V\sqcup W\)
  consisting of its mentions in the \(V\) component and, only when \(i\)
  is hot, its walls in the \(W\) component;
\item
  \(s\) is a \textbf{mole} if it is a spark for fuel overlap;
\item
  the \textbf{slabs} of \(s\) are its parts for fuel overlap.
\end{itemize}

Then:

\begin{enumerate}
\def\labelenumi{\arabic{enumi}.}
\tightlist
\item
  Every slab \(p\) of \(s\) is a mole, and the fuel of its labels
  computed in \(p\) equals the fuel computed in \(s\). Distinct slabs
  have disjoint mention-unions. Distinct slabs that contain a dirty
  label have disjoint wall-unions. (In a mole, containing a dirty label
  is equivalent to all labels being hot.)
\item
  Conversely, let \(S\) be a finite set of moles with pairwise disjoint
  mention-unions, and with pairwise disjoint wall-unions among the
  members that contain a dirty label. Then the slabs of \(\bigcup S\)
  are exactly the members of \(S\).
\item
  By 1 and 2, \(s\mapsto\mathrm{slabs}(s)\) and \(S\mapsto\bigcup S\)
  are inverse bijections between finite sets of labels and such finite
  sets of moles.
\item
  Suppose the walls are supports and, inside a mole \(s\), overlapping
  mentions imply overlapping walls. Then \(s\) is a spark for the walls.
  Thus a family of moles can be estimated by II.2.6 whenever its wall
  supports and weights satisfy that proposition's hypotheses.
\end{enumerate}

\emph{Status:} theorem. Used by III.11 and VII.4. \emph{Lean:}
\href{https://github.com/yang-mills-group/OpenAI-YM-20261028/blob/3456e5e75cf8d4cd558955a5c7638fba2d68f464/YM/Further/SHELL_BIRTH_49/Fusing/Coal.lean\#L123-L126}{\texttt{Kiln}},
\href{https://github.com/yang-mills-group/OpenAI-YM-20261028/blob/3456e5e75cf8d4cd558955a5c7638fba2d68f464/YM/Further/SHELL_BIRTH_49/Fusing/Coal.lean\#L131}{\texttt{Kiln.Hot}},
\href{https://github.com/yang-mills-group/OpenAI-YM-20261028/blob/3456e5e75cf8d4cd558955a5c7638fba2d68f464/YM/Further/SHELL_BIRTH_49/Fusing/Coal.lean\#L132}{\texttt{fuel}},
\href{https://github.com/yang-mills-group/OpenAI-YM-20261028/blob/3456e5e75cf8d4cd558955a5c7638fba2d68f464/YM/Further/SHELL_BIRTH_49/Fusing/Coal.lean\#L133}{\texttt{Mole}},
\href{https://github.com/yang-mills-group/OpenAI-YM-20261028/blob/3456e5e75cf8d4cd558955a5c7638fba2d68f464/YM/Further/SHELL_BIRTH_49/Fusing/Coal.lean\#L165-L169}{\texttt{good}},
\href{https://github.com/yang-mills-group/OpenAI-YM-20261028/blob/3456e5e75cf8d4cd558955a5c7638fba2d68f464/YM/Further/SHELL_BIRTH_49/Fusing/Coal.lean\#L163-L164}{\texttt{fuel\_comp}},
\href{https://github.com/yang-mills-group/OpenAI-YM-20261028/blob/3456e5e75cf8d4cd558955a5c7638fba2d68f464/YM/Further/SHELL_BIRTH_49/Fusing/Coal.lean\#L190-L196}{\texttt{clubs\_away}},
\href{https://github.com/yang-mills-group/OpenAI-YM-20261028/blob/3456e5e75cf8d4cd558955a5c7638fba2d68f464/YM/Further/SHELL_BIRTH_49/Fusing/Coal.lean\#L197-L205}{\texttt{walls\_away}},
\href{https://github.com/yang-mills-group/OpenAI-YM-20261028/blob/3456e5e75cf8d4cd558955a5c7638fba2d68f464/YM/Further/SHELL_BIRTH_49/Fusing/Coal.lean\#L183-L187}{\texttt{takes\_mole}},
\href{https://github.com/yang-mills-group/OpenAI-YM-20261028/blob/3456e5e75cf8d4cd558955a5c7638fba2d68f464/YM/Further/SHELL_BIRTH_49/Fusing/Coal.lean\#L210-L240}{\texttt{Kiln.undo}},
\href{https://github.com/yang-mills-group/OpenAI-YM-20261028/blob/3456e5e75cf8d4cd558955a5c7638fba2d68f464/YM/Further/SHELL_BIRTH_49/Fusing/Coal.lean\#L242-L253}{\texttt{Kiln.attends}}.

\textbf{Proposition II.2.8 (fused gas of forced activities:
definitions).} Fix a gauge group \(G\) and a periodic lattice
\(\Lambda\).

\emph{Data.}

\begin{itemize}
\tightlist
\item
  \textbf{Pipes}: regular addresses \((\mathrm{reg}_i)\); forced
  addresses \((\mathrm{ker}_i)\), each with a support, a nonempty
  inventory contained in it, and an extra test on configurations;
  nonempty ``softer'' supports \((Y_i)\), for \(i\in\mathbb N\).
\item
  \textbf{Tap} \(d\): three families of real-valued functions of the
  configuration, namely regular coefficients \(E_i\), forced
  coefficients \(K_i\) and logarithmic coefficients \(\ell_i\).
\item
  \textbf{Labels} are \(\mathbb N\sqcup\mathbb N\): \emph{left} labels
  are forced kernels, \emph{right} labels are logarithms.

  \begin{itemize}
  \tightlist
  \item
    The \emph{patch} of a left label \(i\) is
    \(\mathrm{supp}(\mathrm{ker}_i)\); that of a right label \(i\) is
    \(Y_i\).
  \item
    The \emph{reach} of a left label is \(K_i(U)\); that of a right
    label is \(e^{\ell_i(U)}-1\).
  \end{itemize}
\item
  A \textbf{pellet} is a finite set of labels that contains at least one
  left label, is a spark for patch overlap, and whose left labels have
  pairwise disjoint kernel supports. Pellets form a countably infinite
  set.
\end{itemize}

\emph{Fused activities.} Fix a activity parameter point \(P\) and a
bijective enumeration \(F\) of the pellets. Eligibility is always tested
at the threshold of \(P\). The fused activities keep the regular part
\((\mathrm{reg}_i,E_i)\). Their forced part is indexed by pellets:

\begin{itemize}
\tightlist
\item
  address of \(F(i)\): support = union of the patches, inventory = union
  of the left inventories, extra test = all left kernels eligible;
\item
  value
  \[\mathrm{solid}_i(U)=\mathbf 1[\text{the address of }F(i)\text{ is eligible at }U]\ \prod_{j\in F(i)}\mathrm{reach}_j(U).\]
\end{itemize}

The unit also defines the hypotheses under which the fused gas is used.
Below, \(B\) is the cover size of \(P\).

\begin{itemize}
\tightlist
\item
  \textbf{Digest condition at \((d,P,\delta,b)\)}:

  \begin{enumerate}
  \def\labelenumi{\arabic{enumi}.}
  \tightlist
  \item
    the unfused activities \((\mathrm{reg},\mathrm{ker},E,K)\) are valid
    at \(P\) with regular and forced bounds \(\delta,b\);
  \item
    \(\ell_i(U)\neq0\) only if \(Y_i\) meets the large-field set of
    \(U\);
  \item
    \(|\ell_i(U)|\le B_i\) for every \(i,U\), with \(\sum_iB_i<\infty\);
  \item
    each \(\ell_i\) is gauge invariant and measurable, and has the same
    value on two configurations that agree on links internal to \(Y_i\)
    and have the same large-field flags on \(Y_i\);
  \item
    lattice affine maps permute the \(\ell_i\) covariantly.
  \end{enumerate}
\item
  \textbf{Fusion bounds at \((d,P,\kappa,\gamma)\)}:

  \begin{enumerate}
  \def\labelenumi{\arabic{enumi}.}
  \tightlist
  \item
    \(\gamma,\kappa\ge0\) and \((2B+1)^4\gamma\le1\);
  \item
    there is a label weight \(f\ge0\), independent of \(U\), with
    \(|\mathrm{reach}|\le e^{-\kappa}f\) on left labels and
    \(|\mathrm{reach}|\le f\) on right labels;
  \item
    \(f\) satisfies the hit condition of II.2.6 for the patches, with
    exponent (support exponent of \(P\)) \(+1\) and bound \(\gamma\).
  \end{enumerate}
\item
  For a parameter set \(T\subseteq\mathbb R\) and taps \(d_a\),
  \(a\in T\), the parametric conditions mean:

  \begin{enumerate}
  \def\labelenumi{\arabic{enumi}.}
  \tightlist
  \item
    for each label \(j\), \(a\mapsto\mathrm{reach}_j(d_a)\) is
    continuous in the supremum over configurations, relative to \(T\);
  \item
    one common choice of \(\kappa,\gamma,f\) satisfies the reach and hit
    bounds for every \(a\in T\);
  \item
    the fused values are continuous in the norm
    \(\sum_i e^{e_P\mathrm{cov}_B(S_i)}\sup_U|K_i(U)|\), where \(e_P\)
    is the support exponent and \(S_i\) the fused support.
  \end{enumerate}
\end{itemize}

\emph{Status:} definitions, plus one theorem (an enumeration exists).
The properties of the fused gas are the work items H01--H04, stated in
\href{https://github.com/yang-mills-group/OpenAI-YM-20261028/blob/3456e5e75cf8d4cd558955a5c7638fba2d68f464/YM/Bid/Kitchen.lean}{\texttt{YM/Bid/Kitchen.lean}}
and proved in III.10:

\begin{itemize}
\tightlist
\item
  \textbf{H01} (under the digest condition and fusion bounds):
  \(\mathrm{gas}_{\mathrm{fused}}(U)=\mathrm{gas}_{\mathrm{unfused}}(U)\exp\sum_i\ell_i(U)\)
  for every configuration \(U\).
\item
  \textbf{H02} (under the digest condition and fusion bounds): the fused
  forced activities satisfy the K-bound \(v\) whenever
  \(e^{-\kappa}(e^{\gamma}-1)\le v\).
\item
  \textbf{H03} (under the digest condition): they form a forced family
  with the lattice moves.
\item
  \textbf{H04}: if the unfused representations form a continuous
  parameter curve, conditions 1 and 2 above imply that the fused
  representations form such a curve. The path data, cutoff and
  parameters \(P\) are fixed along the curve.
\end{itemize}

At startup in III.9, the digest and fusion bounds follow from clip
compatibility, threshold and coupling bounds, and the packet estimates
in the cited results
(\href{https://github.com/yang-mills-group/OpenAI-YM-20261028/blob/3456e5e75cf8d4cd558955a5c7638fba2d68f464/YM/Extra/WILSON_STARTUP/Cake.lean\#L27-L36}{\texttt{Start.digests}},
\href{https://github.com/yang-mills-group/OpenAI-YM-20261028/blob/3456e5e75cf8d4cd558955a5c7638fba2d68f464/YM/Extra/WILSON_STARTUP/Funds.lean\#L120-L123}{\texttt{Loan.buys}}).
\emph{Lean:}
\href{https://github.com/yang-mills-group/OpenAI-YM-20261028/blob/3456e5e75cf8d4cd558955a5c7638fba2d68f464/YM/Ovens/Coffee.lean\#L23-L27}{\texttt{Ovens.Pipes}},
\href{https://github.com/yang-mills-group/OpenAI-YM-20261028/blob/3456e5e75cf8d4cd558955a5c7638fba2d68f464/YM/Ovens/Coffee.lean\#L32-L35}{\texttt{Tap}},
\href{https://github.com/yang-mills-group/OpenAI-YM-20261028/blob/3456e5e75cf8d4cd558955a5c7638fba2d68f464/YM/Ovens/Coffee.lean\#L48}{\texttt{reach}},
\href{https://github.com/yang-mills-group/OpenAI-YM-20261028/blob/3456e5e75cf8d4cd558955a5c7638fba2d68f464/YM/Ovens/Coffee.lean\#L57-L58}{\texttt{contains}},
\href{https://github.com/yang-mills-group/OpenAI-YM-20261028/blob/3456e5e75cf8d4cd558955a5c7638fba2d68f464/YM/Ovens/Coffee.lean\#L72}{\texttt{addr}},
\href{https://github.com/yang-mills-group/OpenAI-YM-20261028/blob/3456e5e75cf8d4cd558955a5c7638fba2d68f464/YM/Ovens/Coffee.lean\#L74-L76}{\texttt{solid}},
\href{https://github.com/yang-mills-group/OpenAI-YM-20261028/blob/3456e5e75cf8d4cd558955a5c7638fba2d68f464/YM/Ovens/Coffee.lean\#L78}{\texttt{fused}},
\href{https://github.com/yang-mills-group/OpenAI-YM-20261028/blob/3456e5e75cf8d4cd558955a5c7638fba2d68f464/YM/Ovens/Coffee.lean\#L81-L90}{\texttt{Digests}},
\href{https://github.com/yang-mills-group/OpenAI-YM-20261028/blob/3456e5e75cf8d4cd558955a5c7638fba2d68f464/YM/Ovens/Coffee.lean\#L94-L98}{\texttt{Buys}},
\href{https://github.com/yang-mills-group/OpenAI-YM-20261028/blob/3456e5e75cf8d4cd558955a5c7638fba2d68f464/YM/Ovens/Coffee.lean\#L105-L107}{\texttt{invariantJet}},
\href{https://github.com/yang-mills-group/OpenAI-YM-20261028/blob/3456e5e75cf8d4cd558955a5c7638fba2d68f464/YM/Ovens/Coffee.lean\#L109-L112}{\texttt{little}},
\href{https://github.com/yang-mills-group/OpenAI-YM-20261028/blob/3456e5e75cf8d4cd558955a5c7638fba2d68f464/YM/Ovens/Coffee.lean\#L115-L117}{\texttt{retakes}},
\href{https://github.com/yang-mills-group/OpenAI-YM-20261028/blob/3456e5e75cf8d4cd558955a5c7638fba2d68f464/YM/Ovens/Coffee.lean\#L121-L133}{\texttt{encoding}};
the obligations
\href{https://github.com/yang-mills-group/OpenAI-YM-20261028/blob/3456e5e75cf8d4cd558955a5c7638fba2d68f464/YM/Bid/Kitchen.lean\#L15-L16}{\texttt{Bid.H01}}--\href{https://github.com/yang-mills-group/OpenAI-YM-20261028/blob/3456e5e75cf8d4cd558955a5c7638fba2d68f464/YM/Bid/Kitchen.lean\#L33-L35}{\texttt{H04}},
proved by
\href{https://github.com/yang-mills-group/OpenAI-YM-20261028/blob/3456e5e75cf8d4cd558955a5c7638fba2d68f464/YM/Work/H01.lean\#L253-L255}{\texttt{Work.H01.answer}},
\href{https://github.com/yang-mills-group/OpenAI-YM-20261028/blob/3456e5e75cf8d4cd558955a5c7638fba2d68f464/YM/Work/H02.lean\#L17-L80}{\texttt{H02.answer}},
\href{https://github.com/yang-mills-group/OpenAI-YM-20261028/blob/3456e5e75cf8d4cd558955a5c7638fba2d68f464/YM/Work/H03.lean\#L332-L334}{\texttt{H03.answer}},
\href{https://github.com/yang-mills-group/OpenAI-YM-20261028/blob/3456e5e75cf8d4cd558955a5c7638fba2d68f464/YM/Work/H04.lean\#L155-L161}{\texttt{H04.answer}}
(III.10).

\paragraph{C. Switching}\label{c.-switching}

\textbf{Proposition II.2.9 (switching).} Let \(\iota\) be finite. For
\(x\in\mathbb R\) and \(g:2^\iota\to\mathbb R\) put
\(\mathrm{val}_x(g)(r)=\sum_{i\subseteq r}x^{|i|}g(i)\). This is the
value at \(x\mathbf 1_r\in\{0,x\}^\iota\) of the multilinear polynomial
\(\sum_ig(i)\prod_{v\in i}X_v\).

\begin{enumerate}
\def\labelenumi{(\alph{enumi})}
\item
  For \(x>0\), \(\mathrm{val}_x\) is a linear bijection of
  \(\mathbb R^{2^\iota}\).
\item
  Let \((f_x)_{x\in[0,1]}\) be functions on \(2^\iota\). Say \((g_x)\)
  \textbf{models} \(f\) if \(\mathrm{val}_x(g_x)=f_x\) for \(x\in[0,1]\)
  and \(g_x\to g_0\) as \(x\downarrow0\). In that case
  \(\mathrm{jet}(f):=\lim_{x\downarrow0}\mathrm{val}_x^{-1}(f_x)\)
  exists and equals \(g_0\).
\item
  A family \((f^u)_u\) is \textbf{paid at cost \(c\)} if there is
  \(b\ge0\) on \(2^\iota\) with \(\sum_ib(i)\le c\) such that, for each
  \(u\), some model \(g^u\) of \(f^u\) satisfies \(|g^u_x(i)|\le b(i)\)
  for all \(x\in[0,1]\). Then:
\end{enumerate}

\begin{itemize}
\tightlist
\item
  \(|f^u_x(r)|\le c\) for every \(u\), \(r\) and \(x\in[0,1]\);
\item
  paid families are closed under finite sums (costs add), finite
  products (costs multiply), scalar multiples (cost multiplied by the
  absolute value), and restriction by zero outside a set of parameters.
  Families of paid functions with summable costs may also be added, and
  their jets sum to the jet of the sum. Products of models obey the rule
  \(X_v^2=xX_v\): \(\mathrm{val}_x\) turns the product
  \((p\ast_xq)(k)=\sum_{i\cup j=k}x^{|i\cap j|}p(i)q(j)\) into the
  pointwise product;
\item
  a monomial \(u\mapsto c(u)\,\mathbf 1[i\subseteq r]\,x^{|i|}\) with
  \(|c|\le b\) is paid at cost \(b\).
\end{itemize}

\begin{enumerate}
\def\labelenumi{(\alph{enumi})}
\setcounter{enumi}{3}
\item
  If \(f\) has a model and \(f_x(r)=f_x(r\setminus\{v\})\) for every
  \(r\) and \(x\in[0,1]\), then \(\mathrm{jet}(f)(i)=0\) whenever
  \(v\in i\). For a paid family on a measurable parameter space, if
  \(u\mapsto f^u_x(r)\) is measurable for each fixed \(x\in[0,1]\) and
  \(r\), then \(u\mapsto\mathrm{jet}(f^u)(i)\) is measurable for every
  \(i\).
\item
  For \(\rho:\iota\to\mathcal A_R\) put
  \(\mathrm{snd}_\rho(f)=\sum_{i\subseteq\iota}\mathrm{jet}(f)(i)\prod_{v\in i}\rho_v\).
\end{enumerate}

\begin{itemize}
\tightlist
\item
  If \((f^u)\) is paid at cost \(d\) and \(\|\rho_v\|_w\le1\) for all
  \(v\), then the coefficients of \(\mathrm{snd}_\rho(f^u)\) are
  dominated, uniformly in \(u\), by a nonnegative array of \(w\)-norm at
  most \(d\).
\item
  If \(\rho_v(\emptyset)=0\) for all \(v\), the coefficient of
  \(\mathrm{snd}_\rho(f)\) on a singleton \(\{p\}\) is
  \(\sum_v\mathrm{jet}(f)(\{v\})\,\rho_v(\{p\})\).
\end{itemize}

\emph{Status:} theorem. For square-zero \(\rho_v\),
\(\mathrm{snd}_\rho\) is multiplicative on functions with models. For
paid families with summable costs, it commutes with sums; for paid
\(f\), it sends \(e^f\) to the exponential series of
\(\mathrm{snd}_\rho(f)\)
(\href{https://github.com/yang-mills-group/OpenAI-YM-20261028/blob/3456e5e75cf8d4cd558955a5c7638fba2d68f464/YM/MergeProof/SwitchCost.lean\#L64-L65}{\texttt{mul\_sounds}},
\href{https://github.com/yang-mills-group/OpenAI-YM-20261028/blob/3456e5e75cf8d4cd558955a5c7638fba2d68f464/YM/MergeProof/SwitchCost.lean\#L81-L87}{\texttt{soundSum}},
\href{https://github.com/yang-mills-group/OpenAI-YM-20261028/blob/3456e5e75cf8d4cd558955a5c7638fba2d68f464/YM/MergeProof/SwitchCost.lean\#L88-L91}{\texttt{soundExp}}).
\emph{Lean:}
\href{https://github.com/yang-mills-group/OpenAI-YM-20261028/blob/3456e5e75cf8d4cd558955a5c7638fba2d68f464/YM/MergeProof/Switch.lean\#L26}{\texttt{Switch.value}},
\href{https://github.com/yang-mills-group/OpenAI-YM-20261028/blob/3456e5e75cf8d4cd558955a5c7638fba2d68f464/YM/MergeProof/Switch.lean\#L33-L42}{\texttt{entire}},
\href{https://github.com/yang-mills-group/OpenAI-YM-20261028/blob/3456e5e75cf8d4cd558955a5c7638fba2d68f464/YM/MergeProof/Switch.lean\#L46}{\texttt{jet}},
\href{https://github.com/yang-mills-group/OpenAI-YM-20261028/blob/3456e5e75cf8d4cd558955a5c7638fba2d68f464/YM/MergeProof/Switch.lean\#L48-L50}{\texttt{Models}},
\href{https://github.com/yang-mills-group/OpenAI-YM-20261028/blob/3456e5e75cf8d4cd558955a5c7638fba2d68f464/YM/MergeProof/Switch.lean\#L58-L63}{\texttt{Models.exact}},
\href{https://github.com/yang-mills-group/OpenAI-YM-20261028/blob/3456e5e75cf8d4cd558955a5c7638fba2d68f464/YM/MergeProof/Switch.lean\#L56}{\texttt{Paid}},
\href{https://github.com/yang-mills-group/OpenAI-YM-20261028/blob/3456e5e75cf8d4cd558955a5c7638fba2d68f464/YM/MergeProof/Switch.lean\#L77-L79}{\texttt{Paid.tasting}},
\href{https://github.com/yang-mills-group/OpenAI-YM-20261028/blob/3456e5e75cf8d4cd558955a5c7638fba2d68f464/YM/MergeProof/Switch.lean\#L88-L93}{\texttt{Paid.add}},
\href{https://github.com/yang-mills-group/OpenAI-YM-20261028/blob/3456e5e75cf8d4cd558955a5c7638fba2d68f464/YM/MergeProof/Switch.lean\#L116-L119}{\texttt{Paid.sum}},
\href{https://github.com/yang-mills-group/OpenAI-YM-20261028/blob/3456e5e75cf8d4cd558955a5c7638fba2d68f464/YM/MergeProof/Switch.lean\#L108-L112}{\texttt{paidClick}},
\href{https://github.com/yang-mills-group/OpenAI-YM-20261028/blob/3456e5e75cf8d4cd558955a5c7638fba2d68f464/YM/MergeProof/SwitchProd.lean\#L10}{\texttt{comb}},
\href{https://github.com/yang-mills-group/OpenAI-YM-20261028/blob/3456e5e75cf8d4cd558955a5c7638fba2d68f464/YM/MergeProof/SwitchProd.lean\#L46-L50}{\texttt{Paid.mul}},
\href{https://github.com/yang-mills-group/OpenAI-YM-20261028/blob/3456e5e75cf8d4cd558955a5c7638fba2d68f464/YM/MergeProof/SwitchProd.lean\#L52-L55}{\texttt{Paid.prod}},
\href{https://github.com/yang-mills-group/OpenAI-YM-20261028/blob/3456e5e75cf8d4cd558955a5c7638fba2d68f464/YM/MergeProof/Digits.lean\#L27-L36}{\texttt{jet\_erase}},
\href{https://github.com/yang-mills-group/OpenAI-YM-20261028/blob/3456e5e75cf8d4cd558955a5c7638fba2d68f464/YM/MergeProof/Digits.lean\#L38-L49}{\texttt{Paid.jet\_meas}},
\href{https://github.com/yang-mills-group/OpenAI-YM-20261028/blob/3456e5e75cf8d4cd558955a5c7638fba2d68f464/YM/MergeProof/SwitchCost.lean\#L33-L36}{\texttt{mapScores}},
\href{https://github.com/yang-mills-group/OpenAI-YM-20261028/blob/3456e5e75cf8d4cd558955a5c7638fba2d68f464/YM/MergeProof/SwitchCost.lean\#L60}{\texttt{sound}},
\href{https://github.com/yang-mills-group/OpenAI-YM-20261028/blob/3456e5e75cf8d4cd558955a5c7638fba2d68f464/YM/MergeProof/Jetty.lean\#L13-L19}{\texttt{Accent.crystal}},
\href{https://github.com/yang-mills-group/OpenAI-YM-20261028/blob/3456e5e75cf8d4cd558955a5c7638fba2d68f464/YM/MergeProof/Jetty.lean\#L29-L35}{\texttt{Accent.quality}}.

\paragraph{D. Gaussian polynomial
calculus}\label{d.-gaussian-polynomial-calculus}

\textbf{Conventions.}

\begin{itemize}
\tightlist
\item
  \(\mathbb C[X]=\mathbb C[X_v:v\in V]\). The variable set \(V\) need
  not be finite; each polynomial uses finitely many variables.
\item
  \(\|p\|_1\) is the sum of the absolute values of the coefficients of
  \(p\).
\item
  For \(b\in\mathbb C^V\), \(\partial_b\) is the derivation with
  \(\partial_bX_v=b_v\).
\item
  A \emph{covariance} is a symmetric \(Q:V\times V\to\mathbb C\); no
  positivity is assumed.
\item
  \(\pi(B,d)=(2B(d+1))^d\).
\end{itemize}

\textbf{Proposition II.2.10 (algebraic Gaussian expectations).} Let
\(Q\) be a covariance. The operators \(X_v+\partial_{Q_{v\cdot}}\) on
\(\mathbb C[X]\) commute; define \(\mathcal G_Qp=p(X+\partial_Q)\,1\).

\begin{enumerate}
\def\labelenumi{(\alph{enumi})}
\tightlist
\item
  \(\mathcal G_Q\) satisfies:
\end{enumerate}

\begin{itemize}
\tightlist
\item
  \(\mathcal G_Q1=1\) and
  \(\mathcal G_Q(X_vp)=X_v\,\mathcal G_Qp+\partial_{Q_{v\cdot}}\mathcal G_Qp\);
\item
  \(\mathcal G_0=\mathrm{id}\) and
  \(\mathcal G_P\mathcal G_Q=\mathcal G_{P+Q}\);
\item
  \(\partial_b\mathcal G_Q=\mathcal G_Q\partial_b\).
\end{itemize}

(Equivalently
\(\mathcal G_Q=\exp(\tfrac12\sum_{v,w}Q_{vw}\partial_v\partial_w)\), so
\(\mathcal G_{-M}\) is Wick ordering with respect to \(M\). On each
polynomial the derivative sum and exponential series have only finitely
many nonzero terms.)

\begin{enumerate}
\def\labelenumi{(\alph{enumi})}
\setcounter{enumi}{1}
\item
  For \(m\in\mathbb C^V\) put \(\mathbb E_{Q,m}(p)=(\mathcal G_Qp)(m)\).
  This is a linear functional with \(\mathbb E_{Q,m}(1)=1\), and it
  satisfies, for every finite set \(s\) and choice of variables \(z_j\),
  the Wick recursion
  \[\mathbb E_{Q,m}\Big(X_v\prod_{j\in s}X_{z_j}\Big)=m_v\,\mathbb E_{Q,m}\Big(\prod_{j\in s}X_{z_j}\Big)+\sum_{j\in s}Q_{vz_j}\,\mathbb E_{Q,m}\Big(\prod_{k\in s\setminus\{j\}}X_{z_k}\Big).\]
\item
  \emph{Change of variables.} For a matrix \(A=(A_{vz})\) with finitely
  supported rows, define the linear substitution \(p\mapsto p\circ A\)
  by \(X_v\mapsto\sum_zA_{vz}X_z\). If \(Q\) is a covariance and \(m\) a
  mean on the new variables, then:
\end{enumerate}

\begin{itemize}
\tightlist
\item
  \(\mathcal G_Q(p\circ A)=(\mathcal G_{AQA^{T}}p)\circ A\);
\item
  \(\mathbb E_{Q,m}(p\circ A)=\mathbb E_{AQA^{T},Am}(p)\).
\end{itemize}

\begin{enumerate}
\def\labelenumi{(\alph{enumi})}
\setcounter{enumi}{3}
\item
  \emph{Identification.} Let \(\xi\) be a random vector in
  \(\mathbb R^V\) under a probability measure. Assume all polynomial
  moments of \(\xi\) are finite and
  \(\mathbb E[\xi_v\,q(\xi)]=\mathbb E[(\partial_{Q_{v\cdot}}q)(\xi)]\)
  for every polynomial \(q\) (Gaussian integration by parts). Then
  \(\mathbb E[q(m+\xi)]=\mathbb E_{Q,m}(q)\) for every
  \(m\in\mathbb R^V\).
\item
  \emph{Bounds.} Let \(B\ge1\), \(|m_v|\le B\), \(|Q_{vw}|\le B\) and
  \(\deg p\le d\). Then \[|\mathbb E_{Q,m}(p)|\le\|p\|_1\,\pi(B,d).\] If
  \((Q',m')\) obeys the same bounds and each corresponding entry of the
  means and covariances differs by at most \(E\ge0\), then
  \(|\mathbb E_{Q,m}(p)-\mathbb E_{Q',m'}(p)|\le\|p\|_1\,E\,\pi(B,d)\).
  In the first bound, and in the second when \(m=m'\), it suffices to
  assume the hypotheses for entries indexed by the variables of \(p\).
\end{enumerate}

\emph{Status:} theorem. Used by III.11 and VII.6. \emph{Lean:}
\href{https://github.com/yang-mills-group/OpenAI-YM-20261028/blob/3456e5e75cf8d4cd558955a5c7638fba2d68f464/YM/Next/RANK/Gaussian.lean\#L21-L23}{\texttt{Normal.Symmetric}},
\href{https://github.com/yang-mills-group/OpenAI-YM-20261028/blob/3456e5e75cf8d4cd558955a5c7638fba2d68f464/YM/Next/RANK/Gaussian.lean\#L40}{\texttt{slit}},
\href{https://github.com/yang-mills-group/OpenAI-YM-20261028/blob/3456e5e75cf8d4cd558955a5c7638fba2d68f464/YM/Next/RANK/Gaussian.lean\#L64-L65}{\texttt{cream\_comm}},
\href{https://github.com/yang-mills-group/OpenAI-YM-20261028/blob/3456e5e75cf8d4cd558955a5c7638fba2d68f464/YM/Next/RANK/Gaussian.lean\#L77}{\texttt{milk}},
\href{https://github.com/yang-mills-group/OpenAI-YM-20261028/blob/3456e5e75cf8d4cd558955a5c7638fba2d68f464/YM/Next/RANK/Gaussian.lean\#L82}{\texttt{feeds}},
\href{https://github.com/yang-mills-group/OpenAI-YM-20261028/blob/3456e5e75cf8d4cd558955a5c7638fba2d68f464/YM/Next/RANK/Gaussian.lean\#L83-L88}{\texttt{skims}},
\href{https://github.com/yang-mills-group/OpenAI-YM-20261028/blob/3456e5e75cf8d4cd558955a5c7638fba2d68f464/YM/Next/RANK/Gaussian.lean\#L89-L94}{\texttt{froze}},
\href{https://github.com/yang-mills-group/OpenAI-YM-20261028/blob/3456e5e75cf8d4cd558955a5c7638fba2d68f464/YM/Next/RANK/Gaussian.lean\#L96-L102}{\texttt{warms}},
\href{https://github.com/yang-mills-group/OpenAI-YM-20261028/blob/3456e5e75cf8d4cd558955a5c7638fba2d68f464/YM/Next/RANK/Gaussian.lean\#L106}{\texttt{blend}},
\href{https://github.com/yang-mills-group/OpenAI-YM-20261028/blob/3456e5e75cf8d4cd558955a5c7638fba2d68f464/YM/Next/RANK/GaussTrack.lean\#L31-L35}{\texttt{spring}},
\href{https://github.com/yang-mills-group/OpenAI-YM-20261028/blob/3456e5e75cf8d4cd558955a5c7638fba2d68f464/YM/Next/RANK/GaussSubst.lean\#L76-L81}{\texttt{prepared}},
\href{https://github.com/yang-mills-group/OpenAI-YM-20261028/blob/3456e5e75cf8d4cd558955a5c7638fba2d68f464/YM/Next/RANK/GaussSubst.lean\#L83-L84}{\texttt{wets}},
\href{https://github.com/yang-mills-group/OpenAI-YM-20261028/blob/3456e5e75cf8d4cd558955a5c7638fba2d68f464/YM/Next/RANK/ParticleMoment.lean\#L21-L25}{\texttt{VMeasure}},
\href{https://github.com/yang-mills-group/OpenAI-YM-20261028/blob/3456e5e75cf8d4cd558955a5c7638fba2d68f464/YM/Next/RANK/ParticleMoment.lean\#L61-L78}{\texttt{payment}},
\href{https://github.com/yang-mills-group/OpenAI-YM-20261028/blob/3456e5e75cf8d4cd558955a5c7638fba2d68f464/YM/Next/RANK/Common.lean\#L36}{\texttt{Common.priced}},
\href{https://github.com/yang-mills-group/OpenAI-YM-20261028/blob/3456e5e75cf8d4cd558955a5c7638fba2d68f464/YM/Next/RANK/Common.lean\#L121-L123}{\texttt{Common.width}},
\href{https://github.com/yang-mills-group/OpenAI-YM-20261028/blob/3456e5e75cf8d4cd558955a5c7638fba2d68f464/YM/Next/RANK/Common.lean\#L124-L126}{\texttt{Common.differs}},
\href{https://github.com/yang-mills-group/OpenAI-YM-20261028/blob/3456e5e75cf8d4cd558955a5c7638fba2d68f464/YM/Next/RANK/Guard.lean\#L44}{\texttt{Tapes.bound\_v}},
\href{https://github.com/yang-mills-group/OpenAI-YM-20261028/blob/3456e5e75cf8d4cd558955a5c7638fba2d68f464/YM/Next/RANK/Guard.lean\#L36-L43}{\texttt{Tapes.differently}}.

\textbf{Proposition II.2.11 (estimates for Gaussian polynomial steps).}
All polynomial degrees below are total degrees.

\begin{enumerate}
\def\labelenumi{(\alph{enumi})}
\item
  \emph{Weighted coefficient norms.} For \(w:V\to\mathbb N\) and
  \(r\ge0\), the quantity
  \(\|p\|_{w,r}=\sum_\alpha r^{w\cdot\alpha}|c_\alpha|\) is
  submultiplicative, and \(\|\cdot\|_1=\|\cdot\|_{0,1}\). It is a norm
  when \(r>0\) and may be a seminorm at \(r=0\). For a linear
  substitution \(X_v\mapsto e_v\) with \(\|e_v\|_{w,r}\le c\) and
  \(c\ge1\), \[\|p\circ e\|_{w,r}\le\|p\|_1\,c^{\deg p}.\]
\item
  \emph{Smoothing near the identity.} If \(0\le e\le1\),
  \(|Q_{vw}|\le e\) and \(\deg p\le D\), then
  \(\|\mathcal G_Qp-p\|_1\le\|p\|_1\,e\,(2(D+2))^D\) and
  \(\deg\mathcal G_Qp\le D\).
\item
  \emph{One step.} Let \(T_{H,Q}p=(\mathcal G_Qp)\circ H\), where
  \(H=(H_v)\) are linear forms in new variables and \(Q\) is a
  covariance.
\end{enumerate}

\begin{itemize}
\tightlist
\item
  If \(C\ge1\), \(\|H_v\|_1\le C\), \(|Q_{vw}|\le C\) and
  \(\deg p\le d\), then \(\|T_{H,Q}p\|_1\le\|p\|_1(2(d+2)C)^d\) and
  \(\deg T_{H,Q}p\le d\).
\item
  For two such pairs, assume \(\varepsilon\ge0\),
  \(\|H_v-H'_v\|_1\le\varepsilon\) and
  \(|Q_{vw}-Q'_{vw}|\le\varepsilon\). The difference bound is
  \(\|p\|_1\,\varepsilon\,(2(d+2)C)^d\)
  (\href{https://github.com/yang-mills-group/OpenAI-YM-20261028/blob/3456e5e75cf8d4cd558955a5c7638fba2d68f464/YM/Next/RANK/Herald.lean\#L125-L126}{\texttt{crossing}}).
\item
  If \(Q\) vanishes between the variables of \(P\) and those of \(R\),
  then \(T_{H,Q}(PR)=T_{H,Q}(P)\,T_{H,Q}(R)\).
\end{itemize}

\begin{enumerate}
\def\labelenumi{(\alph{enumi})}
\setcounter{enumi}{3}
\tightlist
\item
  \emph{Linearization of a rescaled substitution.} Let:
\end{enumerate}

\begin{itemize}
\tightlist
\item
  \(f_v\) be linear forms and \(A_v\) polynomials with no terms of
  degree \(<2\), with \(\|f_v\|_1,\|A_v\|_1\le B\) and \(B\ge1\);
\item
  \(0<t\le1\), and \(\Phi_t\) the algebra map
  \(X_v\mapsto f_v+t^{-1}A_v(tX)\);
\item
  \(K\) is a covariance, \(0\le\varepsilon\le1\),
  \(|K_{vw}|\le\varepsilon\) and \(\deg P\le n\).
\end{itemize}

Then
\[\big\|\Phi_t(P)-(\mathcal G_KP)\circ f\big\|_1\le\|P\|_1\big(n\,tB(2B)^n+\varepsilon(2(n+2))^nB^n\big).\]

\begin{enumerate}
\def\labelenumi{(\alph{enumi})}
\setcounter{enumi}{4}
\tightlist
\item
  \emph{Guards and second moments.} Call \(g\) an
  \textbf{\(\eta\)-guard} for \(\mu\) if \(g\) is measurable,
  \(0\le g\le1\), and \(\int(1-g)\,d\mu\le\eta\).
\end{enumerate}

\begin{itemize}
\tightlist
\item
  For a finite measure \(\mu\), the product of an \(\eta\)-guard and an
  \(\varepsilon\)-guard is an \((\eta+\varepsilon)\)-guard.
\item
  For a probability measure, measurable \(f\) with integrable \(|f|^2\)
  and \(\int|f|^2d\mu\le B\), and an \(\eta\)-guard \(g\):
  \(\big|\int gf\,d\mu-\int f\,d\mu\big|\le\varepsilon'B+\eta/\varepsilon'\)
  for every \(\varepsilon'>0\).
\item
  Under II.2.10(d), if \(\deg p\le d\), \(B\ge1\) and
  \(|m_v|,|Q_{vw}|\le B\) on the variables of \(p\), then
  \(|p(m+\xi)|^2\) is integrable and
  \(\int|p(m+\xi)|^2\le\|p\|_1^2\,\pi(B,2d)\).
\end{itemize}

\begin{enumerate}
\def\labelenumi{(\alph{enumi})}
\setcounter{enumi}{5}
\tightlist
\item
  \emph{Coefficients from values.}
\end{enumerate}

\begin{itemize}
\tightlist
\item
  For finite \(V\), \(d\in\mathbb N\) and \(\varepsilon>0\), there are a
  finite nonempty set \(T\subseteq[-\varepsilon,\varepsilon]^V\) and
  \(C\ge0\) such that \(\|p\|_1\le C\max_{x\in T}|p(x)|\) for all \(p\)
  of degree \(\le d\).
\item
  A polynomial is determined by its values on a real cube of positive
  side length.
\item
  If \(\deg p\le D\), \(B\ge1\), \(v,w\in\mathbb R^V\), \(\delta\ge0\),
  \(|v_i|,|w_i|\le B\) and \(|v_i-w_i|\le\delta\), then
  \(|p(v)-p(w)|\le\|p\|_1\,DB^D\delta\).
\end{itemize}

\emph{Status:} theorem. Used by III.11 and VII.6. \emph{Lean:}
\href{https://github.com/yang-mills-group/OpenAI-YM-20261028/blob/3456e5e75cf8d4cd558955a5c7638fba2d68f464/YM/Next/RANK/Rise.lean\#L13}{\texttt{Rise.cents}},
\href{https://github.com/yang-mills-group/OpenAI-YM-20261028/blob/3456e5e75cf8d4cd558955a5c7638fba2d68f464/YM/Next/RANK/Rise.lean\#L33-L37}{\texttt{Rise.funds}},
\href{https://github.com/yang-mills-group/OpenAI-YM-20261028/blob/3456e5e75cf8d4cd558955a5c7638fba2d68f464/YM/Next/RANK/Rise.lean\#L51-L66}{\texttt{Rise.rise}},
\href{https://github.com/yang-mills-group/OpenAI-YM-20261028/blob/3456e5e75cf8d4cd558955a5c7638fba2d68f464/YM/Next/RANK/Ample.lean\#L79-L96}{\texttt{LiftMark.bounds}},
\href{https://github.com/yang-mills-group/OpenAI-YM-20261028/blob/3456e5e75cf8d4cd558955a5c7638fba2d68f464/YM/Next/RANK/GaussTags.lean\#L113}{\texttt{Normal.pipe}},
\href{https://github.com/yang-mills-group/OpenAI-YM-20261028/blob/3456e5e75cf8d4cd558955a5c7638fba2d68f464/YM/Next/RANK/Herald.lean\#L14-L16}{\texttt{Tapes.Valve}},
\href{https://github.com/yang-mills-group/OpenAI-YM-20261028/blob/3456e5e75cf8d4cd558955a5c7638fba2d68f464/YM/Next/RANK/Herald.lean\#L20-L23}{\texttt{Valve.Fair}},
\href{https://github.com/yang-mills-group/OpenAI-YM-20261028/blob/3456e5e75cf8d4cd558955a5c7638fba2d68f464/YM/Next/RANK/Herald.lean\#L37}{\texttt{crest}},
\href{https://github.com/yang-mills-group/OpenAI-YM-20261028/blob/3456e5e75cf8d4cd558955a5c7638fba2d68f464/YM/Next/RANK/Herald.lean\#L119-L124}{\texttt{crux}},
\href{https://github.com/yang-mills-group/OpenAI-YM-20261028/blob/3456e5e75cf8d4cd558955a5c7638fba2d68f464/YM/Next/RANK/Crust.lean\#L59}{\texttt{bridge}},
\href{https://github.com/yang-mills-group/OpenAI-YM-20261028/blob/3456e5e75cf8d4cd558955a5c7638fba2d68f464/YM/Next/RANK/Pin.lean\#L29-L30}{\texttt{LiftMark.garnish}},
\href{https://github.com/yang-mills-group/OpenAI-YM-20261028/blob/3456e5e75cf8d4cd558955a5c7638fba2d68f464/YM/Next/RANK/Pin.lean\#L66-L84}{\texttt{LiftMark.stock}},
\href{https://github.com/yang-mills-group/OpenAI-YM-20261028/blob/3456e5e75cf8d4cd558955a5c7638fba2d68f464/YM/Next/RANK/Guard.lean\#L69-L73}{\texttt{Guarded}},
\href{https://github.com/yang-mills-group/OpenAI-YM-20261028/blob/3456e5e75cf8d4cd558955a5c7638fba2d68f464/YM/Next/RANK/Guard.lean\#L79-L88}{\texttt{guards\_mul}},
\href{https://github.com/yang-mills-group/OpenAI-YM-20261028/blob/3456e5e75cf8d4cd558955a5c7638fba2d68f464/YM/Next/RANK/Guard.lean\#L108-L125}{\texttt{escaped}},
\href{https://github.com/yang-mills-group/OpenAI-YM-20261028/blob/3456e5e75cf8d4cd558955a5c7638fba2d68f464/YM/Next/RANK/Guard.lean\#L90-L105}{\texttt{squareFits}},
\href{https://github.com/yang-mills-group/OpenAI-YM-20261028/blob/3456e5e75cf8d4cd558955a5c7638fba2d68f464/YM/Next/RANK/SpectrumReal.lean\#L56-L81}{\texttt{Spindle.molds}},
\href{https://github.com/yang-mills-group/OpenAI-YM-20261028/blob/3456e5e75cf8d4cd558955a5c7638fba2d68f464/YM/Next/RANK/SpectrumReal.lean\#L22-L28}{\texttt{Spindle.listen}},
\href{https://github.com/yang-mills-group/OpenAI-YM-20261028/blob/3456e5e75cf8d4cd558955a5c7638fba2d68f464/YM/Next/RANK/Spirals.lean\#L20-L41}{\texttt{Spindle.amove}}.

\paragraph{E. Numbers}\label{e.-numbers}

\textbf{Proposition II.2.12 (scale functions).} For \(s>e\) put
\[y=\log\log s,\quad \Pi=e^{y^2},\quad \theta=s^{-1/2}\Pi,\quad h=e^{y^{3/2}},\quad j=h^{1/5},\quad \omega=e^{-\frac7{10}j},\quad M=\lceil(\log s)^2\rceil,\quad W=M^2,\]
all as functions of \(s\). Here \(\theta\) is the original small-field
threshold. Then:

\begin{enumerate}
\def\labelenumi{\arabic{enumi}.}
\tightlist
\item
  For every \(b\in\mathbb R\) and \(p>0\), eventually
  \(\Pi(s)^b\le s^p\).
\item
  Eventually \(s\theta^2=\Pi^2\), \(s\theta^3\le1\), \(s\theta^4\le1\),
  and \(2\Pi^v\theta^2\le1\) for every \(v\le1\).
\item
  For all \(a>0\), \(k\in\mathbb N\) and \(c\in\mathbb R\), eventually:

  \begin{itemize}
  \tightlist
  \item
    \(W(s)^k\le\Pi(s)^a\) and \(h(s)\le\Pi(s)^a\);
  \item
    \(c\log W(s)\le a\,y(s)^2\);
  \item
    \(y(s)^2\le a\,W(s/2)\) and \(y(s)^2\le a\,j(s/2)\);
  \item
    \(j(s)\le h(s/2)\) and \(y(s)\le y(s/2)+1\).
  \end{itemize}
\item
  For every \(c\in\mathbb R\) and \(a<b\), \(c\,\Pi^a\le\Pi^b\)
  eventually.
\end{enumerate}

Here \(f\ll g\) means \(f/g\to0\) as \(s\to\infty\). For every fixed
\(a,p>0\), \(y^2\ll j\ll h\ll\Pi^a\ll s^p\); also
\(W(s)\sim(\log s)^4\).

\emph{Status:} theorem. Item 2 is used in III.10, III.11 and V.2 (the
per-site vacuum-energy bound
\href{https://github.com/yang-mills-group/OpenAI-YM-20261028/blob/3456e5e75cf8d4cd558955a5c7638fba2d68f464/YM/Vac/Treasury.lean\#L115-L150}{\texttt{treasury\_stage}});
item 3 is used in III.11. \emph{Lean:}
\href{https://github.com/yang-mills-group/OpenAI-YM-20261028/blob/3456e5e75cf8d4cd558955a5c7638fba2d68f464/YM/Vac/Treasury.lean\#L67-L69}{\texttt{Numbers.monitor}},
\href{https://github.com/yang-mills-group/OpenAI-YM-20261028/blob/3456e5e75cf8d4cd558955a5c7638fba2d68f464/YM/Vac/Treasury.lean\#L71}{\texttt{Hat}},
\href{https://github.com/yang-mills-group/OpenAI-YM-20261028/blob/3456e5e75cf8d4cd558955a5c7638fba2d68f464/YM/Vac/Treasury.lean\#L72}{\texttt{P}},
\href{https://github.com/yang-mills-group/OpenAI-YM-20261028/blob/3456e5e75cf8d4cd558955a5c7638fba2d68f464/YM/Vac/Treasury.lean\#L73}{\texttt{theta}},
\href{https://github.com/yang-mills-group/OpenAI-YM-20261028/blob/3456e5e75cf8d4cd558955a5c7638fba2d68f464/YM/Vac/Treasury.lean\#L74-L78}{\texttt{cost}},
\href{https://github.com/yang-mills-group/OpenAI-YM-20261028/blob/3456e5e75cf8d4cd558955a5c7638fba2d68f464/YM/Vac/Treasury.lean\#L83-L108}{\texttt{uniform\_loop}},
\href{https://github.com/yang-mills-group/OpenAI-YM-20261028/blob/3456e5e75cf8d4cd558955a5c7638fba2d68f464/YM/Next/RANK/Reign.lean\#L25}{\texttt{Perch.bill}},
\href{https://github.com/yang-mills-group/OpenAI-YM-20261028/blob/3456e5e75cf8d4cd558955a5c7638fba2d68f464/YM/Next/RANK/Reign.lean\#L28}{\texttt{walls}},
\href{https://github.com/yang-mills-group/OpenAI-YM-20261028/blob/3456e5e75cf8d4cd558955a5c7638fba2d68f464/YM/Next/RANK/Reign.lean\#L30}{\texttt{hot}},
\href{https://github.com/yang-mills-group/OpenAI-YM-20261028/blob/3456e5e75cf8d4cd558955a5c7638fba2d68f464/YM/Next/RANK/Reign.lean\#L31}{\texttt{join}},
\href{https://github.com/yang-mills-group/OpenAI-YM-20261028/blob/3456e5e75cf8d4cd558955a5c7638fba2d68f464/YM/Next/RANK/Reign.lean\#L32}{\texttt{w}},
\href{https://github.com/yang-mills-group/OpenAI-YM-20261028/blob/3456e5e75cf8d4cd558955a5c7638fba2d68f464/YM/Next/RANK/Reign.lean\#L59-L69}{\texttt{logspan}},
\href{https://github.com/yang-mills-group/OpenAI-YM-20261028/blob/3456e5e75cf8d4cd558955a5c7638fba2d68f464/YM/Next/RANK/Reign.lean\#L71-L90}{\texttt{rate\_wall}},
\href{https://github.com/yang-mills-group/OpenAI-YM-20261028/blob/3456e5e75cf8d4cd558955a5c7638fba2d68f464/YM/Next/RANK/Reign.lean\#L94-L108}{\texttt{rate\_inverse}},
\href{https://github.com/yang-mills-group/OpenAI-YM-20261028/blob/3456e5e75cf8d4cd558955a5c7638fba2d68f464/YM/Next/RANK/Reign.lean\#L110-L122}{\texttt{rate\_fire}},
\href{https://github.com/yang-mills-group/OpenAI-YM-20261028/blob/3456e5e75cf8d4cd558955a5c7638fba2d68f464/YM/Next/RANK/Reign.lean\#L127-L137}{\texttt{fire\_stretch}},
\href{https://github.com/yang-mills-group/OpenAI-YM-20261028/blob/3456e5e75cf8d4cd558955a5c7638fba2d68f464/YM/Next/RANK/Spirals.lean\#L127-L134}{\texttt{Spindle.hotter}},
\href{https://github.com/yang-mills-group/OpenAI-YM-20261028/blob/3456e5e75cf8d4cd558955a5c7638fba2d68f464/YM/Next/RANK/NoSpill.lean\#L98-L101}{\texttt{Paint.q\_power}},
\href{https://github.com/yang-mills-group/OpenAI-YM-20261028/blob/3456e5e75cf8d4cd558955a5c7638fba2d68f464/YM/Next/RANK/Planck.lean\#L34-L36}{\texttt{Glimmer.transports}}.

\textbf{Proposition II.2.13 (numerical conditions for the rank
comparison).} Use the scales of II.2.12, and write \(\Pi=\Pi(s)\),
\(y=y(s)\) in formulas involving both \(s\) and \(u\).

\begin{enumerate}
\def\labelenumi{(\alph{enumi})}
\item
  \emph{Shelf conditions.} Fix \(m,Q\in\mathbb N\), \(e,U,V\ge0\) and
  \(a>0\). Put \(A_Q=(2Q+1)^4\) and \(G(u)=e^{-a h(u)^2}\). For all
  sufficiently large \(s\), the following ten inequalities hold
  simultaneously for every \(u\in[s/2,s]\): \[\begin{aligned}
  8&\le W(u),\\
  W(u)^2+2QW(u)&\le M(s)^{10},\\
  64(mU+1)y^2&\le W(u),\\
  64(mU+4\cdot2^m)y^2&\le W(u),\\
  2^m\Pi^{-25/32}+\Pi^{-1}+\Pi^{-2}&\le\Pi^{-3/4},\\
  e^{mUy^2+4\cdot2^m y^2+2j(u)}G(u)+2^m\Pi^{-3}&\le\Pi^{-1},\\
  (2W(u)+1)^4\Pi^{-1}&\le1,\\
  e^{eA_Q}\Pi^{1/8}\Pi^{-31/32}&\le\Pi^{-25/32},\\
  e^{(e+1)A_Q}\Pi^{1/8}\Pi^{1/4}&\le\Pi,\\
  e^{-2j(u)}(e^{\Pi^{-1}}-1)&\le\Pi^{-3/4}\omega(u).
  \end{aligned}\] The fifth line uses \(e^{mVy^2}\Pi^{-mV-2}=\Pi^{-2}\).
  The seventh is the geometric hypothesis of II.2.6 with cover size
  \(W(u)\) and hit bound \(\Pi^{-1}\).
\item
  \emph{Tally conditions.} For the same data, fix \(\alpha\ge1\) and
  \(L>1\), and put
  \[n(L,s)=\left\lceil\frac{10\log M(s)}{\log L}\right\rceil.\] For all
  sufficiently large \(s\),
  \[(\alpha+1)^{n(L,s)}\Pi(s)^{-3/4}\le\Pi(s)^{-1/2},\] and (a) holds
  for every \(u\in[s/2,s]\). Moreover,
  \[G(u)\le\Pi(s)^{-3/4}\omega(v)\qquad(u,v\in[s/2,s]).\]
\item
  \emph{Gaussian prices.} For \(k,n\in\mathbb N\), set
  \(\phi=1/(4(nk+k+1))\). For every fixed \(A\ge1\), eventually in
  \(s\),
  \[n\Pi^{-2}(A+\Pi^{-2})^n\pi(\Pi^\phi,nk)+A^n\Pi^{-2}\pi(\Pi^\phi,nk)\le\Pi^{-3/2}.\]
  Thus the error parameters \(\xi=\Pi^{-2}\), covariance error
  \(\Pi^{-2}\), mean and covariance bound \(\Pi^\phi\), and reference
  norm bound \(A\) satisfy the polynomial comparison condition with
  error \(\Pi^{-3/2}\).
\item
  \emph{Iteration parameters.} A parameter record consists of an era
  \((Q,e,U,V,a,\alpha)\) as above, integers \(t,g,H_0\ge0\), and real
  numbers \(b,c,f,F,K\) satisfying
  \[c\ge1,\quad K\ge1,\quad F\ge0,\quad f+c\le-8,\quad g\ge6,\quad U=V,\quad f=\tfrac1{16}-U,\quad 2V+2\le g-1.\]
  Here \(t\) is the Taylor order, \(g\) the target order, \(H_0\) a
  support parameter, \(b\) an exponent, \(c\) a coefficient exponent,
  and \(F\) an error exponent. Write \(C_d=(2(d+2)K)^d\). At scale \(s\)
  and integer height \(H\), the iteration conditions are
  \[(\alpha+2)^H\le\Pi,\qquad C_t^H\le\Pi,\qquad C_{2t}^H\le\Pi,\qquad \Pi^{2c+1}+\Pi^{2c+2}\le\Pi^{2c+3},\qquad H\le\Pi.\]
  They imply, for every integer \(0\le r\le H\),
  \[(\alpha+2)^r\Pi^{-g}\le\Pi^{1-g}\le\Pi^{-4}.\] For \(b\ge0\), the
  block-size conditions
  \[L>1,\qquad \log L\ge2(b+2)\log q\quad(q\in\{\alpha+2,C_t,C_{2t}\})\]
  hold for all sufficiently large \(L\). For any such fixed \(L\), the
  iteration conditions hold for all sufficiently large \(s\) at height
  \[H(L,s)=\left\lceil\frac{(b+2)y(s)^2}{\log L}\right\rceil\]
  (\href{https://github.com/yang-mills-group/OpenAI-YM-20261028/blob/3456e5e75cf8d4cd558955a5c7638fba2d68f464/YM/Next/RANK/Odometer.lean\#L56-L61}{\texttt{racetrack}}).
\item
  \emph{Moment and fusion errors.} Define \(R_Q(z)=e^{zA_Q}\) and the
  following polynomial error function. For \(B,\varepsilon,w\ge0\) and
  \(D\in\mathbb N\), put \[\begin{aligned}
  q_0&=w\pi(B,D),& d_0&=wDB^D\varepsilon,\\
  q_1&=d_0+\varepsilon w^2\pi(B,2D)+\varepsilon+w\varepsilon\pi(B,D),\\
  q_2&=2q_0d_0+\varepsilon w^4\pi(B,4D)+2\varepsilon+w^2\varepsilon\pi(B,2D),\\
  \mathcal Q(B,\varepsilon,w,D)&=q_1+q_2+(q_0+q_1)q_1+q_1q_0.
  \end{aligned}\] For a parameter record in (d), set \(\tau(s)=3h(s)\).
  The remaining numerical conditions are \[\begin{aligned}
  1\le\tau(s)&\le\Pi-1,& \pi(\tau(s),2t)&\le\Pi,\\
  K\pi(\tau(s),1)&\le\Pi,& K+\Pi^{-F}&\le\Pi,\\
  4R_Q(e+2)\Pi^{-6}&\le\Pi^{-4},\\
  4R_Q(e+2)\mathcal Q(\Pi,\Pi^{-F},\Pi^{2c+3},2t)&\le\Pi^{-g}.
  \end{aligned}\] For every \(u\in[s/2,s]\) they also require
  \[\begin{aligned}
  \tau(u)&\le\tau(s),\\
  G(u)+2\Pi^{-2}R_Q(e+3)&\le\Pi^{-1},\\
  (2W(u)+1)^4\Pi^{-1}&\le1,& e^{\Pi^{-1}}-1&\le1.
  \end{aligned}\] In the second line, \(G(u)=e^{a h(u)^2}G(u)^2\), as in
  the source definition. All these conditions hold for sufficiently
  large \(s\) whenever \[g+10Z+10\le F,\qquad Z=2c+8t+13\]
  (\href{https://github.com/yang-mills-group/OpenAI-YM-20261028/blob/3456e5e75cf8d4cd558955a5c7638fba2d68f464/YM/Next/RANK/Purchased.lean\#L49-L56}{\texttt{funding}}).
\end{enumerate}

\emph{Status:} theorems and numerical definitions. The constants
entering the block-size inequalities in (d) are fixed before \(L\). The
lower bound on \(s\) may depend on all the fixed parameters and on
\(L\). \emph{Lean:}
\href{https://github.com/yang-mills-group/OpenAI-YM-20261028/blob/3456e5e75cf8d4cd558955a5c7638fba2d68f464/YM/Next/RANK/Bird.lean\#L21-L31}{\texttt{Perch.Demand}},
\href{https://github.com/yang-mills-group/OpenAI-YM-20261028/blob/3456e5e75cf8d4cd558955a5c7638fba2d68f464/YM/Next/RANK/Bird.lean\#L39-L49}{\texttt{Shelf}},
\href{https://github.com/yang-mills-group/OpenAI-YM-20261028/blob/3456e5e75cf8d4cd558955a5c7638fba2d68f464/YM/Next/RANK/Bird.lean\#L145-L158}{\texttt{mint}},
\href{https://github.com/yang-mills-group/OpenAI-YM-20261028/blob/3456e5e75cf8d4cd558955a5c7638fba2d68f464/YM/Next/RANK/Stepped.lean\#L37-L40}{\texttt{Era}},
\href{https://github.com/yang-mills-group/OpenAI-YM-20261028/blob/3456e5e75cf8d4cd558955a5c7638fba2d68f464/YM/Next/RANK/Stepped.lean\#L55-L58}{\texttt{Tally}},
\href{https://github.com/yang-mills-group/OpenAI-YM-20261028/blob/3456e5e75cf8d4cd558955a5c7638fba2d68f464/YM/Next/RANK/Stepped.lean\#L81-L89}{\texttt{on\_time}},
\href{https://github.com/yang-mills-group/OpenAI-YM-20261028/blob/3456e5e75cf8d4cd558955a5c7638fba2d68f464/YM/Next/RANK/Planck.lean\#L21}{\texttt{Glimmer.fuel}},
\href{https://github.com/yang-mills-group/OpenAI-YM-20261028/blob/3456e5e75cf8d4cd558955a5c7638fba2d68f464/YM/Next/RANK/Planck.lean\#L39-L62}{\texttt{collected}},
\href{https://github.com/yang-mills-group/OpenAI-YM-20261028/blob/3456e5e75cf8d4cd558955a5c7638fba2d68f464/YM/Next/RANK/Planck.lean\#L68-L69}{\texttt{shed}},
\href{https://github.com/yang-mills-group/OpenAI-YM-20261028/blob/3456e5e75cf8d4cd558955a5c7638fba2d68f464/YM/Next/RANK/Ratios.lean\#L20-L38}{\texttt{Tapes.Purse}},
\href{https://github.com/yang-mills-group/OpenAI-YM-20261028/blob/3456e5e75cf8d4cd558955a5c7638fba2d68f464/YM/Next/RANK/Joule.lean\#L41}{\texttt{Purse.budget}},
\href{https://github.com/yang-mills-group/OpenAI-YM-20261028/blob/3456e5e75cf8d4cd558955a5c7638fba2d68f464/YM/Next/RANK/Ratios.lean\#L64-L69}{\texttt{Racer}},
\href{https://github.com/yang-mills-group/OpenAI-YM-20261028/blob/3456e5e75cf8d4cd558955a5c7638fba2d68f464/YM/Next/RANK/Ratios.lean\#L113-L119}{\texttt{Racer.pos}},
\href{https://github.com/yang-mills-group/OpenAI-YM-20261028/blob/3456e5e75cf8d4cd558955a5c7638fba2d68f464/YM/Next/RANK/Ratios.lean\#L92-L110}{\texttt{Faring}},
\href{https://github.com/yang-mills-group/OpenAI-YM-20261028/blob/3456e5e75cf8d4cd558955a5c7638fba2d68f464/YM/Next/RANK/Ratios.lean\#L61}{\texttt{torque}},
\href{https://github.com/yang-mills-group/OpenAI-YM-20261028/blob/3456e5e75cf8d4cd558955a5c7638fba2d68f464/YM/Next/RANK/Odometer.lean\#L30-L35}{\texttt{Digits}},
\href{https://github.com/yang-mills-group/OpenAI-YM-20261028/blob/3456e5e75cf8d4cd558955a5c7638fba2d68f464/YM/Next/RANK/Odometer.lean\#L41-L42}{\texttt{pulse\_ok}},
\href{https://github.com/yang-mills-group/OpenAI-YM-20261028/blob/3456e5e75cf8d4cd558955a5c7638fba2d68f464/YM/Next/RANK/Prepaid.lean\#L37-L47}{\texttt{Panel}},
\href{https://github.com/yang-mills-group/OpenAI-YM-20261028/blob/3456e5e75cf8d4cd558955a5c7638fba2d68f464/YM/Next/RANK/Canal.lean\#L30}{\texttt{rentQ}},
\href{https://github.com/yang-mills-group/OpenAI-YM-20261028/blob/3456e5e75cf8d4cd558955a5c7638fba2d68f464/YM/Next/RANK/LedgerGate.lean\#L101}{\texttt{load}},
\href{https://github.com/yang-mills-group/OpenAI-YM-20261028/blob/3456e5e75cf8d4cd558955a5c7638fba2d68f464/YM/Next/RANK/Cabal.lean\#L38}{\texttt{capZ}},
\href{https://github.com/yang-mills-group/OpenAI-YM-20261028/blob/3456e5e75cf8d4cd558955a5c7638fba2d68f464/YM/Next/RANK/Cabal.lean\#L39}{\texttt{defect}},
\href{https://github.com/yang-mills-group/OpenAI-YM-20261028/blob/3456e5e75cf8d4cd558955a5c7638fba2d68f464/YM/Next/RANK/Combing.lean\#L51}{\texttt{grant₁}},
\href{https://github.com/yang-mills-group/OpenAI-YM-20261028/blob/3456e5e75cf8d4cd558955a5c7638fba2d68f464/YM/Next/RANK/Combing.lean\#L52}{\texttt{grant₂}},
\href{https://github.com/yang-mills-group/OpenAI-YM-20261028/blob/3456e5e75cf8d4cd558955a5c7638fba2d68f464/YM/Next/RANK/Festoon.lean\#L46-L49}{\texttt{grant}},
\href{https://github.com/yang-mills-group/OpenAI-YM-20261028/blob/3456e5e75cf8d4cd558955a5c7638fba2d68f464/YM/Next/RANK/Fleet.lean\#L62}{\texttt{Paint.trib}},
\href{https://github.com/yang-mills-group/OpenAI-YM-20261028/blob/3456e5e75cf8d4cd558955a5c7638fba2d68f464/YM/Next/RANK/Station.lean\#L38}{\texttt{Station.more}},
\href{https://github.com/yang-mills-group/OpenAI-YM-20261028/blob/3456e5e75cf8d4cd558955a5c7638fba2d68f464/YM/Next/RANK/Station.lean\#L37}{\texttt{Station.bump}}.

\subsubsection{2. What it uses}\label{what-it-uses-2}

\begin{itemize}
\tightlist
\item
  \textbf{III.1, activity format.}

  \begin{itemize}
  \tightlist
  \item
    \href{https://github.com/yang-mills-group/OpenAI-YM-20261028/blob/3456e5e75cf8d4cd558955a5c7638fba2d68f464/YM/RG/Activities.lean\#L62-L66}{\texttt{Activities}}
    consists of regular addresses (a site belonging to the support),
    \href{https://github.com/yang-mills-group/OpenAI-YM-20261028/blob/3456e5e75cf8d4cd558955a5c7638fba2d68f464/YM/RG/Activities.lean\#L42-L47}{forced
    addresses} (support, inventory, extra test) and coefficient
    functions.
  \item
    \href{https://github.com/yang-mills-group/OpenAI-YM-20261028/blob/3456e5e75cf8d4cd558955a5c7638fba2d68f464/YM/RG/Activities.lean\#L54-L55}{Eligibility}
    of a forced address: support ∩ large-field set = inventory, and the
    extra test holds.
  \item
    \href{https://github.com/yang-mills-group/OpenAI-YM-20261028/blob/3456e5e75cf8d4cd558955a5c7638fba2d68f464/YM/RG/Activities.lean\#L106-L111}{\texttt{Valid}}
    requires regularity and vanishing conditions for the regular
    coefficients, gauge invariance, measurability, final-field
    dependence and eligibility for the forced coefficients, affine
    covariance, and the prescribed regular and forced bounds. The
    K-bound \(v\) requires nonnegative envelopes
    \(b_i\ge\sup_U|K_i(U)|\) whose masses have a summable hit bound
    \(v\) at each site. The
    \href{https://github.com/yang-mills-group/OpenAI-YM-20261028/blob/3456e5e75cf8d4cd558955a5c7638fba2d68f464/YM/RG/Activities.lean\#L129}{gas}
    sums products of forced coefficients over finite selections with
    disjoint supports and inventories whose union is exactly the
    large-field set.
  \item
    \href{https://github.com/yang-mills-group/OpenAI-YM-20261028/blob/3456e5e75cf8d4cd558955a5c7638fba2d68f464/YM/RG/Activities.lean\#L59-L60}{\texttt{SumBudget}}
    (nonnegative, summable, sum ≤ bound).
  \item
    \href{https://github.com/yang-mills-group/OpenAI-YM-20261028/blob/3456e5e75cf8d4cd558955a5c7638fba2d68f464/YM/RG/Masks.lean\#L125-L136}{\texttt{BudgetPoint}}
    (order, power, support exponent, cover size \(B\), threshold).
  \item
    The large-field set
    \href{https://github.com/yang-mills-group/OpenAI-YM-20261028/blob/3456e5e75cf8d4cd558955a5c7638fba2d68f464/YM/RG/Masks.lean\#L34-L36}{\texttt{flags}},
    \href{https://github.com/yang-mills-group/OpenAI-YM-20261028/blob/3456e5e75cf8d4cd558955a5c7638fba2d68f464/YM/RG/Masks.lean\#L41}{\texttt{Healthy}},
    \href{https://github.com/yang-mills-group/OpenAI-YM-20261028/blob/3456e5e75cf8d4cd558955a5c7638fba2d68f464/YM/RG/Masks.lean\#L84-L85}{\texttt{UsesFinal}},
    \href{https://github.com/yang-mills-group/OpenAI-YM-20261028/blob/3456e5e75cf8d4cd558955a5c7638fba2d68f464/YM/RG/ActivityBounds.lean\#L12-L13}{\texttt{Dominated}}.
  \item
    II.2.8 uses this format; II.2.6 uses its notion of a nonnegative
    summable bound.
  \end{itemize}
\item
  \textbf{I.1, lattice.} Sites, supports, separation,
  \href{https://github.com/yang-mills-group/OpenAI-YM-20261028/blob/3456e5e75cf8d4cd558955a5c7638fba2d68f464/YM/Lattice/Support.lean\#L43-L44}{\texttt{halo}},
  cover numbers, and the cardinality and intersecting-union bounds
  quoted at the head of §1. The digest condition also uses gauge
  invariance and lattice affine maps.
\item
  \textbf{Standard mathematics.}

  \begin{itemize}
  \tightlist
  \item
    Polymer gases and the Mayer expansion; the Kotecký--Preiss criterion
    and its inductive proof (R. Kotecký, D. Preiss, \emph{Comm. Math.
    Phys.} 103 (1986); R. L. Dobrushin, \emph{AMS Transl.} 177 (1996));
    D. Brydges, \emph{A short course on cluster expansions} (Les Houches
    1984); Glimm--Jaffe, \emph{Quantum Physics}, cluster expansions.
  \item
    Cover-number (connected block-count) weights for polymers in the
    style of Bałaban's renormalization-group papers.
  \item
    Exponential and logarithm on the nilpotent ideal of a commutative
    algebra. Connected parts as logarithms of generating functions;
    Möbius inversion on the Boolean lattice (G.-C. Rota 1964) for the
    triangular inversion in II.2.9.
  \item
    Gaussian integration by parts, Wick recursion and Wick ordering
    (Glimm--Jaffe; B. Simon, \emph{The \(P(\varphi)_2\) Euclidean
    (Quantum) Field Theory}, Ch. I). Gaussian convolution
    \(e^{\frac12\Delta_C}\) on polynomials in the renormalization group
    (D. Brydges, Park City lectures 2009; Brydges--Slade, \emph{J. Stat.
    Phys.} 159 (2015)).
  \item
    Equivalence of norms on finite-dimensional spaces (II.2.11(f)).
  \end{itemize}
\item
  \textbf{Imported bounds and assumptions.}

  \begin{itemize}
  \tightlist
  \item
    The tree bound needs \((2M+1)^4\gamma\le1\) and the exponent \(a+1\)
    in the hit condition. These are hypotheses of each application. In
    the rank comparison they come from II.2.13(a), (e) with \(M=W(u)\)
    and \(\gamma=\Pi(s)^{-1}\); in the fused gas, from the fusion bounds
    with \(M=B\) and exponent \(e_P+1\).
  \item
    Fusion requires the digest and hit bounds in II.2.8. Their startup
    hypotheses are stated in III.9; their verification in every
    subsequent application belongs to the corresponding RG step.
  \item
    Given one compatible set of stage parameters for each index, a
    sequence of numerical records can be chosen with common exponent
    \(b\ge d+6\), common Taylor order \(t\ge d\), common \(f\), and
    coefficient exponents \(c_i\ge bd+d+2\). This is the parameter
    selection used in III.4
    (\href{https://github.com/yang-mills-group/OpenAI-YM-20261028/blob/3456e5e75cf8d4cd558955a5c7638fba2d68f464/YM/Next/RANK/Prepaid.lean\#L91-L109}{\texttt{preSheet}}).
  \end{itemize}
\end{itemize}

\subsubsection{3. Strategy}\label{strategy-2}

Products of local activities are represented in a finite commutative
algebra of square-zero markers. The allowed marker sets determine which
products vanish. The exponential converts sums of nonconstant monomials
into products, and a triangular correction constructs the unique
connected polynomial whose exponential is a given polynomial with
constant term \(1\). The weighted \(\ell^1\) norm controls these
operations without a loss depending on the number of markers. Switching
connects scalar functions of activity parameters to this algebra: for
positive switch size \(x\), coefficients are recovered from values at
\(0/x\) configurations. A convergent model supplies their limit as
\(x\downarrow0\), while a uniform majorant controls that limit. The
relation \(X_v^2=xX_v\) then becomes the square-zero relation.

Connected-component decompositions reduce polymer sums to sums over
packs attached to a support. A Kotecký--Preiss induction bounds the
latter by \(e^{\mathrm{cov}_M(Y)}\). Attaching \(X_i\) costs
\(e^{\mathrm{cov}_M(X_i)}\), paid by the extra unit in the hit exponent
\(a+1\). The geometric input is \(|Y|\le(2M+1)^4\mathrm{cov}_M(Y)\). For
two interactions, mentions always connect, while walls connect only hot
labels. Hotness is unchanged on restriction to a component, so the
component decomposition can be recovered uniquely. The forced-activity
fusion uses ordinary patch components; its density identity and bounds
are proved in III.10.

The Gaussian comparison in III.11 uses algebraic convolution on
polynomials with complex symmetric covariances. Wick recursion gives
coefficient-norm and moment bounds depending on the degree and entrywise
bounds, independently of the number of variables. Gaussian integration
by parts identifies this algebraic functional with an actual expectation
when a measure is supplied. The numerical inequalities then absorb the
errors and iteration losses. There are two iteration lengths:
\(n(L,s)=O(y(s)/\log L)\) and \(H(L,s)=O(y(s)^2/\log L)\) at fixed
\(L\). The first gives losses \(\Pi(s)^{o(1)}\); the second requires the
explicit lower bounds on \(L\) in II.2.13(d). The constants in those
block-size inequalities are chosen before \(L\); all remaining constants
must be fixed before the lower bound on \(s\) is chosen.

\subsubsection{4. Steps}\label{steps-2}

\paragraph{4.1 Marker algebra}\label{marker-algebra}

\textbf{Lemma 4.1 (associativity).} \(\mathcal A_R\) is a commutative
unital ring. \emph{Sketch:}

\begin{enumerate}
\def\labelenumi{\arabic{enumi}.}
\tightlist
\item
  On basis elements, \((\delta_a\delta_b)\delta_c\neq0\) iff \(a,b,c\)
  are pairwise disjoint and \(a\cup b\cup c\in R\).
\item
  This condition is symmetric in \(a,b,c\). Downward closure ensures
  that both intermediate unions are allowed whenever the total union is
  allowed.
\item
  Extend by bilinearity.
\end{enumerate}

\emph{Lean:}
\href{https://github.com/yang-mills-group/OpenAI-YM-20261028/blob/3456e5e75cf8d4cd558955a5c7638fba2d68f464/YM/Marker/Console.lean}{\texttt{Console.lean}}.

\textbf{Lemma 4.2 (norm).} \(\|ab\|_w\le\|a\|_w\|b\|_w\). \emph{Sketch:}
\(\|c\delta_x\cdot c'\delta_y\|_w\le w^x|c|\,w^y|c'|\), with equality
when the product is nonzero, because \(w^{x\cup y}=w^xw^y\) on disjoint
unions. Expand \(a,b\) in the basis and use the triangle inequality.
\emph{Lean:}
\href{https://github.com/yang-mills-group/OpenAI-YM-20261028/blob/3456e5e75cf8d4cd558955a5c7638fba2d68f464/YM/Marker/Weight.lean\#L124-L137}{\texttt{binary}},
\href{https://github.com/yang-mills-group/OpenAI-YM-20261028/blob/3456e5e75cf8d4cd558955a5c7638fba2d68f464/YM/Marker/Weight.lean\#L139-L146}{\texttt{multiplication}}.

\textbf{Lemma 4.3 (nilpotency; norm of \(\mathrm{Exp}\)).} If
\(a(\emptyset)=0\), then \(a^k\) vanishes on sets with fewer than \(k\)
elements. For arbitrary \(a\),
\(\|\mathrm{Exp}(a)\|_w\le\sum_{k\le|V|}\|a\|_w^k/k!\le e^{\|a\|_w}\).
\emph{Sketch:} in \(a^{k+1}=a^k\cdot a\), each factor \(a\) contributes
a nonempty disjoint marker set (induction on \(k\)). The norm bound
follows from Lemma 4.2. Summable families are handled by comparing with
the sup norm (equivalent norms) for convergence, and by passing the
triangle inequality to the limit for the bound. \emph{Lean:}
\href{https://github.com/yang-mills-group/OpenAI-YM-20261028/blob/3456e5e75cf8d4cd558955a5c7638fba2d68f464/YM/Marker/Truncated.lean\#L10-L22}{\texttt{high}},
\href{https://github.com/yang-mills-group/OpenAI-YM-20261028/blob/3456e5e75cf8d4cd558955a5c7638fba2d68f464/YM/Marker/Truncated.lean\#L39-L46}{\texttt{pass\_budget}},
\href{https://github.com/yang-mills-group/OpenAI-YM-20261028/blob/3456e5e75cf8d4cd558955a5c7638fba2d68f464/YM/Marker/Truncated.lean\#L51-L58}{\texttt{budgeted}}.

\textbf{Lemma 4.4 (additivity of \(\mathrm{Exp}\)).} II.2.3(a)--(b).
\emph{Sketch:}

\begin{enumerate}
\def\labelenumi{\arabic{enumi}.}
\tightlist
\item
  Apply the binomial theorem to \((a+b)^n/n!\).
\item
  Terms \(a^ib^j\) with \(i+j>|V|\) vanish by Lemma 4.3, so the
  truncation at \(|V|\) is consistent on both sides.
\item
  For (b): \((c\delta_x)^2=0\) when \(x\neq\emptyset\), and induction
  over \(s\).
\end{enumerate}

A second proof, via the exponential of commuting nilpotent elements in
Mathlib, is in
\href{https://github.com/yang-mills-group/OpenAI-YM-20261028/blob/3456e5e75cf8d4cd558955a5c7638fba2d68f464/YM/Extra/MULTISOURCE_TRAJECTORIES/Algebra.lean}{\texttt{Trajectories.pass\_add}}.
\emph{Lean:}
\href{https://github.com/yang-mills-group/OpenAI-YM-20261028/blob/3456e5e75cf8d4cd558955a5c7638fba2d68f464/YM/Next/RANK/Exponential.lean\#L56-L88}{\texttt{Exponential.pass\_add}},
\href{https://github.com/yang-mills-group/OpenAI-YM-20261028/blob/3456e5e75cf8d4cd558955a5c7638fba2d68f464/YM/Extra/MULTISOURCE_TRAJECTORIES/Algebra.lean\#L124-L134}{\texttt{Trajectories.pass\_sum}}.

\textbf{Lemma 4.5 (logarithm by iteration).} II.2.3(c)--(d).
\emph{Sketch:} say \(u\) vanishes below order \(m\) if \(u(x)=0\) for
\(|x|<m\).

\begin{enumerate}
\def\labelenumi{\arabic{enumi}.}
\tightlist
\item
  \textbf{Key estimate.} If \(c(\emptyset)=0\) and \(u\) vanishes below
  order \(m\ge1\), then
  \(\mathrm{Exp}(c+u)-\mathrm{Exp}(c)-u=(\mathrm{Exp}(c)-1)\,u+\mathrm{Exp}(c)\,(\mathrm{Exp}(u)-1-u)\)
  vanishes below order \(m+1\).
\item
  \textbf{Existence.} With \(u=F-\mathrm{Exp}(C_k)\), the key estimate
  gives that \(F-\mathrm{Exp}(C_{k+1})\) vanishes below order \(k+2\).
  After \(|V|\) steps it vanishes identically.
\item
  \textbf{Uniqueness.} For two solutions \(C,D\), additivity gives
  \(\mathrm{Exp}(C-D)=1\). If \(\mathrm{Exp}(p)=1\) and
  \(p(\emptyset)=0\), then \(\mathrm{Exp}(p)\) agrees with \(1+p\) one
  order beyond the vanishing order of \(p\). So \(p\) vanishes to every
  order, i.e.~\(p=0\).
\item
  \textbf{(d)} follows from \(\mathrm{Exp}(-C)\,\mathrm{Exp}(C)=1\).
\end{enumerate}

\emph{Lean:}
\href{https://github.com/yang-mills-group/OpenAI-YM-20261028/blob/3456e5e75cf8d4cd558955a5c7638fba2d68f464/YM/Next/RANK/Exponential.lean\#L130-L132}{\texttt{roll}},
\href{https://github.com/yang-mills-group/OpenAI-YM-20261028/blob/3456e5e75cf8d4cd558955a5c7638fba2d68f464/YM/Next/RANK/Exponential.lean\#L133-L149}{\texttt{rolls}},
\href{https://github.com/yang-mills-group/OpenAI-YM-20261028/blob/3456e5e75cf8d4cd558955a5c7638fba2d68f464/YM/Next/RANK/Exponential.lean\#L151-L156}{\texttt{roll\_cont}},
\href{https://github.com/yang-mills-group/OpenAI-YM-20261028/blob/3456e5e75cf8d4cd558955a5c7638fba2d68f464/YM/Next/RANK/Exponential.lean\#L161-L174}{\texttt{Expands.uniq}},
\href{https://github.com/yang-mills-group/OpenAI-YM-20261028/blob/3456e5e75cf8d4cd558955a5c7638fba2d68f464/YM/Next/RANK/Exponential.lean\#L184-L185}{\texttt{correct}}.

\textbf{Lemma 4.6 (majorant arrays; first-order coefficients).}

\begin{itemize}
\tightlist
\item
  For nonnegative arrays \(A,B\) on \(R\) define the product
  \((A\ast B)(z)=\sum_{x\sqcup y=z}A(x)B(y)\). Then \(|a|\le A\) and
  \(|b|\le B\) coefficientwise imply \(|ab|\le A\ast B\), and
  \(\sum_zw^z(A\ast B)(z)\le(\sum w^xA(x))(\sum w^yB(y))\). These are
  the uniform-in-configuration majorants used in III.8.
\item
  Every \(p\) decomposes as
  \(p(\emptyset)1+\sum_vp(\{v\})\delta_{\{v\}}+(\text{terms of order}\ge2)\).
\item
  \(\mathrm{Exp}(F)\) and \(F\) have the same singleton coefficients
  when \(F(\emptyset)=0\).
\item
  For a finite family, singleton coefficients satisfy
  \[(\prod_{i\in s}a_i)(\{v\})=\sum_{i\in s}a_i(\{v\})\prod_{j\in s\setminus\{i\}}a_j(\emptyset).\]
  \emph{Sketch:} apply the triangle inequality to each disjoint-union
  coefficient sum. The weights factor on each surviving product.
  Splitting coefficients by marker-set size gives the decomposition. A
  singleton can receive a nonempty marker set from only one factor;
  powers of order at least two therefore make no singleton contribution
  to the exponential.
\end{itemize}

\emph{Lean:}
\href{https://github.com/yang-mills-group/OpenAI-YM-20261028/blob/3456e5e75cf8d4cd558955a5c7638fba2d68f464/YM/Push/GRID_APPROXIMATION/Nil.lean\#L79}{\texttt{Nil.combine}},
\href{https://github.com/yang-mills-group/OpenAI-YM-20261028/blob/3456e5e75cf8d4cd558955a5c7638fba2d68f464/YM/Push/GRID_APPROXIMATION/Nil.lean\#L93-L109}{\texttt{merged\_price}},
\href{https://github.com/yang-mills-group/OpenAI-YM-20261028/blob/3456e5e75cf8d4cd558955a5c7638fba2d68f464/YM/MergeProof/Pitch.lean\#L58-L62}{\texttt{Accent.ballots}},
\href{https://github.com/yang-mills-group/OpenAI-YM-20261028/blob/3456e5e75cf8d4cd558955a5c7638fba2d68f464/YM/MergeProof/Pitch.lean\#L35-L40}{\texttt{whisper}},
\href{https://github.com/yang-mills-group/OpenAI-YM-20261028/blob/3456e5e75cf8d4cd558955a5c7638fba2d68f464/YM/MergeProof/Pitch.lean\#L29-L33}{\texttt{shockProd}}.

\textbf{Lemma 4.7 (repacking).} II.2.4(b). \emph{Sketch:}

\begin{enumerate}
\def\labelenumi{\arabic{enumi}.}
\tightlist
\item
  Sum \(\mathrm{val}\) fibrewise over the slot set \(S=\sigma(s)\).
\item
  For admissible \(S\), the selected sets in the fibre are in bijection
  with sections \(S\ni k\mapsto i_k\in\sigma^{-1}(k)\).
\item
  The sum over sections of \(\prod_kF_{i_k}\) is the limit over finite
  truncations of \(\prod_{k\in S}\sum_{i\in\sigma^{-1}(k)}F_i\).
\item
  The majorant \(\prod_{i\in s}\|F_i\|_w\) is summable over finite sets
  \(s\) because \(\sum_i\|F_i\|_w<\infty\). II.2.2 then gives
  summability of the algebra-valued terms.
\end{enumerate}

\emph{Lean:}
\href{https://github.com/yang-mills-group/OpenAI-YM-20261028/blob/3456e5e75cf8d4cd558955a5c7638fba2d68f464/YM/Extra/MULTISOURCE_TRAJECTORIES/Packets.lean\#L114-L124}{\texttt{Packets.repack}}.

\paragraph{4.2 Components and the tree
bound}\label{components-and-the-tree-bound}

\textbf{Lemma 4.8 (components).} II.2.5. \emph{Sketch:} ``chained in
\(s\)'' is the reflexive--transitive closure of ``supports meet''
restricted to \(s\). Its classes are the parts. A chain cannot leave a
member of a family with disjoint supports, which gives item 2. For item
3, follow a chain from \(j\) to \(i\) and stop immediately before the
first visit to \(i\). \emph{Lean:}
\href{https://github.com/yang-mills-group/OpenAI-YM-20261028/blob/3456e5e75cf8d4cd558955a5c7638fba2d68f464/YM/Ac/Nodes.lean\#L79-L101}{\texttt{separates}},
\href{https://github.com/yang-mills-group/OpenAI-YM-20261028/blob/3456e5e75cf8d4cd558955a5c7638fba2d68f464/YM/Ac/Nodes.lean\#L103-L113}{\texttt{excise}}.

\textbf{Lemma 4.9 (root decomposition).} For finite \(D\subseteq\iota\)
and a support \(Y\), let \(\mathrm{film}(D,Y)=\sum\prod_{i\in s}f_i\),
the sum over packs \(s\subseteq D\) for \(Y\) (including
\(s=\emptyset\)). For \(f\ge0\),
\[\mathrm{film}(D,Y)\le\prod_{i\in D,\ X_i\cap Y\neq\emptyset}\big(1+f_i\,\mathrm{film}(D\setminus\{i\},X_i)\big).\]
\emph{Sketch:}

\begin{enumerate}
\def\labelenumi{\arabic{enumi}.}
\tightlist
\item
  Each part of a pack contains a label meeting \(Y\) (II.2.5, item 4).
  Fix one such label for each part, its \emph{root}; distinct parts have
  distinct roots.
\item
  Removing the root from its part leaves a pack for
  \(X_{\mathrm{root}}\) inside \(D\setminus\{\mathrm{root}\}\) (II.2.5,
  item 3).
\item
  The map \(s\mapsto\) (for each candidate root: nothing, or the
  remainder of the part it roots) is injective, and the weight
  factorizes.
\item
  The right-hand side sums over all such assignments with nonnegative
  terms.
\end{enumerate}

No sum over cuts or orderings occurs, and no label is repeated inside a
part. \emph{Lean:}
\href{https://github.com/yang-mills-group/OpenAI-YM-20261028/blob/3456e5e75cf8d4cd558955a5c7638fba2d68f464/YM/Ac/Table.lean\#L20}{\texttt{Ac.pouch}},
\href{https://github.com/yang-mills-group/OpenAI-YM-20261028/blob/3456e5e75cf8d4cd558955a5c7638fba2d68f464/YM/Ac/Table.lean\#L79}{\texttt{film}},
\href{https://github.com/yang-mills-group/OpenAI-YM-20261028/blob/3456e5e75cf8d4cd558955a5c7638fba2d68f464/YM/Ac/Table.lean\#L81-L106}{\texttt{lookup}}.

\textbf{Lemma 4.10 (cover number of a spark).} If \(s\) is a spark, then
\(\mathrm{cov}_M(X(s))\le\sum_{i\in s}\mathrm{cov}_M(X_i)\).
\emph{Sketch:} enumerate \(s\) so that each new label meets the union of
the previous ones (possible since \(s\) is connected), and use the
intersecting-union bound from I.1 at each step. Subadditivity is not
claimed for disjoint unions. \emph{Lean:}
\href{https://github.com/yang-mills-group/OpenAI-YM-20261028/blob/3456e5e75cf8d4cd558955a5c7638fba2d68f464/YM/Ac/Tree.lean\#L59-L82}{\texttt{Ac.cover}}.

\textbf{Lemma 4.11 (Kotecký--Preiss induction).} Under the hypotheses of
II.2.6, put \(p_i=f_ie^{a\,\mathrm{cov}_M(X_i)}\). Then for every finite
\(D\) and every support \(Y\),
\[\mathrm{film}_p(D,Y)\le e^{\mathrm{cov}_M(Y)}.\] \emph{Sketch:} strong
induction on \(D\). By Lemma 4.9 and the induction hypothesis for
\(D\setminus\{i\}\),
\[\mathrm{film}_p(D,Y)\le\prod_{i\in D,\ X_i\cap Y\neq\emptyset}\big(1+f_ie^{(a+1)\mathrm{cov}_M(X_i)}\big)\le\exp\Big(\sum_{i\in D,\ X_i\cap Y\neq\emptyset}f_ie^{(a+1)\mathrm{cov}_M(X_i)}\Big)\le e^{|Y|\gamma}.\]
The last step sums the hit condition over the sites of \(Y\). Finally
\(|Y|\gamma\le(2M+1)^4\gamma\,\mathrm{cov}_M(Y)\le\mathrm{cov}_M(Y)\).
\emph{Lean:}
\href{https://github.com/yang-mills-group/OpenAI-YM-20261028/blob/3456e5e75cf8d4cd558955a5c7638fba2d68f464/YM/Ac/Tree.lean\#L15-L27}{\texttt{demand}},
\href{https://github.com/yang-mills-group/OpenAI-YM-20261028/blob/3456e5e75cf8d4cd558955a5c7638fba2d68f464/YM/Ac/Tree.lean}{\texttt{Tree.lean}
(\texttt{step}, \texttt{all\_pack})}.

\textbf{Lemma 4.12 (tree bound).} II.2.6. \emph{Sketch:}

\begin{enumerate}
\def\labelenumi{\arabic{enumi}.}
\tightlist
\item
  Take a finite family of sparks whose supports contain \(x\), and let
  \(A\) be the union of their labels.
\item
  Each such spark is a nonempty pack for \(\{x\}\) inside \(A\).
\item
  By Lemma 4.10 and \(a\ge0\),
  \(e^{a\,\mathrm{cov}_M(X(s))}\prod_{i\in s}f_i\le\prod_{i\in s}p_i\).
\item
  The bound \(e^{|Y|\gamma}\) in the proof of Lemma 4.11 with
  \(Y=\{x\}\) (so \(|Y|=1\)) gives
  \(\mathrm{film}_p(A,\{x\})\le e^{\gamma}\).
\item
  Subtract the empty pack. Uniform bounds on finite partial sums give
  summability.
\end{enumerate}

\emph{Lean:}
\href{https://github.com/yang-mills-group/OpenAI-YM-20261028/blob/3456e5e75cf8d4cd558955a5c7638fba2d68f464/YM/Ac/Tree.lean\#L89-L127}{\texttt{Ac.forest}}.

\paragraph{4.3 Two kinds of overlap}\label{two-kinds-of-overlap}

\textbf{Lemma 4.13 (hotness is local).}

\begin{itemize}
\tightlist
\item
  Let \(p\) be a slab of \(s\) and \(i\in p\). Then \(i\) is hot in
  \(p\) iff it is hot in \(s\), and its fuel in \(p\) equals its fuel in
  \(s\).
\item
  In a mole, containing a dirty label is equivalent to all labels being
  hot.
\item
  Hotness spreads along fuel chains.
\end{itemize}

\emph{Sketch:} a mention chain is a fuel chain, so it stays inside the
slab. A wall link joins only hot labels. Moles are also invariant under
label permutations compatible with bijections of mentions and walls that
preserve dirtiness
(\href{https://github.com/yang-mills-group/OpenAI-YM-20261028/blob/3456e5e75cf8d4cd558955a5c7638fba2d68f464/YM/MergeProof/North/Harbor.lean\#L29-L53}{\texttt{Kilner.molecule}}).
\emph{Lean:}
\href{https://github.com/yang-mills-group/OpenAI-YM-20261028/blob/3456e5e75cf8d4cd558955a5c7638fba2d68f464/YM/Further/SHELL_BIRTH_49/Fusing/Coal.lean\#L149-L162}{\texttt{hot\_comp}},
\href{https://github.com/yang-mills-group/OpenAI-YM-20261028/blob/3456e5e75cf8d4cd558955a5c7638fba2d68f464/YM/Further/SHELL_BIRTH_49/Fusing/Coal.lean\#L163-L164}{\texttt{fuel\_comp}},
\href{https://github.com/yang-mills-group/OpenAI-YM-20261028/blob/3456e5e75cf8d4cd558955a5c7638fba2d68f464/YM/Further/SHELL_BIRTH_49/Fusing/Coal.lean\#L183-L187}{\texttt{takes\_mole}}.

\textbf{Lemma 4.14 (recovery).} II.2.7, items 1--3. \emph{Sketch:} let
\(s=\bigcup S\).

\begin{enumerate}
\def\labelenumi{\arabic{enumi}.}
\tightlist
\item
  A mention chain in \(s\) cannot leave a member of \(S\), since
  mention-unions are disjoint. So hotness in \(s\) equals hotness in the
  member, and fuels agree.
\item
  The fuel-unions of distinct members are disjoint: the mention
  components by hypothesis; the wall components because only members
  containing a dirty label have hot labels, and their wall-unions are
  disjoint.
\item
  Apply II.2.5, item 2 to the fuel overlap.
\item
  For slabs of a given set, use Lemma 4.13 to identify their internal
  fuels. Disjoint fuel-unions imply disjoint mention-unions; when both
  slabs contain a dirty label, all their labels are hot, so their
  wall-unions are disjoint as well.
\end{enumerate}

\emph{Lean:}
\href{https://github.com/yang-mills-group/OpenAI-YM-20261028/blob/3456e5e75cf8d4cd558955a5c7638fba2d68f464/YM/Further/SHELL_BIRTH_49/Fusing/Coal.lean\#L210-L240}{\texttt{Kiln.undo}},
\href{https://github.com/yang-mills-group/OpenAI-YM-20261028/blob/3456e5e75cf8d4cd558955a5c7638fba2d68f464/YM/Further/SHELL_BIRTH_49/Fusing/Coal.lean\#L190-L196}{\texttt{clubs\_away}},
\href{https://github.com/yang-mills-group/OpenAI-YM-20261028/blob/3456e5e75cf8d4cd558955a5c7638fba2d68f464/YM/Further/SHELL_BIRTH_49/Fusing/Coal.lean\#L197-L205}{\texttt{walls\_away}}.

\textbf{Lemma 4.15 (moles are wall-sparks).} II.2.7, item 4.
\emph{Sketch:} a fuel overlap is either a mention overlap, which implies
a wall overlap by hypothesis, or a wall overlap between hot labels.
\emph{Lean:}
\href{https://github.com/yang-mills-group/OpenAI-YM-20261028/blob/3456e5e75cf8d4cd558955a5c7638fba2d68f464/YM/Further/SHELL_BIRTH_49/Fusing/Coal.lean\#L242-L253}{\texttt{Kiln.attends}},
\href{https://github.com/yang-mills-group/OpenAI-YM-20261028/blob/3456e5e75cf8d4cd558955a5c7638fba2d68f464/YM/Further/SHELL_BIRTH_49/Fusing/Coal.lean\#L113-L117}{\texttt{Collide.spark\_ac}}.

\textbf{Lemma 4.16 (pellets).} Pellets are countably infinite.
\emph{Sketch:} every singleton \(\{\text{left }i\}\) is a pellet, and
finite subsets of \(\mathbb N\sqcup\mathbb N\) form a countable set.
\emph{Lean:}
\href{https://github.com/yang-mills-group/OpenAI-YM-20261028/blob/3456e5e75cf8d4cd558955a5c7638fba2d68f464/YM/Ovens/Coffee.lean\#L121-L133}{\texttt{Pipes.encoding}}.

\paragraph{4.4 Switching}\label{switching}

\textbf{Lemma 4.17 (triangular inversion).} For \(x>0\),
\(\mathrm{val}_x\) is injective. \emph{Sketch:} strong induction on
\(r\):
\(\mathrm{val}_x(g)(r)=x^{|r|}g(r)+\sum_{i\subsetneq r}x^{|i|}g(i)\) and
\(x^{|r|}\neq0\). This is Möbius inversion on the Boolean lattice,
scaled by \(x^{|i|}\). Hence \(\mathrm{val}_x\) is a linear bijection of
a finite-dimensional space. \emph{Lean:}
\href{https://github.com/yang-mills-group/OpenAI-YM-20261028/blob/3456e5e75cf8d4cd558955a5c7638fba2d68f464/YM/MergeProof/Switch.lean\#L33-L42}{\texttt{entire}},
\href{https://github.com/yang-mills-group/OpenAI-YM-20261028/blob/3456e5e75cf8d4cd558955a5c7638fba2d68f464/YM/MergeProof/Digits.lean}{\texttt{Digits.lean}
(\texttt{vane}, \texttt{decodes})}.

\textbf{Lemma 4.18 (jets; products).}

\begin{itemize}
\tightlist
\item
  If \(g\) models \(f\), then \(\mathrm{val}_x^{-1}(f_x)=g_x\) for
  \(0<x\le1\), so the jet is \(g_0\).
\item
  If \(g,g'\) model \(f,f'\), then \(x\mapsto g_x\ast_xg'_x\) models
  \(ff'\), since
  \(\mathrm{val}_x(p\ast_xq)=\mathrm{val}_x(p)\,\mathrm{val}_x(q)\).
\item
  Bounds: \(|p\ast_xq|\le|p|\ast_1|q|\) for \(x\in[0,1]\), and
  \(\sum(b\ast_1b')=\sum b\cdot\sum b'\).
\end{itemize}

\emph{Lean:}
\href{https://github.com/yang-mills-group/OpenAI-YM-20261028/blob/3456e5e75cf8d4cd558955a5c7638fba2d68f464/YM/MergeProof/Switch.lean\#L58-L63}{\texttt{Models.exact}},
\href{https://github.com/yang-mills-group/OpenAI-YM-20261028/blob/3456e5e75cf8d4cd558955a5c7638fba2d68f464/YM/MergeProof/SwitchProd.lean\#L46-L50}{\texttt{Paid.mul}};
the model product
\href{https://github.com/yang-mills-group/OpenAI-YM-20261028/blob/3456e5e75cf8d4cd558955a5c7638fba2d68f464/YM/MergeProof/SwitchProd.lean\#L15-L25}{\texttt{value\_comb}},
\href{https://github.com/yang-mills-group/OpenAI-YM-20261028/blob/3456e5e75cf8d4cd558955a5c7638fba2d68f464/YM/MergeProof/SwitchProd.lean\#L28-L34}{\texttt{Models.mul}}.

\textbf{Lemma 4.19 (absent markers; measurability).} II.2.9(d).
\emph{Sketch:}
\(\mathrm{val}_x(g)(r)-\mathrm{val}_x(g)(r\setminus\{v\})=\sum_{v\in i\subseteq r}x^{|i|}g(i)\).
If \(f\) does not see \(v\), this vanishes for all \(r\), so
\(g_x(i)=0\) for \(v\in i\) by induction, for every \(0<x\le1\). Let
\(x\downarrow0\). Measurability: the jet is the pointwise limit of
\(\mathrm{val}_x^{-1}(f_x)\), and \(\mathrm{val}_x^{-1}\) is a
continuous linear map. \emph{Lean:}
\href{https://github.com/yang-mills-group/OpenAI-YM-20261028/blob/3456e5e75cf8d4cd558955a5c7638fba2d68f464/YM/MergeProof/Digits.lean\#L27-L36}{\texttt{jet\_erase}},
\href{https://github.com/yang-mills-group/OpenAI-YM-20261028/blob/3456e5e75cf8d4cd558955a5c7638fba2d68f464/YM/MergeProof/Digits.lean\#L38-L49}{\texttt{Paid.jet\_meas}}.

\textbf{Lemma 4.20 (transfer bound).} II.2.9(e). \emph{Sketch:}

\begin{enumerate}
\def\labelenumi{\arabic{enumi}.}
\tightlist
\item
  The coefficient of \(\mathrm{snd}_\rho(f^u)\) at \(z\) is bounded by
  \(\sum_ib(i)\,|(\prod_{v\in i}\rho_v)(z)|\), using
  \(|g_0(i)|\le b(i)\).
\item
  The \(w\)-norm of this majorant array is at most
  \(\sum_ib(i)\,\|\prod_{v\in i}\rho_v\|_w\le\sum_ib(i)\le d\), by
  II.2.2.
\end{enumerate}

\emph{Lean:}
\href{https://github.com/yang-mills-group/OpenAI-YM-20261028/blob/3456e5e75cf8d4cd558955a5c7638fba2d68f464/YM/MergeProof/Jetty.lean\#L13-L19}{\texttt{Accent.crystal}}.

\paragraph{4.5 Gaussian calculus}\label{gaussian-calculus}

\textbf{Lemma 4.21 (construction of \(\mathcal G_Q\)).} II.2.10(a)--(b).
\emph{Sketch:}

\begin{enumerate}
\def\labelenumi{\arabic{enumi}.}
\tightlist
\item
  \([X_v+\partial_{Q_v},X_w+\partial_{Q_w}]=Q_{vw}-Q_{wv}=0\) by
  symmetry, and derivations commute.
\item
  The commutative subalgebra they generate receives an algebra map from
  \(\mathbb C[X]\); apply it to \(1\).
\item
  The semigroup law and commutation with \(\partial_b\) follow by
  induction on polynomials from the recursion. Applying evaluation at
  \(m\) gives the moment functional and its Wick recursion. The
  terminating series \(e^{\Delta_Q/2}\) satisfies the same polynomial
  recursion, giving the equivalent heat-operator formula.
\end{enumerate}

\emph{Lean:}
\href{https://github.com/yang-mills-group/OpenAI-YM-20261028/blob/3456e5e75cf8d4cd558955a5c7638fba2d68f464/YM/Next/RANK/Gaussian.lean\#L64-L65}{\texttt{cream\_comm}},
\href{https://github.com/yang-mills-group/OpenAI-YM-20261028/blob/3456e5e75cf8d4cd558955a5c7638fba2d68f464/YM/Next/RANK/Gaussian.lean\#L73-L74}{\texttt{make}},
\href{https://github.com/yang-mills-group/OpenAI-YM-20261028/blob/3456e5e75cf8d4cd558955a5c7638fba2d68f464/YM/Next/RANK/Gaussian.lean\#L82}{\texttt{feeds}},
\href{https://github.com/yang-mills-group/OpenAI-YM-20261028/blob/3456e5e75cf8d4cd558955a5c7638fba2d68f464/YM/Next/RANK/Gaussian.lean\#L96-L102}{\texttt{warms}},
\href{https://github.com/yang-mills-group/OpenAI-YM-20261028/blob/3456e5e75cf8d4cd558955a5c7638fba2d68f464/YM/Next/RANK/GaussTrack.lean\#L31-L35}{\texttt{spring}}.

\textbf{Lemma 4.22 (size of moments).} II.2.10(e). \emph{Sketch:}

\begin{enumerate}
\def\labelenumi{\arabic{enumi}.}
\tightlist
\item
  By linearity it suffices to treat monomials, each written as
  \(\prod_{j\in s}X_{z_j}\) with \(|s|\le d\). The coefficient sum gives
  the factor \(\|p\|_1\).
\item
  Let \(c_d\) bound the moments of degree \(\le d\). The Wick recursion
  gives \(c_{d+1}\le B\,c_d+d\,B\,c_d\), and \(c_d\le(2B(d+1))^d\) by
  induction.
\item
  For differences, apply \(|xv-yw|\le|x||v-w|+|x-y||w|\) to each term of
  the recursion.
\item
  Localization: restrict \(Q\) and \(m\) to the variables of \(p\) by a
  change of variables.
\end{enumerate}

\emph{Lean:}
\href{https://github.com/yang-mills-group/OpenAI-YM-20261028/blob/3456e5e75cf8d4cd558955a5c7638fba2d68f464/YM/Next/RANK/Common.lean}{\texttt{Common.lean}},
\href{https://github.com/yang-mills-group/OpenAI-YM-20261028/blob/3456e5e75cf8d4cd558955a5c7638fba2d68f464/YM/Next/RANK/Common.lean\#L121-L123}{\texttt{width}},
\href{https://github.com/yang-mills-group/OpenAI-YM-20261028/blob/3456e5e75cf8d4cd558955a5c7638fba2d68f464/YM/Next/RANK/Common.lean\#L124-L126}{\texttt{differs}},
\href{https://github.com/yang-mills-group/OpenAI-YM-20261028/blob/3456e5e75cf8d4cd558955a5c7638fba2d68f464/YM/Next/RANK/Guard.lean}{\texttt{Guard.lean}
(\texttt{pinned})}.

\textbf{Lemma 4.23 (change of variables; factorization).} II.2.10(c) and
II.2.11(c), last clause. \emph{Sketch:} both sides of II.2.10(c) satisfy
the same recursion in \(X_v\). One uses
\(\partial_m(p\circ A)=(\partial_{Am}p)\circ A\) and the fact that
\(\mathcal G_Q\) applied to \((A\text{-linear form})\cdot q\) obeys the
Wick rule with covariance \(AQA^T\). If \(Q\) vanishes between the
variables of \(P\) and \(R\), then
\(\partial_{Q_{v\cdot}}\mathcal G_QR=0\) for every variable \(v\) of
\(P\). The Leibniz rule then gives
\(\mathcal G_Q(PR)=(\mathcal G_QP)(\mathcal G_QR)\); linear substitution
preserves this identity. Related results:

\begin{itemize}
\tightlist
\item
  if \(Q\) is block diagonal, Gaussian expectations of products over
  blocks factorize
  (\href{https://github.com/yang-mills-group/OpenAI-YM-20261028/blob/3456e5e75cf8d4cd558955a5c7638fba2d68f464/YM/Next/RANK/Crust.lean\#L59}{\texttt{bridge}};
  \href{https://github.com/yang-mills-group/OpenAI-YM-20261028/blob/3456e5e75cf8d4cd558955a5c7638fba2d68f464/YM/Next/RANK/BoundParticle.lean\#L29-L36}{\texttt{Particle.tide}}
  for marginals);
\item
  applying Wick ordering separately at each vertex subtracts its chosen
  covariance from the diagonal block
  (\href{https://github.com/yang-mills-group/OpenAI-YM-20261028/blob/3456e5e75cf8d4cd558955a5c7638fba2d68f464/YM/Next/RANK/GaussTags.lean\#L106-L109}{\texttt{lodge}}),
  and two steps compose as \(T_{P,R}T_{H,Q}=T_{HP,\,HRH^T+Q}\)
  (\href{https://github.com/yang-mills-group/OpenAI-YM-20261028/blob/3456e5e75cf8d4cd558955a5c7638fba2d68f464/YM/Next/RANK/GaussTags.lean\#L115-L122}{\texttt{piped}}).
\end{itemize}

\emph{Lean:}
\href{https://github.com/yang-mills-group/OpenAI-YM-20261028/blob/3456e5e75cf8d4cd558955a5c7638fba2d68f464/YM/Next/RANK/GaussSubst.lean\#L76-L81}{\texttt{prepared}},
\href{https://github.com/yang-mills-group/OpenAI-YM-20261028/blob/3456e5e75cf8d4cd558955a5c7638fba2d68f464/YM/Next/RANK/GaussSubst.lean\#L83-L84}{\texttt{wets}}.

\textbf{Lemma 4.24 (identification with a true expectation).}
II.2.10(d). \emph{Sketch:}

\begin{enumerate}
\def\labelenumi{\arabic{enumi}.}
\tightlist
\item
  Substituting \(X\mapsto X+m\) turns integration by parts into the
  recursion
  \(\mathbb E[(m+\xi)_vq(m+\xi)]=m_v\mathbb E[q(m+\xi)]+\mathbb E[(\partial_{Q_v}q)(m+\xi)]\).
\item
  Both \(q\mapsto\mathbb E[q(m+\xi)]\) and \(\mathbb E_{Q,m}\) are
  linear, equal \(1\) on \(1\), and satisfy the Wick recursion on
  monomials.
\item
  Induction on the degree.
\end{enumerate}

\emph{Lean:}
\href{https://github.com/yang-mills-group/OpenAI-YM-20261028/blob/3456e5e75cf8d4cd558955a5c7638fba2d68f464/YM/Next/RANK/ParticleMoment.lean\#L61-L78}{\texttt{payment}}.

\textbf{Lemma 4.25 (step bounds).} II.2.11(a)--(d). \emph{Sketch:}

\begin{itemize}
\tightlist
\item
  \textbf{(a)}: expand in monomials. Submultiplicativity bounds the
  image of a degree-\(k\) monomial by \(c^k\le c^{\deg p}\).
\item
  \textbf{(b)}: induction on the degree, using
  \(\mathcal G_Q(X_vq)-X_vq=X_v(\mathcal G_Qq-q)+\sum_jQ_{vj}\,\mathcal G_Q(\partial_jq)\).
\item
  \textbf{(c)}: the same recursion with \(X_v\) replaced by the linear
  form \(H_v\); each step multiplies the bound by at most
  \(C(1+d)\le2(d+2)C\). For the difference bound, subtract the two
  recursions and split each product difference into the change in one
  factor at a time.
\item
  \textbf{(d)}: write
  \(\Phi_t(P)-(\mathcal G_KP)\circ f=[\Phi_t(P)-P\circ f]+[(P-\mathcal G_KP)\circ f]\).

  \begin{itemize}
  \tightlist
  \item
    First term: \(\|t^{-1}A_v(tX)\|_1\le tB\), since \(A_v\) has no
    terms of degree \(<2\), and
    \(\|U(p)-T(p)\|_1\le\|p\|_1\,d\,B'A^{d}\) for algebra maps that are
    \(B'\)-close on generators (bounded by \(A\)).
  \item
    Second term: (b) followed by (a).
  \end{itemize}
\end{itemize}

\emph{Lean:}
\href{https://github.com/yang-mills-group/OpenAI-YM-20261028/blob/3456e5e75cf8d4cd558955a5c7638fba2d68f464/YM/Next/RANK/Rise.lean\#L51-L66}{\texttt{Rise.rise}},
\href{https://github.com/yang-mills-group/OpenAI-YM-20261028/blob/3456e5e75cf8d4cd558955a5c7638fba2d68f464/YM/Next/RANK/Ample.lean\#L35-L77}{\texttt{LiftMark.puff}},
\href{https://github.com/yang-mills-group/OpenAI-YM-20261028/blob/3456e5e75cf8d4cd558955a5c7638fba2d68f464/YM/Next/RANK/Ample.lean\#L79-L96}{\texttt{LiftMark.bounds}},
\href{https://github.com/yang-mills-group/OpenAI-YM-20261028/blob/3456e5e75cf8d4cd558955a5c7638fba2d68f464/YM/Next/RANK/Herald.lean\#L119-L124}{\texttt{crux}},
\href{https://github.com/yang-mills-group/OpenAI-YM-20261028/blob/3456e5e75cf8d4cd558955a5c7638fba2d68f464/YM/Next/RANK/Instructions.lean\#L29-L53}{\texttt{LiftMark.message}},
\href{https://github.com/yang-mills-group/OpenAI-YM-20261028/blob/3456e5e75cf8d4cd558955a5c7638fba2d68f464/YM/Next/RANK/Pin.lean\#L66-L84}{\texttt{LiftMark.stock}}.

\textbf{Lemma 4.26 (guards; second moment).} II.2.11(e). \emph{Sketch:}

\begin{itemize}
\tightlist
\item
  \(1-gz\le(1-g)+(1-z)\) for \(g,z\in[0,1]\).
\item
  \(|gf-f|=(1-g)|f|\le\varepsilon'|f|^2+(1-g)/\varepsilon'\), using
  \(0\le1-g\le1\); integrate.
\item
  Second moment: \(|p(m+\xi)|^2=(p\bar p)(m+\xi)\), where \(\bar p\)
  conjugates the coefficients. Apply II.2.10(d)--(e) to \(p\bar p\),
  which has degree at most \(2d\) and \(\|p\bar p\|_1\le\|p\|_1^2\).
\end{itemize}

\emph{Lean:}
\href{https://github.com/yang-mills-group/OpenAI-YM-20261028/blob/3456e5e75cf8d4cd558955a5c7638fba2d68f464/YM/Next/RANK/Guard.lean\#L79-L88}{\texttt{guards\_mul}},
\href{https://github.com/yang-mills-group/OpenAI-YM-20261028/blob/3456e5e75cf8d4cd558955a5c7638fba2d68f464/YM/Next/RANK/Guard.lean\#L108-L125}{\texttt{escaped}},
\href{https://github.com/yang-mills-group/OpenAI-YM-20261028/blob/3456e5e75cf8d4cd558955a5c7638fba2d68f464/YM/Next/RANK/Guard.lean\#L90-L105}{\texttt{squareFits}}.

\textbf{Lemma 4.27 (reading coefficients).} II.2.11(f). \emph{Sketch:}

\begin{enumerate}
\def\labelenumi{\arabic{enumi}.}
\tightlist
\item
  Evaluation at all points of the cube is injective on polynomials of
  degree \(\le d\), since a polynomial vanishing on a product of
  infinite sets is zero.
\item
  In finite dimension, a finite subfamily of evaluations is already
  injective: minimize the kernel dimension.
\item
  Take a left inverse and bound its \(\ell^1\) operator norm.
\end{enumerate}

The Lipschitz bound follows by telescoping over monomials. A version
with explicit dimension dependence and constant
\((|V|+1)^dT_{d,\varepsilon}\) is proved in III.11
(\href{https://github.com/yang-mills-group/OpenAI-YM-20261028/blob/3456e5e75cf8d4cd558955a5c7638fba2d68f464/YM/Next/RANK/SpectrumReal.lean\#L145-L156}{\texttt{ball\_count}}).
For polynomials of degree at most \(n\), an \(o(t^n)\) difference at
\(0\) forces equality; for \(C^n\) functions it forces equality of their
Taylor polynomials through order \(n\)
(\href{https://github.com/yang-mills-group/OpenAI-YM-20261028/blob/3456e5e75cf8d4cd558955a5c7638fba2d68f464/YM/Next/RANK/Peano.lean\#L99-L102}{\texttt{Spindle.decipher}},
\href{https://github.com/yang-mills-group/OpenAI-YM-20261028/blob/3456e5e75cf8d4cd558955a5c7638fba2d68f464/YM/Next/RANK/SolveRay.lean\#L160-L167}{\texttt{trill}}).
\emph{Lean:}
\href{https://github.com/yang-mills-group/OpenAI-YM-20261028/blob/3456e5e75cf8d4cd558955a5c7638fba2d68f464/YM/Next/RANK/SpectrumReal.lean\#L22-L28}{\texttt{listen}},
\href{https://github.com/yang-mills-group/OpenAI-YM-20261028/blob/3456e5e75cf8d4cd558955a5c7638fba2d68f464/YM/Next/RANK/SpectrumReal.lean\#L56-L81}{\texttt{molds}},
\href{https://github.com/yang-mills-group/OpenAI-YM-20261028/blob/3456e5e75cf8d4cd558955a5c7638fba2d68f464/YM/Next/RANK/Spirals.lean\#L20-L41}{\texttt{amove}}.

\paragraph{4.6 Numbers}\label{numbers}

\textbf{Lemma 4.28 (loop bounds).} II.2.12, items 1--2. \emph{Sketch:}

\begin{enumerate}
\def\labelenumi{\arabic{enumi}.}
\tightlist
\item
  \((\log\log s)^2=o(\log s)\), so \(b\,y^2\le p\log s\) eventually;
  this is item 1.
\item
  Then \(s\theta^2=\Pi^2\) exactly.
\item
  \(s\theta^3=\Pi^3s^{-1/2}\le1\) by item 1 with \(b=3\),
  \(p=\tfrac12\).
\item
  Item 1 with \(b=1\), \(p=1/2\) gives \(\theta\le1\) eventually, hence
  \(s\theta^4\le s\theta^3\).
\item
  \(s\cdot2\Pi^v\theta^2\le2\Pi^3\le2\sqrt s\le s\) for \(s\ge4\).
\end{enumerate}

\emph{Lean:}
\href{https://github.com/yang-mills-group/OpenAI-YM-20261028/blob/3456e5e75cf8d4cd558955a5c7638fba2d68f464/YM/Vac/Treasury.lean\#L74-L78}{\texttt{cost}},
\href{https://github.com/yang-mills-group/OpenAI-YM-20261028/blob/3456e5e75cf8d4cd558955a5c7638fba2d68f464/YM/Vac/Treasury.lean\#L83-L108}{\texttt{uniform\_loop}}.

\textbf{Lemma 4.29 (comparison of scales).} II.2.12, items 3--4.
\emph{Sketch:}

\begin{itemize}
\tightlist
\item
  \(\log W(s)\le6y(s)\) eventually, so fixed powers of \(W\) are bounded
  by \(\Pi^a\). The ceiling inequalities
  \((\log s)^2\le M(s)<(\log s)^2+1\) also give \(W(s)\sim(\log s)^4\).
  Also \(\log h=y^{3/2}=o(y^2)\).
\item
  If \(a<b\), then \(c\Pi^{a-b}\to0\), proving item 4.
\item
  \(y(s)\le y(s/2)+1\), since \(\log s=\log(s/2)+\log2\).
\item
  \(y^2e^{-y}\to0\) handles \(y(s)^2\le a\,W(s/2)\) and
  \(y(s)^2\le a\,j(s/2)\); for the latter use
  \(j(s/2)=e^{y(s/2)^{3/2}/5}\ge e^{y(s/2)}\) once \(y(s/2)\ge25\).
\item
  Eventually \(y(s)\le2y(s/2)\); then \(2^{3/2}\le5\) gives
  \(j(s)\le h(s/2)\).
\end{itemize}

\emph{Lean:}
\href{https://github.com/yang-mills-group/OpenAI-YM-20261028/blob/3456e5e75cf8d4cd558955a5c7638fba2d68f464/YM/Next/RANK/Reign.lean\#L71-L90}{\texttt{rate\_wall}},
\href{https://github.com/yang-mills-group/OpenAI-YM-20261028/blob/3456e5e75cf8d4cd558955a5c7638fba2d68f464/YM/Next/RANK/Reign.lean\#L94-L108}{\texttt{rate\_inverse}},
\href{https://github.com/yang-mills-group/OpenAI-YM-20261028/blob/3456e5e75cf8d4cd558955a5c7638fba2d68f464/YM/Next/RANK/Reign.lean\#L110-L122}{\texttt{rate\_fire}},
\href{https://github.com/yang-mills-group/OpenAI-YM-20261028/blob/3456e5e75cf8d4cd558955a5c7638fba2d68f464/YM/Next/RANK/Reign.lean\#L127-L137}{\texttt{fire\_stretch}}.

\textbf{Lemma 4.30 (feasibility of the shelf and tally conditions).}
II.2.13(a)--(b). \emph{Sketch:}

\begin{enumerate}
\def\labelenumi{\arabic{enumi}.}
\tightlist
\item
  Use \(y(s)^2\ll W(s/2)\) for the two linear lower bounds on \(W(u)\).
  Since \(W(u)\le W(s)=M(s)^2\), the support-size inequality follows
  once \(M(s)\) is large compared with \(Q\). Fixed powers of \(W\) are
  \(o(\Pi)\), giving the geometric condition.
\item
  The purely exponential terms reduce to fixed constants times powers of
  \(\Pi\) with strictly smaller exponents than the target; for example,
  \(1/8-31/32=-27/32<-25/32\). For the sixth shelf inequality, uniformly
  in \(u\in[s/2,s]\),
  \[\big(mU+4\cdot2^m+2\big)y(s)^2+2j(s)\le a h(u)^2\] eventually. Hence
  its first term is at most \(\Pi^{-2}\). The reserve inequality follows
  from \(j(s/2)\gg y(s)^2\) and \(e^{\Pi^{-1}}-1\le\exp(1)\).
\item
  The ceiling formula gives \(n(L,s)\log(\alpha+1)=O(y(s))\le y(s)^2/4\)
  eventually. The same uniform tail estimate, with coefficients \(3/4\)
  and \(7/10\), gives \(G(u)\le\Pi^{-3/4}\omega(v)\).
\item
  Uniformity uses \(W(s/2)\le W(u)\le W(s)\), \(h(u)\ge h(s/2)\) and
  \(j(u)\le j(s)\). Intersect the finitely many eventual conditions.
\end{enumerate}

\emph{Lean:}
\href{https://github.com/yang-mills-group/OpenAI-YM-20261028/blob/3456e5e75cf8d4cd558955a5c7638fba2d68f464/YM/Next/RANK/Bird.lean\#L145-L158}{\texttt{mint}},
\href{https://github.com/yang-mills-group/OpenAI-YM-20261028/blob/3456e5e75cf8d4cd558955a5c7638fba2d68f464/YM/Next/RANK/Bird.lean\#L67-L77}{\texttt{force}},
\href{https://github.com/yang-mills-group/OpenAI-YM-20261028/blob/3456e5e75cf8d4cd558955a5c7638fba2d68f464/YM/Next/RANK/Stepped.lean\#L81-L89}{\texttt{on\_time}}.

\textbf{Lemma 4.31 (numerical parameter bounds).} II.2.13(c)--(e).
\emph{Sketch:}

\begin{itemize}
\tightlist
\item
  The block-size conditions are finitely many lower bounds on
  \(\log L\). At height \(H(L,s)\), each logarithmic loss is at most
  \(y(s)^2\) for large \(s\).
\item
  For \(r\le H\), \((\alpha+2)^r\le(\alpha+2)^H\le\Pi\); combine this
  with \(g\ge6\).
\item
  For (c), \(\pi(\Pi^\phi,nk)\le(2(nk+1))^{nk}\Pi^{1/4}\) and
  \(\Pi^{-2}\le1\). The left side is therefore at most a fixed constant
  times \(\Pi^{-7/4}\), which is eventually at most \(\Pi^{-3/2}\).
\item
  For (e), fixed-degree moments of \(\tau=3h\) are bounded by
  arbitrarily small positive powers of \(\Pi\). With \(Z=2c+8t+13\), the
  error function in its sixth inequality is \(O(\Pi^{-F+6Z})\). The
  condition \(g+10Z+10\le F\) makes this at most \(\Pi^{-g}\) after
  multiplication by the fixed factor. The remaining tail bound uses
  \(G(u)\le\Pi^{-2}\) uniformly on \([s/2,s]\), and the geometric bound
  follows from \(W(s)^4\ll\Pi\).
\end{itemize}

\emph{Lean:}
\href{https://github.com/yang-mills-group/OpenAI-YM-20261028/blob/3456e5e75cf8d4cd558955a5c7638fba2d68f464/YM/Next/RANK/Odometer.lean\#L41-L42}{\texttt{pulse\_ok}},
\href{https://github.com/yang-mills-group/OpenAI-YM-20261028/blob/3456e5e75cf8d4cd558955a5c7638fba2d68f464/YM/Next/RANK/Ratios.lean\#L113-L119}{\texttt{Racer.pos}},
\href{https://github.com/yang-mills-group/OpenAI-YM-20261028/blob/3456e5e75cf8d4cd558955a5c7638fba2d68f464/YM/Next/RANK/Planck.lean\#L39-L62}{\texttt{collected}},
\href{https://github.com/yang-mills-group/OpenAI-YM-20261028/blob/3456e5e75cf8d4cd558955a5c7638fba2d68f464/YM/Next/RANK/Planck.lean\#L68-L69}{\texttt{shed}}.

\subsubsection{5. Where the difficulty
lies}\label{where-the-difficulty-lies-2}

The tree estimate must preserve the connected cover weight. Cover
numbers are subadditive along intersecting unions, whereas a
disconnected union may require an additional connecting region. The
induction therefore uses connected components and spends exactly one
extra cover exponent on each attachment. It counts finite sets of
labels; it does not estimate Mayer coefficients with arbitrary repeated
labels.

For two interactions, the graph itself depends on the selected set,
through hotness. Unique recovery requires both locality of hotness and
the asymmetric separation conditions: all components separate through
mentions, while only components containing a dirty label separate
through walls.

Switching requires convergence as well as a bound. Bounded coefficients
alone need not have a limit as \(x\downarrow0\). The model supplies
convergence, and the common majorant allows sums and coefficient
estimates to pass to the limit. The Gaussian bounds allow complex
covariances without positivity; their constants grow like \(d^d\) up to
exponential factors at fixed entry bounds. The degree is fixed before
\(L\) and \(s\).

The numerical conditions establish feasibility for fixed parameters.
They do not supply the analytic estimates to which those numbers are
applied. The two iteration lengths in §3 must remain distinct.

The source contains duplicate exponential and component constructions.
Its second trajectory logarithm and
\href{https://github.com/yang-mills-group/OpenAI-YM-20261028/blob/3456e5e75cf8d4cd558955a5c7638fba2d68f464/YM/Next/RANK/Ratios.lean\#L72-L89}{\texttt{Candle.small}}
are absent from the main theorem's dependency set. The latter bounds a
restarted threshold by \(3h(\hat\beta)\beta^{-1/2}\).

\subsubsection{6. Notation}\label{notation-2}

{\def\LTcaptype{none} % do not increment counter
\begin{longtable}[]{@{}
  >{\raggedright\arraybackslash}p{(\linewidth - 4\tabcolsep) * \real{0.3333}}
  >{\raggedright\arraybackslash}p{(\linewidth - 4\tabcolsep) * \real{0.3333}}
  >{\raggedright\arraybackslash}p{(\linewidth - 4\tabcolsep) * \real{0.3333}}@{}}
\toprule\noalign{}
\begin{minipage}[b]{\linewidth}\raggedright
Written proof
\end{minipage} & \begin{minipage}[b]{\linewidth}\raggedright
Lean
\end{minipage} & \begin{minipage}[b]{\linewidth}\raggedright
Meaning
\end{minipage} \\
\midrule\noalign{}
\endhead
\bottomrule\noalign{}
\endlastfoot
\(R\), census &
\href{https://github.com/yang-mills-group/OpenAI-YM-20261028/blob/3456e5e75cf8d4cd558955a5c7638fba2d68f464/YM/Marker/Console.lean\#L22-L26}{\texttt{Marker.Census}}
& finite downward-closed family of allowed marker sets \\
\(\mathcal A_R\) &
\href{https://github.com/yang-mills-group/OpenAI-YM-20261028/blob/3456e5e75cf8d4cd558955a5c7638fba2d68f464/YM/Marker/Console.lean\#L40}{\texttt{Marker.Poly}}
& marker algebra \\
\(c\,\delta_x\) &
\href{https://github.com/yang-mills-group/OpenAI-YM-20261028/blob/3456e5e75cf8d4cd558955a5c7638fba2d68f464/YM/Marker/Console.lean\#L51}{\texttt{Poly.send}}
& monomial with marker set \(x\) \\
\(w^x\), \(\lVert a\rVert_w\) &
\href{https://github.com/yang-mills-group/OpenAI-YM-20261028/blob/3456e5e75cf8d4cd558955a5c7638fba2d68f464/YM/Marker/Weight.lean\#L74}{\texttt{toll}},
\href{https://github.com/yang-mills-group/OpenAI-YM-20261028/blob/3456e5e75cf8d4cd558955a5c7638fba2d68f464/YM/Marker/Weight.lean\#L76}{\texttt{price}}
& weight of a marker set; weighted \(\ell^1\) norm \\
\(\mathrm{Exp}\) &
\href{https://github.com/yang-mills-group/OpenAI-YM-20261028/blob/3456e5e75cf8d4cd558955a5c7638fba2d68f464/YM/Marker/Truncated.lean\#L29}{\texttt{Poly.pass}}
& exponential truncated at order \(\lvert V\rvert\) \\
connected polynomial \(\mathrm{Log}\,F\) &
\href{https://github.com/yang-mills-group/OpenAI-YM-20261028/blob/3456e5e75cf8d4cd558955a5c7638fba2d68f464/YM/Next/RANK/Exponential.lean\#L177}{\texttt{Exponential.connected}},
\href{https://github.com/yang-mills-group/OpenAI-YM-20261028/blob/3456e5e75cf8d4cd558955a5c7638fba2d68f464/YM/Next/RANK/Exponential.lean\#L126}{\texttt{Expands}}
& \(C\) with \(C(\emptyset)=0\), \(\mathrm{Exp}(C)=F\) \\
majorant array &
\href{https://github.com/yang-mills-group/OpenAI-YM-20261028/blob/3456e5e75cf8d4cd558955a5c7638fba2d68f464/YM/Push/GRID_APPROXIMATION/Nil.lean\#L75}{\texttt{Nil.Sheet}}
& nonnegative coefficient bound \\
\(\mathrm{cov}_M(X)\) &
\href{https://github.com/yang-mills-group/OpenAI-YM-20261028/blob/3456e5e75cf8d4cd558955a5c7638fba2d68f464/YM/Lattice/Support.lean\#L69-L71}{\texttt{coverNumber}}
& connected cover number (I.1) \\
spark, \(\mathrm{Parts}(s)\) &
\href{https://github.com/yang-mills-group/OpenAI-YM-20261028/blob/3456e5e75cf8d4cd558955a5c7638fba2d68f464/YM/Ac/Nodes.lean\#L49}{\texttt{Ac.Spark}},
\href{https://github.com/yang-mills-group/OpenAI-YM-20261028/blob/3456e5e75cf8d4cd558955a5c7638fba2d68f464/YM/Ac/Nodes.lean\#L55}{\texttt{Ac.Parts}}
& connected label set; components \\
pack for \(Y\), \(\mathrm{film}(D,Y)\) &
\href{https://github.com/yang-mills-group/OpenAI-YM-20261028/blob/3456e5e75cf8d4cd558955a5c7638fba2d68f464/YM/Ac/Nodes.lean\#L18}{\texttt{Ac.Pack}},
\href{https://github.com/yang-mills-group/OpenAI-YM-20261028/blob/3456e5e75cf8d4cd558955a5c7638fba2d68f464/YM/Ac/Table.lean\#L79}{\texttt{Ac.film}}
& label set attached to \(Y\); its weighted count \\
hit condition &
\href{https://github.com/yang-mills-group/OpenAI-YM-20261028/blob/3456e5e75cf8d4cd558955a5c7638fba2d68f464/YM/Ac/Tree.lean\#L12-L13}{\texttt{Ac.Hit}}
& site-wise summability of \(f\,e^{a\,\mathrm{cov}}\) \\
kiln, hot, fuel, mole, slab &
\href{https://github.com/yang-mills-group/OpenAI-YM-20261028/blob/3456e5e75cf8d4cd558955a5c7638fba2d68f464/YM/Further/SHELL_BIRTH_49/Fusing/Coal.lean\#L123-L126}{\texttt{Kiln}},
\href{https://github.com/yang-mills-group/OpenAI-YM-20261028/blob/3456e5e75cf8d4cd558955a5c7638fba2d68f464/YM/Further/SHELL_BIRTH_49/Fusing/Coal.lean\#L131}{\texttt{Hot}},
\href{https://github.com/yang-mills-group/OpenAI-YM-20261028/blob/3456e5e75cf8d4cd558955a5c7638fba2d68f464/YM/Further/SHELL_BIRTH_49/Fusing/Coal.lean\#L132}{\texttt{fuel}},
\href{https://github.com/yang-mills-group/OpenAI-YM-20261028/blob/3456e5e75cf8d4cd558955a5c7638fba2d68f464/YM/Further/SHELL_BIRTH_49/Fusing/Coal.lean\#L133}{\texttt{Mole}},
\href{https://github.com/yang-mills-group/OpenAI-YM-20261028/blob/3456e5e75cf8d4cd558955a5c7638fba2d68f464/YM/Further/SHELL_BIRTH_49/Fusing/Coal.lean\#L135}{\texttt{slabs}}
& two-interaction polymer data and components \\
pipes, tap, pellet, reach &
\href{https://github.com/yang-mills-group/OpenAI-YM-20261028/blob/3456e5e75cf8d4cd558955a5c7638fba2d68f464/YM/Ovens/Coffee.lean\#L23-L27}{\texttt{Ovens.Pipes}},
\href{https://github.com/yang-mills-group/OpenAI-YM-20261028/blob/3456e5e75cf8d4cd558955a5c7638fba2d68f464/YM/Ovens/Coffee.lean\#L32-L35}{\texttt{Tap}},
\href{https://github.com/yang-mills-group/OpenAI-YM-20261028/blob/3456e5e75cf8d4cd558955a5c7638fba2d68f464/YM/Ovens/Coffee.lean\#L57-L58}{\texttt{contains}},
\href{https://github.com/yang-mills-group/OpenAI-YM-20261028/blob/3456e5e75cf8d4cd558955a5c7638fba2d68f464/YM/Ovens/Coffee.lean\#L48}{\texttt{reach}}
& data and pieces of the fused gas \\
\(\mathrm{val}_x\), model, jet, paid &
\href{https://github.com/yang-mills-group/OpenAI-YM-20261028/blob/3456e5e75cf8d4cd558955a5c7638fba2d68f464/YM/MergeProof/Switch.lean\#L26}{\texttt{Switch.value}},
\href{https://github.com/yang-mills-group/OpenAI-YM-20261028/blob/3456e5e75cf8d4cd558955a5c7638fba2d68f464/YM/MergeProof/Switch.lean\#L48-L50}{\texttt{Models}},
\href{https://github.com/yang-mills-group/OpenAI-YM-20261028/blob/3456e5e75cf8d4cd558955a5c7638fba2d68f464/YM/MergeProof/Switch.lean\#L46}{\texttt{jet}},
\href{https://github.com/yang-mills-group/OpenAI-YM-20261028/blob/3456e5e75cf8d4cd558955a5c7638fba2d68f464/YM/MergeProof/Switch.lean\#L56}{\texttt{Paid}}
& switching objects \\
\(\mathrm{snd}_\rho\) &
\href{https://github.com/yang-mills-group/OpenAI-YM-20261028/blob/3456e5e75cf8d4cd558955a5c7638fba2d68f464/YM/MergeProof/SwitchCost.lean\#L60}{\texttt{Switch.sound}}
& jets inserted into the marker algebra \\
\(\lVert p\rVert_1\), \(\lVert p\rVert_{w,r}\) &
\href{https://github.com/yang-mills-group/OpenAI-YM-20261028/blob/3456e5e75cf8d4cd558955a5c7638fba2d68f464/YM/Next/RANK/Rise.lean\#L47}{\texttt{Rise.coin}},
\href{https://github.com/yang-mills-group/OpenAI-YM-20261028/blob/3456e5e75cf8d4cd558955a5c7638fba2d68f464/YM/Next/RANK/Rise.lean\#L13}{\texttt{Rise.cents}}
& coefficient norm; weighted coefficient seminorm \\
\(\mathcal G_Q\), \(\mathbb E_{Q,m}\) &
\href{https://github.com/yang-mills-group/OpenAI-YM-20261028/blob/3456e5e75cf8d4cd558955a5c7638fba2d68f464/YM/Next/RANK/Gaussian.lean\#L77}{\texttt{Normal.milk}},
\href{https://github.com/yang-mills-group/OpenAI-YM-20261028/blob/3456e5e75cf8d4cd558955a5c7638fba2d68f464/YM/Next/RANK/Gaussian.lean\#L106}{\texttt{Normal.blend}}
& Gaussian smoothing; algebraic expectation \\
\(\partial_b\) &
\href{https://github.com/yang-mills-group/OpenAI-YM-20261028/blob/3456e5e75cf8d4cd558955a5c7638fba2d68f464/YM/Next/RANK/Gaussian.lean\#L40}{\texttt{Normal.slit}}
& derivation \(X_v\mapsto b_v\) \\
\(p\circ A\) &
\href{https://github.com/yang-mills-group/OpenAI-YM-20261028/blob/3456e5e75cf8d4cd558955a5c7638fba2d68f464/YM/Next/RANK/GaussSubst.lean\#L14}{\texttt{Normal.exchange}}
& linear substitution \\
\(T_{H,Q}\) &
\href{https://github.com/yang-mills-group/OpenAI-YM-20261028/blob/3456e5e75cf8d4cd558955a5c7638fba2d68f464/YM/Next/RANK/GaussTags.lean\#L113}{\texttt{Normal.pipe}},
\href{https://github.com/yang-mills-group/OpenAI-YM-20261028/blob/3456e5e75cf8d4cd558955a5c7638fba2d68f464/YM/Next/RANK/Herald.lean\#L18}{\texttt{Valve.read}}
& one Gaussian polynomial step \\
\(\pi(B,d)\) &
\href{https://github.com/yang-mills-group/OpenAI-YM-20261028/blob/3456e5e75cf8d4cd558955a5c7638fba2d68f464/YM/Next/RANK/Common.lean\#L36}{\texttt{Common.priced}}
& \((2B(d+1))^d\) \\
\(\eta\)-guard &
\href{https://github.com/yang-mills-group/OpenAI-YM-20261028/blob/3456e5e75cf8d4cd558955a5c7638fba2d68f464/YM/Next/RANK/Guard.lean\#L69-L73}{\texttt{Guarded}}
& \([0,1]\)-valued cutoff of small defect \\
\(\Pi(s)\), \(\Pi(s)^a\) &
\href{https://github.com/yang-mills-group/OpenAI-YM-20261028/blob/3456e5e75cf8d4cd558955a5c7638fba2d68f464/YM/Vac/Treasury.lean\#L72}{\texttt{Numbers.P}},
\href{https://github.com/yang-mills-group/OpenAI-YM-20261028/blob/3456e5e75cf8d4cd558955a5c7638fba2d68f464/YM/Next/RANK/Reign.lean\#L25}{\texttt{Perch.bill}}
& \(e^{(\log\log s)^2}\) and its powers \\
\(\theta(s)\) &
\href{https://github.com/yang-mills-group/OpenAI-YM-20261028/blob/3456e5e75cf8d4cd558955a5c7638fba2d68f464/YM/Vac/Treasury.lean\#L73}{\texttt{Numbers.theta}}
& \(s^{-1/2}\Pi(s)\), small-field threshold \\
\(h,\ j,\ \omega,\ W\) &
\href{https://github.com/yang-mills-group/OpenAI-YM-20261028/blob/3456e5e75cf8d4cd558955a5c7638fba2d68f464/YM/Next/RANK/Reign.lean\#L30}{\texttt{hot}},
\href{https://github.com/yang-mills-group/OpenAI-YM-20261028/blob/3456e5e75cf8d4cd558955a5c7638fba2d68f464/YM/Next/RANK/Reign.lean\#L31}{\texttt{join}},
\href{https://github.com/yang-mills-group/OpenAI-YM-20261028/blob/3456e5e75cf8d4cd558955a5c7638fba2d68f464/YM/Next/RANK/Reign.lean\#L32}{\texttt{w}},
\href{https://github.com/yang-mills-group/OpenAI-YM-20261028/blob/3456e5e75cf8d4cd558955a5c7638fba2d68f464/YM/Next/RANK/Reign.lean\#L28}{\texttt{walls}}
& \(e^{y^{3/2}}\), \(h^{1/5}\), \(e^{-0.7j}\),
\(\lceil(\log s)^2\rceil^2\) \\
parameter record \((t,g,H_0,b,c,f,F,K)\) &
\href{https://github.com/yang-mills-group/OpenAI-YM-20261028/blob/3456e5e75cf8d4cd558955a5c7638fba2d68f464/YM/Next/RANK/Ratios.lean\#L20-L38}{\texttt{Tapes.Purse}}
& \texttt{taylor}, \texttt{goal}, \texttt{horizon}, \texttt{band},
\texttt{credit}, \texttt{f}, \texttt{foil}, \texttt{L} \\
iteration conditions at height \(H\) &
\href{https://github.com/yang-mills-group/OpenAI-YM-20261028/blob/3456e5e75cf8d4cd558955a5c7638fba2d68f464/YM/Next/RANK/Ratios.lean\#L64-L69}{\texttt{Racer}}
& five inequalities in II.2.13(d) \\
moment and fusion conditions &
\href{https://github.com/yang-mills-group/OpenAI-YM-20261028/blob/3456e5e75cf8d4cd558955a5c7638fba2d68f464/YM/Next/RANK/Ratios.lean\#L92-L110}{\texttt{Faring}}
& inequalities in II.2.13(e) \\
\end{longtable}
}

\subsubsection{7. What remains
unexplained}\label{what-remains-unexplained-2}

\begin{enumerate}
\def\labelenumi{\arabic{enumi}.}
\tightlist
\item
  \textbf{Restarted thresholds.} The original threshold is
  \(s^{-1/2}\Pi(s)\); the restarted Gaussian comparison uses
  \(s^{-1/2}h(s)\) and cover size \(W(s)\)
  (\href{https://github.com/yang-mills-group/OpenAI-YM-20261028/blob/3456e5e75cf8d4cd558955a5c7638fba2d68f464/YM/Next/RANK/Stepped.lean\#L93-L96}{\texttt{Perch.Candle}},
  III.11). Their ratio \(h(s)/\Pi(s)\) tends to zero. The numerical
  estimates here do not establish the existence of restarted
  representations satisfying the required density and trajectory
  identities. That construction must be explained in III.11.
\item
  \textbf{Analytic inputs to the numerical conditions.} The shelf, tally
  and moment conditions are explicit in II.2.13. Their use still
  requires the one-step estimates with the stated supports, covariances
  and constants. The verification of those analytic hypotheses along the
  actual RG trajectories, including every later application of the
  digest and fusion bounds, is outside this unit.
\item
  \textbf{Auxiliary declarations.} The additional regrouping results in
  \texttt{EventBale} and the cost estimates in \texttt{MergeProof/West},
  \texttt{North} and \texttt{East} are not formulated above, apart from
  the cited permutation result. Their full role in III.11 remains to be
  traced.
\end{enumerate}

\subsection{III.1. The block-spin step and the symmetric gauge
fixing}\label{iii.1.-the-block-spin-step-and-the-symmetric-gauge-fixing}

\textbf{Part:} III, Ultraviolet construction (renormalization group).
\textbf{Lean:}
\href{https://github.com/yang-mills-group/OpenAI-YM-20261028/tree/3456e5e75cf8d4cd558955a5c7638fba2d68f464/YM/RG}{\texttt{YM/RG/}}
(Block, Patch, Measurable, Push, Masks, Activities, PathActions,
Bounded, History),
\href{https://github.com/yang-mills-group/OpenAI-YM-20261028/tree/3456e5e75cf8d4cd558955a5c7638fba2d68f464/YM/Branch}{\texttt{YM/Branch/}},
\href{https://github.com/yang-mills-group/OpenAI-YM-20261028/tree/3456e5e75cf8d4cd558955a5c7638fba2d68f464/YM/Pressure}{\texttt{YM/Pressure/}},
\href{https://github.com/yang-mills-group/OpenAI-YM-20261028/tree/3456e5e75cf8d4cd558955a5c7638fba2d68f464/YM/PREC}{\texttt{YM/PREC/}},
\href{https://github.com/yang-mills-group/OpenAI-YM-20261028/tree/3456e5e75cf8d4cd558955a5c7638fba2d68f464/YM/PCalc}{\texttt{YM/PCalc/}},
\href{https://github.com/yang-mills-group/OpenAI-YM-20261028/tree/3456e5e75cf8d4cd558955a5c7638fba2d68f464/YM/Tree}{\texttt{YM/Tree/}},
\href{https://github.com/yang-mills-group/OpenAI-YM-20261028/tree/3456e5e75cf8d4cd558955a5c7638fba2d68f464/YM/Exterior}{\texttt{YM/Exterior/}},
and the small packages \texttt{YM/RG17/},
\texttt{YM/Push/CORRECTED\_COLUMNS/},
\texttt{YM/Planned/CALIBRATED\_MIXED\_HISTORY/},
\texttt{YM/UVNext/ANNULUS\_COMPARISON/}, \texttt{YM/FullVolume/Train/},
\texttt{YM/SectorTools/ACTUAL\_SCALAR\_LEDGER/}. \textbf{Size:} 4,219
lines (693 declarations) in the main-theorem closure.

\textbf{Conventions.} \(N\ge1\) and \(G\subset U(N)\) is a closed
subgroup with probability Haar measure \(dg\). Its Lie algebra
\(\mathfrak g\subset M_N(\mathbb C)\) carries the real inner product
\(\mathrm{Re\,tr}(X^{\dagger}Y)\), and \(\Pi_{\mathfrak g}\) denotes the
orthogonal projection onto \(\mathfrak g\). Matrix norms below are
Frobenius norms. In path actions, Lie-algebra fields, Hessians and the
six works, the physical inner product is
\(N^{-1}\mathrm{Re\,tr}(X^{\dagger}Y)\); write \(|X|=\|X\|/\sqrt N\) for
its norm. This is the coordinate normalization of
\href{https://github.com/yang-mills-group/OpenAI-YM-20261028/blob/3456e5e75cf8d4cd558955a5c7638fba2d68f464/YM/Gauge/Colors.lean\#L56-L62}{\texttt{Gauge.inner\_frob\_color}}.
The curvature field is
\(\mathrm{curv}(u)=\Pi_{\mathfrak g}(u-u^{-1})/2\), with
\(|\mathrm{curv}(u)|\le\mathrm{ch}(u)\). The \emph{chord} is
\(\mathrm{ch}(u)=\|u-1\|/\sqrt N\) and the \emph{price} is
\(\mathrm{pr}(u)=1-\mathrm{Re\,tr}\,u/N\). They satisfy
\(\mathrm{ch}^2=2\,\mathrm{pr}\), both are class functions,
\(\mathrm{ch}(u^{-1})=\mathrm{ch}(u)\le2\),
\(\mathrm{ch}(uv)\le\mathrm{ch}(u)+\mathrm{ch}(v)\), and
\(\|u-v\|^2=2N\,\mathrm{pr}(uv^{-1})\). The Wilson energy is
\(\sum_p\mathrm{pr}(U_p)\).

A step goes from the fine torus \(\Lambda_f=(\mathbb Z/f)^4\) to the
coarse torus \(\Lambda_c=(\mathbb Z/c)^4\), with \(c\ge1\), \(f=cL\) and
\(L\ge3\) odd. The coarse site \(x\) has root \(Lx\) and block
\(Lx+B_{(L-1)/2}\), where \(B_r=\{-r,\dots,r\}^4\subset\mathbb Z^4\).
Configurations are \(U\in G^{\mathrm{links}}\), and \(\mu_f\) is the
product of the probability Haar measures. The gauge action is
\((k\cdot U)_{(y,i)}=k_yU_{(y,i)}k_{y+e_i}^{-1}\). \(T_U(y;w)\) is the
parallel transport along a walk \(w\). \(T_U(y\to z;o)\) is the
transport along the coordinate route that changes the coordinates in the
order given by \(o\in S_4\). For a coarse site \(x\), \(U^{(x)}\) is the
periodic lift of \(U\) to \(\mathbb Z^4\), translated so that \(Lx\)
sits at the origin. All patch constructions take place on
\(\mathbb Z^4\), so a patch may be larger than the torus. For finite
\(X\subset\mathbb Z^4\), \(E_X(U)\) is the sum of the chords of the
links with both ends in \(X\), and \(F_X(U)\) is the sum of the chords
of the plaquettes with all four corners in \(X\); \(F_X\) is gauge
invariant. \(\mathfrak B_4\) is the group of signed permutations of the
axes (384 elements). For \(s\in\mathfrak B_4\) and a point \(b\),
\(s_bU\) is the pullback of \(U\) under \(z\mapsto b+s(z)\); a link
reversed by \(s\) carries the inverse variable. Further,
\(\kappa_{\rm slope}=11C_A/(24\pi^2)\), where \(C_A\) is the adjoint
Casimir, and \(\mathrm{temp}(x)=x^{-1/2}e^{(\log\log x)^2}\) for
\(x>1\).

\textbf{Definition (the block kernel).} A \emph{patch} is \(P=(L,J)\)
with \(L\ge3\) odd and \(J\ge4\). Its radius is \(R=(J+3)L\) and its
points are those of \(B_R\), with centre \(0\). The \emph{kernel
constants} are \(K=(\kappa,\delta,\varepsilon,a)\), all positive:
penalty, frame cutoff, predictor cutoff, axial factor. For a
configuration \(W\) on \(\mathbb Z^4\) the following objects are
defined.

\begin{enumerate}
\def\labelenumi{\arabic{enumi}.}
\tightlist
\item
  \emph{Frame tags and axial frames.} A frame tag is
  \(t=(r,o)\in B_R\times S_4\). The axial frame of \(t\) is
  \(a^t_y=T_W(r\to y;o)\).
\item
  \emph{Minimizing branch.} \(W\in M_\delta\) if
  \(\mathrm{ch}\big(a^{(0,o)}_y\,T_W(y\to z;o')\,(a^{(0,o)}_z)^{-1}\big)<\delta\)
  for all \(o,o'\in S_4\) and \(y,z\in B_R\). In words: in each of the
  24 axial gauges from the centre, every route transport in the patch is
  \(\delta\)-close to \(1\).
\item
  \emph{Score.} For \(h:B_R\to G\),
  \(S_W(h)=\mathrm{avg}_{y,z\in B_R}\mathrm{avg}_{o\in S_4}\,\mathrm{pr}\big(h_yT_W(y\to z;o)h_z^{-1}\big)\).
\item
  \emph{Eligible tags.}
  \(\mathcal E(W)=\{t:\ E_{B_R}(a^t\cdot W)\le a(2R+1)F_{B_R}(W)\}\).
\item
  \emph{Frame and tag mass.} If \(W\in M_\delta\), the frame \(h^W\) is
  a minimizer of \(S_W\) subject to \(h_0=1\), and the tag mass is
  uniform, \(1/(24|B_R|)\). Otherwise \(h^{W,t}_y=(a^t_0)^{-1}a^t_y\)
  (the axial gauge from \(r\), normalized at the centre), with mass
  \(\mathbf 1_{\mathcal E(W)}(t)/|\mathcal E(W)|\).
\item
  \emph{Weight.} \(w\) is the law on \(B_R\) of \(\xi_1+\dots+\xi_J\),
  where the \(\xi_k\) are i.i.d. uniform on \(B_{(L-1)/2}\). It is the
  \(J\)-fold convolution of the uniform weight on a centred block; it is
  \(\mathfrak B_4\)-invariant and supported in \(B_{J(L-1)/2}\).
\item
  \emph{Connector.} For a coarse link \(b=(x,i)\) and frame tags
  \((t_x)_{x\in\Lambda_c}\), set
  \(Y_y=h^{(x)}_y\,T_{U^{(x)}}(y;\ L\text{ steps along }e_i)\,(h^{(x+e_i)}_y)^{-1}\)
  for \(y\in B_R\). Here \(h^{(x)}\) is the frame of \(U^{(x)}\) with
  tag \(t_x\). This is the fine transport from \(Lx+y\) to
  \(L(x+e_i)+y\), written in the frames of the two patches.
\item
  \emph{Prediction.} Let
  \(\Phi_Y(u)=\sum_yw_y\,\mathrm{pr}(Y_yu^{-1})\). If
  \(\min\Phi_Y<\varepsilon\), then \(Q_b\) is a chosen minimizer,
  equivalently a nearest point of \(G\) to the matrix mean
  \(\sum_yw_yY_y\) (an extrinsic, chordal mean); the bond tag then has
  uniform mass \(1/|B_R|\). Otherwise \(Q_b=Y_{r_b}\), with the bond tag
  \(r_b\) drawn with probability \(w_{r_b}\).
\item
  \emph{Kernel.} The tag law is
  \(\lambda(U,t)=\prod_{x}m^{\rm fr}(U^{(x)},t_x)\prod_b m^{\rm pr}(Y_b,r_b)\),
  and
\end{enumerate}

\[\rho_\beta(U,V)=\sum_t\lambda(U,t)\prod_{b}Z_\beta^{-1}e^{-\beta\kappa\,\mathrm{pr}(V_bQ_b(U,t)^{-1})},\qquad Z_\beta=\int_Ge^{-\beta\kappa\,\mathrm{pr}(g)}\,dg .\]

The \emph{branch conditions} for \((P,K,G)\) are:

\begin{itemize}
\tightlist
\item
  (B1) for \(W\in M_\delta\) the constrained minimizer in item 5 is
  unique;
\item
  (B2) whenever \(\min\Phi_Y<\varepsilon\), the minimizer of \(\Phi_Y\)
  is unique;
\item
  (B3) \(\mathcal E(W)\ne\emptyset\) for every \(W\).
\end{itemize}

Compactness gives the chosen minimizers; (B1)--(B2) make the choices
unique on the minimizing branches. Lean:
\href{https://github.com/yang-mills-group/OpenAI-YM-20261028/blob/3456e5e75cf8d4cd558955a5c7638fba2d68f464/YM/RG/Patch.lean\#L29-L34}{\texttt{Patch}},
\href{https://github.com/yang-mills-group/OpenAI-YM-20261028/blob/3456e5e75cf8d4cd558955a5c7638fba2d68f464/YM/RG/Patch.lean\#L37-L45}{\texttt{KernelChoices}},
\href{https://github.com/yang-mills-group/OpenAI-YM-20261028/blob/3456e5e75cf8d4cd558955a5c7638fba2d68f464/YM/RG/Patch.lean\#L109-L110}{\texttt{usesMin}},
\href{https://github.com/yang-mills-group/OpenAI-YM-20261028/blob/3456e5e75cf8d4cd558955a5c7638fba2d68f464/YM/RG/Patch.lean\#L114-L115}{\texttt{score}},
\href{https://github.com/yang-mills-group/OpenAI-YM-20261028/blob/3456e5e75cf8d4cd558955a5c7638fba2d68f464/YM/RG/Patch.lean\#L143-L146}{\texttt{eligible}},
\href{https://github.com/yang-mills-group/OpenAI-YM-20261028/blob/3456e5e75cf8d4cd558955a5c7638fba2d68f464/YM/RG/Patch.lean\#L148-L151}{\texttt{frameMass}},
\href{https://github.com/yang-mills-group/OpenAI-YM-20261028/blob/3456e5e75cf8d4cd558955a5c7638fba2d68f464/YM/RG/Patch.lean\#L153-L155}{\texttt{frame}},
\href{https://github.com/yang-mills-group/OpenAI-YM-20261028/blob/3456e5e75cf8d4cd558955a5c7638fba2d68f464/YM/RG/Patch.lean\#L87-L88}{\texttt{weight}},
\href{https://github.com/yang-mills-group/OpenAI-YM-20261028/blob/3456e5e75cf8d4cd558955a5c7638fba2d68f464/YM/RG/Patch.lean\#L159}{\texttt{bondCost}},
\href{https://github.com/yang-mills-group/OpenAI-YM-20261028/blob/3456e5e75cf8d4cd558955a5c7638fba2d68f464/YM/RG/Patch.lean\#L173-L175}{\texttt{predict}},
\href{https://github.com/yang-mills-group/OpenAI-YM-20261028/blob/3456e5e75cf8d4cd558955a5c7638fba2d68f464/YM/RG/Patch.lean\#L177-L179}{\texttt{predictorMass}},
\href{https://github.com/yang-mills-group/OpenAI-YM-20261028/blob/3456e5e75cf8d4cd558955a5c7638fba2d68f464/YM/RG/Patch.lean\#L182-L185}{\texttt{Branches}},
\href{https://github.com/yang-mills-group/OpenAI-YM-20261028/blob/3456e5e75cf8d4cd558955a5c7638fba2d68f464/YM/RG/Block.lean\#L22-L25}{\texttt{StepData}},
\href{https://github.com/yang-mills-group/OpenAI-YM-20261028/blob/3456e5e75cf8d4cd558955a5c7638fba2d68f464/YM/RG/Block.lean\#L39-L42}{\texttt{connector}},
\href{https://github.com/yang-mills-group/OpenAI-YM-20261028/blob/3456e5e75cf8d4cd558955a5c7638fba2d68f464/YM/RG/Block.lean\#L44-L45}{\texttt{prediction}},
\href{https://github.com/yang-mills-group/OpenAI-YM-20261028/blob/3456e5e75cf8d4cd558955a5c7638fba2d68f464/YM/RG/Block.lean\#L49-L51}{\texttt{likelihood}},
\href{https://github.com/yang-mills-group/OpenAI-YM-20261028/blob/3456e5e75cf8d4cd558955a5c7638fba2d68f464/YM/RG/Block.lean\#L71-L72}{\texttt{density}}.

\subsubsection{1. What this unit proves}\label{what-this-unit-proves-3}

The block-spin kernel is normalized, Borel and covariant under gauge and
lattice symmetries. Its branch cutoffs can be chosen independently of
the patch size. The geometric estimates synchronize neighbouring frames
and bound the predicted coarse curvature by \(L^{-2}\) times the fine
curvature in a surrounding box. Near the trivial configuration the
prediction is analytic, with linearization \(\mathcal M+d_0\mathcal F\).
Tree gauge reduces the step integral to the free links. The unit also
defines effective densities, RG histories and the requirements used in
the later energy and stationary-branch estimates.

\textbf{Proposition III.1.1 (the block kernel is a Markov kernel).}
Assume (B1)--(B3). For every \(\beta\in\mathbb R\):

\begin{itemize}
\tightlist
\item
  \(\rho_\beta\ge0\);
\item
  \((U,V)\mapsto\rho_\beta(U,V)\) is Borel;
\item
  \(\sup\rho_\beta<\infty\);
\item
  \(\int\rho_\beta(U,V)\,d\mu_c(V)=1\) for every \(U\).
\end{itemize}

Consequently, for every Banach-space-valued \(F\in L^1(\mu_f)\) the push
\((\mathcal P_\beta F)(V)=\int\rho_\beta(U,V)F(U)\,d\mu_f(U)\) is
defined for every \(V\). Moreover \(\mathcal P_\beta F\in L^1(\mu_c)\),
\(\mathcal P_\beta\) is linear, and
\(\int\mathcal P_\beta F\,d\mu_c=\int F\,d\mu_f\). This is the operation
that takes level \(m+1\) to level \(m\). \emph{Status:} theorem.
Positivity and integrability in \(V\) hold without (B1)--(B3).
\emph{Lean:}
\href{https://github.com/yang-mills-group/OpenAI-YM-20261028/blob/3456e5e75cf8d4cd558955a5c7638fba2d68f464/YM/RG/Block.lean\#L158-L163}{\texttt{IsKernel}},
\href{https://github.com/yang-mills-group/OpenAI-YM-20261028/blob/3456e5e75cf8d4cd558955a5c7638fba2d68f464/YM/RG/Block.lean\#L203-L208}{\texttt{isKernel}},
\href{https://github.com/yang-mills-group/OpenAI-YM-20261028/blob/3456e5e75cf8d4cd558955a5c7638fba2d68f464/YM/RG/Block.lean\#L77-L79}{\texttt{push}},
\href{https://github.com/yang-mills-group/OpenAI-YM-20261028/blob/3456e5e75cf8d4cd558955a5c7638fba2d68f464/YM/RG/Push.lean\#L33-L35}{\texttt{push\_integrable}},
\href{https://github.com/yang-mills-group/OpenAI-YM-20261028/blob/3456e5e75cf8d4cd558955a5c7638fba2d68f464/YM/RG/Push.lean\#L37-L43}{\texttt{push\_integral}},
\href{https://github.com/yang-mills-group/OpenAI-YM-20261028/blob/3456e5e75cf8d4cd558955a5c7638fba2d68f464/YM/RG/Push.lean\#L45-L50}{\texttt{push\_add}},
\href{https://github.com/yang-mills-group/OpenAI-YM-20261028/blob/3456e5e75cf8d4cd558955a5c7638fba2d68f464/YM/RG/Push.lean\#L52-L55}{\texttt{push\_smul}}.

\textbf{Proposition III.1.2 (exact covariance).} Assume (B1)--(B3). For
all \(\beta\in\mathbb R\), \(U\), \(V\):

\begin{itemize}
\tightlist
\item
  \begin{enumerate}
  \def\labelenumi{(\alph{enumi})}
  \tightlist
  \item
    \(\rho_\beta\big(k\cdot U,(k\circ\mathrm{embed})\cdot V\big)=\rho_\beta(U,V)\)
    for every \(k:\Lambda_f\to G\), where \(\mathrm{embed}(x)=Lx\);
  \end{enumerate}
\item
  (b) \(\rho_\beta(s_{Lb}U,s_bV)=\rho_\beta(U,V)\) for every
  \(s\in\mathfrak B_4\) and \(b\in\Lambda_c\);
\item
  (c) if the kernel constants of a calendar satisfy (B1)--(B3) for every
  patch, then every step of the calendar has \emph{block control}
  (branch conditions together with (a) and (b)), and is a Markov kernel
  by III.1.1.
\end{itemize}

\emph{Status:} theorem. \emph{Lean:}
\href{https://github.com/yang-mills-group/OpenAI-YM-20261028/blob/3456e5e75cf8d4cd558955a5c7638fba2d68f464/YM/Branch/Covariant.lean\#L56-L67}{\texttt{covary}},
\href{https://github.com/yang-mills-group/OpenAI-YM-20261028/blob/3456e5e75cf8d4cd558955a5c7638fba2d68f464/YM/Branch/Covariant.lean\#L50-L54}{\texttt{channel}},
\href{https://github.com/yang-mills-group/OpenAI-YM-20261028/blob/3456e5e75cf8d4cd558955a5c7638fba2d68f464/YM/Branch/Covariant.lean\#L70-L72}{\texttt{Calendar.blocks\_go}},
\href{https://github.com/yang-mills-group/OpenAI-YM-20261028/blob/3456e5e75cf8d4cd558955a5c7638fba2d68f464/YM/RG/History.lean\#L120-L128}{\texttt{Calendar.BlockControl}},
\href{https://github.com/yang-mills-group/OpenAI-YM-20261028/blob/3456e5e75cf8d4cd558955a5c7638fba2d68f464/YM/RG/History.lean\#L158-L159}{\texttt{Levels.transition\_mass}}.

\textbf{Proposition III.1.3 (branch constants depending only on \(G\)).}
Let \(d=d(G)\) be the constant of III.1.5. Let \(r_G>0\) be such that
every \(X\in M_N(\mathbb C)\) at distance \(<r_G\) from \(G\) has a
unique nearest point in \(G\). Put \(\delta_0=d/\sqrt N\) and
\(\varepsilon_0=r_G^2/(2N)\). Then for every \(\kappa>0\),
\(\delta\in(0,\delta_0]\), \(\varepsilon\in(0,\varepsilon_0]\) and
\(a\ge16\), the constants \((\kappa,\delta,\varepsilon,a)\) satisfy
(B1)--(B3) for every patch, that is, for every odd \(L\ge3\) and every
\(J\ge4\). Condition (B3) holds in a stronger form: for every \(W\) and
every ordering \(o\) there is a root \(r\) with \((r,o)\) eligible.
\emph{Status:} theorem. The existence theorem records the case
\((\delta_0,\varepsilon_0,16)\) for a given penalty \(\kappa\).
\emph{Lean:}
\href{https://github.com/yang-mills-group/OpenAI-YM-20261028/blob/3456e5e75cf8d4cd558955a5c7638fba2d68f464/YM/Branch/Pick.lean\#L137-L142}{\texttt{Axial.viable}},
\href{https://github.com/yang-mills-group/OpenAI-YM-20261028/blob/3456e5e75cf8d4cd558955a5c7638fba2d68f464/YM/Branch/Uniform.lean\#L18-L56}{\texttt{RG.frame\_unique}},
\href{https://github.com/yang-mills-group/OpenAI-YM-20261028/blob/3456e5e75cf8d4cd558955a5c7638fba2d68f464/YM/Branch/Uniform.lean\#L58-L63}{\texttt{RG.predictor\_unique}},
\href{https://github.com/yang-mills-group/OpenAI-YM-20261028/blob/3456e5e75cf8d4cd558955a5c7638fba2d68f464/YM/Branch/Pick.lean\#L119-L134}{\texttt{Axial.eligible}}.

\textbf{Proposition III.1.4 (axial gauge controlled by curvature).} Let
\(l\) be a list of distinct axes, \(r\in\mathbb N\), \(o\in B_r\), and
\(U\) any configuration on \(\mathbb Z^4\). Let
\(S=\{y\in B_r:\ y_j=o_j\text{ for }j\notin l\}\) be the slab. Then
there is a root \(x\in S\) such that the axial gauge
\(g_y=T_U(x\to y;\,\text{axes in the order }l)\) satisfies

\[E_S(g\cdot U)\ \le\ 4\,|l|\,(2r+1)\,F_S(U).\]

\emph{Status:} theorem. \emph{Lean:}
\href{https://github.com/yang-mills-group/OpenAI-YM-20261028/blob/3456e5e75cf8d4cd558955a5c7638fba2d68f464/YM/Branch/Pick.lean\#L80-L105}{\texttt{Axial.existence}}.

\textbf{Proposition III.1.5 (least-squares synchronization: uniqueness
up to a global left multiple).} There is \(d=d(G)\in(0,1/4]\) with the
following property. Let \(n,\Lambda\) be finite nonempty sets. Let
\(W:n\times n\times\Lambda\to G\) satisfy \(W_{xxk}=1\) and
\(W_{yx\sigma(k)}=W_{xyk}^{-1}\) for some permutation \(\sigma\) of
\(\Lambda\), and \(\|W_{xyk}-1\|\le d\) for all indices. Then any two
global minimizers \(h,h'\) of
\(\mathrm{cost}(h)=\mathrm{avg}_{x,y,k}\,\mathrm{pr}(h_xW_{xyk}h_y^{-1})\)
over \(G^n\) satisfy \(h'=c\,h\) for a single \(c\in G\). The constant
\(d\) does not depend on \(|n|\) or \(|\Lambda|\). \emph{Status:}
theorem. \emph{Lean:}
\href{https://github.com/yang-mills-group/OpenAI-YM-20261028/blob/3456e5e75cf8d4cd558955a5c7638fba2d68f464/YM/Branch/Synchronize.lean\#L232-L272}{\texttt{Branch.align}},
with
\href{https://github.com/yang-mills-group/OpenAI-YM-20261028/blob/3456e5e75cf8d4cd558955a5c7638fba2d68f464/YM/Branch/Synchronize.lean\#L27-L31}{\texttt{Pairs}}.

\textbf{Proposition III.1.6 (synchronization of the realized frame with
any gauge).} There is \(A=A(G)>0\) with the following property. Let
\(P\) be any patch and \(K\) any kernel constants; no branch condition
is needed. Let \(W\) be a configuration on \(\mathbb Z^4\), \(t\) a
frame tag with \(m^{\rm fr}(W,t)>0\), \(h\) the realized frame, and
\(H:\mathbb Z^4\to G\) any gauge. Then there is \(c\in G\) with

\[\sum_{y\in B_R}\mathrm{ch}\big(c^{-1}h_yH_y^{-1}\big)\ \le\ A\,(2R+1)\Big(E_{B_R}(H\cdot W)+a\,(2R+1)\,F_{B_R}(W)\Big).\]

\emph{Status:} theorem. \emph{Lean:}
\href{https://github.com/yang-mills-group/OpenAI-YM-20261028/blob/3456e5e75cf8d4cd558955a5c7638fba2d68f464/YM/Pressure/Gauge.lean\#L135-L162}{\texttt{Pressure.synchronized}}.

\textbf{Proposition III.1.7 (one gauge for a neighbourhood of blocks).}
Fix \(G\), \(K\) and \(j\ge4\), and put \(\ell_j=2j+15\). Then
\(D=16\,\ell_j\) and \(C=A\,\ell_j(D+a\,\ell_j)\), with \(A\) from
III.1.6, have the following property. Let \(S\) be any step with kernel
constants \(K\) and \(J=j\) (any \(L\), any \(c\)). Let \(U\) be fine,
\(x\) a coarse site, and \(t\) a tag that is \emph{usable} at \(x\),
meaning that \(m^{\rm fr}(U^{(x+y)},t_{x+y})>0\) for all \(y\in B_2\).
Then there are a gauge \(H:\mathbb Z^4\to G\) and pads \(p_y\in G\)
(\(y\in B_2\)) with

\[E_{B_{R'}}(H\cdot U^{(x)})\le D\,L\,\mathrm{port}(x,U),\qquad \sum_{i\in B_R}\mathrm{ch}\big(p_y^{-1}h^{(x+y)}_i\,H_{i+Ly}^{-1}\big)\le C\,L^2\,\mathrm{port}(x,U)\quad(y\in B_2).\]

Here \(R'=R+4L=(J+7)L\) and \(\mathrm{port}(x,U)=F_{B_{R'}}(U^{(x)})\)
is the fine curvature in the macro-box around \(x\). Every tag with
\(\lambda(U,t)>0\) is usable at every coarse site. \emph{Status:}
theorem. \emph{Lean:}
\href{https://github.com/yang-mills-group/OpenAI-YM-20261028/blob/3456e5e75cf8d4cd558955a5c7638fba2d68f464/YM/Pressure/Cost.lean\#L72-L122}{\texttt{Pressure.planting}},
\href{https://github.com/yang-mills-group/OpenAI-YM-20261028/blob/3456e5e75cf8d4cd558955a5c7638fba2d68f464/YM/Pressure/Cost.lean\#L62-L63}{\texttt{usable}},
\href{https://github.com/yang-mills-group/OpenAI-YM-20261028/blob/3456e5e75cf8d4cd558955a5c7638fba2d68f464/YM/Pressure/Cost.lean\#L65-L69}{\texttt{wraps}},
\href{https://github.com/yang-mills-group/OpenAI-YM-20261028/blob/3456e5e75cf8d4cd558955a5c7638fba2d68f464/YM/Pressure/Cost.lean\#L60}{\texttt{port}},
\href{https://github.com/yang-mills-group/OpenAI-YM-20261028/blob/3456e5e75cf8d4cd558955a5c7638fba2d68f464/YM/Pressure/Cost.lean\#L37}{\texttt{macroR}}.

\textbf{Proposition III.1.8 (predicted versus actual coarse curvature).}
For \(G\), \(K\), \(j\) let \(C_8=8\max(1,\varepsilon^{-1})(2C+D)\),
with \(C,D\) from III.1.7. Let \(S\) be any step with constants \(K\)
and \(J=j\). Then for all \(U\) and \(V\), every coarse site \(x\),
every tag \(t\) usable at \(x\) and every plane \(p\):

\[\mathrm{ch}\big(Q(U,t)_{(x,p)}\big)\le \frac{C_8}{L^2}\,\mathrm{port}(x,U),\qquad \mathrm{ch}\big(V_{(x,p)}\big)\le \mathrm{fuel}(Q,V)_{(x,p)}+\frac{C_8}{L^2}\,\mathrm{port}(x,U).\]

Here \(Q_{(x,p)}\) denotes the plaquette holonomy, and
\(\mathrm{fuel}(Q,V)_{(x,p)}=\sum_{e\in\partial(x,p)}\mathrm{ch}(V_eQ_e^{-1})\)
is the sum of the four \emph{wires} of the plaquette, counted with
multiplicity even when edges coincide on a short torus. The macro-box
contains \(6(2R^{\prime})^2(2R^{\prime}+1)^2\) internal plaquettes. Thus
\(L^{-4}\mathrm{port}\) differs from their average chord by a factor
bounded above and below in terms of \(J\) only. The bound therefore says
that the coarse predicted plaquette is at most \(C(G,K,J)\,L^2\) times
the average fine curvature: the area scaling of the non-abelian Stokes
formula, with no further loss in \(L\). \emph{Status:} theorem.
\emph{Lean:}
\href{https://github.com/yang-mills-group/OpenAI-YM-20261028/blob/3456e5e75cf8d4cd558955a5c7638fba2d68f464/YM/Pressure/Cost.lean\#L210-L252}{\texttt{Pressure.qdc}},
\href{https://github.com/yang-mills-group/OpenAI-YM-20261028/blob/3456e5e75cf8d4cd558955a5c7638fba2d68f464/YM/Pressure/Cost.lean\#L189-L190}{\texttt{fuel}},
\href{https://github.com/yang-mills-group/OpenAI-YM-20261028/blob/3456e5e75cf8d4cd558955a5c7638fba2d68f464/YM/Pressure/Cost.lean\#L187}{\texttt{wire}}.

\textbf{Proposition III.1.9 (effective densities and RG histories).} The
following definitions carry four elementary consequences.

\begin{enumerate}
\def\labelenumi{(\alph{enumi})}
\item
  \emph{Budget points, flags, masks.} A budget point is a tuple (integer
  power \(\ge2\), integer order \(\ge6\), support exponent \(>0\), cover
  size \(B\ge1\), threshold \(t\in(0,1]\)). The flagged sites
  \(\mathrm{flags}_t(U)\) are the sites touched by a plaquette of chord
  \(>t\). \(U\) is \emph{healthy} on \(X\) if
  \(X\cap\mathrm{flags}_t(U)=\emptyset\), and the masked function is
  \(\mathbf 1_{\{\text{healthy on }X\}}F\). A function is \emph{regular}
  on \(X\) if it depends only on links in \(X\), is Borel and gauge
  invariant, and at each configuration healthy on \(X\) is \(C^q\) in
  exponential left increments, where \(q\) is the order; its ordered
  left derivatives through order \(q\) must exist. Its \emph{sup bounds}
  control derivatives of orders \(0,\ldots,q\) in directions with
  \(|A_e|\le t(1+\mathrm{dist}(e^-,o))^{\rm power}\), where \(e^-\) is
  the base of the positive link. \emph{Theorem:} a function that reads
  the links in \(X\) and the flags in \(X\) depends only on the links in
  \(\mathrm{inspected}(X)\), the union of \(X\) and the corners of all
  plaquettes meeting \(X\).
\item
  \emph{Activities.} There are regular activities \(E_i\), with an
  address (a site \(o_i\) in a support \(X_i\)), and forced activities
  \(K_i\), with a support \(Y_i\), an inventory
  \(\emptyset\ne I_i\subset Y_i\) and an extra test. \emph{Validity} at
  bounds \(\delta_E,w_K\ge0\) means the following:
\end{enumerate}

\begin{itemize}
\tightlist
\item
  \(E_i\) is regular on \(X_i\) and \(E_i(1)=0\);
\item
  if \(10\,\mathrm{dist}(y,o_i)<n\) for every \(y\in X_i\), where \(n\)
  is the torus side, the Hessian of \(E_i\) at the trivial configuration
  vanishes on affine directions
  \(A_{(y,\mu)}=a_\mu+\sum_\nu\widetilde{(y-o_i)}_\nu d_{\mu\nu}\),
  using the canonical centred lift \(\widetilde{(y-o_i)}\);
  equivalently, \(D^2[E_i(\exp A)]_{A=0}[A,B]=0\) for every two such
  directions \(A,B\);
\item
  \(K_i\) and its extra test are gauge invariant, Borel, and depend only
  on links and flags in \(Y_i\); \(K_i\) vanishes unless
  \(Y_i\cap\mathrm{flags}_t=I_i\) and the test holds;
\item
  covariance under every element of \(\mathfrak B_4\) and every base
  point (index permutations transporting addresses, inventories, tests
  and values);
\item
  there are nonnegative majorants \(e_i,k_i\), with \(e_i\) bounding the
  regular derivatives and \(|K_i|\le k_i\), such that, for every site
  \(x\),
\end{itemize}

\[\sum_{i:o_i=x}e^{\sigma\,\mathrm{cov}_B(X_i)}e_i\le\delta_E,\qquad \sum_{i:x\in Y_i}e^{\sigma\,\mathrm{cov}_B(Y_i)}k_i\le w_K.\]

Here \(\sigma\) is the support exponent and \(\mathrm{cov}_B(X)\) is the
least number of radius-\(B\) balls with connected union covering \(X\);
both series are required to converge.

The \emph{gas} is \(\sum_b\prod_{i\in b}K_i(U)\), summed over finite
families \(b\) with pairwise disjoint supports whose inventories cover
\(\mathrm{flags}_t(U)\) exactly. \emph{Theorem:} the gas equals \(1\)
when \(\mathrm{flags}_t(U)=\emptyset\).

\begin{enumerate}
\def\labelenumi{(\alph{enumi})}
\setcounter{enumi}{2}
\item
  \emph{Leading action.} A clip has \(c_{\rm dir}>16\,c_{\rm lin}>0\)
  and a relaxation length. Path data consist of row series
  \(\ell_p(U)=\sum_kr_{p,k}\,\mathrm{slot}_k(U)\) (weighted
  parallel-transported curvatures) and loop, cubic and quartic atoms.
  Their bounds require exponentially weighted summability, invariance,
  the normalization
  \(c_{\rm dir}\delta_{pq}+\sum_k\hat s_k(p)\hat s_k(q)=\delta_{pq}\) of
  the zero-momentum symbols (the sums of row coefficients at each
  plaquette orientation), and vanishing zeroth and first moments of the
  loop coefficients. The row term at a plaquette based at \(x\) is kept
  exactly when
  \(\mathrm{ch}(U_q)\le t\exp(\mathrm{dist}(x,q^-)/r_{\rm rel})\) for
  every plaquette \(q\), where \(r_{\rm rel}>0\) is the relaxation
  length and \(q^-\) its base. The per-plaquette term is
  \(\mathrm{clip}_t(\mathrm{ch}U_p)+\mathbf 1_{\rm kept}\tfrac12|\ell_p(U)|^2\),
  where \(\mathrm{clip}_t(s)=c_{\rm dir}s^2/2\) for \(s\le t\) and
  \(c_{\rm lin}ts\) for \(s>t\). The atoms are masked by health on their
  path and plaquette footprints. The leading action \(\mathcal S\) is
  the \(\mathfrak B_4\)-average of
  \(\beta\,[\sum_p(\cdot)+\text{cubic}+\text{quartic}]+\text{loop}\). A
  representation
  \(\mathcal R=(\text{path data, clip, activities}, p,\beta,v)\)
  evaluates to
  \(\mathrm{eval}_{\mathcal R}=\exp\big(v|\Lambda|-\mathcal S-\sum_i\mathbf 1_{\rm healthy}E_i\big)\cdot\mathrm{gas}\).
  \emph{Theorem:} if \(\mathcal R\) is lawful (path bounds and valid
  activities), then \(\mathrm{eval}_{\mathcal R}\) is Borel and bounded
  on that finite torus.
\item
  \emph{Histories.}
\end{enumerate}

\begin{itemize}
\tightlist
\item
  A \emph{calendar} holds torus sides \(n_{m,p}\) and factors \(L_m\),
  with \(n_{m+1,p}=L_m n_{m,p}\) and
  \(\mathrm{scale}(m)=\sum_{k<m}\log L_k\), the kernel constants, the
  barrier \(\bar\beta>0\), stage constants, orders \(J_m\ge4\) and
  exponents \(\theta\). It sets
  \(\hat\beta_m=\bar\beta+\kappa_{\rm slope}\,\mathrm{scale}(m)\),
  \(t_m=\mathrm{temp}(\hat\beta_m)\) and cover size
  \(\lceil(\log\hat\beta_m)^2\rceil\).
\item
  Step \(m\) goes from level \(m+1\) to level \(m\), with patch
  \((L_m,J_{m+1})\).
\item
  The \emph{fits} conditions include: \(J_m\) nondecreasing with
  \(J_{m+1}\le J_m+1\) and \(J_m\to\infty\); \(\hat\beta_m>1\);
  \(t_m\le1\); the level-\(m\) torus side at least
  \(100(\text{cover size}_m+1)\) and the fine side of step \(m\) at
  least \(100(R_m+1)\); drift and logarithmic-cost bounds.
\item
  \emph{Levels} means a calendar with its fits conditions and block
  control at every step.
\item
  A history assigns a representation \(\mathcal R_{N,p,m}\) to each mesh
  index \(N\), volume \(p\) and level \(m\le N\). It is \emph{selected}
  if: the path data do not depend on \(p\); each \(\mathcal R_{N,p,m}\)
  has the assigned level-\(m\) clip and budget point and is lawful with
  the level-\(m\) constants and
  \(\beta_{N,p,m}\in[\hat\beta_m/2,2\hat\beta_m]\); \(v_{N,p,N}=0\) and
  \(\mathrm{eval}_{N,p,N}\) is the bare Wilson density;
  \(\mathrm{eval}_{N,p,m}=\mathcal P_{\beta_{N,p,m+1}}\mathrm{eval}_{N,p,m+1}\)
  with
  \(|\beta_{m+1}-\beta_m-\kappa_{\rm slope}\log L_m|\le\mathrm{drift}_{m+1}\);
  \(\beta_{N,p,0}=\bar\beta\); and
  \(|\beta_{N,p,m}-\hat\beta_m|=o(\mathrm{scale}(m))\) uniformly in
  \(N,p\).
\end{itemize}

The lawful bounds at level \(m\) use the stage path width and row and
higher-atom bound, the loop bound
\(2\,\mathrm{monitor}_m^{\theta(J_m)}\), and activity bounds
\(\delta_E=\hat\beta_m^{-11/10}\) and
\(w_K=\exp[-\mathrm{monitor}_m^{1/5}]\), where
\(\mathrm{monitor}_m=\exp[(\log\log\hat\beta_m)^2]\).

The remaining fits conditions are: \(\log L_m/\mathrm{scale}(m)\to0\);
\(L_m\) is the factor assigned to order \(J_{m+1}\), and the assigned
factors are nondecreasing; the derivative orders, powers and support
exponents are nondecreasing along the levels; \(\theta\) is antitone
with \(0<\theta(J_m)\le J_m^{-1}\); and

\[\frac{\mathrm{drift}_{m+1}}{\log L_m}\le J_{m+1}^{-2},\qquad \frac{\log(1+C_{m+1}+\sigma_{m+1}^{-1}+C^{\rm col}_{m+1}+r_{m+1}+s_{m+1})}{\log L_m}\le J_{m+1}^{-2}.\]

Here \(C,\sigma,C^{\rm col},r,s\) are the stage path bound, path width,
column bound, relaxation length and support exponent.

\emph{Theorem:} for a selected history,
\(\int\mathrm{eval}_{N,p,m}\,d\mu=Z_{N,p}(\beta_{N,p,N})\), the Wilson
partition function, for every \(m\le N\).

\emph{Status:} definitions. The four theorems are proved. The
quantitative path and activity bounds needed to construct such histories
belong to III.3--III.10. \emph{Lean:}
\href{https://github.com/yang-mills-group/OpenAI-YM-20261028/blob/3456e5e75cf8d4cd558955a5c7638fba2d68f464/YM/RG/Masks.lean\#L125-L136}{\texttt{BudgetPoint}},
\href{https://github.com/yang-mills-group/OpenAI-YM-20261028/blob/3456e5e75cf8d4cd558955a5c7638fba2d68f464/YM/RG/Masks.lean\#L154-L161}{\texttt{BudgetPoint.Regular}},
\href{https://github.com/yang-mills-group/OpenAI-YM-20261028/blob/3456e5e75cf8d4cd558955a5c7638fba2d68f464/YM/RG/Masks.lean\#L34-L36}{\texttt{flags}},
\href{https://github.com/yang-mills-group/OpenAI-YM-20261028/blob/3456e5e75cf8d4cd558955a5c7638fba2d68f464/YM/RG/Masks.lean\#L108-L110}{\texttt{UsesFinal.allReads}},
\href{https://github.com/yang-mills-group/OpenAI-YM-20261028/blob/3456e5e75cf8d4cd558955a5c7638fba2d68f464/YM/RG/Activities.lean\#L62-L66}{\texttt{Activities}},
\href{https://github.com/yang-mills-group/OpenAI-YM-20261028/blob/3456e5e75cf8d4cd558955a5c7638fba2d68f464/YM/RG/Activities.lean\#L70-L75}{\texttt{RegularFamily}},
\href{https://github.com/yang-mills-group/OpenAI-YM-20261028/blob/3456e5e75cf8d4cd558955a5c7638fba2d68f464/YM/RG/Activities.lean\#L77-L84}{\texttt{ForcedFamily}},
\href{https://github.com/yang-mills-group/OpenAI-YM-20261028/blob/3456e5e75cf8d4cd558955a5c7638fba2d68f464/YM/RG/Activities.lean\#L87-L96}{\texttt{Moves}},
\href{https://github.com/yang-mills-group/OpenAI-YM-20261028/blob/3456e5e75cf8d4cd558955a5c7638fba2d68f464/YM/RG/Activities.lean\#L106-L111}{\texttt{Valid}},
\href{https://github.com/yang-mills-group/OpenAI-YM-20261028/blob/3456e5e75cf8d4cd558955a5c7638fba2d68f464/YM/RG/Activities.lean\#L114-L116}{\texttt{admissible}},
\href{https://github.com/yang-mills-group/OpenAI-YM-20261028/blob/3456e5e75cf8d4cd558955a5c7638fba2d68f464/YM/RG/Activities.lean\#L129}{\texttt{gas}},
\href{https://github.com/yang-mills-group/OpenAI-YM-20261028/blob/3456e5e75cf8d4cd558955a5c7638fba2d68f464/YM/RG/Activities.lean\#L135-L147}{\texttt{gas\_one}},
\href{https://github.com/yang-mills-group/OpenAI-YM-20261028/blob/3456e5e75cf8d4cd558955a5c7638fba2d68f464/YM/RG/PathActions.lean\#L139-L146}{\texttt{Clip}},
\href{https://github.com/yang-mills-group/OpenAI-YM-20261028/blob/3456e5e75cf8d4cd558955a5c7638fba2d68f464/YM/RG/PathActions.lean\#L150-L154}{\texttt{PathData}},
\href{https://github.com/yang-mills-group/OpenAI-YM-20261028/blob/3456e5e75cf8d4cd558955a5c7638fba2d68f464/YM/RG/PathActions.lean\#L173-L185}{\texttt{PathData.Bounds}},
\href{https://github.com/yang-mills-group/OpenAI-YM-20261028/blob/3456e5e75cf8d4cd558955a5c7638fba2d68f464/YM/RG/PathActions.lean\#L221-L224}{\texttt{baseTerm}},
\href{https://github.com/yang-mills-group/OpenAI-YM-20261028/blob/3456e5e75cf8d4cd558955a5c7638fba2d68f464/YM/RG/PathActions.lean\#L236-L238}{\texttt{principal}},
\href{https://github.com/yang-mills-group/OpenAI-YM-20261028/blob/3456e5e75cf8d4cd558955a5c7638fba2d68f464/YM/RG/PathActions.lean\#L241-L242}{\texttt{inChannels}},
\href{https://github.com/yang-mills-group/OpenAI-YM-20261028/blob/3456e5e75cf8d4cd558955a5c7638fba2d68f464/YM/RG/History.lean\#L165-L171}{\texttt{Rep}},
\href{https://github.com/yang-mills-group/OpenAI-YM-20261028/blob/3456e5e75cf8d4cd558955a5c7638fba2d68f464/YM/RG/History.lean\#L179}{\texttt{Rep.eval}},
\href{https://github.com/yang-mills-group/OpenAI-YM-20261028/blob/3456e5e75cf8d4cd558955a5c7638fba2d68f464/YM/RG/History.lean\#L195-L198}{\texttt{Rep.envelope}},
\href{https://github.com/yang-mills-group/OpenAI-YM-20261028/blob/3456e5e75cf8d4cd558955a5c7638fba2d68f464/YM/RG/History.lean\#L64-L73}{\texttt{Calendar}},
\href{https://github.com/yang-mills-group/OpenAI-YM-20261028/blob/3456e5e75cf8d4cd558955a5c7638fba2d68f464/YM/RG/History.lean\#L97-L118}{\texttt{Calendar.Fits}},
\href{https://github.com/yang-mills-group/OpenAI-YM-20261028/blob/3456e5e75cf8d4cd558955a5c7638fba2d68f464/YM/RG/History.lean\#L131-L134}{\texttt{Levels}},
\href{https://github.com/yang-mills-group/OpenAI-YM-20261028/blob/3456e5e75cf8d4cd558955a5c7638fba2d68f464/YM/RG/History.lean\#L141-L156}{\texttt{Levels.budget}},
\href{https://github.com/yang-mills-group/OpenAI-YM-20261028/blob/3456e5e75cf8d4cd558955a5c7638fba2d68f464/YM/RG/History.lean\#L204-L206}{\texttt{History}},
\href{https://github.com/yang-mills-group/OpenAI-YM-20261028/blob/3456e5e75cf8d4cd558955a5c7638fba2d68f464/YM/RG/History.lean\#L234-L240}{\texttt{History.Selected}},
\href{https://github.com/yang-mills-group/OpenAI-YM-20261028/blob/3456e5e75cf8d4cd558955a5c7638fba2d68f464/YM/RG/History.lean\#L242-L249}{\texttt{History.partition\_eq}}.

\textbf{Proposition III.1.10 (the tube identity).} A \emph{tube} for a
step is a tower of transfer kernels together with path data and clips at
the fine level (up) and at the coarse level (down). Fix a path-length
cutoff, a fine threshold \(t\) and a coarse threshold \(s\). From \(V\)
one builds:

\begin{itemize}
\tightlist
\item
  the \emph{vision} \(w\) (the thinned coarse curvature and row data);
\item
  linear maps \(\Phi,\Phi',\mathcal N\), assembled from the adjoint
  action of parallel transports of \(V\) along cut-off routes;
\item
  the fine data \(X_f(U)\), the coarse data \(X_n(V)\), and the noise
  field \(Z_b=\sqrt\kappa\,\mathrm{curv}(V_bQ_b^{-1})\).
\end{itemize}

Use counting inner products and norms on these finite arrays. Let
\(C_f\) and \(C_c\) be the sums of the fine and coarse per-plaquette
crowns, and let \(P=\sum_b\tfrac\kappa2\,\mathrm{ch}(V_bQ_b^{-1})^2\). A
crown is the clipped direct term plus half the kept row's squared
physical norm, with the row truncated at the path cutoff. The totals of
the six local works are

\[\begin{aligned}
W_F&=C_f+\tfrac12\|\Phi w\|^2-\langle\Phi w,X_f\rangle,&
W_I&=P+\tfrac12\|\Phi'w\|^2-\langle\Phi'w,Z\rangle,\\
W_C&=\langle w,X_n\rangle-\tfrac12\|w\|^2-C_c,&
W_L&=\langle w,\Phi^\dagger X_f+\Phi'^\dagger Z-X_n\rangle,\\
W_N&=\tfrac12\|\mathcal Nw\|^2,&
W_D&=\tfrac12\langle w,(1-\Phi^\dagger\Phi-\Phi'^\dagger\Phi'-\mathcal N^\dagger\mathcal N)w\rangle .
\end{aligned}\]

For every path cutoff, thresholds \(t,s\in\mathbb R\) and configurations
\(U,Q,V\),

\[W_F+W_I+W_C+W_N+W_L+W_D=C_f+P-C_c.\]

The fine and coarse data have two components per plaquette:
\(\sqrt{c_{\rm dir}}\,\mathrm{curv}(U_p)\) and the truncated row, using
the respective clip and configuration. The fine transfer uses parallel
transports of the blockwise copy of \(V\); the other two maps use \(V\)
on the coarse torus. The vision caps the direct component by
\(\min(1,s/\mathrm{ch}(V_p))\), taking this factor to be \(1\) at zero
chord, and masks the row by the kept condition. For the local works,
assign each square or inner-product summand to its array coordinate. The
two components at each plaquette are summed together: \(W_F\) is indexed
by fine plaquettes, \(W_I\) by coarse links, \(W_C,W_L,W_D\) by coarse
plaquettes, and \(W_N\) by coarse gain coordinates. The crowns and
penalties are charged at the same indices.

\emph{Status:} theorem (an algebraic identity). \emph{Lean:}
\href{https://github.com/yang-mills-group/OpenAI-YM-20261028/blob/3456e5e75cf8d4cd558955a5c7638fba2d68f464/YM/Pressure/Rows.lean\#L94-L99}{\texttt{Tube}},
\href{https://github.com/yang-mills-group/OpenAI-YM-20261028/blob/3456e5e75cf8d4cd558955a5c7638fba2d68f464/YM/Pressure/Rows.lean\#L145-L163}{\texttt{Tube.locates}},
\href{https://github.com/yang-mills-group/OpenAI-YM-20261028/blob/3456e5e75cf8d4cd558955a5c7638fba2d68f464/YM/Pressure/Rows.lean\#L122-L124}{\texttt{workF}},
\href{https://github.com/yang-mills-group/OpenAI-YM-20261028/blob/3456e5e75cf8d4cd558955a5c7638fba2d68f464/YM/Pressure/Rows.lean\#L138-L140}{\texttt{workD}}.

\textbf{Proposition III.1.11 (energy per seed: the requirement, and its
transfer to a history).}

\emph{The requirement.} Write \(\hat\beta\) for the coarse reference
coupling, \(t=\mathrm{temp}(\hat\beta+\kappa_{\rm slope}\log L)\) for
the fine-level threshold and \(\varrho=\lceil(\log\hat\beta)^2\rceil\)
for the cluster radius.

\begin{itemize}
\item
  A \emph{store} \(A\) is a finite set of fine plaquettes (costs),
  coarse links (aisles) and coarse sites (flags). Its size \(|A|\) is
  the sum of the three cardinalities. Its locations are the parent
  coarse sites of all corners of its fine plaquettes, both endpoints of
  its coarse links, and its coarse sites. Write \(D_r(A)\) for their
  radius-\(r\) neighbourhood, using the maximum-norm distance on the
  coarse torus.
\item
  \((A,U,k,V)\) is a \emph{plot} at \((\hat\beta,W)\) if the following
  conditions hold:

  \begin{itemize}
  \tightlist
  \item
    every fine plaquette whose parent-corner set meets
    \(D_{5\varrho}(A)\) has chord \(\ge t/4\) if it belongs to \(A\),
    and \(\le t/2\) otherwise;
  \item
    every coarse link with an endpoint in \(D_{5\varrho}(A)\) has wire
    \(\mathrm{ch}(V_bQ_b(U,k)^{-1})\ge Wt\) if it belongs to \(A\), and
    \(\le2Wt\) otherwise;
  \item
    on \(D_{5\varrho}(A)\), the coarse sites in \(A\) are exactly
    \(\mathrm{flags}_{\mathrm{temp}(\hat\beta)}(V)\);
  \item
    every site in \(D_{5\varrho}(A)\) has positive frame mass under
    \(k\).
  \end{itemize}
\item
  The \emph{estimator data} fix kernel constants, two stage records, an
  order \(j\), a decay width \(\eta>0\) and a bound \(C\). A tube is
  \emph{compatible} with these data if the kernel and order agree, the
  branch conditions hold, and its transfer kernels have decay bounds
  through \(j-2\) forward fine-coordinate differences: the fine transfer
  row has bound \(C/L^2\), and differences of order \(q\) of the
  background and curvature kernels have bounds \(C/L^{1+q}\) and
  \(C/L^{2+q}\). Its other two rows have bound \(C\). The clips, path
  widths and row and higher-atom bounds are those of the two stages,
  with some finite loop bounds; the path rows must also be attached to
  quadratic metrics satisfying the tower's transfer identities (III.3).
\item
  The \emph{energy bound} at \((L,W,\kappa_0)\) holds if for all
  sufficiently large \(\hat\beta\) the following inequality is true. The
  threshold for \(\hat\beta\) depends on the estimator data, \(L\),
  \(W\) and \(\kappa_0\), not on the torus. The inequality is required
  for every step with factor \(L\) and coarse side
  \(c\ge100(1+\varrho)\), every compatible tube, every forest, and every
  plot with \(U\equiv1\) on its trunks:

  \[\kappa_0\,t^2\,|A|\ \le\ \mathrm{Work}(A).\]

  Here \(\mathrm{Work}(A)\) sums the six local works of III.1.10 with
  path cutoff \(\lfloor\varrho/10^5\rfloor\): a fine plaquette is
  included when its base's parent lies in \(D_{2\varrho}(A)\), and a
  coarse plaquette, link or gain coordinate is included when its base
  lies there.
\item
  The estimator data satisfy the \emph{energy-per-seed requirement} if
  there is \(L_0\) such that for every \(L\ge L_0\) and every
  sufficiently large \(W\) there is \(\kappa_0>0\) for which the energy
  bound holds.
\end{itemize}

\emph{The transfer (theorem).} Let \(B\) be a directed history and
suppose the \emph{ringing} condition holds: for each pair of orders
\(j_{\min}\le i\le j\le i+1\), the inequality holds with some \(\mu>0\)
for every coarse reference coupling above the assigned entry threshold
and wire parameter \(F(j,i)>0\). Here \(j_{\min}=J_0\) is the first
admissible order. Then for each such pair there is \(\mu>0\) with the
following property. For all \(N\), all \(m<N\) with \(J_{m+1}=j\) and
\(J_m=i\), all \(p\), and every plot at \((\hat\beta_m,F(j,i))\) with
\(U\) trivial on the trunks, \[\mu\,t_{m+1}^2\,|A|\le\mathrm{Work}(A).\]
This uses \(\hat\beta_{m+1}=\hat\beta_m+\kappa_{\rm slope}\log L_m\), so
that \(t=\mathrm{temp}(\hat\beta_{m+1})=t_{m+1}\).

\emph{Status:} the requirement is a definition (a hypothesis of other
units). It is proved for every choice of estimator data in III.10
(\href{https://github.com/yang-mills-group/OpenAI-YM-20261028/blob/3456e5e75cf8d4cd558955a5c7638fba2d68f464/YM/Next/FLUCTUATIONS/Grants.lean\#L24}{\texttt{Next.Fluctuations.grants}}).
Its proof applies
\href{https://github.com/yang-mills-group/OpenAI-YM-20261028/blob/3456e5e75cf8d4cd558955a5c7638fba2d68f464/YM/Integrated.lean\#L222-L225}{\texttt{Integrated.pressure\_grants}},
which follows from the work items E01--E09 and from the coarse-row bound
\href{https://github.com/yang-mills-group/OpenAI-YM-20261028/blob/3456e5e75cf8d4cd558955a5c7638fba2d68f464/YM/Bid/Energy.lean\#L148}{\texttt{EBare}},
itself proved by
\href{https://github.com/yang-mills-group/OpenAI-YM-20261028/blob/3456e5e75cf8d4cd558955a5c7638fba2d68f464/YM/Next/FLUCTUATIONS/Bare.lean\#L64-L120}{\texttt{Fluctuations.bare}}.
The ringing condition is derived from these grants by raising the
\(\beta\)-entry thresholds
(\href{https://github.com/yang-mills-group/OpenAI-YM-20261028/blob/3456e5e75cf8d4cd558955a5c7638fba2d68f464/YM/Bid/Energy.lean\#L173-L202}{\texttt{Bid.pricing\_order}},
\href{https://github.com/yang-mills-group/OpenAI-YM-20261028/blob/3456e5e75cf8d4cd558955a5c7638fba2d68f464/YM/Next/FLUCTUATIONS/Grants.lean\#L32-L37}{\texttt{pressure\_pricing}},
III.10). The transfer is a theorem of this unit; its constant is uniform
in \(N,m,p\) for the fixed pair \((j,i)\). \emph{Lean:}
\href{https://github.com/yang-mills-group/OpenAI-YM-20261028/blob/3456e5e75cf8d4cd558955a5c7638fba2d68f464/YM/Pressure/Seeds.lean\#L34-L37}{\texttt{Stores}},
\href{https://github.com/yang-mills-group/OpenAI-YM-20261028/blob/3456e5e75cf8d4cd558955a5c7638fba2d68f464/YM/Pressure/Seeds.lean\#L51-L58}{\texttt{Stores.Plot}},
\href{https://github.com/yang-mills-group/OpenAI-YM-20261028/blob/3456e5e75cf8d4cd558955a5c7638fba2d68f464/YM/Pressure/Seeds.lean\#L85}{\texttt{Stores.assessed}},
\href{https://github.com/yang-mills-group/OpenAI-YM-20261028/blob/3456e5e75cf8d4cd558955a5c7638fba2d68f464/YM/Pressure/Seeds.lean\#L89-L96}{\texttt{Teller}},
\href{https://github.com/yang-mills-group/OpenAI-YM-20261028/blob/3456e5e75cf8d4cd558955a5c7638fba2d68f464/YM/Pressure/Seeds.lean\#L101-L109}{\texttt{Teller.Funds}},
\href{https://github.com/yang-mills-group/OpenAI-YM-20261028/blob/3456e5e75cf8d4cd558955a5c7638fba2d68f464/YM/Pressure/Seeds.lean\#L114-L118}{\texttt{Teller.live}},
\href{https://github.com/yang-mills-group/OpenAI-YM-20261028/blob/3456e5e75cf8d4cd558955a5c7638fba2d68f464/YM/Pressure/Seeds.lean\#L124-L125}{\texttt{Teller.cash}},
\href{https://github.com/yang-mills-group/OpenAI-YM-20261028/blob/3456e5e75cf8d4cd558955a5c7638fba2d68f464/YM/Pressure/Seeds.lean\#L129-L131}{\texttt{Teller.Grants}},
\href{https://github.com/yang-mills-group/OpenAI-YM-20261028/blob/3456e5e75cf8d4cd558955a5c7638fba2d68f464/YM/Pressure/Seeds.lean\#L155-L158}{\texttt{Testing.Stamp}},
\href{https://github.com/yang-mills-group/OpenAI-YM-20261028/blob/3456e5e75cf8d4cd558955a5c7638fba2d68f464/YM/Pressure/Seeds.lean\#L161-L163}{\texttt{Testing.Ringing}},
\href{https://github.com/yang-mills-group/OpenAI-YM-20261028/blob/3456e5e75cf8d4cd558955a5c7638fba2d68f464/YM/Pressure/Seeds.lean\#L166-L181}{\texttt{Testing.notation\_ticket}}.

\textbf{Proposition III.1.12 (kinetic specification of a directed
history).} Fix levels, a horizon of quadratic kernels, a local
exponential-chart density and positive stage constants. Kinetic data
assign, for each \(N,m\), a shift \(\alpha_{N,m}\) and path data, and,
for each volume and forest \(F\), a map \(\psi\) from coarse Lie-algebra
fields to Lie-algebra fields on the free fine links.

To state the stationary requirements, let \(\mathcal A_f^\circ\) be the
unmasked, unclipped fine action with coupling \(1\) and without loop
terms, averaged over signed symmetries. Write \(\mathcal L_f^\circ\) for
the corresponding loop action and \(\mathcal Q_c^\circ\) for the coarse
quadratic part. For a coarse field \(u\) and a free-link field \(b\),
let \(V=\exp u\) and form \(U_F(u,b)\) by putting \(1\) on trunks and
\(\exp(b_e)(\iota V)_e\) on free links, where \(\iota V\) is \(1\) on
internal links and copies \(V\) on links crossing blocks. For each tag
\(t\), set

\[\mathcal I_t(u,b)=\mathcal A_f^\circ(U_F(u,b))+\kappa\sum_e\mathrm{pr}\big(V_eQ_e(U_F(u,b),t)^{-1}\big).\]

The kinetic requirements are:

\begin{itemize}
\tightlist
\item
  \begin{enumerate}
  \def\labelenumi{(\roman{enumi})}
  \tightlist
  \item
    the stages and block factors agree with the fixed choices;
  \end{enumerate}
\item
  (ii) for every \(m\le N\), the path rows are attached to the horizon's
  quadratic kernels and satisfy the stage bounds of III.1.9(c);
\item
  (iii) at level \(N\), the unmasked, unclipped action is analytic and
  has the same derivatives of orders \(0,\ldots,4\) at \(0\) as the
  Wilson energy in exponential coordinates, and every loop coefficient
  vanishes;
\item
  (iv) for \(m<N\), with the prescribed positive stage constant
  \(\eta\), every volume and every forest satisfy
  \[D_b^2\mathcal I_t(0,0)[v,v]\ge\eta\sum_{e\ {\rm free}}|v_e|^2\] for
  every \(t,v\). The map \(\psi\) is analytic at \(0\), with
  \(\psi(0)=0\); each \(\mathcal I_t\) is analytic at \((0,0)\), and
  \(D_b\mathcal I_t(u,\psi(u))=0\) near \(0\). The derivative of
  \(\psi\) is the prescribed linear background, adjusted to tree gauge,
  minus the linearized copy \(\iota_1u\), restricted to free links.
  Explicitly, if \(\mathcal Bu\) is the horizon's linear background, put
  \[\chi_u(y)=\int_{y\to{\rm root}(y)}\mathcal Bu-(\mathcal F\mathcal Bu)({\rm parent}(y)),\qquad D\psi(0)u=(\mathcal Bu+d_0\chi_u-\iota_1u)|_{\rm free},\]
  where the integral follows the forest and \(\mathcal F\) is defined in
  III.1.14. Finally the unmasked coarse action
  \(\mathcal A_c^\circ(\exp u)\) is analytic at \(0\) and has the same
  derivatives as \(\mathcal I_t(u,\psi(u))\) of every order through
  \(4\), for every \(t\);
\item
  (v) for \(m<N\),
  \[|\alpha_{N,m}+\kappa_{\rm slope}\log L_m|+1\le\mathrm{drift}_{m+1};\]
\item
  (vi) for \(m<N\) and every volume and forest, let \(H_t(u)\) be the
  Hessian of \(\mathcal I_t(u,\cdot)\) at \(\psi(u)\), in the counting
  basis of the physical Lie-algebra coordinates, and define
  \[T_t(u)=\mathcal L_f^\circ(U_F(u,\psi(u)))+\tfrac12\log\det H_t(u)-\log\prod_{e\ {\rm free}}J(\psi(u)_e),\qquad D(u)=\alpha_{N,m}\mathcal Q_c^\circ(\exp u)+\mathcal L_c^\circ(\exp u).\]
  For every tag \(t\), \(\det H_t>0\) near \(0\), both \(T_t\) and \(D\)
  are analytic at \(0\), and
  \[|D^2T_t(0)[a,b]-D^2D(0)[a,b]|\le|\Lambda_c|\,\Theta\Lambda e^{-\gamma c}\]
  for any two real plane waves \(a,b\). Each wave is the real part of a
  lattice Fourier character times a complex amplitude whose every
  direction and colour component has modulus at most \(1\). Here
  \(\Theta,\gamma>0\) are the assigned stage constants, \(\Lambda\) is
  the fine loop bound, and \(c\) is the coarse torus side.
\end{itemize}

A \emph{directed history} consists of paid levels (III.2), a horizon
satisfying the decay, positivity and transfer identities of III.3,
kinetic data satisfying (i)--(vi), and a recipe with the same path data.
It also carries a local Haar-chart formula: on some ball
\(B_r\subset\mathfrak g\), the exponential is injective with open image,
\(J\) is positive and analytic, \(J(0)=1\), and the pushforward of
\(c_JJ(X)\,dX\) is Haar measure restricted to that image, for some
\(c_J>0\); \(J\) is Borel globally.

A \emph{move} between a fine representation \(\mathcal R_{m+1}\) and a
coarse representation \(\mathcal R_m\) requires the exact push identity,
the recipe's continuation and retention conditions, and agreement of the
coarse path data, clip, norm parameters and lawful bounds with level
\(m\). It also requires a forced-activity bound
\(\exp[-g_j\,\mathrm{monitor}_{m+1}^2]\), where \(g_j>0\) is prescribed
and \(j=J_{m+1}\), and \[|\beta_m-\beta_{m+1}-\alpha_{N,m}|\le1.\]
Together with (v), this implies the drift bound of III.1.9(d).

\emph{Status:} definitions of requirements. For simple nonabelian
\(\mathfrak g\), the conditional construction theorems
\href{https://github.com/yang-mills-group/OpenAI-YM-20261028/blob/3456e5e75cf8d4cd558955a5c7638fba2d68f464/YM/FullVolume/Roll/Linger.lean\#L23-L51}{\texttt{Roost.companyK}}
and
\href{https://github.com/yang-mills-group/OpenAI-YM-20261028/blob/3456e5e75cf8d4cd558955a5c7638fba2d68f464/YM/Integration/NativeSelection/FundedCompanySelection.lean\#L49-L106}{\texttt{FundedCompanySelection.company}}
supply directed histories with prescribed paid levels as part of the
assembly in III.2. The kinetic requirements themselves do not assert a
coupling reset.

\emph{Lean:}
\href{https://github.com/yang-mills-group/OpenAI-YM-20261028/blob/3456e5e75cf8d4cd558955a5c7638fba2d68f464/YM/PCalc/Chart.lean\#L100-L101}{\texttt{PriceGap}},
\href{https://github.com/yang-mills-group/OpenAI-YM-20261028/blob/3456e5e75cf8d4cd558955a5c7638fba2d68f464/YM/PCalc/Chart.lean\#L88-L91}{\texttt{Critical}},
\href{https://github.com/yang-mills-group/OpenAI-YM-20261028/blob/3456e5e75cf8d4cd558955a5c7638fba2d68f464/YM/PCalc/Chart.lean\#L97-L98}{\texttt{Guessed}},
\href{https://github.com/yang-mills-group/OpenAI-YM-20261028/blob/3456e5e75cf8d4cd558955a5c7638fba2d68f464/YM/PCalc/Chart.lean\#L108-L109}{\texttt{ClassicalMatch}},
\href{https://github.com/yang-mills-group/OpenAI-YM-20261028/blob/3456e5e75cf8d4cd558955a5c7638fba2d68f464/YM/PREC/TreeMatch.lean\#L26-L27}{\texttt{extracted}},
\href{https://github.com/yang-mills-group/OpenAI-YM-20261028/blob/3456e5e75cf8d4cd558955a5c7638fba2d68f464/YM/PREC/Stacks.lean\#L132-L139}{\texttt{Horizon.Suffices}},
\href{https://github.com/yang-mills-group/OpenAI-YM-20261028/blob/3456e5e75cf8d4cd558955a5c7638fba2d68f464/YM/PCalc/Loop.lean\#L74-L77}{\texttt{Kin}},
\href{https://github.com/yang-mills-group/OpenAI-YM-20261028/blob/3456e5e75cf8d4cd558955a5c7638fba2d68f464/YM/PCalc/Loop.lean\#L93-L115}{\texttt{Kin.Planned}},
\href{https://github.com/yang-mills-group/OpenAI-YM-20261028/blob/3456e5e75cf8d4cd558955a5c7638fba2d68f464/YM/PCalc/Loop.lean\#L62-L71}{\texttt{FiniteBound}},
\href{https://github.com/yang-mills-group/OpenAI-YM-20261028/blob/3456e5e75cf8d4cd558955a5c7638fba2d68f464/YM/PCalc/Loop.lean\#L46-L49}{\texttt{thermal}},
\href{https://github.com/yang-mills-group/OpenAI-YM-20261028/blob/3456e5e75cf8d4cd558955a5c7638fba2d68f464/YM/PCalc/Loop.lean\#L42-L44}{\texttt{hessian}},
\href{https://github.com/yang-mills-group/OpenAI-YM-20261028/blob/3456e5e75cf8d4cd558955a5c7638fba2d68f464/YM/PCalc/Loop.lean\#L28-L38}{\texttt{LocalRatio}},
\href{https://github.com/yang-mills-group/OpenAI-YM-20261028/blob/3456e5e75cf8d4cd558955a5c7638fba2d68f464/YM/PCalc/Band.lean\#L26-L36}{\texttt{Directed}},
\href{https://github.com/yang-mills-group/OpenAI-YM-20261028/blob/3456e5e75cf8d4cd558955a5c7638fba2d68f464/YM/PCalc/Band.lean\#L49-L55}{\texttt{Directed.Move}}.

\textbf{Proposition III.1.13 (tree gauge on blocks).} Functions in
(c)--(d) take values in a real Banach space. A \emph{forest} assigns to
every non-root fine site \(y\) (one with
\(y\ne L\cdot\mathrm{parent}(y)\)) a unit step \(\mathrm{up}(y)\) and a
height \(\le4L\), such that the height decreases along \(\mathrm{up}\)
and the parent block does not change. The \emph{trunks} are the links
used by these steps; the others are \emph{free}.

\begin{itemize}
\tightlist
\item
  \begin{enumerate}
  \def\labelenumi{(\alph{enumi})}
  \tightlist
  \item
    For odd \(L\) a forest exists: choose an axis on which the offset
    from the block root is nonzero and step towards the root, with the
    \(\ell^1\) norm of the offset as height.
  \end{enumerate}
\item
  (b) Gauging by
  \(\mathrm{fold}_y=T_U(y\to\text{root along the tree})^{-1}\) makes
  every trunk link equal to \(1\). Moreover \(\mathrm{fold}\) is \(1\)
  at the roots and depends only on the trunk links.
\item
  (c) If \(F\in L^1(\mu_f)\) is invariant under gauge transformations
  equal to \(1\) at the block roots, then
  \(\int F\,d\mu_f=\int F(\mathrm{drape}\,s)\,d\mu_{\rm free}(s)\). Here
  \(\mathrm{drape}\,s\) has trunk links \(1\) and free links \(s\), and
  \(\mu_{\rm free}\) is product Haar measure on the free links.
\item
  (d) Hence, if \(\rho_\beta\) is a kernel satisfying III.1.2(a) and
  \(F\) is gauge invariant and integrable, then
  \((\mathcal P_\beta F)(V)=\int\rho_\beta(\mathrm{drape}\,s,V)\,F(\mathrm{drape}\,s)\,d\mu_{\rm free}(s)\).
\end{itemize}

\emph{Status:} theorem. \emph{Lean:}
\href{https://github.com/yang-mills-group/OpenAI-YM-20261028/blob/3456e5e75cf8d4cd558955a5c7638fba2d68f464/YM/Tree/Pines.lean\#L24-L29}{\texttt{Tree.Forest}},
\href{https://github.com/yang-mills-group/OpenAI-YM-20261028/blob/3456e5e75cf8d4cd558955a5c7638fba2d68f464/YM/Tree/Pines.lean\#L33-L72}{\texttt{Tree.select}},
\href{https://github.com/yang-mills-group/OpenAI-YM-20261028/blob/3456e5e75cf8d4cd558955a5c7638fba2d68f464/YM/Tree/Cut.lean\#L32-L37}{\texttt{Forest.compress}},
\href{https://github.com/yang-mills-group/OpenAI-YM-20261028/blob/3456e5e75cf8d4cd558955a5c7638fba2d68f464/YM/Tree/Cut.lean\#L39-L45}{\texttt{Forest.names}},
\href{https://github.com/yang-mills-group/OpenAI-YM-20261028/blob/3456e5e75cf8d4cd558955a5c7638fba2d68f464/YM/Tree/Cut.lean\#L81-L111}{\texttt{Forest.pack}},
\href{https://github.com/yang-mills-group/OpenAI-YM-20261028/blob/3456e5e75cf8d4cd558955a5c7638fba2d68f464/YM/Tree/Cut.lean\#L114-L120}{\texttt{Forest.pushing}}.

\textbf{Proposition III.1.14 (prediction near the trivial
configuration).} Assume (B1)--(B3). For a \(\mathfrak g\)-valued 1-form
\(v\) on \(\Lambda_f\), let \(U_v=(\exp v_e)_e\). Then for every tag
\(t\) and coarse link \(b=(x,i)\):

\begin{itemize}
\tightlist
\item
  \begin{enumerate}
  \def\labelenumi{(\alph{enumi})}
  \tightlist
  \item
    \(Q_b(U_0,t)=1\);
  \end{enumerate}
\item
  (b) \(v\mapsto Q_b(U_v,t)\) is real analytic near \(v=0\); in a
  sufficiently small neighbourhood all patch frames and bond predictors
  use their minimizing branches;
\item
  (c) the derivatives at \(v=0\) are, componentwise in \(\mathfrak g\),
\end{itemize}

\[\tfrac{d}{ds}h^{(x)}_y(U_{sv})\big|_{0}=\tau_x(y)-\tau_x(0),\qquad \tfrac{d}{ds}Q_b(U_{sv},t)\big|_{0}=(\mathcal Mv+d_0\mathcal Fv)(b).\]

Here \(d_0\varphi(x,i)=\varphi(x+e_i)-\varphi(x)\) and
\(\tau_x(y)=-\mathrm{avg}_{z\in B_R}\mathrm{avg}_{o}\int_{\mathrm{route}(Lx+y\to Lx+z;\,o)}v\),
\((\mathcal Mv)(x,i)=\sum_yw_y\int_{Lx+y}^{Lx+y+Le_i}v\) (a straight
line integral), and \((\mathcal Fv)(x)=\tau_x(0)-\sum_yw_y\tau_x(y)\). -
(d) The linearization \(\mathcal Q=\mathcal M+d_0\mathcal F\) is
independent of \(t\) and maps fine pure gauge to coarse pure gauge:
\(\mathcal Q\,d_0\varphi=d_0(\varphi\circ\mathrm{embed})\).

\emph{Status:} theorem, with the neighbourhood allowed to depend on the
fixed step. Since both tag laws are locally uniform, the kernel itself
is analytic in fine and coarse exponential coordinates near \((1,1)\)
for each fixed \(\beta\). \emph{Lean:}
\href{https://github.com/yang-mills-group/OpenAI-YM-20261028/blob/3456e5e75cf8d4cd558955a5c7638fba2d68f464/YM/PREC/Derivative.lean\#L113-L116}{\texttt{Push.settled}},
\href{https://github.com/yang-mills-group/OpenAI-YM-20261028/blob/3456e5e75cf8d4cd558955a5c7638fba2d68f464/YM/PREC/Derivative.lean\#L118-L124}{\texttt{Push.prediction\_cont}},
\href{https://github.com/yang-mills-group/OpenAI-YM-20261028/blob/3456e5e75cf8d4cd558955a5c7638fba2d68f464/YM/PREC/Derivative.lean\#L38-L68}{\texttt{Push.derivative}},
\href{https://github.com/yang-mills-group/OpenAI-YM-20261028/blob/3456e5e75cf8d4cd558955a5c7638fba2d68f464/YM/PREC/Derivative.lean\#L147-L160}{\texttt{Push.prices}},
\href{https://github.com/yang-mills-group/OpenAI-YM-20261028/blob/3456e5e75cf8d4cd558955a5c7638fba2d68f464/YM/PREC/FrameResponse.lean\#L68}{\texttt{Precision.probe}},
\href{https://github.com/yang-mills-group/OpenAI-YM-20261028/blob/3456e5e75cf8d4cd558955a5c7638fba2d68f464/YM/PREC/FrameResponse.lean\#L92-L96}{\texttt{Precision.root\_gradient}}.

\subsubsection{2. What it uses}\label{what-it-uses-3}

\textbf{From other units.}

\begin{itemize}
\tightlist
\item
  I.1 (lattice, gauge groups). This unit uses:

  \begin{itemize}
  \tightlist
  \item
    the torus and its cover \(\mathbb Z^4\), the block subdivision
    (root, parent), coordinate routes, slabs and layers, parallel
    transport and its gauge rule, the signed symmetries, and the lattice
    forms with \(d\circ d=0\);
  \item
    the gauge group data: chord, price, conjugation invariance,
    probability Haar measure and its two-sided invariance, the
    exponential map;
  \item
    the nearest-point map \(\pi\) of \(G\subset M_N(\mathbb C)\):
    uniqueness of the nearest point within a tubular radius \(r_G\)
    (\href{https://github.com/yang-mills-group/OpenAI-YM-20261028/blob/3456e5e75cf8d4cd558955a5c7638fba2d68f464/YM/Gauge/Nearest.lean\#L141-L153}{\texttt{nearest\_tube}});
    analyticity of \(\pi\) near \(1\) with \(D\pi(1)=\Pi_{\mathfrak g}\)
    (\href{https://github.com/yang-mills-group/OpenAI-YM-20261028/blob/3456e5e75cf8d4cd558955a5c7638fba2d68f464/YM/Gauge/Nearest.lean\#L155-L170}{\texttt{focus\_deriv}});
    the strict-differentiability bound
    \(\|\pi(X)-\pi(Y)-\Pi_{\mathfrak g}(X-Y)\|\le\eta\|X-Y\|\) for
    \(X,Y\) near \(1\)
    (\href{https://github.com/yang-mills-group/OpenAI-YM-20261028/blob/3456e5e75cf8d4cd558955a5c7638fba2d68f464/YM/Gauge/Nearest.lean\#L172-L181}{\texttt{nearest\_approximate}});
    and the first-order condition that \(1\) nearest to \(X\) implies
    \(X-1\perp\mathfrak g\)
    (\href{https://github.com/yang-mills-group/OpenAI-YM-20261028/blob/3456e5e75cf8d4cd558955a5c7638fba2d68f464/YM/Gauge/Nearest.lean\#L43-L63}{\texttt{nearest\_normal\_one}});
  \item
    the Wilson energy, the prior and the partition function; and \(C_A\)
    (\href{https://github.com/yang-mills-group/OpenAI-YM-20261028/blob/3456e5e75cf8d4cd558955a5c7638fba2d68f464/YM/Gauge/Colors.lean\#L180-L182}{\texttt{casimir}}).
  \end{itemize}
\item
  III.16 (analytic tools): the \(\ell^1\) size of lattice offsets used
  for forest heights
  (\href{https://github.com/yang-mills-group/OpenAI-YM-20261028/blob/3456e5e75cf8d4cd558955a5c7638fba2d68f464/YM/Laurent/Eval.lean\#L34}{\texttt{Laurent.size}}).
\item
  III.2 (constants and schedules): the staged choices of kernel and
  analytic constants, block sizes, entry thresholds, barrier and
  windows, with their compatibility conditions
  (\href{https://github.com/yang-mills-group/OpenAI-YM-20261028/blob/3456e5e75cf8d4cd558955a5c7638fba2d68f464/YM/Price/Intent.lean\#L145-L162}{\texttt{Paid}}).
  They enter the statements of III.1.11--III.1.12.
\item
  III.10: the proof of the energy-per-seed requirement (see III.1.11).
  III.2 supplies directed histories (see III.1.12).
\end{itemize}

\textbf{Standard mathematics.}

\begin{itemize}
\tightlist
\item
  Block-spin renormalization for lattice gauge fields, in the
  background-field and axial-gauge setting: T. Bałaban,
  \emph{Propagators and renormalization transformations for lattice
  gauge theories I, II}, Comm. Math. Phys. 95 (1984) 17--40, 96 (1984)
  223--250; \emph{Averaging operations for lattice gauge theories}, CMP
  98 (1985) 17--51; \emph{Spaces of regular gauge field configurations
  on a lattice and gauge fixing conditions}, CMP 99 (1985) 75--102;
  \emph{Renormalization group approach to lattice gauge field theories
  I, II}, CMP 109 (1987) 249--301, 116 (1988) 1--22.
\item
  Axial and tree gauges, and the absence of a Faddeev--Popov factor for
  compact groups: E. Seiler, \emph{Gauge theories as a problem of
  constructive quantum field theory and statistical mechanics}, Lecture
  Notes in Physics 159 (1982). Curvature controls the connection in a
  suitable gauge (continuum analogue): K. Uhlenbeck, CMP 83 (1982)
  31--42.
\item
  Least-squares group synchronization: A. Singer, \emph{Angular
  synchronization by eigenvectors and semidefinite programming}, Appl.
  Comput. Harmon. Anal. 30 (2011) 20--36. Extrinsic means on embedded
  manifolds: R. Bhattacharya, V. Patrangenaru, Ann. Statist. 31 (2003)
  1--29.
\item
  The tubular neighbourhood theorem and the nearest-point projection;
  the real-analytic implicit function theorem (Krantz--Parks, \emph{The
  Implicit Function Theorem}, Birkhäuser 2002), as in
  \href{https://github.com/yang-mills-group/OpenAI-YM-20261028/blob/3456e5e75cf8d4cd558955a5c7638fba2d68f464/YM/PREC/Implicit.lean\#L27-L50}{\texttt{Push.authored}}.
\item
  Continuity of a unique minimizer depending on a parameter, when the
  argmin relation is closed and the target compact
  (\href{https://github.com/yang-mills-group/OpenAI-YM-20261028/blob/3456e5e75cf8d4cd558955a5c7638fba2d68f464/YM/RG/Measurable.lean\#L20-L36}{\texttt{minArg\_cont}}).
  Fubini's theorem, Jensen's inequality, and invariance of Haar measure.
\end{itemize}

\textbf{Imported constants and assumptions.}

\begin{itemize}
\tightlist
\item
  The constants \(d(G)\), \(r_G\) and \(A(G)\), and the clustering
  constant of Lemma 4.14, depend only on \(G\), through \(\pi\). The
  kernel and geometric theorems require no simplicity assumption. The
  construction of directed histories cited in III.1.12 assumes that the
  Lie algebra of \(G\) is simple and nonabelian.
\item
  The kernel constants are chosen before the block factors
  (``pre-\(L\)''). \(L\) odd \(\ge3\) and \(J\ge4\) are part of the
  patch; in a calendar, \(J_m\) and \(L_m\) come from III.2.
\item
  \(\kappa_{\rm slope}=11C_A/(24\pi^2)\) enters only through
  \(\hat\beta_m\) and the identity
  \(\hat\beta_{m+1}=\hat\beta_m+\kappa_{\rm slope}\log L_m\).
\item
  The energy-per-seed requirement (III.1.11) and the kinetic data
  (III.1.12) are hypotheses here, discharged in III.10 and III.2
  respectively.
\end{itemize}

\subsubsection{3. Strategy}\label{strategy-3}

The step follows the block-spin method: fix a frame around each coarse
site, transport across coarse links, average the connectors, and add
independent Haar-normalized Wilson noise. Two rules define the frames on
all configurations. When route holonomies in axial gauges are small,
least-squares synchronization gives a unique frame normalized at the
centre. Otherwise an axial frame is chosen uniformly among tags
satisfying the curvature bound. The predictor likewise uses a unique
nearest point to the weighted matrix mean when the minimum price is
small, and a connector sampled with weight \(w\) otherwise. Uniqueness
and symmetric sampling make these rules covariant. The tag probabilities
sum to \(1\), and the noise has total mass \(1\), so the full kernel is
normalized.

The branch cutoffs are independent of the number of patch sites. For
frame uniqueness, normalize each minimizer by a nearest point to its
centroid. Centroid differences then have zero tangential projection,
making the feedback comparison a contraction. Predictor uniqueness
follows from the tubular neighbourhood of \(G\). Eligible axial frames
exist by induction over coordinate slabs: choose a layer of least
curvature, fix its gauge, and compare the other layers by plaquette
bounds.

For quantitative control, one axial gauge on the macro-box synchronizes
the surrounding patch frames. Each frame defect costs \(O(L^2)\) times
the curvature sum; the averaging weight contributes \(L^{-4}\). Fine
transports cost another \(L\) times a link sum of order \(L\), giving
the same \(L^{-2}\) coarse-curvature estimate. Near the trivial
configuration, the analytic implicit function theorem gives the frame
and predictor derivatives. Tree gauge then reduces integration to free
links. These geometric and analytic results enter the later energy
estimates and stationary-branch analysis.

\subsubsection{4. Steps}\label{steps-3}

\paragraph{4.1 The kernel and its
normalization}\label{the-kernel-and-its-normalization}

\textbf{Lemma 4.1 (the tag law is a probability).} Under (B3),
\(\sum_t\lambda(U,t)=1\). \emph{Sketch:} first sum over the bond tags
with the frame tags fixed. The connectors are then fixed, and each bond
factor sums to \(1\): the mass is uniform on the minimizing branch, and
it is \(w\), with \(\sum w=1\), otherwise. The remaining sum factorizes
over coarse sites, and each frame factor sums to \(1\): the mass is
uniform on all tags, or on the nonempty eligible set. \emph{Lean:}
\href{https://github.com/yang-mills-group/OpenAI-YM-20261028/blob/3456e5e75cf8d4cd558955a5c7638fba2d68f464/YM/RG/Block.lean\#L123-L136}{\texttt{StepData.chosen}},
\href{https://github.com/yang-mills-group/OpenAI-YM-20261028/blob/3456e5e75cf8d4cd558955a5c7638fba2d68f464/YM/RG/Measurable.lean\#L127-L136}{\texttt{frameMass\_sum}},
\href{https://github.com/yang-mills-group/OpenAI-YM-20261028/blob/3456e5e75cf8d4cd558955a5c7638fba2d68f464/YM/RG/Measurable.lean\#L138-L144}{\texttt{predictorMass\_sum}}.

\textbf{Lemma 4.2 (Borel measurability).} Under (B1)--(B2), for each tag
the frames, connectors, predictions and \(\lambda(U,t)\) are Borel
functions of \(U\). \emph{Sketch:} \(M_\delta\) and
\(\{\min\Phi_Y<\varepsilon\}\) are open. The argmin relations are closed
and the targets compact, so a unique minimizer depends continuously on
the parameter. Hence the selected frame is continuous on \(M_\delta\)
and the minimizing predictor is continuous where
\(\min\Phi_Y<\varepsilon\). For a fixed tag, fallback frames are
continuous in \(U\). Eligibility is a closed condition. Piecewise
definitions over Borel sets are Borel. \emph{Lean:}
\href{https://github.com/yang-mills-group/OpenAI-YM-20261028/blob/3456e5e75cf8d4cd558955a5c7638fba2d68f464/YM/RG/Measurable.lean\#L20-L36}{\texttt{minArg\_cont}},
\href{https://github.com/yang-mills-group/OpenAI-YM-20261028/blob/3456e5e75cf8d4cd558955a5c7638fba2d68f464/YM/RG/Measurable.lean\#L50-L53}{\texttt{selected\_cont}},
\href{https://github.com/yang-mills-group/OpenAI-YM-20261028/blob/3456e5e75cf8d4cd558955a5c7638fba2d68f464/YM/RG/Measurable.lean\#L55-L61}{\texttt{best\_cont}},
\href{https://github.com/yang-mills-group/OpenAI-YM-20261028/blob/3456e5e75cf8d4cd558955a5c7638fba2d68f464/YM/RG/Measurable.lean\#L63-L67}{\texttt{min\_open}}.

\textbf{Lemma 4.3 (Proposition III.1.1).} \emph{Sketch:} for each link,
right invariance of Haar measure moves \(Q_b\) to \(1\), so
\(\int Z_\beta^{-1}e^{-\beta\kappa\,\mathrm{pr}(V_bQ_b^{-1})}dV_b=1\).
Fubini over the coarse links then gives
\(\int\rho_\beta\,dV=\sum_t\lambda=1\). The noise is continuous on a
compact space, hence bounded by some \(A_\beta\), and
\(\rho_\beta\le A_\beta\sum_t\lambda=A_\beta\). The properties of the
push follow from Fubini for the bounded kernel times an integrable
\(F\). \emph{Lean:}
\href{https://github.com/yang-mills-group/OpenAI-YM-20261028/blob/3456e5e75cf8d4cd558955a5c7638fba2d68f464/YM/RG/Block.lean\#L146-L155}{\texttt{noise\_mass}},
\href{https://github.com/yang-mills-group/OpenAI-YM-20261028/blob/3456e5e75cf8d4cd558955a5c7638fba2d68f464/YM/RG/Block.lean\#L197-L201}{\texttt{density\_total}},
\href{https://github.com/yang-mills-group/OpenAI-YM-20261028/blob/3456e5e75cf8d4cd558955a5c7638fba2d68f464/YM/RG/Block.lean\#L184-L190}{\texttt{density\_bound}},
\href{https://github.com/yang-mills-group/OpenAI-YM-20261028/blob/3456e5e75cf8d4cd558955a5c7638fba2d68f464/YM/RG/Push.lean\#L17-L24}{\texttt{push\_slice}},
\href{https://github.com/yang-mills-group/OpenAI-YM-20261028/blob/3456e5e75cf8d4cd558955a5c7638fba2d68f464/YM/RG/Push.lean\#L37-L43}{\texttt{push\_integral}}.

\paragraph{4.2 Covariance}\label{covariance}

\textbf{Lemma 4.4 (gauge covariance of the patch data).} Assume
(B1)--(B2). For every \(k:\mathbb Z^4\to G\), \(M_\delta\), the eligible
sets and frame masses are unchanged under \(W\mapsto k\cdot W\), and the
frame becomes \(h_y\mapsto k_0h_yk_y^{-1}\). For a gauge
\(k:\Lambda_f\to G\), the connector of \((x,i)\) becomes
\(k_{Lx}\,Y_y\,k_{L(x+e_i)}^{-1}\). The prediction transforms in the
same way, with unchanged bond mass. \emph{Sketch:}
\(T_{k\cdot W}(y\to z;o)=k_yT_W(y\to z;o)k_z^{-1}\), and ch, pr are
class functions. Then \(S_{k\cdot W}(k_0hk^{-1})=S_W(h)\), the
normalization \(h_0=1\) is preserved, and (B1) identifies the minimizer.
The axial costs are conjugated, hence unchanged.
\(\Phi_{uYv}(uqv)=\Phi_Y(q)\), and (B2) identifies the minimizer.
\emph{Lean:}
\href{https://github.com/yang-mills-group/OpenAI-YM-20261028/blob/3456e5e75cf8d4cd558955a5c7638fba2d68f464/YM/Branch/Transport.lean\#L111-L133}{\texttt{frame\_gauged}},
\href{https://github.com/yang-mills-group/OpenAI-YM-20261028/blob/3456e5e75cf8d4cd558955a5c7638fba2d68f464/YM/Branch/Transport.lean\#L106-L109}{\texttt{min\_gauged}},
\href{https://github.com/yang-mills-group/OpenAI-YM-20261028/blob/3456e5e75cf8d4cd558955a5c7638fba2d68f464/YM/Branch/Predict.lean\#L150-L157}{\texttt{updateGauged}},
\href{https://github.com/yang-mills-group/OpenAI-YM-20261028/blob/3456e5e75cf8d4cd558955a5c7638fba2d68f464/YM/Branch/Predict.lean\#L54-L62}{\texttt{bond\_gauge}},
\href{https://github.com/yang-mills-group/OpenAI-YM-20261028/blob/3456e5e75cf8d4cd558955a5c7638fba2d68f464/YM/Branch/Predict.lean\#L23-L35}{\texttt{predicts}}.

\textbf{Lemma 4.5 (lattice-symmetry covariance of the patch data).}
Assume (B1)--(B2). Let \(s\in\mathfrak B_4\) act about the centre of a
patch. Then:

\begin{itemize}
\tightlist
\item
  \(M_\delta\) is invariant; this is why all 24 orderings enter its
  definition;
\item
  frame tags are permuted by \((r,o)\mapsto(s r,\,\bar s\circ o)\),
  where \(\bar s\in S_4\) is the underlying axis permutation, and
  eligibility and masses correspond;
\item
  the frame at tag \(t\) for \(s_0W\) equals the frame at the
  transformed tag for \(W\), composed with \(s\);
\item
  \(w\) is invariant;
\item
  after reindexing site tags, the connector for the transformed field
  equals the original connector at the image link, composed with \(s\),
  and inverted when \(s\) reverses that link;
\item
  the prediction commutes with both operations.
\end{itemize}

\emph{Sketch:} routes and axial gauges are mapped to routes and axial
gauges, and the score is invariant; (B1) gives the minimizer. Moreover
\(\mathrm{pr}(k^{-1}u)=\mathrm{pr}(ku^{-1})\), so
\(\Phi_{Y^{-1}}(u^{-1})=\Phi_Y(u)\), and (B2) applies. \emph{Lean:}
\href{https://github.com/yang-mills-group/OpenAI-YM-20261028/blob/3456e5e75cf8d4cd558955a5c7638fba2d68f464/YM/Branch/Transport.lean\#L77-L94}{\texttt{frames\_channel}},
\href{https://github.com/yang-mills-group/OpenAI-YM-20261028/blob/3456e5e75cf8d4cd558955a5c7638fba2d68f464/YM/Branch/Transport.lean\#L30-L32}{\texttt{axial\_channel}},
\href{https://github.com/yang-mills-group/OpenAI-YM-20261028/blob/3456e5e75cf8d4cd558955a5c7638fba2d68f464/YM/Branch/Transport.lean\#L34-L38}{\texttt{min\_channel}},
\href{https://github.com/yang-mills-group/OpenAI-YM-20261028/blob/3456e5e75cf8d4cd558955a5c7638fba2d68f464/YM/Branch/Transport.lean\#L73-L75}{\texttt{good\_channel}},
\href{https://github.com/yang-mills-group/OpenAI-YM-20261028/blob/3456e5e75cf8d4cd558955a5c7638fba2d68f464/YM/Branch/Transport.lean\#L135-L141}{\texttt{weight\_invariant}},
\href{https://github.com/yang-mills-group/OpenAI-YM-20261028/blob/3456e5e75cf8d4cd558955a5c7638fba2d68f464/YM/Branch/Predict.lean\#L37-L42}{\texttt{bond\_pull}},
\href{https://github.com/yang-mills-group/OpenAI-YM-20261028/blob/3456e5e75cf8d4cd558955a5c7638fba2d68f464/YM/Branch/Predict.lean\#L44-L52}{\texttt{bond\_inv}},
\href{https://github.com/yang-mills-group/OpenAI-YM-20261028/blob/3456e5e75cf8d4cd558955a5c7638fba2d68f464/YM/Branch/Predict.lean\#L127-L134}{\texttt{StepData.positive}},
\href{https://github.com/yang-mills-group/OpenAI-YM-20261028/blob/3456e5e75cf8d4cd558955a5c7638fba2d68f464/YM/Branch/Predict.lean\#L136-L148}{\texttt{StepData.negative}}.
\emph{(proof read, including
\href{https://github.com/yang-mills-group/OpenAI-YM-20261028/blob/3456e5e75cf8d4cd558955a5c7638fba2d68f464/YM/Branch/Channel.lean\#L117-L127}{\texttt{Axial.channels}})}

\textbf{Lemma 4.6 (Proposition III.1.2).} \emph{Sketch:} (b) The
bijection of tag configurations that relocates frame tags to image
sites, maps them by \(s\), and maps bond tags by \(s\) gives
\(\lambda(s_{Lb}U,t)=\lambda(U,t')\) and \(Q(s_{Lb}U,t)=s_bQ(U,t')\).
The noise is invariant under the simultaneous symmetry, including
inversion of both arguments. Reindex the sum over tags. (a) \(\lambda\)
is invariant and \(Q\mapsto(k\circ\mathrm{embed})\cdot Q\); the noise is
invariant because
\(\mathrm{pr}(kVk'^{-1}(kQk'^{-1})^{-1})=\mathrm{pr}(VQ^{-1})\). (c) is
immediate. \emph{Lean:}
\href{https://github.com/yang-mills-group/OpenAI-YM-20261028/blob/3456e5e75cf8d4cd558955a5c7638fba2d68f464/YM/Branch/Covariant.lean\#L22-L34}{\texttt{mass\_channel}},
\href{https://github.com/yang-mills-group/OpenAI-YM-20261028/blob/3456e5e75cf8d4cd558955a5c7638fba2d68f464/YM/Branch/Covariant.lean\#L13-L20}{\texttt{prediction\_channel}},
\href{https://github.com/yang-mills-group/OpenAI-YM-20261028/blob/3456e5e75cf8d4cd558955a5c7638fba2d68f464/YM/Branch/Covariant.lean\#L36-L48}{\texttt{noise\_channel}},
\href{https://github.com/yang-mills-group/OpenAI-YM-20261028/blob/3456e5e75cf8d4cd558955a5c7638fba2d68f464/YM/Branch/Predict.lean\#L88}{\texttt{withinChoice}},
\href{https://github.com/yang-mills-group/OpenAI-YM-20261028/blob/3456e5e75cf8d4cd558955a5c7638fba2d68f464/YM/Branch/Covariant.lean\#L50-L54}{\texttt{channel}},
\href{https://github.com/yang-mills-group/OpenAI-YM-20261028/blob/3456e5e75cf8d4cd558955a5c7638fba2d68f464/YM/Branch/Covariant.lean\#L56-L67}{\texttt{covary}}.

\paragraph{4.3 Axial gauges and the branch
constants}\label{axial-gauges-and-the-branch-constants}

\textbf{Lemma 4.7 (comparison of parallel layers).} Fix \(r,o,l\) as in
III.1.4 and an axis \(j\notin l\). Let \(S'\) be the slab for \(l+[j]\),
and suppose all \(j\)-links inside \(S'\) equal \(1\). Let \(S_z\) be
the layer of \(S'\) with \(j\)-th coordinate \(z\). Then
\(|E_{S_z}-E_{S_u}|\le4F_{S'}\) for every \(z,u\in\{-r,\ldots,r\}\).
\emph{Sketch:} with trivial \(j\)-links, the \((i,j)\)-plaquette at
\(y\) has holonomy \(U_{(y,i)}U_{(y+e_j,i)}^{-1}\). The inequality
\(|\mathrm{ch}(a)-\mathrm{ch}(b)|\le\mathrm{ch}(ab^{-1})\) is a discrete
Stokes inequality. Summing it over the four axes \(i\) and over \(y\) in
a layer bounds \(|E_{S_{z+1}}-E_{S_z}|\) by \(4\times\) the plaquette
sum at the sites of \(S_z\). Telescoping uses each layer once.
\emph{Lean:}
\href{https://github.com/yang-mills-group/OpenAI-YM-20261028/blob/3456e5e75cf8d4cd558955a5c7638fba2d68f464/YM/Branch/Slide.lean\#L72-L85}{\texttt{variation}},
\href{https://github.com/yang-mills-group/OpenAI-YM-20261028/blob/3456e5e75cf8d4cd558955a5c7638fba2d68f464/YM/Branch/Slide.lean\#L60-L70}{\texttt{squares}},
\href{https://github.com/yang-mills-group/OpenAI-YM-20261028/blob/3456e5e75cf8d4cd558955a5c7638fba2d68f464/YM/Branch/Slide.lean\#L87-L118}{\texttt{slide\_le}},
\href{https://github.com/yang-mills-group/OpenAI-YM-20261028/blob/3456e5e75cf8d4cd558955a5c7638fba2d68f464/YM/Branch/Slide.lean\#L120-L152}{\texttt{window\_le}}.

\textbf{Lemma 4.8 (Proposition III.1.4).} \emph{Sketch:} induction on
\(|l|\); the case \(l=\emptyset\) is trivial. For \(l'=l+[j]\), the slab
\(S'\) is the disjoint union of \(2r+1\) layers \(S_z\), each a slab for
\(l\). Choose the layer \(z_0\) with the least curvature, so that
\((2r+1)F_{S_{z_0}}\le F_{S'}\). By induction there is a root
\(x\in S_{z_0}\) with
\(E_{S_{z_0}}\le4|l|(2r+1)F_{S_{z_0}}\le4|l|F_{S'}\) in its axial gauge
along \(l\). Extend that gauge by moving along \(j\) last. Then all
\(j\)-links in \(S'\) are trivial, and the gauge on \(S_{z_0}\) is
unchanged. By Lemma 4.7,
\[E_{S'}=\sum_zE_{S_z}\le(2r+1)\big(4|l|+4\big)F_{S'}=4|l'|(2r+1)F_{S'}.\]
The choice of the cheapest layer is what closes the induction.
\emph{Lean:}
\href{https://github.com/yang-mills-group/OpenAI-YM-20261028/blob/3456e5e75cf8d4cd558955a5c7638fba2d68f464/YM/Branch/Pick.lean\#L80-L105}{\texttt{existence}},
\href{https://github.com/yang-mills-group/OpenAI-YM-20261028/blob/3456e5e75cf8d4cd558955a5c7638fba2d68f464/YM/Branch/Pick.lean\#L43-L78}{\texttt{total\_slab}},
\href{https://github.com/yang-mills-group/OpenAI-YM-20261028/blob/3456e5e75cf8d4cd558955a5c7638fba2d68f464/YM/Branch/Pick.lean\#L22-L41}{\texttt{vertical}},
\href{https://github.com/yang-mills-group/OpenAI-YM-20261028/blob/3456e5e75cf8d4cd558955a5c7638fba2d68f464/YM/Branch/Slide.lean\#L49-L58}{\texttt{faces\_pins}}.

\textbf{Lemma 4.9 ((B3) for \(a\ge16\)).} For every \(W\) and every
ordering \(o\) there is a root \(r\) with
\(E_{B_R}(a^{(r,o)}\cdot W)\le16(2R+1)F_{B_R}(W)\). \emph{Sketch:} apply
III.1.4 to the whole box with the four axes in the order \(o\); then
\(4\cdot4=16\). \emph{Lean:}
\href{https://github.com/yang-mills-group/OpenAI-YM-20261028/blob/3456e5e75cf8d4cd558955a5c7638fba2d68f464/YM/Branch/Pick.lean\#L119-L134}{\texttt{Axial.eligible}},
\href{https://github.com/yang-mills-group/OpenAI-YM-20261028/blob/3456e5e75cf8d4cd558955a5c7638fba2d68f464/YM/Branch/Pick.lean\#L107-L113}{\texttt{edge\_cost}}.

\textbf{Lemma 4.10 (structure of the synchronization problem).} Let
\(W\) be as in III.1.5, with centroid \(m(h)=\mathrm{avg}_xh_x\) and
\emph{feedback} \(F_x(h)=\mathrm{avg}_{y,k}h_yW_{xyk}^{-1}\), both
computed in \(M_N(\mathbb C)\).

\begin{itemize}
\tightlist
\item
  \begin{enumerate}
  \def\labelenumi{(\alph{enumi})}
  \tightlist
  \item
    \(u\) minimizes \(\sum_iw_i\,\mathrm{pr}(f_iu^{-1})\) (with
    \(w_i\ge0\), \(\sum w_i=1\)) if and only if \(u\) is a nearest point
    of \(G\) to \(\sum_iw_if_i\).
  \end{enumerate}
\item
  (b) The cost is invariant under \(h\mapsto ch\). For every global
  minimizer, \(h_x\) is a nearest point of \(G\) to \(F_x(h)\) for every
  \(x\); this is the self-consistency condition. Changing only \(h_x\)
  combines the two off-diagonal contributions by reflection symmetry.
  The original diagonal contribution is zero; the added diagonal term
  for the competitor is nonnegative. Thus the comparison also holds with
  the self term included in \(F_x(h)\).
\item
  (c) If \(h\) is a global minimizer and \(\|W-1\|\le d\), then
  \(\mathrm{cost}(h)\le\mathrm{cost}(1)\le d^2/(2N)\), so
  \(\mathrm{avg}_{x,y}\|h_x-h_y\|^2\le4d^2\). If moreover \(1\) is a
  nearest point to \(m(h)\), then \(\|m(h)-1\|\le2d\) and
  \(\|F_x(h)-1\|\le3d\).
\item
  (d) If \(\max_x\|p_x-q_x\|\le u\), then \(\|F_x(p)-F_x(q)\|\le u\) and
  \(\|F_x(p)-F_x(q)-(m(p)-m(q))\|\le ud\).
\end{itemize}

\emph{Lean:}
\href{https://github.com/yang-mills-group/OpenAI-YM-20261028/blob/3456e5e75cf8d4cd558955a5c7638fba2d68f464/YM/Branch/Center.lean\#L72-L93}{\texttt{min\_iff}},
\href{https://github.com/yang-mills-group/OpenAI-YM-20261028/blob/3456e5e75cf8d4cd558955a5c7638fba2d68f464/YM/Branch/Center.lean\#L63-L70}{\texttt{price\_diff}},
\href{https://github.com/yang-mills-group/OpenAI-YM-20261028/blob/3456e5e75cf8d4cd558955a5c7638fba2d68f464/YM/Branch/Synchronize.lean\#L64-L98}{\texttt{min\_cov}},
\href{https://github.com/yang-mills-group/OpenAI-YM-20261028/blob/3456e5e75cf8d4cd558955a5c7638fba2d68f464/YM/Branch/Synchronize.lean\#L100-L123}{\texttt{self\_consistent}},
\href{https://github.com/yang-mills-group/OpenAI-YM-20261028/blob/3456e5e75cf8d4cd558955a5c7638fba2d68f464/YM/Branch/Synchronize.lean\#L127-L156}{\texttt{compare\_energy}},
\href{https://github.com/yang-mills-group/OpenAI-YM-20261028/blob/3456e5e75cf8d4cd558955a5c7638fba2d68f464/YM/Branch/Synchronize.lean\#L158-L170}{\texttt{cluster}},
\href{https://github.com/yang-mills-group/OpenAI-YM-20261028/blob/3456e5e75cf8d4cd558955a5c7638fba2d68f464/YM/Branch/Synchronize.lean\#L172-L188}{\texttt{closeness}},
\href{https://github.com/yang-mills-group/OpenAI-YM-20261028/blob/3456e5e75cf8d4cd558955a5c7638fba2d68f464/YM/Branch/Synchronize.lean\#L190-L201}{\texttt{changes}},
\href{https://github.com/yang-mills-group/OpenAI-YM-20261028/blob/3456e5e75cf8d4cd558955a5c7638fba2d68f464/YM/Branch/Synchronize.lean\#L203-L226}{\texttt{weak\_changes}}.

\textbf{Lemma 4.11 (Proposition III.1.5).} \emph{Sketch:} let
\(\delta_1\) be the radius for which the strict-differentiability bound
of \(\pi\) at \(1\) holds with \(\eta=1/4\), and take
\(d=\min(\delta_1/4,1/4)\), so that \(3d<\delta_1\). Left-multiply each
minimizer by the inverse of a nearest point to its centroid. Then \(1\)
is nearest to both centroids, and by 4.10(b) both are still minimizers.
Let \(p,q\) be the normalized minimizers, \(u=\max_x\|p_x-q_x\|\)
attained at \(x\), and \(B=F_x(p)-F_x(q)\). By 4.10(c) both feedbacks
lie within \(3d<\delta_1\) of \(1\), and self-consistency makes
\(p_x,q_x\) nearest points to these feedbacks. Hence
\(\|p_x-q_x-\Pi_{\mathfrak g}B\|\le\frac14\|B\|\le\frac u4\). Since
\(m(p)-1\) and \(m(q)-1\) are normal to \(\mathfrak g\),
\(\Pi_{\mathfrak g}(m(p)-m(q))=0\), and 4.10(d) gives
\(\|\Pi_{\mathfrak g}B\|\le ud\). So \(u\le ud+u/4\le u/2\), hence
\(u=0\) and \(p=q\). Undoing the normalization gives \(h'=ch\).
\emph{Lean:}
\href{https://github.com/yang-mills-group/OpenAI-YM-20261028/blob/3456e5e75cf8d4cd558955a5c7638fba2d68f464/YM/Branch/Synchronize.lean\#L232-L272}{\texttt{align}}.

\textbf{Lemma 4.12 ((B1) and (B2) on patches).} (B1) holds if
\(\delta\le d/\sqrt N\), and (B2) holds if
\(\varepsilon\le r_G^2/(2N)\), for every patch. \emph{Sketch:} (B1).
Take \(n=B_R\), \(\Lambda=S_4\) and \(W_{yzo}=a_yT(y\to z;o)a_z^{-1}\),
where \(a\) is the axial frame from the centre with the identity
ordering. The involution \(\sigma\) reverses the ordering, since the
reversed route is the route with the reversed ordering. On \(M_\delta\),
\(\|W-1\|=\sqrt N\,\mathrm{ch}(W)<d\). Since
\(S_W(h)=\mathrm{cost}(ha^{-1})\) and the cost is left invariant, a
constrained minimizer is a global minimizer. By III.1.5 two of them
differ by a constant, which equals \(1\) at the centre. (B2). If
\(\Phi_Y(u)<r_G^2/(2N)\), Jensen's inequality gives
\(\|\sum w_yY_y-u\|^2\le\sum w_y\|Y_y-u\|^2=2N\Phi_Y(u)<r_G^2\). So the
mean has a unique nearest point, which by 4.10(a) is the unique
minimizer. \emph{Lean:}
\href{https://github.com/yang-mills-group/OpenAI-YM-20261028/blob/3456e5e75cf8d4cd558955a5c7638fba2d68f464/YM/Branch/Uniform.lean\#L18-L56}{\texttt{frame\_unique}},
\href{https://github.com/yang-mills-group/OpenAI-YM-20261028/blob/3456e5e75cf8d4cd558955a5c7638fba2d68f464/YM/Branch/Uniform.lean\#L58-L63}{\texttt{predictor\_unique}},
\href{https://github.com/yang-mills-group/OpenAI-YM-20261028/blob/3456e5e75cf8d4cd558955a5c7638fba2d68f464/YM/Branch/Center.lean\#L95-L108}{\texttt{weighted\_unique}}.

\textbf{Lemma 4.13 (Proposition III.1.3).} Combine 4.9 and 4.12.
\emph{Lean:}
\href{https://github.com/yang-mills-group/OpenAI-YM-20261028/blob/3456e5e75cf8d4cd558955a5c7638fba2d68f464/YM/Branch/Pick.lean\#L137-L142}{\texttt{viable}}.

\paragraph{\texorpdfstring{4.4 Synchronization and the \(L^{-2}\)
estimate}{4.4 Synchronization and the L\^{}\{-2\} estimate}}\label{synchronization-and-the-l-2-estimate}

\textbf{Lemma 4.14 (clustering of minimizers for arbitrary data).} There
is \(C_G>0\) such that for all finite nonempty \(n,\Lambda\) and all
\(W\) satisfying the diagonal and reflection conditions of III.1.5,
every global minimizer \(h\) admits \(y\in G\) with
\(\mathrm{avg}_x\mathrm{ch}(h_xy^{-1})\le C_G\,\mathrm{avg}_{x,z,k}\mathrm{ch}(W_{xzk})\).
\emph{Sketch:} write
\(\mathrm{inc}_x=\mathrm{avg}_{z,k}\mathrm{ch}(W_{xzk})\) and
\(\bar\mu=\mathrm{avg}_x\mathrm{inc}_x\). Since
\(\mathrm{ch}^2\le2\,\mathrm{ch}\),
\(\mathrm{cost}(h)\le\mathrm{cost}(1)\le\bar\mu\), so
\(\mathrm{avg}\|h_x-h_z\|^2\le8N\bar\mu\). Let \(\delta_1\) be the
radius for which the strict-differentiability bound of \(\pi\) at \(1\)
holds with \(\eta=1\).

\begin{itemize}
\tightlist
\item
  If \(\bar\mu<\sigma=\delta_1^2/(32N)\): normalize so that \(1\) is
  nearest to the centroid, which then lies within \(\delta_1/2\) of
  \(1\). Also \(\|F_x-m\|\le\sqrt N\,\mathrm{inc}_x\). For \(x\) with
  \(\mathrm{inc}_x<\delta_1/(2\sqrt N)\), the derivative bound gives
  \(\|h_x-1\|\le2\|F_x-m\|\), that is
  \(\mathrm{ch}(h_x)\le2\,\mathrm{inc}_x\). For the other \(x\), use
  \(\mathrm{ch}\le2\).
\item
  If \(\bar\mu\ge\sigma\): use \(\mathrm{ch}\le2\le(2/\sigma)\bar\mu\).
\end{itemize}

\emph{Lean:}
\href{https://github.com/yang-mills-group/OpenAI-YM-20261028/blob/3456e5e75cf8d4cd558955a5c7638fba2d68f464/YM/Pressure/Single.lean\#L77-L130}{\texttt{Pressure.cluster}},
\href{https://github.com/yang-mills-group/OpenAI-YM-20261028/blob/3456e5e75cf8d4cd558955a5c7638fba2d68f464/YM/Pressure/Single.lean\#L37-L64}{\texttt{price\_square}},
\href{https://github.com/yang-mills-group/OpenAI-YM-20261028/blob/3456e5e75cf8d4cd558955a5c7638fba2d68f464/YM/Pressure/Single.lean\#L66-L75}{\texttt{income\_bounds}}.

\textbf{Lemma 4.15 (route counting).} For every ordering \(o\):
\(\sum_{y,z\in B_r}\mathrm{ch}\,T_U(y\to z;o)\le4\,|B_r|\,(2r+1)\,E_{B_r}(U)\).
\emph{Sketch:} induction over the axes of the route. The last segment,
along the axis \(j\), starts at the point \(\zeta(y,z)\) that has the
already-processed coordinates of \(z\) and the other coordinates of
\(y\). Its chord is at most the sum of the link chords on the whole
\(j\)-line through \(\zeta(y,z)\). The map \((y,z)\mapsto\zeta(y,z)\) is
\(|B_r|\)-to-one onto \(B_r\), and each \(j\)-link lies on one line
containing \(2r+1\) points. \emph{Lean:}
\href{https://github.com/yang-mills-group/OpenAI-YM-20261028/blob/3456e5e75cf8d4cd558955a5c7638fba2d68f464/YM/Pressure/Route.lean\#L107-L130}{\texttt{route\_sum}},
\href{https://github.com/yang-mills-group/OpenAI-YM-20261028/blob/3456e5e75cf8d4cd558955a5c7638fba2d68f464/YM/Pressure/Route.lean\#L50-L54}{\texttt{marginal}},
\href{https://github.com/yang-mills-group/OpenAI-YM-20261028/blob/3456e5e75cf8d4cd558955a5c7638fba2d68f464/YM/Pressure/Route.lean\#L56-L60}{\texttt{slidings}},
\href{https://github.com/yang-mills-group/OpenAI-YM-20261028/blob/3456e5e75cf8d4cd558955a5c7638fba2d68f464/YM/Pressure/Gauge.lean\#L55-L69}{\texttt{Road.revenue}}.

\textbf{Lemma 4.16 (Proposition III.1.6).} \emph{Sketch:}
\emph{Minimizing branch.} Since \(T_{H\cdot W}=H\,T_WH^{-1}\), the
family \(y\mapsto h_yH_y^{-1}\) is a global minimizer of the
synchronization cost for the data \(T_{H\cdot W}(y\to z;o)\). By 4.14
and 4.15 there is \(c\) with
\(\sum_y\mathrm{ch}(h_yH_y^{-1}c^{-1})\le4C_G(2R+1)E_{B_R}(H\cdot W)\).
The cutoff \(\delta\) plays no role here. \emph{Fallback branch.} Here
\(h\) is the normalized axial frame of an eligible tag, so
\(E_{B_R}(h\cdot W)\le a(2R+1)F_{B_R}(W)\). Put \(g=hH^{-1}\); the
pure-gauge configuration \(g\cdot1\) satisfies
\(E(g\cdot1)\le E(h\cdot W)+E(H\cdot W)\) by the triangle inequality.
Apply 4.15 to \(g\cdot1\), whose route transports are \(g_yg_z^{-1}\),
and choose \(v\) with \(\sum_y\mathrm{ch}(g_yg_v^{-1})\) at most its
average over \(v\). Then \(c=g_v\). In both cases \(A=\max(4C_G,4)\),
and \(\mathrm{ch}(c^{-1}x)=\mathrm{ch}(xc^{-1})\). \emph{Lean:}
\href{https://github.com/yang-mills-group/OpenAI-YM-20261028/blob/3456e5e75cf8d4cd558955a5c7638fba2d68f464/YM/Pressure/Gauge.lean\#L135-L162}{\texttt{synchronized}},
\href{https://github.com/yang-mills-group/OpenAI-YM-20261028/blob/3456e5e75cf8d4cd558955a5c7638fba2d68f464/YM/Pressure/Gauge.lean\#L34-L49}{\texttt{Road.strips}},
\href{https://github.com/yang-mills-group/OpenAI-YM-20261028/blob/3456e5e75cf8d4cd558955a5c7638fba2d68f464/YM/Pressure/Gauge.lean\#L75-L88}{\texttt{risks}},
\href{https://github.com/yang-mills-group/OpenAI-YM-20261028/blob/3456e5e75cf8d4cd558955a5c7638fba2d68f464/YM/Pressure/Gauge.lean\#L95-L132}{\texttt{fallback}}.

\textbf{Lemma 4.17 (elementary bounds on the step geometry).}

\begin{itemize}
\tightlist
\item
  \begin{enumerate}
  \def\labelenumi{(\alph{enumi})}
  \tightlist
  \item
    If \(X+v\subset Y\), then the link and plaquette sums over \(X\) of
    the configuration translated by \(v\) are at most those over \(Y\).
  \end{enumerate}
\item
  (b) \(w_y\,L^4\le1\) for all \(y\). One summand \(\xi_1\) is uniform
  on \(L^4\) points and independent of the others; only \(J\ge1\) is
  used.
\item
  (c)
  \(\sum_{y\in B_R}\mathrm{ch}\,T_U(y+v\to y+v+\ell e_j)\le\ell\,E_X(U)\),
  provided all these segments lie in \(X\).
\item
  (d) For all \(\xi,\eta\in G\) and every bond tag,
  \(\mathrm{ch}(\xi^{-1}Q_b\eta)\le2\max(1,\varepsilon^{-1})\sum_yw_y\,\mathrm{ch}(\xi^{-1}Y_y\eta)\).
  On the minimizing branch \(Q_b\) is nearest to the mean
  \(\bar Y=\sum w_yY_y\), so
  \(\|Q_b-\xi\eta^{-1}\|\le2\|\bar Y-\xi\eta^{-1}\|\). On the fallback
  branch,
  \(\varepsilon\le\Phi_Y(\xi\eta^{-1})\le\sum w_y\mathrm{ch}(\xi^{-1}Y_y\eta)\)
  because \(\mathrm{pr}\le\mathrm{ch}\), while \(\mathrm{ch}\le2\).
\end{itemize}

\emph{Lean:}
\href{https://github.com/yang-mills-group/OpenAI-YM-20261028/blob/3456e5e75cf8d4cd558955a5c7638fba2d68f464/YM/Pressure/Macro.lean\#L16-L43}{\texttt{spread}},
\href{https://github.com/yang-mills-group/OpenAI-YM-20261028/blob/3456e5e75cf8d4cd558955a5c7638fba2d68f464/YM/Pressure/Macro.lean\#L48-L79}{\texttt{light}},
\href{https://github.com/yang-mills-group/OpenAI-YM-20261028/blob/3456e5e75cf8d4cd558955a5c7638fba2d68f464/YM/Pressure/Macro.lean\#L81-L100}{\texttt{traverse}},
\href{https://github.com/yang-mills-group/OpenAI-YM-20261028/blob/3456e5e75cf8d4cd558955a5c7638fba2d68f464/YM/Pressure/Macro.lean\#L102-L129}{\texttt{price\_query}}.

\textbf{Lemma 4.18 (Proposition III.1.7).} \emph{Sketch:} apply III.1.4
once, on \(B_{R'}\) around \(Lx\) with all four axes. This gives one
gauge \(H\) with
\(E_{B_{R'}}(H\cdot U^{(x)})\le16(2R'+1)\,\mathrm{port}\le16\,\ell_jL\,\mathrm{port}\).
For \(y\in B_2\), the patch of \(x+y\) is the translate of \(B_R\) by
\(Ly\) and lies inside \(B_{R'}\). Apply III.1.6 to it, with the
translated gauge \(H(\cdot+Ly)\); this gives \(p_y\). Bound its edge and
face sums by those of \(B_{R'}\) using 4.17(a), and use
\(2R+1\le\ell_jL\). Every tag with \(\lambda>0\) is usable, because a
vanishing frame mass would make \(\lambda\) vanish. \emph{Lean:}
\href{https://github.com/yang-mills-group/OpenAI-YM-20261028/blob/3456e5e75cf8d4cd558955a5c7638fba2d68f464/YM/Pressure/Cost.lean\#L72-L122}{\texttt{planting}},
\href{https://github.com/yang-mills-group/OpenAI-YM-20261028/blob/3456e5e75cf8d4cd558955a5c7638fba2d68f464/YM/Pressure/Cost.lean\#L65-L69}{\texttt{wraps}}.

\textbf{Lemma 4.19 (link comparison).} In the setting of III.1.7, for
\(y,y+e_i\in B_2\) with pads \(\xi=p_y\) and \(\eta=p_{y+e_i}\):
\[L^2\,\mathrm{ch}\big(\xi^{-1}Q_{(x+y,i)}\eta\big)\le2\max(1,\varepsilon^{-1})(2C+D)\,\mathrm{port}(x,U).\]
\emph{Sketch:} write
\(\xi^{-1}Y_k\eta=f_k\cdot T_{H\cdot U^{(x)}}(k+Ly\to k+Ly+Le_i)\cdot g_k^{-1}\),
where \(f_k,g_k\) are the frame defects of III.1.7. Then
\(\mathrm{ch}(\xi^{-1}Y_k\eta)\le\mathrm{ch}f_k+\mathrm{ch}T+\mathrm{ch}g_k\).
Multiply by \(w_kL^4\le1\) and sum: the defects contribute
\(\le2CL^2\,\mathrm{port}\), and by 4.17(c) the transports contribute
\(\le L\,E_{B_{R'}}\le DL^2\,\mathrm{port}\). Conclude with 4.17(d).
\emph{Lean:}
\href{https://github.com/yang-mills-group/OpenAI-YM-20261028/blob/3456e5e75cf8d4cd558955a5c7638fba2d68f464/YM/Pressure/Cost.lean\#L129-L179}{\texttt{indicated}}.

\textbf{Lemma 4.20 (Proposition III.1.8).} \emph{Sketch:} conjugate the
plaquette holonomy of \(Q\) at \(x\) by \(p_0\). It factors into four
padded links \(p_y^{-1}Q\,p_{y'}\), with
\(y,y'\in\{0,e_i,e_k,e_i+e_k\}\subset B_1\), and the intermediate pads
cancel. The triangle inequality and 4.19 give
\(\mathrm{ch}\le4\cdot2\max(1,\varepsilon^{-1})(2C+D)L^{-2}\,\mathrm{port}\).
For \(V\), \(\mathrm{ch}(V_{(x,p)}Q_{(x,p)}^{-1})\le\mathrm{fuel}\) by a
telescoping product estimate. \emph{Lean:}
\href{https://github.com/yang-mills-group/OpenAI-YM-20261028/blob/3456e5e75cf8d4cd558955a5c7638fba2d68f464/YM/Pressure/Cost.lean\#L210-L252}{\texttt{qdc}},
\href{https://github.com/yang-mills-group/OpenAI-YM-20261028/blob/3456e5e75cf8d4cd558955a5c7638fba2d68f464/YM/Pressure/Cost.lean\#L192-L207}{\texttt{comp}}.

\paragraph{4.5 Near the trivial
configuration}\label{near-the-trivial-configuration}

\textbf{Lemma 4.21 (Proposition III.1.14).} \emph{Sketch:}

\begin{enumerate}
\def\labelenumi{\arabic{enumi}.}
\tightlist
\item
  \(U_0\) lies in \(M_\delta\), which is open, and at \(v=0\) the frame
  is \(h\equiv1\) by (B1).
\item
  By 4.10(b) the frame solves \(h_0=1\) and
  \(h_y=\pi\big(\mathrm{avg}_{z,o}h_zT_{U_v}(y\to z;o)^{-1}\big)\) for
  \(y\ne0\). This system is real analytic near \((0,1)\). A variation
  \(u\) in the kernel of its partial derivative satisfies \(u_0=0\) and
  \(u_y=\Pi_{\mathfrak g}(\mathrm{avg}\,u)\) for \(y\ne0\). If
  \(n=|B_R|\) and \(A=\Pi_{\mathfrak g}(\mathrm{avg}\,u)\), then
  \(A=(n-1)A/n\), so \(u=0\). Finite dimensionality makes the derivative
  invertible. Continuity of the selected minimizer identifies it with
  the analytic implicit solution.
\item
  Differentiate the fixed-point equation, using
  \(D\pi(1)=\Pi_{\mathfrak g}\) and the fact that the derivative of a
  transport along \(w\) is the line integral of \(v\) along \(w\). This
  gives \(h'_y=\Pi_{\mathfrak g}(\mathrm{avg}\,h')+\tau(y)\); with
  \(h'_0=0\) this is \(h'_y=\tau(y)-\tau(0)\).
\item
  Near \(0\), \(\min\Phi_Y<\varepsilon\) (it is \(0\) at \(v=0\)), so
  \(Q_b=\pi(\sum w_yY_y)\), which is analytic. Its derivative is
  \(\sum_yw_yY'_y=(\mathcal Mv+d_0\mathcal Fv)(b)\).
\item
  For an exact form \(d_0\varphi\), the line integrals telescope; this
  gives (d).
\end{enumerate}

\emph{Lean:}
\href{https://github.com/yang-mills-group/OpenAI-YM-20261028/blob/3456e5e75cf8d4cd558955a5c7638fba2d68f464/YM/PREC/Pred.lean\#L79-L87}{\texttt{hats\_zero}},
\href{https://github.com/yang-mills-group/OpenAI-YM-20261028/blob/3456e5e75cf8d4cd558955a5c7638fba2d68f464/YM/PREC/Pred.lean\#L116-L127}{\texttt{stars}},
\href{https://github.com/yang-mills-group/OpenAI-YM-20261028/blob/3456e5e75cf8d4cd558955a5c7638fba2d68f464/YM/PREC/Pred.lean\#L129-L187}{\texttt{scheme}},
\href{https://github.com/yang-mills-group/OpenAI-YM-20261028/blob/3456e5e75cf8d4cd558955a5c7638fba2d68f464/YM/PREC/Pred.lean\#L44-L47}{\texttt{focused}},
\href{https://github.com/yang-mills-group/OpenAI-YM-20261028/blob/3456e5e75cf8d4cd558955a5c7638fba2d68f464/YM/PREC/Derivative.lean\#L38-L68}{\texttt{derivative}},
\href{https://github.com/yang-mills-group/OpenAI-YM-20261028/blob/3456e5e75cf8d4cd558955a5c7638fba2d68f464/YM/PREC/Derivative.lean\#L95-L111}{\texttt{lips}},
\href{https://github.com/yang-mills-group/OpenAI-YM-20261028/blob/3456e5e75cf8d4cd558955a5c7638fba2d68f464/YM/PREC/Derivative.lean\#L147-L160}{\texttt{prices}},
\href{https://github.com/yang-mills-group/OpenAI-YM-20261028/blob/3456e5e75cf8d4cd558955a5c7638fba2d68f464/YM/PREC/FrameResponse.lean\#L92-L96}{\texttt{root\_gradient}},
\href{https://github.com/yang-mills-group/OpenAI-YM-20261028/blob/3456e5e75cf8d4cd558955a5c7638fba2d68f464/YM/Exterior/Complex.lean\#L99-L116}{\texttt{Exterior.curvature\_linear}}.

\paragraph{4.6 Tree gauge}\label{tree-gauge}

\textbf{Lemma 4.22 (Proposition III.1.13).} \emph{Sketch:} (a) Because
\(L\) is odd, blocks are centred at their roots and the offset \(q\)
lies in \(B_{(L-1)/2}\). A step towards \(0\) along an axis with
\(q_a\ne0\) decreases \(|q|_1\le4L\) and stays in the block. (b)
\(\mathrm{fold}_{y}\,U_{\mathrm{up}(y)}=\mathrm{fold}_{y+\mathrm{up}(y)}\),
by induction on the height. (c) Split the links into trunks and free
links; this split preserves the measure. For fixed trunk values, gauging
by \(\mathrm{fold}\) (trunks set to \(1\)) sends the trunk values to
\(1\). On the free links it acts as
\(s_e\mapsto k_{e^-}s_ek_{e^+}^{-1}\), which preserves product Haar
measure, and \(F\) is unchanged because \(k=1\) at the roots. (d)
\(U\mapsto\rho_\beta(U,V)F(U)\) is invariant under gauge transformations
that are trivial at the roots, by III.1.2(a). \emph{Lean:}
\href{https://github.com/yang-mills-group/OpenAI-YM-20261028/blob/3456e5e75cf8d4cd558955a5c7638fba2d68f464/YM/Tree/Pines.lean\#L33-L72}{\texttt{select}},
\href{https://github.com/yang-mills-group/OpenAI-YM-20261028/blob/3456e5e75cf8d4cd558955a5c7638fba2d68f464/YM/Tree/Cut.lean\#L32-L37}{\texttt{compress}},
\href{https://github.com/yang-mills-group/OpenAI-YM-20261028/blob/3456e5e75cf8d4cd558955a5c7638fba2d68f464/YM/Tree/Cut.lean\#L39-L45}{\texttt{names}},
\href{https://github.com/yang-mills-group/OpenAI-YM-20261028/blob/3456e5e75cf8d4cd558955a5c7638fba2d68f464/YM/Tree/Cut.lean\#L81-L111}{\texttt{pack}},
\href{https://github.com/yang-mills-group/OpenAI-YM-20261028/blob/3456e5e75cf8d4cd558955a5c7638fba2d68f464/YM/Tree/Cut.lean\#L114-L120}{\texttt{pushing}}.

\paragraph{4.7 Formats and energy
accounting}\label{formats-and-energy-accounting}

\textbf{Lemma 4.23 (the theorems of III.1.9).}

\begin{itemize}
\tightlist
\item
  Gas without flags: an admissible family must cover the empty flag set
  by nonempty inventories, so only the empty family occurs.
\item
  Locality: flags on \(X\) depend only on the plaquettes meeting \(X\).
\item
  Envelope: summing the per-site bounds on the finite torus gives
  summable majorants for the regular and forced activities. The gas is
  dominated by the convergent sum of products of the forced majorants.
  Row and atom series have summable majorants as well; the Borel masks
  preserve these bounds. The exponent and gas are therefore Borel and
  bounded.
\item
  Partition function: decreasing induction on \(m\), using
  \(\mathrm{eval}_m=\mathcal P\,\mathrm{eval}_{m+1}\), the identity
  \(\int\mathcal PF=\int F\) (III.1.1, with the kernel property at each
  level), and the initial condition
  \(\mathrm{eval}_N=e^{-\beta\sum\mathrm{pr}}\).
\end{itemize}

\emph{Lean:}
\href{https://github.com/yang-mills-group/OpenAI-YM-20261028/blob/3456e5e75cf8d4cd558955a5c7638fba2d68f464/YM/RG/Activities.lean\#L135-L147}{\texttt{gas\_one}},
\href{https://github.com/yang-mills-group/OpenAI-YM-20261028/blob/3456e5e75cf8d4cd558955a5c7638fba2d68f464/YM/RG/Masks.lean\#L108-L110}{\texttt{UsesFinal.allReads}},
\href{https://github.com/yang-mills-group/OpenAI-YM-20261028/blob/3456e5e75cf8d4cd558955a5c7638fba2d68f464/YM/RG/History.lean\#L195-L198}{\texttt{Rep.envelope}},
\href{https://github.com/yang-mills-group/OpenAI-YM-20261028/blob/3456e5e75cf8d4cd558955a5c7638fba2d68f464/YM/RG/History.lean\#L242-L249}{\texttt{partition\_eq}}.

\textbf{Lemma 4.24 (Proposition III.1.10).} \emph{Sketch:} expand the
six works. The adjointness
\(\langle\Phi q,p\rangle=\langle q,\Phi^{\dagger}p\rangle\) matches the
cross terms of \(W_F\), \(W_I\) with \(W_L\), and the squares
\(\tfrac12\|\Phi w\|^2\), \(\tfrac12\|\Phi'w\|^2\),
\(\tfrac12\|\mathcal Nw\|^2\), \(\tfrac12\|w\|^2\) cancel against
\(W_D\), \(W_N\), \(W_C\). What remains is the right-hand side. III.1.11
requires a lower bound for the corresponding local sum near the seeds;
the identity alone supplies no such bound. \emph{Lean:}
\href{https://github.com/yang-mills-group/OpenAI-YM-20261028/blob/3456e5e75cf8d4cd558955a5c7638fba2d68f464/YM/Pressure/Rows.lean\#L145-L163}{\texttt{Tube.locates}},
\href{https://github.com/yang-mills-group/OpenAI-YM-20261028/blob/3456e5e75cf8d4cd558955a5c7638fba2d68f464/YM/Pressure/Rows.lean\#L86-L89}{\texttt{arity}}.

\textbf{Lemma 4.25 (transfer of the energy per seed, III.1.11).} Assume
the ringing condition and \(j_{\min}\le i\le j\le i+1\). \emph{Sketch:}
the stamp for the pair \((j,i)\) gives \(\mu\) and a \(\beta\)-entry
threshold. In a directed history:

\begin{itemize}
\tightlist
\item
  the tube of step \(m\) is compatible with the estimator data;
\item
  the side condition \(c\ge100(1+\varrho)\) is part of the fits
  conditions;
\item
  \(\hat\beta_m\) exceeds the entry threshold by the price conditions;
\item
  \(\hat\beta_{m+1}=\hat\beta_m+\kappa_{\rm slope}\log L_m\) turns
  \(\mathrm{temp}(\hat\beta_m+\kappa_{\rm slope}\log L)\) into
  \(t_{m+1}\).
\end{itemize}

\emph{Lean:}
\href{https://github.com/yang-mills-group/OpenAI-YM-20261028/blob/3456e5e75cf8d4cd558955a5c7638fba2d68f464/YM/Pressure/Seeds.lean\#L166-L181}{\texttt{notation\_ticket}}.

\subsubsection{5. Where the difficulty
lies}\label{where-the-difficulty-lies-3}

The synchronization cutoff must be independent of the patch size. The
feedback map alone has Lipschitz constant \(1\); normalizing the
centroids removes its tangential average and gives the contraction
factor \(d+\tfrac14\le\tfrac12\).

The curvature estimate must hold on both branches. It combines
clustering of arbitrary minimizers, route multiplicities and
curvature-controlled axial gauges. Losing a power of \(L\) in any of
these comparisons would spoil the area scaling in III.1.8.

Normalization and tree-gauge integration use finite sums and Haar
invariance. Global regularity is more limited: the kernel is Borel,
while analyticity is proved near the trivial configuration. The energy
bound and stationary requirements in III.1.11--III.1.12 need the later
units.

\subsubsection{6. Notation}\label{notation-3}

{\def\LTcaptype{none} % do not increment counter
\begin{longtable}[]{@{}
  >{\raggedright\arraybackslash}p{(\linewidth - 4\tabcolsep) * \real{0.3333}}
  >{\raggedright\arraybackslash}p{(\linewidth - 4\tabcolsep) * \real{0.3333}}
  >{\raggedright\arraybackslash}p{(\linewidth - 4\tabcolsep) * \real{0.3333}}@{}}
\toprule\noalign{}
\begin{minipage}[b]{\linewidth}\raggedright
Written proof
\end{minipage} & \begin{minipage}[b]{\linewidth}\raggedright
Lean
\end{minipage} & \begin{minipage}[b]{\linewidth}\raggedright
Meaning
\end{minipage} \\
\midrule\noalign{}
\endhead
\bottomrule\noalign{}
\endlastfoot
\(\rho_\beta(U,V)\), \(\mathcal P_\beta\) &
\href{https://github.com/yang-mills-group/OpenAI-YM-20261028/blob/3456e5e75cf8d4cd558955a5c7638fba2d68f464/YM/RG/Block.lean\#L71-L72}{\texttt{StepData.density}},
\href{https://github.com/yang-mills-group/OpenAI-YM-20261028/blob/3456e5e75cf8d4cd558955a5c7638fba2d68f464/YM/RG/Block.lean\#L77-L79}{\texttt{push}}
& block kernel and push \\
\(P=(L,J)\), \(R=(J+3)L\) &
\href{https://github.com/yang-mills-group/OpenAI-YM-20261028/blob/3456e5e75cf8d4cd558955a5c7638fba2d68f464/YM/RG/Patch.lean\#L29-L34}{\texttt{Patch}},
\href{https://github.com/yang-mills-group/OpenAI-YM-20261028/blob/3456e5e75cf8d4cd558955a5c7638fba2d68f464/YM/RG/Patch.lean\#L51}{\texttt{radius}}
& block factor, order, patch radius \\
\(K=(\kappa,\delta,\varepsilon,a)\) &
\href{https://github.com/yang-mills-group/OpenAI-YM-20261028/blob/3456e5e75cf8d4cd558955a5c7638fba2d68f464/YM/RG/Patch.lean\#L37-L45}{\texttt{KernelChoices}}
& penalty, frame cutoff, predictor cutoff, axial factor \\
(B1)--(B3) &
\href{https://github.com/yang-mills-group/OpenAI-YM-20261028/blob/3456e5e75cf8d4cd558955a5c7638fba2d68f464/YM/RG/Patch.lean\#L182-L185}{\texttt{Patch.Branches}}
& branch conditions \\
\(M_\delta\) &
\href{https://github.com/yang-mills-group/OpenAI-YM-20261028/blob/3456e5e75cf8d4cd558955a5c7638fba2d68f464/YM/RG/Patch.lean\#L109-L110}{\texttt{usesMin}}
& minimizing branch \\
\(S_W(h)\) &
\href{https://github.com/yang-mills-group/OpenAI-YM-20261028/blob/3456e5e75cf8d4cd558955a5c7638fba2d68f464/YM/RG/Patch.lean\#L114-L115}{\texttt{score}}
& synchronization cost on a patch \\
frame \(h\), tag \(t=(r,o)\) &
\href{https://github.com/yang-mills-group/OpenAI-YM-20261028/blob/3456e5e75cf8d4cd558955a5c7638fba2d68f464/YM/RG/Patch.lean\#L153-L155}{\texttt{frame}},
\texttt{FrameTag} & gauge fixing on a patch; root and axis ordering \\
\(\mathcal E(W)\) &
\href{https://github.com/yang-mills-group/OpenAI-YM-20261028/blob/3456e5e75cf8d4cd558955a5c7638fba2d68f464/YM/RG/Patch.lean\#L143-L146}{\texttt{eligible}}
& eligible axial tags \\
\(w\) &
\href{https://github.com/yang-mills-group/OpenAI-YM-20261028/blob/3456e5e75cf8d4cd558955a5c7638fba2d68f464/YM/RG/Patch.lean\#L87-L88}{\texttt{weight}}
& \(J\)-fold convolution of the block-uniform weight \\
\(Y\), \(\Phi_Y\), \(Q\) &
\href{https://github.com/yang-mills-group/OpenAI-YM-20261028/blob/3456e5e75cf8d4cd558955a5c7638fba2d68f464/YM/RG/Block.lean\#L39-L42}{\texttt{connector}},
\href{https://github.com/yang-mills-group/OpenAI-YM-20261028/blob/3456e5e75cf8d4cd558955a5c7638fba2d68f464/YM/RG/Patch.lean\#L159}{\texttt{bondCost}},
\href{https://github.com/yang-mills-group/OpenAI-YM-20261028/blob/3456e5e75cf8d4cd558955a5c7638fba2d68f464/YM/RG/Block.lean\#L44-L45}{\texttt{prediction}}
& connector, bond cost, prediction \\
\(\lambda(U,t)\) &
\href{https://github.com/yang-mills-group/OpenAI-YM-20261028/blob/3456e5e75cf8d4cd558955a5c7638fba2d68f464/YM/RG/Block.lean\#L49-L51}{\texttt{likelihood}}
& tag law \\
ch, pr &
\href{https://github.com/yang-mills-group/OpenAI-YM-20261028/blob/3456e5e75cf8d4cd558955a5c7638fba2d68f464/YM/Gauge/MatrixGroup.lean\#L115}{\texttt{Gauge.chord}},
\href{https://github.com/yang-mills-group/OpenAI-YM-20261028/blob/3456e5e75cf8d4cd558955a5c7638fba2d68f464/YM/Gauge/MatrixGroup.lean\#L112}{\texttt{Gauge.price}}
& chord distance to \(1\); Wilson plaquette density \\
\(E_X\), \(F_X\) &
\href{https://github.com/yang-mills-group/OpenAI-YM-20261028/blob/3456e5e75cf8d4cd558955a5c7638fba2d68f464/YM/Branch/Slide.lean\#L24}{\texttt{Axial.edges}},
\href{https://github.com/yang-mills-group/OpenAI-YM-20261028/blob/3456e5e75cf8d4cd558955a5c7638fba2d68f464/YM/Branch/Slide.lean\#L32}{\texttt{Axial.faces}}
& link and plaquette chord sums in \(X\) \\
\(m(h)\), \(F_x(h)\) & \texttt{Pairs.centroid}, \texttt{Pairs.feedback}
& centroid and feedback of the synchronization problem \\
\(\pi\) &
\href{https://github.com/yang-mills-group/OpenAI-YM-20261028/blob/3456e5e75cf8d4cd558955a5c7638fba2d68f464/YM/Gauge/Nearest.lean\#L91}{\texttt{Gauge.focus}}
& nearest point of \(G\) (near \(1\)) \\
port, \(R'\) &
\href{https://github.com/yang-mills-group/OpenAI-YM-20261028/blob/3456e5e75cf8d4cd558955a5c7638fba2d68f464/YM/Pressure/Cost.lean\#L60}{\texttt{port}},
\href{https://github.com/yang-mills-group/OpenAI-YM-20261028/blob/3456e5e75cf8d4cd558955a5c7638fba2d68f464/YM/Pressure/Cost.lean\#L37}{\texttt{macroR}}
& fine curvature in the macro-box of radius \((J+7)L\) \\
usable &
\href{https://github.com/yang-mills-group/OpenAI-YM-20261028/blob/3456e5e75cf8d4cd558955a5c7638fba2d68f464/YM/Pressure/Cost.lean\#L62-L63}{\texttt{usable}}
& positive frame masses on the \(5^4\) neighbours \\
toll, wire, fuel &
\href{https://github.com/yang-mills-group/OpenAI-YM-20261028/blob/3456e5e75cf8d4cd558955a5c7638fba2d68f464/YM/Pressure/Rows.lean\#L55}{\texttt{toll}},
\href{https://github.com/yang-mills-group/OpenAI-YM-20261028/blob/3456e5e75cf8d4cd558955a5c7638fba2d68f464/YM/Pressure/Cost.lean\#L187}{\texttt{wire}},
\href{https://github.com/yang-mills-group/OpenAI-YM-20261028/blob/3456e5e75cf8d4cd558955a5c7638fba2d68f464/YM/Pressure/Cost.lean\#L189-L190}{\texttt{fuel}}
& plaquette chord; chord of \(V_bQ_b^{-1}\); sum of four wires \\
\(\mathcal M+d_0\mathcal F\), \(\tau\) &
\href{https://github.com/yang-mills-group/OpenAI-YM-20261028/blob/3456e5e75cf8d4cd558955a5c7638fba2d68f464/YM/PREC/FrameResponse.lean\#L68}{\texttt{probe}},
\href{https://github.com/yang-mills-group/OpenAI-YM-20261028/blob/3456e5e75cf8d4cd558955a5c7638fba2d68f464/YM/PREC/FrameResponse.lean\#L65-L66}{\texttt{mean}},
\href{https://github.com/yang-mills-group/OpenAI-YM-20261028/blob/3456e5e75cf8d4cd558955a5c7638fba2d68f464/YM/PREC/FrameResponse.lean\#L63}{\texttt{fixRoot}},
\href{https://github.com/yang-mills-group/OpenAI-YM-20261028/blob/3456e5e75cf8d4cd558955a5c7638fba2d68f464/YM/PREC/FrameResponse.lean\#L57-L59}{\texttt{tau}}
& linearized prediction (averaging operator) \\
forest, trunks, fold &
\href{https://github.com/yang-mills-group/OpenAI-YM-20261028/blob/3456e5e75cf8d4cd558955a5c7638fba2d68f464/YM/Tree/Pines.lean\#L24-L29}{\texttt{Tree.Forest}},
\href{https://github.com/yang-mills-group/OpenAI-YM-20261028/blob/3456e5e75cf8d4cd558955a5c7638fba2d68f464/YM/Tree/Pines.lean\#L77}{\texttt{trunks}},
\href{https://github.com/yang-mills-group/OpenAI-YM-20261028/blob/3456e5e75cf8d4cd558955a5c7638fba2d68f464/YM/Tree/Cut.lean\#L21}{\texttt{foldAt}}
& tree gauge in blocks \\
budget point, flags &
\href{https://github.com/yang-mills-group/OpenAI-YM-20261028/blob/3456e5e75cf8d4cd558955a5c7638fba2d68f464/YM/RG/Masks.lean\#L125-L136}{\texttt{BudgetPoint}},
\href{https://github.com/yang-mills-group/OpenAI-YM-20261028/blob/3456e5e75cf8d4cd558955a5c7638fba2d68f464/YM/RG/Masks.lean\#L34-L36}{\texttt{flags}}
& norm parameters; large-field sites \\
\(\mathcal R\), eval &
\href{https://github.com/yang-mills-group/OpenAI-YM-20261028/blob/3456e5e75cf8d4cd558955a5c7638fba2d68f464/YM/RG/History.lean\#L165-L171}{\texttt{Rep}},
\href{https://github.com/yang-mills-group/OpenAI-YM-20261028/blob/3456e5e75cf8d4cd558955a5c7638fba2d68f464/YM/RG/History.lean\#L179}{\texttt{Rep.eval}}
& effective density: exponential of the vacuum term minus the action and
regular activities, times the gas \\
\(\hat\beta_m\), \(t_m\) &
\href{https://github.com/yang-mills-group/OpenAI-YM-20261028/blob/3456e5e75cf8d4cd558955a5c7638fba2d68f464/YM/RG/History.lean\#L79}{\texttt{Calendar.hatBeta}},
\href{https://github.com/yang-mills-group/OpenAI-YM-20261028/blob/3456e5e75cf8d4cd558955a5c7638fba2d68f464/YM/RG/History.lean\#L83}{\texttt{Calendar.threshold}}
& reference coupling; small-field threshold
\(\mathrm{temp}(\hat\beta_m)\) \\
tube, works &
\href{https://github.com/yang-mills-group/OpenAI-YM-20261028/blob/3456e5e75cf8d4cd558955a5c7638fba2d68f464/YM/Pressure/Rows.lean\#L94-L99}{\texttt{Tube}},
\href{https://github.com/yang-mills-group/OpenAI-YM-20261028/blob/3456e5e75cf8d4cd558955a5c7638fba2d68f464/YM/Pressure/Rows.lean\#L104}{\texttt{Tube.vision}}
& transfer tower with fine and coarse path data; six local works \\
store, plot, Work &
\href{https://github.com/yang-mills-group/OpenAI-YM-20261028/blob/3456e5e75cf8d4cd558955a5c7638fba2d68f464/YM/Pressure/Seeds.lean\#L34-L37}{\texttt{Stores}},
\href{https://github.com/yang-mills-group/OpenAI-YM-20261028/blob/3456e5e75cf8d4cd558955a5c7638fba2d68f464/YM/Pressure/Seeds.lean\#L51-L58}{\texttt{Stores.Plot}},
\href{https://github.com/yang-mills-group/OpenAI-YM-20261028/blob/3456e5e75cf8d4cd558955a5c7638fba2d68f464/YM/Pressure/Seeds.lean\#L85}{\texttt{assessed}}
& seeds; their large-field characterization; local work \\
estimator data, energy-per-seed requirement &
\href{https://github.com/yang-mills-group/OpenAI-YM-20261028/blob/3456e5e75cf8d4cd558955a5c7638fba2d68f464/YM/Pressure/Seeds.lean\#L89-L96}{\texttt{Teller}},
\href{https://github.com/yang-mills-group/OpenAI-YM-20261028/blob/3456e5e75cf8d4cd558955a5c7638fba2d68f464/YM/Pressure/Seeds.lean\#L129-L131}{\texttt{Teller.Grants}}
& parameters of one energy estimate; energy-per-seed requirement \\
\(t=\mathrm{temp}(\hat\beta+\kappa_{\rm slope}\log L)\) &
\href{https://github.com/yang-mills-group/OpenAI-YM-20261028/blob/3456e5e75cf8d4cd558955a5c7638fba2d68f464/YM/Pressure/Seeds.lean\#L26}{\texttt{fiery}},
\href{https://github.com/yang-mills-group/OpenAI-YM-20261028/blob/3456e5e75cf8d4cd558955a5c7638fba2d68f464/YM/Pressure/Seeds.lean\#L25}{\texttt{temp}},
\href{https://github.com/yang-mills-group/OpenAI-YM-20261028/blob/3456e5e75cf8d4cd558955a5c7638fba2d68f464/YM/Pressure/Seeds.lean\#L23}{\texttt{betaF}}
& fine-level threshold \\
\(\varrho\) &
\href{https://github.com/yang-mills-group/OpenAI-YM-20261028/blob/3456e5e75cf8d4cd558955a5c7638fba2d68f464/YM/Pressure/Seeds.lean\#L22}{\texttt{radiusB}}
& cluster radius \(\lceil(\log\hat\beta)^2\rceil\) \\
kinetic data, directed history &
\href{https://github.com/yang-mills-group/OpenAI-YM-20261028/blob/3456e5e75cf8d4cd558955a5c7638fba2d68f464/YM/PCalc/Loop.lean\#L93-L115}{\texttt{Kin.Planned}},
\href{https://github.com/yang-mills-group/OpenAI-YM-20261028/blob/3456e5e75cf8d4cd558955a5c7638fba2d68f464/YM/PCalc/Band.lean\#L26-L36}{\texttt{Directed}}
& one-loop shift and stationary-branch requirements \\
\end{longtable}
}

\subsubsection{7. What remains
unexplained}\label{what-remains-unexplained-3}

\begin{enumerate}
\def\labelenumi{\arabic{enumi}.}
\tightlist
\item
  \textbf{Energy estimates.} III.1.10 is an identity for the clipped
  direct term and truncated rows. The estimates converting it into
  III.1.11, including bounds on \(W_D\) and the truncation remainders,
  belong to III.10. No sign is asserted here for
  \(1-\Phi^\dagger\Phi-\Phi'^\dagger\Phi'-\mathcal N^\dagger\mathcal N\)
  at arbitrary \(V\).
\item
  \textbf{Construction of the kinetic data.} The stationary conditions
  and the outputs of the two cited construction theorems are identified
  in III.1.12. Their full construction from the estimates of III.3,
  III.6--III.7 and III.9 is outside this unit. The alternative
  constructor
  \href{https://github.com/yang-mills-group/OpenAI-YM-20261028/blob/3456e5e75cf8d4cd558955a5c7638fba2d68f464/YM/Extra/ITERATION_NO_LOSS/Planned.lean\#L37-L55}{\texttt{Extra.Iteration.planned}}
  has the required output but is outside the main-theorem closure.
\item
  \textbf{Choice of neighbourhood and smoothing order.} The curvature
  proof uses pads at plaquette corners in \(B_1\), whereas the stated
  synchronization result supplies them throughout \(B_2\). The need for
  the larger neighbourhood in later estimates is not determined here.
  Likewise, \(w\le L^{-4}\) needs only \(J\ge1\); the later use of the
  stronger patch condition \(J\ge4\) remains to be specified.
\item
  \textbf{Additional files.} The small packages listed in the header are
  not described by the numbered results. Their contributions, and any
  further results from the exterior calculus or quadratic-kernel files
  beyond those cited here, remain to be incorporated into the relevant
  later units.
\end{enumerate}

\subsection{III.2. Choosing the constants: the order of choice, the
level schedule and the shooting
argument}\label{iii.2.-choosing-the-constants-the-order-of-choice-the-level-schedule-and-the-shooting-argument}

\textbf{Part:} III, Ultraviolet construction (renormalization group).
\textbf{Lean:}
\href{https://github.com/yang-mills-group/OpenAI-YM-20261028/tree/3456e5e75cf8d4cd558955a5c7638fba2d68f464/YM/Price}{\texttt{YM/Price/}}
(the constants and the level schedule, 406 lines),
\href{https://github.com/yang-mills-group/OpenAI-YM-20261028/tree/3456e5e75cf8d4cd558955a5c7638fba2d68f464/YM/Shoot}{\texttt{YM/Shoot/}}
(shooting, 228 lines),
\href{https://github.com/yang-mills-group/OpenAI-YM-20261028/tree/3456e5e75cf8d4cd558955a5c7638fba2d68f464/YM/Integration/RG05Preconstruction}{\texttt{YM/Integration/RG05Preconstruction/}},
\href{https://github.com/yang-mills-group/OpenAI-YM-20261028/tree/3456e5e75cf8d4cd558955a5c7638fba2d68f464/YM/Integration/NativeSelection}{\texttt{YM/Integration/NativeSelection/}}
and other files of
\href{https://github.com/yang-mills-group/OpenAI-YM-20261028/tree/3456e5e75cf8d4cd558955a5c7638fba2d68f464/YM/Integration}{\texttt{YM/Integration/}},
\href{https://github.com/yang-mills-group/OpenAI-YM-20261028/tree/3456e5e75cf8d4cd558955a5c7638fba2d68f464/YM/RG08/Calendar}{\texttt{YM/RG08/Calendar/}},
and parts of
\href{https://github.com/yang-mills-group/OpenAI-YM-20261028/tree/3456e5e75cf8d4cd558955a5c7638fba2d68f464/YM/RG08}{\texttt{YM/RG08/}},
\href{https://github.com/yang-mills-group/OpenAI-YM-20261028/tree/3456e5e75cf8d4cd558955a5c7638fba2d68f464/YM/RG12}{\texttt{YM/RG12/}},
\href{https://github.com/yang-mills-group/OpenAI-YM-20261028/tree/3456e5e75cf8d4cd558955a5c7638fba2d68f464/YM/FullVolume}{\texttt{YM/FullVolume/}}
(in particular \texttt{Roll/}) and
\href{https://github.com/yang-mills-group/OpenAI-YM-20261028/tree/3456e5e75cf8d4cd558955a5c7638fba2d68f464/YM/TraceExtension}{\texttt{YM/TraceExtension/}}.
\textbf{Size:} 18,156 lines in the main-theorem closure:

{\def\LTcaptype{none} % do not increment counter
\begin{longtable}[]{@{}
  >{\raggedright\arraybackslash}p{(\linewidth - 2\tabcolsep) * \real{0.5000}}
  >{\raggedright\arraybackslash}p{(\linewidth - 2\tabcolsep) * \real{0.5000}}@{}}
\toprule\noalign{}
\begin{minipage}[b]{\linewidth}\raggedright
part
\end{minipage} & \begin{minipage}[b]{\linewidth}\raggedright
lines
\end{minipage} \\
\midrule\noalign{}
\endhead
\bottomrule\noalign{}
\endlastfoot
proof of the plentiful construction (III.2.5) & about 6,250 \\
proof of the joint same-factor purchase (III.2.4) & about 4,700 \\
proof of the final run (III.2.7) & about 4,000 \\
proof of the sheet and radar purchase & about 2,400 \\
the two directories that carry the mathematics of the schedule
(\texttt{YM/Price/}, \texttt{YM/Shoot/}) & 634 \\
\end{longtable}
}

\textbf{Conventions.} These are those of III.9.

\begin{itemize}
\tightlist
\item
  \textbf{Levels.} Levels \(m=0,1,\dots,N\) run from the coarsest
  (\(m=0\), unit spacing) to the finest (\(m=N\), the Wilson lattice).
\item
  \textbf{Steps and scales.} The step \(m+1\to m\) blocks by an odd
  factor \(L_m\ge3\). The scale is \(s_m=\sum_{i<m}\log L_i\). The
  level-\(m\) torus at volume index \(p\) has side \(n_p e^{s_m}\),
  where \(n_p\) is the side of the level-0 torus. III.9 writes \(L_p\)
  for \(n_p\); the letter \(L\) is kept here for block factors.
\item
  \textbf{Slope.} \(\kappa=11C_G/(24\pi^2)\) (III.9.1), and \(\kappa>0\)
  for simple \(G\).
\item
  \textbf{Reference quantities (III.9.2).}
  \(\hat\beta_m=\bar\beta+\kappa s_m\),
  \(M_m=e^{(\log\log\hat\beta_m)^2}\), the threshold
  \(t_m=\hat\beta_m^{-1/2}M_m\), the cover size
  \(B_m=\lceil(\log\hat\beta_m)^2\rceil\) and the band
  \(\mathfrak B_m=[\hat\beta_m/2,\,2\hat\beta_m]\).
\item
  \textbf{Patch radius.} \(R_m=(J_{m+1}+3)L_m\)
  (\href{https://github.com/yang-mills-group/OpenAI-YM-20261028/blob/3456e5e75cf8d4cd558955a5c7638fba2d68f464/YM/RG/Patch.lean\#L51}{\texttt{Patch.radius}}).
\end{itemize}

The main results concern numerical choices and a conditional shooting
construction:

\begin{enumerate}
\def\labelenumi{\arabic{enumi}.}
\tightlist
\item
  A long list of requirements can be met simultaneously by choosing the
  constants in one fixed order. Each requirement has the form ``this
  inequality holds once the block factor, the coupling or the lattice is
  large enough, given the constants chosen before''.
\item
  There is a level schedule meeting every entrance condition.
\item
  Given the one-step construction as a hypothesis, there is a trajectory
  of effective actions with \(\beta_{N,p,0}=\bar\beta\) (shooting).
\end{enumerate}

The terminology of sheets, prices and purchases records numerical bounds
and their order of choice (glossary, family 1). The unit also contains a
group of analytic lemmas of the Gaussian step, filed here by the
dependency assignment (III.2.9).

\subsubsection{1. What this unit proves}\label{what-this-unit-proves-4}

\textbf{Proposition III.2.1 (the constants, in the order in which they
are chosen).} Let \(G\) be simple. The construction fixes the following
data, in this order.

\begin{enumerate}
\def\labelenumi{\arabic{enumi}.}
\item
  \emph{Sheet (data chosen before the block factors).} It consists of
  the following.

  \begin{itemize}
  \tightlist
  \item
    A kernel choice \(K\) with branches on every patch.
  \item
    For every stage \(j\in\mathbb N\), a stage record \(X_j\) with:

    \begin{itemize}
    \tightlist
    \item
      a clip with relaxation \(r_j\);
    \item
      derivative order \(k_j\ge6\) and spatial power \(q_j\ge2\);
    \item
      support exponent \(a_j>0\), width \(w_j>0\) and path count
      \(\varpi_j>0\);
    \item
      drift allowance \(\omega_j\ge0\) and column bound
      \(\gamma_j\ge1\);
    \item
      the constraint \(r_jw_j\ge12\).
    \end{itemize}
  \item
    A first stage \(j_0\ge4\).
  \item
    Floors \(f_j\in\mathbb N\).
  \end{itemize}

  The sequences \(k_j\), \(q_j\), \(a_j\) are nondecreasing.
\item
  \emph{Stage block factors} \(\ell_j\). They are odd, satisfy
  \(\ell_j\ge\max(3,f_j)\), are nondecreasing in \(j\), and for
  \(j\ge j_0\)
  \[\frac{\omega_j}{\log\ell_j}\le\min\Big(\frac\kappa4,\ \frac1{j^2}\Big),\qquad \frac{\log\mathcal C_j}{\log\ell_j}\le\frac1{j^2},\qquad \mathcal C_j:=1+\varpi_j+w_j^{-1}+\gamma_j+r_j+a_j .\]
\item
  \emph{Smallness exponents} \(\theta_j>0\). They are nonincreasing,
  with \(\theta_j\le1/j\) for \(j\ge j_0\).
\item
  \emph{Delays.} These are the following.

  \begin{itemize}
  \tightlist
  \item
    Entry couplings \(\beta^{\rm ent}_j\in\mathbb R\).
  \item
    Entry scales \(\sigma^{\rm ent}_j(\bar\beta)\). These are functions
    of the stage \emph{and of the barrier} still to be chosen.
  \item
    Volume shells \(S_j\in\mathbb N\) and growth exponents
    \(g_j\in\mathbb N\).
  \end{itemize}

  Delays are ordered componentwise.
\item
  \emph{Barrier} \(\bar\beta>1\). It satisfies
  \(\beta^{\rm ent}_{j_0}\le\bar\beta\) and
  \(x^{-1/2}e^{(\log\log x)^2}\le1\) for all \(x\ge\bar\beta\).
\item
  \emph{Wait.} This consists of:

  \begin{itemize}
  \tightlist
  \item
    readiness scales \(\sigma^{\rm rdy}_j\);
  \item
    windows \(h_j:\mathbb R\to\mathbb N\) with \(h_j(s)/s\to0\) as
    \(s\to\infty\) (for each fixed \(j\), not uniformly);
  \item
    a minimal period \(n_{\min}\).
  \end{itemize}

  Waits are ordered componentwise by readiness scales, windows and
  minimal period.
\end{enumerate}

Items 2, 3 and 5 always exist:

\begin{itemize}
\tightlist
\item
  item 2 exists for every sheet;
\item
  in item 3, \(\theta\) can be chosen below any prescribed positive
  sequence;
\item
  in item 5, \(\bar\beta\) can be chosen above any prescribed real
  number.
\end{itemize}

Item 4 has no further conditions. In item 6, only the stated
sublinearity is required; the readiness scales are real and
\(n_{\min}\in\mathbb N\).

\emph{Status:} items 1--6 are definitions (specifications). There are
three existence theorems. The existence of item 2 uses \(\kappa>0\),
hence the simplicity of \(G\). The later existence theorems
III.2.4--III.2.7 meet specified requirements in this order. \emph{Lean:}
\href{https://github.com/yang-mills-group/OpenAI-YM-20261028/blob/3456e5e75cf8d4cd558955a5c7638fba2d68f464/YM/Price/Intent.lean\#L26-L35}{\texttt{Sheet}},
\href{https://github.com/yang-mills-group/OpenAI-YM-20261028/blob/3456e5e75cf8d4cd558955a5c7638fba2d68f464/YM/RG/History.lean\#L45-L61}{\texttt{Stage}},
\href{https://github.com/yang-mills-group/OpenAI-YM-20261028/blob/3456e5e75cf8d4cd558955a5c7638fba2d68f464/YM/Price/Intent.lean\#L39-L48}{\texttt{Large}},
\href{https://github.com/yang-mills-group/OpenAI-YM-20261028/blob/3456e5e75cf8d4cd558955a5c7638fba2d68f464/YM/Price/Intent.lean\#L50-L53}{\texttt{Rolls}},
\href{https://github.com/yang-mills-group/OpenAI-YM-20261028/blob/3456e5e75cf8d4cd558955a5c7638fba2d68f464/YM/Price/Intent.lean\#L79-L84}{\texttt{Rolls.exists\_rolls}},
\href{https://github.com/yang-mills-group/OpenAI-YM-20261028/blob/3456e5e75cf8d4cd558955a5c7638fba2d68f464/YM/Price/Intent.lean\#L88-L92}{\texttt{Protection}},
\href{https://github.com/yang-mills-group/OpenAI-YM-20261028/blob/3456e5e75cf8d4cd558955a5c7638fba2d68f464/YM/Price/Intent.lean\#L98-L111}{\texttt{Protection.exists\_margin}},
\href{https://github.com/yang-mills-group/OpenAI-YM-20261028/blob/3456e5e75cf8d4cd558955a5c7638fba2d68f464/YM/Price/Intent.lean\#L115-L120}{\texttt{Delays}},
\href{https://github.com/yang-mills-group/OpenAI-YM-20261028/blob/3456e5e75cf8d4cd558955a5c7638fba2d68f464/YM/Push/PRICE_COMPILER/Post.lean\#L34-L38}{\texttt{DelayLe}},
\href{https://github.com/yang-mills-group/OpenAI-YM-20261028/blob/3456e5e75cf8d4cd558955a5c7638fba2d68f464/YM/Price/Intent.lean\#L123-L127}{\texttt{Bar}},
\href{https://github.com/yang-mills-group/OpenAI-YM-20261028/blob/3456e5e75cf8d4cd558955a5c7638fba2d68f464/YM/Price/Length.lean\#L19-L37}{\texttt{Bar.exists\_above}},
\href{https://github.com/yang-mills-group/OpenAI-YM-20261028/blob/3456e5e75cf8d4cd558955a5c7638fba2d68f464/YM/Price/Intent.lean\#L134-L139}{\texttt{Wait}}.

\textbf{Proposition III.2.2 (a level schedule meeting all entrance
conditions).} Let \(G\) be simple and let data 1--6 of III.2.1 be given.
Then there is a level system with the following data:

\begin{itemize}
\tightlist
\item
  a stage index \(J_m\);
\item
  step block factors \(L_m=\ell_{J_{m+1}}\);
\item
  level-0 sides \(n_p=4(M+p+1)\) for some \(M\in\mathbb N\);
\item
  reference couplings \(\hat\beta_m=\bar\beta+\kappa s_m\).
\end{itemize}

It satisfies:

\begin{itemize}
\tightlist
\item
  \begin{enumerate}
  \def\labelenumi{(\alph{enumi})}
  \tightlist
  \item
    \emph{identification:} the kernel, stage records, \(\ell\),
    \(\theta\) and \(\bar\beta\) are those of III.2.1, and
    \(J_0=J_1=j_0\);
  \end{enumerate}
\item
  (b) \emph{entry:} \(\beta^{\rm ent}_{J_{m+1}}\le\hat\beta_m\) for all
  \(m\), and \(\sigma^{\rm ent}_{J_{m+1}}(\bar\beta)\le s_m\) whenever
  \(J_{m+1}>j_0\);
\item
  (c) \emph{volume:} \(n_p\ge n_{\min}\), and
  \(S_j(1+B_m+B_{m+1})^{g_j}\le n_pe^{s_m}\) for all \(m,p\), with
  \(j=J_{m+1}\);
\item
  (d) \emph{windows:} for every \(j\ge j_0\) there is a level \(k\) with
  \(\sigma^{\rm rdy}_j\le s_k\), \(h_j(s_k)\le k\), and \(J_i=j\) for
  \(|i-k|\le h_j(s_k)\);
\item
  (e) \emph{the scale conditions of III.9.2:}

  \begin{itemize}
  \tightlist
  \item
    \(J\) is nondecreasing, \(J_{m+1}\le J_m+1\) and \(J_m\to\infty\);
  \item
    \(k_{J_m}\), \(q_{J_m}\), \(a_{J_m}\) are nondecreasing in \(m\);
  \item
    \(\hat\beta_m>1\) and \(t_m\le1\);
  \item
    \(100(B_m+1)\le n_pe^{s_m}\) and \(100(R_m+1)\le n_pe^{s_{m+1}}\);
  \item
    \(0<\theta_{J_m}\le1/J_m\), and \(\theta\) is nonincreasing;
  \item
    \(\omega_{J_{m+1}}/\log L_m\le J_{m+1}^{-2}\) and
    \(\log\mathcal C_{J_{m+1}}/\log L_m\le J_{m+1}^{-2}\);
  \item
    \(\log L_m/s_m\to0\);
  \end{itemize}
\item
  (f) \emph{block control:} at every step the kernel is well posed,
  symmetric under the lattice symmetries and gauge covariant.
\end{itemize}

The same construction also ensures that every positive integer divides
\(n_p\) for infinitely many \(p\).

\emph{Status:} theorem. The divisibility variant is filed in III.10 and
is the one used in the final assembly (III.2.7). Note that (d) asserts
the existence of one level \(k\) per stage. Taking
\(\sigma^{\rm rdy}_j\) large forces \(k\) to be large. \emph{Lean:}
\href{https://github.com/yang-mills-group/OpenAI-YM-20261028/blob/3456e5e75cf8d4cd558955a5c7638fba2d68f464/YM/Price/Waits.lean\#L137-L239}{\texttt{Wait.ready}},
\href{https://github.com/yang-mills-group/OpenAI-YM-20261028/blob/3456e5e75cf8d4cd558955a5c7638fba2d68f464/YM/Price/Intent.lean\#L145-L162}{\texttt{Paid}},
\href{https://github.com/yang-mills-group/OpenAI-YM-20261028/blob/3456e5e75cf8d4cd558955a5c7638fba2d68f464/YM/RG/History.lean\#L97-L118}{\texttt{Calendar.Fits}},
\href{https://github.com/yang-mills-group/OpenAI-YM-20261028/blob/3456e5e75cf8d4cd558955a5c7638fba2d68f464/YM/RG/History.lean\#L131-L134}{\texttt{Levels}},
\href{https://github.com/yang-mills-group/OpenAI-YM-20261028/blob/3456e5e75cf8d4cd558955a5c7638fba2d68f464/YM/Lattice/Scale.lean\#L129-L130}{\texttt{Schedule.Gentle}},
\href{https://github.com/yang-mills-group/OpenAI-YM-20261028/blob/3456e5e75cf8d4cd558955a5c7638fba2d68f464/YM/Branch/Covariant.lean\#L70-L72}{\texttt{Calendar.blocks\_go}};
variant
\href{https://github.com/yang-mills-group/OpenAI-YM-20261028/blob/3456e5e75cf8d4cd558955a5c7638fba2d68f464/YM/PhysicalNext/RG_NEXT_COMPLETE_REPEATABLE_NONLINEAR_STEP/Pane/Divisors.lean\#L26-L135}{\texttt{Wait.engine}},
\href{https://github.com/yang-mills-group/OpenAI-YM-20261028/blob/3456e5e75cf8d4cd558955a5c7638fba2d68f464/YM/PhysicalNext/RG_NEXT_COMPLETE_REPEATABLE_NONLINEAR_STEP/Pane/Divisors.lean\#L21}{\texttt{AsDivisors}}.

\textbf{Proposition III.2.3 (shooting: a selected history from a
recipe).} Let \(G\) be simple. Fix a level system with
\[\omega_{J_{m+1}}\le\tfrac\kappa4\log L_m\qquad\text{for all } m.\]
This holds for every level system of III.2.2, by III.2.1(2).

\emph{Recipe.} Let a recipe be given: path data \(P^{\mathcal R}_{N,m}\)
and three predicates (a start test, a retention test at each level, a
transition test for consecutive levels).

\emph{Admissible representations.} A level-\(m\) representation \(\rho\)
(III.9: path data, clip, activities, budget point, \(\beta\), \(v\)) is
\emph{admissible for \((N,p)\)} if:

\begin{itemize}
\tightlist
\item
  its path data are \(P^{\mathcal R}_{N,m}\);
\item
  its budget point, clip and polymer bounds are those of level \(m\):
  width \(w_{J_m}\), path count \(\varpi_{J_m}\), loop cap
  \(2M_m^{\theta_{J_m}}\), allowances \(\hat\beta_m^{-11/10}\) and
  \(e^{-M_m^{1/5}}\);
\item
  it passes the retention test.
\end{itemize}

\emph{Continuous families.} A family \(x\mapsto\rho_x\) on an interval
is \emph{continuous} if:

\begin{itemize}
\tightlist
\item
  \(\beta\) and \(v\) are continuous in \(x\);
\item
  the address sets of the regular and forced activities do not depend on
  \(x\);
\item
  differences of regular activities tend to zero in the sum of budget
  masses of their derivative bounds, up to the budget order on the
  healthy domain; differences of forced activities tend to zero in the
  sum of budget masses of their uniform bounds over all configurations.
\end{itemize}

\emph{Hypotheses.} Assume that for every \(N\) and \(p\):

\begin{itemize}
\tightlist
\item
  \textbf{(start)} There is a continuous family
  \(\beta\mapsto\rho_\beta\), \(\beta\in\mathfrak B_N\), of admissible
  level-\(N\) representations. Each \(\rho_\beta\) has coupling
  \(\beta\), \(v=0\), density equal to the Wilson weight at \(\beta\),
  and passes the start test.
\item
  \textbf{(advance)} Let \(m<N\) and \(l\le u\), and let
  \(x\mapsto\rho'_x\), \(x\in[l,u]\), be a continuous family of
  admissible level-\((m+1)\) representations with
  \(\beta(\rho'_x)\in\mathfrak B_{m+1}\) for all \(x\). Then there is a
  continuous family \(x\mapsto\rho_x\) of admissible level-\(m\)
  representations such that, for every \(x\):

  \begin{itemize}
  \tightlist
  \item
    the density of \(\rho_x\) is the block-spin push-forward of the
    density of \(\rho'_x\) at coupling \(\beta(\rho'_x)\);
  \item
    \(|\beta(\rho'_x)-\beta(\rho_x)-\kappa\log L_m|\le\omega_{J_{m+1}}\);
  \item
    the transition test holds.
  \end{itemize}
\end{itemize}

\emph{Conclusion.} There are representations \(\rho_{N,p,m}\),
\(0\le m\le N\), for all \(N,p\), such that:

\begin{itemize}
\tightlist
\item
  every \(\rho_{N,p,m}\) is admissible with
  \(\beta_{N,p,m}\in\mathfrak B_m\);
\item
  \(\rho_{N,p,N}\) is the Wilson weight with \(v=0\);
\item
  consecutive levels are related as in (advance);
\item
  \(\beta_{N,p,0}=\bar\beta\);
\item
  the path data at \((N,m)\) do not depend on \(p\).
\end{itemize}

The resulting history is \emph{selected} in the sense of III.9.3. In
particular it obeys the running law: for every \(\varepsilon>0\) there
is \(m_0\) such that \(|\beta_{N,p,m}-\hat\beta_m|\le\varepsilon s_m\)
for all \(m\ge m_0\), \(N\ge m\) and all \(p\).

\emph{Status:} theorem, conditional on (start) and (advance).

\begin{itemize}
\tightlist
\item
  In the main chain both are verified in III.10
  (\href{https://github.com/yang-mills-group/OpenAI-YM-20261028/blob/3456e5e75cf8d4cd558955a5c7638fba2d68f464/YM/PhysicalNext/RG_NEXT_COMPLETE_REPEATABLE_NONLINEAR_STEP/Graces/Eagles.lean\#L33-L50}{\texttt{Graces.winnow}}):

  \begin{itemize}
  \tightlist
  \item
    (start) from the Wilson start-up curves of III.9
    (\href{https://github.com/yang-mills-group/OpenAI-YM-20261028/blob/3456e5e75cf8d4cd558955a5c7638fba2d68f464/YM/PhysicalNext/RG_NEXT_COMPLETE_REPEATABLE_NONLINEAR_STEP/Graces/Eagles.lean\#L21-L27}{\texttt{hops}});
  \item
    (advance) from the one-step construction with coupling reset
    (\href{https://github.com/yang-mills-group/OpenAI-YM-20261028/blob/3456e5e75cf8d4cd558955a5c7638fba2d68f464/YM/PhysicalNext/RG_NEXT_COMPLETE_REPEATABLE_NONLINEAR_STEP/Mallow/Pediment.lean\#L41-L45}{\texttt{Mallow.faces}}).
  \end{itemize}
\item
  The drift hypothesis is derived from III.2.2 by
  \href{https://github.com/yang-mills-group/OpenAI-YM-20261028/blob/3456e5e75cf8d4cd558955a5c7638fba2d68f464/YM/Shoot/Reach.lean\#L73-L79}{\texttt{Drifting.pay}}.
\item
  A second route, through the step program of III.1/III.9
  (\href{https://github.com/yang-mills-group/OpenAI-YM-20261028/blob/3456e5e75cf8d4cd558955a5c7638fba2d68f464/YM/PCalc/Band.lean\#L65-L72}{\texttt{Directed.Program}}),
  is not used by the main theorem.
\end{itemize}

\emph{Lean:}
\href{https://github.com/yang-mills-group/OpenAI-YM-20261028/blob/3456e5e75cf8d4cd558955a5c7638fba2d68f464/YM/Shoot/Race.lean\#L67-L82}{\texttt{Recipe.races}},
\href{https://github.com/yang-mills-group/OpenAI-YM-20261028/blob/3456e5e75cf8d4cd558955a5c7638fba2d68f464/YM/Shoot/Criteria.lean\#L67-L71}{\texttt{Recipe}},
\href{https://github.com/yang-mills-group/OpenAI-YM-20261028/blob/3456e5e75cf8d4cd558955a5c7638fba2d68f464/YM/Shoot/Criteria.lean\#L74-L79}{\texttt{Lives}},
\href{https://github.com/yang-mills-group/OpenAI-YM-20261028/blob/3456e5e75cf8d4cd558955a5c7638fba2d68f464/YM/Shoot/Criteria.lean\#L51-L57}{\texttt{Curve}},
\href{https://github.com/yang-mills-group/OpenAI-YM-20261028/blob/3456e5e75cf8d4cd558955a5c7638fba2d68f464/YM/Shoot/Criteria.lean\#L88-L91}{\texttt{Recipe.Relays}},
\href{https://github.com/yang-mills-group/OpenAI-YM-20261028/blob/3456e5e75cf8d4cd558955a5c7638fba2d68f464/YM/Shoot/Criteria.lean\#L93-L96}{\texttt{Recipe.Starter}},
\href{https://github.com/yang-mills-group/OpenAI-YM-20261028/blob/3456e5e75cf8d4cd558955a5c7638fba2d68f464/YM/Shoot/Criteria.lean\#L99-L106}{\texttt{Recipe.Issues}},
\href{https://github.com/yang-mills-group/OpenAI-YM-20261028/blob/3456e5e75cf8d4cd558955a5c7638fba2d68f464/YM/Shoot/Reach.lean\#L70-L71}{\texttt{Drifting}},
\href{https://github.com/yang-mills-group/OpenAI-YM-20261028/blob/3456e5e75cf8d4cd558955a5c7638fba2d68f464/YM/RG/History.lean\#L234-L240}{\texttt{History.Selected}}.

\textbf{Proposition III.2.4 (the joint purchase on one emitted factor).}

\emph{Hypotheses.} Assume:

\begin{itemize}
\tightlist
\item
  \(G\) is simple;
\item
  \(K\) is a kernel choice with branches at every patch;
\item
  order data \(v\) are given:

  \begin{itemize}
  \tightlist
  \item
    a first stage \(\ge4\);
  \item
    derivative orders \(k_j\ge\max(6,j)\) and powers
    \(q_j\ge\max(2,j)\), both nondecreasing;
  \item
    support exponents \(a_j\), nondecreasing and \(\ge100\);
  \end{itemize}
\item
  \(Q:\mathbb N\to\mathbb N\) is nondecreasing with \(Q\ge1\);
\item
  the records \(\Omega\)
  (\href{https://github.com/yang-mills-group/OpenAI-YM-20261028/blob/3456e5e75cf8d4cd558955a5c7638fba2d68f464/YM/RemainderComparison/Chron/Age.lean\#L68-L71}{\texttt{Contemple}})
  and \(\Delta\)
  (\href{https://github.com/yang-mills-group/OpenAI-YM-20261028/blob/3456e5e75cf8d4cd558955a5c7638fba2d68f464/YM/RG12/Burst/Wardrobe.lean\#L18-L20}{\texttt{Decal}})
  are given, together with parameters \((\xi,\omega_v,\omega_t)\) and a
  degree-\(300\) interpolation tableau \(\theta_1\) at parameter \(\xi\)
  (\href{https://github.com/yang-mills-group/OpenAI-YM-20261028/blob/3456e5e75cf8d4cd558955a5c7638fba2d68f464/YM/FullVolume/Luna/Cuts.lean\#L18-L24}{\texttt{Tableau}}),
  and the contraction condition
  \href{https://github.com/yang-mills-group/OpenAI-YM-20261028/blob/3456e5e75cf8d4cd558955a5c7638fba2d68f464/YM/FullVolume/Sym/Purse.lean\#L29}{\texttt{sphere}}
  at these parameters; in particular \(\xi\ge2\), \(0<\omega_v\le1\) and
  \(\omega_t>0\);
\item
  \(\pi\) is a predicate on emitted factors (a \emph{prediction}) that
  is prepaid at every sheet and price map.
\end{itemize}

\emph{Conclusion.} There exist:

\begin{itemize}
\tightlist
\item
  a sheet \(s\) with kernel \(K\) and the first stage, orders, powers
  and support exponents prescribed by \(v\);
\item
  a price map \(T\);
\item
  a real sequence \(Y\) and a nondecreasing real sequence \(K'\ge1\);
\item
  a floor \(\tau\ge f^s\).
\end{itemize}

These are followed by the alternation below:

\begin{quote}
for every floor sequence \(k\ge\tau\) (block factors exist for the sheet
\(s_k\), which is \(s\) with floors \(k\)), and every choice of stage
block factors \(\ell\) for \(s_k\), there is a wage threshold \(U\) such
that for every wage \(w\ge U\) and every prescribed postage rate and
exchange datum, there are post-\(L\) cost constants \(F\) such that for
every positive sequence \(\theta^0\) there is a smallness exponent
\(\theta\le\theta^0\) such that for every prescribed menu, auxiliary
record and delay \(T_0\) there is a delay \(V\ge T_0\) such that for
every delay \(V'\ge V\), every barrier \(\bar\beta\) for \(V'\) and
every wait \(W\), there is a wait \(Z\ge W\) such that for every wait
\(W'\ge Z\) and every level system \(a\) that meets the entrance
conditions of \(W'\) (III.2.2(a)--(d), with (e)--(f)) and whose periods
have the divisibility property, there exist a directed datum \(B\) on
\(a\) (kinetic data and a recipe on the level system \(a\)) and an
emitted factor \(X\) such that: - \(X\) records all these choices; -
\(\pi(X)\) holds; - all requirements listed below hold.
\end{quote}

The requirements (with the unit that defines each predicate) are:

{\def\LTcaptype{none} % do not increment counter
\begin{longtable}[]{@{}
  >{\raggedright\arraybackslash}p{(\linewidth - 2\tabcolsep) * \real{0.5000}}
  >{\raggedright\arraybackslash}p{(\linewidth - 2\tabcolsep) * \real{0.5000}}@{}}
\toprule\noalign{}
\begin{minipage}[b]{\linewidth}\raggedright
requirement
\end{minipage} & \begin{minipage}[b]{\linewidth}\raggedright
owner
\end{minipage} \\
\midrule\noalign{}
\endhead
\bottomrule\noalign{}
\endlastfoot
wages, exchanges, menus and auxiliary records equal the prescribed data;
postage rates do not exceed the prescribed rates & III.10
(\href{https://github.com/yang-mills-group/OpenAI-YM-20261028/blob/3456e5e75cf8d4cd558955a5c7638fba2d68f464/YM/PhysicalNext/RG_NEXT_COMPLETE_REPEATABLE_NONLINEAR_STEP/Roost/Foundry.lean\#L26-L50}{\texttt{Factor}}) \\
\href{https://github.com/yang-mills-group/OpenAI-YM-20261028/blob/3456e5e75cf8d4cd558955a5c7638fba2d68f464/YM/FullVolume/Sym/Prices.lean\#L84}{\texttt{Degree}},
\href{https://github.com/yang-mills-group/OpenAI-YM-20261028/blob/3456e5e75cf8d4cd558955a5c7638fba2d68f464/YM/FullVolume/Sym/Shorn.lean\#L32-L35}{\texttt{Savvy}},
\href{https://github.com/yang-mills-group/OpenAI-YM-20261028/blob/3456e5e75cf8d4cd558955a5c7638fba2d68f464/YM/FullVolume/C/Ribbon.lean\#L29}{\texttt{Roofing}},
\href{https://github.com/yang-mills-group/OpenAI-YM-20261028/blob/3456e5e75cf8d4cd558955a5c7638fba2d68f464/YM/FullVolume/Florid/Ready.lean\#L56}{\texttt{PaysW}},
\href{https://github.com/yang-mills-group/OpenAI-YM-20261028/blob/3456e5e75cf8d4cd558955a5c7638fba2d68f464/YM/FullVolume/Luna/Hours.lean\#L43}{\texttt{Spire}},
\href{https://github.com/yang-mills-group/OpenAI-YM-20261028/blob/3456e5e75cf8d4cd558955a5c7638fba2d68f464/YM/FullVolume/Auvi/Entries.lean\#L31}{\texttt{Wild}},
tower bound
\href{https://github.com/yang-mills-group/OpenAI-YM-20261028/blob/3456e5e75cf8d4cd558955a5c7638fba2d68f464/YM/FullVolume/Spiral/Rent.lean\#L25-L27}{\texttt{Tower.Sav}}
with
\href{https://github.com/yang-mills-group/OpenAI-YM-20261028/blob/3456e5e75cf8d4cd558955a5c7638fba2d68f464/YM/FullVolume/Fares/Query.lean\#L45}{\texttt{Canyon}},
\href{https://github.com/yang-mills-group/OpenAI-YM-20261028/blob/3456e5e75cf8d4cd558955a5c7638fba2d68f464/YM/FullVolume/Artes/Events.lean\#L34-L39}{\texttt{Cadre}}
& III.12 \\
\href{https://github.com/yang-mills-group/OpenAI-YM-20261028/blob/3456e5e75cf8d4cd558955a5c7638fba2d68f464/YM/MergeProof/North/Tibetan.lean\#L60}{\texttt{AccordV}}
(with the sequence \(Q\)) & III.11 \\
kinetic data at the top level
(\href{https://github.com/yang-mills-group/OpenAI-YM-20261028/blob/3456e5e75cf8d4cd558955a5c7638fba2d68f464/YM/Planned/RECIPE_STARTUP_PORT/Layout.lean\#L43}{\texttt{AtTop}})
& III.9 \\
amplification
(\href{https://github.com/yang-mills-group/OpenAI-YM-20261028/blob/3456e5e75cf8d4cd558955a5c7638fba2d68f464/YM/FullVolume/Train/Miracles.lean\#L21-L26}{\texttt{Amplify}})
& III.1 \\
support exponents \(\ge100\) at every level & this unit \\
\(\exp((a_j+2)(2\,\mathrm{span}_{ij}+13)^4)\le\ell_i\) for all
\(i,j\in\mathbb N\) with \(j\le i\le j+1\) & this unit \\
a bandage \(b\)
(\href{https://github.com/yang-mills-group/OpenAI-YM-20261028/blob/3456e5e75cf8d4cd558955a5c7638fba2d68f464/YM/PhysicalNext/RG_NEXT_COMPLETE_REPEATABLE_NONLINEAR_STEP/Stage/Iris.lean\#L21-L35}{\texttt{Bandage}},
III.10) with the extras \((\Omega,\Delta,b)\)
(\href{https://github.com/yang-mills-group/OpenAI-YM-20261028/blob/3456e5e75cf8d4cd558955a5c7638fba2d68f464/YM/Integration/RG05Preconstruction/Resources.lean\#L45-L48}{\texttt{Extras}})
& III.15 \\
\end{longtable}
}

Here \(\mathrm{span}_{ij}\in\mathbb N\) is the factor's transition
support parameter; the exponent is given by
\href{https://github.com/yang-mills-group/OpenAI-YM-20261028/blob/3456e5e75cf8d4cd558955a5c7638fba2d68f464/YM/PhysicalNext/V9_JOINT_ENDPOINT_LOAD/Heir/Fusion.lean\#L58}{\texttt{base}}
and
\href{https://github.com/yang-mills-group/OpenAI-YM-20261028/blob/3456e5e75cf8d4cd558955a5c7638fba2d68f464/YM/Planned/ATOM_REGULAR_ESTIMATE/Geometry.lean\#L50}{\texttt{toll}}.

The bandage also supplies the conditional family package of
\href{https://github.com/yang-mills-group/OpenAI-YM-20261028/blob/3456e5e75cf8d4cd558955a5c7638fba2d68f464/YM/FullVolume/Unlatch.lean\#L25-L34}{\texttt{Unlatched}}
with a vacuous additional conclusion; the family hypotheses are those
stated in III.2.5.

\emph{Status:} theorem. The estimates hold on the same emitted factor,
with the displayed order of choice. The hypotheses \(\Omega\),
\(\Delta\), \(\theta_1\), \((\xi,\omega_v,\omega_t)\) and \(v\) are
constructed in the final assembly III.2.7, or supplied there by III.14.
The version without a prediction is
\href{https://github.com/yang-mills-group/OpenAI-YM-20261028/blob/3456e5e75cf8d4cd558955a5c7638fba2d68f464/YM/Integration/RG05Preconstruction/Fares.lean\#L88-L91}{\texttt{richFares}}.
\emph{Lean:}
\href{https://github.com/yang-mills-group/OpenAI-YM-20261028/blob/3456e5e75cf8d4cd558955a5c7638fba2d68f464/YM/Integration/RG05Preconstruction/Fares.lean\#L26-L85}{\texttt{richFaresWithPrediction}},
built on
\href{https://github.com/yang-mills-group/OpenAI-YM-20261028/blob/3456e5e75cf8d4cd558955a5c7638fba2d68f464/YM/Integration/RG05Preconstruction/Expenses.lean\#L133-L144}{\texttt{joinRadarsShipWithPrediction}}
and
\href{https://github.com/yang-mills-group/OpenAI-YM-20261028/blob/3456e5e75cf8d4cd558955a5c7638fba2d68f464/YM/Integration/RG05Preconstruction/Expenses.lean\#L28-L119}{\texttt{joinExpensesShipWithPrediction}}.

\textbf{Proposition III.2.5 (the plentiful construction).} Assume the
hypotheses of III.2.4 and \(Q_i\ge2\). Then its conclusion holds with
the following additional clause. Consider any trajectory compatible with
the emitted factor, any reference levels \(t(d)\) with
\(d+3\le J_{t(d)}\), any section-transfer data and any gauge-door data
between consecutive levels. Assume that the resulting family satisfies
the family-consistency conditions, that its gauge doors satisfy the step
compatibility conditions, and that its section transfers satisfy the
prescribed bounds from level \(t(d)\) whenever \(t(d)\le N\). Then, for
every extraction order \(f\), the family satisfies the complete
weak-continuum requirement
(\href{https://github.com/yang-mills-group/OpenAI-YM-20261028/blob/3456e5e75cf8d4cd558955a5c7638fba2d68f464/YM/MergeProof/East/Gran.lean\#L35}{\texttt{CompleteWeak}},
III.11) and the trace gallery
(\href{https://github.com/yang-mills-group/OpenAI-YM-20261028/blob/3456e5e75cf8d4cd558955a5c7638fba2d68f464/YM/TraceExtension/Chart/Line.lean\#L41-L44}{\texttt{Gallery}},
VII.2). These are the hypotheses and additional conclusion of
\href{https://github.com/yang-mills-group/OpenAI-YM-20261028/blob/3456e5e75cf8d4cd558955a5c7638fba2d68f464/YM/FullVolume/Unlatch.lean\#L25-L34}{\texttt{Unlatched}}.
The floors, delays and waits may increase, and the prediction is
retained. \emph{Status:} theorem. This is where the requirements of the
trace test (Part VII) are bought on the same constants as the
ultraviolet estimates. \emph{Lean:}
\href{https://github.com/yang-mills-group/OpenAI-YM-20261028/blob/3456e5e75cf8d4cd558955a5c7638fba2d68f464/YM/Integration/RG05Preconstruction/Bardo.lean\#L26-L81}{\texttt{plentifulWithPrediction}},
\href{https://github.com/yang-mills-group/OpenAI-YM-20261028/blob/3456e5e75cf8d4cd558955a5c7638fba2d68f464/YM/Integration/RG05Preconstruction/Bardo.lean\#L84-L87}{\texttt{plentiful}}
(vacuous prediction).

\textbf{Proposition III.2.6 (a prediction on the same factor).} Let
\(G\) be simple, let \(s\) carry the estimate package of Lemma 4.16 for
fixed tariff and scope, and let \(\tau>0\). Let \(\pi\) be a prediction
that is \emph{prepaid} at the price map of \(s\), in the following
sense.

\begin{itemize}
\tightlist
\item
  For every choice of targets with the canonical chart, expenses with
  prescribed deficit \(-7/16\), programs, initial wages and baton, there
  is a length floor \(u\).
\item
  For every floor \(l\ge u\), every choice of block factors, wages,
  duration bounds, register, postage, scouts, protection, step records
  and post-\(L\) costs, and every delay \(V\), there is a delay
  \(W\ge V\). The step records must have \(\varepsilon\le1/4\) on
  transitions whose output stage is at least the first stage.
\item
  For every delay \(W'\ge W\), every barrier, wait, pricing record and
  directed datum, and every choice of the remaining factor data, \(\pi\)
  holds on the assembled factor if its shelter and step conditions hold
  and its periods have the divisibility property.
\end{itemize}

Then the complete native purchase of Lemma 4.17 can be carried out so
that, on level systems with divisible periods, \(\pi\) holds on the
factor it emits. The original floor is replaced by its maximum with
\(u\). The delay \(W\) is chosen before the original booth and final
delay. \emph{Status:} theorem. In the main-theorem closure it is used
only with the vacuous prediction
(\href{https://github.com/yang-mills-group/OpenAI-YM-20261028/blob/3456e5e75cf8d4cd558955a5c7638fba2d68f464/YM/Integration/NativeSelection/FundedPredictionSupply.lean\#L16}{\texttt{Always}}):
the only caller of III.2.7 passes it. \emph{Lean:}
\href{https://github.com/yang-mills-group/OpenAI-YM-20261028/blob/3456e5e75cf8d4cd558955a5c7638fba2d68f464/YM/Integration/NativeSelection/FundedCompanySelection.lean\#L49-L106}{\texttt{company}},
\href{https://github.com/yang-mills-group/OpenAI-YM-20261028/blob/3456e5e75cf8d4cd558955a5c7638fba2d68f464/YM/Integration/NativeSelection/FundedCompanySelection.lean\#L30-L42}{\texttt{Prepaid}},
\href{https://github.com/yang-mills-group/OpenAI-YM-20261028/blob/3456e5e75cf8d4cd558955a5c7638fba2d68f464/YM/Integration/NativeSelection/FundedCompanySelection.lean\#L23-L24}{\texttt{Prediction}},
\href{https://github.com/yang-mills-group/OpenAI-YM-20261028/blob/3456e5e75cf8d4cd558955a5c7638fba2d68f464/YM/Integration/NativeSelection/FundedPredictionSupply.lean\#L13-L14}{\texttt{AllPrices}}.

\textbf{Proposition III.2.7 (the complete run).}

\emph{Inputs.} Let \(G\) be simple and \(b>0\) (the kernel penalty). Let
the following records be given: \(\Omega\), \(\Delta\) and a custody
request \(\mathcal H\). Let \(\Lambda\) be a property of level systems
that is \emph{eventually paid} (III.2.8(a)). Let a \emph{demand} be
given:

\begin{itemize}
\tightlist
\item
  for every sheet \(s\) and price map \(q\), a nonempty set of choices
  \(c\);
\item
  for each choice \(c\), a floor \(e_c:\mathbb N\to\mathbb N\);
\item
  a real array \(F_c(n,i)\);
\item
  a real array
  \(F'_c[\ell](https://github.com/scottnarmstrong/yang-mills-informalization/blob/main/outline/units/n,i)\),
  which may depend on the block-factor sequence \(\ell\).
\end{itemize}

\emph{Conclusion.} There is a run \(\mathcal R\) (the record of the
whole construction, III.12) with a campaign such that:

\begin{itemize}
\tightlist
\item
  \begin{enumerate}
  \def\labelenumi{(\alph{enumi})}
  \tightlist
  \item
    the kernel penalty is \(b\);
  \end{enumerate}
\item
  (b) \emph{preselection:} there are \(s,q,c\) and a floor
  \(n\ge\max(f^s,e_c)\) such that:

  \begin{itemize}
  \tightlist
  \item
    \(\mathcal R\) uses the sheet \(s\) with floors \(n\) and the price
    map derived from \(q\);
  \item
    both arrays are paid in the sense of III.2.8(b) on the level system
    of \(\mathcal R\), with \(F'\) evaluated at the actual block factors
    \(\ell^{\mathcal R}\);
  \item
    \(F'_c[\ell^{\mathcal R}](https://github.com/scottnarmstrong/yang-mills-informalization/blob/main/outline/units/0,J_{m+1})\le\hat\beta_m\)
    for all \(m\);
  \end{itemize}
\item
  (c) \(k_j\ge j\) and \(q_j\ge j\) at every stage;
\item
  (d) \(\Lambda\) holds;
\item
  (e) the treasury of III.2.8(d) is funded;
\item
  (f) the common-period purchase is funded;
\item
  (g) the campaign has pillars, the funded, news and read-out records
  hold, and the custody and sky clauses hold.
\end{itemize}

There is also a version in which any prediction prepaid at every sheet
and price map holds on the factor of the run; the intermediate factor
theorem below supplies it. \emph{Status:} theorem. Its only caller in
the main closure is III.14
(\href{https://github.com/yang-mills-group/OpenAI-YM-20261028/blob/3456e5e75cf8d4cd558955a5c7638fba2d68f464/YM/Independent/ConstructionRun.lean\#L14-L28}{\texttt{constructionRun}}),
which passes:

\begin{itemize}
\tightlist
\item
  \(\Lambda\equiv\) true;
\item
  the demand of its native component recipe: a floor and the pre-\(L\)
  and post-\(L\) arrays;
\item
  concrete \(\Omega,\mathcal H,\Delta\);
\item
  the vacuous prediction.
\end{itemize}

\emph{Lean:}
\href{https://github.com/yang-mills-group/OpenAI-YM-20261028/blob/3456e5e75cf8d4cd558955a5c7638fba2d68f464/YM/Integration/NativePreselectedPostnumericPeriodRun.lean\#L99-L114}{\texttt{selectRunWithTreasuryPeriod\_preselected\_quote\_orders\_postnumeric}},
\href{https://github.com/yang-mills-group/OpenAI-YM-20261028/blob/3456e5e75cf8d4cd558955a5c7638fba2d68f464/YM/Integration/NativePreselectedPostnumericPeriodRun.lean\#L31-L46}{\texttt{PostnumericPreselection}},
\href{https://github.com/yang-mills-group/OpenAI-YM-20261028/blob/3456e5e75cf8d4cd558955a5c7638fba2d68f464/YM/Integration/NativePreselectedPostnumericPeriodClock.lean\#L44-L179}{\texttt{allFactorWithTreasuryPeriod\_preselected\_quote\_orders\_postnumeric\_of\_prediction}}.

\textbf{Proposition III.2.8 (purchases that persist under all later
choices).}

\emph{(a) Eventually paid properties.} A property \(U\) of level systems
is \emph{eventually paid} if the following holds. For every sheet, block
factors, smallness exponent and delay \(V\) there is a delay \(X\ge V\)
such that every level system meeting the entrance conditions of any
delay \(W\ge X\) (with any barrier and wait) satisfies \(U\). A
conjunction of two eventually paid properties is eventually paid.

\emph{(b) Diagonal payment of constants.} For every real array
\(F(n,i)\), the following property is eventually paid. At every level
\(m\), \(1\le1+\log\log\hat\beta_m\). For every \(n\) there is \(m_n\)
such that
\[|F(n,i)|\le1+\log\log\hat\beta_m\qquad\text{for all } m\ge m_n \text{ and all } i\le J_{m+1}.\]
Likewise, for every \(F:\mathbb N\to\mathbb R\), the property
``\(F(J_{m+1})\le\hat\beta_m\) for all \(m\)'' is eventually paid.

\emph{(c) Countably many step requests.} Fix a countable index set with
an encoding in \(\mathbb N\), and let \((w_\alpha)\) be a family of
requests indexed by it. Each request consists of a floor
\(\mathrm{pre}_\alpha:\mathbb N\to\mathbb N\) and a cost function of a
stage transition \((j,i)\), its block geometry and a stage record; at
step \(m\) the record is \(X_{J_m}\). Then there is a single floor
\(\mathrm{pre}\), depending only on the family, with the following
property. If \(\ell_j\ge\mathrm{pre}(j)\) for all \(j\), then every
delay can be raised so that every later level system satisfying the
entrance conditions for the enlarged delay, with any barrier and wait,
satisfies, for all \(\alpha\) and \(n\) with
\(\mathrm{code}(\alpha)\le J_n\) and all \(m\ge n\):

\begin{itemize}
\tightlist
\item
  \(\mathrm{pre}_\alpha(J_{m+1})\le L_m\);
\item
  the cost of step \(m\) is at most \(1+\log(2+\log\hat\beta_m)\).
\end{itemize}

\emph{(d) The treasury.} For simple \(G\) and every kernel choice there
is a \emph{plan}. A plan fixes a naming of charges and, for every
request, a dictionary of finite words and spectra. If the block factors
exceed the plan's floor, then every delay can be raised so that every
later level system satisfying its entrance conditions is \emph{treasury
funded}: it meets the step requests of the plan (as in (c)) and the
inverse charge from level 0 on.

\emph{Status:} theorems. Dictionaries are chosen once, before the block
factors, and never reselected. \emph{Lean:}
\href{https://github.com/yang-mills-group/OpenAI-YM-20261028/blob/3456e5e75cf8d4cd558955a5c7638fba2d68f464/YM/R2I/Campaign/Clock.lean\#L36-L38}{\texttt{EventuallyPaid}},
\href{https://github.com/yang-mills-group/OpenAI-YM-20261028/blob/3456e5e75cf8d4cd558955a5c7638fba2d68f464/YM/Integration/NativePreselectedPostnumericPeriodClock.lean\#L31-L39}{\texttt{postnumeric\_quote\_diagonal\_pair\_requests}},
\href{https://github.com/yang-mills-group/OpenAI-YM-20261028/blob/3456e5e75cf8d4cd558955a5c7638fba2d68f464/YM/Integration/NativeDiagonalClockRequest.lean\#L52-L79}{\texttt{eventuallyPaid}},
\href{https://github.com/yang-mills-group/OpenAI-YM-20261028/blob/3456e5e75cf8d4cd558955a5c7638fba2d68f464/YM/Integration/NativeDiagonalClockRequest.lean\#L45-L48}{\texttt{PaidRequests}},
\href{https://github.com/yang-mills-group/OpenAI-YM-20261028/blob/3456e5e75cf8d4cd558955a5c7638fba2d68f464/YM/Integration/NativePostLengthPayment.lean\#L19-L20}{\texttt{StagePaid}},
\href{https://github.com/yang-mills-group/OpenAI-YM-20261028/blob/3456e5e75cf8d4cd558955a5c7638fba2d68f464/YM/RG08/Calendar/Base.lean\#L68-L77}{\texttt{minDelay}},
\href{https://github.com/yang-mills-group/OpenAI-YM-20261028/blob/3456e5e75cf8d4cd558955a5c7638fba2d68f464/YM/RG08/Calendar/Base.lean\#L33-L35}{\texttt{Need}},
\href{https://github.com/yang-mills-group/OpenAI-YM-20261028/blob/3456e5e75cf8d4cd558955a5c7638fba2d68f464/YM/RG08/Calendar/Base.lean\#L43-L45}{\texttt{Affords}},
\href{https://github.com/yang-mills-group/OpenAI-YM-20261028/blob/3456e5e75cf8d4cd558955a5c7638fba2d68f464/YM/Integration/CalendarTreasuryPurchase.lean\#L18-L20}{\texttt{Plan}},
\href{https://github.com/yang-mills-group/OpenAI-YM-20261028/blob/3456e5e75cf8d4cd558955a5c7638fba2d68f464/YM/Integration/CalendarTreasuryPurchase.lean\#L28-L31}{\texttt{plan\_exists}},
\href{https://github.com/yang-mills-group/OpenAI-YM-20261028/blob/3456e5e75cf8d4cd558955a5c7638fba2d68f464/YM/Integration/CalendarTreasuryPurchase.lean\#L42-L64}{\texttt{buy}},
\href{https://github.com/yang-mills-group/OpenAI-YM-20261028/blob/3456e5e75cf8d4cd558955a5c7638fba2d68f464/YM/Integration/CalendarTreasuryPurchase.lean\#L35-L38}{\texttt{TreasuryFunded}}.

\textbf{Proposition III.2.9 (analytic estimates used by the purchases).}
The row and cross norms are componentwise exponentially weighted
\(\ell^1\) bounds on displacement kernels. Multilinear kernels use the
corresponding sum over displacements and colour indices. The colour
spaces are fixed and finite; constants may depend on them. A cross maps
coarse fields to fine fields at block size \(L\). All vanishing orders
refer to the Fourier symbol at zero momentum.

\begin{itemize}
\item
  \emph{Pullback gain.} For \(\sigma,\tau>0\), \(C\in\mathbb R\) and
  \(d\in\mathbb N\), there are \(0<\eta\le\sigma\) and \(K\ge0\) such
  that, for every positive odd integer \(L\) with \(\tau\le L\eta\), the
  following holds. If \(\|A\|_\tau\le C/L\),
  \(\|\nabla^sB\|_\tau\le C/L^{1+|s|}\) for every list of coordinate
  differences with \(|s|\le d\), and \(R\) has \(\|R\|_\sigma\le M\)
  with \(M\ge0\) and vanishes to order \(d\), then
  \(\|A^{\mathsf T}RB\|_\tau\le KML^{2-d}\).
  \href{https://github.com/yang-mills-group/OpenAI-YM-20261028/blob/3456e5e75cf8d4cd558955a5c7638fba2d68f464/YM/RG08/Flat/Pullback.lean\#L56-L110}{\texttt{pullback\_gain}}
\item
  \emph{Contraction of Hessian differences.} For \(\sigma,\tau>0\) and
  \(C\in\mathbb R\), there is \(N\in\mathbb N\) such that, for every
  positive odd \(L\ge N\), crosses \(A,B\) obeying the preceding bounds
  through order \(3\), and a row \(R\) on one-form components with
  \(M\ge0\), \(\|R\|_\sigma\le M\) and vanishing order \(3\), one has
  \(\|A^{\mathsf T}RB\|_\tau\le M L^{-1/2}\).
  \href{https://github.com/yang-mills-group/OpenAI-YM-20261028/blob/3456e5e75cf8d4cd558955a5c7638fba2d68f464/YM/RG08/Flat/KStep.lean\#L130-L146}{\texttt{large\_contraction}}
\item
  \emph{Soft transfer of multilinear kernels.} For \(r,d\in\mathbb N\),
  \(\sigma,w>0\) and \(C\ge0\), there are \(0<\eta\le\sigma\) and
  \(K\ge0\), fixed before \(L\), with the following property. Let \(F\)
  be a real \((r+1)\)-linear coefficient kernel with width \(\sigma\)
  and bound \(B\), supported on tuples whose first site is the origin,
  and vanishing to order \(d\) in the independent momenta. Let \(L\) be
  positive and odd with \((r+2)w\le L\eta\), and let \(Q\) satisfy
  \(\|\nabla^iQ\|_{(r+2)w}\le C/L^{1+|i|}\) for every list \(|i|\le d\).
  The transfer through \(r+1\) copies of \(Q\), followed by
  symmetrization and recentering, has width \(w\) and bound
  \(KBL^4/L^{r+1+d}\).
  \href{https://github.com/yang-mills-group/OpenAI-YM-20261028/blob/3456e5e75cf8d4cd558955a5c7638fba2d68f464/YM/RG08/Ballast/Slide.lean\#L44-L90}{\texttt{flow\_soft}}
\item
  \emph{Division by differences.} Fix a finite coordinate set \(I\),
  \(n\in\mathbb N\) and \(\sigma>0\). There are \(0<\eta\le\sigma\) and
  \(K\ge0\) such that, for every \(A\ge0\) and real kernel
  \(f:\mathbb Z^I\to\mathbb R\) with
  \(\sum_xe^{\sigma|x|_1}|f(x)|\le A\) whose symbol vanishes to order
  \(n\), there are real kernels \(g_b\), indexed by
  \(b:\{1,\ldots,n\}\to I\), with \(\sum_xe^{\eta|x|_1}|g_b(x)|\le KA\)
  and
  \[\widehat f(p)=\sum_b\left(\prod_{a=1}^{n}(e^{ip_{b(a)}}-1)\right)\widehat g_b(p)\]
  for every real momentum \(p\).
  \href{https://github.com/yang-mills-group/OpenAI-YM-20261028/blob/3456e5e75cf8d4cd558955a5c7638fba2d68f464/YM/RG08/Flat/Differences.lean\#L46-L92}{\texttt{real\_difference\_division}}
\item
  \emph{Odd-part cancellation.} Fix a sequence of step scenes whose
  input coupling tends to infinity, leaves \(M_i\), a cutoff parameter
  \(\chi\), chart radii satisfying the logarithm-chart conditions, and
  fluctuation data \(z_i\). Let
  \(H_i=\frac12\log\beta_i^{\mathrm{out}}\) and \(\eta_{*,i}=e^{-H_i}\).
  Assume that the fluctuation bounds, the leaf weights and the
  third-derivative bounds \(D_i\) have upper exponential rate \(0\)
  relative to \(H_i\); that eventually each leaf is a logarithmic term
  given by a masked regular activity whose expression in logarithmic
  coordinates, multiplied by the square of its threshold, has zero first
  derivative at \(0\) and derivatives through order \(3\) bounded by
  \(D_i\) on the unit ball; that both \(z_i\) and \(-z_i\) lie in the
  healthy domain; and that both cutoffs equal \(1\) near the origin in
  the two deformation parameters. If \(P_i(z_i;y)\) denotes the paved
  expansion, then, for every \(\varepsilon>0\), eventually
  \[\big\|D_y^2\big(P_i(z_i;y)+P_i(-z_i;y)\big)\big|_{y=0}\big\|\le\eta_{*,i}^{\,2-\varepsilon}.\]
  Here an upper rate \(a\) means an eventual bound \(e^{rH_i}\) for
  every \(r>a\).
  \href{https://github.com/yang-mills-group/OpenAI-YM-20261028/blob/3456e5e75cf8d4cd558955a5c7638fba2d68f464/YM/RG12/Shadows/Shape.lean\#L54-L70}{\texttt{mirror\_price}}
\item
  \emph{Uniqueness of local read-out.} Fix a representation \(R\), a
  degree-\(d\) frame and a section readable at the origin for the budget
  of \(R\). Let \(P,Q\) be coefficient presentations in the prescribed
  format for the same \(R\): their regular terms have supports
  containing the origin, vanish at the identity and have zero local
  degree-\(d\) projection on nonwrapping supports; their nonzero forced
  terms have supports containing the origin and satisfy the eligibility,
  support, gauge-invariance and measurability conditions. Suppose they
  have equal densities for this section and remainder bounds \(a,b\)
  with the same positive forced-term parameter. For every
  \(m\in\mathbb N\) with \(10(m(4B+2)+2B)<n\), their coefficient vectors
  satisfy \[\|x_P-x_Q\|_{\mathrm{coeff}}\le K_{F,H}e^{-a_*m}(a+b),\]
  where \(n\) is the torus side, \(B\) and \(a_*\) are the cover size
  and support exponent of \(R\), and \(K_{F,H}\) depends on the frame
  and holonomy tests.
  \href{https://github.com/yang-mills-group/OpenAI-YM-20261028/blob/3456e5e75cf8d4cd558955a5c7638fba2d68f464/YM/TraceExtension/Quince/Unum.lean\#L11-L42}{\texttt{unique\_part}}
\end{itemize}

Further linked estimates concern a post-block Hessian bound
(\href{https://github.com/yang-mills-group/OpenAI-YM-20261028/blob/3456e5e75cf8d4cd558955a5c7638fba2d68f464/YM/RG12/Spins/Muse.lean\#L35-L59}{\texttt{design}}),
extraction of cubic and quartic local coefficients
(\href{https://github.com/yang-mills-group/OpenAI-YM-20261028/blob/3456e5e75cf8d4cd558955a5c7638fba2d68f464/YM/FullVolume/Canvas/Reservoir.lean\#L84-L118}{\texttt{tanker}}),
and asymptotic bounds for the trace insertion
(\href{https://github.com/yang-mills-group/OpenAI-YM-20261028/blob/3456e5e75cf8d4cd558955a5c7638fba2d68f464/YM/TraceExtension/Zn/Arc.lean\#L16-L33}{\texttt{cloche}}).
\emph{Status:} theorems, with the hypotheses above. The detailed scene,
coefficient-format and norm interfaces are listed among the open points
in section 7.

\subsubsection{2. What it uses}\label{what-it-uses-4}

\textbf{From other units.}

\begin{itemize}
\tightlist
\item
  I.1: \(\kappa>0\) for simple \(G\)
  (\href{https://github.com/yang-mills-group/OpenAI-YM-20261028/blob/3456e5e75cf8d4cd558955a5c7638fba2d68f464/YM/Gauge/Casimir.lean\#L174-L176}{\texttt{slopeConstant\_pos}}).
\item
  III.1 (interfaces):

  \begin{itemize}
  \tightlist
  \item
    the scale sequence
    (\href{https://github.com/yang-mills-group/OpenAI-YM-20261028/blob/3456e5e75cf8d4cd558955a5c7638fba2d68f464/YM/Lattice/Scale.lean\#L28-L34}{\texttt{Schedule}}),
    the calendar, its conditions and level systems
    (\href{https://github.com/yang-mills-group/OpenAI-YM-20261028/blob/3456e5e75cf8d4cd558955a5c7638fba2d68f464/YM/RG/History.lean\#L64-L73}{\texttt{Calendar}},
    \href{https://github.com/yang-mills-group/OpenAI-YM-20261028/blob/3456e5e75cf8d4cd558955a5c7638fba2d68f464/YM/RG/History.lean\#L97-L118}{\texttt{Fits}},
    \href{https://github.com/yang-mills-group/OpenAI-YM-20261028/blob/3456e5e75cf8d4cd558955a5c7638fba2d68f464/YM/RG/History.lean\#L131-L134}{\texttt{Levels}});
  \item
    histories, transfer and selection
    (\href{https://github.com/yang-mills-group/OpenAI-YM-20261028/blob/3456e5e75cf8d4cd558955a5c7638fba2d68f464/YM/RG/History.lean\#L204-L206}{\texttt{History}},
    \href{https://github.com/yang-mills-group/OpenAI-YM-20261028/blob/3456e5e75cf8d4cd558955a5c7638fba2d68f464/YM/RG/History.lean\#L223-L227}{\texttt{Transfer}},
    \href{https://github.com/yang-mills-group/OpenAI-YM-20261028/blob/3456e5e75cf8d4cd558955a5c7638fba2d68f464/YM/RG/History.lean\#L234-L240}{\texttt{Selected}});
  \item
    block control from the kernel's branches
    (\href{https://github.com/yang-mills-group/OpenAI-YM-20261028/blob/3456e5e75cf8d4cd558955a5c7638fba2d68f464/YM/Branch/Covariant.lean\#L70-L72}{\texttt{blocks\_go}});
  \item
    the directed datum
    (\href{https://github.com/yang-mills-group/OpenAI-YM-20261028/blob/3456e5e75cf8d4cd558955a5c7638fba2d68f464/YM/PCalc/Band.lean\#L26-L36}{\texttt{Directed}}).
  \end{itemize}
\item
  III.9: the reference couplings, threshold and band (III.9.2); the
  meaning of a selected history (III.9.3); the Wilson start-up.
\item
  III.10:

  \begin{itemize}
  \tightlist
  \item
    the verification of the shooting hypotheses
    (\href{https://github.com/yang-mills-group/OpenAI-YM-20261028/blob/3456e5e75cf8d4cd558955a5c7638fba2d68f464/YM/PhysicalNext/RG_NEXT_COMPLETE_REPEATABLE_NONLINEAR_STEP/Graces/Eagles.lean\#L33-L50}{\texttt{Graces.winnow}});
  \item
    the divisible-period schedule
    (\href{https://github.com/yang-mills-group/OpenAI-YM-20261028/blob/3456e5e75cf8d4cd558955a5c7638fba2d68f464/YM/PhysicalNext/RG_NEXT_COMPLETE_REPEATABLE_NONLINEAR_STEP/Pane/Divisors.lean\#L26-L135}{\texttt{Wait.engine}});
  \item
    the emitted factor and bandage records.
  \end{itemize}
\item
  III.9, III.11, III.12, III.15, VII.2: the predicates bought in
  III.2.4--III.2.5 (table in III.2.4) and their suppliers. This unit
  proves only that they hold on one common choice of constants. Their
  analytic content belongs to those units.
\item
  III.12: the run record
  (\href{https://github.com/yang-mills-group/OpenAI-YM-20261028/blob/3456e5e75cf8d4cd558955a5c7638fba2d68f464/YM/FullVolume/Cross.lean\#L30-L48}{\texttt{Run}}).
  III.14: the demand passed to III.2.7.
\item
  III.3, III.6 (probable): the free-kernel and stationary-branch
  estimates from which the sheet is built. The two entry points,
  \href{https://github.com/yang-mills-group/OpenAI-YM-20261028/blob/3456e5e75cf8d4cd558955a5c7638fba2d68f464/YM/FullVolume/Train/Wheels.lean\#L77-L103}{\texttt{reapsY}}
  and
  \href{https://github.com/yang-mills-group/OpenAI-YM-20261028/blob/3456e5e75cf8d4cd558955a5c7638fba2d68f464/YM/FullVolume/Canvas/Mails.lean\#L55-L94}{\texttt{verdant}},
  are themselves filed in this unit.
\item
  III.16: strip and Laurent norms, Taylor jets (for III.2.9).
\end{itemize}

\textbf{Standard mathematics.}

\begin{itemize}
\tightlist
\item
  The intermediate value theorem, in a refined form (Lemma 4.11).
\item
  Elementary asymptotics: \((\log x)^2=o(x)\); polynomials in \(\log s\)
  are \(o(e^s)\).
\item
  Diagonal enumeration of countably many requests.
\item
  For background on the order of choice and bare-coupling tuning in
  constructive RG:

  \begin{itemize}
  \tightlist
  \item
    Choosing constants in a fixed order: first the constants of the
    step, then the block factor \(L\) large depending on them, then the
    coupling large depending on \(L\). See Bałaban's papers on lattice
    gauge theories (Comm. Math. Phys. 1983--1989) and
    Bauerschmidt--Brydges--Slade, \emph{Introduction to a
    renormalisation group method} (LNM 2242, 2019).
  \item
    Tuning the bare coupling as a function of the cutoff so that the
    effective coupling at unit scale is fixed. Compare
    Gawędzki--Kupiainen on the Gross--Neveu and \(\phi^4_4\) models
    (Comm. Math. Phys. 99, 102, 1985).
  \end{itemize}

  Here the tuning is a nested-interval argument on continuous families
  (§4.3).
\end{itemize}

\textbf{Imported constants and hypotheses.}

\begin{itemize}
\tightlist
\item
  \(\kappa\) enters only through \(\kappa>0\) and through the form
  \(\hat\beta_m=\bar\beta+\kappa s_m\). The shift of \(\beta\) by
  \(\kappa\log L\) up to error \(\omega\) is part of the (advance)
  hypothesis of III.2.3, supplied in III.10.
\item
  The hypotheses (start) and (advance) of III.2.3.
\item
  In III.2.7, the records \(\Omega\), \(\mathcal H\), \(\Delta\) and the
  demand come from III.14. Inside the assembly, the kernel with penalty
  \(b\) comes from III.1
  (\href{https://github.com/yang-mills-group/OpenAI-YM-20261028/blob/3456e5e75cf8d4cd558955a5c7638fba2d68f464/YM/Branch/Pick.lean\#L137-L142}{\texttt{Axial.viable}}).
  The tableau and sphere data are constructed in the assembly.
\end{itemize}

\subsubsection{3. Strategy}\label{strategy-4}

Choose the stage records and floors first, then the block factors
\(\ell_j\), post-\(L\) costs and smallness exponents, delays, the
barrier \(\bar\beta\), and finally the wait. Increasing the appropriate
threshold meets each requirement: the block factors dominate the stage
constants, entry couplings pay later costs, and the wait provides
comparison windows and large tori. Persistent requirements combine by
taking maxima of floors and chaining delay and wait increases. This
gives common constants for III.2.4--III.2.7.

The stage index \(J_m\) advances by at most one per level. It enters a
stage only after its entry thresholds are met and stays there long
enough for one prescribed comparison window. At each fixed stage the
volume bill is bounded by a polynomial in \(\log s\), hence eventually
by \(e^s\). Since \(s_m\ge m\log3\), this pays every later volume
demand. One period offset pays the finitely many steps of the first
plateau.

For fixed \(N,p\), start with the Wilson family over
\([\bar\beta+(\kappa/2)s_N,\bar\beta+2\kappa s_N]\subset\mathfrak B_N\).
At each step the drift bound sends the endpoint couplings strictly
beyond the endpoints of the next window. A refined intermediate value
argument selects a smaller parameter interval on which the new family
stays in that window, so the next step applies. At level \(0\) the
window is \(\{\bar\beta\}\). Summing the drifts upward from there and
using \(\omega_{J_{m+1}}\le\log L_m/J_{m+1}^2\) gives the running law
uniformly in depth and volume.

\subsubsection{4. Steps}\label{steps-4}

\paragraph{4.1 Numerical existence
lemmas}\label{numerical-existence-lemmas}

In sections 4.1--4.2, \(G\) is simple and the data are those of III.2.1.

\textbf{Lemma 4.1 (large block factors).} For each stage \(i\) there is
\(M_i\) such that every odd \(n\ge M_i\) satisfies the conditions of
III.2.1(2) at stage \(i\). Hence the stage block factors
\(\ell_i=2\max_{j\le i}M_j+1\) exist. \emph{Sketch:} below the first
stage only oddness and the size floors are required. At later stages,
the remaining conditions have the form \(c/\log n\le\delta\) with
\(\delta\in\{\kappa/4,i^{-2}\}\) and \(\delta>0\). Each holds for large
\(n\) because \(\log n\to\infty\). The running maximum makes \(\ell\)
nondecreasing. \emph{Lean:}
\href{https://github.com/yang-mills-group/OpenAI-YM-20261028/blob/3456e5e75cf8d4cd558955a5c7638fba2d68f464/YM/Price/Intent.lean\#L55-L77}{\texttt{Rolls.eventually\_large}},
\href{https://github.com/yang-mills-group/OpenAI-YM-20261028/blob/3456e5e75cf8d4cd558955a5c7638fba2d68f464/YM/Price/Intent.lean\#L79-L84}{\texttt{Rolls.exists\_rolls}}.

\textbf{Lemma 4.2 (smallness exponents).} For every positive sequence
\(\alpha\), the running minimum
\(\theta_j=\min_{i\le j}\min\big((i+1)^{-1},\alpha_i\big)\) is positive,
nonincreasing, and satisfies \(\theta_j\le\min(\alpha_j,1/(j+1))\).
\emph{Lean:}
\href{https://github.com/yang-mills-group/OpenAI-YM-20261028/blob/3456e5e75cf8d4cd558955a5c7638fba2d68f464/YM/Price/Intent.lean\#L98-L111}{\texttt{Protection.exists\_margin}}.

\textbf{Lemma 4.3 (barrier).} For every \(b\), take
\(\bar\beta=\max(2,x_*,b,\beta^{\rm ent}_{j_0})\), where
\(x^{-1/2}e^{(\log\log x)^2}\le1\) for all \(x\ge x_*\). This is a
barrier. \emph{Sketch:} \((\log y)^2\le\frac12 y\) for large \(y\),
applied with \(y=\log x\). \emph{Lean:}
\href{https://github.com/yang-mills-group/OpenAI-YM-20261028/blob/3456e5e75cf8d4cd558955a5c7638fba2d68f464/YM/Price/Length.lean\#L19-L37}{\texttt{Bar.exists\_above}}.

\textbf{Lemma 4.4 (the volume bill).} Define
\[\mathrm{bill}_j(s)=S_j(1+c+c')^{g_j}+100(c+1)+100\big((j+3)\ell_j+1\big),\]
where \(c=\lceil\log^2(\bar\beta+\kappa s)\rceil\) and
\(c'=\lceil\log^2(\bar\beta+\kappa(s+\log\ell_j))\rceil\). For each
fixed \(j\), \(\mathrm{bill}_j(s)\le e^s\) for all sufficiently large
\(s\). At \(s=s_m\) with \(j=J_{m+1}\) we have \(c=B_m\), \(c'=B_{m+1}\)
and \((j+3)\ell_j=R_m\). So the bill is the sum of the three volume
demands of III.2.2(c),(e). \emph{Sketch:} \(c,c'=O((\log s)^2)\), so the
bill is bounded by a polynomial in \(\log s\). \emph{Lean:}
\href{https://github.com/yang-mills-group/OpenAI-YM-20261028/blob/3456e5e75cf8d4cd558955a5c7638fba2d68f464/YM/Price/Length.lean\#L43-L47}{\texttt{Bar.bill}},
\href{https://github.com/yang-mills-group/OpenAI-YM-20261028/blob/3456e5e75cf8d4cd558955a5c7638fba2d68f464/YM/Price/Length.lean\#L64-L84}{\texttt{Bar.allocated}}.

\textbf{Lemma 4.5 (entering a stage).} For each \(j>j_0\) there is
\(n_j\) with the following property. For every level \(m\ge n_j\) and
every \(s\ge m\log3\):

\begin{itemize}
\tightlist
\item
  \(\beta^{\rm ent}_j\le\bar\beta+\kappa s\);
\item
  \(\sigma^{\rm ent}_j(\bar\beta)\le s\);
\item
  \(s>0\);
\item
  \(\log\ell_j\le s/j\);
\item
  \(\mathrm{bill}_j(s)\le e^s\).
\end{itemize}

\emph{Sketch:} each condition holds for all large \(s\). Take the
largest of the thresholds, and \(n_j\) with \(n_j\log3\) above it.
\emph{Lean:}
\href{https://github.com/yang-mills-group/OpenAI-YM-20261028/blob/3456e5e75cf8d4cd558955a5c7638fba2d68f464/YM/Price/Length.lean\#L90-L106}{\texttt{Bar.enters}}.

\textbf{Lemma 4.6 (windows).} For each \(j\) there is \(c_j\) such that,
for all \(k\ge c_j\) and all \(s\) with \(k\log3\le s\le k\log\ell_j\),
we have \(\sigma^{\rm rdy}_j\le s\) and \(h_j(s)\le k/2\).
\emph{Sketch:} for large \(s\), \(h_j(s)\le s/(2\log\ell_j)\) by
sublinearity, and \(s\le k\log\ell_j\). \emph{Lean:}
\href{https://github.com/yang-mills-group/OpenAI-YM-20261028/blob/3456e5e75cf8d4cd558955a5c7638fba2d68f464/YM/Price/Length.lean\#L109-L127}{\texttt{Wait.lead}}.

\paragraph{4.2 The level schedule}\label{the-level-schedule}

\textbf{Lemma 4.7 (exhaustion).} For any \(b,c:\mathbb N\to\mathbb N\)
there is \(J':\mathbb N\to\mathbb N\) with the following properties:

\begin{itemize}
\tightlist
\item
  \(J'_0=J'_1=0\), \(J'\) is nondecreasing, \(J'_{i+1}\le J'_i+1\) and
  \(J'\to\infty\);
\item
  \(b(J'_{i+1})\le i\) whenever \(J'_{i+1}>0\);
\item
  for every \(i\) there is \(k\in\mathbb N\), \(k\ge c(i)\) with
  \(J'\equiv i\) on \([k-\lfloor k/2\rfloor,\,k+\lfloor k/2\rfloor]\).
\end{itemize}

\emph{Sketch:} set \(a_0=0\) and \(a_{i+1}=a_i+b(i+1)+2k_i+1\), where
\(k_i=c(i)+2a_i+1\). Put \(J'=i\) on \([a_i,a_{i+1})\). Stage \(i+1\)
starts at level \(a_{i+1}>b(i+1)\). The window around \(k_i\) lies in
\([a_i,a_{i+1})\) because \(k_i/2\ge a_i\) and \(\tfrac32k_i<a_{i+1}\).
\emph{Lean:}
\href{https://github.com/yang-mills-group/OpenAI-YM-20261028/blob/3456e5e75cf8d4cd558955a5c7638fba2d68f464/YM/Price/Waits.lean\#L41-L84}{\texttt{Program.exhaust}}.

\textbf{Lemma 4.8 (the schedule meets all entrance conditions; proof of
III.2.2).}

\emph{Construction.}

\begin{enumerate}
\def\labelenumi{\arabic{enumi}.}
\tightlist
\item
  Apply Lemma 4.7 with \(b(0)=0\) and \(b(i)=n_{j_0+i}\) for \(i>0\)
  (Lemma 4.5) and \(c(i)=c_{j_0+i}\) (Lemma 4.6), and set
  \(J_m=j_0+J'_m\), \(L_m=\ell_{J_{m+1}}\) and
  \(s_m=\sum_{i<m}\log L_i\). Then \(m\log3\le s_m\le m\log\ell_{J_m}\),
  using that \(\ell\) and \(J\) are monotone.
\item
  Let \(m_*\) be such that \(J'_m\ge1\) for \(m\ge m_*\). Put
  \(M=n_{\min}+\max_{i\le m_*}\mathrm{bill}_{j_0}(s_i)\) and
  \(n_p=4(M+p+1)\).
\end{enumerate}

\emph{Verification.}

\begin{itemize}
\tightlist
\item
  \emph{Entry.} If \(J_{m+1}=j_0\), then
  \(\beta^{\rm ent}_{j_0}\le\bar\beta\le\hat\beta_m\). Otherwise
  \(m\ge n_{J_{m+1}}\) and \(s_m\ge m\log3\), so Lemma 4.5 applies.
\item
  \emph{Volume off the first plateau.} There
  \(\mathrm{bill}_{J_{m+1}}(s_m)\le e^{s_m}\le n_pe^{s_m}\). This gives
  (c), the cover condition, and the patch condition, the latter because
  sides increase with \(m\).
\item
  \emph{Volume on the first plateau.} There
  \(\mathrm{bill}_{j_0}(s_m)\le M<n_p\).
\item
  \emph{Mesh.} For \(m\ge m_*\), Lemma 4.5 gives
  \(\log L_m/s_m\le1/J_{m+1}\to0\).
\item
  \emph{Windows.} At the level \(k\) given by Lemma 4.7 for stage
  \(j=j_0+i\), \(s_k\) lies in the range of Lemma 4.6. So
  \(h_j(s_k)\le k/2\), and \(J\equiv j\) on \([k-h_j(s_k),k+h_j(s_k)]\).
\item
  \emph{Remaining conditions.} Monotonicity of
  \(k_{J_m},q_{J_m},a_{J_m}\) comes from the sheet. The drift and
  log-cost bounds come from Lemma 4.1, and the bounds on \(\theta\) from
  Lemma 4.2. Block control comes from the kernel's branches.
\item
  \emph{Divisibility variant.} Given \(c\ge1\) and \(z\), the index
  \(p=c(M+z+1)-(M+1)\ge z\) gives \(n_p=4c(M+z+1)\).
\end{itemize}

\emph{Lean:}
\href{https://github.com/yang-mills-group/OpenAI-YM-20261028/blob/3456e5e75cf8d4cd558955a5c7638fba2d68f464/YM/Price/Waits.lean\#L137-L239}{\texttt{Wait.ready}},
\href{https://github.com/yang-mills-group/OpenAI-YM-20261028/blob/3456e5e75cf8d4cd558955a5c7638fba2d68f464/YM/Price/Waits.lean\#L89-L104}{\texttt{Program.book}},
\href{https://github.com/yang-mills-group/OpenAI-YM-20261028/blob/3456e5e75cf8d4cd558955a5c7638fba2d68f464/YM/Price/Waits.lean\#L115-L121}{\texttt{scaled\_lower}},
\href{https://github.com/yang-mills-group/OpenAI-YM-20261028/blob/3456e5e75cf8d4cd558955a5c7638fba2d68f464/YM/Price/Waits.lean\#L123-L130}{\texttt{scaled\_upper}},
\href{https://github.com/yang-mills-group/OpenAI-YM-20261028/blob/3456e5e75cf8d4cd558955a5c7638fba2d68f464/YM/PhysicalNext/RG_NEXT_COMPLETE_REPEATABLE_NONLINEAR_STEP/Pane/Divisors.lean\#L26-L135}{\texttt{Wait.engine}}.

\paragraph{4.3 Shooting}\label{shooting}

\textbf{Lemma 4.9 (running law from the transfer inequality).} Consider
a history on a level system of III.2.2 with \(\beta_{N,p,0}=\bar\beta\)
and
\(|\beta_{N,p,i+1}-\beta_{N,p,i}-\kappa\log L_i|\le\omega_{J_{i+1}}\)
for all \(i<N\). Then it obeys the running law of III.2.3.
\emph{Sketch:}

\begin{enumerate}
\def\labelenumi{\arabic{enumi}.}
\tightlist
\item
  Since \(\hat\beta_{i+1}-\hat\beta_i=\kappa\log L_i\) and
  \(\beta_{N,p,0}=\hat\beta_0\), induction upward from level 0 gives
  \(|\beta_{N,p,i}-\hat\beta_i|\le D_i:=\sum_{k=1}^{i}\omega_{J_k}\).
\item
  Given \(\varepsilon>0\), choose \(i_0\) with
  \(J_{i+1}^{-2}\le\varepsilon/2\) for \(i\ge i_0\). Then
  \(\omega_{J_{i+1}}\le\frac\varepsilon2\log L_i\), so
  \(D_i\le D_{i_0}+\frac\varepsilon2(s_i-s_{i_0})\).
\item
  Finally \(D_{i_0}\le\frac\varepsilon2 s_i\) for large \(i\), because
  \(s_i\to\infty\).
\end{enumerate}

The bound is uniform in \(N\ge i\) and \(p\). \emph{Lean:}
\href{https://github.com/yang-mills-group/OpenAI-YM-20261028/blob/3456e5e75cf8d4cd558955a5c7638fba2d68f464/YM/Shoot/Race.lean\#L23-L53}{\texttt{trajectory}}.

\textbf{Lemma 4.10 (gates).} Put \(\beta^-_m=\bar\beta+\frac\kappa2s_m\)
and \(\beta^+_m=\bar\beta+2\kappa s_m\). Then
\([\beta^-_m,\beta^+_m]\subset\mathfrak B_m\), because \(\bar\beta>0\).
Suppose a level-\(m\) representation \(\rho\) relays a level-\((m+1)\)
representation \(\rho'\), and
\(\omega_{J_{m+1}}\le\frac\kappa4\log L_m\). Then:

\begin{itemize}
\tightlist
\item
  if \(\beta(\rho')=\beta^-_{m+1}\), then
  \(\beta(\rho)\le\beta^-_m-\frac\kappa4\log L_m<\beta^-_m\);
\item
  if \(\beta(\rho')=\beta^+_{m+1}\), then
  \(\beta(\rho)\ge\beta^+_m+\frac{3\kappa}4\log L_m>\beta^+_m\).
\end{itemize}

The drift hypothesis follows from the entrance conditions and
III.2.1(2), since \(L_m=\ell_{J_{m+1}}\). \emph{Lean:}
\href{https://github.com/yang-mills-group/OpenAI-YM-20261028/blob/3456e5e75cf8d4cd558955a5c7638fba2d68f464/YM/Shoot/Reach.lean\#L57}{\texttt{floorBeta}},
\href{https://github.com/yang-mills-group/OpenAI-YM-20261028/blob/3456e5e75cf8d4cd558955a5c7638fba2d68f464/YM/Shoot/Reach.lean\#L58}{\texttt{ceilBeta}},
\href{https://github.com/yang-mills-group/OpenAI-YM-20261028/blob/3456e5e75cf8d4cd558955a5c7638fba2d68f464/YM/Shoot/Reach.lean\#L60-L68}{\texttt{limits}},
\href{https://github.com/yang-mills-group/OpenAI-YM-20261028/blob/3456e5e75cf8d4cd558955a5c7638fba2d68f464/YM/Shoot/Reach.lean\#L81-L90}{\texttt{gates}},
\href{https://github.com/yang-mills-group/OpenAI-YM-20261028/blob/3456e5e75cf8d4cd558955a5c7638fba2d68f464/YM/Shoot/Reach.lean\#L73-L79}{\texttt{Drifting.pay}}.

\textbf{Lemma 4.11 (narrowing).} Let \(x\le y\) and let \(f\) be
continuous on \([x,y]\) with \(f(x)\le d\le u\le f(y)\). Then there are
\(x\le z\le w\le y\) with \(f(z)=d\), \(f(w)=u\) and
\(f([z,w])\subset[d,u]\). \emph{Sketch:} let \(w\) be the least point of
\(f^{-1}(u)\cap[x,y]\) and \(z\) the greatest point of
\(f^{-1}(d)\cap[x,w]\); both sets are nonempty by the intermediate value
theorem and compact by continuity. Suppose some \(c\in[z,w]\) has
\(f(c)<d\). Then the intermediate value theorem on \([c,w]\) gives a
point of \(f^{-1}(d)\) after \(z\), a contradiction. Suppose instead
\(f(c)>u\). Then it gives a point of \(f^{-1}(u)\) before \(w\), again a
contradiction. \emph{Lean:}
\href{https://github.com/yang-mills-group/OpenAI-YM-20261028/blob/3456e5e75cf8d4cd558955a5c7638fba2d68f464/YM/Shoot/Reach.lean\#L18-L49}{\texttt{narrow}}.

\textbf{Lemma 4.12 (nested intervals).} Fix \(N,p\) and assume (start),
(advance) and the drift bound. One can choose intervals
\(I_m=[l_m,u_m]\), \(0\le m\le N\), with \(I_N=[\beta^-_N,\beta^+_N]\)
and \(I_m\subset I_{m+1}\) for \(m<N\), and partial trajectories
\((\rho^x_i)_{m\le i\le N}\) for \(x\in[l_m,u_m]\), such that:

\begin{itemize}
\tightlist
\item
  \(\rho^x_N\) is the start representation at bare coupling \(x\);
\item
  every \(\rho^x_i\) is admissible with
  \(\beta(\rho^x_i)\in\mathfrak B_i\);
\item
  consecutive levels relay;
\item
  \(x\mapsto\rho^x_m\) is continuous;
\item
  \(\beta(\rho^{l_m}_m)=\beta^-_m\) and
  \(\beta(\rho^{u_m}_m)=\beta^+_m\).
\end{itemize}

\emph{Sketch:} downward induction.

\begin{itemize}
\tightlist
\item
  For \(m=N\), restrict the start family to \([\beta^-_N,\beta^+_N]\).
\item
  For \(m+1\to m\), the family at level \(m+1\) has all couplings in
  \(\mathfrak B_{m+1}\), so (advance) yields a continuous family
  \(\rho^x_m\). By Lemma 4.10, \(f(x)=\beta(\rho^x_m)\) satisfies
  \(f(l_{m+1})<\beta^-_m\) and \(f(u_{m+1})>\beta^+_m\).
\item
  Lemma 4.11 gives \([l_m,u_m]\) on which \(f\) maps onto
  \([\beta^-_m,\beta^+_m]\subset\mathfrak B_m\) with the right endpoint
  values.
\end{itemize}

The restriction is needed: (advance) applies only to families whose
couplings lie in the band at every point. Only the continuity of
\(x\mapsto\beta(\rho^x_m)\) enters the intermediate value argument. The
other continuity clauses are carried because (advance) requires them at
the next step. \emph{Lean:}
\href{https://github.com/yang-mills-group/OpenAI-YM-20261028/blob/3456e5e75cf8d4cd558955a5c7638fba2d68f464/YM/Shoot/Reach.lean\#L112-L154}{\texttt{steps}}.

\textbf{Lemma 4.13 (endpoint and assembly; proof of III.2.3).} Since
\(s_0=0\), we have \(\beta^-_0=\beta^+_0=\bar\beta\). So \(x=l_0\) gives
a trajectory from level \(N\) to level 0 with
\(\beta_{N,p,0}=\bar\beta\). Choose one such trajectory for every
\((N,p)\). Then:

\begin{itemize}
\tightlist
\item
  the path data are \(P^{\mathcal R}_{N,m}\), independent of \(p\);
\item
  admissibility gives the budget, clip, polymer bounds and band
  membership;
\item
  (start) gives the Wilson weight with \(v=0\);
\item
  the relays give the transfer identity and the drift inequality;
\item
  Lemma 4.9 gives the running law.
\end{itemize}

Hence the history is selected. \emph{Lean:}
\href{https://github.com/yang-mills-group/OpenAI-YM-20261028/blob/3456e5e75cf8d4cd558955a5c7638fba2d68f464/YM/Shoot/Reach.lean\#L156-L159}{\texttt{successful}},
\href{https://github.com/yang-mills-group/OpenAI-YM-20261028/blob/3456e5e75cf8d4cd558955a5c7638fba2d68f464/YM/Shoot/Reach.lean\#L95-L98}{\texttt{Competes}},
\href{https://github.com/yang-mills-group/OpenAI-YM-20261028/blob/3456e5e75cf8d4cd558955a5c7638fba2d68f464/YM/Shoot/Race.lean\#L67-L82}{\texttt{Recipe.races}}.

\paragraph{4.4 Persistent purchases and their
composition}\label{persistent-purchases-and-their-composition}

\textbf{Lemma 4.14 (composition rules).} Persistent numerical
requirements can be combined in the following forms:

\begin{itemize}
\tightlist
\item
  \begin{enumerate}
  \def\labelenumi{(\roman{enumi})}
  \tightlist
  \item
    a floor \(f\) such that the requirement holds for every later block
    factor;
  \end{enumerate}
\item
  (ii) ``for every delay \(V\) there is \(V'\ge V\) such that every
  later delay, every barrier and wait, and every level system meeting
  those entrance conditions satisfy \(P\)'';
\item
  (iii) the same with waits instead of delays.
\end{itemize}

Floors combine by pointwise maximum. Suppliers of forms (ii) and (iii)
combine by chaining, because the orders on delays and waits are
transitive and the conclusions persist under raising. \emph{Lean:}
\href{https://github.com/yang-mills-group/OpenAI-YM-20261028/blob/3456e5e75cf8d4cd558955a5c7638fba2d68f464/YM/Integration/NativePreselectedPostnumericPeriodClock.lean\#L31-L39}{\texttt{postnumeric\_quote\_diagonal\_pair\_requests}},
and the chains in
\href{https://github.com/yang-mills-group/OpenAI-YM-20261028/blob/3456e5e75cf8d4cd558955a5c7638fba2d68f464/YM/Integration/RG05Preconstruction/Expenses.lean\#L28-L119}{\texttt{joinExpensesShipWithPrediction}}.

\textbf{Lemma 4.15 (diagonal payments; proof of III.2.8).}

\begin{itemize}
\tightlist
\item
  For (b), raise \(\beta^{\rm ent}_j\) to
  \(\max\big(\beta^{\rm ent}_j,\exp\exp(1+\sum_{n,i\le j}|F(n,i)|)\big)\).
  The entrance condition \(\beta^{\rm ent}_{J_{m+1}}\le\hat\beta_m\)
  then bounds the partial sums at stage \(J_{m+1}\) by
  \(\log\log\hat\beta_m\). Since \(J\to\infty\), each fixed \(n\)
  satisfies \(n\le J_{m+1}\) from some level on.
\item
  For the single-array form, raise \(\beta^{\rm ent}_j\) to
  \(\max(\beta^{\rm ent}_j,F(j))\).
\item
  For (c), encode the requests. Take
  \(\mathrm{pre}(j)=\max_{\mathrm{code}(\alpha)\le j}\mathrm{pre}_\alpha(j)\).
  Raise \(\beta^{\rm ent}_j\) until the finitely many costs with
  \(\mathrm{code}(\alpha)\le j\), evaluated on the stage-\(j\) geometry,
  are below \(1+\log(2+\log x)\) for \(x\ge\beta^{\rm ent}_j\).
\item
  For (d), fix the plan before the block factors, apply (c) to its
  requests, and raise the delay once more for the inverse charge, whose
  block floor is zero.
\end{itemize}

\emph{Lean:}
\href{https://github.com/yang-mills-group/OpenAI-YM-20261028/blob/3456e5e75cf8d4cd558955a5c7638fba2d68f464/YM/Integration/NativeDiagonalClockRequest.lean\#L52-L79}{\texttt{eventuallyPaid}},
\href{https://github.com/yang-mills-group/OpenAI-YM-20261028/blob/3456e5e75cf8d4cd558955a5c7638fba2d68f464/YM/RG08/Calendar/Base.lean\#L68-L77}{\texttt{minDelay}},
\href{https://github.com/yang-mills-group/OpenAI-YM-20261028/blob/3456e5e75cf8d4cd558955a5c7638fba2d68f464/YM/Integration/CalendarTreasuryPurchase.lean\#L42-L64}{\texttt{buy}}.

\textbf{Lemma 4.16 (the sheet).} Let \(G\) be simple, let \(K\) have
branches at every patch, let the order data \(v\) satisfy the order
conditions of III.2.4, and fix a sequence \(C_j\ge1\). There are a sheet
\(s\) with kernel \(K\) and the prescribed first stage, orders, powers
and support exponents; two price maps \(O,T\); real sequences
\(\lambda,\gamma\) with \(\gamma_j>0\); and a nondecreasing sequence
\(K'_j\ge1\). For every tariff and scope there is a floor \(f\ge f^s\)
with the following properties.

The sheet satisfies the stationary-path suitability conditions at stages
\(j\ge j_0\) and the budget plans for \(C,K'\). For every floor
\(n\ge f\) and every admissible block-factor sequence \(\ell\), there is
a sequence \(\eta\) satisfying the stationary budget-gap condition. For
every post-\(L\) cost record \(F\) with \(\eta_F=\eta\),
\(\Gamma_F\ge1\) and \(c_F\le\gamma\), every protection, additional
tariff, nonnegative lower bound \(z\) and delay \(A\), there are a delay
\(D\ge A\), a transition-bound array bounded below by \(z\), and two
further transition arrays. Every later delay \(E\ge D\) satisfies the
pricing, ringing and comparison conditions. For every barrier, wait and
level system meeting those entrance conditions, the budget and
multisource conditions hold, and there are a Gaussian horizon and
kinetic data with planned stationary paths, the Wilson top path, the
cadre and amplification conditions, and the retained free and stationary
data specified by the perceived-estimate package.

\emph{Sketch:} the free-kernel and stationary suppliers first fix widths
and preliminary bounds. The sheet construction fixes clips and the
prescribed orders; the startup suppliers enlarge path bounds and floors.
The final profile construction raises drift allowances to dominate the
stationary supplier's bounds and raises column bounds without changing
widths, paths or the prescribed orders. After choosing \(\ell\), the
stationary and multisource suppliers raise the delay and construct the
kinetic data on each admitted level system. The retained data identify
these kinetic paths with the revised stationary paths. \emph{Lean:}
\href{https://github.com/yang-mills-group/OpenAI-YM-20261028/blob/3456e5e75cf8d4cd558955a5c7638fba2d68f464/YM/FullVolume/Roll/Fiscal.lean\#L97-L169}{\texttt{sailsW}},
\href{https://github.com/yang-mills-group/OpenAI-YM-20261028/blob/3456e5e75cf8d4cd558955a5c7638fba2d68f464/YM/FullVolume/Roll/Fiscal.lean\#L63-L79}{\texttt{Perceived}}.
\emph{(proof read at the level of the composition)}

\textbf{Lemma 4.17 (the native purchase and prediction).} Let \(G\) be
simple, fix the perceived-estimate package of Lemma 4.16 at a floor
\(v\), and let \(\tau>0\). There are a floor \(j\ge v\), programs,
initial wages, targets, expenses and a baton, chosen before the block
factors. The targets have the canonical chart and halo \(500\); the
expenses have the prescribed earnings, deficit \(-7/16\) and degree
charge \(0\).

For every floor \(l\ge j\), admissible block factors exist. For every
such choice there is a wage threshold. For every wage above it there are
duration bounds, a register and auxiliary size bounds. For every postage
cap there are postage rates at most that cap and scouts satisfying the
register condition. For every exchange datum and tariff there are reset
prices, reset bounds and post-\(L\) costs. For every positive sequence
\(\theta^0\) there are protection \(\theta\le\theta^0\) and step records
with \(\varepsilon\le1/4\) on transitions whose output stage is at least
the first stage. For every menu and initial delay there are booths; for
every auxiliary record there is an ordered delay above the initial
delay.

Every later delay, barrier, wait and admitted level system satisfies the
budget and multisource conditions and carries a directed datum on that
same system. Its shelter conditions hold at the two pairs of tolerances
\((1/16,1/16)\) and \((\tau,\tau)\), with the third tolerance \(1/8\) in
both cases. It also has the top kinetic data, cadre, amplification, step
conditions, origins and rank-covariance data. The offering condition
additionally assumes divisible periods.

\emph{Sketch:} the shelter supplier gives the initial floor and data. A
further floor pays the register's length requirement; a delay supplier
supplies the step conditions before the booths are chosen. The shelter
output then supplies the remaining data. For III.2.6, combine the
prediction's floor \(u\) with \(j\) by maximum. Insert its delay before
the booth delay; persistence then gives the prediction on the same
emitted factor whenever the divisibility condition holds. \emph{Lean:}
\href{https://github.com/yang-mills-group/OpenAI-YM-20261028/blob/3456e5e75cf8d4cd558955a5c7638fba2d68f464/YM/FullVolume/Roll/Linger.lean\#L23-L51}{\texttt{companyK}}
\emph{(proof read at the level of the composition)};
\href{https://github.com/yang-mills-group/OpenAI-YM-20261028/blob/3456e5e75cf8d4cd558955a5c7638fba2d68f464/YM/Integration/NativeSelection/FundedCompanySelection.lean\#L49-L106}{\texttt{company}}.

\textbf{Lemma 4.18 (proof of III.2.4 and III.2.5).} The proofs proceed
in four steps.

\begin{enumerate}
\def\labelenumi{\arabic{enumi}.}
\item
  \href{https://github.com/yang-mills-group/OpenAI-YM-20261028/blob/3456e5e75cf8d4cd558955a5c7638fba2d68f464/YM/Integration/RG05Preconstruction/Expenses.lean\#L133-L144}{\texttt{joinRadarsShipWithPrediction}}
  takes the sheet from Lemma 4.16 and calls
  \href{https://github.com/yang-mills-group/OpenAI-YM-20261028/blob/3456e5e75cf8d4cd558955a5c7638fba2d68f464/YM/Integration/RG05Preconstruction/Expenses.lean\#L28-L119}{\texttt{joinExpensesShipWithPrediction}}.
  That theorem starts from Lemma 4.17 and threads about thirty further
  suppliers through it:

  \begin{itemize}
  \tightlist
  \item
    floor suppliers, joined into one floor;
  \item
    delay suppliers \(V\le V_1\le\dots\le V_7\le\dots\le w_6\le\dots\),
    chained.
  \end{itemize}
\item
  \href{https://github.com/yang-mills-group/OpenAI-YM-20261028/blob/3456e5e75cf8d4cd558955a5c7638fba2d68f464/YM/Integration/RG05Preconstruction/Fares.lean\#L26-L85}{\texttt{richFaresWithPrediction}}
  adds:

  \begin{itemize}
  \tightlist
  \item
    additional costs and floors before choosing the block factors;
  \item
    a cross-term floor, so the final floor is \((t\vee L\vee\mu)\vee\)
    cross-floor;
  \item
    shells raised to \(\max(S_j,\cdot)\), the mesh and installment
    delays, and a wait enlargement.
  \end{itemize}
\item
  \href{https://github.com/yang-mills-group/OpenAI-YM-20261028/blob/3456e5e75cf8d4cd558955a5c7638fba2d68f464/YM/Integration/RG05Preconstruction/Bardo.lean\#L26-L81}{\texttt{plentifulWithPrediction}}
  adds:

  \begin{itemize}
  \tightlist
  \item
    trace floors and the fixed trace-control floor;
  \item
    trace delays, including a raise of the entry couplings;
  \item
    two wait enlargements.
  \end{itemize}

  It concludes with
  \href{https://github.com/yang-mills-group/OpenAI-YM-20261028/blob/3456e5e75cf8d4cd558955a5c7638fba2d68f464/YM/TraceExtension/Chart/Line.lean\#L49-L56}{\texttt{Trace.Over.gang}}.
\item
  Floors are combined by maximum and delay and wait bounds by
  transitivity, as in Lemma 4.14. The final family conclusion uses the
  hypotheses of III.2.5.
\end{enumerate}

\emph{(proof read at the level of the composition; the suppliers were
not read)}

\textbf{Lemma 4.19 (the final assembly; proof of III.2.7).}
\href{https://github.com/yang-mills-group/OpenAI-YM-20261028/blob/3456e5e75cf8d4cd558955a5c7638fba2d68f464/YM/Integration/NativePreselectedPostnumericPeriodClock.lean\#L44-L179}{\texttt{allFactor…}}
makes the choices in order:

\begin{enumerate}
\def\labelenumi{\arabic{enumi}.}
\tightlist
\item
  \emph{Pre-\(L\) data.}

  \begin{itemize}
  \tightlist
  \item
    Order data: \(j_0=4\), \(k_j=6+j\), \(q_j=2+j\), \(a_j=100\); and
    \(Q_i=i+2\).
  \item
    A kernel with penalty \(b\).
  \item
    A tableau at a parameter \(\xi\ge2\) and sphere data with
    \(\omega_v=1/1024\) and \(0<\omega_t\le1\). The contraction supplier
    chooses \(\omega_t\); its auxiliary constant \(u\ge0\) satisfies
    \(6u\omega_t\le\omega_v\).
  \end{itemize}
\item
  \emph{The plentiful construction.} Apply III.2.5. Pick a choice \(c\)
  of the demand.
\item
  \emph{Block floor and block factors.} The block floor is the maximum
  of five floors: the plentiful floor, a density floor, the treasury
  floor, the period-plan floor and \(e_c\). Fix block factors, wages and
  post-\(L\) costs.
\item
  \emph{Smallness exponent and delays.} Take \(\theta^0\equiv1\) and
  initial data \(\beta^{\rm ent}=\sigma^{\rm ent}=0\), \(S_j=g_j=1\).
  Raise delays in this order: halo, step, thermal requests, \(\Lambda\)
  and both demand arrays together with the all-stage payment of the
  zeroth post-\(L\) row, treasury, period plan, auxiliary budget, silk,
  custody, the delay of III.2.5, the quote bounds, and the header
  bounds.
\item
  \emph{Barrier.} Take \(\bar\beta\ge2\) (Lemma 4.3).
\item
  \emph{Wait.} First apply the wait supplier of III.2.5 to a zero-window
  wait. Separately enlarge a slim wait for density bounds. Choose a
  common enlargement satisfying both waits and the required coverage
  bounds.
\item
  \emph{Level system.} Take the level system of the divisibility variant
  of III.2.2, and the factor from III.2.5.
\end{enumerate}

Finally
\href{https://github.com/yang-mills-group/OpenAI-YM-20261028/blob/3456e5e75cf8d4cd558955a5c7638fba2d68f464/YM/Integration/RG05Preconstruction/PeriodConstruction.lean\#L30-L47}{\texttt{engageWithCustodyFactor}}
builds the run from the factor. \emph{(proof read at the level of the
composition)}

\paragraph{4.5 The analytic lemmas}\label{the-analytic-lemmas}

\textbf{Lemma 4.20 (mechanisms behind III.2.9).} Under the hypotheses of
III.2.9, the estimates follow as follows.

\begin{itemize}
\tightlist
\item
  \emph{Division:} analytic division first gives complex coefficient
  packets on a smaller strip. Taking their real parts and using
  conjugation symmetry preserves the identity at real momenta and gives
  the stated real kernels.
\item
  \emph{Pullback gain:} divide \(R\) into terms with \(d\) forward
  differences and move those differences onto \(B\). The two cross
  bounds contribute \(L^{-2}\), the sum over fine sites contributes
  \(L^4\), and the differences contribute \(L^{-d}\), giving
  \(L^{2-d}\).
\item
  \emph{Contraction:} apply the gain with \(d=3\) and choose
  \(N>\max(\tau/\eta,K^2)\). Then \(K/L\le L^{-1/2}\) for \(L\ge N\).
\item
  \emph{Multilinear transfer:} divide in the independent momenta and
  distribute the \(d\) differences among the \(r+1\) crosses. Their
  bounds contribute \(L^{-(r+1+d)}\); summing the free translation gives
  \(L^4\). The width reserve pays for symmetrization and recentering.
\item
  \emph{Odd-part cancellation:} reflection cancels the first odd
  contribution. The rate calculus controls the remaining second
  derivative with every exponent strictly below \(2\). The second
  derivative of a function composed with \(y\mapsto-y\) equals its
  original second derivative at \(0\).
\item
  \emph{Local read-out:} subtract the presentations. Equality of
  densities expresses the coefficient difference as the difference of
  the regular remainders on the healthy domain. The local projection
  kills nonwrapping supports. A wrapping support has an exponentially
  small contribution, bounded using the side condition; finitely many
  test functionals then control the coefficient norm.
\end{itemize}

\emph{(proofs read at the level of the displayed declarations; the
analytic division, transfer and local-projection inputs are not
independently audited)}

\subsubsection{5. Where the difficulty
lies}\label{where-the-difficulty-lies-4}

\begin{itemize}
\tightlist
\item
  \textbf{Order and persistence.} Each floor must be fixed before the
  block factors that pay it. Costs depending on the chosen factors can
  be paid by later entry thresholds. Delay and wait suppliers must
  remain valid under all subsequent increases; otherwise chaining them
  does not give a common set of constants.
\item
  \textbf{Uniformity in depth and volume.} The constants and schedule
  are common to every \(N,p\). The finite first plateau is handled by
  one period offset \(M\), giving \(n_p=4(M+p+1)\). Later volume
  conditions follow from the exponential growth of the side length.
\item
  \textbf{Continuity during shooting.} Each new interval must keep every
  member of the family inside the next coupling band. The
  nested-interval argument preserves this condition while fixing the
  endpoint couplings. It requires continuity of the full representation
  family at each new step.
\item
  \textbf{Two drift bounds.} The bound
  \(\omega_j\le(\kappa/4)\log\ell_j\) gives the shooting gates. The
  bound \(\omega_j\le\log\ell_j/j^2\) gives the running law. Both are
  paid after the stage constants are fixed.
\item
  \textbf{Eventual payment.} The entry thresholds can absorb any fixed
  real array, including one evaluated after the block factors are
  chosen. Each fixed row is paid eventually; a single level need not pay
  every row. The stronger all-stage condition for the zeroth post-\(L\)
  row is purchased separately.
\end{itemize}

\subsubsection{6. Notation}\label{notation-4}

{\def\LTcaptype{none} % do not increment counter
\begin{longtable}[]{@{}
  >{\raggedright\arraybackslash}p{(\linewidth - 4\tabcolsep) * \real{0.3333}}
  >{\raggedright\arraybackslash}p{(\linewidth - 4\tabcolsep) * \real{0.3333}}
  >{\raggedright\arraybackslash}p{(\linewidth - 4\tabcolsep) * \real{0.3333}}@{}}
\toprule\noalign{}
\begin{minipage}[b]{\linewidth}\raggedright
Written proof
\end{minipage} & \begin{minipage}[b]{\linewidth}\raggedright
Lean
\end{minipage} & \begin{minipage}[b]{\linewidth}\raggedright
Meaning
\end{minipage} \\
\midrule\noalign{}
\endhead
\bottomrule\noalign{}
\endlastfoot
sheet; stage record \(X_j\); first stage \(j_0\); floors \(f_j\) &
\href{https://github.com/yang-mills-group/OpenAI-YM-20261028/blob/3456e5e75cf8d4cd558955a5c7638fba2d68f464/YM/Price/Intent.lean\#L26-L35}{\texttt{Sheet}},
\href{https://github.com/yang-mills-group/OpenAI-YM-20261028/blob/3456e5e75cf8d4cd558955a5c7638fba2d68f464/YM/RG/History.lean\#L45-L61}{\texttt{Stage}}
& data fixed before the block factors \\
\(k_j,q_j,a_j,w_j,\varpi_j,\omega_j,\gamma_j,r_j\) & \texttt{order},
\texttt{power}, \texttt{supportExponent}, \texttt{width},
\texttt{paths}, \texttt{drift}, \texttt{columns},
\texttt{clip.relaxation} & fields of the stage record \\
\(\mathcal C_j\) & (in
\href{https://github.com/yang-mills-group/OpenAI-YM-20261028/blob/3456e5e75cf8d4cd558955a5c7638fba2d68f464/YM/Price/Intent.lean\#L39-L48}{\texttt{Large}})
& \(1+\varpi_j+w_j^{-1}+\gamma_j+r_j+a_j\) \\
\(\ell_j\); \(L_m=\ell_{J_{m+1}}\) &
\href{https://github.com/yang-mills-group/OpenAI-YM-20261028/blob/3456e5e75cf8d4cd558955a5c7638fba2d68f464/YM/Price/Intent.lean\#L50-L53}{\texttt{Rolls}};
\texttt{Schedule.block} & stage and step block factors \\
\(\theta_j\) &
\href{https://github.com/yang-mills-group/OpenAI-YM-20261028/blob/3456e5e75cf8d4cd558955a5c7638fba2d68f464/YM/Price/Intent.lean\#L88-L92}{\texttt{Protection}}
& smallness exponent of the loop cap \\
\(\beta^{\rm ent}_j,\sigma^{\rm ent}_j(\bar\beta),S_j,g_j\) &
\href{https://github.com/yang-mills-group/OpenAI-YM-20261028/blob/3456e5e75cf8d4cd558955a5c7638fba2d68f464/YM/Price/Intent.lean\#L115-L120}{\texttt{Delays}}
& entry thresholds, shells, growth \\
\(\bar\beta\) &
\href{https://github.com/yang-mills-group/OpenAI-YM-20261028/blob/3456e5e75cf8d4cd558955a5c7638fba2d68f464/YM/Price/Intent.lean\#L123-L127}{\texttt{Bar}}
& barrier: coupling at level 0 \\
\(\sigma^{\rm rdy}_j,h_j,n_{\min}\) &
\href{https://github.com/yang-mills-group/OpenAI-YM-20261028/blob/3456e5e75cf8d4cd558955a5c7638fba2d68f464/YM/Price/Intent.lean\#L134-L139}{\texttt{Wait}}
& readiness, windows, minimal period \\
meets the entrance conditions &
\href{https://github.com/yang-mills-group/OpenAI-YM-20261028/blob/3456e5e75cf8d4cd558955a5c7638fba2d68f464/YM/Price/Intent.lean\#L145-L162}{\texttt{Paid}}
& III.2.2(a)--(d) \\
\(\mathrm{bill}_j(s)\) &
\href{https://github.com/yang-mills-group/OpenAI-YM-20261028/blob/3456e5e75cf8d4cd558955a5c7638fba2d68f464/YM/Price/Length.lean\#L43-L47}{\texttt{Bar.bill}}
& required sites per unit length \\
\(J'\), \(J_m=j_0+J'_m\) &
\href{https://github.com/yang-mills-group/OpenAI-YM-20261028/blob/3456e5e75cf8d4cd558955a5c7638fba2d68f464/YM/Price/Waits.lean\#L22-L28}{\texttt{Program}}
& stage index \\
\(n_p=4(M+p+1)\) & \texttt{Program.page} & level-0 side, with
\(M\in\mathbb N\) \\
recipe \(\mathcal R\) &
\href{https://github.com/yang-mills-group/OpenAI-YM-20261028/blob/3456e5e75cf8d4cd558955a5c7638fba2d68f464/YM/Shoot/Criteria.lean\#L67-L71}{\texttt{Recipe}}
& path data and three tests \\
admissible; relays; start &
\href{https://github.com/yang-mills-group/OpenAI-YM-20261028/blob/3456e5e75cf8d4cd558955a5c7638fba2d68f464/YM/Shoot/Criteria.lean\#L74-L79}{\texttt{Lives}},
\href{https://github.com/yang-mills-group/OpenAI-YM-20261028/blob/3456e5e75cf8d4cd558955a5c7638fba2d68f464/YM/Shoot/Criteria.lean\#L88-L91}{\texttt{Relays}},
\href{https://github.com/yang-mills-group/OpenAI-YM-20261028/blob/3456e5e75cf8d4cd558955a5c7638fba2d68f464/YM/Shoot/Criteria.lean\#L93-L96}{\texttt{Starter}}
& \\
continuous family &
\href{https://github.com/yang-mills-group/OpenAI-YM-20261028/blob/3456e5e75cf8d4cd558955a5c7638fba2d68f464/YM/Shoot/Criteria.lean\#L51-L57}{\texttt{Curve}},
\href{https://github.com/yang-mills-group/OpenAI-YM-20261028/blob/3456e5e75cf8d4cd558955a5c7638fba2d68f464/YM/Shoot/Criteria.lean\#L44-L49}{\texttt{deviation}}
& \\
(start)+(advance) &
\href{https://github.com/yang-mills-group/OpenAI-YM-20261028/blob/3456e5e75cf8d4cd558955a5c7638fba2d68f464/YM/Shoot/Criteria.lean\#L99-L106}{\texttt{Recipe.Issues}}
& shooting hypotheses \\
\(\mathfrak B_m\); \([\beta^-_m,\beta^+_m]\) &
\href{https://github.com/yang-mills-group/OpenAI-YM-20261028/blob/3456e5e75cf8d4cd558955a5c7638fba2d68f464/YM/Shoot/Criteria.lean\#L63}{\texttt{band}};
\href{https://github.com/yang-mills-group/OpenAI-YM-20261028/blob/3456e5e75cf8d4cd558955a5c7638fba2d68f464/YM/Shoot/Reach.lean\#L57}{\texttt{floorBeta}},
\href{https://github.com/yang-mills-group/OpenAI-YM-20261028/blob/3456e5e75cf8d4cd558955a5c7638fba2d68f464/YM/Shoot/Reach.lean\#L58}{\texttt{ceilBeta}}
& band; shooting window \\
eventually paid &
\href{https://github.com/yang-mills-group/OpenAI-YM-20261028/blob/3456e5e75cf8d4cd558955a5c7638fba2d68f464/YM/R2I/Campaign/Clock.lean\#L36-L38}{\texttt{EventuallyPaid}}
& III.2.8(a) \\
plan; treasury funded &
\href{https://github.com/yang-mills-group/OpenAI-YM-20261028/blob/3456e5e75cf8d4cd558955a5c7638fba2d68f464/YM/Integration/CalendarTreasuryPurchase.lean\#L18-L20}{\texttt{Plan}},
\href{https://github.com/yang-mills-group/OpenAI-YM-20261028/blob/3456e5e75cf8d4cd558955a5c7638fba2d68f464/YM/Integration/CalendarTreasuryPurchase.lean\#L35-L38}{\texttt{TreasuryFunded}}
& III.2.8(d) \\
emitted factor; prediction &
\href{https://github.com/yang-mills-group/OpenAI-YM-20261028/blob/3456e5e75cf8d4cd558955a5c7638fba2d68f464/YM/PhysicalNext/RG_NEXT_COMPLETE_REPEATABLE_NONLINEAR_STEP/Roost/Foundry.lean\#L26-L50}{\texttt{Factor}},
\href{https://github.com/yang-mills-group/OpenAI-YM-20261028/blob/3456e5e75cf8d4cd558955a5c7638fba2d68f464/YM/Integration/NativeSelection/FundedCompanySelection.lean\#L23-L24}{\texttt{Prediction}}
& record of all choices; predicate on it \\
sheet with floors \(n\) &
\href{https://github.com/yang-mills-group/OpenAI-YM-20261028/blob/3456e5e75cf8d4cd558955a5c7638fba2d68f464/YM/PhysicalNext/UNMASKED_THROUGH_LAUNCH/Merchant.lean\#L36}{\texttt{raise}},
\href{https://github.com/yang-mills-group/OpenAI-YM-20261028/blob/3456e5e75cf8d4cd558955a5c7638fba2d68f464/YM/FullVolume/Roll/Fiscal.lean\#L22-L23}{\texttt{bountyV}}
& \\
order data \(v\) &
\href{https://github.com/yang-mills-group/OpenAI-YM-20261028/blob/3456e5e75cf8d4cd558955a5c7638fba2d68f464/YM/PhysicalNext/UNMASKED_THROUGH_LAUNCH/Dome.lean\#L64-L68}{\texttt{Orders}},
\href{https://github.com/yang-mills-group/OpenAI-YM-20261028/blob/3456e5e75cf8d4cd558955a5c7638fba2d68f464/YM/PhysicalNext/UNMASKED_THROUGH_LAUNCH/Dome.lean\#L70-L77}{\texttt{Draft}},
\href{https://github.com/yang-mills-group/OpenAI-YM-20261028/blob/3456e5e75cf8d4cd558955a5c7638fba2d68f464/YM/PhysicalNext/UNMASKED_THROUGH_LAUNCH/Dome.lean\#L81-L85}{\texttt{WithStage}}
& \\
\end{longtable}
}

\subsubsection{7. What remains
unexplained}\label{what-remains-unexplained-4}

\begin{enumerate}
\def\labelenumi{\arabic{enumi}.}
\tightlist
\item
  \textbf{Analytic content of the purchased conditions.} The
  common-choice statements III.2.4--III.2.7 and their quantifier order
  are verified against the displayed Lean interfaces. The predicates in
  the requirement table, the conditional family package, and the
  campaign's pillars, funding, news, custody and sky conditions are not
  fully expanded into mathematical estimates here. Their equivalence to
  the requirements stated in the owning units remains to be checked.
\item
  \textbf{Underlying suppliers.} The proofs composing the sheet, native,
  joint, plentiful and final-run constructions have been checked at
  their call sites. The proofs of the underlying analytic suppliers,
  including those used by the roughly thirty additions in the joint
  construction, are not independently verified here. The same limitation
  applies to the plan's dictionaries and inverse-charge supplier, and to
  the analytic division and projection results used in III.2.9.
\item
  \textbf{Auxiliary parameters.} \(K'\) is a nondecreasing sequence
  bounded below by \(1\); the tableau is a finite interpolation rule of
  degree \(300\); the sphere condition is a contraction estimate. The
  detailed analytic roles of \(Y\), \(Q\), wages, postage, exchanges,
  menus and the remaining auxiliary records still require identification
  in the step estimates.
\item
  \textbf{Drift normalization.} The sheet construction raises the drift
  allowance to dominate the stationary supplier's sequence. Its
  identification with the bound \(1+D\) used in III.9, including the
  relation between the constants attributed there to III.7 and III.13,
  remains unverified.
\item
  \textbf{Analytic interfaces.} The weighted summability and
  vanishing-order definitions, the first-site support condition in
  multilinear transfer, and the rate interpretation of odd-part
  cancellation have been checked. The full step-scene hypotheses of
  odd-part cancellation and the coefficient-format and readability
  conditions of local read-out still need an independent mathematical
  account. Only the statements of the supplementary Hessian,
  cubic/quartic extraction and trace-insertion estimates linked after
  III.2.9 were checked; their proofs remain outside this review.
\item
  \textbf{Scope of the source review.} The arithmetic construction, the
  divisibility argument, and the shooting proof are verified against the
  source. Other uses of the non-divisible period construction, the
  allocation of the supplementary analytic estimates to outline units,
  and the line-count breakdown given at the outset have not been
  independently checked.
\end{enumerate}

\subsection{III.3. Free lattice kernels and
covariances}\label{iii.3.-free-lattice-kernels-and-covariances}

\textbf{Part:} III, Ultraviolet construction (renormalization group).
\textbf{Lean:}
\href{https://github.com/yang-mills-group/OpenAI-YM-20261028/tree/3456e5e75cf8d4cd558955a5c7638fba2d68f464/YM/Next/FREE_KERNELS}{\texttt{YM/Next/FREE\_KERNELS/}}
(5,450 lines),
\href{https://github.com/yang-mills-group/OpenAI-YM-20261028/tree/3456e5e75cf8d4cd558955a5c7638fba2d68f464/YM/MultiNext/SIGNED_SOURCE_CONTINUATION}{\texttt{YM/MultiNext/SIGNED\_SOURCE\_CONTINUATION/}}
(1,674),
\href{https://github.com/yang-mills-group/OpenAI-YM-20261028/tree/3456e5e75cf8d4cd558955a5c7638fba2d68f464/YM/Further/FREE_COVARIANCE_KP}{\texttt{YM/Further/FREE\_COVARIANCE\_KP/}}
(924),
\href{https://github.com/yang-mills-group/OpenAI-YM-20261028/tree/3456e5e75cf8d4cd558955a5c7638fba2d68f464/YM/UVNext/PRICED_KERNELS}{\texttt{YM/UVNext/PRICED\_KERNELS/}}
(533),
\href{https://github.com/yang-mills-group/OpenAI-YM-20261028/tree/3456e5e75cf8d4cd558955a5c7638fba2d68f464/YM/Further/CORRECTED_LIFT_COEFFICIENTS}{\texttt{YM/Further/CORRECTED\_LIFT\_COEFFICIENTS/}}
(248),
\href{https://github.com/yang-mills-group/OpenAI-YM-20261028/tree/3456e5e75cf8d4cd558955a5c7638fba2d68f464/YM/Push/CUTOFF_PROJECTORS}{\texttt{YM/Push/CUTOFF\_PROJECTORS/}}
(139),
\href{https://github.com/yang-mills-group/OpenAI-YM-20261028/tree/3456e5e75cf8d4cd558955a5c7638fba2d68f464/YM/ProofPorts/GRID_RENEWED_RESERVES}{\texttt{YM/ProofPorts/GRID\_RENEWED\_RESERVES/}}
(30). \textbf{Size:} 8,998 lines in the main-theorem closure.

\textbf{Conventions.}

\begin{itemize}
\tightlist
\item
  \textbf{Tori and forms.} \(\Lambda_n=(\mathbb Z/n\mathbb Z)^4\),
  \(n\ge1\). \(\Omega^r(\Lambda_n)\) is the space of real
  \(r\)-cochains, with the counting inner product.
  \(d_r:\Omega^r\to\Omega^{r+1}\) is the exterior derivative, and
  \(d=d_1\). \(X^*\) denotes the transpose on real cochains and the
  Hermitian adjoint of a symbol at real momentum. Displacements use
  \(|y|_1\); real momenta use the Euclidean norm \(|p|\).
\item
  \textbf{Levels.} In mesh \(N\) the levels are \(i=N,\dots,0\). Level
  \(N\) is the bare lattice. Step \(i<N\) maps level \(i+1\) (fine torus
  \(\Lambda_f\)) to level \(i\) (coarse torus \(\Lambda_c\)), with
  \(f=cL_i\) and \(L_i\ge3\) odd. Level \(i\) has a stage index
  \(J(i)\ge4\), and \(J(i)\le J(i+1)\le J(i)+1\).
\item
  \textbf{Kernels.} A \emph{kernel} is a map
  \(R:\mathbb Z^4\to\mathrm{Hom}(\mathbb R^\beta,\mathbb R^\alpha)\),
  where \(\alpha,\beta\) are sets of face types. It acts on every torus
  by periodisation:
  \((R_n\omega)(x)=\sum_{y\in\mathbb Z^4}R(y)\,\omega(x+y)\). Its symbol
  is \(\hat R(p)=\sum_yR(y)e^{ip\cdot y}\).
\item
  \textbf{Cross kernels.} A \emph{cross kernel} of block size \(L\) also
  carries a fine residue \(r\in(\mathbb Z/L)^4\) in its row index. It
  maps cochains on \(\Lambda_c\) to cochains on \(\Lambda_{cL}\), for
  every \(c\). Its displacements are measured on the coarse lattice.
\item
  \textbf{Strip bound.} \(R\in\mathcal S(\alpha,\Omega)\) means
  \(\sum_ye^{\alpha|y|_1}|R_{ab}(y)|\le\Omega\) for every entry. This
  implies that \(\hat R\) is analytic and bounded on
  \(|\mathrm{Im}\,p_a|<\alpha\). Conversely, a bounded periodic analytic
  symbol has such a weighted coefficient bound on each narrower strip.
\item
  \textbf{Aliases and lattice momenta.} For a coarse momentum \(p\), the
  fine aliases are \(q_\ell=(p+2\pi\ell)/L\), \(\ell\in(\mathbb Z/L)^4\)
  (centred representatives). \(\tau(p)=(e^{ip_a}-1)_{a=1}^4\) is the
  lattice momentum, and \(\sigma_r(p)\) is the symbol of \(d_r\), so
  \(\sigma_0=\tau\).
\item
  \textbf{Complex cubes.} Write
  \(\mathcal C_\theta=\{p\in\mathbb C^4:|\mathrm{Re}\,p_a|<\pi+\theta,\ |\mathrm{Im}\,p_a|<\theta\}\)
  and
  \(\mathcal C_\theta^{\mathbb R}=\mathcal C_\theta\cap\mathbb R^4\).
  Absolute squares in the real-momentum formulas are continued
  holomorphically as \(g(p)g(-p)\) for real-coefficient filters, and
  \(|\tau(p)|^2\) as \(\sum_a(e^{ip_a}-1)(e^{-ip_a}-1)\).
\item
  \textbf{Constants from the kernel choice.} \(\kappa>0\) is the penalty
  of the kernel choice \(K\) (the number \(b\) of the main theorem).
  \(c_i>0\) is the direct coefficient of the clip of level \(i\), that
  is, the coefficient of the bare plaquette term in the quadratic part
  of the action.
\end{itemize}

\subsubsection{1. What this unit proves}\label{what-this-unit-proves-5}

The unit computes the Gaussian part of the block-spin step of III.1 in
closed form, for every history of steps. It proves three things about
it:

\begin{itemize}
\tightlist
\item
  exact algebraic identities on every torus, giving an orthogonal
  splitting of the quadratic form;
\item
  exponential decay, uniform in the block sizes, history length and
  torus, with constants fixed by the kernel choice and stage index, and,
  after repair, by the sheet's direct clip coefficients;
\item
  uniform ellipticity and hypercubic symmetry of the quadratic forms.
\end{itemize}

The constants are fixed before the block sizes (Proposition III.3.1).
Positivity of the rolled kernel requires the lower bound \(L\ge b_j\).

The package also contains:

\begin{itemize}
\tightlist
\item
  the free covariance of the curvature and its comparison with the
  continuum (III.3.6, used in III.11);
\item
  calculus lemmas used in the nonlinear step (III.3.7, used in III.10
  and III.2).
\end{itemize}

\textbf{The identities of one step.} Fix a step with tori
\(\Lambda_f,\Lambda_c\), \(f=cL\). Let
\(Q:\Omega^1(\Lambda_f)\to\Omega^1(\Lambda_c)\) be the linearised block
average of III.1: a weighted average, over a \(J\)-fold smoothing patch,
of line integrals of the fine 1-form along \(L\)-link paths.

\emph{Metric data} on a level is a quintuple
\(\mathcal B=(D,D^{-1},H,H^{-1},R)\) of kernels on 2-forms: a diagonal
multiplier, its inverse, a metric, its inverse, and a root. With the
direct coefficient \(c\) of the level's clip, set
\[V=(\sqrt c\,1,\,R):\Omega^2\to\Omega^2\oplus\Omega^2 .\] Then
\(V^*V=c\,1+R^*R=H\) and \(\langle dA,H\,dA\rangle=|V\,dA|^2\).

\emph{Transfer kernels} are \(\mathcal U=(W,T,M_f,M_n,O)\) with:

\begin{itemize}
\tightlist
\item
  \(W:\Omega^1(\Lambda_c)\to\Omega^1(\Lambda_f)\) and
  \(T:\Omega^2(\Lambda_c)\to\Omega^2(\Lambda_f)\) (cross kernels);
\item
  \(M_f:(\Omega^2)^{\oplus2}(\Lambda_c)\to(\Omega^2)^{\oplus2}(\Lambda_f)\)
  (cross kernel);
\item
  \(M_n:(\Omega^2)^{\oplus2}(\Lambda_c)\to\Omega^1(\Lambda_c)\);
\item
  \(O:(\Omega^2)^{\oplus2}(\Lambda_c)\to\big((\Omega^2)^{\oplus2}\oplus\Omega^3\big)(\Lambda_c)\).
\end{itemize}

With fine data \((\mathcal B,V)\) and coarse data \((\mathcal B',V')\),
the identities are:

\begin{itemize}
\tightlist
\item
  (D1) \(d^*H\,d\,W=Q^*d^*H'd\)
\item
  (D2) \(QW=1-\kappa^{-1}d^*H'd\)
\item
  (D3) \(T^*d=dQ\)
\item
  (D4) \(T\,D'd=D\,d\,W\)
\item
  (R1) \(M_f^*M_f+M_n^*M_n+O^*O=1\)
\item
  (R2) \(\sqrt\kappa\,M_n^*=V'd\)
\item
  (R3) \(M_fV'd=V\,d\,W\)
\item
  (R4) \(M_nV'd=\kappa^{-1/2}d^*H'd\)
\item
  (R5) \(OV'd=0\)
\item
  (R6) \(M_f^*V\,d=\sqrt\kappa\,M_n^*Q\)
\item
  (Δ)
  \(H'\big(H'^{-1}-T^*H^{-1}T-\kappa^{-1}d\,d^*\big)H'=d_2^*P\,d_2\),
  for a 3-form kernel \(P\) on the coarse level.
\end{itemize}

What the identities mean:

\begin{itemize}
\tightlist
\item
  \textbf{(D1)--(D2).} These say that \(A=WB\) is a critical point,
  hence a minimiser, of
  \(A\mapsto\langle dA,H\,dA\rangle+\kappa|QA-B|^2\), and that the
  minimum is \(\langle dB,H'dB\rangle\). So \(W\) is Balaban's
  background-field (minimiser) map, and \(H'\) is the effective metric
  of the coarse level. The value identity is III.1's
  \href{https://github.com/yang-mills-group/OpenAI-YM-20261028/blob/3456e5e75cf8d4cd558955a5c7638fba2d68f464/YM/PREC/TreeMatch.lean\#L93-L103}{\texttt{action}}.
\item
  \textbf{(D3)--(D4).} (D3) makes \(T^*\) a 2-form average compatible
  with \(Q\). (D4) says that \(T\) carries the weighted coarse curvature
  \(D'dB\) to the weighted fine curvature \(D\,dWB\) of the background
  field.
\item
  \textbf{(R1)--(R6): orthogonal splitting.} For
  \(A\in\Omega^1(\Lambda_f)\), \(B\in\Omega^1(\Lambda_c)\) and
  \(Z=A-WB\), (R3), (R4) and (D2) give
  \[\big(V\,dA,\ \sqrt\kappa(B-QA)\big)=(M_f,M_n)V'dB+\big(V\,dZ,\,-\sqrt\kappa\,QZ\big).\]
  The two summands are orthogonal by (R6), and
  \(|(M_f,M_n)V'dB|=|V'dB|\) by (R1) and (R5). Hence
  \[\langle dA,H\,dA\rangle+\kappa|QA-B|^2=\langle dB,H'dB\rangle+\langle dZ,H\,dZ\rangle+\kappa|QZ|^2 .\]
\item
  \textbf{(Δ).} This is the same splitting for inverse metrics. It says
  \(H'^{-1}=T^*H^{-1}T+\kappa^{-1}dd^*+H'^{-1}d_2^*Pd_2H'^{-1}\), and
  the last term sees only \(d_2H'^{-1}\omega\).
\end{itemize}

\textbf{Proposition III.3.1 (catalogue fixed before the block sizes).}
Let \(K\) be a kernel choice. There are numbers \(\alpha_j>0\),
\(\eta_j>0\), \(\Omega_j\in\mathbb R\), \(\lambda_j\ge0\) (amendments),
\(\gamma_j>0\) (clip bounds) and \(b_j\in\mathbb N\) (block thresholds),
\(j\in\mathbb N\), depending only on \(K\), with the following
properties.

\begin{enumerate}
\def\labelenumi{(\alph{enumi})}
\tightlist
\item
  \emph{Metric data of a history.} Let \(j\ge4\), let the clip have
  direct coefficient \(c\le\gamma_j\), and let \(v\) be a history of
  stage index \(j\). Then there is metric data \((D,D^{-1},H,H^{-1},R)\)
  such that:
\end{enumerate}

\begin{itemize}
\tightlist
\item
  \(D\) and \(D^{-1}\) are the explicit kernels of \(v\) with penalty
  \(\kappa\), and \(H^{-1}=D^{-1}+\lambda_j(d_2D^{-1})^*(d_2D^{-1})\)
  (Lemmas 4.4--4.7);
\item
  all five kernels lie in \(\mathcal S(\alpha_j,\Omega_j)\);
\item
  on every torus: \(DD^{-1}=HH^{-1}=1\), \(c\,1+R^*R=H\),
  \(\hat H(0)=1\), \(Hd=Dd\), and \(H\) and \(D\) are symmetric;
\item
  on every torus:
  \(\eta_j|\omega|^2\le\langle\omega,H\omega\rangle\le\Omega_j|\omega|^2\),
  and the same for \(D\);
\item
  for every gauge group \(G\), the form \(\langle A,d^*Hd\,A\rangle\) on
  \(\mathfrak g\)-valued 1-forms is invariant under the signed lattice
  symmetries;
\item
  \(D=1\) if \(v\) is the empty history.
\end{itemize}

\begin{enumerate}
\def\labelenumi{(\alph{enumi})}
\setcounter{enumi}{1}
\tightlist
\item
  \emph{One step.} Let \(j,t\ge4\) with \(t\in\{j,j-1\}\), and let \(P\)
  be a flight of index \((j,t)\): odd block size \(L\ge3\) and fine
  history of index \(j\). Then:
\end{enumerate}

\begin{itemize}
\tightlist
\item
  the explicit transfer kernels \(W,T\) of \(P\) satisfy
  \(\nabla_uW\in\mathcal S(\alpha_j,\Omega_jL^{-1-|u|})\) and
  \(\nabla_uT\in\mathcal S(\alpha_j,\Omega_jL^{-2-|u|})\), for every
  word \(u\) of at most \(j-2\) fine forward differences;
\item
  the rolled 3-form kernel \(\mathcal P\) of \(P\), built with fine
  amendment \(\lambda_j\) and coarse amendment \(\lambda_t\) (Lemma
  4.12), lies in \(\mathcal S(\alpha_j,\Omega_j)\);
\item
  if \(L\ge b_j\), then at every real momentum \(\hat{\mathcal P}(p)\)
  is Hermitian and \(\eta_j\le\hat{\mathcal P}(p)\le\Omega_j\).
\end{itemize}

\emph{Status:} theorem, with no hypothesis beyond the kernel choice. It
uses III.6's square-root theorem to construct \(R\). The catalogue does
not depend on \(G\). \emph{Lean:}
\href{https://github.com/yang-mills-group/OpenAI-YM-20261028/blob/3456e5e75cf8d4cd558955a5c7638fba2d68f464/YM/Next/FREE_KERNELS/Priced.lean\#L29-L33}{\texttt{Catalog}},
\href{https://github.com/yang-mills-group/OpenAI-YM-20261028/blob/3456e5e75cf8d4cd558955a5c7638fba2d68f464/YM/Next/FREE_KERNELS/Priced.lean\#L67-L93}{\texttt{pledges}},
\href{https://github.com/yang-mills-group/OpenAI-YM-20261028/blob/3456e5e75cf8d4cd558955a5c7638fba2d68f464/YM/Next/FREE_KERNELS/Priced.lean\#L22-L24}{\texttt{BottomQuote}},
\href{https://github.com/yang-mills-group/OpenAI-YM-20261028/blob/3456e5e75cf8d4cd558955a5c7638fba2d68f464/YM/Next/FREE_KERNELS/Priced.lean\#L25-L28}{\texttt{HopQuote}},
\href{https://github.com/yang-mills-group/OpenAI-YM-20261028/blob/3456e5e75cf8d4cd558955a5c7638fba2d68f464/YM/Next/FREE_KERNELS/Priced.lean\#L20}{\texttt{Mag}},
\href{https://github.com/yang-mills-group/OpenAI-YM-20261028/blob/3456e5e75cf8d4cd558955a5c7638fba2d68f464/YM/Next/FREE_KERNELS/Ground.lean\#L54-L58}{\texttt{Defined}},
\href{https://github.com/yang-mills-group/OpenAI-YM-20261028/blob/3456e5e75cf8d4cd558955a5c7638fba2d68f464/YM/PREC/Stacks.lean\#L61-L68}{\texttt{Floor}},
\href{https://github.com/yang-mills-group/OpenAI-YM-20261028/blob/3456e5e75cf8d4cd558955a5c7638fba2d68f464/YM/Bid/Craft.lean\#L21-L22}{\texttt{Donate}},
\href{https://github.com/yang-mills-group/OpenAI-YM-20261028/blob/3456e5e75cf8d4cd558955a5c7638fba2d68f464/YM/Bid/Craft.lean\#L38-L40}{\texttt{Funds}}.

\textbf{Proposition III.3.2 (free kernels along an admissible
schedule).} Let \(G\) be a gauge group and \(K\) a kernel choice. There
is a catalogue \((\alpha,\Omega,\eta,\lambda,\gamma,b)\) as in III.3.1
with the following property. For every sheet \(\mathfrak S\) with kernel
\(K\) (the pre-\(L\) constants, III.2) there is a second table
\((\alpha'_j,\Omega'_j,\eta_j)\) such that the conclusion below holds.

Let \(A\) be an RG trajectory with kernel \(K\) and the stage profile of
\(\mathfrak S\), satisfying the \emph{credit condition}:

\begin{itemize}
\tightlist
\item
  (C1) \(c_i\le\gamma_{J(i)}\) for every level \(i\);
\item
  (C2) \(b_{J(i+1)}\le L_i\) for every step \(i\).
\end{itemize}

Then, simultaneously for all \(N\ge0\) and \(0\le i\le N\), there are
two families of metric data and transfer kernels: a pre-repair family
\(\mathsf M=(\mathcal B^{N,i},\mathcal U^{N,i})\) and an output family
\(\mathsf H=(\mathcal B^{N,i},\tilde{\mathcal U}^{N,i})\). They have the
same metric data and the same \(W,T\), and satisfy:

\begin{enumerate}
\def\labelenumi{\arabic{enumi}.}
\tightlist
\item
  \emph{Explicit formulas.} \(\mathcal B^{N,i}\) is the metric data of
  III.3.1(a) for the history of level \(i\) in mesh \(N\), with
  amendment \(\lambda_{J(i)}\). That history is the composite of the
  stage hops of steps \(N-1,\dots,i\), starting from the empty history
  at level \(N\). \(W^{N,i},T^{N,i}\) are the explicit transfer kernels
  of the flight of step \(i\), of index \((J(i+1),J(i))\).
\item
  \emph{Pre-repair conditions (\(\mathsf M\)).}

  \begin{itemize}
  \tightlist
  \item
    \(\mathcal B^{N,i}\in\mathcal S(\alpha_{J(i)},\Omega_{J(i)})\), with
    lower and upper bounds \(\eta_{J(i)},\Omega_{J(i)}\) on every torus
    for the clip of level \(i\), and the symmetry;
  \item
    \(D^{N,N}=1\);
  \item
    \(W,T\) have the derivative bounds of III.3.1(b) at stage
    \(J(i+1)\), with \(J(i+1)-2\) differences;
  \item
    (D1)--(D4) hold on every schedule torus pair of step \(i\);
  \item
    there is a 3-form kernel
    \(P\in\mathcal S(\alpha_{J(i+1)},\Omega_{J(i+1)})\), Hermitian at
    real momenta, with
    \(\langle v,\hat P(p)v\rangle\ge\eta_{J(i+1)}|v|^2\) for real \(p\)
    and \(v\) in the range of \(\sigma_2(p)\), such that (Δ) holds on
    every schedule torus pair.
  \end{itemize}
\item
  \emph{Output conditions (\(\mathsf H\)), with the second table.}

  \begin{itemize}
  \tightlist
  \item
    the metric bounds of III.3.1(a), with strip constants
    \(\alpha'_j,\Omega'_j\) and the same lower bounds \(\eta_j\);
  \item
    \(\tilde M_n,\tilde O\in\mathcal S(\alpha'_{J(i+1)},\Omega'_{J(i+1)})\)
    and
    \(\tilde M_f\in\mathcal S(\alpha'_{J(i+1)},\Omega'_{J(i+1)}L_i^{-2})\),
    and the derivative bounds on \(W,T\) with \(\alpha'_j,\Omega'_j\);
  \item
    (D1)--(D4) and (R1)--(R6) on every schedule torus pair, with the
    clips of levels \(i+1,i\);
  \item
    \(D^{N,N}=1\), and the symmetry of every \(\mathcal B^{N,i}\).
  \end{itemize}
\end{enumerate}

\emph{Status:} theorem. The inputs F06--F09 (III.6) are discharged
inside the proof. For trajectories with the stated kernel and profile,
the additional hypothesis is the credit condition. The variant
\href{https://github.com/yang-mills-group/OpenAI-YM-20261028/blob/3456e5e75cf8d4cd558955a5c7638fba2d68f464/YM/Next/FREE_KERNELS/Calendar.lean\#L108-L124}{\texttt{before\_paid}}
(assigned to III.4) replaces it by:

\begin{itemize}
\tightlist
\item
  two conditions on the sheet: clip coefficients \(\le\gamma_j\), and
  block-size floors \(\ge b_j\), for all stages from the first one on;
\item
  the schedule's payment record (III.2), which gives \(L_i\ge\) floor.
\end{itemize}

III.9 builds the sheet so that these hold: its floors are exactly \(b\)
and its clip coefficients are \(\le\gamma\)
(\href{https://github.com/yang-mills-group/OpenAI-YM-20261028/blob/3456e5e75cf8d4cd558955a5c7638fba2d68f464/YM/PhysicalNext/UNMASKED_THROUGH_LAUNCH/Dome.lean\#L90-L101}{\texttt{housing}}).
The variant also lets the second table be computed from a preliminary
sheet \(\mathfrak S\) and used for any later sheet \(\mathfrak S'\),
provided \(\mathfrak S'\) has the same direct clip coefficients, the
same kernel, a drift at least as large, and logarithmic parameter size
at least as large.

\emph{Lean:}
\href{https://github.com/yang-mills-group/OpenAI-YM-20261028/blob/3456e5e75cf8d4cd558955a5c7638fba2d68f464/YM/Next/FREE_KERNELS/Rack.lean\#L134-L145}{\texttt{retained\_support}};
conditions
\href{https://github.com/yang-mills-group/OpenAI-YM-20261028/blob/3456e5e75cf8d4cd558955a5c7638fba2d68f464/YM/Next/FREE_KERNELS/Rack.lean\#L82-L85}{\texttt{Present}},
\href{https://github.com/yang-mills-group/OpenAI-YM-20261028/blob/3456e5e75cf8d4cd558955a5c7638fba2d68f464/YM/Bid/Craft.lean\#L63-L73}{\texttt{Craft.Rack}},
\href{https://github.com/yang-mills-group/OpenAI-YM-20261028/blob/3456e5e75cf8d4cd558955a5c7638fba2d68f464/YM/Bid/Craft.lean\#L104-L105}{\texttt{Bid.Retained}},
\href{https://github.com/yang-mills-group/OpenAI-YM-20261028/blob/3456e5e75cf8d4cd558955a5c7638fba2d68f464/YM/PREC/Stacks.lean\#L132-L139}{\texttt{Horizon.Suffices}},
\href{https://github.com/yang-mills-group/OpenAI-YM-20261028/blob/3456e5e75cf8d4cd558955a5c7638fba2d68f464/YM/Bid/Station.lean\#L54-L55}{\texttt{Bid.Sworn}},
\href{https://github.com/yang-mills-group/OpenAI-YM-20261028/blob/3456e5e75cf8d4cd558955a5c7638fba2d68f464/YM/Next/FREE_KERNELS/Rack.lean\#L87-L88}{\texttt{Credit}};
discharge
\href{https://github.com/yang-mills-group/OpenAI-YM-20261028/blob/3456e5e75cf8d4cd558955a5c7638fba2d68f464/YM/Next/FREE_KERNELS/Calendar.lean\#L95-L102}{\texttt{credit\_from\_sheet}},
\href{https://github.com/yang-mills-group/OpenAI-YM-20261028/blob/3456e5e75cf8d4cd558955a5c7638fba2d68f464/YM/Next/FREE_KERNELS/Calendar.lean\#L28}{\texttt{Willing}}.

\textbf{Proposition III.3.3 (identities on every torus).} Let a family
satisfy the output conditions of III.3.2(3) for some table. For every
\(N\), \(m<N\) and positive sides \(c,f\) with \(f=cL_m\), (D1)--(D4)
and (R1)--(R6) hold, with the same kernels and the clips of levels
\(m+1,m\). The diagonal identities follow by rigidity (III.10,
\href{https://github.com/yang-mills-group/OpenAI-YM-20261028/blob/3456e5e75cf8d4cd558955a5c7638fba2d68f464/YM/PhysicalNext/RG_NEXT_COMPLETE_REPEATABLE_NONLINEAR_STEP/Pane/Purl.lean\#L20-L53}{\texttt{panes}}).
They also follow from the explicit formulas when the family realises
them
(\href{https://github.com/yang-mills-group/OpenAI-YM-20261028/blob/3456e5e75cf8d4cd558955a5c7638fba2d68f464/YM/UVNext/PRICED_KERNELS/Tori.lean\#L46-L60}{\texttt{universal\_diagonal}}).
Given the diagonal identities, the roll identities follow from the
output conditions by the same rigidity argument. \emph{Status:} theorem.
\emph{Lean:}
\href{https://github.com/yang-mills-group/OpenAI-YM-20261028/blob/3456e5e75cf8d4cd558955a5c7638fba2d68f464/YM/UVNext/PRICED_KERNELS/Tori.lean\#L109-L157}{\texttt{complete\_roll}},
\href{https://github.com/yang-mills-group/OpenAI-YM-20261028/blob/3456e5e75cf8d4cd558955a5c7638fba2d68f464/YM/PREC/Stacks.lean\#L82-L86}{\texttt{Towers.Diagonal}},
\href{https://github.com/yang-mills-group/OpenAI-YM-20261028/blob/3456e5e75cf8d4cd558955a5c7638fba2d68f464/YM/PREC/Stacks.lean\#L88-L96}{\texttt{Towers.Roll}},
\href{https://github.com/yang-mills-group/OpenAI-YM-20261028/blob/3456e5e75cf8d4cd558955a5c7638fba2d68f464/YM/PREC/TreeMatch.lean\#L76-L79}{\texttt{Inputs}}.
Used in III.10.

\textbf{Proposition III.3.4 (positivity of the rolled 3-form kernel).}
Fix a kernel choice \(K\) and \(j,t\ge4\). There are \(C_0,B>0\),
depending only on \((K,j,t)\), with the following property. For every
coarse amendment \(\lambda\ge C_0\) and fine amendment \(\mu\ge0\) there
are \(\eta,\sigma,\Omega>0\) such that, for every flight \(P\) of index
\((j,t)\):

\begin{itemize}
\tightlist
\item
  \(\mathcal P_{\mu,\lambda}(P)\in\mathcal S(\sigma,\Omega)\);
\item
  if \(\mu B\le L^2\), then at every real \(p\),
  \(\hat{\mathcal P}_{\mu,\lambda}(p)\) is Hermitian and
  \(\eta\le\hat{\mathcal P}_{\mu,\lambda}(p)\le\Omega\).
\end{itemize}

\emph{Status:} theorem. \emph{Lean:}
\href{https://github.com/yang-mills-group/OpenAI-YM-20261028/blob/3456e5e75cf8d4cd558955a5c7638fba2d68f464/YM/Next/FREE_KERNELS/Eject.lean\#L99-L140}{\texttt{Flight.frontier}},
\href{https://github.com/yang-mills-group/OpenAI-YM-20261028/blob/3456e5e75cf8d4cd558955a5c7638fba2d68f464/YM/Next/FREE_KERNELS/Eject.lean\#L58-L95}{\texttt{Flight.absorb\_before}}.

\textbf{Proposition III.3.5 (line and box filters).} For positive
integers \(L,M\), define the uncentred line average below; for odd
\(L,M\), also define the centred line and box averages:
\[\mathrm{plank}_M(z)=M^{-1}\sum_{|s|\le(M-1)/2}e^{isz},\qquad \mathrm{rail}_M(z)=M^{-1}\sum_{s=0}^{M-1}e^{isz},\qquad a_L(p)=\prod_a\mathrm{plank}_L(p_a)\]
(centred line average, uncentred line average, box average). Then:

\begin{enumerate}
\def\labelenumi{(\alph{enumi})}
\item
  \(\mathrm{rail}_L(Mz)\,\mathrm{rail}_M(z)=\mathrm{rail}_{LM}(z)\) for
  all positive \(L,M\), and
  \(\mathrm{plank}_L(Mz)\,\mathrm{plank}_M(z)=\mathrm{plank}_{LM}(z)\)
  for odd \(L,M\). For odd \(M\), \(\mathrm{plank}_M\) is even,
  \(2\pi\)-periodic and real on \(\mathbb R\).
\item
  \(|\mathrm{plank}_M(z)|,|\mathrm{rail}_M(z)|\le e^{M|\mathrm{Im}\,z|}\).
\item
  There are absolute constants \(c_*,c_1>0\) such that
  \(|\mathrm{rail}_M(z)|\ge c_*\) for complex \(z\) with
  \(M|z|\le3\pi/2\), and \(|\mathrm{plank}_M(z)|\ge e^{-c_1M/L}\) for
  real \(z\), \(M\le L\) and \(L|z|\le3\pi/2\).
\item
  There is an absolute \(B\) such that
  \(|\mathrm{rail}_L((q_\ell)_a)|,|\mathrm{plank}_L((q_\ell)_a)|\le B/(1+|\ell_a|)\).
  This holds for odd \(L\), every \(p\) in the complex cube
  \(\{|\mathrm{Re}\,p_a|<\tfrac{5\pi}4,\ |\mathrm{Im}\,p_a|<\tfrac\pi4\}\),
  every alias \(\ell\) and every \(a\).
\end{enumerate}

\emph{Status:} theorem. \emph{Lean:}
\href{https://github.com/yang-mills-group/OpenAI-YM-20261028/blob/3456e5e75cf8d4cd558955a5c7638fba2d68f464/YM/Next/FREE_KERNELS/Lines.lean\#L38}{\texttt{plank}},
\href{https://github.com/yang-mills-group/OpenAI-YM-20261028/blob/3456e5e75cf8d4cd558955a5c7638fba2d68f464/YM/Next/FREE_KERNELS/Lines.lean\#L34}{\texttt{rail}},
\href{https://github.com/yang-mills-group/OpenAI-YM-20261028/blob/3456e5e75cf8d4cd558955a5c7638fba2d68f464/YM/Next/FREE_KERNELS/Train.lean\#L18}{\texttt{area}},
\href{https://github.com/yang-mills-group/OpenAI-YM-20261028/blob/3456e5e75cf8d4cd558955a5c7638fba2d68f464/YM/Next/FREE_KERNELS/Lines.lean\#L68-L76}{\texttt{plank\_bind}},
\href{https://github.com/yang-mills-group/OpenAI-YM-20261028/blob/3456e5e75cf8d4cd558955a5c7638fba2d68f464/YM/Next/FREE_KERNELS/Lines.lean\#L51-L62}{\texttt{rail\_bind}},
\href{https://github.com/yang-mills-group/OpenAI-YM-20261028/blob/3456e5e75cf8d4cd558955a5c7638fba2d68f464/YM/Next/FREE_KERNELS/Lines.lean\#L127-L133}{\texttt{plank\_height}},
\href{https://github.com/yang-mills-group/OpenAI-YM-20261028/blob/3456e5e75cf8d4cd558955a5c7638fba2d68f464/YM/Next/FREE_KERNELS/Lines.lean\#L140-L152}{\texttt{rail\_lower}},
\href{https://github.com/yang-mills-group/OpenAI-YM-20261028/blob/3456e5e75cf8d4cd558955a5c7638fba2d68f464/YM/Next/FREE_KERNELS/Lines.lean\#L162-L196}{\texttt{principal\_plank}},
\href{https://github.com/yang-mills-group/OpenAI-YM-20261028/blob/3456e5e75cf8d4cd558955a5c7638fba2d68f464/YM/Next/FREE_KERNELS/Lines.lean\#L200-L245}{\texttt{alias\_rail}}.

\textbf{Proposition III.3.6 (free covariance of the curvature).} Call a
real 2-form kernel \(f\) \emph{firm} with constants
\((\alpha,\Omega,\eta)\) if \(f\in\mathcal S(\alpha,\Omega)\),
\(\hat f(0)=1\), and
\(\eta|\omega|^2\le\langle\omega,f_n\omega\rangle\le\Omega|\omega|^2\)
for every real cochain on every torus. Symmetry is not required. Fix
\(\alpha,\eta>0\) and \(\Omega\in\mathbb R\), and let
\(U=\{p\in\mathbb R^4\setminus\{0\}:|p_a|<3\pi/2\}\). For \(p\in U\)
define:

\begin{itemize}
\tightlist
\item
  \(\Phi_f(p)=\sigma_1(p)^*\hat f(p)\sigma_1(p)+\sigma_0(p)\sigma_0(p)^*\),
  the symbol of \(d_1^*fd_1+d_0d_0^*\) on 1-forms;
\item
  \(G_f(p)=\sigma_1(p)\Phi_f(p)^{-1}\sigma_1(p)^*\), the curvature
  covariance symbol when \(f\) is symmetric and the gauge-fixed
  potential has precision \(d_1^*fd_1+d_0d_0^*\);
\item
  \(G_\infty(p)=s(p)|p|^{-2}s(p)^*\), where \(s(p)\) is wedging with
  \(ip\) on 1-forms: the continuum Feynman-gauge curvature symbol.
\end{itemize}

\begin{enumerate}
\def\labelenumi{(\alph{enumi})}
\item
  There is \(b>0\), depending only on \(\alpha,\Omega,\eta\), such that
  \(\mathrm{Re}\,\langle x,\Phi_f(p)x\rangle\ge b|p|^2|x|^2\) for every
  firm \(f\), \(p\in U\) and complex 1-form \(x\).
\item
  For every \(n\in\mathbb N\) there is \(C_n\ge0\), depending only on
  \(n,\alpha,\Omega,\eta\), such that, for all firm \(f\) and
  \(p\in U\), \(\|D^nG_f(p)\|\le C_n|p|^{-n}\) and
  \(\|D^n(G_f-G_\infty)(p)\|\le C_n|p|^{1-n}\).
\item
  If \(f\) is a scalar kernel with
  \(\sum_xe^{\alpha|x|_1}|f(x)|\le\Omega\), then \(\hat f\) is
  \(C^\infty\) on \(\mathbb R^4\) and \(\|D^n\hat f(p)\|\le A_n\Omega\)
  for every \(p\), where \(A_0=1\) and \(A_n=(n/\alpha)^n\) for
  \(n\ge1\). Here \(D^n\) is the full \(n\)th Fréchet derivative, with
  its multilinear operator norm; this notation also applies to (b).
\end{enumerate}

\emph{Status:} theorems. Part (b) is used in III.11. \emph{Lean:}
\href{https://github.com/yang-mills-group/OpenAI-YM-20261028/blob/3456e5e75cf8d4cd558955a5c7638fba2d68f464/YM/Further/FREE_COVARIANCE_KP/KP/Core.lean\#L115-L125}{\texttt{tank\_full}},
\href{https://github.com/yang-mills-group/OpenAI-YM-20261028/blob/3456e5e75cf8d4cd558955a5c7638fba2d68f464/YM/Further/FREE_COVARIANCE_KP/KP/Alloy.lean\#L110-L120}{\texttt{gem\_quote}},
\href{https://github.com/yang-mills-group/OpenAI-YM-20261028/blob/3456e5e75cf8d4cd558955a5c7638fba2d68f464/YM/Further/FREE_COVARIANCE_KP/KP/Plain.lean\#L101-L133}{\texttt{radial\_prices}},
\href{https://github.com/yang-mills-group/OpenAI-YM-20261028/blob/3456e5e75cf8d4cd558955a5c7638fba2d68f464/YM/Further/FREE_COVARIANCE_KP/KP/Floor.lean\#L17-L21}{\texttt{Firm}},
\href{https://github.com/yang-mills-group/OpenAI-YM-20261028/blob/3456e5e75cf8d4cd558955a5c7638fba2d68f464/YM/Further/FREE_COVARIANCE_KP/KP/Core.lean\#L62}{\texttt{fuel}},
\href{https://github.com/yang-mills-group/OpenAI-YM-20261028/blob/3456e5e75cf8d4cd558955a5c7638fba2d68f464/YM/Further/FREE_COVARIANCE_KP/KP/Alloy.lean\#L15}{\texttt{garnet}},
\href{https://github.com/yang-mills-group/OpenAI-YM-20261028/blob/3456e5e75cf8d4cd558955a5c7638fba2d68f464/YM/Further/FREE_COVARIANCE_KP/KP/Alloy.lean\#L17}{\texttt{pearl}}.

\textbf{Proposition III.3.7 (calculus lemmas for the nonlinear step).}
These do not concern free kernels; they are consumed in III.10 and
III.2.

\begin{enumerate}
\def\labelenumi{(\alph{enumi})}
\item
  \emph{Discrete Taylor expansion with scaled bounds.} Let
  \(b,k\in\mathbb N\) with \(k\le b\). There is \(C=C(b,k)\ge1\) with
  the following property. Let \(L\ge1\), \(M\ge0\), and let
  \(p:\mathbb Z^4\to E\) (\(E\) a real normed vector space) satisfy
  \(\|\nabla_up(x)\|\le ML^{-1-|u|}(1+L^{-1}|x|_1)^b\) for all words
  \(u\) of forward differences with \(|u|\le k\). Then the Newton pieces
  \(P_i(x)=\sum_{|\nu|=i}\binom{x}{\nu}\nabla^\nu p(0)\), of degree at
  most \(i\), satisfy \(\|P_i(x)\|\le CML^{-1-i}(1+|x|_1)^b\) for
  \(i<k\), and the remainder satisfies
  \(\|p(x)-\sum_{i<k}P_i(x)\|\le CML^{-1-k}(1+|x|_1)^b\).
\item
  \emph{Weighted Taylor coefficients.} Let \(d\in\mathbb N\) and let
  \(f:E\to F\) be \(C^d\) at \(0\) between real normed vector spaces.
  Let \(p_v:E\to E\) (\(v\in V\) finite) be continuous linear maps with
  \(\sum_vp_v=1\), and let \(Y_v\ge1\) be integers. Suppose that for
  every \(z\) the function \(t\mapsto f(\sum_vt^{Y_v}p_vz)\) has
  vanishing derivatives of orders \(0,\ldots,d\) at \(0\). Then
  \(D^mf(0)(x_1,\dots,x_m)=0\) whenever there are integers \(k_i\ge1\)
  with \(\sum_ik_i\le d\) and \(p_v(x_i)=0\) for every \(v\) with
  \(Y_v>k_i\).
\item
  \emph{Moments of local gauge-invariant functions.} Fix
  \(b,d\in\mathbb N\) with \(b\le d\), a torus \(\Lambda_n\) with
  \(n\ge2\), an anchor \(o\), and a finite support \(X\) satisfying
  \(10\,\mathrm{dist}(x,o)<n\) for every \(x\in X\), where
  \(\mathrm{dist}\) is the torus \(\ell^\infty\) distance. Let
  \(u:\mathrm{Config}(G,\Lambda_n)\to\mathbb C\) be gauge invariant and
  depend only on links whose two endpoints lie in \(X\), with
  \(u\circ\exp\) of class \(C^b\) at \(0\). For every radial polynomial
  connection jet \(J\) through weight \(d\), let \(U_{J,o}(t)\) be its
  link holonomies at scale \(t\), based at \(o\): a homogeneous
  connection component of degree \(r-1\) has weight \(r\),
  \(2\le r\le d\), and is multiplied by \(t^r\). Assume
  \(\frac{d^k}{dt^k}u(U_{J,o}(t))|_{t=0}=0\) for every \(J\) and
  \(0\le k\le b\). For any \(m\in\mathbb N\), \(p_1,\ldots,p_m\ge1\)
  with \(\sum_i p_i\le b\), and link fields \(w_i\) whose values on
  those links are polynomials of degree at most \(p_i-1\) in the centred
  lift of the initial site relative to \(o\), one has
  \(D^m(u\circ\exp)(0)(w_1,\ldots,w_m)=0\). Values of \(w_i\) on other
  links are unrestricted.
\item
  \emph{Radial decomposition.} Let \(j\in\mathbb N\) and let
  \(Q_a:\mathbb R^4\to\mathfrak g\) (\(a=1,\dots,4\)) be homogeneous
  polynomials of degree \(j+1\). Then \(Q_a=K_a+\partial_ar\), where
  \(r=(j+2)^{-1}\sum_ax_aQ_a\) is homogeneous of degree \(j+2\), and
  \(K\) is homogeneous of degree \(j+1\) with \(\sum_ax_aK_a=0\) (the
  Fock--Schwinger condition).
\item
  \emph{Small-field condition along rays.} There is \(C=C(G)>0\) with
  the following property. Let \(0\le b\le1\), \(\tau\ge Cb^2\), and let
  \(v\) be Lie-algebra link variables with \(|v_x|\le b\) on all links
  within distance \(5\) of \(Y\). If every plaquette touching \(Y\) has
  holonomy within chord distance \(\tau\) of \(1\) for the configuration
  \(\exp(v)\), then the same plaquette bounds hold for \(\exp(sv)\),
  \(0\le s\le1\).
\end{enumerate}

\emph{Status:} theorems. \emph{Lean:}
\href{https://github.com/yang-mills-group/OpenAI-YM-20261028/blob/3456e5e75cf8d4cd558955a5c7638fba2d68f464/YM/MultiNext/SIGNED_SOURCE_CONTINUATION/TownHall.lean\#L31-L90}{\texttt{capital}},
\href{https://github.com/yang-mills-group/OpenAI-YM-20261028/blob/3456e5e75cf8d4cd558955a5c7638fba2d68f464/YM/MultiNext/SIGNED_SOURCE_CONTINUATION/Weighted.lean\#L120-L147}{\texttt{cancel}},
\href{https://github.com/yang-mills-group/OpenAI-YM-20261028/blob/3456e5e75cf8d4cd558955a5c7638fba2d68f464/YM/MultiNext/SIGNED_SOURCE_CONTINUATION/Weighted.lean\#L84-L115}{\texttt{converse}},
\href{https://github.com/yang-mills-group/OpenAI-YM-20261028/blob/3456e5e75cf8d4cd558955a5c7638fba2d68f464/YM/MultiNext/SIGNED_SOURCE_CONTINUATION/PMOM.lean\#L43-L61}{\texttt{wiped}},
\href{https://github.com/yang-mills-group/OpenAI-YM-20261028/blob/3456e5e75cf8d4cd558955a5c7638fba2d68f464/YM/MultiNext/SIGNED_SOURCE_CONTINUATION/Averaging.lean\#L84-L114}{\texttt{locale}},
\href{https://github.com/yang-mills-group/OpenAI-YM-20261028/blob/3456e5e75cf8d4cd558955a5c7638fba2d68f464/YM/MultiNext/SIGNED_SOURCE_CONTINUATION/Safeguard.lean\#L148-L177}{\texttt{strings}}.

\subsubsection{2. What it uses}\label{what-it-uses-5}

\textbf{From other units.}

\begin{itemize}
\tightlist
\item
  \textbf{I.1.} The lattice, the signed (hypercubic) symmetries, and the
  gauge group. In the free-kernel construction the gauge group enters
  only the symmetry statement; III.3.7 also uses its exponential and
  holonomy maps.
\item
  \textbf{III.1.} The block-spin data:

  \begin{itemize}
  \tightlist
  \item
    the block ratio, the patch (odd \(L\ge3\), \(J\ge4\)), the kernel
    choices
    \href{https://github.com/yang-mills-group/OpenAI-YM-20261028/blob/3456e5e75cf8d4cd558955a5c7638fba2d68f464/YM/RG/Patch.lean\#L37-L45}{\texttt{KernelChoices}}
    and the clips
    \href{https://github.com/yang-mills-group/OpenAI-YM-20261028/blob/3456e5e75cf8d4cd558955a5c7638fba2d68f464/YM/RG/PathActions.lean\#L139-L146}{\texttt{Clip}};
  \item
    the linearised average \(Q\)
    (\href{https://github.com/yang-mills-group/OpenAI-YM-20261028/blob/3456e5e75cf8d4cd558955a5c7638fba2d68f464/YM/PREC/FrameResponse.lean\#L65-L66}{\texttt{mean}})
    and the lattice exterior derivative;
  \item
    the trajectory record
    \href{https://github.com/yang-mills-group/OpenAI-YM-20261028/blob/3456e5e75cf8d4cd558955a5c7638fba2d68f464/YM/RG/History.lean\#L131-L134}{\texttt{Levels}}:
    block sizes, stage indices with \(J(i)\le J(i+1)\le J(i)+1\), and
    schedule tori whose sides tend to infinity.
  \end{itemize}

  III.1 also defines the output objects:

  \begin{itemize}
  \tightlist
  \item
    metric data
    \href{https://github.com/yang-mills-group/OpenAI-YM-20261028/blob/3456e5e75cf8d4cd558955a5c7638fba2d68f464/YM/PREC/Stacks.lean\#L35-L40}{\texttt{Bed}},
    \href{https://github.com/yang-mills-group/OpenAI-YM-20261028/blob/3456e5e75cf8d4cd558955a5c7638fba2d68f464/YM/PREC/Stacks.lean\#L49-L51}{\texttt{Bed.full}},
    \href{https://github.com/yang-mills-group/OpenAI-YM-20261028/blob/3456e5e75cf8d4cd558955a5c7638fba2d68f464/YM/PREC/Stacks.lean\#L61-L68}{\texttt{Floor}};
  \item
    transfer kernels
    \href{https://github.com/yang-mills-group/OpenAI-YM-20261028/blob/3456e5e75cf8d4cd558955a5c7638fba2d68f464/YM/PREC/Stacks.lean\#L70-L75}{\texttt{Towers}}
    and their bounds
    \href{https://github.com/yang-mills-group/OpenAI-YM-20261028/blob/3456e5e75cf8d4cd558955a5c7638fba2d68f464/YM/PREC/Stacks.lean\#L111-L114}{\texttt{Towers.Deep}};
  \item
    the identities (D1)--(D4), (R1)--(R6), with \(d^*Hd\) as
    \href{https://github.com/yang-mills-group/OpenAI-YM-20261028/blob/3456e5e75cf8d4cd558955a5c7638fba2d68f464/YM/PREC/TreeMatch.lean\#L68-L69}{\texttt{K}};
  \item
    kernel families
    \href{https://github.com/yang-mills-group/OpenAI-YM-20261028/blob/3456e5e75cf8d4cd558955a5c7638fba2d68f464/YM/PREC/Stacks.lean\#L119-L121}{\texttt{Horizon}},
    the output conditions
    \href{https://github.com/yang-mills-group/OpenAI-YM-20261028/blob/3456e5e75cf8d4cd558955a5c7638fba2d68f464/YM/PREC/Stacks.lean\#L132-L139}{\texttt{Horizon.Suffices}},
    and price tables
    \href{https://github.com/yang-mills-group/OpenAI-YM-20261028/blob/3456e5e75cf8d4cd558955a5c7638fba2d68f464/YM/PREC/Stacks.lean\#L123-L128}{\texttt{MapPrices}}.
  \end{itemize}
\item
  \textbf{III.2.} The sheet
  \href{https://github.com/yang-mills-group/OpenAI-YM-20261028/blob/3456e5e75cf8d4cd558955a5c7638fba2d68f464/YM/Price/Intent.lean\#L26-L35}{\texttt{Sheet}},
  the block sizes
  \href{https://github.com/yang-mills-group/OpenAI-YM-20261028/blob/3456e5e75cf8d4cd558955a5c7638fba2d68f464/YM/Price/Intent.lean\#L50-L53}{\texttt{Rolls}},
  and the payment record
  \href{https://github.com/yang-mills-group/OpenAI-YM-20261028/blob/3456e5e75cf8d4cd558955a5c7638fba2d68f464/YM/Price/Intent.lean\#L145-L162}{\texttt{Paid}}.
\item
  \textbf{III.6.} The four linear-algebra theorems
  \href{https://github.com/yang-mills-group/OpenAI-YM-20261028/blob/3456e5e75cf8d4cd558955a5c7638fba2d68f464/YM/Work/F06.lean\#L64-L128}{\texttt{F06}},
  \href{https://github.com/yang-mills-group/OpenAI-YM-20261028/blob/3456e5e75cf8d4cd558955a5c7638fba2d68f464/YM/Work/F07.lean\#L180-L201}{\texttt{F07}},
  \href{https://github.com/yang-mills-group/OpenAI-YM-20261028/blob/3456e5e75cf8d4cd558955a5c7638fba2d68f464/YM/Work/F08.lean\#L53-L146}{\texttt{F08}},
  \href{https://github.com/yang-mills-group/OpenAI-YM-20261028/blob/3456e5e75cf8d4cd558955a5c7638fba2d68f464/YM/Work/F09.lean\#L41-L150}{\texttt{F09}}.
  Their statements are in Lemma 4.16. F07 is also used in Lemma 4.8.
\item
  \textbf{III.10.}

  \begin{itemize}
  \tightlist
  \item
    The interface conditions
    \href{https://github.com/yang-mills-group/OpenAI-YM-20261028/blob/3456e5e75cf8d4cd558955a5c7638fba2d68f464/YM/Bid/Craft.lean\#L63-L73}{\texttt{Craft.Rack}},
    \href{https://github.com/yang-mills-group/OpenAI-YM-20261028/blob/3456e5e75cf8d4cd558955a5c7638fba2d68f464/YM/Bid/Craft.lean\#L32}{\texttt{Bleeps}},
    \href{https://github.com/yang-mills-group/OpenAI-YM-20261028/blob/3456e5e75cf8d4cd558955a5c7638fba2d68f464/YM/Bid/Craft.lean\#L34-L36}{\texttt{Reaches}},
    \href{https://github.com/yang-mills-group/OpenAI-YM-20261028/blob/3456e5e75cf8d4cd558955a5c7638fba2d68f464/YM/Bid/Craft.lean\#L38-L40}{\texttt{Funds}},
    \href{https://github.com/yang-mills-group/OpenAI-YM-20261028/blob/3456e5e75cf8d4cd558955a5c7638fba2d68f464/YM/Bid/Craft.lean\#L21-L22}{\texttt{Donate}},
    \href{https://github.com/yang-mills-group/OpenAI-YM-20261028/blob/3456e5e75cf8d4cd558955a5c7638fba2d68f464/YM/Bid/Craft.lean\#L53-L58}{\texttt{Replaces}},
    \href{https://github.com/yang-mills-group/OpenAI-YM-20261028/blob/3456e5e75cf8d4cd558955a5c7638fba2d68f464/YM/Bid/Craft.lean\#L104-L105}{\texttt{Bid.Retained}},
    \href{https://github.com/yang-mills-group/OpenAI-YM-20261028/blob/3456e5e75cf8d4cd558955a5c7638fba2d68f464/YM/Bid/Station.lean\#L54-L55}{\texttt{Bid.Sworn}}.
  \item
    The repair theorem
    \href{https://github.com/yang-mills-group/OpenAI-YM-20261028/blob/3456e5e75cf8d4cd558955a5c7638fba2d68f464/YM/Bid/Craft.lean\#L107-L166}{\texttt{Bid.pinning}}
    (Lemma 4.16), and
    \href{https://github.com/yang-mills-group/OpenAI-YM-20261028/blob/3456e5e75cf8d4cd558955a5c7638fba2d68f464/YM/PhysicalNext/RG_NEXT_COMPLETE_REPEATABLE_NONLINEAR_STEP/Pane/Purl.lean\#L20-L53}{\texttt{panes}}
    for III.3.3.
  \end{itemize}
\item
  \textbf{III.16.} Kernels and their action on tori:
  \href{https://github.com/yang-mills-group/OpenAI-YM-20261028/blob/3456e5e75cf8d4cd558955a5c7638fba2d68f464/YM/CellFourier/Circulant.lean\#L22}{\texttt{Row}},
  \href{https://github.com/yang-mills-group/OpenAI-YM-20261028/blob/3456e5e75cf8d4cd558955a5c7638fba2d68f464/YM/CellFourier/Shuttle.lean\#L26}{\texttt{Cross}},
  \href{https://github.com/yang-mills-group/OpenAI-YM-20261028/blob/3456e5e75cf8d4cd558955a5c7638fba2d68f464/YM/CellFourier/Shuttle.lean\#L35}{\texttt{installed}}.
  The weighted norms
  \href{https://github.com/yang-mills-group/OpenAI-YM-20261028/blob/3456e5e75cf8d4cd558955a5c7638fba2d68f464/YM/CellFourier/Circulant.lean\#L26}{\texttt{Row.strip}}
  and
  \href{https://github.com/yang-mills-group/OpenAI-YM-20261028/blob/3456e5e75cf8d4cd558955a5c7638fba2d68f464/YM/Laurent/Eval.lean\#L104}{\texttt{Laurent.Width}},
  where the size of a displacement is its \(\ell^1\) norm. The
  alias-resolved Fourier transform on \(\Lambda_f\), which is unitary up
  to \(f^4\).
\end{itemize}

\textbf{Standard mathematics.}

\begin{itemize}
\tightlist
\item
  Fourier analysis on \(\mathbb Z^4\) and on periodic lattices. Bloch
  (alias) decomposition of operators that commute with coarse
  translations. Periodic strip analyticity gives exponential coefficient
  bounds on narrower strips.
\item
  Gaussian block-spin transformations computed in Fourier space: K.
  Gawędzki, A. Kupiainen, \emph{Massless lattice \(\varphi^4_4\)
  theory}, CMP 99 (1985). Balaban's averaging and background-field
  operators for lattice gauge fields: T. Balaban, CMP 95 and 96 (1984),
  CMP 98 (1985). Exposition: J. Dimock, \emph{The renormalization group
  according to Balaban I}, Rev.~Math. Phys. 25 (2013).
\item
  Lagrange's identity. The Schur-complement (orthogonal) splitting of a
  quadratic form.
\item
  Exactness of the lattice de Rham (Koszul) complex for exponentially
  decaying translation-invariant kernels.
\item
  Newton forward-difference interpolation. Weighted-homogeneous Taylor
  expansions. The Fock--Schwinger (radial) gauge.
\end{itemize}

\textbf{Imported constants and assumptions.}

\begin{itemize}
\tightlist
\item
  The kernel choice \(K\); only the penalty \(\kappa\) enters the free
  kernels.
\item
  The analyticity half-width \(\pi/4\)
  (\href{https://github.com/yang-mills-group/OpenAI-YM-20261028/blob/3456e5e75cf8d4cd558955a5c7638fba2d68f464/YM/Next/FREE_KERNELS/Chip.lean\#L21}{\texttt{zone}})
  and the radius \(3\pi/2\)
  (\href{https://github.com/yang-mills-group/OpenAI-YM-20261028/blob/3456e5e75cf8d4cd558955a5c7638fba2d68f464/YM/Next/FREE_KERNELS/Chip.lean\#L24}{\texttt{domainRadius}}).
  The proved compactness bounds \(c_-|w|\le|e^{iw}-1|\le c_+|w|\) on
  \(|w|\le3\pi/2\)
  (\href{https://github.com/yang-mills-group/OpenAI-YM-20261028/blob/3456e5e75cf8d4cd558955a5c7638fba2d68f464/YM/Next/FREE_KERNELS/Chip.lean\#L82}{\texttt{meter}}).
\item
  The credit condition, a hypothesis of III.3.2. It is discharged
  through
  \href{https://github.com/yang-mills-group/OpenAI-YM-20261028/blob/3456e5e75cf8d4cd558955a5c7638fba2d68f464/YM/Next/FREE_KERNELS/Calendar.lean\#L108-L124}{\texttt{before\_paid}},
  the construction of the sheet in III.9, and the payment record of
  III.2.
\item
  F06--F09 are proved in III.6 and plugged in inside
  \href{https://github.com/yang-mills-group/OpenAI-YM-20261028/blob/3456e5e75cf8d4cd558955a5c7638fba2d68f464/YM/Next/FREE_KERNELS/Rack.lean\#L134-L145}{\texttt{retained\_support}}.
\end{itemize}

\subsubsection{3. Strategy}\label{strategy-5}

The Gaussian block-spin problem is solved in Fourier space. Iterating
the block average and its independent noise gives a diagonal 1-form
covariance symbol. Minimising its precision over gauge directions gives
the effective form \(d^*Dd\) by Lagrange's identity. The
background-field map \(W\) and curvature transfer \(T\) are then
explicit alias sums satisfying (D1)--(D4).

Uniform decay follows from the filters. A stage-\(J\) box filter gives
\(\prod_a(1+|\ell_a|)^{-J}\) alias decay, leaving \(J-2\) fine
differences summable. Along a history, successive dilation ratios
decrease at least geometrically by \(1/3\), so the extra filter products
and accumulated noise stay uniformly bounded. Before estimating inverse
symbols, the principal-alias poles are cancelled. The remaining
denominators are bounded below at real momenta and hence on a thinner
complex cube. These constants depend on the kernel choice and stage,
uniformly in history length and block sizes.

The inverse-metric defect factors through \(d_2\) as \(X'^*CX'\), where
\(X'=d_2D'^{-1}\) and \(C\) need not be positive. Amending the fine and
coarse inverse metrics leaves their quadratic forms on potentials
unchanged and replaces \(C\) by \(C+\lambda_t-\lambda_jS^*S\). The
squared 3-form transfer satisfies \(S^*S=O(L^{-2})\), including the sum
over fine residues. Choose the coarse amendments first, then require
\(\lambda_jB/L^2\le1\). The resulting positive defect yields
\(M_f,M_n,O\) and (R1)--(R6) through III.6 and the repair theorem of
III.10. The repair constants also depend on the sheet's direct clip
coefficients. Rigidity extends the schedule identities to all tori.

\subsubsection{4. Steps}\label{steps-5}

\paragraph{4.1 Filters and histories}\label{filters-and-histories}

\textbf{Lemma 4.1 (filter algebra).} This is III.3.5(a). \emph{Sketch:}
\(M\,\mathrm{rail}_M(z)(e^{iz}-1)=e^{iMz}-1\); apply it twice and cancel
\(e^{iz}-1\) when nonzero. If \(e^{iz}=1\), every uncentred average is
\(1\). For the centred filter, the centring phases add up, since
\(M\,\frac{L-1}2+\frac{M-1}2=\frac{LM-1}2\). \emph{Lean:}
\href{https://github.com/yang-mills-group/OpenAI-YM-20261028/blob/3456e5e75cf8d4cd558955a5c7638fba2d68f464/YM/Next/FREE_KERNELS/Lines.lean\#L51-L62}{\texttt{rail\_bind}},
\href{https://github.com/yang-mills-group/OpenAI-YM-20261028/blob/3456e5e75cf8d4cd558955a5c7638fba2d68f464/YM/Next/FREE_KERNELS/Lines.lean\#L68-L76}{\texttt{plank\_bind}}.

\textbf{Lemma 4.2 (size of filters).} This is III.3.5(b)--(d).
\emph{Sketch:}

\begin{itemize}
\tightlist
\item
  (b) Triangle inequality.
\item
  (c) The compactness constants and Lemma 4.1 give
  \(|\mathrm{rail}_M|\ge c_-/c_+\). On the principal alias, averaging
  \(|e^{isz}-1|\le c_+|sz|\) gives
  \(|\mathrm{plank}_M(z)-1|\le c_+\tfrac{3\pi}2\tfrac ML\), so
  \(|\mathrm{plank}_M|\ge\max(c_*,1-k x)\ge e^{-c_1x}\) with \(x=M/L\).
  The exponent is linear in \(x\) with no prefactor, which is what makes
  products along a history converge.
\item
  (d) \(L\,\mathrm{rail}_L(z)(e^{iz}-1)=e^{ip_a}-1\) at
  \(z=(q_\ell)_a\), and \(L|z|\gtrsim1+|\ell_a|\) when \(\ell_a\ne0\).
\end{itemize}

\emph{Lean:}
\href{https://github.com/yang-mills-group/OpenAI-YM-20261028/blob/3456e5e75cf8d4cd558955a5c7638fba2d68f464/YM/Next/FREE_KERNELS/Lines.lean\#L127-L133}{\texttt{plank\_height}},
\href{https://github.com/yang-mills-group/OpenAI-YM-20261028/blob/3456e5e75cf8d4cd558955a5c7638fba2d68f464/YM/Next/FREE_KERNELS/Lines.lean\#L114-L122}{\texttt{rail\_height}},
\href{https://github.com/yang-mills-group/OpenAI-YM-20261028/blob/3456e5e75cf8d4cd558955a5c7638fba2d68f464/YM/Next/FREE_KERNELS/Lines.lean\#L140-L152}{\texttt{rail\_lower}},
\href{https://github.com/yang-mills-group/OpenAI-YM-20261028/blob/3456e5e75cf8d4cd558955a5c7638fba2d68f464/YM/Next/FREE_KERNELS/Lines.lean\#L154-L155}{\texttt{plank\_lower}},
\href{https://github.com/yang-mills-group/OpenAI-YM-20261028/blob/3456e5e75cf8d4cd558955a5c7638fba2d68f464/YM/Next/FREE_KERNELS/Lines.lean\#L162-L196}{\texttt{principal\_plank}},
\href{https://github.com/yang-mills-group/OpenAI-YM-20261028/blob/3456e5e75cf8d4cd558955a5c7638fba2d68f464/YM/Next/FREE_KERNELS/Lines.lean\#L200-L245}{\texttt{alias\_rail}}.

\textbf{Lemma 4.3 (histories).} A history of stage index \(j\) is either
empty, or a hop \((L,J)\) with \(j\le J\le j+1\) and odd \(L\ge3\),
applied after a history of index \(J\). Its dilation \(M_v\) is the
product of its block sizes. Its cumulative filter is defined by
\(s_\emptyset=1\) and \(s_{(L,J)\circ v}(p)=a_L(M_vp)^J s_v(p)\).
Equivalently, \(s_v=a_{M_v}^{\,j}e_v\), where \(e_v\) is a product of
box filters at the hops where the index drops. Then:

\begin{itemize}
\tightlist
\item
  \(s_v\) is entire, \(2\pi\mathbb Z^4\)-periodic and invariant under
  signed permutations;
\item
  on \(\mathcal C_{\pi/4}\),
  \(|s_v(q_\ell)|\le C_j\prod_a(1+|\ell_a|)^{-j}\) at the aliases
  \(q_\ell=(p+2\pi\ell)/M_v\);
\item
  for \(p\in\mathcal C_{\pi/4}^{\mathbb R}\), \(|s_v(p/M_v)|\ge c_j>0\).
\end{itemize}

Here \(C_j,c_j\) do not depend on \(v\), that is, on the number of hops
or the block sizes. \emph{Sketch:} After \(k\) further hops, a hop of
dilation \(M'\) is evaluated at the momentum scaled by
\(M'/M_v\le3^{-k}\). Lemma 4.2 gives factors \(e^{\pm CM'/M_v}\), and
these products converge. The factor \(a_{M_v}^j\) gives the alias decay
by 4.2(d); the geometric product bounds the extra factors. \emph{Lean:}
\href{https://github.com/yang-mills-group/OpenAI-YM-20261028/blob/3456e5e75cf8d4cd558955a5c7638fba2d68f464/YM/Next/FREE_KERNELS/Train.lean\#L22-L24}{\texttt{Train}},
\href{https://github.com/yang-mills-group/OpenAI-YM-20261028/blob/3456e5e75cf8d4cd558955a5c7638fba2d68f464/YM/Next/FREE_KERNELS/Train.lean\#L36}{\texttt{Train.smooth}},
\href{https://github.com/yang-mills-group/OpenAI-YM-20261028/blob/3456e5e75cf8d4cd558955a5c7638fba2d68f464/YM/Next/FREE_KERNELS/Train.lean\#L43-L48}{\texttt{Train.bind}},
\href{https://github.com/yang-mills-group/OpenAI-YM-20261028/blob/3456e5e75cf8d4cd558955a5c7638fba2d68f464/YM/Next/FREE_KERNELS/Train.lean\#L169-L184}{\texttt{Train.smooth\_upper}},
\href{https://github.com/yang-mills-group/OpenAI-YM-20261028/blob/3456e5e75cf8d4cd558955a5c7638fba2d68f464/YM/Next/FREE_KERNELS/Train.lean\#L186-L194}{\texttt{Train.smooth\_lower}}.

\paragraph{4.2 Explicit metric data}\label{explicit-metric-data}

\textbf{Lemma 4.4 (auxiliary 1-form covariance symbols).} Fix
\(\kappa>0\) and a history \(v\) of dilation \(M\). For a direction
\(a\), let \(\chi_{v,a}(p)=s_v(p)\,\mathrm{rail}_M(p_a)\). For
\(p\in\mathcal C_{\pi/4}^{\mathbb R}\setminus\{0\}\) and
\(q_\ell=(p+2\pi\ell)/M\), define
\[D_a(p)=M^{-2}\sum_{\ell}\frac{|\chi_{v,a}(q_\ell)|^2}{|\tau(q_\ell)|^2}+\kappa^{-1}\sum_{v'}M_{v'}^{-2}\sum_{\ell'}|\chi_{v',a}(q'_{\ell'})|^2 .\]
The second sum runs over the proper final segments \(v'\) of \(v\) (the
last \(k\) hops, \(0\le k<\sharp v\)), with
\(q'_{\ell'}=(p+2\pi\ell')/M_{v'}\). For the empty history,
\(D_a=|\tau|^{-2}\). The formula obeys the exact one-step recursion
\(D'_a(p)=L^{-2}\sum_\ell|a_L(q_\ell)^j\mathrm{rail}_L((q_\ell)_a)|^2D_a(q_\ell)+\kappa^{-1}\),
which is the covariance recursion for an averaged field with independent
noise. These singular symbols are used only away from zero; Lemma 4.6
constructs the regular 2-form symbols. \emph{Lean:}
\href{https://github.com/yang-mills-group/OpenAI-YM-20261028/blob/3456e5e75cf8d4cd558955a5c7638fba2d68f464/YM/Next/FREE_KERNELS/Feed.lean\#L133-L134}{\texttt{Train.levelD}},
\href{https://github.com/yang-mills-group/OpenAI-YM-20261028/blob/3456e5e75cf8d4cd558955a5c7638fba2d68f464/YM/Next/FREE_KERNELS/Feed.lean\#L154-L156}{\texttt{D\_feedback}},
\href{https://github.com/yang-mills-group/OpenAI-YM-20261028/blob/3456e5e75cf8d4cd558955a5c7638fba2d68f464/YM/Next/FREE_KERNELS/Feed.lean\#L131}{\texttt{Train.levelNoise}},
\href{https://github.com/yang-mills-group/OpenAI-YM-20261028/blob/3456e5e75cf8d4cd558955a5c7638fba2d68f464/YM/Next/FREE_KERNELS/Feed.lean\#L62}{\texttt{feed}}.
\emph{(definitions and covariance recursion proof read)}

\textbf{Lemma 4.5 (the diagonal).} At the nonzero real momenta of Lemma
4.4, let \(\rho=\sum_a|\tau_a|^2/D_a\) and
\(D_{\mu\nu}=(\rho\,D_\mu D_\nu)^{-1}\); this is a multiplier diagonal
in the plaquette orientation. Then, with \(F=\mathrm{diag}(D_a)\),
\[\sigma_1^*D\,\sigma_1=F^{-1}-\rho^{-1}F^{-1}\tau\,\tau^*F^{-1}.\] So
\(d^*Dd\) is the precision \(F^{-1}\) minimised over gauge directions.
For the empty history, \(D=1\) (the Wilson quadratic form).
\emph{Sketch:} Lagrange's identity
\(\sum_{\mu<\nu}w_\mu w_\nu|\tau_\mu x_\nu-\tau_\nu x_\mu|^2=(\sum_aw_a|\tau_a|^2)(\sum_aw_a|x_a|^2)-|\sum_aw_a\bar\tau_ax_a|^2\),
with \(w_a=D_a^{-1}\). The weights and \(\rho\) are positive.
\emph{Lean:}
\href{https://github.com/yang-mills-group/OpenAI-YM-20261028/blob/3456e5e75cf8d4cd558955a5c7638fba2d68f464/YM/Next/FREE_KERNELS/Signal.lean\#L103-L120}{\texttt{rotor\_law}},
\href{https://github.com/yang-mills-group/OpenAI-YM-20261028/blob/3456e5e75cf8d4cd558955a5c7638fba2d68f464/YM/Next/FREE_KERNELS/Heart.lean\#L22}{\texttt{pulse}},
\href{https://github.com/yang-mills-group/OpenAI-YM-20261028/blob/3456e5e75cf8d4cd558955a5c7638fba2d68f464/YM/Next/FREE_KERNELS/Heart.lean\#L20}{\texttt{vault}},
\href{https://github.com/yang-mills-group/OpenAI-YM-20261028/blob/3456e5e75cf8d4cd558955a5c7638fba2d68f464/YM/Next/FREE_KERNELS/Curl.lean\#L54}{\texttt{pink}},
\href{https://github.com/yang-mills-group/OpenAI-YM-20261028/blob/3456e5e75cf8d4cd558955a5c7638fba2d68f464/YM/Next/FREE_KERNELS/Curl.lean\#L55}{\texttt{blue}},
\href{https://github.com/yang-mills-group/OpenAI-YM-20261028/blob/3456e5e75cf8d4cd558955a5c7638fba2d68f464/YM/Next/FREE_KERNELS/Heart.lean\#L140-L155}{\texttt{top\_metric}}.

\textbf{Lemma 4.6 (filled symbols, uniform analyticity).} Fix
\(\kappa>0\). Away from \(p=0\) in the real cube, the singular formulas
become regular combinations. With
\(|\tau_M|^2=\sum_aM^2|e^{ip_a/M}-1|^2\):
\[D_a=\beta_a/|\tau_M|^2,\qquad \rho=|\tau_M|^4h,\qquad D_{\mu\nu}=(h\,\beta_\mu\beta_\nu)^{-1}.\]
Here \(\beta_a\) is the principal term \(|\chi_{v,a}(p/M)|^2\) plus
\(|\tau_M|^2\) times the non-principal and noise terms, and \(h\) is a
rational combination of filters whose denominators stay nonzero on a
sufficiently thin cube. The regular formulas give \(\beta_a(0)=h(0)=1\).
For each \(j\ge4\) there are \(\theta,C,c>0\), depending only on
\(\kappa,j\), such that, for all histories of index \(j\), \(\beta_a\)
and \(h\) are analytic and bounded by \(C\) on \(\mathcal C_\theta\),
and \(|\beta_a|,|h|\ge c\) on \(\mathcal C_\theta^{\mathbb R}\). Hence
\(D\) and \(D^{-1}\) lie in a common \(\mathcal S(\alpha,\Omega)\),
uniformly over histories of index \(j\). \emph{Sketch:}

\begin{itemize}
\tightlist
\item
  Upper bounds: Lemma 4.3 summed over the aliases.
\item
  Lower bound on \(\beta_a\): its principal term (4.3 and 4.2(c)).
\item
  Lower bound on \(h\): away from zero use
  \(|\tau_M|^2h=\sum_a|\tau_a|^2/\beta_a\), \(\beta_a\le C\) and
  \(|\tau|^2\asymp|\tau_M|^2\asymp|p|^2\); at zero use \(h(0)=1\).
\item
  Inverses: a lower bound at real momenta plus an analytic bound gives a
  lower bound on a thinner complex cube, so the inverses stay analytic.
\end{itemize}

\emph{Lean:}
\href{https://github.com/yang-mills-group/OpenAI-YM-20261028/blob/3456e5e75cf8d4cd558955a5c7638fba2d68f464/YM/Next/FREE_KERNELS/Bend.lean\#L39}{\texttt{Train.bulb}},
\href{https://github.com/yang-mills-group/OpenAI-YM-20261028/blob/3456e5e75cf8d4cd558955a5c7638fba2d68f464/YM/Next/FREE_KERNELS/Bend.lean\#L43-L44}{\texttt{Train.bend}},
\href{https://github.com/yang-mills-group/OpenAI-YM-20261028/blob/3456e5e75cf8d4cd558955a5c7638fba2d68f464/YM/Next/FREE_KERNELS/Feed.lean\#L192-L201}{\texttt{Train.pole\_split}},
\href{https://github.com/yang-mills-group/OpenAI-YM-20261028/blob/3456e5e75cf8d4cd558955a5c7638fba2d68f464/YM/Next/FREE_KERNELS/Bend.lean\#L88-L100}{\texttt{Train.lift\_bend}},
\href{https://github.com/yang-mills-group/OpenAI-YM-20261028/blob/3456e5e75cf8d4cd558955a5c7638fba2d68f464/YM/Next/FREE_KERNELS/Bend.lean\#L150-L153}{\texttt{Train.recover\_rho}},
\href{https://github.com/yang-mills-group/OpenAI-YM-20261028/blob/3456e5e75cf8d4cd558955a5c7638fba2d68f464/YM/Next/FREE_KERNELS/Bend.lean\#L160-L163}{\texttt{Train.bulb\_firm}},
\href{https://github.com/yang-mills-group/OpenAI-YM-20261028/blob/3456e5e75cf8d4cd558955a5c7638fba2d68f464/YM/Next/FREE_KERNELS/Bend.lean\#L188-L194}{\texttt{Train.bend\_firm}},
\href{https://github.com/yang-mills-group/OpenAI-YM-20261028/blob/3456e5e75cf8d4cd558955a5c7638fba2d68f464/YM/Next/FREE_KERNELS/Dispatch.lean\#L40-L47}{\texttt{heart\_banked}},
\href{https://github.com/yang-mills-group/OpenAI-YM-20261028/blob/3456e5e75cf8d4cd558955a5c7638fba2d68f464/YM/Next/FREE_KERNELS/Moderate.lean\#L105-L106}{\texttt{Moderate}}.

\textbf{Lemma 4.7 (amended metric).} Fix \(\kappa>0\), \(j\ge4\) and
\(\lambda\ge0\). For each history of index \(j\), let
\[H^{-1}=D^{-1}+\lambda\,X^*X,\quad X=d_2D^{-1},\qquad \mathcal T_\lambda=(1+\lambda\,X d_2^*)^{-1}\ \text{on 3-forms}.\]
Then on every torus: \(HH^{-1}=1\), \(Hd=Dd\),
\(XH=\mathcal T_\lambda d_2\), and \(H,\mathcal T_\lambda\) are
symmetric. Moreover \(\hat H(0)=1\), and
\(H,H^{-1},\mathcal T_\lambda,\mathcal T_\lambda^{-1}\) have strip
bounds and two-sided bounds at real momenta, with constants depending
only on \(\kappa,j,\lambda\). The amendment changes \(H\) only away from
exact 2-forms, and leaves \(d^*Hd=d^*Dd\) unchanged. \emph{Sketch:}

\begin{itemize}
\tightlist
\item
  \(XDd=d_2d=0\), which gives \(Hd=Dd\).
\item
  \((1+\lambda Xd_2^*)X=d_2H^{-1}\), which gives
  \(XH=\mathcal T_\lambda d_2\).
\item
  Uniform bounds \(a\,1\le D^{-1}\le A\,1\) and \(\|X\|\le B\) give
  \(a\,1\le H^{-1}\le(A+\lambda B^2)\,1\) and hence two-sided bounds for
  \(H\). The analogous argument applies to \(1+\lambda d_2D^{-1}d_2^*\).
\end{itemize}

\emph{Lean:}
\href{https://github.com/yang-mills-group/OpenAI-YM-20261028/blob/3456e5e75cf8d4cd558955a5c7638fba2d68f464/YM/Next/FREE_KERNELS/Amend.lean\#L30}{\texttt{pledge}},
\href{https://github.com/yang-mills-group/OpenAI-YM-20261028/blob/3456e5e75cf8d4cd558955a5c7638fba2d68f464/YM/Next/FREE_KERNELS/Amend.lean\#L31}{\texttt{amend}},
\href{https://github.com/yang-mills-group/OpenAI-YM-20261028/blob/3456e5e75cf8d4cd558955a5c7638fba2d68f464/YM/Next/FREE_KERNELS/Amend.lean\#L32}{\texttt{bail}},
\href{https://github.com/yang-mills-group/OpenAI-YM-20261028/blob/3456e5e75cf8d4cd558955a5c7638fba2d68f464/YM/Next/FREE_KERNELS/Amend.lean\#L33}{\texttt{treble}},
\href{https://github.com/yang-mills-group/OpenAI-YM-20261028/blob/3456e5e75cf8d4cd558955a5c7638fba2d68f464/YM/Next/FREE_KERNELS/Amend.lean\#L95-L127}{\texttt{obligations}},
\href{https://github.com/yang-mills-group/OpenAI-YM-20261028/blob/3456e5e75cf8d4cd558955a5c7638fba2d68f464/YM/Next/FREE_KERNELS/Amend.lean\#L41-L59}{\texttt{pledge\_warm}},
\href{https://github.com/yang-mills-group/OpenAI-YM-20261028/blob/3456e5e75cf8d4cd558955a5c7638fba2d68f464/YM/Next/FREE_KERNELS/Amend.lean\#L78-L79}{\texttt{insurance}},
\href{https://github.com/yang-mills-group/OpenAI-YM-20261028/blob/3456e5e75cf8d4cd558955a5c7638fba2d68f464/YM/Next/FREE_KERNELS/Amend.lean\#L80-L94}{\texttt{forced\_amend}}.

\textbf{Lemma 4.8 (metric data of every history; III.3.1(a)).} Let
\(\kappa>0\), \(\lambda\ge0\), \(j\ge4\). There are
\(\gamma,\sigma,\eta,C>0\) with the following property. For every clip
with \(c\le\gamma\) and every history of index \(j\), the explicit
\(D,D^{-1},H,H^{-1}\), together with a positive symmetric square root
\(R\) of \(H-c\,1\), satisfy III.3.1(a) with amendment \(\lambda\) and
bounds \(\sigma,\eta,C\). These constants depend only on
\(\kappa,j,\lambda\). \emph{Sketch:}

\begin{itemize}
\tightlist
\item
  \(\gamma\) is half the real-momentum lower bound of \(H\), so
  \(H-c\,1\) has two-sided bounds.
\item
  III.6's square-root theorem gives \(R\in\mathcal S\) with
  \(R^*R=H-c\,1\) on every torus.
\item
  Torus bounds follow from symbol bounds, since \(\Lambda_n\) is
  diagonalised by the momenta \(2\pi\mathbb Z^4/n\).
\item
  The symmetry follows from Lemma 4.3.
\item
  \(D=1\) for the empty history by Lemma 4.5.
\end{itemize}

\emph{Lean:}
\href{https://github.com/yang-mills-group/OpenAI-YM-20261028/blob/3456e5e75cf8d4cd558955a5c7638fba2d68f464/YM/Next/FREE_KERNELS/Ground.lean\#L69-L120}{\texttt{bid\_bed}},
\href{https://github.com/yang-mills-group/OpenAI-YM-20261028/blob/3456e5e75cf8d4cd558955a5c7638fba2d68f464/YM/Next/FREE_KERNELS/Ground.lean\#L54-L58}{\texttt{Defined}},
\href{https://github.com/yang-mills-group/OpenAI-YM-20261028/blob/3456e5e75cf8d4cd558955a5c7638fba2d68f464/YM/Next/FREE_KERNELS/Ground.lean\#L22-L35}{\texttt{full\_cover}},
\href{https://github.com/yang-mills-group/OpenAI-YM-20261028/blob/3456e5e75cf8d4cd558955a5c7638fba2d68f464/YM/Next/FREE_KERNELS/Ground.lean\#L37-L51}{\texttt{low\_ceiling}},
\href{https://github.com/yang-mills-group/OpenAI-YM-20261028/blob/3456e5e75cf8d4cd558955a5c7638fba2d68f464/YM/Next/FREE_KERNELS/Signed.lean\#L65-L70}{\texttt{sworn\_metric}}.

\paragraph{4.3 Transfer kernels of one
flight}\label{transfer-kernels-of-one-flight}

A flight of index \((j,t)\) consists of an odd block size \(L\ge3\), a
fine history \(v\) of index \(j\), and the coarse history (hop \((L,j)\)
after \(v\)) of index \(t\in\{j-1,j\}\).

\textbf{Lemma 4.9 (alias symbols and derivative bounds).} Fix
\(\kappa>0\) and \(j,t\ge4\). There are \(\theta,\alpha,C>0\), depending
only on \(\kappa,j,t\), such that the following holds for every flight
of index \((j,t)\). Its explicit \(W\) and \(T\) are cross kernels; at
this stage \(M_f=M_n=O=0\). Their fibre at coarse momentum \(p\) and
fine residue \(r\) is
\[L^{-e}\sum_\ell e^{iq_\ell\cdot r}w_e(p,\ell),\qquad e=1\ (W),\quad e=2\ (T),\]
with explicit per-alias symbols built from the filters and from
\(\beta_a,h\) of both histories. At the harmonic point \(w_e(0,0)=1\)
and \(w_e(0,\ell)=0\) for \(\ell\ne0\). Each \(w_e(\cdot,\ell)\) is
analytic on \(\mathcal C_\theta\) and bounded by
\(C\prod_a(1+|\ell_a|)^{-j}\), uniformly over flights. Hence, for words
\(u\) with \(|u|\le j-2\),
\(\nabla_uW\in\mathcal S(\alpha,CL^{-1-|u|})\) and
\(\nabla_uT\in\mathcal S(\alpha,CL^{-2-|u|})\). \emph{Sketch:} A fine
difference in direction \(a\) multiplies the alias term by
\(e^{i(q_\ell)_a}-1=O((1+|\ell|)/L)\). The sum
\(\sum_\ell(1+|\ell|)^{|u|}\prod_a(1+|\ell_a|)^{-j}\) converges for
\(|u|\le j-2\). In other words, the order \(j\) of the smoothing pays
for \(j-2\) derivatives. Decay is measured on the coarse lattice, so the
fine-lattice rate is \(\alpha/L\). \emph{Lean:}
\href{https://github.com/yang-mills-group/OpenAI-YM-20261028/blob/3456e5e75cf8d4cd558955a5c7638fba2d68f464/YM/Next/FREE_KERNELS/Primal.lean\#L24-L29}{\texttt{Flight.duo}},
\href{https://github.com/yang-mills-group/OpenAI-YM-20261028/blob/3456e5e75cf8d4cd558955a5c7638fba2d68f464/YM/Next/FREE_KERNELS/Docked.lean\#L103-L110}{\texttt{Flight.pilot\_cost}},
\href{https://github.com/yang-mills-group/OpenAI-YM-20261028/blob/3456e5e75cf8d4cd558955a5c7638fba2d68f464/YM/Next/FREE_KERNELS/Docked.lean\#L53-L71}{\texttt{Flight.ahead\_cost}},
\href{https://github.com/yang-mills-group/OpenAI-YM-20261028/blob/3456e5e75cf8d4cd558955a5c7638fba2d68f464/YM/Next/FREE_KERNELS/Cable.lean\#L97-L103}{\texttt{Flight.wire\_budget}},
\href{https://github.com/yang-mills-group/OpenAI-YM-20261028/blob/3456e5e75cf8d4cd558955a5c7638fba2d68f464/YM/Next/FREE_KERNELS/Strands.lean\#L30-L38}{\texttt{secure\_both}},
\href{https://github.com/yang-mills-group/OpenAI-YM-20261028/blob/3456e5e75cf8d4cd558955a5c7638fba2d68f464/YM/Next/FREE_KERNELS/Strands.lean\#L85-L107}{\texttt{Secured.propels}}.

\textbf{Lemma 4.10 (exact identities for one flight).} Fix a flight of
index \((j,t)\) with \(t\ge4\), a kernel choice \(K\), and any pair of
tori with \(f=cL\). Then (D1)--(D4) hold for:

\begin{itemize}
\tightlist
\item
  \(W\), \(T\) of the flight and \(Q\) of the step;
\item
  \(D\), \(D'\) the diagonals of the fine and coarse histories;
\item
  \(H\), \(H'\) the amended metrics for any amendments \(\ge0\).
\end{itemize}

\emph{Sketch:}

\begin{itemize}
\tightlist
\item
  All operators commute with coarse translations. On the Bloch fibre
  over \(p\), the identities become identities between per-alias
  symbols, for each alias \(q=q_\ell\):

  \begin{itemize}
  \tightlist
  \item
    \((\sigma_1^*D\sigma_1)(q)\,W_\ell=L^{-3}h_1(q)^*(\sigma_1^*D'\sigma_1)(p)\),
    giving (D1);
  \item
    \(\sigma_1(q)^*T_\ell=L^{-3}h_1(q)^*\sigma_1(p)^*\), giving (D3);
  \item
    \(T_\ell D'(p)\sigma_1(p)=D(q)\sigma_1(q)W_\ell\), giving (D4);
  \item
    the alias sum
    \(\sum_\ell Lh_1(q_\ell)W_\ell=1-\kappa^{-1}(\sigma_1^*D'\sigma_1)(p)\),
    giving (D2).
  \end{itemize}
\item
  Here \(W_\ell=L^{-1}w_1(p,\ell)\), \(T_\ell=L^{-2}w_2(p,\ell)\), and
  \(h_r(q)=a_L(q)^j\prod_{a\in\text{face}}\mathrm{rail}_L(q_a)\) is the
  alias symbol of the \(r\)-form average. The forward block average has
  alias row \(Lh_1(q)\), and its adjoint has alias column
  \(L^{-3}h_1(q)^*\) because the fine-to-coarse volume ratio is \(L^4\).
  The form averages intertwine the exterior derivatives:
  \(h_{r+1}(q)\,L\sigma_r(q)=\sigma_r(p)\,h_r(q)\), from
  \(L\,\mathrm{rail}_L(q)(e^{iq}-1)=e^{iLq}-1\).
\item
  The point \(p=0\) is treated separately: the box filter vanishes at
  the non-principal aliases of \(0\).
\item
  The identities pass from the fibres to the tori through the
  alias-resolved Fourier transform.
\item
  Replacing \(D\) by the amended \(H\) in (D1)--(D2) changes nothing,
  since \(d^*Hd=d^*Dd\).
\end{itemize}

\emph{Lean:}
\href{https://github.com/yang-mills-group/OpenAI-YM-20261028/blob/3456e5e75cf8d4cd558955a5c7638fba2d68f464/YM/Next/FREE_KERNELS/Primal.lean\#L63-L118}{\texttt{Flight.tangent}},
\href{https://github.com/yang-mills-group/OpenAI-YM-20261028/blob/3456e5e75cf8d4cd558955a5c7638fba2d68f464/YM/Next/FREE_KERNELS/Rack.lean\#L44-L56}{\texttt{diagonal}},
\href{https://github.com/yang-mills-group/OpenAI-YM-20261028/blob/3456e5e75cf8d4cd558955a5c7638fba2d68f464/YM/Next/FREE_KERNELS/Primal.lean\#L33-L49}{\texttt{Flight.scaled\_tri}},
\href{https://github.com/yang-mills-group/OpenAI-YM-20261028/blob/3456e5e75cf8d4cd558955a5c7638fba2d68f464/YM/Next/FREE_KERNELS/Signal.lean\#L60-L63}{\texttt{hinge\_transport}},
\href{https://github.com/yang-mills-group/OpenAI-YM-20261028/blob/3456e5e75cf8d4cd558955a5c7638fba2d68f464/YM/Next/FREE_KERNELS/ZeroAlias.lean}{\texttt{ZeroAlias.lean}},
\href{https://github.com/yang-mills-group/OpenAI-YM-20261028/blob/3456e5e75cf8d4cd558955a5c7638fba2d68f464/YM/Next/FREE_KERNELS/Pilot.lean\#L46-L47}{\texttt{pilot}},
\href{https://github.com/yang-mills-group/OpenAI-YM-20261028/blob/3456e5e75cf8d4cd558955a5c7638fba2d68f464/YM/Next/FREE_KERNELS/Wire.lean\#L63}{\texttt{wire}},
\href{https://github.com/yang-mills-group/OpenAI-YM-20261028/blob/3456e5e75cf8d4cd558955a5c7638fba2d68f464/YM/Next/FREE_KERNELS/Scripts.lean\#L122-L139}{\texttt{script\_main}},
\href{https://github.com/yang-mills-group/OpenAI-YM-20261028/blob/3456e5e75cf8d4cd558955a5c7638fba2d68f464/YM/Next/FREE_KERNELS/Scripts.lean\#L141-L156}{\texttt{script\_cross}},
\href{https://github.com/yang-mills-group/OpenAI-YM-20261028/blob/3456e5e75cf8d4cd558955a5c7638fba2d68f464/YM/Next/FREE_KERNELS/Scripts.lean\#L179-L211}{\texttt{script\_dual}}.
\emph{(proof read, including the nonzero-alias algebra and harmonic
point)}

\paragraph{4.4 The defect and the
amendment}\label{the-defect-and-the-amendment}

\textbf{Lemma 4.11 (unamended defect).} In the setting of 4.10, let
\(\mathcal V=D'^{-1}-T^*D^{-1}T-\kappa^{-1}dd^*\) on
\(\Omega^2(\Lambda_c)\). Then \(D'\mathcal VD'\) is symmetric, and
\(D'\mathcal VD'd=0\) on every torus. Consequently there is an explicit
symmetric 3-form kernel \(C\), decaying uniformly over flights of index
\((j,t)\), with
\[\mathcal V=(d_2D'^{-1})^*\,C\,(d_2D'^{-1})\quad\text{on every torus.}\]
\emph{Sketch:}

\begin{itemize}
\tightlist
\item
  By (D4), (D3), (D2): \(T^*D^{-1}TD'd=T^*dW=dQW=d-\kappa^{-1}dd^*D'd\).
  Hence \(\mathcal VD'd=0\).
\item
  A symmetric decaying 2-form kernel \(A\) with \(Ad=0\) on every torus
  equals \(d_2^*Bd_2\), for an explicit decaying 3-form kernel \(B\).
  This is a lattice Poincaré lemma for translation-invariant kernels,
  proved by an axis-by-axis contracting homotopy of the Koszul complex
  of \((e^{ip_a}-1)_a\).
\item
  The hypothesis is needed on all tori, not on one: at a single period
  the harmonic (zero-momentum) forms obstruct the factorisation.
\end{itemize}

\emph{Lean:}
\href{https://github.com/yang-mills-group/OpenAI-YM-20261028/blob/3456e5e75cf8d4cd558955a5c7638fba2d68f464/YM/Next/FREE_KERNELS/Sell.lean\#L33-L35}{\texttt{Flight.gap}},
\href{https://github.com/yang-mills-group/OpenAI-YM-20261028/blob/3456e5e75cf8d4cd558955a5c7638fba2d68f464/YM/Next/FREE_KERNELS/Sell.lean\#L158-L187}{\texttt{Flight.dep\_zero}},
\href{https://github.com/yang-mills-group/OpenAI-YM-20261028/blob/3456e5e75cf8d4cd558955a5c7638fba2d68f464/YM/Next/FREE_KERNELS/Lifts.lean\#L121-L143}{\texttt{Double.duplex}},
\href{https://github.com/yang-mills-group/OpenAI-YM-20261028/blob/3456e5e75cf8d4cd558955a5c7638fba2d68f464/YM/Next/FREE_KERNELS/Sell.lean\#L43}{\texttt{Flight.fund}},
\href{https://github.com/yang-mills-group/OpenAI-YM-20261028/blob/3456e5e75cf8d4cd558955a5c7638fba2d68f464/YM/Next/FREE_KERNELS/Sell.lean\#L215-L227}{\texttt{Flight.banked\_dollars}}.
Quantitative control:
\href{https://github.com/yang-mills-group/OpenAI-YM-20261028/blob/3456e5e75cf8d4cd558955a5c7638fba2d68f464/YM/Next/FREE_KERNELS/Lifts.lean\#L145-L154}{\texttt{duplex\_stock}}.

\textbf{Lemma 4.12 (amended defect: identity (Δ)).} In the setting of
Lemma 4.10, take fine amendment \(\mu\ge0\) and coarse amendment
\(\lambda\ge0\). Then (Δ) holds on every pair of tori, with
\(P=\mathcal P_{\mu,\lambda}\) where
\[\mathcal P_{\mu,\lambda}=\mathcal T_\lambda^*\big(C+\lambda\,1-\mu\,S^*S\big)\mathcal T_\lambda .\]
Here \(S\) is the 3-form transfer kernel of the flight, with
\(S\,d_2D'^{-1}=d_2D^{-1}T\), and \(S^*S\) has strip bounds
\(O(L^{-2})\). \emph{Sketch:} Write \(X=d_2D^{-1}\) and
\(X'=d_2D'^{-1}\). Then
\[H'^{-1}-T^*H^{-1}T-\kappa^{-1}dd^*=\mathcal V+\lambda X'^*X'-\mu\,T^*X^*XT .\]
Since \(XT=SX'\), the last term is \(-X'^*(\mu S^*S)X'\); with Lemma
4.11 the whole expression is \(X'^*(C+\lambda-\mu S^*S)X'\). Conjugate
by \(H'\) and use \(X'H'=\mathcal T_\lambda d_2\) (Lemma 4.7).
\emph{Lean:}
\href{https://github.com/yang-mills-group/OpenAI-YM-20261028/blob/3456e5e75cf8d4cd558955a5c7638fba2d68f464/YM/Next/FREE_KERNELS/Eject.lean\#L144-L181}{\texttt{Flight.rebased}},
\href{https://github.com/yang-mills-group/OpenAI-YM-20261028/blob/3456e5e75cf8d4cd558955a5c7638fba2d68f464/YM/Next/FREE_KERNELS/Eject.lean\#L27-L29}{\texttt{Flight.rollFunds}},
\href{https://github.com/yang-mills-group/OpenAI-YM-20261028/blob/3456e5e75cf8d4cd558955a5c7638fba2d68f464/YM/Next/FREE_KERNELS/Eject.lean\#L23-L25}{\texttt{Flight.render}},
\href{https://github.com/yang-mills-group/OpenAI-YM-20261028/blob/3456e5e75cf8d4cd558955a5c7638fba2d68f464/YM/Next/FREE_KERNELS/Sell.lean\#L136-L149}{\texttt{Flight.intake}},
\href{https://github.com/yang-mills-group/OpenAI-YM-20261028/blob/3456e5e75cf8d4cd558955a5c7638fba2d68f464/YM/Next/FREE_KERNELS/Curl.lean\#L90}{\texttt{Flight.sale}},
\href{https://github.com/yang-mills-group/OpenAI-YM-20261028/blob/3456e5e75cf8d4cd558955a5c7638fba2d68f464/YM/Next/FREE_KERNELS/Eject.lean\#L34-L41}{\texttt{Flight.netted}},
\href{https://github.com/yang-mills-group/OpenAI-YM-20261028/blob/3456e5e75cf8d4cd558955a5c7638fba2d68f464/YM/Next/FREE_KERNELS/Rack.lean\#L57-L58}{\texttt{check\_bleep}}.

\textbf{Lemma 4.13 (positivity; III.3.4).} The constants and bounds of
III.3.4 hold uniformly over flights of the fixed index \((j,t)\).
\emph{Sketch:}

\begin{itemize}
\tightlist
\item
  The strip bound on \(C\) bounds its entries at real momenta, so
  \(C\ge-C_1\).
\item
  \(\mu S^*S\le\mu B/L^2\le1\).
\item
  Hence \(C+\lambda-\mu S^*S\ge1\) for \(\lambda\ge C_0=C_1+2\).
\item
  \(\mathcal T_\lambda^{-1}=1+\lambda X'd_2^*\), with \(X'=d_2D'^{-1}\),
  has bounded entries, so \(|x|^2\le W|\mathcal T_\lambda x|^2\), and
  \(\hat{\mathcal P}\ge W^{-1}\).
\item
  The upper bound and the strip bound follow from products of
  strip-bounded kernels.
\end{itemize}

\emph{Lean:}
\href{https://github.com/yang-mills-group/OpenAI-YM-20261028/blob/3456e5e75cf8d4cd558955a5c7638fba2d68f464/YM/Next/FREE_KERNELS/Eject.lean\#L58-L95}{\texttt{Flight.absorb\_before}},
\href{https://github.com/yang-mills-group/OpenAI-YM-20261028/blob/3456e5e75cf8d4cd558955a5c7638fba2d68f464/YM/Next/FREE_KERNELS/Eject.lean\#L99-L140}{\texttt{Flight.frontier}}.

\paragraph{4.5 The catalogue}\label{the-catalogue}

\textbf{Lemma 4.14 (III.3.1).} For every kernel choice \(K\), a
catalogue satisfying III.3.1 exists. \emph{Sketch:}

\begin{itemize}
\tightlist
\item
  \emph{One hop.} For each \((j,t)\) there is \(\lambda_0(j,t)\) with
  the following property: for every coarse amendment
  \(\lambda\ge\lambda_0\) and fine amendment \(\mu\ge0\), there are
  rates, constants and an integer \(b>\mu B(j,t)\) for which III.3.1(b)
  holds. The proof combines Lemma 4.13 with Lemma 4.9 at \(e=1,2\).
\item
  \emph{Assembly.} Set
  \(\lambda_j=\max(\lambda_0(j,j),\lambda_0(j+1,j))\): each stage's
  amendment acts as the coarse amendment of the hops \((j,j)\) and
  \((j+1,j)\), and as the fine amendment of \((j,j)\) and \((j,j-1)\).
  Then:

  \begin{itemize}
  \tightlist
  \item
    \(\alpha_j\) is the minimum of the rates of the bottom quote (Lemma
    4.8 with \(\lambda_j\)) and of the hops \((j,j)\), \((j,j-1)\);
  \item
    \(\Omega_j\) is the corresponding maximum, and \(\eta_j\) the
    corresponding minimum;
  \item
    \(b_j=\max(b(j,j),b(j,j-1))\);
  \item
    weaker constants are always admissible.
  \end{itemize}
\item
  The block size enters only through \(b_j\), which guarantees
  \(\lambda_jB\le L^2\). The stronger choice
  \(L\ge b_j>\lambda_jB(j,t)\) suffices for each admissible \(t\).
\end{itemize}

\emph{Lean:}
\href{https://github.com/yang-mills-group/OpenAI-YM-20261028/blob/3456e5e75cf8d4cd558955a5c7638fba2d68f464/YM/Next/FREE_KERNELS/Priced.lean\#L46-L65}{\texttt{buys}},
\href{https://github.com/yang-mills-group/OpenAI-YM-20261028/blob/3456e5e75cf8d4cd558955a5c7638fba2d68f464/YM/Next/FREE_KERNELS/Priced.lean\#L67-L93}{\texttt{pledges}}.

\paragraph{4.6 Along a schedule: pre-repair, repair,
output}\label{along-a-schedule-pre-repair-repair-output}

\textbf{Lemma 4.15 (pre-repair family).} Fix a catalogue of III.3.1 and
a trajectory with its kernel, satisfying the credit condition. Define:

\begin{itemize}
\tightlist
\item
  \(\mathcal B^{N,i}\): the metric data of Lemma 4.8 for the history of
  level \(i\) in mesh \(N\), with the clip of level \(i\) (allowed by
  (C1));
\item
  \(W^{N,i},T^{N,i}\): those of the flight of step \(i\), of index
  \((J(i+1),J(i))\) (allowed since \(J(i)\le J(i+1)\le J(i)+1\)).
\end{itemize}

This family satisfies the pre-repair conditions of III.3.2(2):

\begin{itemize}
\tightlist
\item
  the bounds, from III.3.1;
\item
  (D1)--(D4), from Lemma 4.10;
\item
  (Δ), with \(P\) the rolled kernel;
\item
  the positivity of \(P\), from III.3.1(b) and (C2).
\end{itemize}

\emph{Lean:}
\href{https://github.com/yang-mills-group/OpenAI-YM-20261028/blob/3456e5e75cf8d4cd558955a5c7638fba2d68f464/YM/Next/FREE_KERNELS/Rack.lean\#L90-L120}{\texttt{shelves}},
\href{https://github.com/yang-mills-group/OpenAI-YM-20261028/blob/3456e5e75cf8d4cd558955a5c7638fba2d68f464/YM/Next/FREE_KERNELS/Rack.lean\#L64-L65}{\texttt{stageHop}},
\href{https://github.com/yang-mills-group/OpenAI-YM-20261028/blob/3456e5e75cf8d4cd558955a5c7638fba2d68f464/YM/Next/FREE_KERNELS/Rack.lean\#L73}{\texttt{location}},
\href{https://github.com/yang-mills-group/OpenAI-YM-20261028/blob/3456e5e75cf8d4cd558955a5c7638fba2d68f464/YM/Next/FREE_KERNELS/Rack.lean\#L74}{\texttt{destination}}.

\textbf{Lemma 4.16 (repair; III.10 with III.6's theorems).} For every
sheet and every table \(X\) there is a table \(X'\) with the following
property. Every pre-repair family for \(X\) (for a trajectory with the
sheet's profile and kernel) has a repaired family with the same metric
data and the same \(W,T\), which satisfies the output conditions for
\(X'\). For step \(i\) the construction is:

\begin{enumerate}
\def\labelenumi{\arabic{enumi}.}
\tightlist
\item
  \textbf{F06} applied to \(P\) from (Δ): a kernel \(P_1\) with
  two-sided bounds at real momenta, a strip bound, and
  \(d_2^*P_1d_2=d_2^*Pd_2\) on every torus.
\item
  \textbf{F07}: a symmetric decaying square root, \(P_1=P_2^*P_2\).
\item
  \textbf{F08}: kernels with the same \(W,T\) and
  \[M_f=V H^{-1}T\,V'^*,\qquad M_n=\kappa^{-1/2}d^*V'^*,\qquad O=\big(1-\Pi,\ P_2\,d_2H'^{-1}V'^*\big),\qquad \Pi=V'H'^{-1}V'^*,\]
  with the strip and derivative bounds of III.3.2(3). The constants
  depend only on \(X\), \(\kappa\) and the sheet's direct clip
  coefficients.
\item
  \textbf{F09} (algebra) gives (R1)--(R6). Since \(V'^*V'=H'\), \(\Pi\)
  is the orthogonal projection onto the range of \(V'\). Then:

  \begin{itemize}
  \tightlist
  \item
    (R1):
    \(M_f^*M_f+M_n^*M_n+O^*O=V'(T^*H^{-1}T+\kappa^{-1}dd^*+H'^{-1}d_2^*P_1d_2H'^{-1})V'^*+1-\Pi=1\),
    by (Δ);
  \item
    (R2) and (R4) follow from \(V'^*V'=H'\);
  \item
    (R3) follows from (D4) and \(H'd=D'd\), \(Hd=Dd\);
  \item
    (R5) follows from \((1-\Pi)V'=0\) and \(d_2d=0\);
  \item
    (R6) follows from (D3).
  \end{itemize}
\end{enumerate}

The second table depends on the sheet only through the direct clip
coefficients, which enter \(V,V'\). It can be chosen with
\(\alpha'_j\le\alpha_j\), \(\Omega'_j\ge\Omega_j\) and unchanged
\(\eta_j\). \emph{Lean:}
\href{https://github.com/yang-mills-group/OpenAI-YM-20261028/blob/3456e5e75cf8d4cd558955a5c7638fba2d68f464/YM/Bid/Craft.lean\#L107-L166}{\texttt{Bid.pinning}},
\href{https://github.com/yang-mills-group/OpenAI-YM-20261028/blob/3456e5e75cf8d4cd558955a5c7638fba2d68f464/YM/Bid/Craft.lean\#L84-L85}{\texttt{F06}},
\href{https://github.com/yang-mills-group/OpenAI-YM-20261028/blob/3456e5e75cf8d4cd558955a5c7638fba2d68f464/YM/Bid/Craft.lean\#L90-L91}{\texttt{F07}},
\href{https://github.com/yang-mills-group/OpenAI-YM-20261028/blob/3456e5e75cf8d4cd558955a5c7638fba2d68f464/YM/Bid/Craft.lean\#L93-L96}{\texttt{F08}},
\href{https://github.com/yang-mills-group/OpenAI-YM-20261028/blob/3456e5e75cf8d4cd558955a5c7638fba2d68f464/YM/Bid/Craft.lean\#L99-L101}{\texttt{F09}},
\href{https://github.com/yang-mills-group/OpenAI-YM-20261028/blob/3456e5e75cf8d4cd558955a5c7638fba2d68f464/YM/Bid/Craft.lean\#L53-L58}{\texttt{Replaces}}.
\emph{(repair proof and F08--F09 proofs read)}

\textbf{Lemma 4.17 (III.3.2 and its variant).} The catalogue and repair
construction give the families and constants asserted in III.3.2,
including its preliminary-sheet variant. \emph{Sketch:} Take the
catalogue (4.14), then the second table from the repair theorem with
III.6's theorems (4.16), then the pre-repair family (4.15), then its
repair. The explicit formulas and the symmetry carry over because the
metric data and \(W,T\) are unchanged. In the variant, the payment
record of the schedule turns sheet-level bounds into (C1)--(C2). A
preliminary sheet is handled by re-running the construction for a copy
of the trajectory with the preliminary profile; this preserves the
histories and the direct clip coefficients. \emph{Lean:}
\href{https://github.com/yang-mills-group/OpenAI-YM-20261028/blob/3456e5e75cf8d4cd558955a5c7638fba2d68f464/YM/Next/FREE_KERNELS/Rack.lean\#L134-L145}{\texttt{retained\_support}},
\href{https://github.com/yang-mills-group/OpenAI-YM-20261028/blob/3456e5e75cf8d4cd558955a5c7638fba2d68f464/YM/Next/FREE_KERNELS/Rack.lean\#L121-L126}{\texttt{retained\_presents}},
\href{https://github.com/yang-mills-group/OpenAI-YM-20261028/blob/3456e5e75cf8d4cd558955a5c7638fba2d68f464/YM/Next/FREE_KERNELS/Calendar.lean\#L108-L124}{\texttt{before\_paid}},
\href{https://github.com/yang-mills-group/OpenAI-YM-20261028/blob/3456e5e75cf8d4cd558955a5c7638fba2d68f464/YM/Next/FREE_KERNELS/Calendar.lean\#L95-L102}{\texttt{credit\_from\_sheet}}.

\paragraph{4.7 All tori}\label{all-tori}

\textbf{Lemma 4.18 (rigidity of decaying kernels).} Suppose two finite
sums of products of decaying kernels, or of cross kernels of the same
block size, define equal matrices on tori \(\Lambda_{c_k}\)
(respectively on pairs \((c_k,c_kL)\)) with \(c_k\to\infty\). Then they
are equal on every torus (pair). \emph{Sketch:} Sums and products of
decaying kernels are decaying kernels, so each side is a single kernel.
A summable kernel is the limit of its periodisations as the period tends
to infinity. Dominated convergence justifies this coefficient limit; the
schedule sides of each fixed step tend to infinity. \emph{Lean:}
\href{https://github.com/yang-mills-group/OpenAI-YM-20261028/blob/3456e5e75cf8d4cd558955a5c7638fba2d68f464/YM/Next/FREE_KERNELS/Bank.lean\#L49-L67}{\texttt{entry\_limit}},
\href{https://github.com/yang-mills-group/OpenAI-YM-20261028/blob/3456e5e75cf8d4cd558955a5c7638fba2d68f464/YM/UVNext/PRICED_KERNELS/Echo.lean\#L77-L81}{\texttt{Echo.all}},
\href{https://github.com/yang-mills-group/OpenAI-YM-20261028/blob/3456e5e75cf8d4cd558955a5c7638fba2d68f464/YM/UVNext/PRICED_KERNELS/Echo.lean\#L105-L116}{\texttt{Echo.allAcross}},
\href{https://github.com/yang-mills-group/OpenAI-YM-20261028/blob/3456e5e75cf8d4cd558955a5c7638fba2d68f464/YM/UVNext/PRICED_KERNELS/Echo.lean\#L38}{\texttt{Echo.Sound}},
\href{https://github.com/yang-mills-group/OpenAI-YM-20261028/blob/3456e5e75cf8d4cd558955a5c7638fba2d68f464/YM/UVNext/PRICED_KERNELS/Echo.lean\#L118-L119}{\texttt{Echo.schedule\_sides}}.

\textbf{Lemma 4.19 (III.3.3).} Every output family of III.3.2 satisfies
(D1)--(D4) and (R1)--(R6) on every torus pair of the same block ratio.
\emph{Sketch:}

\begin{itemize}
\tightlist
\item
  (R1)--(R5) are identities between finite sums of products of kernels,
  with \(M_f\) and \(W\) as cross kernels. Apply 4.18 on the schedule
  tori.
\item
  (R6) contains \(Q\). Given (D3) and (R2) at \((c,f)\), (R6) is
  equivalent to \((V\,d)^*M_f=d^*(T\,V'^*)\), which is again an identity
  between kernels.
\item
  (D1)--(D4) at arbitrary \((c,f)\) follow by the same argument, because
  \(Q\) is itself given by a cross kernel (III.10). They also follow
  directly from Lemma 4.10 for any family realising the explicit
  formulas.
\end{itemize}

\emph{Lean:}
\href{https://github.com/yang-mills-group/OpenAI-YM-20261028/blob/3456e5e75cf8d4cd558955a5c7638fba2d68f464/YM/UVNext/PRICED_KERNELS/Tori.lean\#L109-L157}{\texttt{complete\_roll}},
\href{https://github.com/yang-mills-group/OpenAI-YM-20261028/blob/3456e5e75cf8d4cd558955a5c7638fba2d68f464/YM/UVNext/PRICED_KERNELS/Tori.lean\#L87-L98}{\texttt{orthogonal\_iff}},
\href{https://github.com/yang-mills-group/OpenAI-YM-20261028/blob/3456e5e75cf8d4cd558955a5c7638fba2d68f464/YM/PhysicalNext/RG_NEXT_COMPLETE_REPEATABLE_NONLINEAR_STEP/Pane/Purl.lean\#L20-L53}{\texttt{panes}},
\href{https://github.com/yang-mills-group/OpenAI-YM-20261028/blob/3456e5e75cf8d4cd558955a5c7638fba2d68f464/YM/UVNext/PRICED_KERNELS/Tori.lean\#L46-L60}{\texttt{universal\_diagonal}}.

\paragraph{4.8 Free covariance of the curvature
(III.3.6)}\label{free-covariance-of-the-curvature-iii.3.6}

\textbf{Lemma 4.20 (lower bound).} The coercivity bound of III.3.6(a)
holds uniformly over firm kernels. \emph{Sketch:}

\begin{itemize}
\tightlist
\item
  \(\sigma_1^*\sigma_1+\sigma_0\sigma_0^*=|\tau|^2\,1\) on 1-form
  symbols; this is \(d^*d+dd^*=-\Delta\) in Fourier space.
\item
  A firm \(f\) satisfies
  \(\mathrm{Re}\,\langle x,\hat f(p)x\rangle\ge\eta|x|^2\) at every real
  \(p\): the torus bounds give it at the frequencies
  \(2\pi\mathbb Z^4/n\), and these are dense; conclude by continuity.
\item
  Hence
  \(\mathrm{Re}\,\langle x,\Phi_f(p)x\rangle\ge\min(\eta,1)|\tau(p)|^2|x|^2\).
\item
  \(|\tau(p)|^2\ge c_0|p|^2\) for \(|p_a|\le3\pi/2\), since
  \(e^{ir}\ne1\) for \(0<|r|\le3\pi/2\) (compactness).
\end{itemize}

\emph{Lean:}
\href{https://github.com/yang-mills-group/OpenAI-YM-20261028/blob/3456e5e75cf8d4cd558955a5c7638fba2d68f464/YM/Further/FREE_COVARIANCE_KP/KP/Core.lean\#L41-L57}{\texttt{blade\_id}},
\href{https://github.com/yang-mills-group/OpenAI-YM-20261028/blob/3456e5e75cf8d4cd558955a5c7638fba2d68f464/YM/Further/FREE_COVARIANCE_KP/KP/Core.lean\#L68-L79}{\texttt{ion}},
\href{https://github.com/yang-mills-group/OpenAI-YM-20261028/blob/3456e5e75cf8d4cd558955a5c7638fba2d68f464/YM/Further/FREE_COVARIANCE_KP/KP/Floor.lean\#L94-L100}{\texttt{real\_price}},
\href{https://github.com/yang-mills-group/OpenAI-YM-20261028/blob/3456e5e75cf8d4cd558955a5c7638fba2d68f464/YM/Further/FREE_COVARIANCE_KP/KP/Core.lean\#L115-L125}{\texttt{tank\_full}}.

\textbf{Lemma 4.21 (smooth symbols).} This is III.3.6(c). \emph{Sketch:}
Differentiate the Fourier series termwise and use
\(|y|_1^n\le A_ne^{\alpha|y|_1}\), with \(A_n\) as in III.3.6(c).
\emph{Lean:}
\href{https://github.com/yang-mills-group/OpenAI-YM-20261028/blob/3456e5e75cf8d4cd558955a5c7638fba2d68f464/YM/Further/FREE_COVARIANCE_KP/KP/Plain.lean\#L101-L133}{\texttt{radial\_prices}}.

\textbf{Lemma 4.22 (comparison with the continuum).} This is III.3.6(b).
Let \(S^k(U)\) be the class of families with
\(\|D^ng\|\le C_n|p|^{k-n}\) on \(U\), uniformly over the family. Sums
preserve a fixed order, products add orders, and inversion reverses the
order when the rescaled matrix has a uniform coercivity bound.
\emph{Sketch:}

\begin{itemize}
\tightlist
\item
  \(\tau-ip\in S^2\) (Taylor at \(0\)), and \(\hat f-1\in S^1\) (since
  \(\hat f(0)=1\), and by Lemma 4.21).
\item
  Hence \(\Phi_f\in S^2\) and \(\Phi_f-|p|^2\,1\in S^3\).
\item
  By (a) and the resolvent identity, \(\Phi_f^{-1}\in S^{-2}\) and
  \(\Phi_f^{-1}-|p|^{-2}\,1\in S^{-1}\).
\item
  Hence \(G_f\in S^0\) and \(G_f-G_\infty\in S^1\).
\end{itemize}

\emph{Lean:}
\href{https://github.com/yang-mills-group/OpenAI-YM-20261028/blob/3456e5e75cf8d4cd558955a5c7638fba2d68f464/YM/Further/FREE_COVARIANCE_KP/KP/Alloy.lean\#L110-L120}{\texttt{gem\_quote}},
\href{https://github.com/yang-mills-group/OpenAI-YM-20261028/blob/3456e5e75cf8d4cd558955a5c7638fba2d68f464/YM/Further/FREE_COVARIANCE_KP/KP/Weight.lean\#L39-L41}{\texttt{Weight}},
\href{https://github.com/yang-mills-group/OpenAI-YM-20261028/blob/3456e5e75cf8d4cd558955a5c7638fba2d68f464/YM/Further/FREE_COVARIANCE_KP/KP/Gem.lean\#L17}{\texttt{Value}}.
\emph{(proof read, including the symbol-class product and inverse
estimates)}

\paragraph{4.9 Calculus lemmas (III.3.7)}\label{calculus-lemmas-iii.3.7}

\textbf{Lemma 4.23 (nonlinear calculus).} The statements of III.3.7 hold
under their respective hypotheses. \emph{Sketch:}

\begin{itemize}
\tightlist
\item
  \textbf{(a)} Split at \(|x|_1=L\).

  \begin{itemize}
  \tightlist
  \item
    For \(|x|_1\le L\): localised Newton remainder estimate,
    \(\|p(x)-\sum_{i<k}P_i(x)\|\le M'(1+|x|_1)^k\) with \(M'\) the bound
    on \(k\)-th differences.
  \item
    For \(|x|_1>L\): bound \(p\) and the pieces directly, using
    \(L^{-1-i}(1+|x|)^i\le L^{-1-k}(1+|x|)^b\) when
    \(L^{-1}(1+|x|)\ge1\).
  \end{itemize}
\item
  \textbf{(b)} Substitute \(z\mapsto\sum_vt^{Y_v}p_vz\) in the Taylor
  polynomial of \(f\) and read off each coefficient of \(t^r\) for
  \(r\le d\). Rescaling \(z\) by an independent scalar separates the
  ordinary Taylor degrees, and polarization recovers the multilinear
  coefficients. All weighted components of weight \(\le d\) of every
  Taylor coefficient vanish. A component of weight \(>d\) vanishes on
  the given inputs by the support condition.
\item
  \textbf{(c)} Locality restricts the logarithmic link fields to links
  with both endpoints in \(X\). Gauge invariance identifies \(u\) near
  the identity with its pullback by the map taking a site gauge and a
  radial connection jet to gauged logarithmic holonomies. Part (b) gives
  weighted vanishing for this pullback. The derivative of that map
  covers every polynomial link field of degree at most \(p-1\): part (d)
  supplies the highest degree, line integration changes it only by lower
  degrees, and induction removes the error. Its higher derivatives
  preserve the total weight. Faà di Bruno and induction on the number of
  inputs therefore transfer vanishing to \(u\circ\exp\). Non-wrapping
  and \(n\ge2\) ensure that the centred lifts of adjacent sites differ
  by a unit vector.
\item
  \textbf{(d)} \(r=(j+2)^{-1}\sum_ax_aQ_a\) is homogeneous of degree
  \(j+2\), and Euler's identity gives \(\sum_ax_a(Q_a-\partial_ar)=0\).
\item
  \textbf{(e)} On the unit ball,
  \(\|g(sv)-g(0)\|\le s\|g(v)-g(0)\|+Q(1-s)\|v\|^2\) for the matrix
  \(g\) of the plaquette holonomy of four link variables (second-order
  Taylor bound). Hence the chord distance at \(sv\) is at most
  \(s\tau+(1-s)\tau\).
\end{itemize}

\emph{Lean:} see III.3.7;
\href{https://github.com/yang-mills-group/OpenAI-YM-20261028/blob/3456e5e75cf8d4cd558955a5c7638fba2d68f464/YM/MultiNext/SIGNED_SOURCE_CONTINUATION/Mom.lean\#L43-L52}{\texttt{Maps}},
\href{https://github.com/yang-mills-group/OpenAI-YM-20261028/blob/3456e5e75cf8d4cd558955a5c7638fba2d68f464/YM/MultiNext/SIGNED_SOURCE_CONTINUATION/Routes.lean\#L63-L127}{\texttt{covers}},
\href{https://github.com/yang-mills-group/OpenAI-YM-20261028/blob/3456e5e75cf8d4cd558955a5c7638fba2d68f464/YM/MultiNext/SIGNED_SOURCE_CONTINUATION/Mediator.lean\#L70-L98}{\texttt{attenuation}},
\href{https://github.com/yang-mills-group/OpenAI-YM-20261028/blob/3456e5e75cf8d4cd558955a5c7638fba2d68f464/YM/Source/Local.lean\#L74-L75}{\texttt{testLinks}}.
\emph{(proof read for (a)--(e), including the moment-space definition
and transfer argument)}

\paragraph{4.10 Auxiliary results}\label{auxiliary-results}

Three auxiliary packages concern local observables and equation spaces:

\begin{itemize}
\tightlist
\item
  \textbf{Corrected observables.} An exact identity recovering
  holonomy-jet probes from the coefficients of corrected observables
  (\href{https://github.com/yang-mills-group/OpenAI-YM-20261028/blob/3456e5e75cf8d4cd558955a5c7638fba2d68f464/YM/Further/CORRECTED_LIFT_COEFFICIENTS/Original.lean\#L134-L146}{\texttt{Screen.reproduces}}).
  Used by the rank-covariance part of III.11, by VII.2 and by VII.4.
\item
  \textbf{Equation projections.} Orthogonal projections onto a
  finite-dimensional equation subspace and its complement
  (\href{https://github.com/yang-mills-group/OpenAI-YM-20261028/blob/3456e5e75cf8d4cd558955a5c7638fba2d68f464/YM/Push/CUTOFF_PROJECTORS/Frame.lean\#L104-L105}{\texttt{drop}},
  \href{https://github.com/yang-mills-group/OpenAI-YM-20261028/blob/3456e5e75cf8d4cd558955a5c7638fba2d68f464/YM/Push/CUTOFF_PROJECTORS/Frame.lean\#L108-L109}{\texttt{retain}}).
  The additional criterion
  \href{https://github.com/yang-mills-group/OpenAI-YM-20261028/blob/3456e5e75cf8d4cd558955a5c7638fba2d68f464/YM/Push/CUTOFF_PROJECTORS/Bank.lean\#L161-L169}{\texttt{spectral\_bank}}
  gives spectral separation under construction and comparison
  hypotheses, exact top-jet tests, and approximation by a fixed
  triangular model on finitely many probes.
\item
  \textbf{Local jet vanishing.} Vanishing of test-jet derivatives below
  the engineering degree for non-unit basis observables, assuming their
  local section is readable
  (\href{https://github.com/yang-mills-group/OpenAI-YM-20261028/blob/3456e5e75cf8d4cd558955a5c7638fba2d68f464/YM/ProofPorts/GRID_RENEWED_RESERVES/Speeds.lean}{\texttt{Speeds.lean}},
  \href{https://github.com/yang-mills-group/OpenAI-YM-20261028/blob/3456e5e75cf8d4cd558955a5c7638fba2d68f464/YM/ProofPorts/GRID_RENEWED_RESERVES/Speeds.lean\#L37-L49}{\texttt{canon\_first}}).
  Used in the nonlinear step of III.10.
\end{itemize}

\emph{(statements read; the coefficient identity and jet-vanishing
proofs read)}

\subsubsection{5. Where the difficulty
lies}\label{where-the-difficulty-lies-5}

\begin{itemize}
\tightlist
\item
  \textbf{Uniformity.} The algebraic identities must be accompanied by
  bounds uniform over histories of any length and over all block sizes,
  with the kernel choice and stage index fixed. The mechanism is in
  Lemmas 4.2--4.3:

  \begin{itemize}
  \tightlist
  \item
    geometric dilation along the history;
  \item
    a lower bound on the principal alias, with an exponent linear in
    \(M/L\);
  \item
    \(J\)-fold smoothing, which gives \(\prod_a(1+|\ell_a|)^{-J}\) alias
    decay and therefore \(J-2\) derivatives on \(W,T\).
  \end{itemize}
\item
  \textbf{Positivity is proved only at real momenta.} \(|\tau|^{-2}\) is
  never estimated on a strip. Every symbol is put in a filled form
  (Lemma 4.6) whose denominators are bounded below at real momenta, and
  inverses are taken on a thinner complex cube. The estimates use the
  Banach algebra of exponentially decaying kernels and its analytic
  symbols.
\item
  \textbf{Unusual features.}

  \begin{itemize}
  \tightlist
  \item
    The amendment \(\lambda_j\) (Lemma 4.7) changes the inverse metric.
    It modifies the metric only off exact 2-forms, so \(d^*Hd\) is
    unchanged, and it exists only to make the defect's 3-form factor
    positive.
  \item
    \(L\) enters the catalogue only through \(\lambda_jB\le L^2\).
  \item
    The second table depends on the sheet through its direct clip
    coefficients.
  \item
    The factorisation of the defect (Lemma 4.11) needs the identities on
    all tori, because of the harmonic forms.
  \item
    For explicit families,
    \href{https://github.com/yang-mills-group/OpenAI-YM-20261028/blob/3456e5e75cf8d4cd558955a5c7638fba2d68f464/YM/UVNext/PRICED_KERNELS/Tori.lean\#L46-L60}{\texttt{universal\_diagonal}}
    gives the diagonal identities directly. The output conditions alone
    give them by rigidity.
  \end{itemize}
\item
  \textbf{Units.} Decay is in the \(\ell^1\) norm, on the lattice where
  the kernel acts. For \(W,T,M_f\) that is the coarse lattice, so the
  fine-lattice rate is \(\alpha/L\), with amplitudes \(L^{-1}\),
  \(L^{-2}\) and an extra \(L^{-1}\) per fine difference. These are the
  scalings of Balaban's background field: smooth on scale \(L\).
\item
  \textbf{Statement strength.} III.3.2 is an existence statement whose
  output is a package of bounds and identities. Later units use
  (R1)--(R6), the floors and the strip bounds. The repaired family
  retains the pre-repair metric data and \(W,T\), so the defect identity
  and positivity remain available, although they are not fields of the
  output conditions.
\end{itemize}

\subsubsection{6. Notation}\label{notation-5}

{\def\LTcaptype{none} % do not increment counter
\begin{longtable}[]{@{}
  >{\raggedright\arraybackslash}p{(\linewidth - 4\tabcolsep) * \real{0.3333}}
  >{\raggedright\arraybackslash}p{(\linewidth - 4\tabcolsep) * \real{0.3333}}
  >{\raggedright\arraybackslash}p{(\linewidth - 4\tabcolsep) * \real{0.3333}}@{}}
\toprule\noalign{}
\begin{minipage}[b]{\linewidth}\raggedright
Written proof
\end{minipage} & \begin{minipage}[b]{\linewidth}\raggedright
Lean (glossary)
\end{minipage} & \begin{minipage}[b]{\linewidth}\raggedright
Meaning
\end{minipage} \\
\midrule\noalign{}
\endhead
\bottomrule\noalign{}
\endlastfoot
\(R\), \(\hat R\), \(R_n\); \(\mathcal S(\alpha,\Omega)\) &
\href{https://github.com/yang-mills-group/OpenAI-YM-20261028/blob/3456e5e75cf8d4cd558955a5c7638fba2d68f464/YM/CellFourier/Circulant.lean\#L22}{\texttt{Row}},
\href{https://github.com/yang-mills-group/OpenAI-YM-20261028/blob/3456e5e75cf8d4cd558955a5c7638fba2d68f464/YM/CellFourier/Circulant.lean\#L26}{\texttt{Row.strip}}
& translation-invariant kernel, symbol, action on \(\Lambda_n\); strip
bound \\
cross kernel &
\href{https://github.com/yang-mills-group/OpenAI-YM-20261028/blob/3456e5e75cf8d4cd558955a5c7638fba2d68f464/YM/CellFourier/Shuttle.lean\#L26}{\texttt{Cross}},
\href{https://github.com/yang-mills-group/OpenAI-YM-20261028/blob/3456e5e75cf8d4cd558955a5c7638fba2d68f464/YM/CellFourier/Shuttle.lean\#L35}{\texttt{installed}}
& coarse-to-fine kernel and its action \\
\(\mathrm{plank}_M,\mathrm{rail}_M,a_L,a_L^J\) &
\href{https://github.com/yang-mills-group/OpenAI-YM-20261028/blob/3456e5e75cf8d4cd558955a5c7638fba2d68f464/YM/Next/FREE_KERNELS/Lines.lean\#L38}{\texttt{plank}},
\href{https://github.com/yang-mills-group/OpenAI-YM-20261028/blob/3456e5e75cf8d4cd558955a5c7638fba2d68f464/YM/Next/FREE_KERNELS/Lines.lean\#L34}{\texttt{rail}},
\href{https://github.com/yang-mills-group/OpenAI-YM-20261028/blob/3456e5e75cf8d4cd558955a5c7638fba2d68f464/YM/Next/FREE_KERNELS/Train.lean\#L18}{\texttt{area}},
\href{https://github.com/yang-mills-group/OpenAI-YM-20261028/blob/3456e5e75cf8d4cd558955a5c7638fba2d68f464/YM/Next/FREE_KERNELS/Feed.lean\#L23}{\texttt{massFilter}}
& line, box and \(J\)-fold box filters \\
history \(v\), dilation \(M_v\), filter \(s_v\) &
\href{https://github.com/yang-mills-group/OpenAI-YM-20261028/blob/3456e5e75cf8d4cd558955a5c7638fba2d68f464/YM/Next/FREE_KERNELS/Train.lean\#L22-L24}{\texttt{Train}},
\href{https://github.com/yang-mills-group/OpenAI-YM-20261028/blob/3456e5e75cf8d4cd558955a5c7638fba2d68f464/YM/Next/FREE_KERNELS/Train.lean\#L36}{\texttt{Train.smooth}}
& finite history of block averages \\
flight of index \((j,t)\) &
\href{https://github.com/yang-mills-group/OpenAI-YM-20261028/blob/3456e5e75cf8d4cd558955a5c7638fba2d68f464/YM/Next/FREE_KERNELS/Signal.lean\#L22-L29}{\texttt{Flight}},
\href{https://github.com/yang-mills-group/OpenAI-YM-20261028/blob/3456e5e75cf8d4cd558955a5c7638fba2d68f464/YM/Next/FREE_KERNELS/Signal.lean\#L35}{\texttt{Flight.land}}
& one step: block size, fine history, coarse history \\
\(D_a\); \(\rho\); \(\beta_a\), \(h\) &
\href{https://github.com/yang-mills-group/OpenAI-YM-20261028/blob/3456e5e75cf8d4cd558955a5c7638fba2d68f464/YM/Next/FREE_KERNELS/Feed.lean\#L133-L134}{\texttt{Train.levelD}},
\href{https://github.com/yang-mills-group/OpenAI-YM-20261028/blob/3456e5e75cf8d4cd558955a5c7638fba2d68f464/YM/Next/FREE_KERNELS/Feed.lean\#L154-L156}{\texttt{D\_feedback}};
\href{https://github.com/yang-mills-group/OpenAI-YM-20261028/blob/3456e5e75cf8d4cd558955a5c7638fba2d68f464/YM/Next/FREE_KERNELS/Bend.lean\#L46}{\texttt{Train.rho}};
\href{https://github.com/yang-mills-group/OpenAI-YM-20261028/blob/3456e5e75cf8d4cd558955a5c7638fba2d68f464/YM/Next/FREE_KERNELS/Bend.lean\#L39}{\texttt{Train.bulb}},
\href{https://github.com/yang-mills-group/OpenAI-YM-20261028/blob/3456e5e75cf8d4cd558955a5c7638fba2d68f464/YM/Next/FREE_KERNELS/Bend.lean\#L43-L44}{\texttt{Train.bend}}
& 1-form covariance; gauge factor; filled symbols \\
\(D\), \(D^{-1}\) &
\href{https://github.com/yang-mills-group/OpenAI-YM-20261028/blob/3456e5e75cf8d4cd558955a5c7638fba2d68f464/YM/Next/FREE_KERNELS/Curl.lean\#L54}{\texttt{pink}},
\href{https://github.com/yang-mills-group/OpenAI-YM-20261028/blob/3456e5e75cf8d4cd558955a5c7638fba2d68f464/YM/Next/FREE_KERNELS/Curl.lean\#L55}{\texttt{blue}}
& diagonal 2-form multiplier, inverse \\
\(H^{-1}\), \(H\); \(\mathcal T_\lambda\); \(\lambda\) &
\href{https://github.com/yang-mills-group/OpenAI-YM-20261028/blob/3456e5e75cf8d4cd558955a5c7638fba2d68f464/YM/Next/FREE_KERNELS/Amend.lean\#L30}{\texttt{pledge}},
\href{https://github.com/yang-mills-group/OpenAI-YM-20261028/blob/3456e5e75cf8d4cd558955a5c7638fba2d68f464/YM/Next/FREE_KERNELS/Amend.lean\#L31}{\texttt{amend}};
\href{https://github.com/yang-mills-group/OpenAI-YM-20261028/blob/3456e5e75cf8d4cd558955a5c7638fba2d68f464/YM/Next/FREE_KERNELS/Amend.lean\#L33}{\texttt{treble}}
& amended inverse metric, metric; 3-form factor; amendment \\
metric data \(\mathcal B\), \(V\) &
\href{https://github.com/yang-mills-group/OpenAI-YM-20261028/blob/3456e5e75cf8d4cd558955a5c7638fba2d68f464/YM/PREC/Stacks.lean\#L35-L40}{\texttt{Bed}},
\href{https://github.com/yang-mills-group/OpenAI-YM-20261028/blob/3456e5e75cf8d4cd558955a5c7638fba2d68f464/YM/PREC/Stacks.lean\#L49-L51}{\texttt{Bed.full}},
\href{https://github.com/yang-mills-group/OpenAI-YM-20261028/blob/3456e5e75cf8d4cd558955a5c7638fba2d68f464/YM/PREC/Stacks.lean\#L61-L68}{\texttt{Floor}}
& \((D,D^{-1},H,H^{-1},R)\); \(V=(\sqrt c\,1,R)\) \\
transfer kernels \((W,T,M_f,M_n,O)\) &
\href{https://github.com/yang-mills-group/OpenAI-YM-20261028/blob/3456e5e75cf8d4cd558955a5c7638fba2d68f464/YM/PREC/Stacks.lean\#L70-L75}{\texttt{Towers}}
(\texttt{A0}, \texttt{T}, \texttt{mf}, \texttt{mn}, \texttt{o}),
\href{https://github.com/yang-mills-group/OpenAI-YM-20261028/blob/3456e5e75cf8d4cd558955a5c7638fba2d68f464/YM/Next/FREE_KERNELS/Primal.lean\#L24-L29}{\texttt{Flight.duo}}
& background-field map, curvature transport, isometry blocks \\
(D1)--(D4); (R1)--(R6); (Δ) &
\href{https://github.com/yang-mills-group/OpenAI-YM-20261028/blob/3456e5e75cf8d4cd558955a5c7638fba2d68f464/YM/PREC/Stacks.lean\#L82-L86}{\texttt{Towers.Diagonal}},
\href{https://github.com/yang-mills-group/OpenAI-YM-20261028/blob/3456e5e75cf8d4cd558955a5c7638fba2d68f464/YM/PREC/TreeMatch.lean\#L76-L79}{\texttt{Inputs}};
\href{https://github.com/yang-mills-group/OpenAI-YM-20261028/blob/3456e5e75cf8d4cd558955a5c7638fba2d68f464/YM/PREC/Stacks.lean\#L88-L96}{\texttt{Towers.Roll}};
\href{https://github.com/yang-mills-group/OpenAI-YM-20261028/blob/3456e5e75cf8d4cd558955a5c7638fba2d68f464/YM/Bid/Craft.lean\#L32}{\texttt{Bleeps}}
& identities of one step \\
\(C\), \(S\), \(\mathcal P_{\mu,\lambda}\) &
\href{https://github.com/yang-mills-group/OpenAI-YM-20261028/blob/3456e5e75cf8d4cd558955a5c7638fba2d68f464/YM/Next/FREE_KERNELS/Sell.lean\#L43}{\texttt{Flight.fund}},
\href{https://github.com/yang-mills-group/OpenAI-YM-20261028/blob/3456e5e75cf8d4cd558955a5c7638fba2d68f464/YM/Next/FREE_KERNELS/Curl.lean\#L90}{\texttt{Flight.sale}},
\href{https://github.com/yang-mills-group/OpenAI-YM-20261028/blob/3456e5e75cf8d4cd558955a5c7638fba2d68f464/YM/Next/FREE_KERNELS/Eject.lean\#L27-L29}{\texttt{Flight.rollFunds}}
& unamended 3-form defect, 3-form transport, rolled 3-form kernel \\
derivative bounds &
\href{https://github.com/yang-mills-group/OpenAI-YM-20261028/blob/3456e5e75cf8d4cd558955a5c7638fba2d68f464/YM/Bid/Craft.lean\#L21-L22}{\texttt{Donate}},
\href{https://github.com/yang-mills-group/OpenAI-YM-20261028/blob/3456e5e75cf8d4cd558955a5c7638fba2d68f464/YM/PREC/Stacks.lean\#L111-L114}{\texttt{Towers.Deep}}
& \(\nabla_u\) bounds with \(L^{-1-\lvert u\rvert}\),
\(L^{-2-\lvert u\rvert}\) \\
positive on \(\mathrm{ran}\,\sigma_2\); two-sided at real momenta &
\href{https://github.com/yang-mills-group/OpenAI-YM-20261028/blob/3456e5e75cf8d4cd558955a5c7638fba2d68f464/YM/Bid/Craft.lean\#L34-L36}{\texttt{Reaches}},
\href{https://github.com/yang-mills-group/OpenAI-YM-20261028/blob/3456e5e75cf8d4cd558955a5c7638fba2d68f464/YM/Bid/Craft.lean\#L38-L40}{\texttt{Funds}}
& positivity conditions on 3-form kernels \\
kernel family; pre-repair; output conditions &
\href{https://github.com/yang-mills-group/OpenAI-YM-20261028/blob/3456e5e75cf8d4cd558955a5c7638fba2d68f464/YM/PREC/Stacks.lean\#L119-L121}{\texttt{Horizon}};
\href{https://github.com/yang-mills-group/OpenAI-YM-20261028/blob/3456e5e75cf8d4cd558955a5c7638fba2d68f464/YM/Bid/Craft.lean\#L63-L73}{\texttt{Craft.Rack}};
\href{https://github.com/yang-mills-group/OpenAI-YM-20261028/blob/3456e5e75cf8d4cd558955a5c7638fba2d68f464/YM/PREC/Stacks.lean\#L132-L139}{\texttt{Horizon.Suffices}}
& metric data and transfer kernels for all \((N,i)\) \\
explicit formulas; credit condition &
\href{https://github.com/yang-mills-group/OpenAI-YM-20261028/blob/3456e5e75cf8d4cd558955a5c7638fba2d68f464/YM/Next/FREE_KERNELS/Rack.lean\#L82-L85}{\texttt{Present}};
\href{https://github.com/yang-mills-group/OpenAI-YM-20261028/blob/3456e5e75cf8d4cd558955a5c7638fba2d68f464/YM/Next/FREE_KERNELS/Rack.lean\#L87-L88}{\texttt{Credit}}
& (III.3.2(1)); (C1)--(C2) \\
catalogue; price table &
\href{https://github.com/yang-mills-group/OpenAI-YM-20261028/blob/3456e5e75cf8d4cd558955a5c7638fba2d68f464/YM/Next/FREE_KERNELS/Priced.lean\#L29-L33}{\texttt{Catalog}};
\href{https://github.com/yang-mills-group/OpenAI-YM-20261028/blob/3456e5e75cf8d4cd558955a5c7638fba2d68f464/YM/PREC/Stacks.lean\#L123-L128}{\texttt{MapPrices}}
& constants fixed before \(L\) \\
hypercubic symmetry of \(d^*Hd\) &
\href{https://github.com/yang-mills-group/OpenAI-YM-20261028/blob/3456e5e75cf8d4cd558955a5c7638fba2d68f464/YM/Bid/Station.lean\#L54-L55}{\texttt{Bid.Sworn}}
& invariance under signed permutations \\
firm kernel; \(\Phi_f\); \(G_f\); \(G_\infty\) &
\href{https://github.com/yang-mills-group/OpenAI-YM-20261028/blob/3456e5e75cf8d4cd558955a5c7638fba2d68f464/YM/Further/FREE_COVARIANCE_KP/KP/Floor.lean\#L17-L21}{\texttt{Firm}};
\href{https://github.com/yang-mills-group/OpenAI-YM-20261028/blob/3456e5e75cf8d4cd558955a5c7638fba2d68f464/YM/Further/FREE_COVARIANCE_KP/KP/Core.lean\#L62}{\texttt{fuel}};
\href{https://github.com/yang-mills-group/OpenAI-YM-20261028/blob/3456e5e75cf8d4cd558955a5c7638fba2d68f464/YM/Further/FREE_COVARIANCE_KP/KP/Alloy.lean\#L15}{\texttt{garnet}};
\href{https://github.com/yang-mills-group/OpenAI-YM-20261028/blob/3456e5e75cf8d4cd558955a5c7638fba2d68f464/YM/Further/FREE_COVARIANCE_KP/KP/Alloy.lean\#L17}{\texttt{pearl}}
& curvature covariance, lattice and continuum \\
\end{longtable}
}

\subsubsection{7. What remains
unexplained}\label{what-remains-unexplained-5}

\begin{enumerate}
\def\labelenumi{\arabic{enumi}.}
\tightlist
\item
  \textbf{Gaussian interpretation and zero modes.} Lemma 4.4 verifies
  the covariance recursion away from zero momentum. A measure-level
  derivation requires a stated treatment of gauge directions and
  harmonic modes; neither the auxiliary singular symbols nor III.1's
  \href{https://github.com/yang-mills-group/OpenAI-YM-20261028/blob/3456e5e75cf8d4cd558955a5c7638fba2d68f464/YM/PREC/TreeMatch.lean\#L93-L103}{\texttt{action}}
  alone specifies that Gaussian measure. The identification with an
  iterated Feynman-gauge covariance remains probable.
\item
  \textbf{Change of Gaussian variables.} The quadratic splitting in §1
  follows from the identities. Its use as a change of integration
  variables, including null directions and determinants, belongs to
  III.4 and III.10 and is not established by this unit alone.
\item
  \textbf{Auxiliary spectral criterion.} The finite-test hypotheses of
  the cited cutoff-projector criterion must be supplied at its point of
  use. Its statement does not assert that a trajectory satisfies them.
\item
  \textbf{Further proof detail.} The analytic coefficient-extraction
  machinery underlying the strip bounds and the deeper positivity and
  square-root proofs F06--F07 are inputs here; their full proofs remain
  outside this account. The symbol-class calculations in
  \href{https://github.com/yang-mills-group/OpenAI-YM-20261028/blob/3456e5e75cf8d4cd558955a5c7638fba2d68f464/YM/Further/FREE_COVARIANCE_KP/KP/Gem.lean}{\texttt{Gem.lean}}
  and
  \href{https://github.com/yang-mills-group/OpenAI-YM-20261028/blob/3456e5e75cf8d4cd558955a5c7638fba2d68f464/YM/Further/FREE_COVARIANCE_KP/KP/Weight.lean}{\texttt{Weight.lean}}
  support the covariance comparison, but no position-space continuum
  limit or contact-term statement is asserted in III.3.6.
\end{enumerate}

\subsection{III.4. Gaussian integration: change of Gaussian sector,
determinants, window
law}\label{iii.4.-gaussian-integration-change-of-gaussian-sector-determinants-window-law}

\textbf{Part:} III, Ultraviolet construction (renormalization group).
\textbf{Lean:} about 40 work packages. The main ones are
\href{https://github.com/yang-mills-group/OpenAI-YM-20261028/tree/3456e5e75cf8d4cd558955a5c7638fba2d68f464/YM/SectorNext/GAUSSIAN_SECTOR_CHANGE}{\texttt{YM/SectorNext/GAUSSIAN\_SECTOR\_CHANGE/}}
(exact change of variables, gap, windows),
\href{https://github.com/yang-mills-group/OpenAI-YM-20261028/tree/3456e5e75cf8d4cd558955a5c7638fba2d68f464/YM/AnalyticBridge/RELATIVE_LOGDET_TRACE}{\texttt{YM/AnalyticBridge/RELATIVE\_LOGDET\_TRACE/}}
and
\href{https://github.com/yang-mills-group/OpenAI-YM-20261028/tree/3456e5e75cf8d4cd558955a5c7638fba2d68f464/YM/PhysicalNext/AFTERNOON_INTERACTING_RANK_COMPARISON}{\texttt{YM/PhysicalNext/AFTERNOON\_INTERACTING\_RANK\_COMPARISON/}}
(determinants and the polymer form of a sector),
\href{https://github.com/yang-mills-group/OpenAI-YM-20261028/tree/3456e5e75cf8d4cd558955a5c7638fba2d68f464/YM/AnalyticBridge/ANNULUS_RATIO_INSERTIONS}{\texttt{YM/AnalyticBridge/ANNULUS\_RATIO\_INSERTIONS/}},
\href{https://github.com/yang-mills-group/OpenAI-YM-20261028/tree/3456e5e75cf8d4cd558955a5c7638fba2d68f464/YM/SectorTools/PHYSICAL_GAUSSIAN_INSERTION}{\texttt{YM/SectorTools/PHYSICAL\_GAUSSIAN\_INSERTION/}}
and
\href{https://github.com/yang-mills-group/OpenAI-YM-20261028/tree/3456e5e75cf8d4cd558955a5c7638fba2d68f464/YM/Primitives/CONNECTED_GAUSSIAN_CUMULANTS}{\texttt{YM/Primitives/CONNECTED\_GAUSSIAN\_CUMULANTS/}}
(Gaussian moments),
\href{https://github.com/yang-mills-group/OpenAI-YM-20261028/tree/3456e5e75cf8d4cd558955a5c7638fba2d68f464/YM/AnalyticBridge/SECTOR_AFFINE_JETS}{\texttt{SECTOR\_AFFINE\_JETS/}},
\href{https://github.com/yang-mills-group/OpenAI-YM-20261028/tree/3456e5e75cf8d4cd558955a5c7638fba2d68f464/YM/AnalyticBridge/SECTOR_COVARIANCE_JETS}{\texttt{SECTOR\_COVARIANCE\_JETS/}}
and
\href{https://github.com/yang-mills-group/OpenAI-YM-20261028/tree/3456e5e75cf8d4cd558955a5c7638fba2d68f464/YM/AnalyticBridge/SELECTED_ORDERED_KERNELS}{\texttt{SELECTED\_ORDERED\_KERNELS/}}
(jets),
\href{https://github.com/yang-mills-group/OpenAI-YM-20261028/tree/3456e5e75cf8d4cd558955a5c7638fba2d68f464/YM/AnalyticBridge/LOCAL_GAUSSIAN_LAWS}{\texttt{LOCAL\_GAUSSIAN\_LAWS/}},
\href{https://github.com/yang-mills-group/OpenAI-YM-20261028/tree/3456e5e75cf8d4cd558955a5c7638fba2d68f464/YM/AnalyticBridge/GAUSSIAN_FIELD_ALIGNMENT}{\texttt{GAUSSIAN\_FIELD\_ALIGNMENT/}},
\href{https://github.com/yang-mills-group/OpenAI-YM-20261028/tree/3456e5e75cf8d4cd558955a5c7638fba2d68f464/YM/Push/FREE_RESOLVENT}{\texttt{YM/Push/FREE\_RESOLVENT/}}
and
\href{https://github.com/yang-mills-group/OpenAI-YM-20261028/tree/3456e5e75cf8d4cd558955a5c7638fba2d68f464/YM/RG08/Gaussian}{\texttt{YM/RG08/Gaussian/}}
(locality, infinite-volume kernels),
\href{https://github.com/yang-mills-group/OpenAI-YM-20261028/tree/3456e5e75cf8d4cd558955a5c7638fba2d68f464/YM/SectorNext/LARGE_FIELD_REMAINDERS}{\texttt{YM/SectorNext/LARGE\_FIELD\_REMAINDERS/}},
and the constant-producing statements in \texttt{Next/STATIONARY},
\texttt{MergeProof} and \texttt{PhysicalNext/V9\_JOINT\_ENDPOINT\_LOAD}.
\textbf{Size:} 37,956 lines in 4,737 declarations of the main-theorem
closure.

\textbf{Conventions.} The parameter \(x\) is the coarse inverse-coupling
parameter: in applications \(x=\hat\beta_m\). Throughout:

\begin{itemize}
\tightlist
\item
  \(R(x)=\lceil(\log x)^2\rceil\) is the cluster radius
  (\href{https://github.com/yang-mills-group/OpenAI-YM-20261028/blob/3456e5e75cf8d4cd558955a5c7638fba2d68f464/YM/Pressure/Seeds.lean\#L22}{\texttt{radiusB}}).
\item
  \(t(x)=x^{-1/2}e^{(\log\log x)^2}\) is the small-field threshold
  (\href{https://github.com/yang-mills-group/OpenAI-YM-20261028/blob/3456e5e75cf8d4cd558955a5c7638fba2d68f464/YM/Pressure/Seeds.lean\#L25}{\texttt{temp}}).
\item
  \(\beta_F(x,L)=x+s_G\log L\) is the coupling one level finer.
\item
  \(t_L(x)=t(\beta_F(x,L))\) is the fine-level threshold
  (\href{https://github.com/yang-mills-group/OpenAI-YM-20261028/blob/3456e5e75cf8d4cd558955a5c7638fba2d68f464/YM/Pressure/Seeds.lean\#L26}{\texttt{fiery}}).
\item
  \(g=\beta^{-1/2}\).
\item
  ``For large \(x\)'' means: there is \(x_0\), depending only on the
  constants fixed before it, such that the statement holds for
  \(x\ge x_0\).
\item
  Write \(\|A\|_s\) for the maximum of the row and column sums
  \(\sum_j e^{s\,\mathrm{dist}(i,j)}|A_{ij}|\). Positions lie on the
  coarse torus; several coordinates may have the same position. Matrix
  inequalities refer to quadratic forms. Put
  \(r_x=\lfloor R(x)/60\rfloor\).
\end{itemize}

\subsubsection{1. What this unit proves}\label{what-this-unit-proves-6}

\textbf{Setting.} One block-spin step \(S\) maps a fine torus
\(\Lambda_f\) (side \(f=Lc\), block size \(L\) odd) to a coarse torus
\(\Lambda_c\). The push-forward of the density at fixed coarse field
\(V\) is decomposed (III.1, III.9) into sectors. A sector is given by:

\begin{itemize}
\tightlist
\item
  a store \(A\): the large-field data, made of cost cells, aisles and
  flagged sites, with a partition of unity \(\mathrm{bill}_A\) given by
  an allocation \(D\);
\item
  a choice \(k\) of gauge-fixing tags;
\item
  a spanning forest \(F\) (tree gauge).
\end{itemize}

The integration variables are the free links of \(F\), with product Haar
measure. In each sector the unit splits the free links as follows
(\href{https://github.com/yang-mills-group/OpenAI-YM-20261028/blob/3456e5e75cf8d4cd558955a5c7638fba2d68f464/YM/SectorNext/GAUSSIAN_SECTOR_CHANGE/Shell.lean\#L38-L40}{\texttt{chill}}):

\begin{itemize}
\tightlist
\item
  \textbf{frozen links} \(\mathcal F_A\): a parent block of an endpoint
  lies within distance \(10R(x)\) of \(A\). These are kept as Haar
  variables \(Y\).
\item
  \textbf{active links} \(\mathcal A_A\): all others. These are written
  \(u_e=\exp(a_e)\,\nu_e\), with \(a_e\in\mathfrak g\) and
  \(\nu=\iota V\) the copy of \(V\) on the fine lattice
  (\href{https://github.com/yang-mills-group/OpenAI-YM-20261028/blob/3456e5e75cf8d4cd558955a5c7638fba2d68f464/YM/PCalc/Chart.lean\#L35-L36}{\texttt{replica}}).
\end{itemize}

We write \(n_a=|\mathcal A_A|\dim\mathfrak g\).

On selected stencils, the step energy contains a sum of affine squares
\(Q_Y(a)=\tfrac12\lvert Ma+s(Y)\rvert^2\)
(\href{https://github.com/yang-mills-group/OpenAI-YM-20261028/blob/3456e5e75cf8d4cd558955a5c7638fba2d68f464/YM/SectorNext/GAUSSIAN_SECTOR_CHANGE/Remainder.lean\#L37}{\texttt{guide}}):

\begin{itemize}
\tightlist
\item
  \(M\) is a linear map from active coordinates to rows of linearised
  plaquette curvatures and block constraints;
\item
  \(s(Y)\) depends on the frozen links.
\end{itemize}

From these the unit defines three objects:

\begin{itemize}
\tightlist
\item
  the \textbf{precision} \(W=M^{\top}M\)
  (\href{https://github.com/yang-mills-group/OpenAI-YM-20261028/blob/3456e5e75cf8d4cd558955a5c7638fba2d68f464/YM/SectorNext/GAUSSIAN_SECTOR_CHANGE/Remainder.lean\#L27}{\texttt{Rows.gram}});
\item
  the \textbf{covariance} \(C=W_{\mathcal W}\), the symmetrised windowed
  inverse of \(W\), with windows of coarse radius \(R(x)/20\)
  (\href{https://github.com/yang-mills-group/OpenAI-YM-20261028/blob/3456e5e75cf8d4cd558955a5c7638fba2d68f464/YM/SectorNext/GAUSSIAN_SECTOR_CHANGE/Chosen.lean\#L69-L70}{\texttt{fog}});
\item
  the \textbf{mean} \(m(Y)=-(W^{\mathrm{per}})^{\top}M^{\top}s(Y)\)
  (\href{https://github.com/yang-mills-group/OpenAI-YM-20261028/blob/3456e5e75cf8d4cd558955a5c7638fba2d68f464/YM/SectorNext/GAUSSIAN_SECTOR_CHANGE/Chosen.lean\#L72-L74}{\texttt{meanLog}}).
  This is a window approximation to the minimiser
  \(-W^{-1}M^{\top}s(Y)\); its gradient need not vanish.
\end{itemize}

The unit proves four things:

\begin{itemize}
\tightlist
\item
  an \emph{exact} Gaussian change of variables in every sector of every
  RG step;
\item
  the gap and locality of \(W\) and \(C\), uniformly in volume, store
  and field;
\item
  local expansions of the log-determinants, and a polymer form of each
  sector with large-field suppression;
\item
  Gaussian moment bounds, jet bounds, and the production of all
  post-\(L\) constants.
\end{itemize}

\textbf{Proposition III.4.1 (exact Gaussian change of variables in one
sector).} Fix the following data:

\begin{itemize}
\tightlist
\item
  a step \(S\) and a tube \(T\) satisfying the funding conditions of a
  teller: the prescribed kernel and stage, the multiscale metric bounds,
  the prescribed fine and coarse cutoffs, bounds on both path families,
  and compatible free metrics and towers;
\item
  a representation \(R\) of the fine density whose paths satisfy their
  bounds and whose activities satisfy the validity bounds, with
  \(\beta=\beta_R\);
\item
  \(x\), \(V\), \(k\), \(A\), \(F\), an allocation \(D\) and a cutoff
  profile \(\chi\);
\item
  a chart ratio \(J\), meaning:

  \begin{itemize}
  \tightlist
  \item
    \(J_{\rm sc}>0\) and \(\rho_J>0\), and \(J\) is measurable;
  \item
    \(J\) is real-analytic and positive on a ball
    \(B_{\rho_J}\subset\mathfrak g\);
  \item
    \(J(0)=1\);
  \item
    \(\exp\) is injective on \(B_{\rho_J}\);
  \item
    \(J_{\rm sc}\,J(a)\,da\) is mapped by \(\exp\) onto normalised Haar
    measure restricted to \(\exp B_{\rho_J}\);
  \item
    \(\exp B_{\rho_J}\) is open;
  \end{itemize}
\item
  a number \(d>0\).
\end{itemize}

Let \(C\) be a matrix on the active coordinates and \(m\) a map from
frozen arrays to active coordinates. Assume they are \emph{permitted}
(\href{https://github.com/yang-mills-group/OpenAI-YM-20261028/blob/3456e5e75cf8d4cd558955a5c7638fba2d68f464/YM/SectorNext/GAUSSIAN_SECTOR_CHANGE/Change.lean\#L95-L103}{\texttt{Permitted}}):

\begin{itemize}
\tightlist
\item
  \(C\succ0\), \(\beta>0\), \(m\) measurable, and \(4d<\rho_J\);
\item
  the \emph{standing condition}
  (\href{https://github.com/yang-mills-group/OpenAI-YM-20261028/blob/3456e5e75cf8d4cd558955a5c7638fba2d68f464/YM/SectorNext/GAUSSIAN_SECTOR_CHANGE/Shell.lean\#L51-L55}{\texttt{Standing}}):
  whenever \(\mathrm{bill}_A\ne0\) and the gauge tags are eligible,
  every active link equals \(e^{a_e}\nu_e\) with
  \(\lvert a_e\rvert\le d\).
\end{itemize}

Let
\(\Phi_{A,k}=\mathrm{bill}_A\cdot\mathrm{exposed}_R\cdot\mathrm{enters}_{T,\beta,k,V}\)
be the sector integrand. Here \(\mathrm{exposed}_R\) is the interacting
density of \(R\), and \(\mathrm{enters}\) is the tagged block-spin
kernel times \(e^{-\beta\cdot\text{gap energy}}\). Then the following
hold
(\href{https://github.com/yang-mills-group/OpenAI-YM-20261028/blob/3456e5e75cf8d4cd558955a5c7638fba2d68f464/YM/SectorNext/GAUSSIAN_SECTOR_CHANGE/Change.lean\#L127-L135}{\texttt{Licenses}}):

\begin{itemize}
\tightlist
\item
  \begin{enumerate}
  \def\labelenumi{(\alph{enumi})}
  \tightlist
  \item
    \(\Phi_{A,k}\) is integrable on the slice;
  \end{enumerate}
\item
  (b) for every \(Y\), \(\Psi(Y,\cdot)\in L^1(\gamma_C)\), and
  \(\Psi\in L^1(\mathrm{Haar}^{\mathcal F_A}\otimes\gamma_C)\);
\item
  (c) the identity
  \[\int\Phi_{A,k}\,d\mathrm{Haar}^{F}=\mathrm{pref}_S(\beta)\,J_{\rm sc}^{\,\lvert\mathcal A_A\rvert}\,(2\pi g^2)^{n_a/2}(\det C)^{1/2}\int\mathrm{Haar}^{\mathcal F_A}(dY)\int\gamma_C(dz)\,\Psi(Y,z).\]
\end{itemize}

Here \(\gamma_C=N(0,C)\). Set \(a=m(Y)+gz\), and let \(U\) be the
configuration with frozen links \(Y\), active links \(e^{a_e}\nu_e\) and
tree links \(1\). Then
(\href{https://github.com/yang-mills-group/OpenAI-YM-20261028/blob/3456e5e75cf8d4cd558955a5c7638fba2d68f464/YM/SectorNext/GAUSSIAN_SECTOR_CHANGE/Change.lean\#L63-L71}{\texttt{dyed}})
\[\Psi(Y,z)=\tilde J_d(a)\,\chi_d(a)\,\mathrm{bill}_A\,\mathrm{lik}_k(U)\,\mathrm{exposed}_R(U)\,\exp\!\Big(-\beta\,\overline{\mathrm{res}}(Y,U,a)-\beta Q_Y(m(Y))+\eta\cdot z+\tfrac12 z\cdot(C^{-1}-W)z\Big).\]
In this formula:

\begin{itemize}
\tightlist
\item
  \(\eta=-g^{-1}(Wm(Y)+M^{\top}s(Y))\);
\item
  \(\overline{\mathrm{res}}\) is the orientation-averaged step energy
  minus the affine squares;
\item
  \(\tilde J_d(a)\) is the product of the extended link Jacobians, each
  equal to \(J(a_e)\) for \(\lvert a_e\rvert\le 2d\); \(\chi_d\) is the
  product cutoff, equal to \(1\) when every \(\lvert a_e\rvert\le d\).
\end{itemize}

Equivalently,
\(\mathrm{pref}\cdot\Psi(Y,z)=e^{\frac12 z\cdot C^{-1}z}\,r(Y,m(Y)+gz)\),
where \(r\) is the integrand in Lebesgue log-coordinates
(\href{https://github.com/yang-mills-group/OpenAI-YM-20261028/blob/3456e5e75cf8d4cd558955a5c7638fba2d68f464/YM/SectorNext/GAUSSIAN_SECTOR_CHANGE/Change.lean\#L73-L93}{\texttt{recolored}}).
The identity holds for \emph{every} permitted \((C,m)\). No quantitative
locality assumption is needed.

\emph{Status:} theorem.

\emph{Lean:}
\href{https://github.com/yang-mills-group/OpenAI-YM-20261028/blob/3456e5e75cf8d4cd558955a5c7638fba2d68f464/YM/SectorNext/GAUSSIAN_SECTOR_CHANGE/Change.lean\#L137-L192}{\texttt{authorized\_change}}.

\textbf{Proposition III.4.2 (licence of every committed RG step; the
transfer identity).} Let \(\mathscr R\) be a directed RG run (III.9)
whose delays satisfy the rate conditions and the per-stage chart
conditions
(\href{https://github.com/yang-mills-group/OpenAI-YM-20261028/blob/3456e5e75cf8d4cd558955a5c7638fba2d68f464/YM/SectorNext/GAUSSIAN_SECTOR_CHANGE/Handoff.lean\#L57-L71}{\texttt{Starcash}}).
Let \(m<N\), let \(p\) be a period, and let \(R\) be a representation at
level \(m+1\) on its selected history with
\(\beta_R\in[\hat\beta_{m+1}/2,\,2\hat\beta_{m+1}]\). Write:

\begin{itemize}
\tightlist
\item
  \(S\) for the step from level \(m+1\) to level \(m\);
\item
  \(x=\hat\beta_m\) and \(\beta=\beta_R\);
\item
  \(j=J(m+1)\) and \(i=J(m)\) for the stage indices;
\item
  \(d=a_{ji}\,t(\hat\beta_{m+1})\), and \(w\) for the allocation width
  at this stage pair.
\end{itemize}

Then:

\begin{itemize}
\tightlist
\item
  \(w>0\), \(d>0\), \(\beta>0\), \(E_{ji}>0\), \(2d<\rho_{\rm first}\)
  and \(4d<\rho_J\);
\item
  \(\lvert\log J(t)\rvert\le E_{ji}\lvert t\rvert^2\) for
  \(\lvert t\rvert\le4d\);
\item
  there is \(\kappa>0\) such that for all allocations \(D\) of width
  \(w\), forests \(H\) and profiles \(\chi\):

  \begin{itemize}
  \tightlist
  \item
    \begin{enumerate}
    \def\labelenumi{(\roman{enumi})}
    \tightlist
    \item
      writing \(I_A\) for the signed sector column, its contribution to
      the push-forward, and \(L_A\) for the nonnegative column with
      absolute fine density, after removing the selected packet energy
      and kernel prefactor, both have bounded Borel envelopes and
      \(|I_A(V)|\le\mathrm{pref}(\beta)e^{-\beta\kappa t_L(x)^2|A|}L_A(V)\)
      (\href{https://github.com/yang-mills-group/OpenAI-YM-20261028/blob/3456e5e75cf8d4cd558955a5c7638fba2d68f464/YM/Push/SECTOR_MEASURE/Presented.lean\#L151-L155}{\texttt{PortionCost}});
    \end{enumerate}
  \item
    (ii) for all \(A,k,V\), the chart is admitted
    (\href{https://github.com/yang-mills-group/OpenAI-YM-20261028/blob/3456e5e75cf8d4cd558955a5c7638fba2d68f464/YM/SectorNext/GAUSSIAN_SECTOR_CHANGE/Stamped.lean\#L36-L42}{\texttt{Admitted}});
    when \(\mathrm{flags}(A)=\mathrm{flags}_{t(x)}(V)\), also
    \(C=W_{\mathcal W}\succ0\), \(W\succ0\) and the selected predictor
    rows have analytic derivatives with a common bound; and the
    conclusions of III.4.1 hold with \(C=W_{\mathcal W}\) and the
    windowed mean for every store (for wrong flags trivially, since
    \(\mathrm{bill}_A\equiv0\));
  \item
    (iii) for every \(V\):
    \[\mathcal S_\beta(\rho_R)(V)=\sum_A\mathrm{sieve}(V)\sum_k\mathrm{pref}\,J_{\rm sc}^{\lvert\mathcal A_A\rvert}(2\pi g^2)^{n_a/2}(\det C_{A,k,V})^{1/2}\iint\Psi\,d\gamma_{C}\,d\mathrm{Haar}.\]
  \end{itemize}
\end{itemize}

Here \(\mathcal S_\beta\) is the block-spin push-forward, \(\rho_R\) is
the fine density, and
\(\mathrm{sieve}(V)=\exp(-\beta\cdot\text{orientation-averaged coarse energy of }V)\).

\emph{Status:} theorem, conditional on the run, history and band
hypotheses. III.4.12 supplies the rate and chart conditions; III.9
supplies the selected history.

\emph{Lean:}
\href{https://github.com/yang-mills-group/OpenAI-YM-20261028/blob/3456e5e75cf8d4cd558955a5c7638fba2d68f464/YM/SectorNext/GAUSSIAN_SECTOR_CHANGE/Handoff.lean\#L132-L182}{\texttt{committed}},
\href{https://github.com/yang-mills-group/OpenAI-YM-20261028/blob/3456e5e75cf8d4cd558955a5c7638fba2d68f464/YM/SectorNext/GAUSSIAN_SECTOR_CHANGE/Stamped.lean\#L124-L137}{\texttt{displaced\_wrong}}.

\textbf{Proposition III.4.3 (gap and locality of the actual precision
and covariance).} Fix \(G\), a patch \(P=(L,J)\) and a teller whose
kernel satisfies the branch condition. There are \(v,\eta,b,E>0\) and
\(K\ge0\) (post-\(L\) constants) with the following property. For large
\(x\), take any step with patch \(P\) and coarse side
\(c\ge100(1+R(x))\), any funded tube, and any \(A,F,k,V\) with
\(\mathrm{flags}(A)=\mathrm{flags}_{t(x)}(V)\). Then:

\begin{itemize}
\tightlist
\item
  \begin{enumerate}
  \def\labelenumi{(\roman{enumi})}
  \tightlist
  \item
    \(\eta\lvert a\rvert^2\le a\cdot Wa\) for all active \(a\)
    (\href{https://github.com/yang-mills-group/OpenAI-YM-20261028/blob/3456e5e75cf8d4cd558955a5c7638fba2d68f464/YM/SectorNext/GAUSSIAN_SECTOR_CHANGE/Gap.lean\#L242-L264}{\texttt{Gap.not\_flat}});
  \end{enumerate}
\item
  (ii) \(W\) has exponential off-diagonal decay in the Schur norm:
  \(\sup_i\sum_j e^{v\,\mathrm{dist}(i,j)}\lvert W_{ij}\rvert\le K\),
  and the same for columns, where \(\mathrm{dist}\) is the coarse
  distance between the parent blocks of the coordinates;
\item
  (iii) \(\lVert C-W^{-1}\rVert_0\le E\,e^{-br_x}\), and \(C\succ0\);
\item
  (iv) \(C\le K'\) for some post-\(L\) constant \(K'\)
  (\href{https://github.com/yang-mills-group/OpenAI-YM-20261028/blob/3456e5e75cf8d4cd558955a5c7638fba2d68f464/YM/AnalyticBridge/SHARP_NORMALIZED_TAIL/Apertures.lean\#L43-L55}{\texttt{open\_price}}).
\end{itemize}

The constants do not depend on the volume, the store, the forest or the
field.

\emph{Status:} theorem.

\emph{Lean:}
\href{https://github.com/yang-mills-group/OpenAI-YM-20261028/blob/3456e5e75cf8d4cd558955a5c7638fba2d68f464/YM/SectorNext/GAUSSIAN_SECTOR_CHANGE/Chosen.lean\#L102-L143}{\texttt{opaqueCost}}.

\textbf{Proposition III.4.4 (window law).} Let \(\sigma,\eta>0\),
\(K\ge0\), and fix a nonnegative function \(B:\mathbb R\to\mathbb R\).
There are \(u,E,g>0\) such that for every sufficiently large integer
\(r\) the following holds, uniformly in the torus and the finite
coordinate set. Let the coordinate positions \(p_j\) satisfy
\(\sum_j e^{-v\,\mathrm{dist}(i,p_j)}\le B(v)\) for every coarse site
\(i\) and every \(v>0\)
(\href{https://github.com/yang-mills-group/OpenAI-YM-20261028/blob/3456e5e75cf8d4cd558955a5c7638fba2d68f464/YM/SectorNext/GAUSSIAN_SECTOR_CHANGE/Windows.lean\#L125}{\texttt{Revenue}}).
Let \(A\) be a finite symmetric matrix with:

\begin{itemize}
\tightlist
\item
  Schur decay
  \(\lVert e^{\sigma\,\mathrm{dist}}A\rVert_{\rm Schur}\le K\);
\item
  gap \(\eta\);
\item
  windows \(\mathcal W_i\) that contain the \(r\)-ball around \(i\).
\end{itemize}

Let \(L\) be the windowed inverse,
\(L_{\cdot i}=(A|_{\mathcal W_i})^{-1}e_i\) extended by zero, and
\(A_{\mathcal W}=\tfrac12(L+L^{\top})\)
(\href{https://github.com/yang-mills-group/OpenAI-YM-20261028/blob/3456e5e75cf8d4cd558955a5c7638fba2d68f464/YM/SectorNext/GAUSSIAN_SECTOR_CHANGE/Windows.lean\#L70}{\texttt{clouded}}).
Then:

\begin{itemize}
\tightlist
\item
  \(\lVert A_{\mathcal W}-A^{-1}\rVert_{\rm Schur}\le Ee^{-ur}\);
\item
  \(A_{\mathcal W}\ge g=\eta/(2(K+1)^2)\) and \(A_{\mathcal W}\succ0\).
\end{itemize}

A strengthened form
(\href{https://github.com/yang-mills-group/OpenAI-YM-20261028/blob/3456e5e75cf8d4cd558955a5c7638fba2d68f464/YM/AnalyticBridge/GAUSSIAN_FIELD_ALIGNMENT/Relative.lean\#L72-L118}{\texttt{same\_purchase}})
gives a rate \(s\in(0,\sigma]\) and a constant \(K_1\) such that \(A\),
\(A^{-1}\), \(L\) and \((A_{\mathcal W})^{-1}\) all have \(s\)-weighted
Schur norm \(\le K_1\), and
\(\lVert e^{s\,\mathrm{dist}}(L-A^{-1})\rVert_{\rm Schur}\le K_1e^{-sr}\).

\emph{Status:} theorem. The constants depend only on
\((\sigma,\eta,K,B)\).

\emph{Lean:}
\href{https://github.com/yang-mills-group/OpenAI-YM-20261028/blob/3456e5e75cf8d4cd558955a5c7638fba2d68f464/YM/SectorNext/GAUSSIAN_SECTOR_CHANGE/Windows.lean\#L184-L239}{\texttt{Schur.windows}}.

\textbf{Proposition III.4.5 (log-determinants as local path series).}
Let \(A\succ0\) with windows \(\mathcal W\) (\(i\in\mathcal W_i\)). Let
\(P_u\) be the column-wise windowed inverse of \(A+uI\), and
\(\Delta_u=(A+uI)P_u-I\). Assume the following bounds
(\href{https://github.com/yang-mills-group/OpenAI-YM-20261028/blob/3456e5e75cf8d4cd558955a5c7638fba2d68f464/YM/AnalyticBridge/RELATIVE_LOGDET_TRACE/Walls.lean\#L79-L89}{\texttt{Priced}}),
at a fixed nonnegative Schur rate, with \(C\ge1\) and \(0\le q<1\), for
every \(u\ge0\):

\begin{itemize}
\tightlist
\item
  \(\lVert A\rVert\le C\);
\item
  \(\lVert(A+u)^{-1}\rVert,\ \lVert P_u\rVert\le C/(1+u)\);
\item
  \(\lVert\Delta_u\rVert\le q/(1+u)\);
\item
  \(\lVert S\rVert\le q\), where \(S=AA_{\mathcal W}-1\).
\end{itemize}

Assume also \(A_{\mathcal W}\succ0\). Then
\[\log\det A=\sum_i\Big(\ell_i+\sum_{n\ge0}r_{i,n}\Big),\qquad \log\det(1+S)=\sum_i\sum_{n\ge0}\frac{(-1)^n(S^{n+1})_{ii}}{n+1},\qquad \log\det A_{\mathcal W}=-\log\det A+\log\det(1+S),\]
with:

\begin{itemize}
\tightlist
\item
  \(\ell_i=\int_0^\infty[(1+u)^{-1}-(P_u)_{ii}]\,du=(\log A|_{\mathcal W_i})_{ii}\),
  the single-window letter;
\item
  \(r_{i,n}=-\int_0^\infty(P_u(-\Delta_u)^{n+1})_{ii}\,du\) and
  \(\int\lvert\cdot\rvert\le Cq^{n+1}\).
\end{itemize}

All series converge absolutely. Moreover
(\href{https://github.com/yang-mills-group/OpenAI-YM-20261028/blob/3456e5e75cf8d4cd558955a5c7638fba2d68f464/YM/AnalyticBridge/RELATIVE_LOGDET_TRACE/Alphabet.lean\#L196-L208}{\texttt{labeled}}),
the matrix products expand into closed index paths. A term \(r_{i,n}\)
has one \(P_u\) factor and \(n+1\) defect factors, hence \(n+2\) jumps;
the corresponding trace-series term has \(n+1\) jumps. The absolute
values, summed over the intermediate indices, are bounded by
\(Cq^{n+1}\) and \(q^{n+1}/(n+1)\), respectively.

\emph{Status:} theorem.

\emph{Lean:}
\href{https://github.com/yang-mills-group/OpenAI-YM-20261028/blob/3456e5e75cf8d4cd558955a5c7638fba2d68f464/YM/AnalyticBridge/RELATIVE_LOGDET_TRACE/Series.lean\#L203-L227}{\texttt{Priced.chains}}.

For the relative expansions, the sector regularity conditions are: a
lawful representation; matching fine paths and cutoff; a local chart
with \(d>0\), \(4d<\rho_J\) and \(2d<\rho_{\rm first}\); positive
coupling and sector width; \(0<t_L(x)\le1\); output threshold \(t(x)\);
and the tube's geometric, local-readability, predictor-derivative,
covariance-positivity and gap conditions
(\href{https://github.com/yang-mills-group/OpenAI-YM-20261028/blob/3456e5e75cf8d4cd558955a5c7638fba2d68f464/YM/PhysicalNext/AFTERNOON_INTERACTING_RANK_COMPARISON/Opacity/Desk.lean\#L41-L54}{\texttt{Chips}},
\href{https://github.com/yang-mills-group/OpenAI-YM-20261028/blob/3456e5e75cf8d4cd558955a5c7638fba2d68f464/YM/PhysicalNext/AFTERNOON_INTERACTING_RANK_COMPARISON/Exposure.lean\#L40-L46}{\texttt{Boosted}}).

\textbf{Proposition III.4.6 (relative normalisation of a sector and its
localisation).} Let \(K_F\) be the flat kinetic precision: the Gram
matrix of the move at \(V=1\) with empty store, on all free links of
\(F\).

\begin{enumerate}
\def\labelenumi{(\roman{enumi})}
\item
  If the fine path family satisfies its path bounds and the kernel
  satisfies the patch's branch condition, \(\det K_F\) does not depend
  on the forest or on the gauge tags
  (\href{https://github.com/yang-mills-group/OpenAI-YM-20261028/blob/3456e5e75cf8d4cd558955a5c7638fba2d68f464/YM/SectorNext/GAUSSIAN_SECTOR_CHANGE/Bore.lean\#L303-L305}{\texttt{Native.shared}}).
\item
  Assume:
\end{enumerate}

\begin{itemize}
\tightlist
\item
  the tube is funded, and its fine path family and cutoff are those of
  \(R\);
\item
  the precision certificate of Lemma 4.21 holds;
\item
  \(J\) is a local chart ratio, \(\beta>0\), and the threshold of the
  output budget point is \(t(x)\).
\end{itemize}

Then on every sector point \(p=(A,F,k,V,Y,z)\) with matching flags, the
positive density factor, with the sector masks and polymer product kept
outside it, factorises as
\[\text{freight}(p)=\text{loading}(V)\,e^{\Lambda(p)},\qquad \text{loading}(V)=\text{banker}(V)\,\mathrm{pref}(\beta)\,e^{v_R\lvert\Lambda_f\rvert}\,\mathrm{spare}(\beta)^{\lvert F^{\rm free}\rvert}(\det K_F)^{-1/2}>0.\]
Here
\(\text{freight}=J_{\rm sc}^{\lvert\mathcal A\rvert}(2\pi g^2)^{n_a/2}(\det C)^{1/2}e^{\frac12 z\cdot C^{-1}z}\,\tilde J\cdot(\text{noise})\cdot e^{\mathrm{exponent}_R}\),
and \(\mathrm{spare}(\beta)=(2\pi g^2)^{\dim\mathfrak g/2}J_{\rm sc}\)
is the Gaussian normalisation per link. The exponent is
\[\Lambda(p)=-\beta\,\mathcal E(p)+\eta\cdot z+\tfrac12z\cdot(C^{-1}-W)z+\kappa(p)+\sum_e\log\tilde J(a_e)+(\text{tails}),\]
\[\kappa(p)=-\lvert\mathcal F_A\rvert\log\mathrm{spare}(\beta)-\tfrac12\big(\log\det W-\log\det K_F\big)+\tfrac12\log\det(WC).\]
Here \(\mathcal E\) is the orientation-averaged row energy, including
the subtraction of the affine energy at \(a\) and its addition at \(m\):
\[\mathcal E=E_{\rm noise}(U,V)+E_{\rm fine}^{\rm cut}(U)-E_{\rm coarse}^{\rm cut}(V)-Q_Y(a)+Q_Y(m).\]
The tails are the remaining fine exponent after removing its volume
constant and restoring its cut energy, together with the coarse cutoff
correction. Thus \(\mathcal E\) is not merely a sum of affine squares.

\begin{enumerate}
\def\labelenumi{(\roman{enumi})}
\setcounter{enumi}{2}
\tightlist
\item
  Assume in addition the sector regularity conditions below, a
  nonnegative polymer exponent, and both path-counting bounds of Lemma
  4.21. Then
  \(\kappa(p)=-\lvert\mathcal F_A\rvert\log\mathrm{spare}+\sum_e\mathrm{sect}_e+I+\mathrm{Sh}\)
  (\href{https://github.com/yang-mills-group/OpenAI-YM-20261028/blob/3456e5e75cf8d4cd558955a5c7638fba2d68f464/YM/PhysicalNext/AFTERNOON_INTERACTING_RANK_COMPARISON/Opacity/Shards.lean\#L34-L39}{\texttt{book\_price}}),
  where:
\end{enumerate}

\begin{itemize}
\tightlist
\item
  \(\mathrm{sect}_e=-\tfrac12\sum_{\rm colours}[\ell_e(\text{actual})-\ell_e(\text{flat})]\),
  the difference of the single-window letters of III.4.5 for the
  identity-completed full precision at \((A,V)\) and at
  \((\emptyset,1)\);
\item
  \(I\) is the corresponding difference of path letters;
\item
  \(\mathrm{Sh}=\tfrac12\log\det(1+S)\).
\end{itemize}

\begin{enumerate}
\def\labelenumi{(\roman{enumi})}
\setcounter{enumi}{3}
\tightlist
\item
  Regrouping
  (\href{https://github.com/yang-mills-group/OpenAI-YM-20261028/blob/3456e5e75cf8d4cd558955a5c7638fba2d68f464/YM/PhysicalNext/AFTERNOON_INTERACTING_RANK_COMPARISON/Opacity/Enigma.lean\#L36-L48}{\texttt{rate\_pass}})
  gives \(\kappa=\sum_{o}\mathrm{rent}_o+\mathrm{far}+I+\mathrm{Sh}\):
\end{enumerate}

\begin{itemize}
\tightlist
\item
  \(o\) runs over the cores of \(A\) (the components formed by overlap
  of the \(500R(x)\)-halos of its seeds);
\item
  \(\mathrm{rent}_o\) is the sum of
  \(\mathrm{sect}_e-\log\mathrm{spare}\cdot1_{e\in\mathcal F_A}\) over
  links within \(200R(x)\) of \(o\);
\item
  \(\mathrm{far}\) is the same sum over links farther than \(200R(x)\)
  from \(A\).
\end{itemize}

\emph{Status:} theorem.

\emph{Lean:}
\href{https://github.com/yang-mills-group/OpenAI-YM-20261028/blob/3456e5e75cf8d4cd558955a5c7638fba2d68f464/YM/PhysicalNext/AFTERNOON_INTERACTING_RANK_COMPARISON/Opacity/Prism.lean\#L102-L168}{\texttt{Dyed.cardinal\_spin}},
\href{https://github.com/yang-mills-group/OpenAI-YM-20261028/blob/3456e5e75cf8d4cd558955a5c7638fba2d68f464/YM/PhysicalNext/AFTERNOON_INTERACTING_RANK_COMPARISON/Opacity/Prism.lean\#L50-L54}{\texttt{Dyed.keptEarly}},
\href{https://github.com/yang-mills-group/OpenAI-YM-20261028/blob/3456e5e75cf8d4cd558955a5c7638fba2d68f464/YM/PhysicalNext/AFTERNOON_INTERACTING_RANK_COMPARISON/Opacity/Prism.lean\#L56-L63}{\texttt{Dyed.spoolLift}}.

\textbf{Proposition III.4.7 (polymer form of a sector, with large-field
suppression).} Fix:

\begin{itemize}
\tightlist
\item
  a teller and a patch \(D\) satisfying the branch condition;
\item
  a determinant-jet fund and a local chart ratio \(J\);
\item
  a polymer exponent \(e\ge0\), a chart-width multiplier \(h\ge0\),
  constants \(a,\alpha,\tau>0\) with \(2\alpha\le a/2\), and certificate
  constants \(C_0\ge0\), \(g_0,b_0,d_0>0\), \(\eta_0\in\mathbb R\).
\end{itemize}

Then for large \(x\) the following holds. Let \(y=\mathrm{retime}(x)\),
\(z=\beta_F(x,D.L)\) and \(\mathrm{hot}(s)=e^{(\log\log s)^{3/2}}\).
Eventually \(x\le y\le x^2\) and \(t(y)=x^{-1/2}\mathrm{hot}(x)\). Take
a step with patch \(D\), a sector datum with chart \(J\), and positive
cutoff parameter \(W_0\), subject to the following assumptions:

\begin{itemize}
\tightlist
\item
  the representation is lawful; the tube has the same fine paths and
  cutoff as the representation and satisfies the funded, geometric and
  chart-admission conditions; the chart radius \(d=h\,t_{D.L}(y)\)
  satisfies \(d>0\), \(4d<\rho_J\) and \(2d<\rho_{\rm first}\); the
  output threshold is \(t(y)\);
\item
  the step satisfies the free-metric size condition and
  \(c\ge100(1+R(y))\), and its fine and coarse path families have a
  common path bound;
\item
  \(\beta\in[z/2,2z]\); the fine threshold \(s=z^{-1/2}\mathrm{hot}(z)\)
  also lies in \([t_{D.L}(y)/2,2t_{D.L}(y)]\);
\item
  on each eligible packet the row energy \(\mathcal E\) is at least
  \(a s^2\) times its size, whenever the sector cutoffs are nonzero;
\item
  the support exponents satisfy the reserve inequality with parameters
  \((e,\tau)\); the budget radius \(B_0\) satisfies
  \(25R(y)+30\le B_0\le\mathrm{walls}(y)\), and the fine budget radius
  is at most \(2B_0\);
\item
  the lawful-representation bound can be chosen with activity constant
  \(\Gamma\le\beta_F(y,D.L)^{1/8}\) and smallness parameter
  \(\delta\le y^{-1/2}\), and the fine paths satisfy the required
  row-counting bound;
\item
  the determinant-jet certificate at halo \(20\), exponential order
  \(2\) and derivative order \(1\), and the precision and two
  path-counting certificates of Lemma 4.21 hold at \(y\), with constants
  \((g_0,C_0,\eta_0,b_0,d_0)\);
\item
  uniformly over flag-matched sectors and frozen arrays, \(\|C\|_0\) is
  bounded by a common constant, every component of \(m\) and of
  \(Mm+s(Y)\) has absolute value at most \(1\), and
  \(|\eta_i|\le y^{-4}\); on the chart-cutoff support, \(|z_i|\le y^2\).
\end{itemize}

For such data there is a decomposition
(\href{https://github.com/yang-mills-group/OpenAI-YM-20261028/blob/3456e5e75cf8d4cd558955a5c7638fba2d68f464/YM/PhysicalNext/AFTERNOON_INTERACTING_RANK_COMPARISON/Opacity/Plotter.lean\#L19-L28}{\texttt{Plate}})
\[\Lambda(p)=\sum_{o\in\mathrm{cores}(A)}\mathrm{base}_o(p)+\sum_X\mathrm{aux}_X(p)\]
whenever the flags match and the product of the fine-field, block,
gauge-tag and chart cutoffs is nonzero. The bounds below hold for every
flag-matched \(p\), including points outside that support. In it:

\begin{itemize}
\tightlist
\item
  \(\mathrm{base}_o\) depends only on the data in the
  \(500R(y)\)-neighbourhood of \(o\);
\item
  \(\mathrm{aux}_X\) depends only on the data near the local support
  \(X\);
\item
  both are gauge-invariant and lattice-covariant;
\item
  \(\mathrm{base}_o\le-\alpha\,\mathrm{hot}(z)^2\lvert o\rvert\), where
  \(\lvert o\rvert\) counts cost cells, aisles and flags;
\item
  \(\lvert e^{\mathrm{aux}_X}-1\rvert\le f_X\), with \(f\ge0\)
  satisfying the polymer bound
  \(\sum_{X\ni y'}f_Xe^{e\,N_{B_0}(X)}\le y^{-1/48}\) at every site
  \(y'\), where \(N_{B_0}\) is the cover number at the output budget
  radius
  (\href{https://github.com/yang-mills-group/OpenAI-YM-20261028/blob/3456e5e75cf8d4cd558955a5c7638fba2d68f464/YM/Ac/Tree.lean\#L12-L13}{\texttt{Hit}}).
\end{itemize}

\emph{Status:} theorem.

\emph{Lean:}
\href{https://github.com/yang-mills-group/OpenAI-YM-20261028/blob/3456e5e75cf8d4cd558955a5c7638fba2d68f464/YM/PhysicalNext/AFTERNOON_INTERACTING_RANK_COMPARISON/Opacity/Gush.lean\#L96-L134}{\texttt{foreshadows}},
with hypotheses
\href{https://github.com/yang-mills-group/OpenAI-YM-20261028/blob/3456e5e75cf8d4cd558955a5c7638fba2d68f464/YM/PhysicalNext/AFTERNOON_INTERACTING_RANK_COMPARISON/Opacity/Gush.lean\#L82-L95}{\texttt{Helps}}.

\textbf{Proposition III.4.8 (moments under the sector law).} Fix \(G\).
For every covariance cap \(L>0\) and \(k\in\mathbb N\) there is
\(D\ge1\) with the following property. Take:

\begin{itemize}
\tightlist
\item
  any finite index set \(I\) and any probability space \((\Omega,\mu)\);
\item
  any \(C\succeq0\) on \(\mathbb R^{I\times\dim\mathfrak g}\) with
  \(C\le L\);
\item
  any finite family \((r_i)_{i\in Y}\) of measurable insertions with
  \(w_i,d_i\ge0\),
  \(\lvert r_i(\omega,z)\rvert\le w_ie^{d_i(1+\lvert z_{a_i}\rvert^2+\lvert z_{b_i}\rvert^2)}\),
  and \(\sum_{i:\,x\in\{a_i,b_i\}}d_i\le c=1/(24L)\) at every coordinate
  site \(x\);
\item
  any \(\rho,H\ge0\) and sets \(P,J\subset I\);
\item
  any measurable \(v\) with
  \(\lvert v\rvert\le H(1+\sum_{i\in P}\lvert z_i\rvert^2)^k\) that
  vanishes unless \(\lvert z_i\rvert\ge\rho\) for all \(i\in J\).
\end{itemize}

Then \(f=v\prod_{i\in Y}r_i\in L^1(\mu\otimes\gamma_C)\) and
\[\int\lvert f\rvert\,d(\mu\otimes\gamma_C)\le H\,D\,(1+\lvert P\rvert\dim\mathfrak g)^{2k+1}\,e^{-c\lvert J\rvert\rho^2}\,2^{\dim\mathfrak g\,\lvert J\rvert}\prod_{i\in Y}\big(e^{c}\,4^{\dim\mathfrak g}w_i\big).\]
In a sector, take \(\mu=\mathrm{Haar}^{\mathcal F_A}\) and
\(C=W_{\mathcal W}\); III.4.3 supplies a uniform covariance cap.
Moreover:

\begin{itemize}
\tightlist
\item
  (\href{https://github.com/yang-mills-group/OpenAI-YM-20261028/blob/3456e5e75cf8d4cd558955a5c7638fba2d68f464/YM/AnalyticBridge/ANNULUS_RATIO_INSERTIONS/Ember.lean\#L69-L111}{\texttt{r4}})
  For measurable \(q_i\) on the same product space, let
  \(\lvert q_i\rvert\le x_i^2(1+\lvert z_{a_i}\rvert^2+\lvert z_{b_i}\rvert^2)\)
  with \(0\le x_i\le1\) and \(\sum_{i\ni\text{site}}8x_i\le c\). Then
  \(\int(\prod_i(e^{q_i}-1))^4\le(\prod_iMx_i)^4\), with \(M=M(G,L)\).
\item
  (\href{https://github.com/yang-mills-group/OpenAI-YM-20261028/blob/3456e5e75cf8d4cd558955a5c7638fba2d68f464/YM/AnalyticBridge/ANNULUS_RATIO_INSERTIONS/Echo.lean\#L191-L223}{\texttt{oasis}})
  Fix a patch, a teller satisfying its branch condition and a covariance
  cap \(L>0\). For large \(x\), every funded step with that patch and
  \(c\ge100(1+R(x))\), matching flags, \(0<d\le1\),
  \(2d<\rho_{\rm first}\) and
  \(\beta\in[\beta_F(x,L_{\rm block})/2,2\beta_F(x,L_{\rm block})]\) has
  the fourth-moment insertion bounds above. The single-site exponents
  come from \(\eta\cdot z\) and the pair exponents from
  \(\tfrac12z\cdot(C^{-1}-W)z\). Their weights are \(De^{-uR(x)/480}\)
  for sites and exponentially decaying in distance for pairs; the
  constants are fixed before the threshold in \(x\), which may depend on
  the covariance cap.
\end{itemize}

\emph{Status:} theorem.

\emph{Lean:}
\href{https://github.com/yang-mills-group/OpenAI-YM-20261028/blob/3456e5e75cf8d4cd558955a5c7638fba2d68f464/YM/AnalyticBridge/ANNULUS_RATIO_INSERTIONS/Arc.lean\#L189-L253}{\texttt{same\_law}}.

\textbf{Proposition III.4.9 (Gaussian insertions supported on selected
coordinates).} Fix \(p\ge1\), a finite jet order, \(G\), a patch and a
teller satisfying the branch condition. There are post-\(L\) constants
\(K,M,s,b>0\) and \(B\ge0\) such that, for large \(x\), the following
holds. Take a funded step with that patch, \(c\ge100(1+R(x))\) and
matching flags. Fix any \(g,d,\theta>0\), frozen array and slot set
\(U\). Restrict the linear and quadratic correction to \(U\), obtaining
\(t\cdot z_U+\tfrac12z_U\cdot\mathcal Jz_U\), where \(t=\eta\) and
\(\mathcal J=C^{-1}-W\). Its tilted law is
\(N(\mathcal A t_U,\mathcal A)\), with
\(\mathcal A=(C^{-1}-\mathcal J_U)^{-1}\succ0\). The likelihood ratio of
the \(U\)-marginals satisfies
\[\mathbb E_{\gamma_{C,U}}\big[(dN_U/d\gamma_{C,U})^{p}\big]\le\exp\!\Big(2p\big(\lvert U\rvert\ell^2/(4\theta)+2K\lvert U\rvert(\theta+\varepsilon/2)\big)\Big)\]
whenever \(Kp(\theta+\varepsilon/2)\le1/4\). The parameters are:

\begin{itemize}
\tightlist
\item
  \(\varepsilon=Ke^{-bR(x)}\);
\item
  \(\ell=g^{-1}M^2e^{-sr_x}B\,(t(x)+d)\).
\end{itemize}

The support restriction concerns the Gaussian variable \(z\); the
coefficients may depend on all the background and frozen data. The bound
retains its dependence on \(g\), \(d\) and \(|U|\).

\emph{Status:} theorem.

\emph{Lean:}
\href{https://github.com/yang-mills-group/OpenAI-YM-20261028/blob/3456e5e75cf8d4cd558955a5c7638fba2d68f464/YM/SectorTools/PHYSICAL_GAUSSIAN_INSERTION/Consumer.lean\#L167-L199}{\texttt{Physical.blocks}},
built on
\href{https://github.com/yang-mills-group/OpenAI-YM-20261028/blob/3456e5e75cf8d4cd558955a5c7638fba2d68f464/YM/SectorTools/PHYSICAL_GAUSSIAN_INSERTION/Insertion.lean\#L120-L172}{\texttt{insert\_letter}}.

\textbf{Proposition III.4.10 (jets of the block coefficients and of the
covariance).} Fix \(G\), a patch \(P\), a teller satisfying the branch
condition, and an order \(n\). There are \(r,d,v,\eta,\lambda_0,C_*>0\)
such that for large \(x\) the following holds. Take every funded step
with patch \(P\) and \(c\ge100(1+R(x))\), every \(A,F,V\) with matching
flags, every \(k\), and every logarithmic perturbation \(B\) of the
background with \(\lVert B\rVert\le r\). The store, forest and tags
remain fixed as \(B\) varies. Then:

\begin{itemize}
\item
  the move matrix \(M(B)\) is real-analytic;
\item
  \(M\), \(W=M^{\top}M\) and \(C\) satisfy weighted Schur bounds up to
  order \(n\). For every list \(\ell\) of at most \(n\) coordinate
  directions, with coarse positions \(p_b\), and output and input
  positions \(p_i,p_j\),
  \[\prod_{b\in\ell}e^{d\,\mathrm{dist}(p_b,p_i)}\sum_j e^{d\,\mathrm{dist}(p_i,p_j)}\lvert\partial_\ell X_{ij}(B)\rvert\le C_*,\]
  and the analogous column bound holds with the external weights
  anchored at \(p_j\);
\item
  \(W(B)\ge\eta I\), \(C(B)\ge\lambda_0 I\) and \(C(B)\succ0\);
\item
  for every \(u\ge0\), the full, principal and windowed resolvents have
  the same marked jet bounds with cap \(C_* /(1+u)\) at order zero and
  \(C_* /(1+u)^2\) at positive orders;
\item
  the window defect has marked jet cap \(C_*e^{-vR(x)}/(1+u)\), and the
  difference between the windowed and full resolvents has cap
  \(C_*e^{-vR(x)}/(1+u)^2\).
\end{itemize}

The selected predictor rows are real-analytic, with \(k\)-jets bounded
by constants depending only on \((G,P,K,k)\), on a ball of fixed radius
(\href{https://github.com/yang-mills-group/OpenAI-YM-20261028/blob/3456e5e75cf8d4cd558955a5c7638fba2d68f464/YM/AnalyticBridge/SECTOR_COVARIANCE_JETS/Sounds.lean\#L69-L101}{\texttt{sounds}},
\href{https://github.com/yang-mills-group/OpenAI-YM-20261028/blob/3456e5e75cf8d4cd558955a5c7638fba2d68f464/YM/AnalyticBridge/SECTOR_AFFINE_JETS/Branch.lean\#L140-L219}{\texttt{chosen\_branch}}).
The precision \(W\) is also real-analytic (Lemma 4.27); the covariance
and resolvent conclusions used here are finite-order jet bounds.

\emph{Status:} theorem.

\emph{Lean:}
\href{https://github.com/yang-mills-group/OpenAI-YM-20261028/blob/3456e5e75cf8d4cd558955a5c7638fba2d68f464/YM/AnalyticBridge/SECTOR_AFFINE_JETS/Covariance.lean\#L146-L210}{\texttt{assembly}},
\href{https://github.com/yang-mills-group/OpenAI-YM-20261028/blob/3456e5e75cf8d4cd558955a5c7638fba2d68f464/YM/AnalyticBridge/SECTOR_COVARIANCE_JETS/Sky.lean\#L62-L77}{\texttt{fromTeller}},
\href{https://github.com/yang-mills-group/OpenAI-YM-20261028/blob/3456e5e75cf8d4cd558955a5c7638fba2d68f464/YM/AnalyticBridge/SECTOR_COVARIANCE_JETS/Roof.lean\#L39-L62}{\texttt{coverall}},
\href{https://github.com/yang-mills-group/OpenAI-YM-20261028/blob/3456e5e75cf8d4cd558955a5c7638fba2d68f464/YM/PhysicalNext/PHYSICAL_JOINT_ENDPOINT_LOAD_F16/DLog/Fund.lean\#L120-L127}{\texttt{stores}}.

\textbf{Proposition III.4.11 (affine squares in the large-field
annulus).} Fix a teller, a block size \(L\), \(W_0>0\),
\(K\in\mathbb R\) and an integer halo \(D\ge0\). There are \(C\ge1\) and
an eventually nonnegative function \(e(x)\) with
\((1+R(x))^n\beta_F(x,L)e(x)\to0\) for every \(n\in\mathbb N\), such
that for large \(x\) the following holds. Take a funded step of block
size \(L\) with \(c\ge100(1+R(x))\), stores \(O,A\), a forest \(F\),
fields \(U,V\) and tags \(k\). Assume \(U\) is in tree gauge. In the
\((D+5)R(x)\)-halo of \(O\), assume the following local sector
conditions: fine plaquettes whose support meets that halo and which
belong to \(A\) have toll at least \(t_L(x)/4\), those outside \(A\)
have toll at most \(t_L(x)/2\); block bonds whose support meets that
halo and which belong to \(A\) have discrepancy at least \(W_0t_L(x)\),
those outside have discrepancy at most \(2W_0t_L(x)\); the coarse flags
match \(V\); and the chosen frames have positive mass. Then there is a
direction \(a\), vanishing on the tree links, such that for every pair
of windows \(H,Q\) with \(H_{\rm face}\subseteq Q_{\rm face}\) and
\(H_{\rm noise}\subseteq Q_{\rm noise}\), and with the face, noise and
coarse rows of \(Q\) in \(\{p\notin A^{(2R)},\ p\in O^{(DR)}\}\):

\begin{itemize}
\tightlist
\item
  \(Q\) is charted by \(a\) with constant \(C\);
\item
  the affine series on \(H\) converge, the required curvature and
  predictor derivatives exist with the stated linearisation, and
  \(\mathrm{work}_Q\ge\mathrm{squares}_H(a)-e(x)\lvert O\rvert\);
\item
  for every other \(a'\) charting \(H\) with constant \(K\), the same
  convergence, derivative statements and bound hold with \(a'\).
\end{itemize}

Here \(|O|\) is the number of cost cells, aisles and flags. The
comparison with \(a'\) uses uniqueness of the logarithm on the selected
stencils.

\emph{Status:} theorem.

\emph{Lean:}
\href{https://github.com/yang-mills-group/OpenAI-YM-20261028/blob/3456e5e75cf8d4cd558955a5c7638fba2d68f464/YM/SectorNext/LARGE_FIELD_REMAINDERS/Annulus.lean\#L36-L127}{\texttt{annulus\_pay}},
with
\href{https://github.com/yang-mills-group/OpenAI-YM-20261028/blob/3456e5e75cf8d4cd558955a5c7638fba2d68f464/YM/SectorNext/LARGE_FIELD_REMAINDERS/Clean.lean\#L110-L208}{\texttt{cleans\_at}}
and
\href{https://github.com/yang-mills-group/OpenAI-YM-20261028/blob/3456e5e75cf8d4cd558955a5c7638fba2d68f464/YM/SectorNext/LARGE_FIELD_REMAINDERS/CollarData.lean\#L131-L189}{\texttt{reacts}}.

\textbf{Proposition III.4.12 (production of the step constants).}
Assume:

\begin{itemize}
\tightlist
\item
  \(G\) is simple;
\item
  the kernel choices \(K\) satisfy the branch condition for every patch;
\item
  the initial stage is at least \(4\); the jet orders, position orders
  and support exponents are nondecreasing, the support exponents are
  positive, and at stage \(i\) the two orders are at least \(\max(6,i)\)
  and \(\max(2,i)\), respectively;
\item
  costs \(C_i\ge1\) are given.
\end{itemize}

Then there are pre-\(L\) data with kernel \(K\) and the prescribed
orders, two families of path prices, free-field and slope bounds, and a
cost bound, with the following property. For every tariff and scope
there is a sequence of block-size floors. After any increase of those
floors and any choice of block sizes above them, there is a stagewise
Hessian gap \(\eta\). For every compatible post-\(L\) cost choice (gap
\(\eta\), activity cap at least \(1\), and slope bound below the chosen
cap), protection and initial delays, the delays can be enlarged so that
the rate and chart conditions hold for all further enlargements. On
every paid run with these data there are stationary path families,
including the full stationary estimates in the retained family. For each
adjacent stage pair \((j,i)\) above the initial stage, the per-step
chart conditions hold with \(a_{ji},E_{ji}\) chosen after \(L_j\)
(\href{https://github.com/yang-mills-group/OpenAI-YM-20261028/blob/3456e5e75cf8d4cd558955a5c7638fba2d68f464/YM/SectorNext/GAUSSIAN_SECTOR_CHANGE/Stamped.lean\#L139-L168}{\texttt{sponsored\_permit}}).
\href{https://github.com/yang-mills-group/OpenAI-YM-20261028/blob/3456e5e75cf8d4cd558955a5c7638fba2d68f464/YM/SectorNext/GAUSSIAN_SECTOR_CHANGE/Handoff.lean\#L99-L116}{\texttt{issue\_cosmos}}
replaces each entry value by the maximum of its old value and the
thresholds for \((j,j)\) and \((j,j-1)\), preserving the rate
conditions.

\emph{Status:} theorem under these kernel and order assumptions.

\emph{Lean:}
\href{https://github.com/yang-mills-group/OpenAI-YM-20261028/blob/3456e5e75cf8d4cd558955a5c7638fba2d68f464/YM/MergeProof/Seaside.lean\#L96-L154}{\texttt{sailsV}},
using
\href{https://github.com/yang-mills-group/OpenAI-YM-20261028/blob/3456e5e75cf8d4cd558955a5c7638fba2d68f464/YM/Next/STATIONARY/FullPlanned.lean\#L59-L66}{\texttt{geared}}.

\subsubsection{2. What it uses}\label{what-it-uses-6}

\begin{itemize}
\tightlist
\item
  \textbf{III.1 (block-spin step and gauge fixing):} fine and coarse
  tori with \(f=Lc\), tree-gauge slices, gauge tags, predictor and
  push-forward
  (\href{https://github.com/yang-mills-group/OpenAI-YM-20261028/blob/3456e5e75cf8d4cd558955a5c7638fba2d68f464/YM/RG/Block.lean\#L22-L25}{\texttt{StepData}}).
  The branch condition requires unique frame and predictor minimisers on
  the selected supports and existence of an eligible axial gauge
  (\href{https://github.com/yang-mills-group/OpenAI-YM-20261028/blob/3456e5e75cf8d4cd558955a5c7638fba2d68f464/YM/RG/Patch.lean\#L182-L185}{\texttt{Branches}}).
\item
  \textbf{II.2, III.5, III.8 (representations and activities):} path and
  activity bounds for the fine density
  (\href{https://github.com/yang-mills-group/OpenAI-YM-20261028/blob/3456e5e75cf8d4cd558955a5c7638fba2d68f464/YM/RG/History.lean\#L181-L183}{\texttt{Rep.Lawful}}),
  large-field allocations, and the polymer summability norm of III.4.7.
\item
  \textbf{III.10 (step estimator):} the tube's fine, coarse and noise
  work rows, together with their path bounds, cutoff compatibility and
  free-metric data
  (\href{https://github.com/yang-mills-group/OpenAI-YM-20261028/blob/3456e5e75cf8d4cd558955a5c7638fba2d68f464/YM/Pressure/Seeds.lean\#L101-L109}{\texttt{Funds}}).
  These determine \(M\) and \(s(Y)\).
\item
  \textbf{III.6--III.7 (stationary branch and fourth-order matching):}
  stationary path families and their finite-volume estimates
  (\href{https://github.com/yang-mills-group/OpenAI-YM-20261028/blob/3456e5e75cf8d4cd558955a5c7638fba2d68f464/YM/Next/STATIONARY/Level/Recurse.lean\#L301-L305}{\texttt{FullPaths}}).
  Their production in III.4.12 uses the free-field estimates of III.3.
\item
  \textbf{III.2 and III.9 (schedule and marginal flow):} ordered
  constants, selected histories, the running inverse coupling and its
  band. III.4.2 applies to these data through its directed-run
  hypotheses.
\item
  \textbf{I.1 and III.16 (analytic toolkit):} the matrix gauge group,
  normalised Haar measure, exponential charts, analytic jets and
  finite-dimensional Gaussian measures.
\end{itemize}

The analytic tools are Gaussian affine changes of variables, the Haar
Jacobian in exponential coordinates, Combes--Thomas estimates, the
resolvent integral for \(\log\), Neumann series, Gaussian integration by
parts and Fourier decay from analytic strips. General references are
Helgason, \emph{Differential Geometry, Lie Groups and Symmetric Spaces};
Glimm--Jaffe, \emph{Quantum Physics}; and Katznelson, \emph{An
Introduction to Harmonic Analysis}. Balaban's fluctuation-integral
papers and Dimock's \emph{The renormalization group according to
Balaban} provide related RG constructions; precise bibliographic
locations remain open (§7).

The order of constants is pre-\(L\) data, block sizes, post-\(L\)
bounds, then thresholds in \(x\). Statements with additional parameters
retain those dependencies. The precision and path-counting certificates
come from
\href{https://github.com/yang-mills-group/OpenAI-YM-20261028/blob/3456e5e75cf8d4cd558955a5c7638fba2d68f464/YM/AnalyticBridge/RELATIVE_LOGDET_TRACE/Year.lean\#L113-L197}{\texttt{late\_price}},
and the determinant-jet fund from Lemma 4.21. The packet-energy, chart,
noise and representation hypotheses in III.4.7 are separate assumptions.
The branch and order assumptions of III.4.12, and the rate assumptions
of the scheduling statements in Lemma 4.33, remain explicit.

\subsubsection{3. Strategy}\label{strategy-6}

Keep links within \(10R(x)\) of the large-field store as Haar variables.
Write the others in exponential coordinates about the coarse-field copy
\(\iota V\), and separate the affine squares
\(Q_Y(a)=\tfrac12|Ma+s(Y)|^2\) from the remaining energy. The
substitution \(a=m(Y)+gz\) and division by an arbitrary Gaussian density
give the exact identity of III.4.1. Choose the symmetrised windowed
inverse for \(C\) and the windowed mean for \(m\). The endpoint windows
give finite range, at most \(\lfloor R(x)/20\rfloor\); the linear
correction and \(C^{-1}-W\) remain in the integrand.

A local Poincaré inequality in tree gauge gives the gap of \(W\).
Predictor-derivative comparison permits absorption of the gauge errors
at small threshold. Path decay gives a uniform weighted Schur bound;
Combes--Thomas and the window law then give \(C=W^{-1}+O(e^{-cR(x)})\)
and positivity. The resolvent integral for \(\log\), followed by a
Neumann expansion, expresses the relative Gaussian normalisation as
single-window and path terms. Regroup these near the large-field cores.
Their logarithmic cost is absorbed by the packet energy
\(\beta s^2\gtrsim\mathrm{hot}(z)^2\), leaving exponentially suppressed
core factors and a small polymer remainder on the cutoff support.

Gaussian moment bounds control the exponential corrections, marked jets
control background derivatives, and the annulus estimates recover affine
squares away from the store. The construction follows the
fluctuation-integral method, with window errors retained explicitly.
III.4.12 chooses the constants in the order needed along the run.

\subsubsection{4. Steps}\label{steps-6}

\paragraph{4.1 Exact change of
variables}\label{exact-change-of-variables}

\textbf{Lemma 4.1 (Gaussian density and scaling).} For \(C\succ0\),
\(g>0\), \(m\in\mathbb R^n\), and any Banach-valued \(f\) on
\(\mathbb R^n\), \(f\in L^1\) iff
\(z\mapsto e^{\frac12z\cdot C^{-1}z}f(m+gz)\in L^1(\gamma_C)\), and
\[\int f=g^{n}(2\pi)^{n/2}(\det C)^{1/2}\int e^{\frac12z\cdot C^{-1}z}f(m+gz)\,d\gamma_C.\]
Lean's Bochner integral assigns value zero to a nonintegrable function,
so the identity also holds under that convention without an
integrability premise. \emph{Sketch:} the density of \(\gamma_C\)
follows by pushing \(N(0,1)^{\otimes n}\) forward by \(C^{1/2}\)
(\href{https://github.com/yang-mills-group/OpenAI-YM-20261028/blob/3456e5e75cf8d4cd558955a5c7638fba2d68f464/YM/SectorNext/GAUSSIAN_SECTOR_CHANGE/Gaussian.lean\#L134-L156}{\texttt{cloud\_formula}});
then translate and dilate.

\emph{Lean:}
\href{https://github.com/yang-mills-group/OpenAI-YM-20261028/blob/3456e5e75cf8d4cd558955a5c7638fba2d68f464/YM/SectorNext/GAUSSIAN_SECTOR_CHANGE/Gaussian.lean\#L204-L210}{\texttt{Gaussian.gauged}},
\href{https://github.com/yang-mills-group/OpenAI-YM-20261028/blob/3456e5e75cf8d4cd558955a5c7638fba2d68f464/YM/SectorNext/GAUSSIAN_SECTOR_CHANGE/Counting.lean\#L117-L125}{\texttt{at\_price}}.

\textbf{Lemma 4.2 (exponential chart on links).} For a local chart ratio
\(J\) and base links \(\nu\), the map \(a\mapsto(e^{a_e}\nu_e)_e\) sends
\(\prod_e J_{\rm sc}J(a_e)\,da_e\) on \(B_{\rho_J}^{\mathcal A}\) onto
product Haar measure restricted to the image. Hence
\(\int g\,d\mathrm{Haar}=\int_{B_{\rho_J}^{\mathcal A}}\prod_eJ_{\rm sc}J(a_e)\,g(e^{a}\nu)\,da\)
for strongly measurable \(g\) vanishing off the image, with equivalence
of integrability. In the iterated form on the slice (frozen links Haar,
active links in the chart), assume also that \(g\) is bounded; then
every chart section and the outer integral are integrable.

\emph{Lean:}
\href{https://github.com/yang-mills-group/OpenAI-YM-20261028/blob/3456e5e75cf8d4cd558955a5c7638fba2d68f464/YM/SectorNext/GAUSSIAN_SECTOR_CHANGE/Haar.lean\#L143-L162}{\texttt{log\_transport}},
\href{https://github.com/yang-mills-group/OpenAI-YM-20261028/blob/3456e5e75cf8d4cd558955a5c7638fba2d68f464/YM/SectorNext/GAUSSIAN_SECTOR_CHANGE/Haar.lean\#L167-L184}{\texttt{release\_integral}},
\href{https://github.com/yang-mills-group/OpenAI-YM-20261028/blob/3456e5e75cf8d4cd558955a5c7638fba2d68f464/YM/SectorNext/GAUSSIAN_SECTOR_CHANGE/Haar.lean\#L225-L253}{\texttt{thaw\_integral}}.

\textbf{Lemma 4.3 (regularised integrand).} Assume the branch condition,
a local chart ratio, the standing condition, \(d>0\) and \(4d<\rho_J\).
For a multiplier with a bounded Borel envelope, the chart integrand
equals \(J_{\rm sc}^{\lvert\mathcal A\rvert}\tilde J_d(a)\chi_d(a)\)
times the sector multiplier everywhere. Its sector integral equals
\(J_{\rm sc}^{\lvert\mathcal A\rvert}\iint r\,da\,d\mathrm{Haar}(Y)\),
with section-wise, outer and joint integrability. \emph{Sketch:} on the
support, \(\lvert a_e\rvert\le d\) forces \(\chi_d=1\) and each extended
Jacobian to equal \(J\); off the support both sides vanish
(\href{https://github.com/yang-mills-group/OpenAI-YM-20261028/blob/3456e5e75cf8d4cd558955a5c7638fba2d68f464/YM/SectorNext/GAUSSIAN_SECTOR_CHANGE/Broad.lean\#L115-L120}{\texttt{protected\_pay}}).
The envelope supplies boundedness and measurability.

\emph{Lean:}
\href{https://github.com/yang-mills-group/OpenAI-YM-20261028/blob/3456e5e75cf8d4cd558955a5c7638fba2d68f464/YM/SectorNext/GAUSSIAN_SECTOR_CHANGE/Shell.lean\#L76-L109}{\texttt{sheltered}},
\href{https://github.com/yang-mills-group/OpenAI-YM-20261028/blob/3456e5e75cf8d4cd558955a5c7638fba2d68f464/YM/SectorNext/GAUSSIAN_SECTOR_CHANGE/Shell.lean\#L142-L199}{\texttt{regularized\_price}}.

\textbf{Lemma 4.4 (affine-square splitting).} Under the fine path
bounds, the averaged step energy equals
\(\overline{\mathrm{res}}+Q_Y(a)\)
(\href{https://github.com/yang-mills-group/OpenAI-YM-20261028/blob/3456e5e75cf8d4cd558955a5c7638fba2d68f464/YM/SectorNext/GAUSSIAN_SECTOR_CHANGE/Remainder.lean\#L102-L105}{\texttt{compose\_residual}}).
For any \(C,m,z\) and \(g>0\), with \(\beta=g^{-2}\) and
\(\eta=-g^{-1}(Wm+M^{\top}s(Y))\):
\[-\beta Q_Y(m)+\eta\cdot z+\tfrac12z\cdot(C^{-1}-W)z=\tfrac12z\cdot C^{-1}z-\beta Q_Y(m+gz).\]
\emph{Sketch:} expand \(\tfrac12\lvert M(m+gz)+s\rvert^2\) and use
\(g^{-2}=\beta\).

\emph{Lean:}
\href{https://github.com/yang-mills-group/OpenAI-YM-20261028/blob/3456e5e75cf8d4cd558955a5c7638fba2d68f464/YM/SectorNext/GAUSSIAN_SECTOR_CHANGE/Remainder.lean\#L128-L146}{\texttt{Rows.adjusted}},
\href{https://github.com/yang-mills-group/OpenAI-YM-20261028/blob/3456e5e75cf8d4cd558955a5c7638fba2d68f464/YM/SectorNext/GAUSSIAN_SECTOR_CHANGE/Change.lean\#L73-L93}{\texttt{recolored}}.

\textbf{Lemma 4.5 (proof of III.4.1).} Apply Lemma 4.1 to \(r(Y,\cdot)\)
for each \(Y\), then Lemma 4.4 to identify the Gaussian integrand with
\(\mathrm{pref}\cdot\Psi(Y,\cdot)\). The same change applied to norms
multiplies their Gaussian integral by the positive constant
\(\mathrm{pref}\,(2\pi g^2)^{n_a/2}(\det C)^{1/2}\). Lemma 4.3 and the
Fubini criterion therefore give joint integrability. Inserting the
section identity into the outer Haar integral proves the formula.

\emph{Lean:}
\href{https://github.com/yang-mills-group/OpenAI-YM-20261028/blob/3456e5e75cf8d4cd558955a5c7638fba2d68f464/YM/SectorNext/GAUSSIAN_SECTOR_CHANGE/Change.lean\#L137-L192}{\texttt{authorized\_change}}.

\textbf{Lemma 4.6 (admission of the chart).} Fix a teller, local chart
ratio, block size \(L\) and allocation width \(w>0\). There are
\(D,E>0\) such that for large \(x\), with \(d=D\,t_L(x)\):

\begin{itemize}
\tightlist
\item
  \(0<d\), \(4d<\rho_J\) and \(2d<\rho_{\rm first}\);
\item
  \(\lvert\log J(a)\rvert\le E\lvert a\rvert^2\) for
  \(\lvert a\rvert\le4d\);
\item
  for every funded step of block size \(L\) with \(c\ge100(1+R(x))\) and
  stencil span at most \(R(x)/50\), and every store, forest, allocation
  of width \(w\) and \(V\), each active link (and each frozen link read
  on a selected stencil) in the support of
  \(\mathrm{bill}\cdot\mathrm{lik}\) equals \(e^{a}\nu_e\) with
  \(\lvert a\rvert\le d\).
\end{itemize}

\emph{Sketch:} the allocation gives toll \(\le t_L(x)\) off the cost
cells and block discrepancy \(\le2w\,t_L(x)\) off the aisles. The
axial-gauge estimate bounds the chord distance to \(\nu_e\) by
\(m\,t_L(x)\). The logarithmic chart then gives
\(\lvert a\rvert\le C m\,t_L(x)\).

\emph{Lean:}
\href{https://github.com/yang-mills-group/OpenAI-YM-20261028/blob/3456e5e75cf8d4cd558955a5c7638fba2d68f464/YM/SectorNext/GAUSSIAN_SECTOR_CHANGE/Stamped.lean\#L45-L119}{\texttt{matinee}}.

\textbf{Lemma 4.7 (permits and the licence; proof of III.4.2).} Under
the fixed patch, branch, funding and volume hypotheses, combine Lemma
4.6, III.4.3 and continuity of the windowed mean. For large \(x\),
matching flags and every \(\beta>0\), \((W_{\mathcal W},m)\) is
permitted, \(W\succ0\), and the selected predictor derivatives are
bounded. Enlarge the entry thresholds to cover every admissible stage
pair. Then:

\begin{itemize}
\tightlist
\item
  apply III.4.1 to stores with matching flags;
\item
  wrong flags force \(\mathrm{bill}_A\equiv0\), so all integrals vanish;
\item
  insert both identities into the sector decomposition of the
  push-forward, using tagged kernel times density
  \(=\mathrm{sieve}(V)\cdot\mathrm{exposed}\cdot\mathrm{enters}\).
\end{itemize}

The step size is \(d=a_{ji}t(\hat\beta_{m+1})\), because
\(t_L(\hat\beta_m)=t(\hat\beta_{m+1})\).

\emph{Lean:}
\href{https://github.com/yang-mills-group/OpenAI-YM-20261028/blob/3456e5e75cf8d4cd558955a5c7638fba2d68f464/YM/SectorNext/GAUSSIAN_SECTOR_CHANGE/Stamped.lean\#L139-L168}{\texttt{sponsored\_permit}},
\href{https://github.com/yang-mills-group/OpenAI-YM-20261028/blob/3456e5e75cf8d4cd558955a5c7638fba2d68f464/YM/SectorNext/GAUSSIAN_SECTOR_CHANGE/Handoff.lean\#L132-L182}{\texttt{committed}}.

\textbf{Lemma 4.8 (identity completion).} Assume \(W\ge\eta I\) with
\(\eta>0\) and \(R(x)\ge120\). Complete \(W\) by the identity on the
inactive positions (frozen and tree links), with the corresponding
endpoint windows. The windowed inverse and its symmetrisation are the
identity completions of \(W^{\mathrm{per}}\) and \(W_{\mathcal W}\); the
mean is extended by zero. The determinant is unchanged. This puts the
actual and flat log-determinants on the same full link set.

\emph{Lean:}
\href{https://github.com/yang-mills-group/OpenAI-YM-20261028/blob/3456e5e75cf8d4cd558955a5c7638fba2d68f464/YM/SectorNext/GAUSSIAN_SECTOR_CHANGE/Completion.lean\#L181-L212}{\texttt{Completion.stocked}}.

\paragraph{4.2 Gap, decay and windows}\label{gap-decay-and-windows}

\textbf{Lemma 4.9 (gap of the active precision).} Under the patch,
branch, funding and volume hypotheses of III.4.3, there is
\(\eta=\eta(G,P,\text{teller})>0\) with
\(\eta\lvert a\rvert^2\le a\cdot Wa\) for large \(x\), every
flag-matched sector and every active vector \(a\). \emph{Sketch:}

\begin{itemize}
\tightlist
\item
  Each active link is bounded by a constant times a local sum of
  plaquette jets plus a block jet, computed in a gauge chosen per link.
  The comparison of the predictor's derivative with its value at the
  flat field
  (\href{https://github.com/yang-mills-group/OpenAI-YM-20261028/blob/3456e5e75cf8d4cd558955a5c7638fba2d68f464/YM/SectorNext/GAUSSIAN_SECTOR_CHANGE/Block.lean\#L316-L361}{\texttt{block\_comparison}})
  makes the gauge-rotation errors \(O(\varepsilon)\) once
  \(t(x)\le\delta\).
\item
  Summing gives
  \(\lvert a\rvert^2\le c_1\sum\lvert\text{plaq}\rvert^2+c_2\sum\lvert\text{block}\rvert^2+O(\varepsilon^2)\lvert a\rvert^2\).
\item
  The last term is absorbed, and the energy is bounded by the Gram form.
\end{itemize}

\emph{Lean:}
\href{https://github.com/yang-mills-group/OpenAI-YM-20261028/blob/3456e5e75cf8d4cd558955a5c7638fba2d68f464/YM/SectorNext/GAUSSIAN_SECTOR_CHANGE/Gap.lean\#L242-L264}{\texttt{Gap.not\_flat}}.

the derivative-comparison theorem itself states two things. For links
within the predictor's reach that are \(\delta\)-close to \(1\), the
derivative of \(a\mapsto\) curvature of
\(\mathrm{pred}(e^aU)\mathrm{pred}(U)^{-1}\) at \(0\) is
\(\le B\lVert a\rVert_{\ell^1(\text{stencil})}\). It also differs from
its value at \(U=1\) by \(\le\varepsilon\lVert a\rVert_{\ell^1}\). The
constant \(\delta(\varepsilon)\) comes from continuity and is not
explicit.

\textbf{Lemma 4.10 (Combes--Thomas).} Given \(\sigma,\varepsilon>0\) and
\(M\ge0\), there are \(v\in(0,\sigma]\) and \(C=2/\varepsilon\) such
that every symmetric \(A\) with
\(\lVert e^{\sigma\,\mathrm{dist}}A\rVert_{\rm Schur}\le M\) and gap
\(\varepsilon\) satisfies
\(e^{v\,\mathrm{dist}(i,j)}\lvert A^{-1}_{ij}\rvert\le C\), in any
volume. The same holds for all principal submatrices, since they inherit
the gap. \emph{Sketch:}

\begin{itemize}
\tightlist
\item
  Conjugate \(A\) by \(e^{v\,\mathrm{dist}(\cdot,i)}\) with
  \(v=\min(1,\varepsilon/(2(M+1)))\sigma\).
\item
  The conjugate differs from \(A\) by \(\le\varepsilon/2\) in Schur
  norm, so it keeps gap \(\varepsilon/2\).
\item
  Test it on the weighted \(i\)-th column of \(A^{-1}\).
\end{itemize}

\emph{Lean:}
\href{https://github.com/yang-mills-group/OpenAI-YM-20261028/blob/3456e5e75cf8d4cd558955a5c7638fba2d68f464/YM/SectorNext/GAUSSIAN_SECTOR_CHANGE/Inverse.lean\#L74-L121}{\texttt{Schur.inverse\_tilt}},
\href{https://github.com/yang-mills-group/OpenAI-YM-20261028/blob/3456e5e75cf8d4cd558955a5c7638fba2d68f464/YM/SectorNext/GAUSSIAN_SECTOR_CHANGE/Windows.lean\#L103-L119}{\texttt{Schur.provides}}.

\textbf{Lemma 4.11 (proof of the window law, III.4.4).}

\begin{enumerate}
\def\labelenumi{\arabic{enumi}.}
\tightlist
\item
  \((AL)_{ji}=\delta_{ji}\) for \(j\in\mathcal W_i\), because column
  \(i\) of \(L\) solves the window system. So \(AL-1\) is supported at
  distance \(>r\) and has unweighted Schur norm
  \(\le e^{-ur}\cdot\)const.
\item
  Hence \(L-A^{-1}=A^{-1}(AL-1)\) has the same size, and so does its
  symmetrisation.
\item
  The gap of \(A^{-1}\) is at least \(\eta/(K+1)^2\), because
  \(\lVert A\rVert\le K+1\). For \(r\) large the error is below half of
  it.
\end{enumerate}

\emph{Lean:}
\href{https://github.com/yang-mills-group/OpenAI-YM-20261028/blob/3456e5e75cf8d4cd558955a5c7638fba2d68f464/YM/SectorNext/GAUSSIAN_SECTOR_CHANGE/Windows.lean\#L184-L239}{\texttt{Schur.windows}}.

\textbf{Lemma 4.12 (proof of III.4.3).}

\begin{enumerate}
\def\labelenumi{\arabic{enumi}.}
\tightlist
\item
  Lemma 4.9 gives the gap.
\item
  The path-decay and predictor-derivative bounds give weighted row and
  column bounds for \(M\); multiplying \(M^{\top}M\) gives the weighted
  Schur bound on \(W\).
\item
  The counting function is
  \(B(v)=\dim\mathfrak g\cdot L^4\cdot4\cdot\sum_{z\in\mathbb Z^4}e^{-v|z|_1}\)
  for \(v>0\): each coarse site carries at most \(4L^4\dim\mathfrak g\)
  coordinates.
\item
  The endpoint windows contain \(r_x\)-balls.
\item
  III.4.4 gives (iii) and \(C\succ0\).
\end{enumerate}

\emph{Lean:}
\href{https://github.com/yang-mills-group/OpenAI-YM-20261028/blob/3456e5e75cf8d4cd558955a5c7638fba2d68f464/YM/SectorNext/GAUSSIAN_SECTOR_CHANGE/Chosen.lean\#L102-L143}{\texttt{opaqueCost}}.
The decay certificate
\href{https://github.com/yang-mills-group/OpenAI-YM-20261028/blob/3456e5e75cf8d4cd558955a5c7638fba2d68f464/YM/AnalyticBridge/ANNULUS_RATIO_INSERTIONS/Chord.lean\#L133-L158}{\texttt{holds}}
adds bounds for \(C^{-1}\) and \(\lVert C^{-1}-W\rVert\le Me^{-aR(x)}\).

\textbf{Lemma 4.13 (perturbation of inverses).} Let \(s\ge0\) and
\(A,D\succ0\) with \(\lVert D^{-1}\rVert_s\le K\),
\(\lVert A-D\rVert_s\le\delta\) and \(K\delta\le\frac12\). Then
\(\lVert A^{-1}\rVert_s\le2K\). \emph{Sketch:} from
\(A^{-1}=D^{-1}-D^{-1}(A-D)A^{-1}\), a bound \(t\) on \(A^{-1}\)
improves to \(K+t/2\); iterate from any finite-volume bound and pass to
the limit.

\emph{Lean:}
\href{https://github.com/yang-mills-group/OpenAI-YM-20261028/blob/3456e5e75cf8d4cd558955a5c7638fba2d68f464/YM/SectorTools/PHYSICAL_GAUSSIAN_SWITCH/TableCosts.lean\#L63-L83}{\texttt{inverse\_perturb}}.

\textbf{Lemma 4.14 (infinite-volume kernels).}

\begin{itemize}
\tightlist
\item
  Let \(r>0\), \(M\ge0\). A family of periodic kernels compatible under
  coverings and bounded pointwise by \(Me^{-r\,\mathrm{dist}}\) is, for
  every \(\gamma<r\), exactly the periodisation of a kernel on
  \(\mathbb Z^4\) with strip norm at rate \(\gamma\) bounded by
  \(M\sum_{z\in\mathbb Z^4}e^{-(r-\gamma)|z|_1}\). The proof extracts a
  pointwise limit along the periods \(k!\) and passes compatibility to
  the limit by dominated convergence.
\item
  Given a strip bound with positive rate and a uniform positive gap in
  every period, each nonnegative-shift resolvent is the periodisation of
  a strip-decaying kernel at any rate below the resulting inverse-decay
  margin
  (\href{https://github.com/yang-mills-group/OpenAI-YM-20261028/blob/3456e5e75cf8d4cd558955a5c7638fba2d68f464/YM/Push/FREE_RESOLVENT/Metric.lean\#L30-L37}{\texttt{strip\_resolvent}}).
\item
  Under the analytic-strip inverse bound and real-symmetry condition on
  the symbol, Fourier inversion gives a real inverse kernel with strip
  decay at rate \(\eta/4\)
  (\href{https://github.com/yang-mills-group/OpenAI-YM-20261028/blob/3456e5e75cf8d4cd558955a5c7638fba2d68f464/YM/RG08/Gaussian/PeriodicInverse.lean\#L104-L139}{\texttt{freezer\_row}}).
  Consequently, for each stage \(J\ge4\) there is a rate \(\sigma>0\),
  chosen before the patch, such that every patch at stage \(J\)
  satisfying the branch condition has a strip-decaying inverse of its
  flat gauge-completed Hessian. The cap may depend on the patch; both
  rate and cap are uniform in the permitted kinetic profiles, gauge tags
  and periods
  (\href{https://github.com/yang-mills-group/OpenAI-YM-20261028/blob/3456e5e75cf8d4cd558955a5c7638fba2d68f464/YM/RG08/Gaussian/Completion.lean\#L62-L110}{\texttt{joint\_universal}}).
\end{itemize}

\emph{(proof read for periodic extraction and the uniform flat inverse;
statement only for the other assertions)}

\emph{Lean:}
\href{https://github.com/yang-mills-group/OpenAI-YM-20261028/blob/3456e5e75cf8d4cd558955a5c7638fba2d68f464/YM/Push/FREE_RESOLVENT/Fold.lean\#L139-L180}{\texttt{publish}}.

\textbf{Lemma 4.15 (tree inverse and the free covariance).}

\begin{itemize}
\tightlist
\item
  Gauge-fixing completion: for a symmetric free metric and a step
  satisfying the branch condition, if the gauge-completed joint matrix
  \(\mathcal H\) satisfies \(\mathcal HN=I\), then \(ZNZ^{\top}\)
  inverts the tree matrix and the projected \(N\) is the tree covariance
  (\href{https://github.com/yang-mills-group/OpenAI-YM-20261028/blob/3456e5e75cf8d4cd558955a5c7638fba2d68f464/YM/RG08/Gaussian/TreeInverse.lean\#L64-L115}{\texttt{projecting}},
  pure algebra, proof read).
\item
  Fix a patch, kernel satisfying the branch condition, and positive
  strip and gap bounds. Each free metric satisfying these bounds has a
  strip-decaying operator representing its plane tree-inverse covariance
  at every coarse period at least \(3\). The rate and cap are uniform
  over these metrics. With odd \(L\) and the fine and coarse free-field
  lower bounds and diagonal tower compatibility, the same operator
  represents the conditional covariances of every admissible step with
  that patch. Its range-\(d\) truncation has strip norm at rate \(t/2\)
  at most \(Ae^{-td/2}\)
  (\href{https://github.com/yang-mills-group/OpenAI-YM-20261028/blob/3456e5e75cf8d4cd558955a5c7638fba2d68f464/YM/AnalyticBridge/LOCAL_GAUSSIAN_LAWS/Covers.lean\#L67-L103}{\texttt{earned}},
  \href{https://github.com/yang-mills-group/OpenAI-YM-20261028/blob/3456e5e75cf8d4cd558955a5c7638fba2d68f464/YM/AnalyticBridge/LOCAL_GAUSSIAN_LAWS/Covers.lean\#L114-L136}{\texttt{onePlateau}},
  proof read).
\item
  For a large step, a tube satisfying the fine path bounds, a sector
  with no frozen free links, and an attached, normalised free metric
  whose full kernel has Gram matrix equal to that metric and satisfies
  the free-field symmetry conditions, the windowed resolvent at flat
  background differs from the free tree resolvent by at most
  \(Ce^{-tR(x)}(1+u)^{-2}\) in a positive-rate Schur norm, for every
  \(u\ge0\). The statement applies to the geometric windows used in this
  comparison. At \(u=0\) it also gives precision and log-determinant
  comparison bounds
  (\href{https://github.com/yang-mills-group/OpenAI-YM-20261028/blob/3456e5e75cf8d4cd558955a5c7638fba2d68f464/YM/AnalyticBridge/LOCAL_GAUSSIAN_LAWS/Ideal.lean\#L120-L144}{\texttt{headComparison}},
  \href{https://github.com/yang-mills-group/OpenAI-YM-20261028/blob/3456e5e75cf8d4cd558955a5c7638fba2d68f464/YM/AnalyticBridge/LOCAL_GAUSSIAN_LAWS/HeadSwitch.lean\#L104-L143}{\texttt{freeHead}},
  statement only).
\end{itemize}

\textbf{Lemma 4.16 (covariance and locality of entries).}

\begin{itemize}
\tightlist
\item
  The windowed inverse commutes with site bijections combined with
  orthogonal colour rotations, for every \(u\)
  (\href{https://github.com/yang-mills-group/OpenAI-YM-20261028/blob/3456e5e75cf8d4cd558955a5c7638fba2d68f464/YM/AnalyticBridge/LOCAL_GAUSSIAN_LAWS/ColourBooks.lean\#L163-L192}{\texttt{natural}}).
\item
  If the precision, source and windows are transported together by such
  a map, so are the ordered inverse, windowed covariance and mean; the
  covariance determinant is preserved
  (\href{https://github.com/yang-mills-group/OpenAI-YM-20261028/blob/3456e5e75cf8d4cd558955a5c7638fba2d68f464/YM/AnalyticBridge/LOCAL_GAUSSIAN_LAWS/Binding.lean\#L116-L134}{\texttt{moves}}).
\item
  Under the branch condition, \(R(x)\ge120\) and stencil span at most
  \(R(x)/50\), agreement of the configurations, tags, trees and stores
  on a set containing the \(14R(x)\)-ball about the row position makes
  the corresponding full precision rows equal
  (\href{https://github.com/yang-mills-group/OpenAI-YM-20261028/blob/3456e5e75cf8d4cd558955a5c7638fba2d68f464/YM/AnalyticBridge/RELATIVE_LOGDET_TRACE/Reading.lean\#L194-L212}{\texttt{same\_row}}).
\end{itemize}

\paragraph{4.3 Determinants}\label{determinants}

\textbf{Lemma 4.17 (proof of III.4.5).}

\begin{enumerate}
\def\labelenumi{\arabic{enumi}.}
\tightlist
\item
  On the spectrum of \(A\succ0\),
  \(\log\lambda=\int_0^\infty[(1+u)^{-1}-(\lambda+u)^{-1}]du\), so
  \(\log\det A=\sum_i\int_0^\infty[(1+u)^{-1}-(A+u)^{-1}_{ii}]du\).
\item
  Split \((A+uI)^{-1}=P_u+((A+uI)^{-1}-P_u)\).
\item
  The first part gives \(\ell_i\), bounded by \((C+1)C\) through
  \(\int(1+u)^{-2}=1\).
\item
  Expand the second part as \(P_u\sum_{n\ge1}(-\Delta_u)^n\). Its
  negative enters the logarithm. The \(n\)th correction in III.4.5 has
  absolute integral at most \(Cq^{n+1}\).
\item
  For \(\log\det(1+S)\) use the trace series; \(1+S=AA_{\mathcal W}\).
\end{enumerate}

\emph{Lean:}
\href{https://github.com/yang-mills-group/OpenAI-YM-20261028/blob/3456e5e75cf8d4cd558955a5c7638fba2d68f464/YM/AnalyticBridge/RELATIVE_LOGDET_TRACE/Series.lean\#L203-L227}{\texttt{Priced.chains}},
\href{https://github.com/yang-mills-group/OpenAI-YM-20261028/blob/3456e5e75cf8d4cd558955a5c7638fba2d68f464/YM/AnalyticBridge/RELATIVE_LOGDET_TRACE/Series.lean\#L45-L46}{\texttt{letterBase}}.

\textbf{Lemma 4.18 (proof of III.4.6 (ii)--(iv)).}

\begin{itemize}
\tightlist
\item
  \emph{(ii)} Use
  \(J_{\rm sc}^{\lvert\mathcal A\rvert}(2\pi g^2)^{n_a/2}=\mathrm{spare}^{\lvert F^{\rm free}\rvert}\mathrm{spare}^{-\lvert\mathcal F_A\rvert}\)
  and \(\det C=\det W^{-1}\det(WC)\), with \(\det K_F\) independent of
  the forest. The quadratic part follows from Lemma 4.4, and the coarse
  factors are collected in the coarse normalisation.
\item
  \emph{(iii)} Apply the path-series expansion both to the
  identity-completed full precision at \((A,V)\) (``panel'') and at
  \((\emptyset,1)\) (``frame''), and subtract letter by letter.
\item
  \emph{(iv)} Sum
  \(\mathrm{sect}_e-\log\mathrm{spare}\cdot1_{\mathcal F_A}(e)\) over
  links, split by distance to the cores (\(200R(x)\)).
\end{itemize}

\emph{Lean:}
\href{https://github.com/yang-mills-group/OpenAI-YM-20261028/blob/3456e5e75cf8d4cd558955a5c7638fba2d68f464/YM/PhysicalNext/AFTERNOON_INTERACTING_RANK_COMPARISON/Opacity/Shards.lean\#L34-L39}{\texttt{book\_price}},
\href{https://github.com/yang-mills-group/OpenAI-YM-20261028/blob/3456e5e75cf8d4cd558955a5c7638fba2d68f464/YM/PhysicalNext/AFTERNOON_INTERACTING_RANK_COMPARISON/Opacity/Enigma.lean\#L36-L48}{\texttt{rate\_pass}}.

\textbf{Lemma 4.19 (local bounds on the pieces).} Call a function
\emph{\(\varepsilon\)-small} if it is a sum of local terms \(\phi_X\) on
flag-matched sectors where the cutoff product is nonzero, with
\(\lvert\phi_X\rvert\le b_X\) on every flag-matched sector and
\(\sum_{X\ni y'}b_Xe^{eN_B(X)}\le\varepsilon\) at every site \(y'\).
Here \(B\) is the specified cover radius. Under the corresponding
hypotheses of III.4.7, for large \(x\):

\begin{itemize}
\tightlist
\item
  \begin{enumerate}
  \def\labelenumi{(\alph{enumi})}
  \tightlist
  \item
    At the retimed scale \(y=\mathrm{retime}(x)\), with
    \(z=\beta_F(x,L)\) and \(\beta\in[z/2,2z]\),
    \(\lvert\mathrm{rent}_o\rvert\le\delta\,\mathrm{hot}(z)^2\lvert o\rvert\)
    for any fixed \(\delta>0\)
    (\href{https://github.com/yang-mills-group/OpenAI-YM-20261028/blob/3456e5e75cf8d4cd558955a5c7638fba2d68f464/YM/PhysicalNext/AFTERNOON_INTERACTING_RANK_COMPARISON/Opacity/Enigma.lean\#L90-L114}{\texttt{wages}}).
    Each link contributes \(O(\log y)\), the \(200R(y)\)-halo of a core
    contains \(O((1+R(y))^4\lvert o\rvert)\) links, and
    \(\mathrm{hot}(z)^2\) dominates this logarithmic cost.
  \end{enumerate}
\item
  (b) The far-field term is \(x^{-1/8}\)-small
  (\href{https://github.com/yang-mills-group/OpenAI-YM-20261028/blob/3456e5e75cf8d4cd558955a5c7638fba2d68f464/YM/PhysicalNext/AFTERNOON_INTERACTING_RANK_COMPARISON/Opacity/Enigma.lean\#L71-L89}{\texttt{bills}}).
\item
  (c) \(I\) and \(\mathrm{Sh}\) are \(e^{-dR(x)}\)-small
  (\href{https://github.com/yang-mills-group/OpenAI-YM-20261028/blob/3456e5e75cf8d4cd558955a5c7638fba2d68f464/YM/PhysicalNext/AFTERNOON_INTERACTING_RANK_COMPARISON/Opacity/Shards.lean\#L60-L77}{\texttt{billIt}}).
\item
  (d) The chart terms \(\sum_e\log\tilde J(a_e)\) are
  \(x^{-1/8}\)-small, since \(\lvert\log\tilde J\rvert\le u(4d)^2\),
  \(d\le h t_L(x)\) and \(t_L(x)^2\le x^{-1/4}\) eventually. The linear
  term \(\eta\cdot z\) has the same polymer bound: the noise assumptions
  give \(|\eta_i z_i|\le x^{-2}\) on the cutoff support
  (\href{https://github.com/yang-mills-group/OpenAI-YM-20261028/blob/3456e5e75cf8d4cd558955a5c7638fba2d68f464/YM/PhysicalNext/AFTERNOON_INTERACTING_RANK_COMPARISON/Opacity/Sage.lean\#L82-L114}{\texttt{bolsters}}).
\item
  (e) The remaining path and interaction correction term (path and
  cubic/quartic corrections) is \(x^{-1/24}\)-small
  (\href{https://github.com/yang-mills-group/OpenAI-YM-20261028/blob/3456e5e75cf8d4cd558955a5c7638fba2d68f464/YM/PhysicalNext/AFTERNOON_INTERACTING_RANK_COMPARISON/Opacity/Regale.lean\#L65-L88}{\texttt{Ancestor.abound}}).
\end{itemize}

\emph{(proof read for (a), (b), (d); statement and outline for (c),
(e))}

\textbf{Lemma 4.20 (proof of III.4.7).}

\begin{enumerate}
\def\labelenumi{\arabic{enumi}.}
\tightlist
\item
  On flag-matched sectors,
  \(\Lambda=\sum_o(\mathrm{center}_o+\mathrm{rent}_o)+\mathrm{tenders}\)
  (\href{https://github.com/yang-mills-group/OpenAI-YM-20261028/blob/3456e5e75cf8d4cd558955a5c7638fba2d68f464/YM/PhysicalNext/AFTERNOON_INTERACTING_RANK_COMPARISON/Opacity/Gush.lean\#L27-L35}{\texttt{issued}}).
  The center is \(-\beta\) times the row energy of III.4.6 summed over
  row positions within \(200R(y)\) of \(o\)
  (\href{https://github.com/yang-mills-group/OpenAI-YM-20261028/blob/3456e5e75cf8d4cd558955a5c7638fba2d68f464/YM/PhysicalNext/AFTERNOON_INTERACTING_RANK_COMPARISON/Opacity/Turret.lean\#L74}{\texttt{center}}).
  The remainder consists of the terms of Lemma 4.19 (b)--(e), the row
  energy outside these halos, and the quadratic window correction.
\item
  On the nonzero cutoff support, the packet-energy hypothesis gives
  \(\mathrm{center}_o\le-\beta a s^2\lvert o\rvert\) for each core of
  the store, with \(s=z^{-1/2}\mathrm{hot}(z)\). Since
  \(zs^2=\mathrm{hot}(z)^2\) and \(\beta\ge z/2\), the condition
  \(2\alpha\le a/2\) gives \(2\alpha\,\mathrm{hot}(z)^2\le\beta as^2\).
  The bound is
  \href{https://github.com/yang-mills-group/OpenAI-YM-20261028/blob/3456e5e75cf8d4cd558955a5c7638fba2d68f464/YM/PhysicalNext/AFTERNOON_INTERACTING_RANK_COMPARISON/Opacity/Spreaders.lean\#L53-L61}{\texttt{ranked}}.
  Apply Lemma 4.19 (a) with \(\delta=\alpha\) to obtain
  \(\mathrm{center}_o+\mathrm{rent}_o\le-\alpha\,\mathrm{hot}(z)^2\lvert o\rvert\).
\item
  \href{https://github.com/yang-mills-group/OpenAI-YM-20261028/blob/3456e5e75cf8d4cd558955a5c7638fba2d68f464/YM/PhysicalNext/AFTERNOON_INTERACTING_RANK_COMPARISON/Opacity/Gush.lean\#L52-L65}{\texttt{paltry}}
  uses the lower bound
  \href{https://github.com/yang-mills-group/OpenAI-YM-20261028/blob/3456e5e75cf8d4cd558955a5c7638fba2d68f464/YM/PhysicalNext/AFTERNOON_INTERACTING_RANK_COMPARISON/Opacity/Spreaders.lean\#L82-L89}{\texttt{minimum}}
  to clamp the core term between fixed bounds, preserving its value on
  the nonzero cutoff support and its upper bound on all matching
  sectors.
\item
  The tenders are \(\ell\)-small with \(\ell=4q+y^4\nu+2\nu+a'\), where
  \(q=y^{-1/8}\), \(\nu=y^{-5}\) (from \(e^{-cR(y)}\le y^{-5}\)) and
  \(a'=y^{-1/24}\). Then \(2\ell\le16y^{-1/24}\le y^{-1/48}\) for large
  \(y\) (the power comparison).
\item
  \href{https://github.com/yang-mills-group/OpenAI-YM-20261028/blob/3456e5e75cf8d4cd558955a5c7638fba2d68f464/YM/PhysicalNext/AFTERNOON_INTERACTING_RANK_COMPARISON/Opacity/Petty.lean\#L72-L90}{\texttt{issue\_invoice}}
  symmetrises over lattice motions and converts the small sum into
  \(\lvert e^{\mathrm{aux}_X}-1\rvert\le f_X\).
\end{enumerate}

\textbf{Lemma 4.21 (precision, counting and jet certificates).} Fix
\(G\), a patch and a teller satisfying the branch condition, a halo
multiplier \(l\ge1\) and a polymer exponent \(e\ge0\).

\begin{itemize}
\tightlist
\item
  \href{https://github.com/yang-mills-group/OpenAI-YM-20261028/blob/3456e5e75cf8d4cd558955a5c7638fba2d68f464/YM/AnalyticBridge/RELATIVE_LOGDET_TRACE/Year.lean\#L113-L197}{\texttt{late\_price}}
  gives the precision and path-counting certificates, eventually in
  \(x\):

  \begin{itemize}
  \tightlist
  \item
    positive constants \(v,\eta,b,d\) and \(C\ge1\), chosen before the
    threshold in \(x\), giving the gap, positivity and resolvent bounds
    of III.4.5 for every funded step with that patch and
    \(c\ge100(1+R(x))\), uniformly in matching sectors and every window
    system containing the \(r_x\)-balls, with error factor
    \(e^{-bR(x)}\);
  \item
    cluster counting with factor \(e^{-dR(x)}\).
  \end{itemize}
\item
  For each pair of nonnegative integer orders \((n,k)\) and \(l\ge20\),
  \href{https://github.com/yang-mills-group/OpenAI-YM-20261028/blob/3456e5e75cf8d4cd558955a5c7638fba2d68f464/YM/PhysicalNext/PHYSICAL_JOINT_ENDPOINT_LOAD_F16/DLog/Fund.lean\#L120-L127}{\texttt{stores}}
  chooses a fund whose covariance-jet and determinant-letter bounds hold
  eventually in \(x\). The certificate also includes stencil-span and
  exponential cost bounds. These two results supply the corresponding
  certificates, rather than all the hypotheses of III.4.7.
\end{itemize}

\paragraph{4.4 Gaussian moments and
insertions}\label{gaussian-moments-and-insertions}

\textbf{Lemma 4.22 (exponentials of quadratic forms).} Let
\(x\sim N(0,I)\) and \(\mathcal Q=\sum_iw_i\langle q_i,x\rangle^2\) with
\(w_i\ge0\) and \(\sum_jw_j\lvert\langle q_i,q_j\rangle\rvert\le1/8\)
for each \(i\). Then
\(\mathbb E e^{\mathcal Q}\le e^{2\sum_iw_i\lvert q_i\rvert^2}\).
\emph{Sketch:} Gaussian integration by parts gives
\(F'(t)=\mathbb E[\mathcal Qe^{t\mathcal Q}]\le TF+\tfrac12F'\), so
\((e^{-2Tt}F)'\le0\).

\emph{Lean:}
\href{https://github.com/yang-mills-group/OpenAI-YM-20261028/blob/3456e5e75cf8d4cd558955a5c7638fba2d68f464/YM/Primitives/CONNECTED_GAUSSIAN_CUMULANTS/Trace.lean\#L228-L253}{\texttt{authentic}},
\href{https://github.com/yang-mills-group/OpenAI-YM-20261028/blob/3456e5e75cf8d4cd558955a5c7638fba2d68f464/YM/Primitives/CONNECTED_GAUSSIAN_CUMULANTS/Moment.lean\#L67-L71}{\texttt{amount}}.

\textbf{Lemma 4.23 (polynomial moments).} For \(L>0\) and
\(k\in\mathbb N\) there is \(D(L,k)>0\), independent of \(n\), such that
every \(0\preceq C\le LI\) on \(\mathbb R^n\) satisfies integrability
and \(\mathbb E_{\gamma_C}(1+\lvert z\rvert^2)^k\le D(L,k)(1+n)^{k+1}\).
\emph{Sketch:} dominate a power by an exponential with \(t=1/(4L)\),
bound the sum of squared coordinates by \(n\) times its largest term,
and use \(\mathbb Ee^{tz_x^2}\le2\).

\emph{Lean:}
\href{https://github.com/yang-mills-group/OpenAI-YM-20261028/blob/3456e5e75cf8d4cd558955a5c7638fba2d68f464/YM/AnalyticBridge/ANNULUS_RATIO_INSERTIONS/Buoy.lean\#L104-L152}{\texttt{many}}.

\textbf{Lemma 4.24 (proof of III.4.8).}

\begin{enumerate}
\def\labelenumi{\arabic{enumi}.}
\tightlist
\item
  On \(\{v\ne0\}\),
  \(1\le e^{-c\lvert J\rvert\rho^2}e^{c\sum_J\lvert z_i\rvert^2}\).
\item
  The site budget gives
  \(\sum d_i(1+\lvert z_{a_i}\rvert^2+\lvert z_{b_i}\rvert^2)\le\lvert Y\rvert c+2c\sum_X\lvert z\rvert^2\).
\item
  Use \(ab\le a^2+b^2\) to bound the product of the polynomial score and
  \(e^{3c\sum_U|z|^2}\). The exponential square has expectation at most
  \(2^{\dim\mathfrak g|U|}\), since \(6c=1/(4L)\) and
  \(|U|\le|J|+2|Y|\); Lemma 4.23 bounds the squared polynomial score.
\end{enumerate}

For the fourth-moment bound:
\(\lvert e^q-1\rvert\le xe^{2x(1+\lvert z_a\rvert^2+\lvert z_b\rvert^2)}\),
then apply III.4.8 with \(k=0\). For the coupling-band estimate: the
mean-shift equation writes \(\eta\) through
\(W^{\mathrm{per}}-W^{-1}=O(e^{-uR(x)/60})\), and \(C^{-1}-W\) is
controlled by the inverse-window estimate.

\emph{Lean:}
\href{https://github.com/yang-mills-group/OpenAI-YM-20261028/blob/3456e5e75cf8d4cd558955a5c7638fba2d68f464/YM/AnalyticBridge/ANNULUS_RATIO_INSERTIONS/Arc.lean\#L189-L253}{\texttt{same\_law}},
\href{https://github.com/yang-mills-group/OpenAI-YM-20261028/blob/3456e5e75cf8d4cd558955a5c7638fba2d68f464/YM/AnalyticBridge/ANNULUS_RATIO_INSERTIONS/Ember.lean\#L69-L111}{\texttt{r4}},
\href{https://github.com/yang-mills-group/OpenAI-YM-20261028/blob/3456e5e75cf8d4cd558955a5c7638fba2d68f464/YM/AnalyticBridge/ANNULUS_RATIO_INSERTIONS/Echo.lean\#L191-L223}{\texttt{oasis}}.

\textbf{Lemma 4.25 (Gaussian likelihood ratios).}

\begin{itemize}
\tightlist
\item
  For \(A,D\succ0\), any centres \(a,d\) and any real \(p\), the
  \(p\)-th power of \(dN(a,A)/dN(d,D)\) is \(N(d,D)\)-integrable iff
  \(pA^{-1}+(1-p)D^{-1}\succ0\)
  (\href{https://github.com/yang-mills-group/OpenAI-YM-20261028/blob/3456e5e75cf8d4cd558955a5c7638fba2d68f464/YM/SectorTools/PHYSICAL_GAUSSIAN_INSERTION/Moment.lean\#L183-L185}{\texttt{power\_iff}}).
\item
  For \(A,D\succ0\), if \(D^{-1}-A^{-1}\) is supported in \(S\times S\)
  and \(A^{-1}(d-a)\) on \(S\), the ratio equals the ratio of the
  \(S\)-marginals at \(x_S\)
  (\href{https://github.com/yang-mills-group/OpenAI-YM-20261028/blob/3456e5e75cf8d4cd558955a5c7638fba2d68f464/YM/SectorTools/PHYSICAL_GAUSSIAN_INSERTION/Support.lean\#L178-L207}{\texttt{likelihood\_local}}).
\item
  Let \(A,D\succ0\), \(p\ge1\), \(K,w_i\ge0\) and \(D\le KI\). If the
  change in the quadratic exponent about the reference centre has
  absolute value at most \(b+\sum_iw_iz_i^2\), and \(Kpw_i\le1/4\) for
  every \(i\), then the \(p\)-th likelihood-ratio moment is at most
  \(e^{2p(b+2K\sum_iw_i)}\). Cauchy--Schwarz controls the normaliser
  (\href{https://github.com/yang-mills-group/OpenAI-YM-20261028/blob/3456e5e75cf8d4cd558955a5c7638fba2d68f464/YM/SectorTools/PHYSICAL_GAUSSIAN_INSERTION/Weighted.lean\#L104-L138}{\texttt{weighted}}).
\item
  Let \(0\prec C\le KI\) with \(K>0\), \(\mathcal J\) symmetric,
  \(\varepsilon\ge0\), \(p\ge1\) and \(\theta>0\). For
  \(\lvert z\cdot\mathcal Jz\rvert\le\varepsilon\lvert z\rvert^2\),
  \(\lvert t_i\rvert\le\ell\) on \(S\) and
  \(Kp(\theta+\varepsilon/2)\le1/4\), the inserted covariance is
  positive definite and the moment bound of III.4.9 holds with \(U=S\).
  Positivity follows from \(C^{-1}\ge K^{-1}I>\varepsilon I\); the
  linear part contributes \(e^{p\lvert S\rvert\ell^2/(2\theta)}\) by
  AM--GM
  (\href{https://github.com/yang-mills-group/OpenAI-YM-20261028/blob/3456e5e75cf8d4cd558955a5c7638fba2d68f464/YM/SectorTools/PHYSICAL_GAUSSIAN_INSERTION/Insertion.lean\#L120-L172}{\texttt{insert\_letter}}).
\item
  For symmetric \(A\), integrability of
  \(e^{-\frac12z\cdot Az-b\cdot z+c}\) forces \(A\succ0\)
  (\href{https://github.com/yang-mills-group/OpenAI-YM-20261028/blob/3456e5e75cf8d4cd558955a5c7638fba2d68f464/YM/SectorTools/PHYSICAL_GAUSSIAN_INSERTION/Escape.lean\#L52-L94}{\texttt{posDef\_of\_integrable\_quadratic}}).
\item
  Let \(A,D\succ0\), \(\|A^{-1}\|_s\le M\),
  \(\|A^{-1}-D^{-1}\|_s\le\delta\), \(s\ge0\), and
  \(|\log\det A-\log\det D|\le n\delta\). The ratio \(T(z)\) of the
  \(A\)-Gaussian density at \(z+hu\) to the \(D\)-density at \(z\) obeys
  \[|\log T(z)|\le\tfrac\delta2(n+|z|^2)+\tfrac M2\big(h^2|u|^2+|h|(|u|^2+|z|^2)\big).\]
  Expand the shifted quadratic form and apply the Schur bounds
  (\href{https://github.com/yang-mills-group/OpenAI-YM-20261028/blob/3456e5e75cf8d4cd558955a5c7638fba2d68f464/YM/AnalyticBridge/GAUSSIAN_FIELD_ALIGNMENT/Relative.lean\#L245-L266}{\texttt{priceShift}}).
\end{itemize}

\emph{(proof read for the weighted moment and insertion bounds;
statement only for the other assertions)}

\paragraph{4.5 Jets}\label{jets}

\textbf{Lemma 4.26 (analytic predictor rows).} Fix \(G\), a patch \(P\),
a kernel \(K\) satisfying the branch condition and an order \(k\). There
are \(C\ge0\) and \(r,\delta>0\) such that the following holds. Take a
step with that patch and kernel, matching flags, stencil span at most
\(R(x)/50\), \(\mathrm{allot}(P)+1\le9R(x)\),
\(2(\mathrm{allot}(P)+1)<c\), and \(\mathrm{permit}(P)t(x)\le\delta\).
On every selected block stencil, the predicted group element and its
relative curvature are real-analytic on the \(r\)-ball of mixed coarse
and fine logarithmic perturbations; the curvature has \(k\)-jets bounded
by \(C\). Here \(\mathrm{allot}\) and \(\mathrm{permit}\) are fixed
patch-dependent geometric constants. \emph{Sketch:}

\begin{itemize}
\tightlist
\item
  Transplant the stencil into a fixed reference torus
  (\href{https://github.com/yang-mills-group/OpenAI-YM-20261028/blob/3456e5e75cf8d4cd558955a5c7638fba2d68f464/YM/AnalyticBridge/SECTOR_COVARIANCE_JETS/Press.lean\#L123-L155}{\texttt{impression}}).
\item
  Use a local trimming gauge in which the links have chord
  \(\le\delta\).
\item
  Write the prediction as \(u\cdot\mathrm{paste}(\ldots)\cdot v^{-1}\)
  with maps of norm \(\le1\)
  (\href{https://github.com/yang-mills-group/OpenAI-YM-20261028/blob/3456e5e75cf8d4cd558955a5c7638fba2d68f464/YM/AnalyticBridge/SECTOR_COVARIANCE_JETS/Duplicated.lean\#L52-L101}{\texttt{freshQ}}).
\item
  For large \(x\), all gauge-tag choices give the same predictor in a
  neighbourhood of each \((B,0)\) with \(\|B\|\le r\) on the selected
  stencils; the common predictor has uniform jets.
\end{itemize}

\emph{Lean:}
\href{https://github.com/yang-mills-group/OpenAI-YM-20261028/blob/3456e5e75cf8d4cd558955a5c7638fba2d68f464/YM/AnalyticBridge/SECTOR_COVARIANCE_JETS/Sounds.lean\#L69-L101}{\texttt{sounds}},
\href{https://github.com/yang-mills-group/OpenAI-YM-20261028/blob/3456e5e75cf8d4cd558955a5c7638fba2d68f464/YM/AnalyticBridge/SECTOR_AFFINE_JETS/Branch.lean\#L140-L219}{\texttt{chosen\_branch}}.

\textbf{Lemma 4.27 (precision near the background).} Under the fixed
patch, branch, funding, volume and flag hypotheses of III.4.10,
constants \(r,s,C_*,\eta>0\) precede the threshold in \(x\). On
\(\lVert B\rVert\le r\):

\begin{itemize}
\tightlist
\item
  \(M\) and \(W\) are real-analytic, with weighted jets \(\le C_*\);
\item
  the gap at \(B=0\) comes from Lemma 4.9;
\item
  a Lipschitz perturbation keeps half of it.
\end{itemize}

For \(d\ge0\) with \((1+2k)d\le\text{width}/2\), the move-matrix entries
have marked \(k\)-jets with the corresponding off-diagonal exponential
weight. \emph{(proof read for the precision bound; statement only for
the entry bound)}

\emph{Lean:}
\href{https://github.com/yang-mills-group/OpenAI-YM-20261028/blob/3456e5e75cf8d4cd558955a5c7638fba2d68f464/YM/AnalyticBridge/SECTOR_AFFINE_JETS/Neighborhood.lean\#L122-L159}{\texttt{precision}},
\href{https://github.com/yang-mills-group/OpenAI-YM-20261028/blob/3456e5e75cf8d4cd558955a5c7638fba2d68f464/YM/AnalyticBridge/SECTOR_COVARIANCE_JETS/Quote.lean\#L99-L127}{\texttt{entries}}.

\textbf{Lemma 4.28 (transport along walks).} Fix the jet order and a
bound on the marked jets of the link field. For each \(\kappa>0\) and
endpoint jet cap \(M>0\), there are a positive mark radius and \(N>0\),
uniform in the lattice, walk and scaling parameter \(s\in[0,1]\), such
that adjoint parallel transport moves a marked jet bound from the end of
a walk of length \(\ell\) to its start, at base point \(0\), with cap
multiplied by \(Ne^{\kappa\ell}\)
(\href{https://github.com/yang-mills-group/OpenAI-YM-20261028/blob/3456e5e75cf8d4cd558955a5c7638fba2d68f464/YM/AnalyticBridge/SECTOR_AFFINE_JETS/Visits.lean\#L81-L120}{\texttt{carried}}).
For a strip-decaying kernel of rate \(\gamma>0\), walks with length at
most \(B(1+|z|_1)\), \(B>0\), finite sets of walks below every length
cutoff, and the required endpoint bounds, choosing \(\kappa=\gamma/B\)
gives a uniform marked bound on the truncated sum
(\href{https://github.com/yang-mills-group/OpenAI-YM-20261028/blob/3456e5e75cf8d4cd558955a5c7638fba2d68f464/YM/AnalyticBridge/SECTOR_AFFINE_JETS/Tracks.lean\#L36-L89}{\texttt{follows}}).
If the path family has strip bound \((\sigma,M)\) with \(M\ge0\),
\(d\ge0\), \(\gamma>0\) and \(2\gamma+2kd\le\sigma/2\), selecting only
paths of length at least \(R\) gives marked \(k\)-jets bounded by a
constant times \(Me^{-\sigma R/4}\)
(\href{https://github.com/yang-mills-group/OpenAI-YM-20261028/blob/3456e5e75cf8d4cd558955a5c7638fba2d68f464/YM/AnalyticBridge/SELECTED_ORDERED_KERNELS/TrueTransport.lean\#L257-L299}{\texttt{transport\_rows}}).

\textbf{Lemma 4.29 (cut and window errors).} Assume fine path bounds
with positive rate \(\sigma\) and nonnegative amplitude, and the branch
condition. Near the flat background \(B=0\), the actual sector Gram
matrix equals the reindexed path-cut affine block. If no free links are
frozen, the tree-coordinate identification also identifies their
geometric windowed resolvents
(\href{https://github.com/yang-mills-group/OpenAI-YM-20261028/blob/3456e5e75cf8d4cd558955a5c7638fba2d68f464/YM/AnalyticBridge/LOCAL_GAUSSIAN_LAWS/CutGerm.lean\#L110-L129}{\texttt{notarized}}).
With a large step, \(R(x)\ge120\) and the marked gap and counting bounds
needed for the inverses, the following bounds hold at \(B=0\) for
\(u\ge0\):

\begin{itemize}
\tightlist
\item
  cutting the affine inverse germ at range \(R\) changes its \(k\)-jet
  by \(O(4^kC_0^2e^{-\sigma R/4})\), with \(C_0\propto(1+u)^{-1}\)
  (\href{https://github.com/yang-mills-group/OpenAI-YM-20261028/blob/3456e5e75cf8d4cd558955a5c7638fba2d68f464/YM/Push/FREE_RESOLVENT/AffineResolvent.lean\#L81-L100}{\texttt{shortening}});
\item
  the ordered windowed resolvent differs from the uncut one by a cap
  proportional to
  \((1+u)^{-2}(e^{-(k+1)\gamma r_x}+e^{-\sigma\lfloor R(x)/10^5\rfloor/4})\);
  symmetrisation costs at most a factor \(2\)
  (\href{https://github.com/yang-mills-group/OpenAI-YM-20261028/blob/3456e5e75cf8d4cd558955a5c7638fba2d68f464/YM/AnalyticBridge/LOCAL_GAUSSIAN_LAWS/Taylor.lean\#L73-L108}{\texttt{total}});
\item
  given the marked bounds and positive gaps for both \(W\) and \(C\),
  and the two successive admissible inverse-decay rates, \(C^{-1}-W\)
  has jets through order \(q\) bounded by a constant times
  \(e^{-(q+1)dr_x}\) in the reduced-rate marked norm
  (\href{https://github.com/yang-mills-group/OpenAI-YM-20261028/blob/3456e5e75cf8d4cd558955a5c7638fba2d68f464/YM/AnalyticBridge/SELECTED_ORDERED_KERNELS/SectorWindows.lean\#L139-L165}{\texttt{residual}}).
\end{itemize}

\paragraph{4.6 Large fields and
constants}\label{large-fields-and-constants}

\textbf{Lemma 4.30 (clean fields).}

\begin{itemize}
\tightlist
\item
  Under the funding, block-size, volume, tree-gauge and local sector
  hypotheses of III.4.11, there is \(C\ge1\), independent of the integer
  halo, with the following property. The field
  \(a_e=\mathrm{plog}(U_e\nu_e^{-1})\) vanishes on the tree links. On
  the swept annulus it gives exponential coordinates: within
  \(\lfloor2R(x)/5\rfloor\) of each swept site, \(|a_e|\) and the
  linking are at most \(C\,t_L(x)\). A coarse gauge and its pullback to
  the fine lattice make the corresponding coarse and fine chords at most
  \(C(1+R(x))^5t_L(x)\)
  (\href{https://github.com/yang-mills-group/OpenAI-YM-20261028/blob/3456e5e75cf8d4cd558955a5c7638fba2d68f464/YM/SectorNext/LARGE_FIELD_REMAINDERS/Clean.lean\#L110-L208}{\texttt{cleans\_at}}).
\item
  Given a clean-field constant \(X\ge1\), there are \(C>0\) and
  \(M\ge X\) such that, for large \(x\), the local plaquette curvatures
  and block rows have first-order expansions in \(a\), with linear term
  at most \(Cw\) and remainder at most \(Cw^2\), where
  \(w=M(1+R(x))^5t_L(x)\). Moreover \((1+R(x))^nw\to0\) for every
  \(n\in\mathbb N\); in terms of \(x\), \(w=x^{-1/2+o(1)}\)
  (\href{https://github.com/yang-mills-group/OpenAI-YM-20261028/blob/3456e5e75cf8d4cd558955a5c7638fba2d68f464/YM/SectorNext/LARGE_FIELD_REMAINDERS/CollarData.lean\#L131-L189}{\texttt{reacts}}).
\end{itemize}

\emph{(proof read for the response estimate; statement only for the
clean-field construction)}

\textbf{Lemma 4.31 (affine squares along paths).} Assume fine path
bounds with amplitude \(C\) and cutoff direct coefficient \(d_0\). Fix a
path cutoff \(R_0\in\mathbb N\), thresholds \(t,s\), \(B>0\) and
\(h\ge0\). Suppose the face under consideration has toll at most \(t\)
and satisfies the path-retention condition, its two coarse path rows
have norm at most \(Bh\), and its first-order curvature and curvature
error are at most \(Bh\) and \(Bh^2\). Require the same two curvature
bounds at every face reached by a retained path of length at most
\(R_0\), with chord deviation from the reference field at most \(h\)
along every such path. Then the affine path series converges and work
minus the extracted square is at least \(-4(mB)^2(1+2c_GR_0)h^3\), where
\(m=\sqrt{d_0}+C+1\)
(\href{https://github.com/yang-mills-group/OpenAI-YM-20261028/blob/3456e5e75cf8d4cd558955a5c7638fba2d68f464/YM/SectorNext/LARGE_FIELD_REMAINDERS/AffineRows.lean\#L157-L209}{\texttt{fine\_squares}}).
Under the clean-field and annulus hypotheses of Lemma 4.30, each
fine-face, noise and coarse row at a swept site has error at least
\(-e(x)\), where \(e\) is eventually nonnegative and
\((1+R(x))^n\beta_F(x,L)e(x)\to0\) for every \(n\)
(\href{https://github.com/yang-mills-group/OpenAI-YM-20261028/blob/3456e5e75cf8d4cd558955a5c7638fba2d68f464/YM/SectorNext/LARGE_FIELD_REMAINDERS/Recovery.lean\#L68-L219}{\texttt{coat}}).

\textbf{Lemma 4.32 (proof of III.4.11).}

\begin{enumerate}
\def\labelenumi{\arabic{enumi}.}
\tightlist
\item
  Sum the row bounds of Lemma 4.31 over \(Q\). The counted face, noise
  and coarse rows number at most \(X(1+R(x))^4|O|\), so the total error
  is \(e'(x)|O|\), where \(e'(x)=X(1+R(x))^4e(x)\) still satisfies
  \((1+R(x))^n\beta_F(x,L)e'(x)\to0\) for every \(n\).
\item
  Retain the affine squares selected by \(H\) and discard the other
  squares by positivity. The additional coarse rows in the work are
  nonnegative.
\item
  The extracted squares are unique, because the chart logarithm is
  unique once \(\max(K,C)t_L(x)\) is below the chart radius.
\end{enumerate}

\emph{Lean:}
\href{https://github.com/yang-mills-group/OpenAI-YM-20261028/blob/3456e5e75cf8d4cd558955a5c7638fba2d68f464/YM/SectorNext/LARGE_FIELD_REMAINDERS/Annulus.lean\#L36-L127}{\texttt{annulus\_pay}}.

\textbf{Lemma 4.33 (proof of III.4.12).}

\begin{enumerate}
\def\labelenumi{\arabic{enumi}.}
\tightlist
\item
  Choose the free-field bounds.
\item
  Choose the stationary width caps and path amplitudes, build a sheet
  within them, then choose drift bounds, slope caps and block-size
  floors. The stationary theorem gives a Hessian gap and full stationary
  path estimates only after the free-field, retained-family, paid-run,
  chart and pre-\(L\) compatibility hypotheses have been supplied. Its
  construction combines the one-loop analytic estimate with downward
  induction over levels.
\item
  Build the startup and free-field data for that sheet, preserving the
  required orders and bounds.
\item
  For each \(L\)-sequence the delays are enlarged and the
  chart-threshold enlargement is applied.
\item
  Assemble the stationary path families and their estimates.
\end{enumerate}

Related constant-producing statements:

\begin{itemize}
\tightlist
\item
  The scheduling theorems produce funded delays under their
  transition-rate and compatibility hypotheses;
\item
  The Gaussian comparison theorem gives, for each fixed real moment
  parameter and patch satisfying the branch condition, uniform reverse
  column-switch bounds on sufficiently large funded steps with matching
  flags and the required size condition. It also bounds the source
  \(M^{\top}s(Y)\) by \(C(t(x)+d)\) for every \(d>0\);
\item
  The counting theorem gives a polymer budget \(e^{-cR(x)}\) for fixed
  patch, nonnegative polymer exponent and prescribed integer orders;
  \(c>0\) is fixed by the input bounds.
\end{itemize}

\emph{(proof read for the main constant selection and stationary
theorem; statement only for the ancillary scheduling and comparison
results)}

\emph{Lean:}
\href{https://github.com/yang-mills-group/OpenAI-YM-20261028/blob/3456e5e75cf8d4cd558955a5c7638fba2d68f464/YM/MergeProof/Seaside.lean\#L96-L154}{\texttt{sailsV}},
\href{https://github.com/yang-mills-group/OpenAI-YM-20261028/blob/3456e5e75cf8d4cd558955a5c7638fba2d68f464/YM/Next/STATIONARY/FullPlanned.lean\#L59-L66}{\texttt{geared}},
\href{https://github.com/yang-mills-group/OpenAI-YM-20261028/blob/3456e5e75cf8d4cd558955a5c7638fba2d68f464/YM/PhysicalNext/V9_JOINT_ENDPOINT_LOAD/Calendar.lean\#L133-L138}{\texttt{book}},
\href{https://github.com/yang-mills-group/OpenAI-YM-20261028/blob/3456e5e75cf8d4cd558955a5c7638fba2d68f464/YM/PhysicalNext/RG_NEXT_COMPLETE_REPEATABLE_NONLINEAR_STEP/Instruments.lean\#L69-L86}{\texttt{endures}},
\href{https://github.com/yang-mills-group/OpenAI-YM-20261028/blob/3456e5e75cf8d4cd558955a5c7638fba2d68f464/YM/PhysicalNext/V9_JOINT_ENDPOINT_LOAD/Reverse.lean\#L64-L88}{\texttt{post}},
\href{https://github.com/yang-mills-group/OpenAI-YM-20261028/blob/3456e5e75cf8d4cd558955a5c7638fba2d68f464/YM/PhysicalNext/V9_JOINT_ENDPOINT_LOAD/Folio/Takings.lean\#L99-L116}{\texttt{takings}}.

\subsubsection{5. Where the difficulty
lies}\label{where-the-difficulty-lies-6}

The gap is a gauge-theoretic local Poincaré estimate: the predictor
derivative must be compared with its flat-field value after a different
gauge choice for each active link. Its smallness radius comes from
continuity rather than an explicit formula. Inverse estimates then spend
decay rate, while every bound must remain uniform in the volume, store,
forest, tags and field.

The determinant comparison must accommodate different active dimensions,
hence identity completion and the frozen-link normalisation. The
resulting core costs grow polynomially in \(\log x\); the retimed
large-field threshold supplies the stronger suppression
\(e^{-\alpha\mathrm{hot}(z)^2|o|}\). The final powers \(x^{-1/8}\),
\(x^{-1/24}\) and \(y^{-1/48}\) allow the separate polymer budgets to be
added. Gaussian transport and the affine-square identity are
finite-dimensional; their application still requires the chart support
and integrability conditions.

\subsubsection{6. Notation}\label{notation-6}

{\def\LTcaptype{none} % do not increment counter
\begin{longtable}[]{@{}
  >{\raggedright\arraybackslash}p{(\linewidth - 4\tabcolsep) * \real{0.3333}}
  >{\raggedright\arraybackslash}p{(\linewidth - 4\tabcolsep) * \real{0.3333}}
  >{\raggedright\arraybackslash}p{(\linewidth - 4\tabcolsep) * \real{0.3333}}@{}}
\toprule\noalign{}
\begin{minipage}[b]{\linewidth}\raggedright
Written proof
\end{minipage} & \begin{minipage}[b]{\linewidth}\raggedright
Lean (glossary)
\end{minipage} & \begin{minipage}[b]{\linewidth}\raggedright
Meaning
\end{minipage} \\
\midrule\noalign{}
\endhead
\bottomrule\noalign{}
\endlastfoot
\(\gamma_C\) &
\href{https://github.com/yang-mills-group/OpenAI-YM-20261028/blob/3456e5e75cf8d4cd558955a5c7638fba2d68f464/YM/SectorNext/GAUSSIAN_SECTOR_CHANGE/Gaussian.lean\#L31}{\texttt{Gaussian.cloud}}
& centred Gaussian with covariance \(C\) \\
\(z\cdot C^{-1}z\) &
\href{https://github.com/yang-mills-group/OpenAI-YM-20261028/blob/3456e5e75cf8d4cd558955a5c7638fba2d68f464/YM/SectorNext/GAUSSIAN_SECTOR_CHANGE/Gaussian.lean\#L79}{\texttt{force}}
& precision quadratic form \\
\((2\pi g^2)^{n_a/2}(\det C)^{1/2}\) &
\href{https://github.com/yang-mills-group/OpenAI-YM-20261028/blob/3456e5e75cf8d4cd558955a5c7638fba2d68f464/YM/SectorNext/GAUSSIAN_SECTOR_CHANGE/Counting.lean\#L103-L104}{\texttt{chartCost}}
& Gaussian normalisation \\
\(g=\beta^{-1/2}\) &
\href{https://github.com/yang-mills-group/OpenAI-YM-20261028/blob/3456e5e75cf8d4cd558955a5c7638fba2d68f464/YM/SectorNext/GAUSSIAN_SECTOR_CHANGE/Change.lean\#L35}{\texttt{coupling}}
& coupling of the representation \\
\(J\), \(J_{\rm sc}\), \(\rho_J\) &
\href{https://github.com/yang-mills-group/OpenAI-YM-20261028/blob/3456e5e75cf8d4cd558955a5c7638fba2d68f464/YM/PCalc/Loop.lean\#L23-L26}{\texttt{ChartRatio}},
\href{https://github.com/yang-mills-group/OpenAI-YM-20261028/blob/3456e5e75cf8d4cd558955a5c7638fba2d68f464/YM/PCalc/Loop.lean\#L28-L38}{\texttt{LocalRatio}}
& normalised Haar Jacobian of exp, scale, radius \\
\(\mathcal F_A\), \(\mathcal A_A\) &
\href{https://github.com/yang-mills-group/OpenAI-YM-20261028/blob/3456e5e75cf8d4cd558955a5c7638fba2d68f464/YM/SectorNext/GAUSSIAN_SECTOR_CHANGE/Shell.lean\#L38-L40}{\texttt{chill}}
(\texttt{Frozen}, \texttt{Active}) & free links within \(10R(x)\) of
\(A\) / the others \\
\(\iota V\) &
\href{https://github.com/yang-mills-group/OpenAI-YM-20261028/blob/3456e5e75cf8d4cd558955a5c7638fba2d68f464/YM/PCalc/Chart.lean\#L35-L36}{\texttt{replica}}
& coarse field copied to the fine lattice \\
\(Q_Y(a)\), \(W\) &
\href{https://github.com/yang-mills-group/OpenAI-YM-20261028/blob/3456e5e75cf8d4cd558955a5c7638fba2d68f464/YM/SectorNext/GAUSSIAN_SECTOR_CHANGE/Remainder.lean\#L37}{\texttt{Rows.guide}},
\href{https://github.com/yang-mills-group/OpenAI-YM-20261028/blob/3456e5e75cf8d4cd558955a5c7638fba2d68f464/YM/SectorNext/GAUSSIAN_SECTOR_CHANGE/Remainder.lean\#L27}{\texttt{Rows.gram}}
& affine squares; precision \(M^{\top}M\) \\
\(C=W_{\mathcal W}\), \(m(Y)\) &
\href{https://github.com/yang-mills-group/OpenAI-YM-20261028/blob/3456e5e75cf8d4cd558955a5c7638fba2d68f464/YM/SectorNext/GAUSSIAN_SECTOR_CHANGE/Chosen.lean\#L69-L70}{\texttt{fog}},
\href{https://github.com/yang-mills-group/OpenAI-YM-20261028/blob/3456e5e75cf8d4cd558955a5c7638fba2d68f464/YM/SectorNext/GAUSSIAN_SECTOR_CHANGE/Chosen.lean\#L72-L74}{\texttt{meanLog}}
& windowed covariance; windowed mean \\
\(L\), \(A_{\mathcal W}\) &
\href{https://github.com/yang-mills-group/OpenAI-YM-20261028/blob/3456e5e75cf8d4cd558955a5c7638fba2d68f464/YM/SectorNext/GAUSSIAN_SECTOR_CHANGE/Windows.lean\#L67-L68}{\texttt{Schur.per}},
\href{https://github.com/yang-mills-group/OpenAI-YM-20261028/blob/3456e5e75cf8d4cd558955a5c7638fba2d68f464/YM/SectorNext/GAUSSIAN_SECTOR_CHANGE/Windows.lean\#L70}{\texttt{Schur.clouded}}
& column-wise window inverse; its symmetrisation \\
Schur decay, gap &
\href{https://github.com/yang-mills-group/OpenAI-YM-20261028/blob/3456e5e75cf8d4cd558955a5c7638fba2d68f464/YM/SectorNext/GAUSSIAN_SECTOR_CHANGE/Schur.lean\#L25-L28}{\texttt{Pours}},
\href{https://github.com/yang-mills-group/OpenAI-YM-20261028/blob/3456e5e75cf8d4cd558955a5c7638fba2d68f464/YM/SectorNext/GAUSSIAN_SECTOR_CHANGE/Schur.lean\#L110}{\texttt{Gapped}}
& weighted row and column sums; \(a\cdot Aa\ge\eta\lVert a\rVert^2\) \\
\(\Psi\) &
\href{https://github.com/yang-mills-group/OpenAI-YM-20261028/blob/3456e5e75cf8d4cd558955a5c7638fba2d68f464/YM/SectorNext/GAUSSIAN_SECTOR_CHANGE/Change.lean\#L63-L71}{\texttt{dyed}}
& sector integrand in Gaussian coordinates \\
permitted, licence &
\href{https://github.com/yang-mills-group/OpenAI-YM-20261028/blob/3456e5e75cf8d4cd558955a5c7638fba2d68f464/YM/SectorNext/GAUSSIAN_SECTOR_CHANGE/Change.lean\#L95-L103}{\texttt{Permitted}},
\href{https://github.com/yang-mills-group/OpenAI-YM-20261028/blob/3456e5e75cf8d4cd558955a5c7638fba2d68f464/YM/SectorNext/GAUSSIAN_SECTOR_CHANGE/Change.lean\#L127-L135}{\texttt{Licenses}}
& hypotheses and conclusions of III.4.1 \\
\(\ell_i\), path letters, \(S\) &
\href{https://github.com/yang-mills-group/OpenAI-YM-20261028/blob/3456e5e75cf8d4cd558955a5c7638fba2d68f464/YM/AnalyticBridge/RELATIVE_LOGDET_TRACE/Series.lean\#L45-L46}{\texttt{letterBase}},
\texttt{page}, \texttt{errata},
\href{https://github.com/yang-mills-group/OpenAI-YM-20261028/blob/3456e5e75cf8d4cd558955a5c7638fba2d68f464/YM/AnalyticBridge/RELATIVE_LOGDET_TRACE/Walls.lean\#L35}{\texttt{excess}}
& window log letter; path terms; \(AA_{\mathcal W}-1\) \\
\(\kappa(p)\), \(\Lambda(p)\) &
\href{https://github.com/yang-mills-group/OpenAI-YM-20261028/blob/3456e5e75cf8d4cd558955a5c7638fba2d68f464/YM/PhysicalNext/AFTERNOON_INTERACTING_RANK_COMPARISON/Opacity/Prism.lean\#L50-L54}{\texttt{keptEarly}},
\href{https://github.com/yang-mills-group/OpenAI-YM-20261028/blob/3456e5e75cf8d4cd558955a5c7638fba2d68f464/YM/PhysicalNext/AFTERNOON_INTERACTING_RANK_COMPARISON/Opacity/Prism.lean\#L56-L63}{\texttt{spoolLift}}
& relative log-normalisation; full log-integrand \\
\(\mathrm{spare}(\beta)\) &
\href{https://github.com/yang-mills-group/OpenAI-YM-20261028/blob/3456e5e75cf8d4cd558955a5c7638fba2d68f464/YM/SectorNext/GAUSSIAN_SECTOR_CHANGE/Ledger.lean\#L34}{\texttt{spare}}
& \((2\pi g^2)^{\dim\mathfrak g/2}J_{\rm sc}\) \\
\(\mathrm{rent}_o\), far term, \(I\), \(\mathrm{Sh}\) &
\href{https://github.com/yang-mills-group/OpenAI-YM-20261028/blob/3456e5e75cf8d4cd558955a5c7638fba2d68f464/YM/PhysicalNext/AFTERNOON_INTERACTING_RANK_COMPARISON/Opacity/Enigma.lean\#L32}{\texttt{rents}},
\href{https://github.com/yang-mills-group/OpenAI-YM-20261028/blob/3456e5e75cf8d4cd558955a5c7638fba2d68f464/YM/PhysicalNext/AFTERNOON_INTERACTING_RANK_COMPARISON/Opacity/Enigma.lean\#L35}{\texttt{nightly}},
\href{https://github.com/yang-mills-group/OpenAI-YM-20261028/blob/3456e5e75cf8d4cd558955a5c7638fba2d68f464/YM/PhysicalNext/AFTERNOON_INTERACTING_RANK_COMPARISON/Opacity/Shards.lean\#L31}{\texttt{introTax}},
\href{https://github.com/yang-mills-group/OpenAI-YM-20261028/blob/3456e5e75cf8d4cd558955a5c7638fba2d68f464/YM/PhysicalNext/AFTERNOON_INTERACTING_RANK_COMPARISON/Opacity/Shards.lean\#L32}{\texttt{showTax}}
& pieces of \(\kappa\) \\
\(\varepsilon\)-small &
\href{https://github.com/yang-mills-group/OpenAI-YM-20261028/blob/3456e5e75cf8d4cd558955a5c7638fba2d68f464/YM/PhysicalNext/AFTERNOON_INTERACTING_RANK_COMPARISON/Opacity/Petty.lean\#L28-L29}{\texttt{Petty}}
& sum of local terms with polymer norm \(\le\varepsilon\) \\
\(y^{-1/48}\) &
\href{https://github.com/yang-mills-group/OpenAI-YM-20261028/blob/3456e5e75cf8d4cd558955a5c7638fba2d68f464/YM/PhysicalNext/AFTERNOON_INTERACTING_RANK_COMPARISON/Opacity/Stocks.lean\#L21}{\texttt{garlic}}
& final polymer budget \\
\(R(x)\), \(t(x)\), \(t_L(x)\) & \texttt{radiusB}, \texttt{temp},
\texttt{fiery} & cluster radius, small-field threshold, fine-level
threshold \\
\(M\), \(s(Y)\) & \texttt{move}, \texttt{source} & linear row map and
frozen-link source \\
\(\mathcal E\), \(\mathrm{center}_o\) & \texttt{Milk.levy},
\texttt{Coffer.center} & row energy remainder; its \(-\beta\)-weighted
core sum \\
sector support & \texttt{Dyed.Playing} & nonzero product of field,
block, tag and chart cutoffs \\
matching flags & \texttt{Matched} &
\(\mathrm{flags}(A)=\mathrm{flags}_{t(x)}(V)\) \\
branch condition & \texttt{Patch.Branches} & uniqueness of selected
minimisers and eligible axial gauge \\
sector regularity, noise bounds & \texttt{Chips}, \texttt{noises} &
assumptions in III.4.6--III.4.7 \\
\(c_G\) & \texttt{Gauge.colorScale} & matrix norm factor in the path
estimate \\
price, bill, tax, rent, teller & numerical bounds; \texttt{Teller} & the
step-estimator interface \\
\end{longtable}
}

\subsubsection{7. What remains
unexplained}\label{what-remains-unexplained-6}

\begin{enumerate}
\def\labelenumi{\arabic{enumi}.}
\tightlist
\item
  \textbf{Proof depth.} The supporting proofs marked ``statement only''
  remain incomplete readings: the strip-resolvent and Fourier-inverse
  results in Lemma 4.14; the flat comparison in 4.15; 4.16; the
  path-counting and interaction estimates in 4.19(c),(e); 4.21; parts of
  4.25; the entry bound in 4.27; 4.28--4.29; the clean-field
  construction in 4.30; 4.31; and the ancillary scheduling statements in
  4.33. Their statements and the stated mechanisms have been compared
  with the cited source; these labels do not certify their dependency
  proofs.
\item
  \textbf{Run-level hypotheses.} The packet-energy and support-exponent
  assumptions of III.4.7 are substantive inputs. Their simultaneous
  supply on the selected run, and the transfer from stationary path
  families to the tube's work rows, still need a complete account. The
  precision, determinant-jet and noise bounds have local suppliers, but
  their use does not establish these other premises.
\item
  \textbf{Free-kernel comparison.} Lemma 4.15 compares the flat sector
  under the no-frozen-link and free-metric compatibility conditions. A
  direct statement comparing an arbitrary interacting sector covariance
  with the infinite-volume free operator was not located. The full
  expansion of the free-metric compatibility conditions and the nested
  marked inverse bounds in Lemma 4.29 remains to be written.
\item
  \textbf{Consumers and coverage.} The downstream reset and
  repeatable-step statements have not been checked here. The additional
  assigned packages for Wilson bulk terms, vacuum stability, proof
  interfaces, multisource terms and flat RG data are not covered by
  these steps. The closure-size figures in the opening metadata were
  retained without recounting.
\item
  \textbf{References.} The precise results and bibliographic locations
  in Helgason, Glimm--Jaffe, Katznelson, Balaban and Dimock remain to be
  selected and verified. The cited Lean declarations are the evidence
  for the statements above.
\end{enumerate}

\subsection{III.5. Cluster expansion of the polymer
gas}\label{iii.5.-cluster-expansion-of-the-polymer-gas}

\textbf{Part:} III, Ultraviolet construction (renormalization group).
\textbf{Lean:} work packages
\href{https://github.com/yang-mills-group/OpenAI-YM-20261028/tree/3456e5e75cf8d4cd558955a5c7638fba2d68f464/YM/Push/ABSOLUTE_SUMS}{\texttt{YM/Push/ABSOLUTE\_SUMS/}}
(3,624 lines),
\href{https://github.com/yang-mills-group/OpenAI-YM-20261028/tree/3456e5e75cf8d4cd558955a5c7638fba2d68f464/YM/AnalyticNext/CONNECTED_ACTIVITY_BOUNDS}{\texttt{YM/AnalyticNext/CONNECTED\_ACTIVITY\_BOUNDS/}}
(3,589),
\href{https://github.com/yang-mills-group/OpenAI-YM-20261028/tree/3456e5e75cf8d4cd558955a5c7638fba2d68f464/YM/Push/SECTOR_MEASURE}{\texttt{YM/Push/SECTOR\_MEASURE/}}
(1,153),
\href{https://github.com/yang-mills-group/OpenAI-YM-20261028/tree/3456e5e75cf8d4cd558955a5c7638fba2d68f464/YM/SectorNext/VACUUM_NORMALIZATION}{\texttt{YM/SectorNext/VACUUM\_NORMALIZATION/}}
(935),
\href{https://github.com/yang-mills-group/OpenAI-YM-20261028/tree/3456e5e75cf8d4cd558955a5c7638fba2d68f464/YM/Progress/STATIONARY_SCENES}{\texttt{YM/Progress/STATIONARY\_SCENES/}}
(445). \textbf{Size:} 9,746 lines in 1,097 declarations used by the main
theorem.

\textbf{Conventions.}

\begin{itemize}
\tightlist
\item
  \(\Lambda_n=(\mathbb Z/n\mathbb Z)^4\), with \(n\ge1\), is the
  periodic lattice with the maximum distance \(d\). A \emph{support} is
  a subset of \(\Lambda_n\).
\item
  \(\mathrm{cov}_B(X)\) is the least number of \(d\)-balls of radius
  \(B\) whose union is connected and covers \(X\)
  (\href{https://github.com/yang-mills-group/OpenAI-YM-20261028/blob/3456e5e75cf8d4cd558955a5c7638fba2d68f464/YM/Lattice/Support.lean\#L69-L71}{\texttt{coverNumber}}).
  Two properties are used throughout (I.1):

  \begin{itemize}
  \tightlist
  \item
    \(|X|\le \mathrm{ball}_B\,\mathrm{cov}_B(X)\), with
    \(\mathrm{ball}_B=(2B+1)^4\);
  \item
    \(\mathrm{cov}_B\) is subadditive on unions of overlapping sets.
  \end{itemize}
\item
  A family of supports is \emph{connected} if it is nonempty and its
  overlap graph is connected.
\item
  For \(f\ge0\) on an arbitrary index set, \(\mathrm{SB}(f;C)\) means:
  \(f\) is summable and \(\sum f\le C\)
  (\href{https://github.com/yang-mills-group/OpenAI-YM-20261028/blob/3456e5e75cf8d4cd558955a5c7638fba2d68f464/YM/RG/Activities.lean\#L59-L60}{\texttt{SumBudget}}).
\item
  For polymers \((X_i,w_i)_{i\in\iota}\) with \(w\ge0\), the
  \emph{site-sum condition} \(\mathrm{Hit}_B(w;a,\gamma)\) means
  \(\mathrm{SB}\big(i\mapsto \mathbf 1[x\in X_i]\,w_i\,e^{a\,\mathrm{cov}_B(X_i)};\gamma\big)\)
  for every site \(x\)
  (\href{https://github.com/yang-mills-group/OpenAI-YM-20261028/blob/3456e5e75cf8d4cd558955a5c7638fba2d68f464/YM/Ac/Tree.lean\#L12-L13}{\texttt{Ac.Hit}},
  II.2).
\item
  All weights are real. An upper bound on an infinite sum of nonnegative
  terms includes summability, as in \(\mathrm{SB}\).
\item
  For a relation \(R\) that is not symmetric, the graph used below has
  an edge between distinct occurrences \(a,b\) exactly when both
  \(R(z_a,z_b)\) and \(R(z_b,z_a)\) hold. The coefficient of a
  one-vertex tuple is \(1\).
\end{itemize}

\subsubsection{1. What this unit proves}\label{what-this-unit-proves-7}

The unit consists of five packages.

\begin{itemize}
\tightlist
\item
  \textbf{(A) Mayer expansion.}

  \begin{itemize}
  \tightlist
  \item
    The Mayer expansion of a hard-core polymer gas, with absolute
    bounds, and its re-exponentiation after forced polymers
    (III.5.1--III.5.6).
  \item
    The verification of the smallness condition for gases of large-field
    seeds of a block-spin step (III.5.7--III.5.8).
  \end{itemize}
\item
  \textbf{(B) Vacuum normalization.} A second, graph-free construction
  of the logarithm by deletion identities (III.5.9), and a trace-log
  formula (III.5.10).
\item
  \textbf{(C) Sector decomposition.} The decomposition of one block-spin
  integral into sectors labelled by packets of seeds, with a Gaussian
  gain per seed (III.5.11).
\item
  \textbf{(D) Inverse decay.} A Combes--Thomas bound (III.5.12).
\item
  \textbf{(E) Gaussian bounds.} Windowed inverses (III.5.13), and a
  uniform spectral gap and uniform decay for the precision matrix of the
  small-field Gaussian of a step (III.5.14).
\end{itemize}

III.5.1--III.5.7, III.5.9, III.5.10, III.5.12 and III.5.13 are general
implications. III.5.8(c), III.5.11 and III.5.14 concern the block-spin
step under the funding hypotheses of III.1; III.5.14 also fixes a patch
satisfying its branch conditions.

\textbf{Proposition III.5.1 (Mayer identity for a finite gas).}

\begin{itemize}
\tightlist
\item
  \emph{Data.} Let \(\iota\) be finite, \(R\) a reflexive symmetric
  relation on \(\iota\) (``incompatible''), and \(w:\iota\to\mathbb R\).
  Let \(\Xi(w)=\sum_s\prod_{i\in s}w_i\), the sum over subsets
  \(s\subseteq\iota\) of pairwise compatible labels.
\item
  \emph{Ursell coefficient.} For \(z=(z_0,\dots,z_m)\in\iota^{m+1}\),
  let \(G_R(z)\) be the graph on \(\{0,\dots,m\}\) with an edge
  \(\{a,b\}\), \(a\neq b\), whenever \(R(z_a,z_b)\). Put
  \(\phi^T_R(z)=\sum_G(-1)^{|G|}\), the sum over connected spanning
  subgraphs \(G\subseteq G_R(z)\).
\item
  \emph{Claim.} Let
  \(\ell_m=\frac1{(m+1)!}\sum_{z\in\iota^{m+1}}\phi^T_R(z)\prod_a w_{z_a}\).
  If \(|\ell_m|\le Cq^m\) for all \(m\), with \(C\ge0\) and \(0<q<1\),
  then \(\exp\big(\sum_{m\ge0}\ell_m\big)=\Xi(w)\).
\end{itemize}

\emph{Status:} theorem. \emph{Lean:}
\href{https://github.com/yang-mills-group/OpenAI-YM-20261028/blob/3456e5e75cf8d4cd558955a5c7638fba2d68f464/YM/Push/ABSOLUTE_SUMS/FiniteLog.lean\#L66-L145}{\texttt{Gas.exact\_polynomial}};
coefficient identities
\href{https://github.com/yang-mills-group/OpenAI-YM-20261028/blob/3456e5e75cf8d4cd558955a5c7638fba2d68f464/YM/Push/ABSOLUTE_SUMS/Coef.lean\#L129-L138}{\texttt{Coef.recurrence}},
\href{https://github.com/yang-mills-group/OpenAI-YM-20261028/blob/3456e5e75cf8d4cd558955a5c7638fba2d68f464/YM/Push/ABSOLUTE_SUMS/Partition.lean\#L219-L237}{\texttt{Partition.partition}}.

\textbf{Proposition III.5.2 (tree-graph bound).} Let \(T\) be a graph on
a finite vertex set with a distinguished vertex. Then
\(\big|\sum_G(-1)^{|G|}\big|\), summed over the connected spanning
subgraphs \(G\subseteq T\), is at most the number of spanning trees of
\(T\). In particular \(|\phi^T_R(z)|\) is at most the number of spanning
trees of \(G_R(z)\). \emph{Status:} theorem. Spanning trees are counted
as acyclic parent maps towards the distinguished vertex. \emph{Lean:}
\href{https://github.com/yang-mills-group/OpenAI-YM-20261028/blob/3456e5e75cf8d4cd558955a5c7638fba2d68f464/YM/Push/ABSOLUTE_SUMS/Graph.lean\#L113-L148}{\texttt{Graph.pay}}.

\textbf{Proposition III.5.3 (pinned cluster bounds).}

\begin{itemize}
\item
  \emph{Data.} \(\iota\) is any set; \(X:\iota\to\) supports of
  \(\Lambda_n\); \(B\in\mathbb N\); \(w:\iota\to[0,\infty)\); \(a\ge0\);
  \(R\) is a relation on \(\iota\) with
  \(R(i,j)\Rightarrow X_i\cap X_j\neq\emptyset\). \(R\) need not be
  symmetric.
\item
  \emph{Smallness condition.}
  \[\mathrm{Ready}(w;a,\gamma,\delta,q):\qquad \gamma>0,\quad\delta>0,\quad 4\,\mathrm{ball}_B\,\gamma/\delta\le q<1,\quad \mathrm{Hit}_B(w;\,a+\delta,\,\gamma).\]
\item
  \emph{Notation.} \(X(z)=\bigcup_aX_{z_a}\) for a tuple \(z\).
\end{itemize}

Then, for every site \(x\) and every integer \(m\ge0\):

\begin{itemize}
\tightlist
\item
  \begin{enumerate}
  \def\labelenumi{(\alph{enumi})}
  \tightlist
  \item
    \(\displaystyle\sum_{m'\ge m}\ \sum_{z\in\iota^{m'+1}:\ x\in X(z)}\frac{|\phi^T_R(z)|}{(m'+1)!}\prod_a w_{z_a}\ e^{a\,\mathrm{cov}_B(X(z))}\ \le\ \frac{\gamma\,q^m}{1-q}\).
  \end{enumerate}
\item
  (b) The same bound holds for
  \(\sum_s\prod_{i\in s}\big(w_i e^{a\,\mathrm{cov}_B(X_i)}\big)\). The
  sum runs over finite sets \(s\subseteq\iota\) with \(|s|\ge m+1\)
  whose supports form a connected family containing \(x\).
\item
  (c) (Marked form.) Suppose \(Q\subseteq\iota\) and
  \(\mathrm{SB}\big(i\mapsto\mathbf1[i\in Q]w_ie^{(a+\delta)\mathrm{cov}_B(X_i)};b\big)\).
  Then the sum over \(z\in\iota^{m+1}\) with some \(z_a\in Q\) of
  \(|\phi^T_R(z)|\prod_{j=0}^m\big(w_{z_j}e^{a\,\mathrm{cov}_B(X_{z_j})}\big)/(m+1)!\)
  is at most \(bq^m\).
\end{itemize}

Each displayed bound is a summability statement of the form
\(\mathrm{SB}\). \emph{Status:} theorem. Form (a) is used in III.2,
III.10, III.11, III.12, III.13 and III.15. \emph{Lean:}
\href{https://github.com/yang-mills-group/OpenAI-YM-20261028/blob/3456e5e75cf8d4cd558955a5c7638fba2d68f464/YM/Push/ABSOLUTE_SUMS/Clusters.lean\#L224-L243}{\texttt{Ursell.log\_hit\_tail}},
\href{https://github.com/yang-mills-group/OpenAI-YM-20261028/blob/3456e5e75cf8d4cd558955a5c7638fba2d68f464/YM/Push/ABSOLUTE_SUMS/Distinct.lean\#L245-L276}{\texttt{Trees.touch\_tail}},
\href{https://github.com/yang-mills-group/OpenAI-YM-20261028/blob/3456e5e75cf8d4cd558955a5c7638fba2d68f464/YM/Push/ABSOLUTE_SUMS/Clusters.lean\#L148-L183}{\texttt{Ursell.cluster\_at}},
\href{https://github.com/yang-mills-group/OpenAI-YM-20261028/blob/3456e5e75cf8d4cd558955a5c7638fba2d68f464/YM/Push/ABSOLUTE_SUMS/Clusters.lean\#L185-L192}{\texttt{Ursell.cluster\_tail}};
hypothesis
\href{https://github.com/yang-mills-group/OpenAI-YM-20261028/blob/3456e5e75cf8d4cd558955a5c7638fba2d68f464/YM/Push/ABSOLUTE_SUMS/Forest.lean\#L31-L32}{\texttt{Trees.Ready}}.

\textbf{Proposition III.5.4 (convergence and logarithm of the gas).} Let
\(X,B\) be as in III.5.3, \(R\) reflexive and symmetric with
\(R(i,j)\Rightarrow X_i\cap X_j\neq\emptyset\), \(w:\iota\to\mathbb R\)
and \(a\ge0\). Assume \(\mathrm{Ready}(|w|;a,\gamma,\delta,q)\). Then:

\begin{itemize}
\tightlist
\item
  \(\Xi(w)=\sum_s\prod_{i\in s}w_i\), over finite pairwise compatible
  \(s\subseteq\iota\), converges absolutely, with
  \(\sum_s\prod_{i\in s}|w_i|\le e^{n^4\gamma}\).
\item
  The cluster series
  \(\sum_{m\ge0}\sum_{z\in\iota^{m+1}}\phi^T_R(z)\prod_aw_{z_a}/(m+1)!\)
  converges absolutely. Its absolute sum is at most \(n^4\gamma/(1-q)\),
  and at most \(n^4\gamma q^m\) over the clusters of \(m+1\) labels.
\item
  \(\Xi(w)>0\), and \(\log\Xi(w)\) equals the cluster series.
\end{itemize}

\emph{Status:} theorem. Used in III.10. \emph{Lean:}
\href{https://github.com/yang-mills-group/OpenAI-YM-20261028/blob/3456e5e75cf8d4cd558955a5c7638fba2d68f464/YM/Push/ABSOLUTE_SUMS/Log.lean\#L114-L149}{\texttt{Gas.price\_log}},
\href{https://github.com/yang-mills-group/OpenAI-YM-20261028/blob/3456e5e75cf8d4cd558955a5c7638fba2d68f464/YM/Push/ABSOLUTE_SUMS/FiniteLog.lean\#L53-L57}{\texttt{Gas.absolute}},
\href{https://github.com/yang-mills-group/OpenAI-YM-20261028/blob/3456e5e75cf8d4cd558955a5c7638fba2d68f464/YM/Push/ABSOLUTE_SUMS/FiniteLog.lean\#L35-L51}{\texttt{Gas.absolute\_at}}.

\textbf{Proposition III.5.5 (removing the polymers that meet a region).}
In III.5.4, take for \(R\) the overlap relation (\(R(i,j)\) iff \(i=j\)
or \(X_i\cap X_j\neq\emptyset\)), with nonempty supports. For
\(Y\subseteq\Lambda_n\) let
\(w^Y_i=w_i\,\mathbf 1[X_i\cap Y=\emptyset]\). Then
\[\Xi(w^Y)=\Xi(w)\,e^{-p},\qquad |p|\le\frac{|Y|\,\gamma}{1-q},\] where
\(p\) is the sum of the cluster terms with \(X(z)\cap Y\neq\emptyset\).

\begin{itemize}
\tightlist
\item
  Hence \(e^{-C|Y|}\Xi(w)\le\Xi(w^Y)\le e^{C|Y|}\Xi(w)\) with
  \(C=\gamma/(1-q)\), uniformly in \(n\).
\item
  The clusters of at least \(m+1\) labels that meet \(Y\) contribute at
  most \(|Y|\gamma q^m/(1-q)\) in absolute value.
\end{itemize}

\emph{Status:} theorem. Used in III.10. \emph{Lean:}
\href{https://github.com/yang-mills-group/OpenAI-YM-20261028/blob/3456e5e75cf8d4cd558955a5c7638fba2d68f464/YM/Push/ABSOLUTE_SUMS/Absence.lean\#L69-L95}{\texttt{Gas.restores\_log}},
\href{https://github.com/yang-mills-group/OpenAI-YM-20261028/blob/3456e5e75cf8d4cd558955a5c7638fba2d68f464/YM/Push/ABSOLUTE_SUMS/Absence.lean\#L40-L63}{\texttt{Gas.missing\_tail}}.

\textbf{Proposition III.5.6 (re-exponentiation with forced polymers).}
Fix \(B\in\mathbb N\) and a common lattice \(\Lambda_n\).

\begin{itemize}
\tightlist
\item
  \emph{Data.} ``Unforced'' polymers \((X_i,u_i)_{i\in\iota}\) and
  ``forced'' polymers \((Y_k,d_k)_{k\in\kappa}\), all with nonempty
  supports, \(u,d\ge0\); \(R\) a relation on \(\iota\) contained in the
  overlap relation.
\item
  \emph{Atoms.} The atoms are of two kinds:

  \begin{itemize}
  \tightlist
  \item
    the forced labels \(k\), with support \(Y_k\) and weight \(d_k\);
  \item
    the tuples \(z\in\iota^{m+1}\), with support \(X(z)\) and weight
    \(|\phi^T_R(z)|\prod_au_{z_a}/(m+1)!\).
  \end{itemize}
\item
  \emph{Constants.} \(A\ge0\), \(\delta'>0\), \(t\ge1\),
  \(0\le\sigma<1\), \(\mu\ge0\), \(\eta>0\), \(q'<1\).
\item
  \emph{Hypotheses.} \(\mathrm{Ready}(u;A+\delta',\gamma,\delta,q)\),
  \(\mathrm{Hit}_B(d;A+\delta',c)\), and
  \[\frac{t\,e^{\mu}c}{1-\sigma}+\frac{t\,\gamma}{(1-\sigma)(1-q)}\ \le\ \eta\ \le\ \sigma,\qquad 4\,\mathrm{ball}_B\,\eta/\delta'\le q'.\]
\end{itemize}

Fix a site \(x\) and \(m\ge0\). Consider finite sets \(s\) of atoms with
multiplicities \(n_\alpha\ge1\) such that:

\begin{itemize}
\tightlist
\item
  \(s\) contains at least one forced label;
\item
  \(\sum_\alpha n_\alpha\ge m\);
\item
  the supports of the atoms in \(s\) form a connected family containing
  \(x\).
\end{itemize}

Summed over all such pairs \((s,(n_\alpha)_{\alpha\in s})\), the
quantity
\[e^{A\,\mathrm{cov}_B(\bigcup_{\alpha\in s}\mathrm{supp}\,\alpha)}\ \prod_{\alpha\in s}\frac{\mathrm{wt}(\alpha)^{n_\alpha}}{n_\alpha!}\]
is at most \(t^{-m}e^{-\mu}\,\eta/(1-q')\).

\emph{Status:} theorem. The general form, with an arbitrary selection
predicate, is
\href{https://github.com/yang-mills-group/OpenAI-YM-20261028/blob/3456e5e75cf8d4cd558955a5c7638fba2d68f464/YM/Push/ABSOLUTE_SUMS/Copies.lean\#L135-L169}{\texttt{Multiple.estimates}}.
Used in III.2, III.10, III.11 and III.13. \emph{Lean:}
\href{https://github.com/yang-mills-group/OpenAI-YM-20261028/blob/3456e5e75cf8d4cd558955a5c7638fba2d68f464/YM/Push/ABSOLUTE_SUMS/Omissions.lean\#L48-L100}{\texttt{OmitClusters.restore}}.

\textbf{Proposition III.5.7 (gases of seed packets).}

\emph{Setting.} Let \(S\) be a block-spin step (III.1) from
\(\Lambda_f\) to \(\Lambda_c\) with block side \(L\), and let
\(l,B\in\mathbb N\).

\begin{itemize}
\tightlist
\item
  A \emph{seed} is a fine cell, a coarse bond or a coarse site. Its
  location is the projection of the cell's vertices, the bond's two
  endpoints, or the site itself. Each location lies within distance
  \(1\) of its anchor.
\item
  A \emph{packet} \(A\)
  (\href{https://github.com/yang-mills-group/OpenAI-YM-20261028/blob/3456e5e75cf8d4cd558955a5c7638fba2d68f464/YM/Pressure/Seeds.lean\#L34-L37}{\texttt{Stores}})
  is a finite set of seeds; \(|A|\) is its number of seeds.
\item
  Let \(D_0(A)\) be the union of its seed locations and \(D_r(A)\) its
  radius-\(r\) neighbourhood. The \emph{skins} are the radius-\(lB\)
  neighbourhoods of the individual seed locations; their union is the
  \emph{dome} \(D(A)=D_{lB}(A)\). The packet is \emph{connected} if its
  skins form a connected family. The sector and matrix estimates below
  use \(D_r(A)\) directly, rather than neighbourhoods of \(D(A)\).
\item
  Put
  \[\tau_L=L^4\,\mathrm{card}\{\text{plaquette orientations}\}+5,\qquad K_{\rm hit}=(2lB+3)^4\,\tau_L,\qquad K_{\rm step}=(4lB+5)^4\,\tau_L .\]
\end{itemize}

\textbf{(a) Packet sums.} Assume \(a,C\ge0\),
\(0\le p_A\le C\,e^{(s-g)|A|}\) for every connected \(A\), and
\(\rho:=K_{\rm step}^2\,e^{\,s-g+a(2l+3)^4}\le q<1\). Then for every
coarse site \(x\) and every \(m\ge0\),
\[\sum_{A\ \mathrm{connected},\ x\in D(A),\ |A|\ge m}p_A\,e^{a\,\mathrm{cov}_B(D(A))}\ \le\ K_{\rm hit}\,C\,\frac{q^{\max(1,m)}}{1-q}.\]
The same bound holds for an arbitrary family of variants with
nonnegative weights, each nonzero variant assigned to a connected packet
and supported in its dome, provided the sum of the variant weights of
each packet is at most \(Ce^{(s-g)|A|}\). For variants, the sum is
pinned by the actual support and uses its covering number.

\textbf{(b) Mixed words.}

\begin{itemize}
\tightlist
\item
  \emph{Alphabet.} Add optional polymers \((X_i,w_i)_{i\in\iota}\) on
  \(\Lambda_c\) with \(w\ge0\) and lengths \(\lambda_i\ge0\). A packet
  has footprint \(D(A)\) and length \(e_s|A|\).
\item
  \emph{Hypotheses} (with \(\varepsilon,\delta,\gamma_{\rm root}>0\) and
  \(a,C,e_s,\mu_s\ge0\)):

  \begin{itemize}
  \tightlist
  \item
    \(0\le p_A\le Ce^{(s-g)|A|}\) on connected packets;
  \item
    \begin{enumerate}
    \def\labelenumi{(\alph{enumi})}
    \tightlist
    \item
      holds with \(a+2\varepsilon+\delta\) in place of \(a\),
      \(s+2\varepsilon e_s+\mu_s\) in place of \(s\), and some
      \(q_p<1\);
    \end{enumerate}
  \item
    \(\mathrm{Hit}_B(w\,e^{2\varepsilon\lambda};\,a+2\varepsilon+\delta,\,\gamma_{\rm opt})\);
  \item
    \(K_{\rm hit}Cq_p/(1-q_p)+\gamma_{\rm opt}\le\gamma_{\rm root}\);
  \item
    \(4\,\mathrm{ball}_B\,\gamma_{\rm root}/\delta\le q<1\).
  \end{itemize}
\item
  \emph{Notation.} For a finite set \(s\) of packets and optional
  labels, \(D(s)\) is the union of the footprints, and
  \(\mathrm{ch}(s)=\sum_{\alpha\in s}\big(\mathrm{cov}_B(\text{footprint of }\alpha)+\text{length of }\alpha\big)\).
\item
  \emph{Claim.} For all \(m,k,r\in\mathbb N\), \(T\in\mathbb R\) and
  every site \(x\),
  \[\sum_s\ \left(\prod_{\alpha\in s}\mathrm{wt}(\alpha)\right)\big(1+\mathrm{ch}(s)\big)^k\,e^{a\,\mathrm{cov}_B(D(s))}\ \le\ \frac{k!\,e^{\varepsilon}}{\varepsilon^{k}}\;e^{-\varepsilon T-\mu_s r}\;\frac{\gamma_{\rm root}\,q^m}{1-q},\]
  the sum over connected \(s\) with \(x\in D(s)\), \(|s|\ge m+1\), at
  least \(r\) seeds in total, and \(\mathrm{ch}(s)\ge T\). Packets that
  are not connected have weight \(0\).
\end{itemize}

\emph{Status:} theorem. The hypotheses of (b) are the records
\href{https://github.com/yang-mills-group/OpenAI-YM-20261028/blob/3456e5e75cf8d4cd558955a5c7638fba2d68f464/YM/Push/ABSOLUTE_SUMS/Mixed.lean\#L47-L58}{\texttt{Mixed.Costs.Buy}}
and
\href{https://github.com/yang-mills-group/OpenAI-YM-20261028/blob/3456e5e75cf8d4cd558955a5c7638fba2d68f464/YM/Push/ABSOLUTE_SUMS/Mixed.lean\#L84-L89}{\texttt{Mixed.Accounts}}.
(b) is used in III.2 and III.10--III.13. \emph{Lean:}
\href{https://github.com/yang-mills-group/OpenAI-YM-20261028/blob/3456e5e75cf8d4cd558955a5c7638fba2d68f464/YM/Push/ABSOLUTE_SUMS/Charge.lean\#L66-L105}{\texttt{Packets.packet\_income}},
\href{https://github.com/yang-mills-group/OpenAI-YM-20261028/blob/3456e5e75cf8d4cd558955a5c7638fba2d68f464/YM/Push/ABSOLUTE_SUMS/Charge.lean\#L124-L175}{\texttt{Packets.packet\_variants\_tail}},
\href{https://github.com/yang-mills-group/OpenAI-YM-20261028/blob/3456e5e75cf8d4cd558955a5c7638fba2d68f464/YM/Push/ABSOLUTE_SUMS/Mixed.lean\#L153-L189}{\texttt{Mixed.mixed\_hit}},
\href{https://github.com/yang-mills-group/OpenAI-YM-20261028/blob/3456e5e75cf8d4cd558955a5c7638fba2d68f464/YM/Push/ABSOLUTE_SUMS/Mixed.lean\#L95-L122}{\texttt{Mixed.ready\_vertices}}.

\textbf{Proposition III.5.8 (the numerical conditions hold at large
coupling).}

\emph{Notation.} Fix the gauge group. In (a) and (b), also fix the block
side \(L\) and \(l\in\mathbb N\). For the coarse coupling \(x\) put:

\begin{itemize}
\tightlist
\item
  \(\beta_F=x+b\log L\), with \(b\) a constant of the group;
\item
  \(B=\lceil(\log x)^2\rceil\);
\item
  \(P(s)=e^{(\log\log s)^2}\);
\item
  \(t=\beta_F^{-1/2}P(\beta_F)\), the small-field threshold;
\item
  \(\mathrm{bill}(x)=1+B+(\log\beta_F)^2\);
\item
  for \(\mu,u>0\), \(M=\mu u\,P(\beta_F)^2\).
\end{itemize}

\textbf{(a)} Let \(a,D\ge0\), \(d\in\mathbb N\) and \(z<2\). For all
sufficiently large \(x\) the following holds. Whenever
\(s\le D\,\mathrm{bill}^d P(\beta_F)^z\), \(g\ge M\) and
\(0\le C\le \exp(D\,\mathrm{bill}^dP(\beta_F)^z)\):

\begin{itemize}
\tightlist
\item
  the numerical condition \(\rho\le q<1\) of III.5.7(a) holds with
  \(q=e^{-M/2}\);
\item
  \(K_{\rm hit}Cq^{\max(1,m)}/(1-q)\le e^{-M\max(1,m)/4}\) for every
  \(m\).
\end{itemize}

If \(\beta\ge u\beta_F\), then \(\beta\mu t^2\ge M\).

\textbf{(b)} Let in addition \(\varepsilon,\delta>0\). For all large
\(x\), all \(\beta\ge u\beta_F\), and all \(C,e_s\ge0\),
\(s\in\mathbb R\), and \(\gamma_{\rm opt}\) with
\[C\le e^{\rm fee},\qquad s+2\varepsilon e_s\le{\rm fee}:=D\,\mathrm{bill}^dP(\beta_F)^z,\qquad 0\le\gamma_{\rm opt}\le\delta/(16\,\mathrm{ball}_B),\]
the numerical inequalities of III.5.7(b) hold with:

\begin{itemize}
\tightlist
\item
  gain \(g=\beta\mu t^2\);
\item
  mark \(\mu_s=\mu uP(\beta_F)^2/2\) (half of the guaranteed gain is
  extracted per seed);
\item
  \(q=1/2\);
\item
  \(\gamma_{\rm root}=e^{-\mu uP(\beta_F)^2/8}+\gamma_{\rm opt}\).
\end{itemize}

The packet envelope and optional site-sum estimate in III.5.7(b) remain
hypotheses on the actual weights.

\textbf{(c)} For a step-estimator \(\mathcal T\) (III.1), there is \(n\)
such that for every \(L\ge n\) and all sufficiently large \(W\) there is
\(\mu>0\), chosen before \(l,d,u,a,D,z\), with the following property
for every choice of \(l,d,u,a,D,z\) satisfying (a). For all sufficiently
large \(x\), every funded step with block side \(L\) and coarse side
\(\ge100(1+B)\) satisfies the point charge of III.5.11 with constant
\(\mu\). For every \(\beta\ge u\beta_F\), the numerical conclusions of
(a) hold with \(g=\beta\mu t^2\). The threshold in \(x\) may depend on
the fixed numerical parameters but not on the volume.

\emph{Status:} theorem. (c) takes the energy grant per seed from III.10.
Used in III.10. \emph{Lean:}
\href{https://github.com/yang-mills-group/OpenAI-YM-20261028/blob/3456e5e75cf8d4cd558955a5c7638fba2d68f464/YM/Push/ABSOLUTE_SUMS/Delay.lean\#L165-L214}{\texttt{Delay.fits}},
\href{https://github.com/yang-mills-group/OpenAI-YM-20261028/blob/3456e5e75cf8d4cd558955a5c7638fba2d68f464/YM/Push/ABSOLUTE_SUMS/Delay.lean\#L43-L52}{\texttt{Delay.comparison}},
\href{https://github.com/yang-mills-group/OpenAI-YM-20261028/blob/3456e5e75cf8d4cd558955a5c7638fba2d68f464/YM/Push/ABSOLUTE_SUMS/Buy.lean\#L67-L103}{\texttt{Buy.quoted}},
\href{https://github.com/yang-mills-group/OpenAI-YM-20261028/blob/3456e5e75cf8d4cd558955a5c7638fba2d68f464/YM/Push/ABSOLUTE_SUMS/Grants.lean\#L52-L71}{\texttt{Grant.order}}.

\textbf{Proposition III.5.9 (logarithm by deletion).}

\textbf{(a) Abstract form.}

\begin{itemize}
\item
  \emph{Data.} \(J\) finite; \(X:\iota\to\) nonempty subsets of \(J\)
  (\(\iota\) any set); \(v:\iota\to[0,\infty)\); a weight \(w\) on
  subsets of \(J\) with \(w\ge1\) and \(w(P\cup P')\le w(P)w(P')\)
  whenever \(P\cap P'\neq\emptyset\); \(a\ge0\);
  \(0\le\delta\le1-e^{-a}\).
\item
  \emph{Site condition.} For every \(x\in J\),
  \[\sum_{i:\,x\in X_i}w(X_i)\,v_i\,e^{a|X_i|}\le\delta .\]
\item
  \emph{Gas and its pieces.} For \(u:\iota\to\mathbb R\) with
  \(|u_i|\le v_i\), let \(Z_u(B)=\sum_Q\prod_{i\in Q}u_i\)
  (\(B\subseteq J\)), the sum over finite families \(Q\) of labels with
  pairwise disjoint supports contained in \(B\). Let \(\lambda_u\) be
  the Möbius pieces of \(-\log Z_u\).
\item
  \emph{Claim.}

  \begin{itemize}
  \tightlist
  \item
    \(Z_u(B)>0\) and
    \(Z_u(B)=\exp\big(-\sum_{A\subseteq B}\lambda_u(A)\big)\) for every
    \(B\).
  \item
    \(|\lambda_u(A)|\le\lambda_{-v}(A)\), and
    \(\sum_{A\ni x}w(A)\,\lambda_{-v}(A)\le a\) for every \(x\).
  \item
    For every \(Q\subseteq J\), \(Z_u(J\setminus Q)=Z_u(J)\,e^{l}\),
    with \(l=\sum_{A\cap Q\neq\emptyset}\lambda_u(A)\) and
    \(|l|\le a|Q|\).
  \end{itemize}
\end{itemize}

\textbf{(b) At a coarse field.}

\begin{itemize}
\tightlist
\item
  \emph{Data.} Let \(\iota\) be countable, with nonempty supports
  \(X_i\). Let \(E_i(V)\) (\(i\in\iota\)) be functions of the gauge
  field on \(\Lambda_n\). Each is local in \(X_i\), gauge invariant and
  Borel, with \(|E_i(V)|\le v_i\) whenever no site of \(X_i\) is
  flagged. Let \(\tilde E_i\) be \(E_i\) set to \(0\) when \(X_i\)
  contains a flagged site.
\item
  \emph{Hypotheses.} \(M\in\mathbb N\), \(A,a\ge0\), \(v_i\ge0\),
  \(a(2M+1)^4\le1\), \(\delta\le1-e^{-a}\), and
  \(\sum_{i:\,x\in X_i}e^{(A+1)\mathrm{cov}_M(X_i)}v_i\le\delta\) for
  every \(x\).
\item
  \emph{Claim.} The gas \(Z_V\) of \(\tilde E(V)\) satisfies, for every
  \(V\) and every \(Q\):

  \begin{itemize}
  \tightlist
  \item
    \(Z_V(Q)>0\) and
    \(Z_V(Q)=\exp\big(-\sum_{P\subseteq Q}\lambda_P(V)\big)\);
  \item
    \(|\lambda_P(V)|\le g(P)\), with \(g\) independent of \(V\), and
    \(\sum_{P\ni x}e^{A\,\mathrm{cov}_M(P)}g(P)\le a\);
  \item
    \(Z_V(\Lambda_n\setminus Q)=Z_V(\Lambda_n)e^{l}\) with
    \(l=\sum_{P\cap Q\ne\emptyset}\lambda_P(V)\) and \(|l|\le a|Q|\);
  \item
    there are gauge-invariant Borel functions \(\lambda_P^{\rm raw}\),
    local in the links with both endpoints in \(P\), such that
    \(\lambda_P(V)=\mathbf1[P\text{ contains no flagged site}]\lambda_P^{\rm raw}(V)\).
    Thus the masked coefficients also use the flag tests on \(P\).
  \end{itemize}
\end{itemize}

\emph{Status:} theorem. (b) is used in III.11. \emph{Lean:}
\href{https://github.com/yang-mills-group/OpenAI-YM-20261028/blob/3456e5e75cf8d4cd558955a5c7638fba2d68f464/YM/SectorNext/VACUUM_NORMALIZATION/Gas.lean\#L217-L228}{\texttt{ReLog.catalog}},
\href{https://github.com/yang-mills-group/OpenAI-YM-20261028/blob/3456e5e75cf8d4cd558955a5c7638fba2d68f464/YM/SectorNext/VACUUM_NORMALIZATION/Gas.lean\#L235-L271}{\texttt{ReLog.omitted}},
\href{https://github.com/yang-mills-group/OpenAI-YM-20261028/blob/3456e5e75cf8d4cd558955a5c7638fba2d68f464/YM/SectorNext/VACUUM_NORMALIZATION/Deletion.lean\#L224-L233}{\texttt{ReLog.swept}},
\href{https://github.com/yang-mills-group/OpenAI-YM-20261028/blob/3456e5e75cf8d4cd558955a5c7638fba2d68f464/YM/SectorNext/VACUUM_NORMALIZATION/Bulk.lean\#L157-L183}{\texttt{ReLog.intake}},
\href{https://github.com/yang-mills-group/OpenAI-YM-20261028/blob/3456e5e75cf8d4cd558955a5c7638fba2d68f464/YM/SectorNext/VACUUM_NORMALIZATION/Bulk.lean\#L99-L101}{\texttt{ReLog.retest}},
\href{https://github.com/yang-mills-group/OpenAI-YM-20261028/blob/3456e5e75cf8d4cd558955a5c7638fba2d68f464/YM/SectorNext/VACUUM_NORMALIZATION/Bulk.lean\#L118-L143}{\texttt{ReLog.retain\_mask}}.

\textbf{Proposition III.5.10 (trace-log of a near-identity product).}
Let \(A\) be positive definite and \(C\) symmetric (real matrices), and
\(S=AC-1\) with \(0\le q<1\) and \(\sum_j|S_{ij}|\le q\) for every
\(i\). Then:

\begin{itemize}
\tightlist
\item
  \(\det C>0\) and \(\det(1+S)>0\);
\item
  \(\sum_{n\ge1}\frac{(-1)^{n+1}}{n}\,\mathrm{tr}\,S^n\) converges
  absolutely, and its sum is \(\log\det(1+S)\).
\end{itemize}

\emph{Status:} theorem. Used in III.4. \emph{Lean:}
\href{https://github.com/yang-mills-group/OpenAI-YM-20261028/blob/3456e5e75cf8d4cd558955a5c7638fba2d68f464/YM/SectorNext/VACUUM_NORMALIZATION/Round.lean\#L149-L155}{\texttt{Round.exact\_trace}},
\href{https://github.com/yang-mills-group/OpenAI-YM-20261028/blob/3456e5e75cf8d4cd558955a5c7638fba2d68f464/YM/SectorNext/VACUUM_NORMALIZATION/Round.lean\#L95-L144}{\texttt{Round.twin\_log}}.

\textbf{Proposition III.5.11 (sector decomposition of a block-spin step;
gain per seed).}

\emph{Setting.} Fix a step \(S\) with block side \(L\), a step-estimator
\(\mathcal T\), a tube \(T\) with \(\mathrm{Funds}(S,T)\) (III.1), a
coarse coupling \(x\) with threshold \(t\), and constants
\(W,\kappa>0\).

\emph{Sector of a packet.} Say \((U,k,V)\) satisfies the sector
prerequisites of \(A\) when:

\begin{itemize}
\tightlist
\item
  each fine cell whose projected vertices meet \(D_{5B}(A)\) has
  plaquette chord \(\ge t/4\) if listed in \(A\), and \(\le t/2\)
  otherwise;
\item
  each coarse bond with an endpoint in \(D_{5B}(A)\) has predicted wire
  chord \(\ge Wt\) if listed in \(A\), and \(\le2Wt\) otherwise;
\item
  on \(D_{5B}(A)\), the sites listed in \(A\) are exactly the sites
  flagged by \(V\) at the coarse threshold
  \(t_C=x^{-1/2}e^{(\log\log x)^2}\);
\item
  the frame tag has positive frame mass at each site of \(D_{5B}(A)\).
\end{itemize}

\emph{Point charge.} For every tree-gauge forest \(F\), every fine field
\(U\) equal to \(1\) on the trunks of \(F\), every frame tag \(k\),
every coarse field \(V\) and every packet \(A\) such that \((U,k,V)\)
satisfies these prerequisites:
\[\mathcal E_A(U,V)\ \ge\ \kappa\,t^2\,|A| .\] Here \(\mathcal E_A\) is
the \emph{assessed energy} of \(A\): the sum of the fine,
block-constraint, coarse and cross terms of III.1 with anchors in
\(D_{2B}(A)\), using walks truncated at \(\lfloor B/100000\rfloor\).

\emph{Claim.} Let \(R\) be a lawful representation of the fine density
(III.1) adapted to \(T\): its path data and clip are those of the upper
stage of \(T\), and its threshold is \(t\). Let its coupling satisfy
\(\beta\ge0\). Let \(F\) be any tree-gauge forest and \(D=(D_A)_A\) any
allocation: \(D_A(k,U,V)\ge0\) is Borel, \(\sum_A D_A=1\), and
\(D_A\ne0\) implies the two seed tests above for every fine cell and
coarse bond, together with equality of the flag sets on the whole coarse
lattice. Positive likelihood then implies the local sector
prerequisites. This global support condition is needed to apply the
point charge in every lattice orientation. Then for every coarse field
\(V\)
\[(\mathrm{push}_\beta\,\rho_R)(V)=\sum_A c_A(V),\qquad |c_A(V)|\ \le\ \pi_\beta\;e^{-\beta\kappa t^2|A|}\;\ell_A(V),\]
where:

\begin{itemize}
\tightlist
\item
  \(c_A\) is the \(A\)-th sector output;
\item
  \(\pi_\beta=Z_{\rm bond}(\beta)^{-|\mathrm{Bond}(\Lambda_c)|}\), where
  \(Z_{\rm bond}(\beta)\) is the single-bond penalty integrated against
  probability Haar measure;
\item
  \(\ell_A(V)\) is the remaining sector integral, including the coarse
  factor \(\mathrm{sieve}(V)\), with the orientation-averaged assessed
  energy removed and absolute values on the exposed density;
\item
  \(c_A\) and \(\ell_A\) are bounded Borel functions.
\end{itemize}

\emph{Uniformity.}

\begin{itemize}
\tightlist
\item
  At a committed pair of scales of the run, a single \(\kappa>0\) serves
  all volumes, all steps of that pair, and all lawful representations
  produced by the run whose coupling lies in the admissible band.
\item
  For a fixed step-estimator, choose \(L\) sufficiently large, then
  \(W\) sufficiently large, then \(\kappa>0\). The point charge holds
  for all sufficiently large \(x\) and all funded steps of coarse side
  \(\ge100(1+B)\), uniformly in their volumes.
\end{itemize}

\emph{Status:} theorem, conditional on \(\mathrm{Funds}\) and the point
charge. The point charge is discharged in two ways:

\begin{itemize}
\tightlist
\item
  in general, from the Teller's energy grant (III.10);
\item
  at committed scales, from the committed-scale energy hypothesis of the
  run, through
  \href{https://github.com/yang-mills-group/OpenAI-YM-20261028/blob/3456e5e75cf8d4cd558955a5c7638fba2d68f464/YM/Pressure/Seeds.lean\#L166-L181}{\texttt{notation\_ticket}}
  (III.1).
\end{itemize}

Used in III.4 and III.10. \emph{Lean:}
\href{https://github.com/yang-mills-group/OpenAI-YM-20261028/blob/3456e5e75cf8d4cd558955a5c7638fba2d68f464/YM/Push/SECTOR_MEASURE/Presented.lean\#L157-L210}{\texttt{Presented.portionCost}}
(conclusion
\href{https://github.com/yang-mills-group/OpenAI-YM-20261028/blob/3456e5e75cf8d4cd558955a5c7638fba2d68f464/YM/Push/SECTOR_MEASURE/Presented.lean\#L151-L155}{\texttt{PortionCost}}),
\href{https://github.com/yang-mills-group/OpenAI-YM-20261028/blob/3456e5e75cf8d4cd558955a5c7638fba2d68f464/YM/Push/SECTOR_MEASURE/Consumer.lean\#L42-L65}{\texttt{Consumer.bought}},
\href{https://github.com/yang-mills-group/OpenAI-YM-20261028/blob/3456e5e75cf8d4cd558955a5c7638fba2d68f464/YM/Push/SECTOR_MEASURE/Oriented.lean\#L37-L51}{\texttt{points\_grant}};
definitions
\href{https://github.com/yang-mills-group/OpenAI-YM-20261028/blob/3456e5e75cf8d4cd558955a5c7638fba2d68f464/YM/Push/SECTOR_MEASURE/Cells.lean\#L51-L59}{\texttt{Seeded}},
\href{https://github.com/yang-mills-group/OpenAI-YM-20261028/blob/3456e5e75cf8d4cd558955a5c7638fba2d68f464/YM/Push/SECTOR_MEASURE/Cells.lean\#L62}{\texttt{Alive}},
\href{https://github.com/yang-mills-group/OpenAI-YM-20261028/blob/3456e5e75cf8d4cd558955a5c7638fba2d68f464/YM/Push/SECTOR_MEASURE/Oriented.lean\#L29-L32}{\texttt{PointCharge}},
\href{https://github.com/yang-mills-group/OpenAI-YM-20261028/blob/3456e5e75cf8d4cd558955a5c7638fba2d68f464/YM/Push/SECTOR_MEASURE/Grant.lean\#L28-L29}{\texttt{earn}},
\href{https://github.com/yang-mills-group/OpenAI-YM-20261028/blob/3456e5e75cf8d4cd558955a5c7638fba2d68f464/YM/Push/SECTOR_MEASURE/Profile.lean\#L30-L36}{\texttt{Allocation}}.

\textbf{Proposition III.5.12 (Combes--Thomas bound).} Let \(M\) be a
real symmetric matrix on a finite index set, with positions
\(a_i\in\Lambda_n\). Assume \(M\ge r>0\), \(\eta\ge0\), and
\[\sup_i\sum_je^{\eta d(a_i,a_j)}|M_{ij}|\le C,\qquad \sup_j\sum_ie^{\eta d(a_i,a_j)}|M_{ij}|\le C .\]
If \(0\le\theta\le1\) and \(\theta C\le r/2\), then for every \(u\ge0\)
the matrix \(M+uI\) is invertible and
\[e^{\theta\eta\,d(a_i,a_j)}\,\big|(M+uI)^{-1}_{ij}\big|\le\big(\min(1,r/2)\,(1+u)\big)^{-1}\quad\text{for all }i,j .\]
The hypotheses pass to principal submatrices. \emph{Status:} theorem.
\emph{Lean:}
\href{https://github.com/yang-mills-group/OpenAI-YM-20261028/blob/3456e5e75cf8d4cd558955a5c7638fba2d68f464/YM/Progress/STATIONARY_SCENES/Green.lean\#L188-L251}{\texttt{Green.inverse}},
\href{https://github.com/yang-mills-group/OpenAI-YM-20261028/blob/3456e5e75cf8d4cd558955a5c7638fba2d68f464/YM/Progress/STATIONARY_SCENES/Green.lean\#L93-L113}{\texttt{Green.Bottom.restrict}},
\href{https://github.com/yang-mills-group/OpenAI-YM-20261028/blob/3456e5e75cf8d4cd558955a5c7638fba2d68f464/YM/Progress/STATIONARY_SCENES/Green.lean\#L78-L88}{\texttt{Green.Rows.restrict}}.

\textbf{Proposition III.5.13 (windowed inverses).}

\emph{Setting.}

\begin{itemize}
\tightlist
\item
  A finite index set, positions \(d_j\in\Lambda_n\), and \(U\ge0\) such
  that \(\mathrm{card}\{j:d(y,d_j)\le\rho\}\le U(1+\rho)^4\) for every
  \(y\) and every \(\rho\in\mathbb N\).
\item
  \(A,r,C,\eta,\theta\) as in III.5.12, with \(C\ge0\) and
  \(\eta,\theta>0\); and \(0\le\mu<\nu<\theta\eta\).
\item
  A radius \(R\in\mathbb N\) and windows \(P(j)\ni j\) such that
  \(d(d_i,d_j)\le R\Rightarrow i\in P(j)\).
\end{itemize}

\emph{Objects.}

\begin{itemize}
\tightlist
\item
  \(\mathrm{within}(A)_{ij}=\big((A|_{P(j)})^{-1}\big)_{ij}\) for
  \(i\in P(j)\), and \(0\) otherwise.
\item
  \(E=A\,\mathrm{within}(A)-1\) and
  \(\mathrm{cov}=\tfrac12(\mathrm{within}+\mathrm{within}^T)\).
\item
  \(\Pi_U(\eta)=U\,4!\,(2/\eta)^4\,e^{\eta/2}/(1-e^{-\eta/2})\),
  \(Y=\Pi_U(\theta\eta-\nu)/\min(1,r/2)\), and
  \(s=e^{-(\nu-\mu)R}(CY+1)\).
\end{itemize}

\textbf{(a)} Weighted row and column sums at rate \(\mu\) satisfy:

\begin{itemize}
\tightlist
\item
  \(A^{-1}\) and \(\mathrm{within}(A)\): at most \(Y\);
\item
  \(E\): at most \(s\);
\item
  \(\mathrm{within}(A)-A^{-1}\): at most \(Ys\).
\end{itemize}

At any rate \(0\le\xi\le\eta\), if \(\mathrm{within}-A^{-1}\) has row
and column bounds \(s'\ge0\) and \(\mathrm{within}\) has bounds
\(Y'\ge0\), then \(\mathrm{cov}\), \(\mathrm{cov}-A^{-1}\) and
\(A\,\mathrm{cov}-1\) have bounds \(Y'\), \(s'\) and \(Cs'\),
respectively. If also \(C,Y'>0\) and \(s'\le1/(2C)\), then
\(\mathrm{cov}\ge1/(2C)\) and \(\mathrm{cov}^{-1}\ge(Y')^{-1}\).

\textbf{(b)} Fix positive bounds on the weighted row and column sums,
spectral gap, rate \(\alpha\), and counting constant. There are
\(0<e\le\alpha\) and \(K,c,c_1,c_2>0\), depending only on these data,
such that for all sufficiently large integer \(R\) the following holds
uniformly for all such matrices and windows with room \(R\). At rate
\(e\), the full and windowed inverses of \(tA+(1-t)I\) have bounds
\(K\), and their residual and inverse difference have bounds
\(Ke^{-cR}<1\), for every \(t\in[0,1]\). At \(t=1\), the covariance has
bound \(K\), its difference from \(A^{-1}\) and the residual
\(A\,\mathrm{cov}-I\) have bounds \(Ke^{-cR}\), and
\(\mathrm{cov}\ge c_1\), \(\mathrm{cov}^{-1}\ge c_2\).

\emph{Status:} theorem. \emph{Lean:}
\href{https://github.com/yang-mills-group/OpenAI-YM-20261028/blob/3456e5e75cf8d4cd558955a5c7638fba2d68f464/YM/AnalyticNext/CONNECTED_ACTIVITY_BOUNDS/Windows.lean\#L146-L169}{\texttt{Win.Windows.window\_price}},
\href{https://github.com/yang-mills-group/OpenAI-YM-20261028/blob/3456e5e75cf8d4cd558955a5c7638fba2d68f464/YM/AnalyticNext/CONNECTED_ACTIVITY_BOUNDS/Windows.lean\#L171-L188}{\texttt{covariance\_price}},
\href{https://github.com/yang-mills-group/OpenAI-YM-20261028/blob/3456e5e75cf8d4cd558955a5c7638fba2d68f464/YM/AnalyticNext/CONNECTED_ACTIVITY_BOUNDS/Windows.lean\#L191-L202}{\texttt{covariance\_floor}},
\href{https://github.com/yang-mills-group/OpenAI-YM-20261028/blob/3456e5e75cf8d4cd558955a5c7638fba2d68f464/YM/AnalyticNext/CONNECTED_ACTIVITY_BOUNDS/Birthday.lean\#L106-L178}{\texttt{Win.Promise.receive}}.

\textbf{Proposition III.5.14 (precision matrix of the small-field
Gaussian of a step).}

\emph{Setting.} Fix a step-estimator \(\mathcal T\) and a patch
\(\mathcal D\) (block side \(L\), order \(J\)) satisfying the branch
conditions for its kernel and gauge group. For a step with patch
\(\mathcal D\), tube \(T\), coupling \(x\), tree-gauge forest \(F\),
packet \(A\) and coarse field \(V\), define \(D\) as follows:

\begin{itemize}
\tightlist
\item
  \emph{Rows.} Use the flat plaquette-curl and upper path-sum rows, with
  walks truncated at \(\lfloor B/100000\rfloor\), and the flat
  linearized block-constraint rows.
\item
  \emph{Transport.} Transport the input Lie-algebra field by the replica
  of \(V\) and a fixed path gauge centred at each row, in every lattice
  orientation.
\item
  \emph{Restriction.} Keep rows whose oriented centres lie outside
  \(D_{8B}(A)\), and columns on \emph{active} fine bonds: bonds outside
  the trunks of \(F\) whose projected endpoints lie outside
  \(D_{10B}(A)\). Set all other entries to zero.
\item
  \emph{Normalization.} Write Lie-algebra coordinates and multiply each
  row by the inverse square root of the number of lattice orientations.
\end{itemize}

This is a transported flat linearization. Its equality to the derivative
of the nonlinear cut rows at \(V\) is not asserted.

Let \(\mathrm{gram}=D^TD+\mathrm{diag}(\mathbf 1_{\rm inactive})\), with
positions the coarse parents of the bonds.

\emph{Claim.} There are constants \(r,K,U,\alpha>0\), with \(\alpha\)
half the decay rate of the upper path data, such that the following
holds for all large \(x\). For every step with patch \(\mathcal D\) and
\(\mathrm{Funds}\), coarse side \(\ge100(1+B)\), every \(F\), and every
\(A,V\) for which the flags of \(A\) are the flagged sites of \(V\):

\begin{itemize}
\tightlist
\item
  \(\mathrm{gram}\) is symmetric and \(\mathrm{gram}\ge r\);
\item
  at most \(U(1+\rho)^4\) indices lie within distance \(\rho\) of any
  site;
\item
  its weighted row and column sums at rate \(\alpha\) are at most \(K\);
\item
  \(D_{ij}=0\) when the row centre and column position are separated by
  more than \(\lfloor B/100000\rfloor+R_{\rm pred}(\mathcal D)+2\),
  where \(R_{\rm pred}\) is the patch's predictor range. The range grows
  with \(B\); the weighted bounds are uniform in \(x\).
\end{itemize}

Consequently, for large \(x\), III.5.13(b) applies to the following
windows: a fine bond belongs to the window of an addressed bond exactly
when at least one of its projected endpoints is within distance
\(\lfloor B/4\rfloor\) of at least one projected endpoint of the
addressed bond. Include every color coordinate of each such bond. Their
room is \(\lfloor B/4\rfloor\), and the error is
\(K_{\rm win}e^{-c_{\rm win}\lfloor B/4\rfloor}<1\).

\emph{Status:} theorem, conditional on \(\mathrm{Funds}\) and
\(\mathrm{Branches}\). Used in III.4 and III.11 through the window
estimates below. \emph{Lean:}
\href{https://github.com/yang-mills-group/OpenAI-YM-20261028/blob/3456e5e75cf8d4cd558955a5c7638fba2d68f464/YM/AnalyticNext/CONNECTED_ACTIVITY_BOUNDS/Bedrock.lean\#L142-L158}{\texttt{Aff.bedrock}},
\href{https://github.com/yang-mills-group/OpenAI-YM-20261028/blob/3456e5e75cf8d4cd558955a5c7638fba2d68f464/YM/AnalyticNext/CONNECTED_ACTIVITY_BOUNDS/Bedrock.lean\#L72-L113}{\texttt{Aff.surplus}},
\href{https://github.com/yang-mills-group/OpenAI-YM-20261028/blob/3456e5e75cf8d4cd558955a5c7638fba2d68f464/YM/AnalyticNext/CONNECTED_ACTIVITY_BOUNDS/IndexCost.lean\#L165-L181}{\texttt{Aff.gram\_budget}},
\href{https://github.com/yang-mills-group/OpenAI-YM-20261028/blob/3456e5e75cf8d4cd558955a5c7638fba2d68f464/YM/AnalyticNext/CONNECTED_ACTIVITY_BOUNDS/Warsaw.lean\#L25-L42}{\texttt{Aff.install\_promises}};
definitions
\href{https://github.com/yang-mills-group/OpenAI-YM-20261028/blob/3456e5e75cf8d4cd558955a5c7638fba2d68f464/YM/AnalyticNext/CONNECTED_ACTIVITY_BOUNDS/Affine.lean\#L125-L127}{\texttt{Aff.Dmat}},
\href{https://github.com/yang-mills-group/OpenAI-YM-20261028/blob/3456e5e75cf8d4cd558955a5c7638fba2d68f464/YM/AnalyticNext/CONNECTED_ACTIVITY_BOUNDS/Affine.lean\#L145-L147}{\texttt{Aff.gram}},
\href{https://github.com/yang-mills-group/OpenAI-YM-20261028/blob/3456e5e75cf8d4cd558955a5c7638fba2d68f464/YM/AnalyticNext/CONNECTED_ACTIVITY_BOUNDS/Birthday.lean\#L36-L37}{\texttt{Promise.Served}}.

\textbf{Further outputs.}

\begin{itemize}
\tightlist
\item
  A calculus of budgets \(\mathrm{SB}\)
  (\href{https://github.com/yang-mills-group/OpenAI-YM-20261028/blob/3456e5e75cf8d4cd558955a5c7638fba2d68f464/YM/Push/ABSOLUTE_SUMS/Basic.lean}{\texttt{Basic.lean}}):
  monotonicity, scaling, pull-back, sums over fibres, geometric tails,
  \href{https://github.com/yang-mills-group/OpenAI-YM-20261028/blob/3456e5e75cf8d4cd558955a5c7638fba2d68f464/YM/Push/ABSOLUTE_SUMS/Basic.lean\#L143-L163}{\texttt{budget\_count\_tail}}.
  It is used by about 125 modules elsewhere.
\item
  Site-sum bounds for bundles of walks on the fine lattice, with a tail
  in the true length of the walks
  (\href{https://github.com/yang-mills-group/OpenAI-YM-20261028/blob/3456e5e75cf8d4cd558955a5c7638fba2d68f464/YM/Push/ABSOLUTE_SUMS/Path.lean\#L182-L201}{\texttt{Paths.bundle\_hit}}).
\item
  A weighted row bound for the tree-gauge response kernel
  (\href{https://github.com/yang-mills-group/OpenAI-YM-20261028/blob/3456e5e75cf8d4cd558955a5c7638fba2d68f464/YM/Progress/STATIONARY_SCENES/Response.lean\#L242-L259}{\texttt{Response.response\_price}}).
\item
  Auxiliary Gaussian estimates, listed in Lemma 4.36.
\end{itemize}

\subsubsection{2. What it uses}\label{what-it-uses-7}

\textbf{From other units.}

\begin{itemize}
\tightlist
\item
  \textbf{I.1:} the torus \(\Lambda_n\); the distance \(d\)
  (\href{https://github.com/yang-mills-group/OpenAI-YM-20261028/blob/3456e5e75cf8d4cd558955a5c7638fba2d68f464/YM/Lattice/Site.lean\#L81}{\texttt{separation}});
  \(|X|\le\mathrm{ball}_M\mathrm{cov}_M(X)\)
  (\href{https://github.com/yang-mills-group/OpenAI-YM-20261028/blob/3456e5e75cf8d4cd558955a5c7638fba2d68f464/YM/Corona/Fill.lean\#L129-L131}{\texttt{Corona.pays}});
  subadditivity of \(\mathrm{cov}_M\) on overlapping unions
  (\href{https://github.com/yang-mills-group/OpenAI-YM-20261028/blob/3456e5e75cf8d4cd558955a5c7638fba2d68f464/YM/Corona/Halo.lean\#L14-L30}{\texttt{Corona.meets\_sum}}).
\item
  \textbf{II.2:} the site-sum condition
  \href{https://github.com/yang-mills-group/OpenAI-YM-20261028/blob/3456e5e75cf8d4cd558955a5c7638fba2d68f464/YM/Ac/Tree.lean\#L12-L13}{\texttt{Ac.Hit}};
  connected families
  \href{https://github.com/yang-mills-group/OpenAI-YM-20261028/blob/3456e5e75cf8d4cd558955a5c7638fba2d68f464/YM/Ac/Nodes.lean\#L49}{\texttt{Ac.Spark}};
  \(\mathrm{cov}_M\) of a connected union is at most the sum
  (\href{https://github.com/yang-mills-group/OpenAI-YM-20261028/blob/3456e5e75cf8d4cd558955a5c7638fba2d68f464/YM/Ac/Tree.lean\#L59-L82}{\texttt{Ac.cover}}).
\item
  \textbf{III.1:}

  \begin{itemize}
  \tightlist
  \item
    \(\mathrm{SB}\);
  \item
    the step data: patch, kernel, frame choices, the pushforward
    \href{https://github.com/yang-mills-group/OpenAI-YM-20261028/blob/3456e5e75cf8d4cd558955a5c7638fba2d68f464/YM/RG/Block.lean\#L77-L79}{\texttt{StepData.push}},
    \href{https://github.com/yang-mills-group/OpenAI-YM-20261028/blob/3456e5e75cf8d4cd558955a5c7638fba2d68f464/YM/RG/Patch.lean\#L182-L185}{\texttt{Branches}};
  \item
    lawful representations
    \href{https://github.com/yang-mills-group/OpenAI-YM-20261028/blob/3456e5e75cf8d4cd558955a5c7638fba2d68f464/YM/RG/History.lean\#L181-L183}{\texttt{Rep.Lawful}};
  \item
    packets, domes, sectors and the assessed energy
    (\href{https://github.com/yang-mills-group/OpenAI-YM-20261028/blob/3456e5e75cf8d4cd558955a5c7638fba2d68f464/YM/Pressure/Seeds.lean\#L34-L37}{\texttt{Stores}},
    \href{https://github.com/yang-mills-group/OpenAI-YM-20261028/blob/3456e5e75cf8d4cd558955a5c7638fba2d68f464/YM/Pressure/Seeds.lean\#L51-L58}{\texttt{Plot}},
    \href{https://github.com/yang-mills-group/OpenAI-YM-20261028/blob/3456e5e75cf8d4cd558955a5c7638fba2d68f464/YM/Pressure/Seeds.lean\#L85}{\texttt{assessed}});
  \item
    the step-estimator interface
    \href{https://github.com/yang-mills-group/OpenAI-YM-20261028/blob/3456e5e75cf8d4cd558955a5c7638fba2d68f464/YM/Pressure/Seeds.lean\#L89-L96}{\texttt{Teller}}
    and
    \href{https://github.com/yang-mills-group/OpenAI-YM-20261028/blob/3456e5e75cf8d4cd558955a5c7638fba2d68f464/YM/Pressure/Seeds.lean\#L101-L109}{\texttt{Teller.Funds}},
    with \(\mathrm{Funds}\) supplied for the run by
    \href{https://github.com/yang-mills-group/OpenAI-YM-20261028/blob/3456e5e75cf8d4cd558955a5c7638fba2d68f464/YM/Pressure/Seeds.lean\#L147-L152}{\texttt{Testing.proceeds}};
  \item
    tubes, tree-gauge forests
    (\href{https://github.com/yang-mills-group/OpenAI-YM-20261028/blob/3456e5e75cf8d4cd558955a5c7638fba2d68f464/YM/Tree/Pines.lean\#L24-L29}{\texttt{Tree.Forest}})
    and flags;
  \item
    the constants \(B\), \(t\), \(\beta_F\)
    (\href{https://github.com/yang-mills-group/OpenAI-YM-20261028/blob/3456e5e75cf8d4cd558955a5c7638fba2d68f464/YM/Pressure/Seeds.lean\#L22}{\texttt{radiusB}},
    \href{https://github.com/yang-mills-group/OpenAI-YM-20261028/blob/3456e5e75cf8d4cd558955a5c7638fba2d68f464/YM/Pressure/Seeds.lean\#L26}{\texttt{fiery}},
    \href{https://github.com/yang-mills-group/OpenAI-YM-20261028/blob/3456e5e75cf8d4cd558955a5c7638fba2d68f464/YM/Pressure/Seeds.lean\#L23}{\texttt{betaF}});
  \item
    the point charge at committed scales
    (\href{https://github.com/yang-mills-group/OpenAI-YM-20261028/blob/3456e5e75cf8d4cd558955a5c7638fba2d68f464/YM/Pressure/Seeds.lean\#L166-L181}{\texttt{notation\_ticket}}).
  \end{itemize}
\item
  \textbf{III.10:}

  \begin{itemize}
  \tightlist
  \item
    the step-estimator's energy grant per seed
    (\href{https://github.com/yang-mills-group/OpenAI-YM-20261028/blob/3456e5e75cf8d4cd558955a5c7638fba2d68f464/YM/Next/FLUCTUATIONS/Grants.lean\#L24}{\texttt{Next.Fluctuations.grants}}),
    used for III.5.8(c) and III.5.11;
  \item
    the sector integrals and their reduction
    (\href{https://github.com/yang-mills-group/OpenAI-YM-20261028/blob/3456e5e75cf8d4cd558955a5c7638fba2d68f464/YM/Next/FLUCTUATIONS/Sector.lean\#L229-L259}{\texttt{Sector.peel}},
    \href{https://github.com/yang-mills-group/OpenAI-YM-20261028/blob/3456e5e75cf8d4cd558955a5c7638fba2d68f464/YM/Next/FLUCTUATIONS/Sector.lean\#L71-L72}{\texttt{Sector.prefactor}});
  \item
    the linearized block constraint
    \href{https://github.com/yang-mills-group/OpenAI-YM-20261028/blob/3456e5e75cf8d4cd558955a5c7638fba2d68f464/YM/Next/FLUCTUATIONS/Defs.lean\#L53-L55}{\texttt{zLin}};
  \item
    the growth of \(\beta_F\) and the uniform predictor bound used in
    the spectral gap estimate.
  \end{itemize}
\item
  \textbf{III.6:} a gauge-fixed Poincaré-type inequality
  (\href{https://github.com/yang-mills-group/OpenAI-YM-20261028/blob/3456e5e75cf8d4cd558955a5c7638fba2d68f464/YM/Work/F01.lean\#L63-L108}{\texttt{Work.F01.link\_bound}}).
  A field vanishing on the trunks is bounded on each fine bond by its
  plaquette curls in a box of fixed radius around the block, plus its
  block average. This is the source of the spectral gap \(r\) in
  III.5.14.
\item
  \textbf{III.2:} committed pairs of scales, the construction run and
  its energy hypothesis entering the uniform sector estimate.
\end{itemize}

\textbf{Standard mathematics.} The formal proofs establish the estimates
directly. References for the methods are:

\begin{itemize}
\tightlist
\item
  Mayer expansion, Ursell functions and the Penrose tree partition: R.
  Fernández and A. Procacci,
  \href{https://link.springer.com/article/10.1007/s00220-007-0279-2}{``Cluster
  expansion for abstract polymer models. New bounds from an old
  approach''}, \emph{Commun. Math. Phys.} \textbf{274} (2007), 123--140.
\item
  Polymer convergence criteria: R. Kotecký and D. Preiss,
  \href{https://doi.org/10.1007/BF01211762}{``Cluster expansion for
  abstract polymer models''}, \emph{Commun. Math. Phys.} \textbf{103}
  (1986), 491--498; C. Gruber and H. Kunz,
  \href{https://link.springer.com/article/10.1007/BF01651334}{``General
  properties of polymer systems''}, \emph{Commun. Math. Phys.}
  \textbf{22} (1971), 133--161.
\item
  Lattice-animal counting by spanning walks; the Schur test for matrix
  norms; Möbius inversion on subsets.
\item
  Exponential conjugation for inverse decay: J.-M. Combes and L. Thomas,
  \href{https://link.springer.com/article/10.1007/BF01646473}{``Asymptotic
  behaviour of eigenfunctions for multiparticle Schrödinger
  operators''}, \emph{Commun. Math. Phys.} \textbf{34} (1973), 251--270.
\item
  Large-field sectors and localized propagators belong to Bałaban's
  block-spin method. Here they are implemented through the explicit seed
  allocations and principal inverses stated above.
\end{itemize}

\textbf{Imported constants and assumptions.}

\begin{itemize}
\tightlist
\item
  The constants \(B=\lceil(\log x)^2\rceil\),
  \(t=\beta_F^{-1/2}e^{(\log\log\beta_F)^2}\) and \(\beta_F=x+b\log L\)
  (III.1).
\item
  The halo factor \(l\) is arbitrary in the packet estimates; a
  particular choice belongs to their application.
\item
  The block side \(L\ge n\) and the constant \(W\), chosen large
  (III.2).
\item
  Assumed hypotheses:

  \begin{itemize}
  \tightlist
  \item
    \(\mathrm{Funds}\) and \(\mathrm{Branches}\) (III.1);
  \item
    lawful representations (III.1);
  \item
    the point charge (discharged as in III.5.11);
  \item
    in the auxiliary Gaussian estimates, the coarse-side bound and flag
    equality, together with the compatible quadratic forms specified in
    Lemma 4.36.
  \end{itemize}
\end{itemize}

\subsubsection{3. Strategy}\label{strategy-7}

For a finite gas, inclusion--exclusion writes the partition function as
a sum over tuples. Splitting off the connected component of a
distinguished occurrence gives the coefficient recurrence
\((m+1)g_{m+1}=\sum_{k=0}^m(k+1)\ell_k g_{m-k}\). For the gas polynomial
\(P(t)=\sum_ng_nt^n\) and cluster series \(L(t)=\sum_m\ell_mt^{m+1}\),
this becomes \(P'=L'P\). A geometric bound on the cluster coefficients
justifies differentiation on an interval containing \([0,1]\); then
\(Pe^{-L}=1\). Absolute bounds allow passage to arbitrary label sets
through the directed family of finite subsets.

The bounds combine the tree-graph inequality with a sum over trees. A
parent support \(Y\) meets labels of total inflated weight at most
\(|Y|\gamma\le\mathrm{ball}_B\gamma\,\mathrm{cov}_B(Y)\). The reserve
exponent \(\delta\) absorbs each parent's number of children through
\(x^k\le k!e^x\). Substochastic kernels sum the labels; the sum of the
resulting child factorials over parent maps is at most \(4^mm!\). After
cancellation of the occurrence factorial, the ratio per added vertex is
\(4\,\mathrm{ball}_B\gamma/\delta\). Subadditivity of the covering
number transfers these bounds to cluster unions. The resulting estimates
retain site pins, tails, deleted regions and repeated forced polymers.

For seed packets, a spanning-walk count gives an entropy factor
\(K_{\rm step}^2\) per seed. The point charge supplies an energy gain
\(\beta\kappa t^2\) per seed. Since \(t=\beta_F^{-1/2}P(\beta_F)\) and
\(\beta\ge u\beta_F\), the guaranteed gain is \(\kappa uP(\beta_F)^2\).
At fixed \(L\), this dominates every cost
\(\mathrm{bill}^dP(\beta_F)^z\) with \(z<2\). Half the gain can be
retained as a mark for later estimates. The partition of unity over
global seed records then splits the block-spin pushforward while
preserving a common field, tag array and tree-gauge slice.

The deletion construction obtains another logarithm from the region
recurrence. Its negative-logarithm pieces form a union-convolution
series pinned at a chosen site. A weighted \(\ell^1\) bound closes
induction on the region and compares signed activities with the gas of
negative envelope activities. The final matrix estimates use a
transported flat linearization: a gauge-fixed Poincaré inequality and a
small-background comparison give a uniform spectral gap, while path
bounds give weighted row and column estimates. Exponential conjugation
yields inverse decay, and principal window inverses approximate the full
inverse with error \(O(e^{-cB})\).

\subsubsection{4. Steps}\label{steps-7}

\paragraph{4.1 Algebra of the finite
gas}\label{algebra-of-the-finite-gas}

In Lemmas 4.1--4.4 the label set \(\iota\) is finite, the weights are
real, and \(R\) is reflexive and symmetric. All vertex sets are finite.

\textbf{Lemma 4.1 (gas as a sum over tuples).} Let \(R\) be reflexive
and symmetric, and
\(g_n=\frac1{n!}\sum_{z\in\iota^n}\mathrm{all}(z)\prod_aw_{z_a}\), where
\(\mathrm{all}(z)=\sum_{G\subseteq G_R(z)}(-1)^{|G|}\). Then
\(g_n=\sum_{|s|=n,\ s\ \mathrm{compatible}}\prod_{i\in s}w_i\), and
\(\sum_ng_n=\Xi(w)\). \emph{Sketch:}

\begin{itemize}
\tightlist
\item
  \(\mathrm{all}(z)=\mathbf 1[G_R(z)=\emptyset]\), by the binomial
  identity.
\item
  Since \(R\) is reflexive, an edgeless tuple has distinct entries, and
  so orders a compatible set; each set arises \(n!\) times.
\end{itemize}

\emph{Lean:}
\href{https://github.com/yang-mills-group/OpenAI-YM-20261028/blob/3456e5e75cf8d4cd558955a5c7638fba2d68f464/YM/Push/ABSOLUTE_SUMS/Partition.lean\#L23-L28}{\texttt{Partition.allMass\_eq}},
\href{https://github.com/yang-mills-group/OpenAI-YM-20261028/blob/3456e5e75cf8d4cd558955a5c7638fba2d68f464/YM/Push/ABSOLUTE_SUMS/Gas.lean\#L69-L120}{\texttt{Gas.row}},
\href{https://github.com/yang-mills-group/OpenAI-YM-20261028/blob/3456e5e75cf8d4cd558955a5c7638fba2d68f464/YM/Push/ABSOLUTE_SUMS/Gas.lean\#L128-L137}{\texttt{Gas.sums}}.

\textbf{Lemma 4.2 (splitting off a component).} For any relation \(R\),
any tuple \(z\) on a finite vertex set \(\alpha\), and \(o\in\alpha\):
\[\mathrm{all}(z)=\sum_{s\ni o}\phi^T_R(z|_s)\,\mathrm{all}(z|_{\alpha\setminus s}).\]
\emph{Sketch:} a subgraph is the same thing as three pieces of data:

\begin{itemize}
\tightlist
\item
  the vertex set \(s\) of the component of \(o\);
\item
  a connected spanning graph on \(s\);
\item
  an arbitrary subgraph on the complement, with no edges between \(s\)
  and its complement.
\end{itemize}

The signs multiply. \emph{Lean:}
\href{https://github.com/yang-mills-group/OpenAI-YM-20261028/blob/3456e5e75cf8d4cd558955a5c7638fba2d68f464/YM/Push/ABSOLUTE_SUMS/Partition.lean\#L219-L237}{\texttt{Partition.partition}}.
\emph{(proof read; the bijection lemmas statement only)}

\textbf{Lemma 4.3 (coefficient recurrence).}
\((m+1)\,g_{m+1}=\sum_{k=0}^m(k+1)\,\ell_k\,g_{m-k}\). \emph{Sketch:}

\begin{itemize}
\tightlist
\item
  Sum Lemma 4.2 over all tuples on \(\{o\}\sqcup\{1,\dots,m\}\); the
  tuple sums over \(s\) and over its complement factor.
\item
  There are \(\binom mk\) sets \(s\) of size \(k+1\) containing \(o\).
\item
  Relabelling invariance of \(\phi^T\) finishes the count.
\end{itemize}

\emph{Lean:}
\href{https://github.com/yang-mills-group/OpenAI-YM-20261028/blob/3456e5e75cf8d4cd558955a5c7638fba2d68f464/YM/Push/ABSOLUTE_SUMS/Coef.lean\#L71-L76}{\texttt{Coef.unbased}},
\href{https://github.com/yang-mills-group/OpenAI-YM-20261028/blob/3456e5e75cf8d4cd558955a5c7638fba2d68f464/YM/Push/ABSOLUTE_SUMS/Coef.lean\#L79-L125}{\texttt{Coef.discrete}},
\href{https://github.com/yang-mills-group/OpenAI-YM-20261028/blob/3456e5e75cf8d4cd558955a5c7638fba2d68f464/YM/Push/ABSOLUTE_SUMS/Coef.lean\#L129-L138}{\texttt{Coef.recurrence}}.

\textbf{Lemma 4.4 (III.5.1).} Under the geometric coefficient bound of
III.5.1, \(\exp(\sum_m\ell_m)=\Xi(w)\). \emph{Sketch:}

\begin{itemize}
\tightlist
\item
  Let \(u=(1+q^{-1})/2>1\). On \((-u,u)\), the series
  \(L(t)=\sum\ell_mt^{m+1}\) can be differentiated term by term, because
  \(|(m+1)\ell_mt^m|\le C(m+1)(qu)^m\) is summable.
\item
  \(P(t)=\sum g_nt^n\) is a polynomial, since \(g_n=0\) for
  \(n>|\iota|\).
\item
  Lemma 4.3 and the Cauchy product give \(P'=L'P\). Hence \(H=Pe^{-L}\)
  has \(H'=0\) and \(H(0)=g_0=1\).
\item
  Evaluate at \(t=1\).
\end{itemize}

\emph{Lean:}
\href{https://github.com/yang-mills-group/OpenAI-YM-20261028/blob/3456e5e75cf8d4cd558955a5c7638fba2d68f464/YM/Push/ABSOLUTE_SUMS/FiniteLog.lean\#L66-L145}{\texttt{Gas.exact\_polynomial}}.

\paragraph{4.2 Tree-graph bound}\label{tree-graph-bound}

\textbf{Lemma 4.5 (Penrose partition; III.5.2).} The absolute
connected-graph sum is bounded by the number of spanning trees.
\emph{Sketch:}

\begin{itemize}
\tightlist
\item
  To each connected spanning \(G\) attach its breadth-first tree from
  the distinguished vertex, breaking ties by a fixed enumeration.
\item
  For a fixed tree \(p\), the graphs with tree \(p\) form a Boolean
  interval \([E(p),E(p)\cup F(p)]\). Here \(E(p)\) is the edge set of
  \(p\), and \(F(p)\) consists of the edges of \(T\) that neither change
  the breadth-first layers nor offer an earlier parent.
\item
  The alternating sum over an interval vanishes unless the interval is a
  point: flip one free edge.
\item
  Each fibre therefore contributes at most \(1\), and only trees of
  \(T\) contribute.
\end{itemize}

The bound follows without identifying which trees have nonzero fibres.
\emph{Lean:}
\href{https://github.com/yang-mills-group/OpenAI-YM-20261028/blob/3456e5e75cf8d4cd558955a5c7638fba2d68f464/YM/Push/ABSOLUTE_SUMS/Graph.lean\#L113-L148}{\texttt{Graph.pay}};
breadth-first construction in
\href{https://github.com/yang-mills-group/OpenAI-YM-20261028/blob/3456e5e75cf8d4cd558955a5c7638fba2d68f464/YM/Push/ABSOLUTE_SUMS/Breadth.lean}{\texttt{Breadth.lean}}.
\emph{(proof read; breadth-first construction statement only)}

\paragraph{4.3 Sums over trees}\label{sums-over-trees}

\textbf{Lemma 4.6 (substochastic sums along a forest).} Let
\(\alpha,\beta\) be finite, \(p:\alpha\to\alpha\sqcup\{o\}\) an acyclic
parent map, \(Y\) a fixed root state, and \(X_j\) the state associated
with \(j\in\beta\). If \(K(Y',j)\ge0\) and \(\sum_jK(Y',j)\le1\) for
every state \(Y'\), then
\(\sum_{z:\alpha\to\beta}\prod_iK(X_{z_{p(i)}},z_i)\le1\), with the
first argument interpreted as \(Y\) when \(p(i)=o\). \emph{Sketch:}
induction on the set of summed vertices. Sum out a vertex of maximal
rank first: it is a leaf, and its row sum is \(\le1\). \emph{Lean:}
\href{https://github.com/yang-mills-group/OpenAI-YM-20261028/blob/3456e5e75cf8d4cd558955a5c7638fba2d68f464/YM/Push/ABSOLUTE_SUMS/Parent.lean\#L44-L126}{\texttt{Trees.substochastic}}.

\textbf{Lemma 4.7 (fan-out factorials).}
\(\sum_{p:[m]\to[m]\sqcup\{o\}}\prod_v(\mathrm{card}\{i:p(i)=v\})!=\binom{2m}{m}m!\le4^mm!\).
\emph{Sketch:} with \(m+1\) possible parents fixed, extending a map on
\(n\) vertices multiplies the total factorial weight by \(n+m+1\).
Iterating from \(n=0\) to \(m\) gives \((2m)!/m!\), and
\(\binom{2m}{m}\le4^m\). \emph{Lean:}
\href{https://github.com/yang-mills-group/OpenAI-YM-20261028/blob/3456e5e75cf8d4cd558955a5c7638fba2d68f464/YM/Push/ABSOLUTE_SUMS/Parent.lean\#L191-L193}{\texttt{Trees.parent\_fanout}},
\href{https://github.com/yang-mills-group/OpenAI-YM-20261028/blob/3456e5e75cf8d4cd558955a5c7638fba2d68f464/YM/Push/ABSOLUTE_SUMS/Parent.lean\#L184-L189}{\texttt{fanout\_bound}}.

\textbf{Lemma 4.8 (trees rooted at a support).} Assume \(w\ge0\),
\(\gamma,\delta>0\), \(a\in\mathbb R\), and
\(\mathrm{Hit}_B(w;a+\delta,\gamma)\). For any support \(Y\) and any
finite \(D\subseteq\iota\),
\[\sum_{p}\sum_{z\in D^m}\mathbf 1\big[X_{z_i}\cap X_{\mathrm{parent}(i)}\neq\emptyset\ \forall i\big]\ \frac1{m!}\prod_i\big(w_{z_i}e^{a\,\mathrm{cov}_B(X_{z_i})}\big)\ \le\ e^{\delta\,\mathrm{cov}_B(Y)}\,\big(4\,\mathrm{ball}_B\gamma/\delta\big)^m,\]
where \(p\) runs over acyclic parent maps on the \(m\) labelled vertices
with external root of support \(Y\). \emph{Sketch:}

\begin{itemize}
\tightlist
\item
  Define
  \(K(Y',i)=\mathbf 1[X_i\cap Y'\neq\emptyset]\,w_ie^{(a+\delta)\mathrm{cov}_B(X_i)}\big/\big(\tau\max(\delta\,\mathrm{cov}_B(Y'),1)\big)\)
  with \(\tau=\mathrm{ball}_B\gamma/\delta\). The numerator summed over
  \(i\) is at most
  \(|Y'|\gamma\le\tau\delta\,\mathrm{cov}_B(Y')\le\tau\max(\delta\,\mathrm{cov}_B(Y'),1)\),
  so \(\sum_iK(Y',i)\le1\), also for \(Y'=\emptyset\).
\item
  Rewrite each weight \(w_ie^{a\,\mathrm{cov}}\) as
  \(\tau\max(\delta\,\mathrm{cov}(\mathrm{parent}),1)\,K\,e^{-\delta\,\mathrm{cov}(X_i)}\).
\item
  A parent with \(k\) children contributes
  \(\max(\delta\,\mathrm{cov},1)^k\le k!\,e^{\delta\,\mathrm{cov}}\).
  The exponential at each internal vertex cancels its own
  \(e^{-\delta\,\mathrm{cov}}\), leaving
  \(e^{\delta\,\mathrm{cov}_B(Y)}\) at the root.
\item
  Conclude with Lemmas 4.6 and 4.7.
\end{itemize}

\emph{Lean:}
\href{https://github.com/yang-mills-group/OpenAI-YM-20261028/blob/3456e5e75cf8d4cd558955a5c7638fba2d68f464/YM/Push/ABSOLUTE_SUMS/Weigh.lean\#L192-L228}{\texttt{Trees.trees\_afford}},
\href{https://github.com/yang-mills-group/OpenAI-YM-20261028/blob/3456e5e75cf8d4cd558955a5c7638fba2d68f464/YM/Push/ABSOLUTE_SUMS/Forest.lean\#L46-L56}{\texttt{Trees.forest\_at}}.

\paragraph{4.4 Cluster sums and the
logarithm}\label{cluster-sums-and-the-logarithm}

\textbf{Lemma 4.9 (rooted and marked cluster sums; III.5.3(c)).} Under
III.5.3's hypotheses, the marked sum at order \(m\) is at most \(bq^m\).
\emph{Sketch:}

\begin{itemize}
\tightlist
\item
  \(\frac1{(m+1)!}=\frac1{m+1}\cdot\frac1{m!}\): write the tuple sum as
  the average over the choice of root vertex. A transposition moves the
  root to position \(0\), and \(\phi^T\) is invariant under relabelling.
\item
  At root \(0\), Lemma 4.5 bounds \(|\phi^T(z)|\) by the trees of
  \(G_R(z)\). Since \(R\) is contained in the overlap relation, these
  are trees of supports rooted at \(X_{z_0}\) in the sense of Lemma 4.8.
\item
  The root label is summed against its own weight
  \(w_ie^{(a+\delta)\mathrm{cov}(X_i)}\), which absorbs the factor
  \(e^{\delta\,\mathrm{cov}}\). This gives \(bq^m\).
\end{itemize}

\emph{Lean:}
\href{https://github.com/yang-mills-group/OpenAI-YM-20261028/blob/3456e5e75cf8d4cd558955a5c7638fba2d68f464/YM/Push/ABSOLUTE_SUMS/Clusters.lean\#L148-L183}{\texttt{Ursell.cluster\_at}}.

\textbf{Lemma 4.10 (pinned bounds; III.5.3(a),(b)).} Under III.5.3's
hypotheses, both pinned tails are at most \(\gamma q^m/(1-q)\).
\emph{Sketch:}

\begin{itemize}
\tightlist
\item
  If \(\phi^T_R(z)\neq0\), then \(G_R(z)\) is connected, so
  \(\mathrm{cov}_B(X(z))\le\sum_a\mathrm{cov}_B(X_{z_a})\).
\item
  Apply Lemma 4.9 with \(Q=\{i:x\in X_i\}\), where \(b=\gamma\) by the
  site condition at \(x\), and sum the geometric tail.
\item
  The set version (b) follows in the same way from a count of
  distinct-label forests.
\end{itemize}

\emph{Lean:}
\href{https://github.com/yang-mills-group/OpenAI-YM-20261028/blob/3456e5e75cf8d4cd558955a5c7638fba2d68f464/YM/Push/ABSOLUTE_SUMS/Clusters.lean\#L194-L211}{\texttt{Ursell.connected\_pay}},
\href{https://github.com/yang-mills-group/OpenAI-YM-20261028/blob/3456e5e75cf8d4cd558955a5c7638fba2d68f464/YM/Push/ABSOLUTE_SUMS/Clusters.lean\#L224-L243}{\texttt{Ursell.log\_hit\_tail}},
\href{https://github.com/yang-mills-group/OpenAI-YM-20261028/blob/3456e5e75cf8d4cd558955a5c7638fba2d68f464/YM/Push/ABSOLUTE_SUMS/Distinct.lean\#L245-L276}{\texttt{Trees.touch\_tail}}.
\emph{(proof read; distinct-label forest count statement only)}

\textbf{Lemma 4.11 (from finite to arbitrary label sets; III.5.4).}
Under III.5.4's hypotheses, the gas and cluster series converge
absolutely, and \(\Xi=\exp(\text{cluster series})>0\). \emph{Sketch:}

\begin{itemize}
\tightlist
\item
  Every support is nonempty, so
  \(\sum_i|w_i|e^{(a+\delta)\mathrm{cov}}\le n^4\gamma\). Hence
  \(\sum_s\prod_{i\in s}|w_i|\le e^{n^4\gamma}\), and the cluster series
  has absolute sum \(\le n^4\gamma/(1-q)\).
\item
  On a finite \(D\subseteq\iota\) the hypotheses restrict, and III.5.1
  applies with \(C=n^4\gamma\).
\item
  Let the finite subsets \(D\subseteq\iota\) increase by inclusion and
  use dominated convergence for both series; the dominating budgets do
  not depend on \(D\).
\end{itemize}

\emph{Lean:}
\href{https://github.com/yang-mills-group/OpenAI-YM-20261028/blob/3456e5e75cf8d4cd558955a5c7638fba2d68f464/YM/Push/ABSOLUTE_SUMS/Log.lean\#L29-L41}{\texttt{Gas.finite\_log}},
\href{https://github.com/yang-mills-group/OpenAI-YM-20261028/blob/3456e5e75cf8d4cd558955a5c7638fba2d68f464/YM/Push/ABSOLUTE_SUMS/Log.lean\#L114-L149}{\texttt{Gas.price\_log}},
\href{https://github.com/yang-mills-group/OpenAI-YM-20261028/blob/3456e5e75cf8d4cd558955a5c7638fba2d68f464/YM/Push/ABSOLUTE_SUMS/Product.lean\#L43-L51}{\texttt{Product.everyone\_ne}},
\href{https://github.com/yang-mills-group/OpenAI-YM-20261028/blob/3456e5e75cf8d4cd558955a5c7638fba2d68f464/YM/Push/ABSOLUTE_SUMS/Product.lean\#L53-L64}{\texttt{Product.set\_products}}.

\paragraph{4.5 Deleted regions and forced
polymers}\label{deleted-regions-and-forced-polymers}

\textbf{Lemma 4.12 (III.5.5).} Under III.5.5's hypotheses, deleting
polymers meeting \(Y\) multiplies the gas by \(e^{-p}\), with
\(|p|\le|Y|\gamma/(1-q)\) and the stated cluster tails. \emph{Sketch:}

\begin{itemize}
\tightlist
\item
  The cluster terms of \(w^Y\) are those of \(w\) minus the terms whose
  union meets \(Y\).
\item
  The removed terms are bounded by Lemma 4.10 with \(m=0\), summed over
  \(x\in Y\).
\item
  Apply III.5.4 to both gases; \(|w^Y|\le|w|\), so \(w^Y\) satisfies the
  same hypotheses.
\end{itemize}

\emph{Lean:}
\href{https://github.com/yang-mills-group/OpenAI-YM-20261028/blob/3456e5e75cf8d4cd558955a5c7638fba2d68f464/YM/Push/ABSOLUTE_SUMS/Absence.lean\#L69-L95}{\texttt{Gas.restores\_log}},
\href{https://github.com/yang-mills-group/OpenAI-YM-20261028/blob/3456e5e75cf8d4cd558955a5c7638fba2d68f464/YM/Push/ABSOLUTE_SUMS/Absence.lean\#L40-L63}{\texttt{Gas.missing\_tail}}.

\textbf{Lemma 4.13 (multiplicities on a fixed set).} For a finite set
\(s\), if \(0\le\sigma<1\) and \(0\le w_i\le\sigma\) for \(i\in s\),
then
\(\sum_{(n_i)\in\mathbb N_{\ge1}^{s}}\prod_{i\in s}w_i^{n_i}/n_i!\le\prod_{i\in s}\big(w_i/(1-\sigma)\big)\).
\emph{Sketch:} bound \(w_i^{n_i}/n_i!\) by \(w_i\sigma^{n_i-1}\) and
multiply the geometric sums. \emph{Lean:}
\href{https://github.com/yang-mills-group/OpenAI-YM-20261028/blob/3456e5e75cf8d4cd558955a5c7638fba2d68f464/YM/Push/ABSOLUTE_SUMS/Copies.lean\#L60-L72}{\texttt{Multiple.buy\_at}}.

\textbf{Lemma 4.14 (multiset sums).} Let \(X_i\) be nonempty,
\(w_i\ge0\), \(a\ge0\), \(t\ge1\), \(0\le\sigma<1\), and
\(e_i\in\mathbb R\). Suppose
\(\mathrm{Ready}\big(t\,w\,e^{e}/(1-\sigma);a,\gamma,\delta,q\big)\) and
\(\gamma\le\sigma\). Then the sum over multisets with connected support
through \(x\), satisfying a predicate \(Q\) on which
\(\sum n_ie_i\ge v\), with total multiplicity \(\ge m\), of
\(e^{a\,\mathrm{cov}_B(\cup)}\prod w_i^{n_i}/n_i!\), is at most
\(t^{-m}e^{-v}\gamma/(1-q)\). \emph{Sketch:}

\begin{itemize}
\tightlist
\item
  Each weight \(tw_i e^{e_i}e^{a\,\mathrm{cov}_B(X_i)}\) is at most
  \(\sigma\), by the site condition at a point of its nonempty support.
\item
  Connectedness bounds the union cover by the sum of the individual
  covers; multiplicities are at least one. Lemma 4.13 then sums the
  multiplicities at cost \(1/(1-\sigma)\) per label, and III.5.3(b) sums
  the connected sets.
\item
  The factor \(t^{\sum n}e^{\sum ne}\ge t^me^v\) is extracted on the
  selected multisets.
\end{itemize}

\emph{Lean:}
\href{https://github.com/yang-mills-group/OpenAI-YM-20261028/blob/3456e5e75cf8d4cd558955a5c7638fba2d68f464/YM/Push/ABSOLUTE_SUMS/Copies.lean\#L135-L169}{\texttt{Multiple.estimates}}.

\textbf{Lemma 4.15 (III.5.6).} Under III.5.6's hypotheses, the pinned
multiset sum containing a forced label is at most
\(t^{-m}e^{-\mu}\eta/(1-q')\). \emph{Sketch:}

\begin{itemize}
\tightlist
\item
  The first-level clusters satisfy a site condition with constant
  \(\gamma/(1-q)\) at exponent \(A+\delta'\), by III.5.3(a) with
  \(m=0\); the forced polymers satisfy one with constant \(c\).
\item
  Join the two alphabets, scaling the weights by \(te^\mu/(1-\sigma)\)
  and \(t/(1-\sigma)\) respectively; the resulting constant is at most
  \(\eta\).
\item
  Apply Lemma 4.14 with \(e=\mu\) on forced labels; the presence of a
  forced label gives \(v=\mu\).
\end{itemize}

\emph{Lean:}
\href{https://github.com/yang-mills-group/OpenAI-YM-20261028/blob/3456e5e75cf8d4cd558955a5c7638fba2d68f464/YM/Push/ABSOLUTE_SUMS/Omissions.lean\#L48-L100}{\texttt{OmitClusters.restore}},
\href{https://github.com/yang-mills-group/OpenAI-YM-20261028/blob/3456e5e75cf8d4cd558955a5c7638fba2d68f464/YM/Push/ABSOLUTE_SUMS/Labels.lean\#L83-L96}{\texttt{Labels.together}}.

\paragraph{4.6 Packets and the large-coupling
numerics}\label{packets-and-the-large-coupling-numerics}

\textbf{Lemma 4.16 (animals and covers).} In the setting of III.5.7:

\begin{itemize}
\tightlist
\item
  For every \(k\in\mathbb N\), the number of connected packets with
  \(k\) seeds whose dome contains \(x\) is at most
  \(K_{\rm hit}(K_{\rm step}^2)^k\).
\item
  \(\mathrm{cov}_B(D(A))\le|A|\,(2l+3)^4\) for connected \(A\).
\end{itemize}

\emph{Sketch:}

\begin{itemize}
\tightlist
\item
  A connected packet is the vertex set of a walk of length \(2(k-1)\).
  The walk runs in the graph of seeds whose skins can meet, which has
  degree \(\le K_{\rm step}\). It starts at one of at most
  \(K_{\rm hit}\) seeds whose skin contains \(x\).
\item
  For \(B\ge1\), a skin lies in a ball of radius \(lB+1\le(l+1)B\),
  covered by \((2l+3)^4\) balls of radius \(B\). For \(B=0\), its
  connected location set lies in the radius-\(1\) ball and costs at most
  \(3^4\le(2l+3)^4\) singleton balls. A connected union costs at most
  the sum.
\end{itemize}

\emph{Lean:}
\href{https://github.com/yang-mills-group/OpenAI-YM-20261028/blob/3456e5e75cf8d4cd558955a5c7638fba2d68f464/YM/Push/ABSOLUTE_SUMS/Packet.lean\#L260-L269}{\texttt{Packets.stores\_count}},
\href{https://github.com/yang-mills-group/OpenAI-YM-20261028/blob/3456e5e75cf8d4cd558955a5c7638fba2d68f464/YM/Push/ABSOLUTE_SUMS/Animals.lean\#L122-L137}{\texttt{Animals.animals}},
\href{https://github.com/yang-mills-group/OpenAI-YM-20261028/blob/3456e5e75cf8d4cd558955a5c7638fba2d68f464/YM/Push/ABSOLUTE_SUMS/Packet.lean\#L171-L176}{\texttt{Packets.cover\_support}}.
\emph{(statement only for the animal count; covering proof read)}

\textbf{Lemma 4.17 (III.5.7(a)).} Under III.5.7(a)'s hypotheses, the
packet tail is at most \(K_{\rm hit}Cq^{\max(1,m)}/(1-q)\).
\emph{Sketch:}

\begin{itemize}
\tightlist
\item
  On a connected packet,
  \(p_Ae^{a\,\mathrm{cov}}\le C\big(e^{s-g+a(2l+3)^4}\big)^{|A|}\).
\item
  With the count of Lemma 4.16, sum the geometric series
  (\href{https://github.com/yang-mills-group/OpenAI-YM-20261028/blob/3456e5e75cf8d4cd558955a5c7638fba2d68f464/YM/Push/ABSOLUTE_SUMS/Basic.lean\#L143-L163}{\texttt{budget\_count\_tail}}).
\item
  Connected packets have at least one seed; hence \(\max(1,m)\).
\end{itemize}

\emph{Lean:}
\href{https://github.com/yang-mills-group/OpenAI-YM-20261028/blob/3456e5e75cf8d4cd558955a5c7638fba2d68f464/YM/Push/ABSOLUTE_SUMS/Charge.lean\#L66-L105}{\texttt{Packets.packet\_income}},
\href{https://github.com/yang-mills-group/OpenAI-YM-20261028/blob/3456e5e75cf8d4cd558955a5c7638fba2d68f464/YM/Push/ABSOLUTE_SUMS/Charge.lean\#L42-L53}{\texttt{Packets.cover\_income}}.

\textbf{Lemma 4.18 (III.5.7(b)).} Under III.5.7(b)'s hypotheses, the
mixed-word sum has the bound stated there. \emph{Sketch:}

\begin{itemize}
\tightlist
\item
  \emph{Ready.} Inflate the weights to
  \(p_Ae^{2\varepsilon e_s|A|+\mu_s|A|}\) and
  \(w_ie^{2\varepsilon\lambda_i}\). The joint alphabet then satisfies
  Ready at exponent \(a+2\varepsilon\): packets by Lemma 4.17, optional
  labels by hypothesis, and the constants add up to
  \(\gamma_{\rm root}\).
\item
  \emph{Connected sets.} III.5.3(b) sums the connected sets.
\item
  \emph{Extractions.}

  \begin{itemize}
  \tightlist
  \item
    \((1+\mathrm{ch})^k\le\varepsilon^{-k}k!\,e^{\varepsilon}e^{\varepsilon\,\mathrm{ch}}\);
  \item
    one of the two factors \(e^{\varepsilon\,\mathrm{ch}}\) of the
    inflation gives \(e^{\varepsilon T}\);
  \item
    the marks give \(e^{\mu_sr}\).
  \end{itemize}
\end{itemize}

\emph{Lean:}
\href{https://github.com/yang-mills-group/OpenAI-YM-20261028/blob/3456e5e75cf8d4cd558955a5c7638fba2d68f464/YM/Push/ABSOLUTE_SUMS/Mixed.lean\#L95-L122}{\texttt{Mixed.ready\_vertices}},
\href{https://github.com/yang-mills-group/OpenAI-YM-20261028/blob/3456e5e75cf8d4cd558955a5c7638fba2d68f464/YM/Push/ABSOLUTE_SUMS/Mixed.lean\#L153-L189}{\texttt{Mixed.mixed\_hit}},
\href{https://github.com/yang-mills-group/OpenAI-YM-20261028/blob/3456e5e75cf8d4cd558955a5c7638fba2d68f464/YM/Push/ABSOLUTE_SUMS/Reserves.lean\#L97-L135}{\texttt{Reserves.sets}},
\href{https://github.com/yang-mills-group/OpenAI-YM-20261028/blob/3456e5e75cf8d4cd558955a5c7638fba2d68f464/YM/Push/ABSOLUTE_SUMS/Reserves.lean\#L34-L56}{\texttt{Reserves.reserve\_tail}}.

\textbf{Lemma 4.19 (III.5.8).} For fixed numerical parameters satisfying
III.5.8, all its packet and tree conditions hold for sufficiently large
\(x\), uniformly in volume. \emph{Sketch:}

\begin{itemize}
\tightlist
\item
  \emph{Domination.} \(D\,\mathrm{bill}^d\le P(\beta_F)^e\) eventually,
  for every \(e>0\), since \((\log\beta)^{2d}=e^{2d\log\log\beta}\) is
  eventually below \(e^{e(\log\log\beta)^2}\). Also
  \(K_{\rm hit},K_{\rm step}^2\le e^{M/16}\), because
  \(K_{\rm step}\le\mathrm{const}\cdot\mathrm{bill}^4\).
\item
  \emph{Ratio.} \(\rho\le e^{M/16}e^{M/16-M+M/16}\le e^{-M/2}\).
\item
  \emph{Gain.} \(\beta\ge u\beta_F\) gives
  \(\beta\mu t^2\ge\mu uP(\beta_F)^2\), because
  \(\beta_Ft^2=P(\beta_F)^2\).
\item
  \emph{Assembly.} For (b), subtract the mark \(M/2\) from the gain and
  apply (a) with \(\mu/2\); the remaining gain is at least \(M/2\). For
  (c), use the step-estimator's grant before applying the numerical
  estimates.
\end{itemize}

\emph{Lean:}
\href{https://github.com/yang-mills-group/OpenAI-YM-20261028/blob/3456e5e75cf8d4cd558955a5c7638fba2d68f464/YM/Push/ABSOLUTE_SUMS/Delay.lean\#L74-L101}{\texttt{Delay.monitor\_pays}},
\href{https://github.com/yang-mills-group/OpenAI-YM-20261028/blob/3456e5e75cf8d4cd558955a5c7638fba2d68f464/YM/Push/ABSOLUTE_SUMS/Delay.lean\#L128-L156}{\texttt{Delay.step\_delay}},
\href{https://github.com/yang-mills-group/OpenAI-YM-20261028/blob/3456e5e75cf8d4cd558955a5c7638fba2d68f464/YM/Push/ABSOLUTE_SUMS/Delay.lean\#L165-L214}{\texttt{Delay.fits}},
\href{https://github.com/yang-mills-group/OpenAI-YM-20261028/blob/3456e5e75cf8d4cd558955a5c7638fba2d68f464/YM/Push/ABSOLUTE_SUMS/Buy.lean\#L67-L103}{\texttt{Buy.quoted}},
\href{https://github.com/yang-mills-group/OpenAI-YM-20261028/blob/3456e5e75cf8d4cd558955a5c7638fba2d68f464/YM/Push/ABSOLUTE_SUMS/Grants.lean\#L52-L71}{\texttt{Grant.order}}.
\emph{(proof read; the grant of III.10 statement only)}

\paragraph{4.7 The logarithm by
deletion}\label{the-logarithm-by-deletion}

\textbf{Lemma 4.20 (deletion identity).} For the gas \(Z_u\) of
III.5.9(a) and \(x\in B\):
\[Z_u(B)=Z_u(B\setminus\{x\})+\sum_{P\subseteq B,\ P\ni x}\varphi_u(P)\,Z_u(B\setminus P),\qquad \varphi_u(P)=\sum_{i:\,X_i=P}u_i,\]
and \(Z_u(\emptyset)=1\). \emph{Sketch:}

\begin{itemize}
\tightlist
\item
  At most one member of a disjoint family contains \(x\); removing it
  removes its support.
\item
  The identity holds exactly for finite truncations of \(\iota\) and
  passes to the absolutely convergent limit.
\end{itemize}

\emph{Lean:}
\href{https://github.com/yang-mills-group/OpenAI-YM-20261028/blob/3456e5e75cf8d4cd558955a5c7638fba2d68f464/YM/SectorNext/VACUUM_NORMALIZATION/Gas.lean\#L68-L114}{\texttt{ReLog.stock\_deletion}},
\href{https://github.com/yang-mills-group/OpenAI-YM-20261028/blob/3456e5e75cf8d4cd558955a5c7638fba2d68f464/YM/SectorNext/VACUUM_NORMALIZATION/Gas.lean\#L144-L163}{\texttt{ReLog.shelf}}.

\textbf{Lemma 4.21 (the incidence algebra).} Let \(J\) be finite and let
\(w\) satisfy the weight assumptions of III.5.9(a). For set functions on
\(J\) define:

\begin{itemize}
\tightlist
\item
  \((f*g)(T)=\sum_{P\cup Q=T}f(P)g(Q)\);
\item
  \(\hat f(B)=\sum_{A\subseteq B}f(A)\), so that
  \(\widehat{f*g}=\hat f\hat g\);
\item
  \(\|f\|_w=\sum_Aw(A)|f(A)|\).
\end{itemize}

Then:

\begin{itemize}
\tightlist
\item
  \(\|\cdot\|_w\) is submultiplicative on functions supported on sets
  containing a common point \(x\).
\item
  For such \(f\) with \(\|f\|_w<1\),
  \(\lambda(f)=\sum_{n\ge0}f^{*(n+1)}/(n+1)\) satisfies
  \(\widehat{\lambda(f)}=-\log(1-\hat f)\) and
  \(\|\lambda(f)\|_w\le-\log(1-\|f\|_w)\).
\item
  If \(f\) and \(F\) are supported on sets containing \(x\),
  \(|f|\le F\) pointwise and \(\|F\|_w<1\), then
  \(|\lambda(f)|\le\lambda(F)\).
\end{itemize}

\emph{Sketch:} cumulative sums turn union convolution into
multiplication. Two nonzero factors pinned at \(x\) have intersecting
supports, so the weight inequality gives submultiplicativity. Apply it
to convolution powers and sum \(\sum_{n\ge1}q^n/n=-\log(1-q)\) for
\(q<1\); termwise absolute domination gives the last assertion.

\emph{Lean:}
\href{https://github.com/yang-mills-group/OpenAI-YM-20261028/blob/3456e5e75cf8d4cd558955a5c7638fba2d68f464/YM/SectorNext/VACUUM_NORMALIZATION/Leaves.lean\#L118-L134}{\texttt{ReLog.crop\_glue}},
\href{https://github.com/yang-mills-group/OpenAI-YM-20261028/blob/3456e5e75cf8d4cd558955a5c7638fba2d68f464/YM/SectorNext/VACUUM_NORMALIZATION/Leaves.lean\#L169-L176}{\texttt{ReLog.lamp\_shed}},
\href{https://github.com/yang-mills-group/OpenAI-YM-20261028/blob/3456e5e75cf8d4cd558955a5c7638fba2d68f464/YM/SectorNext/VACUUM_NORMALIZATION/Leaves.lean\#L178-L196}{\texttt{ReLog.lamp\_cropped}},
\href{https://github.com/yang-mills-group/OpenAI-YM-20261028/blob/3456e5e75cf8d4cd558955a5c7638fba2d68f464/YM/SectorNext/VACUUM_NORMALIZATION/Leaves.lean\#L197-L209}{\texttt{ReLog.lamp\_mono}}.

\textbf{Lemma 4.22 (the induction).} Let \(J\) be finite, \(w\) satisfy
III.5.9(a)'s weight assumptions, and \(a\ge0\). Let \(U,z\) be functions
of supports with \(U\ge0\), \(|z|\le U\), and
\(\sum_{P\ni x}w(P)U(P)e^{a|P|}\le1-e^{-a}\) for every \(x\). Suppose
\(Z,Y\) satisfy the deletion identities with activities \(-U,-z\),
respectively, and \(Z(\emptyset)=Y(\emptyset)=1\). Then for every \(B\):

\begin{itemize}
\tightlist
\item
  \(Z(B)>0\) and \(Y(B)>0\);
\item
  \(|\lambda_Y(B)|\le\lambda_Z(B)\);
\item
  \(\sum_{A\subseteq B,\ A\ni x}w(A)\lambda_Z(A)\le a\) for \(x\in B\).
\end{itemize}

\emph{Sketch:} strong induction on \(B\). The empty region is immediate.
For nonempty \(B\), pick \(x\in B\).

\begin{itemize}
\item
  Define \(f_x\) on subsets of \(B\), extended by zero on other
  supports, by Möbius inversion from
  \[\hat f_x(A)=\sum_{P\subseteq A,\ P\ni x}U(P)\exp\!\left(\sum_{T\subseteq A\setminus\{x\},\ T\cap P\neq\emptyset}\lambda_Z(T)\right),\qquad A\subseteq B.\]
  The conditions \(x\in P\subseteq A\) make \(f_x\) supported on sets
  containing \(x\). The deletion identity gives
  \(Z(A)=Z(A\setminus\{x\})(1-\hat f_x(A))\).
\item
  By submultiplicativity of \(w\) and the induction hypothesis at the
  points of \(P\),
  \(\|f_x\|_w\le\sum_{P\ni x}w(P)U(P)e^{a|P|}\le1-e^{-a}\). Hence
  \(Z(B)>0\).
\item
  The pieces of \(-\log Z\) on sets containing \(x\) are
  \(\lambda(f_x)\), with \(\|\lambda(f_x)\|_w\le-\log e^{-a}=a\).
\item
  The same computation for \(Y\), with \(|f^Y_x|\le f_x\), gives the
  domination.
\end{itemize}

\emph{Lean:}
\href{https://github.com/yang-mills-group/OpenAI-YM-20261028/blob/3456e5e75cf8d4cd558955a5c7638fba2d68f464/YM/SectorNext/VACUUM_NORMALIZATION/Deletion.lean\#L180-L211}{\texttt{ReLog.leash}},
\href{https://github.com/yang-mills-group/OpenAI-YM-20261028/blob/3456e5e75cf8d4cd558955a5c7638fba2d68f464/YM/SectorNext/VACUUM_NORMALIZATION/Deletion.lean\#L75-L106}{\texttt{ReLog.matures}},
\href{https://github.com/yang-mills-group/OpenAI-YM-20261028/blob/3456e5e75cf8d4cd558955a5c7638fba2d68f464/YM/SectorNext/VACUUM_NORMALIZATION/Deletion.lean\#L224-L233}{\texttt{ReLog.swept}},
\href{https://github.com/yang-mills-group/OpenAI-YM-20261028/blob/3456e5e75cf8d4cd558955a5c7638fba2d68f464/YM/SectorNext/VACUUM_NORMALIZATION/Gas.lean\#L217-L228}{\texttt{ReLog.catalog}}.

\textbf{Lemma 4.23 (deleting sites).} Under III.5.9(a)'s hypotheses,
with \(Z=Z_u\) and \(\lambda=\lambda_u\),
\(-\log Z(J\setminus Q)=\sum_{A\cap Q=\emptyset}\lambda(A)\), hence
\(l=\sum_{A\cap Q\neq\emptyset}\lambda(A)\), and
\(|l|\le\sum_{x\in Q}\sum_{A\ni x}w(A)\lambda_{-v}(A)\le a|Q|\) because
\(w\ge1\). \emph{Lean:}
\href{https://github.com/yang-mills-group/OpenAI-YM-20261028/blob/3456e5e75cf8d4cd558955a5c7638fba2d68f464/YM/SectorNext/VACUUM_NORMALIZATION/Gas.lean\#L235-L271}{\texttt{ReLog.omitted}}.

\textbf{Lemma 4.24 (at a coarse field; III.5.9(b)).} Under III.5.9(b)'s
hypotheses, its masked logarithm coefficients satisfy the common
envelope and deletion-ratio bounds. \emph{Sketch:}

\begin{itemize}
\tightlist
\item
  \(e^{A\,\mathrm{cov}_M}\) is an admissible weight, since
  \(\mathrm{cov}_M\) is subadditive on overlapping unions.
\item
  \(a(2M+1)^4\le1\) gives \(a|X|\le\mathrm{cov}_M(X)\), which converts
  the condition into III.5.9(a).
\item
  \emph{Masking.} If \(P\) contains a flagged site \(y\), the masked gas
  does not see \(y\), so its Möbius piece at \(P\) vanishes.
\item
  The raw pieces depend only on labels supported in \(P\), so they are
  local in \(P\) and gauge invariant. Countability makes the gas sums
  Borel; finite Möbius inversion preserves this. Multiplying by the
  healthy-region indicator gives the masked coefficients, including the
  external flag tests.
\end{itemize}

\emph{Lean:}
\href{https://github.com/yang-mills-group/OpenAI-YM-20261028/blob/3456e5e75cf8d4cd558955a5c7638fba2d68f464/YM/SectorNext/VACUUM_NORMALIZATION/Bulk.lean\#L157-L183}{\texttt{ReLog.intake}},
\href{https://github.com/yang-mills-group/OpenAI-YM-20261028/blob/3456e5e75cf8d4cd558955a5c7638fba2d68f464/YM/SectorNext/VACUUM_NORMALIZATION/Bulk.lean\#L26-L32}{\texttt{ReLog.native\_rate}},
\href{https://github.com/yang-mills-group/OpenAI-YM-20261028/blob/3456e5e75cf8d4cd558955a5c7638fba2d68f464/YM/SectorNext/VACUUM_NORMALIZATION/Bulk.lean\#L51-L65}{\texttt{ReLog.permits}},
\href{https://github.com/yang-mills-group/OpenAI-YM-20261028/blob/3456e5e75cf8d4cd558955a5c7638fba2d68f464/YM/SectorNext/VACUUM_NORMALIZATION/Bulk.lean\#L118-L143}{\texttt{ReLog.retain\_mask}},
\href{https://github.com/yang-mills-group/OpenAI-YM-20261028/blob/3456e5e75cf8d4cd558955a5c7638fba2d68f464/YM/SectorNext/VACUUM_NORMALIZATION/Bulk.lean\#L104-L114}{\texttt{ReLog.retest\_meas}}.

\textbf{Lemma 4.25 (trace-log; III.5.10).} Under III.5.10's hypotheses,
the trace-log series converges absolutely to \(\log\det(I+S)\), and both
determinants are positive. \emph{Sketch:}

\begin{itemize}
\tightlist
\item
  \(AC=A^{1/2}(A^{1/2}CA^{1/2})A^{-1/2}\) is similar to a symmetric
  matrix.
\item
  Its eigenvalues \(\lambda\) are those of \(1+S\), and
  \(|\lambda-1|\le\|S\|_{\infty\to\infty}\le q\).
\item
  Hence
  \(\log\det(1+S)=\sum\log\lambda=\sum_n(-1)^{n+1}\mathrm{tr}\,S^n/n\),
  and \(\det C=\det(1+S)/\det A>0\).
\end{itemize}

\emph{Lean:}
\href{https://github.com/yang-mills-group/OpenAI-YM-20261028/blob/3456e5e75cf8d4cd558955a5c7638fba2d68f464/YM/SectorNext/VACUUM_NORMALIZATION/Round.lean\#L74-L93}{\texttt{Round.twins}},
\href{https://github.com/yang-mills-group/OpenAI-YM-20261028/blob/3456e5e75cf8d4cd558955a5c7638fba2d68f464/YM/SectorNext/VACUUM_NORMALIZATION/Round.lean\#L95-L144}{\texttt{Round.twin\_log}}.

\paragraph{4.8 The sector decomposition}\label{the-sector-decomposition}

\textbf{Lemma 4.26 (allocations and columns).} An allocation is the
partition of unity with global seed tests specified in III.5.11. It
satisfies \(0\le D_A\le1\); where the likelihood is positive,
\(D_A\ne0\) implies the local sector prerequisites.

\begin{itemize}
\tightlist
\item
  The hard allocation, the indicator of the actual booking, is one
  example.
\item
  A smooth decreasing cutoff \(\chi:\mathbb R\to[0,1]\), equal to \(1\)
  on \((-\infty,1]\) and \(0\) on \([4,\infty)\), gives another
  allocation when \(W,t>0\).
\item
  For every coarse \(V\) and \(\beta\in\mathbb R\), if the step
  satisfies its branch conditions and \(g\) is integrable and gauge
  invariant, then \(\sum_A\mathrm{out}_A(V)=\mathrm{push}_\beta g(V)\).
\end{itemize}

\emph{Sketch:} sum the allocation before integrating, then apply the
tree-gauge integration identity for \(g\) and the gauge-covariant block
kernel.

\emph{Lean:}
\href{https://github.com/yang-mills-group/OpenAI-YM-20261028/blob/3456e5e75cf8d4cd558955a5c7638fba2d68f464/YM/Push/SECTOR_MEASURE/Profile.lean\#L30-L36}{\texttt{Allocation}},
\href{https://github.com/yang-mills-group/OpenAI-YM-20261028/blob/3456e5e75cf8d4cd558955a5c7638fba2d68f464/YM/Push/SECTOR_MEASURE/Profile.lean\#L68-L78}{\texttt{Allocation.hard}},
\href{https://github.com/yang-mills-group/OpenAI-YM-20261028/blob/3456e5e75cf8d4cd558955a5c7638fba2d68f464/YM/Push/SECTOR_MEASURE/Profile.lean\#L95-L101}{\texttt{Chi.canonical}},
\href{https://github.com/yang-mills-group/OpenAI-YM-20261028/blob/3456e5e75cf8d4cd558955a5c7638fba2d68f464/YM/Push/SECTOR_MEASURE/Output.lean\#L102-L114}{\texttt{push\_columns}}.

\textbf{Lemma 4.27 (cut-out identity).} For a representation adapted to
the tube as in III.5.11, let \(\Delta_{\rm cut}(U,V)\) be the
orientation-averaged fine cut-row action minus its coarse counterpart.
Then
\[\text{kernel}\times\rho_R=\mathrm{sieve}(V)\cdot\mathrm{exposed}(U)\cdot\mathrm{kernel}\cdot e^{-\beta\,\Delta_{\rm cut}(U,V)} .\]

\begin{itemize}
\tightlist
\item
  In the input action, split the cut rows of the upper stage, averaged
  over the lattice symmetries, into two parts:

  \begin{itemize}
  \tightlist
  \item
    their coarse counterpart at \(V\), which gives the factor
    \(\mathrm{sieve}(V)\), taken out of the integral;
  \item
    the difference \(\Delta_{\rm cut}(U,V)\), which is kept with the
    block kernel.
  \end{itemize}
\item
  The rest of the action and the polymer gas form
  \(\mathrm{exposed}(U)\).
\end{itemize}

\emph{Lean:}
\href{https://github.com/yang-mills-group/OpenAI-YM-20261028/blob/3456e5e75cf8d4cd558955a5c7638fba2d68f464/YM/Push/SECTOR_MEASURE/Presented.lean\#L40-L56}{\texttt{Presented.expel}},
\href{https://github.com/yang-mills-group/OpenAI-YM-20261028/blob/3456e5e75cf8d4cd558955a5c7638fba2d68f464/YM/Push/SECTOR_MEASURE/Presented.lean\#L66-L78}{\texttt{Presented.use\_gap}}.

\textbf{Lemma 4.28 (energy factorization).} For a step satisfying its
branch conditions,
\[\mathrm{kernel}\cdot e^{-\beta\,\Delta_{\rm cut}}=\pi_\beta\,e^{-\beta\,\overline{\mathcal E}_A}\cdot(\text{likelihood})\,e^{-\beta\,\overline{\mathcal E}^{\rm rest}_A},\]
where bars denote averages over the lattice symmetries, applied to the
rotated packet. On the tree-gauge slice, the global seed tests and
positive likelihood imply the local prerequisites in every orientation.
The point charge therefore gives
\(\overline{\mathcal E}_A\ge\kappa t^2|A|\) wherever \(D_A\ne0\) and the
likelihood is positive; zero-likelihood terms vanish. \emph{Lean:}
\href{https://github.com/yang-mills-group/OpenAI-YM-20261028/blob/3456e5e75cf8d4cd558955a5c7638fba2d68f464/YM/Push/SECTOR_MEASURE/Oriented.lean\#L185-L208}{\texttt{Oriented.enters\_eq}},
\href{https://github.com/yang-mills-group/OpenAI-YM-20261028/blob/3456e5e75cf8d4cd558955a5c7638fba2d68f464/YM/Push/SECTOR_MEASURE/Oriented.lean\#L210-L219}{\texttt{Oriented.mean\_pays}}.

\textbf{Lemma 4.29 (III.5.11).} Under III.5.11's funding, adaptation and
point-charge assumptions, the sector outputs sum to the pushforward and
satisfy its per-seed bound. \emph{Sketch:}

\begin{itemize}
\tightlist
\item
  Express each signed sector output using Lemma 4.27; its sum is the
  pushforward by Lemma 4.26.
\item
  Bound its absolute value by the integral of the absolute exposed
  density. Apply Lemma 4.28 with \(\beta\ge0\) to extract
  \(\pi_\beta e^{-\beta\kappa t^2|A|}\).
\item
  Bounded Borel factors give bounded Borel sector integrals at each
  fixed volume.
\end{itemize}

\emph{Lean:}
\href{https://github.com/yang-mills-group/OpenAI-YM-20261028/blob/3456e5e75cf8d4cd558955a5c7638fba2d68f464/YM/Push/SECTOR_MEASURE/Oriented.lean\#L267-L297}{\texttt{Oriented.turn\_charge}},
\href{https://github.com/yang-mills-group/OpenAI-YM-20261028/blob/3456e5e75cf8d4cd558955a5c7638fba2d68f464/YM/Push/SECTOR_MEASURE/Presented.lean\#L157-L210}{\texttt{Presented.portionCost}}.

\textbf{Lemma 4.30 (one \(\kappa\)).} The uniform point-charge
conclusions of III.5.11 hold. \emph{Sketch:} the general conclusion uses
the step-estimator's energy grant (III.10). At a committed pair of
scales, use the run's energy hypothesis (III.1); its compatibility
conditions identify the representation's paths, clip and threshold with
those of the tube. The constant is chosen before the volume, step and
representation.

\emph{Lean:}
\href{https://github.com/yang-mills-group/OpenAI-YM-20261028/blob/3456e5e75cf8d4cd558955a5c7638fba2d68f464/YM/Push/SECTOR_MEASURE/Oriented.lean\#L37-L51}{\texttt{points\_grant}},
\href{https://github.com/yang-mills-group/OpenAI-YM-20261028/blob/3456e5e75cf8d4cd558955a5c7638fba2d68f464/YM/Push/SECTOR_MEASURE/Consumer.lean\#L42-L65}{\texttt{Consumer.bought}},
\href{https://github.com/yang-mills-group/OpenAI-YM-20261028/blob/3456e5e75cf8d4cd558955a5c7638fba2d68f464/YM/Push/SECTOR_MEASURE/Grant.lean\#L61-L66}{\texttt{charged\_at}}.
\emph{(proof read; the inputs from III.1 and III.10 statement only)}

\paragraph{4.9 Matrices}\label{matrices}

\textbf{Lemma 4.31 (III.5.12).} Under III.5.12's hypotheses, the shifted
inverse has the stated pointwise exponential bound. \emph{Sketch:} fix
\(j\) and put \(s_k=e^{\theta\eta d(a_k,a_j)}\) and
\(w=s\cdot(M+uI)^{-1}e_j\).

\begin{itemize}
\tightlist
\item
  Then \((M+uI+F)w=e_j\), with \(F_{kl}=M_{kl}(s_k/s_l-1)\).
\item
  By convexity of \(\exp\),
  \(|s_k/s_l-1|\le\theta e^{\eta d(a_k,a_l)}\). The Schur test then
  gives \(|\langle w,Fw\rangle|\le\theta C|w|^2\le\frac r2|w|^2\).
\item
  Hence \(\min(1,r/2)(1+u)|w|^2\le w_j\le|w|\), and \(|w_i|\le|w|\)
  gives the claim.
\end{itemize}

\emph{Lean:}
\href{https://github.com/yang-mills-group/OpenAI-YM-20261028/blob/3456e5e75cf8d4cd558955a5c7638fba2d68f464/YM/Progress/STATIONARY_SCENES/Green.lean\#L188-L251}{\texttt{Green.inverse}},
\href{https://github.com/yang-mills-group/OpenAI-YM-20261028/blob/3456e5e75cf8d4cd558955a5c7638fba2d68f464/YM/Progress/STATIONARY_SCENES/Green.lean\#L135-L147}{\texttt{Green.diminishing}},
\href{https://github.com/yang-mills-group/OpenAI-YM-20261028/blob/3456e5e75cf8d4cd558955a5c7638fba2d68f464/YM/Progress/STATIONARY_SCENES/Green.lean\#L149-L168}{\texttt{Green.schur}}.

\textbf{Lemma 4.32 (III.5.13).} Under III.5.13's hypotheses, the window
residual, inverse difference and covariance satisfy its weighted bounds.
\emph{Sketch:}

\begin{itemize}
\tightlist
\item
  \(E_{ij}=0\) for \(i\in P(j)\), because column \(j\) of
  \(\mathrm{within}(A)\) solves the system on the window. So \(E\) lives
  at distances \(>R\).
\item
  Pointwise decay at rate \(\theta\eta\), from III.5.12 applied to \(A\)
  and to its principal submatrices, becomes row-sum decay at rate
  \(\nu\) by the counting bound. This costs \(\Pi_U(\theta\eta-\nu)\).
\item
  The far entries of \(E\) at the slower rate \(\mu\) pay
  \(e^{-(\nu-\mu)R}\).
\item
  Finally \(\mathrm{within}-A^{-1}=A^{-1}E\), and the weighted Schur
  classes are closed under products.
\item
  Part (b) applies (a) uniformly along \(tA+(1-t)\).
\end{itemize}

\emph{Lean:}
\href{https://github.com/yang-mills-group/OpenAI-YM-20261028/blob/3456e5e75cf8d4cd558955a5c7638fba2d68f464/YM/AnalyticNext/CONNECTED_ACTIVITY_BOUNDS/Windows.lean\#L51-L69}{\texttt{Windows.stitched}},
\href{https://github.com/yang-mills-group/OpenAI-YM-20261028/blob/3456e5e75cf8d4cd558955a5c7638fba2d68f464/YM/AnalyticNext/CONNECTED_ACTIVITY_BOUNDS/Windows.lean\#L124-L141}{\texttt{Windows.inverse\_price}},
\href{https://github.com/yang-mills-group/OpenAI-YM-20261028/blob/3456e5e75cf8d4cd558955a5c7638fba2d68f464/YM/AnalyticNext/CONNECTED_ACTIVITY_BOUNDS/Windows.lean\#L146-L169}{\texttt{Windows.window\_price}},
\href{https://github.com/yang-mills-group/OpenAI-YM-20261028/blob/3456e5e75cf8d4cd558955a5c7638fba2d68f464/YM/AnalyticNext/CONNECTED_ACTIVITY_BOUNDS/Schur.lean\#L69-L70}{\texttt{Win.planePrice}}.

\textbf{Lemma 4.33 (spectral gap).} Fix \(\mathcal T,\mathcal D\) as in
III.5.14. There is \(0<r\le1\) such that, for all sufficiently large
\(x\), all funded steps with this patch, coarse side \(\ge100(1+B)\) and
the stated flag equality satisfy \(r|v|^2\le|Dv|^2\) for every \(v\)
supported on active coordinates. \emph{Sketch:}

\begin{itemize}
\tightlist
\item
  Take an active bond. Its projected endpoints lie outside
  \(D_{10B}(A)\), so in an axial gauge the background \(V\) is within
  \(\delta(x)\to0\) of \(1\) near the bond.
\item
  The Poincaré-type inequality of III.6 bounds the bond variable by flat
  curls and block averages near the block. These are the flat rows.
\item
  In this gauge, the transported flat rows differ from the untransported
  flat rows by \(O(\delta)|v|_{\rm loc}\).
\item
  Square and sum with bounded overlap:
  \(|v|^2\le\mathcal V|Dv|^2+o(1)|v|^2\), which gives
  \(r=(2(1+\mathcal V))^{-1}\).
\item
  Then
  \(\langle v,\mathrm{gram}\,v\rangle=|Dv|^2+|v_{\rm inactive}|^2\ge r|v|^2\).
\end{itemize}

\emph{Lean:}
\href{https://github.com/yang-mills-group/OpenAI-YM-20261028/blob/3456e5e75cf8d4cd558955a5c7638fba2d68f464/YM/AnalyticNext/CONNECTED_ACTIVITY_BOUNDS/Bedrock.lean\#L27-L69}{\texttt{Aff.inflow}},
\href{https://github.com/yang-mills-group/OpenAI-YM-20261028/blob/3456e5e75cf8d4cd558955a5c7638fba2d68f464/YM/AnalyticNext/CONNECTED_ACTIVITY_BOUNDS/Bedrock.lean\#L72-L113}{\texttt{Aff.surplus}},
\href{https://github.com/yang-mills-group/OpenAI-YM-20261028/blob/3456e5e75cf8d4cd558955a5c7638fba2d68f464/YM/AnalyticNext/CONNECTED_ACTIVITY_BOUNDS/Bedrock.lean\#L134-L139}{\texttt{Aff.decomposes}}.
\emph{(proof read; per-bond comparison and III.6 inequality statement
only)}

\textbf{Lemma 4.34 (decay).} Fix \(\mathcal T,\mathcal D\) as in
III.5.14. There is \(K>0\), uniform over all funded steps with that
patch, all \(V,A,F\) and every real \(x\), such that \(D\) and its
unmasked version have weighted row and column bounds \(K\) at rate
\(\alpha\) equal to half the upper path-decay rate, and
\(\mathrm{gram}\) has bounds \(K^2+1\). Moreover
\(\mathrm{gram}_{ij}=0\) at distances greater than
\(2(\lfloor B/100000\rfloor+R_{\rm pred}(\mathcal D)+2)\).
\emph{Sketch:}

\begin{itemize}
\tightlist
\item
  Pointwise exponential decay of the row coefficients comes from the
  path-sum bounds of the upper stage and from the block-constraint rows.
\item
  The counting bound converts it into row sums.
\item
  Products and sums give \(\mathrm{gram}\).
\end{itemize}

\emph{Lean:}
\href{https://github.com/yang-mills-group/OpenAI-YM-20261028/blob/3456e5e75cf8d4cd558955a5c7638fba2d68f464/YM/AnalyticNext/CONNECTED_ACTIVITY_BOUNDS/IndexCost.lean\#L165-L181}{\texttt{Aff.gram\_budget}},
\href{https://github.com/yang-mills-group/OpenAI-YM-20261028/blob/3456e5e75cf8d4cd558955a5c7638fba2d68f464/YM/AnalyticNext/CONNECTED_ACTIVITY_BOUNDS/IndexCost.lean\#L120-L159}{\texttt{Aff.stencil}},
\href{https://github.com/yang-mills-group/OpenAI-YM-20261028/blob/3456e5e75cf8d4cd558955a5c7638fba2d68f464/YM/AnalyticNext/CONNECTED_ACTIVITY_BOUNDS/Speed.lean\#L138-L155}{\texttt{Aff.received}}.

\textbf{Lemma 4.35 (III.5.14).} The uniform matrix and window
conclusions of III.5.14 hold. \emph{Sketch:} combine Lemmas 4.33 and
4.34 with the index-counting bound, then apply III.5.13(b). The windows
have room \(\lfloor B/4\rfloor\to\infty\) as \(x\to\infty\).

\emph{Lean:}
\href{https://github.com/yang-mills-group/OpenAI-YM-20261028/blob/3456e5e75cf8d4cd558955a5c7638fba2d68f464/YM/AnalyticNext/CONNECTED_ACTIVITY_BOUNDS/Bedrock.lean\#L142-L158}{\texttt{Aff.bedrock}},
\href{https://github.com/yang-mills-group/OpenAI-YM-20261028/blob/3456e5e75cf8d4cd558955a5c7638fba2d68f464/YM/AnalyticNext/CONNECTED_ACTIVITY_BOUNDS/Warsaw.lean\#L25-L42}{\texttt{Aff.install\_promises}}.

\textbf{Lemma 4.36 (auxiliary Gaussian estimates).} Fix the
step-estimator and patch of III.5.14. The auxiliary Gaussian
(\href{https://github.com/yang-mills-group/OpenAI-YM-20261028/blob/3456e5e75cf8d4cd558955a5c7638fba2d68f464/YM/AnalyticNext/CONNECTED_ACTIVITY_BOUNDS/Mint.lean\#L22-L28}{\texttt{Mint}})
uses the completed precision \(G=D^TD+I_{\rm inactive}\), the affine
constant \(h\) (including the clipped frozen coordinates), source
\(D^Th\), and selected mean \(\mu=-\mathrm{within}(G)^TD^Th\).

For the field estimates below, assume a funded step with this patch,
coarse side \(\ge100(1+B)\), and the flag equality of III.5.14. Write
\(H_{n,p}=B^pt^n\). A \emph{reserve} \(C_{\rm res}>0\) supplies uniform
bounds on the upper and lower row sums and the coarse predicted fields.
A \emph{compatible pair of quadratic forms} means upper and lower forms
attached to the tube, satisfying its diagonal and renormalization
identities, the two quadratic bounds and orientation covariance. The
shift and plateau estimates assume such a pair and the window bounds of
III.5.13(b). Each large-\(x\) threshold is uniform in the remaining
volume and field data.

\begin{itemize}
\tightlist
\item
  \href{https://github.com/yang-mills-group/OpenAI-YM-20261028/blob/3456e5e75cf8d4cd558955a5c7638fba2d68f464/YM/AnalyticNext/CONNECTED_ACTIVITY_BOUNDS/Station.lean\#L222-L241}{\texttt{Aff.escape}}:
  With a compatible pair, for every coarse Lie-algebra field \(b\) and
  every tree-gauge forest, the stationary linear extension solves the
  untruncated flat fine-row and block-constraint equations with the
  predicted coarse linear rows as right-hand sides, in every
  orientation. This is an identity for those linear operators; it does
  not assert that \(D\mu+h=0\).
\item
  \href{https://github.com/yang-mills-group/OpenAI-YM-20261028/blob/3456e5e75cf8d4cd558955a5c7638fba2d68f464/YM/AnalyticNext/CONNECTED_ACTIVITY_BOUNDS/Stint.lean\#L44-L97}{\texttt{Aff.velocity}},
  \href{https://github.com/yang-mills-group/OpenAI-YM-20261028/blob/3456e5e75cf8d4cd558955a5c7638fba2d68f464/YM/AnalyticNext/CONNECTED_ACTIVITY_BOUNDS/Stint.lean\#L100-L121}{\texttt{Aff.shifted}}:
  Assume a reserve. Let \(o\notin D_{20B}(A)\) and let a gauge transform
  of \(V\) agree, on bonds anchored within \(6B\) of \(o\), with
  \(\exp b\), where \(\|b\|_\infty\le H_{1,7}\). Let \(v\) be the
  inverse-gauge transport of the stationary extension minus the replica
  extension, restricted to active coordinates. Then
  \(|\mu_i-v_i|\le H_{2,27}\) whenever the column position is within
  \(2B\) of \(o\). At every row centred outside \(D_{20B}(A)\),
  \(|(D\mu+h)_i|\le H_{2,29}\).
\item
  \href{https://github.com/yang-mills-group/OpenAI-YM-20261028/blob/3456e5e75cf8d4cd558955a5c7638fba2d68f464/YM/AnalyticNext/CONNECTED_ACTIVITY_BOUNDS/Hour.lean\#L55-L116}{\texttt{Jewels.night}}:
  Assume a reserve \(C_{\rm res}\), \(W>0\), and the two numerical
  margins \(C_{\rm res}/L^2\le s_0\) and
  \(C_{\rm res}+s_0\le(W/4)\sqrt{p_0}\), where
  \(s_0=\sqrt{c_{\rm direct}}/64\), \(c_{\rm direct}\) is the upper
  direct quadratic coefficient, and \(p_0>0\) is the kernel penalty.
  There is \(\varepsilon>0\) such that, for all sufficiently large
  \(x\), the following holds when the frozen cutoff scale equals \(Bt\).
  Reconstruct the fine field by inserting \(\exp(\mu+z)\) relative to
  the replica of \(V\) on active bonds and keeping the fixed tree-gauge
  slice field on the other bonds. For any fluctuation \(z\), a fine cell
  whose projected anchor lies outside \(D_{20B}(A)\) has plaquette chord
  \(<t/4\) if \(|z_i|\le\varepsilon t\) on coordinates anchored within
  \(B\) of that anchor. The same test at the anchor of a coarse bond
  outside \(D_{20B}(A)\) gives wire chord \(<Wt\) and makes both
  likelihood cutoffs equal to \(1\).
\item
  \href{https://github.com/yang-mills-group/OpenAI-YM-20261028/blob/3456e5e75cf8d4cd558955a5c7638fba2d68f464/YM/AnalyticNext/CONNECTED_ACTIVITY_BOUNDS/Draught.lean\#L92-L114}{\texttt{Aff.tap\_open}}:
  Fix an order \(p\in\mathbb N\) and \(\delta>0\). For all sufficiently
  large \(x\), if every fine bond anchored within \(B/2\) of a given
  coarse-bond anchor has chord at most \(H_{1,p}\), then the frame and
  connector use their minimizing branches and the predicted coarse-bond
  chord is \(<\delta\), for every tag. The uniqueness of those
  minimizers is part of the branch hypotheses.
\item
  \href{https://github.com/yang-mills-group/OpenAI-YM-20261028/blob/3456e5e75cf8d4cd558955a5c7638fba2d68f464/YM/AnalyticNext/CONNECTED_ACTIVITY_BOUNDS/DCompare.lean\#L190-L207}{\texttt{Aff.perturb}}:
  There is \(K>0\), depending only on the step-estimator and patch, such
  that for every truncation radius \(R\) and every unmasked transported
  row \(L_{R,V}\), with \(s_R=R+R_{\rm pred}(\mathcal D)+2\),
  \(\|(L_{R,V}-L_{R,1})u\|\le K(4s_R+1)\delta C_u\) if
  \(C_u,\delta\ge0\), \(\|u_j\|\le C_u\) on fine bonds anchored within
  \(s_R+1\) of that row, and all coarse-bond chords anchored within
  \(4s_R+1\) are at most \(\delta\). This compares with the flat
  background; it is not a Lipschitz estimate between arbitrary
  backgrounds.
\item
  \href{https://github.com/yang-mills-group/OpenAI-YM-20261028/blob/3456e5e75cf8d4cd558955a5c7638fba2d68f464/YM/AnalyticNext/CONNECTED_ACTIVITY_BOUNDS/Ground.lean\#L51-L130}{\texttt{Aff.arc\_bill}}:
  Assume a reserve and fix \(p\ge7\). For all sufficiently large \(x\),
  take a row centred outside \(D_{2B}(A)\) and a fine Lie-algebra field
  \(u\) bounded by \(H_{1,p}\) on bonds anchored within \(B\). If \(U\)
  agrees there with \(\exp(u)\,\mathrm{replica}(V)\), its nonlinear row
  value \(F(U,V)\) satisfies \(\|F(U,V)\|\le H_{1,p+4}\) and
  \(\|F(U,V)-F(\mathrm{replica}(V),V)-L_Vu\|\le H_{2,2p+6}\), where
  \(L_V\) is the transported flat row at the prescribed truncation. This
  is the error estimate used in place of identifying \(L_V\) with a
  derivative at \(V\).
\item
  \href{https://github.com/yang-mills-group/OpenAI-YM-20261028/blob/3456e5e75cf8d4cd558955a5c7638fba2d68f464/YM/AnalyticNext/CONNECTED_ACTIVITY_BOUNDS/Veil.lean\#L41-L49}{\texttt{Aff.reserves}}:
  For every step-estimator there exists a reserve \(C_{\rm res}\). It is
  chosen before \(L\); for each fixed \(L\), its row and coarse-field
  bounds hold for all sufficiently large \(x\) and all funded steps of
  coarse side \(\ge100(1+B)\).
\end{itemize}

\subsubsection{5. Where the difficulty
lies}\label{where-the-difficulty-lies-7}

The tree-graph argument is standard. The estimates must also retain
arbitrary label sets, repeated occurrences, site pins and covering
weights at the growing scale \(B=\lceil(\log x)^2\rceil\). The entropy
factor \(4(2B+1)^4\) is sufficient because it grows only polynomially in
\(\log x\).

The gain scale is \(P(\beta_F)^2=e^{2(\log\log\beta_F)^2}\). It absorbs
\(\mathrm{bill}^dP(\beta_F)^z\) for fixed \(d\) and \(z<2\), but grows
more slowly than every positive power of \(\beta_F\). The deletion route
gives support-indexed logarithm coefficients and a common
field-independent envelope; the graph route also gives occurrence tails
and multiset bounds. Both give deletion ratios uniform in volume,
although the full logarithm is extensive.

The sector estimate needs global seed records to preserve the local
point-charge prerequisites under lattice orientations. Its remainder
\(\ell_A\) has no volume-uniform bound here. For the Gaussian estimates,
the substantive inputs are the gauge-fixed Poincaré inequality and
control of the transported flat rows away from the flagged region.
Finite range grows with \(B\); uniform exponential row bounds provide
the useful control.

\subsubsection{6. Notation}\label{notation-7}

{\def\LTcaptype{none} % do not increment counter
\begin{longtable}[]{@{}
  >{\raggedright\arraybackslash}p{(\linewidth - 4\tabcolsep) * \real{0.3333}}
  >{\raggedright\arraybackslash}p{(\linewidth - 4\tabcolsep) * \real{0.3333}}
  >{\raggedright\arraybackslash}p{(\linewidth - 4\tabcolsep) * \real{0.3333}}@{}}
\toprule\noalign{}
\begin{minipage}[b]{\linewidth}\raggedright
Written proof
\end{minipage} & \begin{minipage}[b]{\linewidth}\raggedright
Lean
\end{minipage} & \begin{minipage}[b]{\linewidth}\raggedright
Meaning
\end{minipage} \\
\midrule\noalign{}
\endhead
\bottomrule\noalign{}
\endlastfoot
\(\mathrm{SB}(f;C)\) &
\href{https://github.com/yang-mills-group/OpenAI-YM-20261028/blob/3456e5e75cf8d4cd558955a5c7638fba2d68f464/YM/RG/Activities.lean\#L59-L60}{\texttt{SumBudget}}
& \(f\ge0\) summable, \(\sum f\le C\) \\
\(\mathrm{Hit}_B(w;a,\gamma)\) &
\href{https://github.com/yang-mills-group/OpenAI-YM-20261028/blob/3456e5e75cf8d4cd558955a5c7638fba2d68f464/YM/Ac/Tree.lean\#L12-L13}{\texttt{Ac.Hit}}
& site-sum condition \\
\(\mathrm{cov}_B\), \(\mathrm{ball}_B\) &
\href{https://github.com/yang-mills-group/OpenAI-YM-20261028/blob/3456e5e75cf8d4cd558955a5c7638fba2d68f464/YM/Lattice/Support.lean\#L69-L71}{\texttt{coverNumber}},
\href{https://github.com/yang-mills-group/OpenAI-YM-20261028/blob/3456e5e75cf8d4cd558955a5c7638fba2d68f464/YM/Push/ABSOLUTE_SUMS/Animals.lean\#L148}{\texttt{ballCost}}
& connected covering number; \((2B+1)^4\) \\
\(\mathrm{Ready}(w;a,\gamma,\delta,q)\) &
\href{https://github.com/yang-mills-group/OpenAI-YM-20261028/blob/3456e5e75cf8d4cd558955a5c7638fba2d68f464/YM/Push/ABSOLUTE_SUMS/Forest.lean\#L31-L32}{\texttt{Trees.Ready}},
\href{https://github.com/yang-mills-group/OpenAI-YM-20261028/blob/3456e5e75cf8d4cd558955a5c7638fba2d68f464/YM/Push/ABSOLUTE_SUMS/Weigh.lean\#L52}{\texttt{degreePrice}}
& smallness of the cluster expansion \\
\(R\subseteq\) overlap &
\href{https://github.com/yang-mills-group/OpenAI-YM-20261028/blob/3456e5e75cf8d4cd558955a5c7638fba2d68f464/YM/Push/ABSOLUTE_SUMS/Clusters.lean\#L59}{\texttt{Ursell.Tests}},
\href{https://github.com/yang-mills-group/OpenAI-YM-20261028/blob/3456e5e75cf8d4cd558955a5c7638fba2d68f464/YM/Push/ABSOLUTE_SUMS/Ursell.lean\#L22}{\texttt{blocked}}
& incompatibility implies overlap; the overlap relation \\
\(G_R(z)\), \(\phi^T_R(z)\) &
\href{https://github.com/yang-mills-group/OpenAI-YM-20261028/blob/3456e5e75cf8d4cd558955a5c7638fba2d68f464/YM/Push/ABSOLUTE_SUMS/Ursell.lean\#L32-L33}{\texttt{Ursell.pot}},
\href{https://github.com/yang-mills-group/OpenAI-YM-20261028/blob/3456e5e75cf8d4cd558955a5c7638fba2d68f464/YM/Push/ABSOLUTE_SUMS/Ursell.lean\#L35}{\texttt{Ursell.coefficient}}
& incompatibility graph of a tuple; Ursell coefficient \\
\(\Xi(w)\); cluster term; log series &
\href{https://github.com/yang-mills-group/OpenAI-YM-20261028/blob/3456e5e75cf8d4cd558955a5c7638fba2d68f464/YM/Push/ABSOLUTE_SUMS/FiniteLog.lean\#L22}{\texttt{Gas.total}},
\href{https://github.com/yang-mills-group/OpenAI-YM-20261028/blob/3456e5e75cf8d4cd558955a5c7638fba2d68f464/YM/Push/ABSOLUTE_SUMS/FiniteLog.lean\#L24}{\texttt{Gas.term}},
\href{https://github.com/yang-mills-group/OpenAI-YM-20261028/blob/3456e5e75cf8d4cd558955a5c7638fba2d68f464/YM/Push/ABSOLUTE_SUMS/FiniteLog.lean\#L26}{\texttt{Gas.logSum}}
& gas, its cluster terms and their sum \\
\(g_n\), \(\ell_m\) &
\href{https://github.com/yang-mills-group/OpenAI-YM-20261028/blob/3456e5e75cf8d4cd558955a5c7638fba2d68f464/YM/Push/ABSOLUTE_SUMS/Coef.lean\#L19}{\texttt{Coef.gasRow}},
\href{https://github.com/yang-mills-group/OpenAI-YM-20261028/blob/3456e5e75cf8d4cd558955a5c7638fba2d68f464/YM/Push/ABSOLUTE_SUMS/Coef.lean\#L20}{\texttt{Coef.logRow}}
& coefficients of \(\Xi\) and of the log series \\
\(w^Y\), \(p\) &
\href{https://github.com/yang-mills-group/OpenAI-YM-20261028/blob/3456e5e75cf8d4cd558955a5c7638fba2d68f464/YM/Push/ABSOLUTE_SUMS/Absence.lean\#L24}{\texttt{Gas.cut}},
\href{https://github.com/yang-mills-group/OpenAI-YM-20261028/blob/3456e5e75cf8d4cd558955a5c7638fba2d68f464/YM/Push/ABSOLUTE_SUMS/Absence.lean\#L25}{\texttt{Gas.missing}}
& weights off \(Y\); clusters meeting \(Y\) \\
multisets, \(\prod w^{n}/n!\) &
\href{https://github.com/yang-mills-group/OpenAI-YM-20261028/blob/3456e5e75cf8d4cd558955a5c7638fba2d68f464/YM/Push/ABSOLUTE_SUMS/Copies.lean\#L26}{\texttt{Multiple.Reps}},
\href{https://github.com/yang-mills-group/OpenAI-YM-20261028/blob/3456e5e75cf8d4cd558955a5c7638fba2d68f464/YM/Push/ABSOLUTE_SUMS/Copies.lean\#L28-L29}{\texttt{manyW}}
& labels with multiplicities \\
packet \(A\), number of seeds, \(D_r(A)\) &
\href{https://github.com/yang-mills-group/OpenAI-YM-20261028/blob/3456e5e75cf8d4cd558955a5c7638fba2d68f464/YM/Pressure/Seeds.lean\#L34-L37}{\texttt{Stores}},
\href{https://github.com/yang-mills-group/OpenAI-YM-20261028/blob/3456e5e75cf8d4cd558955a5c7638fba2d68f464/YM/Pressure/Seeds.lean\#L46}{\texttt{shares}},
\href{https://github.com/yang-mills-group/OpenAI-YM-20261028/blob/3456e5e75cf8d4cd558955a5c7638fba2d68f464/YM/Pressure/Seeds.lean\#L44}{\texttt{dome}}
& finite set of typed seeds; radius-\(r\) neighbourhood of their
locations \\
connected packet, \(K_{\rm hit}\), \(K_{\rm step}\), \(\tau_L\) &
\href{https://github.com/yang-mills-group/OpenAI-YM-20261028/blob/3456e5e75cf8d4cd558955a5c7638fba2d68f464/YM/Push/ABSOLUTE_SUMS/Packet.lean\#L77}{\texttt{Packets.connected}},
\href{https://github.com/yang-mills-group/OpenAI-YM-20261028/blob/3456e5e75cf8d4cd558955a5c7638fba2d68f464/YM/Push/ABSOLUTE_SUMS/Packet.lean\#L215}{\texttt{hitCost}},
\href{https://github.com/yang-mills-group/OpenAI-YM-20261028/blob/3456e5e75cf8d4cd558955a5c7638fba2d68f464/YM/Push/ABSOLUTE_SUMS/Packet.lean\#L216}{\texttt{stepCost}},
\href{https://github.com/yang-mills-group/OpenAI-YM-20261028/blob/3456e5e75cf8d4cd558955a5c7638fba2d68f464/YM/Push/ABSOLUTE_SUMS/Packet.lean\#L181}{\texttt{typesAt}}
& animal-counting constants \\
\(\rho\), packet tail &
\href{https://github.com/yang-mills-group/OpenAI-YM-20261028/blob/3456e5e75cf8d4cd558955a5c7638fba2d68f464/YM/Push/ABSOLUTE_SUMS/Charge.lean\#L28-L29}{\texttt{ratioCost}},
\href{https://github.com/yang-mills-group/OpenAI-YM-20261028/blob/3456e5e75cf8d4cd558955a5c7638fba2d68f464/YM/Push/ABSOLUTE_SUMS/Charge.lean\#L34-L35}{\texttt{fit}},
\href{https://github.com/yang-mills-group/OpenAI-YM-20261028/blob/3456e5e75cf8d4cd558955a5c7638fba2d68f464/YM/Push/ABSOLUTE_SUMS/Charge.lean\#L37-L38}{\texttt{totalCharge}}
& per-seed ratio; tail of packet sums \\
hypotheses of III.5.7(b) &
\href{https://github.com/yang-mills-group/OpenAI-YM-20261028/blob/3456e5e75cf8d4cd558955a5c7638fba2d68f464/YM/Push/ABSOLUTE_SUMS/Mixed.lean\#L47-L58}{\texttt{Mixed.Costs.Buy}},
\href{https://github.com/yang-mills-group/OpenAI-YM-20261028/blob/3456e5e75cf8d4cd558955a5c7638fba2d68f464/YM/Push/ABSOLUTE_SUMS/Mixed.lean\#L84-L89}{\texttt{Mixed.Accounts}}
& numerical list; weight bounds \\
\(B\), \(\beta_F\), \(t\), \(P\) &
\href{https://github.com/yang-mills-group/OpenAI-YM-20261028/blob/3456e5e75cf8d4cd558955a5c7638fba2d68f464/YM/Pressure/Seeds.lean\#L22}{\texttt{radiusB}},
\href{https://github.com/yang-mills-group/OpenAI-YM-20261028/blob/3456e5e75cf8d4cd558955a5c7638fba2d68f464/YM/Pressure/Seeds.lean\#L23}{\texttt{betaF}},
\href{https://github.com/yang-mills-group/OpenAI-YM-20261028/blob/3456e5e75cf8d4cd558955a5c7638fba2d68f464/YM/Pressure/Seeds.lean\#L26}{\texttt{fiery}},
\href{https://github.com/yang-mills-group/OpenAI-YM-20261028/blob/3456e5e75cf8d4cd558955a5c7638fba2d68f464/YM/Vac/Treasury.lean\#L72}{\texttt{Vac.Numbers.P}}
& \(\lceil(\log x)^2\rceil\); \(x+b\log L\);
\(\beta_F^{-1/2}P(\beta_F)\); \(e^{(\log\log s)^2}\) \\
\(M\), gain \(\beta\mu t^2\) & \texttt{hatWin},
\href{https://github.com/yang-mills-group/OpenAI-YM-20261028/blob/3456e5e75cf8d4cd558955a5c7638fba2d68f464/YM/Push/ABSOLUTE_SUMS/Grants.lean\#L48}{\texttt{Grant.credit}}
& guaranteed and actual gain per seed \\
\(Z(B)\), \(\lambda(A)\) &
\href{https://github.com/yang-mills-group/OpenAI-YM-20261028/blob/3456e5e75cf8d4cd558955a5c7638fba2d68f464/YM/SectorNext/VACUUM_NORMALIZATION/Gas.lean\#L122}{\texttt{fumes}},
\href{https://github.com/yang-mills-group/OpenAI-YM-20261028/blob/3456e5e75cf8d4cd558955a5c7638fba2d68f464/YM/SectorNext/VACUUM_NORMALIZATION/Deletion.lean\#L46}{\texttt{locker}},
\href{https://github.com/yang-mills-group/OpenAI-YM-20261028/blob/3456e5e75cf8d4cd558955a5c7638fba2d68f464/YM/SectorNext/VACUUM_NORMALIZATION/Sprig.lean\#L23-L28}{\texttt{pieces}}
& disjoint-family gas; Möbius pieces of \(-\log Z\) \\
\(f*g\), \(\hat f\), \(\Vert f\Vert_w\), \(\lambda(f)\) &
\href{https://github.com/yang-mills-group/OpenAI-YM-20261028/blob/3456e5e75cf8d4cd558955a5c7638fba2d68f464/YM/SectorNext/VACUUM_NORMALIZATION/Leaves.lean\#L36}{\texttt{glue}},
\href{https://github.com/yang-mills-group/OpenAI-YM-20261028/blob/3456e5e75cf8d4cd558955a5c7638fba2d68f464/YM/SectorNext/VACUUM_NORMALIZATION/Leaves.lean\#L38}{\texttt{shed}},
\href{https://github.com/yang-mills-group/OpenAI-YM-20261028/blob/3456e5e75cf8d4cd558955a5c7638fba2d68f464/YM/SectorNext/VACUUM_NORMALIZATION/Leaves.lean\#L40}{\texttt{crop}},
\href{https://github.com/yang-mills-group/OpenAI-YM-20261028/blob/3456e5e75cf8d4cd558955a5c7638fba2d68f464/YM/SectorNext/VACUUM_NORMALIZATION/Leaves.lean\#L161}{\texttt{lamp}}
& union convolution, cumulative sum, weighted norm, log series \\
hypotheses of III.5.9 &
\href{https://github.com/yang-mills-group/OpenAI-YM-20261028/blob/3456e5e75cf8d4cd558955a5c7638fba2d68f464/YM/SectorNext/VACUUM_NORMALIZATION/Leaves.lean\#L45-L47}{\texttt{Rate}},
\href{https://github.com/yang-mills-group/OpenAI-YM-20261028/blob/3456e5e75cf8d4cd558955a5c7638fba2d68f464/YM/SectorNext/VACUUM_NORMALIZATION/Deletion.lean\#L31-L34}{\texttt{Deletion}},
\href{https://github.com/yang-mills-group/OpenAI-YM-20261028/blob/3456e5e75cf8d4cd558955a5c7638fba2d68f464/YM/SectorNext/VACUUM_NORMALIZATION/Deletion.lean\#L39-L43}{\texttt{Receipt}},
\href{https://github.com/yang-mills-group/OpenAI-YM-20261028/blob/3456e5e75cf8d4cd558955a5c7638fba2d68f464/YM/SectorNext/VACUUM_NORMALIZATION/Gas.lean\#L168-L173}{\texttt{Sufficient}},
\href{https://github.com/yang-mills-group/OpenAI-YM-20261028/blob/3456e5e75cf8d4cd558955a5c7638fba2d68f464/YM/SectorNext/VACUUM_NORMALIZATION/Bulk.lean\#L43-L49}{\texttt{Fund}},
\href{https://github.com/yang-mills-group/OpenAI-YM-20261028/blob/3456e5e75cf8d4cd558955a5c7638fba2d68f464/YM/SectorNext/VACUUM_NORMALIZATION/Bulk.lean\#L72-L77}{\texttt{Tags}}
& weight, deletion identity, site conditions, coefficient properties \\
allocation \(D\); point charge &
\href{https://github.com/yang-mills-group/OpenAI-YM-20261028/blob/3456e5e75cf8d4cd558955a5c7638fba2d68f464/YM/Push/SECTOR_MEASURE/Profile.lean\#L30-L36}{\texttt{Allocation}},
\href{https://github.com/yang-mills-group/OpenAI-YM-20261028/blob/3456e5e75cf8d4cd558955a5c7638fba2d68f464/YM/Push/SECTOR_MEASURE/Oriented.lean\#L29-L32}{\texttt{PointCharge}}
& partition of unity over packets; \(\mathcal E_A\ge\kappa t^2\) per
seed \\
\(\mathcal E_A\); \(c_A\); \(\ell_A\);
\(\pi_\beta e^{-\beta\kappa t^2 k}\) &
\href{https://github.com/yang-mills-group/OpenAI-YM-20261028/blob/3456e5e75cf8d4cd558955a5c7638fba2d68f464/YM/Pressure/Seeds.lean\#L85}{\texttt{assessed}},
\href{https://github.com/yang-mills-group/OpenAI-YM-20261028/blob/3456e5e75cf8d4cd558955a5c7638fba2d68f464/YM/Push/SECTOR_MEASURE/Presented.lean\#L145-L148}{\texttt{cutEntry}},
\href{https://github.com/yang-mills-group/OpenAI-YM-20261028/blob/3456e5e75cf8d4cd558955a5c7638fba2d68f464/YM/Push/SECTOR_MEASURE/Presented.lean\#L140-L143}{\texttt{leftOver}},
\href{https://github.com/yang-mills-group/OpenAI-YM-20261028/blob/3456e5e75cf8d4cd558955a5c7638fba2d68f464/YM/Push/SECTOR_MEASURE/Grant.lean\#L28-L29}{\texttt{earn}}
& sector energy; sector output; remainder; gain \\
weighted Schur classes; \(M\ge r\) &
\href{https://github.com/yang-mills-group/OpenAI-YM-20261028/blob/3456e5e75cf8d4cd558955a5c7638fba2d68f464/YM/Progress/STATIONARY_SCENES/Green.lean\#L26-L27}{\texttt{Green.Rows}},
\href{https://github.com/yang-mills-group/OpenAI-YM-20261028/blob/3456e5e75cf8d4cd558955a5c7638fba2d68f464/YM/Progress/STATIONARY_SCENES/Green.lean\#L29-L30}{\texttt{Green.Both}},
\href{https://github.com/yang-mills-group/OpenAI-YM-20261028/blob/3456e5e75cf8d4cd558955a5c7638fba2d68f464/YM/Progress/STATIONARY_SCENES/Green.lean\#L32-L33}{\texttt{Green.Bottom}}
& row/column sums at rate \(\eta\); lower bound \\
windows, \(\mathrm{within}\), \(E\), \(\mathrm{cov}\), \(\Pi_U\) &
\href{https://github.com/yang-mills-group/OpenAI-YM-20261028/blob/3456e5e75cf8d4cd558955a5c7638fba2d68f464/YM/AnalyticNext/CONNECTED_ACTIVITY_BOUNDS/Windows.lean\#L28-L30}{\texttt{Win.Windows}},
\href{https://github.com/yang-mills-group/OpenAI-YM-20261028/blob/3456e5e75cf8d4cd558955a5c7638fba2d68f464/YM/AnalyticNext/CONNECTED_ACTIVITY_BOUNDS/Windows.lean\#L36-L37}{\texttt{within}},
\href{https://github.com/yang-mills-group/OpenAI-YM-20261028/blob/3456e5e75cf8d4cd558955a5c7638fba2d68f464/YM/AnalyticNext/CONNECTED_ACTIVITY_BOUNDS/Windows.lean\#L39}{\texttt{error}},
\href{https://github.com/yang-mills-group/OpenAI-YM-20261028/blob/3456e5e75cf8d4cd558955a5c7638fba2d68f464/YM/AnalyticNext/CONNECTED_ACTIVITY_BOUNDS/Windows.lean\#L73}{\texttt{cov}},
\href{https://github.com/yang-mills-group/OpenAI-YM-20261028/blob/3456e5e75cf8d4cd558955a5c7638fba2d68f464/YM/AnalyticNext/CONNECTED_ACTIVITY_BOUNDS/Schur.lean\#L69-L70}{\texttt{planePrice}}
& localized inverse and covariance \\
uniform matrix data; window constants &
\href{https://github.com/yang-mills-group/OpenAI-YM-20261028/blob/3456e5e75cf8d4cd558955a5c7638fba2d68f464/YM/AnalyticNext/CONNECTED_ACTIVITY_BOUNDS/Birthday.lean\#L22-L30}{\texttt{Win.Promise}},
\href{https://github.com/yang-mills-group/OpenAI-YM-20261028/blob/3456e5e75cf8d4cd558955a5c7638fba2d68f464/YM/AnalyticNext/CONNECTED_ACTIVITY_BOUNDS/Birthday.lean\#L71-L83}{\texttt{Turns}},
\href{https://github.com/yang-mills-group/OpenAI-YM-20261028/blob/3456e5e75cf8d4cd558955a5c7638fba2d68f464/YM/AnalyticNext/CONNECTED_ACTIVITY_BOUNDS/Birthday.lean\#L91-L101}{\texttt{Turns.Takes}}
& (mass, bottom, rate, counting); bounds of III.5.13(b) \\
\(D\), \(\mathrm{gram}\) &
\href{https://github.com/yang-mills-group/OpenAI-YM-20261028/blob/3456e5e75cf8d4cd558955a5c7638fba2d68f464/YM/AnalyticNext/CONNECTED_ACTIVITY_BOUNDS/Affine.lean\#L125-L127}{\texttt{Aff.Dmat}},
\href{https://github.com/yang-mills-group/OpenAI-YM-20261028/blob/3456e5e75cf8d4cd558955a5c7638fba2d68f464/YM/AnalyticNext/CONNECTED_ACTIVITY_BOUNDS/Affine.lean\#L145-L147}{\texttt{Aff.gram}}
& transported flat row matrix; completed precision matrix \\
transported flat derivative &
\href{https://github.com/yang-mills-group/OpenAI-YM-20261028/blob/3456e5e75cf8d4cd558955a5c7638fba2d68f464/YM/AnalyticNext/CONNECTED_ACTIVITY_BOUNDS/Affine.lean\#L116-L118}{\texttt{Aff.derivative}}
& row operator used to define \(D\) \\
auxiliary Gaussian assumptions &
\href{https://github.com/yang-mills-group/OpenAI-YM-20261028/blob/3456e5e75cf8d4cd558955a5c7638fba2d68f464/YM/AnalyticNext/CONNECTED_ACTIVITY_BOUNDS/Veil.lean\#L34-L39}{\texttt{Aff.Reserve}},
\href{https://github.com/yang-mills-group/OpenAI-YM-20261028/blob/3456e5e75cf8d4cd558955a5c7638fba2d68f464/YM/AnalyticNext/CONNECTED_ACTIVITY_BOUNDS/Station.lean\#L71-L78}{\texttt{Aff.Alongside}},
\href{https://github.com/yang-mills-group/OpenAI-YM-20261028/blob/3456e5e75cf8d4cd558955a5c7638fba2d68f464/YM/AnalyticNext/CONNECTED_ACTIVITY_BOUNDS/Odometer.lean\#L23-L24}{\texttt{Aff.cert}}
& reserve bounds; compatible quadratic forms; volume and flag tests \\
\(H_{n,p}\); comparison field \(v\) &
\href{https://github.com/yang-mills-group/OpenAI-YM-20261028/blob/3456e5e75cf8d4cd558955a5c7638fba2d68f464/YM/AnalyticNext/CONNECTED_ACTIVITY_BOUNDS/Rate.lean\#L27}{\texttt{Rates.allowance}},
\href{https://github.com/yang-mills-group/OpenAI-YM-20261028/blob/3456e5e75cf8d4cd558955a5c7638fba2d68f464/YM/AnalyticNext/CONNECTED_ACTIVITY_BOUNDS/Through.lean\#L29-L30}{\texttt{Aff.unit}}
& \(B^pt^n\); active stationary extension minus replica extension,
transported back by the local gauge \\
\end{longtable}
}

\subsubsection{7. What remains
unexplained}\label{what-remains-unexplained-7}

\begin{enumerate}
\def\labelenumi{\arabic{enumi}.}
\tightlist
\item
  \textbf{Imported analytic inputs.} The step-estimator's energy grant
  in III.10 and the Poincaré inequality in III.6 are used as stated.
  Their proofs, the full per-bond spectral comparison and the origin of
  the run's committed-scale energy hypothesis have not been audited
  here.
\item
  \textbf{Later applications.} The dependency database identifies
  consumers in the units listed above. Their verification of the
  site-sum, per-packet and multiset hypotheses remains outside this
  unit's audit. In particular, the bounds here do not construct
  renormalization outputs or bound the sector remainders uniformly in
  volume.
\item
  \textbf{Unread supporting proofs.} The breadth-first construction, the
  distinct-label forest count, the spanning-walk construction for
  animals and the integration of allocations restricted to sectors were
  used through their statements. In the auxiliary Gaussian chain, the
  statements and the displayed estimates of Lemma 4.36 were checked; its
  supporting local chart and quadratic-form comparison lemmas were not
  audited throughout. These remain proof inputs, rather than additional
  conclusions of the present sketches.
\item
  \textbf{Flag locality.} In III.5.9(b), locality on \(P\) applies to
  the raw coefficients. Their masked versions also inspect flags on
  \(P\); those flags can depend on plaquettes extending outside \(P\). A
  written proof must retain this flag dependence or enlarge the physical
  support before claiming ordinary link locality for the masked
  coefficients.
\item
  \textbf{Verification scope.} This is a source correspondence check. No
  Lean build, axiom computation or rendered-page check was run.
\end{enumerate}

\subsection{III.6. The stationary branch: background field of the
effective
action}\label{iii.6.-the-stationary-branch-background-field-of-the-effective-action}

\textbf{Part:} III, Ultraviolet construction (renormalization group).
\textbf{Lean:} the work items F01--F05 are stated in
\href{https://github.com/yang-mills-group/OpenAI-YM-20261028/blob/3456e5e75cf8d4cd558955a5c7638fba2d68f464/YM/Bid/Station.lean}{\texttt{YM/Bid/Station.lean}}
and F06--F09 in
\href{https://github.com/yang-mills-group/OpenAI-YM-20261028/blob/3456e5e75cf8d4cd558955a5c7638fba2d68f464/YM/Bid/Craft.lean}{\texttt{YM/Bid/Craft.lean}}.
Their proofs are in
\href{https://github.com/yang-mills-group/OpenAI-YM-20261028/tree/3456e5e75cf8d4cd558955a5c7638fba2d68f464/YM/Work}{\texttt{YM/Work/F01}--\texttt{F09}}.
The other packages are: the uniform branch
\href{https://github.com/yang-mills-group/OpenAI-YM-20261028/tree/3456e5e75cf8d4cd558955a5c7638fba2d68f464/YM/Push/STATIONARY_UNIFORM}{\texttt{YM/Push/STATIONARY\_UNIFORM/}};
the jet port
\href{https://github.com/yang-mills-group/OpenAI-YM-20261028/tree/3456e5e75cf8d4cd558955a5c7638fba2d68f464/YM/Planned/STATIONARY_JET_PORT}{\texttt{YM/Planned/STATIONARY\_JET\_PORT/}};
and the iteration step
\href{https://github.com/yang-mills-group/OpenAI-YM-20261028/tree/3456e5e75cf8d4cd558955a5c7638fba2d68f464/YM/Extra/ITERATION_NO_LOSS}{\texttt{YM/Extra/ITERATION\_NO\_LOSS/}}.
Also included are
\href{https://github.com/yang-mills-group/OpenAI-YM-20261028/tree/3456e5e75cf8d4cd558955a5c7638fba2d68f464/YM/Planned/PRIMARY_INTEGRAL_295}{\texttt{YM/Planned/PRIMARY\_INTEGRAL\_295/}},
\href{https://github.com/yang-mills-group/OpenAI-YM-20261028/tree/3456e5e75cf8d4cd558955a5c7638fba2d68f464/YM/Further/PHYSICAL_PERIOD_LIMIT_F933}{\texttt{YM/Further/PHYSICAL\_PERIOD\_LIMIT\_F933/}},
\href{https://github.com/yang-mills-group/OpenAI-YM-20261028/tree/3456e5e75cf8d4cd558955a5c7638fba2d68f464/YM/Planned/ATOM_REGULAR_ESTIMATE}{\texttt{YM/Planned/ATOM\_REGULAR\_ESTIMATE/}},
\href{https://github.com/yang-mills-group/OpenAI-YM-20261028/tree/3456e5e75cf8d4cd558955a5c7638fba2d68f464/YM/Next/ROTATIONS}{\texttt{YM/Next/ROTATIONS/}}
and
\href{https://github.com/yang-mills-group/OpenAI-YM-20261028/tree/3456e5e75cf8d4cd558955a5c7638fba2d68f464/YM/Next/WARD_BIN}{\texttt{YM/Next/WARD\_BIN/}}.
\textbf{Size:} 21,329 lines in the main-theorem closure. Of these:

\begin{itemize}
\tightlist
\item
  the F work items F01--F09 with their auxiliary files: about 3,800;
\item
  the uniform branch: 3,274;
\item
  the jet port: 1,123;
\item
  the iteration step: 1,700;
\item
  the primary-integral package: 3,662;
\item
  the period-limit comparison: 1,262;
\item
  the atom regularity estimate: 474;
\item
  the rotation and Ward material: about 5,900, including the chain of
  results leading to the rotation theorem.
\end{itemize}

\subsubsection{1. What this unit proves}\label{what-this-unit-proves-8}

\textbf{Setting.} The definitions are those of III.1. All lattice sums
use counting measure.

\emph{Lattices, group, distance.}

\begin{itemize}
\tightlist
\item
  One renormalization step goes from a fine torus \(\Lambda_f\) to a
  coarse torus \(\Lambda_c\), with an odd block factor \(L\ge3\) and
  \(f=Lc\), in dimension \(4\). The patch order satisfies \(J\ge4\).
\item
  \(G\) is the gauge group and \(\mathfrak g\) its Lie algebra, normed
  in physical colour coordinates.
\item
  \(\mathrm{chord}(U)=\Vert U-1\Vert/c_G\) and
  \(\mathrm{price}(U)=\tfrac12\mathrm{chord}(U)^2\). With this
  normalization \(\mathrm{price}(e^{tX})=\tfrac12t^2|X|^2+o(t^2)\); the
  price is the Wilson plaquette function up to a constant.
\end{itemize}

\emph{The step datum \(S\).} It consists of:

\begin{itemize}
\tightlist
\item
  a block-spin kernel, with penalty \(\kappa>0\), from the kernel
  choices \(K\);
\item
  a patch, with order \(J\) and macro-radius \(R_S=(J+7)L\);
\item
  a finite set of tags (frame and path choices;
  \href{https://github.com/yang-mills-group/OpenAI-YM-20261028/blob/3456e5e75cf8d4cd558955a5c7638fba2d68f464/YM/RG/Block.lean\#L34}{\texttt{StepData.Choices}});
\item
  for each tag \(k\), a predictor \(U\mapsto Q_k(U)\) from fine
  configurations to coarse configurations.
\end{itemize}

\(P\) denotes the derivative of the predictor at the identity. It acts
on fine \(\mathfrak g\)-valued \(1\)-forms; it is the linearized block
average plus a gradient correction at the block roots. \(\bar Q\)
denotes the linearized block average alone.

\emph{``Larger''.} \(S\) is \emph{larger}
(\href{https://github.com/yang-mills-group/OpenAI-YM-20261028/blob/3456e5e75cf8d4cd558955a5c7638fba2d68f464/YM/Bid/Station.lean\#L64}{\texttt{Larger}})
if three conditions hold:

\begin{itemize}
\tightlist
\item
  the minimizing frame and predictor are unique whenever their
  minimizing branches are selected, and the fallback set of eligible
  axial frames is nonempty; see
  \href{https://github.com/yang-mills-group/OpenAI-YM-20261028/blob/3456e5e75cf8d4cd558955a5c7638fba2d68f464/YM/RG/Patch.lean\#L182-L185}{\texttt{Branches}};
\item
  \(100R_S\le f\);
\item
  \(c>1\).
\end{itemize}

\emph{Forests and the axis chart.}

\begin{itemize}
\tightlist
\item
  A \emph{forest} \(H\) is a spanning forest of the fine torus, made of
  one tree in each block, rooted at the block centre, with tree height
  at most \(4L\)
  (\href{https://github.com/yang-mills-group/OpenAI-YM-20261028/blob/3456e5e75cf8d4cd558955a5c7638fba2d68f464/YM/Tree/Pines.lean\#L24-L29}{\texttt{Forest}}).
  Its bonds are the \emph{tree bonds}, all internal to blocks; the
  remaining bonds are the \emph{free bonds} \(H^{\rm fr}\).
\item
  The \emph{replica} \(\mathcal R(V)\) of a coarse configuration \(V\)
  puts \(V_{(i,\mu)}\) on every fine bond that leaves block \(i\) in
  direction \(\mu\), and \(1\) on bonds internal to a block. Its
  linearization is \(\iota\).
\item
  The \emph{axis chart} has coordinates
  \(x\in\mathfrak g^{\text{coarse bonds}}\) and
  \(a\in\mathfrak g^{H^{\rm fr}}\). It describes the fine field
  \[U^H(x,a)_e=e^{a_e}\,\mathcal R(e^{x})_e\quad(e\in H^{\rm fr}),\qquad U^H(x,a)_e=1\quad(e\ \text{a tree bond}).\]
\end{itemize}

\emph{Clip and path datum.}

\begin{itemize}
\tightlist
\item
  A \emph{clip} \(p\) carries the Wilson coefficient \(c_p>0\).
\item
  A \emph{path datum} \(W\)
  (\href{https://github.com/yang-mills-group/OpenAI-YM-20261028/blob/3456e5e75cf8d4cd558955a5c7638fba2d68f464/YM/RG/PathActions.lean\#L150-L154}{\texttt{PathData}})
  consists of:

  \begin{itemize}
  \tightlist
  \item
    a series \(\ell_W(U)_\pi=\sum_i w_i\,\mathrm{slot}_i(U)\) of
    parallel-transported plaquette variables, attached to each plaquette
    \(\pi\);
  \item
    loop atoms, cubic atoms and quartic atoms: \(\mathrm{Ad}\)-invariant
    polynomials of degree \(2,3,4\) in such slot values.
  \end{itemize}
\item
  \emph{Bounds} \((p,\sigma,C,\Lambda)\) for \(W\)
  (\href{https://github.com/yang-mills-group/OpenAI-YM-20261028/blob/3456e5e75cf8d4cd558955a5c7638fba2d68f464/YM/RG/PathActions.lean\#L173-L185}{\texttt{Bounds}})
  mean four things:

  \begin{itemize}
  \tightlist
  \item
    the coefficients are summable with weight
    \(e^{\sigma\cdot\mathrm{length}}\), with sums at most \(C\) (cubic,
    quartic, \(\ell\)-rows) and \(\Lambda\) (loop atoms);
  \item
    the atoms are \(\mathrm{Ad}\)-invariant;
  \item
    if \(r_{ij}\) is the sum of row-\(i\) coefficients whose plaquette
    type is \(j\), then
    \(c_p\delta_{ij}+\sum_k r_{ki}r_{kj}=\delta_{ij}\);
  \item
    for each pair of plaquette types and colour indices, the loop
    coefficients have zero sum and zero first moment in every component
    of the relative slot offset; these moment series are summable.
  \end{itemize}
\end{itemize}

\emph{The tested action.}

\begin{itemize}
\item
  The \emph{main term}
  (\href{https://github.com/yang-mills-group/OpenAI-YM-20261028/blob/3456e5e75cf8d4cd558955a5c7638fba2d68f464/YM/PCalc/Chart.lean\#L27-L28}{\texttt{mainTerm}})
  is the average over the finite group \(\mathcal S\) of signed
  coordinate permutations, acting by pullback \(s^*U\):
  \[\mathcal M_p(W;U)=\frac1{|\mathcal S|}\sum_{s\in\mathcal S}\Big[\sum_\pi\big(c_p\,\mathrm{price}((s^*U)_\pi)+\tfrac12|\ell_W(s^*U)_\pi|^2\big)+\sum_x\sum_i\big(\mathrm{cub}_i+\mathrm{quart}_i\big)(s^*U)_x\Big].\]
\item
  The \emph{tested action} for tag \(k\)
  (\href{https://github.com/yang-mills-group/OpenAI-YM-20261028/blob/3456e5e75cf8d4cd558955a5c7638fba2d68f464/YM/PCalc/Chart.lean\#L81-L84}{\texttt{tested}})
  is
  \[\mathcal T^H_k(x,a)=\mathcal M_p\big(W;U^H(x,a)\big)+\kappa\sum_{b}\mathrm{price}\big(e^{x_b}\,Q_k(U^H(x,a))_b^{-1}\big),\]
  where \(b\) runs over the coarse bonds. The second term is the
  block-spin penalty.
\end{itemize}

\emph{Beds, towers and the quadratic identities.}

\begin{itemize}
\item
  A \emph{bed} \(B\)
  (\href{https://github.com/yang-mills-group/OpenAI-YM-20261028/blob/3456e5e75cf8d4cd558955a5c7638fba2d68f464/YM/PREC/Stacks.lean\#L35-L40}{\texttt{Bed}})
  is a family of translation-invariant kernels on plaquette \(2\)-forms:
  a metric \(M_B\), an inverse kernel, a diagonal kernel \(D_B\), a dual
  diagonal kernel, and roots \(R_B\). Set \(K_B=d^\top M_Bd\), which
  acts on \(1\)-forms.
\item
  A \emph{floor} \((a,A)\) for \(B\)
  (\href{https://github.com/yang-mills-group/OpenAI-YM-20261028/blob/3456e5e75cf8d4cd558955a5c7638fba2d68f464/YM/PREC/Stacks.lean\#L61-L68}{\texttt{Floor}})
  requires:

  \begin{itemize}
  \tightlist
  \item
    the root identity \(M_B=c_p\,1+R_B^\top R_B\);
  \item
    inverse identities for the metric and diagonal kernels;
  \item
    the Fourier-symbol normalization \(\widehat M_B(0)=I\) and
    \(M_Bd=D_Bd\) on the specified torus;
  \item
    symmetry of \(M_B\) and \(D_B\);
  \item
    the bounds
    \(a|\omega|^2\le\langle\omega,M_B\omega\rangle\le A|\omega|^2\), and
    the same bounds for \(D_B\).
  \end{itemize}
\item
  \(W\) is \emph{attached} to \(B\) if the coefficients of \(\ell_W\),
  summed by offset and plaquette type, are the entries of \(R_B\). \(B\)
  is \emph{sworn} if \(\langle dA,M_BdA\rangle\) is
  \(\mathcal S\)-invariant.
\item
  \emph{Towers} \(T\) for the pair of beds \((B,C)\)
  (\href{https://github.com/yang-mills-group/OpenAI-YM-20261028/blob/3456e5e75cf8d4cd558955a5c7638fba2d68f464/YM/PREC/Stacks.lean\#L82-L86}{\texttt{Towers.Diagonal}})
  provide a linear map \(\mathcal W\) from coarse \(1\)-forms to fine
  \(1\)-forms. Its first clause,
  \href{https://github.com/yang-mills-group/OpenAI-YM-20261028/blob/3456e5e75cf8d4cd558955a5c7638fba2d68f464/YM/PREC/TreeMatch.lean\#L76-L79}{\texttt{Inputs}},
  is
  \[K_B\mathcal W=\bar Q^\top K_C,\qquad \bar Q\mathcal W=1-\kappa^{-1}K_C .\]
\item
  \(\widehat{\mathcal W}_H=\mathcal W+d_0\chi_H\) is the linearized
  gauge transform of \(\mathcal W\) that vanishes on tree bonds
  (\href{https://github.com/yang-mills-group/OpenAI-YM-20261028/blob/3456e5e75cf8d4cd558955a5c7638fba2d68f464/YM/PREC/TreeMatch.lean\#L26-L27}{\texttt{extracted}}).
\item
  \(\mathrm{rel}_H(x)=(\widehat{\mathcal W}_Hx-\iota x)|_{H^{\rm fr}}\)
  (\href{https://github.com/yang-mills-group/OpenAI-YM-20261028/blob/3456e5e75cf8d4cd558955a5c7638fba2d68f464/YM/PCalc/Chart.lean\#L93-L95}{\texttt{relative}}).
\end{itemize}

\emph{Critical branch.} A \emph{critical branch} for \((H,p,W)\)
(\href{https://github.com/yang-mills-group/OpenAI-YM-20261028/blob/3456e5e75cf8d4cd558955a5c7638fba2d68f464/YM/PCalc/Chart.lean\#L88-L91}{\texttt{Critical}})
is a map \(\rho\) from coarse fields to \(\mathfrak g^{H^{\rm fr}}\)
such that:

\begin{itemize}
\tightlist
\item
  \(\rho(0)=0\) and \(\rho\) is real-analytic at \(0\);
\item
  for every tag \(k\), \(\mathcal T^H_k\) is jointly analytic at \(0\);
\item
  for every tag \(k\), \(D_a\mathcal T^H_k(x,\rho(x))=0\) for \(x\) near
  \(0\).
\end{itemize}

\emph{Chart ratio and thermal function.}

\begin{itemize}
\item
  A \emph{local chart ratio} \(\mathscr J\)
  (\href{https://github.com/yang-mills-group/OpenAI-YM-20261028/blob/3456e5e75cf8d4cd558955a5c7638fba2d68f464/YM/PCalc/Loop.lean\#L28-L38}{\texttt{LocalRatio}})
  is a density \(J_{\mathscr J}\) on a ball
  \(\{|X|<r_{\mathscr J}\}\subset\mathfrak g\), where
  \(r_{\mathscr J}>0\) and \(s_{\mathscr J}>0\) is a scale. It
  satisfies:

  \begin{itemize}
  \tightlist
  \item
    \(J_{\mathscr J}\) is measurable, positive and real-analytic on the
    ball, with \(J_{\mathscr J}(0)=1\);
  \item
    \(\exp\) is injective on the ball and its image is open;
  \item
    the pushforward of \(s_{\mathscr J}J_{\mathscr J}(X)\,dX\)
    restricted to the ball is probability Haar measure restricted to its
    exponential image.
  \end{itemize}

  Only these density properties enter the forest comparison.
\item
  The \emph{thermal function} of a branch
  (\href{https://github.com/yang-mills-group/OpenAI-YM-20261028/blob/3456e5e75cf8d4cd558955a5c7638fba2d68f464/YM/PCalc/Loop.lean\#L46-L49}{\texttt{thermal}})
  is
  \[\Theta^H_k(x)=\mathrm{loop}_W\big(U^H(x,\rho(x))\big)+\tfrac12\log\det\mathrm{Hess}_a\mathcal T^H_k(x,\rho(x))-\sum_{e\in H^{\rm fr}}\log J_{\mathscr J}(\rho(x)_e).\]
  The Hessian is taken in the orthonormal colour basis. For a branch
  with a positive vertical gap, the determinant and chart densities are
  positive near \(0\).
\end{itemize}

\textbf{Proposition III.6.1 (second-order jet of the tested action on
the tree slice; F03).} Let \(S\) be larger, \(H\) a forest and \(B\) a
bed with a floor \((a,A)\) on \(\Lambda_f\) for the clip \(p\). Let
\(W\) be a path datum with bounds \((p,\sigma,C,\Lambda)\), attached to
\(B\), with \(B\) sworn on \(\Lambda_f\). Assume that the main term of
\(W\) is analytic along analytic families through the identity (true by
III.6.2). Then for every tag \(k\), \(\mathcal T^H_k\) is real-analytic
at \(0\), with \(\mathcal T^H_k(0)=0\) and \(D\mathcal T^H_k(0)=0\).
Moreover, for all \((x,a),(x',a')\),
\[D^2\mathcal T^H_k(0)\big[(x,a),(x',a')\big]=\langle ds,\,M_B\,ds'\rangle+\kappa\,\langle x-Ps,\;x'-Ps'\rangle,\qquad s=\iota x+a,\ s'=\iota x'+a'.\]
Here \(a\) is extended by \(0\) to tree bonds, and the pairings are
summed over plaquettes or coarse bonds and over colours. \emph{Status:}
theorem. Only the well-posedness clause of ``larger'' is used.
\emph{Lean:}
\href{https://github.com/yang-mills-group/OpenAI-YM-20261028/blob/3456e5e75cf8d4cd558955a5c7638fba2d68f464/YM/Bid/Station.lean\#L78-L81}{\texttt{Bid.F03}},
with conclusions
\href{https://github.com/yang-mills-group/OpenAI-YM-20261028/blob/3456e5e75cf8d4cd558955a5c7638fba2d68f464/YM/Bid/Station.lean\#L47}{\texttt{Prismatic}}
and
\href{https://github.com/yang-mills-group/OpenAI-YM-20261028/blob/3456e5e75cf8d4cd558955a5c7638fba2d68f464/YM/Bid/Station.lean\#L42-L45}{\texttt{Comptroller}};
proof
\href{https://github.com/yang-mills-group/OpenAI-YM-20261028/blob/3456e5e75cf8d4cd558955a5c7638fba2d68f464/YM/Work/F03.lean\#L159-L175}{\texttt{Work.F03.answer}}.

\textbf{Proposition III.6.2 (analyticity of the main term; F02).} Let
\(W\) have bounds \((p,\sigma,C,\Lambda)\). Let \(X\) be a
finite-dimensional real normed space and \(x_0\in X\). Let
\(y\mapsto A(y)\) be a family of fine configurations with \(A(x_0)=1\)
whose link matrices are real-analytic at \(x_0\). Then
\(y\mapsto\mathcal M_p(W;A(y))\) is real-analytic at \(x_0\). The proof
does not use \(A(x_0)=1\): the main term is analytic along analytic
families at every point. \emph{Status:} theorem. It supplies the
hypothesis
\href{https://github.com/yang-mills-group/OpenAI-YM-20261028/blob/3456e5e75cf8d4cd558955a5c7638fba2d68f464/YM/Bid/Station.lean\#L59-L62}{\texttt{Fires}}
of III.6.1. \emph{Lean:}
\href{https://github.com/yang-mills-group/OpenAI-YM-20261028/blob/3456e5e75cf8d4cd558955a5c7638fba2d68f464/YM/Bid/Station.lean\#L73}{\texttt{Bid.F02}},
\href{https://github.com/yang-mills-group/OpenAI-YM-20261028/blob/3456e5e75cf8d4cd558955a5c7638fba2d68f464/YM/Work/F02.lean\#L141-L143}{\texttt{Work.F02.answer}}.

\textbf{Proposition III.6.3 (gap of the vertical Hessian; F01).} Assume
the shell comparison E02 (§2). Fix \(G\), \(L\), \(J\), kernel choices
\(K\) with penalty \(\kappa\), and \(a>0\). Then there is \(\eta>0\)
with the following property. Take sides \(c,f\); a larger step datum
\(S\) with kernel choices \(K\), block factor \(L\) and order \(J\); a
forest \(H\); a clip \(p\); a bed \(B\) with floor \((a,A)\) for some
\(A\); and a path datum \(W\) such that every \(\mathcal T^H_k\) is
analytic at \(0\) with the value, gradient and Hessian specified in
III.6.1. Then for every tag \(r\) and every
\(v\in\mathfrak g^{H^{\rm fr}}\),
\[\eta\sum_{e\in H^{\rm fr}}|v_e|^2\le D_a^2\mathcal T^H_r(0,0)[v,v].\]
One may take
\[\eta=\frac{\min(a,\kappa)}{1+\beta+\gamma},\qquad \beta=2C_{\rm E02}^2\,L^4\cdot4\cdot6\,|X|^2,\qquad \gamma=2L^4,\]
where:

\begin{itemize}
\tightlist
\item
  \(C_{\rm E02}=C_{\rm E02}(G,L,J,K)\) is the constant of E02;
\item
  \(X=\{-R_S,\ldots,R_S\}^4\subset\mathbb Z^4\), where \(R_S=(J+7)L\),
  so \(|X|=(2R_S+1)^4\);
\item
  the factors \(L^4\), \(4\) and \(6\) are the numbers of fine sites per
  block, of axes and of plaquette orientations.
\end{itemize}

So \(\eta\) depends only on \((G,L,J,K,a)\). It does not depend on
\(c\), \(f\), \(H\), \(B\) beyond its floor, or \(W\). \emph{Status:}
theorem, conditional on E02 in its statement. E02 is proved in III.10
(\href{https://github.com/yang-mills-group/OpenAI-YM-20261028/blob/3456e5e75cf8d4cd558955a5c7638fba2d68f464/YM/Work/E02.lean\#L132-L244}{\texttt{Work.E02.answer}}),
and the combination is unconditional in III.6.6. \emph{Lean:}
\href{https://github.com/yang-mills-group/OpenAI-YM-20261028/blob/3456e5e75cf8d4cd558955a5c7638fba2d68f464/YM/Bid/Station.lean\#L68-L71}{\texttt{Bid.F01}},
with conclusion
\href{https://github.com/yang-mills-group/OpenAI-YM-20261028/blob/3456e5e75cf8d4cd558955a5c7638fba2d68f464/YM/PCalc/Chart.lean\#L100-L101}{\texttt{PriceGap}};
proof
\href{https://github.com/yang-mills-group/OpenAI-YM-20261028/blob/3456e5e75cf8d4cd558955a5c7638fba2d68f464/YM/Work/F01.lean\#L266-L311}{\texttt{Work.F01.answer}}.

\textbf{Proposition III.6.4 (existence, tangent and value of the
critical branch; F04).} Let \(S\) be larger, \(H\) a forest, \(B,C\)
beds, \(p\) a clip, \(W\) a path datum, \(a>0\), and \(T\) towers for
\((S,B,C)\) satisfying the diagonal identities. Assume:

\begin{itemize}
\tightlist
\item
  every \(\mathcal T^H_k\) is analytic at \(0\);
\item
  every \(\mathcal T^H_k\) has the value, gradient and Hessian specified
  in III.6.1 with bed \(B\);
\item
  the gap of III.6.3 holds with constant \(a\);
\item
  there are a clip \(q\) and real constants \(v,u\) for which \(C\) has
  floor \((v,u)\) on \(\Lambda_c\).
\end{itemize}

Then there is a critical branch \(\rho\) for \((H,p,W)\) such that:

\begin{itemize}
\tightlist
\item
  (ii) \(D\rho(0)=\mathrm{rel}_H\), built from the tower map
  \(\mathcal W\);
\item
  (iii) for every tag \(k\), \(x\mapsto\mathcal T^H_k(x,\rho(x))\) is
  analytic at \(0\), and its second derivative at \(0\) is
  \(\langle dx,M_C\,dx'\rangle\). Equivalently, its \(2\)-jet is that of
  \(\tfrac12\langle dx,M_Cdx\rangle\).
\end{itemize}

The floor of \(C\) is unused in the proof. Of the diagonal identities,
only the clause \(K_B\mathcal W=\bar Q^\top K_C\),
\(\bar Q\mathcal W=1-\kappa^{-1}K_C\) is used. \emph{Status:} theorem.
It gives a germ only: no control of the size of the neighbourhood.
\emph{Lean:}
\href{https://github.com/yang-mills-group/OpenAI-YM-20261028/blob/3456e5e75cf8d4cd558955a5c7638fba2d68f464/YM/Bid/Station.lean\#L86-L90}{\texttt{Bid.F04}},
with conclusions
\href{https://github.com/yang-mills-group/OpenAI-YM-20261028/blob/3456e5e75cf8d4cd558955a5c7638fba2d68f464/YM/PCalc/Chart.lean\#L88-L91}{\texttt{Critical}},
\href{https://github.com/yang-mills-group/OpenAI-YM-20261028/blob/3456e5e75cf8d4cd558955a5c7638fba2d68f464/YM/PCalc/Chart.lean\#L97-L98}{\texttt{Guessed}}
and
\href{https://github.com/yang-mills-group/OpenAI-YM-20261028/blob/3456e5e75cf8d4cd558955a5c7638fba2d68f464/YM/Bid/Station.lean\#L83-L84}{\texttt{HeatG}};
proof
\href{https://github.com/yang-mills-group/OpenAI-YM-20261028/blob/3456e5e75cf8d4cd558955a5c7638fba2d68f464/YM/Work/F04.lean\#L149-L159}{\texttt{Work.F04.answer}}.

\textbf{Proposition III.6.5 (independence of the gauge-fixing forest;
F05).} Let \(S\) be larger, \(H,H'\) forests, and \(W\) a path datum
with bounds \((p,\sigma,C,\Lambda)\). Let \(\rho,\rho'\) be critical
branches for \(H\) and \(H'\). Assume positive gap constants \(a,b\) for
\(H,H'\), respectively, and let \(\mathscr J\) be a local chart ratio.
Then for every tag \(k\), for \(x\) near \(0\):

\begin{itemize}
\tightlist
\item
  \begin{enumerate}
  \def\labelenumi{(\roman{enumi})}
  \tightlist
  \item
    \(\mathcal T^H_k(x,\rho(x))=\mathcal T^{H'}_k(x,\rho'(x))\);
  \end{enumerate}
\item
  (ii) \(\Theta^H_k(x)=\Theta^{H'}_k(x)\).
\end{itemize}

The gap for \(H\) is unused in the proof. The statement of (ii) allows
an additive constant; the proof gives \(0\). \emph{Status:} theorem.
\emph{Lean:}
\href{https://github.com/yang-mills-group/OpenAI-YM-20261028/blob/3456e5e75cf8d4cd558955a5c7638fba2d68f464/YM/Bid/Station.lean\#L94-L99}{\texttt{Bid.F05}},
\href{https://github.com/yang-mills-group/OpenAI-YM-20261028/blob/3456e5e75cf8d4cd558955a5c7638fba2d68f464/YM/Work/F05.lean\#L60-L67}{\texttt{Work.F05.answer}},
\href{https://github.com/yang-mills-group/OpenAI-YM-20261028/blob/3456e5e75cf8d4cd558955a5c7638fba2d68f464/YM/Work/F05.lean\#L42-L57}{\texttt{Work.F05.same\_thermal}}.

\textbf{Proposition III.6.6 (the stationary branch from F01--F04).} For
every \(G\), \(L\), \(J\), kernel choices \(K\) and \(a>0\) there is
\(\eta>0\) with the following property. Take a larger \(S\) with these
parameters; a forest \(H\); beds \(B,C\); clips \(p,q\); a path datum
\(W\); and towers \(T\). Assume:

\begin{itemize}
\tightlist
\item
  \(T\) satisfies the diagonal identities for \((S,B,C)\);
\item
  \(B\) has floor \((a,A)\) for clip \(p\) on \(\Lambda_f\), and \(C\)
  has floor \((v,u)\) for clip \(q\) on \(\Lambda_c\), for some real
  \(A,v,u\);
\item
  \(W\) has bounds \((p,\sigma,C_0,\Lambda)\) and is attached to \(B\);
\item
  \(B\) is sworn.
\end{itemize}

Then the gap of III.6.3 holds with constant \(\eta\), and there is a
critical branch with properties (ii)--(iii) of III.6.4. \emph{Status:}
theorem.
\href{https://github.com/yang-mills-group/OpenAI-YM-20261028/blob/3456e5e75cf8d4cd558955a5c7638fba2d68f464/YM/Bid/Station.lean\#L108-L113}{\texttt{Bid.cospons}}
derives it from E02 and F01--F04.
\href{https://github.com/yang-mills-group/OpenAI-YM-20261028/blob/3456e5e75cf8d4cd558955a5c7638fba2d68f464/YM/Integrated.lean\#L229-L232}{\texttt{Integrated.stationary\_scene}}
instantiates it with the proofs of E02 and F01--F04, so no work-item
assumption remains. The geometric, floor and path hypotheses above are
retained. \emph{Lean:}
\href{https://github.com/yang-mills-group/OpenAI-YM-20261028/blob/3456e5e75cf8d4cd558955a5c7638fba2d68f464/YM/Bid/Station.lean\#L102-L106}{\texttt{Bid.Scene}}.

\textbf{Proposition III.6.7 (the branch on a volume-independent ball).}

\emph{(a) Uniform implicit function theorem.} Let \(\iota,\tau\) be
finite sets and let \(\ell:\iota\sqcup\tau\to\Lambda\) assign their
coordinates to sites of a four-dimensional torus, with at most \(K_0\)
coordinates per site. Equip the coordinate spaces with the sup norm. Fix
\(C_0,a_0,r,m>0\) and an integer \(n\ge3\). Suppose
\(f:\mathbb R^{\iota\sqcup\tau}\to\mathbb R\) is real-analytic at every
point \(z\) of \(\{\|z\|_\infty<r\}\) and, for \(1\le k\le n\), every
coordinate tuple \(q_1,\ldots,q_k\) and every \(j\le k\),
\[\big|\partial_{q_1}\cdots\partial_{q_k}f(z)\big|\le C_0\exp\!\left(-a_0\sum_{i=1}^k\mathrm{dist}(\ell(q_j),\ell(q_i))\right).\]
Here distance is the periodic max-distance. Assume
\(\nabla_\tau f(0)=0\) and \(D^2_{\tau\tau}f(0)\ge mI\) in counting
coordinates. Then there are positive radii
\(r_{\rm hor}\le r_{\rm zone}\le r_{\rm safe}<r\), depending only on
\((K_0,C_0,a_0,r,m)\), and a map \(g:\mathbb R^\iota\to\mathbb R^\tau\)
with \(g(0)=0\), such that whenever \(\|x\|_\infty<r_{\rm hor}\):

\begin{itemize}
\tightlist
\item
  \(g\) is analytic at \(x\);
\item
  \(\|g(x)\|_\infty\le r_{\rm zone}/2\);
\item
  \(g(x)\) is the unique critical point of \(f(x,\cdot)\) in
  \(\{\|v\|_\infty\le r_{\rm zone}\}\).
\end{itemize}

\emph{(b) Application to the tested action.} Fix \(G\), real constants
\(\mu,C\), a patch order \(J\), kernel choices \(K\), and an integer
\(n\ge3\). There is a choice of auxiliary constants, made before \(L\),
with the following property. For every clip \(p\), real \(\Lambda\),
larger step \(S\) of order \(J\) with kernel choices \(K\), forest \(H\)
and path datum \(W\) with bounds \((p,\mu,C,\Lambda)\), part (a) applies
to the tested action in individual real colour components. It gives a
critical branch \(g\) common to all tags: for each tag \(t\) and
\(\|x\|_\infty<r_{\rm hor}\), \(g(x)\) is the unique critical point of
\(\mathcal T^H_t(x,\cdot)\) in the closed sup-ball of radius
\(r_{\rm zone}\). The radii depend on the auxiliary constants, \(p\) and
\(L\), but not on \(c,f,H,\Lambda\) or \(W\) beyond its bounds. Coarse
coordinates are located at their fine block roots when measuring decay.
\emph{Status:} theorem. \emph{Lean:} (a)
\href{https://github.com/yang-mills-group/OpenAI-YM-20261028/blob/3456e5e75cf8d4cd558955a5c7638fba2d68f464/YM/Push/STATIONARY_UNIFORM/ImplicitRadius.lean\#L179-L215}{\texttt{Receipt.throughout}}
with hypotheses
\href{https://github.com/yang-mills-group/OpenAI-YM-20261028/blob/3456e5e75cf8d4cd558955a5c7638fba2d68f464/YM/Push/STATIONARY_UNIFORM/ImplicitRadius.lean\#L24-L33}{\texttt{Receipt}}
and radii
\href{https://github.com/yang-mills-group/OpenAI-YM-20261028/blob/3456e5e75cf8d4cd558955a5c7638fba2d68f464/YM/Push/STATIONARY_UNIFORM/ImplicitRadius.lean\#L42-L43}{\texttt{reserveX}}.
(b)
\href{https://github.com/yang-mills-group/OpenAI-YM-20261028/blob/3456e5e75cf8d4cd558955a5c7638fba2d68f464/YM/Push/STATIONARY_UNIFORM/Paths.lean\#L253-L294}{\texttt{Before.graph}},
\href{https://github.com/yang-mills-group/OpenAI-YM-20261028/blob/3456e5e75cf8d4cd558955a5c7638fba2d68f464/YM/Push/STATIONARY_UNIFORM/Paths.lean\#L245-L251}{\texttt{Before.Service}},
\href{https://github.com/yang-mills-group/OpenAI-YM-20261028/blob/3456e5e75cf8d4cd558955a5c7638fba2d68f464/YM/Push/STATIONARY_UNIFORM/Paths.lean\#L160-L167}{\texttt{pre\_order}}.

\textbf{Proposition III.6.8 (derivatives along the branch).} Fix an
integer \(N\ge0\), the data and hypotheses of III.6.7(b), auxiliary
constants chosen at order \(N+3\), and any branch \(g\) satisfying its
analytic, size and uniqueness conclusions. Then for every tag the
following hold, with constants depending only on the corresponding
bounds and radii.

\begin{itemize}
\item
  \begin{enumerate}
  \def\labelenumi{(\alph{enumi})}
  \tightlist
  \item
    At every point of the closed joint sup-ball of radius
    \(r_{\rm safe}\ge r_{\rm zone}\), the vertical Hessian, as a matrix
    function of all variables, has a \emph{block package} of order
    \(N\):
  \end{enumerate}

  \begin{itemize}
  \tightlist
  \item
    it is analytic and \(\ge m/2\);
  \item
    its entries and their derivatives up to order \(N+1\) are localized;
  \item
    for every principal sub-block and every shift \(t\ge0\), the shifted
    sub-block is invertible;
  \item
    the inverse of the shifted sub-block is exponentially localized with
    constant \(O((1+t)^{-1})\), has localized derivatives up to order
    \(N+1\), and has \(\ell^2\) norm \(\le(m/2+t)^{-1}\);
  \item
    \(\log\det\) of the shifted sub-block is analytic and its
    derivatives of orders \(1,\ldots,N+1\) are localized, with bounds
    proportional to \((1+t)^{-1}\).
  \end{itemize}
\item
  (b) On the coarse sup-ball of radius \(r_{\rm hor}\), the map \(g\)
  and the graph \(x\mapsto(x,g(x))\) have localized derivatives of
  orders \(1,\ldots,N+2\).
\item
  (c) The classical action \(x\mapsto\mathcal T^H_t(x,g(x))\) has
  localized derivatives of orders \(1,\ldots,N+3\).
\item
  (d) The Hessian along the branch,
  \(x\mapsto\mathrm{Hess}_a\mathcal T^H_t(x,g(x))\), has a block package
  of order \(N\).
\end{itemize}

For a matrix function \(A(z)\), localization of its coefficients means a
bound
\[|\partial_{q_1}\cdots\partial_{q_k}A_{ij}(z)|\le C_k\exp\!\left[-a_k\left(\mathrm{dist}(i,j)+\sum_\nu\mathrm{dist}(i,q_\nu)\right)\right],\]
with coordinates replaced by their assigned sites. For positive-order
derivatives of a map such as \(g\), omit \(j\); for a scalar function
use the bound in III.6.7(a). No volume-independent bound on the value of
the classical action or of \(\log\det\) is asserted.

\emph{Status:} theorem under the hypotheses of III.6.7(b). \emph{Lean:}
\href{https://github.com/yang-mills-group/OpenAI-YM-20261028/blob/3456e5e75cf8d4cd558955a5c7638fba2d68f464/YM/Planned/STATIONARY_JET_PORT/InUse.lean\#L116-L145}{\texttt{Counter.service}},
\href{https://github.com/yang-mills-group/OpenAI-YM-20261028/blob/3456e5e75cf8d4cd558955a5c7638fba2d68f464/YM/Planned/STATIONARY_JET_PORT/InUse.lean\#L75-L88}{\texttt{Counter}},
\href{https://github.com/yang-mills-group/OpenAI-YM-20261028/blob/3456e5e75cf8d4cd558955a5c7638fba2d68f464/YM/Planned/STATIONARY_JET_PORT/Pack.lean\#L65-L84}{\texttt{BlockPack}}.

\textbf{Proposition III.6.9 (one level of the iteration).} Assume the
following.

\begin{itemize}
\tightlist
\item
  The schedule hypotheses of III.2: constants issued before \(L\),
  warm-up, and a paid schedule.
\item
  A horizon of beds \(M_{N,m}\) and towers with the prescribed strip and
  derivative bounds, floors on every torus, diagonal and stack
  identities, and identity diagonal at the cutoff
  (\href{https://github.com/yang-mills-group/OpenAI-YM-20261028/blob/3456e5e75cf8d4cd558955a5c7638fba2d68f464/YM/PREC/Stacks.lean\#L132-L139}{\texttt{Suffices}}).
\item
  All beds are sworn at the periods used. The towers are harmonic: the
  zero-momentum coefficient of their 2-form component is independent of
  the fine residue within a block.
\item
  A gap quote, i.e.~III.6.6 at every level, which the Lean supplies from
  \href{https://github.com/yang-mills-group/OpenAI-YM-20261028/blob/3456e5e75cf8d4cd558955a5c7638fba2d68f464/YM/Extra/ITERATION_NO_LOSS/Post.lean\#L99-L106}{\texttt{Iteration.post}}.
\item
  A local chart ratio.
\item
  The
  \href{https://github.com/yang-mills-group/OpenAI-YM-20261028/blob/3456e5e75cf8d4cd558955a5c7638fba2d68f464/YM/Extra/ITERATION_NO_LOSS/Gates.lean\#L79-L96}{fresh-data
  hypotheses}. For each admissible input datum and each family of
  branches satisfying (i), these hypotheses supply the fresh cubic,
  quartic and loop atoms and their bounds, with the classical and
  thermal matchings each verified on a chosen forest and tag per period.
  These choices may differ between the two matchings. It also supplies
  the cutoff cubic and quartic atoms.
\end{itemize}

Let \(m<N\) and let \(T\) be a path datum \emph{funded} at level \(m+1\)
(\href{https://github.com/yang-mills-group/OpenAI-YM-20261028/blob/3456e5e75cf8d4cd558955a5c7638fba2d68f464/YM/Extra/ITERATION_NO_LOSS/Gates.lean\#L35-L39}{\texttt{Work\ item}}).
That is, the rows of \(T\) enumerate the root-kernel coefficients of
\(M_{N,m+1}\) using the fixed coordinate paths to each offset; \(T\) has
the level-\((m+1)\) bounds; and the sums of absolute cubic and quartic
coefficients weighted by exponential path length are each at most the
level cap. This row convention is stronger than attachment alone. Then
there exist a path datum \(R\) funded at level \(m\), branches
\(\rho_{p,F}\) for every period \(p\) and every forest \(F\), and a
number \(\alpha\) such that:

\begin{itemize}
\tightlist
\item
  \begin{enumerate}
  \def\labelenumi{(\roman{enumi})}
  \tightlist
  \item
    each \(\rho_{p,F}\) satisfies the conclusions of III.6.6 with the
    level-\((m+1)\) gap, its tower map, and the bed \(M_{N,m}\);
  \end{enumerate}
\item
  (ii) \(|\alpha+\kappa_G\log L_m|+1\le\mathrm{drift}_{m+1}\), where
  \(\kappa_G\) is the one-loop slope constant of \(G\);
\item
  (iii) for every \(p\) and \(F\),
  \(x\mapsto\mathcal T^F_k(x,\rho_{p,F}(x))\) (built with \(T\) and the
  clip of level \(m+1\)) has the same jets of orders \(0,\dots,4\) at
  \(0\) as \(x\mapsto\mathcal M_{q}(R;e^x)\) (clip \(q\) of level
  \(m\));
\item
  (iv) for every \(p\) and \(F\), the thermal function is analytic and
  the Hessian determinant is positive near \(0\). Moreover, for any two
  real plane waves \(w_1,w_2\), each the real part of a complex lattice
  plane wave whose coefficient in every axis and colour has modulus at
  most \(1\),
  \[\big|D^2\Theta^F_k(0)[w_1,w_2]-D^2\big(\alpha\,\mathcal M^{(2)}_q(R)+\mathrm{loop}_R\big)(0)[w_1,w_2]\big|\le|\Lambda_c|\,\Gamma\,\lambda_{m+1}\,e^{-\gamma c}.\]
  Here \(\mathcal M^{(2)}_q(R;x)\) denotes the symmetrized
  Wilson-plus-row-square term evaluated on \(e^x\):
  \[\mathcal M^{(2)}_q(R;x)=\frac1{|\mathcal S|}\sum_{s\in\mathcal S}\sum_\pi\left[c_q\,\mathrm{price}((s^*e^x)_\pi)+\tfrac12|\ell_R(s^*e^x)_\pi|^2\right].\]
  It includes all higher-order dependence on \(x\) in these terms; it is
  not their degree-two Taylor polynomial. The constants
  \(\Gamma,\gamma>0\) are chosen after \(L\). All conclusions hold for
  every tag \(k\), and the same \(R,\alpha\) serve every period and
  forest.
\end{itemize}

A companion statement,
\href{https://github.com/yang-mills-group/OpenAI-YM-20261028/blob/3456e5e75cf8d4cd558955a5c7638fba2d68f464/YM/Extra/ITERATION_NO_LOSS/Mechanism.lean\#L35-L48}{\texttt{Iteration.peak}},
gives a funded path datum at the cutoff level \(N\) whose main term has
the jets of order \(\le4\) of the Wilson energy. \emph{Status:} theorem
under the stated schedule, horizon and fresh-data hypotheses. Their
discharge is discussed in §7. \emph{Lean:}
\href{https://github.com/yang-mills-group/OpenAI-YM-20261028/blob/3456e5e75cf8d4cd558955a5c7638fba2d68f464/YM/Extra/ITERATION_NO_LOSS/Mechanism.lean\#L115-L150}{\texttt{Iteration.extend}},
\href{https://github.com/yang-mills-group/OpenAI-YM-20261028/blob/3456e5e75cf8d4cd558955a5c7638fba2d68f464/YM/Extra/ITERATION_NO_LOSS/Mechanism.lean\#L26-L33}{\texttt{Move}},
\href{https://github.com/yang-mills-group/OpenAI-YM-20261028/blob/3456e5e75cf8d4cd558955a5c7638fba2d68f464/YM/PCalc/Chart.lean\#L108-L109}{\texttt{ClassicalMatch}},
\href{https://github.com/yang-mills-group/OpenAI-YM-20261028/blob/3456e5e75cf8d4cd558955a5c7638fba2d68f464/YM/PCalc/Loop.lean\#L62-L71}{\texttt{FiniteBound}}.

\textbf{Further results attached to this unit (statements only; see
§4.7).}

\begin{itemize}
\tightlist
\item
  \emph{Beds and towers (F06--F09).} These are combined in
  \href{https://github.com/yang-mills-group/OpenAI-YM-20261028/blob/3456e5e75cf8d4cd558955a5c7638fba2d68f464/YM/Bid/Craft.lean\#L107-L166}{\texttt{Bid.pinning}}
  (III.10) into a horizon meeting the prescribed bounds. The
  construction retains the supplied beds and the first two tower
  components; the floors and original diagonal identities are inputs,
  and the stack identities are outputs.

  \begin{itemize}
  \tightlist
  \item
    \href{https://github.com/yang-mills-group/OpenAI-YM-20261028/blob/3456e5e75cf8d4cd558955a5c7638fba2d68f464/YM/Work/F06.lean\#L64-L128}{\texttt{F06.answer}}:
    a \(3\)-form kernel that is positive on exact forms can be replaced
    by one that is positive everywhere, without changing
    \(d^\top(\cdot)d\).
  \item
    \href{https://github.com/yang-mills-group/OpenAI-YM-20261028/blob/3456e5e75cf8d4cd558955a5c7638fba2d68f464/YM/Work/F07.lean\#L180-L201}{\texttt{F07.answer}}:
    square roots of positive kernels, with exponential decay.
  \item
    \href{https://github.com/yang-mills-group/OpenAI-YM-20261028/blob/3456e5e75cf8d4cd558955a5c7638fba2d68f464/YM/Work/F08.lean\#L53-L146}{\texttt{F08.answer}}:
    towers with uniform widths.
  \item
    \href{https://github.com/yang-mills-group/OpenAI-YM-20261028/blob/3456e5e75cf8d4cd558955a5c7638fba2d68f464/YM/Work/F09.lean\#L41-L150}{\texttt{F09.answer}}:
    the stack identities.
  \end{itemize}
\item
  \emph{Comparison of two tori}
  (\href{https://github.com/yang-mills-group/OpenAI-YM-20261028/blob/3456e5e75cf8d4cd558955a5c7638fba2d68f464/YM/Further/PHYSICAL_PERIOD_LIMIT_F933/Loom.lean\#L65-L104}{\texttt{Loom.spinning}}).
\item
  \emph{Primary-integral comparisons}
  (\href{https://github.com/yang-mills-group/OpenAI-YM-20261028/blob/3456e5e75cf8d4cd558955a5c7638fba2d68f464/YM/Planned/PRIMARY_INTEGRAL_295/Reaches.lean\#L74-L86}{\texttt{Ramble.billFor}},
  \href{https://github.com/yang-mills-group/OpenAI-YM-20261028/blob/3456e5e75cf8d4cd558955a5c7638fba2d68f464/YM/Planned/PRIMARY_INTEGRAL_295/East.lean\#L112-L139}{\texttt{PPX.toll}},
  \href{https://github.com/yang-mills-group/OpenAI-YM-20261028/blob/3456e5e75cf8d4cd558955a5c7638fba2d68f464/YM/Planned/PRIMARY_INTEGRAL_295/Oval.lean\#L167-L179}{\texttt{Oval.sticker}}).
\item
  \emph{Rotation and Ward material}
  (\href{https://github.com/yang-mills-group/OpenAI-YM-20261028/blob/3456e5e75cf8d4cd558955a5c7638fba2d68f464/YM/Next/ROTATIONS/Route.lean\#L114-L121}{\texttt{Outer.outer}}).
  Mathematically it belongs to VI.1.
\end{itemize}

\subsubsection{2. What it uses}\label{what-it-uses-8}

\begin{itemize}
\item
  \textbf{III.1 (block-spin step and symmetric gauge fixing).}

  \begin{itemize}
  \tightlist
  \item
    The step data, tags and predictors, and the well-posedness
    condition.
  \item
    Spanning forests rooted at block centres.
  \item
    The definitions of the tested action, critical branch, thermal
    function, beds, floors and towers.
  \item
    The identities \(P\widehat{\mathcal W}_H=1-\kappa^{-1}K_C\) and
    \(K_B\widehat{\mathcal W}_H=P^\top K_C\), which follow from the
    tower identities
    (\href{https://github.com/yang-mills-group/OpenAI-YM-20261028/blob/3456e5e75cf8d4cd558955a5c7638fba2d68f464/YM/PREC/TreeMatch.lean\#L81-L91}{\texttt{stationed}}).
  \item
    The change of forest \(\mathrm{slice}(H)\to\mathrm{slice}(H')\), a
    gauge transformation trivial at the roots, preserves product Haar
    measure
    (\href{https://github.com/yang-mills-group/OpenAI-YM-20261028/blob/3456e5e75cf8d4cd558955a5c7638fba2d68f464/YM/PCalc/Profiles.lean\#L57-L77}{\texttt{Profiles.spending}}).
  \item
    The predictor is invariant under gauge transformations trivial at
    the block roots.
  \end{itemize}
\item
  \textbf{III.10 (one nonlinear step).} E02, the shell comparison
  (\href{https://github.com/yang-mills-group/OpenAI-YM-20261028/blob/3456e5e75cf8d4cd558955a5c7638fba2d68f464/YM/Bid/Energy.lean\#L49-L52}{\texttt{Bid.E02}},
  \href{https://github.com/yang-mills-group/OpenAI-YM-20261028/blob/3456e5e75cf8d4cd558955a5c7638fba2d68f464/YM/Bid/Energy.lean\#L44-L47}{\texttt{Cops}}).
  For given \((G,L,J,K)\) there is \(C_{\rm E02}>0\) with the following
  property. For every larger \(S\), forest \(H\), fine configuration
  \(U\) in the slice of \(H\), coarse \(V\), fine bond \(e\) (in block
  \(i\), direction \(\mu\)) and tag \(k\), assume the selected frames
  have positive mass at all coarse sites within periodic offset
  \([-2,2]^4\) of \(i\). Then:
  \[\mathrm{chord}\big(U_e\mathcal R(V)_e^{-1}\big)\le C_{\rm E02}\,\mathrm{port}_i(U)+\mathbf 1_{e\ \mathrm{crossing}}\,\mathrm{wire}(Q_k(U),V)_{(i,\mu)}.\]
  Here \(\mathrm{wire}(Q,V)_b=\mathrm{chord}(V_bQ_b^{-1})\), and
  \(\mathrm{port}_i(U)\le\sum_{z\in X}\sum_\pi\mathrm{chord}(U_\pi)\),
  summed over plaquettes \(\pi\) at \(\mathrm{root}_i+z\). The estimate
  is proved in
  \href{https://github.com/yang-mills-group/OpenAI-YM-20261028/blob/3456e5e75cf8d4cd558955a5c7638fba2d68f464/YM/Work/E02.lean\#L132-L244}{\texttt{Work.E02.answer}}.
  The combination
  \href{https://github.com/yang-mills-group/OpenAI-YM-20261028/blob/3456e5e75cf8d4cd558955a5c7638fba2d68f464/YM/Bid/Station.lean\#L108-L113}{\texttt{Bid.cospons}}
  and the horizon construction
  \href{https://github.com/yang-mills-group/OpenAI-YM-20261028/blob/3456e5e75cf8d4cd558955a5c7638fba2d68f464/YM/Bid/Craft.lean\#L107-L166}{\texttt{Bid.pinning}}
  are filed under III.10.
\item
  \textbf{III.3 (free kernels).} The beds are free lattice kernels in
  Fourier (Bloch/Laurent) form. F06--F09 complete the towers while
  retaining the supplied beds and floors.
\item
  \textbf{III.2 (order of constants).} The constants issued before
  \(L\), the block sizes \(L_m\), the post-\(L\) costs, and the drift
  window, all used in III.6.9. In particular the gap constant is
  post-\(L\) and the constants of III.6.7(b) are chosen before \(L\).
\item
  \textbf{I.1, II.1, II.2.} The lattice, lattice forms and \(d\); the
  group \(G\) and its exponential, chord, price and Haar measure; the
  path data and their bounds.
\item
  \textbf{III.16 (analytic estimates).} Real-analytic functions, Taylor
  jets and iterated derivatives.
\item
  \textbf{Standard mathematics.}

  \begin{itemize}
  \tightlist
  \item
    The analytic implicit and inverse function theorems (Krantz--Parks,
    \emph{The Implicit Function Theorem}, Birkhäuser 2002).
  \item
    Weierstrass's theorem on normally convergent series of holomorphic
    functions.
  \item
    The second-order chain rule at a critical point, so that \(\det\) of
    the Hessian changes by the squared Jacobian.
  \item
    The change-of-variables formula.
  \item
    The Banach fixed-point theorem.
  \item
    The Combes--Thomas estimate (Combes--Thomas, \emph{Comm. Math.
    Phys.} 34 (1973) 251--270).
  \item
    The background-field method in lattice gauge theory: T. Bałaban,
    \emph{Propagators and renormalization transformations for lattice
    gauge theories I, II}, Comm. Math. Phys. 95 (1984), 96 (1984);
    \emph{The variational problem and background fields in
    renormalization group method for lattice gauge theories}, Comm.
    Math. Phys. 102 (1985); \emph{Renormalization group approach to
    lattice gauge field theories I}, Comm. Math. Phys. 109 (1987); J.
    Dimock, \emph{The renormalization group according to Balaban I},
    Rev.~Math. Phys. 25 (2013).
  \end{itemize}
\end{itemize}

\textbf{Imported constants and assumptions.}

\begin{itemize}
\tightlist
\item
  The constant \(C_{\rm E02}(G,L,J,K)\).
\item
  The penalty \(\kappa\) of the kernel choices.
\item
  The floors \((a,A)\) of the beds.
\item
  The one-loop slope \(\kappa_G\), which enters only through the drift
  window of III.6.9 (see III.9).
\item
  The fresh-data, schedule and horizon hypotheses of III.6.9.
\end{itemize}

\subsubsection{3. Strategy}\label{strategy-8}

At fixed coarse field, the renormalization step expands the fine-field
integral around a stationary background. Gauge fixing removes the
blockwise gauge directions. A positive vertical Hessian then gives an
analytic branch, and the tower identities identify its tangent and the
quadratic part of its classical value. The one-loop contribution is the
Hessian determinant corrected by the Haar density and the loop atoms.

The construction uses the penalty
\(\kappa\sum_b\mathrm{price}(e^{x_b}Q_k(U)_b^{-1})\) and fixes the fine
field on a rooted forest inside each block. The shell comparison bounds
the free links by curvature and predictor response. After linearization
and summation, it gives a positive vertical Hessian with a lower bound
independent of the torus volume. The analytic implicit function theorem
produces the branch germ. A change of forest preserves both the
stationary value and the thermal function: the squared Jacobian in the
Hessian determinant cancels the Haar-density ratio.

For a branch on a ball independent of the volume, exponential decay of
the inverse Hessian gives a uniform sup-norm bound. A Newton contraction
produces the branch on that ball; differentiating the stationarity and
resolvent identities controls its localized derivatives. Finally, the
iteration combines inherited atoms with fresh classical and one-loop
data supplied by the hypotheses of III.6.9. This is a penalized
background-field construction in the spirit of Bałaban's method. It
establishes local stationarity and quantitative control near the
identity.

\subsubsection{4. Steps}\label{steps-8}

Lemmas 4.1--4.4 use the hypotheses of III.6.1, Lemmas 4.5--4.7 those of
III.6.3, Lemmas 4.8--4.12 those of III.6.4, and Lemmas 4.13--4.17 those
of III.6.5. Lemmas 4.19--4.20 use III.6.7(a)--(b), respectively, and
Lemmas 4.22--4.24 use III.6.9.

\paragraph{4.1 The second-order jet (F02,
F03)}\label{the-second-order-jet-f02-f03}

\textbf{Lemma 4.1 (holomorphic extension of the main term).} Let \(W\)
have bounds \((p,\sigma,C,\Lambda)\) and let \(U\) be any configuration
on a finite torus. The \(\ell_W\)-series and the sums of cubic and
quartic atoms have holomorphic extensions near the matrix pair
\((U,U^*)\). Their summands are uniformly summable on the ambient ball
of radius \(R_\sigma=e^{\sigma/2}-1\), using the maximum of the linkwise
Frobenius norms. \emph{Sketch:} Treat the link and inverse-link matrices
as independent complex variables. A word of length \(l\) has norm at
most \((1+R_\sigma)^l\|1\|\). A slot of length \(l\) uses words of
length at most \(2l\), so its norm is at most \(C_Ge^{\sigma l}\), with
\(C_G\) independent of \(l\). The coefficient bounds therefore give
normal convergence; a Cauchy estimate on a smaller ball controls the
derivative series. The Wilson term is polynomial and symmetrization is
finite. Composition with \(y\mapsto(A(y),A(y)^*)\), followed by
restriction to real scalars, proves III.6.2.

\emph{Lean:}
\href{https://github.com/yang-mills-group/OpenAI-YM-20261028/blob/3456e5e75cf8d4cd558955a5c7638fba2d68f464/YM/Work/F02/Series.lean\#L57-L74}{\texttt{F02.analytic\_series}},
\href{https://github.com/yang-mills-group/OpenAI-YM-20261028/blob/3456e5e75cf8d4cd558955a5c7638fba2d68f464/YM/Work/F02.lean\#L30-L41}{\texttt{rowExt\_series}},
\href{https://github.com/yang-mills-group/OpenAI-YM-20261028/blob/3456e5e75cf8d4cd558955a5c7638fba2d68f464/YM/Work/F02.lean\#L51-L57}{\texttt{atomExt\_series}},
\href{https://github.com/yang-mills-group/OpenAI-YM-20261028/blob/3456e5e75cf8d4cd558955a5c7638fba2d68f464/YM/Work/F02.lean\#L126-L139}{\texttt{mainTerm\_contDiffAt}}.

\textbf{Lemma 4.2 (the chart and its tangent).} The map
\((x,a)\mapsto U^H(x,a)\) is analytic near \(0\), with \(U^H(0,0)=1\).
The derivative of \(t\mapsto U^H(t(x,a))_e\) at \(0\) is \(s_e\), where
\(s=\iota x+a\). On tree bonds \(s_e=0\), because tree bonds are
internal and \(\iota x\) vanishes there. Consequently
\(\mathrm{incl}(x,a)=(x,\log U^H(x,a))\) is analytic near \(0\), with
\(\exp\circ\log U^H=U^H\) near \(0\) and
\(D\,\mathrm{incl}(0)(x,a)=(x,s)\). \emph{Lean:}
\href{https://github.com/yang-mills-group/OpenAI-YM-20261028/blob/3456e5e75cf8d4cd558955a5c7638fba2d68f464/YM/Work/F03.lean\#L86-L119}{\texttt{F03.live\_rate}},
\href{https://github.com/yang-mills-group/OpenAI-YM-20261028/blob/3456e5e75cf8d4cd558955a5c7638fba2d68f464/YM/Work/F03.lean\#L67-L80}{\texttt{inclusion\_works}},
\href{https://github.com/yang-mills-group/OpenAI-YM-20261028/blob/3456e5e75cf8d4cd558955a5c7638fba2d68f464/YM/Work/F03.lean\#L121-L149}{\texttt{inclusion\_deriv}}.

\textbf{Lemma 4.3 (flat jets).} Let
\(\mathcal F_k(x,A)=\mathcal M_p(W;e^A)+\kappa\sum_b\mathrm{price}(e^{x_b}Q_k(e^A)_b^{-1})\)
be the action in logarithmic coordinates on all fine bonds. Under the
hypotheses of III.6.1:

\begin{itemize}
\tightlist
\item
  \(\mathcal F_k(0)=0\) and \(D\mathcal F_k(0)=0\);
\item
  \(D^2\mathcal F_k(0)[(x,A),(x',A')]=\langle dA,M_BdA'\rangle+\kappa\langle x-PA,x'-PA'\rangle\).
\end{itemize}

\emph{Sketch:} The computation rests on an abstract principle: if \(f\)
is \(C^2\) at \(0\) and, for every \(z\), \(f(tz)/t^2\to\tfrac12Q(z,z)\)
for a symmetric bilinear \(Q\), then \(f(0)=0\), \(Df(0)=0\) and
\(D^2f(0)=Q\). The limits are:

\begin{itemize}
\tightlist
\item
  \(\mathrm{price}((e^{tA})_\pi)/t^2\to\tfrac12|dA_\pi|^2\);
\item
  \(\ell_W(e^{tA})/t\to R_B\,dA\), because \(W\) is attached to \(B\);
\item
  so the quadratic part tends to
  \(\tfrac12\langle dA,(c_p+R_B^\top R_B)dA\rangle=\tfrac12\langle dA,M_BdA\rangle\),
  by the floor identity;
\item
  the cubic and quartic sums are \(o(t^2)\); each slot vanishes at the
  identity and has a linear tangent;
\item
  averaging over \(\mathcal S\) does not change the limit, because \(B\)
  is sworn;
\item
  \(Q_k(e^{tA})\) has derivative \(PA\) at \(t=0\) (this uses
  well-posedness), so the penalty tends to \(\tfrac\kappa2|x-PA|^2\).
\end{itemize}

\emph{Lean:}
\href{https://github.com/yang-mills-group/OpenAI-YM-20261028/blob/3456e5e75cf8d4cd558955a5c7638fba2d68f464/YM/Work/F03/Calculus.lean\#L102-L117}{\texttt{F03.hessian\_of\_limit}},
\href{https://github.com/yang-mills-group/OpenAI-YM-20261028/blob/3456e5e75cf8d4cd558955a5c7638fba2d68f464/YM/Work/F03/Flat.lean\#L23-L38}{\texttt{quadratic\_limit}},
\href{https://github.com/yang-mills-group/OpenAI-YM-20261028/blob/3456e5e75cf8d4cd558955a5c7638fba2d68f464/YM/Work/F03/Flat.lean\#L63-L77}{\texttt{main\_limit}},
\href{https://github.com/yang-mills-group/OpenAI-YM-20261028/blob/3456e5e75cf8d4cd558955a5c7638fba2d68f464/YM/Work/F03/Flat.lean\#L107-L128}{\texttt{noise\_limit}},
\href{https://github.com/yang-mills-group/OpenAI-YM-20261028/blob/3456e5e75cf8d4cd558955a5c7638fba2d68f464/YM/Work/F03/Flat.lean\#L168-L188}{\texttt{flat\_jets}},
\href{https://github.com/yang-mills-group/OpenAI-YM-20261028/blob/3456e5e75cf8d4cd558955a5c7638fba2d68f464/YM/Work/F03/Metric.lean\#L135-L140}{\texttt{metric\_pair}}.
\emph{(proof read; the metric identity and cubic/quartic limit:
statement only)}

\textbf{Lemma 4.4 (transport to the tree slice).} Near \(0\),
\(\mathcal T^H_k=\mathcal F_k\circ\mathrm{incl}\). Since
\(D\mathcal F_k(0)=0\),
\(D^2\mathcal T^H_k(0)[u,u']=D^2\mathcal F_k(0)[D\mathrm{incl}(0)u,D\mathrm{incl}(0)u']\);
no term with \(D^2\mathrm{incl}\) survives. With Lemmas 4.2--4.3 this
gives III.6.1. \emph{Lean:}
\href{https://github.com/yang-mills-group/OpenAI-YM-20261028/blob/3456e5e75cf8d4cd558955a5c7638fba2d68f464/YM/Work/F03.lean\#L151-L156}{\texttt{F03.from\_flat}},
\href{https://github.com/yang-mills-group/OpenAI-YM-20261028/blob/3456e5e75cf8d4cd558955a5c7638fba2d68f464/YM/Work/F03/Calculus.lean\#L120-L142}{\texttt{hessian\_comp}}.

\paragraph{4.2 The gap (F01)}\label{the-gap-f01}

\textbf{Lemma 4.5 (link bound).} Let \(v\in\mathfrak g^{H^{\rm fr}}\)
and \(y=s(0,v)\), i.e.~\(v\) on free bonds and \(0\) on tree bonds.
Assume E02. Then for every fine bond \(e\), in block \(i\) and direction
\(\mu\),
\[|y_e|\le C_{\rm E02}\sum_{z\in X}\sum_{\pi}\big|dy_{(\mathrm{root}_i+z,\pi)}\big|+\big|(Py)_{(i,\mu)}\big| .\]
\emph{Sketch:} Apply E02 with \(V=1\) to \(U_t=e^{ty}\). The
configuration lies in the slice, and for small \(t\) the selected frames
have positive mass. One obtains
\(\mathrm{chord}(U_{t,e})\le C\,\mathrm{port}_i(U_t)+\mathrm{chord}(Q_k(U_t)_{(i,\mu)})\).
Divide by \(|t|\) and let \(t\to0\), using
\(\mathrm{chord}(\gamma(t))/|t|\to|\gamma'(0)|\) along curves through
\(1\). \emph{Lean:}
\href{https://github.com/yang-mills-group/OpenAI-YM-20261028/blob/3456e5e75cf8d4cd558955a5c7638fba2d68f464/YM/Work/F01.lean\#L63-L108}{\texttt{F01.link\_bound}},
\href{https://github.com/yang-mills-group/OpenAI-YM-20261028/blob/3456e5e75cf8d4cd558955a5c7638fba2d68f464/YM/Work/F01.lean\#L22-L28}{\texttt{chord\_limit}}.

\textbf{Lemma 4.6 (summation).} For the extension \(y\) of Lemma 4.5,
\(\sum_e|y_e|^2\le\beta\sum|dy|^2+\gamma\sum|Py|^2\), with
\(\beta,\gamma\) as in III.6.3. \emph{Sketch:}

\begin{itemize}
\tightlist
\item
  Start from \((CZ+Y)^2\le2C^2Z^2+2Y^2\).
\item
  Cauchy--Schwarz over \(X\times\{\text{plaquette orientations}\}\)
  costs \(6|X|\).
\item
  Each coarse site carries \(L^4\) fine sites and \(4\) directions.
\item
  In the re-summation over coarse sites, \(i\mapsto\mathrm{root}_i+z\)
  is injective for fixed \(z\), which costs one more factor \(|X|\).
\end{itemize}

\emph{Lean:}
\href{https://github.com/yang-mills-group/OpenAI-YM-20261028/blob/3456e5e75cf8d4cd558955a5c7638fba2d68f464/YM/Work/F01.lean\#L174-L206}{\texttt{F01.sum\_estimate}},
\href{https://github.com/yang-mills-group/OpenAI-YM-20261028/blob/3456e5e75cf8d4cd558955a5c7638fba2d68f464/YM/Work/F01.lean\#L136-L165}{\texttt{sum\_ports\_sq}}.

\textbf{Lemma 4.7 (the gap).} By III.6.1 on vertical vectors (where
\(x=0\), so the penalty argument is \(-Py\)) and by the floor,
\[D^2_a\mathcal T^H_r(0)[v,v]=\langle dy,M_Bdy\rangle+\kappa|Py|^2\ge a|dy|^2+\kappa|Py|^2 .\]
Hence
\(\eta|v|^2\le\eta(\beta|dy|^2+\gamma|Py|^2)\le a|dy|^2+\kappa|Py|^2\).
This proves III.6.3. \emph{Lean:}
\href{https://github.com/yang-mills-group/OpenAI-YM-20261028/blob/3456e5e75cf8d4cd558955a5c7638fba2d68f464/YM/Work/F01.lean\#L252-L262}{\texttt{F01.hess\_lower}},
\href{https://github.com/yang-mills-group/OpenAI-YM-20261028/blob/3456e5e75cf8d4cd558955a5c7638fba2d68f464/YM/Work/F01.lean\#L218-L230}{\texttt{curl\_floor}},
\href{https://github.com/yang-mills-group/OpenAI-YM-20261028/blob/3456e5e75cf8d4cd558955a5c7638fba2d68f464/YM/Work/F01.lean\#L233-L250}{\texttt{vertical\_hessian}}.

\paragraph{4.3 The branch (F04)}\label{the-branch-f04}

\textbf{Lemma 4.8 (tags agree near the identity).} If \(S\) is well
posed, the predictors \(Q_k\) for different tags coincide on a
neighbourhood of the identity configuration. Hence the functions
\(\mathcal T^H_k\) coincide near \(0\). Near the identity, the frame and
connector rules select their unique minimizers, so the tags do not
affect the predictor. Local logarithms transfer this equality to
arbitrary nearby configurations and then to the axis chart. \emph{Lean:}
\href{https://github.com/yang-mills-group/OpenAI-YM-20261028/blob/3456e5e75cf8d4cd558955a5c7638fba2d68f464/YM/Work/F04/Choices.lean\#L36-L61}{\texttt{F04.prediction\_germ}},
\href{https://github.com/yang-mills-group/OpenAI-YM-20261028/blob/3456e5e75cf8d4cd558955a5c7638fba2d68f464/YM/Work/F04/Choices.lean\#L78-L86}{\texttt{tested\_germ}}.

\textbf{Lemma 4.9 (existence).} The gap makes the vertical Hessian at
\(0\) injective, hence invertible in these finite-dimensional spaces.
The analytic implicit function theorem, applied to
\(D_a\mathcal T^H_{k_0}(x,a)=0\) for one tag \(k_0\), gives an analytic
\(\rho\) with \(\rho(0)=0\). By Lemma 4.8 it is critical for every tag.
\emph{Lean:}
\href{https://github.com/yang-mills-group/OpenAI-YM-20261028/blob/3456e5e75cf8d4cd558955a5c7638fba2d68f464/YM/Work/F04.lean\#L32-L60}{\texttt{F04.invertible}},
\href{https://github.com/yang-mills-group/OpenAI-YM-20261028/blob/3456e5e75cf8d4cd558955a5c7638fba2d68f464/YM/Work/F04.lean\#L62-L86}{\texttt{exists\_branch}}.

\textbf{Lemma 4.10 (test-pair identity).} Assume the tower identities
and the second jet of III.6.1. Then for every \(z=(y,b)\) and coarse
\(x\),
\[D^2\mathcal T^H_k(0)\big[z,(x,\mathrm{rel}_Hx)\big]=\langle dy,M_C\,dx\rangle .\]
\emph{Sketch:}

\begin{itemize}
\tightlist
\item
  \(s(x,\mathrm{rel}_Hx)=\widehat{\mathcal W}_Hx\), since both sides
  vanish on tree bonds.
\item
  The identities of §2 (III.1) give
  \(x-P\widehat{\mathcal W}_Hx=\kappa^{-1}K_Cx\) and
  \(K_B\widehat{\mathcal W}_H=P^\top K_C\).
\item
  So the left side equals
  \(\langle s(z),P^\top K_Cx\rangle+\langle y-Ps(z),K_Cx\rangle=\langle y,K_Cx\rangle\).
\end{itemize}

When \(M_B\) is positive, the tower identities also imply that
\(A=\mathcal Wx\) minimizes
\(\tfrac12\langle dA,M_BdA\rangle+\tfrac\kappa2|x-\bar QA|^2\), with
minimum value \(\tfrac12\langle dx,M_Cdx\rangle\). \emph{Lean:}
\href{https://github.com/yang-mills-group/OpenAI-YM-20261028/blob/3456e5e75cf8d4cd558955a5c7638fba2d68f464/YM/Work/F04/Pairing.lean\#L49-L66}{\texttt{F04.test\_pair}}.

\textbf{Lemma 4.11 (tangent).} Differentiating
\(D_a\mathcal T(x,\rho(x))=0\) at \(0\) gives
\(D^2\mathcal T(0)[(x,D\rho(0)x),(0,v)]=0\). Lemma 4.10 with
\(z=(0,v)\), together with the symmetry of \(D^2\), gives the same
identity with \(\mathrm{rel}_Hx\) in place of \(D\rho(0)x\). Subtracting
and inverting the vertical Hessian gives \(D\rho(0)=\mathrm{rel}_H\).
\emph{Lean:}
\href{https://github.com/yang-mills-group/OpenAI-YM-20261028/blob/3456e5e75cf8d4cd558955a5c7638fba2d68f464/YM/Work/F04.lean\#L88-L126}{\texttt{F04.tangent}}.

\textbf{Lemma 4.12 (value).} Let \(\beta(x)=(x,\rho(x))\). Since
\(D\mathcal T(0)=0\),
\[D^2(\mathcal T\circ\beta)(0)[x,x']=D^2\mathcal T(0)[(x,\mathrm{rel}_Hx),(x',\mathrm{rel}_Hx')]=\langle dx,M_Cdx'\rangle,\]
by Lemma 4.10. This is III.6.4(iii). \emph{Lean:}
\href{https://github.com/yang-mills-group/OpenAI-YM-20261028/blob/3456e5e75cf8d4cd558955a5c7638fba2d68f464/YM/Work/F04.lean\#L128-L147}{\texttt{F04.stationary\_value}}.

\paragraph{4.4 Independence of the forest
(F05)}\label{independence-of-the-forest-f05}

\textbf{Lemma 4.13 (change of forest is a gauge transformation).} For
\((x,a)\) near \(0\) define
\(\mathrm{sw}(x,a)_{e'}=\log\big(\widetilde U_{e'}\,\mathcal R(e^x)_{e'}^{-1}\big)\)
for \(e'\in H'^{\rm fr}\). Here \(\widetilde U\) is the gauge transform
of \(U^H(x,a)\) into the gauge of \(H'\), by a gauge function equal to
\(1\) at the block roots. Then near \(0\):

\begin{itemize}
\tightlist
\item
  \(U^{H'}(x,\mathrm{sw}(x,a))=\widetilde U\);
\item
  \(\mathcal T^H_k(x,a)=\mathcal T^{H'}_k(x,\mathrm{sw}(x,a))\);
\item
  \(\mathrm{loop}_W(U^H(x,a))=\mathrm{loop}_W(U^{H'}(x,\mathrm{sw}(x,a)))\).
\end{itemize}

Gauge invariance holds for each term:

\begin{itemize}
\tightlist
\item
  the atoms are \(\mathrm{Ad}\)-invariant;
\item
  \(\ell_W\) is \(\mathrm{Ad}\)-covariant;
\item
  the price is a class function;
\item
  the predictor is invariant under gauge transformations trivial at the
  roots.
\end{itemize}

\emph{Lean:}
\href{https://github.com/yang-mills-group/OpenAI-YM-20261028/blob/3456e5e75cf8d4cd558955a5c7638fba2d68f464/YM/Work/F05/Actions.lean\#L87-L92}{\texttt{F05.tested\_switch}},
\href{https://github.com/yang-mills-group/OpenAI-YM-20261028/blob/3456e5e75cf8d4cd558955a5c7638fba2d68f464/YM/Work/F05/Actions.lean\#L94-L98}{\texttt{loop\_switch}},
\href{https://github.com/yang-mills-group/OpenAI-YM-20261028/blob/3456e5e75cf8d4cd558955a5c7638fba2d68f464/YM/Work/F05/Actions.lean\#L55-L59}{\texttt{mainTerm\_gauge}},
\href{https://github.com/yang-mills-group/OpenAI-YM-20261028/blob/3456e5e75cf8d4cd558955a5c7638fba2d68f464/YM/Work/F05/Actions.lean\#L66-L75}{\texttt{root\_noise}},
\href{https://github.com/yang-mills-group/OpenAI-YM-20261028/blob/3456e5e75cf8d4cd558955a5c7638fba2d68f464/YM/Work/F05/Transport.lean\#L44-L45}{\texttt{switch}}.

\textbf{Lemma 4.14 (the two graphs correspond).} Near \(0\),
\(\mathrm{sw}(x,\rho(x))=\rho'(x)\). Indeed, \(\mathrm{sw}(x,\rho(x))\)
is a critical point of \(\mathcal T^{H'}_k(x,\cdot)\), by the chain rule
through the local inverse of \(\mathrm{sw}(x,\cdot)\), and both branches
tend to \(0\). Critical branches are unique near \(0\) when the vertical
Hessian at \(0\) is non-degenerate: this is the inverse function theorem
for \((x,a)\mapsto(x,D_a\mathcal T(x,a))\), using the gap of \(H'\).
With Lemma 4.13 this gives III.6.5(i). \emph{Lean:}
\href{https://github.com/yang-mills-group/OpenAI-YM-20261028/blob/3456e5e75cf8d4cd558955a5c7638fba2d68f464/YM/Work/F05/Graphs.lean\#L37-L67}{\texttt{F05.graph\_switch}},
\href{https://github.com/yang-mills-group/OpenAI-YM-20261028/blob/3456e5e75cf8d4cd558955a5c7638fba2d68f464/YM/Work/F05/Calculus.lean\#L112-L150}{\texttt{station\_unique}}.

\textbf{Lemma 4.15 (Hessian determinants).} Write
\(\mathrm{sw}(x,\cdot)=\pi_\epsilon\circ\mathrm{mv}(x,\cdot)\), where
\(\pi_\epsilon\) relabels coordinates along a bijection
\(\epsilon:H^{\rm fr}\to H'^{\rm fr}\); the two forests have the same
number of free bonds. At a critical point the Hessian pulls back by
\(D_a\mathrm{mv}\), so near \(0\)
\[\det\mathrm{Hess}_a\mathcal T^H_k(x,\rho(x))=\big(\det D_a\mathrm{mv}(x,\rho(x))\big)^2\det\mathrm{Hess}_a\mathcal T^{H'}_k(x,\rho'(x)),\]
and the right-hand determinant is positive (by the gap). \emph{Lean:}
\href{https://github.com/yang-mills-group/OpenAI-YM-20261028/blob/3456e5e75cf8d4cd558955a5c7638fba2d68f464/YM/Work/F05/Graphs.lean\#L118-L131}{\texttt{F05.hessian\_ratio}},
\href{https://github.com/yang-mills-group/OpenAI-YM-20261028/blob/3456e5e75cf8d4cd558955a5c7638fba2d68f464/YM/Work/F05/Graphs.lean\#L75-L79}{\texttt{hessian\_pos}},
\href{https://github.com/yang-mills-group/OpenAI-YM-20261028/blob/3456e5e75cf8d4cd558955a5c7638fba2d68f464/YM/Work/F05/Reindex.lean\#L84-L85}{\texttt{move}}.

\textbf{Lemma 4.16 (Haar Jacobian identity).} Near \(0\),
\[\prod_{e\in H^{\rm fr}}J_{\mathscr J}(a_e)=\big|\det D_a\mathrm{mv}(x,a)\big|\prod_{e'\in H'^{\rm fr}}J_{\mathscr J}(\mathrm{sw}(x,a)_{e'}).\]
\emph{Sketch:}

\begin{itemize}
\tightlist
\item
  By the defining property of \(J_{\mathscr J}\) and right invariance of
  Haar measure, \(a\mapsto(e^{a_e}\mathcal R(e^x)_e)_e\) maps the
  density \(\prod_es_{\mathscr J}J_{\mathscr J}(a_e)\,da_e\), restricted
  to the chart, to product Haar measure restricted to the chart image in
  the slice of \(H\).
\item
  The change of forest preserves product Haar measure (III.1).
\item
  By the change-of-variables formula, the two sides therefore have equal
  integrals over every small compact set.
\item
  The forests have the same number of free bonds, so the factors of the
  Haar scale cancel. Continuity and comparison over small closed balls
  give pointwise equality.
\end{itemize}

\emph{Lean:}
\href{https://github.com/yang-mills-group/OpenAI-YM-20261028/blob/3456e5e75cf8d4cd558955a5c7638fba2d68f464/YM/Work/F05/Jacobian.lean\#L129-L142}{\texttt{F05.toll\_switch}}.

\textbf{Lemma 4.17 (cancellation).} If \(B,j,k>0\), \(A=d^2B\) and
\(j=|d|k\), then \(d\ne0\), \(A>0\), and
\(\tfrac12\log A-\log j=\tfrac12\log B-\log k\). With Lemmas 4.13--4.16
and the equality of the loop terms, this gives III.6.5(ii) with constant
\(0\). \emph{Lean:}
\href{https://github.com/yang-mills-group/OpenAI-YM-20261028/blob/3456e5e75cf8d4cd558955a5c7638fba2d68f464/YM/Work/F05.lean\#L42-L57}{\texttt{F05.same\_thermal}}.

\paragraph{4.5 The volume-independent branch and its
jets}\label{the-volume-independent-branch-and-its-jets}

\textbf{Lemma 4.18 (Combes--Thomas bound).} Let \(S\) be a real
symmetric matrix on a finite index set assigned to sites of a
four-dimensional torus, with multiplicity at most \(K_0\). Suppose
\(a_0,M>0\), \(C_0\ge0\), \(|S_{ij}|\le C_0e^{-a_0\mathrm{dist}(i,j)}\),
and \(S\ge MI\). Then
\[|(S^{-1})_{ij}|\le\frac2M e^{-r_{\rm CT}\mathrm{dist}(i,j)},\qquad r_{\rm CT}=\min\left(\frac{a_0}{4},\frac{M}{2(1+4C_0K_0c(a_0/2)/a_0)}\right).\]
Here the site labels are understood in the distances, and
\[c(\alpha)=\frac{16\cdot6!\,e^{2\alpha}}{\alpha^6},\qquad \sum_y e^{-\alpha\mathrm{dist}(x,y)}\le c(\alpha)\quad(\alpha>0),\]
uniformly over the torus side. \emph{Sketch:} Conjugate \(S\) by the
diagonal weight \(e^{r_{\rm CT}\mathrm{dist}(\cdot,j)}\). The difference
has both absolute row and column sums at most \(M/2\). Its quadratic
form is therefore bounded by \((M/2)\|v\|_2^2\), so the conjugated
matrix retains coercivity \(M/2\), although it need not be symmetric.
Applying this inequality to the tilted \(j\)-th inverse column gives the
displayed entry bound.

\emph{Lean:}
\href{https://github.com/yang-mills-group/OpenAI-YM-20261028/blob/3456e5e75cf8d4cd558955a5c7638fba2d68f464/YM/Push/STATIONARY_UNIFORM/Resolvent.lean\#L157-L168}{\texttt{tiltedCharge}},
\href{https://github.com/yang-mills-group/OpenAI-YM-20261028/blob/3456e5e75cf8d4cd558955a5c7638fba2d68f464/YM/Push/STATIONARY_UNIFORM/Resolvent.lean\#L171-L208}{\texttt{inverse\_sealed}}.

\textbf{Lemma 4.19 (uniform implicit function theorem; III.6.7(a)).} In
the notation of III.6.7(a), set \(D=K_0c(a_0)\) and
\[\Pi=1+\frac4m K_0c(r_{\rm CT}),\] where \(r_{\rm CT}\) is the rate of
Lemma 4.18 with gap \(M=m/2\). The radii may be chosen as
\[\begin{aligned}
r_{\rm safe}&=\min\left(\frac r2,\frac{m}{2(1+C_0D^2)}\right),\\
r_{\rm zone}&=\min\left(r_{\rm safe},\frac1{4\Pi(1+C_0D^2)}\right),\\
r_{\rm hor}&=\min\left(r_{\rm zone},\frac{r_{\rm zone}}{4\Pi(1+C_0D)}\right).
\end{aligned}\] \emph{Sketch:} The third-derivative estimate bounds the
change of a Hessian entry by \(r_{\rm safe}C_0D\) times its spatial
decay factor. Summing once more gives absolute row and column sums at
most \(r_{\rm safe}C_0D^2\le m/2\). Thus the vertical Hessian stays
\(\ge(m/2)I\), and its inverse has sup-norm operator bound \(\Pi\)
throughout the closed safety ball. For \(\|x\|_\infty<r_{\rm hor}\), set
\(H_0=D^2_{\tau\tau}f(0)\) and \[Y(v)=v-H_0^{-1}\nabla_\tau f(x,v).\] On
the closed ball of radius \(r_{\rm zone}\),
\(\|DY\|_{\infty\to\infty}\le\Pi r_{\rm zone}C_0D^2\le1/4\), while
\(\|Y(0)\|_\infty\le\Pi\|x\|_\infty C_0D\le r_{\rm zone}/4\). Hence
\(Y\) maps this ball into the ball of half the radius and has a unique
fixed point. The analytic implicit function theorem at each fixed point,
together with uniqueness, makes the resulting branch analytic at every
allowed \(x\). \emph{Lean:}
\href{https://github.com/yang-mills-group/OpenAI-YM-20261028/blob/3456e5e75cf8d4cd558955a5c7638fba2d68f464/YM/Push/STATIONARY_UNIFORM/ImplicitRadius.lean\#L101-L176}{\texttt{Receipt.adjusted}},
\href{https://github.com/yang-mills-group/OpenAI-YM-20261028/blob/3456e5e75cf8d4cd558955a5c7638fba2d68f464/YM/Push/STATIONARY_UNIFORM/ImplicitRadius.lean\#L62-L76}{\texttt{Receipt.snaps}},
\href{https://github.com/yang-mills-group/OpenAI-YM-20261028/blob/3456e5e75cf8d4cd558955a5c7638fba2d68f464/YM/Push/STATIONARY_UNIFORM/ImplicitRadius.lean\#L80-L92}{\texttt{Receipt.priced\_inverse}},
\href{https://github.com/yang-mills-group/OpenAI-YM-20261028/blob/3456e5e75cf8d4cd558955a5c7638fba2d68f464/YM/Push/STATIONARY_UNIFORM/ImplicitRadius.lean\#L179-L215}{\texttt{Receipt.throughout}}.

\textbf{Lemma 4.20 (the tested action satisfies III.6.7(a)).} Take the
auxiliary constants of III.6.7(b). For every clip \(p\), larger \(S\)
with order \(J\) and kernel choices \(K\), forest \(H\) and \(W\) with
bounds \((p,\mu,C,\Lambda)\), the tested action in sup-norm colour
coordinates satisfies:

\begin{itemize}
\tightlist
\item
  analyticity and localized derivatives of orders \(1,\ldots,n\), with
  width \(\mu/(8(n+1)^2(1+L))\), on a ball whose radius is independent
  of \(c\), \(f\) and \(H\);
\item
  the multiplicity bound \(8\dim\mathfrak g\);
\item
  tag-independence on the whole ball and zero vertical gradient at
  \(0\);
\item
  the vertical gap
  \[m=\frac{\min(c_p/|\mathcal S|,\ \kappa)}{1+\beta(L,\ell_L)+2L^4}.\]
\end{itemize}

This is a coercivity estimate for the same vertical Hessian, with no bed
or attachment hypothesis. In it:

\begin{itemize}
\tightlist
\item
  the constant of E02 is replaced by an explicit \(L\)-dependent
  constant \(\ell_L\), obtained by proving the shell comparison directly
  (\href{https://github.com/yang-mills-group/OpenAI-YM-20261028/blob/3456e5e75cf8d4cd558955a5c7638fba2d68f464/YM/Push/STATIONARY_UNIFORM/LCharge.lean\#L145-L251}{\texttt{link\_charge}});
\item
  the floor \(a\) is replaced by \(c_p/|\mathcal S|\), since the
  vertical Hessian is \(\ge(c_p/|\mathcal S|)|dy|^2+\kappa|Py|^2\)
  (\href{https://github.com/yang-mills-group/OpenAI-YM-20261028/blob/3456e5e75cf8d4cd558955a5c7638fba2d68f464/YM/Push/STATIONARY_UNIFORM/HessCost.lean\#L110-L130}{\texttt{wallet}}).
\end{itemize}

\emph{Lean:}
\href{https://github.com/yang-mills-group/OpenAI-YM-20261028/blob/3456e5e75cf8d4cd558955a5c7638fba2d68f464/YM/Push/STATIONARY_UNIFORM/Paths.lean\#L196-L211}{\texttt{Before.receipt}},
\href{https://github.com/yang-mills-group/OpenAI-YM-20261028/blob/3456e5e75cf8d4cd558955a5c7638fba2d68f464/YM/Push/STATIONARY_UNIFORM/HessCost.lean\#L137-L166}{\texttt{symbolic\_gap}},
\href{https://github.com/yang-mills-group/OpenAI-YM-20261028/blob/3456e5e75cf8d4cd558955a5c7638fba2d68f464/YM/Push/STATIONARY_UNIFORM/HessCost.lean\#L132-L133}{\texttt{zeroGap}},
\href{https://github.com/yang-mills-group/OpenAI-YM-20261028/blob/3456e5e75cf8d4cd558955a5c7638fba2d68f464/YM/Push/STATIONARY_UNIFORM/Paths.lean\#L216-L239}{\texttt{Before.full\_pack}}.
\emph{(proof read for the gap and the assembly; the derivative bounds of
the action and the direct shell estimate: statement only)}

\textbf{Lemma 4.21 (derivatives of the inverse and of \(\log\det\);
III.6.8).} Let \(M(y)\) be a finite matrix function analytic at \(y_0\),
with \(M(y_0)\) invertible. Assign its matrix and variable indices to
lattice sites, with matrix-index multiplicity at most \(K_0\). Suppose
its coefficient derivatives through order \(N\) have the localization
bounds of III.6.8 at rate \(a_0>0\), and \(M(y_0)^{-1}\) has
exponentially decaying entries at that rate. Then the derivatives of
\(M^{-1}\) through order \(N\) are localized, with rates and constants
defined recursively from these bounds and \(K_0\). \emph{Sketch:}
Differentiate \[D(M^{-1})=-M^{-1}(DM)M^{-1}.\] The product rule reduces
each new order to lower inverse derivatives and derivatives of \(M\).
Exponential convolution bounds control the internal index sums; each
multiplication reduces the available rate. For a positive determinant,
use \(D\log\det M=\mathrm{tr}(M^{-1}DM)\) and repeat the same estimates.
Under the coercivity hypotheses of III.6.8 this applies to every
principal submatrix with every nonnegative shift. Composing with the
stationary graph gives the estimates along the branch; the joint action
has order \(N+3\), leaving order \(N+1\) after taking two Hessian
derivatives.

\emph{Lean:}
\href{https://github.com/yang-mills-group/OpenAI-YM-20261028/blob/3456e5e75cf8d4cd558955a5c7638fba2d68f464/YM/Planned/STATIONARY_JET_PORT/Lines.lean\#L56-L108}{\texttt{inverse\_lines}},
\href{https://github.com/yang-mills-group/OpenAI-YM-20261028/blob/3456e5e75cf8d4cd558955a5c7638fba2d68f464/YM/Planned/STATIONARY_JET_PORT/Pack.lean\#L88-L108}{\texttt{establish}},
\href{https://github.com/yang-mills-group/OpenAI-YM-20261028/blob/3456e5e75cf8d4cd558955a5c7638fba2d68f464/YM/Planned/STATIONARY_JET_PORT/InUse.lean\#L56-L59}{\texttt{along\_block}},
\href{https://github.com/yang-mills-group/OpenAI-YM-20261028/blob/3456e5e75cf8d4cd558955a5c7638fba2d68f464/YM/Planned/STATIONARY_JET_PORT/InUse.lean\#L95-L112}{\texttt{Counter.background}}.
\emph{(proof read for the inverse induction and the assembly; the
convolution and composition estimates: statement only)}

\paragraph{4.6 One level of the
iteration}\label{one-level-of-the-iteration}

\textbf{Lemma 4.22 (gap quote and branch selection).} Fix a stage index
\(j\). Apply III.6.6 with the block factor and floor bound assigned to
that stage. This gives post-\(L\) gap constants \(\eta_j\)
(\href{https://github.com/yang-mills-group/OpenAI-YM-20261028/blob/3456e5e75cf8d4cd558955a5c7638fba2d68f464/YM/Extra/ITERATION_NO_LOSS/Post.lean\#L99-L106}{\texttt{Iteration.post}}).
At level \(m\), use \(\eta_{J(m+1)}\). Every funded \(T\) at level
\(m+1\) then has, for each period and forest, a branch satisfying
III.6.9(i)
(\href{https://github.com/yang-mills-group/OpenAI-YM-20261028/blob/3456e5e75cf8d4cd558955a5c7638fba2d68f464/YM/Extra/ITERATION_NO_LOSS/Mechanism.lean\#L50-L55}{\texttt{select}}).

\textbf{Lemma 4.23 (the glued datum is funded).} The fresh-data
hypotheses supply fresh cubic and quartic atoms \(t,v\) and fresh loop
atoms \(l\). Let \(R\) be the bundle on the bed \(M_{N,m}\) with:

\begin{itemize}
\tightlist
\item
  cubic atoms: \(t\) joined with the transport of the cubic atoms of
  \(T\) through the towers (``inherited'');
\item
  quartic atoms: \(v\) joined with the inherited quartic atoms;
\item
  loop atoms: \(l\) joined with the inherited loop atoms.
\end{itemize}

Then \(R\) is funded at level \(m\). The inherited cubic and quartic
atoms are charged with relative factors \(L^{-2}\) and \(L^{-4}\), and
the resulting charges stay within \(6\) times the level price.
\emph{Lean:}
\href{https://github.com/yang-mills-group/OpenAI-YM-20261028/blob/3456e5e75cf8d4cd558955a5c7638fba2d68f464/YM/Extra/ITERATION_NO_LOSS/Mechanism.lean\#L70-L112}{\texttt{resume}}.
\emph{(proof read for the assembly; the inherited-atom estimates:
statement only)}

\textbf{Lemma 4.24 (the move).} The fresh one-loop data supply
\(|\alpha+\kappa_G\log L|\le s_{j,k}\), where \(s_{j,k}\) is the
prescribed slope bound for the two stage indices. The schedule gives
\(s_{j,k}+1\le\mathrm{drift}_{m+1}\), proving (ii). The matching (iii)
is proved in three steps:

\begin{itemize}
\tightlist
\item
  On one reference forest per period, orders \(\le2\) follow from
  III.6.4(iii) and the attachment of \(R\) to \(M_{N,m}\)
  (\href{https://github.com/yang-mills-group/OpenAI-YM-20261028/blob/3456e5e75cf8d4cd558955a5c7638fba2d68f464/YM/Extra/ITERATION_NO_LOSS/Calibration.lean\#L66-L94}{\texttt{Calibration.lower}}).
\item
  Orders \(3,4\) split into the inherited part, matched through the
  towers, and the fresh part, supplied by the fresh-data hypotheses.
\item
  Tag-independence and III.6.5 transfer the match to every tag and
  forest
  (\href{https://github.com/yang-mills-group/OpenAI-YM-20261028/blob/3456e5e75cf8d4cd558955a5c7638fba2d68f464/YM/Extra/ITERATION_NO_LOSS/Determinant.lean\#L80-L89}{\texttt{all\_forests}}).
\end{itemize}

The bound (iv) is proved on the reference forest from the fresh
determinant hypothesis, and transferred to every forest by the equality
of thermal functions
(\href{https://github.com/yang-mills-group/OpenAI-YM-20261028/blob/3456e5e75cf8d4cd558955a5c7638fba2d68f464/YM/Extra/ITERATION_NO_LOSS/Determinant.lean\#L91-L103}{\texttt{transplant}}).
\emph{Lean:}
\href{https://github.com/yang-mills-group/OpenAI-YM-20261028/blob/3456e5e75cf8d4cd558955a5c7638fba2d68f464/YM/Extra/ITERATION_NO_LOSS/Mechanism.lean\#L115-L150}{\texttt{extend}}.
\emph{(proof read for the assembly; inherited-atom matching: statement
only)}

\paragraph{4.7 Remaining packages (statements
only)}\label{remaining-packages-statements-only}

\begin{itemize}
\item
  \emph{F06--F09.}

  \begin{itemize}
  \tightlist
  \item
    F06: fix a strip bound \(A\), a positive margin \(v\) and strip
    width \(\eta>0\). A real translation-invariant \(3\)-form kernel
    \(Q\), Hermitian on real Fourier phases and bounded below by \(v\)
    on the image of the \(2\)-form symbol, can be replaced by \(P\)
    bounded above and below on all \(3\)-forms. The replacement
    preserves \(d_2^\top(\cdot)d_2\) on every torus. Its positive lower
    bound, upper bound and strip width depend only on \((A,v,\eta)\).
  \item
    F07: with analogous uniform positivity and strip bounds, a real
    kernel has a symmetric positive square root on every torus, with
    strip width and bounds depending only on the input bounds and the
    finite component space.
  \item
    F08: fix the two clips, penalty, bed strip bounds, a \(3\)-form
    kernel strip bound and a derivative order. Towers whose first two
    components have the prescribed derivative bounds can be completed by
    replacing the remaining components. The completed towers satisfy the
    replacement identities and uniform strip bounds chosen before the
    odd factor \(L\).
  \item
    F09: the replacement identities, floors of the two beds, diagonal
    identities for the original towers and a square-root factorization
    of the defect imply the stack identities for the completed towers.
  \end{itemize}
\item
  \emph{Comparison of two tori}
  (\href{https://github.com/yang-mills-group/OpenAI-YM-20261028/blob/3456e5e75cf8d4cd558955a5c7638fba2d68f464/YM/Further/PHYSICAL_PERIOD_LIMIT_F933/Loom.lean\#L65-L104}{\texttt{Loom.spinning}}).
  Let \(\pi:\beta\to\alpha\) have exactly \(M\) elements in every fibre.
  Assign the indices in \(\beta\) to torus sites with multiplicity at
  most \(K_0\), and suppose distinct elements of a fibre are separated
  by at least \(N\). Let \(F,F'\) be matrix functions on a real normed
  space, analytic at \(0\), with invertible values there and
  \[F(x)E=EF'(x)\quad\text{near }0,\qquad E_{i,j}=\mathbf1_{\pi(i)=j}.\]
  Let \(y,z\) be analytic scalar functions with \(y=Mz\) near \(0\). Fix
  two directions. Suppose the entries of \(F(0)^{-1}\), the two first
  variations and the mixed second variation are bounded by
  \(B_0e^{-d\,\mathrm{dist}}\), \(T_0e^{-d\,\mathrm{dist}}\) and
  \(Q_0e^{-d\,\mathrm{dist}}\), respectively, where \(d>0\) and
  \(B_0,T_0,Q_0\ge0\). The mixed second derivatives of
  \(\tfrac12\log|\det F|-y\) and \(M(\tfrac12\log|\det F'|-z)\) then
  differ by at most \[|\beta|\,C(K_0,T_0,Q_0,B_0,d)e^{-dN/8}.\] For
  positive determinants these are the thermal determinant terms. The
  proof expands the second derivative into two traces. Compatibility
  identifies the small-system trace with sums over covering fibres; the
  error consists of off-diagonal entries between distinct elements of a
  fibre, whose separation is at least \(N\).
\item
  \emph{Primary integral.} This package contains:

  \begin{itemize}
  \tightlist
  \item
    Newton remainder bounds from bounded finite differences;
  \item
    vanishing of a symmetric tensor on lists of pure directions of a
    given total engineering degree when its polynomial component of that
    degree vanishes;
  \item
    under strip bounds for the tower kernel, comparable probe
    amplitudes, equality of the logarithmic probe germs within radius
    \(\chi\), and derivative bounds \(B\) on the true path, a bound
    \(D(1+B)e^{-\sigma\chi/2}\) for the difference of the true and
    forecast probe jets
    (\href{https://github.com/yang-mills-group/OpenAI-YM-20261028/blob/3456e5e75cf8d4cd558955a5c7638fba2d68f464/YM/Planned/PRIMARY_INTEGRAL_295/Reaches.lean\#L74-L86}{\texttt{Ramble.billFor}});
  \item
    triangular structure of the primary transition matrix.
  \end{itemize}

  These comparisons enter the repeatable nonlinear step in III.10.
\item
  \emph{Atom regularity.} Suppose \(t\ge0\), every plaquette incident to
  \(Y\) has chord at most \(t\), the halo of radius \(q+1\) about \(o\)
  lies in \(Y\), and \(2(q+1)<n\). There is a gauge \(g\) such that
  every link \(b\) with \(\mathrm{dist}(o,b.1)\le q\) satisfies
  \(\mathrm{chord}((g\cdot V)_b)\le4t(1+\mathrm{dist}(o,b.1))\).
\item
  \emph{ROTATIONS, WARD\_BIN.} This material is the conclusion of the
  rotation theorem (line-sum property and vanishing of the local
  rotation charge on collinear points) and its chain: moment bounds for
  reflection-positive spectral measures, decay of the Ward integrand,
  and the flux identity. It belongs to VI.1.
\end{itemize}

\subsubsection{5. Where the difficulty
lies}\label{where-the-difficulty-lies-8}

The main estimate is the volume-independent lower bound on the vertical
Hessian. Curvature alone leaves gauge and flat directions; the forest
condition and predictor penalty together control every free link. The
denominator of the gap contains a term proportional to
\(C_{\rm E02}^2L^4|X|^2\), and \(|X|=(2(J+7)L+1)^4\). The gap is chosen
after \(L\); no lower bound uniform in \(L\) is asserted.

A spectral gap gives an \(\ell^2\) inverse bound. To obtain a sup-norm
radius independent of the volume, one also needs the exponential entry
bounds and their summability. The uniform construction uses a direct
shell estimate and the positive Wilson coefficient, so it needs no bed
or attachment hypothesis. It reuses the linearization and counting
estimates of F01, while supplying its own gap and radii.

Forest independence requires both the stationary change-of-coordinates
formula for the Hessian and preservation of product Haar measure. Their
Jacobians cancel exactly, including the Haar normalization scales. The
proof needs only the gap for the target forest. Other unused hypotheses
are the coarse floor in F04, the identity base value in F02, and the
size clauses of ``larger'' in F03--F05.

The conclusions concern a branch near the identity and uniqueness in a
specified sup-ball. Global constrained minimization is not asserted. The
one-loop coefficient and the classical remainders remain inputs to the
iteration through the fresh-data hypotheses in III.6.9.

\subsubsection{6. Notation}\label{notation-8}

{\def\LTcaptype{none} % do not increment counter
\begin{longtable}[]{@{}
  >{\raggedright\arraybackslash}p{(\linewidth - 4\tabcolsep) * \real{0.3333}}
  >{\raggedright\arraybackslash}p{(\linewidth - 4\tabcolsep) * \real{0.3333}}
  >{\raggedright\arraybackslash}p{(\linewidth - 4\tabcolsep) * \real{0.3333}}@{}}
\toprule\noalign{}
\begin{minipage}[b]{\linewidth}\raggedright
Written proof
\end{minipage} & \begin{minipage}[b]{\linewidth}\raggedright
Lean
\end{minipage} & \begin{minipage}[b]{\linewidth}\raggedright
Meaning
\end{minipage} \\
\midrule\noalign{}
\endhead
\bottomrule\noalign{}
\endlastfoot
\(S\), \(\kappa\), \(K\), \(J\), \(L\) &
\href{https://github.com/yang-mills-group/OpenAI-YM-20261028/blob/3456e5e75cf8d4cd558955a5c7638fba2d68f464/YM/RG/Block.lean\#L22-L25}{\texttt{StepData}},
\href{https://github.com/yang-mills-group/OpenAI-YM-20261028/blob/3456e5e75cf8d4cd558955a5c7638fba2d68f464/YM/RG/Patch.lean\#L37-L45}{\texttt{KernelChoices}}
& step datum; penalty; kernel choices; order; block factor \\
larger &
\href{https://github.com/yang-mills-group/OpenAI-YM-20261028/blob/3456e5e75cf8d4cd558955a5c7638fba2d68f464/YM/Bid/Station.lean\#L64}{\texttt{Bid.Larger}}
& well posed, \(100R_S\le f\), \(c>1\) \\
\(Q_k\), \(P\), \(\bar Q\) & \texttt{prediction}, \texttt{probe},
\texttt{mean} & predictor; its linearization; linearized block
average \\
\(H\), \(H^{\rm fr}\) &
\href{https://github.com/yang-mills-group/OpenAI-YM-20261028/blob/3456e5e75cf8d4cd558955a5c7638fba2d68f464/YM/Tree/Pines.lean\#L24-L29}{\texttt{Forest}},
\texttt{Free} & forest gauge; free bonds \\
\(\mathcal R\), \(\iota\) &
\href{https://github.com/yang-mills-group/OpenAI-YM-20261028/blob/3456e5e75cf8d4cd558955a5c7638fba2d68f464/YM/PCalc/Chart.lean\#L35-L36}{\texttt{replica}},
\href{https://github.com/yang-mills-group/OpenAI-YM-20261028/blob/3456e5e75cf8d4cd558955a5c7638fba2d68f464/YM/PCalc/Chart.lean\#L38-L39}{\texttt{linearCopy}}
& coarse field copied onto crossing fine bonds \\
\(U^H(x,a)\), \(a\) &
\href{https://github.com/yang-mills-group/OpenAI-YM-20261028/blob/3456e5e75cf8d4cd558955a5c7638fba2d68f464/YM/PCalc/Chart.lean\#L60-L61}{\texttt{lifting}},
\href{https://github.com/yang-mills-group/OpenAI-YM-20261028/blob/3456e5e75cf8d4cd558955a5c7638fba2d68f464/YM/PCalc/Chart.lean\#L56}{\texttt{AxisChart}}
& fine field in the axis chart \\
\(s\), \(x-Ps\) &
\href{https://github.com/yang-mills-group/OpenAI-YM-20261028/blob/3456e5e75cf8d4cd558955a5c7638fba2d68f464/YM/Bid/Station.lean\#L34-L36}{\texttt{stretch}},
\href{https://github.com/yang-mills-group/OpenAI-YM-20261028/blob/3456e5e75cf8d4cd558955a5c7638fba2d68f464/YM/Bid/Station.lean\#L39}{\texttt{crop}},
\href{https://github.com/yang-mills-group/OpenAI-YM-20261028/blob/3456e5e75cf8d4cd558955a5c7638fba2d68f464/YM/Bid/Station.lean\#L38}{\texttt{bills}}
& linearized fine field; linearized penalty argument \\
\(W\), \(p\), \(c_p\) &
\href{https://github.com/yang-mills-group/OpenAI-YM-20261028/blob/3456e5e75cf8d4cd558955a5c7638fba2d68f464/YM/RG/PathActions.lean\#L150-L154}{\texttt{PathData}},
\href{https://github.com/yang-mills-group/OpenAI-YM-20261028/blob/3456e5e75cf8d4cd558955a5c7638fba2d68f464/YM/RG/PathActions.lean\#L139-L146}{\texttt{Clip}}
& path datum (row series, loop/cubic/quartic atoms); clip; Wilson
coefficient \\
\(\mathcal M_p(W;\cdot)\), \(\mathcal M^{(2)}\) &
\href{https://github.com/yang-mills-group/OpenAI-YM-20261028/blob/3456e5e75cf8d4cd558955a5c7638fba2d68f464/YM/PCalc/Chart.lean\#L27-L28}{\texttt{mainTerm}},
\href{https://github.com/yang-mills-group/OpenAI-YM-20261028/blob/3456e5e75cf8d4cd558955a5c7638fba2d68f464/YM/PCalc/Chart.lean\#L30}{\texttt{secondTerm}}
& symmetrized main term; Wilson-plus-row-square term \\
\(\mathcal T^H_k\) &
\href{https://github.com/yang-mills-group/OpenAI-YM-20261028/blob/3456e5e75cf8d4cd558955a5c7638fba2d68f464/YM/PCalc/Chart.lean\#L81-L84}{\texttt{tested}},
\href{https://github.com/yang-mills-group/OpenAI-YM-20261028/blob/3456e5e75cf8d4cd558955a5c7638fba2d68f464/YM/PCalc/Chart.lean\#L78-L79}{\texttt{noiseTaylor}}
& tested action; penalty term \\
\(B\), \(M_B\), \(R_B\), \(K_B\) &
\href{https://github.com/yang-mills-group/OpenAI-YM-20261028/blob/3456e5e75cf8d4cd558955a5c7638fba2d68f464/YM/PREC/Stacks.lean\#L35-L40}{\texttt{Bed}},
\href{https://github.com/yang-mills-group/OpenAI-YM-20261028/blob/3456e5e75cf8d4cd558955a5c7638fba2d68f464/YM/PREC/TreeMatch.lean\#L68-L69}{\texttt{K}}
& bed; metric; roots; \(d^\top M_Bd\) \\
\(\langle dA,M_B\,dA'\rangle\) &
\href{https://github.com/yang-mills-group/OpenAI-YM-20261028/blob/3456e5e75cf8d4cd558955a5c7638fba2d68f464/YM/Bid/Station.lean\#L31-L32}{\texttt{gaugeFiscal}}
& quadratic form of the bed \\
floor, attached, sworn &
\href{https://github.com/yang-mills-group/OpenAI-YM-20261028/blob/3456e5e75cf8d4cd558955a5c7638fba2d68f464/YM/PREC/Stacks.lean\#L61-L68}{\texttt{Floor}},
\href{https://github.com/yang-mills-group/OpenAI-YM-20261028/blob/3456e5e75cf8d4cd558955a5c7638fba2d68f464/YM/PREC/Stacks.lean\#L44-L46}{\texttt{Attached}},
\href{https://github.com/yang-mills-group/OpenAI-YM-20261028/blob/3456e5e75cf8d4cd558955a5c7638fba2d68f464/YM/Bid/Station.lean\#L54-L55}{\texttt{Sworn}}
& ellipticity of the bed; \(\ell_W\) linearizes to \(R_B\,d\); symmetry
invariance \\
\(T\), \(\mathcal W\), \(\widehat{\mathcal W}_H\) &
\href{https://github.com/yang-mills-group/OpenAI-YM-20261028/blob/3456e5e75cf8d4cd558955a5c7638fba2d68f464/YM/PREC/Stacks.lean\#L82-L86}{\texttt{Towers.Diagonal}},
\href{https://github.com/yang-mills-group/OpenAI-YM-20261028/blob/3456e5e75cf8d4cd558955a5c7638fba2d68f464/YM/PREC/TreeMatch.lean\#L76-L79}{\texttt{Inputs}},
\href{https://github.com/yang-mills-group/OpenAI-YM-20261028/blob/3456e5e75cf8d4cd558955a5c7638fba2d68f464/YM/PREC/TreeMatch.lean\#L26-L27}{\texttt{extracted}}
& towers; linear background map; its forest-gauge version \\
\(\mathrm{rel}_H\) &
\href{https://github.com/yang-mills-group/OpenAI-YM-20261028/blob/3456e5e75cf8d4cd558955a5c7638fba2d68f464/YM/PCalc/Chart.lean\#L93-L95}{\texttt{relative}},
\href{https://github.com/yang-mills-group/OpenAI-YM-20261028/blob/3456e5e75cf8d4cd558955a5c7638fba2d68f464/YM/PCalc/Chart.lean\#L97-L98}{\texttt{Guessed}}
& linear background field in chart coordinates \\
critical branch \(\rho\) &
\href{https://github.com/yang-mills-group/OpenAI-YM-20261028/blob/3456e5e75cf8d4cd558955a5c7638fba2d68f464/YM/PCalc/Chart.lean\#L88-L91}{\texttt{Critical}}
& background field \\
gap \(\eta\) &
\href{https://github.com/yang-mills-group/OpenAI-YM-20261028/blob/3456e5e75cf8d4cd558955a5c7638fba2d68f464/YM/PCalc/Chart.lean\#L100-L101}{\texttt{PriceGap}}
& lower bound of the vertical Hessian at \(0\) \\
second jet; value jet &
\href{https://github.com/yang-mills-group/OpenAI-YM-20261028/blob/3456e5e75cf8d4cd558955a5c7638fba2d68f464/YM/Bid/Station.lean\#L42-L45}{\texttt{Comptroller}},
\href{https://github.com/yang-mills-group/OpenAI-YM-20261028/blob/3456e5e75cf8d4cd558955a5c7638fba2d68f464/YM/Bid/Station.lean\#L83-L84}{\texttt{HeatG}},
\href{https://github.com/yang-mills-group/OpenAI-YM-20261028/blob/3456e5e75cf8d4cd558955a5c7638fba2d68f464/YM/PCalc/Chart.lean\#L105-L106}{\texttt{SameJet}}
& III.6.1; III.6.4(iii); equality of \(j\)-th derivatives \\
\(J_{\mathscr J}\) &
\href{https://github.com/yang-mills-group/OpenAI-YM-20261028/blob/3456e5e75cf8d4cd558955a5c7638fba2d68f464/YM/PCalc/Loop.lean\#L23-L26}{\texttt{ChartRatio}},
\href{https://github.com/yang-mills-group/OpenAI-YM-20261028/blob/3456e5e75cf8d4cd558955a5c7638fba2d68f464/YM/PCalc/Loop.lean\#L28-L38}{\texttt{LocalRatio}}
& Haar density in exponential coordinates \\
\(\Theta^H_k\) &
\href{https://github.com/yang-mills-group/OpenAI-YM-20261028/blob/3456e5e75cf8d4cd558955a5c7638fba2d68f464/YM/PCalc/Loop.lean\#L46-L49}{\texttt{thermal}},
\href{https://github.com/yang-mills-group/OpenAI-YM-20261028/blob/3456e5e75cf8d4cd558955a5c7638fba2d68f464/YM/PCalc/Loop.lean\#L42-L44}{\texttt{hessian}}
& thermal function; Hessian matrix in the colour basis \\
\(\mathrm{sw}\), \(\mathrm{mv}\) &
\href{https://github.com/yang-mills-group/OpenAI-YM-20261028/blob/3456e5e75cf8d4cd558955a5c7638fba2d68f464/YM/Work/F05/Transport.lean\#L44-L45}{\texttt{switch}},
\href{https://github.com/yang-mills-group/OpenAI-YM-20261028/blob/3456e5e75cf8d4cd558955a5c7638fba2d68f464/YM/Work/F05/Reindex.lean\#L84-L85}{\texttt{move}}
& chart change between forests \\
localized derivatives; localized matrix &
\href{https://github.com/yang-mills-group/OpenAI-YM-20261028/blob/3456e5e75cf8d4cd558955a5c7638fba2d68f464/YM/Push/STATIONARY_UNIFORM/Scatter.lean\#L68-L69}{\texttt{Scatter}},
\href{https://github.com/yang-mills-group/OpenAI-YM-20261028/blob/3456e5e75cf8d4cd558955a5c7638fba2d68f464/YM/Push/STATIONARY_UNIFORM/Resolvent.lean\#L41}{\texttt{Sealed}},
\texttt{Cable} & exponential decay of entries and derivatives \\
hypotheses of III.6.7(a) &
\href{https://github.com/yang-mills-group/OpenAI-YM-20261028/blob/3456e5e75cf8d4cd558955a5c7638fba2d68f464/YM/Push/STATIONARY_UNIFORM/ImplicitRadius.lean\#L24-L33}{\texttt{Receipt}}
& analytic, localized, gapped, critical at \(0\) \\
auxiliary constants chosen before \(L\) &
\href{https://github.com/yang-mills-group/OpenAI-YM-20261028/blob/3456e5e75cf8d4cd558955a5c7638fba2d68f464/YM/Push/STATIONARY_UNIFORM/Paths.lean\#L142-L158}{\texttt{Before}}
& radii and costs of the uniform layer \\
block package; jet service &
\href{https://github.com/yang-mills-group/OpenAI-YM-20261028/blob/3456e5e75cf8d4cd558955a5c7638fba2d68f464/YM/Planned/STATIONARY_JET_PORT/Pack.lean\#L65-L84}{\texttt{BlockPack}},
\href{https://github.com/yang-mills-group/OpenAI-YM-20261028/blob/3456e5e75cf8d4cd558955a5c7638fba2d68f464/YM/Planned/STATIONARY_JET_PORT/InUse.lean\#L75-L88}{\texttt{Counter}}
& III.6.8 \\
fresh-data hypotheses &
\href{https://github.com/yang-mills-group/OpenAI-YM-20261028/blob/3456e5e75cf8d4cd558955a5c7638fba2d68f464/YM/Extra/ITERATION_NO_LOSS/Gates.lean\#L79-L96}{\texttt{Rates}}
& cutoff atoms, fresh classical remainders and one-loop matching \\
funded datum; move &
\href{https://github.com/yang-mills-group/OpenAI-YM-20261028/blob/3456e5e75cf8d4cd558955a5c7638fba2d68f464/YM/Extra/ITERATION_NO_LOSS/Gates.lean\#L35-L39}{\texttt{Work\ item}},
\href{https://github.com/yang-mills-group/OpenAI-YM-20261028/blob/3456e5e75cf8d4cd558955a5c7638fba2d68f464/YM/Extra/ITERATION_NO_LOSS/Mechanism.lean\#L26-L33}{\texttt{Move}},
\href{https://github.com/yang-mills-group/OpenAI-YM-20261028/blob/3456e5e75cf8d4cd558955a5c7638fba2d68f464/YM/Extra/ITERATION_NO_LOSS/Gates.lean\#L41-L45}{\texttt{OnShell}}
& III.6.9 \\
\end{longtable}
}

``Price'', ``bill'', ``work item'', ``receipt'', ``toll'' and ``cost''
are numerical bounds and their bookkeeping; they have no economic
meaning.

\subsubsection{7. What remains
unexplained}\label{what-remains-unexplained-8}

\begin{enumerate}
\def\labelenumi{\arabic{enumi}.}
\tightlist
\item
  \textbf{Fresh classical and one-loop inputs.} Proposition III.6.9 is
  conditional on
  \href{https://github.com/yang-mills-group/OpenAI-YM-20261028/blob/3456e5e75cf8d4cd558955a5c7638fba2d68f464/YM/Extra/ITERATION_NO_LOSS/Gates.lean\#L79-L96}{\texttt{Rates}}.
  The adapter
  \href{https://github.com/yang-mills-group/OpenAI-YM-20261028/blob/3456e5e75cf8d4cd558955a5c7638fba2d68f464/YM/Further/ITERATION_RECIPE_PRODUCER/Crown.lean\#L102-L112}{\texttt{IterationRecipe.stocked}}
  supplies the cutoff atoms and retains the classical and quantum
  matching assumptions from
  \href{https://github.com/yang-mills-group/OpenAI-YM-20261028/blob/3456e5e75cf8d4cd558955a5c7638fba2d68f464/YM/Further/ITERATION_RECIPE_PRODUCER/Crown.lean\#L83-L97}{\texttt{Heating}};
  it does not discharge those assumptions. Their discharge in the main
  construction remains unverified here.
\item
  \textbf{Small-field domains.} The comparison between the componentwise
  radii in III.6.7--III.6.8 and the small-field thresholds used
  downstream remains unverified. The relevant branch and derivative
  interfaces are linked in III.6.7--III.6.8.
\item
  \textbf{Iteration over all levels.} The induction beyond one step is
  in
  \href{https://github.com/yang-mills-group/OpenAI-YM-20261028/blob/3456e5e75cf8d4cd558955a5c7638fba2d68f464/YM/Extra/ITERATION_NO_LOSS/Planned.lean\#L18-L35}{\texttt{Iteration.sweep}}
  and
  \href{https://github.com/yang-mills-group/OpenAI-YM-20261028/blob/3456e5e75cf8d4cd558955a5c7638fba2d68f464/YM/FullVolume/Luna/Table.lean\#L166-L172}{\texttt{Luna.uniformZero}}
  (III.12). Their proofs are outside this verification.
\item
  \textbf{Supporting proofs not checked in full.} The remaining unread
  supporting proofs are the metric/path-limit estimates identified after
  Lemma 4.3; the direct shell and action-derivative estimates used in
  Lemma 4.20; the convolution and graph-composition estimates used in
  Lemma 4.21; and the inherited-atom matching and bounds in Lemmas
  4.23--4.24. F06--F09 and the primary-integral, atom-regularity,
  rotation and Ward packages are represented by their statements only.
  The period-comparison trace argument was checked; its geometric
  realization for the actual Hessians was not.
\item
  \textbf{Attached packages.} The primary-integral package feeds the
  repeatable step in III.10; the rotation and Ward material belongs with
  VI.1. The application of the period comparison in
  \href{https://github.com/yang-mills-group/OpenAI-YM-20261028/blob/3456e5e75cf8d4cd558955a5c7638fba2d68f464/YM/Further/PHYSICAL_PERIOD_LIMIT_F933/Fresco.lean\#L98-L148}{\texttt{Fresco.subscribed}}
  remains unverified here. The size figures in the header are the
  inherited dependency-closure counts; they have not been recounted.
\end{enumerate}

\subsection{III.7. Matching the Wilson energy to fourth
order}\label{iii.7.-matching-the-wilson-energy-to-fourth-order}

\textbf{Part:} III, Ultraviolet construction (renormalization group).
\textbf{Lean:}
\href{https://github.com/yang-mills-group/OpenAI-YM-20261028/tree/3456e5e75cf8d4cd558955a5c7638fba2d68f464/YM/Next/STATIONARY}{\texttt{YM/Next/STATIONARY/}},
\href{https://github.com/yang-mills-group/OpenAI-YM-20261028/tree/3456e5e75cf8d4cd558955a5c7638fba2d68f464/YM/WilsonBulk/Spectral}{\texttt{YM/WilsonBulk/Spectral/}},
\href{https://github.com/yang-mills-group/OpenAI-YM-20261028/tree/3456e5e75cf8d4cd558955a5c7638fba2d68f464/YM/UVExtension/Projection}{\texttt{YM/UVExtension/Projection/}},
and a few lines of
\href{https://github.com/yang-mills-group/OpenAI-YM-20261028/tree/3456e5e75cf8d4cd558955a5c7638fba2d68f464/YM/UVExtension/Reward}{\texttt{YM/UVExtension/Reward/}}.
\textbf{Size:} 19,945 lines in the main-theorem closure: 19,391 in
\texttt{Next/STATIONARY}, 286 in \texttt{UVExtension/Projection}, 215 in
\texttt{WilsonBulk/Spectral}, and about 50 elsewhere.

\emph{Conventions.} Levels \(m=0,\dots,N\) run from the coarsest level
(\(m=0\)) to the Wilson lattice (\(m=N\)). Blocking from fine level
\(m+1\) to coarse level \(m\) uses an odd factor \(L_m\). The stage
\(J_m\) is the index of the constants used at level \(m\); it satisfies
\(J_{m+1}\in\{J_m,J_m+1\}\). In one step the fine torus \(\Lambda_f\)
has side \(f=Lc\) and the coarse torus \(\Lambda_c\) has side \(c\); in
general \(|\Lambda_n|=n^4\). Fields are Lie-algebra link variables
\(x\), and configurations are \(U=e^x\).
\(E_W(U)=\sum_p\big(1-d_{\rm rep}^{-1}\mathrm{Re}\,\mathrm{tr}\,U_p\big)\)
is the Wilson energy in the fixed matrix representation of dimension
\(d_{\rm rep}\), \(C_G\) is the adjoint Casimir constant of I.1, and
\(\kappa=11C_G/(24\pi^2)\).

\subsubsection{1. What this unit proves}\label{what-this-unit-proves-9}

At level \(m\) the effective action has the form
\(\beta\,\mathrm{main}_{\mathbf c,R}(U)+\mathrm{loop}_R(U)\), plus
polymer activities (III.10). Here \(R\) is a \emph{path datum}: an
exponentially weighted, explicit sum of gauge-covariant curvature
monomials along lattice paths. This unit constructs the path data and
the coupling shifts.

\begin{itemize}
\tightlist
\item
  \textbf{(A)} At the finest level there is a path datum whose leading
  part has the same Taylor jet as \(E_W\) through order four (III.7.2).
\item
  \textbf{(B)} For one block-spin step, the coarse leading part can be
  chosen as the classical (tree-level) transform of the fine one through
  order four. The second-order Taylor coefficient of the one-loop term
  equals \(\alpha\) times the coarse quadratic form plus coarse two-leg
  atoms, up to an error exponentially small in the side of the torus.
  Moreover \(|\alpha+\kappa\log L|\) is bounded independently of \(L\)
  (III.7.4).
\item
  \textbf{(C)} The value \(\kappa=11C_G/(24\pi^2)\) is computed, not
  assumed. The one-loop determinant of the background-field operator in
  Feynman gauge, regularized in the infrared at scale \(1/L\), has
  \(F^2\)-coefficient \(-\kappa\log L+O(1)\) (III.7.3). Two results
  connect this determinant to the blocked one: an exact Faddeev--Popov
  identity (III.7.6) and uniform invertibility of the completed
  fluctuation operators (III.7.5, §4.6).
\end{itemize}

Iterated over all levels, (A)--(C) give the kinetic data satisfying the
conditions required in III.9. The slope \(\kappa\), a defined constant
in I.1 and III.9.1, is identified here with the one-loop coefficient of
the blocked action.

\textbf{Proposition III.7.1 (path data, beds, Fourier kernels).}

\begin{itemize}
\item
  \emph{Legs and atoms.} A leg is a lattice walk \(\gamma\) from the
  origin of \(\mathbb Z^4\) together with a plaquette orientation \(P\).
  At base point \(x\) it evaluates \(U\) to
  \(\mathrm{Ad}(U_\gamma)\,\mathcal F(U_{\partial p})\in\mathfrak g\).
  Here \(p\) is the plaquette of orientation \(P\) at the end of
  \(\gamma\), \(U_\gamma\) is the transport along \(\gamma\), and
  \(\mathcal F(u)\) is the \(\mathfrak g\)-component of
  \((u-u^{-1})/2\). The length of a leg is \(|\gamma|+4\).
\item
  A \(k\)-atom is a \(k\)-tuple of legs with a coefficient tensor
  \(a:\{1,\dots,\dim\mathfrak g\}^k\to\mathbb R\). It evaluates to
  \(\sum_b a(b)\prod_i(\mathrm{leg}_i)_{b_i}\). It is \emph{invariant}
  if this polynomial is \(\mathrm{Ad}\)-invariant. Its exponential
  weight at width \(\sigma\) is
  \(\|a\|_1e^{\sigma\cdot\mathrm{length}}\).
\item
  \emph{Path datum} \(R\). For each of the six plaquette orientations
  \(P\), a sequence of rows \((\lambda_{P,k},\mathrm{leg}_{P,k})\); and
  sequences of 2-atoms (\emph{loop atoms}), 3-atoms and 4-atoms. The
  rows define
  \(\ell_R(U)_{(x,P)}=\sum_k\lambda_{P,k}\,\mathrm{leg}_{P,k}(U,x)\).
\item
  \emph{Main term.} For a clip \(\mathbf c\) (only \(c_{\rm dir}>0\)
  enters here), whenever the row and atom series converge absolutely:
  \[\mathrm{main}_{\mathbf c,R}(U)=\frac1{384}\sum_{s}\Big[\sum_p\Big(c_{\rm dir}\big(1-d_{\rm rep}^{-1}\mathrm{Re}\,\mathrm{tr}\,U_p\big)+\tfrac12|\ell_R(U)_p|^2\Big)+\sum_x\sum_k\big(\mathcal A^{(3)}_k+\mathcal A^{(4)}_k\big)(U,x)\Big]_{U\to sU}.\]
  This is an average over the 384 signed permutations \(s\) of the axes.
  \(\mathrm{loop}_R\) is the same average of the sum of all loop atoms.
  These are the unclipped analytic germs; the small-field clipping of
  the full action does not enter them.
\item
  \emph{Weight clauses}
  \(R\in\mathrm{Bounds}(\mathbf c;\sigma,C,\Lambda)\):

  \begin{itemize}
  \tightlist
  \item
    \(\sigma>0\), and
    \(\sum_k|\lambda_{P,k}|e^{\sigma\,\mathrm{len}}\le C\) for each
    \(P\);
  \item
    the sums of exponential weights are \(\le C\) for the cubic and
    quartic atoms and \(\le\Lambda\) for the loop atoms;
  \item
    all atoms are invariant;
  \item
    normalization at zero momentum:
    \(c_{\rm dir}\delta_{PQ}+\sum_{P'}\hat\ell_{P'P}(0)\hat\ell_{P'Q}(0)=\delta_{PQ}\);
  \item
    the \emph{moment conditions}: for each pair of leg orientations and
    colours, the loop coefficients have total weight \(0\) and first
    moment \(0\) in the displacement between the two legs (with the
    required summability).
  \end{itemize}
\item
  \emph{Bed} \(B=(D,\tilde D,H,\tilde H,\rho)\): five
  translation-invariant kernels on plaquette fields (diagonal, its
  inverse, metric, its inverse, roots).

  \begin{itemize}
  \tightlist
  \item
    \(R\) is \emph{attached} to \(B\) if its rows, summed by end
    plaquette and displacement, give \(\rho\).
  \item
    \(\mathrm{Floor}(B,\mathbf c,n;\varepsilon,C)\) on \(\Lambda_n\)
    means: \(D\tilde D=1\); \(H\tilde H=1\); \(\Phi^{\mathsf T}\Phi=H\)
    with \(\Phi=(\sqrt{c_{\rm dir}}\,\mathbf 1;\rho)\); \(\hat H(0)=1\);
    \(H\,d_1=D\,d_1\) on link fields; \(H,D\) symmetric with
    \(\varepsilon I\le H,D\le CI\) as quadratic forms.
  \item
    A kernel is in \(\mathrm{strip}(\sigma,C)\) if
    \(\sum_y e^{\sigma|y|_1}|K(y)|\le C\) entrywise.
  \end{itemize}
\item
  \emph{Fourier kernels.} A \((t{+}1)\)-leg kernel assigns a number
  \(\hat K_{b,a}(p_1,\dots,p_t)\) to colours \(b\), directions \(a\) and
  momenta, with \(p_0=-\sum p_j\).

  \begin{itemize}
  \tightlist
  \item
    \(\hat K\) is \emph{strip-analytic of width \(\sigma\) and size
    \(K\)} if it is holomorphic, \(2\pi\)-periodic in each scalar
    momentum coordinate, and bounded by \(K\) on
    \(|\mathrm{Im}\,p_{j,a}|<\sigma\).
  \item
    A function \(F\) on link fields on \(\Lambda_n\) is \emph{recorded}
    by \(\hat K\) with error \(\epsilon\) if its complexified Taylor
    coefficient \(D^{t+1}F(0)/(t+1)!\), evaluated on canonical unit
    plane waves, differs from \(|\Lambda_n|\hat K_{b,a}(p)\) when the
    momenta sum to zero, and from \(0\) otherwise, by at most
    \(\epsilon\). All tuples of lattice momenta are tested.
  \item
    \(F\) \emph{vanishes to order \(s\) along rays} if \(F(tv)=O(t^s)\)
    as \(t\to0\) for every \(v\), with the constant allowed to depend on
    \(v\).
  \end{itemize}
\item
  \emph{\(F^2\)-coefficient.} For a two-leg kernel, \(\alpha_b(\hat K)\)
  is twice the real part of the coefficient of \(s^2\) in
  \(\hat K_{(b,b),(1,1)}(s\,e_0)\). For the quadratic kernel of a bed
  whose metric has positive strip width and \(\hat H(0)=1\), it equals
  \(1\), so \(\alpha\) is measured in units of the classical \(F^2\)
  term.
\end{itemize}

\emph{Status:} definitions. Path data, weight clauses, main term, beds
and Floor are specifications of III.1; the Fourier notions are defined
in this unit. \emph{Lean:}
\href{https://github.com/yang-mills-group/OpenAI-YM-20261028/blob/3456e5e75cf8d4cd558955a5c7638fba2d68f464/YM/RG/PathActions.lean\#L150-L154}{\texttt{RG.PathData}},
\href{https://github.com/yang-mills-group/OpenAI-YM-20261028/blob/3456e5e75cf8d4cd558955a5c7638fba2d68f464/YM/RG/PathActions.lean\#L38-L40}{\texttt{RG.Slot}},
\href{https://github.com/yang-mills-group/OpenAI-YM-20261028/blob/3456e5e75cf8d4cd558955a5c7638fba2d68f464/YM/RG/PathActions.lean\#L84-L86}{\texttt{RG.Atom}},
\href{https://github.com/yang-mills-group/OpenAI-YM-20261028/blob/3456e5e75cf8d4cd558955a5c7638fba2d68f464/YM/RG/PathActions.lean\#L173-L185}{\texttt{PathData.Bounds}},
\href{https://github.com/yang-mills-group/OpenAI-YM-20261028/blob/3456e5e75cf8d4cd558955a5c7638fba2d68f464/YM/RG/PathActions.lean\#L161-L162}{\texttt{metric\_zero}},
\href{https://github.com/yang-mills-group/OpenAI-YM-20261028/blob/3456e5e75cf8d4cd558955a5c7638fba2d68f464/YM/RG/PathActions.lean\#L167-L169}{\texttt{loop\_jets}},
\href{https://github.com/yang-mills-group/OpenAI-YM-20261028/blob/3456e5e75cf8d4cd558955a5c7638fba2d68f464/YM/RG/PathActions.lean\#L139-L146}{\texttt{RG.Clip}},
\href{https://github.com/yang-mills-group/OpenAI-YM-20261028/blob/3456e5e75cf8d4cd558955a5c7638fba2d68f464/YM/PCalc/Chart.lean\#L27-L28}{\texttt{Precision.mainTerm}},
\href{https://github.com/yang-mills-group/OpenAI-YM-20261028/blob/3456e5e75cf8d4cd558955a5c7638fba2d68f464/YM/PREC/Stacks.lean\#L35-L40}{\texttt{Precision.Bed}},
\href{https://github.com/yang-mills-group/OpenAI-YM-20261028/blob/3456e5e75cf8d4cd558955a5c7638fba2d68f464/YM/PREC/Stacks.lean\#L44-L46}{\texttt{Attached}},
\href{https://github.com/yang-mills-group/OpenAI-YM-20261028/blob/3456e5e75cf8d4cd558955a5c7638fba2d68f464/YM/PREC/Stacks.lean\#L61-L68}{\texttt{Floor}},
\href{https://github.com/yang-mills-group/OpenAI-YM-20261028/blob/3456e5e75cf8d4cd558955a5c7638fba2d68f464/YM/Next/STATIONARY/Infrastructure/Jets.lean\#L199}{\texttt{Firm}},
\href{https://github.com/yang-mills-group/OpenAI-YM-20261028/blob/3456e5e75cf8d4cd558955a5c7638fba2d68f464/YM/Next/STATIONARY/Infrastructure/Jets.lean\#L28-L29}{\texttt{Rays}},
\href{https://github.com/yang-mills-group/OpenAI-YM-20261028/blob/3456e5e75cf8d4cd558955a5c7638fba2d68f464/YM/Next/STATIONARY/Lift/Pathways.lean\#L491-L492}{\texttt{Recorded}},
\href{https://github.com/yang-mills-group/OpenAI-YM-20261028/blob/3456e5e75cf8d4cd558955a5c7638fba2d68f464/YM/Next/STATIONARY/Films/LeadingRay.lean\#L699}{\texttt{alphaPay}},
\href{https://github.com/yang-mills-group/OpenAI-YM-20261028/blob/3456e5e75cf8d4cd558955a5c7638fba2d68f464/YM/Next/STATIONARY/Films/LeadingRay.lean\#L700-L708}{\texttt{tip\_fuel}}.

\textbf{Proposition III.7.2 (the Wilson jet at the finest level).} Let
\(G\) be a gauge group, \(\sigma_0>0\) and \(A,Z\ge0\). There are
\(\sigma\in(0,\sigma_0]\) and \(C\ge0\) with the following property.

\emph{Hypotheses:} sides \(n_i\to\infty\); a clip with
\(c_{\rm dir}\le Z\); a bed \(B\) such that

\begin{itemize}
\tightlist
\item
  \(\rho\in\mathrm{strip}(\sigma_0,A)\);
\item
  \(\mathrm{Floor}(B,\mathbf c,n;a,b)\) holds for every side \(n\) and
  some \(a,b\);
\item
  the quadratic gauge form of \(B\) is invariant under signed
  permutations on each \(\Lambda_{n_i}\);
\item
  \(D=1\) on each \(\Lambda_{n_i}\).
\end{itemize}

\emph{Conclusion:} there is a path datum \(Q\) such that:

\begin{enumerate}
\def\labelenumi{\arabic{enumi}.}
\tightlist
\item
  the rows of \(Q\) are the entries of \(\rho\) placed on straight
  lattice routes, \(Q\) is attached to \(B\), and all loop coefficients
  vanish;
\item
  \(Q\in\mathrm{Bounds}(\mathbf c;\tau,k,\Lambda)\) for all
  \(0<\tau\le\sigma\), all \(k\ge\max(C,6Ae^{4\sigma_0})\) and all
  \(\Lambda\ge0\);
\item
  for every \(i\) and every \(j\le4\), the functions
  \(x\mapsto\mathrm{main}_{\mathbf c,Q}(e^x)\) and \(x\mapsto E_W(e^x)\)
  on \(\Lambda_{n_i}\) are real-analytic at \(0\) and have equal
  \(j\)-th derivatives there.
\end{enumerate}

\emph{Status:} theorem. \emph{Lean:}
\href{https://github.com/yang-mills-group/OpenAI-YM-20261028/blob/3456e5e75cf8d4cd558955a5c7638fba2d68f464/YM/Next/STATIONARY/Shape/Scaled.lean\#L270-L305}{\texttt{Stationary.launch}},
\href{https://github.com/yang-mills-group/OpenAI-YM-20261028/blob/3456e5e75cf8d4cd558955a5c7638fba2d68f464/YM/Next/STATIONARY/Level/SliceGate.lean\#L169-L171}{\texttt{rowPath}},
\href{https://github.com/yang-mills-group/OpenAI-YM-20261028/blob/3456e5e75cf8d4cd558955a5c7638fba2d68f464/YM/Next/STATIONARY/Level/SliceGate.lean\#L176}{\texttt{rowPrice}},
\href{https://github.com/yang-mills-group/OpenAI-YM-20261028/blob/3456e5e75cf8d4cd558955a5c7638fba2d68f464/YM/PCalc/Chart.lean\#L105-L106}{\texttt{SameJet}}.

\emph{All stages and all volumes.}

\begin{itemize}
\tightlist
\item
  With stage widths and amplitude bounds given, the statement holds at
  every stage with stage constants \((\sigma_j,A_j)\)
  (\href{https://github.com/yang-mills-group/OpenAI-YM-20261028/blob/3456e5e75cf8d4cd558955a5c7638fba2d68f464/YM/Next/STATIONARY/Level/Recurse.lean\#L101-L105}{\texttt{beckon}},
  \href{https://github.com/yang-mills-group/OpenAI-YM-20261028/blob/3456e5e75cf8d4cd558955a5c7638fba2d68f464/YM/Next/STATIONARY/Level/Recurse.lean\#L93-L99}{\texttt{BasesFine}}).
\item
  Take a compatible history satisfying the thresholds of III.2 and a
  horizon satisfying the free-field conditions of §2, and apply the
  statement to the sequence of periodic sides at level \(N\). The result
  is a top datum
  (\href{https://github.com/yang-mills-group/OpenAI-YM-20261028/blob/3456e5e75cf8d4cd558955a5c7638fba2d68f464/YM/Next/STATIONARY/Level/Recurse.lean\#L220-L225}{\texttt{budding}},
  \href{https://github.com/yang-mills-group/OpenAI-YM-20261028/blob/3456e5e75cf8d4cd558955a5c7638fba2d68f464/YM/Next/STATIONARY/Level/Recurse.lean\#L166}{\texttt{BranchRow}})
  that:

  \begin{itemize}
  \tightlist
  \item
    satisfies the level-\(N\) weight clauses with loop amplitude
    \(\mathrm{loopCap}_N\);
  \item
    matches \(E_W\) through order four on every volume;
  \item
    has vanishing loop coefficients.
  \end{itemize}

  This is the top condition used in III.9.
\end{itemize}

\textbf{Proposition III.7.3 (the one-loop slope).} Let \(G\) be simple.
For \(\sigma,\varepsilon>0\), \(C,Z\ge0\), \(\tau>0\) and \(A\ge0\)
there is \(D\ge0\) with the following property.

\emph{Hypotheses:}

\begin{itemize}
\tightlist
\item
  \(L\ge1\) is odd, \(\mu=1/L\), and \(n_i\to\infty\);
\item
  \(\mathbf c\) is a clip with \(c_{\rm dir}\le Z\);
\item
  \(B\) is a bed with \(H\in\mathrm{strip}(\sigma,C)\),
  \(\mathrm{Floor}(B,\mathbf c,n;\varepsilon,C)\) for every \(n\), and
  gauge form invariant under signed permutations on each
  \(\Lambda_{n_i}\);
\item
  for any \(\Lambda\in\mathbb R\), a datum
  \(R\in\mathrm{Bounds}(\mathbf c;\tau,A,\Lambda)\) is attached to
  \(B\).
\end{itemize}

For a link field \(x\) on \(\Lambda_{n_i}\) and \(U=e^x\) put
\[\Pi_\mu(x)=\tfrac12\log\det\big(\mathcal H_{\mathbf c,R}(U)+d_Ud_U^{*}+\mathcal M_\mu(U)\big)-\log\det\big(d_U^{*}d_U+\mathfrak m_\mu(U)\big).\]
Here:

\begin{itemize}
\tightlist
\item
  \(\mathcal H_{\mathbf c,R}(U)\) is the Hessian at \(y=0\) of
  \(y\mapsto\mathrm{main}_{\mathbf c,R}(e^yU)\);
\item
  \(d_U\) is the covariant difference
  \((d_U\phi)_e=\phi(e_-)-\mathrm{Ad}(U_e)\phi(e_+)\) from adjoint site
  fields to link fields;
\item
  \(\mathfrak m_\mu(U)=\mu^2\big(\mu^2(d_U^*d_U+\mu^2)^{-1}\big)^{9}\);
\item
  \(\mathcal M_\mu\) applies \(\mathfrak m_\mu\) separately to the four
  link components, each regarded as a site field at the initial
  endpoint.
\end{itemize}

\emph{Conclusion:} there are a strip-analytic two-leg kernel \(\hat Y\)
and numbers \(e_i\to0\) such that:

\begin{itemize}
\tightlist
\item
  \begin{enumerate}
  \def\labelenumi{(\alph{enumi})}
  \tightlist
  \item
    \(\hat Y\) records the second-order Taylor coefficient of
    \(\Pi_\mu\) at \(0\) on \(\Lambda_{n_i}\) with error
    \(|\Lambda_{n_i}|\,e_i\), for every \(i\);
  \end{enumerate}
\item
  (b) \(|\alpha_b(\hat Y)+\kappa\log L|\le D\) for every colour \(b\).
\end{itemize}

\(\Pi_\mu\) is the one-loop effective action of the fine leading part in
the background \(U\), in Feynman background gauge: \(d_Ud_U^*\) is the
gauge-fixing term and \(d_U^*d_U\) the ghost operator. The infrared
regulator of mass \(\mu=1/L\) stands in for the block-spin constraint.
\(D\) does not depend on \(L\). The error \(e_i\) decays exponentially
in \(n_i\), at an \(L\)-dependent rate.

\emph{Status:} theorem, for all odd \(L\ge1\): the minimal block factor
produced is \(1\). \emph{Lean:}
\href{https://github.com/yang-mills-group/OpenAI-YM-20261028/blob/3456e5e75cf8d4cd558955a5c7638fba2d68f464/YM/Next/STATIONARY/FullPlanned.lean\#L17-L52}{\texttt{pouring}}
(proves
\href{https://github.com/yang-mills-group/OpenAI-YM-20261028/blob/3456e5e75cf8d4cd558955a5c7638fba2d68f464/YM/Next/STATIONARY/Level/Recurse.lean\#L115}{\texttt{Fluent}},
i.e.~\href{https://github.com/yang-mills-group/OpenAI-YM-20261028/blob/3456e5e75cf8d4cd558955a5c7638fba2d68f464/YM/Next/STATIONARY/Level/Transition.lean\#L20-L23}{\texttt{HorizCalc}}
for all parameters),
\href{https://github.com/yang-mills-group/OpenAI-YM-20261028/blob/3456e5e75cf8d4cd558955a5c7638fba2d68f464/YM/Next/STATIONARY/Spectral/AuxReferences.lean\#L138}{\texttt{PoleTest}},
\href{https://github.com/yang-mills-group/OpenAI-YM-20261028/blob/3456e5e75cf8d4cd558955a5c7638fba2d68f464/YM/Next/STATIONARY/Spectral/AuxReferences.lean\#L46}{\texttt{fineMag}},
\href{https://github.com/yang-mills-group/OpenAI-YM-20261028/blob/3456e5e75cf8d4cd558955a5c7638fba2d68f464/YM/Next/STATIONARY/Spectral/AuxReferences.lean\#L45}{\texttt{rankAux}},
\href{https://github.com/yang-mills-group/OpenAI-YM-20261028/blob/3456e5e75cf8d4cd558955a5c7638fba2d68f464/YM/Next/STATIONARY/Joint/Regulator.lean\#L338}{\texttt{grenade}},
\href{https://github.com/yang-mills-group/OpenAI-YM-20261028/blob/3456e5e75cf8d4cd558955a5c7638fba2d68f464/YM/Gauge/Colors.lean\#L184}{\texttt{Gauge.slopeConstant}}.

\textbf{Proposition III.7.4 (one renormalization step of the leading
part).} Fix a simple \(G\), kernel parameters \(K\) (penalty \(a>0\)), a
stage \(J\), and constants \(\varepsilon,w,\delta>0\) and
\(D,Z,C,Z'\ge0\) for the fine and coarse beds and clips. There are
\(A_*\ge0\) and \(s>0\) (output amplitude and width) such that, for
every input width \(\tau>0\) and amplitude \(A\ge0\), there are:

\begin{itemize}
\tightlist
\item
  \(d\ge0\) (the drift) and \(\gamma>0\) (a rate);
\item
  a minimal block factor \(u\);
\item
  functions \(\Theta_0(L)\) and \(\Lambda_0(L)\),
\end{itemize}

with the following property.

\emph{Hypotheses.}

\begin{itemize}
\tightlist
\item
  \emph{Steps.} Coarse sides \(m_i\to\infty\); block-spin steps
  \(\Lambda_{Lm_i}\to\Lambda_{m_i}\) with one admissible patch of stage
  \(J\) and block factor \(L\ge u\), and kernel \(K\); each coarse side
  is greater than \(1\), and each fine side is at least \(100\) times
  the patch extent.
\item
  \emph{Clips and beds.}

  \begin{itemize}
  \tightlist
  \item
    A fine clip \(\mathbf p\) with \(p_{\rm dir}\le Z\) and a coarse
    clip \(\mathbf q\) with \(q_{\rm dir}\le Z'\).
  \item
    A fine bed \(B\) with \(H,\rho\in\mathrm{strip}(w,D)\),
    \(\mathrm{Floor}(B,\mathbf p,n;\varepsilon,D)\) for every \(n\), and
    signed-permutation-invariant gauge form on the fine tori.
  \item
    A coarse bed \(T\) with \(H_T,\rho_T\in\mathrm{strip}(\delta,C)\),
    \(\mathrm{Floor}(T,\mathbf q,n;\varepsilon',C)\) for every \(n\) and
    some \(\varepsilon'\in\mathbb R\), and signed-permutation-invariant
    gauge form on the coarse tori.
  \item
    A tower with interpolation kernel \(A_0\), satisfying the
    quadratic-form interpolation, curl-transport and diagonal-transport
    identities of III.1 on each step. For every list \(\nu\) of at most
    two fine-coordinate differences, \(\nabla^\nu A_0\) has strip width
    \(w\) and entrywise weighted sum at most \(D/L^{1+|\nu|}\).
  \end{itemize}
\item
  \emph{Fine datum.} \(R\in\mathrm{Bounds}(\mathbf p;\tau,A,\Lambda)\)
  attached to \(B\), with \(\Lambda\ge1\).
\item
  \emph{Stationary branches.} A common \(\eta>0\) and, for each \(i\),
  an axial forest \(F_i\) and a branch \(g_i\) (coarse fields \(\to\)
  free fine variables) such that:

  \begin{itemize}
  \tightlist
  \item
    for every choice of block-kernel data, the second variation in the
    free variables at \((0,0)\) is at least \(\eta\sum_e|v_e|^2\);
  \item
    \(g_i(0)=0\), \(g_i\) and the fine action are real-analytic at
    \(0\), and \(g_i\) is stationary for every choice of block-kernel
    data;
  \item
    on every free edge, \(Dg_i(0)v\) is the free interpolant in the
    chosen axial gauge minus the derivative of the copied coarse
    configuration;
  \item
    the second jet of the tree-level coarse action is the quadratic form
    of \(T\).
  \end{itemize}
\item
  \emph{Chart.} A local Haar chart ratio \(\mathsf J\).
\end{itemize}

\emph{Conclusion.} There are a coarse path datum \(Y\) and
\(\alpha\in\mathbb R\) such that:

\begin{enumerate}
\def\labelenumi{\arabic{enumi}.}
\item
  the rows of \(Y\) are the straight-route rows of \(T\), and \(Y\) is
  attached to \(T\);
\item
  \(Y\in\mathrm{Bounds}(\mathbf q;\eta',A',\Lambda')\) whenever
  \(0<\eta'\le s\), \(A'\ge A_*\) and
  \(\Lambda'\ge\max(\Lambda_0(L),\Lambda/2)\);
\item
  \(|\alpha+\kappa\log L|+1\le d\);
\item
  \emph{classical match:} for every \(i\), every choice \(t\) of
  block-kernel data and every \(j\le4\), the \(j\)-th jets at \(0\) of
  \[x\mapsto S_t(x,g_i(x))\quad\text{and}\quad x\mapsto\mathrm{main}_{\mathbf q,Y}(e^x)\]
  coincide. Here
  \(S_t(x,z)=\mathrm{main}_{\mathbf p,R}(U_{x,z})+a\sum_e\big(1-d_{\rm rep}^{-1}\mathrm{Re}\,\mathrm{tr}(V_e\,\Pi_t(U_{x,z})_e^{-1})\big)\)
  is the fine leading part plus the block-spin penalty. \(U_{x,z}\) is
  the fine configuration lifted from \(V=e^x\) through the forest with
  free variables \(z\), and \(\Pi_t\) is the block predictor;
\item
  \emph{finite bound:} for every \(i\), every \(\Theta\ge\Theta_0(L)\)
  and every \(c'\le\gamma\), consider the one-loop function
  \[\mathrm{th}_{i,t}(x)=\mathrm{loop}_R(U_{x,g_i(x)})+\tfrac12\log\det\mathrm{Hess}_zS_t(x,g_i(x))-\log\textstyle\prod_e\mathsf J(g_i(x)_e).\]
  It is analytic at \(0\), the Hessian determinant is positive near
  \(0\), and for all coarse lattice momenta \(p,p'\) and real plane
  waves \(w_{p,\xi}(x)=\mathrm{Re}(e^{ip\cdot x}\xi)\) and
  \(w_{p',\zeta}\), with every direction-colour component of
  \(\xi,\zeta\) of modulus at most \(1\),
  \[\Big|D^2\mathrm{th}_{i,t}(0)[w_{p,\xi},w_{p',\zeta}]-D^2\big(\alpha\,\mathrm{quad}_{\mathbf q,Y}+\mathrm{loop}_Y\big)(0)[w_{p,\xi},w_{p',\zeta}]\Big|\le|\Lambda_{m_i}|\,\Theta\,e^{-c'm_i}.\]
  Here \(\mathrm{quad}_{\mathbf q,Y}\) is the signed-permutation average
  of the Wilson-plaquette and row-square terms in
  \(\mathrm{main}_{\mathbf q,Y}\), with the cubic and quartic atoms
  omitted. Its Hessian is the coarse quadratic form. The comparison
  function is also analytic at \(0\). The estimate holds for every
  \(t\).
\end{enumerate}

\emph{On the constants.} \(A_*\) and \(s\) are fixed before
\((\tau,A)\). The drift is \(d=d_{\rm cmp}+D_{\rm slope}+1\), where
\(d_{\rm cmp}\) is the comparison constant of §4.8 and \(D_{\rm slope}\)
is the constant of III.7.3. Both may depend on \((\tau,A)\) but are
independent of \(L\). The construction gives
\(\Lambda_0(L)=a_0L^4+b_0\log L+c_0\) and \(\Theta_0(L)=b_1L^4\), with
nonnegative constants depending on the fixed input parameters; both are
\(O(L^4)\).

\emph{Iteration.} Fix positive stage widths and lower bounds for the
free quadratic forms, stage amplitude bounds, kernel parameters \(K\)
and clip bounds. The constants are then chosen in the following order:

\begin{enumerate}
\def\labelenumi{\arabic{enumi}.}
\tightlist
\item
  maximal path widths \(\sigma_j\) and minimal path amplitudes \(A_j\);
\item
  for each positive sequence \(\tau_j\le\sigma_j\) and sequence
  \(A'_j\ge A_j\): drifts \(d_j\ge0\), rates \(\gamma_j>0\) and minimal
  block factors \(u_j\);
\item
  for every compatible choice of stage parameters, with direct clip
  coefficients below their prescribed bounds, drifts at least \(d_j\)
  and block-factor lower bounds at least \(u_j\), and then every
  admissible choice of block factors: volume thresholds and positive
  Hessian lower bounds \(\eta_j\);
\item
  for every admissible protection sequence \(\theta\): entry thresholds
  \(\Delta_j\) for the reference coupling.
\end{enumerate}

For each such choice, take post-step constants \(\Gamma_j\ge1\) and
\(0<c_j\le\gamma_j\) and the Hessian lower bounds \(\eta_j\) just
chosen. The conclusion holds for every history satisfying the
compatibility, size and coupling thresholds of III.2; every free-field
horizon satisfying the strip, Floor, symmetry, top-identity,
interpolation and three-form coercivity conditions of III.10; every
second horizon with the same beds and the same interpolation and
curl-transport kernels; and every local Haar chart ratio. For such data
there are kinetic data \((\alpha_{N,m},P_{N,m},g_{N,m})\) satisfying all
the following conditions:

\begin{itemize}
\tightlist
\item
  every \(P_{N,m}\) is attached to its bed and satisfies the level-\(m\)
  clauses with loop amplitude \(\mathrm{loopCap}_m\);
\item
  \(P_{N,N}\) is the Wilson top datum of III.7.2;
\item
  for every \(m<N\), every volume index and every axial forest: the
  second-variation lower bound \(\eta_{J_{m+1}}\), stationarity, the
  derivative condition and the classical match; also
  \(|\alpha_{N,m}+\kappa\log L_m|+1\le\mathrm{drift}_{m+1}\), and the
  finite bound with constants \((\Gamma_{J_{m+1}},c_{J_{m+1}})\) at loop
  amplitude \(\mathrm{loopCap}_{m+1}\) and also at amplitude \(1\).
\end{itemize}

\emph{Status:}

\begin{itemize}
\tightlist
\item
  The step and iteration are theorems; their output conditions are the
  definitions used in III.9.
\item
  \href{https://github.com/yang-mills-group/OpenAI-YM-20261028/blob/3456e5e75cf8d4cd558955a5c7638fba2d68f464/YM/Next/STATIONARY/Level/Transition.lean\#L92-L182}{\texttt{transition\_cap}}
  proves the step property at every stage using III.7.3 for the input
  constants.
  \href{https://github.com/yang-mills-group/OpenAI-YM-20261028/blob/3456e5e75cf8d4cd558955a5c7638fba2d68f464/YM/Next/STATIONARY/Level/Recurse.lean\#L122-L143}{\texttt{claiming}},
  \href{https://github.com/yang-mills-group/OpenAI-YM-20261028/blob/3456e5e75cf8d4cd558955a5c7638fba2d68f464/YM/Next/STATIONARY/Level/Recurse.lean\#L366-L391}{\texttt{under\_contract}}
  and
  \href{https://github.com/yang-mills-group/OpenAI-YM-20261028/blob/3456e5e75cf8d4cd558955a5c7638fba2d68f464/YM/Next/STATIONARY/FullPlanned.lean\#L59-L66}{\texttt{geared}}
  iterate it.
\item
  The iteration supplies the one-loop slope internally; its remaining
  hypotheses are the free-field, history and Haar-chart conditions just
  listed.
\item
  The coarse-datum construction and induction are also expressed by
  \href{https://github.com/yang-mills-group/OpenAI-YM-20261028/blob/3456e5e75cf8d4cd558955a5c7638fba2d68f464/YM/Next/STATIONARY/Loops/Curvature.lean\#L394-L435}{\texttt{color\_pipe}},
  \href{https://github.com/yang-mills-group/OpenAI-YM-20261028/blob/3456e5e75cf8d4cd558955a5c7638fba2d68f464/YM/Next/STATIONARY/Level/Recurse.lean\#L254-L285}{\texttt{energized}}
  and
  \href{https://github.com/yang-mills-group/OpenAI-YM-20261028/blob/3456e5e75cf8d4cd558955a5c7638fba2d68f464/YM/Next/STATIONARY/Level/Recurse.lean\#L307-L335}{\texttt{plowed}},
  together with the preceding step theorem.
\end{itemize}

\emph{Lean:}
\href{https://github.com/yang-mills-group/OpenAI-YM-20261028/blob/3456e5e75cf8d4cd558955a5c7638fba2d68f464/YM/Next/STATIONARY/Level/Transition.lean\#L86-L90}{\texttt{Transition}},
\href{https://github.com/yang-mills-group/OpenAI-YM-20261028/blob/3456e5e75cf8d4cd558955a5c7638fba2d68f464/YM/Next/STATIONARY/Level/Transition.lean\#L26-L38}{\texttt{Scud}},
\href{https://github.com/yang-mills-group/OpenAI-YM-20261028/blob/3456e5e75cf8d4cd558955a5c7638fba2d68f464/YM/Next/STATIONARY/Level/Transition.lean\#L25}{\texttt{Hull}},
\href{https://github.com/yang-mills-group/OpenAI-YM-20261028/blob/3456e5e75cf8d4cd558955a5c7638fba2d68f464/YM/Next/STATIONARY/Level/Recurse.lean\#L170-L177}{\texttt{Bubble}},
\href{https://github.com/yang-mills-group/OpenAI-YM-20261028/blob/3456e5e75cf8d4cd558955a5c7638fba2d68f464/YM/Next/STATIONARY/Level/Recurse.lean\#L301-L305}{\texttt{FullPaths}},
\href{https://github.com/yang-mills-group/OpenAI-YM-20261028/blob/3456e5e75cf8d4cd558955a5c7638fba2d68f464/YM/PCalc/Loop.lean\#L93-L115}{\texttt{Kin.Planned}},
\href{https://github.com/yang-mills-group/OpenAI-YM-20261028/blob/3456e5e75cf8d4cd558955a5c7638fba2d68f464/YM/PCalc/Chart.lean\#L108-L109}{\texttt{ClassicalMatch}},
\href{https://github.com/yang-mills-group/OpenAI-YM-20261028/blob/3456e5e75cf8d4cd558955a5c7638fba2d68f464/YM/PCalc/Chart.lean\#L81-L84}{\texttt{tested}},
\href{https://github.com/yang-mills-group/OpenAI-YM-20261028/blob/3456e5e75cf8d4cd558955a5c7638fba2d68f464/YM/PCalc/Loop.lean\#L46-L49}{\texttt{thermal}},
\href{https://github.com/yang-mills-group/OpenAI-YM-20261028/blob/3456e5e75cf8d4cd558955a5c7638fba2d68f464/YM/PCalc/Loop.lean\#L62-L71}{\texttt{FiniteBound}},
\href{https://github.com/yang-mills-group/OpenAI-YM-20261028/blob/3456e5e75cf8d4cd558955a5c7638fba2d68f464/YM/PCalc/Loop.lean\#L51}{\texttt{draft}}.

\textbf{Proposition III.7.5 (uniform inverse of the completed free
operator).} Let \(j\) be a stage, \(K\) kernel parameters,
\(\sigma,\varepsilon>0\) and \(\Omega\ge0\). There is \(\eta>0\), and
for every block factor \(L\) a constant \(A_L\ge0\), with the following
property.

For every admissible patch of stage \(j\) and block factor \(L\), and
every bed \(B\) with \(H\in\mathrm{strip}(\sigma/L,\Omega)\) and
\(\mathrm{Floor}(B,\mathbf c,n;\varepsilon,D)\) for all \(n\):

\begin{itemize}
\tightlist
\item
  the inverse \(\mathcal G_B\) of the block-Bloch symbol of the
  completed free operator lies in \(\mathrm{strip}(\eta,A_L)\) on the
  coarse lattice;
\item
  for every step with this patch (any coarse side \(c\), fine side
  \(cL\), any choice of block-kernel data),
  \(\mathcal C_B\,\mathcal G_B=1\).
\end{itemize}

In the rooted coordinates, the joint variables are a fine link field
\(v\) and a coarse site field \(h\). Before gauge completion, the
quadratic form is \[\langle dv,H\,dv\rangle+a\|{-}Qv+d_ch\|^2,\] where
\(Q\) is the linear block average and \(d_c\) the coarse gradient. The
completed operator adds \(F_{1/L}^{\mathsf T}F_{1/L}\), where
\(F_\mu=d_{\!1}^*P_f+\mathfrak m_\mu X\), \(P_f\) projects onto fine
links, and \(X\) is minus the derivative of the logarithm of the fine
rooted frame. Here \(d_{\!1}\) is \(d_U\) at \(U=1\). All matrices use
physical counting bases. The decay rate \(\eta\) is uniform in \(L\);
only the amplitude depends on \(L\).

\emph{Status:} theorem; used in III.4. \emph{Lean:}
\href{https://github.com/yang-mills-group/OpenAI-YM-20261028/blob/3456e5e75cf8d4cd558955a5c7638fba2d68f464/YM/WilsonBulk/Spectral/Backing.lean\#L54-L66}{\texttt{WilsonBulk.uniform\_backing}},
\href{https://github.com/yang-mills-group/OpenAI-YM-20261028/blob/3456e5e75cf8d4cd558955a5c7638fba2d68f464/YM/WilsonBulk/Spectral/Backing.lean\#L18-L20}{\texttt{completedM}},
\href{https://github.com/yang-mills-group/OpenAI-YM-20261028/blob/3456e5e75cf8d4cd558955a5c7638fba2d68f464/YM/WilsonBulk/Spectral/Backing.lean\#L22-L23}{\texttt{jointInverse}}.

\textbf{Proposition III.7.6 (Faddeev--Popov identity for the one-loop
term).} Fix an admissible step, a clip \(\mathbf p\) and a datum
\(R\in\mathrm{Bounds}(\mathbf p;\sigma,A,\Lambda)\) for some
\(\sigma>0\) and \(A,\Lambda\in\mathbb R\). Let \(g(0)=0\) be a
real-analytic stationary branch for an axial forest, stationary for
every choice of block-kernel data, with second variation
\(\ge\eta\sum_e|v_e|^2\) (\(\eta>0\)). Let \(\mu>0\), let \(t\) be a
choice of block-kernel data, and let \(\mathsf J\) be a local Haar chart
ratio. Then for coarse fields \(x\) near \(0\):

\begin{itemize}
\tightlist
\item
  the completed Hessian \(\mathcal P_\mu(x)\) has positive determinant;
\item
  \(\mathrm{th}(x)=\mathrm{loop}_R(U_{x,g(x)})+\tfrac12\log\det\mathcal P_\mu(x)-\log\det\mathcal G_\mu(x)\).
\end{itemize}

\(\mathcal P_\mu(x)\) acts on all fine-link and coarse-frame directions.
It is the Hessian in left exponential coordinates of the joint action
\[\mathscr S_t(V;U,h)=\mathrm{main}_{\mathbf p,R}(U)+a\sum_e\Big(1-d_{\rm rep}^{-1}\mathrm{Re}\,\mathrm{tr}\big(V_e[\mathrm{gauge}_h\Pi_t(U)]_e^{-1}\big)\Big),\]
at \(V=e^x\) and the lifted joint stationary configuration, plus
\(F_\mu^{\mathsf T}F_\mu\). Here \(F_\mu=d_U^*P_f+\mathfrak m_\mu(U)X\),
with the rooted-frame derivative \(X\) as in III.7.5. The joint
stationary configuration combines the lifted fine field \(U_{x,g(x)}\)
with its coarse frame. \(\mathcal G_\mu=d_U^*d_U+\mathfrak m_\mu(U)\) is
the ghost operator at the lifted configuration. The underlying algebraic
identity is \[T^2\det\mathcal P_\mu=\det M\,(\det\mathcal G_\mu)^2,\]
with \(M\) the axial Hessian and \(T\) the Haar Jacobian of the chart.

\emph{Status:} theorem; used in the renormalization and remainder
comparisons of III.4, III.2 and VII.2. \emph{Lean:}
\href{https://github.com/yang-mills-group/OpenAI-YM-20261028/blob/3456e5e75cf8d4cd558955a5c7638fba2d68f464/YM/Next/STATIONARY/Joint/HaarJac.lean\#L177-L202}{\texttt{tribute}},
\href{https://github.com/yang-mills-group/OpenAI-YM-20261028/blob/3456e5e75cf8d4cd558955a5c7638fba2d68f464/YM/Next/STATIONARY/Joint/HaarJac.lean\#L166-L167}{\texttt{bounty}},
\href{https://github.com/yang-mills-group/OpenAI-YM-20261028/blob/3456e5e75cf8d4cd558955a5c7638fba2d68f464/YM/Next/STATIONARY/Joint/Completion.lean\#L59-L60}{\texttt{medals}}.

\subsubsection{2. What it uses}\label{what-it-uses-9}

\begin{itemize}
\item
  \textbf{I.1 (gauge groups, Wilson action).} \(E_W\), the exponential
  chart, \(\mathrm{Ad}\), the Casimir \(C_G\), and \(\kappa\) with
  \(\kappa>0\)
  (\href{https://github.com/yang-mills-group/OpenAI-YM-20261028/blob/3456e5e75cf8d4cd558955a5c7638fba2d68f464/YM/Gauge/Casimir.lean\#L174-L176}{\texttt{slopeConstant\_pos}},
  \href{https://github.com/yang-mills-group/OpenAI-YM-20261028/blob/3456e5e75cf8d4cd558955a5c7638fba2d68f464/YM/Gauge/Colors.lean\#L180-L182}{\texttt{casimir}}).
  Simplicity enters as \(\mathrm{tr}(\mathrm{ad}_b\mathrm{ad}_b)=-C_G\)
  for every unit vector \(b\), and through Schur's lemma.
\item
  \textbf{III.1 (block-spin step).} The specifications of III.7.1. Also:
  steps, patches (block factor \(L\), stage \(J\)), axial forests (tree
  gauge), the fine action plus block penalty, and the stationarity
  notions (stationarity, second-variation lower bound, derivative
  condition, equality of jets, classical match). Further: the one-loop
  function, the finite bound, the kinetic data and the iteration
  conditions, towers, horizons, and local Haar chart ratios.
\item
  \textbf{III.6 (stationary branch).} For every admissible step there is
  a Hessian lower bound \(\eta_j>0\) (chosen after \(L\)) such that
  every admissible fine datum has a stationary branch with:

  \begin{itemize}
  \tightlist
  \item
    second variation \(\ge\eta_j\);
  \item
    derivative given by the free interpolation in the chosen axial
    gauge, after subtracting the copied coarse field;
  \item
    quadratic part equal to the coarse bed.
  \end{itemize}

  This is
  \href{https://github.com/yang-mills-group/OpenAI-YM-20261028/blob/3456e5e75cf8d4cd558955a5c7638fba2d68f464/YM/Integrated.lean\#L229-L232}{\texttt{Integrated.stationary\_scene}},
  packaged here by
  \href{https://github.com/yang-mills-group/OpenAI-YM-20261028/blob/3456e5e75cf8d4cd558955a5c7638fba2d68f464/YM/Next/STATIONARY/Level/SliceGate.lean\#L58-L61}{\texttt{scene\_gate}}.
  Also from III.6: positivity of the Hessian near \(0\)
  (\href{https://github.com/yang-mills-group/OpenAI-YM-20261028/blob/3456e5e75cf8d4cd558955a5c7638fba2d68f464/YM/Work/F05/Graphs.lean\#L75-L79}{\texttt{hessian\_pos}})
  and independence of the one-loop function from the forest
  (\href{https://github.com/yang-mills-group/OpenAI-YM-20261028/blob/3456e5e75cf8d4cd558955a5c7638fba2d68f464/YM/Work/F05.lean\#L42-L57}{\texttt{same\_thermal}}).
\item
  \textbf{III.10 (conditions on beds and steps).}

  \begin{itemize}
  \tightlist
  \item
    \href{https://github.com/yang-mills-group/OpenAI-YM-20261028/blob/3456e5e75cf8d4cd558955a5c7638fba2d68f464/YM/Bid/Station.lean\#L54-L55}{\texttt{Bid.Sworn}}:
    signed-permutation invariance of the bed's gauge form.
  \item
    \href{https://github.com/yang-mills-group/OpenAI-YM-20261028/blob/3456e5e75cf8d4cd558955a5c7638fba2d68f464/YM/Bid/Station.lean\#L83-L84}{\texttt{Bid.HeatG}}:
    quadratic jet of the tree-level action.
  \item
    \href{https://github.com/yang-mills-group/OpenAI-YM-20261028/blob/3456e5e75cf8d4cd558955a5c7638fba2d68f464/YM/Bid/Station.lean\#L64}{\texttt{Bid.Larger}}:
    size of the torus.
  \item
    \href{https://github.com/yang-mills-group/OpenAI-YM-20261028/blob/3456e5e75cf8d4cd558955a5c7638fba2d68f464/YM/Bid/Craft.lean\#L63-L73}{\texttt{Craft.Rack}}
    and
    \href{https://github.com/yang-mills-group/OpenAI-YM-20261028/blob/3456e5e75cf8d4cd558955a5c7638fba2d68f464/YM/Bid/Craft.lean\#L104-L105}{\texttt{Bid.Retained}}:
    uniform strip and Floor bounds, symmetry and three-form coercivity
    for the free-field horizon; and equality of the beds and the two
    interpolation kernels for the second horizon.
  \item
    \href{https://github.com/yang-mills-group/OpenAI-YM-20261028/blob/3456e5e75cf8d4cd558955a5c7638fba2d68f464/YM/Bid/Scope.lean\#L35-L43}{\texttt{Bid.one\_for\_all}}:
    the classical match does not depend on the forest.
  \end{itemize}
\item
  \textbf{III.2 (order of constants).} The admissible stage parameters,
  block factors, protection exponents, entry and volume thresholds, and
  the history compatibility conditions
  \href{https://github.com/yang-mills-group/OpenAI-YM-20261028/blob/3456e5e75cf8d4cd558955a5c7638fba2d68f464/YM/Price/Intent.lean\#L145-L162}{\texttt{Supply.Paid}},
  free-field bounds and constants chosen after the block factors. Also
  the reference coupling
  \href{https://github.com/yang-mills-group/OpenAI-YM-20261028/blob/3456e5e75cf8d4cd558955a5c7638fba2d68f464/YM/RG/History.lean\#L79}{\texttt{β̂\_m}}
  \(=\bar\beta+\kappa s_m\) and the loop cap
  \href{https://github.com/yang-mills-group/OpenAI-YM-20261028/blob/3456e5e75cf8d4cd558955a5c7638fba2d68f464/YM/RG/History.lean\#L86}{\texttt{loopCap\_m}}
  \(=2\exp\big(\theta_{J_m}(\log\log\hat\beta_m)^2\big)\).
\item
  \textbf{II.1.} The local Haar chart ratio of the exponential chart
  (\href{https://github.com/yang-mills-group/OpenAI-YM-20261028/blob/3456e5e75cf8d4cd558955a5c7638fba2d68f464/YM/Next/FOUNDATIONS/Primary.lean\#L84}{\texttt{chart\_ok}}),
  used in the iteration.
\item
  \textbf{III.16, III.3.} Laurent/Bloch calculus on strips (products,
  inverses, Cauchy estimates, Riemann sums versus integrals); the
  block-Bloch (``slab'') representation of fine translation-invariant
  operators on the coarse lattice.
\item
  \textbf{Standard mathematics.}

  \begin{itemize}
  \tightlist
  \item
    Taylor expansion of the plaquette action; polarization (a jet is
    determined by its restrictions to lines).
  \item
    The Koszul complex of the regular sequence \(q_j(p)=e^{ip_j}-1\)
    with Hadamard division in a strip, i.e.~the lattice Poincaré lemma
    in momentum space.
  \item
    Schur's lemma for invariant forms on a simple Lie algebra.
  \item
    The Faddeev--Popov procedure.
  \item
    The background-field method ('t Hooft and Abbott; Peskin--Schroeder,
    §16.6).
  \item
    The one-loop \(\beta\)-function (Gross--Wilczek and Politzer) and
    its orbital/spin split \(-\tfrac{11}3=\tfrac13-4\).
  \item
    Balaban's propagator bounds and background-field renormalization
    group; see also Dimock, \emph{The renormalization group according to
    Balaban}.
  \item
    Rotation-invariant integrals over shells in \(\mathbb R^4\).
  \end{itemize}
\item
  \textbf{Imported constants and assumptions.} No one-loop coefficient
  is assumed. \(\kappa\) is defined in I.1, and III.7.3 proves that it
  is the one-loop slope. The hypotheses carried are those displayed in
  the propositions; the free-field and history assumptions are supplied
  by III.4.
\end{itemize}

\subsubsection{3. Strategy}\label{strategy-9}

\textbf{Classical part.} At every level the quadratic part of the
leading action is fixed by the free-field bed of that level. The cubic
and quartic parts are curvature atoms, determined by Taylor agreement
through order four with a given gauge-invariant function \(E\). At the
top, \(E\) is the Wilson energy. At each step, \(E\) is the tree-level
coarse action \(x\mapsto S_t(x,g(x))\). The atom-emission lemma (§4.3,
Lemma 4.4) serves both cases. Its mechanism has four parts:

\begin{enumerate}
\def\labelenumi{\arabic{enumi}.}
\tightlist
\item
  The Taylor coefficients of a translation sum of local analytic
  functionals have Fourier kernels on \(\mathbb Z^4\) that are
  strip-analytic. Finite-volume coefficients agree with \(|\Lambda_n|\)
  times these kernels up to \(o(|\Lambda_n|)\) along any sequence of
  sides; this is how volume-uniformity enters.
\item
  Gauge invariance makes the kernels transversal in each leg (linearized
  Ward identity).
\item
  Transversal kernels with at least three legs factor through the
  curvature. This is the lattice Poincaré lemma, proved by contracting
  the Koszul complex of the differences \(e^{ip_j}-1\), one axis at a
  time, with Hadamard division. A one-leg solve reduces the width from
  \(\sigma\) to \(\sigma/4\) once, regardless of the four coordinate
  divisions within it. Successive leg solves cause further width losses,
  depending only on the fixed number of legs.
\item
  The real, colour-invariant Laurent coefficients of the factored
  kernels are the coefficients of atoms, and exponential decay is the
  strip bound.
\end{enumerate}

\textbf{One-loop part.}

\begin{enumerate}
\def\labelenumi{\arabic{enumi}.}
\tightlist
\item
  In the axial chart the one-loop term is
  \(\mathrm{loop}+\tfrac12\log\det-\log(\text{Haar Jacobian})\). An
  exact Faddeev--Popov identity turns it into \(\tfrac12\log\det\) of a
  covariantly completed Hessian minus \(\log\det\) of a ghost operator.
\item
  Its second-order Taylor coefficient is a two-leg kernel. Its
  \(F^2\)-coefficient \(\alpha\) is compared, with an \(L\)-independent
  error, to that of the regulated Feynman-gauge background-field
  determinant on the fine lattice.
\item
  The latter is computed. The lattice integrand is the continuum bubble
  integrand \(h(k)\) of the vector loop (weight \(\tfrac12\)) and the
  ghost loop (weight \(-1\)), up to errors
  \(O(|k|^{-3}+\mu^2|k|^{-6})\); its angular average is
  \(-\tfrac{11}3C_G|k|^{-4}\). Integration over the shell
  \(1/L<|k|\le1\) gives \(-\kappa\log L\), and the rest is bounded by
  absolute values.
\item
  After subtracting \(\alpha\) times the bed's quadratic kernel, the
  remaining kernel vanishes to fourth order at zero momentum. It is
  therefore represented by loop atoms satisfying the two moment
  conditions, i.e.~by irrelevant terms. Inherited fine loop atoms
  contribute with a factor \(L^{-2}\).
\end{enumerate}

\textbf{Constants.} The output width and cubic/quartic amplitude of a
step are fixed before the input amplitude. The input enters the output
only with factors \(L^{-1}\) or \(L^{-2}\), which a large \(L\) makes
small. The strip widths and inverse bounds for the completed operators
are uniform in \(L\) in block-scaled norms. The newly generated
loop-atom amplitude is \(O(L^4)\). It is absorbed by
\(\mathrm{loopCap}_m\), which grows with \(\hat\beta_m\).

\textbf{Relation to Balaban.} The construction uses background-field
renormalization with block-spin constraints. Here:

\begin{itemize}
\tightlist
\item
  the effective action is carried by explicit, exponentially weighted
  path data;
\item
  the \(F^2\)-coefficient is obtained by comparison with a regulated
  covariant determinant plus a Fourier computation, with \(O(1)\) error,
  for an arbitrary simple compact gauge group.
\end{itemize}

\subsubsection{4. Steps}\label{steps-9}

\paragraph{4.1 Fourier kernels of local
functionals}\label{fourier-kernels-of-local-functionals}

\textbf{Lemma 4.1 (translation sums of local functionals; limits of
kernels).} (a) Let \(F_k\) be functions of configurations on
\(\mathbb Z^4\). Suppose \(F_k\) depends only on links in a box of size
\(d_k\), has a holomorphic extension with bound \(w_k\) on a polydisc of
radius \(r\) around the identity, uniformly after translation and
periodic restriction, and \(\sum_kw_k\exp(4(t+1)^2\sigma d_k)\le A\),
with \(r,\sigma>0\) and \(w_k\ge0\). Then for every side \(n\) the
translation sum on \(\Lambda_n\) is analytic at \(1\). Its
\((t{+}1)\)-th Taylor coefficient is recorded \emph{exactly} by a kernel
that is strip-analytic of width \(\sigma\) and size \(A\,c(t,r)\),
independent of \(n\). For the Wilson energy, kernels of every order
exist with width \(1\). (b) Suppose a strip-analytic kernel records
\(F_i\) on \(\Lambda_{n_i}\) with error \(|\Lambda_{n_i}|\epsilon_i\),
where \(n_i\to\infty\) and \(\epsilon_i\to0\). If the \(F_i\) are
analytic, invariant near \(0\) under global colour rotations and lattice
gauge transformations, and their Taylor coefficients of positive orders
\(1,\ldots,t\) vanish, then the kernel is transversal in each leg, real
and colour invariant. \emph{Sketch:} (a) Cauchy estimates on the
polydisc bound each local Taylor coefficient; its Fourier support has
size at most \(4(t+1)^2d_k\), which gives the stated exponential factor.
(b) Each property holds exactly for the finite-volume coefficients;
divide by \(|\Lambda_{n_i}|\) and pass to the limit along lattice
momenta approaching any prescribed real momenta. Strip analyticity gives
the required continuity. \emph{Lean:}
\href{https://github.com/yang-mills-group/OpenAI-YM-20261028/blob/3456e5e75cf8d4cd558955a5c7638fba2d68f464/YM/Next/STATIONARY/Shape/LocalOracle.lean\#L312-L348}{\texttt{gathered}},
\href{https://github.com/yang-mills-group/OpenAI-YM-20261028/blob/3456e5e75cf8d4cd558955a5c7638fba2d68f464/YM/Next/STATIONARY/Shape/ClassicalMatch.lean\#L312-L326}{\texttt{Wilson\_stream}},
\href{https://github.com/yang-mills-group/OpenAI-YM-20261028/blob/3456e5e75cf8d4cd558955a5c7638fba2d68f464/YM/Next/STATIONARY/Lift/Ward.lean\#L647-L697}{\texttt{withdrew}},
\href{https://github.com/yang-mills-group/OpenAI-YM-20261028/blob/3456e5e75cf8d4cd558955a5c7638fba2d68f464/YM/Next/STATIONARY/Shape/Vertex.lean\#L371-L379}{\texttt{wedlock}}.
\emph{(proof read at top level; Ward identity input statement only)}

\paragraph{4.2 The lattice Poincaré lemma in momentum
space}\label{the-lattice-poincaruxe9-lemma-in-momentum-space}

\textbf{Lemma 4.2 (one-leg solve).} Let the momentum space have \(d\)
scalar coordinates, including four selected coordinates, and let
\(\delta\) contract differential forms with \(q_j(p)=e^{ip_j}-1\) in
those four coordinates. Let \(\sigma>0\), \(A\ge0\), and let \(f\) be a
form of degree \(r\ge0\) with the following properties:

\begin{itemize}
\tightlist
\item
  each component is strip-analytic of width \(\sigma\) and size \(A\);
\item
  \(\delta f=0\);
\item
  \(f\) vanishes to order \(s\) along rays;
\item
  on real momenta, \(f\) takes values in subspaces unchanged when any of
  the four selected momentum coordinates varies;
\item
  if \(r=0\), \(f\) vanishes where the projected momentum is \(0\).
\end{itemize}

Then \(f=\delta B\), where \(B\) has degree \(r+1\), is strip-analytic
of width \(\sigma/4\) and size \(4A\,\pi_d(\sigma)\) with
\(\pi_d(\sigma)=(4/\sigma)\big(2(1+\sigma/4)/(\sigma/4)\big)^d\),
vanishes to order \(\max(s-1,0)\), and takes values in the same
subspaces on real momenta. \emph{Sketch:} Comb one axis at a time.
\(f(p)-f(p)|_{p_j=0}=q_j(p)D_j(p)\) by Hadamard division
(\href{https://github.com/yang-mills-group/OpenAI-YM-20261028/blob/3456e5e75cf8d4cd558955a5c7638fba2d68f464/YM/Next/STATIONARY/Lift/Koszul.lean\#L88-L128}{\texttt{div\_var}}).
\(D_j\) has width \(\sigma/4\) and size \(A\pi_d(\sigma)\) by Cauchy
estimates, while the restriction \(f|_{p_j=0}\) keeps width \(\sigma\).
After all four axes, no positive-degree component remains. In degree
zero the remainder is \(f\) with all selected momenta set to zero, which
vanishes by hypothesis. The factor \(4\) counts the axes; the width loss
is a single one. For positive vanishing order, division by
\(q_j\sim ip_j\) lowers it by one; order zero remains order zero.
\emph{Lean:}
\href{https://github.com/yang-mills-group/OpenAI-YM-20261028/blob/3456e5e75cf8d4cd558955a5c7638fba2d68f464/YM/Next/STATIONARY/Lift/Koszul.lean\#L571-L587}{\texttt{wall\_lift}}.

\textbf{Lemma 4.3 (transversal kernels are curvature kernels; atoms).}
For every \(t\ge2\) and \(\sigma>0\), there are \(\eta\in(0,\sigma]\)
and \(K\ge0\) such that every \((t{+}1)\)-leg kernel \(f\),
strip-analytic of width \(\sigma\) and size \(A\ge0\) and transversal in
each leg, admits the factorization
\[f(p)=\sum_b\prod_i\hat s(p_i)_{b_ia_i}\,g_b(p),\] for real momenta,
where \(\hat s\) is the symbol of the lattice curl and \(g\) is a
plaquette kernel, strip-analytic of width \(\eta\) and size \(AK\), with
\(\eta,K\) depending on \(t,\sigma\). If \(f\) is also real and colour
invariant, the Laurent coefficients of \(g\) can be taken real, colour
invariant and exponentially decaying. Placed on straight routes, they
define atoms. If \(f\) records an analytic functional invariant under
signed permutations, the signed-permutation average of their translation
sum reproduces that Taylor coefficient up to the recording error.
\emph{Sketch:} Solve the dependent-momentum leg first, then the
independent legs by Lemma 4.2, preserving the subspace constraints from
the earlier solves. Symmetrize by conjugate reflection to obtain real
Laurent coefficients and project the colour tensors onto the invariant
subspace. \emph{Lean:}
\href{https://github.com/yang-mills-group/OpenAI-YM-20261028/blob/3456e5e75cf8d4cd558955a5c7638fba2d68f464/YM/Next/STATIONARY/Lift/ManyLegs.lean\#L533-L549}{\texttt{allleg}},
\href{https://github.com/yang-mills-group/OpenAI-YM-20261028/blob/3456e5e75cf8d4cd558955a5c7638fba2d68f464/YM/Next/STATIONARY/Lift/Pathways.lean\#L193-L201}{\texttt{strokes}},
\href{https://github.com/yang-mills-group/OpenAI-YM-20261028/blob/3456e5e75cf8d4cd558955a5c7638fba2d68f464/YM/Next/STATIONARY/Lift/Ward.lean\#L192-L216}{\texttt{colors\_align}}.
\emph{(proof read at top level)}

\paragraph{4.3 Classical emission and the Wilson
jet}\label{classical-emission-and-the-wilson-jet}

\textbf{Lemma 4.4 (emission of cubic and quartic atoms).} Let
\(\sigma,\sigma_0>0\) and \(A,Z\ge0\). There are \(\tau\in(0,\sigma]\)
and \(C\ge0\) with the following property.

\emph{Hypotheses:}

\begin{itemize}
\tightlist
\item
  sides \(n_i\to\infty\);
\item
  a bed \(B\) with \(\rho\in\mathrm{strip}(\sigma_0,A)\) and Floor on
  the one-site torus, and a clip with \(c_{\rm dir}\le Z\);
\item
  real-analytic functions \(E_i\) of link fields on \(\Lambda_{n_i}\),
  invariant near \(0\) under lattice gauge transformations, global
  colour rotations and signed permutations, such that \(E_i\) minus the
  bare row action of \(B\) vanishes to order \(3\) along rays; the bare
  row action has the straight-route rows of \(\rho\) and no atoms;
\item
  \(M\ge1\) such that the Taylor coefficients of orders \(3\) and \(4\)
  of \(E_i\) are recorded exactly by kernels strip-analytic of width
  \(\sigma\) and size \(M\).
\end{itemize}

\emph{Conclusion:} there are invariant cubic atoms \(X\) and quartic
atoms \(V\) with sums of exponential weights at most \(MC\) at width
\(\tau\) such that, for every choice \(f\) of loop atoms and every
\(i\),
\(x\mapsto E_i(x)-\mathrm{main}_{\mathbf c,(\rho\text{-rows},f,X,V)}(e^x)\)
vanishes to order \(5\) along rays. \emph{Sketch:} At order \(3\),
subtract the full bare row action, whose quadratic jet agrees with that
of \(E\). The residual kernel is transversal (Lemma 4.1(b)) and is
matched by cubic atoms (Lemma 4.3). At order \(4\), subtract the same
bare row action and the cubic atoms; the remainder is transversal and is
matched by quartic atoms. \emph{Lean:}
\href{https://github.com/yang-mills-group/OpenAI-YM-20261028/blob/3456e5e75cf8d4cd558955a5c7638fba2d68f464/YM/Next/STATIONARY/Shape/ClassicalMatch.lean\#L561-L595}{\texttt{assembly}}.

\textbf{Lemma 4.5 (proof of III.7.2).}

\begin{enumerate}
\def\labelenumi{\arabic{enumi}.}
\tightlist
\item
  Take \(E_i=E_W\) on \(\Lambda_{n_i}\). The kernels of orders \(3\) and
  \(4\) come from Lemma 4.1(a).
\item
  Under \(D=1\), Floor and invariance, \(H\) agrees with the identity on
  exact 2-forms. Hence \(E_W\) minus the quadratic term of the bare
  datum (rows of \(\rho\), no atoms) vanishes to order \(3\)
  (\href{https://github.com/yang-mills-group/OpenAI-YM-20261028/blob/3456e5e75cf8d4cd558955a5c7638fba2d68f464/YM/Next/STATIONARY/Shape/ClassicalMatch.lean\#L441-L471}{\texttt{top\_Wilson}}).
\item
  Lemma 4.4 gives \(X,V\), and \(Q\) is formed from the rows of
  \(\rho\), zero loop atoms, \(X\) and \(V\).
\item
  Order-\(5\) vanishing along every ray of the difference of two
  analytic functions gives equality of all derivatives of order \(\le4\)
  (\href{https://github.com/yang-mills-group/OpenAI-YM-20261028/blob/3456e5e75cf8d4cd558955a5c7638fba2d68f464/YM/Next/STATIONARY/Shape/Scaled.lean\#L254-L260}{\texttt{matched}}).
\item
  The weight clauses come from the row prices \(6Ae^{4\sigma_0}\), the
  zero-momentum normalization (Floor on the one-site torus,
  \href{https://github.com/yang-mills-group/OpenAI-YM-20261028/blob/3456e5e75cf8d4cd558955a5c7638fba2d68f464/YM/Next/STATIONARY/Level/SliceGate.lean\#L214-L258}{\texttt{symbol\_match}}),
  and the common atom bound \((1+C_3+C_4)C\), where \(C_3,C_4\) bound
  the Wilson Taylor kernels and \(C\) is the emission constant.
\end{enumerate}

\emph{Lean:}
\href{https://github.com/yang-mills-group/OpenAI-YM-20261028/blob/3456e5e75cf8d4cd558955a5c7638fba2d68f464/YM/Next/STATIONARY/Shape/Scaled.lean\#L270-L305}{\texttt{launch}},
\href{https://github.com/yang-mills-group/OpenAI-YM-20261028/blob/3456e5e75cf8d4cd558955a5c7638fba2d68f464/YM/Next/STATIONARY/Level/Recurse.lean\#L101-L105}{\texttt{beckon}},
\href{https://github.com/yang-mills-group/OpenAI-YM-20261028/blob/3456e5e75cf8d4cd558955a5c7638fba2d68f464/YM/Next/STATIONARY/Level/Recurse.lean\#L220-L225}{\texttt{budding}}.

\textbf{Lemma 4.6 (every admissible leading part is Wilson up to
irrelevant terms).} Fix \(\sigma,\tau>0\), \(\varepsilon\in\mathbb R\)
and \(C,Z,A\ge0\). Let \(n_i\to\infty\), \(c_{\rm dir}\le Z\), and let
\(R\in\mathrm{Bounds}(\mathbf c;\tau,A,\Lambda)\) be attached to a bed
with \(\mathrm{Floor}(B,\mathbf c,n;\varepsilon,C)\) for every \(n\),
invariant gauge form on each \(\Lambda_{n_i}\) and
\(H\in\mathrm{strip}(\sigma,C)\). Then the Taylor kernels of orders
\(2\), \(3\) and \(4\) of \(\mathrm{main}_{\mathbf c,R}-E_W\) vanish at
zero momentum to orders \(3\), \(2\) and \(1\) respectively. Their strip
bounds depend only on \((\sigma,\varepsilon,C,Z,\tau,A)\). In words: all
admissible leading parts share the marginal (dimension-four) part of
\(E_W\). \emph{Sketch:} \(\hat H(0)=1\) makes the quadratic difference
vanish to order \(3\). Ward identities then force the cubic and quartic
differences to vanish to orders \(2\) and \(1\). Apply the same Ward
argument to the invariant cubic and quartic atom sums, whose quadratic
jets vanish. These statements concern their full nonlinear Taylor
coefficients, including transport and curvature expansions. \emph{Lean:}
\href{https://github.com/yang-mills-group/OpenAI-YM-20261028/blob/3456e5e75cf8d4cd558955a5c7638fba2d68f464/YM/Next/STATIONARY/Shape/Scaled.lean\#L363-L403}{\texttt{greyson}},
\href{https://github.com/yang-mills-group/OpenAI-YM-20261028/blob/3456e5e75cf8d4cd558955a5c7638fba2d68f464/YM/Next/STATIONARY/Shape/Scaled.lean\#L349-L358}{\texttt{veneer}},
\href{https://github.com/yang-mills-group/OpenAI-YM-20261028/blob/3456e5e75cf8d4cd558955a5c7638fba2d68f464/YM/Next/STATIONARY/Shape/SoftTransfers.lean\#L59-L103}{\texttt{extra\_curves}},
\href{https://github.com/yang-mills-group/OpenAI-YM-20261028/blob/3456e5e75cf8d4cd558955a5c7638fba2d68f464/YM/Next/STATIONARY/OneLoop/Lattice.lean\#L528-L535}{\texttt{saunter}},
\href{https://github.com/yang-mills-group/OpenAI-YM-20261028/blob/3456e5e75cf8d4cd558955a5c7638fba2d68f464/YM/Next/STATIONARY/OneLoop/Lattice.lean\#L41-L54}{\texttt{tideMark}}.
\emph{(proof read at top level)}

\paragraph{\texorpdfstring{4.4 Two-leg kernels: the \(F^2\)-coefficient
and
irrelevance}{4.4 Two-leg kernels: the F\^{}2-coefficient and irrelevance}}\label{two-leg-kernels-the-f2-coefficient-and-irrelevance}

\textbf{Lemma 4.7 (a two-leg kernel without \(F^2\)-term vanishes to
fourth order).} Let \(G\) be simple. Let \(f\) be a two-leg kernel that
is strip-analytic, transversal, real, colour invariant, symmetric under
exchange of the legs, and invariant under signed permutations. If the
\(p_q^2\)-coefficient of \(f_{(b,b),(w,w)}(p)\), for \(p\) along an axis
\(q\neq w\), vanishes for one colour \(b\), then \(f\) vanishes to order
\(4\) at zero momentum. Consequently, for every bed with positive metric
strip width, \(\hat H(0)=1\) and signed-permutation-invariant quadratic
form, \(f-\alpha_b(f)\,\hat K_B\) vanishes to order \(4\), where
\(\hat K_B\) is the bed's quadratic kernel. \emph{Sketch:}

\begin{itemize}
\tightlist
\item
  Transversality gives vanishing to order \(2\).
\item
  The third-order coefficient is odd under \(p\mapsto-p\), hence \(0\)
  by reflection symmetry.
\item
  By transversality, hypercubic symmetry and Schur's lemma, the
  second-order coefficient is a multiple of
  \((p^2\delta_{aa'}-p_ap_{a'})\delta_{bb'}\), and the multiple is the
  given coefficient. \emph{Lean:}
  \href{https://github.com/yang-mills-group/OpenAI-YM-20261028/blob/3456e5e75cf8d4cd558955a5c7638fba2d68f464/YM/Next/STATIONARY/Films/LeadingRay.lean\#L611-L625}{\texttt{TributeColor.vanish}},
  \href{https://github.com/yang-mills-group/OpenAI-YM-20261028/blob/3456e5e75cf8d4cd558955a5c7638fba2d68f464/YM/Next/STATIONARY/Films/LeadingRay.lean\#L710-L727}{\texttt{subtract\_fee}}.
  \emph{(proof read at top level)}
\end{itemize}

\textbf{Lemma 4.8 (loop atoms are irrelevant).} Let
\(R\in\mathrm{Bounds}(\mathbf c;\sigma,C,\Lambda)\). The two-leg kernel
of the loop atoms, symmetrized over legs and lattice symmetries:

\begin{itemize}
\tightlist
\item
  records \(\mathrm{loop}_R\) exactly on every torus;
\item
  is strip-analytic of width \(\sigma\) and size \(\Lambda\,r(\sigma)\);
\item
  vanishes to order \(4\) at zero momentum.
\end{itemize}

\emph{Sketch:} The kernel is
\(\hat s(-p)^{\mathsf T}\hat K(p)\hat s(p)\) with
\(\hat K(p)=\sum\mathrm{coeff}\,e^{ip\cdot y}\). Here \(\hat s=O(p)\),
and the two moment conditions give \(\hat K(0)=0\) and
\(\nabla\hat K(0)=0\). \emph{Lean:}
\href{https://github.com/yang-mills-group/OpenAI-YM-20261028/blob/3456e5e75cf8d4cd558955a5c7638fba2d68f464/YM/Next/STATIONARY/Shape/Scaled.lean\#L231-L234}{\texttt{filing\_ready}}.

\paragraph{4.5 The one-loop term in covariant
form}\label{the-one-loop-term-in-covariant-form}

\textbf{Lemma 4.9 (proof of III.7.6).} Let \(H_{\rm joint}\) be the
Hessian of the joint action before gauge completion, and let
\(\mathsf R\) select its axial directions. Three determinant identities
combine:

\begin{itemize}
\tightlist
\item
  the completed Hessian factors as
  \(\det\mathcal P_\mu=(\det\mathcal G_\mu)^2\det(\mathsf R^{\mathsf T}H_{\rm joint}\mathsf R)\)
  (\href{https://github.com/yang-mills-group/OpenAI-YM-20261028/blob/3456e5e75cf8d4cd558955a5c7638fba2d68f464/YM/Next/STATIONARY/Joint/Completion.lean\#L491-L502}{\texttt{root\_determines}});
\item
  the axial Hessian satisfies
  \(\det M=(\det\partial s)^2\det(\mathsf R^{\mathsf T}H_{\rm joint}\mathsf R)\),
  where \(s\) is the left-trivializing chart change
  (\href{https://github.com/yang-mills-group/OpenAI-YM-20261028/blob/3456e5e75cf8d4cd558955a5c7638fba2d68f464/YM/Next/STATIONARY/Joint/HaarJac.lean\#L140-L161}{\texttt{certified}});
\item
  the Haar Jacobian is \(T=|\det\partial s|\)
  (\href{https://github.com/yang-mills-group/OpenAI-YM-20261028/blob/3456e5e75cf8d4cd558955a5c7638fba2d68f464/YM/Next/STATIONARY/Joint/HaarJac.lean\#L63-L114}{\texttt{treasury}},
  by comparison of densities on compact sets).
\end{itemize}

Together they give
\(T^2\det\mathcal P_\mu=\det M(\det\mathcal G_\mu)^2\). \(\det M>0\) by
the second-variation lower bound and \(\det\mathcal G_\mu>0\) since
\(\mathcal G_\mu\ge\mathfrak m_\mu>0\); take logarithms.

In addition, the one-loop function is analytic at \(0\)
(\href{https://github.com/yang-mills-group/OpenAI-YM-20261028/blob/3456e5e75cf8d4cd558955a5c7638fba2d68f464/YM/Next/STATIONARY/Level/SliceGate.lean\#L76-L103}{\texttt{born}}),
and a finite bound proved for one forest transfers to any other
(\href{https://github.com/yang-mills-group/OpenAI-YM-20261028/blob/3456e5e75cf8d4cd558955a5c7638fba2d68f464/YM/Next/STATIONARY/Level/SliceGate.lean\#L105-L113}{\texttt{transfer\_bound}}).
\emph{Lean:}
\href{https://github.com/yang-mills-group/OpenAI-YM-20261028/blob/3456e5e75cf8d4cd558955a5c7638fba2d68f464/YM/Next/STATIONARY/Joint/HaarJac.lean\#L177-L202}{\texttt{tribute}}.
\emph{(proof read; determinant factorization inputs statement only)}

\paragraph{4.6 Uniform invertibility of the fluctuation
operators}\label{uniform-invertibility-of-the-fluctuation-operators}

\textbf{Lemma 4.10 (real stability and complex invertibility).} Fix a
stage \(j\), a penalty \(a>0\), a lower bound \(\varepsilon>0\) and a
bound \(A_F\ge0\) for the gauge-completing map \(F\). Let
\[\mathcal C_F(p)=\mathcal Q_B(p)+F(-p)^{\mathsf T}F(p),\] where
\(\mathcal Q_B\) is the block-Bloch symbol of the joint quadratic form
in III.7.5. Suppose \(B\) has metric strip bound \((\sigma/L,\Omega)\)
and Floor lower bound \(\varepsilon\) on every torus, and \(F\) is
holomorphic and periodic in width \(\sigma>0\), with block-scaled norm
at most \(A_F\). On real momenta require
\[F(p)R_{\rm gauge}(p)=\mathcal G_{1/L}(p),\qquad F(-p)^{\mathsf T}=F(p)^*,\]
where \(R_{\rm gauge}\) inserts a fine gauge direction into the joint
variables and \(\mathcal G_{1/L}\) is the free regulated ghost symbol.

\begin{enumerate}
\def\labelenumi{(\alph{enumi})}
\item
  There is \(y>0\), depending only on \((G,j,a,A_F,\varepsilon)\), such
  that \(\|\mathcal C_F(p)v\|\ge y\|v\|\) in the block-scaled
  coordinates for every real \(p\), every \(L\) and every patch of stage
  \(j\).
\item
  If \(0<\sigma\le1\), \(\Omega\ge0\), the patch is admissible, and
  \(F\) represents the rooted gauge-completing map on a sequence of tori
  with coarse sides tending to infinity, then there are
  \(\eta\in(0,\sigma]\) and \(C\ge0\), depending only on
  \((G,j,K,\varepsilon,\sigma,\Omega,A_F)\), such that
  \(\mathcal C_F^{-1}\) is holomorphic in width \(\eta\) and has
  block-scaled norm at most \(C\), uniformly in \(L\).
\end{enumerate}

\emph{Sketch:} The real estimate controls the quotient by gauge
directions using the quadratic form and block penalty; the displayed
identity controls the remaining gauge directions. Strip bounds give a
uniform perturbation estimate from real to nearby complex momenta.
Choose the width so that this perturbation is smaller than the real
lower bound.

\emph{Lean:}
\href{https://github.com/yang-mills-group/OpenAI-YM-20261028/blob/3456e5e75cf8d4cd558955a5c7638fba2d68f464/YM/Next/STATIONARY/Spectral/RealGap.lean\#L647-L664}{\texttt{stadium}},
\href{https://github.com/yang-mills-group/OpenAI-YM-20261028/blob/3456e5e75cf8d4cd558955a5c7638fba2d68f464/YM/Next/STATIONARY/Spectral/ComplexGap.lean\#L330-L353}{\texttt{post\_cert}}.
\emph{(proof read at top level; real quotient and strip perturbation
estimates not expanded)}

\textbf{Lemma 4.11 (all five operators).} Fix the numerical parameters
of III.7.4. There are \(0<\eta\le w_0\le\min(w,\delta,1/128)\) and
\(C\ge0\), depending only on the free-field bounds, \(G,K\) and \(J\),
with the following property. Under the geometric, fine-bed, fine-datum
and interpolation hypotheses of III.7.4, form the actual fluctuation
symbols using the free interpolation \(A_0\) and the block predictor. At
zero background, the following five operators are invertible in width
\(\eta\), with inverse norm at most \(C\) in their respective
block-scaled coordinates, uniformly in \(L\) and in the input path
width, cubic/quartic amplitude and loop amplitude:

\begin{itemize}
\tightlist
\item
  the auxiliary operator;
\item
  its block-aliased copy;
\item
  the completed Hessian;
\item
  the completed Hessian without the contact term;
\item
  the ghost operator.
\end{itemize}

The unaliased free auxiliary operator and the free regulated ghost
operator satisfy corresponding inverse bounds. \emph{Sketch:} Bound the
fine auxiliary and ghost inverses directly in Fourier space and transfer
the auxiliary bound to its aliased copy. Apply Lemma 4.10 to the joint
completion, then use finite-volume compatibility for the completed and
contact-removed symbols. Take minima of widths and maxima of constants.
\emph{Lean:}
\href{https://github.com/yang-mills-group/OpenAI-YM-20261028/blob/3456e5e75cf8d4cd558955a5c7638fba2d68f464/YM/Next/STATIONARY/Films/RankColor.lean\#L25-L68}{\texttt{honored}},
\href{https://github.com/yang-mills-group/OpenAI-YM-20261028/blob/3456e5e75cf8d4cd558955a5c7638fba2d68f464/YM/Next/STATIONARY/Films/RankColor.lean\#L15-L22}{\texttt{GrantF}},
\href{https://github.com/yang-mills-group/OpenAI-YM-20261028/blob/3456e5e75cf8d4cd558955a5c7638fba2d68f464/YM/WilsonBulk/Spectral/Joint.lean\#L44-L69}{\texttt{completion\_choice}}.
\emph{(proof read at top level)}

\textbf{Lemma 4.12 (proof of III.7.5).}

\begin{enumerate}
\def\labelenumi{\arabic{enumi}.}
\tightlist
\item
  Choose \((\eta,C)\) by Lemma 4.10(b) for the regulator family of the
  completed operator.
\item
  For each \(L\), a constant \(W(L)\) converts block-scaled norms into
  the plain strip norm.
\item
  The bound on the inverse gives the strip bound
  \((\eta/2,\ \mathrm{cost}(\eta,CW(L)))\).
\item
  The inverse Fourier coefficients are real by finite-volume
  compatibility. On every torus their Fourier multiplier is the inverse
  symbol, which gives the matrix identity.
\end{enumerate}

\emph{Lean:}
\href{https://github.com/yang-mills-group/OpenAI-YM-20261028/blob/3456e5e75cf8d4cd558955a5c7638fba2d68f464/YM/WilsonBulk/Spectral/Backing.lean\#L54-L66}{\texttt{uniform\_backing}}.

\paragraph{4.7 The slope of the regulated
determinant}\label{the-slope-of-the-regulated-determinant}

\textbf{Lemma 4.13 (the continuum bubble).} Assume \(G\) is simple. Take
an external field of colour \(b\), polarization \(e_1\) and momentum
\(t\,e_0\), and loop momentum \(k\neq0\).

\begin{itemize}
\tightlist
\item
  The vector bubble has orbital vertex \(-2ik_1\) and spin vertex
  \(\pm2it\) in the \((0,1)\) plane.
\item
  Both bubbles carry the colour trace
  \(\mathrm{tr}(\mathrm{ad}_b^2)=-C_G\), which uses simplicity.
\item
  The propagators are \(|k|^{-2}\) and \(|k+te_0|^{-2}\).
\end{itemize}

Writing \(c_{\rm gh}(t)\) for the ghost bubble and \(c_{\rm vec}(t)\)
for the vector bubble,
\[\partial_t^2c_{\rm gh}(0)-\tfrac12\,\partial_t^2c_{\rm vec}(0)=2h(k),\qquad h(k)=4C_G\Big(\frac{k_1^2}{|k|^6}-\frac{4k_1^2k_0^2}{|k|^8}-\frac1{|k|^4}\Big).\]
The terms in \(k_1^2\) come from the orbital coupling: one adjoint
scalar's worth, \(\tfrac12\cdot4\) vector components minus one ghost.
The term \(-|k|^{-4}\) comes from the spin coupling of the vector field.
\emph{Lean:}
\href{https://github.com/yang-mills-group/OpenAI-YM-20261028/blob/3456e5e75cf8d4cd558955a5c7638fba2d68f464/YM/Next/STATIONARY/OneLoop/Bubbles.lean\#L127-L141}{\texttt{haloRate}},
\href{https://github.com/yang-mills-group/OpenAI-YM-20261028/blob/3456e5e75cf8d4cd558955a5c7638fba2d68f464/YM/Next/STATIONARY/OneLoop/Bubbles.lean\#L125}{\texttt{poleVec}}.

\textbf{Lemma 4.14 (the shell integral).} For \(0<\mu\le1\),
\[\frac1{(2\pi)^4}\int_{\mu<|k|\le1}h(k)\,d^4k=-\kappa\log\mu^{-1}.\]
\emph{Sketch:} By rotation invariance
\(\int k_1^2|k|^{-6}=\tfrac14\int|k|^{-4}\), and
\(\int k_1^2k_0^2|k|^{-8}=\tfrac1{24}\int|k|^{-4}\), the latter from
averaging over \(e_i\pm e_j\). Also
\(\int_{\mu<|k|\le1}|k|^{-4}=2\pi^2\log\mu^{-1}\). So \(h\) averages to
\(4C_G(\tfrac14-\tfrac16-1)|k|^{-4}=-\tfrac{11}3C_G|k|^{-4}\): orbital
\(+\tfrac13C_G\), spin \(-4C_G\). Finally
\(\tfrac{11}3C_G\cdot2\pi^2/(2\pi)^4=\kappa\). \emph{Lean:}
\href{https://github.com/yang-mills-group/OpenAI-YM-20261028/blob/3456e5e75cf8d4cd558955a5c7638fba2d68f464/YM/Next/STATIONARY/OneLoop/MomentumIntegral.lean\#L326-L338}{\texttt{poleSpin}},
\href{https://github.com/yang-mills-group/OpenAI-YM-20261028/blob/3456e5e75cf8d4cd558955a5c7638fba2d68f464/YM/Next/STATIONARY/OneLoop/MomentumIntegral.lean\#L270-L275}{\texttt{joy\_two}},
\href{https://github.com/yang-mills-group/OpenAI-YM-20261028/blob/3456e5e75cf8d4cd558955a5c7638fba2d68f464/YM/Next/STATIONARY/OneLoop/MomentumIntegral.lean\#L292-L304}{\texttt{joy\_heights}},
\href{https://github.com/yang-mills-group/OpenAI-YM-20261028/blob/3456e5e75cf8d4cd558955a5c7638fba2d68f464/YM/Next/STATIONARY/OneLoop/MomentumIntegral.lean\#L200-L202}{\texttt{polar\_m}}.

\textbf{Lemma 4.15 (the lattice integrand is close to \(h\)).} Under the
hypotheses of III.7.3, let \(\Phi_\mu\) be the unintegrated
second-derivative kernel of \(\Pi_\mu\), built from the vertex kernels
of Lemma 4.6:
\[\Phi_\mu=\tfrac12\,(\text{vector film})-(\text{ghost film}).\] Let
\(k_b(k)=\tfrac12\partial_s^2\Phi_{\mu,(1,b),(1,b)}(s\,e_0;k)|_{s=0}\)
be its integrand at loop momentum \(k\). The recorded quadratic Taylor
kernel is
\[\hat Y_\mu=\frac1{2(2\pi)^4}\int_{[-\pi,\pi]^4}\Phi_\mu(\,\cdot\,;k)\,d^4k,\]
and
\(\alpha_b(\hat Y_\mu)=(2\pi)^{-4}\int_{[-\pi,\pi]^4}\mathrm{Re}\,k_b(k)\,d^4k\).
There are \(a,b'\ge0\), uniform over all data of III.7.3 and all
\(\mu\in(0,1]\), such that:
\[|k_b(k)|\le\frac{a}{(|k|+\mu)^4}\ \text{on the Brillouin cube},\qquad |k_b(k)-h(k)|\le b'\big(|k|^{-3}+\mu^2|k|^{-6}\big)\ \text{for }|k|\ge\mu.\]
\emph{Sketch:} Each film is (tadpole) \(-\) (bubble). The tadpoles are
\(O(|k|^{-2}+\mu^2|k|^{-6})\). The lattice bubbles differ from the
continuum ones by \(O(|k|^{-3}+\mu^2|k|^{-6})\), using the softness of
Lemma 4.6. Lemma 4.13 identifies the continuum combination. \emph{Lean:}
\href{https://github.com/yang-mills-group/OpenAI-YM-20261028/blob/3456e5e75cf8d4cd558955a5c7638fba2d68f464/YM/Next/STATIONARY/OneLoop/Bubbles.lean\#L278-L292}{\texttt{astronomy}},
\href{https://github.com/yang-mills-group/OpenAI-YM-20261028/blob/3456e5e75cf8d4cd558955a5c7638fba2d68f464/YM/Next/STATIONARY/OneLoop/Bubbles.lean\#L179-L184}{\texttt{Rally.coin}},
\href{https://github.com/yang-mills-group/OpenAI-YM-20261028/blob/3456e5e75cf8d4cd558955a5c7638fba2d68f464/YM/Next/STATIONARY/OneLoop/Bubbles.lean\#L191-L222}{\texttt{Rally.shell}},
\href{https://github.com/yang-mills-group/OpenAI-YM-20261028/blob/3456e5e75cf8d4cd558955a5c7638fba2d68f464/YM/Next/STATIONARY/OneLoop/MomentumIntegral.lean\#L136-L159}{\texttt{polar\_clock}},
\href{https://github.com/yang-mills-group/OpenAI-YM-20261028/blob/3456e5e75cf8d4cd558955a5c7638fba2d68f464/YM/Next/STATIONARY/OneLoop/Lattice.lean\#L333}{\texttt{amena}}.
\emph{(proof read at top level; analytic tadpole and bubble estimates
not expanded)}

\textbf{Lemma 4.16 (from the bounds to the slope).} For \(0<\mu\le1\),
under the strip regularity and two bounds of Lemma 4.15 with
\(a,b'\ge0\),
\[\big|\alpha_b+\kappa\log\mu^{-1}\big|\le(2\pi)^{-4}\big(a\tfrac{\pi^2}2+a(2\pi)^4+3\pi^2b'\big).\]
\emph{Sketch:} Split the cube into three regions.

\begin{itemize}
\tightlist
\item
  On \(|k|\le\mu\): the integrand is \(\le a\mu^{-4}\) and the volume is
  \(\pi^2\mu^4/2\).
\item
  On the shell: compare with \(h\) and use Lemma 4.14; the error is
  \(\int b'(|k|^{-3}+\mu^2|k|^{-6})\le3\pi^2b'\).
\item
  On \(|k|>1\): the integrand is \(\le a\).
\end{itemize}

\emph{Lean:}
\href{https://github.com/yang-mills-group/OpenAI-YM-20261028/blob/3456e5e75cf8d4cd558955a5c7638fba2d68f464/YM/Next/STATIONARY/OneLoop/MomentumIntegral.lean\#L346-L408}{\texttt{absoluteGeode}}.

\textbf{Lemma 4.17 (proof of III.7.3).}

\begin{enumerate}
\def\labelenumi{\arabic{enumi}.}
\tightlist
\item
  Lemma 4.6 and Lemma 4.1 give the vertex kernels of
  \(\mathrm{main}_{\mathbf c,R}=E_W+(\mathrm{main}-E_W)\).
\item
  The film \(\Phi_{1/L}\) represents \(\Pi_{1/L}\)
  (\href{https://github.com/yang-mills-group/OpenAI-YM-20261028/blob/3456e5e75cf8d4cd558955a5c7638fba2d68f464/YM/Next/STATIONARY/OneLoop/Lattice.lean\#L356-L363}{\texttt{gemF}}).
  For fixed \(L\) it is strip-analytic, since the vector and ghost
  symbols are invertible at real momenta, by Floor and the regulator,
  and hence in a strip
  (\href{https://github.com/yang-mills-group/OpenAI-YM-20261028/blob/3456e5e75cf8d4cd558955a5c7638fba2d68f464/YM/Next/STATIONARY/OneLoop/Lattice.lean\#L451-L456}{\texttt{HouseWord.done}}).
\item
  One half of the normalized loop-momentum integral of the
  second-derivative film is the Taylor kernel \(\hat Y\). Comparing the
  torus Riemann sum with the integral gives an error
  \(|\Lambda_{n_i}|K_Le^{-z_Ln_i/4}\)
  (\href{https://github.com/yang-mills-group/OpenAI-YM-20261028/blob/3456e5e75cf8d4cd558955a5c7638fba2d68f464/YM/Next/STATIONARY/Loops/Fourier.lean\#L454-L482}{\texttt{profit}}).
\item
  Lemmas 4.15--4.16 with \(\mu=1/L\) give (b), with
  \(D=(2\pi)^{-4}(\dots)\) independent of \(L\).
\end{enumerate}

\emph{Lean:}
\href{https://github.com/yang-mills-group/OpenAI-YM-20261028/blob/3456e5e75cf8d4cd558955a5c7638fba2d68f464/YM/Next/STATIONARY/FullPlanned.lean\#L17-L52}{\texttt{pouring}}.

\paragraph{4.8 From the regulated determinant to the blocked
determinant}\label{from-the-regulated-determinant-to-the-blocked-determinant}

\textbf{Lemma 4.18 (comparison of \(F^2\)-coefficients).} Assume \(G\)
is simple and the hypotheses of III.7.4, together with the
interpolation, vertex and inverse bounds of Lemma 4.11. Consider the
two-leg films of:

\begin{itemize}
\tightlist
\item
  the medal film: \(\tfrac12\log\det\) of the completed Hessian minus
  \(\log\det\) of the ghost operator, with regulator \(1/L\);
\item
  the core film: the same with the contact term removed;
\item
  the auxiliary film: built from the block-aliased auxiliary operator.
\end{itemize}

Then:

\begin{itemize}
\tightlist
\item
  the medal film is strip-analytic of width \(\eta/2\) and size
  \(L^4a\);
\item
  one half of its normalized loop-momentum integral records the
  determinant part of the one-loop function along the branch on
  \(\Lambda_{m_j}\), with error \(|\Lambda_{m_j}|\,c_Le^{-\eta m_j/8}\);
\item
  \(|\alpha_b(\text{medal})-\alpha_b(\text{auxiliary})|\le d_{\rm cmp}\),
  with \(d_{\rm cmp}\) depending on the stage and vertex bounds,
  including the input width and amplitude, but not on \(L\).
\end{itemize}

\emph{Sketch:}

\begin{itemize}
\tightlist
\item
  For the contact contribution, the tested diagonal entry of the
  integrated kernel vanishes to order \(3\) along the external momentum
  ray. Hence \(\alpha(\text{medal})=\alpha(\text{core})\)
  (\href{https://github.com/yang-mills-group/OpenAI-YM-20261028/blob/3456e5e75cf8d4cd558955a5c7638fba2d68f464/YM/Next/STATIONARY/Films/LeadingRay.lean\#L789-L818}{\texttt{drift\_tip}}).
\item
  The difference core \(-\) auxiliary is strip-analytic with an
  \(L\)-independent bound. Cauchy estimates bound its
  \(F^2\)-coefficient
  (\href{https://github.com/yang-mills-group/OpenAI-YM-20261028/blob/3456e5e75cf8d4cd558955a5c7638fba2d68f464/YM/Next/STATIONARY/Loops/HessianMatch.lean\#L24-L41}{\texttt{boundedTip}}).
\end{itemize}

\emph{Lean:}
\href{https://github.com/yang-mills-group/OpenAI-YM-20261028/blob/3456e5e75cf8d4cd558955a5c7638fba2d68f464/YM/Next/STATIONARY/Loops/HessianMatch.lean\#L67-L96}{\texttt{quoting}},
\href{https://github.com/yang-mills-group/OpenAI-YM-20261028/blob/3456e5e75cf8d4cd558955a5c7638fba2d68f464/YM/Next/STATIONARY/Films/RankColor.lean\#L84}{\texttt{medalCue}},
\href{https://github.com/yang-mills-group/OpenAI-YM-20261028/blob/3456e5e75cf8d4cd558955a5c7638fba2d68f464/YM/Next/STATIONARY/Films/RankColor.lean\#L86}{\texttt{coreCue}},
\href{https://github.com/yang-mills-group/OpenAI-YM-20261028/blob/3456e5e75cf8d4cd558955a5c7638fba2d68f464/YM/Next/STATIONARY/Films/RankColor.lean\#L85}{\texttt{rankCue}}.
\emph{(proof read at top level; contact identities not expanded)}

\textbf{Lemma 4.19 (aliasing preserves \(\alpha\)).} Under the
hypotheses of Lemma 4.18, the \(F^2\)-coefficient of the auxiliary film
on the coarse lattice equals \(\alpha_b(\hat Y)\), where \(\hat Y\) is
the kernel of III.7.3 for \(\Pi_{1/L}\) on the fine tori
\(\Lambda_{Lm_i}\). \emph{Sketch:} The coarse symbol at momentum \(p\)
is the sum over the \(L^4\) aliases of the fine symbol. The nonzero
aliases vanish to order \(4\). The zero alias agrees through order two
with \(L^4L^{-2}\) times the fine kernel at \(p/L\). The additional
\(L^{-2}\) from momentum rescaling gives \(L^4L^{-2}L^{-2}=1\), so the
\(F^2\) coefficient is unchanged. \emph{Lean:}
\href{https://github.com/yang-mills-group/OpenAI-YM-20261028/blob/3456e5e75cf8d4cd558955a5c7638fba2d68f464/YM/Next/STATIONARY/Loops/HessianMatch.lean\#L719-L741}{\texttt{pole\_reduced}},
\href{https://github.com/yang-mills-group/OpenAI-YM-20261028/blob/3456e5e75cf8d4cd558955a5c7638fba2d68f464/YM/Next/STATIONARY/Loops/HessianMatch.lean\#L255-L278}{\texttt{carry\_plan}},
\href{https://github.com/yang-mills-group/OpenAI-YM-20261028/blob/3456e5e75cf8d4cd558955a5c7638fba2d68f464/YM/Next/STATIONARY/Loops/HessianMatch.lean\#L670-L705}{\texttt{fine\_alpha}}.
\emph{(proof read at top level; high-alias and zero-alias estimates
statement only)}

\paragraph{4.9 The step and the
iteration}\label{the-step-and-the-iteration}

\textbf{Lemma 4.20 (construction of the coarse datum).} Assume the
geometric, bed, Haar-chart and stationary-branch hypotheses of III.7.4.
Fix a common Fourier width \(\sigma>0\), coarse-bed width \(w>0\),
coarse metric and root bound \(M\ge0\), and a bound on the coarse direct
clip coefficient. Require \(H_T,\rho_T\in\mathrm{strip}(w,M)\). There
are \(s\in(0,\min(\sigma,w)]\) and \(K,D\ge0\), independent of the
following kernel amplitudes, with this property.

Suppose the tree-level action has exact Taylor kernels of orders \(3\)
and \(4\), of width \(\sigma\) and size \(X\ge1\). Suppose a two-leg
kernel \(\hat A\) of width \(\sigma\) and size \(\Phi\) records the
determinant part with errors \(|\Lambda_{m_i}|e_i\), where \(e_i\ge0\)
and \(e_i\to0\). Suppose a two-leg kernel \(\hat L\) of the same width
and size \(R_0\) records the inherited loop functional exactly and
vanishes to order \(4\) at zero momentum. Set
\(\alpha=\alpha_b(\hat A)\) for one fixed colour \(b\).

There is a coarse datum \(Y\), with the straight-route rows of \(T\) and
attached to \(T\), whose classical jets agree through order four. It
satisfies the weight clauses at every \(0<s'\le s\), every
row/cubic/quartic bound \[A'\ge\max\{6Me^{4w},XK\},\] and every loop
bound \[\Lambda'\ge\big(\Phi+|\alpha|C_T+R_0\big)D,\] where \(C_T\) is
the uniform strip bound for the quadratic kernel of \(T\) at width
\(\min(\sigma,w)\).

For every \(i,\Theta,c',\Omega\), the finite bound of III.7.4 holds with
right-hand side \(|\Lambda_{m_i}|\Theta\Omega e^{-c'm_i}\) whenever
\(B_Ge_i\le\Theta\Omega e^{-c'm_i}\), where
\(B_G=2(4\dim\mathfrak g)^2\).

\emph{Sketch:}

\begin{itemize}
\tightlist
\item
  Apply Lemma 4.4 to the tree-level action. Its quadratic jet agrees
  with the bare row action by the stationary-branch hypotheses.
\item
  The determinant and inherited-loop functionals are analytic and
  invariant under gauge, colour and signed-permutation symmetries. Their
  recorded kernels inherit these properties. Thus
  \(\hat A-\alpha\hat K_T+\hat L\) vanishes to order \(4\) by Lemma 4.7
  and the hypothesis on \(\hat L\).
\item
  Represent this residual kernel by two-leg curvature atoms with zero
  zeroth and first displacement moments
  (\href{https://github.com/yang-mills-group/OpenAI-YM-20261028/blob/3456e5e75cf8d4cd558955a5c7638fba2d68f464/YM/Next/STATIONARY/Loops/Curvature.lean\#L122-L159}{\texttt{reach\_symbol}}).
  The weighted Laurent bound gives the stated loop amplitude.
\item
  The Faddeev--Popov identity converts the determinant recording error
  to the error for the one-loop function. Expanding real waves into
  complex basis waves gives the factor \(B_G\)
  (\href{https://github.com/yang-mills-group/OpenAI-YM-20261028/blob/3456e5e75cf8d4cd558955a5c7638fba2d68f464/YM/Next/STATIONARY/Loops/Curvature.lean\#L281-L324}{\texttt{draft\_pipe}}).
  Classical jet equality follows from order-\(5\) ray vanishing
  (\href{https://github.com/yang-mills-group/OpenAI-YM-20261028/blob/3456e5e75cf8d4cd558955a5c7638fba2d68f464/YM/Next/STATIONARY/Loops/Curvature.lean\#L377-L386}{\texttt{achieved}},
  \href{https://github.com/yang-mills-group/OpenAI-YM-20261028/blob/3456e5e75cf8d4cd558955a5c7638fba2d68f464/YM/Next/STATIONARY/Loops/Curvature.lean\#L343-L376}{\texttt{violet}}).
\end{itemize}

\emph{Lean:}
\href{https://github.com/yang-mills-group/OpenAI-YM-20261028/blob/3456e5e75cf8d4cd558955a5c7638fba2d68f464/YM/Next/STATIONARY/Loops/Curvature.lean\#L394-L435}{\texttt{color\_pipe}}.
\emph{(proof read at top level)}

\textbf{Lemma 4.21 (proof of III.7.4).}

\begin{itemize}
\item
  \emph{Coercivity.} Lemma 4.11 at widths \(\min(w,\delta)\).
\item
  \emph{Tree-level vertices.} Cubic and quartic vertex bounds have the
  form \(a+L^{-1}\ell\) and \(c+L^{-1}r\)
  (\href{https://github.com/yang-mills-group/OpenAI-YM-20261028/blob/3456e5e75cf8d4cd558955a5c7638fba2d68f464/YM/Next/STATIONARY/Shape/VertexWeights.lean\#L357-L382}{\texttt{setWest}}).
  By
  \href{https://github.com/yang-mills-group/OpenAI-YM-20261028/blob/3456e5e75cf8d4cd558955a5c7638fba2d68f464/YM/Next/STATIONARY/Shape/Exchange.lean\#L407-L415}{\texttt{exchangeYield}},
  the order-\(3\) and order-\(4\) kernels of the tree-level action, the
  latter including the one-propagator exchange of two cubic vertices,
  are bounded by \[M=1+\tfrac{a+1}6+\tfrac{(c+1)+3(a+1)^2E}{24}.\] Here
  \(E\) bounds the completed free inverse. The term \(3(a+1)^2E\) bounds
  the three pairings of two cubic vertices joined by one propagator.
\item
  \emph{Inherited loops.} The inherited loop kernel has size
  \(L^{-2}\Lambda C_0\)
  (\href{https://github.com/yang-mills-group/OpenAI-YM-20261028/blob/3456e5e75cf8d4cd558955a5c7638fba2d68f464/YM/Next/STATIONARY/Loops/Curvature.lean\#L66-L92}{\texttt{bequest}}).
\item
  \emph{The slope.} Lemma 4.18 and Lemma 4.19 give
  \(\alpha=\alpha_b(\text{medal film})\) with
  \(|\alpha+\kappa\log L|\le d_{\rm cmp}+D_{\rm slope}\).
\item
  \emph{Assembly.} Lemma 4.20 builds \(Y\).
\item
  \emph{Choice of \(L\).} \(L\ge u\) is chosen so that \(L^{-1}\) times
  the input-dependent constants lies below the output thresholds, and
  \(L^{-2}C_0D_{\rm lift}\le\tfrac14\), where \(D_{\rm lift}\) is the
  loop-representation constant of Lemma 4.20
  (\href{https://github.com/yang-mills-group/OpenAI-YM-20261028/blob/3456e5e75cf8d4cd558955a5c7638fba2d68f464/YM/Next/STATIONARY/Level/Transition.lean\#L42-L49}{\texttt{age\_control}}).
  The loop amplitude is the sum of a newly generated term and an
  inherited term. Their bounds are \(\le\Lambda'/2\) and
  \(\le\Lambda'/2\), respectively, when \(\Lambda'\ge\Lambda_0(L)\) and
  \(\Lambda'\ge\Lambda/2\)
  (\href{https://github.com/yang-mills-group/OpenAI-YM-20261028/blob/3456e5e75cf8d4cd558955a5c7638fba2d68f464/YM/Next/STATIONARY/Level/Transition.lean\#L74-L80}{\texttt{cap\_invoice}}).
\end{itemize}

\emph{Lean:}
\href{https://github.com/yang-mills-group/OpenAI-YM-20261028/blob/3456e5e75cf8d4cd558955a5c7638fba2d68f464/YM/Next/STATIONARY/Level/Transition.lean\#L92-L182}{\texttt{transition\_cap}}.
\emph{(proof read at top level)}

\textbf{Lemma 4.22 (loop caps and the iteration).}

\begin{itemize}
\tightlist
\item
  \emph{Loop caps.} For every \(W_j\) there are thresholds \(\Delta_j\)
  such that, once \(\hat\beta\) at the entry exceeds them,
  \(W_{J_{m+1}}\le\mathrm{loopCap}_m\),
  \(\mathrm{loopCap}_{m+1}\le2\,\mathrm{loopCap}_m\) and
  \(\mathrm{loopCap}_m\ge1\). The history must satisfy the stage,
  block-factor and entry conditions of III.2, with positive protection
  exponents nonincreasing with stage. Use
  \(\hat\beta_{m+1}=\hat\beta_m+\kappa\log L_m\), and
  \((\log\log(s+k))^2-(\log\log s)^2\le3k/s\) for each fixed \(k\ge0\)
  and all sufficiently large \(s\)
  (\href{https://github.com/yang-mills-group/OpenAI-YM-20261028/blob/3456e5e75cf8d4cd558955a5c7638fba2d68f464/YM/Next/STATIONARY/Level/Recurse.lean\#L18-L49}{\texttt{bump}}).
\item
  \emph{Iteration.}

  \begin{itemize}
  \tightlist
  \item
    \href{https://github.com/yang-mills-group/OpenAI-YM-20261028/blob/3456e5e75cf8d4cd558955a5c7638fba2d68f464/YM/Next/STATIONARY/Level/Recurse.lean\#L122-L143}{\texttt{claiming}}
    applies Lemma 4.21 to the stage pairs \((j,j)\) and \((j+1,j)\) and
    takes minima and maxima of the constants.
  \item
    \href{https://github.com/yang-mills-group/OpenAI-YM-20261028/blob/3456e5e75cf8d4cd558955a5c7638fba2d68f464/YM/Next/STATIONARY/Level/Recurse.lean\#L254-L285}{\texttt{energized}}
    produces one step's record at level \(m\).
  \item
    \href{https://github.com/yang-mills-group/OpenAI-YM-20261028/blob/3456e5e75cf8d4cd558955a5c7638fba2d68f464/YM/Next/STATIONARY/Level/Recurse.lean\#L307-L335}{\texttt{plowed}}
    assembles the kinetic data by downward induction from \(m=N\),
    starting from the top datum.
  \item
    \href{https://github.com/yang-mills-group/OpenAI-YM-20261028/blob/3456e5e75cf8d4cd558955a5c7638fba2d68f464/YM/Next/STATIONARY/Level/Recurse.lean\#L366-L391}{\texttt{under\_contract}}
    chooses the volume thresholds so that
    \(F(L)e^{-\gamma n}\le\Gamma e^{-(\gamma/2)n}\)
    (\href{https://github.com/yang-mills-group/OpenAI-YM-20261028/blob/3456e5e75cf8d4cd558955a5c7638fba2d68f464/YM/Next/STATIONARY/Level/Recurse.lean\#L353-L358}{\texttt{waits\_up}}).
  \item
    The final iteration uses III.7.3.
  \end{itemize}
\end{itemize}

\emph{Lean:}
\href{https://github.com/yang-mills-group/OpenAI-YM-20261028/blob/3456e5e75cf8d4cd558955a5c7638fba2d68f464/YM/Next/STATIONARY/Level/Recurse.lean\#L51-L74}{\texttt{clocks}},
\href{https://github.com/yang-mills-group/OpenAI-YM-20261028/blob/3456e5e75cf8d4cd558955a5c7638fba2d68f464/YM/Next/STATIONARY/FullPlanned.lean\#L59-L66}{\texttt{geared}}.

\subsubsection{5. Where the difficulty
lies}\label{where-the-difficulty-lies-9}

\begin{itemize}
\tightlist
\item
  \textbf{The coefficient and its identification.} The continuum bubble
  yields the factor \(11/3\) from the orbital and spin vertices and
  \(\mathrm{tr}(\mathrm{ad}_b^2)=-C_G\). Relating this calculation to
  the blocked determinant also requires the exact Faddeev--Popov
  identity, cancellation of the contact contribution to \(\alpha\), a
  comparison uniform in \(L\), and control of every Fourier alias.
\item
  \textbf{Quantitative curvature factorization.} The lattice Poincaré
  lemma must preserve exponential decay, colour invariance and the
  constraints already imposed on other legs. Each one-leg solve loses a
  factor four in strip width; the number of legs is fixed.
\item
  \textbf{Uniformity in the volume.} A single infinite-volume kernel or
  path datum works on every torus in the chosen sequence. Limits along
  increasingly fine momentum meshes transfer Ward identities and colour
  symmetries to that kernel. In the iteration, the sequence consists of
  the sides at one fixed level as the volume index varies.
\item
  \textbf{Order of constants.} The output width and cubic/quartic
  amplitude precede the input path amplitude. Contributions depending on
  the latter carry a factor \(L^{-1}\); inherited loop atoms carry
  \(L^{-2}\). The new loop amplitude is \(O(L^4)\), and the
  reference-coupling thresholds make the allowed loop amplitude large
  enough to absorb it.
\item
  \textbf{Strength of the conclusion.} The slope estimate has an
  \(O(1)\) error uniform in \(L\). This suffices for III.9 when
  \(\log L\) dominates that error. Polymer activities are absent from
  the argument.
\item
  \textbf{The regulator.} The scalar regulator is
  \(\mu^{20}(\Delta_U+\mu^2)^{-9}\). Its componentwise extension to link
  fields respects gauge transformations and unsigned axis permutations;
  full signed-permutation symmetry is not asserted for this auxiliary
  determinant. The blocked action has the full lattice symmetry.
\end{itemize}

\subsubsection{6. Notation}\label{notation-9}

A determinant film records the second derivative before loop-momentum
integration. Its quadratic Taylor kernel is one half of its normalized
loop-momentum integral.

{\def\LTcaptype{none} % do not increment counter
\begin{longtable}[]{@{}
  >{\raggedright\arraybackslash}p{(\linewidth - 4\tabcolsep) * \real{0.3333}}
  >{\raggedright\arraybackslash}p{(\linewidth - 4\tabcolsep) * \real{0.3333}}
  >{\raggedright\arraybackslash}p{(\linewidth - 4\tabcolsep) * \real{0.3333}}@{}}
\toprule\noalign{}
\begin{minipage}[b]{\linewidth}\raggedright
Written proof
\end{minipage} & \begin{minipage}[b]{\linewidth}\raggedright
Lean
\end{minipage} & \begin{minipage}[b]{\linewidth}\raggedright
Meaning
\end{minipage} \\
\midrule\noalign{}
\endhead
\bottomrule\noalign{}
\endlastfoot
\(R=(\ell,\mathcal A^{(2)},\mathcal A^{(3)},\mathcal A^{(4)})\) &
\href{https://github.com/yang-mills-group/OpenAI-YM-20261028/blob/3456e5e75cf8d4cd558955a5c7638fba2d68f464/YM/RG/PathActions.lean\#L150-L154}{\texttt{RG.PathData}}
& path datum (rows, loop/cubic/quartic atoms) \\
leg, atom &
\href{https://github.com/yang-mills-group/OpenAI-YM-20261028/blob/3456e5e75cf8d4cd558955a5c7638fba2d68f464/YM/RG/PathActions.lean\#L38-L40}{\texttt{RG.Slot}},
\href{https://github.com/yang-mills-group/OpenAI-YM-20261028/blob/3456e5e75cf8d4cd558955a5c7638fba2d68f464/YM/RG/PathActions.lean\#L84-L86}{\texttt{RG.Atom}}
& transported curvature; curvature monomial \\
\(\mathrm{Bounds}(\mathbf c;\sigma,C,\Lambda)\) &
\href{https://github.com/yang-mills-group/OpenAI-YM-20261028/blob/3456e5e75cf8d4cd558955a5c7638fba2d68f464/YM/RG/PathActions.lean\#L173-L185}{\texttt{PathData.Bounds}}
& exponential summability, invariance and moment clauses \\
\(\mathrm{main}_{\mathbf c,R}\), \(\mathrm{loop}_R\) &
\href{https://github.com/yang-mills-group/OpenAI-YM-20261028/blob/3456e5e75cf8d4cd558955a5c7638fba2d68f464/YM/PCalc/Chart.lean\#L27-L28}{\texttt{mainTerm}},
\texttt{loop} & leading part; loop functional \\
clip \(\mathbf c\), \(c_{\rm dir}\) &
\href{https://github.com/yang-mills-group/OpenAI-YM-20261028/blob/3456e5e75cf8d4cd558955a5c7638fba2d68f464/YM/RG/PathActions.lean\#L139-L146}{\texttt{RG.Clip}}
& clipping parameters; \texttt{direct} \\
bed \(B\), Floor, attached &
\href{https://github.com/yang-mills-group/OpenAI-YM-20261028/blob/3456e5e75cf8d4cd558955a5c7638fba2d68f464/YM/PREC/Stacks.lean\#L35-L40}{\texttt{Bed}},
\href{https://github.com/yang-mills-group/OpenAI-YM-20261028/blob/3456e5e75cf8d4cd558955a5c7638fba2d68f464/YM/PREC/Stacks.lean\#L61-L68}{\texttt{Floor}},
\href{https://github.com/yang-mills-group/OpenAI-YM-20261028/blob/3456e5e75cf8d4cd558955a5c7638fba2d68f464/YM/PREC/Stacks.lean\#L44-L46}{\texttt{Attached}}
& free quadratic-form data of a level \\
straight-route rows of \(\rho\) &
\href{https://github.com/yang-mills-group/OpenAI-YM-20261028/blob/3456e5e75cf8d4cd558955a5c7638fba2d68f464/YM/Next/STATIONARY/Level/SliceGate.lean\#L169-L171}{\texttt{rowPath}},
\href{https://github.com/yang-mills-group/OpenAI-YM-20261028/blob/3456e5e75cf8d4cd558955a5c7638fba2d68f464/YM/Next/STATIONARY/Level/SliceGate.lean\#L173-L174}{\texttt{emitBase}}
& rows realizing the bed roots on straight lattice routes \\
strip-analytic \((\sigma,K)\) &
\href{https://github.com/yang-mills-group/OpenAI-YM-20261028/blob/3456e5e75cf8d4cd558955a5c7638fba2d68f464/YM/Next/STATIONARY/Infrastructure/Jets.lean\#L199}{\texttt{Firm}}
& bounded holomorphic periodic kernel \\
vanishing to order \(s\) along rays &
\href{https://github.com/yang-mills-group/OpenAI-YM-20261028/blob/3456e5e75cf8d4cd558955a5c7638fba2d68f464/YM/Next/STATIONARY/Infrastructure/Jets.lean\#L28-L29}{\texttt{Rays}}
& raywise \(O(t^s)\) at \(t=0\) \\
\((t{+}1)\)-leg kernel; records & \texttt{Picture},
\href{https://github.com/yang-mills-group/OpenAI-YM-20261028/blob/3456e5e75cf8d4cd558955a5c7638fba2d68f464/YM/Next/STATIONARY/Lift/Pathways.lean\#L491-L492}{\texttt{Recorded}}
& Fourier kernel of a Taylor coefficient \\
\(\alpha_b(\hat K)\) &
\href{https://github.com/yang-mills-group/OpenAI-YM-20261028/blob/3456e5e75cf8d4cd558955a5c7638fba2d68f464/YM/Next/STATIONARY/Films/LeadingRay.lean\#L699}{\texttt{alphaPay}}
& \(F^2\)-coefficient, normalized by the bed \\
\(\Pi_\mu\) &
\href{https://github.com/yang-mills-group/OpenAI-YM-20261028/blob/3456e5e75cf8d4cd558955a5c7638fba2d68f464/YM/Next/STATIONARY/Spectral/AuxReferences.lean\#L138}{\texttt{PoleTest}},
\href{https://github.com/yang-mills-group/OpenAI-YM-20261028/blob/3456e5e75cf8d4cd558955a5c7638fba2d68f464/YM/Next/STATIONARY/Spectral/AuxReferences.lean\#L46}{\texttt{fineMag}}
& regulated background-field one-loop action \\
\(\mathcal G_\mu\), \(\mathcal P_\mu\) &
\href{https://github.com/yang-mills-group/OpenAI-YM-20261028/blob/3456e5e75cf8d4cd558955a5c7638fba2d68f464/YM/Next/STATIONARY/Joint/Regulator.lean\#L338}{\texttt{grenade}},
\href{https://github.com/yang-mills-group/OpenAI-YM-20261028/blob/3456e5e75cf8d4cd558955a5c7638fba2d68f464/YM/Next/STATIONARY/Joint/Completion.lean\#L59-L60}{\texttt{medals}}
& ghost operator; completed Hessian \\
\(\mu=1/L\) &
\href{https://github.com/yang-mills-group/OpenAI-YM-20261028/blob/3456e5e75cf8d4cd558955a5c7638fba2d68f464/YM/Next/STATIONARY/Spectral/Slabs.lean\#L189}{\texttt{inch}}
& regulator mass \\
\(h(k)\); \(k_b(k)\) &
\href{https://github.com/yang-mills-group/OpenAI-YM-20261028/blob/3456e5e75cf8d4cd558955a5c7638fba2d68f464/YM/Next/STATIONARY/OneLoop/Bubbles.lean\#L125}{\texttt{poleVec}};
\href{https://github.com/yang-mills-group/OpenAI-YM-20261028/blob/3456e5e75cf8d4cd558955a5c7638fba2d68f464/YM/Next/STATIONARY/OneLoop/Bubbles.lean\#L155}{\texttt{keelRay}}
& continuum bubble integrand; lattice integrand \\
bounds of Lemma 4.15 &
\href{https://github.com/yang-mills-group/OpenAI-YM-20261028/blob/3456e5e75cf8d4cd558955a5c7638fba2d68f464/YM/Next/STATIONARY/OneLoop/Bubbles.lean\#L276}{\texttt{Waves1}},
\href{https://github.com/yang-mills-group/OpenAI-YM-20261028/blob/3456e5e75cf8d4cd558955a5c7638fba2d68f464/YM/Next/STATIONARY/OneLoop/Bubbles.lean\#L277}{\texttt{Waves2}}
& global infrared bound and shell comparison bound \\
\((2\pi)^{-4}\) &
\href{https://github.com/yang-mills-group/OpenAI-YM-20261028/blob/3456e5e75cf8d4cd558955a5c7638fba2d68f464/YM/Next/STATIONARY/OneLoop/MomentumIntegral.lean\#L28}{\texttt{cartWeight}}
& normalization of the momentum integral \\
\(S_t(x,z)\); \(\mathrm{th}\) &
\href{https://github.com/yang-mills-group/OpenAI-YM-20261028/blob/3456e5e75cf8d4cd558955a5c7638fba2d68f464/YM/PCalc/Chart.lean\#L81-L84}{\texttt{tested}};
\href{https://github.com/yang-mills-group/OpenAI-YM-20261028/blob/3456e5e75cf8d4cd558955a5c7638fba2d68f464/YM/PCalc/Loop.lean\#L46-L49}{\texttt{thermal}}
& fine action plus penalty; one-loop function \\
classical match; finite bound &
\href{https://github.com/yang-mills-group/OpenAI-YM-20261028/blob/3456e5e75cf8d4cd558955a5c7638fba2d68f464/YM/PCalc/Chart.lean\#L108-L109}{\texttt{ClassicalMatch}};
\href{https://github.com/yang-mills-group/OpenAI-YM-20261028/blob/3456e5e75cf8d4cd558955a5c7638fba2d68f464/YM/PCalc/Loop.lean\#L62-L71}{\texttt{FiniteBound}}
& items 4 and 5 of III.7.4 \\
slope statement &
\href{https://github.com/yang-mills-group/OpenAI-YM-20261028/blob/3456e5e75cf8d4cd558955a5c7638fba2d68f464/YM/Next/STATIONARY/Level/Transition.lean\#L20-L23}{\texttt{HorizCalc}},
\href{https://github.com/yang-mills-group/OpenAI-YM-20261028/blob/3456e5e75cf8d4cd558955a5c7638fba2d68f464/YM/Next/STATIONARY/Level/Recurse.lean\#L115}{\texttt{Fluent}}
& III.7.3 for given / all parameters \\
step property; step record &
\href{https://github.com/yang-mills-group/OpenAI-YM-20261028/blob/3456e5e75cf8d4cd558955a5c7638fba2d68f464/YM/Next/STATIONARY/Level/Transition.lean\#L86-L90}{\texttt{Transition}};
\href{https://github.com/yang-mills-group/OpenAI-YM-20261028/blob/3456e5e75cf8d4cd558955a5c7638fba2d68f464/YM/Next/STATIONARY/Level/Recurse.lean\#L170-L177}{\texttt{Bubble}}
& III.7.4 \\
top datum &
\href{https://github.com/yang-mills-group/OpenAI-YM-20261028/blob/3456e5e75cf8d4cd558955a5c7638fba2d68f464/YM/Next/STATIONARY/Level/Recurse.lean\#L166}{\texttt{BranchRow}}
& III.7.2 at level \(N\) \\
kinetic data, planned conditions &
\href{https://github.com/yang-mills-group/OpenAI-YM-20261028/blob/3456e5e75cf8d4cd558955a5c7638fba2d68f464/YM/PCalc/Loop.lean\#L74-L77}{\texttt{Kin}},
\href{https://github.com/yang-mills-group/OpenAI-YM-20261028/blob/3456e5e75cf8d4cd558955a5c7638fba2d68f464/YM/PCalc/Loop.lean\#L93-L115}{\texttt{Kin.Planned}},
\href{https://github.com/yang-mills-group/OpenAI-YM-20261028/blob/3456e5e75cf8d4cd558955a5c7638fba2d68f464/YM/Next/STATIONARY/Level/Recurse.lean\#L301-L305}{\texttt{FullPaths}}
& output to III.9 \\
invertibility bounds &
\href{https://github.com/yang-mills-group/OpenAI-YM-20261028/blob/3456e5e75cf8d4cd558955a5c7638fba2d68f464/YM/Next/STATIONARY/Films/RankColor.lean\#L15-L22}{\texttt{GrantF}},
\href{https://github.com/yang-mills-group/OpenAI-YM-20261028/blob/3456e5e75cf8d4cd558955a5c7638fba2d68f464/YM/Next/STATIONARY/Loops/UniformMatrix.lean\#L270-L272}{\texttt{Freezer}}
& strip invertibility with inverse bound \\
medal, core, auxiliary films &
\href{https://github.com/yang-mills-group/OpenAI-YM-20261028/blob/3456e5e75cf8d4cd558955a5c7638fba2d68f464/YM/Next/STATIONARY/Films/RankColor.lean\#L84}{\texttt{medalCue}},
\href{https://github.com/yang-mills-group/OpenAI-YM-20261028/blob/3456e5e75cf8d4cd558955a5c7638fba2d68f464/YM/Next/STATIONARY/Films/RankColor.lean\#L86}{\texttt{coreCue}},
\href{https://github.com/yang-mills-group/OpenAI-YM-20261028/blob/3456e5e75cf8d4cd558955a5c7638fba2d68f464/YM/Next/STATIONARY/Films/RankColor.lean\#L85}{\texttt{rankCue}}
& §4.8 \\
\(\hat\beta_m\), \(\mathrm{loopCap}_m\) &
\href{https://github.com/yang-mills-group/OpenAI-YM-20261028/blob/3456e5e75cf8d4cd558955a5c7638fba2d68f464/YM/RG/History.lean\#L79}{\texttt{hatBeta}},
\href{https://github.com/yang-mills-group/OpenAI-YM-20261028/blob/3456e5e75cf8d4cd558955a5c7638fba2d68f464/YM/RG/History.lean\#L86}{\texttt{loopCap}}
& reference coupling; loop-amplitude cap \\
\(\kappa\) &
\href{https://github.com/yang-mills-group/OpenAI-YM-20261028/blob/3456e5e75cf8d4cd558955a5c7638fba2d68f464/YM/Gauge/Colors.lean\#L184}{\texttt{Gauge.slopeConstant}}
& \(11C_G/(24\pi^2)\) \\
\(D^rF(0)/r!\) &
\href{https://github.com/yang-mills-group/OpenAI-YM-20261028/blob/3456e5e75cf8d4cd558955a5c7638fba2d68f464/YM/Next/STATIONARY/Infrastructure/Jets.lean\#L341}{\texttt{noteAt}}
& complexified Taylor coefficient \\
analytic gauge and lattice symmetry &
\href{https://github.com/yang-mills-group/OpenAI-YM-20261028/blob/3456e5e75cf8d4cd558955a5c7638fba2d68f464/YM/Next/STATIONARY/Shape/ClassicalMatch.lean\#L373-L377}{\texttt{Pointwise}}
& hypotheses of atom emission \\
\(\pi_d(\sigma)\) &
\href{https://github.com/yang-mills-group/OpenAI-YM-20261028/blob/3456e5e75cf8d4cd558955a5c7638fba2d68f464/YM/Next/STATIONARY/Lift/Koszul.lean\#L84}{\texttt{priceDivision}}
& one-coordinate division constant in \(d\) momentum coordinates \\
joint free form; gauge completion &
\href{https://github.com/yang-mills-group/OpenAI-YM-20261028/blob/3456e5e75cf8d4cd558955a5c7638fba2d68f464/YM/Next/STATIONARY/Block/Static.lean\#L20}{\texttt{frontR}},
\href{https://github.com/yang-mills-group/OpenAI-YM-20261028/blob/3456e5e75cf8d4cd558955a5c7638fba2d68f464/YM/Next/STATIONARY/Joint/FourPoint.lean\#L221}{\texttt{postroom}},
\href{https://github.com/yang-mills-group/OpenAI-YM-20261028/blob/3456e5e75cf8d4cd558955a5c7638fba2d68f464/YM/Next/STATIONARY/Joint/Regulator.lean\#L499}{\texttt{cannonade}}
& fine-link and coarse-frame operators \\
\end{longtable}
}

\subsubsection{7. What remains
unexplained}\label{what-remains-unexplained-9}

\begin{enumerate}
\def\labelenumi{\arabic{enumi}.}
\tightlist
\item
  \textbf{Analytic estimates below the stated lemmas.} The real quotient
  estimate and complex-strip perturbation estimate used in Lemma 4.10,
  and the analytic tadpole and bubble bounds used in Lemma 4.15, have
  not been checked through every subsidiary proof. Their statements
  support the estimates given here. The determinant-factorization and
  Haar-density lemmas of Lemma 4.9 remain inputs at statement level. The
  contact identities used in Lemma 4.18 and the high-alias and
  zero-alias estimates used in Lemma 4.19 likewise remain inputs at
  statement level.
\item
  \textbf{Choice of regulator power.} The ninth power is part of the
  explicit definition, and the cited estimates use it. A minimal power,
  or a conceptual reason for choosing exactly nine, has not been
  established.
\item
  \textbf{Supply of the global hypotheses.} III.4 supplies the
  free-field horizons and history conditions. The
  \href{https://github.com/yang-mills-group/OpenAI-YM-20261028/blob/3456e5e75cf8d4cd558955a5c7638fba2d68f464/YM/B70FreeStationaryConnection.lean\#L17-L56}{\texttt{free/stationary\ connection}}
  explicitly constructs the two horizons with identical beds and
  interpolation kernels before applying the iteration. The discharge of
  these conditions in every downstream application has not been checked.
\item
  \textbf{Limits of path data.} The projection results support the
  compactness argument in III.13: exponential weights give Hilbert-space
  parameters for curvature coefficients
  (\href{https://github.com/yang-mills-group/OpenAI-YM-20261028/blob/3456e5e75cf8d4cd558955a5c7638fba2d68f464/YM/UVExtension/Projection/Ink.lean\#L79-L99}{\texttt{fitting}});
  coefficientwise convergence with uniform row bounds gives convergence
  of row-action jets
  (\href{https://github.com/yang-mills-group/OpenAI-YM-20261028/blob/3456e5e75cf8d4cd558955a5c7638fba2d68f464/YM/UVExtension/Projection/Lines.lean\#L89-L109}{\texttt{insideLine}});
  compatible source jets and a common positive Hessian lower bound give
  continuity of stationary, tree-level and one-loop jets
  (\href{https://github.com/yang-mills-group/OpenAI-YM-20261028/blob/3456e5e75cf8d4cd558955a5c7638fba2d68f464/YM/UVExtension/Projection/Fumes.lean\#L56-L87}{\texttt{interview}}).
  The complete compactness construction and passage through all its
  quantifiers belong to III.13 and remain outside this unit.
\end{enumerate}

\subsection{III.8. Connected activities and backward dependence
cones}\label{iii.8.-connected-activities-and-backward-dependence-cones}

\textbf{Part:} III, Ultraviolet construction (renormalization group).
\textbf{Lean:}
\href{https://github.com/yang-mills-group/OpenAI-YM-20261028/tree/3456e5e75cf8d4cd558955a5c7638fba2d68f464/YM/RG08/Steam}{\texttt{YM/RG08/Steam/}},
\href{https://github.com/yang-mills-group/OpenAI-YM-20261028/tree/3456e5e75cf8d4cd558955a5c7638fba2d68f464/YM/RG08/Forms}{\texttt{YM/RG08/Forms/}},
\href{https://github.com/yang-mills-group/OpenAI-YM-20261028/tree/3456e5e75cf8d4cd558955a5c7638fba2d68f464/YM/RG08/Skiff}{\texttt{YM/RG08/Skiff/}},
\href{https://github.com/yang-mills-group/OpenAI-YM-20261028/tree/3456e5e75cf8d4cd558955a5c7638fba2d68f464/YM/RG08/Jetty}{\texttt{YM/RG08/Jetty/}},
\href{https://github.com/yang-mills-group/OpenAI-YM-20261028/tree/3456e5e75cf8d4cd558955a5c7638fba2d68f464/YM/Push/GRID_APPROXIMATION}{\texttt{YM/Push/GRID\_APPROXIMATION/}}.
The support packages are listed in §7. \textbf{Size:} 4,147 lines in the
main-theorem closure:

\begin{itemize}
\tightlist
\item
  1,093 lines for the backward cones (\texttt{GRID\_APPROXIMATION});
\item
  1,083 lines for the connected heat expansion (\texttt{Steam} 138,
  \texttt{Forms} 412, \texttt{Skiff} 306, \texttt{Jetty} 227);
\item
  about 1,970 lines in the support packages listed in §7.
\end{itemize}

The unit has two independent halves.

\begin{itemize}
\tightlist
\item
  \textbf{Part A, connected activities.} Bounds at fixed order for a
  recursive heat expansion, organised by powers of the coupling \(g\),
  the heat time \(t\) and the field. For finite index sets and symmetric
  inputs, its unmarked coefficients satisfy the formal equation for
  \(\log\mathbb E_{t\Gamma}e^V\). A second recursion treats one marked
  insertion. Both admit polynomial bounds in an input size \(W\) and
  Lipschitz bounds for differences, uniform in the index set and
  geometry.
\item
  \textbf{Part B, backward dependence cones.} The pull-back of a coarse
  observable through the block-spin kernel of III.1 is its conditional
  expectation given the fine field. It is local in an explicit fine
  window, and pull-backs of observables on disjoint coarse sets multiply
  exactly. Along a history the radius of the window grows by an explicit
  affine map whose increments, in physical units, decay geometrically.
\end{itemize}

\subsubsection{1. What this unit proves}\label{what-this-unit-proves-10}

\paragraph{Part A}\label{part-a}

\emph{Conventions.}

\begin{itemize}
\item
  \(I\) is an arbitrary set (of cells), \(d\) is a pseudometric on
  \(I\), and \(\sigma\ge0\).
\item
  A rank-\(n\) tensor is a function \(F:I^n\to\mathbb R\) with ordered
  slots and no symmetry factors. For a tuple \(x\in I^n\),
  \(\mathrm{diam}_d(x)=\max(0,\max_{i,j}d(x_i,x_j))\), with value \(0\)
  for the empty tuple.
\item
  The \emph{rooted norm} is
  \[\lVert F\rVert_{\sigma}=\max_{1\le i\le n}\ \sup_{z\in I}\ \sum_{x\in I^n,\ x_i=z}|F(x)|\ e^{\sigma\ \mathrm{diam}_d(x)}\ \in[0,\infty].\]
  For \(n=0\) set \(\lVert F\rVert_\sigma=0\): a rooted bound imposes no
  condition on the scalar coefficient. Suprema over empty index sets
  also have value \(0\).
\item
  A \emph{propagator} is any \(\Gamma:I\times I\to\mathbb R\); symmetry
  is not assumed. It is measured by its weighted row and column sums,
  \[[\Gamma]_\sigma=\max\Big(\sup_u\sum_v|\Gamma(u,v)|e^{\sigma d(u,v)},\ \sup_v\sum_u|\Gamma(u,v)|e^{\sigma d(u,v)}\Big).\]
\item
  A pair \((F,F')\) is \emph{\((W,\delta)\)-controlled} if
  \(\lVert F\rVert_\sigma\le W\), \(\lVert F'\rVert_\sigma\le W\) and
  \(\lVert F-F'\rVert_\sigma\le W\delta\). The same definition applies
  to a pair of propagators, with \([\cdot]_\sigma\) in place of
  \(\lVert\cdot\rVert_\sigma\).
\item
  A \emph{polynomial majorant} is a function \(P\) with
  \(0\le P(x)\le Cx^t\) for all \(x\ge1\), for some \(C\ge0\) and
  \(t\in\mathbb N\). The class is closed under sums and products.
\end{itemize}

\emph{Formal potentials.} A formal potential is a family
\(V=(V_{e,n})_{e,n\ge0}\) of rank-\(n\) tensors; \(V_{e,n}\) is the
coefficient of \(g^e\) with \(n\) field variables. In Gaussian
applications the variables may include spectators, on which \(\Gamma\)
vanishes. \(V\) is \emph{eligible} if \(V_{0,n}=0\) for all \(n\), and
\(V_{e,n}\ne0\) implies \(n\le e+2\) and \(e+n\) even. Eligibility does
not require symmetry.

\emph{Heat coefficients.} The heat coefficients
\(V^{(j)}=V^{(j)}(\Gamma,V)\) are defined by \(V^{(0)}=V\) and

\[V^{(j+1)}_{e,n}=\frac1{j+1}\Big[\binom{n+2}2\Delta_\Gamma V^{(j)}_{e,n+2}+\sum_{c=0}^n\sum_{i=0}^j\sum_{a=0}^e\frac{(c+1)(n-c+1)}2\ V^{(i)}_{a,c+1}\star_\Gamma V^{(j-i)}_{e-a,n-c+1}\Big].\qquad(\mathrm{A1})\]

Here \(\Delta_\Gamma\) (contraction) and \(\star_\Gamma\) (join) are
defined in Step A1. For infinite \(I\), the recursions assign value
\(0\) to a sum that is not summable. The bounds below justify the
required sums absolutely for positive-rank unmarked outputs and for all
marked outputs; they make no convergence assertion for unmarked scalar
contractions. The formal heat-equation interpretation for finite \(I\)
and symmetric \(V,\Gamma\) is given in A1.

\textbf{Proposition III.8.1 (bounds at fixed order for the connected
heat coefficients).} Fix \(j,e,n\in\mathbb N\) and put
\[T_j(e,n)=\lbrace(a,r):1\le a\lt e,\ 1\le r\le a+2\rbrace\cup\lbrace(e,r):1\le e,\ n\lt r\le e+2\rbrace,\]
enlarged by \((e,n)\) when \(j=0\), \(e\ge1\) and \(1\le n\le e+2\).
There is a polynomial majorant \(P\) with the following property. Let
\(I\) be a set, \(d\) a pseudometric on it, \(\sigma\ge0\),
\(\delta\ge0\), \(W\ge1\), \(V,V'\) eligible and \(\Gamma,\Gamma'\)
propagators. Suppose \((\Gamma,\Gamma')\) is \((W,\delta)\)-controlled,
and so is \((V_{a,r},V'_{a,r})\) for every \((a,r)\in T_j(e,n)\). Then
\(\big(V^{(j)}_{e,n}(\Gamma,V),\ V^{(j)}_{e,n}(\Gamma',V')\big)\) is
\((P(W),\delta)\)-controlled.

\begin{itemize}
\tightlist
\item
  \(P\) depends on \((j,e,n)\) only. It does not depend on \(I\), \(d\),
  \(\sigma\) or \(\delta\), hence not on the number of cells, the
  period, the scale or the stage at which it is applied.
\item
  For \(j\ge1\), inputs at the same order \(e\) have strictly higher
  rank.
\item
  For \(n=0\) the conclusion is empty.
\end{itemize}

\emph{Status:} theorem. \emph{Lean:}
\href{https://github.com/yang-mills-group/OpenAI-YM-20261028/blob/3456e5e75cf8d4cd558955a5c7638fba2d68f464/YM/RG08/Steam/Routing.lean\#L34-L83}{\texttt{YM.RG08.Steam.navigation}},
input range
\href{https://github.com/yang-mills-group/OpenAI-YM-20261028/blob/3456e5e75cf8d4cd558955a5c7638fba2d68f464/YM/RG08/Steam/Routing.lean\#L14}{\texttt{taps}},
norms
\href{https://github.com/yang-mills-group/OpenAI-YM-20261028/blob/3456e5e75cf8d4cd558955a5c7638fba2d68f464/YM/RG08/Forms/Ground.lean\#L61}{\texttt{Held}},
\href{https://github.com/yang-mills-group/OpenAI-YM-20261028/blob/3456e5e75cf8d4cd558955a5c7638fba2d68f464/YM/RG08/Forms/Lines.lean\#L18-L22}{\texttt{Wire}},
\href{https://github.com/yang-mills-group/OpenAI-YM-20261028/blob/3456e5e75cf8d4cd558955a5c7638fba2d68f464/YM/RG08/Forms/Cot.lean\#L15-L19}{\texttt{Cot}},
\href{https://github.com/yang-mills-group/OpenAI-YM-20261028/blob/3456e5e75cf8d4cd558955a5c7638fba2d68f464/YM/RG08/Forms/Cot.lean\#L21-L24}{\texttt{WireCot}},
majorants
\href{https://github.com/yang-mills-group/OpenAI-YM-20261028/blob/3456e5e75cf8d4cd558955a5c7638fba2d68f464/YM/RG08/Forms/Powers.lean\#L11}{\texttt{Wage}}.

\textbf{Proposition III.8.2 (support; the heat-time series is finite).}
Let \(V\) be eligible and \(\Gamma\) arbitrary.

\begin{itemize}
\tightlist
\item
  Every \(V^{(j)}\) is eligible.
\item
  \(V^{(j)}_{e,n}\ne0\) implies \(n+2j\le3e\).
\item
  For \(j\ge1\), \(V^{(j)}\) is symmetric under permutations of its
  slots.
\end{itemize}

Hence, at fixed \((e,n)\), the heat-time series is zero if \(n\gt3e\)
and otherwise is a polynomial in \(t\) of degree at most
\(\lfloor(3e-n)/2\rfloor\). For every real \(t\) the complete series is
\[\mathcal E_t(V)_{e,n}=\sum_{j=0}^{3e}t^jV^{(j)}_{e,n}=V_{e,n}+\mathcal R_t(V)_{e,n},\qquad\mathcal R_t(V)_{e,n}=\sum_{j=1}^{3e}t^jV^{(j)}_{e,n},\]
and \(\mathcal E_t(V)\) is eligible. \emph{Status:} theorem.
\emph{Lean:}
\href{https://github.com/yang-mills-group/OpenAI-YM-20261028/blob/3456e5e75cf8d4cd558955a5c7638fba2d68f464/YM/RG08/Steam/Valves.lean\#L50-L64}{\texttt{preserves}},
\href{https://github.com/yang-mills-group/OpenAI-YM-20261028/blob/3456e5e75cf8d4cd558955a5c7638fba2d68f464/YM/RG08/Steam/Valves.lean\#L66-L70}{\texttt{live}},
\href{https://github.com/yang-mills-group/OpenAI-YM-20261028/blob/3456e5e75cf8d4cd558955a5c7638fba2d68f464/YM/RG08/Steam/Valves.lean\#L72}{\texttt{off}},
\href{https://github.com/yang-mills-group/OpenAI-YM-20261028/blob/3456e5e75cf8d4cd558955a5c7638fba2d68f464/YM/RG08/Steam/Gears.lean\#L65-L75}{\texttt{coal}},
\href{https://github.com/yang-mills-group/OpenAI-YM-20261028/blob/3456e5e75cf8d4cd558955a5c7638fba2d68f464/YM/RG08/Steam/Valves.lean\#L74}{\texttt{ember}},
\href{https://github.com/yang-mills-group/OpenAI-YM-20261028/blob/3456e5e75cf8d4cd558955a5c7638fba2d68f464/YM/RG08/Steam/Valves.lean\#L75}{\texttt{cinder}},
\href{https://github.com/yang-mills-group/OpenAI-YM-20261028/blob/3456e5e75cf8d4cd558955a5c7638fba2d68f464/YM/RG08/Steam/Valves.lean\#L77-L80}{\texttt{smoulder}},
\href{https://github.com/yang-mills-group/OpenAI-YM-20261028/blob/3456e5e75cf8d4cd558955a5c7638fba2d68f464/YM/RG08/Steam/Valves.lean\#L82-L89}{\texttt{appropriate}}.

\textbf{Proposition III.8.3 (expansion with one marked insertion).}

\emph{Test functions and their norms.} Let \(\mathcal S\) be the real
vector space of smooth real-valued functions on \(\mathbb R^4\) whose
derivatives of every order have polynomial growth. For \(k\in\mathbb N\)
and \(\sigma\gt0\), using the Euclidean norm \(|y|\), put
\[\lVert f\rVert_{k,\sigma}=\max_{|\alpha|\le k}\sup_{y\in\mathbb R^4}|\partial^\alpha f(y)|e^{-\sigma|y|},\]
and, for a real-linear functional \(\lambda\) on \(\mathcal S\),
\(\lVert\lambda\rVert^{\ast}_{k,\sigma}=\sup\lbrace|\lambda(f)|:\lVert f\rVert_{k,\sigma}\le1\rbrace\).

\emph{Marked tensors and the marked norm.} A marked rank-\(n\) tensor
assigns to each \(x\in I^n\) a real-linear functional \(M(x)\) on
\(\mathcal S\). Let \(h:I\to[0,\infty)\) be 1-Lipschitz for \(d\) (a
height, for instance a distance to the insertion), and set
\(h(x)=\max(0,\max_ih(x_i))\) for a tuple \(x\). The marked norm is the
unrooted weighted \(\ell^1\) norm
\[\lVert M\rVert_{h}=\sum_{x\in I^n}\lVert M(x)\rVert^{\ast}_{k,\sigma}\ e^{\sigma h(x)}.\]
Unlike the rooted norm, it gives a genuine bound for \(n=0\). A pair
\((M,M')\) is \((U,\delta)\)-controlled if \(\lVert M\rVert_h\le U\),
\(\lVert M'\rVert_h\le U\) and \(\lVert M-M'\rVert_h\le U\delta\).

\emph{Marked heat coefficients.} For a marked series \(F=(F_{e,n})\)
they are \(M^{(0)}=F\) and
\[M^{(j+1)}_{e,n}=\frac1{j+1}\Big[\binom{n+2}2\Delta_\Gamma M^{(j)}_{e,n+2}+\sum_{c=0}^n\sum_{i=0}^j\sum_{a=0}^e(c+1)(n-c+1)\ M^{(i)}_{a,c+1}\star_\Gamma V^{(j-i)}_{e-a,n-c+1}\Big].\]
The cross term has twice the coefficient of its unmarked counterpart.
The recursive outputs are initially maps on test functions; the bounds
below prove that they are real-linear. Their interpretation as a
normalised Gaussian expectation is discussed in §7.

\begin{enumerate}
\def\labelenumi{(\roman{enumi})}
\tightlist
\item
  \emph{Fixed-order bound.} Fix \(j,e,n\). There is a polynomial
  majorant \(P\) with the following property, for all \(I,d,h\) as
  above, \(\sigma\gt0\), \(k\in\mathbb N\), \(\delta\ge0\), \(W\ge1\),
  \(U\ge0\), eligible \(V,V'\), propagators \(\Gamma,\Gamma'\) and
  marked series \(F,F'\). Assume:
\end{enumerate}

\begin{itemize}
\tightlist
\item
  \((\Gamma,\Gamma')\) is \((W,\delta)\)-controlled;
\item
  the unmarked inputs \((V_{a,r},V'_{a,r})\) with \(1\le a\le e\) and
  \(1\le r\le a+2\) are \((W,\delta)\)-controlled;
\item
  the marked inputs \((F_{a,r},F'_{a,r})\) are \((U,\delta)\)-controlled
  for: \(a\lt e\) (any \(r\)); or \(a=e\) and \(r\gt n\); or \(j=0\) and
  \((a,r)=(e,n)\).
\end{itemize}

Then
\(\big(M^{(j)}_{e,n}(\Gamma,V,F),\ M^{(j)}_{e,n}(\Gamma',V',F')\big)\)
is \((P(W)U,\delta)\)-controlled, with real-linear output values. The
majorant depends only on \((j,e,n)\); in particular it is independent of
\(k,\sigma\) and the height \(h\).

\begin{enumerate}
\def\labelenumi{(\roman{enumi})}
\setcounter{enumi}{1}
\item
  \emph{Support.} Let \(l\in\mathbb N\). If \(V\) is eligible and
  \(F_{a,r}\ne0\) implies \(r\le l+a\) for every \(a,r\), then
  \(M^{(j)}_{e,n}\ne0\) implies \(n+2j\le l+3e\) and \(n\le l+e\).
\item
  \emph{Full positive-time part.} Fix \(l,e,n\). There is a polynomial
  majorant \(P\) such that the pair
  \[\sum_{j=1}^{l+3e}t^jM^{(j)}_{e,n}(\Gamma,V,F),\qquad\sum_{j=1}^{l+3e}t^jM^{(j)}_{e,n}(\Gamma',V',F')\]
  is \((P(W)U,\delta)\)-controlled for every real \(t\) with
  \(|t|\le W\). The hypotheses are those of (i), with marked inputs
  required only for \(a\lt e\), or \(a=e\) and \(r\gt n\). By (ii), this
  is the whole positive-time part when both \(F_{a,r}\) and \(F'_{a,r}\)
  vanish for \(r\gt l+a\), for every \(a,r\).
\end{enumerate}

\emph{Status:} theorem. This is the form in which the main theorem uses
Part A (in III.14). \emph{Lean:}
\href{https://github.com/yang-mills-group/OpenAI-YM-20261028/blob/3456e5e75cf8d4cd558955a5c7638fba2d68f464/YM/RG08/Skiff/Spoke.lean\#L92-L127}{\texttt{journey}},
\href{https://github.com/yang-mills-group/OpenAI-YM-20261028/blob/3456e5e75cf8d4cd558955a5c7638fba2d68f464/YM/RG08/Skiff/Spoke.lean\#L61-L68}{\texttt{full}},
\href{https://github.com/yang-mills-group/OpenAI-YM-20261028/blob/3456e5e75cf8d4cd558955a5c7638fba2d68f464/YM/RG08/Skiff/Shed.lean\#L25-L36}{\texttt{registry}}.
Test space:
\href{https://github.com/yang-mills-group/OpenAI-YM-20261028/blob/3456e5e75cf8d4cd558955a5c7638fba2d68f464/YM/RG08/Jetty/Cube.lean\#L24}{\texttt{Pearl}}.
Norms:
\href{https://github.com/yang-mills-group/OpenAI-YM-20261028/blob/3456e5e75cf8d4cd558955a5c7638fba2d68f464/YM/RG08/Jetty/Bank.lean\#L36}{\texttt{Sings}},
\href{https://github.com/yang-mills-group/OpenAI-YM-20261028/blob/3456e5e75cf8d4cd558955a5c7638fba2d68f464/YM/RG08/Skiff/Pair.lean\#L11-L17}{\texttt{Dyad}},
\href{https://github.com/yang-mills-group/OpenAI-YM-20261028/blob/3456e5e75cf8d4cd558955a5c7638fba2d68f464/YM/RG08/Jetty/Dosh.lean\#L40}{\texttt{purse}},
\href{https://github.com/yang-mills-group/OpenAI-YM-20261028/blob/3456e5e75cf8d4cd558955a5c7638fba2d68f464/YM/RG08/Jetty/Hoard.lean\#L15}{\texttt{beam}},
\href{https://github.com/yang-mills-group/OpenAI-YM-20261028/blob/3456e5e75cf8d4cd558955a5c7638fba2d68f464/YM/RG08/Jetty/Bank.lean\#L15-L17}{\texttt{Loom}},
\href{https://github.com/yang-mills-group/OpenAI-YM-20261028/blob/3456e5e75cf8d4cd558955a5c7638fba2d68f464/YM/RG08/Skiff/Pocket.lean\#L33}{\texttt{fits}}.

\paragraph{Part B}\label{part-b}

\emph{Setting.}

\begin{itemize}
\tightlist
\item
  The gauge group is as in I.1, with its Haar probability measure.
\item
  A block-spin step (III.1) maps a fine periodic lattice to a coarse one
  with odd block factor \(L\ge3\). The coarse site \(x\) sits at the
  fine site \(\iota x\).
\item
  The kernel is built from patches: boxes of radius \(\rho=(J+3)L\)
  around the points \(\iota x\), where the stage satisfies \(J\ge4\).
\item
  \(k_\beta(U,V)\) is the conditional density of the coarse
  configuration \(V\) given the fine configuration \(U\), with respect
  to the product Haar probability measure.
\item
  The pull-back of a bounded measurable complex function \(F\) of the
  coarse configuration is
  \[(\mathcal BF)(U)=\int k_\beta(U,V)F(V)\ dV.\]
\item
  For \(r\in\mathbb N\), \(\Lambda_r(x)\) is the periodic
  \(\ell^\infty\) box of radius \(r\) about \(x\); periodic images are
  allowed. All lattice radii below are nonnegative integers.
\item
  \(F\) is \emph{local on} \(X\) if \(F(V)\) depends only on the links
  of \(V\) with both endpoints in \(X\).
\item
  \(\mathcal I(X)\) is the class of observables that are gauge
  invariant, measurable, bounded and local on \(X\) (IV.1).
\end{itemize}

\emph{Patch hypothesis (III.1).} The minimising frame and the best
predictor are unique whenever their respective branches are used. The
set of admissible fallback frame tags is nonempty for every cover
configuration.

\textbf{Proposition III.8.4 (one step: the pull-back is local).} Assume
the patch hypothesis, and let \(\beta\in\mathbb R\).

\begin{enumerate}
\def\labelenumi{(\roman{enumi})}
\item
  Let \(F\) be a bounded measurable complex function local on a coarse
  set \(X\), and let \(Y\) be a set of fine sites containing
  \(\Lambda_{10\rho}(\iota x)\) for every \(x\in X\). Then
  \(\mathcal BF\) is local on \(Y\).
\item
  Assume moreover that the density is gauge covariant,
  \(k_\beta(U^{g},V^{g\circ\iota})=k_\beta(U,V)\) for every fine gauge
  transformation \(g\). If \(F\in\mathcal I(\Lambda_R(o))\), then
  \(\mathcal BF\in\mathcal I(\Lambda_{LR+10\rho}(\iota o))\), with the
  same sup bound.
\end{enumerate}

\emph{Status:} theorem. III.1 supplies these kernel properties for its
chosen kernel parameters: the patch hypothesis is
\href{https://github.com/yang-mills-group/OpenAI-YM-20261028/blob/3456e5e75cf8d4cd558955a5c7638fba2d68f464/YM/Next/FOUNDATIONS/Primary.lean\#L70-L71}{\texttt{Next.Foundations.branches}}.
Gauge covariance follows from
\href{https://github.com/yang-mills-group/OpenAI-YM-20261028/blob/3456e5e75cf8d4cd558955a5c7638fba2d68f464/YM/Branch/Covariant.lean\#L56-L67}{\texttt{StepData.covary}}.
For each step of a history, both are fields of the block-control data
built in
\href{https://github.com/yang-mills-group/OpenAI-YM-20261028/blob/3456e5e75cf8d4cd558955a5c7638fba2d68f464/YM/Branch/Covariant.lean\#L70-L72}{\texttt{Calendar.blocks\_go}}.
\emph{Lean:}
\href{https://github.com/yang-mills-group/OpenAI-YM-20261028/blob/3456e5e75cf8d4cd558955a5c7638fba2d68f464/YM/Push/GRID_APPROXIMATION/KernelLocal.lean\#L119-L136}{\texttt{back\_local}},
\href{https://github.com/yang-mills-group/OpenAI-YM-20261028/blob/3456e5e75cf8d4cd558955a5c7638fba2d68f464/YM/Push/GRID_APPROXIMATION/KernelLocal.lean\#L156-L169}{\texttt{back\_cylinder}}.

\textbf{Proposition III.8.5 (exact factorisation).} Assume the patch
hypothesis.

\begin{itemize}
\tightlist
\item
  If \(F,G\) are bounded, measurable and local on disjoint coarse sets
  \(X,Y\), then \(\mathcal B(FG)=\mathcal BF\cdot\mathcal BG\) for every
  \(\beta\).
\item
  In particular
  \(\mathcal B\prod_{i\in l}F_i=\prod_{i\in l}\mathcal BF_i\) when
  \(F_i\in\mathcal I(C_i)\) for distinct cells \(C_i\) of a tiling.
\item
  Along a history with the kernel properties stated below, fix \(N,p\)
  and let \(h\le k\le m\). Let \(l\) be a finite set of level-\(h\)
  sites and \(F_i\in\mathcal I(\Lambda_r(\iota_{h\to k}i))\) at level
  \(k\) for \(i\in l\). If \(\mathrm{room}(h,k,r)\) holds, then
  \(\mathcal B_{k\to m}\prod_{i\in l}F_i=\prod_{i\in l}\mathcal B_{k\to m}F_i\).
\item
  The same holds for the product of the tile copies of one
  \(F\in\mathcal I(\Lambda_r(0))\) (translated, reflected, and
  complex-conjugated on odd tiles), the copies on the right being taken
  at level \(m\).
\end{itemize}

\emph{Status:} theorem. It is an identity, not an estimate, and the fine
windows of \(\mathcal BF\) and \(\mathcal BG\) may overlap. \emph{Lean:}
\href{https://github.com/yang-mills-group/OpenAI-YM-20261028/blob/3456e5e75cf8d4cd558955a5c7638fba2d68f464/YM/Push/GRID_APPROXIMATION/KernelLocal.lean\#L141-L150}{\texttt{back\_factors}},
\href{https://github.com/yang-mills-group/OpenAI-YM-20261028/blob/3456e5e75cf8d4cd558955a5c7638fba2d68f464/YM/Push/GRID_APPROXIMATION/Many.lean\#L21-L40}{\texttt{back\_cells}},
\href{https://github.com/yang-mills-group/OpenAI-YM-20261028/blob/3456e5e75cf8d4cd558955a5c7638fba2d68f464/YM/Push/GRID_APPROXIMATION/Many.lean\#L43-L51}{\texttt{backTo\_cells}},
\href{https://github.com/yang-mills-group/OpenAI-YM-20261028/blob/3456e5e75cf8d4cd558955a5c7638fba2d68f464/YM/Push/GRID_APPROXIMATION/Many.lean\#L68-L76}{\texttt{backTo\_copies}}.

\emph{History.}

\begin{itemize}
\tightlist
\item
  Levels \(m=0,1,2,\dots\) get finer as \(m\) grows. Level \(m\) has
  spacing \(a_m\).
\item
  The step from level \(m+1\) to level \(m\) has block factor \(L_m\)
  (odd, \(\ge3\)), so \(a_{m+1}=a_m/L_m\). Its patch radius is
  \(\rho_m=(J_{m+1}+3)L_m\), where \(J_m\ge4\) and \(J_{m+1}\le J_m+1\).
  The inverse couplings are \(\beta_{N,p,m}\).
\item
  Every step satisfies the patch hypothesis, gauge covariance and
  covariance under lattice translations and signed permutations of the
  axes. These are part of the history data supplied by III.1.
\item
  \(\mathcal B_m\) is the pull-back through this step at
  \(\beta_{N,p,m+1}\). Put
  \(\mathcal B_{k\to m}=\mathcal B_{m-1}\circ\dots\circ\mathcal B_k\)
  (with \(\mathcal B_{k\to k}=\mathrm{id}\)), and let
  \(\mathcal B^{\ast}_{k\to m}\) be the composite of the push-forwards
  of III.1 from level \(m\) to level \(k\).
\item
  For a level-\(h\) site \(i\) and \(m\ge h\), \(\iota_{h\to m}i\) is
  its image at level \(m\).
\item
  The cone radius is \(r_{k\to k}(r)=r\),
  \(r_{k\to m+1}(r)=L_m\ r_{k\to m}(r)+10(J_{m+1}+3)L_m\).
\item
  The physical budget is \(\tau_m=30a_m(J_m+5)\).
\item
  \(\mathrm{room}(h,k,r)\) is the condition \(2(a_kr+\tau_k)\lt a_h\).
\end{itemize}

\textbf{Proposition III.8.6 (pull-back along the history).} For all
\(N,p\) and \(k\le m\):

\begin{enumerate}
\def\labelenumi{(\roman{enumi})}
\item
  \(\mathcal B_{k\to m}\) maps bounded measurable functions to bounded
  measurable functions, with \(\sup|\mathcal B_{k\to m}F|\le\sup|F|\).
  It is linear, fixes constants and commutes with complex conjugation.
\item
  For every bounded measurable complex \(F\) at level \(k\) and every
  integrable complex \(g\) at level \(m\),
  \(\int(\mathcal B^{\ast}_{k\to m}g)F=\int g\ \mathcal B_{k\to m}F\)
  and \(\int\mathcal B^{\ast}_{k\to m}g=\int g\). All integrals here use
  the Haar probability measures at their respective levels; the
  push-forward is integrable.
\item
  Let \(h\le k\) and \(i\) a level-\(h\) site. If
  \(F\in\mathcal I(\Lambda_r(\iota_{h\to k}i))\), then
  \(\mathcal B_{k\to m}F\in\mathcal I(\Lambda_{r_{k\to m}(r)}(\iota_{h\to m}i))\).
\item
  On bounded measurable functions, \(\mathcal B_{k\to m}\) commutes with
  signed permutations of the axes followed by translation by
  \(\iota_{h\to\cdot}i\), for \(h\le k\) and each level-\(h\) site
  \(i\), and with the copy maps of the centred level-\(h\) tiling.
\end{enumerate}

\emph{Status:} theorem. The hypotheses on the kernels are supplied by
the history. \emph{Lean:}
\href{https://github.com/yang-mills-group/OpenAI-YM-20261028/blob/3456e5e75cf8d4cd558955a5c7638fba2d68f464/YM/Push/GRID_APPROXIMATION/Iterate.lean\#L24-L25}{\texttt{backTo}},
\href{https://github.com/yang-mills-group/OpenAI-YM-20261028/blob/3456e5e75cf8d4cd558955a5c7638fba2d68f464/YM/Push/GRID_APPROXIMATION/Iterate.lean\#L27-L29}{\texttt{pushTo}},
\href{https://github.com/yang-mills-group/OpenAI-YM-20261028/blob/3456e5e75cf8d4cd558955a5c7638fba2d68f464/YM/Push/GRID_APPROXIMATION/Iterate.lean\#L49-L55}{\texttt{backTo\_bound}},
\href{https://github.com/yang-mills-group/OpenAI-YM-20261028/blob/3456e5e75cf8d4cd558955a5c7638fba2d68f464/YM/Push/GRID_APPROXIMATION/Iterate.lean\#L44-L47}{\texttt{backTo\_env}},
\href{https://github.com/yang-mills-group/OpenAI-YM-20261028/blob/3456e5e75cf8d4cd558955a5c7638fba2d68f464/YM/Push/GRID_APPROXIMATION/Iterate.lean\#L62-L67}{\texttt{backTo\_add}},
\href{https://github.com/yang-mills-group/OpenAI-YM-20261028/blob/3456e5e75cf8d4cd558955a5c7638fba2d68f464/YM/Push/GRID_APPROXIMATION/Iterate.lean\#L57-L60}{\texttt{backTo\_const}},
\href{https://github.com/yang-mills-group/OpenAI-YM-20261028/blob/3456e5e75cf8d4cd558955a5c7638fba2d68f464/YM/Push/GRID_APPROXIMATION/Iterate.lean\#L69-L73}{\texttt{backTo\_star}},
\href{https://github.com/yang-mills-group/OpenAI-YM-20261028/blob/3456e5e75cf8d4cd558955a5c7638fba2d68f464/YM/Push/GRID_APPROXIMATION/Iterate.lean\#L92-L101}{\texttt{pushTo\_dual}},
\href{https://github.com/yang-mills-group/OpenAI-YM-20261028/blob/3456e5e75cf8d4cd558955a5c7638fba2d68f464/YM/Push/GRID_APPROXIMATION/Iterate.lean\#L81-L89}{\texttt{pushTo\_integral}},
\href{https://github.com/yang-mills-group/OpenAI-YM-20261028/blob/3456e5e75cf8d4cd558955a5c7638fba2d68f464/YM/Push/GRID_APPROXIMATION/Iterate.lean\#L105-L115}{\texttt{backTo\_local}},
\href{https://github.com/yang-mills-group/OpenAI-YM-20261028/blob/3456e5e75cf8d4cd558955a5c7638fba2d68f464/YM/Push/GRID_APPROXIMATION/Iterate.lean\#L117-L126}{\texttt{backTo\_orient}},
\href{https://github.com/yang-mills-group/OpenAI-YM-20261028/blob/3456e5e75cf8d4cd558955a5c7638fba2d68f464/YM/Push/GRID_APPROXIMATION/Iterate.lean\#L128-L147}{\texttt{backTo\_copy}}.

\textbf{Proposition III.8.7 (size of the cone in physical units).} For
the history above and all nonnegative integer levels and radii:

\begin{enumerate}
\def\labelenumi{(\roman{enumi})}
\item
  \(\tau_{m+1}\le\frac23\tau_m\); hence \(\tau_m\le(2/3)^m\tau_0\to0\).
\item
  \(a_m\ r_{k\to m}(r)+\tau_m\le a_kr+\tau_k\) for \(k\le m\). Thus an
  observable of lattice radius \(r\) at level \(k\) has, at every finer
  level, a backward cone of physical radius at most \(a_kr+\tau_k\),
  uniformly in \(m\), \(N\) and \(p\).
\item
  If \(h\le k\) and \(\mathrm{room}(h,k,r)\) holds, then for every
  \(m\ge k\) and every level-\(h\) site \(i\), the box
  \(\Lambda_{r_{k\to m}(r)}(\iota_{h\to m}i)\) lies in the cell of \(i\)
  of the centred level-\(h\) tiling at level \(m\). That cell has
  \(L_h\cdots L_{m-1}\) sites per side, i.e.~physical side \(a_h\).
\item
  For fixed \(h\) and \(r\), \(\mathrm{room}(h,k,r)\) holds for all
  sufficiently large \(k\ge h\).
\end{enumerate}

\emph{Status:} theorem. \emph{Lean:}
\href{https://github.com/yang-mills-group/OpenAI-YM-20261028/blob/3456e5e75cf8d4cd558955a5c7638fba2d68f464/YM/Push/GRID_APPROXIMATION/Geometry.lean\#L107}{\texttt{tail}},
\href{https://github.com/yang-mills-group/OpenAI-YM-20261028/blob/3456e5e75cf8d4cd558955a5c7638fba2d68f464/YM/Push/GRID_APPROXIMATION/Geometry.lean\#L110-L121}{\texttt{tail\_next}},
\href{https://github.com/yang-mills-group/OpenAI-YM-20261028/blob/3456e5e75cf8d4cd558955a5c7638fba2d68f464/YM/Push/GRID_APPROXIMATION/Geometry.lean\#L123-L133}{\texttt{tail\_zero}},
\href{https://github.com/yang-mills-group/OpenAI-YM-20261028/blob/3456e5e75cf8d4cd558955a5c7638fba2d68f464/YM/Push/GRID_APPROXIMATION/Geometry.lean\#L135-L145}{\texttt{scales\_next}},
\href{https://github.com/yang-mills-group/OpenAI-YM-20261028/blob/3456e5e75cf8d4cd558955a5c7638fba2d68f464/YM/Push/GRID_APPROXIMATION/Geometry.lean\#L147-L150}{\texttt{scales\_budget}},
\href{https://github.com/yang-mills-group/OpenAI-YM-20261028/blob/3456e5e75cf8d4cd558955a5c7638fba2d68f464/YM/Push/GRID_APPROXIMATION/Geometry.lean\#L153}{\texttt{room}},
\href{https://github.com/yang-mills-group/OpenAI-YM-20261028/blob/3456e5e75cf8d4cd558955a5c7638fba2d68f464/YM/Push/GRID_APPROXIMATION/Geometry.lean\#L160-L169}{\texttt{room\_fits}},
\href{https://github.com/yang-mills-group/OpenAI-YM-20261028/blob/3456e5e75cf8d4cd558955a5c7638fba2d68f464/YM/Push/GRID_APPROXIMATION/Geometry.lean\#L155-L158}{\texttt{plenty}}.

\subsubsection{2. What it uses}\label{what-it-uses-10}

\begin{itemize}
\tightlist
\item
  \textbf{III.1 (block-spin step, history).}

  \begin{itemize}
  \tightlist
  \item
    \emph{The kernel.} \(k_\beta\)
    (\href{https://github.com/yang-mills-group/OpenAI-YM-20261028/blob/3456e5e75cf8d4cd558955a5c7638fba2d68f464/YM/RG/Block.lean\#L71-L72}{\texttt{StepData.density}})
    is a sum over site tags and bond tags of a
    \href{https://github.com/yang-mills-group/OpenAI-YM-20261028/blob/3456e5e75cf8d4cd558955a5c7638fba2d68f464/YM/RG/Block.lean\#L49-L51}{\texttt{likelihood}}
    times a product over bonds of normalised penalty densities. The
    likelihood is a product over sites of frame masses times a product
    over bonds of predictor masses.
  \item
    \emph{Frames and patches.}
    \href{https://github.com/yang-mills-group/OpenAI-YM-20261028/blob/3456e5e75cf8d4cd558955a5c7638fba2d68f464/YM/RG/Patch.lean\#L153-L155}{\texttt{Patch.frame}},
    \href{https://github.com/yang-mills-group/OpenAI-YM-20261028/blob/3456e5e75cf8d4cd558955a5c7638fba2d68f464/YM/RG/Patch.lean\#L148-L151}{\texttt{Patch.frameMass}},
    \href{https://github.com/yang-mills-group/OpenAI-YM-20261028/blob/3456e5e75cf8d4cd558955a5c7638fba2d68f464/YM/RG/Block.lean\#L39-L42}{\texttt{StepData.connector}},
    and the patch radius
    \href{https://github.com/yang-mills-group/OpenAI-YM-20261028/blob/3456e5e75cf8d4cd558955a5c7638fba2d68f464/YM/RG/Patch.lean\#L51}{\texttt{Patch.radius}}
    \(=(J+3)L\).
  \item
    \emph{Patch hypothesis and kernel properties.} The patch hypothesis
    \href{https://github.com/yang-mills-group/OpenAI-YM-20261028/blob/3456e5e75cf8d4cd558955a5c7638fba2d68f464/YM/RG/Patch.lean\#L182-L185}{\texttt{Patch.Branches}}
    is proved in
    \href{https://github.com/yang-mills-group/OpenAI-YM-20261028/blob/3456e5e75cf8d4cd558955a5c7638fba2d68f464/YM/Next/FOUNDATIONS/Primary.lean\#L70-L71}{\texttt{Next.Foundations.branches}}.
    Also used:
    \href{https://github.com/yang-mills-group/OpenAI-YM-20261028/blob/3456e5e75cf8d4cd558955a5c7638fba2d68f464/YM/RG/Measurable.lean\#L127-L136}{\texttt{frameMass\_sum}};
    \href{https://github.com/yang-mills-group/OpenAI-YM-20261028/blob/3456e5e75cf8d4cd558955a5c7638fba2d68f464/YM/RG/Block.lean\#L203-L208}{\texttt{StepData.isKernel}}
    (positivity, measurability, boundedness, mass one); the push-forward
    \href{https://github.com/yang-mills-group/OpenAI-YM-20261028/blob/3456e5e75cf8d4cd558955a5c7638fba2d68f464/YM/RG/Block.lean\#L77-L79}{\texttt{StepData.push}}.
  \item
    \emph{The history.}
    \href{https://github.com/yang-mills-group/OpenAI-YM-20261028/blob/3456e5e75cf8d4cd558955a5c7638fba2d68f464/YM/RG/History.lean\#L204-L206}{\texttt{History}},
    \href{https://github.com/yang-mills-group/OpenAI-YM-20261028/blob/3456e5e75cf8d4cd558955a5c7638fba2d68f464/YM/RG/History.lean\#L131-L134}{\texttt{Levels}},
    \href{https://github.com/yang-mills-group/OpenAI-YM-20261028/blob/3456e5e75cf8d4cd558955a5c7638fba2d68f464/YM/RG/History.lean\#L97-L118}{\texttt{Calendar.Fits}}
    (\(J_{m+1}\le J_m+1\)),
    \href{https://github.com/yang-mills-group/OpenAI-YM-20261028/blob/3456e5e75cf8d4cd558955a5c7638fba2d68f464/YM/RG/History.lean\#L158-L159}{\texttt{Levels.transition\_mass}}.
    The per-step
    \href{https://github.com/yang-mills-group/OpenAI-YM-20261028/blob/3456e5e75cf8d4cd558955a5c7638fba2d68f464/YM/RG/History.lean\#L120-L128}{\texttt{Calendar.BlockControl}}
    holds the patch hypothesis and the lattice-symmetry and gauge
    covariance of each density; it is constructed in
    \href{https://github.com/yang-mills-group/OpenAI-YM-20261028/blob/3456e5e75cf8d4cd558955a5c7638fba2d68f464/YM/Branch/Covariant.lean\#L70-L72}{\texttt{Calendar.blocks\_go}}.
  \end{itemize}
\item
  \textbf{I.1.}

  \begin{itemize}
  \tightlist
  \item
    Periodic lattices, links, Haar measure, gauge transformations,
    transport along routes, and locality
    \href{https://github.com/yang-mills-group/OpenAI-YM-20261028/blob/3456e5e75cf8d4cd558955a5c7638fba2d68f464/YM/Lattice/Path.lean\#L196-L197}{\texttt{LocalOn}}.
  \item
    The spacing schedule
    \href{https://github.com/yang-mills-group/OpenAI-YM-20261028/blob/3456e5e75cf8d4cd558955a5c7638fba2d68f464/YM/Lattice/Scale.lean\#L28-L34}{\texttt{Schedule}}:
    \(L_m\) odd and \(\ge3\);
    \href{https://github.com/yang-mills-group/OpenAI-YM-20261028/blob/3456e5e75cf8d4cd558955a5c7638fba2d68f464/YM/Lattice/Scale.lean\#L103-L104}{\texttt{a\_step}}
    (\(a_{m+1}=a_m/L_m\));
    \href{https://github.com/yang-mills-group/OpenAI-YM-20261028/blob/3456e5e75cf8d4cd558955a5c7638fba2d68f464/YM/Lattice/Scale.lean\#L108-L109}{\texttt{shrinks}}
    (\(a_m\to0\));
    \href{https://github.com/yang-mills-group/OpenAI-YM-20261028/blob/3456e5e75cf8d4cd558955a5c7638fba2d68f464/YM/Lattice/Scale.lean\#L184-L189}{\texttt{embeds\_scale}}
    (\(a_m\ L_h\cdots L_{m-1}=a_h\)).
  \end{itemize}
\item
  \textbf{IV.1.} Tilings, cells and copies; the class
  \href{https://github.com/yang-mills-group/OpenAI-YM-20261028/blob/3456e5e75cf8d4cd558955a5c7638fba2d68f464/YM/Cells/Cut.lean\#L109-L113}{\texttt{Interior}}
  \(=\mathcal I\); and, in Step B7 only, the chessboard root
  \href{https://github.com/yang-mills-group/OpenAI-YM-20261028/blob/3456e5e75cf8d4cd558955a5c7638fba2d68f464/YM/Cells/Tiling.lean\#L226-L227}{\texttt{rootTest}}
  and the chessboard bound
  \href{https://github.com/yang-mills-group/OpenAI-YM-20261028/blob/3456e5e75cf8d4cd558955a5c7638fba2d68f464/YM/Cells/Embedded.lean\#L25-L38}{\texttt{Tiling.planted}}.
  Its reflection-positivity hypotheses are packaged in
  \href{https://github.com/yang-mills-group/OpenAI-YM-20261028/blob/3456e5e75cf8d4cd558955a5c7638fba2d68f464/YM/Cells/Certificate.lean\#L14-L19}{\texttt{Seam}}
  and proved for \(\beta\ge0\) in
  \href{https://github.com/yang-mills-group/OpenAI-YM-20261028/blob/3456e5e75cf8d4cd558955a5c7638fba2d68f464/YM/Cells/Certificate.lean\#L21-L24}{\texttt{laws}}.
\item
  \textbf{Part A is an abstract estimate.} The consumer, III.14,
  supplies \(V\), \(\Gamma\) and the bounds \(W,\delta\). The weighted
  \(\ell^1\) estimates behind the rooted norm come from the support
  packages listed in §7. The covariance used by III.14 is discussed
  there.
\item
  \textbf{Standard mathematics.}

  \begin{itemize}
  \tightlist
  \item
    \emph{Part A.} The formal Hamilton--Jacobi heat equation, often
    called the Polchinski equation, generates connected diagrams by
    contractions and joins. For finite \(I\) and symmetric inputs, put
    \(W_t(\phi)=\log\mathbb E_{t\Gamma}e^{V(g,\phi+\zeta)}\), with
    expectation interpreted by Wick's rule. Then
    \(\partial_tW=\tfrac12\sum_{u,v}\Gamma(u,v)(\partial_u\partial_vW+\partial_uW\,\partial_vW)\).
    Rooted coefficient bounds give a direct inductive estimate for this
    equation. Background references include Polchinski, Nucl. Phys. B
    231 (1984) 269; Brydges--Kennedy, J. Stat. Phys. 48 (1987) 19;
    Salmhofer, \emph{Renormalization: an introduction}, Springer 1999;
    and Brydges, \emph{Lectures on the renormalisation group}, IAS/Park
    City Math. Ser. 16, 2009.
  \item
    \emph{Part B.} Conditional expectation for a Markov kernel; Fubini
    for product measures and for sums over products of probability
    vectors. The model is Bałaban's block averaging (Comm. Math. Phys.
    98 (1985) 17); the kernel here is the probabilistic one of III.1.
  \item
    \emph{Step B7.} The chessboard estimate
    (Fröhlich--Israel--Lieb--Simon, Comm. Math. Phys. 62 (1978) 1).
  \end{itemize}
\item
  \textbf{Imported constants and assumptions.}

  \begin{itemize}
  \tightlist
  \item
    The patch radius \((J+3)L\) (III.1).
  \item
    \(L_m\ge3\) and \(J_{m+1}\le J_m+1\) (I.1, III.1).
  \item
    The patch hypothesis is assumed in III.8.4--III.8.5; gauge
    covariance is also assumed in III.8.4(ii). III.1 supplies these
    properties for its chosen kernels, together with lattice-symmetry
    covariance for iteration.
  \item
    The constants \(10\) (read radius), \(8\) (route length), \(30\) and
    \(5\) (in \(\tau_m\)) are internal.
  \item
    Part A imports no constants.
  \end{itemize}
\end{itemize}

\subsubsection{3. Strategy}\label{strategy-10}

\textbf{Part A.} Each step of (A1) either contracts two slots of one
coefficient or joins two coefficients by a propagator line. The rooted
norm absorbs the internal sums without changing \(\sigma\): the output
diameter is bounded by the input diameters plus the line length, so the
exponential weights multiply. For a marked coefficient, the height
inequality plays the same role, with the marked factor summed over all
its slots.

Induction on heat order gives a polynomial majorant in \(W\).
Telescoping the inputs gives the corresponding difference bound with the
same \(\delta\). Every primitive vertex has positive coupling order and
at most three times that order in slots; every propagator removes two
slots. This gives \(n+2j\le3e\) and finiteness in heat time at each
coupling order. Positive heat order uses only lower coupling orders and
higher ranks at the same coupling order.

\textbf{Part B.} Given the fine field, the kernel first draws
independent site tags, then independent predictor tags on bonds, and
finally independent coarse links conditional on these tags. Integrating
an observable local on \(X\) removes all unused bonds and site tags. The
remaining factors read the fine field within radius \(10\rho\) of
\(\iota X\). Observables on disjoint coarse site sets use disjoint tags
and bonds, so their pull-backs multiply. Under iteration the radius
grows by \(R\mapsto L_mR+10\rho_m\). In physical units the added radius
is at most \(\tau_m/3\), while \(\tau_{m+1}\le2\tau_m/3\). The total
radius therefore stays within \(a_kr+\tau_k\). When this fits in one
level-\(h\) cell, different cells remain separate at every intermediate
step, allowing the product formula to iterate.

\subsubsection{4. Steps}\label{steps-10}

\paragraph{A. Connected heat
coefficients}\label{a.-connected-heat-coefficients}

\textbf{A1. Operations and the recursion.}

\emph{Definitions.} For \(F\) of rank \(m+2\),
\[(\Delta_\Gamma F)(x)=\mathrm{Sym}\sum_{u,v\in I}\Gamma(u,v)F(u,v,x).\]
For \(F\) of rank \(p+1\) and \(G\) of rank \(q+1\), the join has rank
\(p+q\):
\[(F\star_\Gamma G)(x,y)=\mathrm{Sym}\sum_{u,v\in I}F(u,x)\Gamma(u,v)G(v,y).\]
\(\mathrm{Sym}\) averages over permutations of the slots.

\emph{Combinatorial factors.} On a finite index set, write
\(W=\sum_n\langle W_n,\phi^{\otimes n}\rangle\) with symmetric
coefficient tensors and no slot factorials. For symmetric \(\Gamma\),
the differential operators
\(\tfrac12\sum_{u,v}\Gamma(u,v)\partial_u\partial_vW\) and
\(\tfrac12\sum_{u,v}\Gamma(u,v)\partial_uW\,\partial_vW\) give the
factors \(\binom{n+2}2\) and \((c+1)(n-c+1)/2\). For nonsymmetric
tensors (A1) remains a definition, without this identification.

\emph{Vanishing terms and symmetry.} For eligible \(V\) the pair terms
with \(a=0\) or \(a=e\) vanish, because \(V^{(i)}_{0,\cdot}=0\). Every
\(V^{(j)}\) with \(j\ge1\) is symmetric, since every term is
symmetrised.

\emph{Formal heat equation.} For finite \(I\), eligible \(V\) symmetric
in every rank
(\href{https://github.com/yang-mills-group/OpenAI-YM-20261028/blob/3456e5e75cf8d4cd558955a5c7638fba2d68f464/YM/RG08/Glyph/Towers.lean\#L15}{\texttt{waxed}}),
and symmetric \(\Gamma\), III.14 proves:

\begin{itemize}
\tightlist
\item
  the generating series solves the formal Hamilton--Jacobi equation
  (\href{https://github.com/yang-mills-group/OpenAI-YM-20261028/blob/3456e5e75cf8d4cd558955a5c7638fba2d68f464/YM/RG08/Glyph/Tempo.lean\#L43-L73}{\texttt{Glyph.cloudy}});
\item
  its exponential solves the linear heat equation
  (\href{https://github.com/yang-mills-group/OpenAI-YM-20261028/blob/3456e5e75cf8d4cd558955a5c7638fba2d68f464/YM/RG08/Cloud/Arch.lean\#L46-L49}{\texttt{Cloud.Clockwise.linear}});
\item
  summation at time \(t\) gives \(\mathcal E_t(V)\)
  (\href{https://github.com/yang-mills-group/OpenAI-YM-20261028/blob/3456e5e75cf8d4cd558955a5c7638fba2d68f464/YM/RG08/Glyph/Tempo.lean\#L26-L32}{\texttt{Glyph.breaktime}}).
\end{itemize}

These identities support the formal cumulant interpretation. They do not
assert convergence of a Gaussian integral or of the series in \(g\). The
norm estimates use only the recursion.

\emph{Lean:}
\href{https://github.com/yang-mills-group/OpenAI-YM-20261028/blob/3456e5e75cf8d4cd558955a5c7638fba2d68f464/YM/RG08/Forms/Lap.lean\#L17}{\texttt{lap}},
\href{https://github.com/yang-mills-group/OpenAI-YM-20261028/blob/3456e5e75cf8d4cd558955a5c7638fba2d68f464/YM/RG08/Forms/Weld.lean\#L21}{\texttt{bridge}},
\href{https://github.com/yang-mills-group/OpenAI-YM-20261028/blob/3456e5e75cf8d4cd558955a5c7638fba2d68f464/YM/RG08/Forms/Shuffle.lean\#L38}{\texttt{bleach}},
\href{https://github.com/yang-mills-group/OpenAI-YM-20261028/blob/3456e5e75cf8d4cd558955a5c7638fba2d68f464/YM/RG08/Steam/Gears.lean\#L33}{\texttt{unilateral}},
\href{https://github.com/yang-mills-group/OpenAI-YM-20261028/blob/3456e5e75cf8d4cd558955a5c7638fba2d68f464/YM/RG08/Steam/Gears.lean\#L34-L35}{\texttt{pair}},
\href{https://github.com/yang-mills-group/OpenAI-YM-20261028/blob/3456e5e75cf8d4cd558955a5c7638fba2d68f464/YM/RG08/Steam/Gears.lean\#L37-L38}{\texttt{wheels}},
\href{https://github.com/yang-mills-group/OpenAI-YM-20261028/blob/3456e5e75cf8d4cd558955a5c7638fba2d68f464/YM/RG08/Steam/Gears.lean\#L42-L44}{\texttt{bank}},
\href{https://github.com/yang-mills-group/OpenAI-YM-20261028/blob/3456e5e75cf8d4cd558955a5c7638fba2d68f464/YM/RG08/Steam/Gears.lean\#L46}{\texttt{burn}},
\href{https://github.com/yang-mills-group/OpenAI-YM-20261028/blob/3456e5e75cf8d4cd558955a5c7638fba2d68f464/YM/RG08/Steam/Gears.lean\#L55-L62}{\texttt{blaze}},
\href{https://github.com/yang-mills-group/OpenAI-YM-20261028/blob/3456e5e75cf8d4cd558955a5c7638fba2d68f464/YM/RG08/Steam/Gears.lean\#L65-L75}{\texttt{coal}}.

\textbf{A2. Support (Proposition III.8.2).} Since \(\Delta_\Gamma0=0\)
and \(0\star_\Gamma G=F\star_\Gamma0=0\), a nonzero \(V^{(j+1)}_{e,n}\)
has a nonzero predecessor: either \(V^{(j)}_{e,n+2}\ne0\), or both
\(V^{(i)}_{a,c+1}\) and \(V^{(j-i)}_{e-a,n-c+1}\) are nonzero for some
\(c\le n\), \(i\le j\), \(a\le e\). Strong induction on \(j\) gives
\(n+2j\le3e\):

\begin{itemize}
\tightlist
\item
  at \(j=0\), \(n\le e+2\le3e\) because \(e\ge1\);
\item
  the contraction step uses \((n+2)+2j\le3e\);
\item
  the pair step adds \((c+1)+2i\le3a\) and \((n-c+1)+2(j-i)\le3(e-a)\).
\end{itemize}

Parity and \(n\le e+2\) propagate the same way. Splitting off \(j=0\)
gives \(\mathcal E_t=V+\mathcal R_t\). \emph{Lean:}
\href{https://github.com/yang-mills-group/OpenAI-YM-20261028/blob/3456e5e75cf8d4cd558955a5c7638fba2d68f464/YM/RG08/Steam/Valves.lean\#L34-L48}{\texttt{predecessors}},
\href{https://github.com/yang-mills-group/OpenAI-YM-20261028/blob/3456e5e75cf8d4cd558955a5c7638fba2d68f464/YM/RG08/Steam/Valves.lean\#L50-L64}{\texttt{preserves}},
\href{https://github.com/yang-mills-group/OpenAI-YM-20261028/blob/3456e5e75cf8d4cd558955a5c7638fba2d68f464/YM/RG08/Steam/Valves.lean\#L66-L70}{\texttt{live}},
\href{https://github.com/yang-mills-group/OpenAI-YM-20261028/blob/3456e5e75cf8d4cd558955a5c7638fba2d68f464/YM/RG08/Steam/Valves.lean\#L72}{\texttt{off}},
\href{https://github.com/yang-mills-group/OpenAI-YM-20261028/blob/3456e5e75cf8d4cd558955a5c7638fba2d68f464/YM/RG08/Steam/Valves.lean\#L77-L80}{\texttt{smoulder}},
\href{https://github.com/yang-mills-group/OpenAI-YM-20261028/blob/3456e5e75cf8d4cd558955a5c7638fba2d68f464/YM/RG08/Steam/Valves.lean\#L82-L89}{\texttt{appropriate}}.

\textbf{A3. The two contraction estimates.} Let \(\sigma\ge0\), let the
input rooted norms and \([\Gamma]_\sigma\) be finite, and require
positive output rank. Then
\[\lVert\Delta_\Gamma F\rVert_\sigma\le\lVert F\rVert_\sigma\ \sup|\Gamma|\le\lVert F\rVert_\sigma[\Gamma]_\sigma,\qquad\lVert F\star_\Gamma G\rVert_\sigma\le\lVert F\rVert_\sigma[\Gamma]_\sigma\lVert G\rVert_\sigma.\]

\emph{Symmetrisation.} \(\mathrm{Sym}\) does not increase the rooted
norm: it permutes the slots and preserves diameters.

\emph{Contraction.} Fix a slot of \(x\). Then
\(\mathrm{diam}(x)\le\mathrm{diam}(u,v,x)\) and
\(|\Gamma(u,v)|\le\sup|\Gamma|\), and the sum over \((u,v,x)\) with that
slot fixed is a rooted sum of \(F\).

\emph{Join.} The key inequality is
\[\mathrm{diam}(x,y)\le\mathrm{diam}(u,x)+d(u,v)+\mathrm{diam}(v,y).\]
With a slot of \(x\) fixed, sum in three stages:

\begin{enumerate}
\def\labelenumi{\arabic{enumi}.}
\tightlist
\item
  over \(y\) at fixed \(v\): at most \(\lVert G\rVert\), rooted at
  \(v\);
\item
  over \(v\): at most the row sum of \(\Gamma\);
\item
  over \((u,x)\): at most \(\lVert F\rVert\).
\end{enumerate}

A fixed slot of \(y\) is treated symmetrically, with column sums.

\emph{Differences.} Absolute convergence permits telescoping:
\(\Delta_\Gamma F-\Delta_{\Gamma'}F'=\Delta_\Gamma(F-F')+\Delta_{\Gamma-\Gamma'}F'\),
and the join gives three terms. If the tensor pairs are \((P,\delta)\)-
and \((Q,\delta)\)-controlled and \((\Gamma,\Gamma')\) is
\((B,\delta)\)-controlled, with \(P,Q,B,\delta\ge0\), the contraction
pair is \((2PB,\delta)\)-controlled and the join pair is
\((3PBQ,\delta)\)-controlled. At output rank \(0\) these controlled-pair
conclusions are vacuous; telescoping need not be valid for divergent
scalar sums.

\emph{Lean:}
\href{https://github.com/yang-mills-group/OpenAI-YM-20261028/blob/3456e5e75cf8d4cd558955a5c7638fba2d68f464/YM/RG08/Forms/Ground.lean\#L20-L24}{\texttt{Ground}},
\href{https://github.com/yang-mills-group/OpenAI-YM-20261028/blob/3456e5e75cf8d4cd558955a5c7638fba2d68f464/YM/RG08/Forms/Ground.lean\#L61}{\texttt{Held}}
with
\href{https://github.com/yang-mills-group/OpenAI-YM-20261028/blob/3456e5e75cf8d4cd558955a5c7638fba2d68f464/YM/RG08/Forms/Ground.lean\#L78-L84}{\texttt{ensign}},
\href{https://github.com/yang-mills-group/OpenAI-YM-20261028/blob/3456e5e75cf8d4cd558955a5c7638fba2d68f464/YM/RG08/Forms/Lines.lean\#L18-L22}{\texttt{Wire}},
\href{https://github.com/yang-mills-group/OpenAI-YM-20261028/blob/3456e5e75cf8d4cd558955a5c7638fba2d68f464/YM/RG08/Forms/Lap.lean\#L23-L46}{\texttt{unary}},
\href{https://github.com/yang-mills-group/OpenAI-YM-20261028/blob/3456e5e75cf8d4cd558955a5c7638fba2d68f464/YM/RG08/Forms/Lap.lean\#L58}{\texttt{Held.lap}},
\href{https://github.com/yang-mills-group/OpenAI-YM-20261028/blob/3456e5e75cf8d4cd558955a5c7638fba2d68f464/YM/RG08/Forms/Weld.lean\#L26-L42}{\texttt{arc}},
\href{https://github.com/yang-mills-group/OpenAI-YM-20261028/blob/3456e5e75cf8d4cd558955a5c7638fba2d68f464/YM/RG08/Forms/Weld.lean\#L49-L69}{\texttt{root\_sum}},
\href{https://github.com/yang-mills-group/OpenAI-YM-20261028/blob/3456e5e75cf8d4cd558955a5c7638fba2d68f464/YM/RG08/Forms/Weld.lean\#L71-L99}{\texttt{binary\_EN}},
\href{https://github.com/yang-mills-group/OpenAI-YM-20261028/blob/3456e5e75cf8d4cd558955a5c7638fba2d68f464/YM/RG08/Forms/Weld.lean\#L102}{\texttt{Held.bridge}},
\href{https://github.com/yang-mills-group/OpenAI-YM-20261028/blob/3456e5e75cf8d4cd558955a5c7638fba2d68f464/YM/RG08/Forms/Shuffle.lean\#L74-L79}{\texttt{Held.bleach}},
\href{https://github.com/yang-mills-group/OpenAI-YM-20261028/blob/3456e5e75cf8d4cd558955a5c7638fba2d68f464/YM/RG08/Forms/Lap.lean\#L72-L79}{\texttt{lap\_compare}},
\href{https://github.com/yang-mills-group/OpenAI-YM-20261028/blob/3456e5e75cf8d4cd558955a5c7638fba2d68f464/YM/RG08/Forms/Dial.lean\#L13-L24}{\texttt{bridge\_compare}},
\href{https://github.com/yang-mills-group/OpenAI-YM-20261028/blob/3456e5e75cf8d4cd558955a5c7638fba2d68f464/YM/RG08/Forms/Dial.lean\#L26-L29}{\texttt{Held.lap\_diff}},
\href{https://github.com/yang-mills-group/OpenAI-YM-20261028/blob/3456e5e75cf8d4cd558955a5c7638fba2d68f464/YM/RG08/Forms/Dial.lean\#L31-L35}{\texttt{Held.bridge\_diff}},
\href{https://github.com/yang-mills-group/OpenAI-YM-20261028/blob/3456e5e75cf8d4cd558955a5c7638fba2d68f464/YM/RG08/Forms/Cot.lean\#L49-L54}{\texttt{Cot.lap\textquotesingle{}}},
\href{https://github.com/yang-mills-group/OpenAI-YM-20261028/blob/3456e5e75cf8d4cd558955a5c7638fba2d68f464/YM/RG08/Forms/Cot.lean\#L55-L61}{\texttt{Cot.bridge\textquotesingle{}}}.

\textbf{A4. Proof of Proposition III.8.1.} Strong induction on \(j\),
for all \((e,n)\) at once.

\emph{Base, \(j=0\).} For \(n=0\) the claim is vacuous. For \(e=0\), or
\(n\gt e+2\), the coefficient is zero by eligibility. Otherwise
\((e,n)\in T_0(e,n)\) and the bound is \(W\).

\emph{Step, \(j\to j+1\).} Use (A1).

\begin{itemize}
\tightlist
\item
  The contraction term needs the case \((j;e,n+2)\). Its inputs
  (\(a\lt e\), or \(a=e\) with \(r\ge n+2\)) lie in \(T_{j+1}(e,n)\).
\item
  Only pair terms with \(0\lt a\lt e\) survive. Both factors are of
  order less than \(e\), so they need only inputs of order less than
  \(e\).
\end{itemize}

With \(B=W\) for the propagator, A3 gives the polynomial majorant
\[P_{j+1,e,n}(x)=\frac1{j+1}\Big[\binom{n+2}2\ 2x\ P_{j,e,n+2}(x)+\sum_{c=0}^n\sum_{i=0}^j\sum_{a=0}^e\frac{(c+1)(n-c+1)}2\ 3x\ P_{i,a,c+1}(x)P_{j-i,e-a,n-c+1}(x)\Big].\]
\emph{Lean:}
\href{https://github.com/yang-mills-group/OpenAI-YM-20261028/blob/3456e5e75cf8d4cd558955a5c7638fba2d68f464/YM/RG08/Steam/Routing.lean\#L34-L83}{\texttt{navigation}},
\href{https://github.com/yang-mills-group/OpenAI-YM-20261028/blob/3456e5e75cf8d4cd558955a5c7638fba2d68f464/YM/RG08/Steam/Routing.lean\#L14}{\texttt{taps}},
\href{https://github.com/yang-mills-group/OpenAI-YM-20261028/blob/3456e5e75cf8d4cd558955a5c7638fba2d68f464/YM/RG08/Forms/Powers.lean\#L11}{\texttt{Wage}}.

\textbf{A5. Full correction and change of index set (outside the
main-theorem closure).}

\begin{itemize}
\tightlist
\item
  \emph{Full correction.} For fixed \(e,n\), there is a polynomial
  majorant \(P\) such that, under III.8.1's hypotheses on the input
  range \(T_1(e,n)\), the pair
  \((\mathcal R_t(V)_{e,n},\mathcal R_t(V')_{e,n})\) is
  \((P(W),\delta)\)-controlled whenever \(|t|\le W\). Sum A4 over
  \(1\le j\le3e\), using \(|t|^j\le W^j\).
\item
  \emph{Transport.} Let \(n\ge1\), let \((I,d)\) and \((J,D)\) be
  pseudometric spaces, and let \(\sigma,\sigma'\ge0\). For each slot
  \(i\), choose a kernel \(K_i:I\times J\to\mathbb R\), a nonnegative
  weight \(E_i\), and \(b_i\ge0\) bounding both weighted row and column
  sums of \(|K_i|e^{E_i}\). Define
  \((KF)(y)=\sum_{x\in I^n}F(x)\prod_iK_i(x_i,y_i)\). Suppose, for
  distinct slots \(i,j\), that
  \(\sigma'D(y_i,y_j)\le\sigma d(x_i,x_j)+E_i(x_i,y_i)+E_j(x_j,y_j)\)
  whenever the two kernel factors are nonzero. Then
  \(\lVert KF\rVert_{D,\sigma'}\le\lVert F\rVert_{d,\sigma}\prod_i b_i\)
  for every \(F\) with finite rooted norm. Fix an output slot: use its
  column bound to sum its input position, and use row bounds for the
  remaining slots. The geometric inequality bounds the output weight by
  the input weight times all slot weights.
\item
  \emph{Differences.} Suppose \((F,F')\) is \((C,\delta)\)-controlled,
  with \(C,\delta,B\ge0\). If every \(K_i,K'_i\) has weighted row and
  column bounds \(B\), every \(K_i-K'_i\) has bounds \(B\delta\), and
  the same geometric inequality holds for every combination of their
  supports, then \((KF,K'F')\) is \(((n+1)CB^n,\delta)\)-controlled.
  Telescope once in \(F\) and once in each of the \(n\) kernels.
\item
  \emph{Combination.} At \(t=1\), take \(K_i=\mu\) in every slot and
  \(B=W\). For eligible \(V,V'\) and the input bounds above, the
  transported correction has a polynomial bound in \(W\). A bound for
  the transported full potential additionally requires
  \((V_{e,n},V'_{e,n})\) to be \((W,\delta)\)-controlled.
\item
  \emph{Rank \(0\).} The unmarked rooted bounds leave the scalar term
  uncontrolled. On an infinite lattice its sums may diverge.
\end{itemize}

\emph{Lean:}
\href{https://github.com/yang-mills-group/OpenAI-YM-20261028/blob/3456e5e75cf8d4cd558955a5c7638fba2d68f464/YM/RG08/Steam/Routing.lean\#L88-L101}{\texttt{postage}},
\href{https://github.com/yang-mills-group/OpenAI-YM-20261028/blob/3456e5e75cf8d4cd558955a5c7638fba2d68f464/YM/RG08/Forms/Pontoon.lean\#L56}{\texttt{raft}},
\href{https://github.com/yang-mills-group/OpenAI-YM-20261028/blob/3456e5e75cf8d4cd558955a5c7638fba2d68f464/YM/RG08/Forms/Pontoon.lean\#L19-L20}{\texttt{lighthouse}},
\href{https://github.com/yang-mills-group/OpenAI-YM-20261028/blob/3456e5e75cf8d4cd558955a5c7638fba2d68f464/YM/RG08/Forms/Pontoon.lean\#L22-L23}{\texttt{street}},
\href{https://github.com/yang-mills-group/OpenAI-YM-20261028/blob/3456e5e75cf8d4cd558955a5c7638fba2d68f464/YM/RG08/Forms/PontChange.lean\#L38-L43}{\texttt{Alley}},
\href{https://github.com/yang-mills-group/OpenAI-YM-20261028/blob/3456e5e75cf8d4cd558955a5c7638fba2d68f464/YM/RG08/Forms/PontChange.lean\#L72-L89}{\texttt{Cot.raft}},
\href{https://github.com/yang-mills-group/OpenAI-YM-20261028/blob/3456e5e75cf8d4cd558955a5c7638fba2d68f464/YM/RG08/Steam/Ballot.lean\#L34-L48}{\texttt{paving}},
\href{https://github.com/yang-mills-group/OpenAI-YM-20261028/blob/3456e5e75cf8d4cd558955a5c7638fba2d68f464/YM/RG08/Steam/Ballot.lean\#L25}{\texttt{strict}},
\href{https://github.com/yang-mills-group/OpenAI-YM-20261028/blob/3456e5e75cf8d4cd558955a5c7638fba2d68f464/YM/RG08/Steam/Ballot.lean\#L26}{\texttt{allG}}.

\textbf{A6. Marked expansion (Proposition III.8.3).}

\emph{Operations.} For \(M\) marked and \(G\) unmarked, the contraction
\(\Delta_\Gamma M\) is as in A1, and the join is
\((M\star_\Gamma G)(x,y)=\mathrm{Sym}\sum_{u,v}\Gamma(u,v)G(v,y)M(u,x)\).

\emph{Estimates.}
\[\lVert\Delta_\Gamma M\rVert_h\le\sup|\Gamma|\ \lVert M\rVert_h,\qquad\lVert M\star_\Gamma G\rVert_h\le\lVert M\rVert_h[\Gamma]_\sigma\lVert G\rVert_\sigma.\]
The first holds because \(h(x)\le h(u,v,x)\). The second follows from
\(h(x,y)\le h(u,x)+d(u,v)+\mathrm{diam}_d(v,y)\), which uses that \(h\)
is 1-Lipschitz. Sum the marked factor over all its slots, the propagator
by its row bound, and the unmarked factor by its rooted bound at \(v\).
With finite input bounds, the sum of the dual norms is finite.
Evaluation at each test function is therefore absolutely summable and
preserves linearity, including at output rank \(0\). The difference
constants are \(2\) and \(3\), as in A3.

\emph{Induction for (i).} Induct on \(j\) as in A4.

\begin{itemize}
\tightlist
\item
  The contraction term uses the case \((j;e,n+2)\).
\item
  Pair terms with \(a=e\) vanish, because the unmarked factor then has
  order \(0\).
\item
  For \(a\lt e\), the marked factor \(M^{(i)}_{a,c+1}\) is bounded by
  the induction hypothesis, which needs marked inputs of order at most
  \(a\lt e\). The unmarked factor \(V^{(j-i)}_{e-a,n-c+1}\) is bounded
  by III.8.1, whose inputs lie in
  \(\lbrace1\le a'\le e,\ 1\le r\le a'+2\rbrace\).
\end{itemize}

This gives (i) with output \(P(W)U\).

\emph{(ii) and (iii).} For (ii), a nonzero contraction has a marked
predecessor of rank \(n+2\). A nonzero join has marked and unmarked
predecessors whose ranks sum to \(n+2\) and whose coupling orders sum to
\(e\). Add the marked induction bounds to the unmarked bounds of III.8.2
to obtain \(n+2j\le l+3e\) and \(n\le l+e\). For (iii), sum (i) over
\(1\le j\le l+3e\), using \(|t|^j\le W^j\); no bound for the initial
rank-\(n\), order-\(e\) marked coefficient is needed.

\emph{Consistency.} Evaluating a marked tensor at a fixed test function
turns its contraction and join into those of A1.

\emph{Lean:}
\href{https://github.com/yang-mills-group/OpenAI-YM-20261028/blob/3456e5e75cf8d4cd558955a5c7638fba2d68f464/YM/RG08/Jetty/Bank.lean\#L15-L17}{\texttt{Loom}},
\href{https://github.com/yang-mills-group/OpenAI-YM-20261028/blob/3456e5e75cf8d4cd558955a5c7638fba2d68f464/YM/RG08/Jetty/Bank.lean\#L36}{\texttt{Sings}},
\href{https://github.com/yang-mills-group/OpenAI-YM-20261028/blob/3456e5e75cf8d4cd558955a5c7638fba2d68f464/YM/RG08/Jetty/Dosh.lean\#L40}{\texttt{purse}},
\href{https://github.com/yang-mills-group/OpenAI-YM-20261028/blob/3456e5e75cf8d4cd558955a5c7638fba2d68f464/YM/RG08/Jetty/Hoard.lean\#L15}{\texttt{beam}},
\href{https://github.com/yang-mills-group/OpenAI-YM-20261028/blob/3456e5e75cf8d4cd558955a5c7638fba2d68f464/YM/RG08/Jetty/Bank.lean\#L37}{\texttt{Lyrics}},
\href{https://github.com/yang-mills-group/OpenAI-YM-20261028/blob/3456e5e75cf8d4cd558955a5c7638fba2d68f464/YM/RG08/Skiff/Pair.lean\#L11-L17}{\texttt{Dyad}},
\href{https://github.com/yang-mills-group/OpenAI-YM-20261028/blob/3456e5e75cf8d4cd558955a5c7638fba2d68f464/YM/RG08/Skiff/Pocket.lean\#L33}{\texttt{fits}},
\href{https://github.com/yang-mills-group/OpenAI-YM-20261028/blob/3456e5e75cf8d4cd558955a5c7638fba2d68f464/YM/RG08/Skiff/Pocket.lean\#L28}{\texttt{pocket}},
\href{https://github.com/yang-mills-group/OpenAI-YM-20261028/blob/3456e5e75cf8d4cd558955a5c7638fba2d68f464/YM/RG08/Skiff/Pocket.lean\#L75}{\texttt{cloak}},
\href{https://github.com/yang-mills-group/OpenAI-YM-20261028/blob/3456e5e75cf8d4cd558955a5c7638fba2d68f464/YM/RG08/Skiff/Pair.lean\#L39-L48}{\texttt{Dyad.pocket}},
\href{https://github.com/yang-mills-group/OpenAI-YM-20261028/blob/3456e5e75cf8d4cd558955a5c7638fba2d68f464/YM/RG08/Skiff/Pair.lean\#L52-L63}{\texttt{Dyad.cloak}},
\href{https://github.com/yang-mills-group/OpenAI-YM-20261028/blob/3456e5e75cf8d4cd558955a5c7638fba2d68f464/YM/RG08/Skiff/Spoke.lean\#L23}{\texttt{shim}},
\href{https://github.com/yang-mills-group/OpenAI-YM-20261028/blob/3456e5e75cf8d4cd558955a5c7638fba2d68f464/YM/RG08/Skiff/Spoke.lean\#L24}{\texttt{shift}},
\href{https://github.com/yang-mills-group/OpenAI-YM-20261028/blob/3456e5e75cf8d4cd558955a5c7638fba2d68f464/YM/RG08/Skiff/Spoke.lean\#L25-L26}{\texttt{rungs}},
\href{https://github.com/yang-mills-group/OpenAI-YM-20261028/blob/3456e5e75cf8d4cd558955a5c7638fba2d68f464/YM/RG08/Skiff/Spoke.lean\#L31}{\texttt{skate}},
\href{https://github.com/yang-mills-group/OpenAI-YM-20261028/blob/3456e5e75cf8d4cd558955a5c7638fba2d68f464/YM/RG08/Skiff/Spoke.lean\#L38-L40}{\texttt{climbs}},
\href{https://github.com/yang-mills-group/OpenAI-YM-20261028/blob/3456e5e75cf8d4cd558955a5c7638fba2d68f464/YM/RG08/Skiff/Spoke.lean\#L85}{\texttt{across}},
\href{https://github.com/yang-mills-group/OpenAI-YM-20261028/blob/3456e5e75cf8d4cd558955a5c7638fba2d68f464/YM/RG08/Skiff/Spoke.lean\#L92-L127}{\texttt{journey}},
\href{https://github.com/yang-mills-group/OpenAI-YM-20261028/blob/3456e5e75cf8d4cd558955a5c7638fba2d68f464/YM/RG08/Skiff/Spoke.lean\#L52}{\texttt{Lit}},
\href{https://github.com/yang-mills-group/OpenAI-YM-20261028/blob/3456e5e75cf8d4cd558955a5c7638fba2d68f464/YM/RG08/Skiff/Spoke.lean\#L61-L68}{\texttt{full}},
\href{https://github.com/yang-mills-group/OpenAI-YM-20261028/blob/3456e5e75cf8d4cd558955a5c7638fba2d68f464/YM/RG08/Skiff/Spoke.lean\#L71}{\texttt{shaker}},
\href{https://github.com/yang-mills-group/OpenAI-YM-20261028/blob/3456e5e75cf8d4cd558955a5c7638fba2d68f464/YM/RG08/Skiff/Spoke.lean\#L72-L73}{\texttt{cruising}},
\href{https://github.com/yang-mills-group/OpenAI-YM-20261028/blob/3456e5e75cf8d4cd558955a5c7638fba2d68f464/YM/RG08/Skiff/Shed.lean\#L25-L36}{\texttt{registry}},
\href{https://github.com/yang-mills-group/OpenAI-YM-20261028/blob/3456e5e75cf8d4cd558955a5c7638fba2d68f464/YM/RG08/Skiff/Shed.lean\#L20-L21}{\texttt{clocked}},
\href{https://github.com/yang-mills-group/OpenAI-YM-20261028/blob/3456e5e75cf8d4cd558955a5c7638fba2d68f464/YM/RG08/Skiff/Shed.lean\#L22-L23}{\texttt{reel\_stamp}}.

\paragraph{B. Backward dependence
cones}\label{b.-backward-dependence-cones}

\textbf{B1. Product structure of the kernel.} Summing the bond tags out
of the definition in III.1 gives
\[k_\beta(U,V)=\sum_t\prod_xq_x(U,t_x)\prod_bp_b(U,t;V_b).\]

\begin{itemize}
\tightlist
\item
  \(t=(t_x)\) runs over the site (frame) tags, and \(q_x(U,\cdot)\) is
  the frame mass at \(x\).
\item
  \(p_b(U,t;\cdot)\) is a mixture, over predictor tags \(r\) with
  weights summing to one, of normalised penalty densities centred at the
  predictions \(Q_r(Y_b)\), which are built from the connector \(Y_b\)
  of the bond:
  \[p_b(U,t;v)=\sum_rw_r(Y_b)\ Z_\beta^{-1}\exp\big(-\beta\kappa\ \mathrm{price}(v\ Q_r(Y_b)^{-1})\big).\]
  Here \(\kappa\gt0\) is the kernel penalty of III.1, called \(b\) in
  the main theorem.
\end{itemize}

Each \(p_b\) has Haar integral \(1\). The identity \(\sum_rq_x(U,r)=1\)
uses the patch hypothesis (the fallback set is nonempty). Hence
\[\mathcal BF(U)=\sum_t\Big(\prod_xq_x(U,t_x)\Big)S_F(U,t),\qquad S_F(U,t)=\int\prod_bp_b(U,t;V_b)F(V)\ dV.\]
Given \(U\), the site tags are independent; given \(U\) and the tags,
the coarse links are independent. Adjacent bonds share a tag, so the
links are not independent given \(U\) alone. \emph{Lean:}
\href{https://github.com/yang-mills-group/OpenAI-YM-20261028/blob/3456e5e75cf8d4cd558955a5c7638fba2d68f464/YM/Push/GRID_APPROXIMATION/KernelLocal.lean\#L39-L47}{\texttt{density\_parts}},
\href{https://github.com/yang-mills-group/OpenAI-YM-20261028/blob/3456e5e75cf8d4cd558955a5c7638fba2d68f464/YM/Push/GRID_APPROXIMATION/KernelLocal.lean\#L28}{\texttt{noiseLink}},
\href{https://github.com/yang-mills-group/OpenAI-YM-20261028/blob/3456e5e75cf8d4cd558955a5c7638fba2d68f464/YM/Push/GRID_APPROXIMATION/KernelLocal.lean\#L30-L32}{\texttt{perLink}},
\href{https://github.com/yang-mills-group/OpenAI-YM-20261028/blob/3456e5e75cf8d4cd558955a5c7638fba2d68f464/YM/Push/GRID_APPROXIMATION/KernelLocal.lean\#L34-L35}{\texttt{perSite}},
\href{https://github.com/yang-mills-group/OpenAI-YM-20261028/blob/3456e5e75cf8d4cd558955a5c7638fba2d68f464/YM/Push/GRID_APPROXIMATION/KernelLocal.lean\#L51-L55}{\texttt{noiseLink\_mass}},
\href{https://github.com/yang-mills-group/OpenAI-YM-20261028/blob/3456e5e75cf8d4cd558955a5c7638fba2d68f464/YM/Push/GRID_APPROXIMATION/KernelLocal.lean\#L62-L67}{\texttt{perLink\_mass}},
\href{https://github.com/yang-mills-group/OpenAI-YM-20261028/blob/3456e5e75cf8d4cd558955a5c7638fba2d68f464/YM/Push/GRID_APPROXIMATION/KernelLocal.lean\#L83-L85}{\texttt{site\_unit}},
\href{https://github.com/yang-mills-group/OpenAI-YM-20261028/blob/3456e5e75cf8d4cd558955a5c7638fba2d68f464/YM/Push/GRID_APPROXIMATION/KernelLocal.lean\#L87-L97}{\texttt{parts\_back}}.

\textbf{B2. Read sets.}

\begin{enumerate}
\def\labelenumi{(\alph{enumi})}
\tightlist
\item
  \emph{Frames.} The frame and the frame mass at \(x\) depend only on
  \(U\) on links inside \(\Lambda_{10\rho}(\iota x)\). There are two
  branches:
\end{enumerate}

\begin{itemize}
\tightlist
\item
  if a small-field test on transports in the patch passes, the frame is
  the minimiser, normalised at the centre, of an averaged price of
  transports between patch points, and the mass is uniform over all
  tags;
\item
  otherwise the frame is the axial gauge built from the tag, normalised
  at the centre, and the mass is uniform over admissible tags (those
  whose axial cost is bounded by a multiple of the plaquette cost in the
  patch).
\end{itemize}

All transports run along coordinate routes between points of
\(\Lambda_\rho\), of length at most \(8\rho\), so they stay in
\(\Lambda_{9\rho}\). The minimiser is selected by a global choice. Its
uniqueness (patch hypothesis) is what forces two configurations that
agree on the window to select the same frame.

\begin{enumerate}
\def\labelenumi{(\alph{enumi})}
\setcounter{enumi}{1}
\tightlist
\item
  \emph{Connectors.} The connector of the bond from \(x\) to \(x+e_\mu\)
  depends only on the tags at both ends and on \(U\) in
  \(\Lambda_{10\rho}(\iota x)\cup\Lambda_{10\rho}(\iota(x+e_\mu))\). It
  is built from the two frames and from straight transports of length
  \(L\le\rho\).
\end{enumerate}

\emph{Lean:}
\href{https://github.com/yang-mills-group/OpenAI-YM-20261028/blob/3456e5e75cf8d4cd558955a5c7638fba2d68f464/YM/Push/GRID_APPROXIMATION/Cone.lean\#L115-L155}{\texttt{patch\_dep}},
\href{https://github.com/yang-mills-group/OpenAI-YM-20261028/blob/3456e5e75cf8d4cd558955a5c7638fba2d68f464/YM/Push/GRID_APPROXIMATION/Cone.lean\#L85-L99}{\texttt{route\_bound}},
\href{https://github.com/yang-mills-group/OpenAI-YM-20261028/blob/3456e5e75cf8d4cd558955a5c7638fba2d68f464/YM/Push/GRID_APPROXIMATION/Cone.lean\#L67-L83}{\texttt{visits\_bound}},
\href{https://github.com/yang-mills-group/OpenAI-YM-20261028/blob/3456e5e75cf8d4cd558955a5c7638fba2d68f464/YM/Push/GRID_APPROXIMATION/Cone.lean\#L160-L164}{\texttt{site\_dep}},
\href{https://github.com/yang-mills-group/OpenAI-YM-20261028/blob/3456e5e75cf8d4cd558955a5c7638fba2d68f464/YM/Push/GRID_APPROXIMATION/Cone.lean\#L166-L184}{\texttt{connector\_dep}},
\href{https://github.com/yang-mills-group/OpenAI-YM-20261028/blob/3456e5e75cf8d4cd558955a5c7638fba2d68f464/YM/Push/GRID_APPROXIMATION/Cone.lean\#L65}{\texttt{reach}}.

\textbf{B3. Locality (Proposition III.8.4).} Two elementary facts:

\begin{itemize}
\tightlist
\item
  (P1) if a function reads only the coordinates in a set \(s\), its
  integral against a product of probability densities does not change
  when the densities outside \(s\) are changed;
\item
  (P2) the same holds for sums over products of probability vectors.
\end{itemize}

\emph{Proof of (i).} If \(F\) is local on \(X\), it reads only bonds
with both ends in \(X\). By (P1) and B2(b), \(S_F(U,t)\) depends only on
\(t|_X\) and on \(U\) in the windows around \(\iota X\). By (P2), the
tags off \(X\) sum out. By B2(a), \(q_x\) for \(x\in X\) reads \(U\)
only in \(\Lambda_{10\rho}(\iota x)\).

\emph{Proof of (ii).} If \(y\in\Lambda_R(o)\) then
\(\iota y\in\Lambda_{LR}(\iota o)\), and \(\Lambda_{10\rho}\) of a point
of \(\Lambda_{LR}(\iota o)\) lies in \(\Lambda_{LR+10\rho}(\iota o)\).
Gauge invariance follows from the covariance of \(k_\beta\) and the
invariance of Haar measure. The sup bound and measurability hold because
\(k_\beta\) is a probability kernel.

\emph{Lean:}
\href{https://github.com/yang-mills-group/OpenAI-YM-20261028/blob/3456e5e75cf8d4cd558955a5c7638fba2d68f464/YM/Push/GRID_APPROXIMATION/Product.lean\#L19}{\texttt{Reads}},
\href{https://github.com/yang-mills-group/OpenAI-YM-20261028/blob/3456e5e75cf8d4cd558955a5c7638fba2d68f464/YM/Push/GRID_APPROXIMATION/Product.lean\#L88-L96}{\texttt{pi\_ignores}},
\href{https://github.com/yang-mills-group/OpenAI-YM-20261028/blob/3456e5e75cf8d4cd558955a5c7638fba2d68f464/YM/Push/GRID_APPROXIMATION/Product.lean\#L137-L143}{\texttt{sum\_ignores}},
\href{https://github.com/yang-mills-group/OpenAI-YM-20261028/blob/3456e5e75cf8d4cd558955a5c7638fba2d68f464/YM/Push/GRID_APPROXIMATION/KernelLocal.lean\#L100}{\texttt{actualEdges}},
\href{https://github.com/yang-mills-group/OpenAI-YM-20261028/blob/3456e5e75cf8d4cd558955a5c7638fba2d68f464/YM/Push/GRID_APPROXIMATION/KernelLocal.lean\#L107-L114}{\texttt{state\_reads}},
\href{https://github.com/yang-mills-group/OpenAI-YM-20261028/blob/3456e5e75cf8d4cd558955a5c7638fba2d68f464/YM/Push/GRID_APPROXIMATION/KernelLocal.lean\#L119-L136}{\texttt{back\_local}},
\href{https://github.com/yang-mills-group/OpenAI-YM-20261028/blob/3456e5e75cf8d4cd558955a5c7638fba2d68f464/YM/Push/GRID_APPROXIMATION/Cone.lean\#L52-L60}{\texttt{window\_embed}},
\href{https://github.com/yang-mills-group/OpenAI-YM-20261028/blob/3456e5e75cf8d4cd558955a5c7638fba2d68f464/YM/Push/GRID_APPROXIMATION/Cone.lean\#L33-L42}{\texttt{window\_add}},
\href{https://github.com/yang-mills-group/OpenAI-YM-20261028/blob/3456e5e75cf8d4cd558955a5c7638fba2d68f464/YM/Push/GRID_APPROXIMATION/KernelBasic.lean\#L76-L84}{\texttt{back\_natural}},
\href{https://github.com/yang-mills-group/OpenAI-YM-20261028/blob/3456e5e75cf8d4cd558955a5c7638fba2d68f464/YM/Push/GRID_APPROXIMATION/KernelBasic.lean\#L41-L51}{\texttt{back\_bound}},
\href{https://github.com/yang-mills-group/OpenAI-YM-20261028/blob/3456e5e75cf8d4cd558955a5c7638fba2d68f464/YM/Push/GRID_APPROXIMATION/KernelBasic.lean\#L53-L55}{\texttt{back\_env}},
\href{https://github.com/yang-mills-group/OpenAI-YM-20261028/blob/3456e5e75cf8d4cd558955a5c7638fba2d68f464/YM/Push/GRID_APPROXIMATION/KernelLocal.lean\#L156-L169}{\texttt{back\_cylinder}}.

\textbf{B4. Factorisation (Proposition III.8.5).}

\begin{itemize}
\tightlist
\item
  If \(X\cap Y=\emptyset\), the bond sets read by \(F\) and \(G\) are
  disjoint, so \(S_{FG}(U,t)=S_F(U,t)S_G(U,t)\) by Fubini.
\item
  \(S_F\) depends on \(t|_X\) and \(S_G\) on \(t|_Y\), and the
  \(q_x(U,\cdot)\) are probability vectors, so the sum over tags
  factorises.
\item
  Induction over the cells of a tiling gives the finite-product version.
\item
  Along the history, induct on \(m\). By III.8.7(iii), each
  \(\mathcal B_{k\to m}F_i\) lies in \(\mathcal I\) of its own
  level-\(h\) cell, so the one-step product formula applies at every
  step.
\item
  The version for copies follows with III.8.6(iv).
\end{itemize}

\emph{Lean:}
\href{https://github.com/yang-mills-group/OpenAI-YM-20261028/blob/3456e5e75cf8d4cd558955a5c7638fba2d68f464/YM/Push/GRID_APPROXIMATION/Product.lean\#L98-L109}{\texttt{pi\_independent}},
\href{https://github.com/yang-mills-group/OpenAI-YM-20261028/blob/3456e5e75cf8d4cd558955a5c7638fba2d68f464/YM/Push/GRID_APPROXIMATION/Product.lean\#L145-L156}{\texttt{sum\_independent}},
\href{https://github.com/yang-mills-group/OpenAI-YM-20261028/blob/3456e5e75cf8d4cd558955a5c7638fba2d68f464/YM/Push/GRID_APPROXIMATION/KernelLocal.lean\#L141-L150}{\texttt{back\_factors}},
\href{https://github.com/yang-mills-group/OpenAI-YM-20261028/blob/3456e5e75cf8d4cd558955a5c7638fba2d68f464/YM/Push/GRID_APPROXIMATION/Many.lean\#L21-L40}{\texttt{back\_cells}},
\href{https://github.com/yang-mills-group/OpenAI-YM-20261028/blob/3456e5e75cf8d4cd558955a5c7638fba2d68f464/YM/Push/GRID_APPROXIMATION/Many.lean\#L43-L51}{\texttt{backTo\_cells}},
\href{https://github.com/yang-mills-group/OpenAI-YM-20261028/blob/3456e5e75cf8d4cd558955a5c7638fba2d68f464/YM/Push/GRID_APPROXIMATION/Many.lean\#L57-L66}{\texttt{copy\_window}},
\href{https://github.com/yang-mills-group/OpenAI-YM-20261028/blob/3456e5e75cf8d4cd558955a5c7638fba2d68f464/YM/Push/GRID_APPROXIMATION/Many.lean\#L68-L76}{\texttt{backTo\_copies}}.

\textbf{B5. Iteration (Proposition III.8.6).} Induct on \(m-k\), with
\(\mathcal B_{k\to m+1}=\mathcal B_m\circ\mathcal B_{k\to m}\).

\begin{itemize}
\tightlist
\item
  \begin{enumerate}
  \def\labelenumi{(\roman{enumi})}
  \tightlist
  \item
    and (ii) hold for one step because \(k_\beta\) is a real probability
    kernel; (ii) is Fubini.
  \end{enumerate}
\item
  (iii) is III.8.4(ii) at each step. The kernel hypotheses come from the
  block-control data, and \(\iota(\iota_{h\to m}i)=\iota_{h\to m+1}i\).
\item
  (iv) uses the lattice-symmetry covariance of each density.
\end{itemize}

\emph{Lean:}
\href{https://github.com/yang-mills-group/OpenAI-YM-20261028/blob/3456e5e75cf8d4cd558955a5c7638fba2d68f464/YM/Push/GRID_APPROXIMATION/Iterate.lean\#L32-L34}{\texttt{backTo\_next}},
\href{https://github.com/yang-mills-group/OpenAI-YM-20261028/blob/3456e5e75cf8d4cd558955a5c7638fba2d68f464/YM/Push/GRID_APPROXIMATION/Iterate.lean\#L44-L47}{\texttt{backTo\_env}},
\href{https://github.com/yang-mills-group/OpenAI-YM-20261028/blob/3456e5e75cf8d4cd558955a5c7638fba2d68f464/YM/Push/GRID_APPROXIMATION/Iterate.lean\#L49-L55}{\texttt{backTo\_bound}},
\href{https://github.com/yang-mills-group/OpenAI-YM-20261028/blob/3456e5e75cf8d4cd558955a5c7638fba2d68f464/YM/Push/GRID_APPROXIMATION/Iterate.lean\#L62-L67}{\texttt{backTo\_add}},
\href{https://github.com/yang-mills-group/OpenAI-YM-20261028/blob/3456e5e75cf8d4cd558955a5c7638fba2d68f464/YM/Push/GRID_APPROXIMATION/Iterate.lean\#L57-L60}{\texttt{backTo\_const}},
\href{https://github.com/yang-mills-group/OpenAI-YM-20261028/blob/3456e5e75cf8d4cd558955a5c7638fba2d68f464/YM/Push/GRID_APPROXIMATION/Iterate.lean\#L69-L73}{\texttt{backTo\_star}},
\href{https://github.com/yang-mills-group/OpenAI-YM-20261028/blob/3456e5e75cf8d4cd558955a5c7638fba2d68f464/YM/Push/GRID_APPROXIMATION/KernelBasic.lean\#L88-L97}{\texttt{back\_dual}},
\href{https://github.com/yang-mills-group/OpenAI-YM-20261028/blob/3456e5e75cf8d4cd558955a5c7638fba2d68f464/YM/Push/GRID_APPROXIMATION/Iterate.lean\#L81-L89}{\texttt{pushTo\_integral}},
\href{https://github.com/yang-mills-group/OpenAI-YM-20261028/blob/3456e5e75cf8d4cd558955a5c7638fba2d68f464/YM/Push/GRID_APPROXIMATION/Iterate.lean\#L92-L101}{\texttt{pushTo\_dual}},
\href{https://github.com/yang-mills-group/OpenAI-YM-20261028/blob/3456e5e75cf8d4cd558955a5c7638fba2d68f464/YM/Push/GRID_APPROXIMATION/Geometry.lean\#L31-L45}{\texttt{anchor\_next}},
\href{https://github.com/yang-mills-group/OpenAI-YM-20261028/blob/3456e5e75cf8d4cd558955a5c7638fba2d68f464/YM/Push/GRID_APPROXIMATION/Iterate.lean\#L105-L115}{\texttt{backTo\_local}},
\href{https://github.com/yang-mills-group/OpenAI-YM-20261028/blob/3456e5e75cf8d4cd558955a5c7638fba2d68f464/YM/Push/GRID_APPROXIMATION/Iterate.lean\#L117-L126}{\texttt{backTo\_orient}},
\href{https://github.com/yang-mills-group/OpenAI-YM-20261028/blob/3456e5e75cf8d4cd558955a5c7638fba2d68f464/YM/Push/GRID_APPROXIMATION/Iterate.lean\#L128-L147}{\texttt{backTo\_copy}}.

\textbf{B6. Physical size of the cone (Proposition III.8.7).}

\begin{itemize}
\tightlist
\item
  \begin{enumerate}
  \def\labelenumi{(\roman{enumi})}
  \tightlist
  \item
    \(\tau_{m+1}\le30\frac{a_m}3(J_m+6)\le\frac23\tau_m\).
  \end{enumerate}
\item
  (ii) Put \(r'=L_mr+10(J_{m+1}+3)L_m\). Then
  \(a_{m+1}r'=a_mr+10a_m(J_{m+1}+3)\le a_mr+\tau_m/3\). Adding
  \(\tau_{m+1}\le\frac23\tau_m\) gives
  \(a_{m+1}r'+\tau_{m+1}\le a_mr+\tau_m\), and one iterates.
\item
  (iii) The cell of \(i\) at level \(m\) is exactly
  \(\Lambda_{(D-1)/2}(\iota_{h\to m}i)\), with \(D=L_h\cdots L_{m-1}\)
  odd and \(a_mD=a_h\). By (ii) and room,
  \(2a_mr_{k\to m}(r)\lt a_h-2\tau_m\lt a_mD\), so
  \(2r_{k\to m}(r)\lt D\).
\item
  (iv) \(a_k\to0\) and \(\tau_k\to0\).
\end{itemize}

\emph{Lean:}
\href{https://github.com/yang-mills-group/OpenAI-YM-20261028/blob/3456e5e75cf8d4cd558955a5c7638fba2d68f464/YM/Push/GRID_APPROXIMATION/Geometry.lean\#L97}{\texttt{expansion}},
\href{https://github.com/yang-mills-group/OpenAI-YM-20261028/blob/3456e5e75cf8d4cd558955a5c7638fba2d68f464/YM/Push/GRID_APPROXIMATION/Geometry.lean\#L99}{\texttt{grows}},
\href{https://github.com/yang-mills-group/OpenAI-YM-20261028/blob/3456e5e75cf8d4cd558955a5c7638fba2d68f464/YM/Push/GRID_APPROXIMATION/Geometry.lean\#L103-L104}{\texttt{grows\_next}},
\href{https://github.com/yang-mills-group/OpenAI-YM-20261028/blob/3456e5e75cf8d4cd558955a5c7638fba2d68f464/YM/Push/GRID_APPROXIMATION/Geometry.lean\#L110-L121}{\texttt{tail\_next}},
\href{https://github.com/yang-mills-group/OpenAI-YM-20261028/blob/3456e5e75cf8d4cd558955a5c7638fba2d68f464/YM/Push/GRID_APPROXIMATION/Geometry.lean\#L123-L133}{\texttt{tail\_zero}},
\href{https://github.com/yang-mills-group/OpenAI-YM-20261028/blob/3456e5e75cf8d4cd558955a5c7638fba2d68f464/YM/Push/GRID_APPROXIMATION/Geometry.lean\#L135-L145}{\texttt{scales\_next}},
\href{https://github.com/yang-mills-group/OpenAI-YM-20261028/blob/3456e5e75cf8d4cd558955a5c7638fba2d68f464/YM/Push/GRID_APPROXIMATION/Geometry.lean\#L147-L150}{\texttt{scales\_budget}},
\href{https://github.com/yang-mills-group/OpenAI-YM-20261028/blob/3456e5e75cf8d4cd558955a5c7638fba2d68f464/YM/Push/GRID_APPROXIMATION/Geometry.lean\#L49-L85}{\texttt{tiling\_window}},
\href{https://github.com/yang-mills-group/OpenAI-YM-20261028/blob/3456e5e75cf8d4cd558955a5c7638fba2d68f464/YM/Push/GRID_APPROXIMATION/Geometry.lean\#L87-L91}{\texttt{window\_fits}},
\href{https://github.com/yang-mills-group/OpenAI-YM-20261028/blob/3456e5e75cf8d4cd558955a5c7638fba2d68f464/YM/Push/GRID_APPROXIMATION/Geometry.lean\#L160-L169}{\texttt{room\_fits}},
\href{https://github.com/yang-mills-group/OpenAI-YM-20261028/blob/3456e5e75cf8d4cd558955a5c7638fba2d68f464/YM/Push/GRID_APPROXIMATION/Geometry.lean\#L155-L158}{\texttt{plenty}}.

\textbf{B7. Two consequences used downstream.}

\begin{enumerate}
\def\labelenumi{(\alph{enumi})}
\item
  \emph{Replicated differences (used by III.11).} Fix \(N,p\) and
  \(h\le k\le N\). Let \(O\) be a bounded measurable complex function at
  level \(N\), let \(F\in\mathcal I(\Lambda_w(0))\) at level \(k\) with
  \(\mathrm{room}(h,k,w)\), and let \(s\) be a finite set of level-\(h\)
  sites. Then
  \[Z_N\big\langle\prod_{i\in s}\mathrm{copy}_i(\mathcal B_{k\to N}F-O)\big\rangle_N=\sum_{l\subseteq s}\int\mathcal B^{\ast}_{k\to N}\Big[e^{-S_N}\prod_{i\in s\setminus l}(-\mathrm{copy}_iO)\Big]\prod_{i\in l}\mathrm{copy}_iF.\]
  Here \(S_N=\beta_{N,p,N}E_N\) is the Wilson action,
  \(Z_N=\int e^{-S_N}\,dU\), and
  \(\langle Q\rangle_N=Z_N^{-1}\int e^{-S_N}Q\,dU\). The copies on each
  side use the level-\(h\) tiling at the level of their argument; the
  integrals on the right are at level \(k\). Expand the product, commute
  the copies with subtraction, then apply III.8.5 and the duality in
  III.8.6(ii). The old factors stay together inside one push-forward; no
  independence of \(O\) is required.
\item
  \emph{Chessboard telescoping (used by III.13, VII.2).} Fix \(h\le N\),
  \(p\) and a real coupling \(\beta\) for which the chessboard
  inequality of IV.1 holds; the Wilson law has this property for
  \(\beta\ge0\). Let \(T\) be the centred level-\(h\) tiling at level
  \(N\). For a central-cell observable \(Q\), define
  \[\mathrm{root}_T(Q)=\Big|\Big\langle\prod_{a\in T}\mathrm{copy}_aQ\Big\rangle_\beta\Big|^{1/(\sharp T)},\]
  where \(\sharp T\) is the number of cells. Let the index set be
  finite, let \(a_i\) be distinct level-\(h\) sites, and let
  \(O_i,Z_i\in\mathcal I(T_0)\), where \(T_0\) is the central cell. Put
  \(b_i=\iota_{h\to N}a_i\) and let \(\tau_{b_i}\) act on configurations
  by lattice translation. If real bounds \(c_i,e_i\) satisfy
  \(\mathrm{root}_T(O_i)\le c_i\) and
  \(\mathrm{root}_T(Z_i-O_i)\le e_i\), then
  \[\Big|\big\langle\prod_i(Z_i\circ\tau_{b_i})\big\rangle_\beta-\big\langle\prod_i(O_i\circ\tau_{b_i})\big\rangle_\beta\Big|\le\prod_i(e_i+c_i)-\prod_ic_i.\]
  Expand \(\prod_i(O_i+(Z_i-O_i))\) over subsets. The empty-subset term
  cancels, and the chessboard estimate bounds each remaining term by the
  corresponding product of \(c_i\) and \(e_i\). To apply it to literal
  translates, undo the reflection in each copy and, on odd cells,
  conjugate its source. These operations preserve the central-cell class
  and its root. Summing the nonempty-subset bounds gives the right-hand
  side.
\end{enumerate}

\emph{Lean:}
\href{https://github.com/yang-mills-group/OpenAI-YM-20261028/blob/3456e5e75cf8d4cd558955a5c7638fba2d68f464/YM/Push/GRID_APPROXIMATION/Joint.lean\#L102-L128}{\texttt{Joint.restored}},
\href{https://github.com/yang-mills-group/OpenAI-YM-20261028/blob/3456e5e75cf8d4cd558955a5c7638fba2d68f464/YM/Push/GRID_APPROXIMATION/Telescope.lean\#L34-L78}{\texttt{tariff}};
\href{https://github.com/yang-mills-group/OpenAI-YM-20261028/blob/3456e5e75cf8d4cd558955a5c7638fba2d68f464/YM/Push/GRID_APPROXIMATION/Translate.lean\#L104-L109}{\texttt{rooted\_translates}}.

\subsubsection{5. Where the difficulty
lies}\label{where-the-difficulty-lies-10}

\begin{itemize}
\tightlist
\item
  \textbf{Input ranges.} At positive heat order, inputs of coupling
  order \(e\) have strictly higher rank. The purely contracting
  contribution to \(V^{(j)}_{e,n}\) starts at \(V_{e,n+2j}\). This
  permits downward induction in rank at fixed coupling order. The bound
  is uniform in geometry, but its polynomial may depend on \((j,e,n)\);
  no resummation in \(g\) is proved.
\item
  \textbf{Scalar terms.} The unmarked rooted norm gives no scalar
  estimate. The marked norm sums all slots and does control rank \(0\),
  including absolute convergence of its contractions. The same
  \(\sigma\gt0\) occurs in the line, diameter, height and test-function
  weights.
\item
  \textbf{Locality of a selected minimiser.} Agreement on the read
  window makes the minimisation problems identical. Uniqueness then
  makes their selected frames equal. The read radius \(10\rho\) also
  includes fallback tests and connectors; the coordinate routes used by
  frames lie in \(\Lambda_{9\rho}\).
\item
  \textbf{Conditional factorisation.} Disjoint coarse site sets use
  disjoint tags and bonds. Their pull-backs may nevertheless read
  overlapping fine windows. Factorisation therefore concerns the
  conditional kernel, not the fine Gibbs law. Iteration requires
  separation at every intermediate level, supplied by the
  cell-containment bound.
\item
  \textbf{Physical size.} The lattice radius grows by
  \(R\mapsto L_mR+10\rho_m\), while its physical radius stays below
  \(a_kr+30a_k(J_k+5)\). For fixed \(h,r\), this bound eventually fits
  inside a level-\(h\) cell.
\end{itemize}

\subsubsection{6. Notation}\label{notation-10}

{\def\LTcaptype{none} % do not increment counter
\begin{longtable}[]{@{}
  >{\raggedright\arraybackslash}p{(\linewidth - 4\tabcolsep) * \real{0.3333}}
  >{\raggedright\arraybackslash}p{(\linewidth - 4\tabcolsep) * \real{0.3333}}
  >{\raggedright\arraybackslash}p{(\linewidth - 4\tabcolsep) * \real{0.3333}}@{}}
\toprule\noalign{}
\begin{minipage}[b]{\linewidth}\raggedright
Written proof
\end{minipage} & \begin{minipage}[b]{\linewidth}\raggedright
Lean
\end{minipage} & \begin{minipage}[b]{\linewidth}\raggedright
Meaning
\end{minipage} \\
\midrule\noalign{}
\endhead
\bottomrule\noalign{}
\endlastfoot
\(I\), \(d\), \(\mathrm{diam}_d\) &
\href{https://github.com/yang-mills-group/OpenAI-YM-20261028/blob/3456e5e75cf8d4cd558955a5c7638fba2d68f464/YM/RG08/Forms/Ground.lean\#L20-L24}{\texttt{Ground}},
\href{https://github.com/yang-mills-group/OpenAI-YM-20261028/blob/3456e5e75cf8d4cd558955a5c7638fba2d68f464/YM/RG08/Forms/Ground.lean\#L35}{\texttt{Ground.radius}}
& cells with a pseudometric; largest distance within a tuple \\
rank-\(n\) tensor &
\href{https://github.com/yang-mills-group/OpenAI-YM-20261028/blob/3456e5e75cf8d4cd558955a5c7638fba2d68f464/YM/RG08/Forms/Ground.lean\#L17}{\texttt{Tape}},
\href{https://github.com/yang-mills-group/OpenAI-YM-20261028/blob/3456e5e75cf8d4cd558955a5c7638fba2d68f464/YM/RG08/Forms/Ground.lean\#L18}{\texttt{Flake}};
\href{https://github.com/yang-mills-group/OpenAI-YM-20261028/blob/3456e5e75cf8d4cd558955a5c7638fba2d68f464/YM/RG08/Forms/Ground.lean}{file}
& function on \(I^n\), ordered slots \\
\(\lVert F\rVert_\sigma\le p\) &
\href{https://github.com/yang-mills-group/OpenAI-YM-20261028/blob/3456e5e75cf8d4cd558955a5c7638fba2d68f464/YM/RG08/Forms/Ground.lean\#L61}{\texttt{Held}}
& rooted weighted \(\ell^1\) bound (empty at rank \(0\)) \\
\(\Gamma\), \([\Gamma]_\sigma\le B\) &
\href{https://github.com/yang-mills-group/OpenAI-YM-20261028/blob/3456e5e75cf8d4cd558955a5c7638fba2d68f464/YM/RG08/Forms/Lines.lean\#L14}{\texttt{Line}},
\href{https://github.com/yang-mills-group/OpenAI-YM-20261028/blob/3456e5e75cf8d4cd558955a5c7638fba2d68f464/YM/RG08/Forms/Lines.lean\#L18-L22}{\texttt{Wire}}
& propagator; weighted row and column sums \\
\((W,\delta)\)-controlled &
\href{https://github.com/yang-mills-group/OpenAI-YM-20261028/blob/3456e5e75cf8d4cd558955a5c7638fba2d68f464/YM/RG08/Forms/Cot.lean\#L15-L19}{\texttt{Cot}},
\href{https://github.com/yang-mills-group/OpenAI-YM-20261028/blob/3456e5e75cf8d4cd558955a5c7638fba2d68f464/YM/RG08/Forms/Cot.lean\#L21-L24}{\texttt{WireCot}}
& both at most \(W\), difference at most \(W\delta\) \\
polynomial majorant &
\href{https://github.com/yang-mills-group/OpenAI-YM-20261028/blob/3456e5e75cf8d4cd558955a5c7638fba2d68f464/YM/RG08/Forms/Powers.lean\#L11}{\texttt{Wage}}
& \(0\le P(x)\le Cx^t\) for \(x\ge1\) \\
\(V=(V_{e,n})\), eligible &
\href{https://github.com/yang-mills-group/OpenAI-YM-20261028/blob/3456e5e75cf8d4cd558955a5c7638fba2d68f464/YM/RG08/Steam/Gears.lean\#L23}{\texttt{Grain}},
\href{https://github.com/yang-mills-group/OpenAI-YM-20261028/blob/3456e5e75cf8d4cd558955a5c7638fba2d68f464/YM/RG08/Steam/Valves.lean\#L17-L20}{\texttt{Eligible}}
& formal potential \\
\(\Delta_\Gamma\), \(\star_\Gamma\), \(\mathrm{Sym}\) &
\href{https://github.com/yang-mills-group/OpenAI-YM-20261028/blob/3456e5e75cf8d4cd558955a5c7638fba2d68f464/YM/RG08/Forms/Lap.lean\#L17}{\texttt{lap}},
\href{https://github.com/yang-mills-group/OpenAI-YM-20261028/blob/3456e5e75cf8d4cd558955a5c7638fba2d68f464/YM/RG08/Forms/Weld.lean\#L21}{\texttt{bridge}},
\href{https://github.com/yang-mills-group/OpenAI-YM-20261028/blob/3456e5e75cf8d4cd558955a5c7638fba2d68f464/YM/RG08/Forms/Shuffle.lean\#L38}{\texttt{bleach}}
& contraction, join, symmetrisation \\
\(V^{(j)}\) &
\href{https://github.com/yang-mills-group/OpenAI-YM-20261028/blob/3456e5e75cf8d4cd558955a5c7638fba2d68f464/YM/RG08/Steam/Gears.lean\#L46}{\texttt{burn}}
& heat coefficient \\
\(T_j(e,n)\) &
\href{https://github.com/yang-mills-group/OpenAI-YM-20261028/blob/3456e5e75cf8d4cd558955a5c7638fba2d68f464/YM/RG08/Steam/Routing.lean\#L14}{\texttt{taps}}
(\texttt{Residual} is \(T_1\)) & admissible inputs \\
\(\mathcal E_t\), \(\mathcal R_t\) &
\href{https://github.com/yang-mills-group/OpenAI-YM-20261028/blob/3456e5e75cf8d4cd558955a5c7638fba2d68f464/YM/RG08/Steam/Valves.lean\#L74}{\texttt{ember}},
\href{https://github.com/yang-mills-group/OpenAI-YM-20261028/blob/3456e5e75cf8d4cd558955a5c7638fba2d68f464/YM/RG08/Steam/Valves.lean\#L75}{\texttt{cinder}}
& full series; positive-time part \\
\(h\), \(h(x)\) &
\href{https://github.com/yang-mills-group/OpenAI-YM-20261028/blob/3456e5e75cf8d4cd558955a5c7638fba2d68f464/YM/RG08/Jetty/Bank.lean\#L15-L17}{\texttt{Loom}},
\href{https://github.com/yang-mills-group/OpenAI-YM-20261028/blob/3456e5e75cf8d4cd558955a5c7638fba2d68f464/YM/RG08/Jetty/Bank.lean\#L18}{\texttt{Loom.cage}}
& height; its maximum over a tuple \\
\(\lVert f\rVert_{k,\sigma}\), \(\lVert\lambda\rVert^{\ast}_{k,\sigma}\)
&
\href{https://github.com/yang-mills-group/OpenAI-YM-20261028/blob/3456e5e75cf8d4cd558955a5c7638fba2d68f464/YM/RG08/Jetty/Hoard.lean\#L15}{\texttt{beam}},
\href{https://github.com/yang-mills-group/OpenAI-YM-20261028/blob/3456e5e75cf8d4cd558955a5c7638fba2d68f464/YM/RG08/Jetty/Dosh.lean\#L40}{\texttt{purse}}
& weighted \(C^k\) norm of test functions; dual norm \\
\(\lVert M\rVert_h\); \((U,\delta)\)-controlled &
\href{https://github.com/yang-mills-group/OpenAI-YM-20261028/blob/3456e5e75cf8d4cd558955a5c7638fba2d68f464/YM/RG08/Jetty/Bank.lean\#L36}{\texttt{Sings}},
\href{https://github.com/yang-mills-group/OpenAI-YM-20261028/blob/3456e5e75cf8d4cd558955a5c7638fba2d68f464/YM/RG08/Skiff/Pair.lean\#L11-L17}{\texttt{Dyad}}
& marked norm; marked pairs \\
marked \(\Delta_\Gamma\), \(\star_\Gamma\) &
\href{https://github.com/yang-mills-group/OpenAI-YM-20261028/blob/3456e5e75cf8d4cd558955a5c7638fba2d68f464/YM/RG08/Skiff/Pocket.lean\#L28}{\texttt{pocket}},
\href{https://github.com/yang-mills-group/OpenAI-YM-20261028/blob/3456e5e75cf8d4cd558955a5c7638fba2d68f464/YM/RG08/Skiff/Pocket.lean\#L75}{\texttt{cloak}}
& marked contraction and join \\
\(M^{(j)}\); \(\sum_{j\ge1}t^jM^{(j)}\) &
\href{https://github.com/yang-mills-group/OpenAI-YM-20261028/blob/3456e5e75cf8d4cd558955a5c7638fba2d68f464/YM/RG08/Skiff/Spoke.lean\#L31}{\texttt{skate}};
\href{https://github.com/yang-mills-group/OpenAI-YM-20261028/blob/3456e5e75cf8d4cd558955a5c7638fba2d68f464/YM/RG08/Skiff/Spoke.lean\#L71}{\texttt{shaker}}
& marked heat coefficients; positive-time part \\
\(k_\beta\) &
\href{https://github.com/yang-mills-group/OpenAI-YM-20261028/blob/3456e5e75cf8d4cd558955a5c7638fba2d68f464/YM/RG/Block.lean\#L71-L72}{\texttt{StepData.density}}
& conditional density of the coarse field given the fine field \\
\(\mathcal B\) &
\href{https://github.com/yang-mills-group/OpenAI-YM-20261028/blob/3456e5e75cf8d4cd558955a5c7638fba2d68f464/YM/Push/GRID_APPROXIMATION/KernelBasic.lean\#L21-L22}{\texttt{back}}
& one-step pull-back \\
\(q_x\), \(p_b\) &
\href{https://github.com/yang-mills-group/OpenAI-YM-20261028/blob/3456e5e75cf8d4cd558955a5c7638fba2d68f464/YM/Push/GRID_APPROXIMATION/KernelLocal.lean\#L34-L35}{\texttt{perSite}},
\href{https://github.com/yang-mills-group/OpenAI-YM-20261028/blob/3456e5e75cf8d4cd558955a5c7638fba2d68f464/YM/Push/GRID_APPROXIMATION/KernelLocal.lean\#L30-L32}{\texttt{perLink}}
& tag law at a site; link law at a bond \\
\(\Lambda_r(x)\) &
\href{https://github.com/yang-mills-group/OpenAI-YM-20261028/blob/3456e5e75cf8d4cd558955a5c7638fba2d68f464/YM/Push/GRID_APPROXIMATION/Cone.lean\#L22-L23}{\texttt{window}}
& periodic \(\ell^\infty\) box \\
\(10\rho\); \(\rho=(J+3)L\) &
\href{https://github.com/yang-mills-group/OpenAI-YM-20261028/blob/3456e5e75cf8d4cd558955a5c7638fba2d68f464/YM/Push/GRID_APPROXIMATION/Cone.lean\#L65}{\texttt{reach}};
\href{https://github.com/yang-mills-group/OpenAI-YM-20261028/blob/3456e5e75cf8d4cd558955a5c7638fba2d68f464/YM/RG/Patch.lean\#L51}{\texttt{Patch.radius}}
& read radius; patch radius \\
\(\mathcal I(X)\) &
\href{https://github.com/yang-mills-group/OpenAI-YM-20261028/blob/3456e5e75cf8d4cd558955a5c7638fba2d68f464/YM/Cells/Cut.lean\#L109-L113}{\texttt{Interior}}
& gauge invariant, measurable, bounded, local on \(X\) \\
\(\mathcal B_{k\to m}\), \(\mathcal B^{\ast}_{k\to m}\) &
\href{https://github.com/yang-mills-group/OpenAI-YM-20261028/blob/3456e5e75cf8d4cd558955a5c7638fba2d68f464/YM/Push/GRID_APPROXIMATION/Iterate.lean\#L24-L25}{\texttt{backTo}},
\href{https://github.com/yang-mills-group/OpenAI-YM-20261028/blob/3456e5e75cf8d4cd558955a5c7638fba2d68f464/YM/Push/GRID_APPROXIMATION/Iterate.lean\#L27-L29}{\texttt{pushTo}}
& iterated pull-back; push-forward \\
\(\iota_{h\to m}\) &
\href{https://github.com/yang-mills-group/OpenAI-YM-20261028/blob/3456e5e75cf8d4cd558955a5c7638fba2d68f464/YM/Push/GRID_APPROXIMATION/Geometry.lean\#L24}{\texttt{anchor}}
& image of a level-\(h\) site at level \(m\) \\
\(r_{k\to m}(r)\) &
\href{https://github.com/yang-mills-group/OpenAI-YM-20261028/blob/3456e5e75cf8d4cd558955a5c7638fba2d68f464/YM/Push/GRID_APPROXIMATION/Geometry.lean\#L99}{\texttt{grows}},
\href{https://github.com/yang-mills-group/OpenAI-YM-20261028/blob/3456e5e75cf8d4cd558955a5c7638fba2d68f464/YM/Push/GRID_APPROXIMATION/Geometry.lean\#L97}{\texttt{expansion}}
& cone radius in lattice units \\
\(\tau_m\) &
\href{https://github.com/yang-mills-group/OpenAI-YM-20261028/blob/3456e5e75cf8d4cd558955a5c7638fba2d68f464/YM/Push/GRID_APPROXIMATION/Geometry.lean\#L107}{\texttt{tail}}
& physical cone budget \(30a_m(J_m+5)\) \\
\(\mathrm{room}(h,k,r)\) &
\href{https://github.com/yang-mills-group/OpenAI-YM-20261028/blob/3456e5e75cf8d4cd558955a5c7638fba2d68f464/YM/Push/GRID_APPROXIMATION/Geometry.lean\#L153}{\texttt{room}}
& \(2(a_kr+\tau_k)\lt a_h\) \\
\end{longtable}
}

\subsubsection{7. What remains
unexplained}\label{what-remains-unexplained-10}

\begin{enumerate}
\def\labelenumi{\arabic{enumi}.}
\tightlist
\item
  \textbf{Application of the bounds.} The marked bound III.8.3 feeds
  III.14 through
  \href{https://github.com/yang-mills-group/OpenAI-YM-20261028/blob/3456e5e75cf8d4cd558955a5c7638fba2d68f464/YM/RG08/Necks/Prior.lean\#L61-L95}{\texttt{Necks.slushy\_uniform}}.
  The support and finite-sum identities of III.8.2 also enter III.14
  directly. The choices of \(W,\delta,\sigma,h\) and the covariance line
  bounds must be explained there; their derivation from earlier
  covariance estimates has not been traced here. The paired covariance
  and potential bounds are explicit hypotheses of that application.
\item
  \textbf{Gaussian interpretation.} A1 verifies the unmarked formal heat
  equations for finite index sets and symmetric input tensors and
  propagators. For symmetric marked inputs, the intended identity is
  \[M_t(\phi)=\frac{\mathbb E_{t\Gamma}[F(g,\phi+\zeta)e^{V(g,\phi+\zeta)}]}{\mathbb E_{t\Gamma}[e^{V(g,\phi+\zeta)}]}\qquad\text{(probable)},\]
  coefficientwise in \(g\), with expectation interpreted by Wick's rule.
  The marked source announces this identification for later work, but a
  proof was not located. In particular the marked rank-\(0\) coefficient
  should not yet be described as a proved normalised expectation. No
  analytic convergence of the series in \(g\) is asserted.
\item
  \textbf{Infinite-lattice scalars.} The unmarked norm leaves rank \(0\)
  uncontrolled. The scalar recursions are total-sum definitions, whose
  values need not represent convergent infinite-volume sums. A use of
  these scalars requires a separate finite-volume or vacuum-density
  argument.
\item
  \textbf{Scope beyond the stated results.} Step A5 gives bounds outside
  the main-theorem closure; III.14 uses the transported-potential
  definitions without those bounds. The marked transport in
  \texttt{Skiff/Pave.lean} and the further applications in
  \texttt{GRID\_APPROXIMATION/Birth}, \texttt{Nil}, \texttt{Candidates},
  \texttt{Picture}, \texttt{Estimate} and \texttt{Attach*} are not
  developed in this outline. The relation between these formal connected
  coefficients and the polymer activities of III.5 also remains to be
  explained in III.14 (tentative).
\item
  \textbf{Support packages.} About 1,970 assigned lines lie outside the
  two arguments above. Their mathematical roles have only been
  classified from documentation (tentative): \texttt{Carrel} treats path
  templates and their leading curl; \texttt{Camera} treats jets of chart
  germs; \texttt{Room} treats field pull-backs, Jacobians and centred
  Wick operations; \texttt{Rudder} treats symmetric lifts; and
  \texttt{Pert/Bulk} and \texttt{Packets} supply weighted sums.
  \texttt{AnalyticNext/QUARTIC\_MATCH} belongs with the quartic analysis
  of III.7 and III.10. Other assigned packages include \texttt{Isles},
  \texttt{Anchor}, parts of \texttt{Cloud} and \texttt{Glyph},
  \texttt{WilsonBulk/Web} and \texttt{ProofPorts/…}. These packages need
  separate outlines if their results are to be included here.
\end{enumerate}

\subsection{III.9. The marginal flow: the running coupling and its
reset}\label{iii.9.-the-marginal-flow-the-running-coupling-and-its-reset}

\textbf{Part:} III, Ultraviolet construction (renormalization group).
\textbf{Lean:} chapter 21:
\href{https://github.com/yang-mills-group/OpenAI-YM-20261028/tree/3456e5e75cf8d4cd558955a5c7638fba2d68f464/YM/UVNext/COUPLING_RESET}{\texttt{YM/UVNext/COUPLING\_RESET/}},
\href{https://github.com/yang-mills-group/OpenAI-YM-20261028/tree/3456e5e75cf8d4cd558955a5c7638fba2d68f464/YM/AnalyticBridge/VERIFIED_RESET_STEP}{\texttt{YM/AnalyticBridge/VERIFIED\_RESET\_STEP/}},
\href{https://github.com/yang-mills-group/OpenAI-YM-20261028/tree/3456e5e75cf8d4cd558955a5c7638fba2d68f464/YM/Extra/WILSON_STARTUP}{\texttt{YM/Extra/WILSON\_STARTUP/}},
\href{https://github.com/yang-mills-group/OpenAI-YM-20261028/tree/3456e5e75cf8d4cd558955a5c7638fba2d68f464/YM/Push/BETA_TRANSPORT}{\texttt{YM/Push/BETA\_TRANSPORT/}},
\href{https://github.com/yang-mills-group/OpenAI-YM-20261028/tree/3456e5e75cf8d4cd558955a5c7638fba2d68f464/YM/Further/STRONG_STARTUP_CREDIT}{\texttt{YM/Further/STRONG\_STARTUP\_CREDIT/}},
\href{https://github.com/yang-mills-group/OpenAI-YM-20261028/tree/3456e5e75cf8d4cd558955a5c7638fba2d68f464/YM/Push/PRICE_COMPILER}{\texttt{YM/Push/PRICE\_COMPILER/}},
\href{https://github.com/yang-mills-group/OpenAI-YM-20261028/tree/3456e5e75cf8d4cd558955a5c7638fba2d68f464/YM/Extra/MULTISOURCE_BOUNDS}{\texttt{YM/Extra/MULTISOURCE\_BOUNDS/}}.
It also includes five small packages:
\href{https://github.com/yang-mills-group/OpenAI-YM-20261028/tree/3456e5e75cf8d4cd558955a5c7638fba2d68f464/YM/PhysicalNext/UNMASKED_THROUGH_LAUNCH}{\texttt{PhysicalNext/UNMASKED\_THROUGH\_LAUNCH}},
\href{https://github.com/yang-mills-group/OpenAI-YM-20261028/tree/3456e5e75cf8d4cd558955a5c7638fba2d68f464/YM/Planned/RECIPE_STARTUP_PORT}{\texttt{Planned/RECIPE\_STARTUP\_PORT}},
\href{https://github.com/yang-mills-group/OpenAI-YM-20261028/tree/3456e5e75cf8d4cd558955a5c7638fba2d68f464/YM/R2I/Wind}{\texttt{R2I/Wind}},
\href{https://github.com/yang-mills-group/OpenAI-YM-20261028/tree/3456e5e75cf8d4cd558955a5c7638fba2d68f464/YM/Further/THRESHOLD_SOURCE_RESTART}{\texttt{Further/THRESHOLD\_SOURCE\_RESTART}}
and
\href{https://github.com/yang-mills-group/OpenAI-YM-20261028/tree/3456e5e75cf8d4cd558955a5c7638fba2d68f464/YM/PhysicalNext/SERVICE}{\texttt{PhysicalNext/SERVICE}}.
The interfaces this unit fills are in
\href{https://github.com/yang-mills-group/OpenAI-YM-20261028/blob/3456e5e75cf8d4cd558955a5c7638fba2d68f464/YM/RG/History.lean}{\texttt{YM/RG/History.lean}}
and
\href{https://github.com/yang-mills-group/OpenAI-YM-20261028/blob/3456e5e75cf8d4cd558955a5c7638fba2d68f464/YM/PCalc/Band.lean}{\texttt{YM/PCalc/Band.lean}},
both in unit III.1. \textbf{Size:} 7,215 lines in the main-theorem
closure:

{\def\LTcaptype{none} % do not increment counter
\begin{longtable}[]{@{}ll@{}}
\toprule\noalign{}
package & lines \\
\midrule\noalign{}
\endhead
\bottomrule\noalign{}
\endlastfoot
verified reset step & 2,659 \\
Wilson start-up & 2,422 \\
transport of the coupling & 652 \\
coupling reset & 592 \\
multi-source bounds & 344 \\
each remaining package & at most 213 \\
\end{longtable}
}

\textbf{Conventions.}

\begin{itemize}
\tightlist
\item
  \textbf{Levels.} Levels \(m=0,1,\dots,N\) run from the coarsest to the
  finest:

  \begin{itemize}
  \tightlist
  \item
    \(m=0\) is the torus of side \(L_p\) at unit spacing; the sides
    \(L_p\) are multiples of \(4\) and strictly increase with \(p\);
  \item
    \(m=N\) is the Wilson lattice.
  \end{itemize}
\item
  \textbf{Blocking.} The step \(m+1\to m\) blocks by an odd factor
  \(L_m\ge3\). The spacing at level \(m\) is
  \(a_m=(L_0\cdots L_{m-1})^{-1}\), and the scale is
  \(s_m:=\log(1/a_m)=\sum_{i<m}\log L_i\). We write \(\Lambda_m\) for
  the level-\(m\) torus (at volume index \(p\)) and \(|\Lambda_m|\) for
  its number of sites.
\item
  \textbf{Gauge group.} \(G\subset U(n)\) is closed, and its Lie algebra
  \(\mathfrak g\) carries the inner product
  \(\langle X,Y\rangle=-\mathrm{Re}\mathrm{tr}(XY)/n\).
\item
  \textbf{Plaquettes.} \(s(u)=1-\mathrm{Re}\mathrm{tr}u/n\) is the
  plaquette function. \(F_p(U)\in\mathfrak g\) is the orthogonal
  projection of \(\tfrac12(U_p-U_p^{-1})\), and \(S_W=\sum_p s(U_p)\) is
  the Wilson action. Since \(s(e^X)=|X|^2/2+O(|X|^4)\), the inverse
  coupling \(\beta\) equals \(1/g_0^2\) in this normalization.
\end{itemize}

\textbf{Effective actions.} An effective action at level \(m\) is a
\emph{representation} \(\rho=(P,c,E,K,b,\beta,v)\). It evaluates to
\[\rho(U)=\exp\big(v|\Lambda_m|-\beta\,\mathcal K_P(U)-\mathcal L_P(U)-\mathcal E(U)\big)\,\mathcal G(U).\]

\begin{itemize}
\tightlist
\item
  \textbf{The kinetic functional \(\mathcal K_P\).} This is the
  \(\beta\)-coefficient of the leading terms of the path data \(P\):

  \begin{itemize}
  \tightlist
  \item
    a clipped plaquette term;
  \item
    smeared curvature-squared terms;
  \item
    cubic and quartic local polynomials in curvatures along lattice
    paths.
  \end{itemize}

  These are averaged over the point group of the lattice, i.e.~signed
  permutations of the axes.
\item
  \textbf{The other factors.} \(\mathcal L_P\) is the
  \(\beta\)-independent loop functional. \(\mathcal E\) is the sum of
  the regular activities multiplied by their small-field indicators.
  Their weighted norms are summed at each anchor. \(\mathcal G\) is the
  gas of forced activities, whose weighted uniform bounds are summed
  over supports containing each site. \(b=(k,q,a,B,t)\) records the
  derivative order \(k\ge6\), spatial power \(q\ge2\), support exponent
  \(a>0\), cover radius \(B\ge1\) and threshold \(0<t\le1\) of the
  polymer norms.
\item
  \textbf{Normalization.} The zero-momentum quadratic symbol of the
  \(\beta\)-term is the identity on plaquette orientations. The loop
  kernel has vanishing zeroth and first moments. These conditions fix
  the normalization of \(\beta\) within the representation.
\end{itemize}

\textbf{The block-spin step.} \(T^{(m)}_{\beta'}\) denotes one step from
level \(m+1\) to level \(m\) at inverse coupling \(\beta'\):
\[(T^{(m)}_{\beta'}\rho')(V)=\int k_{\beta'}(V,U)\,\rho'(U)\,dU.\] The
kernel \(k_{\beta'}\) is a probability kernel (unit III.1). Densities
are taken relative to product Haar probability measure and are not
normalized.

\subsubsection{1. What this unit proves}\label{what-this-unit-proves-11}

This unit tracks the scalar coefficient \(\beta\) of the normalized
curvature-squared term through the blocking history. Its three results
are:

\begin{itemize}
\tightlist
\item
  \textbf{(a) Start-up.} With the start-up bounds, Wilson top path data
  and recipe acceptance, the Wilson weight at the finest level is
  admissible for every \(\beta\) in the band.
\item
  \textbf{(b) One step.} Given the early presentation, local bounds,
  numerical conditions and recipe acceptance stated below, one
  block-spin step followed by a \emph{reset} of \(\beta\) maps a
  continuous admissible family at level \(m+1\) to one at level \(m\).
  The coupling moves by \(\alpha_{N,m}+d\), where \(\alpha_{N,m}\) is
  the one-loop coefficient of the step and \(|d|\le r_m\le1\).
\item
  \textbf{(c) Thresholds.} The delay-dependent numerical conditions hold
  beyond entry thresholds, provided the earlier constants satisfy the
  compatibility inequalities of III.9.7. These choices respect the order
  of constants in III.2.
\end{itemize}

Combined with the shooting argument of III.2 and the constructions of
III.10 and III.13, this yields the selected history of III.9.3 and the
law \(\beta_{N,p}/\log(1/a_N)\to\kappa\) of the main theorem.

\textbf{Proposition III.9.1 (the slope constant).} Let \(\mathfrak g\)
be simple and nonabelian, and let \((b_i)\) be an orthonormal basis of
\((\mathfrak g,\langle\cdot,\cdot\rangle)\). Put
\[C_G=-\frac1{\dim\mathfrak g}\sum_i\mathrm{tr}_{\mathfrak g}\big((\mathrm{ad}b_i)^2\big),\qquad \kappa=\frac{11\,C_G}{24\pi^2}.\]
Then \(\kappa>0\).

Here \(C_G\) is the adjoint Casimir in the chosen inner product, and
\(\kappa=2b_0\) in the conventional one-loop notation
\(b_0=11C_G/(48\pi^2)\). For \(SU(N)\) in the defining representation,
this normalization gives \(C_G=2N^2\) and \(\kappa=11N^2/(12\pi^2)\).

The invariant form \(-\mathrm{tr}(\mathrm{ad}X\,\mathrm{ad}Y)\) is
\(C_G\langle X,Y\rangle\). Its value at \((X,X)\) is
\(\sum_i|[X,b_i]|^2\), and nonabelianness makes this positive for some
\(X\). Thus \(C_G>0\).

\emph{Status:} definition; positivity is a theorem and uses simplicity.
Inside this unit only two things are used: \(\kappa>0\), and the
identity \(\hat\beta_{m+1}=\hat\beta_m+\kappa\log L_m\). The fact that
the one-loop coefficient of the actual blocked action is
\(-\kappa\log L+O(1)\) is proved in III.7, where \(11/(24\pi^2)\) is the
value of an explicit momentum integral (§2).

\emph{Lean:}
\href{https://github.com/yang-mills-group/OpenAI-YM-20261028/blob/3456e5e75cf8d4cd558955a5c7638fba2d68f464/YM/Gauge/Colors.lean\#L184}{\texttt{YM.Gauge.slopeConstant}},
\href{https://github.com/yang-mills-group/OpenAI-YM-20261028/blob/3456e5e75cf8d4cd558955a5c7638fba2d68f464/YM/Gauge/Colors.lean\#L180-L182}{\texttt{YM.Gauge.casimir}},
\href{https://github.com/yang-mills-group/OpenAI-YM-20261028/blob/3456e5e75cf8d4cd558955a5c7638fba2d68f464/YM/Gauge/Casimir.lean\#L174-L176}{\texttt{YM.Gauge.slopeConstant\_pos}},
\href{https://github.com/yang-mills-group/OpenAI-YM-20261028/blob/3456e5e75cf8d4cd558955a5c7638fba2d68f464/YM/Gauge/Colors.lean\#L71-L75}{\texttt{YM.Gauge.trace\_color}}.

\textbf{Proposition III.9.2 (reference coupling and scale functions).}
Fix \(\kappa>0\), \(\bar\beta>1\) and an infinite sequence of odd block
factors \(L_m\ge3\). Define
\[\hat\beta_m=\bar\beta+\kappa s_m,\quad M_m=e^{(\log\log\hat\beta_m)^2},\quad t_m=\hat\beta_m^{-1/2}M_m,\quad \varepsilon_m=\hat\beta_m^{-11/10},\quad \tilde\varepsilon_m=e^{-M_m^{1/5}},\]
and set \(B_m=\lceil(\log\hat\beta_m)^2\rceil\) and
\(\mathfrak B_m=[\hat\beta_m/2,\,2\hat\beta_m]\). These are, in order:
the reference coupling, the threshold multiplier, the small-field
threshold, the regular and forced allowances, the cover size and the
band.

Then \(\hat\beta_{m+1}=\hat\beta_m+\kappa\log L_m\), \(\hat\beta\) is
nondecreasing, and \(s_m\ge m\log3\), so \(\hat\beta_m\to\infty\).

The scale conditions are \emph{hypotheses} on the scale sequence, not
consequences of the definitions. They are:

\begin{itemize}
\tightlist
\item
  \textbf{Size conditions.} \(\hat\beta_m>1\) and \(t_m\le1\). For every
  volume index \(p\), the level-\(m\) side is at least \(100(B_m+1)\),
  and the level-\((m+1)\) side is at least \(100(R_m+1)\), where \(R_m\)
  is the block-patch radius.
\item
  \textbf{Stage index.} The stage index satisfies \(J_m\ge4\), is
  nondecreasing, \(J_{m+1}\le J_m+1\), and \(J_m\to\infty\). The
  derivative orders, spatial powers and support exponents are
  nondecreasing along the levels. There are nondecreasing stage block
  factors \(\ell_j\) with \(L_m=\ell_{J_{m+1}}\).
\item
  \textbf{Smallness exponent.} \(0<\theta_{J_m}\le1/J_m\), and
  \(\theta\) is antitone.
\item
  \textbf{Gentle scales.} \(\log L_m/s_m\to0\).
\item
  \textbf{Drift and logarithmic costs.} With \(\omega_j\) the drift
  allowance of stage \(j\),
  \[\frac{\omega_{J_{m+1}}}{\log L_m}\le J_{m+1}^{-2},\qquad \frac{\log(1+\text{paths}+\text{width}^{-1}+\text{columns}+\text{relaxation}+\text{support exponent})_{J_{m+1}}}{\log L_m}\le J_{m+1}^{-2}.\]
\end{itemize}

\emph{Status:} definitions (III.1) and elementary theorems. The scale
conditions are part of the scale sequence produced in III.2.

\emph{Lean:}
\href{https://github.com/yang-mills-group/OpenAI-YM-20261028/blob/3456e5e75cf8d4cd558955a5c7638fba2d68f464/YM/RG/History.lean\#L79}{\texttt{Calendar.hatBeta}},
\href{https://github.com/yang-mills-group/OpenAI-YM-20261028/blob/3456e5e75cf8d4cd558955a5c7638fba2d68f464/YM/RG/History.lean\#L83}{\texttt{threshold}},
\href{https://github.com/yang-mills-group/OpenAI-YM-20261028/blob/3456e5e75cf8d4cd558955a5c7638fba2d68f464/YM/RG/History.lean\#L84}{\texttt{allowance}},
\href{https://github.com/yang-mills-group/OpenAI-YM-20261028/blob/3456e5e75cf8d4cd558955a5c7638fba2d68f464/YM/RG/History.lean\#L85}{\texttt{forcedAllowance}},
\href{https://github.com/yang-mills-group/OpenAI-YM-20261028/blob/3456e5e75cf8d4cd558955a5c7638fba2d68f464/YM/RG/History.lean\#L82}{\texttt{coverSize}},
\href{https://github.com/yang-mills-group/OpenAI-YM-20261028/blob/3456e5e75cf8d4cd558955a5c7638fba2d68f464/YM/RG/History.lean\#L97-L118}{\texttt{Calendar.Fits}},
\href{https://github.com/yang-mills-group/OpenAI-YM-20261028/blob/3456e5e75cf8d4cd558955a5c7638fba2d68f464/YM/Shoot/Criteria.lean\#L63}{\texttt{Shoot.band}},
\href{https://github.com/yang-mills-group/OpenAI-YM-20261028/blob/3456e5e75cf8d4cd558955a5c7638fba2d68f464/YM/Next/SCHEDULE/Levels.lean\#L43-L45}{\texttt{Deployment.hat}},
\href{https://github.com/yang-mills-group/OpenAI-YM-20261028/blob/3456e5e75cf8d4cd558955a5c7638fba2d68f464/YM/Extra/WILSON_STARTUP/Fast.lean\#L40-L42}{\texttt{Timing.hat\_mono}},
\href{https://github.com/yang-mills-group/OpenAI-YM-20261028/blob/3456e5e75cf8d4cd558955a5c7638fba2d68f464/YM/Lattice/Scale.lean\#L117-L120}{\texttt{Schedule.divergent}}.

\textbf{Proposition III.9.3 (selected histories).} A \emph{history}
assigns to all \(N,p\) and \(m\le N\) a level-\(m\) representation
\(\rho_{N,p,m}\) with coupling \(\beta_{N,p,m}\). It is \emph{selected}
if the following hold.

\begin{enumerate}
\def\labelenumi{\arabic{enumi}.}
\item
  \textbf{Common path data.} The path data \(P_{N,m}\) of
  \(\rho_{N,p,m}\) do not depend on \(p\).
\item
  \textbf{Shape.} \(\rho_{N,p,m}\) carries the level-\(m\) norm
  parameters and clip. It is lawful at level \(m\):

  \begin{itemize}
  \tightlist
  \item
    its path data obey the stage bounds, with loop cap
    \(2M_m^{\theta_{J_m}}\);
  \item
    its regular activities are local, gauge invariant, vanish at the
    flat field, and have zero Hessian on affine link directions whenever
    their supports do not wrap; their norm bound is \(\varepsilon_m\);
  \item
    its forced activities satisfy the locality, eligibility and symmetry
    conditions of II.2, with norm bound \(\tilde\varepsilon_m\).
  \end{itemize}

  The whole activity family transforms under translations and signed
  permutations of the axes.

  Moreover \(\beta_{N,p,m}\in\mathfrak B_m\).
\item
  \textbf{Start.} \(v_{N,p,N}=0\) and
  \(\rho_{N,p,N}(U)=e^{-\beta_{N,p,N}S_W(U)}\) for all \(U\).
\item
  \textbf{Transfer.} For \(m<N\),
  \(\rho_{N,p,m}=T^{(m)}_{\beta_{N,p,m+1}}\rho_{N,p,m+1}\) at every
  configuration, and
  \[|\beta_{N,p,m+1}-\beta_{N,p,m}-\kappa\log L_m|\le\omega_{J_{m+1}}.\]
\item
  \textbf{End point.} \(\beta_{N,p,0}=\bar\beta\).
\item
  \textbf{Running.} For every \(\varepsilon>0\) there is \(m_0\) with
  \(|\beta_{N,p,m}-\hat\beta_m|\le\varepsilon s_m\) for all
  \(m\ge m_0\), \(N\ge m\) and all \(p\).
\end{enumerate}

\emph{Consequences:}

\begin{itemize}
\tightlist
\item
  \textbf{Partition function.} For all \(m\le N\),
  \(\int\rho_{N,p,m}\,d\mathrm{Haar}=Z_W(\Lambda_{N,p};\beta_{N,p,N})\),
  the Wilson partition function.
\item
  \textbf{Running law.} \(\beta_{N,p,m}/s_m\to\kappa\) uniformly in
  \(N\ge m\) and \(p\), because \(s_m\to\infty\).
\item
  \textbf{Bare coupling.} The bare coupling of the Wilson family,
  \(\beta_{N,p}:=\beta_{N,p,N}\), satisfies
  \(\sup_p|\beta_{N,p}/\log(1/a_N)-\kappa|\to0\). This is (Reg1) of the
  main theorem with \(b_*=\kappa=2b_0\). The continuum running-coupling
  comparison belongs to VII.5.
\end{itemize}

\emph{Status:} (1)--(6) are a specification (III.1). It is proved for
the history of the main theorem in III.13
(\href{https://github.com/yang-mills-group/OpenAI-YM-20261028/blob/3456e5e75cf8d4cd558955a5c7638fba2d68f464/YM/UVExtension/Spire/Spark.lean\#L78-L85}{\texttt{Spire.Uses.exists\_bolt}}),
using the reset steps of III.9.5, the step program from III.10, and the
finite-cutoff construction of Lemma 4.22. The interval shooting of Lemma
4.20 also gives a selected history from the step program. In both
constructions, Lemma 4.21 derives (6) from (4) and (5).

The partition identity follows by induction from probability-kernel
normalization and integrability of lawful representations. The
calibration condition of VIII.1 requires one datum (scales, kinetic data
and recipe) and one family of trails realizing this history, with every
step a reset step for all \(N,p\) and \(m<N\).

\emph{Lean:}
\href{https://github.com/yang-mills-group/OpenAI-YM-20261028/blob/3456e5e75cf8d4cd558955a5c7638fba2d68f464/YM/RG/History.lean\#L234-L240}{\texttt{History.Selected}},
\href{https://github.com/yang-mills-group/OpenAI-YM-20261028/blob/3456e5e75cf8d4cd558955a5c7638fba2d68f464/YM/RG/History.lean\#L213-L218}{\texttt{Conv}},
\href{https://github.com/yang-mills-group/OpenAI-YM-20261028/blob/3456e5e75cf8d4cd558955a5c7638fba2d68f464/YM/RG/History.lean\#L223-L227}{\texttt{Transfer}},
\href{https://github.com/yang-mills-group/OpenAI-YM-20261028/blob/3456e5e75cf8d4cd558955a5c7638fba2d68f464/YM/RG/History.lean\#L229-L230}{\texttt{Running}},
\href{https://github.com/yang-mills-group/OpenAI-YM-20261028/blob/3456e5e75cf8d4cd558955a5c7638fba2d68f464/YM/RG/History.lean\#L242-L249}{\texttt{History.partition\_eq}},
\href{https://github.com/yang-mills-group/OpenAI-YM-20261028/blob/3456e5e75cf8d4cd558955a5c7638fba2d68f464/YM/Spec/Closed.lean\#L58-L61}{\texttt{Spec.Calibrated}}.

\textbf{Proposition III.9.4 (one reset step along a curve).} Fix a datum
\(\mathfrak D\) consisting of three things:

\begin{itemize}
\tightlist
\item
  levels satisfying the scale and kernel conditions of III.9.2 and §2,
  and meeting the entry thresholds of the fixed numerical choices;
\item
  a compatible horizon, a local Haar chart, positive post-block
  constants, and kinetic data \((\alpha_{N,m},P_{N,m})\) satisfying the
  planned one-loop conditions (§2);
\item
  a recipe, with initial, retention and continuation tests.
\end{itemize}

Fix \(N,p\) and \(m<N\), let \(I=[l,u]\) with \(l\le u\), and let
\(x\mapsto\rho'_x\) (\(x\in I\)) be a continuous curve of admissible
level-\((m+1)\) representations.

\begin{itemize}
\tightlist
\item
  \emph{Admissible} means: path data \(P_{N,m+1}\), the level-\((m+1)\)
  norm parameters and clip, lawful with the level-\((m+1)\) constants,
  and accepted by the recipe.
\item
  \emph{Continuous} means: \(\beta\) and \(v\) are continuous; the
  regular and forced addresses, including supports, inventories and
  eligibility tests, are fixed; and the coefficients vary continuously
  in the weighted sum of regular sup-jet norms and forced uniform norms.
  Regular jets are measured only where their supports are small-field,
  at the prescribed derivative order and spatial power.
\end{itemize}

Write \(\beta'_x=\beta(\rho'_x)\) and \(H_x=T^{(m)}_{\beta'_x}\rho'_x\).
Let \((c_m,r_m)\) be the reset constants of the stage pair
\((J_{m+1},J_m)\), with \(c_m>0\) and \(0<r_m\le1\). Let
\(g=g_{J_{m+1}}>0\) be the prescribed fine-stage forced-decay constant.
Assume:

\begin{itemize}
\tightlist
\item
  \textbf{(P) Priced.} \(\eta=\eta_{J_m}\) is a sensor constant (Lemma
  4.7), and \(2(c_m+1)\varepsilon_m\le\eta\,t_m^2\,r_m\).
\item
  \textbf{(V) Probe.} \(V\) is a small-field configuration with
  \(\mathcal K_{P_{N,m}}(V)\ge|\Lambda_m|\eta t_m^2\). Such a \(V\)
  exists by (P).
\item
  \textbf{(I) Income.} There is a family \(F(x,d)\), \(x\in I\),
  \(|d|\le r_m\), of level-\(m\) activities with fixed addresses, valid
  with \((\varepsilon_m,\tilde\varepsilon_m)\). Assume the catalogue
  conditions: each coefficient is continuous in its regular sup-jet norm
  or forced uniform norm, and every convergent parameter sequence admits
  a subsequence with common summable majorants in those weighted norms.
  Moreover \(H_x(1),H_x(V)>0\), and
  \[\Big|(\beta'_x+\alpha_{N,m})\mathcal K(V)+\mathcal L(V)-\log\frac{H_x(1)}{H_x(V)}\Big|\le|\Lambda_m|c_m\varepsilon_m,\]
  \[|\mathcal E_{F(x,b)}(V)-\mathcal E_{F(x,a)}(V)|\le|\Lambda_m|c_m\varepsilon_m|b-a|.\]
\item
  \textbf{(R) Remainder.} Put
  \[\mathrm{shape}_{x,d}(U)=e^{-(\beta'_x+\alpha_{N,m}+d)\mathcal K(U)-\mathcal L(U)-\mathcal E_{F(x,d)}(U)}\,\mathcal G_{F(x,d)}(U).\]
  Require:

  \begin{itemize}
  \tightlist
  \item
    \emph{Shape control:}
    \(|\mathrm{shape}_{x,a}(U)-H_x(U)/H_x(1)|\le C_{x,U}\,|f(x,a)|\) for
    all \(U\) and \(|a|\le r_m\), where \(f\) is the scalar residual of
    Lemma 4.9 and, for each \(x,U\), a finite \(C_{x,U}\ge0\) works for
    every \(|a|\le r_m\);
  \item
    \emph{Forced bound:} the forced activities are bounded by
    \(e^{-g\,M_{m+1}^2}\);
  \item
    \emph{Recipe:} the recipe accepts every exact trial.
  \end{itemize}
\end{itemize}

Conclusion. There is a continuous \(d:I\to[-r_m,r_m]\) with the
following properties. Let \(\rho_x\) have path data \(P_{N,m}\),
activities \(F(x,d(x))\), coupling \(\beta'_x+\alpha_{N,m}+d(x)\) and
\(v_x=|\Lambda_m|^{-1}\log H_x(1)\). Then:

\begin{itemize}
\tightlist
\item
  \(x\mapsto\rho_x\) is a continuous curve;
\item
  \(\rho_x=H_x\) at every configuration;
\item
  \(\rho_x\) is admissible, its forced activities are bounded by
  \(e^{-gM_{m+1}^2}\), and
  \(|\beta(\rho_x)-\beta'_x-\alpha_{N,m}|=|d(x)|\le r_m\le1\).
\end{itemize}

Exact density transfer, recipe continuation, admissibility, the stronger
forced bound and the error bound \(|\beta-\beta'-\alpha_{N,m}|\le1\)
define a \emph{reset step}. Without (R), the same \(d\) is the unique
zero in \([-r_m,r_m]\) of the scalar residual. In that case
\(\rho_x(1)=H_x(1)\) and \(\rho_x(V)=H_x(V)\).

\emph{Status:} theorem.

\emph{Lean:}
\href{https://github.com/yang-mills-group/OpenAI-YM-20261028/blob/3456e5e75cf8d4cd558955a5c7638fba2d68f464/YM/UVNext/COUPLING_RESET/Advance.lean\#L175-L189}{\texttt{CouplingReset.advance}},
\href{https://github.com/yang-mills-group/OpenAI-YM-20261028/blob/3456e5e75cf8d4cd558955a5c7638fba2d68f464/YM/UVNext/COUPLING_RESET/Advance.lean\#L157-L173}{\texttt{correct}},
\href{https://github.com/yang-mills-group/OpenAI-YM-20261028/blob/3456e5e75cf8d4cd558955a5c7638fba2d68f464/YM/UVNext/COUPLING_RESET/Advance.lean\#L78-L85}{\texttt{Income}},
\href{https://github.com/yang-mills-group/OpenAI-YM-20261028/blob/3456e5e75cf8d4cd558955a5c7638fba2d68f464/YM/UVNext/COUPLING_RESET/Advance.lean\#L97-L103}{\texttt{Remainder}},
\href{https://github.com/yang-mills-group/OpenAI-YM-20261028/blob/3456e5e75cf8d4cd558955a5c7638fba2d68f464/YM/UVNext/COUPLING_RESET/Prices.lean\#L66-L69}{\texttt{PricedReset}},
\href{https://github.com/yang-mills-group/OpenAI-YM-20261028/blob/3456e5e75cf8d4cd558955a5c7638fba2d68f464/YM/UVNext/COUPLING_RESET/Advance.lean\#L55-L58}{\texttt{Usable}},
\href{https://github.com/yang-mills-group/OpenAI-YM-20261028/blob/3456e5e75cf8d4cd558955a5c7638fba2d68f464/YM/PCalc/Band.lean\#L49-L55}{\texttt{Directed.Move}}.

\textbf{Proposition III.9.5 (explicit reset and step program).}

\begin{enumerate}
\def\labelenumi{(\alph{enumi})}
\tightlist
\item
  \emph{Explicit correction.} Fix the datum, input curve and reset
  constants of III.9.4, and assume \(G\) is simple and the input
  couplings lie in \(\mathfrak B_{m+1}\). Suppose that, for every
  \(x\in I\), the blocked density has the early presentation
  \[H_x(U)=\exp\Big(\tau_x|\Lambda_m|-\beta'_x\mathcal K(U)-\mathcal L(U)-\sum_i\chi_{X_i}(U)E_{x,i}(U)\Big)\,\mathcal G_x(U)\,e^{-\alpha_{N,m}\mathcal S(U)},\]
  where \(\chi_X\) indicates that \(X\) is small-field and
  \(\mathcal S\) is the masked quadratic receipt of the coarse path
  data. Assume:
\end{enumerate}

\begin{itemize}
\tightlist
\item
  The early regular functions are local, gauge invariant and smooth
  through the required derivative order. Their sup-jet bounds
  \(h_{x,i}\ge0\) satisfy, at every anchor \(o\),
  \[\sum_{i:o_i=o}e^{(a+3)\,\mathrm{cover}_{B_m}(X_i)}h_{x,i}\le H_x^{E},\]
  where \(a\) is the level-\(m\) support exponent. The early forced
  functions satisfy the forced-family conditions and a site-pinned bound
  \(w_x\) at support exponent \(a+1\). The full family transforms under
  translations and the point group.
\item
  The addresses are fixed along \(I\), and the early functions satisfy
  the coefficient-continuity and summable-tail conditions used in
  III.9.4(I).
\item
  The reset is priced; the clip is compatible with the coarse norm
  parameters; and the width and cubic conditions hold. There are
  numerical output bounds \(e_x,k_x\) such that
  \(H_x^E\le c_m\varepsilon_m/4\), the regular and forced reorganization
  estimates fit within \(e_x,k_x\), and these bounds imply both validity
  with \((\varepsilon_m,\tilde\varepsilon_m)\) and the forced bound
  \(e^{-gM_{m+1}^2}\). These are the fee and closure conditions of Lemma
  4.17.
\item
  The recipe accepts the reorganized representation and its continuation
  from \(\rho'_x\).
\end{itemize}

Then the reset correction is the isotropic Hessian coefficient
\(\delta_x\) of the early regular functions (Lemma 4.12):
\[d(x)=\delta_x,\qquad \beta(\rho_x)=\beta'_x+\alpha_{N,m}+\delta_x,\qquad t_m^2|\delta_x|\le H_x^E\le c_m\varepsilon_m/4,\qquad v_x=\tau_x-|\Lambda_m|^{-1}\sum_iE_{x,i}(1).\]
The output is continuous and has all the properties of III.9.4. The
Gaussian-sector formulation additionally requires its integration
licenses and the identity between that integral and the push; a direct
early presentation of the push supplies the equality used in the reset.

\begin{enumerate}
\def\labelenumi{(\alph{enumi})}
\setcounter{enumi}{1}
\tightlist
\item
  \emph{Program.} Assume \(G\) is simple. Suppose the numerical
  conditions above hold at every level, the early presentation and its
  local bounds and continuity are available along every continuous
  admissible input curve with couplings in \(\mathfrak B_{m+1}\), and
  the recipe accepts these outputs. Suppose also that the start-up
  curves below exist and pass the recipe's initial and retention tests.
  Then:
\end{enumerate}

\begin{itemize}
\tightlist
\item
  For all \(N,p\), there is a continuous admissible curve
  \(\beta\mapsto F(\beta)\) on \(\mathfrak B_N\) with coupling
  \(\beta\), vacuum constant \(0\) and density \(e^{-\beta S_W}\),
  passing the initial tests.
\item
  For all \(N,p,m<N\), every continuous admissible input curve on a
  closed interval has, when its couplings lie in \(\mathfrak B_{m+1}\),
  a continuous output curve related to it pointwise by reset steps. The
  finer bound \(|\beta-\beta'-\alpha_{N,m}|\le r_m\) also holds.
\end{itemize}

The output coupling may lie outside \(\mathfrak B_m\); shooting
restricts the parameter interval to restore the band (Lemma 4.20).

\emph{Status:} conditional theorems. III.10 supplies the physical
inputs, including the Gaussian-sector identity and the early
presentation. Its
\href{https://github.com/yang-mills-group/OpenAI-YM-20261028/blob/3456e5e75cf8d4cd558955a5c7638fba2d68f464/YM/PhysicalNext/RG_NEXT_COMPLETE_REPEATABLE_NONLINEAR_STEP/Roost/Foundry.lean\#L59}{Roost.Factor.worksNow}
concludes the step program from these results.

\emph{Lean:}
\href{https://github.com/yang-mills-group/OpenAI-YM-20261028/blob/3456e5e75cf8d4cd558955a5c7638fba2d68f464/YM/AnalyticBridge/VERIFIED_RESET_STEP/Advance.lean\#L64-L90}{\texttt{Return.reset}},
\href{https://github.com/yang-mills-group/OpenAI-YM-20261028/blob/3456e5e75cf8d4cd558955a5c7638fba2d68f464/YM/AnalyticBridge/VERIFIED_RESET_STEP/Payment.lean\#L56-L57}{\texttt{Return.result}},
\href{https://github.com/yang-mills-group/OpenAI-YM-20261028/blob/3456e5e75cf8d4cd558955a5c7638fba2d68f464/YM/AnalyticBridge/VERIFIED_RESET_STEP/Evaluation.lean\#L26-L27}{\texttt{Collection.earlyValue}},
\href{https://github.com/yang-mills-group/OpenAI-YM-20261028/blob/3456e5e75cf8d4cd558955a5c7638fba2d68f464/YM/AnalyticBridge/VERIFIED_RESET_STEP/Downstream.lean\#L85-L88}{\texttt{Use.supplier}},
\href{https://github.com/yang-mills-group/OpenAI-YM-20261028/blob/3456e5e75cf8d4cd558955a5c7638fba2d68f464/YM/UVNext/COUPLING_RESET/Stability.lean\#L67-L70}{\texttt{CouplingReset.program}},
\href{https://github.com/yang-mills-group/OpenAI-YM-20261028/blob/3456e5e75cf8d4cd558955a5c7638fba2d68f464/YM/UVNext/COUPLING_RESET/Stability.lean\#L54-L65}{\texttt{Supplier.ports}},
\href{https://github.com/yang-mills-group/OpenAI-YM-20261028/blob/3456e5e75cf8d4cd558955a5c7638fba2d68f464/YM/PCalc/Band.lean\#L65-L72}{\texttt{Directed.Program}},
\href{https://github.com/yang-mills-group/OpenAI-YM-20261028/blob/3456e5e75cf8d4cd558955a5c7638fba2d68f464/YM/AnalyticBridge/VERIFIED_RESET_STEP/Payment.lean\#L28-L35}{\texttt{Return.Prices}},
\href{https://github.com/yang-mills-group/OpenAI-YM-20261028/blob/3456e5e75cf8d4cd558955a5c7638fba2d68f464/YM/AnalyticBridge/VERIFIED_RESET_STEP/Investor.lean\#L49-L51}{\texttt{Collection.Fees}},
\href{https://github.com/yang-mills-group/OpenAI-YM-20261028/blob/3456e5e75cf8d4cd558955a5c7638fba2d68f464/YM/AnalyticBridge/VERIFIED_RESET_STEP/Comparison.lean\#L72-L74}{\texttt{Collection.Compared}}.

\textbf{Proposition III.9.6 (Wilson start-up).} Let \(c\) be a clip and
\(b\) a set of norm parameters with \(0<t\le1\). They are
\emph{compatible} if \(c_{\rm dir}\le1\), \(B\ge2\), and
\(t\,e^{(B-2)/c_{\rm rel}}\ge2\).

\begin{enumerate}
\def\labelenumi{\arabic{enumi}.}
\tightlist
\item
  \textbf{Exact start.} For every \(\beta\in\mathbb R\) and every
  configuration \(U\), \(\rho^{\rm start}_\beta(U)=e^{-\beta S_W(U)}\).
  Here \(\rho^{\rm start}_\beta\) has:

  \begin{itemize}
  \tightlist
  \item
    the \emph{Wilson top} path data \(P^{\rm top}_c\) (Lemma 4.2);
  \item
    \(v=0\);
  \item
    regular activities \(\beta(1-c_{\rm dir})\times\) (the one-plaquette
    residuals of small-field plaquettes), averaged over the point group;
  \item
    forced activities \(e^{-\beta\,\mathrm{charge}(D,U)}\), one for each
    connected component \(D\) of the large-field region for the padded
    adjacency relation.
  \end{itemize}
\item
  \textbf{Residual bound.} For every gauge group \(G\) and all orders
  \(k\) and spatial powers \(q\) there are \(C>0\) and
  \(\delta\in(0,1]\) with the following property. For every set of norm
  parameters with order \(k\), power \(q\) and threshold \(t\le\delta\),
  every torus, every plaquette \(p\) and every site \(o\) of \(p\), the
  residual \(R_p\) has weighted sup-jet norm at most \(Ct^5\) at \(o\).
\item
  \textbf{Eventual bounds and continuity.} For each fixed stage with
  \(c_{\rm dir}\le1\), there is \(\beta^*\) such that, for every
  realized choice of levels and every \(N\) at that stage with
  \(\hat\beta_N\ge\beta^*\), the start-up conditions hold:
  compatibility, small threshold, and the regular and forced norm bounds
  \(\varepsilon_N,\tilde\varepsilon_N\). The start-up curve is
  continuous on \(\mathfrak B_N\). It is admissible for a given recipe
  if its prescribed top paths are \(P^{\rm top}\), these paths satisfy
  the stage bounds, and its initial and retention tests accept the
  displayed curve.
\item
  \textbf{Installation.} Suppose the old kinetic data satisfy the
  planned conditions, the Wilson top paths satisfy the stage bounds, and
  a revised horizon is supplied whose top quadratic datum is the Wilson
  one, whose lower data agree with the old horizon, and whose initial
  stationary guesses agree. Then the kinetic data can be replaced by
  data for that horizon with \(P_{N,N}=P^{\rm top}\), the same
  \(\alpha\), and every planned condition preserved.
\end{enumerate}

\emph{Status:} theorem.

\emph{Lean:}
\href{https://github.com/yang-mills-group/OpenAI-YM-20261028/blob/3456e5e75cf8d4cd558955a5c7638fba2d68f464/YM/Extra/WILSON_STARTUP/Gas.lean\#L96-L103}{\texttt{WilsonStartup.Start.evaluates}},
\href{https://github.com/yang-mills-group/OpenAI-YM-20261028/blob/3456e5e75cf8d4cd558955a5c7638fba2d68f464/YM/Extra/WILSON_STARTUP/Face.lean\#L76-L101}{\texttt{residual\_price}},
\href{https://github.com/yang-mills-group/OpenAI-YM-20261028/blob/3456e5e75cf8d4cd558955a5c7638fba2d68f464/YM/Extra/WILSON_STARTUP/Cake.lean\#L43-L48}{\texttt{Start.Allows}},
\href{https://github.com/yang-mills-group/OpenAI-YM-20261028/blob/3456e5e75cf8d4cd558955a5c7638fba2d68f464/YM/Extra/WILSON_STARTUP/Clock.lean\#L113-L139}{\texttt{Timing.enough}},
\href{https://github.com/yang-mills-group/OpenAI-YM-20261028/blob/3456e5e75cf8d4cd558955a5c7638fba2d68f464/YM/Planned/RECIPE_STARTUP_PORT/Layout.lean\#L52-L71}{\texttt{Planned.Port.card}},
\href{https://github.com/yang-mills-group/OpenAI-YM-20261028/blob/3456e5e75cf8d4cd558955a5c7638fba2d68f464/YM/Extra/WILSON_STARTUP/Engine.lean\#L62-L82}{\texttt{Change.engine}},
\href{https://github.com/yang-mills-group/OpenAI-YM-20261028/blob/3456e5e75cf8d4cd558955a5c7638fba2d68f464/YM/Extra/WILSON_STARTUP/Ownership.lean\#L20-L23}{\texttt{WilsonStartup.Decent}},
\href{https://github.com/yang-mills-group/OpenAI-YM-20261028/blob/3456e5e75cf8d4cd558955a5c7638fba2d68f464/YM/Extra/WILSON_STARTUP/Match.lean\#L52-L54}{\texttt{Design.Suitable}},
\href{https://github.com/yang-mills-group/OpenAI-YM-20261028/blob/3456e5e75cf8d4cd558955a5c7638fba2d68f464/YM/Extra/WILSON_STARTUP/Engine.lean\#L22-L25}{\texttt{Change.Overlay}}.

\textbf{Proposition III.9.7 (entry thresholds).} Fix a pre-block stage
profile and assume \(G\) is simple. There are sensor constants
\(\eta_j>0\), chosen before the block sizes, with the following
property. Fix the block factors \(\ell_j\), exponents
\(0<\theta_j\le1\), reset constants, forced-decay forecasts and
post-block costs. For each allowed stage transition \(i\le j\le i+1\),
let \(f_{ji}\) be the inherited coefficient and assume
\[8f_{ji}\le q_i\ell_j^{\min},\qquad u_{ji}(x)\longrightarrow0\quad(x\to\infty),\]
where \(q_i=1/[64(1+A_i)]>0\), \(A_i\ge0\) is the coefficient of the
early size in the regular reorganization bound (Lemma 4.17), and
\(\ell_j^{\min}\) is the preassigned block-size floor.

Every entry-threshold record \(V\) has a componentwise increase \(U\)
such that every further increase \(W\ge U\) satisfies the following
conditions beyond its thresholds:

\begin{itemize}
\tightlist
\item
  the reset pricing inequality;
\item
  the cubic inequality of Lemma 4.10;
\item
  clip compatibility when \(c_{\rm dir}\le1\), the width condition and
  the regular and forced reorganization inequalities;
\item
  the inherited inequality of Lemma 4.17.
\end{itemize}

For a realized choice of levels, these conclusions require the same
stage profile, block factors and exponents, and payment of the
thresholds: \(W^{\rm entry}_{J_{m+1}}\le\hat\beta_m\) for every \(m\).
To obtain III.9.5's physical bounds, also assume \(4q_i\le c_{ji}\) and
\(4g_j\le\mathrm{rate}_{ji}\); raising thresholds does not supply these
inequalities or the early presentation. Start-up bounds admit the same
monotone choice of thresholds when \(c_{\rm dir}\le1\) at every used
stage.

\emph{Status:} theorem about numerical conditions, conditional on the
compatibility inequalities stated above.

\emph{Lean:}
\href{https://github.com/yang-mills-group/OpenAI-YM-20261028/blob/3456e5e75cf8d4cd558955a5c7638fba2d68f464/YM/AnalyticBridge/VERIFIED_RESET_STEP/Inherited.lean\#L75-L89}{\texttt{Time.beforeWait}},
\href{https://github.com/yang-mills-group/OpenAI-YM-20261028/blob/3456e5e75cf8d4cd558955a5c7638fba2d68f464/YM/Further/STRONG_STARTUP_CREDIT/Pricing.lean\#L87-L107}{\texttt{StrongStartup.prices}},
\href{https://github.com/yang-mills-group/OpenAI-YM-20261028/blob/3456e5e75cf8d4cd558955a5c7638fba2d68f464/YM/Extra/WILSON_STARTUP/Fast.lean\#L27-L38}{\texttt{Timing.choose}},
\href{https://github.com/yang-mills-group/OpenAI-YM-20261028/blob/3456e5e75cf8d4cd558955a5c7638fba2d68f464/YM/AnalyticBridge/VERIFIED_RESET_STEP/Tables.lean\#L36}{\texttt{Time.credit}}.

\textbf{Proposition III.9.8 (marker-rooted to site-pinned polymer
sums).} Fix \(M,\mu\in\mathbb N\), \(a\in\mathbb R\) and \(c\ge0\). Let
\((X_i)\) be polymers on a four-dimensional torus with weights
\(f_i\ge0\), indexed by an arbitrary set. Fix a finite set of marker
names; each polymer carries a subset of these names, the markers sitting
at sites \(o(v)\), with at most \(\mu\) markers per site. Suppose that
every polymer with \(f_i>0\) contains the site of one of its markers,
and that for every marker \(v\)
\[\sum_{i:\,v\text{ is a marker of }i}e^{(a+1)\,\mathrm{cover}_M(X_i)}f_i\le c .\]
All sums of nonnegative terms here are required to be summable,
including when the index set is uncountable. Then for every site \(x\)
\[\sum_{i:\,x\in X_i}e^{a\,\mathrm{cover}_M(X_i)}f_i\le\mu\,e\,(2(16M+9))^4\,c ,\]
uniformly in the size of the torus.

\emph{Status:} theorem. It is used for observable insertions in III.12
and III.15, not for the coupling flow.

\emph{Lean:}
\href{https://github.com/yang-mills-group/OpenAI-YM-20261028/blob/3456e5e75cf8d4cd558955a5c7638fba2d68f464/YM/Extra/MULTISOURCE_BOUNDS/Root.lean\#L161-L199}{\texttt{Extra.Multi.root\_hit}}.

\textbf{What the rest of the construction uses.}

\begin{itemize}
\tightlist
\item
  The shape, band and transfer clauses of III.9.3.
\item
  The reference coupling \(\hat\beta\) used in the later estimates.
\item
  The step program III.9.5(b), the input of the shooting.
\item
  The finer accuracy
  \(|\beta_{N,p,m}-\beta_{N,p,m+1}-\alpha_{N,m}|\le r_m\).
\item
  III.9.8.
\end{itemize}

\subsubsection{2. What it uses}\label{what-it-uses-11}

\begin{itemize}
\tightlist
\item
  \textbf{III.1 (block-spin step and interfaces).} This unit uses:

  \begin{itemize}
  \tightlist
  \item
    the step \(T^{(m)}_\beta\) with its symmetry, gauge covariance and
    probability normalization
    (\href{https://github.com/yang-mills-group/OpenAI-YM-20261028/blob/3456e5e75cf8d4cd558955a5c7638fba2d68f464/YM/RG/History.lean\#L88-L91}{\texttt{Calendar.op}},
    \href{https://github.com/yang-mills-group/OpenAI-YM-20261028/blob/3456e5e75cf8d4cd558955a5c7638fba2d68f464/YM/RG/History.lean\#L120-L128}{\texttt{BlockControl}},
    \href{https://github.com/yang-mills-group/OpenAI-YM-20261028/blob/3456e5e75cf8d4cd558955a5c7638fba2d68f464/YM/RG/Push.lean\#L37-L43}{\texttt{StepData.push\_integral}});
  \item
    the interface records for histories and selection, the datum
    \(\mathfrak D\)
    (\href{https://github.com/yang-mills-group/OpenAI-YM-20261028/blob/3456e5e75cf8d4cd558955a5c7638fba2d68f464/YM/PCalc/Band.lean\#L26-L36}{\texttt{Directed}}),
    and reset step / program;
  \item
    the conversion of a reset step into the transfer inequality
    (\href{https://github.com/yang-mills-group/OpenAI-YM-20261028/blob/3456e5e75cf8d4cd558955a5c7638fba2d68f464/YM/PCalc/Band.lean\#L57-L63}{\texttt{Directed.links}}).
    This uses \(|\beta-\beta'-\alpha|\le1\) and
    \(|\alpha+\kappa\log L|+1\le\omega\).
  \end{itemize}
\item
  \textbf{II.2 (polymer activities).} Representations, activities, their
  bounds and the large-field masks:
  \href{https://github.com/yang-mills-group/OpenAI-YM-20261028/blob/3456e5e75cf8d4cd558955a5c7638fba2d68f464/YM/RG/History.lean\#L165-L171}{\texttt{Rep}},
  \href{https://github.com/yang-mills-group/OpenAI-YM-20261028/blob/3456e5e75cf8d4cd558955a5c7638fba2d68f464/YM/RG/Activities.lean\#L106-L111}{\texttt{Activities.Valid}},
  \href{https://github.com/yang-mills-group/OpenAI-YM-20261028/blob/3456e5e75cf8d4cd558955a5c7638fba2d68f464/YM/RG/Activities.lean\#L98-L100}{\texttt{EBound}},
  \href{https://github.com/yang-mills-group/OpenAI-YM-20261028/blob/3456e5e75cf8d4cd558955a5c7638fba2d68f464/YM/RG/Activities.lean\#L102-L104}{\texttt{KBound}},
  \href{https://github.com/yang-mills-group/OpenAI-YM-20261028/blob/3456e5e75cf8d4cd558955a5c7638fba2d68f464/YM/RG/PathActions.lean\#L150-L154}{\texttt{PathData}},
  \href{https://github.com/yang-mills-group/OpenAI-YM-20261028/blob/3456e5e75cf8d4cd558955a5c7638fba2d68f464/YM/RG/PathActions.lean\#L173-L185}{\texttt{PathData.Bounds}}.
\item
  \textbf{III.2 (order of constants and shooting).} This unit uses:

  \begin{itemize}
  \tightlist
  \item
    the constants chosen before the block size, and the block sizes
    (\href{https://github.com/yang-mills-group/OpenAI-YM-20261028/blob/3456e5e75cf8d4cd558955a5c7638fba2d68f464/YM/Price/Intent.lean\#L26-L35}{\texttt{Supply.Sheet}},
    \href{https://github.com/yang-mills-group/OpenAI-YM-20261028/blob/3456e5e75cf8d4cd558955a5c7638fba2d68f464/YM/Price/Intent.lean\#L39-L48}{\texttt{Supply.Large}}).
    At each used stage, the block factor \(\ell_j\) satisfies
    \(\omega_j/\log\ell_j\le\kappa/4\) and
    \(\omega_j/\log\ell_j\le j^{-2}\).
  \item
    the record of the scale sequence having met every entry threshold
    (\href{https://github.com/yang-mills-group/OpenAI-YM-20261028/blob/3456e5e75cf8d4cd558955a5c7638fba2d68f464/YM/Price/Intent.lean\#L145-L162}{\texttt{Supply.Paid}});
  \item
    the shooting construction and the running estimate (Lemmas
    4.20--4.21).
  \end{itemize}
\item
  \textbf{III.7 (one-loop coefficient).} The kinetic data satisfy the
  planned conditions
  (\href{https://github.com/yang-mills-group/OpenAI-YM-20261028/blob/3456e5e75cf8d4cd558955a5c7638fba2d68f464/YM/PCalc/Loop.lean\#L93-L115}{\texttt{Kin.Planned}}).
  Among them:

  \begin{itemize}
  \tightlist
  \item
    \(|\alpha_{N,m}+\kappa\log L_m|+1\le\omega_{J_{m+1}}\);
  \item
    the Hessian of the fluctuation term (fine loop functional plus
    \(\tfrac12\log\det\) of the fluctuation Hessian minus the log Haar
    Jacobian) agrees with that of \(\alpha_{N,m}\) times the coarse
    quadratic action plus the coarse loop functional, with an error
    bounded by
    \(|\Lambda_m|\Gamma\Lambda e^{-c\,\mathrm{side}(\Lambda_m)}\) on
    every pair of bounded plane-wave directions
    (\href{https://github.com/yang-mills-group/OpenAI-YM-20261028/blob/3456e5e75cf8d4cd558955a5c7638fba2d68f464/YM/PCalc/Loop.lean\#L62-L71}{\texttt{FiniteBound}});
  \item
    \(P_{N,N}\) matches the Wilson energy to fourth order.
  \end{itemize}

  The bound on \(\alpha\) comes from
  \href{https://github.com/yang-mills-group/OpenAI-YM-20261028/blob/3456e5e75cf8d4cd558955a5c7638fba2d68f464/YM/Next/STATIONARY/FullPlanned.lean\#L17-L52}{\texttt{Stationary.pouring}},
  which, for fixed admissible analytic input parameters, constructs
  one-loop coefficients and constants \(D\ge0,L_{\min}\) with
  \(|\alpha+\kappa\log L|\le D\) for every odd \(L\ge L_{\min}\)
  (\href{https://github.com/yang-mills-group/OpenAI-YM-20261028/blob/3456e5e75cf8d4cd558955a5c7638fba2d68f464/YM/Next/STATIONARY/Level/Transition.lean\#L20-L23}{\texttt{HorizCalc}}).
  The constant is produced by two steps:

  \begin{itemize}
  \tightlist
  \item
    the scalar-minus-half-vector bubble equals \(2\,\mathrm{pole}(p)\),
    where
    \(\mathrm{pole}(p)=4C_G\big(p_1^2|p|^{-6}-4p_1^2p_0^2|p|^{-8}-|p|^{-4}\big)\)
    (\href{https://github.com/yang-mills-group/OpenAI-YM-20261028/blob/3456e5e75cf8d4cd558955a5c7638fba2d68f464/YM/Next/STATIONARY/OneLoop/Bubbles.lean\#L127-L141}{\texttt{haloRate}},
    \href{https://github.com/yang-mills-group/OpenAI-YM-20261028/blob/3456e5e75cf8d4cd558955a5c7638fba2d68f464/YM/Next/STATIONARY/OneLoop/Bubbles.lean\#L125}{\texttt{poleVec}});
  \item
    \((2\pi)^{-4}\int_{\mu\le|p|\le1}\mathrm{pole}=-\kappa\log\mu^{-1}\)
    (\href{https://github.com/yang-mills-group/OpenAI-YM-20261028/blob/3456e5e75cf8d4cd558955a5c7638fba2d68f464/YM/Next/STATIONARY/OneLoop/MomentumIntegral.lean\#L326-L338}{\texttt{poleSpin}}),
    with the remaining errors bounded
    (\href{https://github.com/yang-mills-group/OpenAI-YM-20261028/blob/3456e5e75cf8d4cd558955a5c7638fba2d68f464/YM/Next/STATIONARY/OneLoop/MomentumIntegral.lean\#L346-L408}{\texttt{absoluteGeode}}).
  \end{itemize}

  The angular average of \(\mathrm{pole}\) is
  \(-\tfrac{11}{3}C_G|p|^{-4}\), the classical split \(\tfrac13-4\) into
  orbital and spin parts.
\item
  \textbf{III.13 (kinetic data and history).} Finite-cutoff kinetic data
  are built by downward recursion in \(m\). The limiting one-step
  relation used to compare such data contracts with factor \(\tfrac12\)
  (\href{https://github.com/yang-mills-group/OpenAI-YM-20261028/blob/3456e5e75cf8d4cd558955a5c7638fba2d68f464/YM/UVExtension/Reward/Tides.lean\#L15-L46}{\texttt{Reward.fieldArray}},
  \href{https://github.com/yang-mills-group/OpenAI-YM-20261028/blob/3456e5e75cf8d4cd558955a5c7638fba2d68f464/YM/UVExtension/Reward/AtLimit.lean\#L46-L66}{\texttt{Reward.step\_contracted}},
  \href{https://github.com/yang-mills-group/OpenAI-YM-20261028/blob/3456e5e75cf8d4cd558955a5c7638fba2d68f464/YM/UVExtension/Reward/Tides.lean\#L49-L62}{\texttt{Reward.ties}}).
  The selected history of the main theorem is also built there.
\item
  \textbf{III.10 and III.4 (the nonlinear step).} These provide the
  hypotheses of III.9.5:

  \begin{itemize}
  \tightlist
  \item
    the identification of the blocked density with a Gaussian-sector
    integral
    (\href{https://github.com/yang-mills-group/OpenAI-YM-20261028/blob/3456e5e75cf8d4cd558955a5c7638fba2d68f464/YM/AnalyticBridge/VERIFIED_RESET_STEP/Inputs.lean\#L56-L61}{\texttt{Return.earnings}});
  \item
    its early presentation
    (\href{https://github.com/yang-mills-group/OpenAI-YM-20261028/blob/3456e5e75cf8d4cd558955a5c7638fba2d68f464/YM/AnalyticBridge/VERIFIED_RESET_STEP/Inputs.lean\#L70-L71}{\texttt{Return.atlas}});
  \item
    the local activity bounds
    (\href{https://github.com/yang-mills-group/OpenAI-YM-20261028/blob/3456e5e75cf8d4cd558955a5c7638fba2d68f464/YM/AnalyticBridge/VERIFIED_RESET_STEP/Trays.lean\#L35-L40}{\texttt{Collection.Local}}),
    continuity and recipe acceptance.
  \end{itemize}

  These are delivered by
  \href{https://github.com/yang-mills-group/OpenAI-YM-20261028/blob/3456e5e75cf8d4cd558955a5c7638fba2d68f464/YM/PhysicalNext/RG_NEXT_COMPLETE_REPEATABLE_NONLINEAR_STEP/Rain/Glazing.lean\#L97-L104}{\texttt{Rain.glazing}}
  and
  \href{https://github.com/yang-mills-group/OpenAI-YM-20261028/blob/3456e5e75cf8d4cd558955a5c7638fba2d68f464/YM/PhysicalNext/RG_NEXT_COMPLETE_REPEATABLE_NONLINEAR_STEP/Roost/Foundry.lean\#L26-L50}{\texttt{Roost.Factor}}.
\item
  \textbf{I.1 and III.16 (analytic toolkit).} Lie-group facts and Taylor
  jets of the plaquette function.
\end{itemize}

\textbf{Standard mathematics.}

\begin{itemize}
\tightlist
\item
  Intermediate value theorem, and continuity of the inverse of a
  uniformly increasing family.
\item
  Taylor's formula with uniform constants on compact sets.
\item
  Schur's lemma: an \(\mathrm{Ad}\)-invariant bilinear form on a simple
  \(\mathfrak g\) is a multiple of the inner product.
\item
  The representation theory of the hyperoctahedral group on plaquette
  orientations, used through explicit reflections and permutations.
\item
  Fourier analysis on tori, analytic matrix inversion and exponential
  decay of Fourier coefficients.
\end{itemize}

References:

\begin{itemize}
\tightlist
\item
  One-loop \(\beta\)-function: Gross--Wilczek, \emph{Phys. Rev.~Lett.}
  30 (1973) 1343; Politzer, \emph{Phys. Rev.~Lett.} 30 (1973) 1346.
\item
  The orbital/spin split \(\tfrac13-4\): Nielsen, \emph{Am. J. Phys.} 49
  (1981) 1171; Hughes, \emph{Nucl. Phys. B} 186 (1981) 376.
\item
  Lattice gauge theory: Wilson, \emph{Phys. Rev.~D} 10 (1974) 2445;
  Seiler, \emph{Lecture Notes in Physics} 159 (1982).
\item
  Block-spin renormalization of lattice gauge fields with
  coupling-constant renormalization in four dimensions: Bałaban,
  \emph{Commun. Math. Phys.} 109 (1987) 249; 116 (1988) 1; 119 (1988)
  243; 122 (1989) 175 and 355.
\item
  Exposition of Bałaban's method: Dimock, ``The renormalization group
  according to Balaban'' I--III (\emph{Rev.~Math. Phys.} 25 (2013);
  \emph{J. Math. Phys.} 54 (2013); \emph{Ann. Henri Poincaré} 15
  (2014)).
\item
  Tracking a marginal coupling and choosing the bare parameters so that
  the flow ends at a prescribed value: Bauerschmidt--Brydges--Slade,
  \emph{Introduction to a Renormalisation Group Method}, LNM 2242
  (2019).
\end{itemize}

\textbf{Imported constants and assumptions.}

\begin{itemize}
\tightlist
\item
  \textbf{The drift allowances \(\omega_j\).} These are constants fixed
  before the block size, with \(\omega_j\ge1+D\). The one-loop error
  \(D\) is fixed in III.7 and III.13, and \(\omega\) is then made small
  relative to \(\log L\) by choosing \(L\) large (III.2).
\item
  \textbf{The exponents.} The exponents \(-1/2\), \(-11/10\),
  \((\log\log)^2\) and \(1/5\) in III.9.2 are part of the definitions.
\item
  \textbf{The hypotheses of III.9.5(a).} These are carried here as
  assumptions and discharged in III.10.
\end{itemize}

\subsubsection{3. Strategy}\label{strategy-11}

The nonlinear step gives an early presentation of the blocked density.
Its marginal coefficient has three contributions: the input coupling
\(\beta'\), the one-loop shift \(\alpha_{N,m}\), and the isotropic
Hessian coefficient \(\delta\) of the early regular functions. The reset
sets \(\beta=\beta'+\alpha_{N,m}+\delta\). Subtracting flat values and
quadratic projections, adding compensating path terms, and absorbing
mask discrepancies into the forced gas preserves the full density. The
output activity bounds and \(|\delta|\le r_m\le1\), together with
\(|\alpha_{N,m}+\kappa\log L_m|+1\le\omega\), give the transfer
inequality.

Two configurations identify the reset: the flat field and a small-field
stripe \(V\) with \(\mathcal K(V)\ge|\Lambda|\eta t^2\). Matching at the
flat field fixes \(v=|\Lambda|^{-1}\log H(1)\). Matching the log-ratio
at \(V\) gives a scalar equation \(f(x,d)=0\). The pricing and
activity-motion bounds make \(f\) increasing with slope at least
\(\mathcal K(V)/2\), and give \(|f(x,0)|\le\mathcal K(V)r/2\). Its
unique root lies in \([-r,r]\) and depends continuously on \(x\). Shape
control gives equality at every configuration. For the reorganized
activities, \(f(x,d)=(d-\delta_x)\mathcal K(V)\), so the correction is
explicit.

The interval shooting chains these steps from the Wilson lattice down to
level \(0\). At each level, the transfer inequality forces the two
endpoint couplings to lie on opposite sides of
\([\bar\beta+\kappa s_m/2,\bar\beta+2\kappa s_m]\). Continuity selects a
subinterval whose couplings stay within that window and attain its
endpoints. At level \(0\), a surviving bare parameter gives
\(\beta_0=\bar\beta\). Summing the transfer errors gives the running
law, uniformly in cutoff and volume. III.13 uses compatible
finite-cutoff fixed points to obtain the centered history and its
analytic continuations.

\subsubsection{4. Steps}\label{steps-11}

\paragraph{4.1 Reference scales}\label{reference-scales}

\textbf{Lemma 4.1 (recursion of the reference coupling).}
\(s_{m+1}=s_m+\log L_m\), so
\(\hat\beta_{m+1}=\hat\beta_m+\kappa\log L_m\). \(\hat\beta\) is
nondecreasing if \(\kappa\ge0\). Since \(L_i\ge3\),
\(s_m\ge m\log3\to\infty\), and \(\hat\beta_m\to\infty\) when
\(\kappa>0\).

\emph{Sketch:} \(a_m=(L_0\cdots L_{m-1})^{-1}\) and the definitions.

\emph{Lean:}
\href{https://github.com/yang-mills-group/OpenAI-YM-20261028/blob/3456e5e75cf8d4cd558955a5c7638fba2d68f464/YM/Lattice/Scale.lean\#L124-L127}{\texttt{Schedule.scale\_succ}},
\href{https://github.com/yang-mills-group/OpenAI-YM-20261028/blob/3456e5e75cf8d4cd558955a5c7638fba2d68f464/YM/Next/SCHEDULE/Levels.lean\#L43-L45}{\texttt{Deployment.hat}},
\href{https://github.com/yang-mills-group/OpenAI-YM-20261028/blob/3456e5e75cf8d4cd558955a5c7638fba2d68f464/YM/Extra/WILSON_STARTUP/Fast.lean\#L40-L42}{\texttt{Timing.hat\_mono}},
\href{https://github.com/yang-mills-group/OpenAI-YM-20261028/blob/3456e5e75cf8d4cd558955a5c7638fba2d68f464/YM/Lattice/Scale.lean\#L117-L120}{\texttt{Schedule.divergent}}.

\paragraph{4.2 The Wilson start-up}\label{the-wilson-start-up}

\textbf{Lemma 4.2 (splitting the plaquette action).} Write
\(s(u)=|F(u)|^2/2+Q_4(F(u))+R(u)\). Near the identity, curvature gives
local coordinates on \(G\). Here \(Q_4\) is the diagonal fourth Taylor
coefficient of the plaquette function in those coordinates. The odd
coefficients vanish by inversion symmetry. The residual \(R\), defined
by the displayed identity for every \(u\in G\), has vanishing
derivatives of orders \(0,\ldots,4\) at the identity.

For \(c_{\rm dir}\le1\), the top path data \(P^{\rm top}_c\) consist of:

\begin{itemize}
\tightlist
\item
  the clipped plaquette term, which equals \(c_{\rm dir}\,s(U_p)\) in
  the small-field region;
\item
  the row term \(\sqrt{1-c_{\rm dir}}\,F_p\), contributing
  \((1-c_{\rm dir})|F_p|^2/2\);
\item
  the quartic atoms \((1-c_{\rm dir})Q_4(F_p)\);
\item
  no cubic and no loop terms.
\end{itemize}

They satisfy the zero-momentum normalization
(\(c_{\rm dir}+(1-c_{\rm dir})=1\)). On the small-field region,
\(\beta S_W=\beta\,\mathcal K_{P^{\rm top}}+\beta(1-c_{\rm dir})\sum_pR_p\).

\emph{Sketch:} The derivative of the curvature map at the identity is
the identity, so the analytic inverse function theorem gives curvature
coordinates. Inversion sends \(F\) to \(-F\) and preserves \(s\). Taylor
expansion through order \(4\) gives the quadratic and quartic terms and
the residual. The identity \(s=\tfrac12\mathrm{chord}^2\) fixes the
normalization.

\emph{Lean:}
\href{https://github.com/yang-mills-group/OpenAI-YM-20261028/blob/3456e5e75cf8d4cd558955a5c7638fba2d68f464/YM/Extra/WILSON_STARTUP/Curvature.lean\#L156-L164}{\texttt{convert\_taylor}},
\href{https://github.com/yang-mills-group/OpenAI-YM-20261028/blob/3456e5e75cf8d4cd558955a5c7638fba2d68f464/YM/Extra/WILSON_STARTUP/Curvature.lean\#L166}{\texttt{surplus}},
\href{https://github.com/yang-mills-group/OpenAI-YM-20261028/blob/3456e5e75cf8d4cd558955a5c7638fba2d68f464/YM/Extra/WILSON_STARTUP/Curvature.lean\#L168-L179}{\texttt{surplus\_jet}},
\href{https://github.com/yang-mills-group/OpenAI-YM-20261028/blob/3456e5e75cf8d4cd558955a5c7638fba2d68f464/YM/Extra/WILSON_STARTUP/Design.lean\#L34-L38}{\texttt{Design.top}},
\href{https://github.com/yang-mills-group/OpenAI-YM-20261028/blob/3456e5e75cf8d4cd558955a5c7638fba2d68f464/YM/Extra/WILSON_STARTUP/Design.lean\#L156-L162}{\texttt{Design.principal\_top}}.

\textbf{Lemma 4.3 (exact start).} For a compatible pair \((c,b)\), every
\(\beta\) and every \(U\): the path-data channels of \(P^{\rm top}\) at
\(\beta\), plus the regular activities (residuals of small-field
plaquettes), plus the large-field remainders, add up to
\(\beta S_W(U)\), the sum being averaged over the point group. Moreover
the gas equals \(\exp(-\beta\cdot\text{(large-field remainder)})\).
Hence \(\rho^{\rm start}_\beta=e^{-\beta S_W}\), which is III.9.6(1).

\emph{Sketch:} Fix \(U\). The only nonzero gas term is the family of
connected components of the large-field region for the padded adjacency
relation (cover size \(B\ge2\) and the compatibility condition separate
the padded components). So the gas collapses to a single product
\(\prod_De^{-\beta\,\mathrm{charge}(D,U)}\). The charges sum to the
large-field remainder, and the channel identity accounts for the rest.

\emph{Lean:}
\href{https://github.com/yang-mills-group/OpenAI-YM-20261028/blob/3456e5e75cf8d4cd558955a5c7638fba2d68f464/YM/Extra/WILSON_STARTUP/Channels.lean\#L72-L87}{\texttt{channel\_split}},
\href{https://github.com/yang-mills-group/OpenAI-YM-20261028/blob/3456e5e75cf8d4cd558955a5c7638fba2d68f464/YM/Extra/WILSON_STARTUP/Gas.lean\#L81-L94}{\texttt{Start.gas\_raw}},
\href{https://github.com/yang-mills-group/OpenAI-YM-20261028/blob/3456e5e75cf8d4cd558955a5c7638fba2d68f464/YM/Extra/WILSON_STARTUP/Gas.lean\#L70-L75}{\texttt{Start.issues}},
\href{https://github.com/yang-mills-group/OpenAI-YM-20261028/blob/3456e5e75cf8d4cd558955a5c7638fba2d68f464/YM/Extra/WILSON_STARTUP/Gas.lean\#L96-L103}{\texttt{Start.evaluates}}.

\textbf{Lemma 4.4 (one-plaquette residual).} This is III.9.6(2).

\emph{Sketch:} \(R_p\) is a real-analytic function of the four link
matrices of \(p\). On a small-field plaquette (chord \(\le t\)) its
derivatives of order \(l\le k\) along directions of size
\(\le t(1+4)^q\) are bounded by \(C_lt^5\). This is Taylor's formula
with the flatness of Lemma 4.2 and compactness of \(G\); the factor
\((1+4)^q\) is the spatial weight of sites within distance \(4\). Take
\(C=1+\sum_lC_l\) and the smallest \(\delta\).

\emph{Lean:}
\href{https://github.com/yang-mills-group/OpenAI-YM-20261028/blob/3456e5e75cf8d4cd558955a5c7638fba2d68f464/YM/Extra/WILSON_STARTUP/Face.lean\#L76-L101}{\texttt{residual\_price}},
\href{https://github.com/yang-mills-group/OpenAI-YM-20261028/blob/3456e5e75cf8d4cd558955a5c7638fba2d68f464/YM/Extra/WILSON_STARTUP/Face.lean\#L62-L68}{\texttt{residual\_jet}},
\href{https://github.com/yang-mills-group/OpenAI-YM-20261028/blob/3456e5e75cf8d4cd558955a5c7638fba2d68f464/YM/Extra/WILSON_STARTUP/Face.lean\#L70-L74}{\texttt{healthy\_chord}},
\href{https://github.com/yang-mills-group/OpenAI-YM-20261028/blob/3456e5e75cf8d4cd558955a5c7638fba2d68f464/YM/Extra/WILSON_STARTUP/Residual.lean\#L161-L183}{\texttt{Quad.small\_jet}}.

\textbf{Lemma 4.5 (eventual start-up bounds).} Under the hypotheses of
III.9.6(3), the start-up bounds hold beyond a stage-dependent threshold,
uniformly in \(N,p\) and \(\beta\in\mathfrak B_N\). With the top-path
bounds and recipe approval specified there, the start-up curve is
admissible.

\emph{Sketch:} Put \(x=\hat\beta_N\) and \(M=e^{(\log\log x)^2}\). The
regular bound is a stage-dependent constant times
\(2xt_N^5=2x^{-3/2}M^5=o(x^{-11/10})\). For the forced terms, the proof
uses a decay exponent \(M\) and an additional polymer-counting factor
\(e^\gamma-1\), with \(\gamma=e^{-M/2}\). The connected-polymer estimate
bounds the forced norm by \(e^{-M+1}\le e^{-M^{1/5}}\) for large \(x\).
Compatibility follows from \(B_N\asymp(\log x)^2\) and
\(t_N=x^{-1/2}M\). These estimates give the displayed continuous
start-up family; admissibility also uses its path bounds and the recipe
tests.

\emph{Lean:}
\href{https://github.com/yang-mills-group/OpenAI-YM-20261028/blob/3456e5e75cf8d4cd558955a5c7638fba2d68f464/YM/Extra/WILSON_STARTUP/Clock.lean\#L113-L139}{\texttt{Timing.enough}},
\href{https://github.com/yang-mills-group/OpenAI-YM-20261028/blob/3456e5e75cf8d4cd558955a5c7638fba2d68f464/YM/Extra/WILSON_STARTUP/Ink.lean\#L104-L122}{\texttt{Start.ink}},
\href{https://github.com/yang-mills-group/OpenAI-YM-20261028/blob/3456e5e75cf8d4cd558955a5c7638fba2d68f464/YM/Extra/WILSON_STARTUP/Funds.lean\#L130-L143}{\texttt{Start.Loan.kernel}},
\href{https://github.com/yang-mills-group/OpenAI-YM-20261028/blob/3456e5e75cf8d4cd558955a5c7638fba2d68f464/YM/Planned/RECIPE_STARTUP_PORT/Layout.lean\#L52-L71}{\texttt{Planned.Port.card}}.

\textbf{Lemma 4.6 (installing the Wilson top).} Assume the old fine
paths satisfy the stationary, positive-Hessian, initial-guess,
classical-matching and finite-period one-loop conditions for a step. Let
the new fine paths have the same main-action jets through order \(4\),
and suppose both fine loop kernels vanish. With the same local Haar
chart and block branches, there are new stationary data satisfying those
same conditions, with the same coarse paths and one-loop coefficient.
With the revised horizon and top-path bounds of III.9.6(4), this
replacement gives planned kinetic data at every cutoff, preserving the
stationary data for \(m+1<N\).

\emph{Sketch:} Equality of fourth jets gives equality of the fluctuation
Hessians at the flat field and of their determinant jets through order
\(2\). The analytic implicit function theorem constructs the new
stationary branch; finite-order agreement transfers the initial guess,
classical match and one-loop bound. The old and Wilson top actions both
match the Wilson energy through order \(4\). Only the step \(N\to N-1\)
changes.

\emph{Lean:}
\href{https://github.com/yang-mills-group/OpenAI-YM-20261028/blob/3456e5e75cf8d4cd558955a5c7638fba2d68f464/YM/Extra/WILSON_STARTUP/TopStep.lean\#L160-L180}{\texttt{Change.top\_step}},
\href{https://github.com/yang-mills-group/OpenAI-YM-20261028/blob/3456e5e75cf8d4cd558955a5c7638fba2d68f464/YM/Extra/WILSON_STARTUP/Engine.lean\#L62-L82}{\texttt{Change.engine}}.

\paragraph{4.3 The scalar response}\label{the-scalar-response}

\textbf{Lemma 4.7 (sensor).} Assume \(\mathfrak g\ne\{0\}\), fix a clip
with \(A=c_{\rm dir}>0\), and let \(C\ge0\). Then there is \(\eta>0\),
depending only on these data, with the following property. For every
even torus side, every path datum satisfying the path bounds with cap
\(C\), and every set of norm parameters with \(0<t\le1\), there is a
configuration \(V\) that is small-field everywhere and satisfies
\(\mathcal K(V)\ge|\Lambda|\,\eta\,t^2\).

\emph{Sketch:}

\begin{itemize}
\tightlist
\item
  Choose \(u\in G\) with chord \(h=\tau t\), where
  \(\tau=\min(\varepsilon_G,\tfrac14,A/(512(1+C)))\); such \(u\) exists
  by the intermediate value theorem along a one-parameter subgroup.
\item
  Put \(u\) on the axis-\(0\) links at sites whose coordinate \(1\) is
  even, and \(1\) elsewhere. Every \((0,1)\)-plaquette has chord \(h\),
  and every plaquette has chord \(\le4h\le t\).
\item
  The quadratic term contributes \(\ge|\Lambda|Ah^2/2\). The cubic and
  quartic atoms contribute at least
  \(-|\Lambda|(C(4h)^3+C(4h)^4)\ge-|\Lambda|Ah^2/4\).
\item
  The point-group average permutes plaquettes and preserves these
  bounds. Hence \(\eta=A\tau^2/4\).
\end{itemize}

\emph{Lean:}
\href{https://github.com/yang-mills-group/OpenAI-YM-20261028/blob/3456e5e75cf8d4cd558955a5c7638fba2d68f464/YM/UVNext/COUPLING_RESET/Probe.lean\#L96-L142}{\texttt{CouplingReset.sensor}},
\href{https://github.com/yang-mills-group/OpenAI-YM-20261028/blob/3456e5e75cf8d4cd558955a5c7638fba2d68f464/YM/UVNext/COUPLING_RESET/Probe.lean\#L40-L54}{\texttt{reachable\_chord}}.

\textbf{Lemma 4.8 (monotone response, continuous root).} Let \(I=[l,u]\)
with \(l\le u\), let \(r\ge0\), and let \(f:I\times[-r,r]\to\mathbb R\)
satisfy:

\begin{itemize}
\tightlist
\item
  there is \(\mu>0\) such that \(f(x,b)-f(x,a)\ge\mu(b-a)\) for every
  \(x\in I\) and \(-r\le a\le b\le r\);
\item
  \(f(x,\cdot)\) is continuous for every \(x\in I\), and \(f(\cdot,a)\)
  is continuous for every \(a\in[-r,r]\);
\item
  \(|f(x,0)|\le\mu r\) for every \(x\in I\).
\end{itemize}

Then there is a unique \(d(x)\in[-r,r]\) with \(f(x,d(x))=0\). The
function \(d\) is continuous, and \(|d(x)-a|\le|f(x,a)|/\mu\) for every
\(x\in I\) and \(a\in[-r,r]\).

\emph{Sketch:} \(f(x,-r)\le0\le f(x,r)\), so the intermediate value
theorem gives a zero, unique by strict monotonicity. Moreover
\(|d(y)-d(x)|\le|f(y,d(x))-f(x,d(x))|/\mu\to0\).

\emph{Lean:}
\href{https://github.com/yang-mills-group/OpenAI-YM-20261028/blob/3456e5e75cf8d4cd558955a5c7638fba2d68f464/YM/UVNext/COUPLING_RESET/Root.lean\#L27-L31}{\texttt{Response}},
\href{https://github.com/yang-mills-group/OpenAI-YM-20261028/blob/3456e5e75cf8d4cd558955a5c7638fba2d68f464/YM/UVNext/COUPLING_RESET/Root.lean\#L80-L83}{\texttt{Response.balance\_solve}},
\href{https://github.com/yang-mills-group/OpenAI-YM-20261028/blob/3456e5e75cf8d4cd558955a5c7638fba2d68f464/YM/UVNext/COUPLING_RESET/Root.lean\#L54-L58}{\texttt{Response.continuous}},
\href{https://github.com/yang-mills-group/OpenAI-YM-20261028/blob/3456e5e75cf8d4cd558955a5c7638fba2d68f464/YM/UVNext/COUPLING_RESET/Root.lean\#L60-L70}{\texttt{solve}}.

\textbf{Lemma 4.9 (two-point readout).} Fix path data, a clip, norm
parameters and activities \(F(x,d)\) whose regular functions vanish at
the flat configuration. Let \(V\) be small-field everywhere and suppose
\(H_x(1),H_x(V)>0\). Define
\[f(x,d)=(\beta'_x+\alpha+d)\,\mathcal K(V)+\mathcal L(V)+\mathcal E_{F(x,d)}(V)-\log\frac{H_x(1)}{H_x(V)} .\]
Install the trial with coupling \(\beta'_x+\alpha+d(x)\) and vacuum
constant \(v_x=|\Lambda|^{-1}\log H_x(1)\). If \(f(x,d(x))=0\), the
trial equals \(H_x\) at \(1\) and at \(V\). If moreover the template
satisfies shape control (III.9.4(R)), it equals \(H_x\) at every
configuration.

\emph{Sketch:} At \(U=1\) the kinetic, loop and activity terms vanish
and the gas is \(1\). At small-field \(V\) the gas is \(1\). Shape
control at a zero of \(f\) gives \(\mathrm{shape}=H_x/H_x(1)\).

\emph{Lean:}
\href{https://github.com/yang-mills-group/OpenAI-YM-20261028/blob/3456e5e75cf8d4cd558955a5c7638fba2d68f464/YM/UVNext/COUPLING_RESET/Reset.lean\#L64-L65}{\texttt{Template.residual}},
\href{https://github.com/yang-mills-group/OpenAI-YM-20261028/blob/3456e5e75cf8d4cd558955a5c7638fba2d68f464/YM/UVNext/COUPLING_RESET/Reset.lean\#L99-L105}{\texttt{Template.passes}},
\href{https://github.com/yang-mills-group/OpenAI-YM-20261028/blob/3456e5e75cf8d4cd558955a5c7638fba2d68f464/YM/UVNext/COUPLING_RESET/Reset.lean\#L119-L125}{\texttt{Template.realizes}},
\href{https://github.com/yang-mills-group/OpenAI-YM-20261028/blob/3456e5e75cf8d4cd558955a5c7638fba2d68f464/YM/UVNext/COUPLING_RESET/Reset.lean\#L131-L140}{\texttt{Template.response}}.

\textbf{Lemma 4.10 (surplus from a pre-reset presentation).} Suppose
that on \(V=1\) and on every small-field \(V\)
\[H_x(V)=\exp\big(u_x|\Lambda_m|-\beta'_x\,\mathcal M(V)-\alpha_{N,m}\,\mathcal S(V)-\mathcal L(V)-\mathcal E_x(V)\big),\]
where \(\mathcal M\) is the unmasked main term (quadratic, cubic and
quartic), \(\mathcal S\) its quadratic part, and \(\mathcal E_x\) has
size \(\le c_m\varepsilon_m/4\). Assume the coarse path bounds with cap
\(C\) and the cubic condition
\(|\alpha_{N,m}|\,2Ct_m^3\le c_m\varepsilon_m/2\). Then the surplus
inequality of III.9.4(I) holds. If, in addition, the output activities
satisfy the catalogue and slow-variation assumptions of III.9.4(I), then
all of (I) holds.

\emph{Sketch:}

\begin{itemize}
\tightlist
\item
  On small-field \(V\), \(\mathcal K(V)=\mathcal M(V)\) and the loop
  functional is unmasked.
\item
  \(|\mathcal K(V)-\mathcal S(V)|\le2C|\Lambda|t^3\), since the cubic
  and quartic atoms are \(O(t^3)\) per site.
\item
  The regular activities at \(1\) and at \(V\) cost
  \(2\cdot|\Lambda|c\varepsilon/4\).
\item
  The cubic condition absorbs \(|\alpha|\,|\mathcal K-\mathcal S|\). It
  holds eventually because \(|\alpha|\le\omega+\kappa\log L\) and
  \(t^3\ll\hat\beta^{-11/10}\).
\end{itemize}

\emph{Lean:}
\href{https://github.com/yang-mills-group/OpenAI-YM-20261028/blob/3456e5e75cf8d4cd558955a5c7638fba2d68f464/YM/UVNext/COUPLING_RESET/Regular.lean\#L91-L96}{\texttt{Prereset.income}},
\href{https://github.com/yang-mills-group/OpenAI-YM-20261028/blob/3456e5e75cf8d4cd558955a5c7638fba2d68f464/YM/UVNext/COUPLING_RESET/Regular.lean\#L73-L87}{\texttt{Prereset.value\_income}},
\href{https://github.com/yang-mills-group/OpenAI-YM-20261028/blob/3456e5e75cf8d4cd558955a5c7638fba2d68f464/YM/UVNext/COUPLING_RESET/Precision.lean\#L124-L143}{\texttt{pays\_surplus}},
\href{https://github.com/yang-mills-group/OpenAI-YM-20261028/blob/3456e5e75cf8d4cd558955a5c7638fba2d68f464/YM/UVNext/COUPLING_RESET/Precision.lean\#L52-L69}{\texttt{higher\_estimate}},
\href{https://github.com/yang-mills-group/OpenAI-YM-20261028/blob/3456e5e75cf8d4cd558955a5c7638fba2d68f464/YM/UVNext/COUPLING_RESET/Precision.lean\#L96-L98}{\texttt{CubicGate}},
\href{https://github.com/yang-mills-group/OpenAI-YM-20261028/blob/3456e5e75cf8d4cd558955a5c7638fba2d68f464/YM/UVNext/COUPLING_RESET/Precision.lean\#L89-L94}{\texttt{within\_alpha}}.

\paragraph{4.4 The reset step}\label{the-reset-step}

\textbf{Lemma 4.11 (one step).} This is III.9.4.

\emph{Sketch:}

\begin{itemize}
\tightlist
\item
  (P) and (V) give
  \(|\Lambda_m|(c_m+1)\varepsilon_m\le\tfrac12\mathcal K(V)r_m\).
\item
  Slow motion in \(d\) gives slope \(\ge\mu:=\tfrac12\mathcal K(V)\).
\item
  The surplus inequality and
  \(|\mathcal E(V)|\le|\Lambda|\varepsilon_m\) give
  \(|f(x,0)|\le\mu r_m\).
\item
  \(f\) is continuous in \(x\) by Lemma 4.15.
\item
  Lemma 4.8 gives \(d\); Lemma 4.9 installs the trial; (R) gives the
  identity at all configurations, the forced bound and recipe
  acceptance.
\item
  \(|\beta-\beta'-\alpha|=|d|\le r_m\le1\).
\end{itemize}

\emph{Lean:}
\href{https://github.com/yang-mills-group/OpenAI-YM-20261028/blob/3456e5e75cf8d4cd558955a5c7638fba2d68f464/YM/UVNext/COUPLING_RESET/Advance.lean\#L157-L173}{\texttt{correct}},
\href{https://github.com/yang-mills-group/OpenAI-YM-20261028/blob/3456e5e75cf8d4cd558955a5c7638fba2d68f464/YM/UVNext/COUPLING_RESET/Advance.lean\#L175-L189}{\texttt{advance}},
\href{https://github.com/yang-mills-group/OpenAI-YM-20261028/blob/3456e5e75cf8d4cd558955a5c7638fba2d68f464/YM/UVNext/COUPLING_RESET/Advance.lean\#L121-L137}{\texttt{Income.paidResponse}},
\href{https://github.com/yang-mills-group/OpenAI-YM-20261028/blob/3456e5e75cf8d4cd558955a5c7638fba2d68f464/YM/UVNext/COUPLING_RESET/Advance.lean\#L139-L152}{\texttt{Income.balance}}.

\textbf{Lemma 4.12 (the marginal coefficient of the activities).} Let
\(f\) be a local, gauge-invariant regular function anchored at \(o\),
smooth through the prescribed derivative order, and let \(\mathfrak g\)
be simple. The second derivative of \(f\) at the trivial configuration,
along the canonical affine link test fields centered at \(o\), with
curls \(X\) in orientation \(b\) and \(Y\) in orientation \(b'\), equals
\(\mathrm{coeff}_f(b,b')\langle X,Y\rangle\) (Schur). For the projection
used below, retain this Hessian coefficient when the support enlarged by
the radius-\(2\) halo of \(o\) does not wrap around the torus, and set
the projected coefficient to \(0\) otherwise.

For a family \(A\) of such functions with an anchored regular norm
bound, whose addresses and values transform under translations and the
point group, assume the torus side is at least \(4\) and put
\[\delta(A)=\sum_{f\text{ anchored at }0}q_f(01,01).\] Here \(q_f\)
denotes the projected coefficient just defined. Then:

\begin{itemize}
\tightlist
\item
  for every site \(o\),
  \(\sum_{f\text{ anchored at }o}q_f(b,b')=\delta(A)\,\mathbf 1_{b=b'}\);
\item
  the sum of the projected Hessian forms at \(o\) is
  \(\delta(A)\sum_b|w_b|^2\). Its quadratic energy is
  \(\tfrac12\delta(A)\sum_b|w_b|^2\);
\item
  if \(A\) has regular size \(c\), then \(t^2|\delta(A)|\le c\).
\end{itemize}

\emph{Sketch:} Gauge invariance makes each color bilinear form invariant
under the adjoint action, hence scalar by simplicity. For each pair of
distinct plaquette orientations, reflecting an axis belonging to exactly
one of them reverses the corresponding off-diagonal coefficient; its sum
therefore vanishes. Signed permutations carry each diagonal coefficient
to \((01,01)\). The bound follows from the second derivative in a
direction of size \(t\).

\emph{Lean:}
\href{https://github.com/yang-mills-group/OpenAI-YM-20261028/blob/3456e5e75cf8d4cd558955a5c7638fba2d68f464/YM/AnalyticBridge/VERIFIED_RESET_STEP/Isotropy.lean\#L75}{\texttt{Averaging.delta}},
\href{https://github.com/yang-mills-group/OpenAI-YM-20261028/blob/3456e5e75cf8d4cd558955a5c7638fba2d68f464/YM/AnalyticBridge/VERIFIED_RESET_STEP/Isotropy.lean\#L77-L98}{\texttt{Averaging.isotropic}},
\href{https://github.com/yang-mills-group/OpenAI-YM-20261028/blob/3456e5e75cf8d4cd558955a5c7638fba2d68f464/YM/AnalyticBridge/VERIFIED_RESET_STEP/Isotropy.lean\#L100-L102}{\texttt{Averaging.identity}},
\href{https://github.com/yang-mills-group/OpenAI-YM-20261028/blob/3456e5e75cf8d4cd558955a5c7638fba2d68f464/YM/AnalyticBridge/VERIFIED_RESET_STEP/Isotropy.lean\#L103-L112}{\texttt{Averaging.delta\_bound}},
\href{https://github.com/yang-mills-group/OpenAI-YM-20261028/blob/3456e5e75cf8d4cd558955a5c7638fba2d68f464/YM/AnalyticBridge/VERIFIED_RESET_STEP/Projection.lean\#L34-L40}{\texttt{label\_color}},
\href{https://github.com/yang-mills-group/OpenAI-YM-20261028/blob/3456e5e75cf8d4cd558955a5c7638fba2d68f464/YM/AnalyticBridge/VERIFIED_RESET_STEP/Projection.lean\#L89-L90}{\texttt{projection}}.

\textbf{Lemma 4.13 (exact reorganization).} Assume the early
presentation, local bounds, compatible clip, width, reorganization
inequalities and recipe acceptance of III.9.5(a). Let \(\delta_x\) be
the coefficient of Lemma 4.12. The reorganized representation has the
same paths, coupling \(\beta'_x+\alpha_{N,m}+\delta_x\), and vacuum
constant \(\tau_x-|\Lambda_m|^{-1}\sum_iE_{x,i}(1)\). Its density equals
\(H_x\) at every configuration, and it obeys the level-\(m\) allowances
and the forced bound \(e^{-gM_{m+1}^2}\).

\emph{Sketch:} For an early regular function \(E_i\) with support
\(X_i\) and anchor \(o_i\), put \(Y_i=X_i\cup\mathrm{halo}_2\{o_i\}\).
Let \(Q_i\) be one half of its projected Hessian applied to the local
clover curvatures, and set \(E_i^0=E_i-E_i(1)-Q_i\). Define \(D_i\) to
vanish when \(Y_i\) is small-field, and otherwise by
\[D_i=E_i(1)+\chi_{\{o_i\}}Q_i-\chi_{X_i}E_i.\] Then, at every
configuration,
\[\chi_{X_i}E_i=E_i(1)+\chi_{\{o_i\}}Q_i-D_i+\chi_{Y_i}E_i^0.\] The
projected terms sum to \(\delta_x/2\) times the small-field
clover-curvature square. The regular terms also acquire \(\delta_x\)
times the difference between each quadratic clover proxy and its
original quadratic path term, and \(-(\alpha_{N,m}+\delta_x)\) times
each masked cubic or quartic path term. The remaining mask discrepancies
become logarithmic factors in the forced gas.

More precisely, let \(\mathcal E^{\rm new}_x\) denote the resulting
regular sum and \(\ell_x\) the sum of these logarithmic factors before
they are absorbed into the gas. The activity and path splitting
identities give
\[(\alpha_{N,m}+\delta_x)\mathcal K+\mathcal E^{\rm new}_x-\ell_x
=\alpha_{N,m}\mathcal S+\mathcal E^{\rm early}_x-\sum_iE_{x,i}(1).\]
Absorbing \(e^{\ell_x}\) into the original forced gas proves equality of
the full densities, including at configurations where masks fail. The
reorganization estimates bound the regular terms and the added forced
factors; the numerical closure conditions give the claimed allowances.
Recipe acceptance is a separate hypothesis.

The trial activities are independent of the correction parameter \(d\).
Thus exact equality at \(d=\delta_x\) gives
\(f(x,d)=(d-\delta_x)\mathcal K(V)\) and shape control on
\([-r_m,r_m]\). The unique root of Lemma 4.11 is therefore \(\delta_x\).

\emph{Lean:}
\href{https://github.com/yang-mills-group/OpenAI-YM-20261028/blob/3456e5e75cf8d4cd558955a5c7638fba2d68f464/YM/AnalyticBridge/VERIFIED_RESET_STEP/Evaluation.lean\#L92-L95}{\texttt{Collection.Local.actual}},
\href{https://github.com/yang-mills-group/OpenAI-YM-20261028/blob/3456e5e75cf8d4cd558955a5c7638fba2d68f464/YM/AnalyticBridge/VERIFIED_RESET_STEP/Evaluation.lean\#L70-L82}{\texttt{Local.exponent}},
\href{https://github.com/yang-mills-group/OpenAI-YM-20261028/blob/3456e5e75cf8d4cd558955a5c7638fba2d68f464/YM/AnalyticBridge/VERIFIED_RESET_STEP/Payment.lean\#L77-L84}{\texttt{Return.settled}},
\href{https://github.com/yang-mills-group/OpenAI-YM-20261028/blob/3456e5e75cf8d4cd558955a5c7638fba2d68f464/YM/AnalyticBridge/VERIFIED_RESET_STEP/Inputs.lean\#L74-L94}{\texttt{Return.grants}},
\href{https://github.com/yang-mills-group/OpenAI-YM-20261028/blob/3456e5e75cf8d4cd558955a5c7638fba2d68f464/YM/AnalyticBridge/VERIFIED_RESET_STEP/Advance.lean\#L37-L42}{\texttt{Return.earns}},
\href{https://github.com/yang-mills-group/OpenAI-YM-20261028/blob/3456e5e75cf8d4cd558955a5c7638fba2d68f464/YM/AnalyticBridge/VERIFIED_RESET_STEP/Advance.lean\#L44-L61}{\texttt{Return.remainders}},
\href{https://github.com/yang-mills-group/OpenAI-YM-20261028/blob/3456e5e75cf8d4cd558955a5c7638fba2d68f464/YM/AnalyticBridge/VERIFIED_RESET_STEP/Advance.lean\#L30-L36}{\texttt{Return.physical\_shape}},
\href{https://github.com/yang-mills-group/OpenAI-YM-20261028/blob/3456e5e75cf8d4cd558955a5c7638fba2d68f464/YM/AnalyticBridge/VERIFIED_RESET_STEP/Advance.lean\#L64-L90}{\texttt{Return.reset}},
\href{https://github.com/yang-mills-group/OpenAI-YM-20261028/blob/3456e5e75cf8d4cd558955a5c7638fba2d68f464/YM/AnalyticBridge/VERIFIED_RESET_STEP/Residual.lean\#L114-L120}{\texttt{Residual.splits}},
\href{https://github.com/yang-mills-group/OpenAI-YM-20261028/blob/3456e5e75cf8d4cd558955a5c7638fba2d68f464/YM/AnalyticBridge/VERIFIED_RESET_STEP/Rewrites.lean\#L54-L59}{\texttt{Local.first\_split}},
\href{https://github.com/yang-mills-group/OpenAI-YM-20261028/blob/3456e5e75cf8d4cd558955a5c7638fba2d68f464/YM/AnalyticBridge/VERIFIED_RESET_STEP/Rewrites.lean\#L77-L93}{\texttt{Collection.path\_split}}.
\emph{(proof read, including the reorganized activities in
\href{https://github.com/yang-mills-group/OpenAI-YM-20261028/blob/3456e5e75cf8d4cd558955a5c7638fba2d68f464/YM/AnalyticBridge/VERIFIED_RESET_STEP/Investor.lean\#L112-L128}{\texttt{Local.renew}})}

\textbf{Lemma 4.14 (from per-step inputs to the program).} This is
III.9.5(b). The numerical conditions
(\href{https://github.com/yang-mills-group/OpenAI-YM-20261028/blob/3456e5e75cf8d4cd558955a5c7638fba2d68f464/YM/AnalyticBridge/VERIFIED_RESET_STEP/Downstream.lean\#L41-L49}{\texttt{Use.Machine}}),
the existence of early presentations along every admissible input curve
(\href{https://github.com/yang-mills-group/OpenAI-YM-20261028/blob/3456e5e75cf8d4cd558955a5c7638fba2d68f464/YM/AnalyticBridge/VERIFIED_RESET_STEP/Downstream.lean\#L77-L81}{\texttt{Use.Offering}})
and the start-up curves
(\href{https://github.com/yang-mills-group/OpenAI-YM-20261028/blob/3456e5e75cf8d4cd558955a5c7638fba2d68f464/YM/AnalyticBridge/VERIFIED_RESET_STEP/Downstream.lean\#L83-L84}{\texttt{Use.Origins}})
together give the step program, assuming \(G\) is simple and the machine
uses the compatible recipe and physical bounds of III.9.5.

\emph{Sketch:} For each input curve, Lemma 4.7 provides a probe. Lemma
4.13 provides (I) and (R). Lemma 4.11 then gives the output curve and
the bound \(|\beta-\beta'-\alpha|\le r_m\).

\emph{Lean:}
\href{https://github.com/yang-mills-group/OpenAI-YM-20261028/blob/3456e5e75cf8d4cd558955a5c7638fba2d68f464/YM/AnalyticBridge/VERIFIED_RESET_STEP/Downstream.lean\#L85-L88}{\texttt{Use.supplier}},
\href{https://github.com/yang-mills-group/OpenAI-YM-20261028/blob/3456e5e75cf8d4cd558955a5c7638fba2d68f464/YM/AnalyticBridge/VERIFIED_RESET_STEP/Downstream.lean\#L71-L74}{\texttt{Machine.packet}},
\href{https://github.com/yang-mills-group/OpenAI-YM-20261028/blob/3456e5e75cf8d4cd558955a5c7638fba2d68f464/YM/UVNext/COUPLING_RESET/Stability.lean\#L41-L47}{\texttt{Supplier}},
\href{https://github.com/yang-mills-group/OpenAI-YM-20261028/blob/3456e5e75cf8d4cd558955a5c7638fba2d68f464/YM/UVNext/COUPLING_RESET/Stability.lean\#L67-L70}{\texttt{program}}.

\paragraph{4.5 Continuity of the
transport}\label{continuity-of-the-transport}

\textbf{Lemma 4.15 (joint continuity of the blocked density).} Let
\(x\mapsto\rho'_x\) be a continuous curve of lawful representations with
fixed path data, clip and norm parameters, and assume the block kernel
satisfies the branch conditions of III.1. Then
\((x,V)\mapsto(T^{(m)}_{\beta'_x}\rho'_x)(V)\) is jointly continuous,
and \(x\mapsto\rho'_x(U)\) is continuous for every \(U\).

\emph{Sketch:} Continuity of the activities in the polymer norm gives
uniform convergence of the additive part and of the gas in \(U\). This
yields a local uniform bound on \(\rho'_y\), and dominated convergence
applies under the kernel. The kernel itself is continuous in \(\beta\)
and in the coarse configuration.

\emph{Lean:}
\href{https://github.com/yang-mills-group/OpenAI-YM-20261028/blob/3456e5e75cf8d4cd558955a5c7638fba2d68f464/YM/Push/BETA_TRANSPORT/Curve.lean\#L122-L131}{\texttt{FixedRepData.continuous\_push}},
\href{https://github.com/yang-mills-group/OpenAI-YM-20261028/blob/3456e5e75cf8d4cd558955a5c7638fba2d68f464/YM/Push/BETA_TRANSPORT/Curve.lean\#L93-L97}{\texttt{FixedRepData.continuous\_eval}},
\href{https://github.com/yang-mills-group/OpenAI-YM-20261028/blob/3456e5e75cf8d4cd558955a5c7638fba2d68f464/YM/UVNext/COUPLING_RESET/Catalog.lean\#L39-L45}{\texttt{Catalog.stable}}.

\paragraph{4.6 Entry thresholds}\label{entry-thresholds}

\textbf{Lemma 4.16 (monotone thresholds).} Let \(T_j(x)\) be properties,
one for each stage \(j\), each holding for all large \(x\). Then any
record \(u_j\) can be raised to a nondecreasing record \(m_j\ge u_j\)
such that \(T_j(x)\) holds for all \(x\ge m_j\). For each fine stage
\(j\), combine the two possible coarse targets \(i=j\) and
\(i=\max(j-1,0)\) first.

\emph{Sketch:} Take \(m_j\) to be the prefix maximum of
\(\max(\text{threshold}_i,u_i)\) over \(i\le j\).

\emph{Lean:}
\href{https://github.com/yang-mills-group/OpenAI-YM-20261028/blob/3456e5e75cf8d4cd558955a5c7638fba2d68f464/YM/Push/PRICE_COMPILER/Basic.lean\#L111-L118}{\texttt{collect\_upper}},
\href{https://github.com/yang-mills-group/OpenAI-YM-20261028/blob/3456e5e75cf8d4cd558955a5c7638fba2d68f464/YM/Push/PRICE_COMPILER/Basic.lean\#L126-L131}{\texttt{eventually\_transitions}},
\href{https://github.com/yang-mills-group/OpenAI-YM-20261028/blob/3456e5e75cf8d4cd558955a5c7638fba2d68f464/YM/Push/PRICE_COMPILER/Basic.lean\#L27}{\texttt{StepStages}},
\href{https://github.com/yang-mills-group/OpenAI-YM-20261028/blob/3456e5e75cf8d4cd558955a5c7638fba2d68f464/YM/Push/PRICE_COMPILER/Post.lean\#L34-L38}{\texttt{DelayLe}}.

\textbf{Lemma 4.17 (eventual reset inequalities).} Fix an allowed
transition from fine stage \(j\) to coarse stage \(i\), its block factor
\(L\), and all the constants in III.9.7. Write \(x=\hat\beta_m\),
\(y=x+\kappa\log L\), \(M=e^{(\log\log x)^2}\) and \(q_i\) for the
early-activity fraction. The conclusions are uniform for
\(|\alpha|\le\omega_j+\kappa\log L\) and \(\beta\in[y/2,2y]\). Then all
the delay-dependent reset conditions hold for sufficiently large \(x\).

\emph{Sketch:} The sufficient pricing condition \(2(c+1)\le\eta rM^2\)
holds because \(M\to\infty\). Multiplication by \(x^{-1}\) gives
\(2(c+1)x^{-1}\le\eta rt^2\); since \(\varepsilon=x^{-11/10}\le x^{-1}\)
for \(x\ge1\), this implies the required pricing inequality. The cubic
condition follows from \(t^3/\varepsilon=x^{-2/5}M^3\to0\).

For the local reorganization, its regular bound has the form
\(A_iH+B_i|\alpha|t^3\), with \(A_i,B_i\) fixed before \(L\). The choice
\(H=q_i\varepsilon\), with \(A_iq_i\le1/64\), leaves room for an output
bound \(\varepsilon/32\). The forced estimate controls both the early
forced terms and the added logarithmic factors, at support exponent
\(a+1\), where \(a\) is the coarse support exponent. The width condition
is \(4(a+1)\le\sigma B\), where \(\sigma\) is the path-decay width. It
holds since \(B\to\infty\), as does clip compatibility. The
corresponding forced output bound is \(e^{-R_{ji}(y,\beta)}\). If
\(4g_j\le\mathrm{rate}_{ji}\) and \(\beta\in[y/2,2y]\), its definition
gives \(R_{ji}(y,\beta)\ge g_jM(y)^2\); for large \(x\) it also gives
the ordinary forced allowance at \(x\).

The inherited condition is
\[\big(f_{ji}/L+u_{ji}(x)\big)\varepsilon(y)+\mathcal R_{ji}(x,y,\beta)\le q_i\varepsilon(x),\qquad \beta\in[y/2,2y].\]
Here \(\mathcal R_{ji}\) is the fixed step-estimator residual at
exponent \(\theta_j\); it is \(o(\varepsilon(x))\) uniformly in this
band when \(\theta_j\le1\), and \(u_{ji}(x)\to0\). The pre-block
inequality gives \(f_{ji}/L\le q_i/8\). Together with
\(\varepsilon(y)/\varepsilon(x)\to1\), these bounds prove the
inequality. Lemma 4.16 combines the conditions in the order: inherited
estimates, local reorganization, cubic estimate, reset pricing; the
sensors were fixed before \(L\).

\emph{Lean:}
\href{https://github.com/yang-mills-group/OpenAI-YM-20261028/blob/3456e5e75cf8d4cd558955a5c7638fba2d68f464/YM/UVNext/COUPLING_RESET/Prices.lean\#L56-L62}{\texttt{can\_delay\_reset}},
\href{https://github.com/yang-mills-group/OpenAI-YM-20261028/blob/3456e5e75cf8d4cd558955a5c7638fba2d68f464/YM/UVNext/COUPLING_RESET/Prices.lean\#L71-L89}{\texttt{buys\_reset}},
\href{https://github.com/yang-mills-group/OpenAI-YM-20261028/blob/3456e5e75cf8d4cd558955a5c7638fba2d68f464/YM/UVNext/COUPLING_RESET/Prices.lean\#L46-L49}{\texttt{sensors}},
\href{https://github.com/yang-mills-group/OpenAI-YM-20261028/blob/3456e5e75cf8d4cd558955a5c7638fba2d68f464/YM/UVNext/COUPLING_RESET/Precision.lean\#L100-L106}{\texttt{delay\_cubic}},
\href{https://github.com/yang-mills-group/OpenAI-YM-20261028/blob/3456e5e75cf8d4cd558955a5c7638fba2d68f464/YM/AnalyticBridge/VERIFIED_RESET_STEP/Currency.lean\#L75-L126}{\texttt{Time.cash}},
\href{https://github.com/yang-mills-group/OpenAI-YM-20261028/blob/3456e5e75cf8d4cd558955a5c7638fba2d68f464/YM/AnalyticBridge/VERIFIED_RESET_STEP/Opening.lean\#L48-L59}{\texttt{Time.late}},
\href{https://github.com/yang-mills-group/OpenAI-YM-20261028/blob/3456e5e75cf8d4cd558955a5c7638fba2d68f464/YM/AnalyticBridge/VERIFIED_RESET_STEP/Opening.lean\#L62-L65}{\texttt{Time.open\_shop}},
\href{https://github.com/yang-mills-group/OpenAI-YM-20261028/blob/3456e5e75cf8d4cd558955a5c7638fba2d68f464/YM/AnalyticBridge/VERIFIED_RESET_STEP/Opening.lean\#L84-L119}{\texttt{Time.now}},
\href{https://github.com/yang-mills-group/OpenAI-YM-20261028/blob/3456e5e75cf8d4cd558955a5c7638fba2d68f464/YM/AnalyticBridge/VERIFIED_RESET_STEP/Inherited.lean\#L30-L56}{\texttt{Time.rawGate}},
\href{https://github.com/yang-mills-group/OpenAI-YM-20261028/blob/3456e5e75cf8d4cd558955a5c7638fba2d68f464/YM/AnalyticBridge/VERIFIED_RESET_STEP/Inherited.lean\#L75-L89}{\texttt{Time.beforeWait}}.

\textbf{Lemma 4.18 (joint start-up thresholds).} Assume \(G\) is simple
and \(c_{\rm dir}\le1\) at every used stage. Fix the observable
sections, marker bounds and nonnegative additional start-up requests.
Before choosing block sizes, enlarge the path and column bounds so that
the Wilson top paths satisfy the stage bounds and the column bound is at
least the factorial of the derivative order. Then the block-size floors
and the coupling and scale thresholds can be enlarged so that every
realized choice of levels meeting them has the start-up bounds of Lemma
4.5 and the following additional bounds. Write \(x=\hat\beta_N\),
\(t=t_N\), \(M=M_N\), \(B=B_N\) and \(a\) for the support exponent at
stage \(j=J_N\).

\begin{itemize}
\tightlist
\item
  For every prescribed observable section of degree \(d\le j\), its
  stencil diameter is at most \(B\), and
  \[\tilde\varepsilon_N C_d t^{-s_d}\le\sqrt{\tilde\varepsilon_N},\]
  where \(s_d\) is the total slot order of its section family and
  \(C_d\ge1\) bounds its evaluation on every configuration by \(C_d\)
  times the coefficient \(\ell^1\) norm.
\item
  For \(\beta\in[x/4,2x]\), the start-up regular bound remains
  \(\varepsilon_N\) after increasing the support exponent by any
  preassigned \(\delta_j\ge0\). The connected-polymer estimate retains a
  forced-decay factor \(e^{-\nu M^2}\), where \(\nu=(128h_4)^{-1}>0\)
  and \(h_4=2^4\,4!\) is the order of the lattice point group; its final
  forced bound is at most \(\tilde\varepsilon_N\).
\item
  The mixture estimates have \(B\ge8\), \(x^{-1/2}(k+1)e^{a5^4}\le1\),
  and a connected-polymer bound at most \(1/2\) with support exponent
  increased by \(8\).
\item
  For each preassigned integer observation radius \(r_j\), the
  first-step numerical bound is
  \[(4B+1)^4(2r_jB+1)^4\,8\left[\frac{b_K}{2}+D_G\log\big(4(1+2x)\big)\right]
  \le \frac{x\,t^2}{64h_4},\] where \(b_K>0\) is the kernel penalty and
  \(D_G=\dim_{\mathbb R}M_n(\mathbb C)\).
\end{itemize}

These conclusions persist when the thresholds are increased.

\emph{Sketch:} At a fixed stage only finitely many observable degrees
are requested. At fixed stage, the stencil and cylinder constants are
fixed, while \(B\to\infty\) and \(e^{-M^{1/5}}\) absorbs every fixed
power of \(t^{-1}\). The additional support weights and the first-step
cost grow polynomially in \(\log x\), whereas \(xt^2=M^2\) grows faster
than every such polynomial. Take common thresholds and enlarge the
block-size floors for the cross-level support estimates. The joint
pricing theorem also increases the scale thresholds, while preserving
the prescribed volume and growth data. This establishes numerical
preparation; subsequent nonlinear steps still require their own activity
estimates.

\emph{Lean:}
\href{https://github.com/yang-mills-group/OpenAI-YM-20261028/blob/3456e5e75cf8d4cd558955a5c7638fba2d68f464/YM/Further/STRONG_STARTUP_CREDIT/Pricing.lean\#L87-L107}{\texttt{StrongStartup.prices}},
\href{https://github.com/yang-mills-group/OpenAI-YM-20261028/blob/3456e5e75cf8d4cd558955a5c7638fba2d68f464/YM/Further/STRONG_STARTUP_CREDIT/Pricing.lean\#L47-L64}{\texttt{StrongStartup.before}},
\href{https://github.com/yang-mills-group/OpenAI-YM-20261028/blob/3456e5e75cf8d4cd558955a5c7638fba2d68f464/YM/Extra/MULTISOURCE_BOUNDS/Schedule.lean\#L98-L111}{\texttt{Extra.Multi.pricing}},
\href{https://github.com/yang-mills-group/OpenAI-YM-20261028/blob/3456e5e75cf8d4cd558955a5c7638fba2d68f464/YM/Extra/WILSON_STARTUP/Fast.lean\#L44-L51}{\texttt{Timing.receive}},
\href{https://github.com/yang-mills-group/OpenAI-YM-20261028/blob/3456e5e75cf8d4cd558955a5c7638fba2d68f464/YM/PhysicalNext/UNMASKED_THROUGH_LAUNCH/Dome.lean\#L105-L138}{\texttt{ThroughLaunch.developer}},
\href{https://github.com/yang-mills-group/OpenAI-YM-20261028/blob/3456e5e75cf8d4cd558955a5c7638fba2d68f464/YM/Extra/WILSON_STARTUP/Author.lean\#L34-L36}{\texttt{Mark.Ready}},
\href{https://github.com/yang-mills-group/OpenAI-YM-20261028/blob/3456e5e75cf8d4cd558955a5c7638fba2d68f464/YM/Further/STRONG_STARTUP_CREDIT/Reserve.lean\#L69-L76}{\texttt{StrongStartup.Escrow}},
\href{https://github.com/yang-mills-group/OpenAI-YM-20261028/blob/3456e5e75cf8d4cd558955a5c7638fba2d68f464/YM/Further/STRONG_STARTUP_CREDIT/Sleeves.lean\#L51-L52}{\texttt{StrongStartup.BondFuel}},
\href{https://github.com/yang-mills-group/OpenAI-YM-20261028/blob/3456e5e75cf8d4cd558955a5c7638fba2d68f464/YM/Extra/WILSON_STARTUP/MixRoom.lean\#L75-L78}{\texttt{Mixture.Extras}}.

\paragraph{4.7 Assembly and the running
law}\label{assembly-and-the-running-law}

\textbf{Lemma 4.19 (bounds at the edges of the shooting window).} Assume
\(\kappa>0\) and \(\bar\beta>0\). Put
\(\phi_m=\bar\beta+\tfrac\kappa2s_m\) and
\(\psi_m=\bar\beta+2\kappa s_m\). Then
\([\phi_m,\psi_m]\subset\mathfrak B_m\), and
\(\phi_0=\psi_0=\bar\beta\). Assume
\(\omega_{J_{m+1}}\le\tfrac\kappa4\log L_m\) (guaranteed by the choice
of block sizes). If a step satisfies the transfer inequality, then:

\begin{itemize}
\tightlist
\item
  input coupling \(\phi_{m+1}\) forces output coupling \(<\phi_m\);
\item
  input coupling \(\psi_{m+1}\) forces output coupling \(>\psi_m\).
\end{itemize}

\emph{Sketch:}
\(\beta_m\le\phi_{m+1}-\kappa\log L_m+\omega\le\phi_m-\tfrac\kappa4\log L_m\),
and symmetrically for \(\psi\).

\emph{Lean:}
\href{https://github.com/yang-mills-group/OpenAI-YM-20261028/blob/3456e5e75cf8d4cd558955a5c7638fba2d68f464/YM/Shoot/Reach.lean\#L81-L90}{\texttt{Shoot.gates}},
\href{https://github.com/yang-mills-group/OpenAI-YM-20261028/blob/3456e5e75cf8d4cd558955a5c7638fba2d68f464/YM/Shoot/Reach.lean\#L60-L68}{\texttt{Shoot.limits}},
\href{https://github.com/yang-mills-group/OpenAI-YM-20261028/blob/3456e5e75cf8d4cd558955a5c7638fba2d68f464/YM/Shoot/Reach.lean\#L73-L79}{\texttt{Shoot.Drifting.pay}}.

\textbf{Lemma 4.20 (shooting in the bare coupling).} Suppose the step
program holds and \(G\) is simple. Then for all \(N,p\) there is a trail
\((\rho_{N,p,m})_{m\le N}\) with the following properties:

\begin{itemize}
\tightlist
\item
  it starts with a Wilson start-up \(e^{-\beta S_W}\);
\item
  each level is admissible, with coupling in \([\phi_m,\psi_m]\);
\item
  consecutive levels satisfy the transfer clause;
\item
  \(\beta_{N,p,0}=\bar\beta\).
\end{itemize}

Choosing these trails for all \(N,p\) gives a history satisfying
(1)--(5) of III.9.3. The common-path clause follows because every trail
uses the same prescribed \(P_{N,m}\).

\emph{Sketch:} Decreasing induction on \(m\) carries an interval
\([l,u]\) of bare couplings and a continuous curve of trails.

\begin{itemize}
\tightlist
\item
  At \(m=N\), restrict the start-up curve to \([\phi_N,\psi_N]\).
\item
  At each step, apply the step program, giving an output curve over the
  same interval. By Lemma 4.19 its couplings at \(l\) and \(u\) lie
  below \(\phi_m\) and above \(\psi_m\). A compactness and
  intermediate-value argument gives \([l',u']\subset[l,u]\) whose image
  is \([\phi_m,\psi_m]\), with \(\phi_m,\psi_m\) attained at the ends.
\item
  At \(m=0\) the left end gives \(\beta_0=\bar\beta\).
\end{itemize}

\emph{Lean:}
\href{https://github.com/yang-mills-group/OpenAI-YM-20261028/blob/3456e5e75cf8d4cd558955a5c7638fba2d68f464/YM/Shoot/Reach.lean\#L112-L154}{\texttt{Shoot.steps}},
\href{https://github.com/yang-mills-group/OpenAI-YM-20261028/blob/3456e5e75cf8d4cd558955a5c7638fba2d68f464/YM/Shoot/Reach.lean\#L18-L49}{\texttt{Shoot.narrow}},
\href{https://github.com/yang-mills-group/OpenAI-YM-20261028/blob/3456e5e75cf8d4cd558955a5c7638fba2d68f464/YM/Shoot/Reach.lean\#L156-L159}{\texttt{Shoot.successful}},
\href{https://github.com/yang-mills-group/OpenAI-YM-20261028/blob/3456e5e75cf8d4cd558955a5c7638fba2d68f464/YM/Shoot/Race.lean\#L67-L82}{\texttt{Recipe.races}},
\href{https://github.com/yang-mills-group/OpenAI-YM-20261028/blob/3456e5e75cf8d4cd558955a5c7638fba2d68f464/YM/PhysicalNext/RG_NEXT_COMPLETE_REPEATABLE_NONLINEAR_STEP/Graces/Eagles.lean\#L33-L50}{\texttt{Graces.winnow}}.

\textbf{Lemma 4.21 (running).} For a history on levels satisfying the
drift-smallness and stage-divergence conditions of III.9.2, clauses (4)
and (5) of III.9.3 imply clause (6).

\emph{Sketch:}

\begin{itemize}
\tightlist
\item
  By induction from \(\beta_{N,p,0}=\bar\beta=\hat\beta_0\):
  \(|\beta_{N,p,i}-\hat\beta_i|\le\sum_{k=1}^{i}\omega_{J_k}\).
\item
  Given \(\varepsilon\), pick \(j\) with
  \(J_{k+1}^{-2}\le\varepsilon/2\) for \(k\ge j\), using \(J\to\infty\).
\item
  The drift condition gives
  \(\sum_{k\le i}\omega_{J_k}\le\sum_{k\le j}\omega_{J_k}+\tfrac\varepsilon2(s_i-s_j)\).
\item
  This is \(\le\varepsilon s_i\) for \(i\) large, since
  \(s_i\to\infty\). The bound is uniform in \(N\ge i\) and \(p\).
\end{itemize}

\emph{Lean:}
\href{https://github.com/yang-mills-group/OpenAI-YM-20261028/blob/3456e5e75cf8d4cd558955a5c7638fba2d68f464/YM/Shoot/Race.lean\#L23-L53}{\texttt{Shoot.trajectory}},
\href{https://github.com/yang-mills-group/OpenAI-YM-20261028/blob/3456e5e75cf8d4cd558955a5c7638fba2d68f464/YM/UVExtension/Orchard/First.lean\#L55-L88}{\texttt{Orchard.traject}}.
The second version replaces (5) by the band at level \(0\).

\textbf{Lemma 4.22 (the history of the main theorem).} Fix the
compatible data of III.13, including its analytic continuation bounds
and finite-cutoff solutions. For every \(N,p\) and every real \(b\) with
\(|b|<r_*\), there is a trail in that construction with level-\(0\)
coupling \(\bar\beta+b\). Its bare coupling is unique within the
specified class of finite-cutoff solutions. The coupling and activity
arrays admit analytic continuation to the complex disc \(|b|<r_*\), with
\(r_*>0\) independent of \(N,p\). At \(b=0\), these trails form a
selected history, and every step is a reset step.

\emph{Sketch:} III.13 solves a coupled finite-cutoff fixed-point problem
for all the level couplings. For real offsets its fixed point gives real
trails and the stated endpoint; the comparison theorem gives uniqueness
in this class. At \(b=0\), their admissibility gives the common paths
and level-wise bounds; their reset steps give exact transfer. Lemma 4.21
supplies running. This is the history used in the Wilson family. The
calibration condition records its realized trails and reset steps; it
does not assert uniqueness among all representations of the same
density.

\emph{Lean:}
\href{https://github.com/yang-mills-group/OpenAI-YM-20261028/blob/3456e5e75cf8d4cd558955a5c7638fba2d68f464/YM/UVExtension/Spire/Spark.lean\#L78-L85}{\texttt{Spire.Uses.exists\_bolt}},
\href{https://github.com/yang-mills-group/OpenAI-YM-20261028/blob/3456e5e75cf8d4cd558955a5c7638fba2d68f464/YM/UVExtension/Spire/Spark.lean\#L68}{\texttt{Spire.Uses.hits}},
\href{https://github.com/yang-mills-group/OpenAI-YM-20261028/blob/3456e5e75cf8d4cd558955a5c7638fba2d68f464/YM/UVExtension/Upland/Plinth.lean\#L71-L75}{\texttt{Upland.Arc.exists\_trial}},
\href{https://github.com/yang-mills-group/OpenAI-YM-20261028/blob/3456e5e75cf8d4cd558955a5c7638fba2d68f464/YM/UVExtension/Forest/Glue.lean\#L35-L36}{\texttt{Mast.calibrated}},
\href{https://github.com/yang-mills-group/OpenAI-YM-20261028/blob/3456e5e75cf8d4cd558955a5c7638fba2d68f464/YM/UVExtension/Upland/Plinth.lean\#L55-L58}{\texttt{Upland.Arc.finiteCut}},
\href{https://github.com/yang-mills-group/OpenAI-YM-20261028/blob/3456e5e75cf8d4cd558955a5c7638fba2d68f464/YM/UVExtension/Spire/Orbs.lean\#L41-L42}{\texttt{Spire.Fitted.acts}},
\href{https://github.com/yang-mills-group/OpenAI-YM-20261028/blob/3456e5e75cf8d4cd558955a5c7638fba2d68f464/YM/UVExtension/Spire/Alpine.lean\#L64-L70}{\texttt{Spire.Fitted.unions}}.
\emph{(proof read for the displayed assembly and finite-cutoff
uniqueness theorem; inputs from III.13)}

\paragraph{4.8 Auxiliary results}\label{auxiliary-results-1}

\textbf{Lemma 4.23 (marker-rooted sums).} This is III.9.8.

\emph{Sketch:}

\begin{itemize}
\tightlist
\item
  If \(x\in X_i\) and a marker \(v\) of \(X_i\) sits in \(X_i\), then
  \(\mathrm{dist}(o(v),x)\le\mathrm{cover}_M(X_i)(4M+2)+2M\). So
  \(e^{a\,\mathrm{cover}}\le e^{(a+1)\mathrm{cover}}\,e^{1-\mathrm{dist}/(4M+2)}\).
\item
  Summing over \(i\) marker by marker gives at most \(c\) per marker.
\item
  Summing over markers gives
  \(\le\mu\sum_ye^{1-\mathrm{dist}(y,x)/(4M+2)}\le\mu e(2(16M+9))^4\),
  uniformly in the torus.
\end{itemize}

\emph{Lean:}
\href{https://github.com/yang-mills-group/OpenAI-YM-20261028/blob/3456e5e75cf8d4cd558955a5c7638fba2d68f464/YM/Extra/MULTISOURCE_BOUNDS/Root.lean\#L161-L199}{\texttt{Extra.Multi.root\_hit}},
\href{https://github.com/yang-mills-group/OpenAI-YM-20261028/blob/3456e5e75cf8d4cd558955a5c7638fba2d68f464/YM/Extra/MULTISOURCE_BOUNDS/Root.lean\#L136-L156}{\texttt{torus\_kernel}}.

\textbf{Lemma 4.24 (decay of a shifted Gram inverse).} Fix a finite
channel set \(I\), \(r,\eta>0\) and \(A\ge0\). There are \(C\ge0\) and
\(\tau>0\), depending only on \(I,r,\eta,A\), such that every real
convolution kernel \(P\) on \(\mathbb Z^4\) with
\[\sum_{z\in\mathbb Z^4}e^{\eta|z|_1}|P_{ab}(z)|\le A\qquad(a,b\in I)\]
has a real convolution kernel \(S\) with the corresponding bound \(C\)
at exponent \(\tau\), and, on every torus,
\[S_n\big(rI+P_n^{\mathsf T}P_n\big)=I.\] Here \(P_n,S_n\) are
periodizations. Thus \(S_n=(rI+P_n^{\mathsf T}P_n)^{-1}\), uniformly in
the torus side.

\emph{Sketch:} On real momentum, \(rI+\hat P(k)^*\hat P(k)\) is positive
definite. Compactness in the fixed finite channel space and the weighted
kernel bound give invertibility and a uniform inverse bound in a
narrower complex strip. Fourier coefficient estimates give an
exponentially summable inverse kernel. The discrete Fourier transform
identifies its periodization with the inverse matrix on every torus.

\emph{Lean:}
\href{https://github.com/yang-mills-group/OpenAI-YM-20261028/blob/3456e5e75cf8d4cd558955a5c7638fba2d68f464/YM/Extra/WILSON_STARTUP/Transform.lean\#L86-L91}{\texttt{Polar.deliver}},
\href{https://github.com/yang-mills-group/OpenAI-YM-20261028/blob/3456e5e75cf8d4cd558955a5c7638fba2d68f464/YM/Extra/WILSON_STARTUP/Transform.lean\#L48-L84}{\texttt{Polar.invert\_tori}},
\href{https://github.com/yang-mills-group/OpenAI-YM-20261028/blob/3456e5e75cf8d4cd558955a5c7638fba2d68f464/YM/Extra/WILSON_STARTUP/Polar.lean\#L142-L163}{\texttt{Polar.writes}}.

\subsubsection{5. Where the difficulty
lies}\label{where-the-difficulty-lies-11}

The early presentation with local bounds comes from III.10, the one-loop
coefficient from III.7, and the compatible kinetic data from III.13.
Given those inputs, the reset must preserve the full masked density,
including failed-mask discrepancies and the forced gas. Two-point
matching alone does not prove that identity. Lemma 4.13 supplies it and
proves the output activity bounds.

The marginal coefficient is isolated by gauge symmetry, simplicity and
lattice isotropy. It is a Hessian coefficient, so its quadratic energy
carries a factor \(1/2\). The Wilson start-up has a uniform \(O(t^5)\)
residual even when individual links are far from the identity.
Continuity in the bare parameter must hold in the activity norms despite
the hard configuration masks; shooting depends on this continuity.

The block factors grow. The planned one-loop condition implies
\(\omega_{J_{m+1}}\ge1\), while the drift condition gives
\(\log L_m\ge J_{m+1}^2\to\infty\). The gentle-scale condition requires
\(\log L_m=o(s_m)\). The reset correction is small:
\[|\delta_x|\le\frac{c_m}{4}\hat\beta_m^{-1/10}e^{-2(\log\log\hat\beta_m)^2}.\]
The constants \(c_m\) depend on the stage pair, so this is a fixed-stage
estimate. The separate one-loop error is bounded independently of \(L\)
once its earlier parameters are fixed; no decay of that error in \(L\)
is required. The argument gives \(\beta_m=\bar\beta+\kappa s_m+o(s_m)\)
uniformly in the cutoff and volume. It does not give a bounded remainder
or a two-loop expansion.

The coarsest coupling is prescribed exactly. Shooting chooses the bare
coupling separately for each cutoff and volume and restricts the
surviving parameter intervals to keep every level in its band. The
uniqueness theorem in III.13 applies to its specified continuation
class.

\subsubsection{6. Notation}\label{notation-11}

{\def\LTcaptype{none} % do not increment counter
\begin{longtable}[]{@{}
  >{\raggedright\arraybackslash}p{(\linewidth - 4\tabcolsep) * \real{0.3333}}
  >{\raggedright\arraybackslash}p{(\linewidth - 4\tabcolsep) * \real{0.3333}}
  >{\raggedright\arraybackslash}p{(\linewidth - 4\tabcolsep) * \real{0.3333}}@{}}
\toprule\noalign{}
\begin{minipage}[b]{\linewidth}\raggedright
Written proof
\end{minipage} & \begin{minipage}[b]{\linewidth}\raggedright
Lean (glossary)
\end{minipage} & \begin{minipage}[b]{\linewidth}\raggedright
Meaning
\end{minipage} \\
\midrule\noalign{}
\endhead
\bottomrule\noalign{}
\endlastfoot
\(\kappa=\frac{11C_G}{24\pi^2}=2b_0\) &
\href{https://github.com/yang-mills-group/OpenAI-YM-20261028/blob/3456e5e75cf8d4cd558955a5c7638fba2d68f464/YM/Gauge/Colors.lean\#L184}{\texttt{Gauge.slopeConstant}},
\href{https://github.com/yang-mills-group/OpenAI-YM-20261028/blob/3456e5e75cf8d4cd558955a5c7638fba2d68f464/YM/Gauge/Colors.lean\#L180-L182}{\texttt{Gauge.casimir}}
& slope of \(\beta\) against \(\log(1/a)\) \\
\(s_m\), \(a_m\), \(L_m\) &
\href{https://github.com/yang-mills-group/OpenAI-YM-20261028/blob/3456e5e75cf8d4cd558955a5c7638fba2d68f464/YM/Lattice/Scale.lean\#L115}{\texttt{Schedule.scale}},
\href{https://github.com/yang-mills-group/OpenAI-YM-20261028/blob/3456e5e75cf8d4cd558955a5c7638fba2d68f464/YM/Lattice/Scale.lean\#L99}{\texttt{Schedule.a}},
\texttt{block} & scale \(\log(1/a_m)\), spacing, block factor (odd,
\(\ge3\)) \\
\(\bar\beta\) & \texttt{Calendar.bar} (``barrier'') & coupling at the
coarsest level \\
\(\hat\beta_m\), \(\mathfrak B_m\) &
\href{https://github.com/yang-mills-group/OpenAI-YM-20261028/blob/3456e5e75cf8d4cd558955a5c7638fba2d68f464/YM/RG/History.lean\#L79}{\texttt{Calendar.hatBeta}},
\href{https://github.com/yang-mills-group/OpenAI-YM-20261028/blob/3456e5e75cf8d4cd558955a5c7638fba2d68f464/YM/Shoot/Criteria.lean\#L63}{\texttt{band}}
& reference coupling, band \([\hat\beta/2,2\hat\beta]\) \\
\(t_m\), \(M_m\), \(B_m\) &
\href{https://github.com/yang-mills-group/OpenAI-YM-20261028/blob/3456e5e75cf8d4cd558955a5c7638fba2d68f464/YM/RG/History.lean\#L83}{\texttt{threshold}},
\href{https://github.com/yang-mills-group/OpenAI-YM-20261028/blob/3456e5e75cf8d4cd558955a5c7638fba2d68f464/YM/RG/History.lean\#L81}{\texttt{monitor}},
\href{https://github.com/yang-mills-group/OpenAI-YM-20261028/blob/3456e5e75cf8d4cd558955a5c7638fba2d68f464/YM/RG/History.lean\#L82}{\texttt{coverSize}}
& small-field threshold, its multiplier, cover size \\
\(\varepsilon_m\), \(\tilde\varepsilon_m\) &
\href{https://github.com/yang-mills-group/OpenAI-YM-20261028/blob/3456e5e75cf8d4cd558955a5c7638fba2d68f464/YM/RG/History.lean\#L84}{\texttt{allowance}},
\href{https://github.com/yang-mills-group/OpenAI-YM-20261028/blob/3456e5e75cf8d4cd558955a5c7638fba2d68f464/YM/RG/History.lean\#L85}{\texttt{forcedAllowance}}
& sizes of regular and forced activities \\
\(J_m\), \(\omega_j\), \(\theta_j\) & \texttt{Calendar.J},
\href{https://github.com/yang-mills-group/OpenAI-YM-20261028/blob/3456e5e75cf8d4cd558955a5c7638fba2d68f464/YM/RG/History.lean\#L45-L61}{\texttt{Stage}}
(\texttt{drift}), \texttt{theta} & stage index, drift allowance,
smallness exponent \\
\(\rho=(P,c,E,K,b,\beta,v)\) &
\href{https://github.com/yang-mills-group/OpenAI-YM-20261028/blob/3456e5e75cf8d4cd558955a5c7638fba2d68f464/YM/RG/History.lean\#L165-L171}{\texttt{Rep}}
& effective action (``representation'') \\
\(\mathcal K_P\) &
\href{https://github.com/yang-mills-group/OpenAI-YM-20261028/blob/3456e5e75cf8d4cd558955a5c7638fba2d68f464/YM/Push/BETA_TRANSPORT/Curve.lean\#L22-L25}{\texttt{BetaTransport.kinetic}}
& \(\beta\)-coefficient of the leading terms \\
\(T^{(m)}_\beta\) &
\href{https://github.com/yang-mills-group/OpenAI-YM-20261028/blob/3456e5e75cf8d4cd558955a5c7638fba2d68f464/YM/RG/Block.lean\#L77-L79}{\texttt{StepData.push}}
& one block-spin step \\
\(\beta_{N,p,m}\), \(\rho_{N,p,m}\) & \texttt{History.β},
\texttt{History.rep} & selected history \\
\(\alpha_{N,m}\), \(P_{N,m}\) &
\href{https://github.com/yang-mills-group/OpenAI-YM-20261028/blob/3456e5e75cf8d4cd558955a5c7638fba2d68f464/YM/PCalc/Loop.lean\#L74-L77}{\texttt{Kin}}
(\texttt{alpha}, \texttt{path}) & one-loop shift of the step, path
data \\
\(\mathfrak D\) &
\href{https://github.com/yang-mills-group/OpenAI-YM-20261028/blob/3456e5e75cf8d4cd558955a5c7638fba2d68f464/YM/PCalc/Band.lean\#L26-L36}{\texttt{Directed}}
& scale sequence, kinetic data and recipe \\
reset step &
\href{https://github.com/yang-mills-group/OpenAI-YM-20261028/blob/3456e5e75cf8d4cd558955a5c7638fba2d68f464/YM/PCalc/Band.lean\#L49-L55}{\texttt{Directed.Move}}
& one step with exact transfer and
\(\lvert\beta-\beta'-\alpha\rvert\le1\) \\
step program &
\href{https://github.com/yang-mills-group/OpenAI-YM-20261028/blob/3456e5e75cf8d4cd558955a5c7638fba2d68f464/YM/PCalc/Band.lean\#L65-L72}{\texttt{Directed.Program}}
& start-up curves and reset along all admissible curves \\
\((c_m,r_m)\) &
\href{https://github.com/yang-mills-group/OpenAI-YM-20261028/blob/3456e5e75cf8d4cd558955a5c7638fba2d68f464/YM/UVNext/COUPLING_RESET/Prices.lean\#L38-L43}{\texttt{ResetPrices}}
(\texttt{cost}, \texttt{radius}) & income multiplier, reset window
(\(\le1\)) \\
\(\eta_j\), probe \(V\) &
\href{https://github.com/yang-mills-group/OpenAI-YM-20261028/blob/3456e5e75cf8d4cd558955a5c7638fba2d68f464/YM/UVNext/COUPLING_RESET/Probe.lean\#L92-L94}{\texttt{Sensor}},
\href{https://github.com/yang-mills-group/OpenAI-YM-20261028/blob/3456e5e75cf8d4cd558955a5c7638fba2d68f464/YM/UVNext/COUPLING_RESET/Advance.lean\#L55-L58}{\texttt{Usable}}
& lower bound \(\mathcal K(V)\ge\lvert\Lambda\rvert\eta t^2\) \\
\(f(x,d)\), \(d\) &
\href{https://github.com/yang-mills-group/OpenAI-YM-20261028/blob/3456e5e75cf8d4cd558955a5c7638fba2d68f464/YM/UVNext/COUPLING_RESET/Reset.lean\#L64-L65}{\texttt{Template.residual}}
& scalar residual, coupling correction \\
\(\chi_X\), \(Q_i\), \(E_i^0\), \(D_i\) &
\href{https://github.com/yang-mills-group/OpenAI-YM-20261028/blob/3456e5e75cf8d4cd558955a5c7638fba2d68f464/YM/AnalyticBridge/VERIFIED_RESET_STEP/Residual.lean\#L114-L120}{\texttt{Residual.splits}}
& small-field indicator, local quadratic projection, subtracted regular
function, mask discrepancy \\
\(\delta(A)\) &
\href{https://github.com/yang-mills-group/OpenAI-YM-20261028/blob/3456e5e75cf8d4cd558955a5c7638fba2d68f464/YM/AnalyticBridge/VERIFIED_RESET_STEP/Isotropy.lean\#L75}{\texttt{Averaging.delta}}
& isotropic curvature-squared coefficient of activities \\
early presentation &
\href{https://github.com/yang-mills-group/OpenAI-YM-20261028/blob/3456e5e75cf8d4cd558955a5c7638fba2d68f464/YM/AnalyticBridge/VERIFIED_RESET_STEP/Evaluation.lean\#L26-L27}{\texttt{Collection.earlyValue}},
\href{https://github.com/yang-mills-group/OpenAI-YM-20261028/blob/3456e5e75cf8d4cd558955a5c7638fba2d68f464/YM/AnalyticBridge/VERIFIED_RESET_STEP/Inputs.lean\#L70-L71}{\texttt{Return.atlas}}
& blocked density before the reset \\
\(P^{\rm top}\), \(R_p\), \(Q_4\) &
\href{https://github.com/yang-mills-group/OpenAI-YM-20261028/blob/3456e5e75cf8d4cd558955a5c7638fba2d68f464/YM/Extra/WILSON_STARTUP/Design.lean\#L34-L38}{\texttt{Design.top}},
\href{https://github.com/yang-mills-group/OpenAI-YM-20261028/blob/3456e5e75cf8d4cd558955a5c7638fba2d68f464/YM/Extra/WILSON_STARTUP/Face.lean\#L58}{\texttt{residual}},
\href{https://github.com/yang-mills-group/OpenAI-YM-20261028/blob/3456e5e75cf8d4cd558955a5c7638fba2d68f464/YM/Extra/WILSON_STARTUP/Curvature.lean\#L154}{\texttt{fourth}}
& Wilson top path data, one-plaquette residual, quartic coefficient \\
\(\rho^{\rm start}_\beta\) &
\href{https://github.com/yang-mills-group/OpenAI-YM-20261028/blob/3456e5e75cf8d4cd558955a5c7638fba2d68f464/YM/Bid/Kitchen.lean\#L27-L28}{\texttt{Bid.Cake}}
with \texttt{payload}, \texttt{dictionary} & Wilson start-up
representation \\
compatible \((c,b)\) & \texttt{WilsonStartup.Decent} &
\(c_{\rm dir}\le1\), \(B\ge2\), \(te^{(B-2)/c_{\rm rel}}\ge2\) \\
\(\phi_m,\psi_m\) &
\href{https://github.com/yang-mills-group/OpenAI-YM-20261028/blob/3456e5e75cf8d4cd558955a5c7638fba2d68f464/YM/Shoot/Reach.lean\#L57}{\texttt{floorBeta}},
\href{https://github.com/yang-mills-group/OpenAI-YM-20261028/blob/3456e5e75cf8d4cd558955a5c7638fba2d68f464/YM/Shoot/Reach.lean\#L58}{\texttt{ceilBeta}}
& shooting window \\
numerical bounds and entry thresholds & glossary item 1 & constants in
the polymer norm estimates and lower bounds on \(\hat\beta\) \\
\end{longtable}
}

\subsubsection{7. What remains
unexplained}\label{what-remains-unexplained-11}

\begin{enumerate}
\def\labelenumi{\arabic{enumi}.}
\tightlist
\item
  \textbf{Nonlinear inputs.} The reset is conditional on the early
  presentation, stronger local activity bounds and continuity. Their
  supply is the conclusion of
  \href{https://github.com/yang-mills-group/OpenAI-YM-20261028/blob/3456e5e75cf8d4cd558955a5c7638fba2d68f464/YM/PhysicalNext/RG_NEXT_COMPLETE_REPEATABLE_NONLINEAR_STEP/Rain/Glazing.lean\#L97-L104}{Rain.glazing}
  and
  \href{https://github.com/yang-mills-group/OpenAI-YM-20261028/blob/3456e5e75cf8d4cd558955a5c7638fba2d68f464/YM/PhysicalNext/RG_NEXT_COMPLETE_REPEATABLE_NONLINEAR_STEP/Roost/Foundry.lean\#L56}{Roost.Factor.offering}.
  Their assembly proofs identify the required inputs; the underlying
  nonlinear and Gaussian-sector estimates belong to III.10 and III.4 and
  are not independently verified here.
\item
  \textbf{Activity reorganization.} The explicit splitting and the proof
  of
  \href{https://github.com/yang-mills-group/OpenAI-YM-20261028/blob/3456e5e75cf8d4cd558955a5c7638fba2d68f464/YM/AnalyticBridge/VERIFIED_RESET_STEP/Investor.lean\#L112-L128}{Local.renew}
  give Lemma 4.13. A full written proof should expand the local
  reorganization constants and the connected-polymer estimate that
  absorbs the logarithmic factors. Their detailed norm estimates are
  inputs from II.2 and III.10.
\item
  \textbf{Finite-cutoff uniqueness.} The uniqueness in Lemma 4.22
  concerns the class of solutions specified in III.13, with its fixed
  data and continuation bounds. This construction solves coupled
  fixed-point equations rather than choosing a subinterval by Lemma
  4.20. The displayed assembly and uniqueness theorem have been checked;
  the analytic existence and comparison estimates underneath them remain
  for III.13. Real offsets satisfy \(|b|<r_*\); complex offsets
  parametrize analytic continuations, not real positive densities.
\item
  \textbf{Relaxed histories.} The intermediate history predicate with no
  endpoint restriction is distinct from the selected-history
  specification. Its running estimate uses the band at level \(0\). The
  final centered history satisfies the exact endpoint clause.
\item
  \textbf{Numerical and bibliographic scope.} Lemma 4.18 gives the
  start-up numerical bounds; their use for source polynomials and
  first-step observations requires the norm definitions in III.12 and
  III.15. The cited numerical preparation does not itself construct
  those observables or their nonlinear continuations. The literature
  citations in §2 are background references; their exact bibliographic
  details and comparison with this construction remain unchecked here.
\item
  \textbf{Verification scope.} The correspondence checks use the
  supplied Lean source and link table. No Lean build or independent
  axiom-closure check was performed. The package sizes and main-theorem
  dependency counts are retained from the unit map (tentative).
\end{enumerate}

\subsection{III.10. One full nonlinear renormalization step: bounds on
irrelevant
terms}\label{iii.10.-one-full-nonlinear-renormalization-step-bounds-on-irrelevant-terms}

\textbf{Part:} III, Ultraviolet construction (renormalization group).
\textbf{Lean:}
\href{https://github.com/yang-mills-group/OpenAI-YM-20261028/tree/3456e5e75cf8d4cd558955a5c7638fba2d68f464/YM/Work}{\texttt{YM/Work/}}
(items E01--E09, H01--H04, I01--I04),
\href{https://github.com/yang-mills-group/OpenAI-YM-20261028/tree/3456e5e75cf8d4cd558955a5c7638fba2d68f464/YM/Bid}{\texttt{YM/Bid/}},
\href{https://github.com/yang-mills-group/OpenAI-YM-20261028/tree/3456e5e75cf8d4cd558955a5c7638fba2d68f464/YM/Next/FLUCTUATIONS}{\texttt{YM/Next/FLUCTUATIONS/}},
\href{https://github.com/yang-mills-group/OpenAI-YM-20261028/tree/3456e5e75cf8d4cd558955a5c7638fba2d68f464/YM/AnalyticNext/CONNECTED_ACTIVITY_JETS_E6}{\texttt{YM/AnalyticNext/CONNECTED\_ACTIVITY\_JETS\_E6/}},
\href{https://github.com/yang-mills-group/OpenAI-YM-20261028/tree/3456e5e75cf8d4cd558955a5c7638fba2d68f464/YM/PhysicalNext/RG_NEXT_COMPLETE_REPEATABLE_NONLINEAR_STEP}{\texttt{YM/PhysicalNext/RG\_NEXT\_COMPLETE\_REPEATABLE\_NONLINEAR\_STEP/}},
\href{https://github.com/yang-mills-group/OpenAI-YM-20261028/tree/3456e5e75cf8d4cd558955a5c7638fba2d68f464/YM/PhysicalNext/V9_JOINT_ENDPOINT_LOAD}{\texttt{YM/PhysicalNext/V9\_JOINT\_ENDPOINT\_LOAD/}},
\href{https://github.com/yang-mills-group/OpenAI-YM-20261028/tree/3456e5e75cf8d4cd558955a5c7638fba2d68f464/YM/Pressure}{\texttt{YM/Pressure/}}
(definitions). \textbf{Size:} 49,448 lines in the main-theorem closure.

\textbf{Conventions.} Levels are indexed from coarse to fine, as in
III.9. - \emph{Step and couplings.} One step \(S\) goes from a fine
torus of side \(f=cL\) to a coarse torus of side \(c\), with odd block
factor \(L\ge3\). Write \(x\) for the coarse reference coupling
\(\hat\beta_m\). The fine coupling is
\(\beta_F=x+b_G\log L=\hat\beta_{m+1}\)
(\href{https://github.com/yang-mills-group/OpenAI-YM-20261028/blob/3456e5e75cf8d4cd558955a5c7638fba2d68f464/YM/AnalyticNext/CONNECTED_ACTIVITY_JETS_E6/Intake.lean\#L67-L70}{\texttt{clock\_next}}),
where \(b_G\ge0\) is the one-loop slope of III.9; \(b_G>0\) when the Lie
algebra of \(G\) is simple. - \emph{Thresholds, radius and cutoff.}
\[t_c(x)=x^{-1/2}e^{(\log\log x)^2}\ \text{(coarse field)},\qquad t=t_c(\beta_F)\ \text{(fine field)},\qquad T=\beta_F^{-49/100}\ \text{(large field proper)},\]
with cluster radius \(B=\lceil(\log x)^2\rceil\) and path-length cutoff
\(R=\lfloor B/10^5\rfloor\). - \emph{Sizes.} Since
\(e^{(\log\log\beta)^2}=\beta^{o(1)}\), we have
\(t=\beta_F^{-1/2+o(1)}\ll T\). On the other hand
\(\beta_Ft^2=e^{2(\log\log\beta_F)^2}\) grows faster than every power of
\(\log\beta_F\). - \emph{Field sizes.} \(\chi(g)\) is the chordal
distance from \(g\in G\) to \(1\). For a plaquette \(p\),
\(|U_p|=\chi(U_{\partial p})\). For a coarse bond \(b\),
\(w_b=\chi(V_bQ_b^{-1})\) measures how far the coarse field \(V\) is
from the block-spin predictor \(Q=Q(U,k)\) of the fine field \(U\) with
gauge-fixing tags \(k\) (III.1). - ``For all large \(x\)'' means
\(x\ge x_0\), with \(x_0\) independent of the period and volume. Other
dependencies are specified in each statement.

\subsubsection{1. What this unit proves}\label{what-this-unit-proves-12}

The unit does not construct the effective actions; that is III.9, with
III.5--III.8. It proves the estimates on one nonlinear block-spin step
that the construction consumes: - a deterministic lower bound on the
action gap of the step near large-field regions (III.10.1); - bounds on
the activities produced when the fluctuation field is integrated out,
and on all their derivatives in the coarse field (III.10.2, III.10.3); -
bounds on the Gaussian sector of the step (III.10.4); - the numerical
schedule under which these estimates hold at every level (III.10.5).

III.10.6 collects the start-up fusion and marker estimates. In the
energy argument, constants are chosen before \(L\), after \(L\), or
after \(W\); this order determines which losses they can absorb. The
derivative and Gaussian estimates have their own quantifier orders,
stated below.

\textbf{Data of a step.} - \emph{Request.} A request \(P\)
(\href{https://github.com/yang-mills-group/OpenAI-YM-20261028/blob/3456e5e75cf8d4cd558955a5c7638fba2d68f464/YM/Pressure/Seeds.lean\#L89-L96}{\texttt{Teller}})
consists of: - kernel choices \(K\), with penalty \(a>0\) and cuts; - a
fine stage and a coarse stage, each with clip constants
\(\alpha,\lambda,\rho>0\) satisfying \(\lambda<\alpha/16\)
(\href{https://github.com/yang-mills-group/OpenAI-YM-20261028/blob/3456e5e75cf8d4cd558955a5c7638fba2d68f464/YM/RG/PathActions.lean\#L139-L146}{\texttt{Clip}}),
a path width and a path budget; - a step index \(j\) and constants
\(\eta>0\), \(C\). - \emph{Tube.} A tube \(\mathscr T\)
(\href{https://github.com/yang-mills-group/OpenAI-YM-20261028/blob/3456e5e75cf8d4cd558955a5c7638fba2d68f464/YM/Pressure/Rows.lean\#L94-L99}{\texttt{Tube}})
carries the towers \(\Phi,\Phi',N\) and the fine and coarse path data.
\(\Phi\) maps coarse plaquette coordinates to fine plaquette
coordinates; \(\Phi'\) maps them to coarse-bond (``interior'')
coordinates; \(N\) maps them to gauge coordinates. -
\emph{Admissibility.} A step with tube \((S,\mathscr T)\) is admissible
for \(P\)
(\href{https://github.com/yang-mills-group/OpenAI-YM-20261028/blob/3456e5e75cf8d4cd558955a5c7638fba2d68f464/YM/Pressure/Seeds.lean\#L101-L109}{\texttt{Teller.Funds}})
if: - \(S\) uses \(K\) and index \(j\); - the patch is well posed
(\href{https://github.com/yang-mills-group/OpenAI-YM-20261028/blob/3456e5e75cf8d4cd558955a5c7638fba2d68f464/YM/RG/Patch.lean\#L182-L185}{\texttt{Branches}}:
unique minimizers whenever the frame or predictor selector uses
minimization, and a nonempty set of eligible axial tags for every
configuration); - the scalar tower coefficients have exponentially
weighted absolute sums at rate \(\eta\) bounded by \(C/L^2\) for the
fine kernel and \(C\) for the interior and gauge kernels. Finite
differences of order \(\ell\le j-2\) of the one-form and two-form cross
kernels have bounds \(C/L^{1+\ell}\) and \(C/L^{2+\ell}\), respectively
(\href{https://github.com/yang-mills-group/OpenAI-YM-20261028/blob/3456e5e75cf8d4cd558955a5c7638fba2d68f464/YM/PREC/Stacks.lean\#L111-L114}{\texttt{YM.Precision.Towers.Deep}});
- the clips and path budgets are those of \(P\); - the flat, uncut tower
matrices satisfy the exact identities
\href{https://github.com/yang-mills-group/OpenAI-YM-20261028/blob/3456e5e75cf8d4cd558955a5c7638fba2d68f464/YM/PREC/Stacks.lean\#L88-L96}{\texttt{Towers.Roll}},
among them the resolution of the identity
\[\Phi_0^*\Phi_0+\Phi_0'^*\Phi_0'+N_0^*N_0=1.\] The transported kernels
cut at length \(R\) need not satisfy this identity; their defects are
estimated in §4.5.

\textbf{Seeds.} A store \(A=(A_F,A_I,A_\flat)\)
(\href{https://github.com/yang-mills-group/OpenAI-YM-20261028/blob/3456e5e75cf8d4cd558955a5c7638fba2d68f464/YM/Pressure/Seeds.lean\#L34-L37}{\texttt{Stores}})
consists of finite sets of fine plaquettes (``costs''), coarse bonds
(``aisles'') and coarse sites (``flags''). \(\sharp A\) is the sum of
their cardinalities, including coincident seeds of different types.
\(D_r(A)\) is the coarse \(r\)-neighbourhood of their positions: all
parent sites touched by a fine plaquette, both endpoints of a coarse
bond, and each flag site. \(A\) is \emph{consistent} with \((U,k,V)\) at
\((x,W)\)
(\href{https://github.com/yang-mills-group/OpenAI-YM-20261028/blob/3456e5e75cf8d4cd558955a5c7638fba2d68f464/YM/Pressure/Seeds.lean\#L51-L58}{\texttt{Plot}})
if the following hold: - whenever a parent site touched by \(p\) meets
\(D_{5B}(A)\), \(p\in A_F\) implies \(|U_p|\ge t/4\), and
\(p\notin A_F\) implies \(|U_p|\le t/2\); - whenever an endpoint of
\(b\) meets \(D_{5B}(A)\), \(b\in A_I\) implies \(w_b\ge Wt\), and
\(b\notin A_I\) implies \(w_b\le2Wt\); - \(A_\flat\) agrees with the
flags of \(V\) at threshold \(t_c(x)\) on \(D_{5B}(A)\); - each site tag
based in \(D_{5B}(A)\) has positive frame weight.

Nothing is assumed outside \(D_{5B}(A)\). The \emph{spent} seeds
\(A^{\rm sp}\) are the costs with \(|U_p|>T\) and the aisles with
\(w_b>T\).

\textbf{Action gap.} - \emph{Clipped density.} For a stage with clip
\((\alpha,\lambda,\rho)\) and threshold \(s\), the clipped density of a
plaquette is \(\phi_s(|U_p|)+\mathbf 1_{\rm kept}\tfrac12|\ell_p(U)|^2\)
(\href{https://github.com/yang-mills-group/OpenAI-YM-20261028/blob/3456e5e75cf8d4cd558955a5c7638fba2d68f464/YM/Pressure/Rows.lean\#L68}{\texttt{crown}}).
Here: - \(\phi_s(r)=\alpha r^2/2\) for \(r\le s\) and
\(\phi_s(r)=\lambda sr\) for \(r>s\)
(\href{https://github.com/yang-mills-group/OpenAI-YM-20261028/blob/3456e5e75cf8d4cd558955a5c7638fba2d68f464/YM/RG/PathActions.lean\#L218-L219}{\texttt{clipped}});
- \(\ell_p(U)\) is the path series of the stage, cut at length \(R\); -
``kept'' means \(|U_q|\le s\,e^{{\rm dist}(p,q)/\rho}\) for all \(q\)
(\href{https://github.com/yang-mills-group/OpenAI-YM-20261028/blob/3456e5e75cf8d4cd558955a5c7638fba2d68f464/YM/RG/PathActions.lean\#L210-L211}{\texttt{kept}}).
- \emph{Actions.} \(\mathcal A_F(U)\) is the fine clipped action (fine
stage, threshold \(t\)) and \(\mathcal A_c(V)\) the coarse one (coarse
stage, threshold \(t_c(x)\)). - \emph{Step functional.}
\(\mathcal E_A(U,V)\)
(\href{https://github.com/yang-mills-group/OpenAI-YM-20261028/blob/3456e5e75cf8d4cd558955a5c7638fba2d68f464/YM/Pressure/Seeds.lean\#L85}{\texttt{assessed}})
is the sum of the row contributions of Lemma 4.1 over the rows located
in \(D_{2B}(A)\). Summed over all rows of the torus, the same
contributions give exactly
\[\mathcal A_F(U)+\frac a2\sum_b w_b^2-\mathcal A_c(V).\] So
\(\mathcal E_A\) is the action gap of the step (fine action, plus
block-spin penalty, minus coarse action), localized near the seeds.

\textbf{Proposition III.10.1 (energy per seed).} Let \(P\) be a request.
There is \(L_0\) such that for every block factor \(L\ge L_0\) there is
\(W_0(L)\) such that for every \(W\ge W_0\) there is \(x_0\) with the
following property. Take \(x\ge x_0\) and: - a step \(S\) of block
factor \(L\) with coarse side \(c\ge100(1+B)\), and a tube
\(\mathscr T\) admissible for \(P\); - a spanning forest \(F\) of the
block structure, and a fine field \(U\) equal to \(1\) on the trunks of
\(F\); - a store \(A\), tags \(k\) and a coarse field \(V\) such that
\(A\) is consistent with \((U,k,V)\) at \((x,W)\).

Then
\[\kappa\,t^2\,\sharp A\;\le\;\mathcal E_A(U,V),\qquad \kappa=\tfrac12\min\{\lambda_F/8,\ a/32,\ \alpha_c/64\},\]
where \(\lambda_F\) is the linear clip constant of the fine stage and
\(\alpha_c\) the quadratic clip constant of the coarse stage.

\emph{Status:} theorem under the stated admissibility, tree-gauge and
consistency hypotheses. - The Lean statement
(\href{https://github.com/yang-mills-group/OpenAI-YM-20261028/blob/3456e5e75cf8d4cd558955a5c7638fba2d68f464/YM/Pressure/Seeds.lean\#L124-L125}{\texttt{Teller.cash}},
\href{https://github.com/yang-mills-group/OpenAI-YM-20261028/blob/3456e5e75cf8d4cd558955a5c7638fba2d68f464/YM/Pressure/Seeds.lean\#L129-L131}{\texttt{Teller.Grants}})
asserts only some \(\kappa>0\) chosen after \(L\) and \(W\). The value
above is the one the proof produces (Lemma 4.14), and it is independent
of \(L\) and \(W\). - This deterministic estimate precedes the Gaussian
tail and activity bounds of III.10.2.

\emph{Lean:} proved by
\href{https://github.com/yang-mills-group/OpenAI-YM-20261028/blob/3456e5e75cf8d4cd558955a5c7638fba2d68f464/YM/Next/FLUCTUATIONS/Grants.lean\#L24}{\texttt{Fluctuations.grants}}.
This is
\href{https://github.com/yang-mills-group/OpenAI-YM-20261028/blob/3456e5e75cf8d4cd558955a5c7638fba2d68f464/YM/Integrated.lean\#L222-L225}{\texttt{Integrated.pressure\_grants}}
applied to
\href{https://github.com/yang-mills-group/OpenAI-YM-20261028/blob/3456e5e75cf8d4cd558955a5c7638fba2d68f464/YM/Next/FLUCTUATIONS/Bare.lean\#L64-L120}{\texttt{Fluctuations.bare}},
through
\href{https://github.com/yang-mills-group/OpenAI-YM-20261028/blob/3456e5e75cf8d4cd558955a5c7638fba2d68f464/YM/Bid/Energy.lean\#L155-L158}{\texttt{Bid.issues}}.

\textbf{Proposition III.10.2 (activities of one step).} Let \(P\) be a
request and \(\mathcal T\) target data
(\href{https://github.com/yang-mills-group/OpenAI-YM-20261028/blob/3456e5e75cf8d4cd558955a5c7638fba2d68f464/YM/AnalyticNext/CONNECTED_ACTIVITY_JETS_E6/OptionCost.lean\#L62-L76}{\texttt{Target}}):
constants \(\varepsilon,\delta>0\), \(e\ge0\), a rate
\(r\in(-\tfrac12,0)\), a halo \(h\ge500\) and a local exponential chart
of \(G\), with \(4\varepsilon\) below the support exponent of the fine
stage and \(8\varepsilon\) below both stage widths. Put
\(a_*=e+2\varepsilon+\delta\).

There is \(L_0\) such that for every \(L\ge L_0\) and every large \(W\)
there is \(\mu>0\) with the following property. For every patch \(D\) of
block factor \(L\) satisfying the branch condition there are window
constants \(o\)
(\href{https://github.com/yang-mills-group/OpenAI-YM-20261028/blob/3456e5e75cf8d4cd558955a5c7638fba2d68f464/YM/AnalyticNext/CONNECTED_ACTIVITY_JETS_E6/Issue.lean\#L33-L46}{\texttt{Target.Order}}:
a decay rate, a spectral gap, a shell width \(K\), a Lipschitz constant
\(\gamma\), a drift cap, a plateau \(a_o\) and a trigger level \(\xi\))
and a constant \(C\ge0\) such that the following holds for all large
\(x\), every step \(S\) on \(D\) with \(c\ge100(1+B)\) and every tube
admissible for \(P\)
(\href{https://github.com/yang-mills-group/OpenAI-YM-20261028/blob/3456e5e75cf8d4cd558955a5c7638fba2d68f464/YM/AnalyticNext/CONNECTED_ACTIVITY_JETS_E6/Room.lean\#L25-L28}{\texttt{Works}}).

\begin{itemize}
\item
  \begin{enumerate}
  \def\labelenumi{(\alph{enumi})}
  \tightlist
  \item
    The conclusions
    \href{https://github.com/yang-mills-group/OpenAI-YM-20261028/blob/3456e5e75cf8d4cd558955a5c7638fba2d68f464/YM/AnalyticNext/CONNECTED_ACTIVITY_JETS_E6/Issue.lean\#L51-L71}{\texttt{GrantBodies}}
    hold:
  \end{enumerate}

  \begin{enumerate}
  \def\labelenumi{\arabic{enumi}.}
  \tightlist
  \item
    the bound of III.10.1 holds at this \(x\) with constant \(\mu\)
    (\href{https://github.com/yang-mills-group/OpenAI-YM-20261028/blob/3456e5e75cf8d4cd558955a5c7638fba2d68f464/YM/Push/SECTOR_MEASURE/Oriented.lean\#L29-L32}{\texttt{PointCharge}});
  \item
    if \(U\) is in tree gauge and, within coarse distance \(B/4\) of the
    bond's parent site, site tags have positive frame weight, fine
    plaquettes have size at most \(t/2\), and coarse mismatches are at
    most \(2Wt\), that link lies within \(Kt\) of the replicated coarse
    link; the Lie logarithm is reliable at that radius
    (\href{https://github.com/yang-mills-group/OpenAI-YM-20261028/blob/3456e5e75cf8d4cd558955a5c7638fba2d68f464/YM/AnalyticNext/CONNECTED_ACTIVITY_JETS_E6/Coat.lean\#L25-L29}{\texttt{ShellAccess}});
  \item
    in every sector whose flags are the actual flags of \(V\)
    (\href{https://github.com/yang-mills-group/OpenAI-YM-20261028/blob/3456e5e75cf8d4cd558955a5c7638fba2d68f464/YM/AnalyticNext/CONNECTED_ACTIVITY_JETS_E6/Currents.lean\#L21}{\texttt{manifest}}),
    the Gaussian precision matrix is symmetric, has exponentially
    decaying entries and a positive spectral lower bound, and the
    covariance formed from its local inverses is positive definite
    (\href{https://github.com/yang-mills-group/OpenAI-YM-20261028/blob/3456e5e75cf8d4cd558955a5c7638fba2d68f464/YM/AnalyticNext/CONNECTED_ACTIVITY_JETS_E6/Shutter.lean\#L110-L117}{\texttt{Rotator}});
  \item
    the Gaussian centring drift has norm at most \({\rm cap}\cdot t\) at
    each fine bond. If the tube also satisfies the compatibility
    condition in (b), then, for every \(\beta>0\), a fluctuation
    satisfying \(|z_j|\le\xi t\sqrt\beta\) on all coordinates within
    distance \(B\) of a site farther than \(30B\) from the seeds has
    site and bond tag masses one, \(|U_p|<t/4\) for plaquettes based
    there, and \(w_b<Wt\) for bonds based there;
  \item
    numerical margins hold: \(4a_ot\) is below the chart radius, and,
    for the window's positive inverse-kernel constants \(\tau_0,C_0\)
    and window bound \(Q_0\ge1\), the quantity
    \(\epsilon_B=e^{-\tau_0B}(Q_0C_0+1)\) satisfies \(\epsilon_B<1\) and
    \(Q_0C_0\epsilon_B<1\).
  \end{enumerate}
\item
  (b) The conclusions
  \href{https://github.com/yang-mills-group/OpenAI-YM-20261028/blob/3456e5e75cf8d4cd558955a5c7638fba2d68f464/YM/AnalyticNext/CONNECTED_ACTIVITY_JETS_E6/Issue.lean\#L73-L88}{\texttt{Invoices}}
  hold for every tube satisfying the free-kernel compatibility condition
  \href{https://github.com/yang-mills-group/OpenAI-YM-20261028/blob/3456e5e75cf8d4cd558955a5c7638fba2d68f464/YM/AnalyticNext/CONNECTED_ACTIVITY_JETS_E6/Voyage.lean\#L25-L27}{\texttt{Stamp}},
  and for every input representation \(K\) of the fine-level density as
  in
  \href{https://github.com/yang-mills-group/OpenAI-YM-20261028/blob/3456e5e75cf8d4cd558955a5c7638fba2d68f464/YM/AnalyticNext/CONNECTED_ACTIVITY_JETS_E6/Intake.lean\#L42-L49}{\texttt{Intake}}.
  The input has the fine-stage support exponent, threshold \(t\),
  coupling \(\beta\in[\beta_F/2,2\beta_F]\), the path data and clips of
  \(\mathscr T\), cluster radius \(B(\beta_F)\), and activities obeying
  the incoming bounds of III.9: \(E\)-norm \(\beta_F^{-11/10}\),
  \(K\)-norm \(e^{-M^{1/5}}\) and loop budget \(2M\), where \(M\) is the
  fine monitor. The conclusions are:

  \begin{enumerate}
  \def\labelenumi{\arabic{enumi}.}
  \tightlist
  \item
    the cost record with gain \(\beta\mu t^2\) per seed and spend
    \(C\,{\rm bill}(x)^5\), where \({\rm bill}=1+B+(\log\beta_F)^2\),
    passes the tree-graph summability test
    \href{https://github.com/yang-mills-group/OpenAI-YM-20261028/blob/3456e5e75cf8d4cd558955a5c7638fba2d68f464/YM/Push/ABSOLUTE_SUMS/Mixed.lean\#L47-L58}{\texttt{Buy}};
  \item
    the letters of the new gas carry weights \(k_i,y_i,E\) satisfying a
    rooted Kotecký--Preiss condition with exponent \(\varepsilon\) and
    total \(\beta_F^{\,r}\)
    (\href{https://github.com/yang-mills-group/OpenAI-YM-20261028/blob/3456e5e75cf8d4cd558955a5c7638fba2d68f464/YM/AnalyticNext/CONNECTED_ACTIVITY_JETS_E6/Weights.lean\#L27-L29}{\texttt{Sales}}).
    For a letter with support \(X_i\) and extra length \(d_i\), this
    means nonnegative weights \(u_i\) with
    \[\sum_{i:o\in X_i}u_i\exp\big(a_*N_B(X_i)+2\varepsilon d_i\big)\le\beta_F^r\qquad\text{for every coarse site }o.\]
    Here \(N_B(H)\) is the minimum number of radius-\(B\) balls whose
    connected union contains \(H\). The letters are the forced
    activities of \(K\), the logarithmic letters, and the warning
    letters. Their weights are respectively \(k_i\), \(y_i\) and \(E\)
    times a fixed patch-dependent warning-amplitude bound; warning
    letters have \(d_i=0\);
  \item
    in every manifest sector, the following bounds hold:

    \begin{itemize}
    \tightlist
    \item
      the old activities are bounded by \(k_i\);
    \item
      the logarithmic letters \(\ell_i\) and the differences
      \(e^{\ell_i}-1\) are bounded by \(y_i\);
    \item
      every packet weight satisfies
      \(|{\rm Pack}|\le\exp\big((C\,{\rm bill}^5-\beta\mu t^2)\,\sharp A\big)\)
      (\href{https://github.com/yang-mills-group/OpenAI-YM-20261028/blob/3456e5e75cf8d4cd558955a5c7638fba2d68f464/YM/AnalyticNext/CONNECTED_ACTIVITY_JETS_E6/Revenue.lean\#L23-L24}{\texttt{Pack}});
    \item
      for every set \(v\) of warning addresses, the Gaussian measure of
      the catastrophe set is at most \(E^{\sharp v}\). The catastrophe
      set consists of the fluctuations \(z\) for which, near the anchor
      of each address in \(v\), some coordinate within distance \(2B\)
      satisfies \(|z_j|\ge\xi t\sqrt\beta\). Every fluctuation that
      triggers the warnings on \(v\) lies in that set.
    \end{itemize}
  \end{enumerate}
\end{itemize}

\emph{Status:} theorem. The target data, fine input bounds and
free-kernel compatibility are hypotheses. The input bounds are supplied
along the RG trajectory (III.9); the compatibility construction is
discussed in §7. - \emph{Quantifier order:} \(L_0\), then \(W\), then
\(\mu\) (uniform in \(D\)), then \((o,C)\) for each patch, then \(x\). -
\emph{Size of the per-seed factor.} Since
\(\beta\mu t^2\ge\frac\mu2e^{2(\log\log\beta_F)^2}\) and
\({\rm bill}^5=O((\log\beta_F)^{10})\), the factor
\(e^{C{\rm bill}^5-\beta\mu t^2}\) is eventually smaller than every
power of \(\beta_F\).

\emph{Lean:}
\href{https://github.com/yang-mills-group/OpenAI-YM-20261028/blob/3456e5e75cf8d4cd558955a5c7638fba2d68f464/YM/AnalyticNext/CONNECTED_ACTIVITY_JETS_E6/Issue.lean\#L94-L185}{\texttt{Target.issue}}.
\href{https://github.com/yang-mills-group/OpenAI-YM-20261028/blob/3456e5e75cf8d4cd558955a5c7638fba2d68f464/YM/AnalyticNext/CONNECTED_ACTIVITY_JETS_E6/IssuePlan.lean\#L53-L69}{\texttt{Target.quotesTogether}}
repackages the output as the bundle
\href{https://github.com/yang-mills-group/OpenAI-YM-20261028/blob/3456e5e75cf8d4cd558955a5c7638fba2d68f464/YM/AnalyticNext/CONNECTED_ACTIVITY_JETS_E6/Batch.lean\#L82-L97}{\texttt{Cash}}
(the activity estimates together with witnesses for the letter weights).

\textbf{Proposition III.10.3 (derivatives of the new activities).} Let
\(P\) and \(\mathcal T\) be as in III.10.2. There is \(L_0\) with the
following property. Fix \(L\ge L_0\), a patch \(D\) satisfying the
branch condition, an order \(i\in\mathbb N\), a sector cutoff \(\chi\),
\(W,\mu>0\), window constants \(o\), and chart constants \((A,f)\)
satisfying
\href{https://github.com/yang-mills-group/OpenAI-YM-20261028/blob/3456e5e75cf8d4cd558955a5c7638fba2d68f464/YM/AnalyticNext/CONNECTED_ACTIVITY_JETS_E6/Eider.lean\#L32-L37}{\texttt{YM.ConnectedActivityE6.Blaze.Eider.Canary}}.
The latter require \(A,f>0\), \(\log(\exp u)=u\) for \(|u|\le2A\),
\(\exp(\log(\exp u\,g))=\exp u\,g\) for \(|u|\le A\) and
\(\|g-1\|\le f\), and analyticity of \(M\mapsto\log(\exp u\,M)\) for
\(|u|\le2A\), \(\|M-1\|\le f\). Fix a function \(\mathcal V\) of rate
\(-1/2\) that is eventually nonnegative (Lemma 4.15).

There are \(C\ge0\), functions \(q,\rho,p\) of rates \(-1/2,0,0\),
respectively, and \(l\in\mathbb N\) such that, for all large \(x\), the
following holds for every step on \(D\) satisfying the compatibility
condition, every fine representation \(R\), and every bound bundle \(H\)
of III.10.2 at \((W,x,\mu)\). One has \(8a_ot\le A/2\) and a
derivative-bound structure \(\Xi\)
(\href{https://github.com/yang-mills-group/OpenAI-YM-20261028/blob/3456e5e75cf8d4cd558955a5c7638fba2d68f464/YM/AnalyticNext/CONNECTED_ACTIVITY_JETS_E6/Das.lean\#L27-L46}{\texttt{Utopia}})
with the following properties in every sector at \(x\), of width
\(o.K\), cutoff \(\chi\), and with the actual coarse flags.

\begin{itemize}
\tightlist
\item
  Its nonpacket weights satisfy the rooted bound of III.10.2(b) with
  total \(\beta_F^r\). Its catastrophe probabilities, with trigger
  \(\xi/2\), are at most \((\mathrm{rad}^2)^{\sharp v}\).
\item
  Let \(V(y)\) be a family of coarse fields on a finite-dimensional real
  normed parameter space, based at the sector field \(V(y_0)=V_0\).
  Suppose each link matrix is \(C^i\) at \(y_0\) and
  \[\|D^m V_e(y_0)\|\le c_G\mathcal V(x)^m\qquad(0\le m\le i),\] where
  \(c_G\) is the fixed matrix-norm bound in
  \href{https://github.com/yang-mills-group/OpenAI-YM-20261028/blob/3456e5e75cf8d4cd558955a5c7638fba2d68f464/YM/AnalyticNext/CONNECTED_ACTIVITY_JETS_E6/Tint.lean\#L21-L23}{\texttt{Mode}}.
  Transport the fluctuation coordinates compatibly with this family: the
  physical and Gaussian coordinates represent the same field, their
  physical displacement has all derivatives through order \(i\) bounded
  by \(q(x)\), and, near \(y_0\), the resulting fine field agrees with
  the held fine field on every cold bond and on every hot bond whose
  physical coordinate has norm at most \(8a_ot\)
  (\href{https://github.com/yang-mills-group/OpenAI-YM-20261028/blob/3456e5e75cf8d4cd558955a5c7638fba2d68f464/YM/AnalyticNext/CONNECTED_ACTIVITY_JETS_E6/Eskim.lean\#L86-L95}{\texttt{YM.ConnectedActivityE6.Blaze.Silica.Clutch}}).
  In a region disjoint from the coarse flags, each letter \(a_j\) is
  \(C^i\) at \(y_0\) and satisfies
  \[\|D^m a_j(y_0)\|\le\mathrm{cost}_j\big(\rho(x)\mathrm{span}_j\big)^m\qquad(0\le m\le i)\]
  (\href{https://github.com/yang-mills-group/OpenAI-YM-20261028/blob/3456e5e75cf8d4cd558955a5c7638fba2d68f464/YM/AnalyticNext/CONNECTED_ACTIVITY_JETS_E6/Brush.lean\#L44-L47}{\texttt{Brush}}).
  Packet spans are one; other spans include their support and path
  lengths. The costs use the packet estimate and the nonpacket weights
  before extracting the catastrophe-probability factor.
\end{itemize}

The same \(\Xi\) satisfies an integrated estimate
(\href{https://github.com/yang-mills-group/OpenAI-YM-20261028/blob/3456e5e75cf8d4cd558955a5c7638fba2d68f464/YM/PhysicalNext/RG_NEXT_COMPLETE_REPEATABLE_NONLINEAR_STEP/Daisy/Charm.lean\#L48-L55}{\texttt{Herd}}).
Let \(\mathfrak p\) be a finite molecule with connected support-overlap
graph and total support disjoint from the flags. Require that its packet
stores are connected with pairwise disjoint radius-\(hB\)
neighbourhoods, that its old forced supports are pairwise disjoint, and
that every site in an old forced inventory touches a fine-plaquette seed
in the union of the packet stores. Along a family \(V(y)\) as above, the
prescribed transported integrand is \(C^i\) at \(y_0\) and has every
derivative of order at most \(i\) bounded by an integrable function
\(g\). With \(d\nu=d\lambda_F\,d\gamma\), the product of tree-slice Haar
measure and Gaussian fluctuation measure,
\[\int g\,d\nu\le |p(x)|\prod_{j\in\mathfrak p}\mathrm{fare}_j\,(1+\mathrm{shape}(\mathfrak p))^l.\]
The fares include a factor \(\mathrm{rad}\) for each distinct warning
address. For one measurable, gauge-invariant source letter of amplitude
at most \(w_M\), local with its eligibility tests included, the same
estimate replaces its fare by \(w_M\). A forced source must vanish
outside the original eligibility set; a regular source occupies an
unused regular-letter address. The remaining warning and shape factors
are unchanged.

\emph{Status:} theorem, conditional on the input bound bundle and
compatibility. The constants and functions are uniform in the volume.
\href{https://github.com/yang-mills-group/OpenAI-YM-20261028/blob/3456e5e75cf8d4cd558955a5c7638fba2d68f464/YM/AnalyticNext/CONNECTED_ACTIVITY_JETS_E6/Captains.lean\#L27-L36}{\texttt{Captain}}
and
\href{https://github.com/yang-mills-group/OpenAI-YM-20261028/blob/3456e5e75cf8d4cd558955a5c7638fba2d68f464/YM/PhysicalNext/RG_NEXT_COMPLETE_REPEATABLE_NONLINEAR_STEP/Daisy/Sparks.lean\#L36-L61}{\texttt{rocketsC}}
also control iterated left derivatives in spatially weighted directions.

\emph{Lean:}
\href{https://github.com/yang-mills-group/OpenAI-YM-20261028/blob/3456e5e75cf8d4cd558955a5c7638fba2d68f464/YM/AnalyticNext/CONNECTED_ACTIVITY_JETS_E6/Das.lean\#L49-L102}{\texttt{Utopia.whittle}}
(letter jets),
\href{https://github.com/yang-mills-group/OpenAI-YM-20261028/blob/3456e5e75cf8d4cd558955a5c7638fba2d68f464/YM/PhysicalNext/RG_NEXT_COMPLETE_REPEATABLE_NONLINEAR_STEP/Daisy/Charm.lean\#L58-L179}{\texttt{timing}}
(integrable majorants).

\textbf{Proposition III.10.4 (Gaussian sector of the step).} Fix a patch
\(D\), a request \(Q\) for which \(D\) satisfies the branch condition,
an aperture \(B_0>0\), a sector cutoff \(\chi\), and
\(k,p\in\mathbb N\). There are \(\sigma,d>0\) and \(\pi\ge1\)
(\href{https://github.com/yang-mills-group/OpenAI-YM-20261028/blob/3456e5e75cf8d4cd558955a5c7638fba2d68f464/YM/PhysicalNext/V9_JOINT_ENDPOINT_LOAD/Costs.lean\#L34-L42}{\texttt{Charges}})
such that the following holds for all large \(x\).

Take \(0<\beta\le2\beta_F\), a budget point of power \(p\) and threshold
\(t_c(x)\), a step on \(D\) with \(c\ge100(1+B)\) admissible for \(Q\),
a forest, tags, and a store \(A\) whose flags equal the flags of \(V\)
on the whole coarse torus. Freeze \(A\), the forest and tags while
differentiating. Let \(\Gamma\) be the sector precision matrix,
\(\widetilde\Gamma\) its local approximate inverse, \(m\) the smooth
sector mean at aperture \(B_0t\), and \(h=-g^{-1}\) times its
stationarity residual, with \(g=\beta^{-1/2}\). These are the objects in
\href{https://github.com/yang-mills-group/OpenAI-YM-20261028/blob/3456e5e75cf8d4cd558955a5c7638fba2d68f464/YM/PhysicalNext/V9_JOINT_ENDPOINT_LOAD/Costs.lean\#L60-L74}{\texttt{Value}}.
Then \(\Gamma\) and \(\widetilde\Gamma\) are positive definite, and
\(m\) equals the stationary logarithmic mean for every frozen-link
input.

For a coarse site \(o\) and \(0\le m\le k\), let \(\mathcal D^m\) be an
iterated left derivative in \(m\) directions satisfying
\(|A_e|\le t_c(x)(1+\mathrm{dist}(e,o))^p\). If \(x_i\) is the coarse
location of coordinate \(i\), put \[G_i=(1+\mathrm{dist}(x_i,o))^{pk}.\]
The bounds are
\[|\mathcal D^m\Gamma_{ij}|\le\pi e^{-\sigma\mathrm{dist}(x_i,x_j)}G_i^m,\qquad |\mathcal D^m(\Gamma\widetilde\Gamma-1)_{ij}|\le e^{-dB}\pi e^{-\sigma\mathrm{dist}(x_i,x_j)}G_i^m,\]
\[|\mathcal D^m m_i|\le\pi tG_i^m,\qquad |\mathcal D^m h_i|\le e^{-dB}\pi G_i^m.\]
Thus the spatial exponent at derivative order \(m\) is \(pkm\).

\emph{Status:} theorem.
\href{https://github.com/yang-mills-group/OpenAI-YM-20261028/blob/3456e5e75cf8d4cd558955a5c7638fba2d68f464/YM/PhysicalNext/V9_JOINT_ENDPOINT_LOAD/Supply.lean\#L72-L98}{\texttt{JointLoad.stage}}
also supplies a covariance cap and a trigger-tail estimate; its scales
are given in §7.

\emph{Lean:}
\href{https://github.com/yang-mills-group/OpenAI-YM-20261028/blob/3456e5e75cf8d4cd558955a5c7638fba2d68f464/YM/PhysicalNext/V9_JOINT_ENDPOINT_LOAD/Costs.lean\#L92-L157}{\texttt{JointLoad.capital}}.

\textbf{Proposition III.10.5 (schedule).} - (a) Let \(P\) be a sheet of
pre-\(L\) data (III.2) and \(K\) a price map. There are floors \(k_i\)
with the following property. Take block factors \(L_i\ge k_i\), any
protection, any delays \(V\) and any \(m\ge0\). Then there are delays
\(W\), differing from \(V\) only by larger entry couplings
\(\beta^*_i\), and numbers \(F_{ij}\ge m\), such that for every pair of
stages \(P.\mathrm{first}\le j\le i\le j+1\) the bound of III.10.1 holds
at every coarse coupling \(x\ge\beta^*_i\), for the request with fine
stage \(i\), coarse stage \(j\) and step index \(i\), with block factor
\(L_i\), aisle multiplier \(F_{ij}\) and some \(\mu_{ij}>0\) uniform in
\(x\), period and volume
(\href{https://github.com/yang-mills-group/OpenAI-YM-20261028/blob/3456e5e75cf8d4cd558955a5c7638fba2d68f464/YM/Pressure/Seeds.lean\#L161-L163}{\texttt{Ringing}}).
- (b) If \(G\) is simple, then for every window datum of III.2 (entry
and plateau requirements) there is a level schedule that meets all
requirements of III.2
(\href{https://github.com/yang-mills-group/OpenAI-YM-20261028/blob/3456e5e75cf8d4cd558955a5c7638fba2d68f464/YM/Price/Intent.lean\#L145-L162}{\texttt{Paid}})
and whose periods are divisible: for every side \(c\) and every \(z\)
there is a period index \(p\ge z\) with \(c\mid{\rm period}(p)\). - (c)
Along a directed trajectory satisfying (a), fix an adjacent stage pair
\(i=J(m)\), \(j=J(m+1)\) with \(P.\mathrm{first}\le i\). There is
\(\mu>0\), uniform in the cutoff \(N\), level \(m\) and period, such
that for \(m<N\) at this pair and \(x=\hat\beta_m\), one has
\(\mu\,t_c(\hat\beta_{m+1})^2\,\sharp A\le\mathcal E_A\) whenever the
fine field is in the chosen tree gauge and the store is consistent at
multiplier \(F_{ji}\).

\emph{Status:} theorems. Part (a) is unconditional
(\href{https://github.com/yang-mills-group/OpenAI-YM-20261028/blob/3456e5e75cf8d4cd558955a5c7638fba2d68f464/YM/Next/FLUCTUATIONS/Grants.lean\#L32-L37}{\texttt{pressure\_pricing}}
applies
\href{https://github.com/yang-mills-group/OpenAI-YM-20261028/blob/3456e5e75cf8d4cd558955a5c7638fba2d68f464/YM/Bid/Energy.lean\#L173-L202}{\texttt{Bid.pricing\_order}}
to III.10.1). Divisibility is what lets torus tilings fit in
III.12--III.13.

\emph{Lean:}
\href{https://github.com/yang-mills-group/OpenAI-YM-20261028/blob/3456e5e75cf8d4cd558955a5c7638fba2d68f464/YM/PhysicalNext/RG_NEXT_COMPLETE_REPEATABLE_NONLINEAR_STEP/Pane/Divisors.lean\#L26-L135}{\texttt{Wait.engine}},
\href{https://github.com/yang-mills-group/OpenAI-YM-20261028/blob/3456e5e75cf8d4cd558955a5c7638fba2d68f464/YM/Pressure/Seeds.lean\#L166-L181}{\texttt{notation\_ticket}}.

\textbf{Proposition III.10.6 (start-up fusion H01--H04; markers
I01--I04).}

\begin{enumerate}
\def\labelenumi{(\Alph{enumi})}
\setcounter{enumi}{7}
\tightlist
\item
  \emph{Setting.} Fix a budget point and:
\end{enumerate}

\begin{itemize}
\tightlist
\item
  a \emph{pipe} \(X\): regular and forced addresses, and supports of
  logarithmic letters;
\item
  a \emph{tap} \(d\): regular activities, forced activities \(K_i\) and
  logarithmic letters \(\ell_i\).
\end{itemize}

The \emph{fused} gas
(\href{https://github.com/yang-mills-group/OpenAI-YM-20261028/blob/3456e5e75cf8d4cd558955a5c7638fba2d68f464/YM/Ovens/Coffee.lean\#L78}{\texttt{fused}})
groups the forced activities and the factors \(e^{\ell_i}-1\) into
connected clusters (pellets). Each pellet contains at least one forced
activity, and the forced supports in a pellet are pairwise disjoint.

\emph{Hypotheses.} Fix a bijective enumeration of the pellets. Let
\(\kappa,\gamma\ge0\). - \(d\) is \emph{digested}
(\href{https://github.com/yang-mills-group/OpenAI-YM-20261028/blob/3456e5e75cf8d4cd558955a5c7638fba2d68f464/YM/Ovens/Coffee.lean\#L81-L90}{\texttt{Digests}}):
the gas is valid with specified regular and forced bounds; the
logarithmic family has a summable uniform majorant, and each letter is
measurable, gauge invariant, local with the flag tests included,
vanishes on healthy configurations, and is translation and
lattice-symmetry covariant. - \(d\) is \emph{bought} with constants
\(\kappa,\gamma\)
(\href{https://github.com/yang-mills-group/OpenAI-YM-20261028/blob/3456e5e75cf8d4cd558955a5c7638fba2d68f464/YM/Ovens/Coffee.lean\#L94-L98}{\texttt{Buys}}):
\(|K_i|\le e^{-\kappa}f_i\) and \(|e^{\ell_i}-1|\le f_i\), where
\(f\ge0\) satisfies the rooted Kotecký--Preiss condition at support
exponent one larger than the budget exponent, with total \(\gamma\), and
\((2B+1)^4\gamma\le1\).

\emph{Conclusions.} - (H01) For every configuration \(V\), the fused gas
equals the original gas times \(\exp(\sum_{i=0}^\infty\ell_i(V))\). -
(H02) The fused activities satisfy the \(K\)-bound with any price
\(v\ge e^{-\kappa}(e^\gamma-1)\). - (H03) The fused system is a forced
family and admits the lattice-symmetry moves. - (H04) Fix the paths,
clip, budget point, addresses and pellet enumeration. Suppose the
original representations form a continuous curve on a subset of
\(\mathbb R\), each forced or exponentiated logarithmic letter varies
continuously in the supremum norm over configurations, and a single
nonnegative rooted majorant as above works throughout the parameter set.
Then the fused representations form a continuous curve in the same
representation norm. This is a fixed-volume assertion.

\begin{enumerate}
\def\labelenumi{(\Roman{enumi})}
\tightlist
\item
\end{enumerate}

\begin{itemize}
\tightlist
\item
  (I01) Let \(B\ge8\), let \(F\) be a set of at least two markers with
  pairwise separation at least \(B^2\), and let \(F\subset H\). Then
  \(B\,\sharp F\le64\,N_B(H)\), where \(N_B(H)\) is the minimum number
  of radius-\(B\) balls with connected union covering \(H\). The markers
  lie on a finite periodic lattice with its max-distance; no nonwrapping
  hypothesis is required.
\item
  (I02) Fix a representation with budget order \(n\). Suppose marker
  bands satisfy
  \href{https://github.com/yang-mills-group/OpenAI-YM-20261028/blob/3456e5e75cf8d4cd558955a5c7638fba2d68f464/YM/Bid/Fuels.lean\#L66-L70}{\texttt{YM.Mix.Stubs.Holds}}
  at \(r>0\): valid local addresses, \(r\) times every unit amplitude at
  most one, \(B\ge8\) and separation at least \(B^2\) for each pair of
  distinct markers, factorially weighted derivative majorants, and
  separate rooted budgets at most one for single-marker and
  multiple-marker terms. For the multiple-marker budget one may instead
  use positive weights \(t_v\) and an extra exponent \(\delta\ge0\),
  with \(r\le t_v e^{\delta B/64}\) for every marker \(v\) and a rooted
  budget at most one using support exponent increased by \(\delta\).
  Then, at weight \(r/n!\), the bands satisfy
  \href{https://github.com/yang-mills-group/OpenAI-YM-20261028/blob/3456e5e75cf8d4cd558955a5c7638fba2d68f464/YM/Overlay/Table.lean\#L56-L65}{\texttt{YM.Overlay.Bands.Radial}}:
  positive weights, unit-amplitude bound one, and rooted ordinary
  derivative and forced-entry bounds at most two each.
\item
  (I03) Let a family satisfy the marker schedule condition:
  \(B_{r+1}^2\le2B_r^2\) for every \(r\), and
  \((L_rB_r^2)^4\mathcal F_{r+1}=o(\mathrm{scale}_{r+1})\). Here \(a_r\)
  is the lattice spacing, \(L_r=a_r/a_{r+1}\), and the vacuum bound is
  \[\mathcal F_r=2(2+4P_r)+1+6(\alpha_r+P_r^2)+4d_G\big(\log16+\tfrac12\log\hat\beta_r\big),\]
  with \(P_r\) the stage's path budget, \(\alpha_r\) its direct clip
  coefficient and \(d_G=\dim_{\mathbb R}M_m(\mathbb C)\) for the
  defining matrix representation. This is the hypothesis of
  \href{https://github.com/yang-mills-group/OpenAI-YM-20261028/blob/3456e5e75cf8d4cd558955a5c7638fba2d68f464/YM/Overlay/Prior.lean\#L15-L18}{\texttt{YM.Stock.Tabs}}.
  Fix \(k\ge1\) and levels \(m\le i\) such that \(kB_i^2\le a_m/a_i\)
  (an open stop). For every \(N\ge i\) and every period, the centred
  copies at level \(N\) are the embedded sites from level \(m\); their
  parents agree at consecutive levels from \(m\) to \(N\); and distinct
  embedded sites are separated by at least \(kB_r^2\) for every
  \(i\le r\le N\).
\item
  (I04) Assume \(G\) is simple. For every family, every \(p>0\) and
  \(r\in\mathbb N\), if \({\rm order}_i!\le{\rm columns}_i\) for all
  \(i\), then the widths
  \[w_i=\frac{\hat\beta_i^{-r/2}}{p\,{\rm columns}_i\prod_{m<i}{\rm columns}_{m+1}}\]
  are positive, and for every \(\varepsilon>0\) eventually
  \(5-\log(w_i/{\rm order}_i!)\le\varepsilon\,{\rm scale}_i\).
\end{itemize}

\emph{Status:} theorems. In the main-theorem closure: - H01--H04 are
used by the start-up representation at the finest level; - I01, I03 and
I04 are used by the replicated-marker constructions; - the assemblies
\href{https://github.com/yang-mills-group/OpenAI-YM-20261028/blob/3456e5e75cf8d4cd558955a5c7638fba2d68f464/YM/Bid/Kitchen.lean\#L56-L66}{\texttt{Oust.begins}}
and
\href{https://github.com/yang-mills-group/OpenAI-YM-20261028/blob/3456e5e75cf8d4cd558955a5c7638fba2d68f464/YM/Bid/Fuels.lean\#L146-L174}{\texttt{Bid.installs}}
are not in the closure.

\emph{Lean:} - Statements:
\href{https://github.com/yang-mills-group/OpenAI-YM-20261028/blob/3456e5e75cf8d4cd558955a5c7638fba2d68f464/YM/Bid/Kitchen.lean\#L15-L16}{\texttt{H01}},
\href{https://github.com/yang-mills-group/OpenAI-YM-20261028/blob/3456e5e75cf8d4cd558955a5c7638fba2d68f464/YM/Bid/Kitchen.lean\#L20-L21}{\texttt{H02}},
\href{https://github.com/yang-mills-group/OpenAI-YM-20261028/blob/3456e5e75cf8d4cd558955a5c7638fba2d68f464/YM/Bid/Kitchen.lean\#L23-L24}{\texttt{H03}},
\href{https://github.com/yang-mills-group/OpenAI-YM-20261028/blob/3456e5e75cf8d4cd558955a5c7638fba2d68f464/YM/Bid/Kitchen.lean\#L33-L35}{\texttt{H04}},
\href{https://github.com/yang-mills-group/OpenAI-YM-20261028/blob/3456e5e75cf8d4cd558955a5c7638fba2d68f464/YM/Bid/Fuels.lean\#L124-L126}{\texttt{I01}}--\href{https://github.com/yang-mills-group/OpenAI-YM-20261028/blob/3456e5e75cf8d4cd558955a5c7638fba2d68f464/YM/Bid/Fuels.lean\#L140-L142}{\texttt{I04}}.
- Proofs:
\href{https://github.com/yang-mills-group/OpenAI-YM-20261028/blob/3456e5e75cf8d4cd558955a5c7638fba2d68f464/YM/Work/H01.lean\#L245-L251}{\texttt{H01.gathered}},
\href{https://github.com/yang-mills-group/OpenAI-YM-20261028/blob/3456e5e75cf8d4cd558955a5c7638fba2d68f464/YM/Work/H02.lean\#L17-L80}{\texttt{H02.answer}},
\href{https://github.com/yang-mills-group/OpenAI-YM-20261028/blob/3456e5e75cf8d4cd558955a5c7638fba2d68f464/YM/Work/H03.lean\#L138-L209}{\texttt{H03.fused\_forcedFamily}},
\href{https://github.com/yang-mills-group/OpenAI-YM-20261028/blob/3456e5e75cf8d4cd558955a5c7638fba2d68f464/YM/Work/H04.lean\#L155-L161}{\texttt{H04.answer}},
\href{https://github.com/yang-mills-group/OpenAI-YM-20261028/blob/3456e5e75cf8d4cd558955a5c7638fba2d68f464/YM/Work/I01.lean\#L19-L82}{\texttt{I01.answer}},
\href{https://github.com/yang-mills-group/OpenAI-YM-20261028/blob/3456e5e75cf8d4cd558955a5c7638fba2d68f464/YM/Work/I02.lean\#L159-L180}{\texttt{I02.answer}},
\href{https://github.com/yang-mills-group/OpenAI-YM-20261028/blob/3456e5e75cf8d4cd558955a5c7638fba2d68f464/YM/Work/I03.lean\#L43-L63}{\texttt{I03.answer}},
\href{https://github.com/yang-mills-group/OpenAI-YM-20261028/blob/3456e5e75cf8d4cd558955a5c7638fba2d68f464/YM/Work/I04.lean\#L76-L129}{\texttt{I04.answer}}.

\subsubsection{2. What it uses}\label{what-it-uses-12}

\textbf{From other units.} - \emph{III.1 (block-spin step and gauge
fixing).} - The step data: patch, kernel choices, tree gauge, predictor
\(Q\) and tags. - The definitions of request, admissible tube, store,
consistency and action gap, including the identity
\href{https://github.com/yang-mills-group/OpenAI-YM-20261028/blob/3456e5e75cf8d4cd558955a5c7638fba2d68f464/YM/Pressure/Rows.lean\#L145-L163}{\texttt{Tube.locates}}.
- The branch condition. It is proved for the kernel choices of the
construction by
\href{https://github.com/yang-mills-group/OpenAI-YM-20261028/blob/3456e5e75cf8d4cd558955a5c7638fba2d68f464/YM/Next/FOUNDATIONS/Primary.lean\#L70-L71}{\texttt{Foundations.branches}}
(from
\href{https://github.com/yang-mills-group/OpenAI-YM-20261028/blob/3456e5e75cf8d4cd558955a5c7638fba2d68f464/YM/Branch/Pick.lean\#L137-L142}{\texttt{Axial.viable}})
and is carried by every sheet
(\href{https://github.com/yang-mills-group/OpenAI-YM-20261028/blob/3456e5e75cf8d4cd558955a5c7638fba2d68f464/YM/Price/Intent.lean\#L26-L35}{\texttt{Sheet.has}}).
- With usable tags at the base site, the curvature bound for the
predictor
\href{https://github.com/yang-mills-group/OpenAI-YM-20261028/blob/3456e5e75cf8d4cd558955a5c7638fba2d68f464/YM/Pressure/Cost.lean\#L210-L252}{\texttt{qdc}}:
- \(|Q_p|\le C_QL^{-2}\,{\rm port}(x,U)\); -
\(|V_p|\le\sum_{b\subset\partial p}w_b+C_QL^{-2}\,{\rm port}(x,U)\).

\begin{verbatim}
Here ${\rm port}(x,U)$ is the sum of the fine $|U_q|$ over the macro box of the coarse site $x$. Block averaging reduces curvature by $L^{-2}$.
\end{verbatim}

\begin{itemize}
\tightlist
\item
  \emph{III.2 (constants and schedule):} sheets, block factors,
  protection, delays, bar and waits. The target data precede the block
  choice; the per-step window constants are chosen in III.10.2.
\item
  \emph{III.3 and III.4 (free kernels, Gaussian integration):} the
  towers, their flat identities and exponential decay (III.3 attribution
  probable); the Gaussian tail bound
  \href{https://github.com/yang-mills-group/OpenAI-YM-20261028/blob/3456e5e75cf8d4cd558955a5c7638fba2d68f464/YM/AnalyticNext/CONNECTED_ACTIVITY_JETS_E6/Many.lean\#L89-L101}{\texttt{rectangle}},
  which for covariance at most \(M\) reads
  \(\Pr(|\xi_i|>s\ \forall i\in X)\le2^{|X|}e^{(-\nu s+M\nu^2/2)|X|}\)
  for every \(\nu>0\).
\item
  \emph{III.5 (cluster expansion):} polymer norms, budget points,
  representations of the effective density, Kotecký--Preiss sums.
\item
  \emph{III.6--III.8 (stationary branch, matching, connected
  activities):} the stationary mean and its jets, the quartic matching
  of the path data, and the locality of connected activities.
\item
  \emph{III.9:} the slope \(b_G\), the thresholds, the monitor and the
  incoming allowances. III.9 supplies the fine-level density
  representation and consumes the activity bounds; the norm conversion
  remains to be written out (§7).
\item
  \emph{I.1, II.2, III.16:} lattice, Haar measure, chordal distance,
  polymer activities, Taylor jets.
\end{itemize}

\textbf{Standard mathematics.} The analytic method is Balaban's
lattice-gauge renormalization group. Relevant references are
\emph{Renormalization group approach to lattice gauge field theories},
\href{https://doi.org/10.1007/BF01215223}{I, CMP 109 (1987), 249--301},
\href{https://doi.org/10.1007/BF01239022}{II, CMP 116 (1988), 1--22},
and \emph{Large field renormalization},
\href{https://doi.org/10.1007/BF01257412}{I, CMP 122 (1989), 175--202},
\href{https://doi.org/10.1007/BF01238433}{II, CMP 122 (1989), 355--392}.
The summability estimates use rooted polymer bounds of
\href{https://doi.org/10.1007/BF01211762}{Kotecký--Preiss, CMP 103
(1986), 491--498}, tree-graph estimates, Schur's test, Gaussian
exponential moments, and Taylor estimates on compact charts. Exact
correspondences with Balaban's lemmas remain open (§7).

\textbf{Imported constants and assumptions.} - The request and target
data precede \(L\). The window constants in III.10.2 are chosen after
\(L,W\) and the patch, before the threshold in \(x\). Simplicity of
\(G\) is required by III.10.5(b) and I04. - The following are internal
choices of the proof: the exponent \(49/100\) in \(T\); the factor
\(10^{-5}\) in \(R\); the hysteresis thresholds \(t/4\), \(t/2\),
\(Wt\), \(2Wt\); and the radii \(5B\), \(2B\), \(4B/5\), \(2B/5\),
\(30B\). - The incoming bounds on the fine activities
(\(\beta_F^{-11/10}\), \(e^{-M^{1/5}}\), \(2M\)) are hypotheses supplied
by III.9.

\subsubsection{3. Strategy}\label{strategy-12}

The method combines large-field suppression with Gaussian integration in
the small-field region. Seeds record the fine plaquettes, mismatches and
coarse flags that exceed their thresholds. The deterministic estimate
gives an action cost \(\kappa t^2\) per seed. Since
\(\beta\ge\beta_F/2\), its Boltzmann factor is at most
\(e^{-(\kappa/2)e^{2(\log\log\beta_F)^2}}\). This pays every fixed
polylogarithmic entropy loss and is smaller than every inverse power of
\(\beta_F\), while \(t=\beta_F^{-1/2+o(1)}\) remains perturbatively
small.

An exact completed-square decomposition separates the fine, interior,
coarse and gauge contributions from two kernel defects. Large \(L\)
makes the fine prediction small; large \(W\) separates mismatch seeds
from the interior prediction. Coarse plaquettes above a fixed threshold
lie near spent seeds, whose larger cost \(tT\) pays the affected-cell
losses. Away from those seeds, local gauge fixing and the flat tower
identities reduce the defects to chart errors and exponentially small
kernel tails. Summation over a radius-\(B\) neighbourhood still leaves
an error \(o(t^2)\) per seed.

The resulting gas has packet, old forced, logarithmic and warning
letters. The energy bound and annulus estimate control packets; Gaussian
tails control warnings. Rooted support bounds and tree-graph estimates
sum connected molecules. Derivatives are taken with sector data fixed
and fluctuation coordinates transported compatibly. The Leibniz rule
controls products, and Cauchy--Schwarz retains the square root of the
simultaneous-warning probability when integrating their jets.

\subsubsection{4. Steps}\label{steps-12}

Unless a lemma gives different hypotheses, the energy estimates in
§§4.2--4.6 use the step, tube, forest, fields, tags and consistent store
of III.10.1. A sufficiently large threshold in \(x\) may depend on the
fixed request, \(L,W\), and any additional parameters already chosen,
but never on the period. Row and column bounds refer to sums of operator
norms; vector bounds are pointwise in the coordinate index. Distances
are native torus max-distances.

\paragraph{4.1 The algebraic identity}\label{the-algebraic-identity}

\textbf{Lemma 4.1 (completed squares).} Notation: - \(v=v(V)\) are the
coarse coordinates
(\href{https://github.com/yang-mills-group/OpenAI-YM-20261028/blob/3456e5e75cf8d4cd558955a5c7638fba2d68f464/YM/Pressure/Rows.lean\#L104}{\texttt{vision}}).
For each coarse plaquette \(q\) they are the pair
\((\sqrt{\alpha_c}\,{\rm curv}(V_q),\ \ell_q(V))\), with the first
component multiplied by \(\min(1,t_c/|V_q|)\), interpreted as one when
\(|V_q|=0\), and the second set to zero unless \(q\) is kept. Let
\(d_F,d_c\) denote the clipped densities of §1. - \(X(U)\) and
\(X^c(V)\) are the unthinned fine and coarse coordinates. -
\(Z_b=\sqrt a\,{\rm curv}(V_bQ_b^{-1})\).

The rows are: - \emph{fine plaquette \(p\):}
\(d_F(U,p)+\tfrac12|(\Phi v)_p|^2-\langle(\Phi v)_p,X_p\rangle\); -
\emph{coarse bond \(b\):}
\(\tfrac a2w_b^2+\tfrac12|(\Phi'v)_b|^2-\langle(\Phi'v)_b,Z_b\rangle\);
- \emph{coarse plaquette \(q\):}
\(\langle v_q,X^c_q\rangle-\tfrac12|v_q|^2-d_c(V,q)\); - \emph{gauge
coordinate \(n\):} \(\tfrac12|(Nv)_n|^2\); - \emph{defects at \(q\):} -
the linear-defect contribution
\(\langle v_q,(\Phi^*X+\Phi'^*Z-X^c)_q\rangle\); - the
completeness-defect contribution
\(\tfrac12\langle v_q,(v-(\Phi^*\Phi+\Phi'^*\Phi'+N^*N)v)_q\rangle\).

Their sum over the torus is
\(\mathcal A_F(U)+\frac a2\sum_bw_b^2-\mathcal A_c(V)\). The functional
\(\mathcal E_A\) keeps only the rows located in \(D_{2B}(A)\).

\emph{Proof:} adjointness for the three kernels; all cross terms cancel.
The cancellation is algebraic and works for any coarse coordinate vector
\(v\) if the rows are defined with that vector.

\emph{Lean:}
\href{https://github.com/yang-mills-group/OpenAI-YM-20261028/blob/3456e5e75cf8d4cd558955a5c7638fba2d68f464/YM/Pressure/Rows.lean\#L145-L163}{\texttt{Tube.locates}};
rows
\href{https://github.com/yang-mills-group/OpenAI-YM-20261028/blob/3456e5e75cf8d4cd558955a5c7638fba2d68f464/YM/Pressure/Rows.lean\#L122-L124}{\texttt{workF}},
\href{https://github.com/yang-mills-group/OpenAI-YM-20261028/blob/3456e5e75cf8d4cd558955a5c7638fba2d68f464/YM/Pressure/Rows.lean\#L126-L127}{\texttt{workI}},
\href{https://github.com/yang-mills-group/OpenAI-YM-20261028/blob/3456e5e75cf8d4cd558955a5c7638fba2d68f464/YM/Pressure/Rows.lean\#L129-L131}{\texttt{workC}},
\href{https://github.com/yang-mills-group/OpenAI-YM-20261028/blob/3456e5e75cf8d4cd558955a5c7638fba2d68f464/YM/Pressure/Rows.lean\#L136}{\texttt{workN}},
\href{https://github.com/yang-mills-group/OpenAI-YM-20261028/blob/3456e5e75cf8d4cd558955a5c7638fba2d68f464/YM/Pressure/Rows.lean\#L133-L134}{\texttt{workL}},
\href{https://github.com/yang-mills-group/OpenAI-YM-20261028/blob/3456e5e75cf8d4cd558955a5c7638fba2d68f464/YM/Pressure/Rows.lean\#L138-L140}{\texttt{workD}}.

\paragraph{4.2 Packing and geometry}\label{packing-and-geometry}

\textbf{Lemma 4.2 (packing; E01).} For every \(\delta>0\) there is
\(N\), depending on \(P\) and \(\delta\) but not on \(L\) or \(W\), such
that for all \(L\), all \(W>0\), all large \(x\) and all admissible data
with \(A\) consistent: - (i) for every coarse site \(y\),
\({\rm port}(y,U)=\sum_qm(y,q)|U_q|\) with nonnegative integer charges
\(m\) satisfying \(\sum_ym(y,q)\le N\), and
\(m(y,q)>0\Rightarrow{\rm dist}(y,\bar q)\le N\); - (ii) every coarse
plaquette of \(D_{2B}(A)\) with \(|V_q|>\delta\) (a tide) lies within
\(4B/5\) of \(A^{\rm sp}\) (the hit set); - (iii)
\(\sharp A_\flat\le4\,\sharp\{q\in D_{2B}:|V_q|>t_c\}\); - (iv)
\(\sharp\{\text{tides}\}\le N\sharp A^{\rm sp}\); - (v) the coarse
plaquettes of \(D_{2B}(A)\) number at most \(NB^4\sharp A\), and the hit
set at most \(NB^4\sharp A^{\rm sp}\).

\emph{Proof:} the macro box has radius \((j+7)L\). Counting possible
parent sites gives a bounded total multiplicity for each fine plaquette.
If no spent seed is within distance \(j+8\) of a coarse plaquette,
consistency bounds every relevant input by \(T\); the
predictor-curvature bound then gives \(|V_q|\le C(L,W)T\le\delta\) for
large \(x\). Thus every tide has a nearby spent seed, and \(j+8\le4B/5\)
eventually. A flag is incident to a coarse plaquette above \(t_c\); each
plaquette has at most four incident sites. Four-dimensional ball counts
give the remaining bounds.

\emph{Lean:}
\href{https://github.com/yang-mills-group/OpenAI-YM-20261028/blob/3456e5e75cf8d4cd558955a5c7638fba2d68f464/YM/Bid/Energy.lean\#L40-L41}{\texttt{Bid.E01}}
(\href{https://github.com/yang-mills-group/OpenAI-YM-20261028/blob/3456e5e75cf8d4cd558955a5c7638fba2d68f464/YM/Bid/Energy.lean\#L27-L32}{\texttt{Payments}},
\href{https://github.com/yang-mills-group/OpenAI-YM-20261028/blob/3456e5e75cf8d4cd558955a5c7638fba2d68f464/YM/Bid/Energy.lean\#L35-L38}{\texttt{MacroStore}}),
\href{https://github.com/yang-mills-group/OpenAI-YM-20261028/blob/3456e5e75cf8d4cd558955a5c7638fba2d68f464/YM/Work/E01.lean\#L289-L324}{\texttt{Work.E01.answer}}.

\textbf{Lemma 4.3 (shell comparison; E02).} For every gauge group \(G\),
block factor \(L\), index \(j\) and kernel choices \(K\), there is
\(C>0\) with the following property. Take a step using \(K\), block
factor \(L\) and index \(j\), whose fine side is at least \(100\) times
the macro radius, with \(c>1\); any \(U\) in the tree gauge; any \(V\);
and any fine link \(e\) with usable tags. Then
\[\chi(U_e\tilde V_e^{-1})\le C\,{\rm port}(\bar e,U)+w_{(\bar e,\mu)},\]
where \(\tilde V\) is the replica of \(V\) on fine links and \(\bar e\)
the parent block. The last term is present only when the ends of \(e\)
lie in different blocks.

\emph{Sketch:} close \(e\) by the two tree paths, lift the loop to
\(\mathbb Z^4\), and bound the chord of the product by the sum of chords
along a walk whose length is bounded by the macro budget. For a crossing
link, the predictor estimate controls the discrepancy between the two
block roots; the triangle inequality then adds the coarse mismatch
\(w_b\).

\emph{Lean:}
\href{https://github.com/yang-mills-group/OpenAI-YM-20261028/blob/3456e5e75cf8d4cd558955a5c7638fba2d68f464/YM/Bid/Energy.lean\#L49-L52}{\texttt{Bid.E02}},
\href{https://github.com/yang-mills-group/OpenAI-YM-20261028/blob/3456e5e75cf8d4cd558955a5c7638fba2d68f464/YM/Work/E02.lean\#L132-L244}{\texttt{Work.E02.answer}}.

\paragraph{4.3 Kernel bounds}\label{kernel-bounds}

\textbf{Lemma 4.4 (rows and columns; E03).} There is \(A>0\),
independent of \(L\) and \(W\), such that: - \(\Phi\) has row sums at
most \(A/L^2\) and column sums at most \(AL^2\); - \(\Phi'\) and \(N\)
have row and column sums at most \(A\); - all entries vanish beyond
distance \(B/50\).

Take absolute values of the cover coefficients. Adjoint transport has
norm one, periodization charges each cover term once per row or column,
and the fine residue multiplicity is \(L^4\)
(\href{https://github.com/yang-mills-group/OpenAI-YM-20261028/blob/3456e5e75cf8d4cd558955a5c7638fba2d68f464/YM/Work/E03.lean\#L39-L90}{\texttt{anchored\_of\_terms}}).

\emph{Lean:}
\href{https://github.com/yang-mills-group/OpenAI-YM-20261028/blob/3456e5e75cf8d4cd558955a5c7638fba2d68f464/YM/Bid/Energy.lean\#L67}{\texttt{Bid.E03}}
(\href{https://github.com/yang-mills-group/OpenAI-YM-20261028/blob/3456e5e75cf8d4cd558955a5c7638fba2d68f464/YM/Bid/Energy.lean\#L61-L64}{\texttt{Allowance}}),
\href{https://github.com/yang-mills-group/OpenAI-YM-20261028/blob/3456e5e75cf8d4cd558955a5c7638fba2d68f464/YM/Work/E03.lean\#L228-L256}{\texttt{Work.E03.answer}}.

\textbf{Lemma 4.5 (path series and predictions; E04, given E03).} There
is \(A>0\), independent of \(L\) and \(W\), such that, with \(s=t\) at
the fine stage and \(s=t_c\) at the coarse stage: - the fine and coarse
path series have absolute row and column sums at most \(A\) and range at
most \(B/50\); - \(|\ell_p(U)|\le\sum_q c_{pq}|U_q|\) for the
nonnegative absolute path coefficients \(c_{pq}\), and
\(|\ell_p|\le As\) on kept plaquettes; - \(|v|\le At\),
\(|\Phi v|\le At/L^2\) and \(|\Phi'v|\le At\).

\emph{Proof:} sum the absolute path coefficients with their exponential
length weights. Periodization gives the column count; the kept test
bounds each contributing plaquette by the threshold times the same
exponential weight. Thinning bounds the coarse coordinates, and the row
estimates of Lemma 4.4 bound their predictions. The comparison
\(t_c(x)\le2t\) holds eventually at fixed \(L\).

\emph{Lean:}
\href{https://github.com/yang-mills-group/OpenAI-YM-20261028/blob/3456e5e75cf8d4cd558955a5c7638fba2d68f464/YM/Bid/Energy.lean\#L89-L91}{\texttt{Bid.E04}}
(\href{https://github.com/yang-mills-group/OpenAI-YM-20261028/blob/3456e5e75cf8d4cd558955a5c7638fba2d68f464/YM/Bid/Energy.lean\#L74-L81}{\texttt{LaneCounts}},
\href{https://github.com/yang-mills-group/OpenAI-YM-20261028/blob/3456e5e75cf8d4cd558955a5c7638fba2d68f464/YM/Bid/Energy.lean\#L83-L86}{\texttt{Maximum}}),
\href{https://github.com/yang-mills-group/OpenAI-YM-20261028/blob/3456e5e75cf8d4cd558955a5c7638fba2d68f464/YM/Work/E04.lean\#L282-L307}{\texttt{Work.E04.answer}}.

\paragraph{4.4 Row-by-row lower bounds}\label{row-by-row-lower-bounds}

\textbf{Lemma 4.6 (coarse rows; E05).} There are \(\delta,B_1>0\),
independent of \(L\) and \(W\), such that for every coarse plaquette
with \(\tau=|V_q|\): - \(t_c<\tau\le\delta\) implies row
\(\ge\frac{\alpha_c}4t_c\tau\); - \(\tau\le t_c\) implies row
\(\ge-B_1t_c^4\); - \(\tau>\delta\) implies row
\(\ge-\lambda_ct_c\tau\).

\emph{Sketch:} in matrix space the curvature is the tangential
projection of \(V_q-1\), so \(\chi^2-|{\rm curv}|^2\) is the squared
normal displacement divided by the fixed matrix-to-colour normalization,
which is \(O(\chi^4)\) in a tubular neighbourhood of \(G\). - The path
component cancels whether kept or dropped. For \(\tau\le t_c\) the
direct component is unthinned, and the row equals
\(-\frac{\alpha_c}2(\chi^2-|{\rm curv}|^2)\). - For
\(t_c<\tau\le\delta\) the row is
\(\ge\alpha_ct_c\tau(1-O(\delta^2))-\frac{\alpha_c}2t_c^2-\lambda_ct_c\tau\ge\frac{\alpha_c}4t_c\tau\),
using \(\lambda_c<\alpha_c/16\).

\emph{Lean:}
\href{https://github.com/yang-mills-group/OpenAI-YM-20261028/blob/3456e5e75cf8d4cd558955a5c7638fba2d68f464/YM/Bid/Energy.lean\#L101}{\texttt{Bid.E05}}
(\href{https://github.com/yang-mills-group/OpenAI-YM-20261028/blob/3456e5e75cf8d4cd558955a5c7638fba2d68f464/YM/Bid/Energy.lean\#L94-L97}{\texttt{Crowds}}),
\href{https://github.com/yang-mills-group/OpenAI-YM-20261028/blob/3456e5e75cf8d4cd558955a5c7638fba2d68f464/YM/Work/E05.lean\#L211-L228}{\texttt{Work.E05.answer}}.

\textbf{Lemma 4.7 (fine and interior rows; E06, given E03, E04).} There
is \(L_1\) such that for every \(L\ge L_1\) and every \(W\ge2A/\sqrt a\)
(with \(A\) from E04), for large \(x\), with \(d=\alpha_F/128\):
\[\tfrac{\lambda_F}2\,t\!\!\sum_{p\in A_F,\,|U_p|>t}\!\!|U_p|+d\,t^2\,\sharp\{p\in A_F:|U_p|\le t\}\le\sum\text{fine rows},\qquad\tfrac a8\sum_{b\in A_I}w_b^2\le\sum\text{interior rows}.\]

\emph{Proof:} - \emph{Costs with \(|U_p|\le t\).} The direct part of the
fine row is at least \(\frac12(\sqrt{\alpha_F}|U_p|-|(\Phi v)_p|)^2\).
For a cost, \(|U_p|\ge t/4\), while
\(|(\Phi v)_p|\le At/L^2\le\sqrt{\alpha_F}t/8\) for \(L\ge L_1\); this
gives \(\alpha_Ft^2/128\). - \emph{Costs with \(|U_p|>t\).} The clipped
action is \(\lambda_Ft|U_p|\), and the cross term is at most
\(AL^{-2}\sqrt{\alpha_F}t|U_p|\le\frac{\lambda_F}4t|U_p|\). -
\emph{Path-series part.} It is a complete square on kept plaquettes. On
dropped plaquettes the cross term is paid by the column bounds of E04,
at total cost \(\frac{\lambda_F}4t\sum|U_p|\); this works only after
summation, not row by row. - \emph{Aisles.} On an aisle \(w_b\ge Wt\)
and \(|(\Phi'v)_b|\le At\le\frac{\sqrt a}2w_b\), so the row is at least
\(\frac12(\frac{\sqrt a}2w_b)^2=\frac a8w_b^2\). - \emph{Unlisted rows.}
Interior rows are nonnegative by the same square lower bound. Unlisted
fine direct parts and kept path parts are nonnegative; the dropped path
parts are included in the summed loss above.

\emph{Lean:}
\href{https://github.com/yang-mills-group/OpenAI-YM-20261028/blob/3456e5e75cf8d4cd558955a5c7638fba2d68f464/YM/Bid/Energy.lean\#L110}{\texttt{Bid.E06}}
(\href{https://github.com/yang-mills-group/OpenAI-YM-20261028/blob/3456e5e75cf8d4cd558955a5c7638fba2d68f464/YM/Bid/Energy.lean\#L103-L107}{\texttt{Seals}}),
\href{https://github.com/yang-mills-group/OpenAI-YM-20261028/blob/3456e5e75cf8d4cd558955a5c7638fba2d68f464/YM/Work/E06.lean\#L413-L440}{\texttt{Work.E06.answer}}.

\textbf{Lemma 4.8 (affected cells; E07, given E01, E03, E04).} For every
fixed \(\delta>0\) there is a pre-\(L\) constant \(C_7>0\), and for
every \(L\) and \(W>0\) a constant \(Z\), such that for large \(x\):
\[\text{defects}({\rm hit})+\frac{C_7}{L^2}t\!\!\sum_{p\in A_F,|U_p|>t}\!\!|U_p|+C_7t\sum_{b\in A_I}w_b+Zt^2B^4\sharp A^{\rm sp}\ \ge\ \lambda_ct_c\sum_{\rm tides}|V_q|.\]

\emph{Proof:} 1. By qdc, on the hit set \(|V_q|\le r_0t+H_q\). The macro
excess \(H\) is bounded, through the charges of E01, by a pre-\(L\)
constant times \(L^{-2}\sum_{\rm swell}|U_p|+\sum_{A_I}w_b\). 2. The
defect rows on hit cells are bounded below by Cauchy--Schwarz, using
\(|v|\le At\) and the kernel bounds of E03--E04. The inputs \(X\),
\(Z\), \(X^c\) split into a uniform part, \(O(t)\) per cell, and the
excess. 3. The uniform part is paid by
\(Zt^2\sharp{\rm hit}\le ZNt^2B^4\sharp A^{\rm sp}\); the excess is paid
by the two middle terms.

No cancellation of linear terms is used.

\emph{Lean:}
\href{https://github.com/yang-mills-group/OpenAI-YM-20261028/blob/3456e5e75cf8d4cd558955a5c7638fba2d68f464/YM/Bid/Energy.lean\#L130}{\texttt{Bid.E07}}
(\href{https://github.com/yang-mills-group/OpenAI-YM-20261028/blob/3456e5e75cf8d4cd558955a5c7638fba2d68f464/YM/Bid/Energy.lean\#L121-L128}{\texttt{Lapse}}),
\href{https://github.com/yang-mills-group/OpenAI-YM-20261028/blob/3456e5e75cf8d4cd558955a5c7638fba2d68f464/YM/Work/E07.lean\#L115-L158}{\texttt{Work.E07.answer}}.

\textbf{Lemma 4.9 (unaffected cells; E08, given E01, E02).} For every
\(L\) and \(W>0\) there is \(C\) such that for large \(x\) the following
holds. For every coarse plaquette \(q\) of \(D_{2B}(A)\) outside the hit
set there is a gauge transformation \(z\) with: -
\(\chi(V^z_b)\le h:=CB^{10}T\) for all coarse bonds within \(2B/5\) of
\(q\); - \(\chi(U^{z\circ\,\rm parent}_e)\le h\) for all fine bonds
whose parent is within \(2B/5\).

\emph{Sketch:} there is no spent seed within \(4B/5\) of \(q\), so all
tolls and wires near \(q\) are at most \(T\). An axial gauge on a box of
the coarse cover, with E02 transferring the bound to fine links, loses a
polynomial in \(B\).

\emph{Lean:}
\href{https://github.com/yang-mills-group/OpenAI-YM-20261028/blob/3456e5e75cf8d4cd558955a5c7638fba2d68f464/YM/Bid/Energy.lean\#L140}{\texttt{Bid.E08}}
(\href{https://github.com/yang-mills-group/OpenAI-YM-20261028/blob/3456e5e75cf8d4cd558955a5c7638fba2d68f464/YM/Bid/Energy.lean\#L134-L138}{\texttt{Stitches}}),
\href{https://github.com/yang-mills-group/OpenAI-YM-20261028/blob/3456e5e75cf8d4cd558955a5c7638fba2d68f464/YM/Work/E08.lean\#L342-L344}{\texttt{Work.E08.answer}}.

\paragraph{4.5 Linearization defects on unaffected
cells}\label{linearization-defects-on-unaffected-cells}

\textbf{Lemma 4.10 (flat identities).} Consider the linearized uncut
kernels \(\Phi_0,\Phi_0',N_0\) (scalar matrices built from the towers),
and the linearized inputs for log-fields \(u\) (fine) and \(v\)
(coarse): \(X_0(u)\) (face curls and path series of \(u\)) and
\(Z_0(u,v)=\sqrt{a_{\rm pen}}\,(v-{\rm probe}\,u)\), where
\(a_{\rm pen}\) is the penalty. Under the admissibility identities and
agreement of the beds with the fine and coarse path data, for every
\(u,v\) and every coarse two-layer vector \(w\),
\[\Phi_0^*X_0(u)+\Phi_0'^*Z_0(u,v)-X_0^c(v)=0,\qquad w-(\Phi_0^*\Phi_0+\Phi_0'^*\Phi_0'+N_0^*N_0)w=0.\]
\emph{Proof:} the flat tower identities, applied componentwise to the
scalar matrices and the attached path data.

\emph{Lean:}
\href{https://github.com/yang-mills-group/OpenAI-YM-20261028/blob/3456e5e75cf8d4cd558955a5c7638fba2d68f464/YM/Next/FLUCTUATIONS/Linear.lean\#L89-L136}{\texttt{roll\_l}},
\href{https://github.com/yang-mills-group/OpenAI-YM-20261028/blob/3456e5e75cf8d4cd558955a5c7638fba2d68f464/YM/Next/FLUCTUATIONS/Linear.lean\#L69-L87}{\texttt{roll\_d}}.

\textbf{Lemma 4.11 (chart estimates on the torus).} For a patch \(D\)
satisfying the branch condition there are \(C,\tau,r>0\) with the
following property. For every admissible step on \(D\), all tags, all
\(R\), and all log-fields \(a\), \(b\) with \(\|a\|,\|b\|\le h\le r\),
set \(U=\exp a\) and \(V=\exp b\). Define the defect vectors
\(\mathfrak l=\Phi^*X+(\Phi')^*Z-X^c\) and
\(\mathfrak d(w)=w-(\Phi^*\Phi+(\Phi')^*\Phi'+N^*N)w\). At every coarse
plaquette, componentwise in its two coordinates, -
\(\|\mathfrak l\|\le Ch\,\omega\); -
\(\|\mathfrak d(w)\|\le Cq\,\omega\) whenever \(q\ge0\) and
\(\|w\|\le q\).

Here \(\omega=\omega(R,h,\tau)=(1+R)h+e^{-\tau R}\)
(\href{https://github.com/yang-mills-group/OpenAI-YM-20261028/blob/3456e5e75cf8d4cd558955a5c7638fba2d68f464/YM/Next/FLUCTUATIONS/Defs.lean\#L70}{\texttt{widen}}).

\emph{Proof:} compare every actual object with its linearization, then
subtract Lemma 4.10. - \emph{Kernels:} they are close to their flat
versions
(\href{https://github.com/yang-mills-group/OpenAI-YM-20261028/blob/3456e5e75cf8d4cd558955a5c7638fba2d68f464/YM/Next/FLUCTUATIONS/Kernel.lean\#L264-L298}{\texttt{kernel\_chart\_cost}}).
- \emph{Path-series coordinates}
(\href{https://github.com/yang-mills-group/OpenAI-YM-20261028/blob/3456e5e75cf8d4cd558955a5c7638fba2d68f464/YM/Next/FLUCTUATIONS/Slots.lean\#L83-L158}{\texttt{x\_chart\_cost}}):
- transporting a curvature of size \(O(h)\) along a path of length at
most \(R\) costs \(O(Rh\cdot h)\); - the paths longer than \(R\) are
paid by the weight \(e^{\sigma\,\rm length}\) in the path budget, which
gives \(e^{-(\sigma/d)R}\). - \emph{Predictor coordinates:} quadratic
error
(\href{https://github.com/yang-mills-group/OpenAI-YM-20261028/blob/3456e5e75cf8d4cd558955a5c7638fba2d68f464/YM/Next/FLUCTUATIONS/Prediction.lean\#L243-L267}{\texttt{z\_chart\_cost}}).

\emph{Lean:}
\href{https://github.com/yang-mills-group/OpenAI-YM-20261028/blob/3456e5e75cf8d4cd558955a5c7638fba2d68f464/YM/Next/FLUCTUATIONS/Global.lean\#L76-L148}{\texttt{global\_chart\_cost}}.

\textbf{Lemma 4.12 (interior cost).} For the same \(D\) there are
\(C,\tau,r>0\) with the following property. Suppose that in some gauge
\(z\) both \(U\) and \(V\) are within \(h\le r\) of the identity on all
bonds within distance \(3R+4+{\rm reach}\) of a coarse plaquette \(q\),
where \({\rm reach}\) is the predictor reach and \(h\ge0\). The step
must be on \(D\) and admissible for \(P\); the tags and cut length
\(R\in\mathbb N\) are arbitrary. Then the bounds of Lemma 4.11 hold at
\(q\) for every coarse coordinate vector bounded by a nonnegative
constant.

\emph{Proof:} 1. \emph{Gauge covariance:} the defect vectors transform
by \({\rm Ad}(z_q)\)
(\href{https://github.com/yang-mills-group/OpenAI-YM-20261028/blob/3456e5e75cf8d4cd558955a5c7638fba2d68f464/YM/Next/FLUCTUATIONS/Covariance.lean\#L151-L185}{\texttt{gauge\_gaps}}).
2. \emph{Locality:} the defect vectors depend only on nearby bonds
(\href{https://github.com/yang-mills-group/OpenAI-YM-20261028/blob/3456e5e75cf8d4cd558955a5c7638fba2d68f464/YM/Next/FLUCTUATIONS/Locality.lean\#L151-L174}{\texttt{local\_gaps}}).
3. \emph{Extension:} a configuration that is \(h\)-close to \(1\) on a
ball agrees there with \(\exp\) of a global log-field of norm \(O(h)\)
(\href{https://github.com/yang-mills-group/OpenAI-YM-20261028/blob/3456e5e75cf8d4cd558955a5c7638fba2d68f464/YM/Next/FLUCTUATIONS/Charts.lean\#L89-L106}{\texttt{extend\_interior}}).
4. Apply Lemma 4.11.

\emph{Lean:}
\href{https://github.com/yang-mills-group/OpenAI-YM-20261028/blob/3456e5e75cf8d4cd558955a5c7638fba2d68f464/YM/Next/FLUCTUATIONS/Interior.lean\#L22-L63}{\texttt{interior\_cost}}.

\textbf{Lemma 4.13 (unaffected defects; EBare).} Assume E03, E04 and the
conclusion of E08 at \((P,L,W)\). Then for every \(\varepsilon>0\) and
all large \(x\), the sum of the defect rows over the unhit coarse
plaquettes of \(D_{2B}(A)\) is at least \(-\varepsilon t^2\sharp A\).

\emph{Proof:} 1. Take \(h=CB^{10}T\) from E08 and
\(\omega=\omega(R,h,\tau)\). Lemma 4.12 applies, since
\(3R+4+{\rm reach}\le2B/5\). 2. Cauchy--Schwarz with \(|v|\le At\)
gives, on every unhit cell, defects at least \(-\ell\), where
\(\ell=2(At)(Ch\omega)+(At)(CAt\,\omega)\). 3. There are at most
\(NB^4\sharp A\) such cells (E01 with \(\delta=1\)). 4.
\(NB^4\ell\le\varepsilon t^2\) for large \(x\), because
\(B^n\omega\to0\) and \(B^nh\omega/t\to0\) for every \(n\)
(\href{https://github.com/yang-mills-group/OpenAI-YM-20261028/blob/3456e5e75cf8d4cd558955a5c7638fba2d68f464/YM/Next/FLUCTUATIONS/Scales.lean\#L116-L152}{\texttt{radius\_stitch\_errors}}).
Indeed \(h=\beta_F^{-0.49+o(1)}\), \(h^2/t=\beta_F^{-0.48+o(1)}\), and
\(e^{-\tau R}\le\beta_F^{-1}\) because \(R\asymp(\log x)^2\). 5. If no
patch with the request's data satisfies the branch condition, the
statement is vacuous.

\emph{Lean:}
\href{https://github.com/yang-mills-group/OpenAI-YM-20261028/blob/3456e5e75cf8d4cd558955a5c7638fba2d68f464/YM/Bid/Energy.lean\#L148}{\texttt{Bid.EBare}}
(\href{https://github.com/yang-mills-group/OpenAI-YM-20261028/blob/3456e5e75cf8d4cd558955a5c7638fba2d68f464/YM/Bid/Energy.lean\#L142-L144}{\texttt{Belly}}),
\href{https://github.com/yang-mills-group/OpenAI-YM-20261028/blob/3456e5e75cf8d4cd558955a5c7638fba2d68f464/YM/Next/FLUCTUATIONS/Bare.lean\#L64-L120}{\texttt{Fluctuations.bare}}.

\paragraph{4.6 Assembly of III.10.1}\label{assembly-of-iii.10.1}

\textbf{Lemma 4.14 (assembly; E09).} For a fixed request, the packing
estimate of Lemma 4.2, fine/interior bounds of Lemma 4.7, coarse bounds
of Lemma 4.6 at some \(\delta>0\), affected-cell bound of Lemma 4.8 at
that \(\delta\), and unaffected-cell bound of Lemma 4.13 for every
\(L,W>0\), imply III.10.1.

\emph{Choice of constants:} - \(L\ge\max(L_1,L_7)\), where
\(C_7/L_7^2\le\lambda_F/4\); - \(W\ge\max(1,\ 2A/\sqrt a,\ 16C_7/a)\); -
\(m=\min(\lambda_F/8,\ \alpha_F/128,\ a/32,\ \alpha_c/64)\) and
\(\kappa=m/2\); since \(\lambda_F<\alpha_F/16\), this is the value in
III.10.1; - \(x\) so large that: \(t_c\le\delta\); \(t\le T\);
\(t\le2t_c\); \((4mN+ZB^4)t^2\le m\,tT\);
\(B_1t_c^4NB^4\le\frac m4t^2\); and the unaffected-cell bound holds with
\(\varepsilon=m/4\).

\emph{Proof:} add the rows. 1. \emph{Swell costs:}
\(\frac{\lambda_F}2t|U_p|\). Half of it pays the \(C_7L^{-2}\) term of
Lemma 4.8; the remainder is at least \(2mt^2\), and at least \(m\,tT\)
when the cost is spent. 2. \emph{Other costs:} at least
\(\alpha_Ft^2/128\ge mt^2\). 3. \emph{Aisles:}
\(\frac a8w^2-C_7tw\ge\frac a{16}w^2\ge mt^2\); spent aisles carry
\(m\,tT\). 4. \emph{Tides:} paid by Lemma 4.8. 5. \emph{Coarse
plaquettes in the intermediate range:}
\(\frac{\alpha_c}4t_c|V_q|\ge\frac{\alpha_c}{16}t^2\), which pays
\(4mt^2\) per intermediate plaquette, since
\(\sharp A_\flat\le4(\sharp\text{intermediate}+N\sharp A^{\rm sp})\). 6.
\emph{Small coarse plaquettes:} in total at least
\(-\frac m4t^2\sharp A\). 7. \emph{Gauge rows:} \(\ge0\). 8.
\emph{Defects:} on hit cells they are inside Lemma 4.8; on unhit cells
they total at least \(-\frac m4t^2\sharp A\). 9. \emph{Spent seeds:}
their \(m\,tT\) covers the flag charge against \(A^{\rm sp}\) and the
\(Z\) term.

The total is at least \(\frac m2t^2\sharp A\).

\emph{Lean:}
\href{https://github.com/yang-mills-group/OpenAI-YM-20261028/blob/3456e5e75cf8d4cd558955a5c7638fba2d68f464/YM/Bid/Energy.lean\#L151-L152}{\texttt{Bid.E09}},
\href{https://github.com/yang-mills-group/OpenAI-YM-20261028/blob/3456e5e75cf8d4cd558955a5c7638fba2d68f464/YM/Work/E09.lean\#L113-L159}{\texttt{Work.E09.answer}}
(with
\href{https://github.com/yang-mills-group/OpenAI-YM-20261028/blob/3456e5e75cf8d4cd558955a5c7638fba2d68f464/YM/Work/E09.lean\#L69-L111}{\texttt{assemble}}
and
\href{https://github.com/yang-mills-group/OpenAI-YM-20261028/blob/3456e5e75cf8d4cd558955a5c7638fba2d68f464/YM/Work/E09.lean\#L23-L62}{\texttt{bank}}),
\href{https://github.com/yang-mills-group/OpenAI-YM-20261028/blob/3456e5e75cf8d4cd558955a5c7638fba2d68f464/YM/Bid/Energy.lean\#L155-L158}{\texttt{Bid.issues}}.

\paragraph{4.7 Activities of one step}\label{activities-of-one-step}

The next estimates use the hypotheses of III.10.2. In statements about
packets, logarithmic letters or warnings, the tube satisfies the
compatibility condition, the input satisfies the stated fine-level
bounds, and sectors have the actual coarse flags. Auxiliary window
constants are fixed before the threshold in \(x\).

\textbf{Lemma 4.15 (rate calculus).} A function \(w\) has rate \(a\) if
for every \(v>a\), \(|w(x)|\le\beta_F^{\,v}\) eventually. Products of
functions with rates \(a,b\) have rate \(a+b\); their sums have rate
\(\max(a,b)\). Eventual absolute domination preserves a rate, and
constants have rate \(0\). Thus \(t\) has rate \(-\frac12\), and \(B\)
and \({\rm bill}\) have rate \(0\).

\emph{Lean:}
\href{https://github.com/yang-mills-group/OpenAI-YM-20261028/blob/3456e5e75cf8d4cd558955a5c7638fba2d68f464/YM/AnalyticNext/CONNECTED_ACTIVITY_JETS_E6/Numerics.lean\#L22-L23}{\texttt{Rate}}.

\textbf{Lemma 4.16 (from the action gap to packet weights).} - (a) There
is \(L_0\) such that for all \(L\ge L_0\) and all large \(W\) there is
\(\mu>0\) with two properties. First, the bound of III.10.1 holds with
constant \(\mu\) at all large \(x\)
(\href{https://github.com/yang-mills-group/OpenAI-YM-20261028/blob/3456e5e75cf8d4cd558955a5c7638fba2d68f464/YM/Push/ABSOLUTE_SUMS/Grants.lean\#L52-L71}{\texttt{Grant.order}};
here \(\mu\) is the \(\kappa\) of III.10.1). Second, fix any halo and
nonnegative support exponent, \(d\in\mathbb N\), \(D\ge0\), \(z<2\) and
\(u>0\). For all large \(x\) and every \(\beta\ge u\beta_F\), the gain
\(\beta\mu t^2\) passes the packet summability test for any spend at
most \(D\,\mathrm{bill}^dM^z\) and any nonnegative prefactor at most its
exponential. The total packet charge is exponentially small in the
retained gain. The threshold may depend on these additional parameters.
- (b) Suppose the local sector and tag tests hold through radius
\(300B\), and the physical fluctuation on hot bonds is bounded by a
nonnegative function of rate \(-1/2\) through radius \(304B\). On the
annulus \(2B<\mathrm{dist}\le200B\), every remainder row is at least
\(-E(x)\) for some nonnegative \(E\) of rate \(-3/2\)
(\href{https://github.com/yang-mills-group/OpenAI-YM-20261028/blob/3456e5e75cf8d4cd558955a5c7638fba2d68f464/YM/AnalyticNext/CONNECTED_ACTIVITY_JETS_E6/Annulus.lean\#L164-L185}{\texttt{Coat.annulus}}).
After multiplying \(E\) by the annulus volume and fixed coordinate
multiplicity, it still has rate \(-3/2\). Consequently the assigned
energy is at least \((\mu t^2-E)\sharp A\)
(\href{https://github.com/yang-mills-group/OpenAI-YM-20261028/blob/3456e5e75cf8d4cd558955a5c7638fba2d68f464/YM/AnalyticNext/CONNECTED_ACTIVITY_JETS_E6/AssignedPrice.lean\#L54-L86}{\texttt{Coat.earns}}).
Since \(\beta E\) has rate at most \(-1/2\), the loss is absorbed in
\({\rm bill}^5\). - (c) In every manifest sector,
\(|{\rm Pack}|\le\exp\big((C\,{\rm bill}^5-\beta\mu t^2)\sharp A\big)\)
(\href{https://github.com/yang-mills-group/OpenAI-YM-20261028/blob/3456e5e75cf8d4cd558955a5c7638fba2d68f464/YM/AnalyticNext/CONNECTED_ACTIVITY_JETS_E6/Revenue.lean\#L99-L140}{\texttt{steam\_prices}}).

\emph{((a) and (c) proof read; (b) statement only)}

\textbf{Lemma 4.17 (old and logarithmic letters).} Fix \(-1/2<s<0\) and
any \(u\in\mathbb R\). Assume the sector drift and residual estimates,
with a residual envelope of rate \(-1\), as well as the common
hypotheses of this subsection. For sufficiently large \(L\) and then all
large \(x\), the old forced letters have a rooted bound with total
\(\beta_F^u\), and all logarithmic letters have one with total
\(\beta_F^s\). The latter include inherited regular activities, local
Jacobians and Gaussian determinant corrections. The weights bound both
\(|\ell_i|\) and \(|e^{\ell_i}-1|\). Taking \(u=s\) gives the two totals
used in III.10.2.

\emph{Proof:} transport the incoming forced bounds to coarse supports,
and split the logarithmic alphabet into inherited activities, Jacobians
and Gaussian corrections. Their envelopes have rate \(-1/2\) or smaller.
Choose an intermediate rate strictly between \(-1/2\) and \(s\) to
absorb their finite sum, and use \(|e^z-1|\le2|z|\) once the individual
logarithmic weights are at most one.

\emph{Lean:}
\href{https://github.com/yang-mills-group/OpenAI-YM-20261028/blob/3456e5e75cf8d4cd558955a5c7638fba2d68f464/YM/AnalyticNext/CONNECTED_ACTIVITY_JETS_E6/Quotation.lean\#L47-L92}{\texttt{recommends}}.

\textbf{Lemma 4.18 (window, drift, no new seeds).} For each patch there
are window constants such that, for large \(x\): - every manifest sector
satisfies the Hessian conditions of III.10.2(a)(iii)
(\href{https://github.com/yang-mills-group/OpenAI-YM-20261028/blob/3456e5e75cf8d4cd558955a5c7638fba2d68f464/YM/AnalyticNext/CONNECTED_ACTIVITY_JETS_E6/Shutter.lean\#L119-L144}{\texttt{shoots}});
- the Gaussian centring drift is \(O(t)\) and the residual of its linear
stationarity equation has an envelope of any prescribed rate, in
particular \(-1\)
(\href{https://github.com/yang-mills-group/OpenAI-YM-20261028/blob/3456e5e75cf8d4cd558955a5c7638fba2d68f464/YM/AnalyticNext/CONNECTED_ACTIVITY_JETS_E6/StampQ.lean\#L50-L59}{\texttt{YM.ConnectedActivityE6.Coat.tends\_val}});
- there is \(\xi>0\) with \(\sharp\{\text{colours}\}\cdot\xi\le1\) such
that fluctuations satisfying the radius-\(B\) window test create no
seeds at sites farther than \(30B\) from the given ones
(\href{https://github.com/yang-mills-group/OpenAI-YM-20261028/blob/3456e5e75cf8d4cd558955a5c7638fba2d68f464/YM/AnalyticNext/CONNECTED_ACTIVITY_JETS_E6/Peak.lean\#L92-L188}{\texttt{Coat.ready}}).

\emph{Proof:} exponential row bounds and coercivity give the precision
estimate; the local-inverse error tends to zero and yields a positive
covariance. Applying the approximate inverse to the linear term gives
the centring drift and a small linear residual. For the no-new-seed
assertion, choose \(\xi\) so the Gaussian fluctuation contributes at
most one third of the available margin. Linear cancellation from the
compatibility identities and the small nonlinear remainder pay the other
two thirds. The fine and mismatch prediction bounds then give the strict
thresholds \(t/4\) and \(Wt\).

\textbf{Lemma 4.19 (catastrophes).} For every set \(v\) of warning
addresses, the Gaussian measure of the catastrophe set is at most
\(E^{\sharp v}\), with \(E\) eventually smaller than every power of
\(\beta_F\). The threshold \(\xi t\sqrt\beta\) equals
\(\xi e^{(\log\log\beta_F)^2}(\beta/\beta_F)^{1/2}\) in Gaussian units.

\emph{Proof:} the covariance has a volume-independent Schur bound. Apply
the Gaussian rectangle estimate to simultaneous coordinate excursions
and sum the possible witnesses in each radius-\(2B\) neighbourhood. The
resulting envelope can be taken as \(C B^4 e^{-c M^{1/2}}\), where
\(M=e^{(\log\log\beta_F)^2}\) and \(c>0\). It is smaller than every
inverse power of \(\beta_F\). Repeated warning locations are paid
through the fixed multiplicity bound.

\emph{Lean:}
\href{https://github.com/yang-mills-group/OpenAI-YM-20261028/blob/3456e5e75cf8d4cd558955a5c7638fba2d68f464/YM/AnalyticNext/CONNECTED_ACTIVITY_JETS_E6/ResponsePrice.lean\#L65-L112}{\texttt{Coat.alarm\_bills}}.

\textbf{Lemma 4.20 (assembly of III.10.2).} The preceding estimates
imply III.10.2 with the following parameter choices. 1. The floor on
\(L\) is the sum of the floors from Lemma 4.16(a), the no-new-seed
margin, and Lemma 4.17. 2. \(W\) is taken large for the energy estimate
and the no-new-seed margin; \(\mu\) comes from Lemma 4.16(a) and does
not depend on the patch. 3. For each patch: - \(K,\gamma\) come from the
shell comparison; - the window constants from the precision estimate; -
the drift constant \(C_d\) from the centring estimate, and
\(a_o=\max(C_d+1,\gamma K)\); - \(\xi\) from the no-new-seed estimate; -
\(C\) from the packet estimate; - \(s\) is chosen with \(-\frac12<s<r\).
4. All conditions holding for large \(x\) are merged into one threshold.
The three letter totals have rate \(s<r\), so their sum is eventually at
most \(\beta_F^{\,r}\).

\emph{Lean:}
\href{https://github.com/yang-mills-group/OpenAI-YM-20261028/blob/3456e5e75cf8d4cd558955a5c7638fba2d68f464/YM/AnalyticNext/CONNECTED_ACTIVITY_JETS_E6/Issue.lean\#L94-L185}{\texttt{Target.issue}}.

\paragraph{4.8 Derivatives}\label{derivatives}

\textbf{Lemma 4.21 (letter jets).} Under the hypotheses of III.10.3, the
letter-jet structure exists before integration. Choose \(-1/2<s<r\) and
set \(\mathrm{rad}=\beta_F^s\). The warning amplitude before probability
extraction is a patch constant. The jet scale is a sum
\(\rho(x)=B_0(x)+A_0(x)\) of two rate-zero envelopes, chosen
independently of the input representation and volume.

\emph{Proof:} combine the packet-jet estimate, the old and logarithmic
letter estimates, and the differentiated warning cutoffs, all along
compatible transported fluctuation coordinates. Choose the catastrophe
envelope of rate \(2s-1\); its eventual bound by \(\beta_F^{2s}\) gives
the probability factor \(\mathrm{rad}^2\). The three rooted letter
totals have rate \(s<r\).

\emph{Lean:}
\href{https://github.com/yang-mills-group/OpenAI-YM-20261028/blob/3456e5e75cf8d4cd558955a5c7638fba2d68f464/YM/AnalyticNext/CONNECTED_ACTIVITY_JETS_E6/Das.lean\#L49-L102}{\texttt{Utopia.whittle}}.

\textbf{Lemma 4.22 (molecule integrals).} The letter bounds of Lemma
4.21 imply the integrated majorants of III.10.3, for the compatible
connected molecules specified there.

\emph{Proof:} the Leibniz rule for
\href{https://github.com/yang-mills-group/OpenAI-YM-20261028/blob/3456e5e75cf8d4cd558955a5c7638fba2d68f464/YM/AnalyticNext/CONNECTED_ACTIVITY_JETS_E6/Brush.lean\#L44-L47}{\texttt{Brush}}
adds the letter jet scales. Their sum is bounded by twice the molecule's
shape, giving, for rate-zero transport envelopes \(U,J'\) and fixed
integers \(k,K\ge0\), a pointwise factor
\[(1+|U(x)|+2(\rho(x)+J'(x)))^i(1+\mathrm{shape})^i\,\mathrm{facescale}^{ki}.\]
The bound vanishes outside the catastrophe set for the molecule's
distinct warning addresses. If that set is \(\mathcal C\) and its
address set is \(v\), Cauchy--Schwarz gives
\[\int\mathbf1_{\mathcal C}\,\mathrm{facescale}^{ki}\,d\gamma\le\left(\int\mathrm{facescale}^{2ki}\,d\gamma\right)^{1/2}\gamma(\mathcal C)^{1/2}\le e(x)(1+\mathrm{shape})^K\mathrm{rad}^{\sharp v},\]
where \(e\) has rate zero. Integration over the probability measure on
the tree slice adds no factor. Multiplying by the pointwise coefficient
gives a rate-zero prefactor \(p\) and shape power \(l=K+i\). For a
marked source, remove its ordinary fare and insert its amplitude bound.

\emph{Lean:}
\href{https://github.com/yang-mills-group/OpenAI-YM-20261028/blob/3456e5e75cf8d4cd558955a5c7638fba2d68f464/YM/PhysicalNext/RG_NEXT_COMPLETE_REPEATABLE_NONLINEAR_STEP/Daisy/Charm.lean\#L58-L179}{\texttt{timing}}.

\paragraph{4.9 Start-up fusion and
markers}\label{start-up-fusion-and-markers}

\textbf{Lemma 4.23 (H01--H04).} Under the respective hypotheses in
III.10.6(H), the four fusion conclusions hold. - \emph{H01} is the Mayer
identity \(e^{\sum\ell_i}=\prod_i(1+(e^{\ell_i}-1))\), regrouped with
the forced activities into pellets. Absolute domination of the original
and logarithmic series justifies the product expansion and regrouping;
the argument also proves summability of the fused series. - \emph{H02:}
retain one factor \(e^{-\kappa}\) from a forced letter in every pellet
and use the rooted tree sum \(e^\gamma-1\). - \emph{H03:} transport
every pellet support, inventory and eligibility test under each lattice
symmetry. The local flag tests and the permutations of the original
addresses induce those of the fused family. - \emph{H04:} each pellet is
a fixed finite product; a common summable envelope then gives continuity
of the weighted series by dominated convergence.

\emph{(proof read for H01, H02 and H04; H03 forced-family proof read,
symmetry moves statement only)}

\emph{Lean:}
\href{https://github.com/yang-mills-group/OpenAI-YM-20261028/blob/3456e5e75cf8d4cd558955a5c7638fba2d68f464/YM/Work/H01.lean\#L245-L251}{\texttt{H01.gathered}},
\href{https://github.com/yang-mills-group/OpenAI-YM-20261028/blob/3456e5e75cf8d4cd558955a5c7638fba2d68f464/YM/Work/H02.lean\#L17-L80}{\texttt{H02.answer}},
\href{https://github.com/yang-mills-group/OpenAI-YM-20261028/blob/3456e5e75cf8d4cd558955a5c7638fba2d68f464/YM/Work/H03.lean\#L138-L209}{\texttt{H03.fused\_forcedFamily}},
\href{https://github.com/yang-mills-group/OpenAI-YM-20261028/blob/3456e5e75cf8d4cd558955a5c7638fba2d68f464/YM/Work/H04.lean\#L155-L161}{\texttt{H04.answer}}.

\textbf{Lemma 4.24 (I01--I04).} Under the respective hypotheses in
III.10.6(I), the four marker conclusions hold. - \emph{I01:} traverse
the centres of a minimizing connected \(B\)-ball cover. Distinct markers
force gaps of at least \(B^2-2B\) between the traversal times of their
covering centres. Comparing total traversal length with the number of
balls gives the factor \(64\). - \emph{I02:} I01 converts extra support
decay to an increase of multiple-marker radii. Splitting single-marker
and multiple-marker sums gives total two. Shrinking the radius by the
order factorial converts factorially weighted derivative bounds to
ordinary derivative bounds. - \emph{I03:} centred embeddings compose and
preserve the parent relation. The open-stop inequality persists to later
levels because \(B_{r+1}^2\le2B_r^2\) and each block factor is at least
three. - \emph{I04} follows from
\href{https://github.com/yang-mills-group/OpenAI-YM-20261028/blob/3456e5e75cf8d4cd558955a5c7638fba2d68f464/YM/Work/I04.lean\#L24-L62}{\texttt{columns\_receipt}}:
for every level schedule,
\(\sum_{m<i}\log{\rm columns}_{m+1}=o({\rm scale}_i)\). Combine this
with \(\log{\rm order}_i!\le\log{\rm columns}_i\) and
\(\log\hat\beta_i=o({\rm scale}_i)\).

\emph{Lean:}
\href{https://github.com/yang-mills-group/OpenAI-YM-20261028/blob/3456e5e75cf8d4cd558955a5c7638fba2d68f464/YM/Work/I01.lean\#L19-L82}{\texttt{I01.answer}},
\href{https://github.com/yang-mills-group/OpenAI-YM-20261028/blob/3456e5e75cf8d4cd558955a5c7638fba2d68f464/YM/Work/I02.lean\#L159-L180}{\texttt{I02.answer}},
\href{https://github.com/yang-mills-group/OpenAI-YM-20261028/blob/3456e5e75cf8d4cd558955a5c7638fba2d68f464/YM/Work/I03.lean\#L43-L63}{\texttt{I03.answer}},
\href{https://github.com/yang-mills-group/OpenAI-YM-20261028/blob/3456e5e75cf8d4cd558955a5c7638fba2d68f464/YM/Work/I04.lean\#L76-L129}{\texttt{I04.answer}}.

\paragraph{4.10 The schedule}\label{the-schedule}

\textbf{Lemma 4.25 (pricing order).} Under the hypotheses of
III.10.5(a), the following choices produce its floors and raised entry
couplings. 1. For each pair of adjacent stages, III.10.1 gives a floor
\(b_{ij}\); set \(k_i=\max(b_{ii},b_{i,i-1},\text{sheet floor}_i)\). 2.
For \(L_i\ge k_i\), choose \(W\ge m\) in the admissible range and a
threshold \(x_{ij}\). 3. Raise \(\beta^*_i\) to
\(\max(\beta^*_i,x_{ii},x_{i,i-1})\). 4. The conclusion persists under
further raising
(\href{https://github.com/yang-mills-group/OpenAI-YM-20261028/blob/3456e5e75cf8d4cd558955a5c7638fba2d68f464/YM/Bid/Energy.lean\#L204-L208}{\texttt{recheck}}).

\emph{Lean:}
\href{https://github.com/yang-mills-group/OpenAI-YM-20261028/blob/3456e5e75cf8d4cd558955a5c7638fba2d68f464/YM/Bid/Energy.lean\#L173-L202}{\texttt{Bid.pricing\_order}}.

\textbf{Lemma 4.26 (divisible schedule).} Under the hypotheses of
III.10.5(b), the following construction gives a schedule satisfying all
entry, plateau and divisibility requirements. 1. Entry couplings come
from the bar (using \(G\) simple), leads from the wait data, and a
monotone stage-index function \(J\) with steps at most \(1\) and
\(J\to\infty\) by the finite-request scheduling lemma of III.2. 2. The
scales are \(s_i\); the initial length \(M\) is the minimal length plus
the largest bill over the first stages; the periods are \(4(M+p+1)\).

The checks: - \(\hat\beta_m\ge\) the bar value, and
\(\hat\beta_m\ge\beta^*_{J(m+1)}\); - the scale entries; - the volume
inequalities: \({\rm shells}\cdot(1+B_i+B_{i+1})^{\rm growth}\le\) side,
\(100(B_i+1)\le\) side, and \(100(\text{patch radius}+1)\le\) next side;
- the mesh condition \(\log L_{J(i+1)}/s_i\to0\); - divisibility:
\(p=c(M+z+1)-(M+1)\) gives the period \(4c(M+z+1)\).

\emph{Lean:}
\href{https://github.com/yang-mills-group/OpenAI-YM-20261028/blob/3456e5e75cf8d4cd558955a5c7638fba2d68f464/YM/PhysicalNext/RG_NEXT_COMPLETE_REPEATABLE_NONLINEAR_STEP/Pane/Divisors.lean\#L26-L135}{\texttt{Wait.engine}}.

\subsubsection{5. Where the difficulty
lies}\label{where-the-difficulty-lies-12}

The energy estimate must preserve \(L\)-independent coefficients on the
large fine inputs and \(W\)-independent coefficients on the mismatches.
The fine prediction is \(O(t/L^2)\), so increasing \(L\) absorbs its
cross terms. The interior prediction is \(O(t)\), so increasing \(W\)
absorbs the mismatch cross terms. Losses depending on both \(L,W\) are
charged to spent seeds and absorbed only by increasing \(x\).

On unaffected cells, the transported, truncated kernels have no exact
completeness identity. The flat identity cancels the linear part, while
local gauge fixing, chart estimates and exponential tails control the
remaining defects. This is the unaffected-defect estimate
\href{https://github.com/yang-mills-group/OpenAI-YM-20261028/blob/3456e5e75cf8d4cd558955a5c7638fba2d68f464/YM/Bid/Energy.lean\#L148}{\texttt{Bid.EBare}}.
The constant \(\kappa\) is independent of \(L,W\); the coupling
threshold is not.

The activity and derivative estimates additionally require compatible
quadratic forms, controlled fluctuation transport and probability bounds
for simultaneous warnings. Their product estimates retain one gain for
each distinct warning address before the connected-molecule sum. The
schedule is numerical; start-up fusion and marker estimates serve
separate parts of the construction.

\subsubsection{6. Notation}\label{notation-12}

{\def\LTcaptype{none} % do not increment counter
\begin{longtable}[]{@{}
  >{\raggedright\arraybackslash}p{(\linewidth - 4\tabcolsep) * \real{0.3333}}
  >{\raggedright\arraybackslash}p{(\linewidth - 4\tabcolsep) * \real{0.3333}}
  >{\raggedright\arraybackslash}p{(\linewidth - 4\tabcolsep) * \real{0.3333}}@{}}
\toprule\noalign{}
\begin{minipage}[b]{\linewidth}\raggedright
Written proof
\end{minipage} & \begin{minipage}[b]{\linewidth}\raggedright
Lean (glossary)
\end{minipage} & \begin{minipage}[b]{\linewidth}\raggedright
Meaning
\end{minipage} \\
\midrule\noalign{}
\endhead
\bottomrule\noalign{}
\endlastfoot
\(x=\hat\beta_m\), \(\beta_F=x+b_G\log L\) &
\href{https://github.com/yang-mills-group/OpenAI-YM-20261028/blob/3456e5e75cf8d4cd558955a5c7638fba2d68f464/YM/Pressure/Seeds.lean\#L23}{\texttt{betaF}}
& coarse and fine reference couplings \\
\(t_c(x)\), \(t=t_c(\beta_F)\) &
\href{https://github.com/yang-mills-group/OpenAI-YM-20261028/blob/3456e5e75cf8d4cd558955a5c7638fba2d68f464/YM/Pressure/Seeds.lean\#L25}{\texttt{temp}},
\href{https://github.com/yang-mills-group/OpenAI-YM-20261028/blob/3456e5e75cf8d4cd558955a5c7638fba2d68f464/YM/Pressure/Seeds.lean\#L26}{\texttt{fiery}}
& small-field thresholds; \texttt{fiery} is at the \textbf{fine}
coupling \\
\(T=\beta_F^{-49/100}\) &
\href{https://github.com/yang-mills-group/OpenAI-YM-20261028/blob/3456e5e75cf8d4cd558955a5c7638fba2d68f464/YM/Pressure/Seeds.lean\#L27}{\texttt{largeT}}
& spent-seed threshold \\
\(B\), \(R\) &
\href{https://github.com/yang-mills-group/OpenAI-YM-20261028/blob/3456e5e75cf8d4cd558955a5c7638fba2d68f464/YM/Pressure/Seeds.lean\#L22}{\texttt{radiusB}},
\href{https://github.com/yang-mills-group/OpenAI-YM-20261028/blob/3456e5e75cf8d4cd558955a5c7638fba2d68f464/YM/Pressure/Seeds.lean\#L77}{\texttt{till}}
& cluster radius \(\lceil(\log x)^2\rceil\), path cutoff
\(\lfloor B/10^5\rfloor\) \\
\(\lvert U_p\rvert\), \(w_b\) &
\href{https://github.com/yang-mills-group/OpenAI-YM-20261028/blob/3456e5e75cf8d4cd558955a5c7638fba2d68f464/YM/Pressure/Rows.lean\#L55}{\texttt{toll}},
\href{https://github.com/yang-mills-group/OpenAI-YM-20261028/blob/3456e5e75cf8d4cd558955a5c7638fba2d68f464/YM/Pressure/Cost.lean\#L187}{\texttt{wire}}
& plaquette size; mismatch \(\chi(V_bQ_b^{-1})\) \\
request \(P\); admissible & \texttt{Teller}, \texttt{Teller.Funds} &
data of one call of the step estimator \\
\(A=(A_F,A_I,A_\flat)\), consistent & \texttt{Stores},
\texttt{Stores.Plot} & seeds; two-sided thresholds on \(D_{5B}\) \\
\(A^{\rm sp}\), tides, hit & \texttt{spend}, \texttt{tides},
\texttt{hit} & seeds above \(T\); coarse plaquettes above \(\delta\);
cells near \(A^{\rm sp}\) \\
\(\mathcal E_A\) & \texttt{Stores.assessed} & localized action gap \\
\(v\), \(\Phi,\Phi',N\) & \texttt{vision}, \texttt{Φ},
\texttt{Φ\textquotesingle{}}, \texttt{N} & coarse coordinates;
transported kernels cut at \(R\) \\
\(\Phi_0,\Phi_0',N_0\) & \texttt{flatF}, \texttt{flatI}, \texttt{flatN}
& flat, uncut scalar tower operators \\
\(\omega(R,h,\tau)\) & \texttt{widen} & \((1+R)h+e^{-\tau R}\) \\
rate \(a\) & \texttt{Rate} & \(\le\beta_F^{a+o(1)}\) \\
target data \(\mathcal T\) & \texttt{Target} & \(\varepsilon,\delta,e\),
rate \(r\), halo, chart \\
letter weights, KP condition & \texttt{Sales}, \texttt{Buy} & rooted
Kotecký--Preiss and tree-graph tests \\
new-gas bounds & \texttt{Cash}, \texttt{Utopia}, \texttt{Herd},
\texttt{Captain} & III.10.2 bundle; letter jets; molecule integrals \\
\(N_B(H)\) & \texttt{coverNumber} & minimum connected ball-cover size \\
\(\rho(x)\), \(\mathrm{rad}\) & \texttt{Utopia.active},
\texttt{Utopia.rad} & jet scale of rate zero; square root of a
warning-probability bound \\
Gaussian sector & \texttt{Value}, \texttt{Charges} & precision,
approximate inverse, mean, residual \\
\end{longtable}
}

\subsubsection{7. What remains
unexplained}\label{what-remains-unexplained-12}

\begin{enumerate}
\def\labelenumi{\arabic{enumi}.}
\tightlist
\item
  \textbf{Conversion to the norms of III.9.} The rooted optional-letter
  bound and the packet/tree test are explicit here. Their complete
  conversion to the outgoing \(E\)- and \(K\)-norms of II.2 and III.5,
  including the reset, remains to be written out. The relevant per-step
  assemblies are the activity-bound bundle, its jet version, and the
  molecule-integral bounds.
\item
  \textbf{Free-kernel compatibility along the trajectory.} The condition
  requires fine and coarse beds satisfying the diagonal and flat tower
  identities, invariance under lattice symmetries, attachment to the
  chosen path data, and positive invertibility bounds.
  \href{https://github.com/yang-mills-group/OpenAI-YM-20261028/blob/3456e5e75cf8d4cd558955a5c7638fba2d68f464/YM/PhysicalNext/RG_NEXT_COMPLETE_REPEATABLE_NONLINEAR_STEP/Pane/Rooms.lean\#L39-L41}{\texttt{YM.PhysicalNext.Repeat.Pane.ovation}}
  constructs it on the auxiliary torus from the trajectory's beds,
  assuming their lattice-symmetry property at every period. The proof of
  that symmetry property for the final selected trajectory, and the
  precise III.3 attribution of all tower inputs, have not been checked
  here.
\item
  \textbf{Depth of the analytic proof checks.} The assembly proofs of
  III.10.2--III.10.4 and their defining estimates have been read. The
  underlying stationary coercivity, local-inverse, packet-jet and
  individual letter-jet estimates have not all been read through; their
  detailed analytic proofs remain to be checked. The proof of H03's
  forced-family structure was read; its full lift of lattice-symmetry
  permutations was not checked through. These limitations do not alter
  the conditional statements above.
\item
  \textbf{Additional reset estimates.}
  \href{https://github.com/yang-mills-group/OpenAI-YM-20261028/blob/3456e5e75cf8d4cd558955a5c7638fba2d68f464/YM/PhysicalNext/RG_NEXT_COMPLETE_REPEATABLE_NONLINEAR_STEP/Alabaster/Molt.lean\#L25-L77}{\texttt{meritBar}}
  and
  \href{https://github.com/yang-mills-group/OpenAI-YM-20261028/blob/3456e5e75cf8d4cd558955a5c7638fba2d68f464/YM/PhysicalNext/RG_NEXT_COMPLETE_REPEATABLE_NONLINEAR_STEP/Alabaster/Basin.lean\#L23-L49}{\texttt{rimfire}}
  give summable jet bounds proportional to the input difference, under
  matched-input and derivative-control hypotheses.
  \href{https://github.com/yang-mills-group/OpenAI-YM-20261028/blob/3456e5e75cf8d4cd558955a5c7638fba2d68f464/YM/PhysicalNext/RG_NEXT_COMPLETE_REPEATABLE_NONLINEAR_STEP/Shellac/Sables.lean\#L26-L65}{\texttt{erudite}}
  gives positivity of a residual density and a \(\beta_F^{-23/20}\) jet
  bound on the off-diagonal remainder under matching, finite-bound and
  small-chart hypotheses. Their complete hypotheses and role in III.9's
  reset are not included in this unit's numbered outputs.
\item
  \textbf{Gaussian trigger scales.} In
  \href{https://github.com/yang-mills-group/OpenAI-YM-20261028/blob/3456e5e75cf8d4cd558955a5c7638fba2d68f464/YM/PhysicalNext/V9_JOINT_ENDPOINT_LOAD/Supply.lean\#L72-L98}{\texttt{JointLoad.stage}},
  the scales are \(b(x)=1+B(x)\) and \(p(x)=e^{(\log\log\beta_F)^2}\);
  the trigger estimate absorbs fixed powers of \(b\) into \(e^{-c p^2}\)
  per distinct endpoint. Its full use in the preconstruction has not
  been traced here.
\item
  \textbf{Comparison with Balaban.} The cited papers describe the
  corresponding small-field, cluster and large-field operations. A
  lemma-by-lemma identification, including threshold and normalization
  choices, remains open. The bibliographic dates and volumes above do
  not establish such an identification.
\end{enumerate}

\subsection{III.11. Comparison with the free Gaussian measure:
interacting rank and determinant
budgets}\label{iii.11.-comparison-with-the-free-gaussian-measure-interacting-rank-and-determinant-budgets}

\textbf{Part:} III, Ultraviolet construction (renormalization group).
\textbf{Lean:} chapter 24:
\href{https://github.com/yang-mills-group/OpenAI-YM-20261028/tree/3456e5e75cf8d4cd558955a5c7638fba2d68f464/YM/Next/RANK}{\texttt{YM/Next/RANK/}},
\href{https://github.com/yang-mills-group/OpenAI-YM-20261028/tree/3456e5e75cf8d4cd558955a5c7638fba2d68f464/YM/PhysicalNext/AFTERNOON_INTERACTING_RANK_COMPARISON}{\texttt{YM/PhysicalNext/AFTERNOON\_INTERACTING\_RANK\_COMPARISON/}},
\href{https://github.com/yang-mills-group/OpenAI-YM-20261028/tree/3456e5e75cf8d4cd558955a5c7638fba2d68f464/YM/PhysicalNext/RANK_COVARIANCE_ATLAS}{\texttt{YM/PhysicalNext/RANK\_COVARIANCE\_ATLAS/}},
\href{https://github.com/yang-mills-group/OpenAI-YM-20261028/tree/3456e5e75cf8d4cd558955a5c7638fba2d68f464/YM/MergeProof}{\texttt{YM/MergeProof/}}
(with \texttt{West/}, \texttt{North/}, \texttt{East/}, \texttt{South/}).
\textbf{Size:} 42,556 lines (about 5,300 declarations).

\textbf{Scope.} The unit compares the interacting renormalization-group
(RG) history with the free Maxwell field. It also contains the
stress-flux argument and the label algebra of the Maxwell equation
ideal.

\subsubsection{1. What this unit proves}\label{what-this-unit-proves-13}

The unit has two outputs. The first is the \textbf{rank comparison}. In
the continuum limit, the space \(F_d\) of local fields of degree \(d\)
is isomorphic to the corresponding space of the free Maxwell field,
compatibly with complex conjugation and with the lattice symmetry group.
Here \(F_d\) is the space of labels of degree \(d\) modulo labels whose
correlations vanish. The second output is the \textbf{stress tensor of
the limit}: a nonzero conserved stress tensor that satisfies the
translation and rotation Ward identities and is normalized against the
free Maxwell stress tensor along a scaling sequence. Both outputs rest
on quantitative estimates proved here: a per-site determinant budget,
chessboard bounds for connected parts of corrected observables, and
comparisons of covariances with the free curvature kernel. The trial
systems test a finite separating family of free one- and two-point
correlations.

Notation. Set \(\mu_k=\exp((\log\log\hat\beta_k)^2)\). The centered
estimate and the trial systems use this fixed schedule quantity. \(G\)
is a compact matrix group, assumed simple only where stated; \(\dim G\)
is the dimension of its Lie algebra; \(h\) is normalized Haar measure;
\(\mathrm{chord}(u)\ge0\) is the chordal distance from \(u\) to \(1\);
and \(\mathrm{shell}(r)=\{u\in G:\mathrm{chord}(u)\le r\}\). The torus
of side \(n\ge2\) is \(\Lambda_n=(\mathbb Z/n)^4\), with \(c=n^4\)
sites, \(4c\) bonds and \(6c\) plaquettes. The prior \(\mu_{\rm prior}\)
is product Haar measure on bonds. An RG representation \(R\) on
\(\Lambda_n\) has an inverse coupling \(\beta=\beta(R)\), an exponent
\(\mathrm{exponent}_R\) with constant part \(\mathrm{scalar}(R)\), and
an unnormalized density
\(e^{R}=e^{\mathrm{exponent}_R}\cdot(\text{polymer gas})\).

\textbf{Proposition III.11.1 (determinant budget on the torus, and its
per-site price).} (a) Let \(R\) be a lawful representation on
\(\Lambda_n\), \(n\ge2\), with lawful-class parameters
\((\sigma,A,\Lambda,\delta,w)\), and assume \(e^{R}\ge0\) pointwise. Let
\(0<s\le1\le b\) be a Haar budget (Lemma 4.1). Let \(g>0\) with
\(g\le\mathrm{threshold}(R)\) and \(\beta g^2=1\). Let \(F\) be a
\(\mu_{\rm prior}\)-integrable complex function and \(z\ge0\),
\(X\in\mathbb R\) such that
\(z|F(U)|\le\exp(\mathrm{exponent}_R(U)+cX)\) for every configuration
\(U\). Then
\[\frac{z\,\big|\int F\,d\mu_{\rm prior}\big|}{\int e^{R}\,d\mu_{\rm prior}}\;\le\;\exp\Big(c\,\pi(R,X,g)-4\dim G\,\log g\Big),\]
where the per-site price is
\[\pi=2\,\mathrm{exCost}(A,\Lambda,\delta)+6+\beta\,\mathrm{baseBudget}(A,g)+X+3\,\mathrm{ent}-4\log s,\qquad \mathrm{ent}=\log b-\dim G\,\log\sigma_R,\]
with \(\sigma_R=\min(1,\ell_R)\), where \(q_R\) and \(\ell_R\) are the
quadratic and linear clip coefficients. Here
\(\mathrm{exCost}=\delta+\mathrm{loopsPrice}\) bounds the loop and
additive activity terms per site, and
\(\mathrm{baseBudget}=3(q_R+A^2)g^2\) bounds the remaining action on
configurations whose bonds all lie in \(\mathrm{shell}(g/4)\). (b) Let
\((\mathcal A_m)\) be the levels of an RG schedule for simple \(G\), and
let \(O=(s,b)\) be a Haar budget. For all sufficiently large levels
\(m\), every period \(p\), every representation \(P\) on the level-\(m\)
torus (with \(c_{m,p}\) sites), every \(k\in\mathbb N\), every
\(X\le1+k\), and every \((g,\Lambda,\delta)\) admissible at level \(m\),
the following holds:
\[\pi(P,X,g)-\frac{4\dim G\,\log g}{c_{m,p}}\;\le\;\mathrm{work item}(O,m,k).\]
Admissible at level \(m\) means: same clip as the stage of \(m\),
\(\beta(P)\le4\hat\beta_m\), \(\beta(P)g^2=1\),
\(0<g\le\mathrm{threshold}(P)\le\mathrm{threshold}_m\), \(\delta\le1\),
\(\Lambda\le\mathrm{loopCap}_m\). The work item is
\[\mathrm{work item}=2(8\,\mathrm{paths}_m+2)+6+3(\mathrm{direct}_m+\mathrm{paths}_m^2)+(1+k)+3\,\mathrm{ent}-4\log s+1.\]
In (b), \(\pi\) uses \(A=\mathrm{paths}_m\) and the clip of stage \(m\).
The work item depends on that stage, \(O\) and \(k\), uniformly over the
admissible couplings and periods. \emph{Status:} theorems. \emph{Lean:}
\href{https://github.com/yang-mills-group/OpenAI-YM-20261028/blob/3456e5e75cf8d4cd558955a5c7638fba2d68f464/YM/Next/RANK/Determinant.lean\#L240-L274}{\texttt{determinant}},
\href{https://github.com/yang-mills-group/OpenAI-YM-20261028/blob/3456e5e75cf8d4cd558955a5c7638fba2d68f464/YM/Next/RANK/Determinant.lean\#L225-L227}{\texttt{protection}},
\href{https://github.com/yang-mills-group/OpenAI-YM-20261028/blob/3456e5e75cf8d4cd558955a5c7638fba2d68f464/YM/Next/RANK/Determinant.lean\#L26-L33}{\texttt{HaarBudget}};
(b)
\href{https://github.com/yang-mills-group/OpenAI-YM-20261028/blob/3456e5e75cf8d4cd558955a5c7638fba2d68f464/YM/Next/RANK/StationMeter.lean\#L119-L123}{\texttt{affordable}},
\href{https://github.com/yang-mills-group/OpenAI-YM-20261028/blob/3456e5e75cf8d4cd558955a5c7638fba2d68f464/YM/Next/RANK/StationMeter.lean\#L28-L36}{\texttt{Meets}},
\href{https://github.com/yang-mills-group/OpenAI-YM-20261028/blob/3456e5e75cf8d4cd558955a5c7638fba2d68f464/YM/Next/RANK/StationMeter.lean\#L41-L44}{\texttt{work\ item}}.

\textbf{Proposition III.11.2 (free pair terms).} Fix \(G\),
\(d,n\in\mathbb N\), and two colored plaquette derivative stencils
\(L,S\) based at the origin, with scaling degrees
\(\Delta_L=2+\mathrm{ord}(L)\le d\) and
\(\Delta_S=2+\mathrm{ord}(S)\le d\). There is \(C\ge0\), depending only
on these data, such that the following holds. Let
\(x\in(\mathbb R^4)^n\) have distinct entries and let \(i\ne j\).
Suppose \(D\ge0\) satisfies \(|\mathrm{Re}\,\Gamma_{(i,a),(j,b)}|\le D\)
for every pair of free-current keys \(a,b\) of scaling degree at most
\(d\). Then, for every \(r\ge1\),
\[r\,\Big|r^{\Delta_L+\Delta_S}\,\langle \mathrm{ev}_L\otimes\mathrm{ev}_S,B_r(i,j)\rangle-\mathrm{light}(L,S)(x_i,x_j)\Big|\le CD.\]
Here
\[B_r(i,j)=\sum_{a,b}\mathrm{Re}\,\Gamma_{(i,a),(j,b)}\;r^{-\mathrm{gen}(a)}t_a\otimes r^{-\mathrm{gen}(b)}t_b,\]
where \(t_a\) is the polynomial jet sample of key \(a\),
\(\mathrm{gen}(a)\) is its scaling degree, and \(\mathrm{ev}_L\) applies
the stencil \(L\). The kernel \(\mathrm{light}(L,S)\) is the
corresponding derivative of the free Maxwell curvature two-point
function. \emph{Status:} theorem. \emph{Lean:}
\href{https://github.com/yang-mills-group/OpenAI-YM-20261028/blob/3456e5e75cf8d4cd558955a5c7638fba2d68f464/YM/Next/RANK/StrideCost.lean\#L70-L98}{\texttt{Prayer.fence}},
\href{https://github.com/yang-mills-group/OpenAI-YM-20261028/blob/3456e5e75cf8d4cd558955a5c7638fba2d68f464/YM/Next/RANK/StrideCost.lean\#L64-L67}{\texttt{deck}}.

\textbf{Proposition III.11.3 (window estimate for the free covariance,
uniform in the regulator).} Let \(\alpha,\eta>0\),
\(\Omega\in\mathbb R\) and \(n\in\mathbb N\). Then there is \(C\ge0\)
with the following property. For every regulator \(f\) in the family
\(\mathrm{Firm}(\alpha,\Omega,\eta)\) of III.3, every \(G\), and every
pair of colored plaquette stencils \(a,b\), the frequency chart
\(K=K_f(a,b)\) is measurable, vanishes for \(|\xi|>2\), is \(C^n\) away
from \(0\), and satisfies:

\begin{itemize}
\tightlist
\item
  \(|K(0)|\le C\);
\item
  \(|\partial^iK(\xi)|\,|\xi|^i\le C\) for \(0<|\xi|\le2\) and
  \(i\le n\);
\item
  \(|\partial^i(K-H_{a,b})(\xi)|\,|\xi|^i\le C|\xi|\) for
  \(0<|\xi|\le1\) and \(i\le n\).
\end{itemize}

Here \(\partial^i\) denotes the \(i\)-th Fréchet derivative, measured in
operator norm, and \(H_{a,b}\) is the degree-zero Maxwell curvature
symbol for the plaquettes and colors of \(a,b\), with heat parameter
\((2\pi)^2\). If \(a,b\) have different colors, then \(K=H_{a,b}=0\).
\emph{Status:} theorem. \emph{Lean:}
\href{https://github.com/yang-mills-group/OpenAI-YM-20261028/blob/3456e5e75cf8d4cd558955a5c7638fba2d68f464/YM/PhysicalNext/RANK_COVARIANCE_ATLAS/Window.lean\#L118-L136}{\texttt{window\_prices}},
\href{https://github.com/yang-mills-group/OpenAI-YM-20261028/blob/3456e5e75cf8d4cd558955a5c7638fba2d68f464/YM/Next/RANK/Radian/Signed.lean\#L47-L53}{\texttt{Window}},
\href{https://github.com/yang-mills-group/OpenAI-YM-20261028/blob/3456e5e75cf8d4cd558955a5c7638fba2d68f464/YM/Next/RANK/Radian/Means.lean\#L16-L19}{\texttt{Cast}}.

\textbf{Proposition III.11.4 (frequent good scales produce corrected
trial systems).} Let \(G\) be simple and \(C\) a Wilson family
satisfying the construction record. Let \(d\) be a degree and
\(\mathcal F\) a stage package for degree \(d\) (fixed band
\(b_0\ge\max(6,\,s_d+3)\), with \(s_d\) the label's scaling saving, and
Taylor order \(\ge d\)). Let \(X\) consist of a real basis of
\(F_d^{\rm free}\), starting with the unit, and finitely many generic
one- and two-point tests that separate this space. Choose
representatives of that basis in the label space. Let \(O\) be a Haar
budget, \(J\) a chart ratio, and \(B\) a binding of the RG history of
\(C\) to its price lists. Fix \(w\in\mathbb R\). Assume that at
infinitely many levels \(k\) the following conditions hold
simultaneously:

\begin{itemize}
\tightlist
\item
  the four price conditions: tender, metering with \((O,w)\), and, for
  some common auxiliary parameters \(l,r\), hours with \((J,B,l,r,w)\)
  and vents with \((J,B,l,r)\);
\item
  the lantern-mark condition for \(X\).
\end{itemize}

Then there is a trial system \(Z\) for \(X\) that travels with the
station and satisfies the purchase conditions:

\begin{itemize}
\tightlist
\item
  depths \(k_1<k_2<\dots\) and test scales
  \(s_i=a_{k_i}\mu_{k_i}^{\,b'}\) (fixed \(b'>0\)) strictly decreasing
  to \(0\), with strictly increasing width levels
  \(\mathrm{width}_i\le k_i\) and the stage index constant on
  \([\mathrm{width}_i,k_i]\);
\item
  for each trial \(i\), coefficient vectors \(v_i(N,p)\) of corrected
  observables that are real, bounded uniformly in \((N,p)\), and have
  first entry the reference unit;
\item
  generic test configurations lying strictly inside the centred width
  cell, whose rescalings converge to the prescribed configuration of
  each row;
\item
  for each trial \(i\), eventually in \(N\) and for every period \(p\),
  on the centred tiling of the cutoff torus: for every row, every basis
  column and every nonempty set \(s\) of insertion modes, the chessboard
  root, at coupling \(\beta_{N,p}\), of the \(s\)-connected part of the
  corrected observable (taken relative to its free Maxwell one- and
  two-point profile) is at most \(\gamma_i=\mu_{k_i}^{-1/32}\to0\).
\end{itemize}

\emph{Status:} theorem, with the frequent-scale conditions carried as
hypotheses. In the main chain they are discharged as follows. In
\href{https://github.com/yang-mills-group/OpenAI-YM-20261028/blob/3456e5e75cf8d4cd558955a5c7638fba2d68f464/YM/MergeProof/Powder.lean\#L48-L71}{\texttt{Accent.rated}},
the lantern marks hold eventually in \(k\) by
\href{https://github.com/yang-mills-group/OpenAI-YM-20261028/blob/3456e5e75cf8d4cd558955a5c7638fba2d68f464/YM/MergeProof/SteamW.lean\#L54-L59}{\texttt{bursts}}.
The four price conditions are supplied to the rank assembly by
\href{https://github.com/yang-mills-group/OpenAI-YM-20261028/blob/3456e5e75cf8d4cd558955a5c7638fba2d68f464/YM/MergeProof/Snow.lean\#L55-L60}{\texttt{franchise}}
(this unit) inside the assembly
\href{https://github.com/yang-mills-group/OpenAI-YM-20261028/blob/3456e5e75cf8d4cd558955a5c7638fba2d68f464/YM/SharpGram/Assembly/Flowers.lean\#L36-L66}{\texttt{jointOffers}}.
That assembly is closed by
\href{https://github.com/yang-mills-group/OpenAI-YM-20261028/blob/3456e5e75cf8d4cd558955a5c7638fba2d68f464/YM/SharpGram/Assembly/Flowers.lean\#L75-L130}{\texttt{uv\_log}}
(III.13), whose only hypotheses are that \(G\) is simple and \(b>0\).
\emph{Lean:}
\href{https://github.com/yang-mills-group/OpenAI-YM-20261028/blob/3456e5e75cf8d4cd558955a5c7638fba2d68f464/YM/PhysicalNext/AFTERNOON_INTERACTING_RANK_COMPARISON/Soft/Star.lean\#L64-L66}{\texttt{cast\_off}},
\href{https://github.com/yang-mills-group/OpenAI-YM-20261028/blob/3456e5e75cf8d4cd558955a5c7638fba2d68f464/YM/PhysicalNext/AFTERNOON_INTERACTING_RANK_COMPARISON/Soft/Star.lean\#L68-L69}{\texttt{wheeled}},
\href{https://github.com/yang-mills-group/OpenAI-YM-20261028/blob/3456e5e75cf8d4cd558955a5c7638fba2d68f464/YM/PhysicalNext/AFTERNOON_INTERACTING_RANK_COMPARISON/Soft/Supra.lean\#L84-L95}{\texttt{v\_culminates}},
\href{https://github.com/yang-mills-group/OpenAI-YM-20261028/blob/3456e5e75cf8d4cd558955a5c7638fba2d68f464/YM/PhysicalNext/AFTERNOON_INTERACTING_RANK_COMPARISON/Soft/Supra.lean\#L177-L222}{\texttt{v\_obtain\_purchases}},
\href{https://github.com/yang-mills-group/OpenAI-YM-20261028/blob/3456e5e75cf8d4cd558955a5c7638fba2d68f464/YM/Next/RANK/Bought.lean\#L68-L84}{\texttt{Purchases}}.

\textbf{Proposition III.11.5 (the field spaces of the limit are those of
the free Maxwell field).} Let \(G\) be simple. Write \(\mathcal L_d\)
for abstract labels and \(\mathcal V_d\) for their fixed coordinate
space. Let \(C\) be a Wilson family satisfying the construction record
and the centred estimate: for each \(d\) and \(R,\varepsilon>0\) there
is \(A\ge0\) such that, at every level \(m\ge\mathrm{ref}(d)\),
eventually in \(N\) and for all \(p\), every label of degree \(d\) and
norm \(\le R\) has centred budget \(A\,a_m^{-d-\varepsilon}\). Assume:

\begin{itemize}
\tightlist
\item
  \begin{enumerate}
  \def\labelenumi{(\roman{enumi})}
  \tightlist
  \item
    a projector presentation of the labels: for each \(d\) and cutoff
    \((N,p)\), a projector \(\Pi_{d,N,p}\) on the coordinate label space
    \(\mathcal V_d\) whose kernel is the image of the Maxwell equation
    ideal \(\mathcal E_d\subset\mathcal L_d\) under an invertible
    corrected label map \(\mathcal L_d\to\mathcal V_d\). For each \(d\),
    both are real and equivariant for all sufficiently large \(N\),
    uniformly in \(p\). Assume also affordability and relay-cost bounds
    along every extraction;
  \end{enumerate}
\item
  (ii) the one-step cost package along the gauge-door maps between
  successive RG levels;
\item
  (iii) panels \(Q_d\), the mark exponents \(c_d=2(w_d+1)\), a chart
  ratio and a binding such that, for each \(d\), the four price
  conditions of III.11.4 hold simultaneously at infinitely many levels,
  with \(w=-3\). The panels are specialized to the family using \(c_d\).
\end{itemize}

Then for every extraction order \(f\) there are a refinement \(g\), a
smooth limit table \(T\) with a differentiation plan \(E\), and the
extraction record of IV.1 along \(f\circ g\) (raw limits, bins, spans,
cleaning to \(T\), Euclidean axioms), together with a seed record for
\((C,f\circ g,T)\). The seed record consists of: cutoff, ordered and
generic point limits, Euclidean axioms, cell hypothesis, and corrected
rank widths. Moreover the following holds: there is a real structure on
the free Maxwell table and, for every \(d\), a linear isomorphism
\[\Phi_d:\;F_d(T)\;\xrightarrow{\ \sim\ }\;F_d^{\rm free}=\mathcal L_d/\mathcal E_d\]
that intertwines complex conjugation and the action of the signed
permutations of the four axes. More precisely, there is a real,
equivariant linear isomorphism \(L_d:\mathcal L_d\to\mathcal V_d\)
carrying \(\mathcal E_d\) onto \(\mathrm{null}_d(T)\). Thus
\(\dim F_d(T)=\dim F_d^{\rm free}\); the identification of null spaces
uses this corrected map. \emph{Status:} theorem with hypotheses
(i)--(iii). All three are discharged in the main chain: (i) by
\href{https://github.com/yang-mills-group/OpenAI-YM-20261028/blob/3456e5e75cf8d4cd558955a5c7638fba2d68f464/YM/MergeProof/Eureka.lean\#L22-L30}{\texttt{reception}}
through
\href{https://github.com/yang-mills-group/OpenAI-YM-20261028/blob/3456e5e75cf8d4cd558955a5c7638fba2d68f464/YM/MergeProof/Inns.lean\#L33-L53}{\texttt{chaletrail}}
(this unit), (iii) by
\href{https://github.com/yang-mills-group/OpenAI-YM-20261028/blob/3456e5e75cf8d4cd558955a5c7638fba2d68f464/YM/MergeProof/Snow.lean\#L55-L60}{\texttt{franchise}},
and the binding by
\href{https://github.com/yang-mills-group/OpenAI-YM-20261028/blob/3456e5e75cf8d4cd558955a5c7638fba2d68f464/YM/MergeProof/Snow.lean\#L28-L43}{\texttt{castle}}.
The one-step package (ii) and the centred estimate are also included in
\href{https://github.com/yang-mills-group/OpenAI-YM-20261028/blob/3456e5e75cf8d4cd558955a5c7638fba2d68f464/YM/MergeProof/Almandine.lean\#L24-L36}{\texttt{Chateau}}
and consumed by
\href{https://github.com/yang-mills-group/OpenAI-YM-20261028/blob/3456e5e75cf8d4cd558955a5c7638fba2d68f464/YM/SharpGram/Assembly/Flowers.lean\#L75-L130}{\texttt{uv\_log}}
(III.13). \emph{Lean:}
\href{https://github.com/yang-mills-group/OpenAI-YM-20261028/blob/3456e5e75cf8d4cd558955a5c7638fba2d68f464/YM/MergeProof/Powder.lean\#L48-L71}{\texttt{Accent.rated}},
\href{https://github.com/yang-mills-group/OpenAI-YM-20261028/blob/3456e5e75cf8d4cd558955a5c7638fba2d68f464/YM/Euclid/Comparator.lean\#L135-L139}{\texttt{Maxwell.Ranks}},
\href{https://github.com/yang-mills-group/OpenAI-YM-20261028/blob/3456e5e75cf8d4cd558955a5c7638fba2d68f464/YM/Spec/Closed.lean\#L20-L26}{\texttt{Spec.Seed}},
\href{https://github.com/yang-mills-group/OpenAI-YM-20261028/blob/3456e5e75cf8d4cd558955a5c7638fba2d68f464/YM/MergeProof/Eureka.lean\#L16-L20}{\texttt{Scotia}},
\href{https://github.com/yang-mills-group/OpenAI-YM-20261028/blob/3456e5e75cf8d4cd558955a5c7638fba2d68f464/YM/Extraction/Center.lean\#L38-L40}{\texttt{Wcmain}}.

\textbf{Proposition III.11.6 (the stress tensor of the limit, normalized
against free Maxwell).} (a) Let \(C\) satisfy the construction record
and \(f\) be an extraction order with ordered and generic point limits,
giving a table \(S\). Assume:

\begin{itemize}
\tightlist
\item
  \(S\) satisfies the Euclidean axioms (with spin structure \(H\)) and
  the cell hypothesis;
\item
  \(W\) is a differentiation plan for \(S\) with its inclusion and
  derivative laws;
\item
  \(\dim F_d(S)\le\dim F_d^{\rm free}\) for every \(d\);
\item
  \(\mathcal T\) is a free Maxwell comparison model, with Euclidean
  axioms \(a\), inclusion, derivative and orientation laws, and a
  divergence-free stress;
\item
  the station hypothesis linking \(C\), \(f\) and \(a\) holds;
\item
  the shell estimate of \(C\) holds.
\end{itemize}

Then \(W\) satisfies the orientation law, and there are a stress record
\(I\) (stress tensor \(\mathcal J\), translation parts, spin rows) and a
scaling \(X\) of \(S\) such that:

\begin{itemize}
\tightlist
\item
  \(\mathcal J\neq0\) and \(\mathcal J\) lies in the kernel space of
  \(H\);
\item
  the scaled tables approach the isotropic panel, and the scaling is
  stable and compatible between the spin structures of \(S\) and
  \(\mathcal T\);
\item
  the scaled \(\mathcal J\) is metered against the free stress tensor
  \(\mathcal T.\mathcal J\);
\item
  \(\mathcal J\) is conserved and symmetric and satisfies the
  translation and rotation Ward identities (III.11.7);
\item
  the action law relating \(\mathcal J\) to the free model holds.
\end{itemize}

\begin{enumerate}
\def\labelenumi{(\alph{enumi})}
\setcounter{enumi}{1}
\tightlist
\item
  Let \(G\) be simple, let \(C\) carry the shell estimate, and assume
  the cofinal grid approximation condition and weak completeness of the
  corrected comparison tests for \(f\), together with the seed record of
  III.11.5 for \((C,f,S)\). Then for every differentiation plan \(W\)
  with inclusion and derivative laws there are Maxwell data \(M\) and a
  stress record \(I\) with \(\mathcal J\ne0\) satisfying the UV
  interface record \(\mathrm{UV}(C,f,S,W,M,I)\). \emph{Status:}
  theorems. In (b), the free model comes from
  \href{https://github.com/yang-mills-group/OpenAI-YM-20261028/blob/3456e5e75cf8d4cd558955a5c7638fba2d68f464/YM/Next/MAXWELL/Lit.lean\#L45-L48}{\texttt{maxwell\_supplies}}
  (VI.2). The station hypothesis is built from the seed and weak
  completeness, and the hypothesis
  \href{https://github.com/yang-mills-group/OpenAI-YM-20261028/blob/3456e5e75cf8d4cd558955a5c7638fba2d68f464/YM/Next/UV_STRESS/Invert.lean\#L108}{\texttt{DimAbove}}
  from the seed's corrected rank widths; this is how III.11.5 feeds
  III.11.6. The shell budget is proved by
  \href{https://github.com/yang-mills-group/OpenAI-YM-20261028/blob/3456e5e75cf8d4cd558955a5c7638fba2d68f464/YM/MergeProof/North/Economist.lean\#L23-L32}{\texttt{laurelE}}
  and
  \href{https://github.com/yang-mills-group/OpenAI-YM-20261028/blob/3456e5e75cf8d4cd558955a5c7638fba2d68f464/YM/MergeProof/North/Economist.lean\#L34}{\texttt{fugacious}}
  (this unit) inside
  \href{https://github.com/yang-mills-group/OpenAI-YM-20261028/blob/3456e5e75cf8d4cd558955a5c7638fba2d68f464/YM/SharpGram/Assembly/Flowers.lean\#L75-L130}{\texttt{uv\_log}}.
  \emph{Lean:}
  \href{https://github.com/yang-mills-group/OpenAI-YM-20261028/blob/3456e5e75cf8d4cd558955a5c7638fba2d68f464/YM/MergeProof/Emanate.lean\#L15-L30}{\texttt{Acme.readyA}},
  \href{https://github.com/yang-mills-group/OpenAI-YM-20261028/blob/3456e5e75cf8d4cd558955a5c7638fba2d68f464/YM/MergeProof/Emanate.lean\#L31-L37}{\texttt{Acme.emanate}},
  \href{https://github.com/yang-mills-group/OpenAI-YM-20261028/blob/3456e5e75cf8d4cd558955a5c7638fba2d68f464/YM/MergeProof/West/Ozone.lean\#L16-L19}{\texttt{Aroma}}.
\end{enumerate}

\textbf{Proposition III.11.7 (Ward identity: the flux of the stress
against a Killing field is a local field).} Let \(C\) satisfy the
construction record, let \(f\) be an extraction order with the cell
hypothesis and generic point limits \(S\), and suppose \(S\) satisfies
the Euclidean axioms with real structure \(H\). Let \(W\) satisfy the
derivative laws, and assume the shell estimate of III.11.6. Let:

\begin{itemize}
\tightlist
\item
  \(\mathcal J\) be a stress tensor of \(S\) that is conserved with
  respect to \(W\), and \(v\in F_d(S)\);
\item
  \(U\) be a smooth Killing field on \(\mathbb R^4\), homogeneous of
  degree \(p\) (\(U(ry)=r^pU(y)\) for \(r>0\));
\item
  \(d,D,p\in\mathbb N\) and \(D+3+p=d+4\).
\end{itemize}

Then there is \(B\in F_D(S)\) with the following property. For every
\(n\), all fields \(q_1,\dots,q_n\), every configuration of distinct
points \(x_0,\dots,x_n\), and every window \(\chi\) adapted to \(x\)
(smooth, compactly supported, equal to \(1\) near \(x_0\), with support
avoiding \(x_1,\dots,x_n\)), the flux of \(\mathcal J\) against
\(U(\cdot-x_0)\) and \(\chi\), inserted with \(v\) at \(x_0\) and
\(q_k\) at \(x_k\), equals the correlation of \(B\) at \(x_0\) with the
\(q_k\) at \(x_k\). The sign convention is
\[\mathrm{flux}=-\sum_{i,j}\int_{\mathbb R^4}\partial_i\chi(y)\,U_j(y-x_0)\,\langle\mathcal J_{ij}(y)\,v(x_0)\,q_1(x_1)\cdots q_n(x_n)\rangle\,dy.\]

Nonzero homogeneous Killing fields have \(p=0\) (translations,
\(D=d+1\)) or \(p=1\) (rotations, \(D=d\)). \emph{Status:} theorem.
\emph{Lean:}
\href{https://github.com/yang-mills-group/OpenAI-YM-20261028/blob/3456e5e75cf8d4cd558955a5c7638fba2d68f464/YM/MergeProof/West/Bravo.lean\#L55-L89}{\texttt{annular\_fieldA}},
\href{https://github.com/yang-mills-group/OpenAI-YM-20261028/blob/3456e5e75cf8d4cd558955a5c7638fba2d68f464/YM/MergeProof/West/Blanca.lean\#L49-L120}{\texttt{tenet}},
\href{https://github.com/yang-mills-group/OpenAI-YM-20261028/blob/3456e5e75cf8d4cd558955a5c7638fba2d68f464/YM/MergeProof/West/Orange.lean\#L33-L63}{\texttt{shells\_law}}.

\textbf{Further inputs.}

\begin{itemize}
\tightlist
\item
  \textbf{One-step and root-test estimates for polynomial and
  coefficient data.}
  \href{https://github.com/yang-mills-group/OpenAI-YM-20261028/blob/3456e5e75cf8d4cd558955a5c7638fba2d68f464/YM/MergeProof/Meringue.lean\#L32-L102}{\texttt{single}}
  and
  \href{https://github.com/yang-mills-group/OpenAI-YM-20261028/blob/3456e5e75cf8d4cd558955a5c7638fba2d68f464/YM/MergeProof/Gelatin.lean\#L23-L92}{\texttt{asideP}}
  apply the coefficient bound to already packed coefficients across one
  gauge-door step, with parameters scaled by \(\eta t\). Also in this
  group:
  \href{https://github.com/yang-mills-group/OpenAI-YM-20261028/blob/3456e5e75cf8d4cd558955a5c7638fba2d68f464/YM/Next/RANK/Dock.lean\#L179-L185}{\texttt{grant}},
  \href{https://github.com/yang-mills-group/OpenAI-YM-20261028/blob/3456e5e75cf8d4cd558955a5c7638fba2d68f464/YM/MergeProof/West/Lys.lean\#L40-L60}{\texttt{y\_ready}}
  and
  \href{https://github.com/yang-mills-group/OpenAI-YM-20261028/blob/3456e5e75cf8d4cd558955a5c7638fba2d68f464/YM/Next/RANK/Festoon.lean\#L57-L126}{\texttt{console}}.
\item
  \textbf{Eventual estimates and uniform choices.} These include bounds
  at sufficiently large levels, uniform versions on compact sets of
  distinct-point configurations, and choices of enlarged waiting
  windows:
  \href{https://github.com/yang-mills-group/OpenAI-YM-20261028/blob/3456e5e75cf8d4cd558955a5c7638fba2d68f464/YM/PhysicalNext/AFTERNOON_INTERACTING_RANK_COMPARISON/Chimes.lean\#L113-L153}{\texttt{pedestals}},
  \href{https://github.com/yang-mills-group/OpenAI-YM-20261028/blob/3456e5e75cf8d4cd558955a5c7638fba2d68f464/YM/PhysicalNext/AFTERNOON_INTERACTING_RANK_COMPARISON/Attic.lean\#L157-L203}{\texttt{chamber}},
  \href{https://github.com/yang-mills-group/OpenAI-YM-20261028/blob/3456e5e75cf8d4cd558955a5c7638fba2d68f464/YM/PhysicalNext/AFTERNOON_INTERACTING_RANK_COMPARISON/Morning.lean\#L35-L53}{\texttt{herald}},
  \href{https://github.com/yang-mills-group/OpenAI-YM-20261028/blob/3456e5e75cf8d4cd558955a5c7638fba2d68f464/YM/PhysicalNext/AFTERNOON_INTERACTING_RANK_COMPARISON/Focusing.lean\#L43-L101}{\texttt{scopes}},
  \href{https://github.com/yang-mills-group/OpenAI-YM-20261028/blob/3456e5e75cf8d4cd558955a5c7638fba2d68f464/YM/MergeProof/West/Brite.lean\#L31-L122}{\texttt{on\_hold}},
  \href{https://github.com/yang-mills-group/OpenAI-YM-20261028/blob/3456e5e75cf8d4cd558955a5c7638fba2d68f464/YM/Next/RANK/Promises.lean\#L37-L127}{\texttt{on\_demand}},
  \href{https://github.com/yang-mills-group/OpenAI-YM-20261028/blob/3456e5e75cf8d4cd558955a5c7638fba2d68f464/YM/Next/RANK/Distant.lean\#L69-L100}{\texttt{chamber\_now}},
  \href{https://github.com/yang-mills-group/OpenAI-YM-20261028/blob/3456e5e75cf8d4cd558955a5c7638fba2d68f464/YM/MergeProof/West/Esme.lean\#L42-L57}{\texttt{mural\_cast}},
  \href{https://github.com/yang-mills-group/OpenAI-YM-20261028/blob/3456e5e75cf8d4cd558955a5c7638fba2d68f464/YM/MergeProof/West/Tael.lean\#L58-L90}{\texttt{earnings}},
  \href{https://github.com/yang-mills-group/OpenAI-YM-20261028/blob/3456e5e75cf8d4cd558955a5c7638fba2d68f464/YM/MergeProof/East/Reach.lean\#L27-L42}{\texttt{lava\_v}},
  \href{https://github.com/yang-mills-group/OpenAI-YM-20261028/blob/3456e5e75cf8d4cd558955a5c7638fba2d68f464/YM/MergeProof/East/Gran.lean\#L38-L51}{\texttt{sponsored}},
  \href{https://github.com/yang-mills-group/OpenAI-YM-20261028/blob/3456e5e75cf8d4cd558955a5c7638fba2d68f464/YM/Next/RANK/Starring.lean\#L31-L37}{\texttt{boot}},
  \href{https://github.com/yang-mills-group/OpenAI-YM-20261028/blob/3456e5e75cf8d4cd558955a5c7638fba2d68f464/YM/Next/RANK/Tilt.lean\#L36-L59}{\texttt{forge}}.
  The weak-completeness theorem in this list enlarges the waiting
  window. The order of the choices is part of each statement.
\item
  \textbf{Polynomial and integral identities.} Wick expansion, Taylor
  polynomials in link variables, pushforward of densities as sums over
  sectors, and factorization on disjoint supports; see §4.3.
\end{itemize}

\subsubsection{2. What it uses}\label{what-it-uses-13}

\textbf{Inputs from other units.}

\begin{itemize}
\tightlist
\item
  I.1 (lattice, Wilson measure): the periodic torus, plaquette values,
  the product Haar prior, gauge transformations, the chordal distance on
  \(G\)
  (\href{https://github.com/yang-mills-group/OpenAI-YM-20261028/blob/3456e5e75cf8d4cd558955a5c7638fba2d68f464/YM/Wilson/Law.lean\#L55}{\texttt{prior}},
  \href{https://github.com/yang-mills-group/OpenAI-YM-20261028/blob/3456e5e75cf8d4cd558955a5c7638fba2d68f464/YM/Gauge/MatrixGroup.lean\#L115}{\texttt{chord}}).
\item
  II.1 (field labels): the label spaces \(\mathcal L_d\) with their
  coordinates and real structure; the free Wick realization of labels,
  whose kernel is the Maxwell equation ideal \(\mathcal E_d\); the weave
  and portfolio data used for the corrected label maps
  (\href{https://github.com/yang-mills-group/OpenAI-YM-20261028/blob/3456e5e75cf8d4cd558955a5c7638fba2d68f464/YM/Maxwell/Equation.lean\#L72}{\texttt{equationIdeal}}).
\item
  II.2 (polymer activities, Wilson family, construction record, cells):
  lawful representations and their exponent and gas; the Wilson family
  and its construction record; the tilings and the chessboard root
  (\href{https://github.com/yang-mills-group/OpenAI-YM-20261028/blob/3456e5e75cf8d4cd558955a5c7638fba2d68f464/YM/RG/History.lean\#L181-L183}{\texttt{Rep.Lawful}}).
\item
  III.2 (choice of constants): the levels and stages (\(\hat\beta_m\),
  thresholds, loop caps, clip coefficients, stage index \(J_m\)), the
  pre-\(L\) and post-\(L\) constants, delays and windows; the order of
  choice supplies the eventual estimates and waiting windows listed
  above, including the stage construction
  \href{https://github.com/yang-mills-group/OpenAI-YM-20261028/blob/3456e5e75cf8d4cd558955a5c7638fba2d68f464/YM/MergeProof/North/Expound.lean\#L68-L77}{\texttt{ruese}}.
\item
  III.3, III.4 (free covariances, Gaussian integration): the free
  covariance and its frequency chart, the regulator families with their
  symbol bounds (used in Lemma 4.10), the window-to-kernel transfer
  \href{https://github.com/yang-mills-group/OpenAI-YM-20261028/blob/3456e5e75cf8d4cd558955a5c7638fba2d68f464/YM/Next/RANK/Radian/Signed.lean\#L56-L92}{\texttt{intended}}.
\item
  III.9, III.10 (marginal flow, one nonlinear step): growth of
  \(\hat\beta_m\), the RG history (``course'') from cutoff \(N\) to
  level \(k\) with its exact vacuum identity, the one-step bounds along
  paths, and the waiting-window engine
  \href{https://github.com/yang-mills-group/OpenAI-YM-20261028/blob/3456e5e75cf8d4cd558955a5c7638fba2d68f464/YM/PhysicalNext/RG_NEXT_COMPLETE_REPEATABLE_NONLINEAR_STEP/Pane/Divisors.lean\#L26-L135}{\texttt{Wait.engine}}.
\item
  III.13 (common Wilson family):
  \href{https://github.com/yang-mills-group/OpenAI-YM-20261028/blob/3456e5e75cf8d4cd558955a5c7638fba2d68f464/YM/SharpGram/Assembly/Flowers.lean\#L75-L130}{\texttt{uv\_log}},
  the theorem that closes the chain and supplies every hypothesis of
  III.11.4--III.11.6 on the main path.
\item
  IV.1 (extraction): extraction orders and refinements, raw limits,
  projection inclusion, spans, cleaning to a smooth table, and the
  answers to the work items D01--D06
  (\href{https://github.com/yang-mills-group/OpenAI-YM-20261028/blob/3456e5e75cf8d4cd558955a5c7638fba2d68f464/YM/Next/EXTRACTION/Base.lean\#L42-L51}{\texttt{Base}}).
\item
  VI.1, VI.2 (Ward identities, free Maxwell stress): windows, flux,
  Killing fields, conservation, finite faithfulness of testers (A04),
  the free lighthouse and its stress tensor, the canonical free Maxwell
  data, and the production of the current.
\item
  Consumers: VI.1
  (\href{https://github.com/yang-mills-group/OpenAI-YM-20261028/blob/3456e5e75cf8d4cd558955a5c7638fba2d68f464/YM/Climb/Beams.lean\#L53-L63}{\texttt{Climb.comparison}}
  also derives Maxwell ranks from the action law), VII.4 (OPE columns),
  and VIII.1 (the record of local field laws contains the Maxwell ranks;
  \href{https://github.com/yang-mills-group/OpenAI-YM-20261028/blob/3456e5e75cf8d4cd558955a5c7638fba2d68f464/YM/Org/Input.lean\#L177-L186}{\texttt{Org.UV}}
  is built by III.11.6(b)).
\end{itemize}

\textbf{Standard mathematics.} The chessboard estimate for
reflection-positive measures
(\href{https://www.math.caltech.edu/SimonPapers/93.pdf}{Fröhlich--Israel--Lieb--Simon,
Comm. Math. Phys. 62 (1978)}; Glimm--Jaffe, \emph{Quantum Physics}, 2nd
ed.), together with reflection positivity of the Wilson action
(Osterwalder--Seiler, Ann. Phys. 110 (1978)). Axial (tree) gauge fixing
on a torus and invariance of Haar measure
(\href{https://link.springer.com/book/10.1007/3-540-11559-5}{Seiler,
\emph{Gauge Theories as a Problem of Constructive Quantum Field Theory
and Statistical Mechanics}, Lecture Notes in Physics 159 (1982)}). Haar
volume of small balls in a compact Lie group. Wick (Isserlis) expansion.
Schur complements for Gaussian conditioning. Symbol estimates in
frequency space. Finite-dimensional duality (a vector lies in the range
of a linear map if and only if it is annihilated by every functional
that annihilates the range). Ward identities for a conserved symmetric
stress tensor and Killing fields of \(\mathbb R^4\).

\textbf{Imported constants and assumptions.}

\begin{itemize}
\tightlist
\item
  The Haar-budget constants \((s,b)\) exist by
  \href{https://github.com/yang-mills-group/OpenAI-YM-20261028/blob/3456e5e75cf8d4cd558955a5c7638fba2d68f464/YM/Next/RANK/Haar.lean\#L170-L211}{\texttt{integral\_budget}}
  and
  \href{https://github.com/yang-mills-group/OpenAI-YM-20261028/blob/3456e5e75cf8d4cd558955a5c7638fba2d68f464/YM/Next/RANK/Determinant.lean\#L35-L36}{\texttt{exists\_budget}}.
  They are fixed once for all stages and volumes.
\item
  Hypotheses carried by the outputs: the projector presentation, the
  one-step cost package, the centred budget, the frequent-scale price
  conditions, the binding, the shell budget, Grill and weak
  completeness. On the main path all are discharged before
  \href{https://github.com/yang-mills-group/OpenAI-YM-20261028/blob/3456e5e75cf8d4cd558955a5c7638fba2d68f464/YM/SharpGram/Assembly/Flowers.lean\#L75-L130}{\texttt{uv\_log}}
  (III.13), whose only hypotheses are \(G\) simple and \(b>0\). The
  producers in this unit are
  \href{https://github.com/yang-mills-group/OpenAI-YM-20261028/blob/3456e5e75cf8d4cd558955a5c7638fba2d68f464/YM/MergeProof/Eureka.lean\#L22-L30}{\texttt{reception}}/\href{https://github.com/yang-mills-group/OpenAI-YM-20261028/blob/3456e5e75cf8d4cd558955a5c7638fba2d68f464/YM/MergeProof/Inns.lean\#L33-L53}{\texttt{chaletrail}},
  \href{https://github.com/yang-mills-group/OpenAI-YM-20261028/blob/3456e5e75cf8d4cd558955a5c7638fba2d68f464/YM/MergeProof/Snow.lean\#L55-L60}{\texttt{franchise}},
  \href{https://github.com/yang-mills-group/OpenAI-YM-20261028/blob/3456e5e75cf8d4cd558955a5c7638fba2d68f464/YM/MergeProof/Snow.lean\#L28-L43}{\texttt{castle}},
  \href{https://github.com/yang-mills-group/OpenAI-YM-20261028/blob/3456e5e75cf8d4cd558955a5c7638fba2d68f464/YM/MergeProof/North/Economist.lean\#L23-L32}{\texttt{laurelE}},
  \href{https://github.com/yang-mills-group/OpenAI-YM-20261028/blob/3456e5e75cf8d4cd558955a5c7638fba2d68f464/YM/MergeProof/North/Economist.lean\#L34}{\texttt{fugacious}}
  and
  \href{https://github.com/yang-mills-group/OpenAI-YM-20261028/blob/3456e5e75cf8d4cd558955a5c7638fba2d68f464/YM/MergeProof/South/FinalG.lean\#L76-L83}{\texttt{grilled}}.
\item
  The structures in
  \href{https://github.com/yang-mills-group/OpenAI-YM-20261028/blob/3456e5e75cf8d4cd558955a5c7638fba2d68f464/YM/Next/RANK/LedgerGate.lean}{\texttt{LedgerGate.lean}}
  (one-step ``Estimate'' packages: ignition, local output labellings,
  non-leading terms) are definitions, not theorems. Their docstring
  calls them ``physics inputs, not consequences of the shape/lawful
  structure''. They are conditional one-step interfaces. The restarted
  comparison constructs principal integration data from the binding and
  the one-step estimates; the older interfaces need not all occur in
  that route. Their complete correspondence with the restarted
  interfaces remains open in §7.
\end{itemize}

\subsubsection{3. Strategy}\label{strategy-13}

The field-space comparison has two parts. Algebraic equations give an
upper bound on the dimension of each interacting field space.
Short-distance one- and two-point tests give the reverse bound. A
conjugation of projectors then identifies the real representations of
the signed permutation group.

At each cutoff, the projector kernel is the image of \(\mathcal E_d\)
under an invertible corrected label map. Free probes separate the
on-shell sector, so
\(\dim F_d^{\rm free}=\dim\mathcal V_d-\dim\ker\Pi_{d,N,p}\). The relay
estimates carry these equations to the limit, giving the upper bound.
Along a sequence of good levels, the corrected test moments approach
their free values. Since the free tests separate a basis, the corrected
vectors are eventually linearly independent modulo the limit null space.
This gives the lower bound.

The two bounds imply \(\mathrm{null}_d(T)=\ker\Pi_\infty\) and
\(\dim F_d(T)=\dim F_d^{\rm free}\). An invertible operator close to the
identity conjugates a cutoff projector to \(\Pi_\infty\). Composing it
with the corrected label map identifies the Maxwell ideal with the
interacting null space, preserving conjugation and signed permutations.

The quantitative engine is the chessboard estimate. The product of
reflected copies of a connected part, integrated against the Wilson
measure at cutoff \(N\), is carried by the exact RG history to a coarse
level \(m\). There the number of cells equals the number of sites
\(c_{m,p}\). The flowed observable is bounded pointwise by
\(r^{-|s|q}\exp(\mathrm{exponent}+qX)\) for a set \(s\) of marked modes,
and Proposition III.11.1 bounds the resulting ratio of integrals by
\(e^{c_{m,p}\pi}g^{-4\dim G}\).

The determinant estimate compares two Haar integrals:

\begin{itemize}
\tightlist
\item
  \textbf{Numerator.} Extract the factor \(e^{6c}\) from the plaquette
  estimate. A triangular choice of \(3c-3\) plaquettes, each owning a
  bond, gives an integral equal to
  \((\int e^{-\mathrm{chord}(u)/(g/\sigma_R)}\,dh(u))^{3c-3}\), bounded
  by \((b(g/\sigma_R)^{\dim G})^{3c-3}\).
\item
  \textbf{Vacuum.} Tree gauge leaves \(3c+1\) free bonds. Restricting
  each to a Haar shell of radius \(g/4\) gives the lower bound
  \((s g^{\dim G})^{3c+1}\).
\end{itemize}

The powers of \(g\) cancel up to \(g^{-4\dim G}=\beta^{2\dim G}\), which
does not depend on the volume. After taking the \(c_{m,p}\)-th root, its
logarithmic contribution is at most \(1\), because
\(2\dim G\log\beta\le8\dim G\,\hat\beta_m\le a_m^{-1}\le c_{m,p}\)
eventually. So the root test costs a per-site bound depending only on
the stage parameters. The constants are chosen so that the allowance is
\(e^{2I_k}\) and the error bound is \(e^{-I_k}\), where
\(I_k=\log\mu_k/32\ge\mathrm{work item}\) for the one- and two-point
tests; the root of every nonempty connected part is then
\(\le e^{I_k}/e^{2I_k}=e^{-I_k}\to0\).

All conditions are needed only at \textbf{infinitely many} good levels
\(k\). A diagonal choice of depths then produces the trial sequence
(III.11.4). The covariance splitting identifies the fluctuation law used
in the one-step comparison. The regulator estimates supply the frequency
windows, while the uniform train constants control their propagation;
the tender condition includes the numerical clearance needed for both.

The stress tensor (III.11.6, III.11.7) is built by comparing the limit
with a free Maxwell model along a scaling sequence. The Ward identities
follow from a flux argument on shrinking annuli of radii \(r_l\to0\):

\begin{itemize}
\tightlist
\item
  by conservation the flux does not depend on the window;
\item
  after scaling, the flux is an integral over the annulus
  \(r_l/4<|y-x_0|<2r_l\);
\item
  a bound on correlations in annuli (Lemma 4.23) shows that combinations
  of fluxes annihilating all degree-\(D\) testers tend to \(0\);
\item
  finite-dimensional duality and finite faithfulness of the testers then
  give one local field \(B\) representing the flux.
\end{itemize}

\subsubsection{4. Steps}\label{steps-13}

\paragraph{4.1 Haar estimates and the determinant
budget}\label{haar-estimates-and-the-determinant-budget}

\textbf{Lemma 4.1 (Haar budget).} There are \(0<s\le1\le b\) such that
\(h(\mathrm{shell}(r/4))\ge s\,r^{\dim G}\) for \(0<r\le1\), and
\(\int e^{-\mathrm{chord}(u)/t}\,dh(u)\le b\,t^{\dim G}\) for \(t>0\).
\emph{Sketch:} in exponential coordinates near \(1\), the chord is
comparable to the Lie-algebra norm and the Haar density is bounded above
and below. Compactness extends the upper volume bound
\(h(\mathrm{shell}(r))\le Cr^{\dim G}\) to all \(r>0\). Majorize the
exponential by \(\sum_{k\ge0}e^{-k}1_{\mathrm{shell}((k+1)t)}\) and sum
\(C t^{\dim G}\sum_{k\ge0}(k+1)^{\dim G}e^{-k}\). \emph{Lean:}
\href{https://github.com/yang-mills-group/OpenAI-YM-20261028/blob/3456e5e75cf8d4cd558955a5c7638fba2d68f464/YM/Next/RANK/Haar.lean\#L170-L211}{\texttt{integral\_budget}}.

\textbf{Lemma 4.2 (plaquette elimination).} Let \(n\ge2\) and let
\(\varphi:G\to[0,1]\) be continuous. Then
\(\int\prod_{p}\varphi(U_p)\,d\mu_{\rm prior}(U)\le\big(\int\varphi\,dh\big)^{3c-3}\),
the product running over all \(6c\) plaquettes of \(\Lambda_n\).
\emph{Sketch:} fix the spanning tree in which each site \(x\ne0\) is
joined to \(x-e_{i(x)}\), where \(i(x)\) is the highest axis with
\(x_{i}\ne0\). The \(3c+1\) non-tree bonds are grouped into tiers by
axis. Within each tier, take the wrap bonds first, except one ``handle''
per tier, then the transverse bonds from higher to lower height. This
gives \(3c-3\) swept bonds \(e\), each assigned a plaquette
\(p(e)\ni e\). Triangularity: if \(e\in p(e')\) with \(e'\ne e\), then
\(e'\) has strictly higher priority than \(e\). Integrate the bond of
highest priority first. By left and right invariance of Haar measure,
\(U_{p(e)}\) is Haar-distributed as a function of \(U_e\), and no other
selected plaquette contains \(e\). So
\(\int\varphi(U_{p(e)})\,\Psi\,d\mu=\big(\int\varphi\,dh\big)\int\Psi\,d\mu\).
By induction, the product over the selected plaquettes integrates to
exactly \((\int\varphi)^{3c-3}\); the other plaquette factors are
bounded by \(1\). \emph{Lean:}
\href{https://github.com/yang-mills-group/OpenAI-YM-20261028/blob/3456e5e75cf8d4cd558955a5c7638fba2d68f464/YM/Next/RANK/Axial.lean\#L368-L383}{\texttt{elimination}},
\href{https://github.com/yang-mills-group/OpenAI-YM-20261028/blob/3456e5e75cf8d4cd558955a5c7638fba2d68f464/YM/Next/RANK/Axial.lean\#L347-L365}{\texttt{charged\_integral}},
\href{https://github.com/yang-mills-group/OpenAI-YM-20261028/blob/3456e5e75cf8d4cd558955a5c7638fba2d68f464/YM/Next/RANK/Axial.lean\#L314-L345}{\texttt{pull\_link}},
\href{https://github.com/yang-mills-group/OpenAI-YM-20261028/blob/3456e5e75cf8d4cd558955a5c7638fba2d68f464/YM/Next/RANK/Axial.lean\#L278-L303}{\texttt{independent}}.
\emph{(proof read; the combinatorics of the sweep, i.e.~the count and
triangularity lemmas, statement only)}

\textbf{Lemma 4.3 (vacuum lower bound).} Under the hypotheses of
III.11.1(a),
\(\int e^{R}\,d\mu_{\rm prior}\ge\exp\big(\mathrm{scalar}(R)-c(\mathrm{exCost}+\beta\,\mathrm{baseBudget})\big)\,(s\,g^{\dim G})^{3c+1}\).
\emph{Sketch:} \(e^{R}\) is gauge invariant
(\href{https://github.com/yang-mills-group/OpenAI-YM-20261028/blob/3456e5e75cf8d4cd558955a5c7638fba2d68f464/YM/Next/RANK/Determinant.lean\#L46-L68}{\texttt{gauge\_eval}}),
so its integral equals the integral over the \(3c+1\) free bonds with
tree bonds set to \(1\)
(\href{https://github.com/yang-mills-group/OpenAI-YM-20261028/blob/3456e5e75cf8d4cd558955a5c7638fba2d68f464/YM/Next/RANK/Axial.lean\#L115-L119}{\texttt{free\_count}}).
Restrict to the set where every free bond lies in
\(\mathrm{shell}(g/4)\). Every bond then has chord \(\le g/4\), so every
plaquette has chord \(\le g\). On calm configurations the polymer gas
factor is \(1\), and the exponent is
\(\ge\mathrm{scalar}-c(\mathrm{exCost}+\beta\,\mathrm{baseBudget})\).
Lemma 4.1 with \(r=g\le1\) bounds the Haar volume of the restriction set
from below. Positivity of \(e^R\) handles the complement. \emph{Lean:}
\href{https://github.com/yang-mills-group/OpenAI-YM-20261028/blob/3456e5e75cf8d4cd558955a5c7638fba2d68f464/YM/Next/RANK/Determinant.lean\#L76-L111}{\texttt{vacuum\_bound}}.

\textbf{Lemma 4.4 (exponent bound and numerator).} Under the hypotheses
of III.11.1(a), with \(E=g/\sigma_R\): for every \(U\),
\(\mathrm{exponent}_R(U)\le\mathrm{scalar}+c(\mathrm{exCost}+6)-\sum_p\mathrm{chord}(U_p)/E\).
Consequently
\(z|\int F|\le\exp(\mathrm{scalar}+c(\mathrm{exCost}+6+X))\,(b\,E^{\dim G})^{3c-3}\).
\emph{Sketch:} for the clipped plaquette action at \(\beta=g^{-2}\), one
has \(\mathrm{chord}/E-1\le\beta\cdot(\text{base term})\):

\begin{itemize}
\tightlist
\item
  in the quadratic regime, with \(x=\mathrm{chord}/g\), use
  \(\sigma_R x-1\le\sigma_R x^2/2\) and \(\sigma_R\leq_R\);
\item
  above the clipping threshold \(t\), use \(g\le t\) and
  \(\sigma_R\le\ell_R\) to obtain
  \(\sigma_R\,\mathrm{chord}/g\le\beta\,\ell_R\,t\,\mathrm{chord}\); the
  constant term of the linear branch is nonnegative.
\end{itemize}

The principal part is an average over signed reflections of the
pulled-back configuration. The chord sum is invariant under these, and
the loop and activity corrections are bounded by \(c\) times their
per-site prices. This gives the first bound, with \(6=\) the number of
plaquette orientations accounting for the \(-1\) per plaquette. Then
\(z|F|\le e^{\mathrm{scalar}+c(\dots)}\prod_p e^{-\mathrm{chord}(U_p)/E}\),
and Lemma 4.2 with \(\varphi=e^{-\mathrm{chord}/E}\) together with Lemma
4.1 conclude. \emph{Lean:}
\href{https://github.com/yang-mills-group/OpenAI-YM-20261028/blob/3456e5e75cf8d4cd558955a5c7638fba2d68f464/YM/Next/RANK/Determinant.lean\#L153-L181}{\texttt{exponent\_retains}},
\href{https://github.com/yang-mills-group/OpenAI-YM-20261028/blob/3456e5e75cf8d4cd558955a5c7638fba2d68f464/YM/Next/RANK/Determinant.lean\#L189-L217}{\texttt{numerator}}.

\textbf{Lemma 4.5 (the ratio; III.11.1(a)).} Divide Lemma 4.4 by Lemma
4.3. With \(u=3c-3\) and \(l=3c+1=u+4\), the exponent is
\[c(2\,\mathrm{exCost}+6+\beta\,\mathrm{baseBudget}+X)+u\,\mathrm{ent}-l\log s-4\dim G\log g .\]
Since \(\mathrm{ent}\ge0\), \(\log s\le0\), \(u\le3c\) and \(l\le4c\),
this is at most \(c\pi-4\dim G\log g\). \emph{Lean:}
\href{https://github.com/yang-mills-group/OpenAI-YM-20261028/blob/3456e5e75cf8d4cd558955a5c7638fba2d68f464/YM/Next/RANK/Determinant.lean\#L240-L274}{\texttt{determinant}}.

\textbf{Lemma 4.6 (per-site work item; III.11.1(b)).} Under the
assumptions and uniformly over the parameter ranges of III.11.1(b), for
all sufficiently large \(m\):

\begin{itemize}
\tightlist
\item
  \begin{enumerate}
  \def\labelenumi{(\roman{enumi})}
  \tightlist
  \item
    \(\pi\le\mathrm{work item}-1\);
  \end{enumerate}
\item
  (ii) \(-4\dim G\log g/c_{m,p}\le1\).
\end{itemize}

\emph{Sketch:} (i) uses \(\mathrm{exCost}\le1+(8\,\mathrm{paths}+1)\),
from \(\delta\le1\), \(\Lambda t^2\le1\) and
\(\hat\beta_m t^3,\hat\beta_m t^4\le1\) eventually; and
\(\beta\,\mathrm{baseBudget}=6(\mathrm{direct}+\mathrm{paths}^2)/2\) by
\(\beta g^2=1\). (ii) Since
\(-4\log g=2\log\beta\le2\beta\le8\hat\beta_m\), and \(\hat\beta_m\)
grows linearly in \(\log a_m^{-1}\), one gets
\(8\dim G\,\hat\beta_m\le a_m^{-1}\le c_{m,p}\) eventually. \emph{Lean:}
\href{https://github.com/yang-mills-group/OpenAI-YM-20261028/blob/3456e5e75cf8d4cd558955a5c7638fba2d68f464/YM/Next/RANK/StationMeter.lean\#L119-L123}{\texttt{affordable}},
\href{https://github.com/yang-mills-group/OpenAI-YM-20261028/blob/3456e5e75cf8d4cd558955a5c7638fba2d68f464/YM/Next/RANK/StationMeter.lean\#L46-L79}{\texttt{rolls\_law}}.

\textbf{Lemma 4.7 (chessboard root).} Let \(T\) tile the cutoff torus
into \(q\) cells, where \(q\) is the number of sites of the level-\(m\)
torus. Suppose the exact RG history has terminal density \(P\). Assume
that deleting the additive activities of \(P\) gives a lawful
representation with activity parameter \(0\), and that the deleted
additive term has absolute value at most \(q\delta\). Let \(V\) be a
finite set of modes and let every component \(\Phi_s\) be continuous.
Assume \(r>0\), \(e\ge0\), \(g>0\), \(g\le\mathrm{threshold}(P)\) and
\(\beta(P)g^2=1\). Suppose the local activity bound for the flowed
stamped polynomial, at radius \(r\) and exponents \((e,l,L)\), gives
\[r^{q\,\mathrm{card}\,s}|F_s(U)|\le\exp\big(\mathrm{exponent}_P(U)+q(L+\mathrm{card}\,V\,l)\big)\]
for every terminal configuration \(U\) and every \(s\subseteq V\). Here
\(F_s\) is the flow of the unnormalized Wilson density multiplied by the
product of the reflected copies of \(\Phi_s\). Then
\[\rho_T(\Phi_s)\le\frac{\exp\big(\pi(P,L+\mathrm{card}\,V\,l,g)-4\dim G\log g/q\big)}{r^{\mathrm{card}\,s}},\]
where \(\rho_T\) is the absolute value of the stamped Wilson expectation
to the power \(1/q\).

For the station at level \(k\), assume allocation, per-site
affordability at the width level, \(\mathrm{card}\,V\le M\),
\(l,L\le1\), and \(r=e^{2I}\) with \(I\ge\mathrm{work item}(O,k,M)\).
Then, for every nonempty \(s\), \(\rho_T(\Phi_s)\le e^{-I}\).
\emph{Sketch:} the exact vacuum identity supplies the Wilson partition
function as denominator. Apply III.11.1(a) to the representation with
its additive activities deleted. Comparing both numerator and
denominator costs \(e^{q\delta}\) twice, supplying the \(2\delta\)
already included in \(2\,\mathrm{exCost}\) in \(\pi\). Take the \(q\)-th
root. For the last assertion, allocation keeps the stage constant
between width and \(k\); \(L+\mathrm{card}\,V\,l\le1+M\), and
\(r^{\mathrm{card}\,s}\ge r\). \emph{Lean:}
\href{https://github.com/yang-mills-group/OpenAI-YM-20261028/blob/3456e5e75cf8d4cd558955a5c7638fba2d68f464/YM/PhysicalNext/AFTERNOON_INTERACTING_RANK_COMPARISON/Soft/Denarii.lean\#L49-L85}{\texttt{s\_postage}},
\href{https://github.com/yang-mills-group/OpenAI-YM-20261028/blob/3456e5e75cf8d4cd558955a5c7638fba2d68f464/YM/PhysicalNext/AFTERNOON_INTERACTING_RANK_COMPARISON/Soft/Denarii.lean\#L27-L43}{\texttt{denture}},
\href{https://github.com/yang-mills-group/OpenAI-YM-20261028/blob/3456e5e75cf8d4cd558955a5c7638fba2d68f464/YM/PhysicalNext/AFTERNOON_INTERACTING_RANK_COMPARISON/Soft/Denarii.lean\#L90-L123}{\texttt{coupon}}.
\emph{(proof read; the flow identity and the cell-stamp bookkeeping,
statement only)}

\paragraph{4.2 Free covariances}\label{free-covariances}

\textbf{Lemma 4.8 (III.11.2).} \emph{Sketch:} expand the pair term as
\(\sum_{a,b}\mathrm{Re}\,\Gamma_{ab}\,(r^{\Delta_L}\mathrm{ev}_L(r^{-\mathrm{gen}\,a}t_a))(r^{\Delta_S}\mathrm{ev}_S(r^{-\mathrm{gen}\,b}t_b))\).
Keys of generation below the scaling degree of the stencil evaluate to
\(0\). Keys of generation equal to the scaling degree give exactly the
free curvature kernel
(\href{https://github.com/yang-mills-group/OpenAI-YM-20261028/blob/3456e5e75cf8d4cd558955a5c7638fba2d68f464/YM/Next/RANK/StrideCost.lean\#L64-L67}{\texttt{deck}}).
Every other term carries an extra factor \(r^{-1}\) or smaller. So \(r\)
times the remainder is at most
\(D\sum|\mathrm{ev}_L(t_a)||\mathrm{ev}_S(t_b)|\), summed over the
non-leading pairs. \emph{Lean:}
\href{https://github.com/yang-mills-group/OpenAI-YM-20261028/blob/3456e5e75cf8d4cd558955a5c7638fba2d68f464/YM/Next/RANK/StrideCost.lean\#L70-L98}{\texttt{fence}}.
Its eventual form is the pair penalty
\href{https://github.com/yang-mills-group/OpenAI-YM-20261028/blob/3456e5e75cf8d4cd558955a5c7638fba2d68f464/YM/Next/RANK/EastBell.lean\#L44-L67}{\texttt{Prayer.penalty}}.

\textbf{Lemma 4.9 (measured Gaussian data).} Let \(Q\) be the symmetric
covariance array and \(m\) the mean vector obtained by applying the
normalized stencil and remainder functionals to measured pair and
one-point data. Let \(F\) be the free-current covariance on the stencil
keys, extended by zero on the remainder indices. Suppose \(B\ge1\),
\(T,P,\varepsilon\ge0\), and, entrywise,
\[|m_i|\le PT,\qquad |Q_{ij}-F_{ij}|\le PT^2.\] On the stencil keys
assume the sharper bounds \(|m_i|\le\varepsilon T\) and
\(|Q_{ij}-F_{ij}|\le\varepsilon T^2\). If
\[PT\le B,\qquad |F_{ij}|+PT^2\le B,\qquad \varepsilon T\le e,\qquad \varepsilon T^2\le e,\]
then all measured means and covariances are bounded by \(B\), and their
restriction to the stencil keys differs from \((0,F)\) by at most \(e\)
in both means and covariances. \emph{Sketch:} use the triangle
inequality for \(Q\). In the application, the measured horizons and
bridges satisfy the discrete Bianchi identities and scaled
finite-difference bounds. The stencil normalization \(T\) bounds the
sampling coefficients and the remainder factor \(M^{d-1}\) for addresses
within radius \(M\). These give the displayed broad and sharp bounds.
Positivity of \(Q\) is not required. \emph{Lean:}
\href{https://github.com/yang-mills-group/OpenAI-YM-20261028/blob/3456e5e75cf8d4cd558955a5c7638fba2d68f464/YM/Next/RANK/Prayer.lean\#L74-L111}{\texttt{resolution}}.

\textbf{Lemma 4.10 (III.11.3).} \emph{Sketch:} write \(\lambda=2\pi\)
and let \(T(\zeta)\) be exterior multiplication by \(i\zeta\) from
one-forms to two-forms. In the chart \(\zeta=\lambda\xi\), the heat
symbol is \[e^{-|\zeta|^2}|\zeta|^{-2}(T(\zeta)T(\zeta)^*)_{a,b},\] with
the color Kronecker factor. The regulator estimates in III.3 bound the
chart as a symbol of order \(0\) and its difference from this heat
symbol as a symbol of order \(1\). The chain rule under
\(\xi\mapsto\lambda\xi\) preserves these orders and changes the
constants by fixed powers of \(\lambda\). Sum the constants over the
finitely many plaquette pairs to obtain a bound independent of \(G\) and
the colors. \emph{Lean:}
\href{https://github.com/yang-mills-group/OpenAI-YM-20261028/blob/3456e5e75cf8d4cd558955a5c7638fba2d68f464/YM/PhysicalNext/RANK_COVARIANCE_ATLAS/Window.lean\#L118-L136}{\texttt{window\_prices}},
\href{https://github.com/yang-mills-group/OpenAI-YM-20261028/blob/3456e5e75cf8d4cd558955a5c7638fba2d68f464/YM/PhysicalNext/RANK_COVARIANCE_ATLAS/Window.lean\#L57-L76}{\texttt{heat\_window}}.

\textbf{Lemma 4.11 (covariance splitting and propagation).} (i) Let an
atlas be compatible with the chosen schedule and its retained unmasked
quadratic forms, and assume the atlas symmetry, coercivity, inverse and
crossing laws. Let \(i<N\), choose a tree forest for the step, and
assume the step torus satisfies the size condition for tree inversion.
Then
\[\mathcal C_{i+1}-W_i\mathcal C_iW_i^{\mathsf T}=d_f\Gamma_i d_f^{\mathsf T},\qquad\Gamma_i\succeq0,\]
where \(\mathcal C_i\) is the complete curvature covariance, \(W_i\) the
passage map, \(d_f\) the fine exterior derivative, and \(\Gamma_i\) the
uncut tree covariance, extended by zero on tree bonds. (ii) Fix kernel
choices and stage indices \(j,t\ge4\). There is \(R\ge1\) such that, for
every \(G\), every schedule with these kernel choices, and every
\(N,i,p,d\) satisfying \(J_{i+1}=j\), \(J_i=t\) and \(d-1\le j-2\)
(natural-number subtraction), the difference-motion propagation bound
holds with constant \(R\). \emph{Sketch:} (i) is a Schur-complement
identity on the curl image; the restricted tree inverse gives
positivity. For (ii), use the uniform resolvent estimate for the two
stage indices and sum its six plaquette components.

The table
\href{https://github.com/yang-mills-group/OpenAI-YM-20261028/blob/3456e5e75cf8d4cd558955a5c7638fba2d68f464/YM/PhysicalNext/RANK_COVARIANCE_ATLAS/Gates.lean\#L86}{\texttt{prepared}}
chooses, for each kernel, the uniform window, covariance and propagation
constants; their required smallness is a separate clearance condition.
\emph{Lean:}
\href{https://github.com/yang-mills-group/OpenAI-YM-20261028/blob/3456e5e75cf8d4cd558955a5c7638fba2d68f464/YM/PhysicalNext/RANK_COVARIANCE_ATLAS/Identity.lean\#L198-L213}{\texttt{one\_fold}},
\href{https://github.com/yang-mills-group/OpenAI-YM-20261028/blob/3456e5e75cf8d4cd558955a5c7638fba2d68f464/YM/PhysicalNext/RANK_COVARIANCE_ATLAS/Step.lean\#L106-L116}{\texttt{train\_issue}}.

\textbf{Lemma 4.12 (Wick expansion).} Let \(A\) be a finite index set,
let \(\mathrm{loc}:A\to S\) take values in an ordered location set, and
let \(Z_a\) be formal Gaussian generators with zero mean and symmetric
covariance \(Q\). Assume \(Q(Z_a,Z_b)=0\) whenever
\(\mathrm{loc}(a)=\mathrm{loc}(b)\). For every \(p\subseteq A\), the
formal Wick moment of \(\prod_{a\in p}Z_a\) is the sum, over perfect
matchings of \(p\) joining different locations, of the products of their
covariances. No positivity assumption is imposed on \(Q\).
\emph{Sketch:} remove one index and use the Gaussian moment recursion.
The zero mean removes the singleton term, and the covariance hypothesis
removes pairs at a common location. The remaining choices are in
bijection with the admissible matchings. \emph{Lean:}
\href{https://github.com/yang-mills-group/OpenAI-YM-20261028/blob/3456e5e75cf8d4cd558955a5c7638fba2d68f464/YM/Next/RANK/GaussTrack.lean\#L131-L173}{\texttt{headed}}.

\paragraph{4.3 Polynomial calculus and
identities}\label{polynomial-calculus-and-identities}

\textbf{Lemma 4.13 (polynomial and integral calculus).} Write
\(\|P\|_1\) for the sum of the absolute values of the coefficients of a
polynomial \(P\).

\begin{itemize}
\tightlist
\item
  \begin{enumerate}
  \def\labelenumi{(\roman{enumi})}
  \tightlist
  \item
    Fix derivative rules for \(G\) and an order \(d\). There is
    \(C\ge0\) such that, whenever a local observable \(A\) is regular
    for a parameter set of order at least \(d+1\) and has the prescribed
    derivative bounds \(c\), its pullback \(A\circ\mathcal N\) to a
    link-coordinate chart has a Taylor polynomial \(P\) of degree at
    most \(d\) satisfying
    \(\|P\|_1\le c(\mathrm{card}\,\text{links}+1)^dC\) and
    \(|A(\mathcal N(x))-P(x)|\le c\|x\|^{d+1}\) for \(\|x\|\) within the
    fixed chart radius.
    \href{https://github.com/yang-mills-group/OpenAI-YM-20261028/blob/3456e5e75cf8d4cd558955a5c7638fba2d68f464/YM/Next/RANK/LineBox.lean\#L168-L205}{\texttt{minted}}
  \end{enumerate}
\item
  (ii) For a family \(A_i\) satisfying (i) with bounds \(c_i\), and
  coefficients \(f_i\) such that \(S=\sum_i|f_i|c_i<\infty\), there is a
  degree-\(d\) polynomial \(P\) with
  \(\|P\|_1\le S(\mathrm{card}\,\text{links}+1)^dC\) and remainder at
  most \(S\|x\|^{d+1}\).
  \href{https://github.com/yang-mills-group/OpenAI-YM-20261028/blob/3456e5e75cf8d4cd558955a5c7638fba2d68f464/YM/Next/RANK/Stories.lean\#L99-L112}{\texttt{till}}
\item
  (iii) Fix real normed spaces \(E,V\), with \(V\) a normed algebra, and
  \(D\in\mathbb N\). There are \(C\ge0\) and \(q\in\mathbb N\) such
  that, for every finite list \(f_1,\dots,f_l:E\to V\), point \(x\), and
  \(B\ge1\), if each \(f_i\) is \(C^D\) at \(x\), \(f_i(x)=1\), and
  \(\|\partial^j f_i(x)\|\le B\) for \(j\le D\), then
  \(\|\partial^j\prod_i f_i(x)\|\le C((1+l)B)^q\) for \(j\le D\).
  \href{https://github.com/yang-mills-group/OpenAI-YM-20261028/blob/3456e5e75cf8d4cd558955a5c7638fba2d68f464/YM/Next/RANK/Terrain.lean\#L66-L123}{\texttt{conquer}}
\item
  (iv) For one- or two-insertion polynomial data in the staged
  comparison, assume the common-mask condition, the comparison bounds on
  means and covariances, \(0<\varepsilon\le1\), coefficient bound \(w\),
  degree bound \(D\), the prescribed variable supports, and vanishing of
  inadmissible pair terms. The difference between the resulting
  connected component and its free polynomial comparison is at most
  \(t+z+(u+t)t+tu\), where \(u,t,z\) are respectively the polynomial
  size, one-factor error and two-factor error bounds of that comparison.
  \href{https://github.com/yang-mills-group/OpenAI-YM-20261028/blob/3456e5e75cf8d4cd558955a5c7638fba2d68f464/YM/Next/RANK/Festoon.lean\#L57-L126}{\texttt{console}}
\item
  (v) For a bounded measurable gauge-invariant density \(g\), a
  branching step and an allocation satisfying its chart and
  integrability conditions, the pushforward is exactly the sum over
  matching-flag sectors and tree/tag choices of their weighted
  conditional integrals. If these integrands have a common factor
  \(\Lambda(U)\) and a finite local decomposition with the required
  integrability, the result is \(\Lambda(U)\) times the sum of products
  over pairwise disjoint compatible local components.
  \href{https://github.com/yang-mills-group/OpenAI-YM-20261028/blob/3456e5e75cf8d4cd558955a5c7638fba2d68f464/YM/PhysicalNext/AFTERNOON_INTERACTING_RANK_COMPARISON/Dispenser.lean\#L121-L143}{\texttt{arriving}},
  \href{https://github.com/yang-mills-group/OpenAI-YM-20261028/blob/3456e5e75cf8d4cd558955a5c7638fba2d68f464/YM/PhysicalNext/AFTERNOON_INTERACTING_RANK_COMPARISON/Rendered.lean\#L36-L86}{\texttt{rendered}}
\item
  (vi) Fix a valid coarse configuration and a tree/tag choice. If
  finitely many integrable observables satisfy the locality and
  compatibility conditions of the conditional fluctuation law, and their
  coarse supports are pairwise disjoint, their product is integrable and
  its conditional integral equals the product of their conditional
  integrals.
  \href{https://github.com/yang-mills-group/OpenAI-YM-20261028/blob/3456e5e75cf8d4cd558955a5c7638fba2d68f464/YM/PhysicalNext/AFTERNOON_INTERACTING_RANK_COMPARISON/Sight.lean\#L132-L158}{\texttt{peer\_beams}}
\end{itemize}

\emph{Sketch:} Taylor's theorem gives (i), and absolute coefficient
summability gives (ii). For (iii), split each factor as \(1+(f_i-1)\)
and induct on the derivative order; the highest derivative terms add
rather than multiply. For (iv), the two-point cumulant is the second
moment minus the product of the means; expand the difference and use the
three stated bounds. For (v), apply the allocation identity before
discarding sectors with the wrong flags, then group the local
components. For (vi), the joint local-coordinate marginal is a product
measure, so Fubini gives factorization. \emph{(proof read; local
comparison and allocation inputs statement only)}

\paragraph{4.4 From frequent good scales to trial systems
(III.11.4)}\label{from-frequent-good-scales-to-trial-systems-iii.11.4}

\textbf{Lemma 4.14 (multi-step propagation).} Fix a level \(k\le N\), a
band \(b\), a fitted configuration \(x\) with \(v\) insertions, and an
exact RG history. Let \(m_0\) and \(m_1\) be its width and middle
levels. Assume allocation at \(k\), admissible representations along
\([m_0,m_1]\), a numerical propagation bound for arity \(M\), and
\[\hat\beta_u\in[\hat\beta_k/2,\hat\beta_k]\quad(m_0\le u\le m_1),\qquad v\le M.\]
Assume the instance decomposition for the tested labels: at the middle
level, the marked sources approximate the negative free-profile terms
within the prescribed error, with the required background and source
bounds, and the subsequent steps satisfy the exhibited bounds. Finally
assume each free-profile term has absolute value at most
\(\exp((\log\log\hat\beta_k)^2/4)\).

Then the flow of the corrected stamped observable has a local activity
decomposition at \(m_0\) with radius \(\mathrm{allowance}(M,k)\) and
exponents \((e,1,1)\), where \(e\) is the input source exponent.
\emph{Sketch:} use the middle-level decomposition and the free-profile
bound to initialize the source estimate. The propagation bound controls
its iteration through \([m_0,m_1]\); the coupling window keeps every
intermediate estimate in the same numerical range. The correction
cancels the leading free-profile terms. This is the local decomposition
used in Lemma 4.7. \emph{Lean:}
\href{https://github.com/yang-mills-group/OpenAI-YM-20261028/blob/3456e5e75cf8d4cd558955a5c7638fba2d68f464/YM/PhysicalNext/AFTERNOON_INTERACTING_RANK_COMPARISON/Soft/Supra.lean\#L29-L53}{\texttt{v\_conditioned}}.
\emph{(proof read; helper estimates statement only)}

\textbf{Lemma 4.15 (the chain).} Under the hypotheses of III.11.4:

\begin{enumerate}
\def\labelenumi{\arabic{enumi}.}
\tightlist
\item
  At infinitely many \(k\), the simultaneous price and insertion
  conditions give covariance data, an exact RG history and compatible
  restarted representations, eventually in \(N\) and for every \(p\).
  The uniform covariance and test-point estimates supply the local
  Gaussian comparison.
  (\href{https://github.com/yang-mills-group/OpenAI-YM-20261028/blob/3456e5e75cf8d4cd558955a5c7638fba2d68f464/YM/PhysicalNext/AFTERNOON_INTERACTING_RANK_COMPARISON/Soft/Star.lean\#L64-L66}{\texttt{cast\_off}}
  via
  \href{https://github.com/yang-mills-group/OpenAI-YM-20261028/blob/3456e5e75cf8d4cd558955a5c7638fba2d68f464/YM/PhysicalNext/AFTERNOON_INTERACTING_RANK_COMPARISON/Soft/Star.lean\#L49-L63}{\texttt{backstage}}.)
\item
  Construct a common correction map and polynomial approximations for
  all rows. The one-step comparison supplies the coefficient, support
  and remainder bounds on the same path of representations.
  (\href{https://github.com/yang-mills-group/OpenAI-YM-20261028/blob/3456e5e75cf8d4cd558955a5c7638fba2d68f464/YM/PhysicalNext/AFTERNOON_INTERACTING_RANK_COMPARISON/Soft/Diana.lean\#L57-L104}{\texttt{v\_paid\_factory}},
  \emph{proof read; analytic inputs statement only}.)
\item
  Compare the resulting Gaussian polynomial moments with the free
  profiles. The eventual finite-difference bounds transfer this
  comparison from the chosen representatives to the actual corrected
  columns, giving the instance estimates of Lemma 4.14.
  (\href{https://github.com/yang-mills-group/OpenAI-YM-20261028/blob/3456e5e75cf8d4cd558955a5c7638fba2d68f464/YM/PhysicalNext/AFTERNOON_INTERACTING_RANK_COMPARISON/Soft/Supra.lean\#L84-L95}{\texttt{v\_culminates}},
  \emph{proof read}.)
\item
  Apply Lemma 4.14 to propagate the corrected estimates to the width
  level, retaining the terminal density identity and admissibility
  needed for the root estimate.
  (\href{https://github.com/yang-mills-group/OpenAI-YM-20261028/blob/3456e5e75cf8d4cd558955a5c7638fba2d68f464/YM/PhysicalNext/AFTERNOON_INTERACTING_RANK_COMPARISON/Soft/Supra.lean\#L112-L125}{\texttt{v\_pay\_all}},
  \emph{proof read}.)
\item
  Offers give a trial system. Choose depths \(k_i\) inductively among
  the good levels, with \(k_{i+1}>k_i\), strictly increasing widths, and
  strictly decreasing test scales tending to \(0\). The lifts come from
  \href{https://github.com/yang-mills-group/OpenAI-YM-20261028/blob/3456e5e75cf8d4cd558955a5c7638fba2d68f464/YM/PhysicalNext/AFTERNOON_INTERACTING_RANK_COMPARISON/Soft/Supra.lean\#L149-L173}{\texttt{ticketed}}
  (bounded, real, with first entry the reference unit). The chessboard
  bounds come from Lemma 4.7, using the per-site affordability
  \href{https://github.com/yang-mills-group/OpenAI-YM-20261028/blob/3456e5e75cf8d4cd558955a5c7638fba2d68f464/YM/Next/RANK/Allowance.lean\#L42}{\texttt{banking}},
  which follows from Lemma 4.6.
  (\href{https://github.com/yang-mills-group/OpenAI-YM-20261028/blob/3456e5e75cf8d4cd558955a5c7638fba2d68f464/YM/PhysicalNext/AFTERNOON_INTERACTING_RANK_COMPARISON/Soft/Supra.lean\#L177-L222}{\texttt{v\_obtain\_purchases}},
  \emph{proof read}.)
\end{enumerate}

\emph{Lean:}
\href{https://github.com/yang-mills-group/OpenAI-YM-20261028/blob/3456e5e75cf8d4cd558955a5c7638fba2d68f464/YM/PhysicalNext/AFTERNOON_INTERACTING_RANK_COMPARISON/Soft/Star.lean\#L68-L69}{\texttt{wheeled}}
composes items 2--5.

\textbf{Lemma 4.16 (eventual insertion estimates).} Assume \(G\) is
simple, \(C\) satisfies the construction record, and the one-step cost
package uses the prescribed weight sequence and multiplicities at least
\(1\). Fix a degree-\(d\) column family \(X\) and stage package whose
mark exponent satisfies \(c\ge2(w(G,d)+1)\), where \(w(G,d)=w_d\) is the
prescribed coefficient exponent. Then, for all sufficiently large \(k\),
metering at \(k\) implies the lantern-mark condition for \(X\) at \(k\).
\emph{Sketch:} write \(x=\hat\beta_k\) and
\(b_a(x)=\exp(a(\log\log x)^2)\). For every fixed \(a\), the inequality
\(b_a(x)s\le b_{-8}(x)\) with \(s\ge0\) gives
\[s\le\exp(-(8+a)(\log\log x)^2)\le x\] for sufficiently large \(x\).
Hence \(x^{-c/2}s\le x^{-w(G,d)}\). The one-step estimate and spacing of
the rounded test points turn this coefficient bound into the required
insertion estimate. The threshold is uniform over the tested
coefficients, cutoffs and periods. \emph{Lean:}
\href{https://github.com/yang-mills-group/OpenAI-YM-20261028/blob/3456e5e75cf8d4cd558955a5c7638fba2d68f464/YM/MergeProof/SteamW.lean\#L54-L59}{\texttt{bursts}},
\href{https://github.com/yang-mills-group/OpenAI-YM-20261028/blob/3456e5e75cf8d4cd558955a5c7638fba2d68f464/YM/MergeProof/SteamW.lean\#L45-L52}{\texttt{quaking}}.
\emph{(proof read; analytic inputs statement only)}

\textbf{Lemma 4.17 (discharge of the frequent-scale conditions).} Given
compatible pre-\(L\) price data and one-step bounds for a sheet, with
the stage thresholds below the chosen block lengths and each degree
index eventually present in the stage index sets, a Haar budget,
nonnegative mark exponents and simple \(G\), there are panels \(Q_d\)
and a waiting window extending any given one with the following
property. For every Wilson family \(C\) paid on that window and every
binding with these prices, for every \(d\), the four price conditions
hold at infinitely many levels. \emph{Sketch:} the restarted comparison
and the scale-selection construction produce the panels and the window;
frequency follows because each requested degree index is eventually
present in those stage sets. \emph{Lean:}
\href{https://github.com/yang-mills-group/OpenAI-YM-20261028/blob/3456e5e75cf8d4cd558955a5c7638fba2d68f464/YM/MergeProof/Snow.lean\#L55-L60}{\texttt{franchise}}.
\emph{(proof structure read; helpers statement only)}

\paragraph{4.5 The rank comparison
(III.11.5)}\label{the-rank-comparison-iii.11.5}

\textbf{Lemma 4.18 (equation count).} Under the projector presentation
of III.11.5(i), at every cutoff,
\(\dim\mathcal L_d\le\dim\ker\Pi_{d,N,p}+\dim F_d^{\rm free}\).
\emph{Sketch:} \(\ker\Pi_{d,N,p}\) is the image of \(\mathcal E_d\)
under an injective corrected label map, by the projector presentation.
The free null space is the image of \(\mathcal E_d\) under the label
coordinates. This uses
\href{https://github.com/yang-mills-group/OpenAI-YM-20261028/blob/3456e5e75cf8d4cd558955a5c7638fba2d68f464/YM/Euclid/Comparator.lean\#L102-L118}{\texttt{projectionExact}}
(a label is null in the free table if and only if its on-shell image
lies in the residual) and
\href{https://github.com/yang-mills-group/OpenAI-YM-20261028/blob/3456e5e75cf8d4cd558955a5c7638fba2d68f464/YM/Extra/FREE_DETECTOR_ALGEBRA/Rank.lean\#L38-L54}{\texttt{no\_residual}}
(the residual is \(0\): free probes separate the on-shell sector). Then
\(\dim F_d^{\rm free}+\dim\mathcal E_d=\dim\mathcal L_d\). \emph{Lean:}
\href{https://github.com/yang-mills-group/OpenAI-YM-20261028/blob/3456e5e75cf8d4cd558955a5c7638fba2d68f464/YM/MergeProof/Powder.lean\#L19-L47}{\texttt{savings}},
\href{https://github.com/yang-mills-group/OpenAI-YM-20261028/blob/3456e5e75cf8d4cd558955a5c7638fba2d68f464/YM/Next/RANK/Monitors.lean\#L92-L94}{\texttt{EquationCount}},
\href{https://github.com/yang-mills-group/OpenAI-YM-20261028/blob/3456e5e75cf8d4cd558955a5c7638fba2d68f464/YM/Euclid/Comparator.lean\#L102-L118}{\texttt{projectionExact}},
\href{https://github.com/yang-mills-group/OpenAI-YM-20261028/blob/3456e5e75cf8d4cd558955a5c7638fba2d68f464/YM/Extra/FREE_DETECTOR_ALGEBRA/Rank.lean\#L38-L54}{\texttt{no\_residual}}.
\emph{(proof read; free positivity input statement only)}

\textbf{Lemma 4.19 (short-distance comparison and independence).} Assume
\(C\) satisfies the construction record, \(f\) gives generic point
limits \(S\), and the correlation kernels of \(S\) are continuous on
generic configurations and integrable against every separated test
function. Let \(Z\) satisfy the purchase conditions, and suppose all its
coefficient vectors have iterated limits along \(f\). Then:

\begin{itemize}
\tightlist
\item
  \begin{enumerate}
  \def\labelenumi{(\roman{enumi})}
  \tightlist
  \item
    for every row and every basis column, the correlation of its
    limiting corrected vectors at the trial's physical points differs
    from the free correlation at the rescaled points by at most
    \(P_k\gamma_k\), where \(P_k\) is the polynomial coefficient weight
    of the free profile. The numbers \(P_k\) converge and
    \(\gamma_k\to0\), so these correlations converge to their prescribed
    free test values;
  \end{enumerate}
\item
  (ii) for all sufficiently large \(k\), the limiting vectors of trial
  \(k\) are linearly independent modulo the subspace of labels
  annihilated by every integrated correlation test of \(S\).
\end{itemize}

\emph{Sketch:} expand the ordinary moment using the relative connected
components and apply the chessboard inequality to every nonempty
component. This gives the error \(P_k\gamma_k\) before the cutoff
limits, and the bound passes through those limits. Each row functional
descends to the quotient by the integrated-test null space. The
resulting finite matrix converges to the injective free detector matrix.
Injectivity is open in finite dimensions, giving (ii). \emph{Lean:}
\href{https://github.com/yang-mills-group/OpenAI-YM-20261028/blob/3456e5e75cf8d4cd558955a5c7638fba2d68f464/YM/Next/RANK/Monitors.lean\#L71-L87}{\texttt{measured}};
the comparison is
\href{https://github.com/yang-mills-group/OpenAI-YM-20261028/blob/3456e5e75cf8d4cd558955a5c7638fba2d68f464/YM/Next/RANK/Bought.lean\#L112-L123}{\texttt{Purchases.samples}}/\href{https://github.com/yang-mills-group/OpenAI-YM-20261028/blob/3456e5e75cf8d4cd558955a5c7638fba2d68f464/YM/Next/RANK/Bought.lean\#L125-L133}{\texttt{converge}}
in
\href{https://github.com/yang-mills-group/OpenAI-YM-20261028/blob/3456e5e75cf8d4cd558955a5c7638fba2d68f464/YM/Next/RANK/Bought.lean}{\texttt{Bought.lean}}.
\emph{(proof read; chessboard input statement only)}

\textbf{Lemma 4.20 (seed and null space).} Assume \(C\) satisfies the
construction record. Under a centred port for \(C\) along \(f\), raw
limits \(S\) along that same extraction, relay bins with convergent
projectors, the equation count for every \(d\), and trial systems with
the purchase conditions for every \(d\), there is a refinement along
which the following hold:

\begin{itemize}
\tightlist
\item
  the arrival and spans records;
\item
  a cleaned smooth table \(T\);
\item
  the seed record for \(T\);
\item
  for every \(d\): \(\mathrm{null}_d(T)=\ker\Pi_\infty\), where
  \(\Pi_\infty\) is the limit projector, and
  \(\dim F_d(T)=\dim F_d^{\rm free}\).
\end{itemize}

\emph{Sketch:} the relay estimates give
\(\ker\Pi_\infty\subseteq\mathrm{null}_d(T)\). Lemma 4.19(ii) gives the
reverse dimension inequality, using the equation count and eventual
stability of the projector kernel dimension. Hence the subspaces
coincide, and the field dimensions agree. \emph{Lean:}
\href{https://github.com/yang-mills-group/OpenAI-YM-20261028/blob/3456e5e75cf8d4cd558955a5c7638fba2d68f464/YM/Next/RANK/Beam.lean\#L153-L180}{\texttt{arriving\_seed}},
\href{https://github.com/yang-mills-group/OpenAI-YM-20261028/blob/3456e5e75cf8d4cd558955a5c7638fba2d68f464/YM/Next/RANK/Beam.lean\#L142-L147}{\texttt{equation\_retains}}.
\emph{(proof read; analytic inputs statement only)}

\textbf{Lemma 4.21 (equivariant identification).} Under the projector
presentation of III.11.5(i) and the extraction record for its projectors
and cleaned table \(T\), there are the real,
signed-permutation-equivariant isomorphisms \(\Phi_d\) of III.11.5. The
free table has its canonical real structure and zero residual null
space. \emph{Sketch:} let \(A_i\to\Pi_\infty\) be the projectors along
the extraction, and set
\(Y_i=\Pi_\infty A_i+(1-\Pi_\infty)(1-A_i)\to1\). For large \(i\),
\(Y_i\) is invertible and \(Y_iA_i=\Pi_\infty Y_i\). Since \(A_i\) and
\(\Pi_\infty\) are real and equivariant, so is \(Y_i\). Then
\(L=Y_i\circ(\text{corrected label map})\) maps \(\mathcal E_d\) onto
\(\ker\Pi_\infty=\mathrm{null}_d(T)\)
(\href{https://github.com/yang-mills-group/OpenAI-YM-20261028/blob/3456e5e75cf8d4cd558955a5c7638fba2d68f464/YM/MergeProof/Thrush.lean\#L13-L68}{\texttt{drumroll}}).
The inverse of the induced quotient map, followed by the fixed free
coordinate identification, is \(\Phi_d\). \emph{Lean:}
\href{https://github.com/yang-mills-group/OpenAI-YM-20261028/blob/3456e5e75cf8d4cd558955a5c7638fba2d68f464/YM/MergeProof/Thrush.lean\#L72-L114}{\texttt{thrush}},
\href{https://github.com/yang-mills-group/OpenAI-YM-20261028/blob/3456e5e75cf8d4cd558955a5c7638fba2d68f464/YM/MergeProof/Thrush.lean\#L13-L68}{\texttt{drumroll}}.

\textbf{Lemma 4.22 (III.11.5).} Under the hypotheses of III.11.5, its
conclusion follows by the following steps:

\begin{enumerate}
\def\labelenumi{\arabic{enumi}.}
\tightlist
\item
  extract raw point limits (IV.1) and include the label projections;
\item
  build, for each \(d\), the free detector columns \(X_d\), which form a
  basis of \(F_d^{\rm free}\);
\item
  combine assumption (iii) with Lemma 4.16 and apply III.11.4 to get
  trial systems with the purchase conditions;
\item
  apply Lemma 4.20 with Lemma 4.18;
\item
  assemble the extraction record (using D05);
\item
  conclude with Lemma 4.21.
\end{enumerate}

\emph{Lean:}
\href{https://github.com/yang-mills-group/OpenAI-YM-20261028/blob/3456e5e75cf8d4cd558955a5c7638fba2d68f464/YM/MergeProof/Powder.lean\#L48-L71}{\texttt{rated}}.

\paragraph{4.6 The stress tensor (III.11.6,
III.11.7)}\label{the-stress-tensor-iii.11.6-iii.11.7}

\textbf{Lemma 4.23 (annihilated flux combinations).} Under the
hypotheses of III.11.7, let \(I\) be a finite family of generic tester
configurations on the schedule grids, each with \(x_0=0\), and let
\((g_a)\) be coefficients with \(\sum_ag_a\,\mathrm{test}_a(z)=0\) for
all \(z\in F_D\). Here \(\mathrm{test}_a\) inserts \(z\) at the root
with the remaining fields of tester \(a\); \(\mathrm{flux}_a\) inserts
the flux built from \((\mathcal J,U,v)\) there. Then
\(\sum_ag_a\,\mathrm{flux}_a=0\). \emph{Sketch:} let \(r_l\to0\) be the
radii of the shell budget. Choose a smooth radial bump \(\chi\), equal
to \(1\) on the ball of radius \(1/2\) and supported in the unit ball.
Conservation makes the flux independent of the window, so for large
\(l\) use \(\chi_l(y)=\chi(y/r_l)\). By homogeneity of \(U\), the flux
becomes
\(-\sum_{ij}\int M_{ij}(y)\,r_l^{3+p}\langle\mathcal J_{ij}(r_ly)\cdots\rangle\,dy\),
where \(M_{ij}=\partial_i\chi\cdot U_j\) is integrable and supported in
\(\tfrac12\le|y|\le1\). The annulus bound
\href{https://github.com/yang-mills-group/OpenAI-YM-20261028/blob/3456e5e75cf8d4cd558955a5c7638fba2d68f464/YM/MergeProof/West/Ozone.lean\#L21-L45}{\texttt{zone}}
says that, for annihilating combinations,
\(\sup_{|y|\in(r_l/4,2r_l)}\big|\sum_ag_a\,r_l^{3+p}\langle\mathcal J(y)\cdots\rangle_a\big|\to0\).
So the combination is a constant sequence tending to \(0\). \emph{Lean:}
\href{https://github.com/yang-mills-group/OpenAI-YM-20261028/blob/3456e5e75cf8d4cd558955a5c7638fba2d68f464/YM/MergeProof/West/Blanca.lean\#L49-L120}{\texttt{tenet}},
\href{https://github.com/yang-mills-group/OpenAI-YM-20261028/blob/3456e5e75cf8d4cd558955a5c7638fba2d68f464/YM/MergeProof/West/Ozone.lean\#L21-L45}{\texttt{zone}}.
\emph{(proof read; annulus comparison inputs statement only)}

\textbf{Lemma 4.24 (III.11.7).} \emph{Sketch:}

\begin{itemize}
\tightlist
\item
  for each finite set of tester configurations, Lemma 4.23 and
  finite-dimensional duality
  (\href{https://github.com/yang-mills-group/OpenAI-YM-20261028/blob/3456e5e75cf8d4cd558955a5c7638fba2d68f464/YM/MergeProof/West/Bravo.lean\#L31-L52}{\texttt{obtainable}})
  give \(B\) matching all their fluxes;
\item
  fix a finite separating family of testers (A04), choose \(B\) matching
  it, and compare with the field matching its union with any additional
  tester; separation forces the two fields to agree;
\item
  translation covariance of flux and correlations moves any grid
  configuration to \(x_0=0\);
\item
  continuity of the flux and of the limit kernels, plus density of grid
  configurations, extends the identity to all distinct points.
\end{itemize}

\emph{Lean:}
\href{https://github.com/yang-mills-group/OpenAI-YM-20261028/blob/3456e5e75cf8d4cd558955a5c7638fba2d68f464/YM/MergeProof/West/Bravo.lean\#L55-L89}{\texttt{annular\_fieldA}},
\href{https://github.com/yang-mills-group/OpenAI-YM-20261028/blob/3456e5e75cf8d4cd558955a5c7638fba2d68f464/YM/MergeProof/West/Bravo.lean\#L31-L52}{\texttt{obtainable}}.

\textbf{Lemma 4.25 (stress laws).} Under the hypotheses of III.11.7,
there are linear translation maps \(F_d(S)\to F_{d+1}(S)\) and rotation
maps \(F_d(S)\to F_d(S)\) satisfying the respective flux identities. The
rotation maps also depend real-linearly on the skew matrix.
\emph{Sketch:} apply Lemma 4.24 with constant fields \(U=e_i\) (\(p=0\),
\(D=d+1\)) and with the rotation fields \(U(y)=ky\), \(k\) skew
(\(p=1\), \(D=d\)). Linearity follows from linearity of the flux and
separation of fields by correlations. \emph{Lean:}
\href{https://github.com/yang-mills-group/OpenAI-YM-20261028/blob/3456e5e75cf8d4cd558955a5c7638fba2d68f464/YM/MergeProof/West/Orange.lean\#L33-L63}{\texttt{shells\_law}}.

\textbf{Lemma 4.26 (III.11.6).} Under the respective hypotheses of
III.11.6(a) and (b), the stated stress and ultraviolet records exist.
\emph{Sketch of (a):}

\begin{enumerate}
\def\labelenumi{\arabic{enumi}.}
\tightlist
\item
  D06 gives the orientation law.
\item
  The station equivalences identify the field spaces of \(S\) with those
  of the free model; this uses the hypothesis
  \href{https://github.com/yang-mills-group/OpenAI-YM-20261028/blob/3456e5e75cf8d4cd558955a5c7638fba2d68f464/YM/Next/UV_STRESS/Invert.lean\#L108}{\texttt{DimAbove}}.
\item
  \href{https://github.com/yang-mills-group/OpenAI-YM-20261028/blob/3456e5e75cf8d4cd558955a5c7638fba2d68f464/YM/Next/UV_STRESS/Limit.lean\#L129-L167}{\texttt{produce\_current}}
  (VI.2) gives \((\mathcal J,X)\) with \(\mathcal J\ne0\), the approach,
  stability, the meter and zero current divergence.
\item
  \href{https://github.com/yang-mills-group/OpenAI-YM-20261028/blob/3456e5e75cf8d4cd558955a5c7638fba2d68f464/YM/Next/UV_STRESS/Tensor.lean\#L121-L125}{\texttt{linear\_cons}}
  converts zero current divergence to conservation. The weak comparison
  along the scaling sequence supplies the action law.
\item
  Lemma 4.25 gives the stress laws.
\end{enumerate}

\emph{Sketch of (b):} take free Maxwell data for simple \(G\) and
replace them by the canonical data
(\href{https://github.com/yang-mills-group/OpenAI-YM-20261028/blob/3456e5e75cf8d4cd558955a5c7638fba2d68f464/YM/UVNext/FREE_STRESS_CERTIFICATE/Choose.lean\#L137-L150}{\texttt{FreeStress.fire}}).
Build the station hypothesis from the seed and weak completeness, and
the hypothesis
\href{https://github.com/yang-mills-group/OpenAI-YM-20261028/blob/3456e5e75cf8d4cd558955a5c7638fba2d68f464/YM/Next/UV_STRESS/Invert.lean\#L108}{\texttt{DimAbove}}
from the corrected rank widths. Then apply (a). \emph{Lean:}
\href{https://github.com/yang-mills-group/OpenAI-YM-20261028/blob/3456e5e75cf8d4cd558955a5c7638fba2d68f464/YM/MergeProof/Emanate.lean\#L15-L30}{\texttt{readyA}},
\href{https://github.com/yang-mills-group/OpenAI-YM-20261028/blob/3456e5e75cf8d4cd558955a5c7638fba2d68f464/YM/MergeProof/Emanate.lean\#L31-L37}{\texttt{emanate}}.
\emph{(proof read; helper estimates statement only)}

\subsubsection{5. Where the difficulty
lies}\label{where-the-difficulty-lies-13}

\begin{itemize}
\tightlist
\item
  \textbf{The volume-uniform root test.} A per-site loss of
  \(\log\beta\) in the chessboard root would destroy the comparison,
  because the budget must tend to \(0\) while \(\beta\to\infty\). The
  tree-gauge count of III.11.1 avoids this: it costs only
  \(\beta^{2\dim G}\) in total, not per site. The price for this is a
  crude, but uniform, per-site constant. This estimate uses Haar
  integration rather than a Hessian determinant.
\item
  \textbf{The direction of the rank argument.} The upper bound (equation
  count) is algebraic and holds at every cutoff. The lower bound needs
  the short-distance comparison, at infinitely many scales only, through
  trial sequences. The equivariant isomorphism comes from conjugating
  projectors, not from an explicit basis. The conclusion identifies
  complex representations with compatible conjugations under signed
  permutations.
\item
  \textbf{Frequent rather than eventual scales.} Every price condition
  is required only at infinitely many levels; a diagonal extraction then
  selects the trials. This is weaker than an ``eventually in \(k\)''
  statement, and it is all the rank argument needs.
\item
  \textbf{Uniformity.} The constants must be chosen before the block
  lengths and must control every period. Eventual estimates cannot
  replace bounds on a specified comparison window without the
  corresponding delay.
\item
  \textbf{Connections.} The stress-flux argument supplies the local
  charge fields of VI.1. The equation-ideal calculation uses the free
  label algebra of II.1.
\end{itemize}

\subsubsection{6. Notation}\label{notation-13}

{\def\LTcaptype{none} % do not increment counter
\begin{longtable}[]{@{}
  >{\raggedright\arraybackslash}p{(\linewidth - 4\tabcolsep) * \real{0.3333}}
  >{\raggedright\arraybackslash}p{(\linewidth - 4\tabcolsep) * \real{0.3333}}
  >{\raggedright\arraybackslash}p{(\linewidth - 4\tabcolsep) * \real{0.3333}}@{}}
\toprule\noalign{}
\begin{minipage}[b]{\linewidth}\raggedright
Written proof
\end{minipage} & \begin{minipage}[b]{\linewidth}\raggedright
Lean (glossary)
\end{minipage} & \begin{minipage}[b]{\linewidth}\raggedright
Meaning
\end{minipage} \\
\midrule\noalign{}
\endhead
\bottomrule\noalign{}
\endlastfoot
schedule factor \(\mu_k\) &
\href{https://github.com/yang-mills-group/OpenAI-YM-20261028/blob/3456e5e75cf8d4cd558955a5c7638fba2d68f464/YM/RG/History.lean\#L81}{\texttt{monitor}}
& \(\exp((\log\log\hat\beta_k)^2)\) \\
coefficient exponent \(w_d\) &
\href{https://github.com/yang-mills-group/OpenAI-YM-20261028/blob/3456e5e75cf8d4cd558955a5c7638fba2d68f464/YM/MergeProof/Tender.lean\#L11}{\texttt{winter}}
& degree of the input coefficient family plus \(4\) \\
Haar budget \((s,b)\) &
\href{https://github.com/yang-mills-group/OpenAI-YM-20261028/blob/3456e5e75cf8d4cd558955a5c7638fba2d68f464/YM/Next/RANK/Determinant.lean\#L26-L33}{\texttt{HaarBudget}}
(\texttt{small}, \texttt{big}) & lower constant for shells, upper
constant for \(\int e^{-\mathrm{chord}/t}\) \\
\(\sigma_R\), \(\mathrm{ent}\) &
\href{https://github.com/yang-mills-group/OpenAI-YM-20261028/blob/3456e5e75cf8d4cd558955a5c7638fba2d68f464/YM/Next/RANK/Determinant.lean\#L116}{\texttt{strength}},
\href{https://github.com/yang-mills-group/OpenAI-YM-20261028/blob/3456e5e75cf8d4cd558955a5c7638fba2d68f464/YM/Next/RANK/Determinant.lean\#L220}{\texttt{ent}}
& \(\min(1,\text{linear clip})\); \(\log b-\dim G\log\sigma\) \\
per-site price \(\pi\), work item &
\href{https://github.com/yang-mills-group/OpenAI-YM-20261028/blob/3456e5e75cf8d4cd558955a5c7638fba2d68f464/YM/Next/RANK/Determinant.lean\#L225-L227}{\texttt{protection}},
\href{https://github.com/yang-mills-group/OpenAI-YM-20261028/blob/3456e5e75cf8d4cd558955a5c7638fba2d68f464/YM/Next/RANK/StationMeter.lean\#L41-L44}{\texttt{work\ item}}
& III.11.1 \\
admissible at level \(m\) &
\href{https://github.com/yang-mills-group/OpenAI-YM-20261028/blob/3456e5e75cf8d4cd558955a5c7638fba2d68f464/YM/Next/RANK/StationMeter.lean\#L28-L36}{\texttt{Meets}}
& conditions on \((P,g,\Lambda,\delta)\) in III.11.1(b) \\
chessboard root \(\rho_T\) & \texttt{Cells.Tiling.rootTest} &
\(\lvert E\prod_{\rm cells}\text{copy}\rvert^{1/\mathrm{card\,cells}}\) \\
free curvature kernel & \texttt{light} & derivative of the free Maxwell
curvature two-point kernel between colored stencils \\
scale-\(r\) pair term & \texttt{collect}, \texttt{arc}, \texttt{arbor},
\texttt{Beamer.amount} &
\(\sum\Gamma_{ab}\,r^{-\mathrm{gen}}t_a\otimes r^{-\mathrm{gen}}t_b\);
\(r^{2+\mathrm{ord}}\) \\
Fourier chart rate &
\href{https://github.com/yang-mills-group/OpenAI-YM-20261028/blob/3456e5e75cf8d4cd558955a5c7638fba2d68f464/YM/PhysicalNext/RANK_COVARIANCE_ATLAS/Fourier.lean\#L26}{\texttt{circleRate}}
& \(2\pi\); heat parameter \((2\pi)^2\) \\
window estimate &
\href{https://github.com/yang-mills-group/OpenAI-YM-20261028/blob/3456e5e75cf8d4cd558955a5c7638fba2d68f464/YM/Next/RANK/Radian/Signed.lean\#L47-L53}{\texttt{Window}},
\href{https://github.com/yang-mills-group/OpenAI-YM-20261028/blob/3456e5e75cf8d4cd558955a5c7638fba2d68f464/YM/Next/RANK/Radian/Means.lean\#L16-L19}{\texttt{Cast}}
& symbol bounds of III.11.3 \\
stage package, metering &
\href{https://github.com/yang-mills-group/OpenAI-YM-20261028/blob/3456e5e75cf8d4cd558955a5c7638fba2d68f464/YM/Next/RANK/Concord.lean\#L38-L51}{\texttt{Cohesion}},
\href{https://github.com/yang-mills-group/OpenAI-YM-20261028/blob/3456e5e75cf8d4cd558955a5c7638fba2d68f464/YM/Next/RANK/Concord.lean\#L62-L64}{\texttt{Metering}}
& per-degree stage data; charge and launch at \(k\) \\
price conditions &
\href{https://github.com/yang-mills-group/OpenAI-YM-20261028/blob/3456e5e75cf8d4cd558955a5c7638fba2d68f464/YM/PhysicalNext/RANK_COVARIANCE_ATLAS/Consumer.lean\#L71}{\texttt{Tender}},
\texttt{Metering},
\href{https://github.com/yang-mills-group/OpenAI-YM-20261028/blob/3456e5e75cf8d4cd558955a5c7638fba2d68f464/YM/PhysicalNext/AFTERNOON_INTERACTING_RANK_COMPARISON/Pleas.lean\#L21-L26}{\texttt{ShowHours}},
\href{https://github.com/yang-mills-group/OpenAI-YM-20261028/blob/3456e5e75cf8d4cd558955a5c7638fba2d68f464/YM/PhysicalNext/AFTERNOON_INTERACTING_RANK_COMPARISON/Soft/Postsale.lean\#L19-L21}{\texttt{Vents}},
\href{https://github.com/yang-mills-group/OpenAI-YM-20261028/blob/3456e5e75cf8d4cd558955a5c7638fba2d68f464/YM/PhysicalNext/AFTERNOON_INTERACTING_RANK_COMPARISON/Proceed.lean\#L74-L77}{\texttt{LanternMarks}}
& hypotheses of III.11.4 at a level \(k\) \\
provisions, aggregate, turning, offers &
\href{https://github.com/yang-mills-group/OpenAI-YM-20261028/blob/3456e5e75cf8d4cd558955a5c7638fba2d68f464/YM/PhysicalNext/AFTERNOON_INTERACTING_RANK_COMPARISON/Soft/Diana.lean\#L37}{\texttt{s\_provisions}},
\href{https://github.com/yang-mills-group/OpenAI-YM-20261028/blob/3456e5e75cf8d4cd558955a5c7638fba2d68f464/YM/PhysicalNext/AFTERNOON_INTERACTING_RANK_COMPARISON/Soft/Diana.lean\#L44-L51}{\texttt{s\_empire}},
\href{https://github.com/yang-mills-group/OpenAI-YM-20261028/blob/3456e5e75cf8d4cd558955a5c7638fba2d68f464/YM/PhysicalNext/AFTERNOON_INTERACTING_RANK_COMPARISON/Soft/Supra.lean\#L76-L79}{\texttt{Turning\_}},
\href{https://github.com/yang-mills-group/OpenAI-YM-20261028/blob/3456e5e75cf8d4cd558955a5c7638fba2d68f464/YM/PhysicalNext/AFTERNOON_INTERACTING_RANK_COMPARISON/Soft/Supra.lean\#L110}{\texttt{Offers\_}}
& successive frequent-scale conditions of Lemma 4.15 \\
trial system, purchase conditions & \texttt{Columns.Trials},
\href{https://github.com/yang-mills-group/OpenAI-YM-20261028/blob/3456e5e75cf8d4cd558955a5c7638fba2d68f464/YM/Next/RANK/Bought.lean\#L68-L84}{\texttt{Purchases}},
\href{https://github.com/yang-mills-group/OpenAI-YM-20261028/blob/3456e5e75cf8d4cd558955a5c7638fba2d68f464/YM/PhysicalNext/AFTERNOON_INTERACTING_RANK_COMPARISON/Soft/Supra.lean\#L134-L147}{\texttt{Travels\_}}
& III.11.4 \\
compatible quadratic atlas &
\href{https://github.com/yang-mills-group/OpenAI-YM-20261028/blob/3456e5e75cf8d4cd558955a5c7638fba2d68f464/YM/PhysicalNext/RANK_COVARIANCE_ATLAS/Consumer.lean\#L30-L38}{\texttt{Binding}}
& shape, coercivity, size and path compatibility on the same Wilson
family \\
launch conditions &
\href{https://github.com/yang-mills-group/OpenAI-YM-20261028/blob/3456e5e75cf8d4cd558955a5c7638fba2d68f464/YM/Next/RANK/Ratchet.lean\#L35-L41}{\texttt{Launch}}
& local comparison and propagation conditions at a specified level \\
projector presentation \(\Pi\) &
\href{https://github.com/yang-mills-group/OpenAI-YM-20261028/blob/3456e5e75cf8d4cd558955a5c7638fba2d68f464/YM/MergeProof/Eureka.lean\#L16-L20}{\texttt{Scotia}}
& projectors whose kernels are corrected images of \(\mathcal E_d\) \\
one-step cost package &
\href{https://github.com/yang-mills-group/OpenAI-YM-20261028/blob/3456e5e75cf8d4cd558955a5c7638fba2d68f464/YM/MergeProof/Biscuit.lean\#L29-L34}{\texttt{Scrim}}
& step costs along gauge doors \\
Maxwell equation ideal \(\mathcal E_d\), \(F_d^{\rm free}\) &
\texttt{Maxwell.equationIdeal}, \texttt{Maxwell.freeTable} & labels
vanishing on the free field; free field spaces \\
Maxwell ranks &
\href{https://github.com/yang-mills-group/OpenAI-YM-20261028/blob/3456e5e75cf8d4cd558955a5c7638fba2d68f464/YM/Euclid/Comparator.lean\#L135-L139}{\texttt{Maxwell.Ranks}}
& equivariant isomorphisms \(F_d(T)\cong F_d^{\rm free}\) \\
numerical upper dimension bound &
\href{https://github.com/yang-mills-group/OpenAI-YM-20261028/blob/3456e5e75cf8d4cd558955a5c7638fba2d68f464/YM/Next/UV_STRESS/Invert.lean\#L108}{\texttt{DimAbove}}
& \(\dim F_d(S)\le\dim F_d^{\rm free}\) \\
equation count &
\href{https://github.com/yang-mills-group/OpenAI-YM-20261028/blob/3456e5e75cf8d4cd558955a5c7638fba2d68f464/YM/Next/RANK/Monitors.lean\#L92-L94}{\texttt{EquationCount}}
& Lemma 4.18 \\
seed record &
\href{https://github.com/yang-mills-group/OpenAI-YM-20261028/blob/3456e5e75cf8d4cd558955a5c7638fba2d68f464/YM/Spec/Closed.lean\#L20-L26}{\texttt{Spec.Seed}}
& cutoff, limits, axioms, cell hypothesis, corrected rank widths \\
shell budget &
\href{https://github.com/yang-mills-group/OpenAI-YM-20261028/blob/3456e5e75cf8d4cd558955a5c7638fba2d68f464/YM/MergeProof/West/Ozone.lean\#L16-L19}{\texttt{Aroma}},
\href{https://github.com/yang-mills-group/OpenAI-YM-20261028/blob/3456e5e75cf8d4cd558955a5c7638fba2d68f464/YM/MergeProof/North/Economist.lean\#L20-L21}{\texttt{Economics}}
& shell estimates, open tabs, radii \(r_l\to0\) \\
flux, window & \texttt{flux\ J\ χ\ U\ x\ q\ v}, \texttt{Window} (VI.1) &
flux of the stress against a vector field \\
arbitrarily fine comparison systems &
\href{https://github.com/yang-mills-group/OpenAI-YM-20261028/blob/3456e5e75cf8d4cd558955a5c7638fba2d68f464/YM/Next/UV_STRESS/Station.lean\#L112-L114}{\texttt{Stations}}
& for every order, size and positive tolerances, an admissible station
at an arbitrarily small scale \\
cofinal grid approximation &
\href{https://github.com/yang-mills-group/OpenAI-YM-20261028/blob/3456e5e75cf8d4cd558955a5c7638fba2d68f464/YM/Org/Input.lean\#L121-L123}{\texttt{Grill}}
& approximation systems at cofinally many grid levels \\
weak completeness &
\href{https://github.com/yang-mills-group/OpenAI-YM-20261028/blob/3456e5e75cf8d4cd558955a5c7638fba2d68f464/YM/MergeProof/East/Gran.lean\#L35}{\texttt{CompleteWeak}}
& comparison systems for every fixed arity, degree and positive size
bound \\
free lighthouse & \texttt{Lighthouse}, \texttt{MaxwellData} & free
Maxwell comparison model \\
\end{longtable}
}

\subsubsection{7. What remains
unexplained}\label{what-remains-unexplained-13}

\begin{enumerate}
\def\labelenumi{\arabic{enumi}.}
\tightlist
\item
  \textbf{Analytic inputs to the annulus estimate.} The reduction in
  \href{https://github.com/yang-mills-group/OpenAI-YM-20261028/blob/3456e5e75cf8d4cd558955a5c7638fba2d68f464/YM/MergeProof/West/Ozone.lean\#L21-L45}{\texttt{zone}}
  and the flux argument are checked. The discrete stencil bound,
  interpolation estimate and shell estimates invoked there have not been
  independently checked through their analytic dependencies.
\item
  \textbf{Detailed one-step interfaces.} The principal-data construction
  and the proof assembling the aggregate comparison were checked at
  their application points. The underlying nonlinear integration and
  remainder estimates were not checked in full. The relation between the
  older conditional interfaces in
  \href{https://github.com/yang-mills-group/OpenAI-YM-20261028/blob/3456e5e75cf8d4cd558955a5c7638fba2d68f464/YM/Next/RANK/LedgerGate.lean}{\texttt{LedgerGate.lean}}
  and every restarted replacement remains to be documented; their shared
  names do not identify their hypotheses.
\item
  \textbf{Record definitions.} The stage, launch, hours, vents, relay,
  station and shell hypotheses in the numbered statements refer to the
  linked records. Their complete translation into ordinary mathematical
  estimates remains to be written. Lemmas 4.11, 4.13(iv)--(vi), 4.14 and
  4.17 retain these structural hypotheses; the displayed inequalities
  alone are not their full assumptions.
\item
  \textbf{Auxiliary results and upstream estimates.} The additional
  estimates listed after III.11.7 were checked at statement level,
  including their parameter order. The sweep triangularity in Lemma 4.2
  was checked at statement level. The proofs of the lower-level
  producers of the centred estimate and one-step package, and of the
  schedule-existence theorem in III.10, were not checked in full. The
  source links identify these dependencies; the present outline does not
  independently establish them.
\end{enumerate}

\subsection{III.12. Bounds uniform in the volume and convergence as the
spacing tends to
zero}\label{iii.12.-bounds-uniform-in-the-volume-and-convergence-as-the-spacing-tends-to-zero}

\textbf{Part:} III, Ultraviolet construction (renormalization group).
\textbf{Lean:} chapter 30,
\href{https://github.com/yang-mills-group/OpenAI-YM-20261028/tree/3456e5e75cf8d4cd558955a5c7638fba2d68f464/YM/FullVolume}{\texttt{YM/FullVolume/}}.
\textbf{Size:} 18,973 lines in the main-theorem closure.

\textbf{Scope.} This unit proves volume-uniform bounds and takes the
ultraviolet limit at each fixed volume. It has three outputs:

\begin{itemize}
\tightlist
\item
  \begin{enumerate}
  \def\labelenumi{(\alph{enumi})}
  \tightlist
  \item
    Bounds on lattice correlation functions of composite observables and
    of their finite differences at separated points. The bounds are
    uniform in the volume and eventually uniform in the cutoff.
  \end{enumerate}
\item
  (b) Convergence of these correlation functions as \(N\to\infty\) at
  each fixed volume, with smooth limits.
\item
  (c) Power-series bounds on their dependence on the renormalized
  coupling. The assembly uses these.
\end{itemize}

\textbf{Notation.}

\begin{itemize}
\tightlist
\item
  \(N\) is the mesh index. The lattice spacing is \(a_N=1/D_N\), where
  \(D_N=\ell_0\cdots\ell_{N-1}\) is a product of odd block sizes.
\item
  \(p\) is the volume index. \(\Lambda_{N,p}\) is the torus of physical
  side \(L_p\), with \(L_pD_N\) sites per axis. The block size
  \(\ell_j\) and the physical period \(L_p\) are different quantities.
\item
  \(m,j,k\) are renormalization-group levels. \(\beta_{N,p}\) is the
  bare inverse coupling, \(\mathbb E_{N,p}\) the Wilson expectation, and
  \(\hat\beta_j\) the renormalized inverse coupling at level \(j\).
\item
  \(\mathrm{Lab}_d\) is the finite-dimensional space of labels of degree
  \(d\). \(O_{N,p}(v;o)\) is the lattice representative of
  \(v\in\mathrm{Lab}_d\) at the site \(o\) (II.1).
\item
  \(\mathcal N_d\subset\mathrm{Lab}_d\) is the equation ideal of the
  free Maxwell field, and \(F_d=\mathrm{Lab}_d/\mathcal N_d\).
\item
  \emph{Eventually in the mesh} means: there is \(N_0\) such that the
  statement holds for all \(N\ge N_0\) and all \(p\) (the filter
  \href{https://github.com/yang-mills-group/OpenAI-YM-20261028/blob/3456e5e75cf8d4cd558955a5c7638fba2d68f464/YM/FullVolume/Sluggish.lean\#L13}{\texttt{mesh}}).
  In statements of the form ``for all large \(m\), eventually in the
  mesh,'' \(N_0\) may depend on \(m\).
\item
  A sequence \(w_{N,p}\) in a normed space has \emph{sub-power growth}
  if for every \(\varepsilon>0\), eventually in the mesh,
  \(\|w_{N,p}\|\le a_N^{-\varepsilon}\)
  (\href{https://github.com/yang-mills-group/OpenAI-YM-20261028/blob/3456e5e75cf8d4cd558955a5c7638fba2d68f464/YM/FullVolume/Sluggish.lean\#L37-L38}{\texttt{Sluggish}}).
  The quantity
  \href{https://github.com/yang-mills-group/OpenAI-YM-20261028/blob/3456e5e75cf8d4cd558955a5c7638fba2d68f464/YM/Push/RELAY_DIFFERENCE/Rates.lean\#L40}{\texttt{flight}}
  \((\varepsilon)\) equals \(a_N^{-\varepsilon}\), because
  \href{https://github.com/yang-mills-group/OpenAI-YM-20261028/blob/3456e5e75cf8d4cd558955a5c7638fba2d68f464/YM/Lattice/Scale.lean\#L115}{\texttt{scale}}
  \((N)=\log a_N^{-1}\).
\end{itemize}

\subsubsection{1. What this unit proves}\label{what-this-unit-proves-14}

\textbf{Definitions.}

\begin{itemize}
\tightlist
\item
  A \emph{trial} \(F\) is a family of complex functions \(F_{N,p}(o)\)
  of the gauge field on \(\Lambda_{N,p}\), indexed by a site \(o\). It
  is \emph{local} if, eventually in the mesh, it is translation
  covariant and each \(F_{N,p}(o)\) is gauge invariant, measurable,
  bounded and depending only on links within \(o+X\), for one fixed
  finite \(X\subset\mathbb Z^4\)
  (\href{https://github.com/yang-mills-group/OpenAI-YM-20261028/blob/3456e5e75cf8d4cd558955a5c7638fba2d68f464/YM/Push/RELAY_DIFFERENCE/Stencils.lean\#L20-L22}{\texttt{Local}},
  \href{https://github.com/yang-mills-group/OpenAI-YM-20261028/blob/3456e5e75cf8d4cd558955a5c7638fba2d68f464/YM/FullVolume/Hits.lean\#L29}{\texttt{Stencil}}).
\item
  \emph{Chessboard root at scale \(m\).} For \(m\le N\), \(T_m(o)\) is
  the tiling of \(\Lambda_{N,p}\) by cubes of physical side \(a_m\) with
  one cell centred at \(o\)
  (\href{https://github.com/yang-mills-group/OpenAI-YM-20261028/blob/3456e5e75cf8d4cd558955a5c7638fba2d68f464/YM/Cells/Packing.lean\#L34-L40}{\texttt{center}}).
  Let \(G\) be bounded, measurable and gauge invariant, localized in the
  central cell. Then
  \(R_m(G)=\big|\mathbb E_{N,p}\prod_{c}G^{(c)}\big|^{1/\sharp\text{cells}}\).
  Here \(G^{(c)}\) is obtained by reflecting \(G\) into the cell \(c\)
  through the cell walls, and is complex conjugated after an odd number
  of reflections. The number of cells is
  \(\sharp\text{cells}=(L_p/a_m)^4\)
  (\href{https://github.com/yang-mills-group/OpenAI-YM-20261028/blob/3456e5e75cf8d4cd558955a5c7638fba2d68f464/YM/Cells/Tiling.lean\#L226-L227}{\texttt{rootTest}}).
\item
  \(F\) is \emph{affordable at cost} \(u(m,N,p)\) if \(F\) is local and
  the following holds for all large \(m\), eventually in the mesh:
  \(F_{N,p}(0)\) is localized in the central cell of \(T_m(0)\) and
  \(R_m(F_{N,p}(0))\le u(m,N,p)\).

  \begin{itemize}
  \tightlist
  \item
    \(F\) is \emph{owed at cost} \(v(N,p)\) if it is affordable at a
    cost independent of \(m\).
  \item
    \(F\) is \emph{funded} if it is affordable at a cost \(u(m)\)
    independent of \((N,p)\).
  \end{itemize}

  Lean:
  \href{https://github.com/yang-mills-group/OpenAI-YM-20261028/blob/3456e5e75cf8d4cd558955a5c7638fba2d68f464/YM/FullVolume/Rent.lean\#L19-L22}{\texttt{Affords}},
  \href{https://github.com/yang-mills-group/OpenAI-YM-20261028/blob/3456e5e75cf8d4cd558955a5c7638fba2d68f464/YM/FullVolume/Rent.lean\#L24}{\texttt{Owed}},
  \href{https://github.com/yang-mills-group/OpenAI-YM-20261028/blob/3456e5e75cf8d4cd558955a5c7638fba2d68f464/YM/FullVolume/Rent.lean\#L28}{\texttt{Funded}}.
\item
  \emph{Normalized differences:}
  \((\Delta_\mu F)_{N,p}(o)=a_N^{-1}\big(F_{N,p}(o+e_\mu)-F_{N,p}(o)\big)\),
  and \(\Delta_l=\Delta_{\mu_1}\cdots\Delta_{\mu_k}\) for a list
  \(l=(\mu_1,\dots,\mu_k)\) of axes
  (\href{https://github.com/yang-mills-group/OpenAI-YM-20261028/blob/3456e5e75cf8d4cd558955a5c7638fba2d68f464/YM/FullVolume/Hits.lean\#L16-L17}{\texttt{farther}},
  \href{https://github.com/yang-mills-group/OpenAI-YM-20261028/blob/3456e5e75cf8d4cd558955a5c7638fba2d68f464/YM/FullVolume/Hits.lean\#L19-L21}{\texttt{hops}}).
\item
  A \emph{coin} of degree \(d\) is a sequence \(x=(x_{N,p})\) in
  \(\mathrm{Lab}_d\) whose inverse terminal coordinates have sub-power
  growth
  (\href{https://github.com/yang-mills-group/OpenAI-YM-20261028/blob/3456e5e75cf8d4cd558955a5c7638fba2d68f464/YM/FullVolume/Sluggish.lean\#L146-L147}{\texttt{Coin}}).
  Thus, if \(\tau_{d;N,p}\) is the terminal coordinate isomorphism,
  \(\|\tau_{d;N,p}^{-1}x_{N,p}\|\le a_N^{-\varepsilon}\) eventually for
  every \(\varepsilon>0\). Under the consistency conditions on the
  family, bounded label sequences are coins. \(O(x)\) denotes the trial
  \(O_{N,p}(x_{N,p};\cdot)\).
\item
  A \emph{presentation} is, for each \(d\) and each \((N,p)\), an
  idempotent linear map \(\pi_{N,p}\) on \(\mathrm{Lab}_d\)
  (\href{https://github.com/yang-mills-group/OpenAI-YM-20261028/blob/3456e5e75cf8d4cd558955a5c7638fba2d68f464/YM/Next/EXTRACTION/Relay.lean\#L39-L41}{\texttt{Presentation}},
  IV.1). The kernel condition is an additional hypothesis.
\item
  For a real configuration \(x\), write
  \(\alpha_N(x)_{i\mu}=\lceil x_{i\mu}/a_N\rceil\) for its
  coordinatewise ceiling address
  (\href{https://github.com/yang-mills-group/OpenAI-YM-20261028/blob/3456e5e75cf8d4cd558955a5c7638fba2d68f464/YM/Extraction/Round.lean\#L18}{\texttt{address}}).
  This is the right endpoint convention. Interpolation means multilinear
  interpolation on the corresponding cubes
  (\href{https://github.com/yang-mills-group/OpenAI-YM-20261028/blob/3456e5e75cf8d4cd558955a5c7638fba2d68f464/YM/FullVolume/Panes.lean\#L52-L53}{\texttt{interpolated}}).
\end{itemize}

\textbf{Proposition III.12.1 (kernel bounds uniform in the volume).} Let
\(R\) be a run (§2) for a simple gauge group \(G\), with Wilson family
\(C\). Fix \(n\), labels \(v_i\in\mathrm{Lab}_{d_i}\) (\(1\le i\le n\)),
a finite list \(\lambda\) of pairs (slot \(i\), axis \(\mu\)), and
\(r>0\). Then there are \(B\ge 0\) and \(N_0\) with the following
property. For all \(N\ge N_0\), \(p\) and integer configurations \(x\)
satisfying \(a_N\,\mathrm{dist}(x_i,x_j)\ge r\) for \(i\ne j\),
\[\Big|\Delta_\lambda\,\mathbb E_{N,p}\big[O_{N,p}(v_1;x_1)\cdots O_{N,p}(v_n;x_n)\big]\Big|\le B .\]
Here \(\Delta_{(i,\mu)}\) is the normalized forward difference in
\(x_i\) along \(e_\mu\), and \(\mathrm{dist}\) is the periodic
\(\ell^\infty\) lattice distance. \(N_0\) and \(B\) may depend on
\(R,n,(v_i),\lambda,r\), but not on \(p\). \emph{Status:} theorem, for
every run; the run supplies all inputs (§2). \emph{Lean:}
\href{https://github.com/yang-mills-group/OpenAI-YM-20261028/blob/3456e5e75cf8d4cd558955a5c7638fba2d68f464/YM/FullVolume/Discrete.lean\#L50-L52}{\texttt{FixedBounds}}
(statement),
\href{https://github.com/yang-mills-group/OpenAI-YM-20261028/blob/3456e5e75cf8d4cd558955a5c7638fba2d68f464/YM/FullVolume/Discrete.lean\#L15-L16}{\texttt{gKernel}}
(the differenced kernel),
\href{https://github.com/yang-mills-group/OpenAI-YM-20261028/blob/3456e5e75cf8d4cd558955a5c7638fba2d68f464/YM/FullVolume/Literal.lean\#L62}{\texttt{genuine}}
(theorem),
\href{https://github.com/yang-mills-group/OpenAI-YM-20261028/blob/3456e5e75cf8d4cd558955a5c7638fba2d68f464/YM/FullVolume/Literal.lean\#L42-L61}{\texttt{funding}}
(reduction to III.12.2).

\textbf{Proposition III.12.2 (derivative closure of the chessboard
budget).} Let \(R\) be a run for a simple \(G\). Let \(w=(w_{N,p})\) be
a sequence in \(\mathrm{Lab}_d\) with \(\|w_{N,p}\|\le K\) eventually in
the mesh, for some \(K\ge0\), and let \(l\) be a list of axes. Then
\(\Delta_lO(w)\) is funded: for some \(u(m)\) and all large \(m\),
eventually in the mesh it is localized in the central cell and
\(R_m(\Delta_lO(w)_{N,p}(0))\le u(m)\). \emph{Status:} theorem.
\emph{Lean:}
\href{https://github.com/yang-mills-group/OpenAI-YM-20261028/blob/3456e5e75cf8d4cd558955a5c7638fba2d68f464/YM/FullVolume/Literal.lean\#L24-L33}{\texttt{Run.financed}}.

\textbf{Proposition III.12.3 (closure lemma).} Fix a degree \(d\). Let
\(C\) be a Wilson family satisfying its consistency conditions and the
following hypotheses:

\begin{enumerate}
\def\labelenumi{\arabic{enumi}.}
\tightlist
\item
  the centred budget (§2);
\item
  the scale-selection data of §2;
\item
  sound section data \(V_e\) and linear automorphisms \(W_e\) of the
  abstract labels, for every degree \(e\), satisfying the section,
  symmetry and exact local identities described in §2;
\item
  a presentation \(\pi\) and the full section-comparison data of §2,
  including the quantitative comparison with \(V_e\), identification of
  \(\ker\pi_{N,p}\) with a dressed equation ideal, and a bound
  \(\sup_{N,p}\|\pi_{N,p}\|<\infty\) for each degree;
\item
  degree-\(d\) detector rows satisfying the placement and bounded-lift
  conditions of Lemma 4.6 and the short-distance comparison of Lemma 4.8
  (\href{https://github.com/yang-mills-group/OpenAI-YM-20261028/blob/3456e5e75cf8d4cd558955a5c7638fba2d68f464/YM/FullVolume/Maintain.lean\#L49-L52}{\texttt{VitalRank}}).
\end{enumerate}

Let \(F\) be a funded trial whose values \(F_{N,p}(o)\) are continuous
functions of the gauge field, and let \(l\) be a list of axes. Let \(x\)
be a coin of degree \(d\) such that \(O(x)-\Delta_lF\) is owed at a cost
\(u_{N,p}\to0\) in the mesh. Then \(\Delta_lF\) is funded.
\emph{Status:} theorem under the stated hypotheses. For a run, these
follow from its ties, with the detector comparison obtained from its
short-distance estimates (§2). \emph{Lean:}
\href{https://github.com/yang-mills-group/OpenAI-YM-20261028/blob/3456e5e75cf8d4cd558955a5c7638fba2d68f464/YM/FullVolume/Moneys.lean\#L69-L78}{\texttt{improved}},
\href{https://github.com/yang-mills-group/OpenAI-YM-20261028/blob/3456e5e75cf8d4cd558955a5c7638fba2d68f464/YM/FullVolume/Moneys.lean\#L23-L63}{\texttt{cap}}.

\textbf{Proposition III.12.4 (uniform rank detection).} Assume
hypotheses 1--5 of III.12.3 and fix a list \(l\) of axes. For all large
\(k\) and all sufficiently small \(b>0\) there is \(Z\ge0\) such that,
eventually in the mesh,
\[\|\pi_{N,p}x\|\le Z\,\big\|\mathcal P_{k,b,l;N,p}(\pi_{N,p}x)\big\|\qquad\text{for all }x\in\mathrm{Lab}_d .\]
Here \(\mathcal P_{k,b,l;N,p}:\mathrm{Lab}_d\to\mathbb C^J\) is the
linear detector of Lemma 4.6: correlations with fixed partners, averaged
successively in the directions of \(l\) over lengths between \(b\) and
\(2b\). Repeated axes give repeated averages. \emph{Status:} theorem
(conditional on 1--5). \emph{Lean:}
\href{https://github.com/yang-mills-group/OpenAI-YM-20261028/blob/3456e5e75cf8d4cd558955a5c7638fba2d68f464/YM/FullVolume/Lit.lean\#L51-L91}{\texttt{trained}},
\href{https://github.com/yang-mills-group/OpenAI-YM-20261028/blob/3456e5e75cf8d4cd558955a5c7638fba2d68f464/YM/FullVolume/Windows.lean\#L88}{\texttt{probe}}.

\textbf{Proposition III.12.5 (labels in the equation ideal give
negligible observables).} Assume the consistency conditions on \(C\) and
hypotheses 2--4 of III.12.3. Let \(x\) be a coin of degree \(d\) with
\(\pi_{N,p}x_{N,p}=0\) for all \((N,p)\). Then \(O(x)\) is owed at cost
\(a_N^{1/32}\): it is local, and for all large \(m\), eventually in the
mesh, \(R_m(O(x)_{N,p}(0))\le a_N^{1/32}\). \emph{Status:} theorem.
\emph{Lean:}
\href{https://github.com/yang-mills-group/OpenAI-YM-20261028/blob/3456e5e75cf8d4cd558955a5c7638fba2d68f464/YM/FullVolume/Lit.lean\#L17-L19}{\texttt{sweep}},
\href{https://github.com/yang-mills-group/OpenAI-YM-20261028/blob/3456e5e75cf8d4cd558955a5c7638fba2d68f464/YM/FullVolume/Silent.lean\#L105-L115}{\texttt{banked}}.

\textbf{Proposition III.12.6 (polynomial bounds on the dependence on the
coupling).} Let \(R\) be a run, \(T\) a
\href{https://github.com/yang-mills-group/OpenAI-YM-20261028/blob/3456e5e75cf8d4cd558955a5c7638fba2d68f464/YM/FullVolume/Stage/Crests.lean\#L20-L27}{campaign}
for \(R\) with
\href{https://github.com/yang-mills-group/OpenAI-YM-20261028/blob/3456e5e75cf8d4cd558955a5c7638fba2d68f464/YM/FullVolume/Gables/Vendee.lean\#L22-L27}{pillars}
\(W\), and assume the bounds
\href{https://github.com/yang-mills-group/OpenAI-YM-20261028/blob/3456e5e75cf8d4cd558955a5c7638fba2d68f464/YM/FullVolume/Fares/Oz.lean\#L22-L28}{\texttt{Sky}}.
Let \(\bar\beta\) be the target level-0 coupling and \(w>0\) the fan
width
(\href{https://github.com/yang-mills-group/OpenAI-YM-20261028/blob/3456e5e75cf8d4cd558955a5c7638fba2d68f464/YM/FullVolume/Gain/Resolve.lean\#L72}{\texttt{ask}}).
Fix \(p\), labels \(V_0\) (an \(n\)-tuple) and \(r>0\). Then there are
\(b,m\in\mathbb N\) with \(m\ge1\), \(B\ge0\) and \(c\in(0,1]\) with the
following property.

Take \(m\le j<N\) and \(a\in[-1,1]\). Take a family \(s\mapsto F(s)\),
\(s\in[-1,1]\), of admissible campaign members, with the same block
sequence, sections and labels, whose level-0 inverse coupling at
\((N,p)\) is \(\bar\beta+w\,a\,s\). Use the expectation and
representatives of the attached auxiliary transfer family; their
correlations agree with those of the actual successful trajectory at
these positive cutoffs (Lemma 4.18). Take an \(r\)-separated
configuration \(o\) of sites of \(\Lambda_{j+1,p}\), embedded into
\(\Lambda_{N,p}\). Then the correlation
\[s\longmapsto \mathbb E^{F(cs/2)}_{N,p}\Big[\prod_i O(V_{0,i};o_i)\Big],\qquad s\in[-1,1],\]
equals \(\sum_{i\ge0}P_is^i\) with \(\sum_i|P_i|\le B\,a_{j+1}^{-b}\).
The constants depend only on the campaign data, \(p\), the labels and
\(r\). \emph{Status:} theorem, conditional on the campaign records.
These are supplied on the main-theorem route by the assembly (VIII.1).
\emph{Lean:}
\href{https://github.com/yang-mills-group/OpenAI-YM-20261028/blob/3456e5e75cf8d4cd558955a5c7638fba2d68f464/YM/FullVolume/Fares/Scrap.lean\#L40-L51}{\texttt{Face.financed}},
\href{https://github.com/yang-mills-group/OpenAI-YM-20261028/blob/3456e5e75cf8d4cd558955a5c7638fba2d68f464/YM/FullVolume/Spiral/Dome.lean\#L31-L32}{\texttt{PaidBeam}},
\href{https://github.com/yang-mills-group/OpenAI-YM-20261028/blob/3456e5e75cf8d4cd558955a5c7638fba2d68f464/YM/FullVolume/Ray/Gain.lean\#L22}{\texttt{Ray.sun}}
(\(=[-1,1]\)),
\href{https://github.com/yang-mills-group/OpenAI-YM-20261028/blob/3456e5e75cf8d4cd558955a5c7638fba2d68f464/YM/FullVolume/Relay/Budget.lean\#L34}{\texttt{intervalClock}}
(\(=a_{j+1}^{-1}\) at level \(j+1\)),
\href{https://github.com/yang-mills-group/OpenAI-YM-20261028/blob/3456e5e75cf8d4cd558955a5c7638fba2d68f464/YM/FullVolume/Spiral/Chrome.lean\#L58}{\texttt{slash}},
\href{https://github.com/yang-mills-group/OpenAI-YM-20261028/blob/3456e5e75cf8d4cd558955a5c7638fba2d68f464/YM/FullVolume/Spiral/Chrome.lean\#L65}{\texttt{verso}}.

\textbf{Proposition III.12.7 (Cauchy property in the mesh at fixed
volume).} Fix a simple gauge group and one renormalization-group
construction, with its stationary-path bounds, free amplification,
observable-transfer bounds and common reference sections. Assume the
uniform perturbation estimates for that construction (§2). Fix \(p\).
Let \((C_n)\) be families formed from histories in this construction at
cutoffs \(f(n)\to\infty\). Assume:

\begin{itemize}
\tightlist
\item
  \begin{enumerate}
  \def\labelenumi{(\roman{enumi})}
  \tightlist
  \item
    \href{https://github.com/yang-mills-group/OpenAI-YM-20261028/blob/3456e5e75cf8d4cd558955a5c7638fba2d68f464/YM/FullVolume/Fields/Capital.lean\#L19-L22}{\texttt{Rising}}:
    each history is a selected branch from a level \(k\) on, and agrees
    with a fixed kinetic path. The stationary kinetic paths converge
    with a strict reserve in their admissible bounds. For every fixed
    level \(i\ge k\), the inverse couplings are Cauchy and the polymer
    activities admit Cauchy representatives in the common weighted space
    (\href{https://github.com/yang-mills-group/OpenAI-YM-20261028/blob/3456e5e75cf8d4cd558955a5c7638fba2d68f464/YM/FullVolume/Gables/Honors.lean\#L30-L35}{\texttt{Helmet}}).
  \end{enumerate}
\item
  (ii) For each label degree used below, its reference level \(t(d)\)
  satisfies \(k\le t(d)\) and \(d+3\le J(t(d))\). The transfer data are
  based at the reference section at \(t(d)\) and are obtained by
  composing the branch transfers of Lemma 4.16.
\item
  (iii) For every sufficiently large base level \(r\), the multiscale
  remainder bounds hold for all sufficiently large sequence indices
  \(n\). Their label allowances \(\zeta(d,r)\) are eventually positive
  for each degree used, and their occupancy allowances are at least
  \(1\) at every level.
\end{itemize}

Fix a finite label list \((v_i)_{i=1}^s\) and a configuration
\(x\in(\mathbb R^4)^s\) with \(x_i-x_j\notin L_p\mathbb Z^4\) for
\(i\ne j\). Assume that \(x\)
\href{https://github.com/yang-mills-group/OpenAI-YM-20261028/blob/3456e5e75cf8d4cd558955a5c7638fba2d68f464/YM/FullVolume/Fields/Eves.lean\#L12}{lies
exactly on a common lattice} of some level \(r\): \(x=a_r z\) for an
integer configuration \(z\). Then the correlations of \(C_n\) at cutoff
\(f(n)\), evaluated at the coordinatewise ceiling lattice points of
\(x\), form a Cauchy sequence in \(n\). \emph{Status:} theorem under
these hypotheses. For a run equipped with a campaign, pillars,
fixed-level convergence data and the uniform perturbation estimates
(§2), it applies with \(C_n=C\) and \(f(n)=n\). \emph{Lean:}
\href{https://github.com/yang-mills-group/OpenAI-YM-20261028/blob/3456e5e75cf8d4cd558955a5c7638fba2d68f464/YM/FullVolume/Fields/Seat.lean\#L37-L45}{\texttt{Rising.cinema}},
\href{https://github.com/yang-mills-group/OpenAI-YM-20261028/blob/3456e5e75cf8d4cd558955a5c7638fba2d68f464/YM/FullVolume/Fields/Timber.lean\#L73-L103}{\texttt{Rising.arrives}};
run instance
\href{https://github.com/yang-mills-group/OpenAI-YM-20261028/blob/3456e5e75cf8d4cd558955a5c7638fba2d68f464/YM/FullVolume/Fields/Filled.lean\#L30-L38}{\texttt{Campaign.sees}}.

\textbf{Proposition III.12.8 (positional ultraviolet limit).} Let \(C\)
be a Wilson family and \(e\) a strictly increasing sequence of mesh
indices. Assume:

\begin{itemize}
\tightlist
\item
  the conclusion of III.12.1
  (\href{https://github.com/yang-mills-group/OpenAI-YM-20261028/blob/3456e5e75cf8d4cd558955a5c7638fba2d68f464/YM/FullVolume/Discrete.lean\#L50-L52}{\texttt{FixedBounds}});
\item
  the lattice correlations at coordinatewise ceiling lattice points are
  Cauchy along \(e\) at every configuration satisfying
  \(x_{i\mu}-x_{j\mu}\notin L_p\mathbb Z\) for every \(i\ne j\) and
  every axis \(\mu\)
  (\href{https://github.com/yang-mills-group/OpenAI-YM-20261028/blob/3456e5e75cf8d4cd558955a5c7638fba2d68f464/YM/FullVolume/Tuned.lean\#L56-L58}{\texttt{Purchased}}).
\end{itemize}

Then for every \(p\), \(n\) and \(v_1,\dots,v_n\) there is a function
\(f\) on \((\mathbb R^4)^n\), \(C^\infty\) on
\(\mathrm{NoHit}_p=\{x:\ x_i-x_j\notin L_p\mathbb Z^4\ (i\neq j)\}\)
(\href{https://github.com/yang-mills-group/OpenAI-YM-20261028/blob/3456e5e75cf8d4cd558955a5c7638fba2d68f464/YM/Extraction/Patch.lean\#L25-L26}{\texttt{NoHit}}),
such that:

\begin{itemize}
\tightlist
\item
  \begin{enumerate}
  \def\labelenumi{(\alph{enumi})}
  \tightlist
  \item
    for every smooth compactly supported test function \(t\) whose
    support avoids collisions and lies strictly inside the centred
    period cell in every coordinate, the smeared lattice sums converge
    along \(e\) to \(\int tf\);
  \end{enumerate}
\item
  (b) fix any finite list of coordinate directions (slot, axis), and in
  each direction choose integers \(q_-,q_+\ge0\) with \(q_-+q_+>0\). Use
  the stencil \((F(x+q_+e)-F(x-q_-e))/((q_-+q_+)a_N)\). If \(g\) is the
  corresponding iterated coordinate derivative of \(f\), then:

  \begin{itemize}
  \tightlist
  \item
    along \(e(s_i)\) for every sequence \(s_i\to\infty\), the difference
    quotients at lattice configurations whose physical positions
    converge to \(x\in\mathrm{NoHit}_p\) converge to \(g(x)\)
    (\href{https://github.com/yang-mills-group/OpenAI-YM-20261028/blob/3456e5e75cf8d4cd558955a5c7638fba2d68f464/YM/FullVolume/Strings.lean\#L55-L58}{\texttt{Cluster}});
  \item
    their multilinear interpolations converge locally uniformly on
    \(\mathrm{NoHit}_p\);
  \item
    \(tg\) is integrable and the smeared sums of the differences
    converge to \(\int tg\) for every test \(t\) as in (a).
  \end{itemize}
\end{itemize}

\emph{Status:} theorem under the two displayed hypotheses. For a run
with a campaign, pillars, fixed-level convergence and uniform
perturbation estimates, the full-sequence instance \(e=(0,1,2,\dots)\)
is
\href{https://github.com/yang-mills-group/OpenAI-YM-20261028/blob/3456e5e75cf8d4cd558955a5c7638fba2d68f464/YM/FullVolume/Fields/Filled.lean\#L45}{\texttt{Campaign.realUV}}.
Its pointwise Cauchy hypothesis is derived from III.12.7 in
\href{https://github.com/yang-mills-group/OpenAI-YM-20261028/blob/3456e5e75cf8d4cd558955a5c7638fba2d68f464/YM/FullVolume/Fields/Filled.lean\#L40-L44}{\texttt{Campaign.purchased}}.
\emph{Lean:}
\href{https://github.com/yang-mills-group/OpenAI-YM-20261028/blob/3456e5e75cf8d4cd558955a5c7638fba2d68f464/YM/FullVolume/Profiles.lean\#L25-L36}{\texttt{PositionsUV}},
\href{https://github.com/yang-mills-group/OpenAI-YM-20261028/blob/3456e5e75cf8d4cd558955a5c7638fba2d68f464/YM/FullVolume/Profiles.lean\#L42-L47}{\texttt{positioned}}.
\href{https://github.com/yang-mills-group/OpenAI-YM-20261028/blob/3456e5e75cf8d4cd558955a5c7638fba2d68f464/YM/FullVolume/Profiles.lean\#L52-L57}{\texttt{printed}}
is a variant along an extraction order and is not on the main-theorem
route.

\textbf{Use in the assembly.} The record
\href{https://github.com/yang-mills-group/OpenAI-YM-20261028/blob/3456e5e75cf8d4cd558955a5c7638fba2d68f464/YM/FullVolume/Finale/Program.lean\#L133-L149}{\texttt{Final.Jewel}}
is produced by
\href{https://github.com/yang-mills-group/OpenAI-YM-20261028/blob/3456e5e75cf8d4cd558955a5c7638fba2d68f464/YM/RG17/Joint.lean\#L102-L111}{\texttt{jewelOf17}}.
Its fields take the outputs of this unit:

\begin{itemize}
\tightlist
\item
  The kernel-bound conclusion is III.12.1.
\item
  The full-mesh conclusion is III.12.8 on the full sequence. This is
  (D5); through
  \href{https://github.com/yang-mills-group/OpenAI-YM-20261028/blob/3456e5e75cf8d4cd558955a5c7638fba2d68f464/YM/FullVolume/Finale/Program.lean\#L117-L126}{\texttt{from\_spec}}
  it gives the ordered limit with the full ultraviolet sequence, (D4).
\item
  The analytic-family conclusion uses III.12.6, through the chain of
  §4.4 ending in
  \href{https://github.com/yang-mills-group/OpenAI-YM-20261028/blob/3456e5e75cf8d4cd558955a5c7638fba2d68f464/YM/FullVolume/Finale/Atlas.lean\#L94-L104}{\texttt{Fan.standing}}
  and
  \href{https://github.com/yang-mills-group/OpenAI-YM-20261028/blob/3456e5e75cf8d4cd558955a5c7638fba2d68f464/YM/FullVolume/Finale/Program.lean\#L84-L99}{\texttt{Travel}}.
\end{itemize}

\subsubsection{2. What it uses}\label{what-it-uses-14}

\textbf{The run and its ties.} A
\href{https://github.com/yang-mills-group/OpenAI-YM-20261028/blob/3456e5e75cf8d4cd558955a5c7638fba2d68f464/YM/FullVolume/Cross.lean\#L30-L48}{\texttt{Run}}
(II.2, III.14) fixes the block sizes and numerical bounds, bare history,
interlevel gauge maps and observable transfers. Its degree-\(d\)
reference level satisfies \(J(t(d))\ge d+3\), and its Wilson family
satisfies the consistency conditions.

\href{https://github.com/yang-mills-group/OpenAI-YM-20261028/blob/3456e5e75cf8d4cd558955a5c7638fba2d68f464/YM/FullVolume/Cross.lean\#L73-L75}{\texttt{Run.again}},
proved through
\href{https://github.com/yang-mills-group/OpenAI-YM-20261028/blob/3456e5e75cf8d4cd558955a5c7638fba2d68f464/YM/FullVolume/Forge.lean\#L60-L72}{\texttt{ahead}}
gives an extracted hierarchy and reconstruction after refining any
extraction order, together with
\href{https://github.com/yang-mills-group/OpenAI-YM-20261028/blob/3456e5e75cf8d4cd558955a5c7638fba2d68f464/YM/FullVolume/Cross.lean\#L65-L71}{\texttt{Run.Ties}},
which has these fields:

\begin{itemize}
\tightlist
\item
  The centred budget
  \href{https://github.com/yang-mills-group/OpenAI-YM-20261028/blob/3456e5e75cf8d4cd558955a5c7638fba2d68f464/YM/Extraction/Center.lean\#L38-L40}{\texttt{Wcmain}}
  for the family. For every \(d\), every \(R>0\) and every
  \(\varepsilon>0\) there is \(A\ge0\) such that, for every
  \(m\ge t(d)\), eventually in \(N\), for all \(p\), all
  \(v\in\mathrm{Lab}_d\) with \(\|v\|\le R\) and all sites \(o\):
  \(O_{N,p}(v;o)\) is localized in the cell of \(T_m(o)\) centred at
  \(o\), and its chessboard root is at most \(A\,a_m^{-d-\varepsilon}\)
  (\href{https://github.com/yang-mills-group/OpenAI-YM-20261028/blob/3456e5e75cf8d4cd558955a5c7638fba2d68f464/YM/Extraction/Center.lean\#L21-L23}{\texttt{centeredBudget}}).
\item
  The
  \href{https://github.com/yang-mills-group/OpenAI-YM-20261028/blob/3456e5e75cf8d4cd558955a5c7638fba2d68f464/YM/Overlay/Prior.lean\#L15-L18}{scale-selection
  conditions}.
\item
  Remainder and regularity estimates for observables (III.11).
\item
  Section data \(V_d\), dressings \(W_d\) satisfying
  \href{https://github.com/yang-mills-group/OpenAI-YM-20261028/blob/3456e5e75cf8d4cd558955a5c7638fba2d68f464/YM/Push/RELAY_MARK/Portfolio.lean\#L51-L58}{\texttt{Portfolio}},
  a presentation, and
  \href{https://github.com/yang-mills-group/OpenAI-YM-20261028/blob/3456e5e75cf8d4cd558955a5c7638fba2d68f464/YM/FullVolume/Fire.lean\#L17-L20}{\texttt{Hearthside}}.
\item
  For every limit table and every degree \(d\), detector rows
  (\href{https://github.com/yang-mills-group/OpenAI-YM-20261028/blob/3456e5e75cf8d4cd558955a5c7638fba2d68f464/YM/Cells/Corrected.lean\#L53-L66}{\texttt{Columns.Trials}})
  with the short-distance data
  \href{https://github.com/yang-mills-group/OpenAI-YM-20261028/blob/3456e5e75cf8d4cd558955a5c7638fba2d68f464/YM/Next/RANK/Bought.lean\#L68-L84}{\texttt{Purchases}}.
\end{itemize}

Under the section hypotheses of III.12.3, the section is sound, has the
prescribed polynomial realization and finite regularity bounds, and its
dressing commutes with signed lattice symmetries and conjugation.
Dressed equation labels satisfy the distant-spectator identity of Lemma
4.10.

The
\href{https://github.com/yang-mills-group/OpenAI-YM-20261028/blob/3456e5e75cf8d4cd558955a5c7638fba2d68f464/YM/MergeProof/Canopy.lean\#L30-L32}{section-comparison
estimates} include cutoff-dependent automorphisms \(A_{d;N,p}\) on
terminal coordinates, auxiliary comparison columns, and a sub-power
sequence \(E_d(N)\). Eventually, uniformly in \(p\), \(A_{d;N,p}\) is
the symmetry-averaged map determined by these columns and its inverse
increases coefficient size by at most a factor \(2\). For every
comparison column and every signed or conjugated version, the centred
root of the difference between the auxiliary section evaluation and the
chosen section evaluation is bounded by \[C_d(m,k)E_d(N)a_N^{1/8}.\]
This holds for every admissible auxiliary level \(k<N\) and centred
scale \(m\le N\); \(C_d(m,k)\) is the comparison constant from III.11.
The scale-selection data choose such a \(k\) for each sufficiently large
\(m\). The kernel of \(\pi_{N,p}\) is the image of \(\mathcal N_d\)
under the dressing, threshold rescaling, terminal map and \(A_{d;N,p}\).
The data also retain the projection symmetries and the finite-cutoff
equation and difference comparisons needed by extraction.

Other existence formulations are
\href{https://github.com/yang-mills-group/OpenAI-YM-20261028/blob/3456e5e75cf8d4cd558955a5c7638fba2d68f464/YM/FullVolume/Cross.lean\#L55-L60}{\texttt{Run.ramp}},
\href{https://github.com/yang-mills-group/OpenAI-YM-20261028/blob/3456e5e75cf8d4cd558955a5c7638fba2d68f464/YM/FullVolume/Cross.lean\#L83-L95}{\texttt{realized}},
\href{https://github.com/yang-mills-group/OpenAI-YM-20261028/blob/3456e5e75cf8d4cd558955a5c7638fba2d68f464/YM/FullVolume/Forge.lean\#L28-L49}{\texttt{dagger}},
\href{https://github.com/yang-mills-group/OpenAI-YM-20261028/blob/3456e5e75cf8d4cd558955a5c7638fba2d68f464/YM/FullVolume/Unlatch.lean\#L38-L76}{\texttt{rampant}},
\href{https://github.com/yang-mills-group/OpenAI-YM-20261028/blob/3456e5e75cf8d4cd558955a5c7638fba2d68f464/YM/FullVolume/Finale/Program.lean\#L154-L164}{\texttt{Final.july}}.

\textbf{From other units.}

\begin{itemize}
\tightlist
\item
  \textbf{III.11.}

  \begin{itemize}
  \tightlist
  \item
    The centred budget for the run's family is
    \href{https://github.com/yang-mills-group/OpenAI-YM-20261028/blob/3456e5e75cf8d4cd558955a5c7638fba2d68f464/YM/MergeProof/Gardens.lean\#L19-L40}{\texttt{gardena}},
    via
    \href{https://github.com/yang-mills-group/OpenAI-YM-20261028/blob/3456e5e75cf8d4cd558955a5c7638fba2d68f464/YM/MergeProof/Baskets.lean\#L36-L67}{\texttt{coastwise}}
    and
    \href{https://github.com/yang-mills-group/OpenAI-YM-20261028/blob/3456e5e75cf8d4cd558955a5c7638fba2d68f464/YM/MergeProof/Balloon.lean\#L22-L47}{\texttt{tent}}.
    The centred logarithmic root bound is eventually at most
    \((d+\varepsilon)\log a_i^{-1}\).
  \item
    The short-distance estimates for detector lifts compare their
    cumulant polynomials with the free Maxwell profile in chessboard
    norm. The lifts are bounded and admissibly placed, and the rescaled
    test configurations converge.
  \item
    The remainder comparison for normalized local observables and their
    regularity constants is the input of
    \href{https://github.com/yang-mills-group/OpenAI-YM-20261028/blob/3456e5e75cf8d4cd558955a5c7638fba2d68f464/YM/FullVolume/Paid.lean\#L47-L49}{\texttt{prepared\_rows}}.
  \end{itemize}
\item
  \textbf{II.1.} Labels and representatives. The normalized observables
  \(a_N^dF\)
  (\href{https://github.com/yang-mills-group/OpenAI-YM-20261028/blob/3456e5e75cf8d4cd558955a5c7638fba2d68f464/YM/Push/RELAY_DIFFERENCE/RawTop.lean\#L25-L26}{\texttt{unpriced}})
  and their regularity class
  (\href{https://github.com/yang-mills-group/OpenAI-YM-20261028/blob/3456e5e75cf8d4cd558955a5c7638fba2d68f464/YM/Push/RELAY_DIFFERENCE/RawTop.lean\#L28-L29}{\texttt{clothed}},
  \href{https://github.com/yang-mills-group/OpenAI-YM-20261028/blob/3456e5e75cf8d4cd558955a5c7638fba2d68f464/YM/Push/RELAY_DIFFERENCE/RawTop.lean\#L90-L118}{\texttt{section\_bounds}},
  \href{https://github.com/yang-mills-group/OpenAI-YM-20261028/blob/3456e5e75cf8d4cd558955a5c7638fba2d68f464/YM/Push/RELAY_DIFFERENCE/RawTop.lean\#L133-L151}{\texttt{offmask}}).
  \href{https://github.com/yang-mills-group/OpenAI-YM-20261028/blob/3456e5e75cf8d4cd558955a5c7638fba2d68f464/YM/Push/RELAY_DIFFERENCE/Placement.lean\#L232-L270}{\texttt{offset}}:
  the root on a coarse tiling at any interior offset is bounded by the
  centred root at a finer scale. The matching rate
  \href{https://github.com/yang-mills-group/OpenAI-YM-20261028/blob/3456e5e75cf8d4cd558955a5c7638fba2d68f464/YM/Push/RELAY_DIFFERENCE/Estimate.lean\#L68}{\texttt{gamma}}
  \(=1/8\).
\item
  \textbf{II.2.} The Wilson family, its consistency conditions
  (translation covariance, locality of the representatives, continuity)
  and the run record.
\item
  \textbf{IV.1.} Tilings, chessboard roots and localized observables.
  \href{https://github.com/yang-mills-group/OpenAI-YM-20261028/blob/3456e5e75cf8d4cd558955a5c7638fba2d68f464/YM/Cells/Certificate.lean\#L14-L19}{\texttt{Seam}}
  holds for every \(\beta\ge0\)
  (\href{https://github.com/yang-mills-group/OpenAI-YM-20261028/blob/3456e5e75cf8d4cd558955a5c7638fba2d68f464/YM/Cells/Certificate.lean\#L21-L24}{\texttt{laws}})
  and gives two inequalities:

  \begin{itemize}
  \tightlist
  \item
    the chessboard inequality
    \(|\mathbb E\prod_cG_c^{(c)}|\le\prod_cR(G_c)\);
  \item
    the triangle inequality \(R(G+G')\le R(G)+R(G')\).
  \end{itemize}

  \href{https://github.com/yang-mills-group/OpenAI-YM-20261028/blob/3456e5e75cf8d4cd558955a5c7638fba2d68f464/YM/Cells/Packing.lean\#L81-L84}{\texttt{seamless}}
  gives \(\beta_{N,p}\ge0\). Also:
  \href{https://github.com/yang-mills-group/OpenAI-YM-20261028/blob/3456e5e75cf8d4cd558955a5c7638fba2d68f464/YM/Extraction/Equicontinuity.lean\#L66-L103}{\texttt{stepTotal}}
  (Lipschitz bound in the positions for correlations at separated
  points, derived from centred budgets);
  \href{https://github.com/yang-mills-group/OpenAI-YM-20261028/blob/3456e5e75cf8d4cd558955a5c7638fba2d68f464/YM/Extraction/Tame.lean\#L35-L97}{\texttt{growFrom}}
  (Cauchy on a dense set plus asymptotic equicontinuity gives locally
  uniform convergence).
\item
  \textbf{III.2.} Existence of the section-comparison data and the
  numerical parameter sequence.
\item
  \textbf{III.6, III.8, III.9, III.10.} Histories and selected branches;
  stationary-branch records (used in
  \href{https://github.com/yang-mills-group/OpenAI-YM-20261028/blob/3456e5e75cf8d4cd558955a5c7638fba2d68f464/YM/FullVolume/Fields/Capital.lean\#L31-L35}{\texttt{Rising.flooded}}).
  The running coupling satisfies
  \(\hat\beta_j=\mathrm{const}+\kappa_G\log a_j^{-1}\) with
  \(\kappa_G>0\)
  (\href{https://github.com/yang-mills-group/OpenAI-YM-20261028/blob/3456e5e75cf8d4cd558955a5c7638fba2d68f464/YM/Gauge/Casimir.lean\#L174-L176}{\texttt{slopeConstant\_pos}}).
  The large-field radius grows without bound in \(\hat\beta\)
  (\href{https://github.com/yang-mills-group/OpenAI-YM-20261028/blob/3456e5e75cf8d4cd558955a5c7638fba2d68f464/YM/Work/E01.lean\#L210-L211}{\texttt{radiusB\_tendsto}}).
\item
  \textbf{VIII.1.}
  \href{https://github.com/yang-mills-group/OpenAI-YM-20261028/blob/3456e5e75cf8d4cd558955a5c7638fba2d68f464/YM/FullVolume/Glide/Vac.lean\#L31-L63}{\texttt{Winter}},
  proved by
  \href{https://github.com/yang-mills-group/OpenAI-YM-20261028/blob/3456e5e75cf8d4cd558955a5c7638fba2d68f464/YM/FullVolume/Glide/Vac.lean\#L65-L89}{\texttt{offers}},
  gives fixed-level convergence of the inverse couplings and normalized
  activities for sequences of successful trajectories whose level-0
  couplings converge. It includes the kinetic-path reserve, free
  amplification and analytic versions of the limits. The activity norm
  includes regular derivatives, forced terms and support weights; the
  scalar part of the effective action is excluded. The assembly supplies
  its stationary-path, free-amplification and uniform perturbation
  inputs; the pillars identify the run's transfers with the branch
  transfers. The bare coupling's analytic extension is
  \href{https://github.com/yang-mills-group/OpenAI-YM-20261028/blob/3456e5e75cf8d4cd558955a5c7638fba2d68f464/YM/FullVolume/Finale/Program.lean\#L70-L77}{\texttt{bareTrial}}.
\end{itemize}

\textbf{Standard mathematics.}

\begin{itemize}
\tightlist
\item
  Reflection positivity of the Wilson action: Osterwalder--Seiler,
  \emph{Ann. Phys.} 110 (1978).
\item
  Chessboard estimates: Fröhlich--Israel--Lieb--Simon, \emph{Comm. Math.
  Phys.} 62 (1978) and \emph{J. Stat. Phys.} 22 (1980); Shlosman,
  \emph{Russian Math. Surveys} 41 (1986).
\item
  The chessboard seminorm (triangle inequality for \(R\)):
  Fröhlich--Lieb, \emph{Comm. Math. Phys.} 60 (1978); Biskup,
  \emph{Reflection positivity and phase transitions in lattice spin
  models} (2009).
\item
  Summation by parts. Stability of injectivity of a linear map under
  small perturbations. The dimension count for a linear map of known
  rank.
\item
  Arzelà--Ascoli type arguments. Taylor's theorem with remainder.
\item
  Absolutely convergent power series, uniqueness from their values on a
  real interval, polynomial coefficient estimates and dominated
  convergence for series.
\item
  The renormalization-group control of effective actions uniformly in
  the volume: Balaban, \emph{Comm. Math. Phys.} 1984--1989 (references
  as in III.9).
\end{itemize}

\textbf{Imported constants and assumptions.}

\begin{itemize}
\tightlist
\item
  The exponent \(-d-\varepsilon\) in the centred budget. III.12.2 uses
  it with \(\varepsilon=1\).
\item
  The matching chain: rate \(\gamma=1/8\), giving an owed cost
  \(a_N^{1/16}\) in Lemma 4.14; the equation-ideal estimate gives
  \(a_N^{1/32}\) (III.12.5).
\item
  The tail factor \(e^{-B_j/64}\)
  (\href{https://github.com/yang-mills-group/OpenAI-YM-20261028/blob/3456e5e75cf8d4cd558955a5c7638fba2d68f464/YM/FullVolume/Fields/Sever.lean\#L16}{\texttt{brisk}}),
  where \(B_j\) is the budget radius at level \(j\). Positions must have
  mutual lattice distance at least \(B_j^2\) at level \(j\).
\item
  All numerical thresholds of the run (III.2).
\end{itemize}

\subsubsection{3. Strategy}\label{strategy-14}

The centred budget applies to primaries with bounded physical labels. A
normalized difference carries the divergent factor \(a_N^{-1}\), so it
needs a separate argument. The local regularity estimates and the
matching theorem of III.11 write \(\Delta_lO(w)\) as a
degree-\((d+|l|)\) primary \(O(x)\) with an error of root cost
\(a_N^{1/16}\). The inverse terminal coordinates of \(x\) have only
sub-power control. To bound \(\pi x\), test against fixed partner
observables and average in the differenced directions. Summation by
parts bounds these tests by \((2/b)^{|l|}\) times an undifferenced
correlation bound. At short distances the detector approaches a free
Maxwell matrix that separates the quotient by the equation ideal.
Finite-dimensional stability then bounds \(\pi x\). The auxiliary
equation sections have zero chessboard root by their distant-spectator
identity; section comparison makes \(O(x-\pi x)\) negligible. Thus
\(\Delta_lO(w)\) is funded, and the chessboard inequality bounds
products at separated points (III.12.1--III.12.5).

For convergence, fix the volume and an intermediate level \(j\).
Fixed-level convergence of the couplings, normalized activities and
effective insertions makes the scalar-normalized level-\(j\) correlation
Cauchy as the cutoff grows. The full correlation differs from it by at
most \(2F_0e^{-B_j/64}\) for all sufficiently large cutoffs, where
\(F_0\) is independent of the cutoff and \(j\). Since
\(B_j=\lceil(\log\hat\beta_j)^2\rceil\to\infty\), the full correlations
are Cauchy at configurations lying on a common lattice. Positional
Lipschitz bounds extend convergence from this dense set to every
non-collision configuration. Taylor bounds identify limits of
differences with derivatives, giving smoothness and the smeared limits
(III.12.7--III.12.8).

For dependence on the coupling, vary the level-0 value in a family
parametrized by \(\bar\beta+ws\). After reducing the parameter interval,
the lattice correlations have absolutely convergent power series whose
coefficient norms grow polynomially in \(a_j^{-1}\), uniformly in the
cutoff (III.12.6). Their real level-to-level differences are \(O(a_j)\).
A polynomial coefficient estimate after a further interval reduction
makes the differences summable in coefficient norm. Fixed-level Cauchy
convergence of the series then gives analytic continuum correlations on
a complex disc. This uses the same fixed-level convergence input as the
preceding paragraph, together with the additional analytic estimates.

\subsubsection{4. Steps}\label{steps-14}

\paragraph{4.1 Calculus of the chessboard
budget}\label{calculus-of-the-chessboard-budget}

Throughout §§4.1--4.3, the family satisfies its consistency conditions.
In §4.2, assume hypotheses 1--5 of III.12.3, with the same detector rows
and presentation throughout.

\textbf{Lemma 4.1 (closure properties).} For trials \(F,H\) and costs
\(u,v\):

\begin{itemize}
\tightlist
\item
  \(\mathrm{Aff}(F,u)\) and \(\mathrm{Aff}(H,v)\) imply
  \(\mathrm{Aff}(F\pm H,u+v)\).
\item
  \(\mathrm{Aff}(F,u)\) implies \(\mathrm{Aff}(cF,|c|u)\) for any
  \(c=c_{N,p}\in\mathbb C\).
\item
  \(\mathrm{Aff}(0,0)\), and finite sums behave the same way.
\item
  Affordability is monotone in the cost.
\item
  It is unchanged if \(F\) is modified off a tail of the mesh.
\end{itemize}

Hence funded trials form a vector space. \emph{Sketch:} linear
combinations preserve locality; the root is homogeneous and satisfies
the chessboard triangle inequality. \emph{Lean:}
\href{https://github.com/yang-mills-group/OpenAI-YM-20261028/blob/3456e5e75cf8d4cd558955a5c7638fba2d68f464/YM/FullVolume/Rent.lean\#L63-L79}{\texttt{Affords.sub}},
\href{https://github.com/yang-mills-group/OpenAI-YM-20261028/blob/3456e5e75cf8d4cd558955a5c7638fba2d68f464/YM/FullVolume/Rent.lean\#L80-L87}{\texttt{Affords.scalar}},
\href{https://github.com/yang-mills-group/OpenAI-YM-20261028/blob/3456e5e75cf8d4cd558955a5c7638fba2d68f464/YM/FullVolume/Rent.lean\#L61-L62}{\texttt{Affords.enlarged}},
\href{https://github.com/yang-mills-group/OpenAI-YM-20261028/blob/3456e5e75cf8d4cd558955a5c7638fba2d68f464/YM/FullVolume/Rent.lean\#L107-L110}{\texttt{Affords.recurses}},
\href{https://github.com/yang-mills-group/OpenAI-YM-20261028/blob/3456e5e75cf8d4cd558955a5c7638fba2d68f464/YM/FullVolume/Rent.lean\#L116-L118}{\texttt{Funded.sub}};
negation, addition, the zero trial and finite sums are treated in
\href{https://github.com/yang-mills-group/OpenAI-YM-20261028/blob/3456e5e75cf8d4cd558955a5c7638fba2d68f464/YM/FullVolume/Rent.lean}{the
same file}.

\textbf{Lemma 4.2 (all offsets).} Let \(F\) be affordable at cost \(u\).
For every \(m\) and every integer \(J\ge1\) there is \(k\ge m\) with the
following property, eventually in the mesh. For every tiling of physical
cell side \(a_m\) and every site \(x\) in its central cell at least one
\(J\)-th of the cell side from each wall, \(F_{N,p}(x)\) is localized in
that cell and its root on the tiling is at most \(u(k,N,p)\).
\emph{Sketch:} choose \(k\) with \(3Ja_k\le a_m\) and apply the offset
comparison (II.1), using translation covariance. \emph{Lean:}
\href{https://github.com/yang-mills-group/OpenAI-YM-20261028/blob/3456e5e75cf8d4cd558955a5c7638fba2d68f464/YM/FullVolume/Rent.lean\#L39-L60}{\texttt{Affords.every}}.
\emph{(proof read; offset comparison imported from II.1)}

\textbf{Lemma 4.3 (bounded primaries are funded).} Under the centred
budget, let \(v_{N,p}\in\mathrm{Lab}_d\) satisfy \(\|v_{N,p}\|\le Y\)
eventually in the mesh for some fixed \(Y\). Then \(O(v)\) is funded,
with \(u(m)=A\,a_m^{-d-1}\). \emph{Sketch:} apply the centred budget
with \(\varepsilon=1\) and \(R=\max(1,Y)\), at the offset \(0\).
\emph{Lean:}
\href{https://github.com/yang-mills-group/OpenAI-YM-20261028/blob/3456e5e75cf8d4cd558955a5c7638fba2d68f464/YM/FullVolume/Rent.lean\#L119-L128}{\texttt{Funded.primaries}}.

\textbf{Lemma 4.4 (products at separated points).} Let \(F_1,\dots,F_n\)
be affordable at costs \(u_1,\dots,u_n\), and let \(r>0\). For all large
\(k\), eventually in the mesh, every configuration \(z\) with
\(a_N\,\mathrm{dist}(z_i,z_j)\ge r\) satisfies
\[\Big|\mathbb E_{N,p}\prod_iF_i(z_i)\Big|\le\prod_iu_i(k,N,p).\] If all
\(F_i\) are funded, the left side is bounded uniformly, eventually in
the mesh. \emph{Sketch:}

\begin{itemize}
\tightlist
\item
  Choose \(l\) with \(a_l<r\), and \(k\) with \(6n\,a_k<a_l\).
\item
  Translate the tiling of side \(a_l\) so that each \(z_i\) lies in the
  interior margin of width \(3a_k\) of its cell. Distinct points lie in
  distinct cells, since their distance exceeds the cell side.
\item
  Write each \(F_i(z_i)\) as the reflected copy of an observable
  localized in the central cell. The centred estimate at scale \(k\) and
  the offset comparison bound its root by \(u_i(k,N,p)\).
\item
  Apply the chessboard inequality. The remaining cells carry the
  constant \(1\), whose root is \(1\).
\end{itemize}

\emph{Lean:}
\href{https://github.com/yang-mills-group/OpenAI-YM-20261028/blob/3456e5e75cf8d4cd558955a5c7638fba2d68f464/YM/FullVolume/Cohort.lean\#L39-L66}{\texttt{cohort}},
\href{https://github.com/yang-mills-group/OpenAI-YM-20261028/blob/3456e5e75cf8d4cd558955a5c7638fba2d68f464/YM/FullVolume/Cohort.lean\#L68-L71}{\texttt{FullVolume.financed}}.

\textbf{Lemma 4.5 (III.12.1 from III.12.2).} Suppose \(\Delta_lO(v)\) is
funded for every fixed label \(v\), in every degree, and every list
\(l\). Then the bounds of III.12.1 hold for \(C\). \emph{Sketch:}
induction on the list \(\lambda\) of pairs (slot, axis). Differencing
slot \(i\) differences its insertion
(\href{https://github.com/yang-mills-group/OpenAI-YM-20261028/blob/3456e5e75cf8d4cd558955a5c7638fba2d68f464/YM/FullVolume/Discrete.lean\#L27-L45}{\texttt{actuate}}),
by linearity of expectation. Thus the differenced kernel is a
correlation of \(\Delta_{l_i}O(v_i)\), each funded. Conclude by Lemma
4.4. Then
\href{https://github.com/yang-mills-group/OpenAI-YM-20261028/blob/3456e5e75cf8d4cd558955a5c7638fba2d68f464/YM/FullVolume/Literal.lean\#L62}{\texttt{genuine}}
applies III.12.2 to constant sequences. \emph{Lean:}
\href{https://github.com/yang-mills-group/OpenAI-YM-20261028/blob/3456e5e75cf8d4cd558955a5c7638fba2d68f464/YM/FullVolume/Literal.lean\#L42-L61}{\texttt{funding}}.

\paragraph{4.2 Detectors, rank detection and the closure
lemma}\label{detectors-rank-detection-and-the-closure-lemma}

\textbf{Lemma 4.6 (detectors; summation by parts).} A detector row \(j\)
at test scale \(k\) consists of:

\begin{itemize}
\tightlist
\item
  a test configuration \(y_{k,j}\in(\mathbb R^4)^{1+e_j}\), with
  distinct projections on every coordinate axis and contained in a box
  of side \(a_{\mathrm{width}(k)}\);
\item
  partner labels \(w_{k,j,u}\) (lifted basis labels whose norms, for
  fixed \(k\), are bounded uniformly in \(N,p\) and the finite indices).
\end{itemize}

For a trial \(F\), the \emph{monitor} is the correlation
\(M_j(F)(x)=\mathbb E_{N,p}\big[F(x_0)\prod_uO(w_{k,j,u};x_u)\big]\)
(\href{https://github.com/yang-mills-group/OpenAI-YM-20261028/blob/3456e5e75cf8d4cd558955a5c7638fba2d68f464/YM/FullVolume/Windows.lean\#L14}{\texttt{monitor}}).
The \emph{assay} averages the first slot over \(n_b=\lceil b/a_N\rceil\)
consecutive sites in each direction of \(l\), starting at the prescribed
lattice approximation of \(y_{k,j}\)
(\href{https://github.com/yang-mills-group/OpenAI-YM-20261028/blob/3456e5e75cf8d4cd558955a5c7638fba2d68f464/YM/FullVolume/Windows.lean\#L86}{\texttt{assay}}).
The detector is \(\mathcal P_{k,b,l}(x)_j=\mathrm{assay}_j(O(x))\). For
each fixed \(k,l\), there is \(b_0>0\) such that, for every \(0<b<b_0\),
the following hold:

\begin{itemize}
\tightlist
\item
  \begin{enumerate}
  \def\labelenumi{(\alph{enumi})}
  \tightlist
  \item
    if \(F\) is funded with continuous values, then
    \(\|\mathrm{assay}(\Delta_lF)\|\le v\) eventually in the mesh, for
    some constant \(v\) depending on \(F,k,b,l\);
  \end{enumerate}
\item
  (b) if \(F\) is owed at a cost \(u_{N,p}\to0\), then
  \(\mathrm{assay}(F)\to0\).
\end{itemize}

\emph{Sketch:}

\begin{itemize}
\tightlist
\item
  \href{https://github.com/yang-mills-group/OpenAI-YM-20261028/blob/3456e5e75cf8d4cd558955a5c7638fba2d68f464/YM/FullVolume/Windows.lean\#L43-L51}{\texttt{scores\_hops}}:
  \(M_j(\Delta_lF)\) is the \(a_N\)-normalized lattice difference of
  \(M_j(F)\) in the first slot.
\item
  \href{https://github.com/yang-mills-group/OpenAI-YM-20261028/blob/3456e5e75cf8d4cd558955a5c7638fba2d68f464/YM/FullVolume/Skein.lean\#L40-L45}{\texttt{same\_delta}}:
  averaging \(n\) consecutive normalized differences telescopes to
  \((na_N)^{-1}\big(f(x+ne_\mu)-f(x)\big)\), with \(na_N\in[b,2b]\).
  Iterating gives
  \href{https://github.com/yang-mills-group/OpenAI-YM-20261028/blob/3456e5e75cf8d4cd558955a5c7638fba2d68f464/YM/FullVolume/Skein.lean\#L89-L101}{\texttt{reserve\_abs}}:
  \(|\mathrm{assay}(\Delta_lF)|\le(2/b)^{|l|}\sup|M_j(F)|\) near the
  test point.
\item
  \href{https://github.com/yang-mills-group/OpenAI-YM-20261028/blob/3456e5e75cf8d4cd558955a5c7638fba2d68f464/YM/FullVolume/Maintain.lean\#L86-L101}{\texttt{warded}}:
  near a generic test point, \(|M_j(F)|\le u(w)\cdot A\). This is Lemma
  4.4 applied to \(F\) and the partners, which are funded by Lemma 4.3
  (\href{https://github.com/yang-mills-group/OpenAI-YM-20261028/blob/3456e5e75cf8d4cd558955a5c7638fba2d68f464/YM/FullVolume/Maintain.lean\#L80-L83}{\texttt{siblingsPrice}}).
\item
  For (b), averaging does not increase the bound
  (\href{https://github.com/yang-mills-group/OpenAI-YM-20261028/blob/3456e5e75cf8d4cd558955a5c7638fba2d68f464/YM/FullVolume/Skein.lean\#L77-L87}{\texttt{rounds}}).
\end{itemize}

\emph{Lean:}
\href{https://github.com/yang-mills-group/OpenAI-YM-20261028/blob/3456e5e75cf8d4cd558955a5c7638fba2d68f464/YM/FullVolume/Maintain.lean\#L104-L139}{\texttt{invites}}.

\textbf{Lemma 4.7 (averaging is harmless on primaries).} For every
\(\varepsilon>0\) and \(k\), for small \(b\) and eventually in the mesh,
\(|\mathcal P_{k,b,l}(\beta_c)_j-M_j(O(\beta_c))(y_{k,j})|\le\varepsilon\)
for all rows \(j\) and all lifted basis labels \(\beta_c\).
\emph{Sketch:}
\href{https://github.com/yang-mills-group/OpenAI-YM-20261028/blob/3456e5e75cf8d4cd558955a5c7638fba2d68f464/YM/FullVolume/Pavilion.lean\#L185-L202}{\texttt{lobby}}
gives a radius \(t>0\) such that correlations of labels of bounded norm
vary by less than \(\varepsilon\) within distance \(t\) of the test
configuration. This follows from the positional Lipschitz bound derived
from the centred budget (IV.1), uniformly in \(p\). Then choose
\(2|l|b\le t\). \emph{Lean:}
\href{https://github.com/yang-mills-group/OpenAI-YM-20261028/blob/3456e5e75cf8d4cd558955a5c7638fba2d68f464/YM/FullVolume/Maintain.lean\#L57-L78}{\texttt{probesUp}}.
\emph{(proof read; the single-step Lipschitz bound is an input from
IV.1)}

\textbf{Lemma 4.8 (short distances: the free Maxwell values).} For every
\(\varepsilon>0\), eventually in \(k\), eventually in the mesh,
simultaneously for all rows \(j\) and basis indices \(c\):
\(|M_j(O(\beta_c))(y_{k,j})-\phi_{jc}|<\varepsilon\). Here \(\phi_{jc}\)
is the value of the \(j\)-th test of the free Maxwell table on the
\(c\)-th basis vector of \(F_d\)
(\href{https://github.com/yang-mills-group/OpenAI-YM-20261028/blob/3456e5e75cf8d4cd558955a5c7638fba2d68f464/YM/FullVolume/Maintain.lean\#L21}{\texttt{focus}}).
\emph{Sketch:} the short-distance comparison of III.11 makes the lattice
correlation approach the free Maxwell correlation at the rescaled test
configuration. The latter configurations converge to the prescribed free
tests, whose correlations are continuous there. \emph{Lean:}
\href{https://github.com/yang-mills-group/OpenAI-YM-20261028/blob/3456e5e75cf8d4cd558955a5c7638fba2d68f464/YM/FullVolume/Maintain.lean\#L24-L46}{\texttt{gaugesFull}},
packaged as
\href{https://github.com/yang-mills-group/OpenAI-YM-20261028/blob/3456e5e75cf8d4cd558955a5c7638fba2d68f464/YM/FullVolume/Maintain.lean\#L49-L52}{\texttt{VitalRank}}.
\emph{(proof read; short-distance estimates imported from III.11)}

\textbf{Lemma 4.9 (rank detection; III.12.4).} The matrix
\((\phi_{jc})\) is injective on \(F_d\)
(\href{https://github.com/yang-mills-group/OpenAI-YM-20261028/blob/3456e5e75cf8d4cd558955a5c7638fba2d68f464/YM/FullVolume/Stand.lean\#L73-L83}{\texttt{true\_lower}}):
the free tests separate the quotient. \emph{Sketch:}

\begin{itemize}
\tightlist
\item
  \href{https://github.com/yang-mills-group/OpenAI-YM-20261028/blob/3456e5e75cf8d4cd558955a5c7638fba2d68f464/YM/FullVolume/Stand.lean\#L20-L39}{\texttt{firm}}
  gives \(\delta,L>0\): every linear map whose matrix is
  \(\delta\)-close to \((\phi_{jc})\) satisfies
  \(\|y\|\le L\|\cdot\, y\|\).
\item
  Take \(\varepsilon=\delta/4\). The entries of
  \(\mathcal P(\pi\beta_c)\) are within \(\delta\) of
  \(\phi_{\cdot c}\), by three estimates: Lemma 4.7; Lemma 4.8; and
  \(\mathcal P(\beta_c-\pi\beta_c)\to0\), which holds by Lemma 4.10 and
  Lemma 4.6(b).
\item
  \(\mathrm{rank}\,\pi_{N,p}\le\dim F_d\)
  (\href{https://github.com/yang-mills-group/OpenAI-YM-20261028/blob/3456e5e75cf8d4cd558955a5c7638fba2d68f464/YM/FullVolume/Lit.lean\#L21-L47}{\texttt{weighed}}).
  Indeed \(\ker\pi_{N,p}\) is the dressed image of \(\mathcal N_d\), and
  \(\mathcal N_d\) is the null space of the free table by the exact
  free-field projection.
\item
  \href{https://github.com/yang-mills-group/OpenAI-YM-20261028/blob/3456e5e75cf8d4cd558955a5c7638fba2d68f464/YM/FullVolume/Stand.lean\#L42-L56}{\texttt{rank\_upper}}:
  the map \(y\mapsto\pi(\sum_cy_c\beta_c)\) is injective by the previous
  steps, hence onto \(\mathrm{ran}\,\pi\) by dimension. This gives
  \(\|\pi x\|\le Z\|\mathcal P(\pi x)\|\) with
  \(Z=A\,(\dim F_d)\,V\,L\). Here \(A\) bounds \(\|\pi\|\) and \(V\)
  bounds \(\|\beta_c\|\).
\end{itemize}

\emph{Lean:}
\href{https://github.com/yang-mills-group/OpenAI-YM-20261028/blob/3456e5e75cf8d4cd558955a5c7638fba2d68f464/YM/FullVolume/Lit.lean\#L51-L91}{\texttt{trained}}.

\textbf{Lemma 4.10 (equation terms; III.12.5).} Let \(x\) be a coin with
\(\pi x=0\). Then \(O(x)\) is owed at cost \(a_N^{1/32}\).
\emph{Sketch:} undo the section-comparison automorphism and terminal map
to obtain a terminal-coordinate row \(y\) of sub-power growth whose
associated abstract labels belong to the dressed equation ideal. The
inverse coefficient-size bound in §2 controls this change of
coordinates.

\begin{itemize}
\tightlist
\item
  \href{https://github.com/yang-mills-group/OpenAI-YM-20261028/blob/3456e5e75cf8d4cd558955a5c7638fba2d68f464/YM/FullVolume/Ocean.lean\#L41-L75}{\texttt{ocean}}:
  \(O(x)-H(y)\) is owed at \(a_N^{1/32}\). Here \(H(y)_{N,p}(o)\) is the
  bare lattice evaluation of the section data at \(o\), scaled by
  \(a_N^{-d}\). The proof goes through
  \href{https://github.com/yang-mills-group/OpenAI-YM-20261028/blob/3456e5e75cf8d4cd558955a5c7638fba2d68f464/YM/FullVolume/Paid.lean\#L23-L29}{\texttt{studded}},
  a bound by a constant times a sub-power factor times \(a_N^{1/8}\).
\item
  \href{https://github.com/yang-mills-group/OpenAI-YM-20261028/blob/3456e5e75cf8d4cd558955a5c7638fba2d68f464/YM/FullVolume/Silent.lean\#L61-L90}{\texttt{hole}}:
  \(H(y)\) is owed at cost \(0\). By the section hypotheses, the
  evaluation of a dressed element of \(\mathcal N_d\) is
  \href{https://github.com/yang-mills-group/OpenAI-YM-20261028/blob/3456e5e75cf8d4cd558955a5c7638fba2d68f464/YM/Push/RELAY_MARK/Micro.lean\#L29-L31}{\texttt{MicroAt}}:
  for every \(\beta\neq0\) and every large torus, its Wilson expectation
  against any bounded measurable observable depending only on links at
  lattice distance \(>s\) vanishes. The spectator need not be gauge
  invariant. Once the cell side exceeds \(2s+2\) lattice units, the
  copies in the other cells are such observables. So the chessboard
  product has expectation \(0\) and \(R_m=0\).
\end{itemize}

The auxiliary equation sections are an input from II.1 (§7).
\emph{Lean:}
\href{https://github.com/yang-mills-group/OpenAI-YM-20261028/blob/3456e5e75cf8d4cd558955a5c7638fba2d68f464/YM/FullVolume/Lit.lean\#L17-L19}{\texttt{sweep}},
\href{https://github.com/yang-mills-group/OpenAI-YM-20261028/blob/3456e5e75cf8d4cd558955a5c7638fba2d68f464/YM/FullVolume/Silent.lean\#L105-L115}{\texttt{banked}}.

\textbf{Lemma 4.11 (a priori bound on \(\pi x\)).} Under the hypotheses
of III.12.3, there is a constant \(A\) such that
\(\|\pi_{N,p}x_{N,p}\|\le A\) eventually in the mesh. \emph{Sketch:}

\begin{itemize}
\tightlist
\item
  By Lemma 4.6,
  \(\|\mathcal P(x)\|=\|\mathrm{assay}(O(x))\|\le\|\mathrm{assay}(\Delta_lF)\|+\|\mathrm{assay}(O(x)-\Delta_lF)\|\le c\).
\item
  Put \(X=\pi x\), \(\theta=(1+\|X\|)^{-1}\) and \(z=\theta(x-X)\).
  Since \(|\theta|\le1\), \(\theta x\) is a coin. Since
  \(\|\theta X\|\le1\), \(\theta X\) is a coin too. Thus \(z\) is a coin
  and \(\pi z=0\), so \(Z\|\mathcal P(z)\|\le\frac12\) eventually by
  Lemmas 4.10 and 4.6(b).
\item
  Hence
  \(\|X\|\le Z\|\mathcal P(X)\|\le Z\|\mathcal P(x)\|+Z\|\mathcal P(x-X)\|\le Zc+\tfrac12(1+\|X\|)\).
\end{itemize}

Rearrange to obtain \(A=2Zc+1\). \emph{Lean:}
\href{https://github.com/yang-mills-group/OpenAI-YM-20261028/blob/3456e5e75cf8d4cd558955a5c7638fba2d68f464/YM/FullVolume/Moneys.lean\#L23-L63}{\texttt{cap}}.

\textbf{Lemma 4.12 (closure lemma; III.12.3).} \(O(\pi x)\) is funded by
Lemmas 4.11 and 4.3. \(O(x-\pi x)\) is owed at \(a_N^{1/32}\to0\) by
Lemma 4.10, hence funded: a cost tending to \(0\) is eventually below
\(1\)
(\href{https://github.com/yang-mills-group/OpenAI-YM-20261028/blob/3456e5e75cf8d4cd558955a5c7638fba2d68f464/YM/FullVolume/Moneys.lean\#L65-L67}{\texttt{fixed\_cost}}).
So \(O(x)\) is funded. Finally
\(\Delta_lF=O(x)-\big(O(x)-\Delta_lF\big)\) is funded, by Lemma 4.1.
\emph{Lean:}
\href{https://github.com/yang-mills-group/OpenAI-YM-20261028/blob/3456e5e75cf8d4cd558955a5c7638fba2d68f464/YM/FullVolume/Moneys.lean\#L69-L78}{\texttt{improved}}.

\paragraph{4.3 Differences of primaries
(III.12.2)}\label{differences-of-primaries-iii.12.2}

\textbf{Lemma 4.13 (regularity under normalized differences).} Call
\(F\) \emph{tucked of degree \(d\)} if the following holds. There are a
stencil containing \(0\) and sequences \(c_{N,p},b_{N,p}\) of sub-power
growth. Eventually in the mesh, for every site \(o\),
\(a_N^dF_{N,p}(o)\) is regular on the stencil around \(o\), its
field-derivative bounds are at most \(c_{N,p}\), and its sup norm is at
most \(b_{N,p}\). Then:

\begin{itemize}
\tightlist
\item
  for a simple gauge group, \(O(w)\) is tucked of degree \(d\) for every
  coin \(w\);
\item
  if \(F\) is tucked of degree \(d\) and the constants for shifting the
  field-derivative bounds by one lattice site have sub-power growth,
  then \(\Delta_\mu F\) is tucked of degree \(d+1\). The run's remainder
  data supply these shift bounds.
\end{itemize}

\emph{Sketch:} \(a_N^{d}O(w)\) is the section evaluation at the
threshold-scaled coefficients
(\href{https://github.com/yang-mills-group/OpenAI-YM-20261028/blob/3456e5e75cf8d4cd558955a5c7638fba2d68f464/YM/Push/RELAY_DIFFERENCE/RawTop.lean\#L121-L126}{\texttt{charged}},
II.1); the bounds come from the regularity and supremum estimates for
section evaluations. For the difference,
\(a_N^{d+1}\Delta_\mu F(o)=a_N^dF(o+e_\mu)-a_N^dF(o)\). The stencil
grows by one step, and the derivative bounds pick up the sub-power shift
cost. \emph{Lean:}
\href{https://github.com/yang-mills-group/OpenAI-YM-20261028/blob/3456e5e75cf8d4cd558955a5c7638fba2d68f464/YM/FullVolume/Hits.lean\#L47-L49}{\texttt{Tucked}},
\href{https://github.com/yang-mills-group/OpenAI-YM-20261028/blob/3456e5e75cf8d4cd558955a5c7638fba2d68f464/YM/FullVolume/Hits.lean\#L60-L75}{\texttt{Tucked.up}},
\href{https://github.com/yang-mills-group/OpenAI-YM-20261028/blob/3456e5e75cf8d4cd558955a5c7638fba2d68f464/YM/FullVolume/Hits.lean\#L77-L95}{\texttt{Tucked.farther}},
\href{https://github.com/yang-mills-group/OpenAI-YM-20261028/blob/3456e5e75cf8d4cd558955a5c7638fba2d68f464/YM/FullVolume/Hits.lean\#L97-L101}{\texttt{Tucked.hops}}.

\textbf{Lemma 4.14 (matching a local observable by a primary).} Assume a
simple gauge group, the scale-selection data and the observable
remainder and regularity bounds of III.11 (§2). Let \(F\) be tucked of
degree \(d\), local, and continuous as a function of the gauge field at
the insertion site \(0\), eventually in the mesh. Then there is a coin
\(v\) of degree \(d\) such that \(O(v)-F\) is owed at cost
\(a_N^{1/16}\). \emph{Sketch:} apply the local regularity comparison of
III.11 to \(a_N^dF\). After undoing this normalization, the centred root
of \(O(v)-F\) at scale \(m\) is at most a constant (depending on \(m\)
and an auxiliary level chosen by the scale-selection conditions) times a
sub-power factor times \(a_N^{\gamma}\), with \(\gamma=1/8\).
\href{https://github.com/yang-mills-group/OpenAI-YM-20261028/blob/3456e5e75cf8d4cd558955a5c7638fba2d68f464/YM/FullVolume/Paid.lean\#L35-L46}{\texttt{readyCoin}}
and the sub-power absorption estimate absorb the sub-power factor into
\(a_N^{1/8-1/16}\). \emph{Lean:}
\href{https://github.com/yang-mills-group/OpenAI-YM-20261028/blob/3456e5e75cf8d4cd558955a5c7638fba2d68f464/YM/FullVolume/Paid.lean\#L47-L49}{\texttt{prepared\_rows}}.
\emph{(statement and reduction read; the III.11 inputs statement only)}

\textbf{Lemma 4.15 (proof of III.12.2).} In a run for a simple gauge
group, let \(w\) be eventually bounded in the mesh and set \(F=O(w)\).
Lemma 4.3 funds \(F\). By Lemma 4.13, \(\Delta_lF\) is tucked of degree
\(d+|l|\); it is also local and continuous in the gauge field. Lemma
4.14 gives a coin \(x\) of that degree with \(O(x)-\Delta_lF\) owed at
\(a_N^{1/16}\to0\). The run's ties supply hypotheses 1--5 at degree
\(d+|l|\), so III.12.3 funds \(\Delta_lF\). \emph{Lean:}
\href{https://github.com/yang-mills-group/OpenAI-YM-20261028/blob/3456e5e75cf8d4cd558955a5c7638fba2d68f464/YM/FullVolume/Literal.lean\#L24-L33}{\texttt{Run.financed}}.

\paragraph{4.4 Observable transfer and coupling
dependence}\label{observable-transfer-and-coupling-dependence}

The analytic transfer bounds use III.12.6's campaign data; analytic
continuum convergence also uses III.12.7's fixed-level convergence and
perturbation inputs.

\textbf{Lemma 4.16 (composition of observable transfers).} Fix a simple
gauge group, a branch from level \(k\) to cutoff \(N\), the
observable-transfer bounds and compatible intermediate representations
of §2. Fix \(k\le r\le N\) and \(d+3\le J(r)\). There are degree-\(d\)
transfer data based at the reference section at \(r\), obtained by
composing the one-step transfers along the branch. \emph{Sketch:} a
one-step transfer tool with its capacity estimate acts on an observable
at level \(m+1\) and produces an observable at level \(m\) whose density
is the exact renormalization-group pushforward with the degree-dependent
block normalization. It preserves linearity and symmetry and supplies
coefficient, remainder and response bounds
(\href{https://github.com/yang-mills-group/OpenAI-YM-20261028/blob/3456e5e75cf8d4cd558955a5c7638fba2d68f464/YM/FullVolume/Mark/Wheels.lean\#L38-L247}{\texttt{motor}}).
These steps compose with the prescribed top observable
(\href{https://github.com/yang-mills-group/OpenAI-YM-20261028/blob/3456e5e75cf8d4cd558955a5c7638fba2d68f464/YM/FullVolume/Mark/Vega.lean\#L47-L78}{\texttt{composeT}}).
The branch includes selectedness, available coefficient spaces and the
one-step transfer data
(\href{https://github.com/yang-mills-group/OpenAI-YM-20261028/blob/3456e5e75cf8d4cd558955a5c7638fba2d68f464/YM/FullVolume/Gables/Major.lean\#L30-L33}{\texttt{Branch}});
its estimates provide the required capacities. \emph{Lean:}
\href{https://github.com/yang-mills-group/OpenAI-YM-20261028/blob/3456e5e75cf8d4cd558955a5c7638fba2d68f464/YM/FullVolume/Gables/Major.lean\#L53-L63}{\texttt{Branch.composing}}.
Fixed-level convergence follows from
\href{https://github.com/yang-mills-group/OpenAI-YM-20261028/blob/3456e5e75cf8d4cd558955a5c7638fba2d68f464/YM/FullVolume/Gables/Master.lean\#L50-L53}{\texttt{Branch.stays}},
using the stability of the step maps under two cutoff subsequences
(\href{https://github.com/yang-mills-group/OpenAI-YM-20261028/blob/3456e5e75cf8d4cd558955a5c7638fba2d68f464/YM/FullVolume/Gables/Pole.lean\#L50-L59}{\texttt{Branch.tiers}})
(Lemma 4.19). \emph{(proof read for composition and fixed-level
reduction; one-step capacity estimate statement only)}

\textbf{Lemma 4.17 (power series in the tuning parameter, level by
level).} Assume the campaign data and estimates of III.12.6. Write
\(f\in\mathrm{PB}(t)\) if \(f(s)=\sum_iP_is^i\) on \([-1,1]\) with
\(\sum\|P_i\|\le t\). This class is closed under sums
(\(\mathrm{PB}(t)+\mathrm{PB}(t')\subset\mathrm{PB}(t+t')\)) and under
rescaling \(s\mapsto as\) with \(|a|\le1\).

\begin{itemize}
\tightlist
\item
  The one-step analytic capacity estimates bound the remainder map,
  applied to a coefficient vector \(x\), by \(\mathrm{PB}(\nu\|x\|)\).
  Here \(\nu\) is the prescribed one-step bound in the completed
  remainder norm; the parameter interval has already been reduced by a
  factor \(2\).
\item
  A bound \(\mathrm{PB}(t)\) at level \(m\) propagates to level \(m+1\)
  as \(\mathrm{PB}(\ell_m t)\)
  (\href{https://github.com/yang-mills-group/OpenAI-YM-20261028/blob/3456e5e75cf8d4cd558955a5c7638fba2d68f464/YM/FullVolume/Spiral/Dais.lean\#L64-L80}{\texttt{linked}}),
  or as \(\mathrm{PB}(2v_T\chi_d t)\)
  (\href{https://github.com/yang-mills-group/OpenAI-YM-20261028/blob/3456e5e75cf8d4cd558955a5c7638fba2d68f464/YM/FullVolume/Spiral/Brace.lean\#L61-L78}{\texttt{destination}}).
\item
  Here \(v_T\) is the prescribed transfer bound and \(\chi_d\) the
  finite-dimensional coefficient constant. Since
  \(\prod_{m=0}^j\ell_m=a_{j+1}^{-1}\), iterating gives bounds
  polynomial in \(a_{j+1}^{-1}\).
\end{itemize}

\href{https://github.com/yang-mills-group/OpenAI-YM-20261028/blob/3456e5e75cf8d4cd558955a5c7638fba2d68f464/YM/FullVolume/Relay/Vault.lean\#L83-L98}{\texttt{cruise}}
gives per-label bounds.
\href{https://github.com/yang-mills-group/OpenAI-YM-20261028/blob/3456e5e75cf8d4cd558955a5c7638fba2d68f464/YM/FullVolume/Relay/Income.lean\#L67-L72}{\texttt{polynomial}}
bounds a correlation item at \(s\mapsto cs\) by \(B_0a_{j+1}^{-b_0}\)
for fixed \(B_0,b_0\). Here \(c\) is fixed by
\href{https://github.com/yang-mills-group/OpenAI-YM-20261028/blob/3456e5e75cf8d4cd558955a5c7638fba2d68f464/YM/FullVolume/Relay/Summary.lean\#L62-L68}{\texttt{aperture}}
so that \(|c|\,Q/Z_*\le\frac12\), where \(Q\) bounds the source
contribution and \(Z_*>0\) bounds the scalar-normalized partition
function from below; see also
\href{https://github.com/yang-mills-group/OpenAI-YM-20261028/blob/3456e5e75cf8d4cd558955a5c7638fba2d68f464/YM/FullVolume/Relay/Summary.lean\#L72-L74}{\texttt{storms}}.
\href{https://github.com/yang-mills-group/OpenAI-YM-20261028/blob/3456e5e75cf8d4cd558955a5c7638fba2d68f464/YM/FullVolume/Fares/Scrap.lean\#L40-L51}{\texttt{Face.financed}}
identifies the item for the full mode with the correlation
(\href{https://github.com/yang-mills-group/OpenAI-YM-20261028/blob/3456e5e75cf8d4cd558955a5c7638fba2d68f464/YM/FullVolume/Fields/Seat.lean\#L31-L33}{\texttt{absolute}}),
which gives III.12.6. The weighted analytic bounds for the
parameter-dependent effective densities
(\href{https://github.com/yang-mills-group/OpenAI-YM-20261028/blob/3456e5e75cf8d4cd558955a5c7638fba2d68f464/YM/FullVolume/Relay/Attic.lean\#L50-L88}{\texttt{attic\_close}},
\href{https://github.com/yang-mills-group/OpenAI-YM-20261028/blob/3456e5e75cf8d4cd558955a5c7638fba2d68f464/YM/FullVolume/Relay/Pulse.lean\#L101-L104}{\texttt{many}},
\href{https://github.com/yang-mills-group/OpenAI-YM-20261028/blob/3456e5e75cf8d4cd558955a5c7638fba2d68f464/YM/FullVolume/Nectar/Laid.lean\#L85-L102}{\texttt{program}})
belong to this fan analysis.
\href{https://github.com/yang-mills-group/OpenAI-YM-20261028/blob/3456e5e75cf8d4cd558955a5c7638fba2d68f464/YM/FullVolume/Amp/Blank.lean\#L16-L39}{\texttt{Amp.give}}
and
\href{https://github.com/yang-mills-group/OpenAI-YM-20261028/blob/3456e5e75cf8d4cd558955a5c7638fba2d68f464/YM/FullVolume/Amp/Blank.lean\#L42-L43}{\texttt{Amp.stage}}
extend a multilinear series from a dense isometric image, with geometric
bounds; this builds the series on real jets without complexifying the
field. \emph{(proof read for the power-series calculus, parameter
rescaling, interval selection and final correlation bound; intermediate
analytic estimates statement only)}

\textbf{Lemma 4.18 (analytic continuum correlations).} Assume a campaign
and pillars, the analytic estimates of III.12.6, fixed-level convergence
and the uniform perturbation estimates of §2. Fix \(p\), finitely many
labels, a compact \(K\subset\mathrm{NoHit}_p\) and \(q\in\mathbb N\).
There is a positive parameter width, independent of the cutoff, and an
open neighbourhood \(\Omega\) of \(K\) within \(\mathrm{NoHit}_p\), on
which the limiting kernels are \(C^q\) in position and analytic in the
complex parameter. For each spatial stencil list of length at most
\(q\), the lattice kernels, their multilinear interpolations and their
admissible smeared sums have analytic continuations whose parameter
derivatives of every fixed order converge uniformly on compact subdiscs
and compact spatial subsets of \(\Omega\). \emph{Sketch:}
\href{https://github.com/yang-mills-group/OpenAI-YM-20261028/blob/3456e5e75cf8d4cd558955a5c7638fba2d68f464/YM/FullVolume/Tess/BarsP.lean\#L79-L96}{\texttt{pricedB}}
gives coefficient-norm bounds \(B a_j^{-b}\) for differenced kernels on
a fixed compact set.
\href{https://github.com/yang-mills-group/OpenAI-YM-20261028/blob/3456e5e75cf8d4cd558955a5c7638fba2d68f464/YM/FullVolume/Tess/Lease.lean\#L88-L96}{\texttt{traceS}}
chooses series representing these kernels. Their real values at
consecutive levels differ by \(O(a_j)\).
\href{https://github.com/yang-mills-group/OpenAI-YM-20261028/blob/3456e5e75cf8d4cd558955a5c7638fba2d68f464/YM/FullVolume/Tess/Pier.lean\#L25-L42}{A
polynomial coefficient estimate} after a fixed reduction of the
parameter interval turns this into a coefficient-norm bound
\(C\sqrt{a_j}\), summable in \(j\). For each fixed \(j\),
\href{https://github.com/yang-mills-group/OpenAI-YM-20261028/blob/3456e5e75cf8d4cd558955a5c7638fba2d68f464/YM/FullVolume/Tess/Zoom.lean\#L43-L48}{Cauchy
convergence of the real values} gives Cauchy convergence of the series
in the cutoff after another interval reduction
(\href{https://github.com/yang-mills-group/OpenAI-YM-20261028/blob/3456e5e75cf8d4cd558955a5c7638fba2d68f464/YM/FullVolume/Tess/Stream.lean\#L54-L77}{\texttt{astronomer}}).
Thus the series converge uniformly in position
(\href{https://github.com/yang-mills-group/OpenAI-YM-20261028/blob/3456e5e75cf8d4cd558955a5c7638fba2d68f464/YM/FullVolume/Tess/Tan.lean\#L104-L109}{\texttt{charted}}).
Their absolutely convergent complex power series give analytic limits
and convergence of all parameter derivatives on smaller discs. The width
may depend on \(p,K,q\). The final identification uses the actual
successful Wilson trajectories, whose
\href{https://github.com/yang-mills-group/OpenAI-YM-20261028/blob/3456e5e75cf8d4cd558955a5c7638fba2d68f464/YM/FullVolume/Finale/Real.lean\#L71-L73}{correlations
agree with those of the auxiliary transfer families} at every positive
cutoff. This is the content of the analytic spatial statement
\href{https://github.com/yang-mills-group/OpenAI-YM-20261028/blob/3456e5e75cf8d4cd558955a5c7638fba2d68f464/YM/FullVolume/Finale/Atlas.lean\#L85-L91}{\texttt{Spending}}
and its theorem
\href{https://github.com/yang-mills-group/OpenAI-YM-20261028/blob/3456e5e75cf8d4cd558955a5c7638fba2d68f464/YM/FullVolume/Finale/Atlas.lean\#L94-L104}{\texttt{Fan.standing}},
used by the assembly's analytic family. \emph{(proof read for the
series-convergence and final assembly reductions; spatial derivative
identification statement only)}

\textbf{Additional estimates in this folder.}

\begin{itemize}
\tightlist
\item
  Under the connected-activity hypotheses of III.8, the large-field step
  estimate applies when its rate parameter is at most \(-7/16\)
  (\href{https://github.com/yang-mills-group/OpenAI-YM-20261028/blob/3456e5e75cf8d4cd558955a5c7638fba2d68f464/YM/FullVolume/Florid/Gild.lean\#L57-L101}{\texttt{projects}},
  namespace of the connected-activity bounds, III.8).
\item
  Under the one-step source, degree, threshold and release hypotheses of
  III.10, the primary insertion satisfies its step estimate
  (\href{https://github.com/yang-mills-group/OpenAI-YM-20261028/blob/3456e5e75cf8d4cd558955a5c7638fba2d68f464/YM/FullVolume/Glitz/RockG.lean\#L81-L129}{\texttt{primaryJG}},
  III.10).
\item
  Fix \(g<e\), \(w<w'\), \(w\le1\), a torus with at least four sites per
  axis, and a budget point and clipping parameter satisfying the
  admissibility and kinetic-width conditions. There is \(t>0\) such that
  every \(|a|\le t\) sends a representation with bounds
  \((\sigma,C,\Lambda,g,w)\) to one with bounds
  \((\sigma,C,\Lambda,e,w')\), preserving its evaluation
  (\href{https://github.com/yang-mills-group/OpenAI-YM-20261028/blob/3456e5e75cf8d4cd558955a5c7638fba2d68f464/YM/FullVolume/Shroud/Plateau.lean\#L53-L67}{\texttt{earned}}).
\item
  Comparisons of rank data between two cutoff sequences
  (\href{https://github.com/yang-mills-group/OpenAI-YM-20261028/blob/3456e5e75cf8d4cd558955a5c7638fba2d68f464/YM/FullVolume/Rank/Tracks.lean\#L79}{\texttt{cams}},
  \href{https://github.com/yang-mills-group/OpenAI-YM-20261028/blob/3456e5e75cf8d4cd558955a5c7638fba2d68f464/YM/FullVolume/Rank/Stitch.lean\#L104-L130}{\texttt{linters}}).
\item
  The connected-activity and stationary-branch results used by these
  estimates.
\end{itemize}

These are inputs to the transfer and analytic estimates, not additional
convergence assertions. Their full hypotheses and proofs belong to
III.8--III.10; see §7.

\paragraph{4.5 Convergence in the mesh}\label{convergence-in-the-mesh}

\textbf{Lemma 4.19 (fixed-level convergence).} In the setting of
III.12.7, fix \(m\ge k\) and above the reference levels of all labels
used. Then:

\begin{itemize}
\tightlist
\item
  the inverse couplings and weighted activities are Cauchy
  (\href{https://github.com/yang-mills-group/OpenAI-YM-20261028/blob/3456e5e75cf8d4cd558955a5c7638fba2d68f464/YM/FullVolume/Fields/Capital.lean\#L31-L35}{\texttt{flooded}},
  by the fixed-level convergence input);
\item
  the level-\(m\) coefficients and coins of each label converge
  (\href{https://github.com/yang-mills-group/OpenAI-YM-20261028/blob/3456e5e75cf8d4cd558955a5c7638fba2d68f464/YM/FullVolume/Fields/Capital.lean\#L55-L56}{\texttt{grow}},
  by stability of the composed transfers);
\item
  the effective observable and its polymer parts converge uniformly
  (\href{https://github.com/yang-mills-group/OpenAI-YM-20261028/blob/3456e5e75cf8d4cd558955a5c7638fba2d68f464/YM/FullVolume/Fields/Capital.lean\#L57-L61}{\texttt{hail}}).
  The property
  \href{https://github.com/yang-mills-group/OpenAI-YM-20261028/blob/3456e5e75cf8d4cd558955a5c7638fba2d68f464/YM/FullVolume/Fields/Basin.lean\#L10-L12}{\texttt{Confluence}}
  means eventual uniform boundedness and uniform Cauchy convergence.
\end{itemize}

\emph{Sketch:} the kinetic-path reserve and activity convergence control
the normalized density. Stability of the composed transfers gives Cauchy
coefficients and remainders; evaluation preserves their uniform Cauchy
bounds. \emph{(proof read for these reductions; transfer stability
estimates statement only)}

\textbf{Lemma 4.20 (correlations computed at a fixed level converge).}
Let the level-\(j\) data be as in Lemma 4.19. Assume the effective
insertions and their polymer components are measurable, and that the
scalar-normalized partition function is eventually bounded below by
\(a>0\). Then the level-\(j\) value
\[C_j^I(n)=Z_{j,n}^{-1}\int Q_{j,n}^I(U)\,d\mu_{\mathrm{Haar}}(U)\] is
Cauchy in \(n\) for every fixed label subset \(I\). Here \(Q_{j,n}^I\)
is the coefficient indexed by \(I\) in the scalar-normalized effective
polynomial obtained by retaining the vacuum and singleton source data;
\(Z_{j,n}\) is its vacuum normalization, and \(\mu_{\mathrm{Haar}}\) is
the product Haar probability measure. \emph{Sketch:} uniform convergence
passes through the Haar integral. The normalizations have a positive
limit, so the quotients converge. The additive action scalar cancels.
\emph{Lean:}
\href{https://github.com/yang-mills-group/OpenAI-YM-20261028/blob/3456e5e75cf8d4cd558955a5c7638fba2d68f464/YM/FullVolume/Fields/Ivy.lean\#L37-L39}{\texttt{Confluence.laven}},
\href{https://github.com/yang-mills-group/OpenAI-YM-20261028/blob/3456e5e75cf8d4cd558955a5c7638fba2d68f464/YM/FullVolume/Fields/Timber.lean\#L57-L66}{\texttt{Rising.latter}},
with the uniform bounds
\href{https://github.com/yang-mills-group/OpenAI-YM-20261028/blob/3456e5e75cf8d4cd558955a5c7638fba2d68f464/YM/FullVolume/Fields/Timber.lean\#L51-L53}{\texttt{ships}}
and
\href{https://github.com/yang-mills-group/OpenAI-YM-20261028/blob/3456e5e75cf8d4cd558955a5c7638fba2d68f464/YM/FullVolume/Fields/Timber.lean\#L29-L33}{\texttt{limitVault}}.

\textbf{Lemma 4.21 (tail bound uniform in the cutoff).} In the setting
of III.12.7, fix a non-collision configuration lying exactly on a
lattice of some level. Choose a sufficiently large base level \(r\),
above the reference levels and with distinct insertion sites. Fix
positive label weights so that their weighted coefficient costs are
within the allowances at \(r\); these weights exist by
\href{https://github.com/yang-mills-group/OpenAI-YM-20261028/blob/3456e5e75cf8d4cd558955a5c7638fba2d68f464/YM/FullVolume/Fields/Capital.lean\#L62-L63}{\texttt{creditP}}.
For every fixed subset \(I\) of the labels, there is \(F_0\ge0\),
independent of \(j\) and the cutoff, such that for each sufficiently
large \(j>r\), eventually in \(n\),
\[|p_n^I-C_j^I(n)|\le 2F_0e^{-B_j/64}.\] Here \(p_n^I\) is the product
correlation at cutoff \(f(n)\) and \(C_j^I(n)\) is its scalar-normalized
level-\(j\) approximation. The lower threshold on \(j\) ensures
separation by \(B_j^2\) lattice units, support exponent at least \(2\),
and \(e^{B_j/64}>2\). The threshold on \(n\) may depend on \(j\).
\emph{Sketch:} propagate the exact polynomial density of the inserted
observables down the renormalization group. The weighted remainder
estimate controls its part of degree at least two. Separation forces a
nonzero such term to occupy enough blocks to gain the factor
\(e^{-B_j/64}\). The remaining omission estimate divides by a uniformly
positive scalar-normalized partition function. Uniform bounds on the
effective insertions at \(r\) and \(j\) follow from Lemma 4.19; the
final numerator depends only on the data at \(r\). \emph{Lean:}
\href{https://github.com/yang-mills-group/OpenAI-YM-20261028/blob/3456e5e75cf8d4cd558955a5c7638fba2d68f464/YM/FullVolume/Fields/Peat.lean\#L41-L64}{\texttt{Sever.echoes}}.
\emph{(proof read for the exact-density, remainder and omission
reduction; underlying multiscale estimates are inputs)}

\textbf{Lemma 4.22 (Cauchy property).} Under the hypotheses of Lemma
4.21, \(p_n^I\) is Cauchy for every fixed subset \(I\) of the labels.
\emph{Sketch:} \(B_j=\lceil(\log\hat\beta_j)^2\rceil\), so
\(g(j)=e^{B_j/64}\to\infty\) because
\(\hat\beta_j=\mathrm{const}+\kappa_G\log a_j^{-1}\to\infty\). Choose
\(j\) with \(g(j)>2\) and \(2F_0/g(j)\le\varepsilon/4\). Lemma 4.21
bounds \(|p_n^I-C_j^I(n)|\) by \(\varepsilon/4\) for large \(n\). Lemma
4.20 gives the Cauchy property of \(C_j^I(n)\). Conclude with the
triangle inequality. \emph{Lean:}
\href{https://github.com/yang-mills-group/OpenAI-YM-20261028/blob/3456e5e75cf8d4cd558955a5c7638fba2d68f464/YM/FullVolume/Fields/Timber.lean\#L73-L103}{\texttt{Rising.arrives}}.

\textbf{Lemma 4.23 (III.12.7).} Under the hypotheses of III.12.7, take
the subset containing all labels. The full-subset quantity is the
correlation itself
(\href{https://github.com/yang-mills-group/OpenAI-YM-20261028/blob/3456e5e75cf8d4cd558955a5c7638fba2d68f464/YM/FullVolume/Fields/Seat.lean\#L31-L33}{\texttt{absolute}}).
The positions are \(o_N=\mathrm{wrap}(\alpha_N(x))\). For
\(x\in\mathrm{NoHit}_p\) lying exactly on a lattice of some level, these
satisfy the separation \(B_j^2\) for all large \(j\) by the
lattice-ancestry estimates, and they are injective at the base level.
\emph{Lean:}
\href{https://github.com/yang-mills-group/OpenAI-YM-20261028/blob/3456e5e75cf8d4cd558955a5c7638fba2d68f464/YM/FullVolume/Fields/Seat.lean\#L37-L45}{\texttt{Rising.cinema}}.

\textbf{Lemma 4.24 (all non-collision points, full sequence).} For a run
with a campaign, pillars, fixed-level convergence and uniform
perturbation estimates, III.12.7 holds on the full sequence of cutoffs
(\href{https://github.com/yang-mills-group/OpenAI-YM-20261028/blob/3456e5e75cf8d4cd558955a5c7638fba2d68f464/YM/FullVolume/Fields/Filled.lean\#L30-L38}{\texttt{Campaign.sees}},
with the same cutoff-dependent history family repeated for every
sequence index and \(f(n)=n\)). Configurations lying exactly on some
lattice are dense. III.12.1 gives a local Lipschitz bound
(\href{https://github.com/yang-mills-group/OpenAI-YM-20261028/blob/3456e5e75cf8d4cd558955a5c7638fba2d68f464/YM/FullVolume/Region.lean\#L42-L54}{\texttt{localLip}}),
hence asymptotic equicontinuity on compact subsets of
\(\mathrm{NoHit}_p\)
(\href{https://github.com/yang-mills-group/OpenAI-YM-20261028/blob/3456e5e75cf8d4cd558955a5c7638fba2d68f464/YM/FullVolume/Tuned.lean\#L38-L54}{\texttt{circulating}}).
The dense-set convergence theorem (IV.1) then gives locally uniform
convergence on \(\mathrm{NoHit}_p\), in particular the Cauchy property
on the subset with no coordinate ties. \emph{Lean:}
\href{https://github.com/yang-mills-group/OpenAI-YM-20261028/blob/3456e5e75cf8d4cd558955a5c7638fba2d68f464/YM/FullVolume/Fields/Filled.lean\#L40-L44}{\texttt{Campaign.purchased}}.

\textbf{Lemma 4.25 (positional limits; III.12.8).} Assume the two
hypotheses of III.12.8.

\begin{itemize}
\tightlist
\item
  \href{https://github.com/yang-mills-group/OpenAI-YM-20261028/blob/3456e5e75cf8d4cd558955a5c7638fba2d68f464/YM/FullVolume/Tuned.lean\#L60-L61}{\texttt{established}}:
  the kernels converge locally uniformly on \(\mathrm{NoHit}_p\) to a
  continuous \(f\). The argument is the one of Lemma 4.24.
\item
  \href{https://github.com/yang-mills-group/OpenAI-YM-20261028/blob/3456e5e75cf8d4cd558955a5c7638fba2d68f464/YM/FullVolume/Tuned.lean\#L71-L125}{\texttt{from\_previous}}:
  if the \(l\)-differences converge, so do the differences with one
  additional coordinate direction \(a\). Fix
  \(q=\min(t/2,r/3,\varepsilon/(8(B+1)))>0\), where \(t\) is a
  neighbourhood radius and \(r,B\) are the Taylor constants. Choose
  integers \(n_N\) with \(k_N=n_Na_N\in[q,2q]\) and \(k_N\to q\). If
  \(a_Nx_N\to x\), compare the unit difference \(D_N\) with
  \(W_N=k_N^{-1}(F_N(x_N+n_Ne_a)-F_N(x_N))\). Both endpoint values
  converge
  (\href{https://github.com/yang-mills-group/OpenAI-YM-20261028/blob/3456e5e75cf8d4cd558955a5c7638fba2d68f464/YM/FullVolume/Tuned.lean\#L21-L31}{\texttt{returns}}),
  so \(W_N\) converges. The Taylor estimate
  \(|F_N(x_N+n_Ne_a)-F_N(x_N)-k_ND_N|\le Bk_N^2\)
  (\href{https://github.com/yang-mills-group/OpenAI-YM-20261028/blob/3456e5e75cf8d4cd558955a5c7638fba2d68f464/YM/FullVolume/Region.lean\#L58-L77}{\texttt{taylor}})
  gives \(|W_N-D_N|\le\varepsilon/4\) eventually. Hence \(D_N\) is
  Cauchy.
\item
  Limits of differences are derivatives of \(f\), so
  \(f\in C^\infty(\mathrm{NoHit}_p)\)
  (\href{https://github.com/yang-mills-group/OpenAI-YM-20261028/blob/3456e5e75cf8d4cd558955a5c7638fba2d68f464/YM/FullVolume/Strings.lean\#L90-L92}{\texttt{reachable}}).
\item
  Every fixed forward/backward interval stencil is a finite average of
  translates of a unit forward difference. These translates approach the
  same physical point, so they have the same derivative limit.
  Multilinear interpolation preserves this limit. For admissible tests,
  the corresponding piecewise constant Riemann sums have an integrable
  bound on a fixed compact set, and dominated convergence gives the
  smeared limits.
\end{itemize}

Under the hypotheses of Lemma 4.24, the run instance is
\href{https://github.com/yang-mills-group/OpenAI-YM-20261028/blob/3456e5e75cf8d4cd558955a5c7638fba2d68f464/YM/FullVolume/Fields/Filled.lean\#L45}{\texttt{Campaign.realUV}}:
III.12.8 applied to III.12.1 and Lemma 4.24. \emph{Lean:}
\href{https://github.com/yang-mills-group/OpenAI-YM-20261028/blob/3456e5e75cf8d4cd558955a5c7638fba2d68f464/YM/FullVolume/Profiles.lean\#L42-L47}{\texttt{positioned}}.
\emph{(proof read for equicontinuity, Taylor comparison, derivative
convergence, stencil averaging, interpolation and smearing)}

\subsubsection{5. Where the difficulty
lies}\label{where-the-difficulty-lies-14}

Matching controls only inverse terminal coordinates. Rank detection must
bound the projected physical labels, while the equation-section
comparison controls the projection kernel. These estimates precede every
limit.

The estimates of III.12.1--III.12.5 are uniform in every admitted
volume. The convergence results III.12.7--III.12.8 fix the volume first.
Their thresholds may depend on the period, labels, compact spatial set,
derivative order and test function. The parameter width in Lemma 4.18
may also depend on these data.

All positional limits are off collisions on the torus, including
configurations with coordinate ties and configurations crossing the
boundary of a period cell. Extension across collisions belongs to VII.4.
Every derivative order or stencil is fixed before the cutoff threshold
is chosen.

Sub-power factors are absorbed by the positive error exponents.
Fixed-level convergence uses a positive lower bound on the
scalar-normalized partition function; the additive action scalar need
not converge.

\subsubsection{6. Notation}\label{notation-14}

{\def\LTcaptype{none} % do not increment counter
\begin{longtable}[]{@{}
  >{\raggedright\arraybackslash}p{(\linewidth - 4\tabcolsep) * \real{0.3333}}
  >{\raggedright\arraybackslash}p{(\linewidth - 4\tabcolsep) * \real{0.3333}}
  >{\raggedright\arraybackslash}p{(\linewidth - 4\tabcolsep) * \real{0.3333}}@{}}
\toprule\noalign{}
\begin{minipage}[b]{\linewidth}\raggedright
Written proof
\end{minipage} & \begin{minipage}[b]{\linewidth}\raggedright
Lean (glossary)
\end{minipage} & \begin{minipage}[b]{\linewidth}\raggedright
Meaning
\end{minipage} \\
\midrule\noalign{}
\endhead
\bottomrule\noalign{}
\endlastfoot
eventually in the mesh &
\href{https://github.com/yang-mills-group/OpenAI-YM-20261028/blob/3456e5e75cf8d4cd558955a5c7638fba2d68f464/YM/FullVolume/Sluggish.lean\#L13}{\texttt{mesh}}
& for \(N\ge N_0\) and all \(p\) \\
sub-power growth &
\href{https://github.com/yang-mills-group/OpenAI-YM-20261028/blob/3456e5e75cf8d4cd558955a5c7638fba2d68f464/YM/FullVolume/Sluggish.lean\#L37-L38}{\texttt{Sluggish}},
\href{https://github.com/yang-mills-group/OpenAI-YM-20261028/blob/3456e5e75cf8d4cd558955a5c7638fba2d68f464/YM/Push/RELAY_DIFFERENCE/Rates.lean\#L40}{\texttt{flight}}
& \(\|w\|\le a_N^{-\varepsilon}\) for all \(\varepsilon>0\); flight
\((\varepsilon)=a_N^{-\varepsilon}\) \\
trial \(F\), local & \texttt{Trial},
\href{https://github.com/yang-mills-group/OpenAI-YM-20261028/blob/3456e5e75cf8d4cd558955a5c7638fba2d68f464/YM/FullVolume/Hits.lean\#L29}{\texttt{Stencil}}
& site-indexed lattice observables on a fixed stencil \\
\(T_m(o)\), \(R_m\) &
\href{https://github.com/yang-mills-group/OpenAI-YM-20261028/blob/3456e5e75cf8d4cd558955a5c7638fba2d68f464/YM/Cells/Packing.lean\#L34-L40}{\texttt{center}},
\href{https://github.com/yang-mills-group/OpenAI-YM-20261028/blob/3456e5e75cf8d4cd558955a5c7638fba2d68f464/YM/Cells/Tiling.lean\#L226-L227}{\texttt{rootTest}}
& centred tiling of side \(a_m\); chessboard root \\
affordable, owed, funded &
\href{https://github.com/yang-mills-group/OpenAI-YM-20261028/blob/3456e5e75cf8d4cd558955a5c7638fba2d68f464/YM/FullVolume/Rent.lean\#L19-L22}{\texttt{Affords}},
\href{https://github.com/yang-mills-group/OpenAI-YM-20261028/blob/3456e5e75cf8d4cd558955a5c7638fba2d68f464/YM/FullVolume/Rent.lean\#L24}{\texttt{Owed}},
\href{https://github.com/yang-mills-group/OpenAI-YM-20261028/blob/3456e5e75cf8d4cd558955a5c7638fba2d68f464/YM/FullVolume/Rent.lean\#L28}{\texttt{Funded}}
& centred root bound by \(u(m,N,p)\), \(v(N,p)\), \(u(m)\) \\
\(\Delta_\mu\), \(\Delta_l\) &
\href{https://github.com/yang-mills-group/OpenAI-YM-20261028/blob/3456e5e75cf8d4cd558955a5c7638fba2d68f464/YM/FullVolume/Hits.lean\#L16-L17}{\texttt{farther}},
\href{https://github.com/yang-mills-group/OpenAI-YM-20261028/blob/3456e5e75cf8d4cd558955a5c7638fba2d68f464/YM/FullVolume/Hits.lean\#L19-L21}{\texttt{hops}}
& \(a_N\)-normalized forward differences \\
tucked of degree \(d\) &
\href{https://github.com/yang-mills-group/OpenAI-YM-20261028/blob/3456e5e75cf8d4cd558955a5c7638fba2d68f464/YM/FullVolume/Hits.lean\#L47-L49}{\texttt{Tucked}}
& \(a_N^dF\) regular with sub-power bounds \\
coin; \(O(x)\) &
\href{https://github.com/yang-mills-group/OpenAI-YM-20261028/blob/3456e5e75cf8d4cd558955a5c7638fba2d68f464/YM/FullVolume/Sluggish.lean\#L146-L147}{\texttt{Coin}},
\texttt{Bid.upper} & label sequence with sub-power inverse terminal
coordinates; its representative \\
\(\pi_{N,p}\) &
\href{https://github.com/yang-mills-group/OpenAI-YM-20261028/blob/3456e5e75cf8d4cd558955a5c7638fba2d68f464/YM/Next/EXTRACTION/Relay.lean\#L39-L41}{\texttt{Presentation}},
\href{https://github.com/yang-mills-group/OpenAI-YM-20261028/blob/3456e5e75cf8d4cd558955a5c7638fba2d68f464/YM/FullVolume/Lit.lean\#L50}{\texttt{display}}
& idempotent with kernel the dressed equation ideal \\
\(M_j\), detector \(\mathcal P_{k,b,l}\) &
\href{https://github.com/yang-mills-group/OpenAI-YM-20261028/blob/3456e5e75cf8d4cd558955a5c7638fba2d68f464/YM/FullVolume/Windows.lean\#L14}{\texttt{monitor}},
\href{https://github.com/yang-mills-group/OpenAI-YM-20261028/blob/3456e5e75cf8d4cd558955a5c7638fba2d68f464/YM/FullVolume/Windows.lean\#L86}{\texttt{assay}},
\href{https://github.com/yang-mills-group/OpenAI-YM-20261028/blob/3456e5e75cf8d4cd558955a5c7638fba2d68f464/YM/FullVolume/Windows.lean\#L88}{\texttt{probe}}
& test correlation; iterated directional average \\
\(\phi_{jc}\) &
\href{https://github.com/yang-mills-group/OpenAI-YM-20261028/blob/3456e5e75cf8d4cd558955a5c7638fba2d68f464/YM/FullVolume/Maintain.lean\#L21}{\texttt{focus}}
& free Maxwell value of test \(j\) on basis vector \(c\) \\
centred budget &
\href{https://github.com/yang-mills-group/OpenAI-YM-20261028/blob/3456e5e75cf8d4cd558955a5c7638fba2d68f464/YM/Extraction/Center.lean\#L38-L40}{\texttt{Wcmain}}
& root \(\le A\,a_m^{-d-\varepsilon}\) \\
kernel bounds &
\href{https://github.com/yang-mills-group/OpenAI-YM-20261028/blob/3456e5e75cf8d4cd558955a5c7638fba2d68f464/YM/FullVolume/Discrete.lean\#L50-L52}{\texttt{FixedBounds}}
& III.12.1 \\
run, ties &
\href{https://github.com/yang-mills-group/OpenAI-YM-20261028/blob/3456e5e75cf8d4cd558955a5c7638fba2d68f464/YM/FullVolume/Cross.lean\#L30-L48}{\texttt{Run}},
\href{https://github.com/yang-mills-group/OpenAI-YM-20261028/blob/3456e5e75cf8d4cd558955a5c7638fba2d68f464/YM/FullVolume/Cross.lean\#L65-L71}{\texttt{Run.Ties}}
& record of one construction; its inputs for this unit \\
fan parameter \(s\in[-1,1]\) &
\href{https://github.com/yang-mills-group/OpenAI-YM-20261028/blob/3456e5e75cf8d4cd558955a5c7638fba2d68f464/YM/FullVolume/Ray/Gain.lean\#L22}{\texttt{Ray.sun}},
\href{https://github.com/yang-mills-group/OpenAI-YM-20261028/blob/3456e5e75cf8d4cd558955a5c7638fba2d68f464/YM/FullVolume/Gain/Resolve.lean\#L72}{\texttt{ask}}
& level-0 coupling \(\bar\beta+ws\) \\
\(\mathrm{PB}(t)\) &
\href{https://github.com/yang-mills-group/OpenAI-YM-20261028/blob/3456e5e75cf8d4cd558955a5c7638fba2d68f464/YM/FullVolume/Spiral/Dome.lean\#L31-L32}{\texttt{PaidBeam}}
& power series on \([-1,1]\) with \(\ell^1\) coefficient norm
\(\le t\) \\
\(a_j^{-1}\) &
\href{https://github.com/yang-mills-group/OpenAI-YM-20261028/blob/3456e5e75cf8d4cd558955a5c7638fba2d68f464/YM/FullVolume/Relay/Budget.lean\#L34}{\texttt{intervalClock}}
& \(D_j=\ell_0\cdots\ell_{j-1}\) \\
campaign, pillars, sky &
\href{https://github.com/yang-mills-group/OpenAI-YM-20261028/blob/3456e5e75cf8d4cd558955a5c7638fba2d68f464/YM/FullVolume/Stage/Crests.lean\#L20-L27}{\texttt{Campaign}},
\href{https://github.com/yang-mills-group/OpenAI-YM-20261028/blob/3456e5e75cf8d4cd558955a5c7638fba2d68f464/YM/FullVolume/Gables/Vendee.lean\#L22-L27}{\texttt{Pillars}},
\href{https://github.com/yang-mills-group/OpenAI-YM-20261028/blob/3456e5e75cf8d4cd558955a5c7638fba2d68f464/YM/FullVolume/Fares/Oz.lean\#L22-L28}{\texttt{Sky}}
& records carrying the fan and its bounds \\
fixed-level convergence; function convergence &
\href{https://github.com/yang-mills-group/OpenAI-YM-20261028/blob/3456e5e75cf8d4cd558955a5c7638fba2d68f464/YM/FullVolume/Gables/Honors.lean\#L30-L35}{\texttt{Helmet}},
\href{https://github.com/yang-mills-group/OpenAI-YM-20261028/blob/3456e5e75cf8d4cd558955a5c7638fba2d68f464/YM/FullVolume/Fields/Basin.lean\#L10-L12}{\texttt{Confluence}}
& Cauchy couplings and weighted activities; eventual uniform boundedness
and uniform Cauchy convergence of functions \\
point-Cauchy property &
\href{https://github.com/yang-mills-group/OpenAI-YM-20261028/blob/3456e5e75cf8d4cd558955a5c7638fba2d68f464/YM/FullVolume/Tuned.lean\#L56-L58}{\texttt{Purchased}}
& Cauchy correlations at ceiling addresses when no coordinate
projections tie \\
positional limit &
\href{https://github.com/yang-mills-group/OpenAI-YM-20261028/blob/3456e5e75cf8d4cd558955a5c7638fba2d68f464/YM/FullVolume/Profiles.lean\#L25-L36}{\texttt{PositionsUV}}
& III.12.8; (D5) of the main statement \\
\end{longtable}
}

\subsubsection{7. What remains
unexplained}\label{what-remains-unexplained-14}

\begin{enumerate}
\def\labelenumi{\arabic{enumi}.}
\tightlist
\item
  \textbf{Local matching and equation sections.} Lemma 4.14 imports the
  local-to-primary estimate of III.11; only its reduction to a coin and
  an \(a_N^{1/16}\) root bound is checked here. The construction of the
  auxiliary sections from the Maxwell equation ideal, and its relation
  to a lattice Schwinger--Dyson identity, remain unverified. Lemma 4.10
  uses their exact distant-spectator identity, without identifying it
  with the Wilson action's naive equation of motion.
\item
  \textbf{One-step observable transfer.} The capacity estimate used by
  Lemma 4.16 remains at statement level. Its coefficient and remainder
  constants require the one-step analysis.
\item
  \textbf{Analytic estimates.} The intermediate weighted estimates in
  Lemma 4.17 and the identification of spatial derivatives of the
  analytic limit in Lemma 4.18 remain at statement level. The outer
  series-convergence argument is checked.
\item
  \textbf{Additional estimates in §4.4.} The large-field and
  primary-step results and the two transfer comparisons retain their
  full hypotheses in III.8--III.10. Their proofs, and the detailed
  small-shift lawfulness estimate, are not checked here.
\item
  \textbf{Fixed-level convergence and tails.} The analytic
  trajectory-convergence theorem and the multiscale remainder and
  omission estimates remain imported inputs; their applications are
  checked. The tail argument needs only
  \(B_j=\lceil(\log\hat\beta_j)^2\rceil\to\infty\), not geometric decay
  in \(j\).
\end{enumerate}

\subsection{III.13. Extension to all labels: one common Wilson
family}\label{iii.13.-extension-to-all-labels-one-common-wilson-family}

\textbf{Part:} III, Ultraviolet construction (renormalization group).
\textbf{Lean:}
\href{https://github.com/yang-mills-group/OpenAI-YM-20261028/blob/3456e5e75cf8d4cd558955a5c7638fba2d68f464/YM/SharpGram/Common.lean}{\texttt{YM/SharpGram/Common.lean}},
\href{https://github.com/yang-mills-group/OpenAI-YM-20261028/blob/3456e5e75cf8d4cd558955a5c7638fba2d68f464/YM/SharpGram/Measured.lean}{\texttt{YM/SharpGram/Measured.lean}},
\href{https://github.com/yang-mills-group/OpenAI-YM-20261028/tree/3456e5e75cf8d4cd558955a5c7638fba2d68f464/YM/SharpGram/Assembly}{\texttt{YM/SharpGram/Assembly/}},
\href{https://github.com/yang-mills-group/OpenAI-YM-20261028/tree/3456e5e75cf8d4cd558955a5c7638fba2d68f464/YM/UVExtension}{\texttt{YM/UVExtension/}}.
\textbf{Size:} about 22,100 lines in the unit inventory. Principal
assembly results:
\href{https://github.com/yang-mills-group/OpenAI-YM-20261028/blob/3456e5e75cf8d4cd558955a5c7638fba2d68f464/YM/SharpGram/Assembly/Expenditures.lean\#L31-L79}{\texttt{multitude}},
\href{https://github.com/yang-mills-group/OpenAI-YM-20261028/blob/3456e5e75cf8d4cd558955a5c7638fba2d68f464/YM/SharpGram/Assembly/Flowers.lean\#L75-L130}{\texttt{uv\_log}},
\href{https://github.com/yang-mills-group/OpenAI-YM-20261028/blob/3456e5e75cf8d4cd558955a5c7638fba2d68f464/YM/SharpGram/Assembly/Expenditures.lean\#L84-L93}{\texttt{multitude\_foot}}.

\textbf{Place in the proof.} This unit supplies the ultraviolet
construction for
\href{https://github.com/yang-mills-group/OpenAI-YM-20261028/blob/3456e5e75cf8d4cd558955a5c7638fba2d68f464/YM/SharpGram/Measured.lean\#L11-L36}{\texttt{habitat}}
(Proposition III.13.4). The main theorem
\href{https://github.com/yang-mills-group/OpenAI-YM-20261028/blob/3456e5e75cf8d4cd558955a5c7638fba2d68f464/Construction.lean\#L15-L40}{\texttt{YangMills.constructed}}
does not use it.

\begin{itemize}
\tightlist
\item
  The main theorem gets its Wilson family from the construction run of
  III.14
  (\href{https://github.com/yang-mills-group/OpenAI-YM-20261028/blob/3456e5e75cf8d4cd558955a5c7638fba2d68f464/YM/Independent/ConstructionRun.lean\#L14-L28}{\texttt{constructionRun}},
  \href{https://github.com/yang-mills-group/OpenAI-YM-20261028/blob/3456e5e75cf8d4cd558955a5c7638fba2d68f464/YM/RG17/Joint.lean\#L60-L90}{\texttt{Run.extractible17}}).
\item
  The two constructions rest on the same per-step estimates
  (III.1--III.12) and package them differently.
\end{itemize}

\subsubsection{1. What this unit proves}\label{what-this-unit-proves-15}

\begin{itemize}
\tightlist
\item
  \textbf{The constants are chosen once.} One choice of
  renormalization-group (RG) constants supplies a certified schedule,
  chessboard bounds, the logarithmic OPE package, coverage, weak
  comparison systems and extraction data (§4.1).
\item
  \textbf{Fixed-volume ultraviolet limits.} On every fixed, sufficiently
  large torus, the correlation functions of the renormalized observables
  at separated points converge along the \emph{whole} sequence of
  lattice spacings, together with their lattice derivatives. Couplings
  and normalized effective-action readings depend holomorphically on a
  terminal-coupling shift in a disc independent of the cutoff and
  period. Transfer data and correlations have the analytic extensions
  specified in III.13.7 (§4.2--4.4).
\end{itemize}

\textbf{Conventions.}

\begin{itemize}
\tightlist
\item
  \(G\) is a compact matrix group with simple nonabelian real Lie
  algebra (``\(G\) simple'').
\item
  \(b>0\) is the \emph{penalty} of the block-spin kernel. It is the
  Gaussian block-spin coefficient in the linearized kernel identity. Its
  comparison with Balaban's notation is left to §7.
\item
  \textbf{Levels.} RG levels are \(m=0,1,\dots,N\). Level \(N\) (the
  \emph{cutoff}) is the finest lattice, the Wilson end. Level \(0\) is
  the coarsest. The spacings are \(a_0=1\) and \(a_{m+1}=a_m/L_m\) with
  odd block sizes \(L_m\), and \(\sigma_m=\log(1/a_m)\).
\item
  \textbf{Tori.} The period index is \(p\). The torus \(\Lambda_{N,p}\)
  has spacing \(a_N\) and side \(\ell_p\).
\item
  \textbf{Couplings.} A Wilson family \(C\) has a selected RG trajectory
  with inverse couplings \(\beta_m(N,p)\). Its Wilson coupling is
  \(\beta_{N,p}=\beta_N(N,p)\), its value at level \(0\) at the centre
  of the shooting family is the \emph{barrier} \(\bar\beta\), and
  \(\hat\beta_m\) is the reference running coupling.
\item
  \textbf{Labels and observables.}

  \begin{itemize}
  \tightlist
  \item
    \(\mathcal L_d\) is the space of labels of degree \(d\), and
    \(\mathrm{ref}(d)\) is the reference level of degree \(d\).
  \item
    \(O^{N,p}_v(x)\) is the renormalized lattice observable of a label
    \(v\) at a site \(x\) (II.1), and \(\langle\cdot\rangle_{N,p}\) is
    the Wilson expectation at \(\beta_{N,p}\).
  \end{itemize}
\item
  \textbf{Separated configurations.}
  \(\Omega^n_p=\{x\in(\mathbb R^4)^n:\ x_i\not\equiv x_j \bmod \ell_p\mathbb Z^4 \text{ for } i\neq j\}\).
  A test function is a complex-valued \(C^\infty\) function with compact
  support in the set of pairwise distinct configurations in
  \((\mathbb R^4)^n\). It is \emph{admissible at period \(p\)} if its
  closed support lies in \((-\ell_p/2,\ell_p/2)^{4n}\); this support is
  then a compact subset of \(\Omega^n_p\).
\item
  \textbf{Extraction orders.} An extraction order \(f\) is a pair of
  subsequences (mesh indices, period indices) (II.2), and \(f\circ g\)
  is its refinement by \(g\).
\end{itemize}

\emph{Lean conventions:}
\href{https://github.com/yang-mills-group/OpenAI-YM-20261028/blob/3456e5e75cf8d4cd558955a5c7638fba2d68f464/YM/Euclid/Test.lean\#L18}{\texttt{Test}},
\href{https://github.com/yang-mills-group/OpenAI-YM-20261028/blob/3456e5e75cf8d4cd558955a5c7638fba2d68f464/YM/Extraction/Patch.lean\#L31}{\texttt{Extraction.suitable}}.

\textbf{Definition (chessboard bound in the schedule).} For
\(v\in\mathcal L_d\), levels \(m\le N\) and \(A\ge 0\), the
\emph{centered budget} holds if:

\begin{itemize}
\tightlist
\item
  \(\mathrm{ref}(d)\le N\);
\item
  for every site \(o\) of \(\Lambda_{N,p}\), the observable
  \(O^{N,p}_v(o)\) is supported in the reference cell of the level-\(m\)
  reflection-positivity tiling centred at \(o\);
\item
  its chessboard root
  \(\big|\langle\prod_{\text{cells}}\text{reflected copies}\rangle_{N,p}\big|^{1/\sharp\text{cells}}\)
  is at most \(A\).
\end{itemize}

The family has the \textbf{chessboard bound in the schedule} if for
every \(d\) and all \(R,\varepsilon>0\) there is \(A\ge0\) such that,
for every level \(m\ge\mathrm{ref}(d)\) and all sufficiently large
\(N\), the centered budget holds with constant
\(A\,a_m^{-d-\varepsilon}\) for every period \(p\) and every
\(v\in\mathcal L_d\) with \(\|v\|\le R\).

\emph{Lean:}
\href{https://github.com/yang-mills-group/OpenAI-YM-20261028/blob/3456e5e75cf8d4cd558955a5c7638fba2d68f464/YM/Extraction/Center.lean\#L38-L40}{\texttt{Extraction.Wcmain}},
\href{https://github.com/yang-mills-group/OpenAI-YM-20261028/blob/3456e5e75cf8d4cd558955a5c7638fba2d68f464/YM/Extraction/Center.lean\#L21-L23}{\texttt{centeredBudget}}.

\textbf{Proposition III.13.1 (one family for the ultraviolet and
logarithmic packages).} Let \(G\) be simple and \(b>0\). There is a
Wilson family \(C\) with the following properties.

\begin{enumerate}
\def\labelenumi{\arabic{enumi}.}
\tightlist
\item
  Its kernel penalty is \(b\).
\item
  \textbf{\(C\) is calibrated.} There are a full schedule of constants
  (sheet, block sizes, protection, delays, bar, wait, chart ratio,
  prices, post-cost) and a directed history whose trails generate the
  history of \(C\). Every step from level \(m+1\) to level \(m\)
  satisfies the per-step estimates of III.9--III.12.
\item
  \textbf{\(C\) is consistent.}

  \begin{itemize}
  \tightlist
  \item
    The trajectory is the selected one.
  \item
    The local sections are sound.
  \item
    For every \(d\): \(d+3\le J(\mathrm{ref}(d))\), and \(d+1\le\)
    order\(_m\) and \(d\le\) power\(_m\) for every
    \(m\ge\mathrm{ref}(d)\).
  \item
    Stencils fit: \(20(|x_i|+1)\le\) side\((N,p)\) for every stencil
    point \(x\) of the degree-\(d\) local section, every coordinate
    \(i\), and all \(N,p\) with \(\mathrm{ref}(d)\le N\).
  \item
    For each basis label, the placed local section is regular on its
    footprint and has the prescribed derivative bounds with size
    \(\mathrm{columns}_N\,\theta_N^q\), where \(q\) is its factor degree
    and \(\theta_N\) is the level-\(N\) threshold.
  \item
    The transfer data are calibrated from \(\mathrm{ref}(d)\) on.
  \end{itemize}
\item
  \(C\) has the chessboard bound in the schedule.
\item
  \(C\) is \emph{floral}: it is the centre of an analytic shooting
  family (Proposition III.13.7).
\item
  \(C\) satisfies the three entrance conditions of the OPE:

  \begin{itemize}
  \tightlist
  \item
    the shell estimates of III.11 for every \(d,D,h\) with \(D+h=d+4\):
    there are \(t>1\) and \(j(i)\to\infty\) with
    \(\sigma_{j(i)}/\sigma_i\to t\) for which the prescribed shell
    remainder bounds hold, for every choice of labels and extraction
    order;
  \item
    slow growth of the cover sizes,
    \(\mathrm{cs}_{i+1}^2\le2\,\mathrm{cs}_i^2\), with
    \((L_i\,\mathrm{cs}_i^2)^4\cdot(\text{vacuum cost})\le\varepsilon\,\sigma_{i+1}\)
    eventually;
  \item
    a positive auxiliary radius sequence \(\rho_i\) satisfying
    \(\log\rho_i/\sigma_i\to-1\).
  \end{itemize}
\item
  There are coefficient-map bound data \(W\), allowance indices
  \(r:\mathbb N\to\mathbb N\) and one-step transfer data \(D\) for which
  two properties hold. Write \(K\) for the bound function supplied by
  \(W\).

  \begin{itemize}
  \tightlist
  \item
    \textbf{The logarithmic OPE package.} This contains the transfer and
    remainder estimates of III.11, quantitative triangular label steps,
    auxiliary label comparisons, geometric separation bounds, and the
    sup bounds for products and the additional arity bounds of VII.4. It
    also gives \(L_m\le\hat\beta_m\) and columns\(_m\),
    columns\(_{m+1}\le\hat\beta_m\); \(r\to\infty\); and support
    exponent \(100\) at every level.
  \item
    \textbf{Calibrated label steps.} For each \(d,N,p,m\) with
    \(\mathrm{ref}(d)\le m<N\), there is a linear isomorphism \(T\) such
    that, writing \(x_{m+1},x_m\) for the coefficient maps,
    \[\|x_m(v)-T(V_{L_m}x_{m+1}(v))\|_1\le\mathrm{barunit}(L_m)\,\|\mathrm{core}(x_{m+1}(v))\|_1\qquad(v\in\mathcal L_d),\]
    where \(V_L\) multiplies a component of engineering grade \(e\) by
    \(L^{d-e}\). For an input basis vector of engineering grade \(e\)
    and factor degree \(q\), its only output component of engineering
    grade \(\le e\) is itself, multiplied by
    \((\theta_m/\theta_{m+1})^q\). No component of lower factor degree
    occurs. The map commutes with the core projection and lattice
    symmetries, preserves the unit, and satisfies
    \(\|T\|_1,\|T^{-1}\|_1\le K(J(m+1))\).
  \end{itemize}
\item
  \emph{Coverage} holds. There is \(g(k)\to\infty\) with \(g(k)\le k\)
  and \(a_{g(k)}\le\hat\beta_k\,a_k\) for large \(k\) such that, for
  every \(d\) and all large \(k\), three probe panels of depth \(d\) are
  poised at level \(g(k)\) (VII.1).
\item
  For every extraction order \(F\) satisfying the cell hypotheses
  (IV.1), the grid approximation property holds at arbitrarily fine
  fixed levels, for bounded label degrees and the reflection pairs
  allowed in III.11. Complete weak comparison systems exist for every
  finite demand and degree horizon, every positive radius bound, and
  every real-spin structure on the smooth free comparison table.
\item
  For every extraction order \(f\) there are:

  \begin{itemize}
  \tightlist
  \item
    a refinement \(g\);
  \item
    tables \(A\), \((L_k)\) and \(T\), with a bin \(R\) and a derivative
    plan \(E\), forming extraction base data along \(f\circ g\) (IV.1);
  \item
    a seed \(s\) for \((C,f\circ g,T)\) such that the Maxwell rank
    condition holds for \(T\).
  \end{itemize}
\end{enumerate}

\emph{Status:} theorem. The only hypotheses are \(G\) simple and
\(b>0\).

\emph{Lean:}
\href{https://github.com/yang-mills-group/OpenAI-YM-20261028/blob/3456e5e75cf8d4cd558955a5c7638fba2d68f464/YM/SharpGram/Assembly/Flowers.lean\#L75-L130}{\texttt{uv\_log}}.
The predicates, in the order of the items:
\href{https://github.com/yang-mills-group/OpenAI-YM-20261028/blob/3456e5e75cf8d4cd558955a5c7638fba2d68f464/YM/Spec/Closed.lean\#L58-L61}{\texttt{Spec.Calibrated}},
\href{https://github.com/yang-mills-group/OpenAI-YM-20261028/blob/3456e5e75cf8d4cd558955a5c7638fba2d68f464/YM/Source/Construction.lean\#L44-L62}{\texttt{Family.Makes}},
\href{https://github.com/yang-mills-group/OpenAI-YM-20261028/blob/3456e5e75cf8d4cd558955a5c7638fba2d68f464/YM/UVExtension/Oasis/Flora.lean\#L55-L71}{\texttt{Oasis.Floral}},
\href{https://github.com/yang-mills-group/OpenAI-YM-20261028/blob/3456e5e75cf8d4cd558955a5c7638fba2d68f464/YM/MergeProof/West/Ozone.lean\#L16-L19}{\texttt{Acme.Aroma}},
\href{https://github.com/yang-mills-group/OpenAI-YM-20261028/blob/3456e5e75cf8d4cd558955a5c7638fba2d68f464/YM/OPE/Log/Plans.lean\#L29-L33}{\texttt{LogPackage}},
\href{https://github.com/yang-mills-group/OpenAI-YM-20261028/blob/3456e5e75cf8d4cd558955a5c7638fba2d68f464/YM/SharpGram/Matrix/Leading.lean\#L29-L32}{\texttt{TrueSteps}}
with
\href{https://github.com/yang-mills-group/OpenAI-YM-20261028/blob/3456e5e75cf8d4cd558955a5c7638fba2d68f464/YM/SharpGram/Matrix/Leading.lean\#L19-L27}{\texttt{CalibratedStep}},
\href{https://github.com/yang-mills-group/OpenAI-YM-20261028/blob/3456e5e75cf8d4cd558955a5c7638fba2d68f464/YM/SharpGram/Dense/Passes.lean\#L16-L24}{\texttt{Dense.Coverage}},
\href{https://github.com/yang-mills-group/OpenAI-YM-20261028/blob/3456e5e75cf8d4cd558955a5c7638fba2d68f464/YM/Org/Input.lean\#L121-L123}{\texttt{Org.Grill}},
\href{https://github.com/yang-mills-group/OpenAI-YM-20261028/blob/3456e5e75cf8d4cd558955a5c7638fba2d68f464/YM/MergeProof/East/Gran.lean\#L35}{\texttt{East.CompleteWeak}},
\href{https://github.com/yang-mills-group/OpenAI-YM-20261028/blob/3456e5e75cf8d4cd558955a5c7638fba2d68f464/YM/Next/EXTRACTION/Base.lean\#L42-L51}{\texttt{Next.Extract.Base}},
\href{https://github.com/yang-mills-group/OpenAI-YM-20261028/blob/3456e5e75cf8d4cd558955a5c7638fba2d68f464/YM/Spec/Closed.lean\#L20-L26}{\texttt{Spec.Seed}},
\href{https://github.com/yang-mills-group/OpenAI-YM-20261028/blob/3456e5e75cf8d4cd558955a5c7638fba2d68f464/YM/Euclid/Comparator.lean\#L135-L139}{\texttt{Maxwell.Ranks}}.
The engineering rescaling and leading map are
\href{https://github.com/yang-mills-group/OpenAI-YM-20261028/blob/3456e5e75cf8d4cd558955a5c7638fba2d68f464/YM/Bid/Loan.lean\#L63}{\texttt{Bid.vowel}}
and
\href{https://github.com/yang-mills-group/OpenAI-YM-20261028/blob/3456e5e75cf8d4cd558955a5c7638fba2d68f464/YM/Extra/SOURCE_HISTORIES/LocalDebt.lean\#L77-L81}{\texttt{Heading.leads}}.

\textbf{Proposition III.13.2 (the coordinated family).} Let \(G\) be
simple and \(b>0\). There is one Wilson family \(C\) with the following
properties:

\begin{itemize}
\tightlist
\item
  it has kernel penalty \(b\) and is consistent, calibrated and floral;
\item
  it has the chessboard bound in the schedule;
\item
  it has the fixed-volume limits of III.13.5;
\item
  it has the logarithmic OPE package and the calibrated label steps;
\item
  it has coverage;
\item
  it has complete weak comparison systems along every extraction order
  satisfying the cell hypotheses.
\end{itemize}

In addition, for \textbf{every} extraction order \(F\) there are a
refinement \(g\) and a table \(T\) such that:

\begin{itemize}
\tightlist
\item
  \begin{enumerate}
  \def\labelenumi{(\alph{enumi})}
  \tightlist
  \item
    the ordered limit holds for \(T\) along the \emph{full} mesh
    sequence and the period subsequence of \(F\circ g\);
  \end{enumerate}
\item
  (b) the OPE services of VII.4 hold for \((C,F\circ g,T)\);
\item
  (c) a construction record of II.2 for \((C,F\circ g,T)\) exists.
\end{itemize}

\emph{Status:} theorem. The family is chosen before \(F\).

\emph{Lean:}
\href{https://github.com/yang-mills-group/OpenAI-YM-20261028/blob/3456e5e75cf8d4cd558955a5c7638fba2d68f464/YM/SharpGram/Common.lean\#L16-L37}{\texttt{coordination}},
\href{https://github.com/yang-mills-group/OpenAI-YM-20261028/blob/3456e5e75cf8d4cd558955a5c7638fba2d68f464/YM/Extraction/Limits.lean\#L32-L35}{\texttt{Extraction.Order.Ordered}},
\href{https://github.com/yang-mills-group/OpenAI-YM-20261028/blob/3456e5e75cf8d4cd558955a5c7638fba2d68f464/YM/OPE/ElementaryLog.lean\#L12-L13}{\texttt{OPE.Facilities}},
\href{https://github.com/yang-mills-group/OpenAI-YM-20261028/blob/3456e5e75cf8d4cd558955a5c7638fba2d68f464/YM/Spec/Closed.lean\#L64-L73}{\texttt{Spec.Final}}.

\textbf{Proposition III.13.3 (the common family, all phases).} Let \(G\)
be simple, \(b>0\), and let \(f\) be an extraction order. There are a
family \(C\), a refinement \(g\), a table \(T\), a construction record
\(F\) for \((C,f\circ g,T)\) and a reconstruction record \(\mathcal A\)
for every phase of \(F\) (V.3) such that:

\begin{itemize}
\tightlist
\item
  everything listed in III.13.2 holds for \((C,f\circ g,T)\);
\item
  the reconstruction is \emph{completed} for every phase: Wightman
  realization, surjective comparison maps to the Osterwalder--Schrader
  Hilbert spaces, unique vacuum, light-cone and heat-kernel bounds
  uniform over the phases, trivial commutant, the continued stress
  tensor, continuous dependence of boundary correlations on the phase,
  and the barycentre identities;
\item
  the derivative sector of \(F\) has a prescription \(P\) giving the OPE
  estimates for all fixed finite cluster trees and the corresponding
  estimates with separated spectators (VII.4). These estimates allow an
  arbitrary positive power loss.
\end{itemize}

\emph{Status:} theorem. In the statement \(C\) is quantified after
\(f\). The proof takes \(C\) from III.13.2, so in fact one family serves
all \(f\).

\emph{Lean:}
\href{https://github.com/yang-mills-group/OpenAI-YM-20261028/blob/3456e5e75cf8d4cd558955a5c7638fba2d68f464/YM/SharpGram/Common.lean\#L41-L60}{\texttt{common}},
\href{https://github.com/yang-mills-group/OpenAI-YM-20261028/blob/3456e5e75cf8d4cd558955a5c7638fba2d68f464/YM/EveryPhase/Extension.lean\#L51-L71}{\texttt{EveryPhase.Completed}},
\href{https://github.com/yang-mills-group/OpenAI-YM-20261028/blob/3456e5e75cf8d4cd558955a5c7638fba2d68f464/YM/EveryPhase/Pack.lean\#L117-L132}{\texttt{EveryPhase.Almanac}},
\href{https://github.com/yang-mills-group/OpenAI-YM-20261028/blob/3456e5e75cf8d4cd558955a5c7638fba2d68f464/YM/OPE/Subpower.lean\#L39-L45}{\texttt{OPE.AllShapes}}.

\textbf{Proposition III.13.4 (headline: moving-frame Gram limits on the
common family).} Let \(G\) be simple, \(b>0\), and let \(f\) be an
extraction order. There are \(C,g,T,F,\mathcal A,P\) as in III.13.3 with
all of its conclusions. In addition, coverage holds at the cutoff of
\(F\), and the following holds for every ultrafilter \(\mathcal U\) on
\(\mathbb N\times\mathbb N\) finer than the filter of \(f\circ g\)
(``period index large, then mesh index large'') and every depth \(D\):

\begin{enumerate}
\def\labelenumi{\arabic{enumi}.}
\tightlist
\item
  \emph{Moving-frame portfolios} of depth \(D\) along \(\mathcal U\)
  exist. A portfolio consists of:

  \begin{itemize}
  \tightlist
  \item
    levels \(k'=k'(k)\to\infty\) with \(k'\le k\) and
    \(a_k\le a_{k'}\le\hat\beta_k\,a_k\) for all sufficiently large
    \(k\);
  \item
    three probe panels poised at level \(k'\) for all sufficiently large
    \(k\);
  \item
    label matrices having limits along \(\mathcal U\) for all
    sufficiently large \(k\).
  \end{itemize}
\item
  For every portfolio \(X\) there are linear isomorphisms \(Y_k(d)\)
  from the degree-\(d\) field space of the free Maxwell table onto that
  of \(T\) (\(d\le D\), \(k\in\mathbb N\)), with the following
  properties.

  \begin{itemize}
  \tightlist
  \item
    For all sufficiently large \(k\), they respect the grading through
    depth \(D\).
  \item
    For each \(d\le D\), the Gram forms expressed in these bases are
    bounded and coercive uniformly for all sufficiently large \(k\).
  \item
    For each \(d\le D\), there is \(q_d\in\mathbb N\) such that, for
    every sufficiently large \(k\) and every field \(x\), there are
    representatives \(v\) and \(v_{N,p}\to v\) along \(\mathcal U\),
    projecting to \(x\), whose source costs at level \(k\) are
    eventually at most \(\hat\beta_k^{q_d}a_k^d\|Y_k(d)^{-1}x\|\).
    Conversely, there is \(K_d\ge0\) such that every convergent family
    of representatives with source cost at most \(z\) along any
    nontrivial filter finer than \(\mathcal U\) satisfies
    \(\|Y_k(d)^{-1}x\|\le K_da_k^{-d}z\). Both assertions hold for all
    sufficiently large \(k\).
  \item
    For every \(d\le D\), the Gram form of \(X\) at step \(k\), read in
    the basis \(Y_k(d)\), converges as \(k\to\infty\) to the Gram form
    of the free Maxwell model in degree \(d\).
  \end{itemize}
\end{enumerate}

\emph{Status:} theorem; headline declaration. In the statement \(C\) is
quantified after \(f\), as in III.13.3.

\emph{Lean:}
\href{https://github.com/yang-mills-group/OpenAI-YM-20261028/blob/3456e5e75cf8d4cd558955a5c7638fba2d68f464/YM/SharpGram/Measured.lean\#L11-L36}{\texttt{habitat}},
\href{https://github.com/yang-mills-group/OpenAI-YM-20261028/blob/3456e5e75cf8d4cd558955a5c7638fba2d68f464/YM/SharpGram/Weak/Reading.lean\#L35-L49}{\texttt{Weak.Portfolio}},
\href{https://github.com/yang-mills-group/OpenAI-YM-20261028/blob/3456e5e75cf8d4cd558955a5c7638fba2d68f464/YM/SharpGram/InContext.lean\#L29-L37}{\texttt{PricedIn}}.
\emph{(proof of the wrapper read)}

\textbf{Proposition III.13.5 (fixed-volume limit with \(k\)
derivatives).} Let \(C\) be as in III.13.1, and fix \(d,k\in\mathbb N\).
For all sufficiently large periods \(p\) (the threshold depends on \(d\)
and \(k\)), every \(n\), and all labels \(v_1,\dots,v_n\) of degrees
\(\le d\), there is a function \(g\) on \((\mathbb R^4)^n\) such that,
as \(N\to\infty\) along the whole sequence:

\begin{enumerate}
\def\labelenumi{\arabic{enumi}.}
\tightlist
\item
  \(g\) is continuous on \(\Omega^n_p\). The lattice correlation
  functions \(x\mapsto\langle\prod_iO^{N,p}_{v_i}(x_i)\rangle_{N,p}\),
  extended periodically and then as step functions using coordinatewise
  ceiling projection onto \((a_N\mathbb Z^4)^n\), converge to \(g\)
  uniformly on compact subsets of \(\Omega^n_p\).
\item
  For every word \(u\) of at most \(k\) coordinate directions, the
  step-function extensions of the normalized forward differences
  converge to \(\partial_ug\) uniformly on compact subsets of
  \(\Omega^n_p\). The same limit is obtained by multilinear
  interpolation of any fixed word of normalized coordinate differences
  with positive integer step lengths and fixed integer stencil offsets,
  including forward, backward and centred differences.
\item
  On \(\Omega^n_p\) these iterated directional derivatives exist:
  \(\partial_a(\partial_ug)=\partial_{(a,u)}g\) whenever the word \(u\)
  has length \(<k\).
\item
  For every test function \(t\) admissible at period \(p\), \(tg\) is
  integrable and
  \(a_N^{4n}\sum_r t(r)\,\langle\prod_iO^{N,p}_{v_i}(r_i)\rangle_{N,p}\to\int tg\),
  where the sum is over the centred representatives of the torus
  lattice.
\end{enumerate}

\emph{Status:} theorem.

\emph{Lean:}
\href{https://github.com/yang-mills-group/OpenAI-YM-20261028/blob/3456e5e75cf8d4cd558955a5c7638fba2d68f464/YM/UVExtension/Poppy/Fountain.lean\#L16-L21}{\texttt{Poppy.Margins}},
\href{https://github.com/yang-mills-group/OpenAI-YM-20261028/blob/3456e5e75cf8d4cd558955a5c7638fba2d68f464/YM/UVExtension/Poppy/Summit.lean\#L19-L22}{\texttt{Poppy.bravo}},
\href{https://github.com/yang-mills-group/OpenAI-YM-20261028/blob/3456e5e75cf8d4cd558955a5c7638fba2d68f464/YM/UVExtension/Poppy/Fountain.lean\#L96-L108}{\texttt{Poppy.blossoms}}.

\textbf{Proposition III.13.6 (convergence at separated lattice points).}
Let \(C^{\mathbf s}\) be a member of the shooting family of III.13.7. At
a fixed period \(p\), assume:

\begin{itemize}
\tightlist
\item
  its shifts converge, \(s_{N,p}\to s\) as \(N\to\infty\), with
  \(|s|<r_*\);
\item
  \(C^{\mathbf s}\) has the chessboard bound in the schedule;
\item
  the slow-growth, label-transfer and remainder bounds hold (these are
  supplied by the bundle of III.13.7 for every assignment of shifts).
\end{itemize}

Let \(J\in\mathbb N\), let \(u_1,\dots,u_I\) be labels with
\(\mathrm{ref}(\deg u_i)\le J\), and let \(o_1,\dots,o_I\) be pairwise
distinct sites of the level-\(J\) lattice of the period-\(p\) torus.
Then the sequence
\[\Psi_N=\Big\langle\prod_{i=1}^IO^{N,p}_{u_i}(o_i)\Big\rangle_{N,p},\]
with \(o_i\) read as sites of \(\Lambda_{N,p}\), converges as
\(N\to\infty\). \emph{Status:} theorem. For the family of III.13.1,
\(\mathbf s\equiv0\).

\emph{Lean:}
\href{https://github.com/yang-mills-group/OpenAI-YM-20261028/blob/3456e5e75cf8d4cd558955a5c7638fba2d68f464/YM/UVExtension/Hearth/Table.lean\#L87-L112}{\texttt{Hearth.food}}.

\textbf{Proposition III.13.7 (the family is the centre of an analytic
shooting family).} For the family \(C\) of III.13.1 there are:

\begin{itemize}
\tightlist
\item
  a schedule with its per-step data;
\item
  a shooting structure with an arc;
\item
  a label-transfer bundle;
\item
  reference levels \(\mathrm{ref}\);
\end{itemize}

such that the following holds.

\begin{enumerate}
\def\labelenumi{\arabic{enumi}.}
\tightlist
\item
  \emph{Shooting.} There is \(r_*>0\), depending only on the schedule,
  such that for all \(N,p\) and every real \(s\) with \(|s|<r_*\) there
  is exactly one Wilson coupling whose trajectory under the fixed
  recursion stays in the prescribed scalar coupling bands and reaches
  \(\bar\beta+s\) at level \(0\). For an assignment
  \(\mathbf s=(s_{N,p})\) of real shifts with \(|s_{N,p}|<r_*\), let
  \(C^{\mathbf s}\) be the Wilson family with these trajectories and the
  labels, sections and reference levels of the bundle. Then
  \(C=C^{\mathbf 0}\).
\item
  \emph{Analyticity, uniform in the cutoff.} For \(m\le N\), the
  level-\(m\) coupling \(\beta_m(N,p;s)\) and the normalized activity
  readings of the level-\(m\) effective action extend holomorphically to
  \(|s|<r_*\). They are bounded there by \(3\hat\beta_m\) and by
  \(\mathrm{soft}=\mathrm{rim}/8\) respectively. For fixed \(m,p\) they
  converge as \(N\to\infty\), with all derivatives in \(s\), locally
  uniformly on the disc.
\item
  \emph{Bounds.} \(C\) satisfies the calibration package (with constant
  allowance) and the priced sup bounds. For every assignment
  \(\mathbf s\), \(C^{\mathbf s}\) satisfies the label-transfer and
  remainder bounds and has the chessboard bound in the schedule.
  Extraction base data exist at the centre.
\item
  \emph{Analytic transfer data.} There is \(a>0\), with its disc
  contained in the shooting domain, on which the bundle has holomorphic
  coefficient and transfer data for all \(d,N,p\) with
  \(\mathrm{ref}(d)\le N\). In particular, the inverse terminal label
  maps extend holomorphically there.
\item
  \emph{Analytic correlations.} Fix \(d,m\). For all large \(p\), every
  \(n\) and labels of degree \(\le d\), there are \(b'>0\) and functions
  \(W_N(z,x)\), holomorphic in \(|z|<b'\), and \(f(z,\cdot)\) on
  \((\mathbb R^4)^n\) such that:

  \begin{itemize}
  \tightlist
  \item
    \(W_N(s,\cdot)\) is the lattice correlation function of
    \(C^{\mathbf s}\) whenever \(s_{N,p}=s\) (for every \(N\) with
    \(\mathrm{ref}(\deg v_i)\le N\) for all \(i\));
  \item
    for every complex \(z\) in the disc, \(f(z,\cdot)\) has continuous
    iterated coordinate derivatives through order \(m\) on
    \(\Omega^n_p\). The step-function differences and interpolated
    differences converge to these derivatives, with all derivatives in
    \(z\), uniformly on products of compact subsets of the disc and
    \(\Omega^n_p\). At every real shift these limits agree with the
    fixed-volume limit of the constant-shift family;
  \item
    for admissible \(t\), \(z\mapsto\int t\,f(z,\cdot)\) is holomorphic,
    and the smeared lattice correlations converge to it, with all
    derivatives in \(z\), locally uniformly.
  \end{itemize}
\end{enumerate}

\emph{Status:} theorem (part of III.13.1). The full-mesh ordered limits
in §4.5 use the real shooting trajectories and the bounds of item 3;
they do not use the analytic correlation limits of item 5.

\emph{Lean:}
\href{https://github.com/yang-mills-group/OpenAI-YM-20261028/blob/3456e5e75cf8d4cd558955a5c7638fba2d68f464/YM/UVExtension/Oasis/Flora.lean\#L55-L71}{\texttt{Oasis.Floral}},
\href{https://github.com/yang-mills-group/OpenAI-YM-20261028/blob/3456e5e75cf8d4cd558955a5c7638fba2d68f464/YM/UVExtension/Oasis/Flora.lean\#L41-L49}{\texttt{Oasis.Spring}},
\href{https://github.com/yang-mills-group/OpenAI-YM-20261028/blob/3456e5e75cf8d4cd558955a5c7638fba2d68f464/YM/UVExtension/Upland/Plinth.lean\#L33-L38}{\texttt{Upland.Arc}},
\href{https://github.com/yang-mills-group/OpenAI-YM-20261028/blob/3456e5e75cf8d4cd558955a5c7638fba2d68f464/YM/UVExtension/Mast/Rig.lean\#L29-L37}{\texttt{Mast.Cellar}},
\href{https://github.com/yang-mills-group/OpenAI-YM-20261028/blob/3456e5e75cf8d4cd558955a5c7638fba2d68f464/YM/UVExtension/Mast/Rig.lean\#L60}{\texttt{Cellar.yard}},
\href{https://github.com/yang-mills-group/OpenAI-YM-20261028/blob/3456e5e75cf8d4cd558955a5c7638fba2d68f464/YM/UVExtension/Harvest/Barn.lean\#L20-L23}{\texttt{Harvest.Barn}},
\href{https://github.com/yang-mills-group/OpenAI-YM-20261028/blob/3456e5e75cf8d4cd558955a5c7638fba2d68f464/YM/UVExtension/Gale/Flange.lean\#L43-L50}{\texttt{Gale.Circuit}},
\href{https://github.com/yang-mills-group/OpenAI-YM-20261028/blob/3456e5e75cf8d4cd558955a5c7638fba2d68f464/YM/UVExtension/Beach/Under.lean\#L32-L36}{\texttt{Beach.Shoal}},
\href{https://github.com/yang-mills-group/OpenAI-YM-20261028/blob/3456e5e75cf8d4cd558955a5c7638fba2d68f464/YM/UVExtension/Reports/Prefixes.lean\#L35-L39}{\texttt{Beach.Realizes}}.
\emph{(definitions and producing proofs read)}

\subsubsection{2. What it uses}\label{what-it-uses-15}

Inputs from other units.

\begin{itemize}
\tightlist
\item
  \textbf{I.1.} The Wilson measure and partition function, and the
  one-loop slope \(\kappa_G=11\,C_G/(24\pi^2)\) (\(C_G\) the Casimir).
\item
  \textbf{II.1, II.2.} Labels, local sections, renormalized observables;
  tables, extraction orders, seeds and construction records.
\item
  \textbf{III.1.} Kernel choices with penalty \(b\) for which every
  patch branches (unique minimizing frames and predictors, a non-empty
  axial fallback):
  \href{https://github.com/yang-mills-group/OpenAI-YM-20261028/blob/3456e5e75cf8d4cd558955a5c7638fba2d68f464/YM/Branch/Pick.lean\#L137-L142}{\texttt{Axial.viable}},
  \href{https://github.com/yang-mills-group/OpenAI-YM-20261028/blob/3456e5e75cf8d4cd558955a5c7638fba2d68f464/YM/RG/Patch.lean\#L37-L45}{\texttt{KernelChoices}}.
\item
  \textbf{III.2.} The order of choice. The pre-\(L\) constants form the
  sheet; then come the block sizes, the wages, the protection
  \(\theta\), the delays, the bar and the waits, in this order.
\item
  \textbf{III.4 and III.8.} The relay supply
  \href{https://github.com/yang-mills-group/OpenAI-YM-20261028/blob/3456e5e75cf8d4cd558955a5c7638fba2d68f464/YM/MergeProof/East/Gran.lean\#L29-L33}{\texttt{posters}},
  and the chessboard product estimate
  \href{https://github.com/yang-mills-group/OpenAI-YM-20261028/blob/3456e5e75cf8d4cd558955a5c7638fba2d68f464/YM/Push/GRID_APPROXIMATION/Telescope.lean\#L34-L78}{\texttt{GridApprox.tariff}}.
\item
  \textbf{III.7.} The one-loop annulus calculation.
  \href{https://github.com/yang-mills-group/OpenAI-YM-20261028/blob/3456e5e75cf8d4cd558955a5c7638fba2d68f464/YM/Next/STATIONARY/Level/Transition.lean\#L20-L23}{\texttt{HorizCalc}}
  holds for every simple \(G\) by
  \href{https://github.com/yang-mills-group/OpenAI-YM-20261028/blob/3456e5e75cf8d4cd558955a5c7638fba2d68f464/YM/Next/STATIONARY/FullPlanned.lean\#L17-L52}{\texttt{Stationary.pouring}}
  through
  \href{https://github.com/yang-mills-group/OpenAI-YM-20261028/blob/3456e5e75cf8d4cd558955a5c7638fba2d68f464/YM/Next/STATIONARY/Level/Recurse.lean\#L115}{\texttt{Fluent}}.
  This is where the coefficient of \(\log L\) in the marginal flow comes
  from (Lemma 4.3).
\item
  \textbf{III.9--III.12.} The per-step estimates, the selected
  trajectory and the band, and the directed histories, alloys and
  factors. They enter through the step package
  \href{https://github.com/yang-mills-group/OpenAI-YM-20261028/blob/3456e5e75cf8d4cd558955a5c7638fba2d68f464/YM/UVExtension/Case/Trading.lean\#L22-L50}{\texttt{uv\_commerce}}
  of this unit.
\item
  \textbf{III.11.}

  \begin{itemize}
  \tightlist
  \item
    The calibration package
    \href{https://github.com/yang-mills-group/OpenAI-YM-20261028/blob/3456e5e75cf8d4cd558955a5c7638fba2d68f464/YM/MergeProof/Almandine.lean\#L24-L36}{\texttt{Chateau}}.
    It includes a certified schedule, consistency, transfer and
    remainder estimates, the chessboard bound in the schedule, and
    extraction base data, seeds and Maxwell ranks after refinement of
    every extraction order.
  \item
    The priced sup bounds
    \href{https://github.com/yang-mills-group/OpenAI-YM-20261028/blob/3456e5e75cf8d4cd558955a5c7638fba2d68f464/YM/MergeProof/North/Bonus.lean\#L34-L36}{\texttt{Velour}}.
  \item
    The OPE entrance
    \href{https://github.com/yang-mills-group/OpenAI-YM-20261028/blob/3456e5e75cf8d4cd558955a5c7638fba2d68f464/YM/MergeProof/Emanate.lean\#L31-L37}{\texttt{Acme.emanate}}.
  \item
    The weak comparison systems
    \href{https://github.com/yang-mills-group/OpenAI-YM-20261028/blob/3456e5e75cf8d4cd558955a5c7638fba2d68f464/YM/MergeProof/East/Gran.lean\#L38-L51}{\texttt{sponsored}}.
  \item
    The cost bound
    \href{https://github.com/yang-mills-group/OpenAI-YM-20261028/blob/3456e5e75cf8d4cd558955a5c7638fba2d68f464/YM/MergeProof/North/Economist.lean\#L23-L32}{\texttt{Ribbons.laurelE}}.
  \end{itemize}
\item
  \textbf{III.15, III.16.} Cost accounting for labels, and analytic
  tools.
\item
  \textbf{IV.1.} The chessboard bound in the schedule (as a definition),
  the cell hypotheses, extraction base data, the Maxwell ranks, and the
  Arzelà--Ascoli-type lemma
  \href{https://github.com/yang-mills-group/OpenAI-YM-20261028/blob/3456e5e75cf8d4cd558955a5c7638fba2d68f464/YM/Extraction/Tame.lean\#L35-L97}{\texttt{Extraction.growFrom}}.
\item
  \textbf{V.3.} Reconstruction for every phase:
  \href{https://github.com/yang-mills-group/OpenAI-YM-20261028/blob/3456e5e75cf8d4cd558955a5c7638fba2d68f464/YM/EveryPhase/Extension.lean\#L75-L80}{\texttt{all\_characters}}.
\item
  \textbf{VII.1.}

  \begin{itemize}
  \tightlist
  \item
    Coverage, through
    \href{https://github.com/yang-mills-group/OpenAI-YM-20261028/blob/3456e5e75cf8d4cd558955a5c7638fba2d68f464/YM/SharpGram/Dense/Passes.lean\#L26-L33}{\texttt{Dense.covered}}.
  \item
    The slim-window port
    \href{https://github.com/yang-mills-group/OpenAI-YM-20261028/blob/3456e5e75cf8d4cd558955a5c7638fba2d68f464/YM/SharpGram/Dense/Slender.lean\#L127-L143}{\texttt{slim\_wait\_for}}.
  \item
    The calibrated label steps.
  \item
    The portfolios
    \href{https://github.com/yang-mills-group/OpenAI-YM-20261028/blob/3456e5e75cf8d4cd558955a5c7638fba2d68f464/YM/SharpGram/Weak/Reading.lean\#L66-L67}{\texttt{fromCoverage}}.
  \end{itemize}
\item
  \textbf{VII.4.}

  \begin{itemize}
  \tightlist
  \item
    The logarithmic package and the OPE services.
  \item
    The prescriptions
    \href{https://github.com/yang-mills-group/OpenAI-YM-20261028/blob/3456e5e75cf8d4cd558955a5c7638fba2d68f464/YM/OPE/Splitting/Legacy.lean\#L20-L31}{\texttt{Prescription.chosen}}
    and the shapes.
  \item
    The delay raises
    \href{https://github.com/yang-mills-group/OpenAI-YM-20261028/blob/3456e5e75cf8d4cd558955a5c7638fba2d68f464/YM/OPE/Rough/Gate.lean\#L45-L54}{\texttt{OPE.halo\_delay}}
    and
    \href{https://github.com/yang-mills-group/OpenAI-YM-20261028/blob/3456e5e75cf8d4cd558955a5c7638fba2d68f464/YM/OPE/Log/Plans.lean\#L19-L25}{\texttt{OPE.Log.aux\_delay}}.
  \end{itemize}
\item
  \textbf{VIII.1.} The construction record and
  \href{https://github.com/yang-mills-group/OpenAI-YM-20261028/blob/3456e5e75cf8d4cd558955a5c7638fba2d68f464/YM/MergeProof/Summit.lean\#L17-L30}{\texttt{Acme.glistens}},
  which turns the UV package and the calibration into a record.
\end{itemize}

Standard mathematics used.

\begin{itemize}
\tightlist
\item
  \textbf{Balaban's block-spin renormalization} for lattice gauge theory
  (Balaban, CMP 95--122, 1984--89; Dimock, \emph{The renormalization
  group according to Balaban I--III}). The convergence statements of
  §4.2--4.4 use the per-step estimates together with stability of the
  chosen recursion.
\item
  \textbf{Shooting by contraction} (Banach fixed point), uniform in the
  cutoff.
\item
  \textbf{Vitali's convergence theorem} (Vitali--Porter). Locally
  bounded holomorphic functions that converge on a real segment converge
  locally uniformly, with all derivatives. The Lean proof maps into
  \(\ell^\infty(\mathbb N;Y)\) and applies the identity principle to the
  closed subspace of convergent sequences.
\item
  \textbf{Arzelà--Ascoli}, in the form: asymptotic equicontinuity on
  compacts together with convergence on a dense set gives locally
  uniform convergence.
\item
  \textbf{Dominated convergence.}
\item
  \textbf{Chessboard estimates} for reflection-positive measures
  (Fröhlich--Israel--Lieb--Simon 1978; Glimm--Jaffe, ch.~10).
\item
  \textbf{Ultrafilters.}
\end{itemize}

\textbf{Imported constants and assumptions.} The statements of
III.13.1--III.13.4 have no hypotheses beyond \(G\) simple and \(b>0\).
The one-loop slope \(\kappa_G\) enters the coupling-shift estimate
(Lemma 4.3) through the hypothesis the annulus calculation, which III.7
proves for every simple \(G\). The construction fixes the following
constants:

\begin{itemize}
\tightlist
\item
  the stage data: first stage \(4\), Taylor order \(6+i\), power
  \(2+i\), and support exponent \(100\) at level \(i\);
\item
  the allowances \(Q\equiv1\) (for the floral property) and
  \(Q_2(j)=j+1\) (for the OPE iteration);
\item
  auxiliary data \((40,1,1)\);
\item
  a bar parameter \(2\);
\item
  the target rate \(-7/16\) of the factor.
\end{itemize}

The analytic argument uses \(\mathrm{soft}=\mathrm{rim}/8<1\),
\(\mathrm{rad}=\mathrm{rim}/4\) and the gain
\(e^{(\sigma_k-\sigma_m)/8}\).

\subsubsection{3. Strategy}\label{strategy-15}

\textbf{The order of choice.} Each estimate of III.9--III.12, and each
requirement of the later parts (OPE with logarithmic corrections,
coverage, weak comparison), has the form ``there is a threshold beyond
which the property holds''. It is also stable under further raising: in
the preorders on delays, waits and wages, and under decreasing the
protection. Successive raises meet the finitely many requirements while
preserving this order: sheet → block sizes → wages → protection → delays
→ bar → wait.

Increasing stage parameters (Taylor order \(6+i\), power \(2+i\))
accommodate the countably many label degrees: degree \(d\) is required
only from its fixed reference level \(\mathrm{ref}(d)\) onwards.

\textbf{Shooting and convergence at a fixed level.}

\begin{enumerate}
\def\labelenumi{\arabic{enumi}.}
\tightlist
\item
  The family is the member \(s=0\) of a one-parameter family in which
  the terminal coupling at level \(0\) is \(\bar\beta+s\). For each
  cutoff \(N\) the trajectory is found by shooting: it is the fixed
  point of a contraction with rate \(1/2\), uniformly in \(N\).
\item
  This uniformity, the convergence of the background fields and coupling
  shifts at each level (Lemma 4.7), and an \(\mathrm{soft}^v\) decay of
  the influence of levels far above \(m\) together give convergence, as
  \(N\to\infty\) along the whole sequence, of the coupling and the
  effective action at any fixed level \(m\), for every real \(s\) in the
  shooting disc (Lemma 4.8).
\item
  Holomorphic extension in \(s\) to a disc independent of \(N\) and
  \(p\), with uniform bounds, and Vitali's theorem upgrade this to
  locally uniform convergence of all derivatives in \(s\) on the disc
  (Lemma 4.9).
\end{enumerate}

\textbf{From the effective action to correlations.} A correlation of
renormalized observables at distinct points is handled in three steps.

\begin{enumerate}
\def\labelenumi{\arabic{enumi}.}
\tightlist
\item
  \textbf{RG identity.} For a function of the level-\(k\) variables, its
  conditional average through the block-spin kernels has Wilson
  expectation equal to its expectation in the level-\(k\) effective
  measure. That measure converges as \(N\to\infty\) (Lemmas 4.11--4.12).
\item
  \textbf{Coarse approximation.} Each fine renormalized observable is
  close to the conditional average of a level-\(k\) observable whose
  coefficients converge as \(N\to\infty\). The error is measured in the
  chessboard norm and is small for \(k\) large, because the gain
  \(e^{(\sigma_k-\sigma_m)/8}\) beats the subexponential growth of the
  costs (Lemma 4.13).
\item
  \textbf{Product estimate.} The chessboard product inequality turns
  these one-point approximations into an approximation of the product
  (Lemma 4.14).
\end{enumerate}

A Cauchy argument then gives Proposition III.13.6. Proposition III.13.5
follows by Arzelà--Ascoli: there is convergence on the dense set of grid
points of all levels, and there is equicontinuity from uniform bounds on
lattice differences of order \(k+1\). Those bounds need a non-degeneracy
(``slicing'') of the label snapshots, which is known only for large
periods, so the fixed-volume statement holds for \(p\) large. The
reports then feed this family to the OPE entrance, reconstruction and
Gram machinery defined in other units.

Balaban's estimates are used only through the per-step package; what
this unit adds is the analysis just described and the order of choice.

\subsubsection{4. Steps}\label{steps-15}

\paragraph{4.1 Choosing the constants
once}\label{choosing-the-constants-once}

\textbf{Lemma 4.1 (the step package with all later constants raised).}
Let \(G\) be simple. Fix the purchased sheet, prices, numerical bounds,
scope, tariffs and block-size floor of Lemma 4.2, and assume that, for
every stage pair \(i\le j\le i+1\) with \(i\) beyond the first stage and
every \(d+3\le i\), the local-section and nonconstant-label expansion
costs are bounded by the supplied budget at stage \(j\). Let
\(PB,P_2:\mathbb N\to\mathbb N\) be arbitrary allowance sequences. There
is a floor \(f\ge v_0\) on the block sizes such that, for every
block-size sequence \(L\ge f\), the following holds.

\begin{itemize}
\tightlist
\item
  Block-size rolls \(D\) exist.
\item
  For every roll \(D\) there is a wage threshold \(U\). For every wage
  \(w\ge U\), postage data, exchange data and insertion tariff, there is
  a post-cost \(P\).
\item
  For every sequence \(\theta_0:\mathbb N\to\mathbb R\) with
  \(\theta_0(i)>0\) for all \(i\), there is a protection sequence
  \(\theta\le\theta_0\).
\item
  For all menus, auxiliary data and delays \(T_0\) there is an ordered
  delay \(V\ge T_0\) such that the following holds for every delay
  \(Y\ge V\), every bar, every wait \(I\) and every level system \(a\)
  paid by \(I\) whose spacing schedule has the divisor property. There
  are a directed history \(B\) on the given level system \(a\) and a set
  of per-step data \(R\) for this schedule, with the prescribed wages,
  exchange, menu and auxiliary data, postage rates bounded by the
  requested rates, the prepared insertion data, and:

  \begin{itemize}
  \tightlist
  \item
    label-transfer bounds and connected-activity bounds for iteration;
  \item
    target rates \(\le-7/16\);
  \item
    repeated-insertion bounds for both allowances;
  \item
    startup and smallness data for shooting;
  \item
    continuity bounds for a single RG step;
  \item
    fixed-level limits of the kinematic coupling shifts and paths;
  \item
    covariance comparison data.
  \end{itemize}
\end{itemize}

\emph{Sketch:} The per-step estimates are supplied in the required order
of choice. Successive delay raises add the later requirements. Each
raise is a lemma ``there is \(V'\ge V\) such that every \(Y\ge V'\) has
the property''. Raise the entry-coupling threshold pointwise by taking
its maximum with the shooting threshold. Every subsequent delay increase
preserves the preceding estimates.

\emph{Lean:}
\href{https://github.com/yang-mills-group/OpenAI-YM-20261028/blob/3456e5e75cf8d4cd558955a5c7638fba2d68f464/YM/SharpGram/Assembly/Expenditures.lean\#L31-L79}{\texttt{multitude}},
\href{https://github.com/yang-mills-group/OpenAI-YM-20261028/blob/3456e5e75cf8d4cd558955a5c7638fba2d68f464/YM/UVExtension/Case/Trading.lean\#L22-L50}{\texttt{uv\_commerce}}.
\emph{(proof of \texttt{multitude} read; \texttt{uv\_commerce} read only
as far as its call to \texttt{Roof.uv\_timberT})}

\textbf{Lemma 4.2 (the purchased sheet).} Let \(G\) be simple, \(K\)
kernel choices for which every patch branches, and \(v\) stage data
satisfying the following conditions:

\begin{itemize}
\tightlist
\item
  first stage \(\ge4\);
\item
  Taylor orders monotone and \(\ge\max(6,i)\);
\item
  powers monotone and \(\ge\max(2,i)\);
\item
  support exponents monotone and positive.
\end{itemize}

For every two allowance sequences there are a sheet with kernel \(K\)
and stage data \(v\), prices, a budget sequence and a bound \(W\),
chosen before the tariffs and scopes. For every tariff and scope there
is then a block-size floor above the sheet floor such that Lemma 4.1
applies to every larger block-size sequence.

\emph{Sketch:} Choose the sheet using three inputs:

\begin{itemize}
\tightlist
\item
  the kinematic constant issuer: pre-\(L\) constants chosen so that,
  once \(L\) is revealed, only waiting is needed and the target
  estimates hold;
\item
  the horizon theorem: horizon data for every sheet with kernel \(K\);
\item
  the kinematic convergence of Lemma 4.7.
\end{itemize}

Then apply Lemma 4.1.

\emph{Lean:}
\href{https://github.com/yang-mills-group/OpenAI-YM-20261028/blob/3456e5e75cf8d4cd558955a5c7638fba2d68f464/YM/SharpGram/Assembly/Expenditures.lean\#L84-L93}{\texttt{multitude\_foot}},
\href{https://github.com/yang-mills-group/OpenAI-YM-20261028/blob/3456e5e75cf8d4cd558955a5c7638fba2d68f464/YM/UVExtension/Case/Wing.lean\#L95-L154}{\texttt{sailsV}},
\href{https://github.com/yang-mills-group/OpenAI-YM-20261028/blob/3456e5e75cf8d4cd558955a5c7638fba2d68f464/YM/UVExtension/Reward/Clocks.lean\#L43-L108}{\texttt{office}},
\href{https://github.com/yang-mills-group/OpenAI-YM-20261028/blob/3456e5e75cf8d4cd558955a5c7638fba2d68f464/YM/UVExtension/Free/Lease.lean\#L171-L182}{\texttt{profiles}}.
\emph{(proof of \texttt{multitude\_foot} read; the others statement
only)}

\textbf{Lemma 4.3 (the one-loop running of the marginal coupling).} Let
\(G\) be simple. Fix kernel choices, a stage \(J\), a colour direction,
positive tolerances \(\varepsilon,w,\delta\), and nonnegative
\(D,Z,C,Z'\). There are \(\kappa\ge0\) and \(s>0\) with the following
property. For all \(\tau,\mathrm{pace}>0\) and \(A\ge0\) satisfying the
annulus calculation of III.7 with parameters
\((G,w,\varepsilon,D,Z,\tau,A)\), there are:

\begin{itemize}
\tightlist
\item
  \(d\ge0\) and \(\gamma>0\);
\item
  an integer \(u\);
\item
  functions \(\beta(L)\) and \(\Lambda_0(L)\);
\end{itemize}

such that every sequence of step data with stage \(J\), block size
\(L\ge u\), diverging coarse side lengths, and the comparison hypotheses
of III.7 admits the following data. Those hypotheses include input and
output covariance strip bounds, input path bounds, admissible forests
and charts, spectral gap and criticality estimates, heat bounds, direct
clip bounds \(Z,Z'\), a positive gap parameter and amplitude
\(\Lambda\ge1\). The output consists of:

\begin{itemize}
\item
  background path data \(Y\);
\item
  a number \(\alpha\), the shift of the marginal coupling read from the
  actual quadratic form of the step in a colour direction, with
  \[|\alpha+\kappa_G\log L|+1\le d;\]
\item
  equality of the classical Taylor coefficients through order \(4\), and
  a finite-volume thermal comparison estimate. For every
  \(\Theta\ge\beta(L)\) and \(c'\le\gamma\), on a coarse torus with
  \(\nu\) sites in each coordinate, the second derivatives differ by at
  most \(\nu^4\Theta e^{-c'\nu}\) on the prescribed unit-bounded wave
  probes. The output path has its announced bounds whenever
  \(0<\eta'\le s\), \(A'\ge\kappa\) and
  \(\Lambda'\ge\max(\Lambda_0(L),\mathrm{pace}\,\Lambda)\).
\end{itemize}

The annulus calculation supplies a comparison quadratic form whose
coupling shift differs from \(-\kappa_G\log L\) by a constant
independent of \(L\). Its conclusion applies to all odd \(L\) beyond a
threshold under the input path, covariance, floor and compatibility
assumptions of III.7.

\emph{Sketch:} \(\alpha\) is compared with the coupling shift of the
annulus comparison form, and the two bounds are added. The constant
\(d\) is chosen before \(L\), so the error is bounded uniformly in
\(L\).

\emph{Lean:}
\href{https://github.com/yang-mills-group/OpenAI-YM-20261028/blob/3456e5e75cf8d4cd558955a5c7638fba2d68f464/YM/UVExtension/Reward/Heal.lean\#L24-L114}{\texttt{ceilings}},
\href{https://github.com/yang-mills-group/OpenAI-YM-20261028/blob/3456e5e75cf8d4cd558955a5c7638fba2d68f464/YM/UVExtension/Reward/Choices.lean\#L137-L178}{\texttt{colorReady}},
\href{https://github.com/yang-mills-group/OpenAI-YM-20261028/blob/3456e5e75cf8d4cd558955a5c7638fba2d68f464/YM/Next/STATIONARY/Level/Transition.lean\#L20-L23}{\texttt{HorizCalc}}.
\emph{(the part of the proof of \texttt{ceilings} that bounds \(\alpha\)
read; \texttt{colorReady} statement only)}

\textbf{Lemma 4.4 (the bundle and the joint offers).} Let \(G\) be
simple, \(K\) kernel choices for which every patch branches, \(v\)
satisfying the stage conditions of Lemma 4.2 and having support
exponents \(\ge100\), and \(Q,Q_2:\mathbb N\to\mathbb N\) monotone with
values \(\ge1\). There are:

\begin{itemize}
\tightlist
\item
  a sheet with kernel \(K\) and stage data \(v\);
\item
  prices, a budget sequence and a bound \(K'\);
\item
  a floor \(\tau\ge\mathrm{floor}_L\);
\end{itemize}

such that for every block-size sequence \(k\ge\tau\) rolls exist, and
the following holds through the quantifier chain of Lemma 4.1. In that
chain the waits receive one more raise \(Z\), of readiness, minimal
length and window. Every level system paid by this raised wait and
satisfying the divisor property then carries:

\begin{itemize}
\tightlist
\item
  a directed history and a factor;
\item
  a shooting structure \(\mathcal S\) with an arc;
\item
  a label-transfer bundle over \(\mathcal S\);
\end{itemize}

and the family \(C^{\mathbf 0}\) built from the bundle has the following
properties:

\begin{itemize}
\tightlist
\item
  the calibration package for both \(Q\) and \(Q_2\);
\item
  slow growth of the cover sizes;
\item
  the priced sup bounds (three variants);
\item
  the bundle conditions of III.13.7.3;
\item
  calibrated label steps;
\item
  the leases, and target rates \(\le-7/16\).
\end{itemize}

In addition:

\begin{itemize}
\tightlist
\item
  there is a radius \(r>0\) such that the disc of radius \(r\) lies in
  the analytic domain and the transfer data extend holomorphically to it
  for all \(d,N,p\) with \(\mathrm{ref}(d)\le N\);
\item
  complete weak comparison systems exist for every extraction order.
\end{itemize}

\emph{Sketch:} Lemma 4.2 supplies the schedule. Increase the block-size
floor to ensure
\(\mathrm{rate}(L)(1+2\,\mathrm{columns})\,\mathrm{orbitCost}\le\mathrm{barunit}(L)\),
then construct the bundle. The relay supply and a further wait increase
give weak comparison systems. The leases and target rates give the
common disc for holomorphic transfer.

\emph{Lean:}
\href{https://github.com/yang-mills-group/OpenAI-YM-20261028/blob/3456e5e75cf8d4cd558955a5c7638fba2d68f464/YM/SharpGram/Assembly/SameFactory.lean\#L60-L137}{\texttt{havens\_log}},
\href{https://github.com/yang-mills-group/OpenAI-YM-20261028/blob/3456e5e75cf8d4cd558955a5c7638fba2d68f464/YM/SharpGram/Assembly/Flowers.lean\#L36-L66}{\texttt{jointOffers}},
\href{https://github.com/yang-mills-group/OpenAI-YM-20261028/blob/3456e5e75cf8d4cd558955a5c7638fba2d68f464/YM/UVExtension/Gale/Flange.lean\#L57-L67}{\texttt{Gale.sheaf}}.
\emph{(proof of \texttt{jointOffers} read; statement and opening of
\texttt{havens\_log} read)}

\textbf{Lemma 4.5 (instantiation).} Proposition III.13.1 follows from
Lemma 4.4 with the following choices.

\begin{itemize}
\tightlist
\item
  \textbf{Kernel.} The kernel choices have penalty \(b\) (III.1). The
  stage data are \((4,\,6+i,\,2+i,\,100)\), and the allowances are
  \(Q\equiv1\), \(Q_2(j)=j+1\).
\item
  \textbf{Block sizes.} Take the maximum of the floor of Lemma 4.4 and
  the relay floor.
\item
  \textbf{Post-\(L\) data.} The menu is canonical, the auxiliary data
  are \((40,1,1)\), the wage is the threshold itself, and
  \(\theta_0\equiv1\).
\item
  \textbf{Delays.} Zero delays are raised by the geometric delay bound
  and the logarithmic delay bound; the latter gives \(L_m\), columns
  \(\le\hat\beta_m\).
\item
  \textbf{Bar and wait.} A bar above the delays with parameter \(2\) is
  chosen. A wait with slim window is thickened by the slim-window
  theorem so that it carries a certified station port, then joined with
  the wait of Lemma 4.4. The schedule existence theorem gives a paid
  level system.
\end{itemize}

Then:

\begin{itemize}
\tightlist
\item
  coverage follows from the certified station port;
\item
  floral from Lemma 4.18;
\item
  the OPE entrance conditions from the shell estimate, the slow growth
  of cover sizes and the radius asymptotics;
\item
  the logarithmic package is assembled from the label-transfer and
  remainder bounds, the priced bounds, the calibrated steps, the budget,
  \(r\to\infty\) (from the depth of the fits) and support exponent
  \(100\).
\end{itemize}

\emph{Lean:}
\href{https://github.com/yang-mills-group/OpenAI-YM-20261028/blob/3456e5e75cf8d4cd558955a5c7638fba2d68f464/YM/SharpGram/Assembly/Flowers.lean\#L75-L130}{\texttt{uv\_log}},
\href{https://github.com/yang-mills-group/OpenAI-YM-20261028/blob/3456e5e75cf8d4cd558955a5c7638fba2d68f464/YM/Branch/Pick.lean\#L137-L142}{\texttt{Axial.viable}}.

\paragraph{4.2 Shooting, and the limit at a fixed
level}\label{shooting-and-the-limit-at-a-fixed-level}

\textbf{Lemma 4.6 (shooting).} Let \(r_*>0\) be the radius of the arc;
it depends only on the schedule. For all \(N,p\) and real \(|s|<r_*\):

\begin{itemize}
\tightlist
\item
  there is exactly one \(\beta\in\mathbb R\) for which the fixed startup
  and step recursion has a trajectory staying in all prescribed scalar
  bands and satisfying \(\beta_N=\beta\), \(\beta_0=\bar\beta+s\);
\item
  the realized trajectory satisfies the band and step conditions at
  every level, and its level-\(0\) coupling is \(\bar\beta+s\);
\item
  \(\beta\) is the value at \(s\) of the holomorphic function
  \(\beta_N(N,p;\cdot)\).
\end{itemize}

\emph{Lean:}
\href{https://github.com/yang-mills-group/OpenAI-YM-20261028/blob/3456e5e75cf8d4cd558955a5c7638fba2d68f464/YM/UVExtension/Upland/Plinth.lean\#L71-L75}{\texttt{Arc.exists\_trial}},
\href{https://github.com/yang-mills-group/OpenAI-YM-20261028/blob/3456e5e75cf8d4cd558955a5c7638fba2d68f464/YM/UVExtension/Oasis/Flora.lean\#L51-L53}{\texttt{Oasis.spring}},
\href{https://github.com/yang-mills-group/OpenAI-YM-20261028/blob/3456e5e75cf8d4cd558955a5c7638fba2d68f464/YM/UVExtension/Spire/Spark.lean\#L37-L41}{\texttt{Spire.Uses}}.
\emph{(the wrappers read; existence and uniqueness come from the fitted
structure, statement only)}

\textbf{Lemma 4.7 (kinematic convergence at each level).} Fix a paid
level system, the horizon and chart data, a colour direction, the small
positive pulse parameter, the one-step kinematic estimates and the
limiting covariance data supplied by Lemma 4.2. There are background
paths and coupling shifts \(\alpha_N(m)\) satisfying all kinematic step
estimates such that, for every level \(m\):

\begin{itemize}
\tightlist
\item
  \(\alpha_N(m)\) converges as \(N\to\infty\);
\item
  there are a limiting background path and \(w_m\) strictly larger than
  the level-\(m\) width such that, for every \(u<w_m\), the path
  coefficients converge in the exponentially weighted sum norm of width
  \(u\). Their path supports agree with those of the limit for all
  sufficiently large \(N\).
\end{itemize}

\emph{Lean:}
\href{https://github.com/yang-mills-group/OpenAI-YM-20261028/blob/3456e5e75cf8d4cd558955a5c7638fba2d68f464/YM/UVExtension/Reward/Lights.lean\#L74-L78}{\texttt{Reward.work\ item}},
\href{https://github.com/yang-mills-group/OpenAI-YM-20261028/blob/3456e5e75cf8d4cd558955a5c7638fba2d68f464/YM/UVExtension/Reward/Tides.lean\#L103-L116}{\texttt{Reward.turning}}.
\emph{(proof of \texttt{work\ item} read; \texttt{turning} statement
only)}

\textbf{Lemma 4.8 (ultraviolet convergence at a fixed level, real
shifts).} Fix the shooting structure and assume its one-step continuity
property and the kinematic convergence of Lemma 4.7; both are supplied
by the arc. For every real \(s\) in the analytic domain, every level
\(m\) and every period \(p\), two quantities converge as \(N\to\infty\)
along the whole sequence: the effective activities at level \(m\), and
the level-\(m\) coupling data of the trajectory ending at
\(\bar\beta+s\).

\emph{Sketch:}

\begin{enumerate}
\def\labelenumi{\arabic{enumi}.}
\tightlist
\item
  \textbf{Contraction.} The trajectory with cutoff \(N\) is the fixed
  point of a map on coupling sequences. Its \(k\)-th iterate from \(0\)
  lies within \(\mathrm{rad}\cdot2^{1-k}\) of the fixed point, uniformly
  in \(N\).
\item
  \textbf{Stability.} If the level-\((m+1)\) data converge as
  \(N\to\infty\), so do the level-\(m\) data and the claimed coupling
  shift of the step. By induction on \(k\), every iterate therefore
  converges level by level. The coupling shifts converge by Lemma 4.7.
\item
  \textbf{Decay of far-ultraviolet influence.} The level-\(m\)
  activities depend on the levels above \(m+v\) only up to
  \(\mathrm{soft}^v\) with \(\mathrm{soft}<1\), uniformly in \(N\).
\item
  \textbf{Conclusion.} An \(\varepsilon/4\) argument closes each step:
  in a complete space, a sequence that can be approximated uniformly by
  convergent sequences converges.
\end{enumerate}

\emph{Lean:}
\href{https://github.com/yang-mills-group/OpenAI-YM-20261028/blob/3456e5e75cf8d4cd558955a5c7638fba2d68f464/YM/UVExtension/Camp/Pear.lean\#L70-L71}{\texttt{Camp.artistry}},
\href{https://github.com/yang-mills-group/OpenAI-YM-20261028/blob/3456e5e75cf8d4cd558955a5c7638fba2d68f464/YM/UVExtension/Camp/Pear.lean\#L62-L68}{\texttt{Camp.climb}},
\href{https://github.com/yang-mills-group/OpenAI-YM-20261028/blob/3456e5e75cf8d4cd558955a5c7638fba2d68f464/YM/UVExtension/Camp/Pear.lean\#L31-L35}{\texttt{Camp.income}},
\href{https://github.com/yang-mills-group/OpenAI-YM-20261028/blob/3456e5e75cf8d4cd558955a5c7638fba2d68f464/YM/UVExtension/Camp/Buckles.lean\#L40-L41}{\texttt{Camp.Sign}},
\href{https://github.com/yang-mills-group/OpenAI-YM-20261028/blob/3456e5e75cf8d4cd558955a5c7638fba2d68f464/YM/UVExtension/Camp/Buckles.lean\#L43-L49}{\texttt{Camp.trim}}.
\emph{(proof read; the contraction rate and the \(\mathrm{soft}^v\)
decay are taken from statements)}

\textbf{Lemma 4.9 (Vitali, and analyticity uniform in the cutoff).} Let
\(r>0\) and let \(F_N\) be holomorphic on \(|z|<r\) with values in a
complex Banach space, uniformly bounded, and convergent at every real
\(z\in(-r,r)\). Then there is a holomorphic \(f\) on the disc such that
every iterated derivative of \(F_N\) converges to that of \(f\), locally
uniformly. The lemma applies to \(\beta_m(N,p;\cdot)\) (bounded by
\(3\hat\beta_m\)) and to the normalized level-\(m\) activity readings
(bounded by \(\mathrm{soft}\)). With Lemma 4.8 this gives III.13.7.2.

\emph{Lean:}
\href{https://github.com/yang-mills-group/OpenAI-YM-20261028/blob/3456e5e75cf8d4cd558955a5c7638fba2d68f464/YM/UVExtension/Camp/Veto.lean\#L64-L69}{\texttt{Camp.viva}},
\href{https://github.com/yang-mills-group/OpenAI-YM-20261028/blob/3456e5e75cf8d4cd558955a5c7638fba2d68f464/YM/UVExtension/Camp/Veto.lean\#L62}{\texttt{Camp.SummitUV}},
\href{https://github.com/yang-mills-group/OpenAI-YM-20261028/blob/3456e5e75cf8d4cd558955a5c7638fba2d68f464/YM/UVExtension/Upland/Plinth.lean\#L62-L70}{\texttt{Arc.uv}}.

\paragraph{4.3 Renormalized observables}\label{renormalized-observables}

\textbf{Lemma 4.10 (label transfer).} Fix the shooting structure,
coefficient bundle and reference levels of Lemma 4.4.

\begin{itemize}
\tightlist
\item
  \textbf{Coefficient data.} For each degree \(d\), period \(p\), level
  \(i\), cutoff \(N\) and shift \(s\), the bundle provides the
  coefficient data of the degree-\(d\) renormalized observables at level
  \(i\).
\item
  \textbf{Bounds.} For every assignment of shifts, degree \(d\) and
  level \(m\ge\mathrm{ref}(d)\), the coefficient histories satisfy the
  supplied transfer and remainder bounds. Also
  \(d+3\le J(\mathrm{ref}(d))\).
\item
  \textbf{Continuity in the shift.} If the shifts at period \(p\)
  converge to \(s\), the level-\(i\) data converge for
  \(i\ge\mathrm{ref}(d)\).
\item
  \textbf{Transfer data.} The transfer data of \(C^{\mathbf s}\) are
  read from these coefficient data and are calibrated from
  \(\mathrm{ref}(d)\) on.
\item
  \textbf{Composition.} In a selected history with sound local sections,
  bounds on the coefficient maps and their inverses, and compatible
  one-step transfer estimates at every \(r\le m<N\), a valid top
  coefficient at level \(N\) extends to a multiscale history down to
  \(r\). It retains that top coefficient and satisfies the one-step
  transfer, remainder and matching identities. The one-step construction
  transports a prescribed coefficient; downward induction composes it
  across the required levels.
\end{itemize}

\emph{Lean:}
\href{https://github.com/yang-mills-group/OpenAI-YM-20261028/blob/3456e5e75cf8d4cd558955a5c7638fba2d68f464/YM/UVExtension/Mast/Rig.lean\#L29-L37}{\texttt{Mast.Cellar}},
\href{https://github.com/yang-mills-group/OpenAI-YM-20261028/blob/3456e5e75cf8d4cd558955a5c7638fba2d68f464/YM/UVExtension/Mast/Rig.lean\#L59}{\texttt{Cellar.gate}},
\href{https://github.com/yang-mills-group/OpenAI-YM-20261028/blob/3456e5e75cf8d4cd558955a5c7638fba2d68f464/YM/UVExtension/Lacing/Circles.lean\#L13-L47}{\texttt{Lacing.compose}},
\href{https://github.com/yang-mills-group/OpenAI-YM-20261028/blob/3456e5e75cf8d4cd558955a5c7638fba2d68f464/YM/UVExtension/Lacing/Reprise.lean\#L45-L258}{\texttt{repriseQ}},
\href{https://github.com/yang-mills-group/OpenAI-YM-20261028/blob/3456e5e75cf8d4cd558955a5c7638fba2d68f464/YM/UVExtension/Lacing/Chunk.lean\#L34-L138}{\texttt{acclaim}},
\href{https://github.com/yang-mills-group/OpenAI-YM-20261028/blob/3456e5e75cf8d4cd558955a5c7638fba2d68f464/YM/UVExtension/Lacing/Couplet.lean\#L57-L89}{\texttt{antilles}}.
\emph{(\texttt{Cellar} definition read; the Lacing lemmas statement
only)}

\paragraph{4.4 Correlations at fixed
volume}\label{correlations-at-fixed-volume}

\textbf{Lemma 4.11 (the RG identity).} Let the history be selected,
\(m\le N\), and \(F\) a bounded measurable function of the level-\(m\)
configuration. Let \(B_{N\to m}F\) be its conditional average obtained
by iterating the block-spin kernels from level \(m\) to level \(N\).
Then
\[\big\langle B_{N\to m}F\big\rangle_{N,p}=\frac{\int\rho_{m,N,p}(U)F(U)\,d\mathrm{Haar}(U)}{\int\rho_{m,N,p}(U)\,d\mathrm{Haar}(U)},\]
where \(\rho_{m,N,p}\) is the effective density with its
field-independent scalar factor removed.

\emph{Sketch:} Iterating the one-step density identity gives the
effective measure at level \(m\). Duality with the conditional averaging
kernels and equality of partition functions give the formula; the
removed scalar factor cancels in the ratio.

\emph{Lean:}
\href{https://github.com/yang-mills-group/OpenAI-YM-20261028/blob/3456e5e75cf8d4cd558955a5c7638fba2d68f464/YM/UVExtension/Hearth/Dine.lean\#L42-L52}{\texttt{Hearth.scaledK}},
\href{https://github.com/yang-mills-group/OpenAI-YM-20261028/blob/3456e5e75cf8d4cd558955a5c7638fba2d68f464/YM/UVExtension/Hearth/Dine.lean\#L40}{\texttt{Hearth.tasting}},
\href{https://github.com/yang-mills-group/OpenAI-YM-20261028/blob/3456e5e75cf8d4cd558955a5c7638fba2d68f464/YM/Push/GRID_APPROXIMATION/KernelBasic.lean\#L21-L22}{\texttt{GridApprox.back}},
\href{https://github.com/yang-mills-group/OpenAI-YM-20261028/blob/3456e5e75cf8d4cd558955a5c7638fba2d68f464/YM/Push/GRID_APPROXIMATION/Iterate.lean\#L24-L25}{\texttt{GridApprox.backTo}}.

\textbf{Lemma 4.12 (convergence of the effective measure at a fixed
level).} Fix the arc and its coefficient bundle. At a fixed period
\(p\), let the shifts satisfy \(s_{N,p}\to s\) with \(s\) real in the
disc. If \(F_N\) are bounded measurable functions of the level-\(m\)
configuration and form a uniformly Cauchy sequence, then the level-\(m\)
effective expectations of \(F_N\) converge as \(N\to\infty\).

\emph{Sketch:} The effective densities with their scalar factors removed
converge uniformly, by stability of the trajectory family. Their
integrals have a uniform positive lower bound. The effective
expectations are ratios of convergent integrals.

\emph{Lean:}
\href{https://github.com/yang-mills-group/OpenAI-YM-20261028/blob/3456e5e75cf8d4cd558955a5c7638fba2d68f464/YM/UVExtension/Hearth/Dine.lean\#L65-L72}{\texttt{Hearth.arrival}},
\href{https://github.com/yang-mills-group/OpenAI-YM-20261028/blob/3456e5e75cf8d4cd558955a5c7638fba2d68f464/YM/UVExtension/Hearth/Dine.lean\#L58-L62}{\texttt{Hearth.seating}},
\href{https://github.com/yang-mills-group/OpenAI-YM-20261028/blob/3456e5e75cf8d4cd558955a5c7638fba2d68f464/YM/UVExtension/Hearth/Trail.lean\#L82-L84}{\texttt{strippedOn}},
\href{https://github.com/yang-mills-group/OpenAI-YM-20261028/blob/3456e5e75cf8d4cd558955a5c7638fba2d68f464/YM/UVExtension/Hearth/Trail.lean\#L85-L97}{\texttt{zPositive}},
\href{https://github.com/yang-mills-group/OpenAI-YM-20261028/blob/3456e5e75cf8d4cd558955a5c7638fba2d68f464/YM/UVExtension/Forest/Coda.lean\#L47-L52}{\texttt{Mast.signs}}.
\emph{(proof of \texttt{arrival} read; the three inputs statement only)}

\textbf{Lemma 4.13 (coarse approximation of renormalized observables).}
Fix the arc and coefficient bundle. At period \(p\), assume the shifts
converge to a real shift in the shooting disc, and assume the transfer
and remainder estimates and slow-growth bounds of III.13.6. Let
\(J\le m\le k\), with \(\mathrm{cs}_m^2\le a_J/a_m\). Require enough
separation for the supports of the degree-\(d\) coarse observables:
\(2(a_kr_d+\mathrm{tail}_k)<a_J\), where \(r_d\) is their fixed support
radius and \(\mathrm{tail}_k=30a_k(J(k)+5)\) bounds the reach of the
later kernels. Let \(u\) be a degree-\(d\) label with
\(\mathrm{ref}(d)\le m\), and let \(w>0\) satisfy
\[w\,a_m^{-d}\cdot4(1+\mathrm{columns}_k)\,e^{-(\sigma_k-\sigma_m)/8}\le\tfrac12\,\mathrm{fly}_{d,m}.\]
Then there are level-\(k\) observables \(g_N=\sum_jc_j(N)\,\Lambda_j\)
with the following properties:

\begin{itemize}
\tightlist
\item
  the \(\Lambda_j\) are fixed level-\(k\) observables, and the
  coefficients \(c_j(N)\) converge as \(N\to\infty\);
\item
  for large \(N\), the chessboard root, at level \(J\), of
  \(B_{N\to k}g_N-O^{N,p}_u(0)\) is at most
  \(\mathrm{stock}_k(u)\cdot\mathrm{stakes}(J,m)/w\).
\end{itemize}

For fixed \(d,m,J,u\) and every \(\zeta,\varepsilon>0\), all
sufficiently large \(k\) admit \(w>0\) with
\(wa_m^{-d}4(1+\mathrm{columns}_k)e^{-(\sigma_k-\sigma_m)/8}\le\zeta/2\)
and \(\mathrm{stock}_k(u)\mathrm{stakes}(J,m)/w\le\varepsilon\). Apply
this with \(\zeta=\mathrm{fly}_{d,m}\). The product of the two
requirements tends to \(0\), because columns and stock grow
subexponentially in \(\sigma_k\) while \(e^{(\sigma_k-\sigma_m)/8}\)
grows exponentially.

\emph{Lean:}
\href{https://github.com/yang-mills-group/OpenAI-YM-20261028/blob/3456e5e75cf8d4cd558955a5c7638fba2d68f464/YM/UVExtension/Hearth/Spices.lean\#L64-L68}{\texttt{Hearth.peeling}},
\href{https://github.com/yang-mills-group/OpenAI-YM-20261028/blob/3456e5e75cf8d4cd558955a5c7638fba2d68f464/YM/UVExtension/Hearth/Table.lean\#L37-L59}{\texttt{Hearth.seatingM}}.
\emph{(proofs read; the level-\(k\) representation itself is statement
only)}

\textbf{Lemma 4.14 (chessboard product estimate).} Assume the Wilson
measure satisfies the reflection-positivity and translation properties
for the level-\(h\) tiling. Let \(O_i,Z_i\) be bounded measurable
observables supported in its reference cell, with chessboard roots
\(\|O_i\|\le c_i\) and \(\|Z_i-O_i\|\le e_i\), and let \(a_i\) be
distinct sites of the level-\(h\) lattice. Then
\[\Big|\Big\langle\prod_iZ_i(\cdot+a_i)\Big\rangle-\Big\langle\prod_iO_i(\cdot+a_i)\Big\rangle\Big|\le\prod_i(e_i+c_i)-\prod_ic_i.\]

\emph{Sketch:} Expand \(Z_i=O_i+(Z_i-O_i)\) and subtract the term
containing only the \(O_i\). Apply the chessboard inequality to each
remaining product on distinct cells; summing gives the displayed bound.

\emph{Lean:}
\href{https://github.com/yang-mills-group/OpenAI-YM-20261028/blob/3456e5e75cf8d4cd558955a5c7638fba2d68f464/YM/Push/GRID_APPROXIMATION/Telescope.lean\#L34-L78}{\texttt{GridApprox.tariff}},
applied in
\href{https://github.com/yang-mills-group/OpenAI-YM-20261028/blob/3456e5e75cf8d4cd558955a5c7638fba2d68f464/YM/UVExtension/Hearth/Table.lean\#L63-L80}{\texttt{Hearth.banquet}}.
\emph{(\texttt{tariff} statement only; \texttt{banquet} proof read)}

\emph{Proof of Proposition III.13.6.}

\begin{enumerate}
\def\labelenumi{\arabic{enumi}.}
\tightlist
\item
  The chessboard bound in the schedule gives, for each \(i\), a bound
  \(t_i\) on the chessboard root of \(O_{u_i}\) at level \(J\).
\item
  Fix \(q>0\). Choose \(k\) large and weights \(w_i\) by Lemma 4.13, so
  that each error is \(\le q\). The support separation makes conditional
  averaging commute with the product of translated coarse observables.
  The RG identity (Lemma 4.11) and the chessboard product estimate
  (Lemma 4.14) then give a sequence \(F_N\) such that, for large \(N\),
  \(|F_N-\Psi_N|\le\Phi(q):=\prod(q+t_i)-\prod t_i\). Here \(F_N\) is
  the level-\(k\) effective expectation of the product of the translated
  \(g^{(i)}_N\), and it converges by Lemma 4.12.
\item
  Since \(\Phi\) is continuous with \(\Phi(0)=0\), one can choose \(q\)
  with \(\Phi(q)<\varepsilon/4\). The triangle inequality through two
  terms of the convergent approximating sequence shows that \((\Psi_N)\)
  is Cauchy.
\end{enumerate}

\emph{Lean:}
\href{https://github.com/yang-mills-group/OpenAI-YM-20261028/blob/3456e5e75cf8d4cd558955a5c7638fba2d68f464/YM/UVExtension/Hearth/Table.lean\#L87-L112}{\texttt{Hearth.food}},
\href{https://github.com/yang-mills-group/OpenAI-YM-20261028/blob/3456e5e75cf8d4cd558955a5c7638fba2d68f464/YM/UVExtension/Hearth/Table.lean\#L82-L84}{\texttt{Hearth.closet}}.

\textbf{Lemma 4.15 (slicing for large periods).}

\emph{Definition.} \emph{Slicing}\((d,p)\) holds if there are:

\begin{itemize}
\tightlist
\item
  finitely many test shots, admissible at period \(p\);
\item
  a constant \(a>0\);
\item
  a subspace \(U\subset\mathcal L_d\) of codimension
  \(\mathrm{charge}(G,d)\), the dimension of the degree-\(d\) Maxwell
  equation ideal;
\end{itemize}

such that, for large \(N\),
\(a\|x\|\le\sum_i\|\mathrm{snapshot}_i(N,p)x\|\) for \(x\in U\).

\emph{Statement.} Assume a fixed presentation identifying the
finite-cutoff kernels of the label projections with the transported
Maxwell equation ideals. Suppose every admissible snapshot at every
fixed period converges, and every extraction order has a refinement with
base data whose projections agree with this presentation. Then
Slicing\((d,p)\) holds for all large \(p\). For the bundle of III.13.7
the conclusion takes this form: for each \(d\) and all large \(p\) there
is \(\delta>0\) such that Slicing holds for every assignment of shifts
with \(|s_{N,p}|\le\delta\) eventually.

\emph{Sketch:} By contradiction. A sequence of bad periods defines an
extraction order. Along a refinement, finite detection on the range of
the limiting projection gives Slicing at all large indices.

\emph{Lean:}
\href{https://github.com/yang-mills-group/OpenAI-YM-20261028/blob/3456e5e75cf8d4cd558955a5c7638fba2d68f464/YM/UVExtension/Harvest/Oats.lean\#L59-L61}{\texttt{Harvest.rootwards}},
\href{https://github.com/yang-mills-group/OpenAI-YM-20261028/blob/3456e5e75cf8d4cd558955a5c7638fba2d68f464/YM/UVExtension/Harvest/Oats.lean\#L34-L36}{\texttt{Harvest.Slicing}},
\href{https://github.com/yang-mills-group/OpenAI-YM-20261028/blob/3456e5e75cf8d4cd558955a5c7638fba2d68f464/YM/UVExtension/Harvest/Barn.lean\#L27-L30}{\texttt{Barn.ripe}}.

\textbf{Lemma 4.16 (proof of Proposition III.13.5).} Fix \(d,k\) and
\(p\) beyond the slicing threshold for the degrees up to \(d+k+1\). Use
the arc and bundle hypotheses of Lemma 4.4, zero shifts, and labels of
degree \(\le d\). Then:

\begin{enumerate}
\def\labelenumi{\arabic{enumi}.}
\tightlist
\item
  \textbf{Dense set.} On the points of \(\Omega^n_p\) whose coordinates
  lie on the level-\(m\) grid for some \(m\ge\max\mathrm{ref}\) (a dense
  set), the lattice correlations converge by III.13.6.
\item
  \textbf{Uniform difference bounds.} Slicing and the label-transfer
  bounds give bounds on lattice differences of order \(\le k+1\),
  uniform in \(N\) on compact subsets of \(\Omega^n_p\). Hence the
  correlations and their differences of order \(\le k\) are
  asymptotically equicontinuous there.
\item
  \textbf{Arzelà--Ascoli.} It gives locally uniform convergence to a
  continuous \(g\), and to continuous limits of the difference
  quotients. These limits are the iterated derivatives of \(g\).
\item
  \textbf{Smeared limits.} For admissible tests, the smeared
  correlations converge by dominated convergence, with domination by the
  uniform period-\(p\) bound.
\end{enumerate}

\emph{Lean:}
\href{https://github.com/yang-mills-group/OpenAI-YM-20261028/blob/3456e5e75cf8d4cd558955a5c7638fba2d68f464/YM/UVExtension/Poppy/Fountain.lean\#L96-L108}{\texttt{Poppy.blossoms}},
\href{https://github.com/yang-mills-group/OpenAI-YM-20261028/blob/3456e5e75cf8d4cd558955a5c7638fba2d68f464/YM/UVExtension/Poppy/Fountain.lean\#L40-L69}{\texttt{Poppy.coveringR}},
\href{https://github.com/yang-mills-group/OpenAI-YM-20261028/blob/3456e5e75cf8d4cd558955a5c7638fba2d68f464/YM/Extraction/Tame.lean\#L35-L97}{\texttt{Extraction.growFrom}},
\href{https://github.com/yang-mills-group/OpenAI-YM-20261028/blob/3456e5e75cf8d4cd558955a5c7638fba2d68f464/YM/UVExtension/Poppy/Sparrow.lean\#L42-L48}{\texttt{Poppy.nourishing}},
\href{https://github.com/yang-mills-group/OpenAI-YM-20261028/blob/3456e5e75cf8d4cd558955a5c7638fba2d68f464/YM/UVExtension/Poppy/Sparrow.lean\#L61-L62}{\texttt{Poppy.mature}};
difference bounds:
\href{https://github.com/yang-mills-group/OpenAI-YM-20261028/blob/3456e5e75cf8d4cd558955a5c7638fba2d68f464/YM/UVExtension/Fallow/Plow.lean\#L66-L70}{\texttt{Fallow.reaped}}.
\emph{(proofs read through the stated transfer estimates)}

\textbf{Lemma 4.17 (analytic dependence of the correlations on the
shift).} Assume the arc, bundle bounds, paid schedule, analytic transfer
disc, leases and target rate bound of Lemma 4.18. Fix \(d,m\), a period
\(p\) with slicing stable for all degrees through \(d+m+1\) for
sufficiently small shifts, and finitely many labels of degree \(\le d\).
There is \(b'>0\) such that, for all sufficiently large \(N\), the
lattice correlations extend holomorphically to \(|z|<b'\). On every
compact subset of \(\Omega^n_p\), their lattice differences through
order \(m\) have bounds uniform in \(N\) and \(z\) throughout this disc.
At arbitrary lattice configurations the bound supplied at each cutoff
may depend on \(N\).

\emph{Sketch:} Slicing gives uniform real difference bounds through one
higher order. The analytic transfer estimates and complex correlation
bounds permit the Banach-valued Vitali argument on each separated
compact set. The resulting limits have the spatial derivatives and
smeared limits of III.13.7.5. Shrinking the disc once more incorporates
every cutoff at which all reference levels fit. At each fixed level and
period, the normalized effective densities have bounded holomorphic
extensions; after shrinking the disc, the moduli of their partition
integrals are bounded away from zero uniformly over the allowed cutoffs.

\emph{Lean:}
\href{https://github.com/yang-mills-group/OpenAI-YM-20261028/blob/3456e5e75cf8d4cd558955a5c7638fba2d68f464/YM/UVExtension/Beach/Rolls.lean\#L37-L54}{\texttt{Beach.inlet}},
\href{https://github.com/yang-mills-group/OpenAI-YM-20261028/blob/3456e5e75cf8d4cd558955a5c7638fba2d68f464/YM/UVExtension/Beach/Under.lean\#L43-L63}{\texttt{Beach.sparksP}},
\href{https://github.com/yang-mills-group/OpenAI-YM-20261028/blob/3456e5e75cf8d4cd558955a5c7638fba2d68f464/YM/UVExtension/Reports/Prefixes.lean\#L78-L112}{\texttt{Beach.keep\_prefix}},
\href{https://github.com/yang-mills-group/OpenAI-YM-20261028/blob/3456e5e75cf8d4cd558955a5c7638fba2d68f464/YM/UVExtension/Pier/Plates.lean\#L68-L87}{\texttt{Pier.roomR}},
\href{https://github.com/yang-mills-group/OpenAI-YM-20261028/blob/3456e5e75cf8d4cd558955a5c7638fba2d68f464/YM/UVExtension/Pier/Heath.lean\#L28-L50}{\texttt{Pier.mounted}}.
\emph{(proofs of \texttt{sparksP} and \texttt{keep\_prefix} read;
\texttt{inlet}, \texttt{roomR} and \texttt{mounted} statement only)}

\paragraph{4.5 Reports}\label{reports}

\textbf{Lemma 4.18 (the centre is floral).} Assume:

\begin{itemize}
\tightlist
\item
  an arc and the bundle conditions;
\item
  a paid family;
\item
  a disc in the analytic domain, with holomorphic transfer data for all
  \(d,N,p\) with \(\mathrm{ref}(d)\le N\);
\item
  the leases and target rates \(\le-7/16\);
\item
  the calibration package with \(Q\equiv1\), and the priced sup bounds.
\end{itemize}

Then the member \(s=0\) is floral.

\emph{Sketch:} The shooting and analyticity properties come from the arc
(Lemmas 4.6, 4.9). For large \(p\), the analytic correlation data come
from Lemma 4.17.

\emph{Lean:}
\href{https://github.com/yang-mills-group/OpenAI-YM-20261028/blob/3456e5e75cf8d4cd558955a5c7638fba2d68f464/YM/UVExtension/Oasis/Flora.lean\#L81-L87}{\texttt{Oasis.survives}}.

\textbf{Lemma 4.19 (mesh limits along the full sequence).} For a floral
family \(C\), an extraction order \(f\) and a table \(T\) with a
construction record for \((C,f,T)\), the order (full mesh sequence,
periods of \(f\)) is ordered for \(T\).

\emph{Sketch:} At each fixed period there are tables whose smeared
limits are taken along the full mesh sequence; they come from
Proposition III.13.6 with shift \(0\), through the fixed-period smearing
theorem. The record supplies the limit along the mesh subsequence of
\(f\), and uniqueness of limits identifies the two.

\emph{Lean:}
\href{https://github.com/yang-mills-group/OpenAI-YM-20261028/blob/3456e5e75cf8d4cd558955a5c7638fba2d68f464/YM/UVExtension/Reports/Ultraviolet.lean\#L50-L58}{\texttt{Reports.full\_tests}},
\href{https://github.com/yang-mills-group/OpenAI-YM-20261028/blob/3456e5e75cf8d4cd558955a5c7638fba2d68f464/YM/UVExtension/Reports/Ultraviolet.lean\#L62-L69}{\texttt{Reports.useFull}}.

\textbf{Lemma 4.20 (assembly of the reports).}

\begin{itemize}
\item
  \textbf{III.13.2.} For each extraction order \(F\):

  \begin{itemize}
  \tightlist
  \item
    item 10 of III.13.1 gives \(g\), the base data and a seed;
  \item
    the seed gives the cell hypotheses, hence grid approximants and weak
    comparison systems;
  \item
    the OPE entrance theorem gives Maxwell data, infrared data with
    non-zero \(J\), and the UV package;
  \item
    the reconstruction assembly theorem gives a construction record;
  \item
    Lemma 4.19 gives the full-mesh ordering;
  \item
    the OPE services are assembled from the seed, the logarithmic
    package and the base data.
  \end{itemize}

  Proposition III.13.5 supplies the fixed-volume limits.
\item
  \textbf{III.13.3.} Apply the every-phase reconstruction theorem to the
  record, and the coherent-prescription theorem to the derivative
  sector.
\item
  \textbf{III.13.4.}

  \begin{itemize}
  \tightlist
  \item
    Coverage gives portfolios along any finer ultrafilter: take the
    levels \(g(k)\) of item 8 of III.13.1, and the limits along
    \(\mathcal U\).
  \item
    The OPE services give the bases \(Y_k\) and the convergence of the
    Gram forms (the Gram comparison theorem, through the grading theorem
    and the Gram limit theorem).
  \end{itemize}
\end{itemize}

\emph{Lean:}
\href{https://github.com/yang-mills-group/OpenAI-YM-20261028/blob/3456e5e75cf8d4cd558955a5c7638fba2d68f464/YM/SharpGram/Common.lean\#L16-L37}{\texttt{coordination}},
\href{https://github.com/yang-mills-group/OpenAI-YM-20261028/blob/3456e5e75cf8d4cd558955a5c7638fba2d68f464/YM/SharpGram/Common.lean\#L41-L60}{\texttt{common}},
\href{https://github.com/yang-mills-group/OpenAI-YM-20261028/blob/3456e5e75cf8d4cd558955a5c7638fba2d68f464/YM/SharpGram/Measured.lean\#L11-L36}{\texttt{habitat}},
\href{https://github.com/yang-mills-group/OpenAI-YM-20261028/blob/3456e5e75cf8d4cd558955a5c7638fba2d68f464/YM/SharpGram/Weak/Reading.lean\#L66-L67}{\texttt{fromCoverage}},
\href{https://github.com/yang-mills-group/OpenAI-YM-20261028/blob/3456e5e75cf8d4cd558955a5c7638fba2d68f464/YM/SharpGram/InContext.lean\#L65-L71}{\texttt{attains}}.
\emph{(proofs of the three reports and of \texttt{fromCoverage} read;
\texttt{attains} read to the level of its two inputs)}

\textbf{Lemma 4.21 (further estimates used in choosing the step
constants).}

\begin{itemize}
\tightlist
\item
  \textbf{Connected-activity targets.} Take a target with an expense
  bound. There is \(M\) such that for every side \(L\ge M\),
  \(W,\mu>0\), \(a\ge0\), indices \(j,k,N\), \(0\le\xi<\mathrm{exp}\)
  and \(e>0\) with \(\mathrm{rate}+e<\eta<0\), the prescribed
  connected-activity estimate holds for all sufficiently large \(x\),
  with weight \(\beta_F(x,L)^e\) and every target order with parameter
  \(W\). This bound supplies the steepness requirement in Lemma 4.1.
\item
  \textbf{Ursell sum.} Assume the support compatibility tests,
  nonnegative weights, a nonnegative support exponent, the
  cluster-expansion readiness bound with \(0\le q<1\), and a rooted
  input budget \(B_0\). The averaged marked Ursell sum over clusters,
  weighted by the cost of the union of their supports, has budget
  \(B_0/(1-q)\) (III.5 in the form of III.8).
\item
  \textbf{Uniform rate.} For a connected-activity target there is a side
  threshold such that, after fixing a larger side, a branching patch,
  cutoff function, positive \(W,\mu\), arity and order bounds,
  \(0<c<\mathrm{exp}\) and \(e>0\) with \(\mathrm{rate}+e<\gamma<0\),
  there is a rate function \(\vartheta(x)\) with
  \(|\vartheta(x)|\le\beta_F(x,L)^{\gamma+\varepsilon}\) eventually for
  every \(\varepsilon>0\). Eventually in \(x\), it bounds the prescribed
  analytic costs uniformly over admissible tubes, representatives and
  budget points with exponent at most \(c\) and the specified order,
  power, threshold and radius bounds.
\item
  \textbf{Jet realization.} Let the real variable space be finite
  dimensional, and fix a point \(x\) and jet order \(j\). Suppose a
  determining parameter set carries \(C^j\) germs at \(x\), and the
  prescribed jet coefficients extend holomorphically along every
  holomorphic curve in the parameter domain, agreeing with those germs
  on the determining set. There is a family of \(C^j\) germs at \(x\)
  realizing these coefficients throughout that domain, holomorphic along
  such curves together with its jets through order \(j\).
\end{itemize}

\emph{Lean:}
\href{https://github.com/yang-mills-group/OpenAI-YM-20261028/blob/3456e5e75cf8d4cd558955a5c7638fba2d68f464/YM/UVExtension/Linden/Seasons.lean\#L38-L85}{\texttt{Linden.lives}},
\href{https://github.com/yang-mills-group/OpenAI-YM-20261028/blob/3456e5e75cf8d4cd558955a5c7638fba2d68f464/YM/UVExtension/Copse/Spring.lean\#L106-L114}{\texttt{Copse.tangled}},
\href{https://github.com/yang-mills-group/OpenAI-YM-20261028/blob/3456e5e75cf8d4cd558955a5c7638fba2d68f464/YM/UVExtension/Concert/Vines.lean\#L22-L50}{\texttt{Concert.basisE}},
\href{https://github.com/yang-mills-group/OpenAI-YM-20261028/blob/3456e5e75cf8d4cd558955a5c7638fba2d68f464/YM/UVExtension/Lamppost/Blossom.lean\#L20-L70}{\texttt{Lamppost.forge}}.

\subsubsection{5. Where the difficulty
lies}\label{where-the-difficulty-lies-15}

Fixed-level convergence requires control of arbitrarily many ultraviolet
levels. The contraction estimate and the \(\mathrm{soft}^v\) bound give
two uniform approximations by convergent sequences. For correlations,
Lemma 4.13 balances exponentially improving coarse approximation against
subexponential coefficient costs; the chessboard inequality then
controls products.

The derivative limit holds only for sufficiently large periods, with
threshold depending on the degree and derivative order. Slicing gives a
lower bound on a subspace of the required dimension. The plain smeared
limit, used in Lemma 4.19, holds at every fixed period.

The coupling and activity analyticity holds in one disc for all cutoffs
and periods. The analytic correlation disc may depend on the finite
label family and the period. The construction fixes support exponent
\(100\) and target rate \(-7/16\); no alternative rate is asserted here.

The family in III.13.2 is chosen before the extraction order. III.13.3
and III.13.4 state a weaker quantifier order, but their proofs use that
same choice. The main theorem follows the separate construction route of
III.14; equality of the families produced by the two routes is not
asserted.

\subsubsection{6. Notation}\label{notation-15}

{\def\LTcaptype{none} % do not increment counter
\begin{longtable}[]{@{}
  >{\raggedright\arraybackslash}p{(\linewidth - 4\tabcolsep) * \real{0.3333}}
  >{\raggedright\arraybackslash}p{(\linewidth - 4\tabcolsep) * \real{0.3333}}
  >{\raggedright\arraybackslash}p{(\linewidth - 4\tabcolsep) * \real{0.3333}}@{}}
\toprule\noalign{}
\begin{minipage}[b]{\linewidth}\raggedright
Written proof
\end{minipage} & \begin{minipage}[b]{\linewidth}\raggedright
Lean
\end{minipage} & \begin{minipage}[b]{\linewidth}\raggedright
Meaning
\end{minipage} \\
\midrule\noalign{}
\endhead
\bottomrule\noalign{}
\endlastfoot
\(C\) &
\href{https://github.com/yang-mills-group/OpenAI-YM-20261028/blob/3456e5e75cf8d4cd558955a5c7638fba2d68f464/YM/Source/Construction.lean\#L24-L30}{\texttt{Family}}
& Wilson family: history, holonomy data, label frames, local sections,
reference levels, transfer data \\
\(\beta_{N,p}\) & \texttt{Family.beta} & Wilson coupling at cutoff
\(N\), period \(p\) \\
consistent &
\href{https://github.com/yang-mills-group/OpenAI-YM-20261028/blob/3456e5e75cf8d4cd558955a5c7638fba2d68f464/YM/Source/Construction.lean\#L44-L62}{\texttt{Family.Makes}}
& selected trajectory, sound sections, stage and stencil conditions,
priced bounds, calibrated transfer \\
calibrated &
\href{https://github.com/yang-mills-group/OpenAI-YM-20261028/blob/3456e5e75cf8d4cd558955a5c7638fba2d68f464/YM/Spec/Closed.lean\#L58-L61}{\texttt{Spec.Calibrated}}
& history generated by a full schedule with per-step estimates \\
\(\mathrm{ref}(d)\) & \texttt{Family.ref} & reference level of degree
\(d\) \\
penalty \(b\) & \texttt{KernelChoices.penalty} & coefficient of the
Gaussian block-spin term \\
sheet, rolls, wages, \(\theta\), delays, bar, wait, paid &
\texttt{Sheet}, \texttt{Rolls}, \texttt{Return.Wages},
\texttt{Protection}, \texttt{Delays}, \texttt{Bar}, \texttt{Wait},
\texttt{Paid} & constants in the order of choice \\
alloy, factor, directed history & \texttt{Alloy}, \texttt{Roost.Factor},
\texttt{Directed} & a schedule of RG steps with their bounds \\
shooting structure, arc, \(r_*\) & \texttt{Spire.Uses},
\texttt{Upland.Arc}, \texttt{Arc.radius} & trajectories ending at
\(\bar\beta+s\); disc of analyticity \\
shift \(s\), assignment \(\mathbf s\) & \texttt{Orchard.Rose},
\texttt{Orchard.Keys}, \texttt{Mast.roseA} (\(\mathbf s=0\)) & real
shift of the terminal coupling \\
label-transfer bundle, \(C^{\mathbf s}\) & \texttt{Mast.Cellar},
\texttt{Cellar.yard} & coefficient and transfer data for every shift
assignment; the associated family \\
floral & \texttt{Oasis.Floral}, \texttt{Oasis.Spring} & III.13.7 \\
bundle conditions & \texttt{Harvest.Barn} & Scrim bounds and scheduled
chessboard bound for all shifts; base data at the centre \\
chessboard bound in the schedule & \texttt{Extraction.Wcmain} & §1
definition \\
fixed-volume limit with \(k\) derivatives & \texttt{Poppy.Margins} &
III.13.5 \\
slicing & \texttt{Harvest.Slicing} & Lemma 4.15 \\
logarithmic OPE package & \texttt{OPE.Log.LogPackage} & III.13.1.7 \\
calibrated label steps & \texttt{SharpGram.TrueSteps} & III.13.1.7 \\
coverage & \texttt{Dense.Coverage} & III.13.1.8 \\
OPE entrance conditions & \texttt{Acme.Aroma} & III.13.1.6 \\
grid approximants, weak comparison systems & \texttt{Org.Grill},
\texttt{East.CompleteWeak} & III.13.1.9 \\
OPE services & \texttt{OPE.Facilities} & services of VII.4 with a log
package \\
construction record, reconstruction for every phase &
\texttt{Spec.Final}, \texttt{EveryPhase.Almanac},
\texttt{EveryPhase.Completed} & II.2, V.3 \\
portfolio & \texttt{Weak.Portfolio} & moving-frame probe data along an
ultrafilter \\
effective expectation & \texttt{Hearth.tasting} & normalized expectation
against the effective density with its scalar factor removed \\
\(B_{N\to m}\) &
\href{https://github.com/yang-mills-group/OpenAI-YM-20261028/blob/3456e5e75cf8d4cd558955a5c7638fba2d68f464/YM/Push/GRID_APPROXIMATION/Iterate.lean\#L24-L25}{\texttt{GridApprox.backTo}}
& iterated conditional averaging through the block-spin kernels \\
\(V_L\) &
\href{https://github.com/yang-mills-group/OpenAI-YM-20261028/blob/3456e5e75cf8d4cd558955a5c7638fba2d68f464/YM/Bid/Loan.lean\#L63}{\texttt{Bid.vowel}}
& engineering rescaling by \(L^{d-e}\) on grade \(e\) \\
\(\mathrm{stock}_k(u)\) &
\href{https://github.com/yang-mills-group/OpenAI-YM-20261028/blob/3456e5e75cf8d4cd558955a5c7638fba2d68f464/YM/MergeProof/South/Dine.lean\#L27}{\texttt{stockK}}
& coefficient bound with subexponential growth in \(\sigma_k\) \\
\(\mathrm{stakes}(J,m)\) &
\href{https://github.com/yang-mills-group/OpenAI-YM-20261028/blob/3456e5e75cf8d4cd558955a5c7638fba2d68f464/YM/MergeProof/Strudel.lean\#L21}{\texttt{stakes}}
& fixed-scale chessboard comparison cost \\
\(\mathrm{fly}_{d,m}\) &
\href{https://github.com/yang-mills-group/OpenAI-YM-20261028/blob/3456e5e75cf8d4cd558955a5c7638fba2d68f464/YM/MergeProof/Stove.lean\#L37}{\texttt{Bandage.fly}}
& positive allowance in the label-transfer estimate \\
\(\mathrm{soft}\), \(\mathrm{rad}\), gain & \texttt{Camp.soft},
\texttt{Camp.rad}, \texttt{Bid.gainChart} & \(\mathrm{rim}/8\),
\(\mathrm{rim}/4\), \(e^{(\sigma_k-\sigma_m)/8}\) \\
\end{longtable}
}

\subsubsection{7. What remains
unexplained}\label{what-remains-unexplained-15}

\begin{enumerate}
\def\labelenumi{\arabic{enumi}.}
\tightlist
\item
  \textbf{Quantitative packages.} The shell estimates in the OPE
  entrance condition, the full transfer, remainder, sup and weak
  comparison bounds, and the one-step comparison assumptions in Lemma
  4.3 retain their contracts from III.7, III.11 and VII.4. Their
  complete expansion into conventional notation remains to be written
  there. The summaries here do not replace those contracts.
\item
  \textbf{Proofs not fully checked.} The detailed constant issuers
  underlying Lemmas 4.1--4.4, the fitted recursion behind Lemma 4.6, the
  contraction and far-level decay estimates in Lemma 4.8, the one-step
  label estimates behind Lemma 4.10, and the auxiliary estimates in
  Lemma 4.21 are used as stated. Their internal proofs have not all been
  checked. The difference estimate in Lemma 4.16 is traced through
  slicing and the transfer bounds; its preceding one-step estimates
  remain inputs.
\item
  \textbf{Two construction routes.} No identification of the families
  from III.13 and III.14, or transfer of III.13.5 to the III.14 family,
  is established here. A common witness satisfying both lists would
  require a separate argument.
\item
  \textbf{Moving-frame estimates.} The basis construction and its
  comparison with the free Gram form belong to VII.1/VII.3. Lemma 4.20
  uses those estimates; their proofs are not reproduced here.
\item
  \textbf{The penalty.} Its role as the Gaussian block-spin coefficient
  is fixed by the kernel choices. Its precise normalization relative to
  Balaban's parameter \(a\) has not been verified.
\item
  \textbf{Other reports.}
  \href{https://github.com/yang-mills-group/OpenAI-YM-20261028/blob/3456e5e75cf8d4cd558955a5c7638fba2d68f464/YM/UVExtension/Reports/Ultraviolet.lean\#L101-L112}{\texttt{Reports.ultraviolet}},
  \href{https://github.com/yang-mills-group/OpenAI-YM-20261028/blob/3456e5e75cf8d4cd558955a5c7638fba2d68f464/YM/UVExtension/Oasis/Fruition.lean\#L33-L41}{\texttt{Oasis.fruition}}
  give additional reports. They are not needed for the assembly of
  III.13.1--III.13.4. The same applies to the offset and analytic-tuning
  test wrappers and the separate fixed-volume existence theorem.
\end{enumerate}

\subsection{III.14. The construction run: flat representations and the
ultraviolet
rate}\label{iii.14.-the-construction-run-flat-representations-and-the-ultraviolet-rate}

\textbf{Part:} III, Ultraviolet construction (renormalization group).
\textbf{Lean:}

\begin{itemize}
\tightlist
\item
  the endpoint
  \href{https://github.com/yang-mills-group/OpenAI-YM-20261028/blob/3456e5e75cf8d4cd558955a5c7638fba2d68f464/YM/Independent/ConstructionRun.lean}{\texttt{YM/Independent/ConstructionRun.lean}};
\item
  the native component recipe
  \href{https://github.com/yang-mills-group/OpenAI-YM-20261028/blob/3456e5e75cf8d4cd558955a5c7638fba2d68f464/YM/Integration/NativeComponentRecipe.lean}{\texttt{YM/Integration/NativeComponentRecipe.lean}}
  and
  \href{https://github.com/yang-mills-group/OpenAI-YM-20261028/blob/3456e5e75cf8d4cd558955a5c7638fba2d68f464/YM/Integration/ActualRunNativeComponentRecipe.lean}{\texttt{ActualRunNativeComponentRecipe.lean}};
\item
  the flat-representation chain
  \href{https://github.com/yang-mills-group/OpenAI-YM-20261028/tree/3456e5e75cf8d4cd558955a5c7638fba2d68f464/YM/RG08/Flat}{\texttt{YM/RG08/Flat/}},
  \href{https://github.com/yang-mills-group/OpenAI-YM-20261028/tree/3456e5e75cf8d4cd558955a5c7638fba2d68f464/YM/RG08/Packets}{\texttt{YM/RG08/Packets/}}
  and
  \href{https://github.com/yang-mills-group/OpenAI-YM-20261028/tree/3456e5e75cf8d4cd558955a5c7638fba2d68f464/YM/RG08/Rate}{\texttt{YM/RG08/Rate/}};
\item
  the local-estimate package of
  \href{https://github.com/yang-mills-group/OpenAI-YM-20261028/blob/3456e5e75cf8d4cd558955a5c7638fba2d68f464/YM/RG12/Burst/Wardrobe.lean}{\texttt{YM/RG12/Burst/Wardrobe.lean}};
\item
  the packages that define the numerical floors of the recipe:
  \href{https://github.com/yang-mills-group/OpenAI-YM-20261028/tree/3456e5e75cf8d4cd558955a5c7638fba2d68f464/YM/RG08/Redo}{\texttt{YM/RG08/Redo/}},
  \texttt{Yamen/}, \texttt{Sideshow/}, \texttt{Necks/},
  \texttt{Rudder/}, \texttt{Room/}, \texttt{Weald/}, \texttt{Wisp/},
  \texttt{Trill/}, \texttt{Halo/}, \texttt{Line/}, \ldots{} (§4.6).
\end{itemize}

\textbf{Size:} 14,290 lines in the main-theorem closure (2,509
declarations).

The unit selects the construction run and proves a quantitative limit
theorem for its Gaussian quadratic part.

\begin{itemize}
\tightlist
\item
  \textbf{The endpoint.} For a compact matrix gauge group \(G\) with
  simple nonabelian Lie algebra and a penalty \(b>0\), it selects a
  \emph{run}: the Wilson family of the main theorem with its scale and
  extraction data. The run satisfies the \emph{funded} condition of
  III.15 and carries the \emph{native component recipe}.
\item
  \textbf{The Gaussian limit.} At each fixed stage, the order-zero
  quadratic kernel converges as the ultraviolet cutoff is removed, with
  a power rate in the ratio of lattice spacings (Proposition III.14.3).
  This is one of three components of the recipe.
\end{itemize}

\textbf{Use in the main theorem.} VIII.1
(\href{https://github.com/yang-mills-group/OpenAI-YM-20261028/blob/3456e5e75cf8d4cd558955a5c7638fba2d68f464/YM/Nonperturbative/Continuum.lean\#L28-L48}{\texttt{continuum\_with\_position\_tests}})
uses the run, its penalty, the pillars and the funded property. It
discards the attached-recipe conjunct.

The selection theorem of III.2 accepts any nonempty family of demands.
Here the demand is a recipe, with lower bounds on the block factors and
numerical requests on the running coupling. The Gaussian limit is
therefore an additional property of the selected run; VIII.1 does not
invoke it separately. Replacing this selection by one with fewer demands
is a separate question (§7).

\subsubsection{1. What this unit proves}\label{what-this-unit-proves-16}

\textbf{Conventions.} We use the level systems and sufficiency condition
of III.1. Throughout, \(G\) is a compact matrix gauge group; ``simple''
means that its real Lie algebra is simple and nonabelian.

\begin{itemize}
\tightlist
\item
  \textbf{Level system.} A level system \(\mathcal L\) fixes stages
  \(i\in\mathbb N\) (finer as \(i\) grows), odd block factors
  \(L_i\ge3\) for the step \(i+1\to i\), and spacings \(a_0=1\),
  \(a_{i+1}=a_i/L_i\), with log-scales \(s_i=\log(1/a_i)\). Its
  smoothing indices satisfy \(J_i\le J_{i+1}\le J_i+1\). Write
  \(\ell_j\) for the block factor assigned to smoothing index \(j\), so
  \(L_i=\ell_{J_{i+1}}\). The level system also fixes the block-spin
  kernel with penalty \(b>0\) and the reference inverse coupling
  \(\hat\beta_i=\bar\beta+s_Gs_i>1\), where \(s_G=11C_G/(24\pi^2)\)
  (\href{https://github.com/yang-mills-group/OpenAI-YM-20261028/blob/3456e5e75cf8d4cd558955a5c7638fba2d68f464/YM/RG/History.lean\#L79}{\texttt{hatBeta}},
  \href{https://github.com/yang-mills-group/OpenAI-YM-20261028/blob/3456e5e75cf8d4cd558955a5c7638fba2d68f464/YM/Gauge/Colors.lean\#L184}{\texttt{slopeConstant}}).
  The remaining compatibility conditions are those of III.1.
\item
  \textbf{Cutoff.} ``Cutoff \(N\)'' means the lattice of spacing
  \(a_N\), renormalized through stages \(N,N-1,\dots\).
\item
  \textbf{Horizon.} A horizon \(H\) assigns plaquette-metric and
  transfer data to each cutoff and stage. Write \(M^\infty_{N,i}\) for
  its translation-invariant metric kernel on \(\mathbb Z^4\),
  \(\widehat M_{N,i}(\xi)\) for its Fourier symbol, and
  \(M^{\mathbb T}_{N,i}\) for its periodization. On a torus
  \(\mathbb T\), the Gaussian quadratic form on 1-forms has matrix
  \(\mathcal K^{\mathbb T}_{N,i}=d^\top M^{\mathbb T}_{N,i}d\).
\item
  \textbf{Price map and sufficiency.} A price map
  \(P=(\alpha_j,\Omega_j,\eta_j)_j\) specifies positive strip widths
  \(\alpha_j\), upper bounds \(\Omega_j\) and positive coercivity bounds
  \(\eta_j\). The full sufficiency condition of III.1 includes strip and
  derivative bounds for all metric and transfer data, inverse and
  factorization identities, symmetry and coercivity on every torus, the
  background-field and fluctuation identities between consecutive
  stages, and the bare terminal condition. In particular, a sufficient
  horizon satisfies \(\widehat M_{N,i}(0)=1\) and
  \(\mathcal K^{\mathbb T}_{N,N}=d^\top d\).
\item
  \textbf{Position-space kernel.}
  \(\mathcal O^H_{N,i}=d^\top M^\infty_{N,i}d\) is the quadratic kernel
  on the links of \(\mathbb Z^4\)
  (\href{https://github.com/yang-mills-group/OpenAI-YM-20261028/blob/3456e5e75cf8d4cd558955a5c7638fba2d68f464/YM/RG08/Flat/Curl.lean\#L18}{\texttt{output}}).
  This is the Hessian convention: no factor \(1/2\) is included in
  \(\mathcal O\).
\item
  \textbf{Kernel norms.} Let \(x(u)\) be the base point of a link or
  plaquette \(u\). For link kernels, put
  \(\lVert F\rVert_\sigma=\sup_u\sup_\phi\sum_{y\in\mathbb Z^4}|F(u,(y,\phi))|e^{\sigma|x(u)-y|_1}\).
  For plaquette kernels, put
  \(\lVert B\rVert^{\rm pl}_\sigma=\sup_u\sum_{v\in\mathbb Z^4\times\mathcal F}|B(u,v)|e^{\sigma|x(u)-x(v)|_1}\),
  where \(\mathcal F\) is the set of plaquette orientations. The latter
  sum includes all target orientations.
\item
  \textbf{Age.} \(\Lambda_{k,i}=5+\log\hat\beta_k+|s_i-s_k|\) is the age
  of stage \(i\) seen from a reporting scale \(k\).
\end{itemize}

\textbf{Proposition III.14.1 (the construction run).} Let \(G\) be
simple and \(b>0\). There exist a run \(R\), a campaign \(T\) for \(R\)
and pillars \(W\) for \((R,T)\) such that:

\begin{enumerate}
\def\labelenumi{\arabic{enumi}.}
\tightlist
\item
  the penalty of the block-spin kernel of \(R\) equals \(b\);
\item
  \(R\) satisfies the full \emph{funded} condition of III.15 relative to
  the local-estimate package \(\Delta_{10}\). Its thirteen clauses (§6)
  include the prescribed local estimates, \(Q(i)\to\infty\), support
  exponent \(100\) at every stage, and the extraction, OPE and thermal
  conditions. The selection uses \(Q(i)=i+2\);
\item
  the native component recipe is \emph{attached} to \(R\). There are a
  sheet \(s\), a price map \(q\), a recipe \(C\) for \(s\) (Proposition
  III.14.2) and floors \(n\ge\max(n_s,E_C)\), where \(n_s\) is the
  sheet's floor and \(E_C\) the combined recipe floor of Lemma 4.14,
  such that:

  \begin{itemize}
  \tightlist
  \item
    \(R\) is built on \(s\) with floors raised to \(n\), with the
    primary price map \(P_{s,q}\) of III.12;
  \item
    both countable request families, before and after the block factors
    are chosen, satisfy the eventual bounds of Lemma 4.16; the recipe's
    post-length bill \(B_C\) also satisfies
    \(B_C(\ell,J_{m+1})\le\hat\beta_m\) for every \(m\);
  \item
    the realized recipe (Lemma 4.14) holds on \(R\)'s own levels and
    horizon.
  \end{itemize}
\end{enumerate}

\emph{Status:} theorem. A run is fixed by choice as
\href{https://github.com/yang-mills-group/OpenAI-YM-20261028/blob/3456e5e75cf8d4cd558955a5c7638fba2d68f464/YM/Independent/ConstructionRun.lean\#L30-L31}{\texttt{fundedRun}}.

\emph{Lean:}
\href{https://github.com/yang-mills-group/OpenAI-YM-20261028/blob/3456e5e75cf8d4cd558955a5c7638fba2d68f464/YM/Independent/ConstructionRun.lean\#L14-L28}{\texttt{YM.Independent.constructionRun}};
\href{https://github.com/yang-mills-group/OpenAI-YM-20261028/blob/3456e5e75cf8d4cd558955a5c7638fba2d68f464/YM/R2I/Campaign/Accepted.lean\#L60-L73}{\texttt{R2I.Funded}};
\href{https://github.com/yang-mills-group/OpenAI-YM-20261028/blob/3456e5e75cf8d4cd558955a5c7638fba2d68f464/YM/RG12/Burst/Wardrobe.lean\#L63}{\texttt{garment10}};
\href{https://github.com/yang-mills-group/OpenAI-YM-20261028/blob/3456e5e75cf8d4cd558955a5c7638fba2d68f464/YM/FullVolume/Gables/Vendee.lean\#L22-L27}{\texttt{Campaign.Pillars}};
\href{https://github.com/yang-mills-group/OpenAI-YM-20261028/blob/3456e5e75cf8d4cd558955a5c7638fba2d68f464/YM/Integration/ActualRunNativeComponentRecipe.lean\#L56-L69}{\texttt{ActualRunNativeComponentRecipe.Attached}}.

\textbf{Proposition III.14.2 (existence of the native component
recipe).} Let \(G\) be simple and let \(s\) be a sheet (the kernel
choice and stage profile, fixed before any block factor). There exist a
price map \(X\) (the \emph{secondary} price map) and floors
\(E^{\flat},E^{\rm sw},E^{\rm du}:\mathbb N\to\mathbb N\) with the
following properties, for every level system \(\mathcal L\) with the
kernel and stage profile of \(s\).

\begin{itemize}
\tightlist
\item
  \textbf{(flat representation)} If \(L_i\ge E^\flat(J_{i+1})\) for all
  \(i\), then for every \(e\) satisfying
  \(\Omega^X_{J_i}\le\log\hat\beta_i\) for \(i\ge e\), there are kernels
  and constants satisfying clauses (a)--(c) of Proposition III.14.3 with
  this \(X\), simultaneously for every sufficient horizon.
\item
  \textbf{(point-group invariance)} If \(L_i\ge E^{\rm sw}(J_{i+1})\)
  for all \(i\), then for every horizon \(H\) sufficient for some price
  map, every \(i\le N\) and every torus, the colour-summed bilinear form
  \(\sum_a\langle A^a,\mathcal K^{\mathbb T}_{N,i}B^a\rangle\) is
  invariant under simultaneous pullback of \(A,B\) by signed
  permutations of the axes fixing the origin. Pullback reverses the sign
  when a link orientation is reversed.
\item
  \textbf{(background-field identities)} If
  \(L_i\ge E^{\rm du}(J_{i+1})\) for all \(i\), then for every such
  \(H\), every \(m<N\) and every torus in the level system, the
  canonical free background-field map of III.3 satisfies the relations
  of Proposition III.14.5 for the fine and coarse metrics
  \((M^{\mathbb T_f}_{N,m+1},M^{\mathbb T_c}_{N,m})\). It therefore
  gives a minimizer of the Gaussian variational problem for the actual
  quadratic forms.
\end{itemize}

\emph{Status:} theorem.

\emph{Lean:}
\href{https://github.com/yang-mills-group/OpenAI-YM-20261028/blob/3456e5e75cf8d4cd558955a5c7638fba2d68f464/YM/Integration/NativeComponentRecipe.lean\#L29-L51}{\texttt{NativeComponentRecipe.Choice}},
\href{https://github.com/yang-mills-group/OpenAI-YM-20261028/blob/3456e5e75cf8d4cd558955a5c7638fba2d68f464/YM/Integration/NativeComponentRecipe.lean\#L54-L59}{\texttt{available}};
the three components follow from
\href{https://github.com/yang-mills-group/OpenAI-YM-20261028/blob/3456e5e75cf8d4cd558955a5c7638fba2d68f464/YM/RG08/Flat/Entrance.lean\#L28-L50}{\texttt{Flat.entrance\_gate}},
\href{https://github.com/yang-mills-group/OpenAI-YM-20261028/blob/3456e5e75cf8d4cd558955a5c7638fba2d68f464/YM/RG08/Cast/Profile.lean\#L16-L34}{\texttt{Cast.assured}}
(with
\href{https://github.com/yang-mills-group/OpenAI-YM-20261028/blob/3456e5e75cf8d4cd558955a5c7638fba2d68f464/YM/Bid/Station.lean\#L54-L55}{\texttt{Bid.Sworn}})
and
\href{https://github.com/yang-mills-group/OpenAI-YM-20261028/blob/3456e5e75cf8d4cd558955a5c7638fba2d68f464/YM/RG08/Aura/Publisher.lean\#L28-L52}{\texttt{Aura.dual\_support}}
(with
\href{https://github.com/yang-mills-group/OpenAI-YM-20261028/blob/3456e5e75cf8d4cd558955a5c7638fba2d68f464/YM/RG08/Aura/Publisher.lean\#L24-L26}{\texttt{Aura.licensed}}).

\textbf{Proposition III.14.3 (ultraviolet rate for the Gaussian
quadratic part).} Let \(G\) be simple and \(s\) a sheet. There exist a
price map \(X\) and a floor \(E:\mathbb N\to\mathbb N\) with the
following property. Let \(\mathcal L\) be a level system with the kernel
and stage profile of \(s\) and with \(L_i\ge E(J_{i+1})\) for all
\(i\ge0\). Let \(e\in\mathbb N\) satisfy
\(\Omega^X_{J_i}\le\log\hat\beta_i\) for \(i\ge e\). Then there are
kernels \(F_i\) (\(i\ge e\)) on the links of \(\mathbb Z^4\) and
constants \(A,B\ge0\), \(c>0\), \(d,l\in\mathbb N\) such that the
following holds for every horizon \(H\) of \(\mathcal L\) sufficient for
some price map, every reporting scale \(k\in\mathbb N\) and every
\(i\ge e\). Write \(w_j=\min(1,\alpha^X_j,\alpha^X_{j+1})\).

\begin{itemize}
\tightlist
\item
  \begin{enumerate}
  \def\labelenumi{(\alph{enumi})}
  \tightlist
  \item
    \(\lVert F_i\rVert_{w_{J_i}}\le B\,\Lambda_{k,i}^{\,l}\);
  \end{enumerate}
\item
  (b) \(\mathcal O^H_{n,i}(u,v)\to F_i(u,v)\) as \(n\to\infty\), for all
  links \(u,v\);
\item
  (c)
  \(\lVert\mathcal O^H_{N,i}-F_i\rVert_{w_{J_i}}\le A\,\Lambda_{k,i}^{\,d}\,(a_N/a_i)^{c}\)
  for all \(N\ge i\).
\end{itemize}

One may take \(c=1/4\), \(d=l=1\),
\(A=34\sigma_0|\mathcal F|(1+s_G)/(1-3^{-1/8})\) and
\(B=\sigma_0|\mathcal F|(1+s_G)\), where
\(\sigma_0=|\mathcal F|^2\lambda_d^2\) and \(\lambda_d\) is a width-one
strip bound for the lattice exterior derivative. These constants depend
only on \(G\) and the lattice, independently of the sheet, level system,
entry stage and horizon.

\emph{Status:} theorem. This is an order-zero estimate (see
\href{https://github.com/yang-mills-group/OpenAI-YM-20261028/blob/3456e5e75cf8d4cd558955a5c7638fba2d68f464/YM/RG08/Flat/Mesh.lean}{\texttt{Mesh.lean}}).

\emph{Lean:}
\href{https://github.com/yang-mills-group/OpenAI-YM-20261028/blob/3456e5e75cf8d4cd558955a5c7638fba2d68f464/YM/RG08/Flat/Entrance.lean\#L13-L21}{\texttt{Flat.Quorum}}
(the three clauses),
\href{https://github.com/yang-mills-group/OpenAI-YM-20261028/blob/3456e5e75cf8d4cd558955a5c7638fba2d68f464/YM/RG08/Flat/Entrance.lean\#L28-L50}{\texttt{Flat.entrance\_gate}}
(existence),
\href{https://github.com/yang-mills-group/OpenAI-YM-20261028/blob/3456e5e75cf8d4cd558955a5c7638fba2d68f464/YM/RG08/Flat/Mesh.lean\#L107-L164}{\texttt{Flat.supplies}},
\href{https://github.com/yang-mills-group/OpenAI-YM-20261028/blob/3456e5e75cf8d4cd558955a5c7638fba2d68f464/YM/RG08/Flat/Mesh.lean\#L75-L101}{\texttt{Flat.meshed}},
\href{https://github.com/yang-mills-group/OpenAI-YM-20261028/blob/3456e5e75cf8d4cd558955a5c7638fba2d68f464/YM/RG08/Flat/Mesh.lean\#L30}{\texttt{Flat.clock}},
\href{https://github.com/yang-mills-group/OpenAI-YM-20261028/blob/3456e5e75cf8d4cd558955a5c7638fba2d68f464/YM/RG08/Flat/Length.lean\#L14}{\texttt{Flat.curve}},
\href{https://github.com/yang-mills-group/OpenAI-YM-20261028/blob/3456e5e75cf8d4cd558955a5c7638fba2d68f464/YM/RG08/Flat/Length.lean\#L19}{\texttt{Flat.Delay}}.

\textbf{Proposition III.14.4 (flat representations at all stages,
abstract form).} Let \(P\) be a price map, \(\theta>0\), and
\(w:\mathbb N\to(0,\infty)\). Let
\((\mathcal L_x,H_x,r_x,N_x,n_x,k_x,Z_x)_{x\in X}\) be a family of level
systems with horizons, integers \(r_x\le N_x\le n_x\), reporting scales
\(k_x\in\mathbb N\) and offsets \(Z_x\ge0\). Put
\(\Lambda_x(i)=2+Z_x+|s_i-s_{k_x}|\), with the scale data understood to
belong to \(\mathcal L_x\). Fix \(A_0\ge0\) and \(d_0\in\mathbb N\),
independent of \(x\). Assume:

\begin{itemize}
\tightlist
\item
  each \(H_x\) is sufficient for \(P\), and its plaquette-transfer
  kernels at cutoff \(N_x\) and stages \(r_x\le i<N_x\) are the
  canonical free kernels of III.3;
\item
  \(w_j\le\alpha_j\) for all \(j\), and for \(r_x\le i<N_x\):

  \begin{itemize}
  \tightlist
  \item
    \(w_{J_i}\le\min(1,\alpha_{J_i},\alpha_{J_{i+1}})\);
  \item
    \(2w_{J_i}\le L_i\,w_{J_{i+1}}\);
  \item
    \(q_i(w_{J_{i+1}})\le L_i^{-\theta}\), where \(q_i\) is the one-step
    factor of Lemma 4.3;
  \end{itemize}
\item
  \(2|\mathcal F|\Omega_{J_{N_x}}\le A_0\Lambda_x(N_x)^{d_0}\) for every
  \(x\).
\end{itemize}

Then there are \(A\ge0\), \(c>0\), \(d\in\mathbb N\) and kernels
\(K_{x,i}\) such that, for all \(x\) and \(r_x\le i\le N_x\),
\(K_{x,i}\) is a flat representation (Definition 4.1) of
\(\mathcal K_{n_x,i}-\mathcal K_{N_x,i}\) at stage \(i\), of width
\(w_{J_i}\) and cost \(A\,\Lambda_x(i)^d(a_{N_x}/a_i)^c\). One may take
\(c=\min(1,\theta/2)\) and \(d=d_0\), with \(A\) depending only on
\(A_0,d_0,\theta\).

\emph{Status:} theorem.

\emph{Lean:}
\href{https://github.com/yang-mills-group/OpenAI-YM-20261028/blob/3456e5e75cf8d4cd558955a5c7638fba2d68f464/YM/RG08/Flat/World.lean\#L165-L174}{\texttt{Flat.World.overlap}},
with
\href{https://github.com/yang-mills-group/OpenAI-YM-20261028/blob/3456e5e75cf8d4cd558955a5c7638fba2d68f464/YM/RG08/Flat/World.lean\#L82-L91}{\texttt{World}},
\href{https://github.com/yang-mills-group/OpenAI-YM-20261028/blob/3456e5e75cf8d4cd558955a5c7638fba2d68f464/YM/RG08/Flat/World.lean\#L101-L110}{\texttt{World.Funding}},
\href{https://github.com/yang-mills-group/OpenAI-YM-20261028/blob/3456e5e75cf8d4cd558955a5c7638fba2d68f464/YM/RG08/Flat/World.lean\#L112-L115}{\texttt{World.Provides}}.

\textbf{Proposition III.14.5 (Schur identity for two cutoffs, and the
flat step).}

\emph{Identity.} Fix one RG step: a fine torus \(\mathbb T_f\) with
\(L\) times the side of the coarse torus \(\mathbb T_c\), the block
average \(Q\) and penalty \(b>0\). Say that \((W;M,M')\) satisfies the
background-field relations when
\[\mathcal K(M)\,W=Q^\top\mathcal K(M'),\qquad QW=1-b^{-1}\mathcal K(M').\]
If \(M\) is symmetric and positive semidefinite, these relations make
\(WA'\) a minimizer of \(\langle A,\mathcal K(M)A\rangle+b|A'-QA|^2\),
with minimum \(\langle A',\mathcal K(M')A'\rangle\). For the algebraic
identity below, positivity is unnecessary: if \((W_1;M_1,M_1')\) and
\((W_2;M_2,M_2')\) satisfy the relations and \(M_2,M_2'\) are symmetric,
then
\[\mathcal K(M_2')-\mathcal K(M_1')=W_2^\top\big(\mathcal K(M_2)-\mathcal K(M_1)\big)W_1 .\]

\emph{Flat step.} Let \(\tau>0\), \(r\ge0\), and let \(B\) be a neutral,
translation-invariant plaquette kernel with
\(\lVert B\rVert^{\rm pl}_\tau\le r\), whose periodization represents
\(\mathcal K(M_2)-\mathcal K(M_1)\) on \(\mathbb T_f\). Suppose \(T,U\)
are order-two strip-bounded curvature-transport kernels with constants
\(q,p\ge0\) in width \(\omega\). Assume \(T(x+Lv,y+v)=T(x,y)\) and the
same identity for \(U\), that their periodizations \(C_2,C_1\) satisfy
\(dW_2=C_2d\), \(dW_1=C_1d\), and that \(\sum_zU((x,\varphi),(z,\psi))\)
is independent of \(x\) for each pair of orientations. If
\(0\le\sigma\le\omega\) and \(2\sigma\le L\tau\), then
\[B'(y,z)=\sum_{x,w}T(x,y)B(x,w)U(w,z),\qquad \lVert B'\rVert^{\rm pl}_\sigma\le\frac{2|\mathcal F|^2qp}{\tau L}r.\]
The kernel \(B'\) is neutral and translation invariant, and its
periodization represents \(\mathcal K(M_2')-\mathcal K(M_1')\) on
\(\mathbb T_c\).

\emph{Sketch:} transpose the second pair of background-field relations
and substitute the first pair; the mixed
\(b^{-1}\mathcal K(M_2')\mathcal K(M_1')\) terms cancel. The
intertwining identities then give the asserted curvature representative.
Neutrality permits subtraction of the second transport kernel at the
first base point, yielding the norm estimate in Lemma 4.3.

\emph{Status:} theorem. No inverse of the curvature form is used.
Preservation of neutrality requires the separate constant-row-sum
condition.

\emph{Lean:}
\href{https://github.com/yang-mills-group/OpenAI-YM-20261028/blob/3456e5e75cf8d4cd558955a5c7638fba2d68f464/YM/RG08/Packets/FlatStep.lean\#L24-L38}{\texttt{Flat.poised}}
(identity),
\href{https://github.com/yang-mills-group/OpenAI-YM-20261028/blob/3456e5e75cf8d4cd558955a5c7638fba2d68f464/YM/RG08/Packets/FlatStep.lean\#L50-L78}{\texttt{Flat.toward}}
(step),
\href{https://github.com/yang-mills-group/OpenAI-YM-20261028/blob/3456e5e75cf8d4cd558955a5c7638fba2d68f464/YM/RG08/Packets/Lines.lean\#L87-L113}{\texttt{Packets.push\_null}},
\href{https://github.com/yang-mills-group/OpenAI-YM-20261028/blob/3456e5e75cf8d4cd558955a5c7638fba2d68f464/YM/RG08/Packets/Lines.lean\#L84-L85}{\texttt{neutralCoin}}.

\textbf{Proposition III.14.6 (rate lemma for fine-to-coarse
recursions).} For each \(x\in X\), fix a scale system with \(L_i\ge3\),
stages \(r_x\le i\le t_x\) and age \(\Lambda_x(i)=2+Z_x+|s_i-s_{k_x}|\),
where \(Z_x\ge0\). Fix \(\theta,c_0>0\), \(A_0\ge0\) and
\(d_0\in\mathbb N\), uniformly in \(x\). Suppose \(f_x(i),g_x(i)\ge0\)
on these stage ranges and

\begin{itemize}
\tightlist
\item
  \(f_x(i)\le L_i^{-\theta}f_x(i+1)+g_x(i)\) for \(r_x\le i<t_x\), and
  \(f_x(t_x)\le g_x(t_x)\);
\item
  \(g_x(i)\le A_0\Lambda_x(i)^{d_0}(a_{t_x}/a_i)^{c_0}\) for
  \(r_x\le i\le t_x\).
\end{itemize}

Then \(f_x(i)\le A\Lambda_x(i)^{d_0}(a_{t_x}/a_i)^\mu\) throughout each
range, where \(\mu=\min(c_0,\theta/2)>0\). One may take
\(A=A_0(K+1)^{d_0}/(1-3^{-\theta'/2})\), with \(\theta'=\theta-\mu>0\)
and \(K=2(d_0+1)/\theta'\).

\emph{Status:} theorem.

\emph{Lean:}
\href{https://github.com/yang-mills-group/OpenAI-YM-20261028/blob/3456e5e75cf8d4cd558955a5c7638fba2d68f464/YM/RG08/Rate/Ascent.lean\#L118-L134}{\texttt{Rate.Stair.fade\_solve}},
\href{https://github.com/yang-mills-group/OpenAI-YM-20261028/blob/3456e5e75cf8d4cd558955a5c7638fba2d68f464/YM/RG08/Rate/Ascent.lean\#L56-L116}{\texttt{poly\_solve}},
\href{https://github.com/yang-mills-group/OpenAI-YM-20261028/blob/3456e5e75cf8d4cd558955a5c7638fba2d68f464/YM/RG08/Rate/Ascent.lean\#L26-L54}{\texttt{slope}},
\href{https://github.com/yang-mills-group/OpenAI-YM-20261028/blob/3456e5e75cf8d4cd558955a5c7638fba2d68f464/YM/RG08/Rate/Shadows.lean\#L90-L97}{\texttt{Fade.cap}}.

\textbf{Further content.} The auxiliary estimates in §4.6 supply
numerical requirements for the recipe. They include the one-step transit
estimate
\href{https://github.com/yang-mills-group/OpenAI-YM-20261028/blob/3456e5e75cf8d4cd558955a5c7638fba2d68f464/YM/RG08/Sideshow/Orbit.lean\#L81-L190}{\texttt{harbor}},
curvature-coefficient lifts
\href{https://github.com/yang-mills-group/OpenAI-YM-20261028/blob/3456e5e75cf8d4cd558955a5c7638fba2d68f464/YM/RG08/Rudder/Sym/Oaks.lean\#L65-L81}{\texttt{arc}}
and
\href{https://github.com/yang-mills-group/OpenAI-YM-20261028/blob/3456e5e75cf8d4cd558955a5c7638fba2d68f464/YM/RG08/Spar/Overhead.lean\#L39-L49}{\texttt{warehouse}},
and formal Gaussian identities
\href{https://github.com/yang-mills-group/OpenAI-YM-20261028/blob/3456e5e75cf8d4cd558955a5c7638fba2d68f464/YM/RG08/Room/Time.lean\#L54-L81}{\texttt{bottle}},
\href{https://github.com/yang-mills-group/OpenAI-YM-20261028/blob/3456e5e75cf8d4cd558955a5c7638fba2d68f464/YM/RG08/Room/Glass.lean\#L98-L111}{\texttt{still}}
and
\href{https://github.com/yang-mills-group/OpenAI-YM-20261028/blob/3456e5e75cf8d4cd558955a5c7638fba2d68f464/YM/RG08/Room/Alight.lean\#L86-L123}{\texttt{absorbed}}.

\subsubsection{2. What it uses}\label{what-it-uses-16}

\textbf{From other units.}

\begin{itemize}
\tightlist
\item
  \textbf{III.1 (block-spin step and its linearization).}

  \begin{itemize}
  \tightlist
  \item
    Level systems: stages, spacings, block factors, the smoothing
    indices \(J\), and \(\hat\beta_i\).
  \item
    Horizons and their sufficiency for a price map, and the metric data
    with kinetic normalization \(\widehat M(0)=1\).
  \item
    The background-field relations, and the identity
    \(\mathcal K(M')=W^\top\mathcal K(M)W+b(1-QW)^\top(1-QW)\).
  \item
    Links:
    \href{https://github.com/yang-mills-group/OpenAI-YM-20261028/blob/3456e5e75cf8d4cd558955a5c7638fba2d68f464/YM/RG/History.lean\#L131-L134}{\texttt{Levels}},
    \href{https://github.com/yang-mills-group/OpenAI-YM-20261028/blob/3456e5e75cf8d4cd558955a5c7638fba2d68f464/YM/PREC/Stacks.lean\#L119-L121}{\texttt{Horizon}},
    \href{https://github.com/yang-mills-group/OpenAI-YM-20261028/blob/3456e5e75cf8d4cd558955a5c7638fba2d68f464/YM/PREC/Stacks.lean\#L132-L139}{\texttt{Horizon.Suffices}},
    \href{https://github.com/yang-mills-group/OpenAI-YM-20261028/blob/3456e5e75cf8d4cd558955a5c7638fba2d68f464/YM/PREC/Stacks.lean\#L123-L128}{\texttt{MapPrices}},
    \href{https://github.com/yang-mills-group/OpenAI-YM-20261028/blob/3456e5e75cf8d4cd558955a5c7638fba2d68f464/YM/PREC/Stacks.lean\#L61-L68}{\texttt{Floor}},
    \href{https://github.com/yang-mills-group/OpenAI-YM-20261028/blob/3456e5e75cf8d4cd558955a5c7638fba2d68f464/YM/PREC/TreeMatch.lean\#L68-L69}{\texttt{Precision.K}},
    \href{https://github.com/yang-mills-group/OpenAI-YM-20261028/blob/3456e5e75cf8d4cd558955a5c7638fba2d68f464/YM/PREC/TreeMatch.lean\#L76-L79}{\texttt{Inputs}},
    \href{https://github.com/yang-mills-group/OpenAI-YM-20261028/blob/3456e5e75cf8d4cd558955a5c7638fba2d68f464/YM/PREC/TreeMatch.lean\#L93-L103}{\texttt{action}}.
  \end{itemize}
\item
  \textbf{III.2 (choosing the constants).}

  \begin{itemize}
  \tightlist
  \item
    Sheets:
    \href{https://github.com/yang-mills-group/OpenAI-YM-20261028/blob/3456e5e75cf8d4cd558955a5c7638fba2d68f464/YM/Price/Intent.lean\#L26-L35}{\texttt{Sheet}}.
  \item
    Selection with a preselected demand:
    \href{https://github.com/yang-mills-group/OpenAI-YM-20261028/blob/3456e5e75cf8d4cd558955a5c7638fba2d68f464/YM/Integration/NativePreselectedPostnumericPeriodRun.lean\#L99-L114}{\texttt{selectRunWithTreasuryPeriod\_preselected\_quote\_orders\_postnumeric}}
    and
    \href{https://github.com/yang-mills-group/OpenAI-YM-20261028/blob/3456e5e75cf8d4cd558955a5c7638fba2d68f464/YM/Integration/NativePreselectedPostnumericPeriodRun.lean\#L31-L46}{\texttt{PostnumericPreselection}}.
    The quantified conclusion used here is Lemma 4.16.
  \end{itemize}
\item
  \textbf{III.3 (free kernels).}
  \href{https://github.com/yang-mills-group/OpenAI-YM-20261028/blob/3456e5e75cf8d4cd558955a5c7638fba2d68f464/YM/Next/FREE_KERNELS/Rack.lean\#L134-L145}{\texttt{retained\_support}}
  gives, for each kernel choice, the following:

  \begin{itemize}
  \tightlist
  \item
    a clip profile \(\gamma\) and a floor \(n\);
  \item
    for every sheet with that kernel, a price map \(X\) such that every
    level system with clip parameters \(\le\gamma\) and
    \(L_i\ge n(J_{i+1})\) (the ``credit'',
    \href{https://github.com/yang-mills-group/OpenAI-YM-20261028/blob/3456e5e75cf8d4cd558955a5c7638fba2d68f464/YM/Next/FREE_KERNELS/Rack.lean\#L87-L88}{\texttt{Credit}})
    carries a horizon sufficient for \(X\), whose transfer data are the
    canonical free data
    (\href{https://github.com/yang-mills-group/OpenAI-YM-20261028/blob/3456e5e75cf8d4cd558955a5c7638fba2d68f464/YM/Next/FREE_KERNELS/Rack.lean\#L74}{\texttt{destination}})
    and whose forms are point-group invariant.
  \end{itemize}
\item
  \textbf{III.8 (kernel norms).} Weighted row and column norms, and
  order-\(m\) strip bounds for kernels from fine to coarse lattices:
  \href{https://github.com/yang-mills-group/OpenAI-YM-20261028/blob/3456e5e75cf8d4cd558955a5c7638fba2d68f464/YM/Pert/Bulk/Atlas.lean\#L28-L29}{\texttt{RowCost}},
  \href{https://github.com/yang-mills-group/OpenAI-YM-20261028/blob/3456e5e75cf8d4cd558955a5c7638fba2d68f464/YM/Pert/Bulk/Atlas.lean\#L37-L42}{\texttt{Atlas}},
  the carrier map
  \href{https://github.com/yang-mills-group/OpenAI-YM-20261028/blob/3456e5e75cf8d4cd558955a5c7638fba2d68f464/YM/Pert/Bulk/Inherit.lean\#L24}{\texttt{carrier}}.
\item
  \textbf{III.12.} The structures
  \href{https://github.com/yang-mills-group/OpenAI-YM-20261028/blob/3456e5e75cf8d4cd558955a5c7638fba2d68f464/YM/FullVolume/Cross.lean\#L30-L48}{\texttt{Run}}
  and
  \href{https://github.com/yang-mills-group/OpenAI-YM-20261028/blob/3456e5e75cf8d4cd558955a5c7638fba2d68f464/YM/FullVolume/Gables/Vendee.lean\#L22-L27}{\texttt{Pillars}},
  and the primary price map \(P_{s,q}\)
  (\href{https://github.com/yang-mills-group/OpenAI-YM-20261028/blob/3456e5e75cf8d4cd558955a5c7638fba2d68f464/YM/FullVolume/Roll/Fiscal.lean\#L18}{\texttt{raging}}).
\item
  \textbf{III.15.}

  \begin{itemize}
  \tightlist
  \item
    The funded property
    \href{https://github.com/yang-mills-group/OpenAI-YM-20261028/blob/3456e5e75cf8d4cd558955a5c7638fba2d68f464/YM/R2I/Campaign/Accepted.lean\#L60-L73}{\texttt{Funded}}.
  \item
    The predicate
    \href{https://github.com/yang-mills-group/OpenAI-YM-20261028/blob/3456e5e75cf8d4cd558955a5c7638fba2d68f464/YM/R2I/Campaign/Clock.lean\#L36-L38}{\texttt{EventuallyPaid}}:
    a property of level systems is eventually paid if it can be enforced
    by raising the entry thresholds of \(\beta\). Here it is applied to
    the trivial property.
  \end{itemize}
\item
  \textbf{I.1.} The group data and the one-loop slope \(s_G\).
\item
  \textbf{Use in VIII.1.}
  \href{https://github.com/yang-mills-group/OpenAI-YM-20261028/blob/3456e5e75cf8d4cd558955a5c7638fba2d68f464/YM/Nonperturbative/Continuum.lean\#L28-L48}{\texttt{continuum\_with\_position\_tests}}
  applies
  \href{https://github.com/yang-mills-group/OpenAI-YM-20261028/blob/3456e5e75cf8d4cd558955a5c7638fba2d68f464/YM/R2I/Campaign/Accepted.lean\#L150-L167}{\texttt{R2IGoals.fromFunded}}
  to the pillars and funded property.
\end{itemize}

\textbf{Standard mathematics.}

\begin{itemize}
\tightlist
\item
  Gaussian block-spin transformation with background field (Balaban,
  \emph{Propagators and renormalization transformations for lattice
  gauge theories I, II}, CMP 95 (1984), CMP 96 (1984); \emph{Averaging
  operations for lattice gauge theories}, CMP 98 (1985);
  Gawędzki--Kupiainen, \emph{Block spin renormalization group for dipole
  gas and \((\nabla\varphi)^4\)}, Ann. Phys. 147 (1983)). The coarse
  quadratic form is the Schur complement obtained by minimizing over the
  fine field.
\item
  Exponentially weighted \(\ell^1\) kernel norms and their behaviour
  under composition and blocking, as in Balaban's propagator estimates.
\item
  A discrete Grönwall argument for \(f_i\le\lambda_if_{i+1}+g_i\).
\item
  Uniqueness of a translation-invariant kernel of finite weighted row
  norm from its periodizations on tori of diverging size (dominated
  convergence).
\end{itemize}

\textbf{Imported constants and assumptions.}

\begin{itemize}
\tightlist
\item
  The endpoint assumes only simplicity of \(G\) and \(b>0\).
\item
  The local-estimate package \(\Delta_{10}\) is defined in this unit, in
  \href{https://github.com/yang-mills-group/OpenAI-YM-20261028/blob/3456e5e75cf8d4cd558955a5c7638fba2d68f464/YM/RG12/Burst/Wardrobe.lean\#L63}{\texttt{Wardrobe.lean}}.
  It is the pair of the predicate
  \href{https://github.com/yang-mills-group/OpenAI-YM-20261028/blob/3456e5e75cf8d4cd558955a5c7638fba2d68f464/YM/RG12/Burst/Wardrobe.lean\#L30-L34}{\texttt{Attic}}
  and the lemma
  \href{https://github.com/yang-mills-group/OpenAI-YM-20261028/blob/3456e5e75cf8d4cd558955a5c7638fba2d68f464/YM/RG12/Burst/Wardrobe.lean\#L36-L61}{\texttt{annex}}:

  \begin{itemize}
  \tightlist
  \item
    the predicate combines one-step connected-activity bounds, tree-sum
    smallness and bounds through any fixed finite range of marks and
    degrees;
  \item
    for a target and expense with expense exponent at most \(-7/16\),
    there is a block-factor threshold such that, for every larger block
    factor, every \(W,\mu>0\), admissible order, stage profile and
    finite range, the predicate holds for all sufficiently large \(x\).
    The threshold for \(x\) may depend on those fixed data.
  \end{itemize}
\item
  Inside the flat chain every hypothesis is discharged in this unit or
  in III.3:

  \begin{itemize}
  \tightlist
  \item
    sufficiency, free transfer data and point-group invariance come from
    III.3 via the shadow (Lemma 4.11);
  \item
    the width conditions follow from the floors (Lemma 4.6);
  \item
    the delay condition follows from the post-length bill (Lemma 4.14).
  \end{itemize}
\end{itemize}

\subsubsection{3. Strategy}\label{strategy-16}

The endpoint applies the selection theorem of III.2 to the demand for a
recipe. First choose the sheet, primary price map, recipe and combined
floor \(E_C\). Next choose the block factors \(\ell_j\) above that
floor. Finally raise the coupling thresholds to pay the bill
\(B_C(\ell,j)\) and the countable numerical requests. The selection
retains the resulting inequalities on the run itself, so Lemma 4.15
realizes the recipe on that run's levels and horizon.

For the Gaussian estimate, compare two cutoffs \(N\le n\) at a common
stage \(i\). At stage \(N\), their metric difference
\(M^\infty_{n,N}-M^\infty_{N,N}\) represents the quadratic-form
difference and has zero row sums in each orientation, since both Fourier
symbols equal the identity at zero momentum. The mixed Schur identity
transports this representative through the two background-field maps. A
transport kernel has row scale \(L^{-2}\), column scale \(L^2\) and
first-difference scale \(L^{-3}\). Neutrality permits subtraction in one
leg, gaining a factor \(L^{-1}\); the distance introduced by subtraction
is absorbed using half of the input exponential weight. The block-factor
floor makes the one-step factor at most \(L_i^{-1/2}\). A discrete
Grönwall estimate absorbs polynomial growth in the age and gives the
rate \((a_N/a_i)^{1/4}\). Passing from curvature representatives to link
kernels and then taking a Cauchy limit gives Proposition III.14.3.

The argument carries representatives modulo kernels annihilated by
\(d^\top(\cdot)d\), without inverting the degenerate curvature form. The
plaquette metrics do have inverses, used in the transport kernels. The
Gaussian quadratic forms are independent of the sufficient horizon,
because the terminal form and the Schur recursion determine them. This
permits the estimate to be proved on a shadow level system with altered
clip parameters, where III.3 supplies canonical free transfer data, and
then transferred to the original horizon.

\subsubsection{4. Steps}\label{steps-16}

\paragraph{4.1 Flat representations and their
propagation}\label{flat-representations-and-their-propagation}

\textbf{Definition 4.1 (flat representation).} Let \(H\) be a horizon of
\(\mathcal L\), \(N\le n\), \(i\le N\), and \(\sigma,p\ge0\). A kernel
\(K\) on plaquettes \(\mathbb Z^4\times\mathcal F\) is a \emph{flat
representation} of \(\mathcal K_{n,i}-\mathcal K_{N,i}\) at stage \(i\),
of width \(\sigma\) and cost \(p\), if:

\begin{itemize}
\tightlist
\item
  \(\lVert K\rVert^{\rm pl}_\sigma\le p\);
\item
  \(\sum_{u\in\mathbb Z^4}K(x,(u,\varphi))=0\) for every \(x\) and every
  orientation \(\varphi\) (\emph{neutral});
\item
  \(K(x+u,y+u)=K(x,y)\);
\item
  for every torus \(\mathbb T\) of the level system at stage \(i\), the
  periodization \(K^{\mathbb T}\) satisfies
  \(d^\top K^{\mathbb T}d=\mathcal K^{\mathbb T}_{n,i}-\mathcal K^{\mathbb T}_{N,i}\).
\end{itemize}

\emph{Lean:}
\href{https://github.com/yang-mills-group/OpenAI-YM-20261028/blob/3456e5e75cf8d4cd558955a5c7638fba2d68f464/YM/RG08/Flat/World.lean\#L36-L40}{\texttt{Flat.FlatRep}},
\href{https://github.com/yang-mills-group/OpenAI-YM-20261028/blob/3456e5e75cf8d4cd558955a5c7638fba2d68f464/YM/RG08/Packets/FlatStep.lean\#L43-L44}{\texttt{Sigil}},
\href{https://github.com/yang-mills-group/OpenAI-YM-20261028/blob/3456e5e75cf8d4cd558955a5c7638fba2d68f464/YM/Pert/Bulk/Stacks.lean\#L14}{\texttt{Folded}},
\href{https://github.com/yang-mills-group/OpenAI-YM-20261028/blob/3456e5e75cf8d4cd558955a5c7638fba2d68f464/YM/Pert/Bulk/NullPair.lean\#L9}{\texttt{neutral}},
\href{https://github.com/yang-mills-group/OpenAI-YM-20261028/blob/3456e5e75cf8d4cd558955a5c7638fba2d68f464/YM/RG08/Packets/Torus.lean\#L18}{\texttt{town}}.

\textbf{Lemma 4.2 (the finest stage).} Let \(H\) be sufficient for
\(P\), \(N\le n\), and \(0<\sigma\le\alpha_{J_N}\). Then
\(M^\infty_{n,N}-M^\infty_{N,N}\) is a flat representation at stage
\(N\) of width \(\sigma\) and cost \(2|\mathcal F|\Omega_{J_N}\).

\emph{Sketch:} each metric has plaquette norm
\(\le|\mathcal F|\Omega_{J_N}\) by its strip bound. The difference is
neutral because \(\widehat M(0)=1\) for both metrics. The torus identity
is \(\mathcal K(M_{n,N})-\mathcal K(M_{N,N})=d^\top(M_{n,N}-M_{N,N})d\).

\emph{Lean:} finest-stage lemma (private,
\href{https://github.com/yang-mills-group/OpenAI-YM-20261028/blob/3456e5e75cf8d4cd558955a5c7638fba2d68f464/YM/RG08/Flat/World.lean}{\texttt{World.lean}}),
via
\href{https://github.com/yang-mills-group/OpenAI-YM-20261028/blob/3456e5e75cf8d4cd558955a5c7638fba2d68f464/YM/RG08/Packets/Present.lean\#L52-L78}{\texttt{turn\_class}};
cost
\href{https://github.com/yang-mills-group/OpenAI-YM-20261028/blob/3456e5e75cf8d4cd558955a5c7638fba2d68f464/YM/RG08/Flat/World.lean\#L24}{\texttt{topTab}}.

\textbf{Lemma 4.3 (one step to a coarser stage).} Let \(H\) be
sufficient for \(P\), with \(i<N\le n\). Assume the plaquette-transfer
kernel of \(H\) at \((N,i)\) is the canonical free kernel. Let \(K\) be
a flat representation at stage \(i+1\) of width \(u>0\) and cost
\(p\ge0\), and let \(0<\sigma\le\min(1,\alpha_{J_i},\alpha_{J_{i+1}})\)
with \(2\sigma\le L_iu\). Then
\[K'(y,z)=\sum_{x,w}T_{n,i}(x,y)\,K(x,w)\,T_{N,i}(w,z)\] is a flat
representation at stage \(i\) of width \(\sigma\) and cost
\(q_i(u)\,p\), where:

\begin{itemize}
\tightlist
\item
  \(T_{m,i}=(M^\infty_{m,i+1})^{-1}\mathcal T_{m,i}M^\infty_{m,i}\) is
  the curvature-transport kernel at cutoff \(m\), with
  \(\mathcal T_{m,i}\) the installed plaquette-transfer kernel and
  products denoting kernel composition;
\item
  \(q_i(u)=2|\mathcal F|^2h_i^2/(uL_i)\);
\item
  \(h_i=\exp(4)|\mathcal F|^2\Omega_{J_{i+1}}^2\Omega_{J_i}\).
\end{itemize}

\emph{Sketch:} the four properties come from four inputs.

\begin{itemize}
\tightlist
\item
  \emph{Algebra.} The torus identity follows from the Schur identity and
  from the intertwining \(dW_m=C_md\)
  (\href{https://github.com/yang-mills-group/OpenAI-YM-20261028/blob/3456e5e75cf8d4cd558955a5c7638fba2d68f464/YM/Pert/Bulk/Inherit.lean\#L32-L34}{\texttt{inherited\_precision}}).
  These give
  \(\mathcal K_{n,i}-\mathcal K_{N,i}=d^\top C_n^\top K^{\mathbb T}C_N d\),
  and \(C_n^\top K^{\mathbb T}C_N\) is the periodization of \(K'\)
  (\href{https://github.com/yang-mills-group/OpenAI-YM-20261028/blob/3456e5e75cf8d4cd558955a5c7638fba2d68f464/YM/RG08/Packets/Pushed.lean\#L116-L122}{\texttt{town\_hall}}).
\item
  \emph{Bounds on the transport kernels.} Both are order-2 strip-bounded
  with constant \(h_i\)
  (\href{https://github.com/yang-mills-group/OpenAI-YM-20261028/blob/3456e5e75cf8d4cd558955a5c7638fba2d68f464/YM/Pert/Bulk/Coin.lean\#L40}{\texttt{response}},
  \href{https://github.com/yang-mills-group/OpenAI-YM-20261028/blob/3456e5e75cf8d4cd558955a5c7638fba2d68f464/YM/Pert/Bulk/Keys.lean\#L14}{\texttt{allotment}}).
\item
  \emph{Neutrality.} The free plaquette-transfer kernel has a
  zero-momentum symbol independent of the fine residue class, and
  therefore position-independent row sums. Composition with the
  translation-invariant metric kernels preserves this property for
  \(T_{N,i}\)
  (\href{https://github.com/yang-mills-group/OpenAI-YM-20261028/blob/3456e5e75cf8d4cd558955a5c7638fba2d68f464/YM/RG08/Packets/Swing.lean\#L75}{\texttt{travels}}).
  Thus constants of each orientation go to constants and neutrality
  passes to \(K'\)
  (\href{https://github.com/yang-mills-group/OpenAI-YM-20261028/blob/3456e5e75cf8d4cd558955a5c7638fba2d68f464/YM/RG08/Packets/Swing.lean\#L77-L104}{\texttt{invariant}}).
\item
  \emph{Norm.} Neutrality replaces \(T_{N,i}(w,z)\) by
  \(T_{N,i}(w,z)-T_{N,i}(\bar w,z)\), where \(\bar w\) has the base
  point of \(x\) and the orientation of \(w\). In the weighted row
  estimate this costs \(rh_i/L_i^3\), where \(r\) is the \(\ell^1\)
  distance between the base points of \(x,w\). The inequality
  \(re^{\sigma r/L_i}\le(2/u)e^{ur}\) for \(r\ge0\) follows from
  \(2\sigma\le L_iu\). The column estimate for \(T_{n,i}\) contributes
  \(h_iL_i^2\); summing orientations gives the factor
  \(2|\mathcal F|^2h_i^2/(uL_i)\).
\end{itemize}

\emph{Lean:} one-step lemma (private,
\href{https://github.com/yang-mills-group/OpenAI-YM-20261028/blob/3456e5e75cf8d4cd558955a5c7638fba2d68f464/YM/RG08/Flat/World.lean}{\texttt{World.lean}});
\href{https://github.com/yang-mills-group/OpenAI-YM-20261028/blob/3456e5e75cf8d4cd558955a5c7638fba2d68f464/YM/RG08/Packets/FlatStep.lean\#L50-L78}{\texttt{toward}},
\href{https://github.com/yang-mills-group/OpenAI-YM-20261028/blob/3456e5e75cf8d4cd558955a5c7638fba2d68f464/YM/RG08/Packets/FlatStep.lean\#L24-L38}{\texttt{poised}},
\href{https://github.com/yang-mills-group/OpenAI-YM-20261028/blob/3456e5e75cf8d4cd558955a5c7638fba2d68f464/YM/RG08/Packets/Lines.lean\#L87-L113}{\texttt{push\_null}},
\href{https://github.com/yang-mills-group/OpenAI-YM-20261028/blob/3456e5e75cf8d4cd558955a5c7638fba2d68f464/YM/RG08/Flat/World.lean\#L23}{\texttt{quote}}.

\textbf{Lemma 4.4 (induction over stages).} Under the hypotheses of
Proposition III.14.4 there are flat representations \(K_{x,i}\)
(\(r_x\le i\le N_x\)) of width \(w_{J_i}\) and cost
\(c_x(i)=2|\mathcal F|\Omega_{J_{N_x}}(a_{N_x}/a_i)^\theta\).

\emph{Sketch:} the costs are defined explicitly, so that
\(c_x(i)=L_i^{-\theta}c_x(i+1)\). The proof is downward induction on
\(N_x-i\):

\begin{itemize}
\tightlist
\item
  start with Lemma 4.2;
\item
  apply Lemma 4.3 with \(u=w_{J_{i+1}}\), \(\sigma=w_{J_i}\);
\item
  use \(q_i(w_{J_{i+1}})\le L_i^{-\theta}\) and monotonicity of flat
  representations in the cost.
\end{itemize}

The representatives for all \((x,i)\) are then fixed by choice.

\emph{Lean:} downward-induction lemma (private,
\href{https://github.com/yang-mills-group/OpenAI-YM-20261028/blob/3456e5e75cf8d4cd558955a5c7638fba2d68f464/YM/RG08/Flat/World.lean}{\texttt{World.lean}}).

\textbf{Lemma 4.5 (the rate; proof of Proposition III.14.4).} Write
\(g_x(i)=2|\mathcal F|\Omega_{J_{N_x}}\mathbf 1_{i=N_x}\). The costs of
Lemma 4.4 satisfy the recursion of Proposition III.14.6 with forcing
\(g\). By the last hypothesis of III.14.4, \(g\) has the required form
with \(c_0=1\) (it vanishes below the top). Proposition III.14.6 then
gives cost \(\le A\Lambda_x(i)^{d_0}(a_{N_x}/a_i)^{\min(1,\theta/2)}\).

\emph{Proof of III.14.6.}

\begin{enumerate}
\def\labelenumi{\arabic{enumi}.}
\tightlist
\item
  Multiply by \((a_i/a_t)^{\mu}\); the recursion keeps its form with
  \(\theta'=\theta-\mu\ge\theta/2\).
\item
  Prove \(\tilde f_x(i)\le\tilde B\,(K+\Lambda_x(i))^{d_0}\) by downward
  induction, with \(\tilde B=A_0/(1-3^{-\theta'/2})\).
\item
  The step uses the polynomial estimate: since
  \(\Lambda_x(i+1)\le\Lambda_x(i)+\log L_i\), one has
  \((K+\Lambda_x(i+1))^{d}L_i^{-\theta'}\le(K+\Lambda_x(i))^{d}L_i^{-\theta'/2}\)
  when \(d\le K\theta'/2\). Then \(L_i^{-\theta'/2}\le3^{-\theta'/2}\),
  and \(\tilde B3^{-\theta'/2}+A_0=\tilde B\).
\end{enumerate}

Finally, \(\Lambda_x(i)\ge2\) gives
\((K+\Lambda_x(i))^{d_0}\le(K+1)^{d_0}\Lambda_x(i)^{d_0}\). Undoing the
tilt gives the stated value of \(A\).

\emph{Lean:}
\href{https://github.com/yang-mills-group/OpenAI-YM-20261028/blob/3456e5e75cf8d4cd558955a5c7638fba2d68f464/YM/RG08/Flat/World.lean\#L165-L174}{\texttt{World.overlap}},
\href{https://github.com/yang-mills-group/OpenAI-YM-20261028/blob/3456e5e75cf8d4cd558955a5c7638fba2d68f464/YM/RG08/Rate/Ascent.lean\#L118-L134}{\texttt{fade\_solve}},
\href{https://github.com/yang-mills-group/OpenAI-YM-20261028/blob/3456e5e75cf8d4cd558955a5c7638fba2d68f464/YM/RG08/Rate/Ascent.lean\#L56-L116}{\texttt{poly\_solve}},
\href{https://github.com/yang-mills-group/OpenAI-YM-20261028/blob/3456e5e75cf8d4cd558955a5c7638fba2d68f464/YM/RG08/Rate/Ascent.lean\#L26-L54}{\texttt{slope}},
\href{https://github.com/yang-mills-group/OpenAI-YM-20261028/blob/3456e5e75cf8d4cd558955a5c7638fba2d68f464/YM/RG08/Rate/Shadows.lean\#L83-L87}{\texttt{Fade.of\_top}},
\href{https://github.com/yang-mills-group/OpenAI-YM-20261028/blob/3456e5e75cf8d4cd558955a5c7638fba2d68f464/YM/RG08/Rate/Stair.lean\#L79-L82}{\texttt{drop\_length}}.

\paragraph{4.2 Choice of widths and passage to position
space}\label{choice-of-widths-and-passage-to-position-space}

\textbf{Lemma 4.6 (choice of widths).} Let \(G\) be simple, \(X\) a
price map and \(w_j=\min(1,\alpha^X_j,\alpha^X_{j+1})\). Put
\[\beta_{j,j'}=2|\mathcal F|^2h_{j,j'}^2/w_j,\qquad h_{j,j'}=\exp(4)|\mathcal F|^2(\Omega^X_j)^2\Omega^X_{j'},\qquad E^\flat(j)=\lceil 2/w_j+\beta_{j,j}^2+\beta_{j,j-1}^2\rceil.\]
Here \(j-1\) is truncated at zero. Suppose \(L_i\ge E^\flat(J_{i+1})\)
and \(\Omega^X_{J_i}\le\log\hat\beta_i\) for every \(i\ge e\). With
\(q_i\) computed from \(X\), the following hold for \(i\ge e\):

\begin{itemize}
\tightlist
\item
  \(w_{J_i}\le\min(1,\alpha^X_{J_i},\alpha^X_{J_{i+1}})\);
\item
  \(2w_{J_i}\le L_iw_{J_{i+1}}\);
\item
  \(q_i(w_{J_{i+1}})\le L_i^{-1/2}\);
\item
  \(\Omega^X_{J_j}\le(1+s_G)\Lambda_{k,j}\) for all \(k\) and all
  \(j\ge e\).
\end{itemize}

\emph{Sketch:} since \(J_i\in\{J_{i+1},J_{i+1}-1\}\), the constant
\(h_i\) is one of \(h_{j,j},h_{j,j-1}\) with \(j=J_{i+1}\). Then
\(q_i(w_j)=\beta/L_i\le L_i^{-1/2}\) because \(\beta^2\le L_i\). Next,
\(2w_{J_i}\le2\le L_iw_j\). For the age cap: \(\hat\beta\) is affine in
\(s\), so \(\log\hat\beta_j\le\log\hat\beta_k+s_G|s_j-s_k|\).

\emph{Lean:}
\href{https://github.com/yang-mills-group/OpenAI-YM-20261028/blob/3456e5e75cf8d4cd558955a5c7638fba2d68f464/YM/RG08/Flat/Length.lean\#L64-L92}{\texttt{spacing\_choice}},
\href{https://github.com/yang-mills-group/OpenAI-YM-20261028/blob/3456e5e75cf8d4cd558955a5c7638fba2d68f464/YM/RG08/Flat/Length.lean\#L38-L62}{\texttt{universeClock}},
\href{https://github.com/yang-mills-group/OpenAI-YM-20261028/blob/3456e5e75cf8d4cd558955a5c7638fba2d68f464/YM/RG08/Flat/Length.lean\#L29}{\texttt{work\ items}},
\href{https://github.com/yang-mills-group/OpenAI-YM-20261028/blob/3456e5e75cf8d4cd558955a5c7638fba2d68f464/YM/RG08/Flat/Length.lean\#L25}{\texttt{entrance}},
\href{https://github.com/yang-mills-group/OpenAI-YM-20261028/blob/3456e5e75cf8d4cd558955a5c7638fba2d68f464/YM/RG08/Flat/Length.lean\#L21-L23}{\texttt{before}},
\href{https://github.com/yang-mills-group/OpenAI-YM-20261028/blob/3456e5e75cf8d4cd558955a5c7638fba2d68f464/YM/RG08/Flat/Mesh.lean\#L32-L41}{\texttt{Spacing}}.

\textbf{Lemma 4.7 (the family estimate).} Assume the conclusions of
Lemma 4.6 and let \(H\) be sufficient for \(X\). Index the family by
\(x=(k,N,n)\) with \(k\in\mathbb N\) and \(e\le N\le n\), and put
\(r_x=e\), \(\theta=1/2\) and \(Z_x=3+\log\hat\beta_k\). Then the
numerical hypotheses of Proposition III.14.4 hold with
\(A_0=2|\mathcal F|(1+s_G)\) and \(d_0=1\). If the plaquette-transfer
data at every \((N,i)\) with \(e\le i<N\) are the canonical free data,
all its hypotheses hold.

\emph{Sketch:} the first three width estimates are exactly the required
one-step inequalities. The age bound gives
\(2|\mathcal F|\Omega_{J_N}\le A_0\Lambda_{k,N}\) uniformly in
\(k,N,n\).

\emph{Lean:}
\href{https://github.com/yang-mills-group/OpenAI-YM-20261028/blob/3456e5e75cf8d4cd558955a5c7638fba2d68f464/YM/RG08/Flat/Mesh.lean\#L45-L54}{\texttt{berth}},
\href{https://github.com/yang-mills-group/OpenAI-YM-20261028/blob/3456e5e75cf8d4cd558955a5c7638fba2d68f464/YM/RG08/Flat/Mesh.lean\#L56-L68}{\texttt{Spacing.funded}}.

\textbf{Lemma 4.8 (from a representative to the position-space kernel).}
Let \(0\le\sigma,\varepsilon\le1\) and \(p,C\ge0\). Suppose both metric
kernels are strip-bounded in width \(\sigma\) with constant \(p\), and
that \(K\) represents \(\mathcal K_{n,i}-\mathcal K_{N,i}\) on tori
whose sides tend to infinity, with
\(\lVert K\rVert^{\rm pl}_\varepsilon\le C\). Then
\[\mathcal O_{n,i}-\mathcal O_{N,i}=d^\top K\,d\quad\text{on the links of }\mathbb Z^4,\qquad \lVert\mathcal O_{n,i}-\mathcal O_{N,i}\rVert_\varepsilon\le\sigma_0C,\qquad \sigma_0=|\mathcal F|^2\lambda_d^2,\]
where \(\lambda_d\) is a width-one strip bound for \(d\).

\emph{Sketch:} \(\mathcal O_{n,i}-\mathcal O_{N,i}\) and \(d^\top Kd\)
have finite row norms and the same periodizations on the given tori.
Dominated convergence as the torus side tends to infinity identifies the
kernels. Composition with the two exterior derivatives gives the factor
\(\sigma_0\).

\emph{Lean:}
\href{https://github.com/yang-mills-group/OpenAI-YM-20261028/blob/3456e5e75cf8d4cd558955a5c7638fba2d68f464/YM/RG08/Flat/Curl.lean\#L101-L115}{\texttt{resolve}},
\href{https://github.com/yang-mills-group/OpenAI-YM-20261028/blob/3456e5e75cf8d4cd558955a5c7638fba2d68f464/YM/RG08/Flat/Curl.lean\#L89-L90}{\texttt{output\_cost}},
\href{https://github.com/yang-mills-group/OpenAI-YM-20261028/blob/3456e5e75cf8d4cd558955a5c7638fba2d68f464/YM/RG08/Flat/Curl.lean\#L66}{\texttt{sumToken}}.

\textbf{Lemma 4.9 (rate for pairs of cutoffs).} Under the conclusions of
Lemma 4.6, with \(H\) sufficient for \(X\) and with canonical free
plaquette-transfer data at all \((N,i)\), \(e\le i<N\), the constant
\(A\) of Proposition III.14.3 satisfies
\[\lVert\mathcal O_{n,i}-\mathcal O_{N,i}\rVert_{w_{J_i}}\le A\,\Lambda_{k,i}\,(a_N/a_i)^{1/4}\qquad\text{for all }k\text{ and }e\le i\le N\le n.\]

\emph{Sketch:} apply Proposition III.14.4 to the family of Lemma 4.7,
then Lemma 4.8.

\emph{Lean:}
\href{https://github.com/yang-mills-group/OpenAI-YM-20261028/blob/3456e5e75cf8d4cd558955a5c7638fba2d68f464/YM/RG08/Flat/Mesh.lean\#L75-L101}{\texttt{meshed}}.

\textbf{Lemma 4.10 (the limit kernel).} Under the same hypotheses there
are kernels \(F_i\) satisfying (a)--(c) of Proposition III.14.3, in the
same width \(w_{J_i}\).

\emph{Sketch:} for fixed \(i\), Lemma 4.9 gives a Cauchy modulus
\(A\Lambda_{k,i}(a_n/a_i)^{1/4}\to0\) in the weighted row norm.
Completeness of the scalar entries gives a pointwise limit, and
finite-sum bounds followed by monotone convergence preserve the row
estimates
(\href{https://github.com/yang-mills-group/OpenAI-YM-20261028/blob/3456e5e75cf8d4cd558955a5c7638fba2d68f464/YM/RG08/Packets/Limit.lean\#L51-L59}{\texttt{solvent}}).
This gives norm convergence and (c). For \(n\ge i\),
\(\lVert\mathcal O_{n,i}\rVert_{w_{J_i}}\le\sigma_0|\mathcal F|\Omega_{J_i}\le\sigma_0|\mathcal F|(1+s_G)\Lambda_{k,i}\),
which passes to the limit and gives (a). Uniqueness of pointwise limits
makes \(F_i\) independent of \(k\).

\emph{Lean:}
\href{https://github.com/yang-mills-group/OpenAI-YM-20261028/blob/3456e5e75cf8d4cd558955a5c7638fba2d68f464/YM/RG08/Flat/Mesh.lean\#L107-L164}{\texttt{supplies}}.

\paragraph{4.3 Uniqueness of the Gaussian part; the three recipe
components}\label{uniqueness-of-the-gaussian-part-the-three-recipe-components}

\textbf{Lemma 4.11 (the Gaussian quadratic forms do not depend on the
horizon).} Let \(H,H'\) be horizons, each sufficient for some price map,
on level systems \(\mathcal L,\mathcal L'\) that differ at most in the
clip parameters of their stages; for example \(\mathcal L'\) may be the
\emph{shadow} of \(\mathcal L\), with clip profile replaced by
\(\gamma\). Then
\(\mathcal K^{\mathbb T}_{N,i}(H)=\mathcal K^{\mathbb T}_{N,i}(H')\) for
all \(i\le N\) and all tori, and
\(\mathcal O^H_{N,i}=\mathcal O^{H'}_{N,i}\).

\emph{Sketch:} downward induction from \(i=N\), where both forms equal
\(d^\top d\). If the forms agree at stage \(i+1\), the Schur identity
gives
\(\mathcal K_{N,i}(H)-\mathcal K_{N,i}(H')=W^\top\cdot0\cdot W'=0\); the
clip parameters do not enter the step data. This proves equality on the
tori of the level system. Lemma 4.8 then identifies the position-space
kernels; periodizing that equality gives equality on every torus.

\emph{Lean:}
\href{https://github.com/yang-mills-group/OpenAI-YM-20261028/blob/3456e5e75cf8d4cd558955a5c7638fba2d68f464/YM/RG08/Flat/Agreement.lean\#L23-L36}{\texttt{same\_path}},
\href{https://github.com/yang-mills-group/OpenAI-YM-20261028/blob/3456e5e75cf8d4cd558955a5c7638fba2d68f464/YM/RG08/Flat/Agreement.lean\#L38-L43}{\texttt{same\_output}},
\href{https://github.com/yang-mills-group/OpenAI-YM-20261028/blob/3456e5e75cf8d4cd558955a5c7638fba2d68f464/YM/RG08/Flat/Agreement.lean\#L19-L21}{\texttt{same\_terminal}},
\href{https://github.com/yang-mills-group/OpenAI-YM-20261028/blob/3456e5e75cf8d4cd558955a5c7638fba2d68f464/YM/RG08/Flat/Shadow.lean\#L30-L50}{\texttt{shadow}},
\href{https://github.com/yang-mills-group/OpenAI-YM-20261028/blob/3456e5e75cf8d4cd558955a5c7638fba2d68f464/YM/RG08/Flat/Shadow.lean\#L52-L75}{\texttt{same\_shadow}}.

\textbf{Lemma 4.12 (the flat component; proof of Proposition III.14.3).}
For a simple \(G\) and a sheet \(s\), take the clip profile \(\gamma\)
and floor \(E_{\rm free}\) from III.3, and obtain \(X\) by applying its
sufficiency construction to the sheet with clip profile \(\gamma\). Set
\(E=\max(E^\flat,E_{\rm free})\). Given \(\mathcal L\) and \(e\) as in
Proposition III.14.3, the shadow \(\mathcal L'\) meets the assumptions
of III.3, so it carries a horizon \(Y\) sufficient for \(X\) with
canonical free transfer data. Apply Lemmas 4.6--4.10 to
\((\mathcal L',Y)\); the delay condition is unchanged, since
\(\hat\beta\) and the scale data are unchanged. Transfer (b) and (c) to
every sufficient horizon \(H\) of \(\mathcal L\) by Lemma 4.11. The
floor \(E\) is required at every stage because III.3 requires its floor
everywhere.

\emph{Lean:}
\href{https://github.com/yang-mills-group/OpenAI-YM-20261028/blob/3456e5e75cf8d4cd558955a5c7638fba2d68f464/YM/RG08/Flat/Entrance.lean\#L28-L50}{\texttt{entrance\_gate}},
\href{https://github.com/yang-mills-group/OpenAI-YM-20261028/blob/3456e5e75cf8d4cd558955a5c7638fba2d68f464/YM/RG08/Flat/Shadow.lean\#L21-L23}{\texttt{rearrange}}.

\textbf{Lemma 4.13 (the other two components of the recipe).} For each
sheet, the point-group and background-field conclusions of Proposition
III.14.2 hold above suitable floors. Both use the shadow construction of
Lemma 4.12. Their floors may be chosen to be \(E_{\rm free}\); the flat
component also requires \(E^\flat\).

\begin{itemize}
\tightlist
\item
  \emph{Point-group invariance.} III.3 supplies point-group invariance
  for the shadow horizon, and the forms coincide by Lemma 4.11.
\item
  \emph{Background-field identities.} On the shadow the free
  background-field map satisfies the relations of Proposition III.14.5,
  because its transfer data are the free data. The shadow's quadratic
  forms equal those of \(H\) (Lemma 4.11), and the free data of
  \(\mathcal L\) and of its shadow coincide. Hence the same relations
  hold for \(H\).
\end{itemize}

\emph{Lean:}
\href{https://github.com/yang-mills-group/OpenAI-YM-20261028/blob/3456e5e75cf8d4cd558955a5c7638fba2d68f464/YM/RG08/Cast/Profile.lean\#L16-L34}{\texttt{Cast.assured}},
\href{https://github.com/yang-mills-group/OpenAI-YM-20261028/blob/3456e5e75cf8d4cd558955a5c7638fba2d68f464/YM/RG08/Aura/Publisher.lean\#L28-L52}{\texttt{Aura.dual\_support}},
\href{https://github.com/yang-mills-group/OpenAI-YM-20261028/blob/3456e5e75cf8d4cd558955a5c7638fba2d68f464/YM/Integration/NativeComponentRecipe.lean\#L54-L59}{\texttt{available}}.

\paragraph{4.4 Realizing the recipe on a
run}\label{realizing-the-recipe-on-a-run}

\textbf{Lemma 4.14 (realization).} Let \(s\) be a sheet, \(q\) a price
map, \(C\) a recipe for \(s\) with secondary price map \(X\), and
\(\mathcal L\) a level system with the kernel and profile of \(s\). Let
\(H\) be sufficient for the primary price map \(P_{s,q}\). Let \(E_C\)
be the maximum of \(E^\flat,E^{\rm sw},E^{\rm du}\) and the two
numerical floors in §4.6, and let \(B_C(\ell,j)\) be the maximum of the
two post-length bills there. Assume

\begin{itemize}
\tightlist
\item
  \(\ell_j\ge E_C(j)\) for every smoothing index \(j\);
\item
  \(\hat\beta_m\ge B_C(\ell,J_{m+1})\) for every stage \(m\).
\end{itemize}

Then on \((\mathcal L,H)\):

\begin{itemize}
\tightlist
\item
  both numerical floor and bill conditions hold;
\item
  the background-field identities and point-group invariance hold;
\item
  Proposition III.14.3 holds with \(e=0\).
\end{itemize}

\emph{Sketch:} the floor bounds apply to \(L_i=\ell_{J_{i+1}}\) and to
all smoothing-index requirements. For the delay, the second bill is the
exponential of a finite bound \(S_i\) for the actual patch, with
\(\Omega^X_q\le S_i\) for every \(q\le J_{i+1}\). Its payment gives
\(S_i\le\log\hat\beta_i\)
(\href{https://github.com/yang-mills-group/OpenAI-YM-20261028/blob/3456e5e75cf8d4cd558955a5c7638fba2d68f464/YM/RG08/Yamen/Penny.lean\#L56-L64}{\texttt{Yamen.hatch}}).
Since \(J_i\le J_{i+1}\), the delay holds with \(e=0\). Apply the three
components of \(C\); the fixed chart of III.3 supplies the remaining
local-chart condition.

\emph{Lean:}
\href{https://github.com/yang-mills-group/OpenAI-YM-20261028/blob/3456e5e75cf8d4cd558955a5c7638fba2d68f464/YM/Integration/NativeComponentRecipe.lean\#L132-L182}{\texttt{realize}},
\href{https://github.com/yang-mills-group/OpenAI-YM-20261028/blob/3456e5e75cf8d4cd558955a5c7638fba2d68f464/YM/Integration/NativeComponentRecipe.lean\#L117-L128}{\texttt{Realized}},
\href{https://github.com/yang-mills-group/OpenAI-YM-20261028/blob/3456e5e75cf8d4cd558955a5c7638fba2d68f464/YM/Integration/NativeComponentRecipe.lean\#L61-L65}{\texttt{floor}},
\href{https://github.com/yang-mills-group/OpenAI-YM-20261028/blob/3456e5e75cf8d4cd558955a5c7638fba2d68f464/YM/Integration/NativeComponentRecipe.lean\#L75-L79}{\texttt{post}},
\href{https://github.com/yang-mills-group/OpenAI-YM-20261028/blob/3456e5e75cf8d4cd558955a5c7638fba2d68f464/YM/Integration/NativeComponentRecipe.lean\#L71-L72}{\texttt{numeric}}.

\textbf{Lemma 4.15 (binding to the run).} Let \(s,q,C\) be as in Lemma
4.14 and \(n:\mathbb N\to\mathbb N\) with \(n\ge E_C\). Let \(R\) be a
run whose sheet is \(s\) with floors raised to \(n\), whose price map is
\(P_{s,q}\), and whose levels pay \(B_C\). Then the recipe is realized
on \(R\)'s own levels and horizon.

\emph{Sketch:} the sheet and price-map identities identify the kernel,
profile and sufficient horizon required in Lemma 4.14. The run's
block-factor inequalities give \(\ell_j\ge n(j)\ge E_C(j)\). Apply that
lemma. The retained selection data of III.2 contain exactly these
identities and inequalities.

\emph{Lean:}
\href{https://github.com/yang-mills-group/OpenAI-YM-20261028/blob/3456e5e75cf8d4cd558955a5c7638fba2d68f464/YM/Integration/ActualRunNativeComponentRecipe.lean\#L25-L54}{\texttt{on\_run}},
\href{https://github.com/yang-mills-group/OpenAI-YM-20261028/blob/3456e5e75cf8d4cd558955a5c7638fba2d68f464/YM/Integration/ActualRunNativeComponentRecipe.lean\#L71-L77}{\texttt{of\_preselection}}.

\paragraph{4.5 The run}\label{the-run}

\textbf{Lemma 4.16 (selection with a preselected demand; III.2).} Fix
the comparison and custody data used in III.2. Let \(G\) be simple,
\(b>0\), \(\Lambda\) a property of level systems that can be enforced by
raising coupling-entry thresholds, and \(\Delta\) a local-estimate
package. For every sheet \(s\) and price map \(q\), suppose a nonempty
demand set \(\mathcal D(s,q)\) is given. For each
\(C\in\mathcal D(s,q)\), prescribe an integer floor \(E_C(j)\), real
requests \(U_C(r,j)\), and real post-length requests \(V_C(\ell,r,j)\)
depending on the block-factor sequence \(\ell\).

Then there are a run \(R\), a campaign \(T\), a sheet \(s\), a price map
\(q\), a choice \(C\in\mathcal D(s,q)\) and an integer floor
\(n\ge\max(n_s,E_C)\) such that:

\begin{itemize}
\tightlist
\item
  the penalty is \(b\), \(\Lambda\) holds on the levels of \(R\),
  pillars exist, and \(R\) is funded for \(\Delta\);
\item
  \(R\) has the sheet \(s\) with floor raised to \(n\), and price map
  \(P_{s,q}\);
\item
  \(\log\log\hat\beta_m\ge0\) for every \(m\);
\item
  for each of the request families \(F=U_C\) and
  \(F(r,j)=V_C(\ell,r,j)\), and every fixed \(r\), there is \(m_0\) such
  that \(|F(r,j)|\le1+\log\log\hat\beta_m\) whenever \(m\ge m_0\) and
  \(j\le J_{m+1}\);
\item
  \(V_C(\ell,0,J_{m+1})\le\hat\beta_m\) for every \(m\).
\end{itemize}

The comparison and custody data have canonical choices depending only on
\(G\), used in Lemma 4.17. The additional conclusions about orders,
periods and extraction in III.2 are not needed here.

\emph{Lean:}
\href{https://github.com/yang-mills-group/OpenAI-YM-20261028/blob/3456e5e75cf8d4cd558955a5c7638fba2d68f464/YM/Integration/NativePreselectedPostnumericPeriodRun.lean\#L99-L114}{\texttt{selectRunWithTreasuryPeriod\_preselected\_quote\_orders\_postnumeric}};
the plain version is
\href{https://github.com/yang-mills-group/OpenAI-YM-20261028/blob/3456e5e75cf8d4cd558955a5c7638fba2d68f464/YM/R2I/Campaign/Accepted.lean\#L75-L84}{\texttt{selectRun}}.
\emph{(statement only; the construction behind it belongs to III.2)}

\textbf{Lemma 4.17 (assembly of the construction run).} Proposition
III.14.1 holds.

\emph{Sketch:} use the canonical comparison and custody data and apply
Lemma 4.16 with \(\Lambda\) identically true, \(\Delta=\Delta_{10}\),
and \(\mathcal D(s,q)\) the set of recipes for \(s\). Proposition
III.14.2 makes this set nonempty. Use floor \(E_C\), reciprocal-width
requests \(U_C(r,j)=\rho_C(r,j)^{-1}\), where \(\rho_C(r,j)>0\) is the
recipe width for demand \(r\) at smoothing index \(j\), and
\(V_C(\ell,r,j)=B_C(\ell,j)\), independent of \(r\). Lemma 4.15 realizes
the selected recipe on the run.

\emph{Lean:}
\href{https://github.com/yang-mills-group/OpenAI-YM-20261028/blob/3456e5e75cf8d4cd558955a5c7638fba2d68f464/YM/Independent/ConstructionRun.lean\#L14-L28}{\texttt{constructionRun}},
\href{https://github.com/yang-mills-group/OpenAI-YM-20261028/blob/3456e5e75cf8d4cd558955a5c7638fba2d68f464/YM/Integration/ActualRunNativeComponentRecipe.lean\#L21-L23}{\texttt{postnumeric}}.

\paragraph{4.6 The packages behind the numerical
floors}\label{the-packages-behind-the-numerical-floors}

The combined recipe floor includes
\href{https://github.com/yang-mills-group/OpenAI-YM-20261028/blob/3456e5e75cf8d4cd558955a5c7638fba2d68f464/YM/RG08/Redo/Scaling.lean\#L46}{\texttt{rail}},
and its post-length bill uses
\href{https://github.com/yang-mills-group/OpenAI-YM-20261028/blob/3456e5e75cf8d4cd558955a5c7638fba2d68f464/YM/RG08/Redo/Scaling.lean\#L65}{\texttt{receipt}}/\href{https://github.com/yang-mills-group/OpenAI-YM-20261028/blob/3456e5e75cf8d4cd558955a5c7638fba2d68f464/YM/RG08/Redo/Scaling.lean\#L67}{\texttt{after}}.
The second numerical component uses
\href{https://github.com/yang-mills-group/OpenAI-YM-20261028/blob/3456e5e75cf8d4cd558955a5c7638fba2d68f464/YM/RG08/Yamen/Pledge.lean\#L36}{\texttt{Yamen.need}}
and
\href{https://github.com/yang-mills-group/OpenAI-YM-20261028/blob/3456e5e75cf8d4cd558955a5c7638fba2d68f464/YM/RG08/Yamen/Penny.lean\#L48}{\texttt{Yamen.stamp}}.

The supporting estimates below supply finite widths, numerical floors
and post-length bounds. The widths also enter the reciprocal-width
requests of Lemma 4.17. These are summaries under the hypotheses of the
linked results; the complete auxiliary statements are not reproduced
here.

\begin{itemize}
\tightlist
\item
  \textbf{One-step transit and graded differences.}
  \href{https://github.com/yang-mills-group/OpenAI-YM-20261028/blob/3456e5e75cf8d4cd558955a5c7638fba2d68f464/YM/RG08/Sideshow/Orbit.lean\#L81-L190}{\texttt{harbor}}:
  for admissible parameters there are a constant \(A\ge1\) and a growth
  function \(J\) of controlled degree bounding a one-step transit of
  paired marked kernel families. Supporting results:
  \href{https://github.com/yang-mills-group/OpenAI-YM-20261028/blob/3456e5e75cf8d4cd558955a5c7638fba2d68f464/YM/RG08/Necks/Elect.lean\#L78-L140}{\texttt{DocketBill}},
  \href{https://github.com/yang-mills-group/OpenAI-YM-20261028/blob/3456e5e75cf8d4cd558955a5c7638fba2d68f464/YM/RG08/Necks/Prior.lean\#L61-L95}{\texttt{slushy\_uniform}},
  \href{https://github.com/yang-mills-group/OpenAI-YM-20261028/blob/3456e5e75cf8d4cd558955a5c7638fba2d68f464/YM/RG08/Necks/Lob.lean\#L32-L124}{\texttt{conquer}},
  \href{https://github.com/yang-mills-group/OpenAI-YM-20261028/blob/3456e5e75cf8d4cd558955a5c7638fba2d68f464/YM/RG08/Necks/Tides.lean\#L45-L88}{\texttt{drive}}
  (graded differences of two kernel families with scale-independent
  constants).
\item
  \textbf{Curvature-coefficient lifts with linear cost.} Closed
  coefficients satisfying the required grading, symmetry and vanishing
  conditions admit curvature lifts, with reduced width and a bound
  linear in the input size. The bilinear case needs an additional
  zero-momentum condition:
  \href{https://github.com/yang-mills-group/OpenAI-YM-20261028/blob/3456e5e75cf8d4cd558955a5c7638fba2d68f464/YM/RG08/Rudder/FinalLegs.lean\#L49-L74}{\texttt{line\_legs}},
  \href{https://github.com/yang-mills-group/OpenAI-YM-20261028/blob/3456e5e75cf8d4cd558955a5c7638fba2d68f464/YM/RG08/Rudder/Loop/Empty.lean\#L81-L107}{\texttt{doubleLift}},
  \href{https://github.com/yang-mills-group/OpenAI-YM-20261028/blob/3456e5e75cf8d4cd558955a5c7638fba2d68f464/YM/RG08/Rudder/Sym/Oaks.lean\#L65-L81}{\texttt{arc}},
  \href{https://github.com/yang-mills-group/OpenAI-YM-20261028/blob/3456e5e75cf8d4cd558955a5c7638fba2d68f464/YM/RG08/Spar/Overhead.lean\#L39-L49}{\texttt{warehouse}},
  \href{https://github.com/yang-mills-group/OpenAI-YM-20261028/blob/3456e5e75cf8d4cd558955a5c7638fba2d68f464/YM/AnalyticNext/QUARTIC_MATCH/Cells/LiftOne.lean\#L198-L206}{\texttt{arm}},
  \href{https://github.com/yang-mills-group/OpenAI-YM-20261028/blob/3456e5e75cf8d4cd558955a5c7638fba2d68f464/YM/AnalyticNext/QUARTIC_MATCH/Cells/Resolve.lean\#L158-L178}{\texttt{multi}}.
\item
  \textbf{Formal Gaussian change of variables.} Near-identity
  substitutions preserve the formal Gaussian functional after Jacobian
  correction. The supporting results include marked-variable vanishing
  and compatibility of connected logarithms:
  \href{https://github.com/yang-mills-group/OpenAI-YM-20261028/blob/3456e5e75cf8d4cd558955a5c7638fba2d68f464/YM/RG08/Room/Time.lean\#L54-L81}{\texttt{bottle}},
  \href{https://github.com/yang-mills-group/OpenAI-YM-20261028/blob/3456e5e75cf8d4cd558955a5c7638fba2d68f464/YM/RG08/Room/Glass.lean\#L98-L111}{\texttt{still}},
  \href{https://github.com/yang-mills-group/OpenAI-YM-20261028/blob/3456e5e75cf8d4cd558955a5c7638fba2d68f464/YM/RG08/Room/Alight.lean\#L86-L123}{\texttt{absorbed}},
  \href{https://github.com/yang-mills-group/OpenAI-YM-20261028/blob/3456e5e75cf8d4cd558955a5c7638fba2d68f464/YM/RG08/Room/Refresh.lean\#L87-L106}{\texttt{impress}}.
\item
  \textbf{Formal Ward identities and gauge variations.} Change of
  variables with an insertion, a finite-torus Ward identity, and bounds
  on the response to adding gradients:
  \href{https://github.com/yang-mills-group/OpenAI-YM-20261028/blob/3456e5e75cf8d4cd558955a5c7638fba2d68f464/YM/RG08/Wisp/Chord.lean\#L48-L76}{\texttt{stirred}},
  \href{https://github.com/yang-mills-group/OpenAI-YM-20261028/blob/3456e5e75cf8d4cd558955a5c7638fba2d68f464/YM/RG08/Weald/Pitch.lean\#L80-L98}{\texttt{welcomed}},
  \href{https://github.com/yang-mills-group/OpenAI-YM-20261028/blob/3456e5e75cf8d4cd558955a5c7638fba2d68f464/YM/RG08/Flint/Hunger.lean\#L77-L91}{\texttt{spark}}.
\item
  \textbf{Representatives and evaluation.} Gauge-equivalent
  parametrizations give equal block-integral germs and rescaled
  coefficients under the stated regularity conditions; the
  block-integral comparison also assumes equality of the quadratic
  germs. A related formula evaluates completed coefficients by averaging
  over signed lattice symmetries:
  \href{https://github.com/yang-mills-group/OpenAI-YM-20261028/blob/3456e5e75cf8d4cd558955a5c7638fba2d68f464/YM/RG08/Atlas/Furnace.lean\#L32-L51}{\texttt{blossom}},
  \href{https://github.com/yang-mills-group/OpenAI-YM-20261028/blob/3456e5e75cf8d4cd558955a5c7638fba2d68f464/YM/RG08/Aggregate/Voice.lean\#L84-L94}{\texttt{switch\_rep}},
  \href{https://github.com/yang-mills-group/OpenAI-YM-20261028/blob/3456e5e75cf8d4cd558955a5c7638fba2d68f464/YM/RG08/Carrel/Works.lean\#L43-L52}{\texttt{hatch}}.
\item
  \textbf{Marked-kernel receipts.} Bounds on marked kernels including
  the exact leading term, with constants announced before the block
  factor:
  \href{https://github.com/yang-mills-group/OpenAI-YM-20261028/blob/3456e5e75cf8d4cd558955a5c7638fba2d68f464/YM/RG08/Trill/Swift.lean\#L40-L96}{\texttt{podium}},
  \href{https://github.com/yang-mills-group/OpenAI-YM-20261028/blob/3456e5e75cf8d4cd558955a5c7638fba2d68f464/YM/RG08/Halo/Kumi.lean\#L63-L106}{\texttt{baton}}.
\item
  \textbf{Line inverse in cell coordinates.} An algebraic inverse on the
  constrained subspace, followed by strip bounds and the finite-torus
  inverse identity:
  \href{https://github.com/yang-mills-group/OpenAI-YM-20261028/blob/3456e5e75cf8d4cd558955a5c7638fba2d68f464/YM/RG08/Line/Solid.lean\#L21-L38}{\texttt{jointProjection}},
  \href{https://github.com/yang-mills-group/OpenAI-YM-20261028/blob/3456e5e75cf8d4cd558955a5c7638fba2d68f464/YM/RG08/Line/Recipe.lean\#L60-L83}{\texttt{preheat}},
  \href{https://github.com/yang-mills-group/OpenAI-YM-20261028/blob/3456e5e75cf8d4cd558955a5c7638fba2d68f464/YM/RG08/Line/Complex.lean\#L108-L138}{\texttt{transcribe}}.
\end{itemize}

\subsubsection{5. Where the difficulty
lies}\label{where-the-difficulty-lies-16}

\begin{itemize}
\tightlist
\item
  \textbf{Degenerate quadratic forms.} Gauge zero modes prevent
  inversion of \(d^\top Md\). The mixed Schur identity propagates
  curvature representatives without this inversion.
\item
  \textbf{Neutrality.} Zero row sums are required separately in every
  target orientation. Kinetic normalization gives the initial
  cancellation; the constant-row-sum identity for the free transfer
  kernel preserves it. Scaled translation covariance alone does not
  suffice.
\item
  \textbf{The exponent.} The width floor reduces the \(L^{-1}\) gain to
  \(L^{-1/2}\), and the recursion estimate uses half of the remaining
  exponent to absorb age growth. The resulting exponent \(1/4\) is not
  claimed to be optimal.
\item
  \textbf{Order of choices.} The flat floor is fixed from the price map
  before the block factors. It involves the constants at indices
  \(j,j-1\) and widths at \(j,j+1\). The post-length bill then enforces
  the delay condition through the coupling thresholds.
\item
  \textbf{Scope.} The estimate controls the Gaussian link kernel
  \(d^\top M^\infty d\). It gives no convergence rate for the
  undetermined part of the plaquette metric, positive-order coefficients
  or source kernels. The formal recipe is not an identification with
  physical marked observables.
\end{itemize}

\subsubsection{6. Notation}\label{notation-16}

{\def\LTcaptype{none} % do not increment counter
\begin{longtable}[]{@{}
  >{\raggedright\arraybackslash}p{(\linewidth - 4\tabcolsep) * \real{0.3333}}
  >{\raggedright\arraybackslash}p{(\linewidth - 4\tabcolsep) * \real{0.3333}}
  >{\raggedright\arraybackslash}p{(\linewidth - 4\tabcolsep) * \real{0.3333}}@{}}
\toprule\noalign{}
\begin{minipage}[b]{\linewidth}\raggedright
Written proof
\end{minipage} & \begin{minipage}[b]{\linewidth}\raggedright
Lean
\end{minipage} & \begin{minipage}[b]{\linewidth}\raggedright
Meaning
\end{minipage} \\
\midrule\noalign{}
\endhead
\bottomrule\noalign{}
\endlastfoot
run \(R\), campaign \(T\), pillars \(W\) &
\href{https://github.com/yang-mills-group/OpenAI-YM-20261028/blob/3456e5e75cf8d4cd558955a5c7638fba2d68f464/YM/FullVolume/Cross.lean\#L30-L48}{\texttt{Run}},
\texttt{Campaign},
\href{https://github.com/yang-mills-group/OpenAI-YM-20261028/blob/3456e5e75cf8d4cd558955a5c7638fba2d68f464/YM/FullVolume/Gables/Vendee.lean\#L22-L27}{\texttt{Pillars}}
& Wilson family with its scale and extraction data; campaign comparison
data; selected histories from stage \(0\) for all \((N,p)\), with local
step data and accounting identities \\
funded &
\href{https://github.com/yang-mills-group/OpenAI-YM-20261028/blob/3456e5e75cf8d4cd558955a5c7638fba2d68f464/YM/R2I/Campaign/Accepted.lean\#L60-L73}{\texttt{R2I.Funded}}
& 13 clauses: \texttt{localBudget}, \texttt{attentive},
\texttt{logCopy}, \texttt{locrian}, \texttt{wild}, \texttt{silk},
\texttt{poly}, \texttt{fixed}, \texttt{tic}, \texttt{halo},
\texttt{paced}, \texttt{stoves}, \texttt{hasSky} \\
\(\Delta_{10}\) &
\href{https://github.com/yang-mills-group/OpenAI-YM-20261028/blob/3456e5e75cf8d4cd558955a5c7638fba2d68f464/YM/RG12/Burst/Wardrobe.lean\#L63}{\texttt{garment10}},
\href{https://github.com/yang-mills-group/OpenAI-YM-20261028/blob/3456e5e75cf8d4cd558955a5c7638fba2d68f464/YM/RG12/Burst/Wardrobe.lean\#L18-L20}{\texttt{Decal}}
& local one-step estimate package with its eventual-validity lemma \\
sheet \(s\) &
\href{https://github.com/yang-mills-group/OpenAI-YM-20261028/blob/3456e5e75cf8d4cd558955a5c7638fba2d68f464/YM/Price/Intent.lean\#L26-L35}{\texttt{Sheet}}
& kernel choice, stage profile, floors; fixed before \(L\) \\
price map \(P=(\alpha_j,\Omega_j,\eta_j)\) &
\href{https://github.com/yang-mills-group/OpenAI-YM-20261028/blob/3456e5e75cf8d4cd558955a5c7638fba2d68f464/YM/PREC/Stacks.lean\#L123-L128}{\texttt{MapPrices}}
& strip width, upper bound and coercivity bound for each smoothing
index \\
horizon \(H\), sufficient for \(P\) &
\href{https://github.com/yang-mills-group/OpenAI-YM-20261028/blob/3456e5e75cf8d4cd558955a5c7638fba2d68f464/YM/PREC/Stacks.lean\#L119-L121}{\texttt{Horizon}},
\href{https://github.com/yang-mills-group/OpenAI-YM-20261028/blob/3456e5e75cf8d4cd558955a5c7638fba2d68f464/YM/PREC/Stacks.lean\#L132-L139}{\texttt{Suffices}}
& metrics \(M_{N,i}\) and transfer data, with strip bounds and Schur
relations \\
\(\mathcal K^{\mathbb T}_{N,i}=d^\top M^{\mathbb T}_{N,i}d\) &
\href{https://github.com/yang-mills-group/OpenAI-YM-20261028/blob/3456e5e75cf8d4cd558955a5c7638fba2d68f464/YM/PREC/TreeMatch.lean\#L68-L69}{\texttt{Precision.K}},
\texttt{quadratic} & Gaussian quadratic form on a torus \\
\(\mathcal O_{N,i}\) &
\href{https://github.com/yang-mills-group/OpenAI-YM-20261028/blob/3456e5e75cf8d4cd558955a5c7638fba2d68f464/YM/RG08/Flat/Curl.lean\#L18}{\texttt{output}},
\href{https://github.com/yang-mills-group/OpenAI-YM-20261028/blob/3456e5e75cf8d4cd558955a5c7638fba2d68f464/YM/RG08/Flat/Curl.lean\#L15-L16}{\texttt{writ}}
& position-space kernel \(d^\top M^\infty d\) on links \\
background-field relations &
\href{https://github.com/yang-mills-group/OpenAI-YM-20261028/blob/3456e5e75cf8d4cd558955a5c7638fba2d68f464/YM/PREC/TreeMatch.lean\#L76-L79}{\texttt{Inputs}}
& \(\mathcal K W=Q^\top\mathcal K'\), \(QW=1-b^{-1}\mathcal K'\) \\
curvature transport \(T_{m,i}\) &
\href{https://github.com/yang-mills-group/OpenAI-YM-20261028/blob/3456e5e75cf8d4cd558955a5c7638fba2d68f464/YM/Pert/Bulk/Keys.lean\#L15}{\texttt{fine\_share}},
\href{https://github.com/yang-mills-group/OpenAI-YM-20261028/blob/3456e5e75cf8d4cd558955a5c7638fba2d68f464/YM/Pert/Bulk/Coin.lean\#L40}{\texttt{response}}
& \((M^\infty_{m,i+1})^{-1}\mathcal T_{m,i}M^\infty_{m,i}\) \\
\(h_i\), \(q_i(u)\) &
\href{https://github.com/yang-mills-group/OpenAI-YM-20261028/blob/3456e5e75cf8d4cd558955a5c7638fba2d68f464/YM/RG08/Flat/World.lean\#L22}{\texttt{header}},
\href{https://github.com/yang-mills-group/OpenAI-YM-20261028/blob/3456e5e75cf8d4cd558955a5c7638fba2d68f464/YM/RG08/Flat/World.lean\#L23}{\texttt{quote}},
\href{https://github.com/yang-mills-group/OpenAI-YM-20261028/blob/3456e5e75cf8d4cd558955a5c7638fba2d68f464/YM/RG08/Packets/Lines.lean\#L84-L85}{\texttt{neutralCoin}}
& \(\exp(4)\lvert\mathcal F\rvert^2\Omega_{J_{i+1}}^2\Omega_{J_i}\);
\(2\lvert\mathcal F\rvert^2h_i^2/(uL_i)\) \\
flat representation &
\href{https://github.com/yang-mills-group/OpenAI-YM-20261028/blob/3456e5e75cf8d4cd558955a5c7638fba2d68f464/YM/RG08/Flat/World.lean\#L36-L40}{\texttt{FlatRep}},
\href{https://github.com/yang-mills-group/OpenAI-YM-20261028/blob/3456e5e75cf8d4cd558955a5c7638fba2d68f464/YM/RG08/Packets/FlatStep.lean\#L43-L44}{\texttt{Sigil}}
& Definition 4.1 \\
weighted norms &
\href{https://github.com/yang-mills-group/OpenAI-YM-20261028/blob/3456e5e75cf8d4cd558955a5c7638fba2d68f464/YM/Pert/Bulk/Stacks.lean\#L14}{\texttt{Folded}},
\href{https://github.com/yang-mills-group/OpenAI-YM-20261028/blob/3456e5e75cf8d4cd558955a5c7638fba2d68f464/YM/Pert/Bulk/Atlas.lean\#L28-L29}{\texttt{RowCost}},
\href{https://github.com/yang-mills-group/OpenAI-YM-20261028/blob/3456e5e75cf8d4cd558955a5c7638fba2d68f464/YM/Pert/Bulk/Atlas.lean\#L37-L42}{\texttt{Atlas}}
& plaquette norm \(\lVert K\rVert^{\rm pl}_\sigma\); link row norm
\(\lVert F\rVert_\sigma\); row, column and first-difference bounds of
order \(m\) \\
neutral, translation invariant &
\href{https://github.com/yang-mills-group/OpenAI-YM-20261028/blob/3456e5e75cf8d4cd558955a5c7638fba2d68f464/YM/Pert/Bulk/NullPair.lean\#L9}{\texttt{neutral}},
\href{https://github.com/yang-mills-group/OpenAI-YM-20261028/blob/3456e5e75cf8d4cd558955a5c7638fba2d68f464/YM/RG08/Packets/Torus.lean\#L18}{\texttt{town}}
& zero symbol at zero momentum; \(K(x+u,y+u)=K(x,y)\) \\
\(w_j\), \(\Lambda_{k,i}\) &
\href{https://github.com/yang-mills-group/OpenAI-YM-20261028/blob/3456e5e75cf8d4cd558955a5c7638fba2d68f464/YM/RG08/Flat/Length.lean\#L14}{\texttt{curve}},
\href{https://github.com/yang-mills-group/OpenAI-YM-20261028/blob/3456e5e75cf8d4cd558955a5c7638fba2d68f464/YM/RG08/Flat/Mesh.lean\#L30}{\texttt{clock}},
\href{https://github.com/yang-mills-group/OpenAI-YM-20261028/blob/3456e5e75cf8d4cd558955a5c7638fba2d68f464/YM/RG08/Flat/Length.lean\#L17}{\texttt{Lambda\textquotesingle{}}}
& \(\min(1,\alpha_j,\alpha_{j+1})\);
\(5+\log\hat\beta_k+\lvert s_i-s_k\rvert\) \\
delay condition &
\href{https://github.com/yang-mills-group/OpenAI-YM-20261028/blob/3456e5e75cf8d4cd558955a5c7638fba2d68f464/YM/RG08/Flat/Length.lean\#L19}{\texttt{Delay}}
& \(\Omega_{J_i}\le\log\hat\beta_i\) for \(i\ge e\) \\
flat rate (III.14.3) &
\href{https://github.com/yang-mills-group/OpenAI-YM-20261028/blob/3456e5e75cf8d4cd558955a5c7638fba2d68f464/YM/RG08/Flat/Entrance.lean\#L13-L21}{\texttt{Quorum}}
& clauses (a)--(c) \\
family and numerical assumptions &
\href{https://github.com/yang-mills-group/OpenAI-YM-20261028/blob/3456e5e75cf8d4cd558955a5c7638fba2d68f464/YM/RG08/Flat/World.lean\#L82-L91}{\texttt{World}},
\href{https://github.com/yang-mills-group/OpenAI-YM-20261028/blob/3456e5e75cf8d4cd558955a5c7638fba2d68f464/YM/RG08/Flat/World.lean\#L101-L110}{\texttt{Funding}},
\href{https://github.com/yang-mills-group/OpenAI-YM-20261028/blob/3456e5e75cf8d4cd558955a5c7638fba2d68f464/YM/RG08/Flat/Mesh.lean\#L32-L41}{\texttt{Spacing}}
& hypotheses of III.14.4; their form on a common level system \\
stairs, drop, weight &
\href{https://github.com/yang-mills-group/OpenAI-YM-20261028/blob/3456e5e75cf8d4cd558955a5c7638fba2d68f464/YM/RG08/Rate/Stair.lean\#L24-L30}{\texttt{Stair}},
\href{https://github.com/yang-mills-group/OpenAI-YM-20261028/blob/3456e5e75cf8d4cd558955a5c7638fba2d68f464/YM/RG08/Rate/Stair.lean\#L66}{\texttt{drop}},
\href{https://github.com/yang-mills-group/OpenAI-YM-20261028/blob/3456e5e75cf8d4cd558955a5c7638fba2d68f464/YM/RG08/Rate/Stair.lean\#L108}{\texttt{Fade}}
& index family; \(L_i^{-\theta}\); age polynomial times
\((a_t/a_i)^c\) \\
shadow &
\href{https://github.com/yang-mills-group/OpenAI-YM-20261028/blob/3456e5e75cf8d4cd558955a5c7638fba2d68f464/YM/RG08/Flat/Shadow.lean\#L30-L50}{\texttt{shadow}}
& same level system with clip profile replaced \\
recipe \(C\), floor \(E_C\), bill \(B_C\) &
\href{https://github.com/yang-mills-group/OpenAI-YM-20261028/blob/3456e5e75cf8d4cd558955a5c7638fba2d68f464/YM/Integration/NativeComponentRecipe.lean\#L29-L51}{\texttt{Choice}},
\href{https://github.com/yang-mills-group/OpenAI-YM-20261028/blob/3456e5e75cf8d4cd558955a5c7638fba2d68f464/YM/Integration/NativeComponentRecipe.lean\#L61-L65}{\texttt{floor}},
\href{https://github.com/yang-mills-group/OpenAI-YM-20261028/blob/3456e5e75cf8d4cd558955a5c7638fba2d68f464/YM/Integration/NativeComponentRecipe.lean\#L75-L79}{\texttt{post}}
& three components with pre-length floors; combined floor \(E_C\) and
post-length bill \(B_C\) \\
point-group invariance (``sworn'') &
\href{https://github.com/yang-mills-group/OpenAI-YM-20261028/blob/3456e5e75cf8d4cd558955a5c7638fba2d68f464/YM/Bid/Station.lean\#L54-L55}{\texttt{Bid.Sworn}}
& colour-summed bilinear form invariant under signed axis
permutations \\
background-field identities (``licensed'') &
\href{https://github.com/yang-mills-group/OpenAI-YM-20261028/blob/3456e5e75cf8d4cd558955a5c7638fba2d68f464/YM/RG08/Aura/Publisher.lean\#L24-L26}{\texttt{Aura.licensed}}
& free background map satisfies the background-field identities for the
actual forms \\
``prices'', ``bill'', ``floor'', ``numeric'', ``post'' & --- & numerical
bounds on block factors and the reference coupling \\
\end{longtable}
}

\subsubsection{7. What remains
unexplained}\label{what-remains-unexplained-16}

\begin{enumerate}
\def\labelenumi{\arabic{enumi}.}
\tightlist
\item
  \textbf{Selection in III.2.} Lemma 4.16 is used as an input; the
  underlying construction of the scale and extraction data is not
  audited here. Its two auxiliary inputs have canonical choices.
  Removing the recipe demand would give a separate existence statement,
  not an identification of the resulting run with the fixed run of
  Proposition III.14.1.
\item
  \textbf{The funded condition.} Its thirteen clauses are retained in
  full, rather than replaced by their summary in Proposition III.14.1.
  The route through
  \href{https://github.com/yang-mills-group/OpenAI-YM-20261028/blob/3456e5e75cf8d4cd558955a5c7638fba2d68f464/YM/R2I/Campaign/Accepted.lean\#L150-L167}{\texttt{fromFunded}}
  uses the clauses for extraction, comparison, OPE and asymptotic
  control; their detailed mathematical meanings and proofs belong to
  III.15 and later units. In particular, that route does not project the
  local-estimate clause involving \(\Delta_{10}\).
\item
  \textbf{Auxiliary estimates in §4.6.} Only their statements and
  numerical uses are summarized. Their full admissibility conditions and
  proofs, including the curvature-lift, marked-kernel and formal Ward
  estimates, remain outside this verification. The summaries are not
  substitutes for those hypotheses.
\item
  \textbf{Inputs from III.3 and III.8.} The free-kernel existence
  theorem and the underlying lattice kernel estimates are used at their
  stated interfaces. The proofs of those upstream results are not
  audited here.
\end{enumerate}

\subsection{III.15. Cost accounting for observables and the remainder
comparison}\label{iii.15.-cost-accounting-for-observables-and-the-remainder-comparison}

\textbf{Part:} III, Ultraviolet construction (renormalization group).
\textbf{Lean:} chapter 32:
\href{https://github.com/yang-mills-group/OpenAI-YM-20261028/tree/3456e5e75cf8d4cd558955a5c7638fba2d68f464/YM/RemainderComparison}{\texttt{YM/RemainderComparison/}}
(packages Chron, Cadenza, Halo, Muse, Solo, Lyric, Herbal, Cascade,
Sleet, Access, Portal, Aux and about twenty smaller ones),
\href{https://github.com/yang-mills-group/OpenAI-YM-20261028/tree/3456e5e75cf8d4cd558955a5c7638fba2d68f464/YM/Progress/DIRECTED_RECIPE}{\texttt{YM/Progress/DIRECTED\_RECIPE/}},
\href{https://github.com/yang-mills-group/OpenAI-YM-20261028/tree/3456e5e75cf8d4cd558955a5c7638fba2d68f464/YM/Next/SCHEDULE}{\texttt{YM/Next/SCHEDULE/}},
\href{https://github.com/yang-mills-group/OpenAI-YM-20261028/tree/3456e5e75cf8d4cd558955a5c7638fba2d68f464/YM/RG05Readout}{\texttt{YM/RG05Readout/}},
\href{https://github.com/yang-mills-group/OpenAI-YM-20261028/tree/3456e5e75cf8d4cd558955a5c7638fba2d68f464/YM/Integration/ChronPreconstruction}{\texttt{YM/Integration/ChronPreconstruction/}},
\href{https://github.com/yang-mills-group/OpenAI-YM-20261028/tree/3456e5e75cf8d4cd558955a5c7638fba2d68f464/YM/Integration/RG05Preconstruction}{\texttt{YM/Integration/RG05Preconstruction/}},
and small parts of \texttt{YM/R2I/}, \texttt{YM/RG12/},
\texttt{YM/OPE/Therm/}. \textbf{Size:} 7,382 lines in 1,053 declarations
used by the main theorem. About 4,200 of these lines are the launch
theorem and its prerequisites.

\textbf{Overview.} The unit has three independent parts, plus some
record types.

\begin{itemize}
\tightlist
\item
  \textbf{(A) Pair-step pricing} (\texttt{DIRECTED\_RECIPE}). Two
  numerical facts used by the RG step of III.9 and III.10. A hard family
  of polymer activities is merged with the Mayer expansion of the
  exponential of soft local terms. The merged family keeps a weighted
  polymer bound and carries an explicit Gaussian large-field factor
  (III.15.1--2).
\item
  \textbf{(B) Derivative bounds for local observables}
  (\texttt{SCHEDULE/Wages}). On the small-field region, a transported
  curvature slot has left derivatives of every order \(\le K\) bounded
  by a constant times \(t\). A \(k\)-leg monomial in slots has them
  bounded by a constant times \(t^k\). Here \(t\) is the small-field
  threshold (III.15.3--4).
\item
  \textbf{(C) Launch data for the remainder comparison}
  (\texttt{RemainderComparison}). For connected-activity targets with
  rates below \(-\tfrac38\), a sufficiently large block side permits the
  numerical conditions of a comparison step. After the auxiliary data
  are fixed, one chooses the comparison functions and then an entry
  coupling above which those conditions hold (III.15.5--7).
\item
  \textbf{(D) Records and numerical definitions.} These include the
  construction-run bounds, the additional-expense and readout records,
  and the RG05 readout prices. Their scope is described in §4.6.
\end{itemize}

\textbf{Status of the parts for the main theorem.}

\begin{itemize}
\tightlist
\item
  \begin{enumerate}
  \def\labelenumi{(\Alph{enumi})}
  \tightlist
  \item
    and (B) are consumed by III.2, III.9 and III.10.
  \end{enumerate}
\item
  (C) supplies additional comparison data in the main-theorem dependency
  chain:

  \begin{itemize}
  \tightlist
  \item
    III.2's assembly lemma
    \href{https://github.com/yang-mills-group/OpenAI-YM-20261028/blob/3456e5e75cf8d4cd558955a5c7638fba2d68f464/YM/Integration/RG05Preconstruction/Expenses.lean\#L28-L119}{\texttt{YM.PhysicalNext.Repeat.Legato.R2I.RG05Joint.joinExpensesShipWithPrediction}}
    calls the launch and delay-enlargement theorems to build the
    additional conjunct
    \href{https://github.com/yang-mills-group/OpenAI-YM-20261028/blob/3456e5e75cf8d4cd558955a5c7638fba2d68f464/YM/RemainderComparison/Chron/Rig.lean\#L21-L22}{\texttt{YM.PhysicalNext.Repeat.Legato.Stagehands.News}},
    which records a family of quarries and their stagewise bounds.
  \item
    That conjunct is passed unchanged up to the selected-run theorem
    \href{https://github.com/yang-mills-group/OpenAI-YM-20261028/blob/3456e5e75cf8d4cd558955a5c7638fba2d68f464/YM/Integration/NativePreselectedPostnumericPeriodRun.lean\#L99-L114}{\texttt{YM.PhysicalNext.Repeat.Legato.R2I.RG05Joint.selectRunWithTreasuryPeriod\_preselected\_quote\_orders\_postnumeric}}
    of III.14.
  \item
    In the main-theorem dependency chain,
    \href{https://github.com/yang-mills-group/OpenAI-YM-20261028/blob/3456e5e75cf8d4cd558955a5c7638fba2d68f464/YM/Independent/ConstructionRun.lean\#L14-L28}{\texttt{YM.Independent.constructionRun}}
    discards it. The retained construction-run conditions form a
    separate record.
  \item
    Part C imposes additional lower bounds on block sides and entry
    couplings. Its comparison conclusions are absent from the final
    construction-run statement.
  \end{itemize}
\end{itemize}

\textbf{Conventions.}

\begin{itemize}
\tightlist
\item
  \(x\) denotes the coarse renormalized inverse coupling \(\hat\beta\).
\item
  \(\beta_F(x,L)=x+b_G\log L\) is the inverse coupling at the next finer
  scale
  (\href{https://github.com/yang-mills-group/OpenAI-YM-20261028/blob/3456e5e75cf8d4cd558955a5c7638fba2d68f464/YM/Pressure/Seeds.lean\#L23}{\texttt{YM.Pressure.betaF}});
  \(b_G\) is the slope constant of \(G\), positive when \(G\) is simple.
\end{itemize}

For \(\beta>0\) and \(x>1\), use the scale functions

\[g(\beta)=\beta^{-1/2},\qquad p(x)=e^{(\log\log x)^2},\qquad t(x)=g(x)\,p(x),\qquad \mathcal B(x)=\lceil(\log x)^2\rceil .\]

\begin{itemize}
\tightlist
\item
  \(t\) is the small-field threshold
  (\href{https://github.com/yang-mills-group/OpenAI-YM-20261028/blob/3456e5e75cf8d4cd558955a5c7638fba2d68f464/YM/Pressure/Seeds.lean\#L25}{\texttt{YM.Pressure.temp}},
  equal to
  \href{https://github.com/yang-mills-group/OpenAI-YM-20261028/blob/3456e5e75cf8d4cd558955a5c7638fba2d68f464/YM/Next/SCHEDULE/Dues.lean\#L59}{\texttt{YM.Next.Scheduling.Instrument.t}}
  for the plain instrument).
\item
  \(\mathcal B\) is the cluster radius
  (\href{https://github.com/yang-mills-group/OpenAI-YM-20261028/blob/3456e5e75cf8d4cd558955a5c7638fba2d68f464/YM/Pressure/Seeds.lean\#L22}{\texttt{YM.Pressure.radiusB}},
  the same formula as
  \href{https://github.com/yang-mills-group/OpenAI-YM-20261028/blob/3456e5e75cf8d4cd558955a5c7638fba2d68f464/YM/Next/SCHEDULE/Clocks.lean\#L38}{\texttt{YM.Next.Scheduling.Clock.cover}}).
\item
  The band of \(y\) is \([y/4,4y]\)
  (\href{https://github.com/yang-mills-group/OpenAI-YM-20261028/blob/3456e5e75cf8d4cd558955a5c7638fba2d68f464/YM/Next/SCHEDULE/Dues.lean\#L99}{\texttt{YM.Next.Scheduling.inBand}}).
  On it, \(g(y)/2\le g(\beta)\le 2g(y)\)
  (\href{https://github.com/yang-mills-group/OpenAI-YM-20261028/blob/3456e5e75cf8d4cd558955a5c7638fba2d68f464/YM/Next/SCHEDULE/Dues.lean\#L105-L111}{\texttt{YM.Next.Scheduling.band\_gap}}).
\end{itemize}

Growth classes, at fixed \(L\), as \(x\to\infty\):

\begin{itemize}
\tightlist
\item
  \textbf{Rate} \(\le a\)
  (\href{https://github.com/yang-mills-group/OpenAI-YM-20261028/blob/3456e5e75cf8d4cd558955a5c7638fba2d68f464/YM/AnalyticNext/CONNECTED_ACTIVITY_JETS_E6/Numerics.lean\#L22-L23}{\texttt{YM.ConnectedActivityE6.Rate}}):
  for every \(v>a\), eventually \(|w(x)|\le\beta_F(x,L)^v\).
\item
  \textbf{Reserve}
  (\href{https://github.com/yang-mills-group/OpenAI-YM-20261028/blob/3456e5e75cf8d4cd558955a5c7638fba2d68f464/YM/RemainderComparison/Access/Geometry.lean\#L25-L26}{\texttt{YM.RemainderComparison.Reserve}}):
  \(|w(x)|\le M(x)^k\) eventually for some \(k\in\mathbb N\), where
  \(M(x)=p(\beta_F(x,L))\)
  (\href{https://github.com/yang-mills-group/OpenAI-YM-20261028/blob/3456e5e75cf8d4cd558955a5c7638fba2d68f464/YM/Push/ABSOLUTE_SUMS/Delay.lean\#L24}{\texttt{YM.Push.AbsoluteSums.Delay.fineMonitor}}).
  Every reserve has rate at most \(0\).
\item
  \textbf{Gear} \(a\)
  (\href{https://github.com/yang-mills-group/OpenAI-YM-20261028/blob/3456e5e75cf8d4cd558955a5c7638fba2d68f464/YM/RemainderComparison/Gear/Class.lean\#L25}{\texttt{YM.RemainderComparison.Gear}}):
  \(\beta_F^{-a}w\) is a reserve, equivalently
  \(|w(x)|\le\beta_F(x,L)^a M(x)^k\) eventually for some
  \(k\in\mathbb N\). This implies rate at most \(a\).
\end{itemize}

Weighted hit bound
(\href{https://github.com/yang-mills-group/OpenAI-YM-20261028/blob/3456e5e75cf8d4cd558955a5c7638fba2d68f464/YM/Ac/Tree.lean\#L12-L13}{\texttt{YM.Ac.Hit}}):
for supports \(X_i\), weights \(f_i\ge0\), scale \(M\) and exponent
\(a\), write \(\mathrm{Hit}_M(f;a)\le\gamma\) if

\[\sum_{i:\,z\in X_i} f_i\,e^{a N_M(X_i)}\le\gamma\quad\text{for every site }z,\]
where \(N_M(X)\) is the minimal number of \(M\)-balls with connected
union covering \(X\).

The quote is \(\mathrm{quote}(\kappa,\gamma)=e^{-\kappa}(e^{\gamma}-1)\)
(\href{https://github.com/yang-mills-group/OpenAI-YM-20261028/blob/3456e5e75cf8d4cd558955a5c7638fba2d68f464/YM/Bid/Kitchen.lean\#L18}{\texttt{YM.Bid.quote}}).

In the code, a price or bill is a numerical bound; ``financed'' means
that the bound satisfies the specified growth conditions.

\subsubsection{1. What this unit proves}\label{what-this-unit-proves-17}

\paragraph{Part A: pair-step pricing}\label{part-a-pair-step-pricing}

\textbf{Proposition III.15.1 (pair-step rates).} Let \(\rho>0\),
\(C\ge0\), \(r>\rho\), \(d\in\mathbb N\), \(e\in\mathbb R\) and \(q>0\).
For \(y>1\) and \(\beta>0\) put

\[s(y,\beta)=\frac{t(y)}{g(\beta)}=p(y)\sqrt{\beta/y},\quad \mathrm{mild}(y)=C(1+\mathcal B(y))^{d}p(y)^{e}g(y)^{q},\]
\[\mathrm{strong}(y,\beta)=C e^{-r s^2},\quad \kappa(y,\beta)=\rho s^2 .\]
For all sufficiently large \(y>1\), the following holds for every
\(\beta\in[y/4,4y]\) and every integer \(0\le B\le\mathcal B(y)\). With
\(a=\mathrm{strong}(y,\beta)\), \(b=\mathrm{mild}(y)\) and
\(\gamma=e^{\kappa}a+2b\):

\[\kappa\ge0,\quad a\ge0,\quad 0\le b\le1,\quad (2B+1)^4\gamma\le1,\quad \mathrm{quote}(\kappa,\gamma)\le e^{-\kappa}.\]
These are the rate conditions with target \(w=e^{-\kappa}\). The
threshold in \(y\) depends only on \((\rho,C,r,d,e,q)\). \emph{Status:}
theorem. \emph{Lean:}
\href{https://github.com/yang-mills-group/OpenAI-YM-20261028/blob/3456e5e75cf8d4cd558955a5c7638fba2d68f464/YM/Progress/DIRECTED_RECIPE/Economy.lean\#L76-L139}{\texttt{YM.Progress.DirectedRecipe.Payments.Pairing.currency}};
data
\href{https://github.com/yang-mills-group/OpenAI-YM-20261028/blob/3456e5e75cf8d4cd558955a5c7638fba2d68f464/YM/Progress/DIRECTED_RECIPE/Economy.lean\#L28-L36}{\texttt{YM.Progress.DirectedRecipe.Payments.Pairing}};
conditions
\href{https://github.com/yang-mills-group/OpenAI-YM-20261028/blob/3456e5e75cf8d4cd558955a5c7638fba2d68f464/YM/Progress/DIRECTED_RECIPE/Tender.lean\#L33-L39}{\texttt{YM.Progress.DirectedRecipe.Payments.Rates}}
and
\href{https://github.com/yang-mills-group/OpenAI-YM-20261028/blob/3456e5e75cf8d4cd558955a5c7638fba2d68f464/YM/Progress/DIRECTED_RECIPE/Tender.lean\#L32}{\texttt{YM.Progress.DirectedRecipe.Payments.inflation}}.
\emph{Used in:}

\begin{itemize}
\tightlist
\item
  III.9, in
  \href{https://github.com/yang-mills-group/OpenAI-YM-20261028/blob/3456e5e75cf8d4cd558955a5c7638fba2d68f464/YM/AnalyticBridge/VERIFIED_RESET_STEP/Currency.lean\#L75-L126}{\texttt{YM.AnalyticBridge.VerifiedReset.Time.cash}},
  with \(d=e=0\), \(q=\tfrac12\), \(y=x+b_G\log L\) and
  \(B=\mathcal B(x)\).
\item
  III.2, in
  \href{https://github.com/yang-mills-group/OpenAI-YM-20261028/blob/3456e5e75cf8d4cd558955a5c7638fba2d68f464/YM/FullVolume/Sym/Purse.lean\#L72-L85}{\texttt{YM.PhysicalNext.Repeat.Jams.reunite}}.
\end{itemize}

\textbf{Proposition III.15.2 (merging a hard and a soft family).} Fix a
parameter point \(P\) with derivative order \(K\ge6\), direction power
\(P_0\ge2\), cluster scale \(B_P\ge1\), support exponent \(\sigma>0\)
and threshold \(0<t_P\le1\). Consider a countably indexed regular family
\(E_i\), a \textbf{hard family} of forced activities \(F_i\) with
supports \(X_i\), and a \textbf{soft family} of real local exponential
corrections \(\ell_i\) with nonempty supports \(X'_i\). Assume:

\begin{enumerate}
\def\labelenumi{\arabic{enumi}.}
\tightlist
\item
  The regular terms satisfy the regular-family conditions at \(P\):
  regularity on their supports, vanishing at the flat field, and
  vanishing affine Hessian on nonwrapping supports. The forced terms and
  their eligibility tests are gauge invariant, measurable and determined
  by the links and final flags on their supports; a nonzero forced term
  satisfies its eligibility test. The regular and forced families admit
  the prescribed translation/reflection permutations, including the
  support and inventory transformations. Each \(\ell_i\) is gauge
  invariant, measurable and determined by the links and final flags on
  \(X'_i\); it vanishes on configurations small-field on \(X'_i\). The
  soft family is also covariant under translations and reflections.
\item
  The regular family has bound \(\delta\): it admits nonnegative bounds
  on all left derivatives through order \(K\) on its small-field
  supports, in the directions specified by \(P\), whose weighted sum at
  each anchor is at most \(\delta\).
\item
  For every \(i\) and every configuration \(U\), \(|F_i(U)|\le p_i\) and
  \(|\ell_i(U)|\le q_i\), with \(p_i,q_i\ge0\).
\item
  \(\mathrm{Hit}_{B_P}(p;\sigma+1)\le a\) on \((X_i)\) and
  \(\mathrm{Hit}_{B_P}(q;\sigma+1)\le b\) on \((X'_i)\).
\item
  The rate inequalities of III.15.1 hold with \(B=B_P\), with the final
  inequality replaced by \(\mathrm{quote}(\kappa,e^\kappa a+2b)\le w\).
\end{enumerate}

Then the fusion input condition with \((\delta,w)\) holds. Explicitly:

\begin{itemize}
\tightlist
\item
  the original regular and forced families retain bounds \(\delta\) and
  \(a\), respectively; the soft corrections admit a summable uniform
  majorant;
\item
  the merged family \((F_i)\sqcup(e^{\ell_i}-1)\), with weights
  \(e^{\kappa}p_i\) and \(2q_i\), satisfies
  \(|F_i|\le e^{-\kappa}\cdot e^{\kappa}p_i\) and
  \(|e^{\ell_i}-1|\le 2q_i\);
\item
  that family has
  \(\mathrm{Hit}_{B_P}(\cdot\,;\sigma+1)\le\gamma:=e^{\kappa}a+2b\),
  with \((2B_P+1)^4\gamma\le1\);
\item
  \(\mathrm{quote}(\kappa,\gamma)\le w\).
\end{itemize}

\emph{Status:} theorem. The only analytic input is \(|e^z-1|\le2|z|\)
for \(|z|\le1\). Through the fusion lemmas of III.10
(\texttt{Work.H01}--\texttt{H03}, used in
\href{https://github.com/yang-mills-group/OpenAI-YM-20261028/blob/3456e5e75cf8d4cd558955a5c7638fba2d68f464/YM/Progress/DIRECTED_RECIPE/Batch.lean\#L63-L68}{\texttt{YM.Progress.DirectedRecipe.Ingredients.gates}}),
the conclusion gives the fused polymer family with norm bound \(\le w\).
With \(w=e^{-\kappa}\) from III.15.1, this is the Gaussian large-field
factor \(\exp(-\rho\,(\beta/y)\,p(y)^2)\). \emph{Lean:}
\href{https://github.com/yang-mills-group/OpenAI-YM-20261028/blob/3456e5e75cf8d4cd558955a5c7638fba2d68f464/YM/Progress/DIRECTED_RECIPE/Tender.lean\#L124-L133}{\texttt{YM.Progress.DirectedRecipe.Payments.tender}};
\href{https://github.com/yang-mills-group/OpenAI-YM-20261028/blob/3456e5e75cf8d4cd558955a5c7638fba2d68f464/YM/Progress/DIRECTED_RECIPE/Tender.lean\#L42-L43}{\texttt{YM.Progress.DirectedRecipe.Payments.hardCost}},
\href{https://github.com/yang-mills-group/OpenAI-YM-20261028/blob/3456e5e75cf8d4cd558955a5c7638fba2d68f464/YM/Progress/DIRECTED_RECIPE/Tender.lean\#L45-L46}{\texttt{YM.Progress.DirectedRecipe.Payments.softCost}},
\href{https://github.com/yang-mills-group/OpenAI-YM-20261028/blob/3456e5e75cf8d4cd558955a5c7638fba2d68f464/YM/Progress/DIRECTED_RECIPE/Tender.lean\#L48-L49}{\texttt{YM.Progress.DirectedRecipe.Payments.rooted}},
\href{https://github.com/yang-mills-group/OpenAI-YM-20261028/blob/3456e5e75cf8d4cd558955a5c7638fba2d68f464/YM/Progress/DIRECTED_RECIPE/Tender.lean\#L52-L62}{\texttt{YM.Progress.DirectedRecipe.Payments.Fittings}},
\href{https://github.com/yang-mills-group/OpenAI-YM-20261028/blob/3456e5e75cf8d4cd558955a5c7638fba2d68f464/YM/Progress/DIRECTED_RECIPE/Batch.lean\#L32-L33}{\texttt{YM.Progress.DirectedRecipe.Ingredients}}.
\emph{Used in:} III.9
(\href{https://github.com/yang-mills-group/OpenAI-YM-20261028/blob/3456e5e75cf8d4cd558955a5c7638fba2d68f464/YM/AnalyticBridge/VERIFIED_RESET_STEP/Investor.lean\#L112-L128}{\texttt{YM.AnalyticBridge.VerifiedReset.Collection.Local.renew}})
and III.10
(\href{https://github.com/yang-mills-group/OpenAI-YM-20261028/blob/3456e5e75cf8d4cd558955a5c7638fba2d68f464/YM/PhysicalNext/RG_NEXT_COMPLETE_REPEATABLE_NONLINEAR_STEP/Marlin/Pylon.lean\#L38-L49}{\texttt{YM.PhysicalNext.Repeat.Marlin.cannon}}).

\paragraph{Part B: derivative bounds for local
observables}\label{part-b-derivative-bounds-for-local-observables}

A \emph{slot} \(s\) is a lattice walk \(\omega\) from the origin,
followed by a plaquette \(\pi\) (two axes) at the endpoint of \(\omega\)
(\href{https://github.com/yang-mills-group/OpenAI-YM-20261028/blob/3456e5e75cf8d4cd558955a5c7638fba2d68f464/YM/RG/PathActions.lean\#L38-L40}{\texttt{YM.RG.Slot}}).
Its length is \(|s|=|\omega|+4\), and its value at a site \(o\) is the
transported curvature

\[s(U,o)=\mathrm{Ad}(U_{\omega})\,\mathcal C(U_{\pi})\in\mathfrak g ,\]
where \(\mathcal C\) is the colour projection of the anti-Hermitian
part.

A \emph{\(k\)-leg atom} \(a\) has slots \(s_1,\dots,s_k\) and a real
coefficient tensor \(c_b\) indexed by colour multi-indices
\(b\in\{1,\dots,\dim\mathfrak g\}^k\)
(\href{https://github.com/yang-mills-group/OpenAI-YM-20261028/blob/3456e5e75cf8d4cd558955a5c7638fba2d68f464/YM/RG/PathActions.lean\#L84-L86}{\texttt{YM.RG.Atom}}).
Its value is \(a(U,o)=\sum_b c_b\prod_i s_i(U,o)^{b_i}\), its size is
\(|a|=\sum_b|c_b|\), and its length is \(\mathrm{len}(a)=\sum_i|s_i|\).

Fix a parameter point with derivative order \(K\ge6\), power \(P\ge2\)
and threshold \(t\in(0,1]\). Directions are colour-valued increments
\(A\) of the left derivative \(\partial_A\). A direction is
\emph{admissible at \(o\)} if

\[\|A_e\|\le t\,(1+\mathrm{dist}(e,o))^{P}\quad\text{for every link }e.\]
Here \(\mathrm{dist}(e,o)\) is the periodic maximum distance from the
basepoint of \(e\) to \(o\)
(\href{https://github.com/yang-mills-group/OpenAI-YM-20261028/blob/3456e5e75cf8d4cd558955a5c7638fba2d68f464/YM/RG/Masks.lean\#L141-L142}{\texttt{YM.RG.BudgetPoint.dirBound}}).

A configuration is \emph{small-field on \(X\)} (healthy,
\href{https://github.com/yang-mills-group/OpenAI-YM-20261028/blob/3456e5e75cf8d4cd558955a5c7638fba2d68f464/YM/RG/Masks.lean\#L41}{\texttt{YM.RG.Healthy}})
if every plaquette touching a site of \(X\) has chord distance \(\le t\)
from the identity. Let \(\kappa_G>0\) be the colour scale of \(G\) and
\(\|1\|\) the Frobenius norm of the identity matrix.

\textbf{Proposition III.15.3 (slot derivative bound).} Let
\(\ell\in\mathbb N\) and \(s\) be a slot with \(|s|\le\ell\), let \(S\)
be a finite set of points of \(\mathbb Z^4\) containing the stencil of
\(s\), and let \(o\) be a site. Then \(U\mapsto s(U,o)\) is \(K\) times
left-differentiable. Let \(U\) be small-field on \(o+S\) (wrapped on the
torus), and let \(A_1,\dots,A_m\) be admissible directions at \(o\) with
\(m\le K\). Then

\[\big\|\partial_{A_1}\cdots\partial_{A_m}s(U,o)\big\|\le \mathrm{slotCost}(P,K,\ell)\,t,\]
\[\mathrm{slotCost}(P,K,\ell)=\big(1+\kappa_G^{-1}\|1\|\big)\,\mathrm{margin}(P,\ell)^{K},\qquad \mathrm{margin}(P,\ell)=2(\kappa_G+1)\,2^P(\ell+1)^{P+1}.\]
The constant depends on \(P,K,\ell,G\) only, not on \(t\), \(o\) or the
volume. At \(m=0\) the stronger order-zero bound is \(\|s(U,o)\|\le t\).
\emph{Status:} theorem. \emph{Lean:}
\href{https://github.com/yang-mills-group/OpenAI-YM-20261028/blob/3456e5e75cf8d4cd558955a5c7638fba2d68f464/YM/Next/SCHEDULE/SlotPay.lean\#L108-L159}{\texttt{YM.Next.Scheduling.Wages.slotPays}};
\href{https://github.com/yang-mills-group/OpenAI-YM-20261028/blob/3456e5e75cf8d4cd558955a5c7638fba2d68f464/YM/Next/SCHEDULE/SlotPay.lean\#L95}{\texttt{YM.Next.Scheduling.Wages.slotCost}},
\href{https://github.com/yang-mills-group/OpenAI-YM-20261028/blob/3456e5e75cf8d4cd558955a5c7638fba2d68f464/YM/Next/SCHEDULE/SlotPay.lean\#L94}{\texttt{YM.Next.Scheduling.Wages.margin}},
\href{https://github.com/yang-mills-group/OpenAI-YM-20261028/blob/3456e5e75cf8d4cd558955a5c7638fba2d68f464/YM/Work/G07/Calculus.lean\#L124-L129}{\texttt{YM.Work.G07.LeftBound}}.
\emph{Used in:} III.9
(\href{https://github.com/yang-mills-group/OpenAI-YM-20261028/blob/3456e5e75cf8d4cd558955a5c7638fba2d68f464/YM/AnalyticBridge/VERIFIED_RESET_STEP/Quadratic.lean\#L64-L83}{\texttt{YM.AnalyticBridge.VerifiedReset.Low.plain\_bill}})
and III.15.4.

\textbf{Proposition III.15.4 (bound for \(k\)-leg atoms).} Let
\(k\in\mathbb N\), let \(a\) be a \(k\)-leg atom, \(o\) a site, and
\(U\) small-field on the footprint \(o+\mathrm{stencil}(a)\). Then
\(U\mapsto a(U,o)\) is \(K\) times left-differentiable, and every left
derivative of order \(m\le K\) in admissible directions at \(o\) is
bounded by

\[\mathrm{price}(a)=|a|\,2^K\big(2^K\,\mathrm{slotCost}(P,K,\mathrm{len}(a))\,t\big)^{k}=\mathrm{start}(G,k,P,K)\;t^{k}\;|a|\;(\mathrm{len}(a)+1)^{(P+1)Kk},\]
where
\(\mathrm{start}(G,k,P,K)=2^K\big(2^K\mathrm{slotCost}(P,K,0)\big)^k\).
\emph{Status:} theorem. \emph{Lean:}
\href{https://github.com/yang-mills-group/OpenAI-YM-20261028/blob/3456e5e75cf8d4cd558955a5c7638fba2d68f464/YM/Next/SCHEDULE/TensorCost.lean\#L56-L71}{\texttt{YM.Next.Scheduling.Wages.fin}},
\href{https://github.com/yang-mills-group/OpenAI-YM-20261028/blob/3456e5e75cf8d4cd558955a5c7638fba2d68f464/YM/Next/SCHEDULE/TensorCost.lean\#L52-L54}{\texttt{YM.Next.Scheduling.Wages.price}},
\href{https://github.com/yang-mills-group/OpenAI-YM-20261028/blob/3456e5e75cf8d4cd558955a5c7638fba2d68f464/YM/Next/SCHEDULE/TensorCost.lean\#L85-L89}{\texttt{YM.Next.Scheduling.Wages.scan}},
\href{https://github.com/yang-mills-group/OpenAI-YM-20261028/blob/3456e5e75cf8d4cd558955a5c7638fba2d68f464/YM/Next/SCHEDULE/TensorCost.lean\#L78}{\texttt{YM.Next.Scheduling.Wages.startup}}.
\emph{Used in:} III.9
(\href{https://github.com/yang-mills-group/OpenAI-YM-20261028/blob/3456e5e75cf8d4cd558955a5c7638fba2d68f464/YM/AnalyticBridge/VERIFIED_RESET_STEP/Higher.lean\#L35-L41}{\texttt{YM.AnalyticBridge.VerifiedReset.High.paid}})
and III.10.

\paragraph{Part C: launch data for the remainder
comparison}\label{part-c-launch-data-for-the-remainder-comparison}

A \emph{connected-activity target} \(Q\)
(\href{https://github.com/yang-mills-group/OpenAI-YM-20261028/blob/3456e5e75cf8d4cd558955a5c7638fba2d68f464/YM/AnalyticNext/CONNECTED_ACTIVITY_JETS_E6/OptionCost.lean\#L62-L76}{\texttt{YM.ConnectedActivityE6.Target}},
III.10) carries:

\begin{itemize}
\tightlist
\item
  a rate \(\rho_Q\in(-\tfrac12,0)\);
\item
  an exponent \(\mathrm{e}_Q\ge0\);
\item
  small parameters \(\varepsilon_Q,\delta_Q>0\) with \(4\varepsilon_Q\)
  below the top support exponent;
\item
  a halo \(h_Q\ge500\);
\item
  a chart table with its local-ratio property. The target also requires
  \(8\varepsilon_Q\) below both input chart widths.
\end{itemize}

Its relative exponent is
\(\mathrm{rel}_Q=\mathrm e_Q+2\varepsilon_Q+\delta_Q\).

Stage pairs \((j,i)\) are those with \(i\in\{j,j-1\}\)
(\href{https://github.com/yang-mills-group/OpenAI-YM-20261028/blob/3456e5e75cf8d4cd558955a5c7638fba2d68f464/YM/Push/PRICE_COMPILER/Basic.lean\#L27}{\texttt{YM.Push.PriceCompiler.StepStages}}).
The stage profile \(P_j\) has support exponent \(\sigma_j\).

\textbf{Proposition III.15.5 (launch data).} Let \(G\) be simple. Fix:

\begin{itemize}
\tightlist
\item
  a sheet with stage profiles \(P_j\) and a kernel choice;
\item
  a compatible price table and targets \(R_{ji}\) with
  \(\tfrac38+\rho_{R_{ji}}<0\) for all \(j,i\in\mathbb N\);
\item
  tolerances \(h_{ji}>0\);
\item
  libraries \(C_{ji}\) (of the order of \(P_i\));
\item
  mark numbers \(r_{ji}\).
\end{itemize}

Then there are thresholds \(f:\mathbb N\to\mathbb N\) and constants
\(C'_{ji}\ge1\) such that:

\begin{enumerate}
\def\labelenumi{\arabic{enumi}.}
\tightlist
\item
  For all \(j,i\), the \emph{quotation bound} holds for
  \((C_{ji},\mathrm{power}(P_i),r_{ji})\) with fine exponent
  \(\sigma_j\), coarse exponent \(\sigma_i+2\) and constant \(C'_{ji}\).
  More precisely, take an RG step of side \(L\), fine and coarse
  parameter points \(B,P\), and supports \(X,Y\) such that \(Y\) lies in
  the \(r_{ji}\mathcal B(x)\)-neighbourhood of the coarse image of
  \(X\). Assume \(P\) has the power of \(P_i\), \(B\) has support
  exponent \(\sigma_j\), \(B_P=\mathcal B(x)\), \(B_B\le2B_P\),
  \(t_P\le2t_B\), and the acceptance condition in (c) below. For every
  \(u\ge0\) satisfying the local analytic smallness condition for the
  return fee and every \(\tau\ge0\),
  \[e^{(\sigma_i+2)N_{B_P}(Y)}\,\tau\,\mathrm{returnFee}\le C'_{ji}L^{-5}e^{\sigma_jN_{B_B}(X)}\tau.\]
\item
  For every stage pair \((j,i)\) and every block side \(L\ge f_j\):

  \begin{itemize}
  \tightlist
  \item
    \begin{enumerate}
    \def\labelenumii{(\alph{enumii})}
    \tightlist
    \item
      the quay condition (III.15.6) holds for \(R_{ji}\) at side \(L\);
    \end{enumerate}
  \item
    (b)
    \((\mathrm{rel}_{R_{ji}}+1)\,10\,c_{h+5}/L+2\varepsilon_{R_{ji}}\le\sigma_j\),
    where \(c_{h+5}\) is the pre-cost of halo \(h_{R_{ji}}+5\);
  \item
    (c) the acceptance condition
    \(20(\sigma_i+2)\,\mathrm{toll}(r_{ji})\le L\sigma_j\);
  \item
    (d) \(C'_{ji}/L\le h_{ji}\).
  \end{itemize}
\end{enumerate}

\emph{Status:} theorem. The construction run calls it with rates
\(<-\tfrac7{16}\); these are supplied by III.2's company selection, and
\(-\tfrac7{16}\) implies the hypothesis. \emph{Lean:}
\href{https://github.com/yang-mills-group/OpenAI-YM-20261028/blob/3456e5e75cf8d4cd558955a5c7638fba2d68f464/YM/RemainderComparison/Chron/Age.lean\#L98-L111}{\texttt{YM.RemainderComparison.Chron.sails}},
\href{https://github.com/yang-mills-group/OpenAI-YM-20261028/blob/3456e5e75cf8d4cd558955a5c7638fba2d68f464/YM/RemainderComparison/Chron/Age.lean\#L94-L95}{\texttt{YM.RemainderComparison.Chron.Pledge}},
\href{https://github.com/yang-mills-group/OpenAI-YM-20261028/blob/3456e5e75cf8d4cd558955a5c7638fba2d68f464/YM/PhysicalNext/V9_JOINT_ENDPOINT_LOAD/Heir/Choirs.lean\#L23-L27}{\texttt{YM.PhysicalNext.JointLoad.Heir.Release.Quotation}},
\href{https://github.com/yang-mills-group/OpenAI-YM-20261028/blob/3456e5e75cf8d4cd558955a5c7638fba2d68f464/YM/PhysicalNext/V9_JOINT_ENDPOINT_LOAD/Heir/Fusion.lean\#L59}{\texttt{YM.PhysicalNext.JointLoad.Heir.Release.accepts}},
\href{https://github.com/yang-mills-group/OpenAI-YM-20261028/blob/3456e5e75cf8d4cd558955a5c7638fba2d68f464/YM/RemainderComparison/Lyric/Orbit.lean\#L35-L36}{\texttt{YM.RemainderComparison.Lyric.radiusCost}},
\href{https://github.com/yang-mills-group/OpenAI-YM-20261028/blob/3456e5e75cf8d4cd558955a5c7638fba2d68f464/YM/RemainderComparison/Chron/Age.lean\#L68-L71}{\texttt{YM.RemainderComparison.Chron.Contemple}}
(tolerances).

The \emph{heart condition} on rays for a stage \(E'\) and exponent
\(\mathrm e_Q\) consists of the support inequalities

\[0\le\sigma_*,\quad \sigma_{E'}\le\sigma_*,\quad \sigma_*+1\le a_*,\quad \varepsilon_*>0,\quad a_*+3\varepsilon_*\le\mathrm e_Q,\]
and the order inequalities

\[\mathrm{ord}(E')\le J_*,\quad \mathrm{ord}(E')+(p_*+1)+p_*+2\le n_*,\quad p_*+1\le s_*,\quad k_*+1\le s_*.\]
The starred quantities are the support exponents, slack and requested
orders carried by the rays.

\textbf{Proposition III.15.6 (the quay for large sides).} Let \(G\) be
simple, \(K\) a kernel choice, and \(Q\) a target with
\(\tfrac38+\rho_Q<0\). There is \(n\) such that for every block side
\(L\ge n\) the following holds. Take any:

\begin{itemize}
\tightlist
\item
  patch \(p\) of side \(L\) satisfying the branching condition;
\item
  \(W,\mu>0\), and order datum at \(W\);
\item
  songbook and scrollwork at side \(L\), and \(\chi\);
\item
  stages \(E,E'\), and rays with the heart property for
  \((E',\mathrm e_Q)\);
\item
  library \(C\) of the order of \(E'\);
\item
  record with parameter \(\varepsilon\le\tfrac14\);
\item
  \(j\in\mathbb N\), reserve \(f\), and \(\eta,\delta>0\).
\end{itemize}

Then there exist a quarry \(D\) (a record of bounding functions of
\(x\), with \(d(D)>0\) and paint supply) and a threshold \(P_1\) with:

\begin{itemize}
\tightlist
\item
  \(D\) is financed at rate \(\rho_Q\);
\item
  for every \(x\ge P_1\), the poster comparison and chart conditions and
  the instructions hold. For every
  \(\beta\in[\beta_F(x,L)/4,4\beta_F(x,L)]\) and every parameter point
  \(Z\) tempered to \(E'\) at \(x\), the active, resonance and tolerance
  conditions hold.
\end{itemize}

Here \emph{tempered} means that \(Z\) has cluster radius
\(\mathcal B(x)\), threshold \(t(x)\), the order, power and support
exponent of \(E'\), and threshold at most \(2t(\beta_F(x,L))\). The
active conditions comprise the local analytic estimates, readiness of
the spatial restrictions, compatibility estimates and the prescribed
tail bounds. The block-side threshold \(n\) is independent of the
auxiliary data; the quarry and \(P_1\) may depend on all of them. The
quarry is a fixed record of functions of \(x\). \emph{Status:} theorem.
\emph{Lean:}
\href{https://github.com/yang-mills-group/OpenAI-YM-20261028/blob/3456e5e75cf8d4cd558955a5c7638fba2d68f464/YM/RemainderComparison/Cadenza/Syllabus.lean\#L47-L58}{\texttt{YM.RemainderComparison.Cadenza.buoys}},
\href{https://github.com/yang-mills-group/OpenAI-YM-20261028/blob/3456e5e75cf8d4cd558955a5c7638fba2d68f464/YM/RemainderComparison/Cadenza/Syllabus.lean\#L46}{\texttt{YM.RemainderComparison.Cadenza.Quay}},
\href{https://github.com/yang-mills-group/OpenAI-YM-20261028/blob/3456e5e75cf8d4cd558955a5c7638fba2d68f464/YM/RemainderComparison/Cadenza/Syllabus.lean\#L42-L44}{\texttt{YM.RemainderComparison.Cadenza.Allot}},
\href{https://github.com/yang-mills-group/OpenAI-YM-20261028/blob/3456e5e75cf8d4cd558955a5c7638fba2d68f464/YM/RemainderComparison/Cadenza/Syllabus.lean\#L35-L41}{\texttt{YM.RemainderComparison.Cadenza.Coaxes}},
\href{https://github.com/yang-mills-group/OpenAI-YM-20261028/blob/3456e5e75cf8d4cd558955a5c7638fba2d68f464/YM/RemainderComparison/Cadenza/Sheriff.lean\#L34-L39}{\texttt{YM.RemainderComparison.Cadenza.Quarry.Financed}}.

\textbf{Proposition III.15.7 (using the launch data).} Fix a sheet,
compatible target and expense tables \(R,X\), chosen block sizes
\(L_j\), a protection, a register and a firmament. The register supplies
\(W_{ji},\mu_{ji}\) and order data, with positivity for stage pairs. The
firmament gives, for each stage pair, finitely many \emph{request
slots}; each carries a scrollwork, a songbook, grade and order data, a
reserve, positive constants, and an extra condition that holds
eventually for every financed quarry. Fix also tolerances, libraries
\(C_{ji}\), mark numbers \(r_{ji}\) and records \(u_{ji}\). Assume:

\begin{itemize}
\tightlist
\item
  launch data as in III.15.5 with thresholds \(k_j\le L_j\);
\item
  \(\varepsilon(u_{ji})\le\tfrac14\) for stage pairs with \(i\) at least
  the first stage;
\item
  \(\sigma_i+3<\mathrm e_{R_{ji}}\) for all \(j,i\).
\end{itemize}

Then for every initial delay datum \(T_0\) there are a delay datum
\(W\ge T_0\) and a quarry \(V_{jig}\) for every request slot \(g\),
indexed at all \((j,i)\), with the following properties for every delay
datum \(U\ge W\):

\begin{itemize}
\tightlist
\item
  \begin{enumerate}
  \def\labelenumi{(\alph{enumi})}
  \tightlist
  \item
    launch data with \(f_j\le L_j\);
  \end{enumerate}
\item
  (b) for every stage pair \((j,i)\) with \(j\ge4\), \(i\) at least the
  first stage, and the patch branching, and for every slot \(g\): the
  slot's rays have the heart property; \(d(V_{jig})>0\); paint supply
  holds; the allotment of III.15.6 holds with threshold
  \(\beta^*_j(U)\); and the slot's extra condition holds for
  \(x\ge\beta^*_j(U)\);
\item
  (c) for every stage pair with \(i\) at least the first stage, every
  request slot \(g\), every \(x\ge\beta^*_j(U)\) and every torus side
  \(c\ge \mathrm{shells}_j(U)\,(1+\mathcal B(x))^{\mathrm{growth}_j(U)}\),
  the slot's scope is tiled at \((x,c)\). This means \(2z_g(x)<c\) and
  \(4\,\mathrm{halo}_g(x)+16\le c\), where \(z_g\) and
  \(\mathrm{halo}_g\) belong to its scrollwork.
\end{itemize}

The order on delay data increases entry couplings, scale-entry
functions, shell counts and growth exponents componentwise. The quarries
are chosen before \(U\).

\emph{Status:} theorem. \emph{Lean:}
\href{https://github.com/yang-mills-group/OpenAI-YM-20261028/blob/3456e5e75cf8d4cd558955a5c7638fba2d68f464/YM/RemainderComparison/Chron/Corners.lean\#L33-L70}{\texttt{YM.RemainderComparison.Chron.grows}},
\href{https://github.com/yang-mills-group/OpenAI-YM-20261028/blob/3456e5e75cf8d4cd558955a5c7638fba2d68f464/YM/RemainderComparison/Chron/Corners.lean\#L23-L28}{\texttt{YM.RemainderComparison.Chron.Pillar}},
\href{https://github.com/yang-mills-group/OpenAI-YM-20261028/blob/3456e5e75cf8d4cd558955a5c7638fba2d68f464/YM/RemainderComparison/Chron/Age.lean\#L64-L66}{\texttt{YM.RemainderComparison.Chron.Firmament}},
\href{https://github.com/yang-mills-group/OpenAI-YM-20261028/blob/3456e5e75cf8d4cd558955a5c7638fba2d68f464/YM/RemainderComparison/Chron/Age.lean\#L49-L63}{\texttt{YM.RemainderComparison.Chron.Glint}}.
\emph{Use of III.15.5 and III.15.7:} III.2
(\href{https://github.com/yang-mills-group/OpenAI-YM-20261028/blob/3456e5e75cf8d4cd558955a5c7638fba2d68f464/YM/Integration/RG05Preconstruction/Expenses.lean\#L28-L119}{\texttt{YM.PhysicalNext.Repeat.Legato.R2I.RG05Joint.joinExpensesShipWithPrediction}}),
whose comparison conjunct reaches the selected-run theorem of III.14 and
is discarded by the final construction-run theorem (see the overview).

\subsubsection{2. What it uses}\label{what-it-uses-17}

\textbf{From other units.}

\begin{itemize}
\tightlist
\item
  III.1: the scale functions \(t\), \(\mathcal B\) and \(\beta_F\);
  slots and atoms.
\item
  II.1: ordered left-derivative calculus; the order-zero slot bound
  \href{https://github.com/yang-mills-group/OpenAI-YM-20261028/blob/3456e5e75cf8d4cd558955a5c7638fba2d68f464/YM/Work/G07/Stencil.lean\#L165-L176}{\texttt{YM.Work.G07.slot\_bound}}.
\item
  II.2: weighted hit bounds.
\item
  III.5: the logarithmic tail bound of the cluster expansion, used in
  Lemma 4.15.
\item
  III.10, a large set of inputs:

  \begin{itemize}
  \tightlist
  \item
    the growth classes, connected-activity targets and step records;
  \item
    the three fusion theorems, which turn the fusion input condition
    into a fused activity bound;
  \item
    the quotation constants
    \href{https://github.com/yang-mills-group/OpenAI-YM-20261028/blob/3456e5e75cf8d4cd558955a5c7638fba2d68f464/YM/PhysicalNext/V9_JOINT_ENDPOINT_LOAD/Heir/Choirs.lean\#L28-L29}{\texttt{YM.PhysicalNext.JointLoad.Heir.Release.quoting}}
    (including the factor \(L^{-5}\));
  \item
    the per-vertex weight
    \href{https://github.com/yang-mills-group/OpenAI-YM-20261028/blob/3456e5e75cf8d4cd558955a5c7638fba2d68f464/YM/PhysicalNext/RG_NEXT_COMPLETE_REPEATABLE_NONLINEAR_STEP/Height/Highway.lean\#L25}{\texttt{YM.PhysicalNext.Repeat.Height.grow}}
    \(=\beta_F^{3/8}\) and the criterion
    \href{https://github.com/yang-mills-group/OpenAI-YM-20261028/blob/3456e5e75cf8d4cd558955a5c7638fba2d68f464/YM/PhysicalNext/RG_NEXT_COMPLETE_REPEATABLE_NONLINEAR_STEP/Height/Highway.lean\#L30-L35}{\texttt{YM.PhysicalNext.Repeat.Height.pressureYoung}};
  \item
    the local derivative and connected-activity estimates, including the
    rate-\(\rho_Q\) bound on the tree sum.
  \end{itemize}
\item
  III.2: the targets \(R_{ji}\) with rates \(<-\tfrac7{16}\) and
  \(\sigma_i+3<\mathrm e_{R_{ji}}\) (the expense records of the company
  selection). Also the block-size and delay order into which III.15.5
  and III.15.7 are inserted.
\end{itemize}

\textbf{Standard mathematics.}

\begin{itemize}
\tightlist
\item
  Cluster and polymer expansions with weighted per-site bounds
  (Kotecký--Preiss, Commun. Math. Phys. 103 (1986); Brydges, \emph{A
  short course on cluster expansions}, Les Houches 1984).
\item
  The small-field/large-field decomposition and Mayer expansion used in
  the RG step. The \(L^{-5}\) estimate is an input from III.10, with its
  explicit support and analytic hypotheses.
\item
  The Leibniz rule and the chain rule for products of group elements
  along a path.
\item
  \(|e^z-1|\le2|z|\) for \(|z|\le1\).
\item
  Elementary asymptotics: \(e^{a(\log\log y)^d}\) grows faster than
  every power of \(\log y\) when \(a>0\) and \(d>1\), and
  \(e^{-c\,p(y)^a}\) decays faster than every inverse power of \(y\)
  when \(a,c>0\).
\end{itemize}

\textbf{Imported constants and assumptions.}

\begin{itemize}
\tightlist
\item
  Simplicity of \(G\). In Part C it enters through the local
  single-occurrence estimate used in Lemma 4.16. Parts (A) and (B) do
  not use it.
\item
  \(\tfrac38+\rho<0\) for the targets. The value \(\tfrac38\) is the
  exponent of the per-vertex tree weight \(\beta_F^{3/8}\) of III.10. In
  Lemma 4.17 it is used to make that weight times the activity summable.
  Since targets have \(\rho_Q>-\tfrac12\), the admissible window is
  \((-\tfrac12,-\tfrac38)\); the run uses \((-\tfrac12,-\tfrac7{16})\).
\item
  The tolerance table \(h_{ji}\) and the firmament are inputs. In the
  run they come from a fixed default
  (\href{https://github.com/yang-mills-group/OpenAI-YM-20261028/blob/3456e5e75cf8d4cd558955a5c7638fba2d68f464/YM/Integration/RiderResourceDynamics.lean\#L24-L25}{\texttt{YM.Integration.RiderResourceDynamics.defaultDynamics}}).
\end{itemize}

\subsubsection{3. Strategy}\label{strategy-17}

Part A separates the numerical loss from the fusion argument. On the
coupling band, the hard bound is \(Ce^{-rs^2}\) with \(s\ge p(y)/2\).
Retaining \(e^{-\rho s^2}\) leaves \(Ce^{-(r-\rho)s^2}\) to absorb the
support-counting factor. The soft bound decays because its negative
power of \(y\) dominates its powers of \(p(y)\) and \(\mathcal B(y)\).
The Mayer expansion then combines the hard activities with
\(e^{\ell_i}-1\); their weighted bounds add to \(e^\kappa a+2b\). Both
families may be supported on large fields. The terminology ``hard'' and
``soft'' distinguishes their estimates.

Part B represents a transported curvature slot as the projection of the
single group word \(\omega\cdot\partial\pi\cdot\omega^{-1}\). Unitarity
bounds each differentiated product without an exponential loss in path
length. Each derivative contributes one direction norm, at most
\(t(1+2\ell)^P\), and has at most \(2\ell\) possible link occurrences.
For positive derivative order, \(t^m\le t\); at order zero, the
small-field plaquette bound gives the factor \(t\). The Leibniz rule
gives \(t^k\) for a product of \(k\) slots, with a polynomial loss in
total path length.

Part C first chooses a block-side threshold and then, for each permitted
side and each fixed set of auxiliary data, chooses comparison functions
and an entry coupling. The connected-activity bounds have rate
\(\rho_Q\); the extra tree weight \(\beta_F^{3/8}\) still leaves
negative rate. A marked input family of size
\(\delta\le\beta_F^{-11/10}\) gives an anchored output bound
\((C'/L+o(1))\delta\), subject to the support and local analytic
conditions in Lemma 4.16. Clusters with at least one marked occurrence
and \(m\) further occurrences outside the inserted leaves gain an
additional factor \(S'(x)^m\) of negative rate. At each fixed expansion
order, the tail bounds are small enough to absorb any prescribed power
of \(g(\beta)\). Finally, finitely many requests at each stage are met
by enlarging the entry coupling and spatial-size thresholds.

\subsubsection{4. Steps}\label{steps-17}

\paragraph{4.1 Elementary asymptotics}\label{elementary-asymptotics}

\textbf{Lemma 4.1.} For every \(k\in\mathbb N\), as \(x\to\infty\):

\begin{enumerate}
\def\labelenumi{\arabic{enumi}.}
\tightlist
\item
  \(p(x)^{a}g(x)^{q}\to0\) for all \(a\in\mathbb R\) and \(q>0\);
\item
  \((1+\mathcal B(x))^k/p_d(x)^a\to0\) for \(d>1\) and \(a>0\), where
  \(p_d(x)=e^{(\log\log x)^d}\);
\item
  \(g(x)^{q}e^{-c\,p_d(x)^a}\to0\) for every real \(q\), \(c,a>0\) and
  \(d>1\).
\end{enumerate}

\emph{Sketch:} substitute \(u=\log\log x\), so that \(\log x=e^u\); each
claim reduces to a polynomial or exponential in \(u\) against \(e^{u}\)
or \(e^{u^d}\). \emph{Lean:}
\href{https://github.com/yang-mills-group/OpenAI-YM-20261028/blob/3456e5e75cf8d4cd558955a5c7638fba2d68f464/YM/Next/SCHEDULE/Clocks.lean}{\texttt{Clocks.lean}};
\href{https://github.com/yang-mills-group/OpenAI-YM-20261028/blob/3456e5e75cf8d4cd558955a5c7638fba2d68f464/YM/Next/SCHEDULE/Clocks.lean\#L100-L106}{\texttt{YM.Next.Scheduling.Clock.light}},
\href{https://github.com/yang-mills-group/OpenAI-YM-20261028/blob/3456e5e75cf8d4cd558955a5c7638fba2d68f464/YM/Next/SCHEDULE/Clocks.lean\#L110-L120}{\texttt{YM.Next.Scheduling.Clock.ceiling\_cheaper}},
\href{https://github.com/yang-mills-group/OpenAI-YM-20261028/blob/3456e5e75cf8d4cd558955a5c7638fba2d68f464/YM/Next/SCHEDULE/Clocks.lean\#L141-L152}{\texttt{YM.Next.Scheduling.Clock.burnt}}.

\paragraph{4.2 Pair-step rates}\label{pair-step-rates}

\textbf{Lemma 4.2 (III.15.1).} Under the hypotheses of III.15.1, the
rate conditions hold eventually, uniformly over the coupling band and
\(B\le\mathcal B(y)\).

\emph{Proof sketch:}

\begin{enumerate}
\def\labelenumi{\arabic{enumi}.}
\tightlist
\item
  Put \(\eta=(r-\rho)/4\). On the band, \(g(\beta)\le2g(y)\), so
  \(s\ge p(y)/2\). Hence
  \(e^{\kappa}\,\mathrm{strong}=Ce^{-(r-\rho)s^2}\le h(y):=Ce^{-\eta p(y)^2}\).
\item
  Put \(v=h+2\,\mathrm{mild}\). For every \(k\),
  \((1+\mathcal B)^k v\to0\):

  \begin{itemize}
  \tightlist
  \item
    for \(h\), write
    \((1+\mathcal B)^k e^{-\eta p^2}=\big[(1+\mathcal B)^k g\big]\cdot\big[g^{-1}e^{-\eta p^2}\big]\)
    and use Lemma 4.1(1, 2, 3);
  \item
    for \(\mathrm{mild}\), write
    \((1+\mathcal B)^{d+k}p^e g^q=\big[(1+\mathcal B)^{d+k}/p\big]\cdot\big[p^{1+e}g^q\big]\).
  \end{itemize}
\item
  With \(k=4\): \((2B+1)^4\gamma\le2^4(1+\mathcal B)^4v\le1\)
  eventually.
\item
  \(e^{v}-1\le1\) eventually gives
  \(\mathrm{quote}(\kappa,\gamma)\le e^{-\kappa}\), and
  \(\mathrm{mild}\to0\) gives \(\mathrm{mild}\le1\).
\end{enumerate}

All thresholds are independent of \(\beta\) and of
\(B\le\mathcal B(y)\). \emph{Lean:}
\href{https://github.com/yang-mills-group/OpenAI-YM-20261028/blob/3456e5e75cf8d4cd558955a5c7638fba2d68f464/YM/Progress/DIRECTED_RECIPE/Economy.lean\#L76-L139}{\texttt{YM.Progress.DirectedRecipe.Payments.Pairing.currency}},
the two private asymptotic estimates in
\href{https://github.com/yang-mills-group/OpenAI-YM-20261028/blob/3456e5e75cf8d4cd558955a5c7638fba2d68f464/YM/Progress/DIRECTED_RECIPE/Economy.lean}{\texttt{Economy.lean}}.
The same file also defines the record of these rate conditions along the
chosen sequence of scales.

\paragraph{4.3 Merging hard and soft
families}\label{merging-hard-and-soft-families}

\textbf{Lemma 4.3 (blend).} Let \(p,q\ge0\) satisfy
\(\mathrm{Hit}_{B}(p;\sigma+1)\le a\) on \((X_i)\) and
\(\mathrm{Hit}_{B}(q;\sigma+1)\le b\) on \((X'_i)\). Then the weights
\((e^{\kappa}p_i)\sqcup(2q_i)\) on the disjoint union satisfy
\(\mathrm{Hit}_B(\cdot\,;\sigma+1)\le e^{\kappa}a+2b\).

\emph{Sketch:} multiply a hit bound by a constant, then add two summable
families. \emph{Lean:}
\href{https://github.com/yang-mills-group/OpenAI-YM-20261028/blob/3456e5e75cf8d4cd558955a5c7638fba2d68f464/YM/Progress/DIRECTED_RECIPE/Tender.lean\#L78-L86}{\texttt{YM.Progress.DirectedRecipe.Payments.give}},
\href{https://github.com/yang-mills-group/OpenAI-YM-20261028/blob/3456e5e75cf8d4cd558955a5c7638fba2d68f464/YM/Progress/DIRECTED_RECIPE/Tender.lean\#L88}{\texttt{YM.Progress.DirectedRecipe.Payments.blend}},
\href{https://github.com/yang-mills-group/OpenAI-YM-20261028/blob/3456e5e75cf8d4cd558955a5c7638fba2d68f464/YM/Progress/DIRECTED_RECIPE/Tender.lean\#L98-L105}{\texttt{YM.Progress.DirectedRecipe.Payments.merges}}.

\textbf{Lemma 4.4 (activity bounds of the merged family).} Under the
nonempty-support and positive-exponent conventions of III.15.2, assume
\(p_i,q_i\ge0\), \(|F_i|\le p_i\), \(|\ell_i|\le q_i\), the soft hit
bound \(b\), and \(b\le1\). Then
\(\|F_i\|\le e^{-\kappa}(e^{\kappa}p_i)\) and \(|e^{\ell_i}-1|\le2q_i\).

\emph{Sketch:}

\begin{enumerate}
\def\labelenumi{\arabic{enumi}.}
\tightlist
\item
  Dropping one unit of support exponent gives a sum bound for the
  weights \(e^{\sigma N(X'_i)}q_i\).
\item
  A single term is at most the sum, so
  \(|\ell_i|\le q_i\le e^{\sigma N(X'_i)}q_i\le b\le1\).
\item
  Then \(|e^{z}-1|\le2|z|\) for \(|z|\le1\).
\end{enumerate}

\emph{Lean:}
\href{https://github.com/yang-mills-group/OpenAI-YM-20261028/blob/3456e5e75cf8d4cd558955a5c7638fba2d68f464/YM/Progress/DIRECTED_RECIPE/Tender.lean\#L107-L117}{\texttt{YM.Progress.DirectedRecipe.Payments.shares}},
\href{https://github.com/yang-mills-group/OpenAI-YM-20261028/blob/3456e5e75cf8d4cd558955a5c7638fba2d68f464/YM/Progress/DIRECTED_RECIPE/Tender.lean\#L64-L76}{\texttt{YM.Progress.DirectedRecipe.Payments.drop\_reserve}}.

\textbf{Lemma 4.5 (III.15.2).} Under the hypotheses of III.15.2, the
fusion input condition holds, with witness
\((\kappa,\ e^{\kappa}a+2b,\ a)\).

\emph{Sketch:} the original-family estimate combines the structural
assumptions, the regular bound \(\delta\), the hard sup bounds and the
reduction of the support exponent by one. The merged-family estimate
uses \(\gamma\ge0\) and \(\kappa\ge0\), the capital inequality, Lemma
4.3 and Lemma 4.4. The quote inequality is the last rate condition.
\emph{Lean:}
\href{https://github.com/yang-mills-group/OpenAI-YM-20261028/blob/3456e5e75cf8d4cd558955a5c7638fba2d68f464/YM/Progress/DIRECTED_RECIPE/Tender.lean\#L124-L133}{\texttt{YM.Progress.DirectedRecipe.Payments.tender}}.

\paragraph{4.4 Slots and atoms}\label{slots-and-atoms}

\textbf{Lemma 4.6 (a slot is one group word).} For every slot \(s\),
site \(o\) and configuration \(U\),

\[s(U,o)=\Pi\big(U_{\omega\cdot\partial\pi\cdot\omega^{-1}}\big),\qquad \Pi(A)=\mathcal P_G\big(\tfrac12(A-A^{*})\big),\qquad \|\Pi(A)\|\le\kappa_G^{-1}\|A\|.\]
Here \(\mathcal P_G\) is the colour-coordinate projection with the
normalization of \(G\).

\emph{Sketch:} \(\mathcal C(hgh^{-1})=\mathrm{Ad}(h)\,\mathcal C(g)\).
The detour has length \(2|\omega|+4\le2|s|\). \emph{Lean:}
\href{https://github.com/yang-mills-group/OpenAI-YM-20261028/blob/3456e5e75cf8d4cd558955a5c7638fba2d68f464/YM/Next/SCHEDULE/SlotPay.lean\#L64-L65}{\texttt{YM.Next.Scheduling.Wages.detour}},
\href{https://github.com/yang-mills-group/OpenAI-YM-20261028/blob/3456e5e75cf8d4cd558955a5c7638fba2d68f464/YM/Next/SCHEDULE/SlotPay.lean\#L67-L68}{\texttt{YM.Next.Scheduling.Wages.projector}},
\href{https://github.com/yang-mills-group/OpenAI-YM-20261028/blob/3456e5e75cf8d4cd558955a5c7638fba2d68f464/YM/Next/SCHEDULE/SlotPay.lean\#L70-L80}{\texttt{YM.Next.Scheduling.Wages.hangs}},
\href{https://github.com/yang-mills-group/OpenAI-YM-20261028/blob/3456e5e75cf8d4cd558955a5c7638fba2d68f464/YM/Next/SCHEDULE/SlotPay.lean\#L82-L88}{\texttt{YM.Next.Scheduling.Wages.norm\_projector}}.

\textbf{Lemma 4.7 (derivatives of a word).} Let \(w\) be a product of
link variables and inverses with \(n\) link uses, and let \(T\ge0\).
Write \(\iota_G\) for the colour-to-matrix map. Suppose the matrix
representative of every increment satisfies
\(\|\iota_G(A_e)\|=\kappa_G\|A_e\|\le T\) on every used link \(e\), for
every direction in the list \(\vec A\). Then

\[\|\partial_{\vec A}U_w\|\le\|1\|\,(Tn)^{|\vec A|}.\]

\emph{Sketch:} the Leibniz rule; each derivative replaces one factor by
its left increment. \emph{Lean:}
\href{https://github.com/yang-mills-group/OpenAI-YM-20261028/blob/3456e5e75cf8d4cd558955a5c7638fba2d68f464/YM/Next/SCHEDULE/GridCost.lean\#L184-L197}{\texttt{YM.Next.Scheduling.Wages.Path.bill}}.

\textbf{Lemma 4.8 (III.15.3).} Under the hypotheses of III.15.3, the
slot derivative bound holds.

\emph{Sketch:}

\begin{enumerate}
\def\labelenumi{\arabic{enumi}.}
\tightlist
\item
  Every used link basepoint of the detour is within distance \(2\ell\)
  of \(o\). Admissible directions therefore give
  \(T=\kappa_G t(1+2\ell)^P\), and
  \(T\cdot2\ell\le\mathrm{margin}\cdot t\).
\item
  Lemma 4.7, then \(\Pi\), give
  \(\kappa_G^{-1}\|1\|(\mathrm{margin}\,t)^m\).
\item
  Since \(\mathrm{margin}\ge1\), \(t\le1\) and \(1\le m\le K\), this is
  at most \(\kappa_G^{-1}\|1\|\,\mathrm{margin}^K t\).
\item
  For \(m=0\), use the order-zero small-field slot bound.
\end{enumerate}

\emph{Lean:}
\href{https://github.com/yang-mills-group/OpenAI-YM-20261028/blob/3456e5e75cf8d4cd558955a5c7638fba2d68f464/YM/Next/SCHEDULE/SlotPay.lean\#L108-L159}{\texttt{YM.Next.Scheduling.Wages.slotPays}}.

\textbf{Lemma 4.9 (III.15.4).} Under the hypotheses of III.15.4, the
atom derivative bound and its closed form hold.

\emph{Sketch:}

\begin{enumerate}
\def\labelenumi{\arabic{enumi}.}
\tightlist
\item
  Extract colour components by an orthonormal projection.
\item
  Apply the Leibniz rule to the product of the \(k\) leg components: two
  derivative bounds \(c_1,c_2\) up to order \(K\) give \(2^Kc_1c_2\).
\item
  Multiply by \(c_b\) and sum over \(b\).
\item
  The closed form follows from
  \(\mathrm{slotCost}(P,K,\ell)=\mathrm{slotCost}(P,K,0)(\ell+1)^{(P+1)K}\).
\end{enumerate}

\emph{Lean:}
\href{https://github.com/yang-mills-group/OpenAI-YM-20261028/blob/3456e5e75cf8d4cd558955a5c7638fba2d68f464/YM/Next/SCHEDULE/TensorCost.lean\#L56-L71}{\texttt{YM.Next.Scheduling.Wages.fin}},
\href{https://github.com/yang-mills-group/OpenAI-YM-20261028/blob/3456e5e75cf8d4cd558955a5c7638fba2d68f464/YM/Next/SCHEDULE/TensorCost.lean\#L85-L89}{\texttt{YM.Next.Scheduling.Wages.scan}}.

\paragraph{4.5 The Chron chain (Part C)}\label{the-chron-chain-part-c}

The following conventions specify the conditional estimates in this
chain. An \emph{admissible step record} consists of step data with the
prescribed patch, a stamped tube, an input representation \(R\), and its
connected-activity hypotheses at \((W,x,\mu,L,Q)\). A \emph{marked
input} \((E,w,\delta)\) is a set \(E\) of regular input indices with
derivative bounds \(w_i\ge0\), zero outside \(E\), satisfying

\[\sum_{i:\,o_i=z}e^{\sigma_RN_{B_R}(X_i)}w_i\le\delta\qquad\text{for every fine site }z.\]
Here \(o_i,X_i\) are the input anchors and supports, and the derivative
bounds use the input parameter point. In particular, \(\delta\ge0\).

For a finite family of inserted leaves with supports \(Z_\alpha\),
nonnegative coefficients \(l_\alpha\), attached root sets \(A_\alpha\)
and leaf weights \(\omega_\alpha\), the \emph{leaf-scale condition} with
parameters \((n,\sigma,Y)\) requires \(Y\ge0\), at most \(n\) roots over
each coarse site, and

\[\sum_{\alpha:\,v\in A_\alpha}e^{\sigma N_{B_R}(Z_\alpha)}\omega_\alpha\le Y\qquad\text{for every root }v.\]
The weights \(\omega_\alpha\) include the leaf norm and \(l_\alpha\),
and each positive weight must have an attached root lying in its leaf
support. The \emph{radius condition at \(\sigma\)} is
\((\mathrm{rel}_Q+1)10c_{h_Q+5}/L+2\varepsilon_Q\le\sigma\).

A cluster together with chosen leaf insertions has \(N\) \emph{remaining
occurrences}: the sum, over its vertices, of the numbers of input
occurrences outside those insertions. It is \emph{marked} if at least
one remaining occurrence comes from \(E\). Thus \(N\ge1\) for a marked
cluster, but the other remaining occurrences need not be marked. The
insertion count is the total number of leaves, counted with
multiplicity. These are the two counts used below.

\textbf{Lemma 4.10 (quotation and acceptance).} Let \(t,u\) be stage
profiles, \(C\) a library of order \(\mathrm{ord}(u)\),
\(r'\in\mathbb N\) and \(h>0\). There are \(E\ge1\) and \(k\) such that:

\begin{itemize}
\tightlist
\item
  the quotation bound holds with exponents \((\sigma_t,\sigma_u+2)\) and
  constant \(E\) under the support and analytic hypotheses of III.15.5;
\item
  for every \(L\ge k\): \(20(\sigma_u+2)\mathrm{toll}(r')\le L\sigma_t\)
  and \(E/L\le h\).
\end{itemize}

\emph{Sketch:} the quotation theorem of III.10 gives an
\(L\)-independent \(E\). Then \(E/L\to0\) and \(L\sigma_t\to\infty\).
The return-fee estimate gives the factor \(L^{-5}\). Summing over the
\(L^4\) fine sites of a block leaves \(E/L\), by the anchored-sum
estimate. \emph{Lean:} inner step of
\href{https://github.com/yang-mills-group/OpenAI-YM-20261028/blob/3456e5e75cf8d4cd558955a5c7638fba2d68f464/YM/RemainderComparison/Chron/Age.lean\#L98-L111}{\texttt{YM.RemainderComparison.Chron.sails}};
\href{https://github.com/yang-mills-group/OpenAI-YM-20261028/blob/3456e5e75cf8d4cd558955a5c7638fba2d68f464/YM/PhysicalNext/V9_JOINT_ENDPOINT_LOAD/Heir/Choirs.lean\#L28-L29}{\texttt{YM.PhysicalNext.JointLoad.Heir.Release.quoting}}.

\textbf{Lemma 4.11 (radius cost).} If \(b>0\) and
\(b+2\varepsilon_Q\le v\), then
\((\mathrm{rel}_Q+1)10\,c_{h_Q+5}/L+2\varepsilon_Q\le v\) for all large
\(L\). In III.15.5, take \(b=\varepsilon_Q\): the target condition gives
\(3\varepsilon_Q<4\varepsilon_Q<\sigma_j\).

\emph{Sketch:} multiply the desired inequality by \(L>0\) and choose
\(L\) so large that \((\mathrm{rel}_Q+1)10c_{h_Q+5}\le bL\).
\emph{Lean:}
\href{https://github.com/yang-mills-group/OpenAI-YM-20261028/blob/3456e5e75cf8d4cd558955a5c7638fba2d68f464/YM/RemainderComparison/Lyric/Orbit.lean\#L37-L42}{\texttt{YM.RemainderComparison.Lyric.bidLift}}.

\textbf{Lemma 4.12 (comparison functions).} Fix a target \(Q\), a patch,
stages \(E,E'\), a kernel choice, an order datum at \(W\), a songbook, a
scrollwork, \(\mu>0\), and rays satisfying the heart condition for
\((E',\mathrm e_Q)\). There is a poster \(D\) such that \(d(D)>0\), its
paint-supply condition holds, and its radius, charge and four comparison
functions are reserves. Eventually in \(x\), it satisfies the
conditional comparison statement, the chart-radius estimates and the
axis-chart existence statement required by the poster condition.

\emph{Sketch:} obtain the radius and chart estimates from the chart
construction, and choose the paint constant. The comparison-function
theorem applies when the chart parameter is positive; otherwise the
chart hypotheses are inconsistent. It requires positive comparison
parameters \(a,d\), a subpower bound on the volume function, reserves
for the radius, charge and derivative control, paint supply, and orders
satisfying \(j+(r+1)+r+2\le n\) and \(n+N\le k\). The ray heart
condition supplies these order inequalities. It produces four reserve
functions, uniformly over the subsequent step and chart choices. Its
conclusion is conditional on local derivative bounds, the chart
estimates and coefficient matching: derivatives through order \(j\) of
the differences of the expansion coefficients through degree \(r\) have
the prescribed bound. The proof combines finite-order exponential
bounds, analytic expansion estimates and the assumed coefficient match.
\emph{Lean:}
\href{https://github.com/yang-mills-group/OpenAI-YM-20261028/blob/3456e5e75cf8d4cd558955a5c7638fba2d68f464/YM/RemainderComparison/Cadenza/Metro.lean\#L29-L51}{\texttt{YM.RemainderComparison.Cadenza.Poster.quorum}},
\href{https://github.com/yang-mills-group/OpenAI-YM-20261028/blob/3456e5e75cf8d4cd558955a5c7638fba2d68f464/YM/RemainderComparison/Herbal/Orbit.lean\#L41-L56}{\texttt{YM.RemainderComparison.Herbal.compile}},
\href{https://github.com/yang-mills-group/OpenAI-YM-20261028/blob/3456e5e75cf8d4cd558955a5c7638fba2d68f464/YM/RemainderComparison/Herbal/Orbit.lean\#L33-L39}{\texttt{YM.RemainderComparison.Herbal.Airstrip}},
\href{https://github.com/yang-mills-group/OpenAI-YM-20261028/blob/3456e5e75cf8d4cd558955a5c7638fba2d68f464/YM/RemainderComparison/Portal/Wire.lean\#L80-L81}{\texttt{YM.RemainderComparison.Portal.JMatch}}
\emph{(proof read; proofs of the analytic expansion estimates not
expanded)}.

\textbf{Lemma 4.13 (connected-activity structures).} Fix an input
specification and target \(Q\). There is \(l\) such that, for every
\(L\ge l\), branching patch, order \(j\), auxiliary choice \(\chi\),
\(W,\mu>0\) and order datum at \(W\), one can choose constants and
functions before \(x\) as follows: \(d,v\) have gear \(-\tfrac12\);
\(B,r,T'\) have rate at most \(0\); and \(p\) is a reserve, with
\(T',p\ge0\). Eventually in \(x\), every admissible step record admits a
connected-activity structure with the heavy, derivative and marked-input
estimates specified by its three accompanying conditions.

\emph{Sketch:} combine the connected-activity existence theorem with the
derivative estimates for its local factors and the marked-input
estimate. The growth-class rules keep \(d,v\) at gear \(-\tfrac12\) and
the remaining stated functions at rate at most \(0\). The structure is
chosen separately for each step record; the bounding functions are
common to all those records. \emph{Lean:}
\href{https://github.com/yang-mills-group/OpenAI-YM-20261028/blob/3456e5e75cf8d4cd558955a5c7638fba2d68f464/YM/RemainderComparison/Halo/Nexus.lean\#L12-L57}{\texttt{YM.RemainderComparison.Halo.seeds}}
\emph{(proof read; analytic inputs from III.10)}.

\textbf{Lemma 4.14 (remaining-occurrence bound).} There is a block-side
threshold depending only on the input specification and \(Q\) such that,
for every \(L\) above it and every branching patch, the following holds.
Fix the remaining data of Lemma 4.13 and \(a\ge0\), \(\varepsilon>0\)
with \(a+2\varepsilon\le\mathrm e_Q\). One can choose bounding functions
such that, for every fixed pair of orders \(I,J\), reserve \(Y\) and
nonnegative integer-valued function \(n\) whose real-valued cast is a
reserve, there are a reserve \(V\) and a function \(S'\) of gear \(q<0\)
with the following property. Eventually in \(x\), every admissible step
record admits a connected-activity structure for which the following
holds. For every finite leaf family satisfying the leaf-scale condition
\((n(x),\sigma,Y(x))\) and the radius condition at \(\sigma\), every
marked input with \(\delta\le\beta_F^{-11/10}\), and every
\(m\in\mathbb N\), the bills restricted to marked clusters with
\(N\ge1+m\) remaining occurrences and insertion count at most \(I\)
satisfy

\[\mathrm{Hit}_{\mathcal B(x)}(\mathrm{bill};a)\le V(x)S'(x)^m\delta.\]
The derivative order is \(j\) and the support-polynomial order in the
bill is \(J\). The bill is formed from the local majorant supplied by
the chosen structure.

\emph{Sketch:} increase the weights of the remaining occurrences by a
factor whose inverse has negative gear. On the indicated clusters the
inflated product contains at least one marked factor and \(m\) further
factors. Apply the weighted cluster bound and divide by those factors.
The leaf-scale and radius conditions control the inserted leaves; the
insertion cap \(I\) controls the accompanying polynomial losses.
\emph{Lean:}
\href{https://github.com/yang-mills-group/OpenAI-YM-20261028/blob/3456e5e75cf8d4cd558955a5c7638fba2d68f464/YM/RemainderComparison/Halo/Tome.lean\#L33-L65}{\texttt{YM.RemainderComparison.Halo.ramp}}
\emph{(proof read; weighted-cluster estimates used as inputs)}.

\textbf{Lemma 4.15 (restriction to a cluster predicate).} Fix an
admissible step record, a connected-activity structure, a nonnegative
local majorant satisfying its derivative condition, nonnegative leaf
coefficients and nonnegative vertex weights \(\gamma\). Let \(a\ge0\),
\(e_1,\delta>0\) and \(q'<1\). Suppose the polymer majorant has hit
bound \(e_1\) at exponent \(a+2\delta\), and its degree price with slack
\(\delta\) is at most \(q'\). If a cluster predicate \(V\) and \(v\ge0\)
satisfy \(1\le v\prod\gamma\) on \(V\), then the restricted bills have
hit bound at exponent \(a\) at most

\[v A_\delta(x;j,m)e_1(1-q')^{-1},\qquad A_\delta(x;j,m)=[4(2\mathcal B(x)+1)]^{mj}(mj)!\,\delta^{-mj}e^\delta.\]
The bill uses derivative order \(j\) and support-polynomial order \(m\).

\emph{Sketch:} the logarithmic cluster-tail estimate gives the geometric
sum \(e_1/(1-q')\). Absorb \((1+N)^{mj}\) into \(e^{\delta N}\), using
\(N^{mj}e^{-\delta N}\le(mj)!\delta^{-mj}\) with the corresponding shift
by \(1\). The spatial factor is \([4(2\mathcal B(x)+1)]^{mj}\). Finally
insert \(v\prod\gamma\ge1\) on the retained clusters. \emph{Lean:}
\href{https://github.com/yang-mills-group/OpenAI-YM-20261028/blob/3456e5e75cf8d4cd558955a5c7638fba2d68f464/YM/RemainderComparison/Lyric/Salary.lean\#L49-L76}{\texttt{YM.RemainderComparison.Lyric.gathers}},
\href{https://github.com/yang-mills-group/OpenAI-YM-20261028/blob/3456e5e75cf8d4cd558955a5c7638fba2d68f464/YM/RemainderComparison/Lyric/Salary.lean\#L82-L100}{\texttt{YM.RemainderComparison.Lyric.run}}.

\textbf{Lemma 4.16 (anchored bound for marked outputs).} There is a
block-side threshold depending only on the input specification and \(Q\)
such that the following holds for every \(L\) above it and every
branching patch. Fix the remaining data of Lemma 4.13 and \(a\ge0\),
\(\varepsilon>0\) with \(a+3\varepsilon\le\mathrm e_Q\). The
connected-activity functions can be chosen as follows. Fix orders
\(J,r,I,n,P_0\) with \(2\le n\le j\), assume \(G\) simple, and fix an
auxiliary choice, a reserve \(Y\) and a nonnegative integer-valued
function \(i\) whose cast is a reserve, a record with parameter at most
\(\tfrac14\), and a library of order \(n\). There are a reserve \(V\)
and a function \(V'\) of some rate \(q'<0\) such that, eventually in
\(x\), every admissible step record has:

\begin{itemize}
\tightlist
\item
  the marked hit bound \(V(x)\delta\) at exponent \(a+\varepsilon\), for
  leaf families with scale \((i(x),\sigma,Y(x))\) satisfying the radius
  condition, marked inputs \(\delta\le\beta_F^{-11/10}\), and insertion
  count at most \(I\);
\item
  an anchored output bound \((C_0/L+V'(x))\delta\) at exponent \(a\),
  provided the output parameter point has radius \(\mathcal B(x)\),
  threshold \(t(x)\) and power \(P_0\le J\), the library's compatibility
  condition holds, \(r\ge h_Q+5\), and the acceptance, quotation and
  local analytic smallness conditions hold with quotation constant
  \(C_0\).
\end{itemize}

The anchored bound includes regularity and all left derivatives through
the output order on the small-field supports. The proof gives
\(q'=\max(q,-\tfrac12)\), where \(q<0\) is the gear exponent in Lemma
4.14.

\emph{Sketch:} separate clusters with a single remaining occurrence from
those with at least two, and separate supports whose covering number
exceeds \(\mathcal B(x)\). The single-occurrence contribution uses the
quotation bound and the sum over \(L^4\) sites. Lemma 4.14 bounds the
other contributions by \(V S'\) and \(e^{-\varepsilon\mathcal B}V\). A
further local comparison error has rate at most \(-\tfrac12\). Their sum
defines \(V'\). \emph{Lean:}
\href{https://github.com/yang-mills-group/OpenAI-YM-20261028/blob/3456e5e75cf8d4cd558955a5c7638fba2d68f464/YM/RemainderComparison/Solo/Barque.lean\#L21-L61}{\texttt{YM.RemainderComparison.Solo.atlas}}
\emph{(proof read; local singleton estimate used as an input)}.

\textbf{Lemma 4.17 (tail bounds).} Fix \(L\), \(\mu>0\), \(a\ge0\),
\(b>0\) with \(a+2b\le\mathrm e_Q\), and \(\tfrac38+\rho_Q<0\). Fix
orders \(j,n,k,s,J\), functions \(B,p_0\) of rate at most \(0\), and a
nonnegative integer-valued function \(m\) whose cast has rate at most
\(0\). Put \(G_*(x)=\beta_F(x,L)^{3/8}\) and \(\tau(x)=G_*(x)^5\). Let
\(I(x)\) be the additional root-activity bound formed from \(m,p_0\),
and let \(\Sigma(x)\) be the resulting tree-sum majorant with derivative
factor \(1+2^jB(x)\). Eventually in \(x\):

\begin{itemize}
\tightlist
\item
  the degree prices of \(G_*((\text{root activity})+I)\) and of
  \(\Sigma\) are at most \(\tfrac12\), with slacks \(\delta_Q\) and
  \(b\), respectively;
\item
  \(G_*\ge1\), \(I\) is below the prescribed root allowance, and
  \(G_*^{5k+2s}\le e^{\mathrm{mark}}\) and
  \(G_*^{5k+2s}\le e^{\varepsilon_Q\lfloor\mathcal B/100000\rfloor}\);
\item
  for every \(\beta\) in the band of \(\beta_F\),
  \[\tau^kG_*^{-(5k+2s)}A_b(x;j,J)\,2\Sigma(x)\le g(\beta)^s,\]
  \[\tau^k e^{-\mathrm{mark}}\,2\mathrm{pledge}(x)\le e^{-\mathrm{mark}/2}g(\beta)^s;\]
\item
  \(e^{-\mathrm{mark}/2}\le e^{-4p(x)^{1/5}}\).
\end{itemize}

\emph{Sketch:} the additional root term has rate at most \(-\tfrac32\).
Thus \(\Sigma\) has rate at most \(\tfrac38+\rho_Q<0\), and the
polynomial spatial factors have rate at most \(0\). These facts give the
two degree-price inequalities. In the first displayed tail bound the
factors \(\tau^k\) and \(G_*^{-5k}\) cancel; \(G_*^{-2s}\) supplies
\(\beta_F^{-3s/4}\), sufficient for \(g(\beta)^s\) uniformly on the
band. The large marking exponent absorbs the remaining powers in the
second bound. \emph{Lean:}
\href{https://github.com/yang-mills-group/OpenAI-YM-20261028/blob/3456e5e75cf8d4cd558955a5c7638fba2d68f464/YM/RemainderComparison/Lyric/Forecast.lean\#L93-L120}{\texttt{YM.RemainderComparison.Lyric.epochs}},
\href{https://github.com/yang-mills-group/OpenAI-YM-20261028/blob/3456e5e75cf8d4cd558955a5c7638fba2d68f464/YM/RemainderComparison/Lyric/Forecast.lean\#L80-L89}{\texttt{YM.RemainderComparison.Lyric.Premium}}.

\textbf{Lemma 4.18 (simultaneous derivative and tail estimates).}
Suppose \(\tfrac38+\rho_Q<0\). There is \(\ell\) such that, for every
\(L\ge\ell\), branching patch, \(W,\mu>0\), order datum, auxiliary
choice, orders \(j,J,k,s\), and \(a\ge0\), \(\varepsilon,\Delta>0\) with
\(a+2\varepsilon\le\mathrm e_Q\), the following holds. Fix a reserve
\(Y\), a nonnegative integer-valued function \(n\) whose cast is a
reserve, and a function \(p_0\) of any fixed finite rate. There are
reserves \(A,B,C\) and a function \(E\) of rate \(q<0\) such that,
eventually in \(x\), each admissible step record admits three
connected-activity structures with these properties:

\begin{itemize}
\tightlist
\item
  the first has the derivative estimates, mass condition and tail bounds
  of Lemma 4.17, and the required depth bound for \(G_*^{5k+2s}\);
\item
  \(A(x)\ge1\), and the full bill from the second has hit bound \(B(x)\)
  at exponent \(a\) for leaves whose coefficients have been divided by
  \(A(x)\), whenever the original leaf-scale condition
  \((n(x),\sigma,Y(x))\) and radius condition hold;
\item
  the third gives hit bound \(C(x)E(x)^i\delta\) at exponent \(a\) for
  marked clusters with \(N\ge1+i\), insertion count at most \(k\), and
  \(\delta\le\beta_F^{-11/10}\), under the same leaf-scale and radius
  conditions;
\item
  writing
  \(\pi_\Delta(x)=e^{-\Delta\lfloor\mathcal B(x)/100000\rfloor}\), the
  crop inequalities hold:
  \[\pi_\Delta p_0\le\beta_F^{-11/10},\qquad C\pi_\Delta p_0\le g(\beta_H)^s,\qquad \pi_\Delta A^kB\le g(\beta_H)^s.\]
\end{itemize}

The full and marked majorants also satisfy their local derivative
conditions. Here \(\beta_H\) is the coupling of the step record; the
last bounds are for each fixed \(s\).

\emph{Sketch:} use separate connected-activity structures for the mass
bound, the full bill and the marked bill. Combine the full-bill
estimate, Lemma 4.14, Lemma 4.17 and the decay of \(\pi_\Delta\) against
every function of finite rate. The three structures need not coincide.
\emph{Lean:}
\href{https://github.com/yang-mills-group/OpenAI-YM-20261028/blob/3456e5e75cf8d4cd558955a5c7638fba2d68f464/YM/RemainderComparison/Muse/Baton.lean\#L80-L96}{\texttt{YM.RemainderComparison.Muse.generation}},
\href{https://github.com/yang-mills-group/OpenAI-YM-20261028/blob/3456e5e75cf8d4cd558955a5c7638fba2d68f464/YM/RemainderComparison/Muse/Baton.lean\#L63-L77}{\texttt{YM.RemainderComparison.Muse.Roles}}
\emph{(proof read; connected-activity and full-bill estimates used as
inputs)}.

\textbf{Lemma 4.19 (extending the comparison functions).} Suppose \(G\)
is simple and \(\tfrac38+\rho_Q<0\). There is \(\ell\) such that, for
every branching patch of side at least \(\ell\), \(W,\mu>0\), order
datum, songbook, scrollwork, auxiliary choice, mark number \(r\), stage
\(E'\), rays with the heart property, library of order
\(\mathrm{ord}(E')\), record with parameter at most \(\tfrac14\), and
financed poster \(D\), there is a quarry \(F\) extending \(D\), chosen
before \(x\), such that:

\begin{itemize}
\tightlist
\item
  its new functions \(Y\), spin and \(e\) are reserves, its
  connected-activity bound \(R_b\) has rate at most \(\rho_Q\), and its
  marked-output correction \(h\) has some negative rate;
\item
  eventually in \(x\), every parameter point tempered to \(E'\) and
  every compatible admissible step record has the required local
  analytic bounds and a connected-activity structure with tree-sum bound
  at most \(R_b(x)\);
\item
  its marked-output bound is \((C_0/L+h(x))\delta\), under all the
  marked-input, compatibility, acceptance, quotation and local analytic
  hypotheses of Lemma 4.16, with \(r\ge h_Q+5\).
\end{itemize}

\emph{Sketch:} take the new reserve functions and the finite-rate tail
function from Lemma 4.18. The tree-sum bound is twice the polynomial
spatial factor times the connected-activity majorant, whose rate is at
most \(\rho_Q\). Take \(h=V'\) from Lemma 4.16 and apply that lemma with
an empty leaf family. The additional tail function is a field of the
quarry but carries no separate growth assertion in its financing
condition. \emph{Lean:}
\href{https://github.com/yang-mills-group/OpenAI-YM-20261028/blob/3456e5e75cf8d4cd558955a5c7638fba2d68f464/YM/RemainderComparison/Cadenza/Sheriff.lean\#L64-L80}{\texttt{YM.RemainderComparison.Cadenza.rollBefore}},
\href{https://github.com/yang-mills-group/OpenAI-YM-20261028/blob/3456e5e75cf8d4cd558955a5c7638fba2d68f464/YM/RemainderComparison/Cadenza/Sheriff.lean\#L57-L60}{\texttt{YM.RemainderComparison.Cadenza.Instructions}},
\href{https://github.com/yang-mills-group/OpenAI-YM-20261028/blob/3456e5e75cf8d4cd558955a5c7638fba2d68f464/YM/RemainderComparison/Cadenza/Sheriff.lean\#L47-L51}{\texttt{YM.RemainderComparison.Cadenza.Quarry.Staked}},
\href{https://github.com/yang-mills-group/OpenAI-YM-20261028/blob/3456e5e75cf8d4cd558955a5c7638fba2d68f464/YM/RemainderComparison/Cadenza/Sheriff.lean\#L44-L45}{\texttt{YM.RemainderComparison.Cadenza.Quarry.Bequest}}.

\textbf{Lemma 4.20 (III.15.6).} Under the hypotheses of III.15.6, the
quay condition holds for all sufficiently large block sides.

\emph{Sketch:} take \(n\) from 4.19. For given data, 4.12 gives a
poster, and 4.19 extends it to a quarry. The remaining step conditions
hold eventually in \(x\) by their local analytic estimates and the
scrollwork growth condition. Take \(P_1\) beyond all of these
thresholds. \emph{Lean:}
\href{https://github.com/yang-mills-group/OpenAI-YM-20261028/blob/3456e5e75cf8d4cd558955a5c7638fba2d68f464/YM/RemainderComparison/Cadenza/Syllabus.lean\#L47-L58}{\texttt{YM.RemainderComparison.Cadenza.buoys}}
\emph{(proof read; individual local analytic estimates not expanded)}.

\textbf{Lemma 4.21 (III.15.5).} Under the hypotheses of III.15.5, the
launch data exist.

\emph{Sketch:} for each \((j,i)\), take \(k_{ji}\) (4.10), \(g_{ji}\)
(4.20) and \(n_{ji}\) (4.11). Set

\[f_j=\max_{i\in\{j,j-1\}}g_{ji}+\max_{i\in\{j,j-1\}}k_{ji}+\max_{i\in\{j,j-1\}}n_{ji}\]
and take \(C'_{ji}=E_{ji}\) from 4.10. \emph{Lean:}
\href{https://github.com/yang-mills-group/OpenAI-YM-20261028/blob/3456e5e75cf8d4cd558955a5c7638fba2d68f464/YM/RemainderComparison/Chron/Age.lean\#L98-L111}{\texttt{YM.RemainderComparison.Chron.sails}}.

\textbf{Lemma 4.22 (III.15.7).} Under the hypotheses of III.15.7, the
quarries and a common enlargement of the delay data exist.

\emph{Sketch:}

\begin{enumerate}
\def\labelenumi{\arabic{enumi}.}
\tightlist
\item
  For each stage pair and slot, apply the pledge at \(L=L_j\) (allowed
  since \(k_j\le L_j\)) to the slot's data. The chosen rays have the
  heart property because \(\sigma_i+3<\mathrm e_{R_{ji}}\). This gives a
  quarry \(V_{jig}\) and a threshold. The slot's extra condition holds
  eventually.
\item
  Set
  \(\mathrm{shells}_j=\max(\mathrm{shells}_j(T_0),\max_{i\in\{j,j-1\},g}(4\ell_g+50))\)
  and \(\mathrm{growth}_j=\max(\mathrm{growth}_j(T_0),1)\). The inner
  maximum is \(0\) when there are no request slots. Tiling follows from
  the scrollwork growth conditions.
\item
  Since there are finitely many slots and two stage pairs per \(j\),
  choose \(\beta^*_j\) beyond all thresholds.
\end{enumerate}

\emph{Lean:}
\href{https://github.com/yang-mills-group/OpenAI-YM-20261028/blob/3456e5e75cf8d4cd558955a5c7638fba2d68f464/YM/RemainderComparison/Chron/Corners.lean\#L33-L70}{\texttt{YM.RemainderComparison.Chron.grows}}.

\paragraph{4.6 Records and other fragments (Part
D)}\label{records-and-other-fragments-part-d}

The construction-run record collects local estimates, growth of the copy
count, bounds for the selected representations, and the conditions
required for later OPE and comparison arguments. Its fields are
established when the run is assembled. The additional-expense records
collect the hypotheses needed on the same selected run. The readout
record extends three primitive conditions and adds a growth condition, a
bound on block sizes by an iterated logarithm of the coupling, and
cross-stage comparison inequalities.

The unit also contains numerical definitions and lemmas for the RG05
readout, the scale functions, and the instrument and postage records.
These are not all mere projections: some prove asymptotic and transport
estimates. The readout formulas and the remaining RG12 and R2I fragments
are outside the detailed statements above; see §7.

\emph{Lean:}
\href{https://github.com/yang-mills-group/OpenAI-YM-20261028/blob/3456e5e75cf8d4cd558955a5c7638fba2d68f464/YM/R2I/Campaign/Accepted.lean\#L60-L73}{\texttt{YM.PhysicalNext.Repeat.Legato.R2I.Funded}},
\href{https://github.com/yang-mills-group/OpenAI-YM-20261028/blob/3456e5e75cf8d4cd558955a5c7638fba2d68f464/YM/R2I/Campaign/Expenses.lean\#L35-L45}{\texttt{YM.PhysicalNext.Repeat.Legato.R2I.Extras}},
\href{https://github.com/yang-mills-group/OpenAI-YM-20261028/blob/3456e5e75cf8d4cd558955a5c7638fba2d68f464/YM/Integration/RG05Preconstruction/Resources.lean\#L41-L43}{\texttt{YM.Integration.RG05Preconstruction.PrimitiveExtras}},
\href{https://github.com/yang-mills-group/OpenAI-YM-20261028/blob/3456e5e75cf8d4cd558955a5c7638fba2d68f464/YM/Integration/RG05Preconstruction/Resources.lean\#L28-L39}{\texttt{YM.Integration.RG05Preconstruction.ReadoutPurchase}}.

\subsubsection{5. Where the difficulty
lies}\label{where-the-difficulty-lies-17}

The active outputs III.15.1--4 require uniform estimates in the coupling
band and the volume. In III.15.1, the threshold must work for every
\(\beta\) in the band and every \(B\le\mathcal B(y)\). In III.15.3--4,
the constants depend on derivative order and path length but not on
\(t\) or the torus side. The single-word representation and unitarity
give that independence.

Part C adds the conditional comparison bounds to the selected run, but
those bounds are discarded in the final construction-run theorem. Its
main estimates concern an inherited marked family: the anchored bound
\((C'/L+o(1))\delta\) and the geometric gain per additional remaining
occurrence. Their application requires the leaf-scale, radius, quotation
and local analytic conditions. The gain counts all remaining
occurrences, with at least one marked occurrence.

The exponent \(\tfrac38\) is the tree-weight exponent inherited from
III.10. Requiring \(\rho_Q<-\tfrac38\) makes the relevant tree-sum
majorant decay; the run's choice \(\rho_Q<-\tfrac7{16}\) leaves a margin
greater than \(\tfrac1{16}\). The large-side threshold is independent of
the later auxiliary data. The comparison functions and entry-coupling
threshold may depend on those data, and the estimates then hold for all
larger \(x\).

\subsubsection{6. Notation}\label{notation-17}

{\def\LTcaptype{none} % do not increment counter
\begin{longtable}[]{@{}
  >{\raggedright\arraybackslash}p{(\linewidth - 4\tabcolsep) * \real{0.3333}}
  >{\raggedright\arraybackslash}p{(\linewidth - 4\tabcolsep) * \real{0.3333}}
  >{\raggedright\arraybackslash}p{(\linewidth - 4\tabcolsep) * \real{0.3333}}@{}}
\toprule\noalign{}
\begin{minipage}[b]{\linewidth}\raggedright
Written proof
\end{minipage} & \begin{minipage}[b]{\linewidth}\raggedright
Lean (glossary)
\end{minipage} & \begin{minipage}[b]{\linewidth}\raggedright
Meaning
\end{minipage} \\
\midrule\noalign{}
\endhead
\bottomrule\noalign{}
\endlastfoot
\(g(\beta)\) &
\href{https://github.com/yang-mills-group/OpenAI-YM-20261028/blob/3456e5e75cf8d4cd558955a5c7638fba2d68f464/YM/Next/SCHEDULE/Clocks.lean\#L40}{\texttt{YM.Next.Scheduling.Clock.gap}},
\texttt{gStep} & \(\beta^{-1/2}\) \\
\(p(x)\) &
\href{https://github.com/yang-mills-group/OpenAI-YM-20261028/blob/3456e5e75cf8d4cd558955a5c7638fba2d68f464/YM/Next/SCHEDULE/Dues.lean\#L55}{\texttt{YM.Next.Scheduling.Instrument.P}}
(plain),
\href{https://github.com/yang-mills-group/OpenAI-YM-20261028/blob/3456e5e75cf8d4cd558955a5c7638fba2d68f464/YM/Next/SCHEDULE/Clocks.lean\#L42}{\texttt{YM.Next.Scheduling.Clock.watch}}
\(2\) & \(e^{(\log\log x)^2}\) \\
\(t(x)\) &
\href{https://github.com/yang-mills-group/OpenAI-YM-20261028/blob/3456e5e75cf8d4cd558955a5c7638fba2d68f464/YM/Pressure/Seeds.lean\#L25}{\texttt{YM.Pressure.temp}},
\href{https://github.com/yang-mills-group/OpenAI-YM-20261028/blob/3456e5e75cf8d4cd558955a5c7638fba2d68f464/YM/Next/SCHEDULE/Dues.lean\#L59}{\texttt{YM.Next.Scheduling.Instrument.t}}
& small-field threshold \(g(x)p(x)\) \\
\(\mathcal B(x)\) &
\href{https://github.com/yang-mills-group/OpenAI-YM-20261028/blob/3456e5e75cf8d4cd558955a5c7638fba2d68f464/YM/Pressure/Seeds.lean\#L22}{\texttt{YM.Pressure.radiusB}},
\href{https://github.com/yang-mills-group/OpenAI-YM-20261028/blob/3456e5e75cf8d4cd558955a5c7638fba2d68f464/YM/Next/SCHEDULE/Clocks.lean\#L38}{\texttt{YM.Next.Scheduling.Clock.cover}}
& cluster radius \(\lceil(\log x)^2\rceil\) \\
\(\beta_F(x,L)\) &
\href{https://github.com/yang-mills-group/OpenAI-YM-20261028/blob/3456e5e75cf8d4cd558955a5c7638fba2d68f464/YM/Pressure/Seeds.lean\#L23}{\texttt{YM.Pressure.betaF}}
& \(x+b_G\log L\) \\
band of \(y\) &
\href{https://github.com/yang-mills-group/OpenAI-YM-20261028/blob/3456e5e75cf8d4cd558955a5c7638fba2d68f464/YM/Next/SCHEDULE/Dues.lean\#L99}{\texttt{YM.Next.Scheduling.inBand}}
& \([y/4,4y]\) \\
\(\mathrm{Rate}_L(a)\) &
\href{https://github.com/yang-mills-group/OpenAI-YM-20261028/blob/3456e5e75cf8d4cd558955a5c7638fba2d68f464/YM/AnalyticNext/CONNECTED_ACTIVITY_JETS_E6/Numerics.lean\#L22-L23}{\texttt{YM.ConnectedActivityE6.Rate}}
& \(\lvert w\rvert\le\beta_F^{v}\) eventually, for every \(v>a\) \\
reserve; \(M(x)\) &
\href{https://github.com/yang-mills-group/OpenAI-YM-20261028/blob/3456e5e75cf8d4cd558955a5c7638fba2d68f464/YM/RemainderComparison/Access/Geometry.lean\#L25-L26}{\texttt{YM.RemainderComparison.Reserve}},
\href{https://github.com/yang-mills-group/OpenAI-YM-20261028/blob/3456e5e75cf8d4cd558955a5c7638fba2d68f464/YM/Push/ABSOLUTE_SUMS/Delay.lean\#L24}{\texttt{YM.Push.AbsoluteSums.Delay.fineMonitor}}
& \(\lvert w\rvert\le M^k\) with \(M=p(\beta_F)\) \\
gear \(a\) &
\href{https://github.com/yang-mills-group/OpenAI-YM-20261028/blob/3456e5e75cf8d4cd558955a5c7638fba2d68f464/YM/RemainderComparison/Gear/Class.lean\#L25}{\texttt{YM.RemainderComparison.Gear}}
& \(\beta_F^{-a}w\) is a reserve \\
\(\mathrm{Hit}_M(f;a)\le\gamma\) &
\href{https://github.com/yang-mills-group/OpenAI-YM-20261028/blob/3456e5e75cf8d4cd558955a5c7638fba2d68f464/YM/Ac/Tree.lean\#L12-L13}{\texttt{YM.Ac.Hit}}
& weighted per-site polymer bound \\
\(\mathrm{quote}(\kappa,\gamma)\) &
\href{https://github.com/yang-mills-group/OpenAI-YM-20261028/blob/3456e5e75cf8d4cd558955a5c7638fba2d68f464/YM/Bid/Kitchen.lean\#L18}{\texttt{YM.Bid.quote}}
& \(e^{-\kappa}(e^\gamma-1)\) \\
rate conditions \((B,w,a,b,\kappa)\) &
\href{https://github.com/yang-mills-group/OpenAI-YM-20261028/blob/3456e5e75cf8d4cd558955a5c7638fba2d68f464/YM/Progress/DIRECTED_RECIPE/Tender.lean\#L33-L39}{\texttt{YM.Progress.DirectedRecipe.Payments.Rates}}
& the five inequalities of III.15.1 \\
mild, strong, \(\kappa\) &
\href{https://github.com/yang-mills-group/OpenAI-YM-20261028/blob/3456e5e75cf8d4cd558955a5c7638fba2d68f464/YM/Progress/DIRECTED_RECIPE/Economy.lean\#L28-L36}{\texttt{YM.Progress.DirectedRecipe.Payments.Pairing}}
(\texttt{mild}, \texttt{strong}, \texttt{kGoal}) & power-decaying and
Gaussian pair-step costs \\
hard and soft family &
\href{https://github.com/yang-mills-group/OpenAI-YM-20261028/blob/3456e5e75cf8d4cd558955a5c7638fba2d68f464/YM/Progress/DIRECTED_RECIPE/Tender.lean\#L42-L43}{\texttt{YM.Progress.DirectedRecipe.Payments.hardCost}},
\href{https://github.com/yang-mills-group/OpenAI-YM-20261028/blob/3456e5e75cf8d4cd558955a5c7638fba2d68f464/YM/Progress/DIRECTED_RECIPE/Tender.lean\#L45-L46}{\texttt{YM.Progress.DirectedRecipe.Payments.softCost}}
& forced activities and local exponential corrections \\
Ingredients &
\href{https://github.com/yang-mills-group/OpenAI-YM-20261028/blob/3456e5e75cf8d4cd558955a5c7638fba2d68f464/YM/Progress/DIRECTED_RECIPE/Batch.lean\#L32-L33}{\texttt{YM.Progress.DirectedRecipe.Ingredients}}
& input to the fusion of III.10 \\
slot \(s\), atom \(a\) & \texttt{Slot}, \texttt{Atom} (III.1) &
transported curvature; \(k\)-leg monomial \\
small-field on \(X\) &
\href{https://github.com/yang-mills-group/OpenAI-YM-20261028/blob/3456e5e75cf8d4cd558955a5c7638fba2d68f464/YM/RG/Masks.lean\#L41}{\texttt{YM.RG.Healthy}}
& plaquettes touching \(X\) within chord \(t\) of \(1\) \\
admissible direction &
\href{https://github.com/yang-mills-group/OpenAI-YM-20261028/blob/3456e5e75cf8d4cd558955a5c7638fba2d68f464/YM/RG/Masks.lean\#L141-L142}{\texttt{YM.RG.BudgetPoint.dirBound}}
& \(\lVert A_e\rVert\le t(1+\mathrm{dist}(e,o))^P\) \\
\(\mathrm{slotCost}\), \(\mathrm{margin}\), \(\mathrm{price}\) &
\href{https://github.com/yang-mills-group/OpenAI-YM-20261028/blob/3456e5e75cf8d4cd558955a5c7638fba2d68f464/YM/Next/SCHEDULE/SlotPay.lean\#L95}{\texttt{YM.Next.Scheduling.Wages.slotCost}},
\href{https://github.com/yang-mills-group/OpenAI-YM-20261028/blob/3456e5e75cf8d4cd558955a5c7638fba2d68f464/YM/Next/SCHEDULE/SlotPay.lean\#L94}{\texttt{YM.Next.Scheduling.Wages.margin}},
\href{https://github.com/yang-mills-group/OpenAI-YM-20261028/blob/3456e5e75cf8d4cd558955a5c7638fba2d68f464/YM/Next/SCHEDULE/TensorCost.lean\#L52-L54}{\texttt{YM.Next.Scheduling.Wages.price}}
& derivative constants of III.15.3--4 \\
target \(Q\); \(\rho_Q\), \(\mathrm e_Q\), \(\varepsilon_Q\) &
\href{https://github.com/yang-mills-group/OpenAI-YM-20261028/blob/3456e5e75cf8d4cd558955a5c7638fba2d68f464/YM/AnalyticNext/CONNECTED_ACTIVITY_JETS_E6/OptionCost.lean\#L62-L76}{\texttt{YM.ConnectedActivityE6.Target}}
& connected-activity target (rate, exponent, \ldots) \\
\(\beta_F^{3/8}\) &
\href{https://github.com/yang-mills-group/OpenAI-YM-20261028/blob/3456e5e75cf8d4cd558955a5c7638fba2d68f464/YM/PhysicalNext/RG_NEXT_COMPLETE_REPEATABLE_NONLINEAR_STEP/Height/Highway.lean\#L25}{\texttt{YM.PhysicalNext.Repeat.Height.grow}}
& per-vertex tree weight (III.10) \\
pledge &
\href{https://github.com/yang-mills-group/OpenAI-YM-20261028/blob/3456e5e75cf8d4cd558955a5c7638fba2d68f464/YM/RemainderComparison/Chron/Age.lean\#L94-L95}{\texttt{YM.RemainderComparison.Chron.Pledge}}
& launch conditions (1)--(2) of III.15.5 \\
quay &
\href{https://github.com/yang-mills-group/OpenAI-YM-20261028/blob/3456e5e75cf8d4cd558955a5c7638fba2d68f464/YM/RemainderComparison/Cadenza/Syllabus.lean\#L46}{\texttt{YM.RemainderComparison.Cadenza.Quay}}
& all step conditions available at side \(L\) \\
quarry, financed &
\href{https://github.com/yang-mills-group/OpenAI-YM-20261028/blob/3456e5e75cf8d4cd558955a5c7638fba2d68f464/YM/RemainderComparison/Cadenza/Sheriff.lean\#L21-L27}{\texttt{YM.RemainderComparison.Cadenza.Quarry}},
\href{https://github.com/yang-mills-group/OpenAI-YM-20261028/blob/3456e5e75cf8d4cd558955a5c7638fba2d68f464/YM/RemainderComparison/Cadenza/Sheriff.lean\#L34-L39}{\texttt{YM.RemainderComparison.Cadenza.Quarry.Financed}}
& bounding functions of \(x\) with reserve/rate bounds \\
Pillar &
\href{https://github.com/yang-mills-group/OpenAI-YM-20261028/blob/3456e5e75cf8d4cd558955a5c7638fba2d68f464/YM/RemainderComparison/Chron/Corners.lean\#L23-L28}{\texttt{YM.RemainderComparison.Chron.Pillar}}
& pledge plus per-slot allotments and tilings \\
radius cost &
\href{https://github.com/yang-mills-group/OpenAI-YM-20261028/blob/3456e5e75cf8d4cd558955a5c7638fba2d68f464/YM/RemainderComparison/Lyric/Orbit.lean\#L35-L36}{\texttt{YM.RemainderComparison.Lyric.radiusCost}}
& \((\mathrm{rel}_Q+1)10c_{h+5}/L+2\varepsilon_Q\le v\) \\
comparison estimate &
\href{https://github.com/yang-mills-group/OpenAI-YM-20261028/blob/3456e5e75cf8d4cd558955a5c7638fba2d68f464/YM/RemainderComparison/Herbal/Orbit.lean\#L33-L39}{\texttt{YM.RemainderComparison.Herbal.Airstrip}},
\href{https://github.com/yang-mills-group/OpenAI-YM-20261028/blob/3456e5e75cf8d4cd558955a5c7638fba2d68f464/YM/RemainderComparison/Portal/Wire.lean\#L80-L81}{\texttt{YM.RemainderComparison.Portal.JMatch}}
& interacting vs.~model, with jets matched to order \(r\) \\
heart condition &
\href{https://github.com/yang-mills-group/OpenAI-YM-20261028/blob/3456e5e75cf8d4cd558955a5c7638fba2d68f464/YM/RemainderComparison/Cadenza/Cord.lean\#L50-L59}{\texttt{YM.RemainderComparison.Cadenza.Rays.Heart}}
& support and order inequalities stated before III.15.6 \\
\(B_P,t_P\) &
\href{https://github.com/yang-mills-group/OpenAI-YM-20261028/blob/3456e5e75cf8d4cd558955a5c7638fba2d68f464/YM/RG/Masks.lean\#L125-L136}{\texttt{YM.RG.BudgetPoint}}
& cluster radius and threshold of a parameter point \\
support enlargement &
\href{https://github.com/yang-mills-group/OpenAI-YM-20261028/blob/3456e5e75cf8d4cd558955a5c7638fba2d68f464/YM/PhysicalNext/V9_JOINT_ENDPOINT_LOAD/Heir/Depth.lean\#L44}{\texttt{YM.PhysicalNext.JointLoad.Heir.Walled}}
& \(Y\) lies within distance \(r\mathcal B(x)\) of the coarse image of
\(X\) \\
local analytic smallness &
\href{https://github.com/yang-mills-group/OpenAI-YM-20261028/blob/3456e5e75cf8d4cd558955a5c7638fba2d68f464/YM/PhysicalNext/V9_JOINT_ENDPOINT_LOAD/Heir/Fusion.lean\#L25-L26}{\texttt{YM.PhysicalNext.JointLoad.Heir.Lute.Plented}}
& numerical hypothesis on the return fee \\
library compatibility &
\href{https://github.com/yang-mills-group/OpenAI-YM-20261028/blob/3456e5e75cf8d4cd558955a5c7638fba2d68f464/YM/PhysicalNext/V9_JOINT_ENDPOINT_LOAD/Heir/Sonata.lean\#L66-L75}{\texttt{YM.PhysicalNext.JointLoad.Heir.Plays}}
& registered local estimates for the library \\
marked input &
\href{https://github.com/yang-mills-group/OpenAI-YM-20261028/blob/3456e5e75cf8d4cd558955a5c7638fba2d68f464/YM/RemainderComparison/Halo/Cut.lean\#L26-L29}{\texttt{YM.RemainderComparison.Halo.Presents}}
& derivative bounds on an inherited regular subfamily \\
leaf-scale condition &
\href{https://github.com/yang-mills-group/OpenAI-YM-20261028/blob/3456e5e75cf8d4cd558955a5c7638fba2d68f464/YM/RemainderComparison/Lyric/Flux.lean\#L54-L59}{\texttt{YM.RemainderComparison.Lyric.Scale}}
& nonnegative weights, attached roots, rooted bound and occupancy \\
\(N\); marked cluster &
\href{https://github.com/yang-mills-group/OpenAI-YM-20261028/blob/3456e5e75cf8d4cd558955a5c7638fba2d68f464/YM/RemainderComparison/Halo/Strike.lean\#L39}{\texttt{YM.RemainderComparison.Halo.gross}}
& remaining occurrences after removing leaf insertions \\
insertion count &
\href{https://github.com/yang-mills-group/OpenAI-YM-20261028/blob/3456e5e75cf8d4cd558955a5c7638fba2d68f464/YM/RemainderComparison/Lyric/Plectrum.lean\#L33}{\texttt{YM.RemainderComparison.Lyric.notes}}
& number of inserted leaves, with multiplicity \\
\(A_b(x;j,J)\) &
\href{https://github.com/yang-mills-group/OpenAI-YM-20261028/blob/3456e5e75cf8d4cd558955a5c7638fba2d68f464/YM/AnalyticNext/CONNECTED_ACTIVITY_JETS_E6/Armada.lean\#L25}{\texttt{YM.ConnectedActivityE6.Blaze.March.collars}}
& polynomial spatial factor in Lemma 4.15 \\
\(\pi_\Delta(x)\) &
\href{https://github.com/yang-mills-group/OpenAI-YM-20261028/blob/3456e5e75cf8d4cd558955a5c7638fba2d68f464/YM/RemainderComparison/Sleet/Flint.lean\#L42}{\texttt{YM.RemainderComparison.Sleet.pigment}}
& \(e^{-\Delta\lfloor\mathcal B(x)/100000\rfloor}\) \\
tempered point &
\href{https://github.com/yang-mills-group/OpenAI-YM-20261028/blob/3456e5e75cf8d4cd558955a5c7638fba2d68f464/YM/RemainderComparison/Cadenza/Cord.lean\#L63-L69}{\texttt{YM.RemainderComparison.Cadenza.Temper}}
& agreement of scale, threshold and stage parameters \\
comparison structure &
\href{https://github.com/yang-mills-group/OpenAI-YM-20261028/blob/3456e5e75cf8d4cd558955a5c7638fba2d68f464/YM/RemainderComparison/Portal/Lamps.lean\#L51-L58}{\texttt{YM.RemainderComparison.Portal.Pan}}
& analytic expansion data and coefficient match \\
\(\iota_G\); \(\mathcal P_G\) &
\href{https://github.com/yang-mills-group/OpenAI-YM-20261028/blob/3456e5e75cf8d4cd558955a5c7638fba2d68f464/YM/Gauge/Colors.lean\#L64}{\texttt{YM.Gauge.colorMat}},
\href{https://github.com/yang-mills-group/OpenAI-YM-20261028/blob/3456e5e75cf8d4cd558955a5c7638fba2d68f464/YM/Gauge/Colors.lean\#L96-L97}{\texttt{YM.Gauge.projectColor}}
& colour-to-matrix map; normalized colour projection \\
\(I(x)\); \(\Sigma(x)\) &
\href{https://github.com/yang-mills-group/OpenAI-YM-20261028/blob/3456e5e75cf8d4cd558955a5c7638fba2d68f464/YM/RemainderComparison/Lyric/Forecast.lean\#L31}{\texttt{YM.RemainderComparison.Lyric.ink}},
\href{https://github.com/yang-mills-group/OpenAI-YM-20261028/blob/3456e5e75cf8d4cd558955a5c7638fba2d68f464/YM/RemainderComparison/Lyric/Store.lean\#L62-L63}{\texttt{YM.RemainderComparison.Lyric.salaries}}
& additional root term; tree-sum majorant in Lemma 4.17 \\
\(\mathrm{pledge}(x)\) &
\href{https://github.com/yang-mills-group/OpenAI-YM-20261028/blob/3456e5e75cf8d4cd558955a5c7638fba2d68f464/YM/PhysicalNext/RG_NEXT_COMPLETE_REPEATABLE_NONLINEAR_STEP/Daisy/Corral.lean\#L26}{\texttt{YM.PhysicalNext.Repeat.Daisy.pledgeQ}}
& tail coefficient in Lemma 4.17 \\
marked cluster &
\href{https://github.com/yang-mills-group/OpenAI-YM-20261028/blob/3456e5e75cf8d4cd558955a5c7638fba2d68f464/YM/RemainderComparison/Halo/Strike.lean\#L37}{\texttt{YM.RemainderComparison.Halo.dyed}}
& a remaining occurrence belongs to the inherited family \\
\end{longtable}
}

\subsubsection{7. What remains
unexplained}\label{what-remains-unexplained-17}

\begin{enumerate}
\def\labelenumi{\arabic{enumi}.}
\tightlist
\item
  \textbf{Scope of Part C.} The selected-run theorem retains its
  comparison data as an additional conjunct; the final construction-run
  theorem discards that conjunct. A written proof using only the final
  run properties can omit these comparison conclusions. An all-order
  expansion using them would also have to supply the
  coefficient-matching hypotheses; existence of the comparison functions
  alone does not supply that match.
\item
  \textbf{Local analytic conditions.} The quotation bound and Lemmas
  4.12, 4.16 and 4.19 retain the source's support-enlargement,
  library-compatibility and analytic-smallness conditions. Their full
  numerical definitions, and the local analytic, readiness,
  compatibility, resonance and tolerance conditions in III.15.6, remain
  to be written out. In particular, the local analytic condition
  attached to a funded step has not been translated into an independent
  mathematical statement.
\item
  \textbf{Provider proofs.} The assembly proofs in §4 have been read.
  The deeper proofs of the local singleton estimate, weighted-cluster
  bounds, full-bill estimates and analytic expansion estimates have not
  all been checked here. The contraction and occurrence-count statements
  are consequences of the cited conditional estimates, not unconditional
  bounds for arbitrary local terms.
\item
  \textbf{Historical interpretation.} The estimate \(L^{-5}\) per fine
  site and its conversion to \(L^{-1}\) per coarse block are explicit.
  Identifying it with a particular dimension-five operator is not
  justified by those inequalities alone. The reason for the choice
  \(\tfrac38\) belongs to III.10's tree estimates.
\item
  \textbf{Other fragments.} The construction-run, additional-expense and
  readout record statements have been read. The RG05 readout price
  formulas and their proofs, and the remaining RG12 and R2I fragments,
  are not covered by the detailed verification. The thermodynamic
  fragment also includes eventual numerical estimates and a
  delay-enlargement lemma; it is not solely a collection of definitions.
\end{enumerate}

\subsection{III.16. Analytic toolkit}\label{iii.16.-analytic-toolkit}

\textbf{Part:} III, Ultraviolet construction (renormalization group).
\textbf{Lean:} eight folders, none a numbered chapter of its own:
\href{https://github.com/yang-mills-group/OpenAI-YM-20261028/tree/3456e5e75cf8d4cd558955a5c7638fba2d68f464/YM/Complex}{\texttt{YM/Complex/}},
\href{https://github.com/yang-mills-group/OpenAI-YM-20261028/tree/3456e5e75cf8d4cd558955a5c7638fba2d68f464/YM/TimeKernel}{\texttt{YM/TimeKernel/}},
\href{https://github.com/yang-mills-group/OpenAI-YM-20261028/tree/3456e5e75cf8d4cd558955a5c7638fba2d68f464/YM/Analysis}{\texttt{YM/Analysis/}},
\href{https://github.com/yang-mills-group/OpenAI-YM-20261028/tree/3456e5e75cf8d4cd558955a5c7638fba2d68f464/YM/Laurent}{\texttt{YM/Laurent/}},
\href{https://github.com/yang-mills-group/OpenAI-YM-20261028/tree/3456e5e75cf8d4cd558955a5c7638fba2d68f464/YM/CellFourier}{\texttt{YM/CellFourier/}},
\href{https://github.com/yang-mills-group/OpenAI-YM-20261028/tree/3456e5e75cf8d4cd558955a5c7638fba2d68f464/YM/Work/F07}{\texttt{YM/Work/F07/}},
\href{https://github.com/yang-mills-group/OpenAI-YM-20261028/tree/3456e5e75cf8d4cd558955a5c7638fba2d68f464/YM/Maxwell}{\texttt{YM/Maxwell/}},
\href{https://github.com/yang-mills-group/OpenAI-YM-20261028/tree/3456e5e75cf8d4cd558955a5c7638fba2d68f464/YM/Labels}{\texttt{YM/Labels/}}.
\textbf{Size:} 5,318 physical source lines in 41 files, plus elementary
helper lemmas in other folders.

The results fall into three groups.

\begin{itemize}
\tightlist
\item
  \textbf{(A) Traces of kernel products.} A Hölder inequality and
  Cauchy's estimate give a dimension-free Lipschitz bound for moving
  insertions between time slices. Quadrature passes the bound from
  matrices to kernels; IV.1 uses it for equicontinuity.
\item
  \textbf{(B) Analytic and Fourier estimates.} Holomorphic symbols
  correspond to exponentially summable Laurent coefficients. Fourier and
  Bloch identities diagonalize periodized kernels, and a binomial series
  supplies a matrix square root. Taylor coefficient identities and
  several-variable analyticity support these arguments.
\item
  \textbf{(C) Maxwell fields and labels.} Radial-gauge potentials
  correspond to closed curvature jets. Restriction to the Maxwell shell
  defines the free contractions and Wick hierarchy. Gauge-invariant
  polynomial labels admit coherent complex bases with real coefficients.
\end{itemize}

\subsubsection{1. What this unit proves}\label{what-this-unit-proves-18}

The contraction kernel and Wick hierarchy are definitions whose
properties are proved here.

\emph{Notation for (A).} - Matrices are complex \(\iota\times\iota\)
matrices with \(\iota\) finite, and \(\|\cdot\|\) is the operator norm.
- \(\|b\|_{2q}:=\big(\mathrm{tr}((b^*b)^q)\big)^{1/(2q)}\) is the
Schatten \(2q\)-norm. - For \(A\succeq0\) with spectral decomposition
\(A=\sum_\lambda\lambda P_\lambda\), put
\(A^z:=\sum_{\lambda>0}e^{z\log\lambda}P_\lambda\) for
\(z\in\mathbb C\). This is entire in \(z\). It vanishes on \(\ker A\)
for every \(z\), so \(A^0\) is the projection onto the range, and it
equals the ordinary power \(A^m\) for every integer \(m\ge1\). -
Products \(\prod_{k=1}^n\) are ordered by \(k\).

\textbf{Proposition III.16.1 (integrating a positive-semidefinite
kernel).}

\emph{Hypotheses.} - \(X\) is a compact metric space with a finite Borel
measure \(\mu\), and \(\mathcal A,\mathcal B\) are sets. -
\(K:\mathcal A\times\mathcal A\to\mathbb C\) is positive semidefinite:
for every \(m\) and all \(a_1,\dots,a_m\in\mathcal A\) (repetitions
allowed), the matrix \((K(a_i,a_j))_{i,j}\) is Hermitian and positive
semidefinite. - \(g:\mathcal B\times X\to\mathcal A\) is such that
\((x,y)\mapsto K(g(b,x),g(c,y))\) is continuous on \(X\times X\) for all
\(b,c\in\mathcal B\).

\emph{Conclusion.} The kernel
\[\tilde K(b,c)=\int_X\!\int_X K\big(g(b,x),g(c,y)\big)\,d\mu(y)\,d\mu(x)\]
is positive semidefinite on \(\mathcal B\).

\emph{Status:} theorem. \emph{Lean:}
\href{https://github.com/yang-mills-group/OpenAI-YM-20261028/blob/3456e5e75cf8d4cd558955a5c7638fba2d68f464/YM/Analysis/Projected.lean\#L44-L70}{\texttt{YM.Analysis.NNKernel.accumulate}};
the notion of positive semidefinite kernel is
\href{https://github.com/yang-mills-group/OpenAI-YM-20261028/blob/3456e5e75cf8d4cd558955a5c7638fba2d68f464/YM/Analysis/Nonnegative.lean\#L18-L19}{\texttt{NNKernel}}.

\textbf{Proposition III.16.2 (Hölder inequality for traces with complex
exponents).}

\emph{Hypotheses.} - \(A\succeq0\) and \(b_1,\dots,b_n\) are complex
\(\iota\times\iota\) matrices. - \(q,M\ge1\) are integers, and
\(Z_A:=\mathrm{tr}\,A^M>0\). - \(v\in\mathbb C^n\) satisfies
\(\mathrm{Re}\,v_k\ge0\) for all \(k\) and
\(\frac{n}{2q}+\frac1M\sum_k\mathrm{Re}\,v_k=1\).

\emph{Conclusion.}
\[\Big|\mathrm{tr}\prod_{k=1}^n b_kA^{v_k}\Big|\le\prod_k\|b_k\|_{2q}\cdot Z_A^{\,1-n/(2q)}.\]

For real \(v_k>0\), the corresponding Schatten exponent is \(M/v_k\) and
\(\|A^{v_k}\|_{M/v_k}=Z_A^{v_k/M}\). At \(v_k=0\), the factor is the
range projection, whose operator norm is \(1\) since \(Z_A>0\). The
bound has no dimension factor.

\emph{Status:} theorem. \emph{Lean:}
\href{https://github.com/yang-mills-group/OpenAI-YM-20261028/blob/3456e5e75cf8d4cd558955a5c7638fba2d68f464/YM/Complex/Gaps.lean\#L124-L165}{\texttt{YM.Holder.estimation}}.
The left side is
\href{https://github.com/yang-mills-group/OpenAI-YM-20261028/blob/3456e5e75cf8d4cd558955a5c7638fba2d68f464/YM/Complex/Gaps.lean\#L108-L109}{\texttt{total}}
and the right side is
\href{https://github.com/yang-mills-group/OpenAI-YM-20261028/blob/3456e5e75cf8d4cd558955a5c7638fba2d68f464/YM/Complex/Gaps.lean\#L111-L112}{\texttt{cost}}.

\textbf{Proposition III.16.3 (Lipschitz bound in the gaps, matrices; key
result).}

\emph{Hypotheses.} - \(A\succeq0\) and \(b_1,\dots,b_n\) are as in
III.16.2. - \(q,w\ge1\) are integers, \(M=2qw\), and
\(\mathrm{tr}\,A^M>0\). - \(Y\in\mathbb N^n\) satisfies
\(nw+\sum_kY_k=M\). - For some \(i\ne j\) and some \(\ell\in\mathbb N\),
the vector \(Y'\in\mathbb N^n\) moves \(\ell\) units from gap \(j\) to
gap \(i\): \(Y'_i=Y_i+\ell\), \(Y'_j+\ell=Y_j\), and \(Y'_k=Y_k\)
otherwise. - \(r>0\), \(Y_k>2r\) for every \(k\), and \(\ell\le r\).

\emph{Conclusion.}
\[\Big|\mathrm{tr}\prod_k b_kA^{Y'_k}-\mathrm{tr}\prod_k b_kA^{Y_k}\Big|\le\frac{\ell}{r}\,\prod_k\|b_k\|_{2q}\,\big(\mathrm{tr}\,A^M\big)^{1-n/(2q)}.\]

The condition \(i\ne j\) forces \(n\ge2\).

\emph{Status:} theorem. \emph{Lean:}
\href{https://github.com/yang-mills-group/OpenAI-YM-20261028/blob/3456e5e75cf8d4cd558955a5c7638fba2d68f464/YM/Complex/Gaps.lean\#L170-L229}{\texttt{YM.Holder.finite\_lip}}.

\textbf{Proposition III.16.4 (Lipschitz bound in the gaps, kernels; key
result).}

\emph{Setting.} \(X\) is a compact metric space with a finite Borel
measure \(\mu\). For \(L\ge1\) and continuous
\(k_1,\dots,k_L:X\times X\to\mathbb C\), write
\[\mathrm{Tr}(k_1\cdots k_L)=\int_{X^L}\prod_{l=1}^Lk_l(x_l,x_{l+1})\,d\mu(x_1)\cdots d\mu(x_L),\qquad x_{L+1}=x_1,\]
and write \(k^*(x,y)=\overline{k(y,x)}\) for the adjoint kernel. Powers
inside \(\mathrm{Tr}\) denote repetitions in the cyclic word.

\emph{Hypotheses.} - \(K\) is a continuous positive-semidefinite kernel,
and \(B_1,\dots,B_n\) are continuous kernels. - \(q,w\ge1\), \(M=2qw\),
and \(\mathrm{Tr}(K^M)=Z\) with \(Z>0\) real. - \(Y,Y',\ell,r\) are as
in III.16.3. - \(c_1,\dots,c_n\) are real numbers with
\(\big|\mathrm{Tr}\big((B_k^*B_k)^q\big)\big|^{1/(2q)}\le c_k\,Z^{1/(2q)}\)
for every \(k\).

\emph{Conclusion.}
\[\big|\mathrm{Tr}(B_1K^{Y'_1}\cdots B_nK^{Y'_n})-\mathrm{Tr}(B_1K^{Y_1}\cdots B_nK^{Y_n})\big|\le Z\Big(\prod_kc_k\Big)\frac{\ell}{r}.\]

\emph{Status:} theorem. \emph{Lean:}
\href{https://github.com/yang-mills-group/OpenAI-YM-20261028/blob/3456e5e75cf8d4cd558955a5c7638fba2d68f464/YM/TimeKernel/Unit.lean\#L76-L106}{\texttt{YM.TimeKernel.passage}},
with
\href{https://github.com/yang-mills-group/OpenAI-YM-20261028/blob/3456e5e75cf8d4cd558955a5c7638fba2d68f464/YM/TimeKernel/Trace.lean\#L76}{\texttt{closed}}
(the cyclic integral),
\href{https://github.com/yang-mills-group/OpenAI-YM-20261028/blob/3456e5e75cf8d4cd558955a5c7638fba2d68f464/YM/TimeKernel/Unit.lean\#L17-L18}{\texttt{along}}
(the word with gaps) and
\href{https://github.com/yang-mills-group/OpenAI-YM-20261028/blob/3456e5e75cf8d4cd558955a5c7638fba2d68f464/YM/TimeKernel/Unit.lean\#L20}{\texttt{multiply}}
(the column word \((B^*B)^q\)).

\emph{Notation for (B).} - \(\alpha\) is a finite index set; in the
applications it is the four axes. - \(|x|_1=\sum_a|x_a|\) for
\(x\in\mathbb Z^\alpha\), and \(e^{ip\cdot x}=\exp(i\sum_ap_ax_a)\) for
\(p\in\mathbb C^\alpha\). -
\(S_\eta=\{p\in\mathbb C^\alpha:|\mathrm{Im}\,p_a|<\eta\ \forall a\}\).
- For \(f:\mathbb Z^\alpha\to\mathbb C\), the symbol is
\(\hat f(p)=\sum_xf(x)e^{ip\cdot x}\).

\textbf{Proposition III.16.5 (holomorphy in a strip and exponentially
weighted Laurent coefficients).} Let \(\eta>0\).

\begin{enumerate}
\def\labelenumi{(\alph{enumi})}
\tightlist
\item
  Let \(u\) be holomorphic on \(S_\eta\) and \(2\pi\)-periodic there:
  \(u(p+2\pi x)=u(p)\) for \(p\in S_\eta\) and \(x\in\mathbb Z^\alpha\).
  For \(p\in S_\eta\) let
  \(\hat u_p(x)=\int_{[0,1]^\alpha}u(p+2\pi\theta)e^{-2\pi i\theta\cdot x}d\theta\),
  and put \(\hat u=\hat u_0\). Then:
\end{enumerate}

\begin{itemize}
\tightlist
\item
  the integral \(\int u(p+2\pi\theta)\,d\theta\) does not depend on
  \(p\in S_\eta\);
\item
  \(\hat u_p(x)=e^{ip\cdot x}\hat u(x)\).
\end{itemize}

If moreover \(|u|\le K\) on \(S_\eta\), then: - \(|\hat u_p(x)|\le K\);
- for \(H\ge0\) and \(b>0\) with \(H+b<\eta\),
\(\sum_xe^{H|x|_1}|\hat u(x)|\le K\big(2(1+b)/b\big)^{\mathrm{card}\,\alpha}\);
- \(u(p)=\sum_x\hat u(x)e^{ip\cdot x}\) on \(S_\eta\), with absolute
convergence.

\begin{enumerate}
\def\labelenumi{(\alph{enumi})}
\setcounter{enumi}{1}
\tightlist
\item
  Let \(\sum_xe^{\eta|x|_1}|f(x)|\le W\). Then:
\end{enumerate}

\begin{itemize}
\tightlist
\item
  \(|\hat f(p)|\le W\) whenever \(|\mathrm{Im}\,p_a|\le\eta\) for all
  \(a\);
\item
  for \(0<\sigma<\eta\), \(\hat f\) is holomorphic on \(S_\sigma\) with
  \(\|D\hat f(p)\|\le W/(\eta-\sigma)\) (operator norm for the sup norm
  on \(\mathbb C^\alpha\)), so
  \(|\hat f(p)-\hat f(p')|\le W\|p-p'\|_\infty/(\eta-\sigma)\) on
  \(S_\sigma\);
\item
  for real \(p\), the symbol of \(\mathrm{Re}\,f\) equals
  \(\tfrac12\big(\hat f(p)+\overline{\hat f(-p)}\big)\).
\end{itemize}

\emph{Status:} theorems. \emph{Lean:} - for (a):
\href{https://github.com/yang-mills-group/OpenAI-YM-20261028/blob/3456e5e75cf8d4cd558955a5c7638fba2d68f464/YM/Laurent/Rune.lean\#L91-L99}{\texttt{market\_cap}},
\href{https://github.com/yang-mills-group/OpenAI-YM-20261028/blob/3456e5e75cf8d4cd558955a5c7638fba2d68f464/YM/Laurent/Turn.lean\#L103-L112}{\texttt{full\_turn}},
\href{https://github.com/yang-mills-group/OpenAI-YM-20261028/blob/3456e5e75cf8d4cd558955a5c7638fba2d68f464/YM/Laurent/Turn.lean\#L114-L118}{\texttt{coat\_price}},
\href{https://github.com/yang-mills-group/OpenAI-YM-20261028/blob/3456e5e75cf8d4cd558955a5c7638fba2d68f464/YM/Laurent/Art.lean\#L61-L87}{\texttt{Vane.issues}},
\href{https://github.com/yang-mills-group/OpenAI-YM-20261028/blob/3456e5e75cf8d4cd558955a5c7638fba2d68f464/YM/Laurent/Art.lean\#L90-L104}{\texttt{Vane.paints}};
- for (b):
\href{https://github.com/yang-mills-group/OpenAI-YM-20261028/blob/3456e5e75cf8d4cd558955a5c7638fba2d68f464/YM/Laurent/Eval.lean\#L131-L133}{\texttt{Width.estimate}},
\href{https://github.com/yang-mills-group/OpenAI-YM-20261028/blob/3456e5e75cf8d4cd558955a5c7638fba2d68f464/YM/Work/F07/Read.lean\#L55-L95}{\texttt{readDeriv}},
\href{https://github.com/yang-mills-group/OpenAI-YM-20261028/blob/3456e5e75cf8d4cd558955a5c7638fba2d68f464/YM/Work/F07/Read.lean\#L97-L101}{\texttt{readDiff}},
\href{https://github.com/yang-mills-group/OpenAI-YM-20261028/blob/3456e5e75cf8d4cd558955a5c7638fba2d68f464/YM/Work/F07/Read.lean\#L119-L133}{\texttt{readRealPart}}.

\textbf{Proposition III.16.6 (Fourier and Bloch decomposition on tori).}

\emph{Setting.} - All torus side lengths \(n,c,f,L\) are positive
integers, and \(\Lambda_n=(\mathbb Z/n\mathbb Z)^4\). For
\(p\in\Lambda_n\), represented in \(\{0,\dots,n-1\}^4\), let
\(k_p=2\pi p/n\), and let \(F_n=(e^{ik_p\cdot x})_{x,p\in\Lambda_n}\). -
\(\alpha,\beta\) are finite index sets, and the real kernel
\(R=(R_{ab}(x))\), \(x\in\mathbb Z^4\), satisfies
\(\sum_x|R_{ab}(x)|<\infty\) for every \(a\in\alpha\), \(b\in\beta\). -
Its periodization on \(\Lambda_n\) is
\(R_n\big((y,a),(y',b)\big)=\sum_{x:\,y+x\equiv y'}R_{ab}(x)\), and its
symbol is \(\hat R(k)_{ab}=\sum_xR_{ab}(x)e^{ik\cdot x}\).

\begin{enumerate}
\def\labelenumi{(\alph{enumi})}
\item
  \(F_nF_n^*=F_n^*F_n=n^4I\), and
  \(R_n(F_n\otimes I_\beta)=(F_n\otimes I_\alpha)\,\mathrm{diag}_{p\in\Lambda_n}\hat R(k_p)\).
  This holds for every \(p\), including \(p=0\).
\item
  Let \(f=cL\). For \(p\in\Lambda_c\) and \(l\in\Lambda_L\) let
  \(k_{p,l}=2\pi(p+cl)/f\). These are the \(L^4\) fine momenta \(k\)
  with \(Lk\equiv2\pi p/c\) modulo \(2\pi\), i.e.~the aliases of the
  coarse momentum \(2\pi p/c\). The fine Fourier matrix
  \(\Phi=(e^{ik_{p,l}\cdot y})_{y\in\Lambda_f,(p,l)}\) satisfies:
\end{enumerate}

\begin{itemize}
\tightlist
\item
  \(\Phi\Phi^*=f^4I\);
\item
  \(R_f(\Phi\otimes I_\beta)=(\Phi\otimes I_\alpha)\,\mathrm{diag}_{p\in\Lambda_c}\,\mathrm{diag}_{l\in\Lambda_L}\hat R(k_{p,l})\).
\end{itemize}

\begin{enumerate}
\def\labelenumi{(\alph{enumi})}
\setcounter{enumi}{2}
\tightlist
\item
  For every \(p\in\mathbb C^4\), the alias block
  \(E_L(p)=(e^{i(p+2\pi l)\cdot x/L})_{x,l\in\Lambda_L}\) satisfies
  \(E_L(\bar p)^*E_L(p)=L^4I\).
\end{enumerate}

\emph{Status:} theorems. \emph{Lean:}
\href{https://github.com/yang-mills-group/OpenAI-YM-20261028/blob/3456e5e75cf8d4cd558955a5c7638fba2d68f464/YM/CellFourier/Mesh.lean\#L117-L119}{\texttt{gram}},
\href{https://github.com/yang-mills-group/OpenAI-YM-20261028/blob/3456e5e75cf8d4cd558955a5c7638fba2d68f464/YM/CellFourier/Mesh.lean\#L121-L127}{\texttt{coGram}},
\href{https://github.com/yang-mills-group/OpenAI-YM-20261028/blob/3456e5e75cf8d4cd558955a5c7638fba2d68f464/YM/CellFourier/Circulant.lean\#L88-L104}{\texttt{circulant}},
\href{https://github.com/yang-mills-group/OpenAI-YM-20261028/blob/3456e5e75cf8d4cd558955a5c7638fba2d68f464/YM/CellFourier/Shuttle.lean\#L66-L75}{\texttt{readingsGram}},
\href{https://github.com/yang-mills-group/OpenAI-YM-20261028/blob/3456e5e75cf8d4cd558955a5c7638fba2d68f464/YM/CellFourier/Shuttle.lean\#L110-L129}{\texttt{aligns}},
\href{https://github.com/yang-mills-group/OpenAI-YM-20261028/blob/3456e5e75cf8d4cd558955a5c7638fba2d68f464/YM/CellFourier/Alias.lean\#L41-L53}{\texttt{straightens}}.

\textbf{Proposition III.16.7 (analytic square root near a spectral
band).} Let \(m>0\) and let the finite matrix index set be nonempty. For
complex matrices \(X\) with \(\|X/m-1\|<1\), put
\(\Phi_m(X):=\sqrt m\sum_{k\ge0}\binom{1/2}{k}(X/m-1)^k\). - \(\Phi_m\)
is holomorphic on this ball. - \(\Phi_m\) is bounded on each smaller
ball \(\|X/m-1\|\le r<1\). - \(\Phi_m\) commutes with entrywise complex
conjugation. - If \(X\) is Hermitian and \(v\le X\le m\) as quadratic
forms, with \(v>0\), then \(\Phi_m(X)\) is Hermitian, \(\Phi_m(X)^2=X\),
and \(\sqrt v\le\Phi_m(X)\le\sqrt m\).

\emph{Status:} theorem. \emph{Lean:}
\href{https://github.com/yang-mills-group/OpenAI-YM-20261028/blob/3456e5e75cf8d4cd558955a5c7638fba2d68f464/YM/Work/F07/MatrixRoot.lean\#L185-L198}{\texttt{binRoot\_on\_band}},
\href{https://github.com/yang-mills-group/OpenAI-YM-20261028/blob/3456e5e75cf8d4cd558955a5c7638fba2d68f464/YM/Work/F07/MatrixRoot.lean\#L127-L130}{\texttt{binRoot\_deriv}},
\href{https://github.com/yang-mills-group/OpenAI-YM-20261028/blob/3456e5e75cf8d4cd558955a5c7638fba2d68f464/YM/Work/F07/MatrixRoot.lean\#L116-L123}{\texttt{binSum\_bound}},
\href{https://github.com/yang-mills-group/OpenAI-YM-20261028/blob/3456e5e75cf8d4cd558955a5c7638fba2d68f464/YM/Work/F07/MatrixRoot.lean\#L138-L143}{\texttt{binRoot\_conj}}.

\emph{Notation for (C).} - Fix a compact matrix gauge group \(G\). Its
colour space \(\mathfrak g\) is the real Lie algebra in coordinates
orthonormal for
\(-(\mathrm{dim}\,\rho)^{-1}\mathrm{Re}\,\mathrm{tr}(XY)\), where
\(\rho\) is the fixed defining representation. - A homogeneous
polynomial on spacetime is identified with its symmetric multilinear
Taylor coefficient. - For \(j\ge0\): - \(\mathcal A_j\) is the space of
\(4\)-tuples \((A_l)\) of \(\mathfrak g\)-valued homogeneous polynomials
of degree \(j+1\) on \(\mathbb R^4\) with \(\sum_lx_lA_l(x)=0\) (radial,
or Fock--Schwinger, gauge); - \(\mathcal C_j\) is the space of
\(\mathfrak g\)-valued homogeneous polynomials \(F_{il}\) of degree
\(j\) with \(F_{il}=-F_{li}\) and
\(\partial_aF_{bc}+\partial_bF_{ca}+\partial_cF_{ab}=0\); -
\(\mathcal S_j\subset\mathcal C_j\), the Maxwell shell, adds the
condition \(\sum_a\partial_aF_{ab}=0\).

\textbf{Proposition III.16.8 (free Maxwell field).}

\begin{enumerate}
\def\labelenumi{(\alph{enumi})}
\item
  \emph{(Poincaré lemma in radial gauge.)} The curl
  \(A\mapsto F_{il}=\partial_iA_l-\partial_lA_i\) is a linear
  isomorphism \(\mathcal A_j\to\mathcal C_j\). Its inverse is
  \(A_i(x)=(j+2)^{-1}\sum_kx_kF_{ki}(x)\).
\item
  \emph{(Taylor coefficients of solutions.)} Let \(U\subset\mathbb R^4\)
  be open, and let the \(\mathfrak g\)-valued field \(f=(f_{ab})\) be
  smooth and antisymmetric on \(U\) and satisfy the Bianchi and Maxwell
  equations there. Then for every \(x\in U\) and \(j\ge0\), the \(j\)-th
  Taylor coefficient of \(f\) at \(x\) lies in \(\mathcal S_j\).
\item
  \emph{(Free two-point function.)} The function
  \(G_0(x)=1/(4\pi^2|x|^2)\) is harmonic on
  \(\mathbb R^4\setminus\{0\}\). Put
  \[C_{ab,km}=\partial_a\psi_{b,km}-\partial_b\psi_{a,km},\qquad\psi_{b,km}=\delta_{kb}\partial_mG_0-\delta_{mb}\partial_kG_0.\]
  This is the field-strength two-point function
  \(\langle F_{ab}(x)F_{km}(0)\rangle\) of one colour component of the
  free field in Feynman gauge. Then:
\end{enumerate}

\begin{itemize}
\tightlist
\item
  for fixed \((k,m)\), \((a,b)\mapsto C_{ab,km}\) is smooth and
  antisymmetric on \(x\ne0\) and solves the Bianchi and Maxwell
  equations there;
\item
  the same holds in \((k,m)\) for fixed \((a,b)\);
\item
  \(C_{ab,km}(x)=C_{km,ab}(x)=C_{ab,km}(-x)\).
\end{itemize}

\begin{enumerate}
\def\labelenumi{(\alph{enumi})}
\setcounter{enumi}{3}
\tightlist
\item
  \emph{(One contraction between jets.)} Let \(j,l\ge0\), and let \(p\)
  and \(q\) be real linear forms on \(\mathcal S_j\) and
  \(\mathcal S_l\) (probes). Write \(c_jh(z)=D^jh(z)/j!\), and let
  \((e_c)\) be the orthonormal colour basis. Choose linear extensions
  \(P,Q\) of \(p,q\) to the unrestricted two-index multilinear tensors.
  Define the colour-valued field
  \[R_{P,km}(z)=\sum_cP\big((c_j(C_{ab,km}e_c)(z))_{a,b}\big)e_c,\] and,
  for \(z\ne0\), put
  \[\kappa(p,q)(z)=(-1)^lQ\big((c_lR_{P,km}(z))_{k,m}\big).\] This
  definition is independent of both extensions. It is bilinear in
  \((p,q)\), smooth for \(z\ne0\), and satisfies
  \(\kappa(p,q)(z)=\kappa(q,p)(-z)\).
\end{enumerate}

For the coordinate probes \(p(F)=F_{ab}[u]_c\) and \(q(F)=F_{km}[v]_d\),
with \(u\in(\mathbb R^4)^j\) and \(v\in(\mathbb R^4)^l\),
\[\kappa(p,q)(z)=\delta_{cd}(-1)^l c_l\big(t\mapsto c_jC_{ab,km}(t)[u]\big)(z)[v].\]
Equivalently, at \(z=x-y\) this is the mixed Taylor coefficient at
\((x,y)\) of \(\delta_{cd}C_{ab,km}(x-y)\). Thus \(\kappa\) is the free
covariance of the corresponding jet observables, with Taylor
coefficients used at both endpoints.

\emph{Status:} theorems in (a)--(c); definition and theorems in (d).
\emph{Lean:} - (a):
\href{https://github.com/yang-mills-group/OpenAI-YM-20261028/blob/3456e5e75cf8d4cd558955a5c7638fba2d68f464/YM/Maxwell/Exterior.lean\#L218-L224}{\texttt{potential}},
\href{https://github.com/yang-mills-group/OpenAI-YM-20261028/blob/3456e5e75cf8d4cd558955a5c7638fba2d68f464/YM/Maxwell/Exterior.lean\#L108-L128}{\texttt{closedCurl}},
\href{https://github.com/yang-mills-group/OpenAI-YM-20261028/blob/3456e5e75cf8d4cd558955a5c7638fba2d68f464/YM/Maxwell/Exterior.lean\#L163-L195}{\texttt{curl\_inverse}},
\href{https://github.com/yang-mills-group/OpenAI-YM-20261028/blob/3456e5e75cf8d4cd558955a5c7638fba2d68f464/YM/Maxwell/Exterior.lean\#L197-L214}{\texttt{potential\_recover}};
- (b):
\href{https://github.com/yang-mills-group/OpenAI-YM-20261028/blob/3456e5e75cf8d4cd558955a5c7638fba2d68f464/YM/Maxwell/Shell.lean\#L73-L122}{\texttt{WaveOn.jets}};
- (c):
\href{https://github.com/yang-mills-group/OpenAI-YM-20261028/blob/3456e5e75cf8d4cd558955a5c7638fba2d68f464/YM/Maxwell/PotentialKernel.lean\#L69-L77}{\texttt{harmonic}},
\href{https://github.com/yang-mills-group/OpenAI-YM-20261028/blob/3456e5e75cf8d4cd558955a5c7638fba2d68f464/YM/Maxwell/Curls.lean\#L82-L98}{\texttt{Bare.equation}},
\href{https://github.com/yang-mills-group/OpenAI-YM-20261028/blob/3456e5e75cf8d4cd558955a5c7638fba2d68f464/YM/Maxwell/Curls.lean\#L111-L123}{\texttt{Bare.equation\_right}},
\href{https://github.com/yang-mills-group/OpenAI-YM-20261028/blob/3456e5e75cf8d4cd558955a5c7638fba2d68f464/YM/Maxwell/Curls.lean\#L104-L109}{\texttt{Bare.exchange}},
\href{https://github.com/yang-mills-group/OpenAI-YM-20261028/blob/3456e5e75cf8d4cd558955a5c7638fba2d68f464/YM/Maxwell/Curls.lean\#L100-L102}{\texttt{Bare.even}};
- (d):
\href{https://github.com/yang-mills-group/OpenAI-YM-20261028/blob/3456e5e75cf8d4cd558955a5c7638fba2d68f464/YM/Maxwell/Leg.lean\#L161-L162}{\texttt{legKernel}},
\href{https://github.com/yang-mills-group/OpenAI-YM-20261028/blob/3456e5e75cf8d4cd558955a5c7638fba2d68f464/YM/Maxwell/Leg.lean\#L148-L159}{\texttt{descends}},
\href{https://github.com/yang-mills-group/OpenAI-YM-20261028/blob/3456e5e75cf8d4cd558955a5c7638fba2d68f464/YM/Maxwell/Bilinear.lean\#L24-L47}{\texttt{pairing}},
\href{https://github.com/yang-mills-group/OpenAI-YM-20261028/blob/3456e5e75cf8d4cd558955a5c7638fba2d68f464/YM/Maxwell/Bilinear.lean\#L159-L169}{\texttt{exchange}},
\href{https://github.com/yang-mills-group/OpenAI-YM-20261028/blob/3456e5e75cf8d4cd558955a5c7638fba2d68f464/YM/Maxwell/Bilinear.lean\#L72-L87}{\texttt{raw\_mixed}},
\href{https://github.com/yang-mills-group/OpenAI-YM-20261028/blob/3456e5e75cf8d4cd558955a5c7638fba2d68f464/YM/Maxwell/Bilinear.lean\#L90-L105}{\texttt{reads\_mixed}}.

\textbf{Proposition III.16.9 (field labels).}

\emph{Setting.} Let \(d\ge0\) be an integer.

\begin{itemize}
\tightlist
\item
  Fix real bases of the spaces \(\mathcal A_j\).
\item
  \(\mathrm{Pol}\) is the algebra of complex polynomials in the
  resulting coordinates of a jet sequence \((A_j)_{j\ge0}\).
\item
  A coordinate from \(\mathcal A_j\) has weight \(j+2\), and a monomial
  has the sum of the weights of its factors.
\item
  \(\mathrm{Label}_d\subset\mathrm{Pol}\) consists of the polynomials of
  weight \(\le d\) that are invariant under the simultaneous adjoint
  action of \(G\) on all \(A_j\). Invariance is imposed on the whole
  polynomial.
\end{itemize}

\begin{enumerate}
\def\labelenumi{(\alph{enumi})}
\item
  A polynomial that vanishes on all real jet sequences is zero.
\item
  There are finitely many monomials of weight \(\le d\), so
  \(\mathrm{Label}_d\) is finite-dimensional.
\item
  Every element of \(\mathrm{Label}_d\) is the sum of its bihomogeneous
  components (weight, polynomial degree), both indices \(\le d\).
\end{enumerate}

\begin{itemize}
\tightlist
\item
  Each component is again a label.
\item
  Taking components commutes with complex conjugation of coefficients
  and with the action of \(O(4)\) on jets.
\item
  For \(s\in O(4)\), set \((T_sA_j)_i(x)=\sum_ks_{ik}(A_j)_k(s^{-1}x)\).
  The induced label action is \((s\cdot P)(A)=P(T_{s^{-1}}A)\); it
  satisfies \((st)\cdot P=s\cdot(t\cdot P)\).
\end{itemize}

\begin{enumerate}
\def\labelenumi{(\alph{enumi})}
\setcounter{enumi}{3}
\tightlist
\item
  For every integer \(d\ge0\), \(\mathrm{Label}_d\) has a complex basis
  (a frame) \(e_0=1,e_1,\dots,e_{N_d}\) of bihomogeneous labels with
  real coefficients. These bases can be chosen simultaneously so that,
  under the natural inclusion, every basis vector at level \(d\) occurs
  at level \(d+1\). No ordering beyond \(e_0=1\) is prescribed.
\end{enumerate}

\emph{Status:} theorems. \emph{Lean:}
\href{https://github.com/yang-mills-group/OpenAI-YM-20261028/blob/3456e5e75cf8d4cd558955a5c7638fba2d68f464/YM/Labels/Polynomial.lean\#L38-L49}{\texttt{eval\_ext}},
\href{https://github.com/yang-mills-group/OpenAI-YM-20261028/blob/3456e5e75cf8d4cd558955a5c7638fba2d68f464/YM/Labels/Polynomial.lean\#L63-L94}{\texttt{cap\_finite}},
\href{https://github.com/yang-mills-group/OpenAI-YM-20261028/blob/3456e5e75cf8d4cd558955a5c7638fba2d68f464/YM/Labels/Polynomial.lean\#L114-L117}{\texttt{labels\_finite}},
\href{https://github.com/yang-mills-group/OpenAI-YM-20261028/blob/3456e5e75cf8d4cd558955a5c7638fba2d68f464/YM/Labels/Grading.lean\#L87-L105}{\texttt{decompose}},
\href{https://github.com/yang-mills-group/OpenAI-YM-20261028/blob/3456e5e75cf8d4cd558955a5c7638fba2d68f464/YM/Labels/Grading.lean\#L66-L71}{\texttt{gradePart\_bar}},
\href{https://github.com/yang-mills-group/OpenAI-YM-20261028/blob/3456e5e75cf8d4cd558955a5c7638fba2d68f464/YM/Labels/Grading.lean\#L73-L74}{\texttt{gradePart\_turn}},
\href{https://github.com/yang-mills-group/OpenAI-YM-20261028/blob/3456e5e75cf8d4cd558955a5c7638fba2d68f464/YM/Labels/Action.lean\#L82-L88}{\texttt{rotation\_law}},
\href{https://github.com/yang-mills-group/OpenAI-YM-20261028/blob/3456e5e75cf8d4cd558955a5c7638fba2d68f464/YM/Labels/Basis.lean\#L50-L57}{\texttt{Frame}},
\href{https://github.com/yang-mills-group/OpenAI-YM-20261028/blob/3456e5e75cf8d4cd558955a5c7638fba2d68f464/YM/Labels/Basis.lean\#L107-L111}{\texttt{frames}}.

\textbf{Proposition III.16.10 (free Wick hierarchy on labels).}

\emph{Definitions.} Let \(n\ge0\), let \(d_1,\dots,d_n\ge0\) be
integers, and let \(x_1,\dots,x_n\in\mathbb R^4\) be pairwise distinct.
- For Wick monomials in the coordinates of shell jets, the \emph{free
moment} is the sum, over all pairings of the factor occurrences in which
partners sit at different points, of the products of the kernels
\(\kappa(\cdot,\cdot)(x_i-x_{i'})\) of the paired coordinate forms. It
is extended multilinearly to Wick polynomials. - The \emph{free
correlation} of labels \(P_k\in\mathrm{Label}_{d_k}\) is the moment of
their restrictions to on-shell jets, \(A_j=\mathrm{curl}^{-1}F_j\) with
\(F_j\in\mathcal S_j\).

\emph{Properties.} - The free correlation is smooth on the configuration
space. - It vanishes if one label vanishes on all on-shell jets,
i.e.~lies in the equation ideal. - It is invariant under permuting the
points together with the labels. More generally, embedding the index set
into a larger configuration of distinct points preserves the correlation
when each added point carries the label \(1\). - The moment with every
insertion equal to \(1\) is \(1\), including at \(n=0\). For constant
insertions \(a_1,\dots,a_n\in\mathbb C\), it is \(\prod_i a_i\).

\emph{Status:} definitions and theorems. \emph{Lean:}
\href{https://github.com/yang-mills-group/OpenAI-YM-20261028/blob/3456e5e75cf8d4cd558955a5c7638fba2d68f464/YM/Maxwell/Wick.lean\#L126-L127}{\texttt{Wick.moment}},
\href{https://github.com/yang-mills-group/OpenAI-YM-20261028/blob/3456e5e75cf8d4cd558955a5c7638fba2d68f464/YM/Maxwell/Wick.lean\#L145-L147}{\texttt{Wick.correlation}},
\href{https://github.com/yang-mills-group/OpenAI-YM-20261028/blob/3456e5e75cf8d4cd558955a5c7638fba2d68f464/YM/Maxwell/Wick.lean\#L152-L154}{\texttt{Wick.smooth}},
\href{https://github.com/yang-mills-group/OpenAI-YM-20261028/blob/3456e5e75cf8d4cd558955a5c7638fba2d68f464/YM/Maxwell/Wick.lean\#L149-L150}{\texttt{Wick.kills\_equation}},
\href{https://github.com/yang-mills-group/OpenAI-YM-20261028/blob/3456e5e75cf8d4cd558955a5c7638fba2d68f464/YM/Maxwell/Equation.lean\#L74-L79}{\texttt{ideal\_meaning}},
\href{https://github.com/yang-mills-group/OpenAI-YM-20261028/blob/3456e5e75cf8d4cd558955a5c7638fba2d68f464/YM/Maxwell/Reorder.lean\#L112-L124}{\texttt{Wick.exchangeSites}},
\href{https://github.com/yang-mills-group/OpenAI-YM-20261028/blob/3456e5e75cf8d4cd558955a5c7638fba2d68f464/YM/Maxwell/Reorder.lean\#L103-L110}{\texttt{Wick.monExtension}},
\href{https://github.com/yang-mills-group/OpenAI-YM-20261028/blob/3456e5e75cf8d4cd558955a5c7638fba2d68f464/YM/Maxwell/Reorder.lean\#L126-L144}{\texttt{Wick.adjoin}},
\href{https://github.com/yang-mills-group/OpenAI-YM-20261028/blob/3456e5e75cf8d4cd558955a5c7638fba2d68f464/YM/Maxwell/Reorder.lean\#L146-L153}{\texttt{Wick.constants}}.

Other units also quote, by name, the general lemmas of §4: the algebra
of positive kernels (4.1), quadrature (4.2), the calculus of Taylor
coefficients (4.19), and analyticity in several variables and weak
analyticity (4.20).

\subsubsection{2. What it uses}\label{what-it-uses-18}

\begin{itemize}
\tightlist
\item
  \textbf{I.1 (lattice, symmetries, gauge groups).}

  \begin{itemize}
  \tightlist
  \item
    The torus \(\Lambda_n\), its cover \(\mathbb Z^4\), and the
    coarse/fine subdivision maps:
    \href{https://github.com/yang-mills-group/OpenAI-YM-20261028/blob/3456e5e75cf8d4cd558955a5c7638fba2d68f464/YM/Lattice/Site.lean\#L19}{\texttt{Lattice.Site}},
    \href{https://github.com/yang-mills-group/OpenAI-YM-20261028/blob/3456e5e75cf8d4cd558955a5c7638fba2d68f464/YM/Lattice/Site.lean\#L35-L38}{\texttt{Lattice.wrap}},
    \href{https://github.com/yang-mills-group/OpenAI-YM-20261028/blob/3456e5e75cf8d4cd558955a5c7638fba2d68f464/YM/Lattice/Subdivision.lean\#L22-L24}{\texttt{Lattice.Ratio}}.
  \item
    \(\mathbb R^4\) with its axes, and \(O(4)\):
    \href{https://github.com/yang-mills-group/OpenAI-YM-20261028/blob/3456e5e75cf8d4cd558955a5c7638fba2d68f464/YM/Space/Configurations.lean\#L24}{\texttt{E4}},
    \href{https://github.com/yang-mills-group/OpenAI-YM-20261028/blob/3456e5e75cf8d4cd558955a5c7638fba2d68f464/YM/Space/Orthogonal.lean\#L20}{\texttt{Ortho}}.
  \item
    The gauge-group data, namely the colour space and the adjoint
    action:
    \href{https://github.com/yang-mills-group/OpenAI-YM-20261028/blob/3456e5e75cf8d4cd558955a5c7638fba2d68f464/YM/Gauge/MatrixGroup.lean\#L23-L26}{\texttt{Gauge}},
    \href{https://github.com/yang-mills-group/OpenAI-YM-20261028/blob/3456e5e75cf8d4cd558955a5c7638fba2d68f464/YM/Gauge/Colors.lean\#L43}{\texttt{Gauge.Color}},
    \href{https://github.com/yang-mills-group/OpenAI-YM-20261028/blob/3456e5e75cf8d4cd558955a5c7638fba2d68f464/YM/Gauge/Colors.lean\#L114-L117}{\texttt{Gauge.adjoint}}.
  \item
    The configuration space of distinct points:
    \href{https://github.com/yang-mills-group/OpenAI-YM-20261028/blob/3456e5e75cf8d4cd558955a5c7638fba2d68f464/YM/Space/Configurations.lean\#L63}{\texttt{conf}}.
  \item
    All of these are definitions.
  \end{itemize}
\item
  \textbf{III.1:} the summability condition ``nonnegative, summable, sum
  \(\le W\)'',
  \href{https://github.com/yang-mills-group/OpenAI-YM-20261028/blob/3456e5e75cf8d4cd558955a5c7638fba2d68f464/YM/RG/Activities.lean\#L59-L60}{\texttt{RG.SumBudget}},
  used to define the weighted norms of (B). It is a definition.
\item
  \textbf{Mathlib:}

  \begin{itemize}
  \tightlist
  \item
    Hadamard's three-lines theorem;
  \item
    Cauchy's estimate for the derivative on a circle, and the mean value
    inequality on convex sets;
  \item
    the spectral theorem for Hermitian matrices;
  \item
    the Schur product theorem;
  \item
    dominated convergence;
  \item
    Cauchy's theorem for rectangles and the Cauchy integral formula;
  \item
    multivariable Fourier inversion for continuous torus functions with
    absolutely summable Fourier coefficients;
  \item
    the binomial series in Banach algebras;
  \item
    symmetry of second derivatives;
  \item
    the identity theorem for polynomials on infinite sets.
  \end{itemize}
\item
  \textbf{Standard references:} B. Simon, \emph{Trace Ideals and Their
  Applications}, for Schatten inequalities; Berg--Christensen--Ressel,
  \emph{Harmonic Analysis on Semigroups}, for positive kernels;
  Katznelson, \emph{An Introduction to Harmonic Analysis}, for torus
  Fourier series; Reed--Simon, \emph{Methods of Modern Mathematical
  Physics}, vol.~I, for operator functional calculus; Hörmander,
  \emph{An Introduction to Complex Analysis in Several Variables}, and
  Rudin, \emph{Functional Analysis}, for holomorphy; Glimm--Jaffe,
  \emph{Quantum Physics}, for Wick products and reflection positivity.
\item
  \textbf{Imported constants and assumptions.} All analytic assumptions
  appear in the statements above. The numerical factors are explicit.
  The parameters \(q,w,M,r\) in (A), the strip widths and torus sizes in
  (B), and the label cutoff \(d\) are chosen in the applications.
\end{itemize}

\textbf{Uses in other units.} The principal applications are listed
below; the scope of the longer dependency lists is noted in §7. -
\textbf{III.16.1--III.16.4:} the external applications of this trace
argument are in IV.1. - III.16.1 gives the positivity of the Wilson
transfer kernel:
\href{https://github.com/yang-mills-group/OpenAI-YM-20261028/blob/3456e5e75cf8d4cd558955a5c7638fba2d68f464/YM/Slice/Positive.lean\#L81-L91}{\texttt{Slice.step\_positive}},
Lemma IV.1.f. - III.16.4 gives the time-direction Lipschitz bound:
\href{https://github.com/yang-mills-group/OpenAI-YM-20261028/blob/3456e5e75cf8d4cd558955a5c7638fba2d68f464/YM/Slice/Calendar.lean\#L72-L102}{\texttt{Slice.Time.powering}},
Lemmas IV.1.e--f. - Lemma 4.1 (the algebra of positive kernels) is also
used in III.2, III.4, III.6, III.10, III.11, III.13 and VIII.1. -
\textbf{III.16.5:} III.7, III.3, III.6, III.2, III.4, III.9, III.10,
III.12--III.15, VII.1, VII.2. - \textbf{III.16.6:} III.7, III.3, III.4,
III.2, III.6, III.14, III.10, III.13, III.9, III.11, III.1, III.5,
VII.1, VII.2. - \textbf{III.16.7:} III.6, in the square-root theorem
described in §3. - \textbf{III.16.8 and III.16.10:} VI.2, VII.1, VII.3,
VI.1, III.11, II.1, VII.2, IV.1. - \textbf{III.16.9:} II.1, III.11,
VII.2, III.12, III.13, III.6, VI.2, VI.1, III.10, III.2, VII.3, VII.1,
VII.4, VIII.1. - \textbf{Lemma 4.20:} III.6 and V.2 use
\href{https://github.com/yang-mills-group/OpenAI-YM-20261028/blob/3456e5e75cf8d4cd558955a5c7638fba2d68f464/YM/Analysis/Polydisc.lean\#L222-L238}{\texttt{finiteD}};
III.13, V.1 and V.2 use
\href{https://github.com/yang-mills-group/OpenAI-YM-20261028/blob/3456e5e75cf8d4cd558955a5c7638fba2d68f464/YM/Analysis/Weak.lean\#L21-L50}{\texttt{lift\_weak}}.

\subsubsection{3. Strategy}\label{strategy-18}

\textbf{(A) Moving an insertion.} Transfer-kernel correlations have the
form \(\mathrm{Tr}(B_1K^{Y_1}\cdots B_nK^{Y_n})/Z\). Moving one
insertion transfers \(\ell\) units between two gaps. For finite
matrices, continue the exponents along \(Y+z(e_i-e_j)\). The trace is
entire, and the Hölder balance is unchanged. III.16.2 bounds it by a
constant \(C\) on \(|z|\le2r\); Cauchy's estimate gives \(|F'|\le C/r\)
on \(|z|\le r\), hence a change of at most \(C\ell/r\). The
dimension-free bound passes to kernels by sampling every factor and
adjoint with the same partition. In IV.1, for nonnegative Wilson
coupling, positivity follows from III.16.1 and the positive-kernel
algebra; \(Z\) is the partition function. The condition \(M=2qw\) makes
the reflected column \((B^*B)^q\) fill the time period, so chessboard
estimates bound its normalized trace. The normalized difference is at
most \((\prod_kc_k)\ell/r\); uniformity in volume and spacing requires
corresponding control of \(c_k\) and \(\ell/r\).

\textbf{(B) Kernels and symbols.} Exponentially weighted summability
gives bounded holomorphic symbols. Moving the Fourier integration torus
in the imaginary directions gives the converse, with a narrower
exponential weight. Discrete Fourier identities then transfer symbol
identities to every finite torus, including the zero momentum. For the
square-root construction in III.6, take a real square matrix kernel
\(Q\) with \(\sum_xe^{\eta|x|_1}|Q_{ab}(x)|\le A\) for every entry,
\(\eta>0\), and Hermitian real-momentum symbols satisfying
\(vI\le\hat Q(p)\le AI\), where \(0<v\le A\). The derivative bound,
together with the norm conversion from entrywise bounds to the matrix
operator norm, keeps \(\hat Q/A-I\) in a ball of radius less than \(1\)
on a thinner strip. III.16.7 gives a bounded holomorphic root there. Its
Laurent coefficients have an exponential bound; taking their real parts
preserves the positive square-root symbol at real momenta because
\(\hat Q(-p)=\overline{\hat Q(p)}\). Fourier diagonalization yields
\(P_n^{\mathsf T}P_n=Q_n\) with \(P_n\) symmetric on every torus. The
parameter choices and assembly belong to III.6.

\textbf{(C) Free fields and labels.} For homogeneous closed curvature
jets, the radial-gauge formula \(A_i=(j+2)^{-1}\sum_kx_kF_{ki}\) inverts
the curl. Imposing the Maxwell equation gives the shell. The explicit
two-curl kernel derived from \(G_0(x)=(4\pi^2|x|^2)^{-1}\) solves both
equations in each leg away from the origin, so its Taylor coefficients
define pairings on shell probes. Wick sums of these pairings give the
free hierarchy on polynomial labels. Bihomogeneous decomposition and
conjugation provide coherent complex bases with real coefficients. This
hierarchy is the comparison object for short-distance limits,
stress-tensor identities, the trace anomaly and nontriviality in
VI.1--VI.2 and VII.1--VII.3; its kernels are used at distinct points.

\subsubsection{4. Steps}\label{steps-18}

\paragraph{4.1 Positive kernels and
quadrature}\label{positive-kernels-and-quadrature}

\textbf{Lemma 4.1 (algebra of positive-semidefinite kernels).}

\emph{Statement.} - Generators: the class contains the Gram kernels
\(\overline{f(x)}f(y)\) and the constant \(1\). - Closure: it is closed
under multiplication by constants \(c\ge0\), finite sums, pointwise
(Schur) products, nonnegative integer powers, pullbacks
\(K(\varphi(x),\varphi(y))\), pointwise limits, and \(K\mapsto e^K\). -
Criterion: a kernel is positive semidefinite if and only if it is
Hermitian and \(\sum_{i,j}\bar c_iK(x_i,x_j)c_j\ge0\) for all finite
families. - Integrated form: on a compact metric space with a finite
Borel measure, for continuous positive-semidefinite \(K\) and bounded
measurable \(f\), \(\iint K(x,y)\overline{f(x)}f(y)\,d\mu\,d\mu\ge0\).

\emph{Sketch.} - Products follow from the Schur product theorem. -
\(e^K\) is the limit of the partial sums \(\sum_{m<N}K^m/m!\). - Limits
are handled by the quadratic-form criterion. Hermitian symmetry passes
to the limit through the \(2\times2\) Gram matrices. - For the
integrated form, replace \(K\) by its values at the tags of a fine
partition (Lemma 4.2). The integral becomes
\(\sum_{i,j}\overline{\int_{\text{cell }i}f}\,K(x_i,x_j)\int_{\text{cell }j}f\ge0\),
and dominated convergence finishes the proof.

\emph{Lean:}
\href{https://github.com/yang-mills-group/OpenAI-YM-20261028/blob/3456e5e75cf8d4cd558955a5c7638fba2d68f464/YM/Analysis/Nonnegative.lean\#L24-L28}{\texttt{gram}},
\href{https://github.com/yang-mills-group/OpenAI-YM-20261028/blob/3456e5e75cf8d4cd558955a5c7638fba2d68f464/YM/Analysis/Nonnegative.lean\#L33-L34}{\texttt{scaled}},
\href{https://github.com/yang-mills-group/OpenAI-YM-20261028/blob/3456e5e75cf8d4cd558955a5c7638fba2d68f464/YM/Analysis/Nonnegative.lean\#L35-L36}{\texttt{plus}},
\href{https://github.com/yang-mills-group/OpenAI-YM-20261028/blob/3456e5e75cf8d4cd558955a5c7638fba2d68f464/YM/Analysis/Nonnegative.lean\#L37-L38}{\texttt{times}},
\href{https://github.com/yang-mills-group/OpenAI-YM-20261028/blob/3456e5e75cf8d4cd558955a5c7638fba2d68f464/YM/Analysis/Nonnegative.lean\#L49-L52}{\texttt{pow}},
\href{https://github.com/yang-mills-group/OpenAI-YM-20261028/blob/3456e5e75cf8d4cd558955a5c7638fba2d68f464/YM/Analysis/Nonnegative.lean\#L54-L55}{\texttt{titled}},
\href{https://github.com/yang-mills-group/OpenAI-YM-20261028/blob/3456e5e75cf8d4cd558955a5c7638fba2d68f464/YM/Analysis/Nonnegative.lean\#L69-L75}{\texttt{criterion}},
\href{https://github.com/yang-mills-group/OpenAI-YM-20261028/blob/3456e5e75cf8d4cd558955a5c7638fba2d68f464/YM/Analysis/Nonnegative.lean\#L80-L88}{\texttt{closed}},
\href{https://github.com/yang-mills-group/OpenAI-YM-20261028/blob/3456e5e75cf8d4cd558955a5c7638fba2d68f464/YM/Analysis/Nonnegative.lean\#L90-L99}{\texttt{exp}},
\href{https://github.com/yang-mills-group/OpenAI-YM-20261028/blob/3456e5e75cf8d4cd558955a5c7638fba2d68f464/YM/Analysis/MomentPD.lean\#L19-L62}{\texttt{integrated}}.

\textbf{Lemma 4.2 (quadrature).}

\emph{Statement.} Let \(X\) be a compact metric space with a finite
Borel measure \(\mu\). - For each \(k\ge0\) there is a finite measurable
partition of \(X\) with assigned tags such that every point lies within
\(1/(k+1)\) of its cell's tag. Empty cells and zero masses are allowed;
a tag need not lie in its cell. - Let \(m_i=\mu(\text{cell }i)\). For
every finite index set \(I\) and every continuous \(F:X^I\to E\) with
\(E\) a Banach space, the Riemann sums
\(\sum_{i\in\text{cells}^I}\prod_{j\in I}m_{i_j}\,F\big((x_{i_j})_{j\in I}\big)\)
converge to \(\int F\,d\mu^{\otimes I}\).

\emph{Sketch.} - The partitions: cover \(X\) by finitely many balls of
radius \(1/(k+1)\) and disjointify. - The convergence: the Riemann sum
equals \(\int F\circ\pi_k\,d\mu^{\otimes I}\), where \(\pi_k\) sends
every point to the tag of its cell. Since \(\pi_k\to\mathrm{id}\)
pointwise and \(F\) is bounded, dominated convergence applies. - The
statement concerns each fixed \(F\) and finite product; it asserts no
uniform rate over families of integrands.

\emph{Lean:}
\href{https://github.com/yang-mills-group/OpenAI-YM-20261028/blob/3456e5e75cf8d4cd558955a5c7638fba2d68f464/YM/Analysis/Quadrature.lean\#L22-L26}{\texttt{Quad}},
\href{https://github.com/yang-mills-group/OpenAI-YM-20261028/blob/3456e5e75cf8d4cd558955a5c7638fba2d68f464/YM/Analysis/Quadrature.lean\#L40-L63}{\texttt{refine\_}},
\href{https://github.com/yang-mills-group/OpenAI-YM-20261028/blob/3456e5e75cf8d4cd558955a5c7638fba2d68f464/YM/Analysis/Quadrature.lean\#L65-L66}{\texttt{shrink}},
\href{https://github.com/yang-mills-group/OpenAI-YM-20261028/blob/3456e5e75cf8d4cd558955a5c7638fba2d68f464/YM/Analysis/Quadrature.lean\#L95-L114}{\texttt{counting}},
\href{https://github.com/yang-mills-group/OpenAI-YM-20261028/blob/3456e5e75cf8d4cd558955a5c7638fba2d68f464/YM/Analysis/Quadrature.lean\#L119-L133}{\texttt{approximation}},
\href{https://github.com/yang-mills-group/OpenAI-YM-20261028/blob/3456e5e75cf8d4cd558955a5c7638fba2d68f464/YM/Analysis/Projected.lean\#L19-L38}{\texttt{approx}}.

\textbf{Lemma 4.3 (proof of III.16.1).}

\emph{Sketch.} - For a partition, define
\(G_Q(b,c)=\sum_{i,j}m_im_jK(g(b,x_i),g(c,x_j))\). - \(G_Q\) is
Hermitian. For each finite family of elements \(b\in\mathcal B\) and
coefficients \(r_b\), its quadratic form
\(\sum_{b,c}\bar r_bG_Q(b,c)r_c\) is the quadratic form of \(K\) at the
points \(g(b,x_i)\), indexed by \((b,i)\), with coefficients \(r_bm_i\).
Hence \(G_Q\) is positive semidefinite. - By Lemma 4.2 in two variables,
\(G_Q\to\tilde K\) pointwise. - Positive-semidefinite kernels are closed
under pointwise limits (Lemma 4.1).

\emph{Lean:}
\href{https://github.com/yang-mills-group/OpenAI-YM-20261028/blob/3456e5e75cf8d4cd558955a5c7638fba2d68f464/YM/Analysis/Projected.lean\#L44-L70}{\texttt{accumulate}}.

\paragraph{4.2 The finite trace
inequality}\label{the-finite-trace-inequality}

\textbf{Lemma 4.4 (three-lines Hölder lemma).}

\emph{Definition.} A \emph{strip path} is an entire map
\(V:\mathbb C\to M_\iota(\mathbb C)\) with: - \(\|V(t)\|\le1\) for
\(0\le\mathrm{Re}\,t\le1\); - trace norm
\(\sup_{\|C\|\le1}|\mathrm{tr}\,V(t)C|\le1\) for \(\mathrm{Re}\,t=1\).

\emph{Statement.} If \(V_1,\dots,V_n\) are strip paths and
\(h_k\in\mathbb C\) satisfy \(\mathrm{Re}\,h_k\in[0,1]\) and
\(\sum_k\mathrm{Re}\,h_k=1\), then
\(|\mathrm{tr}\,V_1(h_1)\cdots V_n(h_n)|\le1\).

\emph{Sketch.} Prove \(|\mathrm{tr}(c\,V_1(h_1)\cdots V_n(h_n))|\le1\)
for every contraction \(c\), by induction on \(n\). The case \(n=0\) is
excluded by the balance condition. In the induction step let
\(t=\mathrm{Re}\,h_1\). - \emph{Case \(t=1\).} The other factors are
contractions. Apply the trace-norm bound of \(V_1\) after a cyclic
permutation. - \emph{Case \(t<1\).} Put
\(E(x)=\mathrm{tr}\big(c\,V_1(x+i\,\mathrm{Im}\,h_1)\prod_{k\ge2}V_k(x_k(x))\big)\)
with \(x_k(x)=\mathrm{Re}\,h_k\,\frac{1-x}{1-t}+i\,\mathrm{Im}\,h_k\). -
\(E\) is entire and bounded on \(0\le\mathrm{Re}\,x\le1\). The bound
depends on \(\mathrm{card}\,\iota\), but it is used only qualitatively.
- On \(\mathrm{Re}\,x=0\), the first factor is a contraction absorbed
into \(c\), and the remaining real parts sum to \(1\); the induction
hypothesis gives \(|E|\le1\). - On \(\mathrm{Re}\,x=1\), the factors
\(V_k(x_k(x))\) with \(k\ge2\) are contractions, and the trace-norm
bound of \(V_1\) gives \(|E|\le1\). - The three-lines theorem gives
\(|E(t)|\le1\), and \(E(t)\) is the trace to be bounded.

\emph{Lean:}
\href{https://github.com/yang-mills-group/OpenAI-YM-20261028/blob/3456e5e75cf8d4cd558955a5c7638fba2d68f464/YM/Complex/Holder.lean\#L25-L29}{\texttt{Flex}},
\href{https://github.com/yang-mills-group/OpenAI-YM-20261028/blob/3456e5e75cf8d4cd558955a5c7638fba2d68f464/YM/Complex/Holder.lean\#L73-L126}{\texttt{holder\_core}},
\href{https://github.com/yang-mills-group/OpenAI-YM-20261028/blob/3456e5e75cf8d4cd558955a5c7638fba2d68f464/YM/Complex/Holder.lean\#L130-L132}{\texttt{convexStrip}}.

\textbf{Lemma 4.5 (Schatten factorization).} For every matrix \(b\) and
integer \(q\ge1\) there is a strip path \(V\) with
\(b=\|b\|_{2q}\,V(1/(2q))\).

\emph{Sketch.} - Write \(b^*b=U\,\mathrm{diag}(s_i^2)\,U^*\). Then
\(b=W\,\mathrm{diag}(s)\,U^*\) with \(\|W\|\le1\), where
\(W=bU\,\mathrm{diag}(s)^{-1}\) on the support of \(s\). - If \(b=0\),
take \(V=0\). Otherwise let \(f_i=s_i^{2q}/\sum_js_j^{2q}\) and
\(V(t)=W\,\mathrm{diag}(f_i^{\,t})\,U^*\), with the convention
\(0^t=0\). - Since \(f_i\le1\), \(\|V(t)\|\le1\) on the closed strip. -
For \(\mathrm{Re}\,t=1\),
\(|\mathrm{tr}\,V(t)C|\le\sum_if_i|(U^*CW)_{ii}|\le1\). - At
\(t=1/(2q)\), \(f_i^{1/(2q)}=s_i/\|b\|_{2q}\).

\emph{Lean:}
\href{https://github.com/yang-mills-group/OpenAI-YM-20261028/blob/3456e5e75cf8d4cd558955a5c7638fba2d68f464/YM/Complex/Columns.lean\#L183-L214}{\texttt{even\_cost}},
\href{https://github.com/yang-mills-group/OpenAI-YM-20261028/blob/3456e5e75cf8d4cd558955a5c7638fba2d68f464/YM/Complex/Columns.lean\#L114-L156}{\texttt{imageForm}},
\href{https://github.com/yang-mills-group/OpenAI-YM-20261028/blob/3456e5e75cf8d4cd558955a5c7638fba2d68f464/YM/Complex/Columns.lean\#L60-L81}{\texttt{weighted}},
\href{https://github.com/yang-mills-group/OpenAI-YM-20261028/blob/3456e5e75cf8d4cd558955a5c7638fba2d68f464/YM/Complex/Columns.lean\#L38}{\texttt{specPow}}.

\textbf{Lemma 4.6 (complex powers).} For \(A\succeq0\) and integer
\(M\ge1\) there is a strip path \(P\) with \(A^z=Z_A^{z/M}P(z/M)\) for
all \(z\in\mathbb C\). Here \(Z_A^w=e^{w\log Z_A}\) if \(Z_A>0\), and
\(Z_A^w=0\) otherwise. Moreover \(A^z\) equals the ordinary power
\(A^m\) at every integer \(z=m\ge1\).

\emph{Sketch.} If \(Z_A=0\), positivity gives \(A=0\) and one may take
\(P=0\). Otherwise take
\(P(w)=U\,\mathrm{diag}\big((\lambda_i^M/Z_A)^w\big)U^*\). The
normalized eigenvalues sum to \(1\), so the strip and trace-norm bounds
follow as in Lemma 4.5.

\emph{Lean:}
\href{https://github.com/yang-mills-group/OpenAI-YM-20261028/blob/3456e5e75cf8d4cd558955a5c7638fba2d68f464/YM/Complex/Gaps.lean\#L22-L24}{\texttt{leap}},
\href{https://github.com/yang-mills-group/OpenAI-YM-20261028/blob/3456e5e75cf8d4cd558955a5c7638fba2d68f464/YM/Complex/Gaps.lean\#L26-L32}{\texttt{leap\_at}},
\href{https://github.com/yang-mills-group/OpenAI-YM-20261028/blob/3456e5e75cf8d4cd558955a5c7638fba2d68f464/YM/Complex/Gaps.lean\#L39-L63}{\texttt{flexPower}}.

\textbf{Lemma 4.7 (proof of III.16.2).}

\emph{Sketch.} - Write \(b_k=\|b_k\|_{2q}V_k(1/(2q))\) (Lemma 4.5) and
\(A^{v_k}=Z_A^{v_k/M}P(v_k/M)\) (Lemma 4.6). - The scalar prefactors
have modulus
\(\prod_k\|b_k\|_{2q}\,Z_A^{\sum\mathrm{Re}v_k/M}=\prod_k\|b_k\|_{2q}\,Z_A^{1-n/(2q)}\).
- The \(2n\) strip paths \(V_1,P,V_2,P,\dots\) are evaluated at
exponents whose real parts, \(1/(2q)\) and \(\mathrm{Re}\,v_k/M\), lie
in \([0,1]\) and sum to \(1\). - Lemma 4.4 bounds the remaining trace by
\(1\).

\emph{Lean:}
\href{https://github.com/yang-mills-group/OpenAI-YM-20261028/blob/3456e5e75cf8d4cd558955a5c7638fba2d68f464/YM/Complex/Gaps.lean\#L124-L165}{\texttt{estimation}}.

\textbf{Lemma 4.8 (proof of III.16.3).}

\emph{Sketch.} - Let \(F(z)=\mathrm{tr}\prod_kb_kA^{W_k(z)}\) with
\(W(z)=Y+z(e_i-e_j)\), and put
\(C=(\prod_k\|b_k\|_{2q})(\mathrm{tr}\,A^M)^{1-n/(2q)}\). The function
\(F\) is entire. - The relation \(nw+\sum Y_k=M\) is the balance
\(n/(2q)+\sum_kY_k/M=1\), and \(\sum_k\mathrm{Re}\,W_k(z)=\sum_kY_k\). -
On \(|z|\le2r\), \(\mathrm{Re}\,W_k(z)\ge Y_k-2r>0\), so III.16.2 gives
\(|F|\le C\) there. - For \(|x|\le r\), the circle of radius \(r\) about
\(x\) lies in this disc, and Cauchy's estimate gives \(|F'(x)|\le C/r\).
- The mean value inequality on the convex disc \(|z|\le r\), between
\(0\) and \(\ell\), gives the claim. - \(F(0)\) and \(F(\ell)\) are the
two traces, because all \(Y_k,Y'_k\ge1\).

\emph{Lean:}
\href{https://github.com/yang-mills-group/OpenAI-YM-20261028/blob/3456e5e75cf8d4cd558955a5c7638fba2d68f464/YM/Complex/Gaps.lean\#L170-L229}{\texttt{finite\_lip}}.

\paragraph{4.3 From matrices to kernels}\label{from-matrices-to-kernels}

\textbf{Lemma 4.9 (sampled traces).} Let \(L\ge1\). For a partition with
masses \(m_i\) and tags \(x_i\), sample a kernel as
\(k^Q=(\sqrt{m_i}\,k(x_i,x_j)\sqrt{m_j})_{ij}\). Then: -
\((k^*)^Q=(k^Q)^*\); - \(k^Q\) is positive semidefinite whenever \(k\)
is; -
\(\mathrm{tr}(k_1^Q\cdots k_L^Q)=\sum_{i_1,\dots,i_L}\prod_lm_{i_l}\prod_lk_l(x_{i_l},x_{i_{l+1}})\),
because every vertex contributes two square-root weights; - this Riemann
sum converges to \(\mathrm{Tr}(k_1\cdots k_L)\) for continuous \(k_l\)
along the partitions of Lemma 4.2.

\emph{Lean:}
\href{https://github.com/yang-mills-group/OpenAI-YM-20261028/blob/3456e5e75cf8d4cd558955a5c7638fba2d68f464/YM/TimeKernel/Trace.lean\#L81}{\texttt{sample}},
\href{https://github.com/yang-mills-group/OpenAI-YM-20261028/blob/3456e5e75cf8d4cd558955a5c7638fba2d68f464/YM/TimeKernel/Trace.lean\#L90-L94}{\texttt{safe}},
\href{https://github.com/yang-mills-group/OpenAI-YM-20261028/blob/3456e5e75cf8d4cd558955a5c7638fba2d68f464/YM/TimeKernel/Trace.lean\#L96-L104}{\texttt{measureWeight}},
\href{https://github.com/yang-mills-group/OpenAI-YM-20261028/blob/3456e5e75cf8d4cd558955a5c7638fba2d68f464/YM/TimeKernel/Trace.lean\#L109-L119}{\texttt{learns}}.

\textbf{Lemma 4.10 (proof of III.16.4).}

\emph{Sketch.} Apply III.16.3, for each partition, to \(A=K^Q\) and
\(b_l=B_l^Q\). - By Lemma 4.9, \(\mathrm{tr}A^M\to Z>0\), so the
hypothesis \(\mathrm{tr}A^M>0\) holds for fine partitions. -
\(\|b_l\|_{2q}\to|\mathrm{Tr}((B_l^*B_l)^q)|^{1/(2q)}\), because the
column word is sampled with the same partition. - The two gap traces
converge to the cyclic integrals. - Passing to the limit at fixed
\((q,w,Y,Y',\ell,r)\) gives the bound
\(\frac{\ell}{r}\prod_l|\mathrm{Tr}((B_l^*B_l)^q)|^{1/(2q)}\,Z^{1-n/(2q)}\).
- The column hypothesis turns this into \(Z\prod_lc_l\,\ell/r\).

\emph{Lean:}
\href{https://github.com/yang-mills-group/OpenAI-YM-20261028/blob/3456e5e75cf8d4cd558955a5c7638fba2d68f464/YM/TimeKernel/Unit.lean\#L76-L106}{\texttt{passage}},
\href{https://github.com/yang-mills-group/OpenAI-YM-20261028/blob/3456e5e75cf8d4cd558955a5c7638fba2d68f464/YM/TimeKernel/Unit.lean\#L45-L58}{\texttt{goes}},
\href{https://github.com/yang-mills-group/OpenAI-YM-20261028/blob/3456e5e75cf8d4cd558955a5c7638fba2d68f464/YM/TimeKernel/Unit.lean\#L60-L72}{\texttt{totality}}.

\paragraph{4.4 Holomorphy in a strip and Laurent
coefficients}\label{holomorphy-in-a-strip-and-laurent-coefficients}

\textbf{Lemma 4.11 (torus integrals inside the strip).} These are the
first three claims of III.16.5(a).

\emph{Sketch.} - The coefficient bound: \(\hat u_p(x)\) is an average of
\(u\) over the torus \(p+2\pi[0,1]^\alpha\) against a unimodular
character. - Independence of \(p\): change one coordinate at a time. A
real shift leaves a full-period integral unchanged. For an imaginary
shift in coordinate \(i\), apply Cauchy's theorem to a rectangle of
horizontal width \(2\pi\) between heights \(\mathrm{Im}\,p_i\) and
\(\mathrm{Im}\,p'_i\). Periodicity cancels the vertical sides, so the
horizontal integrals agree. - The identity
\(\hat u_p(x)=e^{ip\cdot x}\hat u(x)\) is independence applied to
\(u(p)e^{-ip\cdot x}\).

\emph{Lean:}
\href{https://github.com/yang-mills-group/OpenAI-YM-20261028/blob/3456e5e75cf8d4cd558955a5c7638fba2d68f464/YM/Laurent/Rune.lean\#L91-L99}{\texttt{market\_cap}},
\href{https://github.com/yang-mills-group/OpenAI-YM-20261028/blob/3456e5e75cf8d4cd558955a5c7638fba2d68f464/YM/Laurent/Turn.lean\#L103-L112}{\texttt{full\_turn}},
\href{https://github.com/yang-mills-group/OpenAI-YM-20261028/blob/3456e5e75cf8d4cd558955a5c7638fba2d68f464/YM/Laurent/Turn.lean\#L114-L118}{\texttt{coat\_price}}.

\textbf{Lemma 4.12 (weighted \(\ell^1\) bound).} Under the holomorphy
and periodicity hypotheses of III.16.5(a), if \(|u|\le K\) on
\(S_\eta\), \(H\ge0\), \(b>0\), and \(H+b<\eta\), then
\(\sum_xe^{H|x|_1}|\hat u(x)|\le K(2(1+b)/b)^{\mathrm{card}\,\alpha}\).

\emph{Sketch.} - For given \(x\), take \(p\) with
\(\mathrm{Im}\,p_a=-(H+b)\,\mathrm{sign}\,x_a\). Then
\(|e^{ip\cdot x}|=e^{(H+b)|x|_1}\), and Lemma 4.11 gives
\(|\hat u(x)|\le Ke^{-(H+b)|x|_1}\). - Then
\(\sum_{x\in\mathbb Z^\alpha}e^{-b|x|_1}\le\big(2\sum_{k\ge0}e^{-bk}\big)^{\mathrm{card}\,\alpha}\le(2(1+b)/b)^{\mathrm{card}\,\alpha}\).

\emph{Lean:}
\href{https://github.com/yang-mills-group/OpenAI-YM-20261028/blob/3456e5e75cf8d4cd558955a5c7638fba2d68f464/YM/Laurent/Art.lean\#L61-L87}{\texttt{Vane.issues}},
\href{https://github.com/yang-mills-group/OpenAI-YM-20261028/blob/3456e5e75cf8d4cd558955a5c7638fba2d68f464/YM/Laurent/Art.lean\#L39-L57}{\texttt{Width.compactPay}},
\href{https://github.com/yang-mills-group/OpenAI-YM-20261028/blob/3456e5e75cf8d4cd558955a5c7638fba2d68f464/YM/Laurent/Art.lean\#L13-L35}{\texttt{pins}}.

\textbf{Lemma 4.13 (Laurent representation).} Under the same holomorphy
and periodicity hypotheses, if \(|u|\le K\) on \(S_\eta\), then
\(u(p)=\sum_x\hat u(x)e^{ip\cdot x}\) on \(S_\eta\).

\emph{Sketch.} - By Lemma 4.12, with
\(H=\max(\{0\}\cup\{|\mathrm{Im}\,p_a|:a\in\alpha\})\) and
\(b=(\eta-H)/2\), the Fourier series of
\(\theta\mapsto u(p+2\pi\theta)\) is absolutely summable. - Fourier
inversion for the continuous shifted function identifies this series
with the function. - Evaluate at \(\theta=0\) and use
\(\hat u_p(x)=e^{ip\cdot x}\hat u(x)\).

\emph{Lean:}
\href{https://github.com/yang-mills-group/OpenAI-YM-20261028/blob/3456e5e75cf8d4cd558955a5c7638fba2d68f464/YM/Laurent/Art.lean\#L90-L104}{\texttt{Vane.paints}}.

\textbf{Lemma 4.14 (Laurent sums).} This is III.16.5(b).

\emph{Sketch.} - Differentiate term by term, using
\(\|D_pe^{ip\cdot x}\|\le|x|_1e^{\sigma|x|_1}\le(\eta-\sigma)^{-1}e^{\eta|x|_1}\)
on \(S_\sigma\). - The derivative series is uniformly summable on the
open convex strip, which gives holomorphy and the derivative bound. -
The Lipschitz bound follows from the mean value inequality. - The
real-part identity follows from
\(\overline{e^{-ip\cdot x}}=e^{ip\cdot x}\) for real \(p\).

\emph{Lean:}
\href{https://github.com/yang-mills-group/OpenAI-YM-20261028/blob/3456e5e75cf8d4cd558955a5c7638fba2d68f464/YM/Laurent/Eval.lean\#L131-L133}{\texttt{Width.estimate}},
\href{https://github.com/yang-mills-group/OpenAI-YM-20261028/blob/3456e5e75cf8d4cd558955a5c7638fba2d68f464/YM/Laurent/Eval.lean\#L48-L56}{\texttt{reward}},
\href{https://github.com/yang-mills-group/OpenAI-YM-20261028/blob/3456e5e75cf8d4cd558955a5c7638fba2d68f464/YM/Work/F07/Read.lean\#L55-L95}{\texttt{readDeriv}},
\href{https://github.com/yang-mills-group/OpenAI-YM-20261028/blob/3456e5e75cf8d4cd558955a5c7638fba2d68f464/YM/Work/F07/Read.lean\#L97-L101}{\texttt{readDiff}},
\href{https://github.com/yang-mills-group/OpenAI-YM-20261028/blob/3456e5e75cf8d4cd558955a5c7638fba2d68f464/YM/Work/F07/Read.lean\#L119-L133}{\texttt{readRealPart}}.

\paragraph{4.5 Fourier analysis on tori}\label{fourier-analysis-on-tori}

\textbf{Lemma 4.15 (orthogonality).}
\(\sum_{x\in\Lambda_n}e^{i(k_p-k_{p'})\cdot x}=n^4\delta_{pp'}\), and
hence \(F_nF_n^*=F_n^*F_n=n^4I\).

\emph{Sketch.} For \(p\ne0\) some unit step \(e_a\) has
\(e^{ik_p\cdot e_a}\ne1\). Translation invariance of the sum then forces
it to vanish.

\emph{Lean:}
\href{https://github.com/yang-mills-group/OpenAI-YM-20261028/blob/3456e5e75cf8d4cd558955a5c7638fba2d68f464/YM/CellFourier/Mesh.lean\#L98-L104}{\texttt{sum\_cell}},
\href{https://github.com/yang-mills-group/OpenAI-YM-20261028/blob/3456e5e75cf8d4cd558955a5c7638fba2d68f464/YM/CellFourier/Mesh.lean\#L106-L110}{\texttt{orthogonal}},
\href{https://github.com/yang-mills-group/OpenAI-YM-20261028/blob/3456e5e75cf8d4cd558955a5c7638fba2d68f464/YM/CellFourier/Mesh.lean\#L117-L119}{\texttt{gram}},
\href{https://github.com/yang-mills-group/OpenAI-YM-20261028/blob/3456e5e75cf8d4cd558955a5c7638fba2d68f464/YM/CellFourier/Mesh.lean\#L121-L127}{\texttt{coGram}}.

\textbf{Lemma 4.16 (circulant diagonalization; III.16.6(a)).} Under the
hypotheses of III.16.6, periodization is diagonalized by the Fourier
matrix. For real \(p\), the symbol is \(2\pi\)-periodic and satisfies
\(\hat R(-\bar p)=\overline{\hat R(p)}\) entrywise. These identities
also hold at complex momenta wherever the series converge absolutely.

\emph{Sketch.} For fixed \(a,b,p\), summing the periodized kernel
against the Fourier column gives
\(\sum_zR_{ab}(z)e^{ik_p\cdot(y+z)}=e^{ik_p\cdot y}\hat R(k_p)_{ab}\).
Absolute convergence permits regrouping by residues. Periodicity follows
from \(e^{2\pi i m\cdot z}=1\), and conjugation follows term by term
because \(R\) is real.

\emph{Lean:}
\href{https://github.com/yang-mills-group/OpenAI-YM-20261028/blob/3456e5e75cf8d4cd558955a5c7638fba2d68f464/YM/CellFourier/Circulant.lean\#L88-L104}{\texttt{circulant}},
\href{https://github.com/yang-mills-group/OpenAI-YM-20261028/blob/3456e5e75cf8d4cd558955a5c7638fba2d68f464/YM/CellFourier/Circulant.lean\#L49-L50}{\texttt{Row.periods}},
\href{https://github.com/yang-mills-group/OpenAI-YM-20261028/blob/3456e5e75cf8d4cd558955a5c7638fba2d68f464/YM/CellFourier/Circulant.lean\#L39-L45}{\texttt{Row.conjugate}}.

\textbf{Lemma 4.17 (aliases; III.16.6(b)--(c)).}

\emph{Sketch.} - The fine Fourier matrix factors as the coarse Fourier
matrix composed with the alias blocks \(E_L(2\pi p/c)\). -
\(E_L(\bar p)^*E_L(p)=L^4I\) is Lemma 4.15 on \(\Lambda_L\). The common
factor \(e^{ip\cdot x/L}\) cancels against its conjugate taken at
\(\bar p\). - The intertwining relation is Lemma 4.16 on the fine torus,
using \(e^{ik_{p,l}\cdot fz}=1\).

\emph{Lean:}
\href{https://github.com/yang-mills-group/OpenAI-YM-20261028/blob/3456e5e75cf8d4cd558955a5c7638fba2d68f464/YM/CellFourier/Alias.lean\#L20}{\texttt{aliasAt}},
\href{https://github.com/yang-mills-group/OpenAI-YM-20261028/blob/3456e5e75cf8d4cd558955a5c7638fba2d68f464/YM/CellFourier/Alias.lean\#L30-L31}{\texttt{aliasBlock}},
\href{https://github.com/yang-mills-group/OpenAI-YM-20261028/blob/3456e5e75cf8d4cd558955a5c7638fba2d68f464/YM/CellFourier/Alias.lean\#L41-L53}{\texttt{straightens}},
\href{https://github.com/yang-mills-group/OpenAI-YM-20261028/blob/3456e5e75cf8d4cd558955a5c7638fba2d68f464/YM/CellFourier/Alias.lean\#L65-L69}{\texttt{energy}},
\href{https://github.com/yang-mills-group/OpenAI-YM-20261028/blob/3456e5e75cf8d4cd558955a5c7638fba2d68f464/YM/CellFourier/Shuttle.lean\#L46}{\texttt{readings}},
\href{https://github.com/yang-mills-group/OpenAI-YM-20261028/blob/3456e5e75cf8d4cd558955a5c7638fba2d68f464/YM/CellFourier/Shuttle.lean\#L66-L75}{\texttt{readingsGram}},
\href{https://github.com/yang-mills-group/OpenAI-YM-20261028/blob/3456e5e75cf8d4cd558955a5c7638fba2d68f464/YM/CellFourier/Shuttle.lean\#L110-L129}{\texttt{aligns}}.

\paragraph{4.6 Matrix square root}\label{matrix-square-root}

\textbf{Lemma 4.18 (= III.16.7).}

\emph{Sketch.} - Diagonalize \(X=U\,\mathrm{diag}(s)\,U^*\) with
\(v\le s_i\le m\). Then \(X/m-1\) has eigenvalues of modulus
\(\le1-v/m<1\). - The scalar identity
\(\sum_k\binom{1/2}{k}t^k=\sqrt{1+t}\) for \(|t|<1\), applied to the
eigenvalues, gives \(\Phi_m(X)=U\,\mathrm{diag}(\sqrt{s_i})\,U^*\). -
The binomial series in the Banach algebra of matrices has radius of
convergence at least \(1\), giving holomorphy. Continuity and
compactness give boundedness on each smaller closed ball. Its real
scalar coefficients imply commutation with entrywise conjugation.

\emph{Lean:}
\href{https://github.com/yang-mills-group/OpenAI-YM-20261028/blob/3456e5e75cf8d4cd558955a5c7638fba2d68f464/YM/Work/F07/MatrixRoot.lean\#L185-L198}{\texttt{binRoot\_on\_band}},
\href{https://github.com/yang-mills-group/OpenAI-YM-20261028/blob/3456e5e75cf8d4cd558955a5c7638fba2d68f464/YM/Work/F07/MatrixRoot.lean\#L104-L114}{\texttt{binSum\_deriv}},
\href{https://github.com/yang-mills-group/OpenAI-YM-20261028/blob/3456e5e75cf8d4cd558955a5c7638fba2d68f464/YM/Work/F07/MatrixRoot.lean\#L21-L24}{\texttt{Band}}.

\paragraph{4.7 Taylor coefficients and
analyticity}\label{taylor-coefficients-and-analyticity}

\textbf{Lemma 4.19 (Taylor coefficients).} Let \(E,F\) be real normed
spaces, let \(j,l\ge0\) be integers, let \(f:E\to F\), and put
\(c_jf(x)=D^jf(x)/j!\).

\begin{itemize}
\tightlist
\item
  If \(f\) is \(C^j\) at \(x\), then \(D^jf(x)\) is symmetric.
\item
  Two symmetric continuous \(j\)-linear forms with the same diagonal are
  equal.
\item
  If \(U\subset E\) is open, \(f\) is smooth on \(U\), \(x\in U\),
  \(u\in E^j\), and \(v\in E^l\), then
  \[c_j\big(y\mapsto c_lf(y)[v]\big)(x)[u]=c_l\big(y\mapsto c_jf(y)[u]\big)(x)[v].\]
\item
  If \(f\) is \(C^j\) at \(-x\), then
  \(c_j(f(-\cdot))(x)=(-1)^jc_jf(-x)\). If \(f\) is \(C^j\) at \(x-y\),
  then \(c_j(y'\mapsto f(x-y'))(y)=(-1)^jc_jf(x-y)\).
\item
  If \(f\) is \(C^{j+1}\) at \(x\), then, for \(u,z\in E\),
  \[D_u\big(t\mapsto c_{j+1}f(x)[t,\dots,t]\big)(z)=c_j(D_uf)(x)[z,\dots,z].\]
\end{itemize}

\emph{Sketch.} Induct on \(j\) for symmetry, using symmetry of the
second derivative to exchange the first two slots and the induction
hypothesis for the remaining slots. Equality on diagonals then implies
equality of the forms by differentiating the diagonal and inducting on
the arity. For commutation, peel off one directional derivative at a
time. The chain rule gives the reflection identities. Differentiating a
symmetric diagonal contributes \(j+1\) equal terms; the factorials give
the last identity.

\emph{Lean:}
\href{https://github.com/yang-mills-group/OpenAI-YM-20261028/blob/3456e5e75cf8d4cd558955a5c7638fba2d68f464/YM/Analysis/SmoothJets.lean\#L71-L134}{\texttt{sym\_Jet}},
\href{https://github.com/yang-mills-group/OpenAI-YM-20261028/blob/3456e5e75cf8d4cd558955a5c7638fba2d68f464/YM/Analysis/SmoothJets.lean\#L54-L60}{\texttt{commuteD}},
\href{https://github.com/yang-mills-group/OpenAI-YM-20261028/blob/3456e5e75cf8d4cd558955a5c7638fba2d68f464/YM/Analysis/Symmetric.lean\#L62-L77}{\texttt{ext\_diag}},
\href{https://github.com/yang-mills-group/OpenAI-YM-20261028/blob/3456e5e75cf8d4cd558955a5c7638fba2d68f464/YM/Analysis/Homogeneous.lean\#L65-L66}{\texttt{coeff}},
\href{https://github.com/yang-mills-group/OpenAI-YM-20261028/blob/3456e5e75cf8d4cd558955a5c7638fba2d68f464/YM/Analysis/Exchange.lean\#L36-L69}{\texttt{exchangeJet}},
\href{https://github.com/yang-mills-group/OpenAI-YM-20261028/blob/3456e5e75cf8d4cd558955a5c7638fba2d68f464/YM/Analysis/Exchange.lean\#L72-L99}{\texttt{reflectionJet}},
\href{https://github.com/yang-mills-group/OpenAI-YM-20261028/blob/3456e5e75cf8d4cd558955a5c7638fba2d68f464/YM/Analysis/Exchange.lean\#L102-L108}{\texttt{oppositeCoeff}},
\href{https://github.com/yang-mills-group/OpenAI-YM-20261028/blob/3456e5e75cf8d4cd558955a5c7638fba2d68f464/YM/Analysis/Homogeneous.lean\#L71-L78}{\texttt{slot\_lower}}.

\textbf{Lemma 4.20 (several-variable and weak analyticity).} - A map
between finite-dimensional complex normed spaces that is
complex-differentiable on a neighbourhood of \(x\) is analytic at \(x\).
- Let \(E\) be a complex Banach space and \(f:\mathbb C\to E\)
continuous on an open set \(U\). If \(L\circ f\) is holomorphic on \(U\)
for every \(L\) in a family of continuous complex linear functionals
that separates points, then \(f\) is analytic on \(U\).

\emph{Sketch.} - First claim: choose coordinates and a closed polydisc
contained in the differentiability neighbourhood. Successive
applications of the one-variable Cauchy formula express each scalar
component as a continuous linear functional applied to the product of
the Cauchy kernels and the fixed boundary values. Each kernel is
analytic in the interior variables, since its denominator is invertible
in the Banach algebra of continuous functions on the angle torus. Their
finite product and the resulting integral are analytic. Coordinate
changes recover the assertion for the original spaces. - Second claim:
the Cauchy integral formula holds after applying each \(L\), the
separating family lifts it to \(E\), and the Cauchy power series
converges.

\emph{Lean:}
\href{https://github.com/yang-mills-group/OpenAI-YM-20261028/blob/3456e5e75cf8d4cd558955a5c7638fba2d68f464/YM/Analysis/Polydisc.lean\#L222-L238}{\texttt{finiteD}},
\href{https://github.com/yang-mills-group/OpenAI-YM-20261028/blob/3456e5e75cf8d4cd558955a5c7638fba2d68f464/YM/Analysis/Weak.lean\#L21-L50}{\texttt{lift\_weak}}.

\paragraph{4.8 Free Maxwell field}\label{free-maxwell-field}

\textbf{Lemma 4.21 (Poincaré lemma in radial gauge; III.16.8(a)).}

\emph{Sketch.} - The curl is closed by symmetry of second derivatives. -
Inverse on \(\mathcal C_j\): for \(F\in\mathcal C_j\) and
\(A_i=(j+2)^{-1}x_kF_{ki}\), compute
\(\partial_iA_l-\partial_lA_i=(j+2)^{-1}\big(2F_{il}+x_k(\partial_iF_{kl}-\partial_lF_{ki})\big)\).
By Bianchi this equals
\((j+2)^{-1}\big(2F_{il}+x_k\partial_kF_{il}\big)\), and by Euler's
identity this is \(F_{il}\). - The map lands in radial gauge, since
\(x_ix_kF_{ki}=0\) by antisymmetry. - Inverse on \(\mathcal A_j\):
differentiating \(x_kA_k=0\) gives \(x_k\partial_iA_k=-A_i\). Hence
\(x_kF_{ki}=(j+1)A_i+A_i=(j+2)A_i\).

\emph{Lean:}
\href{https://github.com/yang-mills-group/OpenAI-YM-20261028/blob/3456e5e75cf8d4cd558955a5c7638fba2d68f464/YM/Maxwell/Exterior.lean\#L97-L98}{\texttt{takeCurl}},
\href{https://github.com/yang-mills-group/OpenAI-YM-20261028/blob/3456e5e75cf8d4cd558955a5c7638fba2d68f464/YM/Maxwell/Exterior.lean\#L108-L128}{\texttt{closedCurl}},
\href{https://github.com/yang-mills-group/OpenAI-YM-20261028/blob/3456e5e75cf8d4cd558955a5c7638fba2d68f464/YM/Maxwell/Exterior.lean\#L130-L154}{\texttt{contracted}},
\href{https://github.com/yang-mills-group/OpenAI-YM-20261028/blob/3456e5e75cf8d4cd558955a5c7638fba2d68f464/YM/Maxwell/Exterior.lean\#L163-L195}{\texttt{curl\_inverse}},
\href{https://github.com/yang-mills-group/OpenAI-YM-20261028/blob/3456e5e75cf8d4cd558955a5c7638fba2d68f464/YM/Maxwell/Exterior.lean\#L197-L214}{\texttt{potential\_recover}},
\href{https://github.com/yang-mills-group/OpenAI-YM-20261028/blob/3456e5e75cf8d4cd558955a5c7638fba2d68f464/YM/Maxwell/Exterior.lean\#L218-L224}{\texttt{potential}}.

\textbf{Lemma 4.22 (the shell; III.16.8(b)).} - \(\mathcal S_j\) is a
subspace of \(\mathcal C_j\). - The inverse curl restricted to the shell
is injective:
\href{https://github.com/yang-mills-group/OpenAI-YM-20261028/blob/3456e5e75cf8d4cd558955a5c7638fba2d68f464/YM/Maxwell/Shell.lean\#L46-L47}{\texttt{gaugeJet}},
\href{https://github.com/yang-mills-group/OpenAI-YM-20261028/blob/3456e5e75cf8d4cd558955a5c7638fba2d68f464/YM/Maxwell/Shell.lean\#L49-L50}{\texttt{leg\_injective}}.

\emph{Sketch for III.16.8(b).} By Lemma 4.19, for \(j\ge1\) the
derivatives of the \(j\)-th homogeneous Taylor polynomial are the
\((j-1)\)-th Taylor coefficients of the derivatives; the degree-zero
case is constant. The Bianchi and Maxwell combinations of \(f\) vanish
near \(x\), so the corresponding combinations of Taylor coefficients
vanish.

\emph{Lean:}
\href{https://github.com/yang-mills-group/OpenAI-YM-20261028/blob/3456e5e75cf8d4cd558955a5c7638fba2d68f464/YM/Maxwell/Shell.lean\#L18-L42}{\texttt{onShell}},
\href{https://github.com/yang-mills-group/OpenAI-YM-20261028/blob/3456e5e75cf8d4cd558955a5c7638fba2d68f464/YM/Maxwell/Shell.lean\#L57-L62}{\texttt{WaveOn}},
\href{https://github.com/yang-mills-group/OpenAI-YM-20261028/blob/3456e5e75cf8d4cd558955a5c7638fba2d68f464/YM/Maxwell/Shell.lean\#L73-L122}{\texttt{WaveOn.jets}}.

\textbf{Lemma 4.23 (the free kernel; III.16.8(c)).}

\emph{Sketch.} -
\(\partial_i\partial_jG_0=-2N\delta_{ij}|x|^{-4}+8Nx_ix_j|x|^{-6}\) with
\(N=(4\pi^2)^{-1}\). The trace vanishes in dimension four. -
\(\psi_{\cdot,km}\) is divergence-free and harmonic, so
\(\sum_a\partial_aC_{ab,km}=\Delta\psi_b-\partial_b\,\mathrm{div}\,\psi=0\).
- Bianchi holds because \(C\) is a curl. - Exchange and evenness follow
from the explicit Hessian.

\emph{Lean:}
\href{https://github.com/yang-mills-group/OpenAI-YM-20261028/blob/3456e5e75cf8d4cd558955a5c7638fba2d68f464/YM/Maxwell/PotentialKernel.lean\#L22}{\texttt{Green}},
\href{https://github.com/yang-mills-group/OpenAI-YM-20261028/blob/3456e5e75cf8d4cd558955a5c7638fba2d68f464/YM/Maxwell/PotentialKernel.lean\#L49-L66}{\texttt{second\_Green}},
\href{https://github.com/yang-mills-group/OpenAI-YM-20261028/blob/3456e5e75cf8d4cd558955a5c7638fba2d68f464/YM/Maxwell/PotentialKernel.lean\#L69-L77}{\texttt{harmonic}},
\href{https://github.com/yang-mills-group/OpenAI-YM-20261028/blob/3456e5e75cf8d4cd558955a5c7638fba2d68f464/YM/Maxwell/Curls.lean\#L17-L18}{\texttt{curlOpp}},
\href{https://github.com/yang-mills-group/OpenAI-YM-20261028/blob/3456e5e75cf8d4cd558955a5c7638fba2d68f464/YM/Maxwell/Curls.lean\#L20}{\texttt{bareCurl}},
\href{https://github.com/yang-mills-group/OpenAI-YM-20261028/blob/3456e5e75cf8d4cd558955a5c7638fba2d68f464/YM/Maxwell/Curls.lean\#L82-L98}{\texttt{Bare.equation}},
\href{https://github.com/yang-mills-group/OpenAI-YM-20261028/blob/3456e5e75cf8d4cd558955a5c7638fba2d68f464/YM/Maxwell/Curls.lean\#L104-L109}{\texttt{Bare.exchange}}.

\textbf{Lemma 4.24 (jet pairing; III.16.8(d)).}

\emph{Sketch.} - By III.16.8(b)--(c), the Taylor coefficients of \(C\)
in either leg lie on the shell. Hence only the restrictions of \(p\) and
\(q\) matter. - Symmetry: it suffices to check probes that read single
mixed Taylor coefficients, because these span the dual of the shell. For
those, the claim follows from: - commutation of Taylor coefficients of
different orders (Lemma 4.19); - the reflection rule, with sign
\((-1)^l\); - \(C_{ab,km}(x)=C_{km,ab}(-x)\).

\emph{Lean:}
\href{https://github.com/yang-mills-group/OpenAI-YM-20261028/blob/3456e5e75cf8d4cd558955a5c7638fba2d68f464/YM/Maxwell/Leg.lean\#L64}{\texttt{Probe}},
\href{https://github.com/yang-mills-group/OpenAI-YM-20261028/blob/3456e5e75cf8d4cd558955a5c7638fba2d68f464/YM/Maxwell/Leg.lean\#L145-L146}{\texttt{contraction}},
\href{https://github.com/yang-mills-group/OpenAI-YM-20261028/blob/3456e5e75cf8d4cd558955a5c7638fba2d68f464/YM/Maxwell/Leg.lean\#L148-L159}{\texttt{descends}},
\href{https://github.com/yang-mills-group/OpenAI-YM-20261028/blob/3456e5e75cf8d4cd558955a5c7638fba2d68f464/YM/Maxwell/Leg.lean\#L168-L172}{\texttt{kernel\_smooth}},
\href{https://github.com/yang-mills-group/OpenAI-YM-20261028/blob/3456e5e75cf8d4cd558955a5c7638fba2d68f464/YM/Maxwell/Bilinear.lean\#L159-L169}{\texttt{exchange}}.

\textbf{Lemma 4.25 (Wick hierarchy; III.16.10).}

\emph{Definitions.} - \emph{Pairings} are fixed-point-free involutions
of the factor occurrences whose partners sit at different points; no
self-contractions are allowed. - The \emph{moment} is defined on the
monomial basis and extended multilinearly. - Labels enter through the
substitution \(A_j=\mathrm{curl}^{-1}F_j\), \(F_j\in\mathcal S_j\). -
The \emph{equation ideal} is the kernel of this substitution. It
consists of the labels that vanish on every on-shell jet sequence.

\emph{Sketch of the properties.} - Vanishing on the equation ideal:
multilinearity. - Relabelling and insertion of \(1\): added unit
monomials have no occurrences. An injection of the remaining indices
induces bijections of occurrences and pairings. If the orientation of an
edge reverses, the exchange identity for \(\kappa\) preserves its value.
Extend from monomials by multilinearity. - Smoothness: each polynomial
has finitely many monomials and each monomial finitely many pairings, so
the result is a finite sum of products of smooth kernels off the
diagonal. When all insertions are \(1\), there are no occurrences and
exactly one empty pairing, whose product is \(1\).

\emph{Lean:}
\href{https://github.com/yang-mills-group/OpenAI-YM-20261028/blob/3456e5e75cf8d4cd558955a5c7638fba2d68f464/YM/Maxwell/Wick.lean\#L24}{\texttt{Wick.entry}},
\href{https://github.com/yang-mills-group/OpenAI-YM-20261028/blob/3456e5e75cf8d4cd558955a5c7638fba2d68f464/YM/Maxwell/Wick.lean\#L75-L76}{\texttt{Wick.Occ.Matching}},
\href{https://github.com/yang-mills-group/OpenAI-YM-20261028/blob/3456e5e75cf8d4cd558955a5c7638fba2d68f464/YM/Maxwell/Wick.lean\#L107}{\texttt{Wick.monomial}},
\href{https://github.com/yang-mills-group/OpenAI-YM-20261028/blob/3456e5e75cf8d4cd558955a5c7638fba2d68f464/YM/Maxwell/Equation.lean\#L47-L50}{\texttt{reduceEq}},
\href{https://github.com/yang-mills-group/OpenAI-YM-20261028/blob/3456e5e75cf8d4cd558955a5c7638fba2d68f464/YM/Maxwell/Equation.lean\#L66-L67}{\texttt{observable}},
\href{https://github.com/yang-mills-group/OpenAI-YM-20261028/blob/3456e5e75cf8d4cd558955a5c7638fba2d68f464/YM/Maxwell/Equation.lean\#L72}{\texttt{equationIdeal}},
and the links of III.16.10.

\paragraph{4.9 Labels}\label{labels}

\textbf{Lemma 4.26 (= III.16.9).}

\emph{Sketch.} - (a) A complex polynomial that vanishes on
\(\mathbb R^N\) is zero. - (b) A coordinate's exponent in a monomial of
weight \(\le d\) is at most \(d\), and only finitely many coordinates
have weight \(\le d\). - (c) Bihomogeneous components of an invariant
polynomial are invariant, because the adjoint action is linear in each
\(A_j\) and preserves bidegree. They are real if the label is real. -
(d) Real bihomogeneous labels span \(\mathrm{Label}_d\) over
\(\mathbb C\), by taking real and imaginary parts of its bihomogeneous
components. A basis containing \(1\) can be extracted, and the bases are
extended inductively from level \(d\) to level \(d+1\). - Each frame
gives coordinates on \(\mathrm{Label}_d\). For every real \(t\),
rescaling the coordinate of \(e_i\) by \(t^{\deg e_i}\) (polynomial
degree) commutes with the action of \(O(4)\), since that action
preserves each polynomial degree:
\href{https://github.com/yang-mills-group/OpenAI-YM-20261028/blob/3456e5e75cf8d4cd558955a5c7638fba2d68f464/YM/Labels/Reference.lean\#L58-L61}{\texttt{Frame.scaled}},
\href{https://github.com/yang-mills-group/OpenAI-YM-20261028/blob/3456e5e75cf8d4cd558955a5c7638fba2d68f464/YM/Labels/Reference.lean\#L96-L109}{\texttt{Frame.scaling}}.

\emph{Lean:}
\href{https://github.com/yang-mills-group/OpenAI-YM-20261028/blob/3456e5e75cf8d4cd558955a5c7638fba2d68f464/YM/Labels/Jet.lean\#L22-L37}{\texttt{radial}},
\href{https://github.com/yang-mills-group/OpenAI-YM-20261028/blob/3456e5e75cf8d4cd558955a5c7638fba2d68f464/YM/Labels/Polynomial.lean\#L38-L49}{\texttt{eval\_ext}},
\href{https://github.com/yang-mills-group/OpenAI-YM-20261028/blob/3456e5e75cf8d4cd558955a5c7638fba2d68f464/YM/Labels/Polynomial.lean\#L63-L94}{\texttt{cap\_finite}},
\href{https://github.com/yang-mills-group/OpenAI-YM-20261028/blob/3456e5e75cf8d4cd558955a5c7638fba2d68f464/YM/Labels/Grading.lean\#L35-L48}{\texttt{component\_mem}},
\href{https://github.com/yang-mills-group/OpenAI-YM-20261028/blob/3456e5e75cf8d4cd558955a5c7638fba2d68f464/YM/Labels/Grading.lean\#L87-L105}{\texttt{decompose}},
\href{https://github.com/yang-mills-group/OpenAI-YM-20261028/blob/3456e5e75cf8d4cd558955a5c7638fba2d68f464/YM/Labels/Basis.lean\#L83-L86}{\texttt{frame\_exists}},
\href{https://github.com/yang-mills-group/OpenAI-YM-20261028/blob/3456e5e75cf8d4cd558955a5c7638fba2d68f464/YM/Labels/Basis.lean\#L107-L111}{\texttt{frames}},
\href{https://github.com/yang-mills-group/OpenAI-YM-20261028/blob/3456e5e75cf8d4cd558955a5c7638fba2d68f464/YM/Labels/Action.lean\#L82-L88}{\texttt{rotation\_law}}.

\subsubsection{5. Where the difficulty
lies}\label{where-the-difficulty-lies-18}

The analytic issue in III.16.3--III.16.4 is to retain a bound
independent of matrix dimension while varying the exponents and passing
to integral kernels. The three-lines argument uses dimension only for
qualitative boundedness on the strip. Arbitrary insertions are allowed:
their Schatten norms supply the entire insertion cost. Positivity and a
positive cyclic trace are required of the transfer kernel, and
quadrature is taken at fixed exponents and period.

The complex-power convention matters at zero: \(A^0\) is the range
projection. The strict inequalities \(Y_k>2r\) ensure that the integer
endpoints agree with ordinary positive powers. Fourier inversion loses
strip width and the explicit factor \((2(1+b)/b)^4\); differentiating a
Laurent sum costs \(1/(\eta-\sigma)\). The Bloch identities retain all
momenta, and at complex momentum the alias identity uses the adjoint
evaluated at \(\bar p\).

For Maxwell jets, the conventions are Taylor coefficients \(D^j/j!\),
real orthonormal colour coordinates, the normalization
\((4\pi^2)^{-1}\), and the sign \((-1)^l\) at the second endpoint. The
equation subspace is defined by restriction to shell jets; no
presentation by invariant generators is needed. The chosen label bases
are coherent as sets of polynomials, rather than as prescribed ordered
lists.

\subsubsection{6. Notation}\label{notation-18}

{\def\LTcaptype{none} % do not increment counter
\begin{longtable}[]{@{}
  >{\raggedright\arraybackslash}p{(\linewidth - 4\tabcolsep) * \real{0.3333}}
  >{\raggedright\arraybackslash}p{(\linewidth - 4\tabcolsep) * \real{0.3333}}
  >{\raggedright\arraybackslash}p{(\linewidth - 4\tabcolsep) * \real{0.3333}}@{}}
\toprule\noalign{}
\begin{minipage}[b]{\linewidth}\raggedright
Written proof
\end{minipage} & \begin{minipage}[b]{\linewidth}\raggedright
Lean
\end{minipage} & \begin{minipage}[b]{\linewidth}\raggedright
Meaning
\end{minipage} \\
\midrule\noalign{}
\endhead
\bottomrule\noalign{}
\endlastfoot
positive-semidefinite kernel &
\href{https://github.com/yang-mills-group/OpenAI-YM-20261028/blob/3456e5e75cf8d4cd558955a5c7638fba2d68f464/YM/Analysis/Nonnegative.lean\#L18-L19}{\texttt{NNKernel}}
& all finite Gram matrices Hermitian and positive semidefinite \\
tagged partitions \(Q_k\), masses \(m_i\) &
\href{https://github.com/yang-mills-group/OpenAI-YM-20261028/blob/3456e5e75cf8d4cd558955a5c7638fba2d68f464/YM/Analysis/Quadrature.lean\#L22-L26}{\texttt{Quad}},
\href{https://github.com/yang-mills-group/OpenAI-YM-20261028/blob/3456e5e75cf8d4cd558955a5c7638fba2d68f464/YM/Analysis/Quadrature.lean\#L65-L66}{\texttt{Quad.shrink}}
& finite measurable partition with tags; mesh \(1/(k+1)\) \\
\(\lVert b\rVert_{2q}\) &
\href{https://github.com/yang-mills-group/OpenAI-YM-20261028/blob/3456e5e75cf8d4cd558955a5c7638fba2d68f464/YM/Complex/Gaps.lean\#L106}{\texttt{paysNorm}},
\href{https://github.com/yang-mills-group/OpenAI-YM-20261028/blob/3456e5e75cf8d4cd558955a5c7638fba2d68f464/YM/Complex/Gaps.lean\#L105}{\texttt{rate}}
\(=1/(2q)\) & Schatten norm via \(\mathrm{tr}((b^*b)^q)\) \\
strip path &
\href{https://github.com/yang-mills-group/OpenAI-YM-20261028/blob/3456e5e75cf8d4cd558955a5c7638fba2d68f464/YM/Complex/Holder.lean\#L25-L29}{\texttt{Flex}}
& entire; contraction on the strip; trace norm \(\le1\) on
\(\mathrm{Re}=1\) \\
\(A^z\) &
\href{https://github.com/yang-mills-group/OpenAI-YM-20261028/blob/3456e5e75cf8d4cd558955a5c7638fba2d68f464/YM/Complex/Gaps.lean\#L22-L24}{\texttt{leap}},
\href{https://github.com/yang-mills-group/OpenAI-YM-20261028/blob/3456e5e75cf8d4cd558955a5c7638fba2d68f464/YM/Complex/Columns.lean\#L38}{\texttt{specPow}}
& complex power, \(0\) on \(\ker A\) \\
\(Z_A=\mathrm{tr}A^M\) &
\href{https://github.com/yang-mills-group/OpenAI-YM-20261028/blob/3456e5e75cf8d4cd558955a5c7638fba2d68f464/YM/Complex/Gaps.lean\#L37}{\texttt{Holder.Q}}
& absolute trace, equal to the trace for \(A\succeq0\) \\
sides of III.16.2 &
\href{https://github.com/yang-mills-group/OpenAI-YM-20261028/blob/3456e5e75cf8d4cd558955a5c7638fba2d68f464/YM/Complex/Gaps.lean\#L108-L109}{\texttt{total}},
\href{https://github.com/yang-mills-group/OpenAI-YM-20261028/blob/3456e5e75cf8d4cd558955a5c7638fba2d68f464/YM/Complex/Gaps.lean\#L111-L112}{\texttt{cost}}
& trace product and Hölder bound \\
\(q,w,M,\ell,K,Z,c_k,Y'\) & \texttt{p}, \texttt{width}, \texttt{N},
\texttt{m}, \texttt{T}, \texttt{R}, \texttt{K}, \texttt{Z} & Lean
parameters of \texttt{finite\_lip}, \texttt{passage} \\
\(\mathrm{Tr}(k_1\cdots k_L)\) &
\href{https://github.com/yang-mills-group/OpenAI-YM-20261028/blob/3456e5e75cf8d4cd558955a5c7638fba2d68f464/YM/TimeKernel/Trace.lean\#L76}{\texttt{TimeKernel.closed}}
& cyclic integral \\
\(k^Q\), \(k^*\) &
\href{https://github.com/yang-mills-group/OpenAI-YM-20261028/blob/3456e5e75cf8d4cd558955a5c7638fba2d68f464/YM/TimeKernel/Trace.lean\#L81}{\texttt{sample}},
\href{https://github.com/yang-mills-group/OpenAI-YM-20261028/blob/3456e5e75cf8d4cd558955a5c7638fba2d68f464/YM/TimeKernel/Trace.lean\#L83}{\texttt{op}}
& sampled matrix; adjoint kernel \\
\(B_1K^{Y_1}\cdots\); \((B^*B)^q\) &
\href{https://github.com/yang-mills-group/OpenAI-YM-20261028/blob/3456e5e75cf8d4cd558955a5c7638fba2d68f464/YM/TimeKernel/Unit.lean\#L17-L18}{\texttt{along}},
\href{https://github.com/yang-mills-group/OpenAI-YM-20261028/blob/3456e5e75cf8d4cd558955a5c7638fba2d68f464/YM/TimeKernel/Unit.lean\#L20}{\texttt{multiply}}
& gap word; column word \\
\(S_\eta\); \(e^{ip\cdot x}\); \(\lvert x\rvert_1\) &
\href{https://github.com/yang-mills-group/OpenAI-YM-20261028/blob/3456e5e75cf8d4cd558955a5c7638fba2d68f464/YM/Laurent/Eval.lean\#L43}{\texttt{strip}},
\href{https://github.com/yang-mills-group/OpenAI-YM-20261028/blob/3456e5e75cf8d4cd558955a5c7638fba2d68f464/YM/Laurent/Eval.lean\#L39}{\texttt{mode}},
\href{https://github.com/yang-mills-group/OpenAI-YM-20261028/blob/3456e5e75cf8d4cd558955a5c7638fba2d68f464/YM/Laurent/Eval.lean\#L34}{\texttt{size}}
& strip, Fourier character and lattice one-norm \\
\(\sum e^{\eta\lvert x\rvert_1}\lvert f(x)\rvert\le W\) &
\href{https://github.com/yang-mills-group/OpenAI-YM-20261028/blob/3456e5e75cf8d4cd558955a5c7638fba2d68f464/YM/Laurent/Eval.lean\#L104}{\texttt{Width}}
& weighted norm bound \\
\(\hat f\) &
\href{https://github.com/yang-mills-group/OpenAI-YM-20261028/blob/3456e5e75cf8d4cd558955a5c7638fba2d68f464/YM/Laurent/Eval.lean\#L123}{\texttt{Laurent.read}},
\href{https://github.com/yang-mills-group/OpenAI-YM-20261028/blob/3456e5e75cf8d4cd558955a5c7638fba2d68f464/YM/CellFourier/Circulant.lean\#L30-L31}{\texttt{Row.read}}
& symbol (Laurent series) \\
periodic holomorphic \(u\) on \(S_\eta\); \(\hat u_p\) &
\href{https://github.com/yang-mills-group/OpenAI-YM-20261028/blob/3456e5e75cf8d4cd558955a5c7638fba2d68f464/YM/Laurent/Rune.lean\#L18-L21}{\texttt{Vane}},
\href{https://github.com/yang-mills-group/OpenAI-YM-20261028/blob/3456e5e75cf8d4cd558955a5c7638fba2d68f464/YM/Laurent/Rune.lean\#L80}{\texttt{market}}
& strip hypotheses and shifted Fourier coefficients \\
kernel \(R\); periodization \(R_n\) &
\href{https://github.com/yang-mills-group/OpenAI-YM-20261028/blob/3456e5e75cf8d4cd558955a5c7638fba2d68f464/YM/CellFourier/Circulant.lean\#L22}{\texttt{Bloch.Row}},
\href{https://github.com/yang-mills-group/OpenAI-YM-20261028/blob/3456e5e75cf8d4cd558955a5c7638fba2d68f464/YM/CellFourier/Circulant.lean\#L33-L34}{\texttt{Row.acts}},
\href{https://github.com/yang-mills-group/OpenAI-YM-20261028/blob/3456e5e75cf8d4cd558955a5c7638fba2d68f464/YM/CellFourier/Circulant.lean\#L26}{\texttt{Row.strip}}
& real kernel over \(\mathbb Z^4\) with internal indices \\
\(k_p\); \(F_n\) &
\href{https://github.com/yang-mills-group/OpenAI-YM-20261028/blob/3456e5e75cf8d4cd558955a5c7638fba2d68f464/YM/CellFourier/Mesh.lean\#L18}{\texttt{mon}},
\href{https://github.com/yang-mills-group/OpenAI-YM-20261028/blob/3456e5e75cf8d4cd558955a5c7638fba2d68f464/YM/CellFourier/Mesh.lean\#L20}{\texttt{cell}},
\href{https://github.com/yang-mills-group/OpenAI-YM-20261028/blob/3456e5e75cf8d4cd558955a5c7638fba2d68f464/YM/CellFourier/Mesh.lean\#L113}{\texttt{F}}
& torus momentum, character and Fourier matrix \\
alias momenta; \(E_L(p)\) &
\href{https://github.com/yang-mills-group/OpenAI-YM-20261028/blob/3456e5e75cf8d4cd558955a5c7638fba2d68f464/YM/CellFourier/Alias.lean\#L20}{\texttt{aliasAt}},
\href{https://github.com/yang-mills-group/OpenAI-YM-20261028/blob/3456e5e75cf8d4cd558955a5c7638fba2d68f464/YM/CellFourier/Alias.lean\#L30-L31}{\texttt{aliasBlock}}
& fine momenta over one coarse momentum and their Fourier block \\
\(\Phi_m\); \(v\le X\le m\) &
\href{https://github.com/yang-mills-group/OpenAI-YM-20261028/blob/3456e5e75cf8d4cd558955a5c7638fba2d68f464/YM/Work/F07/MatrixRoot.lean\#L125}{\texttt{binRoot}},
\href{https://github.com/yang-mills-group/OpenAI-YM-20261028/blob/3456e5e75cf8d4cd558955a5c7638fba2d68f464/YM/Work/F07/MatrixRoot.lean\#L21-L24}{\texttt{Band}}
& analytic root and quadratic-form bounds \\
\(c_jf(x)=D^jf(x)/j!\) &
\href{https://github.com/yang-mills-group/OpenAI-YM-20261028/blob/3456e5e75cf8d4cd558955a5c7638fba2d68f464/YM/Analysis/Homogeneous.lean\#L65-L66}{\texttt{Tensor.coeff}}
& Taylor coefficient \\
\(\partial_i\) &
\href{https://github.com/yang-mills-group/OpenAI-YM-20261028/blob/3456e5e75cf8d4cd558955a5c7638fba2d68f464/YM/Maxwell/Exterior.lean\#L20}{\texttt{Maxwell.grad}}
& partial derivative along axis \(i\) \\
\(\mathcal A_j\) &
\href{https://github.com/yang-mills-group/OpenAI-YM-20261028/blob/3456e5e75cf8d4cd558955a5c7638fba2d68f464/YM/Labels/Jet.lean\#L39}{\texttt{Labels.Taylor}}
(\href{https://github.com/yang-mills-group/OpenAI-YM-20261028/blob/3456e5e75cf8d4cd558955a5c7638fba2d68f464/YM/Labels/Jet.lean\#L22-L37}{\texttt{radial}})
& radial-gauge potentials, degree \(j+1\), weight \(j+2\) \\
\(\mathcal C_j\); \(\mathcal S_j\) &
\href{https://github.com/yang-mills-group/OpenAI-YM-20261028/blob/3456e5e75cf8d4cd558955a5c7638fba2d68f464/YM/Maxwell/Exterior.lean\#L63-L95}{\texttt{formSpace}}
(\href{https://github.com/yang-mills-group/OpenAI-YM-20261028/blob/3456e5e75cf8d4cd558955a5c7638fba2d68f464/YM/Maxwell/Exterior.lean\#L57-L61}{\texttt{ClosedJet}});
\href{https://github.com/yang-mills-group/OpenAI-YM-20261028/blob/3456e5e75cf8d4cd558955a5c7638fba2d68f464/YM/Maxwell/Shell.lean\#L44}{\texttt{Shell}}
& closed curvature jets; Maxwell shell \\
curl and its inverse &
\href{https://github.com/yang-mills-group/OpenAI-YM-20261028/blob/3456e5e75cf8d4cd558955a5c7638fba2d68f464/YM/Maxwell/Exterior.lean\#L97-L98}{\texttt{takeCurl}},
\href{https://github.com/yang-mills-group/OpenAI-YM-20261028/blob/3456e5e75cf8d4cd558955a5c7638fba2d68f464/YM/Maxwell/Exterior.lean\#L130-L154}{\texttt{contracted}}
& homogeneous radial-gauge identification \\
\(G_0\); \(C_{ab,km}\) &
\href{https://github.com/yang-mills-group/OpenAI-YM-20261028/blob/3456e5e75cf8d4cd558955a5c7638fba2d68f464/YM/Maxwell/PotentialKernel.lean\#L22}{\texttt{Green}},
\href{https://github.com/yang-mills-group/OpenAI-YM-20261028/blob/3456e5e75cf8d4cd558955a5c7638fba2d68f464/YM/Maxwell/Curls.lean\#L20}{\texttt{bareCurl}}
& \(G_0=(4\pi^2\lvert x\rvert^2)^{-1}\); free field-strength two-point
function \\
probe; \(\kappa(p,q)\) &
\href{https://github.com/yang-mills-group/OpenAI-YM-20261028/blob/3456e5e75cf8d4cd558955a5c7638fba2d68f464/YM/Maxwell/Leg.lean\#L64}{\texttt{Probe}},
\href{https://github.com/yang-mills-group/OpenAI-YM-20261028/blob/3456e5e75cf8d4cd558955a5c7638fba2d68f464/YM/Maxwell/Leg.lean\#L161-L162}{\texttt{legKernel}}
& linear form on \(\mathcal S_j\); one-contraction kernel \\
free moment; free correlation &
\href{https://github.com/yang-mills-group/OpenAI-YM-20261028/blob/3456e5e75cf8d4cd558955a5c7638fba2d68f464/YM/Maxwell/Wick.lean\#L126-L127}{\texttt{Wick.moment}},
\href{https://github.com/yang-mills-group/OpenAI-YM-20261028/blob/3456e5e75cf8d4cd558955a5c7638fba2d68f464/YM/Maxwell/Wick.lean\#L145-L147}{\texttt{Wick.correlation}}
& Wick pairing sum and its restriction to labels \\
equation ideal &
\href{https://github.com/yang-mills-group/OpenAI-YM-20261028/blob/3456e5e75cf8d4cd558955a5c7638fba2d68f464/YM/Maxwell/Equation.lean\#L72}{\texttt{equationIdeal}}
& labels vanishing on all on-shell jets \\
\(\mathrm{Pol}\); \(\mathrm{Label}_d\); weight &
\href{https://github.com/yang-mills-group/OpenAI-YM-20261028/blob/3456e5e75cf8d4cd558955a5c7638fba2d68f464/YM/Labels/Polynomial.lean\#L29}{\texttt{Labels.Pol}},
\href{https://github.com/yang-mills-group/OpenAI-YM-20261028/blob/3456e5e75cf8d4cd558955a5c7638fba2d68f464/YM/Labels/Polynomial.lean\#L112}{\texttt{Labels.Label}},
\href{https://github.com/yang-mills-group/OpenAI-YM-20261028/blob/3456e5e75cf8d4cd558955a5c7638fba2d68f464/YM/Labels/Polynomial.lean\#L55}{\texttt{height}}
& polynomial algebra, invariant filtered subspace and engineering
weight \\
frame &
\href{https://github.com/yang-mills-group/OpenAI-YM-20261028/blob/3456e5e75cf8d4cd558955a5c7638fba2d68f464/YM/Labels/Basis.lean\#L50-L57}{\texttt{Frame}},
\href{https://github.com/yang-mills-group/OpenAI-YM-20261028/blob/3456e5e75cf8d4cd558955a5c7638fba2d68f464/YM/Labels/Basis.lean\#L107-L111}{\texttt{frames}}
& complex basis of real bihomogeneous polynomials, with \(e_0=1\),
coherent in \(d\) \\
\end{longtable}
}

\subsubsection{7. What remains
unexplained}\label{what-remains-unexplained-18}

\begin{enumerate}
\def\labelenumi{\arabic{enumi}.}
\tightlist
\item
  \textbf{Additional lemmas and dependency lists.} The elementary helper
  lemmas assigned to this unit outside the eight folders have not been
  checked individually. The longer unit lists in §2 have not been
  independently verified entry by entry and should be treated as guides
  to placement.
\item
  \textbf{Norms and parameter choices.} The square-root argument uses
  the matrix operator norm; its entrywise Laurent bounds require a
  finite-dimensional norm conversion. In III.6 the strip is narrowed to
  \(\tau=\min(\eta/2,v/(2CD))\), with \(D=2A/\eta\) and \(C\) a
  norm-conversion constant, and then to \(\tau/3\). IV.1 chooses
  \(q,w,r\) and sums successive insertion moves. Those choices must
  accompany any claim of uniformity in volume or spacing.
\item
  \textbf{Placement.} The trace estimate belongs with IV.1, Laurent and
  Bloch analysis with III.2--III.7, labels with II.1, and Maxwell jets
  and Wick correlations with VI.1--VI.2 and VII.1--VII.3. The full
  square-root theorem is assembled in III.6.
\end{enumerate}

\subsection{IV.1. Infinite volume: extraction of a subsequence and the
Euclidean axioms of the
limit}\label{iv.1.-infinite-volume-extraction-of-a-subsequence-and-the-euclidean-axioms-of-the-limit}

\textbf{Part:} IV, infinite volume (chapter 10). \textbf{Lean:}
\href{https://github.com/yang-mills-group/OpenAI-YM-20261028/tree/3456e5e75cf8d4cd558955a5c7638fba2d68f464/YM/Extraction}{\texttt{YM/Extraction}},
\href{https://github.com/yang-mills-group/OpenAI-YM-20261028/tree/3456e5e75cf8d4cd558955a5c7638fba2d68f464/YM/Cells}{\texttt{YM/Cells}},
\href{https://github.com/yang-mills-group/OpenAI-YM-20261028/tree/3456e5e75cf8d4cd558955a5c7638fba2d68f464/YM/Slice}{\texttt{YM/Slice}},
\href{https://github.com/yang-mills-group/OpenAI-YM-20261028/tree/3456e5e75cf8d4cd558955a5c7638fba2d68f464/YM/Euclid}{\texttt{YM/Euclid}},
\href{https://github.com/yang-mills-group/OpenAI-YM-20261028/tree/3456e5e75cf8d4cd558955a5c7638fba2d68f464/YM/Dash}{\texttt{YM/Dash}},
\href{https://github.com/yang-mills-group/OpenAI-YM-20261028/tree/3456e5e75cf8d4cd558955a5c7638fba2d68f464/YM/Next/EXTRACTION}{\texttt{YM/Next/EXTRACTION}},
\href{https://github.com/yang-mills-group/OpenAI-YM-20261028/tree/3456e5e75cf8d4cd558955a5c7638fba2d68f464/YM/Lane}{\texttt{YM/Lane}},
and the proofs of the work items
\href{https://github.com/yang-mills-group/OpenAI-YM-20261028/blob/3456e5e75cf8d4cd558955a5c7638fba2d68f464/YM/Work/D01.lean}{D01},
\href{https://github.com/yang-mills-group/OpenAI-YM-20261028/blob/3456e5e75cf8d4cd558955a5c7638fba2d68f464/YM/Work/D02.lean}{D02},
\href{https://github.com/yang-mills-group/OpenAI-YM-20261028/blob/3456e5e75cf8d4cd558955a5c7638fba2d68f464/YM/Work/D03.lean}{D03},
\href{https://github.com/yang-mills-group/OpenAI-YM-20261028/blob/3456e5e75cf8d4cd558955a5c7638fba2d68f464/YM/Work/D04.lean}{D04},
\href{https://github.com/yang-mills-group/OpenAI-YM-20261028/blob/3456e5e75cf8d4cd558955a5c7638fba2d68f464/YM/Work/D05.lean}{D05},
\href{https://github.com/yang-mills-group/OpenAI-YM-20261028/blob/3456e5e75cf8d4cd558955a5c7638fba2d68f464/YM/Work/D06.lean}{D06}.
Their statements are in
\href{https://github.com/yang-mills-group/OpenAI-YM-20261028/blob/3456e5e75cf8d4cd558955a5c7638fba2d68f464/YM/Bid/Drops.lean}{\texttt{YM/Bid/Drops.lean}},
which is filed with III.10. \textbf{Size:} 7,444 lines in 1,049
declarations used by the main theorem. Other material in these folders
concerns stress tensors, Maxwell comparison and label quotients (Parts
V--VII), or corrected fields (III.11--III.12).

\subsubsection{1. What this unit proves}\label{what-this-unit-proves-19}

\textbf{Setting and notation.} Let \(G\) be a compact matrix gauge group
and \(C\) a Wilson family (II.2, III.13). It has:

\begin{itemize}
\item
  lattice spacings \(a_N\downarrow0\), with \(a_0=1\),
  \(a_{N+1}=a_N/b_N\) and \(b_N\ge3\) odd;
\item
  periods \(L_p\uparrow\infty\), each divisible by \(4\);
\item
  couplings \(\beta_{N,p}\);
\item
  the periodic lattice
  \(\Lambda_{N,p}=(\mathbb Z/(L_p/a_N)\mathbb Z)^4\), whose sites are
  placed in the box \((-L_p/2,L_p/2]^4\);
\item
  the normalized Wilson expectation \(\langle\cdot\rangle_{N,p}\).
\end{itemize}

\textbf{Labels.} Labels at degree \(d\) (polynomial weight at most
\(d\)) form a finite-dimensional complex space \(\mathrm{Label}_d\),
given by coefficients in a real basis, so that \(v\mapsto\bar v\) is
defined. We use the supremum norm on coefficient vectors.
\(\mathrm{Label}_d\) carries:

\begin{itemize}
\item
  a representation \(\rho_d\) of the group \(H_4\) of signed coordinate
  permutations;
\item
  a unit label \(1_d\);
\item
  for each axis \(a\) the reflected label
  \(\theta_av=\rho_d(s_a)\bar v\), where \(s_a\) reverses axis \(a\).
\end{itemize}

\textbf{Observables and correlations.} A label \(v\in\mathrm{Label}_d\)
has a renormalized lattice observable \(I_{N,p}(v,r)\) at each site
\(r\), linear in \(v\). Covariance, reality and the unit identity are
required for \(N\ge\mathrm{ref}(d)\). Write \(\widehat r\) for the
physical position of a site in \((-L_p/2,L_p/2]^4\). For a label list
\(v=(v_1,\dots,v_n)\),

\[\mathcal C_{N,p}(v;r)=\Big\langle\prod_{i=1}^nI_{N,p}(v_i,r_i)\Big\rangle_{N,p},\qquad \mathcal C_{N,p}(v;t)=a_N^{4n}\sum_{r\in\Lambda_{N,p}^n}t(\widehat r)\,\mathcal C_{N,p}(v;r).\]

For a lifted tuple \(r\in(\mathbb Z^4)^n\), \(\mathcal C_{N,p}(v;r)\)
means evaluation after reduction modulo \(L_p/a_N\); its lifted physical
position is \(a_Nr\).

\textbf{Configuration sets.}

\begin{itemize}
\item
  \(\mathrm{Conf}_n\): configurations with \(x_i\ne x_j\) for
  \(i\ne j\).
\item
  \(\mathrm{Gen}_n\): configurations with \(x_i^a\neq x_j^a\) for all
  \(i\ne j\) and all axes \(a\). This set is open and dense, and its
  complement is Lebesgue-null.
\item
  On the torus: \(\mathrm{Conf}^{(p)}_n\) requires that for each
  \(i\ne j\) some \(x_i^a-x_j^a\notin L_p\mathbb Z\);
  \(\mathrm{Gen}^{(p)}_n\) requires all
  \(x_i^a-x_j^a\notin L_p\mathbb Z\).
\end{itemize}

\textbf{Test functions.} \(\mathcal D_n=C_c^\infty(\mathrm{Conf}_n)\). A
test function \(t\) is \emph{admissible at volume \(p\)} if
\(\mathrm{supp}t\subset(-L_p/2,L_p/2)^{4n}\).

\textbf{Extractions.} An \emph{extraction} \(f\) is a pair of strictly
increasing sequences \((N_j)\) and \((p_k)\); a refinement \(f\circ g\)
passes to subsequences of both. A property of \((N_j,p_k)\) holds
\emph{eventually along ordered tails} if it holds for \(k\ge k_0\) and
then \(j\ge j_0(k)\). Convergence along ordered tails means that, for
each neighbourhood of the limit, membership holds eventually in this
sense. The iterated limits in IV.1.1 (c) also require an actual spacing
limit at every admissible fixed volume.

\textbf{Tables.} A \emph{table} \(S\) assigns to each arity \(n\),
degrees \(d_1,\dots,d_n\) and point \(x\) a multilinear form
\(S(\cdot\,;x)\) on
\(\mathrm{Label}_{d_1}\times\dots\times\mathrm{Label}_{d_n}\). Its
values off \(\mathrm{Conf}_n\) play no role.

\textbf{Standing hypotheses.}

\begin{itemize}
\item
  \textbf{(H1)} The family satisfies its specification (§2, input 1).
\item
  \textbf{(H2)} The fixed-volume centered budgets, the power budgets
  along ordered tails, reflection positivity and the chessboard
  inequalities hold along \(f\), with the quantifiers in §2, input 2.
\end{itemize}

\textbf{Proposition IV.1.1 (subsequential raw limits).} Assume (H1) and
(H2) for an extraction \(f\). Then there are a refinement
\(f\circ g=((N_j),(p_k))\), a table \(S\) and tables \(P_p\)
(\(p\in\mathbb N\)) such that the following hold for every \(n\) and
every list of labels \(v\).

\begin{enumerate}
\def\labelenumi{(\alph{enumi})}
\tightlist
\item
  \emph{(Spacing limit at fixed volume.)} For every \(p\),
  \(P_p(v;\cdot)\) is continuous on \(\mathrm{Gen}^{(p)}_n\). For every
  compact \(K\subset\mathrm{Gen}^{(p)}_n\) and \(\varepsilon>0\), for
  all sufficiently large \(j\) and all lifted tuples
  \(r\in(\mathbb Z^4)^n\) with \(a_{N_j}r\in K\),
\end{enumerate}

\[|\mathcal C_{N_j,p}(v;r)-P_p(v;a_{N_j}r)|<\varepsilon.\]

\begin{enumerate}
\def\labelenumi{(\alph{enumi})}
\setcounter{enumi}{1}
\item
  \emph{(Volume limit.)} \(S(v;\cdot)\) is continuous on
  \(\mathrm{Gen}_n\), and \(P_{p_k}(v;\cdot)\to S(v;\cdot)\) uniformly
  on compact subsets of \(\mathrm{Gen}_n\).
\item
  \emph{(Ordered smeared limit.)} For every \(t\in\mathcal D_n\):
\end{enumerate}

\begin{itemize}
\item
  \(t\) is admissible at \(p_k\) for \(k\) large;
\item
  for every admissible \(p\) the limit
  \(V_p=\lim_j\mathcal C_{N_j,p}(v;t)\) exists (in the proof,
  \(V_p=\int tP_p(v;\cdot)\));
\item
  \(V_{p_k}\to\int t\,S(v;\cdot)\).
\end{itemize}

\begin{enumerate}
\def\labelenumi{(\alph{enumi})}
\setcounter{enumi}{3}
\tightlist
\item
  \(t\,S(v;\cdot)\) is integrable for every \(t\in\mathcal D_n\),
  i.e.~\(S\) is locally integrable on \(\mathrm{Conf}_n\).
\end{enumerate}

\emph{Status:} theorem. (H2) holds for every \(f\) as soon as the
volume-uniform budget of §2, input 2, holds
(\href{https://github.com/yang-mills-group/OpenAI-YM-20261028/blob/3456e5e75cf8d4cd558955a5c7638fba2d68f464/YM/Extraction/Center.lean\#L42-L56}{\texttt{Wcmain.projects}}).
Part (d) uses the uniform bound of Lemma IV.1.d.

\emph{Lean:}

\begin{itemize}
\item
  existence:
  \href{https://github.com/yang-mills-group/OpenAI-YM-20261028/blob/3456e5e75cf8d4cd558955a5c7638fba2d68f464/YM/Extraction/Held.lean\#L136-L137}{\texttt{Extraction.crops}};
\item
  conclusion:
  \href{https://github.com/yang-mills-group/OpenAI-YM-20261028/blob/3456e5e75cf8d4cd558955a5c7638fba2d68f464/YM/Extraction/Lipschitz.lean\#L156-L160}{\texttt{RawLimits}},
  with fields
  \href{https://github.com/yang-mills-group/OpenAI-YM-20261028/blob/3456e5e75cf8d4cd558955a5c7638fba2d68f464/YM/Extraction/Limits.lean\#L57-L62}{\texttt{Order.Meshes}},
  \href{https://github.com/yang-mills-group/OpenAI-YM-20261028/blob/3456e5e75cf8d4cd558955a5c7638fba2d68f464/YM/Extraction/Limits.lean\#L64-L66}{\texttt{Order.VolumePoints}},
  \href{https://github.com/yang-mills-group/OpenAI-YM-20261028/blob/3456e5e75cf8d4cd558955a5c7638fba2d68f464/YM/Extraction/Limits.lean\#L32-L35}{\texttt{Order.Ordered}}
  and
  \href{https://github.com/yang-mills-group/OpenAI-YM-20261028/blob/3456e5e75cf8d4cd558955a5c7638fba2d68f464/YM/Extraction/Orders.lean\#L61-L64}{\texttt{Order.limits}};
\item
  part (d):
  \href{https://github.com/yang-mills-group/OpenAI-YM-20261028/blob/3456e5e75cf8d4cd558955a5c7638fba2d68f464/YM/Extraction/Integrate.lean\#L115-L153}{\texttt{Order.ordinary}};
\item
  hypothesis (H2):
  \href{https://github.com/yang-mills-group/OpenAI-YM-20261028/blob/3456e5e75cf8d4cd558955a5c7638fba2d68f464/YM/Extraction/Center.lean\#L30-L34}{\texttt{CenterPort}}.
\end{itemize}

\textbf{Proposition IV.1.2 (limits of label derivatives).} Let \(S\) be
an ordered limit along \(f\) (IV.1.1 (c)) that is locally integrable
(IV.1.1 (d)), with (H1) and (H2) along \(f\). Let \(\mathrm{Null}_d(S)\)
be the space of \(u\in\mathrm{Label}_d\) with
\(\int tS(\dots,u,\dots)=0\) for every arity \(n\), slot, spectator
label list and \(t\in\mathcal D_n\). Assume the relay data of §2, input
3:

\begin{itemize}
\item
  projections \(\pi_d(N,p)\) on \(\mathrm{Label}_d\) converging along
  ordered tails in operator norm to a projection \(\pi_d\);
\item
  the relay properties for every \(d\);
\item
  the dimension count \(\dim\mathrm{Null}_d(S)\le\dim\ker\pi_d\).
\end{itemize}

Then \(\ker\pi_d=\mathrm{Null}_d(S)\). Moreover there are linear maps
\(\kappa_{d,e}:\mathrm{Label}_d\to\mathrm{Label}_e\) (\(d\le e\)) and
\(\tau_{d,a}:\mathrm{Label}_d\to\mathrm{Label}_{d+1}\) (\(a\) an axis)
such that, for all \(n\), slots \(i\), labels \(v\) and
\(t\in\mathcal D_n\),

\[\int t\,S(v_{i\to\kappa v_i};\cdot)=\int t\,S(v;\cdot),\qquad \int t\,S(v_{i\to\tau_av_i};\cdot)=-\int(\partial_{x_i^a}t)\,S(v;\cdot).\]

Here \(v_{i\to w}\) is \(v\) with the \(i\)-th label replaced by \(w\).
In other words,
\(\partial_{x_i^a}S(v;\cdot)=S(v_{i\to\tau_av_i};\cdot)\) in
\(\mathcal D'(\mathrm{Conf}_n)\).

\emph{Status:} theorem, conditional on the relay data and the dimension
count. Their suppliers in III.10--III.11 are listed in §2, input 3.

\emph{Lean:}
\href{https://github.com/yang-mills-group/OpenAI-YM-20261028/blob/3456e5e75cf8d4cd558955a5c7638fba2d68f464/YM/Work/D04.lean\#L212-L240}{\texttt{Work.D04.answer}}
(conclusion
\href{https://github.com/yang-mills-group/OpenAI-YM-20261028/blob/3456e5e75cf8d4cd558955a5c7638fba2d68f464/YM/Dash/Inert.lean\#L55-L60}{\texttt{Spans}}),
\href{https://github.com/yang-mills-group/OpenAI-YM-20261028/blob/3456e5e75cf8d4cd558955a5c7638fba2d68f464/YM/Dash/Inert.lean\#L41-L42}{\texttt{Bin.exactC}},
\href{https://github.com/yang-mills-group/OpenAI-YM-20261028/blob/3456e5e75cf8d4cd558955a5c7638fba2d68f464/YM/Dash/Inert.lean\#L71-L87}{\texttt{lifts\_model}}
(maps
\href{https://github.com/yang-mills-group/OpenAI-YM-20261028/blob/3456e5e75cf8d4cd558955a5c7638fba2d68f464/YM/Dash/Inert.lean\#L63-L65}{\texttt{MonitorOps}},
identities
\href{https://github.com/yang-mills-group/OpenAI-YM-20261028/blob/3456e5e75cf8d4cd558955a5c7638fba2d68f464/YM/Dash/Inert.lean\#L67-L69}{\texttt{MonitorOps.Wakes}}),
\href{https://github.com/yang-mills-group/OpenAI-YM-20261028/blob/3456e5e75cf8d4cd558955a5c7638fba2d68f464/YM/Dash/Inert.lean\#L89-L96}{\texttt{Wakes.through\_ties}}.

\textbf{Proposition IV.1.3 (smoothness of the limit).} Let \(S\) be as
in IV.1.1, along an extraction for which IV.1.2 applies. Then there is a
table \(T\) with the following properties.

\begin{itemize}
\item
  \(T=S\) on \(\mathrm{Gen}_n\), and
  \(T(v;\cdot)\in C^\infty(\mathrm{Conf}_n)\) for every \(n\) and \(v\).
\item
  IV.1.1 (a)--(d) hold for \(T\) along the same extraction and with the
  same \(P_p\).
\item
  \(T\) satisfies the basic identities (E1) below.
\item
  The maps \(\kappa\) and \(\tau\) descend to linear maps
  \(\mathrm{rise}_{d,e}:F_d\to F_e\) and
  \(\partial_{d,a}:F_d\to F_{d+1}\) on the quotients
  \(F_d=\mathrm{Label}_d/\mathrm{Null}_d(T)\). They satisfy, pointwise
  on \(\mathrm{Conf}_n\),
\end{itemize}

\[T(\dots,\mathrm{rise}\,u,\dots)=T(\dots,u,\dots),\qquad T(\dots,\partial_au_i,\dots)=\partial_{x_i^a}T(\dots,u_i,\dots).\]

Such a pair of maps is a \emph{derivative plan} with its inclusion and
derivative identities.

On \(\mathrm{Conf}_n\), the smooth ordered limit along a given
extraction is unique
(\href{https://github.com/yang-mills-group/OpenAI-YM-20261028/blob/3456e5e75cf8d4cd558955a5c7638fba2d68f464/YM/Extraction/Limits.lean\#L47-L50}{\texttt{Order.determined}}).

\emph{Status:} theorem.

\emph{Lean:}
\href{https://github.com/yang-mills-group/OpenAI-YM-20261028/blob/3456e5e75cf8d4cd558955a5c7638fba2d68f464/YM/Org/Input.lean\#L101-L109}{\texttt{Org.Arrives.serves}}
(filed with VIII.1; hypotheses
\href{https://github.com/yang-mills-group/OpenAI-YM-20261028/blob/3456e5e75cf8d4cd558955a5c7638fba2d68f464/YM/Org/Input.lean\#L83-L86}{\texttt{Org.Arrives}},
conclusion
\href{https://github.com/yang-mills-group/OpenAI-YM-20261028/blob/3456e5e75cf8d4cd558955a5c7638fba2d68f464/YM/Org/Input.lean\#L90-L96}{\texttt{Org.Cleaner}}),
\href{https://github.com/yang-mills-group/OpenAI-YM-20261028/blob/3456e5e75cf8d4cd558955a5c7638fba2d68f464/YM/Dash/Slew.lean\#L26-L58}{\texttt{Dash.awake}},
\href{https://github.com/yang-mills-group/OpenAI-YM-20261028/blob/3456e5e75cf8d4cd558955a5c7638fba2d68f464/YM/Extraction/Lipschitz.lean\#L194-L202}{\texttt{RawLimits.ready}},
\href{https://github.com/yang-mills-group/OpenAI-YM-20261028/blob/3456e5e75cf8d4cd558955a5c7638fba2d68f464/YM/Euclid/Differential.lean\#L20-L22}{\texttt{DiffPlan}}
with
\href{https://github.com/yang-mills-group/OpenAI-YM-20261028/blob/3456e5e75cf8d4cd558955a5c7638fba2d68f464/YM/Euclid/Differential.lean\#L58-L60}{\texttt{Inclusions}}
and
\href{https://github.com/yang-mills-group/OpenAI-YM-20261028/blob/3456e5e75cf8d4cd558955a5c7638fba2d68f464/YM/Euclid/Differential.lean\#L61-L64}{\texttt{Derivatives}}.

\textbf{Proposition IV.1.4 (the Euclidean hierarchy).} Let \(T\) be a
table and let \(f\) be an extraction. Assume (H1), IV.1.1 (a)--(c) hold
for \(T\) and \(T(v;\cdot)\in C^\infty(\mathrm{Conf}_n)\) for every
arity \(n\) and label list \(v\). Then \(T\) satisfies (E1)--(E5), with
all points in \(\mathrm{Conf}_n\):

\begin{itemize}
\item
  \textbf{(E1) Basic identities.} \(T\) is invariant under simultaneous
  permutation of points and labels; \(T=1\) for \(n=0\); and
  \(T((1_d,v);x)=T(v;(x_2,\dots,x_{n+1}))\).
\item
  \textbf{(E2) Real spin.} \(T(\rho(s)v;sx)=T(v;x)\) for \(s\in H_4\),
  and \(T(\bar v;x)=\overline{T(v;x)}\).
\item
  \textbf{(E3) Homogeneity.} \(T(v;x+s)=T(v;x)\) for every
  \(s\in\mathbb R^4\).
\item
  \textbf{(E4) Reflection positivity at points.} Fix an axis \(a\) and
  words \(w_1,\dots,w_m\), each a finite list of (point, label) pairs
  with distinct points and all \(x^a>0\). Let \(\Theta_aw\) be the word
  in reverse order with \(x^a\mapsto-x^a\) and labels \(\theta_a\). Then
  for all \(c\in\mathbb C^m\) the following expression is real and
  nonnegative:
\end{itemize}

\[\sum_{i,j}\bar c_ic_j\,T(\Theta_aw_i\,w_j)\ge0.\]

\begin{itemize}
\tightlist
\item
  \textbf{(E5) Smeared reflection positivity.} The same holds for tests:
  replace \(w_i\) by pairs (labels \(v^{(i)}\), test
  \(t_i\in\mathcal D_{n_i}\) supported in \(\{x_k^a>0\ \forall k\}\))
  and \(T(\Theta_aw_iw_j)\) by
  \(\iint\overline{t_i(y)}\,t_j(z)\,T(\Theta_a(v^{(i)},y)\,(v^{(j)},z))\).
  These integrals converge absolutely, and the Gram matrix is positive
  semidefinite.
\end{itemize}

If in addition \textbf{(E7)} (polynomial bounds on derivatives) holds,
then \(T\) satisfies the hierarchy (E1)--(E7), where:

\begin{itemize}
\item
  \textbf{(E6)} is smoothness on \(\mathrm{Conf}_n\);
\item
  \textbf{(E7)} says: for every arity \(n\), degree list, label list
  \(v\) and order \(j\ge0\), there are \(A\in\mathbb N\) and \(C\ge0\)
  such that \(\|D^j_xT(v;x)\|\le Cr^{-A}\) whenever \(0<r\le1\) and
  \(|x_i-x_\ell|_\infty\ge r\) for every \(i\ne\ell\). Here \(D^j\) is
  the iterated real Fréchet derivative with its operator norm on
  \((\mathbb R^4)^n\). The constants are independent of \(r\) and \(x\).
\end{itemize}

\emph{Status:} theorem. (E6) and (E7) are hypotheses of the assembly.
(E6) is supplied by IV.1.3; (E7) is supplied by IV.1.5 together with the
promotion lemma B10 of V.1.

\emph{Lean:}

\begin{itemize}
\item
  assembly:
  \href{https://github.com/yang-mills-group/OpenAI-YM-20261028/blob/3456e5e75cf8d4cd558955a5c7638fba2d68f464/YM/Extraction/Harvest.lean\#L62-L65}{\texttt{Order.assembles}},
  producing
  \href{https://github.com/yang-mills-group/OpenAI-YM-20261028/blob/3456e5e75cf8d4cd558955a5c7638fba2d68f464/YM/Euclid/Words.lean\#L141-L148}{\texttt{Table.Hierarchy}};
\item
  fields:
  \href{https://github.com/yang-mills-group/OpenAI-YM-20261028/blob/3456e5e75cf8d4cd558955a5c7638fba2d68f464/YM/Euclid/Table.lean\#L75-L78}{\texttt{Basic}},
  \href{https://github.com/yang-mills-group/OpenAI-YM-20261028/blob/3456e5e75cf8d4cd558955a5c7638fba2d68f464/YM/Euclid/Table.lean\#L80-L84}{\texttt{RealSpin}},
  \href{https://github.com/yang-mills-group/OpenAI-YM-20261028/blob/3456e5e75cf8d4cd558955a5c7638fba2d68f464/YM/Euclid/Table.lean\#L86}{\texttt{Homogeneous}},
  \href{https://github.com/yang-mills-group/OpenAI-YM-20261028/blob/3456e5e75cf8d4cd558955a5c7638fba2d68f464/YM/Euclid/Words.lean\#L123-L124}{\texttt{RP}},
  \href{https://github.com/yang-mills-group/OpenAI-YM-20261028/blob/3456e5e75cf8d4cd558955a5c7638fba2d68f464/YM/Euclid/Words.lean\#L136-L138}{\texttt{RPSmear}},
  \href{https://github.com/yang-mills-group/OpenAI-YM-20261028/blob/3456e5e75cf8d4cd558955a5c7638fba2d68f464/YM/Euclid/Table.lean\#L88-L89}{\texttt{Smooth}},
  \href{https://github.com/yang-mills-group/OpenAI-YM-20261028/blob/3456e5e75cf8d4cd558955a5c7638fba2d68f464/YM/Euclid/Table.lean\#L97-L99}{\texttt{Capped}};
\item
  call site:
  \href{https://github.com/yang-mills-group/OpenAI-YM-20261028/blob/3456e5e75cf8d4cd558955a5c7638fba2d68f464/YM/Org/Input.lean\#L113-L114}{\texttt{Org.Cleaner.stay}}.
\end{itemize}

\textbf{Proposition IV.1.5 (polynomial bound at short distance; D05).}
Assume (H1) and (H2) along \(f\). Let \(T(v;\cdot)\) be smooth on
\(\mathrm{Conf}_n\) for every arity and label list, and suppose \(T\)
satisfies IV.1.1 (a)--(b) along \(f\). Then for every arity \(n\),
degree list \((d_i)\) and label list \(v\) there is \(C\ge0\),
independent of \(r\) and \(x\), such that

\[|T(v;x)|\le C\,r^{-A}\quad\text{whenever }0<r\le1\text{ and }|x_i-x_\ell|_\infty\ge r\text{ for every }i\ne\ell,\qquad A=4n\big(1+\max_id_i\big).\]

The maximum of an empty degree list is taken to be \(0\). The displayed
exponent follows by choosing the largest degree in the proof.

\emph{Status:} theorem. Together with the derivative plan of IV.1.3, the
lemma B10 of V.1 upgrades this to (E7).

\emph{Lean:}
\href{https://github.com/yang-mills-group/OpenAI-YM-20261028/blob/3456e5e75cf8d4cd558955a5c7638fba2d68f464/YM/Work/D05.lean\#L160-L180}{\texttt{Work.D05.answer}}
(statement
\href{https://github.com/yang-mills-group/OpenAI-YM-20261028/blob/3456e5e75cf8d4cd558955a5c7638fba2d68f464/YM/Bid/Drops.lean\#L70-L71}{\texttt{Bid.D05}},
conclusion
\href{https://github.com/yang-mills-group/OpenAI-YM-20261028/blob/3456e5e75cf8d4cd558955a5c7638fba2d68f464/YM/Bid/Phase.lean\#L97-L98}{\texttt{Bid.OneCap}});
promotion
\href{https://github.com/yang-mills-group/OpenAI-YM-20261028/blob/3456e5e75cf8d4cd558955a5c7638fba2d68f464/YM/Bid/Phase.lean\#L101-L102}{\texttt{Bid.B10}},
proved by
\href{https://github.com/yang-mills-group/OpenAI-YM-20261028/blob/3456e5e75cf8d4cd558955a5c7638fba2d68f464/YM/Work/B10.lean\#L95-L105}{\texttt{Work.B10.answer}}.

\textbf{Proposition IV.1.6 (the derivative plan respects the symmetries;
D06).} Let \(T\) be smooth on \(\mathrm{Conf}_n\) for every arity and
label list, permutation symmetric and with real spin, and let
\((\mathrm{rise},\partial)\) be a derivative plan for \(T\) satisfying
its inclusion and derivative identities. Use \(\rho_d\) also for the
induced action on \(F_d\). For every \(d\le e\), axis and \(s\in H_4\):

\begin{itemize}
\item
  \(\mathrm{rise}_{d,e}\) and every \(\partial_a\) commute with complex
  conjugation;
\item
  \(\mathrm{rise}\) commutes with \(\rho(s)\) for every \(s\in H_4\);
\item
  \(\rho_{d+1}(s)\,\partial_b=\sum_a s_{ab}\,\partial_a\,\rho_d(s)\),
  where \((s_{ab})\) is the signed permutation matrix of \(s\).
\end{itemize}

\emph{Status:} theorem.

\emph{Lean:}
\href{https://github.com/yang-mills-group/OpenAI-YM-20261028/blob/3456e5e75cf8d4cd558955a5c7638fba2d68f464/YM/Work/D06.lean\#L98-L135}{\texttt{Work.D06.compat}},
\href{https://github.com/yang-mills-group/OpenAI-YM-20261028/blob/3456e5e75cf8d4cd558955a5c7638fba2d68f464/YM/Work/D06.lean\#L137-L139}{\texttt{Work.D06.answer}},
\href{https://github.com/yang-mills-group/OpenAI-YM-20261028/blob/3456e5e75cf8d4cd558955a5c7638fba2d68f464/YM/Euclid/Differential.lean\#L66-L72}{\texttt{DiffPlan.Orientation}}.

\textbf{Remarks.}

\begin{enumerate}
\def\labelenumi{\arabic{enumi}.}
\item
  \textbf{The package handed downstream} is the bundle
  \href{https://github.com/yang-mills-group/OpenAI-YM-20261028/blob/3456e5e75cf8d4cd558955a5c7638fba2d68f464/YM/Next/EXTRACTION/Base.lean\#L42-L51}{\texttt{Next.Extract.Base}}.
  It contains (H1), (H2), the relay data, the raw limit \(S\) with local
  integrability and the spanning property, the smooth table \(T\) with
  its plan, the bound of IV.1.5 and the hierarchy of IV.1.4. It also
  records \(\ker\pi_d=\mathrm{Null}_d(T)\)
  (\href{https://github.com/yang-mills-group/OpenAI-YM-20261028/blob/3456e5e75cf8d4cd558955a5c7638fba2d68f464/YM/Next/EXTRACTION/Base.lean\#L72-L74}{\texttt{Base.null}}).
  On the main theorem's path it is assembled in III.11
  (\href{https://github.com/yang-mills-group/OpenAI-YM-20261028/blob/3456e5e75cf8d4cd558955a5c7638fba2d68f464/YM/MergeProof/Powder.lean\#L48-L71}{\texttt{MergeProof.Accent.rated}})
  as follows:

  \begin{itemize}
  \item
    extract with
    \href{https://github.com/yang-mills-group/OpenAI-YM-20261028/blob/3456e5e75cf8d4cd558955a5c7638fba2d68f464/YM/Extraction/Held.lean\#L136-L137}{\texttt{crops}};
  \item
    refine further so that the projections converge
    (\href{https://github.com/yang-mills-group/OpenAI-YM-20261028/blob/3456e5e75cf8d4cd558955a5c7638fba2d68f464/YM/Next/EXTRACTION/Relay.lean\#L86-L93}{\texttt{include\_projections}});
  \item
    convert the relay estimates
    (\href{https://github.com/yang-mills-group/OpenAI-YM-20261028/blob/3456e5e75cf8d4cd558955a5c7638fba2d68f464/YM/Next/EXTRACTION/Relay.lean\#L178-L181}{\texttt{Projected.ready}}).
  \item
    include the rank-detector limits along a further common refinement
    and obtain the dimension comparison
    (\href{https://github.com/yang-mills-group/OpenAI-YM-20261028/blob/3456e5e75cf8d4cd558955a5c7638fba2d68f464/YM/Next/RANK/Beam.lean\#L153-L180}{\texttt{Beam.arriving\_seed}}),
    then apply IV.1.2--IV.1.5. All earlier limits are retained.
  \end{itemize}
\item
  \textbf{The extraction is subsequential in both indices.} It gives no
  independence from the volume subsequence. The public form of the
  theorem
  (\href{https://github.com/yang-mills-group/OpenAI-YM-20261028/blob/3456e5e75cf8d4cd558955a5c7638fba2d68f464/PublicWilson.lean\#L85-L90}{\texttt{WilsonFamily.Converges}},
  via
  \href{https://github.com/yang-mills-group/OpenAI-YM-20261028/blob/3456e5e75cf8d4cd558955a5c7638fba2d68f464/HierarchyBridge.lean\#L44-L56}{\texttt{wilsonLimit}})
  needs the ordered limit along the \emph{full} sequence of spacings.
  The upgrade is proved in III.13 by
  \href{https://github.com/yang-mills-group/OpenAI-YM-20261028/blob/3456e5e75cf8d4cd558955a5c7638fba2d68f464/YM/UVExtension/Reports/Ultraviolet.lean\#L62-L69}{\texttt{UVExtension.Reports.useFull}},
  using
  \href{https://github.com/yang-mills-group/OpenAI-YM-20261028/blob/3456e5e75cf8d4cd558955a5c7638fba2d68f464/YM/UVExtension/Reports/Ultraviolet.lean\#L50-L58}{\texttt{full\_tests}}.
  At each admissible fixed volume, uniqueness identifies the
  full-sequence smeared limit with its subsequential limit; the volume
  limit and the table \(T\) are retained.
\end{enumerate}

\subsubsection{2. What it uses}\label{what-it-uses-19}

\textbf{Inputs from other units.}

\begin{enumerate}
\def\labelenumi{\arabic{enumi}.}
\item
  \textbf{Specification of the family.} The definition is in II.1--II.2.
  It is proved for the constructed family in III.13
  (\href{https://github.com/yang-mills-group/OpenAI-YM-20261028/blob/3456e5e75cf8d4cd558955a5c7638fba2d68f464/YM/UVExtension/Forest/Glue.lean\#L35-L36}{\texttt{UVExtension.Mast.calibrated}}).
  Its consequences used here are the following; covariance, reality and
  the unit identity require \(N\ge\mathrm{ref}(d)\):

  \begin{itemize}
  \item
    \(\beta_{N,p}\ge0\);
  \item
    \(I_{N,p}(v,x)\) is continuous, gauge invariant and linear in \(v\);
  \item
    \(I_{N,p}(v,x)\) depends only on links inside a fixed finite stencil
    translated to \(x\); for each stencil offset \(z\) and axis \(a\),
    \(20(|z^a|+1)\le L_p/a_N\) whenever \(N\ge\mathrm{ref}(d)\);
  \item
    for \(N\ge\mathrm{ref}(d)\), the terminal transfer \(B_{d,N,p}\)
    from cutoff coordinates to reference coordinates is normalized as
    follows: the coefficient trajectory of \(v\) has
    \(x_{\mathrm{ref}(d)}=v\), \(x_N=B_{d,N,p}^{-1}v\), and
    \(\|x_{m+1}\|_1\le\mathrm{column}_{m+1}\|x_m\|_1\) for
    \(\mathrm{ref}(d)\le m<N\);
  \item
    translation covariance;
  \item
    \(H_4\)-covariance: \(I(\rho(s)v,sx+b)\) is \(I(v,x)\) composed with
    the pull-back of the lattice symmetry;
  \item
    reality: \(I(\bar v,x)=\overline{I(v,x)}\);
  \item
    \(I(1_d,x)=1\) for \(N\ge\mathrm{ref}(d)\). \emph{Lean:}
    \href{https://github.com/yang-mills-group/OpenAI-YM-20261028/blob/3456e5e75cf8d4cd558955a5c7638fba2d68f464/YM/Source/Construction.lean\#L44-L62}{\texttt{Family.Makes}},
    \href{https://github.com/yang-mills-group/OpenAI-YM-20261028/blob/3456e5e75cf8d4cd558955a5c7638fba2d68f464/YM/Source/Construction.lean\#L42}{\texttt{Family.O}},
    \href{https://github.com/yang-mills-group/OpenAI-YM-20261028/blob/3456e5e75cf8d4cd558955a5c7638fba2d68f464/YM/Source/Construction.lean\#L88-L98}{\texttt{rotated}},
    \href{https://github.com/yang-mills-group/OpenAI-YM-20261028/blob/3456e5e75cf8d4cd558955a5c7638fba2d68f464/YM/Source/Construction.lean\#L74-L86}{\texttt{reality}},
    \href{https://github.com/yang-mills-group/OpenAI-YM-20261028/blob/3456e5e75cf8d4cd558955a5c7638fba2d68f464/YM/Source/Construction.lean\#L100-L102}{\texttt{constant}}.
  \end{itemize}
\item
  \textbf{Centered chessboard budgets} (interface item 2). A sufficient
  volume-uniform input is the following: for every degree \(d\), every
  \(R>0\) and every \(\varepsilon>0\) there is \(A\ge0\) such that for
  every \(m\ge\mathrm{ref}(d)\) there is \(N_0\ge m\) with the following
  property. For all \(N\ge N_0\), all \(p\), all
  \(v\in\mathrm{Label}_d\) with \(\|v\|\le R\), and all sites \(x\):

  \begin{itemize}
  \item
    \(I_{N,p}(v,x)\) is an interior function of the base cell of the
    centered tiling \(\mathcal T_m(x)\);
  \item
    \(\|I_{N,p}(v,x)\|_{\mathcal T_m(x)}\le A\,a_m^{-d-\varepsilon}\).
  \end{itemize}

  This stronger volume-uniform input implies \textbf{(H2)} along every
  extraction \(f\). Precisely, (H2) requires:

  \begin{itemize}
  \item
    for every fixed \(p\) and label \(v\), at every sufficiently fine
    scale \(m\) there is \(A\ge0\) such that the centered interior
    condition and root bound
    \(\|I_{N_j,p}(v,x)\|_{\mathcal T_m(x)}\le A\) hold for all \(x\) and
    all large \(j\), with \(N_j\ge\max(m,\mathrm{ref}(d))\) for
    \(v\in\mathrm{Label}_d\);
  \item
    reflection positivity, the chessboard inequality and subadditivity
    of roots at every \((N,p)\); and, for every \(d\), \(R>0\) and
    \(\varepsilon>0\), numbers \(m_0\ge\mathrm{ref}(d)\) and \(A\ge0\)
    such that for every \(m\ge m_0\) and \(\eta>0\), eventually along
    ordered tails, the centered interior condition and
    \(\|I_{N,p}(v,x)\|_{\mathcal T_m(x)}\le A\,a_m^{-d-\varepsilon}+\eta\)
    hold for all \(\|v\|\le R\) and all sites \(x\), with \(N\ge m\).
  \end{itemize}

  On the main theorem's path it is supplied by
  \href{https://github.com/yang-mills-group/OpenAI-YM-20261028/blob/3456e5e75cf8d4cd558955a5c7638fba2d68f464/YM/MergeProof/Gardens.lean\#L19-L40}{\texttt{MergeProof.Accent.gardena}}
  (III.11). The variant
  \href{https://github.com/yang-mills-group/OpenAI-YM-20261028/blob/3456e5e75cf8d4cd558955a5c7638fba2d68f464/YM/Org/Input.lean\#L75-L76}{\texttt{Org.Assembly.wcmain}}
  is not on that path. \emph{Lean:}
  \href{https://github.com/yang-mills-group/OpenAI-YM-20261028/blob/3456e5e75cf8d4cd558955a5c7638fba2d68f464/YM/Extraction/Center.lean\#L38-L40}{\texttt{Wcmain}},
  \href{https://github.com/yang-mills-group/OpenAI-YM-20261028/blob/3456e5e75cf8d4cd558955a5c7638fba2d68f464/YM/Extraction/Center.lean\#L21-L23}{\texttt{centeredBudget}},
  \href{https://github.com/yang-mills-group/OpenAI-YM-20261028/blob/3456e5e75cf8d4cd558955a5c7638fba2d68f464/YM/Cells/Packing.lean\#L34-L40}{\texttt{Cells.center}},
  \href{https://github.com/yang-mills-group/OpenAI-YM-20261028/blob/3456e5e75cf8d4cd558955a5c7638fba2d68f464/YM/Cells/Packing.lean\#L68-L79}{\texttt{CellHyp}},
  \href{https://github.com/yang-mills-group/OpenAI-YM-20261028/blob/3456e5e75cf8d4cd558955a5c7638fba2d68f464/YM/Extraction/Center.lean\#L58-L66}{\texttt{CenterPort.book}},
  \href{https://github.com/yang-mills-group/OpenAI-YM-20261028/blob/3456e5e75cf8d4cd558955a5c7638fba2d68f464/YM/Extraction/Center.lean\#L25-L28}{\texttt{CenteredBook}}.
\item
  \textbf{Relay data} (interface item 3, in the form used here;
  III.10--III.11). A family \(Q_{N,p}(x)\) is \emph{chessboard-small} if
  there is a real rate \(y(N,p)\to0\) along ordered tails with the
  following property. For every integer \(q\ge0\) and lower scale \(m\),
  there is \(k\ge m\) such that, eventually along ordered tails with
  \(N\ge k\), the bound \(\|Q_{N,p}(x)\|_{\mathcal T}\le y(N,p)\) and
  the interior condition hold for every translated bond tiling
  \(\mathcal T\) of side \(a_k\) and every \(x\) in its base cell with
  lattice coordinates \(z^a\) relative to the base corner satisfying

  \[M\le z^a,\qquad z^a+M+1\le a_k/a_N,\qquad M=\Big\lfloor\frac{a_k/a_N}{8(q+1)}\Big\rfloor.\]
  Thus the error bound must hold for translated cells with prescribed
  relative interior margins. Centered cells alone do not suffice.

  A label family \(w(N,p)\) has \emph{subpolynomial growth} if, for
  every \(\varepsilon>0\),

  \[\|B_{d,N,p}^{-1}w(N,p)\|\le a_N^{-\varepsilon}\] eventually along
  ordered tails. For each degree \(d\) the relay data are:

  \begin{itemize}
  \item
    projections \(\pi_d(N,p)\) on \(\mathrm{Label}_d\), bounded at each
    fixed \(p\) and along ordered tails;
  \item
    \emph{pace}: for \(v\in\mathrm{Label}_d\) and each axis \(a\) there
    are degree-\((d+1)\) labels \(w(N,p)\) of subpolynomial growth such
    that
    \(I_{N,p}(w(N,p),\cdot)-a_N^{-1}\big[I_{N,p}(v,\cdot+e_a)-I_{N,p}(v,\cdot)\big]\)
    is chessboard-small;
  \end{itemize}

  Here \(e_a\) is one lattice step, of physical length \(a_N\).

  \begin{itemize}
  \item
    \emph{transit}: the same, with a label of degree \(d'\le d\)
    approximated at degree \(d\);
  \item
    \emph{mark}: if \(w(N,p)\in\ker\pi_d(N,p)\) has subpolynomial
    growth, then \(I_{N,p}(w(N,p),\cdot)\) is chessboard-close to
    measurable observables, bounded at each cutoff and site, whose
    correlations with arbitrary spectator labels vanish exactly at every
    fixed positive separation, eventually along ordered tails;
  \item
    the dimension count \(\dim\mathrm{Null}_d\le\dim\ker\pi_d\).
  \end{itemize}

  Supplied in III.11 by
  \href{https://github.com/yang-mills-group/OpenAI-YM-20261028/blob/3456e5e75cf8d4cd558955a5c7638fba2d68f464/YM/MergeProof/Eureka.lean\#L16-L20}{\texttt{MergeProof.Scotia}},
  \href{https://github.com/yang-mills-group/OpenAI-YM-20261028/blob/3456e5e75cf8d4cd558955a5c7638fba2d68f464/YM/MergeProof/Powder.lean\#L19-L47}{\texttt{savings}}
  and
  \href{https://github.com/yang-mills-group/OpenAI-YM-20261028/blob/3456e5e75cf8d4cd558955a5c7638fba2d68f464/YM/Next/RANK/Beam.lean\#L153-L180}{\texttt{Beam.arriving\_seed}}.
  \emph{Lean:}
  \href{https://github.com/yang-mills-group/OpenAI-YM-20261028/blob/3456e5e75cf8d4cd558955a5c7638fba2d68f464/YM/Next/EXTRACTION/Relay.lean\#L39-L41}{\texttt{Presentation}},
  \href{https://github.com/yang-mills-group/OpenAI-YM-20261028/blob/3456e5e75cf8d4cd558955a5c7638fba2d68f464/YM/Next/EXTRACTION/Relay.lean\#L46-L48}{\texttt{Affordable}},
  \href{https://github.com/yang-mills-group/OpenAI-YM-20261028/blob/3456e5e75cf8d4cd558955a5c7638fba2d68f464/YM/Next/EXTRACTION/Relay.lean\#L115-L119}{\texttt{RelayCosts}},
  \href{https://github.com/yang-mills-group/OpenAI-YM-20261028/blob/3456e5e75cf8d4cd558955a5c7638fba2d68f464/YM/Bid/Drops.lean\#L46-L52}{\texttt{Bid.RelayBin}},
  \href{https://github.com/yang-mills-group/OpenAI-YM-20261028/blob/3456e5e75cf8d4cd558955a5c7638fba2d68f464/YM/Bid/Drops.lean\#L39}{\texttt{vet}},
  \href{https://github.com/yang-mills-group/OpenAI-YM-20261028/blob/3456e5e75cf8d4cd558955a5c7638fba2d68f464/YM/Bid/Drops.lean\#L25}{\texttt{close}},
  \href{https://github.com/yang-mills-group/OpenAI-YM-20261028/blob/3456e5e75cf8d4cd558955a5c7638fba2d68f464/YM/Bid/Drops.lean\#L18-L19}{\texttt{dollar}},
  \href{https://github.com/yang-mills-group/OpenAI-YM-20261028/blob/3456e5e75cf8d4cd558955a5c7638fba2d68f464/YM/Bid/Drops.lean\#L43-L44}{\texttt{heighten}},
  \href{https://github.com/yang-mills-group/OpenAI-YM-20261028/blob/3456e5e75cf8d4cd558955a5c7638fba2d68f464/YM/Dash/Inert.lean\#L39}{\texttt{Bin.shortage}}.
\item
  \textbf{Analytic toolkit (III.16).} The Lipschitz bound for traces of
  products of a positive kernel (Lemma IV.1.e). \emph{Lean:}
  \href{https://github.com/yang-mills-group/OpenAI-YM-20261028/blob/3456e5e75cf8d4cd558955a5c7638fba2d68f464/YM/TimeKernel/Unit.lean\#L76-L106}{\texttt{TimeKernel.passage}},
  \href{https://github.com/yang-mills-group/OpenAI-YM-20261028/blob/3456e5e75cf8d4cd558955a5c7638fba2d68f464/YM/Complex/Gaps.lean\#L170-L229}{\texttt{Holder.finite\_lip}},
  \href{https://github.com/yang-mills-group/OpenAI-YM-20261028/blob/3456e5e75cf8d4cd558955a5c7638fba2d68f464/YM/Complex/Gaps.lean\#L124-L165}{\texttt{Holder.estimation}},
  \href{https://github.com/yang-mills-group/OpenAI-YM-20261028/blob/3456e5e75cf8d4cd558955a5c7638fba2d68f464/YM/Complex/Holder.lean\#L130-L132}{\texttt{Holder.convexStrip}},
  \href{https://github.com/yang-mills-group/OpenAI-YM-20261028/blob/3456e5e75cf8d4cd558955a5c7638fba2d68f464/YM/TimeKernel/Trace.lean\#L109-L119}{\texttt{TimeKernel.learns}}.
\item
  \textbf{Scale schedule (I.1, III.2, V.1).}

  \begin{itemize}
  \item
    Gentleness, \(\log b_m/\log a_m^{-1}\to0\), gives
    \(a_{m+1}\ge c\,a_m^2\)
    (\href{https://github.com/yang-mills-group/OpenAI-YM-20261028/blob/3456e5e75cf8d4cd558955a5c7638fba2d68f464/YM/Lattice/Scale.lean\#L129-L130}{\texttt{Schedule.Gentle}},
    \href{https://github.com/yang-mills-group/OpenAI-YM-20261028/blob/3456e5e75cf8d4cd558955a5c7638fba2d68f464/YM/Lattice/Scale.lean\#L133-L162}{\texttt{mesh\_inequality}},
    \href{https://github.com/yang-mills-group/OpenAI-YM-20261028/blob/3456e5e75cf8d4cd558955a5c7638fba2d68f464/YM/OS/Halves.lean\#L29-L39}{\texttt{OS.thickness}}).
    This is used in D05.
  \item
    The column constants of the renormalization group have
    \(\sum_{m\le N}\log(\text{column}_m)=o(\log a_N^{-1})\). This is
    used in D01.
  \end{itemize}
\item
  \textbf{Promotion of derivative bounds (V.1).} The lemma B10 (see
  IV.1.5).
\end{enumerate}

Interface item 4 (whole-sequence convergence at fixed volume) enters
only in Remark 2: the extraction itself passes to a spacing subsequence.
Interface item 5 (corrected fields) is not used either:
\href{https://github.com/yang-mills-group/OpenAI-YM-20261028/blob/3456e5e75cf8d4cd558955a5c7638fba2d68f464/YM/Cells/Packing.lean\#L68-L79}{\texttt{CellHyp}}
consists only of the chessboard inequality and the budget for
\(I_{N,p}\). The corrected-field definitions serve III.11--III.12.

\textbf{Standard mathematics.}

\begin{itemize}
\item
  \emph{Reflection positivity of the Wilson action, site and bond
  reflections.} K. Osterwalder, E. Seiler, Ann. Phys. 110 (1978)
  440--471. Axioms: Osterwalder--Schrader, CMP 31 (1973), 42 (1975).
\item
  \emph{Chessboard estimates.}

  \begin{itemize}
  \tightlist
  \item
    J. Glimm, A. Jaffe, T. Spencer, CMP 45 (1975) 203--216;
  \item
    J. Fröhlich, R. Israel, E. Lieb, B. Simon, CMP 62 (1978) 1--34 and
    J. Stat. Phys. 22 (1980) 297--347;
  \item
    M. Biskup, \emph{Reflection positivity and phase transitions in
    lattice spin models}, LNM 1970 (2009), §5.
  \end{itemize}
\item
  \emph{Transfer matrices of reflection-positive lattice theories.}
  Glimm--Jaffe, \emph{Quantum Physics}, 2nd ed.
\item
  \emph{Hölder's inequality for Schatten norms and complex interpolation
  (Hadamard three lines).} B. Simon, \emph{Trace Ideals and Their
  Applications}; Reed--Simon II, §IX.4. Cauchy's estimate for
  holomorphic functions.
\item
  \emph{Compactness and integration.} Arzelà--Ascoli; Tychonoff for
  countable products; dominated convergence.
\item
  \emph{Regularity.} A distribution whose localizations have rapidly
  decreasing Fourier transforms is \(C^\infty\) (Hörmander, \emph{The
  Analysis of Linear Partial Differential Operators I}, Ch. 7).
\end{itemize}

\textbf{Imported constants and hypotheses.}

\begin{itemize}
\item
  Numbers: the budget constants \(A(d,R,\varepsilon)\) and thresholds
  \(N_0(m,d,R,\varepsilon)\) (input 2); the reference levels
  \(\mathrm{ref}(d)\); the gentleness constant \(c\) of input 5. The
  exponent \(4n(1+\max d_i)\) of IV.1.5 comes from choosing
  \(\varepsilon=1\) and from applying \(a_{m+1}\ge ca_m^2\) twice.
\item
  Hypotheses: (H1), (H2), and the relay data and dimension count (input
  3). The promotion theorem B10 is proved in V.1.
\end{itemize}

\subsubsection{3. Strategy}\label{strategy-19}

Chessboard estimates bound correlations at distinct points uniformly in
the cutoff and volume. To obtain equicontinuity, fix an axis and express
a correlation as a cyclic trace of transfer kernels and observable
kernels supported in slabs. Moving one insertion transfers length
between two gaps. Finite-dimensional complex interpolation bounds the
trace for complex gap shifts; Cauchy's estimate gives a Lipschitz bound.
The required column moments are bounded by the same chessboard budgets.
This argument needs separation in every coordinate. Two applications of
Arzelà--Ascoli, with countable basis indices, give spacing limits at
every fixed volume and then a volume limit on \(\mathrm{Gen}_n\). Bounds
on all of \(\mathrm{Conf}_n\) supply domination for smeared limits.

Smoothness uses the label calculus. The relay estimates approximate each
lattice forward difference quotient by a degree-raised observable.
Summation by parts computes its weak limit, and chessboard-small errors
have vanishing smeared limits. The projections and dimension count
identify the null labels. A finite-dimensional compactness argument then
replaces the varying approximating labels by fixed linear maps, giving
distributional derivatives as kernels with raised labels. Iteration
yields rapid Fourier decay after localization, hence a smooth
representative on all distinct-point configurations.

The exact lattice symmetries pass first at generic grid points and then
by continuity. Reflection positivity passes for joint tests supported in
a half-space; normalized bump functions give its pointwise form. A
two-scale choice of cells yields a polynomial bound at short distances,
and the derivative identities promote it to bounds of every order. These
arguments assemble the Euclidean hierarchy.

\subsubsection{4. Steps}\label{steps-19}

\paragraph{4.1 Reflection positivity and chessboard
estimates}\label{reflection-positivity-and-chessboard-estimates}

\textbf{Lemma IV.1.a (reflection positivity of the Wilson measure).} Let
\(\Lambda=(\mathbb Z/n\mathbb Z)^4\) with \(n\) even. Let \(a\) be an
axis, \(b\in\mathbb Z/n\), and \(\Theta\) either the site reflection
\(x^a\mapsto2b-x^a\) or the bond reflection \(x^a\mapsto2b-1-x^a\)
(other coordinates fixed), acting on configurations by pull-back. Let
\(\Lambda_+\) be the half-torus \(0<(x^a-b)\bmod n<n/2\) (site case) or
\(0\le(x^a-b)\bmod n<n/2\) (bond case), with residues in
\(\{0,\dots,n-1\}\). For every \(\beta\ge0\) and every bounded,
measurable, gauge-invariant \(F\) depending only on links in
\(\Lambda_+\),

\[\langle F\,\overline{F\circ\Theta}\rangle_\beta\ge0.\]

\emph{Proof sketch.} This is the Osterwalder--Seiler argument.
Reparametrize configurations by the links in the two fixed planes, the
links crossing them, and the two open halves; this map preserves the
measure.

\begin{itemize}
\item
  \emph{Site reflections.} Given the plane links, the halves are
  independent and exchanged by \(\Theta\), so the expectation is a
  weighted integral of \(|\mathbb E(F\mid\text{plane links})|^2\) with
  nonnegative weights.
\item
  \emph{Bond reflections.} Integrating the crossing links gives for each
  plane the one-step kernel \(K_\beta(X,Y)\). It is positive
  semidefinite because \(\mathrm{Re}\,\mathrm{tr}(XY^*)\) is a sum of
  Gram kernels, and exponentials, products and averages of positive
  semidefinite kernels stay positive semidefinite (Schur). The product
  of the two boundary kernels is positive semidefinite, so its quadratic
  integral is nonnegative.
\end{itemize}

\emph{Lean:}
\href{https://github.com/yang-mills-group/OpenAI-YM-20261028/blob/3456e5e75cf8d4cd558955a5c7638fba2d68f464/YM/Slice/Reflection.lean\#L98-L114}{\texttt{Slice.reflected\_positive}},
\href{https://github.com/yang-mills-group/OpenAI-YM-20261028/blob/3456e5e75cf8d4cd558955a5c7638fba2d68f464/YM/Cells/Cut.lean\#L116-L117}{\texttt{Cells.Positive}},
\href{https://github.com/yang-mills-group/OpenAI-YM-20261028/blob/3456e5e75cf8d4cd558955a5c7638fba2d68f464/YM/Cells/Cut.lean\#L109-L113}{\texttt{Cells.Interior}},
\href{https://github.com/yang-mills-group/OpenAI-YM-20261028/blob/3456e5e75cf8d4cd558955a5c7638fba2d68f464/YM/Cells/Cut.lean\#L40-L45}{\texttt{Cells.Cut}},
\href{https://github.com/yang-mills-group/OpenAI-YM-20261028/blob/3456e5e75cf8d4cd558955a5c7638fba2d68f464/YM/Slice/Gate.lean\#L94-L102}{\texttt{Slice.portPos}},
\href{https://github.com/yang-mills-group/OpenAI-YM-20261028/blob/3456e5e75cf8d4cd558955a5c7638fba2d68f464/YM/Slice/Positive.lean\#L81-L91}{\texttt{Slice.step\_positive}}.

\textbf{Lemma IV.1.b (chessboard estimate).} Let \(\mathcal T\) be a
tiling of \(\Lambda\) into \(M^4\) congruent cells (\(M\) even), with
site or bond seams. For a cell \(c\) let \(\theta_c\) be the composition
of seam reflections carrying the base cell to \(c\). Put
\(\mathrm{copy}_cF=F\circ\theta_c\), or its complex conjugate when
\(\theta_c\) reverses an odd number of axes. For \(F\) interior to the
base cell (bounded, measurable, gauge invariant, depending only on links
inside it) put
\(\|F\|_{\mathcal T}=|\langle\prod_c\mathrm{copy}_cF\rangle_\beta|^{1/M^4}\).
If \(\beta\ge0\), then for interior \(F_c\)

\[\Big|\Big\langle\prod_c\mathrm{copy}_cF_c\Big\rangle_\beta\Big|\le\prod_c\|F_c\|_{\mathcal T},\]

and \(\|\cdot\|_{\mathcal T}\) is subadditive. By (H1) this holds for
every \((N,p)\) of the family.

\emph{Proof sketch.} This is the Fröhlich--Israel--Lieb--Simon argument.

\begin{itemize}
\item
  \emph{Fold.} For an assignment \(A\) of functions to cells, put
  \(\Phi(A)=|\langle\prod_c\mathrm{copy}_c A_c\rangle_\beta|\).
  Cauchy--Schwarz for a band cut gives
  \(\Phi(\mathrm{assign}(f,g))^2\le\Phi(\mathrm{assign}(f,f))\Phi(\mathrm{assign}(g,g))\),
  where the two assignments supply the two halves.
\item
  \emph{Maximization.} Among assignments of finitely many functions to
  cells, the maximum of \(|\langle\prod\rangle|\) is attained by a
  constant assignment, because a fold-closed property satisfied by some
  assignment is satisfied by a constant one.
\item
  \emph{Normalization.} Apply this to \(F_c/(\|F_c\|+\epsilon)\) and let
  \(\epsilon\to0\).
\item
  \emph{Subadditivity.} Expand \(\prod_c(F+F')\) into \(2^{M^4}\)
  assignments and apply the estimate to each.
\end{itemize}

\emph{Lean:}
\href{https://github.com/yang-mills-group/OpenAI-YM-20261028/blob/3456e5e75cf8d4cd558955a5c7638fba2d68f464/YM/Cells/Certificate.lean\#L14-L19}{\texttt{Cells.Seam}},
\href{https://github.com/yang-mills-group/OpenAI-YM-20261028/blob/3456e5e75cf8d4cd558955a5c7638fba2d68f464/YM/Cells/Certificate.lean\#L21-L24}{\texttt{Cells.laws}},
\href{https://github.com/yang-mills-group/OpenAI-YM-20261028/blob/3456e5e75cf8d4cd558955a5c7638fba2d68f464/YM/Cells/Chess.lean\#L80-L105}{\texttt{Tiling.tiles}},
\href{https://github.com/yang-mills-group/OpenAI-YM-20261028/blob/3456e5e75cf8d4cd558955a5c7638fba2d68f464/YM/Cells/Chess.lean\#L32-L63}{\texttt{Tiling.maximize}},
\href{https://github.com/yang-mills-group/OpenAI-YM-20261028/blob/3456e5e75cf8d4cd558955a5c7638fba2d68f464/YM/Cells/Polar.lean\#L94-L102}{\texttt{Tiling.fold\_both}},
\href{https://github.com/yang-mills-group/OpenAI-YM-20261028/blob/3456e5e75cf8d4cd558955a5c7638fba2d68f464/YM/Cells/Flood.lean\#L110-L128}{\texttt{Tiling.existence}},
\href{https://github.com/yang-mills-group/OpenAI-YM-20261028/blob/3456e5e75cf8d4cd558955a5c7638fba2d68f464/YM/Cells/Chess.lean\#L108-L131}{\texttt{Tiling.adds}},
\href{https://github.com/yang-mills-group/OpenAI-YM-20261028/blob/3456e5e75cf8d4cd558955a5c7638fba2d68f464/YM/Cells/Tiling.lean\#L226-L227}{\texttt{Tiling.rootTest}},
\href{https://github.com/yang-mills-group/OpenAI-YM-20261028/blob/3456e5e75cf8d4cd558955a5c7638fba2d68f464/YM/Cells/Tiling.lean\#L221-L224}{\texttt{Tiling.copy}},
\href{https://github.com/yang-mills-group/OpenAI-YM-20261028/blob/3456e5e75cf8d4cd558955a5c7638fba2d68f464/YM/Cells/Embedded.lean\#L25-L38}{\texttt{Tiling.planted}},
\href{https://github.com/yang-mills-group/OpenAI-YM-20261028/blob/3456e5e75cf8d4cd558955a5c7638fba2d68f464/YM/Cells/Packing.lean\#L81-L84}{\texttt{Cells.seamless}}.

\textbf{Lemma IV.1.c (roots of uncentered observables).} Assume (H1),
\(m\le k\le N\) and \(N\ge\mathrm{ref}(d)\). Let \(\mathcal X\) be a
bond tiling of \(\Lambda_{N,p}\) with cells of side \(a_m\), and let
\(u\in\mathrm{Label}_d\) have all stencil offsets of coordinate size
\(<a_k/a_N\), and assume that for every image \(gu\),
\(g\in H_4\times\{1,\text{conj}\}\), the centered roots satisfy
\(\|I(gu,y)\|_{\mathcal T_k(y)}\le A\) for all \(y\), with the centered
interior conditions. If \(o\) is at distance \(\ge3a_k\) from the
boundary of the base cell of \(\mathcal X\), then
\(\|I(u,o)\|_{\mathcal X}\le A\).

\emph{Proof sketch.}

\begin{itemize}
\item
  \emph{Centred case.} If \(o\) is the centre of a fine cell aligned
  with \(\mathcal X\), the \(\mathcal X\)-copies of \(I(u,o)\) are
  copies, in the fine centered tiling, of observables \(I(gu,y)\) (by
  \(H_4\)-covariance and reality) sitting in distinct fine cells. The
  chessboard estimate in \(\mathcal T_k\), with \(1\) in the remaining
  cells, gives
  \(\|I(u,o)\|_{\mathcal X}^{\sharp \mathcal X}\le A^{\sharp \mathcal X}\).
\item
  \emph{Log-convexity.} Under unit translations,
  \(\|F\|_{\mathcal X}^2\le\|F\circ\tau_{-e}\|_{\mathcal X}\,\|F\circ\tau_{e}\|_{\mathcal X}\).
  This follows from translation invariance and the chessboard estimate
  with alternating translates, since reversed cells see the opposite
  shift. A nonnegative log-convex sequence is bounded by its values at
  the endpoints.
\item
  \emph{Interpolation.} Move between fine-cell centres one axis at a
  time.
\end{itemize}

\emph{Lean:}
\href{https://github.com/yang-mills-group/OpenAI-YM-20261028/blob/3456e5e75cf8d4cd558955a5c7638fba2d68f464/YM/Extraction/Around.lean\#L95-L129}{\texttt{Around.interior\_from}},
\href{https://github.com/yang-mills-group/OpenAI-YM-20261028/blob/3456e5e75cf8d4cd558955a5c7638fba2d68f464/YM/Extraction/Around.lean\#L57-L93}{\texttt{Around.swept}},
\href{https://github.com/yang-mills-group/OpenAI-YM-20261028/blob/3456e5e75cf8d4cd558955a5c7638fba2d68f464/YM/Cells/Slide.lean\#L89-L101}{\texttt{Tiling.pathway}},
\href{https://github.com/yang-mills-group/OpenAI-YM-20261028/blob/3456e5e75cf8d4cd558955a5c7638fba2d68f464/YM/Cells/Slide.lean\#L60-L86}{\texttt{Tiling.trip}},
\href{https://github.com/yang-mills-group/OpenAI-YM-20261028/blob/3456e5e75cf8d4cd558955a5c7638fba2d68f464/YM/Extraction/Around.lean\#L42-L50}{\texttt{Spin.uses}}.

\paragraph{4.2 Boundedness at separated
points}\label{boundedness-at-separated-points}

\textbf{Lemma IV.1.d (uniform bound).} Assume (H1). Let
\(x_1,\dots,x_n\in\Lambda_{N,p}\) be such that every pair has some
coordinate at cyclic physical distance \(\ge r\). Let \(m\le k\le N\)
with \(a_m<r\) and \(6n\,a_k<a_m\), and let \(A\ge0\) bound the centered
roots at scale \(k\) of all images \(gv_i\),
\(g\in H_4\times\{1,\text{conj}\}\), with the centered interior
conditions. Assume \(N\ge\max_i\mathrm{ref}(d_i)\) and that all stencil
offsets of these labels have coordinate size \(<a_k/a_N\). Then
\(|\mathcal C_{N,p}(v;x)|\le A^n\). Under (H1)--(H2) this yields the
bound \textbf{(B)}:

\begin{itemize}
\item
  \textbf{(B1)} for each \(p\) and compact
  \(K\subset\mathrm{Conf}^{(p)}_n\), \(|\mathcal C_{N_j,p}|\) is bounded
  at lattice points of \(K\) for \(j\) large;
\item
  \textbf{(B2)} for compact \(K\subset\mathrm{Conf}_n\), it is bounded
  eventually along ordered tails.
\end{itemize}

\emph{Proof sketch.}

\begin{itemize}
\item
  By pigeonhole, some translate of the tiling with cells of side \(a_m\)
  puts every \(x_i\) at distance \(\ge3a_k\) from all cell boundaries
  (possible because \(n\cdot2\cdot3a_k<a_m\)).
\item
  Distinct points lie in distinct cells because \(a_m<r\).
\item
  Each \(I(v_i,x_i)\) is the copy of some \(I(g_iv_i,y_i)\) with \(y_i\)
  in the base cell.
\item
  Apply the chessboard estimate with \(1\) elsewhere, then Lemma IV.1.c.
\item
  For (B): by compactness a single \(r>0\) works on \(K\), and for large
  \(p\) there is no wrap-around. Choose \(m\), then \(k\), then take the
  budgets from (H2).
\end{itemize}

\emph{Lean:}
\href{https://github.com/yang-mills-group/OpenAI-YM-20261028/blob/3456e5e75cf8d4cd558955a5c7638fba2d68f464/YM/Extraction/Held.lean\#L73-L114}{\texttt{Around.fits\_budget}},
\href{https://github.com/yang-mills-group/OpenAI-YM-20261028/blob/3456e5e75cf8d4cd558955a5c7638fba2d68f464/YM/Extraction/Held.lean\#L120-L132}{\texttt{Around.size\_done}}
(conclusion
\href{https://github.com/yang-mills-group/OpenAI-YM-20261028/blob/3456e5e75cf8d4cd558955a5c7638fba2d68f464/YM/Extraction/Integrate.lean\#L25-L30}{\texttt{Bound}}),
\href{https://github.com/yang-mills-group/OpenAI-YM-20261028/blob/3456e5e75cf8d4cd558955a5c7638fba2d68f464/YM/Cells/Proximity.lean\#L53-L66}{\texttt{Tiling.fits\_together}},
\href{https://github.com/yang-mills-group/OpenAI-YM-20261028/blob/3456e5e75cf8d4cd558955a5c7638fba2d68f464/YM/Extraction/Gaps.lean\#L63-L80}{\texttt{Gaps.some}},
\href{https://github.com/yang-mills-group/OpenAI-YM-20261028/blob/3456e5e75cf8d4cd558955a5c7638fba2d68f464/YM/Extraction/Gaps.lean\#L129-L137}{\texttt{Gaps.someGrows}}.

\paragraph{4.3 Equicontinuity}\label{equicontinuity}

\textbf{Lemma IV.1.e (Lipschitz bound for traces; proved in III.16).}
Let \(X\) be a compact metric space with a finite Borel measure, \(K\) a
continuous positive semidefinite kernel, and \(B_1,\dots,B_n\)
continuous kernels, with \(n\ge2\). Let \(q,w\ge1\) be integers, put
\(M=2qw\), and assume \(Z=\mathrm{Tr}\,K^M>0\). Take gaps
\(Y,Y'\in\mathbb N^n\) with \(nw+\sum_kY_k=M\). Suppose \(Y'\) transfers
an integer \(\ell\ge0\) units from one gap to a different gap,
\(Y_k>2r\) for every \(k\), and \(0\le\ell\le r\) with \(r>0\). If

\[\|B_k\|_{2q}:=\bigl(\mathrm{Tr}(B_k^*B_k)^q\bigr)^{1/(2q)}\le c_kZ^{1/(2q)},\]

then

\[\Big|\mathrm{Tr}\prod_kB_kK^{Y'_k}-\mathrm{Tr}\prod_kB_kK^{Y_k}\Big|\le Z\Big(\prod_kc_k\Big)\frac{\ell}{r}.\]

Here traces of kernel products are their cyclic integrals.

\emph{Proof sketch.} Approximate the cyclic integrals by finite matrices
\((A,b_k)\) using one quadrature for all kernels and their adjoints. For
the positive semidefinite matrix \(A\), use the entire spectral family
equal to \(e^{z\log\lambda}\) on positive eigenvalues and zero on
\(\ker A\) for every \(z\). It agrees with \(A^m\) at positive integers
\(m\). The trace with gap vector \(W(z)=Y+z(e_i-e_j)\) is entire. On
\(|z|\le2r\), all real gap parts are nonnegative and

\[\frac{n}{2q}+\sum_k\frac{\mathrm{Re}\,W_k(z)}{M}=1.\]

Writing \(Z_A=\mathrm{Tr}\,A^M\), complex interpolation gives the bound
\(\prod_k\|b_k\|_{2q}\,Z_A^{1-n/(2q)}\). Cauchy's estimate bounds the
derivative on \(|z|\le r\) by this quantity divided by \(r\). Integrate
from \(0\) to \(\ell\) and pass to the kernel limit; the partition
traces and column moments converge under the common quadrature.

\emph{Lean:}
\href{https://github.com/yang-mills-group/OpenAI-YM-20261028/blob/3456e5e75cf8d4cd558955a5c7638fba2d68f464/YM/TimeKernel/Unit.lean\#L76-L106}{\texttt{TimeKernel.passage}},
\href{https://github.com/yang-mills-group/OpenAI-YM-20261028/blob/3456e5e75cf8d4cd558955a5c7638fba2d68f464/YM/Complex/Gaps.lean\#L170-L229}{\texttt{Holder.finite\_lip}},
\href{https://github.com/yang-mills-group/OpenAI-YM-20261028/blob/3456e5e75cf8d4cd558955a5c7638fba2d68f464/YM/Complex/Gaps.lean\#L124-L165}{\texttt{Holder.estimation}},
\href{https://github.com/yang-mills-group/OpenAI-YM-20261028/blob/3456e5e75cf8d4cd558955a5c7638fba2d68f464/YM/Complex/Gaps.lean\#L22-L24}{\texttt{Holder.leap}}.

\textbf{Lemma IV.1.f (the Wilson transfer matrix).} Fix \(\beta\ge0\),
an axis \(a\) and \(\Lambda_{N,p}\). The one-step kernel \(K_\beta\) on
slice configurations, with the links in direction \(a\) integrated out
and half of each slice action assigned to either side, is continuous and
positive semidefinite. With \(M=L_p/a_N\), \(\mathrm{Tr}\,K_\beta^M=Z\)
is the partition function. Every continuous observable local to a slab
of width \(1\le w<M\) lattice steps has a slab representation, whose
weighted integration gives a continuous kernel \(B\). Products of such
observables, all of width \(w\) and separated by gaps \(Y\), satisfy

\[\Big\langle\prod_kF_k\Big\rangle_{\beta}=\frac{\mathrm{Tr}\prod_kB_kK_\beta^{Y_k}}{Z},\qquad nw+\sum_kY_k=M.\]

If \(M=2qw\) with \(q\ge1\), the normalized column bounds
\(\bigl(\mathrm{Tr}(B_k^*B_k)^q/Z\bigr)^{1/(2q)}\le c_k\) therefore give
the difference bound \((\prod_kc_k)\ell/r\) under the gap and shift
hypotheses of Lemma IV.1.e.

\emph{Proof sketch.} Fill the complement of the slab with a fixed
configuration to express the observable as a function of its slab links.
Integrate the interior links with the Wilson weights. Fubini's theorem
identifies the product with the cyclic kernel integral. Apply Lemma
IV.1.e and divide by \(Z>0\).

\emph{Lean:}
\href{https://github.com/yang-mills-group/OpenAI-YM-20261028/blob/3456e5e75cf8d4cd558955a5c7638fba2d68f464/YM/Slice/Positive.lean\#L81-L91}{\texttt{Slice.step\_positive}},
\href{https://github.com/yang-mills-group/OpenAI-YM-20261028/blob/3456e5e75cf8d4cd558955a5c7638fba2d68f464/YM/Slice/Calendar.lean\#L34-L38}{\texttt{Time.partition\_seen}},
\href{https://github.com/yang-mills-group/OpenAI-YM-20261028/blob/3456e5e75cf8d4cd558955a5c7638fba2d68f464/YM/Slice/Calendar.lean\#L72-L102}{\texttt{Time.powering}},
\href{https://github.com/yang-mills-group/OpenAI-YM-20261028/blob/3456e5e75cf8d4cd558955a5c7638fba2d68f464/YM/Slice/Resident.lean\#L74-L87}{\texttt{Time.selects}},
\href{https://github.com/yang-mills-group/OpenAI-YM-20261028/blob/3456e5e75cf8d4cd558955a5c7638fba2d68f464/YM/Slice/Ordered.lean\#L68-L115}{\texttt{Time.organized}}.

\textbf{Lemma IV.1.g (chessboard bound for the column norm).} Assume
(H1), \(m\le N\) and \(N\ge\mathrm{ref}(d_v)\) for a label \(v\) of
degree \(d_v\). Put \(d=a_m/a_N\). Suppose \(F=I_{N,p}(v,x)\) and
\(J=I_{N,p}(\theta_av,x)\) are interior to the base cell of
\(\mathcal T_m(x)\) with roots \(\le c\). Then \(F\) has a slab
representation of width \(2d\) lattice steps. Its kernel \(B\) and the
kernel of \(J\) are adjoints. Writing \(L_p/a_N=4qd\), with \(q\ge1\),
and letting \(Z\) be the partition function, one has

\[\left(\frac{\mathrm{Tr}(B^*B)^q}{Z}\right)^{1/(2q)}\le c.\]

\emph{Proof sketch.} The observable \(J\) equals
\(\overline{F\circ\Theta}\) for the site reflection through \(x^a\).
Reversing the slab gives the adjoint kernel. The column consists of
\(2q\) alternating copies of \(J\) and \(F\), placed in every other cell
of the row of \(\mathcal T_m(x)\) in direction \(a\). Apply the
chessboard estimate with \(1\) in all other cells and take the
\(1/(2q)\)-th power.

\emph{Lean:}
\href{https://github.com/yang-mills-group/OpenAI-YM-20261028/blob/3456e5e75cf8d4cd558955a5c7638fba2d68f464/YM/Cells/Present.lean\#L76-L111}{\texttt{Tiling.enough}},
\href{https://github.com/yang-mills-group/OpenAI-YM-20261028/blob/3456e5e75cf8d4cd558955a5c7638fba2d68f464/YM/Extraction/Difference.lean\#L51-L87}{\texttt{Difference.reading}},
\href{https://github.com/yang-mills-group/OpenAI-YM-20261028/blob/3456e5e75cf8d4cd558955a5c7638fba2d68f464/YM/Extraction/Difference.lean\#L37-L49}{\texttt{Difference.stars}}.

\textbf{Lemma IV.1.h (Lipschitz bound in one coordinate).} Assume (H1),
\(m\le N\) and \(N\ge\max_i\mathrm{ref}(d_i)\). Let \(n\ge2\) lattice
points have pairwise cyclic physical \(a\)-gaps \(>2a_m+2\rho\), with
\(\rho>0\). Assume the centered interior conditions and let \(c_i\)
bound the roots of \(I(v_i,x_i)\) and \(I(\theta_av_i,x_i)\) on
\(\mathcal T_m(x_i)\). Moving one point by \(h\in a_N\mathbb N\),
\(0\le h\le\rho\), in direction \(+e_a\) changes
\(\mathcal C_{N,p}(v;x)\) by at most \((\prod_ic_i)\,h/\rho\).

\emph{Proof sketch.} Put \(d=a_m/a_N\) and reorder the points cyclically
in the \(a\)-direction (permutation symmetry). Represent each observable
by a brick of width \(2d\); the gaps are then \(>2\rho/a_N\) in lattice
units. Combine Lemmas IV.1.f and IV.1.g.

\emph{Lean:}
\href{https://github.com/yang-mills-group/OpenAI-YM-20261028/blob/3456e5e75cf8d4cd558955a5c7638fba2d68f464/YM/Extraction/Difference.lean\#L90-L145}{\texttt{Difference.primitive}},
\href{https://github.com/yang-mills-group/OpenAI-YM-20261028/blob/3456e5e75cf8d4cd558955a5c7638fba2d68f464/YM/Extraction/Difference.lean\#L24-L35}{\texttt{Difference.chrono}}.

\textbf{Lemma IV.1.i (equicontinuity).} Under (H1)--(H2), for every
arity \(n\):

\begin{itemize}
\item
  \textbf{(Q1)} for every \(p\), compact
  \(K\subset\mathrm{Gen}^{(p)}_n\), label list \(v\) and \(\epsilon>0\),
  there is \(\delta>0\) such that, for all large \(j\), the lifted
  lattice correlations at any \(x,y\in K\) with \(|x-y|<\delta\) differ
  by less than \(\epsilon\);
\item
  \textbf{(Q2)} for every \(\epsilon>0\), the same assertion holds for
  every compact \(K\subset\mathrm{Gen}_n\), eventually along ordered
  tails, with one \(\delta\) independent of the volume.
\end{itemize}

Quantitatively, let \(R>0\) bound below all cyclic coordinate gaps on a
closed neighbourhood of \(K\), and let \(A\) bound the centered roots of
the \(v_i\) and all \(\theta_av_i\) at a scale \(a_m<R/4\). Then
\(|\Delta\mathcal C|\le16nA^n|x-y|/R\) for \(|x-y|\) sufficiently small.

\emph{Proof sketch.} Telescope through at most \(4n\) coordinate moves
inside the neighbourhood and apply Lemma IV.1.h with \(\rho=R/4\). For a
negative move, interchange its endpoints. For \(n=1\) use translation
invariance; \(n=0\) is constant.

\emph{Lean:}
\href{https://github.com/yang-mills-group/OpenAI-YM-20261028/blob/3456e5e75cf8d4cd558955a5c7638fba2d68f464/YM/Extraction/Equicontinuity.lean\#L181-L214}{\texttt{Difference.completed}}
(conclusion
\href{https://github.com/yang-mills-group/OpenAI-YM-20261028/blob/3456e5e75cf8d4cd558955a5c7638fba2d68f464/YM/Extraction/Lipschitz.lean\#L26-L30}{\texttt{TransferLip}}),
\href{https://github.com/yang-mills-group/OpenAI-YM-20261028/blob/3456e5e75cf8d4cd558955a5c7638fba2d68f464/YM/Extraction/Equicontinuity.lean\#L115-L163}{\texttt{comparable}},
\href{https://github.com/yang-mills-group/OpenAI-YM-20261028/blob/3456e5e75cf8d4cd558955a5c7638fba2d68f464/YM/Extraction/Equicontinuity.lean\#L66-L103}{\texttt{stepTotal}},
\href{https://github.com/yang-mills-group/OpenAI-YM-20261028/blob/3456e5e75cf8d4cd558955a5c7638fba2d68f464/YM/Extraction/Equicontinuity.lean\#L51-L64}{\texttt{permits}},
\href{https://github.com/yang-mills-group/OpenAI-YM-20261028/blob/3456e5e75cf8d4cd558955a5c7638fba2d68f464/YM/Extraction/Gaps.lean\#L82-L93}{\texttt{Gaps.all}},
\href{https://github.com/yang-mills-group/OpenAI-YM-20261028/blob/3456e5e75cf8d4cd558955a5c7638fba2d68f464/YM/Extraction/Gaps.lean\#L119-L127}{\texttt{Gaps.allGrows}}.

\paragraph{4.4 Compactness and ordered
limits}\label{compactness-and-ordered-limits}

\textbf{Lemma IV.1.j (Arzelà--Ascoli for countably many families).} Let
\((Y_i)\) be a countable family of locally compact separable metric
spaces, \(U_i\subset Y_i\) open, and \(T^{(j)}_i:Y_i\to\mathbb C\).
Assume that \(T^{(j)}_i(y)\) is eventually bounded for each
\(y\in U_i\), and that the family is eventually equicontinuous on
compact subsets of \(U_i\). Then one strictly increasing sequence
\((t_j)\) makes every \(T^{(t_j)}_i\) converge locally uniformly on
\(U_i\) to a continuous limit.

\emph{Proof sketch.} Take countable dense sets and replace the finitely
many out-of-bound values at each dense point by a point in its bounding
ball. By Tychonoff, the countable product of closed balls is
sequentially compact, which gives a single subsequence converging at all
dense points. Equicontinuity upgrades this to locally uniform
convergence.

\emph{Lean:}
\href{https://github.com/yang-mills-group/OpenAI-YM-20261028/blob/3456e5e75cf8d4cd558955a5c7638fba2d68f464/YM/Extraction/Tame.lean\#L100-L114}{\texttt{localAscoli}},
\href{https://github.com/yang-mills-group/OpenAI-YM-20261028/blob/3456e5e75cf8d4cd558955a5c7638fba2d68f464/YM/Extraction/Simultaneous.lean\#L21-L32}{\texttt{simultaneous}},
\href{https://github.com/yang-mills-group/OpenAI-YM-20261028/blob/3456e5e75cf8d4cd558955a5c7638fba2d68f464/YM/Extraction/Tame.lean\#L20-L21}{\texttt{oscillation}}.

\textbf{Lemma IV.1.k (two-stage extraction).} The bound (B) and
equicontinuity (Q) along \(f\) imply IV.1.1 (a)--(b) along a refinement.

\emph{Proof sketch.}

\begin{itemize}
\item
  \emph{Indexing.} Index by
  \((p,n,\text{degrees},\text{basis vectors})\), a countable set. Limits
  are taken on basis labels and extended multilinearly.
\item
  \emph{Spacing.} Use the step extensions of the lattice correlators on
  \(\mathrm{Gen}^{(p)}_n\); the extension moves points by at most
  \(4na_N\). Equicontinuity comes from (Q1) and the bound from (B1).
  Lemma IV.1.j gives one subsequence of spacings for all \(p\).
\item
  \emph{Volume.} Use the functions \(P_{p_k}\) on \(\mathrm{Gen}_n\).
  Equicontinuity and the bound come from (Q2) and (B2) passed to
  \(j\to\infty\); this uses that a compact subset of \(\mathrm{Gen}_n\)
  lies in \(\mathrm{Gen}^{(p)}_n\) for large \(p\). Apply Lemma IV.1.j
  again.
\end{itemize}

\emph{Lean:}
\href{https://github.com/yang-mills-group/OpenAI-YM-20261028/blob/3456e5e75cf8d4cd558955a5c7638fba2d68f464/YM/Extraction/Lipschitz.lean\#L166-L184}{\texttt{first\_surface}},
\href{https://github.com/yang-mills-group/OpenAI-YM-20261028/blob/3456e5e75cf8d4cd558955a5c7638fba2d68f464/YM/Extraction/Lipschitz.lean\#L78-L105}{\texttt{setMeshes}},
\href{https://github.com/yang-mills-group/OpenAI-YM-20261028/blob/3456e5e75cf8d4cd558955a5c7638fba2d68f464/YM/Extraction/Lipschitz.lean\#L51-L68}{\texttt{each\_oscillates}},
\href{https://github.com/yang-mills-group/OpenAI-YM-20261028/blob/3456e5e75cf8d4cd558955a5c7638fba2d68f464/YM/Extraction/Lipschitz.lean\#L116-L138}{\texttt{volume\_oscillates}},
\href{https://github.com/yang-mills-group/OpenAI-YM-20261028/blob/3456e5e75cf8d4cd558955a5c7638fba2d68f464/YM/Extraction/Stencil.lean\#L69-L70}{\texttt{fitted}}.

\textbf{Lemma IV.1.l (ordered smeared limits and local integrability).}
The bound (B) together with IV.1.1 (a)--(b) implies IV.1.1 (c)--(d).

\emph{Proof sketch.} Dominated convergence, applied twice.

\begin{itemize}
\item
  \emph{Fixed volume.} At an admissible \(p\), the lattice sum is the
  integral of a step function. It converges almost everywhere (on
  \(\mathrm{Gen}^{(p)}_n\)) to \(tP_p\) and is dominated by
  \(\sup|t|\cdot D\) on the box, using (B1) on
  \(\mathrm{supp}t\subset\mathrm{Conf}^{(p)}_n\).
\item
  \emph{Volume limit.} \(|P_{p_k}|\le D\) on
  \(\mathrm{supp}t\cap\mathrm{Gen}_n\) for large \(k\), by (B2) in the
  limit \(j\to\infty\), and \(P_{p_k}\to S\) on the full-measure set
  \(\mathrm{Gen}_n\).
\end{itemize}

IV.1.1 is Lemmas IV.1.d, IV.1.i, IV.1.k and IV.1.l combined.

\emph{Lean:}
\href{https://github.com/yang-mills-group/OpenAI-YM-20261028/blob/3456e5e75cf8d4cd558955a5c7638fba2d68f464/YM/Extraction/Integrate.lean\#L155-L158}{\texttt{Order.poses}},
\href{https://github.com/yang-mills-group/OpenAI-YM-20261028/blob/3456e5e75cf8d4cd558955a5c7638fba2d68f464/YM/Extraction/Integrate.lean\#L115-L153}{\texttt{Order.ordinary}},
\href{https://github.com/yang-mills-group/OpenAI-YM-20261028/blob/3456e5e75cf8d4cd558955a5c7638fba2d68f464/YM/Extraction/Integrate.lean\#L66-L95}{\texttt{Order.intraperiod}},
\href{https://github.com/yang-mills-group/OpenAI-YM-20261028/blob/3456e5e75cf8d4cd558955a5c7638fba2d68f464/YM/Extraction/Integrate.lean\#L98-L113}{\texttt{Order.ring\_bound}},
\href{https://github.com/yang-mills-group/OpenAI-YM-20261028/blob/3456e5e75cf8d4cd558955a5c7638fba2d68f464/YM/Extraction/Ties.lean\#L44-L49}{\texttt{almostGeneric}};
\href{https://github.com/yang-mills-group/OpenAI-YM-20261028/blob/3456e5e75cf8d4cd558955a5c7638fba2d68f464/YM/Extraction/Held.lean\#L136-L137}{\texttt{crops}}.

\paragraph{4.5 Label derivatives}\label{label-derivatives}

\textbf{Lemma IV.1.m (one-slot difference quotients; D03).} Assume (H1),
IV.1.1 (c)--(d) and (B) along \(f\). For every arity \(n\), label list
\(v\), slot \(i\), axis \(a\) and \(t\in\mathcal D_n\), put \(u=v_i\)
and replace \(I(u,x_i)\) by the forward difference quotient
\(a_N^{-1}[I_{N,p}(u,x_i+e_a)-I_{N,p}(u,x_i)]\). The resulting smeared
lattice correlation converges along ordered tails to
\(-\int(\partial_{x_i^a}t)\,S(v;\cdot)\).

\emph{Proof sketch.}

\begin{itemize}
\item
  Summation by parts on the torus, which is exact, moves the difference
  onto the test as a backward difference quotient.
\item
  For \(p\) large and then \(N\) large the shifted test is evaluated
  without wrap-around.
\item
  The backward quotient \([t(x-a_Ne_{i,a})-t(x)]/a_N\) converges
  uniformly to \(-\partial_{x_i^a}t\), where \(e_{i,a}\) is the
  coordinate vector for slot \(i\) and axis \(a\). The error is at most
  \(L_p^{4n}\cdot\sup|\text{residual}|\cdot D\), with \(D\) from (B2) on
  a compact neighbourhood of \(\mathrm{supp}t\).
\item
  The main term converges by (c) applied to \(\partial_{x_i^a}t\).
\end{itemize}

\emph{Lean:}
\href{https://github.com/yang-mills-group/OpenAI-YM-20261028/blob/3456e5e75cf8d4cd558955a5c7638fba2d68f464/YM/Work/D03.lean\#L244-L249}{\texttt{Work.D03.answer}}
(statement
\href{https://github.com/yang-mills-group/OpenAI-YM-20261028/blob/3456e5e75cf8d4cd558955a5c7638fba2d68f464/YM/Bid/Drops.lean\#L61-L62}{\texttt{Bid.D03}}),
\href{https://github.com/yang-mills-group/OpenAI-YM-20261028/blob/3456e5e75cf8d4cd558955a5c7638fba2d68f464/YM/Work/D03.lean\#L71-L77}{\texttt{on\_lattice}},
\href{https://github.com/yang-mills-group/OpenAI-YM-20261028/blob/3456e5e75cf8d4cd558955a5c7638fba2d68f464/YM/Work/D03.lean\#L185-L242}{\texttt{error\_sum}},
\href{https://github.com/yang-mills-group/OpenAI-YM-20261028/blob/3456e5e75cf8d4cd558955a5c7638fba2d68f464/YM/Bid/Drops.lean\#L43-L44}{\texttt{heighten}}.

\textbf{Lemma IV.1.n (chessboard-small observables are invisible; D02).}
Assume (H1)--(H2). Let \(Q_{N,p}\) be chessboard-small with rate
\(y(N,p)\to0\). For every \(\delta>0\), arity, slot and fixed spectator
label list, there is \(A_\delta\ge0\) such that, eventually along
ordered tails, the correlation with \(Q\) in that slot has modulus at
most \(A_\delta|y(N,p)|\) at every configuration separated by \(\delta\)
in the periodic physical supremum norm. Consequently every smeared
lattice correlation with \(Q\) in one slot tends to \(0\) along ordered
tails.

\emph{Proof sketch.} The first claim is the chessboard argument of Lemma
IV.1.d with \(Q\) in one cell. For the second, write the smeared sum as
the integral of a step function supported in a fixed ball and bounded by
\(\sup|t|\,A_\delta|y|\).

\emph{Lean:}
\href{https://github.com/yang-mills-group/OpenAI-YM-20261028/blob/3456e5e75cf8d4cd558955a5c7638fba2d68f464/YM/Lane/Lease.lean\#L57-L124}{\texttt{Lease.take}},
\href{https://github.com/yang-mills-group/OpenAI-YM-20261028/blob/3456e5e75cf8d4cd558955a5c7638fba2d68f464/YM/Lane/Lease.lean\#L20-L22}{\texttt{Lease.rents}},
\href{https://github.com/yang-mills-group/OpenAI-YM-20261028/blob/3456e5e75cf8d4cd558955a5c7638fba2d68f464/YM/Lane/Lease.lean\#L54-L55}{\texttt{Lease.gain}},
\href{https://github.com/yang-mills-group/OpenAI-YM-20261028/blob/3456e5e75cf8d4cd558955a5c7638fba2d68f464/YM/Work/D02.lean\#L24-L84}{\texttt{Work.D02.answer}}
(statement
\href{https://github.com/yang-mills-group/OpenAI-YM-20261028/blob/3456e5e75cf8d4cd558955a5c7638fba2d68f464/YM/Bid/Drops.lean\#L57-L58}{\texttt{Bid.D02}}).

\textbf{Lemma IV.1.o (bounded labels have subpolynomial cutoff
coordinates; D01).} Under (H1), every degree-\(d\) label family
\(w(N,p)\) bounded along ordered tails has subpolynomial growth in the
sense of §2, input 3.

\emph{Proof sketch.} The inverse terminal transfer is bounded, up to a
fixed degree-dependent norm factor, by the product of the column
constants up to level \(N\), and their logarithms sum to
\(o(\log a_N^{-1})\).

\emph{Lean:}
\href{https://github.com/yang-mills-group/OpenAI-YM-20261028/blob/3456e5e75cf8d4cd558955a5c7638fba2d68f464/YM/Work/D01.lean\#L127-L134}{\texttt{Work.D01.answer}}
(statement
\href{https://github.com/yang-mills-group/OpenAI-YM-20261028/blob/3456e5e75cf8d4cd558955a5c7638fba2d68f464/YM/Bid/Drops.lean\#L54}{\texttt{Bid.D01}}),
\href{https://github.com/yang-mills-group/OpenAI-YM-20261028/blob/3456e5e75cf8d4cd558955a5c7638fba2d68f464/YM/Work/D01.lean\#L92-L125}{\texttt{terminal\_ball}},
\href{https://github.com/yang-mills-group/OpenAI-YM-20261028/blob/3456e5e75cf8d4cd558955a5c7638fba2d68f464/YM/Work/D01.lean\#L34-L63}{\texttt{columnLogs\_le}}.

\textbf{Lemma IV.1.p (spanning; D04).} Assume (H1), (H2), IV.1.1
(c)--(d), and the relay data and dimension count for every \(d\). Then:

\begin{enumerate}
\def\labelenumi{\arabic{enumi}.}
\item
  for each fixed \(u\in\mathrm{Label}_d\), the smeared lattice
  correlations of \(u-\pi_d(N,p)u\) (in any slot, with any spectators
  and tests) tend to \(0\) along ordered tails; together with the
  dimension count this gives \(\ker\pi_d=\mathrm{Null}_d(S)\);
\item
  for \(d\le e\) and \(u\in\mathrm{Label}_d\) there are degree-\(e\)
  labels \(w(N,p)\) whose smeared correlations converge along ordered
  tails to those of \(u\) in the limit;
\item
  for \(u\in\mathrm{Label}_d\) and each axis \(a\) there are
  degree-\((d+1)\) labels \(w(N,p)\) whose smeared correlations converge
  along ordered tails to \(-\int(\partial_{x_i^a}t)\,S(v;\cdot)\).
\end{enumerate}

In 2 and 3, one family \(w(N,p)\) works for every arity, slot, spectator
list and test. It may be chosen with subpolynomial growth, and the
normalized remainders \((w-\pi w)/(1+\|\pi w\|)\) also have vanishing
smeared limits.

\emph{Proof sketch.}

\begin{itemize}
\item
  \begin{enumerate}
  \def\labelenumi{(\arabic{enumi})}
  \tightlist
  \item
    \(u-\pi(N,p)u\) is bounded and lies in \(\ker\pi(N,p)\). By Lemma
    IV.1.o it has subpolynomial growth, so \emph{mark} applies: it is
    chessboard-close to an observable with exactly vanishing separated
    correlations. Lemma IV.1.n then shows the smeared correlations
    vanish. The one-slot smear maps converge in operator norm, since
    their domains are finite-dimensional. This allows passage to varying
    labels and gives \(\ker\pi_d\subset\mathrm{Null}_d(S)\); the
    dimension count gives equality.
  \end{enumerate}
\item
  (2)--(3) \emph{Transit}, resp. \emph{pace}, gives labels whose
  observables are chessboard-close to the embedded label, resp. to the
  forward difference quotient. Lemma IV.1.n transfers the smeared
  limits, which Lemma IV.1.m computes for the difference quotient.
\item
  For the normalized remainder, dividing \(w\) by \(1+\|\pi w\|\)
  preserves subpolynomial growth, while \(\pi w/(1+\|\pi w\|)\) is
  bounded and Lemma IV.1.o applies. Their difference lies in the
  projection kernel, so \emph{mark} and Lemma IV.1.n apply.
\end{itemize}

\emph{Lean:}
\href{https://github.com/yang-mills-group/OpenAI-YM-20261028/blob/3456e5e75cf8d4cd558955a5c7638fba2d68f464/YM/Work/D04.lean\#L212-L240}{\texttt{Work.D04.answer}}
(statement
\href{https://github.com/yang-mills-group/OpenAI-YM-20261028/blob/3456e5e75cf8d4cd558955a5c7638fba2d68f464/YM/Bid/Drops.lean\#L65-L67}{\texttt{Bid.D04}}),
\href{https://github.com/yang-mills-group/OpenAI-YM-20261028/blob/3456e5e75cf8d4cd558955a5c7638fba2d68f464/YM/Dash/Inert.lean\#L55-L60}{\texttt{Spans}},
\href{https://github.com/yang-mills-group/OpenAI-YM-20261028/blob/3456e5e75cf8d4cd558955a5c7638fba2d68f464/YM/Dash/Inert.lean\#L24-L26}{\texttt{Bin.Upper}},
\href{https://github.com/yang-mills-group/OpenAI-YM-20261028/blob/3456e5e75cf8d4cd558955a5c7638fba2d68f464/YM/Dash/Inert.lean\#L41-L42}{\texttt{Bin.exactC}}.

\textbf{Lemma IV.1.q (fixed derivative labels).} The conclusions of
Lemma IV.1.p yield the maps \(\kappa\) and \(\tau\) of IV.1.2.

\emph{Proof sketch.} This is a finite-dimensional compactness argument.

\begin{itemize}
\item
  Take a cluster point \((c,z)\) of
  \(\bigl((1+\|\pi w\|)^{-1},\,\pi w/(1+\|\pi w\|)\bigr)\) along the
  ordered-tail filter. Use one refining ultrafilter for this
  finite-dimensional pair; all spectator tests then retain their limits.
\item
  Idempotence and convergence of the projections give \(\pi z=z\). For
  every spectator test, the limits give
  \(\langle S,z\rangle=c\cdot(\text{target})\).
\item
  \(c\ne0\): otherwise \(z\in\mathrm{ran}\pi\cap\ker\pi=\{0\}\) (since
  \(\ker\pi=\mathrm{Null}\)), contradicting \(\|z\|+c=1\). So \(z/c\) is
  a fixed label with the target kernels.
\item
  Define \(\kappa\) and \(\tau\) this way on a basis and extend
  linearly.
\end{itemize}

The identity for \(\tau\) is the distributional relation
\(\partial_{x_i^a}S_v=S_{v_{i\to\tau_av_i}}\).

\emph{Lean:}
\href{https://github.com/yang-mills-group/OpenAI-YM-20261028/blob/3456e5e75cf8d4cd558955a5c7638fba2d68f464/YM/Dash/Inert.lean\#L71-L87}{\texttt{lifts\_model}},
\href{https://github.com/yang-mills-group/OpenAI-YM-20261028/blob/3456e5e75cf8d4cd558955a5c7638fba2d68f464/YM/Dash/Limits.lean\#L69-L111}{\texttt{Dash.brace}},
\href{https://github.com/yang-mills-group/OpenAI-YM-20261028/blob/3456e5e75cf8d4cd558955a5c7638fba2d68f464/YM/Dash/Inert.lean\#L89-L96}{\texttt{Wakes.through\_ties}},
\href{https://github.com/yang-mills-group/OpenAI-YM-20261028/blob/3456e5e75cf8d4cd558955a5c7638fba2d68f464/YM/Dash/Scope.lean\#L108}{\texttt{Dash.monitor}}.

\paragraph{4.6 Smoothness}\label{smoothness}

\textbf{Lemma IV.1.r (distributional regularity).} Let \(K_w\), with
\(w\) ranging over label lists of arity \(n\), be locally integrable on
\(\mathrm{Conf}_n\) and continuous on \(\mathrm{Gen}_n\). Assume that
for every \(w\), slot \(i\) and axis \(a\) there is \(w'\) with
\(\partial_{x_i^a}K_w=K_{w'}\) in \(\mathcal D'(\mathrm{Conf}_n)\). Then
every \(K_w\) agrees almost everywhere on \(\mathrm{Conf}_n\), and
everywhere on \(\mathrm{Gen}_n\), with a function in
\(C^\infty(\mathrm{Conf}_n)\).

\emph{Proof sketch.}

\begin{itemize}
\item
  For \(\varphi\in\mathcal D_n\):
  \(\xi_{i,a}\,\widehat{\varphi K_w}(\xi)\) is a constant times
  \(\widehat{\varphi K_{w'}}+\widehat{(\partial_{i,a}\varphi)K_w}\).
\item
  By induction on \(m\), \(|\xi_{i,a}|^m|\widehat{\varphi K_w}|\) is
  bounded for all \(m\) and all coordinates. So
  \(\widehat{\varphi K_w}\) decays rapidly, and \(\varphi K_w\) equals a
  smooth function almost everywhere.
\item
  Take \(\varphi=1\) near an arbitrary point of \(\mathrm{Conf}_n\),
  generic or not.
\item
  The local smooth representatives agree on \(\mathrm{Gen}_n\), which
  has full measure, hence on overlaps; glue them. Multilinearity in the
  labels follows from the same density argument.
\end{itemize}

\emph{Lean:}
\href{https://github.com/yang-mills-group/OpenAI-YM-20261028/blob/3456e5e75cf8d4cd558955a5c7638fba2d68f464/YM/Dash/Slew.lean\#L26-L58}{\texttt{Dash.awake}},
\href{https://github.com/yang-mills-group/OpenAI-YM-20261028/blob/3456e5e75cf8d4cd558955a5c7638fba2d68f464/YM/Dash/Regular.lean\#L91-L113}{\texttt{aggregate}},
\href{https://github.com/yang-mills-group/OpenAI-YM-20261028/blob/3456e5e75cf8d4cd558955a5c7638fba2d68f464/YM/Dash/Regular.lean\#L62-L71}{\texttt{repair}},
\href{https://github.com/yang-mills-group/OpenAI-YM-20261028/blob/3456e5e75cf8d4cd558955a5c7638fba2d68f464/YM/Dash/Twist.lean\#L98-L115}{\texttt{arises}},
\href{https://github.com/yang-mills-group/OpenAI-YM-20261028/blob/3456e5e75cf8d4cd558955a5c7638fba2d68f464/YM/Dash/Twist.lean\#L94-L95}{\texttt{Hidden}},
\href{https://github.com/yang-mills-group/OpenAI-YM-20261028/blob/3456e5e75cf8d4cd558955a5c7638fba2d68f464/YM/Dash/Twist.lean\#L58-L59}{\texttt{currents}},
\href{https://github.com/yang-mills-group/OpenAI-YM-20261028/blob/3456e5e75cf8d4cd558955a5c7638fba2d68f464/YM/Dash/Slew.lean\#L17-L18}{\texttt{closeBy}}.

\textbf{Lemma IV.1.s (the smooth version is a raw limit with a
derivative plan).} Let \(T\) be given by Lemma IV.1.r applied to \(S\)
and the maps of IV.1.2. Then \(T\) satisfies the conclusions of IV.1.3.

\emph{Proof sketch.}

\begin{itemize}
\item
  \(T\) satisfies IV.1.1 (a)--(d) with the same \(P_p\), because the
  volume limit only sees \(\mathrm{Gen}_n\), and (c) follows from Lemma
  IV.1.l;
\item
  (E1) holds by Lemma IV.1.t;
\item
  the kernel identities pass from \(S\) to \(T\), since the smeared
  values are the same;
\item
  \(\kappa\) and \(\tau\) preserve null labels: their smeared kernels
  are respectively the original kernel and its distributional
  derivative. Smoothness turns these into pointwise identities, so the
  maps descend to the quotients \(F_d\).
\end{itemize}

\emph{Lean:}
\href{https://github.com/yang-mills-group/OpenAI-YM-20261028/blob/3456e5e75cf8d4cd558955a5c7638fba2d68f464/YM/Org/Input.lean\#L101-L109}{\texttt{Org.Arrives.serves}},
\href{https://github.com/yang-mills-group/OpenAI-YM-20261028/blob/3456e5e75cf8d4cd558955a5c7638fba2d68f464/YM/Extraction/Lipschitz.lean\#L194-L202}{\texttt{RawLimits.ready}},
\href{https://github.com/yang-mills-group/OpenAI-YM-20261028/blob/3456e5e75cf8d4cd558955a5c7638fba2d68f464/YM/Euclid/Null.lean\#L74}{\texttt{Table.Field}}.

\paragraph{4.7 Euclidean axioms, bounds and
assembly}\label{euclidean-axioms-bounds-and-assembly}

\textbf{Lemma IV.1.t (symmetries and basic identities).} If \(T\) is
smooth on \(\mathrm{Conf}_n\) and is a point limit (IV.1.1 (a)--(b))
along \(f\), and (H1) holds, then \(T\) satisfies (E1)--(E3).

\emph{Proof sketch.}

\begin{itemize}
\item
  \emph{Grid convergence.} For a point \(x\) of the level-\(k\) grid
  \(a_k\mathbb Z^{4n}\) with \(x\in\mathrm{Gen}_n\), use its reduction
  in \(\Lambda_{N,p}\) for \(N\ge k\) and sufficiently large \(p\). The
  lattice correlations there converge to \(T(v;x)\) along ordered tails.
\item
  \emph{Lattice identities.} Permutation symmetry, \(H_4\)-covariance
  about the origin, translations by grid vectors, reality and the unit
  label hold exactly on the lattice, by (H1) and the invariance of the
  Wilson measure.
\item
  \emph{Extension.} Equality on the dense set of generic grid points
  extends to \(\mathrm{Gen}_n\) by continuity, and then to
  \(\mathrm{Conf}_n\), since both sides are continuous there and
  \(\mathrm{Gen}_n\) has full measure. Real translations are
  approximated by grid translations.
\end{itemize}

\emph{Lean:}
\href{https://github.com/yang-mills-group/OpenAI-YM-20261028/blob/3456e5e75cf8d4cd558955a5c7638fba2d68f464/YM/Extraction/Inherited.lean\#L151-L193}{\texttt{Order.limits\_basic}},
\href{https://github.com/yang-mills-group/OpenAI-YM-20261028/blob/3456e5e75cf8d4cd558955a5c7638fba2d68f464/YM/Extraction/Inherited.lean\#L124-L149}{\texttt{Order.spin}},
\href{https://github.com/yang-mills-group/OpenAI-YM-20261028/blob/3456e5e75cf8d4cd558955a5c7638fba2d68f464/YM/Extraction/Inherited.lean\#L83-L112}{\texttt{Order.homogeneous}},
\href{https://github.com/yang-mills-group/OpenAI-YM-20261028/blob/3456e5e75cf8d4cd558955a5c7638fba2d68f464/YM/Extraction/Inherited.lean\#L39-L44}{\texttt{Order.inPatch}},
\href{https://github.com/yang-mills-group/OpenAI-YM-20261028/blob/3456e5e75cf8d4cd558955a5c7638fba2d68f464/YM/Extraction/Inherited.lean\#L48-L51}{\texttt{Order.equal}}.

\textbf{Lemma IV.1.u (reflection positivity of the limit).} If \(T\) is
smooth, symmetric and an ordered limit along \(f\), and (H1) holds, then
(E5) holds; (E5) implies (E4).

\emph{Proof sketch.}

\begin{itemize}
\item
  \emph{Lattice side.} For a test supported in \(\{x^a>0\}\), take \(p\)
  large enough that it is admissible. The smeared lattice observable
  \(a_N^{4n}\sum_rt(\widehat r)\prod I(v_k,r_k)\) is then, for \(N\)
  large, interior to the half-torus \(0<x^a<L_p/2\). Its reflection
  through the site plane \(x^a=0\) carries labels \(\theta_a\), because
  \(I(\theta_av,\Theta x)=\overline{I(v,x)\circ\Theta}\). The lattice
  Gram matrix of such observables is the smeared lattice correlation of
  \(\overline{t_i}\otimes t_j\) pulled back by reflection of the first
  configuration. Symmetry permits reversal of its entries. The resulting
  joint test has compact support in \(\mathrm{Conf}_{n_i+n_j}\), since
  the two configurations lie strictly on opposite sides of the plane.
  Lemma IV.1.a makes the lattice Gram matrix positive semidefinite.
\item
  \emph{Limit.} Smoothness gives absolute integrability on the compact
  joint supports. Pass to the ordered limit using (c).
\item
  \emph{(E4) from (E5).} Use bump functions concentrating at the points,
  together with continuity of \(T\).
\end{itemize}

\emph{Lean:}
\href{https://github.com/yang-mills-group/OpenAI-YM-20261028/blob/3456e5e75cf8d4cd558955a5c7638fba2d68f464/YM/Extraction/Harvest.lean\#L46-L60}{\texttt{Order.realPositive}},
\href{https://github.com/yang-mills-group/OpenAI-YM-20261028/blob/3456e5e75cf8d4cd558955a5c7638fba2d68f464/YM/Extraction/Pairs.lean\#L104-L125}{\texttt{Banks.at\_cut}},
\href{https://github.com/yang-mills-group/OpenAI-YM-20261028/blob/3456e5e75cf8d4cd558955a5c7638fba2d68f464/YM/Extraction/Banks.lean\#L94-L102}{\texttt{Banks.kept}},
\href{https://github.com/yang-mills-group/OpenAI-YM-20261028/blob/3456e5e75cf8d4cd558955a5c7638fba2d68f464/YM/Extraction/Banks.lean\#L19}{\texttt{Banks.plane}},
\href{https://github.com/yang-mills-group/OpenAI-YM-20261028/blob/3456e5e75cf8d4cd558955a5c7638fba2d68f464/YM/Euclid/RPPeak.lean\#L126-L131}{\texttt{Pairing.at\_point}},
\href{https://github.com/yang-mills-group/OpenAI-YM-20261028/blob/3456e5e75cf8d4cd558955a5c7638fba2d68f464/YM/Euclid/RPPeak.lean\#L16-L33}{\texttt{Pairing.nearGram}}.

\textbf{Lemma IV.1.v (the polynomial bound).} Under the hypotheses of
IV.1.5, its stated exponent \(4n(1+\max_i d_i)\) is valid.

\emph{Proof sketch.} Set \(D=1+\max_i d_i\). The centered power budgets
for all transformed labels can be bounded by \(B a_k^{-D}\) for every
\(k\ge M\), with fixed \(B\ge0\) and \(M\).

\begin{itemize}
\item
  Choose the first \(m\ge M\) with \(a_m<r/2\), and the first \(k\ge M\)
  with \(a_k<a_m/[6(n+1)]\). The scale inequality, including the
  possible endpoint \(m=M\) or \(k=M\), gives \(a_m\ge c(r/2)^2\) and
  then \(a_k\ge c'r^4\).
\item
  Lemma IV.1.d bounds the lattice correlations by \((B a_k^{-D})^n\).
  Pass first to the spacing limit and then to the volume limit at
  generic points with separation \(>r/2\).
\item
  Continuity and density extend the bound to every \(r\)-separated
  configuration. Since \(a_k\ge c'r^4\), the bound is at most
  \(C r^{-4Dn}\), with \(C\) independent of positions and \(r\).
\end{itemize}

\emph{Lean:}
\href{https://github.com/yang-mills-group/OpenAI-YM-20261028/blob/3456e5e75cf8d4cd558955a5c7638fba2d68f464/YM/Work/D05.lean\#L160-L180}{\texttt{Work.D05.answer}},
\href{https://github.com/yang-mills-group/OpenAI-YM-20261028/blob/3456e5e75cf8d4cd558955a5c7638fba2d68f464/YM/Work/D05.lean\#L43-L68}{\texttt{image\_power}},
\href{https://github.com/yang-mills-group/OpenAI-YM-20261028/blob/3456e5e75cf8d4cd558955a5c7638fba2d68f464/YM/Work/D05.lean\#L132-L158}{\texttt{scales}},
\href{https://github.com/yang-mills-group/OpenAI-YM-20261028/blob/3456e5e75cf8d4cd558955a5c7638fba2d68f464/YM/Work/D05.lean\#L95-L128}{\texttt{separated\_bound}}.

\textbf{Lemma IV.1.w (orientation).} Under the hypotheses of IV.1.6, its
four symmetry identities hold on the quotient spaces \(F_d\).

\emph{Proof sketch.} A class in \(F_d\) is determined by its kernels
against all spectators. Compare the kernels of both sides using (E2) and
the chain rule for \(x\mapsto sx\) and for complex conjugation.

\emph{Lean:}
\href{https://github.com/yang-mills-group/OpenAI-YM-20261028/blob/3456e5e75cf8d4cd558955a5c7638fba2d68f464/YM/Work/D06.lean\#L98-L135}{\texttt{Work.D06.compat}}.

\textbf{Lemma IV.1.x (assembly).} If \(T\) is an ordered limit and a
point limit along \(f\), (H1) holds, \(T\) is smooth and (E7) holds,
then \(T\) satisfies the hierarchy (E1)--(E7).

\emph{Proof sketch.} Lemmas IV.1.t and IV.1.u give (E1)--(E5);
(E6)--(E7) are assumed. When (H2) and the derivative plan of IV.1.3 are
also available, IV.1.5 and B10 supply (E7). B10 inducts on the
derivative order: each coordinate derivative is a kernel with one raised
label, and the operator norm is bounded by a finite sum of coordinate
norms. The base case is IV.1.5. On the main path, compactness of the
bounded projection matrices gives a common refinement before forming the
bundle of Remark 1.

\emph{Lean:}
\href{https://github.com/yang-mills-group/OpenAI-YM-20261028/blob/3456e5e75cf8d4cd558955a5c7638fba2d68f464/YM/Extraction/Harvest.lean\#L62-L65}{\texttt{Order.assembles}},
\href{https://github.com/yang-mills-group/OpenAI-YM-20261028/blob/3456e5e75cf8d4cd558955a5c7638fba2d68f464/YM/Org/Input.lean\#L113-L114}{\texttt{Org.Cleaner.stay}},
\href{https://github.com/yang-mills-group/OpenAI-YM-20261028/blob/3456e5e75cf8d4cd558955a5c7638fba2d68f464/YM/Bid/Drops.lean\#L91-L97}{\texttt{Bid.healthy}},
\href{https://github.com/yang-mills-group/OpenAI-YM-20261028/blob/3456e5e75cf8d4cd558955a5c7638fba2d68f464/YM/Work/B10.lean\#L72-L90}{\texttt{Work.B10.capped\_tags}},
\href{https://github.com/yang-mills-group/OpenAI-YM-20261028/blob/3456e5e75cf8d4cd558955a5c7638fba2d68f464/YM/Next/EXTRACTION/Relay.lean\#L86-L93}{\texttt{include\_projections}},
\href{https://github.com/yang-mills-group/OpenAI-YM-20261028/blob/3456e5e75cf8d4cd558955a5c7638fba2d68f464/YM/Next/EXTRACTION/Relay.lean\#L178-L181}{\texttt{Projected.ready}},
\href{https://github.com/yang-mills-group/OpenAI-YM-20261028/blob/3456e5e75cf8d4cd558955a5c7638fba2d68f464/YM/Next/EXTRACTION/Base.lean\#L42-L51}{\texttt{Next.Extract.Base}}.

\subsubsection{5. Where the difficulty
lies}\label{where-the-difficulty-lies-19}

The main estimates are the transfer-kernel Lipschitz bound and the
control of uncentered chessboard roots. Complex interpolation and
Cauchy's estimate turn reflection positivity into equicontinuity without
prior derivative bounds. Translation log-convexity allows the centered
budgets to control observables at arbitrary interior positions.

The regularity argument depends on more than compactness: the relay
estimates supply degree-raised approximants to differences, and the rank
comparison identifies precisely which labels vanish. The normalized
remainder estimate in Lemma IV.1.p is needed to obtain fixed derivative
labels. Fourier regularity, passage of symmetries by density, and
positivity of limiting Gram matrices then use standard arguments.

Three distinctions matter:

\begin{itemize}
\item
  Equicontinuity is proved on \(\mathrm{Gen}_n\), whereas bounds and the
  smooth representative extend to \(\mathrm{Conf}_n\).
\item
  The maps \(\tau_{d,a}\) are chosen by compactness. Their kernel
  identities determine them only modulo \(\mathrm{Null}_{d+1}\); the
  induced maps on \(F_d\) are intrinsic.
\item
  The hierarchy here has signed coordinate symmetries and reflection
  positivity at coordinate planes. Translation invariance moves those
  planes. Continuous rotation symmetry is a later result.
\end{itemize}

\subsubsection{6. Notation}\label{notation-19}

{\def\LTcaptype{none} % do not increment counter
\begin{longtable}[]{@{}
  >{\raggedright\arraybackslash}p{(\linewidth - 4\tabcolsep) * \real{0.3333}}
  >{\raggedright\arraybackslash}p{(\linewidth - 4\tabcolsep) * \real{0.3333}}
  >{\raggedright\arraybackslash}p{(\linewidth - 4\tabcolsep) * \real{0.3333}}@{}}
\toprule\noalign{}
\begin{minipage}[b]{\linewidth}\raggedright
object
\end{minipage} & \begin{minipage}[b]{\linewidth}\raggedright
code name
\end{minipage} & \begin{minipage}[b]{\linewidth}\raggedright
proposed notation
\end{minipage} \\
\midrule\noalign{}
\endhead
\bottomrule\noalign{}
\endlastfoot
Wilson family; spacing, period, coupling &
\href{https://github.com/yang-mills-group/OpenAI-YM-20261028/blob/3456e5e75cf8d4cd558955a5c7638fba2d68f464/YM/Source/Construction.lean\#L24-L30}{\texttt{Family}}
(\texttt{C.S.a}, \texttt{C.S.period}, \texttt{C.beta}) & \(C\); \(a_N\),
\(L_p\), \(\beta_{N,p}\) \\
specification of the family &
\href{https://github.com/yang-mills-group/OpenAI-YM-20261028/blob/3456e5e75cf8d4cd558955a5c7638fba2d68f464/YM/Source/Construction.lean\#L44-L62}{\texttt{Family.Makes}}
& (H1) \\
label, register of labels &
\href{https://github.com/yang-mills-group/OpenAI-YM-20261028/blob/3456e5e75cf8d4cd558955a5c7638fba2d68f464/YM/Euclid/Register.lean\#L88}{\texttt{Register.Token}},
\href{https://github.com/yang-mills-group/OpenAI-YM-20261028/blob/3456e5e75cf8d4cd558955a5c7638fba2d68f464/YM/Extraction/LatticeFamily.lean\#L20}{\texttt{Family.book}}
& \(v\in\mathrm{Label}_d\) \\
unit label; reflected label &
\href{https://github.com/yang-mills-group/OpenAI-YM-20261028/blob/3456e5e75cf8d4cd558955a5c7638fba2d68f464/YM/Euclid/Register.lean\#L90}{\texttt{Register.one}},
\href{https://github.com/yang-mills-group/OpenAI-YM-20261028/blob/3456e5e75cf8d4cd558955a5c7638fba2d68f464/YM/Euclid/Register.lean\#L93}{\texttt{Register.reflectedToken}}
& \(1_d\); \(\theta_av=\rho_d(s_a)\bar v\) \\
renormalized lattice observable &
\href{https://github.com/yang-mills-group/OpenAI-YM-20261028/blob/3456e5e75cf8d4cd558955a5c7638fba2d68f464/YM/Source/Construction.lean\#L42}{\texttt{Family.O}}
& \(I_{N,p}(v,x)\) \\
lattice correlation; smeared &
\href{https://github.com/yang-mills-group/OpenAI-YM-20261028/blob/3456e5e75cf8d4cd558955a5c7638fba2d68f464/YM/Extraction/LatticeFamily.lean\#L68-L69}{\texttt{Family.numeric}},
\href{https://github.com/yang-mills-group/OpenAI-YM-20261028/blob/3456e5e75cf8d4cd558955a5c7638fba2d68f464/YM/Extraction/Limits.lean\#L26-L27}{\texttt{Order.smears}}
& \(\mathcal C_{N,p}(v;x)\), \(\mathcal C_{N,p}(v;t)\) \\
configuration sets &
\href{https://github.com/yang-mills-group/OpenAI-YM-20261028/blob/3456e5e75cf8d4cd558955a5c7638fba2d68f464/YM/Space/Configurations.lean\#L63}{\texttt{conf}},
\href{https://github.com/yang-mills-group/OpenAI-YM-20261028/blob/3456e5e75cf8d4cd558955a5c7638fba2d68f464/YM/Space/Configurations.lean\#L66}{\texttt{generic}},
\href{https://github.com/yang-mills-group/OpenAI-YM-20261028/blob/3456e5e75cf8d4cd558955a5c7638fba2d68f464/YM/Extraction/Patch.lean\#L25-L26}{\texttt{NoHit}},
\href{https://github.com/yang-mills-group/OpenAI-YM-20261028/blob/3456e5e75cf8d4cd558955a5c7638fba2d68f464/YM/Extraction/Patch.lean\#L28-L29}{\texttt{Cyclic}}
& \(\mathrm{Conf}_n\), \(\mathrm{Gen}_n\), \(\mathrm{Conf}^{(p)}_n\),
\(\mathrm{Gen}^{(p)}_n\) \\
test functions; admissible at \(p\) &
\href{https://github.com/yang-mills-group/OpenAI-YM-20261028/blob/3456e5e75cf8d4cd558955a5c7638fba2d68f464/YM/Euclid/Test.lean\#L18}{\texttt{Test}},
\href{https://github.com/yang-mills-group/OpenAI-YM-20261028/blob/3456e5e75cf8d4cd558955a5c7638fba2d68f464/YM/Extraction/Patch.lean\#L31}{\texttt{suitable}}
& \(\mathcal D_n\); \(\mathrm{supp}t\subset(-L_p/2,L_p/2)^{4n}\) \\
extraction; refinement; ordered sense &
\href{https://github.com/yang-mills-group/OpenAI-YM-20261028/blob/3456e5e75cf8d4cd558955a5c7638fba2d68f464/YM/Extraction/Orders.lean\#L34-L36}{\texttt{Order}},
\href{https://github.com/yang-mills-group/OpenAI-YM-20261028/blob/3456e5e75cf8d4cd558955a5c7638fba2d68f464/YM/Extraction/Orders.lean\#L40}{\texttt{refine}},
\href{https://github.com/yang-mills-group/OpenAI-YM-20261028/blob/3456e5e75cf8d4cd558955a5c7638fba2d68f464/YM/Extraction/Orders.lean\#L47}{\texttt{late}},
\href{https://github.com/yang-mills-group/OpenAI-YM-20261028/blob/3456e5e75cf8d4cd558955a5c7638fba2d68f464/YM/Extraction/Inherited.lean\#L27}{\texttt{fuse}}
& \(f=((N_j),(p_k))\), \(f\circ g\), ``\(k\) large, then \(j\)
large'' \\
table &
\href{https://github.com/yang-mills-group/OpenAI-YM-20261028/blob/3456e5e75cf8d4cd558955a5c7638fba2d68f464/YM/Euclid/Table.lean\#L55-L56}{\texttt{Table}}
& \(S(v;x)\), \(T(v;x)\) \\
raw limits &
\href{https://github.com/yang-mills-group/OpenAI-YM-20261028/blob/3456e5e75cf8d4cd558955a5c7638fba2d68f464/YM/Extraction/Lipschitz.lean\#L156-L160}{\texttt{RawLimits}},
\href{https://github.com/yang-mills-group/OpenAI-YM-20261028/blob/3456e5e75cf8d4cd558955a5c7638fba2d68f464/YM/Extraction/Limits.lean\#L68}{\texttt{PointLimits}}
& IV.1.1 (a)--(c); \(P_p\) \\
tiling, copy, chessboard root &
\href{https://github.com/yang-mills-group/OpenAI-YM-20261028/blob/3456e5e75cf8d4cd558955a5c7638fba2d68f464/YM/Cells/Tiling.lean\#L19-L25}{\texttt{Tiling}},
\href{https://github.com/yang-mills-group/OpenAI-YM-20261028/blob/3456e5e75cf8d4cd558955a5c7638fba2d68f464/YM/Cells/Tiling.lean\#L221-L224}{\texttt{copy}},
\href{https://github.com/yang-mills-group/OpenAI-YM-20261028/blob/3456e5e75cf8d4cd558955a5c7638fba2d68f464/YM/Cells/Tiling.lean\#L226-L227}{\texttt{rootTest}}
& \(\mathcal T\), \(\mathrm{copy}_c\), \(\|F\|_{\mathcal T}\) \\
centered tiling at \(x\), cells of side \(a_m\) &
\href{https://github.com/yang-mills-group/OpenAI-YM-20261028/blob/3456e5e75cf8d4cd558955a5c7638fba2d68f464/YM/Cells/Packing.lean\#L34-L40}{\texttt{Cells.center}}
& \(\mathcal T_m(x)\) \\
interior function of a cell &
\href{https://github.com/yang-mills-group/OpenAI-YM-20261028/blob/3456e5e75cf8d4cd558955a5c7638fba2d68f464/YM/Cells/Cut.lean\#L109-L113}{\texttt{Interior}}
& \(F\in\mathcal A(\text{cell})\) \\
reflection positivity and chessboard inequality &
\href{https://github.com/yang-mills-group/OpenAI-YM-20261028/blob/3456e5e75cf8d4cd558955a5c7638fba2d68f464/YM/Cells/Certificate.lean\#L14-L19}{\texttt{Seam}}
& (RP), (CB) \\
budgets &
\href{https://github.com/yang-mills-group/OpenAI-YM-20261028/blob/3456e5e75cf8d4cd558955a5c7638fba2d68f464/YM/Extraction/Center.lean\#L38-L40}{\texttt{Wcmain}},
\href{https://github.com/yang-mills-group/OpenAI-YM-20261028/blob/3456e5e75cf8d4cd558955a5c7638fba2d68f464/YM/Extraction/Center.lean\#L30-L34}{\texttt{CenterPort}}
& input 2; (H2) \\
uniform bound; equicontinuity &
\href{https://github.com/yang-mills-group/OpenAI-YM-20261028/blob/3456e5e75cf8d4cd558955a5c7638fba2d68f464/YM/Extraction/Integrate.lean\#L25-L30}{\texttt{Bound}},
\href{https://github.com/yang-mills-group/OpenAI-YM-20261028/blob/3456e5e75cf8d4cd558955a5c7638fba2d68f464/YM/Extraction/Lipschitz.lean\#L26-L30}{\texttt{TransferLip}}
& (B1)--(B2), (E1)--(E2) \\
forward difference quotient &
\href{https://github.com/yang-mills-group/OpenAI-YM-20261028/blob/3456e5e75cf8d4cd558955a5c7638fba2d68f464/YM/Bid/Drops.lean\#L43-L44}{\texttt{heighten}}
& \(a_N^{-1}[I(v,\cdot+e_a)-I(v,\cdot)]\) \\
chessboard-small; subpolynomial growth &
\href{https://github.com/yang-mills-group/OpenAI-YM-20261028/blob/3456e5e75cf8d4cd558955a5c7638fba2d68f464/YM/Bid/Drops.lean\#L25}{\texttt{close}},
\href{https://github.com/yang-mills-group/OpenAI-YM-20261028/blob/3456e5e75cf8d4cd558955a5c7638fba2d68f464/YM/Bid/Drops.lean\#L18-L19}{\texttt{dollar}}
& --- \\
projections; relay data &
\href{https://github.com/yang-mills-group/OpenAI-YM-20261028/blob/3456e5e75cf8d4cd558955a5c7638fba2d68f464/YM/Dash/Scope.lean\#L119}{\texttt{Dash.Bin}},
\href{https://github.com/yang-mills-group/OpenAI-YM-20261028/blob/3456e5e75cf8d4cd558955a5c7638fba2d68f464/YM/Bid/Drops.lean\#L46-L52}{\texttt{RelayBin}}
& \(\pi_d(N,p)\to\pi_d\); pace, transit, mark \\
null labels; label quotient &
\href{https://github.com/yang-mills-group/OpenAI-YM-20261028/blob/3456e5e75cf8d4cd558955a5c7638fba2d68f464/YM/Dash/Scope.lean\#L108}{\texttt{Dash.monitor}},
\href{https://github.com/yang-mills-group/OpenAI-YM-20261028/blob/3456e5e75cf8d4cd558955a5c7638fba2d68f464/YM/Euclid/Null.lean\#L34}{\texttt{Table.null}},
\href{https://github.com/yang-mills-group/OpenAI-YM-20261028/blob/3456e5e75cf8d4cd558955a5c7638fba2d68f464/YM/Euclid/Null.lean\#L74}{\texttt{Table.Field}}
& \(\mathrm{Null}_d\), \(F_d=\mathrm{Label}_d/\mathrm{Null}_d\) \\
derivative labels; derivative plan &
\href{https://github.com/yang-mills-group/OpenAI-YM-20261028/blob/3456e5e75cf8d4cd558955a5c7638fba2d68f464/YM/Dash/Inert.lean\#L63-L65}{\texttt{MonitorOps}},
\href{https://github.com/yang-mills-group/OpenAI-YM-20261028/blob/3456e5e75cf8d4cd558955a5c7638fba2d68f464/YM/Euclid/Differential.lean\#L20-L22}{\texttt{DiffPlan}}
& \(\kappa_{d,e}\), \(\tau_{d,a}\); \(\mathrm{rise}_{d,e}\),
\(\partial_{d,a}\) \\
Euclidean hierarchy &
\href{https://github.com/yang-mills-group/OpenAI-YM-20261028/blob/3456e5e75cf8d4cd558955a5c7638fba2d68f464/YM/Euclid/Words.lean\#L141-L148}{\texttt{Table.Hierarchy}}
& (E1)--(E7) \\
zeroth-order polynomial bound &
\href{https://github.com/yang-mills-group/OpenAI-YM-20261028/blob/3456e5e75cf8d4cd558955a5c7638fba2d68f464/YM/Bid/Phase.lean\#L97-L98}{\texttt{Bid.OneCap}}
& IV.1.5 \\
bundle handed downstream &
\href{https://github.com/yang-mills-group/OpenAI-YM-20261028/blob/3456e5e75cf8d4cd558955a5c7638fba2d68f464/YM/Next/EXTRACTION/Base.lean\#L42-L51}{\texttt{Next.Extract.Base}}
& --- \\
\end{longtable}
}

\subsubsection{7. What remains
unexplained}\label{what-remains-unexplained-19}

\begin{enumerate}
\def\labelenumi{\arabic{enumi}.}
\item
  \textbf{Ultraviolet relay and rank estimates.} The interfaces in
  \href{https://github.com/yang-mills-group/OpenAI-YM-20261028/blob/3456e5e75cf8d4cd558955a5c7638fba2d68f464/YM/MergeProof/Eureka.lean\#L16-L20}{\texttt{MergeProof.Scotia}}
  and
  \href{https://github.com/yang-mills-group/OpenAI-YM-20261028/blob/3456e5e75cf8d4cd558955a5c7638fba2d68f464/YM/MergeProof/Powder.lean\#L19-L47}{\texttt{savings}}
  were checked, together with their use in the dimension argument. The
  upstream proofs of the matching estimates and interacting rank
  comparison in III.10--III.11 were not audited here. In particular, the
  dimension inequality is a separate input; it does not follow from the
  centered root budgets.
\item
  \textbf{Volume-uniform budget.} The supplier is
  \href{https://github.com/yang-mills-group/OpenAI-YM-20261028/blob/3456e5e75cf8d4cd558955a5c7638fba2d68f464/YM/MergeProof/Gardens.lean\#L19-L40}{\texttt{gardena}}
  (III.11). Its proof reduces to earlier estimates on the same family,
  but their relation to the volume-uniform bounds assigned to III.12 was
  not traced in full.
\item
  \textbf{Fixed-volume ultraviolet limit.} The full-spacing upgrade
  required by
  \href{https://github.com/yang-mills-group/OpenAI-YM-20261028/blob/3456e5e75cf8d4cd558955a5c7638fba2d68f464/HierarchyBridge.lean\#L44-L56}{\texttt{wilsonLimit}}
  is located and explained in Remark 2. The upstream construction of the
  full fixed-volume kernels used by that upgrade was not audited here.
\item
  \textbf{Alternative assembly.}
  \href{https://github.com/yang-mills-group/OpenAI-YM-20261028/blob/3456e5e75cf8d4cd558955a5c7638fba2d68f464/YM/Next/EXTRACTION/Base.lean\#L79-L92}{\texttt{Next.Extract.finish}}
  and
  \href{https://github.com/yang-mills-group/OpenAI-YM-20261028/blob/3456e5e75cf8d4cd558955a5c7638fba2d68f464/YM/Next/EXTRACTION/Base.lean\#L111-L122}{\texttt{common}}
  give another conditional route to the same bundle. They require relay
  costs and detector estimates in addition to the centered budgets.
  Their hypotheses were checked; they are not additional proofs of the
  ultraviolet inputs.
\end{enumerate}

\subsection{V.1. Decomposition of the limit into
phases}\label{v.1.-decomposition-of-the-limit-into-phases}

\textbf{Part:} V, phases and reconstruction. \textbf{Lean:}
\href{https://github.com/yang-mills-group/OpenAI-YM-20261028/tree/3456e5e75cf8d4cd558955a5c7638fba2d68f464/YM/OS}{\texttt{YM/OS}},
\href{https://github.com/yang-mills-group/OpenAI-YM-20261028/tree/3456e5e75cf8d4cd558955a5c7638fba2d68f464/YM/Phase}{\texttt{YM/Phase}},
\href{https://github.com/yang-mills-group/OpenAI-YM-20261028/tree/3456e5e75cf8d4cd558955a5c7638fba2d68f464/YM/Ground}{\texttt{YM/Ground}},
\href{https://github.com/yang-mills-group/OpenAI-YM-20261028/tree/3456e5e75cf8d4cd558955a5c7638fba2d68f464/YM/Heat}{\texttt{YM/Heat}},
\href{https://github.com/yang-mills-group/OpenAI-YM-20261028/tree/3456e5e75cf8d4cd558955a5c7638fba2d68f464/YM/SNAG}{\texttt{YM/SNAG}},
\href{https://github.com/yang-mills-group/OpenAI-YM-20261028/tree/3456e5e75cf8d4cd558955a5c7638fba2d68f464/YM/Paley}{\texttt{YM/Paley}},
\href{https://github.com/yang-mills-group/OpenAI-YM-20261028/tree/3456e5e75cf8d4cd558955a5c7638fba2d68f464/YM/Spectral}{\texttt{YM/Spectral}},
the work items B01--B18 in
\href{https://github.com/yang-mills-group/OpenAI-YM-20261028/tree/3456e5e75cf8d4cd558955a5c7638fba2d68f464/YM/Work}{\texttt{YM/Work}},
and the assembly files
\href{https://github.com/yang-mills-group/OpenAI-YM-20261028/blob/3456e5e75cf8d4cd558955a5c7638fba2d68f464/YM/Bid/Phase.lean}{\texttt{Bid/Phase.lean}},
\href{https://github.com/yang-mills-group/OpenAI-YM-20261028/blob/3456e5e75cf8d4cd558955a5c7638fba2d68f464/YM/Bid/Fetch.lean}{\texttt{Bid/Fetch.lean}},
\href{https://github.com/yang-mills-group/OpenAI-YM-20261028/blob/3456e5e75cf8d4cd558955a5c7638fba2d68f464/YM/Bid/Market.lean}{\texttt{Bid/Market.lean}},
\href{https://github.com/yang-mills-group/OpenAI-YM-20261028/blob/3456e5e75cf8d4cd558955a5c7638fba2d68f464/YM/Org/Phase.lean}{\texttt{Org/Phase.lean}}.
\textbf{Size:} about 5,700 lines in declarations used by the main
theorem, including about 2,700 in B01--B18, plus the assembly files.

\subsubsection{1. What this unit proves}\label{what-this-unit-proves-20}

\textbf{Setting and notation.}

\begin{itemize}
\tightlist
\item
  \textbf{Tables and words.} \(S\) is the limit table of IV.1 (the mixed
  state \(S^{\mathrm{mix}}\) of the main theorem). A \emph{label} is a
  pair \((d,v)\): its degree \(d\in\mathbb N\) and a vector \(v\) in the
  finite-dimensional reference space \(V_d\). A \emph{word}
  \(w=((x_1,v_1),\dots,(x_n,v_n))\) is a finite list of points of
  \(\mathbb R^4\) with labels; \(|w|=n\), and \(S(w)=S(x;v)\) is its
  value, multilinear in the labels.
\item
  \textbf{Time, slabs, buffers.} Time is the last coordinate \(t(x)\),
  with unit vector \(e_t\); \(y\perp e_t\) denotes a spatial vector. For
  a different coordinate axis \(a\), replace \(t(x)\) by \(x^a\) in the
  following definitions.

  \begin{itemize}
  \tightlist
  \item
    A word is \emph{separated} if its points are distinct, and
    \emph{positive} if in addition all \(t(x_i)>0\).
  \item
    It lies in the \emph{slab} \((0,d)\) if \(0<t(x_i)<d\) for all
    \(i\).
  \item
    It is \emph{\(r\)-buffered} in that slab if its points are at mutual
    sup-distance \(\ge r\) and \(r\le t(x_i)\le d-r\).
  \end{itemize}
\item
  \textbf{Operations on words.} \(w+a\) is the translate,
  \(\tau_dw=w+de_t\), \(w\cdot w'\) is concatenation, and \(\bar w\)
  conjugates only the labels.
\item
  \textbf{Two reflections.}

  \begin{itemize}
  \tightlist
  \item
    \(\Theta w\), the OS reflection: reverse the list, reflect the
    points in \(\{t=0\}\), and apply to the labels the antilinear map
    (conjugation followed by the time-mirror action).
  \item
    \(\vartheta w\), the linear time reflection: reflect the points and
    apply the time-mirror action to the labels, with neither conjugation
    nor reversal.
  \item
    For real labels, \(\Theta w\) is a permutation of \(\vartheta w\).
  \end{itemize}
\item
  \textbf{Null spaces.} \(N_d(S)\subset V_d\) consists of labels whose
  correlations with every separated family of spectators vanish. Write
  \(F_d(S)=V_d/N_d(S)\) and \(q_d\colon V_d\to F_d(S)\) for the quotient
  map. A table \(S'\) \emph{respects} \(S\) if \(N_d(S)\subset N_d(S')\)
  for all \(d\).
\item
  \textbf{Field algebra (VI.1).} It consists of label derivatives
  \(\partial_j\colon F_d(S)\to F_{d+1}(S)\), inclusions
  \(F_d(S)\to F_e(S)\) for \(d\le e\), an action \(R\) of \(O(4)\) on
  these spaces, and a stress tensor \(J_{ab}\in F_4(S)\). In a word,
  \(\partial_jv\) and \(R(K)v\) mean representatives of
  \(\partial_jq_d(v)\) and \(R(K)q_d(v)\). All tables and operators to
  which this notation is applied factor through the relevant quotients.
  Covariance means \(S(Kx_1,\dots,Kx_n;R(K)v_1,\dots,R(K)v_n)=S(x;v)\)
  for every \(K\in O(4)\) and separated configuration
  \(x\in\mathrm{Conf}_n\).
\item
  \textbf{Caps.} For each fixed arity, degree list, label list and
  derivative order \(j\), there are \(A\in\mathbb N\) and \(C\ge0\) such
  that \(\|D^jS(x;v)\|\le Cr^{-A}\) whenever \(0<r\le1\) and the points
  have mutual sup-distance at least \(r\).
\item
  \textbf{The OS space.} The OS space \(\mathcal H\) of \(S\) for an
  axis \(a\) is the completion of the span of \(a\)-positive words
  modulo null vectors, with \(\langle b,b'\rangle=S(\Theta b\cdot b')\).
  \(|b\rangle\) denotes the vector of \(b\), and
  \(\Omega=|\emptyset\rangle\), with \(\|\Omega\|=1\).
\end{itemize}

Nowhere does the unit assume clustering, a unique vacuum or a spectral
gap.

\textbf{Proposition V.1.1 (chambers and ground space).} Let \(S\)
satisfy the Euclidean hierarchy and let \(a\) be a coordinate axis.
Then:

\begin{itemize}
\tightlist
\item
  On the OS space \(\mathcal H\) there is a strongly continuous
  semigroup \((E_t)_{t\ge0}\) of self-adjoint contractions with
  \(E_t|b\rangle=|b+te_a\rangle\).
\item
  There is a strongly continuous unitary group \((U_y)_{y\perp e_a}\)
  commuting with it, with \(U_y|b\rangle=|b+y\rangle\).
\item
  The triple is a \emph{chamber}. Let
  \(\mathcal H_0=\{\psi:E_t\psi=\psi\ \forall t\}\) be its \emph{ground
  space} and \(P_0\) the orthogonal projection onto it. Then
  \(\Omega\in\mathcal H_0\), and \(E_t\psi\to P_0\psi\) as
  \(t\to\infty\) for every \(\psi\).
\end{itemize}

\emph{Status:} theorem. \emph{Lean:}
\href{https://github.com/yang-mills-group/OpenAI-YM-20261028/blob/3456e5e75cf8d4cd558955a5c7638fba2d68f464/YM/OS/Dynamics.lean\#L106-L107}{\texttt{OS.realization}},
\href{https://github.com/yang-mills-group/OpenAI-YM-20261028/blob/3456e5e75cf8d4cd558955a5c7638fba2d68f464/YM/OS/Dynamics.lean\#L63-L104}{\texttt{OS.implement}},
\href{https://github.com/yang-mills-group/OpenAI-YM-20261028/blob/3456e5e75cf8d4cd558955a5c7638fba2d68f464/YM/OS/Core.lean\#L38-L78}{\texttt{OS.Semigroup.emerges}};
definitions
\href{https://github.com/yang-mills-group/OpenAI-YM-20261028/blob/3456e5e75cf8d4cd558955a5c7638fba2d68f464/YM/OS/Dynamics.lean\#L57-L61}{\texttt{OS.Chamber}},
\href{https://github.com/yang-mills-group/OpenAI-YM-20261028/blob/3456e5e75cf8d4cd558955a5c7638fba2d68f464/YM/OS/Flow.lean\#L94-L106}{\texttt{OS.Semigroup}},
\href{https://github.com/yang-mills-group/OpenAI-YM-20261028/blob/3456e5e75cf8d4cd558955a5c7638fba2d68f464/YM/OS/Core.lean\#L14}{\texttt{Semigroup.ground}}.

\textbf{Proposition V.1.2 (spectral calculus in a chamber).} In a
chamber, write formally \(E_t=e^{-tH}\).

\begin{itemize}
\tightlist
\item
  \begin{enumerate}
  \def\labelenumi{(\alph{enumi})}
  \tightlist
  \item
    There is a strongly continuous unitary group \(s\mapsto e^{isH}\),
    determined by \(e^{isH}E_1\psi=E_1^{\,1-is}\psi\). Here \(E_1^z\),
    for \(\mathrm{Re}\,z>0\), is the continuous functional calculus of
    the positive contraction \(E_1\). The map
    \(V(x)=e^{i\,x^aH}U_{x_\perp}\), with \(x_\perp=x-x^ae_a\), is a
    strongly continuous unitary representation of \(\mathbb R^4\).
  \end{enumerate}
\item
  (b) There is a projection-valued measure \(P\) on \(\mathbb R^4\) with
  \(\langle\psi,V(x)\psi\rangle=\int e^{ip\cdot x}\,d\mu_\psi(p)\),
  where \(\mu_\psi(B)=\langle\psi,P(B)\psi\rangle\) for Borel \(B\).
\item
  (c) Every \(\mu_\psi\) is supported in \(\{p^a\ge0\}\).
\item
  (d) For \(\mathrm{Im}\,z\ge0\):
  \[\langle\psi,V(x)\,e^{i(\mathrm{Re}\,z)H}E_{\mathrm{Im}\,z}\,\psi\rangle=\int e^{ip\cdot x}\,e^{ip^az}\,d\mu_\psi(p).\]
\end{itemize}

\emph{Status:} theorem. It holds for every semigroup with the properties
of V.1.1. No gap is assumed or proved. \emph{Lean:}
\href{https://github.com/yang-mills-group/OpenAI-YM-20261028/blob/3456e5e75cf8d4cd558955a5c7638fba2d68f464/YM/Heat/Imaginary.lean\#L49}{\texttt{Semigroup.edge}},
\href{https://github.com/yang-mills-group/OpenAI-YM-20261028/blob/3456e5e75cf8d4cd558955a5c7638fba2d68f464/YM/Heat/Lift.lean\#L50-L58}{\texttt{Semigroup.unit}},
\href{https://github.com/yang-mills-group/OpenAI-YM-20261028/blob/3456e5e75cf8d4cd558955a5c7638fba2d68f464/YM/SNAG/Escape.lean\#L52-L119}{\texttt{Spectral.Orbit.total}},
\href{https://github.com/yang-mills-group/OpenAI-YM-20261028/blob/3456e5e75cf8d4cd558955a5c7638fba2d68f464/YM/SNAG/Escape.lean\#L50}{\texttt{Orbit.P}},
\href{https://github.com/yang-mills-group/OpenAI-YM-20261028/blob/3456e5e75cf8d4cd558955a5c7638fba2d68f464/YM/Heat/Lift.lean\#L102-L104}{\texttt{Semigroup.nonneg}},
\href{https://github.com/yang-mills-group/OpenAI-YM-20261028/blob/3456e5e75cf8d4cd558955a5c7638fba2d68f464/YM/Heat/Lift.lean\#L108-L114}{\texttt{Semigroup.heat}}.

\textbf{Proposition V.1.3 (slab operators).} Assume the hierarchy and
the bench (V.1.7). Then every chamber of \(S\) for the time axis carries
a unique family of bounded operators \(T_d(w)\), one for each \(d>0\)
and each word \(w\) in the slab \((0,d)\), with
\[T_d(w)\,|b\rangle=|w\cdot\tau_d b\rangle\qquad\text{for every positive word }b .\]
- \textbf{All-copy bound.} For every \(D\in\mathbb N\) and \(\rho>0\)
there are \(C>0\) and \(\kappa\in\mathbb N\) such that
\[\|T_d(w)\|\le C^{|w|}\big(|w|/r\big)^{\kappa|w|}\] for every \(d>0\),
every \(r\in(0,1]\), and every word \(w\) that is \(r\)-buffered in
\((0,d)\) with degrees \(\le D\) and label norms \(\le\rho\). The right
side is \(1\) for the empty word.

\begin{itemize}
\tightlist
\item
  \textbf{Smoothness.} If moreover \(S\) satisfies the derivative
  identities of the plan \(\partial\), then for fixed \(d\) and labels
  the map \(x\mapsto T_d(w(x))\) is \(C^\infty\) in operator norm on the
  open set where \(w(x)\) lies in the slab. Its derivatives are
  \(\partial_{x_p^j}T_d(w(x))=T_d(w'(x))\), where \(w'\) replaces the
  label \(v_p\) by \(\partial_jv_p\).
\end{itemize}

\emph{Status:} theorem, given the bench. More generally, smoothness
follows for any slab family with the all-copy bound on a time-axis
chamber of a table \(S'\) that satisfies the hierarchy, respects \(S\),
and satisfies the derivative identities of the plan on \(S\).
\emph{Lean:}
\href{https://github.com/yang-mills-group/OpenAI-YM-20261028/blob/3456e5e75cf8d4cd558955a5c7638fba2d68f464/YM/Work/B01.lean\#L50-L71}{\texttt{Work.B01.answer}},
\href{https://github.com/yang-mills-group/OpenAI-YM-20261028/blob/3456e5e75cf8d4cd558955a5c7638fba2d68f464/YM/Work/B02.lean\#L17-L30}{\texttt{Work.B02.answer}};
definitions
\href{https://github.com/yang-mills-group/OpenAI-YM-20261028/blob/3456e5e75cf8d4cd558955a5c7638fba2d68f464/YM/OS/Slabs.lean\#L101-L104}{\texttt{Spread}},
\href{https://github.com/yang-mills-group/OpenAI-YM-20261028/blob/3456e5e75cf8d4cd558955a5c7638fba2d68f464/YM/OS/Slabs.lean\#L111-L113}{\texttt{Spread.AllCopy}},
\href{https://github.com/yang-mills-group/OpenAI-YM-20261028/blob/3456e5e75cf8d4cd558955a5c7638fba2d68f464/YM/OS/Slabs.lean\#L153-L162}{\texttt{Spread.Laws}}.

\textbf{Proposition V.1.4 (the compressed vacuum algebra).} Assume the
hierarchy and \(O(4)\)-covariance of \(S\). Fix a time-axis chamber and
any family \(T_d(w)\) with the defining identity of V.1.3. Then:

\begin{itemize}
\tightlist
\item
  \begin{enumerate}
  \def\labelenumi{(\alph{enumi})}
  \tightlist
  \item
    \(U_y\psi=\psi\) for every \(\psi\in\mathcal H_0\) and every spatial
    \(y\).
  \end{enumerate}
\item
  (b) For a separated word \(w\), the operator
  \(M(w):=P_0\,T_d(w+s)|_{\mathcal H_0}\) does not depend on the choice
  of \(s\in\mathbb R^4\) and \(d>0\) with \(w+s\) in the slab \((0,d)\).
  For positive \(b\), \(M(b)\Omega=P_0|b\rangle\). For every separated
  \(w\), \(\langle\Omega,M(w)\Omega\rangle=S(w)\).
\item
  (c) Cluster limit. Let \(q_1,\dots,q_m\) be separated words and
  \(a_1(n),\dots,a_m(n)\in\mathbb R^4\) with
  \(t(a_j(n)-a_i(n))\to+\infty\) for \(i<j\); the spatial components are
  arbitrary. Then for large \(n\) the word
  \((q_1+a_1(n))\cdots(q_m+a_m(n))\) is separated, and its value under
  \(S\) tends to \(\langle\Omega,M(q_1)\cdots M(q_m)\Omega\rangle\).
\item
  (d) The norm-closed \(*\)-algebra
  \(\mathcal A\subset B(\mathcal H_0)\) generated by the \(M(w)\) is
  commutative, \(\mathcal A\Omega\) is dense in \(\mathcal H_0\), and
  \(p\in\mathcal A\) with \(p\Omega=0\) implies \(p=0\).
\item
  (e) Suppose in addition that the family satisfies the smoothness part
  of V.1.3. Then:

  \begin{itemize}
  \tightlist
  \item
    \(M\) is invariant under permutations of the word and under all
    translations; \(M(\emptyset)=1\); a unit label may be dropped
    together with its point; and \(M(w)^*=M(\Theta w)\).
  \item
    For \(x\in\mathrm{Conf}_n\), the map \(v\mapsto M(x;v)\) is
    multilinear and factors through
    \(F_{d_1}(S)\times\cdots\times F_{d_n}(S)\).
  \item
    \(x\mapsto M(x;v)\) is \(C^\infty\) in norm on \(\mathrm{Conf}_n\).
  \item
    \(\mathcal A\) is norm-separable.
  \end{itemize}
\end{itemize}

\emph{Status:} theorem. \emph{Lean:}
\href{https://github.com/yang-mills-group/OpenAI-YM-20261028/blob/3456e5e75cf8d4cd558955a5c7638fba2d68f464/YM/Ground/Lean.lean\#L185-L188}{\texttt{Ground.inert}},
\href{https://github.com/yang-mills-group/OpenAI-YM-20261028/blob/3456e5e75cf8d4cd558955a5c7638fba2d68f464/YM/Ground/Compress.lean\#L111-L119}{\texttt{Ground.overline}},
\href{https://github.com/yang-mills-group/OpenAI-YM-20261028/blob/3456e5e75cf8d4cd558955a5c7638fba2d68f464/YM/Ground/Compress.lean\#L128}{\texttt{Ground.almost}},
\href{https://github.com/yang-mills-group/OpenAI-YM-20261028/blob/3456e5e75cf8d4cd558955a5c7638fba2d68f464/YM/Phase/Mixture.lean\#L59-L63}{\texttt{Spread.insertValue}},
\href{https://github.com/yang-mills-group/OpenAI-YM-20261028/blob/3456e5e75cf8d4cd558955a5c7638fba2d68f464/YM/Ground/Moments.lean\#L138-L148}{\texttt{Ground.arcs}},
\href{https://github.com/yang-mills-group/OpenAI-YM-20261028/blob/3456e5e75cf8d4cd558955a5c7638fba2d68f464/YM/Ground/Permute.lean\#L149-L161}{\texttt{Ground.adjunction}},
\href{https://github.com/yang-mills-group/OpenAI-YM-20261028/blob/3456e5e75cf8d4cd558955a5c7638fba2d68f464/YM/Ground/Compress.lean\#L145-L162}{\texttt{Ground.full}},
\href{https://github.com/yang-mills-group/OpenAI-YM-20261028/blob/3456e5e75cf8d4cd558955a5c7638fba2d68f464/YM/Ground/Spectrum.lean\#L115-L124}{\texttt{Ground.attains}},
\href{https://github.com/yang-mills-group/OpenAI-YM-20261028/blob/3456e5e75cf8d4cd558955a5c7638fba2d68f464/YM/Work/B03.lean\#L405-L415}{\texttt{Work.B03.answer}},
\href{https://github.com/yang-mills-group/OpenAI-YM-20261028/blob/3456e5e75cf8d4cd558955a5c7638fba2d68f464/YM/Work/B12.lean\#L81-L92}{\texttt{Work.B12.answer}};
definitions
\href{https://github.com/yang-mills-group/OpenAI-YM-20261028/blob/3456e5e75cf8d4cd558955a5c7638fba2d68f464/YM/Ground/Compress.lean\#L89-L91}{\texttt{Ground.M}},
\href{https://github.com/yang-mills-group/OpenAI-YM-20261028/blob/3456e5e75cf8d4cd558955a5c7638fba2d68f464/YM/Phase/Mixture.lean\#L54}{\texttt{Spread.algebra}},
\href{https://github.com/yang-mills-group/OpenAI-YM-20261028/blob/3456e5e75cf8d4cd558955a5c7638fba2d68f464/YM/Bid/Phase.lean\#L45-L56}{\texttt{Bid.Theme}}.

\textbf{Proposition V.1.5 (decomposition into phases).} Assume the
setting of V.1.4(a)--(e). Then:

\begin{itemize}
\tightlist
\item
  There are a compact metrizable space \(X\) (the Gelfand spectrum of
  \(\mathcal A\)), a Borel probability measure \(\nu\) on \(X\) with
  \(\nu(O)>0\) for every nonempty open \(O\), and characters
  \(\chi_\xi\) (\(\xi\in X\)).
\item
  \(\xi\mapsto\chi_\xi\) is a bijection onto the characters of
  \(\mathcal A\); each \(\xi\mapsto\chi_\xi(p)\) is continuous; and
  \(\int_X\chi_\xi(p)\,d\nu(\xi)=\langle\Omega,p\Omega\rangle\) for
  every \(p\in\mathcal A\).
\item
  For each \(\xi\), the formula \(S^\xi(x;v)=\chi_\xi(M(x;v))\) on
  \(\mathrm{Conf}_n\) (and \(0\) off it) defines a table.
\item
  \(S^\xi(w)=\chi_\xi(M(w))\) for every separated \(w\);
  \(\xi\mapsto S^\xi(w)\) is continuous; and the barycenter identity
  holds:
  \[\int_XS^\xi(w)\,d\nu(\xi)=S(w)\qquad\text{for every separated word }w.\]
\end{itemize}

\emph{Status:} theorem. \emph{Lean:}
\href{https://github.com/yang-mills-group/OpenAI-YM-20261028/blob/3456e5e75cf8d4cd558955a5c7638fba2d68f464/YM/Ground/Spectrum.lean\#L126-L140}{\texttt{Ground.resolution}},
\href{https://github.com/yang-mills-group/OpenAI-YM-20261028/blob/3456e5e75cf8d4cd558955a5c7638fba2d68f464/YM/Work/B07.lean\#L12-L13}{\texttt{Work.B07.state}},
\href{https://github.com/yang-mills-group/OpenAI-YM-20261028/blob/3456e5e75cf8d4cd558955a5c7638fba2d68f464/YM/Phase/Mixture.lean\#L113-L117}{\texttt{Resolution.barycenter}};
definitions
\href{https://github.com/yang-mills-group/OpenAI-YM-20261028/blob/3456e5e75cf8d4cd558955a5c7638fba2d68f464/YM/Phase/Mixture.lean\#L65-L72}{\texttt{Spread.Resolution}},
\href{https://github.com/yang-mills-group/OpenAI-YM-20261028/blob/3456e5e75cf8d4cd558955a5c7638fba2d68f464/YM/Phase/Mixture.lean\#L78-L86}{\texttt{Resolution.Out}},
\href{https://github.com/yang-mills-group/OpenAI-YM-20261028/blob/3456e5e75cf8d4cd558955a5c7638fba2d68f464/YM/Phase/Mixture.lean\#L89-L91}{\texttt{Resolution.receives}}.

\textbf{Proposition V.1.6 (every phase satisfies the hierarchy and the
field laws).} In addition to V.1.3--V.1.5 (with the smoothness laws for
the plan of the field algebra), assume:

\begin{itemize}
\tightlist
\item
  the Seed of \(S\) (IV.1);
\item
  the grid hypothesis (§2);
\item
  the laws of the field algebra for \(S\) (VI.1).
\end{itemize}

Then two operator identities hold before any character is applied:

\begin{itemize}
\tightlist
\item
  (e) \emph{phase reflection}: \(M(w)=M(\vartheta w)\) for every
  separated word with arbitrary complex labels;
\item
  (f) \(M(w)=M(Kw)\) for every separated \(w\) and every \(K\in O(4)\),
  where \(K\) acts on points and, through \(R(K)\), on labels.
\end{itemize}

As a consequence, for \textbf{every} \(\xi\in X\):

\begin{itemize}
\item
  \begin{enumerate}
  \def\labelenumi{(\alph{enumi})}
  \tightlist
  \item
    \(S^\xi\) satisfies the Euclidean hierarchy:
  \end{enumerate}

  \begin{itemize}
  \tightlist
  \item
    permutation symmetry, \(S^\xi(\emptyset)=1\), and unit labels;
  \item
    invariance under signed permutations of the axes, and reality
    \(S^\xi(x;\bar v)=\overline{S^\xi(x;v)}\);
  \item
    translation invariance;
  \item
    reflection positivity for every coordinate axis, in point and
    smeared form;
  \item
    smoothness on \(\mathrm{Conf}_n\), and caps on all derivatives.
  \end{itemize}
\item
  (b) \(S^\xi\) respects \(S\) and is \(O(4)\)-covariant for the same
  action \(R\). It satisfies the inclusion and derivative identities of
  the plan and the stress-tensor Ward identities of VI.1.
\item
  (c) Uniform cap: for every \(D\in\mathbb N\) and \(\rho>0\) there are
  \(K>0\) and \(\kappa\in\mathbb N\) with
  \[|S^\xi(w)|\le K^{|w|}\big(|w|/r\big)^{\kappa|w|}\] for all \(\xi\),
  all \(r\in(0,1]\), and all words with degrees \(\le D\), label norms
  \(\le\rho\) and points at mutual sup-distance \(\ge r\).
\item
  (d) \(N_d(S^\xi)=N_d(S)\) for every \(d\), hence
  \(F_d(S^\xi)=F_d(S)\).
\end{itemize}

\emph{Status:} theorem. \emph{Lean:}
\href{https://github.com/yang-mills-group/OpenAI-YM-20261028/blob/3456e5e75cf8d4cd558955a5c7638fba2d68f464/YM/Work/B05.lean\#L90-L115}{\texttt{Work.B05.answer}},
\href{https://github.com/yang-mills-group/OpenAI-YM-20261028/blob/3456e5e75cf8d4cd558955a5c7638fba2d68f464/YM/Work/B06.lean\#L238-L242}{\texttt{B06.answer}},
\href{https://github.com/yang-mills-group/OpenAI-YM-20261028/blob/3456e5e75cf8d4cd558955a5c7638fba2d68f464/YM/Work/B07.lean\#L119-L124}{\texttt{B07.answer}},
\href{https://github.com/yang-mills-group/OpenAI-YM-20261028/blob/3456e5e75cf8d4cd558955a5c7638fba2d68f464/YM/Work/B08.lean\#L206-L209}{\texttt{B08.answer}},
\href{https://github.com/yang-mills-group/OpenAI-YM-20261028/blob/3456e5e75cf8d4cd558955a5c7638fba2d68f464/YM/Work/B09.lean\#L12-L15}{\texttt{B09.answer}},
\href{https://github.com/yang-mills-group/OpenAI-YM-20261028/blob/3456e5e75cf8d4cd558955a5c7638fba2d68f464/YM/Work/B10.lean\#L95-L105}{\texttt{B10.answer}},
\href{https://github.com/yang-mills-group/OpenAI-YM-20261028/blob/3456e5e75cf8d4cd558955a5c7638fba2d68f464/YM/Work/B11.lean\#L75-L89}{\texttt{B11.answer}},
\href{https://github.com/yang-mills-group/OpenAI-YM-20261028/blob/3456e5e75cf8d4cd558955a5c7638fba2d68f464/YM/Phase/Picks.lean\#L128-L158}{\texttt{Retain.lift\_detect}};
definitions
\href{https://github.com/yang-mills-group/OpenAI-YM-20261028/blob/3456e5e75cf8d4cd558955a5c7638fba2d68f464/YM/Org/Phase.lean\#L19-L20}{\texttt{Phase.reflects}},
\href{https://github.com/yang-mills-group/OpenAI-YM-20261028/blob/3456e5e75cf8d4cd558955a5c7638fba2d68f464/YM/Bid/Phase.lean\#L72-L80}{\texttt{Bid.Garland}},
\href{https://github.com/yang-mills-group/OpenAI-YM-20261028/blob/3456e5e75cf8d4cd558955a5c7638fba2d68f464/YM/Org/Phase.lean\#L25-L32}{\texttt{Phase.Scripts}}.

\textbf{Proposition V.1.7 (bench).} For every \(D\in\mathbb N\) and
\(\rho>0\) there are \(C_0>0\) and \(\kappa\in\mathbb N\) with the
following property. Let \(d>0\), \(r\in(0,1]\), and let
\(w\neq\emptyset\) be \(r\)-buffered in \((0,d)\), with degrees
\(\le D\) and label norms \(\le\rho\). Then there are a lattice level
\(b\) and a width \(s\le d\) with \(s\in a_b\mathbb N\) and \(w\) in the
slab \((0,s)\), such that every positive word \(v\) has a constant
\(B_v\) with
\[\big|S(\Theta v_n\cdot v_n)\big|\le B_v\,A^{4n}\ \ (n\ge0),\qquad A=C_0^{|w|}\big(|w|/r\big)^{\kappa|w|},\qquad v_0=v,\ \ v_{n+1}=\tau_s(\Theta w)\cdot\tau_s(w\cdot\tau_s v_n).\]
Here \(v_n\) is \(v\) preceded by \(2n\) consecutive slabs of width
\(s\), holding alternately \(w\) and its reflection. Thus
\(|v_n\rangle=(T_s(w)^*T_s(w))^n|v\rangle\) once \(T_s(w)\) exists, and
\(S(\Theta v_n\cdot v_n)\) is a reflection-symmetric train of \(4n\)
copies of \(w\).

\emph{Status:} a hypothesis of this unit, required of every limit table
with a Seed. It is proved in VII.4
(\href{https://github.com/yang-mills-group/OpenAI-YM-20261028/blob/3456e5e75cf8d4cd558955a5c7638fba2d68f464/YM/Next/PHASE/Seats.lean\#L154-L160}{\texttt{Next.Phase.seats}})
from the cutoff family, the hierarchy, the point limits and the cell
hypothesis of the Seed.

\begin{itemize}
\tightlist
\item
  \textbf{Chessboard estimate.} The cell hypothesis bounds every signed
  or reflected field of degree \(\le D\) by \(J\,a_k^{-(D+1)}\) in
  chessboard norm. The scale can be chosen with \(a_k\ge c\,(r/|w|)^4\),
  where \(c>0\) depends only on the fixed family and \((D,\rho)\). A
  finer scale handles the exterior word. The resulting train bound is
  \(B^{2|v|}(Ja_k^{-(D+1)})^{4n|w|}\)
  (\href{https://github.com/yang-mills-group/OpenAI-YM-20261028/blob/3456e5e75cf8d4cd558955a5c7638fba2d68f464/YM/Next/PHASE/Trains.lean\#L273-L304}{\texttt{Next.Phase.travel\_price}}).
  Thus one may take \(\kappa=4(D+1)\); the exterior constant is
  independent of \(n\).
\end{itemize}

\emph{Lean:}
\href{https://github.com/yang-mills-group/OpenAI-YM-20261028/blob/3456e5e75cf8d4cd558955a5c7638fba2d68f464/YM/Bid/Phase.lean\#L21-L24}{\texttt{Bid.Bench}},
\href{https://github.com/yang-mills-group/OpenAI-YM-20261028/blob/3456e5e75cf8d4cd558955a5c7638fba2d68f464/YM/Bid/Phase.lean\#L27}{\texttt{Bid.Seats}},
\href{https://github.com/yang-mills-group/OpenAI-YM-20261028/blob/3456e5e75cf8d4cd558955a5c7638fba2d68f464/YM/OS/Train.lean\#L70-L71}{\texttt{OS.QRead}},
\href{https://github.com/yang-mills-group/OpenAI-YM-20261028/blob/3456e5e75cf8d4cd558955a5c7638fba2d68f464/YM/Next/PHASE/Seats.lean\#L119-L151}{\texttt{Next.Phase.replicas}},
\href{https://github.com/yang-mills-group/OpenAI-YM-20261028/blob/3456e5e75cf8d4cd558955a5c7638fba2d68f464/YM/Next/PHASE/Seats.lean\#L55-L77}{\texttt{Next.Phase.rates}}.

\textbf{Proposition V.1.8 (vacuum uniqueness in almost every phase).}
Let \((X,\nu,\chi)\) and \(S^\xi\) be as in V.1.5, with every \(S^\xi\)
satisfying the hierarchy. For each \(\xi\) fix a time-axis chamber of
\(S^\xi\), with vacuum \(\Omega_\xi\) and ground space
\(\mathcal H_0^\xi\). Then:

\begin{itemize}
\tightlist
\item
  \(\mathcal H_0^\xi=\mathbb C\,\Omega_\xi\) for \(\nu\)-almost every
  \(\xi\); in particular such phases exist.
\item
  Every time-axis chamber of every \(S^\xi\) carries a slab family
  \(T^\xi_d(w)\) with \(\|T^\xi_d(w)\|\le\|T_d(w)\|\) for each valid
  slab word. If the original family has the all-copy bound, so does this
  family. Under V.1.6 it also has the smoothness laws.
\end{itemize}

\emph{Status:} theorem. Uniqueness for every phase is proved later, in
V.3
(\href{https://github.com/yang-mills-group/OpenAI-YM-20261028/blob/3456e5e75cf8d4cd558955a5c7638fba2d68f464/YM/EveryPhase/Vacuum.lean\#L146-L153}{\texttt{EveryPhase.every}}),
by an argument outside this unit. \emph{Lean:}
\href{https://github.com/yang-mills-group/OpenAI-YM-20261028/blob/3456e5e75cf8d4cd558955a5c7638fba2d68f464/YM/Phase/Unique.lean\#L56-L86}{\texttt{Cluster.exp\_real}},
\href{https://github.com/yang-mills-group/OpenAI-YM-20261028/blob/3456e5e75cf8d4cd558955a5c7638fba2d68f464/YM/Phase/Unique.lean\#L91-L110}{\texttt{Cluster.lifts}},
\href{https://github.com/yang-mills-group/OpenAI-YM-20261028/blob/3456e5e75cf8d4cd558955a5c7638fba2d68f464/YM/Phase/Unique.lean\#L111-L114}{\texttt{Cluster.selects}},
\href{https://github.com/yang-mills-group/OpenAI-YM-20261028/blob/3456e5e75cf8d4cd558955a5c7638fba2d68f464/YM/Work/B13.lean\#L91-L105}{\texttt{Work.B13.answer}},
\href{https://github.com/yang-mills-group/OpenAI-YM-20261028/blob/3456e5e75cf8d4cd558955a5c7638fba2d68f464/YM/Work/B14.lean\#L145-L148}{\texttt{Work.B14.answer}};
definition
\href{https://github.com/yang-mills-group/OpenAI-YM-20261028/blob/3456e5e75cf8d4cd558955a5c7638fba2d68f464/YM/Phase/Mixture.lean\#L35}{\texttt{Chamber.UniqueHeat}}.

\textbf{Proposition V.1.9 (selected-phase certificate).} Let \(C\) be a
Wilson family with extraction order \(f\) and limit table \(S\). Assume:

\begin{itemize}
\tightlist
\item
  the Seed of \(S\);
\item
  the grid hypothesis;
\item
  a field algebra on \(S\) with its laws.
\end{itemize}

Then there are:

\begin{itemize}
\tightlist
\item
  a time-axis chamber of \(S\) whose ground space is fixed pointwise by
  spatial translations;
\item
  a slab family on it with the all-copy bound and the smoothness laws;
\item
  a decomposition \((X,\nu,\chi_\xi,S^\xi)\) as in V.1.5, with every
  \(S^\xi\) satisfying V.1.6(a), (b), (d);
\item
  a phase \(\xi_*\), with a time-axis chamber of \(S^{\xi_*}\) whose
  ground space is \(\mathbb C\Omega_{\xi_*}\), and a slab family on that
  chamber with the smoothness laws.
\end{itemize}

\emph{Status:} theorem.

\begin{itemize}
\tightlist
\item
  \textbf{Assembly.} The bench theorem and B01--B18 give the slab and
  phase constructions
  (\href{https://github.com/yang-mills-group/OpenAI-YM-20261028/blob/3456e5e75cf8d4cd558955a5c7638fba2d68f464/YM/Bid/Market.lean\#L10-L36}{\texttt{Bid.Looms}});
  the selection argument produces the certificate
  (\href{https://github.com/yang-mills-group/OpenAI-YM-20261028/blob/3456e5e75cf8d4cd558955a5c7638fba2d68f464/YM/Org/Phase.lean\#L60-L76}{\texttt{Phase.Loom.cleanse}}).
\item
  \textbf{Use in the main theorem.} The main theorem applies it in
  \href{https://github.com/yang-mills-group/OpenAI-YM-20261028/blob/3456e5e75cf8d4cd558955a5c7638fba2d68f464/YM/Pert/Construction.lean\#L21-L31}{\texttt{Pert.withFinal}}
  and
  \href{https://github.com/yang-mills-group/OpenAI-YM-20261028/blob/3456e5e75cf8d4cd558955a5c7638fba2d68f464/YM/MergeProof/Summit.lean\#L17-L30}{\texttt{MergeProof.Acme.glistens}}.
  The result is the \texttt{phase} field of
  \href{https://github.com/yang-mills-group/OpenAI-YM-20261028/blob/3456e5e75cf8d4cd558955a5c7638fba2d68f464/YM/Spec/Closed.lean\#L64-L73}{\texttt{Spec.Final}}.
\end{itemize}

\emph{Lean:} definitions
\href{https://github.com/yang-mills-group/OpenAI-YM-20261028/blob/3456e5e75cf8d4cd558955a5c7638fba2d68f464/YM/Org/Phase.lean\#L50-L54}{\texttt{Phase.Loom}},
\href{https://github.com/yang-mills-group/OpenAI-YM-20261028/blob/3456e5e75cf8d4cd558955a5c7638fba2d68f464/YM/Org/Phase.lean\#L42-L46}{\texttt{Phase.Stall}},
\href{https://github.com/yang-mills-group/OpenAI-YM-20261028/blob/3456e5e75cf8d4cd558955a5c7638fba2d68f464/YM/Org/Phase.lean\#L34-L40}{\texttt{Phase.Vendor}},
\href{https://github.com/yang-mills-group/OpenAI-YM-20261028/blob/3456e5e75cf8d4cd558955a5c7638fba2d68f464/YM/Phase/Mixture.lean\#L127-L134}{\texttt{OS.Selection}},
\href{https://github.com/yang-mills-group/OpenAI-YM-20261028/blob/3456e5e75cf8d4cd558955a5c7638fba2d68f464/YM/Phase/Mixture.lean\#L136-L148}{\texttt{Selection.Certificate}},
\href{https://github.com/yang-mills-group/OpenAI-YM-20261028/blob/3456e5e75cf8d4cd558955a5c7638fba2d68f464/YM/Spec/Closed.lean\#L45-L47}{\texttt{Spec.Materializes}}.

\subsubsection{2. What it uses}\label{what-it-uses-20}

\begin{itemize}
\tightlist
\item
  \textbf{IV.1 (limit table and Seed).}

  \begin{itemize}
  \tightlist
  \item
    The limit table \(S\) with the Euclidean hierarchy: Lean
    \href{https://github.com/yang-mills-group/OpenAI-YM-20261028/blob/3456e5e75cf8d4cd558955a5c7638fba2d68f464/YM/Euclid/Words.lean\#L141-L148}{\texttt{Table.Hierarchy}},
    with fields basic laws, real-spin, translations, point and smeared
    reflection positivity, smoothness, caps.
  \item
    The Seed, Lean
    \href{https://github.com/yang-mills-group/OpenAI-YM-20261028/blob/3456e5e75cf8d4cd558955a5c7638fba2d68f464/YM/Spec/Closed.lean\#L20-L26}{\texttt{Spec.Seed}}.
    Its components:

    \begin{itemize}
    \tightlist
    \item
      \begin{enumerate}
      \def\labelenumi{(\roman{enumi})}
      \tightlist
      \item
        the cutoff family;
      \end{enumerate}
    \item
      (ii) the ordered limit;
    \item
      (iii) point limits at generic points
      (\href{https://github.com/yang-mills-group/OpenAI-YM-20261028/blob/3456e5e75cf8d4cd558955a5c7638fba2d68f464/YM/Extraction/Limits.lean\#L68}{\texttt{PointLimits}});
    \item
      (iv) the hierarchy;
    \item
      (v) the \emph{cell hypothesis}
      (\href{https://github.com/yang-mills-group/OpenAI-YM-20261028/blob/3456e5e75cf8d4cd558955a5c7638fba2d68f464/YM/Cells/Packing.lean\#L68-L79}{\texttt{CellHyp}}):
      reflection-positivity seams at every cutoff; and, for every degree
      \(d\) and \(\varepsilon,\rho>0\), a threshold \(m_0\) and
      \(A\ge0\) such that for every \(m\ge m_0\) and \(\eta>0\), late in
      the ordered cutoff, every degree-\(d\) label of norm at most
      \(\rho\) is supported inside its centered cell and has chessboard
      root at most \(A\,a_m^{-d-\varepsilon}+\eta\), uniformly over cell
      centers;
    \item
      (vi) the \emph{corrected-field data}
      (\href{https://github.com/yang-mills-group/OpenAI-YM-20261028/blob/3456e5e75cf8d4cd558955a5c7638fba2d68f464/YM/Cells/Corrected.lean\#L138}{\texttt{RankWidths}},
      \href{https://github.com/yang-mills-group/OpenAI-YM-20261028/blob/3456e5e75cf8d4cd558955a5c7638fba2d68f464/YM/Cells/Corrected.lean\#L38-L46}{\texttt{Columns}},
      \href{https://github.com/yang-mills-group/OpenAI-YM-20261028/blob/3456e5e75cf8d4cd558955a5c7638fba2d68f464/YM/Cells/Corrected.lean\#L98-L118}{\texttt{Trials.Valid}}).
      For each degree \(d\): finitely many one- or two-point test rows
      and a basis of the free-field space of degree \(d\), for which the
      free test map is injective; and lattice ``corrected'' fields at
      levels \(k\) whose classes span \(F_d(S)\) for all sufficiently
      large \(k\), whose one-field differences from free values and
      centered two-field products minus free cumulants have chessboard
      bounds \(\mathrm{gap}_k+\eta\) late in the cutoff for every
      \(\eta>0\), with \(\mathrm{gap}_k\to0\). The free values use the
      rescaled test points.
    \end{itemize}
  \item
    The free table
    (\href{https://github.com/yang-mills-group/OpenAI-YM-20261028/blob/3456e5e75cf8d4cd558955a5c7638fba2d68f464/YM/Cells/Corrected.lean\#L22}{\texttt{Cells.standard}}).
  \item
    Where these are used: (i)--(iv) in B04, B16 and the bench; (v) in
    the bench, B15 and the grid hypothesis; (vi) in B18 and
    \texttt{lift\_detect}.
  \end{itemize}
\item
  \textbf{VII.4 (the bench).} V.1.7,
  i.e.~\href{https://github.com/yang-mills-group/OpenAI-YM-20261028/blob/3456e5e75cf8d4cd558955a5c7638fba2d68f464/YM/Next/PHASE/Seats.lean\#L154-L160}{\texttt{Next.Phase.seats}}.
\item
  \textbf{III.11 (the grid hypothesis)}
  \href{https://github.com/yang-mills-group/OpenAI-YM-20261028/blob/3456e5e75cf8d4cd558955a5c7638fba2d68f464/YM/Org/Input.lean\#L121-L123}{\texttt{Org.Grill}}.
  Proved from bounded ultraviolet approximants and the cell hypothesis
  by
  \href{https://github.com/yang-mills-group/OpenAI-YM-20261028/blob/3456e5e75cf8d4cd558955a5c7638fba2d68f464/YM/MergeProof/South/FinalG.lean\#L76-L83}{\texttt{MergeProof.Accent.grilled}}.

  \begin{itemize}
  \tightlist
  \item
    Setup: take a degree bound \(d\) and a level \(m\). There is a level
    \(k\ge m\) such that the following holds for every word \(w\) with
    real labels of degree \(\le d\) at generic points of the level-\(k\)
    grid.
  \item
    Let \(F\) be the lattice product of the fields of \(w\), and
    \(\vartheta F\) the product of the mirrored fields at the reflected
    points.
  \item
    The lattice variables \(f_j\) are the translates of
    \(F-\vartheta F\) by \(t\,j\) grid steps in every coordinate.
  \item
    Conclusion: for each approximation index \(\varkappa\) the \(f_j\)
    admit real measurable approximants, bounded by \(R_\varkappa\)
    uniformly in the cutoff, whose pair integrals, for \(i<j\), differ
    from \(\mathbb E[\mathrm{Re}(f_if_j)]\) by at most
    \(\delta_\varkappa\to0\). For a fixed grid word, \(R_\varkappa\) and
    \(\delta_\varkappa\) are independent of the translation, copy
    indices and cutoffs. The pair estimates hold for all sufficiently
    large \(t\), then late in the cutoff
    (\href{https://github.com/yang-mills-group/OpenAI-YM-20261028/blob/3456e5e75cf8d4cd558955a5c7638fba2d68f464/YM/Phase/Approx.lean\#L63-L67}{\texttt{Grid.Covering}}).
    It is used only in B05.
  \end{itemize}
\item
  \textbf{VI.1 (field algebra and its laws on \(S\))}
  (\href{https://github.com/yang-mills-group/OpenAI-YM-20261028/blob/3456e5e75cf8d4cd558955a5c7638fba2d68f464/YM/Spec/Closed.lean\#L28-L31}{\texttt{Spec.Algebra}},
  \href{https://github.com/yang-mills-group/OpenAI-YM-20261028/blob/3456e5e75cf8d4cd558955a5c7638fba2d68f464/YM/Spec/Closed.lean\#L34-L43}{\texttt{Spec.Locals}}).

  \begin{itemize}
  \tightlist
  \item
    \(O(4)\)-covariance of \(S\)
    (\href{https://github.com/yang-mills-group/OpenAI-YM-20261028/blob/3456e5e75cf8d4cd558955a5c7638fba2d68f464/YM/Euclid/Stress.lean\#L46-L48}{\texttt{Rotation.Oriented}})
    is used in V.1.4(a), (d), in \texttt{Ground.resolution}, and in B06,
    B07.
  \item
    The rotation laws
    (\href{https://github.com/yang-mills-group/OpenAI-YM-20261028/blob/3456e5e75cf8d4cd558955a5c7638fba2d68f464/YM/Euclid/Stress.lean\#L31-L42}{\texttt{Rotation.Laws}}):
    group law, and \(R(s)\) equal to the mirror action of \(S\) for
    signed permutations \(s\). Used in B06, B07.
  \item
    The derivative identities are used in B02 and B09; the inclusion
    identities and stress-tensor laws
    (\href{https://github.com/yang-mills-group/OpenAI-YM-20261028/blob/3456e5e75cf8d4cd558955a5c7638fba2d68f464/YM/Euclid/Stress.lean\#L90-L97}{\texttt{StressLaws}})
    in B09.
  \item
    \textbf{Order of dependence.} VI.1 uses the chamber construction and
    the bounded heat functional calculus from the beginning of this
    unit, as well as elementary geometry and norm estimates. It does not
    use the compressed algebra, phase reflection or phase selection.
    Those constructions then use the \(O(4)\) action supplied by VI.1.
  \end{itemize}
\item
  \textbf{V.2.} The Gram/completion construction of the OS space
  (\href{https://github.com/yang-mills-group/OpenAI-YM-20261028/blob/3456e5e75cf8d4cd558955a5c7638fba2d68f464/YM/Hilbert/Kolmogorov.lean\#L34-L38}{\texttt{Kernel.Gram}},
  with \texttt{by\_bound}, \texttt{by\_bud}).
\item
  \textbf{I.1, II.1.} Lattice, generic points, label registers, the
  Wilson family.
\end{itemize}

\textbf{Standard mathematics.}

\begin{itemize}
\tightlist
\item
  \textbf{Osterwalder--Schrader reconstruction} and the doubling bound
  for the semigroup: Osterwalder--Schrader, CMP 31 (1973), 42 (1975);
  Glimm--Jaffe, \emph{Quantum Physics}, Ch. 19 (probable).
\item
  \textbf{Chessboard estimates:} Fröhlich--Israel--Lieb--Simon, CMP 62
  (1978).
\item
  \textbf{Pure phases.} The classical analogue is the decomposition of a
  Nelson--Symanzik positive theory into pure phases, Fröhlich, Ann.
  Phys. 97 (1976) (probable), and the decomposition of states along
  abelian algebras, Bratteli--Robinson I, Ch. 4.
\item
  \textbf{Operator algebras and measure theory:}

  \begin{itemize}
  \tightlist
  \item
    Gelfand--Naimark for commutative C*-algebras and spectral permanence
    (Bratteli--Robinson I, §2.1--2.2);
  \item
    Riesz--Markov (Reed--Simon I, Thm. IV.14);
  \item
    continuous functional calculus;
  \item
    Hahn--Banach;
  \item
    dominated convergence.
  \end{itemize}
\item
  \textbf{Harmonic analysis:} the Stone--Naimark--Ambrose--Godement
  theorem (Folland, \emph{A Course in Abstract Harmonic Analysis}, §4.4)
  and a Paley--Wiener support theorem for bounded holomorphic functions
  in a half-plane (Reed--Simon II, §IX.3) (probable).
\item
  \textbf{Geometry:} Cartan--Dieudonné (reflections generate \(O(4)\)).
\end{itemize}

\textbf{Imported constants and assumptions.}

\begin{itemize}
\tightlist
\item
  No numerical constant is fixed in the unit.
\item
  The all-copy bound uses the constants of the bench; one may choose
  \(\kappa=4(D+1)\). These constants may depend on the fixed family and
  \((D,\rho)\), but not on the slab width, buffer, word length or number
  of copies.
\item
  The hypotheses are: the hierarchy and the Seed (IV.1), the grid
  hypothesis (III.11), the bench (VII.4), and the field-algebra laws
  (VI.1).
\end{itemize}

\subsubsection{3. Strategy}\label{strategy-20}

\textbf{The limit table as a mixture.} The table \(S\) is obtained by
ordered limits of finite-volume Wilson correlations. Its OS semigroup
may have a multidimensional ground space \(\mathcal H_0\). The unit
decomposes the correlations through an algebra acting on that space.

\begin{itemize}
\tightlist
\item
  \textbf{The algebra.} Inserting a separated word into a time slab
  gives a bounded operator; compressing it to \(\mathcal H_0\) gives
  \(M(w)\). These operators generate a commutative C*-algebra
  \(\mathcal A\), for which \(\Omega\) is cyclic and separating.
\item
  \textbf{Phases.} The Gelfand spectrum \(X\) of \(\mathcal A\), with
  the measure \(\nu\) representing
  \(p\mapsto\langle\Omega,p\Omega\rangle\), is the space of phases. Each
  character gives a table, \(S^\xi=\chi_\xi\circ M\), and
  \(S=\int S^\xi\,d\nu\).
\item
  \textbf{Comparison with the classical case.} The algebra describes
  observables at infinite time separation. It replaces the invariant
  observables used in the ergodic decomposition of a Euclidean measure.
  Here the construction uses only the OS space of the limit
  correlations. Operator identities give the hierarchy and field laws
  for every \(\xi\); vacuum uniqueness is first obtained almost
  everywhere.
\end{itemize}

\textbf{Four inputs.} Four inputs make this work, and each is used where
the classical argument would use the measure.

\begin{itemize}
\tightlist
\item
  \textbf{(i) Boundedness.} A chessboard estimate controls alternating
  trains of \(4n\) copies with constants independent of \(n\). A power
  argument converts this estimate into the slab-operator bound
  \(C^{|w|}(|w|/r)^{\kappa|w|}\).
\item
  \textbf{(ii) \(O(4)\)-covariance} (VI.1). It shows that ground states
  are invariant under spatial translations, which makes \(M(w)\) well
  defined. Spatial displacements are rotated into time displacements,
  which are asymptotically invisible. It also makes \(\mathcal A\)
  commutative: exchanging the time axis with a spatial axis reverses the
  time order of two clusters inside a four-cluster limit.
\item
  \textbf{(iii) Phase reflection} \(M(w)=M(\vartheta w)\). Together with
  \(M(w)^*=M(\Theta w)\), this makes every phase value on real labels
  real. Put \(D=M(w)-M(\vartheta w)\) for such a word. Then \(D^*=-D\),
  so \(\langle\Omega,D^2\Omega\rangle=-\|D\Omega\|^2\). Bounded real
  lattice approximants give the reverse inequality by positivity of the
  variance of many far-separated copies. Thus \(D=0\). Conjugating time
  reflection gives invariance of \(M\) under \(O(4)\).
\item
  \textbf{(iv) Chessboard detectors.} For the corrected fields they show
  that the one- and two-point functions of every phase, on a spanning
  family of fields, are close to free-field values for which the test
  map is injective. Hence every phase has the null spaces of \(S\).
\end{itemize}

\textbf{Selecting a phase.} For a positive word \(b\), the quantity
\(\|P_0^\xi|b\rangle\|^2-|S^\xi(b)|^2\ge0\) has \(\nu\)-integral zero.
This follows by dominated convergence along \(E_s\) and the barycenter
identity. One countable family of words is total in every phase chamber,
so a single null set suffices, and outside it every phase has a
one-dimensional ground space. The spectral calculus of V.1.2 (SNAG,
positivity of the energy from analyticity of the heat semigroup) is
built without assuming a gap and is used downstream (V.2, V.3).

\subsubsection{4. Steps}\label{steps-20}

\paragraph{4.1 OS space, chambers, spectral
calculus}\label{os-space-chambers-spectral-calculus}

\textbf{Lemma V.1.a (chamber).} Under the hierarchy, V.1.1 holds except
for the convergence \(E_t\to P_0\). \emph{Proof sketch.}

\begin{itemize}
\tightlist
\item
  \textbf{Contraction (OS doubling trick).} For \(p\) in the span of
  positive words,
  \(\|E_tp\|^2=\langle p,E_{2t}p\rangle\le\|p\|\sup_s\|E_sp\|\). The
  supremum \(c\) is finite, because
  \(s\mapsto S(\Theta b\cdot\tau_sb')\) is bounded for \(s\ge0\): the
  points stay at a fixed positive distance, and the zeroth cap applies.
  Hence \(c^2\le\|p\|c\), so \(\|E_tp\|\le\|p\|\).
\item
  \textbf{Self-adjointness and unitarity.} Self-adjointness of \(E_t\)
  and unitarity of \(U_y\) follow from translation invariance.
\item
  \textbf{Strong continuity.} It follows from continuity of the Gram
  functions (smoothness) together with the uniform bounds.
\end{itemize}

\emph{Lean:}
\href{https://github.com/yang-mills-group/OpenAI-YM-20261028/blob/3456e5e75cf8d4cd558955a5c7638fba2d68f464/YM/OS/Flow.lean\#L70-L89}{\texttt{OS.contract}},
\href{https://github.com/yang-mills-group/OpenAI-YM-20261028/blob/3456e5e75cf8d4cd558955a5c7638fba2d68f464/YM/OS/Dynamics.lean\#L63-L104}{\texttt{OS.implement}}.

\textbf{Lemma V.1.b (ground projection).} For a semigroup of
self-adjoint contractions, \(E_t\psi\to P_0\psi\) for every \(\psi\).
\emph{Proof sketch.}

\begin{itemize}
\tightlist
\item
  For \(u\le v\),
  \(\langle E_u\psi,E_v\psi\rangle=\|E_{(u+v)/2}\psi\|^2\ge\|E_v\psi\|^2\).
  Hence \(\|E_u\psi-E_v\psi\|^2\le\|E_u\psi\|^2-\|E_v\psi\|^2\), and the
  decreasing limit of \(\|E_t\psi\|^2\) makes \(E_t\psi\) Cauchy.
\item
  The limit is fixed by every \(E_s\), and \(\psi\) minus the limit is
  orthogonal to \(\mathcal H_0\), because
  \(\langle E_t\psi,u\rangle=\langle\psi,u\rangle\) for
  \(u\in\mathcal H_0\).
\item
  No spectral theorem is used.
\end{itemize}

\emph{Lean:}
\href{https://github.com/yang-mills-group/OpenAI-YM-20261028/blob/3456e5e75cf8d4cd558955a5c7638fba2d68f464/YM/OS/Core.lean\#L38-L78}{\texttt{Semigroup.emerges}}.

\textbf{Lemma V.1.c (real-time group and SNAG).} V.1.2(a), (b) hold.
\emph{Proof sketch.}

\begin{itemize}
\tightlist
\item
  \textbf{Real-time group.} \(E_1=E_{1/2}^2\) is positive, with spectrum
  in \([0,1]\). The functional calculus gives \(E_1^z\), holomorphic for
  \(\mathrm{Re}\,z>0\), and \(E_1^t=E_t\) for \(t>0\) by positive roots
  and continuity. If \(E_t\psi=0\), analyticity makes \(E_s\psi=0\) for
  every \(s>0\), and continuity at zero gives \(\psi=0\).
  Self-adjointness then makes the range of \(E_1\) dense. Since
  \(\|E_1^{1-is}\psi\|=\|E_1\psi\|\), the map
  \(E_1\psi\mapsto E_1^{1-is}\psi\) extends to an isometry. The power
  laws give the group law and inverse at \(-s\); continuity on the dense
  range gives strong continuity.
\item
  \textbf{SNAG.} Integrate \(V\) against Fourier transforms of Schwartz
  functions to obtain a \(*\)-representation of continuous functions on
  the one-point compactification of \(\mathbb R^4\). Riesz--Markov gives
  a projection-valued measure there. The point at infinity carries no
  mass: evaluate at \(x=0\) along an approximate identity.
\end{itemize}

\emph{Lean:}
\href{https://github.com/yang-mills-group/OpenAI-YM-20261028/blob/3456e5e75cf8d4cd558955a5c7638fba2d68f464/YM/Heat/Extend.lean\#L42-L65}{\texttt{Semigroup.roots\_heat}},
\href{https://github.com/yang-mills-group/OpenAI-YM-20261028/blob/3456e5e75cf8d4cd558955a5c7638fba2d68f464/YM/Heat/Imaginary.lean\#L33-L46}{\texttt{Semigroup.pillar}},
\href{https://github.com/yang-mills-group/OpenAI-YM-20261028/blob/3456e5e75cf8d4cd558955a5c7638fba2d68f464/YM/Heat/Imaginary.lean\#L49}{\texttt{Semigroup.edge}},
\href{https://github.com/yang-mills-group/OpenAI-YM-20261028/blob/3456e5e75cf8d4cd558955a5c7638fba2d68f464/YM/Heat/Lift.lean\#L50-L58}{\texttt{Semigroup.unit}},
\href{https://github.com/yang-mills-group/OpenAI-YM-20261028/blob/3456e5e75cf8d4cd558955a5c7638fba2d68f464/YM/SNAG/Escape.lean\#L52-L119}{\texttt{Spectral.Orbit.total}},
\href{https://github.com/yang-mills-group/OpenAI-YM-20261028/blob/3456e5e75cf8d4cd558955a5c7638fba2d68f464/YM/Spectral/Projection.lean\#L19-L29}{\texttt{Spectral.PVM}}.

\textbf{Lemma V.1.d (spectrum condition and Laplace formula).} V.1.2(c),
(d) hold. \emph{Proof sketch.}

\begin{itemize}
\tightlist
\item
  \textbf{The function \(J\).} Fix \(\psi\) and \(x\). The function
  \(J(z)=\langle\psi,V(x)e^{i(\mathrm{Re}\,z)H}E_{\mathrm{Im}\,z}\psi\rangle\)
  is continuous on \(\mathrm{Im}\,z\ge0\). In the interior it equals
  \(\langle\psi,V(x)E_1^{-iz}\psi\rangle\), so it is holomorphic there.
  It is bounded by \(\|\psi\|^2\). Its boundary values are the Fourier
  transform of \(e^{ip\cdot x}d\mu_\psi(p)\) in the variable \(p^a\).
\item
  \textbf{Support theorem.} Let \(\mu\) be a finite positive Borel
  measure on \(\mathbb R^4\), let \(\ell\colon\mathbb R^4\to\mathbb R\)
  be continuous, and let \(F\) be bounded and holomorphic in the upper
  half-plane, continuous up to the boundary, with
  \(F(b)=\int e^{i\ell(y)b}d\mu(y)\). Then
  \(\mathrm{supp}\,\mu\subset\{\ell\ge0\}\). To see this, integrate
  \(F\) against the Gaussian \(e^{-ibz-\varepsilon z^2}\) on horizontal
  lines and shift the contour upward. For \(b<0\), choosing height
  \(-b/(2\varepsilon)\) bounds the mass where \(|\ell(y)-b|<|b|/2\) by
  \(C e^{-3b^2/(16\varepsilon)}\). Letting \(\varepsilon\downarrow0\)
  excludes negative \(\ell\) from the support. Apply this with \(x=0\)
  in \(J\) and \(\ell(p)=p^a\).
\item
  \textbf{Laplace formula.} Both sides of (d) are bounded and
  holomorphic with the same boundary values, so they agree.
\end{itemize}

\emph{Lean:}
\href{https://github.com/yang-mills-group/OpenAI-YM-20261028/blob/3456e5e75cf8d4cd558955a5c7638fba2d68f464/YM/Heat/Lift.lean\#L75-L93}{\texttt{Semigroup.altitude}},
\href{https://github.com/yang-mills-group/OpenAI-YM-20261028/blob/3456e5e75cf8d4cd558955a5c7638fba2d68f464/YM/Heat/Lift.lean\#L102-L104}{\texttt{Semigroup.nonneg}},
\href{https://github.com/yang-mills-group/OpenAI-YM-20261028/blob/3456e5e75cf8d4cd558955a5c7638fba2d68f464/YM/Heat/Lift.lean\#L108-L114}{\texttt{Semigroup.heat}},
\href{https://github.com/yang-mills-group/OpenAI-YM-20261028/blob/3456e5e75cf8d4cd558955a5c7638fba2d68f464/YM/Paley/Half.lean\#L63-L113}{\texttt{Paley.supports}},
\href{https://github.com/yang-mills-group/OpenAI-YM-20261028/blob/3456e5e75cf8d4cd558955a5c7638fba2d68f464/YM/Paley/Mapped.lean\#L28-L57}{\texttt{Paley.heat}},
\href{https://github.com/yang-mills-group/OpenAI-YM-20261028/blob/3456e5e75cf8d4cd558955a5c7638fba2d68f464/YM/Paley/Line.lean\#L31-L34}{\texttt{Paley.Air}}.

\paragraph{4.2 Slab operators}\label{slab-operators}

\textbf{Lemma V.1.e (power argument).}

\begin{itemize}
\tightlist
\item
  \textbf{Setting.} Let \(\Delta>0\), let \(A,B\) be linear maps on a
  complex vector space, and let \(Y\) be a linear map into a Hilbert
  space, with \(\langle YAv,Yw\rangle=\langle Yv,YBw\rangle\) for all
  \(v,w\).
\item
  \textbf{Statement.} If every \(v\) has a constant \(K_v\) with
  \(\|Y(BA)^nv\|\le K_v\Delta^{2n}\) for all \(n\), then
  \(\|YAv\|\le\Delta\|Yv\|\).
\end{itemize}

\emph{Proof sketch.}

\begin{itemize}
\tightlist
\item
  Put \(\delta_n=\|Y(BA)^nv\|/\Delta^{2n}\). Then
  \(\delta_n^2\le\|Yv\|\,\delta_{2n}\), so \(\sup_n\delta_n\le\|Yv\|\).
\item
  Hence \(\|YAv\|^2=\langle Yv,YBAv\rangle\le\Delta^2\|Yv\|^2\).
\item
  Apply this with \(A\) = insertion of \(w\), \(B\) = insertion of
  \(\tau_s\Theta w\), and \(Y\) = the OS map. The bench gives
  \(K_v=\sqrt{B_v}\) for single words; for finite sums use the triangle
  inequality. Extend by continuity.
\end{itemize}

\emph{Lean:}
\href{https://github.com/yang-mills-group/OpenAI-YM-20261028/blob/3456e5e75cf8d4cd558955a5c7638fba2d68f464/YM/OS/Train.lean\#L20-L60}{\texttt{Drive.against}},
\href{https://github.com/yang-mills-group/OpenAI-YM-20261028/blob/3456e5e75cf8d4cd558955a5c7638fba2d68f464/YM/OS/Train.lean\#L95-L114}{\texttt{OS.size}},
\href{https://github.com/yang-mills-group/OpenAI-YM-20261028/blob/3456e5e75cf8d4cd558955a5c7638fba2d68f464/YM/OS/Train.lean\#L116-L123}{\texttt{OS.producer}}.

\textbf{Lemma V.1.f (existence and all-copy bound; B01).} The first part
of V.1.3 holds. \emph{Proof sketch.}

\begin{itemize}
\tightlist
\item
  Let \(w\) be \(r\)-buffered in \((0,d)\). The bench gives \(s\le d\)
  with the train bound at width \(s\). V.1.e then gives a bounded
  operator \(T'\) with \(\|T'\|\le A\) and
  \(T'|b\rangle=|w\cdot\tau_sb\rangle\).
\item
  Put \(T_d(w)=T'E_{d-s}\), so \(\|T_d(w)\|\le A\); for the empty word,
  \(T_d(\emptyset)=E_d\).
\item
  Every word in a slab is buffered for some \((D,\rho,r)\). Operators
  agreeing on all positive-word vectors coincide, so the family is
  unique on valid slab words.
\end{itemize}

\emph{Lean:}
\href{https://github.com/yang-mills-group/OpenAI-YM-20261028/blob/3456e5e75cf8d4cd558955a5c7638fba2d68f464/YM/Work/B01.lean\#L50-L71}{\texttt{Work.B01.answer}}.

\textbf{Lemma V.1.g (smoothness laws; B02).}

\begin{itemize}
\tightlist
\item
  \textbf{Hypotheses.} \(S'\) is a table respecting \(S\), both
  satisfying the hierarchy, and \(S'\) satisfies the derivative
  identities of the plan of \(S\). The slab family on a chamber of
  \(S'\) has the all-copy bound.
\item
  \textbf{Conclusion.} The family is \(C^\infty\) in operator norm, with
  derivatives given by differentiated labels.
\end{itemize}

\emph{Proof sketch.}

\begin{itemize}
\tightlist
\item
  \textbf{Local bound on the candidate derivative.} The candidate
  derivative is the operator with the label differentiated. It is
  locally bounded: an \(r\)-buffered configuration is \(r/2\)-buffered
  on a neighbourhood, so the all-copy bound applies there.
\item
  \textbf{Matrix elements.} Between positive-word vectors, the
  derivative identity of \(S'\) gives the candidate's matrix elements as
  derivatives.
\item
  \textbf{From matrix elements to norm.} A mean-value estimate with
  constants uniform in the observing vectors gives first Lipschitz
  continuity in norm, then norm continuity of the candidate (apply the
  same to the next derivative), then Fréchet differentiability.
  Induction on the order gives \(C^\infty\).
\end{itemize}

\emph{Lean:}
\href{https://github.com/yang-mills-group/OpenAI-YM-20261028/blob/3456e5e75cf8d4cd558955a5c7638fba2d68f464/YM/Work/B02.lean\#L17-L30}{\texttt{Work.B02.answer}},
\href{https://github.com/yang-mills-group/OpenAI-YM-20261028/blob/3456e5e75cf8d4cd558955a5c7638fba2d68f464/YM/Work/B02/Analysis.lean\#L134-L153}{\texttt{B02.norm\_remainder}},
\href{https://github.com/yang-mills-group/OpenAI-YM-20261028/blob/3456e5e75cf8d4cd558955a5c7638fba2d68f464/YM/Work/B02/Slab.lean}{\texttt{B02/Slab}}.

\paragraph{4.3 The compressed algebra}\label{the-compressed-algebra}

\textbf{Lemma V.1.h (ground states are spatially invariant).} Under the
hierarchy and \(O(4)\)-covariance, \(U_y\psi=\psi\) for all
\(\psi\in\mathcal H_0\). \emph{Proof sketch.} Let \(p,q\) be positive
words, \(y\) spatial, and \(u=y+e_t\).

\begin{itemize}
\tightlist
\item
  \textbf{Pure time shifts are invisible.} For any \(d(s)\) with
  \(t(d(s))\to\infty\) and any fixed \(\lambda\),
  \(S(\Theta p\cdot(q+d(s)+\lambda e_t))-S(\Theta p\cdot(q+d(s)))\to0\).
  Both terms are matrix elements
  \(\langle p',U_{z(s)}E_{\tau(s)+\lambda}q'\rangle\) and
  \(\langle p',U_{z(s)}E_{\tau(s)}q'\rangle\), with \(\tau(s)\to\infty\)
  and arbitrary spatial \(z(s)\), and \(E_\tau q'\) converges (V.1.b).
\item
  \textbf{Rotate the spatial displacement into time.} Choose a
  hyperplane reflection \(K\) with \(Ku=|u|e_t\) and \(t(Ke_t)>0\).
  Covariance turns
  \(S(\Theta p\cdot(q+u+se_t))-S(\Theta p\cdot(q+se_t))\) into a pure
  time shift by \(|u|\) at displacement \(sKe_t\), so it tends to \(0\).
\item
  \textbf{Conclusion.} Combining the two,
  \(\langle p,U_yE_sq\rangle-\langle p,E_sq\rangle\to0\). Hence
  \(U_yP_0|q\rangle=P_0|q\rangle\), and this extends to all of
  \(\mathcal H_0\) by density.
\end{itemize}

\emph{Lean:}
\href{https://github.com/yang-mills-group/OpenAI-YM-20261028/blob/3456e5e75cf8d4cd558955a5c7638fba2d68f464/YM/Ground/Lean.lean\#L185-L188}{\texttt{Ground.inert}},
\href{https://github.com/yang-mills-group/OpenAI-YM-20261028/blob/3456e5e75cf8d4cd558955a5c7638fba2d68f464/YM/Ground/Lean.lean\#L153-L183}{\texttt{Ground.rays}},
\href{https://github.com/yang-mills-group/OpenAI-YM-20261028/blob/3456e5e75cf8d4cd558955a5c7638fba2d68f464/YM/Ground/Lean.lean\#L135-L150}{\texttt{Ground.across\_axis}},
\href{https://github.com/yang-mills-group/OpenAI-YM-20261028/blob/3456e5e75cf8d4cd558955a5c7638fba2d68f464/YM/Ground/Lean.lean\#L109-L133}{\texttt{Ground.overtops}}.

\textbf{Lemma V.1.i (the compression and its identities; B03).} V.1.4(b)
and (e), except separability, hold. \emph{Proof sketch.}

\begin{itemize}
\tightlist
\item
  \textbf{Independence of the shift.} For \(u\ge0\) the defining
  identity gives \(T_{d+u}(w)=T_d(w)E_u\) and
  \(T_{u+d}(\tau_uw)=E_uT_d(w)\); spatial translations give
  \(U_yT_d(w)=T_d(w+y)U_y\). Together with \(P_0E_u=P_0\) and V.1.h,
  these identities make \(M(w)\) independent of the shift and width.
\item
  \textbf{Values.} \(\langle\Omega,T_d(w)\Omega\rangle=S(w)\) because
  \(\Omega=|\emptyset\rangle\); translation invariance removes the
  shift.
\item
  \textbf{Identities.} Pairing against positive-word vectors gives
  \(\langle p,T_d(w)q\rangle=S(\Theta p\cdot w\cdot\tau_dq)\) with
  separated exterior words. Permutation symmetry, unit labels,
  multilinearity and dependence on \(F_d(S)\) therefore hold before
  compression.
\item
  \textbf{Adjoint and smoothness.} \(T_d(w)^*=T_d(\tau_d\Theta w)\) by
  translation invariance. Smoothness: near a configuration one shift and
  width serve (the slab condition is open), compression is bounded
  linear, and V.1.g applies.
\end{itemize}

\emph{Lean:}
\href{https://github.com/yang-mills-group/OpenAI-YM-20261028/blob/3456e5e75cf8d4cd558955a5c7638fba2d68f464/YM/Ground/Compress.lean\#L71-L87}{\texttt{Ground.at\_any}},
\href{https://github.com/yang-mills-group/OpenAI-YM-20261028/blob/3456e5e75cf8d4cd558955a5c7638fba2d68f464/YM/Ground/Compress.lean\#L111-L119}{\texttt{Ground.overline}},
\href{https://github.com/yang-mills-group/OpenAI-YM-20261028/blob/3456e5e75cf8d4cd558955a5c7638fba2d68f464/YM/Ground/Compress.lean\#L38-L43}{\texttt{Ground.pins}},
\href{https://github.com/yang-mills-group/OpenAI-YM-20261028/blob/3456e5e75cf8d4cd558955a5c7638fba2d68f464/YM/Ground/Compress.lean\#L129-L143}{\texttt{Ground.slim}},
\href{https://github.com/yang-mills-group/OpenAI-YM-20261028/blob/3456e5e75cf8d4cd558955a5c7638fba2d68f464/YM/Work/B03.lean\#L405-L415}{\texttt{Work.B03.answer}},
\href{https://github.com/yang-mills-group/OpenAI-YM-20261028/blob/3456e5e75cf8d4cd558955a5c7638fba2d68f464/YM/Bid/Phase.lean\#L41-L43}{\texttt{Bid.depot}}.

\textbf{Lemma V.1.j (cluster limits).} V.1.4(c) holds. \emph{Proof
sketch.}

\begin{itemize}
\tightlist
\item
  \textbf{Reduction to positive clusters.} Shift each cluster by a fixed
  vector to make it positive; this does not change \(M\).
\item
  \textbf{Induction on \(m\).} Shift the union so that \(q_1\) sits in
  its slab. Its vector is then
  \(T_d(q_1)U_{y_n}E_{u_n}|\text{rest}_n\rangle\) with \(u_n\to\infty\).
\item
  \textbf{Key convergence.} If \(\psi_n\to\psi\) and \(u_n\to\infty\),
  then \(U_{y_n}E_{u_n}\psi_n\to P_0\psi\) for arbitrary spatial
  \(y_n\), by V.1.b and V.1.h. Inductively the vector converges to
  \(T_d(q_1)M(q_2)\cdots M(q_m)\Omega\), and pairing with \(\Omega\)
  inserts \(P_0\).
\end{itemize}

\emph{Lean:}
\href{https://github.com/yang-mills-group/OpenAI-YM-20261028/blob/3456e5e75cf8d4cd558955a5c7638fba2d68f464/YM/Ground/Moments.lean\#L138-L148}{\texttt{Ground.arcs}},
\href{https://github.com/yang-mills-group/OpenAI-YM-20261028/blob/3456e5e75cf8d4cd558955a5c7638fba2d68f464/YM/Ground/Moments.lean\#L41-L93}{\texttt{Ground.future\_train}},
\href{https://github.com/yang-mills-group/OpenAI-YM-20261028/blob/3456e5e75cf8d4cd558955a5c7638fba2d68f464/YM/Ground/Moments.lean\#L26-L39}{\texttt{Ground.receives}}.

\textbf{Lemma V.1.k (commutativity, cyclicity, separation).} V.1.4(d)
holds. \emph{Proof sketch.} Write
\(\langle q_1,\dots,q_m\rangle=\langle\Omega,M(q_1)\cdots M(q_m)\Omega\rangle\).

\begin{itemize}
\tightlist
\item
  \textbf{Simultaneous rotation.} If \(K\in O(4)\) maps some \(u\) with
  \(t(u)>0\) to a vector with positive time, then
  \(\langle q_1,\dots,q_m\rangle=\langle Kq_1,\dots,Kq_m\rangle\). Place
  cluster \(j\) at \(n\,j\,u\), apply V.1.j to both orders, and use
  covariance of \(S\) at finite \(n\).
\item
  \textbf{Exchange of two clusters.} Let \(k\) exchange the time axis
  with the first spatial axis. Place four clusters at \(0\),
  \(n(e_t+2e_1)\), \(n(2e_t+e_1)\) and \(n(3e_t+3e_1)\), and apply
  \(k\): the two middle clusters exchange their time order. With the
  simultaneous-rotation step for \(u=e_t+e_1\), this gives
  \(\langle x,y,w,z\rangle=\langle x,w,y,z\rangle\).
\item
  \textbf{Commutativity of generators.} The vectors
  \(P_0|f\rangle=M(f)\Omega\) for positive \(f\) have dense linear span
  in \(\mathcal H_0\), and \(M(f)^*=M(\Theta f)\). Thus
  \(\langle M(f)\Omega,M(u)M(v)M(p)\Omega\rangle=\langle\Theta f,u,v,p\rangle\)
  is symmetric in \(u,v\). Testing on both dense spans proves
  \(M(u)M(v)=M(v)M(u)\).
\item
  \textbf{Extension to \(\mathcal A\).} The generators are closed under
  adjoints, so \(\mathcal A\) is commutative.
\item
  \textbf{Cyclicity and separation.}
  \(\mathcal A\Omega\supset P_0\{|b\rangle\}\) is dense. If
  \(p\Omega=0\), then \(p\,a\Omega=a\,p\Omega=0\) for all
  \(a\in\mathcal A\), so \(p=0\).
\end{itemize}

\emph{Lean:}
\href{https://github.com/yang-mills-group/OpenAI-YM-20261028/blob/3456e5e75cf8d4cd558955a5c7638fba2d68f464/YM/Ground/Permute.lean\#L45-L61}{\texttt{Ground.same\_up}},
\href{https://github.com/yang-mills-group/OpenAI-YM-20261028/blob/3456e5e75cf8d4cd558955a5c7638fba2d68f464/YM/Ground/Permute.lean\#L72-L93}{\texttt{Ground.exchanges}},
\href{https://github.com/yang-mills-group/OpenAI-YM-20261028/blob/3456e5e75cf8d4cd558955a5c7638fba2d68f464/YM/Ground/Permute.lean\#L95-L136}{\texttt{Ground.gens}},
\href{https://github.com/yang-mills-group/OpenAI-YM-20261028/blob/3456e5e75cf8d4cd558955a5c7638fba2d68f464/YM/Ground/Permute.lean\#L149-L161}{\texttt{Ground.adjunction}},
\href{https://github.com/yang-mills-group/OpenAI-YM-20261028/blob/3456e5e75cf8d4cd558955a5c7638fba2d68f464/YM/Ground/Compress.lean\#L145-L162}{\texttt{Ground.full}},
\href{https://github.com/yang-mills-group/OpenAI-YM-20261028/blob/3456e5e75cf8d4cd558955a5c7638fba2d68f464/YM/Ground/Spectrum.lean\#L115-L124}{\texttt{Ground.attains}}.

\textbf{Lemma V.1.l (separability; Gelfand--Riesz; B12).} \(\mathcal A\)
is norm-separable, and the first part of V.1.5 holds. \emph{Proof
sketch.}

\begin{itemize}
\tightlist
\item
  \textbf{Separability.} By multilinearity each generator is a finite
  combination of \(M(x;e)\) with basis labels \(e\). For fixed arity,
  degrees and basis indices, \(x\mapsto M(x;e)\) is continuous on the
  separable set \(\mathrm{Conf}_n\). A countable union of separable
  images, finite products, adjoints and closure stay separable.
\item
  \textbf{Gelfand--Riesz.} Gelfand gives \(\mathcal A\cong C(X)\), and a
  countable dense sequence makes \(X\) metrizable. The vacuum state is
  represented by a probability measure \(\nu\).
\item
  \textbf{Full support.} If \(\nu(O)=0\) for some open \(O\ni\xi_0\), a
  continuous bump supported in \(O\) corresponds to an element \(p\ne0\)
  with \(\|p\Omega\|^2=\int|\hat p|^2d\nu=0\), contradicting separation.
\end{itemize}

\emph{Lean:}
\href{https://github.com/yang-mills-group/OpenAI-YM-20261028/blob/3456e5e75cf8d4cd558955a5c7638fba2d68f464/YM/Work/B12.lean\#L81-L92}{\texttt{Work.B12.answer}},
\href{https://github.com/yang-mills-group/OpenAI-YM-20261028/blob/3456e5e75cf8d4cd558955a5c7638fba2d68f464/YM/Ground/Spectrum.lean\#L126-L140}{\texttt{Ground.resolution}},
\href{https://github.com/yang-mills-group/OpenAI-YM-20261028/blob/3456e5e75cf8d4cd558955a5c7638fba2d68f464/YM/Ground/Spectrum.lean\#L34-L41}{\texttt{Ground.tiny}},
\href{https://github.com/yang-mills-group/OpenAI-YM-20261028/blob/3456e5e75cf8d4cd558955a5c7638fba2d68f464/YM/Ground/Spectrum.lean\#L43-L66}{\texttt{Ground.supports}}.

\paragraph{4.4 Phase reflection and rotation
invariance}\label{phase-reflection-and-rotation-invariance}

\textbf{Lemma V.1.m (grid passage; B04).}

\begin{itemize}
\tightlist
\item
  \textbf{Setting.} Assume the cutoff laws and generic point limits of
  the Seed. Let \(w\) have real labels at generic points of the
  level-\(k\) grid, and let \(w_j\) be its translate by \(t\,j\) grid
  steps in every coordinate.
\item
  \textbf{Statement.} For \(i<j\) and all large \(t\), the lattice pair
  expectation \(\mathbb E_{N,p}[\mathrm{Re}(f_if_j)]\) of §2 converges,
  along the ordered limit, to
  \[e_{ij}(t)=\mathrm{Re}\big[S(w_i\cdot w_j)-S(w_i\cdot(\vartheta w)_j)-S((\vartheta w)_i\cdot w_j)+S((\vartheta w)_i\cdot(\vartheta w)_j)\big].\]
\end{itemize}

\emph{Proof sketch.} Expand the product into four point moments of the
Wilson family. For large \(t\) the union of two translated grid words is
generic, and the point limits of the Seed apply term by term. The
reflection acts before the shift, so the shift stays outside it.
\emph{Lean:}
\href{https://github.com/yang-mills-group/OpenAI-YM-20261028/blob/3456e5e75cf8d4cd558955a5c7638fba2d68f464/YM/Work/B04.lean\#L168-L177}{\texttt{Work.B04.answer}},
\href{https://github.com/yang-mills-group/OpenAI-YM-20261028/blob/3456e5e75cf8d4cd558955a5c7638fba2d68f464/YM/Phase/Approx.lean\#L78-L79}{\texttt{Grid.Expects}},
\href{https://github.com/yang-mills-group/OpenAI-YM-20261028/blob/3456e5e75cf8d4cd558955a5c7638fba2d68f464/YM/Phase/Approx.lean\#L51-L54}{\texttt{Grid.event}}.

\textbf{Lemma V.1.n (phase reflection; B05).}

\begin{itemize}
\tightlist
\item
  \textbf{Hypotheses.} The Seed, the grid hypothesis,
  \(O(4)\)-covariance, V.1.h and V.1.i.
\item
  \textbf{Statement.} \(M(w)=M(\vartheta w)\) for every separated word
  with arbitrary labels.
\end{itemize}

\emph{Proof sketch.}

\begin{itemize}
\tightlist
\item
  \textbf{(1) Real labels at generic grid points.} Here \(\Theta w\) is
  a permutation of \(\vartheta w\), so \(M(w)^*=M(\vartheta w)\). Hence
  \(D:=M(w)-M(\vartheta w)\) satisfies \(D^*=-D\) and
  \(a:=\langle\Omega,D^2\Omega\rangle=-\|D\Omega\|^2\le0\).
\item
  \textbf{(2) The variance bound.} Fix \(q\) and \(\varkappa\). For all
  large \(t\), the approximants \(g_0,\dots,g_q\) satisfy, late in the
  cutoff,
  \[0\le\int\Big(\sum_{j=0}^qg_j\Big)^2\le(q+1)R_\varkappa^2+2\sum_{0\le i<j\le q}\big(\mathbb E_{N,p}[\mathrm{Re}(f_if_j)]+\delta_\varkappa\big).\]
  First take the ordered cutoff limit using V.1.m. Then let
  \(t\to\infty\); V.1.j gives \(e_{ij}(t)\to a\), so
  \[0\le(q+1)\big(R_\varkappa^2+q(a+\delta_\varkappa)\big).\] Let
  \(q\to\infty\), then \(\varkappa\to\infty\). This gives \(a\ge0\),
  hence \(D\Omega=0\) and \(D=0\) by V.1.k.
\item
  \textbf{(3) Extension to all configurations.} Approximate a generic
  configuration by grid configurations of finer levels; the grid
  hypothesis provides a cofinal set of levels. Both sides are continuous
  (V.1.i), so equality holds on the co-null set of generic
  configurations, and then on \(\mathrm{Conf}_n\).
\item
  \textbf{(4) Complex labels.} Extend by multilinearity, since real
  basis vectors span.
\end{itemize}

\emph{Lean:}
\href{https://github.com/yang-mills-group/OpenAI-YM-20261028/blob/3456e5e75cf8d4cd558955a5c7638fba2d68f464/YM/Work/B05.lean\#L90-L115}{\texttt{Work.B05.answer}},
\href{https://github.com/yang-mills-group/OpenAI-YM-20261028/blob/3456e5e75cf8d4cd558955a5c7638fba2d68f464/YM/Work/B05.lean\#L48-L88}{\texttt{B05.real\_depot}},
\href{https://github.com/yang-mills-group/OpenAI-YM-20261028/blob/3456e5e75cf8d4cd558955a5c7638fba2d68f464/YM/Phase/Square.lean\#L38-L50}{\texttt{Grid.horizon}},
\href{https://github.com/yang-mills-group/OpenAI-YM-20261028/blob/3456e5e75cf8d4cd558955a5c7638fba2d68f464/YM/Phase/Square.lean\#L56-L63}{\texttt{Grid.estimates}},
\texttt{Grid.Covering.square} in
\href{https://github.com/yang-mills-group/OpenAI-YM-20261028/blob/3456e5e75cf8d4cd558955a5c7638fba2d68f464/YM/Phase/Square.lean}{\texttt{Phase/Square.lean}},
\href{https://github.com/yang-mills-group/OpenAI-YM-20261028/blob/3456e5e75cf8d4cd558955a5c7638fba2d68f464/YM/Phase/Approx.lean\#L33}{\texttt{Grid.anti}}.

\textbf{Lemma V.1.o (invariance under \(O(4)\); B06).} Given phase
reflection and the rotation laws, \(M(w)=M(Kw)\) for all \(K\in O(4)\).
\emph{Proof sketch.}

\begin{itemize}
\tightlist
\item
  \textbf{Products.} The set of \(K\) with \(M(\cdot)=M(K\,\cdot)\) is
  closed under products. It contains the time reflection, because by the
  rotation law \(R(\vartheta)\) is the mirror action on \(F_d(S)\) and
  V.1.n applies.
\item
  \textbf{Two-point functions.}
  \(\langle p,x\rangle=\langle Kp,Kx\rangle\) for every \(K\). This is
  V.1.k's simultaneous-rotation step. If \(t(Ke_t)>0\) use \(u=e_t\); if
  \(t(Ke_t)=0\) use \(u=e_t+K^{-1}e_t\); if \(t(Ke_t)<0\) compose with
  the time reflection.
\item
  \textbf{Conjugates of the time reflection.} Rotating a pair back by
  \(J^{-1}\), reflecting the second factor and rotating forward gives
  \(\langle p,x\rangle=\langle p,J\vartheta J^{-1}x\rangle\) for all
  \(p\). Since \(M(\Theta p)=M(p)^*\), this gives
  \(M(x)\Omega=M(J\vartheta J^{-1}x)\Omega\), hence equality by
  separation.
\item
  \textbf{Generation.} Every hyperplane reflection is such a conjugate,
  and reflections generate \(O(4)\) (Cartan--Dieudonné). No separate
  treatment of time--space rotations is needed.
\end{itemize}

\emph{Lean:}
\href{https://github.com/yang-mills-group/OpenAI-YM-20261028/blob/3456e5e75cf8d4cd558955a5c7638fba2d68f464/YM/Work/B06.lean\#L238-L242}{\texttt{Work.B06.answer}},
auxiliary lemmas in
\href{https://github.com/yang-mills-group/OpenAI-YM-20261028/blob/3456e5e75cf8d4cd558955a5c7638fba2d68f464/YM/Work/B06.lean}{\texttt{Work/B06.lean}}.

\paragraph{4.5 The phase tables}\label{the-phase-tables}

\textbf{Lemma V.1.p (phase tables, symmetries, barycenter; B07).}

\begin{itemize}
\tightlist
\item
  \textbf{Definition.} \(S^\xi(x;v)=\chi_\xi(M(x;v))\) is a table.
\item
  \textbf{Properties.}

  \begin{itemize}
  \tightlist
  \item
    It is basic: permutation symmetry, \(S^\xi(\emptyset)=1\), unit
    labels.
  \item
    It is translation invariant, respects \(S\) (labels in \(N_d(S)\)
    give \(M=0\)), and is \(O(4)\)-covariant (V.1.o).
  \item
    It is real-spin. Signed permutations act through \(R\) by the
    rotation law. Reality:
    \(S^\xi(x;\bar v)=\chi_\xi(M(\bar w))=\chi_\xi(M(\vartheta\bar w))=\chi_\xi(M(\Theta w))=\overline{\chi_\xi(M(w))}\),
    using V.1.n and a permutation.
  \item
    It is smooth: \(\chi_\xi\) is extended by Hahn--Banach to a bounded
    functional on \(B(\mathcal H_0)\) and composed with the smooth map
    \(x\mapsto M(x;v)\).
  \end{itemize}
\item
  \textbf{Barycenter.}
  \(\int S^\xi(w)\,d\nu=\langle\Omega,M(w)\Omega\rangle=S(w)\), by V.1.l
  and V.1.i.
\end{itemize}

\emph{Lean:}
\href{https://github.com/yang-mills-group/OpenAI-YM-20261028/blob/3456e5e75cf8d4cd558955a5c7638fba2d68f464/YM/Work/B07.lean\#L119-L124}{\texttt{Work.B07.answer}},
\href{https://github.com/yang-mills-group/OpenAI-YM-20261028/blob/3456e5e75cf8d4cd558955a5c7638fba2d68f464/YM/Work/B07.lean\#L83-L106}{\texttt{B07.state\_spin}},
\href{https://github.com/yang-mills-group/OpenAI-YM-20261028/blob/3456e5e75cf8d4cd558955a5c7638fba2d68f464/YM/Work/B07.lean\#L108-L117}{\texttt{B07.state\_smooth}},
\href{https://github.com/yang-mills-group/OpenAI-YM-20261028/blob/3456e5e75cf8d4cd558955a5c7638fba2d68f464/YM/Phase/Mixture.lean\#L113-L117}{\texttt{Resolution.barycenter}}.

\textbf{Lemma V.1.q (reflection positivity in every phase; B08).} Every
\(S^\xi\) is reflection positive for every axis, in point and smeared
form. \emph{Proof sketch.}

\begin{itemize}
\tightlist
\item
  \textbf{Time axis.} Given finitely many positive words \(w_i\), choose
  \(d\) so they all lie in \((0,d)\). By V.1.i,
  \(M(\Theta w_i\cdot w_j)=Y_i^*Y_j\) with
  \(Y_i=T_d(w_i)|_{\mathcal H_0}\colon\mathcal H_0\to\mathcal H\).
  Therefore \(\sum\bar c_ic_jM(\Theta w_i\cdot w_j)=Y^*Y\ge0\), where
  \(Y=\sum c_iY_i\). Spectral permanence makes its character value
  nonnegative.
\item
  \textbf{Other axes.} Real-spin exchanges the axis with time.
\item
  \textbf{Smeared form.} Realize the positive point kernel in a Hilbert
  space. A compactly supported test gives the integral of its word
  vectors against the test function. Smoothness makes these
  vector-valued functions continuous and integrable; the smeared
  quadratic form is the squared norm of a finite sum of these integrals.
\end{itemize}

\emph{Lean:}
\href{https://github.com/yang-mills-group/OpenAI-YM-20261028/blob/3456e5e75cf8d4cd558955a5c7638fba2d68f464/YM/Work/B08.lean\#L206-L209}{\texttt{Work.B08.answer}},
\href{https://github.com/yang-mills-group/OpenAI-YM-20261028/blob/3456e5e75cf8d4cd558955a5c7638fba2d68f464/YM/Work/B08.lean\#L27-L33}{\texttt{B08.char\_nonneg}},
\href{https://github.com/yang-mills-group/OpenAI-YM-20261028/blob/3456e5e75cf8d4cd558955a5c7638fba2d68f464/YM/Work/B08.lean\#L193-L203}{\texttt{B08.smear}}.

\textbf{Lemma V.1.r (derivative, inclusion and stress laws in every
phase; B09).} Every \(S^\xi\) satisfies the inclusion and derivative
identities of the plan and the stress-tensor Ward identities.
\emph{Proof sketch.}

\begin{itemize}
\tightlist
\item
  \textbf{The evaluator.} There is a bounded functional \(\varphi_\xi\)
  on \(B(\mathcal H)\) with
  \(S^\xi(x;v)=\varphi_\xi\big(T_D(w(x)+s)\big)\) on a padded
  neighbourhood of each configuration; it is \(\chi_\xi\) extended by
  Hahn--Banach, composed with compression.
\item
  \textbf{Derivatives and inclusions.} The identities hold among slab
  operators. Derivatives come from V.1.g; pairing with positive-word
  vectors reduces inclusions to the identities of \(S\).
\item
  \textbf{Stress laws.} Choose one shift \(s\) and width \(D\) that put
  the configuration and the support of a smooth compactly supported
  window \(\chi\) in \((0,D)\), with \(\chi=1\) near the distinguished
  point \(x_0\) and all spectators outside its support. The operators
  \(-\sum_a\int\partial_a\chi(y)\,T_D((J_{ab}(y)B(x_0)\cdots)+s)\,dy\)
  are Bochner integrals. The shell weight vanishes near every possible
  collision, which gives integrability. Pairing with positive-word
  vectors adds only exterior points outside the window, so the Ward
  identities of \(S\) apply. The angular law uses the additional weight
  \((g(y-x_0))_b\) for a skew matrix \(g\). Apply \(\varphi_\xi\) to
  both operator identities.
\end{itemize}

\emph{Lean:}
\href{https://github.com/yang-mills-group/OpenAI-YM-20261028/blob/3456e5e75cf8d4cd558955a5c7638fba2d68f464/YM/Work/B09.lean\#L12-L15}{\texttt{Work.B09.answer}},
\href{https://github.com/yang-mills-group/OpenAI-YM-20261028/blob/3456e5e75cf8d4cd558955a5c7638fba2d68f464/YM/Work/B09/Core.lean\#L114-L129}{\texttt{B09.has\_evaluator}},
\href{https://github.com/yang-mills-group/OpenAI-YM-20261028/blob/3456e5e75cf8d4cd558955a5c7638fba2d68f464/YM/Work/B09/Shell.lean\#L147-L215}{\texttt{B09.Evaluates.flux}}.

\textbf{Lemma V.1.s (caps in every phase; B10, B11).}

\begin{itemize}
\tightlist
\item
  \begin{enumerate}
  \def\labelenumi{(\roman{enumi})}
  \tightlist
  \item
    Uniform zeroth cap: V.1.6(c).
  \end{enumerate}
\item
  (ii) A smooth basic table that respects \(S\), satisfies the inclusion
  and derivative identities, and has the zeroth cap for every fixed list
  of labels, has caps on all derivatives.
\end{itemize}

\emph{Proof sketch.}

\begin{itemize}
\tightlist
\item
  \begin{enumerate}
  \def\labelenumi{(\roman{enumi})}
  \tightlist
  \item
    Translate \(w\) into a slab with buffer \(r\) and unit padding. Then
    \(|\chi_\xi(M(w))|\le\|M(w)\|\le\|T_d(w)\|\), and the all-copy bound
    applies; the constants depend only on \((D,\rho)\).
  \end{enumerate}
\item
  (ii) A coordinate derivative replaces one label by its derivative
  label, raising its degree by one. Induct on the order, and bound each
  multilinear derivative norm by the sum over its finitely many
  coordinate directions. The constants may depend on the arity, labels
  and derivative order, but not on the positions or \(r\).
\end{itemize}

\emph{Lean:}
\href{https://github.com/yang-mills-group/OpenAI-YM-20261028/blob/3456e5e75cf8d4cd558955a5c7638fba2d68f464/YM/Work/B11.lean\#L75-L89}{\texttt{Work.B11.answer}},
\href{https://github.com/yang-mills-group/OpenAI-YM-20261028/blob/3456e5e75cf8d4cd558955a5c7638fba2d68f464/YM/Work/B11.lean\#L60-L73}{\texttt{B11.phase\_cap}},
\href{https://github.com/yang-mills-group/OpenAI-YM-20261028/blob/3456e5e75cf8d4cd558955a5c7638fba2d68f464/YM/Work/B10.lean\#L95-L105}{\texttt{Work.B10.answer}},
\href{https://github.com/yang-mills-group/OpenAI-YM-20261028/blob/3456e5e75cf8d4cd558955a5c7638fba2d68f464/YM/Work/B10.lean\#L72-L90}{\texttt{B10.capped\_tags}},
\href{https://github.com/yang-mills-group/OpenAI-YM-20261028/blob/3456e5e75cf8d4cd558955a5c7638fba2d68f464/YM/Bid/Phase.lean\#L97-L98}{\texttt{Bid.OneCap}}.

\paragraph{4.6 Null spaces of the
phases}\label{null-spaces-of-the-phases}

\textbf{Lemma V.1.t (moment bounds give bounds in every phase).} Let
\(\tau\ge0\) and \(p\in\mathcal A\). If
\(\langle\Omega,(p^*p)^n\Omega\rangle\le\tau^{2n}\) for every
\(n\in\mathbb N\), then \(|\chi_\xi(p)|\le\tau\) for every \(\xi\).
\emph{Proof sketch.}

\begin{itemize}
\tightlist
\item
  The hypothesis says \(\int|\chi_\xi(p)|^{2n}d\nu\le\tau^{2n}\) for all
  \(n\).
\item
  If \(|\chi_{\xi_0}(p)|>\tau'>\tau\), then continuity gives an open
  \(O\ni\xi_0\) on which \(|\chi_\xi(p)|>\tau'\).
\item
  Full support gives \(\nu(O)>0\), so \(\nu(O)\tau'^{2n}\le\tau^{2n}\)
  for all \(n\), a contradiction.
\end{itemize}

\emph{Lean:}
\href{https://github.com/yang-mills-group/OpenAI-YM-20261028/blob/3456e5e75cf8d4cd558955a5c7638fba2d68f464/YM/Phase/Detect.lean\#L22-L58}{\texttt{Retain.bit}},
\href{https://github.com/yang-mills-group/OpenAI-YM-20261028/blob/3456e5e75cf8d4cd558955a5c7638fba2d68f464/YM/Phase/Detect.lean\#L21}{\texttt{Retain.Tall}}.

\textbf{Lemma V.1.u (detectors; B15--B18).}

\begin{itemize}
\tightlist
\item
  \textbf{The detectors.} Assume the Seed and phase reflection. Fix a
  degree \(d\), an index \(k\) of its corrected-field data, and a test
  row with a basis index. Let \(F\) be a one-point word at a prescribed
  grid position, with free value \(s\) at the rescaled position. For a
  two-point row, let \(G\) be its second one-point word, \(s'\) its free
  value and \(\gamma\) the free cumulant. There are detectors in
  \(\mathcal A\):

  \begin{itemize}
  \tightlist
  \item
    the one-point detector \(\Delta_1=M(F)-s\cdot1\);
  \item
    the two-point detector
    \(\Delta_2=M(F\cdot G)-s\,M(G)-s'M(F)+(ss'-\gamma)\cdot1\).
  \end{itemize}
\item
  \textbf{Statement.} For every slot there is a one-point detector; for
  a two-point row there is also the two-point detector. Each satisfies
  \(\langle\Omega,(\Delta^*\Delta)^n\Omega\rangle\le\mathrm{gap}_k^{2n}\)
  for every \(n\in\mathbb N\), where \(\mathrm{gap}_k\ge0\) and
  \(\mathrm{gap}_k\to0\).
\end{itemize}

\emph{Proof sketch.} Each detector is the cluster element
\(c_E=\sum_i\lambda_iM(\text{word}_i)\) of a finite lattice
\emph{episode} \(E\), whose chessboard budget is \(\mathrm{gap}_k\)
(B18, from the corrected-field data). Then:

\begin{itemize}
\tightlist
\item
  \textbf{B15 (cutoff level).} The absolute lattice expectation of the
  product of \(2n\) diagonal translates of the episode function,
  alternately conjugated, is at most
  \((\mathrm{gap}_k+\varepsilon)^{2n}\) late in the cutoff. This is the
  chessboard estimate.
\item
  \textbf{B16 (limit table).} For fixed \(n\) and all sufficiently large
  translations, expand the product into finitely many monomials. Their
  points are generic: within each cluster this is part of the data, and
  between clusters each coordinate tie excludes at most one translation
  parameter. Point limits apply to each monomial, including the
  conjugated cutoff labels.
\item
  \textbf{B17 (vacuum moments).} Let the translation tend to infinity
  and apply V.1.j to each ordered monomial, using \(M(\bar w)=M(w)^*\)
  from V.1.i and V.1.n. The limit is
  \(\langle\Omega,(c_E^*c_E)^n\Omega\rangle\). Finally let
  \(\varepsilon\downarrow0\).
\end{itemize}

\emph{Lean:}
\href{https://github.com/yang-mills-group/OpenAI-YM-20261028/blob/3456e5e75cf8d4cd558955a5c7638fba2d68f464/YM/Work/B15.lean\#L108-L118}{\texttt{Work.B15.answer}},
\href{https://github.com/yang-mills-group/OpenAI-YM-20261028/blob/3456e5e75cf8d4cd558955a5c7638fba2d68f464/YM/Work/B16.lean\#L185-L190}{\texttt{B16.answer}},
\href{https://github.com/yang-mills-group/OpenAI-YM-20261028/blob/3456e5e75cf8d4cd558955a5c7638fba2d68f464/YM/Work/B17.lean\#L71-L90}{\texttt{B17.answer}},
\href{https://github.com/yang-mills-group/OpenAI-YM-20261028/blob/3456e5e75cf8d4cd558955a5c7638fba2d68f464/YM/Work/B18.lean\#L130-L133}{\texttt{B18.answer}},
\href{https://github.com/yang-mills-group/OpenAI-YM-20261028/blob/3456e5e75cf8d4cd558955a5c7638fba2d68f464/YM/Bid/Fetch.lean\#L96-L100}{\texttt{Bid.episode}};
definitions
\href{https://github.com/yang-mills-group/OpenAI-YM-20261028/blob/3456e5e75cf8d4cd558955a5c7638fba2d68f464/YM/Bid/Fetch.lean\#L18-L27}{\texttt{Bid.Episode}},
\href{https://github.com/yang-mills-group/OpenAI-YM-20261028/blob/3456e5e75cf8d4cd558955a5c7638fba2d68f464/YM/Bid/Fetch.lean\#L38-L40}{\texttt{Episode.Sold}},
\href{https://github.com/yang-mills-group/OpenAI-YM-20261028/blob/3456e5e75cf8d4cd558955a5c7638fba2d68f464/YM/Phase/Detect.lean\#L73}{\texttt{Retain.Δone}},
\href{https://github.com/yang-mills-group/OpenAI-YM-20261028/blob/3456e5e75cf8d4cd558955a5c7638fba2d68f464/YM/Phase/Detect.lean\#L75-L77}{\texttt{Retain.Δtwo}},
\href{https://github.com/yang-mills-group/OpenAI-YM-20261028/blob/3456e5e75cf8d4cd558955a5c7638fba2d68f464/YM/Phase/Detect.lean\#L81-L83}{\texttt{Retain.Funds}}.

\textbf{Lemma V.1.v (retention of null spaces).} Assume the
corrected-field data and detector bounds of V.1.u. If a table \(S'\) is
symmetric, respects \(S\), and satisfies \(S'(w)=\chi_\xi(M(w))\) for
every separated word and some fixed \(\xi\), then \(N_d(S')=N_d(S)\) for
every \(d\). \emph{Proof sketch.}

\begin{itemize}
\tightlist
\item
  \textbf{Comparison maps.} Let \(F_d^{\mathrm{free}}\) be the
  free-field space in the corrected-field data. Define
  \(U_k\colon F_d^{\mathrm{free}}\to V_d\) by sending its chosen basis
  to the corrected labels. Compose \(U_k\), the quotient map for \(S'\),
  and the finite family of \(S'\)-probes to obtain
  \(f_k\colon F_d^{\mathrm{free}}\to\mathbb C^m\).
\item
  \textbf{Injectivity.} V.1.t and V.1.u give one-point errors at most
  \(\mathrm{gap}_k\) and two-point errors at most
  \(\mathrm{gap}_k(1+|s|+|s'|)\). The free singleton values converge and
  are bounded, and the rescaled test points converge. Thus \(f_k\)
  converges in norm to the injective free test map. Injectivity is open
  in this finite-dimensional operator space, so \(f_k\) is injective for
  large \(k\).
\item
  \textbf{Conclusion.} For all large \(k\), \(q_dU_k\) is also
  surjective by the spanning hypothesis. If \(v\in N_d(S')\), choose
  \(z\in F_d^{\mathrm{free}}\) with \(q_dU_kz=q_dv\). Since \(S'\)
  respects \(S\), \(U_kz\in N_d(S')\), hence \(f_kz=0\). Injectivity
  gives \(z=0\), so \(q_dv=0\). The reverse inclusion is the respect
  hypothesis.
\end{itemize}

\emph{Lean:}
\href{https://github.com/yang-mills-group/OpenAI-YM-20261028/blob/3456e5e75cf8d4cd558955a5c7638fba2d68f464/YM/Phase/Picks.lean\#L128-L158}{\texttt{Retain.lift\_detect}},
\href{https://github.com/yang-mills-group/OpenAI-YM-20261028/blob/3456e5e75cf8d4cd558955a5c7638fba2d68f464/YM/Phase/Picks.lean\#L62-L108}{\texttt{Retain.result}},
\href{https://github.com/yang-mills-group/OpenAI-YM-20261028/blob/3456e5e75cf8d4cd558955a5c7638fba2d68f464/YM/Phase/Detect.lean\#L100-L117}{\texttt{Retain.pair\_near}}.

\paragraph{4.7 Choosing a phase}\label{choosing-a-phase}

\textbf{Lemma V.1.w (slab operators in a phase chamber; B13).} For every
\(\xi\), every time-axis chamber of \(S^\xi\) carries a slab family with
\(\|T^\xi_d(w)\|\le\|T_d(w)\|\). \emph{Proof sketch.}

\begin{itemize}
\tightlist
\item
  \textbf{The train in the phase.} In the phase chamber the train Gram
  is \(S^\xi(\Theta v_n\cdot v_n)=\chi_\xi(Y(v_n)^*Y(v_n))\), where
  \(Y(v_n)=(T_d(w)^*T_d(w))^nY(v)\).
\item
  \textbf{Bench in the phase.} Hence
  \(|S^\xi(\Theta v_n\cdot v_n)|\le\|\chi_\xi\|\,\|Y(v)\|^2\,\Delta^{4n}\)
  for every \(\Delta\ge\|T_d(w)\|\). This is the bench in the phase,
  with a constant depending on \(v\).
\item
  \textbf{Conclusion.} V.1.e gives norm \(\le\Delta\) for every
  \(\Delta>\|T_d(w)\|\); the constant depending on \(v\) drops out. No
  operator inequality between the chambers is used.
\end{itemize}

\emph{Lean:}
\href{https://github.com/yang-mills-group/OpenAI-YM-20261028/blob/3456e5e75cf8d4cd558955a5c7638fba2d68f464/YM/Work/B13.lean\#L91-L105}{\texttt{Work.B13.answer}},
\href{https://github.com/yang-mills-group/OpenAI-YM-20261028/blob/3456e5e75cf8d4cd558955a5c7638fba2d68f464/YM/Work/B13.lean\#L51-L72}{\texttt{B13.buys}}.

\textbf{Lemma V.1.x (one countable family; B14).} There is a single
sequence \((b_k)\) of positive words, depending only on the label
register and the time axis, whose vectors have dense linear span in
every OS Hilbert space realizing the positive-word Gram form of a
symmetric smooth table. \emph{Proof sketch.}

\begin{itemize}
\tightlist
\item
  \textbf{Parametrization.} For fixed arity and degrees, positive words
  are parametrized by a separable space: positive separated
  configurations times label coefficients. Choose countable dense
  subsets, together with the empty word.
\item
  \textbf{Continuity.} Inner products of the vectors are values of the
  table at combined configurations, which are continuous in the points
  (smoothness) and in the coefficients (finite-dimensional
  multilinearity). So the vector map is continuous, and the closed span
  of the countable family contains all vectors.
\end{itemize}

\emph{Lean:}
\href{https://github.com/yang-mills-group/OpenAI-YM-20261028/blob/3456e5e75cf8d4cd558955a5c7638fba2d68f464/YM/Work/B14.lean\#L119-L143}{\texttt{Work.B14.universal}},
\href{https://github.com/yang-mills-group/OpenAI-YM-20261028/blob/3456e5e75cf8d4cd558955a5c7638fba2d68f464/YM/Work/B14.lean\#L145-L148}{\texttt{B14.answer}}.

\textbf{Lemma V.1.y (almost every phase has a unique ground state).}

\begin{itemize}
\tightlist
\item
  \textbf{Statement.} In the setting of V.1.8, for each positive \(b\),
  \(P^\xi_0|b\rangle_\xi=S^\xi(b)\,\Omega_\xi\) for \(\nu\)-almost every
  \(\xi\). This proves the uniqueness assertion of V.1.8.
\end{itemize}

\emph{Proof sketch.}

\begin{itemize}
\tightlist
\item
  \textbf{The variance identity.} Let
  \(\omega_\xi=\|P_0^\xi|b\rangle_\xi\|^2=\lim_{n\to\infty}S^\xi(\Theta\tau_nb\cdot\tau_nb)\).
  For \(n\in\mathbb N\) these functions of \(\xi\) are continuous,
  nonnegative, decreasing in \(n\) and dominated by their value at
  \(n=0\), which is integrable on compact \(X\). Their integrals are the
  corresponding quantities for \(S\). Dominated convergence gives
  \(\int\omega_\xi\,d\nu=\|P_0|b\rangle\|^2=\|M(b)\Omega\|^2=\int|\chi_\xi(M(b))|^2d\nu\).
\item
  \textbf{Almost-everywhere vanishing.} Since
  \(\|P_0^\xi|b\rangle-S^\xi(b)\Omega_\xi\|^2=\omega_\xi-|S^\xi(b)|^2\ge0\)
  has integral zero, it vanishes a.e.
\item
  \textbf{Uniqueness.} Apply this to the family \((b_k)\) of V.1.x; the
  union of the countably many null sets is null. Outside it, \(P_0^\xi\)
  maps a dense set into \(\mathbb C\Omega_\xi\), so
  \(\mathcal H_0^\xi=\mathbb C\Omega_\xi\). Since \(\nu(X)=1\), such
  \(\xi\) exist; the choice is not canonical.
\end{itemize}

\emph{Lean:}
\href{https://github.com/yang-mills-group/OpenAI-YM-20261028/blob/3456e5e75cf8d4cd558955a5c7638fba2d68f464/YM/Phase/Unique.lean\#L56-L86}{\texttt{Cluster.exp\_real}},
\href{https://github.com/yang-mills-group/OpenAI-YM-20261028/blob/3456e5e75cf8d4cd558955a5c7638fba2d68f464/YM/Phase/Unique.lean\#L41-L47}{\texttt{Cluster.champ}},
\href{https://github.com/yang-mills-group/OpenAI-YM-20261028/blob/3456e5e75cf8d4cd558955a5c7638fba2d68f464/YM/Phase/Unique.lean\#L31-L36}{\texttt{Cluster.fits}},
\href{https://github.com/yang-mills-group/OpenAI-YM-20261028/blob/3456e5e75cf8d4cd558955a5c7638fba2d68f464/YM/Phase/Unique.lean\#L91-L110}{\texttt{Cluster.lifts}},
\href{https://github.com/yang-mills-group/OpenAI-YM-20261028/blob/3456e5e75cf8d4cd558955a5c7638fba2d68f464/YM/Phase/Unique.lean\#L111-L114}{\texttt{Cluster.selects}}.

\textbf{Lemma V.1.z (assembly).} V.1.9 holds. \emph{Proof sketch.}

\begin{itemize}
\tightlist
\item
  \textbf{Slab and phase constructions.} B01 and B02 give the original
  slab family with its laws. V.1.h gives spatial invariance; B03 gives
  the compression identities; B04 and B05 give phase reflection;
  B15--B18 give the detector bounds; and B12 gives separability. For any
  decomposition, B06 and B07 construct the tables, and B08--B11 give
  their hierarchy and field laws. B13 and B02 give slab families in the
  phase chambers, and B14 supplies a common countable total family.
\item
  \textbf{The certificate.} Choose a chamber for every phase, select
  \(\xi_*\) by V.1.y, and apply V.1.v to obtain \(N_d(S^\xi)=N_d(S)\)
  for every \(\xi,d\).
\end{itemize}

\emph{Lean:}
\href{https://github.com/yang-mills-group/OpenAI-YM-20261028/blob/3456e5e75cf8d4cd558955a5c7638fba2d68f464/YM/Bid/Market.lean\#L10-L36}{\texttt{Bid.Looms}},
\href{https://github.com/yang-mills-group/OpenAI-YM-20261028/blob/3456e5e75cf8d4cd558955a5c7638fba2d68f464/YM/Org/Phase.lean\#L60-L76}{\texttt{Phase.Loom.cleanse}},
\href{https://github.com/yang-mills-group/OpenAI-YM-20261028/blob/3456e5e75cf8d4cd558955a5c7638fba2d68f464/YM/Pert/Construction.lean\#L21-L31}{\texttt{Pert.withFinal}}.

\subsubsection{5. Where the difficulty
lies}\label{where-the-difficulty-lies-20}

\textbf{Distinctive arguments.}

\begin{itemize}
\tightlist
\item
  \textbf{The decomposition uses the OS space of the limit
  correlations.} Three properties require separate arguments:

  \begin{itemize}
  \tightlist
  \item
    spatial invariance of the ground states (V.1.h);
  \item
    commutativity of \(\mathcal A\) (V.1.k);
  \item
    reality of phase values (V.1.n).
  \end{itemize}
\item
  \textbf{Lattice input after the limit.} It enters at three places
  only: the bench (V.1.7), phase reflection (V.1.n, through the grid
  hypothesis), and the detectors (V.1.u).
\item
  \textbf{Phase reflection is the least standard argument.} It is a
  variance-positivity statement for far-separated copies of a real
  lattice variable. It is equivalent to the claim that the abstract
  characters assign real values to real observables.
\item
  \textbf{Detectors.} They require quantitative UV input: one- and
  two-point chessboard deviations \(\mathrm{gap}_k\to0\) of corrected
  fields from free values. These estimates preserve the field null
  spaces in every phase.
\end{itemize}

\textbf{Routine but long.}

\begin{itemize}
\tightlist
\item
  B02, B03, B09, B12, B14, B16, B18: smooth dependence on positions,
  label quotients, finite polynomial limits, separability.
\item
  \texttt{SNAG}, \texttt{Paley}, \texttt{Heat} (about 1,100 lines)
  reprove textbook theorems in an explicit form adapted to the library.
\end{itemize}

\textbf{Unusual or worth checking.} 1. \textbf{Selection.} The selected
phase comes from a measure-theoretic existence argument and is not
canonical. No finite bound on the number of phases is asserted. V.3
later proves unique ground states for every phase. 2. \textbf{Sub-width
in the bench.} The bench gives the train bound only at a
lattice-compatible width \(s\le d\). The remaining width is absorbed by
the contraction \(E_{d-s}\) (V.1.f). 3. \textbf{Which symmetries are
used where.} - Commutativity uses covariance only under the exchange of
the time axis with one spatial axis. This is a lattice symmetry, already
contained in the real-spin part of the hierarchy. - Spatial invariance
of \(\mathcal H_0\) needs arbitrary hyperplane reflections, i.e.~the
continuum \(O(4)\) action of VI.1. In the Lean, both steps draw on that
action. - Both use improper elements of \(O(4)\) (reflections), not only
rotations. 4. \textbf{Scope of the certificate.} Its unique-ground-state
assertion concerns the time-axis chamber of the selected phase. It also
records the hierarchy, field laws and null-space equality for every
phase. The continuation linked to the selected Wightman realization uses
that selected table. 5. \textbf{Smoothness hypothesis in B02.} B02 needs
the derivative identities of the table on which the slab family lives.
For phases these come from B09, which uses the slab laws of \(S\), so
the two applications of B02 are ordered.

\subsubsection{6. Notation}\label{notation-20}

{\def\LTcaptype{none} % do not increment counter
\begin{longtable}[]{@{}
  >{\raggedright\arraybackslash}p{(\linewidth - 4\tabcolsep) * \real{0.3333}}
  >{\raggedright\arraybackslash}p{(\linewidth - 4\tabcolsep) * \real{0.3333}}
  >{\raggedright\arraybackslash}p{(\linewidth - 4\tabcolsep) * \real{0.3333}}@{}}
\toprule\noalign{}
\begin{minipage}[b]{\linewidth}\raggedright
object
\end{minipage} & \begin{minipage}[b]{\linewidth}\raggedright
code name
\end{minipage} & \begin{minipage}[b]{\linewidth}\raggedright
proposed notation
\end{minipage} \\
\midrule\noalign{}
\endhead
\bottomrule\noalign{}
\endlastfoot
limit table, value on a word & \texttt{Table}, \texttt{S.expect\ w} &
\(S=S^{\mathrm{mix}}\), \(S(w)\) \\
OS reflection; linear time reflection of a word &
\texttt{reflect\ w\ timeAxis};
\href{https://github.com/yang-mills-group/OpenAI-YM-20261028/blob/3456e5e75cf8d4cd558955a5c7638fba2d68f464/YM/Phase/Approx.lean\#L33}{\texttt{Grid.anti}}
& \(\Theta w\); \(\vartheta w\) \\
OS space, vector of a positive word & \texttt{GNS.H}, \texttt{GNS.plant}
& \(\mathcal H\), \(\lvert b\rangle\) \\
chamber &
\href{https://github.com/yang-mills-group/OpenAI-YM-20261028/blob/3456e5e75cf8d4cd558955a5c7638fba2d68f464/YM/OS/Dynamics.lean\#L57-L61}{\texttt{Chamber}}
& \((\mathcal H,E_t,U_y)\) \\
ground space, projection, vacuum &
\href{https://github.com/yang-mills-group/OpenAI-YM-20261028/blob/3456e5e75cf8d4cd558955a5c7638fba2d68f464/YM/OS/Core.lean\#L14}{\texttt{Semigroup.ground}},
\texttt{orth},
\href{https://github.com/yang-mills-group/OpenAI-YM-20261028/blob/3456e5e75cf8d4cd558955a5c7638fba2d68f464/YM/Phase/Mixture.lean\#L32-L33}{\texttt{omega0}}
& \(\mathcal H_0\), \(P_0\), \(\Omega\) \\
real-time group; 4-parameter group; spectral measure &
\href{https://github.com/yang-mills-group/OpenAI-YM-20261028/blob/3456e5e75cf8d4cd558955a5c7638fba2d68f464/YM/Heat/Imaginary.lean\#L49}{\texttt{edge}};
\href{https://github.com/yang-mills-group/OpenAI-YM-20261028/blob/3456e5e75cf8d4cd558955a5c7638fba2d68f464/YM/Heat/Lift.lean\#L50-L58}{\texttt{unit}};
\href{https://github.com/yang-mills-group/OpenAI-YM-20261028/blob/3456e5e75cf8d4cd558955a5c7638fba2d68f464/YM/SNAG/Escape.lean\#L50}{\texttt{Orbit.P}}
& \(e^{isH}\); \(V(x)\); \(P(dp)\) \\
slab, buffer &
\href{https://github.com/yang-mills-group/OpenAI-YM-20261028/blob/3456e5e75cf8d4cd558955a5c7638fba2d68f464/YM/OS/Slabs.lean\#L19}{\texttt{Slab}},
\href{https://github.com/yang-mills-group/OpenAI-YM-20261028/blob/3456e5e75cf8d4cd558955a5c7638fba2d68f464/YM/OS/Slabs.lean\#L21-L22}{\texttt{buffer}}
& slab \((0,d)\), \(r\)-buffered \\
slab operators (``Spread'') and all-copy bound &
\href{https://github.com/yang-mills-group/OpenAI-YM-20261028/blob/3456e5e75cf8d4cd558955a5c7638fba2d68f464/YM/OS/Slabs.lean\#L101-L104}{\texttt{Spread}},
\texttt{place},
\href{https://github.com/yang-mills-group/OpenAI-YM-20261028/blob/3456e5e75cf8d4cd558955a5c7638fba2d68f464/YM/OS/Slabs.lean\#L111-L113}{\texttt{AllCopy}},
\href{https://github.com/yang-mills-group/OpenAI-YM-20261028/blob/3456e5e75cf8d4cd558955a5c7638fba2d68f464/YM/OS/Slabs.lean\#L153-L162}{\texttt{Laws}}
& \(T_d(w)\); \(C^n(n/r)^{\kappa n}\) \\
bench, train bound &
\href{https://github.com/yang-mills-group/OpenAI-YM-20261028/blob/3456e5e75cf8d4cd558955a5c7638fba2d68f464/YM/Bid/Phase.lean\#L21-L24}{\texttt{Bench}},
\href{https://github.com/yang-mills-group/OpenAI-YM-20261028/blob/3456e5e75cf8d4cd558955a5c7638fba2d68f464/YM/Bid/Phase.lean\#L27}{\texttt{Seats}},
\href{https://github.com/yang-mills-group/OpenAI-YM-20261028/blob/3456e5e75cf8d4cd558955a5c7638fba2d68f464/YM/OS/Train.lean\#L70-L71}{\texttt{QRead}}
& bench (V.1.7) \\
spatial invariance of \(\mathcal H_0\) &
\href{https://github.com/yang-mills-group/OpenAI-YM-20261028/blob/3456e5e75cf8d4cd558955a5c7638fba2d68f464/YM/Ground/Compress.lean\#L20}{\texttt{Steady}}
& V.1.4(a) \\
compression; its multilinear form; algebra &
\href{https://github.com/yang-mills-group/OpenAI-YM-20261028/blob/3456e5e75cf8d4cd558955a5c7638fba2d68f464/YM/Ground/Compress.lean\#L89-L91}{\texttt{M}},
\href{https://github.com/yang-mills-group/OpenAI-YM-20261028/blob/3456e5e75cf8d4cd558955a5c7638fba2d68f464/YM/Phase/Mixture.lean\#L46}{\texttt{atom}};
\href{https://github.com/yang-mills-group/OpenAI-YM-20261028/blob/3456e5e75cf8d4cd558955a5c7638fba2d68f464/YM/Bid/Phase.lean\#L41-L43}{\texttt{depot}};
\href{https://github.com/yang-mills-group/OpenAI-YM-20261028/blob/3456e5e75cf8d4cd558955a5c7638fba2d68f464/YM/Phase/Mixture.lean\#L54}{\texttt{Spread.algebra}}
& \(M(w)\); \(M(x;v)\); \(\mathcal A\) \\
identities of \(M\) (``Theme'') &
\href{https://github.com/yang-mills-group/OpenAI-YM-20261028/blob/3456e5e75cf8d4cd558955a5c7638fba2d68f464/YM/Bid/Phase.lean\#L45-L56}{\texttt{Theme}}
& V.1.4(e) \\
phase reflection &
\href{https://github.com/yang-mills-group/OpenAI-YM-20261028/blob/3456e5e75cf8d4cd558955a5c7638fba2d68f464/YM/Org/Phase.lean\#L19-L20}{\texttt{Phase.reflects}}
& \(M(w)=M(\vartheta w)\) \\
decomposition (``Resolution'') &
\href{https://github.com/yang-mills-group/OpenAI-YM-20261028/blob/3456e5e75cf8d4cd558955a5c7638fba2d68f464/YM/Phase/Mixture.lean\#L65-L72}{\texttt{Spread.Resolution}},
\texttt{X}, \texttt{μ}, \texttt{character} & \(X\), \(\nu\),
\(\chi_\xi\) \\
phase tables; their properties & \texttt{ensemble},
\href{https://github.com/yang-mills-group/OpenAI-YM-20261028/blob/3456e5e75cf8d4cd558955a5c7638fba2d68f464/YM/Bid/Phase.lean\#L72-L80}{\texttt{Garland}},
\href{https://github.com/yang-mills-group/OpenAI-YM-20261028/blob/3456e5e75cf8d4cd558955a5c7638fba2d68f464/YM/Org/Phase.lean\#L25-L32}{\texttt{Scripts}}
& \(S^\xi\) \\
\(S^\xi\) receives \(\chi_\xi\) &
\href{https://github.com/yang-mills-group/OpenAI-YM-20261028/blob/3456e5e75cf8d4cd558955a5c7638fba2d68f464/YM/Phase/Mixture.lean\#L89-L91}{\texttt{Resolution.receives}}
& \(S^\xi(w)=\chi_\xi(M(w))\) \\
grid hypothesis, bounded approximants &
\href{https://github.com/yang-mills-group/OpenAI-YM-20261028/blob/3456e5e75cf8d4cd558955a5c7638fba2d68f464/YM/Org/Input.lean\#L121-L123}{\texttt{Org.Grill}},
\href{https://github.com/yang-mills-group/OpenAI-YM-20261028/blob/3456e5e75cf8d4cd558955a5c7638fba2d68f464/YM/Phase/Approx.lean\#L63-L67}{\texttt{Grid.Covering}}
& grid hypothesis \\
lattice episode, cluster element, budget &
\href{https://github.com/yang-mills-group/OpenAI-YM-20261028/blob/3456e5e75cf8d4cd558955a5c7638fba2d68f464/YM/Bid/Fetch.lean\#L18-L27}{\texttt{Episode}},
\href{https://github.com/yang-mills-group/OpenAI-YM-20261028/blob/3456e5e75cf8d4cd558955a5c7638fba2d68f464/YM/Bid/Fetch.lean\#L33}{\texttt{cluster}},
\href{https://github.com/yang-mills-group/OpenAI-YM-20261028/blob/3456e5e75cf8d4cd558955a5c7638fba2d68f464/YM/Bid/Fetch.lean\#L38-L40}{\texttt{Sold}}
& \(E\), \(c_E\) \\
detectors; moment bound &
\href{https://github.com/yang-mills-group/OpenAI-YM-20261028/blob/3456e5e75cf8d4cd558955a5c7638fba2d68f464/YM/Phase/Detect.lean\#L73}{\texttt{Δone}},
\href{https://github.com/yang-mills-group/OpenAI-YM-20261028/blob/3456e5e75cf8d4cd558955a5c7638fba2d68f464/YM/Phase/Detect.lean\#L75-L77}{\texttt{Δtwo}};
\href{https://github.com/yang-mills-group/OpenAI-YM-20261028/blob/3456e5e75cf8d4cd558955a5c7638fba2d68f464/YM/Phase/Detect.lean\#L21}{\texttt{Tall}},
\href{https://github.com/yang-mills-group/OpenAI-YM-20261028/blob/3456e5e75cf8d4cd558955a5c7638fba2d68f464/YM/Phase/Detect.lean\#L81-L83}{\texttt{Funds}}
& \(\Delta_1,\Delta_2\) \\
countable dense family &
\href{https://github.com/yang-mills-group/OpenAI-YM-20261028/blob/3456e5e75cf8d4cd558955a5c7638fba2d68f464/YM/Phase/Unique.lean\#L88-L89}{\texttt{Cluster.horizon}}
& \((b_k)\) \\
unique ground state &
\href{https://github.com/yang-mills-group/OpenAI-YM-20261028/blob/3456e5e75cf8d4cd558955a5c7638fba2d68f464/YM/Phase/Mixture.lean\#L35}{\texttt{UniqueHeat}}
& \(\mathcal H_0^\xi=\mathbb C\Omega_\xi\) \\
selected phase and certificate &
\href{https://github.com/yang-mills-group/OpenAI-YM-20261028/blob/3456e5e75cf8d4cd558955a5c7638fba2d68f464/YM/Phase/Mixture.lean\#L127-L134}{\texttt{Selection}},
\texttt{atPoint},
\href{https://github.com/yang-mills-group/OpenAI-YM-20261028/blob/3456e5e75cf8d4cd558955a5c7638fba2d68f464/YM/Phase/Mixture.lean\#L136-L148}{\texttt{Certificate}},
\href{https://github.com/yang-mills-group/OpenAI-YM-20261028/blob/3456e5e75cf8d4cd558955a5c7638fba2d68f464/YM/Spec/Closed.lean\#L45-L47}{\texttt{Materializes}}
& \(\xi_*\); V.1.9 \\
bundles of obligations &
\href{https://github.com/yang-mills-group/OpenAI-YM-20261028/blob/3456e5e75cf8d4cd558955a5c7638fba2d68f464/YM/Org/Phase.lean\#L50-L54}{\texttt{Loom}},
\href{https://github.com/yang-mills-group/OpenAI-YM-20261028/blob/3456e5e75cf8d4cd558955a5c7638fba2d68f464/YM/Org/Phase.lean\#L42-L46}{\texttt{Stall}},
\href{https://github.com/yang-mills-group/OpenAI-YM-20261028/blob/3456e5e75cf8d4cd558955a5c7638fba2d68f464/YM/Org/Phase.lean\#L34-L40}{\texttt{Vendor}}
& assembly assumptions \\
\end{longtable}
}

The measure \(\nu\) corresponds to \(\mu\) in the Lean.

\subsubsection{7. What remains
unexplained}\label{what-remains-unexplained-20}

\begin{enumerate}
\def\labelenumi{\arabic{enumi}.}
\tightlist
\item
  \textbf{Ultraviolet inputs.} The chessboard geometry below
  \href{https://github.com/yang-mills-group/OpenAI-YM-20261028/blob/3456e5e75cf8d4cd558955a5c7638fba2d68f464/YM/Next/PHASE/Trains.lean\#L273-L304}{\texttt{Next.Phase.travel\_price}}
  belongs to VII.4. Its train estimate and the subsequent scale choice
  give the bench used here; the supporting packing argument remains
  outside this unit. The bounded-approximation input to
  \href{https://github.com/yang-mills-group/OpenAI-YM-20261028/blob/3456e5e75cf8d4cd558955a5c7638fba2d68f464/YM/MergeProof/South/FinalG.lean\#L76-L83}{\texttt{Accent.grilled}}
  belongs to III.11. Its output has the cutoff-independent bounds
  required in V.1.n, but the underlying ultraviolet estimates and the
  production of the corrected-field data remain inputs.
\item
  \textbf{Order of limits.} In V.1.n, fix the number of copies and
  approximation index, then fix a sufficiently large translation, take
  the ordered cutoff limit, and let the translation tend to infinity.
  Only then let the number of copies tend to infinity and the
  approximation error tend to zero. The detector argument similarly
  takes the cutoff limit before the large-translation limit and sends
  its budget slack to zero last.
\item
  \textbf{Later vacuum uniqueness.} V.1.8 gives uniqueness almost
  everywhere, which suffices to select \(\xi_*\). The stronger result
  for every phase,
  \href{https://github.com/yang-mills-group/OpenAI-YM-20261028/blob/3456e5e75cf8d4cd558955a5c7638fba2d68f464/YM/EveryPhase/Vacuum.lean\#L146-L153}{\texttt{EveryPhase.every}},
  uses reconstruction in V.3; its proof is outside this unit.
\item
  \textbf{Placement of the spectral calculus.} V.1.2 supplies V.2 and
  V.3 and can be placed at the start of reconstruction in the written
  proof. The \(O(4)\) action from VI.1 is needed only after the
  preliminary chamber and heat constructions.
\item
  \textbf{Bibliography.} The reference locations marked ``probable'' in
  §2 need bibliographic verification.
\end{enumerate}

\subsection{V.2. Osterwalder--Schrader reconstruction of a Wightman
theory}\label{v.2.-osterwalderschrader-reconstruction-of-a-wightman-theory}

\textbf{Part:} V, phases and reconstruction.

\textbf{Lean:}

\begin{itemize}
\tightlist
\item
  The work items A16 and C01--C29.

  \begin{itemize}
  \tightlist
  \item
    Statements:
    \href{https://github.com/yang-mills-group/OpenAI-YM-20261028/blob/3456e5e75cf8d4cd558955a5c7638fba2d68f464/YM/Bid/Corners.lean}{\texttt{Bid/Corners}},
    \href{https://github.com/yang-mills-group/OpenAI-YM-20261028/blob/3456e5e75cf8d4cd558955a5c7638fba2d68f464/YM/Bid/Tar.lean}{\texttt{Bid/Tar}},
    \href{https://github.com/yang-mills-group/OpenAI-YM-20261028/blob/3456e5e75cf8d4cd558955a5c7638fba2d68f464/YM/Bid/Grow.lean}{\texttt{Bid/Grow}}.
  \item
    Proofs:
    \href{https://github.com/yang-mills-group/OpenAI-YM-20261028/tree/3456e5e75cf8d4cd558955a5c7638fba2d68f464/YM/Work}{\texttt{YM/Work}},
    files C01--C29, subfolders C07/ and C13/.
  \item
    Collected by
    \href{https://github.com/yang-mills-group/OpenAI-YM-20261028/blob/3456e5e75cf8d4cd558955a5c7638fba2d68f464/YM/Integrated.lean\#L182-L212}{\texttt{Integrated.reconstructionArcs}}.
  \end{itemize}
\item
  The folders
  \href{https://github.com/yang-mills-group/OpenAI-YM-20261028/tree/3456e5e75cf8d4cd558955a5c7638fba2d68f464/YM/Boundary}{\texttt{YM/Boundary}},
  \href{https://github.com/yang-mills-group/OpenAI-YM-20261028/tree/3456e5e75cf8d4cd558955a5c7638fba2d68f464/YM/Fields}{\texttt{YM/Fields}},
  \href{https://github.com/yang-mills-group/OpenAI-YM-20261028/tree/3456e5e75cf8d4cd558955a5c7638fba2d68f464/YM/Hilbert}{\texttt{YM/Hilbert}},
  \href{https://github.com/yang-mills-group/OpenAI-YM-20261028/tree/3456e5e75cf8d4cd558955a5c7638fba2d68f464/YM/Compact}{\texttt{YM/Compact}},
  \href{https://github.com/yang-mills-group/OpenAI-YM-20261028/tree/3456e5e75cf8d4cd558955a5c7638fba2d68f464/YM/Distribution}{\texttt{YM/Distribution}}
  and
  \href{https://github.com/yang-mills-group/OpenAI-YM-20261028/tree/3456e5e75cf8d4cd558955a5c7638fba2d68f464/YM/Spark}{\texttt{YM/Spark}}.
  Spark is the Schwartz-space calculus used by C05, C08, C12, C14 and
  C25.
\item
  Material filed elsewhere but described here because it completes the
  reconstruction (§4.3, §4.6, §4.8):

  \begin{itemize}
  \tightlist
  \item
    The boundary-value and germ lemmas of
    \href{https://github.com/yang-mills-group/OpenAI-YM-20261028/tree/3456e5e75cf8d4cd558955a5c7638fba2d68f464/YM/Shift}{\texttt{YM/Shift}},
    filed partly under VI.1 by the declaration table.
  \item
    The assembly in
    \href{https://github.com/yang-mills-group/OpenAI-YM-20261028/blob/3456e5e75cf8d4cd558955a5c7638fba2d68f464/YM/Next/BOUNDARY/Stage.lean}{\texttt{Next/BOUNDARY/Stage.lean}}
    and
    \href{https://github.com/yang-mills-group/OpenAI-YM-20261028/blob/3456e5e75cf8d4cd558955a5c7638fba2d68f464/YM/Org/Theatre.lean}{\texttt{Org/Theatre.lean}},
    filed under VIII.1.
  \end{itemize}
\end{itemize}

\textbf{Size:} 6,456 lines in 705 declarations.

{\def\LTcaptype{none} % do not increment counter
\begin{longtable}[]{@{}ll@{}}
\toprule\noalign{}
part & lines \\
\midrule\noalign{}
\endhead
\bottomrule\noalign{}
\endlastfoot
proofs collected by reconstructionArcs & 4,279 \\
Fields & 438 \\
Compact & 218 \\
Distribution & 189 \\
Spark & 120 \\
Hilbert & 85 \\
Boundary & 66 \\
Overlay and Vac (ultraviolet material) & 558 \\
\end{longtable}
}

Three bundles in Integrated.lean are also ultraviolet material placed
here by the dependency cut; they are listed at the end of §1.

\subsubsection{1. What this unit proves}\label{what-this-unit-proves-21}

\emph{Conventions.}

\begin{itemize}
\tightlist
\item
  \textbf{Points.} A point of \(\mathbb R^4\) is \(x=(\vec x,x^0)\);
  time is the last coordinate, as in the code. A configuration is
  \(x=(x_1,\dots,x_n)\in\mathbb R^{4n}\).
\item
  \textbf{Pairing.} \(\langle\cdot,\cdot\rangle\) is the
  \emph{Euclidean} inner product of \(\mathbb R^{4n}\), extended
  complex-bilinearly. Dual cones are taken with respect to this pairing.
  Hilbert-space inner products are antilinear in the first argument.
\item
  \textbf{Cones and groups.} \(V_+=\{\lvert\vec x\rvert<x^0\}\), with
  closure \(\bar V_+\). \(\mathcal L^\uparrow_+\) is the proper
  orthochronous Lorentz group. \(\Lambda_a(t)\) is the boost of rapidity
  \(t\) along the spatial axis \(a\). The Poincaré group \(\mathcal P\)
  acts on tests by \((g_\ast f)(x)=f(g^{-1}x_1,\dots,g^{-1}x_n)\).
\item
  \textbf{Heights.} Heights \(y\in\mathbb R^n\) may be negative.
  \(E_n=\{y_1<\dots<y_n\}\) is the set of ordered heights, and
  \(\Sigma_n=\{z\in\mathbb C^n:\mathrm{Im}\,z\in E_n\}\) is the strict
  region. We write
  \(\underline y=((\vec0,y_1),\dots,(\vec0,y_n))\in\mathbb R^{4n}\).
\item
  \textbf{Fourier transform.}
  \(\hat f(p)=\int e^{-i\langle x,p\rangle}f(x)\,dx\),
  \(\check f(p)=(2\pi)^{-4n}\int e^{i\langle p,x\rangle}f(x)\,dx\), and
  \((\mathcal FW)(f)=W(\hat f)\).
\end{itemize}

Unless stated otherwise, arities \(n,m\) range over \(\mathbb N\),
including \(0\), and label tuples range over \(\mathcal F^n\). Supports
of functions are closed supports.

\emph{Tables and the OS space.}

\begin{itemize}
\tightlist
\item
  \(S\) and \(T\) are tables on one label register (glossary).
  \(N_d(Q)\) denotes the raw degree-\(d\) labels whose correlations in
  \(Q\) vanish with arbitrary other insertions. The quotient label space
  of \(S\) is \(F_d=\mathrm{Label}_d/N_d(S)\), and \(\mathcal F\) is
  their union, into which \(F_d\) embeds by \(\mathrm{pin}_d\).
\item
  \(T(x;v)\) is the Euclidean \(n\)-point function of \(T\) at a
  non-coincident configuration \(x\), read on tags \(v_k=(d_k,v_k)\) of
  \(S\). For a tuple of tags, \(\mathrm{pin}\,v\) means
  \((\mathrm{pin}_{d_k}v_k)_k\); the insertion into \(\mathcal F\) is
  implicit in \(\psi_v\). It is meaningful when \(T\) \emph{respects}
  \(S\), i.e.~\(N_d(S)\subseteq N_d(T)\).
\item
  \(\mathcal H\) is the OS Hilbert space of \(T\) for the time axis. It
  carries the vacuum \(\Omega\), a Hamiltonian \(H\ge0\) with
  holomorphic contraction semigroup \(e^{-\zeta H}\)
  (\(\mathrm{Re}\,\zeta>0\)), the spatial translations \(U(\vec a)\),
  and a unitary space-time translation group.
\item
  A \emph{spread} \(A\) gives slab operators \(A_d(w)\) for words \(w\)
  in a time slab of width \(d\) (V.1).
\item
  A differential plan \(P\) consists of linear maps
  \(\iota_{d,e}:F_d\to F_e\) for \(d\le e\) and derivative maps
  \(\partial_{d,a}:F_d\to F_{d+1}\). Its inclusions are \emph{compatible
  with a table} \(Q\) if replacing any tag \((d,v)\) by
  \((e,\iota_{d,e}v)\) leaves \(Q(x;v)\) unchanged at every
  non-coincident configuration, with arbitrary other tags. Its
  derivatives are compatible with \(Q\) if inserting \(\partial_{d,a}v\)
  differentiates the corresponding correlation in that slot and
  coordinate. Writing \(C_d\) for label conjugation and \(M_d(s)\) for a
  signed coordinate symmetry, \(P\) is \emph{oriented} if its inclusions
  commute with \(C_d,M_d(s)\), its derivatives commute with conjugation,
  and \[M_{d+1}(s)\partial_{d,b}=\sum_a s_{ab}\,\partial_{d,a}M_d(s).\]
  The induced star satisfies
  \(J\mathrm{pin}_dv=\mathrm{pin}_d(M_d(r_0)C_dv)\), where \(r_0\) is
  time reflection.
\end{itemize}

\emph{Continuations and boundary families.}

\begin{itemize}
\tightlist
\item
  A \textbf{continuation} \(F\) is a family of maps \(F_n(z;\vec x)\),
  \(z\in\mathbb C^n\), \(\vec x\in(\mathbb R^3)^n\), each multilinear on
  \(\mathcal F^n\). We write

  \begin{itemize}
  \tightlist
  \item
    \(F^y_n(x)=F_n(x^0+iy;\vec x)\), with real times \(x^0_k\) and
    heights \(y_k\);
  \item
    \(\psi_v(y;f)=\int F^y_n(x)(v)f(x)\,dx\).
  \end{itemize}
\item
  A \textbf{boundary family} \(W\) is a family of multilinear maps
  \(W_n:\mathcal F^n\to\mathcal S'(\mathbb R^{4n})\).

  \begin{itemize}
  \tightlist
  \item
    For a word \(b=((u_1,f_1),\dots,(u_n,f_n))\), with
    \(u_k\in\mathcal F\) and \(f_k\in\mathcal S(\mathbb R^4)\), we set
    \(W(b)=W_n(u)(f_1\otimes\dots\otimes f_n)\).
  \item
    For an antilinear map \(J\) on \(\mathcal F\) we set
    \(b^{\ast}=((Ju_n,\bar f_n),\dots,(Ju_1,\bar f_1))\) and
    \(\langle b,c\rangle_W=W(b^{\ast}c)\).
  \end{itemize}
\end{itemize}

The thirty reconstruction statements are conditional results for
arbitrary tables, OS data and continuations. VIII.1 applies them to the
mixed limit table and to a selected phase (§4.8).

\textbf{Proposition V.2.1 (the reconstruction bundle).} The statements
A16 and C01--C29 (V.2.2--V.2.8) hold simultaneously. Each retains its
own hypotheses. \emph{Status:} theorem. \emph{Lean:}
\href{https://github.com/yang-mills-group/OpenAI-YM-20261028/blob/3456e5e75cf8d4cd558955a5c7638fba2d68f464/YM/Integrated.lean\#L182-L212}{\texttt{Integrated.reconstructionArcs}},
\href{https://github.com/yang-mills-group/OpenAI-YM-20261028/blob/3456e5e75cf8d4cd558955a5c7638fba2d68f464/YM/Bid/Arena.lean\#L15-L45}{\texttt{Bid.Arcs}}.

\textbf{Proposition V.2.2 (strict-time continuation; C01--C03).} For the
tables \(S,T\), the OS chamber \(V\) with Hilbert space \(\mathcal H\),
a differential plan \(P\) and a spread \(A\), assume the slab hypotheses
(H):

\begin{itemize}
\tightlist
\item
  \(S\) and \(T\) satisfy the Euclidean hierarchy;
\item
  \(T\) respects \(S\);
\item
  the inclusions of \(P\) are compatible with \(S\) and \(T\), and its
  derivatives with \(T\);
\item
  the spread \(A\) obeys the slab bound (SB) (§2), is smooth in the
  insertion positions on every valid slab, and differentiating an
  insertion replaces its label by the corresponding derivative under
  \(P\).
\end{itemize}

Then:

\begin{enumerate}
\def\labelenumi{(\alph{enumi})}
\tightlist
\item
  (C01) There are linear maps
  \(B_p:\mathcal F\to\mathcal B(\mathcal H)\), \(p\ge0\), such that:
\end{enumerate}

\begin{itemize}
\tightlist
\item
  for \(p>0\), \(B_p(\mathrm{pin}_dv)\) is the slab operator of width
  \(2p\) of the one-point word made of the representative of \(v\)
  placed at \((\vec 0,p)\);
\item
  for every \(u\in\mathcal F\) there are \(C>0\) and \(a\in\mathbb N\)
  with \(\|B_p(u)\|\le C\,p^{-a}\) for \(0<p<1\).
\end{itemize}

\begin{enumerate}
\def\labelenumi{(\alph{enumi})}
\setcounter{enumi}{1}
\tightlist
\item
  (C02) Let \(B\) be such a family. Put
  \[G_p(w,\vec a)=e^{-(\mathrm{Im}\,w-2p)H}\,e^{i\,\mathrm{Re}\,w\,H}\,U(\vec a),\]
  \[M_p(z;\vec x)(v)=\big\langle\Omega,\,B_p(v_1)\,G_p(z_2-z_1,\vec x_2-\vec x_1)\,B_p(v_2)\cdots G_p(z_n-z_{n-1},\vec x_n-\vec x_{n-1})\,B_p(v_n)\Omega\big\rangle .\]
  For \(n=0\) set \(M_p=\|\Omega\|^2\). Call \(p\) \emph{admissible} for
  \(z\) if \(0<p<1\) and \(\mathrm{Im}(z_j-z_i)>4p\) for \(i<j\). Then:
\end{enumerate}

\begin{itemize}
\tightlist
\item
  \(M_p=M_q\) whenever \(p\) and \(q\) are both admissible for \(z\);
\item
  \(M_p(iy;\vec x)(\mathrm{pin}\,v)=T\big((\vec x_k,y_k)_k;v\big)\) for
  \(y\in E_n\) and \(p\) admissible.
\end{itemize}

\begin{enumerate}
\def\labelenumi{(\alph{enumi})}
\setcounter{enumi}{2}
\tightlist
\item
  (C03, from (b) alone) There is a continuation \(F\) with
  \(F_n(z;\vec x)=M_p(z;\vec x)\) for \(z\in\Sigma_n\) and every
  admissible \(p\), satisfying:
\end{enumerate}

\begin{itemize}
\tightlist
\item
  (SC) for each fixed label tuple \(v\),
  \((z,\vec x)\mapsto F_n(z;\vec x)(v)\) is continuous on
  \(\Sigma_n\times(\mathbb R^3)^n\) and holomorphic in \(z\in\Sigma_n\),
  and \(F_n(iy;\vec x)(\mathrm{pin}\,v)=T((\vec x_k,y_k)_k;v)\) for
  \(y\in E_n\);
\item
  (PG) for all \(n\) and \(v\) there are \(m\in\mathbb N\) and \(C\ge0\)
  with \(\lvert F^y_n(x)(v)\rvert\le C\,k^{-m}\) for all \(x\), all
  \(k\in(0,1]\), and all \(y\) with \(y_j-y_i\ge k\) for \(i<j\).
\end{itemize}

\emph{Status:} theorem. The derivative and spatial regularity hypotheses
remain part of the stated input package. \emph{Lean:}

\begin{itemize}
\tightlist
\item
  \href{https://github.com/yang-mills-group/OpenAI-YM-20261028/blob/3456e5e75cf8d4cd558955a5c7638fba2d68f464/YM/Work/C01.lean\#L99-L117}{\texttt{C01.answer}},
  \href{https://github.com/yang-mills-group/OpenAI-YM-20261028/blob/3456e5e75cf8d4cd558955a5c7638fba2d68f464/YM/Work/C02.lean\#L168-L177}{\texttt{C02.answer}},
  \href{https://github.com/yang-mills-group/OpenAI-YM-20261028/blob/3456e5e75cf8d4cd558955a5c7638fba2d68f464/YM/Work/C03.lean\#L245-L248}{\texttt{C03.answer}}.
\item
  Statements
  \href{https://github.com/yang-mills-group/OpenAI-YM-20261028/blob/3456e5e75cf8d4cd558955a5c7638fba2d68f464/YM/Bid/Corners.lean\#L37-L38}{\texttt{Bid.C01}},
  \href{https://github.com/yang-mills-group/OpenAI-YM-20261028/blob/3456e5e75cf8d4cd558955a5c7638fba2d68f464/YM/Bid/Corners.lean\#L47-L48}{\texttt{Bid.C02}},
  \href{https://github.com/yang-mills-group/OpenAI-YM-20261028/blob/3456e5e75cf8d4cd558955a5c7638fba2d68f464/YM/Bid/Corners.lean\#L56-L57}{\texttt{Bid.C03}}.
\item
  Hypothesis and conclusion structures:
  \href{https://github.com/yang-mills-group/OpenAI-YM-20261028/blob/3456e5e75cf8d4cd558955a5c7638fba2d68f464/YM/Bid/Corners.lean\#L23-L30}{\texttt{Pantry}},
  \href{https://github.com/yang-mills-group/OpenAI-YM-20261028/blob/3456e5e75cf8d4cd558955a5c7638fba2d68f464/YM/Bid/Corners.lean\#L33-L35}{\texttt{Piece}},
  \href{https://github.com/yang-mills-group/OpenAI-YM-20261028/blob/3456e5e75cf8d4cd558955a5c7638fba2d68f464/YM/Bid/Corners.lean\#L40-L45}{\texttt{Slice}},
  \href{https://github.com/yang-mills-group/OpenAI-YM-20261028/blob/3456e5e75cf8d4cd558955a5c7638fba2d68f464/YM/Bid/Corners.lean\#L50-L54}{\texttt{Direct}}.
\item
  Compatibility definitions:
  \href{https://github.com/yang-mills-group/OpenAI-YM-20261028/blob/3456e5e75cf8d4cd558955a5c7638fba2d68f464/YM/Euclid/Differential.lean\#L58-L60}{\texttt{DiffPlan.Inclusions}},
  \href{https://github.com/yang-mills-group/OpenAI-YM-20261028/blob/3456e5e75cf8d4cd558955a5c7638fba2d68f464/YM/Euclid/Differential.lean\#L66-L72}{\texttt{DiffPlan.Orientation}}.
\item
  Definitions
  \href{https://github.com/yang-mills-group/OpenAI-YM-20261028/blob/3456e5e75cf8d4cd558955a5c7638fba2d68f464/YM/Fields/Tall.lean\#L19-L20}{\texttt{Small}},
  \href{https://github.com/yang-mills-group/OpenAI-YM-20261028/blob/3456e5e75cf8d4cd558955a5c7638fba2d68f464/YM/Fields/Tall.lean\#L49-L51}{\texttt{Small.matrix}},
  \href{https://github.com/yang-mills-group/OpenAI-YM-20261028/blob/3456e5e75cf8d4cd558955a5c7638fba2d68f464/YM/Fields/Tall.lean\#L62-L63}{\texttt{Small.rated}},
  \href{https://github.com/yang-mills-group/OpenAI-YM-20261028/blob/3456e5e75cf8d4cd558955a5c7638fba2d68f464/YM/Fields/Tall.lean\#L66-L67}{\texttt{Small.Defined}},
  \href{https://github.com/yang-mills-group/OpenAI-YM-20261028/blob/3456e5e75cf8d4cd558955a5c7638fba2d68f464/YM/Boundary/Contour.lean\#L57-L58}{\texttt{Waves}},
  \href{https://github.com/yang-mills-group/OpenAI-YM-20261028/blob/3456e5e75cf8d4cd558955a5c7638fba2d68f464/YM/Boundary/Contour.lean\#L71-L77}{\texttt{Waves.Region}},
  \href{https://github.com/yang-mills-group/OpenAI-YM-20261028/blob/3456e5e75cf8d4cd558955a5c7638fba2d68f464/YM/Shift/Ridge.lean\#L72-L73}{\texttt{Waves.Breathes}}.
\end{itemize}

\textbf{Proposition V.2.3 (Cauchy--Riemann identities and tempered
boundary values; C04).} Let \(F\) satisfy (SC) and (PG).

\begin{enumerate}
\def\labelenumi{(\alph{enumi})}
\tightlist
\item
  (C04) Three properties hold:
\end{enumerate}

\begin{itemize}
\tightlist
\item
  \begin{enumerate}
  \def\labelenumi{(\Roman{enumi})}
  \tightlist
  \item
    \(x\mapsto F^y_n(x)(v)f(x)\) is integrable for \(y\in E_n\) and
    \(f\in\mathcal S(\mathbb R^{4n})\).
  \end{enumerate}
\item
  (HC) \(y\mapsto\psi_v(y;f)\) is continuous on \(E_n\).
\item
  (CR) For \(a\in\mathbb R^n\) and \(y+sa\in E_n\),
  \[\tfrac{d}{dt}\psi_v(y+ta;f)\big|_{t=s}=\psi_v\big(y+sa;-i\,\partial_{\underline a}f\big),\qquad \partial_{\underline a}=\textstyle\sum_ka_k\,\partial/\partial x^0_k .\]
\end{itemize}

\begin{enumerate}
\def\labelenumi{(\alph{enumi})}
\setcounter{enumi}{1}
\tightlist
\item
  If \(F\) satisfies (I), (HC), (CR) and (PG), there is a boundary
  family \(W\) with (BV): the integrands defining \(\psi_v(y;f)\) are
  integrable for \(y\in E_n\), and \(\psi_v(y;f)\to W_n(v)(f)\) as
  \(y\to0\) within \(E_n\), for all \(n\), \(v\) and \(f\). In
  particular, the radial boundary limit (BV\(_{\rm ray}\)) holds:
  \(\psi_v(\varepsilon u;f)\to W_n(v)(f)\) as \(\varepsilon\downarrow0\)
  for each \(u\in E_n\).
\end{enumerate}

\emph{Status:} theorem. Part (b) lives in Shift/Breath and Shift/Osc,
which are filed under VI.1. It is used through
\href{https://github.com/yang-mills-group/OpenAI-YM-20261028/blob/3456e5e75cf8d4cd558955a5c7638fba2d68f464/YM/Org/Theatre.lean\#L29-L31}{\texttt{Show.Rhythm.opens}}
(VIII.1). The boundary family is an output of the reconstruction.
\emph{Lean:}

\begin{itemize}
\tightlist
\item
  \href{https://github.com/yang-mills-group/OpenAI-YM-20261028/blob/3456e5e75cf8d4cd558955a5c7638fba2d68f464/YM/Work/C04.lean\#L221-L222}{\texttt{C04.answer}},
  \href{https://github.com/yang-mills-group/OpenAI-YM-20261028/blob/3456e5e75cf8d4cd558955a5c7638fba2d68f464/YM/Bid/Corners.lean\#L59-L60}{\texttt{Bid.C04}}.
\item
  \href{https://github.com/yang-mills-group/OpenAI-YM-20261028/blob/3456e5e75cf8d4cd558955a5c7638fba2d68f464/YM/Shift/Ridge.lean\#L83-L84}{\texttt{Waves.Bounds}},
  \href{https://github.com/yang-mills-group/OpenAI-YM-20261028/blob/3456e5e75cf8d4cd558955a5c7638fba2d68f464/YM/Shift/Ridge.lean\#L80-L81}{\texttt{Waves.Lines}},
  \href{https://github.com/yang-mills-group/OpenAI-YM-20261028/blob/3456e5e75cf8d4cd558955a5c7638fba2d68f464/YM/Shift/Ridge.lean\#L86}{\texttt{Waves.Along}}.
\item
  \href{https://github.com/yang-mills-group/OpenAI-YM-20261028/blob/3456e5e75cf8d4cd558955a5c7638fba2d68f464/YM/Shift/Osc.lean\#L83-L86}{\texttt{Waves.jointly}},
  \href{https://github.com/yang-mills-group/OpenAI-YM-20261028/blob/3456e5e75cf8d4cd558955a5c7638fba2d68f464/YM/Boundary/Contour.lean\#L96-L100}{\texttt{Waves.NonTangentialFree}},
  \href{https://github.com/yang-mills-group/OpenAI-YM-20261028/blob/3456e5e75cf8d4cd558955a5c7638fba2d68f464/YM/Boundary/Contour.lean\#L89-L93}{\texttt{Waves.Opens}},
  \href{https://github.com/yang-mills-group/OpenAI-YM-20261028/blob/3456e5e75cf8d4cd558955a5c7638fba2d68f464/YM/Boundary/Contour.lean\#L102-L112}{\texttt{Waves.exposed}}.
\end{itemize}

\textbf{Proposition V.2.4 (Laplace identity and spectral support; C05,
C06).} Let \(F\) satisfy (SC), (PG), (CR) and (HC), and have boundary
family \(W\) (BV).

\begin{enumerate}
\def\labelenumi{(\alph{enumi})}
\item
  (C05; Laplace identity (L)) For \(y\in E_n\) and compactly supported
  \(f\in\mathcal S(\mathbb R^{4n})\),
  \[W_n(v)(\hat f)=\psi_v\big(y;\widehat{e^{\langle\underline y,\cdot\rangle}f}\big).\]
\item
  (C06, from (SC), (PG), (I) and (a)) \(\mathcal FW_n(v)\) is supported
  in
  \[\mathcal P_n=\Big\{p:\textstyle\sum_ky_kp^0_k\ge0\text{ for every nondecreasing }y\Big\}=\Big\{p:\textstyle\sum_kp^0_k=0,\ \sum_{k\ge j}p^0_k\ge0\ \forall j\Big\}.\]
\item
  (Lorentz covariance) If \(W\) also satisfies (Cov) (V.2.5), then
  \(\mathcal FW_n(v)\) is supported in
  \[\mathcal V_n=\bigcap_{\Lambda\in\mathcal L^\uparrow_+}\tilde\Lambda^{-1}\mathcal P_n ,\]
  where \(\tilde\Lambda=\mathsf r\Lambda\mathsf r\) is the dual action
  (\(\mathsf r\) reverses the spatial components). Every
  \(p\in\mathcal V_n\) has \(\sum_kp_k=0\) and
  \(\sum_{k\ge j}p_k\in\bar V_+\) for all \(j\). This is the
  cone-support hypothesis of C14--C16 and C19.
\end{enumerate}

\emph{Status:} theorem. (SC) is not used in C05 or C06, and (I) is not
used in C06. \emph{Lean:}

\begin{itemize}
\tightlist
\item
  \href{https://github.com/yang-mills-group/OpenAI-YM-20261028/blob/3456e5e75cf8d4cd558955a5c7638fba2d68f464/YM/Work/C05.lean\#L205-L230}{\texttt{C05.answer}},
  \href{https://github.com/yang-mills-group/OpenAI-YM-20261028/blob/3456e5e75cf8d4cd558955a5c7638fba2d68f464/YM/Work/C06.lean\#L136-L157}{\texttt{C06.answer}}.
\item
  Statements
  \href{https://github.com/yang-mills-group/OpenAI-YM-20261028/blob/3456e5e75cf8d4cd558955a5c7638fba2d68f464/YM/Bid/Corners.lean\#L62-L63}{\texttt{Bid.C05}},
  \href{https://github.com/yang-mills-group/OpenAI-YM-20261028/blob/3456e5e75cf8d4cd558955a5c7638fba2d68f464/YM/Bid/Corners.lean\#L65-L66}{\texttt{Bid.C06}}.
\item
  \href{https://github.com/yang-mills-group/OpenAI-YM-20261028/blob/3456e5e75cf8d4cd558955a5c7638fba2d68f464/YM/Shift/Sol.lean\#L101-L102}{\texttt{Waves.during}},
  \href{https://github.com/yang-mills-group/OpenAI-YM-20261028/blob/3456e5e75cf8d4cd558955a5c7638fba2d68f464/YM/Shift/Sol.lean\#L106}{\texttt{Surface.cone\_tests}},
  \href{https://github.com/yang-mills-group/OpenAI-YM-20261028/blob/3456e5e75cf8d4cd558955a5c7638fba2d68f464/YM/Shift/Closing.lean\#L20}{\texttt{half}},
  \href{https://github.com/yang-mills-group/OpenAI-YM-20261028/blob/3456e5e75cf8d4cd558955a5c7638fba2d68f464/YM/Shift/Closing.lean\#L23}{\texttt{coneM}},
  \href{https://github.com/yang-mills-group/OpenAI-YM-20261028/blob/3456e5e75cf8d4cd558955a5c7638fba2d68f464/YM/Shift/Closing.lean\#L106-L121}{\texttt{tents}},
  \href{https://github.com/yang-mills-group/OpenAI-YM-20261028/blob/3456e5e75cf8d4cd558955a5c7638fba2d68f464/YM/Bid/Tar.lean\#L39}{\texttt{Bid.cones}}.
\end{itemize}

\textbf{Proposition V.2.5 (Poincaré covariance; A16, C07, C08, C09,
C28).}

\emph{Data.}

\begin{itemize}
\tightlist
\item
  A rotation structure \(R\) on \(S\) (VI.1): an action of \(O(4)\) on
  the \(F_d\), with Lie action \(\rho\) of \(\mathfrak{so}(4)\). It is
  subject to Euclidean covariance (the functions \(T(x;v)\) are
  \(O(4)\)-covariant) and the rotation laws of VI.1.
\item
  A spin structure \(K\): representations \(K(\Lambda)\) of
  \(\mathcal L^\uparrow_+\) on \(\mathcal F\) and on each \(F_d\), with
  spin compatibility with \(R\), relative to the real structure of
  \(S\), which consists of:

  \begin{itemize}
  \tightlist
  \item
    compatibility with \(\mathrm{pin}_d\);
  \item
    the representation property and continuity on each
    finite-dimensional level;
  \item
    \(\frac{d}{dt}K_d(\Lambda_a(t))\vert_{t=0}=-i\rho(e_0\wedge e_a)\);
  \item
    \(K=R\) on spatial rotations;
  \item
    commutation with the reflection adjoint.
  \end{itemize}
\end{itemize}

\begin{enumerate}
\def\labelenumi{(\alph{enumi})}
\item
  (A16) \(\mathcal L^\uparrow_+\) is generated, as an abstract group, by
  the spatial rotations and the coordinate boosts.
\item
  (C08, no hypotheses) For all \(n\) and
  \(f\in\mathcal S(\mathbb R^{4n})\):
\end{enumerate}

\begin{itemize}
\tightlist
\item
  \(g\mapsto g_\ast f\) is continuous from \(\mathcal P\) into
  \(\mathcal S(\mathbb R^{4n})\);
\item
  \(t^{-1}(\Lambda_a(t)_\ast f-f)\to-D_af\) in \(\mathcal S\) as
  \(t\to0\), where
  \(D_a=\sum_k(x^a_k\partial_{x^0_k}+x^0_k\partial_{x^a_k})\).
\end{itemize}

\begin{enumerate}
\def\labelenumi{(\alph{enumi})}
\setcounter{enumi}{2}
\item
  (C07) Assume \(S\) and \(T\) satisfy the hierarchy, the inclusions of
  \(P\) are compatible with \(S\), \(T\) respects \(S\), and Euclidean
  covariance, the rotation laws of VI.1 and spin compatibility hold.
  Then every \(F\) with (SC), (PG), (I), (CR) and (HC) satisfies the
  boost identity (B\(_y\)): for all \(y\in E_n\), tags \(v\), spatial
  axes \(a\) and \(f\in\mathcal S(\mathbb R^{4n})\),
  \[\sum_{k=1}^n\psi_{v^{(k)}}(y;f)=\psi_v\big(y;D_{a,y}f\big),\qquad D_{a,y}=\sum_k\Big(x^a_k\,\partial_{x^0_k}+(x^0_k+iy_k)\,\partial_{x^a_k}\Big).\]
  Here \(v^{(k)}\) is \(v\) with \(v_k\) replaced by
  \(-i\rho(e_0\wedge e_a)v_k\). Among the table, plan and spin inputs,
  the boost argument uses the Euclidean hierarchy and rotation data.
\item
  (C09) Assume A16, C08 and the same table, plan, rotation and spin
  hypotheses as in (c). If \(F\) satisfies (SC), (PG) and (B\(_y\)) and
  has boundary family \(W\) (BV), then \(W\) satisfies (Cov):
\end{enumerate}

\begin{itemize}
\tightlist
\item
  \(W_n(v)(f(\cdot-a))=W_n(v)(f)\) for \(a\in\mathbb R^4\);
\item
  \(W_n(K(\Lambda)v)(\Lambda_\ast f)=W_n(v)(f)\) for all
  \(\Lambda\in\mathcal L^\uparrow_+\).
\end{itemize}

\begin{enumerate}
\def\labelenumi{(\alph{enumi})}
\setcounter{enumi}{4}
\tightlist
\item
  (C28) Assume C08, \(S\) symmetric, the inclusions of \(P\) compatible
  with \(S\), spin compatibility, and \(J\) the induced star. Let
  \(g=(a,\Lambda)\) act on words letterwise by
  \((u,f)\mapsto(K(\Lambda)u,g_\ast f)\). Then (Cov) implies:
\end{enumerate}

\begin{itemize}
\tightlist
\item
  (invariance of the word Gram form)
  \(\langle g\cdot b,g\cdot c\rangle_W=\langle b,c\rangle_W\);
\item
  (continuity of word matrix elements)
  \(g\mapsto\langle b,g\cdot c\rangle_W\) is continuous.
\end{itemize}

\emph{Status:} theorem. \emph{Lean:}

\begin{itemize}
\tightlist
\item
  \href{https://github.com/yang-mills-group/OpenAI-YM-20261028/blob/3456e5e75cf8d4cd558955a5c7638fba2d68f464/YM/Work/A16.lean\#L91-L101}{\texttt{A16.answer}}
  (statement
  \href{https://github.com/yang-mills-group/OpenAI-YM-20261028/blob/3456e5e75cf8d4cd558955a5c7638fba2d68f464/YM/Bid/Angular.lean\#L116}{\texttt{Bid.A16}},
  \href{https://github.com/yang-mills-group/OpenAI-YM-20261028/blob/3456e5e75cf8d4cd558955a5c7638fba2d68f464/YM/Bid/Angular.lean\#L114-L115}{\texttt{beamGens}}).
\item
  \href{https://github.com/yang-mills-group/OpenAI-YM-20261028/blob/3456e5e75cf8d4cd558955a5c7638fba2d68f464/YM/Work/C07.lean\#L53-L66}{\texttt{C07.answer}},
  \href{https://github.com/yang-mills-group/OpenAI-YM-20261028/blob/3456e5e75cf8d4cd558955a5c7638fba2d68f464/YM/Work/C08.lean\#L285-L286}{\texttt{C08.answer}},
  \href{https://github.com/yang-mills-group/OpenAI-YM-20261028/blob/3456e5e75cf8d4cd558955a5c7638fba2d68f464/YM/Work/C09.lean\#L403-L442}{\texttt{C09.answer}},
  \href{https://github.com/yang-mills-group/OpenAI-YM-20261028/blob/3456e5e75cf8d4cd558955a5c7638fba2d68f464/YM/Work/C28.lean\#L134-L143}{\texttt{C28.answer}}.
\item
  Statements
  \href{https://github.com/yang-mills-group/OpenAI-YM-20261028/blob/3456e5e75cf8d4cd558955a5c7638fba2d68f464/YM/Bid/Corners.lean\#L95-L97}{\texttt{Bid.C07}},
  \href{https://github.com/yang-mills-group/OpenAI-YM-20261028/blob/3456e5e75cf8d4cd558955a5c7638fba2d68f464/YM/Bid/Corners.lean\#L99-L101}{\texttt{Bid.C08}},
  \href{https://github.com/yang-mills-group/OpenAI-YM-20261028/blob/3456e5e75cf8d4cd558955a5c7638fba2d68f464/YM/Bid/Corners.lean\#L103-L106}{\texttt{Bid.C09}},
  \href{https://github.com/yang-mills-group/OpenAI-YM-20261028/blob/3456e5e75cf8d4cd558955a5c7638fba2d68f464/YM/Bid/Grow.lean\#L50-L52}{\texttt{Bid.C28}}.
\item
  Definitions
  \href{https://github.com/yang-mills-group/OpenAI-YM-20261028/blob/3456e5e75cf8d4cd558955a5c7638fba2d68f464/YM/Bid/Corners.lean\#L92-L93}{\texttt{Boosty}},
  \href{https://github.com/yang-mills-group/OpenAI-YM-20261028/blob/3456e5e75cf8d4cd558955a5c7638fba2d68f464/YM/Bid/Corners.lean\#L84-L85}{\texttt{inclined}},
  \href{https://github.com/yang-mills-group/OpenAI-YM-20261028/blob/3456e5e75cf8d4cd558955a5c7638fba2d68f464/YM/Bid/Corners.lean\#L87-L88}{\texttt{kick}},
  \href{https://github.com/yang-mills-group/OpenAI-YM-20261028/blob/3456e5e75cf8d4cd558955a5c7638fba2d68f464/YM/Shift/Closing.lean\#L83-L85}{\texttt{Urs}},
  \href{https://github.com/yang-mills-group/OpenAI-YM-20261028/blob/3456e5e75cf8d4cd558955a5c7638fba2d68f464/YM/Fields/Star.lean\#L26-L35}{\texttt{Spins.Descends}},
  \href{https://github.com/yang-mills-group/OpenAI-YM-20261028/blob/3456e5e75cf8d4cd558955a5c7638fba2d68f464/YM/Fields/Motion.lean\#L23-L24}{\texttt{Surface.Transforms}},
  \href{https://github.com/yang-mills-group/OpenAI-YM-20261028/blob/3456e5e75cf8d4cd558955a5c7638fba2d68f464/YM/Fields/Motion.lean\#L26-L27}{\texttt{Surface.Strong}}.
\item
  Rotation laws:
  \href{https://github.com/yang-mills-group/OpenAI-YM-20261028/blob/3456e5e75cf8d4cd558955a5c7638fba2d68f464/YM/Euclid/Stress.lean\#L31-L42}{\texttt{Table.Rotation.Laws}}.
\end{itemize}

\textbf{Proposition V.2.6 (Fourier--Laplace continuation in the forward
tube; C10--C16).}

\emph{Definitions.}

\begin{itemize}
\item
  Tube and cones:

  \begin{itemize}
  \tightlist
  \item
    \(V^+_n=\{\eta\in\mathbb R^{4n}:\eta_{k+1}-\eta_k\in V_+,\ 1\le k<n\}\)
    (the open imaginary-part cone);
  \item
    \(\bar V^+_n\) is the analogue with \(\bar V_+\) (the closed
    imaginary-part cone);
  \item
    \(\mathcal T_n=\{z\in\mathbb C^{4n}:\mathrm{Im}\,z\in V^+_n\}\) (the
    tube). This is the forward tube in absolute coordinates; there is no
    condition on \(\mathrm{Im}\,z_1\).
  \end{itemize}
\item
  A \emph{cutoff} \(\chi_n\) is a real smooth function whose derivatives
  are all polynomially bounded, equal to \(1\) on the
  \(1\)-neighbourhood of \(\mathcal V_n\), with closed support contained
  in its open \(3\)-neighbourhood.
\item
  \emph{Fourier--Laplace continuation:}
  \[\mathcal W_n(v)(z)=(2\pi)^{-4n}\,(\mathcal FW_n(v))\big(e^{i\langle\cdot,z\rangle}\chi_n\big),\qquad z\in\mathcal T_n .\]
\item
  \emph{Partial boundary values:}
  \(W_n(v)[\eta](https://github.com/scottnarmstrong/yang-mills-informalization/blob/main/outline/units/f)=(\mathcal FW_n(v))\big(e^{-\langle\eta,\cdot\rangle}\chi_n\,\check f\big)\),
  formally \(W_n(v)(\cdot+i\eta)\).
\end{itemize}

\emph{Statements.}

\begin{enumerate}
\def\labelenumi{(\alph{enumi})}
\item
  (C10, cone duality) For \(p\in\mathcal V_n\) and
  \(\eta\in\bar V^+_n\), \(\langle p,\eta\rangle\ge0\). For
  \(\eta\in V^+_n\) there is \(c>0\) with
  \(\langle p,u\rangle\ge c\lvert p\rvert\) for all \(p\in\mathcal V_n\)
  and all \(u\) near \(\eta\).
\item
  (C11) Cutoffs exist.
\item
  (C12) Given (a), every cutoff satisfies the damping properties (D):
\end{enumerate}

\begin{itemize}
\tightlist
\item
  \(e^{-\langle\eta,\cdot\rangle}\chi_n\) has temperate growth for
  \(\eta\in\bar V^+_n\);
\item
  \(\eta\mapsto e^{-\langle\eta,\cdot\rangle}\chi_nf\) is continuous
  from \(\bar V^+_n\) into \(\mathcal S\).
\end{itemize}

\begin{enumerate}
\def\labelenumi{(\alph{enumi})}
\setcounter{enumi}{3}
\item
  (C13) Given (a), \(\mathcal W_n(v)\) is holomorphic on
  \(\mathcal T_n\).
\item
  (C14, the integral representation) Given (c), for \(\eta\in V^+_n\)
  the identity
  \(W_n(v)[\eta](https://github.com/scottnarmstrong/yang-mills-informalization/blob/main/outline/units/f)=\int\mathcal W_n(v)(x+i\eta)f(x)\,dx\)
  holds, with integrable integrand.
\item
  (C15) Suppose:
\end{enumerate}

\begin{itemize}
\item
  (c) holds and \(\mathrm{supp}\mathcal FW\subseteq\mathcal V_n\);
\item
  (d) and (e) hold;
\item
  \(F\) satisfies (SC), (PG), (I) and the Laplace identity V.2.4(a).

  Then \(\mathcal W_n(v)(x+i\underline y)=F^y_n(x)(v)\) for all
  \(y\in E_n\) and \(x\in\mathbb R^{4n}\). Together with (d) and (e)
  this is (TC).
\end{itemize}

\begin{enumerate}
\def\labelenumi{(\alph{enumi})}
\setcounter{enumi}{6}
\tightlist
\item
  (C16) Given \(\mathrm{supp}\mathcal FW\subseteq\mathcal V_n\) and
  (Cov), \(\mathcal W_n(v)(z)=\mathcal W_n(K(\Lambda)v)(\Lambda z)\) for
  \(z\in\mathcal T_n\) and \(\Lambda\in\mathcal L^\uparrow_+\), with
  \(\Lambda\) acting on real and imaginary parts.
\end{enumerate}

\emph{Status:} theorem. Several listed hypotheses are not used (§5).
\emph{Lean:}

\begin{itemize}
\tightlist
\item
  \href{https://github.com/yang-mills-group/OpenAI-YM-20261028/blob/3456e5e75cf8d4cd558955a5c7638fba2d68f464/YM/Work/C10.lean\#L71}{\texttt{C10}},
  \href{https://github.com/yang-mills-group/OpenAI-YM-20261028/blob/3456e5e75cf8d4cd558955a5c7638fba2d68f464/YM/Work/C11.lean\#L57-L100}{\texttt{C11}},
  \href{https://github.com/yang-mills-group/OpenAI-YM-20261028/blob/3456e5e75cf8d4cd558955a5c7638fba2d68f464/YM/Work/C12.lean\#L150-L169}{\texttt{C12}},
  \href{https://github.com/yang-mills-group/OpenAI-YM-20261028/blob/3456e5e75cf8d4cd558955a5c7638fba2d68f464/YM/Work/C13.lean\#L127-L132}{\texttt{C13}},
  \href{https://github.com/yang-mills-group/OpenAI-YM-20261028/blob/3456e5e75cf8d4cd558955a5c7638fba2d68f464/YM/Work/C14.lean\#L257-L260}{\texttt{C14}},
  \href{https://github.com/yang-mills-group/OpenAI-YM-20261028/blob/3456e5e75cf8d4cd558955a5c7638fba2d68f464/YM/Work/C15.lean\#L120-L131}{\texttt{C15}},
  \href{https://github.com/yang-mills-group/OpenAI-YM-20261028/blob/3456e5e75cf8d4cd558955a5c7638fba2d68f464/YM/Work/C16.lean\#L75-L91}{\texttt{C16}}
  (all answer theorems).
\item
  Statements
  \href{https://github.com/yang-mills-group/OpenAI-YM-20261028/blob/3456e5e75cf8d4cd558955a5c7638fba2d68f464/YM/Bid/Tar.lean\#L20}{\texttt{Bid.C10}},
  \href{https://github.com/yang-mills-group/OpenAI-YM-20261028/blob/3456e5e75cf8d4cd558955a5c7638fba2d68f464/YM/Bid/Tar.lean\#L22}{\texttt{C11}},
  \href{https://github.com/yang-mills-group/OpenAI-YM-20261028/blob/3456e5e75cf8d4cd558955a5c7638fba2d68f464/YM/Bid/Tar.lean\#L24}{\texttt{C12}},
  \href{https://github.com/yang-mills-group/OpenAI-YM-20261028/blob/3456e5e75cf8d4cd558955a5c7638fba2d68f464/YM/Bid/Tar.lean\#L27-L29}{\texttt{C13}},
  \href{https://github.com/yang-mills-group/OpenAI-YM-20261028/blob/3456e5e75cf8d4cd558955a5c7638fba2d68f464/YM/Bid/Tar.lean\#L41-L42}{\texttt{C14}},
  \href{https://github.com/yang-mills-group/OpenAI-YM-20261028/blob/3456e5e75cf8d4cd558955a5c7638fba2d68f464/YM/Bid/Tar.lean\#L44-L46}{\texttt{C15}},
  \href{https://github.com/yang-mills-group/OpenAI-YM-20261028/blob/3456e5e75cf8d4cd558955a5c7638fba2d68f464/YM/Bid/Tar.lean\#L53-L54}{\texttt{C16}}.
\item
  Definitions
  \href{https://github.com/yang-mills-group/OpenAI-YM-20261028/blob/3456e5e75cf8d4cd558955a5c7638fba2d68f464/YM/Bid/Tar.lean\#L16-L18}{\texttt{Wide}},
  \href{https://github.com/yang-mills-group/OpenAI-YM-20261028/blob/3456e5e75cf8d4cd558955a5c7638fba2d68f464/YM/Bid/Tar.lean\#L33-L37}{\texttt{Income}},
  \href{https://github.com/yang-mills-group/OpenAI-YM-20261028/blob/3456e5e75cf8d4cd558955a5c7638fba2d68f464/YM/Bid/Tar.lean\#L50-L51}{\texttt{wave\_turn}},
  \href{https://github.com/yang-mills-group/OpenAI-YM-20261028/blob/3456e5e75cf8d4cd558955a5c7638fba2d68f464/YM/Shift/Sol.lean\#L60-L64}{\texttt{Blanket}},
  \href{https://github.com/yang-mills-group/OpenAI-YM-20261028/blob/3456e5e75cf8d4cd558955a5c7638fba2d68f464/YM/Shift/Sol.lean\#L70-L72}{\texttt{Blanket.Modified}},
  \href{https://github.com/yang-mills-group/OpenAI-YM-20261028/blob/3456e5e75cf8d4cd558955a5c7638fba2d68f464/YM/Shift/Sol.lean\#L78-L80}{\texttt{Blanket.partialD}},
  \href{https://github.com/yang-mills-group/OpenAI-YM-20261028/blob/3456e5e75cf8d4cd558955a5c7638fba2d68f464/YM/Shift/Sol.lean\#L109-L111}{\texttt{Surface.Aurora}},
  \href{https://github.com/yang-mills-group/OpenAI-YM-20261028/blob/3456e5e75cf8d4cd558955a5c7638fba2d68f464/YM/Shift/Sol.lean\#L116}{\texttt{Aurora.emit}},
  \href{https://github.com/yang-mills-group/OpenAI-YM-20261028/blob/3456e5e75cf8d4cd558955a5c7638fba2d68f464/YM/Shift/Sol.lean\#L118-L124}{\texttt{Aurora.FullBody}},
  \href{https://github.com/yang-mills-group/OpenAI-YM-20261028/blob/3456e5e75cf8d4cd558955a5c7638fba2d68f464/YM/Shift/Sol.lean\#L45}{\texttt{tube}},
  \href{https://github.com/yang-mills-group/OpenAI-YM-20261028/blob/3456e5e75cf8d4cd558955a5c7638fba2d68f464/YM/Shift/Closing.lean\#L70}{\texttt{shores}},
  \href{https://github.com/yang-mills-group/OpenAI-YM-20261028/blob/3456e5e75cf8d4cd558955a5c7638fba2d68f464/YM/Shift/Closing.lean\#L71}{\texttt{fields}}.
\end{itemize}

\textbf{Proposition V.2.7 (locality; C17--C21, C27).}

\emph{Setting.}

\begin{itemize}
\tightlist
\item
  Let \(n\ge2\) and let \(Q=(i,i+1)\) be an adjacent pair. \(s\) is the
  transposition of \(i\) and \(i+1\), and \(\sigma\) swaps slots \(i\)
  and \(i+1\) of real, imaginary or complex configurations.
\item
  \(\Gamma_Q=\{x:x_{i+1}-x_i\text{ spacelike}\}\).
\item
  A height profile \(y\) is a \emph{spectator profile} for \(Q\) if
  \(y_i=y_{i+1}\) and \(y_k<y_l\) for all other pairs \(k<l\). Then
  \(\underline y\) and \(\sigma\underline y\) lie in \(\bar V^+_n\), and
  for every real \(x\) the point \(x+i\underline y\) lies in the
  closures of \(\mathcal T_n\) and of \(\sigma^{-1}\mathcal T_n\).
\item
  \(\Lambda_e(\zeta)\) is the complexified boost of rapidity
  \(\zeta\in\mathbb C\) along a unit spatial direction \(e\).
\end{itemize}

\emph{Statements.}

\begin{enumerate}
\def\labelenumi{(\alph{enumi})}
\tightlist
\item
  (C17) There are spectator profiles \(y_m\to0\) with the following
  property. For every \(m\) and every \(x\in\Gamma_Q\) there are a unit
  direction \(e\), \(t>0\), and neighbourhoods \(U_{\rm o},U_{\rm t}\)
  of \(x+i\underline{y_m}\) such that, for every \(b\in[0,t]\):
\end{enumerate}

\begin{itemize}
\item
  \(\Lambda_e(ib)w\in\mathcal T_n\) for all
  \(w\in U_{\rm o}\cap\mathcal T_n\);
\item
  \(\sigma\Lambda_e(-ib)w\in\mathcal T_n\) for all
  \(w\in U_{\rm t}\cap\sigma^{-1}\mathcal T_n\);
\item
  the same holds at \(b=t\) for \(w=x+i\underline{y_m}\) itself.

  The profiles are \(y_m=(m+1)^{-1}d\), where \(d_k=k-1\) except that
  \(d_{i+1}=i-1\). Thus the entries at \(i\) and \(i+1\) agree; all
  later entries keep their original values.
\end{itemize}

\begin{enumerate}
\def\labelenumi{(\alph{enumi})}
\setcounter{enumi}{1}
\item
  (C18) Under spin compatibility, \(t\mapsto K_d(\Lambda_e(t))\) extends
  to an entire function of \(t\in\mathbb C\), for each \(d\) and \(e\).
\item
  (C19) Assume (TC), cone support V.2.4(c), spin compatibility, \(S\)
  symmetric, the inclusions of \(P\) compatible with \(S\), (b), and the
  real covariance V.2.6(g). Let \(v\) be labels, \(Q\) a pair, and \(y\)
  a spectator profile admitting the paths of (a). Then every
  \(x\in\Gamma_Q\) carries two germs:
\end{enumerate}

\begin{itemize}
\tightlist
\item
  \(g^{\rm o}_x\), holomorphic near \(x+i\underline y\) and equal to
  \(\mathcal W_n(v)\) on \(\mathcal T_n\) near that point;
\item
  \(g^{\rm t}_x\), equal to \(\mathcal W_n(v\circ s)\circ\sigma\) on
  \(\sigma^{-1}\mathcal T_n\) near that point.
\end{itemize}

\begin{enumerate}
\def\labelenumi{(\alph{enumi})}
\setcounter{enumi}{3}
\item
  (C20) Assume \(S\) and \(T\) symmetric, the inclusions of \(P\)
  compatible with \(S\), \(T\) smooth, \(T\) respects \(S\), (SC) and
  (TC). For any such germs and any \(x\in\Gamma_Q\) with all
  \(x^0_k=0\), \(g^{\rm o}_x=g^{\rm t}_x\) near \(x+i\underline y\).
\item
  (C21) Assume (D) and (TC). Each \(x\in\Gamma_Q\) has a radius \(r>0\)
  such that, for compactly supported \(f\) with support in \(B(x,r)\):
  \[W_n(v)[\underline y](https://github.com/scottnarmstrong/yang-mills-informalization/blob/main/outline/units/f)=\int g^{\rm o}_x(\xi+i\underline y)f(\xi)\,d\xi,\qquad W_n(v\circ s)[\sigma\underline y](https://github.com/scottnarmstrong/yang-mills-informalization/blob/main/outline/units/f\circ\sigma)=\int g^{\rm t}_x(\xi+i\underline y)f(\xi)\,d\xi.\]
\item
  (Boundary exchange) Suppose:
\end{enumerate}

\begin{itemize}
\item
  \(\mathrm{supp}\mathcal FW_n(u)\subseteq\mathcal V_n\) for all \(u\);
\item
  the cutoff satisfies (D);
\item
  for a sequence of spectator profiles \(y_m\to0\) there are germs
  satisfying (d) and (e).

  Then \(W_n(v)(f)=W_n(v\circ s)(f\circ\sigma)\) for every
  \(f\in\mathcal S\) with \(\mathrm{supp}f\subset\Gamma_Q\). When this
  holds for every arity, label tuple and adjacent pair, \(W\) satisfies
  \emph{joint locality}.
\end{itemize}

\begin{enumerate}
\def\labelenumi{(\alph{enumi})}
\setcounter{enumi}{6}
\tightlist
\item
  (C27) Joint locality implies \emph{word locality}:
  \(W(a\,(u,f)(u',g)\,c)=W(a\,(u',g)(u,f)\,c)\) for all words \(a\),
  \(c\), labels \(u\), \(u'\), and tests \(f\), \(g\) such that
  \(\mathrm{supp}f-\mathrm{supp}g\) consists of spacelike vectors.
\end{enumerate}

\emph{Status:} theorem. \emph{Lean:}

\begin{itemize}
\tightlist
\item
  \href{https://github.com/yang-mills-group/OpenAI-YM-20261028/blob/3456e5e75cf8d4cd558955a5c7638fba2d68f464/YM/Work/C17.lean\#L277-L282}{\texttt{C17.answer}},
  \href{https://github.com/yang-mills-group/OpenAI-YM-20261028/blob/3456e5e75cf8d4cd558955a5c7638fba2d68f464/YM/Work/C18.lean\#L58-L79}{\texttt{C18.answer}},
  \href{https://github.com/yang-mills-group/OpenAI-YM-20261028/blob/3456e5e75cf8d4cd558955a5c7638fba2d68f464/YM/Work/C19.lean\#L213-L226}{\texttt{C19.answer}},
  \href{https://github.com/yang-mills-group/OpenAI-YM-20261028/blob/3456e5e75cf8d4cd558955a5c7638fba2d68f464/YM/Work/C20.lean\#L242-L307}{\texttt{C20.answer}},
  \href{https://github.com/yang-mills-group/OpenAI-YM-20261028/blob/3456e5e75cf8d4cd558955a5c7638fba2d68f464/YM/Work/C21.lean\#L66-L87}{\texttt{C21.answer}},
  \href{https://github.com/yang-mills-group/OpenAI-YM-20261028/blob/3456e5e75cf8d4cd558955a5c7638fba2d68f464/YM/Work/C27.lean\#L87-L92}{\texttt{C27.answer}}.
\item
  Statements
  \href{https://github.com/yang-mills-group/OpenAI-YM-20261028/blob/3456e5e75cf8d4cd558955a5c7638fba2d68f464/YM/Bid/Tar.lean\#L68-L69}{\texttt{Bid.C17}},
  \href{https://github.com/yang-mills-group/OpenAI-YM-20261028/blob/3456e5e75cf8d4cd558955a5c7638fba2d68f464/YM/Bid/Tar.lean\#L75-L76}{\texttt{C18}},
  \href{https://github.com/yang-mills-group/OpenAI-YM-20261028/blob/3456e5e75cf8d4cd558955a5c7638fba2d68f464/YM/Bid/Tar.lean\#L78-L81}{\texttt{C19}},
  \href{https://github.com/yang-mills-group/OpenAI-YM-20261028/blob/3456e5e75cf8d4cd558955a5c7638fba2d68f464/YM/Bid/Tar.lean\#L84-L86}{\texttt{C20}},
  \href{https://github.com/yang-mills-group/OpenAI-YM-20261028/blob/3456e5e75cf8d4cd558955a5c7638fba2d68f464/YM/Bid/Tar.lean\#L89-L90}{\texttt{C21}},
  \href{https://github.com/yang-mills-group/OpenAI-YM-20261028/blob/3456e5e75cf8d4cd558955a5c7638fba2d68f464/YM/Bid/Grow.lean\#L48}{\texttt{C27}}.
\item
  Path definitions:
  \href{https://github.com/yang-mills-group/OpenAI-YM-20261028/blob/3456e5e75cf8d4cd558955a5c7638fba2d68f464/YM/Bid/Tar.lean\#L57-L60}{\texttt{augment}},
  \href{https://github.com/yang-mills-group/OpenAI-YM-20261028/blob/3456e5e75cf8d4cd558955a5c7638fba2d68f464/YM/Bid/Tar.lean\#L63-L66}{\texttt{travelable}},
  \href{https://github.com/yang-mills-group/OpenAI-YM-20261028/blob/3456e5e75cf8d4cd558955a5c7638fba2d68f464/YM/Bid/Tar.lean\#L72-L73}{\texttt{Weighty}}.
\item
  Pair definitions:
  \href{https://github.com/yang-mills-group/OpenAI-YM-20261028/blob/3456e5e75cf8d4cd558955a5c7638fba2d68f464/YM/Shift/Swap.lean\#L19-L22}{\texttt{NearPair}},
  \href{https://github.com/yang-mills-group/OpenAI-YM-20261028/blob/3456e5e75cf8d4cd558955a5c7638fba2d68f464/YM/Shift/Swap.lean\#L49}{\texttt{gap}},
  \href{https://github.com/yang-mills-group/OpenAI-YM-20261028/blob/3456e5e75cf8d4cd558955a5c7638fba2d68f464/YM/Shift/Swap.lean\#L53}{\texttt{spectator}},
  \href{https://github.com/yang-mills-group/OpenAI-YM-20261028/blob/3456e5e75cf8d4cd558955a5c7638fba2d68f464/YM/Shift/Swap.lean\#L56-L61}{\texttt{Fleck}}.
\item
  Germ definitions:
  \href{https://github.com/yang-mills-group/OpenAI-YM-20261028/blob/3456e5e75cf8d4cd558955a5c7638fba2d68f464/YM/Shift/Exchange.lean\#L33-L36}{\texttt{Panes}},
  \href{https://github.com/yang-mills-group/OpenAI-YM-20261028/blob/3456e5e75cf8d4cd558955a5c7638fba2d68f464/YM/Shift/Exchange.lean\#L40-L41}{\texttt{sameGround}},
  \href{https://github.com/yang-mills-group/OpenAI-YM-20261028/blob/3456e5e75cf8d4cd558955a5c7638fba2d68f464/YM/Shift/Exchange.lean\#L62-L66}{\texttt{Forms}}.
\item
  Locality lemmas:
  \href{https://github.com/yang-mills-group/OpenAI-YM-20261028/blob/3456e5e75cf8d4cd558955a5c7638fba2d68f464/YM/Shift/Exchange.lean\#L43-L56}{\texttt{Panes.reaches}},
  \href{https://github.com/yang-mills-group/OpenAI-YM-20261028/blob/3456e5e75cf8d4cd558955a5c7638fba2d68f464/YM/Shift/Exchange.lean\#L68-L90}{\texttt{Panes.forgets}},
  \href{https://github.com/yang-mills-group/OpenAI-YM-20261028/blob/3456e5e75cf8d4cd558955a5c7638fba2d68f464/YM/Shift/Exchange.lean\#L92-L112}{\texttt{NearPair.actor}},
  \href{https://github.com/yang-mills-group/OpenAI-YM-20261028/blob/3456e5e75cf8d4cd558955a5c7638fba2d68f464/YM/Shift/Exchange.lean\#L115-L116}{\texttt{Actor}},
  \href{https://github.com/yang-mills-group/OpenAI-YM-20261028/blob/3456e5e75cf8d4cd558955a5c7638fba2d68f464/YM/Fields/Stems.lean\#L33-L34}{\texttt{Surface.Regional}}.
\end{itemize}

\textbf{Proposition V.2.8 (real-time vectors, Gram limits, nulls;
C22--C26, C29).}

\emph{Heated vectors.} For \(n\ge1\), take \(v\in\mathcal F^n\), heights
\(y\in E_n\) with positive entries, \(p\in(0,1)\) satisfying \(4p<y_k\)
and \(4p<y_j-y_i\) for \(i<j\), and real points \(r\in\mathbb R^{4n}\).
Set \(z_k=r^0_k+iy_k\). The heated vector is
\[\Psi_p(v;y;r)=e^{-(y_1-p)H}e^{ir^0_1H}U(\vec r_1)\,B_p(v_1)\,G_p(z_2-z_1,\vec r_2-\vec r_1)\,B_p(v_2)\cdots G_p(z_n-z_{n-1},\vec r_n-\vec r_{n-1})\,B_p(v_n)\Omega .\]
For \(n=0\) set \(\Psi_p=\Omega\). For a word
\(b=((u_1,f_1),\ldots,(u_n,f_n))\) and \(t>0\), the heat-cup vector is
\[\Phi_t(b)=\int f_1(x_1)\cdots f_n(x_n)\,\Psi_{p(t)}\big(u;(t,2t,\dots,nt);x\big)\,dx,\qquad p(t)=\min(t,1)/8 .\]

\begin{enumerate}
\def\labelenumi{(\alph{enumi})}
\item
  (C22, the pointwise Gram identity) Assume (H), \(J\) the induced star,
  the outputs of V.2.2(a)--(c), the orientation of \(P\), and \(T\)
  respects \(S\). Take \(x\in E_n\) and \(y\in E_m\) with positive
  entries, and widths \(t\), \(u\) admissible with all heights exceeding
  \(4t\) (resp. \(4u\)). Then
  \[\langle\Psi_t(v;x;r),\Psi_u(w;y;q)\rangle=F^{h}_{n+m}(\varrho)\,(Jv_n,\dots,Jv_1,w_1,\dots,w_m),\]
  with \(h=(-x_n,\dots,-x_1,y_1,\dots,y_m)\) and
  \(\varrho=(r_n,\dots,r_1,q_1,\dots,q_m)\).
\item
  (C23, integrability and translation covariance of heat cups) Assume
  (H), the slab-insertion properties V.2.2(a), and spin compatibility.
  For every word \(b\) and \(t>0\), the integrand of \(\Phi_t(b)\) is
  Bochner integrable, and \(\Phi_t(g\cdot b)=U(a)\Phi_t(b)\) for
  translations \(g=(a,1)\), where \(U(a)\) is the unitary translation of
  \(\mathcal H\).
\item
  (C24, the heat-cup Gram limit) Assume (H), (SC), (BV), integrability
  and translation covariance of heat cups and the pointwise Gram
  identity. Then
  \(\langle\Phi_s(b),\Phi_{s'}(c)\rangle\to\langle b,c\rangle_W\) as
  \((s,s')\to(0^+,0^+)\) independently.
\item
  (C25) For every \(n\ge0\), tensor products span a dense subspace of
  \(\mathcal S(\mathbb R^{4n})\).
\item
  (C26, label faithfulness) Assume C25, \(N_d(T)=N_d(S)\) for all \(d\),
  the hierarchies for \(S,T\), compatibility of the inclusions of \(P\)
  with both tables, (SC) and (TC). A label \(u\in\mathcal F\) with
  \(W(a\,(u,f)\,c)=0\) for all words \(a\), \(c\) and all \(f\) is zero.
  This concerns single labels, not words.
\item
  (C29, no hypotheses) For every boundary family \(W\), every antilinear
  map \(J\), all words \(a,c\) and all labels \(u,u'\), the following
  hold:
\end{enumerate}

\begin{itemize}
\tightlist
\item
  (word bilinearity) \((u,f)\mapsto W(a\,(u,f)\,c)\) is bilinear;
\item
  (continuity in the two tests)
  \((f,g)\mapsto\langle (u,f)a,(u',g)c\rangle_W\) is continuous on
  \(\mathcal S(\mathbb R^4)^2\).
\end{itemize}

\emph{Status:} theorem. \emph{Lean:}

\begin{itemize}
\tightlist
\item
  \href{https://github.com/yang-mills-group/OpenAI-YM-20261028/blob/3456e5e75cf8d4cd558955a5c7638fba2d68f464/YM/Work/C22.lean\#L348-L354}{\texttt{C22.answer}},
  \href{https://github.com/yang-mills-group/OpenAI-YM-20261028/blob/3456e5e75cf8d4cd558955a5c7638fba2d68f464/YM/Work/C23.lean\#L218-L220}{\texttt{C23.answer}},
  \href{https://github.com/yang-mills-group/OpenAI-YM-20261028/blob/3456e5e75cf8d4cd558955a5c7638fba2d68f464/YM/Work/C24.lean\#L146-L153}{\texttt{C24.answer}},
  \href{https://github.com/yang-mills-group/OpenAI-YM-20261028/blob/3456e5e75cf8d4cd558955a5c7638fba2d68f464/YM/Work/C25.lean\#L135-L138}{\texttt{C25.answer}},
  \href{https://github.com/yang-mills-group/OpenAI-YM-20261028/blob/3456e5e75cf8d4cd558955a5c7638fba2d68f464/YM/Work/C26.lean\#L77-L92}{\texttt{C26.answer}},
  \href{https://github.com/yang-mills-group/OpenAI-YM-20261028/blob/3456e5e75cf8d4cd558955a5c7638fba2d68f464/YM/Work/C29.lean\#L110}{\texttt{C29.answer}}.
\item
  Statements
  \href{https://github.com/yang-mills-group/OpenAI-YM-20261028/blob/3456e5e75cf8d4cd558955a5c7638fba2d68f464/YM/Bid/Grow.lean\#L29-L31}{\texttt{Bid.C22}},
  \href{https://github.com/yang-mills-group/OpenAI-YM-20261028/blob/3456e5e75cf8d4cd558955a5c7638fba2d68f464/YM/Bid/Grow.lean\#L33-L34}{\texttt{C23}},
  \href{https://github.com/yang-mills-group/OpenAI-YM-20261028/blob/3456e5e75cf8d4cd558955a5c7638fba2d68f464/YM/Bid/Grow.lean\#L36-L38}{\texttt{C24}},
  \href{https://github.com/yang-mills-group/OpenAI-YM-20261028/blob/3456e5e75cf8d4cd558955a5c7638fba2d68f464/YM/Bid/Grow.lean\#L42}{\texttt{C25}},
  \href{https://github.com/yang-mills-group/OpenAI-YM-20261028/blob/3456e5e75cf8d4cd558955a5c7638fba2d68f464/YM/Bid/Grow.lean\#L44-L46}{\texttt{C26}},
  \href{https://github.com/yang-mills-group/OpenAI-YM-20261028/blob/3456e5e75cf8d4cd558955a5c7638fba2d68f464/YM/Bid/Grow.lean\#L54}{\texttt{C29}}.
\item
  Definitions
  \href{https://github.com/yang-mills-group/OpenAI-YM-20261028/blob/3456e5e75cf8d4cd558955a5c7638fba2d68f464/YM/Bid/Grow.lean\#L19-L22}{\texttt{Hull}},
  \href{https://github.com/yang-mills-group/OpenAI-YM-20261028/blob/3456e5e75cf8d4cd558955a5c7638fba2d68f464/YM/Bid/Grow.lean\#L24-L27}{\texttt{Roof}},
  \href{https://github.com/yang-mills-group/OpenAI-YM-20261028/blob/3456e5e75cf8d4cd558955a5c7638fba2d68f464/YM/Fields/Tall.lean\#L53-L56}{\texttt{Small.shot}},
  \href{https://github.com/yang-mills-group/OpenAI-YM-20261028/blob/3456e5e75cf8d4cd558955a5c7638fba2d68f464/YM/Fields/Tall.lean\#L92-L93}{\texttt{Small.heatCup}},
  \href{https://github.com/yang-mills-group/OpenAI-YM-20261028/blob/3456e5e75cf8d4cd558955a5c7638fba2d68f464/YM/Fields/Tall.lean\#L96-L100}{\texttt{Small.Lamp}},
  \href{https://github.com/yang-mills-group/OpenAI-YM-20261028/blob/3456e5e75cf8d4cd558955a5c7638fba2d68f464/YM/Fields/Star.lean\#L97}{\texttt{Surface.kernel}},
  \href{https://github.com/yang-mills-group/OpenAI-YM-20261028/blob/3456e5e75cf8d4cd558955a5c7638fba2d68f464/YM/Fields/Stems.lean\#L38}{\texttt{Persists}},
  \href{https://github.com/yang-mills-group/OpenAI-YM-20261028/blob/3456e5e75cf8d4cd558955a5c7638fba2d68f464/YM/Fields/Stems.lean\#L26-L27}{\texttt{Symbol}},
  \href{https://github.com/yang-mills-group/OpenAI-YM-20261028/blob/3456e5e75cf8d4cd558955a5c7638fba2d68f464/YM/Fields/Stems.lean\#L29-L30}{\texttt{Cam}}.
\end{itemize}

\textbf{Proposition V.2.9 (the target: a Wightman structure realizing
\(W\)).}

A \emph{Wightman structure} consists of:

\begin{itemize}
\tightlist
\item
  families \(K(\Lambda)\) of linear automorphisms of \(\mathcal F\) and
  of each \(F_d\);
\item
  a field system: a Hilbert space \(\mathcal H_W\), a subspace
  \(\mathcal D\), a vacuum \(\Omega_W\in\mathcal D\), and fields
  \(\varphi(u,f)\in\mathrm{End}(\mathcal D)\), bilinear in
  \((u,f)\in\mathcal F\times\mathcal S(\mathbb R^4)\);
\item
  an antilinear map \(J:\mathcal F\to\mathcal F\);
\item
  a unitary, strongly continuous representation \(U\) of \(\mathcal P\);
\item
  a projection-valued measure \(\mathsf E\) on \(\mathbb R^4\).
\end{itemize}

It \emph{realizes} \(W\) if:

\begin{itemize}
\tightlist
\item
  \(J\) is the induced star and \(K\) satisfies spin compatibility;
\item
  totality and temperedness: \(\|\Omega_W\|=1\); \(\mathcal D\) is dense
  and equals the span of the word vectors \(\varphi(b)\Omega_W\);
  \(f\mapsto\varphi(u,f)\psi\) is continuous;
\item
  the adjoint identity:
  \(\langle\varphi(u,f)\psi,\chi\rangle=\langle\psi,\varphi(Ju,\bar f)\chi\rangle\)
  on \(\mathcal D\);
\item
  vacuum expectations:
  \(\langle\Omega_W,\varphi(b)\Omega_W\rangle=W(b)\);
\item
  local commutativity: \([\varphi(u,f),\varphi(u',g)]=0\) on
  \(\mathcal D\) when \(\mathrm{supp}f-\mathrm{supp}g\) is spacelike;
\item
  field faithfulness: for \(u\ne0\) some \(\varphi(u,f)\ne0\);
\item
  covariance: \(U(g)\) preserves \(\mathcal D\) and fixes \(\Omega_W\),
  and \(U(g)\varphi(u,f)\psi=\varphi(K(\Lambda)u,g_\ast f)U(g)\psi\);
\item
  \(\mathsf E\) is the spectral measure of the translations and vanishes
  off \(\bar V_+\);
\item
  every translation-invariant vector is a multiple of \(\Omega_W\).
\end{itemize}

The \emph{field axioms without vacuum uniqueness} (VIII.1) are the same
list without the last clause. \emph{Status:} definition (a
specification). It is produced in VIII.1 by the construction of V.2.10.
Consequence:
\href{https://github.com/yang-mills-group/OpenAI-YM-20261028/blob/3456e5e75cf8d4cd558955a5c7638fba2d68f464/YM/Fields/System.lean\#L62-L70}{\texttt{Fields.general}}
and
\href{https://github.com/yang-mills-group/OpenAI-YM-20261028/blob/3456e5e75cf8d4cd558955a5c7638fba2d68f464/YM/Fields/System.lean\#L82-L84}{\texttt{Fields.sees}}
show that a realization has Gram form
\(\langle\varphi(b)\Omega_W,\varphi(c)\Omega_W\rangle=\langle b,c\rangle_W\),
so \(W\) is positive on words. \emph{Lean:}
\href{https://github.com/yang-mills-group/OpenAI-YM-20261028/blob/3456e5e75cf8d4cd558955a5c7638fba2d68f464/YM/Fields/System.lean\#L117-L122}{\texttt{Wightman}},
\href{https://github.com/yang-mills-group/OpenAI-YM-20261028/blob/3456e5e75cf8d4cd558955a5c7638fba2d68f464/YM/Fields/System.lean\#L124-L135}{\texttt{Wightman.Realizes}},
\href{https://github.com/yang-mills-group/OpenAI-YM-20261028/blob/3456e5e75cf8d4cd558955a5c7638fba2d68f464/YM/Fields/System.lean\#L22-L26}{\texttt{Fields}},
\href{https://github.com/yang-mills-group/OpenAI-YM-20261028/blob/3456e5e75cf8d4cd558955a5c7638fba2d68f464/YM/Fields/System.lean\#L41-L45}{\texttt{Fields.Total}},
\href{https://github.com/yang-mills-group/OpenAI-YM-20261028/blob/3456e5e75cf8d4cd558955a5c7638fba2d68f464/YM/Fields/System.lean\#L111-L115}{\texttt{Symmetries}},
\href{https://github.com/yang-mills-group/OpenAI-YM-20261028/blob/3456e5e75cf8d4cd558955a5c7638fba2d68f464/YM/Fields/System.lean\#L92-L97}{\texttt{Action}},
\href{https://github.com/yang-mills-group/OpenAI-YM-20261028/blob/3456e5e75cf8d4cd558955a5c7638fba2d68f464/YM/Next/BOUNDARY/Core.lean\#L34-L44}{\texttt{FieldCore}}.

\textbf{Proposition V.2.10 (field realization of \(W\)).} Fix rotation
data \(R\) and a real structure on \(S\). Let \(J\) be the induced
involutive star and assume spin compatibility. Let \(W\) satisfy word
bilinearity, continuity in the two tests, word locality, label
faithfulness, continuity of word matrix elements and invariance of the
word Gram form. Suppose vectors \(\Phi_s(b)\in\mathcal H\), \(s>0\),
satisfy the Gram limit V.2.8(c), \(\Phi_s(\varnothing)\to\Omega\) as
\(s\downarrow0\), and, for every word \(b\) and translation \(a\),
\(\Phi_s(a\cdot b)=U(a)\Phi_s(b)\) for all sufficiently small \(s>0\).
Assume \(T\) satisfies the normalized, unital, permutation-symmetric
Euclidean laws. Then there is a Wightman structure satisfying every
clause of V.2.9 except possibly vacuum uniqueness. If the common fixed
vectors of all \(e^{-tH}\), \(t>0\), are multiples of \(\Omega\), it
realizes \(W\). \emph{Status:} theorem. The route used by the main
theorem is
\href{https://github.com/yang-mills-group/OpenAI-YM-20261028/blob/3456e5e75cf8d4cd558955a5c7638fba2d68f464/YM/Next/BOUNDARY/Stage.lean\#L167-L184}{\texttt{Next.BOUNDARY.establishing}}
followed by
\href{https://github.com/yang-mills-group/OpenAI-YM-20261028/blob/3456e5e75cf8d4cd558955a5c7638fba2d68f464/YM/Next/BOUNDARY/Stage.lean\#L217-L221}{\texttt{Compatible.vacuum}}
(VIII.1).
\href{https://github.com/yang-mills-group/OpenAI-YM-20261028/blob/3456e5e75cf8d4cd558955a5c7638fba2d68f464/YM/Fields/Lo.lean\#L85-L95}{\texttt{Boundary.harvest}}
states the same in one lemma and is used only by an alternative route.
The ingredients are in §4.8.

\emph{Not reconstruction.} Integrated.lean also proves three ultraviolet
bundles, and the folders Overlay and Vac are ultraviolet material; they
belong to Part III (III.6, III.10--III.13):

\begin{itemize}
\tightlist
\item
  \href{https://github.com/yang-mills-group/OpenAI-YM-20261028/blob/3456e5e75cf8d4cd558955a5c7638fba2d68f464/YM/Integrated.lean\#L216-L219}{\texttt{source\_sections}}
  (G01--G07);
\item
  \href{https://github.com/yang-mills-group/OpenAI-YM-20261028/blob/3456e5e75cf8d4cd558955a5c7638fba2d68f464/YM/Integrated.lean\#L222-L225}{\texttt{pressure\_grants}}
  (E01--E09, given
  \href{https://github.com/yang-mills-group/OpenAI-YM-20261028/blob/3456e5e75cf8d4cd558955a5c7638fba2d68f464/YM/Bid/Energy.lean\#L148}{\texttt{EBare}}).
  The energy hypothesis is proved by
  \href{https://github.com/yang-mills-group/OpenAI-YM-20261028/blob/3456e5e75cf8d4cd558955a5c7638fba2d68f464/YM/Next/FLUCTUATIONS/Bare.lean\#L64-L120}{\texttt{Next.Fluctuations.bare}};
\item
  \href{https://github.com/yang-mills-group/OpenAI-YM-20261028/blob/3456e5e75cf8d4cd558955a5c7638fba2d68f464/YM/Integrated.lean\#L229-L232}{\texttt{stationary\_scene}}
  (E02, F01--F04).
\end{itemize}

\subsubsection{2. What it uses}\label{what-it-uses-21}

\textbf{Inputs from other units.}

\begin{enumerate}
\def\labelenumi{\arabic{enumi}.}
\tightlist
\item
  \textbf{The Euclidean hierarchy} (IV.1; for each phase, V.1):
  \href{https://github.com/yang-mills-group/OpenAI-YM-20261028/blob/3456e5e75cf8d4cd558955a5c7638fba2d68f464/YM/Euclid/Words.lean\#L141-L148}{\texttt{Table.Hierarchy}}.
  It consists of permutation symmetry, normalization, the unit-insertion
  law, compatibility with signed symmetries and conjugation, translation
  homogeneity, reflection positivity on words and on smeared joint
  tests, smoothness off coincidences, and capped polynomial bounds near
  collisions. Its uses include:

  \begin{itemize}
  \tightlist
  \item
    permutation symmetry of \(S\) and \(T\) (C02, C20, C22, C26);
  \item
    smoothness of \(T\) off coincidences (C07, C20, C26);
  \item
    translation homogeneity of \(T\) (C02, C09);
  \item
    the reflection adjoint on labels (the star \(J\));
  \item
    the null spaces (null-space compatibility, C26).
  \end{itemize}
\item
  \textbf{The OS space and the slab operators} (V.1, the B-items; files
  YM/OS, YM/Heat).

  \begin{itemize}
  \tightlist
  \item
    \href{https://github.com/yang-mills-group/OpenAI-YM-20261028/blob/3456e5e75cf8d4cd558955a5c7638fba2d68f464/YM/OS/Dynamics.lean\#L57-L61}{\texttt{OS.Chamber}}
    provides:

    \begin{itemize}
    \tightlist
    \item
      the GNS space with vacuum \(\Omega\);
    \item
      the semigroup \(e^{-tH}\) and its holomorphic extension
      \(e^{-\zeta H}\), a contraction for \(\mathrm{Re}\,\zeta>0\), with
      the semigroup law and \(e^{-tH}\Omega=\Omega\) (Heat/Rows,
      Heat/Extend);
    \item
      spatial translations;
    \item
      a unitary space-time translation group whose spectral measure lies
      in \(\{p^0\ge0\}\);
    \item
      the ground space, i.e.~the invariant vectors.
    \end{itemize}
  \item
    \href{https://github.com/yang-mills-group/OpenAI-YM-20261028/blob/3456e5e75cf8d4cd558955a5c7638fba2d68f464/YM/OS/Slabs.lean\#L101-L104}{\texttt{OS.Spread}}:
    slab operators \(A_d(w)\) acting on GNS vectors by concatenation of
    words.
  \item
    The laws
    \href{https://github.com/yang-mills-group/OpenAI-YM-20261028/blob/3456e5e75cf8d4cd558955a5c7638fba2d68f464/YM/OS/Slabs.lean\#L153-L162}{\texttt{Spread.Laws}}.
    Among them is the slab bound (SB),
    \href{https://github.com/yang-mills-group/OpenAI-YM-20261028/blob/3456e5e75cf8d4cd558955a5c7638fba2d68f464/YM/OS/Slabs.lean\#L111-L113}{\texttt{AllCopy}}:
    for each degree bound \(D\) and representative bound \(R\) there are
    \(C>0\) and \(a\in\mathbb N\) such that
    \[\|A_d(w)\|\le C^{\lvert w\rvert}\big(\lvert w\rvert/r\big)^{a\lvert w\rvert}\]
    for every word \(w\) whose labels have degree at most \(D\) and
    representative norm at most \(R\), in a slab of width \(d\), with
    pairwise distances in the maximum norm and distances from both time
    boundaries at least \(r\in(0,1]\). This is the only quantitative
    input of the unit; it gives the bound \(p^{-a}\) of C01.
  \item
    VIII.1 obtains (H) from the phase data, both for the mixed table
    \(S=T\) and for the selected phase.
  \end{itemize}
\item
  \textbf{Rotation, spin and star data} (VI.1).

  \begin{itemize}
  \tightlist
  \item
    \(R\) with Euclidean covariance and rotation laws (A10--A15).
  \item
    \(K\) with spin compatibility (A16--A17).
  \item
    A16 itself, whose proof
    \href{https://github.com/yang-mills-group/OpenAI-YM-20261028/blob/3456e5e75cf8d4cd558955a5c7638fba2d68f464/YM/Work/A16.lean}{\texttt{Work/A16}}
    is a field of the bundle.
  \item
    The induced star \(J\) (the induced-star identity, from the
    orientation of the differential plan).
  \end{itemize}
\item
  \textbf{The label structures} of I.1 and II.1: the gauge group, the
  register, the quotient labels, \(\mathrm{pin}\), and the common level
  of finitely many labels.
\end{enumerate}

\textbf{Standard mathematics.}

\begin{itemize}
\tightlist
\item
  \emph{Osterwalder--Schrader reconstruction:} Osterwalder--Schrader,
  \emph{Axioms for Euclidean Green's functions} I, II, CMP 31 (1973), 42
  (1975); Glimm--Jaffe, \emph{Quantum Physics}, Ch. 6 and 19.

  \begin{itemize}
  \tightlist
  \item
    Here the continuation is built from operators on the OS space, so
    the Laplace-transform estimates of OS II are not needed.
  \end{itemize}
\item
  \emph{Boundary values of holomorphic functions of polynomial growth as
  distributions of finite order:} Hörmander, \emph{ALPDO I}, Thm.
  3.1.15; Vladimirov, \emph{Methods of the Theory of Generalized
  Functions}; Streater--Wightman, Thm. 2-10.
\item
  \emph{Fourier--Laplace transforms of distributions with support in a
  cone, and holomorphy in tubes:} Reed--Simon II, Thm. IX.16;
  Streater--Wightman, Thm. 2-6.
\item
  \emph{Locality:} Jost points and the complex Lorentz group (Jost,
  \emph{The General Theory of Quantized Fields}; Streater--Wightman,
  Thms. 2-11, 2-12). Locality from the permutation symmetry of the
  Schwinger functions follows OS I.

  \begin{itemize}
  \tightlist
  \item
    A holomorphic germ is determined by its restriction to a maximal
    totally real subspace (elementary, by power series).
  \end{itemize}
\item
  \emph{Hilbert-space tools:} the Kolmogorov (GNS) construction;
  Riesz--Markov; Stone--Weierstrass; the spectral measure of a
  continuous unitary representation of \(\mathbb R^4\)
  (Stone--Naimark--Ambrose--Godement); Reed--Simon I, Ch. VII--VIII.
\item
  \emph{Other standard facts:}

  \begin{itemize}
  \tightlist
  \item
    density of tensor products in \(\mathcal S(\mathbb R^{4n})\)
    (Reed--Simon I, §V.3);
  \item
    continuous one-parameter groups in finite dimension are
    exponentials;
  \item
    generation of \(\mathcal L^\uparrow_+\) by rotations and boosts.
  \end{itemize}
\end{itemize}

\textbf{Imported constants and assumptions.}

\begin{itemize}
\tightlist
\item
  There are no numerical constants beyond conventions: the width
  \(p<1\), gaps \(>4p\), and the choice \(p=\min(k,1)/8\). The exponents
  \(a\) and constants \(C\) of C01 are those of (SB) (V.1).
\item
  The hypotheses not discharged here are:

  \begin{itemize}
  \tightlist
  \item
    (H);
  \item
    the data \(R\), \(K\), \(J\) with their properties;
  \item
    \(N_d(T)=N_d(S)\) for every \(d\);
  \item
    one-dimensionality of the heat-fixed space.
  \end{itemize}

  VIII.1 supplies the table, slab and symmetry data from V.1 and VI.1.
  One-dimensionality of the heat-fixed space is supplied for the
  selected phase.
\item
  The boundary family \(W\), (BV) and cone support are produced by
  V.2.3(b) and V.2.4(c) from the outputs of C03, C04, C06 and C09.
\end{itemize}

\subsubsection{3. Strategy}\label{strategy-21}

The slab bound controls a field insertion \(B_p(v)\) between two heat
factors. The relation \(B_{p+u}=e^{-uH}B_pe^{-uH}\) makes products of
these insertions independent of the admissible width. Their vacuum
matrix elements give a continuation to strictly ordered imaginary times,
with polynomial growth as the time differences shrink. Cauchy--Riemann
identities and Taylor's formula along \((0,1,\ldots,n-1)\) then give
tempered boundary values for arbitrary approaches through ordered
heights. A Laplace identity yields total energy zero and nonnegative
suffix energies.

Euclidean covariance gives Poincaré covariance of the boundary family:
translations and spatial rotations pass through a one-variable identity
theorem; infinitesimal rotations in a time--space plane give the boost
identity, which is integrated along the real boost groups. Lorentz
covariance improves the spectral support to \(\mathcal V_n\).
Fourier--Laplace transformation extends the continuation to the forward
tube. Locality follows from two holomorphic germs near each spacelike
adjacent pair: imaginary boosts construct the germs, Euclidean
permutation symmetry identifies their jets on the Wick slice, and
continuation along a rescaling of real times propagates their agreement.
Partial boundary values then give exchange of adjacent spacelike fields.

Finally, inner products of heated OS vectors equal the continuation at a
glued, reflected configuration. Their limit gives positivity of the word
kernel and an isometric embedding of its Kolmogorov space into the OS
Hilbert space. Left multiplication defines the fields on the common word
domain. The embedding intertwines translations, transferring positive
energy; Lorentz covariance gives the forward-cone spectrum condition. A
one-dimensional OS heat-fixed space gives vacuum uniqueness. These are
the operator and local-germ forms of Osterwalder--Schrader
reconstruction used here.

\subsubsection{4. Steps}\label{steps-21}

\paragraph{4.1 Tools}\label{tools}

\textbf{Lemma V.2.a (Kolmogorov).} Let \(K\) be a positive semidefinite
kernel on a set \(X\). There are a Hilbert space and vectors
\((\psi_x)_{x\in X}\) with \(\langle\psi_x,\psi_y\rangle=K(x,y)\) and
dense span. \emph{Proof sketch.} Complete the finitely supported
functions modulo the null space of \(\sum\bar f(x)g(y)K(x,y)\).
\emph{Lean:}
\href{https://github.com/yang-mills-group/OpenAI-YM-20261028/blob/3456e5e75cf8d4cd558955a5c7638fba2d68f464/YM/Hilbert/Kolmogorov.lean\#L67-L92}{\texttt{Kernel.host}},
\href{https://github.com/yang-mills-group/OpenAI-YM-20261028/blob/3456e5e75cf8d4cd558955a5c7638fba2d68f464/YM/Hilbert/Kolmogorov.lean\#L34-L38}{\texttt{Kernel.Gram}}.

\textbf{Lemma V.2.b (projection-valued measure).} Let \(X\) be a compact
Hausdorff space and \(A:C(X)\to\mathcal B(\mathcal H)\) a unital complex
\(\ast\)-algebra homomorphism. There is a strongly countably additive
projection-valued measure on the Borel sets of \(X\) representing \(A\).
\emph{Proof sketch.}

\begin{enumerate}
\def\labelenumi{\arabic{enumi}.}
\tightlist
\item
  For each \(\psi\), the positive functional
  \(f\mapsto\langle\psi,A(f)\psi\rangle\) gives a Riesz--Markov measure
  \(\mu_\psi\).
\item
  \(f\mapsto A(f)\psi\) extends to an isometry
  \(L^2(\mu_\psi)\to\mathcal H\).
\item
  \(E(s)\psi\) is the image of \(1_s\); for disjoint Borel sets, the
  corresponding indicator sums converge in \(L^2(\mu_\psi)\), giving
  strong countable additivity.
\end{enumerate}

It is used (V.1, SNAG/Escape) on the one-point compactification of
\(\mathbb R^4\) to build the spectral measure of the translations.
\emph{Lean:}
\href{https://github.com/yang-mills-group/OpenAI-YM-20261028/blob/3456e5e75cf8d4cd558955a5c7638fba2d68f464/YM/Compact/PVM.lean\#L51-L62}{\texttt{Spectral.Rep.ships}}
\emph{(proof read for additivity; Compact/Borel statement only)}.

\textbf{Lemma V.2.c (Schwartz calculus; tensor density, C25).}

\begin{enumerate}
\def\labelenumi{\arabic{enumi}.}
\tightlist
\item
  Let \(f_\lambda,g\in\mathcal S(\mathbb R^d)\). If every jet of
  \(f_\lambda\) converges locally uniformly to the corresponding jet of
  \(g\), and every Schwartz seminorm of \(f_\lambda\) is eventually
  bounded uniformly in \(\lambda\), then \(f_\lambda\to g\) in
  \(\mathcal S\).
\item
  Let \(K\) be compact metric, \(\mu\) a finite Borel measure, and
  \(u\mapsto f_u\) continuous from \(K\) into
  \(\mathcal S(\mathbb R^d)\). If \(g\in\mathcal S\) satisfies
  \(g(x)=\int_K f_u(x)\,d\mu(u)\), then the sums
  \(\sum_j\mu(E_j)f_{u_j}\) over finite measurable partitions, with
  \(u_j\in E_j\) and \(\max_j\mathrm{diam}(E_j)\to0\), converge to \(g\)
  in \(\mathcal S\).
\item
  The linear span of tensor products is dense in
  \(\mathcal S(\mathbb R^{4n})\).
\end{enumerate}

\emph{Proof sketch.} For 3:

\begin{itemize}
\tightlist
\item
  reduce to compactly supported \(v\);
\item
  choose a product cutoff equal to \(1\) on the support;
\item
  replace \(\hat v\) by compactly supported cutoffs
  \(\hat v_N\to\hat v\) in \(\mathcal S\);
\item
  the fixed product cutoff times
  \((2\pi)^{-4n}\int\hat v_N(\xi)e^{i\langle\xi,\cdot\rangle}\,d\xi\)
  converges to \(v\) by Fourier inversion;
\item
  approximate each compact integral by quadrature sums of products of
  cutoff plane waves, using 2. \emph{Lean:}
  \href{https://github.com/yang-mills-group/OpenAI-YM-20261028/blob/3456e5e75cf8d4cd558955a5c7638fba2d68f464/YM/Spark/Limit.lean\#L28-L51}{\texttt{Spark.ready}},
  \href{https://github.com/yang-mills-group/OpenAI-YM-20261028/blob/3456e5e75cf8d4cd558955a5c7638fba2d68f464/YM/Spark/Pile.lean\#L77-L114}{\texttt{Spark.records}};
  \href{https://github.com/yang-mills-group/OpenAI-YM-20261028/blob/3456e5e75cf8d4cd558955a5c7638fba2d68f464/YM/Work/C25.lean\#L109-L124}{\texttt{C25.compact\_mem}},
  \href{https://github.com/yang-mills-group/OpenAI-YM-20261028/blob/3456e5e75cf8d4cd558955a5c7638fba2d68f464/YM/Work/C25.lean\#L85-L107}{\texttt{C25.cut\_port\_mem}}.
\end{itemize}

\paragraph{4.2 The strict-time
continuation}\label{the-strict-time-continuation}

\textbf{Lemma V.2.d (slab insertions; C01).} As in V.2.2(a). \emph{Proof
sketch.}

\begin{itemize}
\tightlist
\item
  \textbf{Construction.} For each degree, \(B_p\) is defined on
  representatives by the one-point slab operator.
\item
  \textbf{Linearity and descent to quotient labels.} An operator on
  \(\mathcal H\) is determined by its matrix elements between GNS
  vectors. These matrix elements are Schwinger functions of \(T\) on the
  word \emph{reflected word, centre insertion, raised word}; they are
  linear in the inserted label and vanish on \(S\)-null labels because
  \(T\) respects \(S\). The degrees are glued using compatibility with
  the inclusions. For \(p=0\) the map is set to \(0\).
\item
  \textbf{Bound.} (SB) applies to a word of length 1 in a slab of width
  \(2p\) with buffer \(r=p\le1\) and gives
  \(\|B_p(\mathrm{pin}_dw)\|\le C(1/p)^a\). The constants depend on the
  degree and on the norm of the representative.
\end{itemize}

\emph{Lean:}
\href{https://github.com/yang-mills-group/OpenAI-YM-20261028/blob/3456e5e75cf8d4cd558955a5c7638fba2d68f464/YM/Work/C01.lean\#L71-L80}{\texttt{C01.single}},
\href{https://github.com/yang-mills-group/OpenAI-YM-20261028/blob/3456e5e75cf8d4cd558955a5c7638fba2d68f464/YM/Work/C01.lean\#L51-L57}{\texttt{C01.inner\_raw}},
\href{https://github.com/yang-mills-group/OpenAI-YM-20261028/blob/3456e5e75cf8d4cd558955a5c7638fba2d68f464/YM/Work/C01.lean\#L91-L97}{\texttt{C01.at\_width}},
\href{https://github.com/yang-mills-group/OpenAI-YM-20261028/blob/3456e5e75cf8d4cd558955a5c7638fba2d68f464/YM/Work/C01.lean\#L99-L117}{\texttt{C01.answer}}.

\textbf{Lemma V.2.e (Euclidean values and width independence; C02).} As
in V.2.2(b). \emph{Proof sketch.}

\begin{itemize}
\tightlist
\item
  \textbf{Euclidean values.} At imaginary times each factor \(G_p\) is
  the heat semigroup at a positive real time. By induction from the
  right, the vector \(B_p(v_k)G_p\cdots B_p(v_n)\Omega\) is the GNS
  vector of a positive-time Euclidean word. At the end,
  \(\langle\Omega,\text{vector}\rangle\) is the Schwinger function, and
  translation invariance removes the offset.
\item
  \textbf{Width independence.}

  \begin{itemize}
  \tightlist
  \item
    First, \(B_{p+u}=e^{-uH}B_pe^{-uH}\).
  \item
    Second, \(G_p(w)=e^{-uH}G_{p+u}(w)e^{-uH}\).
  \item
    The factors \(e^{-uH}\) telescope, and \(e^{-uH}\Omega=\Omega\)
    removes the last one.
  \end{itemize}
\end{itemize}

Only gaps \(>2p\) are needed; the factor \(4\) in ``admissible'' is
slack. \emph{Lean:}
\href{https://github.com/yang-mills-group/OpenAI-YM-20261028/blob/3456e5e75cf8d4cd558955a5c7638fba2d68f464/YM/Work/C02.lean\#L68-L97}{\texttt{C02.train\_pure}},
\href{https://github.com/yang-mills-group/OpenAI-YM-20261028/blob/3456e5e75cf8d4cd558955a5c7638fba2d68f464/YM/Work/C02.lean\#L118-L128}{\texttt{C02.pad\_bar}},
\href{https://github.com/yang-mills-group/OpenAI-YM-20261028/blob/3456e5e75cf8d4cd558955a5c7638fba2d68f464/YM/Work/C02.lean\#L156-L166}{\texttt{C02.matrix\_pad}}.

\textbf{Lemma V.2.f (holomorphy and polynomial growth; C03).} As in
V.2.2(c). \emph{Proof sketch.}

\begin{itemize}
\tightlist
\item
  \textbf{Definition.} \(F_n(z)=M_p(z)\) with \(p=\min(k,1)/8\), where
  \(k>0\) is any lower bound for the imaginary-time differences when
  \(n\ge2\); take \(k=1\) for \(n=0,1\). By V.2.e this does not depend
  on the admissible \(p\), and admissibility is an open condition.
\item
  \textbf{Holomorphy and continuity.} These follow from those of
  \(\zeta\mapsto e^{-\zeta H}\) in norm, combined with strong continuity
  of \(U\), through the finite product.
\item
  \textbf{Growth.} For \(0<k\le1\), take \(p=k/8\). The factors \(G_p\)
  are contractions, so
  \(\lvert F\rvert\le\|\Omega\|^2\prod_j\|B_p(v_j)\|\le\|\Omega\|^2\big(\prod_jC_j8^{a_j}\big)k^{-\sum_ja_j}\).
  This gives (PG) with \(m=\sum a_k\).
\end{itemize}

\emph{Lean:}
\href{https://github.com/yang-mills-group/OpenAI-YM-20261028/blob/3456e5e75cf8d4cd558955a5c7638fba2d68f464/YM/Work/C03.lean\#L218-L226}{\texttt{C03.defined\_exists}},
\href{https://github.com/yang-mills-group/OpenAI-YM-20261028/blob/3456e5e75cf8d4cd558955a5c7638fba2d68f464/YM/Work/C03.lean\#L228-L243}{\texttt{C03.defined\_region}},
\href{https://github.com/yang-mills-group/OpenAI-YM-20261028/blob/3456e5e75cf8d4cd558955a5c7638fba2d68f464/YM/Work/C03.lean\#L200-L216}{\texttt{C03.defined\_power}}.

\paragraph{4.3 Boundary values and spectral
support}\label{boundary-values-and-spectral-support}

\textbf{Lemma V.2.g (C04).} As in V.2.3(a). \emph{Proof sketch.}

\begin{itemize}
\tightlist
\item
  (PG) gives a bound on \(F^y\) that is uniform in \(x\) and locally
  uniform in \(y\in E_n\); this yields integrability and continuity by
  dominated convergence.
\item
  On each one-variable slice \(w\mapsto F(\cdot+w\underline a)\),
  Cauchy--Riemann gives
  \(\partial_tF(y+ta)=i\,\partial_{\underline a}F\).
\item
  Cauchy estimates on discs of radius uniform in \(x\) allow
  differentiation under the integral.
\item
  Integration by parts moves \(\partial_{\underline a}\) onto \(f\).
\end{itemize}

\emph{Lean:}
\href{https://github.com/yang-mills-group/OpenAI-YM-20261028/blob/3456e5e75cf8d4cd558955a5c7638fba2d68f464/YM/Work/C04.lean\#L131-L144}{\texttt{C04.cr\_height}},
\href{https://github.com/yang-mills-group/OpenAI-YM-20261028/blob/3456e5e75cf8d4cd558955a5c7638fba2d68f464/YM/Work/C04.lean\#L146-L171}{\texttt{C04.deriv\_bound\_near}},
\href{https://github.com/yang-mills-group/OpenAI-YM-20261028/blob/3456e5e75cf8d4cd558955a5c7638fba2d68f464/YM/Work/C04.lean\#L173-L219}{\texttt{C04.lines}}.

\textbf{Lemma V.2.h (boundary values).} As in V.2.3(b). For every
nondecreasing \(y\), including \(y=0\), the boundary value is given by
an explicit formula. Let \(\varrho=(0,1,\dots,n-1)\) and
\(D=-i\partial_{\underline\varrho}\), and let \(m\) be the exponent of
(PG). Then
\[W^y_n(v)(f)=\sum_{k=0}^m\frac{(-1)^k}{k!}\,\psi_v(y+\varrho;D^kf)-\int_0^1\frac{(-s)^m}{m!}\,\psi_v(y+s\varrho;D^{m+1}f)\,ds .\]
\emph{Proof sketch.}

\begin{itemize}
\tightlist
\item
  For \(y\in E_n\), Taylor's formula with integral remainder along
  \(t\mapsto y+t\varrho\) (by (CR)) shows \(W^y=\psi(y;\cdot)\).
\item
  For nondecreasing \(y\) the profile \(y+s\varrho\) has all gaps
  \(\ge s\). So (PG) bounds the remainder integrand by
  \(C\|D^{m+1}f\|_{L^1}/m!\), uniformly in \(y\).
\item
  Hence:

  \begin{itemize}
  \tightlist
  \item
    \(y\mapsto W^y(f)\) is continuous on the closed set of nondecreasing
    heights, by dominated convergence;
  \item
    \(\lvert W^y(f)\rvert\) is bounded by finitely many \(L^1\) norms of
    \(D^kf\), so \(W^y\) is tempered.
  \end{itemize}
\item
  Multilinearity in \(v\) follows from uniqueness of limits. At \(y=0\)
  this gives (BV).
\end{itemize}

\emph{Lean:}
\href{https://github.com/yang-mills-group/OpenAI-YM-20261028/blob/3456e5e75cf8d4cd558955a5c7638fba2d68f464/YM/Shift/Breath.lean\#L27-L28}{\texttt{Waves.desired}},
\href{https://github.com/yang-mills-group/OpenAI-YM-20261028/blob/3456e5e75cf8d4cd558955a5c7638fba2d68f464/YM/Shift/Breath.lean\#L50-L66}{\texttt{Waves.meant}},
\href{https://github.com/yang-mills-group/OpenAI-YM-20261028/blob/3456e5e75cf8d4cd558955a5c7638fba2d68f464/YM/Shift/Breath.lean\#L78-L101}{\texttt{Waves.maintains}},
\href{https://github.com/yang-mills-group/OpenAI-YM-20261028/blob/3456e5e75cf8d4cd558955a5c7638fba2d68f464/YM/Shift/Osc.lean\#L27-L36}{\texttt{Waves.cost}},
\href{https://github.com/yang-mills-group/OpenAI-YM-20261028/blob/3456e5e75cf8d4cd558955a5c7638fba2d68f464/YM/Shift/Osc.lean\#L44-L80}{\texttt{Waves.seeks}}.

\textbf{Lemma V.2.i (Laplace identity; C05).} As in V.2.4(a).
\emph{Proof sketch.} Put
\(A(s)=\psi_v\big(sy;\widehat{e^{s\langle\underline y,\cdot\rangle}f}\big)\)
for \(s>0\).

\begin{enumerate}
\def\labelenumi{\arabic{enumi}.}
\tightlist
\item
  \textbf{\(A\) is constant.} By (CR) along \(y\), and since
  \(\frac{d}{ds}\widehat{e^{s\langle\underline y,\cdot\rangle}f}=i\,\partial_{\underline y}\widehat{e^{s\langle\underline y,\cdot\rangle}f}\)
  in \(\mathcal S\) (here compact support is needed), \(A'(s)=0\). The
  product rule is justified by the uniform bound
  \(\lvert\psi_{y'}(g)\rvert\le\mathrm{toll}(g)\) of V.2.h.
\item
  \textbf{The limit.} As \(s\to0^+\), \(A(s)\to W_n(v)(\hat f)\), by the
  radial limit and the same bound.
\end{enumerate}

Only the radial limit is used. \emph{Lean:}
\href{https://github.com/yang-mills-group/OpenAI-YM-20261028/blob/3456e5e75cf8d4cd558955a5c7638fba2d68f464/YM/Work/C05.lean\#L182-L201}{\texttt{C05.hasDerivAt\_zero\_psi}},
\href{https://github.com/yang-mills-group/OpenAI-YM-20261028/blob/3456e5e75cf8d4cd558955a5c7638fba2d68f464/YM/Work/C05.lean\#L157-L177}{\texttt{C05.tendsto\_psi}}.

\textbf{Lemma V.2.j (support in \(\mathcal P_n\); C06).} As in V.2.4(b).
\emph{Proof sketch.} Let \(y\) be nondecreasing and \(f\) compactly
supported in \(\{p:\langle\underline y,p\rangle<-1\}\).

\begin{enumerate}
\def\labelenumi{\arabic{enumi}.}
\tightlist
\item
  Apply V.2.i at the ordered heights \(Y(t)=ty+\varrho\), whose gaps are
  all \(\ge1\).
\item
  (PG) with \(k=1\) gives
  \(\lvert\psi_{Y(t)}(\hat g)\rvert\le C\|\hat g\|_{L^1}\).
\item
  Absorb \(e^{\langle\underline\varrho,\cdot\rangle}\) into the compact
  test. Every seminorm of the remaining factor
  \(e^{t\langle\underline y,\cdot\rangle}\) times that test is bounded
  by a polynomial in \(t\) times \(e^{-t}\), and tends to zero.
\item
  Hence \(W(\hat f)=0\). Rescaling \(y\) removes the \(-1\).
\end{enumerate}

Constant profiles of both signs force \(\sum_kp^0_k=0\). \emph{Lean:}
\href{https://github.com/yang-mills-group/OpenAI-YM-20261028/blob/3456e5e75cf8d4cd558955a5c7638fba2d68f464/YM/Work/C06.lean\#L101-L134}{\texttt{C06.patch}},
\href{https://github.com/yang-mills-group/OpenAI-YM-20261028/blob/3456e5e75cf8d4cd558955a5c7638fba2d68f464/YM/Work/C06.lean\#L38-L86}{\texttt{C06.lift\_falls}}.

\paragraph{4.4 Poincaré covariance}\label{poincaruxe9-covariance}

\textbf{Lemma V.2.k (Poincaré action on \(\mathcal S\); C08).} As in
V.2.5(b). \emph{Proof sketch.}

\begin{itemize}
\tightlist
\item
  \textbf{Continuity.} The seminorm estimate
  \(p_{i,k}(f\circ(A(\cdot-b)))\le(\|A^{-1}\|+\lvert b\rvert)^i\|A\|^k2^i\max_{(i',k')\le(i,k)}p_{i',k'}(f)\)
  combines with locally uniform convergence of jets (V.2.c, item 1).
\item
  \textbf{Derivative.} The pointwise derivative in \(t\) of
  \(\Lambda_a(t)_\ast f\) is \(-\Lambda_a(t)_\ast D_af\). The slope
  error is written as \(\int_0^1(\cdot)\,du\) (fundamental theorem of
  calculus) and controlled in \(\mathcal S\) by the continuity just
  proved.
\end{itemize}

\emph{Lean:}
\href{https://github.com/yang-mills-group/OpenAI-YM-20261028/blob/3456e5e75cf8d4cd558955a5c7638fba2d68f464/YM/Work/C08.lean\#L31-L59}{\texttt{C08.size\_affine\_le}},
\href{https://github.com/yang-mills-group/OpenAI-YM-20261028/blob/3456e5e75cf8d4cd558955a5c7638fba2d68f464/YM/Work/C08.lean\#L179-L185}{\texttt{C08.continuous\_gpush}},
\href{https://github.com/yang-mills-group/OpenAI-YM-20261028/blob/3456e5e75cf8d4cd558955a5c7638fba2d68f464/YM/Work/C08.lean\#L127-L161}{\texttt{C08.probe\_slope\_of\_pointwise}}.

\textbf{Lemma V.2.l (boost identity at strict heights; C07).} As in
V.2.5(c). \emph{Proof sketch.}

\begin{enumerate}
\def\labelenumi{\arabic{enumi}.}
\tightlist
\item
  \textbf{Euclidean input.} Differentiating the \(O(4)\) covariance of
  \(T\) along \(\exp(s\,e_0\wedge e_a)\) at a non-coincident
  configuration gives the rotation Ward identity: the derivative of
  \(T(x;v)\) along the rotation field equals
  \(-\sum_kT(x;\dots,\rho(e_0\wedge e_a)v_k,\dots)\).
\item
  \textbf{The analytic function \(E(z)\).} Fix \(f\) with compact
  support and introduce a complex parameter \(z\) and the times
  \(z\,x^0_k+iy_k\). Let \(E(z)\) be the integral of
  \[F\cdot\Big(\sum_k\big[x^a_k\partial_{x^0_k}+z(zx^0_k+iy_k)\partial_{x^a_k}\big]f\Big)-z\sum_kF(v^{(k)})f .\]
  The times stay strictly ordered on \(\mathrm{supp}f\) for \(z\) in a
  convex open set containing \(\mathbb R\), and \(E\) is holomorphic
  there.
\item
  \textbf{\(E\) vanishes on a segment of the imaginary axis.} At
  \(z=iw\), \(w>0\) small, the configuration is Euclidean, with times
  \(y_k+wx^0_k\). There the vector field is \(w\) times the pull-back of
  the \((0,a)\)-rotation field, and \(F\) is the Schwinger function. By
  step 1 the integrand is the divergence of the compactly supported
  vector field \(fT\) times the divergence-free rotation field, so
  \(E(iw)=0\).
\item
  \textbf{Conclusion.} The identity theorem gives \(E(1)=0\), which is
  (B\(_y\)) for compactly supported \(f\). Density (cutoffs) and the
  continuity of \(f\mapsto\psi(y;f)\) (V.2.h at fixed \(y\)) give all
  \(f\).
\end{enumerate}

The term \(iy_k\partial_{x^a_k}\) is the imaginary part of the complex
time multiplying the spatial derivative. \emph{Lean:}
\href{https://github.com/yang-mills-group/OpenAI-YM-20261028/blob/3456e5e75cf8d4cd558955a5c7638fba2d68f464/YM/Work/C07/Euclid.lean\#L104-L125}{\texttt{C07.oriented\_differential}},
\href{https://github.com/yang-mills-group/OpenAI-YM-20261028/blob/3456e5e75cf8d4cd558955a5c7638fba2d68f464/YM/Work/C07/Integrate.lean\#L47-L75}{\texttt{C07.integral\_orbit}},
\href{https://github.com/yang-mills-group/OpenAI-YM-20261028/blob/3456e5e75cf8d4cd558955a5c7638fba2d68f464/YM/Work/C07/Integrate.lean\#L126-L162}{\texttt{C07.error\_imag}},
\href{https://github.com/yang-mills-group/OpenAI-YM-20261028/blob/3456e5e75cf8d4cd558955a5c7638fba2d68f464/YM/Work/C07/Continue.lean\#L128-L148}{\texttt{C07.error\_one}},
\href{https://github.com/yang-mills-group/OpenAI-YM-20261028/blob/3456e5e75cf8d4cd558955a5c7638fba2d68f464/YM/Work/C07/Continue.lean\#L150-L170}{\texttt{C07.compact\_boost}},
\href{https://github.com/yang-mills-group/OpenAI-YM-20261028/blob/3456e5e75cf8d4cd558955a5c7638fba2d68f464/YM/Work/C07.lean\#L53-L66}{\texttt{C07.answer}}.

\textbf{Lemma V.2.m (real Poincaré covariance; C09, A16).} As in
V.2.5(a), (d). \emph{Proof sketch.}

\begin{itemize}
\tightlist
\item
  \textbf{Translations by \(a=(\vec a,a^0)\).} For fixed \(y\in E_n\),
  \(F^y_n(x+a)=F^y_n(x)\). This follows from a one-variable identity
  theorem in \(c\) for the times
  \(i\,\mathrm{Im}\,z_k+c(\mathrm{Re}\,z_k+a^0)\): at \(c=ib\) the times
  are imaginary, and Euclidean translation invariance by
  \((\vec a,ba^0)\) applies. Change variables and let \(y\to0\) (BV).
\item
  \textbf{Spatial rotations.} The same argument applies, using
  \(O(3)\subset O(4)\) covariance and \(K=R\) on rotations.
\item
  \textbf{Boosts.}

  \begin{itemize}
  \tightlist
  \item
    Passing (B\(_y\)) to \(y\to0\) removes the term
    \(\sum iy_k\partial_{x^a_k}\) and gives
    \(\sum_kW(v^{(k)})(f)=W(v)(D_af)\).
  \item
    The function \(h(t)=W_n(K(\Lambda_a(t))v)(\Lambda_a(t)_\ast f)\) has
    derivative \(0\) at \(t=0\), by C08(b), the boost derivative, and
    finite-dimensional expansion of the labels.
  \item
    By the group law \(h'\) vanishes everywhere, so \(h\) is constant.
  \end{itemize}
\item
  \textbf{The Lorentz group.} The \(\Lambda\) for which covariance holds
  form a subgroup containing the rotations and the coordinate boosts.
  A16 makes it all of \(\mathcal L^\uparrow_+\). The proof of A16 is
  purely group-theoretic: up to three boosts bring a future unit vector
  to rest, and its stabilizer consists of rotations.
\end{itemize}

\emph{Lean:}
\href{https://github.com/yang-mills-group/OpenAI-YM-20261028/blob/3456e5e75cf8d4cd558955a5c7638fba2d68f464/YM/Work/C09.lean\#L33-L115}{\texttt{C09.time\_compare}},
\href{https://github.com/yang-mills-group/OpenAI-YM-20261028/blob/3456e5e75cf8d4cd558955a5c7638fba2d68f464/YM/Work/C09.lean\#L123-L148}{\texttt{C09.restricted\_translation}},
\href{https://github.com/yang-mills-group/OpenAI-YM-20261028/blob/3456e5e75cf8d4cd558955a5c7638fba2d68f464/YM/Work/C09.lean\#L168-L199}{\texttt{C09.restricted\_spatial}},
\href{https://github.com/yang-mills-group/OpenAI-YM-20261028/blob/3456e5e75cf8d4cd558955a5c7638fba2d68f464/YM/Work/C09.lean\#L242-L279}{\texttt{C09.boost\_corner}},
\href{https://github.com/yang-mills-group/OpenAI-YM-20261028/blob/3456e5e75cf8d4cd558955a5c7638fba2d68f464/YM/Work/C09.lean\#L374-L399}{\texttt{C09.boost\_invariant}},
\href{https://github.com/yang-mills-group/OpenAI-YM-20261028/blob/3456e5e75cf8d4cd558955a5c7638fba2d68f464/YM/Work/A16.lean\#L91-L101}{\texttt{A16.answer}}.

\textbf{Lemma V.2.n (word covariance and cone support; C28).} As in
V.2.5(e) and V.2.4(c). \emph{Proof sketch.}

\begin{itemize}
\tightlist
\item
  \textbf{Gram invariance.} The induced star commutes with the
  letterwise action by spin compatibility. Pushed tensor tests are
  \(g_\ast\) of the tensor test; then apply (Cov).
\item
  \textbf{Continuity.} Use C08(b, first part), continuity of
  \(\Lambda\mapsto K(\Lambda)\) on a finite-dimensional level, and
  continuity of \(W_n(v)\).
\item
  \textbf{Lorentz improvement of spectral support.} The Fourier
  transform intertwines \(\Lambda\) with \(\tilde\Lambda\), so
  covariance carries the support condition \(\mathcal P_n\) to
  \(\tilde\Lambda^{-1}\mathcal P_n\) for every \(\Lambda\). The
  properties of \(\mathcal V_n\) follow from half-space support applied
  to step profiles and to constant profiles.
\end{itemize}

\emph{Lean:}
\href{https://github.com/yang-mills-group/OpenAI-YM-20261028/blob/3456e5e75cf8d4cd558955a5c7638fba2d68f464/YM/Work/C28.lean\#L134-L143}{\texttt{C28.answer}},
\href{https://github.com/yang-mills-group/OpenAI-YM-20261028/blob/3456e5e75cf8d4cd558955a5c7638fba2d68f464/YM/Shift/Closing.lean\#L106-L121}{\texttt{tents}}.

\paragraph{4.5 Continuation to the full
tube}\label{continuation-to-the-full-tube}

\textbf{Lemma V.2.o (cone geometry, cutoff, damping; C10--C12).} As in
V.2.6(a)--(c). \emph{Proof sketch.}

\begin{itemize}
\tightlist
\item
  \textbf{C10.} Summation by parts gives
  \(\langle p,\eta\rangle=\sum_k\langle P_{k+1},\eta_{k+1}-\eta_k\rangle\),
  where \(P_j=\sum_{l\ge j}p_l\) and \(P_1=0\). Each term is \(\ge0\)
  because \(\bar V_+\) is self-dual for the Euclidean pairing. For the
  local bound: if \(B(\eta,r)\subset V^+_n\) and \(p\ne0\), then for
  \(u\in B(\eta,r/2)\) the point \(u-\tfrac r2p/\lvert p\rvert\) lies in
  \(\bar V^+_n\), so \(\langle p,u\rangle\ge\tfrac r2\lvert p\rvert\).
\item
  \textbf{C11.} Take \(\chi_n\) = (a bump of integral \(1\) and outer
  radius \(1/2\)) \(\ast\) (indicator of the open \(2\)-neighbourhood of
  \(\mathcal V_n\)).
\item
  \textbf{C12.} By C10, \(-\langle\eta,p\rangle\le3\lvert\eta\rvert\) on
  the \(3\)-neighbourhood. Leibniz bounds then hold uniformly for
  bounded \(\eta\in\bar V^+_n\), and continuity into \(\mathcal S\)
  follows by V.2.c, item 1.
\end{itemize}

\emph{Lean:}
\href{https://github.com/yang-mills-group/OpenAI-YM-20261028/blob/3456e5e75cf8d4cd558955a5c7638fba2d68f464/YM/Work/C10.lean\#L25-L35}{\texttt{C10.telescope}},
\href{https://github.com/yang-mills-group/OpenAI-YM-20261028/blob/3456e5e75cf8d4cd558955a5c7638fba2d68f464/YM/Work/C10.lean\#L48-L69}{\texttt{C10.cone\_bound}},
\href{https://github.com/yang-mills-group/OpenAI-YM-20261028/blob/3456e5e75cf8d4cd558955a5c7638fba2d68f464/YM/Work/C11.lean\#L57-L100}{\texttt{C11.answer}},
\href{https://github.com/yang-mills-group/OpenAI-YM-20261028/blob/3456e5e75cf8d4cd558955a5c7638fba2d68f464/YM/Work/C12.lean\#L68-L76}{\texttt{C12.cone\_estimate}},
\href{https://github.com/yang-mills-group/OpenAI-YM-20261028/blob/3456e5e75cf8d4cd558955a5c7638fba2d68f464/YM/Work/C12.lean\#L84-L108}{\texttt{C12.attenuate\_estimate}}.

\textbf{Lemma V.2.p (holomorphy in the tube; C13).} As in V.2.6(d).
\emph{Proof sketch.}

\begin{enumerate}
\def\labelenumi{\arabic{enumi}.}
\tightlist
\item
  \textbf{Exponential reserve.} For \(z\in\mathcal T_n\), on
  \(\mathrm{supp}\chi_n\) one has
  \(\mathrm{Re}\,(i\langle p,z\rangle)=-\langle p,\mathrm{Im}\,z\rangle\le A-c\lvert p\rvert\),
  by C10 and distance \(\le3\). So \(e^{i\langle\cdot,z\rangle}\chi_n\)
  has finite seminorms weighted by \(e^{c\lvert p\rvert/2}\).
\item
  \textbf{Differentiability.} Write
  \(L_z(p)=e^{i\langle p,z\rangle}\chi_n(p)\). For each Schwartz
  seminorm \(q\), exponential decay gives
  \(q(L_{z+h}-L_z-i\langle\cdot,h\rangle L_z)\le C_q|h|^2\) for all
  sufficiently small \(h\). Thus \(z\mapsto L_z\) is complex Fréchet
  differentiable into \(\mathcal S\).
\item
  \textbf{Analyticity.} Compose with the continuous functional
  \(\mathcal FW\). Complex differentiability on an open set gives
  analyticity.
\end{enumerate}

The argument applies to every Schwartz test with the prescribed
exponential form; C13 asserts the existence of such tests and the
resulting holomorphy. \emph{Lean:}
\href{https://github.com/yang-mills-group/OpenAI-YM-20261028/blob/3456e5e75cf8d4cd558955a5c7638fba2d68f464/YM/Work/C13.lean\#L61-L89}{\texttt{C13.lens\_reserve}},
\href{https://github.com/yang-mills-group/OpenAI-YM-20261028/blob/3456e5e75cf8d4cd558955a5c7638fba2d68f464/YM/Work/C13.lean\#L113-L125}{\texttt{C13.analytic\_interior}};
\href{https://github.com/yang-mills-group/OpenAI-YM-20261028/blob/3456e5e75cf8d4cd558955a5c7638fba2d68f464/YM/Work/C13/Decay.lean\#L93-L123}{\texttt{ExpReserve.of\_cut}},
\href{https://github.com/yang-mills-group/OpenAI-YM-20261028/blob/3456e5e75cf8d4cd558955a5c7638fba2d68f464/YM/Work/C13/Variation.lean\#L150-L193}{\texttt{hasFDerivAt\_shift}}.

\textbf{Lemma V.2.q (partial boundary values as integrals; agreement
with \(F\); C14, C15).} As in V.2.6(e), (f). \emph{Proof sketch.}

\begin{itemize}
\tightlist
\item
  \textbf{C14.} For real \(x\) the Fourier--Laplace test at \(x+i\eta\)
  is the modulation \(e^{i\langle\cdot,x\rangle}g\), with
  \(g=e^{-\langle\eta,\cdot\rangle}\chi_n\). Hence
  \(\int\mathcal W_n(v)(x+i\eta)f(x)\,dx=(\mathcal FW)(g\check f)\). The
  distribution is moved inside the \(x\)-integral via Riemann sums in
  \(\mathcal S\), first for compactly supported \(f\), then by density.
  Integrability comes from the polynomial growth of
  \(x\mapsto\mathcal W(x+i\eta)\). Cone support is not needed.
\item
  \textbf{C15.}

  \begin{enumerate}
  \def\labelenumi{\arabic{enumi}.}
  \tightlist
  \item
    Since \(\underline y\in V^+_n\) and
    \(\mathrm{supp}\mathcal FW\subseteq\mathcal V_n\), the multipliers
    \(e^{-\langle\underline y,\cdot\rangle}\chi_n\) and
    \(e^{-\langle\underline y,\cdot\rangle}\) agree near the support.
  \item
    So V.2.i gives \(W_n(v)[\underline y]=\psi_v(y;\cdot)\) on compactly
    supported tests, hence on all tests (\(\psi(y;\cdot)\) is tempered
    by (I) and (PG)).
  \item
    With C14, \(\int\mathcal W(x+i\underline y)f=\int F^y(x)f\) for all
    \(f\). Both integrands are continuous, so they are equal.
  \end{enumerate}
\end{itemize}

\emph{Lean:}
\href{https://github.com/yang-mills-group/OpenAI-YM-20261028/blob/3456e5e75cf8d4cd558955a5c7638fba2d68f464/YM/Work/C14.lean\#L133-L141}{\texttt{C14.modulation\_eq}},
\href{https://github.com/yang-mills-group/OpenAI-YM-20261028/blob/3456e5e75cf8d4cd558955a5c7638fba2d68f464/YM/Work/C14.lean\#L233-L253}{\texttt{C14.emit\_integral}},
\href{https://github.com/yang-mills-group/OpenAI-YM-20261028/blob/3456e5e75cf8d4cd558955a5c7638fba2d68f464/YM/Work/C15.lean\#L60-L74}{\texttt{C15.blanketed\_eq}},
\href{https://github.com/yang-mills-group/OpenAI-YM-20261028/blob/3456e5e75cf8d4cd558955a5c7638fba2d68f464/YM/Work/C15.lean\#L83-L99}{\texttt{C15.partialD\_psi}}.

\textbf{Lemma V.2.r (Lorentz covariance in the tube; C16).} As in
V.2.6(g). \emph{Proof sketch.}

\begin{itemize}
\tightlist
\item
  \(\mathcal V_n\) is invariant under the dual action.
\item
  The Fourier--Laplace test at \(\Lambda z\), composed with
  \(\tilde\Lambda\), differs from the lens at \(z\) only away from
  \(\mathcal V_n\): both cutoffs equal \(1\) there, and the phases agree
  by duality.
\item
  Cone support, the intertwining of \(\mathcal F\) with \(\Lambda\) and
  \(\tilde\Lambda\), and (Cov) conclude.
\end{itemize}

\emph{Lean:}
\href{https://github.com/yang-mills-group/OpenAI-YM-20261028/blob/3456e5e75cf8d4cd558955a5c7638fba2d68f464/YM/Work/C16.lean\#L45-L73}{\texttt{C16.lens\_support}},
\href{https://github.com/yang-mills-group/OpenAI-YM-20261028/blob/3456e5e75cf8d4cd558955a5c7638fba2d68f464/YM/Work/C16.lean\#L75-L91}{\texttt{C16.answer}}.

\paragraph{4.6 Locality}\label{locality}

\textbf{Lemma V.2.s (imaginary boosts at a spectator profile; C17).} As
in V.2.7(a). \emph{Proof sketch.}

\begin{enumerate}
\def\labelenumi{\arabic{enumi}.}
\tightlist
\item
  \textbf{Spectator profiles.} The profiles \(y_m=d/(m+1)\) are
  spectator profiles. Monotone heights give
  \(\underline{y_m}\in\bar V^+_n\), and \(\sigma\) fixes
  \(\underline{y_m}\). The closure statements follow by approaching
  through strictly ordered heights.
\item
  \textbf{Choice of direction.} Let \(\xi=x_{i+1}-x_i\) be spacelike, so
  \(\vec\xi\ne0\), and take \(e=\vec\xi/\lvert\vec\xi\rvert\). Then
  \[\mathrm{Im}\big(\Lambda_e(ib)w\big)_k=\mathrm{Im}\,w_k+\sin b\cdot c_e(\mathrm{Re}\,w_k)+(\cos b-1)\,s_e(\mathrm{Im}\,w_k),\]
  with \(c_e(\vec X,t)=(te,e\cdot\vec X)\) and
  \(s_e(\vec X,t)=((e\cdot\vec X)e,t)\).
\item
  \textbf{The pair difference.} Its imaginary part is \(0\) at \(b=0\),
  and \(c_e(\xi)=(\xi^0e,\lvert\vec\xi\rvert)\in V_+\) exactly because
  \(\lvert\xi^0\rvert<\lvert\vec\xi\rvert\) (the Jost condition). So the
  difference enters \(V_+\) for small \(b>0\).
\item
  \textbf{The other consecutive differences.} They are already in
  \(V_+\) and stay there by continuity, uniformly near the base point.
\item
  \textbf{The swapped tube.} After the swap the pair difference is
  \(-\xi\), and the opposite rapidity \(-ib\) compensates.
\end{enumerate}

For \(n=2\) the profile is identically \(0\), so the germs sit at real
points. \emph{Lean:}
\href{https://github.com/yang-mills-group/OpenAI-YM-20261028/blob/3456e5e75cf8d4cd558955a5c7638fba2d68f464/YM/Work/C17.lean\#L247-L275}{\texttt{C17.height\_path}},
\href{https://github.com/yang-mills-group/OpenAI-YM-20261028/blob/3456e5e75cf8d4cd558955a5c7638fba2d68f464/YM/Work/C17.lean\#L44-L77}{\texttt{C17.fleck\_of\_spectator}},
\href{https://github.com/yang-mills-group/OpenAI-YM-20261028/blob/3456e5e75cf8d4cd558955a5c7638fba2d68f464/YM/Work/C17.lean\#L79-L92}{\texttt{C17.choose\_direction}},
\href{https://github.com/yang-mills-group/OpenAI-YM-20261028/blob/3456e5e75cf8d4cd558955a5c7638fba2d68f464/YM/Work/C17.lean\#L126-L171}{\texttt{C17.enters}},
\href{https://github.com/yang-mills-group/OpenAI-YM-20261028/blob/3456e5e75cf8d4cd558955a5c7638fba2d68f464/YM/Work/C17.lean\#L173-L241}{\texttt{C17.local\_paths}}.

\textbf{Lemma V.2.t (complexified labels; the two germs; C18, C19).} As
in V.2.7(b), (c). \emph{Proof sketch.}

\begin{itemize}
\tightlist
\item
  \textbf{C18.} A continuous one-parameter group on a finite-dimensional
  space is \(e^{tX}\) (via \(\int_0^t\) and invertibility for small
  \(t\)), and \(e^{\zeta X}\) is entire.
\item
  \textbf{C19.}

  \begin{enumerate}
  \def\labelenumi{\arabic{enumi}.}
  \tightlist
  \item
    Put all labels at one level \(d\) and write \(K_{d,e}(\zeta)\) for
    the entire extension from C18. Define the germ at
    \(x+i\underline y\) by
    \(w\mapsto\mathcal W_n(\mathrm{pin}_d K_{d,e}(it)v)(\Lambda_e(it)w)\).
    It is holomorphic near \(x+i\underline y\) because
    \(\Lambda_e(it)(x+i\underline y)\in\mathcal T_n\) and
    \(\Lambda_e(it)\) is complex linear.
  \item
    Fix \(w\in\mathcal T_n\) near the point. The function
    \(\zeta\mapsto\mathcal W_n(\mathrm{pin}_d K_{d,e}(\zeta)v)(\Lambda_e(\zeta)w)\)
    is holomorphic near the segment \(i[0,t]\), by the path conditions,
    and constant for real \(\zeta\), by V.2.6(g).
  \item
    The identity theorem makes it constant on the segment, so the germ
    equals \(\mathcal W_n(v)\) on \(\mathcal T_n\) near
    \(x+i\underline y\).
  \item
    The tilted germ uses \(-it\), the labels \(v\circ s\), and
    \(\sigma\).
  \end{enumerate}
\end{itemize}

\emph{Lean:}
\href{https://github.com/yang-mills-group/OpenAI-YM-20261028/blob/3456e5e75cf8d4cd558955a5c7638fba2d68f464/YM/Work/C18.lean\#L12-L44}{\texttt{C18.eq\_exponential}},
\href{https://github.com/yang-mills-group/OpenAI-YM-20261028/blob/3456e5e75cf8d4cd558955a5c7638fba2d68f464/YM/Work/C18.lean\#L46-L54}{\texttt{C18.entire}},
\href{https://github.com/yang-mills-group/OpenAI-YM-20261028/blob/3456e5e75cf8d4cd558955a5c7638fba2d68f464/YM/Work/C19.lean\#L161-L211}{\texttt{C19.paneAlong}}.

\textbf{Lemma V.2.u (agreement at Euclidean points; C20).} As in
V.2.7(d). \emph{Proof sketch.}

\begin{enumerate}
\def\labelenumi{\arabic{enumi}.}
\tightlist
\item
  \textbf{The Wick slice.} Let \(M(\vec x,\tau)=(\vec x,i\tau)\). Its
  complex span is \(\mathbb C^{4n}\). If \(x\) has zero times, then
  \(M(\vec x,y)=x+i\underline y\).
\item
  \textbf{The original germ.} On the open set of strictly ordered
  \(\tau\), \(M\) maps into \(\mathcal T_n\), and
  \(g^{\rm o}\circ M=\mathcal W\circ M=F\) at imaginary times. That
  equals the Schwinger function \(T((\vec x_k,\tau_k);v)\) (V.2.6(f),
  (SC)).
\item
  \textbf{The tilted germ.} On the swapped-ordered set,
  \(g^{\rm t}\circ M\) equals the same Schwinger function, by
  permutation symmetry of \(T\).
\item
  \textbf{Smoothness at the base point.} The point \((\vec x,y)\) is
  non-coincident: the pair has equal times but
  \(\vec x_i\ne\vec x_{i+1}\) by spacelikeness. So \(T\) is smooth
  there.
\item
  \textbf{Jets.} Both open sets accumulate at \((\vec x,y)\). Hence all
  real jets of \(g^{\rm o}\circ M\) and \(g^{\rm t}\circ M\) at that
  point equal those of \(T\).
\item
  \textbf{Conclusion.} Real jets on a slice that spans over
  \(\mathbb C\) determine all complex derivatives, so
  \(g^{\rm o}=g^{\rm t}\) near \(x+i\underline y\).
\end{enumerate}

No Schwarz reflection and no edge-of-the-wedge theorem are used.
\emph{Lean:}
\href{https://github.com/yang-mills-group/OpenAI-YM-20261028/blob/3456e5e75cf8d4cd558955a5c7638fba2d68f464/YM/Work/C20.lean\#L34-L46}{\texttt{C20.jets\_on\_open\_side}},
\href{https://github.com/yang-mills-group/OpenAI-YM-20261028/blob/3456e5e75cf8d4cd558955a5c7638fba2d68f464/YM/Work/C20.lean\#L67-L87}{\texttt{C20.complex\_germ}},
\href{https://github.com/yang-mills-group/OpenAI-YM-20261028/blob/3456e5e75cf8d4cd558955a5c7638fba2d68f464/YM/Work/C20.lean\#L230-L240}{\texttt{C20.emit\_table}}.

\textbf{Lemma V.2.v (local realization and exchange; C21).} As in
V.2.7(e), (f). \emph{Proof sketch.}

\begin{itemize}
\tightlist
\item
  \textbf{C21.} As \(\eta\to\underline y\) within \(V^+_n\),
  \(\int\mathcal W(\xi+i\eta)f(\xi)\,d\xi\to W[\underline y](https://github.com/scottnarmstrong/yang-mills-informalization/blob/main/outline/units/f)\),
  by C14 and continuity of \(\eta\mapsto W[\eta]\) on \(\bar V^+_n\).
  Near \(x+i\underline y\) the germ is bounded and equals \(\mathcal W\)
  on \(\mathcal T_n\), so dominated convergence gives the integral of
  the germ. The swapped case is the same after a change of variables.
\item
  \textbf{Propagation.} For \(x\in\Gamma_Q\), the path
  \(x\mapsto(\vec x,tx^0)\), \(t\in[0,1]\), stays in \(\Gamma_Q\).
  Agreement of the two germs is an open and closed condition along this
  path, so it passes from the zero-time point to \(x\).
\item
  \textbf{Local equality.} On a small ball around each point of
  \(\Gamma_Q\) both partial boundary values are integrals of the same
  germ, so their difference vanishes on \(\Gamma_Q\).
\item
  \textbf{Exchange.} Let \(m\to\infty\). Then
  \(W[\underline{y_m}]\to W[0]=W\), by continuity in \(\eta\) and
  because \(\chi_n=1\) near \(\mathcal V_n\), and similarly on the
  swapped side.
\end{itemize}

\emph{Lean:}
\href{https://github.com/yang-mills-group/OpenAI-YM-20261028/blob/3456e5e75cf8d4cd558955a5c7638fba2d68f464/YM/Work/C21.lean\#L15-L56}{\texttt{C21.realize}},
\href{https://github.com/yang-mills-group/OpenAI-YM-20261028/blob/3456e5e75cf8d4cd558955a5c7638fba2d68f464/YM/Shift/Exchange.lean\#L43-L56}{\texttt{Panes.reaches}},
\href{https://github.com/yang-mills-group/OpenAI-YM-20261028/blob/3456e5e75cf8d4cd558955a5c7638fba2d68f464/YM/Shift/Exchange.lean\#L68-L90}{\texttt{Panes.forgets}},
\href{https://github.com/yang-mills-group/OpenAI-YM-20261028/blob/3456e5e75cf8d4cd558955a5c7638fba2d68f464/YM/Shift/Exchange.lean\#L92-L112}{\texttt{NearPair.actor}};
\href{https://github.com/yang-mills-group/OpenAI-YM-20261028/blob/3456e5e75cf8d4cd558955a5c7638fba2d68f464/YM/Shift/Swap.lean\#L101-L130}{\texttt{Pane.propagate}},
\href{https://github.com/yang-mills-group/OpenAI-YM-20261028/blob/3456e5e75cf8d4cd558955a5c7638fba2d68f464/YM/Shift/Sol.lean\#L85-L96}{\texttt{Blanket.flatD}}.

\textbf{Lemma V.2.w (locality for words; C27).} As in V.2.7(g).
\emph{Proof sketch.} If \(\mathrm{supp}f-\mathrm{supp}g\) is spacelike,
the tensor test of the word \(a\,(u,f)(u',g)\,c\) is supported in
\(\Gamma_Q\), with \(Q\) the slots of \((u,f)\) and \((u',g)\).
Composing with \(\sigma\) swaps the two factors. Joint locality applies
to all Schwartz tests supported in \(\Gamma_Q\), so no density argument
is needed. The pair may be at any position in a word of any length
\(\ge2\). \emph{Lean:}
\href{https://github.com/yang-mills-group/OpenAI-YM-20261028/blob/3456e5e75cf8d4cd558955a5c7638fba2d68f464/YM/Work/C27.lean\#L12-L19}{\texttt{C27.coordinate\_tsupport}},
\href{https://github.com/yang-mills-group/OpenAI-YM-20261028/blob/3456e5e75cf8d4cd558955a5c7638fba2d68f464/YM/Work/C27.lean\#L28-L39}{\texttt{C27.exchange}}.

\paragraph{4.7 Real-time vectors, Gram limits,
nulls}\label{real-time-vectors-gram-limits-nulls}

\textbf{Lemma V.2.x (Gram form of heated vectors; C22).} As in V.2.8(a).
\emph{Proof sketch.}

\begin{enumerate}
\def\labelenumi{\arabic{enumi}.}
\tightlist
\item
  \textbf{Imaginary times.} At imaginary times each heated vector is the
  GNS vector of a positive-time Euclidean word.

  \begin{itemize}
  \tightlist
  \item
    Their inner product is the reflection-positive form, i.e.~the
    Schwinger function of \emph{reflected first word, second word}:
    times negated, order reversed, labels starred.
  \item
    By (SC) this is \(F\) at \(i\cdot h\). Here the star \(J\) appears
    because it is the induced reflection adjoint, and null-space
    compatibility is used.
  \end{itemize}
\item
  \textbf{Continuation in real times.} Introduce one complex parameter
  \(c\) and the times \(iy_k+c\,r^0_k\). Both sides are holomorphic on a
  strip \(\lvert\mathrm{Im}\,c\rvert<\delta\); the left side is
  antiholomorphic in the first slot, handled by conjugation.
\item
  \textbf{Conclusion.} The two sides agree at \(c=i\tau\) for small
  \(\tau\), hence at \(c=1\).
\end{enumerate}

The width formula of C03 is not used. \emph{Lean:}
\href{https://github.com/yang-mills-group/OpenAI-YM-20261028/blob/3456e5e75cf8d4cd558955a5c7638fba2d68f464/YM/Work/C22.lean\#L70-L107}{\texttt{C22.euclideanRef}},
\href{https://github.com/yang-mills-group/OpenAI-YM-20261028/blob/3456e5e75cf8d4cd558955a5c7638fba2d68f464/YM/Work/C22.lean\#L170-L192}{\texttt{C22.onImag}},
\href{https://github.com/yang-mills-group/OpenAI-YM-20261028/blob/3456e5e75cf8d4cd558955a5c7638fba2d68f464/YM/Work/C22.lean\#L274-L336}{\texttt{C22.pointJoined}}.

\textbf{Lemma V.2.y (heat cups and their Gram limit; C23, C24).} As in
V.2.8(b), (c). \emph{Proof sketch.}

\begin{itemize}
\tightlist
\item
  \textbf{C23.}

  \begin{itemize}
  \tightlist
  \item
    \(x\mapsto\Psi(\cdot;x)\) is continuous and bounded by
    \(\|\Omega\|\prod\|B_p(v_k)\|\), since the propagators are
    contractions. So the integrand is Bochner integrable.
  \item
    Translation covariance comes from the difference structure of
    \(\Psi\) and \(U(a)\Omega=\Omega\).
  \end{itemize}
\item
  \textbf{C24.}

  \begin{enumerate}
  \def\labelenumi{\arabic{enumi}.}
  \tightlist
  \item
    By the pointwise Gram identity and Fubini,
    \(\langle\Phi_s(b),\Phi_{s'}(c)\rangle\) is the integral of \(F\) at
    heights \(h=(-ns,\dots,-s,s',\dots,ms')\) against the test
    \(f_b^{\ast}\otimes f_c\).
  \item
    As \((s,s')\to0\), \(h\to0\) within \(E_{n+m}\) along a
    two-parameter family with no fixed direction. The unrestricted
    boundary limit applies to this family, whose two scales vary
    independently.
  \item
    (BV) gives the limit \(W(b^{\ast}c)\).
  \end{enumerate}
\end{itemize}

\emph{Lean:}
\href{https://github.com/yang-mills-group/OpenAI-YM-20261028/blob/3456e5e75cf8d4cd558955a5c7638fba2d68f464/YM/Work/C23.lean\#L53-L83}{\texttt{C23.shot\_control}},
\href{https://github.com/yang-mills-group/OpenAI-YM-20261028/blob/3456e5e75cf8d4cd558955a5c7638fba2d68f464/YM/Work/C23.lean\#L207-L216}{\texttt{C23.heatCup\_shift}},
\href{https://github.com/yang-mills-group/OpenAI-YM-20261028/blob/3456e5e75cf8d4cd558955a5c7638fba2d68f464/YM/Work/C24.lean\#L61-L73}{\texttt{C24.toward}},
\href{https://github.com/yang-mills-group/OpenAI-YM-20261028/blob/3456e5e75cf8d4cd558955a5c7638fba2d68f464/YM/Work/C24.lean\#L118-L144}{\texttt{C24.integrated}}.

\textbf{Lemma V.2.z (nulls; word insertion; C26, C29).} As in V.2.8(e),
(f). \emph{Proof sketch.}

\begin{itemize}
\tightlist
\item
  \textbf{C26.} Let \(u=\mathrm{pin}_d[w]\) be a marker.

  \begin{enumerate}
  \def\labelenumi{\arabic{enumi}.}
  \tightlist
  \item
    By C25, every \(W_n(\dots,u,\dots)\) vanishes as a distribution.
  \item
    By V.2.6(f), \(F^y=0\) for ordered \(y\), so the Schwinger functions
    with \(w\) in some slot vanish at time-ordered non-coincident
    points.
  \item
    Symmetry and continuity give vanishing at all non-coincident points.
    So \(w\in N_d(T)=N_d(S)\), and \(u=0\).
  \end{enumerate}
\item
  \textbf{C29.}

  \begin{itemize}
  \tightlist
  \item
    word bilinearity is the multilinearity of \(W_n\) composed with that
    of \(\otimes\).
  \item
    continuity in the two tests follows from continuity of \(W_n(v)\),
    of \((f_1,\dots,f_n)\mapsto\otimes f_k\), and of conjugation.
  \end{itemize}
\end{itemize}

\emph{Lean:}
\href{https://github.com/yang-mills-group/OpenAI-YM-20261028/blob/3456e5e75cf8d4cd558955a5c7638fba2d68f464/YM/Work/C26.lean\#L23-L38}{\texttt{C26.marker\_on\_joint}},
\href{https://github.com/yang-mills-group/OpenAI-YM-20261028/blob/3456e5e75cf8d4cd558955a5c7638fba2d68f464/YM/Work/C26.lean\#L40-L51}{\texttt{C26.strict\_zero}},
\href{https://github.com/yang-mills-group/OpenAI-YM-20261028/blob/3456e5e75cf8d4cd558955a5c7638fba2d68f464/YM/Work/C26.lean\#L55-L75}{\texttt{C26.read\_zero}},
\href{https://github.com/yang-mills-group/OpenAI-YM-20261028/blob/3456e5e75cf8d4cd558955a5c7638fba2d68f464/YM/Work/C29.lean\#L110}{\texttt{C29.answer}}.

\paragraph{4.8 Assembly (VIII.1) and the field
realization}\label{assembly-viii.1-and-the-field-realization}

The reconstruction is assembled in
\href{https://github.com/yang-mills-group/OpenAI-YM-20261028/blob/3456e5e75cf8d4cd558955a5c7638fba2d68f464/YM/Next/BOUNDARY/Stage.lean}{\texttt{Next/BOUNDARY/Stage.lean}},
in this order.

\begin{enumerate}
\def\labelenumi{\arabic{enumi}.}
\tightlist
\item
  \href{https://github.com/yang-mills-group/OpenAI-YM-20261028/blob/3456e5e75cf8d4cd558955a5c7638fba2d68f464/YM/Next/BOUNDARY/Stage.lean\#L39-L44}{\texttt{rhythm}}
  chains C01 → C02 → C03 → C04.

  \begin{itemize}
  \tightlist
  \item
    \href{https://github.com/yang-mills-group/OpenAI-YM-20261028/blob/3456e5e75cf8d4cd558955a5c7638fba2d68f464/YM/Org/Theatre.lean\#L29-L31}{\texttt{Rhythm.opens}}
    (V.2.h) produces \(W\).
  \item
    \href{https://github.com/yang-mills-group/OpenAI-YM-20261028/blob/3456e5e75cf8d4cd558955a5c7638fba2d68f464/YM/Next/BOUNDARY/Stage.lean\#L52-L78}{\texttt{at\_contour}}
    shows that any prescribed continuation--boundary pair for \(T\)
    satisfies (BV).
  \end{itemize}
\item
  \href{https://github.com/yang-mills-group/OpenAI-YM-20261028/blob/3456e5e75cf8d4cd558955a5c7638fba2d68f464/YM/Next/BOUNDARY/Stage.lean\#L82-L90}{\texttt{illuminated}}
  applies C07, C09 (with A16 and C08), C05 and C06.
  \href{https://github.com/yang-mills-group/OpenAI-YM-20261028/blob/3456e5e75cf8d4cd558955a5c7638fba2d68f464/YM/Next/BOUNDARY/Stage.lean\#L92}{\texttt{spectra}}
  uses Lorentz improvement of spectral support to get cone support.
\item
  \href{https://github.com/yang-mills-group/OpenAI-YM-20261028/blob/3456e5e75cf8d4cd558955a5c7638fba2d68f464/YM/Next/BOUNDARY/Stage.lean\#L94-L95}{\texttt{blankets}}
  (C10, C12),
  \href{https://github.com/yang-mills-group/OpenAI-YM-20261028/blob/3456e5e75cf8d4cd558955a5c7638fba2d68f464/YM/Next/BOUNDARY/Stage.lean\#L97-L99}{\texttt{visible}}
  (C13) and
  \href{https://github.com/yang-mills-group/OpenAI-YM-20261028/blob/3456e5e75cf8d4cd558955a5c7638fba2d68f464/YM/Next/BOUNDARY/Stage.lean\#L101-L111}{\texttt{neon\_any}}
  (C13--C15) give (TC).
  \href{https://github.com/yang-mills-group/OpenAI-YM-20261028/blob/3456e5e75cf8d4cd558955a5c7638fba2d68f464/YM/Next/BOUNDARY/Unique.lean\#L187-L203}{\texttt{emit\_eq}}
  shows that the continuation on \(\mathcal T_n\) does not depend on the
  cutoff.
\item
  \href{https://github.com/yang-mills-group/OpenAI-YM-20261028/blob/3456e5e75cf8d4cd558955a5c7638fba2d68f464/YM/Next/BOUNDARY/Stage.lean\#L122-L138}{\texttt{windows}}
  combines C17, C19 (with C18 and C16), C20 and C21.
  \href{https://github.com/yang-mills-group/OpenAI-YM-20261028/blob/3456e5e75cf8d4cd558955a5c7638fba2d68f464/YM/Org/Theatre.lean\#L64-L67}{\texttt{Show.uses}}
  then gets joint locality from the boundary exchange lemma and Lorentz
  improvement of spectral support.
\item
  \href{https://github.com/yang-mills-group/OpenAI-YM-20261028/blob/3456e5e75cf8d4cd558955a5c7638fba2d68f464/YM/Next/BOUNDARY/Stage.lean\#L142-L162}{\texttt{revealing}}
  applies C22, C26 with C25, and C24 with C23.
\item
  \href{https://github.com/yang-mills-group/OpenAI-YM-20261028/blob/3456e5e75cf8d4cd558955a5c7638fba2d68f464/YM/Next/BOUNDARY/Stage.lean\#L167-L184}{\texttt{establishing}}
  applies C28, C29 and C27, then builds the field system (step 7).
\item
  The field system is built (V.2.10) in the following order.

  \begin{enumerate}
  \def\labelenumii{\arabic{enumii}.}
  \tightlist
  \item
    The heat cups converge: the two-variable Gram limit makes
    \(s\mapsto\Phi_s(b)\) Cauchy
    (\href{https://github.com/yang-mills-group/OpenAI-YM-20261028/blob/3456e5e75cf8d4cd558955a5c7638fba2d68f464/YM/Fields/Lo.lean\#L19-L31}{\texttt{norms\_kernel}},
    \href{https://github.com/yang-mills-group/OpenAI-YM-20261028/blob/3456e5e75cf8d4cd558955a5c7638fba2d68f464/YM/Fields/Lo.lean\#L55-L60}{\texttt{Cups.out}}).
    The limit vectors \(\Phi_0(b)\) satisfy
    \(\langle\Phi_0(b),\Phi_0(c)\rangle=\langle b,c\rangle_W\) and
    \(\Phi_0(\varnothing)=\Omega\).
  \item
    Hence \(\langle\cdot,\cdot\rangle_W\) is positive
    (\href{https://github.com/yang-mills-group/OpenAI-YM-20261028/blob/3456e5e75cf8d4cd558955a5c7638fba2d68f464/YM/Fields/Lo.lean\#L66-L71}{\texttt{Outfield.safe}}),
    and V.2.a gives its Kolmogorov space.
  \item
    Fields act by left multiplication of words. The identity
    \(\langle (u,f)b,c\rangle_W=\langle b,(Ju,\bar f)c\rangle_W\) shows
    that a null linear combination of words remains null after left
    multiplication, so the fields are well defined
    (\href{https://github.com/yang-mills-group/OpenAI-YM-20261028/blob/3456e5e75cf8d4cd558955a5c7638fba2d68f464/YM/Fields/Stems.lean\#L120-L137}{\texttt{Soil.starter}}).
  \item
    The same identity gives the formal adjoint. The other properties of
    \(W\) give the remaining clauses of V.2.9: temperedness (continuity
    in the two tests); local commutativity (word locality); field
    faithfulness (label faithfulness).
  \item
    A unitary representation \(U\) of \(\mathcal P\) is defined by
    \(U(g)[b]=[g\cdot b]\), using invariance of the word Gram form and
    continuity of word matrix elements
    (\href{https://github.com/yang-mills-group/OpenAI-YM-20261028/blob/3456e5e75cf8d4cd558955a5c7638fba2d68f464/YM/Fields/Motion.lean\#L72-L87}{\texttt{Soil.mobilize}}).
  \item
    \([b]\mapsto\Phi_0(b)\) is an isometry into \(\mathcal H\) that
    intertwines translations and sends the vacuum to \(\Omega\)
    (\href{https://github.com/yang-mills-group/OpenAI-YM-20261028/blob/3456e5e75cf8d4cd558955a5c7638fba2d68f464/YM/Fields/Motion.lean\#L31-L53}{\texttt{universal\_embed}},
    \href{https://github.com/yang-mills-group/OpenAI-YM-20261028/blob/3456e5e75cf8d4cd558955a5c7638fba2d68f464/YM/Fields/Lo.lean\#L76-L82}{\texttt{Outfield.daylight}}).
  \item
    \textbf{Spectrum condition}
    (\href{https://github.com/yang-mills-group/OpenAI-YM-20261028/blob/3456e5e75cf8d4cd558955a5c7638fba2d68f464/YM/Fields/Spectrum.lean\#L108-L114}{\texttt{Action.spectrum}}).
    The isometric intertwiner identifies each scalar spectral measure
    with that of its image in \(\mathcal H\), by equality of
    characteristic functions. These measures are supported in
    \(\{p^0\ge0\}\) by positivity of \(H\). Lorentz covariance of \(U\)
    upgrades this to support in \(\bar V_+\).
  \item
    \textbf{Vacuum uniqueness}
    (\href{https://github.com/yang-mills-group/OpenAI-YM-20261028/blob/3456e5e75cf8d4cd558955a5c7638fba2d68f464/YM/Fields/Spectrum.lean\#L116-L121}{\texttt{Action.only}}).
    Under one-dimensionality of the heat-fixed space, an invariant
    vector is mapped to the ground space of \(\mathcal H\), which is
    \(\mathbb C\Omega\).
  \end{enumerate}
\item
  \href{https://github.com/yang-mills-group/OpenAI-YM-20261028/blob/3456e5e75cf8d4cd558955a5c7638fba2d68f464/YM/Next/BOUNDARY/Stage.lean\#L217-L221}{\texttt{Compatible.vacuum}}
  adds uniqueness for the selected phase.
  \href{https://github.com/yang-mills-group/OpenAI-YM-20261028/blob/3456e5e75cf8d4cd558955a5c7638fba2d68f464/YM/Next/BOUNDARY/Stage.lean\#L247-L255}{\texttt{after\_materializes}}
  runs the chain for the mixed table (\(S=T\), without vacuum
  uniqueness) and for the selected phase.
\end{enumerate}

\subsubsection{5. Where the difficulty
lies}\label{where-the-difficulty-lies-21}

The boost argument must control moving labels and tests without imposing
a topology on the full label union. All labels are first placed in one
finite-dimensional level; Schwartz estimates then justify the derivative
and its boundary limit. The locality argument needs equality of complete
jets on the Wick slice. Equality of values at zero real times would not
determine the two germs. Smoothness of the Euclidean correlations at
non-coincident configurations, including equal pair times, supplies
those jets.

The Taylor formula in V.2.h gives a bound uniform over all nondecreasing
heights. This uniformity permits varying tests in the Laplace identity
and independent heating scales in the Gram limit. Positivity of the word
kernel then suffices for the field construction; surjectivity of the
embedding into the OS space is not required.

Several formal hypotheses are stronger than the corresponding proof
needs. The slab-insertion construction uses the slab bound but not the
derivative laws of the spread. C05 and C06 use the tested identities and
growth bound without Euclidean root values. C14 uses the damping
properties without spectral support; C16 uses spectral support and
Lorentz covariance without agreement with the strict-time continuation.
These distinctions do not remove the hypotheses from the statements
above.

The pairing is Euclidean, so momenta transform by
\(\tilde\Lambda=\mathsf r\Lambda\mathsf r\). The tube uses absolute
coordinates and imposes no condition on its first imaginary point. For
\(n=1\), it is \(\mathbb C^4\) and \(\mathcal V_1=\{0\}\). The
polynomial bound is per label tuple, with no uniformity in the arity.
The chosen value \(B_0=0\) is outside the admissible widths. The Laplace
identity is stated on compactly supported tests.

\subsubsection{6. Notation}\label{notation-21}

{\def\LTcaptype{none} % do not increment counter
\begin{longtable}[]{@{}
  >{\raggedright\arraybackslash}p{(\linewidth - 4\tabcolsep) * \real{0.3333}}
  >{\raggedright\arraybackslash}p{(\linewidth - 4\tabcolsep) * \real{0.3333}}
  >{\raggedright\arraybackslash}p{(\linewidth - 4\tabcolsep) * \real{0.3333}}@{}}
\toprule\noalign{}
\begin{minipage}[b]{\linewidth}\raggedright
object
\end{minipage} & \begin{minipage}[b]{\linewidth}\raggedright
code name
\end{minipage} & \begin{minipage}[b]{\linewidth}\raggedright
proposed notation
\end{minipage} \\
\midrule\noalign{}
\endhead
\bottomrule\noalign{}
\endlastfoot
OS Hilbert space of \(T\), time axis &
\href{https://github.com/yang-mills-group/OpenAI-YM-20261028/blob/3456e5e75cf8d4cd558955a5c7638fba2d68f464/YM/OS/Dynamics.lean\#L57-L61}{\texttt{OS.Chamber}}
& \(\mathcal H\), \(\Omega\), \(H\), \(U(\vec a)\) \\
slab operators of a spread; slab bound &
\href{https://github.com/yang-mills-group/OpenAI-YM-20261028/blob/3456e5e75cf8d4cd558955a5c7638fba2d68f464/YM/OS/Slabs.lean\#L101-L104}{\texttt{OS.Spread}},
\href{https://github.com/yang-mills-group/OpenAI-YM-20261028/blob/3456e5e75cf8d4cd558955a5c7638fba2d68f464/YM/OS/Slabs.lean\#L111-L113}{\texttt{AllCopy}}
& \(A_d(w)\); (SB) \\
hypothesis package of C01 &
\href{https://github.com/yang-mills-group/OpenAI-YM-20261028/blob/3456e5e75cf8d4cd558955a5c7638fba2d68f464/YM/Bid/Corners.lean\#L23-L30}{\texttt{Bid.Pantry}}
& (H) \\
labels of \(S\); tags; insertion & \texttt{S.Union}, \texttt{S.Tag},
\texttt{pin} & \(\mathcal F\), \((d,v)\), \(\mathrm{pin}_d\) \\
Euclidean function of \(T\) on labels of \(S\) & \texttt{S.useAt\ T} &
\(T(x;v)\) \\
single insertions between heat factors &
\href{https://github.com/yang-mills-group/OpenAI-YM-20261028/blob/3456e5e75cf8d4cd558955a5c7638fba2d68f464/YM/Fields/Tall.lean\#L19-L20}{\texttt{Small}}
(\texttt{bar}) & \(B_p(v)\) \\
propagator between insertions &
\href{https://github.com/yang-mills-group/OpenAI-YM-20261028/blob/3456e5e75cf8d4cd558955a5c7638fba2d68f464/YM/Fields/Tall.lean\#L33}{\texttt{Small.rip}}
& \(G_p(w,\vec a)\) \\
strict-time matrix element &
\href{https://github.com/yang-mills-group/OpenAI-YM-20261028/blob/3456e5e75cf8d4cd558955a5c7638fba2d68f464/YM/Fields/Tall.lean\#L49-L51}{\texttt{Small.matrix}}
& \(M_p(z;\vec x)\) \\
ordered heights; strict region &
\href{https://github.com/yang-mills-group/OpenAI-YM-20261028/blob/3456e5e75cf8d4cd558955a5c7638fba2d68f464/YM/Boundary/Contour.lean\#L27}{\texttt{east}},
\href{https://github.com/yang-mills-group/OpenAI-YM-20261028/blob/3456e5e75cf8d4cd558955a5c7638fba2d68f464/YM/Boundary/Contour.lean\#L28}{\texttt{strict}}
& \(E_n\), \(\Sigma_n\) \\
heights placed as times &
\href{https://github.com/yang-mills-group/OpenAI-YM-20261028/blob/3456e5e75cf8d4cd558955a5c7638fba2d68f464/YM/Shift/Ridge.lean\#L24}{\texttt{dawn}}
& \(\underline y\) \\
continuation; restriction; smearing &
\href{https://github.com/yang-mills-group/OpenAI-YM-20261028/blob/3456e5e75cf8d4cd558955a5c7638fba2d68f464/YM/Boundary/Contour.lean\#L57-L58}{\texttt{Waves}},
\texttt{restricted},
\href{https://github.com/yang-mills-group/OpenAI-YM-20261028/blob/3456e5e75cf8d4cd558955a5c7638fba2d68f464/YM/Shift/Ridge.lean\#L75-L76}{\texttt{psi}}
& \(F_n\), \(F^y_n\), \(\psi_v(y;f)\) \\
strict continuation; integrability; height continuity; Cauchy--Riemann &
\texttt{Region}, \texttt{Bounds}, \texttt{Along}, \texttt{Lines} & (SC),
(I), (HC), (CR) \\
polynomial growth in the inverse time separation & \texttt{Breathes} &
(PG) \\
boundary family (Wightman functions) &
\href{https://github.com/yang-mills-group/OpenAI-YM-20261028/blob/3456e5e75cf8d4cd558955a5c7638fba2d68f464/YM/Boundary/Contour.lean\#L54-L55}{\texttt{Surface}}
& \(W_n(v)\); \(W(b)\) \\
boundary-value hypotheses & \texttt{Opens}, \texttt{NonTangentialFree} &
(BV\(_{\rm ray}\)), (BV) \\
Laplace identity; half-space support & \texttt{during},
\texttt{cone\_tests}, \texttt{half} & (L); \(\mathcal P_n\) \\
Wightman spectral set & \texttt{coneM}, \texttt{cones} &
\(\mathcal V_n\) \\
imaginary-part cones; tube & \texttt{fields}, \texttt{shores},
\texttt{tube} & \(V^+_n\), \(\bar V^+_n\), \(\mathcal T_n\) \\
cutoff; damping properties; partial boundary value & \texttt{Blanket},
\texttt{Modified}, \texttt{partialD} & \(\chi_n\); (D);
\(W_n(v)[\eta]\) \\
Fourier--Laplace continuation &
\href{https://github.com/yang-mills-group/OpenAI-YM-20261028/blob/3456e5e75cf8d4cd558955a5c7638fba2d68f464/YM/Shift/Sol.lean\#L116}{\texttt{Aurora.emit}},
\texttt{FullBody} & \(\mathcal W_n(v)(z)\); (TC) \\
rotation structure; spin structure & \texttt{Rotation} (\(R\),
\texttt{lie}),
\href{https://github.com/yang-mills-group/OpenAI-YM-20261028/blob/3456e5e75cf8d4cd558955a5c7638fba2d68f464/YM/Fields/Star.lean\#L22-L24}{\texttt{Spins}}
(\texttt{rail}, \texttt{shell}) & \(R\), \(\rho\); \(K(\Lambda)\) \\
boost identity; Poincaré covariance & \texttt{Boosty}; \texttt{Urs} &
(B\(_y\)); (Cov) \\
induced star; word star; kernel &
\href{https://github.com/yang-mills-group/OpenAI-YM-20261028/blob/3456e5e75cf8d4cd558955a5c7638fba2d68f464/YM/Fields/Star.lean\#L37}{\texttt{StarOp}},
\texttt{rev}, \texttt{kernel} & \(J\); \(b^{\ast}\);
\(\langle b,c\rangle_W\) \\
adjacent pair; spacelike set; swap &
\href{https://github.com/yang-mills-group/OpenAI-YM-20261028/blob/3456e5e75cf8d4cd558955a5c7638fba2d68f464/YM/Shift/Swap.lean\#L19-L22}{\texttt{NearPair}},
\texttt{gap}, \texttt{switch}/\texttt{name} & \(Q=(i,i+1)\),
\(\Gamma_Q\), \(\sigma\) \\
spectator profile; complex boost & \texttt{spectator}/\texttt{Fleck};
\texttt{augment} & \(y\); \(\Lambda_e(\zeta)\) \\
germs at a spacelike pair &
\href{https://github.com/yang-mills-group/OpenAI-YM-20261028/blob/3456e5e75cf8d4cd558955a5c7638fba2d68f464/YM/Shift/Exchange.lean\#L33-L36}{\texttt{Panes}}
& \(g^{\rm o}_x\), \(g^{\rm t}_x\) \\
heated vector; heat cup & \texttt{shot}, \texttt{heatCup} &
\(\Psi_p(v;y;r)\); \(\Phi_t(b)\) \\
reconstructed theory &
\href{https://github.com/yang-mills-group/OpenAI-YM-20261028/blob/3456e5e75cf8d4cd558955a5c7638fba2d68f464/YM/Fields/System.lean\#L117-L122}{\texttt{Boundary.Wightman}}
& \((\mathcal H_W,\mathcal D,\Omega_W,\varphi,J,U,\mathsf E)\) \\
Kolmogorov space of \(\langle\cdot,\cdot\rangle_W\) &
\href{https://github.com/yang-mills-group/OpenAI-YM-20261028/blob/3456e5e75cf8d4cd558955a5c7638fba2d68f464/YM/Fields/Stems.lean\#L20}{\texttt{Soil}}
& \(\mathcal H_W\) \\
the bundle &
\href{https://github.com/yang-mills-group/OpenAI-YM-20261028/blob/3456e5e75cf8d4cd558955a5c7638fba2d68f464/YM/Bid/Arena.lean\#L15-L45}{\texttt{Bid.Arcs}}
& (A16), (C01)--(C29) \\
\end{longtable}
}

\subsubsection{7. What remains
unexplained}\label{what-remains-unexplained-21}

\begin{enumerate}
\def\labelenumi{\arabic{enumi}.}
\tightlist
\item
  \textbf{Upstream slab estimate.} The constants in (SB) are independent
  of the slab width and word length for fixed degree bound \(D\) and
  representative bound \(R\); they may depend on \(D,R\). VIII.1 obtains
  the mixed and selected slab laws from the phase-data package. The
  upstream proof of those laws, beyond their use in reconstruction, is
  outside this check.
\item
  \textbf{Imported analytic results.} The holomorphic heat-semigroup
  construction and the OS translation spectral measure are used by their
  stated properties. The construction of the compact-model Borel
  functional calculus, and the estimates underlying complex
  differentiation of the Fourier--Laplace tests, have not been checked
  in full. Their exact uses are identified in V.2.b, V.2.f and V.2.p.
\item
  \textbf{Location of the arguments.} The boundary-value construction,
  the spectral-support improvement and the germ exchange lemmas are
  filed under VI.1; the field realization and application to phase data
  are filed under VIII.1. They belong with the reconstruction argument
  in a written proof.
\item
  \textbf{Vacuum uniqueness.} The construction for the mixed table
  supplies the field axioms without vacuum uniqueness. The selected
  phase has uniqueness because its OS heat-fixed space is
  one-dimensional. No uniqueness assertion for the mixed table follows
  here.
\item
  \textbf{References.} The bibliographic pointers in §2 identify the
  relevant standard methods; their theorem numbering has not been
  verified.
\end{enumerate}

\subsection{V.3. Reconstruction in every
phase}\label{v.3.-reconstruction-in-every-phase}

\textbf{Part:} V, phases and reconstruction. \textbf{Lean:}
\href{https://github.com/yang-mills-group/OpenAI-YM-20261028/tree/3456e5e75cf8d4cd558955a5c7638fba2d68f464/YM/EveryPhase}{\texttt{YM/EveryPhase}}
(34 files). Its hypotheses are recorded in
\href{https://github.com/yang-mills-group/OpenAI-YM-20261028/blob/3456e5e75cf8d4cd558955a5c7638fba2d68f464/YM/Spec/Closed.lean\#L20-L73}{\texttt{YM/Spec/Closed.lean}}
and
\href{https://github.com/yang-mills-group/OpenAI-YM-20261028/blob/3456e5e75cf8d4cd558955a5c7638fba2d68f464/YM/Phase/Mixture.lean\#L65-L148}{\texttt{YM/Phase/Mixture.lean}}.
\textbf{Size:} 5,268 lines.

\begin{itemize}
\tightlist
\item
  The 436 declarations occupy about 4,180 lines; 419 of them (about
  4,070 lines) occur in the main theorem's dependency graph.
\item
  About half the folder (15 files, about 2,600 lines) continues the
  stress tensor (§4.7).
\end{itemize}

\subsubsection{1. What this unit proves}\label{what-this-unit-proves-22}

\textbf{Setting.} Fix a Wilson family, an extraction and the limit table
\(S=S^{\rm mix}\), together with a record
\href{https://github.com/yang-mills-group/OpenAI-YM-20261028/blob/3456e5e75cf8d4cd558955a5c7638fba2d68f464/YM/Spec/Closed.lean\#L64-L73}{\texttt{Final}}
for them. The unit uses the following parts of it.

\begin{itemize}
\tightlist
\item
  The Euclidean axioms of \(S^{\rm mix}\) (IV.1).
\item
  The phase decomposition of V.1:

  \begin{itemize}
  \tightlist
  \item
    a compact metric space \(X\) with a Borel probability measure
    \(\nu\) of full support;
  \item
    the commutative C\(^*\)-algebra \(\mathfrak A\) generated by the
    time-slab operators of \(S^{\rm mix}\), compressed to the ground
    space \(\mathcal D\) of the Osterwalder--Schrader (OS) space
    \(\mathcal H^{\rm OS}\) of \(S^{\rm mix}\);
  \item
    continuous characters \(\chi_\xi\) of \(\mathfrak A\), indexed
    bijectively by \(X\), with \(\xi\mapsto\chi_\xi(M)\) continuous for
    every \(M\in\mathfrak A\) and
    \(\int_X\chi_\xi(M)\,d\nu=\langle\omega,M\omega\rangle\);
  \item
    spatial translations acting trivially on \(\mathcal D\);
  \item
    tables \(S^\xi\) satisfying the Euclidean axioms, with
    \(S^\xi(w)=\chi_\xi(M_w)\) for every word \(w\) of pairwise distinct
    points. Here \(M_w\in\mathfrak A\) is the compressed slab operator
    of a time translate of \(w\).
  \end{itemize}
\item
  The relativistic theory \(\mathfrak W_{\xi_*}\) of V.2 at the selected
  phase \(\xi_*\), and the analytic continuations of \(S^{\rm mix}\) and
  \(S^{\xi_*}\).
\item
  The field algebra of VI.1 (derivatives, rotations, the Euclidean
  stress tensor \(T^{\rm E}_{ab}\in F_4\)), and the Euclidean Ward
  identities in every phase.
\end{itemize}

\textbf{Notation.}

\begin{itemize}
\tightlist
\item
  \textbf{OS side.} For each \(\xi\), \(\mathcal H^{\rm OS}_\xi\) is an
  OS space of \(S^\xi\) for the time axis. It carries the transfer
  semigroup \(e^{-sH^{\rm OS}_\xi}\) and the Euclidean vacuum
  \(\omega_\xi\). Its ground space is
  \(\mathcal D_\xi=\{v:\ e^{-sH^{\rm OS}_\xi}v=v\ \text{for all } s\ge0\}\).
\item
  \textbf{Wightman theory.}
  \(\mathfrak W_\xi=(\mathcal H_\xi,\mathfrak D_\xi,\Omega_\xi,\phi_\xi,U_\xi,P_\xi)\)
  consists of:

  \begin{itemize}
  \tightlist
  \item
    a Hilbert space \(\mathcal H_\xi\) and a dense common invariant
    domain \(\mathfrak D_\xi\) spanned by word vectors;
  \item
    a unit vacuum \(\Omega_\xi\);
  \item
    fields \(\phi_\xi(a,f)\) for labels \(a\in\mathbb F\) (the common
    label space generated by the \(F_d\)) and
    \(f\in\mathcal S(\mathbb R^4)\);
  \item
    a strongly continuous unitary Poincaré representation \(U_\xi\) and
    an energy--momentum spectral measure \(P_\xi\).
  \end{itemize}
\item
  \textbf{Time and energy.} Time is the fourth coordinate. \(H_\xi\)
  generates the time translations, and \(P^0_\xi=P_\xi(\{p^4=0\})\).
\item
  \textbf{Words.} A \emph{finite sum of words} is
  \(A=\sum_kc_k\,\phi(a_{k1},f_{k1})\cdots\phi(a_{kn_k},f_{kn_k})\). It
  is \emph{compactly supported} if every
  \(f_{kj}\in C^\infty_c(\mathbb R^4)\). \(A^\dagger\) is its formal
  adjoint (conjugate coefficients and tests, apply the star operation to
  labels, and reverse each word). The spectral measure of
  \(A\Omega_\xi\) is
  \(\mu^\xi_A(B)=\langle A\Omega_\xi,P_\xi(B)A\Omega_\xi\rangle\).
\item
  \textbf{Continuations.} \(\mathsf F^\xi_n(a;z,\vec x)\) is the
  analytic continuation of \(S^\xi\) to complex times with
  \(\mathrm{Im}z_1<\dots<\mathrm{Im}z_n\). \(W^\xi_n(a;g)\), for
  \(g\in\mathcal S((\mathbb R^4)^n)\), are its boundary distributions.
  \(W^{\rm mix}_n\) are those of the continuation of \(S^{\rm mix}\)
  (V.2).
\end{itemize}

\textbf{Proposition V.3.1 (reconstruction in every phase).} There is a
family indexed by \(X\) with the following data for each \(\xi\in X\):

\begin{itemize}
\tightlist
\item
  an OS space \(\mathcal H^{\rm OS}_\xi\) of \(S^\xi\) with its
  time-slab operators;
\item
  a continuation \(\mathsf F^\xi\) of \(S^\xi\) with boundary
  distributions \(W^\xi\);
\item
  a Wightman theory \(\mathfrak W_\xi\) on the label space
  \(\mathbb F\), with the same star operation and the same Poincaré
  action on labels as \(\mathfrak W_{\xi_*}\);
\item
  an isometry \(J_\xi:\mathcal H_\xi\to\mathcal H^{\rm OS}_\xi\) with
  \(J_\xi\Omega_\xi=\omega_\xi\). It intertwines the space--time
  translations and sends word vectors to the boundary values of the
  corresponding Euclidean vectors at imaginary times (Lemma V.3.b).
\end{itemize}

At \(\xi=\xi_*\) these are the data of V.1 and V.2. For \textbf{every}
\(\xi\in X\), with no exceptional set, the following hold.

\begin{enumerate}
\def\labelenumi{\arabic{enumi}.}
\item
  \emph{Wightman axioms.}

  \begin{itemize}
  \tightlist
  \item
    The vacuum expectation values of products of fields are \(W^\xi\).
  \item
    The Hilbert-space adjoint of \(\phi_\xi(a,f)\) extends
    \(\phi_\xi(a^*,\bar f)\).
  \item
    Fields with spacelike separated supports commute on
    \(\mathfrak D_\xi\).
  \item
    \(U_\xi\) implements Poincaré covariance.
  \item
    The joint spectrum of the translations lies in the closed forward
    cone \(\overline V_+\).
  \item
    For every nonzero label \(a\), there is a test function \(f\) for
    which \(\phi_\xi(a,f)\ne0\) on \(\mathfrak D_\xi\).
  \item
    Every translation-invariant vector is a multiple of \(\Omega_\xi\).
  \end{itemize}
\item
  \emph{Surjectivity.} \(J_\xi\) is unitary.
\item
  \emph{Euclidean vacuum.} \(\mathcal D_\xi=\mathbb C\omega_\xi\).
\item
  \emph{Zero energy.} For every \(x\in\mathcal H_\xi\), \(P^0_\xi x=x\)
  if and only if \(x\in\mathbb C\Omega_\xi\).
\item
  \emph{Low-energy bound.} For every compactly supported finite sum of
  words \(A\) there is \(M\ge0\) such that
  \(\mu^\xi_A(\{0<p^4\le u\})\le Mu^2\) for all \(\xi\in X\) and
  \(u\in(0,1]\).
\item
  \emph{Heat decay.} For every such \(A\) there is \(M\ge0\) such that,
  for all \(\xi\in X\) and \(s\ge1\),
  \[0\le\big\langle A\Omega_\xi,(e^{-sH_\xi}-P^0_\xi)A\Omega_\xi\big\rangle\le Ms^{-2}.\]
\item
  \emph{Irreducibility.} For a bounded complex-linear operator \(B\) on
  \(\mathcal H_\xi\), the following are equivalent:
  \(B\in\mathbb C\,1\); and, for every \(a\), \(f\) and
  \(u,v\in\mathfrak D_\xi\), one has
  \(u\in\mathrm{dom}\,\phi_\xi(a,f)^*\) and
  \(\langle\phi_\xi(a,f)^*u,Bv\rangle=\langle u,B\phi_\xi(a,f)v\rangle\).
\item
  \emph{Stress tensor.} The conclusions (T0)--(T6) of Proposition V.3.y
  hold for \(\mathfrak W_\xi\).
\item
  \emph{Dependence on the phase.}

  \begin{itemize}
  \tightlist
  \item
    For all \(n\), \(a\in\mathbb F^n\) and
    \(g\in\mathcal S((\mathbb R^4)^n)\), the map
    \(\xi\mapsto W^\xi_n(a;g)\) is continuous and
    \(\int_XW^\xi_n(a;g)\,d\nu(\xi)=W^{\rm mix}_n(a;g)\).
  \item
    For every word \(w\) of pairwise distinct points,
    \(\int_XS^\xi(w)\,d\nu(\xi)=S^{\rm mix}(w)\).
  \end{itemize}
\end{enumerate}

In item 6, \(e^{-sH_\xi}\) is defined as
\(J_\xi^*e^{-sH^{\rm OS}_\xi}J_\xi\). Since \(J_\xi\) intertwines the
spectral projections (V.3.b), this equals \(\int e^{-sp^4}\,dP_\xi(p)\).
The projection subtracted in item 6 is the full zero-energy projection.
Items 5 and 6 are therefore proved before item 4 and independently of
it.

\emph{Status:} theorem, for an arbitrary record
\href{https://github.com/yang-mills-group/OpenAI-YM-20261028/blob/3456e5e75cf8d4cd558955a5c7638fba2d68f464/YM/Spec/Closed.lean\#L64-L73}{\texttt{Final}}.
The record is constructed in the proof of the main theorem (VIII.1), so
the statement is unconditional there.

\begin{itemize}
\tightlist
\item
  The detailed main theorem states V.3.1 as its clause (D7), via
  \href{https://github.com/yang-mills-group/OpenAI-YM-20261028/blob/3456e5e75cf8d4cd558955a5c7638fba2d68f464/YM/Independent/Core.lean\#L18-L34}{\texttt{Independent.Core}}.
\item
  The public theorem uses the Euclidean barycenter formula of item 9 for
  clause (R12), via
  \href{https://github.com/yang-mills-group/OpenAI-YM-20261028/blob/3456e5e75cf8d4cd558955a5c7638fba2d68f464/HierarchyBridge.lean\#L79-L133}{\texttt{publicReconstruction}}.
\item
  Item 1 and the averaged Wightman distributions of item 9 enter
  \href{https://github.com/yang-mills-group/OpenAI-YM-20261028/blob/3456e5e75cf8d4cd558955a5c7638fba2d68f464/YM/Nonperturbative/HeadlineSpec.lean\#L75-L81}{\texttt{Headline.Reconstruction}}.
\item
  No unit uses items 2--8 as hypotheses.
\end{itemize}

\emph{Lean:}
\href{https://github.com/yang-mills-group/OpenAI-YM-20261028/blob/3456e5e75cf8d4cd558955a5c7638fba2d68f464/YM/EveryPhase/Extension.lean\#L75-L80}{\texttt{all\_characters}},
\href{https://github.com/yang-mills-group/OpenAI-YM-20261028/blob/3456e5e75cf8d4cd558955a5c7638fba2d68f464/YM/EveryPhase/Extension.lean\#L51-L71}{\texttt{Completed}}
(the statement; see §6 for the fields),
\href{https://github.com/yang-mills-group/OpenAI-YM-20261028/blob/3456e5e75cf8d4cd558955a5c7638fba2d68f464/YM/EveryPhase/Extension.lean\#L19-L22}{\texttt{member}}
(weak commutant),
\href{https://github.com/yang-mills-group/OpenAI-YM-20261028/blob/3456e5e75cf8d4cd558955a5c7638fba2d68f464/YM/EveryPhase/Pack.lean\#L117-L132}{\texttt{Almanac}}
(the data),
\href{https://github.com/yang-mills-group/OpenAI-YM-20261028/blob/3456e5e75cf8d4cd558955a5c7638fba2d68f464/YM/EveryPhase/Pack.lean\#L157-L194}{\texttt{almanac}}
(existence).

\textbf{Proposition V.3.2 (no exceptional phase).} For every
\(\xi\in X\) and \(x\in\mathcal H_\xi\),
\[P^0_\xi x=\langle\Omega_\xi,x\rangle\,\Omega_\xi.\] \emph{Status:}
theorem (Lemma V.3.r). Items 3, 4 and 7 of V.3.1 are derived from it.
\emph{Lean:}
\href{https://github.com/yang-mills-group/OpenAI-YM-20261028/blob/3456e5e75cf8d4cd558955a5c7638fba2d68f464/YM/EveryPhase/Vacuum.lean\#L115-L144}{\texttt{all\_plant}},
\href{https://github.com/yang-mills-group/OpenAI-YM-20261028/blob/3456e5e75cf8d4cd558955a5c7638fba2d68f464/YM/EveryPhase/Vacuum.lean\#L158-L162}{\texttt{boundary\_zero\_energy}}.

\textbf{Remark V.3.3 (recorded consequences, not used by the main
theorem).} Under V.3.1, for every \(\xi\in X\) and
\(x\in\mathcal H_\xi\):

\begin{itemize}
\tightlist
\item
  \begin{enumerate}
  \def\labelenumi{(\alph{enumi})}
  \tightlist
  \item
    \(e^{-sH_\xi}x=x\) for all \(s\ge0\) if and only if
    \(x\in\mathbb C\Omega_\xi\);
  \end{enumerate}
\item
  (b) \(x\) is invariant under all translations if and only if
  \(x\in\mathbb C\Omega_\xi\);
\item
  (c) for every degree \(d\) and degree-\(d\) token coefficient vector
  \(v\), \(v\) is null for \(S^\xi\) if and only if the field associated
  to its class in \(F_d\) vanishes on every test function;
\item
  (d) reflection positivity holds for every affine hyperplane. More
  precisely, let \(O\in O(4)\), \(b\in\mathbb R^4\) and
  \(a\in\{1,2,3,4\}\). For any finite family of words \(w_i\) with
  distinct points in \(\{x:(O^{-1}(x-b))^a>0\}\),
  \[\sum_{i,j}\overline{c_i}c_j\,S^\xi(\theta_{O,b,a}w_i\,w_j)\ge0\qquad(c_i\in\mathbb C).\]
  Here \(\theta_{O,b,a}\) reverses the word, reflects its points across
  the hyperplane, and acts on labels by transporting the \(a\)-axis
  reflection through \(O\).
\end{itemize}

\emph{Lean:}
\href{https://github.com/yang-mills-group/OpenAI-YM-20261028/blob/3456e5e75cf8d4cd558955a5c7638fba2d68f464/YM/EveryPhase/Extension.lean\#L100-L102}{\texttt{fixedHeat}},
\href{https://github.com/yang-mills-group/OpenAI-YM-20261028/blob/3456e5e75cf8d4cd558955a5c7638fba2d68f464/YM/EveryPhase/Extension.lean\#L103-L106}{\texttt{exactly\_fixed}},
\href{https://github.com/yang-mills-group/OpenAI-YM-20261028/blob/3456e5e75cf8d4cd558955a5c7638fba2d68f464/YM/EveryPhase/Extension.lean\#L90-L98}{\texttt{recognizes}},
\href{https://github.com/yang-mills-group/OpenAI-YM-20261028/blob/3456e5e75cf8d4cd558955a5c7638fba2d68f464/YM/EveryPhase/Extension.lean\#L84-L88}{\texttt{allPlanes}}.

\subsubsection{2. What it uses}\label{what-it-uses-22}

\textbf{From other units.}

\begin{itemize}
\tightlist
\item
  \textbf{IV.1 (infinite volume).} The Euclidean axioms of
  \(S^{\rm mix}\): translation homogeneity, reflection positivity,
  smoothness. Lean:
  \href{https://github.com/yang-mills-group/OpenAI-YM-20261028/blob/3456e5e75cf8d4cd558955a5c7638fba2d68f464/YM/Spec/Closed.lean\#L20-L26}{\texttt{Seed}}.
\item
  \textbf{V.1 (decomposition into phases).} The data:

  \begin{itemize}
  \tightlist
  \item
    the triple \((X,\nu,\chi)\) with the properties of §1:
    \href{https://github.com/yang-mills-group/OpenAI-YM-20261028/blob/3456e5e75cf8d4cd558955a5c7638fba2d68f464/YM/Phase/Mixture.lean\#L65-L72}{\texttt{Resolution}},
    \href{https://github.com/yang-mills-group/OpenAI-YM-20261028/blob/3456e5e75cf8d4cd558955a5c7638fba2d68f464/YM/Phase/Mixture.lean\#L78-L86}{\texttt{Resolution.Out}},
    \href{https://github.com/yang-mills-group/OpenAI-YM-20261028/blob/3456e5e75cf8d4cd558955a5c7638fba2d68f464/YM/Phase/Mixture.lean\#L89-L91}{\texttt{Resolution.receives}};
  \item
    for every \(\xi\): the Euclidean axioms of \(S^\xi\), null spaces
    equal to those of \(S^{\rm mix}\), the inclusion, derivative and
    rotation laws, and the Euclidean Ward identities of the stress
    tensor;
  \item
    the slab laws of \(S^{\rm mix}\), including its copy bounds (below):
    \href{https://github.com/yang-mills-group/OpenAI-YM-20261028/blob/3456e5e75cf8d4cd558955a5c7638fba2d68f464/YM/Phase/Mixture.lean\#L127-L134}{\texttt{Selection}},
    \href{https://github.com/yang-mills-group/OpenAI-YM-20261028/blob/3456e5e75cf8d4cd558955a5c7638fba2d68f464/YM/Phase/Mixture.lean\#L136-L148}{\texttt{Selection.Certificate}}.
  \end{itemize}
\item
  \textbf{V.1, tools.}

  \begin{itemize}
  \tightlist
  \item
    The OS space of a table:
    \href{https://github.com/yang-mills-group/OpenAI-YM-20261028/blob/3456e5e75cf8d4cd558955a5c7638fba2d68f464/YM/OS/Dynamics.lean\#L106-L107}{\texttt{OS.realization}}.
  \item
    The multiple-reflection bound for slab operators:
    \href{https://github.com/yang-mills-group/OpenAI-YM-20261028/blob/3456e5e75cf8d4cd558955a5c7638fba2d68f464/YM/OS/Train.lean\#L116-L123}{\texttt{OS.producer}}.
  \item
    A countable family of positive-time words whose OS vectors are total
    in the OS space of every symmetric smooth table:
    \href{https://github.com/yang-mills-group/OpenAI-YM-20261028/blob/3456e5e75cf8d4cd558955a5c7638fba2d68f464/YM/Work/B14.lean\#L119-L143}{\texttt{Work.B14.universal}}.
  \item
    Continuity of OS vectors in the points:
    \href{https://github.com/yang-mills-group/OpenAI-YM-20261028/blob/3456e5e75cf8d4cd558955a5c7638fba2d68f464/YM/Work/B14.lean\#L99-L107}{\texttt{Work.B14.plant\_cont}}.
  \item
    The full slab laws from the copy bounds:
    \href{https://github.com/yang-mills-group/OpenAI-YM-20261028/blob/3456e5e75cf8d4cd558955a5c7638fba2d68f464/YM/Work/B02.lean\#L17-L30}{\texttt{Work.B02.answer}}.
  \item
    Almost-everywhere uniqueness of the Euclidean vacuum (Lemma V.3.c):
    \href{https://github.com/yang-mills-group/OpenAI-YM-20261028/blob/3456e5e75cf8d4cd558955a5c7638fba2d68f464/YM/Phase/Unique.lean\#L91-L110}{\texttt{Cluster.lifts}}.
  \item
    The Euclidean barycenter formula:
    \href{https://github.com/yang-mills-group/OpenAI-YM-20261028/blob/3456e5e75cf8d4cd558955a5c7638fba2d68f464/YM/Phase/Mixture.lean\#L113-L117}{\texttt{Resolution.barycenter}}.
  \end{itemize}
\item
  \textbf{V.2 (OS reconstruction).}

  \begin{itemize}
  \tightlist
  \item
    The continuations of \(S^{\rm mix}\) and of \(S^{\xi_*}\), and the
    theory \(\mathfrak W_{\xi_*}\) with all Wightman axioms:
    \href{https://github.com/yang-mills-group/OpenAI-YM-20261028/blob/3456e5e75cf8d4cd558955a5c7638fba2d68f464/YM/Spec/Closed.lean\#L50-L54}{\texttt{Relativistic}}.
  \item
    Compatibility of both continuations with the slab operators, which
    includes the polynomial bound on one-point slab operators used in
    V.3.d:
    \href{https://github.com/yang-mills-group/OpenAI-YM-20261028/blob/3456e5e75cf8d4cd558955a5c7638fba2d68f464/YM/EveryPhase/Compatibility.lean\#L26-L29}{\texttt{onFinal}}.
  \item
    The estimates of work items C01--C29 on the continuation and its
    partial boundary values (used in V.3.d, V.3.g, V.3.i).
  \end{itemize}
\item
  \textbf{Reconstruction machinery from V.2:}

  \begin{itemize}
  \tightlist
  \item
    Wightman data lacking only vacuum uniqueness:
    \href{https://github.com/yang-mills-group/OpenAI-YM-20261028/blob/3456e5e75cf8d4cd558955a5c7638fba2d68f464/YM/Next/BOUNDARY/Core.lean\#L34-L44}{\texttt{FieldCore}}.
  \item
    The boundary isometry and its properties:
    \href{https://github.com/yang-mills-group/OpenAI-YM-20261028/blob/3456e5e75cf8d4cd558955a5c7638fba2d68f464/YM/Next/BOUNDARY/Core.lean\#L156-L161}{\texttt{Attachment}},
    \href{https://github.com/yang-mills-group/OpenAI-YM-20261028/blob/3456e5e75cf8d4cd558955a5c7638fba2d68f464/YM/Next/BOUNDARY/Core.lean\#L198-L204}{\texttt{Attachment.measures}},
    \href{https://github.com/yang-mills-group/OpenAI-YM-20261028/blob/3456e5e75cf8d4cd558955a5c7638fba2d68f464/YM/Next/BOUNDARY/Core.lean\#L207-L246}{\texttt{Attachment.projections}},
    \href{https://github.com/yang-mills-group/OpenAI-YM-20261028/blob/3456e5e75cf8d4cd558955a5c7638fba2d68f464/YM/Next/BOUNDARY/Core.lean\#L191-L194}{\texttt{Attachment.with\_heat}}.
  \item
    The constructions run at each phase:
    \href{https://github.com/yang-mills-group/OpenAI-YM-20261028/blob/3456e5e75cf8d4cd558955a5c7638fba2d68f464/YM/Next/BOUNDARY/Stage.lean\#L46-L49}{\texttt{contour}},
    \href{https://github.com/yang-mills-group/OpenAI-YM-20261028/blob/3456e5e75cf8d4cd558955a5c7638fba2d68f464/YM/Next/BOUNDARY/Stage.lean\#L200-L215}{\texttt{compatible}},
    \href{https://github.com/yang-mills-group/OpenAI-YM-20261028/blob/3456e5e75cf8d4cd558955a5c7638fba2d68f464/YM/Next/BOUNDARY/Stage.lean\#L167-L184}{\texttt{establishing}}.
  \end{itemize}
\item
  \textbf{VI.1 (Ward identities, rotations, stress tensor).} Derivative
  plan, rotations and Euclidean stress tensor:
  \href{https://github.com/yang-mills-group/OpenAI-YM-20261028/blob/3456e5e75cf8d4cd558955a5c7638fba2d68f464/YM/Spec/Closed.lean\#L28-L31}{\texttt{Algebra}}.
  Their laws (conservation and symmetry of \(T^{\rm E}\), tensor
  covariance, rotation laws):
  \href{https://github.com/yang-mills-group/OpenAI-YM-20261028/blob/3456e5e75cf8d4cd558955a5c7638fba2d68f464/YM/Spec/Closed.lean\#L34-L43}{\texttt{Locals}}.
\end{itemize}

\textbf{Standard mathematics.}

\begin{itemize}
\tightlist
\item
  OS reconstruction: Osterwalder--Schrader, CMP 31 (1973) and CMP 42
  (1975); Glimm--Jaffe, \emph{Quantum Physics}, Ch. 6 and 19.
\item
  Boundedness of transfer operators by iterated reflection and the
  Cauchy--Schwarz inequality: Glimm--Jaffe, Ch. 10;
  Fröhlich--Israel--Lieb--Simon, CMP 62 (1978).
\item
  Decomposition of a state along an abelian algebra of the vacuum
  sector, used in V.1 and in V.3.c: Bratteli--Robinson, vol.~1, Ch. 4.
\item
  Spectral theorem for unitary representations of \(\mathbb R^4\),
  Bochner's theorem, Laplace transforms of positive measures:
  Reed--Simon I--II.
\item
  Boundary values of holomorphic functions of polynomial growth and
  their Taylor representation: Streater--Wightman, Ch. 2; Reed--Simon
  II, §IX.3.
\item
  Complex analysis: Vitali--Montel theorem, identity theorem, Morera's
  theorem, Liouville's theorem.
\item
  Uniqueness of the vacuum and irreducibility of the field algebra:
  Streater--Wightman, Ch. 4; Jost, \emph{The General Theory of Quantized
  Fields}.
\end{itemize}

\textbf{Imported constants and assumptions.} The record
\href{https://github.com/yang-mills-group/OpenAI-YM-20261028/blob/3456e5e75cf8d4cd558955a5c7638fba2d68f464/YM/Spec/Closed.lean\#L64-L73}{\texttt{Final}}
is the only hypothesis. Its quantitative inputs are as follows.

\begin{itemize}
\item
  \textbf{(i) Copy bounds of \(S^{\rm mix}\)} (V.1, from the chessboard
  estimates). For every degree bound \(D\) and coefficient bound \(R\)
  there are \(C>0\) and \(Q\in\mathbb N\) with
  \[\|\Pi(d,w)\|\le C^{|w|}(|w|/r)^{Q|w|}.\] This holds for every word
  \(w\) of \(|w|\) points with degrees \(\le D\) and coefficients
  \(\le R\), lying in the slab \([r,d-r]\), with points \(r\)-separated
  and \(0<r\le1\).
\item
  \textbf{(ii) One-point slab operators} (V.2). For each label
  \(a\in\mathbb F\) there are \(Q_a>0\) and \(N_a\in\mathbb N\) such
  that the one-point slab operator of \(S^{\rm mix}\), of thickness
  \(2p\) with \(a\) at \((\vec0,p)\), has norm at most \(Q_a p^{-N_a}\)
  for \(0<p<1\). The exponent \(N_a\) is a bound, not an asserted
  scaling dimension.
\end{itemize}

The one-point bound is
\href{https://github.com/yang-mills-group/OpenAI-YM-20261028/blob/3456e5e75cf8d4cd558955a5c7638fba2d68f464/YM/Fields/Tall.lean\#L62-L63}{\texttt{Boundary.Small.rated}}.

The unit's own constants are explicit:

\begin{itemize}
\tightlist
\item
  \(R=4r+1\) in V.3.m;
\item
  \(c(R,M)=16\pi\,|B_1|\,|B_R|\,M/c_3\) in V.3.k, where \(|B_r|\) is the
  volume of the ball of radius \(r\) in \(\mathbb R^3\) and
  \(c_3=(2\pi)^3\) for the Fourier convention in V.3.k;
\item
  the factor \(2\) in V.3.l;
\item
  \(\Gamma(3)=2\) in V.3.o.
\end{itemize}

The Fourier convention and constant are
\href{https://github.com/yang-mills-group/OpenAI-YM-20261028/blob/3456e5e75cf8d4cd558955a5c7638fba2d68f464/YM/Distribution/Tests.lean\#L84-L90}{\texttt{Boundary.Probe.hat\_eq}},
\href{https://github.com/yang-mills-group/OpenAI-YM-20261028/blob/3456e5e75cf8d4cd558955a5c7638fba2d68f464/YM/Distribution/Tests.lean\#L48}{\texttt{Boundary.Probe.pref}}.

\subsubsection{3. Strategy}\label{strategy-22}

V.2 reconstructs a relativistic theory at a selected phase \(\xi_*\).
The same construction applies to every \(S^\xi\) before vacuum
uniqueness is known. The norm of each phase's slab operator is bounded
by the corresponding norm for \(S^{\rm mix}\) (V.3.a): repeated
reflection gives a power bound, and the fixed norm of the character
disappears on taking roots. This gives polynomial continuation bounds
uniform in \(\xi\). Character continuity at imaginary times propagates
to complex times by Vitali's theorem, then to boundary distributions by
a Taylor formula with an integrable remainder. Integration over phases
commutes with both passages (V.3.d--V.3.h).

Vacuum uniqueness holds almost everywhere by the phase decomposition
(V.3.c). To extend it to every phase, fix a compactly supported finite
sum of words \(A\) and set
\[r_\xi(s)=\langle A\Omega_\xi,e^{-sH_\xi}A\Omega_\xi\rangle-|\langle\Omega_\xi,A\Omega_\xi\rangle|^2.\]
Locality makes the vacuum commutator vanish at large spatial separation
for \(|t|\le1\). Its time-averaged odd part is the spatial Fourier
transform of a positive measure, weighted by \((1-\cos E)/E\). A
momentum bump of radius \(u\) bounds its shell mass by \(O(u^2)\): the
bump contributes \(u^3\), while the weight contributes \(u\). Dyadic
summation and a Laplace transform give decay \(Ms^{-2}\), uniformly in
\(\xi\), after subtracting the full zero-energy projection
(V.3.k--V.3.o). Almost everywhere this is the vacuum projection. Since
\(r_\xi(s)\) is continuous, the same estimate holds everywhere by full
support of \(\nu\). Letting \(s\to\infty\) identifies the zero-energy
projection with the vacuum projection at every phase (V.3.q--V.3.r).
This needs no continuity of the zero-energy projection itself.

Separately, positive-height Euclidean vectors are images of Wightman
vectors with exponentially weighted frequency tests, so \(J_\xi\) is
onto (V.3.i--V.3.j). Vacuum uniqueness therefore holds on both sides. A
Liouville argument gives irreducibility. Finally, continuation of the
Euclidean Ward identities gives the local stress-charge commutator
(V.3.u--V.3.y). The removal of the exceptional set uses uniform
low-energy control from locality, rather than a mass gap.

\subsubsection{4. Steps}\label{steps-22}

\paragraph{4.1 Slab bounds and the family of
reconstructions}\label{slab-bounds-and-the-family-of-reconstructions}

\textbf{Lemma V.3.a (slab operators of a phase are dominated).} Fix
\(\xi\in X\). Let \(\mathcal K\) be any OS space of \(S^\xi\) for the
time axis, and \(\Pi_\xi\) any family of slab operators on it. Then for
every thickness \(d>0\) and every word \(w\) of pairwise distinct points
in the slab \((0,d)\), \[\|\Pi_\xi(d,w)\|\le\|\Pi(d,w)\|,\] where
\(\Pi\) are the slab operators of \(S^{\rm mix}\) on
\(\mathcal H^{\rm OS}\). Here \(\Pi(d,w)\) maps the OS vector of a
positive-time word \(v\) to that of \(w\) followed by \(v\) translated
by \(d\).

\emph{Sketch.}

\begin{enumerate}
\def\labelenumi{\arabic{enumi}.}
\item
  Let \(u=\Pi(d,w)\) and \(D>\|u\|\). For a positive-time word \(v\),
  let \(X(v)\) be the slab operator of \(v\) restricted to the ground
  space \(\mathcal D\).
\item
  The compression of the doubled word \(\theta v\,v\) is
  \(X(v)^*X(v)\in\mathfrak A\). Let \(r_n\) be \(v\) with \(n\)
  reflected pairs of copies of \(w\) inserted. Then
  \(X(r_n)=(u^*u)^nX(v)\).
\item
  Hence the diagonal Gram entry of \(r_n\) in \(S^\xi\) satisfies
  \[|S^\xi(\theta r_n\,r_n)|=|\chi_\xi(X(r_n)^*X(r_n))|\le\|\chi_\xi\|\,\|X(v)\|^2D^{4n}.\]
\item
  The power argument (iterated Cauchy--Schwarz on the span of OS
  vectors) turns this into an operator on \(\mathcal K\) of norm at most
  \(D\) that inserts \(w\). The fixed factor
  \(\|\chi_\xi\|\,\|X(v)\|^2\) disappears in the limit.
\item
  Slab operators are determined by their action on OS vectors, so this
  operator is \(\Pi_\xi(d,w)\). Finally let \(D\downarrow\|u\|\).
\end{enumerate}

\emph{Lean:}
\href{https://github.com/yang-mills-group/OpenAI-YM-20261028/blob/3456e5e75cf8d4cd558955a5c7638fba2d68f464/YM/EveryPhase/Pack.lean\#L95-L100}{\texttt{sameBound}},
with
\href{https://github.com/yang-mills-group/OpenAI-YM-20261028/blob/3456e5e75cf8d4cd558955a5c7638fba2d68f464/YM/EveryPhase/Pack.lean\#L32-L50}{\texttt{x\_train}},
\href{https://github.com/yang-mills-group/OpenAI-YM-20261028/blob/3456e5e75cf8d4cd558955a5c7638fba2d68f464/YM/EveryPhase/Pack.lean\#L53-L74}{\texttt{buys}},
\href{https://github.com/yang-mills-group/OpenAI-YM-20261028/blob/3456e5e75cf8d4cd558955a5c7638fba2d68f464/YM/EveryPhase/Pack.lean\#L76-L91}{\texttt{actual}},
and
\href{https://github.com/yang-mills-group/OpenAI-YM-20261028/blob/3456e5e75cf8d4cd558955a5c7638fba2d68f464/YM/OS/Train.lean\#L116-L123}{\texttt{OS.producer}}.

\textbf{Lemma V.3.b (the family of reconstructions).} For every \(\xi\)
there are:

\begin{itemize}
\tightlist
\item
  \begin{enumerate}
  \def\labelenumi{(\roman{enumi})}
  \tightlist
  \item
    an OS space \(\mathcal H^{\rm OS}_\xi\) of \(S^\xi\) with slab
    operators \(\Pi_\xi\) satisfying the slab laws (copy bounds,
    smoothness in the positions, derivative rule), with the copy
    constants of \(S^{\rm mix}\), hence uniformly in \(\xi\);
  \end{enumerate}
\item
  (ii) a continuation \(\mathsf F^\xi\) of \(S^\xi\) with non-tangential
  boundary values \(W^\xi\);
\item
  (iii) a Wightman theory \(\mathfrak W_\xi\) with all properties of
  V.3.1(1) except uniqueness of the vacuum, with the star operation and
  label action of \(\mathfrak W_{\xi_*}\);
\item
  (iv) an isometry \(J_\xi\) with \(J_\xi\Omega_\xi=\omega_\xi\) and
  \(J_\xi U_\xi(a)=U^{\rm OS}_\xi(a)J_\xi\) for \(a\in\mathbb R^4\),
  where \(U^{\rm OS}_\xi\) is the OS translation group. It is uniquely
  determined by the following boundary property: Euclidean word vectors
  at positive imaginary heights, smeared with a joint Schwartz test
  \(g\), converge in Hilbert norm to \(J_\xi\Phi_\xi(a;g)\) as the
  heights and slab widths tend to zero with the positive margins of
  V.3.i. Here \(\Phi_\xi\) is the continuous extension of word vectors
  to joint tests.
\end{itemize}

At \(\xi_*\) these are the data of V.1--V.2. Moreover:

\begin{itemize}
\tightlist
\item
  \(J_\xi P_\xi(B)=P^{\rm OS}_\xi(B)J_\xi\) for every Borel set
  \(B\subset\mathbb R^4\), so \(\mu^\xi_x\) is the spectral measure of
  \(J_\xi x\);
\item
  if \(\mathcal D_\xi=\mathbb C\omega_\xi\), then \(\mathfrak W_\xi\)
  satisfies all of V.3.1(1).
\end{itemize}

\emph{Sketch.}

\begin{itemize}
\tightlist
\item
  For \(\xi\ne\xi_*\), the OS space is built from \(S^\xi\) and the slab
  operators are those of V.3.a.
\item
  Their copy bounds are those of \(S^{\rm mix}\), by V.3.a; the
  remaining slab laws follow from the copy bounds.
\item
  The continuation, the compatibility data and a theory without vacuum
  uniqueness come from the reconstruction machinery of V.2. The boundary
  property in (iv) determines the isometry on a dense span of word
  vectors.
\item
  At \(\xi_*\) the existing data are taken, so the construction does not
  alter V.2.
\item
  The intertwining of the spectral projections follows from intertwining
  the translations, by comparing characteristic functions.
\item
  The last assertion holds because a translation-invariant \(x\) has
  \(J_\xi x\in\mathcal D_\xi\) and \(J_\xi\) is injective.
\end{itemize}

\emph{Lean:}
\href{https://github.com/yang-mills-group/OpenAI-YM-20261028/blob/3456e5e75cf8d4cd558955a5c7638fba2d68f464/YM/EveryPhase/Pack.lean\#L134-L155}{\texttt{storeys}},
\href{https://github.com/yang-mills-group/OpenAI-YM-20261028/blob/3456e5e75cf8d4cd558955a5c7638fba2d68f464/YM/EveryPhase/Pack.lean\#L157-L194}{\texttt{almanac}},
\href{https://github.com/yang-mills-group/OpenAI-YM-20261028/blob/3456e5e75cf8d4cd558955a5c7638fba2d68f464/YM/EveryPhase/Pack.lean\#L196-L199}{\texttt{Almanac.priced}},
\href{https://github.com/yang-mills-group/OpenAI-YM-20261028/blob/3456e5e75cf8d4cd558955a5c7638fba2d68f464/YM/Next/BOUNDARY/Core.lean\#L207-L246}{\texttt{Attachment.projections}},
\href{https://github.com/yang-mills-group/OpenAI-YM-20261028/blob/3456e5e75cf8d4cd558955a5c7638fba2d68f464/YM/Next/BOUNDARY/Core.lean\#L191-L194}{\texttt{Attachment.with\_heat}}.
\emph{(proof read for these;
\href{https://github.com/yang-mills-group/OpenAI-YM-20261028/blob/3456e5e75cf8d4cd558955a5c7638fba2d68f464/YM/Next/BOUNDARY/Stage.lean\#L46-L49}{\texttt{contour}},
\href{https://github.com/yang-mills-group/OpenAI-YM-20261028/blob/3456e5e75cf8d4cd558955a5c7638fba2d68f464/YM/Next/BOUNDARY/Stage.lean\#L200-L215}{\texttt{compatible}},
\href{https://github.com/yang-mills-group/OpenAI-YM-20261028/blob/3456e5e75cf8d4cd558955a5c7638fba2d68f464/YM/Next/BOUNDARY/Stage.lean\#L167-L184}{\texttt{establishing}}
and
\href{https://github.com/yang-mills-group/OpenAI-YM-20261028/blob/3456e5e75cf8d4cd558955a5c7638fba2d68f464/YM/Work/B02.lean\#L17-L30}{\texttt{Work.B02.answer}}
also read; underlying analytic inputs from V.2)}

\textbf{Lemma V.3.c (unique Euclidean vacuum for almost every phase;
from V.1).} For \(\nu\)-almost every \(\xi\),
\(\mathcal D_\xi=\mathbb C\omega_\xi\).

\emph{Sketch (a variance identity).} For a positive-time word \(b\), let
\(\Psi_\xi(b)\in\mathcal H^{\rm OS}_\xi\) be its OS vector and \(Q_\xi\)
the projection onto \(\mathcal D_\xi\).

\begin{enumerate}
\def\labelenumi{\arabic{enumi}.}
\tightlist
\item
  \(\|Q_\xi\Psi_\xi(b)\|^2=\lim_{s\to\infty}S^\xi(\theta b_s\,b_s)\),
  where \(b_s\) is \(b\) translated by \(s\) in time. The terms are
  character values, continuous in \(\xi\), nonincreasing in \(s\) and
  dominated by the value at \(s=0\).
\item
  By dominated convergence and the barycenter identity,
  \(\int\|Q_\xi\Psi_\xi(b)\|^2d\nu=\|Q\Psi(b)\|^2\).
\item
  On the other hand \(S^\xi(b)=\chi_\xi(M_b)\) and
  \(M_b\omega=Q\Psi(b)\). Multiplicativity of characters gives
  \(\int|S^\xi(b)|^2d\nu=\int\chi_\xi(M_b^*M_b)\,d\nu=\|Q\Psi(b)\|^2\).
\item
  Hence the nonnegative function
  \(\|Q_\xi\Psi_\xi(b)\|^2-|S^\xi(b)|^2=\|Q_\xi\Psi_\xi(b)-S^\xi(b)\omega_\xi\|^2\)
  has integral zero, so \(Q_\xi\Psi_\xi(b)\in\mathbb C\omega_\xi\)
  almost everywhere.
\item
  Let \(b\) run over a countable family of words whose OS vectors are
  total in every OS space
  (\href{https://github.com/yang-mills-group/OpenAI-YM-20261028/blob/3456e5e75cf8d4cd558955a5c7638fba2d68f464/YM/Work/B14.lean\#L119-L143}{\texttt{Work.B14.universal}}).
\end{enumerate}

\emph{Lean:}
\href{https://github.com/yang-mills-group/OpenAI-YM-20261028/blob/3456e5e75cf8d4cd558955a5c7638fba2d68f464/YM/EveryPhase/Pack.lean\#L208-L212}{\texttt{Almanac.almostEvery}},
\href{https://github.com/yang-mills-group/OpenAI-YM-20261028/blob/3456e5e75cf8d4cd558955a5c7638fba2d68f464/YM/Phase/Unique.lean\#L91-L110}{\texttt{Cluster.lifts}},
\href{https://github.com/yang-mills-group/OpenAI-YM-20261028/blob/3456e5e75cf8d4cd558955a5c7638fba2d68f464/YM/Phase/Unique.lean\#L56-L86}{\texttt{Cluster.exp\_real}}.

\paragraph{4.2 Continuity in the phase and
barycenter}\label{continuity-in-the-phase-and-barycenter}

\textbf{Lemma V.3.d (polynomial bound, uniform in the phase).} For every
\(n\) and \(a\in\mathbb F^n\) there are \(m\in\mathbb N\) and
\(C_0\ge0\) with the following property. For all \(\xi\), all real
configurations \(x\), and all heights \(y\) with \(y_j-y_i\ge k\) for
\(i<j\), where \(k\in(0,1]\),
\[|\mathsf F^\xi_n(a;\,t(x)+iy,\,\vec x)|\le C_0\,k^{-m}.\]

\emph{Sketch.}

\begin{itemize}
\tightlist
\item
  \(\mathsf F^\xi_n\) at such points is a matrix element, between
  \(\omega_\xi\) and \(\omega_\xi\), of a product of contractions
  (complex-time evolutions, spatial translations) and one-point slab
  operators of thickness of order \(k\).
\item
  Each one-point slab operator is bounded by the corresponding one for
  \(S^{\rm mix}\) (V.3.a), which is polynomial in \(1/k\) (input (ii) of
  §2).
\item
  Choose slab half-width \(p=k/8\). With the constants of §2(ii), take
  \(m=\sum_jN_{a_j}\) and \(C_0=\prod_j(Q_{a_j}8^{N_{a_j}})\).
\end{itemize}

\emph{Lean:}
\href{https://github.com/yang-mills-group/OpenAI-YM-20261028/blob/3456e5e75cf8d4cd558955a5c7638fba2d68f464/YM/EveryPhase/Pages.lean\#L19-L46}{\texttt{power}}.
\emph{(proof read; the matrix-element representation of the continuation
comes from V.2)}

\textbf{Lemma V.3.e (strict heights).} For all \(n\), \(a\), real \(x\)
and \(y_1<\dots<y_n\), the map
\(\xi\mapsto\mathsf F^\xi_n(a;t(x)+iy,\vec x)\) is continuous and its
\(\nu\)-integral is \(\mathsf F^{\rm mix}_n(a;t(x)+iy,\vec x)\).

\emph{Sketch.}

\begin{enumerate}
\def\labelenumi{\arabic{enumi}.}
\tightlist
\item
  At purely imaginary times \(z_j=iy_j\), \(\mathsf F^\xi_n\) is
  \(S^\xi\) of a word of distinct points, hence a character value. It is
  therefore continuous in \(\xi\), with barycenter \(S^{\rm mix}\).
\item
  Along the complex line \(c\mapsto iy+c\,t(x)\), for \(c\) in a strip
  \(|\mathrm{Im}c|<\delta\) (where the heights stay ordered with gaps
  \(\ge k/2\)), the functions are holomorphic and bounded uniformly in
  \(\xi\) by V.3.d.
\item
  By Vitali's theorem, convergence along nets \(\xi'\to\xi\) on the
  imaginary segment through \(0\) propagates to the whole strip.
\item
  The integral over \(X\) is holomorphic in \(c\). It agrees with the
  unselected function on that segment, hence on the strip by the
  identity theorem. Evaluate at \(c=1\).
\end{enumerate}

\emph{Lean:}
\href{https://github.com/yang-mills-group/OpenAI-YM-20261028/blob/3456e5e75cf8d4cd558955a5c7638fba2d68f464/YM/EveryPhase/Pages.lean\#L152-L170}{\texttt{imaginary\_evaluation}},
\href{https://github.com/yang-mills-group/OpenAI-YM-20261028/blob/3456e5e75cf8d4cd558955a5c7638fba2d68f464/YM/EveryPhase/Pages.lean\#L193-L228}{\texttt{strict\_evaluation}},
\href{https://github.com/yang-mills-group/OpenAI-YM-20261028/blob/3456e5e75cf8d4cd558955a5c7638fba2d68f464/YM/EveryPhase/Pages.lean\#L114-L146}{\texttt{integral\_holomorphic}},
\href{https://github.com/yang-mills-group/OpenAI-YM-20261028/blob/3456e5e75cf8d4cd558955a5c7638fba2d68f464/YM/EveryPhase/Vitali.lean\#L67-L89}{\texttt{propagate\_limit}}.

\textbf{Lemma V.3.f (tested strict heights).} For \(y_1<\dots<y_n\) and
\(g\in\mathcal S((\mathbb R^4)^n)\), define
\(\psi^\xi(a,y,g)=\int\mathsf F^\xi_n(a;t(x)+iy,\vec x)\,g(x)\,dx\).
Then \(\xi\mapsto\psi^\xi(a,y,g)\) is continuous, and
\(\int\psi^\xi\,d\nu=\psi^{\rm mix}\).

\emph{Sketch.} Dominated convergence with dominator \(C_0k^{-m}|g|\)
(V.3.d), and Fubini.

\emph{Lean:}
\href{https://github.com/yang-mills-group/OpenAI-YM-20261028/blob/3456e5e75cf8d4cd558955a5c7638fba2d68f464/YM/EveryPhase/Barycenter.lean\#L18-L42}{\texttt{eager}}.

\textbf{Lemma V.3.g (continuity and barycenter of the Wightman
distributions).} Item 9 of V.3.1, first part.

\emph{Sketch.} Let \(\rho=(0,1,\dots,n-1)\) and
\(D=-i\sum_{j=1}^n(j-1)\partial_{t_j}\) on joint test functions. The
Cauchy--Riemann equations and integration by parts give
\(\frac{d}{dt}\psi(a,y+t\rho,g)=\psi(a,y+t\rho,Dg)\).

\begin{enumerate}
\def\labelenumi{\arabic{enumi}.}
\item
  With \(m\) as in V.3.d, the boundary value has the Taylor
  representation
  \[W^\xi_n(a;g)=\sum_{k=0}^m\frac{(-1)^k}{k!}\,\psi^\xi(a,\rho,D^kg)-\int_0^1\frac{(-t)^m}{m!}\,\psi^\xi(a,t\rho,D^{m+1}g)\,dt .\]
\item
  By V.3.d the integrand is bounded by \(C_0\|D^{m+1}g\|_{L^1}\),
  uniformly in \(t\) and \(\xi\).
\item
  Continuity then follows from V.3.f and dominated convergence in \(t\).
\item
  For the barycenter: \(\psi^\xi(a,y,g)\to W^\xi_n(a;g)\) as \(y\to0\)
  among ordered heights, with a bound uniform in \(\xi\) and \(y\) given
  by the same representation. Integrate over \(X\) by dominated
  convergence, and use V.3.f at each \(y\).
\end{enumerate}

\emph{Lean:}
\href{https://github.com/yang-mills-group/OpenAI-YM-20261028/blob/3456e5e75cf8d4cd558955a5c7638fba2d68f464/YM/EveryPhase/Barycenter.lean\#L48-L73}{\texttt{word\_continuity}},
\href{https://github.com/yang-mills-group/OpenAI-YM-20261028/blob/3456e5e75cf8d4cd558955a5c7638fba2d68f464/YM/EveryPhase/Barycenter.lean\#L75-L90}{\texttt{barycenter}},
\href{https://github.com/yang-mills-group/OpenAI-YM-20261028/blob/3456e5e75cf8d4cd558955a5c7638fba2d68f464/YM/Shift/Breath.lean\#L27-L28}{\texttt{Waves.desired}},
\href{https://github.com/yang-mills-group/OpenAI-YM-20261028/blob/3456e5e75cf8d4cd558955a5c7638fba2d68f464/YM/Shift/Breath.lean\#L50-L66}{\texttt{Waves.meant}}.
The derivative convention is
\href{https://github.com/yang-mills-group/OpenAI-YM-20261028/blob/3456e5e75cf8d4cd558955a5c7638fba2d68f464/YM/Shift/Ridge.lean\#L78}{\texttt{Boundary.Waves.clock}}.

\textbf{Lemma V.3.h (inner products).} For finite sums of words
\(A,A'\):

\begin{itemize}
\tightlist
\item
  \(\xi\mapsto\langle A\Omega_\xi,A'\Omega_\xi\rangle\) is continuous on
  \(X\);
\item
  \(M_A:=\sup_\xi\|A\Omega_\xi\|^2=\sup_\xi\mu^\xi_A(\mathbb R^4)<\infty\).
\end{itemize}

\emph{Sketch.}
\(\langle\phi_\xi(w)\Omega_\xi,\phi_\xi(w')\Omega_\xi\rangle=W^\xi(w^\dagger w')\),
with the same star operation for all \(\xi\). Apply V.3.g, then use
compactness of \(X\).

\emph{Lean:}
\href{https://github.com/yang-mills-group/OpenAI-YM-20261028/blob/3456e5e75cf8d4cd558955a5c7638fba2d68f464/YM/EveryPhase/Light.lean\#L179-L183}{\texttt{word\_pair}},
\href{https://github.com/yang-mills-group/OpenAI-YM-20261028/blob/3456e5e75cf8d4cd558955a5c7638fba2d68f464/YM/EveryPhase/Light.lean\#L185-L195}{\texttt{sum\_pair}},
\href{https://github.com/yang-mills-group/OpenAI-YM-20261028/blob/3456e5e75cf8d4cd558955a5c7638fba2d68f464/YM/EveryPhase/Light.lean\#L197-L206}{\texttt{word\_bound}}.

\paragraph{4.3 Surjectivity of the boundary
isometry}\label{surjectivity-of-the-boundary-isometry}

\textbf{Lemma V.3.i (Euclidean vectors at positive heights are in the
range).} Fix \(n\ge1\), \(\xi\in X\), labels \(a\in\mathbb F^n\),
heights \(0<y_1<\dots<y_n\) and \(p\in(0,1)\) with \(4p<y_1\) and
\(4p<y_{j+1}-y_j\) for \(1\le j<n\).

\begin{itemize}
\tightlist
\item
  Let \(B_{\xi,p}(a)\) be the centered one-point slab operator of
  thickness \(2p\). For \(\mathrm{Im}z>\ell\) put
  \[R_{\xi,\ell}(z,\vec x)=e^{-(\mathrm{Im}z-\ell)H^{\rm OS}_\xi}e^{i\mathrm{Re}zH^{\rm OS}_\xi}U^{\rm OS}_\xi(\vec x).\]
\end{itemize}

For a real configuration \(q\), write \(z_j=q_j^4+iy_j\) and define
\[\Psi^\xi_p(a;q+iy)=R_{\xi,p}(z_1,\vec q_1)B_{\xi,p}(a_1)\left(\prod_{j=2}^n R_{\xi,2p}(z_j-z_{j-1},\vec q_j-\vec q_{j-1})B_{\xi,p}(a_j)\right)\omega_\xi,\]
where the factors are ordered by increasing \(j\) from left to right;
the product is empty when \(n=1\).

\begin{itemize}
\item
  Let \(h\in C^\infty_c((\mathbb R^4)^n)\) and use the convention
  \(g(q)=\widehat h(q)=\int e^{-i\sum_j p_j\cdot q_j}h(p)\,dp\). Set
  \(g_y=\widehat{e^{-\sum_jy_jp_j^4}h}\). Then
  \[\int g(q)\,\Psi^\xi_p(a;q+iy)\,dq=J_\xi\,\Phi_\xi(a;g_y).\] Here
  \(\Phi_\xi(a;\cdot)\) is the joint-test extension from V.3.b. Compact
  support of \(h\) ensures that the exponential multiplier preserves the
  Schwartz class.
\item
  Consequently every \(\Psi^\xi_p(a;q+iy)\) lies in
  \(\mathrm{ran}\,J_\xi\).
\end{itemize}

\emph{Sketch.}

\begin{enumerate}
\def\labelenumi{\arabic{enumi}.}
\tightlist
\item
  Call the two sides \(v\) and \(J_\xi\Phi\). Their squared norms and
  their inner product, \(\|v\|^2\), \(\langle v,J_\xi\Phi\rangle\) and
  \(\|J_\xi\Phi\|^2\), are each expressed through partial boundary
  values of \(W^\xi\) at heights \(y\).
\item
  On tests of compact frequency support these partial boundary values
  equal \(W^\xi\) tested against the damped test functions, so the three
  numbers coincide and \(\|v-J_\xi\Phi\|=0\).
\item
  Density of compact-frequency tests and continuity in \(g\) extend the
  membership in \(\mathrm{ran}\,J_\xi\) to all Schwartz \(g\).
\item
  Testing against all \(g\), and using continuity in \(q\), gives the
  pointwise statement. The range is closed since \(J_\xi\) is an
  isometry.
\end{enumerate}

\emph{Lean:}
\href{https://github.com/yang-mills-group/OpenAI-YM-20261028/blob/3456e5e75cf8d4cd558955a5c7638fba2d68f464/YM/EveryPhase/Hot.lean\#L29-L75}{\texttt{vector\_cooled}},
\href{https://github.com/yang-mills-group/OpenAI-YM-20261028/blob/3456e5e75cf8d4cd558955a5c7638fba2d68f464/YM/EveryPhase/Hot.lean\#L99-L109}{\texttt{vector\_in}},
\href{https://github.com/yang-mills-group/OpenAI-YM-20261028/blob/3456e5e75cf8d4cd558955a5c7638fba2d68f464/YM/EveryPhase/Hot.lean\#L111-L124}{\texttt{shot\_in}},
\href{https://github.com/yang-mills-group/OpenAI-YM-20261028/blob/3456e5e75cf8d4cd558955a5c7638fba2d68f464/YM/EveryPhase/Frequency.lean\#L103-L111}{\texttt{boundary\_tilt}},
\href{https://github.com/yang-mills-group/OpenAI-YM-20261028/blob/3456e5e75cf8d4cd558955a5c7638fba2d68f464/YM/EveryPhase/Frequency.lean\#L83}{\texttt{chilled}}.
\emph{(proof read; V.2's partial-boundary identities statement only)}

\textbf{Lemma V.3.j (surjectivity).} Item 2 of V.3.1.

\emph{Sketch.}

\begin{enumerate}
\def\labelenumi{\arabic{enumi}.}
\tightlist
\item
  The empty word gives \(\omega_\xi=J_\xi\Omega_\xi\). Take a nonempty
  positive-time Euclidean word with strictly increasing times. Its OS
  vector is a vector \(\Psi^\xi_p\) of V.3.i (real times \(0\), heights
  equal to the Euclidean times), so it lies in the range.
\item
  Take a general positive-time word. By symmetry its points may be
  reordered so that the times are nondecreasing, say \(r\). Then
  \(r+u\rho\) is strictly increasing for \(u>0\).
\item
  Its OS vector is the limit, as \(u\downarrow0\), of OS vectors in the
  range, by continuity of OS vectors in the points. It therefore lies in
  the range, because the range is closed.
\item
  These vectors span a dense subspace, so \(\mathrm{ran}\,J_\xi\) is
  dense and closed.
\end{enumerate}

Vacuum uniqueness is not used.

\emph{Lean:}
\href{https://github.com/yang-mills-group/OpenAI-YM-20261028/blob/3456e5e75cf8d4cd558955a5c7638fba2d68f464/YM/EveryPhase/Surjective.lean\#L178-L223}{\texttt{onto}},
\href{https://github.com/yang-mills-group/OpenAI-YM-20261028/blob/3456e5e75cf8d4cd558955a5c7638fba2d68f464/YM/EveryPhase/Surjective.lean\#L160-L176}{\texttt{ordered\_covered}},
\href{https://github.com/yang-mills-group/OpenAI-YM-20261028/blob/3456e5e75cf8d4cd558955a5c7638fba2d68f464/YM/EveryPhase/Surjective.lean\#L119-L128}{\texttt{bright\_shot}},
\href{https://github.com/yang-mills-group/OpenAI-YM-20261028/blob/3456e5e75cf8d4cd558955a5c7638fba2d68f464/YM/EveryPhase/Projection.lean\#L20-L25}{\texttt{reaches}}.

\paragraph{4.4 Low-energy bound and heat
decay}\label{low-energy-bound-and-heat-decay}

\textbf{Lemma V.3.k (time-averaged commutator).} Let \(\rho,\sigma\) be
finite positive measures on \(\mathbb R^4\) supported in
\(\overline V_+\), with \(\rho(\mathbb R^4)+\sigma(\mathbb R^4)\le M\).
Put
\[\hat\rho(\vec x,t)=\int e^{i(\vec p\cdot\vec x+p^4t)}\,d\rho(p),\qquad a=\hat\rho(\cdot)-\hat\sigma(-\,\cdot),\qquad D(\vec x)=\int_0^1\frac{a(\vec x,t)-a(\vec x,-t)}{2i}\,dt .\]
Then \(|D|\le M\). Write
\(\mathcal F f(x)=\int e^{-i\vec p\cdot x}f(\vec p)\,d\vec p\), so
\(c_3=(2\pi)^3\). If moreover \(D(\vec x)=0\) for \(|\vec x|>R\), then
for \(0<u\le1\)
\[\rho(\{u/2<p^4\le u\})\le c(R,M)\,u^2,\qquad c(R,M)=16\pi\,|B_1|\,|B_R|\,M/c_3 .\]

\emph{Sketch.}

\begin{enumerate}
\def\labelenumi{\arabic{enumi}.}
\tightlist
\item
  By Fubini, \(D\) is the spatial Fourier transform of
  \(\rho^\sharp+\sigma^\sharp\). Here \(\rho^\sharp\) is the image under
  \(p\mapsto\vec p\) of \(\pi(p^4)\,d\rho(p)\), with
  \(\pi(E)=\int_0^1\sin(Et)\,dt\), equal to \((1-\cos E)/E\) for \(E>0\)
  and to \(0\) at \(E=0\), and \(\sigma^\sharp\) is the analogous image
  of \(\sigma\) under \(p\mapsto-\vec p\).
\item
  On \(\overline V_+\) one has \(\pi(p^4)\ge0\) at all energies, and
  \(\pi(E)\ge E/\pi\) for \(0\le E\le1\).
\item
  Test \(D\) against \(\mathcal F f\), where \(f\) is a smooth
  spatial-momentum bump with \(0\le f\le1\), equal to \(1\) on the ball
  of radius \(u\) and supported in the ball of radius \(2u\). Then
  \(|\mathcal F f|\le|B_{2u}|=8|B_1|u^3\). Since \(D\) vanishes outside
  \(B_R\), the pairing is at most \(8|B_1||B_R|M\,u^3\).
\item
  The pairing equals
  \(c_3\big(\int f\,d\rho^\sharp+\int f\,d\sigma^\sharp\big)\), which is
  at least \(c_3\,\tfrac{u}{2\pi}\,\rho(\{u/2<p^4\le u\})\). On that
  shell \(|\vec p|\le p^4\le u\), so \(f=1\) there, and
  \(\pi(p^4)\ge u/(2\pi)\).
\item
  Thus the exponent \(2\) is \(3\) (spatial dimension, from the volume
  of the momentum ball) minus \(1\) (first-order vanishing of the
  weight).
\end{enumerate}

\emph{Lean:}
\href{https://github.com/yang-mills-group/OpenAI-YM-20261028/blob/3456e5e75cf8d4cd558955a5c7638fba2d68f464/YM/EveryPhase/Light.lean\#L43-L129}{\texttt{annulus}},
\href{https://github.com/yang-mills-group/OpenAI-YM-20261028/blob/3456e5e75cf8d4cd558955a5c7638fba2d68f464/YM/EveryPhase/Light.lean\#L16-L25}{\texttt{averaged\_bound}},
\href{https://github.com/yang-mills-group/OpenAI-YM-20261028/blob/3456e5e75cf8d4cd558955a5c7638fba2d68f464/YM/EveryPhase/Light.lean\#L35-L36}{\texttt{fee}},
\href{https://github.com/yang-mills-group/OpenAI-YM-20261028/blob/3456e5e75cf8d4cd558955a5c7638fba2d68f464/YM/EveryPhase/Average.lean\#L149-L160}{\texttt{two\_average}},
\href{https://github.com/yang-mills-group/OpenAI-YM-20261028/blob/3456e5e75cf8d4cd558955a5c7638fba2d68f464/YM/EveryPhase/Average.lean\#L113-L147}{\texttt{one\_average}},
\href{https://github.com/yang-mills-group/OpenAI-YM-20261028/blob/3456e5e75cf8d4cd558955a5c7638fba2d68f464/YM/EveryPhase/Average.lean\#L35-L44}{\texttt{pulse\_positive}},
\href{https://github.com/yang-mills-group/OpenAI-YM-20261028/blob/3456e5e75cf8d4cd558955a5c7638fba2d68f464/YM/EveryPhase/Average.lean\#L46-L59}{\texttt{pulse\_lower}}.

\textbf{Lemma V.3.l (dyadic summation).} Let \(\rho\) be a finite
measure and \(C\ge0\) with \(\rho(\{u/2<p^4\le u\})\le Cu^2\) for all
\(u\in(0,1]\). Then \(\rho(\{0<p^4\le u\})\le2Cu^2\) for \(u\in(0,1]\).

\emph{Sketch.} Cover \((0,u]\) by the shells with \(u_n=u2^{-n}\). Then
\(\sum_nCu^2\,4^{-n}\le\frac43Cu^2\); the bound \(2\) follows already
from the coarser majorant \(\sum_n2^{-n}=2\).

\emph{Lean:}
\href{https://github.com/yang-mills-group/OpenAI-YM-20261028/blob/3456e5e75cf8d4cd558955a5c7638fba2d68f464/YM/EveryPhase/Light.lean\#L135-L168}{\texttt{dyadic}}.

\textbf{Lemma V.3.m (locality at large distance, uniformly in the
phase).} Let \(A\) be a compactly supported finite sum of words, and
\(r\ge0\) a bound for \(|t|\) and \(|\vec x|\) on the supports of all
its test functions. Put \(R=4r+1\). Then for all \(\xi\), \(|t|\le1\)
and \(|\vec x|>R\):
\[\hat\mu^\xi_A(\vec x,t)-\hat\mu^\xi_{A^\dagger}(-\vec x,-t)=\langle\Omega_\xi,[A^\dagger,\alpha_{(\vec x,t)}A]\,\Omega_\xi\rangle=0,\]
where \(\alpha\) denotes translation of the test functions. Hence the
function \(D\) of V.3.k, built from \(\rho=\mu^\xi_A\) and
\(\sigma=\mu^\xi_{A^\dagger}\), vanishes for \(|\vec x|>R\).

\emph{Sketch.}

\begin{itemize}
\tightlist
\item
  The first identity follows from covariance and the adjoint relation.
\item
  The supports of \(A^\dagger\) and \(\alpha_{(\vec x,t)}A\) are
  spacelike separated, so the commutator vanishes by locality.
\item
  \(R\) depends only on the supports, and translating test functions is
  the same operation for every \(\xi\).
\end{itemize}

\emph{Lean:}
\href{https://github.com/yang-mills-group/OpenAI-YM-20261028/blob/3456e5e75cf8d4cd558955a5c7638fba2d68f464/YM/EveryPhase/Light.lean\#L208-L219}{\texttt{late}},
\href{https://github.com/yang-mills-group/OpenAI-YM-20261028/blob/3456e5e75cf8d4cd558955a5c7638fba2d68f464/YM/EveryPhase/Polynomials.lean\#L146-L160}{\texttt{distant}},
\href{https://github.com/yang-mills-group/OpenAI-YM-20261028/blob/3456e5e75cf8d4cd558955a5c7638fba2d68f464/YM/EveryPhase/Polynomials.lean\#L175-L207}{\texttt{comm\_expect}},
\href{https://github.com/yang-mills-group/OpenAI-YM-20261028/blob/3456e5e75cf8d4cd558955a5c7638fba2d68f464/YM/EveryPhase/Polynomials.lean\#L106-L119}{\texttt{op\_comm}}.

\textbf{Proposition V.3.n (low-energy bound).} Item 5 of V.3.1, with
\(M=2c(R,M_A+M_{A^\dagger})\).

\emph{Sketch.} Apply V.3.k to \(\rho=\mu^\xi_A\) and
\(\sigma=\mu^\xi_{A^\dagger}\). These are supported in \(\overline V_+\)
by the spectrum condition, with mass bounds from V.3.h and \(R\) from
V.3.m. Then apply V.3.l.

\emph{Lean:}
\href{https://github.com/yang-mills-group/OpenAI-YM-20261028/blob/3456e5e75cf8d4cd558955a5c7638fba2d68f464/YM/EveryPhase/Light.lean\#L221-L232}{\texttt{light\_bound}}.

\textbf{Proposition V.3.o (heat decay with the full zero-energy
projection).} Item 6 of V.3.1, with
\(M=2\max(M_A,2c(R,M_A+M_{A^\dagger}))\).

\emph{Sketch.}

\begin{enumerate}
\def\labelenumi{\arabic{enumi}.}
\item
  By the spectral theorem and
  \(\mathrm{supp}\,\mu^\xi_A\subset\{p^4\ge0\}\):
  \[\langle A\Omega_\xi,(e^{-sH_\xi}-P^0_\xi)A\Omega_\xi\rangle=\int_{\{p^4>0\}}e^{-sp^4}\,d\mu^\xi_A\ \ge0 .\]
\item
  Let \(\rho\) be a finite measure with \(\rho(\{0<p^4\le t\})\le Ct^2\)
  for all \(t>0\); for \(t>1\) this follows from the total-mass bound.
  Write \(e^{-sE}=\int_E^\infty se^{-st}\,dt\). Fubini gives
  \[\int_{\{p^4>0\}}e^{-sp^4}d\rho=\int_0^\infty se^{-st}\rho(\{0<p^4\le t\})\,dt\le Cs\int_0^\infty t^2e^{-st}dt=\Gamma(3)\,C\,s^{-2}.\]
\end{enumerate}

\emph{Lean:}
\href{https://github.com/yang-mills-group/OpenAI-YM-20261028/blob/3456e5e75cf8d4cd558955a5c7638fba2d68f464/YM/EveryPhase/Projection.lean\#L145-L150}{\texttt{residual\_pair}},
\href{https://github.com/yang-mills-group/OpenAI-YM-20261028/blob/3456e5e75cf8d4cd558955a5c7638fba2d68f464/YM/EveryPhase/Decay.lean\#L22-L75}{\texttt{decay}},
\href{https://github.com/yang-mills-group/OpenAI-YM-20261028/blob/3456e5e75cf8d4cd558955a5c7638fba2d68f464/YM/EveryPhase/Decay.lean\#L173-L192}{\texttt{inverse\_square}},
\href{https://github.com/yang-mills-group/OpenAI-YM-20261028/blob/3456e5e75cf8d4cd558955a5c7638fba2d68f464/YM/EveryPhase/Projection.lean\#L97-L100}{\texttt{bHeat\_pair}}.

\textbf{Lemma V.3.p (continuity of the heat form).} For every finite sum
of words \(A\) and \(s>0\),
\(\xi\mapsto\langle A\Omega_\xi,e^{-sH_\xi}A\Omega_\xi\rangle=\int e^{-sp^4}\,d\mu^\xi_A\)
is continuous.

\emph{Sketch.}

\begin{enumerate}
\def\labelenumi{\arabic{enumi}.}
\tightlist
\item
  The Fourier transform
  \(\hat\mu^\xi_A(p)=\langle A\Omega_\xi,(\alpha_pA)\Omega_\xi\rangle\)
  is continuous in \(\xi\) by V.3.h, and bounded by \(M_A\).
\item
  Hence \(\int h\,d\mu^\xi_A\) is continuous in \(\xi\) for every
  Schwartz \(h\), by Parseval and dominated convergence.
\item
  On \(\overline V_+\) outside the ball of radius \(2N\) one has
  \(p^4\ge N\). Multiplying \(e^{-sp^4}\) by a bump equal to \(1\) on
  that ball changes the integral by at most \(M_Ae^{-sN}\), uniformly in
  \(\xi\).
\item
  A uniform limit of continuous functions is continuous.
\end{enumerate}

\emph{Lean:}
\href{https://github.com/yang-mills-group/OpenAI-YM-20261028/blob/3456e5e75cf8d4cd558955a5c7638fba2d68f464/YM/EveryPhase/Decay.lean\#L113-L170}{\texttt{uniform\_heat}},
\href{https://github.com/yang-mills-group/OpenAI-YM-20261028/blob/3456e5e75cf8d4cd558955a5c7638fba2d68f464/YM/EveryPhase/Decay.lean\#L82-L88}{\texttt{char\_cont}},
\href{https://github.com/yang-mills-group/OpenAI-YM-20261028/blob/3456e5e75cf8d4cd558955a5c7638fba2d68f464/YM/EveryPhase/Decay.lean\#L93-L111}{\texttt{test\_cont}}.

\paragraph{4.5 Every phase}\label{every-phase}

\textbf{Lemma V.3.q (vacuum-subtracted decay at every phase).} For every
compactly supported finite sum of words \(A\) there is \(M\ge0\) such
that \(0\le r_\xi(s)\le Ms^{-2}\) for all \(\xi\in X\) and \(s\ge1\),
with \(r_\xi(s)\) as in §3.

\emph{Sketch.}

\begin{enumerate}
\def\labelenumi{\arabic{enumi}.}
\tightlist
\item
  \(r_\xi(s)\) is continuous in \(\xi\) by V.3.p and V.3.h;
  \(\Omega_\xi\) is the vector of the empty word.
\item
  Where \(\mathcal D_\xi=\mathbb C\omega_\xi\), which is \(\nu\)-almost
  everywhere by V.3.c, \(J_\xi\) intertwines the zero-energy
  projections. Then \(P^0_\xi x=\langle\Omega_\xi,x\rangle\Omega_\xi\),
  so
  \(r_\xi(s)=\langle A\Omega_\xi,(e^{-sH_\xi}-P^0_\xi)A\Omega_\xi\rangle\in[0,Ms^{-2}]\)
  by V.3.o.
\item
  The set \(\{\xi:\ r_\xi(s)\in[0,Ms^{-2}]\}\) is closed and of full
  measure. Since \(\nu\) charges every nonempty open set, it is \(X\).
\end{enumerate}

\emph{Lean:}
\href{https://github.com/yang-mills-group/OpenAI-YM-20261028/blob/3456e5e75cf8d4cd558955a5c7638fba2d68f464/YM/EveryPhase/Vacuum.lean\#L22-L39}{\texttt{proj\_eq}},
\href{https://github.com/yang-mills-group/OpenAI-YM-20261028/blob/3456e5e75cf8d4cd558955a5c7638fba2d68f464/YM/EveryPhase/Vacuum.lean\#L42-L46}{\texttt{univ\_of\_closed}},
\href{https://github.com/yang-mills-group/OpenAI-YM-20261028/blob/3456e5e75cf8d4cd558955a5c7638fba2d68f464/YM/EveryPhase/Vacuum.lean\#L51-L78}{\texttt{extends\_decay}}.

\textbf{Lemma V.3.r (zero-energy projection = vacuum projection).}
Proposition V.3.2.

\emph{Sketch.}

\begin{enumerate}
\def\labelenumi{\arabic{enumi}.}
\tightlist
\item
  \(e^{-sH_\xi}\to P^0_\xi\) strongly as \(s\to\infty\), so V.3.q gives
  \(\langle x,P^0_\xi x\rangle=|\langle\Omega_\xi,x\rangle|^2\) for
  \(x=A\Omega_\xi\).
\item
  Since \(P^0_\xi\) is an orthogonal projection with
  \(P^0_\xi\Omega_\xi=\Omega_\xi\), this gives
  \(\|P^0_\xi x-\langle\Omega_\xi,x\rangle\Omega_\xi\|^2=\|P^0_\xi x\|^2-|\langle\Omega_\xi,x\rangle|^2=0\).
\item
  The identity extends from compactly supported to Schwartz test
  functions by cutoffs converging in \(\mathcal S\) and continuity of
  the word vectors in the joint test function.
\item
  Bounded operators that agree on all word vectors are equal.
\end{enumerate}

\emph{Lean:}
\href{https://github.com/yang-mills-group/OpenAI-YM-20261028/blob/3456e5e75cf8d4cd558955a5c7638fba2d68f464/YM/EveryPhase/Vacuum.lean\#L16-L20}{\texttt{groundvac}},
\href{https://github.com/yang-mills-group/OpenAI-YM-20261028/blob/3456e5e75cf8d4cd558955a5c7638fba2d68f464/YM/EveryPhase/Projection.lean\#L102-L109}{\texttt{bHeat\_limit}},
\href{https://github.com/yang-mills-group/OpenAI-YM-20261028/blob/3456e5e75cf8d4cd558955a5c7638fba2d68f464/YM/EveryPhase/Vacuum.lean\#L80-L112}{\texttt{cluster\_zero}},
\href{https://github.com/yang-mills-group/OpenAI-YM-20261028/blob/3456e5e75cf8d4cd558955a5c7638fba2d68f464/YM/EveryPhase/Vacuum.lean\#L115-L144}{\texttt{all\_plant}}.

\textbf{Lemma V.3.s (Euclidean vacuum and Wightman axioms at every
phase).} Items 1, 3 and 4 of V.3.1.

\emph{Sketch.}

\begin{enumerate}
\def\labelenumi{\arabic{enumi}.}
\tightlist
\item
  Let \(v\in\mathcal D_\xi\). By V.3.j, \(v=J_\xi x\) for some \(x\).
\item
  The projection onto \(\mathcal D_\xi\) equals the zero-energy spectral
  projection of the OS translation group: both are the strong limit of
  \(e^{-sH^{\rm OS}_\xi}\).
\item
  Intertwining with V.3.2 gives
  \(v=J_\xi P^0_\xi x=\langle\Omega_\xi,x\rangle\,\omega_\xi\). This is
  item 3.
\item
  Item 1 then follows from V.3.b, and item 4 from V.3.2.
\end{enumerate}

\emph{Lean:}
\href{https://github.com/yang-mills-group/OpenAI-YM-20261028/blob/3456e5e75cf8d4cd558955a5c7638fba2d68f464/YM/EveryPhase/Vacuum.lean\#L146-L153}{\texttt{every}},
\href{https://github.com/yang-mills-group/OpenAI-YM-20261028/blob/3456e5e75cf8d4cd558955a5c7638fba2d68f464/YM/EveryPhase/Projection.lean\#L77-L86}{\texttt{reflection\_projection}},
\href{https://github.com/yang-mills-group/OpenAI-YM-20261028/blob/3456e5e75cf8d4cd558955a5c7638fba2d68f464/YM/EveryPhase/Vacuum.lean\#L155-L156}{\texttt{boundary\_realizes}},
\href{https://github.com/yang-mills-group/OpenAI-YM-20261028/blob/3456e5e75cf8d4cd558955a5c7638fba2d68f464/YM/EveryPhase/Vacuum.lean\#L158-L162}{\texttt{boundary\_zero\_energy}}.

\paragraph{4.6 Irreducibility}\label{irreducibility}

\textbf{Proposition V.3.t (trivial weak commutant).} Item 7 of V.3.1.

\emph{Sketch.} The weak commutant is first rewritten with
\(\phi_\xi(a^*,\bar f)\) in place of the Hilbert adjoint, since the
former is the restriction of the latter to \(\mathfrak D_\xi\). The
converse direction is the adjoint relation. Now let \(B\) be weakly
central and write \(\Omega=\Omega_\xi\).

\begin{enumerate}
\def\labelenumi{\arabic{enumi}.}
\item
  For word vectors \(u=\phi(w)\Omega\) and \(v=\phi(w^\dagger)\Omega\),
  and \(\mathrm{Im}z\ge0\), set
  \[p(z)=\langle J_\xi u,e^{izH^{\rm OS}_\xi}J_\xi B\Omega\rangle,\qquad q(z)=\langle J_\xi B^*\Omega,e^{izH^{\rm OS}_\xi}J_\xi v\rangle .\]
  Both are bounded, continuous, and holomorphic for \(\mathrm{Im}z>0\).
  On the real axis \(p(t)=\langle u,U_\xi(t)B\Omega\rangle\), and at
  \(z=is\) one has \(p(is)=\langle u,e^{-sH_\xi}B\Omega\rangle\).
\item
  Covariance and weak centrality give \(p(t)=q(-t)\) for real \(t\).
\item
  Gluing \(p\) on the upper half-plane to \(q(-\,\cdot)\) on the lower
  one gives a bounded continuous function, holomorphic off
  \(\mathbb R\). It is entire by Morera's theorem and constant by
  Liouville's theorem.
\item
  Hence \(\langle u,e^{-sH_\xi}B\Omega\rangle=\langle u,B\Omega\rangle\)
  for all \(s\ge0\). Letting \(s\to\infty\) and using V.3.2 gives
  \(B\Omega=a\Omega\) with \(a=\langle\Omega,B\Omega\rangle\).
\item
  Weak centrality then gives
  \(\langle\phi(c)\Omega,B\phi(b)\Omega\rangle=\langle\phi(b^\dagger)\phi(c)\Omega,B\Omega\rangle=a\langle\phi(c)\Omega,\phi(b)\Omega\rangle\)
  for all words, so \(B=a\,1\).
\end{enumerate}

The proof uses V.3.2 and the reconstruction data of V.3.b. It requires
neither surjectivity of \(J_\xi\) nor invariance of \(\mathfrak D_\xi\)
under \(B\).

\emph{Lean:}
\href{https://github.com/yang-mills-group/OpenAI-YM-20261028/blob/3456e5e75cf8d4cd558955a5c7638fba2d68f464/YM/EveryPhase/Irreducible.lean\#L118-L193}{\texttt{commutant}},
\href{https://github.com/yang-mills-group/OpenAI-YM-20261028/blob/3456e5e75cf8d4cd558955a5c7638fba2d68f464/YM/EveryPhase/Irreducible.lean\#L56-L84}{\texttt{saturated}},
\href{https://github.com/yang-mills-group/OpenAI-YM-20261028/blob/3456e5e75cf8d4cd558955a5c7638fba2d68f464/YM/EveryPhase/Irreducible.lean\#L22-L54}{\texttt{paste}},
\href{https://github.com/yang-mills-group/OpenAI-YM-20261028/blob/3456e5e75cf8d4cd558955a5c7638fba2d68f464/YM/EveryPhase/Irreducible.lean\#L195-L205}{\texttt{precise}},
\href{https://github.com/yang-mills-group/OpenAI-YM-20261028/blob/3456e5e75cf8d4cd558955a5c7638fba2d68f464/YM/EveryPhase/Irreducible.lean\#L207-L216}{\texttt{irreducible}}.

\paragraph{4.7 The continued stress
tensor}\label{the-continued-stress-tensor}

\textbf{Conventions.}

\begin{itemize}
\tightlist
\item
  Indices \(a,b\in\{1,2,3,4\}\), with \(4\) the time index.
\item
  \(\eta=\mathrm{diag}(-1,-1,-1,1)\), \(\zeta_4=-i\), \(\zeta_j=1\).
\item
  The Minkowski stress labels are
  \(T^{\rm M}_{ab}=-\zeta_a\zeta_bT^{\rm E}_{ab}\), that is
  \(T^{\rm M}_{44}=T^{\rm E}_{44}\), \(T^{\rm M}_{4j}=iT^{\rm E}_{4j}\),
  \(T^{\rm M}_{jk}=-T^{\rm E}_{jk}\). Write
  \(T_{ab}(f)=\phi_\xi(T^{\rm M}_{ab},f)\).
\item
  A Killing field is \(K(x)=k+Vx\), where \(k\in\mathbb R^4\) and
  \(V^{\mathsf T}\eta+\eta V=0\). Write \(\Sigma_K\) for the
  infinitesimal label action of \(V\) (VI.1). Its Euclidean continuation
  is \(K_{\rm E}^b(y)=\zeta_b(k^b+\sum_c V_{bc}\kappa_c y^c)\), where
  \(\kappa_4=i\) and \(\kappa_j=1\) for \(j\le3\). The Euclidean
  derivative label is denoted by \(\partial_b^{\rm E}B\)
  (\href{https://github.com/yang-mills-group/OpenAI-YM-20261028/blob/3456e5e75cf8d4cd558955a5c7638fba2d68f464/YM/EveryPhase/Killing.lean\#L20-L23}{\texttt{EveryPhase.Stress.Killing}}).
\item
  The charge is \[Q_K(f,c)=\sum_bT_{4b}\big(K^b\,(c\otimes f)\big),\]
  for a time profile \(f\) and a spatial cutoff \(c\).
\end{itemize}

\textbf{Lemma V.3.u (Euclidean contact identity and current jump).} Fix
\(\xi\in X\), distinct Euclidean points \(x_1,\ldots,x_{n+1}\), a
distinguished label \(B\in F_d\) at \(x_k\), arbitrary spectator labels,
and a Killing field \(K\). Let
\(\chi\in C^\infty_c(\mathbb R^4;\mathbb R)\) equal \(1\) near \(x_k\),
with every other \(x_j\) outside \(\mathrm{supp}\,\chi\). Then the
following integral is absolutely convergent and
\[-\int\sum_{a,b}\partial_a\chi(y)\,K^b_{\rm E}(y)\,S^\xi(\ldots,T^{\rm E}_{ab}(y),\ldots)\,dy=S^\xi(\ldots,\Sigma_KB(x_k),\ldots)+\sum_bK^b_{\rm E}(x_k)\,S^\xi(\ldots,\partial_b^{\rm E}B(x_k),\ldots).\]
If the exterior Euclidean times are strictly ordered and
\(c\in C^\infty_c(\mathbb R^3;\mathbb R)\) equals \(1\) near
\(\vec x_k\), the difference of the currents
\(\sum_b\int c(\vec y)K_{\rm E}^b(\vec y,t)S^\xi(\ldots,T^{\rm E}_{4b}(\vec y,t),\ldots)\,d\vec y\),
with the stress insertion immediately after \(B\) at
\(t=x_k^4+\varepsilon\) and immediately before \(B\) at
\(t=x_k^4-\varepsilon\), tends to the right side as
\(\varepsilon\downarrow0\).

\emph{Sketch.} Decompose \(V\) into a spatial rotation and a boost. The
translation and rotation Ward identities give the displayed formula,
including the Wick factors. Choose \(\chi(y)=c(\vec y)h(y^4-x_k^4)\)
with a fixed smooth time cutoff supported between the neighboring
exterior times. Conservation identifies the time derivative of the
integrated current with the flux through \(\nabla c\). The fundamental
theorem of calculus, applied on both sides of \(x_k^4\), gives the jump.
For complex exterior times, the spacelike continuation germs in (T0)
give equal limits for the two spatial fluxes; bounded derivatives give
joint current limits. Vitali's theorem then continues the jump identity
from imaginary exterior times. Only spatially integrated currents have a
jump limit here.

\emph{Lean:}
\href{https://github.com/yang-mills-group/OpenAI-YM-20261028/blob/3456e5e75cf8d4cd558955a5c7638fba2d68f464/YM/EveryPhase/Shell.lean\#L128-L176}{\texttt{euclidean\_shell}},
\href{https://github.com/yang-mills-group/OpenAI-YM-20261028/blob/3456e5e75cf8d4cd558955a5c7638fba2d68f464/YM/EveryPhase/Pure.lean\#L43-L115}{\texttt{Stress.pure\_clock}},
\href{https://github.com/yang-mills-group/OpenAI-YM-20261028/blob/3456e5e75cf8d4cd558955a5c7638fba2d68f464/YM/EveryPhase/Through.lean\#L129-L142}{\texttt{Stress.through\_jump}}.
The cutoff condition is
\href{https://github.com/yang-mills-group/OpenAI-YM-20261028/blob/3456e5e75cf8d4cd558955a5c7638fba2d68f464/YM/Euclid/Stress.lean\#L72-L73}{\texttt{Table.Window}}.

\textbf{Lemma V.3.v (Minkowski jump identity).} Fix \(\xi\), a Killing
field \(K\), labels at \(n+1\) points and a distinguished index \(k\).
Let \(M\in C^\infty_c((\mathbb R^4)^{n+1})\),
\(f\in C^\infty_c(\mathbb R)\) and
\(c\in C^\infty_c(\mathbb R^3;\mathbb R)\). For every
\(x\in\mathrm{supp}\,M\), assume that \(c=1\) near \(\vec x_k\) and
\[|s-x_k^4|<|\vec y-\vec x_k|\qquad(s\in\mathrm{supp}\,f,\ \vec y\in\mathrm{supp}\,\partial_jc,\ j=1,2,3).\]
Then
\[\sum_b\zeta_b\Big[W^\xi(\ldots,a_k,\mathcal T_b,\ldots)-W^\xi(\ldots,\mathcal T_b,a_k,\ldots)\Big]=\Big(\int f\Big)\Big[W^\xi(\ldots,\Sigma_Ka_k,\ldots;M)+\sum_b\zeta_bW^\xi(\ldots,\partial_b^{\rm E}a_k,\ldots;K^b_kM)\Big].\]
The left side is tested jointly against \(M\) in the original variables
and \(K^b(y)c(\vec y)f(y^4)\) in the inserted stress variable \(y\).
Thus \(\mathcal T_b\) denotes the Euclidean label \(T^{\rm E}_{4b}\)
inserted immediately after, respectively before, the \(k\)-th field. On
the right, \(K^b_kM(x)=K^b(x_k)M(x)\).

\emph{Sketch.}

\begin{itemize}
\tightlist
\item
  Use the common sequence of tied heights from (T0). At each height,
  integrate the continued current jump of V.3.u against \(M(x)f(t)\).
  Joint convergence and compactness of the test supports give a common
  integrable bound; the right side is independent of \(t\), producing
  \(\int f\).
\item
  Pass to the limit with the continuity of the partial boundary values.
\end{itemize}

\emph{Lean:}
\href{https://github.com/yang-mills-group/OpenAI-YM-20261028/blob/3456e5e75cf8d4cd558955a5c7638fba2d68f464/YM/EveryPhase/Decree.lean\#L208-L242}{\texttt{Stress.verdict}},
\href{https://github.com/yang-mills-group/OpenAI-YM-20261028/blob/3456e5e75cf8d4cd558955a5c7638fba2d68f464/YM/EveryPhase/Decree.lean\#L77-L164}{\texttt{Stress.partial\_jump}},
\href{https://github.com/yang-mills-group/OpenAI-YM-20261028/blob/3456e5e75cf8d4cd558955a5c7638fba2d68f464/YM/EveryPhase/Decree.lean\#L73-L75}{\texttt{Stress.Jury}}.

\textbf{Lemma V.3.w (Ward identity with arbitrary spectators).} In V.3.v
take \(\int f=1\) and a product test with a compactly supported factor
\(m\) at the distinguished field. The resulting identity holds with
arbitrary Schwartz spectator words before and after the inserted pair.
The cutoff conditions need hold only for points of \(\mathrm{supp}\,m\).

\emph{Sketch.} Cut off the spectator test functions. Apply V.3.v. Pass
to the limit by continuity of \(W^\xi\) on joint test functions.

\emph{Lean:}
\href{https://github.com/yang-mills-group/OpenAI-YM-20261028/blob/3456e5e75cf8d4cd558955a5c7638fba2d68f464/YM/EveryPhase/Court.lean\#L147-L163}{\texttt{Stress.Schwartz\_shells}},
\href{https://github.com/yang-mills-group/OpenAI-YM-20261028/blob/3456e5e75cf8d4cd558955a5c7638fba2d68f464/YM/EveryPhase/Court.lean\#L77-L117}{\texttt{Stress.compact\_shells}}.

\textbf{Lemma V.3.x (the charge generates the transformation).} Fix
\(\xi\in X\), \(B\in F_d\) and a Killing field \(K\). Let
\(\mathcal O\subset\mathbb R^4\) be compact,
\(m\in C^\infty_c(\mathbb R^4)\) with
\(\mathrm{supp}\,m\subset\mathcal O\), and
\(f\in C^\infty_c(\mathbb R;\mathbb R)\) with \(\int f=1\). Let
\(c\in C^\infty_c(\mathbb R^3;\mathbb C)\) equal \(1\) on an open ball
such that, for every \(s\in\mathrm{supp}\,f\), every point
\((\vec y,s)\) over its complement is strictly spacelike to every point
of \(\mathcal O\). Then, on \(\mathfrak D_\xi\),
\[i[Q_K(f,c),\phi_\xi(B,m)]=\phi_\xi(\Sigma_KB,m)-\sum_{a=1}^4\phi_\xi\bigl(B,\partial_a(K^am)\bigr).\]
The derivative includes the affine coefficient \(K^a\). Equivalently,
the last term is \(\sum_a\zeta_a\phi_\xi(\partial_a^{\rm E}B,K^am)\).

\emph{Sketch.}

\begin{itemize}
\tightlist
\item
  Replace \(c\) by its real part. The difference of the two charges
  commutes with \(\phi_\xi(B,m)\) by locality, since it is supported off
  the ball.
\item
  Then take matrix elements between word vectors and apply V.3.w.
\end{itemize}

\emph{Lean:}
\href{https://github.com/yang-mills-group/OpenAI-YM-20261028/blob/3456e5e75cf8d4cd558955a5c7638fba2d68f464/YM/EveryPhase/Charge.lean\#L143-L177}{\texttt{Stress.local\_action}},
\href{https://github.com/yang-mills-group/OpenAI-YM-20261028/blob/3456e5e75cf8d4cd558955a5c7638fba2d68f464/YM/EveryPhase/Charge.lean\#L104-L137}{\texttt{Stress.measure\_identity}},
\href{https://github.com/yang-mills-group/OpenAI-YM-20261028/blob/3456e5e75cf8d4cd558955a5c7638fba2d68f464/YM/EveryPhase/Charge.lean\#L80-L81}{\texttt{Stress.Charge}},
\href{https://github.com/yang-mills-group/OpenAI-YM-20261028/blob/3456e5e75cf8d4cd558955a5c7638fba2d68f464/YM/EveryPhase/Charge.lean\#L14-L17}{\texttt{Stress.Season}},
\href{https://github.com/yang-mills-group/OpenAI-YM-20261028/blob/3456e5e75cf8d4cd558955a5c7638fba2d68f464/YM/EveryPhase/Charge.lean\#L37-L40}{\texttt{Stress.Reservation}}.

\textbf{Proposition V.3.y (the stress tensor of \(\mathfrak W_\xi\)).}
For every \(\xi\in X\) the following hold. Statements (T1)--(T4) hold
for all \(f\in\mathcal S(\mathbb R^4)\), \(a,b\in\{1,2,3,4\}\) and
\(u,v,x\in\mathfrak D_\xi\).

\begin{itemize}
\tightlist
\item
  \textbf{(T0)} There are holomorphic extensions of \(\mathsf F^\xi\) to
  the full forward tube and continuous tested partial boundary values,
  equal to \(W^\xi\) at zero height. For each adjacent pair there is a
  sequence of tied imaginary heights tending to zero, independent of the
  labels, at which the two field orders have equal continuation germs at
  spacelike separation. The tube consists of complex configurations with
  consecutive imaginary differences in the open forward cone; partial
  boundary values are taken on the corresponding closed cone. These data
  are used in V.3.u--V.3.v.
\item
  \textbf{(T1)} \(U_\xi(g)\mathfrak D_\xi\subset\mathfrak D_\xi\) for
  every proper orthochronous Poincaré transformation \(g\), and
  \(U_\xi(g)T_{ab}(f)x=\sum_{i,j}\Lambda(g)_{ia}\Lambda(g)_{jb}\,T_{ij}(f\circ g^{-1})\,U_\xi(g)x\).
\item
  \textbf{(T2)} \(T_{ab}=T_{ba}\).
\item
  \textbf{(T3)}
  \(\langle T_{ab}(f)u,v\rangle=\langle u,T_{ab}(\bar f)v\rangle\).
\item
  \textbf{(T4)} \(\sum_a\eta_{aa}\,T_{ab}(\partial_af)=0\).
\item
  \textbf{(T5)} The commutator identity of V.3.x.
\item
  \textbf{(T6)} Under the cutoff and profile conditions of V.3.x, for
  every finite sum of words \(A\) whose test functions are supported in
  \(\mathcal O\), \(i[Q_K(f,c),A]\) is the sum obtained by replacing one
  factor at a time by the right side of V.3.x and retaining the original
  order of the other factors. The formula holds for any choice of
  representatives of the labels in the spaces \(F_d\).
\end{itemize}

\emph{Sketch.}

\begin{itemize}
\tightlist
\item
  (T0) is supplied by the full-tube and partial-boundary reconstruction
  of V.2, with one common sequence of exchange heights for all labels
  (\href{https://github.com/yang-mills-group/OpenAI-YM-20261028/blob/3456e5e75cf8d4cd558955a5c7638fba2d68f464/YM/EveryPhase/Cross.lean\#L76-L83}{\texttt{EveryPhase.Stress.Prospect}},
  \href{https://github.com/yang-mills-group/OpenAI-YM-20261028/blob/3456e5e75cf8d4cd558955a5c7638fba2d68f464/YM/EveryPhase/Decree.lean\#L167-L171}{\texttt{EveryPhase.Stress.Gallery}}).
\item
  (T1)--(T4) are the tensor covariance, symmetry, reality and
  conservation of \(T^{\rm E}\) (VI.1), read through the Wick factors
  and the derivative rule for fields.
\item
  (T5) is V.3.x; (T6) follows from (T5) by the derivation property of
  commutators.
\end{itemize}

\emph{Lean:}
\href{https://github.com/yang-mills-group/OpenAI-YM-20261028/blob/3456e5e75cf8d4cd558955a5c7638fba2d68f464/YM/EveryPhase/Implement.lean\#L81-L92}{\texttt{Stress.Force}},
\href{https://github.com/yang-mills-group/OpenAI-YM-20261028/blob/3456e5e75cf8d4cd558955a5c7638fba2d68f464/YM/EveryPhase/Implement.lean\#L94-L116}{\texttt{Stress.force}},
\href{https://github.com/yang-mills-group/OpenAI-YM-20261028/blob/3456e5e75cf8d4cd558955a5c7638fba2d68f464/YM/EveryPhase/Implement.lean\#L64-L78}{\texttt{Stress.Generated}},
\href{https://github.com/yang-mills-group/OpenAI-YM-20261028/blob/3456e5e75cf8d4cd558955a5c7638fba2d68f464/YM/EveryPhase/Tensor.lean\#L210-L215}{\texttt{Stress.tensor\_lorentz}},
\href{https://github.com/yang-mills-group/OpenAI-YM-20261028/blob/3456e5e75cf8d4cd558955a5c7638fba2d68f464/YM/EveryPhase/Tensor.lean\#L217-L220}{\texttt{Stress.tensor\_herm}},
\href{https://github.com/yang-mills-group/OpenAI-YM-20261028/blob/3456e5e75cf8d4cd558955a5c7638fba2d68f464/YM/EveryPhase/Tensor.lean\#L227-L237}{\texttt{Stress.tensor\_conserved}}.

\paragraph{4.8 Assembly}\label{assembly}

\textbf{Proof of V.3.1.} Take the family of V.3.b. Then:

\begin{itemize}
\tightlist
\item
  item 1 is V.3.s; item 2 is V.3.j; item 3 is V.3.s; item 4 is V.3.2;
\item
  item 5 is V.3.n; item 6 is V.3.o; item 7 is V.3.t; item 8 is V.3.y;
\item
  item 9: the Wightman part is V.3.g; the Euclidean part is V.1's
  barycenter formula, which translates a word of distinct points into a
  slab using the homogeneity of \(S^{\rm mix}\).
\end{itemize}

\emph{Lean:}
\href{https://github.com/yang-mills-group/OpenAI-YM-20261028/blob/3456e5e75cf8d4cd558955a5c7638fba2d68f464/YM/EveryPhase/Extension.lean\#L75-L80}{\texttt{all\_characters}},
\href{https://github.com/yang-mills-group/OpenAI-YM-20261028/blob/3456e5e75cf8d4cd558955a5c7638fba2d68f464/YM/Phase/Mixture.lean\#L113-L117}{\texttt{Resolution.barycenter}}.

\subsubsection{5. Where the difficulty
lies}\label{where-the-difficulty-lies-22}

The delicate points are the common bounds near the analytic boundary,
the recovery of positive-height Euclidean vectors from joint frequency
tests, and the joint current limits in the stress-tensor argument. The
last passage uses equal continuation germs only at spacelike separation;
the current itself is spatially integrated. It gives a charge commutator
for every admissible time profile, without defining a stress field at a
fixed real time.

The low-energy estimate needs only locality, the spectrum condition and
the two vector-norm bounds. Its positive weight avoids differentiating
spectral integrals or imposing generator-domain bounds. The exponent
\(2\) is compatible with the massless scalar measure, since
\(\int_{|\vec p|\le u}d^3p/|\vec p|\sim u^2\). Thus this estimate does
not imply a mass gap. For removal of the exceptional set, any decay rate
tending to zero uniformly in \(\xi\) would suffice.

The remaining integrations, dyadic summation and assembly are routine
once their uniform bounds have been established.

\emph{What looks unusual.}

\begin{itemize}
\tightlist
\item
  Item 8 quantifies over all admissible time profiles and spatial
  cutoffs. Admissible profiles and cutoffs exist
  (\href{https://github.com/yang-mills-group/OpenAI-YM-20261028/blob/3456e5e75cf8d4cd558955a5c7638fba2d68f464/YM/EveryPhase/Charge.lean\#L24-L35}{\texttt{Stress.accepts}},
  \href{https://github.com/yang-mills-group/OpenAI-YM-20261028/blob/3456e5e75cf8d4cd558955a5c7638fba2d68f464/YM/EveryPhase/Charge.lean\#L42-L59}{\texttt{Stress.reservable}});
  this existence assertion is separate from the commutator conclusion.
  The main theorem does not use it or the auxiliary results of Remark
  V.3.3.
\item
  The reconstruction agrees with V.2 at \(\xi_*\). The subsequent proof
  of vacuum uniqueness applies to every phase.
\item
  The Euclidean half of item 9 restates V.1, and the public clause (R12)
  is derived from it.
\item
  Nothing here or in the main theorem decides whether \(X\) is a single
  point. Each \(\mathfrak W_\xi\) is irreducible, while \(S^{\rm mix}\)
  is the \(\nu\)-mixture of the \(S^\xi\).
\end{itemize}

\subsubsection{6. Notation}\label{notation-22}

{\def\LTcaptype{none} % do not increment counter
\begin{longtable}[]{@{}
  >{\raggedright\arraybackslash}p{(\linewidth - 4\tabcolsep) * \real{0.3333}}
  >{\raggedright\arraybackslash}p{(\linewidth - 4\tabcolsep) * \real{0.3333}}
  >{\raggedright\arraybackslash}p{(\linewidth - 4\tabcolsep) * \real{0.3333}}@{}}
\toprule\noalign{}
\begin{minipage}[b]{\linewidth}\raggedright
object
\end{minipage} & \begin{minipage}[b]{\linewidth}\raggedright
code name
\end{minipage} & \begin{minipage}[b]{\linewidth}\raggedright
proposed notation
\end{minipage} \\
\midrule\noalign{}
\endhead
\bottomrule\noalign{}
\endlastfoot
phases, measure, selected phase & \texttt{F.α.space.X},
\texttt{F.α.space.μ}, \texttt{F.α.atPoint} & \(X\), \(\nu\),
\(\xi_*\) \\
characters of the ground algebra & \texttt{F.α.space.character},
\href{https://github.com/yang-mills-group/OpenAI-YM-20261028/blob/3456e5e75cf8d4cd558955a5c7638fba2d68f464/YM/Phase/Mixture.lean\#L54}{\texttt{Spread.algebra}}
& \(\chi_\xi\) on \(\mathfrak A\) \\
unselected table; table of phase \(\xi\) & \texttt{S};
\texttt{F.α.ensemble\ ξ} & \(S^{\rm mix}\); \(S^\xi\) \\
OS space and slab operators of \(S^{\rm mix}\) & \texttt{F.α.room},
\texttt{F.α.unselected} & \(\mathcal H^{\rm OS}\), \(\Pi(d,w)\) \\
ground space, Euclidean vacuum & \texttt{Chamber.D},
\texttt{Chamber.omega0} & \(\mathcal D_\xi\), \(\omega_\xi\) \\
one-dimensional ground space &
\href{https://github.com/yang-mills-group/OpenAI-YM-20261028/blob/3456e5e75cf8d4cd558955a5c7638fba2d68f464/YM/Phase/Mixture.lean\#L35}{\texttt{Chamber.UniqueHeat}}
& \(\mathcal D_\xi=\mathbb C\omega_\xi\) \\
family of reconstructions &
\href{https://github.com/yang-mills-group/OpenAI-YM-20261028/blob/3456e5e75cf8d4cd558955a5c7638fba2d68f464/YM/EveryPhase/Pack.lean\#L117-L132}{\texttt{Almanac}}
&
\((\mathcal H^{\rm OS}_\xi,\Pi_\xi,\mathsf F^\xi,\mathfrak W_\xi,J_\xi)_{\xi\in X}\) \\
OS space, slab operators of \(S^\xi\) & \texttt{L.chambers\ ξ},
\texttt{L.spreads\ ξ} & \(\mathcal H^{\rm OS}_\xi\), \(\Pi_\xi(d,w)\) \\
continuation and boundary values & \texttt{L.contours\ ξ}
(\texttt{waves}, \texttt{world}) & \(\mathsf F^\xi_n\), \(W^\xi_n\) \\
Wightman theory of phase \(\xi\) & \texttt{L.actors\ ξ} &
\(\mathfrak W_\xi=(\mathcal H_\xi,\mathfrak D_\xi,\Omega_\xi,\phi_\xi,U_\xi,P_\xi)\) \\
boundary isometry & \texttt{L.J\ ξ}, \texttt{L.links\ ξ}
(\href{https://github.com/yang-mills-group/OpenAI-YM-20261028/blob/3456e5e75cf8d4cd558955a5c7638fba2d68f464/YM/Next/BOUNDARY/Core.lean\#L156-L161}{\texttt{Attachment}})
& \(J_\xi\) \\
Wightman data without vacuum uniqueness &
\href{https://github.com/yang-mills-group/OpenAI-YM-20261028/blob/3456e5e75cf8d4cd558955a5c7638fba2d68f464/YM/Next/BOUNDARY/Core.lean\#L34-L44}{\texttt{FieldCore}}
& --- \\
heat semigroup on \(\mathcal H_\xi\) &
\href{https://github.com/yang-mills-group/OpenAI-YM-20261028/blob/3456e5e75cf8d4cd558955a5c7638fba2d68f464/YM/EveryPhase/Projection.lean\#L94-L95}{\texttt{bHeat}}
& \(e^{-sH_\xi}=J_\xi^*e^{-sH^{\rm OS}_\xi}J_\xi\) \\
zero-energy projection & \texttt{momentum.region\ zeroE} &
\(P^0_\xi\) \\
finite sum of words; compact supports; adjoint; vector &
\href{https://github.com/yang-mills-group/OpenAI-YM-20261028/blob/3456e5e75cf8d4cd558955a5c7638fba2d68f464/YM/EveryPhase/Polynomials.lean\#L14}{\texttt{SumOfWords}},
\href{https://github.com/yang-mills-group/OpenAI-YM-20261028/blob/3456e5e75cf8d4cd558955a5c7638fba2d68f464/YM/EveryPhase/Polynomials.lean\#L85}{\texttt{finite}},
\href{https://github.com/yang-mills-group/OpenAI-YM-20261028/blob/3456e5e75cf8d4cd558955a5c7638fba2d68f464/YM/EveryPhase/Polynomials.lean\#L42}{\texttt{reverseP}},
\href{https://github.com/yang-mills-group/OpenAI-YM-20261028/blob/3456e5e75cf8d4cd558955a5c7638fba2d68f464/YM/EveryPhase/Polynomials.lean\#L55}{\texttt{stands}}
& \(A\), \(A^\dagger\), \(A\Omega_\xi\) \\
spectral measure of \(A\Omega_\xi\) &
\href{https://github.com/yang-mills-group/OpenAI-YM-20261028/blob/3456e5e75cf8d4cd558955a5c7638fba2d68f464/YM/EveryPhase/Light.lean\#L173-L174}{\texttt{νletter}}
& \(\mu^\xi_A\) \\
low-energy region; shell &
\href{https://github.com/yang-mills-group/OpenAI-YM-20261028/blob/3456e5e75cf8d4cd558955a5c7638fba2d68f464/YM/EveryPhase/Light.lean\#L132}{\texttt{lowE}},
\href{https://github.com/yang-mills-group/OpenAI-YM-20261028/blob/3456e5e75cf8d4cd558955a5c7638fba2d68f464/YM/EveryPhase/Light.lean\#L40}{\texttt{ringE}}
& \(\{0<p^4\le u\}\); \(\{u/2<p^4\le u\}\) \\
vacuum commutator function &
\href{https://github.com/yang-mills-group/OpenAI-YM-20261028/blob/3456e5e75cf8d4cd558955a5c7638fba2d68f464/YM/EveryPhase/Polynomials.lean\#L171-L173}{\texttt{delta}}
& \(\langle\Omega_\xi,[A^\dagger,\alpha_pA]\Omega_\xi\rangle\) \\
time-averaged odd part; weight &
\href{https://github.com/yang-mills-group/OpenAI-YM-20261028/blob/3456e5e75cf8d4cd558955a5c7638fba2d68f464/YM/EveryPhase/Average.lean\#L88-L89}{\texttt{averaged}},
\texttt{pulse} & \(D(\vec x)\); \(\pi(0)=0\), \(\pi(E)=(1-\cos E)/E\)
for \(E>0\) \\
strict-height values; tested values & \texttt{waves.restricted},
\texttt{waves.psi} & \(\mathsf F^\xi_n(a;t(x)+iy,\vec x)\);
\(\psi^\xi(a,y,g)\) \\
weak commutant &
\href{https://github.com/yang-mills-group/OpenAI-YM-20261028/blob/3456e5e75cf8d4cd558955a5c7638fba2d68f464/YM/EveryPhase/Extension.lean\#L19-L22}{\texttt{member}},
\href{https://github.com/yang-mills-group/OpenAI-YM-20261028/blob/3456e5e75cf8d4cd558955a5c7638fba2d68f464/YM/EveryPhase/Irreducible.lean\#L93-L96}{\texttt{weaklyCentral}}
& --- \\
Minkowski stress labels and fields &
\href{https://github.com/yang-mills-group/OpenAI-YM-20261028/blob/3456e5e75cf8d4cd558955a5c7638fba2d68f464/YM/EveryPhase/Tensor.lean\#L27}{\texttt{Stress.M}},
\href{https://github.com/yang-mills-group/OpenAI-YM-20261028/blob/3456e5e75cf8d4cd558955a5c7638fba2d68f464/YM/EveryPhase/Tensor.lean\#L206}{\texttt{Stress.fieldM}}
& \(T^{\rm M}_{ab}\), \(T_{ab}(f)\) \\
Killing field; charge; Lie derivative of a field &
\texttt{Stress.Killing},
\href{https://github.com/yang-mills-group/OpenAI-YM-20261028/blob/3456e5e75cf8d4cd558955a5c7638fba2d68f464/YM/EveryPhase/Charge.lean\#L80-L81}{\texttt{Stress.Charge}},
\href{https://github.com/yang-mills-group/OpenAI-YM-20261028/blob/3456e5e75cf8d4cd558955a5c7638fba2d68f464/YM/EveryPhase/Charge.lean\#L86-L88}{\texttt{Stress.action}}
& \(K\), \(Q_K(f,c)\), \(\phi(\Sigma_KB)+K^a\partial_a\phi(B)\) \\
stress record &
\href{https://github.com/yang-mills-group/OpenAI-YM-20261028/blob/3456e5e75cf8d4cd558955a5c7638fba2d68f464/YM/EveryPhase/Implement.lean\#L81-L92}{\texttt{Stress.Force}}
(fields appointment, orbital, symmetric, hermitian, conserve,
commutator, rule) & (T0)--(T6) \\
complete record &
\href{https://github.com/yang-mills-group/OpenAI-YM-20261028/blob/3456e5e75cf8d4cd558955a5c7638fba2d68f464/YM/EveryPhase/Extension.lean\#L51-L71}{\texttt{Completed}}
(fields isTheory, into\_onto, euclid, mink, light, heat, commutant,
continued\_stress, depends, at\_boundary, at\_words) & items 1, 2, 3, 4,
5, 6, 7, 8, 9, 9, 9 of V.3.1 \\
\end{longtable}
}

\subsubsection{7. What remains
unexplained}\label{what-remains-unexplained-22}

\begin{enumerate}
\def\labelenumi{\arabic{enumi}.}
\tightlist
\item
  \textbf{Analytic inputs from V.2.} The constructions
  \href{https://github.com/yang-mills-group/OpenAI-YM-20261028/blob/3456e5e75cf8d4cd558955a5c7638fba2d68f464/YM/Next/BOUNDARY/Stage.lean\#L46-L49}{\texttt{contour}},
  \href{https://github.com/yang-mills-group/OpenAI-YM-20261028/blob/3456e5e75cf8d4cd558955a5c7638fba2d68f464/YM/Next/BOUNDARY/Stage.lean\#L200-L215}{\texttt{compatible}},
  \href{https://github.com/yang-mills-group/OpenAI-YM-20261028/blob/3456e5e75cf8d4cd558955a5c7638fba2d68f464/YM/Next/BOUNDARY/Stage.lean\#L167-L184}{\texttt{establishing}}
  and their use here were checked in source. The underlying C01--C29
  analytic estimates, including the partial-boundary identities used in
  V.3.i and the exchange continuations in (T0), remain inputs from V.2;
  their proofs are not audited here. The local statements and proof
  sketches have been checked against the declarations that consume them.
\item
  \textbf{Partial-boundary data.} (T0) states the existence and
  compatibility needed for the current-jump argument. A detailed proof
  must specify the admissible tube faces and the spacelike neighborhoods
  of the tied pair; it must not replace integrated current limits by
  pointwise limits of the stress insertion.
\end{enumerate}

\subsection{VI.1. Ward identities, the rotation action and the conserved
stress
tensor}\label{vi.1.-ward-identities-the-rotation-action-and-the-conserved-stress-tensor}

\textbf{Part:} VI, symmetry: stress tensor and Ward identities.
\textbf{Lean:}

\begin{itemize}
\tightlist
\item
  the work items A01--A17, stated in
  \href{https://github.com/yang-mills-group/OpenAI-YM-20261028/blob/3456e5e75cf8d4cd558955a5c7638fba2d68f464/YM/Bid/Angular.lean}{\texttt{YM/Bid/Angular.lean}},
  combined in
  \href{https://github.com/yang-mills-group/OpenAI-YM-20261028/blob/3456e5e75cf8d4cd558955a5c7638fba2d68f464/YM/Bid/Town.lean}{\texttt{YM/Bid/Town.lean}}
  and proved in
  \href{https://github.com/yang-mills-group/OpenAI-YM-20261028/tree/3456e5e75cf8d4cd558955a5c7638fba2d68f464/YM/Work}{\texttt{YM/Work/A01}--\texttt{A17}};
\item
  the collinear flux estimate, in
  \href{https://github.com/yang-mills-group/OpenAI-YM-20261028/tree/3456e5e75cf8d4cd558955a5c7638fba2d68f464/YM/Next/ROTATIONS}{\texttt{YM/Next/ROTATIONS}}
  with
  \href{https://github.com/yang-mills-group/OpenAI-YM-20261028/tree/3456e5e75cf8d4cd558955a5c7638fba2d68f464/YM/Axis}{\texttt{YM/Axis}};
\item
  the analytic continuation, in
  \href{https://github.com/yang-mills-group/OpenAI-YM-20261028/tree/3456e5e75cf8d4cd558955a5c7638fba2d68f464/YM/Ward}{\texttt{YM/Ward}};
\item
  the scaling comparison, in
  \href{https://github.com/yang-mills-group/OpenAI-YM-20261028/tree/3456e5e75cf8d4cd558955a5c7638fba2d68f464/YM/Climb}{\texttt{YM/Climb}};
\item
  the descent of the rotation action, in
  \href{https://github.com/yang-mills-group/OpenAI-YM-20261028/tree/3456e5e75cf8d4cd558955a5c7638fba2d68f464/YM/Spin}{\texttt{YM/Spin}};
\item
  the assembled record, in
  \href{https://github.com/yang-mills-group/OpenAI-YM-20261028/blob/3456e5e75cf8d4cd558955a5c7638fba2d68f464/YM/Org/Twist.lean}{\texttt{YM/Org/Twist.lean}};
\item
  in part,
  \href{https://github.com/yang-mills-group/OpenAI-YM-20261028/tree/3456e5e75cf8d4cd558955a5c7638fba2d68f464/YM/Push/STRESS_FLUX}{\texttt{YM/Push/STRESS\_FLUX}},
  \href{https://github.com/yang-mills-group/OpenAI-YM-20261028/tree/3456e5e75cf8d4cd558955a5c7638fba2d68f464/YM/Shift}{\texttt{YM/Shift}}
  and
  \href{https://github.com/yang-mills-group/OpenAI-YM-20261028/tree/3456e5e75cf8d4cd558955a5c7638fba2d68f464/YM/Next/BOUNDARY}{\texttt{YM/Next/BOUNDARY}}.
\end{itemize}

\textbf{Size (approximate):}

\begin{itemize}
\tightlist
\item
  7,704 lines are assigned to this unit and used by the main theorem. Of
  these, about 900 (\texttt{Shift} and most of \texttt{Next/BOUNDARY})
  are reconstruction material for V.2, and about 1,000
  (\texttt{STRESS\_FLUX}) are free-field algebra used in VII.2 (Remark
  in §4.10).
\item
  About 6,100 further lines of \texttt{Next/ROTATIONS} carry the
  collinear flux estimate. The declaration table files them under III.6,
  but they belong here, and the III.6 note says so too.
\end{itemize}

\textbf{Conventions.}

\begin{itemize}
\tightlist
\item
  \emph{Table.} \(S\) is the limit table of IV.1. For labels
  \(v_i\in F_{d_i}\) and \(x\in\mathrm{Conf}_n\) (pairwise distinct
  points of \(\mathbb R^4\)), \(S(v_1,\dots,v_n;x_1,\dots,x_n)\) is a
  correlation kernel, smooth on \(\mathrm{Conf}_n\).
\item
  \emph{Label spaces.} \(F_d\) is the finite-dimensional complex space
  of labels of degree \(\le d\) modulo null labels. It has inclusions
  \(\iota=\iota_{d\to e}\), a unit \(\mathbf 1_d\), the action
  \(M_\sigma\) of the signed permutation matrices
  \(\sigma\in B_4\subset O(4)\), and an antilinear conjugation
  \(\Gamma\).
\item
  \emph{Euclidean axioms}:

  \begin{itemize}
  \tightlist
  \item
    permutation symmetry, normalisation of the empty correlation,
    deletion of unit insertions, and translation invariance;
  \item
    covariance under \(B_4\) and \(\Gamma\);
  \item
    reflection positivity in the coordinate hyperplanes, for point words
    and for smeared tests;
  \item
    smoothness on configuration space;
  \item
    for each fixed arity, labels and derivative order \(j\), constants
    \(C\ge0\) and \(A\in\mathbb N\) such that \(\|D^jS(x)\|\le Cr^{-A}\)
    whenever \(0<r\le1\) and the points are separated by at least \(r\)
    in the maximum norm. The constants are uniform in the positions.
  \end{itemize}
\item
  \emph{Differential plan.} Complex linear maps
  \(\partial_a:F_d\to F_{d+1}\) satisfying three laws:

  \begin{itemize}
  \tightlist
  \item
    the \emph{derivative law}
    \(S(\dots,\partial_av_k,\dots;x)=\partial_{x_k^a}S(\dots,v_k,\dots;x)\);
  \item
    the \emph{inclusion law}
    \(S(\dots,\iota v_k,\dots)=S(\dots,v_k,\dots)\);
  \item
    \emph{orientation}: \(\Gamma\) and \(M_\sigma\) commute with the
    inclusions, \(\partial\) commutes with \(\Gamma\), and
    \(M_\sigma\partial_b=\sum_a\sigma_{ab}\,\partial_aM_\sigma\).
    Commutation of \(\partial\) with the inclusions follows from the
    inclusion and derivative laws.
  \end{itemize}
\item
  \emph{Rotations.} \(\mathfrak{so}(4)\) is the space of real skew
  \(4\times4\) matrices with \([K,K']=KK'-K'K\).
  \(K^{(ba)}=E_{ba}-E_{ab}\) generates the rotation field
  \(x_ae_b-x_be_a\).
\item
  \emph{Spin rows.} A system of spin rows is a family of
  \(\mathbb R\)-linear maps
  \(T:\mathfrak{so}(4)\to\mathrm{End}_{\mathbb C}F_d\).
\item
  \emph{Angular sum.} For spin rows \(T\) and maps
  \(P_b:F_d\to F_{d+1}\),
\end{itemize}

\[\mathcal W_K[T,P](https://github.com/scottnarmstrong/yang-mills-informalization/blob/main/outline/units/q;x)=\sum_{i=1}^n\Big[S(\dots,T_Kq_i,\dots;x)+\sum_b(Kx_i)_b\,S(\dots,P_bq_i,\dots;x)\Big],\qquad \mathcal W_K[T]:=\mathcal W_K[T,\partial].\]
By the derivative law, \(\mathcal W_K[T]=0\) says
\(\frac{d}{dt}\big|_{0}S(q;e^{tK}x)=-\sum_iS(\dots,T_Kq_i,\dots;x)\).
This is infinitesimal rotation covariance with spin part \(T\). -
\emph{Lorentz group.} Axes are numbered \(0,1,2,3\); axis \(3\) is time
and \(0,1,2\) are spatial. Write \(e_t=e_3\). The axis \(0\) used for
the continuation in §4.4 has no distinguished physical role. - All
identities hold at separated points. No contact term is evaluated.

\subsubsection{1. What this unit proves}\label{what-this-unit-proves-23}

\textbf{Inputs.}

\begin{itemize}
\tightlist
\item
  The table and its differential plan.
\item
  An unnormalised stress \(J=(J_{ab})\), \(J_{ab}\in F_4\), with complex
  linear \emph{raw translation maps} \(p_a:F_d\to F_{d+1}\) and real
  linear \emph{raw spin rows}
  \(Y:\mathfrak{so}(4)\to\mathrm{End}\,F_d\). These satisfy the shell
  identities below with unknown right-hand sides.
\item
  Free-Maxwell comparison data: a free table with its differential plan
  \(\partial^{\rm free}\), rotation action \(R^{\rm free}\) with
  generators \(\mathcal L^{\rm free}_K\), stress \(J^{\rm free}\) and
  Ward constant \(\kappa\ne0\).
\item
  A weak scaling comparison of \(S\) with the free table.
\end{itemize}

The table and differential plan come from IV.1. The stress, shell maps
and comparison data come from VI.2, using ultraviolet inputs as
specified in §2. Every result below assumes the Euclidean axioms for
\(S\), unless it is a statement solely about groups or operators.

\textbf{Outputs.} Given the differential plan and the comparison inputs,
the unit constructs:

\begin{itemize}
\tightlist
\item
  an action of \(O(4)\) on every \(F_d\) whose generators are
  \(c^{-1}Y\);
\item
  \(O(4)\)-covariance of \(S\);
\item
  a conserved, symmetric, real, \(O(4)\)-tensorial stress \(J'\)
  satisfying the Ward identity (13) of the assessment and its rotation
  analogue;
\item
  a Lorentz action on the labels, compatible with the rotation action.
\end{itemize}

\textbf{Shells and stress laws.} A \emph{window} for
\(x=(x_0,\dots,x_n)\in\mathrm{Conf}_{n+1}\) is a real function
\(\chi\in C^\infty_c(\mathbb R^4)\) that equals \(1\) near \(x_0\) and
whose closed support excludes \(x_1,\dots,x_n\). For \(B\in F_d\) and
tags \(q\) at \(x_1,\dots,x_n\), the two shells are

\[\Phi^{\rm tr}_\chi(b)=-\sum_a\int\partial_a\chi(y)\,S\big(J_{ab}(y),B(x_0),q\big)\,dy,\qquad \Phi^{\rm rot}_\chi(K)=-\sum_{a,b}\int\partial_a\chi(y)\,\big(K(y-x_0)\big)_b\,S\big(J_{ab}(y),B(x_0),q\big)\,dy .\]

The integrands vanish near all insertions. Given maps
\(P_b:F_d\to F_{d+1}\) and \(A:\mathfrak{so}(4)\to\mathrm{End}\,F_d\),
the triple \((J;P,A)\) satisfies the \emph{stress laws} if

\[\Phi^{\rm tr}_\chi(b)=S(P_bB(x_0),q)\qquad\text{and}\qquad\Phi^{\rm rot}_\chi(K)=S(A_KB(x_0),q)\]

for all windows, all \(B\) and all \(q\)
(\href{https://github.com/yang-mills-group/OpenAI-YM-20261028/blob/3456e5e75cf8d4cd558955a5c7638fba2d68f464/YM/Euclid/Stress.lean\#L90-L97}{\texttt{Table.StressLaws}},
\href{https://github.com/yang-mills-group/OpenAI-YM-20261028/blob/3456e5e75cf8d4cd558955a5c7638fba2d68f464/YM/Euclid/Stress.lean\#L79-L81}{\texttt{shellGrad}},
\href{https://github.com/yang-mills-group/OpenAI-YM-20261028/blob/3456e5e75cf8d4cd558955a5c7638fba2d68f464/YM/Euclid/Stress.lean\#L84-L86}{\texttt{shellSpin}},
\href{https://github.com/yang-mills-group/OpenAI-YM-20261028/blob/3456e5e75cf8d4cd558955a5c7638fba2d68f464/YM/Euclid/Stress.lean\#L72-L73}{\texttt{Window}}).

\textbf{Raw hypotheses (R)}
(\href{https://github.com/yang-mills-group/OpenAI-YM-20261028/blob/3456e5e75cf8d4cd558955a5c7638fba2d68f464/YM/Climb/Beams.lean\#L66-L69}{\texttt{Climb.Pre}}):

\begin{itemize}
\tightlist
\item
  \(J\) is \emph{conserved}: symmetric, \(\Gamma\)-real, and
  \(\sum_a\partial_aJ_{ab}=0\) in \(F_5\)
  (\href{https://github.com/yang-mills-group/OpenAI-YM-20261028/blob/3456e5e75cf8d4cd558955a5c7638fba2d68f464/YM/Euclid/Stress.lean\#L59-L62}{\texttt{Conserved}});
\item
  \(J\) is \emph{raw-tensorial}:
  \(M_\sigma J_{ab}=\sum_{ij}\sigma_{ia}\sigma_{jb}J_{ij}\) for
  \(\sigma\in B_4\)
  (\href{https://github.com/yang-mills-group/OpenAI-YM-20261028/blob/3456e5e75cf8d4cd558955a5c7638fba2d68f464/YM/Euclid/Stress.lean\#L68-L70}{\texttt{RawTensor}});
\item
  \((J;p,Y)\) satisfies the stress laws.
\end{itemize}

\textbf{Proposition VI.1.1 (collinear rotation Ward identity and line
sums).}

\emph{Hypotheses.}

\begin{itemize}
\tightlist
\item
  \(S\) satisfies the Euclidean axioms.
\item
  \(\partial\) satisfies the derivative law and orientation.
\item
  \((J,p,Y)\) satisfies (R).
\end{itemize}

\emph{Conclusions.}

\begin{itemize}
\tightlist
\item
  \begin{enumerate}
  \def\labelenumi{(\roman{enumi})}
  \tightlist
  \item
    Let \(x\in\mathrm{Conf}_n\) have all its points on a coordinate axis
    \(\mathbb Re_l\). Then for every \(K\in\mathfrak{so}(4)\) and all
    tags \(q\):
    \(\mathcal W_K[Y,p](https://github.com/scottnarmstrong/yang-mills-informalization/blob/main/outline/units/q;x)=0\).
  \end{enumerate}
\item
  (ii) (Line sums.) Let \(x\in\mathrm{Conf}_n\) lie on \(\mathbb Re_l\),
  let \(j\neq l\), and let \(v_1,\dots,v_n\in F_d\). Then
  \(\sum_iS(v_1,\dots,p_jv_i,\dots,v_n;x)=0\).
\end{itemize}

\emph{Status:} theorem.

\begin{itemize}
\tightlist
\item
  It answers the request
  \href{https://github.com/yang-mills-group/OpenAI-YM-20261028/blob/3456e5e75cf8d4cd558955a5c7638fba2d68f464/YM/Bid/Town.lean\#L13-L15}{\texttt{Bid.Miles}}
  through
  \href{https://github.com/yang-mills-group/OpenAI-YM-20261028/blob/3456e5e75cf8d4cd558955a5c7638fba2d68f464/YM/Next/ROTATIONS/Route.lean\#L123}{\texttt{Outer.miles}}.
\item
  The analytic input is the exterior flux estimate (Lemma VI.1.e). It
  uses reflection positivity and the polynomial bounds, but no mass gap
  or clustering.
\end{itemize}

\emph{Lean:}
\href{https://github.com/yang-mills-group/OpenAI-YM-20261028/blob/3456e5e75cf8d4cd558955a5c7638fba2d68f464/YM/Next/ROTATIONS/Route.lean\#L114-L121}{\texttt{Outer.outer}},
\href{https://github.com/yang-mills-group/OpenAI-YM-20261028/blob/3456e5e75cf8d4cd558955a5c7638fba2d68f464/YM/Bid/Angular.lean\#L35-L36}{\texttt{Bid.LineSum}},
\href{https://github.com/yang-mills-group/OpenAI-YM-20261028/blob/3456e5e75cf8d4cd558955a5c7638fba2d68f464/YM/Climb/Beams.lean\#L72-L74}{\texttt{Climb.localTurn}}.

\textbf{Proposition VI.1.2 (normalisation of the shell maps:
\(p=c\,\partial\)).}

\emph{Hypotheses.}

\begin{itemize}
\tightlist
\item
  \(\partial\) satisfies all three laws.
\item
  The ultraviolet record holds, in particular (R) and the weak
  comparison law against the free data (input 4 of §2).
\item
  The free-Maxwell data hold (input 3 of §2).
\end{itemize}

\emph{Conclusion.} There is a real \(c\ne0\), independent of the degree
and coordinate axis, such that:

\begin{itemize}
\tightlist
\item
  \(p_a=c\,\partial_a\) on every \(F_d\), for every \(a\);
\item
  the rescaled rows \(D:=c^{-1}Y\) satisfy

  \begin{itemize}
  \tightlist
  \item
    (C) \(\partial_bD_K+\sum_iK_{ib}\,\partial_i=D_K\,\partial_b\) on
    every \(F_d\);
  \item
    (L)
    \(\mathcal W_K[D](https://github.com/scottnarmstrong/yang-mills-informalization/blob/main/outline/units/q;x)=0\)
    for every configuration \(x\) on a coordinate axis.
  \end{itemize}
\end{itemize}

\emph{Status:} theorem.

\emph{Lean:}
\href{https://github.com/yang-mills-group/OpenAI-YM-20261028/blob/3456e5e75cf8d4cd558955a5c7638fba2d68f464/YM/Bid/Town.lean\#L49-L59}{\texttt{Bid.pressured}},
\href{https://github.com/yang-mills-group/OpenAI-YM-20261028/blob/3456e5e75cf8d4cd558955a5c7638fba2d68f464/YM/Bid/Town.lean\#L61-L78}{\texttt{Bid.formed}};
the record is
\href{https://github.com/yang-mills-group/OpenAI-YM-20261028/blob/3456e5e75cf8d4cd558955a5c7638fba2d68f464/YM/Climb/Beams.lean\#L83-L87}{\texttt{Climb.Crystal}}.

\textbf{Proposition VI.1.3 (rotation Ward identity at all separated
configurations).}

\emph{Hypotheses.}

\begin{itemize}
\tightlist
\item
  \(S\) satisfies the Euclidean axioms.
\item
  \(\partial\) satisfies the derivative law.
\item
  The spin rows \(T\) satisfy (C) and (L).
\end{itemize}

\emph{Conclusion.}
\(\mathcal W_K[T](https://github.com/scottnarmstrong/yang-mills-informalization/blob/main/outline/units/q;x)=0\)
for all \(n\), all \(x\in\mathrm{Conf}_n\), all tags \(q\) and all
\(K\).

\emph{Status:} theorem.

\emph{Lean:}
\href{https://github.com/yang-mills-group/OpenAI-YM-20261028/blob/3456e5e75cf8d4cd558955a5c7638fba2d68f464/YM/Ward/Axis.lean\#L130-L139}{\texttt{Ward.global}};
the vanishing is
\href{https://github.com/yang-mills-group/OpenAI-YM-20261028/blob/3456e5e75cf8d4cd558955a5c7638fba2d68f464/YM/Ward/Walk.lean\#L45}{\texttt{Ward.Still}};
the hypotheses are
\href{https://github.com/yang-mills-group/OpenAI-YM-20261028/blob/3456e5e75cf8d4cd558955a5c7638fba2d68f464/YM/Ward/Sum.lean\#L29-L30}{\texttt{Ward.commutes}}
and
\href{https://github.com/yang-mills-group/OpenAI-YM-20261028/blob/3456e5e75cf8d4cd558955a5c7638fba2d68f464/YM/Ward/Sum.lean\#L34-L35}{\texttt{Ward.Collinear}};
the sum is
\href{https://github.com/yang-mills-group/OpenAI-YM-20261028/blob/3456e5e75cf8d4cd558955a5c7638fba2d68f464/YM/Ward/Sum.lean\#L25-L27}{\texttt{Ward.angular}}.

\textbf{Proposition VI.1.4 (\(\mathfrak{so}(4)\) brackets; A08).}

\emph{Hypotheses.}

\begin{itemize}
\tightlist
\item
  \(\partial\) satisfies the inclusion and derivative laws.
\item
  \(T\) satisfies the conclusion of VI.1.3.
\item
  \(T\) is \emph{compatible}
  (\href{https://github.com/yang-mills-group/OpenAI-YM-20261028/blob/3456e5e75cf8d4cd558955a5c7638fba2d68f464/YM/Bid/Angular.lean\#L16-L20}{\texttt{Bid.Hyper}}):
  it commutes with \(\Gamma\) and with \(\iota\), kills \(\mathbf 1\),
  and \(M_\sigma T_K=T_{\sigma K\sigma^{-1}}M_\sigma\).
\item
  \emph{Internal row erasure} holds
  (\href{https://github.com/yang-mills-group/OpenAI-YM-20261028/blob/3456e5e75cf8d4cd558955a5c7638fba2d68f464/YM/Org/Input.lean\#L149-L152}{\texttt{Org.erases}}).
  That is, for some \(N\) and all \(d\ge N\), let
  \(\Delta:\mathfrak{so}(4)\times\mathfrak{so}(4)\to\mathrm{End}\,F_d\)
  be a real bilinear family of complex linear endomorphisms with these
  properties
  (\href{https://github.com/yang-mills-group/OpenAI-YM-20261028/blob/3456e5e75cf8d4cd558955a5c7638fba2d68f464/YM/Org/Input.lean\#L141-L146}{\texttt{Org.Winds}}):

  \begin{itemize}
  \tightlist
  \item
    it is covariant under signed permutations of determinant \(1\);
  \item
    it is \(\Gamma\)-real;
  \item
    it kills \(\mathbf 1\);
  \item
    it is alternating;
  \item
    for every \(e\le d\) and \(u,v\), \(\Delta(u,v)\) preserves
    \(\iota_{e\to d}(F_e)\).
  \end{itemize}

  If moreover \(\sum_iS(w_1,\dots,\Delta(u,v)w_i,\dots,w_m;x)=0\) for
  all \(m\), all \(x\in\mathrm{Conf}_m\), all \(u,v\in\mathfrak{so}(4)\)
  and all \(w\in F_d^m\), then \(\Delta=0\).
\end{itemize}

\emph{Conclusion.} \([T_u,T_v]=T_{[u,v]}\) on every \(F_d\).

\emph{Status:} theorem. Erasure for \(S\) is itself proved (A07, Lemma
VI.1.q).

\emph{Lean:}
\href{https://github.com/yang-mills-group/OpenAI-YM-20261028/blob/3456e5e75cf8d4cd558955a5c7638fba2d68f464/YM/Work/A08.lean\#L232-L247}{\texttt{Work.A08.answer}}
for
\href{https://github.com/yang-mills-group/OpenAI-YM-20261028/blob/3456e5e75cf8d4cd558955a5c7638fba2d68f464/YM/Bid/Angular.lean\#L76-L78}{\texttt{Bid.A08}};
the conclusion is
\href{https://github.com/yang-mills-group/OpenAI-YM-20261028/blob/3456e5e75cf8d4cd558955a5c7638fba2d68f464/YM/Bid/Angular.lean\#L74}{\texttt{Bid.Brackets}}.

\textbf{Proposition VI.1.5 (descent of the free rotation action).} Under
the hypotheses of VI.1.2, let \(c\) be its normalisation constant and
put \(D=c^{-1}Y\).

\emph{Conclusion.} There are complex linear isomorphisms
\(V_d:F^{\rm free}_d\to F_d\) with these properties:

\begin{itemize}
\tightlist
\item
  \(V_dM^{\rm free}_\sigma=M_\sigma V_d\) and
  \(V_d\Gamma^{\rm free}=\Gamma V_d\);
\item
  \(D_KV_d=V_d\,\mathcal L^{\rm free}_K\) for all
  \(K\in\mathfrak{so}(4)\).
\end{itemize}

Hence \(R_d(k):=V_d\,R^{\rm free}_d(k)\,V_d^{-1}\), \(k\in O(4)\), is a
continuous representation of \(O(4)\) on \(F_d\). It equals \(M_\sigma\)
on \(B_4\), commutes with \(\Gamma\), and has infinitesimal generator
\(D\). Call such an action, with its specified generators, a
\emph{rotation port}. The Ward identity of VI.1.3 holds for \(D\).

\emph{Status:} theorem. The \(V_d\) are chosen independently for each
\(d\). Compatibility with the inclusions comes later (VI.1.6).

\emph{Lean:}
\href{https://github.com/yang-mills-group/OpenAI-YM-20261028/blob/3456e5e75cf8d4cd558955a5c7638fba2d68f464/YM/Climb/Beams.lean\#L90-L98}{\texttt{Climb.descends}},
\href{https://github.com/yang-mills-group/OpenAI-YM-20261028/blob/3456e5e75cf8d4cd558955a5c7638fba2d68f464/YM/Spin/Lands.lean\#L26-L47}{\texttt{Spin.produces}},
\href{https://github.com/yang-mills-group/OpenAI-YM-20261028/blob/3456e5e75cf8d4cd558955a5c7638fba2d68f464/YM/Spin/Sheets.lean\#L134-L157}{\texttt{Spin.meet}},
\href{https://github.com/yang-mills-group/OpenAI-YM-20261028/blob/3456e5e75cf8d4cd558955a5c7638fba2d68f464/YM/Bid/Town.lean\#L124-L157}{\texttt{Bid.services}};
the transported action is
\href{https://github.com/yang-mills-group/OpenAI-YM-20261028/blob/3456e5e75cf8d4cd558955a5c7638fba2d68f464/YM/Spin/Lands.lean\#L22-L24}{\texttt{Spin.morph}}
and a port is
\href{https://github.com/yang-mills-group/OpenAI-YM-20261028/blob/3456e5e75cf8d4cd558955a5c7638fba2d68f464/YM/Spin/Lands.lean\#L13-L19}{\texttt{Spin.Port}}.

\textbf{Proposition VI.1.6 (covariance and the rotation laws; A14).}

\emph{Hypotheses.}

\begin{itemize}
\tightlist
\item
  \(\partial\) satisfies the inclusion and derivative laws.
\item
  \(R\) is a port with generators \(T\).
\item
  \(\mathcal W_K[T]=0\) on all of \(\mathrm{Conf}\).
\end{itemize}

\emph{Conclusions.}

\begin{itemize}
\tightlist
\item
  \(R\) is \emph{oriented}:
  \(S(R(k)v_1,\dots,R(k)v_n;kx_1,\dots,kx_n)=S(v;x)\) for all
  \(k\in O(4)\) and all \(x\in\mathrm{Conf}_n\).
\item
  \(R\) satisfies the \emph{rotation laws}: it fixes \(\mathbf 1\); it
  commutes with \(\iota\); it acts on derivatives as a vector,
  \(R_{d+1}(k)\,\partial_av=\sum_ik_{ia}\,\partial_iR_d(k)v\); and it
  keeps the port properties.
\end{itemize}

\emph{Status:} theorem. It uses A13: \(O(4)\) is generated by \(B_4\)
and the one-parameter groups \(e^{tK}\).

\emph{Lean:}
\href{https://github.com/yang-mills-group/OpenAI-YM-20261028/blob/3456e5e75cf8d4cd558955a5c7638fba2d68f464/YM/Work/A14.lean\#L417-L420}{\texttt{Work.A14.answer}}
for
\href{https://github.com/yang-mills-group/OpenAI-YM-20261028/blob/3456e5e75cf8d4cd558955a5c7638fba2d68f464/YM/Bid/Angular.lean\#L105-L107}{\texttt{Bid.A14}};
\href{https://github.com/yang-mills-group/OpenAI-YM-20261028/blob/3456e5e75cf8d4cd558955a5c7638fba2d68f464/YM/Euclid/Stress.lean\#L31-L42}{\texttt{Rotation.Laws}},
\href{https://github.com/yang-mills-group/OpenAI-YM-20261028/blob/3456e5e75cf8d4cd558955a5c7638fba2d68f464/YM/Euclid/Stress.lean\#L46-L48}{\texttt{Rotation.Oriented}};
\href{https://github.com/yang-mills-group/OpenAI-YM-20261028/blob/3456e5e75cf8d4cd558955a5c7638fba2d68f464/YM/Work/A13.lean\#L191-L200}{\texttt{Work.A13.answer}}.

\textbf{Proposition VI.1.7 (the averaged stress tensor and the Ward
identity (13); A15).}

\emph{Hypotheses.}

\begin{itemize}
\tightlist
\item
  \(\partial\) satisfies the derivative law.
\item
  \(R\) is a rotation action satisfying the rotation laws, and \(R\) is
  oriented.
\item
  \(p=c\,\partial\) with \(c\neq0\), and \((J,p,Y)\) satisfies (R).
\item
  The generators of \(R\) are \(c^{-1}Y\).
\end{itemize}

\emph{Conclusion.} There is a stress \(J'\) with these properties:

\begin{itemize}
\tightlist
\item
  it is conserved, i.e.~symmetric, real and divergence-free;
\item
  it is \emph{tensorial}: \(R_4(k)J'_{ab}=\sum_{ij}k_{ia}k_{jb}J'_{ij}\)
  for all \(k\in O(4)\);
\item
  \((J';\partial,c^{-1}Y)\) satisfies the stress laws.
\end{itemize}

The translation law is the Ward identity (13):

\[-\sum_a\int\partial_a\chi(y)\,S\big(J'_{ab}(y),B(x_0),q\big)\,dy=\frac{\partial}{\partial x_0^b}S\big(B(x_0),q\big).\]

\emph{Status:} theorem. It uses A10.

\emph{Lean:}
\href{https://github.com/yang-mills-group/OpenAI-YM-20261028/blob/3456e5e75cf8d4cd558955a5c7638fba2d68f464/YM/Work/A15.lean\#L187-L192}{\texttt{Work.A15.answer}}
for
\href{https://github.com/yang-mills-group/OpenAI-YM-20261028/blob/3456e5e75cf8d4cd558955a5c7638fba2d68f464/YM/Bid/Angular.lean\#L109-L112}{\texttt{Bid.A15}};
\href{https://github.com/yang-mills-group/OpenAI-YM-20261028/blob/3456e5e75cf8d4cd558955a5c7638fba2d68f464/YM/Euclid/Stress.lean\#L64-L66}{\texttt{Table.Tensorial}}.

\textbf{Proposition VI.1.8 (Lorentz action on the labels).} For
(b)--(c), assume the hypotheses of VI.1.2 and choose \(V_d\), \(R\) and
\(J'\) as in VI.1.5--7.

\emph{(a)} \(SO^+(1,3)\) is generated by the spatial rotations and the
coordinate boosts (A16).

\emph{(b)} Let \(B^{\rm free}\) be the Lorentz action on the free
labels, which is part of the free data. Let
\(B_d(\Lambda)=V_dB^{\rm free}_d(\Lambda)V_d^{-1}\). Then
\(B_e(\Lambda)\,\iota=\iota\,B_d(\Lambda)\) for all \(d\le e\) and
\(\Lambda\in SO^+(1,3)\) (A17).

\emph{(c)} The compatible actions \(B_d\) define a representation on the
union of the label spaces. For every degree \(d\) and spatial axis
\(i\):

\begin{itemize}
\tightlist
\item
  the action on \(F_d\) is continuous;
\item
  \(\frac{d}{ds}\big|_{s=0}B_d(\mathrm{boost}_i(s))v=-i\,D_{K^{(3i)}}v\)
  for every \(v\in F_d\);
\item
  \(B_d\) equals \(R_d\) on spatial rotations;
\item
  \(B_d\) commutes with the time-reflection adjoint
  \(M_{\theta_t}\Gamma\), where \(\theta_t\) reverses the time
  coordinate.
\end{itemize}

\emph{Status:} theorems.

\emph{Lean:}
\href{https://github.com/yang-mills-group/OpenAI-YM-20261028/blob/3456e5e75cf8d4cd558955a5c7638fba2d68f464/YM/Work/A16.lean\#L91-L101}{\texttt{Work.A16.answer}}
for
\href{https://github.com/yang-mills-group/OpenAI-YM-20261028/blob/3456e5e75cf8d4cd558955a5c7638fba2d68f464/YM/Bid/Angular.lean\#L116}{\texttt{Bid.A16}};
\href{https://github.com/yang-mills-group/OpenAI-YM-20261028/blob/3456e5e75cf8d4cd558955a5c7638fba2d68f464/YM/Work/A17.lean\#L88-L101}{\texttt{Work.A17.answer}}
for
\href{https://github.com/yang-mills-group/OpenAI-YM-20261028/blob/3456e5e75cf8d4cd558955a5c7638fba2d68f464/YM/Bid/Angular.lean\#L118-L119}{\texttt{Bid.A17}};
\href{https://github.com/yang-mills-group/OpenAI-YM-20261028/blob/3456e5e75cf8d4cd558955a5c7638fba2d68f464/YM/Spin/Floor.lean\#L24-L25}{\texttt{Spin.wand}},
\href{https://github.com/yang-mills-group/OpenAI-YM-20261028/blob/3456e5e75cf8d4cd558955a5c7638fba2d68f464/YM/Spin/Floor.lean\#L39-L74}{\texttt{Spin.reach}},
\href{https://github.com/yang-mills-group/OpenAI-YM-20261028/blob/3456e5e75cf8d4cd558955a5c7638fba2d68f464/YM/Fields/Star.lean\#L26-L35}{\texttt{Boundary.Spins.Descends}}.

\textbf{Proposition VI.1.9 (the assembled record).} Let the free data,
the infrared triple \((J,p,Y)\) and the ultraviolet record be as in §2
(inputs 3--4). Then the six assembly clauses hold: compatibility of the
raw maps, their normalisation, descent of rotations, covariance,
normalisation and averaging of the stress, and compatibility of the
Lorentz action with inclusions.

Hence, if \(\partial\) satisfies the inclusion and derivative laws,
there exist \(c\), \((V_d)\) and \(J'\) such that the following local
field laws hold:

\begin{itemize}
\tightlist
\item
  the normalisation of VI.1.2 for \(c\);
\item
  \(V_d\) intertwines \(M\) and \(\Gamma\);
\item
  \(R=V R^{\rm free}V^{-1}\) with generators \(c^{-1}Y\) is a port;
\item
  the three plan laws; the rotation laws and orientation of \(R\); each
  \(F_d\) is isomorphic to \(F^{\rm free}_d\) compatibly with \(M\) and
  \(\Gamma\); and \(J'\) is conserved, tensorial and satisfies the
  stress laws with \((\partial,c^{-1}Y)\).
\end{itemize}

There is also a Lorentz action descending from \(R\) (VI.1.8(c)).

\emph{Status:} theorem conditional on the inputs in §2. Their production
uses VI.2 and the ultraviolet units; VIII.1 assembles the resulting
local field laws with the phase and reconstruction data.

\emph{Lean:}
\href{https://github.com/yang-mills-group/OpenAI-YM-20261028/blob/3456e5e75cf8d4cd558955a5c7638fba2d68f464/YM/Next/ROTATIONS.lean\#L126-L131}{\texttt{Outer.commission}},
\href{https://github.com/yang-mills-group/OpenAI-YM-20261028/blob/3456e5e75cf8d4cd558955a5c7638fba2d68f464/YM/Bid/Town.lean\#L159-L168}{\texttt{Bid.station}},
\href{https://github.com/yang-mills-group/OpenAI-YM-20261028/blob/3456e5e75cf8d4cd558955a5c7638fba2d68f464/YM/Org/Twist.lean\#L29-L43}{\texttt{Org.Commission}},
\href{https://github.com/yang-mills-group/OpenAI-YM-20261028/blob/3456e5e75cf8d4cd558955a5c7638fba2d68f464/YM/Org/Twist.lean\#L45-L53}{\texttt{Commission.picks}},
\href{https://github.com/yang-mills-group/OpenAI-YM-20261028/blob/3456e5e75cf8d4cd558955a5c7638fba2d68f464/YM/Org/Twist.lean\#L19-L24}{\texttt{Org.Tab}},
\href{https://github.com/yang-mills-group/OpenAI-YM-20261028/blob/3456e5e75cf8d4cd558955a5c7638fba2d68f464/YM/Spec/Closed.lean\#L34-L43}{\texttt{Spec.Locals}}.

\subsubsection{2. What it uses}\label{what-it-uses-23}

\textbf{Inputs from other units.}

\begin{enumerate}
\def\labelenumi{\arabic{enumi}.}
\tightlist
\item
  \textbf{IV.1 (limit table and differential plan).}

  \begin{itemize}
  \tightlist
  \item
    The Euclidean axioms
    \href{https://github.com/yang-mills-group/OpenAI-YM-20261028/blob/3456e5e75cf8d4cd558955a5c7638fba2d68f464/YM/Euclid/Words.lean\#L141-L148}{\texttt{Table.Hierarchy}}.
  \item
    The plan laws
    \href{https://github.com/yang-mills-group/OpenAI-YM-20261028/blob/3456e5e75cf8d4cd558955a5c7638fba2d68f464/YM/Euclid/Differential.lean\#L61-L64}{\texttt{DiffPlan.Derivatives}},
    \href{https://github.com/yang-mills-group/OpenAI-YM-20261028/blob/3456e5e75cf8d4cd558955a5c7638fba2d68f464/YM/Euclid/Differential.lean\#L58-L60}{\texttt{Inclusions}}
    and
    \href{https://github.com/yang-mills-group/OpenAI-YM-20261028/blob/3456e5e75cf8d4cd558955a5c7638fba2d68f464/YM/Euclid/Differential.lean\#L66-L72}{\texttt{Orientation}}.
  \item
    Also: \(\partial_a\partial_b=\partial_b\partial_a\); injectivity of
    \(\iota\); the unique lift of compatible maps to the union of label
    spaces.
  \end{itemize}
\item
  \textbf{V.1 (OS chambers).}

  \begin{itemize}
  \tightlist
  \item
    For each coordinate axis there is the Osterwalder--Schrader Hilbert
    space of reflection in \(\{x_a=0\}\)
    (\href{https://github.com/yang-mills-group/OpenAI-YM-20261028/blob/3456e5e75cf8d4cd558955a5c7638fba2d68f464/YM/OS/Dynamics.lean\#L57-L61}{\texttt{OS.Chamber}}).
    It carries a self-adjoint contraction semigroup \(e^{-tH}\), the
    unitary translations parallel to the hyperplane, and their joint
    spectral measure.
  \item
    It is built from the axioms by
    \href{https://github.com/yang-mills-group/OpenAI-YM-20261028/blob/3456e5e75cf8d4cd558955a5c7638fba2d68f464/YM/OS/Dynamics.lean\#L106-L107}{\texttt{OS.realization}},
    so it is not an assumption.
  \end{itemize}
\item
  \textbf{VI.2 (free-Maxwell data)}
  \href{https://github.com/yang-mills-group/OpenAI-YM-20261028/blob/3456e5e75cf8d4cd558955a5c7638fba2d68f464/YM/Org/Input.lean\#L159-L170}{\texttt{Org.MaxwellData}}:

  \begin{itemize}
  \tightlist
  \item
    a free table with its axioms;
  \item
    free comparison data
    (\href{https://github.com/yang-mills-group/OpenAI-YM-20261028/blob/3456e5e75cf8d4cd558955a5c7638fba2d68f464/YM/Climb/Beams.lean\#L24-L29}{\texttt{Climb.Lighthouse}}),
    i.e.~\((\partial^{\rm free},R^{\rm free},J^{\rm free},\kappa\ne0)\),
    with laws
    \href{https://github.com/yang-mills-group/OpenAI-YM-20261028/blob/3456e5e75cf8d4cd558955a5c7638fba2d68f464/YM/Climb/Beams.lean\#L34-L42}{\texttt{Goes}}:

    \begin{itemize}
    \tightlist
    \item
      plan laws;
    \item
      rotation laws and orientation;
    \item
      stress laws for
      \((J^{\rm free};\kappa\partial^{\rm free},\kappa\mathcal L^{\rm free})\);
    \item
      \(\partial^{\rm free}\) is nonzero on \(F^{\rm free}_4\);
    \item
      at every sufficiently large degree, finitely many collinear
      readouts whose common kernel within the free vector-map space is
      \(\mathbb R\partial^{\rm free}\);
    \end{itemize}
  \item
    a free Lorentz action \(B^{\rm free}\) descending from
    \(R^{\rm free}\);
  \item
    the perfect commutator relations among the free generators;
  \item
    Casimir projection polynomials;
  \item
    free row erasure.
  \end{itemize}

  The free-data producer
  \href{https://github.com/yang-mills-group/OpenAI-YM-20261028/blob/3456e5e75cf8d4cd558955a5c7638fba2d68f464/YM/Next/MAXWELL/Lit.lean\#L45-L48}{\texttt{Maxwell.Free.maxwell\_supplies}}
  and its stress replacement
  \href{https://github.com/yang-mills-group/OpenAI-YM-20261028/blob/3456e5e75cf8d4cd558955a5c7638fba2d68f464/YM/UVNext/FREE_STRESS_CERTIFICATE/Choose.lean\#L137-L150}{\texttt{FreeStress.fire}}
  are supplied in VI.2. A12 and A10+A09 give the Casimir polynomials and
  the generator relations. These free-data producers assume \(G\)
  simple.
\item
  \textbf{VI.2 and III.11 (raw stress, shell maps and weak comparison).}

  \begin{itemize}
  \tightlist
  \item
    \href{https://github.com/yang-mills-group/OpenAI-YM-20261028/blob/3456e5e75cf8d4cd558955a5c7638fba2d68f464/YM/Org/Input.lean\#L172-L175}{\texttt{Org.IR}}
    is the triple \((J,p,Y)\).
  \item
    \href{https://github.com/yang-mills-group/OpenAI-YM-20261028/blob/3456e5e75cf8d4cd558955a5c7638fba2d68f464/YM/Org/Input.lean\#L177-L186}{\texttt{Org.UV}}
    contains the lattice and extraction data, corrected-field bounds and
    Euclidean axioms for \(S\); orientation of \(\partial\); (R); the
    \emph{weak action law}; and grid approximations.
  \item
    The weak action law
    \href{https://github.com/yang-mills-group/OpenAI-YM-20261028/blob/3456e5e75cf8d4cd558955a5c7638fba2d68f464/YM/Climb/Beams.lean\#L50-L51}{\texttt{Climb.ActLaw}}
    says: for every arity bound \(N\), degree bound \(d\), derivative
    bound \(m\) and \(R\ge1\), there is a \emph{scaling system}
    \(X=(r_k\downarrow0,\ \xi_k>0,\ U_{k,e}:F^{\rm free}_e\to F_e)\)
    (\href{https://github.com/yang-mills-group/OpenAI-YM-20261028/blob/3456e5e75cf8d4cd558955a5c7638fba2d68f464/YM/Climb/Cells.lean\#L24-L31}{\texttt{Climb.Scaling}})
    with a \emph{track}
    (\href{https://github.com/yang-mills-group/OpenAI-YM-20261028/blob/3456e5e75cf8d4cd558955a5c7638fba2d68f464/YM/Climb/Beams.lean\#L44-L48}{\texttt{Climb.Track}}):

    \begin{itemize}
    \tightlist
    \item
      every \(U_{k,e}\) is a complex linear isomorphism; for
      \(e\le\max(d,4)\) it intertwines \(M\), \(\Gamma\) and the unit,
      and these maps respect inclusions within that degree range;
    \item
      for \(n\le N\), labels of degree at most \(d\), and \(j\le m\),
      \(D_y^jS(U_kv;r_ky)\to D^jS^{\rm free}(v;y)\) uniformly on the
      compact set \(\|y\|\le R\) with pairwise maximum-norm separation
      at least \(R^{-1}\);
    \item
      \(\xi_k^{-1}r_k^4\,U_{k,4}^{-1}J_{ab}\to J^{\rm free}_{ab}\) for
      every \(a,b\).
    \end{itemize}
  \item
    The scaling system may depend on all four bounds. Any finite
    collection of requests can be combined by increasing them.
  \item
    The conserved-current construction produces the raw stress; the
    annular-field construction supplies its shell maps. The immediate
    suppliers are
    \href{https://github.com/yang-mills-group/OpenAI-YM-20261028/blob/3456e5e75cf8d4cd558955a5c7638fba2d68f464/YM/MergeProof/Emanate.lean\#L15-L30}{\texttt{Acme.readyA}}
    and
    \href{https://github.com/yang-mills-group/OpenAI-YM-20261028/blob/3456e5e75cf8d4cd558955a5c7638fba2d68f464/YM/MergeProof/West/Orange.lean\#L33-L63}{\texttt{Acme.shells\_law}}.
    They provide the weak comparison and shell laws for this same
    stress.
  \end{itemize}
\item
  \textbf{VIII.1 (assembly).} Combines the local field laws obtained
  here with the phase decomposition and reconstruction.
\end{enumerate}

\textbf{Standard mathematics.}

\begin{itemize}
\tightlist
\item
  Ward identities for a conserved symmetric stress tensor: the flux of
  \(J_{ab}V_b\) through a closed shell, for a Killing field \(V\)
  (Noether; e.g.~Itzykson--Zuber §2-4).
\item
  Osterwalder--Schrader reconstruction and the transfer operator along a
  coordinate axis (Osterwalder--Schrader 1973/1975; Glimm--Jaffe,
  \emph{Quantum Physics}, ch.~6 and 19). Spectral theorem and Laplace
  representations.
\item
  Analytic vectors (Nelson, Ann. of Math. 1959). Holomorphic contraction
  semigroups (Reed--Simon II). Formal uniqueness at a noncharacteristic
  line (as proved in VI.1.m). Identity theorem for real-analytic
  functions.
\item
  Representation theory:

  \begin{itemize}
  \tightlist
  \item
    Haar measure and Weyl's unitarian trick;
  \item
    the Casimir element;
  \item
    McWeeny purification \(x\mapsto3x^2-2x^3\) of approximate
    idempotents (Rev.~Mod. Phys. 1960).
  \end{itemize}
\item
  Groups and ODEs:

  \begin{itemize}
  \tightlist
  \item
    Cartan--Dieudonné (\(O(n)\) is generated by reflections; Mathlib);
  \item
    Givens rotations;
  \item
    generation of \(SO^+(1,3)\) by rotations and boosts;
  \item
    Picard--Lindelöf uniqueness.
  \end{itemize}
\end{itemize}

\textbf{Imported constants and assumptions.}

\begin{itemize}
\tightlist
\item
  \(\kappa\neq0\) comes from the chosen free data. The initial free-data
  constructor has \(\kappa=-\tau\); the stress replacement used in
  assembly supplies a common nonzero Ward constant without identifying
  its numerical value.
\item
  \(c\) is determined by \(p=c\,\partial\). Along scaling systems
  satisfying the degree-\(4\) requests in VI.1.g, \(c/\xi_k\to\kappa\)
  (Lemma VI.1.u).
\item
  The flux estimate uses \(\lambda_k=e^k\) and the exponents
  \(100,150,300,3000\). These are internal choices.
\item
  Inputs 3--4 are hypotheses of the assembly theorem; their existence is
  supplied by the units cited above.
\end{itemize}

\subsubsection{3. Strategy}\label{strategy-23}

Conservation and the raw shell laws make the maps \(p\) and \(Y\)
compatible with inclusions, conjugation and signed permutations. Moving
the center inside a fixed window gives their commutators with
derivatives. To obtain a global rotation identity on a coordinate line,
test the stress against a compact vector field that equals the rotation
field near the insertions. Its cutoff has widths \(\lambda^{100}\) in
the rotation plane, \(\lambda^{150}\) in the line direction if that
direction is transverse to the plane, and \(\lambda^{300}\) in the
remaining directions. Two derivative terms cancel by symmetry of the
stress. Reflection positivity and spectral estimates make the remaining
exterior flux \(O(\lambda^{-1})\), so the collinear angular sum
vanishes. Translation along the line then gives the transverse
translation sums.

The free comparison identifies \(p=c\partial\): finitely many collinear
tests have a one-dimensional kernel in the free vector-map space, and
this rank condition persists along a sufficiently accurate scaling
system. Dividing by \(c\) gives the collinear Ward identity for
\(D=c^{-1}Y\). Its derivatives determine the transverse Taylor series of
each two-point function from an explicit analytic comparison function
built from radial correlations. Positivity turns the resulting bounds on
Gram-kernel jets into factorial bounds for derivatives of the
corresponding OS vectors. Correlations are therefore analytic when one
insertion moves while remaining extreme in some coordinate. The Ward
identity propagates along finite chains of these moves; density gives it
at every separated configuration.

Differentiating the global Ward identity twice gives the
\(\mathfrak{so}(4)\) brackets, after row erasure is transported from the
free theory. A relative-Casimir polynomial and McWeeny iteration produce
exact intertwiners with the free rotation generators. Transporting the
free action through them gives \(O(4)\) actions on each degree. The Ward
identity implies covariance, which in turn gives compatibility with
inclusions, the unit and derivatives. Haar averaging of \(c^{-1}J\)
gives the tensorial stress \(J'\) without changing its Ward laws. The
free Lorentz action is transported through the same intertwiners;
uniqueness for the boost ODEs makes it compatible across degrees.

\subsubsection{4. Steps}\label{steps-23}

\paragraph{4.1 Local structure of the raw shell maps
(A01--A03)}\label{local-structure-of-the-raw-shell-maps-a01a03}

\textbf{Lemma VI.1.a (A01: compatibility).} Assume (R), the derivative
law and the inclusion law. Then:

\begin{itemize}
\tightlist
\item
  \(p\) lies in the \emph{colour space}
  (\href{https://github.com/yang-mills-group/OpenAI-YM-20261028/blob/3456e5e75cf8d4cd558955a5c7638fba2d68f464/YM/Climb/Cells.lean\#L89-L91}{\texttt{Climb.ColorLaw}}):
  \(p_i\mathbf 1=0\); \(p\) commutes with \(\iota\) and \(\Gamma\);
  \(M_\sigma p_i=\sum_j\sigma_{ji}\,p_jM_\sigma\).
\item
  \(Y\) is compatible
  (\href{https://github.com/yang-mills-group/OpenAI-YM-20261028/blob/3456e5e75cf8d4cd558955a5c7638fba2d68f464/YM/Bid/Angular.lean\#L16-L20}{\texttt{Bid.Hyper}}).
\end{itemize}

\emph{Proof sketch.} Each identity is tested against all words in a
fixed window. The shells are fluxes of \(J\) against
\(\partial\chi\otimes V\), with \(V\equiv e_i\) for \(p\) and
\(V(y)=K(y-x_0)\) for \(Y\).

\begin{itemize}
\tightlist
\item
  \emph{Unit.} If the label at the centre is \(\mathbf 1\), the unit
  axiom deletes that insertion. Integrating by parts against \(\chi\)
  gives \(\int\chi\,\mathrm{div}\big(\sum_bV_bS(J_{ab}(y),\dots)\big)\),
  and this vanishes. Conservation handles one term. The Killing property
  (\(\partial_aV_b\) antisymmetric) together with \(J_{ab}=J_{ba}\)
  handles the other.
\item
  \emph{Mirrors.} Substitute \(y\mapsto\sigma y\) and use the window
  \(\chi\circ\sigma^{-1}\) at \(\sigma x\). Then use \(B_4\)-covariance
  of \(S\), raw tensoriality of \(J\), and
  \(\sigma K(\sigma^{-1}\cdot)=(\sigma K\sigma^{-1})\cdot\).
\item
  \emph{Reality.} Reality of \(S\) and \(J\), and the weights are real.
\item
  \emph{Inclusions.} The inclusion law inside the integrand.
\item
  The orientation hypothesis is not used.
\end{itemize}

\emph{Lean:}
\href{https://github.com/yang-mills-group/OpenAI-YM-20261028/blob/3456e5e75cf8d4cd558955a5c7638fba2d68f464/YM/Work/A01.lean\#L115-L117}{\texttt{Work.A01.answer}},
\href{https://github.com/yang-mills-group/OpenAI-YM-20261028/blob/3456e5e75cf8d4cd558955a5c7638fba2d68f464/YM/Work/A01.lean\#L27-L76}{\texttt{translation\_color}},
\href{https://github.com/yang-mills-group/OpenAI-YM-20261028/blob/3456e5e75cf8d4cd558955a5c7638fba2d68f464/YM/Work/A01.lean\#L78-L112}{\texttt{angular\_hyper}},
\href{https://github.com/yang-mills-group/OpenAI-YM-20261028/blob/3456e5e75cf8d4cd558955a5c7638fba2d68f464/YM/Work/A01/Conserve.lean\#L154-L181}{\texttt{cutoff\_zero}},
\href{https://github.com/yang-mills-group/OpenAI-YM-20261028/blob/3456e5e75cf8d4cd558955a5c7638fba2d68f464/YM/Work/A01/Conserve.lean\#L183-L203}{\texttt{flux\_unit}},
\href{https://github.com/yang-mills-group/OpenAI-YM-20261028/blob/3456e5e75cf8d4cd558955a5c7638fba2d68f464/YM/Work/A01/Mirror.lean\#L161-L194}{\texttt{flux\_mirror}}.

\textbf{Lemma VI.1.b (A02: moving the centre).} Assume (R) and the
derivative law. Then

\[\partial_bY_Ku+\sum_cK_{cb}\,p_cu=Y_K\,\partial_bu\qquad\text{and}\qquad\partial_b\,p_i=p_i\,\partial_b\]

on every \(F_d\)
(\href{https://github.com/yang-mills-group/OpenAI-YM-20261028/blob/3456e5e75cf8d4cd558955a5c7638fba2d68f464/YM/Bid/Angular.lean\#L27-L29}{\texttt{Bid.ShellCom}}).

\emph{Proof sketch.} Keep \(\chi\) fixed and move the centre to
\(x_0+te_b\); for small \(t\), \(\chi\) is still a window. Differentiate
both shell laws at \(t=0\).

\begin{itemize}
\tightlist
\item
  The right-hand sides give \(S(\partial_bY_KB)\) and
  \(S(\partial_bp_iB)\), by the derivative law.
\item
  On the left, the kernel's \(t\)-derivative replaces \(B\) by
  \(\partial_bB\). The weight \(K(y-x_0-te_b)\) contributes
  \(-\sum_cK_{cb}\) times the translation shell.
\item
  Differentiation under the integral is dominated on
  \(\mathrm{supp}\,\nabla\chi\times[-r,r]\), so no near-diagonal
  behaviour is needed.
\end{itemize}

\emph{Lean:}
\href{https://github.com/yang-mills-group/OpenAI-YM-20261028/blob/3456e5e75cf8d4cd558955a5c7638fba2d68f464/YM/Work/A02.lean\#L272-L307}{\texttt{Work.A02.answer}},
\href{https://github.com/yang-mills-group/OpenAI-YM-20261028/blob/3456e5e75cf8d4cd558955a5c7638fba2d68f464/YM/Work/A02.lean\#L213-L259}{\texttt{deriv\_shellSpin}},
\href{https://github.com/yang-mills-group/OpenAI-YM-20261028/blob/3456e5e75cf8d4cd558955a5c7638fba2d68f464/YM/Work/A02.lean\#L37-L48}{\texttt{stable}}.

\textbf{Lemma VI.1.c (A03: transversality).} Assume the plan laws, the
colour law for \(p\), the second identity of VI.1.b, and the line sums
of VI.1.1(ii). Then \(p\) and \(\partial\) are \emph{transverse}
(\href{https://github.com/yang-mills-group/OpenAI-YM-20261028/blob/3456e5e75cf8d4cd558955a5c7638fba2d68f464/YM/Climb/Peaks.lean\#L57}{\texttt{Climb.transverse}}).
That is: for \(x\in\mathrm{Conf}_n\) on \(\mathbb Re_a\), \(b\neq a\)
and \(v\in F_d^n\), the function
\(y\mapsto\sum_kS(v_1,\dots,P_bv_k,\dots,v_n;y)\) has vanishing jet of
every order at \(x\), for \(P=p\) and for \(P=\partial\).

\emph{Proof sketch.}

\begin{itemize}
\tightlist
\item
  The inclusion law extends the line sums to labels of mixed degree.
\item
  By \(\partial p=p\,\partial\) and the derivative law, a derivative of
  such a row sum in a direction \(e_{k,c}\) is the same kind of row sum
  with \(\partial_c\) inserted at slot \(k\). Induction on the order,
  quantified over all words, reduces every jet to a line sum at \(x\).
\item
  For \(\partial\), the line sum is translation invariance, and the
  colour law is
  \href{https://github.com/yang-mills-group/OpenAI-YM-20261028/blob/3456e5e75cf8d4cd558955a5c7638fba2d68f464/YM/Climb/Cells.lean\#L93-L96}{\texttt{Climb.native}},
  which uses orientation.
\end{itemize}

\emph{Lean:}
\href{https://github.com/yang-mills-group/OpenAI-YM-20261028/blob/3456e5e75cf8d4cd558955a5c7638fba2d68f464/YM/Work/A03.lean\#L197-L202}{\texttt{Work.A03.answer}},
\href{https://github.com/yang-mills-group/OpenAI-YM-20261028/blob/3456e5e75cf8d4cd558955a5c7638fba2d68f464/YM/Work/A03.lean\#L51-L80}{\texttt{all\_line}},
\href{https://github.com/yang-mills-group/OpenAI-YM-20261028/blob/3456e5e75cf8d4cd558955a5c7638fba2d68f464/YM/Work/A03.lean\#L83-L103}{\texttt{row\_D}},
\href{https://github.com/yang-mills-group/OpenAI-YM-20261028/blob/3456e5e75cf8d4cd558955a5c7638fba2d68f464/YM/Work/A03.lean\#L152-L186}{\texttt{full\_jet}}.

\paragraph{4.2 The collinear Ward
identity}\label{the-collinear-ward-identity}

\textbf{Lemma VI.1.d (flux identity).} Assume (R) and the derivative
law. Let \(x\in\mathrm{Conf}_n\) and \(K\in\mathfrak{so}(4)\). Let \(F\)
be a smooth compactly supported vector field that equals \(y\mapsto Ky\)
near every \(x_i\). Then

\[\mathcal W_K[Y,p](https://github.com/scottnarmstrong/yang-mills-informalization/blob/main/outline/units/q;x)=-\int\sum_{a,b}\partial_aF_b(y)\,S\big(J_{ab}(y),q;x\big)\,dy .\]

\emph{Proof sketch.}

\begin{itemize}
\tightlist
\item
  Choose disjoint windows \(\chi_i\) around the \(x_i\), and split
  \(F=(F-u)+u\) with \(u=\sum_i\chi_i\,K(\cdot)\).
\item
  The flux of \(F-u\) vanishes: it is a sum of constant-field fluxes
  against a cutoff supported away from all points, which are zero by
  conservation (as in VI.1.a).
\item
  For \(u\), write
  \(\partial_a(\chi_iKy)_b=\partial_a\chi_i\,(Ky)_b+\chi_iK_{ba}\). The
  second term pairs to zero with the symmetric \(J\).
\item
  In the first term, write \(Ky=K(y-x_i)+Kx_i\). These produce, at
  \(x_i\), the rotation shell and \((Kx_i)_b\) times the translation
  shell. Use permutation symmetry to place \(x_i\) in the first slot.
\item
  The shell laws turn these into the two terms of \(\mathcal W\).
\end{itemize}

\emph{Lean:}
\href{https://github.com/yang-mills-group/OpenAI-YM-20261028/blob/3456e5e75cf8d4cd558955a5c7638fba2d68f464/YM/Next/ROTATIONS/River.lean\#L84-L113}{\texttt{Outer.flux\_value}},
\href{https://github.com/yang-mills-group/OpenAI-YM-20261028/blob/3456e5e75cf8d4cd558955a5c7638fba2d68f464/YM/Next/ROTATIONS/River.lean\#L40-L64}{\texttt{sands}},
\href{https://github.com/yang-mills-group/OpenAI-YM-20261028/blob/3456e5e75cf8d4cd558955a5c7638fba2d68f464/YM/Next/ROTATIONS/River.lean\#L66-L82}{\texttt{paid\_zero}}.

\textbf{Lemma VI.1.e (exterior flux estimate).} Fix data and cutoff
field as follows.

\begin{itemize}
\tightlist
\item
  Fix axes \(a\neq b\) and a line direction \(l\). Put
  \(\lambda=\lambda_k=e^k\), with \(k\to\infty\).
\item
  Let \(\eta\) be an even smooth bump equal to \(1\) on \(|s|\le1\) and
  zero for \(|s|\ge2\), and put \(\ell(s)=s\eta(s)\). Then \(\ell\) is
  odd and equals \(s\) on \(|s|<1\). Its derivative \(\ell'\) is even,
  compactly supported and equal to \(1\) on \(|s|<1\).
\item
  Put
  \(\mathrm{petal}_{ab}(y)=\lambda^{100}\,\ell(\lambda^{-100}y_a)\prod_{j\ne a}\ell'(\lambda^{c_j}y_j)\),
  with \(c_j=-100\) for \(j\in\{a,b\}\), \(c_l=-150\) if
  \(l\notin\{a,b\}\), and \(c_j=-300\) otherwise.
\end{itemize}

\emph{Hypotheses.}

\begin{itemize}
\tightlist
\item
  \(J\) is conserved and raw-tensorial.
\item
  \(\partial\) satisfies the derivative law and orientation.
\item
  \(x\in\mathrm{Conf}_n\) lies on \(\mathbb Re_l\).
\item
  \(d\neq a\).
\end{itemize}

\emph{Conclusion.} For every fixed tuple of tags \(q\), the integrand is
absolutely integrable for all sufficiently large \(k\), and
\(\int\partial_d\,\mathrm{petal}_{ab}(y)\,S\big(J_{db}(y),q;x\big)\,dy=O(\lambda^{-1})\).
The constant may depend on the tags, configuration, axes and stress, but
not on \(k\). The integrand is supported in the exterior shell
\(\lambda^{-c_d}\le|y_d|\le2\lambda^{-c_d}\).

\emph{Proof sketch.}

\begin{itemize}
\tightlist
\item
  \emph{Cluster form.} Off the insertions,
  \(S(J_{db}(y),q;x)=\langle\bar B;J_{db}(y)\rangle\). Here \(B\) is the
  reflected word of the collinear insertions. Its insertion points have
  transverse coordinates \(0\)
  (\href{https://github.com/yang-mills-group/OpenAI-YM-20261028/blob/3456e5e75cf8d4cd558955a5c7638fba2d68f464/YM/Next/ROTATIONS/Yard.lean\#L84-L103}{\texttt{collimated}}).
\item
  \emph{Parity.} Decompose \(B\) into components of definite reflection
  parity in each axis. Only the component odd in \(a\) and \(b\)
  survives integration against the test
  (\href{https://github.com/yang-mills-group/OpenAI-YM-20261028/blob/3456e5e75cf8d4cd558955a5c7638fba2d68f464/YM/Next/ROTATIONS/Yard.lean}{\texttt{Yard.lean}}:
  \href{https://github.com/yang-mills-group/OpenAI-YM-20261028/blob/3456e5e75cf8d4cd558955a5c7638fba2d68f464/YM/Next/ROTATIONS/Yard.lean\#L49-L77}{\texttt{pick\_price}}).
\item
  \emph{Laplace representation.} Represent the correlation through the
  OS space of axis \(d\). Because the test vanishes for
  \(|y_d|<\lambda^{-c_d}\), the \(y_d\)-integral becomes a Laplace
  transform. Its time factor satisfies
  \(|\mathrm{Lap}(E)|\le C\lambda^{c_d}/(E+\lambda^{c_d})\), paired with
  the polarised spectral measure \(\nu_d(B,J_{db})\) of the transfer
  operator
  (\href{https://github.com/yang-mills-group/OpenAI-YM-20261028/tree/3456e5e75cf8d4cd558955a5c7638fba2d68f464/YM/Axis}{\texttt{YM/Axis}};
  \href{https://github.com/yang-mills-group/OpenAI-YM-20261028/blob/3456e5e75cf8d4cd558955a5c7638fba2d68f464/YM/Next/ROTATIONS/Vanishing.lean\#L68-L120}{\texttt{exterior\_size}}).
\item
  \emph{Outside the low box.} Where \(E>1\) or some
  \(|p_j|>\lambda^{1+c_j}\), the contribution is
  \(O(\lambda^{-\infty})\). This uses only Schwartz decay of the
  profiles and the zeroth moment.
\item
  \emph{Inside the box.} A weighted Cauchy--Schwarz (Young) inequality
  splits the bound between \(\nu_d(B)\) and \(\nu_d(J_{db})\).
\item
  \emph{Shear case \(d\notin\{a,b\}\).}

  \begin{itemize}
  \tightlist
  \item
    The line cost of the cluster is
    \(\int_{E\le1}(E+\lambda^{-3000})^{-1}\prod_jh_j\,d\nu_d(B)=O(\lambda^{1})\)
    for \(d=l\), and \(O(\lambda^{150})\) otherwise.
  \item
    The point cost
    \(\int_{E\le1}(E+\lambda^{c_d})^{-1}\prod_jh_j\,d\nu_d(J)=O(\lambda^{\delta})\)
    holds for every \(\delta>0\).
  \item
    Here \(h_j\) are fixed nonnegative even smooth bumps, rescaled in
    \(L^1\) at widths \(\lambda^{1+c_j}\); products run over \(j\ne d\).
  \item
    Before multiplication by the prefactor \(\lambda^{100}\) in the
    cutoff, the exponent is at most \(-146\).
  \end{itemize}
\item
  \emph{Diagonal case \(d=b\).}

  \begin{itemize}
  \tightlist
  \item
    \(J_{bb}\) is even in every axis, so the mixed-parity point estimate
    does not apply. Split at \(E=\lambda^{-3000}\).
  \item
    Above the split, the cluster's line cost is paired with the
    \emph{mass estimate}
    \(\int_{E\le1}EW^4(E+W)^{-4}\prod h_j\,d\nu_b(J_{bb})=O(\lambda^{\delta}W^4)\)
    for every \(\delta>0\), at \(W=\lambda^{-99}\).
  \item
    Below the split, the packet estimate gains \(\lambda^{-1500}\); the
    line estimate then suffices with a polynomial bound on the stress
    measure.
  \end{itemize}
\item
  \emph{Point cost.} Proved by a \emph{sign alternation}. The same
  Euclidean integral is a positive spectral pairing along one axis and
  minus a positive pairing along an axis of opposite parity.
  Moment-subtracted profiles have Laplace responses bounded below by
  \(c/(E+w)\) down to \(E=0\). Walking through the axes, with orders
  chosen backwards, moves every profile to scale \(\ge1\) at subpower
  cost. These estimates also control the contribution at zero energy and
  require no spectral gap.
\item
  \emph{Moment bounds.} For an axis-pointlike weight,
  \(\int e^{-tE}|p_j|^{2m}\,d\nu\le Ct^{-D}\) for \(0<t\le1\) and
  \(j\ne d\). The polynomial kernel bounds give this estimate by finite
  differences. Shifted cluster weights have analogous damped moment
  estimates.
\item
  \emph{Mass estimate.} Uses conservation in tested form. Tested
  conservation relates the contact-normalised two-point symbols by
  \(B=p_i^2A+D_B\) and \(C=p_i^2B+D_C\). Contact terms vanish on tests
  with enough zero moments; the remaining error estimates enter after
  normalisation by the band widths. Positivity gives alternating signs
  for the three residual band differences. Comparing the bands \(W\) and
  \(2W\) bounds the positive residual, which dominates the displayed
  mass integral.
\item
  The mass estimate uses conservation, the derivative law and
  orientation. Raw tensoriality supplies the parities. Combining the two
  spectral regions gives \(O(\lambda^{-101})\) before the prefactor
  \(\lambda^{100}\), and hence the stated \(O(\lambda^{-1})\) bound.
\end{itemize}

\emph{Lean:}
\href{https://github.com/yang-mills-group/OpenAI-YM-20261028/blob/3456e5e75cf8d4cd558955a5c7638fba2d68f464/YM/Next/ROTATIONS/Yard.lean\#L105-L114}{\texttt{Outer.blows}},
\href{https://github.com/yang-mills-group/OpenAI-YM-20261028/blob/3456e5e75cf8d4cd558955a5c7638fba2d68f464/YM/Next/ROTATIONS/Moments.lean\#L228-L243}{\texttt{point\_moment}},
\href{https://github.com/yang-mills-group/OpenAI-YM-20261028/blob/3456e5e75cf8d4cd558955a5c7638fba2d68f464/YM/Next/ROTATIONS/Switch.lean\#L145-L160}{\texttt{alternate}},
\href{https://github.com/yang-mills-group/OpenAI-YM-20261028/blob/3456e5e75cf8d4cd558955a5c7638fba2d68f464/YM/Next/ROTATIONS/Pivot.lean\#L59-L72}{\texttt{budget\_lock}};
the cutoff profile is defined in
\href{https://github.com/yang-mills-group/OpenAI-YM-20261028/blob/3456e5e75cf8d4cd558955a5c7638fba2d68f464/YM/Next/ROTATIONS/Patch.lean}{\texttt{Patch.lean}};
the line and point estimates are in
\href{https://github.com/yang-mills-group/OpenAI-YM-20261028/blob/3456e5e75cf8d4cd558955a5c7638fba2d68f464/YM/Next/ROTATIONS/Tunnel.lean}{\texttt{Tunnel.lean}}
and the mass estimate is in
\href{https://github.com/yang-mills-group/OpenAI-YM-20261028/blob/3456e5e75cf8d4cd558955a5c7638fba2d68f464/YM/Next/ROTATIONS/Mass.lean}{\texttt{Mass.lean}}.
\emph{(proof read for the reduction and combination of the spectral
estimates; supporting analytic inputs remain as listed in §7)}

\textbf{Lemma VI.1.f (collinear identity and line sums).} This is
Proposition VI.1.1.

\emph{Proof sketch.}

\begin{itemize}
\tightlist
\item
  Take \(F=\mathrm{vane}_{ab}\), with \(F_b=\mathrm{petal}_{ab}\),
  \(F_a=-\mathrm{petal}_{ba}\) and \(F_d=0\) otherwise. For sufficiently
  large cutoff index it equals the Killing field of \(K^{(ba)}\) near
  every \(x_i\).
\item
  In the flux of VI.1.d, the diagonal terms
  \(\partial_a\mathrm{petal}_{ab}\,J_{ab}\) and
  \(\partial_b\mathrm{petal}_{ba}\,J_{ba}\) cancel identically, since
  \(\partial_a\mathrm{petal}_{ab}=\partial_b\mathrm{petal}_{ba}\) and
  \(J\) is symmetric. What remains is a finite sum of integrals of
  VI.1.e.
\item
  Hence
  \(\mathcal W_{K^{(ba)}}[Y,p](https://github.com/scottnarmstrong/yang-mills-informalization/blob/main/outline/units/q;x)\),
  which does not depend on \(n\), is \(O(e^{-n})\), so it is \(0\).
\item
  \(\mathcal W_K[Y,p]\) is \(\mathbb R\)-linear in \(K\), so it vanishes
  for all \(K\).
\item
  Translating the configuration by \(e_l\) along its line changes
  \(\mathcal W_{K^{(jl)}}[Y,p]\) by exactly
  \(\sum_iS(\dots,p_jq_i,\dots)\). Both values vanish, which gives the
  line sums.
\end{itemize}

\emph{Lean:}
\href{https://github.com/yang-mills-group/OpenAI-YM-20261028/blob/3456e5e75cf8d4cd558955a5c7638fba2d68f464/YM/Next/ROTATIONS/Route.lean\#L31-L48}{\texttt{draught\_river}},
\href{https://github.com/yang-mills-group/OpenAI-YM-20261028/blob/3456e5e75cf8d4cd558955a5c7638fba2d68f464/YM/Next/ROTATIONS/Route.lean\#L49-L56}{\texttt{drained}},
\href{https://github.com/yang-mills-group/OpenAI-YM-20261028/blob/3456e5e75cf8d4cd558955a5c7638fba2d68f464/YM/Next/ROTATIONS/Route.lean\#L57-L67}{\texttt{axisTurn}},
\href{https://github.com/yang-mills-group/OpenAI-YM-20261028/blob/3456e5e75cf8d4cd558955a5c7638fba2d68f464/YM/Next/ROTATIONS/Route.lean\#L83-L96}{\texttt{fullTurn}},
\href{https://github.com/yang-mills-group/OpenAI-YM-20261028/blob/3456e5e75cf8d4cd558955a5c7638fba2d68f464/YM/Next/ROTATIONS/Route.lean\#L98-L108}{\texttt{stagger}},
\href{https://github.com/yang-mills-group/OpenAI-YM-20261028/blob/3456e5e75cf8d4cd558955a5c7638fba2d68f464/YM/Next/ROTATIONS/Route.lean\#L21-L27}{\texttt{vane\_home}}.

\paragraph{\texorpdfstring{4.3 Normalisation \(p=c\,\partial\)
(A04--A06)}{4.3 Normalisation p=c\textbackslash,\textbackslash partial (A04--A06)}}\label{normalisation-pcpartial-a04a06}

\textbf{Lemma VI.1.g (scaling limits; A04, A05).} Assume the inclusion
and derivative laws, (R), and the free data of §2. For each degree
\(d\), there are bounds \(d_0,N_0,m_0,R_0\), with \(R_0\ge1\), such that
every track whose arity, degree, derivative and size bounds are at least
\(N_0,d_0,m_0,R_0\), respectively has the following limits. Write \(X\)
for its scaling system and define
\(\mathrm{aisle}_{k,d}(P)_a=U_{k,d+1}^{-1}P_aU_{k,d}\). At this fixed
degree:

\begin{itemize}
\tightlist
\item
  \(r_k\,\mathrm{aisle}_{k,d}(\partial)\to\partial^{\rm free}\);
\item
  \((r_k/\xi_k)\,\mathrm{aisle}_{k,d}(p)\to\kappa\,\partial^{\rm free}\);
\item
  \(\xi_k^{-1}U_{k,d}^{-1}Y_KU_{k,d}\to\kappa\,\mathcal L^{\rm free}_K\).
\end{itemize}

\emph{Proof sketch.}

\begin{itemize}
\tightlist
\item
  \emph{First limit.} Test against a finite family of free correlations
  that is faithful on \(F^{\rm free}_{d+1}\). The chain rule gives the
  factor \(r_k\). Then use first-order jet convergence, and invert by
  finite-dimensional linear algebra.
\item
  \emph{Second and third limits.} Write the shell law at \(r_kx\) with
  the window \(\chi(\cdot/r_k)\). Substituting \(y=r_kz\) produces
  \(r_k^4\cdot r_k^{-1}\). For the translation shell, multiplication by
  \(r_k/\xi_k\) gives the stress factor \(\xi_k^{-1}r_k^4J\). For the
  rotation shell, the weight \(K(y-r_kx_0)=r_kK(z-x_0)\) supplies the
  additional factor, so the multiplier is \(\xi_k^{-1}\). The stress
  convergence and uniform kernel convergence on the compact shell
  supports give the free shell laws. Finite faithful tests recover the
  two operator limits.
\end{itemize}

\emph{Lean:}
\href{https://github.com/yang-mills-group/OpenAI-YM-20261028/blob/3456e5e75cf8d4cd558955a5c7638fba2d68f464/YM/Work/A04.lean\#L163-L183}{\texttt{Work.A04.answer}},
\href{https://github.com/yang-mills-group/OpenAI-YM-20261028/blob/3456e5e75cf8d4cd558955a5c7638fba2d68f464/YM/Work/A05.lean\#L95-L121}{\texttt{Work.A05.answer}},
\href{https://github.com/yang-mills-group/OpenAI-YM-20261028/blob/3456e5e75cf8d4cd558955a5c7638fba2d68f464/YM/Work/A05/Window.lean\#L158-L172}{\texttt{integrate\_sheet\_scale}},
\href{https://github.com/yang-mills-group/OpenAI-YM-20261028/blob/3456e5e75cf8d4cd558955a5c7638fba2d68f464/YM/Work/A05/Read.lean\#L70-L81}{\texttt{recover\_on}};
the requests are
\href{https://github.com/yang-mills-group/OpenAI-YM-20261028/blob/3456e5e75cf8d4cd558955a5c7638fba2d68f464/YM/Bid/Angular.lean\#L43-L46}{\texttt{Bid.Watches}}
and
\href{https://github.com/yang-mills-group/OpenAI-YM-20261028/blob/3456e5e75cf8d4cd558955a5c7638fba2d68f464/YM/Bid/Angular.lean\#L48-L49}{\texttt{Bid.ConLap}}.

\textbf{Lemma VI.1.h (A06 and the identification \(p=c\,\partial\)).}
Assume the hypotheses of VI.1.2 and the transversality conclusions of
VI.1.c.

\emph{Setting.}

\begin{itemize}
\tightlist
\item
  The free data supply \(d_*\) such that, for every \(d\ge d_*\), there
  is a finite \emph{selection} of collinear transverse readouts on
  \(F^{\rm free}\): \(m\)-th jets of row sums at configurations on a
  coordinate axis.
\item
  These readouts vanish, within the colour space of the free table,
  exactly on \(\mathbb R\,\partial^{\rm free}\)
  (\href{https://github.com/yang-mills-group/OpenAI-YM-20261028/blob/3456e5e75cf8d4cd558955a5c7638fba2d68f464/YM/Climb/Tests.lean\#L101-L104}{\texttt{Climb.Ens}},
  \href{https://github.com/yang-mills-group/OpenAI-YM-20261028/blob/3456e5e75cf8d4cd558955a5c7638fba2d68f464/YM/Climb/Tests.lean\#L95-L99}{\texttt{Select}}).
\end{itemize}

\emph{Claims.}

\begin{itemize}
\tightlist
\item
  (A06) Along any scaling system meeting a sufficiently large finite
  request for these readouts, the actual rescaled readouts converge to
  the free ones. They annihilate \(\mathrm{aisle}_{k,d}(p)\) and
  \(\mathrm{aisle}_{k,d}(\partial)\) exactly at every \(k\). The reason
  is that a readout of \(\mathrm{aisle}_{k,d}(P)\) equals \(r_k^m\)
  times an actual collinear jet of \(P\), which vanishes by VI.1.c.
\item
  Consequently there is a real \(c\neq0\) with \(p_a=c\,\partial_a\) on
  all \(F_d\).
\end{itemize}

\emph{Proof sketch.}

\begin{itemize}
\tightlist
\item
  \emph{Persistence lemma.} Let \(\mathcal C\) be the fixed real space
  of free vector maps and suppose \(Z_k\to Z\) with
  \(\ker Z\cap\mathcal C=\mathbb R x\), \(x\ne0\). Choose a real linear
  functional \(\ell\) with \(\ell(x)\ne0\). The map \((Z,\ell)\) is
  injective on \(\mathcal C\); injectivity persists under small
  perturbations in finite dimension. Thus \(\ker Z_k\cap\mathcal C\) has
  dimension at most one for large \(k\). This gives
  \(\mathrm{aisle}_{k,d}(p)=c_d\,\mathrm{aisle}_{k,d}(\partial)\), hence
  \(p=c_d\partial\) on \(F_d\) for large \(d\).
\item
  The constants agree across \(d\), and lower degrees follow, by
  injectivity of \(\iota\) and the colour law.
\item
  \(p\neq0\) and \(\partial\ne0\) by VI.1.g and
  \(\kappa\partial^{\rm free}\neq0\).
\end{itemize}

\emph{Lean:}
\href{https://github.com/yang-mills-group/OpenAI-YM-20261028/blob/3456e5e75cf8d4cd558955a5c7638fba2d68f464/YM/Work/A06.lean\#L236-L256}{\texttt{Work.A06.answer}},
\href{https://github.com/yang-mills-group/OpenAI-YM-20261028/blob/3456e5e75cf8d4cd558955a5c7638fba2d68f464/YM/Climb/Peaks.lean\#L98-L111}{\texttt{Climb.extends\_form}},
\href{https://github.com/yang-mills-group/OpenAI-YM-20261028/blob/3456e5e75cf8d4cd558955a5c7638fba2d68f464/YM/Bid/Town.lean\#L49-L59}{\texttt{Bid.pressured}}.

\textbf{Lemma VI.1.i (the normalised Ward data).} With \(c\) as in
VI.1.h, \(D=c^{-1}Y\) satisfies (C) and (L). This proves Proposition
VI.1.2.

\emph{Proof sketch.} Divide the first identity of VI.1.b, and the
identity of VI.1.1(i), by \(c\), using \(p=c\,\partial\).

\emph{Lean:}
\href{https://github.com/yang-mills-group/OpenAI-YM-20261028/blob/3456e5e75cf8d4cd558955a5c7638fba2d68f464/YM/Bid/Town.lean\#L61-L78}{\texttt{Bid.formed}},
\href{https://github.com/yang-mills-group/OpenAI-YM-20261028/blob/3456e5e75cf8d4cd558955a5c7638fba2d68f464/YM/Bid/Town.lean\#L33-L39}{\texttt{Bid.witnessing}}.

\paragraph{4.4 Continuation off the
line}\label{continuation-off-the-line}

Throughout this subsection \(S\) satisfies the Euclidean axioms, \(T\)
satisfies (C) and (L), and \(\partial\) satisfies the derivative law.

\textbf{Lemma VI.1.j (derivatives of angular sums).}

\begin{itemize}
\tightlist
\item
  \(\partial_{x_k^a}\mathcal W_K[T](https://github.com/scottnarmstrong/yang-mills-informalization/blob/main/outline/units/q;x)=\mathcal W_K[T](https://github.com/scottnarmstrong/yang-mills-informalization/blob/main/outline/units/q';x)\),
  where \(q'\) is \(q\) with \(q_k\) replaced by \(\partial_aq_k\). In
  particular
  \(\mathcal W_K[T](https://github.com/scottnarmstrong/yang-mills-informalization/blob/main/outline/units/q;\cdot)\)
  is \(C^\infty\) on \(\mathrm{Conf}\).
\item
  Hence, if
  \(\mathcal W_K[T](https://github.com/scottnarmstrong/yang-mills-informalization/blob/main/outline/units/\cdot\,;y)=0\)
  for all tags at one point \(y\), then every derivative of
  \(\mathcal W_K[T](https://github.com/scottnarmstrong/yang-mills-informalization/blob/main/outline/units/q;\cdot)\)
  at \(y\) vanishes, for every \(q\).
\end{itemize}

\emph{Proof sketch.} Product rule. The derivative of \((Kx_k)_b\) in
\(x_k^a\) is \(K_{ba}\). Use
\(\partial_a\partial_b=\partial_b\partial_a\), and (C) absorbs the extra
term. The second statement is induction on the order, with \(q\)
generalised.

\emph{Lean:}
\href{https://github.com/yang-mills-group/OpenAI-YM-20261028/blob/3456e5e75cf8d4cd558955a5c7638fba2d68f464/YM/Ward/Sum.lean\#L92-L136}{\texttt{Ward.traverses}},
\href{https://github.com/yang-mills-group/OpenAI-YM-20261028/blob/3456e5e75cf8d4cd558955a5c7638fba2d68f464/YM/Ward/Sum.lean\#L138-L143}{\texttt{Ward.ac}},
\href{https://github.com/yang-mills-group/OpenAI-YM-20261028/blob/3456e5e75cf8d4cd558955a5c7638fba2d68f464/YM/Ward/Motion.lean\#L31-L42}{\texttt{Ward.jet\_line}}.

\textbf{Lemma VI.1.k (OS vectors of one insertion).} Let
\(\mathcal H_a\) be the OS space for reflection in \(\{x_a=0\}\). Write
\(\Theta_a\) for the corresponding reflection of words, including the
label adjoint \(M_{\theta_a}\Gamma\). For \(x_a>0\) let
\(\mathrm{spot}(A,x)\in\mathcal H_a\) be the vector of the one-point
word \(A(x)\).

\begin{itemize}
\tightlist
\item
  \begin{enumerate}
  \def\labelenumi{(\roman{enumi})}
  \tightlist
  \item
    For any word \(w\) supported in \(\{y_a>0\}\) with distinct
    insertion points, its OS vector \([w]\) satisfies
    \(\langle[w],\mathrm{spot}(A,x)\rangle=S(A(x),\Theta_aw)\).
  \end{enumerate}
\item
  (ii) \(\mathrm{spot}(A,x+se_a)=e^{-sH}\mathrm{spot}(A,x)\) for
  \(s\ge0\), and \(\mathrm{spot}(A,x+v)=U(v)\,\mathrm{spot}(A,x)\) for
  \(v\perp e_a\).
\item
  (iii) For \(b\neq a\), \(s\mapsto\mathrm{spot}(A,x+se_b)\) is
  \(C^\infty\) with derivative \(\mathrm{spot}(\partial_bA,\cdot)\). The
  weak derivative comes from the derivative law and is upgraded to a
  strong one.
\item
  (iv) \(t\mapsto\mathrm{spot}(A,x+te_a)\) is real-analytic for
  \(x_a+t>0\), because \(e^{-zH}\) is holomorphic on
  \(\mathrm{Re}\,z>0\).
\item
  (v) \(t\mapsto S(A(te_a),B(0))\) is real-analytic for \(t>0\).
\end{itemize}

\emph{Lean:}
\href{https://github.com/yang-mills-group/OpenAI-YM-20261028/blob/3456e5e75cf8d4cd558955a5c7638fba2d68f464/YM/Ward/Plants.lean\#L55}{\texttt{Ward.spot}},
\href{https://github.com/yang-mills-group/OpenAI-YM-20261028/blob/3456e5e75cf8d4cd558955a5c7638fba2d68f464/YM/Ward/Plants.lean\#L64-L73}{\texttt{link}},
\href{https://github.com/yang-mills-group/OpenAI-YM-20261028/blob/3456e5e75cf8d4cd558955a5c7638fba2d68f464/YM/Ward/Plants.lean\#L102-L118}{\texttt{responds}},
\href{https://github.com/yang-mills-group/OpenAI-YM-20261028/blob/3456e5e75cf8d4cd558955a5c7638fba2d68f464/YM/Ward/Plants.lean\#L120-L133}{\texttt{blooms}},
\href{https://github.com/yang-mills-group/OpenAI-YM-20261028/blob/3456e5e75cf8d4cd558955a5c7638fba2d68f464/YM/Ward/Plants.lean\#L137-L158}{\texttt{advancing}},
\href{https://github.com/yang-mills-group/OpenAI-YM-20261028/blob/3456e5e75cf8d4cd558955a5c7638fba2d68f464/YM/Ward/Ket.lean\#L19-L41}{\texttt{Analysis.represents}},
\href{https://github.com/yang-mills-group/OpenAI-YM-20261028/blob/3456e5e75cf8d4cd558955a5c7638fba2d68f464/YM/Ward/Genuine.lean\#L54-L68}{\texttt{Ward.radial}}.

\textbf{Lemma VI.1.l (transverse jets of the two-point function).} Fix
\(a\neq b\) and \(K=K^{(ba)}\). For a pair of labels
\(\pi=(\pi_1,\pi_2)\) put
\(\Phi_\pi(x_1,x_2)=S\big(\pi_1(x_1e_a+x_2e_b),\pi_2(0)\big)\) on
\(\{x_1>0\}\). Define

\[(\mathcal L\Phi)_\pi=(x_1\partial_2-x_2\partial_1)\Phi_\pi+\Phi_{(T_K\pi_1,\pi_2)}+\Phi_{(\pi_1,T_K\pi_2)} .\]

Then:

\begin{itemize}
\tightlist
\item
  \((\mathcal L\Phi)_\pi(x)=\mathcal W_K[T]\big(\pi;(x_1e_a+x_2e_b,0)\big)\);
\item
  \(\partial_2^n(\mathcal L\Phi)_\pi(t,0)=0\) for all \(n\), all \(t>0\)
  and all \(\pi\).
\end{itemize}

\emph{Proof sketch.} The first identity is the definition of
\(\mathcal W\), by the derivative law. At \(x_2=0\) both points lie on
\(\mathbb Re_a\), so (L) gives \(\mathcal W=0\) for all labels. VI.1.j
then gives all \(x_2\)-derivatives. This supplies the collinear jets
needed for the analytic comparison.

\emph{Lean:}
\href{https://github.com/yang-mills-group/OpenAI-YM-20261028/blob/3456e5e75cf8d4cd558955a5c7638fba2d68f464/YM/Ward/Rows.lean\#L73-L97}{\texttt{Ward.sigma}},
\href{https://github.com/yang-mills-group/OpenAI-YM-20261028/blob/3456e5e75cf8d4cd558955a5c7638fba2d68f464/YM/Ward/Rows.lean\#L99-L116}{\texttt{Ward.realized}}.

\textbf{Lemma VI.1.m (comparison function and formal uniqueness).} Let
\((r,\theta)\) be polar coordinates of \((x_1,x_2)\) with \(x_1>0\). Put

\[Q_\pi(x)=S\big((e^{-\theta T_K}\pi_1)(re_a),\,(e^{-\theta T_K}\pi_2)(0)\big).\]

\begin{itemize}
\item
  \(Q\) is real-analytic on \(\{x_1>0\}\), by VI.1.k(v),
  finite-dimensionality of the label spaces, and the matrix exponential.
\item
  \(\mathcal LQ=0\) identically, since \(\theta\) has derivative \(1\)
  and \(r\) has derivative \(0\) along the rotation field.
\item
  \(Q=\Phi\) on \(\{x_2=0\}\).
\item
  (Formal uniqueness.) Let \(L\) be smooth on the half-plane with
  \(L_\pi(t,0)=0\) and \(\partial_2^n(\mathcal LL)_\pi(t,0)=0\) for all
  \(n,\pi\) and \(t>0\). Then \(\partial_2^nL_\pi(t,0)=0\) for all
  \(n\).

  Applied to \(L=\Phi-Q\): the \(x_2\)-Taylor series of \(\Phi\) at
  \((t,0)\) is that of the analytic function \(Q\).
\end{itemize}

\emph{Proof sketch.} Strong induction on \(n\), for all \(\pi\) at once.
For \(n\ge1\), at \((t,0)\),

\[t\,\partial_2^{n+1}L_\pi=-\partial_2^n\big(L_{(T_K\pi_1,\pi_2)}+L_{(\pi_1,T_K\pi_2)}\big)+n\,\partial_1\partial_2^{n-1}L_\pi .\]

For \(n=0\) the transport equation directly gives
\(t\partial_2L_\pi=-L_{(T_K\pi_1,\pi_2)}-L_{(\pi_1,T_K\pi_2)}=0\). In
the displayed recurrence, the first term vanishes by induction, and the
last term vanishes because \(\partial_2^{n-1}L_\pi\) vanishes on the
whole half-line. Since \(t\neq0\), the line is noncharacteristic.

\emph{Lean:}
\href{https://github.com/yang-mills-group/OpenAI-YM-20261028/blob/3456e5e75cf8d4cd558955a5c7638fba2d68f464/YM/Ward/Circles.lean\#L89-L90}{\texttt{Ward.Q}},
\href{https://github.com/yang-mills-group/OpenAI-YM-20261028/blob/3456e5e75cf8d4cd558955a5c7638fba2d68f464/YM/Ward/Circles.lean\#L115-L128}{\texttt{htQ}},
\href{https://github.com/yang-mills-group/OpenAI-YM-20261028/blob/3456e5e75cf8d4cd558955a5c7638fba2d68f464/YM/Ward/Circles.lean\#L130-L174}{\texttt{solves}},
\href{https://github.com/yang-mills-group/OpenAI-YM-20261028/blob/3456e5e75cf8d4cd558955a5c7638fba2d68f464/YM/Ward/PDE.lean\#L121-L141}{\texttt{PDE.sure}}.

\textbf{Lemma VI.1.n (analytic Gram kernel).}

\begin{itemize}
\tightlist
\item
  (Abstract.) Let \(f:\mathbb R\to\mathcal H\) be \(C^\infty\) with
  \(\langle f(s),f(s')\rangle=A(s'-s)\), where \(A\) has at \(0\) the
  Taylor series of a real-analytic germ. Then \(f\) is real-analytic on
  \(\mathbb R\).
\item
  (Application.) Hence, for \(x_a>0\) and \(b\neq a\),
  \(s\mapsto\mathrm{spot}(A,x+se_b)\) is real-analytic on \(\mathbb R\).
\end{itemize}

\emph{Proof sketch.}

\begin{itemize}
\tightlist
\item
  \(\|f^{(m)}(s)\|^2=|A^{(2m)}(0)|\le C\rho^{-2m}(2m)!\) uniformly in
  \(s\). With \((2m)!\le4^m(m!)^2\) this gives
  \(\|f^{(m)}\|\le\sqrt C(\rho/2)^{-m}m!\) everywhere, and the Taylor
  remainder finishes.
\item
  For the application: by reflection and translation invariance,
  \(\langle f(s),f(s')\rangle=\Phi_{(A,\Theta_a A)}(2x_a,s'-s)\). VI.1.m
  identifies its Taylor series at \(0\) with that of
  \(s\mapsto Q_{(A,\Theta_a A)}(2x_a,s)\).
\item
  Uniformity in \(s\), which comes from positivity and translation
  invariance, is why a jet at a single point suffices.
\end{itemize}

\emph{Lean:}
\href{https://github.com/yang-mills-group/OpenAI-YM-20261028/blob/3456e5e75cf8d4cd558955a5c7638fba2d68f464/YM/Ward/Ket.lean\#L76-L96}{\texttt{Analysis.blossoms}},
\href{https://github.com/yang-mills-group/OpenAI-YM-20261028/blob/3456e5e75cf8d4cd558955a5c7638fba2d68f464/YM/Ward/Ket.lean\#L46-L74}{\texttt{bud}},
\href{https://github.com/yang-mills-group/OpenAI-YM-20261028/blob/3456e5e75cf8d4cd558955a5c7638fba2d68f464/YM/Ward/Rate.lean\#L93-L102}{\texttt{builds}},
\href{https://github.com/yang-mills-group/OpenAI-YM-20261028/blob/3456e5e75cf8d4cd558955a5c7638fba2d68f464/YM/Ward/Genuine.lean\#L72-L81}{\texttt{Ward.fires}},
\href{https://github.com/yang-mills-group/OpenAI-YM-20261028/blob/3456e5e75cf8d4cd558955a5c7638fba2d68f464/YM/Ward/Genuine.lean\#L33-L43}{\texttt{swing}},
\href{https://github.com/yang-mills-group/OpenAI-YM-20261028/blob/3456e5e75cf8d4cd558955a5c7638fba2d68f464/YM/Ward/Circles.lean\#L179-L209}{\texttt{Ward.pure}}.

\textbf{Lemma VI.1.o (analyticity at an extreme point).} Suppose that
after moving \(x_k\) to \(x_k+se_b\), the point \(k\) is strictly
latest, or strictly earliest, in coordinate \(a\). Then
\(u\mapsto S(q;x\text{ with }x_k+ue_b)\) is real-analytic at \(s\), for
every \(b\) (including \(b=a\)). Hence \(\mathcal W_K[T]\) is
real-analytic along such a move.

\emph{Proof sketch.}

\begin{itemize}
\tightlist
\item
  Translate in \(e_a\) so that the other points lie in \(\{y_a<0\}\) and
  \(x_k+se_b\) in \(\{y_a>0\}\).
\item
  Then
  \(S=\langle\Theta(\text{word of the others}),\mathrm{spot}(q_k,\cdot)\rangle\),
  a bounded functional of one vector. This is analytic by VI.1.n for
  \(b\neq a\), and by VI.1.k(iv) for \(b=a\).
\item
  The ``earliest'' case reduces to the ``latest'' case by the reflection
  \(\Theta_a\) with mirrored labels.
\end{itemize}

\emph{Lean:}
\href{https://github.com/yang-mills-group/OpenAI-YM-20261028/blob/3456e5e75cf8d4cd558955a5c7638fba2d68f464/YM/Ward/Extreme.lean\#L23-L69}{\texttt{Ward.higher}},
\href{https://github.com/yang-mills-group/OpenAI-YM-20261028/blob/3456e5e75cf8d4cd558955a5c7638fba2d68f464/YM/Ward/Extreme.lean\#L102-L121}{\texttt{last\_adv}},
\href{https://github.com/yang-mills-group/OpenAI-YM-20261028/blob/3456e5e75cf8d4cd558955a5c7638fba2d68f464/YM/Ward/Extreme.lean\#L131-L153}{\texttt{first\_adv}},
\href{https://github.com/yang-mills-group/OpenAI-YM-20261028/blob/3456e5e75cf8d4cd558955a5c7638fba2d68f464/YM/Ward/Walk.lean\#L15-L26}{\texttt{Ward.channel}}.

\textbf{Lemma VI.1.p (walks; proof of VI.1.3).} A \emph{walk} is a
finite chain of segments. Each segment moves one point \(x_k\) along
\(e_b\), and on the whole segment \(x_k\) is strictly earliest or latest
in some coordinate \(a\).

\begin{itemize}
\tightlist
\item
  If
  \(\mathcal W_K[T](https://github.com/scottnarmstrong/yang-mills-informalization/blob/main/outline/units/\cdot\,;y)=0\)
  (all tags) at the end \(y\) of a walk, then the same holds at its
  start.
\item
  Every configuration with pairwise distinct \(0\)-coordinates is joined
  by a walk to a configuration on \(\mathbb Re_0\).
\item
  Such configurations are dense in \(\mathrm{Conf}_n\).
\end{itemize}

\emph{Proof sketch.}

\begin{itemize}
\tightlist
\item
  \emph{Propagation.} On a segment, \(\mathcal W\) is analytic (VI.1.o),
  and all its derivatives vanish at the endpoint (VI.1.j). By the
  identity theorem it vanishes on the segment.
\item
  \emph{Geometry.} Take the earliest point in coordinate \(0\). Move it
  beyond all others in coordinate \(1\), then beyond all others in
  coordinate \(0\), then set its coordinates \(1,2,3\) to \(0\), keeping
  it latest in coordinate \(0\). Induct on the remaining points.
\item
  \emph{Density.} Add \(t(0,1,\dots,n-1)\) to the \(0\)-coordinates.
\item
  \emph{Conclusion.} (L) holds on \(\mathbb Re_0\), so \(\mathcal W=0\)
  on a dense set. \(\mathcal W\) is continuous on \(\mathrm{Conf}\)
  (VI.1.j).
\end{itemize}

\emph{Lean:}
\href{https://github.com/yang-mills-group/OpenAI-YM-20261028/blob/3456e5e75cf8d4cd558955a5c7638fba2d68f464/YM/Ward/Walk.lean\#L60-L71}{\texttt{Ward.stood}},
\href{https://github.com/yang-mills-group/OpenAI-YM-20261028/blob/3456e5e75cf8d4cd558955a5c7638fba2d68f464/YM/Ward/Walk.lean\#L75-L78}{\texttt{Walk}},
\href{https://github.com/yang-mills-group/OpenAI-YM-20261028/blob/3456e5e75cf8d4cd558955a5c7638fba2d68f464/YM/Ward/Walk.lean\#L79-L84}{\texttt{along}},
\href{https://github.com/yang-mills-group/OpenAI-YM-20261028/blob/3456e5e75cf8d4cd558955a5c7638fba2d68f464/YM/Ward/Axis.lean\#L36-L53}{\texttt{park1}},
\href{https://github.com/yang-mills-group/OpenAI-YM-20261028/blob/3456e5e75cf8d4cd558955a5c7638fba2d68f464/YM/Ward/Axis.lean\#L57-L74}{\texttt{pushes}},
\href{https://github.com/yang-mills-group/OpenAI-YM-20261028/blob/3456e5e75cf8d4cd558955a5c7638fba2d68f464/YM/Ward/Axis.lean\#L86-L108}{\texttt{line}},
\href{https://github.com/yang-mills-group/OpenAI-YM-20261028/blob/3456e5e75cf8d4cd558955a5c7638fba2d68f464/YM/Ward/Axis.lean\#L130-L139}{\texttt{global}}.

\paragraph{4.5 Brackets (A07, A08)}\label{brackets-a07-a08}

\textbf{Lemma VI.1.q (A07: erasure for the limit table).} The weak
action law, with its alignment of \(M\), \(\Gamma\), units and
inclusions, transports free row erasure to \(S\) and \(\partial\).

\emph{Proof sketch.}

\begin{itemize}
\tightlist
\item
  For \(d\ge N\), the free bilinear families of the kind in VI.1.4 form
  a finite-dimensional real space \(\mathcal M\). Free erasure means
  that the free row tests separate \(\mathcal M\), so finitely many of
  them do.
\item
  Along a scaling system, the transported actual rows converge to these,
  so they are injective on \(\mathcal M\) at some scale \(k\).
\item
  An actual family \(\Delta\) with vanishing row sums, pulled back by
  \(U_k\), lies in \(\mathcal M\), because \(U_k\) respects \(M\),
  \(\Gamma\), \(\mathbf 1\) and \(\iota\). Its rows are actual row sums
  at \(r_kx\), which vanish. So \(\Delta=0\).
\end{itemize}

\emph{Lean:}
\href{https://github.com/yang-mills-group/OpenAI-YM-20261028/blob/3456e5e75cf8d4cd558955a5c7638fba2d68f464/YM/Work/A07.lean\#L287-L289}{\texttt{Work.A07.answer}},
\href{https://github.com/yang-mills-group/OpenAI-YM-20261028/blob/3456e5e75cf8d4cd558955a5c7638fba2d68f464/YM/Work/A07.lean\#L243-L285}{\texttt{transfer}}.

\textbf{Lemma VI.1.r (A08: brackets).} This is Proposition VI.1.4.

\emph{Proof sketch.}

\begin{itemize}
\tightlist
\item
  Let \(\mathcal O_u\) be the orbital derivative along
  \(x_i\mapsto ux_i\) for all points at once. VI.1.3 says
  \(\mathcal O_uS(q)=-\sum_iS(\dots,T_uq_i,\dots)\) on the open set
  \(\mathrm{Conf}\).
\item
  Apply it twice:
  \(\mathcal O_u\mathcal O_vS=\sum_{i,j}S(\dots T_uq_j\dots T_vq_i\dots)\).
  Linear vector fields satisfy
  \([\mathcal O_u,\mathcal O_v]=-\mathcal O_{[u,v]}\).
\item
  Subtracting, the off-diagonal terms cancel. So the defect
  \(\Delta(u,v)=[T_u,T_v]-T_{[u,v]}\) has vanishing row sums.
\item
  By compatibility of \(T\) and
  \(\sigma[u,v]\sigma^{-1}=[\sigma u\sigma^{-1},\sigma v\sigma^{-1}]\),
  \(\Delta\) is a family of the kind in VI.1.4. Erasure gives
  \(\Delta=0\) for \(d\ge N\).
\item
  Lower degrees follow by injecting into degree \(\max(d,N)\).
\end{itemize}

\emph{Lean:}
\href{https://github.com/yang-mills-group/OpenAI-YM-20261028/blob/3456e5e75cf8d4cd558955a5c7638fba2d68f464/YM/Work/A08.lean\#L232-L247}{\texttt{Work.A08.answer}},
\href{https://github.com/yang-mills-group/OpenAI-YM-20261028/blob/3456e5e75cf8d4cd558955a5c7638fba2d68f464/YM/Work/A08.lean\#L62-L68}{\texttt{orbital\_comm}},
\href{https://github.com/yang-mills-group/OpenAI-YM-20261028/blob/3456e5e75cf8d4cd558955a5c7638fba2d68f464/YM/Work/A08.lean\#L91-L107}{\texttt{orbit\_twice}},
\href{https://github.com/yang-mills-group/OpenAI-YM-20261028/blob/3456e5e75cf8d4cd558955a5c7638fba2d68f464/YM/Work/A08.lean\#L192-L205}{\texttt{defect\_winds}},
\href{https://github.com/yang-mills-group/OpenAI-YM-20261028/blob/3456e5e75cf8d4cd558955a5c7638fba2d68f464/YM/Work/A08.lean\#L209-L230}{\texttt{defect\_ward}}.

\paragraph{4.6 Descent of the free rotation action
(A09--A12)}\label{descent-of-the-free-rotation-action-a09a12}

Fix \(d\). Index the generators by ordered pairs \(\pi=(i,j)\), with
\(K_\pi=E_{ij}-E_{ji}\) and \(C_\pi=\mathcal L^{\rm free}(K_\pi)\) on
\(F^{\rm free}_d\). For tuples \(B,C\) of operators on
\(F^{\rm free}_d\) define

\begin{itemize}
\tightlist
\item
  \(\mathrm{ad}_{B,C}(X)=BX-XC\) on \(\mathrm{End}\,F^{\rm free}_d\);
\item
  \(\Omega(B,C)=\sum_\pi\mathrm{ad}_{B_\pi,C_\pi}^2\), the relative
  Casimir.
\end{itemize}

A tuple satisfies the \emph{perfect relations} if
\(X_{ij}=[X_{ik},X_{kj}]\) for a fixed \(k\notin\{i,j\}\). Any
\(\mathfrak{so}(4)\) representation does (A09).

\textbf{Lemma VI.1.s (A12: the intertwiner projection).} There is a real
polynomial \(z\) such that \(P=z(\Omega(C,C))\) satisfies:

\begin{itemize}
\tightlist
\item
  \(P^2=P\);
\item
  \(P(\mathrm{id})=\mathrm{id}\);
\item
  \(\mathrm{ad}_{C_\pi,C_\pi}P=0\) for all \(\pi\).
\end{itemize}

So \(P\) projects onto the commutant of the \(C_\pi\), along
\(\mathrm{range}\,\Omega\).

\emph{Proof sketch.}

\begin{itemize}
\tightlist
\item
  Haar averaging over \(O(4)\) gives an inner product on
  \(\mathrm{End}\,F^{\rm free}_d\) in which every \(\mathrm{ad}\,C_\pi\)
  is skew (unitarian trick). So \(\Omega\) is self-adjoint,
  \(\ker\Omega=\bigcap_\pi\ker\mathrm{ad}\,C_\pi\), and
  \(\ker\Omega^2=\ker\Omega\).
\item
  Cayley--Hamilton, with the powers of \(X\) removed, gives a real \(z\)
  with \(z(0)=1\) and \(\Omega z(\Omega)=0\).
\end{itemize}

\emph{Lean:}
\href{https://github.com/yang-mills-group/OpenAI-YM-20261028/blob/3456e5e75cf8d4cd558955a5c7638fba2d68f464/YM/Work/A12.lean\#L323-L336}{\texttt{Work.A12.answer}},
\href{https://github.com/yang-mills-group/OpenAI-YM-20261028/blob/3456e5e75cf8d4cd558955a5c7638fba2d68f464/YM/Spin/Shape.lean\#L68-L71}{\texttt{Spin.Works}},
\href{https://github.com/yang-mills-group/OpenAI-YM-20261028/blob/3456e5e75cf8d4cd558955a5c7638fba2d68f464/YM/Spin/Shape.lean\#L57}{\texttt{wave}}.

\textbf{Lemma VI.1.t (approximate to exact intertwiners).}

\begin{itemize}
\tightlist
\item
  \begin{enumerate}
  \def\labelenumi{(\roman{enumi})}
  \tightlist
  \item
    \emph{Purification.} In a Banach algebra, let \(Z_n\to z\) with
    \(z^2=z\). There are \(C>0\) and \(n_0\) such that, for
    \(n\ge n_0\), the iterates of \(\varphi(x)=3x^2-2x^3\) starting at
    \(Z_n\) converge to idempotents \(P_n\to z\), with
    \(\|P_n-Z_n\|\le2C\|Z_n^2-Z_n\|\). This uses
    \(\varphi(x)-x=(1-2x)(x^2-x)\) and
    \(\mathrm{def}(\varphi x)=(4\,\mathrm{def}\,x-3)(\mathrm{def}\,x)^2\),
    where \(\mathrm{def}\,x=x^2-x\).
  \end{enumerate}
\item
  (ii) \emph{Norm gap.} If operators \(D_\pi\) satisfy the perfect
  relations and all \(\|D_\pi\|<\tfrac12\), then all \(D_\pi=0\), since
  \(\max\|D\|\le2\max\|D\|^2\).
\item
  (iii) \emph{Exact intertwiners.} Assume:

  \begin{itemize}
  \tightlist
  \item
    tuples \(B_n\to C\), with \(B_n\) (eventually) and \(C\) satisfying
    the perfect relations;
  \item
    \(P=z(\Omega(C,C))\) is as in VI.1.s;
  \item
    there is a closed additive subgroup
    \(\mathcal F\subset\mathrm{End}_{\mathbb C}F^{\rm free}_d\)
    containing \(\mathrm{id}\) such that, eventually,
    \(z(\Omega(B_n,C))\) preserves \(\mathcal F\) and commutes with
    every \(\mathrm{ad}_{B_{n,\pi},C_\pi}\).
  \end{itemize}

  Then for large \(n\) there is an invertible \(f\in\mathcal F\) with
  \(B_{n,\pi}f=fC_\pi\) for all \(\pi\).
\end{itemize}

\emph{Proof sketch of (iii).}

\begin{itemize}
\tightlist
\item
  Work in \(\mathrm{End}(\mathrm{End}\,F^{\rm free}_d)\) with
  \(D_\pi(n)=\mathrm{ad}_{B_{n,\pi},C_\pi}\). These satisfy the perfect
  relations and converge to \(\mathrm{ad}_{C_\pi,C_\pi}\).
\item
  \(z(\Omega(B_n,C))\to P\). Purify it to idempotents \(Q_n\). Every
  iterate preserves \(\mathcal F\) and commutes with \(D_\pi(n)\). The
  preserving operators form a closed subring of
  \(\mathrm{End}(\mathrm{End}\,F^{\rm free}_d)\), so these properties
  pass to \(Q_n\).
\item
  The compressions \(D_\pi(n)Q_n\) satisfy the perfect relations and
  tend in norm to \(\mathrm{ad}_{C_\pi,C_\pi}P=0\). By (ii) they vanish
  for large \(n\).
\item
  Then \(f_n=Q_n(\mathrm{id})\to P(\mathrm{id})=\mathrm{id}\), so
  \(f_n\) is invertible. It lies in \(\mathcal F\), and
  \(\mathrm{ad}_{B_{n,\pi},C_\pi}f_n=0\).
\item
  No complete reducibility of the representation \(D=c^{-1}Y\) on
  \(F_d\), and no Lie integration for it, is used.
\end{itemize}

\emph{Lean:}
\href{https://github.com/yang-mills-group/OpenAI-YM-20261028/blob/3456e5e75cf8d4cd558955a5c7638fba2d68f464/YM/Spin/Isolates.lean\#L25-L69}{\texttt{Spin.rectify}},
\href{https://github.com/yang-mills-group/OpenAI-YM-20261028/blob/3456e5e75cf8d4cd558955a5c7638fba2d68f464/YM/Spin/Isolates.lean\#L72-L98}{\texttt{stable}},
\href{https://github.com/yang-mills-group/OpenAI-YM-20261028/blob/3456e5e75cf8d4cd558955a5c7638fba2d68f464/YM/Spin/Close.lean\#L27-L38}{\texttt{enough}},
\href{https://github.com/yang-mills-group/OpenAI-YM-20261028/blob/3456e5e75cf8d4cd558955a5c7638fba2d68f464/YM/Spin/Close.lean\#L42-L70}{\texttt{onto}},
\href{https://github.com/yang-mills-group/OpenAI-YM-20261028/blob/3456e5e75cf8d4cd558955a5c7638fba2d68f464/YM/Spin/Shape.lean\#L77-L93}{\texttt{realize}},
\href{https://github.com/yang-mills-group/OpenAI-YM-20261028/blob/3456e5e75cf8d4cd558955a5c7638fba2d68f464/YM/Spin/Shape.lean\#L74-L75}{\texttt{Tracks}}.

\textbf{Lemma VI.1.u (descent of rotations; proof of VI.1.5).} Assume
the hypotheses of VI.1.2. The normalised rows \(D=c^{-1}Y\) satisfy the
global Ward identity by VI.1.3.

\begin{itemize}
\tightlist
\item
  \textbf{Choice of \(U_n\).} Choose one scaling system meeting the
  requests of VI.1.g at degrees \(d\) and \(4\), and aligned through
  degree \(d\). Put \(U_n=U_{n,d}\).
\item
  \textbf{Convergence.} At degree \(4\), \(p=c\,\partial\) and VI.1.g
  give \(c/\xi_k\to\kappa\). Hence
  \(U_n^{-1}DU_n\to\mathcal L^{\rm free}\).
\item
  \textbf{Inputs of VI.1.t(iii).} These are the \emph{tilt}
  (\href{https://github.com/yang-mills-group/OpenAI-YM-20261028/blob/3456e5e75cf8d4cd558955a5c7638fba2d68f464/YM/Spin/Sheets.lean\#L116-L122}{\texttt{Spin.Tilt}}):

  \begin{itemize}
  \tightlist
  \item
    \(U_n\) intertwines \(M\) and \(\Gamma\) (track);
  \item
    \(D\) satisfies the perfect relations, by A08 (with erasure from
    A07) and A09;
  \item
    the free data supply these relations and the Casimir polynomial of
    VI.1.s for \(\mathcal L^{\rm free}\);
  \item
    let \(\mathcal F\subset\mathrm{End}_{\mathbb C}F^{\rm free}_d\) be
    the closed real subspace of maps commuting with the free \(M\) and
    \(\Gamma\). The operators \(z(\Omega(U_n^{-1}DU_n,C))\) preserve
    \(\mathcal F\) (A11: real coefficients and cancellation of the
    signed-permutation factors in the sum of squares) and commute with
    the relative generators (A09).
  \end{itemize}
\item
  \textbf{Exact intertwiner.} VI.1.t gives \(f_n\). Put \(V_d=U_nf_n\).
  Then \(V_d\) intertwines \(M\) and \(\Gamma\), and
  \(D_{K_\pi}V_d=V_dC_\pi\). By linearity this holds for all \(K\).
\item
  \textbf{Port.} \(R_d=V_dR^{\rm free}_dV_d^{-1}\) inherits continuity,
  the group law, \(R(\sigma)=M_\sigma\), reality, and the generator
  \(D\).
\end{itemize}

\emph{Lean:}
\href{https://github.com/yang-mills-group/OpenAI-YM-20261028/blob/3456e5e75cf8d4cd558955a5c7638fba2d68f464/YM/Bid/Town.lean\#L124-L157}{\texttt{Bid.services}},
\href{https://github.com/yang-mills-group/OpenAI-YM-20261028/blob/3456e5e75cf8d4cd558955a5c7638fba2d68f464/YM/Bid/Town.lean\#L80-L90}{\texttt{Bid.currencies}},
\href{https://github.com/yang-mills-group/OpenAI-YM-20261028/blob/3456e5e75cf8d4cd558955a5c7638fba2d68f464/YM/Bid/Town.lean\#L102-L122}{\texttt{Bid.imaginary}},
\href{https://github.com/yang-mills-group/OpenAI-YM-20261028/blob/3456e5e75cf8d4cd558955a5c7638fba2d68f464/YM/Work/A09.lean\#L85-L90}{\texttt{Work.A09.answer}},
\href{https://github.com/yang-mills-group/OpenAI-YM-20261028/blob/3456e5e75cf8d4cd558955a5c7638fba2d68f464/YM/Work/A11.lean\#L199-L206}{\texttt{Work.A11.answer}},
\href{https://github.com/yang-mills-group/OpenAI-YM-20261028/blob/3456e5e75cf8d4cd558955a5c7638fba2d68f464/YM/Spin/Sheets.lean\#L55-L64}{\texttt{Spin.fine}},
\href{https://github.com/yang-mills-group/OpenAI-YM-20261028/blob/3456e5e75cf8d4cd558955a5c7638fba2d68f464/YM/Spin/Sheets.lean\#L134-L157}{\texttt{Spin.meet}},
\href{https://github.com/yang-mills-group/OpenAI-YM-20261028/blob/3456e5e75cf8d4cd558955a5c7638fba2d68f464/YM/Spin/Lands.lean\#L26-L47}{\texttt{Spin.produces}},
\href{https://github.com/yang-mills-group/OpenAI-YM-20261028/blob/3456e5e75cf8d4cd558955a5c7638fba2d68f464/YM/Climb/Beams.lean\#L90-L98}{\texttt{Climb.descends}}.

\paragraph{\texorpdfstring{4.7 Integration to \(O(4)\) (A13,
A14)}{4.7 Integration to O(4) (A13, A14)}}\label{integration-to-o4-a13-a14}

\textbf{Lemma VI.1.v (A13).} \(O(4)\) is generated by \(B_4\) and the
one-parameter groups \(e^{tK}\), \(K\in\mathfrak{so}(4)\).

\emph{Proof sketch.} \(O(4)\) is generated by hyperplane reflections
(Cartan--Dieudonné). The reflection in \(v^\perp\) equals
\(g^{-1}R_0g\), where \(R_0\) is a coordinate mirror and \(g\) is a
product of plane rotations \(e^{t(E_{ij}-E_{ji})}\) that turn \(v\) onto
\(e_0\) (Givens rotations).

\emph{Lean:}
\href{https://github.com/yang-mills-group/OpenAI-YM-20261028/blob/3456e5e75cf8d4cd558955a5c7638fba2d68f464/YM/Work/A13.lean\#L191-L200}{\texttt{Work.A13.answer}},
\href{https://github.com/yang-mills-group/OpenAI-YM-20261028/blob/3456e5e75cf8d4cd558955a5c7638fba2d68f464/YM/Work/A13.lean\#L151-L174}{\texttt{reflection\_conjugate}}.

\textbf{Lemma VI.1.w (A14: covariance and laws).} This is Proposition
VI.1.6.

\emph{Proof sketch.}

\begin{itemize}
\tightlist
\item
  The set of \(k\) with \(S(R(k)v;kx)=S(v;x)\) for all \(v,x\) is a
  subgroup of \(O(4)\).
\item
  It contains \(B_4\), since \(R(\sigma)=M_\sigma\) and \(S\) is
  \(B_4\)-covariant.
\item
  It contains each \(e^{tK}\). The function
  \(f(t)=S(R(e^{tK})v;e^{tK}x)\) is differentiable, and by the group law
  \(f'(t)\) is \(\mathcal W_K[T]\) at the moved data, which is \(0\). So
  \(f\) is constant.
\item
  By A13 the subgroup is all of \(O(4)\).
\item
  The rotation laws:

  \begin{itemize}
  \tightlist
  \item
    \(R(k)\mathbf 1=\mathbf 1\): apply covariance, delete the unit
    insertion, apply covariance again;
  \item
    commutation with \(\iota\): covariance plus the inclusion law
    (\href{https://github.com/yang-mills-group/OpenAI-YM-20261028/blob/3456e5e75cf8d4cd558955a5c7638fba2d68f464/YM/Spin/Lands.lean\#L50-L69}{\texttt{Spin.team}});
  \item
    the vector law: differentiate the covariance identity and use the
    chain rule for \(x\mapsto kx\);
  \item
    the remaining fields are those of the port.
  \end{itemize}
\end{itemize}

\emph{Lean:}
\href{https://github.com/yang-mills-group/OpenAI-YM-20261028/blob/3456e5e75cf8d4cd558955a5c7638fba2d68f464/YM/Work/A14.lean\#L277-L292}{\texttt{Work.A14.oriented}},
\href{https://github.com/yang-mills-group/OpenAI-YM-20261028/blob/3456e5e75cf8d4cd558955a5c7638fba2d68f464/YM/Work/A14.lean\#L245-L261}{\texttt{cov\_parameter}},
\href{https://github.com/yang-mills-group/OpenAI-YM-20261028/blob/3456e5e75cf8d4cd558955a5c7638fba2d68f464/YM/Work/A14.lean\#L310-L325}{\texttt{unit\_fixed}},
\href{https://github.com/yang-mills-group/OpenAI-YM-20261028/blob/3456e5e75cf8d4cd558955a5c7638fba2d68f464/YM/Work/A14.lean\#L364-L398}{\texttt{vectorLaw}},
\href{https://github.com/yang-mills-group/OpenAI-YM-20261028/blob/3456e5e75cf8d4cd558955a5c7638fba2d68f464/YM/Work/A14.lean\#L403-L415}{\texttt{laws\_of\_oriented}}.

\paragraph{4.8 The averaged stress tensor
(A15)}\label{the-averaged-stress-tensor-a15}

\textbf{Lemma VI.1.x (A15: Haar averaging).} Under the hypotheses of
VI.1.7, put \(J_0=c^{-1}J\), so that \((J_0;\partial,c^{-1}Y)\)
satisfies the stress laws. Define

\[J'=\int_{O(4)}k\cdot J_0\,d\mu(k),\qquad (k\cdot J)_{ab}=\sum_{ij}k_{ai}k_{bj}\,R_4(k)J_{ij},\]

with \(\mu\) normalised Haar measure on \(O(4)\). Then \(J'\) satisfies
Proposition VI.1.7.

\emph{Proof sketch.}

\begin{itemize}
\tightlist
\item
  \(k\cdot\) is a linear group action on the finite-dimensional space of
  stresses. Fixed points are exactly the tensorial stresses, and
  invariance of \(\mu\) makes \(J'\) fixed.
\item
  Each \(k\cdot J_0\) is conserved: symmetry by the index structure,
  reality by the reality law, divergence by the vector law.
\item
  Each \(k\cdot J_0\) satisfies the same stress laws. Substitute
  \(y\mapsto ky\) and use the window \(\chi\circ k\) at \(k^{-1}x\);
  pull the labels back by \(R(k^{-1})\); use covariance of \(S\).

  \begin{itemize}
  \tightlist
  \item
    The translation weight \(e_b\) becomes \(k^{-1}e_b\). The right side
    is independent of \(k\) by the vector law.
  \item
    The weight \(K(y-x_0)\) becomes \(k^{-1}Kk\) centred at
    \(k^{-1}x_0\). The right side is independent of \(k\) by A10:
    \(R(k)\,\mathcal L_K=\mathcal L_{kKk^{-1}}\,R(k)\).
  \end{itemize}
\item
  All the conditions are linear, so they pass to the integral.
\end{itemize}

\emph{Lean:}
\href{https://github.com/yang-mills-group/OpenAI-YM-20261028/blob/3456e5e75cf8d4cd558955a5c7638fba2d68f464/YM/Work/A15.lean\#L187-L192}{\texttt{Work.A15.answer}},
\href{https://github.com/yang-mills-group/OpenAI-YM-20261028/blob/3456e5e75cf8d4cd558955a5c7638fba2d68f464/YM/Work/A15.lean\#L22-L41}{\texttt{normalize\_flux}},
\href{https://github.com/yang-mills-group/OpenAI-YM-20261028/blob/3456e5e75cf8d4cd558955a5c7638fba2d68f464/YM/Work/A15.lean\#L86}{\texttt{averageStress}},
\href{https://github.com/yang-mills-group/OpenAI-YM-20261028/blob/3456e5e75cf8d4cd558955a5c7638fba2d68f464/YM/Work/A15.lean\#L117-L118}{\texttt{tensorial\_average}},
\href{https://github.com/yang-mills-group/OpenAI-YM-20261028/blob/3456e5e75cf8d4cd558955a5c7638fba2d68f464/YM/Work/A15.lean\#L134-L154}{\texttt{conserved\_average}},
\href{https://github.com/yang-mills-group/OpenAI-YM-20261028/blob/3456e5e75cf8d4cd558955a5c7638fba2d68f464/YM/Work/A15.lean\#L156-L183}{\texttt{stressLaws\_average}},
\href{https://github.com/yang-mills-group/OpenAI-YM-20261028/blob/3456e5e75cf8d4cd558955a5c7638fba2d68f464/YM/Work/A15/Transfer.lean\#L124-L168}{\texttt{rotate\_flux}},
\href{https://github.com/yang-mills-group/OpenAI-YM-20261028/blob/3456e5e75cf8d4cd558955a5c7638fba2d68f464/YM/Work/A10.lean\#L113-L115}{\texttt{Work.A10.answer}}.

\paragraph{4.9 Lorentz group, and
assembly}\label{lorentz-group-and-assembly}

\textbf{Lemma VI.1.y (A16, A17 and the descended Lorentz action).} This
is Proposition VI.1.8.

\emph{Proof sketch.}

\begin{itemize}
\tightlist
\item
  \emph{(a) Generation.} Let \(\Lambda\in SO^+(1,3)\). Then
  \(\Lambda e_t\) is a future unit timelike vector. At each step, apply
  a coordinate boost with rapidity \(-\mathrm{artanh}(y_i/y_t)\) to the
  current vector \(y\). This kills its \(i\)-th spatial component and
  preserves the other spatial components. After three steps the vector
  is \(e_t\). The resulting Lorentz transformation fixes \(e_t\), so it
  is a spatial rotation.
\item
  \emph{(b) Compatibility with inclusions.} Induct over the generators.

  \begin{itemize}
  \tightlist
  \item
    Spatial rotations: here \(B\) is the transported rotation action,
    which commutes with \(\iota\) by VI.1.6.
  \item
    Boosts: \(t\mapsto B_e(\mathrm{boost}_i(t))\iota v\) and
    \(t\mapsto\iota B_d(\mathrm{boost}_i(t))v\) both solve \(f'=G f\)
    with \(G=-i\,D_{K^{(3i)}}\), and they agree at \(t=0\). Here \(G\)
    commutes with \(\iota\), which follows by differentiating the
    inclusion law of \(R\). Uniqueness for the linear ODE finishes.
  \end{itemize}
\item
  \emph{(c) Descent.} Compatibility with inclusions and the unique lift
  to the union of label spaces give the representation. Its descent
  properties transfer from the free action through \(V\)
  (\href{https://github.com/yang-mills-group/OpenAI-YM-20261028/blob/3456e5e75cf8d4cd558955a5c7638fba2d68f464/YM/Spin/Floor.lean\#L30-L37}{\texttt{Spin.transmits}}).
\end{itemize}

\emph{Lean:}
\href{https://github.com/yang-mills-group/OpenAI-YM-20261028/blob/3456e5e75cf8d4cd558955a5c7638fba2d68f464/YM/Work/A16.lean\#L82-L89}{\texttt{exists\_rest}},
\href{https://github.com/yang-mills-group/OpenAI-YM-20261028/blob/3456e5e75cf8d4cd558955a5c7638fba2d68f464/YM/Work/A17.lean\#L70-L84}{\texttt{run\_boost}},
\href{https://github.com/yang-mills-group/OpenAI-YM-20261028/blob/3456e5e75cf8d4cd558955a5c7638fba2d68f464/YM/Work/A17.lean\#L57-L68}{\texttt{rise\_generator}},
\href{https://github.com/yang-mills-group/OpenAI-YM-20261028/blob/3456e5e75cf8d4cd558955a5c7638fba2d68f464/YM/Spin/Floor.lean\#L21-L22}{\texttt{Spin.run}},
\href{https://github.com/yang-mills-group/OpenAI-YM-20261028/blob/3456e5e75cf8d4cd558955a5c7638fba2d68f464/YM/Spin/Floor.lean\#L39-L74}{\texttt{Spin.reach}}.

\textbf{Lemma VI.1.z (assembly).} Under the hypotheses of VI.1.9, the
preceding constructions give its local field laws and Lorentz action.

\emph{Proof sketch.} A01--A03 and the collinear flux argument give the
raw compatibility laws. VI.1.h--i give \(c\) and the normalised Ward
data. VI.1.3 and VI.1.u give the intertwiners, rotation port and global
Ward identity. A13--A14 give covariance and the rotation laws; A10 and
A15 give the conserved tensorial stress with the required shells. Each
\(F_d\) has the free comparison isomorphism by any aligned scaling
system at that degree. A16--A17 give compatibility of the Lorentz
actions, and the unique lift to the union gives the final
representation.

\emph{Lean:}
\href{https://github.com/yang-mills-group/OpenAI-YM-20261028/blob/3456e5e75cf8d4cd558955a5c7638fba2d68f464/YM/Bid/Town.lean\#L159-L168}{\texttt{Bid.station}},
\href{https://github.com/yang-mills-group/OpenAI-YM-20261028/blob/3456e5e75cf8d4cd558955a5c7638fba2d68f464/YM/Org/Twist.lean\#L45-L53}{\texttt{Commission.picks}},
\href{https://github.com/yang-mills-group/OpenAI-YM-20261028/blob/3456e5e75cf8d4cd558955a5c7638fba2d68f464/YM/Climb/Beams.lean\#L53-L63}{\texttt{Climb.comparison}}.

\paragraph{4.10 Free-field stress
algebra}\label{free-field-stress-algebra}

Assume \(G\) is simple. The free stress satisfies the following
finite-dimensional algebraic facts:

\begin{itemize}
\tightlist
\item
  adjoint-invariant linear forms vanish, and invariant bilinear forms
  are multiples of the inner product.
\item
  Gauge-invariant composites of weight \(\le5\) have explicit normal
  forms.
\item
  The real symmetric \(B_4\)-covariant \(2\)-tensors of weight at most
  \(4\) form a space of dimension \(\le4\), spanned by the
  four-parameter family
\end{itemize}

\[T_{ab}=\delta_{ab}\Big(w_0+w_1\sum_i\langle F_{ai},F_{ai}\rangle+w_2\sum_{ij}\langle F_{ij},F_{ij}\rangle\Big)+(1-\delta_{ab})\,w_3\sum_i\langle F_{ai},F_{bi}\rangle .\]
Here \(F_{ab}\) is the free curvature and \(w\in\mathbb R^4\). The
Maxwell stress is \(w=(0,1,-\tfrac14,1)\). - Wick's theorem computes the
free one- and two-point tests. Centring and conservation force
\(w=w_3\cdot(0,1,-\tfrac14,1)\).

These facts feed VII.2's statement that, in the free table, the
divergence map from centred, real, symmetric, \(B_4\)-tensorial stresses
in \((F^{\rm free}_4)^{4\times4}\) to real \(B_4\)-vector fields in
\((F^{\rm free}_5)^4\) is onto and has a one-dimensional kernel
(\href{https://github.com/yang-mills-group/OpenAI-YM-20261028/blob/3456e5e75cf8d4cd558955a5c7638fba2d68f464/YM/Further/PHYSICAL_PERIOD_LIMIT_F933/FreeRanks.lean\#L159-L165}{\texttt{free\_rank}}).
They are not used in the Ward chain above.

\emph{Lean:}
\href{https://github.com/yang-mills-group/OpenAI-YM-20261028/blob/3456e5e75cf8d4cd558955a5c7638fba2d68f464/YM/Push/STRESS_FLUX/Color.lean\#L79-L91}{\texttt{bilinear\_invariant}},
\href{https://github.com/yang-mills-group/OpenAI-YM-20261028/blob/3456e5e75cf8d4cd558955a5c7638fba2d68f464/YM/Push/STRESS_FLUX/Counts.lean\#L129-L167}{\texttt{finrank\_two}},
\href{https://github.com/yang-mills-group/OpenAI-YM-20261028/blob/3456e5e75cf8d4cd558955a5c7638fba2d68f464/YM/Push/STRESS_FLUX/Counts.lean\#L190-L231}{\texttt{finrank\_one}},
\href{https://github.com/yang-mills-group/OpenAI-YM-20261028/blob/3456e5e75cf8d4cd558955a5c7638fba2d68f464/YM/Push/STRESS_FLUX/Polynomial.lean\#L202-L214}{\texttt{stress\_tensor}},
\href{https://github.com/yang-mills-group/OpenAI-YM-20261028/blob/3456e5e75cf8d4cd558955a5c7638fba2d68f464/YM/Push/STRESS_FLUX/FreeTests.lean\#L177-L184}{\texttt{div\_kernel}},
\href{https://github.com/yang-mills-group/OpenAI-YM-20261028/blob/3456e5e75cf8d4cd558955a5c7638fba2d68f464/YM/Push/STRESS_FLUX/FreeTests.lean\#L227-L235}{\texttt{stress\_bijective}}.
\emph{(proof of the divergence test read; polynomial classification and
dimension bounds: statement only)}

\subsubsection{5. Where the difficulty
lies}\label{where-the-difficulty-lies-23}

The exterior flux estimate (VI.1.e) bears most of the analytic burden.
Reflection positivity along axes of opposite parity gives the subpower
estimate without a gap. Diagonal stress components require a separate
conservation argument, since they have no odd reflection parity. The
widths \(\lambda^{100},\lambda^{150},\lambda^{300}\) and the split at
\(\lambda^{-3000}\) make the powers in these estimates compatible.

Continuation off the line (VI.1.j--p) requires more than a smooth
function with vanishing jets. The collinear Ward identities determine
the transverse jets of the two-point function; Hilbert-space positivity
turns those jets into uniform factorial derivative bounds and hence
analyticity.

The descent (VI.1.s--u) constructs exact intertwiners from the free
comparison. Lie brackets alone would give an action of a covering group
after integration; they would not by themselves identify the required
\(O(4)\) action or its signed-permutation subgroup. Covariance
subsequently forces compatibility between degrees.

The remaining steps use compactly supported integration by parts,
finite-dimensional tests, group generation, Haar averaging and
uniqueness for linear ODEs.

\subsubsection{6. Notation}\label{notation-23}

{\def\LTcaptype{none} % do not increment counter
\begin{longtable}[]{@{}
  >{\raggedright\arraybackslash}p{(\linewidth - 4\tabcolsep) * \real{0.3333}}
  >{\raggedright\arraybackslash}p{(\linewidth - 4\tabcolsep) * \real{0.3333}}
  >{\raggedright\arraybackslash}p{(\linewidth - 4\tabcolsep) * \real{0.3333}}@{}}
\toprule\noalign{}
\begin{minipage}[b]{\linewidth}\raggedright
object
\end{minipage} & \begin{minipage}[b]{\linewidth}\raggedright
code name
\end{minipage} & \begin{minipage}[b]{\linewidth}\raggedright
proposed notation
\end{minipage} \\
\midrule\noalign{}
\endhead
\bottomrule\noalign{}
\endlastfoot
label spaces, tags, inclusions, unit &
\href{https://github.com/yang-mills-group/OpenAI-YM-20261028/blob/3456e5e75cf8d4cd558955a5c7638fba2d68f464/YM/Euclid/Null.lean\#L74}{\texttt{Field\ d}},
\href{https://github.com/yang-mills-group/OpenAI-YM-20261028/blob/3456e5e75cf8d4cd558955a5c7638fba2d68f464/YM/Euclid/Lift.lean\#L66}{\texttt{Tag}},
\href{https://github.com/yang-mills-group/OpenAI-YM-20261028/blob/3456e5e75cf8d4cd558955a5c7638fba2d68f464/YM/Euclid/Differential.lean\#L20-L22}{\texttt{rise}},
\href{https://github.com/yang-mills-group/OpenAI-YM-20261028/blob/3456e5e75cf8d4cd558955a5c7638fba2d68f464/YM/Euclid/Null.lean\#L100}{\texttt{unit}}
& \(F_d\), \((d,v)\), \(\iota\), \(\mathbf 1\) \\
signed permutations, conjugation &
\href{https://github.com/yang-mills-group/OpenAI-YM-20261028/blob/3456e5e75cf8d4cd558955a5c7638fba2d68f464/YM/Space/Index.lean\#L23-L26}{\texttt{Signed}},
\href{https://github.com/yang-mills-group/OpenAI-YM-20261028/blob/3456e5e75cf8d4cd558955a5c7638fba2d68f464/YM/Euclid/Null.lean\#L138-L140}{\texttt{mirror}},
\href{https://github.com/yang-mills-group/OpenAI-YM-20261028/blob/3456e5e75cf8d4cd558955a5c7638fba2d68f464/YM/Euclid/Null.lean\#L142-L144}{\texttt{conj}}
& \(B_4\), \(M_\sigma\), \(\Gamma\) \\
differential plan and its laws &
\href{https://github.com/yang-mills-group/OpenAI-YM-20261028/blob/3456e5e75cf8d4cd558955a5c7638fba2d68f464/YM/Euclid/Differential.lean\#L20-L22}{\texttt{DiffPlan}},
\href{https://github.com/yang-mills-group/OpenAI-YM-20261028/blob/3456e5e75cf8d4cd558955a5c7638fba2d68f464/YM/Euclid/Differential.lean\#L61-L64}{\texttt{Derivatives}},
\href{https://github.com/yang-mills-group/OpenAI-YM-20261028/blob/3456e5e75cf8d4cd558955a5c7638fba2d68f464/YM/Euclid/Differential.lean\#L58-L60}{\texttt{Inclusions}},
\href{https://github.com/yang-mills-group/OpenAI-YM-20261028/blob/3456e5e75cf8d4cd558955a5c7638fba2d68f464/YM/Euclid/Differential.lean\#L66-L72}{\texttt{Orientation}}
& \(\partial_a\); derivative, inclusion, orientation laws \\
stress, window, shells &
\href{https://github.com/yang-mills-group/OpenAI-YM-20261028/blob/3456e5e75cf8d4cd558955a5c7638fba2d68f464/YM/Euclid/Stress.lean\#L57}{\texttt{Table.Stress}},
\href{https://github.com/yang-mills-group/OpenAI-YM-20261028/blob/3456e5e75cf8d4cd558955a5c7638fba2d68f464/YM/Euclid/Stress.lean\#L72-L73}{\texttt{Window}},
\href{https://github.com/yang-mills-group/OpenAI-YM-20261028/blob/3456e5e75cf8d4cd558955a5c7638fba2d68f464/YM/Euclid/Stress.lean\#L79-L81}{\texttt{shellGrad}},
\href{https://github.com/yang-mills-group/OpenAI-YM-20261028/blob/3456e5e75cf8d4cd558955a5c7638fba2d68f464/YM/Euclid/Stress.lean\#L84-L86}{\texttt{shellSpin}}
& \(J_{ab}\in F_4\), \(\chi\), \(\Phi^{\rm tr}_\chi\),
\(\Phi^{\rm rot}_\chi\) \\
raw shell maps; raw hypotheses &
\href{https://github.com/yang-mills-group/OpenAI-YM-20261028/blob/3456e5e75cf8d4cd558955a5c7638fba2d68f464/YM/Climb/Cells.lean\#L58}{\texttt{Climb.Parts}}
(\(p\)),
\href{https://github.com/yang-mills-group/OpenAI-YM-20261028/blob/3456e5e75cf8d4cd558955a5c7638fba2d68f464/YM/Ward/Sum.lean\#L17}{\texttt{Ward.SpinRows}}
(\(Y\));
\href{https://github.com/yang-mills-group/OpenAI-YM-20261028/blob/3456e5e75cf8d4cd558955a5c7638fba2d68f464/YM/Climb/Beams.lean\#L66-L69}{\texttt{Climb.Pre}}
& \(p_a\), \(Y_K\); (R) \\
angular sum; its vanishing &
\href{https://github.com/yang-mills-group/OpenAI-YM-20261028/blob/3456e5e75cf8d4cd558955a5c7638fba2d68f464/YM/Climb/Beams.lean\#L72-L74}{\texttt{localTurn}},
\href{https://github.com/yang-mills-group/OpenAI-YM-20261028/blob/3456e5e75cf8d4cd558955a5c7638fba2d68f464/YM/Ward/Sum.lean\#L25-L27}{\texttt{Ward.angular}};
\href{https://github.com/yang-mills-group/OpenAI-YM-20261028/blob/3456e5e75cf8d4cd558955a5c7638fba2d68f464/YM/Ward/Walk.lean\#L45}{\texttt{Still}},
\href{https://github.com/yang-mills-group/OpenAI-YM-20261028/blob/3456e5e75cf8d4cd558955a5c7638fba2d68f464/YM/Ward/Sum.lean\#L34-L35}{\texttt{Collinear}},
\href{https://github.com/yang-mills-group/OpenAI-YM-20261028/blob/3456e5e75cf8d4cd558955a5c7638fba2d68f464/YM/Ward/Sum.lean\#L29-L30}{\texttt{commutes}}
& \(\mathcal W_K[Y,p]\), \(\mathcal W_K[T]\); (L), (C) \\
line sums, transversality &
\href{https://github.com/yang-mills-group/OpenAI-YM-20261028/blob/3456e5e75cf8d4cd558955a5c7638fba2d68f464/YM/Bid/Angular.lean\#L35-L36}{\texttt{LineSum}},
\href{https://github.com/yang-mills-group/OpenAI-YM-20261028/blob/3456e5e75cf8d4cd558955a5c7638fba2d68f464/YM/Climb/Peaks.lean\#L57}{\texttt{Climb.transverse}}
& line sums, transverse \\
colour space and law &
\href{https://github.com/yang-mills-group/OpenAI-YM-20261028/blob/3456e5e75cf8d4cd558955a5c7638fba2d68f464/YM/Climb/Cells.lean\#L66-L87}{\texttt{Colors}},
\href{https://github.com/yang-mills-group/OpenAI-YM-20261028/blob/3456e5e75cf8d4cd558955a5c7638fba2d68f464/YM/Climb/Cells.lean\#L89-L91}{\texttt{ColorLaw}}
& colour space \\
normalised data &
\href{https://github.com/yang-mills-group/OpenAI-YM-20261028/blob/3456e5e75cf8d4cd558955a5c7638fba2d68f464/YM/Climb/Beams.lean\#L83-L87}{\texttt{Climb.Crystal}}
& \(p=c\,\partial\), \(D=c^{-1}Y\) \\
free data &
\href{https://github.com/yang-mills-group/OpenAI-YM-20261028/blob/3456e5e75cf8d4cd558955a5c7638fba2d68f464/YM/Org/Input.lean\#L159-L170}{\texttt{MaxwellData}},
\href{https://github.com/yang-mills-group/OpenAI-YM-20261028/blob/3456e5e75cf8d4cd558955a5c7638fba2d68f464/YM/Climb/Cells.lean\#L20}{\texttt{Climb.Bel}},
\href{https://github.com/yang-mills-group/OpenAI-YM-20261028/blob/3456e5e75cf8d4cd558955a5c7638fba2d68f464/YM/Climb/Beams.lean\#L24-L29}{\texttt{Lighthouse}},
\href{https://github.com/yang-mills-group/OpenAI-YM-20261028/blob/3456e5e75cf8d4cd558955a5c7638fba2d68f464/YM/Climb/Beams.lean\#L34-L42}{\texttt{Goes}}
&
\(F^{\rm free}_d,\partial^{\rm free},R^{\rm free},\mathcal L^{\rm free},J^{\rm free},\kappa\) \\
scaling system, track, weak law &
\href{https://github.com/yang-mills-group/OpenAI-YM-20261028/blob/3456e5e75cf8d4cd558955a5c7638fba2d68f464/YM/Climb/Cells.lean\#L24-L31}{\texttt{Climb.Scaling}},
\href{https://github.com/yang-mills-group/OpenAI-YM-20261028/blob/3456e5e75cf8d4cd558955a5c7638fba2d68f464/YM/Climb/Beams.lean\#L44-L48}{\texttt{Track}},
\href{https://github.com/yang-mills-group/OpenAI-YM-20261028/blob/3456e5e75cf8d4cd558955a5c7638fba2d68f464/YM/Climb/Beams.lean\#L50-L51}{\texttt{ActLaw}},
\href{https://github.com/yang-mills-group/OpenAI-YM-20261028/blob/3456e5e75cf8d4cd558955a5c7638fba2d68f464/YM/Climb/Tests.lean\#L49-L57}{\texttt{aisle}}
& \((r_k,\xi_k,U_{k,d})\), \(\mathrm{aisle}_{k,d}\) \\
cutoff field &
\href{https://github.com/yang-mills-group/OpenAI-YM-20261028/blob/3456e5e75cf8d4cd558955a5c7638fba2d68f464/YM/Next/ROTATIONS/Patch.lean\#L72}{\texttt{petal}},
\href{https://github.com/yang-mills-group/OpenAI-YM-20261028/blob/3456e5e75cf8d4cd558955a5c7638fba2d68f464/YM/Next/ROTATIONS/Route.lean\#L14}{\texttt{vane}},
\href{https://github.com/yang-mills-group/OpenAI-YM-20261028/blob/3456e5e75cf8d4cd558955a5c7638fba2d68f464/YM/Next/ROTATIONS/Budget.lean\#L15}{\texttt{fuel}}
& \(\mathrm{petal}_{ab}\), \(\lambda=e^k\) \\
OS vector of one insertion &
\href{https://github.com/yang-mills-group/OpenAI-YM-20261028/blob/3456e5e75cf8d4cd558955a5c7638fba2d68f464/YM/Ward/Plants.lean\#L55}{\texttt{spot}},
\href{https://github.com/yang-mills-group/OpenAI-YM-20261028/blob/3456e5e75cf8d4cd558955a5c7638fba2d68f464/YM/OS/Dynamics.lean\#L57-L61}{\texttt{OS.Chamber}}
& \(\mathrm{spot}(A,x)\in\mathcal H_a\) \\
row erasure &
\href{https://github.com/yang-mills-group/OpenAI-YM-20261028/blob/3456e5e75cf8d4cd558955a5c7638fba2d68f464/YM/Org/Input.lean\#L149-L152}{\texttt{erases}},
\href{https://github.com/yang-mills-group/OpenAI-YM-20261028/blob/3456e5e75cf8d4cd558955a5c7638fba2d68f464/YM/Org/Input.lean\#L141-L146}{\texttt{Winds}}
& erasure \\
compatibility of spin rows &
\href{https://github.com/yang-mills-group/OpenAI-YM-20261028/blob/3456e5e75cf8d4cd558955a5c7638fba2d68f464/YM/Bid/Angular.lean\#L16-L20}{\texttt{Bid.Hyper}}
& compatible \\
perfect relations, Casimir projection &
\href{https://github.com/yang-mills-group/OpenAI-YM-20261028/blob/3456e5e75cf8d4cd558955a5c7638fba2d68f464/YM/Spin/Sheets.lean\#L35}{\texttt{Spin.rel}},
\href{https://github.com/yang-mills-group/OpenAI-YM-20261028/blob/3456e5e75cf8d4cd558955a5c7638fba2d68f464/YM/Spin/Shape.lean\#L57}{\texttt{wave}},
\href{https://github.com/yang-mills-group/OpenAI-YM-20261028/blob/3456e5e75cf8d4cd558955a5c7638fba2d68f464/YM/Spin/Shape.lean\#L68-L71}{\texttt{Works}},
\href{https://github.com/yang-mills-group/OpenAI-YM-20261028/blob/3456e5e75cf8d4cd558955a5c7638fba2d68f464/YM/Spin/Shape.lean\#L21}{\texttt{effect}}
& perfect relations, \(z(\Omega(B,C))\), \(\mathrm{ad}_{B,C}\) \\
tilt, intertwiners &
\href{https://github.com/yang-mills-group/OpenAI-YM-20261028/blob/3456e5e75cf8d4cd558955a5c7638fba2d68f464/YM/Spin/Sheets.lean\#L116-L122}{\texttt{Spin.Tilt}},
\href{https://github.com/yang-mills-group/OpenAI-YM-20261028/blob/3456e5e75cf8d4cd558955a5c7638fba2d68f464/YM/Spin/Sheets.lean\#L50-L52}{\texttt{realH}},
\href{https://github.com/yang-mills-group/OpenAI-YM-20261028/blob/3456e5e75cf8d4cd558955a5c7638fba2d68f464/YM/Spin/Sheets.lean\#L44}{\texttt{change}}
& \(U_n\), \(V_d\) \\
port, transported action &
\href{https://github.com/yang-mills-group/OpenAI-YM-20261028/blob/3456e5e75cf8d4cd558955a5c7638fba2d68f464/YM/Spin/Lands.lean\#L13-L19}{\texttt{Spin.Port}},
\href{https://github.com/yang-mills-group/OpenAI-YM-20261028/blob/3456e5e75cf8d4cd558955a5c7638fba2d68f464/YM/Spin/Lands.lean\#L22-L24}{\texttt{Spin.morph}},
\href{https://github.com/yang-mills-group/OpenAI-YM-20261028/blob/3456e5e75cf8d4cd558955a5c7638fba2d68f464/YM/Org/Twist.lean\#L17}{\texttt{Org.mill}}
& \(R_d(k)=V_dR^{\rm free}_d(k)V_d^{-1}\) \\
rotation laws, covariance &
\href{https://github.com/yang-mills-group/OpenAI-YM-20261028/blob/3456e5e75cf8d4cd558955a5c7638fba2d68f464/YM/Euclid/Stress.lean\#L31-L42}{\texttt{Rotation.Laws}},
\href{https://github.com/yang-mills-group/OpenAI-YM-20261028/blob/3456e5e75cf8d4cd558955a5c7638fba2d68f464/YM/Euclid/Stress.lean\#L46-L48}{\texttt{Rotation.Oriented}}
& rotation laws; \(S(R(k)v;kx)=S(v;x)\) \\
conserved, tensorial; average &
\href{https://github.com/yang-mills-group/OpenAI-YM-20261028/blob/3456e5e75cf8d4cd558955a5c7638fba2d68f464/YM/Euclid/Stress.lean\#L59-L62}{\texttt{Conserved}},
\href{https://github.com/yang-mills-group/OpenAI-YM-20261028/blob/3456e5e75cf8d4cd558955a5c7638fba2d68f464/YM/Euclid/Stress.lean\#L64-L66}{\texttt{Tensorial}},
\href{https://github.com/yang-mills-group/OpenAI-YM-20261028/blob/3456e5e75cf8d4cd558955a5c7638fba2d68f464/YM/Work/A15.lean\#L86}{\texttt{averageStress}}
& \(J'=\int k\cdot c^{-1}J\,d\mu\) \\
Lorentz action on labels &
\href{https://github.com/yang-mills-group/OpenAI-YM-20261028/blob/3456e5e75cf8d4cd558955a5c7638fba2d68f464/YM/Spin/Floor.lean\#L21-L22}{\texttt{Spin.run}},
\href{https://github.com/yang-mills-group/OpenAI-YM-20261028/blob/3456e5e75cf8d4cd558955a5c7638fba2d68f464/YM/Spin/Floor.lean\#L24-L25}{\texttt{wand}},
\href{https://github.com/yang-mills-group/OpenAI-YM-20261028/blob/3456e5e75cf8d4cd558955a5c7638fba2d68f464/YM/Fields/Star.lean\#L22-L24}{\texttt{Boundary.Spins}},
\href{https://github.com/yang-mills-group/OpenAI-YM-20261028/blob/3456e5e75cf8d4cd558955a5c7638fba2d68f464/YM/Fields/Star.lean\#L26-L35}{\texttt{Descends}}
& \(B_d(\Lambda)\) \\
record &
\href{https://github.com/yang-mills-group/OpenAI-YM-20261028/blob/3456e5e75cf8d4cd558955a5c7638fba2d68f464/YM/Org/Twist.lean\#L29-L43}{\texttt{Org.Commission}},
\href{https://github.com/yang-mills-group/OpenAI-YM-20261028/blob/3456e5e75cf8d4cd558955a5c7638fba2d68f464/YM/Org/Twist.lean\#L19-L24}{\texttt{Org.Tab}},
\href{https://github.com/yang-mills-group/OpenAI-YM-20261028/blob/3456e5e75cf8d4cd558955a5c7638fba2d68f464/YM/Spec/Closed.lean\#L34-L43}{\texttt{Spec.Locals}}
& \((c,V,J')\); local laws \\
\end{longtable}
}

\subsubsection{7. What remains
unexplained}\label{what-remains-unexplained-23}

\begin{enumerate}
\def\labelenumi{\arabic{enumi}.}
\tightlist
\item
  \textbf{Supporting analysis in the exterior flux estimate (probable).}
  The reductions to spectral integrals, the point and line estimates,
  the diagonal mass estimate and their combination in VI.1.e were
  checked against the cited proofs. The proofs of the remote-gap
  profiles, the odd temporal variant, the contact extension on product
  tests, and the finite-difference and interpolation estimates were not
  checked in full. These are the files \texttt{GapTest},
  \texttt{Nodal}/\texttt{Mode}/\texttt{Odd},
  \texttt{Collision}/\texttt{Fold}/\texttt{Journal}, \texttt{MomentFD}
  and \texttt{Nodes}. In particular, the construction and normalisation
  of the contact extension remain inputs to the mass-estimate sketch.
\item
  \textbf{Free-field stress algebra (probable).} The statements in §4.10
  and the proof that centring and the two divergence tests select the
  Maxwell line were checked. The polynomial normal-form and
  dimension-count proofs were not checked in full. The map in the cited
  rank theorem has centred stresses as its domain; centring cannot be
  omitted from the one-dimensional-kernel claim.
\item
  \textbf{Production of the inputs.} The immediate producers of the raw
  stress, shell maps, weak comparison and free data are identified in
  §2. Their ultraviolet hypotheses, including the annular estimates and
  finite comparison systems, are not discharged in this unit. The unused
  lattice shell-law interfaces in \texttt{STRESS\_FLUX} are not the
  source of (R).
\item
  \textbf{Normalisation and centring.} The Ward coefficient of \(J'\) is
  exactly \(1\). Centring is not part of the conclusion of VI.1.7; the
  centred trace is formed in VII.2. The stress replacement used before
  the weak comparison supplies \(\kappa\ne0\) without the numerical
  identity \(\kappa=-\tau\). Along the scaling systems used in VI.1.u,
  positivity of \(\xi_k\) and \(c/\xi_k\to\kappa\ne0\) imply that \(c\)
  and \(\kappa\) have the same sign. Neither sign is specified by this
  unit's hypotheses.
\item
  \textbf{Scope.} The folder ownership and line counts in the header
  have not been independently recounted. The assignment of
  \texttt{Shift} and most of \texttt{Next/BOUNDARY} to V.2 is tentative.
  The statements above are source comparisons; Lean elaboration and the
  transitive axiom dependencies were not checked.
\end{enumerate}

\subsection{VI.2. The free Maxwell stress tensor and the local Ward
charge}\label{vi.2.-the-free-maxwell-stress-tensor-and-the-local-ward-charge}

\textbf{Part:} VI, symmetry: stress tensor and Ward identities.
\textbf{Lean:}
\href{https://github.com/yang-mills-group/OpenAI-YM-20261028/tree/3456e5e75cf8d4cd558955a5c7638fba2d68f464/YM/Next/MAXWELL}{\texttt{YM/Next/MAXWELL}},
\href{https://github.com/yang-mills-group/OpenAI-YM-20261028/tree/3456e5e75cf8d4cd558955a5c7638fba2d68f464/YM/UVNext/FREE_STRESS_CERTIFICATE}{\texttt{YM/UVNext/FREE\_STRESS\_CERTIFICATE}},
\href{https://github.com/yang-mills-group/OpenAI-YM-20261028/tree/3456e5e75cf8d4cd558955a5c7638fba2d68f464/YM/Next/UV_STRESS}{\texttt{YM/Next/UV\_STRESS}},
\href{https://github.com/yang-mills-group/OpenAI-YM-20261028/tree/3456e5e75cf8d4cd558955a5c7638fba2d68f464/YM/Pert}{\texttt{YM/Pert}}.
\textbf{Size:} 14,179 lines used by the main theorem (Maxwell 8.2k,
free-stress classification 3.1k, UV stress 2.6k, Pert 0.2k).

\subsubsection{1. What this unit proves}\label{what-this-unit-proves-24}

The free Maxwell theory supplies the comparison model for the stress
tensor of the Wilson limit. Its local Ward identities fix the sign and
normalization of the free stress. For a simple gauge group, a
finite-dimensional calculation identifies its centred conserved tensors
of degree at most four as a line. Short-distance comparison transfers
this line to the interacting limit; the translation Ward identity then
fixes the normalization of the interacting stress.

Throughout, \(G\) is a gauge group as in I.1, \(\mathfrak g\) is its Lie
algebra, and \(c=1,\dots,\dim\mathfrak g\) indexes an orthonormal colour
basis for the invariant inner product. Spacetime indices range over
\(0,1,2,3\). Inner products of \(\mathfrak g\)-valued forms include the
colour sum.

A \emph{Maxwell wave} on an open set \(U\subset\mathbb R^4\) is a smooth
\(\mathfrak g\)-valued antisymmetric field \(f_{ab}\) satisfying
\[\partial_af_{bc}+\partial_bf_{ca}+\partial_cf_{ab}=0,\qquad (\mathrm{Div}\,f)_b:=\sum_a\partial_af_{ab}=0.\]
Here \(\mathrm{Div}=-d^*\) with the usual Euclidean codifferential
convention. The space \(\mathrm{Sh}_j\) consists of the homogeneous
Taylor coefficients of order \(j\) satisfying these linear equations,
not the cumulative jets through order \(j\). The algebra \(\mathcal W\)
consists of complex polynomials in coordinates on these spaces; a
coordinate from \(\mathrm{Sh}_j\) has weight \(j+2\). Its \emph{shell
sector} \(\mathcal P_d\) is the image of labels of degree at most \(d\)
under substitution of on-shell Taylor coefficients. Equivalently, it
consists of the \(\mathrm{Ad}(G)\)-invariant polynomials of weight at
most \(d\)
(\href{https://github.com/yang-mills-group/OpenAI-YM-20261028/blob/3456e5e75cf8d4cd558955a5c7638fba2d68f464/YM/Next/MAXWELL/Algebra.lean\#L108-L116}{\texttt{sector}}).
A \emph{probe} of order \(j\) is a real linear functional on
\(\mathrm{Sh}_j\); \(p(f)(x)\) denotes its value on the order-\(j\)
Taylor coefficient of \(f\) at \(x\).

The curvature propagator is
\[\langle F^c_{im}(x)F^{c'}_{ab}(y)\rangle_0=\delta^{cc'}\mathcal C_{im,ab}(x-y),\qquad G_0(x)=\frac1{4\pi^2|x|^2},\qquad-\Delta G_0=\delta,\]
where \(\mathcal C\) is the curl--curl of \(G_0\). The generated field
\(f^{p,x}(y)=\langle p(F)(x)F(y)\rangle_0\) is a Maxwell wave off \(x\).
Correlations of Wick polynomials at distinct sites sum over complete
pairings of legs between different sites; there are no contractions
within a site.

Write
\[T^{\rm free}_{ab}=\sum_{c,i}F^c_{ai}F^c_{bi}-\frac14\delta_{ab}\sum_{c,i,k}F^c_{ik}F^c_{ik}\in\mathcal P_4.\]
Its polarization and the current associated with a vector field \(v\)
are
\[\Theta_{ab}(f,g)=\sum_i\bigl(\langle f_{ai},g_{bi}\rangle+\langle f_{bi},g_{ai}\rangle\bigr)-\frac12\delta_{ab}\sum_{i,k}\langle f_{ik},g_{ik}\rangle,\qquad j^v_a(f,g)=\sum_bv_b\Theta_{ab}(f,g).\]
Thus \(\Theta(f,f)=2T^{\rm free}(f)\). Both tensors are symmetric and
traceless. For a real cutoff \(\chi\in C_c^\infty(\mathbb R^4)\), the
outward flux is \(\Phi_\chi(j)=-\sum_a\int\partial_a\chi\,j_a\).

The free table on a register \(L\) is \(\mathcal M=\mathrm{Bel}(L)\),
with filtered field spaces \(\mathcal F^{\rm free}_d\). The Wilson limit
table \(S\) has spaces \(\mathcal F^S_d\). For a lighthouse
\(\mathcal T\), write \(T_0\) for its free stress and \(\kappa\) for its
Ward constant. Let \(\mathrm{Conf}_n\) denote the distinct-point
configurations. For \(R>0\), set
\[K_{R,n}=\{y\in(\mathbb R^4)^n:\ |y_i-y_j|_\infty\ge R^{-1}\ (i\ne j),\ \|y\|\le R\},\]
where \(\|y\|\) is the Euclidean norm of the full configuration.

\textbf{Proposition VI.2.1 (the free Maxwell table).} For every register
\(L\), the free table \(\mathcal M=\mathrm{Bel}(L)\) has the following
properties.

\begin{enumerate}
\def\labelenumi{(\alph{enumi})}
\item
  For every \(d\ge0\), substitution of on-shell Taylor coefficients
  induces a complex linear isomorphism
  \(\iota_d:\mathcal F^{\rm free}_d\to\mathcal P_d\). A label is null
  exactly when its polynomial vanishes on every on-shell Taylor tower.
  At distinct points, field correlations equal the Wick moments of their
  images under \(\iota\).
\item
  The table satisfies the Euclidean hierarchy of IV.1: permutation
  symmetry, the unit laws, covariance under signed coordinate
  permutations, reality, translation invariance, point and smeared
  reflection positivity, and smoothness on \(\mathrm{Conf}_n\). The
  correlations are in fact real-analytic there. For each fixed tuple of
  fields and derivative order, there are constants \(A\ge0\) and
  \(C\ge0\) bounding that derivative by \(Cr^{-A}\) on configurations
  with mutual sup-distances at least \(r\), for \(0<r\le1\).
\item
  Differentiation of Taylor coefficients induces maps
  \(\partial^{\rm free}_a:\mathcal F^{\rm free}_d\to\mathcal F^{\rm free}_{d+1}\)
  satisfying the inclusion, derivative and orientation laws. In
  particular,
  \[\langle(\partial^{\rm free}_a B)(x_0)\cdots\rangle^{\rm free}=\partial_{x_0^a}\langle B(x_0)\cdots\rangle^{\rm free}.\]
  The action of \(O(4)\) on coefficients induces actions \(R_g\) on all
  field spaces and generators \(\rho(\omega)\) for
  \(\omega\in\mathfrak{so}(4)\). These satisfy the rotation laws of
  VI.1: continuity, the group law, extension of signed permutations,
  reality, preservation of the unit and inclusions, vector covariance of
  derivatives, and the infinitesimal-generator identity. All free
  correlations are rotation invariant.
\item
  The fields \(\iota_4^{-1}(T^{\rm free}_{ab})\), also denoted
  \(T^{\rm free}_{ab}\), are symmetric, real, traceless and centred, and
  satisfy \(\sum_a\partial^{\rm free}_aT^{\rm free}_{ab}=0\) in
  \(\mathcal F^{\rm free}_5\).
\end{enumerate}

\emph{Status:} theorem for every gauge group in the setting of I.1;
simplicity is not required. \emph{Lean:}
\href{https://github.com/yang-mills-group/OpenAI-YM-20261028/blob/3456e5e75cf8d4cd558955a5c7638fba2d68f464/YM/Next/MAXWELL/System.lean\#L106-L114}{\texttt{Maxwell.Free.hierarchy}},
\href{https://github.com/yang-mills-group/OpenAI-YM-20261028/blob/3456e5e75cf8d4cd558955a5c7638fba2d68f464/YM/Next/MAXWELL/System.lean\#L33-L35}{\texttt{fields}},
\href{https://github.com/yang-mills-group/OpenAI-YM-20261028/blob/3456e5e75cf8d4cd558955a5c7638fba2d68f464/YM/Next/MAXWELL/Caps.lean\#L90-L113}{\texttt{no\_residual}},
\href{https://github.com/yang-mills-group/OpenAI-YM-20261028/blob/3456e5e75cf8d4cd558955a5c7638fba2d68f464/YM/Next/MAXWELL/System.lean\#L80-L83}{\texttt{scheme}},
\href{https://github.com/yang-mills-group/OpenAI-YM-20261028/blob/3456e5e75cf8d4cd558955a5c7638fba2d68f464/YM/Next/MAXWELL/System.lean\#L153-L170}{\texttt{gear}},
\href{https://github.com/yang-mills-group/OpenAI-YM-20261028/blob/3456e5e75cf8d4cd558955a5c7638fba2d68f464/YM/Next/MAXWELL/Beacon.lean\#L54-L61}{\texttt{gear\_laws},
\texttt{gear\_motion}},
\href{https://github.com/yang-mills-group/OpenAI-YM-20261028/blob/3456e5e75cf8d4cd558955a5c7638fba2d68f464/YM/Next/MAXWELL/System.lean\#L172}{\texttt{original}},
\href{https://github.com/yang-mills-group/OpenAI-YM-20261028/blob/3456e5e75cf8d4cd558955a5c7638fba2d68f464/YM/Next/MAXWELL/System.lean\#L176-L179}{\texttt{realStress}}.

\textbf{Proposition VI.2.2 (local Ward charge).} Let
\(U\subset\mathbb R^4\) be open, \(x\in U\), and \(f\) a Maxwell wave on
\(U\). Let \(v\) be a smooth conformal Killing field on \(\mathbb R^4\):
\[\partial_av_b=-\partial_bv_a\quad(a\ne b),\qquad \partial_0v_0=\partial_1v_1=\partial_2v_2=\partial_3v_3.\]
Let \(\chi\in C_c^\infty(U;\mathbb R)\) equal one near \(x\). Since
\(df=0\), its Lie derivative is
\[(\mathcal L_vf)_{il}=\sum_b\bigl[\partial_i(v_bf_{bl})-\partial_l(v_bf_{bi})\bigr]=(d\iota_vf)_{il}.\]

\begin{enumerate}
\def\labelenumi{(\alph{enumi})}
\item
  For every zero-order component probe \(F^c_{il}\),
  \[\Phi_\chi\bigl(j^v(f^{F^c_{il},x},f)\bigr)=-(\mathcal L_vf)^c_{il}(x).\]
\item
  If \(\mathcal L_vf\) is also a Maxwell wave on \(U\), then, for every
  order \(j\ge0\) and every probe \(p\) of that order,
  \[\Phi_\chi\bigl(j^v(f^{p,x},f)\bigr)=-p(\mathcal L_vf)(x).\] For
  translations \(v=e_a\), this is
  \(-\partial_a[p(f)](https://github.com/scottnarmstrong/yang-mills-informalization/blob/main/outline/units/x)\).
  For rotations \(v(y)=\omega(y-x)\), it is \(-(\omega\cdot p)(f)(x)\),
  where \(\omega\cdot p\) is the infinitesimal label action. For
  dilation \(v(y)=y-x\), it is \(-(j+2)p(f)(x)\).
\item
  For any two Maxwell waves \(f,g\) on a neighbourhood of
  \(\mathrm{supp}\,\chi\), the flux \(\Phi_\chi(j^v(f,g))\) vanishes. It
  also vanishes for two fields \(f^{p,x},f^{q,x}\) generated at the same
  point, provided \(v\) is homogeneous about \(x\) of degree
  \(e\in\{0,1\}\).
\end{enumerate}

\emph{Status:} theorem, with the propagator, polarization and
outward-flux conventions above. \emph{Lean:}
\href{https://github.com/yang-mills-group/OpenAI-YM-20261028/blob/3456e5e75cf8d4cd558955a5c7638fba2d68f464/YM/Next/MAXWELL/Contact.lean\#L137-L156}{\texttt{Maxwell.Free.local\_force}},
\href{https://github.com/yang-mills-group/OpenAI-YM-20261028/blob/3456e5e75cf8d4cd558955a5c7638fba2d68f464/YM/Next/MAXWELL/Force.lean\#L131-L186}{\texttt{jet\_force}},
\href{https://github.com/yang-mills-group/OpenAI-YM-20261028/blob/3456e5e75cf8d4cd558955a5c7638fba2d68f464/YM/Next/MAXWELL/Affine.lean\#L249-L255}{\texttt{drift\_flux}},
\href{https://github.com/yang-mills-group/OpenAI-YM-20261028/blob/3456e5e75cf8d4cd558955a5c7638fba2d68f464/YM/Next/MAXWELL/Angular.lean\#L181-L195}{\texttt{angular\_flux}},
\href{https://github.com/yang-mills-group/OpenAI-YM-20261028/blob/3456e5e75cf8d4cd558955a5c7638fba2d68f464/YM/Next/MAXWELL/Affine.lean\#L220-L247}{\texttt{radial\_flux}},
\href{https://github.com/yang-mills-group/OpenAI-YM-20261028/blob/3456e5e75cf8d4cd558955a5c7638fba2d68f464/YM/Next/MAXWELL/Current.lean\#L239-L250}{\texttt{sourcefree\_flux}},
\href{https://github.com/yang-mills-group/OpenAI-YM-20261028/blob/3456e5e75cf8d4cd558955a5c7638fba2d68f464/YM/Next/MAXWELL/Repeat.lean\#L63-L128}{\texttt{two\_flux\_zero}};
definitions
\href{https://github.com/yang-mills-group/OpenAI-YM-20261028/blob/3456e5e75cf8d4cd558955a5c7638fba2d68f464/YM/Next/MAXWELL/Current.lean\#L140-L143}{\texttt{Killing}},
\href{https://github.com/yang-mills-group/OpenAI-YM-20261028/blob/3456e5e75cf8d4cd558955a5c7638fba2d68f464/YM/Next/MAXWELL/Regulator.lean\#L49}{\texttt{variation}},
\href{https://github.com/yang-mills-group/OpenAI-YM-20261028/blob/3456e5e75cf8d4cd558955a5c7638fba2d68f464/YM/Next/MAXWELL/Kernel.lean\#L101-L102}{\texttt{pulse}},
\href{https://github.com/yang-mills-group/OpenAI-YM-20261028/blob/3456e5e75cf8d4cd558955a5c7638fba2d68f464/YM/Next/MAXWELL/Current.lean\#L187}{\texttt{singleFlux}}.

\textbf{Proposition VI.2.3 (Ward identities of the free table).} Let
\(n\ge0\), \(x\in\mathrm{Conf}_{n+1}\), and
\(p_0,\dots,p_n\in\mathcal W\). Let
\(\chi\in C_c^\infty(\mathbb R^4;\mathbb R)\) equal one near \(x_0\),
with \(x_1,\dots,x_n\notin\mathrm{supp}\,\chi\). Take one of the pairs
\[(v,D)=(e_a,\partial_a),\qquad (\omega(y-x_0),\rho(\omega)),\qquad (y-x_0,\varepsilon),\]
where \(a\) is an axis, \(\omega\in\mathfrak{so}(4)\), and
\(\varepsilon\) multiplies each monomial by its weight. The operators
act on \(\mathcal W\) by derivations. Then
\[-\sum_{a,b}\int\partial_a\chi(y)\,v_b(y)\,\langle T^{\rm free}_{ab}(y)p_0(x_0)\cdots p_n(x_n)\rangle_0\,dy=-\langle(Dp_0)(x_0)p_1(x_1)\cdots p_n(x_n)\rangle_0.\]
Consequently, for every \(\tau\in\mathbb R\), the stress
\(\tau T^{\rm free}\) satisfies the translation and rotation Ward
identities with rows \(-\tau\partial^{\rm free}\) and \(-\tau\rho\). The
stress normalized to give \(+\partial^{\rm free}\) in the translation
identity is \(-T^{\rm free}\).

\emph{Status:} theorem. \emph{Lean:}
\href{https://github.com/yang-mills-group/OpenAI-YM-20261028/blob/3456e5e75cf8d4cd558955a5c7638fba2d68f464/YM/Next/MAXWELL/Taxes.lean\#L177-L200}{\texttt{Maxwell.Free.moment\_force}},
\href{https://github.com/yang-mills-group/OpenAI-YM-20261028/blob/3456e5e75cf8d4cd558955a5c7638fba2d68f464/YM/Next/MAXWELL/Taxes.lean\#L34-L60}{\texttt{straight\_motion},
\texttt{rotational\_motion}, \texttt{radial\_motion}},
\href{https://github.com/yang-mills-group/OpenAI-YM-20261028/blob/3456e5e75cf8d4cd558955a5c7638fba2d68f464/YM/Next/MAXWELL/Beacon.lean\#L74-L136}{\texttt{beacon\_signal}},
\href{https://github.com/yang-mills-group/OpenAI-YM-20261028/blob/3456e5e75cf8d4cd558955a5c7638fba2d68f464/YM/Euclid/Stress.lean\#L90-L97}{\texttt{Table.StressLaws}}.

\textbf{Proposition VI.2.4 (the free stress line).} Suppose \(G\) is
simple. Let \(W\) be a derivative plan on the free table satisfying the
derivative and orientation laws. The derivative law determines its
derivative maps uniquely as \(\partial^{\rm free}\). Define the
following real vector spaces:

\begin{itemize}
\tightlist
\item
  \(\mathcal Z\) consists of symmetric real tensors
  \(J\in(\mathcal F^{\rm free}_4)^{4\times4}\) that are covariant under
  signed coordinate permutations and centred:
  \(\langle J_{ab}\rangle=0\).
\item
  \(\mathcal V\) consists of real vectors in
  \((\mathcal F^{\rm free}_5)^4\) covariant under signed coordinate
  permutations.
\end{itemize}

The divergence \(\mathrm{div}:\mathcal Z\to\mathcal V\),
\((\mathrm{div}J)_b=\sum_a W_aJ_{ab}\), is surjective, and
\[\dim_{\mathbb R}\mathcal Z=3,\qquad \dim_{\mathbb R}\mathcal V=2,\qquad \ker\mathrm{div}=\mathbb R T^{\rm free}.\]
Thus a lighthouse with such a derivative plan and stress
\(\tau T^{\rm free}\), \(\tau\ne0\), satisfies (FD): its divergence is
onto, its kernel has real dimension one, and its stress is a nonzero
centred conserved tensor.

\emph{Status:} theorem. \emph{Lean:}
\href{https://github.com/yang-mills-group/OpenAI-YM-20261028/blob/3456e5e75cf8d4cd558955a5c7638fba2d68f464/YM/UVNext/FREE_STRESS_CERTIFICATE/Certificate.lean\#L43-L47}{\texttt{FreeStress.rank\_on}},
\href{https://github.com/yang-mills-group/OpenAI-YM-20261028/blob/3456e5e75cf8d4cd558955a5c7638fba2d68f464/YM/UVNext/FREE_STRESS_CERTIFICATE/Certificate.lean\#L73-L79}{\texttt{line\_dim}},
\href{https://github.com/yang-mills-group/OpenAI-YM-20261028/blob/3456e5e75cf8d4cd558955a5c7638fba2d68f464/YM/UVNext/FREE_STRESS_CERTIFICATE/Certificate.lean\#L86-L94}{\texttt{lighting}},
\href{https://github.com/yang-mills-group/OpenAI-YM-20261028/blob/3456e5e75cf8d4cd558955a5c7638fba2d68f464/YM/UVNext/FREE_STRESS_CERTIFICATE/Certificate.lean\#L125-L132}{\texttt{certificate}},
\href{https://github.com/yang-mills-group/OpenAI-YM-20261028/blob/3456e5e75cf8d4cd558955a5c7638fba2d68f464/YM/Next/UV_STRESS/Dimension.lean\#L153-L158}{\texttt{FreeDiv}},
\href{https://github.com/yang-mills-group/OpenAI-YM-20261028/blob/3456e5e75cf8d4cd558955a5c7638fba2d68f464/YM/Next/UV_STRESS/Tensor.lean\#L75-L80}{\texttt{KernelSpace}},
\href{https://github.com/yang-mills-group/OpenAI-YM-20261028/blob/3456e5e75cf8d4cd558955a5c7638fba2d68f464/YM/Next/UV_STRESS/Tensor.lean\#L87-L97}{\texttt{Vectors}},
\href{https://github.com/yang-mills-group/OpenAI-YM-20261028/blob/3456e5e75cf8d4cd558955a5c7638fba2d68f464/YM/Next/UV_STRESS/Tensor.lean\#L117-L119}{\texttt{currentDiv}}.

\textbf{Proposition VI.2.5 (Maxwell data).} Suppose \(G\) is simple, let
\(C\) be a Wilson family for \(G\), and let
\(\tau\in\mathbb R\setminus\{0\}\).

\begin{enumerate}
\def\labelenumi{(\alph{enumi})}
\tightlist
\item
  The lighthouse
  \(\mathcal T_\tau=(\partial^{\rm free},R,\tau T^{\rm free},\kappa=-\tau)\)
  satisfies all lighthouse laws of VI.1. These include the derivative,
  rotation and Ward laws of VI.2.1--3, a nonzero derivative in degree
  four, and the following selection property. There is \(N\) such that
  for every \(d\ge N\), finitely many derivatives of Ward sums at
  collinear configurations detect the line
  \(\mathbb R\partial^{\rm free}\) among complex linear maps
  \(t_a:\mathcal F^{\rm free}_d\to\mathcal F^{\rm free}_{d+1}\)
  satisfying:
\end{enumerate}

\begin{itemize}
\tightlist
\item
  they kill the unit, are real, and transform as a vector under signed
  permutations;
\item
  for every \(e\le d\), \(t_a\) maps the image of
  \(\mathcal F^{\rm free}_e\) into the image of
  \(\mathcal F^{\rm free}_{e+1}\).
\end{itemize}

For component \(t_b\), a diagram is an iterated directional derivative
of
\[x\longmapsto\sum_i\langle w_1(x_1)\cdots(t_bw_i)(x_i)\cdots w_n(x_n)\rangle^{\rm free},\qquad w_i\in\mathcal F^{\rm free}_d,\]
evaluated at a distinct-point configuration on a coordinate axis
perpendicular to \(e_b\). The selected diagrams all vanish exactly when
\(t\) is a real multiple of the derivative tuple.

\begin{enumerate}
\def\labelenumi{(\alph{enumi})}
\setcounter{enumi}{1}
\item
  There exists a Maxwell datum for \(C\) with lighthouse
  \(\mathcal T_\tau\). It also contains the free Euclidean hierarchy,
  the Casimir polynomial data of VI.1, compatible representations of the
  proper orthochronous Lorentz group on the field spaces and their
  union, continuous in each degree, and the following Ward-vanishing
  property. For every \(d\ge4\), let \(T(u,v)\) be a real-bilinear
  family of complex linear endomorphisms of \(\mathcal F^{\rm free}_d\),
  with \(u,v\in\mathfrak{so}(4)\). Suppose it is alternating, real,
  equivariant under orientation-preserving signed permutations, kills
  the unit, and preserves every lower-degree image. If
  \[\sum_i\langle w_1(x_1)\cdots(T(u,v)w_i)(x_i)\cdots w_m(x_m)\rangle=0\]
  for every \(m\ge0\), every \(x\in\mathrm{Conf}_m\), every
  \(w_i\in\mathcal F^{\rm free}_d\) and every \(u,v\), then \(T=0\). The
  Lorentz representation agrees with \(R\) on spatial rotations,
  commutes with time-reflected conjugation, and has boost generator
  \(-i\rho\) in each time--space plane.
\item
  Every Maxwell datum can be replaced by one with the same derivative
  plan, rotations, hierarchy and Lorentz representation, stress
  \(\tau T^{\rm free}\), full rotation covariance of that stress, and
  (FD). The Ward constant of any Maxwell datum with stress
  \(\tau T^{\rm free}\) is \(-\tau\).
\end{enumerate}

\emph{Status:} theorem. The construction used in III.11 chooses
\(\tau=1\) and hence \(\kappa=-1\). \emph{Lean:}
\href{https://github.com/yang-mills-group/OpenAI-YM-20261028/blob/3456e5e75cf8d4cd558955a5c7638fba2d68f464/YM/Next/MAXWELL/Lit.lean\#L29-L34}{\texttt{Maxwell.Free.herald}},
\href{https://github.com/yang-mills-group/OpenAI-YM-20261028/blob/3456e5e75cf8d4cd558955a5c7638fba2d68f464/YM/Next/MAXWELL/Lit.lean\#L36-L42}{\texttt{shines}},
\href{https://github.com/yang-mills-group/OpenAI-YM-20261028/blob/3456e5e75cf8d4cd558955a5c7638fba2d68f464/YM/Next/MAXWELL/Lit.lean\#L45-L48}{\texttt{maxwell\_supplies}},
\href{https://github.com/yang-mills-group/OpenAI-YM-20261028/blob/3456e5e75cf8d4cd558955a5c7638fba2d68f464/YM/Next/MAXWELL/AwayTeam.lean\#L48-L118}{\texttt{erasing}},
\href{https://github.com/yang-mills-group/OpenAI-YM-20261028/blob/3456e5e75cf8d4cd558955a5c7638fba2d68f464/YM/Next/MAXWELL/Scenes.lean\#L125-L132}{\texttt{scanner}},
\href{https://github.com/yang-mills-group/OpenAI-YM-20261028/blob/3456e5e75cf8d4cd558955a5c7638fba2d68f464/YM/Next/MAXWELL/Descent.lean\#L179-L181}{\texttt{launched}},
\href{https://github.com/yang-mills-group/OpenAI-YM-20261028/blob/3456e5e75cf8d4cd558955a5c7638fba2d68f464/YM/UVNext/FREE_STRESS_CERTIFICATE/Choose.lean\#L137-L150}{\texttt{FreeStress.fire}},
\href{https://github.com/yang-mills-group/OpenAI-YM-20261028/blob/3456e5e75cf8d4cd558955a5c7638fba2d68f464/YM/UVNext/FREE_STRESS_CERTIFICATE/Pressure.lean\#L167-L171}{\texttt{names}},
\href{https://github.com/yang-mills-group/OpenAI-YM-20261028/blob/3456e5e75cf8d4cd558955a5c7638fba2d68f464/YM/UVNext/FREE_STRESS_CERTIFICATE/Pressure.lean\#L158-L161}{\texttt{required}};
records
\href{https://github.com/yang-mills-group/OpenAI-YM-20261028/blob/3456e5e75cf8d4cd558955a5c7638fba2d68f464/YM/Org/Input.lean\#L159-L170}{\texttt{Org.MaxwellData}},
\href{https://github.com/yang-mills-group/OpenAI-YM-20261028/blob/3456e5e75cf8d4cd558955a5c7638fba2d68f464/YM/Org/Input.lean\#L149-L152}{\texttt{Org.erases}},
\href{https://github.com/yang-mills-group/OpenAI-YM-20261028/blob/3456e5e75cf8d4cd558955a5c7638fba2d68f464/YM/Climb/Beams.lean\#L24-L29}{\texttt{Climb.Lighthouse}},
\href{https://github.com/yang-mills-group/OpenAI-YM-20261028/blob/3456e5e75cf8d4cd558955a5c7638fba2d68f464/YM/Climb/Beams.lean\#L34-L42}{\texttt{Lighthouse.Goes}}.

\textbf{Proposition VI.2.6 (strict reflection positivity).} Fix a
coordinate axis \(a\), and let \(\theta_a\) reflect in \(y_a=0\). Let
\(\vartheta_a\) act on polynomials by the corresponding jet reflection
and coefficient conjugation.

\begin{enumerate}
\def\labelenumi{(\alph{enumi})}
\item
  For every \(x\) with \(x_a>0\) and every nonzero \(p\in\mathcal W\),
  \[\langle(\vartheta_ap)(\theta_ax)p(x)\rangle_0>0.\]
\item
  For every finite set of distinct coordinate probes of the spaces
  \(\mathrm{Sh}_j\), the reflected propagator Gram matrix at
  \((\theta_ax,x)\) is positive definite.
\item
  The free table is reflection positive for finite families of words in
  the positive half-space, with distinct sites within each word.
\item
  For every \(d\ge0\), a polynomial \(p\in\mathcal P_d\) whose
  correlations with every finite family of spectators from
  \(\mathcal P_d\) vanish at every distinct-point configuration is zero.
\end{enumerate}

\emph{Status:} theorem. \emph{Lean:}
\href{https://github.com/yang-mills-group/OpenAI-YM-20261028/blob/3456e5e75cf8d4cd558955a5c7638fba2d68f464/YM/Next/MAXWELL/Core.lean\#L94-L110}{\texttt{detects}},
\href{https://github.com/yang-mills-group/OpenAI-YM-20261028/blob/3456e5e75cf8d4cd558955a5c7638fba2d68f464/YM/Next/MAXWELL/Gram.lean\#L82-L105}{\texttt{coordGram}},
\href{https://github.com/yang-mills-group/OpenAI-YM-20261028/blob/3456e5e75cf8d4cd558955a5c7638fba2d68f464/YM/Next/MAXWELL/Caps.lean\#L137-L155}{\texttt{positive\_table}},
\href{https://github.com/yang-mills-group/OpenAI-YM-20261028/blob/3456e5e75cf8d4cd558955a5c7638fba2d68f464/YM/Next/MAXWELL/Statistic.lean\#L80-L99}{\texttt{ghosts\_out}}.

\textbf{Proposition VI.2.7 (a conserved stress on the limit table).} Let
\(C\) be a Wilson family satisfying the construction conditions of II.2.
Let \(F\) be an extraction order with ordered smeared convergence and
generic point limits to a table \(S\) with Euclidean hierarchy \(H\).
Let \(W\) be a derivative plan on \(S\) satisfying the inclusion and
derivative laws; its orientation law then follows from IV.1. Let \(H_0\)
be the free Euclidean hierarchy, and \(\mathcal T\) a lighthouse whose
plan satisfies the inclusion, derivative and orientation laws and for
which (FD) holds. Assume the station property of VI.2.8 and
\[\dim_{\mathbb C}\mathcal F^S_d\le\dim_{\mathbb C}\mathcal F^{\rm free}_d\qquad(d\ge0).\]

\begin{enumerate}
\def\labelenumi{(\alph{enumi})}
\item
  There exists a scaling \(X=(r_k,\xi_k,U_{k,d})\), where \(r_k>0\)
  decreases strictly to zero, \(\xi_k>0\), and
  \(U_{k,d}:\mathcal F^{\rm free}_d\to\mathcal F^S_d\) is a complex
  linear isomorphism for every \(k,d\). It is \emph{approaching}: for
  each fixed tuple \(v=(v_i)\) of free fields, the difference
  \[K_k(y)=\langle U_kv_1(r_ky_1)\cdots U_kv_n(r_ky_n)\rangle^S-\langle v_1(y_1)\cdots v_n(y_n)\rangle^{\rm free}\]
  satisfies, for every compact \(K\subset\mathrm{Conf}_n\),
  \[\sup_{\substack{t\in C_c^\infty(\mathrm{Conf}_n;\mathbb C)\\ \mathrm{supp}\,t\subset K,\ \|t\|_\infty\le1}}\left|\int t(y)K_k(y)\,dy\right|\longrightarrow0.\]
  It is \emph{stable}: the maps eventually commute with conjugation and
  signed permutations in each degree, preserve the unit in each degree,
  and intertwine each fixed inclusion of degrees. The initial choice may
  have \(\xi_k=1\).
\item
  For every approaching stable scaling, the divergence
  \(\mathrm{div}^S:\mathcal Z^S\to\mathcal V^S\) is onto and has a
  one-dimensional real kernel. These spaces are defined as in VI.2.4,
  using \(S\) and \(W\).
\item
  There exist a nonzero \(J\in\mathcal Z^S\) with \(\mathrm{div}^S J=0\)
  and an approaching stable refinement \(Y\) of \(X\), obtained by
  taking a subsequence and changing \(\xi_k\), such that
  \[\xi_k^{-1}U_{k,4}^{-1}(r_k^4J_{ab})\longrightarrow T_{0,ab}\qquad\text{for all }a,b.\]
\item
  For every finite arity bound \(N\), degree bound \(d\), derivative
  order \(m\), and \(R\ge1\), a tail of \(Y\) is a track: it intertwines
  symmetries, units and inclusions through degree \(\max(d,4)\),
  satisfies (c), and has convergence of all correlation derivatives of
  orders \(0\le j\le m\) uniformly on \(K_{R,n}\) for \(n\le N\). The
  convergence is for each fixed tuple of fields of degrees at most
  \(d\); its tail may depend on that tuple and the error tolerance. In
  particular, \(J\) satisfies the comparison law of VI.1.
\end{enumerate}

\emph{Status:} theorem under these hypotheses. III.11 supplies the
station estimates; the rank-corrected widths of the construction give
the dimension bound; VI.2.4 gives (FD) for the canonical datum. To
obtain the infrared and ultraviolet data used later, III.11 also
supplies the shell estimates needed for the Ward rows of \(J\).
\emph{Lean:}
\href{https://github.com/yang-mills-group/OpenAI-YM-20261028/blob/3456e5e75cf8d4cd558955a5c7638fba2d68f464/YM/Next/UV_STRESS/Invert.lean\#L132-L181}{\texttt{Stations.equivalences}},
\href{https://github.com/yang-mills-group/OpenAI-YM-20261028/blob/3456e5e75cf8d4cd558955a5c7638fba2d68f464/YM/Next/UV_STRESS/Dimension.lean\#L165-L195}{\texttt{differential\_rank}},
\href{https://github.com/yang-mills-group/OpenAI-YM-20261028/blob/3456e5e75cf8d4cd558955a5c7638fba2d68f464/YM/Next/UV_STRESS/Limit.lean\#L129-L167}{\texttt{produce\_current}},
\href{https://github.com/yang-mills-group/OpenAI-YM-20261028/blob/3456e5e75cf8d4cd558955a5c7638fba2d68f464/YM/Next/UV_STRESS/Limit.lean\#L197-L202}{\texttt{focused}},
\href{https://github.com/yang-mills-group/OpenAI-YM-20261028/blob/3456e5e75cf8d4cd558955a5c7638fba2d68f464/YM/Next/UV_STRESS/Invert.lean\#L110-L115}{\texttt{cap\_given}};
records
\href{https://github.com/yang-mills-group/OpenAI-YM-20261028/blob/3456e5e75cf8d4cd558955a5c7638fba2d68f464/YM/Climb/Cells.lean\#L24-L31}{\texttt{Climb.Scaling}},
\href{https://github.com/yang-mills-group/OpenAI-YM-20261028/blob/3456e5e75cf8d4cd558955a5c7638fba2d68f464/YM/Climb/Cells.lean\#L54-L55}{\texttt{Scaling.meter}},
\href{https://github.com/yang-mills-group/OpenAI-YM-20261028/blob/3456e5e75cf8d4cd558955a5c7638fba2d68f464/YM/Next/UV_STRESS/Rates.lean\#L52-L53}{\texttt{Approaches}},
\href{https://github.com/yang-mills-group/OpenAI-YM-20261028/blob/3456e5e75cf8d4cd558955a5c7638fba2d68f464/YM/Next/UV_STRESS/Invert.lean\#L121-L125}{\texttt{Stable}},
\href{https://github.com/yang-mills-group/OpenAI-YM-20261028/blob/3456e5e75cf8d4cd558955a5c7638fba2d68f464/YM/Climb/Beams.lean\#L50-L51}{\texttt{ActLaw}},
\href{https://github.com/yang-mills-group/OpenAI-YM-20261028/blob/3456e5e75cf8d4cd558955a5c7638fba2d68f464/YM/Org/Input.lean\#L172-L175}{\texttt{Org.IR}},
\href{https://github.com/yang-mills-group/OpenAI-YM-20261028/blob/3456e5e75cf8d4cd558955a5c7638fba2d68f464/YM/Org/Input.lean\#L177-L186}{\texttt{Org.UV}}.

\textbf{Proposition VI.2.8 (weak comparison stations).} Fix
\(C,F,S,H,H_0\) as in VI.2.7. A \emph{station} of depth \(D\) consists
of \(r>0\) and complex linear lifts
\(A_{N,p,d}:\mathcal F^{\rm free}_d\to\mathrm{Ref}_d\) for \(d\le D\) at
every lattice cutoff \((N,p)\).

It is \emph{admissible} for arity bound \(n_{\max}\), radius \(R\ge1\)
and tolerance \(\varepsilon>0\) if:

\begin{itemize}
\tightlist
\item
  the lifts are eventually uniformly bounded along the extraction order,
  real and signed-covariant, and preserve the degree-zero unit;
\item
  for every \(n\le n_{\max}\), every tuple \(v_i\) of degrees at most
  \(D\), every smooth complex test \(t\) supported in \(K_{R,n}\), every
  \(M>0\) with \(\|t\|_\infty\le M\), and every \(\eta>0\), eventually
  along the ordered cutoff tails,
  \[\left|\mathcal S_{N,p}\bigl(Av;r^{-4n}t(\cdot/r)\bigr)-\int t(y)\langle v_1(y_1)\cdots v_n(y_n)\rangle^{\rm free}\,dy\right|\le\varepsilon M\prod_i\|v_i\|+\eta.\]
  Here \(\mathcal S_{N,p}\) is the lattice-smeared correlation of IV.1.
  The cutoff tail may depend on \(v,t,M,\eta\).
\end{itemize}

The \emph{station property} requires an admissible station with
\(r\le\delta\) for every \(D,n_{\max}\), \(R\ge1\) and
\(\varepsilon,\delta>0\).

\begin{enumerate}
\def\labelenumi{(\alph{enumi})}
\item
  An admissible station yields a panel
  \(U_d:\mathcal F^{\rm free}_d\to\mathcal F^S_d\) for \(d\le D\). It is
  real, signed-covariant, preserves the degree-zero unit, and satisfies
  \[\left|\int t(y)\bigl[\langle Uv_1(ry_1)\cdots Uv_n(ry_n)\rangle^S-\langle v_1(y_1)\cdots v_n(y_n)\rangle^{\rm free}\bigr]dy\right|\le\varepsilon M\prod_i\|v_i\|\]
  for the same tuples and tests.
\item
  Quantitative grid, strip and finite-difference estimates satisfying
  the balance conditions in VI.2.s, for every finite arity bound, degree
  bound and \(R\ge1\), imply the station property.
\end{enumerate}

\emph{Status:} the station and balance conditions are definitions; (a)
and (b) are theorems. III.11 supplies the balanced estimates through its
complete weak-comparison result. \emph{Lean:}
\href{https://github.com/yang-mills-group/OpenAI-YM-20261028/blob/3456e5e75cf8d4cd558955a5c7638fba2d68f464/YM/Next/UV_STRESS/Station.lean\#L68-L71}{\texttt{Station}},
\href{https://github.com/yang-mills-group/OpenAI-YM-20261028/blob/3456e5e75cf8d4cd558955a5c7638fba2d68f464/YM/Next/UV_STRESS/Station.lean\#L89-L97}{\texttt{Station.Admissible}},
\href{https://github.com/yang-mills-group/OpenAI-YM-20261028/blob/3456e5e75cf8d4cd558955a5c7638fba2d68f464/YM/Next/UV_STRESS/Station.lean\#L112-L114}{\texttt{Stations}},
\href{https://github.com/yang-mills-group/OpenAI-YM-20261028/blob/3456e5e75cf8d4cd558955a5c7638fba2d68f464/YM/Next/UV_STRESS/Station.lean\#L132-L180}{\texttt{harvested}},
\href{https://github.com/yang-mills-group/OpenAI-YM-20261028/blob/3456e5e75cf8d4cd558955a5c7638fba2d68f464/YM/Next/UV_STRESS/Granular.lean\#L30-L34}{\texttt{Granular}},
\href{https://github.com/yang-mills-group/OpenAI-YM-20261028/blob/3456e5e75cf8d4cd558955a5c7638fba2d68f464/YM/Next/UV_STRESS/Granular.lean\#L134-L140}{\texttt{realizes\_stations}},
\href{https://github.com/yang-mills-group/OpenAI-YM-20261028/blob/3456e5e75cf8d4cd558955a5c7638fba2d68f464/YM/MergeProof/East/Gran.lean\#L35}{\texttt{MergeProof.East.CompleteWeak}}.

\textbf{Proposition VI.2.9 (the normalized stress and its short-distance
law).} Let \(C\) be a Wilson family, \(F\) an extraction order, \(S\) a
table with derivative plan \(W\), and \(M\) a Maxwell datum with
lighthouse \(\mathcal T\) and Ward constant \(\kappa\ne0\). Let
\(I=(J_0,p^I,Y^I)\) and let \(X\) be a scaling. Assume the
ultraviolet-data conditions of VIII.1 for \((C,F,S,W,M,I)\). In
particular, they give a Euclidean hierarchy on \(S\), the orientation
law for \(W\), a symmetric real conserved signed-covariant \(J_0\) with
Ward rows \((p^I,Y^I)\), and tracks comparing \(J_0\) to \(T_0\). Assume
also the inclusion and derivative laws for \(W\), and the following
additional properties: \(J_0\) is a nonzero member of \(\mathcal Z^S\);
\(\mathcal T\) satisfies (FD); and \(X\) is approaching and stable, with
\[\xi_k^{-1}U_{k,4}^{-1}(r_k^4J_{0,ab})\longrightarrow T_{0,ab}.\]

Let \(J\in(\mathcal F^S_4)^{4\times4}\) be symmetric, real, conserved
and covariant under signed permutations. Suppose it satisfies the
translation Ward identity with rows \(W\) and the rotation Ward identity
with some rotation rows. Then:

\begin{enumerate}
\def\labelenumi{(\alph{enumi})}
\item
  \(U_{k,4}^{-1}(r_k^4J_{ab})\to\kappa^{-1}T_{0,ab}\) for every \(a,b\).
\item
  For every finite tuple of free fields \(q_i\), every assignment of
  slots to either \(U_kq_i\) or \(r_k^4J_{a_ib_i}\), and every
  \(j\ge0\), the resulting correlations at \(r_ky_i\) converge with
  their order-\(j\) derivatives uniformly on compact subsets of
  \(\mathrm{Conf}_n\). The limit is the free correlation with each
  stress slot replaced by \(\kappa^{-1}T_{0,a_ib_i}\).
\item
  For the canonical datum with plan \(\partial^{\rm free}\) and stress
  \(T^{\rm free}\), \(\kappa=-1\). Thus the Ward-normalized physical
  stress converges to \(-T^{\rm free}\) without the factor
  \(\xi_k^{-1}\).
\end{enumerate}

\emph{Status:} theorem. VI.1 supplies a stress \(J\) satisfying these
Ward laws. The mixed-correlation conclusion follows from (a) and VI.2.u;
the linked mixed-correlation theorem states its specialization to that
stress. \emph{Lean:}
\href{https://github.com/yang-mills-group/OpenAI-YM-20261028/blob/3456e5e75cf8d4cd558955a5c7638fba2d68f464/YM/Pert/Weak.lean\#L118-L144}{\texttt{Pert.meter\_physical}},
\href{https://github.com/yang-mills-group/OpenAI-YM-20261028/blob/3456e5e75cf8d4cd558955a5c7638fba2d68f464/YM/Pert/Lifted.lean\#L85-L124}{\texttt{Pert.mixed\_leading}},
\href{https://github.com/yang-mills-group/OpenAI-YM-20261028/blob/3456e5e75cf8d4cd558955a5c7638fba2d68f464/YM/Pert/Lifted.lean\#L65-L74}{\texttt{kappa\_eq}},
\href{https://github.com/yang-mills-group/OpenAI-YM-20261028/blob/3456e5e75cf8d4cd558955a5c7638fba2d68f464/YM/Pert/Lifted.lean\#L38-L49}{\texttt{common\_stress}},
\href{https://github.com/yang-mills-group/OpenAI-YM-20261028/blob/3456e5e75cf8d4cd558955a5c7638fba2d68f464/YM/Pert/Weak.lean\#L22-L32}{\texttt{Pert.Supplier}},
\href{https://github.com/yang-mills-group/OpenAI-YM-20261028/blob/3456e5e75cf8d4cd558955a5c7638fba2d68f464/YM/Org/Twist.lean\#L19-L24}{\texttt{Org.Tab}}.

\subsubsection{2. What it uses}\label{what-it-uses-24}

\textbf{From other units.}

\begin{itemize}
\tightlist
\item
  III.16 supplies on-shell Taylor spaces, label frames, substitution
  into curvature polynomials, and the cross-site Wick pairing formula:
  \href{https://github.com/yang-mills-group/OpenAI-YM-20261028/blob/3456e5e75cf8d4cd558955a5c7638fba2d68f464/YM/Maxwell/Wick.lean\#L126-L127}{\texttt{YM.Maxwell.Wick.moment}};
  \href{https://github.com/yang-mills-group/OpenAI-YM-20261028/blob/3456e5e75cf8d4cd558955a5c7638fba2d68f464/YM/Maxwell/Equation.lean\#L66-L67}{\texttt{YM.Maxwell.observable}}.
\item
  IV.1 supplies tables, quotient fields, ordered cutoff limits and
  lattice smearing. The derivative and inclusion laws give the
  orientation law:
  \href{https://github.com/yang-mills-group/OpenAI-YM-20261028/blob/3456e5e75cf8d4cd558955a5c7638fba2d68f464/YM/Work/D06.lean\#L137-L139}{\texttt{YM.Work.D06.answer}}.
  Its rank-corrected widths give the upper dimension bound used here:
  \href{https://github.com/yang-mills-group/OpenAI-YM-20261028/blob/3456e5e75cf8d4cd558955a5c7638fba2d68f464/YM/Next/UV_STRESS/Invert.lean\#L110-L115}{\texttt{YM.Next.UVStress.cap\_given}}.
\item
  V.1 gives smeared reflection positivity from point reflection
  positivity and derivative bounds from order-zero bounds:
  \href{https://github.com/yang-mills-group/OpenAI-YM-20261028/blob/3456e5e75cf8d4cd558955a5c7638fba2d68f464/YM/Work/B08.lean\#L193-L203}{\texttt{YM.Work.B08.smear}};
  \href{https://github.com/yang-mills-group/OpenAI-YM-20261028/blob/3456e5e75cf8d4cd558955a5c7638fba2d68f464/YM/Work/B10.lean\#L95-L105}{\texttt{YM.Work.B10.answer}}.
\item
  VI.1 supplies local windows, finite separating correlation tests,
  collinear Ward diagrams, and the derivative, rotation and track laws.
  Rotation laws follow from invariance:
  \href{https://github.com/yang-mills-group/OpenAI-YM-20261028/blob/3456e5e75cf8d4cd558955a5c7638fba2d68f464/YM/Work/A14.lean\#L403-L415}{\texttt{YM.Work.A14.laws\_of\_oriented}}.
  It also supplies the Casimir relations and annihilating polynomials:
  \href{https://github.com/yang-mills-group/OpenAI-YM-20261028/blob/3456e5e75cf8d4cd558955a5c7638fba2d68f464/YM/Work/A09.lean\#L85-L90}{\texttt{YM.Work.A09.answer}};
  \href{https://github.com/yang-mills-group/OpenAI-YM-20261028/blob/3456e5e75cf8d4cd558955a5c7638fba2d68f464/YM/Work/A10.lean\#L113-L115}{\texttt{YM.Work.A10.answer}};
  \href{https://github.com/yang-mills-group/OpenAI-YM-20261028/blob/3456e5e75cf8d4cd558955a5c7638fba2d68f464/YM/Work/A12.lean\#L323-L336}{\texttt{YM.Work.A12.answer}}.
\item
  III.10 assembles Maxwell data from these components:
  \href{https://github.com/yang-mills-group/OpenAI-YM-20261028/blob/3456e5e75cf8d4cd558955a5c7638fba2d68f464/YM/Bid/Scope.lean\#L24-L29}{\texttt{YM.Bid.masts}}.
\item
  III.11 supplies balanced weak-comparison estimates and the shell
  estimates for the interacting Ward identities. It combines these with
  the present stress-line construction to obtain the infrared and
  ultraviolet data:
  \href{https://github.com/yang-mills-group/OpenAI-YM-20261028/blob/3456e5e75cf8d4cd558955a5c7638fba2d68f464/YM/MergeProof/East/Gran.lean\#L35}{\texttt{YM.MergeProof.East.CompleteWeak}};
  \href{https://github.com/yang-mills-group/OpenAI-YM-20261028/blob/3456e5e75cf8d4cd558955a5c7638fba2d68f464/YM/MergeProof/Emanate.lean\#L15-L30}{\texttt{YM.MergeProof.Acme.readyA}};
  \href{https://github.com/yang-mills-group/OpenAI-YM-20261028/blob/3456e5e75cf8d4cd558955a5c7638fba2d68f464/YM/MergeProof/Emanate.lean\#L31-L37}{\texttt{YM.MergeProof.Acme.emanate}}.
\item
  VII.5 supplies the comparison of stress and derivative normalizations
  along tracks and the nontriviality test for a Ward-normalized stress:
  \href{https://github.com/yang-mills-group/OpenAI-YM-20261028/blob/3456e5e75cf8d4cd558955a5c7638fba2d68f464/YM/Extra/WEAK_SCALING_ACTLAW/Orientation.lean\#L43-L54}{\texttt{YM.Extra.WeakScaling.tracks\_ratio}};
  \href{https://github.com/yang-mills-group/OpenAI-YM-20261028/blob/3456e5e75cf8d4cd558955a5c7638fba2d68f464/YM/Extra/WEAK_SCALING_ACTLAW/Criteria.lean\#L64-L70}{\texttt{YM.Extra.WeakScaling.flux\_nontrivial}}.
\item
  VIII.1 defines the Maxwell, infrared and ultraviolet data:
  \href{https://github.com/yang-mills-group/OpenAI-YM-20261028/blob/3456e5e75cf8d4cd558955a5c7638fba2d68f464/YM/Org/Input.lean\#L159-L170}{\texttt{YM.Org.MaxwellData}};
  \href{https://github.com/yang-mills-group/OpenAI-YM-20261028/blob/3456e5e75cf8d4cd558955a5c7638fba2d68f464/YM/Org/Input.lean\#L172-L175}{\texttt{YM.Org.IR}};
  \href{https://github.com/yang-mills-group/OpenAI-YM-20261028/blob/3456e5e75cf8d4cd558955a5c7638fba2d68f464/YM/Org/Input.lean\#L177-L186}{\texttt{YM.Org.UV}}.
\end{itemize}

\textbf{Standard mathematics.} Wick's theorem and the bosonic Fock inner
product (Glimm--Jaffe, ch.~6); reflection positivity of Gaussian fields
(Osterwalder--Schrader, 1973/75); the Euclidean Green function and
Cartan's formula; integration by parts and approximate identities;
homogeneous harmonic functions and the Kelvin transform;
finite-dimensional rank--nullity, openness of injectivity and
surjectivity, compactness of spheres, and uniqueness for linear ODEs.
The passage from uniform order-zero test estimates to smooth convergence
uses closure under gradients, localization and Fourier inversion
(VI.2.u).

\textbf{Imported constants and assumptions.} Simplicity of \(G\) enters
the stress classification and the rigidity arguments for the Maxwell
data. It is essential for the one-dimensional stress line: for
\(\mathfrak g=\mathfrak g_1\oplus\mathfrak g_2\), the two factors have
separately invariant conserved Maxwell stresses. The propagator
normalization is \(c_0=(4\pi^2)^{-1}\). The explicit correlation tests
give the integers \(-256,-512,128,-128\) in VI.2.l and the factor
\(256\dim\mathfrak g\,c_0^2\) in VI.2.n. The station estimates and
dimension bound enter as hypotheses of the transfer argument, with their
discharge stated in VI.2.7--8.

\subsubsection{3. Strategy}\label{strategy-24}

The free table is realized by invariant polynomials in on-shell
curvature coefficients with Wick correlations. Translation and rotation
covariance, analyticity and polynomial bounds follow from the
propagator. Reflection positivity uses a Fock-space factorization. For
two polynomials at mirror points, the reflected Gram matrix of their
finitely many coordinate probes is positive definite. Factoring it as
\(L^*L\) identifies their correlation with a Fock inner product after an
invertible linear substitution. Strict positivity eliminates further
null fields and identifies the free quotient with \(\mathcal P_d\).

For the Ward identity, replace the singular propagator by a smooth curl
source. Its divergence is a derivative of an approximate identity. Two
integrations by parts and removal of the regulator give the local charge
\(-\mathcal L_vf\). Differentiation in the source point extends the
charge to every Taylor order. Wick expansion then leaves three
possibilities: both stress legs meet the marked site, giving zero by
homogeneity; one meets that site and one a spectator, giving the desired
derivation; or both meet spectators, giving a source-free current inside
the cutoff.

The free stress line is computed in finite-dimensional spaces. Signed
covariance allows three centred symmetric tensors in degree four and at
most two vectors in degree five. Two explicit correlation tests prove
that divergence has rank two. Its kernel is therefore the line generated
by \(T^{\rm free}\). A translation Ward test fixes the constant
\(\kappa=-\tau\) for stress \(\tau T^{\rm free}\).

Stations give comparison maps at scales tending to zero. Correct them
recursively to intertwine inclusions. Finite separating tests make the
corrected maps eventually injective; the dimension bound makes them
isomorphisms. Integration by parts and the derivative law show that
\(r_kU_k^{-1}W_aU_k\to\partial^{\rm free}_a\). The correlation
differences, including those with convergent field coefficients, are
closed under gradients, so the uniform test estimates improve to smooth
local convergence.

Transport the interacting divergence to the free spaces, recentering
tensors after transport. It converges to free divergence, so
surjectivity and rank--nullity transfer the one-dimensional kernel to
\(S\). Normalize a nonzero vector in this kernel to a sphere in the free
stress space and pass to a subsequence. Its limit is a nonzero centred
conserved free tensor, hence a multiple of \(T^{\rm free}\); a single
sign choice fixes the limit on the entire comparison sequence.

Finally, a Ward-normalized stress \(J\) on \(S\) has centred part
proportional to the constructed stress:
\(c(J-\langle J\rangle\mathbf1)=J_0\). Comparing the derivative and
stress Ward identities along tracks gives \(c/\xi_k\to\kappa\). Thus
\(U_k^{-1}(r_k^4J)\to\kappa^{-1}T_0=-T^{\rm free}\) in the canonical
normalization. This obtains the interacting stress from the
finite-dimensional conservation law, then fixes its scale by the Ward
identity.

\subsubsection{4. Steps}\label{steps-24}

\paragraph{4.1 The free table}\label{the-free-table}

\textbf{Lemma VI.2.a (field spaces and correlations).} For every \(d\),
the null labels are exactly the kernel of substitution into
\(\mathcal P_d\), and \(\iota_d:\mathcal F^{\rm free}_d\to\mathcal P_d\)
is an isomorphism carrying correlations to Wick moments.

\emph{Sketch.} The free table is defined by substitution followed by the
Wick moment. A residual polynomial has zero correlations with every
family of shell-sector spectators. A nonzero polynomial has positive
two-point correlation with its own reflected conjugate by VI.2.b, which
remains in the same sector. Thus the residual subspace is zero.
\emph{Lean:}
\href{https://github.com/yang-mills-group/OpenAI-YM-20261028/blob/3456e5e75cf8d4cd558955a5c7638fba2d68f464/YM/Next/MAXWELL/System.lean\#L33-L35}{\texttt{fields}},
\href{https://github.com/yang-mills-group/OpenAI-YM-20261028/blob/3456e5e75cf8d4cd558955a5c7638fba2d68f464/YM/Next/MAXWELL/Caps.lean\#L90-L113}{\texttt{no\_residual}},
\href{https://github.com/yang-mills-group/OpenAI-YM-20261028/blob/3456e5e75cf8d4cd558955a5c7638fba2d68f464/YM/Euclid/Comparator.lean\#L23-L24}{\texttt{freeTable}},
\href{https://github.com/yang-mills-group/OpenAI-YM-20261028/blob/3456e5e75cf8d4cd558955a5c7638fba2d68f464/YM/Maxwell/Wick.lean\#L126-L127}{\texttt{Wick.moment}}.

\textbf{Lemma VI.2.b (strict reflection positivity).} For \(x_a>0\), the
pairing \(\langle(\vartheta_ap)(\theta_ax)q(x)\rangle_0\) is a
positive-definite Hermitian form on \(\mathcal W\).

\emph{Sketch.} Let \(s\) be the finite union of coordinate variables
appearing in \(p\) and \(q\). The reflected propagator Gram matrix
\(M_s\) is positive definite; factor it as \(M_s=L^*L\). The Wick
pairing between the two sites then equals
\(\langle T_Lp,T_Lq\rangle_{\rm Fock}\), where \(T_L\) substitutes
\(e_i\mapsto\sum_jL_{ji}e_j\). The Fock form makes monomials orthogonal
with squared norms \(\alpha!\). Since \(L\) is invertible, the form is
definite. Taking \(q=p\) gives strict positivity; testing a residual
polynomial against its reflected conjugate gives VI.2.6(d). \emph{Lean:}
\href{https://github.com/yang-mills-group/OpenAI-YM-20261028/blob/3456e5e75cf8d4cd558955a5c7638fba2d68f464/YM/Next/MAXWELL/Core.lean\#L94-L110}{\texttt{detects}},
\href{https://github.com/yang-mills-group/OpenAI-YM-20261028/blob/3456e5e75cf8d4cd558955a5c7638fba2d68f464/YM/Next/MAXWELL/Core.lean\#L44-L92}{\texttt{core}},
\href{https://github.com/yang-mills-group/OpenAI-YM-20261028/blob/3456e5e75cf8d4cd558955a5c7638fba2d68f464/YM/Next/MAXWELL/Fock.lean\#L188-L206}{\texttt{Fock.halves}},
\href{https://github.com/yang-mills-group/OpenAI-YM-20261028/blob/3456e5e75cf8d4cd558955a5c7638fba2d68f464/YM/Next/MAXWELL/Gram.lean\#L82-L105}{\texttt{coordGram}},
\href{https://github.com/yang-mills-group/OpenAI-YM-20261028/blob/3456e5e75cf8d4cd558955a5c7638fba2d68f464/YM/Next/MAXWELL/Statistic.lean\#L80-L99}{\texttt{ghosts\_out}}.
\emph{(\texttt{detects}, \texttt{core} and \texttt{ghosts\_out}: proof
read; reflected Gram positivity and Fock pairing factorization:
statement checked)}

\textbf{Lemma VI.2.c (the Euclidean hierarchy).} The free table is
reflection positive, real-analytic on distinct-point configurations,
translation invariant, and has polynomial bounds there, including bounds
for all derivatives.

\emph{Sketch.} For words with distinct sites in a positive half-space,
the multi-site Fock factorization makes the matrix of reflected
correlations positive semidefinite. A leg kernel between probes of
orders \(j,l\) is analytic and harmonic off zero, homogeneous of degree
\(-(4+j+l)\), and becomes analytic on all of \(\mathbb R^4\) after
multiplication by \(|x|^{2(3+j+l)}\). Its bound on a sphere and
homogeneity give a polynomial bound away from zero. Finite Wick sums
inherit analyticity and bounds. The derivative law gives bounds for all
derivatives; V.1 gives smeared positivity. The propagator depends only
on point differences, which gives translation invariance. \emph{Lean:}
\href{https://github.com/yang-mills-group/OpenAI-YM-20261028/blob/3456e5e75cf8d4cd558955a5c7638fba2d68f464/YM/Next/MAXWELL/Caps.lean\#L137-L155}{\texttt{positive\_table}},
\href{https://github.com/yang-mills-group/OpenAI-YM-20261028/blob/3456e5e75cf8d4cd558955a5c7638fba2d68f464/YM/Next/MAXWELL/Studio.lean\#L141-L158}{\texttt{cabinets\_positive}},
\href{https://github.com/yang-mills-group/OpenAI-YM-20261028/blob/3456e5e75cf8d4cd558955a5c7638fba2d68f464/YM/Next/MAXWELL/Pole.lean\#L238-L260}{\texttt{kernel\_pole}},
\href{https://github.com/yang-mills-group/OpenAI-YM-20261028/blob/3456e5e75cf8d4cd558955a5c7638fba2d68f464/YM/Next/MAXWELL/Caps.lean\#L69-L75}{\texttt{cap\_moment}},
\href{https://github.com/yang-mills-group/OpenAI-YM-20261028/blob/3456e5e75cf8d4cd558955a5c7638fba2d68f464/YM/Next/MAXWELL/System.lean\#L106-L114}{\texttt{hierarchy}}.
\emph{(proof read; multi-site Fock factorization used as an input)}

\textbf{Lemma VI.2.d (derivatives, rotations and conservation).} The
Taylor derivations and rotation action give VI.2.1(c), and
\(T^{\rm free}\) satisfies VI.2.1(d).

\emph{Sketch.} The directional derivative raises the Taylor order by one
and the weight by one. Applying it at one insertion differentiates
precisely the propagator legs incident to that insertion. Rotation
covariance follows from the natural jet action and invariance of the
propagator; VI.1 then gives the rotation laws.

To prove conservation as a polynomial identity, evaluate on an arbitrary
Taylor tower. Only its orders zero and one occur, so realize them by a
wave \(f=A_0+A_1\) with homogeneous components of degrees zero and one.
The evaluated divergence is
\(\frac12\sum_a\partial_a\Theta_{ab}(f,f)=\sum_i\langle(\mathrm{Div}\,f)_i,f_{bi}\rangle=0\).
Evaluation on all towers separates polynomials. Symmetry, reality,
tracelessness and centring follow from the displayed stress polynomial
and the absence of one-site contractions. \emph{Lean:}
\href{https://github.com/yang-mills-group/OpenAI-YM-20261028/blob/3456e5e75cf8d4cd558955a5c7638fba2d68f464/YM/Next/MAXWELL/System.lean\#L80-L83}{\texttt{scheme}},
\href{https://github.com/yang-mills-group/OpenAI-YM-20261028/blob/3456e5e75cf8d4cd558955a5c7638fba2d68f464/YM/Next/MAXWELL/System.lean\#L153-L170}{\texttt{gear}},
\href{https://github.com/yang-mills-group/OpenAI-YM-20261028/blob/3456e5e75cf8d4cd558955a5c7638fba2d68f464/YM/Next/MAXWELL/System.lean\#L176-L179}{\texttt{realStress}},
\href{https://github.com/yang-mills-group/OpenAI-YM-20261028/blob/3456e5e75cf8d4cd558955a5c7638fba2d68f464/YM/Next/MAXWELL/Stress.lean\#L76-L117}{\texttt{pol\_conservation}},
\href{https://github.com/yang-mills-group/OpenAI-YM-20261028/blob/3456e5e75cf8d4cd558955a5c7638fba2d68f464/YM/Next/MAXWELL/Stress.lean\#L24-L25}{\texttt{polStress}}.

\paragraph{4.2 The local Ward charge}\label{the-local-ward-charge}

\textbf{Lemma VI.2.e (divergence of the current).} For smooth closed
two-forms \(f,g\) on an open set \(U\) and a smooth conformal Killing
field \(v\),
\[\sum_a\partial_aj^v_a(f,g)=\sum_bv_b\sum_i\bigl(\langle(\mathrm{Div}\,f)_i,g_{bi}\rangle+\langle f_{bi},(\mathrm{Div}\,g)_i\rangle\bigr).\]
In particular, the flux vanishes when both fields are Maxwell waves near
the cutoff support.

\emph{Sketch.} Apply the product rule. The contraction of \(\nabla v\)
with \(\Theta\) vanishes because \(\Theta\) is symmetric and traceless
and the symmetric part of \(\nabla v\) is scalar. The remaining identity
follows by the Bianchi equations. Integrating against the cutoff proves
the flux statement. \emph{Lean:}
\href{https://github.com/yang-mills-group/OpenAI-YM-20261028/blob/3456e5e75cf8d4cd558955a5c7638fba2d68f464/YM/Next/MAXWELL/Current.lean\#L151-L181}{\texttt{current\_div}},
\href{https://github.com/yang-mills-group/OpenAI-YM-20261028/blob/3456e5e75cf8d4cd558955a5c7638fba2d68f464/YM/Next/MAXWELL/Current.lean\#L239-L250}{\texttt{sourcefree\_flux}}.

\textbf{Lemma VI.2.f (regulated source).} Under the hypotheses of
VI.2.2, fix \(i,l,c\) and \(t>0\). Set
\[G_t(y)=c_0(t^2+|y|^2)^{-1},\qquad \rho_t(y)=-\Delta G_t(y)=8c_0t^2(t^2+|y|^2)^{-3},\]
and let \(R^t=d\alpha^t\,e_c\), with
\(\alpha^t_b=\delta_{ib}\partial_lG_t(\cdot-x)-\delta_{lb}\partial_iG_t(\cdot-x)\).
Then \(R^t\) is closed and
\[(\mathrm{Div}\,R^t)_b=(\delta_{lb}\partial_i\rho_t-\delta_{ib}\partial_l\rho_t)(\cdot-x)e_c.\]
Writing \(h_a=(\iota_vf)^c_a\), its flux is
\[\Phi_\chi(j^v(R^t,f))=\int\rho_t(y-x)\bigl[\partial_l(\chi h_i)-\partial_i(\chi h_l)\bigr](https://github.com/scottnarmstrong/yang-mills-informalization/blob/main/outline/units/y)\,dy.\]

\emph{Sketch.} Apply VI.2.e and integrate \(\chi\,\mathrm{div}\,j^v\) by
parts. Insert the displayed divergence of \(R^t\) and integrate by parts
again to move derivatives from \(\rho_t\) to \(\chi h\). \emph{Lean:}
the regulated-flux identity in
\href{https://github.com/yang-mills-group/OpenAI-YM-20261028/blob/3456e5e75cf8d4cd558955a5c7638fba2d68f464/YM/Next/MAXWELL/Contact.lean\#L56-L98}{\texttt{Contact.lean}},
\href{https://github.com/yang-mills-group/OpenAI-YM-20261028/blob/3456e5e75cf8d4cd558955a5c7638fba2d68f464/YM/Next/MAXWELL/Regulator.lean\#L147-L148}{\texttt{regulated}}.

\textbf{Lemma VI.2.g (removing the regulator).} The regulated flux
extends continuously to \(t=0\) and gives VI.2.2(a).

\emph{Sketch.} On the compact support of \(\nabla\chi\), the source
point is excluded and \(R^t\) is jointly continuous down to zero, where
it equals \(f^{F^c_{il},x}\). The functions \(\rho_t\) form an
approximate identity. Hence VI.2.f tends to
\((\partial_lh_i-\partial_ih_l)(x)=-(\mathcal L_vf)^c_{il}(x)\), since
\(\chi=1\) near \(x\). \emph{Lean:}
\href{https://github.com/yang-mills-group/OpenAI-YM-20261028/blob/3456e5e75cf8d4cd558955a5c7638fba2d68f464/YM/Next/MAXWELL/Contact.lean\#L137-L156}{\texttt{local\_force}}.

\textbf{Lemma VI.2.h (charges of higher-order probes).} Under the
hypotheses of VI.2.2(b), the flux identity holds in every order.
Translations, rotations and dilation give its three stated
specializations.

\emph{Sketch.} Probes of order \(j+1\) are spanned by directional
derivatives of probes of order \(j\). Keep \(v\) and \(\chi\) fixed
while differentiating in the source point. The generated fields satisfy
\(\frac d{ds}f^{p,x+su}=f^{\partial_up,x+su}\), where \(\partial_up\) is
the probe representing the derivative of \(p(f)\). Compact support away
from the source justifies differentiation under the flux integral. The
right side differentiates as a probe evaluation because
\(\mathcal L_vf\) is a Maxwell wave.

For a centred affine motion, the order-\(j\) coefficient of
\(\mathcal L_vf\) at the centre depends only on the order-\(j\)
coefficient of \(f\). For rotations it is the derivative of the inverse
action on solution coefficients: if that action is \(Q_g\), then
\((\omega\cdot p)(A)=\frac d{dt}|_{t=0}p(Q_{\exp(-t\omega)}A)\). For
dilation, Euler's identity on the homogeneous wave gives the factor
\(j+2\). \emph{Lean:}
\href{https://github.com/yang-mills-group/OpenAI-YM-20261028/blob/3456e5e75cf8d4cd558955a5c7638fba2d68f464/YM/Next/MAXWELL/Force.lean\#L131-L186}{\texttt{jet\_force}},
\href{https://github.com/yang-mills-group/OpenAI-YM-20261028/blob/3456e5e75cf8d4cd558955a5c7638fba2d68f464/YM/Next/MAXWELL/Affine.lean\#L249-L255}{\texttt{drift\_flux}},
\href{https://github.com/yang-mills-group/OpenAI-YM-20261028/blob/3456e5e75cf8d4cd558955a5c7638fba2d68f464/YM/Next/MAXWELL/Angular.lean\#L181-L195}{\texttt{angular\_flux}},
\href{https://github.com/yang-mills-group/OpenAI-YM-20261028/blob/3456e5e75cf8d4cd558955a5c7638fba2d68f464/YM/Next/MAXWELL/Affine.lean\#L220-L247}{\texttt{radial\_flux}}.

\textbf{Lemma VI.2.i (two sources at one point).} With \(\chi,v\) as in
VI.2.2, probes \(p,q\) of orders \(j,l\), and \(v\) homogeneous about
\(x\) of degree \(e\in\{0,1\}\), the flux
\(Z=\Phi_\chi(j^v(f^{p,x},f^{q,x}))\) vanishes.

\emph{Sketch.} The current is divergence-free off \(x\), so changing the
cutoff while keeping it equal to one near \(x\) preserves its flux.
Replace it by \(\chi(x+(\cdot-x)/2)\). The two generated fields have
homogeneities \(-(4+j)\) and \(-(4+l)\), so a change of variables gives
\(Z=2^{3+e-(8+j+l)}Z\). The factor is strictly less than one.
\emph{Lean:}
\href{https://github.com/yang-mills-group/OpenAI-YM-20261028/blob/3456e5e75cf8d4cd558955a5c7638fba2d68f464/YM/Next/MAXWELL/Repeat.lean\#L63-L128}{\texttt{two\_flux\_zero}}.

\textbf{Lemma VI.2.j (Ward identities for Wick polynomials).} The three
charges of VI.2.h imply VI.2.3 for every tuple of Wick polynomials.

\emph{Sketch.} More generally, take a conformal Killing field \(v\) and
a polynomial derivation \(D\) for which the two-source flux at the
marked site vanishes and
\[\Phi_\chi(j^v(f^{p,x_0},f^{q,y}))=-\langle(Dp)(x_0)q(y)\rangle_0\qquad(y\notin\mathrm{supp}\,\chi).\]
The three motions satisfy these conditions by VI.2.h--i. For monomials,
expand the two stress legs in the Wick moment. The integrand is one half
the sum over ordered pairs of extracted legs of their polarized current
times the remaining Wick moment. Two marked legs give zero by VI.2.i;
two spectator legs give zero by VI.2.e. The mixed pairs give \(D\) on
each marked leg, with the two orders cancelling the factor one half. The
Leibniz rule reassembles the result as \(Dp_0\). Multilinearity gives
the identity for all polynomials, and \(\iota\) carries it to
field-space Ward laws. \emph{Lean:}
\href{https://github.com/yang-mills-group/OpenAI-YM-20261028/blob/3456e5e75cf8d4cd558955a5c7638fba2d68f464/YM/Next/MAXWELL/Taxes.lean\#L115-L172}{\texttt{fees}},
\href{https://github.com/yang-mills-group/OpenAI-YM-20261028/blob/3456e5e75cf8d4cd558955a5c7638fba2d68f464/YM/Next/MAXWELL/Taxes.lean\#L177-L200}{\texttt{moment\_force}},
\href{https://github.com/yang-mills-group/OpenAI-YM-20261028/blob/3456e5e75cf8d4cd558955a5c7638fba2d68f464/YM/Next/MAXWELL/Taxes.lean\#L24-L28}{\texttt{Motion}},
\href{https://github.com/yang-mills-group/OpenAI-YM-20261028/blob/3456e5e75cf8d4cd558955a5c7638fba2d68f464/YM/Next/MAXWELL/Beacon.lean\#L74-L136}{\texttt{beacon\_signal}}.

\paragraph{4.3 The free stress line and the Ward
constant}\label{the-free-stress-line-and-the-ward-constant}

\textbf{Lemma VI.2.k (the tensor and vector spaces).} Assume \(G\) is
simple. Put
\[Q_{ak,bj}=\sum_cF^c_{ak}F^c_{bj},\qquad P_{ab}=\sum_kQ_{ak,bk},\qquad N=\sum_{b,k}Q_{bk,bk},\]
\[A_{ab}=\delta_{ab}P_{aa},\qquad B_{ab}=\delta_{ab}(N-2P_{aa}),\qquad C_{ab}=(1-\delta_{ab})P_{ab}.\]
Then \(A,B,C\) form a real basis of \(\mathcal Z\), and
\(\dim_{\mathbb R}\mathcal V\le2\).

\emph{Sketch.} The simple-group invariant calculation through weight
four, followed by sign changes and coordinate permutations, leaves two
diagonal coefficients and one off-diagonal coefficient. Centring removes
the scalar unit term, and reality makes the coefficients real. The
tensors above have zero one-point functions. Their independence follows
from the tests in VI.2.l together with a normalized two-point test of
the \(01\) component against \(Q_{02,12}\), whose values on \((A,B,C)\)
are \((0,0,-16)\).

For the vector bound, the Bianchi and Maxwell equations reduce the
zeroth component to the real span of two fixed
curvature--first-derivative combinations. Signed permutations determine
every other component from it. Thus projection to the zeroth component
embeds \(\mathcal V\) into a space of dimension at most two.
\emph{Lean:}
\href{https://github.com/yang-mills-group/OpenAI-YM-20261028/blob/3456e5e75cf8d4cd558955a5c7638fba2d68f464/YM/UVNext/FREE_STRESS_CERTIFICATE/StressSpace.lean\#L115-L122}{\texttt{exhaust}},
\href{https://github.com/yang-mills-group/OpenAI-YM-20261028/blob/3456e5e75cf8d4cd558955a5c7638fba2d68f464/YM/UVNext/FREE_STRESS_CERTIFICATE/Certificate.lean\#L56-L71}{\texttt{mixed\_space}},
\href{https://github.com/yang-mills-group/OpenAI-YM-20261028/blob/3456e5e75cf8d4cd558955a5c7638fba2d68f464/YM/UVNext/FREE_STRESS_CERTIFICATE/VRank.lean\#L102-L116}{\texttt{capped\_v}}.
\emph{(proof read; invariant spanning calculations used as inputs)}

\textbf{Lemma VI.2.l (rank of free divergence).} Under the hypotheses of
VI.2.4, let \(s_1,s_2\) test the zeroth component of a vector at \(e_0\)
against \(Q_{01,01}(0)\) and \(Q_{12,12}(0)\), respectively, taking real
parts and dividing by \(\dim\mathfrak g\,c_0^2\). Then
\[(s_1,s_2)(\mathrm{div}A)=(-256,0),\qquad(s_1,s_2)(\mathrm{div}B)=(0,-512),\qquad(s_1,s_2)(\mathrm{div}C)=(128,-128).\]

\emph{Sketch.} Differentiate the explicit two-point curvature
correlations and evaluate the finite index sums. The first two columns
make \(s\circ\mathrm{div}\) onto \(\mathbb R^2\). By VI.2.k, \(s\) is
therefore an isomorphism and divergence is onto. Rank--nullity gives a
one-dimensional kernel. The columns also show that
\(\frac12A-\frac14B+C\) lies in it, and this tensor is
\(P-\frac14\delta N=T^{\rm free}\). \emph{Lean:}
\href{https://github.com/yang-mills-group/OpenAI-YM-20261028/blob/3456e5e75cf8d4cd558955a5c7638fba2d68f464/YM/UVNext/FREE_STRESS_CERTIFICATE/Pivots.lean\#L96-L99}{\texttt{east\_column}},
\href{https://github.com/yang-mills-group/OpenAI-YM-20261028/blob/3456e5e75cf8d4cd558955a5c7638fba2d68f464/YM/UVNext/FREE_STRESS_CERTIFICATE/Pivots.lean\#L100-L103}{\texttt{west\_column}},
\href{https://github.com/yang-mills-group/OpenAI-YM-20261028/blob/3456e5e75cf8d4cd558955a5c7638fba2d68f464/YM/UVNext/FREE_STRESS_CERTIFICATE/Pivots.lean\#L104-L107}{\texttt{cross\_column}},
\href{https://github.com/yang-mills-group/OpenAI-YM-20261028/blob/3456e5e75cf8d4cd558955a5c7638fba2d68f464/YM/UVNext/FREE_STRESS_CERTIFICATE/Certificate.lean\#L33-L41}{\texttt{scope\_exact}},
\href{https://github.com/yang-mills-group/OpenAI-YM-20261028/blob/3456e5e75cf8d4cd558955a5c7638fba2d68f464/YM/UVNext/FREE_STRESS_CERTIFICATE/Certificate.lean\#L43-L47}{\texttt{rank\_on}},
\href{https://github.com/yang-mills-group/OpenAI-YM-20261028/blob/3456e5e75cf8d4cd558955a5c7638fba2d68f464/YM/UVNext/FREE_STRESS_CERTIFICATE/Certificate.lean\#L73-L79}{\texttt{line\_dim}},
\href{https://github.com/yang-mills-group/OpenAI-YM-20261028/blob/3456e5e75cf8d4cd558955a5c7638fba2d68f464/YM/UVNext/FREE_STRESS_CERTIFICATE/Certificate.lean\#L86-L94}{\texttt{lighting}},
\href{https://github.com/yang-mills-group/OpenAI-YM-20261028/blob/3456e5e75cf8d4cd558955a5c7638fba2d68f464/YM/UVNext/FREE_STRESS_CERTIFICATE/Certificate.lean\#L125-L132}{\texttt{certificate}}.
\emph{(proof read; underlying differentiated propagator table used as an
input)}

\textbf{Lemma VI.2.m (canonical choice of stress).} For any Maxwell
datum for a simple \(G\) and any \(\tau\ne0\), one can retain its
derivative plan, rotations, hierarchy and Lorentz representation while
replacing its stress by \(\tau T^{\rm free}\) and obtaining (FD) and
full rotation covariance.

\emph{Sketch.} Divide the original stress by its Ward constant to obtain
Ward rows \((\partial,\rho)\). Average over \(O(4)\) with Haar measure
and take the compatible real part. The Ward identities imply symmetry
and conservation of the resulting tensor. Centre it by subtracting its
expectation times the unit; this preserves its Ward laws and covariance.
The centred tensor \(j_0\) is nonzero, since otherwise its translation
Ward law would make the derivative vanish in degree four. By VI.2.l,
\(j_0=cT^{\rm free}\) with \(c\ne0\). Scaling it by \(\tau/c\) gives the
desired stress with Ward constant \(\tau/c\); VI.2.n identifies this
constant as \(-\tau\). \emph{Lean:}
\href{https://github.com/yang-mills-group/OpenAI-YM-20261028/blob/3456e5e75cf8d4cd558955a5c7638fba2d68f464/YM/UVNext/FREE_STRESS_CERTIFICATE/Choose.lean\#L77-L134}{\texttt{ignite}},
\href{https://github.com/yang-mills-group/OpenAI-YM-20261028/blob/3456e5e75cf8d4cd558955a5c7638fba2d68f464/YM/UVNext/FREE_STRESS_CERTIFICATE/Choose.lean\#L41-L51}{\texttt{every\_light}},
\href{https://github.com/yang-mills-group/OpenAI-YM-20261028/blob/3456e5e75cf8d4cd558955a5c7638fba2d68f464/YM/UVNext/FREE_STRESS_CERTIFICATE/Choose.lean\#L137-L150}{\texttt{fire}}.
\emph{(proof read; Haar averaging and Ward-implied conservation and
symmetry: statement checked)}

\textbf{Lemma VI.2.n (the Ward constant).} If a Maxwell datum for a
simple \(G\) has stress \(\tau T^{\rm free}\), its Ward constant is
\(\kappa=-\tau\).

\emph{Sketch.} Test its translation Ward identity in direction zero with
\(Q_{01,01}\) at \(0\) and at \(-e_0\). Wick expansion leaves one stress
leg at each insertion. The single-contraction flux is
\(-\partial_0u_{01}(0)\) for a smooth Maxwell wave \(u\) on the cutoff
support. Evaluating the remaining propagators gives flux
\(256\tau\dim\mathfrak g\,c_0^2\). The derivative side is
\(-256\kappa\dim\mathfrak g\,c_0^2\). Since simplicity implies
\(\dim\mathfrak g>0\), the constants agree only if \(\kappa=-\tau\).

There is also a structural proof in the canonical case. If the plan is
\(\partial^{\rm free}\) and the stress is \(T^{\rm free}\), VI.2.3 gives
Ward constant \(-1\). Comparing translation Ward identities for two
constants, using separating correlations, makes their difference times
the derivative tuple in degree four zero. That tuple is nonzero, so the
constant is unique. \emph{Lean:}
\href{https://github.com/yang-mills-group/OpenAI-YM-20261028/blob/3456e5e75cf8d4cd558955a5c7638fba2d68f464/YM/UVNext/FREE_STRESS_CERTIFICATE/Pressure.lean\#L132-L151}{\texttt{negated}},
\href{https://github.com/yang-mills-group/OpenAI-YM-20261028/blob/3456e5e75cf8d4cd558955a5c7638fba2d68f464/YM/UVNext/FREE_STRESS_CERTIFICATE/Pressure.lean\#L158-L161}{\texttt{required}},
\href{https://github.com/yang-mills-group/OpenAI-YM-20261028/blob/3456e5e75cf8d4cd558955a5c7638fba2d68f464/YM/UVNext/FREE_STRESS_CERTIFICATE/Contact.lean\#L106-L124}{\texttt{single\_cross}},
\href{https://github.com/yang-mills-group/OpenAI-YM-20261028/blob/3456e5e75cf8d4cd558955a5c7638fba2d68f464/YM/UVNext/FREE_STRESS_CERTIFICATE/TripleMeter.lean\#L63-L80}{\texttt{battery}},
\href{https://github.com/yang-mills-group/OpenAI-YM-20261028/blob/3456e5e75cf8d4cd558955a5c7638fba2d68f464/YM/UVNext/FREE_STRESS_CERTIFICATE/Pressure.lean\#L96-L118}{\texttt{at\_tower}},
\href{https://github.com/yang-mills-group/OpenAI-YM-20261028/blob/3456e5e75cf8d4cd558955a5c7638fba2d68f464/YM/Pert/Lifted.lean\#L65-L74}{\texttt{kappa\_eq}},
\href{https://github.com/yang-mills-group/OpenAI-YM-20261028/blob/3456e5e75cf8d4cd558955a5c7638fba2d68f464/YM/Pert/Lifted.lean\#L57-L63}{\texttt{same\_translation}}.

\paragraph{4.4 The remaining Maxwell
data}\label{the-remaining-maxwell-data}

\textbf{Lemma VI.2.o (selection of the derivative).} For simple \(G\),
there is \(d_0\) such that spectators in \(\mathcal P_{d_0}\) detect
every shell-sector polynomial, of any weight. For every
\(d\ge\max(d_0,4)\), the admissible vector maps in VI.2.5(a) whose Ward
diagrams vanish form the line \(\mathbb R\partial^{\rm free}\). Finitely
many collinear diagrams suffice.

\emph{Sketch.} For detection, induct on the weight. Mixed contractions
lower the weight and constrain derivatives in the highest-order
variables. An explicit colour matrix has nonzero determinant as a matrix
of polynomials, so these constraints eliminate the high-order variables.
The remaining bounded Taylor orders are detected by a fixed finite
exterior degree.

For selection, the low-degree classification first gives a scalar
multiple of the derivative on degree four. Subtract that multiple. The
remainder kills the free stress; the dilation Ward identity and the
detection result make it commute with the weight derivation. Vector
covariance makes it anticommute with total coordinate inversion. Since
inversion acts on weight \(w\) by \((-1)^w\), the remainder is zero.
Analyticity transports vanishing of all derivatives at collinear
configurations to arbitrary configurations. Finite-dimensionality
reduces the intersection of all diagram kernels to a finite
intersection. \emph{Lean:}
\href{https://github.com/yang-mills-group/OpenAI-YM-20261028/blob/3456e5e75cf8d4cd558955a5c7638fba2d68f464/YM/Next/MAXWELL/Hidden.lean\#L113-L137}{\texttt{exposed}},
\href{https://github.com/yang-mills-group/OpenAI-YM-20261028/blob/3456e5e75cf8d4cd558955a5c7638fba2d68f464/YM/Next/MAXWELL/Hidden.lean\#L66-L87}{\texttt{matrix\_meal}},
\href{https://github.com/yang-mills-group/OpenAI-YM-20261028/blob/3456e5e75cf8d4cd558955a5c7638fba2d68f464/YM/Next/MAXWELL/Magnitude.lean\#L142-L181}{\texttt{charges\_quiet}},
\href{https://github.com/yang-mills-group/OpenAI-YM-20261028/blob/3456e5e75cf8d4cd558955a5c7638fba2d68f464/YM/Next/MAXWELL/Scenes.lean\#L92-L123}{\texttt{finite\_lenses}},
\href{https://github.com/yang-mills-group/OpenAI-YM-20261028/blob/3456e5e75cf8d4cd558955a5c7638fba2d68f464/YM/Next/MAXWELL/Scenes.lean\#L125-L132}{\texttt{scanner}},
\href{https://github.com/yang-mills-group/OpenAI-YM-20261028/blob/3456e5e75cf8d4cd558955a5c7638fba2d68f464/YM/Next/MAXWELL/Born.lean\#L51-L70}{\texttt{born}}.
\emph{(proof read; low-degree vector-map classification and
bounded-order detection used as inputs)}

\textbf{Lemma VI.2.p (the Ward-vanishing test).} The endomorphism family
in VI.2.5(b) vanishes for every \(d\ge4\).

\emph{Sketch.} Write the family on polynomials. Exact degree-four
two-point correlations imply that each endomorphism commutes with
coordinate quarter turns on the degree-four image. Equivariance then
makes its alternating bilinear dependence on \(u,v\in\mathfrak{so}(4)\)
invariant under simultaneous conjugation by those turns. The finite
quarter-turn calculation forces this alternating form to vanish. In
particular, the family kills \(T^{\rm free}\).

Apply the local rotational Ward identity to its vanishing Ward sums. The
commutator with every \(\rho(\omega)\) has zero correlations with
degree-\(d\) spectators and still lies in \(\mathcal P_d\); VI.2.6(d)
makes it zero. Uniqueness for the linear rotation ODE gives commutation
with quarter turns on all of \(\mathcal F^{\rm free}_d\). The same
alternating-form calculation now gives \(T=0\). \emph{Lean:}
\href{https://github.com/yang-mills-group/OpenAI-YM-20261028/blob/3456e5e75cf8d4cd558955a5c7638fba2d68f464/YM/Next/MAXWELL/AwayTeam.lean\#L48-L118}{\texttt{erasing}},
\href{https://github.com/yang-mills-group/OpenAI-YM-20261028/blob/3456e5e75cf8d4cd558955a5c7638fba2d68f464/YM/Next/MAXWELL/Statistic.lean\#L104-L150}{\texttt{commute\_ghost}},
\href{https://github.com/yang-mills-group/OpenAI-YM-20261028/blob/3456e5e75cf8d4cd558955a5c7638fba2d68f464/YM/Next/MAXWELL/AwayTeam.lean\#L29-L46}{\texttt{complete\_turn}},
\href{https://github.com/yang-mills-group/OpenAI-YM-20261028/blob/3456e5e75cf8d4cd558955a5c7638fba2d68f464/YM/Next/MAXWELL/Depth.lean\#L162-L181}{\texttt{low\_motion}}.

\textbf{Lemma VI.2.q (assembly of Maxwell data).} For simple \(G\) and
\(\tau\ne0\), the free lighthouse and the remaining data satisfy
VI.2.5(a),(b).

\emph{Sketch.} VI.2.d gives the derivative and rotation laws; VI.2.j
gives the Ward laws. The derivative is nonzero in every degree
\(d\ge4\): differentiate a nonconstant invariant quadratic curvature
polynomial and use faithfulness of the associated colour pairing. VI.2.o
supplies the selection property. The Lorentz action on the free
polynomial fields descends compatibly to their union, with boost
generator \(-i\rho\) and the prescribed spatial rotations. Combine these
results with the Casimir statements of VI.1 and the Ward-vanishing test
VI.2.p. \emph{Lean:}
\href{https://github.com/yang-mills-group/OpenAI-YM-20261028/blob/3456e5e75cf8d4cd558955a5c7638fba2d68f464/YM/Next/MAXWELL/Lit.lean\#L36-L42}{\texttt{shines}},
\href{https://github.com/yang-mills-group/OpenAI-YM-20261028/blob/3456e5e75cf8d4cd558955a5c7638fba2d68f464/YM/Next/MAXWELL/Lit.lean\#L45-L48}{\texttt{maxwell\_supplies}},
\href{https://github.com/yang-mills-group/OpenAI-YM-20261028/blob/3456e5e75cf8d4cd558955a5c7638fba2d68f464/YM/Bid/Scope.lean\#L24-L29}{\texttt{Bid.masts}},
\href{https://github.com/yang-mills-group/OpenAI-YM-20261028/blob/3456e5e75cf8d4cd558955a5c7638fba2d68f464/YM/Next/MAXWELL/Descent.lean\#L173-L177}{\texttt{voyage}}.
\emph{(assembly and nonzero derivative: proof read; Lorentz and Casimir
data: statement checked)}

\paragraph{4.5 Comparison with the Wilson
limit}\label{comparison-with-the-wilson-limit}

\textbf{Lemma VI.2.r (panels from a station).} An admissible station
under the hypotheses of VI.2.8 yields the panel and estimate of
VI.2.8(a).

\emph{Sketch.} The bounded lift maps lie in a fixed finite-dimensional
space. Pass to a convergent subfilter and project their limits to the
quotient fields of \(S\). Reality, signed covariance and the degree-zero
unit identity pass to the limit. Ordered smeared convergence also holds
for convergent label coefficients, by finite-dimensional multilinearity.
The change of variables in the dilated test then turns the admissibility
estimate into the panel estimate. The subfilter proves an estimate for
the same continuum table; it does not change the physical extraction
order. \emph{Lean:}
\href{https://github.com/yang-mills-group/OpenAI-YM-20261028/blob/3456e5e75cf8d4cd558955a5c7638fba2d68f464/YM/Next/UV_STRESS/Station.lean\#L132-L180}{\texttt{harvested}},
\href{https://github.com/yang-mills-group/OpenAI-YM-20261028/blob/3456e5e75cf8d4cd558955a5c7638fba2d68f464/YM/Next/UV_STRESS/Vary.lean\#L62-L71}{\texttt{calibrated\_smear}}.

\textbf{Lemma VI.2.s (balanced estimates give stations).} Fix the
extraction hypotheses of VI.2.8 and finite bounds \(N,D\ge0\) and
\(R\ge1\). Suppose stations satisfy \(r_k\to0\) and the eventual lift
bounds and symmetries of VI.2.8. Let their normalized grid spacings be
\(s_k>0\) and their strip widths \(\sigma_k>0\). For each basis word
\(w\) of arity \(n\le N\) and degrees at most \(D\), suppose there are
errors \(e_k\to0\) and eventually nonnegative costs \(A_k,B_k\) with
\[\sigma_k\to0,\qquad4ns_k\le\sigma_k\ \text{eventually},\qquad A_k^n\sigma_k\to0,\qquad B_k^ns_k/\sigma_k\to0.\]
For every sufficiently large fixed \(k\) and every convergent subfilter
of the lifts along the extraction order, let \(Q_k\) be the rescaled
continuum correlation of their projected limits. Write \(Q_0\) for the
free correlation of \(w\) and
\[E_\sigma=\{y:\ |y_i^a-y_j^a|>\sigma\ \text{for every }i\ne j,a\}.\]
Assume the following bounds for every such limit:
\[|Q_k|\le A_k^n\quad\text{on }K_{R,n},\qquad \|DQ_k\|\le4nB_k^n/\sigma_k\quad\text{on }K_{2R,n}\cap E_{4\sigma_k},\]
\[|Q_k-Q_0|\le e_k\quad\text{on }(s_k\mathbb Z)^{4n}\cap K_{2R,n}\cap E_0.\]
These bounds for all \(N,D,R\) imply the station property. The grid,
bulk and finite-difference assumptions of the balanced estimate in
III.11 imply precisely these bounds for every cluster limit.

\emph{Sketch.} The excluded band \(K_{R,n}\setminus E_{8\sigma_k}\) has
volume \(O_{n,R}(\sigma_k)\), so its contribution is
\(O_{n,R}((A_k^n+\|Q_0\|_\infty)\sigma_k)\). Outside the band, rounding
to the grid stays in \(E_{4\sigma_k}\) and \(K_{2R,n}\) for large \(k\).
Interpolation costs \(O_n(B_k^ns_k/\sigma_k)\) plus the free
correlation's modulus of continuity. All these terms and \(e_k\) tend to
zero. A contradiction subfilter with convergent lifts converts the bound
into the ordered cutoff estimate. Expansion in the finite product basis,
with coefficient sum bounded by a constant times \(\prod_i\|v_i\|\),
gives admissibility for arbitrary tuples. Costs may depend on the basis
word; only finitely many words occur for the fixed bounds. \emph{Lean:}
\href{https://github.com/yang-mills-group/OpenAI-YM-20261028/blob/3456e5e75cf8d4cd558955a5c7638fba2d68f464/YM/Next/UV_STRESS/Glide.lean\#L102-L129}{\texttt{cutoff\_balance}},
\href{https://github.com/yang-mills-group/OpenAI-YM-20261028/blob/3456e5e75cf8d4cd558955a5c7638fba2d68f464/YM/Next/UV_STRESS/Glide.lean\#L72-L99}{\texttt{pays\_film}},
\href{https://github.com/yang-mills-group/OpenAI-YM-20261028/blob/3456e5e75cf8d4cd558955a5c7638fba2d68f464/YM/Next/UV_STRESS/Glide.lean\#L24-L31}{\texttt{Balances}},
\href{https://github.com/yang-mills-group/OpenAI-YM-20261028/blob/3456e5e75cf8d4cd558955a5c7638fba2d68f464/YM/Next/UV_STRESS/Fringe.lean\#L38-L42}{\texttt{Wheel.Spots}},
\href{https://github.com/yang-mills-group/OpenAI-YM-20261028/blob/3456e5e75cf8d4cd558955a5c7638fba2d68f464/YM/Next/UV_STRESS/Granular.lean\#L121-L128}{\texttt{supply\_wheel}},
\href{https://github.com/yang-mills-group/OpenAI-YM-20261028/blob/3456e5e75cf8d4cd558955a5c7638fba2d68f464/YM/Next/UV_STRESS/Granular.lean\#L134-L140}{\texttt{realizes\_stations}}.

\textbf{Lemma VI.2.t (comparison isomorphisms).} Under the hypotheses of
VI.2.7, there is an approaching stable scaling, initially with
\(\xi_k=1\).

\emph{Sketch.} Choose stations of depth and arity bound \(k\), radius
bound \(k+1\), tolerance \((k+1)^{-1}\), and strictly decreasing radii
\(r_k\le(k+1)^{-1}\). Their panels approach in the test sense of
VI.2.7(a). For each consecutive free inclusion choose a real,
signed-covariant left inverse \(R_d\). Starting with
\(Q_{k,0}=U_{k,0}\), set
\[Q_{k,d+1}=\mathrm{inc}^S Q_{k,d}R_d+U_{k,d+1}(1-\mathrm{inc}^{\rm free}R_d).\]
These maps intertwine inclusions exactly. Expansion into mixed tuples of
lower-degree fields and complementary fields preserves the comparison
estimates. Finitely many separating joint tests make \(Q_{k,d}\)
eventually injective. The dimension bound then makes it bijective. In
each degree, fill the finitely many earlier indices with any one of
these isomorphisms. This preserves eventual comparison, symmetry and
inclusion identities and yields isomorphisms for every \(k,d\).
\emph{Lean:}
\href{https://github.com/yang-mills-group/OpenAI-YM-20261028/blob/3456e5e75cf8d4cd558955a5c7638fba2d68f464/YM/Next/UV_STRESS/Rates.lean\#L27-L41}{\texttt{Stations.ranging}},
\href{https://github.com/yang-mills-group/OpenAI-YM-20261028/blob/3456e5e75cf8d4cd558955a5c7638fba2d68f464/YM/Next/UV_STRESS/Rates.lean\#L55-L73}{\texttt{Stations.build\_curve}},
\href{https://github.com/yang-mills-group/OpenAI-YM-20261028/blob/3456e5e75cf8d4cd558955a5c7638fba2d68f464/YM/Next/UV_STRESS/Arrange.lean\#L29-L72}{\texttt{signed\_section}},
\href{https://github.com/yang-mills-group/OpenAI-YM-20261028/blob/3456e5e75cf8d4cd558955a5c7638fba2d68f464/YM/Next/UV_STRESS/Arrange.lean\#L78-L81}{\texttt{corrected}},
\href{https://github.com/yang-mills-group/OpenAI-YM-20261028/blob/3456e5e75cf8d4cd558955a5c7638fba2d68f464/YM/Next/UV_STRESS/Invert.lean\#L88-L103}{\texttt{MeterTest.injects}},
\href{https://github.com/yang-mills-group/OpenAI-YM-20261028/blob/3456e5e75cf8d4cd558955a5c7638fba2d68f464/YM/Next/UV_STRESS/Invert.lean\#L132-L181}{\texttt{Stations.equivalences}},
\href{https://github.com/yang-mills-group/OpenAI-YM-20261028/blob/3456e5e75cf8d4cd558955a5c7638fba2d68f464/YM/Next/UV_STRESS/Invert.lean\#L108}{\texttt{DimAbove}}.

\textbf{Lemma VI.2.u (smooth convergence).} With the hierarchy and
derivative laws of VI.2.7, every approaching scaling satisfies
\[r_kU_{k,d+1}^{-1}W_aU_{k,d}\longrightarrow\partial^{\rm free}_a\] in
each finite-dimensional operator space. Its rescaled correlations
converge with all derivatives uniformly on compact subsets of
\(\mathrm{Conf}_n\). This also holds for free field tuples \(v_k\to v\)
with fixed degrees.

\emph{Sketch.} The chain rule identifies a position derivative of a
rescaled correlation with the correlation obtained by inserting the
displayed transported derivative. Integrate against tests and move the
derivative to the test. Approaching convergence and finitely many
separating tests then determine the limit of the operator.

Finite-dimensional multilinearity extends the test estimates to
convergent field tuples. The resulting correlation differences are
closed under gradients, by operator convergence. After localization,
their Fourier transforms tend uniformly to zero: multiplication of a
fixed test by any Fourier phase preserves its sup norm. Repeated
integration by parts gives arbitrarily high decay in frequency, using
the same estimates for gradient descendants. Fourier inversion gives
uniform convergence; repeating this for descendants gives every
derivative order. \emph{Lean:}
\href{https://github.com/yang-mills-group/OpenAI-YM-20261028/blob/3456e5e75cf8d4cd558955a5c7638fba2d68f464/YM/Next/UV_STRESS/Focus.lean\#L22-L48}{\texttt{scaled\_identity}},
\href{https://github.com/yang-mills-group/OpenAI-YM-20261028/blob/3456e5e75cf8d4cd558955a5c7638fba2d68f464/YM/Next/UV_STRESS/Focus.lean\#L67-L92}{\texttt{bend\_limit}},
\href{https://github.com/yang-mills-group/OpenAI-YM-20261028/blob/3456e5e75cf8d4cd558955a5c7638fba2d68f464/YM/Next/UV_STRESS/Focus.lean\#L158-L181}{\texttt{closed\_path}},
\href{https://github.com/yang-mills-group/OpenAI-YM-20261028/blob/3456e5e75cf8d4cd558955a5c7638fba2d68f464/YM/Next/UV_STRESS/Focus.lean\#L183-L196}{\texttt{travels\_jet}},
\href{https://github.com/yang-mills-group/OpenAI-YM-20261028/blob/3456e5e75cf8d4cd558955a5c7638fba2d68f464/YM/Next/UV_STRESS/Local.lean\#L152-L181}{\texttt{GradientFamily.jets}}.

\textbf{Lemma VI.2.v (rank transfer).} Every approaching stable scaling
under VI.2.7 transfers surjectivity and the one-dimensional divergence
kernel to \(S\).

\emph{Sketch.} For large \(k\), transport a free centred tensor by
\(U_{k,4}\) and then recenter it in \(S\); this defines a real
isomorphism \(C_k:\mathcal Z^{\rm free}\to\mathcal Z^S\). The inverse
transports back and recentres; preservation of the unit makes the two
composites identities. Let \(V_k:\mathcal V^{\rm free}\to\mathcal V^S\)
be transport by \(U_{k,5}\). Then
\[Z_k=V_k^{-1}\circ(r_k\,\mathrm{div}^S)\circ C_k\longrightarrow\mathrm{div}^{\rm free}.\]
The centring term disappears on differentiation because the derivative
kills the unit. Compose \(Z_k\) with a right inverse of free divergence.
The result tends to the identity on \(\mathcal V^{\rm free}\), hence is
eventually invertible. Thus \(\mathrm{div}^S\) is onto. The two
isomorphisms identify the dimensions of domain and target with the free
dimensions, so rank--nullity gives kernel dimension one. \emph{Lean:}
\href{https://github.com/yang-mills-group/OpenAI-YM-20261028/blob/3456e5e75cf8d4cd558955a5c7638fba2d68f464/YM/Next/UV_STRESS/Dimension.lean\#L165-L195}{\texttt{differential\_rank}},
\href{https://github.com/yang-mills-group/OpenAI-YM-20261028/blob/3456e5e75cf8d4cd558955a5c7638fba2d68f464/YM/Next/UV_STRESS/Dimension.lean\#L77-L82}{\texttt{centerEquiv}}.

\textbf{Lemma VI.2.w (one current along one comparison sequence).} Under
VI.2.7, there are a nonzero centred conserved \(J\) and a refinement
\(Y\) satisfying VI.2.7(c),(d).

\emph{Sketch.} Choose \(B\ne0\) in \(\ker\mathrm{div}^S\) and normalize
its entire transported stress tensor:
\[\xi_k=\frac{\|U_{k,4}^{-1}(r_k^4B)\|}{\|T_0\|},\qquad \beta_k=\xi_k^{-1}U_{k,4}^{-1}(r_k^4B).\]
A subsequence converges on the sphere of radius \(\|T_0\|\). The limit
is symmetric, real and signed-covariant by closedness. Convergence of
one-point functions makes it centred. VI.2.u and the equation
\(\sum_a(r_kU_{k,5}^{-1}W_aU_{k,4})\beta_{k,ab}=0\) make it conserved.
By (FD), it equals \(cT_0\) for a real \(c\) with \(|c|=1\). Choose
\(J=B\) or \(-B\) once so the limit is \(T_0\). For each finite set of
track requirements, stability holds on a tail and VI.2.u supplies the
required derivatives. This tail is a track for the same \(J\).
\emph{Lean:}
\href{https://github.com/yang-mills-group/OpenAI-YM-20261028/blob/3456e5e75cf8d4cd558955a5c7638fba2d68f464/YM/Next/UV_STRESS/Limit.lean\#L129-L167}{\texttt{produce\_current}},
\href{https://github.com/yang-mills-group/OpenAI-YM-20261028/blob/3456e5e75cf8d4cd558955a5c7638fba2d68f464/YM/Next/UV_STRESS/Limit.lean\#L53-L68}{\texttt{subsector}},
\href{https://github.com/yang-mills-group/OpenAI-YM-20261028/blob/3456e5e75cf8d4cd558955a5c7638fba2d68f464/YM/Next/UV_STRESS/Limit.lean\#L106-L124}{\texttt{on\_axis}},
\href{https://github.com/yang-mills-group/OpenAI-YM-20261028/blob/3456e5e75cf8d4cd558955a5c7638fba2d68f464/YM/Next/UV_STRESS/Limit.lean\#L88-L104}{\texttt{center\_expect}},
\href{https://github.com/yang-mills-group/OpenAI-YM-20261028/blob/3456e5e75cf8d4cd558955a5c7638fba2d68f464/YM/Next/UV_STRESS/Limit.lean\#L173-L194}{\texttt{focused\_track}},
\href{https://github.com/yang-mills-group/OpenAI-YM-20261028/blob/3456e5e75cf8d4cd558955a5c7638fba2d68f464/YM/Next/UV_STRESS/Limit.lean\#L197-L202}{\texttt{focused}}.

\textbf{Lemma VI.2.x (physical normalization).} Under VI.2.9, there is
\(c\ne0\) such that
\[c(J-\langle J\rangle\mathbf1)=J_0,\qquad c/\xi_k\longrightarrow\kappa.\]
Consequently VI.2.9(a)--(c) hold and \(\xi_k\to c/\kappa\ne0\).

\emph{Sketch.} The centred stress \(J_c\) lies in
\(\ker\mathrm{div}^S\). It is nonzero: otherwise the translation Ward
identity would make the derivative tuple in degree four zero. Its
transported limit would then contradict the nonzero free derivative. The
kernel is a real line, so \(cJ_c=J_0\) for some \(c\ne0\). Comparing
their translation Ward identities gives \(p^I=cW\) in degree four.

Along tracks, the transported derivatives tend to
\(\partial^{\rm free}\), whereas the transported Ward rows of \(J_0\),
divided by \(\xi_k\), tend to \(\kappa\partial^{\rm free}\). Since the
free derivative is nonzero, a separating linear functional gives
\(c/\xi_k\to\kappa\). Therefore
\[U_{k,4}^{-1}(r_k^4J_c)=\frac{\xi_k}{c}\bigl[\xi_k^{-1}U_{k,4}^{-1}(r_k^4J_0)\bigr]\longrightarrow\kappa^{-1}T_0.\]
The unit part tends to zero because \(r_k^4\to0\) and the maps
eventually preserve the unit. VI.2.u applies to these convergent
preimages in every combination of slots, giving the mixed-correlation
limit. VI.2.n and the equality of the two displayed free-stress
constructions give the canonical limit \(-T^{\rm free}\). \emph{Lean:}
\href{https://github.com/yang-mills-group/OpenAI-YM-20261028/blob/3456e5e75cf8d4cd558955a5c7638fba2d68f464/YM/Pert/Weak.lean\#L82-L116}{\texttt{centered\_line}},
\href{https://github.com/yang-mills-group/OpenAI-YM-20261028/blob/3456e5e75cf8d4cd558955a5c7638fba2d68f464/YM/Pert/Weak.lean\#L62-L66}{\texttt{ratio}},
\href{https://github.com/yang-mills-group/OpenAI-YM-20261028/blob/3456e5e75cf8d4cd558955a5c7638fba2d68f464/YM/Pert/Weak.lean\#L118-L144}{\texttt{meter\_physical}},
\href{https://github.com/yang-mills-group/OpenAI-YM-20261028/blob/3456e5e75cf8d4cd558955a5c7638fba2d68f464/YM/Pert/Lifted.lean\#L85-L124}{\texttt{mixed\_leading}},
\href{https://github.com/yang-mills-group/OpenAI-YM-20261028/blob/3456e5e75cf8d4cd558955a5c7638fba2d68f464/YM/Extra/WEAK_SCALING_ACTLAW/Orientation.lean\#L43-L54}{\texttt{tracks\_ratio}}.

\subsubsection{5. Where the difficulty
lies}\label{where-the-difficulty-lies-24}

The transfer to \(S\) depends on exact field-space isomorphisms, not
merely comparison maps. Finite separating tests and the upper dimension
bound provide them. The derivative law then turns the uniform test
estimates into convergence of transported derivatives and smooth
correlations. Recentring is necessary before transferring divergence: a
field-space isomorphism need not preserve the expectation functional.

The normalization is a separate assertion. Conservation determines a
line, while the Ward identity determines a point on it. Comparing that
identity along tracks removes the initially unknown factor \(\xi_k\) and
fixes the negative free-stress limit.

Two algebraic rigidity results support VI.1: detection of
arbitrary-weight fields by bounded-degree spectators, and vanishing of
alternating internal Ward rows. They require more than reflection
positivity alone. The three-dimensional tensor space also reflects the
chosen symmetry group: signed coordinate permutations permit three
coefficients, and conservation reduces them to one. The local current
calculation admits conformal Killing fields, although the polynomial
Ward argument uses only translations, rotations and dilation.

\subsubsection{6. Notation}\label{notation-24}

{\def\LTcaptype{none} % do not increment counter
\begin{longtable}[]{@{}
  >{\raggedright\arraybackslash}p{(\linewidth - 4\tabcolsep) * \real{0.3333}}
  >{\raggedright\arraybackslash}p{(\linewidth - 4\tabcolsep) * \real{0.3333}}
  >{\raggedright\arraybackslash}p{(\linewidth - 4\tabcolsep) * \real{0.3333}}@{}}
\toprule\noalign{}
\begin{minipage}[b]{\linewidth}\raggedright
object
\end{minipage} & \begin{minipage}[b]{\linewidth}\raggedright
code name
\end{minipage} & \begin{minipage}[b]{\linewidth}\raggedright
proposed notation
\end{minipage} \\
\midrule\noalign{}
\endhead
\bottomrule\noalign{}
\endlastfoot
free Maxwell table &
\href{https://github.com/yang-mills-group/OpenAI-YM-20261028/blob/3456e5e75cf8d4cd558955a5c7638fba2d68f464/YM/Euclid/Comparator.lean\#L23-L24}{\texttt{Maxwell.freeTable}},
\href{https://github.com/yang-mills-group/OpenAI-YM-20261028/blob/3456e5e75cf8d4cd558955a5c7638fba2d68f464/YM/Climb/Cells.lean\#L20}{\texttt{Climb.Bel}},
\texttt{Maxwell.Free.Model} & \(\mathcal M\), fields
\(\mathcal F^{\rm free}_d\) \\
homogeneous on-shell Taylor coefficients, polynomial algebra &
\href{https://github.com/yang-mills-group/OpenAI-YM-20261028/blob/3456e5e75cf8d4cd558955a5c7638fba2d68f464/YM/Maxwell/Shell.lean\#L44}{\texttt{Shell}},
\href{https://github.com/yang-mills-group/OpenAI-YM-20261028/blob/3456e5e75cf8d4cd558955a5c7638fba2d68f464/YM/Maxwell/Equation.lean\#L24}{\texttt{WickPol}}
& \(\mathrm{Sh}_j\), \(\mathcal W\) \\
shell sector &
\href{https://github.com/yang-mills-group/OpenAI-YM-20261028/blob/3456e5e75cf8d4cd558955a5c7638fba2d68f464/YM/Euclid/Comparator.lean\#L64}{\texttt{shellSector}}
& \(\mathcal P_d\) \\
field ↔ polynomial &
\href{https://github.com/yang-mills-group/OpenAI-YM-20261028/blob/3456e5e75cf8d4cd558955a5c7638fba2d68f464/YM/Next/MAXWELL/System.lean\#L33-L35}{\texttt{fields}},
\href{https://github.com/yang-mills-group/OpenAI-YM-20261028/blob/3456e5e75cf8d4cd558955a5c7638fba2d68f464/YM/Next/MAXWELL/System.lean\#L39}{\texttt{address}}
& \(\iota_d\) \\
Wick moment &
\href{https://github.com/yang-mills-group/OpenAI-YM-20261028/blob/3456e5e75cf8d4cd558955a5c7638fba2d68f464/YM/Maxwell/Wick.lean\#L126-L127}{\texttt{Wick.moment}}
& \(\langle p_1(x_1)\cdots p_n(x_n)\rangle_0\) \\
probe evaluation, generated field, leg kernel &
\href{https://github.com/yang-mills-group/OpenAI-YM-20261028/blob/3456e5e75cf8d4cd558955a5c7638fba2d68f464/YM/Next/MAXWELL/Kernel.lean\#L22-L23}{\texttt{sample}},
\href{https://github.com/yang-mills-group/OpenAI-YM-20261028/blob/3456e5e75cf8d4cd558955a5c7638fba2d68f464/YM/Next/MAXWELL/Kernel.lean\#L101-L102}{\texttt{pulse}},
\href{https://github.com/yang-mills-group/OpenAI-YM-20261028/blob/3456e5e75cf8d4cd558955a5c7638fba2d68f464/YM/Next/MAXWELL/Kernel.lean\#L123}{\texttt{K}}
& \(p(f)(x)\), \(f^{p,x}\), \(\langle p(x)q(y)\rangle_0\) \\
derivative plan, rotations &
\href{https://github.com/yang-mills-group/OpenAI-YM-20261028/blob/3456e5e75cf8d4cd558955a5c7638fba2d68f464/YM/Next/MAXWELL/System.lean\#L80-L83}{\texttt{scheme}},
\href{https://github.com/yang-mills-group/OpenAI-YM-20261028/blob/3456e5e75cf8d4cd558955a5c7638fba2d68f464/YM/Next/MAXWELL/System.lean\#L153-L170}{\texttt{gear}},
\href{https://github.com/yang-mills-group/OpenAI-YM-20261028/blob/3456e5e75cf8d4cd558955a5c7638fba2d68f464/YM/Next/MAXWELL/Derives.lean\#L27-L28}{\texttt{rewired}}
(\texttt{glide}, \texttt{wheel}, \texttt{metric}) &
\(\partial^{\rm free}_a\); \(R_g\), \(\rho(\omega)\); \(\varepsilon\) \\
free Maxwell stress &
\href{https://github.com/yang-mills-group/OpenAI-YM-20261028/blob/3456e5e75cf8d4cd558955a5c7638fba2d68f464/YM/Next/MAXWELL/Stress.lean\#L24-L25}{\texttt{polStress}},
\href{https://github.com/yang-mills-group/OpenAI-YM-20261028/blob/3456e5e75cf8d4cd558955a5c7638fba2d68f464/YM/Next/MAXWELL/System.lean\#L172}{\texttt{original}},
\href{https://github.com/yang-mills-group/OpenAI-YM-20261028/blob/3456e5e75cf8d4cd558955a5c7638fba2d68f464/YM/UVNext/FREE_STRESS_CERTIFICATE/Certificate.lean\#L82}{\texttt{canonicalStress}},
\href{https://github.com/yang-mills-group/OpenAI-YM-20261028/blob/3456e5e75cf8d4cd558955a5c7638fba2d68f464/YM/UVNext/FREE_STRESS_CERTIFICATE/Certificate.lean\#L111}{\texttt{torch}}
& \(T^{\rm free}\), \(\tau T^{\rm free}\) \\
polarized stress, current, flux &
\href{https://github.com/yang-mills-group/OpenAI-YM-20261028/blob/3456e5e75cf8d4cd558955a5c7638fba2d68f464/YM/Next/MAXWELL/Current.lean\#L31-L33}{\texttt{cross}},
\href{https://github.com/yang-mills-group/OpenAI-YM-20261028/blob/3456e5e75cf8d4cd558955a5c7638fba2d68f464/YM/Next/MAXWELL/Current.lean\#L145}{\texttt{current}},
\href{https://github.com/yang-mills-group/OpenAI-YM-20261028/blob/3456e5e75cf8d4cd558955a5c7638fba2d68f464/YM/Next/MAXWELL/Current.lean\#L184-L185}{\texttt{exportFlux}}
& \(\Theta_{ab}\), \(j^v\), \(\Phi_\chi\) \\
conformal Killing field, Lie derivative &
\href{https://github.com/yang-mills-group/OpenAI-YM-20261028/blob/3456e5e75cf8d4cd558955a5c7638fba2d68f464/YM/Next/MAXWELL/Current.lean\#L140-L143}{\texttt{Killing}},
\href{https://github.com/yang-mills-group/OpenAI-YM-20261028/blob/3456e5e75cf8d4cd558955a5c7638fba2d68f464/YM/Next/MAXWELL/Regulator.lean\#L49}{\texttt{variation}}
& \(v\), \(\mathcal L_vf=d\iota_vf\) \\
centred stresses, vectors, divergence &
\href{https://github.com/yang-mills-group/OpenAI-YM-20261028/blob/3456e5e75cf8d4cd558955a5c7638fba2d68f464/YM/Next/UV_STRESS/Tensor.lean\#L75-L80}{\texttt{KernelSpace}},
\href{https://github.com/yang-mills-group/OpenAI-YM-20261028/blob/3456e5e75cf8d4cd558955a5c7638fba2d68f464/YM/Next/UV_STRESS/Tensor.lean\#L87-L97}{\texttt{Vectors}},
\href{https://github.com/yang-mills-group/OpenAI-YM-20261028/blob/3456e5e75cf8d4cd558955a5c7638fba2d68f464/YM/Next/UV_STRESS/Tensor.lean\#L117-L119}{\texttt{currentDiv}}
& \(\mathcal Z\), \(\mathcal V\), \(\mathrm{div}\) \\
free divergence property &
\href{https://github.com/yang-mills-group/OpenAI-YM-20261028/blob/3456e5e75cf8d4cd558955a5c7638fba2d68f464/YM/Next/UV_STRESS/Dimension.lean\#L153-L158}{\texttt{FreeDiv}}
& (FD) \\
lighthouse, Maxwell data &
\href{https://github.com/yang-mills-group/OpenAI-YM-20261028/blob/3456e5e75cf8d4cd558955a5c7638fba2d68f464/YM/Climb/Beams.lean\#L24-L29}{\texttt{Climb.Lighthouse}},
\href{https://github.com/yang-mills-group/OpenAI-YM-20261028/blob/3456e5e75cf8d4cd558955a5c7638fba2d68f464/YM/Next/MAXWELL/Lit.lean\#L29-L34}{\texttt{herald}},
\href{https://github.com/yang-mills-group/OpenAI-YM-20261028/blob/3456e5e75cf8d4cd558955a5c7638fba2d68f464/YM/Org/Input.lean\#L159-L170}{\texttt{Org.MaxwellData}}
& \(\mathcal T=(\partial^0,R^0,T_0,\kappa)\), \(M\) \\
null-test space, Ward vanishing &
\href{https://github.com/yang-mills-group/OpenAI-YM-20261028/blob/3456e5e75cf8d4cd558955a5c7638fba2d68f464/YM/Next/MAXWELL/Statistic.lean\#L65-L76}{\texttt{ghost}},
\href{https://github.com/yang-mills-group/OpenAI-YM-20261028/blob/3456e5e75cf8d4cd558955a5c7638fba2d68f464/YM/Next/MAXWELL/Statistic.lean\#L40-L42}{\texttt{silence}}
& \(\mathcal N_d\); \(\sum_i\langle\cdots T p_i\cdots\rangle=0\) \\
station, admissibility, estimates &
\href{https://github.com/yang-mills-group/OpenAI-YM-20261028/blob/3456e5e75cf8d4cd558955a5c7638fba2d68f464/YM/Next/UV_STRESS/Station.lean\#L68-L71}{\texttt{Station}},
\href{https://github.com/yang-mills-group/OpenAI-YM-20261028/blob/3456e5e75cf8d4cd558955a5c7638fba2d68f464/YM/Next/UV_STRESS/Station.lean\#L112-L114}{\texttt{Stations}},
\href{https://github.com/yang-mills-group/OpenAI-YM-20261028/blob/3456e5e75cf8d4cd558955a5c7638fba2d68f464/YM/Next/UV_STRESS/Granular.lean\#L30-L34}{\texttt{Granular}},
\href{https://github.com/yang-mills-group/OpenAI-YM-20261028/blob/3456e5e75cf8d4cd558955a5c7638fba2d68f464/YM/Next/UV_STRESS/Glide.lean\#L24-L31}{\texttt{Balances}}
& \((r,A_{N,p})\) \\
panel, scaling, rescaled derivative &
\href{https://github.com/yang-mills-group/OpenAI-YM-20261028/blob/3456e5e75cf8d4cd558955a5c7638fba2d68f464/YM/Next/UV_STRESS/Station.lean\#L118}{\texttt{Panel}},
\href{https://github.com/yang-mills-group/OpenAI-YM-20261028/blob/3456e5e75cf8d4cd558955a5c7638fba2d68f464/YM/Climb/Cells.lean\#L24-L31}{\texttt{Climb.Scaling}},
\href{https://github.com/yang-mills-group/OpenAI-YM-20261028/blob/3456e5e75cf8d4cd558955a5c7638fba2d68f464/YM/Next/UV_STRESS/Focus.lean\#L19-L20}{\texttt{tiltScale}}
& \(U\); \((r_k,\xi_k,U_k)\); \(\mathrm{tilt}_k\) \\
uniform test convergence, stability, stress limit &
\href{https://github.com/yang-mills-group/OpenAI-YM-20261028/blob/3456e5e75cf8d4cd558955a5c7638fba2d68f464/YM/Next/UV_STRESS/Rates.lean\#L52-L53}{\texttt{Approaches}},
\href{https://github.com/yang-mills-group/OpenAI-YM-20261028/blob/3456e5e75cf8d4cd558955a5c7638fba2d68f464/YM/Next/UV_STRESS/Weak.lean\#L24-L27}{\texttt{Fades}},
\href{https://github.com/yang-mills-group/OpenAI-YM-20261028/blob/3456e5e75cf8d4cd558955a5c7638fba2d68f464/YM/Next/UV_STRESS/Invert.lean\#L121-L125}{\texttt{Stable}},
\href{https://github.com/yang-mills-group/OpenAI-YM-20261028/blob/3456e5e75cf8d4cd558955a5c7638fba2d68f464/YM/Climb/Cells.lean\#L54-L55}{\texttt{Scaling.meter}}
& --- \\
\end{longtable}
}

\subsubsection{7. What remains
unexplained}\label{what-remains-unexplained-24}

\begin{enumerate}
\def\labelenumi{\arabic{enumi}.}
\tightlist
\item
  \textbf{Inputs not checked through all subsidiary proofs.} The
  reflected Gram positivity and Fock pairing factorization in VI.2.b
  were checked at statement level. The Haar-averaging and Ward-implied
  conservation and symmetry results used in VI.2.m were also checked at
  statement level. The invariant spanning results behind VI.2.k, the
  differentiated propagator table behind VI.2.l, and the low-degree
  vector-map classification and bounded-order detection behind VI.2.o
  remain inputs to the sketches. Their direct applications and
  hypotheses were checked, but their full derivations were not
  independently reconstructed.
\item
  \textbf{Lorentz and Casimir data.} The representation, descent and
  generator statements were checked, together with their use in
  assembly. The complexified-rotation construction of the representation
  and the Casimir proofs A09--A12 were not checked in full.
\item
  \textbf{Renormalization-group estimates.} The conversion of balanced
  grid, bulk and finite-difference estimates into stations was checked
  here. Their production in III.11, including all parameter choices in
  the complete weak-comparison result and the shell estimates needed for
  the interacting Ward rows, remains to be verified in that unit. The
  excluded region here is precisely a band around coordinate-projection
  ties; it is not a neighbourhood of the full collision diagonal.
\item
  \textbf{Short-distance normalization.} The comparison of derivative
  and stress rows uses the track estimates of VI.1. The scalar-ratio
  argument and its application here were checked; the underlying track
  Ward estimates A04--A05 remain inputs. The mixed-correlation theorem
  is stated in Lean for the stress assembled in VI.1; VI.2.9(b) gives
  the general consequence of the same convergent-preimage argument.
\end{enumerate}

\subsection{VII.1. Short distances: convergence to free Maxwell
fields}\label{vii.1.-short-distances-convergence-to-free-maxwell-fields}

\textbf{Part:} VII, identification as Yang--Mills. \textbf{Lean:} the
chain of
\href{https://github.com/yang-mills-group/OpenAI-YM-20261028/tree/3456e5e75cf8d4cd558955a5c7638fba2d68f464/YM/Pert}{\texttt{YM/Pert}}
(folders \texttt{Fourier}, \texttt{Align}, \texttt{Warp}, \texttt{Rig},
the jet stencils of \texttt{Lattice}, the statement
\texttt{Position/Crest.finite\_leading}, and \texttt{Dense}),
\href{https://github.com/yang-mills-group/OpenAI-YM-20261028/tree/3456e5e75cf8d4cd558955a5c7638fba2d68f464/YM/Push/FINITE_DETECTION}{\texttt{YM/Push/FINITE\_DETECTION}},
\href{https://github.com/yang-mills-group/OpenAI-YM-20261028/tree/3456e5e75cf8d4cd558955a5c7638fba2d68f464/YM/SharpGram/Weak}{\texttt{YM/SharpGram/Weak}},
\href{https://github.com/yang-mills-group/OpenAI-YM-20261028/tree/3456e5e75cf8d4cd558955a5c7638fba2d68f464/YM/SharpGram/Dense}{\texttt{YM/SharpGram/Dense}},
\href{https://github.com/yang-mills-group/OpenAI-YM-20261028/tree/3456e5e75cf8d4cd558955a5c7638fba2d68f464/YM/RG08}{\texttt{YM/RG08}}
(about thirty small sub-folders),
\href{https://github.com/yang-mills-group/OpenAI-YM-20261028/tree/3456e5e75cf8d4cd558955a5c7638fba2d68f464/YM/PertExtension}{\texttt{YM/PertExtension}},
\href{https://github.com/yang-mills-group/OpenAI-YM-20261028/tree/3456e5e75cf8d4cd558955a5c7638fba2d68f464/YM/OPE/Levy}{\texttt{YM/OPE/Levy}}.
\textbf{Size:} 8,117 lines. They divide as follows:

\begin{itemize}
\tightlist
\item
  3,275 lines in the chain ending in \texttt{finite\_leading};
\item
  1,006 in \texttt{FINITE\_DETECTION};
\item
  616 in \texttt{SharpGram/Weak};
\item
  about 2,250 in \texttt{RG08} (\texttt{Dart} 538, \texttt{Ballast} 338,
  \texttt{Vale} 315, the rest under 160 each);
\item
  316 in \texttt{SharpGram/Dense}, 310 in \texttt{PertExtension}, 108 in
  \texttt{SharpGram/Matrix}, 95 in \texttt{Pert/Fourier} and 63 in
  \texttt{OPE/Levy};
\item
  the remainder in packages of fewer than 40 lines.
\end{itemize}

The main result identifies the order-zero lattice correlations with the
free Maxwell hierarchy. The short-distance limit of the interacting
correlations is stated in
\href{https://github.com/yang-mills-group/OpenAI-YM-20261028/blob/3456e5e75cf8d4cd558955a5c7638fba2d68f464/YM/Pert/FirstTarget.lean\#L26-L39}{\texttt{FirstExtension}}
(VI.2), established by
\href{https://github.com/yang-mills-group/OpenAI-YM-20261028/blob/3456e5e75cf8d4cd558955a5c7638fba2d68f464/YM/Pert/FirstTarget.lean\#L57-L81}{\texttt{firstTarget}}
in VIII.1. The results are:

\begin{itemize}
\tightlist
\item
  the order-zero term of lattice perturbation theory converges, as the
  spacing tends to zero, to the Wick-ordered free Maxwell hierarchy
  evaluated on the on-shell images of the labels (VII.1.1--VII.1.3);
\item
  the free hierarchy detects every label outside the Maxwell equation
  ideal through a single two-point function at a fixed pair of points
  (VII.1.4--VII.1.5);
\item
  using this detection, the measuring charts of the construction can be
  corrected to charts that are exactly compatible with the filtration,
  the symmetries and the unit (VII.1.6);
\item
  estimates on backward windows and local expansions used in the rank
  comparison (VII.1.7).
\end{itemize}

\subsubsection{1. What this unit proves}\label{what-this-unit-proves-25}

\textbf{Setting.}

\emph{Gauge group and continuum objects.}

\begin{itemize}
\tightlist
\item
  \(G\) is a compact matrix gauge group whose Lie algebra
  \(\mathfrak g\) is real simple and nonabelian, with an invariant inner
  product and an orthonormal basis of colours.
\item
  A \emph{key} \(a=(j,\alpha)\) is a coordinate of the real radial-gauge
  potential space \(T_j\): homogeneous \(\mathfrak g\)-valued 1-forms of
  polynomial degree \(j+1\) with \(x\cdot A(x)=0\). Curl identifies
  \(T_j\) with the closed homogeneous curvature jets of degree \(j\).
  Each key has engineering weight \(j+2\).
\item
  \(\mathrm{Pol}\) is the complex polynomial algebra in the keys.
  \(V_d\) consists of its \(\mathrm{Ad}(G)\)-invariant polynomials of
  total engineering weight at most \(d\). The free field space \(F_d\)
  is the quotient of \(V_d\) by the free null space (VI.2).
\item
  \(\mathrm{Sh}_j\) is the space of curvature tensors, symmetric in the
  \(j\) Taylor slots and antisymmetric in the two form indices, whose
  homogeneous fields solve \(dF=d{*}F=0\). Taylor coefficients are
  normalized by \(1/j!\). Write \(\kappa:\mathrm{Sh}_j\to T_j\) for the
  inverse radial-gauge curl map
  (\href{https://github.com/yang-mills-group/OpenAI-YM-20261028/blob/3456e5e75cf8d4cd558955a5c7638fba2d68f464/YM/Maxwell/Shell.lean\#L46-L47}{\texttt{gaugeJet}}).
\item
  \(\rho:\mathrm{Pol}\to \mathrm{WickPol}=\mathbb C[\text{shell coordinates}]\)
  restricts a polynomial to on-shell jets
  (\href{https://github.com/yang-mills-group/OpenAI-YM-20261028/blob/3456e5e75cf8d4cd558955a5c7638fba2d68f464/YM/Maxwell/Equation.lean\#L47-L50}{\texttt{reduceEq}}).
\item
  The \emph{equation ideal} \(\mathcal I_d=V_d\cap\ker\rho\) consists of
  labels that vanish on every on-shell Taylor tower
  (\href{https://github.com/yang-mills-group/OpenAI-YM-20261028/blob/3456e5e75cf8d4cd558955a5c7638fba2d68f464/YM/Maxwell/Equation.lean\#L72}{\texttt{equationIdeal}}).
\item
  \(\omega_d=\rho\circ\mathrm{sec}_d:F_d\to\mathrm{WickPol}\) is the
  on-shell image of a label. Here \(\mathrm{sec}_d\) is the
  representative polynomial of VII.3, obtained by projecting radial
  curvature coordinates onto the Maxwell shell
  (\href{https://github.com/yang-mills-group/OpenAI-YM-20261028/blob/3456e5e75cf8d4cd558955a5c7638fba2d68f464/YM/Pert/Lattice/Bundles.lean\#L35}{\texttt{sections}},
  \href{https://github.com/yang-mills-group/OpenAI-YM-20261028/blob/3456e5e75cf8d4cd558955a5c7638fba2d68f464/YM/Pert/Lattice/Lifts.lean\#L55-L60}{\texttt{fixedPolynomial}}).
\end{itemize}

\emph{The free Maxwell hierarchy (III.16).}

\begin{itemize}
\tightlist
\item
  \(\Gamma(x)=1/(4\pi^2|x|^2)\)
  (\href{https://github.com/yang-mills-group/OpenAI-YM-20261028/blob/3456e5e75cf8d4cd558955a5c7638fba2d68f464/YM/Maxwell/PotentialKernel.lean\#L22}{\texttt{Green}}).
\item
  The kernel
  \href{https://github.com/yang-mills-group/OpenAI-YM-20261028/blob/3456e5e75cf8d4cd558955a5c7638fba2d68f464/YM/Maxwell/Curls.lean\#L20}{\texttt{bareCurl}}
  is
  \(\delta_{bk}\partial_a\partial_m\Gamma-\delta_{bm}\partial_a\partial_k\Gamma-\delta_{ak}\partial_b\partial_m\Gamma+\delta_{am}\partial_b\partial_k\Gamma\).
  This is \(\langle F_{ab}(x)F_{km}(y)\rangle\) at \(x-y\) for \(F=dA\),
  where \(A\) is the Feynman-gauge free field with
  \(\langle A^c_\mu(x)A^{c'}_\nu(y)\rangle=\delta^{cc'}\delta_{\mu\nu}\Gamma(x-y)\).
\item
  For linear forms \(P\) on \(\mathrm{Sh}_j\) and \(Q\) on
  \(\mathrm{Sh}_l\), \(K(P,Q)(x-y)\) is \(P\) applied to the \(j\)-jet
  in \(x\) and \(Q\) to the \(l\)-jet in \(y\) of this kernel,
  colour-diagonal
  (\href{https://github.com/yang-mills-group/OpenAI-YM-20261028/blob/3456e5e75cf8d4cd558955a5c7638fba2d68f464/YM/Maxwell/Leg.lean\#L161-L162}{\texttt{legKernel}}).
  The value is independent of extensions of \(P,Q\) to full curvature
  tensors, because the kernel solves the free equations in each
  variable.
\item
  For a key \(a\), write \(x_a\) for its real coordinate functional on
  \(T_j\) and \(\kappa_a=x_a\circ\kappa\), and write
  \(K_{ab}=K(\kappa_a,\kappa_b)\). The Taylor coefficients in both legs
  include their factorial denominators; differentiation at the second
  endpoint contributes \((-1)^l\).
\item
  The \emph{Wick moment} \(\mathcal W_x(q_1,\dots,q_r)\), for
  \(x\in\mathrm{Conf}_r\) and \(q_i\in\mathrm{WickPol}\), is multilinear
  (\href{https://github.com/yang-mills-group/OpenAI-YM-20261028/blob/3456e5e75cf8d4cd558955a5c7638fba2d68f464/YM/Maxwell/Wick.lean\#L126-L127}{\texttt{Wick.moment}}).
  On monomials it is the sum, over perfect matchings of the variable
  occurrences that never pair two occurrences at the same point, of the
  product of the entries \(K(\cdot,\cdot)(x_i-x_j)\). It is the
  correlation of the Wick-ordered products \(:\!q_i\!:(x_i)\) of the
  Gaussian jet field with covariance \(K\).
\item
  The free hierarchy of the labels is
  \(x\mapsto\mathcal W_x(\omega_d p_1,\dots,\omega_d p_r)\)
  (\href{https://github.com/yang-mills-group/OpenAI-YM-20261028/blob/3456e5e75cf8d4cd558955a5c7638fba2d68f464/YM/Ward/Sum.lean\#L23}{\texttt{Ward.disc}}
  of the free table).
\end{itemize}

\emph{The lattice prescription at order zero (definitions from VII.3).}

\begin{itemize}
\item
  \(\Lambda_N=(\mathbb Z/N)^4\) with \(N\) even, \(N\ge 4\).
\item
  The field space \(\mathcal A_N\) consists of \(\mathfrak g\)-valued
  link variables that are odd under the point reflection:
  \(A(-x-e_i,i)=-A(x,i)\), i.e.~\(A\) is invariant under \(x\mapsto -x\)
  as a 1-form
  (\href{https://github.com/yang-mills-group/OpenAI-YM-20261028/blob/3456e5e75cf8d4cd558955a5c7638fba2d68f464/YM/Pert/Lattice/Geometry.lean\#L107-L111}{\texttt{LinkP}}).
\item
  The quadratic form
  \(\mathsf Q_N(A)=\tfrac12\big(\sum_{x,i<j}|\mathrm{curl}_{ij}A(x)|^2+\sum_x|\mathrm{div}A(x)|^2\big)\)
  is the Hessian at the identity of the Wilson action plus the Feynman
  gauge-fixing term \(\tfrac14\sum|\mathrm{div}A|^2\)
  (\href{https://github.com/yang-mills-group/OpenAI-YM-20261028/blob/3456e5e75cf8d4cd558955a5c7638fba2d68f464/YM/Pert/Lattice/Precision.lean\#L76-L78}{\texttt{form}},
  \href{https://github.com/yang-mills-group/OpenAI-YM-20261028/blob/3456e5e75cf8d4cd558955a5c7638fba2d68f464/YM/Pert/Lattice/Action.lean\#L46}{\texttt{action}}).
\item
  \(\mathsf Q_N\) is positive definite on \(\mathcal A_N\): a closed and
  divergence-free field is constant, and a reflection-odd constant field
  is zero. \(\mu_N\) is the centred Gaussian with density proportional
  to \(\exp(-\mathsf Q_N(A)/2)\)
  (\href{https://github.com/yang-mills-group/OpenAI-YM-20261028/blob/3456e5e75cf8d4cd558955a5c7638fba2d68f464/YM/Pert/Lattice/Measure.lean\#L31}{\texttt{data}}).
\item
  Expectations are formal power series in the coupling. The interaction
  exponent has order-\(n\) coefficient equal to minus the
  order-\((n+2)\) action coefficient for \(n\ge1\), and zero constant
  coefficient. Its exponential is multiplied by the Taylor series of the
  Haar Jacobian and Faddeev--Popov density. Each jet-coordinate series
  uses its order-\((n+1)\) Taylor coefficient at order \(n\)
  (\href{https://github.com/yang-mills-group/OpenAI-YM-20261028/blob/3456e5e75cf8d4cd558955a5c7638fba2d68f464/YM/Pert/Lattice/Series.lean\#L61}{\texttt{weight}},
  \href{https://github.com/yang-mills-group/OpenAI-YM-20261028/blob/3456e5e75cf8d4cd558955a5c7638fba2d68f464/YM/Pert/Lattice/Readers.lean\#L32}{\texttt{seriesLex}}).
  Order zero is therefore a \(\mu_N\)-expectation;
  \(\langle\cdot\rangle^{(0)}\) denotes the order-zero coefficient.
\item
  \(\ell^h_{x}:\mathcal A_N\to T_j\) is the radial-gauge image of the
  linear lattice curvature jet at site \(x\)
  (\href{https://github.com/yang-mills-group/OpenAI-YM-20261028/blob/3456e5e75cf8d4cd558955a5c7638fba2d68f464/YM/Pert/Lattice/TopLinear.lean\#L38}{\texttt{slope}}).
  It is obtained from the \(j\)-fold forward difference of the plaquette
  curl divided by \(j!\) by symmetrization and the radial-gauge map,
  scaled by \(h^{-(j+2)}\) and averaged over the hyperoctahedral group
  \(W_4\) of signed axis permutations. \(\ell^h_{x,a}\) is its \(a\)-th
  coordinate.
\item
  The order-zero lattice representative of a polynomial
  \(p\in\mathrm{Pol}\) using only keys of rank at most \(d\), at \(o\)
  is \[R_h(p,o)(A)=\mathbb E_\xi\, p\big(\ell^h_o(A)+i\xi\big),\] where
  \(\xi\) is an independent centred Gaussian on the keys of rank
  \(\le d\) with covariance
  \[S_h(a,b)=(2\pi)^{-4}\int_{[-\pi,\pi]^4}\hat\Delta(k)^{-1}\sum_c \lambda_a(c,k)\lambda_b(c,k)\,dk .\]
  Here \(\hat\Delta(k)=\sum_i4\sin^2(k_i/2)\), \(c\) runs over
  cosine/sine, axis and colour channels, and \(\lambda_a(c,\cdot)\) is
  \(\ell^h_{0,a}\) applied to the corresponding plane wave on
  \(\mathbb Z^4\). So \(S_h\) is the covariance of the jet stencils at
  one site for the Feynman-gauge Gaussian on the infinite lattice
  \(\mathbb Z^4\), and \(R_h(p,o)=\,:\!p(\ell^h_o)\!:_{S_h}\) is
  Wick-ordered with respect to the infinite-lattice self-contraction
  (\href{https://github.com/yang-mills-group/OpenAI-YM-20261028/blob/3456e5e75cf8d4cd558955a5c7638fba2d68f464/YM/Pert/Lattice/Blank.lean\#L41}{\texttt{column}},
  \href{https://github.com/yang-mills-group/OpenAI-YM-20261028/blob/3456e5e75cf8d4cd558955a5c7638fba2d68f464/YM/Pert/Lattice/Symbol.lean\#L50}{\texttt{link}}).
\item
  \textbf{Regulators.} For an integer \(L\ge2\) set \(h_m=L^{-m}\) and
  \(N_m=4L^{m^2+m+1}\). Lattice points are reported relative to the
  centre \(c_m=(N_m/4,\dots,N_m/4)\), through \(o\mapsto h_m(o-c_m)\)
  with representatives in \((-N_m/2,N_m/2]^4\).
\item
  \textbf{Profile.} For \(f:(\mathbb R^4)^r\to\mathbb C\) and
  \(p\in F_d^{\,r}\),
  \[\Pi_m(f;p)=h_m^{4r}\sum_{o\in\Lambda_{N_m}^r}f\big(h_m(o-c_m)\big)\,\Big\langle\prod_{j}R_{h_m}(\mathrm{sec}_d\,p_j,o_j)\Big\rangle^{(0)}_{N_m}\]
  (\href{https://github.com/yang-mills-group/OpenAI-YM-20261028/blob/3456e5e75cf8d4cd558955a5c7638fba2d68f464/YM/Pert/Lattice/Bundles.lean\#L55-L56}{\texttt{profile}},
  \href{https://github.com/yang-mills-group/OpenAI-YM-20261028/blob/3456e5e75cf8d4cd558955a5c7638fba2d68f464/YM/Pert/Lattice/Bundles.lean\#L41-L42}{\texttt{row}}).
\item
  \textbf{Itinerary.} An itinerary is a sequence \((N_i,h_i)\) with
  \(N_i\) even, \(h_i>0\), \(N_ih_i\) a positive integer tending to
  \(\infty\), \(h_i\to0\), and \(N_ih_i^k\to\infty\) for every
  \(k\in\mathbb N\)
  (\href{https://github.com/yang-mills-group/OpenAI-YM-20261028/blob/3456e5e75cf8d4cd558955a5c7638fba2d68f464/YM/Pert/Align/Burst.lean\#L50-L59}{\texttt{Itinerary}}).
  Sites \(o_i\in\Lambda_{N_i}\) are \emph{homing} to \(z\) if
  \(h_i(o_i-c_{N_i})\to z\)
  (\href{https://github.com/yang-mills-group/OpenAI-YM-20261028/blob/3456e5e75cf8d4cd558955a5c7638fba2d68f464/YM/Pert/Align/Schedule.lean\#L53}{\texttt{Home}}).
\end{itemize}

\textbf{Proposition VII.1.1 (regulator removal at order zero).} Fix an
integer \(L\ge2\), \(d,r\in\mathbb N\), label frames, and
\(p_1,\dots,p_r\in F_d\). Let \(f\) be continuous with compact support
and \(\mathrm{supp}\,f\subset\mathrm{Conf}_r\), where support is closed.
Then
\[\lim_{m\to\infty}\Pi_m(f;p)=\int f(x)\,\mathcal W_x(\omega_dp_1,\dots,\omega_dp_r)\,dx .\]
In board form: let \(Q\) be a position board whose profiles are smooth
test functions compactly supported in \(\mathrm{Conf}\) (``clean
profiles''). Then the order-zero board readings converge to the free
readings, slot by slot and for every tuple of basis labels.
\emph{Status:} theorem. There is no hypothesis on the construction; it
concerns lattice Gaussian integrals and the continuum kernel only.
\emph{Lean:}
\href{https://github.com/yang-mills-group/OpenAI-YM-20261028/blob/3456e5e75cf8d4cd558955a5c7638fba2d68f464/YM/Pert/Rig/Linked.lean\#L59-L69}{\texttt{Rig.zerothProfile}}.
Board form:
\href{https://github.com/yang-mills-group/OpenAI-YM-20261028/blob/3456e5e75cf8d4cd558955a5c7638fba2d68f464/YM/Pert/Position/Crest.lean\#L13-L16}{\texttt{Position.finite\_leading}},
using the definitions
\href{https://github.com/yang-mills-group/OpenAI-YM-20261028/blob/3456e5e75cf8d4cd558955a5c7638fba2d68f464/YM/Pert/Position/Board.lean\#L98}{\texttt{finite}},
\href{https://github.com/yang-mills-group/OpenAI-YM-20261028/blob/3456e5e75cf8d4cd558955a5c7638fba2d68f464/YM/Pert/Position/Board.lean\#L103-L104}{\texttt{white}}
and
\href{https://github.com/yang-mills-group/OpenAI-YM-20261028/blob/3456e5e75cf8d4cd558955a5c7638fba2d68f464/YM/Pert/Position/Actual.lean\#L13}{\texttt{CleanProfiles}}
(all VII.3).

\textbf{Proposition VII.1.2 (pointwise convergence of order-zero
correlations).} Let \((N_i,h_i)\) be an itinerary with \(4\mid N_i\).
Let \(z\in\mathrm{Conf}_r\), and let \(o_{i,j}\in\Lambda_{N_i}\) be
homing to \(z_j\) for each \(j\). Then for all \(p_1,\dots,p_r\in F_d\),
\[\Big\langle\prod_jR_{h_i}(\mathrm{sec}_dp_j,o_{i,j})\Big\rangle^{(0)}_{N_i}\longrightarrow\mathcal W_z(\omega_dp_1,\dots,\omega_dp_r).\]
\emph{Status:} theorem. \emph{Lean:}
\href{https://github.com/yang-mills-group/OpenAI-YM-20261028/blob/3456e5e75cf8d4cd558955a5c7638fba2d68f464/YM/Pert/Warp/Horizons.lean\#L65-L68}{\texttt{Warp.row\_limit}},
through
\href{https://github.com/yang-mills-group/OpenAI-YM-20261028/blob/3456e5e75cf8d4cd558955a5c7638fba2d68f464/YM/Pert/Warp/Horizons.lean\#L50-L63}{\texttt{poly\_limit}}
and
\href{https://github.com/yang-mills-group/OpenAI-YM-20261028/blob/3456e5e75cf8d4cd558955a5c7638fba2d68f464/YM/Pert/Warp/Approach.lean\#L19-L32}{\texttt{earth}}.

\textbf{Proposition VII.1.3 (corrected two-leg covariance).} Let
\((N_i,h_i)\) be an itinerary with \(4\mid N_i\). Let \(\iota\) be any
index set, let \(z:\iota\to\mathbb R^4\) be injective, and for every
\(j\in\iota\) let \(o_{i,j}\) be sites homing to \(z_j\). For indices
\(j,l\in\iota\) and arbitrary keys \(a,b\), define
\[\kappa_h(j,a;l,b)=\mathrm{Cov}_{\mu_N}\big(\ell^h_{o_j,a},\ell^h_{o_l,b}\big)-\delta_{jl}\,S_h(a,b).\]
Then
\[\kappa_{h_i}(j,a;l,b)\to\begin{cases}K_{ab}(z_j-z_l), & j\ne l,\\ 0, & j=l.\end{cases}\]
\emph{Status:} theorem. \emph{Lean:}
\href{https://github.com/yang-mills-group/OpenAI-YM-20261028/blob/3456e5e75cf8d4cd558955a5c7638fba2d68f464/YM/Pert/Align/Drivers.lean\#L39-L44}{\texttt{Align.corrected\_limit}},
from
\href{https://github.com/yang-mills-group/OpenAI-YM-20261028/blob/3456e5e75cf8d4cd558955a5c7638fba2d68f464/YM/Pert/Align/Cancel.lean\#L80-L88}{\texttt{orbit\_limit}},
\href{https://github.com/yang-mills-group/OpenAI-YM-20261028/blob/3456e5e75cf8d4cd558955a5c7638fba2d68f464/YM/Pert/Align/Cancel.lean\#L89-L95}{\texttt{orbit\_contact}}
and
\href{https://github.com/yang-mills-group/OpenAI-YM-20261028/blob/3456e5e75cf8d4cd558955a5c7638fba2d68f464/YM/Pert/Align/Cancel.lean\#L108-L113}{\texttt{shell\_exact}}.

\textbf{Proposition VII.1.4 (two-point detection).} Let
\(e_{\mathrm t}\) be the unit vector of the time axis, \(\sigma\) the
time reflection \(x_{\mathrm t}\mapsto-x_{\mathrm t}\) acting on labels,
and \(\bar v\) the complex conjugate label. If
\(v\in V_d\setminus\mathcal I_d\), then
\[\mathrm{Re}\;\mathcal W_{(e_{\mathrm t},0)}\big(\rho v,\ \rho(\sigma\bar v)\big)>0 .\]
More generally, let \(\Phi\) and \(\Psi\) be the algebra maps sending
shell coordinates to the diagonalizing and reflected probe variables of
Lemma VII.1.m. For every nonzero \(q\in\mathrm{WickPol}\),
\(\mathrm{Re}\,\mathcal W_{(e_{\mathrm t},0)}(\Phi q,\Psi\bar q)>0\).
Here \(\Phi\) is surjective, and \(\rho(\sigma\bar v)=\Psi(\bar q)\)
whenever \(\rho v=\Phi q\). \emph{Status:} theorem, about the free
kernel only. No Wilson family and no reflection positivity of a measure
is involved. \emph{Lean:}
\href{https://github.com/yang-mills-group/OpenAI-YM-20261028/blob/3456e5e75cf8d4cd558955a5c7638fba2d68f464/YM/Push/FINITE_DETECTION/TwoPoint.lean\#L151-L155}{\texttt{Free.two\_detect}},
\href{https://github.com/yang-mills-group/OpenAI-YM-20261028/blob/3456e5e75cf8d4cd558955a5c7638fba2d68f464/YM/Push/FINITE_DETECTION/TwoPoint.lean\#L132-L149}{\texttt{Free.positive}},
\href{https://github.com/yang-mills-group/OpenAI-YM-20261028/blob/3456e5e75cf8d4cd558955a5c7638fba2d68f464/YM/Push/FINITE_DETECTION/TwoPoint.lean\#L26}{\texttt{detectors}}.

\textbf{Proposition VII.1.5 (Kelvin formulas for the free kernel).}

\begin{enumerate}
\def\labelenumi{(\alph{enumi})}
\item
  Let \(A\) be a real symmetric \(j\)-tensor on \(\mathbb R^4\) whose
  diagonal \(x\mapsto A(x,\dots,x)\) is harmonic. Then for \(x\ne0\)
  \[\frac1{j!}A(\partial,\dots,\partial)\,\Gamma(x)=\frac{(-2)^j}{4\pi^2}\,|x|^{-2(j+1)}A(x,\dots,x).\]
\item
  Let \(Y\in\mathrm{Sh}_j\) be \(\mathfrak g\)-valued, with homogeneous
  Maxwell solution \(Y^\natural(x)=Y(x,\dots,x)\). Contract the first
  leg of the curvature kernel at \(x\) with \(Y\), using its rank-\(j\)
  Taylor coefficient and summing over colours, ordered form indices and
  jet indices, without a factor \(1/2\). For \(x\ne0\) the result equals
  \[\gamma_j\,(\iota^*Y^\natural)(x),\qquad \gamma_j=\frac{2^j(j+1)(j+2)}{\pi^2}.\]
  Here \(\iota(x)=-x/|x|^2\) and \(\iota^*\) is the pullback of 2-forms:
  \[(\iota^*F)_{km}(x)=|x|^{-4}\left[B_{km}-2|x|^{-2}\left(x_k\sum_a x_aB_{am}-x_m\sum_a x_aB_{ak}\right)\right],\qquad B_{ab}=F_{ab}(\iota x).\]
  \emph{Status:} theorem. \emph{Lean:}
  \href{https://github.com/yang-mills-group/OpenAI-YM-20261028/blob/3456e5e75cf8d4cd558955a5c7638fba2d68f464/YM/Push/FINITE_DETECTION/Harmonic.lean\#L157-L180}{\texttt{Kelvin.harmonic\_Green}}
  (with
  \href{https://github.com/yang-mills-group/OpenAI-YM-20261028/blob/3456e5e75cf8d4cd558955a5c7638fba2d68f464/YM/Push/FINITE_DETECTION/Harmonic.lean\#L83}{\texttt{R}}
  \(=|x|^{-2}\)), and
  \href{https://github.com/yang-mills-group/OpenAI-YM-20261028/blob/3456e5e75cf8d4cd558955a5c7638fba2d68f464/YM/Push/FINITE_DETECTION/Response.lean\#L172-L176}{\texttt{Kelvin.response\_eq}}
  (with
  \href{https://github.com/yang-mills-group/OpenAI-YM-20261028/blob/3456e5e75cf8d4cd558955a5c7638fba2d68f464/YM/Push/FINITE_DETECTION/Response.lean\#L64-L65}{\texttt{inverted}}
  and
  \href{https://github.com/yang-mills-group/OpenAI-YM-20261028/blob/3456e5e75cf8d4cd558955a5c7638fba2d68f464/YM/Push/FINITE_DETECTION/Response.lean\#L92}{\texttt{gain}}).
\end{enumerate}

\textbf{Proposition VII.1.6 (corrected measuring charts).} Fix a Wilson
family \(C\) with a realization of its defining conditions, a
construction record \(U\) with limit table \(S\) and extraction order
\(f\), and OPE table data as in VII.4. Let \(D\in\mathbb N\), and let
\(u\) be a nontrivial filter on \(\mathbb N\times\mathbb N\) finer than
the extraction filter of \(f\).

Assume a weak measuring portfolio \(X\) through degree \(D\). Its data
include the observable bounds, Haar bounds, a Haar chart ratio, three
price panels, and a constant dominating the label costs through \(D\).
It supplies an index \(g(k)\to\infty\), finite-regulator linear label
maps, and limiting maps \(\Phi_{k,d}\) on the reference label-coordinate
spaces. Eventually in \(k\):

\begin{itemize}
\tightlist
\item
  \(g(k)\le k\) and
  \(\mathrm{scale}(k)\le\mathrm{scale}(g(k))+\log\hat\beta_k\);
\item
  the finite-regulator maps satisfy the simultaneous rank-comparison
  conditions through \(D\): every degree is available at step \(g(k)\),
  the block ratio is below the monitor bound, the maps are tested and
  equivariant, and all three panels supply the required observable and
  remainder estimates;
\item
  their rescaled read-out maps converge to \(\Phi_{k,d}\) along the same
  filter \(u\), for every \(d\le D\) and every reference label.
\end{itemize}

Let \(\lambda_d\) be the fixed equivariant representative of the free
quotient in reference label coordinates, and let \(\pi^S_d\) be
projection to \(S_d\). Set \(D_{k,d}=\pi^S_d\Phi_{k,d}\lambda_d\).
Define \(\chi_k(d)=D_{k,d}\) when this map is bijective, and choose a
fixed isomorphism \(F_d\to S_d\) otherwise. Then there are complex
linear isomorphisms \(p_k(d):F_d\to S_d\), for every
\(k,d\in\mathbb N\), such that:

\begin{enumerate}
\def\labelenumi{\arabic{enumi}.}
\item
  For all sufficiently large \(k\), simultaneously for \(d\le D\),
  \(p_k(d)\) commutes with signed permutations and conjugation and
  preserves the unit. For all \(d\le e\le D\),
  \[p_k(e)\circ\iota^F_{d,e}=\iota^S_{d,e}\circ p_k(d).\]
\item
  For every \(d\le D\) and \(v\in F_d\),
  \[\chi_k(d)^{-1}p_k(d)v\longrightarrow v.\]
\end{enumerate}

The simultaneous finite-regulator conditions are
\href{https://github.com/yang-mills-group/OpenAI-YM-20261028/blob/3456e5e75cf8d4cd558955a5c7638fba2d68f464/YM/MergeProof/East/Plume.lean\#L23-L27}{\texttt{Poised}},
and the representative is
\href{https://github.com/yang-mills-group/OpenAI-YM-20261028/blob/3456e5e75cf8d4cd558955a5c7638fba2d68f464/YM/Further/CORRECTED_LIFT_COEFFICIENTS/Basis.lean\#L115-L117}{\texttt{ladder}}.

\emph{Status:} theorem under the stated construction and portfolio
hypotheses. Lemma VII.1.q supplies the required two-point convergence.
\emph{Lean:}
\href{https://github.com/yang-mills-group/OpenAI-YM-20261028/blob/3456e5e75cf8d4cd558955a5c7638fba2d68f464/YM/SharpGram/Weak/Compatible.lean\#L140-L159}{\texttt{Weak.Portfolio.tune\_grades}},
with
\href{https://github.com/yang-mills-group/OpenAI-YM-20261028/blob/3456e5e75cf8d4cd558955a5c7638fba2d68f464/YM/SharpGram/Weak/Reading.lean\#L35-L49}{\texttt{Portfolio}}
and
\href{https://github.com/yang-mills-group/OpenAI-YM-20261028/blob/3456e5e75cf8d4cd558955a5c7638fba2d68f464/YM/SharpGram/Weak/Compatible.lean\#L76-L79}{\texttt{chart}}.

\textbf{Proposition VII.1.7 (backward windows and local expansions).}

\begin{enumerate}
\def\labelenumi{(\alph{enumi})}
\item
  Fix the RG constants of III.2, Haar bounds, a Haar chart ratio
  satisfying the local ratio bounds, and a relay of admissible price
  panels with activation bounds \(u_j\le L_j\). Each degree is
  eventually active; its three panels satisfy the comparison
  inequalities and the analytic bounds required for the rank estimates.
  Let \(U\) be a wait with window
  \[w_j(t)=\left\lfloor\frac{\log(\bar\beta+c_Gt)}{\log L_j}\right\rfloor,\]
  where \(c_G>0\) is the coupling slope. There is a wait \(W\) with no
  smaller minimum stage length or entry thresholds and with the same
  window, together with panels and cost bounds for every degree, having
  the following property. For every realized Wilson family satisfying
  \(W\), the prescribed observable bounds and the
  polynomial-representative conditions, and every \(D\), all
  sufficiently late steps \(k\) whose scale exceeds the stage entry
  threshold and whose backward window has length at most \(k\) and lies
  in one stage admit simultaneous finite-regulator rank-comparison data
  through \(D\), as in VII.1.6. The returned panels satisfy the three
  comparison inequalities at every degree, and each cost bound dominates
  the label costs through that degree.
\item
  Let \(U\) have the window in (a), and let \(V\) have no smaller
  stage-length bounds, entry thresholds or windows. Assume, for every
  stage \(j\), that the entry threshold of \(V\) at \(j\) is at least
  that of \(U\) at \(j+1\), and that \(\bar\beta+c_Gt\ge c_G\) above the
  \(V\) threshold at \(j\). Every RG trajectory satisfying \(V\) admits
  \(g(k)\to\infty\) such that eventually \(g(k)\le k\), the step
  \(g(k)\) is above its \(U\) entry threshold, its backward \(U\) window
  has length at most \(g(k)\) and lies in one stage, and
  \[\mathrm{scale}(k)\le\mathrm{scale}(g(k))+\log\hat\beta_k.\] These
  threshold conditions are
  \href{https://github.com/yang-mills-group/OpenAI-YM-20261028/blob/3456e5e75cf8d4cd558955a5c7638fba2d68f464/YM/SharpGram/Dense/Cover.lean\#L27-L28}{\texttt{LateEnough}}.
\item
  Fix a finite index set \(I\), a marker subset \(T\subset I\), an
  anchor \(p\in T\), \(D,k,E,q\in\mathbb N\), a support parameter
  \(R\in\mathbb R\), and \(v>0\), with \(q\le k+E\). There are
  \(0<w\le v\) and \(C_0\ge0\), independent of the layout and its bound
  \(B\), with the following property. Suppose a layout has marker
  weights \(a_i\in\mathbb N\) with \(\sum_{i\in T}a_i=k\), depends only
  on the occupied markers, vanishes on the empty marker set, and
  satisfies the prescribed tensor symmetry, rank and analytic
  realization conditions
  (\href{https://github.com/yang-mills-group/OpenAI-YM-20261028/blob/3456e5e75cf8d4cd558955a5c7638fba2d68f464/YM/RG08/Creek/Step.lean\#L16-L20}{\texttt{Sense}}).
  Suppose also that every coefficient of expansion order \(i\le E\), at
  every leg rank, has the anchored \(v\)-weighted absolute-sum bound
  \(B\). Translate the root to \(x_p\) and shift the order by \(q\),
  setting coefficients with \(q+n<k\) to zero. Through leg rank
  \(k+E-q\), the resulting coefficients are exactly the Taylor
  coefficients of a finite sum of curvature expressions of ranks
  \[k\mathbin{\dot-}q,\ldots,k+E-q,\] where
  \(k\mathbin{\dot-}q=\max(k-q,0)\). Each expression depends only on
  \(T\), retains the required symmetries and analytic realization, and
  has the anchored absolute-sum bound \(BC_0\) with weight
  \(\exp(w\,\mathrm{girth})\), including the exponential weight on
  relative marker positions.
\end{enumerate}

\emph{Status:} theorems. \emph{Lean:}
\href{https://github.com/yang-mills-group/OpenAI-YM-20261028/blob/3456e5e75cf8d4cd558955a5c7638fba2d68f464/YM/SharpGram/Dense/Slender.lean\#L127-L143}{\texttt{Dense.slim\_wait\_for}},
\href{https://github.com/yang-mills-group/OpenAI-YM-20261028/blob/3456e5e75cf8d4cd558955a5c7638fba2d68f464/YM/SharpGram/Dense/Cover.lean\#L155-L165}{\texttt{Dense.nearby}},
\href{https://github.com/yang-mills-group/OpenAI-YM-20261028/blob/3456e5e75cf8d4cd558955a5c7638fba2d68f464/YM/RG08/Dart/Seams.lean\#L34-L45}{\texttt{Dart.coin\_pre}}.

\subsubsection{2. What it uses}\label{what-it-uses-25}

\textbf{From other units.}

\begin{itemize}
\tightlist
\item
  \emph{VII.3 (the lattice prescription and position boards):}

  \begin{itemize}
  \tightlist
  \item
    the reflection-odd field space, the Feynman-gauge form, the Gaussian
    \(\mu_N\) and the coupling series;
  \item
    the jet stencils
    (\href{https://github.com/yang-mills-group/OpenAI-YM-20261028/blob/3456e5e75cf8d4cd558955a5c7638fba2d68f464/YM/Pert/Lattice/TopLinear.lean\#L33-L35}{\texttt{wideLinear}}),
    the Wick-subtraction covariance \(S_h\)
    (\href{https://github.com/yang-mills-group/OpenAI-YM-20261028/blob/3456e5e75cf8d4cd558955a5c7638fba2d68f464/YM/Pert/Lattice/Symbol.lean\#L58-L59}{\texttt{shellCov}})
    and the on-shell representatives
    \href{https://github.com/yang-mills-group/OpenAI-YM-20261028/blob/3456e5e75cf8d4cd558955a5c7638fba2d68f464/YM/Pert/Lattice/Lifts.lean\#L78}{\texttt{physicalLabel}};
  \item
    rows, profiles and the board objects
    \href{https://github.com/yang-mills-group/OpenAI-YM-20261028/blob/3456e5e75cf8d4cd558955a5c7638fba2d68f464/YM/Pert/Position/Board.lean\#L54-L59}{\texttt{Conditions}},
    \href{https://github.com/yang-mills-group/OpenAI-YM-20261028/blob/3456e5e75cf8d4cd558955a5c7638fba2d68f464/YM/Pert/Position/Board.lean\#L98}{\texttt{finite}}
    and
    \href{https://github.com/yang-mills-group/OpenAI-YM-20261028/blob/3456e5e75cf8d4cd558955a5c7638fba2d68f464/YM/Pert/Position/Board.lean\#L103-L104}{\texttt{white}}.
  \end{itemize}
\item
  \emph{III.16 (free Maxwell analysis):} the shells,
  \href{https://github.com/yang-mills-group/OpenAI-YM-20261028/blob/3456e5e75cf8d4cd558955a5c7638fba2d68f464/YM/Maxwell/Equation.lean\#L47-L50}{\texttt{reduceEq}},
  \href{https://github.com/yang-mills-group/OpenAI-YM-20261028/blob/3456e5e75cf8d4cd558955a5c7638fba2d68f464/YM/Maxwell/Equation.lean\#L72}{\texttt{equationIdeal}},
  \(\Gamma\), the curvature kernel
  \href{https://github.com/yang-mills-group/OpenAI-YM-20261028/blob/3456e5e75cf8d4cd558955a5c7638fba2d68f464/YM/Maxwell/Leg.lean\#L161-L162}{\texttt{legKernel}}
  and the Wick moment
  \href{https://github.com/yang-mills-group/OpenAI-YM-20261028/blob/3456e5e75cf8d4cd558955a5c7638fba2d68f464/YM/Maxwell/Wick.lean\#L126-L127}{\texttt{Wick.moment}}.
\item
  \emph{VI.2 and VI.1:} the free table and its label spaces \(F_d\)
  (VI.2); the hierarchy as a function of positions,
  \href{https://github.com/yang-mills-group/OpenAI-YM-20261028/blob/3456e5e75cf8d4cd558955a5c7638fba2d68f464/YM/Ward/Sum.lean\#L23}{\texttt{Ward.disc}}
  (VI.1).
\item
  \emph{II.1 and I.1:} labels, keys and Taylor jets (II.1); the lattice
  and the signed permutations \(W_4\) (I.1).
\item
  \emph{For VII.1.6:}

  \begin{itemize}
  \tightlist
  \item
    the construction record (VIII.1) and the limit table \(S\);
  \item
    the OPE table work items, of which only the system part is used
    (VII.4);
  \item
    the chart-correction machinery
    \href{https://github.com/yang-mills-group/OpenAI-YM-20261028/blob/3456e5e75cf8d4cd558955a5c7638fba2d68f464/YM/Extra/WEAK_SCALING_ACTLAW/Adjust.lean\#L35-L55}{\texttt{retraction}},
    \href{https://github.com/yang-mills-group/OpenAI-YM-20261028/blob/3456e5e75cf8d4cd558955a5c7638fba2d68f464/YM/Extra/WEAK_SCALING_ACTLAW/Adjust.lean\#L61-L66}{\texttt{adjusted}},
    \href{https://github.com/yang-mills-group/OpenAI-YM-20261028/blob/3456e5e75cf8d4cd558955a5c7638fba2d68f464/YM/Extra/WEAK_SCALING_ACTLAW/Adjust.lean\#L137-L140}{\texttt{rotateChart}}
    and
    \href{https://github.com/yang-mills-group/OpenAI-YM-20261028/blob/3456e5e75cf8d4cd558955a5c7638fba2d68f464/YM/Extra/WEAK_SCALING_ACTLAW/Adjust.lean\#L171-L172}{\texttt{Graded}}
    (VII.5);
  \item
    the rank-comparison error bound and its decay
    \href{https://github.com/yang-mills-group/OpenAI-YM-20261028/blob/3456e5e75cf8d4cd558955a5c7638fba2d68f464/YM/Next/RANK/Allowance.lean\#L34-L37}{\texttt{Ship.diminishes}}
    (III.11), and the chessboard-to-moment estimate for observables
    \href{https://github.com/yang-mills-group/OpenAI-YM-20261028/blob/3456e5e75cf8d4cd558955a5c7638fba2d68f464/YM/Next/RANK/Finite.lean\#L242-L250}{\texttt{Read.full\_cost}}
    (III.11);
  \item
    the lemma
    \href{https://github.com/yang-mills-group/OpenAI-YM-20261028/blob/3456e5e75cf8d4cd558955a5c7638fba2d68f464/YM/Work/A04.lean\#L34-L57}{\texttt{tendsto\_of\_fitted}}
    (VI.1): if \(f_k,g:E\to F\) are linear maps between
    finite-dimensional complex normed spaces, \(f_k\to g\) pointwise,
    \(g\) is injective and \(f_k(w_k)\to g(z)\), then \(w_k\to z\).
  \end{itemize}
\item
  \emph{For VII.1.7:} stage-length bounds, entry thresholds, backward
  windows, and the RG scale sequence (III.2); price panels and the
  linear checks on their exponents (II.2).
\end{itemize}

\textbf{Standard mathematics.}

\begin{itemize}
\tightlist
\item
  Gaussian moments and Wick's theorem in recursive form; Wick ordering
  as \(:\!p(\varphi)\!:\,=\mathbb E_\xi\,p(\varphi+i\xi)\). References:
  Glimm--Jaffe, \emph{Quantum Physics}, §6.3 and §8.6; B. Simon,
  \emph{The \(P(\varphi)_2\) Euclidean (Quantum) Field Theory}, Ch. I.
\item
  Fourier analysis on \((\mathbb Z/N)^4\); asymptotics of the lattice
  Green function and its finite differences (Lawler--Limic, \emph{Random
  Walk: A Modern Introduction}, Ch. 4); power counting for lattice
  Feynman integrals (T. Reisz, \emph{Commun. Math. Phys.} 116 (1988)
  81).
\item
  Riemann sums for continuous compactly supported functions; dominated
  convergence.
\item
  Derivatives of \(|x|^{2-n}\) against harmonic polynomials (Hobson's
  formula; E. W. Hobson, \emph{The Theory of Spherical and Ellipsoidal
  Harmonics}, 1931); the Kelvin transform (Axler--Bourdon--Ramey,
  \emph{Harmonic Function Theory}, Ch. 4); invariance of the source-free
  Maxwell equations under inversion (Bateman, Cunningham, 1910).
\end{itemize}

\textbf{Imported constants and assumptions.}

\begin{itemize}
\tightlist
\item
  VII.1.1--VII.1.5 use the stated gauge, label and regulator data and no
  assumptions on an interacting construction.
\item
  VII.1.6 imports the construction record, the portfolio estimates at
  step \(g(k)\), the error bound of III.11, which tends to \(0\), and
  the OPE work items.
\item
  VII.1.7 imports the RG scale sequence (stage lengths \(L_j\), the
  scale and \(\hat\beta_k\)) from III.2.
\end{itemize}

\subsubsection{3. Strategy}\label{strategy-25}

At order zero, the normalized perturbative expectation is Gaussian. Wick
ordering introduces an independent imaginary Gaussian for each
insertion, so the pair contraction is the torus stencil covariance minus
the infinite-lattice self-covariance at a common insertion. Fourier
analysis identifies the contractions at separated points with the
Maxwell kernel. At one insertion, the remaining finite-volume error
tends to zero because the period grows faster than every inverse power
of the spacing. The lattice moments and the continuum Wick moments
satisfy the same deletion recursion; induction on the number of
distinguishable legs gives VII.1.2. Compactness and dominated
convergence then give the tested limit VII.1.1.

For detection, inversion and translation turn a homogeneous Maxwell
solution of degree \(j\) into a field whose Taylor expansion at the
detector begins at order \(j\). Adding the corresponding lower-rank
probes gives a triangular change of variables. In an orthogonal shell
basis, the reflected one-leg pairing becomes diagonal with positive
entries. Wick's rule then makes distinct monomials orthogonal and gives
every monomial a positive squared norm, proving VII.1.4.

The interacting comparison uses the rank estimates of III.11. They imply
convergence of all mixed two-point readings in one portfolio. Free
detection makes the raw charts eventually invertible and forces their
filtration defects to tend to zero. Fixed equivariant retractions of the
free inclusions remove these defects recursively, giving the exactly
compatible charts of VII.1.6. Proposition VII.1.7 supplies the
backward-window and local-expansion estimates used in this construction.

\subsubsection{4. Steps}\label{steps-25}

\paragraph{4.1 The free lattice
covariance}\label{the-free-lattice-covariance}

\textbf{Lemma VII.1.a (covariance of link forms).} For links
\(a=(x,i)\), \(b=(y,j)\) and colours \(u,t\),
\[\mathrm{Cov}_{\mu_N}\big(\langle A(a),u\rangle,\langle A(b),t\rangle\big)=\langle u,t\rangle\,\mathbf 1_{i=j}\,\big[g_N(y-x)-g_N(y+x+e_i)\big].\]
Here \(g_N\) is the mean-zero torus Green function,
\(-\Delta g_N=\delta_0-N^{-4}\) and \(\sum_xg_N=0\), written as a
Fourier sum over \((\mathbb Z/N)^4\) with the zero mode omitted
(\href{https://github.com/yang-mills-group/OpenAI-YM-20261028/blob/3456e5e75cf8d4cd558955a5c7638fba2d68f464/YM/Pert/Fourier/Inversion.lean\#L19}{\texttt{toroid}},
\href{https://github.com/yang-mills-group/OpenAI-YM-20261028/blob/3456e5e75cf8d4cd558955a5c7638fba2d68f464/YM/Pert/Fourier/Inversion.lean\#L46-L52}{\texttt{lap\_eq}}).
\emph{Proof sketch.} Summation by parts gives
\(\mathsf Q_N(A,B)=-\tfrac12\sum\langle A(x,i),\Delta B_i(x)\rangle\),
the Feynman-gauge identity curl² + div² = \(-A\cdot\Delta A\)
(\href{https://github.com/yang-mills-group/OpenAI-YM-20261028/blob/3456e5e75cf8d4cd558955a5c7638fba2d68f464/YM/Pert/Fourier/Adjoint.lean\#L62-L71}{\texttt{total}}).
The response to a unit forcing at \(a\) is \(g_N(\cdot-x)e_i\) minus the
same function at the image link. The constants \(N^{-4}\) cancel between
the two images, and oddness of the test field doubles the pairing
(\href{https://github.com/yang-mills-group/OpenAI-YM-20261028/blob/3456e5e75cf8d4cd558955a5c7638fba2d68f464/YM/Pert/Fourier/Source.lean\#L40-L54}{\texttt{inverse\_eq}}).
\emph{Lean:}
\href{https://github.com/yang-mills-group/OpenAI-YM-20261028/blob/3456e5e75cf8d4cd558955a5c7638fba2d68f464/YM/Pert/Fourier/Source.lean\#L55-L59}{\texttt{Orb.covar\_link}}.

\textbf{Lemma VII.1.b (separated legs).} Let \((N_i,h_i)\) be an
itinerary. For a torus site \(x\), write \([x]_{N}\) for its centred
lift in \((-N/2,N/2]^4\). Let \(x_i,y_i\) satisfy
\(h_i[x_i-y_i+p]_{N_i}\to z\ne0\) and
\(\|h_i[x_i+y_i+p]_{N_i}\|\to\infty\) for every fixed
\(p\in\mathbb Z^4\). Then for real linear forms \(P\) on \(T_j\) and
\(Q\) on \(T_k\),
\[\mathrm{Cov}_{\mu_{N_i}}\big(P\circ\ell^{h_i}_{x_i},Q\circ\ell^{h_i}_{y_i}\big)\to K(P\circ\kappa,Q\circ\kappa)(z).\]
The image condition holds automatically when both families are homing
and \(4\mid N_i\), since then each coordinate of \(x_i+y_i\) is
\(N_i/2+O(h_i^{-1})\) modulo \(N_i\)
(\href{https://github.com/yang-mills-group/OpenAI-YM-20261028/blob/3456e5e75cf8d4cd558955a5c7638fba2d68f464/YM/Pert/Align/Schedule.lean\#L90-L103}{\texttt{finale\_home}}).
\emph{Proof sketch.}

\begin{enumerate}
\def\labelenumi{\arabic{enumi}.}
\tightlist
\item
  Average over \(W_4\) and expand the forms in component functionals. By
  Lemma VII.1.a each component pairing becomes \(h^{-(|D|+2)}\) times a
  forward difference \(\nabla^+_D\) of \(g_N\), evaluated at \(x-y\)
  (direct term) and at \(x+y+e\) (image term). The Fourier symbol of
  \(\nabla^+_D\) is \(\prod_{k\in D}(e^{ip_k}-1)\).
\item
  \emph{Direct term}
  (\href{https://github.com/yang-mills-group/OpenAI-YM-20261028/blob/3456e5e75cf8d4cd558955a5c7638fba2d68f464/YM/Pert/Fourier/Return.lean\#L95-L101}{\texttt{joint\_casual}}).
  At physical momentum scale \(s\), set \(t=hs\) and write lattice
  momentum as \(tk\). The unscaled Fourier multiplier becomes
  \(N^{-4}t^{|D|-2}\) times a rescaled symbol; after multiplication by
  \(h^{-(|D|+2)}\), the prefactor is \((Nh)^{-4}s^{|D|-2}\). The
  rescaled symbol is jointly smooth near \((0,k)\) for \(k\ne0\) and
  equals \(\prod_{l\in D}(ik_l)/|k|^2\) at \(t=0\). Split momentum space
  with a smooth bump into a low-momentum piece, dyadic annuli and a
  tail.

  \begin{itemize}
  \tightlist
  \item
    Each annulus is a Riemann sum of a compactly supported rescaled
    integrand and converges to its continuum integral.
  \item
    For separations bounded below by \(R>0\), the tail is bounded by
    \(C(D,R)/s\) uniformly in \(N\) once the mesh is fine enough
    relative to \(hs\)
    (\href{https://github.com/yang-mills-group/OpenAI-YM-20261028/blob/3456e5e75cf8d4cd558955a5c7638fba2d68f464/YM/Pert/Fourier/Power.lean\#L63-L70}{\texttt{upper\_tail}}).
    The proof uses \(|D|+3\) summations by parts in momentum and the
    lower bound \(|e^{2\pi i[x]_f/N}-1|\gtrsim\|h[x]_N\|/(Nh)\) for a
    suitable coordinate \(f\).
  \item
    The low piece is \(O(s^{|D|+2})\).
  \item
    The same bounds hold for the continuum integral after heat
    regularization. A \(3\varepsilon\) argument concludes
    (\href{https://github.com/yang-mills-group/OpenAI-YM-20261028/blob/3456e5e75cf8d4cd558955a5c7638fba2d68f464/YM/Pert/Fourier/Return.lean\#L56-L81}{\texttt{wrap\_tube}}).
  \item
    Besides \(h\to0\), this part requires only \(Nh\to\infty\), rather
    than faster volume growth.
  \end{itemize}
\item
  \emph{Image term.} The same decomposition gives \(0\) when the norm of
  the centred image displacement tends to infinity in physical units
  (\href{https://github.com/yang-mills-group/OpenAI-YM-20261028/blob/3456e5e75cf8d4cd558955a5c7638fba2d68f464/YM/Pert/Fourier/Image.lean\#L22-L31}{\texttt{joint\_far}}).
\item
  The limit is identified with \(K\) on rotated on-shell jets, and
  \(W_4\)-averaging returns \(K(P\circ\kappa,Q\circ\kappa)\).
\end{enumerate}

\emph{Lean:}
\href{https://github.com/yang-mills-group/OpenAI-YM-20261028/blob/3456e5e75cf8d4cd558955a5c7638fba2d68f464/YM/Pert/Align/Bilateral.lean\#L75-L83}{\texttt{Align.chart\_limit}},
\href{https://github.com/yang-mills-group/OpenAI-YM-20261028/blob/3456e5e75cf8d4cd558955a5c7638fba2d68f464/YM/Pert/Align/Cancel.lean\#L80-L88}{\texttt{Align.orbit\_limit}}.
\emph{(proof read through the Fourier bounds; remaining analytic inputs
in §7)}

\textbf{Lemma VII.1.c (coincident legs).} Let \((N_i,h_i)\) be an
itinerary and let \(x_i\in\Lambda_{N_i}\) satisfy
\(\|h_i[2x_i+p]_{N_i}\|\to\infty\) for every fixed \(p\in\mathbb Z^4\).
For real linear forms \(P\) on \(T_j\) and \(Q\) on \(T_k\),
\[\mathrm{Cov}_{\mu_{N_i}}\big(P\circ\ell^{h_i}_{x_i},Q\circ\ell^{h_i}_{x_i}\big)-\mathcal S_{h_i}(P,Q)\to0,\]
where \(\mathcal S_h\) is the \(\mathbb Z^4\) covariance of the same
stencils. In particular \(S_h(a,b)=\mathcal S_h(x_a,x_b)\)
(\href{https://github.com/yang-mills-group/OpenAI-YM-20261028/blob/3456e5e75cf8d4cd558955a5c7638fba2d68f464/YM/Pert/Align/Cancel.lean\#L108-L113}{\texttt{shell\_exact}}).
\emph{Proof sketch.} The image term vanishes as in Lemma VII.1.b. The
direct term is a finite combination of values \(g_N(a)\) at fixed
lattice offsets \(a\), compared with the \(\mathbb Z^4\) Green function
\(G_{\mathbb Z^4}(a)\) and multiplied by \(h^{-(j+k+4)}\). The estimate
for each fixed offset \(a\)
(\href{https://github.com/yang-mills-group/OpenAI-YM-20261028/blob/3456e5e75cf8d4cd558955a5c7638fba2d68f464/YM/Pert/Fourier/Fee.lean\#L84-L112}{\texttt{toroid\_cost}})
gives constants \(C,A\ge0\), independent of \(N\) and \(r\), such that
for \(N\) even and \(0<r\le1\), \(8/N\le r\),
\[|g_N(a)-G_{\mathbb Z^4}(a)|\le C\,r^2+A/(N r^4).\] It is proved by
viewing \(g_N(a)\) as a Riemann sum of
\(\cos(2\pi a\cdot q)/\hat\Delta(2\pi q)\):

\begin{itemize}
\tightlist
\item
  near \(q=0\) both terms are \(O(|q|^{-2})\), whose integral over
  \(\{|q|\le r\}\) is \(O(r^2)\);
\item
  elsewhere the integrand is \(O(r^{-4})\)-Lipschitz and the grid points
  move by \(O(1/N)\).
\end{itemize}

With \(r=h^{n+1}\), where \(n=j+k+4\) is the power of \(h^{-1}\) to be
beaten, the error times \(h^{-n}\) is at most
\(Ch^{n+2}+A/(Nh^{5n+4})\to0\). This comparison uses the condition
\(N_ih_i^k\to\infty\) for every \(k\)
(\href{https://github.com/yang-mills-group/OpenAI-YM-20261028/blob/3456e5e75cf8d4cd558955a5c7638fba2d68f464/YM/Pert/Fourier/Visits.lean\#L29-L46}{\texttt{strike\_two}}).
\emph{Lean:}
\href{https://github.com/yang-mills-group/OpenAI-YM-20261028/blob/3456e5e75cf8d4cd558955a5c7638fba2d68f464/YM/Pert/Align/Bilateral.lean\#L85-L89}{\texttt{Align.chart\_return}},
\href{https://github.com/yang-mills-group/OpenAI-YM-20261028/blob/3456e5e75cf8d4cd558955a5c7638fba2d68f464/YM/Pert/Align/Cancel.lean\#L89-L95}{\texttt{Align.orbit\_contact}}.

\textbf{Lemma VII.1.d (VII.1.3).} This is Lemma VII.1.b for \(j\ne l\)
(with
\href{https://github.com/yang-mills-group/OpenAI-YM-20261028/blob/3456e5e75cf8d4cd558955a5c7638fba2d68f464/YM/Pert/Align/Schedule.lean\#L71-L77}{\texttt{scene\_home}}
and
\href{https://github.com/yang-mills-group/OpenAI-YM-20261028/blob/3456e5e75cf8d4cd558955a5c7638fba2d68f464/YM/Pert/Align/Schedule.lean\#L90-L103}{\texttt{finale\_home}})
and Lemma VII.1.c for \(j=l\). The covariance matrix
\(\mathsf Q_N^{-1}\) evaluated on the stencil vectors equals the orbital
pairing
(\href{https://github.com/yang-mills-group/OpenAI-YM-20261028/blob/3456e5e75cf8d4cd558955a5c7638fba2d68f464/YM/Pert/Align/Drivers.lean\#L31}{\texttt{key\_orb}}).
\emph{Lean:}
\href{https://github.com/yang-mills-group/OpenAI-YM-20261028/blob/3456e5e75cf8d4cd558955a5c7638fba2d68f464/YM/Pert/Align/Drivers.lean\#L39-L44}{\texttt{corrected\_limit}}.

\paragraph{4.2 Order zero as a Gaussian
recursion}\label{order-zero-as-a-gaussian-recursion}

\textbf{Lemma VII.1.e (order zero is a joint Gaussian expectation).} Fix
\(h>0\), an even period \(N\ge4\), insertion sites \(o_1,\ldots,o_r\),
and \(d\in\mathbb N\). Let \(s\) be a finite set of distinguishable
occurrences, with insertion map \(\tau:s\to\{1,\ldots,r\}\) and key map
\(a:s\to\{\text{keys of rank at most }d\}\). Repeated keys are allowed.
Set \(\mathrm{root}(s)_j=\prod_{\nu:\tau(\nu)=j}x_{a(\nu)}\). Then
\[\Big\langle\prod_jR_h(\mathrm{root}(s)_j,o_j)\Big\rangle^{(0)}=E_h(s):=\mathbb E\Big[\prod_{\nu\in s}\big(\ell^h_{o_{\tau(\nu)},a(\nu)}+i\,\xi^{(\tau(\nu))}_{a(\nu)}\big)\Big],\]
under \(\mu_N\otimes\bigotimes_{j=1}^r\mathcal N(0,S_h)\). \emph{Proof
sketch.} The weight and the vacuum both have constant coefficient \(1\),
so the order-zero coefficient of the normalized expectation is a
\(\mu_N\)-integral. The constant coefficient of each lattice jet is its
linear part \(\ell^h\). Fubini over the independent auxiliary copies
finishes the proof. \emph{Lean:}
\href{https://github.com/yang-mills-group/OpenAI-YM-20261028/blob/3456e5e75cf8d4cd558955a5c7638fba2d68f464/YM/Pert/Warp/Initial.lean\#L39-L56}{\texttt{Warp.amplitude\_zero}}.
The pair covariance of two legs is \(\kappa_h\)
(\href{https://github.com/yang-mills-group/OpenAI-YM-20261028/blob/3456e5e75cf8d4cd558955a5c7638fba2d68f464/YM/Pert/Warp/Engine.lean\#L59-L63}{\texttt{node\_wage}}):
the auxiliary copies contribute \(i^2S_h=-S_h\) for legs assigned to the
same insertion and \(0\) otherwise, even if two insertion sites happen
to coincide.

\textbf{Lemma VII.1.f (symmetric deletion identity).} Let \(\mu\) be a
centred Gaussian on a finite-dimensional real space (possibly
degenerate) and \(P_a\), \(a\in s\), real linear functionals with values
in \(\mathbb C\). Then
\[|s|\,\mathbb E\prod_{a\in s}P_a=\sum_{a\in s}\sum_{b\in s\setminus\{a\}}\mathbb E[P_aP_b]\;\mathbb E\prod_{c\in s\setminus\{a,b\}}P_c.\]
The products are bilinear, without complex conjugation. For the legs of
Lemma VII.1.e, write
\(\kappa_h(\nu,\mu)=\kappa_h(\tau(\nu),a(\nu);\tau(\mu),a(\mu))\); then
\(|s|E_h(s)=\sum_{\nu\ne\mu}\kappa_h(\nu,\mu)E_h(s\setminus\{\nu,\mu\})\).
\emph{Proof sketch.} For powers of a single real functional \(L\) of
variance \(v\), \(\mathbb E L^{n+1}=nv\,\mathbb EL^{n-1}\) for
\(n\ge1\), with \(\mathbb EL=0\), read off from the moment generating
function
(\href{https://github.com/yang-mills-group/OpenAI-YM-20261028/blob/3456e5e75cf8d4cd558955a5c7638fba2d68f464/YM/Pert/Warp/Diagonal.lean\#L15-L19}{\texttt{ev\_next}}).
The multilinear case follows by polarization
(\href{https://github.com/yang-mills-group/OpenAI-YM-20261028/blob/3456e5e75cf8d4cd558955a5c7638fba2d68f464/YM/Pert/Warp/Symmetric.lean\#L73-L77}{\texttt{packet\_cancel}}).
The complex case follows by splitting into real and imaginary parts.
\emph{Lean:}
\href{https://github.com/yang-mills-group/OpenAI-YM-20261028/blob/3456e5e75cf8d4cd558955a5c7638fba2d68f464/YM/Pert/Warp/Exchange.lean\#L68-L91}{\texttt{exchange\_s}},
\href{https://github.com/yang-mills-group/OpenAI-YM-20261028/blob/3456e5e75cf8d4cd558955a5c7638fba2d68f464/YM/Pert/Warp/Approach.lean\#L12-L18}{\texttt{Warp.eq\_cancel}}.

\textbf{Lemma VII.1.g (the continuum recursion).} For
\(z\in\mathrm{Conf}_r\) set
\(c_z(s)=\mathcal W_z(\rho\circ\mathrm{root}(s))\). Then for every
\(\nu\in s\),
\[c_z(s)=\sum_{\mu\in s\setminus\{\nu\}}\mathbf 1_{\tau(\mu)\ne\tau(\nu)}\,K_{a(\nu),a(\mu)}\big(z_{\tau(\nu)}-z_{\tau(\mu)}\big)\,c_z(s\setminus\{\nu,\mu\}),\]
and summing over \(\nu\in s\) gives the analogue of Lemma VII.1.f.
\emph{Proof sketch.} Removing one on-shell variable \(\rho(x_{a(\nu)})\)
at the point \(z_{\tau(\nu)}\) from a Wick moment amounts to applying,
at every other point, the derivation that replaces one variable by its
pairing with \(a(\nu)\)
(\href{https://github.com/yang-mills-group/OpenAI-YM-20261028/blob/3456e5e75cf8d4cd558955a5c7638fba2d68f464/YM/Pert/Warp/Rowdy.lean\#L45-L49}{\texttt{take\_line}}).
On a monomial the derivation produces the sum over its variables
(\href{https://github.com/yang-mills-group/OpenAI-YM-20261028/blob/3456e5e75cf8d4cd558955a5c7638fba2d68f464/YM/Pert/Warp/Rowdy.lean\#L59-L61}{\texttt{D\_kill}}).
There are no same-point pairings. \emph{Lean:}
\href{https://github.com/yang-mills-group/OpenAI-YM-20261028/blob/3456e5e75cf8d4cd558955a5c7638fba2d68f464/YM/Pert/Warp/Shards.lean\#L26-L48}{\texttt{cut\_one}},
\href{https://github.com/yang-mills-group/OpenAI-YM-20261028/blob/3456e5e75cf8d4cd558955a5c7638fba2d68f464/YM/Pert/Warp/Shards.lean\#L49}{\texttt{cut\_cancel}}.

\textbf{Lemma VII.1.h (convergence of leg products).} In the setting of
VII.1.2, \(E_{h_i}(s)\to c_z(s)\) for every finite set of legs.
\emph{Proof sketch.} The proof is by strong induction on \(|s|\).

\begin{itemize}
\tightlist
\item
  For \(s=\emptyset\) both sides equal \(1\) (probability measures).
\item
  Otherwise divide the identities of Lemmas VII.1.f and VII.1.g by
  \(|s|\). Use VII.1.3 for the coefficients and the induction hypothesis
  for the shorter products.
\end{itemize}

\emph{Lean:}
\href{https://github.com/yang-mills-group/OpenAI-YM-20261028/blob/3456e5e75cf8d4cd558955a5c7638fba2d68f464/YM/Pert/Warp/Approach.lean\#L19-L32}{\texttt{Warp.earth}}.

\textbf{Lemma VII.1.i (monomials, polynomials, labels; VII.1.2).}

\begin{itemize}
\tightlist
\item
  Every tuple of monomials involving only keys of rank at most \(d\) is
  \(\mathrm{root}(s)\) for a suitable set of distinguishable
  occurrences, so Lemma VII.1.h applies.
\item
  Multilinearity extends the limit to arbitrary polynomials involving
  only keys of rank at most \(d\); no bound on their number of factors
  is needed.
\item
  The representatives \(\mathrm{sec}_d\) have engineering weight at most
  \(d\), hence use only keys of rank at most \(d\). Composing gives the
  label limit, with \(\rho\circ\mathrm{sec}_d=\omega_d\).
\end{itemize}

\emph{Lean:}
\href{https://github.com/yang-mills-group/OpenAI-YM-20261028/blob/3456e5e75cf8d4cd558955a5c7638fba2d68f464/YM/Pert/Warp/Horizons.lean\#L31-L40}{\texttt{mono\_single}},
\href{https://github.com/yang-mills-group/OpenAI-YM-20261028/blob/3456e5e75cf8d4cd558955a5c7638fba2d68f464/YM/Pert/Warp/Horizons.lean\#L41-L47}{\texttt{mono\_limit}},
\href{https://github.com/yang-mills-group/OpenAI-YM-20261028/blob/3456e5e75cf8d4cd558955a5c7638fba2d68f464/YM/Pert/Warp/Horizons.lean\#L50-L63}{\texttt{poly\_limit}},
\href{https://github.com/yang-mills-group/OpenAI-YM-20261028/blob/3456e5e75cf8d4cd558955a5c7638fba2d68f464/YM/Pert/Warp/Horizons.lean\#L65-L68}{\texttt{row\_limit}}.

\paragraph{4.3 Tested sums}\label{tested-sums}

\textbf{Lemma VII.1.j (Riemann sums with continuous convergence).}
Hypotheses:

\begin{itemize}
\tightlist
\item
  centred torus grids of positive integer side \(N_m\) and spacing
  \(t_m=1/n_m\), \(n_m\in\mathbb N_{>0}\), with \(t_m\to0\) and
  \(t_mN_m\to\infty\);
\item
  kernels \(\mathcal K_m\) on \(r\)-tuples of grid points;
\item
  \(f\) continuous with compact support in an open set \(\omega\);
\item
  along every cofinal reindexing \(\sigma:\mathbb N\to\mathbb N\), and
  every sequence of tuples on the reindexed grids whose embeddings
  converge to \(z\in\omega\), the corresponding kernels converge to
  \(\mathcal G(z)\).
\end{itemize}

Writing \(\jmath_m(o)=t_m[o-c_{N_m}]_{N_m}\) for the embedding, one has
\(t_m^{4r}\sum_o f(\jmath_m(o))\mathcal K_m(o)\to\int f\mathcal G\).
\emph{Proof sketch.}

\begin{itemize}
\tightlist
\item
  The kernels are eventually bounded uniformly on \(\mathrm{supp}\,f\)
  (\href{https://github.com/yang-mills-group/OpenAI-YM-20261028/blob/3456e5e75cf8d4cd558955a5c7638fba2d68f464/YM/Pert/Rig/Moving.lean\#L68-L86}{\texttt{bounded\_fee}}).
  Otherwise there are points of \(\mathrm{supp}\,f\) along which they
  blow up, and a convergent subsequence of these points contradicts the
  hypothesis.
\item
  Write the sum as the integral of a step function supported in a fixed
  ball. Apply dominated convergence; the box eventually contains the
  support.
\end{itemize}

\emph{Lean:}
\href{https://github.com/yang-mills-group/OpenAI-YM-20261028/blob/3456e5e75cf8d4cd558955a5c7638fba2d68f464/YM/Pert/Rig/Tested.lean\#L23-L52}{\texttt{Rig.grid\_limit}},
with
\href{https://github.com/yang-mills-group/OpenAI-YM-20261028/blob/3456e5e75cf8d4cd558955a5c7638fba2d68f464/YM/Pert/Rig/Moving.lean\#L50-L51}{\texttt{Arrives}}.

\textbf{Lemma VII.1.k (VII.1.1).} Apply Lemma VII.1.j with
\(\omega=\mathrm{Conf}_r\), \(\mathcal K_m\) the order-zero correlations
and \(\mathcal G=\mathcal W(\omega_dp)\), along the regulator sequence
\(h_m=L^{-m}\), \(N_m=4L^{m^2+m+1}\)
(\href{https://github.com/yang-mills-group/OpenAI-YM-20261028/blob/3456e5e75cf8d4cd558955a5c7638fba2d68f464/YM/Pert/Rig/Linked.lean\#L14-L21}{\texttt{routine}}).

\begin{itemize}
\tightlist
\item
  The arrival property is VII.1.2, applied to each cofinal reindexing of
  the regulator sequence, which is again an itinerary with \(4\mid N\).
\item
  The tested sum is the order-zero coefficient of the profile.
\item
  The identity \(\rho\circ\mathrm{sec}_d=\omega_d\) identifies the
  continuum moment with the free hierarchy.
\item
  The board form applies this slot by slot to basis labels; clean
  profiles are smooth test functions with compact support in
  \(\mathrm{Conf}\)
  (\href{https://github.com/yang-mills-group/OpenAI-YM-20261028/blob/3456e5e75cf8d4cd558955a5c7638fba2d68f464/YM/Euclid/Test.lean\#L18}{\texttt{Test}}).
\end{itemize}

\emph{Lean:}
\href{https://github.com/yang-mills-group/OpenAI-YM-20261028/blob/3456e5e75cf8d4cd558955a5c7638fba2d68f464/YM/Pert/Rig/Linked.lean\#L59-L69}{\texttt{Rig.zerothProfile}},
\href{https://github.com/yang-mills-group/OpenAI-YM-20261028/blob/3456e5e75cf8d4cd558955a5c7638fba2d68f464/YM/Pert/Position/Crest.lean\#L13-L16}{\texttt{Position.finite\_leading}}.

\textbf{Downstream use (not part of this unit).} The assembly in VIII.1
uses the tested order-zero limit in two ways:

\begin{itemize}
\tightlist
\item
  \href{https://github.com/yang-mills-group/OpenAI-YM-20261028/blob/3456e5e75cf8d4cd558955a5c7638fba2d68f464/YM/Pert/Position/Crest.lean\#L50-L52}{\texttt{credible}}
  subtracts it from
  \href{https://github.com/yang-mills-group/OpenAI-YM-20261028/blob/3456e5e75cf8d4cd558955a5c7638fba2d68f464/YM/Pert/Position/Actual.lean\#L16-L19}{\texttt{measured\_tends}}
  (which assumes the interacting short-distance comparison of VI.2).
  This shows that measured scaled means minus order-zero lattice
  readings tend to \(0\).
\item
  \href{https://github.com/yang-mills-group/OpenAI-YM-20261028/blob/3456e5e75cf8d4cd558955a5c7638fba2d68f464/YM/Pert/Position/Crest.lean\#L36-L43}{\texttt{calibrate}}
  solves the finite position equations when the leading Jacobian is
  invertible.
\end{itemize}

The position-comparison data are described in VII.3.

\paragraph{4.4 Detection on the free
side}\label{detection-on-the-free-side}

\textbf{Lemma VII.1.l (Wick rule under linear substitution).} Fix
\(x\in\mathrm{Conf}_r\) and a finite list of variable occurrences
assigned to its insertions. Replace each occurrence \(a\) by a finite
complex linear combination \(u_a\) of shell coordinates, allowing
different Taylor ranks. Extend \(K\) to these combinations by complex
bilinearity. The Wick moment of the resulting products equals
\[\sum_{\text{matchings}}\prod_{\text{pairs}}K\big(u_a,u_{a'}\big)(x_{i(a)}-x_{i(a')}).\]
Only matchings between distinct insertions are included. \emph{Proof
sketch.} Expand the linear forms in the coordinate basis. Multilinearity
reduces the formula to the defining Wick sum on coordinate monomials.
\emph{Lean:}
\href{https://github.com/yang-mills-group/OpenAI-YM-20261028/blob/3456e5e75cf8d4cd558955a5c7638fba2d68f464/YM/Push/FINITE_DETECTION/Polarize.lean\#L108-L119}{\texttt{Free.linear\_substitution}}.

\textbf{Lemma VII.1.m (diagonal one-leg kernel).} With the bilinear
extension of \(K\) from Lemma VII.1.l, choose in each \(\mathrm{Sh}_j\)
a basis \((Y_a)\) orthogonal for full component contraction, and set
\(\lambda_a=\gamma_j\langle Y_a,Y_a\rangle>0\). For
\(Y\in\mathrm{Sh}_j\), let \(Y_i\) be the order-\(i\) Taylor coefficient
of \(Y^\natural\) at \(e_{\mathrm t}\) and let
\(\langle Y_i,\cdot\rangle\) denote the full-contraction probe. Define
\[\epsilon_a=\sum_{i=0}^j\frac{\gamma_j}{\gamma_i}\langle (Y_a)_i,\cdot\rangle,\qquad \epsilon^\sigma_a=\epsilon_a\circ\sigma.\]
Then
\[K(\epsilon_a,\epsilon^\sigma_b)(e_{\mathrm t})=\delta_{ab}\lambda_a.\]
\emph{Proof sketch.} Expand the translated field
\(x\mapsto Y^\natural(e_{\mathrm t}+x)\) into its homogeneous Taylor
pieces and apply the Kelvin formula to each piece. Since
\(e_{\mathrm t}+\iota(e_{\mathrm t})=0\), its inverted field has no
Taylor coefficient at \(e_{\mathrm t}\) below order \(j\); its
order-\(j\) coefficient is \((-1)^j\) times the time-reflected tensor
\(Y\). The weighted sum defining \(\epsilon_a\) produces exactly
\(\gamma_j\) times this inverted translated field. The second-endpoint
sign cancels \((-1)^j\) on the diagonal. Thus pairings with lower ranks
vanish and equal-rank pairings are \(\gamma_j\) times full contraction.
Symmetry gives vanishing for the opposite rank ordering, and the chosen
orthogonal bases give the displayed matrix.

\emph{Lean:}
\href{https://github.com/yang-mills-group/OpenAI-YM-20261028/blob/3456e5e75cf8d4cd558955a5c7638fba2d68f464/YM/Push/FINITE_DETECTION/Variables.lean\#L114-L118}{\texttt{Free.pairing\_var}},
with
\href{https://github.com/yang-mills-group/OpenAI-YM-20261028/blob/3456e5e75cf8d4cd558955a5c7638fba2d68f464/YM/Push/FINITE_DETECTION/Centered.lean\#L70-L73}{\texttt{center\_lower}}
and
\href{https://github.com/yang-mills-group/OpenAI-YM-20261028/blob/3456e5e75cf8d4cd558955a5c7638fba2d68f464/YM/Push/FINITE_DETECTION/Centered.lean\#L80-L101}{\texttt{center\_diag}}.

\textbf{Lemma VII.1.n (orthogonality and positivity; VII.1.4).} For
finitely supported exponent vectors \(\alpha,\beta\) in the probe
variables,
\[\mathcal W_{(e_{\mathrm t},0)}(\Phi x^\alpha,\Psi x^\beta)=\delta_{\alpha\beta}s_\alpha,\qquad s_\alpha=\prod_a\alpha_a!\lambda_a^{\alpha_a}>0.\]
The empty monomial has \(s_0=1\). \emph{Proof sketch.}

\begin{itemize}
\tightlist
\item
  By Lemmas VII.1.l and VII.1.m only matchings that pair each variable
  at \(e_{\mathrm t}\) with the same variable at \(0\) survive.
\item
  For equal exponent vectors, there are \(\prod_a\alpha_a!\) surviving
  matchings, each of amplitude \(\prod_a\lambda_a^{\alpha_a}\).
\item
  For \(q\ne0\), expanding \(q=\sum c_\alpha x^\alpha\) gives
  \(\mathrm{Re}=\sum_\alpha|c_\alpha|^2s_\alpha>0\).
\item
  For a label \(v\):

  \begin{itemize}
  \tightlist
  \item
    \(\Phi\) is onto
    (\href{https://github.com/yang-mills-group/OpenAI-YM-20261028/blob/3456e5e75cf8d4cd558955a5c7638fba2d68f464/YM/Push/FINITE_DETECTION/Variables.lean\#L137-L142}{\texttt{forward\_full}}),
    so \(\rho v=\Phi q\) for some \(q\);
  \item
    \(\rho(\sigma\bar v)=\Psi\bar q\)
    (\href{https://github.com/yang-mills-group/OpenAI-YM-20261028/blob/3456e5e75cf8d4cd558955a5c7638fba2d68f464/YM/Push/FINITE_DETECTION/Variables.lean\#L165-L173}{\texttt{reversed\_observable}});
  \item
    \(q\ne0\) because \(v\notin\mathcal I_d=V_d\cap\ker\rho\).
  \end{itemize}
\end{itemize}

\emph{Lean:}
\href{https://github.com/yang-mills-group/OpenAI-YM-20261028/blob/3456e5e75cf8d4cd558955a5c7638fba2d68f464/YM/Push/FINITE_DETECTION/TwoPoint.lean\#L57-L79}{\texttt{pair\_test}},
\href{https://github.com/yang-mills-group/OpenAI-YM-20261028/blob/3456e5e75cf8d4cd558955a5c7638fba2d68f464/YM/Push/FINITE_DETECTION/TwoPoint.lean\#L109-L115}{\texttt{wick\_orth}},
\href{https://github.com/yang-mills-group/OpenAI-YM-20261028/blob/3456e5e75cf8d4cd558955a5c7638fba2d68f464/YM/Push/FINITE_DETECTION/TwoPoint.lean\#L132-L149}{\texttt{positive}},
\href{https://github.com/yang-mills-group/OpenAI-YM-20261028/blob/3456e5e75cf8d4cd558955a5c7638fba2d68f464/YM/Push/FINITE_DETECTION/TwoPoint.lean\#L151-L155}{\texttt{two\_detect}}.

\textbf{Lemma VII.1.o (Kelvin formulas; VII.1.5).} \emph{Proof sketch
for (a).} The case \(j=0\) is the definition of \(\Gamma\). For
\(j\ge1\), use
\[\tfrac1{j!}A(\partial)=\tfrac{1}{j}\sum_i\partial_i\,\tfrac1{(j-1)!}A(e_i,\cdot)(\partial).\]
Each contraction \(A(e_i,\cdot)\) is again symmetric with harmonic
diagonal. Then use \(\partial_i|x|^{-2j}=-2j\,x_i|x|^{-2(j+1)}\) and the
vanishing of \(\sum_i\partial_iA(e_i,x,\dots,x)\) (harmonicity). The
contraction identity \(\sum_ix_iA(e_i,x,\dots)=A(x,\dots,x)\) closes the
step. \emph{Proof sketch for (b).} Work colour by colour. Each component
\(Y_{ab}^\natural\) is harmonic (closed and coclosed), so (a) applies to
every term of the curvature kernel. The terms are recombined using the
divergence equation and
\(\partial_k(x\lrcorner Y^\natural)_m-\partial_m(x\lrcorner Y^\natural)_k=(j+2)Y^\natural_{km}\)
(Bianchi plus Euler). \emph{Lean:}
\href{https://github.com/yang-mills-group/OpenAI-YM-20261028/blob/3456e5e75cf8d4cd558955a5c7638fba2d68f464/YM/Push/FINITE_DETECTION/Harmonic.lean\#L157-L180}{\texttt{harmonic\_Green}},
\href{https://github.com/yang-mills-group/OpenAI-YM-20261028/blob/3456e5e75cf8d4cd558955a5c7638fba2d68f464/YM/Push/FINITE_DETECTION/Response.lean\#L172-L176}{\texttt{response\_eq}}.

\textbf{Lemma VII.1.p (finite detection).} Let \(\beta\) be a basis of
\(F_d\). The map
\(v\mapsto\big(\mathcal W_{(e_{\mathrm t},0)}(\omega_dv,\omega_d\beta_i)\big)_i\)
is injective on \(F_d\). A label whose two-point functions with all
labels vanish at the detector pair is null in the free table.
\emph{Proof sketch.} If all basis probes vanish, linearity gives
vanishing against every second label. For a representative of \(v\),
take its time-reflected conjugate as the second label and apply VII.1.4.
Conversely, \(\rho v=0\) makes every free correlation containing that
label vanish. Thus the free null space in \(V_d\) is \(\mathcal I_d\).
\emph{Lean:}
\href{https://github.com/yang-mills-group/OpenAI-YM-20261028/blob/3456e5e75cf8d4cd558955a5c7638fba2d68f464/YM/Push/FINITE_DETECTION/Test.lean\#L21-L32}{\texttt{free\_criterion}},
\href{https://github.com/yang-mills-group/OpenAI-YM-20261028/blob/3456e5e75cf8d4cd558955a5c7638fba2d68f464/YM/Push/FINITE_DETECTION/Test.lean\#L34-L42}{\texttt{free\_null}},
\href{https://github.com/yang-mills-group/OpenAI-YM-20261028/blob/3456e5e75cf8d4cd558955a5c7638fba2d68f464/YM/Push/FINITE_DETECTION/Test.lean\#L77-L85}{\texttt{free\_injective}}.
In the same spirit,
\href{https://github.com/yang-mills-group/OpenAI-YM-20261028/blob/3456e5e75cf8d4cd558955a5c7638fba2d68f464/YM/OPE/Levy/Panel.lean\#L39-L51}{\texttt{OPE.Levy.free\_positive}}
shows that the reflected free two-point form at the detector pair is
positive on \(F_d\setminus0\).

For bounded complex linear maps \(P:E\to E\), \(X:V\to E\), \(A:E\to Z\)
and \(H:V\to Z\) between complex normed spaces, the perturbation lemma
\href{https://github.com/yang-mills-group/OpenAI-YM-20261028/blob/3456e5e75cf8d4cd558955a5c7638fba2d68f464/YM/Push/FINITE_DETECTION/Matrix.lean\#L38-L46}{\texttt{test\_projected}}
is not used by the main theorem. It states that a lower bound
\(a\|x\|\le\|Hx\|\) for every \(x\in V\) and
\(\|H-AX\|+\|(A-AP)X\|\le\delta\) imply \((a-\delta)\|x\|\le\|APXx\|\).

\paragraph{4.5 Corrected charts and
bookkeeping}\label{corrected-charts-and-bookkeeping}

\textbf{Lemma VII.1.q (two-point convergence at moving detectors).}
Under the hypotheses of VII.1.6, let \(\delta_k=r_k y_k\) be the
portfolio detector pair, where \(r_k>0\), \(r_k\to0\), and
\(y_k\to(e_{\mathrm t},0)\). For fixed reference labels \(v_0,v_1\) of
respective weights at most \(d_0,d_1\le D\), put
\(\widehat v_{k,j}=\pi^S_{d_j}\Phi_{k,d_j}v_j\). Then
\[S_{d_0,d_1}(\delta_k;\widehat v_{k,0},\widehat v_{k,1})\longrightarrow\mathcal W_{(e_{\mathrm t},0)}(\rho v_0,\rho v_1),\]
where \(\rho v_j\) means the on-shell image of the label polynomial
represented by \(v_j\). \emph{Proof sketch.} The finite-regulator
estimate
\href{https://github.com/yang-mills-group/OpenAI-YM-20261028/blob/3456e5e75cf8d4cd558955a5c7638fba2d68f464/YM/Next/RANK/Finite.lean\#L242-L250}{\texttt{Read.full\_cost}}
converts bounds on all nonempty marker components into a moment bound.
The portfolio supplies these component bounds, and passage to the common
extraction subfilter gives
\href{https://github.com/yang-mills-group/OpenAI-YM-20261028/blob/3456e5e75cf8d4cd558955a5c7638fba2d68f464/YM/SharpGram/Weak/Inspect.lean\#L80-L99}{\texttt{actualRead}}.
The resulting comparison
\href{https://github.com/yang-mills-group/OpenAI-YM-20261028/blob/3456e5e75cf8d4cd558955a5c7638fba2d68f464/YM/SharpGram/Weak/Reading.lean\#L94-L105}{\texttt{compareC}}
bounds the difference between the interacting moment at \(\delta_k\) and
the free moment at \(y_k\) by \(C\varepsilon_{g(k)}\), with
\(\varepsilon_{g(k)}\to0\) by III.11. Continuity of the free kernel at
the detector pair gives the limit. \emph{Lean:}
\href{https://github.com/yang-mills-group/OpenAI-YM-20261028/blob/3456e5e75cf8d4cd558955a5c7638fba2d68f464/YM/SharpGram/Weak/Reading.lean\#L107-L124}{\texttt{Portfolio.twoC}}.
\emph{(proof read through the rank-comparison input)}

\textbf{Lemma VII.1.r (raw charts are eventually bijective and nearly
compatible).} Under the hypotheses of VII.1.6:

\begin{enumerate}
\def\labelenumi{\arabic{enumi}.}
\tightlist
\item
  The read-out maps \(D_{k,d}\) are eventually bijective for every
  \(d\le D\), so \(\chi_k(d)=D_{k,d}\). The real part of the reflected
  measured form converges in operator norm to a free form bounded below
  by \(\delta\|v\|^2\) for some \(\delta>0\). It is therefore bounded
  below by \((\delta/2)\|v\|^2\) eventually. This proves injectivity;
  the comparator isomorphism in the construction record gives equality
  of dimensions and hence surjectivity
  (\href{https://github.com/yang-mills-group/OpenAI-YM-20261028/blob/3456e5e75cf8d4cd558955a5c7638fba2d68f464/YM/SharpGram/Weak/Pairing.lean\#L75-L76}{\texttt{freely}},
  \href{https://github.com/yang-mills-group/OpenAI-YM-20261028/blob/3456e5e75cf8d4cd558955a5c7638fba2d68f464/YM/SharpGram/Weak/Pairing.lean\#L110-L128}{\texttt{certified}},
  \href{https://github.com/yang-mills-group/OpenAI-YM-20261028/blob/3456e5e75cf8d4cd558955a5c7638fba2d68f464/YM/SharpGram/Weak/Compatible.lean\#L84-L87}{\texttt{chartDraw}}).
\item
  For every \(d<D\) and \(v\in F_d\),
  \(\chi_k(d+1)^{-1}\iota^S\chi_k(d)v\to\iota^Fv\)
  (\href{https://github.com/yang-mills-group/OpenAI-YM-20261028/blob/3456e5e75cf8d4cd558955a5c7638fba2d68f464/YM/SharpGram/Weak/Compatible.lean\#L98-L121}{\texttt{stepFits}}).
  Both sides are probed by two-point functions against
  \(\chi_k(d+1)\beta_i\), where \((\beta_i)\) is a basis of \(F_{d+1}\);
  Lemma VII.1.q gives convergence of the probes, and injectivity of the
  free probe (Lemma VII.1.p), through the finite-dimensional fitting
  lemma of §2, gives convergence of the arguments.
\item
  Eventually each \(\chi_k(d)\) commutes with conjugation and signed
  permutations
  (\href{https://github.com/yang-mills-group/OpenAI-YM-20261028/blob/3456e5e75cf8d4cd558955a5c7638fba2d68f464/YM/SharpGram/Weak/Compatible.lean\#L124-L136}{\texttt{symmetrical}}),
  and it preserves the unit.
\end{enumerate}

\textbf{Lemma VII.1.s (VII.1.6).} Under the hypotheses of VII.1.6:

\begin{enumerate}
\def\labelenumi{\arabic{enumi}.}
\item
  Take a retraction \(P_d:F_{d+1}\to F_d\) of the free inclusion that is
  equivariant and commutes with conjugation (averaging over \(W_4\) and
  symmetrizing with the conjugation).
\item
  For fixed \(k\), define recursively \(A_0=\chi_k(0)\) and
  \[A_{d+1}=\iota^S A_dP_d+\chi_k(d+1)(1-\iota^FP_d).\]
\item
  For each \(d\le D\), the operators \(\chi_k(d)^{-1}A_d\) tend to the
  identity by induction using Lemma VII.1.r(2). They are eventually
  invertible. Set \(p_k(d)=A_d\) when it is bijective and
  \(p_k(d)=\chi_k(d)\) otherwise.
\item
  Exact compatibility follows from the construction together with
  VII.1.r(3), and \(\chi_k^{-1}p_k\to1\) from VII.1.r(2).
\end{enumerate}

\emph{Lean:}
\href{https://github.com/yang-mills-group/OpenAI-YM-20261028/blob/3456e5e75cf8d4cd558955a5c7638fba2d68f464/YM/SharpGram/Weak/Compatible.lean\#L140-L159}{\texttt{tune\_grades}}.

\textbf{Lemma VII.1.t (VII.1.7).}

\begin{itemize}
\tightlist
\item
  \begin{enumerate}
  \def\labelenumi{(\alph{enumi})}
  \tightlist
  \item
    \href{https://github.com/yang-mills-group/OpenAI-YM-20261028/blob/3456e5e75cf8d4cd558955a5c7638fba2d68f464/YM/SharpGram/Dense/Slender.lean\#L127-L143}{\texttt{slim\_wait\_for}}.
    The required minimum window is \(o(\log\hat\beta)\) at each fixed
    stage, while the prescribed window is
    \(\lfloor\log\hat\beta/\log L_j\rfloor\). Raise entry thresholds to
    fit the finitely many active estimates inside it. Repeat for the
    three panels and raise the thresholds further to ensure that the
    block ratio is below the monitor bound. The station estimates then
    give the simultaneous data in VII.1.7(a).
  \end{enumerate}
\item
  (b)
  \href{https://github.com/yang-mills-group/OpenAI-YM-20261028/blob/3456e5e75cf8d4cd558955a5c7638fba2d68f464/YM/SharpGram/Dense/Cover.lean\#L155-L165}{\texttt{nearby}}.
  If the backward window at \(k\) stays in its stage, take \(g(k)=k\).
  Otherwise take the last step of the preceding stage. Within one stage
  the logarithmic window grows by at most one per level. The threshold
  conditions ensure that this last step is usable; the number of
  intervening levels times \(\log L_j\) is at most \(\log\hat\beta_k\).
  The selected stage tends to infinity, so \(g(k)\to\infty\).
\item
  (c)
  \href{https://github.com/yang-mills-group/OpenAI-YM-20261028/blob/3456e5e75cf8d4cd558955a5c7638fba2d68f464/YM/RG08/Dart/Seams.lean\#L34-L45}{\texttt{coin\_pre}}.
  Put \(z=k+E-q\) and first reduce the weight to \(\sigma=v/(z+2)\) to
  pay for re-rooting. Apply
  \href{https://github.com/yang-mills-group/OpenAI-YM-20261028/blob/3456e5e75cf8d4cd558955a5c7638fba2d68f464/YM/RG08/Dart/Verge.lean\#L62-L86}{\texttt{arose}}
  from rank \(\max(k-q,0)\) through \(z\). Inductively lift the lowest
  remaining coefficient to a curvature expression, subtract its Taylor
  expansion, and continue. At each of these finitely many steps, reduce
  the exponential weight and enlarge the constant to control the
  subtraction. The analytic realization hypothesis permits the lift at
  every step.
\end{itemize}

Parts (a) and (b) supply the common portfolio used in VII.1.6; the
assembly is in III.13. Part (c) supplies local curvature expansions for
the RG estimates.

\emph{(proof read through the stated station and curvature-lifting
inputs)}

\subsubsection{5. Where the difficulty
lies}\label{where-the-difficulty-lies-25}

The main analytic work is the control of finite differences of the torus
Green function, uniformly as the spacing shrinks and the volume grows.
The Fourier argument must retain oscillatory cancellation; an absolute
bound on the full momentum integral would not give the required
derivative limits.

Wick ordering uses the infinite-lattice covariance. Consequently the
same-insertion contraction is not exactly zero at finite volume. For
each fixed stencil offset, the available estimate is \(Cr^2+A/(Nr^4)\),
multiplied by \(h^{-n}\). Taking \(r=h^{n+1}\) gives
\(Ch^{n+2}+A/(Nh^{5n+4})\). The chosen period handles every fixed \(n\)
at once. This is a sufficient growth condition; no optimal volume rate
is asserted.

The reflection-odd field space removes constant zero modes but
introduces an image term. Reporting around the quarter period sends that
image to infinite physical distance. Detection requires a separate
algebraic argument: the translated inversion formula and the triangular
probe substitution establish strict positivity on every nonzero on-shell
polynomial, including constants. The chart correction then uses
finite-dimensional perturbation theory and the interacting estimates of
III.11.

\subsubsection{6. Notation}\label{notation-25}

{\def\LTcaptype{none} % do not increment counter
\begin{longtable}[]{@{}
  >{\raggedright\arraybackslash}p{(\linewidth - 4\tabcolsep) * \real{0.3333}}
  >{\raggedright\arraybackslash}p{(\linewidth - 4\tabcolsep) * \real{0.3333}}
  >{\raggedright\arraybackslash}p{(\linewidth - 4\tabcolsep) * \real{0.3333}}@{}}
\toprule\noalign{}
\begin{minipage}[b]{\linewidth}\raggedright
object
\end{minipage} & \begin{minipage}[b]{\linewidth}\raggedright
code name
\end{minipage} & \begin{minipage}[b]{\linewidth}\raggedright
proposed notation
\end{minipage} \\
\midrule\noalign{}
\endhead
\bottomrule\noalign{}
\endlastfoot
reflection-odd link fields, Feynman-gauge form, Gaussian &
\href{https://github.com/yang-mills-group/OpenAI-YM-20261028/blob/3456e5e75cf8d4cd558955a5c7638fba2d68f464/YM/Pert/Lattice/Geometry.lean\#L107-L111}{\texttt{LinkP}},
\href{https://github.com/yang-mills-group/OpenAI-YM-20261028/blob/3456e5e75cf8d4cd558955a5c7638fba2d68f464/YM/Pert/Lattice/Precision.lean\#L76-L78}{\texttt{form}},
\href{https://github.com/yang-mills-group/OpenAI-YM-20261028/blob/3456e5e75cf8d4cd558955a5c7638fba2d68f464/YM/Pert/Lattice/Measure.lean\#L31}{\texttt{data}}
& \(\mathcal A_N\), \(\mathsf Q_N\), \(\mu_N\) \\
torus Green function (mean zero) &
\href{https://github.com/yang-mills-group/OpenAI-YM-20261028/blob/3456e5e75cf8d4cd558955a5c7638fba2d68f464/YM/Pert/Fourier/Inversion.lean\#L19}{\texttt{toroid}}
& \(g_N\) \\
linear lattice curvature jet &
\href{https://github.com/yang-mills-group/OpenAI-YM-20261028/blob/3456e5e75cf8d4cd558955a5c7638fba2d68f464/YM/Pert/Lattice/TopLinear.lean\#L38}{\texttt{slope}}
& \(\ell^h_x\) \\
\(\mathbb Z^4\) self-covariance of the jets (Wick subtraction) &
\href{https://github.com/yang-mills-group/OpenAI-YM-20261028/blob/3456e5e75cf8d4cd558955a5c7638fba2d68f464/YM/Pert/Lattice/Symbol.lean\#L50}{\texttt{link}},
\href{https://github.com/yang-mills-group/OpenAI-YM-20261028/blob/3456e5e75cf8d4cd558955a5c7638fba2d68f464/YM/Pert/Lattice/Symbol.lean\#L58-L59}{\texttt{shellCov}}
& \(S_h\) \\
Wick-ordered representative &
\href{https://github.com/yang-mills-group/OpenAI-YM-20261028/blob/3456e5e75cf8d4cd558955a5c7638fba2d68f464/YM/Pert/Lattice/Blank.lean\#L41}{\texttt{column}}
& \(R_h(p,o)\) \\
corrected and limit two-leg covariances &
\href{https://github.com/yang-mills-group/OpenAI-YM-20261028/blob/3456e5e75cf8d4cd558955a5c7638fba2d68f464/YM/Pert/Align/Drivers.lean\#L34-L35}{\texttt{corrected}},
\href{https://github.com/yang-mills-group/OpenAI-YM-20261028/blob/3456e5e75cf8d4cd558955a5c7638fba2d68f464/YM/Pert/Align/Drivers.lean\#L36}{\texttt{indirect}}
& \(\kappa_h\), \(K_{ab}\) \\
itinerary; homing sites &
\href{https://github.com/yang-mills-group/OpenAI-YM-20261028/blob/3456e5e75cf8d4cd558955a5c7638fba2d68f464/YM/Pert/Align/Burst.lean\#L50-L59}{\texttt{Itinerary}},
\href{https://github.com/yang-mills-group/OpenAI-YM-20261028/blob/3456e5e75cf8d4cd558955a5c7638fba2d68f464/YM/Pert/Align/Schedule.lean\#L53}{\texttt{Home}}
& \((N_i,h_i)\); \(h_i(o_i-c_{N_i})\to z\) \\
regulator sequence, centre, profile &
\href{https://github.com/yang-mills-group/OpenAI-YM-20261028/blob/3456e5e75cf8d4cd558955a5c7638fba2d68f464/YM/Pert/Lattice/Bundles.lean\#L46}{\texttt{topSpacing}},
\href{https://github.com/yang-mills-group/OpenAI-YM-20261028/blob/3456e5e75cf8d4cd558955a5c7638fba2d68f464/YM/Pert/Lattice/Bundles.lean\#L49}{\texttt{period}},
\href{https://github.com/yang-mills-group/OpenAI-YM-20261028/blob/3456e5e75cf8d4cd558955a5c7638fba2d68f464/YM/Pert/Lattice/Bundles.lean\#L52}{\texttt{center}},
\href{https://github.com/yang-mills-group/OpenAI-YM-20261028/blob/3456e5e75cf8d4cd558955a5c7638fba2d68f464/YM/Pert/Lattice/Bundles.lean\#L55-L56}{\texttt{profile}}
& \(h_m\), \(N_m\), \(c_m\), \(\Pi_m(f;p)\) \\
shell, on-shell reduction, equation ideal & \texttt{Shell},
\href{https://github.com/yang-mills-group/OpenAI-YM-20261028/blob/3456e5e75cf8d4cd558955a5c7638fba2d68f464/YM/Maxwell/Equation.lean\#L47-L50}{\texttt{reduceEq}},
\href{https://github.com/yang-mills-group/OpenAI-YM-20261028/blob/3456e5e75cf8d4cd558955a5c7638fba2d68f464/YM/Maxwell/Equation.lean\#L72}{\texttt{equationIdeal}}
& \(\mathrm{Sh}_j\), \(\rho\), \(\mathcal I_d\) \\
Green function, curvature jet kernel &
\href{https://github.com/yang-mills-group/OpenAI-YM-20261028/blob/3456e5e75cf8d4cd558955a5c7638fba2d68f464/YM/Maxwell/PotentialKernel.lean\#L22}{\texttt{Green}},
\href{https://github.com/yang-mills-group/OpenAI-YM-20261028/blob/3456e5e75cf8d4cd558955a5c7638fba2d68f464/YM/Maxwell/Leg.lean\#L161-L162}{\texttt{legKernel}}
& \(\Gamma\), \(K(P,Q)\) \\
Wick moment &
\href{https://github.com/yang-mills-group/OpenAI-YM-20261028/blob/3456e5e75cf8d4cd558955a5c7638fba2d68f464/YM/Maxwell/Wick.lean\#L126-L127}{\texttt{Wick.moment}}
& \(\mathcal W_x\) \\
detector pair, time reflection &
\href{https://github.com/yang-mills-group/OpenAI-YM-20261028/blob/3456e5e75cf8d4cd558955a5c7638fba2d68f464/YM/Push/FINITE_DETECTION/TwoPoint.lean\#L26}{\texttt{detectors}},
\texttt{Kelvin.σ} & \((e_{\mathrm t},0)\), \(\sigma\) \\
inversion pullback, gain &
\href{https://github.com/yang-mills-group/OpenAI-YM-20261028/blob/3456e5e75cf8d4cd558955a5c7638fba2d68f464/YM/Push/FINITE_DETECTION/Response.lean\#L64-L65}{\texttt{inverted}},
\href{https://github.com/yang-mills-group/OpenAI-YM-20261028/blob/3456e5e75cf8d4cd558955a5c7638fba2d68f464/YM/Push/FINITE_DETECTION/Response.lean\#L92}{\texttt{gain}}
& \(\iota^*\), \(\gamma_j\) \\
weak measuring portfolio, raw chart, corrected chart &
\href{https://github.com/yang-mills-group/OpenAI-YM-20261028/blob/3456e5e75cf8d4cd558955a5c7638fba2d68f464/YM/SharpGram/Weak/Reading.lean\#L35-L49}{\texttt{Portfolio}},
\href{https://github.com/yang-mills-group/OpenAI-YM-20261028/blob/3456e5e75cf8d4cd558955a5c7638fba2d68f464/YM/SharpGram/Weak/Compatible.lean\#L76-L79}{\texttt{chart}},
\texttt{tune\_grades} & \(X\), \(\chi_k(d)\), \(p_k(d)\) \\
\end{longtable}
}

\subsubsection{7. What remains
unexplained}\label{what-remains-unexplained-25}

\begin{enumerate}
\def\labelenumi{\arabic{enumi}.}
\item
  \textbf{Analytic inputs below the Fourier limits.} The symbol
  smoothness, higher-difference estimate, continuum annular decay, tail
  estimate and fixed-offset quadrature estimate support Lemmas
  VII.1.b--c.~The remaining lower-level estimates for periodic annular
  sums, the low-momentum bound and the annular Riemann-sum comparison
  have not been checked line by line. Their uniformity in the period and
  spacing remains an input to the sketches. (probable)
\item
  \textbf{Imported interacting estimates.} The two-point limit follows
  from the simultaneous portfolio conditions and the rank-comparison
  error bound. The full analytic proofs supplying those conditions,
  including the station estimates and curvature-lifting estimates used
  in VII.1.7, are outside this verification. (probable)
\item
  \textbf{Other assigned packages.} The remaining RG local-expansion
  packages and the perturbative-extension packages are not described
  here; no proof outline for their internal estimates is supplied. Their
  placement relative to III.13, III.15 and VII.6 remains to be settled.
  (tentative)
\item
  \textbf{Scope of the comparison.} The Gaussian \(\mu_N\) is defined
  independently of the RG Gaussian measures. Proposition VII.1.1
  requires no identification with them. The interacting short-distance
  theorem and its combination with the tested order-zero limit are the
  linked results in VIII.1; no all-orders perturbative comparison
  follows from this unit.
\item
  \textbf{References and size estimates.} The literature citations and
  attributed source-line counts have not been independently verified.
  (tentative)
\end{enumerate}

\subsection{VII.2. The trace of the stress tensor and the trace
anomaly}\label{vii.2.-the-trace-of-the-stress-tensor-and-the-trace-anomaly}

\textbf{Part:} VII, identification as Yang--Mills. \textbf{Lean:}
\href{https://github.com/yang-mills-group/OpenAI-YM-20261028/tree/3456e5e75cf8d4cd558955a5c7638fba2d68f464/YM/TraceExtension}{\texttt{YM/TraceExtension}}
(namespaces \texttt{YM.TraceExtension} and \texttt{YM.Trace}). The
folder splits as follows:

\begin{itemize}
\tightlist
\item
  the centered trace and its spectral measure: the top-level files
  \texttt{Centered}, \texttt{WilsonTrace}, \texttt{LogScaling},
  \texttt{ConeLaplace}, \texttt{Dyadic}, \texttt{Belfry},
  \texttt{Agora}, \texttt{Endpoint};
\item
  the renormalization-group proof of the anomaly: \texttt{Chart},
  \texttt{Slate}, \texttt{Shelves}, \texttt{Terrace}, \texttt{Margin},
  \texttt{Loft}, \texttt{Lamp}, \texttt{Det}, \texttt{Perfect},
  \texttt{Anchor}, \texttt{Cask}, \texttt{Shepherd}, with the auxiliary
  folders \texttt{Zn}, \texttt{Weir}, \texttt{Pills}, \texttt{Turret},
  \texttt{Wicket}, \texttt{Fen}, \texttt{Swift}, \texttt{Filch},
  \texttt{Quince}, \texttt{Pi}, \texttt{Xo};
\item
  the short-distance law: \texttt{Verse} and \texttt{Foundation};
\item
  the exclusions: \texttt{Wick} and \texttt{Fock};
\item
  a position-space variant: \texttt{Atlas}, \texttt{Pith},
  \texttt{Sundial}, \texttt{Link}.
\end{itemize}

The unit also contains
\href{https://github.com/yang-mills-group/OpenAI-YM-20261028/blob/3456e5e75cf8d4cd558955a5c7638fba2d68f464/YM/Further/PHYSICAL_PERIOD_LIMIT_F933/FreeRanks.lean}{\texttt{FreeRanks.lean}},
the specification
\href{https://github.com/yang-mills-group/OpenAI-YM-20261028/blob/3456e5e75cf8d4cd558955a5c7638fba2d68f464/YM/Nonperturbative/HeadlineSpec.lean}{\texttt{HeadlineSpec.lean}},
\href{https://github.com/yang-mills-group/OpenAI-YM-20261028/blob/3456e5e75cf8d4cd558955a5c7638fba2d68f464/YM/Integration/ManuscriptTraceTests.lean}{\texttt{ManuscriptTraceTests.lean}}
and small pieces of \texttt{YM/Integration} and
\texttt{YM/RG08/Calendar}.

\textbf{Size:} 20,681 lines. The argument headed by
\href{https://github.com/yang-mills-group/OpenAI-YM-20261028/blob/3456e5e75cf8d4cd558955a5c7638fba2d68f464/YM/TraceExtension/Chart/Time.lean\#L63-L68}{\texttt{Store\_l}}
and the argument headed by
\href{https://github.com/yang-mills-group/OpenAI-YM-20261028/blob/3456e5e75cf8d4cd558955a5c7638fba2d68f464/YM/TraceExtension/Margin/Stand.lean\#L23-L35}{\texttt{agenda}}
are the two main RG arguments.

\textbf{Conventions.} \(G\) is a compact matrix group with simple
nonabelian Lie algebra \(\mathfrak g\); \(\dim\mathfrak g\) is the real
dimension. \(C_A\) is the adjoint Casimir of I.1, defined by
\(-\mathrm{tr}(\mathrm{ad}_x\mathrm{ad}_y)=C_A\langle x,y\rangle\). The
constants are
\[\kappa_G=\frac{11\,C_A}{24\pi^2},\qquad b_0=\frac{\kappa_G}{2}=\frac{11\,C_A}{48\pi^2},\qquad d_2=-\frac{\kappa_G}{4}=-\frac{b_0}{2},\qquad A_G=12\Big(\frac{d_2}{b_0}\Big)^2\frac{\dim\mathfrak g}{\pi^4}=\frac{3\dim\mathfrak g}{\pi^4}.\]
\(\kappa_G\) is the slope of the coupling schedule
\(\hat\beta_k=\bar\beta+\kappa_G\log(1/a_k)\) of the renormalization
group (RG) at level \(k\) (II.2, III.9). Write \(g(r)^2=\beta(r)^{-1}\)
for the running coupling in free-field normalization.

\begin{itemize}
\tightlist
\item
  Unless other hypotheses are stated, fix the reconstructed theory
  supplied by a final record \(F\) (VIII.1), over a Wilson family \(C\),
  an extraction order \(f\) and a limit table \(S\).
\item
  \(F_d\) is the space of degree-\(d\) field labels;
  \(\langle\cdot\rangle_{\xi_*}\) is the one-point function of the
  selected phase.
\item
  \(T=(T_{\mu\nu})\) is the Ward-normalized stress tensor of VI.1.
\item
  \(\mathsf F\) is the free Maxwell table (the two constructions listed
  in §6 coincide by definition): \(\dim\mathfrak g\) independent free
  Maxwell fields, with propagator
  \(\langle A A\rangle=1/(4\pi^2|x|^2)\).
\item
  \(\mathcal O\in\mathsf F_4\) is the Wick-ordered label
  \(\tfrac14\sum_{a,\mu,\nu}F^a_{\mu\nu}F^a_{\mu\nu}\), and
  \(T^{\rm free}\) is the traceless free Maxwell stress tensor.
\item
  \(G_\Theta(r)=S^{\xi_*}\big(\Theta(-\tfrac r2e_0)\,\Theta(\tfrac r2e_0)\big)\)
  is the time-axis two-point function of the selected phase.
\end{itemize}

\subsubsection{1. What this unit proves}\label{what-this-unit-proves-26}

The unit defines the centered trace \(\Theta\) and proves four results
about it:

\begin{itemize}
\tightlist
\item
  the trace anomaly at the level of fields: at scale \(r\), \(\Theta\)
  equals \(-\kappa_G\beta(r)^{-1}\mathcal O+o(\beta(r)^{-1})\), where
  \(\beta(r)\sim\kappa_G\log(1/r)\) is the running inverse coupling;
\item
  the resulting short-distance law
  \(r^8\log^2(1/r)\,G_\Theta(r)\to 3\dim\mathfrak g/\pi^4\), and the
  corresponding law for the spectral measure;
\item
  the uniqueness (rank one) of the conserved stress tensor;
\item
  two exclusions of free-field descriptions.
\end{itemize}

Most of the lines prove the field-level anomaly, Proposition VII.2.4.
That proof is a renormalization-group argument: a lattice Ward identity
is transported through the block-spin steps, and a one-loop computation
fixes the coefficient.

\textbf{Proposition VII.2.1 (the centered trace).} Let
\[\Theta:=\sum_\mu T_{\mu\mu}-\Big\langle\sum_\mu T_{\mu\mu}\Big\rangle_{\xi_*}\mathbf 1\in F_4 .\]
(a) \(\langle\Theta\rangle_{\xi_*}=0\). \(\Theta\) is real, invariant
under every \(K\in O(4)\) for the rotation action of VI.1, and fixed by
the time-reflection adjoint. The same holds for the trace centered in
any phase table respecting \(S\).

\begin{enumerate}
\def\labelenumi{(\alph{enumi})}
\setcounter{enumi}{1}
\item
  In the reconstructed Wightman theory of the selected phase, the pinned
  field label \(\iota(\Theta)\) is fixed by the conjugation \(J\) and by
  the Lorentz group. Moreover
  \(\langle\Omega,\Phi_\Theta(f)\Omega\rangle=0\) for every Schwartz
  \(f\).
\item
  For every Schwartz test \(f\), as an identity on the common invariant
  field domain,
  \(\Phi_\Theta(f)=\Phi_{\sum_\mu T_{\mu\mu}}(f)-\big\langle\sum_\mu T_{\mu\mu}\big\rangle_{\xi_*}\big(\int f\big)\,\mathrm{id}\).
\item
  If \(r^8\log^2(1/r)\,G_\Theta(r)\to A>0\), then
  \(\iota(\Theta)\neq0\).
\end{enumerate}

\emph{Status:} theorem. It uses the conservation, reality and
tensoriality of \(T\) and the rotation laws (VI.1), and the
reconstruction (V.2--V.3). \emph{Lean:}

\begin{itemize}
\tightlist
\item
  definitions:
  \href{https://github.com/yang-mills-group/OpenAI-YM-20261028/blob/3456e5e75cf8d4cd558955a5c7638fba2d68f464/YM/TraceExtension/Centered.lean\#L17}{\texttt{stressTrace}},
  \href{https://github.com/yang-mills-group/OpenAI-YM-20261028/blob/3456e5e75cf8d4cd558955a5c7638fba2d68f464/YM/TraceExtension/Centered.lean\#L63}{\texttt{centerIn}},
  \href{https://github.com/yang-mills-group/OpenAI-YM-20261028/blob/3456e5e75cf8d4cd558955a5c7638fba2d68f464/YM/TraceExtension/Centered.lean\#L64}{\texttt{theta}},
  \href{https://github.com/yang-mills-group/OpenAI-YM-20261028/blob/3456e5e75cf8d4cd558955a5c7638fba2d68f464/YM/TraceExtension/WilsonTrace.lean\#L48}{\texttt{center}};
\item
  (a):
  \href{https://github.com/yang-mills-group/OpenAI-YM-20261028/blob/3456e5e75cf8d4cd558955a5c7638fba2d68f464/YM/TraceExtension/Centered.lean\#L66-L67}{\texttt{centerIn\_vac}},
  \href{https://github.com/yang-mills-group/OpenAI-YM-20261028/blob/3456e5e75cf8d4cd558955a5c7638fba2d68f464/YM/TraceExtension/Centered.lean\#L69-L71}{\texttt{trace\_real}},
  \href{https://github.com/yang-mills-group/OpenAI-YM-20261028/blob/3456e5e75cf8d4cd558955a5c7638fba2d68f464/YM/TraceExtension/Centered.lean\#L73-L74}{\texttt{trace\_turn}},
  \href{https://github.com/yang-mills-group/OpenAI-YM-20261028/blob/3456e5e75cf8d4cd558955a5c7638fba2d68f464/YM/TraceExtension/Centered.lean\#L76-L78}{\texttt{trace\_adj}};
\item
  (b):
  \href{https://github.com/yang-mills-group/OpenAI-YM-20261028/blob/3456e5e75cf8d4cd558955a5c7638fba2d68f464/YM/TraceExtension/WilsonTrace.lean\#L58-L59}{\texttt{self\_adjoint}},
  \href{https://github.com/yang-mills-group/OpenAI-YM-20261028/blob/3456e5e75cf8d4cd558955a5c7638fba2d68f464/YM/TraceExtension/WilsonTrace.lean\#L61-L63}{\texttt{pinned\_scalar}}
  with
  \href{https://github.com/yang-mills-group/OpenAI-YM-20261028/blob/3456e5e75cf8d4cd558955a5c7638fba2d68f464/YM/TraceExtension/Centered.lean\#L80-L106}{\texttt{scalar\_rail}},
  \href{https://github.com/yang-mills-group/OpenAI-YM-20261028/blob/3456e5e75cf8d4cd558955a5c7638fba2d68f464/YM/TraceExtension/WilsonTrace.lean\#L65-L67}{\texttt{pinned\_centered}}
  with
  \href{https://github.com/yang-mills-group/OpenAI-YM-20261028/blob/3456e5e75cf8d4cd558955a5c7638fba2d68f464/YM/TraceExtension/WilsonTrace.lean\#L16-L41}{\texttt{vev\_boundary}};
\item
  (c):
  \href{https://github.com/yang-mills-group/OpenAI-YM-20261028/blob/3456e5e75cf8d4cd558955a5c7638fba2d68f464/YM/TraceExtension/Link/Identity.lean\#L158-L161}{\texttt{center\_field}};
\item
  (d):
  \href{https://github.com/yang-mills-group/OpenAI-YM-20261028/blob/3456e5e75cf8d4cd558955a5c7638fba2d68f464/YM/TraceExtension/Sundial/Grand.lean\#L23-L27}{\texttt{nonzero\_strength}};
\item
  public form: the zero-expectation and field-identification clauses of
  \href{https://github.com/yang-mills-group/OpenAI-YM-20261028/blob/3456e5e75cf8d4cd558955a5c7638fba2d68f464/YM/Nonperturbative/HeadlineSpec.lean\#L84-L120}{\texttt{TraceExclusions}},
  i.e.~(T1).
\end{itemize}

\textbf{Proposition VII.2.2 (spectral representation).} There is a
unique positive measure \(\mu_\Theta\) on \(\mathbb R^4\) with the
following properties.

\begin{itemize}
\tightlist
\item
  It has tempered growth: there are \(C_\mu\ge0\) and
  \(k\in\mathbb N_0\) such that
  \(\mu_\Theta\{p:p_0\le R\}\le C_\mu R^k\) for every \(R\ge1\).
\item
  It is supported in the closed forward cone \(\bar V_+\) and is Lorentz
  invariant.
\item
  \(e^{-sp_0}\mu_\Theta\) is finite for \(s>0\), and
  \(G_\Theta(s)=\int e^{-sp_0}\,d\mu_\Theta(p)\) for \(s>0\).
\item
  For every Schwartz \(\varphi\), the spectral measure of the
  translation group in \(\Phi_\Theta(\varphi)\Omega\) is
  \(|\hat\varphi|^2\mu_\Theta\).
\item
  For every Schwartz test \(\phi\) on \((\mathbb R^4)^2\), the Wightman
  function \(W_2(\Theta,\Theta)(\phi)\) equals
  \(\int\widehat\phi^{\,\rm wave}\,d\mu_\Theta\).
\end{itemize}

In particular \(G_\Theta\) is real, nonnegative and nonincreasing on
\((0,\infty)\).

\emph{Status:} theorem: the Källén--Lehmann representation, obtained
from the spectral theorem for the translations of the reconstruction
(V.2--V.3). \emph{Lean:}
\href{https://github.com/yang-mills-group/OpenAI-YM-20261028/blob/3456e5e75cf8d4cd558955a5c7638fba2d68f464/YM/TraceExtension/WilsonTrace.lean\#L83-L94}{\texttt{TraceSpectrum}},
\href{https://github.com/yang-mills-group/OpenAI-YM-20261028/blob/3456e5e75cf8d4cd558955a5c7638fba2d68f464/YM/TraceExtension/WilsonTrace.lean\#L99-L119}{\texttt{traceSpectrum\_exists}},
\href{https://github.com/yang-mills-group/OpenAI-YM-20261028/blob/3456e5e75cf8d4cd558955a5c7638fba2d68f464/YM/TraceExtension/Verse/Cape.lean\#L13-L21}{\texttt{monotone\_axis}}.

\textbf{Proposition VII.2.3 (one-loop certificates along every filter).}
Let \(G\) be simple and let \(C\) be the Wilson family issued by run
data satisfying the RG construction hypotheses of §2, item 4, including
the chosen bounded block-spin run. Let \(\mathcal L\) be any filter on
\(\mathbb N\times\mathbb N\) along which both coordinates tend to
infinity. Then there are work items \(t=(t_{k,N,p})\), that is, level-5
lattice data at RG level \(k\) for the pair \((N,p)\) of spacing and
volume indices, with the following property. For all large \(k\) and
\(\mathcal L\)-eventually in \((N,p)\), the lattice certificate of §6
holds at \((k,N,p)\). In particular the Ward dart \(J=t_{k,N,p}.J\)
(Lemma VII.2.d), a lattice stress tensor with its Ward identity at level
\(k\), satisfies for every color index \(q\)
\[\big|\mathrm{midPay}_q(J)-(-12\kappa_G)\big|\le\hat\beta_k^{-100}.\]
Here \(\mathrm{midPay}_q\) is the real coefficient selector applied to
the Haar divergence of the displacement field rooted at the origin; its
exact definition is
\href{https://github.com/yang-mills-group/OpenAI-YM-20261028/blob/3456e5e75cf8d4cd558955a5c7638fba2d68f464/YM/TraceExtension/Loft/Leaf.lean\#L38}{\texttt{midPay}}.

\emph{Status:} theorem. It is used by III.2 to choose the run and as a
hypothesis of VII.2.4. \emph{Lean:}
\href{https://github.com/yang-mills-group/OpenAI-YM-20261028/blob/3456e5e75cf8d4cd558955a5c7638fba2d68f464/YM/TraceExtension/Margin/Stand.lean\#L23-L35}{\texttt{agenda}},
\href{https://github.com/yang-mills-group/OpenAI-YM-20261028/blob/3456e5e75cf8d4cd558955a5c7638fba2d68f464/YM/TraceExtension/Margin/Stand.lean\#L16}{\texttt{Grant}},
\href{https://github.com/yang-mills-group/OpenAI-YM-20261028/blob/3456e5e75cf8d4cd558955a5c7638fba2d68f464/YM/TraceExtension/Margin/Ansel.lean\#L30-L40}{\texttt{WorthB}}.

\textbf{Proposition VII.2.4 (the gallery property: the anomaly at the
level of fields).} Assume the following:

\begin{itemize}
\tightlist
\item
  \(C\) satisfies the Wilson consistency conditions of II.1, and \(G\)
  is simple;
\item
  the quantitative RG and scaling-sequence hypotheses of §2, items 4 and
  9, hold; the auxiliary iteration depths \(n_k\) are eventually at
  least \(2\), and a sequence \(z_d\) bounds every generator budget of
  degree at most \(d\) and satisfies the degree-\(d\) scaling estimates;
\item
  \(f\) is an extraction order and \(\mathcal L\) is an ultrafilter on
  \(\mathbb N^2\) finer than the fuse filter of \(f\);
\item
  the work items \(t\) of VII.2.3 satisfy the certificate along
  \(\mathcal L\).
\end{itemize}

Then the following property holds. For every limit table \(S\) along
\(f\) with ordered Wilson convergence, a Euclidean hierarchy, point
limits, the cell hypothesis and corrected rank estimates, and every
field algebra on \(S\) satisfying the derivative, rotation, conservation
and Ward laws, there exist:

\begin{itemize}
\tightlist
\item
  limits \(Z\) of the RG data along \(\mathcal L\):
  \(Z.\beta(k)=\lim_{\mathcal L}\beta_{N,p}(k)\), the renormalized
  inverse coupling at level \(k\), with
  \(\hat\beta_k/2\le Z.\beta(k)\le2\hat\beta_k\), together with the
  limits of the label transfers;
\item
  a free Maxwell hierarchy \(\mathsf a\) on \(\mathsf F\),
\end{itemize}

such that the following two statements hold.

\begin{enumerate}
\def\labelenumi{(\roman{enumi})}
\tightlist
\item
  Let \((r_i,U_i)\) be a scaling sequence of \(S\) that approaches
  \(\mathsf a\), with the RG lifts and pointwise lattice limits of
  VII.6. Write \(k_i\to\infty\) for its base levels and \(\ell_i\) for
  the later readout levels selected by these data. Here
  \(r_i\downarrow0\), the \(U_i:\mathsf F_d\to F_d\) are complex-linear
  identifications that eventually preserve the unit, reality and signed
  coordinate actions in each degree, and their rescaled correlations
  approach the free hierarchy. Set \(w_i:=Z.\beta(\ell_i)\). Then:
\end{enumerate}

\begin{itemize}
\tightlist
\item
  \(w_i\to\infty\) and \(\log(1/r_i)/w_i\to\kappa_G^{-1}\);
\item
  \(r_i^4\,U_i^{-1}\mathring T_{\mu\nu}\to T^{\rm free}_{\mu\nu}\),
  where \(\mathring T\) is \(T\) centered in \(S\);
\item
  for every phase table \(T'\) with a hierarchy respecting \(S\), there
  are real sequences \(d_i,e_i\) such that the trace \(\Theta_{T'}\)
  centered in \(T'\) satisfies, eventually,
  \[r_i^4\,U_i^{-1}\Theta_{T'}=e_i\,\mathbf 1+4d_i\,\mathcal O,\qquad w_i d_i\to d_2,\quad w_ie_i\to0,\]
  hence
  \(w_i\,r_i^4U_i^{-1}\Theta_{T'}\to4d_2\,\mathcal O=-\kappa_G\,\mathcal O\).
\end{itemize}

\begin{enumerate}
\def\labelenumi{(\roman{enumi})}
\setcounter{enumi}{1}
\tightlist
\item
  Every sequence \(s_j>0\) with \(s_j\to0\) has a subsequence that is
  the scale sequence of a scaling sequence satisfying (i).
\end{enumerate}

Thus the rescaled trace has leading term \(-\kappa_G g(r)^2\mathcal O\).
In canonically normalized field-strength notation, this is the one-loop
term \(\frac{\beta_{\mathrm{RG}}(g)}{2g}F^a_{\mu\nu}F^a_{\mu\nu}\), with
\(\beta_{\mathrm{RG}}(g)=-b_0g^3+o(g^3)\). The precise field statement
is the sequence limit above.

\emph{Status:} theorem. Its hypotheses are data of the construction run.
They are discharged in III.2: in
\href{https://github.com/yang-mills-group/OpenAI-YM-20261028/blob/3456e5e75cf8d4cd558955a5c7638fba2d68f464/YM/Integration/RG05Preconstruction/Bardo.lean}{\texttt{plentifulWithPrediction}},
for each \(f\) one takes an ultrafilter \(\mathcal L\) finer than the
fuse filter of \(f\), the work items of VII.2.3, and the gallery
assembly. The result is recorded as
\href{https://github.com/yang-mills-group/OpenAI-YM-20261028/blob/3456e5e75cf8d4cd558955a5c7638fba2d68f464/YM/R2I/Campaign/Accepted.lean\#L93-L95}{\texttt{Funded.gallery}}
(VIII.1). \emph{Lean:}
\href{https://github.com/yang-mills-group/OpenAI-YM-20261028/blob/3456e5e75cf8d4cd558955a5c7638fba2d68f464/YM/TraceExtension/Chart/Line.lean\#L49-L56}{\texttt{gang}},
\href{https://github.com/yang-mills-group/OpenAI-YM-20261028/blob/3456e5e75cf8d4cd558955a5c7638fba2d68f464/YM/TraceExtension/Chart/Line.lean\#L41-L44}{\texttt{Gallery}},
\href{https://github.com/yang-mills-group/OpenAI-YM-20261028/blob/3456e5e75cf8d4cd558955a5c7638fba2d68f464/YM/TraceExtension/Chart/Time.lean\#L52-L59}{\texttt{Scribe}},
\href{https://github.com/yang-mills-group/OpenAI-YM-20261028/blob/3456e5e75cf8d4cd558955a5c7638fba2d68f464/YM/TraceExtension/Chart/Time.lean\#L63-L68}{\texttt{Store\_l}}.

\textbf{Proposition VII.2.5 (short-distance law and spectral law).} Let
\(G\) be simple, and suppose the gallery property holds for \((C,f)\).
Then:

\begin{enumerate}
\def\labelenumi{(\alph{enumi})}
\item
  \(\displaystyle\lim_{r\downarrow0}r^8\log^2(1/r)\,G_\Theta(r)=A_G=\frac{3\dim\mathfrak g}{\pi^4}>0\),
  as a limit in \(\mathbb C\).
\item
  As \(\rho\to\infty\),
  \[\rho^{-8}\log^2\!\rho\int\varphi(p/\rho)\,d\mu_\Theta(p)\to\frac{A_G}{1536\pi}\int_{V_+}(p^2)^2\,\varphi(p)\,d^4p,\]
  for continuous compactly supported \(\varphi\) and for Schwartz
  \(\varphi\). Here \(p^2=p_0^2-|\vec p|^2\); the measure on the right
  has Laplace transform \(A_Gs^{-8}\). More explicitly, put
  \(\phi_\rho(x,y)=\rho^8\phi(\rho x,\rho y)\) and
  \(\widehat\phi^{\,\mathrm{wave}}(p)=\int e^{i p\cdot(y-x)}\phi(x,y)\,d^4x\,d^4y\),
  in the Fourier convention of VII.2.2. Then
  \(\rho^{-8}\log^2\rho\,W_2(\Theta,\Theta)(\phi_\rho)\) tends to the
  same cone integral with \(\varphi=\widehat\phi^{\,\mathrm{wave}}\).
\end{enumerate}

\emph{Status:} theorem under the gallery property; (b) uses only (a) and
VII.2.2. The final short-distance and spectral conclusions are filed in
VIII.1. \emph{Lean:}

\begin{itemize}
\tightlist
\item
  (a):
  \href{https://github.com/yang-mills-group/OpenAI-YM-20261028/blob/3456e5e75cf8d4cd558955a5c7638fba2d68f464/YM/TraceExtension/Verse/Cape.lean\#L77-L88}{\texttt{outside}},
  \href{https://github.com/yang-mills-group/OpenAI-YM-20261028/blob/3456e5e75cf8d4cd558955a5c7638fba2d68f464/YM/TraceExtension/LogScaling.lean\#L74-L75}{\texttt{EuclideanAsymptotic}},
  \href{https://github.com/yang-mills-group/OpenAI-YM-20261028/blob/3456e5e75cf8d4cd558955a5c7638fba2d68f464/YM/TraceExtension/Verse/Day.lean\#L41}{\texttt{award}},
  \href{https://github.com/yang-mills-group/OpenAI-YM-20261028/blob/3456e5e75cf8d4cd558955a5c7638fba2d68f464/YM/TraceExtension/Verse/Cape.lean\#L23-L27}{\texttt{award\_pos}};
\item
  (b):
  \href{https://github.com/yang-mills-group/OpenAI-YM-20261028/blob/3456e5e75cf8d4cd558955a5c7638fba2d68f464/YM/TraceExtension/LogScaling.lean\#L151-L157}{\texttt{compact\_limit}},
  \href{https://github.com/yang-mills-group/OpenAI-YM-20261028/blob/3456e5e75cf8d4cd558955a5c7638fba2d68f464/YM/TraceExtension/LogScaling.lean\#L163-L167}{\texttt{spectral\_limit}},
  \href{https://github.com/yang-mills-group/OpenAI-YM-20261028/blob/3456e5e75cf8d4cd558955a5c7638fba2d68f464/YM/TraceExtension/LogScaling.lean\#L203-L213}{\texttt{relative\_limit}},
  \href{https://github.com/yang-mills-group/OpenAI-YM-20261028/blob/3456e5e75cf8d4cd558955a5c7638fba2d68f464/YM/TraceExtension/LogScaling.lean\#L216-L222}{\texttt{joint\_limit}};
\item
  bundles:
  \href{https://github.com/yang-mills-group/OpenAI-YM-20261028/blob/3456e5e75cf8d4cd558955a5c7638fba2d68f464/YM/TraceExtension/Belfry.lean\#L16-L30}{\texttt{Crest}}
  and
  \href{https://github.com/yang-mills-group/OpenAI-YM-20261028/blob/3456e5e75cf8d4cd558955a5c7638fba2d68f464/YM/TraceExtension/Belfry.lean\#L31-L37}{\texttt{crest\_ready}},
  \href{https://github.com/yang-mills-group/OpenAI-YM-20261028/blob/3456e5e75cf8d4cd558955a5c7638fba2d68f464/YM/TraceExtension/Endpoint.lean\#L67-L82}{\texttt{HighEnergy}};
\item
  public form: the amplitude identity and
  \href{https://github.com/yang-mills-group/OpenAI-YM-20261028/blob/3456e5e75cf8d4cd558955a5c7638fba2d68f464/TraceBridge.lean\#L35-L44}{\texttt{publicTrace}},
  i.e.~(T2), (T3).
\end{itemize}

\textbf{Proposition VII.2.6 (the conserved stress tensor is unique).}
Let \(\mathcal K(S)\) be the real space of symmetric tensors
\((B_{\mu\nu})\) in \(F_4\) that are real, centered with respect to
\(S\), and covariant under signed coordinate permutations. For a
derivative plan satisfying the derivative laws and the orientation
condition, let \(\mathrm{div}:\mathcal K(S)\to\mathcal V(S)\) be the
current divergence, where \(\mathcal V(S)\) consists of real
degree-\(5\) vectors covariant under signed coordinate permutations.

\begin{enumerate}
\def\labelenumi{(\alph{enumi})}
\item
  (free) For \(G\) simple, the free Maxwell hierarchy on \(\mathsf F\)
  and a derivative plan satisfying the derivative laws and the
  orientation condition: \(\mathrm{div}\) is onto and
  \(\dim\ker\mathrm{div}=1\).
\item
  (interacting) In the setting of VII.2.4(i), \(\dim\ker\mathrm{div}=1\)
  on \(\mathcal K(S)\); moreover \(\dim\mathcal K(S)=3\) and the target
  of \(\mathrm{div}\) is two-dimensional.
\end{enumerate}

\emph{Status:} theorem. (b) is the uniqueness statement of the
assessment, §7.1. It is used in VII.2.4 to identify \(T\) with the
tensor constructed from the RG. \emph{Lean:}
\href{https://github.com/yang-mills-group/OpenAI-YM-20261028/blob/3456e5e75cf8d4cd558955a5c7638fba2d68f464/YM/Further/PHYSICAL_PERIOD_LIMIT_F933/FreeRanks.lean\#L159-L165}{\texttt{free\_rank}}
(from
\href{https://github.com/yang-mills-group/OpenAI-YM-20261028/blob/3456e5e75cf8d4cd558955a5c7638fba2d68f464/YM/Further/PHYSICAL_PERIOD_LIMIT_F933/FreeRanks.lean\#L83-L113}{\texttt{free\_closure}}),
\href{https://github.com/yang-mills-group/OpenAI-YM-20261028/blob/3456e5e75cf8d4cd558955a5c7638fba2d68f464/YM/TraceExtension/Chart/Masts.lean\#L41-L51}{\texttt{ranking}},
the rank clause of
\href{https://github.com/yang-mills-group/OpenAI-YM-20261028/blob/3456e5e75cf8d4cd558955a5c7638fba2d68f464/YM/TraceExtension/Chart/Masts.lean\#L53-L55}{\texttt{ChoiceK}}.

\textbf{Proposition VII.2.7 (no finite Wick profile).} Suppose
\(r^8\log^2(1/r)\,G_\Theta(r)\to A\) with \(A>0\).

\begin{enumerate}
\def\labelenumi{(\alph{enumi})}
\item
  Let
  \(W(r)=\sum_\alpha c_\alpha\prod_\beta K_{m_{\alpha\beta},q_{\alpha\beta},h_{\alpha\beta}}(r)\)
  be any finite combination, with \(c_\alpha\in\mathbb C\), where
  \[K_{m,q,h}(r)=\int_{\mathbb R^3}\frac{E_m(p)^q\,h(p)}{2E_m(p)}\,e^{-rE_m(p)}\,dp,\qquad E_m=\sqrt{m^2+|p|^2},\]
  with \(m\ge0\), \(q\in\mathbb N_0\), and \(h:\mathbb R^3\to\mathbb R\)
  continuous and positively homogeneous of degree \(j\in\mathbb N_0\).
  Then for every \(\delta>0\) there is \(r\in(0,\delta)\) with
  \(G_\Theta(r)\ne W(r)\).
\item
  No such collection reproduces the Wightman two-point function of
  \(\Theta\) through its mass-shell expression on all product tests
  \(\varphi\otimes\psi\) whose wave transform has compact support. The
  same holds for the Wick contraction of two finite polynomial fields
  from any palette of free-field modes.
\end{enumerate}

\emph{Status:} theorem. (a) uses only the law; (b) also uses VII.2.2.
The internal class is the public class of Definition 1.15 field by field
(\href{https://github.com/yang-mills-group/OpenAI-YM-20261028/blob/3456e5e75cf8d4cd558955a5c7638fba2d68f464/TraceBridge.lean\#L19-L23}{\texttt{profile\_eq}}).
With VII.2.5 this gives (T4). The proof is also described in VII.3.
\emph{Lean:}
\href{https://github.com/yang-mills-group/OpenAI-YM-20261028/blob/3456e5e75cf8d4cd558955a5c7638fba2d68f464/YM/TraceExtension/Wick/Refute.lean\#L75-L80}{\texttt{euclidean\_excludes}},
\href{https://github.com/yang-mills-group/OpenAI-YM-20261028/blob/3456e5e75cf8d4cd558955a5c7638fba2d68f464/YM/TraceExtension/Wick/Refute.lean\#L86-L105}{\texttt{polynomial\_excludes}},
\href{https://github.com/yang-mills-group/OpenAI-YM-20261028/blob/3456e5e75cf8d4cd558955a5c7638fba2d68f464/YM/TraceExtension/Wick/Palette.lean\#L130-L133}{\texttt{wick\_excludes}},
\href{https://github.com/yang-mills-group/OpenAI-YM-20261028/blob/3456e5e75cf8d4cd558955a5c7638fba2d68f464/YM/TraceExtension/Wick/PolynomialField.lean\#L23-L29}{\texttt{Mode}},
\href{https://github.com/yang-mills-group/OpenAI-YM-20261028/blob/3456e5e75cf8d4cd558955a5c7638fba2d68f464/YM/TraceExtension/Wick/PolynomialField.lean\#L90-L92}{\texttt{Track}}.

\textbf{Proposition VII.2.8 (no relatively local realization in
free-photon Fock space).} Suppose the law of VII.2.7 holds with \(A>0\).
For every \(m\in\mathbb N_0\), let \(\mathcal F_m\) be the Fock space of
\(m\) species of free photons, realized by symmetric transverse
(``physical'') wave functions of on-shell momenta. Let
\(U:\mathcal H\to\mathcal F_m\) be a complex-linear isometry from the
reconstructed space of the selected phase with:

\begin{itemize}
\tightlist
\item
  \(U\Omega=\Omega_{\mathcal F}\);
\item
  image in physical waves;
\item
  \(U\) intertwines the translations with multiplication by
  \(e^{ib\cdot p_{\rm tot}}\).
\end{itemize}

Then there are a compactly supported smooth scalar test \(f\) and
compactly supported smooth tests \(h_c:\mathbb R^4\to\mathbb R^3\),
\(1\le c\le m\), each causally separated from \(f\), such that
\[\langle W(-h)\Omega_{\mathcal F},\,U\Phi_\Theta(f)\Omega\rangle\neq\langle U\Phi_\Theta(\bar f)\Omega,\,W(h)\Omega_{\mathcal F}\rangle ,\]
where \(W(h)\Omega_{\mathcal F}\) is the coherent state defined in §6.
Equality for all such pairs would hold if \(\Phi_\Theta(f)\) commuted
with the free Weyl operators localized spacelike to \(f\).

The same conclusion holds in two variants:

\begin{itemize}
\tightlist
\item
  with any bounded self-adjoint operator fixing \(U\mathcal H\) applied
  to both coherent vectors;
\item
  with an average over a probability family of maps that preserve inner
  products of physical waves and fix \(U\mathcal H\); for every \(r>0\),
  this variant has tests supported in open sets \(Q,W\) with
  \(\bar Q\subset B(0,r)\) and \(\bar Q-\bar W\) spacelike, with \(f\)
  supported in \(Q\) and every component of \(h\) in \(W\).
\end{itemize}

\emph{Status:} theorem. These are the free-photon and averaged
free-photon clauses of the trace-exclusion specification; the final
conclusions are filed in VIII.1. \emph{Lean:}
\href{https://github.com/yang-mills-group/OpenAI-YM-20261028/blob/3456e5e75cf8d4cd558955a5c7638fba2d68f464/YM/TraceExtension/Fock/Identify.lean\#L23-L27}{\texttt{Comparison}},
\href{https://github.com/yang-mills-group/OpenAI-YM-20261028/blob/3456e5e75cf8d4cd558955a5c7638fba2d68f464/YM/TraceExtension/Fock/Identify.lean\#L32-L35}{\texttt{Nfl}},
\href{https://github.com/yang-mills-group/OpenAI-YM-20261028/blob/3456e5e75cf8d4cd558955a5c7638fba2d68f464/YM/TraceExtension/Fock/Conclude.lean\#L46-L49}{\texttt{excludes}},
\href{https://github.com/yang-mills-group/OpenAI-YM-20261028/blob/3456e5e75cf8d4cd558955a5c7638fba2d68f464/YM/TraceExtension/Fock/Conclude.lean\#L33-L44}{\texttt{saturate}},
\href{https://github.com/yang-mills-group/OpenAI-YM-20261028/blob/3456e5e75cf8d4cd558955a5c7638fba2d68f464/YM/TraceExtension/Fock/Conclude.lean\#L51-L70}{\texttt{excludes\_projected}},
\href{https://github.com/yang-mills-group/OpenAI-YM-20261028/blob/3456e5e75cf8d4cd558955a5c7638fba2d68f464/YM/TraceExtension/Fock/Occult.lean\#L12-L21}{\texttt{smell}},
\href{https://github.com/yang-mills-group/OpenAI-YM-20261028/blob/3456e5e75cf8d4cd558955a5c7638fba2d68f464/YM/TraceExtension/Fock/Regions.lean\#L36-L38}{\texttt{regions\_excludes}}.

\textbf{Proposition VII.2.9 (position-space coefficient;
supplementary).} For every simple \(G\), every block-spin kernel
parameter \(b>0\) and every prescribed extraction order \(p\), there are
a consistent Wilson family \(C\) with that kernel parameter and a
refinement \(f\) of \(p\) for which final records exist and the
following holds for every final record \(F\) over \((C,f)\).

There are a free hierarchy, coupling limits \(Z\) and an ultrafilter
such that admissible degree-\(4\) position-space frames exist. For every
such pair of frames \(P(s),P'(s):\mathsf F_4\to F_4\), with a window
cutoff \(\chi\), there are real functions \(\kappa,d,e\) such that, for
all sufficiently small \(s>0\),
\[P'(s)^{-1}(s^4T_{\mu\nu})=\kappa(s)T^{\rm free}_{\mu\nu}+d(s)\delta_{\mu\nu}\mathcal O+e(s)\delta_{\mu\nu}\mathbf1,\qquad \kappa(s)\to1.\]
Admissibility means that \(P\) has the position-space Gram
normalization, branch and angular alignment, and convergence of all
prescribed finite off-diagonal jets; \(P'\) is obtained from \(P\) by
positive rescaling of the centered scalar sector and its centered
complement, fixing the unit, with the two shell tests normalized to
their free values. These conditions are specified by the frame predicate
linked below.

For every coupling function \(g\) that is eventually positive and
satisfies \(\log(1/s)g(s)^2\to\kappa_G^{-1}\),
\[\frac{d(s)}{g(s)^2}\to d_2.\] The coupling \(g_X\) defined by the
position-space tests satisfies \(g_X/g_W\to1\), where \(g_W\) is the
Wilson coupling from \(Z\), and obeys the stronger estimate
\[d(s)-d_2g_X(s)^2=O(g_X(s)^4).\] Here the free scalar is exactly
\(\mathcal O\). Centering gives the exact identity
\(s^4\Theta=4d(s)\,\mathrm{center}_{\xi_*}(P'(s)\mathcal O)\) for
sufficiently small \(s>0\).

\emph{Proof sketch.} Subsequence compactness gives position-space frames
at every sufficiently small scale. The two positive sector rescalings
tend to \(1\) and impose the shell normalizations. Put \(B_0\ne0\) for
the common normalized shell value and \(Q(s)\) for the remaining shell
readout, which has a finite limit. Conservation and the Ward identity
give \(d=d_2g_X^2\kappa\) and \(\kappa-1=-dQ/B_0\). Since
\(d/g_X^2\to d_2\), these identities imply \(\kappa-1=O(g_X^2)\) and
then \(d-d_2g_X^2=O(g_X^4)\). The Wilson comparison is established
before the definition of \(g_X\) by computing the same divergence
diagnostic in Wilson normalization.

\emph{Status:} theorem for the family and refinement supplied by the
supplementary construction. \emph{Lean:}
\href{https://github.com/yang-mills-group/OpenAI-YM-20261028/blob/3456e5e75cf8d4cd558955a5c7638fba2d68f464/YM/TraceExtension/Sundial/Noon.lean\#L50-L83}{\texttt{meterQ}}
\emph{(proof read; frame-existence inputs statement only)};
\href{https://github.com/yang-mills-group/OpenAI-YM-20261028/blob/3456e5e75cf8d4cd558955a5c7638fba2d68f464/YM/TraceExtension/Endpoint.lean\#L155-L161}{\texttt{main}},
\href{https://github.com/yang-mills-group/OpenAI-YM-20261028/blob/3456e5e75cf8d4cd558955a5c7638fba2d68f464/YM/TraceExtension/Endpoint.lean\#L29-L34}{\texttt{XScheme}},
\href{https://github.com/yang-mills-group/OpenAI-YM-20261028/blob/3456e5e75cf8d4cd558955a5c7638fba2d68f464/YM/TraceExtension/Sundial/Shade.lean\#L20-L22}{\texttt{ShellNorm}};
the stress clause of
\href{https://github.com/yang-mills-group/OpenAI-YM-20261028/blob/3456e5e75cf8d4cd558955a5c7638fba2d68f464/YM/TraceExtension/Endpoint.lean\#L99-L120}{\texttt{Result}},
\href{https://github.com/yang-mills-group/OpenAI-YM-20261028/blob/3456e5e75cf8d4cd558955a5c7638fba2d68f464/YM/TraceExtension/Endpoint.lean\#L37-L64}{\texttt{FieldComparison}},
\href{https://github.com/yang-mills-group/OpenAI-YM-20261028/blob/3456e5e75cf8d4cd558955a5c7638fba2d68f464/YM/TraceExtension/Atlas/Crash.lean\#L33-L36}{\texttt{odd\_limit}},
\href{https://github.com/yang-mills-group/OpenAI-YM-20261028/blob/3456e5e75cf8d4cd558955a5c7638fba2d68f464/YM/TraceExtension/Pith/Wall.lean\#L13-L14}{\texttt{Measures}}.

\subsubsection{2. What it uses}\label{what-it-uses-26}

\textbf{Inputs from other units.}

\begin{enumerate}
\def\labelenumi{\arabic{enumi}.}
\item
  \emph{Gauge data} (I.1): the Casimir \(C_A>0\) and \(\kappa_G>0\)
  (\href{https://github.com/yang-mills-group/OpenAI-YM-20261028/blob/3456e5e75cf8d4cd558955a5c7638fba2d68f464/YM/Gauge/Colors.lean\#L184}{\texttt{slopeConstant}});
  simplicity; a nontrivial color space.
\item
  \emph{Coupling schedule} (II.2, III.9):
  \(\hat\beta_k=\bar\beta+\kappa_G\log(1/a_k)\), with
  \(\hat\beta_k\to\infty\)
  (\href{https://github.com/yang-mills-group/OpenAI-YM-20261028/blob/3456e5e75cf8d4cd558955a5c7638fba2d68f464/YM/TraceExtension/Terrace/Lyre.lean\#L8}{\texttt{hrate}}).
  The actual couplings follow the schedule:
  \(Z.\beta(k)/\hat\beta_k\to1\)
  (\href{https://github.com/yang-mills-group/OpenAI-YM-20261028/blob/3456e5e75cf8d4cd558955a5c7638fba2d68f464/YM/TraceExtension/Terrace/Dusk.lean\#L89-L95}{\texttt{gamut}},
  \href{https://github.com/yang-mills-group/OpenAI-YM-20261028/blob/3456e5e75cf8d4cd558955a5c7638fba2d68f464/YM/TraceExtension/Terrace/Bang.lean\#L23-L37}{\texttt{Barry}}).
\item
  \emph{The one-loop integral} (III.7): with infrared mass
  \(\mu\in(0,1]\), the one-loop momentum integral has coefficient
  \(-\kappa_G\log(1/\mu)\) up to a bounded error
  (\href{https://github.com/yang-mills-group/OpenAI-YM-20261028/blob/3456e5e75cf8d4cd558955a5c7638fba2d68f464/YM/Next/STATIONARY/OneLoop/MomentumIntegral.lean\#L346-L408}{\texttt{absoluteGeode}}).
\item
  \emph{The RG step} (III.4--III.15):

  \begin{itemize}
  \tightlist
  \item
    the stationary background, criticality and fourth-order classical
    matching of actions (III.6--III.7):
    \href{https://github.com/yang-mills-group/OpenAI-YM-20261028/blob/3456e5e75cf8d4cd558955a5c7638fba2d68f464/YM/Next/STATIONARY/Joint/Gauge.lean\#L219-L220}{\texttt{encamp}};
  \item
    cluster and polymer bounds (III.5, III.8) and determinant budgets
    (III.4, III.11);
  \item
    the bounded block-spin run, eventual level estimates, classical and
    determinant bounds, support conditions, and small-field and
    large-field estimates. The exact run predicates are listed in §6.
  \end{itemize}

  These hypotheses are discharged when the run is chosen in III.2.
\item
  \emph{Infinite volume} (IV.1): the extraction order and its fuse
  filter, point limits,
  \href{https://github.com/yang-mills-group/OpenAI-YM-20261028/blob/3456e5e75cf8d4cd558955a5c7638fba2d68f464/YM/Cells/Packing.lean\#L68-L79}{\texttt{CellHyp}},
  and the chessboard root tests on centered chambers.
\item
  \emph{Phases and reconstruction} (V.1--V.3): the selected phase; the
  analytic continuation and its boundary values; the spectral measure of
  the translations; the Lorentz action and the conjugation \(J\). All of
  this is packaged in
  \href{https://github.com/yang-mills-group/OpenAI-YM-20261028/blob/3456e5e75cf8d4cd558955a5c7638fba2d68f464/YM/Spec/Closed.lean\#L64-L73}{\texttt{Final}}.
\item
  \emph{Field algebra and stress tensor} (VI.1): the derivative plan,
  the rotation laws (including that rotations and boosts generate the
  Lorentz group), and the stress tensor with its conservation,
  tensoriality and Ward identities
  (\href{https://github.com/yang-mills-group/OpenAI-YM-20261028/blob/3456e5e75cf8d4cd558955a5c7638fba2d68f464/YM/Spec/Closed.lean\#L34-L43}{\texttt{Locals}}).
\item
  \emph{Free Maxwell model} (VI.2, VII.1, VII.3):

  \begin{itemize}
  \tightlist
  \item
    the free table \(\mathsf F\) and its Wick calculus;
  \item
    the free stress tensors
    (\href{https://github.com/yang-mills-group/OpenAI-YM-20261028/blob/3456e5e75cf8d4cd558955a5c7638fba2d68f464/YM/TraceExtension/Shelves/Pear.lean\#L50}{\texttt{materialB}},
    \href{https://github.com/yang-mills-group/OpenAI-YM-20261028/blob/3456e5e75cf8d4cd558955a5c7638fba2d68f464/YM/TraceExtension/Chart/Auth.lean\#L30-L34}{\texttt{Litpay}});
  \item
    the normalization
    \(\langle\mathcal O(x)\mathcal O(y)\rangle=3\dim\mathfrak g/(\pi^4|x-y|^8)\)
    (\href{https://github.com/yang-mills-group/OpenAI-YM-20261028/blob/3456e5e75cf8d4cd558955a5c7638fba2d68f464/YM/TraceExtension/Foundation/Constant.lean\#L38-L52}{\texttt{wick\_norm}},
    proved by an explicit curl computation);
  \item
    convergence of rescaled two-point functions at lattice-rounded
    points
    (\href{https://github.com/yang-mills-group/OpenAI-YM-20261028/blob/3456e5e75cf8d4cd558955a5c7638fba2d68f464/YM/TraceExtension/Verse/Round.lean\#L51-L63}{\texttt{rainbow}}).
  \end{itemize}
\item
  \emph{Scaling sequences} (VII.6): scaling sequences approaching the
  free hierarchy
  (\href{https://github.com/yang-mills-group/OpenAI-YM-20261028/blob/3456e5e75cf8d4cd558955a5c7638fba2d68f464/YM/TraceExtension/Prom/Zoom/Running.lean\#L19-L23}{\texttt{StarPath}},
  \href{https://github.com/yang-mills-group/OpenAI-YM-20261028/blob/3456e5e75cf8d4cd558955a5c7638fba2d68f464/YM/TraceExtension/Terrace/Gondola.lean\#L18-L21}{\texttt{Gondola}})
  exist along a subsequence of any positive null sequence
  (\href{https://github.com/yang-mills-group/OpenAI-YM-20261028/blob/3456e5e75cf8d4cd558955a5c7638fba2d68f464/YM/TraceExtension/Prom/Zoom/Running.lean\#L25-L31}{\texttt{pays\_sky}},
  \href{https://github.com/yang-mills-group/OpenAI-YM-20261028/blob/3456e5e75cf8d4cd558955a5c7638fba2d68f464/YM/TraceExtension/Terrace/Gondola.lean\#L22-L28}{\texttt{gravity}}).
\end{enumerate}

\textbf{Standard mathematics.}

\begin{itemize}
\tightlist
\item
  The trace-anomaly mechanism (Adler--Collins--Duncan, Phys. Rev.~D 15
  (1977); Collins--Duncan--Joglekar, Phys. Rev.~D 16 (1977); Nielsen,
  Nucl. Phys. B 120 (1977)) and the background-field one-loop beta
  function ('t Hooft 1973; Abbott, Nucl. Phys. B 185 (1981)).
\item
  The Källén--Lehmann representation and the spectral theorem for the
  translation group (Reed--Simon II, §IX.8; Glimm--Jaffe, ch.~19).
\item
  The identity theorem for holomorphic functions in tube domains.
\item
  Laplace transforms of positive measures, continuity theorems and
  Tauberian monotonicity (Feller II, §XIII.1--5; Widder, \emph{The
  Laplace Transform}).
\item
  Fock space, Weyl operators and coherent states (Reed--Simon II, §X.7;
  Bratteli--Robinson II, §5.2).
\item
  Simultaneous eigenvector decomposition for commuting operators
  annihilated by polynomials with distinct real roots.
\item
  Lattice difference operators and finite-difference stencils of order
  8.
\end{itemize}

\textbf{Imported constants and assumptions.} \(\kappa_G\), \(b_0\) and
\(d_2\) are defined from \(C_A\)
(\href{https://github.com/yang-mills-group/OpenAI-YM-20261028/blob/3456e5e75cf8d4cd558955a5c7638fba2d68f464/YM/TraceExtension/Chart/Nectar.lean\#L60}{\texttt{b0}}),
so \(\kappa_G=2b_0\) and \(d_2=-b_0/2\) are identities by definition
(\href{https://github.com/yang-mills-group/OpenAI-YM-20261028/blob/3456e5e75cf8d4cd558955a5c7638fba2d68f464/YM/TraceExtension/Chart/Nectar.lean\#L61}{\texttt{beta\_eq}},
\href{https://github.com/yang-mills-group/OpenAI-YM-20261028/blob/3456e5e75cf8d4cd558955a5c7638fba2d68f464/YM/TraceExtension/Chart/Nectar.lean\#L63}{\texttt{force\_b}}).
The value \(\kappa_G\) enters the proofs in two independent places:

\begin{itemize}
\tightlist
\item
  as the slope of the coupling schedule (input 2);
\item
  as the logarithmic coefficient of the one-loop integral (input 3).
\end{itemize}

The amplitude \(3\dim\mathfrak g/\pi^4\) depends on these two being
equal (§5). The gallery property is a hypothesis of VII.2.5; it is
discharged for the constructed family in III.2, so the trace data must
be carried by the family chosen in the ultraviolet construction.

\subsubsection{3. Strategy}\label{strategy-26}

\emph{Physical picture and the approach taken.} Classically the
Yang--Mills stress tensor is traceless. The running coupling produces a
nonzero trace after renormalization. It characterizes \(T\) by two
properties: uniqueness (the conserved admissible tensors form a line,
VII.2.6) and the Ward normalization. At each scale \(r_i\) it builds
three lattice tensors \(T_i^{(B)}=\mathrm{tub}_i(B)\) from the reference
tensors \(B=L,\Delta,\Xi\): the free Maxwell stress \(L\), the
pure-trace tensor \(\Delta_{\mu\nu}=\delta_{\mu\nu}\mathcal O\), and a
third traceless tensor \(\Xi\).

The candidate is
\(X_i=w_iT_i^{(L)}+c_{i,0}T_i^{(\Xi)}+c_{i,1}T_i^{(\Delta)}\), where
\(w_i\sim\beta(r_i)\). The suitably rescaled divergence of the canonical
part \(w_iT_i^{(L)}\), tested by \(\rho\), tends to \(-\kappa_G\).
Conservation forces \(c_{i,1}\to-\kappa_G/4=d_2\). Uniqueness makes
\(X_i\) a multiple of the centered stress tensor; the Ward normalization
makes that multiple tend to \(1\). The weighted trace therefore
satisfies \(w_i\,\mathrm{tr}(r_i^4U_i^{-1}X_i)\to4d_2\mathcal O\).
Combining this with \(w_i\sim\kappa_G\log(1/r_i)\) gives the leading
rescaled trace \(-\mathcal O/\log(1/r_i)\).

\emph{The one-loop input.} The value \(-\kappa_G\) comes from a lattice
Ward identity at a single RG level, a ``dart''. A dart is a lattice
stress tensor whose divergence equals the first variation of the action
along a displacement field, through fourth order in the field. Darts at
a given level are obtained from finer levels by block-spin transport.
The transport preserves the Ward identity, because the first variations
of the fine and coarse effective actions agree to fourth order; errors
contract by \(\frac12\) per step. The level-5 ``payoff'' of the dart is
then compared with the infrared-regulated one-loop effective action,
\[\tfrac12\log\det(\text{gauge Hessian}+\mu^2)-\log\det(\text{ghost operator}+\mu^2).\]
Its normalized quadratic momentum coefficient has the logarithmic term
\(-\kappa_G\log(1/\mu)\) per plane, \(12\) ordered planes in total. A
mean-value argument in \(\log\mu\) turns this pole into the limit
\(-12\kappa_G\) for the payoff. This is Balaban's renormalization-group
framework (III) together with a background-field one-loop computation;
what is specific here is that the anomaly coefficient is extracted from
the Ward identity carried along the RG flow.

\emph{From fields to the law, and the exclusions.} By free-field
convergence (VII.1), the two-point function of
\(w_ir_i^4U_i^{-1}\Theta\) converges to
\((4d_2)^2\cdot3\dim\mathfrak g/(\pi^4|x|^8)\), and
\(\log(1/r_i)/w_i\to1/\kappa_G\). Rotation and translation invariance
move the points to the time axis. Monotonicity of \(G_\Theta\) in the
separation removes the lattice rounding of points. Finally, every
positive null sequence has a good subsequence, which upgrades the
subsequential limit to \(r\downarrow0\). A Tauberian argument transfers
the law to the spectral measure. The exclusions use only the law and
\(\mu_\Theta\):

\begin{itemize}
\tightlist
\item
  every finite Wick profile expands in integer powers and nonnegative
  powers of logarithms, which are incompatible with
  \(r^{-8}\log^{-2}(1/r)\);
\item
  a relatively local Fock realization forces each \(n\)-photon amplitude
  to split into eigencomponents of the commuting momentum-doubling maps,
  with total scaling degree at most \(3\); this bounds
  \(\mu_\Theta\{p_0\le R\}\) by \(R^7\), against \(R^8/\log^2R\).
\end{itemize}

\subsubsection{4. Steps}\label{steps-26}

\paragraph{4.1 The centered trace and its spectral
measure}\label{the-centered-trace-and-its-spectral-measure}

\textbf{Lemma VII.2.a (centering and symmetry).} The map
\(v\mapsto\langle v\rangle\) is linear with
\(\langle\mathbf 1\rangle=1\), so \(\langle\Theta\rangle=0\). Reality
and \(O(4)\) invariance follow from those of \(T\): contracting the
tensorial law with \(KK^{\mathsf T}=1\) gives
\(\sum_aK_{ia}K_{ja}=\delta_{ij}\). Time-reflection invariance follows
from reality and rotation invariance. For Lorentz invariance of
\(\iota(\Theta)\), the Lorentz group is generated by spatial rotations
and axis boosts. Along a boost,
\(t\mapsto\mathrm{boost}(t)\iota(\Theta)\) solves the linear ODE
\(\dot y=-iX_{0a}y\), and the generator annihilates \(\iota(\Theta)\)
because \(\Theta\) is rotation invariant. Uniqueness of solutions of the
ODE gives invariance. \emph{Lean:}
\href{https://github.com/yang-mills-group/OpenAI-YM-20261028/blob/3456e5e75cf8d4cd558955a5c7638fba2d68f464/YM/TraceExtension/Centered.lean\#L73-L74}{\texttt{trace\_turn}},
\href{https://github.com/yang-mills-group/OpenAI-YM-20261028/blob/3456e5e75cf8d4cd558955a5c7638fba2d68f464/YM/TraceExtension/Centered.lean\#L80-L106}{\texttt{scalar\_rail}}.

\textbf{Lemma VII.2.b (zero Wightman vacuum expectation).} For a
centered label, the one-point function continued to complex times is
holomorphic and vanishes on the Euclidean slice, so it vanishes
identically by the identity theorem. Its boundary value, the Wightman
one-point function, is therefore \(0\). \emph{Lean:}
\href{https://github.com/yang-mills-group/OpenAI-YM-20261028/blob/3456e5e75cf8d4cd558955a5c7638fba2d68f464/YM/TraceExtension/WilsonTrace.lean\#L16-L41}{\texttt{vev\_boundary}}.

\textbf{Lemma VII.2.c (spectral measure).} The boundary realization
(V.2) assigns a positive spectral measure to the label
\(\iota(\Theta)\). Smearing the field with \(\varphi\) gives a
Hilbert-space vector whose translation spectral measure is
\(|\hat\varphi|^2\mu_\Theta\).

\begin{itemize}
\tightlist
\item
  The spectrum condition gives support in \(\bar V_+\); the boundary
  estimates give tempered growth and the polynomial energy budget.
\item
  Lorentz invariance comes from Lemma VII.2.a.
\item
  The Laplace formula comes from continuing \(G_\Theta(2s)\) to the
  boundary pairing at imaginary time
  (\href{https://github.com/yang-mills-group/OpenAI-YM-20261028/blob/3456e5e75cf8d4cd558955a5c7638fba2d68f464/YM/TraceExtension/WilsonTrace.lean\#L69-L78}{\texttt{continuation}}).
\item
  Uniqueness holds because \(\mu\) is determined by the integrals
  \(\int|\hat\varphi|^2d\mu\).
\end{itemize}

Monotonicity holds because \(e^{-sp_0}\) decreases in \(s\) on
\(\bar V_+\). \emph{Lean:}
\href{https://github.com/yang-mills-group/OpenAI-YM-20261028/blob/3456e5e75cf8d4cd558955a5c7638fba2d68f464/YM/TraceExtension/WilsonTrace.lean\#L99-L119}{\texttt{traceSpectrum\_exists}},
\href{https://github.com/yang-mills-group/OpenAI-YM-20261028/blob/3456e5e75cf8d4cd558955a5c7638fba2d68f464/YM/TraceExtension/Verse/Cape.lean\#L13-L21}{\texttt{monotone\_axis}}.
\emph{(proof read at the top level; the boundary spectral lemmas are
inputs from V.2)}

\paragraph{4.2 Ward darts and their transport through the
RG}\label{ward-darts-and-their-transport-through-the-rg}

\textbf{Lemma VII.2.d (sound darts and the block-spin step).} A dart on
a torus of side \(f\) consists of site-indexed functions and vector
fields: a tensor \(J_{ab}\), a displacement field \(Y_a\) on link
configurations, scalar entries \(E_a,Q_a\), and scalar entries
\(\mathrm{tw}_{abi}\). It is \emph{sound} if:

\begin{itemize}
\tightlist
\item
  the scalar entries are analytic and gauge invariant near the identity
  configuration, with zero constant and linear terms; the displacement
  fields are analytic and gauge covariant and vanish at the identity;
\item
  \(\sum_aJ_{aa}=\sum_a\partial_aQ_a\) and
  \(J_{ab}-J_{ba}=\sum_i\partial_i\mathrm{tw}_{abi}\), with lattice
  differences \(\partial\);
\item
  the Ward relation holds: \(\sum_a\partial_aJ_{ab}-E_b\) agrees through
  fourth order at zero field with the first variation of the action in
  the direction \(Y_b\).
\end{itemize}

The block-spin map takes a dart at the fine level to one at the coarse
level, by block averaging of inherited jets scaled by \(L^k\) with
\(k=3,4,5\). It preserves soundness, under the criticality, branch,
bound, positive price-gap and classical-matching hypotheses of
III.6--III.7. The mechanism is fourth-order matching of first
variations: along an analytic coarse direction \(Y\) with \(Y(0)=0\),
the derivative of the coarse action at the stationary background in the
direction \(Y\) agrees through fourth order with the derivative of the
fine action in the lifted direction. This follows by the chain rule
through the stationary background map. \emph{Lean:}
\href{https://github.com/yang-mills-group/OpenAI-YM-20261028/blob/3456e5e75cf8d4cd558955a5c7638fba2d68f464/YM/TraceExtension/Perfect/Badge.lean\#L22-L27}{\texttt{Dart}},
\href{https://github.com/yang-mills-group/OpenAI-YM-20261028/blob/3456e5e75cf8d4cd558955a5c7638fba2d68f464/YM/TraceExtension/Perfect/Badge.lean\#L82-L90}{\texttt{Dart.Sound}},
\href{https://github.com/yang-mills-group/OpenAI-YM-20261028/blob/3456e5e75cf8d4cd558955a5c7638fba2d68f464/YM/TraceExtension/Perfect/Badge.lean\#L75-L80}{\texttt{Dart.down}},
\href{https://github.com/yang-mills-group/OpenAI-YM-20261028/blob/3456e5e75cf8d4cd558955a5c7638fba2d68f464/YM/TraceExtension/Perfect/Badge.lean\#L110-L141}{\texttt{Sound.balance}};
\href{https://github.com/yang-mills-group/OpenAI-YM-20261028/blob/3456e5e75cf8d4cd558955a5c7638fba2d68f464/YM/TraceExtension/Perfect/Wicker.lean\#L37-L46}{\texttt{aroma\_transfer}},
\href{https://github.com/yang-mills-group/OpenAI-YM-20261028/blob/3456e5e75cf8d4cd558955a5c7638fba2d68f464/YM/TraceExtension/Perfect/Relay.lean\#L143-L160}{\texttt{recoup}}.

\textbf{Lemma VII.2.e (darts with prescribed precision).} Under the
bounded-run and eventual level hypotheses of VII.2.3, for all large
levels \(k\), every \(\varepsilon>0\) and every filter along which
\(N\to\infty\) and the lattice side at each fixed level tends to
infinity, the following holds eventually. There is a sound dart at level
\(k\) for the actual configuration with these properties:

\begin{itemize}
\tightlist
\item
  its jets obey the size bounds;
\item
  its tested second derivatives are within \(\varepsilon\) of the
  tree-level values: total differential weight at most \(4\) for \(J\)
  and at most \(5\) for \(Y\); the derivatives of \(E\) of orders
  \(2\)--\(4\) have size at most \(\varepsilon\).
\end{itemize}

\emph{Proof sketch.} Choose \(i\) with
\(2^{-(i+1)}C(\hat\beta_k)<\varepsilon\), where \(C(\hat\beta_k)\)
bounds the finitely many relevant derivative orders. Darts at level
\(k\) are obtained by transporting darts from level \(k+i+1\) through
\(i+1\) steps. Each step multiplies the deviation from tree level by at
most \(\frac12\). The resulting geometric bound is converted into the
required estimates. \emph{Lean:}
\href{https://github.com/yang-mills-group/OpenAI-YM-20261028/blob/3456e5e75cf8d4cd558955a5c7638fba2d68f464/YM/TraceExtension/Perfect/Pallas.lean\#L46-L64}{\texttt{treasury}},
\href{https://github.com/yang-mills-group/OpenAI-YM-20261028/blob/3456e5e75cf8d4cd558955a5c7638fba2d68f464/YM/TraceExtension/Perfect/Pallas.lean\#L38-L40}{\texttt{cash\_viz}};
\href{https://github.com/yang-mills-group/OpenAI-YM-20261028/blob/3456e5e75cf8d4cd558955a5c7638fba2d68f464/YM/TraceExtension/Perfect/Flo.lean\#L41-L46}{\texttt{through}},
\href{https://github.com/yang-mills-group/OpenAI-YM-20261028/blob/3456e5e75cf8d4cd558955a5c7638fba2d68f464/YM/TraceExtension/Perfect/Native.lean\#L81-L88}{\texttt{pay\_step}}.

\paragraph{4.3 The one-loop slope}\label{the-one-loop-slope}

\textbf{Lemma VII.2.f (logarithmic pole).} Let
\(\Gamma_\mu(U)=\frac12\log\det(\mathcal H(U)+MM^{\mathsf T}+\mu^2)-\log\det(\mathcal G(U)+\mu^2)\)
be the infrared-regulated one-loop effective action of a single RG
level. Here \(\mathcal H\) is the Hessian of the action,
\(MM^{\mathsf T}\) the gauge-fixing term, and \(\mathcal G\) the ghost
operator. Let \(R_\mu=\sum_{d\ne i}y_{di}(\Gamma_\mu)\), where
\(y_{di}\) is the quadratic color-\(q\), polarization-\(i\) momentum
test in direction \(d\), divided by the number of sites. It applies the
centered second difference to the quadratic field derivative at the
opposite momenta \(\pm h e_d\), with the factor \(h^{-2}\), where
\(h=2\pi/f\) for torus side \(f\). For a sequence of lattices of side
\(\to\infty\) satisfying the uniform covariance, resolvent and action
bounds of the one-loop comparison (the predicates in §6), there is
\(C'\) such that for every \(\mu\in(0,1]\), eventually,
\[\big|\mathrm{Re}\,R_\mu+12\,\kappa_G\log(1/\mu)\big|\le C'.\]
\emph{Proof sketch.} By symmetry all \(12\) planes contribute equally.
The \((0,1)\) coefficient converges, along the sequence, to a momentum
integral whose \(\log(1/\mu)\) coefficient is \(-\kappa_G\) up to a
bounded error
(\href{https://github.com/yang-mills-group/OpenAI-YM-20261028/blob/3456e5e75cf8d4cd558955a5c7638fba2d68f464/YM/Next/STATIONARY/OneLoop/MomentumIntegral.lean\#L346-L408}{\texttt{absoluteGeode}},
III.7). \emph{Lean:}
\href{https://github.com/yang-mills-group/OpenAI-YM-20261028/blob/3456e5e75cf8d4cd558955a5c7638fba2d68f464/YM/TraceExtension/Loft/Samen.lean\#L23-L42}{\texttt{slopeL}},
\href{https://github.com/yang-mills-group/OpenAI-YM-20261028/blob/3456e5e75cf8d4cd558955a5c7638fba2d68f464/YM/TraceExtension/Det/Port.lean\#L62}{\texttt{marke}},
\href{https://github.com/yang-mills-group/OpenAI-YM-20261028/blob/3456e5e75cf8d4cd558955a5c7638fba2d68f464/YM/TraceExtension/Loft/Fluent.lean\#L9}{\texttt{ridge}}
\emph{(proof read;
\href{https://github.com/yang-mills-group/OpenAI-YM-20261028/blob/3456e5e75cf8d4cd558955a5c7638fba2d68f464/YM/TraceExtension/Loft/Skis.lean\#L12-L24}{\texttt{effect}},
\href{https://github.com/yang-mills-group/OpenAI-YM-20261028/blob/3456e5e75cf8d4cd558955a5c7638fba2d68f464/YM/TraceExtension/Loft/Drifts.lean\#L32-L35}{\texttt{banana}},
\href{https://github.com/yang-mills-group/OpenAI-YM-20261028/blob/3456e5e75cf8d4cd558955a5c7638fba2d68f464/YM/TraceExtension/Loft/Tips.lean\#L32-L47}{\texttt{payday}},
\href{https://github.com/yang-mills-group/OpenAI-YM-20261028/blob/3456e5e75cf8d4cd558955a5c7638fba2d68f464/YM/TraceExtension/Loft/Gondola.lean\#L31-L35}{\texttt{jury}}
statement only; the momentum selector is defined in §6)}.

\textbf{Lemma VII.2.g (from the pole to a value: mean-value
extraction).} The function \(R\) satisfies
\(\mu\,\partial_\mu R_\mu=-\mathrm{seaD}_\mu\), where
\(\mathrm{seaD}_\mu=\mu^2\sum_{d\ne i}u_{di}(\mu)\) and \(u_{di}\)
collects the same normalized momentum tests of
\(-\mathrm{tr}(\mathcal H+MM^{\mathsf T}+\mu^2)^{-1}+2\,\mathrm{tr}(\mathcal G+\mu^2)^{-1}\)
(from
\(\partial_\mu\log\det(F+\mu^2)=2\mu\,\mathrm{tr}(F+\mu^2)^{-1}\)). By
Lemma VII.2.f, \(x\mapsto\mathrm{Re}\,R_{e^x}-12\kappa_Gx\) is bounded
by \(C'\) eventually for each fixed \(x\le0\). For every \(a<b\le0\),
the mean value theorem on \([a,b]\) then gives \(\mu_w\in[e^a,e^b]\)
with \(|\mathrm{Re}\,\mathrm{seaD}_{\mu_w}-(-12\kappa_G)|\le 2C'/(b-a)\)
eventually. \emph{Lean:}
\href{https://github.com/yang-mills-group/OpenAI-YM-20261028/blob/3456e5e75cf8d4cd558955a5c7638fba2d68f464/YM/TraceExtension/Loft/Fluent.lean\#L47-L60}{\texttt{ridge\_deriv}},
\href{https://github.com/yang-mills-group/OpenAI-YM-20261028/blob/3456e5e75cf8d4cd558955a5c7638fba2d68f464/YM/TraceExtension/Loft/Awake.lean\#L13-L30}{\texttt{coastMark}}.

\textbf{Lemma VII.2.h (the payoff converges to \(-12\kappa_G\)).} Under
the operator bounds of VII.2.f, take translation-covariant sound darts
\(B_w\) on sides \(g(w)\to\infty\), with a fixed positive decay
parameter, uniform jet bounds and precision errors
\(0\le\varepsilon_w\le1\) tending to \(0\). Let \(X_w\) be the summed
normalized momentum test of the Haar divergence of their rooted
displacement fields. Then \(\mathrm{Re}\,X_w\to-12\kappa_G\); the
difference between the finite-momentum test and the Taylor selector,
including its quadratic correction, tends to \(0\), so
\(\mathrm{midPay}_q(B_w)\to-12\kappa_G\). \emph{Proof sketch.} For every
fixed \(0<b'\le c\le1\) and every choice \(\mu_w\in[b',c]\), the Ward
relation gives \(X_w=\mathrm{seaD}_{\mu_w}+\mu_w^2K_w(\mu_w)+o(1)\) with
\(|\mu_w^2K_w(\mu_w)|\le B'\mu_w\) eventually, where \(B'\) is
independent of the fixed regulator interval. Here \(K_w\) is a trace of
resolvents paired with the linearized displacement and ghost terms. The
regulated Ward identity yields the decomposition; its other
action-variation and precision terms vanish in the period limit. Choose
\(u\le0\) with \(B'e^u\le\varepsilon/4\), then \(C\) with
\(2C'/C\le\varepsilon/4\), and apply Lemma VII.2.g on \([u-C,u]\).
\emph{Lean:}
\href{https://github.com/yang-mills-group/OpenAI-YM-20261028/blob/3456e5e75cf8d4cd558955a5c7638fba2d68f464/YM/TraceExtension/Loft/Awake.lean\#L32-L53}{\texttt{awake}},
\href{https://github.com/yang-mills-group/OpenAI-YM-20261028/blob/3456e5e75cf8d4cd558955a5c7638fba2d68f464/YM/TraceExtension/Loft/Leaf.lean\#L40-L52}{\texttt{midAwake}};
\href{https://github.com/yang-mills-group/OpenAI-YM-20261028/blob/3456e5e75cf8d4cd558955a5c7638fba2d68f464/YM/TraceExtension/Lamp/Pious.lean\#L97-L127}{\texttt{saga}},
\href{https://github.com/yang-mills-group/OpenAI-YM-20261028/blob/3456e5e75cf8d4cd558955a5c7638fba2d68f464/YM/TraceExtension/Loft/Peak.lean\#L15-L29}{\texttt{crest\_v}}
\emph{(proof read; their analytic inputs statement only)};
\href{https://github.com/yang-mills-group/OpenAI-YM-20261028/blob/3456e5e75cf8d4cd558955a5c7638fba2d68f464/YM/TraceExtension/Lamp/Pine.lean\#L11-L39}{\texttt{fawn}}.

\textbf{Lemma VII.2.i (uniform certificates; proof of VII.2.3).}

\begin{enumerate}
\def\labelenumi{\arabic{enumi}.}
\tightlist
\item
  At a fixed level \(k\), fix a positive decay parameter and a uniform
  jet bound. For every \(\delta>0\) there are a precision tolerance
  \(\eta>0\) and a volume threshold \(T\) such that every
  translation-covariant dart satisfying those bounds with precision
  \(\eta\), with \(N>k\) and \(p\ge T\), has
  \(|\mathrm{midPay}-(-12\kappa_G)|\le\delta\). The proof is by
  contradiction: a sequence of counterexamples with \(p\to\infty\)
  contradicts Lemma VII.2.h.
\item
  Average the darts of Lemma VII.2.e over lattice translations. The
  bounds are preserved, so this gives darts with payoff accuracy
  \(\varepsilon\) for all large \(k\), eventually along every filter for
  which \(N,p\to\infty\).
\item
  With \(\varepsilon=\hat\beta_k^{-100}\) and the remaining stage,
  support, cluster and coefficient bounds, the work items are chosen by
  choice along \(\mathcal L\).
\end{enumerate}

\emph{Lean:}
\href{https://github.com/yang-mills-group/OpenAI-YM-20261028/blob/3456e5e75cf8d4cd558955a5c7638fba2d68f464/YM/TraceExtension/Loft/Prise.lean\#L13-L28}{\texttt{wire\_mid}},
\href{https://github.com/yang-mills-group/OpenAI-YM-20261028/blob/3456e5e75cf8d4cd558955a5c7638fba2d68f464/YM/TraceExtension/Loft/Prise.lean\#L31-L44}{\texttt{dewWt}},
\href{https://github.com/yang-mills-group/OpenAI-YM-20261028/blob/3456e5e75cf8d4cd558955a5c7638fba2d68f464/YM/TraceExtension/Margin/Stand.lean\#L23-L35}{\texttt{agenda}}.
\emph{(proof read; stage and certificate-supply inputs statement only)}

\paragraph{4.4 From the slope to the coefficient of the
trace}\label{from-the-slope-to-the-coefficient-of-the-trace}

\textbf{Notation.} Fix the run data and a scaling sequence as in
VII.2.4(i). All statements in §§4.4--4.5 have these hypotheses. Write
\(\hat\beta_i=\hat\beta_{\ell_i}\),
\(E_i=\exp(\frac15(\log\log\hat\beta_i)^2)\) and
\(R_i=\exp(\frac16(\log\log\hat\beta_i)^2)\).

\begin{itemize}
\tightlist
\item
  \emph{Reference tensors.} The centered, real degree-\(4\) free tensors
  covariant under signed coordinate permutations form a family
  \(\mathrm{st}(w)\), \(w\in\mathbb R^3\), with
  \(\mathrm{tr}\,\mathrm{st}(w)=4(w_0+4w_1)\mathcal O\)
  (\href{https://github.com/yang-mills-group/OpenAI-YM-20261028/blob/3456e5e75cf8d4cd558955a5c7638fba2d68f464/YM/TraceExtension/Foundation/Dilate.lean\#L67-L69}{\texttt{mechanical\_trace}}).
  The three reference tensors are
  \(L=\mathrm{st}(-(1,-\frac14,1))=T^{\rm free}\),
  \(\Delta=\mathrm{st}(0,\frac14,0)=\delta_{\mu\nu}\mathcal O\) and
  \(\Xi=\mathrm{st}(0,0,1)\).
\item
  \emph{The functional \(\rho\).} On raw degree-\(5\) divergence labels,
  before quotienting by the Maxwell equations,
  \(\rho(\nabla_{\rm raw}\mathrm{st}(w))=-4(w_0+4w_1)\)
  (\href{https://github.com/yang-mills-group/OpenAI-YM-20261028/blob/3456e5e75cf8d4cd558955a5c7638fba2d68f464/YM/TraceExtension/Foundation/Dilate.lean\#L81-L84}{\texttt{rho\_vector}}).
  It annihilates the Maxwell equation direction. Below, \(\rho(B)\)
  abbreviates \(\rho(\nabla_{\rm raw}B)\), so \(\rho(L)=\rho(\Xi)=0\)
  and \(\rho(\Delta)=-4\). Explicitly, \(\Xi_{\mu\mu}=0\) and
  \(\Xi_{\mu\nu}=\sum_{a,c}:F^c_{\mu a}F^c_{\nu a}:\) for \(\mu\ne\nu\).
\item
  \emph{Lattice-built tensors.} \(\mathrm{tub}_i(B)\in\mathcal K(S)\),
  for \(B\in\{L,\Delta,\Xi\}\), are built from the work items at the
  readout level \(\ell_i\).
\item
  \emph{Rescaled divergence.} Eventually,
  \(K_i:=w_ir_i^5\,U_i^{-1}\circ\mathrm{div}\)
  (\href{https://github.com/yang-mills-group/OpenAI-YM-20261028/blob/3456e5e75cf8d4cd558955a5c7638fba2d68f464/YM/TraceExtension/Shelves/Lore.lean\#L40-L44}{\texttt{Optics}}).
\item
  \emph{Rescaling of tensors.} \(\mathrm{Ph}_i:=w_ir_i^4U_i^{-1}\).
\end{itemize}

\textbf{Lemma VII.2.j (the divergence vector of the dart).} Let
\(\Sigma_k\) be the level-5 coefficient vector of the divergence of the
dart's stress tensor at level \(k\), normalized by the coupling limits
and taken along \(\mathcal L\). Then eventually
\(\|\Sigma_k\|\le c\,e^{\frac16(\log\log\hat\beta_k)^2}\), which grows
more slowly than every positive power of \(\hat\beta_k\), and
\(\rho(\Sigma_k)\to-\kappa_G\). \emph{Proof sketch.} At each \((k,N,p)\)
the revised readout decomposes into the payoff term
\(\mathrm{midPay}/12\) plus terms bounded by a fixed multiple of
\(e^{-(\log\log\hat\beta_k)^2}\) and \(\hat\beta_k^{-30}\). Passing to
the ultrafilter limit adds a coefficient-transfer error
\(O(e^{-\frac14(\log\log\hat\beta_k)^2})\). The certificate of VII.2.3
then gives \(|\mathrm{midPay}/12+\kappa_G|\le\hat\beta_k^{-100}/12\).
For the three reference tensors, the analogous raw vectors converge to
\(\nabla_{\rm raw}L\), \(\nabla_{\rm raw}\Delta\) and
\(\nabla_{\rm raw}\Xi\), at rate
\(o(e^{-\frac15(\log\log\hat\beta)^2})\) after applying \(\rho\).
\emph{Lean:}
\href{https://github.com/yang-mills-group/OpenAI-YM-20261028/blob/3456e5e75cf8d4cd558955a5c7638fba2d68f464/YM/TraceExtension/Slate/Pony.lean\#L30-L65}{\texttt{penumbra}},
\href{https://github.com/yang-mills-group/OpenAI-YM-20261028/blob/3456e5e75cf8d4cd558955a5c7638fba2d68f464/YM/TraceExtension/Margin/Beans.lean\#L16-L42}{\texttt{brass}};
\href{https://github.com/yang-mills-group/OpenAI-YM-20261028/blob/3456e5e75cf8d4cd558955a5c7638fba2d68f464/YM/TraceExtension/Slate/Limit.lean\#L33-L53}{\texttt{tiara}},
\href{https://github.com/yang-mills-group/OpenAI-YM-20261028/blob/3456e5e75cf8d4cd558955a5c7638fba2d68f464/YM/TraceExtension/Slate/Limit.lean\#L54-L62}{\texttt{diadem}},
\href{https://github.com/yang-mills-group/OpenAI-YM-20261028/blob/3456e5e75cf8d4cd558955a5c7638fba2d68f464/YM/TraceExtension/Slate/Spite.lean\#L32-L46}{\texttt{article}}.

\textbf{Lemma VII.2.k (divergences at scale \(r_i\)).} There are linear
maps \(s_i\) from the fixed raw coefficient space to
\(\mathcal V(\mathsf F)\), converging to the free projection \(s\), and
functionals \(\mu_i\) on \(\mathcal V(\mathsf F)\) with
\(\|\mu_i\|\le c\) eventually, such that
\(E_i\|\mu_i\circ s_i-\rho\|\to0\) in operator norm, where \(\rho\) is
expressed in the lattice coefficient coordinates. With these:

\begin{itemize}
\tightlist
\item
  \(w_iK_i(\mathrm{tub}_iL)-s_i(\Sigma_{\ell_i})\to0\), hence
  \(\mu_i\big(w_iK_i(\mathrm{tub}_iL)\big)\to-\kappa_G\);
\item
  \(\|w_iK_i(\mathrm{tub}_iL)\|\le CR_i\) eventually;
\item
  for \(B=\Delta,\Xi\):
  \(K_i(\mathrm{tub}_iB)\to\mathrm{div}^{\rm free}B\) and
  \(E_i\big(\mu_i(K_i\mathrm{tub}_iB)-\rho(B)\big)\to0\).
\end{itemize}

\emph{Inputs.}

\begin{itemize}
\tightlist
\item
  The lattice-to-continuum remainders are \(o(\hat\beta_i^{-p})\) for
  every \(p\in\mathbb N_0\). This rests on the stage estimates and on
  the cluster bounds, which are the RG estimates of III.5 and III.8 run
  once more on the trace observables.
\item
  Lattice observables are identified with continuum labels by two tools.
  Pullback along centered chambers: if a lattice observable differs from
  the local representative of \(v_k\) by a chamber-interior term with
  fast chessboard root, then its correlations with displaced spectators
  agree with those of the limit of \(v_k\) up to a fast error.
  Eighth-order difference stencils: the lattice stencil of an observable
  converges to the continuum derivative of its limit.
\end{itemize}

\emph{Lean:}
\href{https://github.com/yang-mills-group/OpenAI-YM-20261028/blob/3456e5e75cf8d4cd558955a5c7638fba2d68f464/YM/TraceExtension/Chart/Dray.lean\#L45-L55}{\texttt{pairingD}},
\href{https://github.com/yang-mills-group/OpenAI-YM-20261028/blob/3456e5e75cf8d4cd558955a5c7638fba2d68f464/YM/TraceExtension/Chart/Dray.lean\#L57-L61}{\texttt{sizeQ\textquotesingle{}}},
\href{https://github.com/yang-mills-group/OpenAI-YM-20261028/blob/3456e5e75cf8d4cd558955a5c7638fba2d68f464/YM/TraceExtension/Chart/Dray.lean\#L63-L77}{\texttt{transferX}},
\href{https://github.com/yang-mills-group/OpenAI-YM-20261028/blob/3456e5e75cf8d4cd558955a5c7638fba2d68f464/YM/TraceExtension/Chart/Token.lean\#L61-L68}{\texttt{beater}},
\href{https://github.com/yang-mills-group/OpenAI-YM-20261028/blob/3456e5e75cf8d4cd558955a5c7638fba2d68f464/YM/TraceExtension/Chart/Token.lean\#L75-L88}{\texttt{mere}}.
The following are statement only:
\href{https://github.com/yang-mills-group/OpenAI-YM-20261028/blob/3456e5e75cf8d4cd558955a5c7638fba2d68f464/YM/TraceExtension/Terrace/Lease.lean\#L27-L30}{\texttt{Prosper}},
\href{https://github.com/yang-mills-group/OpenAI-YM-20261028/blob/3456e5e75cf8d4cd558955a5c7638fba2d68f464/YM/TraceExtension/Slate/Garnet.lean\#L62-L68}{\texttt{rungs}},
\href{https://github.com/yang-mills-group/OpenAI-YM-20261028/blob/3456e5e75cf8d4cd558955a5c7638fba2d68f464/YM/TraceExtension/Slate/Garnet.lean\#L59-L61}{\texttt{legacy}},
\href{https://github.com/yang-mills-group/OpenAI-YM-20261028/blob/3456e5e75cf8d4cd558955a5c7638fba2d68f464/YM/TraceExtension/Slate/Phantom.lean\#L37-L50}{\texttt{hullage}},
\href{https://github.com/yang-mills-group/OpenAI-YM-20261028/blob/3456e5e75cf8d4cd558955a5c7638fba2d68f464/YM/TraceExtension/Shelves/Lore.lean\#L57-L59}{\texttt{gapC}},
\href{https://github.com/yang-mills-group/OpenAI-YM-20261028/blob/3456e5e75cf8d4cd558955a5c7638fba2d68f464/YM/TraceExtension/Terrace/Ell.lean\#L59-L82}{\texttt{crestship}},
\href{https://github.com/yang-mills-group/OpenAI-YM-20261028/blob/3456e5e75cf8d4cd558955a5c7638fba2d68f464/YM/TraceExtension/Terrace/Toil.lean\#L29-L56}{\texttt{touring}},
\href{https://github.com/yang-mills-group/OpenAI-YM-20261028/blob/3456e5e75cf8d4cd558955a5c7638fba2d68f464/YM/TraceExtension/Terrace/Heli.lean\#L39-L51}{\texttt{trod}},
\href{https://github.com/yang-mills-group/OpenAI-YM-20261028/blob/3456e5e75cf8d4cd558955a5c7638fba2d68f464/YM/TraceExtension/Terrace/Terrain.lean\#L26-L51}{\texttt{small\_means}},
\href{https://github.com/yang-mills-group/OpenAI-YM-20261028/blob/3456e5e75cf8d4cd558955a5c7638fba2d68f464/YM/TraceExtension/Terrace/Mortal.lean\#L37-L43}{\texttt{expire}},
\href{https://github.com/yang-mills-group/OpenAI-YM-20261028/blob/3456e5e75cf8d4cd558955a5c7638fba2d68f464/YM/TraceExtension/Cask/Common.lean\#L17-L54}{\texttt{lay}},
\href{https://github.com/yang-mills-group/OpenAI-YM-20261028/blob/3456e5e75cf8d4cd558955a5c7638fba2d68f464/YM/TraceExtension/Shepherd/Spindle.lean\#L46-L51}{\texttt{welfare}},
\href{https://github.com/yang-mills-group/OpenAI-YM-20261028/blob/3456e5e75cf8d4cd558955a5c7638fba2d68f464/YM/TraceExtension/Shepherd/Folly.lean\#L38-L61}{\texttt{garnet}},
\href{https://github.com/yang-mills-group/OpenAI-YM-20261028/blob/3456e5e75cf8d4cd558955a5c7638fba2d68f464/YM/TraceExtension/Cask/Votive.lean\#L42-L67}{\texttt{quaff}},
\href{https://github.com/yang-mills-group/OpenAI-YM-20261028/blob/3456e5e75cf8d4cd558955a5c7638fba2d68f464/YM/TraceExtension/Anchor/Book.lean\#L33-L76}{\texttt{from\_center}},
\href{https://github.com/yang-mills-group/OpenAI-YM-20261028/blob/3456e5e75cf8d4cd558955a5c7638fba2d68f464/YM/TraceExtension/Anchor/Derivative.lean\#L80-L115}{\texttt{axis\_advance}},
\href{https://github.com/yang-mills-group/OpenAI-YM-20261028/blob/3456e5e75cf8d4cd558955a5c7638fba2d68f464/YM/TraceExtension/Anchor/Bleed.lean\#L81-L144}{\texttt{differentiation}}.

\textbf{Lemma VII.2.l (the coefficients of the improvement terms).}
There are \(c_i\in\mathbb R^2\) and \(C\ge0\) such that, eventually:

\begin{itemize}
\tightlist
\item
  \(\|c_i\|\le C\,R_i\);
\item
  \(f_i:c\mapsto c_0K_i(\mathrm{tub}_i\Xi)+c_1K_i(\mathrm{tub}_i\Delta)\)
  is bijective;
\item
  \(w_iK_i(\mathrm{tub}_iL)+f_i(c_i)=0\).
\end{itemize}

Moreover \(c_{i,1}\to-\kappa_G/4=d_2\). \emph{Proof sketch.}

\begin{enumerate}
\def\labelenumi{\arabic{enumi}.}
\tightlist
\item
  The maps \(f_i\) converge to
  \(c\mapsto c_0\,\mathrm{div}\,\Xi+c_1\,\mathrm{div}\,\Delta\), which
  is bijective for simple \(G\). Hence \(f_i\) is eventually bijective
  with uniformly bounded inverse.
\item
  Set \(c_i=-f_i^{-1}\big(w_iK_i(\mathrm{tub}_iL)\big)\).
\item
  Apply \(\mu_i\) to the balance relation and use Lemma VII.2.k:
  \[4c_{i,1}=\mu_i\big(w_iK_i\mathrm{tub}_iL\big)+c_{i,0}\,\mu_i(K_i\mathrm{tub}_i\Xi)+c_{i,1}\big(\mu_i(K_i\mathrm{tub}_i\Delta)+4\big)\to-\kappa_G .\]
  The last two terms tend to \(0\) because \(|c_{i,j}|\le CR_i\) and
  \(R_i/E_i\to0\).
\end{enumerate}

\emph{Lean:}
\href{https://github.com/yang-mills-group/OpenAI-YM-20261028/blob/3456e5e75cf8d4cd558955a5c7638fba2d68f464/YM/TraceExtension/Chart/Laurel.lean\#L53-L78}{\texttt{formulate}},
\href{https://github.com/yang-mills-group/OpenAI-YM-20261028/blob/3456e5e75cf8d4cd558955a5c7638fba2d68f464/YM/TraceExtension/Chart/Limits.lean\#L45-L58}{\texttt{wend}},
\href{https://github.com/yang-mills-group/OpenAI-YM-20261028/blob/3456e5e75cf8d4cd558955a5c7638fba2d68f464/YM/TraceExtension/Chart/Limits.lean\#L68-L80}{\texttt{quill}}.

\textbf{Lemma VII.2.m (the conserved candidate).} Let
\(X_i:=w_i\,\mathrm{tub}_iL+c_{i,0}\,\mathrm{tub}_i\Xi+c_{i,1}\,\mathrm{tub}_i\Delta\).
Then:

\begin{itemize}
\tightlist
\item
  \(X_i\) is conserved eventually, since \(K_iX_i=0\) and \(K_i\)
  vanishes exactly on conserved tensors;
\item
  \(r_i^4U_i^{-1}X_i\to T^{\rm free}\), because
  \(\mathrm{Ph}_i\,\mathrm{tub}_iL\to T^{\rm free}\) and
  \(c_i/w_i\to0\);
\item
  \(w_i\,\mathrm{tr}\big(r_i^4U_i^{-1}X_i\big)\to4d_2\,\mathcal O\),
  because for \(B\in\{L,\Xi\}\) the trace of
  \(\mathrm{Ph}_i\,\mathrm{tub}_iB\) is \(o(\hat\beta^{-p})\) for every
  \(p\), while
  \(\mathrm{tr}\,\mathrm{Ph}_i\,\mathrm{tub}_i\Delta\to4\mathcal O\).
\end{itemize}

\emph{Lean:}
\href{https://github.com/yang-mills-group/OpenAI-YM-20261028/blob/3456e5e75cf8d4cd558955a5c7638fba2d68f464/YM/TraceExtension/Chart/Stage.lean\#L64-L65}{\texttt{serene}},
\href{https://github.com/yang-mills-group/OpenAI-YM-20261028/blob/3456e5e75cf8d4cd558955a5c7638fba2d68f464/YM/TraceExtension/Chart/Stage.lean\#L69-L80}{\texttt{yields}},
\href{https://github.com/yang-mills-group/OpenAI-YM-20261028/blob/3456e5e75cf8d4cd558955a5c7638fba2d68f464/YM/TraceExtension/Chart/Stage.lean\#L82-L95}{\texttt{blossoms}},
\href{https://github.com/yang-mills-group/OpenAI-YM-20261028/blob/3456e5e75cf8d4cd558955a5c7638fba2d68f464/YM/TraceExtension/Chart/Stage.lean\#L51-L53}{\texttt{sum\_delta}};
\href{https://github.com/yang-mills-group/OpenAI-YM-20261028/blob/3456e5e75cf8d4cd558955a5c7638fba2d68f464/YM/TraceExtension/Shelves/Lore.lean\#L51-L55}{\texttt{aprilA}};
\href{https://github.com/yang-mills-group/OpenAI-YM-20261028/blob/3456e5e75cf8d4cd558955a5c7638fba2d68f464/YM/TraceExtension/Shelves/Pear.lean\#L65-L71}{\texttt{nearing\_h}},
\href{https://github.com/yang-mills-group/OpenAI-YM-20261028/blob/3456e5e75cf8d4cd558955a5c7638fba2d68f464/YM/TraceExtension/Shelves/Pear.lean\#L73-L77}{\texttt{muteP}}
\emph{(proof read at the top; inputs statement only)}.

\textbf{Lemma VII.2.n (identification with \(T\); proof of VII.2.6).}

\begin{itemize}
\tightlist
\item
  \emph{Rank.} The free side gives \(\mathrm{div}\) onto with
  one-dimensional kernel; this is reduced by linear algebra to the free
  closure map. Eventually
  \(K_i\circ(c\mapsto c_0\mathrm{tub}_i\Xi+c_1\mathrm{tub}_i\Delta)\) is
  onto the two-dimensional target, and \(\ker K_i=\ker\mathrm{div}\).
  Since \(\dim\mathcal K(S)=\dim\mathcal K(\mathsf F)\), a dimension
  count gives \(\dim\ker\mathrm{div}=1\) on \(\mathcal K(S)\).
\item
  \emph{Identification.} Hence eventually \(X_i=v_i\mathring T\) for
  real \(v_i\).
\item
  \emph{Normalization.} Choose a free observer \(o\) and a label \(u\)
  with \(o(\partial_b u)\ne0\). The rescaled Ward flux
  \(r_i\,\mathrm{flux}_b(\mathring T)\) tested against \(U_iu\)
  converges to \(o(\partial_bu)\), and the same flux for \(X_i\)
  converges to the free value. Comparing the two limits gives
  \(v_i\to1\).
\item
  \emph{Conclusion.} \(r_i^4U_i^{-1}\mathring T\to T^{\rm free}\) and
  \(w_i\,\mathrm{tr}(r_i^4U_i^{-1}\mathring T)\to4d_2\mathcal O\).
\end{itemize}

\emph{Lean:}
\href{https://github.com/yang-mills-group/OpenAI-YM-20261028/blob/3456e5e75cf8d4cd558955a5c7638fba2d68f464/YM/TraceExtension/Chart/Masts.lean\#L41-L51}{\texttt{ranking}},
\href{https://github.com/yang-mills-group/OpenAI-YM-20261028/blob/3456e5e75cf8d4cd558955a5c7638fba2d68f464/YM/TraceExtension/Chart/Masts.lean\#L58-L65}{\texttt{reveals}},
\href{https://github.com/yang-mills-group/OpenAI-YM-20261028/blob/3456e5e75cf8d4cd558955a5c7638fba2d68f464/YM/TraceExtension/Chart/North.lean\#L28-L65}{\texttt{year\_spec}},
\href{https://github.com/yang-mills-group/OpenAI-YM-20261028/blob/3456e5e75cf8d4cd558955a5c7638fba2d68f464/YM/Further/PHYSICAL_PERIOD_LIMIT_F933/FreeRanks.lean\#L159-L165}{\texttt{free\_rank}};
\href{https://github.com/yang-mills-group/OpenAI-YM-20261028/blob/3456e5e75cf8d4cd558955a5c7638fba2d68f464/YM/Further/PHYSICAL_PERIOD_LIMIT_F933/FreeRanks.lean\#L83-L113}{\texttt{free\_closure}};
the Ward-flux comparison inputs \emph{(statement only; VI.1)}.

\textbf{Lemma VII.2.o (the field-level anomaly; proof of VII.2.4).}

\begin{enumerate}
\def\labelenumi{\arabic{enumi}.}
\tightlist
\item
  \emph{Centering is invisible at short distance.} For any phase table
  \(T'\), \(\Theta_{T'}=\mathrm{tr}\,\mathring T-x\,\mathbf 1\) for some
  constant \(x\). Since \(w_ir_i^4\to0\), the term \(x\mathbf 1\)
  disappears after rescaling. This gives
  \(w_ir_i^4U_i^{-1}\Theta_{T'}\to4d_2\mathcal O\).
\item
  \emph{Form of \(\Theta\).} \(\Theta_{T'}\) is a real \(O(4)\)-scalar.
  Eventually \(U_i^{-1}\) maps it into the real degree-\(4\) free labels
  fixed by signed coordinate permutations, which are spanned by
  \(\mathbf 1\) and \(\mathcal O\). This gives the stated form with
  coefficients \(e_i,d_i\). The two basis labels are linearly
  independent; applying their coordinate functionals to the limit gives
  \(w_ie_i\to0\) and \(w_id_i\to d_2\).
\item
  \emph{Clock.} \(\log(1/r_i)/w_i\to\kappa_G^{-1}\) holds because
  \(r_i=a_{k_i}m_i\) with \(\log m_i=o(\hat\beta_{k_i})\),
  \(\hat\beta_k=\bar\beta+\kappa_G\log(1/a_k)\), and
  \(w_i/\hat\beta_{k_i}\to1\).
\item
  \emph{Gallery property.} The free Maxwell comparison supplies the
  hierarchy \(\mathsf a\) and its Ward-normalized stress tensor;
  compactness along the ultrafilter supplies \(Z\). The preceding
  conclusions apply to every compatible scaling sequence. The
  subsequence statement comes from VII.6.
\end{enumerate}

\emph{Lean:}
\href{https://github.com/yang-mills-group/OpenAI-YM-20261028/blob/3456e5e75cf8d4cd558955a5c7638fba2d68f464/YM/TraceExtension/Chart/Nectar.lean\#L79-L84}{\texttt{july}},
\href{https://github.com/yang-mills-group/OpenAI-YM-20261028/blob/3456e5e75cf8d4cd558955a5c7638fba2d68f464/YM/TraceExtension/Chart/Nectar.lean\#L69-L72}{\texttt{shift\_unit}},
\href{https://github.com/yang-mills-group/OpenAI-YM-20261028/blob/3456e5e75cf8d4cd558955a5c7638fba2d68f464/YM/TraceExtension/Chart/Time.lean\#L41-L51}{\texttt{leaderB}},
\href{https://github.com/yang-mills-group/OpenAI-YM-20261028/blob/3456e5e75cf8d4cd558955a5c7638fba2d68f464/YM/TraceExtension/Chart/Orr.lean\#L42-L62}{\texttt{zoom\_BQ}},
\href{https://github.com/yang-mills-group/OpenAI-YM-20261028/blob/3456e5e75cf8d4cd558955a5c7638fba2d68f464/YM/TraceExtension/Chart/Orr.lean\#L73-L76}{\texttt{band\_cost}},
\href{https://github.com/yang-mills-group/OpenAI-YM-20261028/blob/3456e5e75cf8d4cd558955a5c7638fba2d68f464/YM/TraceExtension/Chart/Time.lean\#L63-L68}{\texttt{Store\_l}},
\href{https://github.com/yang-mills-group/OpenAI-YM-20261028/blob/3456e5e75cf8d4cd558955a5c7638fba2d68f464/YM/TraceExtension/Chart/Line.lean\#L49-L56}{\texttt{gang}};
\href{https://github.com/yang-mills-group/OpenAI-YM-20261028/blob/3456e5e75cf8d4cd558955a5c7638fba2d68f464/YM/TraceExtension/Chart/Scalar.lean\#L72-L81}{\texttt{payload}},
\href{https://github.com/yang-mills-group/OpenAI-YM-20261028/blob/3456e5e75cf8d4cd558955a5c7638fba2d68f464/YM/TraceExtension/Terrace/Dusk.lean\#L72-L75}{\texttt{styling}}.

\paragraph{4.5 The short-distance law}\label{the-short-distance-law}

\textbf{Lemma VII.2.p (law along an anomaly sequence at fixed
configurations).} Let \(x=(x_0,x_1)\) have distinct coordinates in every
axis and lie in the interior of the fixed scaling window with bound
\(100\). Let \(x^{(i)}\to x\) be its lattice roundings at level \(i\).
Then
\[r_i^8\log^2(1/r_i)\;G_\Theta\big(r_i|x_0^{(i)}-x_1^{(i)}|\big)\to\frac{A_G}{|x_0-x_1|^8}.\]
\emph{Proof sketch.}

\begin{enumerate}
\def\labelenumi{\arabic{enumi}.}
\tightlist
\item
  Write \(U_i(a_i)=w_ir_i^4\Theta\), where
  \(a_i=U_i^{-1}(w_ir_i^4\Theta)\).
\item
  The two-point function in degree 4 is quadratic,
  \(\langle(pu)(pu)\rangle=p^2\langle uu\rangle\).
\item
  By \(O(4)\) and translation invariance, any pair of points is moved to
  the time axis, so the left side equals
  \((\log(1/r_i)/w_i)^2\,S^{\xi_*}\big(U_ia_i(r_ix_0^{(i)})\,U_ia_i(r_ix_1^{(i)})\big)\).
\item
  Since \(a_i\to4d_2\mathcal O\), the free-field convergence and the
  normalization of \(\mathcal O\) give the limit
  \(\kappa_G^{-2}(4d_2)^2\cdot3\dim\mathfrak g/(\pi^4|x_0-x_1|^8)=A_G/|x_0-x_1|^8\).
\end{enumerate}

\emph{Lean:}
\href{https://github.com/yang-mills-group/OpenAI-YM-20261028/blob/3456e5e75cf8d4cd558955a5c7638fba2d68f464/YM/TraceExtension/Verse/Day.lean\#L50-L69}{\texttt{tagline}},
\href{https://github.com/yang-mills-group/OpenAI-YM-20261028/blob/3456e5e75cf8d4cd558955a5c7638fba2d68f464/YM/TraceExtension/Verse/Day.lean\#L18-L37}{\texttt{axis\_shape}},
\href{https://github.com/yang-mills-group/OpenAI-YM-20261028/blob/3456e5e75cf8d4cd558955a5c7638fba2d68f464/YM/TraceExtension/Foundation/Constant.lean\#L38-L52}{\texttt{wick\_norm}};
\href{https://github.com/yang-mills-group/OpenAI-YM-20261028/blob/3456e5e75cf8d4cd558955a5c7638fba2d68f464/YM/TraceExtension/Verse/Round.lean\#L51-L63}{\texttt{rainbow}},
\href{https://github.com/yang-mills-group/OpenAI-YM-20261028/blob/3456e5e75cf8d4cd558955a5c7638fba2d68f464/YM/TraceExtension/Verse/Round.lean\#L16-L23}{\texttt{rounds}}
\emph{(proof read at the top)}.

\textbf{Lemma VII.2.q (removing the rounding; the full limit).}

\begin{enumerate}
\def\labelenumi{\arabic{enumi}.}
\item
  \emph{Removing the rounding.} Take
  \(x=(0,(\frac s2,\frac s2,\frac s2,\frac s2))\) with
  \(s\in(\frac12,2)\), so that \(|x_0-x_1|=s\). Since \(G_\Theta\) is
  nonincreasing (VII.2.2) and the rounded separations tend to \(s\):

  \begin{itemize}
  \tightlist
  \item
    \(\liminf_i r_i^8\log^2(1/r_i)\,G_\Theta(r_i)\ge A_Gs^{-8}\) for
    \(1<s<2\);
  \item
    \(\limsup_i r_i^8\log^2(1/r_i)\,G_\Theta(r_i)\le A_Gs^{-8}\) for
    \(1/2<s<1\).
  \end{itemize}

  Letting \(s\to1\) gives the limit \(A_G\) along \((r_i)\).
\item
  \emph{The full limit.} If the law failed as \(r\downarrow0\), there
  would be a null sequence avoiding a neighbourhood of \(A_G\). By
  VII.2.4(ii), a subsequence satisfies the field-level anomaly, which
  contradicts step 1.
\end{enumerate}

\emph{Lean:}
\href{https://github.com/yang-mills-group/OpenAI-YM-20261028/blob/3456e5e75cf8d4cd558955a5c7638fba2d68f464/YM/TraceExtension/Verse/Cape.lean\#L52-L76}{\texttt{border}},
\href{https://github.com/yang-mills-group/OpenAI-YM-20261028/blob/3456e5e75cf8d4cd558955a5c7638fba2d68f464/YM/TraceExtension/Verse/Cape.lean\#L77-L88}{\texttt{outside}}.

\textbf{Lemma VII.2.r (Tauberian step; proof of VII.2.5(b)).} Suppose
VII.2.2 holds and \(r^8\log^2(1/r)G_\Theta(r)\to A\ge0\). Let
\(\nu_\rho:=\rho^{-8}\log^2\rho\;\mu_\Theta(\rho\,\cdot)\). By Lorentz
invariance, the Laplace transforms \(\int e^{-u\cdot p}d\nu_\rho\) for
\(u\) in the open forward cone reduce to time-axis values. The law then
gives their convergence to \(A(u^2)^{-4}\), the Laplace transform of
\(\frac{A}{1536\pi}(p^2)^2\mathbf 1_{V_+}d^4p\). The law also gives a
uniform energy budget \(\nu_\rho\{p_0\le R\}\le CR^8\) for all \(R\ge1\)
and all sufficiently large \(\rho\), with one constant \(C\). Laplace
convergence gives convergence on continuous compactly supported tests.
The uniform polynomial growth bound controls the tails of Schwartz tests
and extends the limit to that class. \emph{Lean:}
\href{https://github.com/yang-mills-group/OpenAI-YM-20261028/blob/3456e5e75cf8d4cd558955a5c7638fba2d68f464/YM/TraceExtension/LogScaling.lean\#L93-L106}{\texttt{transform\_limit}},
\href{https://github.com/yang-mills-group/OpenAI-YM-20261028/blob/3456e5e75cf8d4cd558955a5c7638fba2d68f464/YM/TraceExtension/LogScaling.lean\#L120-L141}{\texttt{rescaled\_energy}},
\href{https://github.com/yang-mills-group/OpenAI-YM-20261028/blob/3456e5e75cf8d4cd558955a5c7638fba2d68f464/YM/TraceExtension/LogScaling.lean\#L151-L157}{\texttt{compact\_limit}},
\href{https://github.com/yang-mills-group/OpenAI-YM-20261028/blob/3456e5e75cf8d4cd558955a5c7638fba2d68f464/YM/TraceExtension/LogScaling.lean\#L216-L222}{\texttt{joint\_limit}};
\href{https://github.com/yang-mills-group/OpenAI-YM-20261028/blob/3456e5e75cf8d4cd558955a5c7638fba2d68f464/YM/TraceExtension/LaplaceSchwartz.lean\#L89-L99}{\texttt{laplace\_schwartz}},
\href{https://github.com/yang-mills-group/OpenAI-YM-20261028/blob/3456e5e75cf8d4cd558955a5c7638fba2d68f464/YM/TraceExtension/ConeLaplace.lean\#L137-L150}{\texttt{cone\_laplace}}.

\paragraph{4.6 Exclusions}\label{exclusions}

\textbf{Lemma VII.2.s (power--logarithm expansions cannot carry the
law).} Suppose \(f=u+O(1)\) as \(s\downarrow0\), where \(u\) is a finite
real combination of \(s^a\log^k s\) with \(a\in\mathbb Z\) and
\(k\in\mathbb N_0\), and suppose \(s^8\log^2(1/s)\,f(s)\to A\). Then
\(A=0\). \emph{Proof sketch.} The \(O(1)\) part contributes
\(s^8\log^2(1/s)\cdot O(1)\to0\). Along \(s=2^{-n}\),
\(s^8\log^2(1/s)\,u(s)\) is a combination of
\(c_t\,\rho_t^{\,n}n^{k_t+2}\) with \(\rho_t=2^{-(8+a_t)}\). Every term
has a factor \(n^{k_t+2}\) with \(k_t+2\ge2\), so no term tends to a
nonzero constant. Collect equal pairs \((a_t,k_t)\), discard zero
coefficients, and order the remaining terms lexicographically by
\((\rho_t,k_t+2)\). If the leading term has \(\rho_t<1\), all terms tend
to \(0\). Otherwise the leading term diverges and dominates the others,
so a finite limit is impossible unless its coefficient vanishes.
\emph{Lean:}
\href{https://github.com/yang-mills-group/OpenAI-YM-20261028/blob/3456e5e75cf8d4cd558955a5c7638fba2d68f464/YM/TraceExtension/Wick/Log.lean\#L96-L108}{\texttt{no\_powerlog}}.

\textbf{Lemma VII.2.t (Wick profiles expand; proof of VII.2.7).}

\begin{itemize}
\tightlist
\item
  \emph{Expansion.} Each mode kernel \(K_{m,q,h}\) has an expansion to
  all orders in \(\{s^a\log^ks\}\). The class is closed under products
  and sums, so \(\mathrm{Re}\,W\) has the form required in Lemma
  VII.2.s; part (a) follows. Equality on \((0,\delta)\) is formulated as
  eventual equality as \(s\downarrow0\).
\item
  \emph{Part (b).} If the Wightman function equals the mass-shell
  expression on all products with compactly supported wave transform,
  their wave transforms include every compactly supported Schwartz
  momentum test, so \(\int\psi\,d\mu_\Theta\) equals the mass-shell
  expression for every compactly supported \(\psi\). Compactly supported
  Schwartz \(w_n\to e^{-sp_0}\), dominated by \(e^{-sp_0}\), together
  with dominated convergence on both sides, give \(G_\Theta(s)=W(s)\)
  for \(s>0\). The palette version reduces to the collection version.
\end{itemize}

\emph{Lean:}
\href{https://github.com/yang-mills-group/OpenAI-YM-20261028/blob/3456e5e75cf8d4cd558955a5c7638fba2d68f464/YM/TraceExtension/Wick/Expansion.lean\#L47}{\texttt{Expands}},
\href{https://github.com/yang-mills-group/OpenAI-YM-20261028/blob/3456e5e75cf8d4cd558955a5c7638fba2d68f464/YM/TraceExtension/Wick/PolynomialField.lean\#L116-L125}{\texttt{moments\_expand}},
\href{https://github.com/yang-mills-group/OpenAI-YM-20261028/blob/3456e5e75cf8d4cd558955a5c7638fba2d68f464/YM/TraceExtension/Wick/Refute.lean\#L18-L22}{\texttt{real\_profile}},
\href{https://github.com/yang-mills-group/OpenAI-YM-20261028/blob/3456e5e75cf8d4cd558955a5c7638fba2d68f464/YM/TraceExtension/Wick/Refute.lean\#L62-L73}{\texttt{identifies\_euclid}},
\href{https://github.com/yang-mills-group/OpenAI-YM-20261028/blob/3456e5e75cf8d4cd558955a5c7638fba2d68f464/YM/TraceExtension/Wick/Refute.lean\#L86-L105}{\texttt{polynomial\_excludes}}.
\emph{(proof read; the single-mode expansion not read)}

\textbf{Lemma VII.2.u (Fock sectors; proof of VII.2.8).} Assume the
positive short-distance law of VII.2.8 and a comparison \(U\) satisfying
its relative-locality identity.

\begin{enumerate}
\def\labelenumi{\arabic{enumi}.}
\tightlist
\item
  \emph{Sector decomposition.} By translation covariance,
  \(U\Phi_\Theta(\varphi)\Omega=\hat\varphi(p_{\rm tot})\,Y\) for fixed
  measurable sector amplitudes \(Y=(Y_{n,s})\), which need not form a
  square-integrable vector before smearing. Hence \(\mu_\Theta\) is the
  image of \(\sum_{n,s}|Y_{n,s}|^2\,d(\text{photon measures})\) under
  total momentum, and
  \(\mu_\Theta\{p_0\le R\}=\sum_{n,s}\mathrm{mass}(Y_{n,s};\{\sum|k_j|\le R\})\).
\item
  \emph{Evaporation.} The mass of \(Y_{n,s}\) on the box
  \(\{|k_j|\le2^J\text{ for all }j\}\) is at most
  \(\mu_\Theta\{p_0\le(n+1)2^J\}\). By the law, \(2^{-8J}\) times this
  tends to \(0\).
\item
  \emph{Dilation polynomials.} Relative locality gives equality of the
  forward and adjoint sector pairings when each test is supported
  spacelike to the origin. Approximate identities in time show that the
  Cauchy data of their difference are supported at the spatial origin. A
  distribution supported at one point has finite jet order. If \(D_j\)
  doubles the \(j\)th momentum, these jet bounds and mass-shell
  uniqueness give, for some \(K\),
  \(\prod_{a=0}^{K-1}(D_j-2^{a-1})Y_{n,s}=0\) in every slot. The roots
  are distinct and the \(D_j\) commute, so \(Y_{n,s}\) is a finite sum
  of joint eigencomponents.
\item
  \emph{Finitely many sectors.} A joint eigencomponent with slot degrees
  \(a_j\) has box mass proportional to \(R^{2\sum_j a_j}\) at dyadic
  radii \(R\). Evaporation excludes \(\sum_j a_j\ge4\). A component with
  any \(a_j=0\) vanishes by local square integrability near zero
  momentum. Thus every nonzero component has \(a_j\ge1\), so
  \(Y_{n,s}=0\) for \(n\ge4\). The remaining components have box mass at
  most \(C_{n,s}2^{6i}\) at radius \(2^i\).
\item
  \emph{Contradiction.} Summing over the finitely many sectors with
  \(n\le3\) gives \(\mu_\Theta\{p_0\le2^k\}\le D\,2^{7k}\). Multiply
  this bound by \(R^{-8}\log^2R\): it tends to zero. The spectral law
  instead gives a positive limit against a nonnegative compact test
  supported in \(\{0<p_0<1,\ p^2>0\}\), a contradiction.
\end{enumerate}

\emph{Lean:}
\href{https://github.com/yang-mills-group/OpenAI-YM-20261028/blob/3456e5e75cf8d4cd558955a5c7638fba2d68f464/YM/TraceExtension/Fock/Form.lean\#L27-L39}{\texttt{form\_exists}},
\href{https://github.com/yang-mills-group/OpenAI-YM-20261028/blob/3456e5e75cf8d4cd558955a5c7638fba2d68f464/YM/TraceExtension/Fock/Weight.lean\#L76-L78}{\texttt{parts}},
\href{https://github.com/yang-mills-group/OpenAI-YM-20261028/blob/3456e5e75cf8d4cd558955a5c7638fba2d68f464/YM/TraceExtension/Fock/Conclude.lean\#L15-L20}{\texttt{burn}},
\href{https://github.com/yang-mills-group/OpenAI-YM-20261028/blob/3456e5e75cf8d4cd558955a5c7638fba2d68f464/YM/TraceExtension/Fock/Conclude.lean\#L22-L27}{\texttt{consumes}},
\href{https://github.com/yang-mills-group/OpenAI-YM-20261028/blob/3456e5e75cf8d4cd558955a5c7638fba2d68f464/YM/TraceExtension/Fock/Conclude.lean\#L29-L32}{\texttt{finite\_particles}},
\href{https://github.com/yang-mills-group/OpenAI-YM-20261028/blob/3456e5e75cf8d4cd558955a5c7638fba2d68f464/YM/TraceExtension/Fock/Conclude.lean\#L33-L44}{\texttt{saturate}},
\href{https://github.com/yang-mills-group/OpenAI-YM-20261028/blob/3456e5e75cf8d4cd558955a5c7638fba2d68f464/YM/TraceExtension/Fock/SpectrumGap.lean\#L108-L112}{\texttt{local\_eigen}};
\href{https://github.com/yang-mills-group/OpenAI-YM-20261028/blob/3456e5e75cf8d4cd558955a5c7638fba2d68f464/YM/TraceExtension/Fock/VacuumLocal.lean\#L95-L139}{\texttt{local\_position}},
\href{https://github.com/yang-mills-group/OpenAI-YM-20261028/blob/3456e5e75cf8d4cd558955a5c7638fba2d68f464/YM/TraceExtension/Fock/TimeData.lean\#L86-L115}{\texttt{cauchy\_point}},
\href{https://github.com/yang-mills-group/OpenAI-YM-20261028/blob/3456e5e75cf8d4cd558955a5c7638fba2d68f464/YM/TraceExtension/Fock/Branches.lean\#L133-L136}{\texttt{binds}},
\href{https://github.com/yang-mills-group/OpenAI-YM-20261028/blob/3456e5e75cf8d4cd558955a5c7638fba2d68f464/YM/TraceExtension/Fock/Eigen.lean\#L44-L83}{\texttt{splitParts}},
\href{https://github.com/yang-mills-group/OpenAI-YM-20261028/blob/3456e5e75cf8d4cd558955a5c7638fba2d68f464/YM/TraceExtension/Dyadic.lean\#L23-L38}{\texttt{log\_high}},
\href{https://github.com/yang-mills-group/OpenAI-YM-20261028/blob/3456e5e75cf8d4cd558955a5c7638fba2d68f464/YM/TraceExtension/Dyadic.lean\#L65-L82}{\texttt{does\_not\_fit}}
\emph{(proof read; the point-support and analytic uniqueness inputs are
not expanded here)}.

\subsubsection{5. Where the difficulty
lies}\label{where-the-difficulty-lies-26}

\emph{Hard.} Three pieces carry the mathematics, and all three are
needed for the exact constant.

\begin{itemize}
\tightlist
\item
  The transport of the lattice Ward identity through the RG with
  fourth-order matching (VII.2.d--e).
\item
  The one-loop pole with coefficient exactly \(\kappa_G\) (VII.2.f,
  using III.7), and its conversion into a value by the mean-value
  argument (VII.2.g--h).
\item
  The control of lattice-to-continuum remainders to
  \(o(\hat\beta^{-p})\) for every \(p\) (VII.2.k). The trace must be
  controlled after multiplication by the inverse coupling. The separate
  coefficient errors have the slower exponential rates displayed in
  VII.2.k. The bounds on \(c_i\) grow like
  \(e^{\frac16(\log\log\hat\beta)^2}\) and are beaten by decay factors
  \(e^{-\frac15(\log\log\hat\beta)^2}\).
\end{itemize}

The identification of the explicit conserved tensor with \(T\) (VII.2.n)
is short once the rank-one statement and the Ward normalization are
available.

\emph{Routine.}

\begin{itemize}
\tightlist
\item
  Centering and the symmetries of \(\Theta\).
\item
  The passage from subsequential limits at rounded points to the full
  limit (monotonicity and the subsequence principle).
\item
  The Tauberian step.
\item
  The Wick exclusion.
\end{itemize}

The Fock exclusion also needs a distributional argument: relative
locality yields point-supported Cauchy data, which force polynomial
identities for momentum doubling. The growth comparison alone does not
supply this step.

\emph{Unusual or worth checking.}

\begin{enumerate}
\def\labelenumi{\arabic{enumi}.}
\item
  \emph{Logic of dependence.} The law is derived from the gallery
  property, which is established while the ultraviolet run is chosen
  (III.2): the same Wilson family must carry the trace certificates. The
  exclusions are proved from the law, not independently. Their public
  form combines the two.
\item
  \emph{The cancellation of \(b_0\) is a consistency check, not a
  tautology.} The identities \(\kappa_G=2b_0\) and \(d_2=-\kappa_G/4\)
  are definitions. The content is:

  \begin{itemize}
  \tightlist
  \item
    \(c_{i,1}\to-\kappa_G/4\), where \(\kappa_G\) comes from the
    one-loop integral (VII.2.f);
  \item
    \(\log(1/r)/w\to1/\kappa_G\), where \(\kappa_G\) is the slope of the
    coupling schedule.
  \end{itemize}

  Only because both are \(\kappa_G\) is the amplitude
  \(A_G=(4d_2/\kappa_G)^2\cdot3\dim\mathfrak g/\pi^4=3\dim\mathfrak g/\pi^4\),
  independent of \(C_A\). Had the one-loop slope differed from the
  schedule, the amplitude would differ by the square of their ratio.
\item
  \emph{Geometry of the points.} The law is first obtained at the
  generic configuration \((0,\frac s2(1,1,1,1))\), as the lattice
  machinery requires distinct coordinates in each axis. It is then
  rotated to the time axis, which uses the full \(O(4)\) invariance of
  \(\Theta\) from VI.1.
\item
  \emph{Precision.} The payoff certificate is accurate to
  \(\hat\beta_k^{-100}\), much more than the argument needs.
\item
  \emph{Scope of the exclusions.} The Wick class is exactly the public
  finite-profile class. The conclusion does not exclude continuous-mass
  generalized free fields. The Fock comparison is tested through vacuum
  matrix elements with coherent states; it gives an obstruction to
  relative locality with free photons under the stated embedding
  hypotheses.
\end{enumerate}

\subsubsection{6. Notation}\label{notation-26}

{\def\LTcaptype{none} % do not increment counter
\begin{longtable}[]{@{}
  >{\raggedright\arraybackslash}p{(\linewidth - 4\tabcolsep) * \real{0.3333}}
  >{\raggedright\arraybackslash}p{(\linewidth - 4\tabcolsep) * \real{0.3333}}
  >{\raggedright\arraybackslash}p{(\linewidth - 4\tabcolsep) * \real{0.3333}}@{}}
\toprule\noalign{}
\begin{minipage}[b]{\linewidth}\raggedright
object
\end{minipage} & \begin{minipage}[b]{\linewidth}\raggedright
code name
\end{minipage} & \begin{minipage}[b]{\linewidth}\raggedright
notation
\end{minipage} \\
\midrule\noalign{}
\endhead
\bottomrule\noalign{}
\endlastfoot
trace of the stress tensor &
\href{https://github.com/yang-mills-group/OpenAI-YM-20261028/blob/3456e5e75cf8d4cd558955a5c7638fba2d68f464/YM/TraceExtension/Centered.lean\#L17}{\texttt{stressTrace}}
& \(\sum_\mu T_{\mu\mu}\) \\
one-point function in a table &
\href{https://github.com/yang-mills-group/OpenAI-YM-20261028/blob/3456e5e75cf8d4cd558955a5c7638fba2d68f464/YM/TraceExtension/Centered.lean\#L34}{\texttt{euclideanVEV}}
& \(\langle v\rangle_{\xi}\) \\
centered trace, its pinned label &
\href{https://github.com/yang-mills-group/OpenAI-YM-20261028/blob/3456e5e75cf8d4cd558955a5c7638fba2d68f464/YM/TraceExtension/Centered.lean\#L64}{\texttt{theta}},
\href{https://github.com/yang-mills-group/OpenAI-YM-20261028/blob/3456e5e75cf8d4cd558955a5c7638fba2d68f464/YM/TraceExtension/WilsonTrace.lean\#L48}{\texttt{center}},
\href{https://github.com/yang-mills-group/OpenAI-YM-20261028/blob/3456e5e75cf8d4cd558955a5c7638fba2d68f464/YM/TraceExtension/WilsonTrace.lean\#L49}{\texttt{pinned}}
& \(\Theta\), \(\Theta_{T'}\), \(\iota(\Theta)\) \\
time-axis two-point function &
\href{https://github.com/yang-mills-group/OpenAI-YM-20261028/blob/3456e5e75cf8d4cd558955a5c7638fba2d68f464/YM/TraceExtension/WilsonTrace.lean\#L53}{\texttt{euclidean}}
& \(G_\Theta(r)\) \\
spectral measure &
\href{https://github.com/yang-mills-group/OpenAI-YM-20261028/blob/3456e5e75cf8d4cd558955a5c7638fba2d68f464/YM/TraceExtension/WilsonTrace.lean\#L83-L94}{\texttt{TraceSpectrum}}
& \(\mu_\Theta\) \\
cone measure &
\href{https://github.com/yang-mills-group/OpenAI-YM-20261028/blob/3456e5e75cf8d4cd558955a5c7638fba2d68f464/YM/TraceExtension/ConeLaplace.lean\#L57}{\texttt{coneMeasure}}
& \(\frac{A}{1536\pi}(p^2)^2\mathbf 1_{V_+}d^4p\) \\
scale factor &
\href{https://github.com/yang-mills-group/OpenAI-YM-20261028/blob/3456e5e75cf8d4cd558955a5c7638fba2d68f464/YM/TraceExtension/LogScaling.lean\#L20}{\texttt{Zlog}},
\href{https://github.com/yang-mills-group/OpenAI-YM-20261028/blob/3456e5e75cf8d4cd558955a5c7638fba2d68f464/YM/TraceExtension/Verse/Day.lean\#L39}{\texttt{clue}}
& \(\rho^{-8}\log^2\rho\); \(r^8\log^2(1/r)\) \\
slope, one-loop coefficient, force &
\href{https://github.com/yang-mills-group/OpenAI-YM-20261028/blob/3456e5e75cf8d4cd558955a5c7638fba2d68f464/YM/Gauge/Colors.lean\#L184}{\texttt{slopeConstant}},
\href{https://github.com/yang-mills-group/OpenAI-YM-20261028/blob/3456e5e75cf8d4cd558955a5c7638fba2d68f464/YM/TraceExtension/Chart/Nectar.lean\#L60}{\texttt{b0}},
\href{https://github.com/yang-mills-group/OpenAI-YM-20261028/blob/3456e5e75cf8d4cd558955a5c7638fba2d68f464/YM/TraceExtension/Chart/Nectar.lean\#L62}{\texttt{d2}}
& \(\kappa_G\), \(b_0\), \(d_2\) \\
amplitude; payoff constant &
\href{https://github.com/yang-mills-group/OpenAI-YM-20261028/blob/3456e5e75cf8d4cd558955a5c7638fba2d68f464/YM/TraceExtension/Verse/Day.lean\#L41}{\texttt{award}};
\href{https://github.com/yang-mills-group/OpenAI-YM-20261028/blob/3456e5e75cf8d4cd558955a5c7638fba2d68f464/YM/TraceExtension/Loft/Awake.lean\#L9}{\texttt{landPay}}
& \(A_G\); \(-12\kappa_G\) \\
free table; \(\frac14F^2\) label &
\href{https://github.com/yang-mills-group/OpenAI-YM-20261028/blob/3456e5e75cf8d4cd558955a5c7638fba2d68f464/YM/Climb/Cells.lean\#L20}{\texttt{Climb.Bel}},
\href{https://github.com/yang-mills-group/OpenAI-YM-20261028/blob/3456e5e75cf8d4cd558955a5c7638fba2d68f464/YM/Cells/Corrected.lean\#L22}{\texttt{Cells.standard}};
\href{https://github.com/yang-mills-group/OpenAI-YM-20261028/blob/3456e5e75cf8d4cd558955a5c7638fba2d68f464/YM/TraceExtension/Chart/Stage.lean\#L50}{\texttt{OScale}},
\href{https://github.com/yang-mills-group/OpenAI-YM-20261028/blob/3456e5e75cf8d4cd558955a5c7638fba2d68f464/YM/TraceExtension/Foundation/Dilate.lean\#L66}{\texttt{square4}}
& \(\mathsf F\); \(\mathcal O\) \\
free tensors, divergence functional &
\href{https://github.com/yang-mills-group/OpenAI-YM-20261028/blob/3456e5e75cf8d4cd558955a5c7638fba2d68f464/YM/TraceExtension/Shelves/Pear.lean\#L50}{\texttt{materialB}},
\href{https://github.com/yang-mills-group/OpenAI-YM-20261028/blob/3456e5e75cf8d4cd558955a5c7638fba2d68f464/YM/TraceExtension/Slate/April.lean\#L27}{\texttt{BornZ}},
\href{https://github.com/yang-mills-group/OpenAI-YM-20261028/blob/3456e5e75cf8d4cd558955a5c7638fba2d68f464/YM/TraceExtension/Foundation/Dilate.lean\#L78-L80}{\texttt{rho}},
\href{https://github.com/yang-mills-group/OpenAI-YM-20261028/blob/3456e5e75cf8d4cd558955a5c7638fba2d68f464/YM/TraceExtension/Terrace/Lease.lean\#L22}{\texttt{Lease}}
& \(T^{\rm free}=L\), \(\Delta\), \(\Xi\); \(\rho\) \\
scaling sequence; coupling limits &
\href{https://github.com/yang-mills-group/OpenAI-YM-20261028/blob/3456e5e75cf8d4cd558955a5c7638fba2d68f464/YM/Climb/Cells.lean\#L24-L31}{\texttt{Climb.Scaling}}
(\(r_i\), \(U_i\));
\href{https://github.com/yang-mills-group/OpenAI-YM-20261028/blob/3456e5e75cf8d4cd558955a5c7638fba2d68f464/YM/TraceExtension/Terrace/Dusk.lean\#L65-L67}{\texttt{Arch}},
\href{https://github.com/yang-mills-group/OpenAI-YM-20261028/blob/3456e5e75cf8d4cd558955a5c7638fba2d68f464/YM/TraceExtension/Chart/Dray.lean\#L42}{\texttt{BQ}}
& \((r_i,U_i)\); \(Z.\beta(k)\), \(w_i\) \\
rescaled fields &
\href{https://github.com/yang-mills-group/OpenAI-YM-20261028/blob/3456e5e75cf8d4cd558955a5c7638fba2d68f464/YM/TraceExtension/Chart/Nectar.lean\#L64}{\texttt{ZoomValue}},
\href{https://github.com/yang-mills-group/OpenAI-YM-20261028/blob/3456e5e75cf8d4cd558955a5c7638fba2d68f464/YM/TraceExtension/Chart/Stage.lean\#L55}{\texttt{Year}},
\texttt{Phase4},
\href{https://github.com/yang-mills-group/OpenAI-YM-20261028/blob/3456e5e75cf8d4cd558955a5c7638fba2d68f464/YM/TraceExtension/Shelves/Lore.lean\#L40-L44}{\texttt{Optics}}
& \(U_i^{-1}(w_ir_i^4\,\cdot)\), \(r_i^4U_i^{-1}\), \(\mathrm{Ph}_i\),
\(K_i\) \\
lattice-built tensors; candidate &
\href{https://github.com/yang-mills-group/OpenAI-YM-20261028/blob/3456e5e75cf8d4cd558955a5c7638fba2d68f464/YM/TraceExtension/Chart/Token.lean\#L59}{\texttt{tub}},
\href{https://github.com/yang-mills-group/OpenAI-YM-20261028/blob/3456e5e75cf8d4cd558955a5c7638fba2d68f464/YM/TraceExtension/Chart/Laurel.lean\#L48}{\texttt{pollen}};
\href{https://github.com/yang-mills-group/OpenAI-YM-20261028/blob/3456e5e75cf8d4cd558955a5c7638fba2d68f464/YM/TraceExtension/Chart/Stage.lean\#L61}{\texttt{Cloak}}
& \(\mathrm{tub}_i(B)\); \(X_i\) \\
divergence vector of the dart &
\href{https://github.com/yang-mills-group/OpenAI-YM-20261028/blob/3456e5e75cf8d4cd558955a5c7638fba2d68f464/YM/TraceExtension/Slate/Spite.lean\#L49}{\texttt{Shine5}}
& \(\Sigma_k\) \\
lattice Ward dart; precision &
\href{https://github.com/yang-mills-group/OpenAI-YM-20261028/blob/3456e5e75cf8d4cd558955a5c7638fba2d68f464/YM/TraceExtension/Perfect/Badge.lean\#L22-L27}{\texttt{Dart}},
\href{https://github.com/yang-mills-group/OpenAI-YM-20261028/blob/3456e5e75cf8d4cd558955a5c7638fba2d68f464/YM/TraceExtension/Perfect/Badge.lean\#L82-L90}{\texttt{Dart.Sound}};
\href{https://github.com/yang-mills-group/OpenAI-YM-20261028/blob/3456e5e75cf8d4cd558955a5c7638fba2d68f464/YM/TraceExtension/Perfect/Pallas.lean\#L15-L23}{\texttt{Viz}}
& dart \((J,Y,E,Q,\mathrm{tw})\) \\
regulated one-loop action and coefficient &
\href{https://github.com/yang-mills-group/OpenAI-YM-20261028/blob/3456e5e75cf8d4cd558955a5c7638fba2d68f464/YM/TraceExtension/Det/Port.lean\#L62}{\texttt{marke}},
\href{https://github.com/yang-mills-group/OpenAI-YM-20261028/blob/3456e5e75cf8d4cd558955a5c7638fba2d68f464/YM/TraceExtension/Loft/Fluent.lean\#L9}{\texttt{ridge}},
\href{https://github.com/yang-mills-group/OpenAI-YM-20261028/blob/3456e5e75cf8d4cd558955a5c7638fba2d68f464/YM/TraceExtension/Loft/Awake.lean\#L10}{\texttt{seaD}}
& \(\Gamma_\mu\), \(R_\mu\), \(\mathrm{seaD}_\mu\) \\
payoff of a dart &
\href{https://github.com/yang-mills-group/OpenAI-YM-20261028/blob/3456e5e75cf8d4cd558955a5c7638fba2d68f464/YM/TraceExtension/Loft/Leaf.lean\#L38}{\texttt{midPay}},
\href{https://github.com/yang-mills-group/OpenAI-YM-20261028/blob/3456e5e75cf8d4cd558955a5c7638fba2d68f464/YM/TraceExtension/Loft/Leaf.lean\#L37}{\texttt{Holt}}
& \(\mathrm{midPay}_q\) \\
work items and certificate &
\href{https://github.com/yang-mills-group/OpenAI-YM-20261028/blob/3456e5e75cf8d4cd558955a5c7638fba2d68f464/YM/TraceExtension/Margin/Stand.lean\#L15}{\texttt{Work\ items}},
\href{https://github.com/yang-mills-group/OpenAI-YM-20261028/blob/3456e5e75cf8d4cd558955a5c7638fba2d68f464/YM/TraceExtension/Margin/Stand.lean\#L16}{\texttt{Grant}},
\href{https://github.com/yang-mills-group/OpenAI-YM-20261028/blob/3456e5e75cf8d4cd558955a5c7638fba2d68f464/YM/TraceExtension/Margin/Ansel.lean\#L30-L40}{\texttt{WorthB}}
& certificates along \(\mathcal L\) \\
slowly varying rates &
\href{https://github.com/yang-mills-group/OpenAI-YM-20261028/blob/3456e5e75cf8d4cd558955a5c7638fba2d68f464/YM/TraceExtension/Zn/Have.lean\#L10}{\texttt{Wave}},
\href{https://github.com/yang-mills-group/OpenAI-YM-20261028/blob/3456e5e75cf8d4cd558955a5c7638fba2d68f464/YM/TraceExtension/Terrace/Look.lean\#L10}{\texttt{AgeK}}
& \(e^{a(\log\log\hat\beta_k)^2}\) \\
field-level anomaly; its bundle &
\href{https://github.com/yang-mills-group/OpenAI-YM-20261028/blob/3456e5e75cf8d4cd558955a5c7638fba2d68f464/YM/TraceExtension/Chart/Time.lean\#L52-L59}{\texttt{Scribe}},
\href{https://github.com/yang-mills-group/OpenAI-YM-20261028/blob/3456e5e75cf8d4cd558955a5c7638fba2d68f464/YM/TraceExtension/Chart/Line.lean\#L41-L44}{\texttt{Gallery}}
& Scribe property; gallery property \\
short-distance law &
\href{https://github.com/yang-mills-group/OpenAI-YM-20261028/blob/3456e5e75cf8d4cd558955a5c7638fba2d68f464/YM/TraceExtension/LogScaling.lean\#L74-L75}{\texttt{EuclideanAsymptotic}}
& \(r^8\log^2(1/r)G_\Theta(r)\to A\) \\
Wick modes and tracks &
\href{https://github.com/yang-mills-group/OpenAI-YM-20261028/blob/3456e5e75cf8d4cd558955a5c7638fba2d68f464/YM/TraceExtension/Wick/PolynomialField.lean\#L23-L29}{\texttt{Mode}},
\href{https://github.com/yang-mills-group/OpenAI-YM-20261028/blob/3456e5e75cf8d4cd558955a5c7638fba2d68f464/YM/TraceExtension/Wick/PolynomialField.lean\#L90-L92}{\texttt{Track}}
& \(K_{m,q,h}\), \(\prod_\beta K\) \\
Fock comparison; relative locality &
\href{https://github.com/yang-mills-group/OpenAI-YM-20261028/blob/3456e5e75cf8d4cd558955a5c7638fba2d68f464/YM/TraceExtension/Fock/Identify.lean\#L23-L27}{\texttt{Comparison}},
\href{https://github.com/yang-mills-group/OpenAI-YM-20261028/blob/3456e5e75cf8d4cd558955a5c7638fba2d68f464/YM/TraceExtension/Fock/Identify.lean\#L32-L35}{\texttt{Nfl}},
\href{https://github.com/yang-mills-group/OpenAI-YM-20261028/blob/3456e5e75cf8d4cd558955a5c7638fba2d68f464/YM/TraceExtension/Fock/Space.lean\#L111}{\texttt{Ambient}},
\href{https://github.com/yang-mills-group/OpenAI-YM-20261028/blob/3456e5e75cf8d4cd558955a5c7638fba2d68f464/YM/TraceExtension/Fock/Weyl.lean\#L119}{\texttt{coherent}}
& \(U:\mathcal H\to\mathcal F_m\); \(W(h)\Omega_{\mathcal F}\) \\
output bundles &
\href{https://github.com/yang-mills-group/OpenAI-YM-20261028/blob/3456e5e75cf8d4cd558955a5c7638fba2d68f464/YM/Nonperturbative/HeadlineSpec.lean\#L84-L120}{\texttt{TraceExclusions}},
\href{https://github.com/yang-mills-group/OpenAI-YM-20261028/blob/3456e5e75cf8d4cd558955a5c7638fba2d68f464/YM/TraceExtension/Belfry.lean\#L16-L30}{\texttt{Crest}},
\href{https://github.com/yang-mills-group/OpenAI-YM-20261028/blob/3456e5e75cf8d4cd558955a5c7638fba2d68f464/YM/TraceExtension/Endpoint.lean\#L99-L120}{\texttt{Result}},
\href{https://github.com/yang-mills-group/OpenAI-YM-20261028/blob/3456e5e75cf8d4cd558955a5c7638fba2d68f464/YM/Integration/ManuscriptTraceTests.lean\#L62-L75}{\texttt{Tests}}
& (T1)--(T4) and internal variants \\
\end{longtable}
}

For the coefficient selectors, let \(M_{iq}^H(p_1,p_2)\) be the
quadratic field derivative of \(H\) at zero, complexified and evaluated
on two plane waves of polarization \(i\) and color \(q\), using the
rooted trigonometric interpolation. Write \([t^j]\) for the Taylor
coefficient, including \(1/j!\). For the rooted displacement \(Y_b(0)\),
set \(D_b(z)=\mathrm{tr}(\mathrm{tidal}(Y_b(0))(\exp z))\), with the
exponential applied linkwise. Let \(\Lambda\) be the set of lattice
sites and let \(\mathrm{i}=\sqrt{-1}\). Then the payoff is exactly
\[\mathrm{midPay}_q=\mathrm{Re}\sum_{d\ne i}\left(\mathrm{i}\sum_b[u^2v]M_{iq}^{D_b}(u e_d,-u e_d+v e_b)-\frac12[t^2]M_{iq}^{D_i}(t e_d,-t e_d)\right).\]
The cubic momentum coefficient is extracted from the quadratic field
derivative. The one-loop selector is
\[y_{di}(H)=\frac{M_{iq}^H(h e_d,-h e_d)-2M_{iq}^H(0,0)+M_{iq}^H(-h e_d,h e_d)}{2h^2\,\mathrm{card}\,\Lambda},\qquad h=2\pi/f.\]

\emph{Lean:}
\href{https://github.com/yang-mills-group/OpenAI-YM-20261028/blob/3456e5e75cf8d4cd558955a5c7638fba2d68f464/YM/TraceExtension/Det/Port.lean\#L69}{\texttt{lush}},
\href{https://github.com/yang-mills-group/OpenAI-YM-20261028/blob/3456e5e75cf8d4cd558955a5c7638fba2d68f464/YM/TraceExtension/Lamp/Cake.lean\#L63}{\texttt{molt}},
\href{https://github.com/yang-mills-group/OpenAI-YM-20261028/blob/3456e5e75cf8d4cd558955a5c7638fba2d68f464/YM/TraceExtension/Lamp/Keg.lean\#L11}{\texttt{neap}},
\href{https://github.com/yang-mills-group/OpenAI-YM-20261028/blob/3456e5e75cf8d4cd558955a5c7638fba2d68f464/YM/TraceExtension/Lamp/Lock.lean\#L20}{\texttt{pia}},
\href{https://github.com/yang-mills-group/OpenAI-YM-20261028/blob/3456e5e75cf8d4cd558955a5c7638fba2d68f464/YM/TraceExtension/Lamp/Taps.lean\#L31-L34}{\texttt{Deb}},
\href{https://github.com/yang-mills-group/OpenAI-YM-20261028/blob/3456e5e75cf8d4cd558955a5c7638fba2d68f464/YM/TraceExtension/Lamp/Lock.lean\#L21}{\texttt{cob}},
\href{https://github.com/yang-mills-group/OpenAI-YM-20261028/blob/3456e5e75cf8d4cd558955a5c7638fba2d68f464/YM/TraceExtension/Loft/Tips.lean\#L12}{\texttt{yard}}.

For the Fock tests, \(h=(h_c^a)_{1\le c\le m,\,1\le a\le3}\) is a list
of real scalar spacetime tests. Its one-particle wave is
\[\eta_c^a(k)=|k|\sum_b\left(\delta_{ab}-\frac{k_a k_b}{|k|^2}\right)\widehat h_c^b(|k|,k),\qquad d\nu(k)=\frac{d^3k}{2|k|}.\]
The coherent vector has \(n\)-particle component
\(e^{-\|\eta\|^2/2}\mathrm{i}^n\eta^{\otimes n}/\sqrt{n!}\). The value
at \(k=0\) is irrelevant to \(\nu\).

The run predicates used in VII.2.3 are \texttt{Bolt}, \texttt{Bandage},
\texttt{Nest}, \texttt{Vigil}, \texttt{Acclaim} and \texttt{supportQ};
those used in VII.2.4 are \texttt{PostT}, \texttt{Stock.Tabs},
\texttt{Scrim}, \texttt{Sceptre}, \texttt{supportQ}, \texttt{Moonlit}
and \texttt{auraQ}, together with \texttt{Binding}, \texttt{HaarBudget},
\texttt{ChartRatio}, the generator bounds and \texttt{FiresSky}. The
one-loop hypotheses in VII.2.f--h are \texttt{Cred} and \texttt{Flare}.
Their parameters and dependencies are those of the fixed run and
operator data.

\subsubsection{7. What remains
unexplained}\label{what-remains-unexplained-26}

\begin{enumerate}
\def\labelenumi{\arabic{enumi}.}
\tightlist
\item
  \textbf{Transport and one-loop estimates.} The statements of the
  fourth-order first-variation transfer and the quantitative contraction
  through successive RG steps were checked, but their proofs were not
  read in full. The same limitation applies to the covariance and
  resolvent estimates feeding the one-loop pole. The proof sketches use
  these as inputs from III.4--III.11.
\item
  \textbf{Remainder estimates.} The stage, cluster and chamber-pullback
  estimates listed after VII.2.k were checked as statements, including
  their scale, degree, support and separation restrictions. Their
  detailed proofs, and the construction of the fast bounds behind the
  trace estimate, remain to be expanded. The slowly decaying coefficient
  errors and the remainders smaller than every inverse power are
  distinct estimates.
\item
  \textbf{Wick expansion.} The finite-profile reduction and the
  incompatibility with the trace law were checked. The all-orders
  expansion of a single massive mode is used as a theorem; its proof was
  not read.
\item
  \textbf{Distributional inputs to the Fock exclusion.} The reductions
  through remote tests, equal-time data, momentum-doubling identities,
  eigenvector splitting and dyadic growth were read. The finite-order
  theorem for point-supported distributions and the analytic uniqueness
  results used in those reductions remain inputs; their full proofs were
  not audited here.
\item
  \textbf{Position-space frames.} The scalar identification, frame
  quantifiers, shell normalization and the \(O(g_X^4)\) deduction were
  checked. The full existence proof for the position-space Gram and
  angular alignment, and the detailed formulas for the two shell
  observables, remain to be expanded from the frame definitions.
  Proposition VII.2.9 concerns the family and extraction refinement
  supplied by the supplementary construction.
\item
  \textbf{Unit boundaries.} The declaration table places the spectral
  measure, the final short-distance steps and the Fock conclusion in
  VIII.1, the scaling-sequence existence results in VII.6, and the
  logarithmic momentum integral in III.7. Their mathematical roles are
  described here. The remaining integration and RG bookkeeping records
  are outside the proof sketches in this unit.
\end{enumerate}

\subsection{VII.3. The theory is not free: Gram positivity and exclusion
of free-field
profiles}\label{vii.3.-the-theory-is-not-free-gram-positivity-and-exclusion-of-free-field-profiles}

\textbf{Part:} VII, identification as Yang--Mills. \textbf{Lean:}
\href{https://github.com/yang-mills-group/OpenAI-YM-20261028/tree/3456e5e75cf8d4cd558955a5c7638fba2d68f464/YM/Extra/FREE_DETECTOR_ALGEBRA}{\texttt{YM/Extra/FREE\_DETECTOR\_ALGEBRA}},
\href{https://github.com/yang-mills-group/OpenAI-YM-20261028/tree/3456e5e75cf8d4cd558955a5c7638fba2d68f464/YM/Pert/Position}{\texttt{YM/Pert/Position}},
\href{https://github.com/yang-mills-group/OpenAI-YM-20261028/tree/3456e5e75cf8d4cd558955a5c7638fba2d68f464/YM/Pert/Rigid}{\texttt{YM/Pert/Rigid}},
\href{https://github.com/yang-mills-group/OpenAI-YM-20261028/tree/3456e5e75cf8d4cd558955a5c7638fba2d68f464/YM/Pert/Lattice}{\texttt{YM/Pert/Lattice}},
and a few declarations of
\href{https://github.com/yang-mills-group/OpenAI-YM-20261028/tree/3456e5e75cf8d4cd558955a5c7638fba2d68f464/YM/Push/LOOP_COEFFICIENT}{\texttt{YM/Push/LOOP\_COEFFICIENT}}.
Proposition VII.3.7 is in
\href{https://github.com/yang-mills-group/OpenAI-YM-20261028/tree/3456e5e75cf8d4cd558955a5c7638fba2d68f464/YM/TraceExtension/Wick}{\texttt{YM/TraceExtension/Wick}}
and
\href{https://github.com/yang-mills-group/OpenAI-YM-20261028/blob/3456e5e75cf8d4cd558955a5c7638fba2d68f464/TraceBridge.lean}{\texttt{TraceBridge.lean}}.
\textbf{Size:} about 3,800 declaration lines on the main dependency path
in the listed folders, excluding the trace-exclusion files.

The reflected Gram form of the free Maxwell Wick algebra is strictly
positive. It identifies the free field spaces with on-shell polynomial
labels and supplies injective two-point tests (VII.3.1--VII.3.5). At
degree four, rigidity gives finite position equations with invertible
linearization and a leading comparison with lattice perturbation theory
(VII.3.6). Independently, the trace anomaly excludes every finite Wick
profile (VII.3.7). The lattice series used in the position comparison
are defined in §4.7.

\subsubsection{1. What this unit proves}\label{what-this-unit-proves-27}

\textbf{Setting.}

Write
\(\mathrm{Conf}_n=\{(x_1,\dots,x_n)\in(\mathbb R^4)^n:x_i\ne x_j\text{ for }i\ne j\}\).

\emph{Coefficients and jets.}

\begin{itemize}
\tightlist
\item
  \(\mathfrak g\) is the Lie algebra of the gauge group, with its
  invariant inner product. \(\mathfrak g=0\) is allowed in
  VII.3.1--VII.3.5.
\item
  For \(j\ge0\), \(\mathrm{Sol}_j\) is the real space of
  \(\mathfrak g\)-valued 2-forms \(S=(S_{\mu\nu})\) on \(\mathbb R^4\)
  that satisfy:

  \begin{itemize}
  \tightlist
  \item
    the coefficients are homogeneous polynomials of degree \(j\);
  \item
    \(dS=0\) and \(\delta S=0\), where \(\delta\) is the codifferential
    (\href{https://github.com/yang-mills-group/OpenAI-YM-20261028/blob/3456e5e75cf8d4cd558955a5c7638fba2d68f464/YM/Extra/FREE_DETECTOR_ALGEBRA/Solutions.lean\#L115}{\texttt{Sol}},
    \href{https://github.com/yang-mills-group/OpenAI-YM-20261028/blob/3456e5e75cf8d4cd558955a5c7638fba2d68f464/YM/Maxwell/Shell.lean\#L18-L42}{\texttt{onShell}}).
  \end{itemize}
\item
  Such \(S\) are harmonic
  (\href{https://github.com/yang-mills-group/OpenAI-YM-20261028/blob/3456e5e75cf8d4cd558955a5c7638fba2d68f464/YM/Extra/FREE_DETECTOR_ALGEBRA/Solutions.lean\#L129}{\texttt{Sol.harmonic}}).
\item
  \(\mathrm{Sol}_j\) is the space of on-shell \(j\)-jets: the \(j\)-th
  Taylor coefficient of a solution of the free Maxwell equations lies in
  \(\mathrm{Sol}_j\).
\item
  \(S(x)\) denotes the polynomial 2-form.
\item
  Write \(T_jF=(j!)^{-1}D^jF\) for a Taylor coefficient. The inner
  product \(\langle S,T\rangle\) sums the color inner products over both
  ordered form indices and all ordered derivative indices; there is no
  factor \(1/2\) for the antisymmetric form indices.
\item
  \(\mathrm{Tow}_n:=\bigoplus_{j<n}\mathrm{Sol}_j\)
  (\href{https://github.com/yang-mills-group/OpenAI-YM-20261028/blob/3456e5e75cf8d4cd558955a5c7638fba2d68f464/YM/Extra/FREE_DETECTOR_ALGEBRA/PositiveLeg.lean\#L69}{\texttt{ProbeTower}}).
\end{itemize}

\emph{Wick polynomials.}

\begin{itemize}
\tightlist
\item
  Fix a basis of each \(\mathrm{Sol}_j\). Put
  \(\mathcal W=\mathbb C[X_s]_{s\in\mathsf S}\) with
  \(\mathsf S=\{(j,a)\}\)
  (\href{https://github.com/yang-mills-group/OpenAI-YM-20261028/blob/3456e5e75cf8d4cd558955a5c7638fba2d68f464/YM/Maxwell/Equation.lean\#L24}{\texttt{WickPol}},
  \href{https://github.com/yang-mills-group/OpenAI-YM-20261028/blob/3456e5e75cf8d4cd558955a5c7638fba2d68f464/YM/Maxwell/Equation.lean\#L22}{\texttt{Slot}}).
\item
  \(X_{(j,a)}\) is the \(a\)-th coordinate of the \(j\)-th on-shell
  Taylor coefficient of the free curvature at a point.
\item
  For \(p\in\mathrm{Tow}_n\), \(\ell_p\in\mathcal W\) is the real linear
  polynomial \(A\mapsto\sum_{j<n}\langle p_j,A_j\rangle\)
  (\href{https://github.com/yang-mills-group/OpenAI-YM-20261028/blob/3456e5e75cf8d4cd558955a5c7638fba2d68f464/YM/Extra/FREE_DETECTOR_ALGEBRA/Symbol.lean\#L77}{\texttt{digit}},
  \href{https://github.com/yang-mills-group/OpenAI-YM-20261028/blob/3456e5e75cf8d4cd558955a5c7638fba2d68f464/YM/Extra/FREE_DETECTOR_ALGEBRA/Script.lean\#L16-L23}{\texttt{code}}).
\end{itemize}

\emph{Free moments.}

\begin{itemize}
\tightlist
\item
  \(C_{st}(x-y)\) is the free covariance of jet coordinate \(s\) at
  \(x\) with jet coordinate \(t\) at \(y\)
  (\href{https://github.com/yang-mills-group/OpenAI-YM-20261028/blob/3456e5e75cf8d4cd558955a5c7638fba2d68f464/YM/Maxwell/Wick.lean\#L24}{\texttt{Wick.entry}},
  \href{https://github.com/yang-mills-group/OpenAI-YM-20261028/blob/3456e5e75cf8d4cd558955a5c7638fba2d68f464/YM/Maxwell/Leg.lean\#L161-L162}{\texttt{legKernel}}).
  It is smooth off the diagonal and built from derivatives of
  \(\Gamma(z)=(4\pi^2|z|^2)^{-1}\).
\item
  For \(x\in\mathrm{Conf}_n\), the free moment
  \(\langle P_1(x_1)\cdots P_n(x_n)\rangle_0\) is multilinear
  (\href{https://github.com/yang-mills-group/OpenAI-YM-20261028/blob/3456e5e75cf8d4cd558955a5c7638fba2d68f464/YM/Maxwell/Wick.lean\#L126-L127}{\texttt{Wick.moment}}).
  On monomials it is the sum, over perfect matchings of the variable
  occurrences with no pair inside one factor, of
  \(\prod_{\rm pairs}C_{st}(x_i-x_k)\).
\item
  These are the Wick-ordered correlations of \(\dim\mathfrak g\)
  independent free Maxwell curvatures. Multiplication in \(\mathcal W\)
  therefore represents Wick multiplication at one insertion.
\end{itemize}

\emph{Reflected pairing.}

\begin{itemize}
\tightlist
\item
  For \(z\ne0\), \(\mathcal R_zv=v-2\langle v,z\rangle|z|^{-2}z\) is the
  reflection in \(z^\perp\)
  (\href{https://github.com/yang-mills-group/OpenAI-YM-20261028/blob/3456e5e75cf8d4cd558955a5c7638fba2d68f464/YM/Extra/FREE_DETECTOR_ALGEBRA/Forms.lean\#L103}{\texttt{face}}).
  Orthogonal transformations act on all tensor indices of jets and on
  polynomials by \((R\cdot P)(A)=P(R^{-1}\cdot A)\)
  (\href{https://github.com/yang-mills-group/OpenAI-YM-20261028/blob/3456e5e75cf8d4cd558955a5c7638fba2d68f464/YM/Extra/FREE_DETECTOR_ALGEBRA/Symbol.lean\#L54-L55}{\texttt{pullSol}}).
\item
  \(\bar P\) conjugates the coefficients of \(P\)
  (\href{https://github.com/yang-mills-group/OpenAI-YM-20261028/blob/3456e5e75cf8d4cd558955a5c7638fba2d68f464/YM/Extra/FREE_DETECTOR_ALGEBRA/Symbol.lean\#L46}{\texttt{barPol}}).
\item
  For \(x\ne y\)
  (\href{https://github.com/yang-mills-group/OpenAI-YM-20261028/blob/3456e5e75cf8d4cd558955a5c7638fba2d68f464/YM/Extra/FREE_DETECTOR_ALGEBRA/Duplex.lean\#L54}{\texttt{duplex}},
  \href{https://github.com/yang-mills-group/OpenAI-YM-20261028/blob/3456e5e75cf8d4cd558955a5c7638fba2d68f464/YM/Extra/FREE_DETECTOR_ALGEBRA/Duplex.lean\#L49-L52}{\texttt{aligned}}):
\end{itemize}

\[\langle P,Q\rangle_{x,y}:=\big\langle P(x)\,(\mathcal R_{x-y}\bar Q)(y)\big\rangle_0 .\]

\begin{itemize}
\tightlist
\item
  This is the local two-point Osterwalder--Schrader pairing for
  reflection in the perpendicular bisector of \([x,y]\), with the
  convention that the first argument is linear and the second
  conjugate-linear.
\end{itemize}

\emph{Tower Gram form.}

\begin{itemize}
\tightlist
\item
  On \(\mathrm{Tow}_n\) the restriction to linear polynomials is
  \(B_{x-y}(p,q):=\langle\ell_p,\ell_q\rangle_{x,y}\). It is real and
  depends only on \(z=x-y\).
\item
  Explicitly, \(B_z(p,q)=\sum_{j,k<n}\mathrm{scan}_z(p_j,q_k)\), where
  \(\mathrm{scan}_z(P,Q)\) is the free covariance of
  \(\langle P,T_jF\rangle(z)\) with
  \(\langle\mathcal R_zQ,T_kF\rangle(0)\)
  (\href{https://github.com/yang-mills-group/OpenAI-YM-20261028/blob/3456e5e75cf8d4cd558955a5c7638fba2d68f464/YM/Extra/FREE_DETECTOR_ALGEBRA/PositiveLeg.lean\#L71-L77}{\texttt{towerB}},
  \href{https://github.com/yang-mills-group/OpenAI-YM-20261028/blob/3456e5e75cf8d4cd558955a5c7638fba2d68f464/YM/Extra/FREE_DETECTOR_ALGEBRA/PositiveLeg.lean\#L29-L32}{\texttt{scan}}).
\item
  The identification of the two descriptions is
  \href{https://github.com/yang-mills-group/OpenAI-YM-20261028/blob/3456e5e75cf8d4cd558955a5c7638fba2d68f464/YM/Extra/FREE_DETECTOR_ALGEBRA/Script.lean\#L83-L94}{\texttt{crash\_test}}
  together with
  \href{https://github.com/yang-mills-group/OpenAI-YM-20261028/blob/3456e5e75cf8d4cd558955a5c7638fba2d68f464/YM/Extra/FREE_DETECTOR_ALGEBRA/Duplex.lean\#L60-L66}{\texttt{join\_crash}}.
\end{itemize}

\textbf{Proposition VII.3.1 (strict reflected positivity of the free
Wick algebra).} Let \(x\ne y\) in \(\mathbb R^4\). (a) For every
\(P\in\mathcal W\): \(\langle P,P\rangle_{x,y}>0\) if and only if
\(P\ne0\). (b)
\(\overline{\langle P,Q\rangle_{x,y}}=\langle Q,P\rangle_{x,y}\) for all
\(P,Q\in\mathcal W\). (c) For every finite \(\mathbb C\)-linearly
independent family \((V_i)\) in \(\mathcal W\), the matrix
\((\langle V_j,V_i\rangle_{x,y})_{ij}\) is positive definite. No
simplicity or positive-dimension assumption is needed, and the
polynomial degree is unrestricted. Self-pairings are real, so the
inequality in (a) is a real inequality. \emph{Status:} theorem. Only (a)
lies on the path to the main theorem; (b) and (c) are proved but not
used. \emph{Lean:} (a)
\href{https://github.com/yang-mills-group/OpenAI-YM-20261028/blob/3456e5e75cf8d4cd558955a5c7638fba2d68f464/YM/Extra/FREE_DETECTOR_ALGEBRA/WickPositive.lean\#L47-L53}{\texttt{reflected\_positive}};
(b)
\href{https://github.com/yang-mills-group/OpenAI-YM-20261028/blob/3456e5e75cf8d4cd558955a5c7638fba2d68f464/YM/Extra/FREE_DETECTOR_ALGEBRA/WickPositive.lean\#L57-L76}{\texttt{reflected\_symm}};
(c)
\href{https://github.com/yang-mills-group/OpenAI-YM-20261028/blob/3456e5e75cf8d4cd558955a5c7638fba2d68f464/YM/Extra/FREE_DETECTOR_ALGEBRA/WickPositive.lean\#L80-L97}{\texttt{reflected\_pd}}.

\textbf{Proposition VII.3.2 (tower Gram form and Fock structure).} Let
\(x\ne y\), \(z=x-y\), \(n\in\mathbb N\). (a) \(B_z\) is symmetric and
positive definite on \(\mathrm{Tow}_n\). More precisely, there is a
block-unitriangular automorphism \(\sigma_z\) of \(\mathrm{Tow}_n\)
(Lemma VII.3.f) such that

\[B_z(\sigma_zp,\sigma_zq)=\sum_{j<n}\pi_j(z)\,\langle p_j,q_j\rangle,\qquad \pi_j(z)=\frac{(j+1)(j+2)\,2^j}{\pi^2\,|z|^{2j+4}}>0 .\]

\begin{enumerate}
\def\labelenumi{(\alph{enumi})}
\setcounter{enumi}{1}
\tightlist
\item
  Let \(b_1,\dots,b_N\in\mathrm{Tow}_n\) satisfy
  \(B_z(b_i,b_k)=d_i\delta_{ik}\), with any \(d_i\in\mathbb R\). Let
  \(\Phi_b:\mathbb C[Y_1,\dots,Y_N]\to\mathcal W\) be the algebra
  homomorphism with \(Y_i\mapsto\ell_{b_i}\)
  (\href{https://github.com/yang-mills-group/OpenAI-YM-20261028/blob/3456e5e75cf8d4cd558955a5c7638fba2d68f464/YM/Extra/FREE_DETECTOR_ALGEBRA/Script.lean\#L25}{\texttt{encode}}).
  Then
\end{enumerate}

\[\langle\Phi_bp,\Phi_bq\rangle_{x,y}=\sum_\alpha p_\alpha\,\overline{q_\alpha}\;\alpha!\,d^\alpha,\qquad \alpha!\,d^\alpha:=\prod_i\alpha_i!\,d_i^{\alpha_i}.\]

\emph{Status:} theorem. \emph{Lean:} (a)
\href{https://github.com/yang-mills-group/OpenAI-YM-20261028/blob/3456e5e75cf8d4cd558955a5c7638fba2d68f464/YM/Extra/FREE_DETECTOR_ALGEBRA/PositiveLeg.lean\#L90-L94}{\texttt{towerB\_sym}},
\href{https://github.com/yang-mills-group/OpenAI-YM-20261028/blob/3456e5e75cf8d4cd558955a5c7638fba2d68f464/YM/Extra/FREE_DETECTOR_ALGEBRA/PositiveLeg.lean\#L150-L154}{\texttt{towerB\_pos}},
\href{https://github.com/yang-mills-group/OpenAI-YM-20261028/blob/3456e5e75cf8d4cd558955a5c7638fba2d68f464/YM/Extra/FREE_DETECTOR_ALGEBRA/PositiveLeg.lean\#L126-L133}{\texttt{shiftAll\_gram}},
\href{https://github.com/yang-mills-group/OpenAI-YM-20261028/blob/3456e5e75cf8d4cd558955a5c7638fba2d68f464/YM/Extra/FREE_DETECTOR_ALGEBRA/PositiveLeg.lean\#L24}{\texttt{price}},
\href{https://github.com/yang-mills-group/OpenAI-YM-20261028/blob/3456e5e75cf8d4cd558955a5c7638fba2d68f464/YM/Extra/FREE_DETECTOR_ALGEBRA/PositiveLeg.lean\#L25}{\texttt{price\_pos}};
(b)
\href{https://github.com/yang-mills-group/OpenAI-YM-20261028/blob/3456e5e75cf8d4cd558955a5c7638fba2d68f464/YM/Extra/FREE_DETECTOR_ALGEBRA/PowerGram.lean\#L100-L116}{\texttt{duplex\_script}},
\href{https://github.com/yang-mills-group/OpenAI-YM-20261028/blob/3456e5e75cf8d4cd558955a5c7638fba2d68f464/YM/Extra/FREE_DETECTOR_ALGEBRA/PowerGram.lean\#L60-L85}{\texttt{aligned\_script}}.

\textbf{Proposition VII.3.3 (Wick recursion; characterization of the
free moments).} Let \(n\ge1\), \(x\in\mathrm{Conf}_n\),
\(W\in\mathcal W^n\), \(1\le i\le n\) and \(s\in\mathsf S\). Let
\(\partial_s^{y,y'}\) be the \(\mathbb C\)-linear derivation of
\(\mathcal W\) with \(\partial_s^{y,y'}X_t=C_{st}(y-y')\). Then

\[\big\langle W_1(x_1)\cdots(W_iX_s)(x_i)\cdots W_n(x_n)\big\rangle_0=\sum_{k\ne i}\big\langle W_1(x_1)\cdots(\partial_s^{x_i,x_k}W_k)(x_k)\cdots W_n(x_n)\big\rangle_0 ,\]

where on the right the factor at \(x_i\) is \(W_i\). Conversely, for any
\(n\ge0\) and \(x\in\mathrm{Conf}_n\), a \(\mathbb C\)-multilinear map
\(T:\mathcal W^n\to\mathbb C\) with \(T(1,\dots,1)=1\) satisfying this
recursion for every \(i,s,W\) is the free moment at \(x\). Consequences:

\begin{itemize}
\tightlist
\item
  \(\langle\cdots\rangle_0\) is invariant under any simultaneous
  \(O(4)\) rotation of points and polynomials.
\item
  \(\langle\bar P_1(x_1)\cdots\bar P_n(x_n)\rangle_0=\overline{\langle P_1(x_1)\cdots P_n(x_n)\rangle_0}\).
\end{itemize}

\emph{Status:} theorem. \emph{Lean:}
\href{https://github.com/yang-mills-group/OpenAI-YM-20261028/blob/3456e5e75cf8d4cd558955a5c7638fba2d68f464/YM/Extra/FREE_DETECTOR_ALGEBRA/Removal.lean\#L91-L108}{\texttt{Pair.remove}},
\href{https://github.com/yang-mills-group/OpenAI-YM-20261028/blob/3456e5e75cf8d4cd558955a5c7638fba2d68f464/YM/Extra/FREE_DETECTOR_ALGEBRA/Removal.lean\#L25-L26}{\texttt{Pair.extractVar}};
characterization
\href{https://github.com/yang-mills-group/OpenAI-YM-20261028/blob/3456e5e75cf8d4cd558955a5c7638fba2d68f464/YM/Extra/FREE_DETECTOR_ALGEBRA/Gaussian.lean\#L24-L62}{\texttt{byAnnihilators}};
consequences
\href{https://github.com/yang-mills-group/OpenAI-YM-20261028/blob/3456e5e75cf8d4cd558955a5c7638fba2d68f464/YM/Extra/FREE_DETECTOR_ALGEBRA/Gaussian.lean\#L81-L101}{\texttt{wick\_rotate}},
\href{https://github.com/yang-mills-group/OpenAI-YM-20261028/blob/3456e5e75cf8d4cd558955a5c7638fba2d68f464/YM/Extra/FREE_DETECTOR_ALGEBRA/Gaussian.lean\#L111-L137}{\texttt{wick\_bar}}.

\textbf{Proposition VII.3.4 (the free table is faithful on on-shell
labels).}

Setting:

\begin{itemize}
\tightlist
\item
  \(\mathrm{Label}_d\) is the space of complex polynomials in radial
  potential jets that are invariant under the adjoint action of the
  gauge group and have engineering weight at most \(d\); an order-\(j\)
  slot has weight \(j+2\) (II.1, III.16;
  \href{https://github.com/yang-mills-group/OpenAI-YM-20261028/blob/3456e5e75cf8d4cd558955a5c7638fba2d68f464/YM/Labels/Polynomial.lean\#L112}{\texttt{Label}}).
\item
  \(\omega_d:\mathrm{Label}_d\to\mathcal W\) evaluates a label on free
  on-shell jets
  (\href{https://github.com/yang-mills-group/OpenAI-YM-20261028/blob/3456e5e75cf8d4cd558955a5c7638fba2d68f464/YM/Maxwell/Equation.lean\#L66-L67}{\texttt{observable}}).
\item
  \(\mathrm{sh}_d:=\omega_d(\mathrm{Label}_d)\)
  (\href{https://github.com/yang-mills-group/OpenAI-YM-20261028/blob/3456e5e75cf8d4cd558955a5c7638fba2d68f464/YM/Euclid/Comparator.lean\#L64}{\texttt{shellSector}}).
\item
  The free table
  (\href{https://github.com/yang-mills-group/OpenAI-YM-20261028/blob/3456e5e75cf8d4cd558955a5c7638fba2d68f464/YM/Euclid/Comparator.lean\#L23-L24}{\texttt{freeTable}},
  IV.1) has correlations
  \(\langle\omega(a_1)(x_1)\cdots\omega(a_n)(x_n)\rangle_0\).
\item
  Its field space \(E_d\) (the glossary's \(F_d\), for the free table)
  is the space of degree-\(d\) reference labels modulo null labels.
\end{itemize}

Statements: (a) For every \(d\ge0\), let \(P\in\mathrm{sh}_d\). Suppose
\(\langle P(x_0)Q_1(x_1)\cdots Q_n(x_n)\rangle_0=0\) for every
\(n\ge0\), every choice of degrees \(d_1,\dots,d_n\ge0\), all
\(Q_i\in\mathrm{sh}_{d_i}\) and all \(x\in\mathrm{Conf}_{n+1}\). Then
\(P=0\). Hence:

\begin{itemize}
\tightlist
\item
  the null labels of the free table at degree \(d\) are exactly the
  labels with \(\omega_d(a)=0\);
\item
  \(\omega_d\) induces a linear isomorphism \(E_d\cong\mathrm{sh}_d\).
  Write \(\varpi_d:E_d\hookrightarrow\mathcal W\) for the resulting
  injection.
\end{itemize}

\begin{enumerate}
\def\labelenumi{(\alph{enumi})}
\setcounter{enumi}{1}
\item
  For every basis \((b_i)_{i\in I}\) of \(E_d\) and every \(x\ne y\),
  the map
  \(v\mapsto\big(\langle\varpi v(x)\,\varpi b_i(y)\rangle_0\big)_{i\in I}\)
  is injective.
\item
  Fix a Wilson family \(C\), any table \(S\) on its register, and a pair
  satisfying \(x_\mu\ne y_\mu\) for each coordinate \(\mu\), for example
  \((x,y)=(0,(1,1,1,1))\). For every \(d\ge0\) there is a column system
  (IV.1) for \(E_d\). It consists of:
\end{enumerate}

\begin{itemize}
\tightlist
\item
  a complex basis fixed by conjugation, with the unit first, whose
  elements are images of frame labels with recorded weight and
  polynomial length;
\item
  the pair readings of (b) against this basis at \((x,y)\), which are
  injective.
\end{itemize}

\emph{Status:} theorem. Part (a) enters the rank comparisons of III.11
and III.12; (c) supplies columns for III.11 and III.2. The embedding
\(\varpi_d\) is used in VII.1. \emph{Lean:} (a)
\href{https://github.com/yang-mills-group/OpenAI-YM-20261028/blob/3456e5e75cf8d4cd558955a5c7638fba2d68f464/YM/Extra/FREE_DETECTOR_ALGEBRA/Rank.lean\#L38-L54}{\texttt{no\_residual}}
(with the definitions
\href{https://github.com/yang-mills-group/OpenAI-YM-20261028/blob/3456e5e75cf8d4cd558955a5c7638fba2d68f464/YM/Euclid/Comparator.lean\#L84-L85}{\texttt{residual}},
\href{https://github.com/yang-mills-group/OpenAI-YM-20261028/blob/3456e5e75cf8d4cd558955a5c7638fba2d68f464/YM/Euclid/Comparator.lean\#L67-L72}{\texttt{ProbeFamily}}),
\href{https://github.com/yang-mills-group/OpenAI-YM-20261028/blob/3456e5e75cf8d4cd558955a5c7638fba2d68f464/YM/Extra/FREE_DETECTOR_ALGEBRA/Rank.lean\#L58-L59}{\texttt{free\_null}},
\href{https://github.com/yang-mills-group/OpenAI-YM-20261028/blob/3456e5e75cf8d4cd558955a5c7638fba2d68f464/YM/Extra/FREE_DETECTOR_ALGEBRA/Rank.lean\#L64-L65}{\texttt{exactField}},
\href{https://github.com/yang-mills-group/OpenAI-YM-20261028/blob/3456e5e75cf8d4cd558955a5c7638fba2d68f464/YM/Extra/FREE_DETECTOR_ALGEBRA/Rank.lean\#L70-L71}{\texttt{writes}};
(b)
\href{https://github.com/yang-mills-group/OpenAI-YM-20261028/blob/3456e5e75cf8d4cd558955a5c7638fba2d68f464/YM/Extra/FREE_DETECTOR_ALGEBRA/Columns.lean\#L103-L115}{\texttt{scanner\_inj}};
(c)
\href{https://github.com/yang-mills-group/OpenAI-YM-20261028/blob/3456e5e75cf8d4cd558955a5c7638fba2d68f464/YM/Extra/FREE_DETECTOR_ALGEBRA/Vane.lean\#L25-L38}{\texttt{vane}},
\href{https://github.com/yang-mills-group/OpenAI-YM-20261028/blob/3456e5e75cf8d4cd558955a5c7638fba2d68f464/YM/Extra/FREE_DETECTOR_ALGEBRA/Vane.lean\#L41-L42}{\texttt{vane\_possible}},
\href{https://github.com/yang-mills-group/OpenAI-YM-20261028/blob/3456e5e75cf8d4cd558955a5c7638fba2d68f464/YM/Extra/FREE_DETECTOR_ALGEBRA/GradedBasis.lean\#L18-L27}{\texttt{RStage}},
\href{https://github.com/yang-mills-group/OpenAI-YM-20261028/blob/3456e5e75cf8d4cd558955a5c7638fba2d68f464/YM/Extra/FREE_DETECTOR_ALGEBRA/GradedBasis.lean\#L41-L49}{\texttt{recorded}}.

\textbf{Proposition VII.3.5 (covariance of the free table).} For every
gauge group in the setting above, the free table has the following
symmetries:

\begin{itemize}
\tightlist
\item
  its correlations are invariant under the signed coordinate
  permutations, acting on points and labels;
\item
  conjugating the labels conjugates the correlations.
\end{itemize}

In addition, \(O(4)\) acts on each \(E_d\) by complex-linear
isomorphisms satisfying \(\varpi_d(R\cdot v)=R\cdot(\varpi_dv)\).
\emph{Status:} theorem. The proof uses the \(O(4)\)-covariance of the
free kernel, not any reconstruction theorem. \emph{Lean:}
\href{https://github.com/yang-mills-group/OpenAI-YM-20261028/blob/3456e5e75cf8d4cd558955a5c7638fba2d68f464/YM/Extra/FREE_DETECTOR_ALGEBRA/LabelTurn.lean\#L80-L91}{\texttt{table\_spin}},
\href{https://github.com/yang-mills-group/OpenAI-YM-20261028/blob/3456e5e75cf8d4cd558955a5c7638fba2d68f464/YM/Extra/FREE_DETECTOR_ALGEBRA/Rank.lean\#L106-L110}{\texttt{rideO}},
\href{https://github.com/yang-mills-group/OpenAI-YM-20261028/blob/3456e5e75cf8d4cd558955a5c7638fba2d68f464/YM/Extra/FREE_DETECTOR_ALGEBRA/Rank.lean\#L99-L101}{\texttt{ride}}.

\textbf{Proposition VII.3.6 (rigidity at degree four; position
equations).} Parts (a) and (c) assume that \(\mathfrak g\) is simple and
nonabelian. Part (b) holds for any gauge group in the setting above.

\emph{Buffers.} For a degree \(d\), a \emph{buffer}
(\href{https://github.com/yang-mills-group/OpenAI-YM-20261028/blob/3456e5e75cf8d4cd558955a5c7638fba2d68f464/YM/Pert/Position/Board.lean\#L31-L40}{\texttt{buffers}})
is a \(\mathbb C\)-linear map \(B:E_d\to E_d\) such that:

\begin{itemize}
\tightlist
\item
  \(B\) kills the unit;
\item
  for each \(e<d\), \(B\) preserves the subspace of fields of degree
  \(\le e\);
\item
  \(B\) commutes with conjugation and with the signed coordinate
  permutations.
\end{itemize}

Buffers form a finite-dimensional real space \(\mathrm{Buf}_d\).

\emph{Stationarity.} \(B\) is \emph{stationary} if

\[\sum_{j=1}^r\langle v_1(x_1)\cdots(Bv_j)(x_j)\cdots v_r(x_r)\rangle_0=0\quad\text{for all } r,\ v\in E_d^r,\ x\in\mathrm{Conf}_r .\]

Equivalently, the change of identification \(v\mapsto(1+tB)v\) of
degree-\(d\) fields leaves all free correlations unchanged to first
order in \(t\).

\begin{enumerate}
\def\labelenumi{(\alph{enumi})}
\tightlist
\item
  \emph{Rigidity.} Every stationary buffer of degree \(4\) vanishes. In
  fact \(B\in\mathrm{Buf}_4\) vanishes as soon as both of the following
  hold, with \(e\) the unit vector of coordinate axis \(0\):
\end{enumerate}

\begin{itemize}
\tightlist
\item
  \(\langle Bv\rangle_0=0\) for all \(v\);
\item
  \(\langle(Bv)(e)\,w(0)\rangle_0+\langle v(e)\,(Bw)(0)\rangle_0=0\) for
  all \(v,w\).
\end{itemize}

For \(d<4\), \(\mathrm{Buf}_d=0\).

\begin{enumerate}
\def\labelenumi{(\alph{enumi})}
\setcounter{enumi}{1}
\tightlist
\item
  \emph{Square systems.} Suppose every stationary buffer of degree \(d\)
  vanishes. Then there are finitely many \emph{sections}
  \((r_i,p_i,f_i)_{1\le i\le k}\), with \(r_i\ge0\),
  \(p_i\in E_d^{r_i}\) and
  \(f_i\in C_c^\infty(\mathrm{Conf}_{r_i};\mathbb C)\), and a
  real-linear map \(g:\mathbb R^k\to\mathbb R^m\) with
  \(m=\dim_{\mathbb R}\mathrm{Buf}_d\), such that the map
\end{enumerate}

\[B\longmapsto g\Big(\mathrm{Re}\int f_i(x)\sum_j\big\langle p_{i1}(x_1)\cdots(Bp_{ij})(x_j)\cdots\big\rangle_0\,dx\Big)_{i\le k}\]

is an isomorphism \(\mathrm{Buf}_d\to\mathbb R^m\).

\begin{itemize}
\tightlist
\item
  This map is the derivative in the buffer direction, at the free data
  and \(B=0\), of the \emph{position chart}
  \((w,B)\mapsto g\big(\mathrm{Re}\,w_i((1+B)p_i)\big)_i\). Here \(w_i\)
  ranges over complex-multilinear forms on \(E_d^{r_i}\).
\item
  The system is called \emph{square}, and \emph{clean} (all profiles are
  test functions).
\item
  For simple nonabelian \(\mathfrak g\) and \(d=4\) such a system
  exists.
\end{itemize}

\begin{enumerate}
\def\labelenumi{(\alph{enumi})}
\setcounter{enumi}{2}
\tightlist
\item
  \emph{Leading position comparison at degree four.} Hypotheses:
\end{enumerate}

\begin{itemize}
\tightlist
\item
  \(b\) is a final record over the Wilson family, with limit table
  \(S\).
\item
  \(X=(r_k,U_{k,d},\xi_k)\) is a scaling sequence: \(r_k>0\) decreases
  strictly to \(0\), \(\xi_k>0\), and \(U_{k,d}:E_d\to F_d(S)\) is a
  complex-linear isomorphism for every \(k,d\). Write \(U_k=U_{k,4}\).
\item
  Along \(X\), the leading free-Maxwell comparison of VI.2 holds,
  including eventual compatibility of the frames with the unit,
  conjugation, signed permutations and degree inclusions.
\item
  \(Q\) is a clean square degree-four system. The map \(\varphi\)
  satisfies the four conditions in VII.3.r: a local analytic implicit
  solution, convergence of the order-zero lattice data, the formal
  position identities, and an actual leading comparison along \(X\).
\end{itemize}

Put \(B_k:=\varphi(m_Q(k))\), where \(m_Q(k)\) are the tested rescaled
correlations \(\int f(x)\langle U_kv_1(r_kx_1)\cdots\rangle_S\,dx\) on
basis fields. Then:

\begin{itemize}
\tightlist
\item
  \(B_k\to0\);
\item
  eventually \(1+B_k\) is invertible, and the corrected frames
  \(U_k(1+B_k)\) satisfy the position equations of \(Q\) exactly;
\item
  eventually \(U_k(1+B_k)\) commutes with conjugation and signed
  permutations, maps the free unit to the unit of \(S\), and is
  compatible with the frames of degrees \(\le3\) under the inclusions
  \(E_d\subset E_4\);
\item
  for every clean degree-four system \(R\) and every integer \(L\ge2\),
\end{itemize}

\[m_R^{B_k}(k)-\mathrm{chart}_R\bigl(\pi_R^{(0)}(L,k),\varphi(\pi_Q^{(0)}(L,k))\bigr)\longrightarrow0.\]

Here \(m_R^{B_k}(k)\) is the tested reading after replacing each field
\(v\) by \(U_k(1+B_k)v\), and \(\pi_R^{(0)}(L,k)\) is the array of
order-zero lattice readings on basis fields at spacing \(L^{-k}\).

These conclusions, together with the comparisons for all clean systems
described in VII.3.r, give clause (D15) of the detailed main theorem.

\emph{Status:}

\begin{itemize}
\tightlist
\item
  \begin{enumerate}
  \def\labelenumi{(\alph{enumi})}
  \tightlist
  \item
    and (b): theorems.
  \end{enumerate}
\item
  (c): theorem under the leading comparison hypothesis (VI.2; discharged
  by
  \href{https://github.com/yang-mills-group/OpenAI-YM-20261028/blob/3456e5e75cf8d4cd558955a5c7638fba2d68f464/YM/Pert/FirstTarget.lean\#L57-L81}{\texttt{firstTarget}})
  and the four conditions on \(\varphi\) (discharged by
  \href{https://github.com/yang-mills-group/OpenAI-YM-20261028/blob/3456e5e75cf8d4cd558955a5c7638fba2d68f464/YM/Pert/Position/Crest.lean\#L36-L43}{\texttt{calibrate}},
  VIII.1).
\item
  VIII.1 assembles this with the existence of a clean square system:
  simplicity of \(\mathfrak g\) and the leading free comparison imply
  the full degree-four position conclusion.
\end{itemize}

\emph{Lean:} (a)
\href{https://github.com/yang-mills-group/OpenAI-YM-20261028/blob/3456e5e75cf8d4cd558955a5c7638fba2d68f464/YM/Pert/Rigid/Camera.lean\#L107-L113}{\texttt{Rigid.frozen\_four}},
\href{https://github.com/yang-mills-group/OpenAI-YM-20261028/blob/3456e5e75cf8d4cd558955a5c7638fba2d68f464/YM/Pert/Rigid/Camera.lean\#L114-L117}{\texttt{Rigid.first\_frozen}}
(unused), definition
\href{https://github.com/yang-mills-group/OpenAI-YM-20261028/blob/3456e5e75cf8d4cd558955a5c7638fba2d68f464/YM/Pert/Position/Detect.lean\#L102}{\texttt{Frozen}};
(b)
\href{https://github.com/yang-mills-group/OpenAI-YM-20261028/blob/3456e5e75cf8d4cd558955a5c7638fba2d68f464/YM/Pert/Position/Detect.lean\#L131-L147}{\texttt{square\_of\_frozen}},
\href{https://github.com/yang-mills-group/OpenAI-YM-20261028/blob/3456e5e75cf8d4cd558955a5c7638fba2d68f464/YM/Pert/Position/Board.lean\#L54-L59}{\texttt{Conditions}},
\href{https://github.com/yang-mills-group/OpenAI-YM-20261028/blob/3456e5e75cf8d4cd558955a5c7638fba2d68f464/YM/Pert/Position/Board.lean\#L69}{\texttt{chart}},
\href{https://github.com/yang-mills-group/OpenAI-YM-20261028/blob/3456e5e75cf8d4cd558955a5c7638fba2d68f464/YM/Pert/Position/Board.lean\#L103-L104}{\texttt{white}},
\href{https://github.com/yang-mills-group/OpenAI-YM-20261028/blob/3456e5e75cf8d4cd558955a5c7638fba2d68f464/YM/Pert/Position/Board.lean\#L106}{\texttt{square}},
\href{https://github.com/yang-mills-group/OpenAI-YM-20261028/blob/3456e5e75cf8d4cd558955a5c7638fba2d68f464/YM/Pert/Position/Four.lean\#L13-L15}{\texttt{square\_four}};
(c)
\href{https://github.com/yang-mills-group/OpenAI-YM-20261028/blob/3456e5e75cf8d4cd558955a5c7638fba2d68f464/YM/Pert/Position/Four.lean\#L32-L43}{\texttt{residency}},
\href{https://github.com/yang-mills-group/OpenAI-YM-20261028/blob/3456e5e75cf8d4cd558955a5c7638fba2d68f464/YM/Pert/Position/Four.lean\#L23-L30}{\texttt{Resident}},
\href{https://github.com/yang-mills-group/OpenAI-YM-20261028/blob/3456e5e75cf8d4cd558955a5c7638fba2d68f464/YM/Pert/Position/Four.lean\#L45-L47}{\texttt{FourMeters}},
\href{https://github.com/yang-mills-group/OpenAI-YM-20261028/blob/3456e5e75cf8d4cd558955a5c7638fba2d68f464/YM/Pert/Position/Crest.lean\#L45-L48}{\texttt{Crest}},
\href{https://github.com/yang-mills-group/OpenAI-YM-20261028/blob/3456e5e75cf8d4cd558955a5c7638fba2d68f464/YM/Pert/Position/Crest.lean\#L26-L34}{\texttt{BoardAttains}},
\href{https://github.com/yang-mills-group/OpenAI-YM-20261028/blob/3456e5e75cf8d4cd558955a5c7638fba2d68f464/YM/Pert/Position/Four.lean\#L49-L53}{\texttt{four\_on}},
\href{https://github.com/yang-mills-group/OpenAI-YM-20261028/blob/3456e5e75cf8d4cd558955a5c7638fba2d68f464/YM/Pert/Position/Four.lean\#L56-L63}{\texttt{normalized\_four}}.

\textbf{Proposition VII.3.7 (exclusion of finite Wick profiles; assigned
to VII.2 and VIII.1).}

Definitions:

\begin{itemize}
\tightlist
\item
  \(G_\Theta(r)\) is the two-point function of the centered trace
  \(\Theta\) of the stress tensor in the selected phase, at the points
  \(\mp\frac r2e_{\rm t}\) on the time axis
  (\href{https://github.com/yang-mills-group/OpenAI-YM-20261028/blob/3456e5e75cf8d4cd558955a5c7638fba2d68f464/ChallengeSpec.lean\#L149}{\texttt{traceTwo}}).
\item
  A \emph{mode} is \((\mathsf m,j,q,h)\) with \(\mathsf m\ge0\),
  \(j,q\in\mathbb N\), and \(h:\mathbb R^3\to\mathbb R\) continuous with
  \(h(tp)=t^jh(p)\) for every \(t>0\) and \(p\in\mathbb R^3\).
\item
  For \(r>0\), its integrated kernel is
  \(K_{\mathsf m,q,h}(r)=\int_{\mathbb R^3}E^q\,h(p)\,e^{-rE}(2E)^{-1}dp\),
  with \(E=\sqrt{\mathsf m^2+|p|^2}\).
\item
  A \emph{Wick profile} is a finite sum
  \(W=\sum_\alpha c_\alpha\prod_\beta K_{\alpha\beta}\) with
  \(c_\alpha\in\mathbb C\); empty products equal \(1\)
  (\href{https://github.com/yang-mills-group/OpenAI-YM-20261028/blob/3456e5e75cf8d4cd558955a5c7638fba2d68f464/ChallengeSpec.lean\#L152-L158}{\texttt{Mode}},
  \href{https://github.com/yang-mills-group/OpenAI-YM-20261028/blob/3456e5e75cf8d4cd558955a5c7638fba2d68f464/ChallengeSpec.lean\#L159-L161}{\texttt{Mode.kernel}},
  \href{https://github.com/yang-mills-group/OpenAI-YM-20261028/blob/3456e5e75cf8d4cd558955a5c7638fba2d68f464/ChallengeSpec.lean\#L162-L164}{\texttt{Track}},
  \href{https://github.com/yang-mills-group/OpenAI-YM-20261028/blob/3456e5e75cf8d4cd558955a5c7638fba2d68f464/ChallengeSpec.lean\#L165-L166}{\texttt{wickProfile}};
  statement.md, Definition 1.15).
\end{itemize}

\begin{enumerate}
\def\labelenumi{(\alph{enumi})}
\tightlist
\item
  Suppose \(r^8\log^2(1/r)\,G_\Theta(r)\to A\) as \(r\to0^+\), with
  \(A>0\). Then for every Wick profile \(W\) and every \(\delta>0\)
  there is \(r\in(0,\delta)\) with \(G_\Theta(r)\ne W(r)\). Only this
  asymptotic law of \(G_\Theta\) is used.
\item
  Under the same asymptotic hypothesis, in the selected reconstructed
  theory, no finite sum of the corresponding mass-shell contraction
  integrals reproduces the Wightman two-point function of \(\Theta\) on
  every product \(\bar f\otimes g\) of Schwartz tests for which
  \(\psi_{f,g}=\overline{\widehat f}\,\widehat g\) is compactly
  supported, using the positive-sign Fourier convention. A product of
  mode kernels corresponds to testing \(\psi_{f,g}\) at the sum of their
  positive-energy on-shell momenta, with the product of their densities
  \(E^qh(p)/(2E)\)
  (\href{https://github.com/yang-mills-group/OpenAI-YM-20261028/blob/3456e5e75cf8d4cd558955a5c7638fba2d68f464/YM/TraceExtension/Momentum.lean\#L26-L35}{\texttt{waveTest\_product}}).
  Thus, for every proposed finite sum, some such pair \(f,g\) gives
  different readings.
\end{enumerate}

With \(A=3\dim\mathfrak g/\pi^4\) from VII.2, (a) is clause (T4) of the
main theorem. \emph{Status:} theorem. (a) is assigned to VII.2 and (b)
to VIII.1. \emph{Lean:} (a)
\href{https://github.com/yang-mills-group/OpenAI-YM-20261028/blob/3456e5e75cf8d4cd558955a5c7638fba2d68f464/YM/TraceExtension/Wick/Refute.lean\#L75-L80}{\texttt{Wick.euclidean\_excludes}},
\href{https://github.com/yang-mills-group/OpenAI-YM-20261028/blob/3456e5e75cf8d4cd558955a5c7638fba2d68f464/YM/TraceExtension/Wick/Log.lean\#L96-L108}{\texttt{Wick.no\_powerlog}};
public form
\href{https://github.com/yang-mills-group/OpenAI-YM-20261028/blob/3456e5e75cf8d4cd558955a5c7638fba2d68f464/TraceBridge.lean\#L35-L44}{\texttt{publicTrace}}
for
\href{https://github.com/yang-mills-group/OpenAI-YM-20261028/blob/3456e5e75cf8d4cd558955a5c7638fba2d68f464/ChallengeSpec.lean\#L168-L175}{\texttt{Hierarchy.TraceTest}}.
(b)
\href{https://github.com/yang-mills-group/OpenAI-YM-20261028/blob/3456e5e75cf8d4cd558955a5c7638fba2d68f464/YM/TraceExtension/Wick/Refute.lean\#L86-L105}{\texttt{Wick.polynomial\_excludes}}.
The asymptotic law and related polynomial-field exclusions occur in
\href{https://github.com/yang-mills-group/OpenAI-YM-20261028/blob/3456e5e75cf8d4cd558955a5c7638fba2d68f464/YM/Nonperturbative/HeadlineSpec.lean\#L84-L120}{\texttt{Headline.TraceExclusions}}.

The Euclidean conclusion concerns the single function \(G_\Theta\) along
one axis. The excluded class allows arbitrary complex coefficients and
continuous homogeneous numerators, but only finitely many modes and
products. It does not include general continuous-mass generalized free
fields or infinite Wick expansions.

\subsubsection{2. What it uses}\label{what-it-uses-27}

\textbf{Inputs from other units.}

\begin{enumerate}
\def\labelenumi{\arabic{enumi}.}
\tightlist
\item
  \textbf{III.16 (analytic toolkit).}

  \begin{itemize}
  \tightlist
  \item
    The free curvature kernel, Wick moments, on-shell jet spaces, label
    evaluation map and real homogeneous reference frames defined above.
  \item
    The response identity
    \href{https://github.com/yang-mills-group/OpenAI-YM-20261028/blob/3456e5e75cf8d4cd558955a5c7638fba2d68f464/YM/Extra/FREE_DETECTOR_ALGEBRA/CurlDetect.lean\#L28-L31}{\texttt{replyTest}}:
    for \(P\in\mathrm{Sol}_j\) and \(Q\in\mathrm{Sol}_k\), the
    covariance of \(\langle P,T_jF\rangle(z)\) with
    \(\langle Q,T_kF\rangle(0)\) equals
    \((-1)^k\langle Q,T_k[R_P](https://github.com/scottnarmstrong/yang-mills-informalization/blob/main/outline/units/z)\rangle\).
    Here \(R_P\) is the curvature field produced by \(P\) and \(T_k\) is
    the \(k\)-th Taylor coefficient.
  \item
    Used in VII.3.1--VII.3.5.
  \end{itemize}
\item
  \textbf{IV.1.} The free table, its on-shell image and spectator tests;
  correlation tables, null and quotient field spaces, real conjugation
  and signed-permutation symmetry, column systems and generic
  configurations.
\item
  \textbf{VI.2 (free Maxwell stress tensor).} For simple nonabelian
  \(\mathfrak g\):

  \begin{itemize}
  \tightlist
  \item
    \(E_d=\mathbb C\cdot1\) for \(d<4\)
    (\href{https://github.com/yang-mills-group/OpenAI-YM-20261028/blob/3456e5e75cf8d4cd558955a5c7638fba2d68f464/YM/UVNext/FREE_STRESS_CERTIFICATE/Size.lean\#L46-L49}{\texttt{below}},
    \href{https://github.com/yang-mills-group/OpenAI-YM-20261028/blob/3456e5e75cf8d4cd558955a5c7638fba2d68f464/YM/UVNext/FREE_STRESS_CERTIFICATE/Size.lean\#L51-L54}{\texttt{dim\_before}}).
  \item
    \(E_4\) has the basis consisting of the unit and the 21 invariant
    quadratics \(\langle F_P,F_Q\rangle\), one for each unordered pair
    \(\{P,Q\}\) of coordinate planes. Its dual functionals are built
    from the one-point function and the two-point functions at
    separation \(e\)
    (\href{https://github.com/yang-mills-group/OpenAI-YM-20261028/blob/3456e5e75cf8d4cd558955a5c7638fba2d68f464/YM/UVNext/FREE_STRESS_CERTIFICATE/Size.lean\#L105-L110}{\texttt{stand}},
    \href{https://github.com/yang-mills-group/OpenAI-YM-20261028/blob/3456e5e75cf8d4cd558955a5c7638fba2d68f464/YM/UVNext/FREE_STRESS_CERTIFICATE/Size.lean\#L82-L94}{\texttt{separates}},
    \href{https://github.com/yang-mills-group/OpenAI-YM-20261028/blob/3456e5e75cf8d4cd558955a5c7638fba2d68f464/YM/UVNext/FREE_STRESS_CERTIFICATE/Size.lean\#L78-L80}{\texttt{portLine}},
    \href{https://github.com/yang-mills-group/OpenAI-YM-20261028/blob/3456e5e75cf8d4cd558955a5c7638fba2d68f464/YM/UVNext/FREE_STRESS_CERTIFICATE/Size.lean\#L69-L75}{\texttt{up\_to}}).
  \item
    The free curvature two-point function at \(e_t\) is diagonal in
    planes, with entries \(\kappa_t(P)=-\pi^{-2}\) if \(t\in P\) and
    \(+\pi^{-2}\) otherwise
    (\href{https://github.com/yang-mills-group/OpenAI-YM-20261028/blob/3456e5e75cf8d4cd558955a5c7638fba2d68f464/YM/UVNext/FREE_STRESS_CERTIFICATE/Diagonal.lean\#L39-L67}{\texttt{bare\_diagonal}},
    \href{https://github.com/yang-mills-group/OpenAI-YM-20261028/blob/3456e5e75cf8d4cd558955a5c7638fba2d68f464/YM/UVNext/FREE_STRESS_CERTIFICATE/Diagonal.lean\#L34}{\texttt{coat}}).
  \item
    Also from VI.2:
    \href{https://github.com/yang-mills-group/OpenAI-YM-20261028/blob/3456e5e75cf8d4cd558955a5c7638fba2d68f464/YM/Pert/FirstTarget.lean\#L26-L39}{\texttt{FirstExtension}}.
    Along \(X\):

    \begin{itemize}
    \tightlist
    \item
      the rescaled correlations in the frames \(U_k\) converge in every
      \(C^j\) on compacta of \(\mathrm{Conf}_n\) to those of the free
      table, with the stress tensor compared to the Maxwell stress
      tensor;
    \item
      the centered real symmetric degree-four tensors that transform
      tensorially under signed permutations and have zero divergence
      form a one-dimensional real space;
    \item
      derivatives and rotation generators converge.
    \end{itemize}
  \item
    This leading comparison is a hypothesis of VII.3.6(c).
  \end{itemize}
\item
  \textbf{VI.1.} Scaling sequences of radii, stress normalizations and
  degreewise complex-linear frame changes.
\item
  \textbf{VII.1.}
  \href{https://github.com/yang-mills-group/OpenAI-YM-20261028/blob/3456e5e75cf8d4cd558955a5c7638fba2d68f464/YM/Pert/Rig/Linked.lean\#L59-L69}{\texttt{Rig.zerothProfile}}:
  the order-zero coefficient of the lattice perturbative profile
  converges to the tested free reading. It supplies the lattice limit in
  VII.3.r.
\item
  \textbf{VIII.1.}

  \begin{itemize}
  \tightlist
  \item
    The final construction data, including the selected phase and its
    reconstructed theory.
  \item
    \href{https://github.com/yang-mills-group/OpenAI-YM-20261028/blob/3456e5e75cf8d4cd558955a5c7638fba2d68f464/YM/Pert/Position/Crest.lean\#L36-L43}{\texttt{calibrate}}
    and
    \href{https://github.com/yang-mills-group/OpenAI-YM-20261028/blob/3456e5e75cf8d4cd558955a5c7638fba2d68f464/YM/Pert/Position/Crest.lean\#L50-L52}{\texttt{credible}}:
    the four implicit-position conditions hold for every clean square
    system, and the leading lattice comparison holds for every clean
    system.
  \item
    \href{https://github.com/yang-mills-group/OpenAI-YM-20261028/blob/3456e5e75cf8d4cd558955a5c7638fba2d68f464/YM/Pert/Position/Implicit.lean\#L34-L35}{\texttt{Reports.of\_square}}
    (implicit function theorem).
  \end{itemize}
\item
  \textbf{VII.2.} For VII.3.7:

  \begin{itemize}
  \tightlist
  \item
    the short-distance law
    \href{https://github.com/yang-mills-group/OpenAI-YM-20261028/blob/3456e5e75cf8d4cd558955a5c7638fba2d68f464/YM/TraceExtension/LogScaling.lean\#L74-L75}{\texttt{EuclideanAsymptotic}}
    with constant \(3\dim\mathfrak g/\pi^4\);
  \item
    the expansion lemmas
    \href{https://github.com/yang-mills-group/OpenAI-YM-20261028/blob/3456e5e75cf8d4cd558955a5c7638fba2d68f464/YM/TraceExtension/Wick/Massive.lean\#L122-L131}{\texttt{moment\_expands}},
    \href{https://github.com/yang-mills-group/OpenAI-YM-20261028/blob/3456e5e75cf8d4cd558955a5c7638fba2d68f464/YM/TraceExtension/Wick/PolynomialField.lean\#L116-L125}{\texttt{Track.moments\_expand}},
    \href{https://github.com/yang-mills-group/OpenAI-YM-20261028/blob/3456e5e75cf8d4cd558955a5c7638fba2d68f464/YM/TraceExtension/Wick/Refute.lean\#L18-L22}{\texttt{real\_profile}}.
  \end{itemize}
\item
  \textbf{V.3 and VII.2.} The selected trace field has a positive-energy
  spectral measure \(\mu\), its two-point function is the Laplace
  transform of \(\mu\), and its Wightman two-point distribution is the
  corresponding momentum-space integral. Exponential damping is
  integrable for every positive time
  (\href{https://github.com/yang-mills-group/OpenAI-YM-20261028/blob/3456e5e75cf8d4cd558955a5c7638fba2d68f464/YM/TraceExtension/WilsonTrace.lean\#L83-L94}{\texttt{TraceSpectrum}},
  \href{https://github.com/yang-mills-group/OpenAI-YM-20261028/blob/3456e5e75cf8d4cd558955a5c7638fba2d68f464/YM/TraceExtension/WilsonTrace.lean\#L99-L119}{\texttt{traceSpectrum\_exists}}).
  This is used only in VII.3.7(b).
\item
  \textbf{I.1.} The lattice, the Wilson action and the Haar measure (for
  §4.7).
\end{enumerate}

\textbf{Standard mathematics.}

\begin{itemize}
\tightlist
\item
  Wick's theorem, Gaussian integration by parts, and Fock-space
  orthogonality of Wick monomials: Glimm--Jaffe, \emph{Quantum Physics},
  ch.~6; B. Simon, \emph{The \(P(\phi)_2\) Euclidean (Quantum) Field
  Theory}, ch.~I; S. Janson, \emph{Gaussian Hilbert Spaces}, ch.~3.
\item
  Reflection positivity of the free field: Osterwalder--Schrader, CMP 31
  (1973) and 42 (1975); Glimm--Jaffe, ch.~6. It is not invoked: strict
  positivity is proved by explicit diagonalization.
\item
  Hobson's formula for \(h(\partial)|x|^{-2}\) in four dimensions, \(h\)
  harmonic and homogeneous: Hobson, \emph{The Theory of Spherical and
  Ellipsoidal Harmonics}, ch.~III; Stein--Weiss, \emph{Introduction to
  Fourier Analysis on Euclidean Spaces}, ch.~IV.
\item
  Invariance of the four-dimensional Maxwell equations under inversion:
  Bateman (1910), Cunningham (1910).
\item
  Existence of orthogonal bases for symmetric bilinear forms; the
  real-analytic implicit function theorem; characters of the group
  \((\mathbb Z/2)^4\) of axis reflections.
\item
  Small-\(r\) expansions of Laplace integrals with logarithms, as for
  modified Bessel functions: Watson, \emph{A Treatise on the Theory of
  Bessel Functions}, ch.~III.
\item
  Lattice perturbation theory with Faddeev--Popov gauge fixing:
  Montvay--Münster, \emph{Quantum Fields on a Lattice}, ch.~3--4.
\end{itemize}

\textbf{Imported constants and assumptions.}

\begin{itemize}
\tightlist
\item
  Constants used:

  \begin{itemize}
  \tightlist
  \item
    \(\Gamma=1/(4\pi^2|z|^2)\);
  \item
    \(s_j=(j+1)(j+2)(-2)^j/\pi^2\), the code's
    \href{https://github.com/yang-mills-group/OpenAI-YM-20261028/blob/3456e5e75cf8d4cd558955a5c7638fba2d68f464/YM/Extra/FREE_DETECTOR_ALGEBRA/Kelvin.lean\#L46}{\texttt{strength}};
  \item
    \(\pi_j(z)=(-1)^js_j|z|^{-2j-4}\);
  \item
    \(\kappa_t(P)=\mp\pi^{-2}\) (VI.2).
  \end{itemize}
\item
  Hypotheses: the leading free comparison (VI.2) and the four
  implicit-position conditions (VIII.1) for VII.3.6(c); the
  short-distance law with \(A>0\) for VII.3.7 (VII.2).
\item
  Simplicity of \(\mathfrak g\) enters VII.3.6 through VI.2: invariant
  linear color forms vanish, invariant bilinear color forms are
  multiples of the inner product, and \(\dim\mathfrak g>0\).
\end{itemize}

\subsubsection{3. Strategy}\label{strategy-27}

The reflected covariance couples different jet orders. Hobson's formula
identifies the field of an order-\(j\) multipole with \(s_j\) times its
Kelvin transform. Translating the polynomial multipole by the negative
inverted separation and rescaling its Taylor components produces a
triangular change of probes whose response has no Taylor terms below
order \(j\) at the second point. The reflected covariance becomes
diagonal, with positive pivots \(\pi_j\). Wick recursion extends this to
polynomial monomials, with factorial norms \(\alpha!d^\alpha\).

Strict positivity detects every nonzero on-shell label. It identifies
the free quotient and gives injective two-point columns. At degree four,
the explicit quadratic basis reduces buffer rigidity to a coefficient
calculation: two-point stationarity gives weighted antisymmetry, while
signed permutations give symmetry or zero. Compactly supported tests
convert this rigidity into finitely many equations with invertible
derivative. The implicit function theorem solves them along the scaling
sequence.

Independently, finite Wick profiles have expansions in integer powers of
distance and nonnegative powers of its logarithm. Their real parts
cannot have the trace field's leading term \(A r^{-8}(\log(1/r))^{-2}\).
Comparison along \(r=2^{-n}\) proves the contradiction. The spectral
representation transfers equality of tested Wightman two-point functions
to equality of the Euclidean profiles.

\subsubsection{4. Steps}\label{steps-27}

\paragraph{4.1 The Wick algebra}\label{the-wick-algebra}

\textbf{Lemma VII.3.a (Wick recursion).} Proposition VII.3.3.
\emph{Proof sketch.}

\begin{itemize}
\tightlist
\item
  By multilinearity it suffices to treat monomials. Occurrences of
  variables are distinguishable copies.
\item
  Adjoin one occurrence of \(X_s\) at \(x_i\). Matchings of the enlarged
  system correspond bijectively to pairs (a partner occurrence of some
  \(X_t\) at some \(x_k\) with \(k\ne i\), a matching of the reduced
  system). Here the reduced system has one copy of \(X_t\) removed at
  \(x_k\).
\item
  The multiplicity of \(X_t\) in \(W_k\) is the coefficient produced by
  \(\partial_s\).
\item
  Uniqueness: induct on the number of non-unit entries and on the
  polynomial structure. The recursion and the normalization
  \(T(1,\dots,1)=1\) determine \(T\).
\item
  Rotated and conjugated moments satisfy the same recursion, by
  covariance and reality of the kernel. This gives \(O(4)\)-invariance
  and reality.
\end{itemize}

\emph{Lean:}
\href{https://github.com/yang-mills-group/OpenAI-YM-20261028/blob/3456e5e75cf8d4cd558955a5c7638fba2d68f464/YM/Extra/FREE_DETECTOR_ALGEBRA/Removal.lean\#L91-L108}{\texttt{Pair.remove}},
\href{https://github.com/yang-mills-group/OpenAI-YM-20261028/blob/3456e5e75cf8d4cd558955a5c7638fba2d68f464/YM/Extra/FREE_DETECTOR_ALGEBRA/Removal.lean\#L64-L86}{\texttt{Pair.remove\_mon}},
bijections in
\href{https://github.com/yang-mills-group/OpenAI-YM-20261028/blob/3456e5e75cf8d4cd558955a5c7638fba2d68f464/YM/Extra/FREE_DETECTOR_ALGEBRA/Match.lean}{\texttt{Match.lean}}
and
\href{https://github.com/yang-mills-group/OpenAI-YM-20261028/blob/3456e5e75cf8d4cd558955a5c7638fba2d68f464/YM/Extra/FREE_DETECTOR_ALGEBRA/Occurrences.lean}{\texttt{Occurrences.lean}};
\href{https://github.com/yang-mills-group/OpenAI-YM-20261028/blob/3456e5e75cf8d4cd558955a5c7638fba2d68f464/YM/Extra/FREE_DETECTOR_ALGEBRA/Gaussian.lean\#L24-L62}{\texttt{byAnnihilators}},
\href{https://github.com/yang-mills-group/OpenAI-YM-20261028/blob/3456e5e75cf8d4cd558955a5c7638fba2d68f464/YM/Extra/FREE_DETECTOR_ALGEBRA/Gaussian.lean\#L81-L101}{\texttt{wick\_rotate}},
\href{https://github.com/yang-mills-group/OpenAI-YM-20261028/blob/3456e5e75cf8d4cd558955a5c7638fba2d68f464/YM/Extra/FREE_DETECTOR_ALGEBRA/Gaussian.lean\#L111-L137}{\texttt{wick\_bar}}.

\textbf{Lemma VII.3.b (two-point adjoint relation).} For
\(p\in\mathrm{Tow}_n\) and \(x\ne y\), let \(D^{x,y}_p\) be the
derivation of \(\mathcal W\) sending each variable at \(y\) to its free
covariance with \(\ell_p\) at \(x\)
(\href{https://github.com/yang-mills-group/OpenAI-YM-20261028/blob/3456e5e75cf8d4cd558955a5c7638fba2d68f464/YM/Extra/FREE_DETECTOR_ALGEBRA/Script.lean\#L81}{\texttt{crash}}).
For all \(P,Q\in\mathcal W\) and \(q\in\mathrm{Tow}_n\):

\begin{itemize}
\tightlist
\item
  \(\langle (P\ell_p)(x)\,Q(y)\rangle_0=\langle P(x)\,(D^{x,y}_pQ)(y)\rangle_0\);
\item
  \(D^{x,y}_p(\mathcal R_{x-y}\ell_q)=B_{x-y}(p,q)\), a constant.
\end{itemize}

\emph{Proof sketch.} Apply VII.3.a with \(n=2\): the sum over \(k\ne i\)
has a single term. On linear polynomials, the resulting contraction is
exactly \(B_{x-y}(p,q)\). \emph{Lean:}
\href{https://github.com/yang-mills-group/OpenAI-YM-20261028/blob/3456e5e75cf8d4cd558955a5c7638fba2d68f464/YM/Extra/FREE_DETECTOR_ALGEBRA/Duplex.lean\#L60-L66}{\texttt{join\_crash}},
\href{https://github.com/yang-mills-group/OpenAI-YM-20261028/blob/3456e5e75cf8d4cd558955a5c7638fba2d68f464/YM/Extra/FREE_DETECTOR_ALGEBRA/Script.lean\#L96-L100}{\texttt{crash\_relation}},
\href{https://github.com/yang-mills-group/OpenAI-YM-20261028/blob/3456e5e75cf8d4cd558955a5c7638fba2d68f464/YM/Extra/FREE_DETECTOR_ALGEBRA/Script.lean\#L83-L94}{\texttt{crash\_test}}.

\textbf{Lemma VII.3.c (exhaustion by towers).} For every
\(P\in\mathcal W\) there is \(n_0\ge0\) with the following property. For
all \(n\ge n_0\) and every real basis \((b_i)_{1\le i\le N}\) of
\(\mathrm{Tow}_n\), one has \(P=\Phi_b(p)\) for some
\(p\in\mathbb C[Y_1,\dots,Y_N]\). \emph{Proof sketch.}

\begin{itemize}
\tightlist
\item
  Each variable \(X_{(j,a)}\) equals \(\ell_r\), where \(r\) is the
  Riesz dual of the \(a\)-th coordinate functional, placed in order
  \(j\)
  (\href{https://github.com/yang-mills-group/OpenAI-YM-20261028/blob/3456e5e75cf8d4cd558955a5c7638fba2d68f464/YM/Extra/FREE_DETECTOR_ALGEBRA/Symbol.lean\#L95-L96}{\texttt{digit\_coeff}}).
\item
  \(\ell_r\) is a real linear combination of the \(\ell_{b_i}\) once
  \(n>j\).
\item
  Induct on the structure of \(P\): constants, sums, and products with
  one variable.
\end{itemize}

\emph{Lean:}
\href{https://github.com/yang-mills-group/OpenAI-YM-20261028/blob/3456e5e75cf8d4cd558955a5c7638fba2d68f464/YM/Extra/FREE_DETECTOR_ALGEBRA/Script.lean\#L51-L68}{\texttt{from\_large}},
\href{https://github.com/yang-mills-group/OpenAI-YM-20261028/blob/3456e5e75cf8d4cd558955a5c7638fba2d68f464/YM/Extra/FREE_DETECTOR_ALGEBRA/Script.lean\#L37-L46}{\texttt{from\_basis}}.

\paragraph{4.2 The tower Gram form}\label{the-tower-gram-form}

\textbf{Lemma VII.3.d (Hobson's formula in \(\mathbb R^4\)).} Let
\(S\in\mathrm{Sol}_j\) and \(x\ne0\). Let \(S(\partial)\) be the
differential operator with symbol \(S\), normalized with Taylor
coefficients (a factor \(1/j!\)). Then

\[S(\partial)\Gamma(x)=(-2)^j\,|x|^{-2j}\,S(x)\,\Gamma(x).\]

\emph{Proof sketch.} Induct on \(j\). In the step from \(j\) to \(j+1\),
take \(S\in\mathrm{Sol}_{j+1}\) and let \(S_a\in\mathrm{Sol}_j\) be its
partial contraction with \(e_a\). Taylor normalization gives
\(S(\partial)=\frac1{j+1}\sum_a\partial_aS_a(\partial)\)
(\href{https://github.com/yang-mills-group/OpenAI-YM-20261028/blob/3456e5e75cf8d4cd558955a5c7638fba2d68f464/YM/Extra/FREE_DETECTOR_ALGEBRA/Solutions.lean\#L140-L159}{\texttt{Sol.freeze}}).
Apply the induction hypothesis and differentiate:

\begin{itemize}
\tightlist
\item
  the term with the divergence of the partial contractions vanishes,
  because \(S\) is harmonic
  (\href{https://github.com/yang-mills-group/OpenAI-YM-20261028/blob/3456e5e75cf8d4cd558955a5c7638fba2d68f464/YM/Extra/FREE_DETECTOR_ALGEBRA/Solutions.lean\#L164-L172}{\texttt{Sol.freeze\_trace}});
\item
  the radial derivative gives \(-2(j+1)|x|^{-2(j+2)}\sum_ax_aS_a(x)\),
  and \(\sum_ax_aS_a(x)=S(x)\)
  (\href{https://github.com/yang-mills-group/OpenAI-YM-20261028/blob/3456e5e75cf8d4cd558955a5c7638fba2d68f464/YM/Extra/FREE_DETECTOR_ALGEBRA/Harmonic.lean\#L109-L119}{\texttt{freeze\_diagonal}}).
  The factor \(j+1\) cancels, leaving the required factor \(-2\).
\end{itemize}

\emph{Lean:}
\href{https://github.com/yang-mills-group/OpenAI-YM-20261028/blob/3456e5e75cf8d4cd558955a5c7638fba2d68f464/YM/Extra/FREE_DETECTOR_ALGEBRA/Harmonic.lean\#L122-L152}{\texttt{shell\_green}}.

\textbf{Lemma VII.3.e (the field of a multipole is a Kelvin transform).}
For \(P\in\mathrm{Sol}_j\) and \(z\ne0\), the curvature field \(R_P\)
produced by \(P\) satisfies

\[R_P(z)=s_j\,|z|^{-4}\,(\mathcal R_z\otimes\mathcal R_z)\,P(\check z)=s_j\,|z|^{-2j-4}\,(\mathcal R_z\otimes\mathcal R_z)\,P(z),\qquad \check z=z/|z|^2 .\]

That is, \(R_P=s_j\,I^\ast P\), the pull-back of the 2-form \(P\) under
the inversion \(I\), whose differential is \(|z|^{-2}\mathcal R_z\).
\emph{Proof sketch.}

\begin{itemize}
\tightlist
\item
  The curvature kernel is a curl-curl of \(\Gamma\). Its contraction
  with \(P\) is \((R_P)_{im}=2(\partial_m(D_P)_i-\partial_i(D_P)_m)\),
  where \((D_P)_i=\sum_a\partial_a(P_{ai}(\partial)\Gamma)\). By
  VII.3.d,
\end{itemize}

\[D_P(x)=-2(j+1)(-2)^j(4\pi^2)^{-1}|x|^{-2(j+2)}\,\iota_xP(x)\]

(\href{https://github.com/yang-mills-group/OpenAI-YM-20261028/blob/3456e5e75cf8d4cd558955a5c7638fba2d68f464/YM/Extra/FREE_DETECTOR_ALGEBRA/Kelvin.lean\#L59-L77}{\texttt{divP\_value}}).

\begin{itemize}
\tightlist
\item
  Its curl is computed with
  \(\partial_i(\iota_xP)_m-\partial_m(\iota_xP)_i=(j+2)P_{im}\). This
  uses Bianchi and Euler; see
  \href{https://github.com/yang-mills-group/OpenAI-YM-20261028/blob/3456e5e75cf8d4cd558955a5c7638fba2d68f464/YM/Extra/FREE_DETECTOR_ALGEBRA/Kelvin.lean\#L28-L44}{\texttt{radial\_curl}}.
\item
  Collecting terms gives
  \(s_j|x|^{-2(j+2)}\big(P_{im}-2|x|^{-2}(x_i(\iota_xP)_m-x_m(\iota_xP)_i)\big)\),
  which is the reflected tensor.
\end{itemize}

\emph{Lean:}
\href{https://github.com/yang-mills-group/OpenAI-YM-20261028/blob/3456e5e75cf8d4cd558955a5c7638fba2d68f464/YM/Extra/FREE_DETECTOR_ALGEBRA/Kelvin.lean\#L82-L109}{\texttt{curl\_kelvin}},
\href{https://github.com/yang-mills-group/OpenAI-YM-20261028/blob/3456e5e75cf8d4cd558955a5c7638fba2d68f464/YM/Extra/FREE_DETECTOR_ALGEBRA/Kelvin.lean\#L111-L113}{\texttt{reply\_kelvin}},
\href{https://github.com/yang-mills-group/OpenAI-YM-20261028/blob/3456e5e75cf8d4cd558955a5c7638fba2d68f464/YM/Extra/FREE_DETECTOR_ALGEBRA/LegAct.lean\#L112-L122}{\texttt{reply\_radial}},
\href{https://github.com/yang-mills-group/OpenAI-YM-20261028/blob/3456e5e75cf8d4cd558955a5c7638fba2d68f464/YM/Extra/FREE_DETECTOR_ALGEBRA/Kelvin.lean\#L51-L55}{\texttt{strength\_sign}}.

\textbf{Lemma VII.3.f (shifted towers).} For \(x\ne0\), \(j<n\) and
\(P\in\mathrm{Sol}_j\), define
\(\sigma_x(P):=(Q_0,\dots,Q_j,0,\dots)\in\mathrm{Tow}_n\) with

\[Q_k:=(s_j/s_k)\,T_k[P](-\check x).\]

Here
\(T_k[P](https://github.com/scottnarmstrong/yang-mills-informalization/blob/main/outline/units/a)\in\mathrm{Sol}_k\)
is the order-\(k\) Taylor coefficient of \(P\) at \(a\). In particular
\(Q_j=P\). For every \(y\ne0\), the total field of the shifted tower is

\[\textstyle\sum_kR_{Q_k}(y)=s_j|y|^{-4}(\mathcal R_y\otimes\mathcal R_y)P(\check y-\check x).\]

It vanishes to order \(j\) at \(y=x\):

\[T_k\big[\textstyle\sum R_{Q_i}\big](https://github.com/scottnarmstrong/yang-mills-informalization/blob/main/outline/units/x)=0\ \ (k<j),\qquad T_j\big[\textstyle\sum R_{Q_i}\big](https://github.com/scottnarmstrong/yang-mills-informalization/blob/main/outline/units/x)=s_j|x|^{-2(j+2)}\,\mathcal R_xP .\]

\emph{Proof sketch.}

\begin{itemize}
\tightlist
\item
  Since
  \(\sum_kT_k[P](https://github.com/scottnarmstrong/yang-mills-informalization/blob/main/outline/units/a)(w)=P(w+a)\)
  and \(R\) is linear, VII.3.e gives the formula for the total field.
\item
  Along the line \(y=x+tv\):
  \(\check y-\check x=t\,|x+tv|^{-2}(\mathcal R_xv-t|v|^2\check x)\),
  which is smooth in \(t\).
\item
  By homogeneity the field is \(t^j\) times a smooth function whose
  value at \(t=0\) is \(s_j|x|^{-2(j+2)}(\mathcal R_xP)(v)\).
\item
  Taylor coefficients of symmetric jets are determined by their diagonal
  values.
\end{itemize}

\emph{Lean:}
\href{https://github.com/yang-mills-group/OpenAI-YM-20261028/blob/3456e5e75cf8d4cd558955a5c7638fba2d68f464/YM/Extra/FREE_DETECTOR_ALGEBRA/Shift.lean\#L20-L21}{\texttt{shifts}},
\href{https://github.com/yang-mills-group/OpenAI-YM-20261028/blob/3456e5e75cf8d4cd558955a5c7638fba2d68f464/YM/Extra/FREE_DETECTOR_ALGEBRA/Translation.lean\#L104-L113}{\texttt{towards}},
\href{https://github.com/yang-mills-group/OpenAI-YM-20261028/blob/3456e5e75cf8d4cd558955a5c7638fba2d68f464/YM/Extra/FREE_DETECTOR_ALGEBRA/Translation.lean\#L118-L119}{\texttt{field\_expands}},
\href{https://github.com/yang-mills-group/OpenAI-YM-20261028/blob/3456e5e75cf8d4cd558955a5c7638fba2d68f464/YM/Extra/FREE_DETECTOR_ALGEBRA/Shift.lean\#L27-L32}{\texttt{fields\_shift}},
\href{https://github.com/yang-mills-group/OpenAI-YM-20261028/blob/3456e5e75cf8d4cd558955a5c7638fba2d68f464/YM/Extra/FREE_DETECTOR_ALGEBRA/Shift.lean\#L68-L81}{\texttt{signal\_kelvin}},
\href{https://github.com/yang-mills-group/OpenAI-YM-20261028/blob/3456e5e75cf8d4cd558955a5c7638fba2d68f464/YM/Extra/FREE_DETECTOR_ALGEBRA/Shift.lean\#L40-L56}{\texttt{invert\_line}},
\href{https://github.com/yang-mills-group/OpenAI-YM-20261028/blob/3456e5e75cf8d4cd558955a5c7638fba2d68f464/YM/Extra/FREE_DETECTOR_ALGEBRA/Shift.lean\#L135-L139}{\texttt{signal\_lower}},
\href{https://github.com/yang-mills-group/OpenAI-YM-20261028/blob/3456e5e75cf8d4cd558955a5c7638fba2d68f464/YM/Extra/FREE_DETECTOR_ALGEBRA/Shift.lean\#L141-L147}{\texttt{signal\_top}},
\href{https://github.com/yang-mills-group/OpenAI-YM-20261028/blob/3456e5e75cf8d4cd558955a5c7638fba2d68f464/YM/Extra/FREE_DETECTOR_ALGEBRA/PositiveLeg.lean\#L98-L99}{\texttt{newLeg}}.

\textbf{Lemma VII.3.g (diagonal Gram form in shifted towers).} For
\(x\ne0\), \(i,m<n\), \(P\in\mathrm{Sol}_i\) and \(Q\in\mathrm{Sol}_m\):

\[B_x(\sigma_xP,\sigma_xQ)=\begin{cases}\pi_i(x)\langle P,Q\rangle & i=m,\\ 0 & i\ne m.\end{cases}\]

Hence \(B_x(\sigma_xp,\sigma_xq)=\sum_i\pi_i(x)\langle p_i,q_i\rangle\),
where \(\sigma_xp:=\sum_i\sigma_x(p_i)\). \emph{Proof sketch.}

\begin{itemize}
\tightlist
\item
  By the response identity (III.16), for \(q\in\mathrm{Sol}_k\) placed
  in order \(k\):
  \(B_x(\sigma_xP,q)=(-1)^k\big\langle\mathcal R_xq,\ T_k[\textstyle\sum R_{Q_i}](https://github.com/scottnarmstrong/yang-mills-informalization/blob/main/outline/units/x)\big\rangle\).
\item
  By VII.3.f this is \(0\) for \(k<i\), and
  \((-1)^is_i|x|^{-2(i+2)}\langle\mathcal R_xq,\mathcal R_xP\rangle=\pi_i(x)\langle P,q\rangle\)
  for \(k=i\), since \(\mathcal R_x\) is orthogonal.
\item
  If \(m<i\), expand \(\sigma_xQ\) into components of orders \(\le m\);
  each gives zero. The case \(m>i\) follows by symmetry.
\item
  If \(m=i\), only the top component \(Q\) of \(\sigma_xQ\) contributes.
\end{itemize}

\emph{Lean:}
\href{https://github.com/yang-mills-group/OpenAI-YM-20261028/blob/3456e5e75cf8d4cd558955a5c7638fba2d68f464/YM/Extra/FREE_DETECTOR_ALGEBRA/PositiveLeg.lean\#L49-L57}{\texttt{slide\_sum}},
\href{https://github.com/yang-mills-group/OpenAI-YM-20261028/blob/3456e5e75cf8d4cd558955a5c7638fba2d68f464/YM/Extra/FREE_DETECTOR_ALGEBRA/PositiveLeg.lean\#L59-L60}{\texttt{slide\_below}},
\href{https://github.com/yang-mills-group/OpenAI-YM-20261028/blob/3456e5e75cf8d4cd558955a5c7638fba2d68f464/YM/Extra/FREE_DETECTOR_ALGEBRA/PositiveLeg.lean\#L62-L67}{\texttt{slide\_top}},
\href{https://github.com/yang-mills-group/OpenAI-YM-20261028/blob/3456e5e75cf8d4cd558955a5c7638fba2d68f464/YM/Extra/FREE_DETECTOR_ALGEBRA/PositiveLeg.lean\#L105-L111}{\texttt{newLeg\_below}},
\href{https://github.com/yang-mills-group/OpenAI-YM-20261028/blob/3456e5e75cf8d4cd558955a5c7638fba2d68f464/YM/Extra/FREE_DETECTOR_ALGEBRA/PositiveLeg.lean\#L113-L121}{\texttt{newLeg\_diag}},
\href{https://github.com/yang-mills-group/OpenAI-YM-20261028/blob/3456e5e75cf8d4cd558955a5c7638fba2d68f464/YM/Extra/FREE_DETECTOR_ALGEBRA/PositiveLeg.lean\#L126-L133}{\texttt{shiftAll\_gram}}.

\textbf{Lemma VII.3.h (VII.3.2(a)).} For every \(z\ne0\) and \(n\ge0\),
\(B_z\) is symmetric and positive definite on \(\mathrm{Tow}_n\).
\emph{Proof sketch.}

\begin{itemize}
\tightlist
\item
  Symmetry: the kernel is symmetric under exchanging the two probes and
  replacing \(z\) by \(-z\). It is also \(O(4)\)-covariant
  (\href{https://github.com/yang-mills-group/OpenAI-YM-20261028/blob/3456e5e75cf8d4cd558955a5c7638fba2d68f464/YM/Extra/FREE_DETECTOR_ALGEBRA/LegAct.lean\#L156-L160}{\texttt{leg\_cov}}).
  Apply this with \(\mathcal R_zz=-z\) and \(\mathcal R_z^2=1\).
\item
  Positivity:

  \begin{itemize}
  \tightlist
  \item
    \(\sigma_z\) is injective, since \(B_z(\sigma_zp,\sigma_zp)>0\) for
    \(p\ne0\) by VII.3.g;
  \item
    hence it is bijective, by finite dimension (it is also visibly block
    unitriangular);
  \item
    so \(B_z(p,p)=\sum\pi_j\lVert p'_j\rVert^2>0\) for
    \(p=\sigma_zp'\ne0\).
  \end{itemize}
\end{itemize}

\emph{Lean:}
\href{https://github.com/yang-mills-group/OpenAI-YM-20261028/blob/3456e5e75cf8d4cd558955a5c7638fba2d68f464/YM/Extra/FREE_DETECTOR_ALGEBRA/PositiveLeg.lean\#L37-L47}{\texttt{scan\_sym}},
\href{https://github.com/yang-mills-group/OpenAI-YM-20261028/blob/3456e5e75cf8d4cd558955a5c7638fba2d68f464/YM/Extra/FREE_DETECTOR_ALGEBRA/PositiveLeg.lean\#L90-L94}{\texttt{towerB\_sym}},
\href{https://github.com/yang-mills-group/OpenAI-YM-20261028/blob/3456e5e75cf8d4cd558955a5c7638fba2d68f464/YM/Extra/FREE_DETECTOR_ALGEBRA/PositiveLeg.lean\#L135-L141}{\texttt{shiftAll\_pos}},
\href{https://github.com/yang-mills-group/OpenAI-YM-20261028/blob/3456e5e75cf8d4cd558955a5c7638fba2d68f464/YM/Extra/FREE_DETECTOR_ALGEBRA/PositiveLeg.lean\#L143-L146}{\texttt{shiftDual}},
\href{https://github.com/yang-mills-group/OpenAI-YM-20261028/blob/3456e5e75cf8d4cd558955a5c7638fba2d68f464/YM/Extra/FREE_DETECTOR_ALGEBRA/PositiveLeg.lean\#L150-L154}{\texttt{towerB\_pos}}.

\paragraph{4.3 From towers to all Wick
polynomials}\label{from-towers-to-all-wick-polynomials}

\textbf{Lemma VII.3.i (Fock orthogonality).} Let \(b\) be as in
VII.3.2(b), and let \(\alpha,\beta\in\mathbb N^N\). Then, without
conjugation,

\[\big\langle\Phi_b(Y^\alpha)(x)\,(\mathcal R_{x-y}\Phi_b(Y^\beta))(y)\big\rangle_0=\delta_{\alpha\beta}\,\alpha!\,d^\alpha .\]

\emph{Proof sketch.}

\begin{itemize}
\tightlist
\item
  Since the \(b_k\) are orthogonal, VII.3.b gives
  \(D^{x,y}_{b_i}\,\mathcal R\Phi_b(q)=d_i\,\mathcal R\Phi_b(\partial_{Y_i}q)\)
  (\href{https://github.com/yang-mills-group/OpenAI-YM-20261028/blob/3456e5e75cf8d4cd558955a5c7638fba2d68f464/YM/Extra/FREE_DETECTOR_ALGEBRA/Duplex.lean\#L88-L102}{\texttt{crash\_basis}}).
\item
  Thus transferring one factor \(Y_i\) from the first insertion
  differentiates the second polynomial and multiplies by \(d_i\)
  (\href{https://github.com/yang-mills-group/OpenAI-YM-20261028/blob/3456e5e75cf8d4cd558955a5c7638fba2d68f464/YM/Extra/FREE_DETECTOR_ALGEBRA/Duplex.lean\#L104-L109}{\texttt{aligned\_step}}).
\item
  Induct on \(\alpha\), one factor at a time
  (\href{https://github.com/yang-mills-group/OpenAI-YM-20261028/blob/3456e5e75cf8d4cd558955a5c7638fba2d68f464/YM/Extra/FREE_DETECTOR_ALGEBRA/PowerGram.lean\#L44-L52}{\texttt{tally\_step}}).
\item
  The base cases are the unit-unit reading \(1\) and the zero reading of
  the unit against a nonconstant monomial
  (\href{https://github.com/yang-mills-group/OpenAI-YM-20261028/blob/3456e5e75cf8d4cd558955a5c7638fba2d68f464/YM/Extra/FREE_DETECTOR_ALGEBRA/Duplex.lean\#L81-L84}{\texttt{aligned\_unit}},
  \href{https://github.com/yang-mills-group/OpenAI-YM-20261028/blob/3456e5e75cf8d4cd558955a5c7638fba2d68f464/YM/Extra/FREE_DETECTOR_ALGEBRA/Duplex.lean\#L76-L79}{\texttt{aligned\_one\_right}}).
\end{itemize}

This is the factorial orthogonality of Wick monomials. \emph{Lean:}
\href{https://github.com/yang-mills-group/OpenAI-YM-20261028/blob/3456e5e75cf8d4cd558955a5c7638fba2d68f464/YM/Extra/FREE_DETECTOR_ALGEBRA/PowerGram.lean\#L60-L85}{\texttt{aligned\_script}}.

\textbf{Lemma VII.3.j (VII.3.2(b)).} The monomials \(\Phi_b(Y^\alpha)\)
are real
(\href{https://github.com/yang-mills-group/OpenAI-YM-20261028/blob/3456e5e75cf8d4cd558955a5c7638fba2d68f464/YM/Extra/FREE_DETECTOR_ALGEBRA/PowerGram.lean\#L90-L95}{\texttt{bar\_mono}}).
Sesquilinear extension of VII.3.i then gives
\(\langle\Phi_bp,\Phi_bq\rangle_{x,y}=\sum_\alpha p_\alpha\overline{q_\alpha}\alpha!d^\alpha\).
\emph{Lean:}
\href{https://github.com/yang-mills-group/OpenAI-YM-20261028/blob/3456e5e75cf8d4cd558955a5c7638fba2d68f464/YM/Extra/FREE_DETECTOR_ALGEBRA/PowerGram.lean\#L100-L116}{\texttt{duplex\_script}}.

\textbf{Lemma VII.3.k (VII.3.1).} \emph{Proof sketch.}

\begin{itemize}
\tightlist
\item
  \begin{enumerate}
  \def\labelenumi{(\alph{enumi})}
  \tightlist
  \item
    Given \(P\ne0\), take \(n\) from VII.3.c. Take a
    \(B_{x-y}\)-orthogonal basis of \(\mathrm{Tow}_n\); it exists for a
    symmetric form, and \(d_i>0\) by VII.3.h
    (\href{https://github.com/yang-mills-group/OpenAI-YM-20261028/blob/3456e5e75cf8d4cd558955a5c7638fba2d68f464/YM/Extra/FREE_DETECTOR_ALGEBRA/WickPositive.lean\#L20-L29}{\texttt{upright}}).
    Then \(P=\Phi_b(p)\) with \(p\ne0\), and
    \(\langle P,P\rangle=\sum|p_\alpha|^2\alpha!d^\alpha>0\).
  \end{enumerate}
\item
  (b) Write \(P\) and \(Q\) in a common basis; the formula of VII.3.j is
  Hermitian.
\item
  (c) Apply (a) to \(\sum a_iV_i\).
\end{itemize}

\emph{Lean:}
\href{https://github.com/yang-mills-group/OpenAI-YM-20261028/blob/3456e5e75cf8d4cd558955a5c7638fba2d68f464/YM/Extra/FREE_DETECTOR_ALGEBRA/WickPositive.lean\#L47-L53}{\texttt{reflected\_positive}},
\href{https://github.com/yang-mills-group/OpenAI-YM-20261028/blob/3456e5e75cf8d4cd558955a5c7638fba2d68f464/YM/Extra/FREE_DETECTOR_ALGEBRA/WickPositive.lean\#L38-L44}{\texttt{script\_pos}},
\href{https://github.com/yang-mills-group/OpenAI-YM-20261028/blob/3456e5e75cf8d4cd558955a5c7638fba2d68f464/YM/Extra/FREE_DETECTOR_ALGEBRA/WickPositive.lean\#L57-L76}{\texttt{reflected\_symm}},
\href{https://github.com/yang-mills-group/OpenAI-YM-20261028/blob/3456e5e75cf8d4cd558955a5c7638fba2d68f464/YM/Extra/FREE_DETECTOR_ALGEBRA/WickPositive.lean\#L80-L97}{\texttt{reflected\_pd}}.

\paragraph{4.4 Faithfulness and covariance of the free
table}\label{faithfulness-and-covariance-of-the-free-table}

\textbf{Lemma VII.3.l (VII.3.4(a)).} \emph{Proof sketch.}

\begin{itemize}
\tightlist
\item
  \(\mathrm{sh}_d\) is stable under conjugation and under \(O(4)\),
  because \(\omega_d\) intertwines the label operations with
  \(\bar{\ }\) and with substitution
  (\href{https://github.com/yang-mills-group/OpenAI-YM-20261028/blob/3456e5e75cf8d4cd558955a5c7638fba2d68f464/YM/Extra/FREE_DETECTOR_ALGEBRA/Rank.lean\#L25-L26}{\texttt{sector\_bar}},
  \href{https://github.com/yang-mills-group/OpenAI-YM-20261028/blob/3456e5e75cf8d4cd558955a5c7638fba2d68f464/YM/Extra/FREE_DETECTOR_ALGEBRA/Rank.lean\#L27-L28}{\texttt{sector\_rot}},
  \href{https://github.com/yang-mills-group/OpenAI-YM-20261028/blob/3456e5e75cf8d4cd558955a5c7638fba2d68f464/YM/Extra/FREE_DETECTOR_ALGEBRA/LabelTurn.lean\#L56}{\texttt{observe\_rot}},
  \href{https://github.com/yang-mills-group/OpenAI-YM-20261028/blob/3456e5e75cf8d4cd558955a5c7638fba2d68f464/YM/Extra/FREE_DETECTOR_ALGEBRA/LabelTurn.lean\#L58}{\texttt{observe\_bar}}).
\item
  So for \(P\in\mathrm{sh}_d\) the one-element family
  \(Q_1=\mathcal R_{x-y}\bar P\) at \((x,y)=(0,e_{\rm t})\) is
  admissible. Its test value is \(\langle P,P\rangle_{0,e_{\rm t}}\)
  (\href{https://github.com/yang-mills-group/OpenAI-YM-20261028/blob/3456e5e75cf8d4cd558955a5c7638fba2d68f464/YM/Extra/FREE_DETECTOR_ALGEBRA/Rank.lean\#L31-L36}{\texttt{point\_probe}}),
  which is positive by VII.3.1 if \(P\ne0\).
\item
  The description of the null labels then follows from the identity
  between free-table readings and these tests (IV.1).
\end{itemize}

\emph{Lean:}
\href{https://github.com/yang-mills-group/OpenAI-YM-20261028/blob/3456e5e75cf8d4cd558955a5c7638fba2d68f464/YM/Extra/FREE_DETECTOR_ALGEBRA/Rank.lean\#L38-L54}{\texttt{no\_residual}},
\href{https://github.com/yang-mills-group/OpenAI-YM-20261028/blob/3456e5e75cf8d4cd558955a5c7638fba2d68f464/YM/Extra/FREE_DETECTOR_ALGEBRA/Rank.lean\#L58-L59}{\texttt{free\_null}},
\href{https://github.com/yang-mills-group/OpenAI-YM-20261028/blob/3456e5e75cf8d4cd558955a5c7638fba2d68f464/YM/Extra/FREE_DETECTOR_ALGEBRA/Rank.lean\#L64-L65}{\texttt{exactField}}.

\textbf{Lemma VII.3.m (VII.3.4(b), (c)).} \emph{Proof sketch of (b).}

\begin{itemize}
\tightlist
\item
  Suppose \(v\ne w\) have the same readings against a basis. By
  bilinearity, \(u=v-w\) has zero reading against every field.
\item
  In particular \(u\) has zero reading against
  \(\mathcal R_{x-y}\cdot\bar u\in E_d\)
  (\href{https://github.com/yang-mills-group/OpenAI-YM-20261028/blob/3456e5e75cf8d4cd558955a5c7638fba2d68f464/YM/Extra/FREE_DETECTOR_ALGEBRA/Columns.lean\#L96-L99}{\texttt{faceTest\_eq}}).
\item
  That reading is \(\langle\varpi u,\varpi u\rangle_{x,y}>0\), a
  contradiction.
\end{itemize}

\emph{Proof sketch of (c).}

\begin{itemize}
\tightlist
\item
  Extend the unit to a complex basis chosen among the images of frame
  labels. These images are fixed by conjugation
  (\href{https://github.com/yang-mills-group/OpenAI-YM-20261028/blob/3456e5e75cf8d4cd558955a5c7638fba2d68f464/YM/Extra/FREE_DETECTOR_ALGEBRA/Columns.lean\#L29-L69}{\texttt{labeled}},
  \href{https://github.com/yang-mills-group/OpenAI-YM-20261028/blob/3456e5e75cf8d4cd558955a5c7638fba2d68f464/YM/Extra/FREE_DETECTOR_ALGEBRA/Columns.lean\#L71-L83}{\texttt{frameRows}}).
\item
  Its elements carry their weighted degree, with an order-\(j\) jet
  weighted \(j+2\), and their length.
\item
  The pair readings are injective by (b).
\end{itemize}

\emph{Lean:}
\href{https://github.com/yang-mills-group/OpenAI-YM-20261028/blob/3456e5e75cf8d4cd558955a5c7638fba2d68f464/YM/Extra/FREE_DETECTOR_ALGEBRA/Columns.lean\#L103-L115}{\texttt{scanner\_inj}},
\href{https://github.com/yang-mills-group/OpenAI-YM-20261028/blob/3456e5e75cf8d4cd558955a5c7638fba2d68f464/YM/Extra/FREE_DETECTOR_ALGEBRA/GradedBasis.lean\#L41-L49}{\texttt{recorded}},
\href{https://github.com/yang-mills-group/OpenAI-YM-20261028/blob/3456e5e75cf8d4cd558955a5c7638fba2d68f464/YM/Extra/FREE_DETECTOR_ALGEBRA/Vane.lean\#L25-L38}{\texttt{vane}}.

\textbf{Lemma VII.3.n (VII.3.5).} \emph{Proof sketch.}

\begin{itemize}
\tightlist
\item
  \(\omega\) intertwines the label action of \(O(4)\) with substitution,
  via the identity relating rotation of curvature jets to rotation of
  potential jets
  (\href{https://github.com/yang-mills-group/OpenAI-YM-20261028/blob/3456e5e75cf8d4cd558955a5c7638fba2d68f464/YM/Extra/FREE_DETECTOR_ALGEBRA/LabelTurn.lean\#L43-L46}{\texttt{radial\_turn}}).
  It intertwines label conjugation with \(\bar{\ }\).
\item
  The free moments are \(O(4)\)-invariant and real (VII.3.3).
\item
  The \(O(4)\) action on \(E_d\) is transported through
  \(E_d\cong\mathrm{sh}_d\).
\end{itemize}

\emph{Lean:}
\href{https://github.com/yang-mills-group/OpenAI-YM-20261028/blob/3456e5e75cf8d4cd558955a5c7638fba2d68f464/YM/Extra/FREE_DETECTOR_ALGEBRA/LabelTurn.lean\#L80-L91}{\texttt{table\_spin}},
\href{https://github.com/yang-mills-group/OpenAI-YM-20261028/blob/3456e5e75cf8d4cd558955a5c7638fba2d68f464/YM/Extra/FREE_DETECTOR_ALGEBRA/Rank.lean\#L99-L101}{\texttt{ride}}.

\paragraph{4.5 Degree four: rigidity and the position
equations}\label{degree-four-rigidity-and-the-position-equations}

\textbf{Lemma VII.3.o (coordinates on \(E_4\)).} Assume that
\(\mathfrak g\) is simple and nonabelian. Let \(e\) be the unit vector
of coordinate axis \(0\). The basis and dual functionals below are from
VI.2.

\emph{Basis.} Quadratic products below are Wick-ordered. The space
\(E_4\) has the basis \(\{\star_c\}\),
\(c\in\{\varnothing\}\cup\{\text{faces}\}\), where:

\begin{itemize}
\tightlist
\item
  \(\star_\varnothing=1\);
\item
  \(\star_{\{P,Q\}}=\langle F_P,F_Q\rangle\), the invariant contraction
  of two curvature components, for \(P\le Q\) two of the six coordinate
  planes;
\item
  so \(\dim E_4=22\).
\end{itemize}

\emph{Dual functionals.} The dual basis is
\(\lambda_\varnothing(v)=\langle v\rangle_0\) and

\[\lambda_{\{P,Q\}}(v)=\frac{m(P,Q)}{\dim\mathfrak g\;\kappa(P)\kappa(Q)}\,\big\langle v(e)\,\langle F_P,F_Q\rangle(0)\big\rangle_0 ,\]

with \(m(P,Q)=\tfrac12\) if \(P=Q\) and \(1\) otherwise, and
\(\kappa(P)=\kappa_0(P)=\mp\pi^{-2}\).

\emph{Two-point formulas.} For a quadratic index \(c=\{P,Q\}\), put
\(\rho(c)=\dim\mathfrak g\,\kappa(P)\kappa(Q)/m(P,Q)\ne0\). For every
\(v\in E_4\):

\begin{itemize}
\tightlist
\item
  \(\langle v(e)\star_c(0)\rangle_0=\rho(c)\lambda_c(v)\);
\item
  if \(\lambda_\varnothing(v)=0\), then also
  \(\langle\star_c(e)v(0)\rangle_0=\rho(c)\lambda_c(v)\).
\end{itemize}

\emph{Lean:}
\href{https://github.com/yang-mills-group/OpenAI-YM-20261028/blob/3456e5e75cf8d4cd558955a5c7638fba2d68f464/YM/Pert/Rigid/Camera.lean\#L50-L52}{\texttt{price\_tag}},
\href{https://github.com/yang-mills-group/OpenAI-YM-20261028/blob/3456e5e75cf8d4cd558955a5c7638fba2d68f464/YM/Pert/Rigid/Camera.lean\#L59-L69}{\texttt{backwards}},
\href{https://github.com/yang-mills-group/OpenAI-YM-20261028/blob/3456e5e75cf8d4cd558955a5c7638fba2d68f464/YM/Pert/Rigid/Camera.lean\#L39}{\texttt{rate\_ne}},
\href{https://github.com/yang-mills-group/OpenAI-YM-20261028/blob/3456e5e75cf8d4cd558955a5c7638fba2d68f464/YM/Pert/Rigid/Camera.lean\#L30-L36}{\texttt{big\_weight}};
inputs
\href{https://github.com/yang-mills-group/OpenAI-YM-20261028/blob/3456e5e75cf8d4cd558955a5c7638fba2d68f464/YM/UVNext/FREE_STRESS_CERTIFICATE/Size.lean\#L105-L110}{\texttt{stand}},
\href{https://github.com/yang-mills-group/OpenAI-YM-20261028/blob/3456e5e75cf8d4cd558955a5c7638fba2d68f464/YM/UVNext/FREE_STRESS_CERTIFICATE/Size.lean\#L82-L94}{\texttt{separates}}.

\textbf{Lemma VII.3.p (VII.3.6(a)).} Let \(B\in\mathrm{Buf}_4\) satisfy
the one- and two-point conditions of VII.3.6(a). Put
\(\mathrm{coeff}(c,c')=\lambda_c(B\star_{c'})\). Then every coefficient
vanishes, and so \(B=0\). \emph{Proof sketch.}

\begin{itemize}
\tightlist
\item
  \textbf{Unit row and column.}

  \begin{itemize}
  \tightlist
  \item
    \(\mathrm{coeff}(\varnothing,\cdot)=0\) by the one-point condition
    (\href{https://github.com/yang-mills-group/OpenAI-YM-20261028/blob/3456e5e75cf8d4cd558955a5c7638fba2d68f464/YM/Pert/Rigid/Camera.lean\#L73-L81}{\texttt{alone}}).
  \item
    \(\mathrm{coeff}(\cdot,\varnothing)=0\) because \(B1=0\).
  \end{itemize}
\item
  \textbf{Antisymmetry.} The two-point condition with \(v=\star_c\),
  \(w=\star_{c'}\) gives, by VII.3.o,
\end{itemize}

\[\rho(c')\,\mathrm{coeff}(c',c)+\rho(c)\,\mathrm{coeff}(c,c')=0\]

(\href{https://github.com/yang-mills-group/OpenAI-YM-20261028/blob/3456e5e75cf8d4cd558955a5c7638fba2d68f464/YM/Pert/Rigid/Camera.lean\#L82-L89}{\texttt{cancel}}).

\begin{itemize}
\tightlist
\item
  \textbf{Different reflection characters.} Reflection of axis \(a\)
  multiplies \(\star_{\{P,Q\}}\) by \((-1)^{[a\in P]+[a\in Q]}\)
  (\href{https://github.com/yang-mills-group/OpenAI-YM-20261028/blob/3456e5e75cf8d4cd558955a5c7638fba2d68f464/YM/Pert/Rigid/Pallet.lean\#L33-L34}{\texttt{draw\_mask}}).
  If two quadratic indices have different characters of
  \((\mathbb Z/2)^4\), some reflection separates them, and equivariance
  of \(B\) gives \(\mathrm{coeff}=0\)
  (\href{https://github.com/yang-mills-group/OpenAI-YM-20261028/blob/3456e5e75cf8d4cd558955a5c7638fba2d68f464/YM/Pert/Rigid/Spin.lean\#L63-L67}{\texttt{unpaint}}).
\item
  \textbf{Equal reflection characters.} A finite check over the 21 faces
  (\href{https://github.com/yang-mills-group/OpenAI-YM-20261028/blob/3456e5e75cf8d4cd558955a5c7638fba2d68f464/YM/Pert/Rigid/Chess.lean\#L29-L31}{\texttt{havePath}},
  decided by computation) finds an involutive axis permutation
  exchanging \(\star_c\) and \(\star_{c'}\), up to a common sign. It
  also shows the two faces are both diagonal or both off-diagonal.

  \begin{itemize}
  \tightlist
  \item
    Equivariance then gives the symmetry
    \(\mathrm{coeff}(c',c)=\mathrm{coeff}(c,c')\)
    (\href{https://github.com/yang-mills-group/OpenAI-YM-20261028/blob/3456e5e75cf8d4cd558955a5c7638fba2d68f464/YM/Pert/Rigid/Spin.lean\#L69-L84}{\texttt{reflect\_path}}).
  \item
    It also gives equal rates, \(\rho(c)=\rho(c')\)
    (\href{https://github.com/yang-mills-group/OpenAI-YM-20261028/blob/3456e5e75cf8d4cd558955a5c7638fba2d68f464/YM/Pert/Rigid/Camera.lean\#L48-L49}{\texttt{common\_rate}}).
  \item
    Combined with the antisymmetry:
    \(2\rho(c)\,\mathrm{coeff}(c,c')=0\).
  \end{itemize}
\item
  \textbf{Degrees \(d<4\):} \(E_d=\mathbb C\cdot1\), and buffers kill
  the unit.
\end{itemize}

\emph{Lean:}
\href{https://github.com/yang-mills-group/OpenAI-YM-20261028/blob/3456e5e75cf8d4cd558955a5c7638fba2d68f464/YM/Pert/Rigid/Camera.lean\#L90-L105}{\texttt{coeff\_barren}},
\href{https://github.com/yang-mills-group/OpenAI-YM-20261028/blob/3456e5e75cf8d4cd558955a5c7638fba2d68f464/YM/Pert/Rigid/Camera.lean\#L107-L113}{\texttt{frozen\_four}},
\href{https://github.com/yang-mills-group/OpenAI-YM-20261028/blob/3456e5e75cf8d4cd558955a5c7638fba2d68f464/YM/Pert/Rigid/Camera.lean\#L114-L117}{\texttt{first\_frozen}},
\href{https://github.com/yang-mills-group/OpenAI-YM-20261028/blob/3456e5e75cf8d4cd558955a5c7638fba2d68f464/YM/Pert/Rigid/Spin.lean\#L60-L62}{\texttt{orient\_transfer}},
\href{https://github.com/yang-mills-group/OpenAI-YM-20261028/blob/3456e5e75cf8d4cd558955a5c7638fba2d68f464/YM/Pert/Rigid/Pallet.lean\#L74-L80}{\texttt{wearer}}.

\textbf{Lemma VII.3.q (VII.3.6(b)).} \emph{Proof sketch.}

\begin{itemize}
\tightlist
\item
  The \emph{feedback} of a section \((r,p,f)\) is the \(B\)-derivative
  at \(0\) of \(\mathrm{Re}\int f\langle(1+B)p_1\cdots\rangle_0\)
  (\href{https://github.com/yang-mills-group/OpenAI-YM-20261028/blob/3456e5e75cf8d4cd558955a5c7638fba2d68f464/YM/Pert/Position/Detect.lean\#L58}{\texttt{feedback}}).
  It is a linear functional on \(\mathrm{Buf}_d\).
\item
  If all feedbacks vanish at \(B\), then \(B\) is stationary
  (\href{https://github.com/yang-mills-group/OpenAI-YM-20261028/blob/3456e5e75cf8d4cd558955a5c7638fba2d68f464/YM/Pert/Position/Detect.lean\#L103-L112}{\texttt{detected}}):

  \begin{itemize}
  \tightlist
  \item
    the profiles \(f\) and \(if\) give the real and imaginary parts
    (\href{https://github.com/yang-mills-group/OpenAI-YM-20261028/blob/3456e5e75cf8d4cd558955a5c7638fba2d68f464/YM/Pert/Position/Detect.lean\#L97-L100}{\texttt{feedback\_twist}});
  \item
    the integrand is continuous on \(\mathrm{Conf}_r\)
    (\href{https://github.com/yang-mills-group/OpenAI-YM-20261028/blob/3456e5e75cf8d4cd558955a5c7638fba2d68f464/YM/Pert/Position/Detect.lean\#L19-L20}{\texttt{visible\_cont}}),
    so it vanishes pointwise.
  \end{itemize}
\item
  The span of all feedbacks lies in the finite-dimensional dual of
  \(\mathrm{Buf}_d\). Pick \(k\) sections whose feedbacks span the same
  subspace; their joint kernel is \(0\)
  (\href{https://github.com/yang-mills-group/OpenAI-YM-20261028/blob/3456e5e75cf8d4cd558955a5c7638fba2d68f464/YM/Pert/Position/Detect.lean\#L115-L129}{\texttt{finite\_detectors}}).
\item
  Take a left inverse \(g'\) of the resulting injection
  \(\mathrm{Buf}_d\to\mathbb R^k\), and put \(g=H\circ g'\) with
  \(H:\mathrm{Buf}_d\cong\mathbb R^m\) a coordinate isomorphism.
\item
  The chart derivative is then \(H\)
  (\href{https://github.com/yang-mills-group/OpenAI-YM-20261028/blob/3456e5e75cf8d4cd558955a5c7638fba2d68f464/YM/Pert/Position/Detect.lean\#L82-L93}{\texttt{chart\_deriv\_formula}}).
\end{itemize}

\emph{Lean:}
\href{https://github.com/yang-mills-group/OpenAI-YM-20261028/blob/3456e5e75cf8d4cd558955a5c7638fba2d68f464/YM/Pert/Position/Detect.lean\#L131-L147}{\texttt{square\_of\_frozen}}.

\textbf{Lemma VII.3.r (VII.3.6(c)).}

\emph{The data.}

\begin{itemize}
\tightlist
\item
  \(m_Q(k)\) is the tested correlation of the limit theory at scale
  \(r_k\) in the frames \(U_k\), on basis fields
  (\href{https://github.com/yang-mills-group/OpenAI-YM-20261028/blob/3456e5e75cf8d4cd558955a5c7638fba2d68f464/YM/Pert/Position/Actual.lean\#L15}{\texttt{measured}},
  \href{https://github.com/yang-mills-group/OpenAI-YM-20261028/blob/3456e5e75cf8d4cd558955a5c7638fba2d68f464/YM/Pert/Position/RawMean.lean\#L19}{\texttt{exactMean}}).
\item
  The leading free comparison, with derivative order \(0\) and no stress
  slots, gives \(m_Q(k)\to w_Q\), the free tested readings
  (\href{https://github.com/yang-mills-group/OpenAI-YM-20261028/blob/3456e5e75cf8d4cd558955a5c7638fba2d68f464/YM/Pert/Position/Actual.lean\#L16-L19}{\texttt{measured\_tends}},
  \href{https://github.com/yang-mills-group/OpenAI-YM-20261028/blob/3456e5e75cf8d4cd558955a5c7638fba2d68f464/YM/Pert/Position/RawMean.lean\#L42-L53}{\texttt{exact\_tape}}).
\end{itemize}

\emph{The hypothesis on \(\varphi\).} It consists of four conditions:

\begin{itemize}
\tightlist
\item
  \(\varphi\) is a real-analytic implicit solution of
  \(\mathrm{chart}_Q(w,\varphi(w))=\mathrm{chart}_Q(w_Q,0)\) near
  \(w_Q\), with \(\varphi(w_Q)=0\). Locally the solution set is the
  graph of \(\varphi\)
  (\href{https://github.com/yang-mills-group/OpenAI-YM-20261028/blob/3456e5e75cf8d4cd558955a5c7638fba2d68f464/YM/Pert/Position/Implicit.lean\#L29-L33}{\texttt{Reports}}).
\item
  For every integer \(L\ge2\), \(\pi^{(0)}_Q(L,m)\to w_Q\) as
  \(m\to\infty\).
\item
  For every \(L\ge2\) and \(p,j\ge0\), eventually in \(m\), all
  order-\(j\) coupling derivatives at \(0\) of the position equations
  evaluated on the order-\(p\) lattice truncation and its image under
  \(\varphi\) equal those of the constant free value. The threshold may
  depend on \(L,p,j\).
\item
  Some buffers \(A_k\to0\) have \(1+A_k\) invertible and solve the
  position equations eventually along \(X\). For every clean degree-four
  system \(R\) and \(L\ge2\), their corrected tested readings differ by
  \(o(1)\) from
  \(\mathrm{chart}_R(\pi_R^{(0)}(L,k),\varphi(\pi_Q^{(0)}(L,k)))\).
\end{itemize}

\emph{Proof sketch.}

\begin{itemize}
\tightlist
\item
  \(B_k=\varphi(m_Q(k))\to\varphi(w_Q)=0\).
\item
  By local uniqueness of the implicit solution, \(B_k=A_k\) eventually.
  This transfers exact solvability and the leading comparison from
  \(A_k\) to \(B_k\).
\item
  The required symmetries and unit identity follow from the eventual
  compatibility of \(U_k\) in VI.2 and the fact that buffers commute
  with conjugation and signed permutations and kill the unit
  (\href{https://github.com/yang-mills-group/OpenAI-YM-20261028/blob/3456e5e75cf8d4cd558955a5c7638fba2d68f464/YM/Pert/Position/Frame.lean\#L39-L46}{\texttt{adjusted\_own}},
  \href{https://github.com/yang-mills-group/OpenAI-YM-20261028/blob/3456e5e75cf8d4cd558955a5c7638fba2d68f464/YM/Pert/Position/Frame.lean\#L47-L49}{\texttt{adjusted\_unit}}).
\item
  Compatibility with degrees \(\le3\): \(1+B\) is the identity on
  \(E_d\subset E_4\) for \(d<4\), because \(E_d=\mathbb C\cdot1\)
  (\href{https://github.com/yang-mills-group/OpenAI-YM-20261028/blob/3456e5e75cf8d4cd558955a5c7638fba2d68f464/YM/Pert/Position/Frame.lean\#L55-L61}{\texttt{motor\_small}}).
\end{itemize}

The full degree-four position conclusion consists of:

\begin{itemize}
\tightlist
\item
  for every degree, every clean system and every \(L\ge2\), the
  uncorrected tested data minus the order-zero lattice data at spacing
  \(L^{-k}\) tend to \(0\); every clean square system admits a map
  \(\varphi\) satisfying the four conditions above;
\item
  a clean square degree-four system \(Q\) and such a map \(\varphi\) for
  which VII.3.6(c) holds.
\end{itemize}

\emph{Lean:}
\href{https://github.com/yang-mills-group/OpenAI-YM-20261028/blob/3456e5e75cf8d4cd558955a5c7638fba2d68f464/YM/Pert/Position/Four.lean\#L32-L43}{\texttt{residency}},
\href{https://github.com/yang-mills-group/OpenAI-YM-20261028/blob/3456e5e75cf8d4cd558955a5c7638fba2d68f464/YM/Pert/Position/Frame.lean\#L29-L30}{\texttt{adjusted}},
\href{https://github.com/yang-mills-group/OpenAI-YM-20261028/blob/3456e5e75cf8d4cd558955a5c7638fba2d68f464/YM/Pert/Position/Actual.lean\#L42}{\texttt{testedMean}},
\href{https://github.com/yang-mills-group/OpenAI-YM-20261028/blob/3456e5e75cf8d4cd558955a5c7638fba2d68f464/YM/Pert/Position/Board.lean\#L98}{\texttt{finite}},
\href{https://github.com/yang-mills-group/OpenAI-YM-20261028/blob/3456e5e75cf8d4cd558955a5c7638fba2d68f464/YM/Pert/Position/Implicit.lean\#L22}{\texttt{readJet}};
packaged by
\href{https://github.com/yang-mills-group/OpenAI-YM-20261028/blob/3456e5e75cf8d4cd558955a5c7638fba2d68f464/YM/Pert/Position/Four.lean\#L49-L53}{\texttt{four\_on}}
(VIII.1).

\paragraph{4.6 Exclusion of Wick profiles
(VII.3.7)}\label{exclusion-of-wick-profiles-vii.3.7}

\textbf{Lemma VII.3.s (expansions; VII.2).} Let \(\mathcal B\) be the
real span of \(r^a(\log r)^k\) for \(a\in\mathbb Z\), \(k\in\mathbb N\),
on \(r>0\). Every integrated mode kernel, and each of the real and
imaginary parts of every Wick profile, has an expansion to all orders in
\(\mathcal B\): for every \(N\in\mathbb Z\) there is \(u\in\mathcal B\)
with \(f-u=O(r^N)\) as \(r\to0^+\). \emph{Proof sketch.}

\begin{itemize}
\tightlist
\item
  \textbf{Reduction to radial moments.} In polar coordinates,
  \(K_{\mathsf m,q,h}=\tfrac12\big(\int_{S^2}h\big)M_{j,q}\), where
\end{itemize}

\[M_{j,k}(s)=\int_0^\infty t^{j+2}E^{k-1}e^{-sE}dt,\qquad E=\sqrt{\mathsf m^2+t^2}.\]

\begin{itemize}
\tightlist
\item
  \textbf{Relations among the moments.}

  \begin{itemize}
  \tightlist
  \item
    \(M_{j,k}'=-M_{j,k+1}\);
  \item
    \(M_{j,k+2}=\mathsf m^2M_{j,k}+M_{j+2,k}\);
  \item
    the Euler identity \(sM_{j+2,0}=(j+3)M_{j,1}\), from integration by
    parts.
  \item
    Together they give \((s^{j+3}M_{j,1})'=-\mathsf m^2s^{j+3}M_{j,0}\).
  \end{itemize}
\item
  \textbf{Bootstrap.}

  \begin{itemize}
  \tightlist
  \item
    Start from \(M_{j,0}=O(s^{-j-2})\).
  \item
    If \(M_{j,0}\) has an expansion with remainder \(O(s^k)\),
    \(k\ge-j-2\), integrating the weighted equation gives one for
    \(M_{j,1}\) with remainder \(O(s^{k+1})\).
  \item
    Integrating \(M_{j,0}'=-M_{j,1}\) then improves the known remainder
    for \(M_{j,0}\) by at least one power. When the remainder to be
    integrated has order \(s^{-1}\), first weaken it to \(O(s^{-2})\);
    this still suffices for the induction.
  \item
    The explicit power-log terms have primitives in \(\mathcal B\);
    integrating \(s^{-1}(\log s)^k\) raises the logarithmic power to
    \(k+1\).
  \end{itemize}
\item
  Products and sums preserve the class.
\end{itemize}

\emph{Lean:}
\href{https://github.com/yang-mills-group/OpenAI-YM-20261028/blob/3456e5e75cf8d4cd558955a5c7638fba2d68f464/YM/TraceExtension/Wick/Expansion.lean\#L47}{\texttt{Expands}},
\href{https://github.com/yang-mills-group/OpenAI-YM-20261028/blob/3456e5e75cf8d4cd558955a5c7638fba2d68f464/YM/TraceExtension/Wick/Massive.lean\#L69-L88}{\texttt{euler}},
\href{https://github.com/yang-mills-group/OpenAI-YM-20261028/blob/3456e5e75cf8d4cd558955a5c7638fba2d68f464/YM/TraceExtension/Wick/Massive.lean\#L99-L105}{\texttt{derived}},
\href{https://github.com/yang-mills-group/OpenAI-YM-20261028/blob/3456e5e75cf8d4cd558955a5c7638fba2d68f464/YM/TraceExtension/Wick/Massive.lean\#L106-L121}{\texttt{basic\_expands}},
\href{https://github.com/yang-mills-group/OpenAI-YM-20261028/blob/3456e5e75cf8d4cd558955a5c7638fba2d68f464/YM/TraceExtension/Wick/Massive.lean\#L122-L131}{\texttt{moment\_expands}},
\href{https://github.com/yang-mills-group/OpenAI-YM-20261028/blob/3456e5e75cf8d4cd558955a5c7638fba2d68f464/YM/TraceExtension/Wick/PolynomialField.lean\#L116-L125}{\texttt{Track.moments\_expand}},
\href{https://github.com/yang-mills-group/OpenAI-YM-20261028/blob/3456e5e75cf8d4cd558955a5c7638fba2d68f464/YM/TraceExtension/Wick/Refute.lean\#L18-L22}{\texttt{real\_profile}},
\href{https://github.com/yang-mills-group/OpenAI-YM-20261028/blob/3456e5e75cf8d4cd558955a5c7638fba2d68f464/YM/TraceExtension/Wick/Gain.lean\#L109-L113}{\texttt{primitives}}.

\textbf{Lemma VII.3.t (no power-log law; VII.2).} Let
\(f:(0,\infty)\to\mathbb R\) and \(A\in\mathbb R\). Suppose \(f=u+O(1)\)
with \(u\in\mathcal B\) as \(s\to0^+\), and \(s^8\log^2(1/s)f(s)\to A\).
Then \(A=0\). \emph{Proof sketch.}

\begin{itemize}
\tightlist
\item
  The \(O(1)\) part contributes \(o(1)\).
\item
  Collect equal power-log terms and discard zero coefficients. Along
  \(s=2^{-n}\), \(s^8\log^2(1/s)u(s)=\sum_tc_t\rho_t^n n^{k_t+2}\), with
  \(c_t\ne0\), \(\rho_t=2^{-(8+a_t)}\) and the pairs \((\rho_t,k_t)\)
  distinct. The empty sum has limit \(0\).
\item
  Let \(t^\ast\) be the lexicographically largest term.

  \begin{itemize}
  \tightlist
  \item
    If \(\rho_{t^\ast}<1\), every term tends to \(0\), so \(A=0\).
  \item
    Otherwise its positive growth factor
    \(\rho_{t^\ast}^n n^{k_{t^\ast}+2}\) diverges and dominates those of
    all remaining terms. Dividing by this factor, the left side tends to
    \(0\) while the right side tends to \(c_{t^\ast}\ne0\), a
    contradiction.
  \end{itemize}
\end{itemize}

\emph{Lean:}
\href{https://github.com/yang-mills-group/OpenAI-YM-20261028/blob/3456e5e75cf8d4cd558955a5c7638fba2d68f464/YM/TraceExtension/Wick/Log.lean\#L96-L108}{\texttt{no\_powerlog}}.

\textbf{Lemma VII.3.u (VII.3.7).} \emph{Proof sketch.}

\begin{itemize}
\tightlist
\item
  \begin{enumerate}
  \def\labelenumi{(\alph{enumi})}
  \tightlist
  \item
    If \(G_\Theta=W\) on \((0,\delta)\), apply VII.3.t to
    \(\mathrm{Re}\,W\) and the real part of the hypothesis. This gives
    \(A=0\), a contradiction.
  \end{enumerate}
\item
  (b) Write \(G_\Theta(s)=\int e^{-sp^0}d\mu(p)\), using the spectral
  measure \(\mu\) of \(\Theta\).

  \begin{itemize}
  \tightlist
  \item
    Every compactly supported momentum test \(\psi\) has the form
    \(\overline{\widehat f}\,\widehat g\): choose \(|\widehat f|=1\) on
    its support and \(\widehat g=\widehat f\,\psi\). Thus the
    hypothesized product-test identities give equality on all compactly
    supported momentum tests.
  \item
    Approximate \(e^{-sp^0}\) by compactly supported Schwartz tests
    dominated by it
    (\href{https://github.com/yang-mills-group/OpenAI-YM-20261028/blob/3456e5e75cf8d4cd558955a5c7638fba2d68f464/YM/TraceExtension/Wick/Refute.lean\#L24-L38}{\texttt{approximate}}).
  \item
    Exponential damping is integrable for the trace spectral measure and
    for every mode product. Dominated convergence on both sides
    identifies \(G_\Theta\) with the profile of the description on
    \((0,\infty)\)
    (\href{https://github.com/yang-mills-group/OpenAI-YM-20261028/blob/3456e5e75cf8d4cd558955a5c7638fba2d68f464/YM/TraceExtension/Wick/Refute.lean\#L40-L54}{\texttt{reads}},
    \href{https://github.com/yang-mills-group/OpenAI-YM-20261028/blob/3456e5e75cf8d4cd558955a5c7638fba2d68f464/YM/TraceExtension/Wick/Refute.lean\#L62-L73}{\texttt{identifies\_euclid}}).
  \item
    Then apply (a).
  \end{itemize}
\item
  The public form uses the same mode kernels and identifies the
  asymptotic constant with \(3\dim\mathfrak g/\pi^4\).
\end{itemize}

\emph{Lean:}
\href{https://github.com/yang-mills-group/OpenAI-YM-20261028/blob/3456e5e75cf8d4cd558955a5c7638fba2d68f464/YM/TraceExtension/Wick/Refute.lean\#L75-L80}{\texttt{euclidean\_excludes}},
\href{https://github.com/yang-mills-group/OpenAI-YM-20261028/blob/3456e5e75cf8d4cd558955a5c7638fba2d68f464/YM/TraceExtension/Wick/Refute.lean\#L86-L105}{\texttt{polynomial\_excludes}},
\href{https://github.com/yang-mills-group/OpenAI-YM-20261028/blob/3456e5e75cf8d4cd558955a5c7638fba2d68f464/TraceBridge.lean\#L35-L44}{\texttt{publicTrace}}.

\paragraph{4.7 The formal lattice perturbation
theory}\label{the-formal-lattice-perturbation-theory}

\textbf{Definitions VII.3.v.} Fix a finite periodic four-dimensional
lattice. The perturbation theory uses inversion-odd link potentials,

\[A(-x-e_\mu,\mu)=-A(x,\mu).\]

Its definitions are:

\textbf{Gauge-fixed linearization and Gaussian measure.}

The links are \(U=\exp(A)\). The action \(S\) is the finite-regulator
Wilson term plus the quadratic Feynman gauge-fixing term
\(\tfrac14\sum_x|\mathrm{div}\,A(x)|^2\)
(\href{https://github.com/yang-mills-group/OpenAI-YM-20261028/blob/3456e5e75cf8d4cd558955a5c7638fba2d68f464/YM/Pert/Lattice/Action.lean\#L46}{\texttt{action}}).

The Haar Jacobian \(J\) and the normalized Faddeev--Popov determinant
\(\mathrm{FP}\) are smooth at \(0\), with \(J(0)=\mathrm{FP}(0)=1\). The
latter is the determinant of the gauge-condition derivative on
inversion-even site parameters fixed to zero at the root, divided by its
flat value
(\href{https://github.com/yang-mills-group/OpenAI-YM-20261028/blob/3456e5e75cf8d4cd558955a5c7638fba2d68f464/YM/Pert/Lattice/GaugeFix.lean\#L94}{\texttt{Orb.FP}}).
The matrix determinant, logarithmic determinant and inverse estimates
come from the matrix calculus listed in the header.

The centered Gaussian probability has covariance equal to the inverse
Hessian of this gauge-fixed action at \(0\)
(\href{https://github.com/yang-mills-group/OpenAI-YM-20261028/blob/3456e5e75cf8d4cd558955a5c7638fba2d68f464/YM/Pert/Lattice/Measure.lean\#L31}{\texttt{Orb.data}}).

\textbf{Formal series in the coupling.}

Write \(S_j(A)=(j!)^{-1}D^jS(0)[A,\dots,A]\). After the rescaling
\(A\mapsto gA\), the formal weight is

\[w(g,A)=\exp\!\left(-\sum_{i\ge1}g^iS_{i+2}(A)\right)(J\,\mathrm{FP})(gA),\]

where the last factor denotes its formal Taylor series at \(g=0\).

\(\mathbb E_0\) is the Gaussian expectation, taken coefficientwise.

The normalized expectation is
\(\langle f\rangle_{\rm pert}=\mathbb E_0(wf)/\mathbb E_0(w)\)
(\href{https://github.com/yang-mills-group/OpenAI-YM-20261028/blob/3456e5e75cf8d4cd558955a5c7638fba2d68f464/YM/Pert/Lattice/Series.lean\#L61}{\texttt{Orb.weight}},
\href{https://github.com/yang-mills-group/OpenAI-YM-20261028/blob/3456e5e75cf8d4cd558955a5c7638fba2d68f464/YM/Pert/Lattice/Series.lean\#L74}{\texttt{Orb.expect}},
\href{https://github.com/yang-mills-group/OpenAI-YM-20261028/blob/3456e5e75cf8d4cd558955a5c7638fba2d68f464/YM/Pert/Lattice/Series.lean\#L85}{\texttt{Orb.normalize}}).

Every coefficient is continuous with at most polynomial growth, hence
Gaussian-integrable. Since \(w(0,A)=1\), the denominator has constant
coefficient \(1\) and is invertible as a formal series.

\textbf{Lattice readers of labels.} Substitute the lattice jet series
plus \(iZ\) into each polynomial label, where \(Z\) is an independent
centered real Gaussian jet with the prescribed covariance, and average
in \(Z\) separately within each insertion. Multiply these averaged
insertions and take the normalized formal expectation. The separate
auxiliary averages implement Wick subtraction without auxiliary
contractions between different insertions
(\href{https://github.com/yang-mills-group/OpenAI-YM-20261028/blob/3456e5e75cf8d4cd558955a5c7638fba2d68f464/YM/Pert/Lattice/Blank.lean\#L41}{\texttt{Full.column}},
\href{https://github.com/yang-mills-group/OpenAI-YM-20261028/blob/3456e5e75cf8d4cd558955a5c7638fba2d68f464/YM/Pert/Lattice/Blank.lean\#L87}{\texttt{Full.amplitude}}).

\textbf{Profiles.} The profile is the lattice Riemann sum

\[\textstyle\sum_o h^{4r}f(\iota_h(o))\,\mathrm{amplitude}(o),\qquad h=L^{-m},\ \text{side }4L^{m^2+m+1},\]

where \(o\) ranges over \(r\) lattice sites and \(\iota_h\) takes the
chosen integer lifts relative to the quarter-side origin and multiplies
them by \(h\)
(\href{https://github.com/yang-mills-group/OpenAI-YM-20261028/blob/3456e5e75cf8d4cd558955a5c7638fba2d68f464/YM/Pert/Lattice/Bundles.lean\#L55-L56}{\texttt{Full.profile}}).

The order-zero profiles enter VII.3.r through the convergence theorem of
VII.1. Their higher coefficients enter the formal position identities
and the degree-four parity assertion assembled in VIII.1.

\href{https://github.com/yang-mills-group/OpenAI-YM-20261028/blob/3456e5e75cf8d4cd558955a5c7638fba2d68f464/YM/Pert/Lattice/Together.lean\#L58-L61}{\texttt{Orb.finite\_agreement}}
shows that retaining the interaction series, density series and
insertions only through coupling order \(k\) preserves every normalized
coefficient through order \(k\). The required action Taylor coefficients
extend through order \(k+2\).

\subsubsection{5. Where the difficulty
lies}\label{where-the-difficulty-lies-27}

The main calculation is the simultaneous diagonalization of all jet
orders (VII.3.d--h). The Taylor normalization and the sum over ordered
form indices determine the Kelvin constants and positive pivots;
changing either convention changes the formulas. Recentring at the
inverted separation eliminates all lower Taylor coefficients, so
positivity follows from a triangular change of basis.

At degree four, the two-point form is indefinite before reflection.
Rigidity uses its nonzero diagonal rates together with
signed-permutation symmetry: equal reflection characters force symmetric
coefficients and equal rates, while the two-point condition forces
weighted antisymmetry. Full rotational equivariance of the buffer is
unnecessary.

The position comparison is leading-order only. Its formal identities at
every coupling order follow from the local implicit equation at a finite
regulator; they give no all-order comparison with the interacting
theory. The free quotient and its injective columns likewise give no
interacting rank bound without the comparison estimates of III.11.

The non-freeness argument is independent of Gram positivity. Its
analytic input is the power-log expansion of the massive Laplace
moments, including the remainder estimates needed for integration. It
excludes inverse-logarithmic leading behavior; logarithms with
nonnegative powers already occur in free massive kernels.

\subsubsection{6. Notation}\label{notation-27}

{\def\LTcaptype{none} % do not increment counter
\begin{longtable}[]{@{}
  >{\raggedright\arraybackslash}p{(\linewidth - 4\tabcolsep) * \real{0.3333}}
  >{\raggedright\arraybackslash}p{(\linewidth - 4\tabcolsep) * \real{0.3333}}
  >{\raggedright\arraybackslash}p{(\linewidth - 4\tabcolsep) * \real{0.3333}}@{}}
\toprule\noalign{}
\begin{minipage}[b]{\linewidth}\raggedright
object
\end{minipage} & \begin{minipage}[b]{\linewidth}\raggedright
code name
\end{minipage} & \begin{minipage}[b]{\linewidth}\raggedright
notation
\end{minipage} \\
\midrule\noalign{}
\endhead
\bottomrule\noalign{}
\endlastfoot
on-shell \(j\)-jets (homogeneous Maxwell solutions) &
\href{https://github.com/yang-mills-group/OpenAI-YM-20261028/blob/3456e5e75cf8d4cd558955a5c7638fba2d68f464/YM/Extra/FREE_DETECTOR_ALGEBRA/Solutions.lean\#L115}{\texttt{Sol}}
& \(\mathrm{Sol}_j\) \\
towers of multipoles &
\href{https://github.com/yang-mills-group/OpenAI-YM-20261028/blob/3456e5e75cf8d4cd558955a5c7638fba2d68f464/YM/Extra/FREE_DETECTOR_ALGEBRA/PositiveLeg.lean\#L69}{\texttt{ProbeTower}}
& \(\mathrm{Tow}_n\) \\
Wick polynomials &
\href{https://github.com/yang-mills-group/OpenAI-YM-20261028/blob/3456e5e75cf8d4cd558955a5c7638fba2d68f464/YM/Maxwell/Equation.lean\#L24}{\texttt{WickPol}},
\href{https://github.com/yang-mills-group/OpenAI-YM-20261028/blob/3456e5e75cf8d4cd558955a5c7638fba2d68f464/YM/Maxwell/Equation.lean\#L22}{\texttt{Slot}}
& \(\mathcal W=\mathbb C[X_s]\) \\
linear polynomial of a tower &
\href{https://github.com/yang-mills-group/OpenAI-YM-20261028/blob/3456e5e75cf8d4cd558955a5c7638fba2d68f464/YM/Extra/FREE_DETECTOR_ALGEBRA/Script.lean\#L16-L23}{\texttt{code}},
\href{https://github.com/yang-mills-group/OpenAI-YM-20261028/blob/3456e5e75cf8d4cd558955a5c7638fba2d68f464/YM/Extra/FREE_DETECTOR_ALGEBRA/Symbol.lean\#L77}{\texttt{digit}}
& \(\ell_p\) \\
encoding in a basis of a tower &
\href{https://github.com/yang-mills-group/OpenAI-YM-20261028/blob/3456e5e75cf8d4cd558955a5c7638fba2d68f464/YM/Extra/FREE_DETECTOR_ALGEBRA/Script.lean\#L25}{\texttt{encode}}
& \(\Phi_b\) \\
free moment &
\href{https://github.com/yang-mills-group/OpenAI-YM-20261028/blob/3456e5e75cf8d4cd558955a5c7638fba2d68f464/YM/Maxwell/Wick.lean\#L126-L127}{\texttt{Wick.moment}},
\href{https://github.com/yang-mills-group/OpenAI-YM-20261028/blob/3456e5e75cf8d4cd558955a5c7638fba2d68f464/YM/Extra/FREE_DETECTOR_ALGEBRA/Duplex.lean\#L28}{\texttt{join}}
& \(\langle P_1(x_1)\cdots\rangle_0\) \\
reflection in \(z^\perp\) &
\href{https://github.com/yang-mills-group/OpenAI-YM-20261028/blob/3456e5e75cf8d4cd558955a5c7638fba2d68f464/YM/Extra/FREE_DETECTOR_ALGEBRA/Forms.lean\#L103}{\texttt{face}},
\href{https://github.com/yang-mills-group/OpenAI-YM-20261028/blob/3456e5e75cf8d4cd558955a5c7638fba2d68f464/YM/Extra/FREE_DETECTOR_ALGEBRA/Symbol.lean\#L54-L55}{\texttt{pullSol}},
\href{https://github.com/yang-mills-group/OpenAI-YM-20261028/blob/3456e5e75cf8d4cd558955a5c7638fba2d68f464/YM/Extra/FREE_DETECTOR_ALGEBRA/Covariance.lean\#L159-L163}{\texttt{solO}}
& \(\mathcal R_z\) \\
reflected pairing &
\href{https://github.com/yang-mills-group/OpenAI-YM-20261028/blob/3456e5e75cf8d4cd558955a5c7638fba2d68f464/YM/Extra/FREE_DETECTOR_ALGEBRA/Duplex.lean\#L54}{\texttt{duplex}},
\href{https://github.com/yang-mills-group/OpenAI-YM-20261028/blob/3456e5e75cf8d4cd558955a5c7638fba2d68f464/YM/Extra/FREE_DETECTOR_ALGEBRA/Duplex.lean\#L49-L52}{\texttt{aligned}}
& \(\langle P,Q\rangle_{x,y}\) \\
tower Gram form &
\href{https://github.com/yang-mills-group/OpenAI-YM-20261028/blob/3456e5e75cf8d4cd558955a5c7638fba2d68f464/YM/Extra/FREE_DETECTOR_ALGEBRA/PositiveLeg.lean\#L71-L77}{\texttt{towerB}},
\href{https://github.com/yang-mills-group/OpenAI-YM-20261028/blob/3456e5e75cf8d4cd558955a5c7638fba2d68f464/YM/Extra/FREE_DETECTOR_ALGEBRA/PositiveLeg.lean\#L29-L32}{\texttt{scan}}
& \(B_z\) \\
inversion; inverted point &
\href{https://github.com/yang-mills-group/OpenAI-YM-20261028/blob/3456e5e75cf8d4cd558955a5c7638fba2d68f464/YM/Extra/FREE_DETECTOR_ALGEBRA/LegAct.lean\#L110}{\texttt{inversePoint}},
\href{https://github.com/yang-mills-group/OpenAI-YM-20261028/blob/3456e5e75cf8d4cd558955a5c7638fba2d68f464/YM/Extra/FREE_DETECTOR_ALGEBRA/Harmonic.lean\#L25}{\texttt{rr}}
& \(I(z)=\check z=z/\lvert z\rvert^{2}\) \\
field of a multipole &
\href{https://github.com/yang-mills-group/OpenAI-YM-20261028/blob/3456e5e75cf8d4cd558955a5c7638fba2d68f464/YM/Extra/FREE_DETECTOR_ALGEBRA/CurlDetect.lean\#L24}{\texttt{reply}}
& \(R_P=s_jI^{\ast}P\) \\
Kelvin constant; pivot &
\href{https://github.com/yang-mills-group/OpenAI-YM-20261028/blob/3456e5e75cf8d4cd558955a5c7638fba2d68f464/YM/Extra/FREE_DETECTOR_ALGEBRA/Kelvin.lean\#L46}{\texttt{strength}},
\href{https://github.com/yang-mills-group/OpenAI-YM-20261028/blob/3456e5e75cf8d4cd558955a5c7638fba2d68f464/YM/Extra/FREE_DETECTOR_ALGEBRA/PositiveLeg.lean\#L24}{\texttt{price}}
& \(s_j\), \(\pi_j(z)\) \\
shifted tower &
\href{https://github.com/yang-mills-group/OpenAI-YM-20261028/blob/3456e5e75cf8d4cd558955a5c7638fba2d68f464/YM/Extra/FREE_DETECTOR_ALGEBRA/Shift.lean\#L20-L21}{\texttt{shifts}},
\href{https://github.com/yang-mills-group/OpenAI-YM-20261028/blob/3456e5e75cf8d4cd558955a5c7638fba2d68f464/YM/Extra/FREE_DETECTOR_ALGEBRA/PositiveLeg.lean\#L98-L99}{\texttt{newLeg}},
\href{https://github.com/yang-mills-group/OpenAI-YM-20261028/blob/3456e5e75cf8d4cd558955a5c7638fba2d68f464/YM/Extra/FREE_DETECTOR_ALGEBRA/PositiveLeg.lean\#L123-L124}{\texttt{shiftAll}}
& \(\sigma_z\) \\
two-point annihilator &
\href{https://github.com/yang-mills-group/OpenAI-YM-20261028/blob/3456e5e75cf8d4cd558955a5c7638fba2d68f464/YM/Extra/FREE_DETECTOR_ALGEBRA/Script.lean\#L81}{\texttt{crash}},
\href{https://github.com/yang-mills-group/OpenAI-YM-20261028/blob/3456e5e75cf8d4cd558955a5c7638fba2d68f464/YM/Extra/FREE_DETECTOR_ALGEBRA/Removal.lean\#L25-L26}{\texttt{extractVar}}
& \(D^{x,y}_p\), \(\partial^{y,y'}_s\) \\
Fock weight &
\href{https://github.com/yang-mills-group/OpenAI-YM-20261028/blob/3456e5e75cf8d4cd558955a5c7638fba2d68f464/YM/Extra/FREE_DETECTOR_ALGEBRA/PowerGram.lean\#L41}{\texttt{tally}}
& \(\alpha!\,d^\alpha\) \\
free-field map of labels; shell sector &
\href{https://github.com/yang-mills-group/OpenAI-YM-20261028/blob/3456e5e75cf8d4cd558955a5c7638fba2d68f464/YM/Maxwell/Equation.lean\#L66-L67}{\texttt{observable}},
\href{https://github.com/yang-mills-group/OpenAI-YM-20261028/blob/3456e5e75cf8d4cd558955a5c7638fba2d68f464/YM/Euclid/Comparator.lean\#L64}{\texttt{shellSector}}
& \(\omega_d\), \(\mathrm{sh}_d\) \\
field space of the free table; embedding &
\href{https://github.com/yang-mills-group/OpenAI-YM-20261028/blob/3456e5e75cf8d4cd558955a5c7638fba2d68f464/YM/Euclid/Null.lean\#L74}{\texttt{(freeTable\ F).Field\ d}},
\href{https://github.com/yang-mills-group/OpenAI-YM-20261028/blob/3456e5e75cf8d4cd558955a5c7638fba2d68f464/YM/Pert/Position/Board.lean\#L11}{\texttt{Position.E}},
\href{https://github.com/yang-mills-group/OpenAI-YM-20261028/blob/3456e5e75cf8d4cd558955a5c7638fba2d68f464/YM/Extra/FREE_DETECTOR_ALGEBRA/Rank.lean\#L70-L71}{\texttt{writes}}
& \(E_d\), \(\varpi_d\) \\
buffers; corrected frame &
\href{https://github.com/yang-mills-group/OpenAI-YM-20261028/blob/3456e5e75cf8d4cd558955a5c7638fba2d68f464/YM/Pert/Position/Board.lean\#L31-L40}{\texttt{buffers}},
\href{https://github.com/yang-mills-group/OpenAI-YM-20261028/blob/3456e5e75cf8d4cd558955a5c7638fba2d68f464/YM/Pert/Position/Board.lean\#L45}{\texttt{motor}},
\href{https://github.com/yang-mills-group/OpenAI-YM-20261028/blob/3456e5e75cf8d4cd558955a5c7638fba2d68f464/YM/Pert/Position/Frame.lean\#L29-L30}{\texttt{adjusted}}
& \(\mathrm{Buf}_d\), \(U_k(1+B)\) \\
stationary buffers vanish &
\href{https://github.com/yang-mills-group/OpenAI-YM-20261028/blob/3456e5e75cf8d4cd558955a5c7638fba2d68f464/YM/Pert/Position/Detect.lean\#L102}{\texttt{Frozen}},
\href{https://github.com/yang-mills-group/OpenAI-YM-20261028/blob/3456e5e75cf8d4cd558955a5c7638fba2d68f464/YM/Pert/Rigid/Camera.lean\#L71-L72}{\texttt{stationery}}
& rigidity \\
basis and dual basis of \(E_4\) &
\href{https://github.com/yang-mills-group/OpenAI-YM-20261028/blob/3456e5e75cf8d4cd558955a5c7638fba2d68f464/YM/UVNext/FREE_STRESS_CERTIFICATE/Size.lean\#L105-L110}{\texttt{stand}},
\href{https://github.com/yang-mills-group/OpenAI-YM-20261028/blob/3456e5e75cf8d4cd558955a5c7638fba2d68f464/YM/Pert/Rigid/Pallet.lean\#L15}{\texttt{star}},
\href{https://github.com/yang-mills-group/OpenAI-YM-20261028/blob/3456e5e75cf8d4cd558955a5c7638fba2d68f464/YM/UVNext/FREE_STRESS_CERTIFICATE/Size.lean\#L78-L80}{\texttt{portLine}},
\href{https://github.com/yang-mills-group/OpenAI-YM-20261028/blob/3456e5e75cf8d4cd558955a5c7638fba2d68f464/YM/Pert/Rigid/Spin.lean\#L59}{\texttt{coeff}}
& \(\star_c\), \(\lambda_c\), \(\mathrm{coeff}(c,c')\) \\
rate &
\href{https://github.com/yang-mills-group/OpenAI-YM-20261028/blob/3456e5e75cf8d4cd558955a5c7638fba2d68f464/YM/Pert/Rigid/Camera.lean\#L37}{\texttt{rate}}
& \(\rho(c)\) \\
test system; free readings; chart &
\href{https://github.com/yang-mills-group/OpenAI-YM-20261028/blob/3456e5e75cf8d4cd558955a5c7638fba2d68f464/YM/Pert/Position/Board.lean\#L54-L59}{\texttt{Conditions}},
\href{https://github.com/yang-mills-group/OpenAI-YM-20261028/blob/3456e5e75cf8d4cd558955a5c7638fba2d68f464/YM/Pert/Position/Board.lean\#L103-L104}{\texttt{white}},
\href{https://github.com/yang-mills-group/OpenAI-YM-20261028/blob/3456e5e75cf8d4cd558955a5c7638fba2d68f464/YM/Pert/Position/Board.lean\#L69}{\texttt{chart}},
\href{https://github.com/yang-mills-group/OpenAI-YM-20261028/blob/3456e5e75cf8d4cd558955a5c7638fba2d68f464/YM/Pert/Position/Board.lean\#L106}{\texttt{square}}
& \(Q\), \(w_Q\), \(\mathrm{chart}_Q\) \\
measured tested readings &
\href{https://github.com/yang-mills-group/OpenAI-YM-20261028/blob/3456e5e75cf8d4cd558955a5c7638fba2d68f464/YM/Pert/Position/Actual.lean\#L15}{\texttt{measured}},
\href{https://github.com/yang-mills-group/OpenAI-YM-20261028/blob/3456e5e75cf8d4cd558955a5c7638fba2d68f464/YM/Pert/Position/RawMean.lean\#L19}{\texttt{exactMean}},
\href{https://github.com/yang-mills-group/OpenAI-YM-20261028/blob/3456e5e75cf8d4cd558955a5c7638fba2d68f464/YM/Pert/Position/Actual.lean\#L42}{\texttt{testedMean}}
& \(m_Q(k)\) \\
leading position comparison &
\href{https://github.com/yang-mills-group/OpenAI-YM-20261028/blob/3456e5e75cf8d4cd558955a5c7638fba2d68f464/YM/Pert/Position/Four.lean\#L45-L47}{\texttt{FourMeters}},
\href{https://github.com/yang-mills-group/OpenAI-YM-20261028/blob/3456e5e75cf8d4cd558955a5c7638fba2d68f464/YM/Pert/Position/Four.lean\#L23-L30}{\texttt{Resident}},
\href{https://github.com/yang-mills-group/OpenAI-YM-20261028/blob/3456e5e75cf8d4cd558955a5c7638fba2d68f464/YM/Pert/Position/Crest.lean\#L45-L48}{\texttt{Crest}},
\href{https://github.com/yang-mills-group/OpenAI-YM-20261028/blob/3456e5e75cf8d4cd558955a5c7638fba2d68f464/YM/Pert/Position/Crest.lean\#L26-L34}{\texttt{BoardAttains}}
& (D15) \\
mode, track, Wick profile &
\href{https://github.com/yang-mills-group/OpenAI-YM-20261028/blob/3456e5e75cf8d4cd558955a5c7638fba2d68f464/ChallengeSpec.lean\#L152-L158}{\texttt{Mode}},
\href{https://github.com/yang-mills-group/OpenAI-YM-20261028/blob/3456e5e75cf8d4cd558955a5c7638fba2d68f464/ChallengeSpec.lean\#L162-L164}{\texttt{Track}},
\href{https://github.com/yang-mills-group/OpenAI-YM-20261028/blob/3456e5e75cf8d4cd558955a5c7638fba2d68f464/ChallengeSpec.lean\#L165-L166}{\texttt{wickProfile}},
\href{https://github.com/yang-mills-group/OpenAI-YM-20261028/blob/3456e5e75cf8d4cd558955a5c7638fba2d68f464/YM/TraceExtension/Wick/Refute.lean\#L16}{\texttt{atPoint}}
& \(K_{\mathsf m,q,h}\), \(W\) \\
span of \(r^a(\log r)^k\); full expansion &
\href{https://github.com/yang-mills-group/OpenAI-YM-20261028/blob/3456e5e75cf8d4cd558955a5c7638fba2d68f464/YM/TraceExtension/Wick/Expansion.lean\#L23}{\texttt{board}},
\href{https://github.com/yang-mills-group/OpenAI-YM-20261028/blob/3456e5e75cf8d4cd558955a5c7638fba2d68f464/YM/TraceExtension/Wick/Expansion.lean\#L47}{\texttt{Expands}},
\href{https://github.com/yang-mills-group/OpenAI-YM-20261028/blob/3456e5e75cf8d4cd558955a5c7638fba2d68f464/YM/TraceExtension/Wick/Expansion.lean\#L46}{\texttt{AtOrder}}
& \(\mathcal B\) \\
trace two-point function &
\href{https://github.com/yang-mills-group/OpenAI-YM-20261028/blob/3456e5e75cf8d4cd558955a5c7638fba2d68f464/ChallengeSpec.lean\#L149}{\texttt{traceTwo}},
\href{https://github.com/yang-mills-group/OpenAI-YM-20261028/blob/3456e5e75cf8d4cd558955a5c7638fba2d68f464/YM/TraceExtension/WilsonTrace.lean\#L53}{\texttt{euclidean}}
& \(G_\Theta(r)\) \\
perturbative weight, expectation, normalization &
\href{https://github.com/yang-mills-group/OpenAI-YM-20261028/blob/3456e5e75cf8d4cd558955a5c7638fba2d68f464/YM/Pert/Lattice/Series.lean\#L61}{\texttt{Orb.weight}},
\href{https://github.com/yang-mills-group/OpenAI-YM-20261028/blob/3456e5e75cf8d4cd558955a5c7638fba2d68f464/YM/Pert/Lattice/Series.lean\#L74}{\texttt{expect}},
\href{https://github.com/yang-mills-group/OpenAI-YM-20261028/blob/3456e5e75cf8d4cd558955a5c7638fba2d68f464/YM/Pert/Lattice/Series.lean\#L85}{\texttt{normalize}}
& \(w\), \(\mathbb E_0\), \(\langle\cdot\rangle_{\rm pert}\) \\
lattice perturbative profile &
\href{https://github.com/yang-mills-group/OpenAI-YM-20261028/blob/3456e5e75cf8d4cd558955a5c7638fba2d68f464/YM/Pert/Lattice/Bundles.lean\#L55-L56}{\texttt{Full.profile}},
\href{https://github.com/yang-mills-group/OpenAI-YM-20261028/blob/3456e5e75cf8d4cd558955a5c7638fba2d68f464/YM/Pert/Position/Board.lean\#L98}{\texttt{Position.finite}}
& \(\pi^{(n)}_Q(L,m)\) \\
\end{longtable}
}

\subsubsection{7. What remains
unexplained}\label{what-remains-unexplained-27}

\begin{enumerate}
\def\labelenumi{\arabic{enumi}.}
\tightlist
\item
  \textbf{Image versus intrinsic invariant algebra.} The statements use
  \(\mathrm{sh}_d=\omega_d(\mathrm{Label}_d)\) exactly. An intrinsic
  characterization as all gauge-invariant on-shell polynomials of weight
  at most \(d\) would require a separate surjectivity argument; no such
  characterization is needed here.
\item
  \textbf{Interacting rank estimates.} The free columns enter the
  construction of determinant trials in III.11, and the free-null
  description enters the upper-rank comparison. The estimates that turn
  these into interacting rank bounds and compatible interacting field
  spaces are inputs from III.11 and III.12; their proofs are outside
  this unit's verification.
\item
  \textbf{Continuum Wick subtraction.} The auxiliary imaginary Gaussian
  average in VII.3.v is the exact lattice definition. Its comparison
  with continuum Wick-ordered correlations uses VII.1's order-zero
  profile theorem. The underlying regulator-removal and Wick-subtraction
  estimates are accepted as VII.1 inputs; their proofs are not checked
  here.
\item
  \textbf{Further non-freeness statements.} The larger trace-exclusion
  package also contains exclusions involving embeddings into Maxwell
  Fock spaces and invariant averages. Their proofs are not covered by
  VII.3.7, which treats finite mass-shell contraction profiles and their
  tested two-point identities.
\end{enumerate}

\subsection{VII.4. Operator product expansion and contact
extensions}\label{vii.4.-operator-product-expansion-and-contact-extensions}

\textbf{Part:} VII, identification as Yang--Mills. \textbf{Lean:}
\href{https://github.com/yang-mills-group/OpenAI-YM-20261028/tree/3456e5e75cf8d4cd558955a5c7638fba2d68f464/YM/OPE}{\texttt{YM/OPE}}.
Its folders, by role:

\begin{itemize}
\tightlist
\item
  prescription and Reader: \texttt{Splitting}, \texttt{Orbits};
\item
  bounds in the ground algebra: \texttt{Eye}, \texttt{Cards},
  \texttt{Exchange};
\item
  lattice cluster estimate: \texttt{Sweep}, \texttt{Replicas},
  \texttt{Grouping}, \texttt{Columns}, \texttt{Prices};
\item
  passage to arbitrary configurations: \texttt{Schemes},
  \texttt{Interpolate}, \texttt{Height}, \texttt{Rate};
\item
  trees: \texttt{Nesting}, \texttt{Trees}, \texttt{Single},
  \texttt{Subpower}.
\end{itemize}

Also
\href{https://github.com/yang-mills-group/OpenAI-YM-20261028/tree/3456e5e75cf8d4cd558955a5c7638fba2d68f464/YM/Next/PHASE}{\texttt{YM/Next/PHASE}},
\href{https://github.com/yang-mills-group/OpenAI-YM-20261028/tree/3456e5e75cf8d4cd558955a5c7638fba2d68f464/YM/Contacts}{\texttt{YM/Contacts}}
(\texttt{Normal}, \texttt{Score}, \texttt{Layers}, \texttt{Harvest},
\texttt{Solstice}, \texttt{Transport}),
\href{https://github.com/yang-mills-group/OpenAI-YM-20261028/tree/3456e5e75cf8d4cd558955a5c7638fba2d68f464/YM/MultipleContact}{\texttt{YM/MultipleContact}}
and
\href{https://github.com/yang-mills-group/OpenAI-YM-20261028/tree/3456e5e75cf8d4cd558955a5c7638fba2d68f464/YM/SharpGram}{\texttt{YM/SharpGram}}.
The public statements are
\href{https://github.com/yang-mills-group/OpenAI-YM-20261028/blob/3456e5e75cf8d4cd558955a5c7638fba2d68f464/HierarchyBridge.lean\#L68-L77}{\texttt{publicOPE}}
and
\href{https://github.com/yang-mills-group/OpenAI-YM-20261028/blob/3456e5e75cf8d4cd558955a5c7638fba2d68f464/ContactBridge.lean\#L14-L49}{\texttt{contactDistribution}}.

\textbf{Size:} about 10,200 lines in 2,035 declarations assigned to this
unit:

\begin{itemize}
\tightlist
\item
  \texttt{OPE} about 4,600, mostly in \texttt{Grouping},
  \texttt{Replicas}, \texttt{Columns}, \texttt{Trees}, \texttt{Sweep};
\item
  \texttt{Next/PHASE} 840;
\item
  \texttt{Contacts} about 3,400;
\item
  \texttt{MultipleContact} about 290;
\item
  \texttt{SharpGram} about 250.
\end{itemize}

\textbf{Conventions.}

\begin{itemize}
\tightlist
\item
  \(G\) is the compact gauge group, with simple Lie algebra, and
  \(\mathcal W\) is the Wilson family of the construction. Its lattice
  levels \(k=0,1,2,\dots\) have spacings \(a_k\downarrow0\). Larger
  \(k\) means finer. The sequence is \emph{gentle}:
  \(\log(a_k/a_{k+1})/\log(1/a_k)\to0\).
\item
  \(I_{N,p}(v,x)\) is the renormalized lattice observable of a lattice
  label (token) \(v\) at the site \(x\). \(\|F\|_{\mathcal T_k}\) is the
  chessboard root of a lattice function \(F\) in the centred tiling of
  scale \(a_k\). Both follow IV.1.
\item
  \(S\) is the limit table along the refined extraction order. It has
  field spaces \(F_d\) and the inclusions
  \(\iota_{m,n}\colon F_m\to F_n\) (\(m\le n\)) of the field algebra. A
  label list is \(a=(a_i)_{i\in I}\) with \(a_i\in F_{d_i}\), and its
  total degree is \(w(a)=\sum_i d_i\).
\item
  \(\mathcal A\) is the ground algebra of V.1: the commutative
  C*-algebra generated by the compressed slab operators \(M(u)\) of
  separated words \(u\). Its vacuum is \(\Omega\), its Gelfand spectrum
  is the phase space \(X\) with measure \(\nu\) of full support, and its
  characters are \(\chi_\xi\).

  \begin{itemize}
  \tightlist
  \item
    \(S^\xi(u)=\chi_\xi(M(u))\) and
    \(S^{\mathrm{mix}}(u)=\langle\Omega,M(u)\Omega\rangle\).
  \item
    Hence, for \(A\in\mathcal A\) and \(c\ge0\): \(\|A\|\le c\) if and
    only if \(|s(A)|\le c\) for every state
    \(s\in X\sqcup\{\mathrm{mix}\}\)
    (\href{https://github.com/yang-mills-group/OpenAI-YM-20261028/blob/3456e5e75cf8d4cd558955a5c7638fba2d68f464/YM/OPE/Eye/Views.lean\#L67-L68}{\texttt{glance\_contract}},
    \href{https://github.com/yang-mills-group/OpenAI-YM-20261028/blob/3456e5e75cf8d4cd558955a5c7638fba2d68f464/YM/OPE/Cards/Pressure.lean\#L89-L92}{\texttt{Cards.pressure}}).
  \end{itemize}
\item
  Clusters and clearance follow Definition 1.13 of the statement.

  \begin{itemize}
  \tightlist
  \item
    \((x,o)\) is an \emph{\((r,M,\eta)\)-cluster} if \(|x_i-o|\le Mr\)
    and \(|x_i-x_{i'}|\ge\eta r\) for \(i\ne i'\).
  \item
    \(y\in(\mathbb R^4)^J\) is \emph{\(c\)-clear of \(o\)} if
    \(|o-y_j|\ge c\) and \(|y_j-y_{j'}|\ge c\) for \(j\ne j'\).
  \item
    \(S^s(a,x;b,y)\) and \(M(a,x;b,y)\) denote the concatenated word.
  \end{itemize}
\item
  ``For all sufficiently small \(r\)'' means: there is \(r_0>0\),
  depending only on the data quantified before the phrase, such that the
  assertion holds for all \(r\in(0,r_0)\) and all configurations and
  states quantified after it.
\end{itemize}

\subsubsection{1. What this unit proves}\label{what-this-unit-proves-28}

\textbf{Proposition VII.4.1 (separated OPE; properties (O1)--(O2) of the
main theorem).}

\emph{Setting.}

\begin{itemize}
\tightlist
\item
  \((\mathcal W,S,F)\) is a family, limit table and final record
  satisfying the RG hypotheses of Remark VII.4.7, denoted (RG).
\item
  \(\mathcal P\) is any coherent prescription (VII.4.6).
\item
  \(\mathcal R\) is the Reader built from \(\mathcal P\) and the
  rotation and reality structures of \(F\) (Lemma VII.4.c). It provides
  strata \(E_j\subset F_j\), bases fixed by conjugation
  \((O_{j,\kappa})_{\kappa\in\mathrm{Ch}_j}\) of \(E_j\), and
  coefficients \(c_{j,\kappa}(a;x,o)\).
\end{itemize}

\emph{Conclusions.}

\begin{itemize}
\tightlist
\item
  \textbf{(O1)} Fix a finite \(I\), labels \(a\), \(M\ge0\), \(\eta>0\),
  \(\varepsilon>0\), \(j\) and \(\kappa\in\mathrm{Ch}_j\). For all
  sufficiently small \(r\) and every \((r,M,\eta)\)-cluster \((x,o)\):
  \(|c_{j,\kappa}(a;x,o)|\le r^{\,j-w(a)-\varepsilon}\).
\item
  \textbf{(O2)} Fix finite \(I,J\), labels \(a,b\), \(M\ge0\),
  \(\eta,c,\varepsilon>0\) and \(D\in\mathbb N\). For all sufficiently
  small \(r\), every \((r,M,\eta)\)-cluster \((x,o)\) and every \(y\)
  that is \(c\)-clear of \(o\):

  \begin{itemize}
  \tightlist
  \item
    the points of \((x,y)\) are pairwise distinct, and so are those of
    \((o,y)\);
  \item
    the remainder satisfies
    \[\Big\|M(a,x;b,y)-\sum_{j\le D}\sum_{\kappa\in\mathrm{Ch}_j}c_{j,\kappa}(a;x,o)\,M\big((j,O_{j,\kappa}),o;\,b,y\big)\Big\|_{\mathcal A}\le r^{\,D+1-w(a)-\varepsilon};\]
  \item
    hence
    \(\big|S^s(a,x;b,y)-\sum_{j\le D}\sum_\kappa c_{j,\kappa}(a;x,o)\,S^s((j,O_{j,\kappa}),o;b,y)\big|\le r^{D+1-w(a)-\varepsilon}\)
    for every \(s\in X\sqcup\{\mathrm{mix}\}\).
  \end{itemize}
\end{itemize}

The coefficients do not depend on \(b\), \(y\), \(s\) or \(D\). The
thresholds do not depend on \(x\), \(o\), \(y\) or \(s\).

\emph{Status:} theorem for every Reader on the given final record under
(RG). The main construction supplies the coherent prescription and (RG)
in VIII.1.

\emph{Lean:}

\begin{itemize}
\tightlist
\item
  \href{https://github.com/yang-mills-group/OpenAI-YM-20261028/blob/3456e5e75cf8d4cd558955a5c7638fba2d68f464/YM/OPE/Orbits/Kernel.lean\#L62-L64}{\texttt{Orbits.Foundation}}:
  clause \texttt{coefficients} is (O1); clause \texttt{discrepancy}
  gives injectivity, the bound in \(\mathcal A\) and the bound in every
  state.
\item
  \href{https://github.com/yang-mills-group/OpenAI-YM-20261028/blob/3456e5e75cf8d4cd558955a5c7638fba2d68f464/YM/OPE/Orbits/Kernel.lean\#L67-L68}{\texttt{Orbits.produce}},
  \href{https://github.com/yang-mills-group/OpenAI-YM-20261028/blob/3456e5e75cf8d4cd558955a5c7638fba2d68f464/YM/OPE/Orbits/Separated.lean\#L64-L69}{\texttt{coefficient\_separated}},
  \href{https://github.com/yang-mills-group/OpenAI-YM-20261028/blob/3456e5e75cf8d4cd558955a5c7638fba2d68f464/YM/OPE/Orbits/Separated.lean\#L71-L85}{\texttt{quantitative\_separated}}.
\item
  The error element is
  \href{https://github.com/yang-mills-group/OpenAI-YM-20261028/blob/3456e5e75cf8d4cd558955a5c7638fba2d68f464/YM/OPE/Orbits/System.lean\#L61-L63}{\texttt{algebraError}}.
\item
  Headline:
  \href{https://github.com/yang-mills-group/OpenAI-YM-20261028/blob/3456e5e75cf8d4cd558955a5c7638fba2d68f464/YM/Nonperturbative/HeadlineSpec.lean\#L24-L48}{\texttt{SeparatedOPE}},
  clauses \texttt{coefficient} and \texttt{remainder}, the latter via
  \href{https://github.com/yang-mills-group/OpenAI-YM-20261028/blob/3456e5e75cf8d4cd558955a5c7638fba2d68f464/YM/Independent/Core.lean\#L49-L63}{\texttt{Independent.Core.point}}.
\item
  Public form:
  \href{https://github.com/yang-mills-group/OpenAI-YM-20261028/blob/3456e5e75cf8d4cd558955a5c7638fba2d68f464/HierarchyBridge.lean\#L68-L77}{\texttt{publicOPE}}
  into
  \href{https://github.com/yang-mills-group/OpenAI-YM-20261028/blob/3456e5e75cf8d4cd558955a5c7638fba2d68f464/ChallengeSpec.lean\#L70-L80}{\texttt{Hierarchy.OperatorProductExpansion}}.
\end{itemize}

\textbf{Proposition VII.4.2 (integrated remainder; item (D8)).} Fix the
data of (O2), a measurable space \(K\) and \(B\ge0\). For all
sufficiently small \(r\) the following holds.

\emph{Hypotheses.} Let \(\mu\) be a positive measure on \(K\),
\(w\in L^1(\mu;\mathbb C)\) with \(\int|w|\,d\mu\le B\), and
\(k\mapsto(x(k),o(k),y(k))\) an a.e.-measurable configuration map.
Assume that \(\mu\)-a.e. on \(\{w\ne0\}\), \((x,o)\) is an
\((r,M,\eta)\)-cluster and \(y\) is \(c\)-clear of \(o\).

\emph{Conclusion.} Take for \(s(k)\) either the mixed state, or
\(\xi(k)\) for an arbitrary a.e.-measurable phase map
\(\xi\colon K\to X\). Then:

\begin{itemize}
\tightlist
\item
  \(w\,S^{s}(a,x;b,y)\), \(w\) times the truncated sum, and \(w\) times
  their difference are integrable;
\item
  the integrals satisfy
  \[\Big|\int_K w(k)\Big[S^{s(k)}(a,x;b,y)-\sum_{j\le D}\sum_\kappa c_{j,\kappa}(a;x,o)\,S^{s(k)}((j,O_{j,\kappa}),o;b,y)\Big]d\mu(k)\Big|\le B\,r^{D+1-w(a)-\varepsilon}.\]
\end{itemize}

In particular, for every probability measure \(\nu'\) on \(X\), the
\(\nu'\)-mixture of the phases satisfies the bound of (O2)
(\href{https://github.com/yang-mills-group/OpenAI-YM-20261028/blob/3456e5e75cf8d4cd558955a5c7638fba2d68f464/YM/OPE/Orbits/Kernel.lean\#L52-L58}{\texttt{probability\_mixtures}}).
The same pointwise threshold works for every \(K\) and \(B\).

\emph{Status:} theorem.

\emph{Lean:}
\href{https://github.com/yang-mills-group/OpenAI-YM-20261028/blob/3456e5e75cf8d4cd558955a5c7638fba2d68f464/YM/OPE/Orbits/Profile.lean\#L98-L108}{\texttt{estimates\_weighted}},
\href{https://github.com/yang-mills-group/OpenAI-YM-20261028/blob/3456e5e75cf8d4cd558955a5c7638fba2d68f464/YM/OPE/Orbits/Profile.lean\#L53-L57}{\texttt{JointEstimate}},
\href{https://github.com/yang-mills-group/OpenAI-YM-20261028/blob/3456e5e75cf8d4cd558955a5c7638fba2d68f464/YM/OPE/Orbits/Profile.lean\#L17}{\texttt{FitsOn}};
clause \texttt{profiles} of \texttt{SeparatedOPE}, through
\texttt{Core.joint}.

\textbf{Proposition VII.4.3 (nested OPE, i.e.~composition of
coefficients over trees; item (D10)).}

\emph{Trees and plots.}

\begin{itemize}
\tightlist
\item
  \(t\) is a finite rooted tree whose root is an inner node, and \(v\)
  labels its leaves.
\item
  \(w_p\) is the total degree of the leaves below the inner node \(p\).
\item
  A \emph{plot} \(z\) assigns points to the leaves and an anchor \(z_p\)
  to each inner node.
\item
  Radii \(s_p>0\) are attached to the inner nodes, with
  \(\rho_p=s_p/s_{\mathrm{parent}(p)}\) for non-root \(p\).
\end{itemize}

\emph{Admissibility.} \((s,z)\) is \emph{\((r,M,\eta)\)-admissible} if:

\begin{itemize}
\tightlist
\item
  \(s_{\mathrm{root}}=r\) and \(s_p\le s_{\mathrm{parent}(p)}\);
\item
  at every inner node \(p\), all positions of the subtree lie within
  \(Ms_p\) of \(z_p\);
\item
  positions in different children of \(p\) are at distance
  \(\ge\eta s_p\).
\end{itemize}

\emph{Costs.} Put
\[\mathrm{toll}_{m}(s)=r^{\,m-w(v)-\varepsilon}\prod_{p\ne\mathrm{root}}\rho_p^{-w_p-\varepsilon},\qquad \mathrm{nest}_D(s)=\sum_{p\ne\mathrm{root}}\rho_p^{\,D_p+1}.\]

\emph{Conclusions.} Fix \(m\), \(\kappa\in\mathrm{Ch}_m\), truncation
degrees \(D_p\) at the inner nodes, \(M\ge0\) and
\(\eta,\varepsilon>0\). For all sufficiently small \(r\) and every
admissible \((s,z)\):

\begin{itemize}
\tightlist
\item
  \begin{enumerate}
  \def\labelenumi{(\alph{enumi})}
  \tightlist
  \item
    \(|c_{m,\kappa}(v;z_{\mathrm{leaves}},z_{\mathrm{root}})|\le\mathrm{toll}_m(s)\).
  \end{enumerate}
\item
  (b)
  \(|c_{m,\kappa}(v;\cdot)-c^{(D)}_{m,\kappa}(v;z)|\le\mathrm{toll}_m(s)\,\mathrm{nest}_D(s)\).
  Define \(c^{(D)}\) recursively: each inner child \(q\) contributes a
  channel \((j_q,\kappa_q)\) with \(j_q\le D_q\), its composed
  coefficient, and the output field \(O_{j_q,\kappa_q}\) at \(z_q\).
  Each leaf child contributes its original label and the factor \(1\).
  Multiply these factors by the root coefficient for the resulting child
  labels, and sum over the inner children's channels. The root output is
  fixed at \((m,\kappa)\); \(D_{\mathrm{root}}\) is unused.
\item
  (c) Fix a finite spectator list \(b\) and \(\delta>0\) before taking
  \(r\) small. If \(y\) is at distance \(\ge\delta\) from every point of
  the plot, with mutual distances \(\ge\delta\), then the degree-\(m\)
  remainder of VII.4.1 for the leaves at \(z_{\mathrm{root}}\) has
  \(\mathcal A\)-norm, and absolute value in every state, at most
  \(\mathrm{toll}_{m+1}(s)\).
\item
  (d) Fix a measurable space \(K\). For each small \(r\), keep the radii
  \(s\) fixed throughout the integration. Let \(\mu\) be a positive
  measure on \(K\) and \(q\in L^1(\mu;\mathbb C)\) with
  \(\int|q|\,d\mu\le B\). Let the plot, and in (c) the spectators, be
  a.e.-measurable maps satisfying the stated geometry a.e. where
  \(q\ne0\). The direct and composed coefficients are integrable after
  multiplication by \(q\); their integral and the difference of their
  integrals satisfy (a)--(b) with an additional factor \(B\). The tested
  remainder has the integrated conclusion of VII.4.2 with bound
  \(B\,\mathrm{toll}_{m+1}(s)\), for the mixed state or any
  a.e.-measurable phase map. The thresholds are uniform in \(s,\mu,q,B\)
  and the configuration maps.
\end{itemize}

\emph{Status:} theorem for every Reader on the final record under (RG).

\emph{Lean:}

\begin{itemize}
\tightlist
\item
  \href{https://github.com/yang-mills-group/OpenAI-YM-20261028/blob/3456e5e75cf8d4cd558955a5c7638fba2d68f464/YM/OPE/Nesting/Composite.lean\#L33-L35}{\texttt{Composite.token}},
  \href{https://github.com/yang-mills-group/OpenAI-YM-20261028/blob/3456e5e75cf8d4cd558955a5c7638fba2d68f464/YM/OPE/Nesting/Composite.lean\#L93-L95}{\texttt{Composite.nShot}}
  distinguish unchanged leaf inputs from inner-node expansions.
\item
  \href{https://github.com/yang-mills-group/OpenAI-YM-20261028/blob/3456e5e75cf8d4cd558955a5c7638fba2d68f464/YM/Nonperturbative/HeadlineSpec.lean\#L128-L138}{\texttt{NestedOPE}}
  (clauses \texttt{direct}, \texttt{tested}, \texttt{joint},
  \texttt{jointLaw}).
\item
  \href{https://github.com/yang-mills-group/OpenAI-YM-20261028/blob/3456e5e75cf8d4cd558955a5c7638fba2d68f464/YM/OPE/Nesting/Endpoint.lean\#L17-L27}{\texttt{Nesting.Uniform}},
  \href{https://github.com/yang-mills-group/OpenAI-YM-20261028/blob/3456e5e75cf8d4cd558955a5c7638fba2d68f464/YM/OPE/Nesting/Endpoint.lean\#L31-L36}{\texttt{Nesting.uniform}}.
\item
  \href{https://github.com/yang-mills-group/OpenAI-YM-20261028/blob/3456e5e75cf8d4cd558955a5c7638fba2d68f464/YM/OPE/Nesting/Quantified.lean\#L15-L18}{\texttt{Accept}},
  \href{https://github.com/yang-mills-group/OpenAI-YM-20261028/blob/3456e5e75cf8d4cd558955a5c7638fba2d68f464/YM/OPE/Nesting/Actual.lean\#L18-L19}{\texttt{Tree.toll}},
  \href{https://github.com/yang-mills-group/OpenAI-YM-20261028/blob/3456e5e75cf8d4cd558955a5c7638fba2d68f464/YM/OPE/Nesting/Weights.lean\#L102-L103}{\texttt{Tree.nestPrice}},
  \href{https://github.com/yang-mills-group/OpenAI-YM-20261028/blob/3456e5e75cf8d4cd558955a5c7638fba2d68f464/YM/OPE/Nesting/Composite.lean\#L115}{\texttt{Composite.stacked}},
  \href{https://github.com/yang-mills-group/OpenAI-YM-20261028/blob/3456e5e75cf8d4cd558955a5c7638fba2d68f464/YM/OPE/Nesting/Integration.lean\#L66-L71}{\texttt{CoefIntegral}}.
\end{itemize}

\textbf{Proposition VII.4.4 (contact extensions; items (D2),
(D11)--(D13)).}

\emph{Setting.}

\begin{itemize}
\tightlist
\item
  \(\mathcal M\) is a system of contact devices on the final record. Its
  prescription is the Reader's
  (\(\mathcal M.\mathrm{basis}=\mathcal P\); this is (D11)) and its
  rotation action is that of \(F\).
\item
  (RG) holds.
\item
  For labels \(a_1,\dots,a_k\) of degrees \(d_i\) and a state \(s\):

  \begin{itemize}
  \tightlist
  \item
    let \(\mathrm{cl}(n)>0\) be the clearance of the devices up to
    degree \(n\);
  \item
    put \(B_0=\min(1,\mathrm{cl}(w(a)))\) and fix a subtraction length
    \(s_0\in(0,B_0/2]\); the public form uses \(s_0=B_0/4\);
  \item
    let
    \(\mathcal D^s_a(f)=\int\langle\mathrm{full}\rangle_s(f;x)\,dx\),
    the kernel being given by the forest formula of Lemma VII.4.r.
  \end{itemize}
\end{itemize}

\emph{Conclusions.}

\begin{itemize}
\tightlist
\item
  \begin{enumerate}
  \def\labelenumi{(\alph{enumi})}
  \tightlist
  \item
    \(\mathcal D^s_a\) is a distribution of order at most
    \((2+w(a))\,k\). For every compact \(K\) there is \(C_K\) with
    \(\int|\langle\mathrm{full}\rangle_s(f;x)|\,dx\le C_K\max_{|\alpha|\le(2+w(a))k}\|\partial^\alpha f\|_\infty\)
    for all \(f\) supported in \(K\).
  \end{enumerate}
\item
  (b) \textbf{(D2)} \(\mathcal D^s_a(f)=\int f(x)\,S^s(a;x)\,dx\) for
  every \(f\in C_c^\infty\) supported in \(\mathrm{Conf}_k\).
\item
  (c) The same construction gives distributions on affine slices
  \(o+P(E)\), where \(P\) is a linear isometry whose range contains the
  normal space of the diagonal of a group \(g\subset\{1,\dots,k\}\), and
  \(o\) lies on that diagonal.
\item
  (d) \textbf{(D12)} For the admissible stage families described below,
  the pairing equals the sum over subtraction choices and dyadic shells.
  Every shell kernel is integrable, the sum of their \(L^1\) norms is
  finite, and the iterated forest sums give the same pairing.
\item
  (e) \textbf{(D12)} For the single-block partition \(Q=\{g\}\) with
  \(g\ne\emptyset\), use pure dilation in its normal space \(K_g\), with
  \(R\in(0,1]\). Put \(q_g=4(|g|-1)\) and
  \(f_R(x)=R^{-q_g}f(R^{-1}x_{K_g}+x_{K_g^\perp})\). If
  \(\mathrm{supp}f\) lies in the rescaled stage and
  \(\|f\|_{C^N}\le H\), where \(N=(2+\sum_{i=1}^k d_i)(|g|-1)\), then
  for every \(\varepsilon>0\) the integral of the absolute kernel is at
  most \(C H R^{\,u-w(a_g)-\varepsilon}\). Here \(u=0\) for the full
  resolution, \(u=n\) for the weight-\(n\) coefficient and \(u=n+1\) for
  the remainder after degree \(n\). For the last two bounds the stage
  must also lie within clearance \(\mathrm{cl}(n)\) of the
  \(g\)-diagonal. For \(|g|>1\), tests supported away from that diagonal
  give the same pairing as the punctured expression, which already
  includes proper-subcluster subtractions.
\item
  (f) \textbf{(D13)} The device strata have bases fixed by conjugation,
  and the basis in degree 0 contains the unit field.
\item
  (g) \textbf{(D13)} For every positive sequence \((f_n)\) there are
  thresholds \(0<t_n\le f_n\) with \(4t_{n+1}<t_n\), \(t_n\to0\) and
  \(2(n+1)t_n\le\mathrm{cl}(n)\).
\end{itemize}

In (d)--(e), an admissible stage family consists of a partition \(Q\) of
a group \(g\), affine placements \(o+P(E)\), positive subtraction
lengths \(s\), and block scales \(R_i\in(0,1]\). Each \(o\) lies on the
block diagonals, and the range of \(P\) contains their normal spaces. A
common ceiling \(B\le\min(1,\mathrm{cl}(\sum_i d_i))\) satisfies
\(2s\le B\). At each placement a common linear frame acts by \(R_i\) on
every block normal space and preserves each orthogonal complement. The
closed stages lie in a fixed ball in frame coordinates and within normal
distance \(B\) of each block diagonal; for nonsingleton blocks
\(R_i\le C_1s\) with fixed \(C_1\). Original and spectator labels have
uniform norm bounds. Spectator positions and labels are continuous on
each stage and constant in each block's normal directions. All centroids
of nonempty subgroups of different blocks, and all spectators, are
uniformly separated by some \(\delta>0\). These conditions are the
geometric and field hypotheses for the summation and scaling estimates;
their constants may depend on the fixed family.

Uniqueness, reflection positivity on overlapping supports, contact Ward
identities for arbitrary subtractions and real-time products are not
asserted.

\emph{Status:} theorem
(\href{https://github.com/yang-mills-group/OpenAI-YM-20261028/blob/3456e5e75cf8d4cd558955a5c7638fba2d68f464/YM/Contacts/Solstice/Result.lean\#L65-L69}{\texttt{provided}}).
In the main theorem the devices are assembled in
\href{https://github.com/yang-mills-group/OpenAI-YM-20261028/blob/3456e5e75cf8d4cd558955a5c7638fba2d68f464/YM/Nonperturbative/Continuum.lean\#L14-L24}{\texttt{contacts\_on}}
(VIII.1), with detectors from
\href{https://github.com/yang-mills-group/OpenAI-YM-20261028/blob/3456e5e75cf8d4cd558955a5c7638fba2d68f464/YM/MultipleContact/Detector.lean\#L28-L35}{\texttt{detector\_exists}}.
(D2) appears in the detailed theorem
\href{https://github.com/yang-mills-group/OpenAI-YM-20261028/blob/3456e5e75cf8d4cd558955a5c7638fba2d68f464/Construction.lean\#L15-L40}{\texttt{constructed}},
not in the public joint existence statement.

\emph{Lean:}

\begin{itemize}
\tightlist
\item
  Stage hypotheses:
  \href{https://github.com/yang-mills-group/OpenAI-YM-20261028/blob/3456e5e75cf8d4cd558955a5c7638fba2d68f464/YM/Contacts/Harvest/Route.lean\#L24-L45}{\texttt{Route}},
  \href{https://github.com/yang-mills-group/OpenAI-YM-20261028/blob/3456e5e75cf8d4cd558955a5c7638fba2d68f464/YM/Contacts/Harvest/Reader.lean\#L57}{\texttt{Route.Authority}},
  \href{https://github.com/yang-mills-group/OpenAI-YM-20261028/blob/3456e5e75cf8d4cd558955a5c7638fba2d68f464/YM/Contacts/Solstice/Strict.lean\#L48-L52}{\texttt{Backdrop.Eligible}};
  output clearance:
  \href{https://github.com/yang-mills-group/OpenAI-YM-20261028/blob/3456e5e75cf8d4cd558955a5c7638fba2d68f464/YM/Contacts/Solstice/Projection.lean\#L28}{\texttt{sees}};
  derivative order:
  \href{https://github.com/yang-mills-group/OpenAI-YM-20261028/blob/3456e5e75cf8d4cd558955a5c7638fba2d68f464/YM/Contacts/Harvest/Aperture.lean\#L25}{\texttt{stairs}}.
\item
  \href{https://github.com/yang-mills-group/OpenAI-YM-20261028/blob/3456e5e75cf8d4cd558955a5c7638fba2d68f464/YM/Contacts/Solstice/Result.lean\#L52-L63}{\texttt{MultipleLight}},
  with clauses \texttt{slope} =
  \href{https://github.com/yang-mills-group/OpenAI-YM-20261028/blob/3456e5e75cf8d4cd558955a5c7638fba2d68f464/YM/Contacts/Solstice/Result.lean\#L33-L40}{\texttt{Sunshine}},
  \texttt{moment} =
  \href{https://github.com/yang-mills-group/OpenAI-YM-20261028/blob/3456e5e75cf8d4cd558955a5c7638fba2d68f464/YM/Contacts/Solstice/Result.lean\#L23-L31}{\texttt{Resolution}},
  and \texttt{distribution} =
  \href{https://github.com/yang-mills-group/OpenAI-YM-20261028/blob/3456e5e75cf8d4cd558955a5c7638fba2d68f464/YM/Contacts/Solstice/Venue.lean\#L30-L31}{\texttt{Charge}}
  of order
  \href{https://github.com/yang-mills-group/OpenAI-YM-20261028/blob/3456e5e75cf8d4cd558955a5c7638fba2d68f464/YM/Contacts/Solstice/Skies.lean\#L48}{\texttt{totalTestOrder}}.
\item
  \href{https://github.com/yang-mills-group/OpenAI-YM-20261028/blob/3456e5e75cf8d4cd558955a5c7638fba2d68f464/YM/Contacts/Transport/Register.lean\#L47-L51}{\texttt{SystemBasis}},
  \href{https://github.com/yang-mills-group/OpenAI-YM-20261028/blob/3456e5e75cf8d4cd558955a5c7638fba2d68f464/YM/Contacts/Solstice/Gates.lean\#L23-L28}{\texttt{Thresholds}}.
\item
  Public form:
  \href{https://github.com/yang-mills-group/OpenAI-YM-20261028/blob/3456e5e75cf8d4cd558955a5c7638fba2d68f464/ContactBridge.lean\#L14-L49}{\texttt{contactDistribution}}
  for
  \href{https://github.com/yang-mills-group/OpenAI-YM-20261028/blob/3456e5e75cf8d4cd558955a5c7638fba2d68f464/ContactSpec.lean\#L12-L14}{\texttt{Hierarchy.Contacts}}.
\end{itemize}

\textbf{Proposition VII.4.5 (slab-operator bound; the bench V.1.7).}

\emph{Hypotheses.} \((\mathcal W,f,S)\) is a seed:

\begin{itemize}
\tightlist
\item
  \(\mathcal W\) has the construction properties;
\item
  the extraction is ordered and has generic point limits;
\item
  \(S\) satisfies the Euclidean hierarchy;
\item
  the corrected field estimates of IV.1 hold;
\item
  the cell hypothesis holds: reflection positivity and the chessboard
  inequality at every cutoff, and centred bounds
  \(\|I_{N,p}(v,x)\|_{\mathcal T_m}\le A\,a_m^{-d-\varepsilon}+\eta\)
  late in the extraction.
\end{itemize}

\emph{Conclusion.} For every \(D\in\mathbb N\) and \(R>0\) there are
\(C_0>0\) and \(P\in\mathbb N\) with the following property. Let
\(d>0\), \(r\in(0,1]\), and let \(w\neq\emptyset\) be a word with these
properties:

\begin{itemize}
\tightlist
\item
  its tokens have degree \(\le D\) and coefficient norm \(\le R\);
\item
  it lies in the time slab \((0,d)\) and is \(r\)-buffered there.
\end{itemize}

Then there are a level \(b\) and a width \(0<s\le d\) with
\(s\in a_b\mathbb N\) and \(w\) in the slab \((0,s)\), such that every
word \(v\) supported at positive times has a constant \(B_v\) with
\[|S(\Theta v_n\cdot v_n)|\le B_v\,\rho^{4n}\quad(n\in\mathbb N),\qquad \rho=C_0^{|w|}\,(|w|/r)^{P|w|}.\]
Here \(v_n\) is \(v\) preceded by \(2n\) consecutive slabs of width
\(s\) that hold alternately \(w\) and its reflection. In V.1 an iterated
Schwarz inequality turns this into \(\|T_s(w)\|\le\rho\). One may take
\(P=4(D+1)\).

\emph{Status:} theorem:
\href{https://github.com/yang-mills-group/OpenAI-YM-20261028/blob/3456e5e75cf8d4cd558955a5c7638fba2d68f464/YM/Next/PHASE/Seats.lean\#L154-L160}{\texttt{Next.Phase.seats}}
proves
\href{https://github.com/yang-mills-group/OpenAI-YM-20261028/blob/3456e5e75cf8d4cd558955a5c7638fba2d68f464/YM/Bid/Phase.lean\#L27}{\texttt{Bid.Seats}}.

\begin{itemize}
\tightlist
\item
  It enters the phase construction of V.1 through
  \href{https://github.com/yang-mills-group/OpenAI-YM-20261028/blob/3456e5e75cf8d4cd558955a5c7638fba2d68f464/YM/Next/PHASE/Assembly.lean\#L69-L73}{\texttt{Next.Phase.vacuum}}.
\item
  The OPE proofs of this unit do not use it directly; they use the
  algebra \(\mathcal A\) it helps to build.
\end{itemize}

\emph{Lean:}
\href{https://github.com/yang-mills-group/OpenAI-YM-20261028/blob/3456e5e75cf8d4cd558955a5c7638fba2d68f464/YM/Bid/Phase.lean\#L21-L24}{\texttt{Bid.Bench}},
\href{https://github.com/yang-mills-group/OpenAI-YM-20261028/blob/3456e5e75cf8d4cd558955a5c7638fba2d68f464/YM/OS/Train.lean\#L70-L71}{\texttt{OS.QRead}},
\href{https://github.com/yang-mills-group/OpenAI-YM-20261028/blob/3456e5e75cf8d4cd558955a5c7638fba2d68f464/YM/OS/Slabs.lean\#L109}{\texttt{Spread.radius}};
steps in Lemmas VII.4.o--p.

\textbf{Proposition VII.4.6 (coherent prescription).}

\emph{Setting.}

\begin{itemize}
\tightlist
\item
  A \emph{probe} is a finite separated configuration of tokens; its
  first point is the anchor, and the other points are taken relative to
  it.
\item
  \(\mathrm{sp}_n\colon F_n\to\mathbb C^{\{\text{probes}\}}\) sends
  \(\phi\) to the correlations of \(\phi\) at the anchor with the probe
  tokens. This map is injective, by the definition of \(F_n\) as a
  quotient by null fields, and compatible with the inclusions.
\item
  A \emph{mode} is a pair \(i=(n,\alpha)\), where \(\alpha\) indexes a
  basis of \(F_n/\iota(F_{n-1})\), with \(F_{-1}=0\).
\end{itemize}

\emph{Conclusion.} There are, for every mode, a vector \(u_i\in F_n\)
(needle) and a finitely supported family \(\lambda_i\) of complex
weights on probes (sheet), with the following property. Write
\(\lambda_i^{(m)}=\sum_p\lambda_i(p)\,\mathrm{sp}_m(\cdot)(p)\) for the
induced functional (test) on \(F_m\). Then for every \(n\):

\begin{itemize}
\tightlist
\item
  \(\lambda_i^{(n)}(\iota u_j)=\delta_{ij}\) for all modes \(i,j\) of
  degree \(\le n\);
\item
  \(\phi=\sum_{\deg i\le n}\lambda_i^{(n)}(\phi)\,\iota u_i\) for all
  \(\phi\in F_n\).
\end{itemize}

Tests of degree \(>n\) vanish on \(F_n\). One prescription serves every
truncation degree.

\emph{Status:} theorem.

\emph{Lean:}
\href{https://github.com/yang-mills-group/OpenAI-YM-20261028/blob/3456e5e75cf8d4cd558955a5c7638fba2d68f464/YM/OPE/Splitting/Basis.lean\#L44-L46}{\texttt{Prescription}},
\href{https://github.com/yang-mills-group/OpenAI-YM-20261028/blob/3456e5e75cf8d4cd558955a5c7638fba2d68f464/YM/OPE/Splitting/Basis.lean\#L59-L61}{\texttt{SignedOff}},
\href{https://github.com/yang-mills-group/OpenAI-YM-20261028/blob/3456e5e75cf8d4cd558955a5c7638fba2d68f464/YM/OPE/Splitting/Basis.lean\#L63}{\texttt{Coherent}},
\href{https://github.com/yang-mills-group/OpenAI-YM-20261028/blob/3456e5e75cf8d4cd558955a5c7638fba2d68f464/YM/OPE/Splitting/Legacy.lean\#L20-L31}{\texttt{Prescription.chosen}},
\href{https://github.com/yang-mills-group/OpenAI-YM-20261028/blob/3456e5e75cf8d4cd558955a5c7638fba2d68f464/YM/OPE/Splitting/Legacy.lean\#L32-L35}{\texttt{highVan}}.

\textbf{Remark VII.4.7 (hypotheses and order of choices).}

\emph{The order of choices.} In the main theorem
(\href{https://github.com/yang-mills-group/OpenAI-YM-20261028/blob/3456e5e75cf8d4cd558955a5c7638fba2d68f464/YM/Nonperturbative/HeadlineSpec.lean\#L143-L155}{\texttt{WilsonExistence}}),
the family is chosen first, for \(G\) with simple Lie algebra and any
penalty \(b>0\). Then, for every extraction order, a refinement, the
table, the final record, the Reader and the devices are chosen together.

\emph{The hypotheses (RG)}
(\href{https://github.com/yang-mills-group/OpenAI-YM-20261028/blob/3456e5e75cf8d4cd558955a5c7638fba2d68f464/YM/OPE/Exchange/Arrangement.lean\#L56-L60}{\texttt{Services}})
bundle four items:

\begin{itemize}
\tightlist
\item
  the grid-approximant property;
\item
  simplicity of the Lie algebra of \(G\);
\item
  the per-step and cluster RG estimates;
\item
  the continuum extraction record.
\end{itemize}

For the table of the construction these hypotheses are proved by
\href{https://github.com/yang-mills-group/OpenAI-YM-20261028/blob/3456e5e75cf8d4cd558955a5c7638fba2d68f464/YM/RG17/Joint.lean\#L60-L90}{\texttt{Run.extractible17}}
(VIII.1, from Part III outputs; §2 input 6).

\emph{Separate existence theorems.}
\href{https://github.com/yang-mills-group/OpenAI-YM-20261028/blob/3456e5e75cf8d4cd558955a5c7638fba2d68f464/YM/OPE/Single.lean\#L32-L37}{\texttt{through\_final}}
and the raw tree theorem allow any penalty \(b>0\).
\href{https://github.com/yang-mills-group/OpenAI-YM-20261028/blob/3456e5e75cf8d4cd558955a5c7638fba2d68f464/YM/OPE/Orbits/Kernel.lean\#L70-L77}{\texttt{Orbits.foundation}},
the single-cluster existence theorem, the covariant tree existence
theorem and
\href{https://github.com/yang-mills-group/OpenAI-YM-20261028/blob/3456e5e75cf8d4cd558955a5c7638fba2d68f464/YM/Contacts/Solstice/Result.lean\#L78-L79}{\texttt{celestial}}
choose \(b=1\). The main theorem instead applies the OPE estimates to
its already chosen family.

\emph{Further results.} The same construction also gives:

\begin{itemize}
\tightlist
\item
  the OPE in the raw (complex, non-covariant) needle basis
  (\href{https://github.com/yang-mills-group/OpenAI-YM-20261028/blob/3456e5e75cf8d4cd558955a5c7638fba2d68f464/YM/OPE/Single.lean\#L40-L51}{\texttt{SingleSubpower}});
\item
  tree versions for every coherent prescription
  (\href{https://github.com/yang-mills-group/OpenAI-YM-20261028/blob/3456e5e75cf8d4cd558955a5c7638fba2d68f464/YM/OPE/Subpower.lean\#L39-L45}{\texttt{AllShapes}},
  via
  \href{https://github.com/yang-mills-group/OpenAI-YM-20261028/blob/3456e5e75cf8d4cd558955a5c7638fba2d68f464/YM/OPE/Subpower.lean\#L47-L48}{\texttt{Services.everything}}).
\end{itemize}

\emph{What is not claimed.} The coefficients depend on the prescription
and the Reader. Independence of the extraction subsequence is not
claimed (Assessment §7.5).

\subsubsection{2. What it uses}\label{what-it-uses-28}

\textbf{Inputs from other units}

\begin{enumerate}
\def\labelenumi{\arabic{enumi}.}
\tightlist
\item
  \emph{The final record and its seed} (IV.1, V.1--V.3, VI.1; assembled
  in VIII.1). This comprises:

  \begin{itemize}
  \tightlist
  \item
    the cutoff family with the construction properties (II.2, III.13);
  \item
    the hierarchy (Euclidean axioms) of \(S\), and generic point limits
    on grid configurations;
  \item
    the field algebra: inclusions \(\iota\), derivative identities,
    rotation action and its laws, reality structure.
  \end{itemize}

  Lean:
  \href{https://github.com/yang-mills-group/OpenAI-YM-20261028/blob/3456e5e75cf8d4cd558955a5c7638fba2d68f464/YM/Spec/Closed.lean\#L20-L26}{\texttt{Seed}},
  \texttt{Locals}, \texttt{Final}.
\item
  \emph{Reflection positivity and centred bounds} (IV.1, interface item
  2; I.1). Two parts:

  \begin{itemize}
  \tightlist
  \item
    the chessboard inequality at every cutoff;
  \item
    for every degree \(d\) and \(\varepsilon,R>0\), there are \(m_0\)
    and \(A\ge0\) such that, for every \(m\ge m_0\) and \(\eta>0\),
    eventually along the extraction every \(I_{N,p}(v,x)\) with
    \(\|v\|\le R\) is interior to its centred cell at scale \(a_m\),
    with root \(\le A\,a_m^{-d-\varepsilon}+\eta\), uniformly in \(v\)
    and the site \(x\).
  \end{itemize}

  Lean:
  \href{https://github.com/yang-mills-group/OpenAI-YM-20261028/blob/3456e5e75cf8d4cd558955a5c7638fba2d68f464/YM/Cells/Packing.lean\#L68-L79}{\texttt{CellHyp}}.
\item
  \emph{The ground algebra and phases} (V.1). This comprises:

  \begin{itemize}
  \tightlist
  \item
    \(\mathcal A\) commutative, with \(X\) its Gelfand spectrum and
    \(\nu\) of full support;
  \item
    \href{https://github.com/yang-mills-group/OpenAI-YM-20261028/blob/3456e5e75cf8d4cd558955a5c7638fba2d68f464/YM/Work/B17.lean\#L71-L90}{\texttt{B17.answer}}:
    \(\langle\Omega,(u^*u)^n\Omega\rangle\) is the limit, as the spacing
    \(\to\infty\), of the continuum correlation of \(2n\)
    time-translates of \(u\), alternately conjugated;
  \item
    \href{https://github.com/yang-mills-group/OpenAI-YM-20261028/blob/3456e5e75cf8d4cd558955a5c7638fba2d68f464/YM/Phase/Detect.lean\#L22-L58}{\texttt{Retain.bit}}:
    if \(c\ge0\) and \(\langle\Omega,(u^*u)^n\Omega\rangle\le c^{2n}\)
    for all \(n\), then \(|\chi_\xi(u)|\le c\) for all \(\xi\);
  \item
    the selection, the every-phase tables, and their smoothness.
  \end{itemize}
\item
  \emph{The derivative identity} (VI.1, IV.1): the derivative in one
  point of a correlation is the correlation with the derivative label,
  of degree one higher. It is used in Lemma VII.4.i.
\item
  \emph{Haar measure on \(O(4)\) and the rotation laws} (VI.1). They are
  used for the covariant Reader.
\item
  \emph{The RG hypotheses (RG).} The estimates come from Part III and
  are collected in VIII.1:

  \begin{itemize}
  \tightlist
  \item
    \href{https://github.com/yang-mills-group/OpenAI-YM-20261028/blob/3456e5e75cf8d4cd558955a5c7638fba2d68f464/YM/Org/Input.lean\#L121-L123}{\texttt{Org.Grill}}:
    for every degree bound \(d\) and level \(m\) there is \(k\ge m\)
    such that conjugation-fixed token tuples of degree \(\le d\) on
    generic level-\(k\) grid configurations admit coverings that work
    along the extraction. It is supplied in III.11 and used in the a
    priori bounds (Lemma VII.4.e).
  \item
    \href{https://github.com/yang-mills-group/OpenAI-YM-20261028/blob/3456e5e75cf8d4cd558955a5c7638fba2d68f464/YM/OPE/Replicas/Budget.lean\#L13-L16}{\texttt{Package}}
    =
    \href{https://github.com/yang-mills-group/OpenAI-YM-20261028/blob/3456e5e75cf8d4cd558955a5c7638fba2d68f464/YM/OPE/Rough/Factory.lean\#L18-L27}{\texttt{TableTickets}},
    together with cluster product cuts for every arity. It supplies
    per-step RG bounds, coefficient transport and column estimates, halo
    summability, polymer representations, and per-cell scale choices.
    These are the Part III inputs to Lemmas VII.4.f--g, collected in
    VIII.1.
  \item
    The continuum extraction record (IV.1,
    \href{https://github.com/yang-mills-group/OpenAI-YM-20261028/blob/3456e5e75cf8d4cd558955a5c7638fba2d68f464/YM/Next/EXTRACTION/Base.lean}{\texttt{Next/EXTRACTION/Base.lean}}):
    the extraction to the smooth continuum table. It is supplied with
    the extraction in III.11 and used in the grid estimates (Lemma
    VII.4.h).
  \end{itemize}
\end{enumerate}

\textbf{Standard mathematics}

\begin{itemize}
\tightlist
\item
  Commutative C*-algebras: the Gelfand isometry and norm \(=\sup\) over
  characters (Bratteli--Robinson I, §2.3.5).
\item
  Reflection positivity and chessboard estimates
  (Fröhlich--Israel--Lieb--Simon 1978; Fröhlich--Simon 1977;
  Glimm--Jaffe, Ch. 6 and 18).
\item
  The iterated Schwarz inequality for transfer operators on an invariant
  pre-Hilbert domain: exponential bounds on the iterates of an adjoint
  pair give \(\|Av\|\le D\|v\|\). This is used in V.1 (Glimm--Jaffe, Ch.
  6).
\item
  Block-spin renormalization and polymer representations (Bałaban;
  Brydges, \emph{Lectures on the renormalisation group};
  Gawędzki--Kupiainen).
\item
  The Wilson--Zimmermann OPE (Wilson 1969; Zimmermann 1973) and its
  axiomatic form (Hollands--Wald 2008). The conclusions here are
  finite-order asymptotic bounds with \(\varepsilon\)-loss.
\item
  Extension of distributions across diagonals with Taylor subtraction of
  order given by the scaling degree (Epstein--Glaser 1973;
  Brunetti--Fredenhagen 2000; Hörmander I, Thm 3.2.3--3.2.4). Forest
  formulas (Zimmermann 1969).
\item
  Normalized Haar measure on \(O(4)\), finite-dimensional duality,
  Taylor's formula with remainder, and the mean value theorem.
\end{itemize}

\textbf{Imported constants and assumptions}

\begin{itemize}
\tightlist
\item
  \(\varepsilon>0\) is an arbitrary exponent loss. The estimates do not
  determine logarithmic corrections.
\item
  Constants of the cluster estimate (Lemma VII.4.f):

  \begin{itemize}
  \tightlist
  \item
    the gain margin \(1/16\), from the RG transport factor
    \((a_k/a_m)^{1/8}\);
  \item
    the loss \((a_m/a_j)^{4|I|}\), where the exponent \(4\) is
    bookkeeping slack (\(L_i^4\ge6L_i\));
  \item
    subpower factors in the per-cell, halo and insertion bounds,
    absorbed by \(\varepsilon\).
  \end{itemize}
\item
  The constant \(P=4(D+1)\) of VII.4.5 comes from degree-\(d\) bounds of
  the form \(a_k^{-(d+1)}\) and \(a_k\ge c\,(r/|w|)^4\).
\end{itemize}

\subsubsection{3. Strategy}\label{strategy-28}

A coherent system of probes reads every output degree from the mixed
correlation table. For mode \(i\), the raw coefficient is
\(\lambda_i(a;x,o)=\sum_p\lambda_i(p)S(a,x-o;\text{probe }p)\), with
each reading set to zero when the combined configuration is not
separated
(\href{https://github.com/yang-mills-group/OpenAI-YM-20261028/blob/3456e5e75cf8d4cd558955a5c7638fba2d68f464/YM/OPE/Splitting/Detector.lean\#L80}{\texttt{Prescription.ope}}).
Biorthogonality makes these readings exact on single fields
(\href{https://github.com/yang-mills-group/OpenAI-YM-20261028/blob/3456e5e75cf8d4cd558955a5c7638fba2d68f464/YM/OPE/Splitting/Observe.lean\#L27-L31}{\texttt{sumAll}}).
Averaging over \(O(4)\), symmetrizing the input and output conjugations,
and removing lower-degree components gives the Reader's real covariant
strata. These choices are fixed before the spectator labels, state and
truncation degree.

The RG estimate replaces a lattice cluster by one label \(Q\) of
prescribed degree \(t\). Its remainder gains \((a_m/a_k)^{t+1/16}\),
while the label and remainder both pay the factor
\(a_j^{-w}(a_m/a_j)^{4|I|}\). Along the scale induction, nearby groups
fuse and separated groups pass through the next block-spin step. After
extraction, exact reconstruction of \(Q\) by the probes yields the grid
OPE. Reflection positivity bounds decks of widely translated copies by
products of chessboard roots, treating a cluster remainder as one
factor. The deck limit gives all vacuum moments of \(u^*u\) in the
ground algebra; full support of the phase measure and the Gelfand
isometry then give \(\|u\|_{\mathcal A}\). Each spectator costs only its
own one-field bound. The resulting norm estimate controls every phase
and the mixed state with the same coefficients.

Exponential rates organize the choices of lattice scales. A downward
induction in total degree extends the grid bounds to arbitrary
configurations: each derivative raises the degree by one, while the
distance to the fine grid recovers one power. At sufficiently high
degree the crude bound already suffices to start the induction. Bringing
the auxiliary scale heights together gives an arbitrary
\(\varepsilon\)-loss; refinement arguments turn these rates into uniform
small-radius estimates. Tree expansions repeat this replacement
recursively, leaving leaf labels unchanged. Contacts use a forest
recursion: each retained weight receives a normal Taylor subtraction of
its own order, and the remainder gains enough radial power for
integration. Partial collisions are handled by the same recursion on
finer partitions. Separately, the slab-operator bound follows by
applying the two-scale chessboard estimate to a reflection-positive
train and using iterated Schwarz to remove the fixed exterior cost.

\subsubsection{4. Steps}\label{steps-28}

\paragraph{4.1 Rates and sequences}\label{rates-and-sequences}

\textbf{Lemma VII.4.a (from sequences and rates to uniform small-\(r\)
bounds).}

\begin{itemize}
\tightlist
\item
  \begin{enumerate}
  \def\labelenumi{(\roman{enumi})}
  \tightlist
  \item
    Let \(P(r,z)\) and \(Q(r,z)\) be predicates. Suppose that for every
    choice \(r\mapsto z_r\) and every non-trivial refinement of
    \(r\to0^+\) along which \(P(r,z_r)\) holds eventually, a further
    refinement makes \(Q(r,z_r)\) hold eventually. Then, for all
    sufficiently small \(r\), \(P(r,z)\) implies \(Q(r,z)\) for every
    \(z\).
  \end{enumerate}
\item
  (ii) \emph{Rates.} A quantity has upper rate \(\kappa\) along
  \(T\to\infty\) if \(|f|\le e^{\kappa'T}\) eventually, for every
  \(\kappa'>\kappa\). With \(T=-\log r\) this means
  \(|f|\le r^{-\kappa-\varepsilon}\) eventually, for every
  \(\varepsilon>0\). Below, ``rate'' means upper rate unless an exact
  rate is specified.
\item
  (iii) \emph{Clusters and spectators.} Under \(T=-\log r\), an
  \((r,M,\eta)\)-cluster has extent of upper rate \(-1\) and reciprocal
  separation of upper rate \(1\). For \(c\)-clear spectators, both the
  reciprocal mutual separation and the reciprocal distance to the centre
  have upper rate \(0\).
\item
  (iv) \emph{Heights.} A sequence of levels \(m\to\infty\) has height
  \(c\) if \(\log(1/a_m)/T\to c\). Every height \(c\ge0\) is attained
  for a gentle spacing sequence.
\end{itemize}

\emph{Proof sketch.} For (i), choose a bad configuration at each bad
radius and refine to a non-trivial filter supported on those radii. The
assumed further refinement is a contradiction. Parts (ii)--(iii) follow
from \(r=e^{-T}\). For (iv), choose the first level whose logarithmic
spacing exceeds \(cT\) when \(c>0\), and \(\sqrt T\) when \(c=0\);
gentleness makes the relative overshoot tend to zero.

\emph{Lean:}
\href{https://github.com/yang-mills-group/OpenAI-YM-20261028/blob/3456e5e75cf8d4cd558955a5c7638fba2d68f464/YM/OPE/Exchange/Regime.lean\#L17-L22}{\texttt{refine\_geometry}},
\href{https://github.com/yang-mills-group/OpenAI-YM-20261028/blob/3456e5e75cf8d4cd558955a5c7638fba2d68f464/YM/OPE/Rate.lean\#L23-L24}{\texttt{RateLE}},
\href{https://github.com/yang-mills-group/OpenAI-YM-20261028/blob/3456e5e75cf8d4cd558955a5c7638fba2d68f464/YM/OPE/Exchange/Regime.lean\#L35-L37}{\texttt{loglife}},
\href{https://github.com/yang-mills-group/OpenAI-YM-20261028/blob/3456e5e75cf8d4cd558955a5c7638fba2d68f464/YM/OPE/Exchange/Regime.lean\#L38-L42}{\texttt{room\_regime}},
\href{https://github.com/yang-mills-group/OpenAI-YM-20261028/blob/3456e5e75cf8d4cd558955a5c7638fba2d68f464/YM/OPE/Exchange/Regime.lean\#L43-L47}{\texttt{free\_clear}},
\href{https://github.com/yang-mills-group/OpenAI-YM-20261028/blob/3456e5e75cf8d4cd558955a5c7638fba2d68f464/YM/OPE/Height.lean\#L18-L20}{\texttt{Height}},
\href{https://github.com/yang-mills-group/OpenAI-YM-20261028/blob/3456e5e75cf8d4cd558955a5c7638fba2d68f464/YM/OPE/Height.lean\#L58-L64}{\texttt{found}}.

\paragraph{4.2 Prescription and Reader}\label{prescription-and-reader}

\textbf{Lemma VII.4.b (extension one level at a time).} A prescription
that is biorthogonal and reconstructing on \(F_n\) extends to one with
the same property on \(F_{n+1}\), without changing the needles and
sheets of degree \(\le n\).

\emph{Proof sketch.}

\begin{itemize}
\tightlist
\item
  Let
  \(B=\sum_{\deg i\le n}\lambda^{(n+1)}_i(\cdot)\,\iota_{\deg i,n}u_i\colon F_{n+1}\to F_n\),
  a left inverse of \(\iota_{n,n+1}\).
\item
  Take any linear section of \(F_{n+1}\to F_{n+1}/\iota F_n\), and
  subtract from it its projection \(\iota\circ B\). This gives a lift
  \(\Phi\) whose image lies in the kernel of all old tests.
\item
  The new needles are \(\Phi\) applied to a basis of the quotient.
\item
  The new sheets realize the dual coordinates composed with the quotient
  map. Such sheets exist because the probe functionals span the dual of
  \(F_{n+1}\)
  (\href{https://github.com/yang-mills-group/OpenAI-YM-20261028/blob/3456e5e75cf8d4cd558955a5c7638fba2d68f464/YM/OPE/Splitting/Basis.lean\#L37-L42}{\texttt{reachReads}},
  by faithfulness).
\item
  A recursive choice, collected along \(n\), gives a coherent
  prescription.
\end{itemize}

\emph{Lean:}
\href{https://github.com/yang-mills-group/OpenAI-YM-20261028/blob/3456e5e75cf8d4cd558955a5c7638fba2d68f464/YM/OPE/Splitting/Needle.lean\#L44-L53}{\texttt{startRows}},
\href{https://github.com/yang-mills-group/OpenAI-YM-20261028/blob/3456e5e75cf8d4cd558955a5c7638fba2d68f464/YM/OPE/Splitting/Needle.lean\#L57-L92}{\texttt{nextRows}},
\href{https://github.com/yang-mills-group/OpenAI-YM-20261028/blob/3456e5e75cf8d4cd558955a5c7638fba2d68f464/YM/OPE/Splitting/Legacy.lean\#L20-L31}{\texttt{Prescription.chosen}}.

\textbf{Lemma VII.4.c (the Reader).} Fix a coherent prescription.

\emph{Construction.}

\begin{itemize}
\item
  The \emph{raw vector} of degree \(n\) is
  \(\sigma_n(a;x,o)=\sum_{\deg i\le n}\lambda_i(a;x,o)\,\iota u_i\in F_n\).
\item
  Define the covariant average
  \[A_n(a;x,o)=\int_{O(4)}\rho_n(g)\,\sigma_n(\rho(g^{-1})a;g^{-1}x,g^{-1}o)\,dg,\]
  where \(\rho(g^{-1})\) acts on every input label. Set \(A_n=0\)
  outside the domain where all rotated probes through degree \(n\) stay
  clear. The conjugation-symmetric reading is
  \[\bar\sigma_n(a;x,o)=\tfrac12\big(A_n(a;x,o)+\overline{A_n(\bar a;x,o)}\big).\]
\item
  Define recursively
  \(\pi_n=\mathrm{id}-\sum_{i<n}\iota\circ\pi_i\circ\ell_{i,n}\) on
  \(F_n\), where \(\ell_{i,n}(\phi)=\bar\sigma_i(\phi;0,0)\).
\item
  The \emph{cut tests} are \(\tau_{n,q}=\pi_n\circ\ell_{n,q}\), and the
  \emph{stratum} is \(E_n=\bigcap_{i<n}\ker\tau_{i,n}\).
\item
  The \emph{fresh vector} is \(\pi_n\bar\sigma_n(a;x,o)\in E_n\).
\item
  \((O_{n,\kappa})\) is a basis fixed by conjugation of \(E_n\), and
  \(c_{n,\kappa}(a;x,o)\) are the coordinates of the fresh vector in it.
\end{itemize}

\emph{Properties.}

\begin{itemize}
\tightlist
\item
  \(\tau_{n,q}=0\) for \(q<n\).
\item
  \(\sum_{i\le n}\iota\tau_{i,q}=\iota_{q,n}\) for \(q\le n\).
\item
  The fresh vectors and cut tests intertwine rotations and conjugation;
  simultaneous translations and permutations of the inputs preserve the
  fresh vectors. Coordinates transform in the corresponding channel
  representations.
\item
  The coefficients are smooth where the probes are clear.
\end{itemize}

\emph{Proof sketch.}

\begin{itemize}
\tightlist
\item
  On a single field of degree \(j\le n\) at the origin, the degree-\(n\)
  averaged reading is \(\iota_{j,n}\)
  (\href{https://github.com/yang-mills-group/OpenAI-YM-20261028/blob/3456e5e75cf8d4cd558955a5c7638fba2d68f464/YM/OPE/Orbits/Real.lean\#L81-L84}{\texttt{checkReal\_fixed}});
  this uses coherence and the rotation laws.
\item
  The identities for \(\tau\) then follow by strong induction on \(n\).
\item
  A complex finite-dimensional space with an antilinear involution has a
  basis of fixed vectors.
\end{itemize}

\emph{Lean:}
\href{https://github.com/yang-mills-group/OpenAI-YM-20261028/blob/3456e5e75cf8d4cd558955a5c7638fba2d68f464/YM/OPE/Orbits/Real.lean\#L43-L44}{\texttt{checkValue}},
\href{https://github.com/yang-mills-group/OpenAI-YM-20261028/blob/3456e5e75cf8d4cd558955a5c7638fba2d68f464/YM/OPE/Orbits/Cuts.lean\#L17-L23}{\texttt{Reader}},
\href{https://github.com/yang-mills-group/OpenAI-YM-20261028/blob/3456e5e75cf8d4cd558955a5c7638fba2d68f464/YM/OPE/Orbits/Average.lean\#L22}{\texttt{averageV}},
\href{https://github.com/yang-mills-group/OpenAI-YM-20261028/blob/3456e5e75cf8d4cd558955a5c7638fba2d68f464/YM/OPE/Orbits/Real.lean\#L37-L39}{\texttt{checkReal}},
\href{https://github.com/yang-mills-group/OpenAI-YM-20261028/blob/3456e5e75cf8d4cd558955a5c7638fba2d68f464/YM/OPE/Orbits/Cuts.lean\#L49}{\texttt{knife}},
\href{https://github.com/yang-mills-group/OpenAI-YM-20261028/blob/3456e5e75cf8d4cd558955a5c7638fba2d68f464/YM/OPE/Orbits/Cuts.lean\#L51}{\texttt{cutTest}},
\href{https://github.com/yang-mills-group/OpenAI-YM-20261028/blob/3456e5e75cf8d4cd558955a5c7638fba2d68f464/YM/OPE/Orbits/Cuts.lean\#L66-L68}{\texttt{reconstruct}},
\href{https://github.com/yang-mills-group/OpenAI-YM-20261028/blob/3456e5e75cf8d4cd558955a5c7638fba2d68f464/YM/OPE/Orbits/Cuts.lean\#L70-L79}{\texttt{kills\_lower}},
\href{https://github.com/yang-mills-group/OpenAI-YM-20261028/blob/3456e5e75cf8d4cd558955a5c7638fba2d68f464/YM/OPE/Orbits/Cuts.lean\#L83}{\texttt{stratum}},
\href{https://github.com/yang-mills-group/OpenAI-YM-20261028/blob/3456e5e75cf8d4cd558955a5c7638fba2d68f464/YM/OPE/Orbits/Fresh.lean\#L15}{\texttt{fresh}},
\href{https://github.com/yang-mills-group/OpenAI-YM-20261028/blob/3456e5e75cf8d4cd558955a5c7638fba2d68f464/YM/OPE/Orbits/Stars.lean\#L81}{\texttt{Reader.basis}},
\href{https://github.com/yang-mills-group/OpenAI-YM-20261028/blob/3456e5e75cf8d4cd558955a5c7638fba2d68f464/YM/OPE/Orbits/System.lean\#L18-L19}{\texttt{coef}}.

\paragraph{4.3 Norms in the ground
algebra}\label{norms-in-the-ground-algebra}

\textbf{Lemma VII.4.d (norm from decks).} Let \(u\in\mathcal A\) come
from a finite combination of separated grid words. It has lattice
versions \(U_{N,p}\) at the cutoffs.

\emph{Hypothesis.} Let \(c\ge0\). The lattice labels converge to those
defining \(u\) along a non-trivial refinement of the extraction. For
every \(n\) and all sufficiently large translation spacings \(t\),
eventually along that refinement the lattice moment of the deck
\(\prod_{l<2n}U^{(\ast)}_{N,p}\circ\tau_{2tl}\) (time translates, even
ones conjugated) has modulus \(\le c^{2n}\).

\emph{Conclusion.} \(\|u\|_{\mathcal A}\le c\).

\emph{Proof sketch.}

\begin{itemize}
\tightlist
\item
  The deck moments converge to the continuum deck correlation.
\item
  As \(t\to\infty\), that correlation tends to
  \(\langle\Omega,(u^*u)^n\Omega\rangle=\int|\chi_\xi(u)|^{2n}d\nu\)
  (B17, V.1).
\item
  Since \(\nu\) has full support, the bound for all \(n\) gives
  \(|\chi_\xi(u)|\le c\) for every \(\xi\).
\item
  By Gelfand, \(\|u\|=\sup_\xi|\chi_\xi(u)|\).
\end{itemize}

\emph{Lean:}
\href{https://github.com/yang-mills-group/OpenAI-YM-20261028/blob/3456e5e75cf8d4cd558955a5c7638fba2d68f464/YM/OPE/Cards/Pressure.lean\#L98-L104}{\texttt{Cards.carrying}},
\href{https://github.com/yang-mills-group/OpenAI-YM-20261028/blob/3456e5e75cf8d4cd558955a5c7638fba2d68f464/YM/OPE/Cards/Pressure.lean\#L89-L92}{\texttt{Cards.pressure}}.
\emph{(proof read; the geometric inputs to \texttt{B17.answer} are taken
from V.1)}

\textbf{Lemma VII.4.e (a priori and replacement bounds in
\(\mathcal A\)).} Assume (RG).

\emph{Setting.} Main points are at mutual distance \(>2a_b\). Exterior
points are at mutual distance \(>2a_l\), and at distance \(>2a_b\) from
the main points. The levels satisfy \(b\le f\le l\le n\) and \(f\le q\),
with \[|I|\cdot10a_f<a_b,\qquad |J|\cdot6a_n<a_l.\] Main tokens have
centred bounds \(B_i\) at level \(q\), and exterior tokens have bounds
\(D_j\) at level \(n\).

\emph{Conclusions.}

\begin{itemize}
\item
  \begin{enumerate}
  \def\labelenumi{(\roman{enumi})}
  \tightlist
  \item
    \(\|M(\text{main}\cup\text{exterior})\|_{\mathcal A}\le\prod_iB_i\prod_jD_j\).
  \end{enumerate}
\item
  (ii) For a cluster on a fixed grid with distinct coordinates on every
  axis, suppose lattice labels \(Q_{N,p}\) of degree \(t\) converge to
  \(Q\) along a non-trivial refinement of the extraction. Suppose
  \(R_{N,p}=\prod_i I_{N,p}(v_i,x_i)-I_{N,p}(Q_{N,p},0)\) is interior to
  its centred cell at level \(q\) with root \(\le\alpha\) along that
  refinement. Require \(b\) no finer than the grid, the origin to be
  \(2a_b\)-clear of the exterior, and both the cluster--exterior and
  origin--exterior configurations to be separated. Use the cell-fitting
  inequalities above with one main factor. Then
  \[\big\|M(v,x;\text{ext})-M((t,Q),0;\text{ext})\big\|_{\mathcal A}\le\alpha\prod_jD_j.\]
\item
  (iii) In rate form, let \(0\le\alpha\le\beta\). If reciprocal
  main--main and main--exterior distances have upper rate \(\alpha\),
  and reciprocal exterior--exterior distances have upper rate \(\beta\),
  then \(M\) has upper rate \(w(\text{main})\alpha+w(\text{ext})\beta\),
  uniformly in states.
\end{itemize}

\emph{Proof sketch.}

\begin{itemize}
\tightlist
\item
  (i)--(ii): the two-scale chessboard inequality is applied to decks of
  \(2n\) translates. Main factors occupy distinct cells of the coarse
  tiling. Exterior factors, possibly several per coarse cell, are
  grouped and bounded in a finer aligned tiling. This gives the deck
  bound \((\prod B\prod D)^{2n}\) or \((\alpha\prod D)^{2n}\).
\item
  Lemma VII.4.d then gives the norm bound at grid configurations.
  Smoothness of correlations extends it to the limit configuration.
\item
  (iii): take bounds \(a_k^{-d-\delta}\) at heights slightly above the
  separation rates.
\end{itemize}

\emph{Lean:}
\href{https://github.com/yang-mills-group/OpenAI-YM-20261028/blob/3456e5e75cf8d4cd558955a5c7638fba2d68f464/YM/OPE/Cards/Literal.lean\#L107-L127}{\texttt{Cards.plant}},
\href{https://github.com/yang-mills-group/OpenAI-YM-20261028/blob/3456e5e75cf8d4cd558955a5c7638fba2d68f464/YM/OPE/Cards/Mix.lean\#L57-L94}{\texttt{Cards.mixed}},
\href{https://github.com/yang-mills-group/OpenAI-YM-20261028/blob/3456e5e75cf8d4cd558955a5c7638fba2d68f464/YM/OPE/Cards/Mix.lean\#L45-L51}{\texttt{Cards.Hall}},
\href{https://github.com/yang-mills-group/OpenAI-YM-20261028/blob/3456e5e75cf8d4cd558955a5c7638fba2d68f464/YM/OPE/Eye/Civic.lean\#L32-L46}{\texttt{civic}},
\href{https://github.com/yang-mills-group/OpenAI-YM-20261028/blob/3456e5e75cf8d4cd558955a5c7638fba2d68f464/YM/OPE/Eye/Civic.lean\#L63-L87}{\texttt{viewing\_admits}},
\href{https://github.com/yang-mills-group/OpenAI-YM-20261028/blob/3456e5e75cf8d4cd558955a5c7638fba2d68f464/YM/OPE/Eye/Admire.lean\#L47-L62}{\texttt{mirrorMatch}},
\href{https://github.com/yang-mills-group/OpenAI-YM-20261028/blob/3456e5e75cf8d4cd558955a5c7638fba2d68f464/YM/OPE/Eye/Fast.lean\#L22-L31}{\texttt{admitted}},
\href{https://github.com/yang-mills-group/OpenAI-YM-20261028/blob/3456e5e75cf8d4cd558955a5c7638fba2d68f464/YM/OPE/Eye/Fast.lean\#L49-L54}{\texttt{phase\_budget}},
\href{https://github.com/yang-mills-group/OpenAI-YM-20261028/blob/3456e5e75cf8d4cd558955a5c7638fba2d68f464/YM/OPE/Eye/Fast.lean\#L55-L57}{\texttt{viewer\_cost}}.
\emph{(proof read; \texttt{TileDoc.global} statement only)}

\paragraph{4.4 The lattice cluster
estimate}\label{the-lattice-cluster-estimate}

\textbf{Lemma VII.4.f (lattice cluster replacement).} Fix the RG
representation, product-cut and transport data of §2 for output degree
\(t\). Let \(I\) be finite and nonempty. Levels
\(J\le k\le m<m+1\le j<N\) are ordered from coarse to fine. Require
\(\mathrm{ref}(t)\le k\) and \(\mathrm{ref}(d_i)\le j\) for every input
degree \(d_i\).

\emph{Hypotheses.}

\begin{itemize}
\item
  \emph{Geometry.} Lattice observables \(f_i=I_{N,p}(s_i,o_i)\) of
  degrees \(d_i\) sit at level-\(j\) sites whose level-\((m+1)\)
  ancestor is the centre cell. They are pairwise at least \(B_j^2\)
  level-\(j\) units apart, where \(B_j\) is the polymer scale. Require
  \[(B_{m+1}^2+1)a_{m+1}\le a_J,\qquad \mathfrak c_k^2a_k\le a_J,\qquad 2(a_m\ell_t+\tau_m)<a_J.\]
  Here \(\mathfrak c_k\) is the supplied cover count, \(\ell_t\) is the
  output stencil radius plus \(4\), and \(\tau_m=30a_m(\jmath_m+5)\) is
  the cone-tail bound with supplied aperture parameter \(\jmath_m\). All
  \(f_i\) and all \(I_{N,p}(z,0)\) are interior to the centre chamber.
\item
  \emph{Weights.} Choose \(\alpha_i>0\) with
  \(\alpha_i\,\mathrm{paid}_j(s_i)\le\zeta(d_i,j)\). Here
  \(\mathrm{paid}_j\) is the insertion-size bound and \(\zeta\) the
  allowed RG weight. Choose \(0<\alpha\le1\) below every halo allowance
  from \(m\) through \(j-1\). With \(H=|I|\), require
  \[0<A\le\alpha^H\left(\frac{a_j}{a_{m+1}}\right)^{4H}\left(\frac{\alpha}{8}\right)^{H^2}\left(\frac{\delta_j\varphi_j}{8}\right)^H\prod_i\alpha_i.\]
\item
  \emph{Output weights.} Let \(\delta_m,x_m,\varphi_m,\zeta(t,k)\) and
  \(c_m\) be the supplied RG constants, with \(c_m\) the column count.
  Write \(\mathfrak m_{k,m}(b)\) for the weight obtained by transporting
  \(b>0\) from \(k\) to \(m\), and \(A_{k,t}\) for the output transport
  map. Require \(e>0\) and
  \[\mathfrak m_{k,m}(b)\le\frac{\delta_m A}{16x_m},\qquad 8(1+c_m)a_{m+1}^{t}\frac{4x_m}{A}\le e,\]
  \[32(1+c_m)b\,\|A_{k,t}\|\frac{e}{(a_k/a_m)^{1/8}}\le\frac{\zeta(t,k)}2.\]
\end{itemize}

\emph{Conclusion.} There is a lattice label \(z\) of degree \(t\) with
these properties:

\begin{itemize}
\tightlist
\item
  its coefficient size at level \(m\) is at most \(2e\);
\item
  \(R=\prod_if_i-I_{N,p}(z,0)\) is interior to the centre chamber;
\item
  \(\|R\|_{\mathcal T_J}\le2\,\mathrm{stakes}(J,k)/b\).
\end{itemize}

Here
\(\mathrm{stakes}(J,k)=\exp\big((a_J/a_k)^4\mathrm{treasury}_k+5\big)\)
is the supplied per-cell chessboard factor. In Lemma VII.4.g the two
coarse levels are chosen so that this factor has upper rate \(0\).

\emph{Proof sketch.}

\begin{itemize}
\tightlist
\item
  \emph{Transport.} From \(N\) down to \(j\), each insertion is
  transported separately, since the points are separated.
\item
  \emph{Induction from \(j\) down to \(m+1\).} The carried state is a
  polymer representation of the RG-pushed generating function, with the
  points grouped and the representatives spaced. At each step one of two
  things happens:

  \begin{itemize}
  \tightlist
  \item
    if the representatives stay spaced, one block-spin step is taken,
    and the weights lose at most a factor \(L_i^4\) per point;
  \item
    otherwise a maximal spaced set is kept and the remaining
    representatives are fused into it. Enlarging polymer supports to
    reach a representative costs a bounded factor in halo summability.
    Each fusion pays a subpower toll.
  \end{itemize}
\item
  \emph{The weights.} Their surviving product contains
  \((a_j/a_{m+1})^{4|I|}(\alpha/8)^{|I|(|I|-s)}\) when \(s\) groups
  remain. Its reciprocal contributes to the final error bound.
\item
  \emph{Final step.} At level \(m+1\), fuse the remaining groups into
  one variable per cell copy. The RG step to level \(m\) splits the
  result into a coefficient \(Q_0\) of degree \(\le t\), with its
  coefficient and remainder bounds, and an irrelevant remainder.
  Transport to the coarse level gives a factor \((a_m/a_k)^{t+1}\) in
  the remainder bound.
\item
  The constant \(2\) comes from two halves, each
  \(\le\mathrm{stakes}/b\), joined by the triangle inequality for roots.
\end{itemize}

\emph{Lean:}
\href{https://github.com/yang-mills-group/OpenAI-YM-20261028/blob/3456e5e75cf8d4cd558955a5c7638fba2d68f464/YM/OPE/Sweep/Camp.lean\#L54-L94}{\texttt{joins}},
\href{https://github.com/yang-mills-group/OpenAI-YM-20261028/blob/3456e5e75cf8d4cd558955a5c7638fba2d68f464/YM/OPE/Sweep/Induct.lean\#L93-L117}{\texttt{inductionLive}},
\href{https://github.com/yang-mills-group/OpenAI-YM-20261028/blob/3456e5e75cf8d4cd558955a5c7638fba2d68f464/YM/OPE/Sweep/Ballet.lean\#L35-L76}{\texttt{Ballet}},
\href{https://github.com/yang-mills-group/OpenAI-YM-20261028/blob/3456e5e75cf8d4cd558955a5c7638fba2d68f464/YM/OPE/Replicas/Antique.lean\#L33-L46}{\texttt{ranch}},
\href{https://github.com/yang-mills-group/OpenAI-YM-20261028/blob/3456e5e75cf8d4cd558955a5c7638fba2d68f464/YM/OPE/Grouping/Shift.lean\#L54-L93}{\texttt{hitB}}.
\emph{(statement read; proof read in outline, at the level of the calls
in \texttt{Ballet}, \texttt{inductionLive}, \texttt{Live.gathers},
\texttt{bequeathed})}

\textbf{Lemma VII.4.g (replacement rates and the limit label).} Assume
(RG), fix a nonempty finite token list of total degree \(w\), and take
heights \(0\le b<a<c\) for the coarse, output and insertion levels. Fix
a target degree \(t\), and put \(Z'=w\,c+4|I|(c-a)\).

\emph{Conclusion.} There are rate functions \(g\) and \(o\) with these
properties:

\begin{itemize}
\tightlist
\item
  \(g\) has exact rate \(Z'-ta\);
\item
  \(o\) has upper rate \(Z'-(t+\tfrac1{16})(a-b)\);
\item
  eventually, the replacement of Lemma VII.4.f has coefficient size
  \(\le g\) and remainder root \(\le o\), whenever its cluster and
  interior conditions hold.
\end{itemize}

Up to subpower factors, the coefficient size is bounded by
\(a_j^{-w}(a_m/a_j)^{4|I|}a_m^{t}\), and the remainder root by
\(a_j^{-w}(a_m/a_j)^{4|I|}(a_m/a_k)^{t+1/16}\). For each fixed
configuration satisfying the cluster and interior conditions, bounded
coefficient size gives a convergent refinement of the extraction in the
finite-dimensional label space. Its limit \(Q\) retains the remainder
bound.

\emph{Proof sketch.}

\begin{itemize}
\tightlist
\item
  Apply Lemma VII.4.f with all parameters exponential in \(T\), and
  check its size conditions as inequalities between exponents.
\item
  The loss \(4|I|(c-a)\) is the rate of \((a_{m+1}/a_j)^{4|I|}\), the
  reciprocal of the surviving weight factor.
\item
  The binding condition is the transport condition, which needs exponent
  \(<(t+\tfrac18)(a-b)-Z'\). The value \(\tfrac1{16}\) leaves strict
  slack.
\item
  Subpower factors have rate \(0\).
\end{itemize}

\emph{Lean:}
\href{https://github.com/yang-mills-group/OpenAI-YM-20261028/blob/3456e5e75cf8d4cd558955a5c7638fba2d68f464/YM/OPE/Prices/Circuit.lean\#L64-L87}{\texttt{funding}},
\href{https://github.com/yang-mills-group/OpenAI-YM-20261028/blob/3456e5e75cf8d4cd558955a5c7638fba2d68f464/YM/OPE/Prices/Circuit.lean\#L52-L59}{\texttt{PriceAt}},
\href{https://github.com/yang-mills-group/OpenAI-YM-20261028/blob/3456e5e75cf8d4cd558955a5c7638fba2d68f464/YM/OPE/Eye/Remit.lean\#L47-L65}{\texttt{price\_plot}},
\href{https://github.com/yang-mills-group/OpenAI-YM-20261028/blob/3456e5e75cf8d4cd558955a5c7638fba2d68f464/YM/OPE/Eye/Remit.lean\#L42-L44}{\texttt{Survival}},
\href{https://github.com/yang-mills-group/OpenAI-YM-20261028/blob/3456e5e75cf8d4cd558955a5c7638fba2d68f464/YM/OPE/Eye/Extract.lean\#L25-L32}{\texttt{for\_plot}}.

\paragraph{4.5 The OPE at grid
configurations}\label{the-ope-at-grid-configurations}

\textbf{Lemma VII.4.h (grid estimates).}

\emph{Setting.} Assume (RG) and fix a coherent prescription and a
nonempty finite token list \(a\). Heights satisfy \(0<\beta<a_0<x<d\),
for levels \(k,m,j\) at heights \(\beta,a_0,d\). The cluster consists of
level-\(j\) grid points with distinct coordinates on every axis, extent
of upper rate \(-x\) about the origin and reciprocal separation of upper
rate \(x\). Spectator distances from the origin have reciprocal upper
rate \(0\le c<\beta\), and their mutual distances have reciprocal upper
rate \(u\ge c\). Put \(\mathrm{tax}=w(a)d+4|I|(d-a_0)\).

\emph{Conclusions.}

\begin{itemize}
\tightlist
\item
  each fixed raw coefficient \(\lambda_i(a;\text{cluster},0)\) has upper
  rate \(\mathrm{tax}-(\deg i)a_0\);
\item
  for every \(N\in\mathbb N\), the correlation with the spectators minus
  \(\sum_{\deg i\le N}\lambda_i\,S^s(u_i,0;\text{spect.})\) has upper
  rate \(\mathrm{tax}-(N+1)(a_0-c)+w(b)u\), uniformly in the state
  \(s\).
\end{itemize}

\emph{Proof sketch.} For the remainder choose \(q\ge N\) with
\((N+1)a_0\le(q+\tfrac1{16})(a_0-\beta)\); for a coefficient choose
\(q\ge\deg i\) with the analogous inequality. Let \(Q\) be the
degree-\(q\) label of Lemma VII.4.g. Then:

\begin{itemize}
\tightlist
\item
  \begin{enumerate}
  \def\labelenumi{(\arabic{enumi})}
  \tightlist
  \item
    by Lemma VII.4.e(ii), replacing the cluster by \(Q\) at \(0\) costs
    rate \(\mathrm{tax}-(q+\tfrac1{16})(a_0-\beta)+w(b)u\) in
    \(\mathcal A\);
  \end{enumerate}
\item
  (2) the same with the probes in the role of spectators:
  \(|\lambda_i(\text{cluster})-\lambda_i(Q)|\) has the rate of the
  remainder root;
\item
  (3) for \(\deg i\le q\), the coefficient of \(Q\) is bounded by its
  coefficient size times \(a_m^{-(q-\deg i)}\), up to subpower factors,
  and hence has rate \(\mathrm{tax}-(\deg i)a_0\);
\item
  (4) exact reconstruction:
  \(S(Q,0;\cdot)=\sum_{\deg i\le q}\lambda_i(Q)\,S(u_i,0;\cdot)\);
\item
  (5) the modes of degree in \((N,q]\) contribute at most rate
  \(\mathrm{tax}-(N+1)(a_0-c)+w(b)u\), using the upper rate
  \((\deg i)c+w(b)u\) for \(S^s(u_i,0;\text{spect.})\) from Lemma
  VII.4.e(iii).
\end{itemize}

\emph{Lean:}
\href{https://github.com/yang-mills-group/OpenAI-YM-20261028/blob/3456e5e75cf8d4cd558955a5c7638fba2d68f464/YM/OPE/Schemes/Ladder.lean\#L34-L42}{\texttt{grid\_coeff}},
\href{https://github.com/yang-mills-group/OpenAI-YM-20261028/blob/3456e5e75cf8d4cd558955a5c7638fba2d68f464/YM/OPE/Schemes/Ladder.lean\#L53-L73}{\texttt{grid\_remainder}},
\href{https://github.com/yang-mills-group/OpenAI-YM-20261028/blob/3456e5e75cf8d4cd558955a5c7638fba2d68f464/YM/OPE/Eye/Extract.lean\#L64-L66}{\texttt{see\_cost}},
\href{https://github.com/yang-mills-group/OpenAI-YM-20261028/blob/3456e5e75cf8d4cd558955a5c7638fba2d68f464/YM/OPE/Schemes/Read.lean\#L68-L80}{\texttt{det\_cost}},
\href{https://github.com/yang-mills-group/OpenAI-YM-20261028/blob/3456e5e75cf8d4cd558955a5c7638fba2d68f464/YM/OPE/Schemes/Read.lean\#L37-L52}{\texttt{at\_grade}},
\href{https://github.com/yang-mills-group/OpenAI-YM-20261028/blob/3456e5e75cf8d4cd558955a5c7638fba2d68f464/YM/OPE/Splitting/Observe.lean\#L27-L31}{\texttt{sumAll}}.

\paragraph{4.6 Arbitrary configurations and
exponents}\label{arbitrary-configurations-and-exponents}

\textbf{Lemma VII.4.i (derivative bootstrap).} Fix a finite index set
\(I\) and levels \(j\) of height \(d>x_0\). A configuration path
\(x(T)\) lies in the height-\(x_0\) corridor if its extent about the
origin has upper rate \(-x_0\) and its reciprocal separation has upper
rate \(x_0\). Let \(h_T(v;x)\) be functions for every token list \(v\)
on \(I\). Suppose the following hold for every such list and every
configuration path in this corridor, with constants \(c_0,d_0\)
independent of the list:

\begin{itemize}
\tightlist
\item
  (rough) \(h_T(v;x(T))\) has upper rate \(x_0w(v)+c_0\);
\item
  (grid) on level-\(j\) grid paths with distinct coordinates on every
  axis, \(h_T(v;x(T))\) has upper rate \(dw(v)+d_0\);
\item
  (derivative) eventually \(h_T(v;\cdot)\) is differentiable at
  \(x(T)\), and its derivative in each point and coordinate equals
  \(h_T\) with that token replaced by its derivative token, of degree
  one higher.
\end{itemize}

Then \(h\) has rate \(d\,w(v)+d_0\) on the whole corridor.

\emph{Proof sketch.}

\begin{itemize}
\tightlist
\item
  Use induction downward on \(w(v)\).
\item
  For \(w(v)\) large, \((d-x_0)w(v)\ge c_0-d_0\), so the rough bound
  suffices.
\item
  Otherwise choose a grid path whose distance from \(x(T)\) has upper
  rate \(-d\). All intervening points stay in the same corridor. The
  mean value theorem and the induction bound for the derivative give
  rate \(-d+d(w(v)+1)+d_0=dw(v)+d_0\) for the change.
\end{itemize}

\emph{Lean:}
\href{https://github.com/yang-mills-group/OpenAI-YM-20261028/blob/3456e5e75cf8d4cd558955a5c7638fba2d68f464/YM/OPE/Interpolate/Corridor.lean\#L34-L47}{\texttt{interpolated}},
\href{https://github.com/yang-mills-group/OpenAI-YM-20261028/blob/3456e5e75cf8d4cd558955a5c7638fba2d68f464/YM/OPE/Interpolate/Leap.lean\#L39-L49}{\texttt{cloud\_jump}},
\href{https://github.com/yang-mills-group/OpenAI-YM-20261028/blob/3456e5e75cf8d4cd558955a5c7638fba2d68f464/YM/OPE/Interpolate/Leap.lean\#L20-L27}{\texttt{leap\_cost}}.
\emph{(proof read for \texttt{interpolated})}

\textbf{Lemma VII.4.j (raw coefficient and remainder rates).}

\emph{Setting.} Assume (RG) and fix a coherent prescription and finite
token lists \(a,b\). Along \(T\to\infty\), the cluster has extent of
upper rate \(-x\) around \(o\) and reciprocal separation of upper rate
\(x>0\). Spectator distances from \(o\) have reciprocal upper rate
\(0\le y<x\), and their mutual distances have reciprocal upper rate
\(u\ge y\).

\emph{Conclusions.}

\begin{itemize}
\tightlist
\item
  For each fixed mode \(i\), \(\lambda_i(a;\cdot,o)\) has upper rate
  \((w(a)-\deg i)x\).
\item
  For every \(N\), the remainder of the raw truncation has upper rate
  \(w(a)x-(N+1)(x-y)+w(b)u\), uniformly in states and in \(\mathcal A\).
\item
  In powers, with \(r=e^{-xT}\), spectator distance \(\rho\) and
  spectator separation \(\rho'\): the remainder is
  \(\lesssim r^{-w(a)}(r/\rho)^{N+1}\rho'^{-w(b)}\) up to
  \(\varepsilon\)-losses.
\end{itemize}

\emph{Proof sketch.}

\begin{itemize}
\tightlist
\item
  Apply Lemma VII.4.i to Lemma VII.4.h, with grid height \(d\) and
  auxiliary height \(a_0\). The rough bound is Lemma VII.4.e(iii).
\item
  This gives the bounds of Lemma VII.4.h with
  \(\mathrm{tax}=wd+4|I|(d-a_0)\). Let \(a_0\uparrow x\) and
  \(d\downarrow x\); continuity of the exponents and availability of all
  intermediate heights give the claimed rates. Empty clusters are
  handled by the exact unit identities.
\item
  Translation invariance moves the centre to \(o\).
\item
  The bound in \(\mathcal A\) follows from the bounds in every state by
  the norm characterization in the conventions.
\end{itemize}

\emph{Lean:}
\href{https://github.com/yang-mills-group/OpenAI-YM-20261028/blob/3456e5e75cf8d4cd558955a5c7638fba2d68f464/YM/OPE/Schemes/Corridor.lean\#L22-L25}{\texttt{coefficient\_rough}},
\href{https://github.com/yang-mills-group/OpenAI-YM-20261028/blob/3456e5e75cf8d4cd558955a5c7638fba2d68f464/YM/OPE/Schemes/Corridor.lean\#L37-L41}{\texttt{common\_rate}},
\href{https://github.com/yang-mills-group/OpenAI-YM-20261028/blob/3456e5e75cf8d4cd558955a5c7638fba2d68f464/YM/OPE/Schemes/Corridor.lean\#L68-L74}{\texttt{early\_rate}},
\href{https://github.com/yang-mills-group/OpenAI-YM-20261028/blob/3456e5e75cf8d4cd558955a5c7638fba2d68f464/YM/OPE/Schemes/Travel.lean\#L24-L30}{\texttt{group\_coeff}},
\href{https://github.com/yang-mills-group/OpenAI-YM-20261028/blob/3456e5e75cf8d4cd558955a5c7638fba2d68f464/YM/OPE/Schemes/Travel.lean\#L31-L41}{\texttt{group\_remainder}},
\href{https://github.com/yang-mills-group/OpenAI-YM-20261028/blob/3456e5e75cf8d4cd558955a5c7638fba2d68f464/YM/OPE/Exchange/Arrangement.lean\#L49-L52}{\texttt{ground\_cost}},
\href{https://github.com/yang-mills-group/OpenAI-YM-20261028/blob/3456e5e75cf8d4cd558955a5c7638fba2d68f464/YM/OPE/Exchange/Arrangement.lean\#L26-L27}{\texttt{groundError}},
\href{https://github.com/yang-mills-group/OpenAI-YM-20261028/blob/3456e5e75cf8d4cd558955a5c7638fba2d68f464/YM/OPE/Exchange/Arrangement.lean\#L68-L71}{\texttt{arrangement\_cost}}.

\paragraph{4.7 The Reader and the separated
OPE}\label{the-reader-and-the-separated-ope}

\textbf{Lemma VII.4.k (transfer to the Reader).} In the regime of Lemma
VII.4.j:

\begin{itemize}
\tightlist
\item
  the fresh vector of degree \(n\) has rate \((w(a)-n)x\), and hence so
  does each \(c_{n,\kappa}\);
\item
  the error element of VII.4.1 has rate \(w(a)x-(D+1)(x-y)+w(b)u\) in
  \(\mathcal A\).
\end{itemize}

\emph{Proof sketch.}

\begin{itemize}
\tightlist
\item
  Compare the cluster with its raw truncation at a horizon \(H\), placed
  at \(o\). Inside every rotated probe correlation, the difference has
  rate \((w-H-1)x\) by Lemma VII.4.j.
\item
  The raw estimate holds uniformly over all rotations, including the
  conjugate input reading. Integrating over normalized Haar measure and
  applying the fixed linear projections preserve its upper rate. Thus
  the fresh vector differs from \(\tau_{n,H}\) of the raw vector by
  upper rate \((w-H-1)x\).
\item
  In \(\tau_{n,n}\) of the raw vector, only modes of degree \(n\)
  survive, with rate \((w-n)x\).
\item
  For the remainder, write the error as the raw error plus
  \(\sum_{j\le D}\) (the correlation of
  \(\tau_{j,D}(\text{raw})-\text{fresh}_j\)), and bound each term.
\end{itemize}

\emph{Lean:}
\href{https://github.com/yang-mills-group/OpenAI-YM-20261028/blob/3456e5e75cf8d4cd558955a5c7638fba2d68f464/YM/OPE/Orbits/Fee.lean\#L45-L50}{\texttt{fresh\_bridge}},
\href{https://github.com/yang-mills-group/OpenAI-YM-20261028/blob/3456e5e75cf8d4cd558955a5c7638fba2d68f464/YM/OPE/Orbits/Gain.lean\#L21-L30}{\texttt{fresh\_price}},
\href{https://github.com/yang-mills-group/OpenAI-YM-20261028/blob/3456e5e75cf8d4cd558955a5c7638fba2d68f464/YM/OPE/Orbits/Gain.lean\#L57-L72}{\texttt{balance\_rate}},
\href{https://github.com/yang-mills-group/OpenAI-YM-20261028/blob/3456e5e75cf8d4cd558955a5c7638fba2d68f464/YM/OPE/Orbits/Separated.lean\#L17-L22}{\texttt{coefficient\_rate}},
\href{https://github.com/yang-mills-group/OpenAI-YM-20261028/blob/3456e5e75cf8d4cd558955a5c7638fba2d68f464/YM/OPE/Orbits/Separated.lean\#L30-L35}{\texttt{algebra\_rate}}.
\emph{(proof read; \texttt{fit\_old}, \texttt{discrepancy\_rate\_old}
statement only)}

\textbf{Lemma VII.4.l (proof of VII.4.1).} For (O1) and (O2), apply
Lemma VII.4.a with \(T=-\log r\).

\begin{itemize}
\tightlist
\item
  An \((r,M,\eta)\)-cluster sits at height \(x=1\).
\item
  Fixed clearance \(c\) gives \(y=u=0\).
\item
  So Lemma VII.4.k gives rates \(w-n\) and \(w-(D+1)\). These are the
  powers \(r^{n-w-\varepsilon}\) and \(r^{D+1-w-\varepsilon}\),
  uniformly in configurations and states.
\end{itemize}

\emph{Lean:}
\href{https://github.com/yang-mills-group/OpenAI-YM-20261028/blob/3456e5e75cf8d4cd558955a5c7638fba2d68f464/YM/OPE/Orbits/Separated.lean\#L64-L69}{\texttt{coefficient\_separated}},
\href{https://github.com/yang-mills-group/OpenAI-YM-20261028/blob/3456e5e75cf8d4cd558955a5c7638fba2d68f464/YM/OPE/Orbits/Separated.lean\#L71-L85}{\texttt{quantitative\_separated}},
\href{https://github.com/yang-mills-group/OpenAI-YM-20261028/blob/3456e5e75cf8d4cd558955a5c7638fba2d68f464/YM/OPE/Orbits/Kernel.lean\#L67-L68}{\texttt{produce}};
assembly in
\href{https://github.com/yang-mills-group/OpenAI-YM-20261028/blob/3456e5e75cf8d4cd558955a5c7638fba2d68f464/YM/R2I/Campaign/Accepted.lean\#L150-L167}{\texttt{R2IGoals.fromFunded}}.

\paragraph{4.8 Integrated and nested
forms}\label{integrated-and-nested-forms}

\textbf{Lemma VII.4.m (integrated forms; VII.4.2).}

\emph{Proof sketch.}

\begin{itemize}
\tightlist
\item
  For small \(r\), every admissible configuration lies in an open set on
  which the correlations and the coefficients are continuous. There, the
  raw correlation is bounded by \(r^{-w-\varepsilon}\) and the remainder
  by \(r^{D+1-w-\varepsilon}\), in every state.
\item
  So the integrands are measurable and dominated by \(|w|\) times a
  constant, and the integrals are bounded by \(B\) times the pointwise
  bound.
\item
  For any probability measure on \(X\), continuity of each fixed
  separated correlation in \(\xi\) permits termwise integration of the
  finite truncation; the same pointwise bound then gives the mixture
  estimate.
\end{itemize}

\emph{Lean:}
\href{https://github.com/yang-mills-group/OpenAI-YM-20261028/blob/3456e5e75cf8d4cd558955a5c7638fba2d68f464/YM/OPE/Orbits/Profile.lean\#L98-L108}{\texttt{estimates\_weighted}},
\href{https://github.com/yang-mills-group/OpenAI-YM-20261028/blob/3456e5e75cf8d4cd558955a5c7638fba2d68f464/YM/OPE/Orbits/Kernel.lean\#L52-L58}{\texttt{probability\_mixtures}}.

\textbf{Lemma VII.4.n (nested OPE; VII.4.3).}

\emph{Proof sketch.}

\begin{itemize}
\tightlist
\item
  Along a sequence, put \(T=\sum_p\log(1/s_p)\), summing over the inner
  nodes. The normalized logarithms lie in a compact simplex, so a
  refinement gives heights \(a_p\ge0\) with \(\sum_pa_p=1\). This also
  covers scales shrinking faster than every power of the root radius.
\item
  Then induct over the tree.
\item
  At a fork, children whose rate equals the parent's are merged into the
  parent cluster.
\item
  Each other child is replaced by its own expansion at an auxiliary
  horizon high enough that its remainder is below any prescribed rate.
\item
  The replacements are telescoped in a fixed triangular order, so that
  each switch sees only outputs already fixed.
\item
  Summing the node costs gives (a). The composition error gains the
  smallest value of \((D_p+1)(a_p-a_{\mathrm{parent}(p)})\) over
  non-root inner nodes. The extra \(\varepsilon\)-loss converts this
  rate estimate into (b). If there are no such nodes, the composed
  coefficient equals the direct coefficient exactly.
\item
  (c) and (d) follow as in Lemmas VII.4.k--m.
\end{itemize}

\emph{Lean:}
\href{https://github.com/yang-mills-group/OpenAI-YM-20261028/blob/3456e5e75cf8d4cd558955a5c7638fba2d68f464/YM/OPE/Nesting/Fusion.lean\#L97-L104}{\texttt{Fusion.tariff}},
\href{https://github.com/yang-mills-group/OpenAI-YM-20261028/blob/3456e5e75cf8d4cd558955a5c7638fba2d68f464/YM/OPE/Nesting/Quantified.lean}{\texttt{Quantified.lean}}
(\href{https://github.com/yang-mills-group/OpenAI-YM-20261028/blob/3456e5e75cf8d4cd558955a5c7638fba2d68f464/YM/OPE/Nesting/Quantified.lean\#L43-L51}{\texttt{Limits.coe\_power}},
\href{https://github.com/yang-mills-group/OpenAI-YM-20261028/blob/3456e5e75cf8d4cd558955a5c7638fba2d68f464/YM/OPE/Nesting/Quantified.lean\#L54-L71}{\texttt{Limits.tested\_power}}),
\href{https://github.com/yang-mills-group/OpenAI-YM-20261028/blob/3456e5e75cf8d4cd558955a5c7638fba2d68f464/YM/OPE/Nesting/Geometry.lean\#L21-L23}{\texttt{Tree.Tapered}},
\href{https://github.com/yang-mills-group/OpenAI-YM-20261028/blob/3456e5e75cf8d4cd558955a5c7638fba2d68f464/YM/OPE/Nesting/Grove.lean\#L99-L103}{\texttt{Tree.Relative}}.
\emph{(proof read for \texttt{coe\_power}, \texttt{tested\_power},
\texttt{tariff}; \texttt{supplies} proof read; \texttt{sequential} and
\texttt{high} statement only)}

\paragraph{4.9 The slab bound of V.1}\label{the-slab-bound-of-v.1}

\textbf{Lemma VII.4.o (point and train bounds).} Take levels
\(m\le k\le l\le b\).

\begin{itemize}
\item
  \begin{enumerate}
  \def\labelenumi{(\alph{enumi})}
  \tightlist
  \item
    \emph{Point bound.} Assume:
  \end{enumerate}

  \begin{itemize}
  \tightlist
  \item
    all points are separated; main--main and main--exterior
    sup-distances exceed \(a_m\), and exterior--exterior sup-distances
    exceed \(a_l\);
  \item
    main points are congruent mod \(a_m\) to one of \(P\) fixed points;
  \item
    \(P\cdot10a_k<a_m\) and \(\sharp(\text{ext})\cdot6a_b<a_l\);
  \item
    \(A,B\ge0\); main tokens have one-field bounds \(A\) at level \(k\)
    for all signed and reflected variants, and exterior tokens have
    bounds \(B\) at level \(b\), eventually along the extraction;
  \item
    the cutoff family has the construction properties, \(S\) is smooth
    on separated configurations, and generic point limits hold along
    this extraction.
  \end{itemize}

  Then \(|S(x;v)|\le A^{\sharp\mathrm{main}}B^{\sharp\mathrm{ext}}\).
\item
  (b) \emph{Train bound.} Assume the Euclidean hierarchy and the same
  cutoff and point-limit hypotheses. Let \(w\) be \(r\)-buffered in a
  slab of width \(d\in a_m\mathbb N\), with \(a_m<r\). Let \(v\) be a
  word supported at positive times with buffer larger than \(a_l\).
  Require \(2|w|\cdot10a_k<a_m\) and \(2|v|\cdot6a_b<a_l\), and the
  one-field bounds in (a). Then for every \(q\in\mathbb N\) the Gram
  pairing satisfies \[|S(\Theta v_q\cdot v_q)|\le B^{2|v|}A^{4q|w|}.\]
\end{itemize}

\emph{Proof sketch.}

\begin{itemize}
\tightlist
\item
  \begin{enumerate}
  \def\labelenumi{(\alph{enumi})}
  \tightlist
  \item
    The two-scale chessboard estimate at the cutoff: main factors are in
    separated coarse cells, and exterior groups in a finer tiling, with
    no restriction on how many exterior factors share a coarse cell. A
    translation of less than one exterior cell fits all the margins.
    Point limits at generic configurations and continuity give the
    continuum bound.
  \end{enumerate}
\item
  (b) The Gram pairing is the correlation of \(4q|w|\) main and \(2|v|\)
  exterior insertions at separated points. Apply (a).
\end{itemize}

\emph{Lean:}
\href{https://github.com/yang-mills-group/OpenAI-YM-20261028/blob/3456e5e75cf8d4cd558955a5c7638fba2d68f464/YM/Next/PHASE/Roots.lean\#L125-L179}{\texttt{cutoff\_bound}},
\href{https://github.com/yang-mills-group/OpenAI-YM-20261028/blob/3456e5e75cf8d4cd558955a5c7638fba2d68f464/YM/Next/PHASE/Passage.lean\#L187-L229}{\texttt{point\_bound}},
\href{https://github.com/yang-mills-group/OpenAI-YM-20261028/blob/3456e5e75cf8d4cd558955a5c7638fba2d68f464/YM/Next/PHASE/Passage.lean\#L84-L87}{\texttt{Gapped}},
\href{https://github.com/yang-mills-group/OpenAI-YM-20261028/blob/3456e5e75cf8d4cd558955a5c7638fba2d68f464/YM/Next/PHASE/Trains.lean\#L273-L304}{\texttt{travel\_price}}.

\textbf{Lemma VII.4.p (uniform constants; proof of VII.4.5).}

\emph{Proof sketch.}

\begin{itemize}
\tightlist
\item
  \emph{One-field bounds.} For degrees \(\le D\) and norms \(\le R\)
  there are \(F\) and \(J>0\) such that, for \(k\ge F\), every such
  token, together with all its signed and reflected variants, has
  one-field bound \(J\,a_k^{-(D+1)}\).
\item
  \emph{Choice of level.} For \(r\in(0,1]\), take the first \(m\ge F\)
  with \(a_m<r/32\). Gentleness gives a uniform lower bound
  \(a_m\ge c_1r^2\). A second level choice gives \(k\ge m\) with
  \(2|w|\cdot10a_k<a_m\) and \(a_k\ge c(r/|w|)^4\). Hence
  \(A^{|w|}\le\rho=C_0^{|w|}(|w|/r)^{4(D+1)|w|}\) with
  \(C_0=Jc^{-(D+1)}\).
\item
  \emph{Exterior.} For each exterior word \(v\), choose its own levels
  and bound \(X_v\). Lemma VII.4.o(b) then gives the bound of VII.4.5
  with \(B_v=X_v^{2|v|}\).
\item
  \emph{Width.} A sub-width \(s\le d\) covered at level \(m\) is chosen
  before the exterior.
\end{itemize}

\emph{Lean:}
\href{https://github.com/yang-mills-group/OpenAI-YM-20261028/blob/3456e5e75cf8d4cd558955a5c7638fba2d68f464/YM/Next/PHASE/Seats.lean\#L55-L77}{\texttt{rates}},
\href{https://github.com/yang-mills-group/OpenAI-YM-20261028/blob/3456e5e75cf8d4cd558955a5c7638fba2d68f464/YM/Next/PHASE/Seats.lean\#L119-L151}{\texttt{replicas}},
\href{https://github.com/yang-mills-group/OpenAI-YM-20261028/blob/3456e5e75cf8d4cd558955a5c7638fba2d68f464/YM/Next/PHASE/Seats.lean\#L154-L160}{\texttt{seats}}.
In V.1:
\href{https://github.com/yang-mills-group/OpenAI-YM-20261028/blob/3456e5e75cf8d4cd558955a5c7638fba2d68f464/YM/OS/Train.lean\#L20-L60}{\texttt{Drive.against}},
\href{https://github.com/yang-mills-group/OpenAI-YM-20261028/blob/3456e5e75cf8d4cd558955a5c7638fba2d68f464/YM/OS/Train.lean\#L116-L123}{\texttt{producer}}.

\paragraph{4.10 Contact extensions}\label{contact-extensions}

\textbf{Lemma VII.4.q (Taylor deviation on shells).} Let \(E\) be a
finite-dimensional Euclidean space, \(K\subset E\) a subspace,
\(m,n\in\mathbb N\), \(B\ge1\), and \(U\) a compact subset of
\(\{\tfrac12\le|x_K|\le2\}\). Write \(D_a x=ax_K+x_{K^\perp}\). There is
\(X\ge1\), depending only on \(E,K,U,m,n,B\), with the following
property. Let \(a,r,s>0\) and let \(A\colon E\to E\) be a linear
isomorphism satisfying \((Ax)_K=rx_K\) for every \(x\) and \(Ax=rx\) for
\(x\in K\). Let \(f\) be smooth and assume \(ra/s\le B\), \(r\le Bs\),
and
\[\mathrm{supp}(f\circ A)\subset\{x:|x_K|\le B/2\},\qquad \|f\circ A\|_{C^{n+m+1}}\le H.\]
Then the cut-off normal Taylor operator satisfies
\[\big\|(f-T^K_{m,s}f)\circ A\circ D_a\big\|_{C^n(U)}\le XH\begin{cases}a^{m+1},&a\le1,\\a^m,&a>1.\end{cases}\]

\emph{Proof sketch.} Taylor's formula with remainder, together with jet
bounds for the smooth cut-off.

\emph{Lean:}
\href{https://github.com/yang-mills-group/OpenAI-YM-20261028/blob/3456e5e75cf8d4cd558955a5c7638fba2d68f464/YM/Contacts/Normal/ShellTest.lean\#L43-L72}{\texttt{compact\_deviation}}.

\textbf{Lemma VII.4.r (the subtraction formula).}

\emph{Notation.}

\begin{itemize}
\tightlist
\item
  \(\varphi_i\) is the field \(a_i\) at \(x_i\).
\item
  For a group \(h\subset\{1,\dots,k\}\), \(|h|\ge2\): \(x^\perp_h\) is
  the component of \(x\) normal to the \(h\)-diagonal, \(\bar x_h\) the
  centroid, and \(L_h=\max\big(0,\ \sum_{i\in h}d_i+1-4(|h|-1)\big)\).
\item
  \(T^h_m f(x)=\chi(|x_h^\perp|^2/s_0^2)\sum_{l\le m}\frac1{l!}D^lf(x-x^\perp_h)[x^\perp_h,\dots,x^\perp_h]\),
  with \(\chi=1\) near \(0\) and \(\chi(|x^\perp_h|^2/s_0^2)=0\) for
  \(|x^\perp_h|\ge2s_0\).
\item
  \(R^h_w\) replaces a field polynomial by the weight-\(w\) component of
  its OPE projection, placed at \(\bar x_h\). The device detectors use
  finite probe sets, a linear decoder, an \(O(4)\)-average and
  simultaneous input and output conjugation, followed by the device's
  graded splitting.
\end{itemize}

\emph{Definition.} Set \(C_{\{i\}}=\varphi_i\otimes1\) and, for
\(|h|\ge2\),
\[C_h=-\sum_{w<L_h}R^h_w(\Pi_h)\cdot T^h_{L_h-1-w},\qquad \Pi_h=\sum_{\pi}\prod_{\beta\in\pi}C_\beta,\qquad \mathrm{full}_h=\Pi_h+C_h,\]
where \(\pi\) runs over partitions of \(h\) into at least two blocks.
For a singleton the full expression is \(\varphi_i\otimes1\); for the
empty group it is the unit. The kernel of \(\mathcal D^s_a\) at \(x\) is
the \(s\)-expectation of \(\mathrm{full}_{\{1..k\}}\) applied to \(f\).
Expanding the recursion gives a Zimmermann-type forest sum. The child
expressions are constructed recursively first. On the test function, the
parent Taylor operator acts first, followed by the operators already
present in the child expression; nested operators need not commute.

\emph{Per-group form.}
\(\big[\Pi_h-\sum_{w<L_h}R^h_w\Pi_h\big]f+\sum_{w<L_h}R^h_w(\Pi_h)\,(1-T^h_{L_h-1-w})f\).

\begin{itemize}
\tightlist
\item
  Put \(d_h=\sum_{i\in h}d_i\) and \(q_h=4(|h|-1)\). After retaining
  weights \(w<L_h\), the remainder has radial order
  \(O(|x^\perp_h|^{L_h-d_h-\varepsilon})\), with \(L_h-d_h\ge1-q_h\).
\item
  A retained weight-\(w\) term has coefficient order
  \(O(|x^\perp_h|^{w-d_h-\varepsilon})\) and test deviation
  \(O(|x^\perp_h|^{L_h-w})\).
\item
  For \(0<\varepsilon<1\) these radial powers are integrable in \(q_h\)
  normal dimensions. Angular sectors containing partial collisions
  require the subcluster induction of Lemma VII.4.s.
\item
  Each weight-\(w\) coefficient has scaling degree at most \(d_h-w\),
  and receives its own Taylor subtraction.
\item
  If \(\sum_{i\in h}d_i+1\le4(|h|-1)\), nothing is subtracted.
\end{itemize}

\emph{Lean:}
\href{https://github.com/yang-mills-group/OpenAI-YM-20261028/blob/3456e5e75cf8d4cd558955a5c7638fba2d68f464/YM/Contacts/Score/Tape.lean\#L56-L60}{\texttt{takeRight}},
\href{https://github.com/yang-mills-group/OpenAI-YM-20261028/blob/3456e5e75cf8d4cd558955a5c7638fba2d68f464/YM/Contacts/Score/Formula.lean}{\texttt{Score/Formula.lean}},
\href{https://github.com/yang-mills-group/OpenAI-YM-20261028/blob/3456e5e75cf8d4cd558955a5c7638fba2d68f464/YM/Contacts/Normal/Subtractions.lean\#L68}{\texttt{Normal.taylor}},
\href{https://github.com/yang-mills-group/OpenAI-YM-20261028/blob/3456e5e75cf8d4cd558955a5c7638fba2d68f464/YM/Contacts/Solstice/Venue.lean\#L59}{\texttt{majorKernel}},
\href{https://github.com/yang-mills-group/OpenAI-YM-20261028/blob/3456e5e75cf8d4cd558955a5c7638fba2d68f464/YM/Contacts/Tomography.lean\#L17-L26}{\texttt{Devices}},
\href{https://github.com/yang-mills-group/OpenAI-YM-20261028/blob/3456e5e75cf8d4cd558955a5c7638fba2d68f464/YM/Contacts/Coherent.lean\#L137-L139}{\texttt{Devices.clearance}}.
\emph{(definitions read)}

\textbf{Lemma VII.4.s (absolute bound; summation and scaling).} Fix the
devices, labels, state and subtraction length as in VII.4.4. For every
compact \(U\subset(\mathbb R^4)^k\) there is \(C>0\) such that
\[\int|\langle\mathrm{full}\rangle_s(f;x)|\,dx\le C\,\|f\|_{C^{(2+w(a))k}}\quad\text{whenever }\mathrm{supp}f\subset U.\]
The summation formulas and scaling bounds hold under the stage-family
hypotheses of VII.4.4(d)--(e).

\emph{Proof sketch.}

\begin{itemize}
\tightlist
\item
  \emph{Patching.} Cover \(U\) by finitely many patches. On each patch
  the relevant partition is the set of coincidence classes of its
  centre. A smooth partition of unity and the Leibniz rule combine the
  patch bounds.
\item
  \emph{Induction.} On a patch, use strong induction on \(|h|\) minus
  the number of blocks.

  \begin{itemize}
  \tightlist
  \item
    Base case: the separated correlation times \(f\), bounded pointwise.
  \item
    Inductive step: split into the remainder arm and the weight-\(w\)
    arms. Cut into dyadic normal shells \(|x^\perp_h|\sim2^{-n}\) and
    angular fans. Rescale each sector to a finer partition, and apply
    the induction hypothesis.
  \end{itemize}
\item
  \emph{OPE input.} The raw coefficient rate and the remainder rate in
  \(\mathcal A\) from Lemma VII.4.j.
\item
  \emph{Shell exponent.} Including normal volume, the effective exponent
  is \(4(|h|-1)+L_h-d_h\ge1\):

  \begin{itemize}
  \tightlist
  \item
    the remainder gains \(L_h\);
  \item
    a weight-\(w\) arm gains \(w\) from the coefficient and \(L_h-w\)
    from Lemma VII.4.q.
  \end{itemize}
\item
  \emph{Summing.} Choose \(\gamma>0\) with
  \((3^{|g|-|Q|}+1)\gamma\le1\), where \(|Q|\) is the number of blocks.
  The accumulated exponent losses leave a factor \(2^{-n\gamma}\) in the
  shell sum. The corresponding geometric series converges.
\end{itemize}

\emph{Lean:}
\href{https://github.com/yang-mills-group/OpenAI-YM-20261028/blob/3456e5e75cf8d4cd558955a5c7638fba2d68f464/YM/Contacts/Solstice/Skies.lean\#L59-L79}{\texttt{glues\_total}},
\href{https://github.com/yang-mills-group/OpenAI-YM-20261028/blob/3456e5e75cf8d4cd558955a5c7638fba2d68f464/YM/Contacts/Solstice/Venue.lean\#L62-L63}{\texttt{global\_charge}},
\href{https://github.com/yang-mills-group/OpenAI-YM-20261028/blob/3456e5e75cf8d4cd558955a5c7638fba2d68f464/YM/Contacts/Solstice/Result.lean\#L44-L45}{\texttt{commissions}},
\href{https://github.com/yang-mills-group/OpenAI-YM-20261028/blob/3456e5e75cf8d4cd558955a5c7638fba2d68f464/YM/Contacts/Solstice/Scaling.lean\#L61}{\texttt{Claim}},
\href{https://github.com/yang-mills-group/OpenAI-YM-20261028/blob/3456e5e75cf8d4cd558955a5c7638fba2d68f464/YM/Contacts/Solstice/Puncture.lean\#L53}{\texttt{Annulus}}.
\emph{(proof read in outline)}

\textbf{Lemma VII.4.t (distribution and restriction).}

\begin{itemize}
\tightlist
\item
  A kernel with the bound of Lemma VII.4.s defines a continuous linear
  functional on \(C_c^\infty\), of order \(\le(2+w(a))k\).
\item
  With the full group, \(P=\mathrm{id}\), \(o=0\),
  \(B_0=\min(1,\mathrm{cl}(w(a)))\) and \(s_0=B_0/4\), the hypotheses of
  Lemma VII.4.s hold on all of \((\mathbb R^4)^k\).
\item
  If \(\mathrm{supp}f\subset\mathrm{Conf}_k\), the Taylor terms evaluate
  jets of \(f\) on diagonals, where \(f\) vanishes with all its
  derivatives. So the kernel equals \(f(x)S^s(a;x)\).
\end{itemize}

\emph{Lean:}
\href{https://github.com/yang-mills-group/OpenAI-YM-20261028/blob/3456e5e75cf8d4cd558955a5c7638fba2d68f464/YM/Contacts/Solstice/Venue.lean\#L30-L31}{\texttt{Charge}},
\href{https://github.com/yang-mills-group/OpenAI-YM-20261028/blob/3456e5e75cf8d4cd558955a5c7638fba2d68f464/YM/Contacts/Transport/Blend.lean\#L70-L71}{\texttt{Transport.harmonic}},
\href{https://github.com/yang-mills-group/OpenAI-YM-20261028/blob/3456e5e75cf8d4cd558955a5c7638fba2d68f464/ContactBridge.lean\#L14-L49}{\texttt{contactDistribution}}.

\textbf{Lemma VII.4.u (real bases and thresholds).}

\begin{itemize}
\tightlist
\item
  The device strata have bases fixed by the conjugation, and the basis
  in degree 0 contains the unit. Proof: real and imaginary parts span,
  and a real independent set extends to a basis.
\item
  Given \(f_n>0\), a valid choice of thresholds is
  \[g_n=\min\big(f_n,\ 2^{-n},\ \mathrm{cl}(n)/(2(n+1))\big),\qquad t_0=g_0,\qquad t_{n+1}=\min(t_n/8,\ g_{n+1}).\]
  The thresholds define smooth scalar cutoffs. At a length
  \(0<\ell\le t_n\), the cutoff-weighted projection agrees with the
  original projection through degree \(n\). These cutoffs do not enter
  \(\mathcal D^s_a\).
\end{itemize}

\emph{Lean:}
\href{https://github.com/yang-mills-group/OpenAI-YM-20261028/blob/3456e5e75cf8d4cd558955a5c7638fba2d68f464/YM/Contacts/Transport/Register.lean\#L52-L61}{\texttt{real\_basis}},
\href{https://github.com/yang-mills-group/OpenAI-YM-20261028/blob/3456e5e75cf8d4cd558955a5c7638fba2d68f464/YM/Contacts/Solstice/Gates.lean\#L30-L36}{\texttt{installs}}.

\subsubsection{5. Where the difficulty
lies}\label{where-the-difficulty-lies-28}

\emph{Routine but long.}

\begin{itemize}
\tightlist
\item
  The rate bookkeeping (Lemmas VII.4.a, j).
\item
  The \(O(4)\)-covariant Reader (Lemmas VII.4.c, k).
\item
  The tree combinatorics (Lemma VII.4.n).
\item
  The geometry of normal shells, angular sectors and compact patches.
\end{itemize}

\emph{Main estimates.}

\begin{enumerate}
\def\labelenumi{\arabic{enumi}.}
\tightlist
\item
  \textbf{The OPE comes from the RG.} The basic estimate (Lemmas
  VII.4.f--g) is a lattice statement: a cluster of renormalized
  observables equals one observable of degree \(t\) plus a remainder
  that gains \((a_m/a_k)^{t+1/16}\), at a cost \((a_m/a_j)^{4|I|}\). The
  margin \(\tfrac1{16}\) is set by the transport gain
  \((a_k/a_m)^{1/8}\) of Part III's RG, not by the contraction
  \((a_m/a_k)^{t+1}\) of the irrelevant part. This is where Part III
  enters.
\item
  \textbf{Uniformity in states and spectators comes from the C*-norm.}
  The norm of \(\mathcal A\) is computed by chessboard bounds on decks
  of translated copies, together with Gelfand and the full support of
  \(\nu\) (Lemmas VII.4.d--e). The spectators then cost only their own
  one-field bounds, at their own scale, and the bound holds in every
  phase at once.
\item
  \textbf{The derivative bootstrap} (Lemma VII.4.i). Each derivative
  costs one power of the fine spacing, so a first-order comparison with
  the grid gains nothing. The induction therefore starts from high
  degree, where the crude bound is already strong enough.
\item
  \textbf{One coefficient system for every truncation degree.} It comes
  from a coherent prescription (VII.4.6) and the strata of the Reader.
\end{enumerate}

\emph{Unusual, weaker or stronger than expected.}

\begin{itemize}
\tightlist
\item
  \emph{Exponents.} The estimates have an arbitrary \(\varepsilon\)-loss
  and do not determine logarithmic corrections. The degree \(w\) is the
  sum of the label degrees; the statement does not identify it with a
  scaling dimension.
\item
  \emph{Coefficients.} They depend on the prescription and the Reader,
  and nothing relates them to perturbation theory (VII.5, VII.6).
\item
  \emph{Spectator distance.} The remainder bound holds with
  power-separated spectators, \((r/\rho)^{N+1}\). The main theorem uses
  only a fixed clearance.
\item
  \emph{Contacts.} The contact extension is an existence statement. Its
  order \((2+w(a))k\) is crude. Its OPE projection uses the device
  detectors, with the same prescription as the Reader but separate probe
  sets.
\item
  \emph{Interface to Part III.} The cluster replacement requires the
  numerical and geometric hypotheses in Lemma VII.4.f.
\end{itemize}

\subsubsection{6. Notation}\label{notation-28}

{\def\LTcaptype{none} % do not increment counter
\begin{longtable}[]{@{}
  >{\raggedright\arraybackslash}p{(\linewidth - 4\tabcolsep) * \real{0.3333}}
  >{\raggedright\arraybackslash}p{(\linewidth - 4\tabcolsep) * \real{0.3333}}
  >{\raggedright\arraybackslash}p{(\linewidth - 4\tabcolsep) * \real{0.3333}}@{}}
\toprule\noalign{}
\begin{minipage}[b]{\linewidth}\raggedright
object
\end{minipage} & \begin{minipage}[b]{\linewidth}\raggedright
code name
\end{minipage} & \begin{minipage}[b]{\linewidth}\raggedright
proposed notation
\end{minipage} \\
\midrule\noalign{}
\endhead
\bottomrule\noalign{}
\endlastfoot
upper and exact exponential rates &
\href{https://github.com/yang-mills-group/OpenAI-YM-20261028/blob/3456e5e75cf8d4cd558955a5c7638fba2d68f464/YM/OPE/Rate.lean\#L23-L24}{\texttt{RateLE}},
\texttt{RateIs} & \(f\preceq e^{\kappa T}\); \(\log f/T\to\kappa\) with
\(f>0\) eventually \\
height of a level sequence &
\href{https://github.com/yang-mills-group/OpenAI-YM-20261028/blob/3456e5e75cf8d4cd558955a5c7638fba2d68f464/YM/OPE/Height.lean\#L18-L20}{\texttt{Height}}
& \(m\to\infty\), \(\log(1/a_m)/T\to c\) \\
cluster, clearance &
\href{https://github.com/yang-mills-group/OpenAI-YM-20261028/blob/3456e5e75cf8d4cd558955a5c7638fba2d68f464/YM/OPE/Exchange/Regime.lean\#L24-L26}{\texttt{WithinRoom}},
\href{https://github.com/yang-mills-group/OpenAI-YM-20261028/blob/3456e5e75cf8d4cd558955a5c7638fba2d68f464/YM/OPE/Exchange/Regime.lean\#L28-L30}{\texttt{ClearOf}}
& \((r,M,\eta)\)-cluster; \(c\)-clear \\
total degree & \texttt{weight}, \texttt{sigma} & \(w(a)\) \\
needles, sheets, tests &
\href{https://github.com/yang-mills-group/OpenAI-YM-20261028/blob/3456e5e75cf8d4cd558955a5c7638fba2d68f464/YM/OPE/Splitting/Basis.lean\#L44-L46}{\texttt{Prescription}}
& \(u_i\), \(\lambda_i\), \(\lambda_i^{(n)}\) \\
raw coefficient &
\href{https://github.com/yang-mills-group/OpenAI-YM-20261028/blob/3456e5e75cf8d4cd558955a5c7638fba2d68f464/YM/OPE/Splitting/Detector.lean\#L80}{\texttt{Prescription.ope}}
& \(\lambda_i(a;x,o)\) \\
Reader &
\href{https://github.com/yang-mills-group/OpenAI-YM-20261028/blob/3456e5e75cf8d4cd558955a5c7638fba2d68f464/YM/OPE/Orbits/Cuts.lean\#L17-L23}{\texttt{Orbits.Reader}}
& \(\mathcal R\) \\
lower-degree removal, cut test, stratum & \texttt{knife},
\texttt{cutTest}, \texttt{stratum} & \(\pi_n\), \(\tau_{n,q}\),
\(E_n\) \\
channels, actors, coefficients & \texttt{channels}, \texttt{actor},
\texttt{coef} & \(\mathrm{Ch}_j\), \(O_{j,\kappa}\),
\(c_{j,\kappa}(a;x,o)\) \\
ground algebra, states & \texttt{unselected.algebra}, \texttt{glance},
\texttt{Viewers} & \(\mathcal A\); \(s\in X\sqcup\{\mathrm{mix}\}\) \\
error element &
\href{https://github.com/yang-mills-group/OpenAI-YM-20261028/blob/3456e5e75cf8d4cd558955a5c7638fba2d68f464/YM/OPE/Orbits/System.lean\#L61-L63}{\texttt{algebraError}},
\href{https://github.com/yang-mills-group/OpenAI-YM-20261028/blob/3456e5e75cf8d4cd558955a5c7638fba2d68f464/YM/OPE/Exchange/Arrangement.lean\#L26-L27}{\texttt{groundError}}
& \(\mathcal E_D(a,x;b,y)\) \\
chessboard root with interior condition &
\href{https://github.com/yang-mills-group/OpenAI-YM-20261028/blob/3456e5e75cf8d4cd558955a5c7638fba2d68f464/YM/OPE/Cards/Literal.lean\#L28-L31}{\texttt{Cards.PaidAt}}
& \(\|F\|_{\mathcal T_k}\le A\) \\
deck of translates & \texttt{deck}, \texttt{Episode} &
\(\prod_{l<2n}U^{(*)}\circ\tau_{2tl}\) \\
limit label of a cluster &
\href{https://github.com/yang-mills-group/OpenAI-YM-20261028/blob/3456e5e75cf8d4cd558955a5c7638fba2d68f464/YM/OPE/Eye/Remit.lean\#L42-L44}{\texttt{Survival}}
& \(Q\) \\
coefficient size of a label & \texttt{quota} & \(\|z\|_m\) \\
RG insertion, step and column bounds & \texttt{paid}, \texttt{ζ},
\texttt{δ}, \texttt{φ}, \texttt{x}, \texttt{columns}, \texttt{alto},
\texttt{march}, \texttt{gainChart} & \(\mathrm{paid}_j\), \(\zeta\),
\(\delta\), \(\varphi\), \(x\), \(c_m\), \(A_{k,t}\),
\(\mathfrak m_{k,m}\), \((a_k/a_m)^{1/8}\) \\
cell and stencil fitting data & \texttt{coverSize}, \texttt{lensScope},
\texttt{GridApprox.tail}, \texttt{clock.J} & \(\mathfrak c_k\),
\(\ell_t\), \(\tau_m\), \(\jmath_m\) \\
RG package &
\href{https://github.com/yang-mills-group/OpenAI-YM-20261028/blob/3456e5e75cf8d4cd558955a5c7638fba2d68f464/YM/OPE/Exchange/Arrangement.lean\#L56-L60}{\texttt{Services}}
& (RG) \\
tree plot, admissibility &
\href{https://github.com/yang-mills-group/OpenAI-YM-20261028/blob/3456e5e75cf8d4cd558955a5c7638fba2d68f464/YM/OPE/Nesting/Quantified.lean\#L15-L18}{\texttt{Accept}},
\texttt{Plot} & \((s,z)\) admissible \\
tree costs &
\href{https://github.com/yang-mills-group/OpenAI-YM-20261028/blob/3456e5e75cf8d4cd558955a5c7638fba2d68f464/YM/OPE/Nesting/Actual.lean\#L18-L19}{\texttt{Tree.toll}},
\href{https://github.com/yang-mills-group/OpenAI-YM-20261028/blob/3456e5e75cf8d4cd558955a5c7638fba2d68f464/YM/OPE/Nesting/Weights.lean\#L102-L103}{\texttt{nestPrice}}
& \(\mathrm{toll}_m\), \(\mathrm{nest}_D\) \\
one-field bound, separation &
\href{https://github.com/yang-mills-group/OpenAI-YM-20261028/blob/3456e5e75cf8d4cd558955a5c7638fba2d68f464/YM/Next/PHASE/Roots.lean\#L25-L26}{\texttt{Next.Phase.Single}},
\href{https://github.com/yang-mills-group/OpenAI-YM-20261028/blob/3456e5e75cf8d4cd558955a5c7638fba2d68f464/YM/Next/PHASE/Passage.lean\#L84-L87}{\texttt{Gapped}}
& \(\|I(v,\cdot)\|_{\mathcal T_k}\le A\) \\
train, quantum read & \texttt{ride},
\href{https://github.com/yang-mills-group/OpenAI-YM-20261028/blob/3456e5e75cf8d4cd558955a5c7638fba2d68f464/YM/OS/Train.lean\#L70-L71}{\texttt{QRead}}
& \(v_n\); bench bound \(\rho\) \\
contact devices, clearance &
\href{https://github.com/yang-mills-group/OpenAI-YM-20261028/blob/3456e5e75cf8d4cd558955a5c7638fba2d68f464/YM/Contacts/Tomography.lean\#L17-L26}{\texttt{Contacts.Devices}},
\texttt{clearance} & \(\mathcal M\), \(\mathrm{cl}(n)\) \\
cumulant, punctured, full & \texttt{cumulant}, \texttt{punctured},
\texttt{full} & \(C_h\), \(\Pi_h\), \(\mathrm{full}_h\) \\
normal Taylor operator &
\href{https://github.com/yang-mills-group/OpenAI-YM-20261028/blob/3456e5e75cf8d4cd558955a5c7638fba2d68f464/YM/Contacts/Normal/Subtractions.lean\#L68}{\texttt{Normal.taylor}}
& \(T^h_m\) \\
distribution of finite order &
\href{https://github.com/yang-mills-group/OpenAI-YM-20261028/blob/3456e5e75cf8d4cd558955a5c7638fba2d68f464/YM/Contacts/Solstice/Venue.lean\#L30-L31}{\texttt{Charge}},
\texttt{totalTestOrder} & \(\mathcal D^s_a\), order \((2+w)k\) \\
contact thresholds &
\href{https://github.com/yang-mills-group/OpenAI-YM-20261028/blob/3456e5e75cf8d4cd558955a5c7638fba2d68f464/YM/Contacts/Solstice/Gates.lean\#L23-L28}{\texttt{Thresholds}}
& \(t_n\) \\
\end{longtable}
}

\subsubsection{7. What remains
unexplained}\label{what-remains-unexplained-28}

\begin{enumerate}
\def\labelenumi{\arabic{enumi}.}
\tightlist
\item
  \textbf{RG inputs below the cluster replacement.} Lemmas VII.4.f--g
  identify the scalar inequalities and the scale induction. The polymer
  norm estimates inside the one-step representation, fusion and final
  remainder transport have not been checked line by line. The dependence
  of their constants, and the complete Part III proofs of these inputs,
  remain to be reconciled with the stated RG package.
\item
  \textbf{Proofs used through their statements.} The full geometry
  behind the vacuum deck limit and the two-scale continuum transfer
  belongs to V.1. The rotated comparison estimates used in Lemma
  VII.4.k, and the triangular replacement and high-horizon lemmas used
  in VII.4.n, have not been checked internally. Their uses and exponent
  arithmetic agree with their stated conclusions.
\item
  \textbf{Contact sector induction.} The global patching argument,
  subtraction recursion, shell Taylor bound and scaling normalization
  have been checked. The finer angular decompositions and the full
  induction through partially coincident strata remain checked only at
  the level of their theorem calls. Lemma VII.4.r's radial estimate
  alone does not settle these partial collisions.
\item
  \textbf{Contact projection versus Reader coefficients.} The devices
  share the Reader's prescription but use separate detectors and a
  separate graded splitting. An identity between their weight-\(w\) rows
  and \(\sum_\kappa c_{w,\kappa}O_{w,\kappa}\) has not been located. No
  such identity is used in the contact statement above.
\item
  \textbf{Limits of the conclusions.} The results do not assert
  uniqueness of coefficients or contact extensions, independence of the
  extraction subsequence, convergence as \(D\to\infty\), reflection
  positivity or Ward identities for the extensions, or agreement with
  perturbative coefficients.
\end{enumerate}

\subsection{VII.5. Asymptotic freedom of the running
coupling}\label{vii.5.-asymptotic-freedom-of-the-running-coupling}

\textbf{Part:} VII, identification as Yang--Mills. \textbf{Lean:}
\href{https://github.com/yang-mills-group/OpenAI-YM-20261028/tree/3456e5e75cf8d4cd558955a5c7638fba2d68f464/YM/Pert/Asymptotics}{\texttt{YM/Pert/Asymptotics}}
(the running coupling). Auxiliary declarations are in
\href{https://github.com/yang-mills-group/OpenAI-YM-20261028/tree/3456e5e75cf8d4cd558955a5c7638fba2d68f464/YM/Extra/WEAK_SCALING_ACTLAW}{\texttt{YM/Extra/WEAK\_SCALING\_ACTLAW}}
(files \texttt{Orientation}, \texttt{Criteria}, \texttt{Adjust},
\texttt{Analysis}, \texttt{Station}), from
\href{https://github.com/yang-mills-group/OpenAI-YM-20261028/tree/3456e5e75cf8d4cd558955a5c7638fba2d68f464/YM/TraceExtension/Foundation}{\texttt{YM/TraceExtension/Foundation}}
(files \texttt{Caps}, \texttt{Dilate}), and one definition from
\href{https://github.com/yang-mills-group/OpenAI-YM-20261028/tree/3456e5e75cf8d4cd558955a5c7638fba2d68f464/YM/TraceExtension/Soot}{\texttt{YM/TraceExtension/Soot}}.
\textbf{Size:} 232 lines in 43 declarations: 32 lines in
\texttt{Pert/Asymptotics} (5 declarations), 162 in
\texttt{WEAK\_SCALING\_ACTLAW} (24), 37 in
\texttt{TraceExtension/Foundation} (13), 1 in
\texttt{TraceExtension/Soot} (1).

\subsubsection{1. What this unit proves}\label{what-this-unit-proves-29}

\textbf{(A) The running coupling (5 declarations).} The inverse squared
coupling of the constructed renormalization-group trajectory grows
logarithmically, uniformly in the lattice level and the volume
(VII.5.1). This gives the regularity clause (Reg1) of the main theorem
with slope \(b_*=2b_0\) (VII.5.2). It also gives the law
\(g^2\log(1/a)\to1/(2b_0)\) (VII.5.3).

\textbf{(B) Auxiliary results (38 declarations).}

\begin{itemize}
\tightlist
\item
  VII.5.4: the normalisation of the stress tensor along free-field
  scaling sequences (used by VI.2).
\item
  VII.5.5: exactly compatible frames, built from approximately
  compatible ones (used by VII.1).
\item
  VII.5.6: two elementary convergence lemmas (used by III.12).
\item
  VII.5.7--VII.5.8: the definitions of the classical stress-tensor
  algebra and a read-out predicate (used by VI.1, VII.2, VIII.1).
\end{itemize}

\textbf{Setting for (A).}

\begin{itemize}
\item
  \emph{Gauge group.} \(G\subset U(q)\), \(q\ge1\), is a compact matrix
  group with Lie algebra \(\mathfrak g\) and inner product
  \(\langle X,Y\rangle_{\mathfrak g}=-\frac1q\mathrm{Re}\,\mathrm{tr}(XY)\).
  For an orthonormal basis \((e_i)\), define
  \(C_{\mathfrak g}:=-(\dim\mathfrak g)^{-1}\sum_i\mathrm{tr}_{\mathfrak g}((\mathrm{ad}\,e_i)^2)\),
  with value \(0\) when \(\mathfrak g=0\). Here and below, simple means
  real simple and nonabelian. If \(\mathfrak g\) is simple, its Killing
  form equals
  \(-C_{\mathfrak g}\langle\cdot,\cdot\rangle_{\mathfrak g}\). For
  \(SU(q)\subset U(q)\) with \(q\ge2\), this gives
  \(C_{\mathfrak g}=2q^2\). Put

  \[\kappa_G:=\frac{11\,C_{\mathfrak g}}{24\pi^2},\qquad b_0:=\frac{11\,C_{\mathfrak g}}{48\pi^2}=\tfrac12\kappa_G .\]
\item
  \emph{Scales.} The block factors \(b_k\ge3\) are odd. Set
  \(D_m=\prod_{k<m}b_k\), \(a_m=1/D_m\) and

  \[s_m:=\log(1/a_m)=\sum_{k<m}\log b_k .\]

  Then \(s_m\to\infty\).
\item
  \emph{Trajectory.} For a lattice level \(N\), a volume index \(p\) and
  a stage \(m\le N\), \(\beta_{N,p,m}\) is the coefficient of the
  normalised kinetic term in the effective action at spacing \(a_m\).
  The represented density is obtained from the Wilson weight by \(N-m\)
  block-spin steps.

  \begin{itemize}
  \tightlist
  \item
    \(\beta_{N,p}:=\beta_{N,p,N}\) is the bare coupling of the main
    theorem.
  \item
    \(\beta_{N,p,0}=\bar\beta>0\) for every \(N,p\). This is the
    renormalization condition at unit scale.
  \item
    The reference line is \(\widehat\beta_m:=\bar\beta+\kappa_G s_m\).
    One has \(\beta_{N,p,m}\in[\widehat\beta_m/2,\,2\widehat\beta_m]\).
  \item
    The running coupling is \(g_{N,p,m}:=\beta_{N,p,m}^{-1/2}\).
  \end{itemize}
\end{itemize}

The \textbf{running clause} (R) of the construction conditions reads:
for every \(\varepsilon>0\) there is \(m_0\) such that

\[|\beta_{N,p,m}-\widehat\beta_m|\le\varepsilon\,s_m\quad\text{for all } m\ge m_0,\ N\ge m,\ p\in\mathbb N .\]

\textbf{Proposition VII.5.1 (uniform logarithmic growth of the inverse
squared coupling).} Let the family satisfy its construction conditions
(II.2). For every \(\varepsilon>0\) there is \(m_0\ge1\) such that

\[\Bigl|\frac{\beta_{N,p,m}}{s_m}-2b_0\Bigr|<\varepsilon\qquad\text{for all } m\ge m_0,\ N\ge m,\ p\in\mathbb N .\]

\emph{Status:} theorem. Of the construction conditions only (R) is used.
Simplicity of \(\mathfrak g\) is not needed. \emph{Lean:}
\href{https://github.com/yang-mills-group/OpenAI-YM-20261028/blob/3456e5e75cf8d4cd558955a5c7638fba2d68f464/YM/Pert/Asymptotics/RawAxis.lean\#L15-L30}{\texttt{YM.Pert.Asymptotics.zeroUp\_bound}},
with
\href{https://github.com/yang-mills-group/OpenAI-YM-20261028/blob/3456e5e75cf8d4cd558955a5c7638fba2d68f464/YM/Pert/Asymptotics/RawAxis.lean\#L13}{\texttt{b0}}
and
\href{https://github.com/yang-mills-group/OpenAI-YM-20261028/blob/3456e5e75cf8d4cd558955a5c7638fba2d68f464/YM/Pert/Asymptotics/RawAxis.lean\#L14}{\texttt{b0\_normalizes}}.

\textbf{Corollary VII.5.2 ((Reg1) for the constructed family).} Let
\(\mathfrak g\) be simple and let the family satisfy its construction
conditions. Then its Wilson family satisfies (Reg1) with slope
\(b_*=2b_0=11C_{\mathfrak g}/(24\pi^2)>0\): for every \(\varepsilon>0\)
there is \(N_0\ge1\) with

\[\Bigl|\frac{\beta_{N,p}}{\log(1/a_N)}-b_*\Bigr|<\varepsilon\quad\text{for all } N\ge N_0 \text{ and all } p.\]

\emph{Status:} theorem, proved in VIII.1 as the first clause of the
regularity of the constructed family. It is VII.5.1 at \(m=N\), together
with \(C_{\mathfrak g}>0\) (I.1). This is how the main theorem uses part
(A). The main theorem asserts the existence of a positive slope without
specifying its value. \emph{Lean:}
\href{https://github.com/yang-mills-group/OpenAI-YM-20261028/blob/3456e5e75cf8d4cd558955a5c7638fba2d68f464/HierarchyBridge.lean\#L16-L25}{\texttt{YangMills.regular}}
(running clause of
\href{https://github.com/yang-mills-group/OpenAI-YM-20261028/blob/3456e5e75cf8d4cd558955a5c7638fba2d68f464/ChallengeSpec.lean\#L12-L22}{\texttt{YangMills.WilsonRegularity}}),
via
\href{https://github.com/yang-mills-group/OpenAI-YM-20261028/blob/3456e5e75cf8d4cd558955a5c7638fba2d68f464/YM/Pert/Asymptotics/RawAxis.lean\#L15-L30}{\texttt{zeroUp\_bound}}.

\textbf{Proposition VII.5.3 (asymptotic freedom of the running
coupling).} Let \(\mathfrak g\) be simple and let the family satisfy its
construction conditions. Let \((N_i)\) and \((k_i)\) be sequences of
lattice levels and volume indices with \(i\le N_i\) for all large \(i\).
Then

\[\lim_{i\to\infty}g_{N_i,k_i,i}\,\sqrt{\log(1/a_i)}=\frac1{\sqrt{2b_0}},\qquad\text{equivalently}\qquad g_{N_i,k_i,i}^2\,\log(1/a_i)\longrightarrow\frac1{2b_0}.\]

\begin{itemize}
\tightlist
\item
  For \(N_i=i\) this is the law of the bare coupling,
  \(g_0(a)^2\sim\bigl(2b_0\log(1/a)\bigr)^{-1}\).
\item
  For \(N_i\ge i\) it is the law of the effective coupling at spacing
  \(a_i\) of any finer lattice. No condition is placed on the volumes
  \(k_i\).
\end{itemize}

\emph{Status:} theorem; its square gives the second limit. No other
declaration cites it. \emph{Lean:}
\href{https://github.com/yang-mills-group/OpenAI-YM-20261028/blob/3456e5e75cf8d4cd558955a5c7638fba2d68f464/YM/Pert/Asymptotics/RawAxis.lean\#L36-L48}{\texttt{YM.Pert.Asymptotics.comparison}},
\href{https://github.com/yang-mills-group/OpenAI-YM-20261028/blob/3456e5e75cf8d4cd558955a5c7638fba2d68f464/YM/Pert/Asymptotics/RawAxis.lean\#L35}{\texttt{regulatorRun}}.

\textbf{Setting for (B).}

\begin{itemize}
\tightlist
\item
  \emph{Spaces and maps.}

  \begin{itemize}
  \tightlist
  \item
    \(\mathcal F_d\) is the finite-dimensional complex space of labels
    of degree at most \(d\) in a correlation table \(S\) carrying a
    Euclidean hierarchy \(H\) (IV.1).
  \item
    \(\mathcal F^\circ_d\) is the corresponding space of the free
    Maxwell comparison table, with its hierarchy \(H^\circ\), built on
    the same reference labels (VI.1).
  \item
    A \emph{derivative plan} (IV.1) on either table consists of
    complex-linear inclusions \(\iota_{d\to e}\) (\(d\le e\)) and
    complex-linear partial derivatives
    \(\partial_{d,a}:\mathcal F_d\to\mathcal F_{d+1}\), \(a=1,\dots,4\).
  \item
    The inclusion axiom says that inserting \(\iota_{d\to e}v\) in a
    separated correlation gives the same value as inserting \(v\). It
    implies injectivity, \(\iota_{d\to d}=1\), transitivity and
    preservation of the unit. The derivative axiom identifies insertion
    of \(\partial_{d,a}v\) with differentiation of the correlation in
    that insertion point.
  \item
    The orientation axioms say that inclusions commute with the
    hyperoctahedral group \(\Gamma_4\) of signed permutations (\(384\)
    elements) and with conjugation. Derivatives commute with conjugation
    and transform as vectors under \(\Gamma_4\).
  \end{itemize}
\item
  \emph{Stress data.} A \emph{candidate stress tensor} is a family
  \(J=(J_{ab})\) in \(\mathcal F_4\). Its \emph{local Ward identities}
  (IV.1, VI.1) involve a translation part
  \(p=(p_{d,a}:\mathcal F_d\to\mathcal F_{d+1})\) and a rotation part
  \(Y\).
\item
  \emph{Free Maxwell data} (VI.2, VIII.1): the comparison data specified
  in §2, including a derivative plan \(W\), a rotation action, a tensor
  \(J^\circ\) and \(\kappa\ne0\). The plan satisfies the inclusion,
  derivative and orientation axioms; \(J^\circ\) satisfies the local
  Ward identities with translation part \(\kappa\,\partial^W\) and
  rotation part \(\kappa\) times the infinitesimal rotation action; and
  the family \(\partial^W_4=(\partial^W_{4,a})_a\) is nonzero.
\item
  \emph{Scaling sequences.} A \emph{scaling sequence} (VI.1) consists of
  frames \(U_k(d):\mathcal F^\circ_d\to\mathcal F_d\) (complex-linear
  isomorphisms), positive, strictly decreasing scales \(r_k\to0\) and
  stress normalisations \(\xi_k>0\).
\item
  \emph{Tracking.} A scaling sequence \emph{tracks} \((V,W,J)\) at level
  \((n,d,j,R)\) (VI.1) if three conditions hold:

  \begin{itemize}
  \tightlist
  \item
    the frames are exactly compatible with inclusions, signed
    permutations, conjugation and units up to degree \(\max(d,4)\);
  \item
    for each fixed list of at most \(n\) fields of degrees at most
    \(d\), the correlations transported by \(U_k\) and evaluated at
    \(r_kx\) converge, with all Fréchet derivatives of orders
    \(0,\ldots,j\), uniformly on the window
    \(\{x:\|x\|\le R,\ \|x_a-x_b\|_\infty\ge R^{-1}\ \text{for all }a\ne b\}\),
    to the free correlations;
  \item
    \(\xi_k^{-1}r_k^4\,U_k(4)^{-1}J_{ab}\to J^\circ_{ab}\).
  \end{itemize}
\end{itemize}

\textbf{Proposition VII.5.4 (normalisation of the stress tensor along
free-field scaling sequences).} Fix the two hierarchies and the free
Maxwell comparison data above. Let \(V\) be a derivative plan on \(S\)
satisfying the inclusion and derivative axioms. Let \(J\) be symmetric,
real and \(\Gamma_4\)-covariant, with \(\sum_a\partial^V_{4,a}J_{ab}=0\)
for every \(b\), and let \((J,p,Y)\) satisfy the local Ward identities.

\begin{enumerate}
\def\labelenumi{(\alph{enumi})}
\item
  Suppose \(p_{4,a}=c\,\partial^V_{4,a}\) for \(a=1,\dots,4\) and some
  \(c\in\mathbb R\). Then there are \(D,N,M\in\mathbb N\) and
  \(R_0\ge1\) such that \(c/\xi_k\to\kappa\) for every scaling sequence
  that tracks \((V,W,J)\) at a level \((n,d,j,R)\) with \(n\ge N\),
  \(d\ge D\), \(j\ge M\) and \(R\ge R_0\).
\item
  If \(p\neq0\), then \(J\neq0\).
\end{enumerate}

\emph{Status:} theorems. Part (a) uses A04 and A05 (VI.1); part (b)
requires only the hierarchy and local Ward identities. \emph{Lean:}
\href{https://github.com/yang-mills-group/OpenAI-YM-20261028/blob/3456e5e75cf8d4cd558955a5c7638fba2d68f464/YM/Extra/WEAK_SCALING_ACTLAW/Orientation.lean\#L43-L54}{\texttt{YM.Extra.WeakScaling.tracks\_ratio}},
\href{https://github.com/yang-mills-group/OpenAI-YM-20261028/blob/3456e5e75cf8d4cd558955a5c7638fba2d68f464/YM/Extra/WEAK_SCALING_ACTLAW/Orientation.lean\#L28-L39}{\texttt{scale\_limit}},
\href{https://github.com/yang-mills-group/OpenAI-YM-20261028/blob/3456e5e75cf8d4cd558955a5c7638fba2d68f464/YM/Extra/WEAK_SCALING_ACTLAW/Criteria.lean\#L64-L70}{\texttt{flux\_nontrivial}}.
Used in VI.2 by
\href{https://github.com/yang-mills-group/OpenAI-YM-20261028/blob/3456e5e75cf8d4cd558955a5c7638fba2d68f464/YM/Pert/Weak.lean\#L62-L66}{\texttt{YM.Pert.ratio}}
and
\href{https://github.com/yang-mills-group/OpenAI-YM-20261028/blob/3456e5e75cf8d4cd558955a5c7638fba2d68f464/YM/Pert/Weak.lean\#L82-L116}{\texttt{YM.Pert.centered\_line}}.

\textbf{Proposition VII.5.5 (exactly compatible frames).} Suppose \(S\)
and its free comparison table carry Euclidean hierarchies. Let \(V\) and
\(W\) be derivative plans on these tables satisfying the inclusion and
orientation axioms.

\begin{enumerate}
\def\labelenumi{(\alph{enumi})}
\item
  For every \(d\) there is a \(\mathbb C\)-linear map
  \(P_d:\mathcal F^\circ_{d+1}\to\mathcal F^\circ_d\) with
  \(P_d\,\iota^W_{d\to d+1}=1\) that commutes with \(\Gamma_4\) and with
  conjugation.
\item
  Let the \(P_d\) be as in (a), fix \(D\in\mathbb N\), and let
  \(A_d:\mathcal F^\circ_d\to\mathcal F_d\) be complex-linear
  isomorphisms for every \(d\). Define recursively
\end{enumerate}

\[X_0:=A_0,\qquad X_{d+1}:=\iota^V_{d\to d+1}\,X_d\,P_d+A_{d+1}\bigl(1-\iota^W_{d\to d+1}P_d\bigr).\]

Suppose that for every \(d\le D\) the map \(A_d\) commutes with
\(\Gamma_4\) and conjugation and maps unit to unit. Then:

\begin{itemize}
\tightlist
\item
  \(X_e\,\iota^W_{d\to e}=\iota^V_{d\to e}\,X_d\) for all
  \(d\le e\le D\);
\item
  for \(d\le D\), each \(X_d\) commutes with \(\Gamma_4\) and
  conjugation and maps unit to unit.
\end{itemize}

\begin{enumerate}
\def\labelenumi{(\alph{enumi})}
\setcounter{enumi}{2}
\tightlist
\item
  Let \(A^{(k)}_d\) be complex-linear isomorphisms such that, for every
  \(d<D\) and every \(v\),
\end{enumerate}

\[(A^{(k)}_{d+1})^{-1}\,\iota^V_{d\to d+1}\,A^{(k)}_d\,v\longrightarrow\iota^W_{d\to d+1}v\qquad(k\to\infty).\]

Let \(X^{(k)}\) be built from \(A^{(k)}\) as in (b). Then for every
\(d\le D\):

\begin{itemize}
\tightlist
\item
  \((A^{(k)}_d)^{-1}X^{(k)}_d\to1\) in operator norm;
\item
  \(X^{(k)}_d\) is bijective for all large \(k\);
\item
  the isomorphisms \(\widetilde X^{(k)}_d\) (equal to \(X^{(k)}_d\) when
  it is bijective and to \(A^{(k)}_d\) otherwise) satisfy
  \((A^{(k)}_d)^{-1}\widetilde X^{(k)}_dv\to v\) for every \(v\).
\end{itemize}

If the \(A^{(k)}_d\) also preserve the symmetries and units in (b),
then, for all sufficiently large \(k\), the frames
\(\widetilde X^{(k)}_d\) are exactly compatible simultaneously in all
degrees \(d\le D\).

\emph{Status:} theorem (finite-dimensional algebra and analysis).
\emph{Lean:}
\href{https://github.com/yang-mills-group/OpenAI-YM-20261028/blob/3456e5e75cf8d4cd558955a5c7638fba2d68f464/YM/Extra/WEAK_SCALING_ACTLAW/Adjust.lean\#L35-L55}{\texttt{retraction}},
\href{https://github.com/yang-mills-group/OpenAI-YM-20261028/blob/3456e5e75cf8d4cd558955a5c7638fba2d68f464/YM/Extra/WEAK_SCALING_ACTLAW/Adjust.lean\#L61-L66}{\texttt{adjusted}},
\href{https://github.com/yang-mills-group/OpenAI-YM-20261028/blob/3456e5e75cf8d4cd558955a5c7638fba2d68f464/YM/Extra/WEAK_SCALING_ACTLAW/Adjust.lean\#L72-L76}{\texttt{Exact}},
\href{https://github.com/yang-mills-group/OpenAI-YM-20261028/blob/3456e5e75cf8d4cd558955a5c7638fba2d68f464/YM/Extra/WEAK_SCALING_ACTLAW/Adjust.lean\#L80-L102}{\texttt{adjusted\_exact}},
\href{https://github.com/yang-mills-group/OpenAI-YM-20261028/blob/3456e5e75cf8d4cd558955a5c7638fba2d68f464/YM/Extra/WEAK_SCALING_ACTLAW/Adjust.lean\#L110-L111}{\texttt{Near}},
\href{https://github.com/yang-mills-group/OpenAI-YM-20261028/blob/3456e5e75cf8d4cd558955a5c7638fba2d68f464/YM/Extra/WEAK_SCALING_ACTLAW/Adjust.lean\#L113-L129}{\texttt{perturb\_limit}},
\href{https://github.com/yang-mills-group/OpenAI-YM-20261028/blob/3456e5e75cf8d4cd558955a5c7638fba2d68f464/YM/Extra/WEAK_SCALING_ACTLAW/Adjust.lean\#L131-L135}{\texttt{eventually\_unit}},
\href{https://github.com/yang-mills-group/OpenAI-YM-20261028/blob/3456e5e75cf8d4cd558955a5c7638fba2d68f464/YM/Extra/WEAK_SCALING_ACTLAW/Adjust.lean\#L137-L140}{\texttt{rotateChart}},
\href{https://github.com/yang-mills-group/OpenAI-YM-20261028/blob/3456e5e75cf8d4cd558955a5c7638fba2d68f464/YM/Extra/WEAK_SCALING_ACTLAW/Adjust.lean\#L144-L148}{\texttt{relative}},
\href{https://github.com/yang-mills-group/OpenAI-YM-20261028/blob/3456e5e75cf8d4cd558955a5c7638fba2d68f464/YM/Extra/WEAK_SCALING_ACTLAW/Adjust.lean\#L171-L172}{\texttt{Graded}}.
Used in VII.1 by
\href{https://github.com/yang-mills-group/OpenAI-YM-20261028/blob/3456e5e75cf8d4cd558955a5c7638fba2d68f464/YM/SharpGram/Weak/Compatible.lean\#L140-L159}{\texttt{YM.SharpGram.Weak.Portfolio.tune\_grades}}.
Exact compatibility is also used in III.13 and VIII.1.

\textbf{Proposition VII.5.6 (finite-dimensional convergence).}

\begin{enumerate}
\def\labelenumi{(\alph{enumi})}
\item
  Let \(f:E\to X\) be an injective linear map between finite-dimensional
  real normed spaces. If \(f(a_y)\to f(e)\) along a filter, then
  \(a_y\to e\).
\item
  Let \(f_y,g:(\mathbb R^4)^n\to\mathbb C\) be continuous real-linear
  maps, and fix a filter on the index set of \(y\). If
  \(f_y(e_{i,a})\to g(e_{i,a})\) at each of the \(4n\) coordinate
  vectors, then \(f_y\to g\) in operator norm along that filter.
\end{enumerate}

\emph{Status:} theorem. \emph{Lean:}
\href{https://github.com/yang-mills-group/OpenAI-YM-20261028/blob/3456e5e75cf8d4cd558955a5c7638fba2d68f464/YM/Extra/WEAK_SCALING_ACTLAW/Analysis.lean\#L44-L51}{\texttt{norm\_fit}},
\href{https://github.com/yang-mills-group/OpenAI-YM-20261028/blob/3456e5e75cf8d4cd558955a5c7638fba2d68f464/YM/Extra/WEAK_SCALING_ACTLAW/Analysis.lean\#L63-L70}{\texttt{norm\_fit\_deriv}},
\href{https://github.com/yang-mills-group/OpenAI-YM-20261028/blob/3456e5e75cf8d4cd558955a5c7638fba2d68f464/YM/Extra/WEAK_SCALING_ACTLAW/Analysis.lean\#L55-L60}{\texttt{components}}.
Used in III.12 by
\href{https://github.com/yang-mills-group/OpenAI-YM-20261028/blob/3456e5e75cf8d4cd558955a5c7638fba2d68f464/YM/FullVolume/Region.lean\#L28-L30}{\texttt{YM.FullVolume.stretching}}.

\textbf{Definition VII.5.7 (the classical stress-tensor algebra).}
Labels of degree \(d\) are gauge-invariant polynomials of weight at most
\(d\) in radial Taylor coefficients at the base point. In the formulas
below, \(F\) and \(\partial F\) denote the linear curls of the degree-0
and degree-1 Taylor coefficients used in the free Maxwell comparison.
The curvature has weight 2 and each derivative adds 1 (II.1).
\(\langle\cdot,\cdot\rangle\) denotes the invariant inner product on
\(\mathfrak g\). These are polynomial identities, before passing to the
quotient by null fields; they do not multiply limiting fields at
coincident points.

\begin{itemize}
\item
  \emph{Weight 4.} The square is
  \(\mathfrak s:=\frac14\sum_{i,j}\langle F_{ij},F_{ij}\rangle\), summed
  over ordered pairs. The candidate stress tensors of VI.1 form the
  family

  \[T^{(t)}_{ab}=\delta_{ab}\Bigl(t_0\mathbf 1+t_1\sum_i\langle F_{ai},F_{ai}\rangle+t_2\sum_{i,j}\langle F_{ij},F_{ij}\rangle\Bigr)+(1-\delta_{ab})\,t_3\sum_i\langle F_{ai},F_{bi}\rangle .\]

  The \emph{canonical} tensor is \(T^{\mathrm{can}}:=T^{(0,1,-1/4,1)}\),
  that is,
  \(T^{\mathrm{can}}_{ab}=\sum_i\langle F_{ai},F_{bi}\rangle-\delta_{ab}\,\mathfrak s\).
\item
  \emph{Weight 5.} Define the vectors

  \[\mu^1_b:=\sum_i\langle F_{bi},\partial_bF_{bi}\rangle=\tfrac12\partial_b\sum_i|F_{bi}|^2,\qquad \mu^2_b:=\sum_{i,j}\langle F_{ij},\partial_bF_{ij}\rangle=\tfrac12\partial_b\sum_{i,j}|F_{ij}|^2,\]

  and the equation-of-motion vector
  \(E_b:=\sum_{a,i}\langle F_{bi},\partial_aF_{ai}\rangle\) (VI.1). Let
  \(\mathcal R_1\) be the real space of real, \(\Gamma_4\)-covariant
  vectors of degree at most 5 (VI.1). Define
  \(\beta,\delta:\mathbb R^3\to\mathcal R_1\) by

  \[\beta(w):=w_0\mu^1+w_1\mu^2+w_2E,\qquad \delta(w):=\beta\bigl(2(w_0-w_2),\ 2w_1+\tfrac12w_2,\ w_2\bigr).\]
\end{itemize}

The definitions have two accompanying lemmas:

\begin{itemize}
\tightlist
\item
  \(\mu^1,\mu^2\in\mathcal R_1\);
\item
  every degree-0 radial jet is realised by a constant-curvature free
  Maxwell field, with an affine potential.
\end{itemize}

\emph{Status:} definitions and two lemmas. \emph{Lean:}
\href{https://github.com/yang-mills-group/OpenAI-YM-20261028/blob/3456e5e75cf8d4cd558955a5c7638fba2d68f464/YM/TraceExtension/Foundation/Dilate.lean\#L66}{\texttt{square4}},
\href{https://github.com/yang-mills-group/OpenAI-YM-20261028/blob/3456e5e75cf8d4cd558955a5c7638fba2d68f464/YM/TraceExtension/Foundation/Dilate.lean\#L70}{\texttt{raw}},
\href{https://github.com/yang-mills-group/OpenAI-YM-20261028/blob/3456e5e75cf8d4cd558955a5c7638fba2d68f464/YM/TraceExtension/Foundation/Caps.lean\#L16}{\texttt{momentum1}},
\href{https://github.com/yang-mills-group/OpenAI-YM-20261028/blob/3456e5e75cf8d4cd558955a5c7638fba2d68f464/YM/TraceExtension/Foundation/Caps.lean\#L14}{\texttt{momentum2}},
\href{https://github.com/yang-mills-group/OpenAI-YM-20261028/blob/3456e5e75cf8d4cd558955a5c7638fba2d68f464/YM/TraceExtension/Foundation/Caps.lean\#L28-L34}{\texttt{momentum1\_mem}},
\href{https://github.com/yang-mills-group/OpenAI-YM-20261028/blob/3456e5e75cf8d4cd558955a5c7638fba2d68f464/YM/TraceExtension/Foundation/Caps.lean\#L25-L27}{\texttt{momentum2\_mem}},
\href{https://github.com/yang-mills-group/OpenAI-YM-20261028/blob/3456e5e75cf8d4cd558955a5c7638fba2d68f464/YM/TraceExtension/Foundation/Caps.lean\#L21-L24}{\texttt{momentum2\_turn}},
\href{https://github.com/yang-mills-group/OpenAI-YM-20261028/blob/3456e5e75cf8d4cd558955a5c7638fba2d68f464/YM/TraceExtension/Foundation/Caps.lean\#L63-L66}{\texttt{basisV}},
\href{https://github.com/yang-mills-group/OpenAI-YM-20261028/blob/3456e5e75cf8d4cd558955a5c7638fba2d68f464/YM/TraceExtension/Foundation/Dilate.lean\#L33}{\texttt{bareVector}}
(with
\href{https://github.com/yang-mills-group/OpenAI-YM-20261028/blob/3456e5e75cf8d4cd558955a5c7638fba2d68f464/YM/TraceExtension/Foundation/Dilate.lean\#L32}{\texttt{mat3}},
\href{https://github.com/yang-mills-group/OpenAI-YM-20261028/blob/3456e5e75cf8d4cd558955a5c7638fba2d68f464/YM/TraceExtension/Foundation/Dilate.lean\#L31}{\texttt{proj3}}),
\href{https://github.com/yang-mills-group/OpenAI-YM-20261028/blob/3456e5e75cf8d4cd558955a5c7638fba2d68f464/YM/TraceExtension/Foundation/Caps.lean\#L43-L50}{\texttt{shell\_affine}},
\href{https://github.com/yang-mills-group/OpenAI-YM-20261028/blob/3456e5e75cf8d4cd558955a5c7638fba2d68f464/YM/TraceExtension/Foundation/Caps.lean\#L17-L20}{\texttt{sumPerm}}.
The facts proved about these objects in VI.1 and VII.2 are listed in
§4.5.

\textbf{Definition VII.5.8 (sharp read-out bound).} Fix the family, a
nondecreasing control function \(K:\mathbb N\to\mathbb R\) with
\(K(j)\ge1\) for every \(j\), a degree \(d\), stages
\(\ell_d\le k\le N\), a volume index \(p\) and \(u\in\mathbb R\). Write
\(\|c\|_1:=\sum_\alpha|c_\alpha|\) for the coefficient norm, and let
\(P\) remove the unit coordinate. Here \(\ell_d\) is the reference stage
fixed by the family. The predicate quantifies over every \(w\ge0\) and
all insertion presentations \(A_i\) with coefficient vectors \(c_i\),
remainder costs \(q_i\), inserted densities \(\mathcal D_i\) and
insertion points at the origin. They must satisfy:

\begin{itemize}
\item
  At each \(\ell_d\le i\le k\), \(A_i\) has the admissible regular and
  forced remainders of III.12, with total remainder cost at most
  \(q_i\), using the square root of the stage-\(i\) forced allowance.
\item
  \(c_{\ell_d}=0\). Only this coefficient vector must vanish at the
  reference stage.
\item
  The full inserted densities obey
  \(\mathcal D_i=b_i^d\,\mathcal T^{(i)}_{\beta_{N,p,i+1}}\mathcal D_{i+1}\)
  for \(\ell_d\le i<k\).
\item
  Put \((S_i c)_\alpha:=b_i^{d-d_\alpha}c_\alpha\), where \(d_\alpha\)
  is the weight of the basis label, and \(K_i:=K(J(i+1))\). For each
  \(\ell_d\le i<k\), \(q_i,q_{i+1}\ge0\) and there is a complex-linear
  automorphism \(B_i\) commuting with \(P\) such that
  \(\|B_i c\|_1,\|B_i^{-1}c\|_1\le K_i\|c\|_1\) for every \(c\), and

  \[\|c_i-B_iS_ic_{i+1}\|_1\le b_i^{-3/8}(\|Pc_{i+1}\|_1+q_{i+1}),\qquad q_i\le K_i\|PS_ic_{i+1}\|_1+b_i^{-3/8}q_{i+1}.\]
\item
  The terminal cost satisfies \(q_k\le w+n_k(w+\|c_k\|_1)\), where
  \(w\ge0\) and \(n_k\) is the stage-\(k\) column number.
\end{itemize}

The sharp read-out predicate is the assertion that every such chain
satisfies

\[u\,\|c_k\|_1\le4(1+n_k)w.\]

\emph{Status:} definition. The admissibility and remainder conditions
are those of
\href{https://github.com/yang-mills-group/OpenAI-YM-20261028/blob/3456e5e75cf8d4cd558955a5c7638fba2d68f464/YM/Source/Insertion.lean\#L52-L61}{\texttt{Format}},
\href{https://github.com/yang-mills-group/OpenAI-YM-20261028/blob/3456e5e75cf8d4cd558955a5c7638fba2d68f464/YM/Source/Insertion.lean\#L63-L68}{\texttt{remBound}};
the complete chain conditions are
\href{https://github.com/yang-mills-group/OpenAI-YM-20261028/blob/3456e5e75cf8d4cd558955a5c7638fba2d68f464/YM/Bid/Loan.lean\#L115-L127}{\texttt{Duet.Woke}}
and
\href{https://github.com/yang-mills-group/OpenAI-YM-20261028/blob/3456e5e75cf8d4cd558955a5c7638fba2d68f464/YM/Bid/Loan.lean\#L65-L69}{\texttt{direction}}.
\emph{Lean:}
\href{https://github.com/yang-mills-group/OpenAI-YM-20261028/blob/3456e5e75cf8d4cd558955a5c7638fba2d68f464/YM/TraceExtension/Soot/Skein.lean\#L26}{\texttt{YM.MergeProof.Accent.SharpRead}}.
Used by:
\href{https://github.com/yang-mills-group/OpenAI-YM-20261028/blob/3456e5e75cf8d4cd558955a5c7638fba2d68f464/YM/TraceExtension/Terrace/Knell.lean\#L62-L66}{\texttt{YM.Trace.Terrace.Sceptre}}
and the cost lemmas of VII.2.

\subsubsection{2. What it uses}\label{what-it-uses-29}

\textbf{Inputs from other units.}

\begin{itemize}
\tightlist
\item
  \emph{I.1 (gauge group and scales).}

  \begin{itemize}
  \tightlist
  \item
    \(C_{\mathfrak g}\)
    (\href{https://github.com/yang-mills-group/OpenAI-YM-20261028/blob/3456e5e75cf8d4cd558955a5c7638fba2d68f464/YM/Gauge/Colors.lean\#L180-L182}{\texttt{YM.Gauge.casimir}})
    and \(\kappa_G\), which is \emph{defined} as
    \(11C_{\mathfrak g}/(24\pi^2)\)
    (\href{https://github.com/yang-mills-group/OpenAI-YM-20261028/blob/3456e5e75cf8d4cd558955a5c7638fba2d68f464/YM/Gauge/Colors.lean\#L184}{\texttt{slopeConstant}}).
  \item
    For simple \(\mathfrak g\), the Killing-form identity in §1 is
    \href{https://github.com/yang-mills-group/OpenAI-YM-20261028/blob/3456e5e75cf8d4cd558955a5c7638fba2d68f464/YM/Gauge/Casimir.lean\#L140-L158}{\texttt{trace\_adjoint}}.
  \item
    Positivity \(C_{\mathfrak g}>0\) and \(\kappa_G>0\) for simple
    \(\mathfrak g\)
    (\href{https://github.com/yang-mills-group/OpenAI-YM-20261028/blob/3456e5e75cf8d4cd558955a5c7638fba2d68f464/YM/Gauge/Casimir.lean\#L160-L172}{\texttt{casimir\_pos}},
    \href{https://github.com/yang-mills-group/OpenAI-YM-20261028/blob/3456e5e75cf8d4cd558955a5c7638fba2d68f464/YM/Gauge/Casimir.lean\#L174-L176}{\texttt{slopeConstant\_pos}}).
  \item
    \(a_m=1/D_m\)
    (\href{https://github.com/yang-mills-group/OpenAI-YM-20261028/blob/3456e5e75cf8d4cd558955a5c7638fba2d68f464/YM/Lattice/Scale.lean\#L99}{\texttt{Schedule.a}});
    \(s_m\) is \emph{defined} as \(\log(1/a_m)\)
    (\href{https://github.com/yang-mills-group/OpenAI-YM-20261028/blob/3456e5e75cf8d4cd558955a5c7638fba2d68f464/YM/Lattice/Scale.lean\#L115}{\texttt{Schedule.scale}});
    and \(s_m\to\infty\)
    (\href{https://github.com/yang-mills-group/OpenAI-YM-20261028/blob/3456e5e75cf8d4cd558955a5c7638fba2d68f464/YM/Lattice/Scale.lean\#L117-L120}{\texttt{Schedule.divergent}}).
  \end{itemize}
\item
  \emph{II.2 (construction conditions)}
  \href{https://github.com/yang-mills-group/OpenAI-YM-20261028/blob/3456e5e75cf8d4cd558955a5c7638fba2d68f464/YM/Source/Construction.lean\#L44-L62}{\texttt{YM.Family.Makes}}.
  Its first field says that the trajectory is a selected history
  (III.1). The bare coupling is
  \href{https://github.com/yang-mills-group/OpenAI-YM-20261028/blob/3456e5e75cf8d4cd558955a5c7638fba2d68f464/YM/Source/Construction.lean\#L38}{\texttt{Family.beta}},
  that is, \(\beta_{N,p,N}\).
\item
  \emph{III.1 (selected history)}
  \href{https://github.com/yang-mills-group/OpenAI-YM-20261028/blob/3456e5e75cf8d4cd558955a5c7638fba2d68f464/YM/RG/History.lean\#L234-L240}{\texttt{YM.RG.History.Selected}}.
  It contains:

  \begin{itemize}
  \tightlist
  \item
    \(\beta_{N,p,0}=\bar\beta\);
  \item
    the per-step relation
    \(|\beta_{N,p,m+1}-\beta_{N,p,m}-\kappa_G\log b_m|\le\delta_{m+1}\)
    (\href{https://github.com/yang-mills-group/OpenAI-YM-20261028/blob/3456e5e75cf8d4cd558955a5c7638fba2d68f464/YM/RG/History.lean\#L223-L227}{\texttt{Transfer}});
  \item
    the band \([\widehat\beta_m/2,2\widehat\beta_m]\)
    (\href{https://github.com/yang-mills-group/OpenAI-YM-20261028/blob/3456e5e75cf8d4cd558955a5c7638fba2d68f464/YM/RG/History.lean\#L213-L218}{\texttt{Conv}});
  \item
    the running clause (R)
    (\href{https://github.com/yang-mills-group/OpenAI-YM-20261028/blob/3456e5e75cf8d4cd558955a5c7638fba2d68f464/YM/RG/History.lean\#L229-L230}{\texttt{Running}}),
    with \(\widehat\beta_m\) given by
    \href{https://github.com/yang-mills-group/OpenAI-YM-20261028/blob/3456e5e75cf8d4cd558955a5c7638fba2d68f464/YM/RG/History.lean\#L79}{\texttt{hatBeta}}.
  \end{itemize}
\item
  \emph{III.2 (choice of constants, shooting).}

  \begin{itemize}
  \tightlist
  \item
    (R) is derived in
    \href{https://github.com/yang-mills-group/OpenAI-YM-20261028/blob/3456e5e75cf8d4cd558955a5c7638fba2d68f464/YM/Shoot/Race.lean\#L23-L53}{\texttt{YM.RG.Shoot.trajectory}}.
    The inputs are \(\beta_{N,p,0}=\bar\beta\), the per-step relation,
    and \(\delta_{m+1}/\log b_m\le J(m+1)^{-2}\), where
    \(J(m)\to\infty\) is the stage index of the scale sequence (Lemma
    VII.5.b).
  \item
    The history with \(\beta_{N,p,0}=\bar\beta\) is produced by shooting
    (\href{https://github.com/yang-mills-group/OpenAI-YM-20261028/blob/3456e5e75cf8d4cd558955a5c7638fba2d68f464/YM/Shoot/Race.lean\#L67-L82}{\texttt{Recipe.races}},
    \href{https://github.com/yang-mills-group/OpenAI-YM-20261028/blob/3456e5e75cf8d4cd558955a5c7638fba2d68f464/YM/Shoot/Reach.lean\#L156-L159}{\texttt{successful}}).
  \item
    The per-step relation follows from the coupling reset (III.9). If
    the one-loop shift \(\alpha_{N,m}\) satisfies
    \(|\alpha_{N,m}+\kappa_G\log b_m|+1\le\delta_{m+1}\) and the reset
    changes \(\beta\) by \(\alpha_{N,m}\) with error at most \(1\), the
    triangle inequality gives the transfer bound. This is
    \href{https://github.com/yang-mills-group/OpenAI-YM-20261028/blob/3456e5e75cf8d4cd558955a5c7638fba2d68f464/YM/PCalc/Band.lean\#L57-L63}{\texttt{Directed.links}},
    using
    \href{https://github.com/yang-mills-group/OpenAI-YM-20261028/blob/3456e5e75cf8d4cd558955a5c7638fba2d68f464/YM/PCalc/Loop.lean\#L93-L115}{\texttt{Kin.Planned}}.
  \end{itemize}
\item
  \emph{VIII.1:} uses VII.5.2 in the assembly of the main theorem.
\item
  \emph{For VII.5.4.}

  \begin{itemize}
  \tightlist
  \item
    The work items A04
    (\href{https://github.com/yang-mills-group/OpenAI-YM-20261028/blob/3456e5e75cf8d4cd558955a5c7638fba2d68f464/YM/Work/A04.lean\#L163-L183}{\texttt{answer}})
    and A05
    (\href{https://github.com/yang-mills-group/OpenAI-YM-20261028/blob/3456e5e75cf8d4cd558955a5c7638fba2d68f464/YM/Work/A05.lean\#L95-L121}{\texttt{answer}}),
    both proved in VI.1. Along tracking scaling sequences of sufficient
    arity, degree, derivative order and window size, they give:
    \(r_k\,\partial^V_d\) conjugated by the frames converges to
    \(\partial^W_d\); and \((r_k/\xi_k)\,p_d\) conjugated by the frames
    converges to \(\kappa\,\partial^W_d\).
  \item
    The tracking notion
    (\href{https://github.com/yang-mills-group/OpenAI-YM-20261028/blob/3456e5e75cf8d4cd558955a5c7638fba2d68f464/YM/Climb/Beams.lean\#L44-L48}{\texttt{Track}})
    and the hypotheses on \(J\)
    (\href{https://github.com/yang-mills-group/OpenAI-YM-20261028/blob/3456e5e75cf8d4cd558955a5c7638fba2d68f464/YM/Climb/Beams.lean\#L66-L69}{\texttt{Pre}}),
    from VI.1.
  \item
    The Ward identities
    (\href{https://github.com/yang-mills-group/OpenAI-YM-20261028/blob/3456e5e75cf8d4cd558955a5c7638fba2d68f464/YM/Euclid/Stress.lean\#L90-L97}{\texttt{StressLaws}})
    and derivative plans
    (\href{https://github.com/yang-mills-group/OpenAI-YM-20261028/blob/3456e5e75cf8d4cd558955a5c7638fba2d68f464/YM/Euclid/Differential.lean\#L20-L22}{\texttt{DiffPlan}}),
    from IV.1.
  \item
    The free Maxwell data
    (\href{https://github.com/yang-mills-group/OpenAI-YM-20261028/blob/3456e5e75cf8d4cd558955a5c7638fba2d68f464/YM/Org/Input.lean\#L159-L170}{\texttt{MaxwellData}},
    \href{https://github.com/yang-mills-group/OpenAI-YM-20261028/blob/3456e5e75cf8d4cd558955a5c7638fba2d68f464/YM/Climb/Beams.lean\#L34-L42}{\texttt{Lighthouse.Goes}}).
    Besides the data in §1, the full package includes finite separating
    tests and the boundary, polynomial and internal-rotation conditions
    of VI.2. The ratio argument uses the free hierarchy, derivative and
    Ward identities, and the nonzero degree-4 derivative.
  \end{itemize}
\item
  \emph{For VII.5.5:} the inclusion and orientation axioms of derivative
  plans (IV.1), and the free comparison table
  (\href{https://github.com/yang-mills-group/OpenAI-YM-20261028/blob/3456e5e75cf8d4cd558955a5c7638fba2d68f464/YM/Spin/Sheets.lean\#L38}{\texttt{Spin.comparator}},
  VI.1).
\item
  \emph{For VII.5.7:}

  \begin{itemize}
  \tightlist
  \item
    the label polynomials \(\langle F_{ab},F_{cd}\rangle\)
    (\href{https://github.com/yang-mills-group/OpenAI-YM-20261028/blob/3456e5e75cf8d4cd558955a5c7638fba2d68f464/YM/Push/STRESS_FLUX/Polynomial.lean\#L50-L51}{\texttt{pair}})
    and \(\langle F_{ab},\partial_iF_{cd}\rangle\)
    (\href{https://github.com/yang-mills-group/OpenAI-YM-20261028/blob/3456e5e75cf8d4cd558955a5c7638fba2d68f464/YM/Push/STRESS_FLUX/Polynomial.lean\#L57-L58}{\texttt{pairDeriv}});
  \item
    \(E\)
    (\href{https://github.com/yang-mills-group/OpenAI-YM-20261028/blob/3456e5e75cf8d4cd558955a5c7638fba2d68f464/YM/Push/STRESS_FLUX/Polynomial.lean\#L105}{\texttt{eom}}),
    the family \(T^{(t)}\)
    (\href{https://github.com/yang-mills-group/OpenAI-YM-20261028/blob/3456e5e75cf8d4cd558955a5c7638fba2d68f464/YM/Push/STRESS_FLUX/Polynomial.lean\#L177-L181}{\texttt{stress}},
    \href{https://github.com/yang-mills-group/OpenAI-YM-20261028/blob/3456e5e75cf8d4cd558955a5c7638fba2d68f464/YM/Push/STRESS_FLUX/Polynomial.lean\#L187}{\texttt{weyl}})
    and \(\mathcal R_1\)
    (\href{https://github.com/yang-mills-group/OpenAI-YM-20261028/blob/3456e5e75cf8d4cd558955a5c7638fba2d68f464/YM/Push/STRESS_FLUX/Counts.lean\#L53}{\texttt{RawOne}}),
    all from VI.1;
  \item
    the ideal of labels vanishing on free Maxwell solutions
    (\href{https://github.com/yang-mills-group/OpenAI-YM-20261028/blob/3456e5e75cf8d4cd558955a5c7638fba2d68f464/YM/Maxwell/Equation.lean\#L72}{\texttt{Maxwell.equationIdeal}}),
    from III.16.
  \end{itemize}
\item
  \emph{For VII.5.8:} the comparison-chain definitions of III.2 and the
  insertion estimates of III.12. For the constructed family,
  \href{https://github.com/yang-mills-group/OpenAI-YM-20261028/blob/3456e5e75cf8d4cd558955a5c7638fba2d68f464/YM/TraceExtension/Terrace/Earn.lean\#L22-L52}{\texttt{attain}}
  (III.2) supplies the read-out field of the trace-analysis assumptions
  under its scale and one-step hypotheses.
\end{itemize}

\textbf{Standard mathematics.}

\begin{itemize}
\tightlist
\item
  \emph{Asymptotic freedom.} Gross--Wilczek and Politzer (Phys.
  Rev.~Lett. 30, 1973). In the conventional normalisation for \(SU(q)\),
  with \(\beta=2q/g^2\), the one-loop law is
  \(1/g^2\sim2b_0^{\mathrm{conv}}\log(1/a)\) with
  \(b_0^{\mathrm{conv}}=11q/(48\pi^2)\).

  \begin{itemize}
  \tightlist
  \item
    Here \(\beta=1/g_0^2\) in the normalisation
    \(\langle X,Y\rangle_{\mathfrak g}=-\frac1q\mathrm{Re}\,\mathrm{tr}(XY)\),
    and \(C_{\mathfrak g}=2q^2\). So
    \(b_*=2b_0=11q^2/(12\pi^2)=2q\cdot2b_0^{\mathrm{conv}}\), which
    agrees with the conventional law.
  \item
    For \(SU(2)\), \(b_*=11/(3\pi^2)\).
  \item
    For the relation between the bare coupling and the lattice spacing
    see Creutz, \emph{Quarks, Gluons and Lattices} (1983).
  \end{itemize}
\item
  \emph{Block-spin RG for lattice gauge fields.} Balaban's
  renormalization group (Comm. Math. Phys. 109 (1987), 116 (1988), and
  related papers) provides the background for the marginal-flow
  estimate. The specific input here is the per-step bound of III.9; (R)
  follows by summation.
\item
  \emph{Exactly compatible frames.} Averaging over a finite group gives
  equivariant projections (Maschke). In finite dimension, injective
  linear maps form an open set, and linear maps are continuous.
\item
  \emph{Classical stress tensor.} The symmetric gauge-field stress
  tensor \(T^{\mathrm{can}}\) is conserved on the equations of motion
  and is traceless in four dimensions (classical scale invariance;
  Callan--Coleman--Jackiw, Ann. Phys. 59, 1970). The breaking of
  tracelessness in the quantum theory is the trace anomaly
  (Adler--Collins--Duncan 1977; Collins--Duncan--Joglekar 1977), treated
  in VII.2.
\end{itemize}

\textbf{Imported constants and assumptions.}

\begin{itemize}
\tightlist
\item
  The value \(\kappa_G=11C_{\mathfrak g}/(24\pi^2)\) is a definition
  (I.1), and \(b_0:=\kappa_G/2\), so \(\kappa_G=2b_0\) holds by
  unfolding. That \(\kappa_G\) is the actual slope of the trajectory is
  the content of the per-step relation, proved in Part III.
\item
  \(\bar\beta>0\) is the fixed reference coupling (III.2). The limiting
  slope is independent of \(\bar\beta\); the convergence thresholds may
  depend on it.
\item
  VII.5.1 assumes the construction conditions. VII.5.2--VII.5.3
  additionally assume simplicity. The auxiliary hypotheses are stated in
  §1.
\end{itemize}

\subsubsection{3. Strategy}\label{strategy-29}

\textbf{Part (A).} The coupling law follows from the one-step bound by
summation and division by the logarithmic scale.

\begin{enumerate}
\def\labelenumi{\arabic{enumi}.}
\item
  \emph{Marginal flow (III.9).} Blocking from stage \(m+1\) to stage
  \(m\) gives
  \(\beta_{N,p,m}=\beta_{N,p,m+1}-\kappa_G\log b_m+e_{N,p,m}\), with
  \(|e_{N,p,m}|\le\delta_{m+1}\) and
  \(\delta_{m+1}/\log b_m\le J(m+1)^{-2}\to0\). The reset uses the
  one-loop estimate of III.7.
\item
  \emph{Shooting and summation (III.2).} For each \(N,p\) the bare
  coupling \(\beta_{N,p,N}\) is chosen so that the flow reaches
  \(\beta_{N,p,0}=\bar\beta\); this is the renormalization condition at
  unit scale. Summing the per-step relation from step 0 to step \(m\)
  gives

  \[|\beta_{N,p,m}-\widehat\beta_m|\le\sum_{i<m}\delta_{i+1}.\]

  Since \(\delta_{i+1}=o(\log b_i)\) and
  \(s_m=\sum_{i<m}\log b_i\to\infty\), the right side is \(o(s_m)\) by a
  Cesàro argument. The bound is uniform in \(N\ge m\) and in \(p\),
  because the per-step bounds are.
\item
  \emph{Division by \(s_m\) (this unit).} Since \(\bar\beta/s_m\to0\),
  one gets \(\beta_{N,p,m}/s_m\to\kappa_G=2b_0\), uniformly. Taking
  \(m=N\) gives (Reg1); taking square roots gives the coupling law.
\end{enumerate}

The result concerns the actual inverse coupling of the trajectory,
including its bare endpoint. It does not identify this quantity with a
coupling defined from continuum correlations.

\textbf{Part (B).} The auxiliary arguments use finite-dimensional limits
and equivariant linear algebra.

\begin{itemize}
\tightlist
\item
  \emph{VII.5.4} compares two convergent sequences in a
  finite-dimensional space. Along a free-field scaling sequence, the
  rescaled derivative \(r_k\partial^V_4\) tends to
  \(\partial^W_4\neq0\). The rescaled flux
  \((r_k/\xi_k)p_4=(c/\xi_k)\,r_k\partial^V_4\) tends to
  \(\kappa\,\partial^W_4\). Hence the scalar \(c/\xi_k\) tends to
  \(\kappa\).
\item
  \emph{VII.5.5} uses equivariant averaging and a recursive correction:

  \begin{enumerate}
  \def\labelenumi{\arabic{enumi}.}
  \tightlist
  \item
    an equivariant retraction of each inclusion, obtained by averaging
    over \(\Gamma_4\) and conjugation;
  \item
    a recursive correction, degree by degree, that turns approximately
    compatible frames into exactly compatible ones;
  \item
    a perturbation argument showing the correction is eventually
    invertible and asymptotically equal to the original frames.
  \end{enumerate}
\item
  \emph{VII.5.7} is the algebra of stress tensors quadratic in \(F\)
  under the hyperoctahedral group. On the lattice only \(\Gamma_4\)
  acts, so the candidate family has four parameters and its divergences
  have three coordinates. For simple \(\mathfrak g\), conservation on
  free Maxwell fields singles out the canonical tensor up to scale and
  addition of a constant multiple of \(\delta_{ab}\mathbf1\). The
  canonical tensor is traceless.
\end{itemize}

\subsubsection{4. Steps}\label{steps-29}

\paragraph{4.1 The running coupling}\label{the-running-coupling}

\textbf{Lemma VII.5.a (normalisation).} \(2b_0=\kappa_G\). \emph{Proof.}
Unfold both definitions. \emph{Lean:}
\href{https://github.com/yang-mills-group/OpenAI-YM-20261028/blob/3456e5e75cf8d4cd558955a5c7638fba2d68f464/YM/Pert/Asymptotics/RawAxis.lean\#L14}{\texttt{b0\_normalizes}}.

\textbf{Lemma VII.5.b (the running clause; input from III.2).} Let
\(\delta:\mathbb N\to[0,\infty)\) and \(J:\mathbb N\to\mathbb N\)
satisfy \(J(m)\ge1\) and \(J(m)\to\infty\). With the scales and
reference line of §1, suppose:

\begin{itemize}
\tightlist
\item
  \(\beta_{N,p,0}=\bar\beta\) for all \(N,p\);
\item
  \(|\beta_{N,p,m+1}-\beta_{N,p,m}-\kappa_G\log b_m|\le\delta_{m+1}\)
  for \(m<N\);
\item
  \(\delta_{m+1}/\log b_m\le J(m+1)^{-2}\) for every \(m\).
\end{itemize}

Then (R) holds.

\emph{Proof sketch.}

\begin{enumerate}
\def\labelenumi{\arabic{enumi}.}
\tightlist
\item
  Since \(s_{m+1}=s_m+\log b_m\), induction on \(m\) gives
  \(|\beta_{N,p,m}-\widehat\beta_m|\le\Delta_m:=\sum_{i<m}\delta_{i+1}\)
  for \(m\le N\).
\item
  Given \(\varepsilon\), choose \(j\) with
  \(J(i+1)^{-2}\le\varepsilon/2\) for \(i\ge j\). Then
  \(\Delta_m\le\Delta_j+\frac\varepsilon2(s_m-s_j)\) for \(m\ge j\).
\item
  For \(m\) large, \(\Delta_j\le\frac\varepsilon2 s_m\), hence
  \(\Delta_m\le\varepsilon s_m\).
\end{enumerate}

The constants and thresholds do not involve \(N\) or \(p\). \emph{Lean:}
\href{https://github.com/yang-mills-group/OpenAI-YM-20261028/blob/3456e5e75cf8d4cd558955a5c7638fba2d68f464/YM/Shoot/Race.lean\#L23-L53}{\texttt{YM.RG.Shoot.trajectory}}
(III.2).

\textbf{Lemma VII.5.c (= Proposition VII.5.1).} \emph{Proof sketch.}

\begin{enumerate}
\def\labelenumi{\arabic{enumi}.}
\item
  Apply (R) with \(\varepsilon/4\). Because \(s_m\to\infty\), one may
  also require \(s_m\ge4(|\bar\beta|+1)/\varepsilon\).
\item
  Then \(s_m>0\) and \(|\bar\beta|<\frac\varepsilon4s_m\).
\item
  With \(\kappa_G=2b_0\), the triangle inequality gives

  \[|\beta_{N,p,m}-2b_0s_m|\le|\beta_{N,p,m}-\widehat\beta_m|+|\bar\beta|<\varepsilon s_m .\]
\item
  Divide by \(s_m\).
\end{enumerate}

\emph{Lean:}
\href{https://github.com/yang-mills-group/OpenAI-YM-20261028/blob/3456e5e75cf8d4cd558955a5c7638fba2d68f464/YM/Pert/Asymptotics/RawAxis.lean\#L15-L30}{\texttt{zeroUp\_bound}}.

\textbf{Lemma VII.5.d (= Corollary VII.5.2).} \emph{Proof sketch.} The
slope is \(2b_0\), which is positive because \(C_{\mathfrak g}>0\) for
simple \(\mathfrak g\). Given \(\varepsilon\), take \(N_0:=m_0\) from
VII.5.1 and specialise to \(m=N\). The bare coupling of the public
family is \(\beta_{N,p,N}\), and its \(\log(1/a_N)\) is \(s_N\) by
definition. \emph{Lean:}
\href{https://github.com/yang-mills-group/OpenAI-YM-20261028/blob/3456e5e75cf8d4cd558955a5c7638fba2d68f464/HierarchyBridge.lean\#L16-L25}{\texttt{YangMills.regular}}
(VIII.1).

\textbf{Lemma VII.5.e (= Proposition VII.5.3).} \emph{Proof sketch.}

\begin{enumerate}
\def\labelenumi{\arabic{enumi}.}
\tightlist
\item
  For each \(\varepsilon\), eventually both \(i\ge m_0(\varepsilon)\)
  and \(i\le N_i\). So VII.5.1 gives \(\beta_{N_i,k_i,i}/s_i\to2b_0\).
\item
  Since \(\kappa_G>0\) for simple \(\mathfrak g\), \(2b_0>0\). Then
  \(x\mapsto\sqrt x\) is continuous at \(2b_0\) and \(y\mapsto y^{-1}\)
  is continuous at \(\sqrt{2b_0}\ne0\). Hence
  \((\beta_{N_i,k_i,i}/s_i)^{-1/2}\to(2b_0)^{-1/2}\).
\item
  For all sufficiently large \(i\), both \(s_i\) and
  \(\beta_{N_i,k_i,i}\) are positive, so
  \(\beta_{N_i,k_i,i}^{-1/2}s_i^{1/2}=(\beta_{N_i,k_i,i}/s_i)^{-1/2}\),
  and \(s_i=\log(1/a_i)\) by definition.
\end{enumerate}

\emph{Lean:}
\href{https://github.com/yang-mills-group/OpenAI-YM-20261028/blob/3456e5e75cf8d4cd558955a5c7638fba2d68f464/YM/Pert/Asymptotics/RawAxis.lean\#L36-L48}{\texttt{comparison}}.

\paragraph{4.2 Normalisation of the stress tensor
(VII.5.4)}\label{normalisation-of-the-stress-tensor-vii.5.4}

\textbf{Lemma VII.5.f (ratio lemma).} Let \(Z\) be a finite-dimensional
real normed space. If \(X_k\to x\neq0\) in \(Z\) and \(f_kX_k\to c\,x\)
with \(f_k\in\mathbb R\), then \(f_k\to c\). \emph{Proof sketch.} Choose
a linear functional \(\ell\) with \(\ell(x)\neq0\). Then
\(\ell(X_k)\ne0\) eventually, and \(f_k=\ell(f_kX_k)/\ell(X_k)\to c\).
\emph{Lean:}
\href{https://github.com/yang-mills-group/OpenAI-YM-20261028/blob/3456e5e75cf8d4cd558955a5c7638fba2d68f464/YM/Extra/WEAK_SCALING_ACTLAW/Orientation.lean\#L28-L39}{\texttt{scale\_limit}}.

\textbf{Lemma VII.5.g (= VII.5.4(a)).} \emph{Proof sketch.}

\begin{enumerate}
\def\labelenumi{\arabic{enumi}.}
\tightlist
\item
  A04 at degree 4 gives thresholds beyond which
  \(X_k:=r_k\,U_k(5)^{-1}\partial^V_4U_k(4)\to x:=\partial^W_4\).
\item
  A05 at degree 4, applied to \((J,p,Y)\), gives thresholds beyond which
  \((r_k/\xi_k)\,U_k(5)^{-1}p_4U_k(4)\to\kappa\,\partial^W_4\).
\item
  Take the componentwise maximum of the two sets of thresholds.
\item
  Substituting \(p_4=c\,\partial^V_4\) in step 2 gives
  \((c/\xi_k)X_k\to\kappa x\), where \(x\ne0\) is part of the free data.
  Lemma VII.5.f gives \(c/\xi_k\to\kappa\).
\end{enumerate}

The transport is \(U_k(d+1)^{-1}(\,\cdot\,)U_k(d)\), as defined by
\href{https://github.com/yang-mills-group/OpenAI-YM-20261028/blob/3456e5e75cf8d4cd558955a5c7638fba2d68f464/YM/Climb/Tests.lean\#L44-L45}{\texttt{carry}}
and
\href{https://github.com/yang-mills-group/OpenAI-YM-20261028/blob/3456e5e75cf8d4cd558955a5c7638fba2d68f464/YM/Climb/Tests.lean\#L49-L57}{\texttt{aisle}}.

\emph{Lean:}
\href{https://github.com/yang-mills-group/OpenAI-YM-20261028/blob/3456e5e75cf8d4cd558955a5c7638fba2d68f464/YM/Extra/WEAK_SCALING_ACTLAW/Orientation.lean\#L43-L54}{\texttt{tracks\_ratio}}.

\textbf{Lemma VII.5.h (= VII.5.4(b)).} \emph{Proof sketch.} Suppose
\(J=0\). The translation Ward identity writes every correlation with an
insertion of \(p_{d,a}v\) as a flux of \(J\) through a shell, which is
then \(0\). Correlations separate the elements of \(\mathcal F_{d+1}\),
so \(p=0\). \emph{Lean:}
\href{https://github.com/yang-mills-group/OpenAI-YM-20261028/blob/3456e5e75cf8d4cd558955a5c7638fba2d68f464/YM/Extra/WEAK_SCALING_ACTLAW/Criteria.lean\#L64-L70}{\texttt{flux\_nontrivial}}.

\emph{Remark (sign of \(c\kappa\)).} Suppose in addition that tracking
sequences exist at every level. Then \(c\kappa>0\): since \(\xi_k>0\),
\(c/\xi_k\) has the sign of \(c\), and its limit is \(\kappa\ne0\). This
is
\href{https://github.com/yang-mills-group/OpenAI-YM-20261028/blob/3456e5e75cf8d4cd558955a5c7638fba2d68f464/YM/Extra/WEAK_SCALING_ACTLAW/Orientation.lean\#L56-L70}{\texttt{kappa\_sign}}.
It is not assigned to this unit and is not used by the main theorem.
VI.2 uses only (a), and obtains \(c\neq0\) from \(\kappa\ne0\).

\paragraph{4.3 Exactly compatible frames
(VII.5.5)}\label{exactly-compatible-frames-vii.5.5}

\textbf{Lemma VII.5.i (equivariant retraction, = VII.5.5(a)).}
\emph{Proof sketch.}

\begin{enumerate}
\def\labelenumi{\arabic{enumi}.}
\tightlist
\item
  \(\iota^W\) is injective, so it has a linear left inverse \(R\).
\item
  Average over \(\Gamma_4\):
  \(Q:=|\Gamma_4|^{-1}\sum_{s\in\Gamma_4}s^{-1}Rs\). Since \(\iota^W\)
  commutes with \(\Gamma_4\), \(Q\) is again a left inverse, and \(Q\)
  is \(\Gamma_4\)-equivariant.
\item
  Let \(\theta\) be the conjugation, an antilinear involution. Put
  \(P:=\frac12(Q+\theta Q\theta)\). Then \(P\) is \(\mathbb C\)-linear.
  It is a left inverse, since \(\iota^W\) commutes with \(\theta\). It
  commutes with \(\theta\). It stays \(\Gamma_4\)-equivariant, since
  \(\theta\) commutes with \(\Gamma_4\).
\end{enumerate}

\emph{Lean:}
\href{https://github.com/yang-mills-group/OpenAI-YM-20261028/blob/3456e5e75cf8d4cd558955a5c7638fba2d68f464/YM/Extra/WEAK_SCALING_ACTLAW/Adjust.lean\#L35-L55}{\texttt{retraction}},
\href{https://github.com/yang-mills-group/OpenAI-YM-20261028/blob/3456e5e75cf8d4cd558955a5c7638fba2d68f464/YM/Extra/WEAK_SCALING_ACTLAW/Adjust.lean\#L27-L32}{\texttt{averaged\_cov}}.

\textbf{Lemma VII.5.j (exactness of the correction, = VII.5.5(b)).}
\emph{Proof sketch.}

\begin{enumerate}
\def\labelenumi{\arabic{enumi}.}
\tightlist
\item
  \emph{One step.}
  \(X_{d+1}\iota^Wv=\iota^VX_dP_d\iota^Wv+A_{d+1}(\iota^Wv-\iota^WP_d\iota^Wv)=\iota^VX_dv\).
\item
  \emph{Intertwining for \(d\le e\).} Induction on \(e\), using
  \(\iota_{d\to e+1}=\iota_{e\to e+1}\iota_{d\to e}\) and
  \(\iota_{d\to d}=1\).
\item
  \emph{Symmetries.} Induction on \(d\). Let \(g\) be a signed
  permutation or the conjugation. Then \(g\) commutes with \(P_d\),
  \(\iota^W\), \(\iota^V\), \(A_{d+1}\) and, by induction, with \(X_d\).
  Hence it commutes with both terms of \(X_{d+1}\).
\item
  \emph{Units.} \(\mathbf 1_{d+1}=\iota^W\mathbf 1_d\), so
  \(X_{d+1}\mathbf 1_{d+1}=\iota^VX_d\mathbf 1_d=\iota^V\mathbf 1_d=\mathbf 1_{d+1}\).
\end{enumerate}

\emph{Lean:}
\href{https://github.com/yang-mills-group/OpenAI-YM-20261028/blob/3456e5e75cf8d4cd558955a5c7638fba2d68f464/YM/Extra/WEAK_SCALING_ACTLAW/Adjust.lean\#L78-L79}{\texttt{adjusted\_succ}},
\href{https://github.com/yang-mills-group/OpenAI-YM-20261028/blob/3456e5e75cf8d4cd558955a5c7638fba2d68f464/YM/Extra/WEAK_SCALING_ACTLAW/Adjust.lean\#L80-L102}{\texttt{adjusted\_exact}}.

\textbf{Lemma VII.5.k (asymptotic identity).} Under the hypotheses of
VII.5.5(c), for \(d\le D\),
\(\Pi^{(k)}_d:=(A^{(k)}_d)^{-1}X^{(k)}_d\to1\) in operator norm.
\emph{Proof sketch.}

\begin{enumerate}
\def\labelenumi{\arabic{enumi}.}
\item
  Put \(\Phi^{(k)}_d:=(A^{(k)}_{d+1})^{-1}\iota^VA^{(k)}_d\). Then

  \[\Pi^{(k)}_{d+1}v=\Phi^{(k)}_d\,\Pi^{(k)}_d P_dv+(v-\iota^WP_dv),\qquad \Pi^{(k)}_0=1 .\]
\item
  \(\Phi^{(k)}_d\to\iota^W\) pointwise by hypothesis, hence in norm
  (finite dimension).
\item
  By induction, \(\Pi^{(k)}_dP_dv\to P_dv\).
\item
  Therefore \(\Pi^{(k)}_{d+1}v\to\iota^WP_dv+v-\iota^WP_dv=v\). Finite
  dimensionality converts pointwise convergence to operator-norm
  convergence.
\end{enumerate}

\emph{Lean:}
\href{https://github.com/yang-mills-group/OpenAI-YM-20261028/blob/3456e5e75cf8d4cd558955a5c7638fba2d68f464/YM/Extra/WEAK_SCALING_ACTLAW/Adjust.lean\#L106-L108}{\texttt{perturb}},
\href{https://github.com/yang-mills-group/OpenAI-YM-20261028/blob/3456e5e75cf8d4cd558955a5c7638fba2d68f464/YM/Extra/WEAK_SCALING_ACTLAW/Adjust.lean\#L113-L129}{\texttt{perturb\_limit}}.

\textbf{Lemma VII.5.l (bijectivity and the new frames, VII.5.5(c)).}
\emph{Proof sketch.}

\begin{enumerate}
\def\labelenumi{\arabic{enumi}.}
\tightlist
\item
  \(\Pi^{(k)}_d\to1\) in norm, and injective operators form an open set.
  So \(\Pi^{(k)}_d\), and hence \(X^{(k)}_d\), is injective for large
  \(k\).
\item
  \(\dim\mathcal F^\circ_d=\dim\mathcal F_d\) (by \(A^{(k)}_d\)), so
  \(X^{(k)}_d\) is bijective.
\item
  Hence \(\widetilde X^{(k)}_d=X^{(k)}_d\) eventually, and
  \((A^{(k)}_d)^{-1}\widetilde X^{(k)}_dv\to v\) for every \(v\).
\item
  Intersect the eventual sets for \(d=0,\ldots,D\). If the original
  frames satisfy the symmetry and unit conditions of (b), all corrected
  frames then agree with the exactly compatible maps \(X^{(k)}_d\) on
  this degree range.
\end{enumerate}

\emph{Lean:}
\href{https://github.com/yang-mills-group/OpenAI-YM-20261028/blob/3456e5e75cf8d4cd558955a5c7638fba2d68f464/YM/Extra/WEAK_SCALING_ACTLAW/Adjust.lean\#L131-L135}{\texttt{eventually\_unit}},
\href{https://github.com/yang-mills-group/OpenAI-YM-20261028/blob/3456e5e75cf8d4cd558955a5c7638fba2d68f464/YM/Extra/WEAK_SCALING_ACTLAW/Adjust.lean\#L142-L143}{\texttt{rotate\_eq}},
\href{https://github.com/yang-mills-group/OpenAI-YM-20261028/blob/3456e5e75cf8d4cd558955a5c7638fba2d68f464/YM/Extra/WEAK_SCALING_ACTLAW/Adjust.lean\#L144-L148}{\texttt{relative}}.

\emph{Use.} In VII.1 the frames of the free-field scaling sequence
intertwine the inclusions only asymptotically. VII.5.5 replaces them by
frames that are exactly graded for all large \(k\) and asymptotically
equal to the original ones, as tracking requires
(\href{https://github.com/yang-mills-group/OpenAI-YM-20261028/blob/3456e5e75cf8d4cd558955a5c7638fba2d68f464/YM/SharpGram/Weak/Compatible.lean\#L140-L159}{\texttt{tune\_grades}}).

\paragraph{4.4 Finite-dimensional convergence
(VII.5.6)}\label{finite-dimensional-convergence-vii.5.6}

\textbf{Lemma VII.5.m (= Proposition VII.5.6).} \emph{Proof sketch.}

\begin{itemize}
\tightlist
\item
  \begin{enumerate}
  \def\labelenumi{(\alph{enumi})}
  \tightlist
  \item
    A linear left inverse \(g\) of \(f\) is continuous in finite
    dimension, so \(a_y=gf(a_y)\to gf(e)=e\).
  \end{enumerate}
\item
  (b) Every point of \((\mathbb R^4)^n\) is
  \(\sum_{i,a}v_{i,a}e_{i,a}\). Hence evaluation at the \(4n\)
  coordinate vectors is an injective real-linear map from the space of
  complex-valued real-linear maps to \(\mathbb C^{4n}\), and (a)
  applies.
\end{itemize}

\emph{Lean:}
\href{https://github.com/yang-mills-group/OpenAI-YM-20261028/blob/3456e5e75cf8d4cd558955a5c7638fba2d68f464/YM/Extra/WEAK_SCALING_ACTLAW/Analysis.lean\#L44-L51}{\texttt{norm\_fit}},
\href{https://github.com/yang-mills-group/OpenAI-YM-20261028/blob/3456e5e75cf8d4cd558955a5c7638fba2d68f464/YM/Extra/WEAK_SCALING_ACTLAW/Analysis.lean\#L55-L60}{\texttt{components}},
\href{https://github.com/yang-mills-group/OpenAI-YM-20261028/blob/3456e5e75cf8d4cd558955a5c7638fba2d68f464/YM/Extra/WEAK_SCALING_ACTLAW/Analysis.lean\#L63-L70}{\texttt{norm\_fit\_deriv}}.

\paragraph{4.5 The classical stress-tensor algebra
(VII.5.7)}\label{the-classical-stress-tensor-algebra-vii.5.7}

\textbf{Lemma VII.5.n (covariance of \(\mu^1,\mu^2\)).} For a signed
permutation \(t\) with signs \(\varepsilon_t(\cdot)\):

\[t\cdot\mu^2_a=\varepsilon_t(a)\,\mu^2_{t(a)},\qquad t\cdot\mu^1_a=\varepsilon_t(a)\,\mu^1_{t(a)} .\]

Both are real. \emph{Proof sketch.} \(t\) multiplies
\(\langle F_{ij},\partial_bF_{ij}\rangle\) by
\(\varepsilon(i)\varepsilon(j)\,\varepsilon(b)\varepsilon(i)\varepsilon(j)=\varepsilon(b)\)
and relabels the indices. The double sum over \((i,j)\) is invariant
under reindexing by \(t\). For \(\mu^1_b\), the sign is
\(\varepsilon(b)^3\varepsilon(i)^2=\varepsilon(b)\), and the remaining
index is reindexed in the same way. Reality follows from the real inner
products. \emph{Lean:}
\href{https://github.com/yang-mills-group/OpenAI-YM-20261028/blob/3456e5e75cf8d4cd558955a5c7638fba2d68f464/YM/TraceExtension/Foundation/Caps.lean\#L21-L24}{\texttt{momentum2\_turn}},
\href{https://github.com/yang-mills-group/OpenAI-YM-20261028/blob/3456e5e75cf8d4cd558955a5c7638fba2d68f464/YM/TraceExtension/Foundation/Caps.lean\#L28-L34}{\texttt{momentum1\_mem}},
\href{https://github.com/yang-mills-group/OpenAI-YM-20261028/blob/3456e5e75cf8d4cd558955a5c7638fba2d68f464/YM/TraceExtension/Foundation/Caps.lean\#L25-L27}{\texttt{momentum2\_mem}},
\href{https://github.com/yang-mills-group/OpenAI-YM-20261028/blob/3456e5e75cf8d4cd558955a5c7638fba2d68f464/YM/TraceExtension/Foundation/Caps.lean\#L17-L20}{\texttt{sumPerm}}.

\textbf{Lemma VII.5.o (affine shells).} Every degree-0 radial jet is
realised by a free Maxwell field with constant curvature and an affine
potential. \emph{Proof sketch.} The associated potential has constant
curvature, so \(\sum_a\partial_aF_{ab}=0\). \emph{Lean:}
\href{https://github.com/yang-mills-group/OpenAI-YM-20261028/blob/3456e5e75cf8d4cd558955a5c7638fba2d68f464/YM/TraceExtension/Foundation/Caps.lean\#L43-L50}{\texttt{shell\_affine}}.
VII.2 uses this, with its degree-1 analogue
\href{https://github.com/yang-mills-group/OpenAI-YM-20261028/blob/3456e5e75cf8d4cd558955a5c7638fba2d68f464/YM/TraceExtension/Foundation/Caps.lean\#L52-L61}{\texttt{shell\_quadratic}},
to build test jets on Maxwell solutions.

\textbf{Consequences in VI.1 and VII.2.} Write
\(T^{(0,w)}:=T^{(0,w_0,w_1,w_2)}\) for \(w\in\mathbb R^3\).

\begin{enumerate}
\def\labelenumi{\arabic{enumi}.}
\item
  \emph{Trace (VI.1).}
  \(\sum_aT^{(0,w)}_{aa}=4(w_0+4w_1)\,\mathfrak s\). In particular
  \(T^{\mathrm{can}}\) is traceless
  (\href{https://github.com/yang-mills-group/OpenAI-YM-20261028/blob/3456e5e75cf8d4cd558955a5c7638fba2d68f464/YM/TraceExtension/Foundation/Dilate.lean\#L67-L69}{\texttt{mechanical\_trace}},
  \href{https://github.com/yang-mills-group/OpenAI-YM-20261028/blob/3456e5e75cf8d4cd558955a5c7638fba2d68f464/YM/TraceExtension/Foundation/Dilate.lean\#L71-L72}{\texttt{conformal}}).
\item
  \emph{Divergence (VII.2).}
  \(\delta(w)_b=\sum_a\partial_aT^{(0,w)}_{ab}\)
  (\href{https://github.com/yang-mills-group/OpenAI-YM-20261028/blob/3456e5e75cf8d4cd558955a5c7638fba2d68f464/YM/TraceExtension/Foundation/Dilate.lean\#L38-L52}{\texttt{euler}}).
  The key identity is the Bianchi consequence

  \[\sum_{a,i}\langle F_{ai},\partial_aF_{bi}\rangle=\tfrac12\mu^2_b .\]

  With it, the shear part
  \(\sum_{a\ne b}\partial_a\sum_i\langle F_{ai},F_{bi}\rangle\) equals
  \(E_b+\frac12\mu^2_b-2\mu^1_b\).
\item
  \emph{Coordinates (VII.2).} For simple \(\mathfrak g\), \(\beta\) and
  \(\delta\) are bijections \(\mathbb R^3\to\mathcal R_1\)
  (\href{https://github.com/yang-mills-group/OpenAI-YM-20261028/blob/3456e5e75cf8d4cd558955a5c7638fba2d68f464/YM/TraceExtension/Foundation/Caps.lean\#L112-L124}{\texttt{coordinates1}},
  \href{https://github.com/yang-mills-group/OpenAI-YM-20261028/blob/3456e5e75cf8d4cd558955a5c7638fba2d68f464/YM/TraceExtension/Foundation/Dilate.lean\#L54-L63}{\texttt{bareCoordinates}}).
  Under the same simplicity hypothesis, \(\beta(w)\) vanishes on free
  Maxwell solutions if and only if \(w_0=w_1=0\)
  (\href{https://github.com/yang-mills-group/OpenAI-YM-20261028/blob/3456e5e75cf8d4cd558955a5c7638fba2d68f464/YM/TraceExtension/Foundation/Caps.lean\#L96-L110}{\texttt{through\_eq}}).
\item
  \emph{Conservation of the canonical tensor (VII.2).}
  \(\delta(1,-\tfrac14,1)=E\): \(T^{\mathrm{can}}\) is conserved on the
  equations of motion
  (\href{https://github.com/yang-mills-group/OpenAI-YM-20261028/blob/3456e5e75cf8d4cd558955a5c7638fba2d68f464/YM/TraceExtension/Foundation/Dilate.lean\#L73-L74}{\texttt{eom\_vector\_raw}}).
\item
  \emph{The functional \(\rho\) (VII.2).} For simple \(\mathfrak g\),
  define \(\rho:\mathcal R_1\to\mathbb R\) by
  \(\rho(\delta(w)):=-4(w_0+4w_1)\). Then
  \(\sum_aT^{(0,w)}_{aa}=-\rho(\delta(w))\,\mathfrak s\). The functional
  \(\rho\) vanishes on vectors whose components vanish on free Maxwell
  solutions
  (\href{https://github.com/yang-mills-group/OpenAI-YM-20261028/blob/3456e5e75cf8d4cd558955a5c7638fba2d68f464/YM/TraceExtension/Foundation/Dilate.lean\#L78-L80}{\texttt{rho}},
  \href{https://github.com/yang-mills-group/OpenAI-YM-20261028/blob/3456e5e75cf8d4cd558955a5c7638fba2d68f464/YM/TraceExtension/Foundation/Dilate.lean\#L85-L88}{\texttt{rho\_eom}}).
\end{enumerate}

For simple \(\mathfrak g\), a tensor in this four-parameter family is
conserved on free Maxwell fields precisely when it is a multiple of
\(T^{\mathrm{can}}\) plus a constant multiple of
\(\delta_{ab}\mathbf1\). With the unit term removed, its trace is
determined by its divergence through \(\rho\). VII.2 uses this
functional in
\href{https://github.com/yang-mills-group/OpenAI-YM-20261028/blob/3456e5e75cf8d4cd558955a5c7638fba2d68f464/YM/TraceExtension/Terrace/Lease.lean\#L22}{\texttt{Lease}}
and its limiting comparison
\href{https://github.com/yang-mills-group/OpenAI-YM-20261028/blob/3456e5e75cf8d4cd558955a5c7638fba2d68f464/YM/TraceExtension/Chart/Token.lean\#L75-L88}{\texttt{mere}}.

\subsubsection{5. Where the difficulty
lies}\label{where-the-difficulty-lies-29}

The logarithmic coupling law uses only the uniform running clause and
elementary real analysis. The analytic work lies in the upstream
one-loop estimate, the coupling reset and the construction of a
trajectory reaching \(\bar\beta\) at unit scale. The auxiliary results
use finite-dimensional linear algebra, equivariant averaging and
scalar-ratio limits.

The result is uniform both in the volume index \(p\) and in the finer
lattice level \(N\ge m\). Thus VII.5.3 allows arbitrary volume
sequences. If inverse couplings converge as \(N,p\to\infty\) at fixed
\(m\), their limits inherit the non-strict bound
\(|\beta/s_m-2b_0|\le\varepsilon\) for \(m\ge m_0(\varepsilon)\).
Existence of those limits and comparison with a coupling defined by
continuum correlations require separate results.

Only the leading logarithm is established here, for the coefficients of
the lattice effective actions. No two-loop correction or
\(\Lambda\)-parameter is asserted. The main theorem uses VII.5.1 through
(Reg1), retaining only the existence of a positive slope. Proposition
VII.5.3 is not used in the assembly.

The 38 auxiliary declarations serve other parts of the proof: stress
normalisation in VI.2, frame correction in VII.1, convergence in III.12,
and classical stress algebra and read-out estimates in VI.1 and VII.2.

\subsubsection{6. Notation}\label{notation-29}

{\def\LTcaptype{none} % do not increment counter
\begin{longtable}[]{@{}
  >{\raggedright\arraybackslash}p{(\linewidth - 4\tabcolsep) * \real{0.3333}}
  >{\raggedright\arraybackslash}p{(\linewidth - 4\tabcolsep) * \real{0.3333}}
  >{\raggedright\arraybackslash}p{(\linewidth - 4\tabcolsep) * \real{0.3333}}@{}}
\toprule\noalign{}
\begin{minipage}[b]{\linewidth}\raggedright
object
\end{minipage} & \begin{minipage}[b]{\linewidth}\raggedright
code name
\end{minipage} & \begin{minipage}[b]{\linewidth}\raggedright
proposed notation
\end{minipage} \\
\midrule\noalign{}
\endhead
\bottomrule\noalign{}
\endlastfoot
Casimir constant &
\href{https://github.com/yang-mills-group/OpenAI-YM-20261028/blob/3456e5e75cf8d4cd558955a5c7638fba2d68f464/YM/Gauge/Colors.lean\#L180-L182}{\texttt{YM.Gauge.casimir}}
& \(C_{\mathfrak g}\) (\(=2q^2\) for \(SU(q)\)) \\
slope constant &
\href{https://github.com/yang-mills-group/OpenAI-YM-20261028/blob/3456e5e75cf8d4cd558955a5c7638fba2d68f464/YM/Gauge/Colors.lean\#L184}{\texttt{YM.Gauge.slopeConstant}}
& \(\kappa_G=11C_{\mathfrak g}/(24\pi^2)\) \\
one-loop coefficient &
\href{https://github.com/yang-mills-group/OpenAI-YM-20261028/blob/3456e5e75cf8d4cd558955a5c7638fba2d68f464/YM/Pert/Asymptotics/RawAxis.lean\#L13}{\texttt{b0}}
& \(b_0=\kappa_G/2\); slope \(b_*=2b_0\) \\
lattice spacing at stage \(m\) &
\href{https://github.com/yang-mills-group/OpenAI-YM-20261028/blob/3456e5e75cf8d4cd558955a5c7638fba2d68f464/YM/Lattice/Scale.lean\#L99}{\texttt{Schedule.a}}
& \(a_m=1/D_m\) \\
logarithmic scale &
\href{https://github.com/yang-mills-group/OpenAI-YM-20261028/blob/3456e5e75cf8d4cd558955a5c7638fba2d68f464/YM/Lattice/Scale.lean\#L115}{\texttt{Schedule.scale}}
& \(s_m=\log(1/a_m)\) \\
coupling at stage \(m\) & \texttt{C.path.β\ N\ p\ m} &
\(\beta_{N,p,m}\); bare coupling \(\beta_{N,p}=\beta_{N,p,N}\) \\
reference line &
\href{https://github.com/yang-mills-group/OpenAI-YM-20261028/blob/3456e5e75cf8d4cd558955a5c7638fba2d68f464/YM/RG/History.lean\#L79}{\texttt{hatBeta}}
& \(\widehat\beta_m=\bar\beta+\kappa_Gs_m\) \\
running clause &
\href{https://github.com/yang-mills-group/OpenAI-YM-20261028/blob/3456e5e75cf8d4cd558955a5c7638fba2d68f464/YM/RG/History.lean\#L229-L230}{\texttt{Running}}
& (R) \\
running coupling &
\href{https://github.com/yang-mills-group/OpenAI-YM-20261028/blob/3456e5e75cf8d4cd558955a5c7638fba2d68f464/YM/Pert/Asymptotics/RawAxis.lean\#L35}{\texttt{regulatorRun}}
& \(g_{N,p,m}=\beta_{N,p,m}^{-1/2}\) \\
free comparison table &
\href{https://github.com/yang-mills-group/OpenAI-YM-20261028/blob/3456e5e75cf8d4cd558955a5c7638fba2d68f464/YM/Extra/WEAK_SCALING_ACTLAW/Station.lean\#L28}{\texttt{Free}}
& \(\mathcal F^\circ_d\) \\
equivariant retraction &
\href{https://github.com/yang-mills-group/OpenAI-YM-20261028/blob/3456e5e75cf8d4cd558955a5c7638fba2d68f464/YM/Extra/WEAK_SCALING_ACTLAW/Adjust.lean\#L35-L55}{\texttt{retraction}}
& \(P_d\) \\
corrected frames &
\href{https://github.com/yang-mills-group/OpenAI-YM-20261028/blob/3456e5e75cf8d4cd558955a5c7638fba2d68f464/YM/Extra/WEAK_SCALING_ACTLAW/Adjust.lean\#L61-L66}{\texttt{adjusted}},
\href{https://github.com/yang-mills-group/OpenAI-YM-20261028/blob/3456e5e75cf8d4cd558955a5c7638fba2d68f464/YM/Extra/WEAK_SCALING_ACTLAW/Adjust.lean\#L137-L140}{\texttt{rotateChart}}
& \(X_d\), \(\widetilde X^{(k)}_d\) \\
exact compatibility up to degree \(D\) &
\href{https://github.com/yang-mills-group/OpenAI-YM-20261028/blob/3456e5e75cf8d4cd558955a5c7638fba2d68f464/YM/Extra/WEAK_SCALING_ACTLAW/Adjust.lean\#L72-L76}{\texttt{Exact}},
\href{https://github.com/yang-mills-group/OpenAI-YM-20261028/blob/3456e5e75cf8d4cd558955a5c7638fba2d68f464/YM/Extra/WEAK_SCALING_ACTLAW/Adjust.lean\#L171-L172}{\texttt{Graded}}
& ``exactly graded'' \\
asymptotic intertwining &
\href{https://github.com/yang-mills-group/OpenAI-YM-20261028/blob/3456e5e75cf8d4cd558955a5c7638fba2d68f464/YM/Extra/WEAK_SCALING_ACTLAW/Adjust.lean\#L110-L111}{\texttt{Near}}
& hypothesis of VII.5.5(c) \\
square of the curvature &
\href{https://github.com/yang-mills-group/OpenAI-YM-20261028/blob/3456e5e75cf8d4cd558955a5c7638fba2d68f464/YM/TraceExtension/Foundation/Dilate.lean\#L66}{\texttt{square4}}
& \(\mathfrak s=\frac14\sum_{i,j}\langle F_{ij},F_{ij}\rangle\) \\
canonical stress tensor &
\href{https://github.com/yang-mills-group/OpenAI-YM-20261028/blob/3456e5e75cf8d4cd558955a5c7638fba2d68f464/YM/TraceExtension/Foundation/Dilate.lean\#L70}{\texttt{raw}}
& \(T^{\mathrm{can}}\) \\
weight-5 vectors &
\href{https://github.com/yang-mills-group/OpenAI-YM-20261028/blob/3456e5e75cf8d4cd558955a5c7638fba2d68f464/YM/TraceExtension/Foundation/Caps.lean\#L16}{\texttt{momentum1}},
\href{https://github.com/yang-mills-group/OpenAI-YM-20261028/blob/3456e5e75cf8d4cd558955a5c7638fba2d68f464/YM/TraceExtension/Foundation/Caps.lean\#L14}{\texttt{momentum2}}
& \(\mu^1\), \(\mu^2\) \\
coordinates on \(\mathcal R_1\) &
\href{https://github.com/yang-mills-group/OpenAI-YM-20261028/blob/3456e5e75cf8d4cd558955a5c7638fba2d68f464/YM/TraceExtension/Foundation/Caps.lean\#L63-L66}{\texttt{basisV}}
& \(\beta(w)\) \\
divergence of \(T^{(0,w)}\) &
\href{https://github.com/yang-mills-group/OpenAI-YM-20261028/blob/3456e5e75cf8d4cd558955a5c7638fba2d68f464/YM/TraceExtension/Foundation/Dilate.lean\#L33}{\texttt{bareVector}}
& \(\delta(w)\) \\
sharp read-out bound &
\href{https://github.com/yang-mills-group/OpenAI-YM-20261028/blob/3456e5e75cf8d4cd558955a5c7638fba2d68f464/YM/TraceExtension/Soot/Skein.lean\#L26}{\texttt{SharpRead}}
& Definition VII.5.8 \\
\end{longtable}
}

\subsubsection{7. What remains
unexplained}\label{what-remains-unexplained-29}

\begin{enumerate}
\def\labelenumi{\arabic{enumi}.}
\tightlist
\item
  \textbf{The upstream one-loop calculation.} The transfer bound is
  traced through the coupling reset and the one-loop shift bound in §2.
  The evaluation of the momentum integral giving
  \(11C_{\mathfrak g}/(24\pi^2)\) belongs to III.7; its analytic proof
  is not verified here. The coupling results in this unit take the
  construction conditions as hypotheses.
\item
  \textbf{The coupling after the limits.} The assessment's Theorem 7.2
  concerns \(g_W(r)\) after removal of the initial cutoff and a
  subsequential infinite-volume limit, together with comparison to a
  coupling read from continuum correlations. VII.5.3 proves the
  lattice-stage law. A declaration establishing the additional limits
  and comparison has not been identified here (tentative). The files
  \texttt{Beacon}, \texttt{Hymnal} and \texttt{Musicology} in
  \texttt{YM/RemainderComparison} are possible starting points.
\item
  \textbf{Existence of the comparison inputs.} VII.5.4 assumes the free
  Maxwell data of VI.2 and a Euclidean hierarchy on \(S\); it gives a
  conclusion for every tracking sequence meeting its finite thresholds.
  It does not construct such sequences. Their existence belongs to VI.1
  and VII.1. Likewise, the supply of the read-out bound in III.2 uses
  the hypotheses of the theorem linked in §2; the underlying one-step
  estimates are outside this unit.
\end{enumerate}

\subsection{VII.6. Scaling variants and fixed-field
estimates}\label{vii.6.-scaling-variants-and-fixed-field-estimates}

\textbf{Part:} VII, identification as Yang--Mills. \textbf{Lean:}
\href{https://github.com/yang-mills-group/OpenAI-YM-20261028/tree/3456e5e75cf8d4cd558955a5c7638fba2d68f464/YM/TraceExtension/Prom}{\texttt{YM/TraceExtension/Prom/}}
contains the ``Sky'' files, 133 files with ``dynamic stops''. Other
outputs are
\href{https://github.com/yang-mills-group/OpenAI-YM-20261028/blob/3456e5e75cf8d4cd558955a5c7638fba2d68f464/YM/Trace2/Transfer.lean\#L31-L37}{\texttt{Trace.first\_grade}}
and
\href{https://github.com/yang-mills-group/OpenAI-YM-20261028/blob/3456e5e75cf8d4cd558955a5c7638fba2d68f464/YM/UVNext/FREE_STRESS_CERTIFICATE/Size.lean\#L51-L54}{\texttt{FreeStress.dim\_before}}.
Also included are
\href{https://github.com/yang-mills-group/OpenAI-YM-20261028/blob/3456e5e75cf8d4cd558955a5c7638fba2d68f464/YM/SharpGram/Clauses.lean\#L126-L129}{\texttt{SharpGram.inletFirst}}
and its helpers in
\href{https://github.com/yang-mills-group/OpenAI-YM-20261028/blob/3456e5e75cf8d4cd558955a5c7638fba2d68f464/YM/SharpGram/Clauses.lean}{\texttt{YM/SharpGram/Clauses.lean}}.
\textbf{Size:} 4,205 lines in 374 declarations:

\begin{itemize}
\tightlist
\item
  about 3,600 lines in the chain of
  \href{https://github.com/yang-mills-group/OpenAI-YM-20261028/blob/3456e5e75cf8d4cd558955a5c7638fba2d68f464/YM/TraceExtension/Prom/Zoom/Running.lean\#L25-L31}{\texttt{pays\_sky}};
\item
  about 560 further lines in the dynamic-stop files;
\item
  43 lines in the chain of the supply theorem.
\end{itemize}

The unit has two parts.

\begin{itemize}
\tightlist
\item
  \textbf{(A) Scaling sequences at prescribed radii.} The ultraviolet
  construction (III.11) compares the lattice theory with the free field
  at reading scales \(\ell(b,k)=a_k\mu_k^b\). Varying the band \(b\)
  over a unit interval makes these scales cover every sufficiently small
  radius. The stage estimates remain valid across one change of block
  factor. They produce an approaching scaling of the limit theory along
  a subsequence of any prescribed positive null sequence.
\item
  \textbf{(B) Degrees below four.} In the limit theory every field of
  degree at most \(3\) is a multiple of the unit. It has a constant
  family of lattice representatives whose stage-\(k\) coefficient size
  is at most \(\hat\beta_k\) on an ordered cutoff tail, for every
  sufficiently large \(k\).
\end{itemize}

\subsubsection{1. What this unit proves}\label{what-this-unit-proves-30}

\textbf{Setting.}

\emph{The family and its stages.}

\begin{itemize}
\item
  \(G\) is a compact matrix group with simple Lie algebra
  \(\mathfrak g\), and \(C\) is a Wilson family with its construction
  record (II.2, III.13).
\item
  Its RG stages \(k\) have spacings \(a_k\downarrow0\) and block factors
  \(L_k=a_k/a_{k+1}\ge2\).
\item
  The running inverse coupling is
  \(\hat\beta_k=\bar\beta+c_G\log a_k^{-1}\), where \(c_G>0\). Put
  \[\mu_k=\exp\big((\log\log\hat\beta_k)^2\big),\] which grows faster
  than every positive power of \(\log\hat\beta_k\) and slower than every
  positive power of \(\hat\beta_k\) (III.11).
\item
  For a real band \(b\), the \emph{reading scale} of stage \(k\) is
  \(\ell(b,k)=a_k\mu_k^{\,b}\). For fixed \(b\) it tends to \(0\).
\end{itemize}

\emph{The limit theory and the free reference.}

\begin{itemize}
\tightlist
\item
  \(F\) is an extraction order and \(\Phi_F\) its filter of ordered
  cutoff tails on pairs \((N,p)\) (IV.1).
\item
  \(S\) is a limit table of \(C\) along \(F\) (ordered smeared limits
  and generic point limits), with a Euclidean hierarchy \(H\).
\item
  \(\mathcal F^S_d=\mathrm{Label}_d/\mathrm{Null}_d(S)\) is the filtered
  field space of degree at most \(d\),
  \(\pi_d:\mathrm{Label}_d\to\mathcal F^S_d\) is the projection, and
  \(\mathbf 1_d\) is the unit.
\item
  The reference is the free Maxwell table of VI.2: gauge-invariant Wick
  polynomials in a free \(\mathfrak g\)-valued field. Its field spaces
  are \(\mathcal F^{\rm free}_d\), its unit is
  \(\mathbf 1^{\rm free}_d\) and its hierarchy is \(H_0\).
\item
  Correlation kernels on \(\mathrm{Conf}_n\) are written
  \(\langle\phi_1(x_1)\cdots\phi_n(x_n)\rangle^S\) and
  \(\langle\cdots\rangle^{\rm free}\).
\end{itemize}

\emph{Scalings.} A \emph{scaling} (VI.1) consists of

\begin{itemize}
\tightlist
\item
  radii \(r_k>0\) strictly decreasing to \(0\),
\item
  constants \(\xi_k>0\),
\item
  complex-linear isomorphisms
  \(U_k(d):\mathcal F^{\rm free}_d\to\mathcal F^S_d\).
\end{itemize}

It is \emph{approaching} (VI.2.7(a)) if the following holds for every
\(n\) and all \(v_i\in\mathcal F^{\rm free}_{d_i}\). Put
\[K_k(x)=\langle U_kv_1(r_kx_1)\cdots U_kv_n(r_kx_n)\rangle^S-\langle v_1(x_1)\cdots v_n(x_n)\rangle^{\rm free}.\]
The kernels \(K_k\) are smooth on \(\mathrm{Conf}_n\) and tend to zero
locally with an order-zero bound: for every compact
\(K\subset\mathrm{Conf}_n\) and all \(M,\varepsilon>0\), eventually
\(|\int tK_k|\le\varepsilon\) for every smooth test function \(t\)
supported in \(K\) and satisfying \(\|t\|_\infty\le M\).

\emph{Fixed data.}

\begin{itemize}
\tightlist
\item
  \(B\) is a binding of the RG history to its price lists, \(O\) a Haar
  budget and \(J\) a chart ratio (III.11).
\item
  (S) is the one-step cost package of III.11.
\item
  \(c_d=2(w_d+1)\) are the mark exponents of III.11, where \(w_d\) is
  the polynomial degree of the degree-\(d\) label frame plus \(4\).
\item
  Write \(J_k\) for the price index of stage \(k\), \(\Lambda_j\) for
  its block-factor list, and \(T_d\) for the reference stage of degree
  \(d\). Thus \(L_k=\Lambda_{J_{k+1}}\). Stage packages, poised stages,
  the uniform-band property \(\mathrm{Fire}_d(z)\), vanes and checked
  vanes are defined in §4.
\end{itemize}

\textbf{Proposition VII.6.1 (approaching scalings along a subsequence of
any null sequence).} Assume:

\begin{itemize}
\tightlist
\item
  \begin{enumerate}
  \def\labelenumi{(\roman{enumi})}
  \tightlist
  \item
    \(G\) is simple and (S) holds for \(C\);
  \end{enumerate}
\item
  (ii) \((z_d)\) are reals with \(z_d\ge c_i\) for all \(i\le d\);
\item
  (iii) \(\mathrm{Fire}_d(z_d)\) holds for every \(d\) (Definition
  VII.6.d);
\item
  (iv) \(\dim\mathcal F^S_d\le\dim\mathcal F^{\rm free}_d\) for every
  \(d\);
\item
  (v) \(\mathcal L\) is an ultrafilter on \(\mathbb N^2\) finer than
  \(\Phi_F\).
\end{itemize}

Let \(s_i>0\) with \(s_i\to0\). Then there are a strictly increasing
\(f:\mathbb N\to\mathbb N\) and a scaling with \(r_k=s_{f(k)}\) and
\(\xi_k=1\) such that:

\begin{itemize}
\tightlist
\item
  \begin{enumerate}
  \def\labelenumi{(\alph{enumi})}
  \tightlist
  \item
    the scaling is approaching;
  \end{enumerate}
\item
  (b) for every \(d\) and all large \(k\), \(U_k(d)\) is real and
  covariant under signed coordinate permutations between \(H_0\) and
  \(H\), and \(U_k(d)\mathbf 1^{\rm free}_d=\mathbf 1_d\);
\item
  (c) there are vanes \(V_k\) of depth \(k\) with the following
  properties:

  \begin{itemize}
  \tightlist
  \item
    their stages tend to infinity;
  \item
    \(V_k\) is checked at level \(k\);
  \item
    its radius is \(\ell(V_k)=r_k\);
  \item
    for every \(d\) and all large \(k\), the full lift tuple of \(V_k\)
    in degrees \(\le k\) converges along \(\mathcal L\) to a tuple
    \(Y_k\), and \(U_k(d)=\pi_d\circ Y_{k,d}\).
  \end{itemize}
\end{itemize}

\emph{Status:} theorem, with (iii) and (iv) carried as hypotheses. Both
are supplied inside the construction run, where VII.2 applies this
result:

\begin{itemize}
\tightlist
\item
  (iii) is produced by the uniform-band construction (III.2);
\item
  (iv) is the rank-width bound (VI.2), applied to the corrected rank
  widths of the seed (IV.1).
\end{itemize}

\emph{Lean:}
\href{https://github.com/yang-mills-group/OpenAI-YM-20261028/blob/3456e5e75cf8d4cd558955a5c7638fba2d68f464/YM/TraceExtension/Prom/Zoom/Running.lean\#L25-L31}{\texttt{pays\_sky}}.
Conclusions (a)--(c) are the structure
\href{https://github.com/yang-mills-group/OpenAI-YM-20261028/blob/3456e5e75cf8d4cd558955a5c7638fba2d68f464/YM/TraceExtension/Prom/Zoom/Running.lean\#L19-L23}{\texttt{StarPath}}.

\textbf{Proposition VII.6.2 (checked vanes at prescribed radii).} Assume
(i)--(iii) of VII.6.1. Use the ordered and generic point limits from the
setting. Only smoothness of \(S\), rather than its full hierarchy, is
needed here; the free hierarchy \(H_0\) is still assumed. Let \(s_i>0\)
with \(s_i\to0\). Then there are

\begin{itemize}
\tightlist
\item
  a strictly increasing \(p:\mathbb N\to\mathbb N\) with \(s\circ p\)
  strictly decreasing, and
\item
  for each \(d\in\mathbb N\), a vane \(V_d\) of depth \(d\) with stage
  \(\ge d\), checked at level \(d\), with \(\ell(V_d)=s_{p(d)}\).
\end{itemize}

\emph{Status:} theorem.

\emph{Lean:}
\href{https://github.com/yang-mills-group/OpenAI-YM-20261028/blob/3456e5e75cf8d4cd558955a5c7638fba2d68f464/YM/TraceExtension/Prom/Zoom/Boats.lean\#L75-L87}{\texttt{wakesSky}},
with
\href{https://github.com/yang-mills-group/OpenAI-YM-20261028/blob/3456e5e75cf8d4cd558955a5c7638fba2d68f464/YM/TraceExtension/Prom/Zoom/Boats.lean\#L65-L69}{\texttt{Vane}}
and
\href{https://github.com/yang-mills-group/OpenAI-YM-20261028/blob/3456e5e75cf8d4cd558955a5c7638fba2d68f464/YM/TraceExtension/Prom/Zoom/Boats.lean\#L42-L47}{\texttt{TimberSky}}.

\textbf{Proposition VII.6.3 (stage-\(k\) chessboard bound with a moving
band).} Assume (i) of VII.6.1, let \(d\in\mathbb N\), and let
\(z\ge c_i\) for \(i\le d\). Fix the following data:

\begin{itemize}
\tightlist
\item
  \(\Delta\in\mathbb N\), \(e,u_0\in\mathbb R\), and stage packages
  \(P_k\) of degree \(\le d\) with Taylor order \(\Delta\), exponent
  \(e\) and bands \(b_k\le u_0\);
\item
  \(n\ge2\) and \(m+1\le n\);
\item
  labels \(g_i\in\mathrm{Label}_{d_i}\) with \(d_i\le d\), for
  \(0\le i\le m\);
\item
  a compact \(K\subset\mathrm{Conf}_{m+1}\).
\end{itemize}

Then the following holds for all large \(k\). Take:

\begin{itemize}
\tightlist
\item
  a stage-\(k\) cost record \(L\) of \(P_k\) with exponent \(-3\), mark
  exponent \(z\) and arity bound \(n\);
\item
  \(N>k\), a period \(p\), and a stage-\(k\) RG representative \(V\)
  compatible with \(L\) along the same RG history;
\item
  \(x\) whose rounding to the stage-\(k\) lattice, in units of
  \(\ell(b_k,k)\), lies in \(K\) and fits the stopping stage;
\item
  stage-\(k\) polynomial lift data \(T_i\) of Taylor order \(\Delta\)
  for \(V\), and labels \(f_i\in\mathrm{Label}_{d_i}\), satisfying the
  rescaling equation: after applying \(T_i\) and the coupling dilation
  of VII.6.n, the result is \(g_i\) weighted by \(\ell(b_k,k)/a_k\).
\end{itemize}

Then for every nonempty subset \(s\subset\{0,\ldots,m\}\),
\[\rho_{\mathcal T}(\Phi_s)\ \le\ \mathrm{bud}_n(k):=\mu_k^{-1/(16n)}.\]
To define \(\Phi_s\), introduce commuting markers \(\eta_i\) with
\(\eta_i^2=0\). If \(O_i\) is the lattice probe with label \(f_i\), and
\(D\) is the free cumulant polynomial of the \(g_i\) at the rounded,
rescaled points, then
\[\Phi_s=[\eta_s]\left(e^{-D}\prod_{i=0}^{m}(1+\eta_i O_i)\right),\qquad \eta_s=\prod_{i\in s}\eta_i.\]
The exponential is a finite polynomial in the marker algebra.
\(\rho_{\mathcal T}\) is the chessboard root at coupling \(\beta_{N,p}\)
over the tiling by cells of the width stage \(m_0(b_k,k)\).

\emph{Status:} theorem. The threshold in \(k\) depends on
\((d,z,\Delta,e,u_0,g,K,n)\) and on the bands only through \(u_0\). This
uniformity is the difference with the fixed-band
\href{https://github.com/yang-mills-group/OpenAI-YM-20261028/blob/3456e5e75cf8d4cd558955a5c7638fba2d68f464/YM/MergeProof/West/Tael.lean\#L58-L90}{\texttt{earnings}}
of III.11.

\emph{Lean:}
\href{https://github.com/yang-mills-group/OpenAI-YM-20261028/blob/3456e5e75cf8d4cd558955a5c7638fba2d68f464/YM/TraceExtension/Prom/Front/West/Tael.lean\#L56-L91}{\texttt{earningsSky}}.
The form for poised families is
\href{https://github.com/yang-mills-group/OpenAI-YM-20261028/blob/3456e5e75cf8d4cd558955a5c7638fba2d68f464/YM/TraceExtension/Prom/Orion/Reach.lean\#L34-L50}{\texttt{lava\_vSky}}.

\textbf{Proposition VII.6.4 (fixed-field supply in degrees \(\le3\)).}
Let \(Y\) be a final record for \((C,F,S)\) (VIII.1) and \(G\) simple.
Let \(\varphi\) be a filter on \(\mathbb N^2\) finer than \(\Phi_F\),
let \(d\le3\) and let \(x\in\mathrm{Label}_d\). Then \(\pi_dx\) is
\emph{supplied along \(\varphi\)}. That is, there is \(p\in\mathbb N\)
such that for all large \(k\) there are \(y\in\mathrm{Label}_d\) and
labels \(q_{N,p'}\in\mathrm{Label}_d\) with:

\begin{itemize}
\tightlist
\item
  \(\pi_dy=\pi_dx\);
\item
  \(q\to y\) along \(\varphi\);
\item
  \(\mathrm{quota}_k^{N,p'}(q_{N,p'})\le\hat\beta_k^{\,p}\) eventually
  along \(\varphi\).
\end{itemize}

Here \(\mathrm{quota}_k^{N,p'}(u)\) is the coefficient size of \(u\) at
cutoff \((N,p')\), written in stage-\(k\) coordinates
(\href{https://github.com/yang-mills-group/OpenAI-YM-20261028/blob/3456e5e75cf8d4cd558955a5c7638fba2d68f464/YM/Next/UV_STRESS/Least.lean\#L20-L21}{\texttt{quota}}).
One can take \(p=1\) and \(q\) constant.

\emph{Status:} theorem.

\emph{Lean:}
\href{https://github.com/yang-mills-group/OpenAI-YM-20261028/blob/3456e5e75cf8d4cd558955a5c7638fba2d68f464/YM/SharpGram/Clauses.lean\#L126-L129}{\texttt{inletFirst}},
\href{https://github.com/yang-mills-group/OpenAI-YM-20261028/blob/3456e5e75cf8d4cd558955a5c7638fba2d68f464/YM/SharpGram/Clauses.lean\#L18-L21}{\texttt{Supplies}}.

\textbf{Proposition VII.6.5 (fields of degree below four).} Let \(G\) be
simple and let \(d\in\mathbb N\) satisfy \(d<4\).

\begin{itemize}
\tightlist
\item
  \begin{enumerate}
  \def\labelenumi{(\alph{enumi})}
  \tightlist
  \item
    \(\dim\mathcal F^{\rm free}_d=1\).
  \end{enumerate}
\item
  (b) If \(Y\) is a final record for \((C,F,S)\), then every
  \(u\in\mathcal F^S_d\) is of the form \(t\,\mathbf 1_d\) with
  \(t\in\mathbb C\).
\end{itemize}

\emph{Status:} theorem. Used by VII.6.4 and by
\href{https://github.com/yang-mills-group/OpenAI-YM-20261028/blob/3456e5e75cf8d4cd558955a5c7638fba2d68f464/YM/Trace2/Transfer.lean\#L94-L111}{\texttt{vacuum\_truncation}}.

\emph{Lean:} (a)
\href{https://github.com/yang-mills-group/OpenAI-YM-20261028/blob/3456e5e75cf8d4cd558955a5c7638fba2d68f464/YM/UVNext/FREE_STRESS_CERTIFICATE/Size.lean\#L51-L54}{\texttt{dim\_before}};
(b)
\href{https://github.com/yang-mills-group/OpenAI-YM-20261028/blob/3456e5e75cf8d4cd558955a5c7638fba2d68f464/YM/Trace2/Transfer.lean\#L31-L37}{\texttt{first\_grade}}.

\subsubsection{2. What it uses}\label{what-it-uses-30}

\textbf{From other units.}

\begin{itemize}
\tightlist
\item
  \textbf{III.11 (interacting rank).} The stage estimates and numerical
  bounds:

  \begin{itemize}
  \tightlist
  \item
    The fixed-band estimates underlying the moving-band argument:

    \begin{itemize}
    \tightlist
    \item
      \href{https://github.com/yang-mills-group/OpenAI-YM-20261028/blob/3456e5e75cf8d4cd558955a5c7638fba2d68f464/YM/MergeProof/West/Tael.lean\#L58-L90}{\texttt{earnings}},
      \href{https://github.com/yang-mills-group/OpenAI-YM-20261028/blob/3456e5e75cf8d4cd558955a5c7638fba2d68f464/YM/MergeProof/East/Reach.lean\#L27-L42}{\texttt{lava\_v}};
    \item
      \href{https://github.com/yang-mills-group/OpenAI-YM-20261028/blob/3456e5e75cf8d4cd558955a5c7638fba2d68f464/YM/MergeProof/East/Plume.lean\#L23-L27}{\texttt{Poised}},
      \href{https://github.com/yang-mills-group/OpenAI-YM-20261028/blob/3456e5e75cf8d4cd558955a5c7638fba2d68f464/YM/MergeProof/East/Medal.lean\#L36-L76}{\texttt{aspect}},
      \href{https://github.com/yang-mills-group/OpenAI-YM-20261028/blob/3456e5e75cf8d4cd558955a5c7638fba2d68f464/YM/MergeProof/East/Largo.lean\#L21-L29}{\texttt{podium}}.
    \end{itemize}
  \item
    The reading-scale arithmetic:

    \begin{itemize}
    \tightlist
    \item
      \href{https://github.com/yang-mills-group/OpenAI-YM-20261028/blob/3456e5e75cf8d4cd558955a5c7638fba2d68f464/YM/Next/RANK/Station.lean\#L49}{\texttt{Station.gap}},
      \href{https://github.com/yang-mills-group/OpenAI-YM-20261028/blob/3456e5e75cf8d4cd558955a5c7638fba2d68f464/YM/Next/RANK/Station.lean\#L47}{\texttt{mag}};
    \item
      the rounding
      \href{https://github.com/yang-mills-group/OpenAI-YM-20261028/blob/3456e5e75cf8d4cd558955a5c7638fba2d68f464/YM/Next/RANK/Tuning.lean\#L24}{\texttt{tested}};
    \item
      the rates
      \href{https://github.com/yang-mills-group/OpenAI-YM-20261028/blob/3456e5e75cf8d4cd558955a5c7638fba2d68f464/YM/Next/RANK/Station.lean\#L55-L56}{\texttt{hat\_rate}}
      and
      \href{https://github.com/yang-mills-group/OpenAI-YM-20261028/blob/3456e5e75cf8d4cd558955a5c7638fba2d68f464/YM/Next/RANK/Station.lean\#L60-L88}{\texttt{gap\_rate}}.
    \end{itemize}
  \item
    The budget
    \href{https://github.com/yang-mills-group/OpenAI-YM-20261028/blob/3456e5e75cf8d4cd558955a5c7638fba2d68f464/YM/Next/RANK/Allowance.lean\#L28}{\texttt{Ship.budget}}
    and the pattern of the root bound
    \href{https://github.com/yang-mills-group/OpenAI-YM-20261028/blob/3456e5e75cf8d4cd558955a5c7638fba2d68f464/YM/Next/RANK/Allowance.lean\#L95-L128}{\texttt{Ship.rebate}}
    (III.11, Lemma 4.7).
  \item
    The step from root bounds to a pointwise comparison,
    \href{https://github.com/yang-mills-group/OpenAI-YM-20261028/blob/3456e5e75cf8d4cd558955a5c7638fba2d68f464/YM/Next/RANK/Finite.lean\#L242-L250}{\texttt{Read.full\_cost}}.
  \item
    The one-step cost package
    \href{https://github.com/yang-mills-group/OpenAI-YM-20261028/blob/3456e5e75cf8d4cd558955a5c7638fba2d68f464/YM/MergeProof/Biscuit.lean\#L29-L34}{\texttt{Scrim}}.
  \item
    The band inequalities
    \href{https://github.com/yang-mills-group/OpenAI-YM-20261028/blob/3456e5e75cf8d4cd558955a5c7638fba2d68f464/YM/MergeProof/East/Medal.lean\#L23-L28}{\texttt{Checks}}
    and the vane lift
    \href{https://github.com/yang-mills-group/OpenAI-YM-20261028/blob/3456e5e75cf8d4cd558955a5c7638fba2d68f464/YM/MergeProof/East/Pint.lean\#L25-L31}{\texttt{wok}}.
  \end{itemize}
\item
  \textbf{III.2 (choice of constants).} Production of the uniform-band
  property by
  \href{https://github.com/yang-mills-group/OpenAI-YM-20261028/blob/3456e5e75cf8d4cd558955a5c7638fba2d68f464/YM/TraceExtension/Prom/Order/State.lean\#L33-L50}{\texttt{embalms}},
  \href{https://github.com/yang-mills-group/OpenAI-YM-20261028/blob/3456e5e75cf8d4cd558955a5c7638fba2d68f464/YM/TraceExtension/Prom/Orion/Plume.lean\#L30-L53}{\texttt{stagingSky}}
  and the package extension of III.2. Also the construction runs
  \href{https://github.com/yang-mills-group/OpenAI-YM-20261028/blob/3456e5e75cf8d4cd558955a5c7638fba2d68f464/YM/Integration/RG05Preconstruction/Bardo.lean\#L26-L81}{\texttt{plentifulWithPrediction}},
  \href{https://github.com/yang-mills-group/OpenAI-YM-20261028/blob/3456e5e75cf8d4cd558955a5c7638fba2d68f464/YM/R2I/Campaign/Bardo.lean\#L24-L77}{\texttt{plentiful}}
  and
  \href{https://github.com/yang-mills-group/OpenAI-YM-20261028/blob/3456e5e75cf8d4cd558955a5c7638fba2d68f464/YM/TraceExtension/Made/Bardo.lean\#L25-L76}{\texttt{merano}}.
  These runs:

  \begin{itemize}
  \tightlist
  \item
    choose \(z_d=\sum_{i\le d}c_i\);
  \item
    choose an ultrafilter \(\mathcal L\) finer than \(\Phi_F\);
  \item
    invoke the uniform-band construction and then the gallery theorem of
    VII.2.
  \end{itemize}
\item
  \textbf{III.9, III.10, III.11.} The RG history, background fields and
  one-step bounds used in VII.6.g.
\item
  \textbf{IV.1.}

  \begin{itemize}
  \tightlist
  \item
    Extraction orders and limit tables.
  \item
    The chessboard estimate (Lemma IV.1.b; reflection positivity of the
    Wilson action).
  \item
    The corrected rank widths
    \href{https://github.com/yang-mills-group/OpenAI-YM-20261028/blob/3456e5e75cf8d4cd558955a5c7638fba2d68f464/YM/Cells/Corrected.lean\#L138}{\texttt{Cells.RankWidths}}
    that give hypothesis (iv).
  \end{itemize}
\item
  \textbf{VI.1, VI.2.} The scaling record
  \href{https://github.com/yang-mills-group/OpenAI-YM-20261028/blob/3456e5e75cf8d4cd558955a5c7638fba2d68f464/YM/Climb/Cells.lean\#L24-L31}{\texttt{Climb.Scaling}}.
  The free Maxwell table
  \href{https://github.com/yang-mills-group/OpenAI-YM-20261028/blob/3456e5e75cf8d4cd558955a5c7638fba2d68f464/YM/Euclid/Comparator.lean\#L23-L24}{\texttt{freeTable}}.
  The rest of the station machinery:

  \begin{itemize}
  \tightlist
  \item
    stations and admissibility
    \href{https://github.com/yang-mills-group/OpenAI-YM-20261028/blob/3456e5e75cf8d4cd558955a5c7638fba2d68f464/YM/Next/UV_STRESS/Station.lean\#L89-L97}{\texttt{Station.Admissible}}
    (VI.2.8);
  \item
    balanced estimates give stations:
    \href{https://github.com/yang-mills-group/OpenAI-YM-20261028/blob/3456e5e75cf8d4cd558955a5c7638fba2d68f464/YM/Next/UV_STRESS/Granular.lean\#L30-L34}{\texttt{Granular}},
    \href{https://github.com/yang-mills-group/OpenAI-YM-20261028/blob/3456e5e75cf8d4cd558955a5c7638fba2d68f464/YM/Next/UV_STRESS/Granular.lean\#L121-L128}{\texttt{supply\_wheel}}
    (VI.2.s);
  \item
    \href{https://github.com/yang-mills-group/OpenAI-YM-20261028/blob/3456e5e75cf8d4cd558955a5c7638fba2d68f464/YM/Next/UV_STRESS/Rates.lean\#L52-L53}{\texttt{Approaches}}
    and
    \href{https://github.com/yang-mills-group/OpenAI-YM-20261028/blob/3456e5e75cf8d4cd558955a5c7638fba2d68f464/YM/Next/UV_STRESS/Weak.lean\#L24-L27}{\texttt{Fades}};
  \item
    eventual injectivity
    \href{https://github.com/yang-mills-group/OpenAI-YM-20261028/blob/3456e5e75cf8d4cd558955a5c7638fba2d68f464/YM/Next/UV_STRESS/Invert.lean\#L88-L103}{\texttt{MeterTest.injects}};
  \item
    the dimension cap
    \href{https://github.com/yang-mills-group/OpenAI-YM-20261028/blob/3456e5e75cf8d4cd558955a5c7638fba2d68f464/YM/Next/UV_STRESS/Invert.lean\#L108}{\texttt{DimAbove}}
    and its proof
    \href{https://github.com/yang-mills-group/OpenAI-YM-20261028/blob/3456e5e75cf8d4cd558955a5c7638fba2d68f464/YM/Next/UV_STRESS/Invert.lean\#L110-L115}{\texttt{cap\_given}};
  \item
    the fixed-radius analogue
    \href{https://github.com/yang-mills-group/OpenAI-YM-20261028/blob/3456e5e75cf8d4cd558955a5c7638fba2d68f464/YM/Next/UV_STRESS/Invert.lean\#L132-L181}{\texttt{Stations.equivalences}}
    (VI.2.t);
  \item
    the free labels of degree \(<4\):
    \href{https://github.com/yang-mills-group/OpenAI-YM-20261028/blob/3456e5e75cf8d4cd558955a5c7638fba2d68f464/YM/UVNext/FREE_STRESS_CERTIFICATE/Size.lean\#L29-L44}{\texttt{just\_one}},
    \href{https://github.com/yang-mills-group/OpenAI-YM-20261028/blob/3456e5e75cf8d4cd558955a5c7638fba2d68f464/YM/UVNext/FREE_STRESS_CERTIFICATE/Size.lean\#L46-L49}{\texttt{below}},
    \href{https://github.com/yang-mills-group/OpenAI-YM-20261028/blob/3456e5e75cf8d4cd558955a5c7638fba2d68f464/YM/UVNext/FREE_STRESS_CERTIFICATE/Low.lean\#L133-L143}{\texttt{span\_four}}.
  \end{itemize}
\item
  \textbf{VIII.1.} The final record
  \href{https://github.com/yang-mills-group/OpenAI-YM-20261028/blob/3456e5e75cf8d4cd558955a5c7638fba2d68f464/YM/Spec/Closed.lean\#L64-L73}{\texttt{Spec.Final}}.
  Part B uses three of its fields:

  \begin{itemize}
  \tightlist
  \item
    the construction record;
  \item
    the hierarchy laws, giving \(\mathbf 1_d\ne0\);
  \item
    the comparator, a linear isomorphism
    \(\mathcal F^S_d\cong\mathcal F^{\rm free}_d\) for each \(d\).
  \end{itemize}
\item
  \textbf{VII.2 (consumer).}
  \href{https://github.com/yang-mills-group/OpenAI-YM-20261028/blob/3456e5e75cf8d4cd558955a5c7638fba2d68f464/YM/TraceExtension/Chart/Line.lean\#L49-L56}{\texttt{gang}}
  proves the last clause of
  \href{https://github.com/yang-mills-group/OpenAI-YM-20261028/blob/3456e5e75cf8d4cd558955a5c7638fba2d68f464/YM/TraceExtension/Chart/Line.lean\#L41-L44}{\texttt{Gallery}}
  from VII.6.1. That clause says: for every null sequence \(s\) there
  are a subsequence and a scaling with radii \(s\circ f\) carrying the
  trace-comparison witness.
\end{itemize}

\textbf{Standard mathematics.}

\begin{itemize}
\tightlist
\item
  Reflection positivity of the Wilson action (Osterwalder--Seiler, Ann.
  Phys. 110 (1978)) and chessboard estimates
  (Fröhlich--Israel--Lieb--Simon, CMP 62 (1978)).
\item
  Bałaban's block-spin renormalization group with background fields
  (CMP, 1983--1989). It is used only through the stage estimates of
  III.9--III.11.
\item
  The intermediate value theorem and recursive (diagonal) choice of
  subsequences.
\item
  Ultrafilter limits in finite-dimensional spaces (bounded sequences
  converge along any ultrafilter).
\item
  For simple \(\mathfrak g\): the adjoint representation has no nonzero
  invariant vectors, and invariant symmetric bilinear forms on
  \(\mathfrak g\) are multiples of the Killing form. Both are used only
  through VI.2.
\end{itemize}

\textbf{Imported constants and assumptions.}

\begin{itemize}
\tightlist
\item
  Hypotheses carried by VII.6.1: (S), the uniform-band property, the
  dimension cap and the ultrafilter. Their discharge is described above.
\item
  The budget \(\mathrm{bud}_n(k)=\exp(-\log\mu_k/(16n))\). The
  cost-record exponent \(-3\).
\item
  The mark exponents \(c_d=2(w_d+1)\)
  (\href{https://github.com/yang-mills-group/OpenAI-YM-20261028/blob/3456e5e75cf8d4cd558955a5c7638fba2d68f464/YM/MergeProof/SteamW.lean\#L18}{\texttt{generator}},
  \href{https://github.com/yang-mills-group/OpenAI-YM-20261028/blob/3456e5e75cf8d4cd558955a5c7638fba2d68f464/YM/MergeProof/Tender.lean\#L11}{\texttt{winter}}).
\item
  Bands at least \(d+6\), and the inequalities between three bands in
  VII.6.d.
\item
  Arity bound \(d+2\) for degree \(d\)
  (\href{https://github.com/yang-mills-group/OpenAI-YM-20261028/blob/3456e5e75cf8d4cd558955a5c7638fba2d68f464/YM/MergeProof/East/Gran.lean\#L25}{\texttt{trays}}).
\item
  In the stopping rule: the margin \(\mu_k^2\) and the depth
  \(\lceil(\log\hat\beta_k)^2\rceil^{10}\).
\item
  In the checked condition: tolerance \(1/(j+1)\) with window radius
  \(j+1\), and window radius \(100\) for the two-point bound.
\item
  The error \(\hat\beta_i^{-1/32}\) in Lemma VII.6.j.
\end{itemize}

\subsubsection{3. Strategy}\label{strategy-30}

For fixed \(b\), the reading scales \(\ell(b,k)\) form a discrete null
sequence. The fixed-band argument of III.11 also requires a plateau of
constant block factor across its RG window. These restrictions do not
ensure that its radii meet a prescribed null sequence. VII.2 needs such
a sequence to test the trace anomaly at arbitrary short distances.

Let the band vary over \([b_0,b_0+1]\), with \(b_0\ge0\). Eventually
\(\mu_k\le\mu_{k+1}\) and \(L_k\le\mu_{k+1}\), so
\[\ell(b_0,k)=a_{k+1}L_k\mu_k^{b_0}\le a_{k+1}\mu_{k+1}^{b_0+1}=\ell(b_0+1,k+1).\]
The reading intervals therefore overlap and cover all sufficiently small
positive radii. Continuity in the band gives exact equality with each
target radius. The stage estimates must now be uniform over a unit
interval of bands, and their stopping rule must allow a change of block
factor. The uniform-band input comes from III.2; the modified estimates
occupy most of this unit.

The chessboard bound gives a pointwise comparison with free
correlations. Two auxiliary bands control the strip near the coincidence
sets and the Riemann-sum error on its complement. These bounds yield
admissible stations at the target radii. A diagonal choice makes their
depth and arity increase while their tolerance tends to zero. Their
bounded lift tuples converge along a single ultrafilter; projection to
the limit field spaces gives approaching linear maps. Free test
functionals make these maps eventually injective, and the dimension cap
makes them isomorphisms. For degrees below four, a separate argument
transports the one-dimensional free field spaces through the final
comparator. Multiples of the unit then have constant lattice
representatives with the required coefficient bound.

\subsubsection{4. Steps}\label{steps-30}

\paragraph{4.1 Reading scales and dynamic
stops}\label{reading-scales-and-dynamic-stops}

\textbf{Lemma VII.6.a (stopping stages and fit).} Fix a band \(b\) and a
stage \(k\). The middle stage \(m_1(b,k)\) is the largest \(j\le k\) for
which \(a_j\ge\mu_k^2\ell(b,k)\), with value \(0\) if no such \(j\)
exists. Put \(Q_k=\lceil(\log\hat\beta_k)^2\rceil\). The width stage
\(m_0(b,k)\) is the largest \(j\le m_1\) for which
\(a_j\ge Q_k^{10}a_{m_1}\), again with value \(0\) if the set is empty.

The pair \((b,k)\) is \emph{allowed} when both defining inequalities
hold at the selected stages,
\[\min(\Lambda_{J_k},\Lambda_{J_k-1})\ge2,\qquad J_m\ge J_k-1\quad(m_0\le m\le k).\]
Monotonicity also gives \(J_m\le J_k\). Thus each block factor in this
window is at least \(\lambda_k=\min(\Lambda_{J_k},\Lambda_{J_k-1})\),
and \[\lambda_k^{j-m}\le a_m/a_j\qquad(m_0\le m\le j\le k).\]

Round \(x\) to the stage-\(k\) lattice in units of \(\ell(b,k)\), and
multiply by \(\ell(b,k)\) to obtain its physical configuration. It
\emph{fits} if that configuration is generic and every coordinate has
absolute value at most \(a_{m_1}/8\). For fixed \(b>0\) and generic
\(x\), it fits for every sufficiently large allowed \(k\). More
uniformly, for each \(R\in\mathbb R\) and arity \(n\), every generic
rescaled lattice configuration of norm at most \(R\) fits for
sufficiently large allowed \(k\), uniformly in \(b\).

\emph{Proof sketch.} The stopping indices are maxima of finite sets. The
allowed conditions give the two scale margins. Since
\(J_m\in\{J_k-1,J_k\}\), multiplying the lower block-factor bound gives
the ratio estimate. Rounding converges for a fixed positive band; the
margin \(\mu_k^2\to\infty\) then absorbs each fixed configuration. For
rescaled lattice configurations, rounding is exact, so the same margin
absorbs \(R\) uniformly in the band.

On an allowed plateau of III.11, the maximum definitions agree with its
ceiling-based stopping stages: \(a_j/a_k=L^{k-j}\) turns the defining
inequalities into the same ceiling inequalities.

Inputs:
\href{https://github.com/yang-mills-group/OpenAI-YM-20261028/blob/3456e5e75cf8d4cd558955a5c7638fba2d68f464/YM/TraceExtension/Prom/Hill/Tuning.lean\#L20-L38}{\texttt{readySky}},
\href{https://github.com/yang-mills-group/OpenAI-YM-20261028/blob/3456e5e75cf8d4cd558955a5c7638fba2d68f464/YM/TraceExtension/Prom/Orion/Suds.lean\#L34-L48}{\texttt{open\_readSky}}.

\emph{Lean:}
\href{https://github.com/yang-mills-group/OpenAI-YM-20261028/blob/3456e5e75cf8d4cd558955a5c7638fba2d68f464/YM/TraceExtension/Prom/Hill/Station.lean\#L15-L16}{\texttt{gainSky}},
\href{https://github.com/yang-mills-group/OpenAI-YM-20261028/blob/3456e5e75cf8d4cd558955a5c7638fba2d68f464/YM/TraceExtension/Prom/Hill/Station.lean\#L31}{\texttt{middleSky}},
\href{https://github.com/yang-mills-group/OpenAI-YM-20261028/blob/3456e5e75cf8d4cd558955a5c7638fba2d68f464/YM/TraceExtension/Prom/Hill/Station.lean\#L32}{\texttt{widthSky}},
\href{https://github.com/yang-mills-group/OpenAI-YM-20261028/blob/3456e5e75cf8d4cd558955a5c7638fba2d68f464/YM/TraceExtension/Prom/Hill/Station.lean\#L42-L46}{\texttt{AllowedSky}},
\href{https://github.com/yang-mills-group/OpenAI-YM-20261028/blob/3456e5e75cf8d4cd558955a5c7638fba2d68f464/YM/TraceExtension/Prom/Hill/Tuning.lean\#L15-L17}{\texttt{FitSky}};
fixed versions
\href{https://github.com/yang-mills-group/OpenAI-YM-20261028/blob/3456e5e75cf8d4cd558955a5c7638fba2d68f464/YM/Next/RANK/Station.lean\#L44}{\texttt{middle}},
\href{https://github.com/yang-mills-group/OpenAI-YM-20261028/blob/3456e5e75cf8d4cd558955a5c7638fba2d68f464/YM/Next/RANK/Station.lean\#L93-L96}{\texttt{Allowed}}.
\emph{(proof read; stopping definitions read)}

\textbf{Lemma VII.6.b (overlapping reading intervals).} Let \(G\) be
simple, \(b\ge0\), and assume
\(\max(\Lambda_{J_k},\Lambda_{J_k-1})\le\mu_k\) for all sufficiently
large \(k\). Then eventually \(\ell(b,k)\le\ell(b+1,k+1)\). For every
\(J\in\mathbb N\) there is \(r>0\) such that each \(s\in(0,r]\) equals
\(\ell(b+c,k)\) for some \(k\ge J\) and \(c\in[0,1]\).

\emph{Proof sketch.} Eventually \(\hat\beta_k\ge e\) and
\(\mu_k\le\mu_{k+1}\). The factor \(L_k\) is among the two factors
controlled at stage \(k+1\), so \(L_k\le\mu_{k+1}\). This gives overlap.
Choose a sufficiently large \(i\ge J\) and put \(r=\ell(b,i)\). For
\(0<s\le r\), take the least \(k>i\) with \(\ell(b,k)\le s\). Then
\[\ell(b,k)\le s\le\ell(b,k-1)\le\ell(b+1,k).\] Continuity of
\(c\mapsto\ell(b+c,k)\) gives \(c\).

\emph{Lean:}
\href{https://github.com/yang-mills-group/OpenAI-YM-20261028/blob/3456e5e75cf8d4cd558955a5c7638fba2d68f464/YM/TraceExtension/Prom/Orion/Ramp.lean\#L14-L18}{\texttt{tide\_mode}},
\href{https://github.com/yang-mills-group/OpenAI-YM-20261028/blob/3456e5e75cf8d4cd558955a5c7638fba2d68f464/YM/TraceExtension/Prom/Orion/Ramp.lean\#L19-L26}{\texttt{near\_skies}},
\href{https://github.com/yang-mills-group/OpenAI-YM-20261028/blob/3456e5e75cf8d4cd558955a5c7638fba2d68f464/YM/TraceExtension/Prom/Orion/Ramp.lean\#L27-L31}{\texttt{gap\_curve}},
\href{https://github.com/yang-mills-group/OpenAI-YM-20261028/blob/3456e5e75cf8d4cd558955a5c7638fba2d68f464/YM/TraceExtension/Prom/Orion/Ramp.lean\#L32-L43}{\texttt{gap\_account}}.

\textbf{Lemma VII.6.c (a null sequence is realized at increasing
stages).} Under the hypotheses of VII.6.b, let \(s_j>0\) with
\(s_j\to0\), and let \(J\in\mathbb N\). Then there are strictly
increasing \(K,P:\mathbb N\to\mathbb N\) and \(c_j\in[0,1]\) with
\(P_j\ge J\) and \(s_{K_j}=\ell(b+c_j,P_j)\).

\emph{Proof sketch.} Recursion: apply VII.6.b with threshold
\(\max(J,P_{j-1})+1\), and pick \(K_j>K_{j-1}\) with \(s_{K_j}\le r\).

\emph{Lean:}
\href{https://github.com/yang-mills-group/OpenAI-YM-20261028/blob/3456e5e75cf8d4cd558955a5c7638fba2d68f464/YM/TraceExtension/Prom/Orion/Ramp.lean\#L45-L53}{\texttt{gap\_orders}}.

\paragraph{4.2 The uniform-band
property}\label{the-uniform-band-property}

\textbf{Definition VII.6.d (packages, poised stages and the uniform-band
property).} A package of degree \(\le D\) consists of a band
\(b\ge D+6\), a Taylor order \(\theta\ge D\), an exponent \(e\), and a
purse of numerical cost parameters at each price index. Every purse has
band \(b\), Taylor order \(\theta\), exponent \(e\), and credit at least
\(bD+D+2\).

A finite family of packages is \emph{poised at stage \(k\)}, with arity
bound \(n\), mark exponent \(z\) and tables
\(R_{N,p,d}\in\mathrm{End}_{\mathbb C}(\mathrm{Label}_d)\) for
\(d\le D\), when:

\begin{itemize}
\tightlist
\item
  \(T_d\le k\) for every \(d\le D\);
\item
  \(\max(\Lambda_{J_k},\Lambda_{J_k-1})\le\mu_k\);
\item
  for every \(N>k\) and \(p\), the tables send the reference unit label
  to its frame unit vector, preserve the real structure, commute with
  signed coordinate permutations, and approximate the actual stage-\(k\)
  coordinate maps to relative coefficient-size error at most \(1/2\);
\item
  the sharper relative error bound \(\mu_k^{-1/2}\) also holds;
\item
  for each package there is a stage-\(k\) cost record with exponent
  \(-3\), mark exponent \(z\) and arity bound \(n\); for every \(N>k\)
  and \(p\), it has a compatible RG representative and polynomial lift
  data realizing the same tables.
\end{itemize}

The cost record also contains the allowed stopping stages of VII.6.a,
numerical and background-field bounds along the RG window, room bounds
for each arity at most \(n\), and the one-step estimates used in
VII.6.g.

The \emph{uniform-band property} \(\mathrm{Fire}_d(z)\) asserts that
there are package families \(P_s^{(i)}\), \(i=0,1,2\), base bands
\(b_0,b_1,b_2\), and \(\kappa<0\), satisfying
\[b_1+6<b_0,\quad b_2+5<b_0+\kappa,\quad (b_0-b_1)d(d+2)+\kappa<0,\quad (b_0-b_2)d(d+2)-b_0-\kappa<0.\]
For \(s\in[0,1]\), the bands are \(b_i+s\), and each family's Taylor
order and exponent are independent of \(s\). There is a single \(k_0\)
such that, for every \(k\ge k_0\) and every \(s\in[0,1]\), the triple is
poised at stage \(k\) with arity bound \(d+2\) and mark exponent \(z\),
for some tables.

\emph{Production in III.2.} Starting from its package and block-size
requests, a local chart ratio and nonnegative mark exponents, the
construction enlarges a waiting window so that every family paid on the
enlarged window, with the required covariance of its coordinate maps and
a binding to the same prices, satisfies \(\mathrm{Fire}_d(z_d)\) for
every \(d\). Packages chosen at bands \(b_i+1\) are lowered to \(b_i+s\)
without changing their Taylor orders or exponents. The poised-stage
threshold is uniform in \(s\in[0,1]\).

\emph{Lean:}
\href{https://github.com/yang-mills-group/OpenAI-YM-20261028/blob/3456e5e75cf8d4cd558955a5c7638fba2d68f464/YM/Next/RANK/Prepaid.lean\#L37-L47}{\texttt{Tapes.Panel}},
\href{https://github.com/yang-mills-group/OpenAI-YM-20261028/blob/3456e5e75cf8d4cd558955a5c7638fba2d68f464/YM/TraceExtension/Prom/Orion/Plume.lean\#L23-L28}{\texttt{PoisedSky}},
\href{https://github.com/yang-mills-group/OpenAI-YM-20261028/blob/3456e5e75cf8d4cd558955a5c7638fba2d68f464/YM/TraceExtension/Prom/Front/West/Lys.lean\#L38-L45}{\texttt{LegacyCostSky}},
\href{https://github.com/yang-mills-group/OpenAI-YM-20261028/blob/3456e5e75cf8d4cd558955a5c7638fba2d68f464/YM/TraceExtension/Prom/Order/State.lean\#L29-L32}{\texttt{FiresSky}},
\href{https://github.com/yang-mills-group/OpenAI-YM-20261028/blob/3456e5e75cf8d4cd558955a5c7638fba2d68f464/YM/MergeProof/East/Medal.lean\#L23-L28}{\texttt{Checks}},
\href{https://github.com/yang-mills-group/OpenAI-YM-20261028/blob/3456e5e75cf8d4cd558955a5c7638fba2d68f464/YM/MergeProof/East/Gran.lean\#L14-L22}{\texttt{raises}},
\href{https://github.com/yang-mills-group/OpenAI-YM-20261028/blob/3456e5e75cf8d4cd558955a5c7638fba2d68f464/YM/TraceExtension/Prom/Order/State.lean\#L33-L50}{\texttt{embalms}}.
\emph{(definitions read; production proof read)}

\textbf{Lemma VII.6.e (leaps).} Let \(G\) be simple and assume
\(\mathrm{Fire}_d(z)\) with base bands \(b_0,b_1,b_2\). Then \(b_0>5\),
and the following holds for every \(s_i>0\) with \(s_i\to0\). There are:

\begin{itemize}
\tightlist
\item
  a strictly increasing \(f\);
\item
  strictly increasing stages \(k_i\ge k_0\);
\item
  offsets \(e_k\in[0,1]\);
\item
  tables \(Y_i\);
\end{itemize}

such that the triple \(P_{e_{k_i}}\) is poised at stage \(k_i\) with
\(Y_i\), and \(s_{f(i)}=\ell(b_0+e_{k_i},k_i)\).

\emph{Proof sketch.}

\begin{enumerate}
\def\labelenumi{\arabic{enumi}.}
\tightlist
\item
  \(b_0\ge d+6\).
\item
  The poised condition gives
  \(\max(\Lambda_{J_k},\Lambda_{J_k-1})\le\mu_k\), which is the
  hypothesis of VII.6.b.
\item
  Apply VII.6.c with base \(b_0\) and threshold \(k_0\). Put
  \(e_{P_j}=c_j\) and \(e=0\) elsewhere.
\item
  \(\mathrm{Fire}_d\) at \(s=c_j\) gives poised data at stage
  \(P_j\ge k_0\). Here the threshold must not depend on \(s\).
\end{enumerate}

\emph{Lean:}
\href{https://github.com/yang-mills-group/OpenAI-YM-20261028/blob/3456e5e75cf8d4cd558955a5c7638fba2d68f464/YM/TraceExtension/Prom/Order/State.lean\#L68-L84}{\texttt{fliesSky}},
\href{https://github.com/yang-mills-group/OpenAI-YM-20261028/blob/3456e5e75cf8d4cd558955a5c7638fba2d68f464/YM/TraceExtension/Prom/Order/State.lean\#L56-L60}{\texttt{LeapSky}},
\href{https://github.com/yang-mills-group/OpenAI-YM-20261028/blob/3456e5e75cf8d4cd558955a5c7638fba2d68f464/YM/TraceExtension/Prom/Orion/Reach.lean\#L21-L24}{\texttt{ArcSky}}.

\paragraph{4.3 Stage estimates with a moving
band}\label{stage-estimates-with-a-moving-band}

\textbf{Lemma VII.6.f (= Proposition VII.6.3).}

\emph{Proof sketch.} Follow the fixed-band argument of III.11 with three
uniform estimates and the modified stopping stages.

\begin{enumerate}
\def\labelenumi{\arabic{enumi}.}
\tightlist
\item
  The local probe bounds hold uniformly for bands \(\le u_0\).
\item
  The polynomial lifts of the \(g_i\) have bounds independent of the
  band.
\item
  The mark exponent \(z\ge c_{d_i}\) converts normalized coefficient
  bounds on the \(f_i\) into source-expansion bounds
  (\href{https://github.com/yang-mills-group/OpenAI-YM-20261028/blob/3456e5e75cf8d4cd558955a5c7638fba2d68f464/YM/MergeProof/SteamW.lean\#L19-L41}{\texttt{tremor}}).
\item
  Per-site affordability holds from some stage \(u\) on. It is needed at
  the width stage \(m_0\).
\item
  \(u\le m_0\) for large \(k\). If \(m_0<u\), the bracket condition of
  VII.6.a would give \(J_k-1\le J_{m_0}\le J_u\), contradicting
  \(J_k>J_u+1\). In the fixed version the plateau is used here, with
  \(J_k>J_u\).
\item
  Lemma VII.6.g then gives the bound.
\end{enumerate}

\emph{Lean:}
\href{https://github.com/yang-mills-group/OpenAI-YM-20261028/blob/3456e5e75cf8d4cd558955a5c7638fba2d68f464/YM/TraceExtension/Prom/Front/West/Tael.lean\#L56-L91}{\texttt{earningsSky}},
\href{https://github.com/yang-mills-group/OpenAI-YM-20261028/blob/3456e5e75cf8d4cd558955a5c7638fba2d68f464/YM/TraceExtension/Prom/Front/West/Tael.lean\#L40-L42}{\texttt{ActumSky}}.

\textbf{Lemma VII.6.g (root bound at one stage).} Fix a cost record of
degree bound \(D\), band \(b\ge0\), purse exponent \(e\), mark exponent
\(z\) and arity bound \(u\ge2\). Fix \(N\ge k\) and a period \(p\), let
\(d_i\le D\) for \(0\le i\le m\), and assume \(m+1\le u\), \(T_D\le k\)
and \(T_{d_i}\le k\). Take an RG history from the cutoff to \(m_0(b,k)\)
satisfying the cost record, the per-site affordability bound at \(m_0\),
the required local work items, and the implication from the stage-\(k\)
seal to its covariance model. Let \(g_i,f_i\in\mathrm{Label}_{d_i}\) and
polynomial lift data satisfy the rescaling equation of VII.6.3. Assume
the configuration fits, its rounded position belongs to the specified
compact set, and the uniform local probe bounds hold there. Write
\(\gamma\) for the purse credit and \(M_{k,N,p,d}\) for the normalized
stage-coordinate coefficient size multiplied by the stage multiplicity
bound. Finally assume:

\begin{itemize}
\tightlist
\item
  \(\mu_k^e\,M_{k,N,p,d_i}(f_i)\le\mu_k^{-8}\) for each probe;
\item
  its lifted polynomial has coefficient sum at most \(\mu_k^\gamma\);
\item
  separation of the probe seats implies the source-expansion bound with
  weight \(g_k^z\), where \(g_k=\hat\beta_k^{-1/2}\).
\end{itemize}

Then \(\rho_{\mathcal T}(\Phi_s)\le\mathrm{bud}_u(k)\) for every
nonempty \(s\).

\emph{Proof sketch.} The history gives compatible background fields and
marker estimates at each stage. The covariance model supplies the free
comparison data, and VII.6.h supplies the source expansion at the middle
stage. Propagating that expansion to the width stage gives a marker
estimate with allowance \(A=e^{2I}\), where \(I=\log\mu_k/(16u)\). The
per-site affordability and allocated cost imply
\[\rho_{\mathcal T}(\Phi_s)\le e^I/A^{|s|}\le e^{-I}=\mathrm{bud}_u(k).\]

Root estimate:
\href{https://github.com/yang-mills-group/OpenAI-YM-20261028/blob/3456e5e75cf8d4cd558955a5c7638fba2d68f464/YM/TraceExtension/Prom/Port/Soft/Denarii.lean\#L31-L62}{\texttt{couponSky}};
fixed-band counterpart:
\href{https://github.com/yang-mills-group/OpenAI-YM-20261028/blob/3456e5e75cf8d4cd558955a5c7638fba2d68f464/YM/Next/RANK/Allowance.lean\#L95-L128}{\texttt{Ship.rebate}}.

\emph{Lean:}
\href{https://github.com/yang-mills-group/OpenAI-YM-20261028/blob/3456e5e75cf8d4cd558955a5c7638fba2d68f464/YM/TraceExtension/Prom/Front/West/Lys.lean\#L51-L71}{\texttt{y\_readySky}},
\href{https://github.com/yang-mills-group/OpenAI-YM-20261028/blob/3456e5e75cf8d4cd558955a5c7638fba2d68f464/YM/TraceExtension/Prom/Hill/Allowance.lean\#L26-L28}{\texttt{yardSky}},
\href{https://github.com/yang-mills-group/OpenAI-YM-20261028/blob/3456e5e75cf8d4cd558955a5c7638fba2d68f464/YM/TraceExtension/Prom/Hill/Allowance.lean\#L40-L41}{\texttt{invoiceSky}},
\href{https://github.com/yang-mills-group/OpenAI-YM-20261028/blob/3456e5e75cf8d4cd558955a5c7638fba2d68f464/YM/Next/RANK/Allowance.lean\#L28}{\texttt{Ship.budget}}.

\textbf{Lemma VII.6.h (source expansion at the middle stage).} Fix the
history, polynomial lift data and labels of VII.6.g. Assume its launch
and survivor data, the compatible background-field record, the room
bound for \(m+1\) probes at depth \(k-m_1\), and the leading and
nonleading marker estimates along the window. If the configuration fits,
its probe seats are separated, and the normalized coefficient,
polynomial coefficient and source-expansion bounds of VII.6.g hold, then
the marked probe product has the required source expansion at \(m_1\).

More precisely, for every admissible middle-stage numerical bound, there
are finite source and background families and a distinguished source for
each cell and nonempty mode subset. They reproduce the propagated marked
product exactly. Each distinguished source has the prescribed support
and mode set; on a healthy background it differs from the corresponding
free-cumulant expression by at most \(\mu_k^{-3/2}\), after the
permitted covariance adjustment.

\emph{Proof sketch.} The polynomial coefficient bound controls the probe
family. The launch and survivor data give the marked-source
representation. The room, fit and marker estimates propagate it to the
middle stage, where the support and comparison bounds give the stated
expansion.

\emph{Lean:}
\href{https://github.com/yang-mills-group/OpenAI-YM-20261028/blob/3456e5e75cf8d4cd558955a5c7638fba2d68f464/YM/TraceExtension/Prom/Front/West/Delft.lean\#L30-L38}{\texttt{attireSky}}.

\textbf{Lemma VII.6.i (assembly of weighted words).} In a fixed one-step
stage configuration, assume the background is healthy, the room and
one-step bounds hold with size constant \(B\ge1\), and
\(0<\varepsilon\le1\). Let \(P_u\) be polynomials in the mode variables
invariant under simultaneous gauge rotations, each of total degree at
most \(D\) and coefficient sum at most \(r\ge0\), with all variables in
the permitted probe regions. Let \(S\) be a finite list of pairs
\((c,w)\), where \(c\in\mathbb R\) and \(w\) is a word of allowed modes,
each supported in the selected cell. Then the assembled chart integrands
are integrable, and
\[\left|\mathcal A\left(\sum_{(c,w)\in S}c\prod_{u\in w}P_u\right)-\mathcal R\left(\sum_{(c,w)\in S}c\prod_{u\in w}P_u\right)\right|\le\sum_{(c,w)\in S}|c|\,E(B,\varepsilon,r,D,|w|).\]
Here \(\mathcal A\) is the chart assembly, \(\mathcal R\) is the
corresponding reference reading with the same cell-orientation test, and
\(E\) is the one-word error bound from III.11.

\emph{Proof sketch.} Induct on the list. Integrability is preserved by
addition and scalar multiplication. Linearity and the triangle
inequality add the one-word errors, with the factor \(|c|\) for each
word.

\emph{Lean:}
\href{https://github.com/yang-mills-group/OpenAI-YM-20261028/blob/3456e5e75cf8d4cd558955a5c7638fba2d68f464/YM/TraceExtension/Prom/Front/West/Karin.lean\#L34-L57}{\texttt{sweepSky}}.

\textbf{Lemma VII.6.j (one step at a change of block factor).} Fix a
step \(i+1\to i\) with \(i<N\), a coarse-field pose \((o,W,V,v)\), a
width multiplier \(w\in\mathbb N\), and its auxiliary block-spin
choices, forest and fluctuation data. Assume the one-step coupling
bounds, with the larger of \(\Omega_{J_{i+1}},\Omega_{J_i}\) and the
smaller of \(\alpha_{J_{i+1}},\alpha_{J_i}\). Assume also the
small-field and Hessian-control bounds on the patch. Put
\(x=\hat\beta_i\), let \(\widetilde x\) be the auxiliary reference
coupling used in the stationary-field data, let \(y\) be the fine
reference coupling, and set \(L=wB(y)\), where \(B\) is the block radius
of III.11.

Let \(r\) be a fine-lattice site. Assume the fluctuation guards on the
enlarged patch about \(r\),
\[2(L+1)<\mathrm{side}_{i+1,p},\qquad 100(1+R_B(\widetilde x))\le\mathrm{side}_{i,p},\]
and, for every guarded fine bond \(e\), require its parent to lie inside
the pose with the margin
\[\mathrm{dist}(o,\mathrm{parent}(e))+B(x)+2\le W.\] Then the fine
configuration \(U\) constructed by the background-field map is healthy
at the small-field threshold \(\tau(y)\) on the radius-\(L\) patch. For
every \(f\) and gauge transformation \(t\) with \(t\cdot f=U\), and
every coloured incidence \(j\) with \(\mathrm{dist}(r,j)+2\le L\),
\[\left|\mathcal O_j(f,t;r,g)-\bigl(\mathrm{Re}\,m_j+(DZ)_j\bigr)\right|\le\hat\beta_i^{-1/32}.\]
Here \(R_B(s)=\lceil(\log s)^2\rceil\), \(g\) is the step coupling, and
\(\mathcal O_j\) is the colour component of \(g^{-1}\) times the
plaquette curvature, in tree gauge rooted at \(r\) and rotated by the
adjoint action of \(t(r)\). The mean \(m\) and linear fluctuation
correction \(DZ\) use the same coordinates.

\emph{Proof sketch.} The Hessian-control bound and guards give link
estimates throughout the enlarged patch. A four-bond estimate makes
every tested plaquette small, proving health. The link estimate and
stationary-field curl identity give the displayed error on incidences
two lattice units inside the patch.

\emph{Lean:}
\href{https://github.com/yang-mills-group/OpenAI-YM-20261028/blob/3456e5e75cf8d4cd558955a5c7638fba2d68f464/YM/TraceExtension/Prom/Port/Rail.lean\#L95-L127}{\texttt{laminateSky}},
used by
\href{https://github.com/yang-mills-group/OpenAI-YM-20261028/blob/3456e5e75cf8d4cd558955a5c7638fba2d68f464/YM/TraceExtension/Prom/Port/Emulation.lean\#L33-L71}{\texttt{screensSky}}.

\textbf{Lemma VII.6.k (poised families supply the stage bound).} Assume
(S), simplicity of \(G\), and \(z\ge c_i\) for \(i\le D\). Let \(P_k\)
be a finite family of degree-\(D\) packages whose bands are
\(q_\alpha+e_k\), with \(e_k\in[0,1]\), and whose Taylor orders and
exponents are independent of \(k\). Fix a package index \(\alpha\),
labels \(y_i\in\mathrm{Label}_{d_i}\) with \(d_i\le D\), a compact
\(K\subset\mathrm{Conf}_{m+1}\), and \(2\le n\), \(m+1\le n\). For every
sufficiently large \(k\), every table with which \(P_k\) is poised,
every \(N>k\) and \(p\), and every fitted configuration whose rounded
position belongs to \(K\), the bound of VII.6.3 holds for the probes
obtained by the vane lift from \(y_i\).

\emph{Proof sketch.} Apply VII.6.3 with band ceiling \(q_\alpha+1\). The
poised data supply the cost record, RG representative and polynomial
lifts. Inverting the charted table gives the rescaling equation.

\emph{Lean:}
\href{https://github.com/yang-mills-group/OpenAI-YM-20261028/blob/3456e5e75cf8d4cd558955a5c7638fba2d68f464/YM/TraceExtension/Prom/Orion/Reach.lean\#L34-L50}{\texttt{lava\_vSky}}.

\textbf{Lemma VII.6.l (pointwise comparison).} Under the hypotheses of
VII.6.k, fix an extraction order \(F\), a package index \(\alpha\), a
strip exponent \(\chi\), \(R>0\), and a free word of arity \(m\le n\)
with degrees at most \(D\). Assume the free table is smooth. There is
\(C\) such that, for every sufficiently large \(k\) and every poised
table, the lattice \(m\)-point function of the lifted probes differs
from the free function by at most \(C\,\mathrm{bud}_n(k)\) in the
ordered upper-limit sense, at each generic stage-\(k\) lattice
configuration in the rescaled window of radius \(2R\).

Thus, for each configuration and each \(\eta>0\), the difference is
eventually at most \(C\,\mathrm{bud}_n(k)+\eta\) along \(F\); no common
cutoff threshold over all configurations is asserted.

\emph{Proof sketch.} VII.6.k bounds every nonempty recentred part. The
chessboard inequality and moment--cumulant identity bound the difference
of moments by the budget times the sum of absolute coefficients of the
free moment polynomial. Smoothness bounds that finite sum on the compact
window. At arity zero both moments equal \(1\).

Inputs:
\href{https://github.com/yang-mills-group/OpenAI-YM-20261028/blob/3456e5e75cf8d4cd558955a5c7638fba2d68f464/YM/Next/RANK/Finite.lean\#L242-L250}{\texttt{Read.full\_cost}}.

\emph{Lean:}
\href{https://github.com/yang-mills-group/OpenAI-YM-20261028/blob/3456e5e75cf8d4cd558955a5c7638fba2d68f464/YM/TraceExtension/Prom/Orion/Suds.lean\#L51-L95}{\texttt{sightSky}}.

\textbf{Lemma VII.6.m (balanced estimates).} Under the hypotheses of
VII.6.k, suppose three base bands \(q_x,q_y,q_w\) and \(c<0\) satisfy
\[q_y+6<q_x,\quad q_w+5<q_x+c,\quad(q_x-q_y)Dn+c<0,\quad(q_x-q_w)Dn-q_x-c<0.\]
Assume the free hierarchy is given. Let \(f(i)\to\infty\), and let
\(Y_i\) be tables for which \(P_{f(i)}\) is poised. Then the stations at
stage \(f(i)\), main band \(q_x+e_{f(i)}\) and strip exponent \(c\)
satisfy the balanced estimates of VI.2 for arity at most \(n\) and every
window radius \(R>0\).

\emph{Proof sketch.} Write \(k=f(i)\). The pointwise error is
\(O(\mathrm{bud}_n(k))\to0\). The auxiliary package at band \(q_y+e_k\)
supplies the strip bound \[A_k=O\bigl(\mu_k^{(q_x-q_y)D}\bigr),\] and
the package at \(q_w+e_k\) supplies the bound used for interpolation off
the strip, \[B_k=O\bigl(\mu_k^{(q_x-q_w)D}\bigr).\] Each auxiliary
package supplies both types of local bound; its designated role is
chosen in the balancing argument. Their stopping stages fit inside the
main geometric window by the first two inequalities. With
\(\sigma_k=\mu_k^c\) and normalized spacing
\(h_k=\mu_k^{-(q_x+e_k)}\le\mu_k^{-q_x}\), the last two inequalities
give \[A_k^n\sigma_k\to0,\qquad B_k^n h_k/\sigma_k\to0.\] The radii tend
to zero, and the lift bounds and symmetries hold on ordered tails.
VI.2's station theorem gives eventual admissibility with any positive
tolerance.

Inputs:
\href{https://github.com/yang-mills-group/OpenAI-YM-20261028/blob/3456e5e75cf8d4cd558955a5c7638fba2d68f464/YM/TraceExtension/Prom/Orion/Latch.lean\#L52-L77}{\texttt{thirdSky}},
\href{https://github.com/yang-mills-group/OpenAI-YM-20261028/blob/3456e5e75cf8d4cd558955a5c7638fba2d68f464/YM/TraceExtension/Prom/Orion/Pint.lean\#L37-L57}{\texttt{w\_rightSky}},
\href{https://github.com/yang-mills-group/OpenAI-YM-20261028/blob/3456e5e75cf8d4cd558955a5c7638fba2d68f464/YM/Next/UV_STRESS/Granular.lean\#L121-L128}{\texttt{supply\_wheel}}.

\emph{Lean:}
\href{https://github.com/yang-mills-group/OpenAI-YM-20261028/blob/3456e5e75cf8d4cd558955a5c7638fba2d68f464/YM/TraceExtension/Prom/Orion/Medal.lean\#L39-L90}{\texttt{aspectSky}},
\href{https://github.com/yang-mills-group/OpenAI-YM-20261028/blob/3456e5e75cf8d4cd558955a5c7638fba2d68f464/YM/TraceExtension/Prom/Orion/Largo.lean\#L24-L34}{\texttt{podiumSky}}.

\paragraph{4.4 Checked vanes}\label{checked-vanes}

\textbf{Definition VII.6.n (vanes and checked vanes).} A \emph{vane of
depth \(D\)} consists of a stage \(k\), band \(b\), strip exponent
\(c\), and complex-linear tables \(R_{N,p,d}\) on \(\mathrm{Label}_d\)
for \(d\le D\).

Its station has radius \(\ell(b,k)\), strip width \(\mu_k^c\), and lifts
\[A_{N,p,d}=R_{N,p,d}^{-1}\circ D_{k,N,p,d}^{-1}\circ W_{\mu_k^b,d}\circ\iota_d.\]
Here \(\iota_d:\mathcal F^{\rm free}_d\to\mathrm{Label}_d\) is the
canonical equivariant section, \(W\) applies the position scaling
weight, and \(D\) is the coupling dilation. On a frame vector of
bidegree \((h,q)\), \(W_{\rho,d}\) multiplies by \(\rho^h\), where \(h\)
is the engineering degree, while \(D_{k,N,p,d}\) multiplies by
\(s_{k,N,p}^{q}\), where \(q\) is the number of curvature factors and
\[s_{k,N,p}=\frac{\beta_{N,p;k}^{-1/2}}{\hat\beta_k^{-1/2}\mu_k}.\] Here
\(\beta_{N,p;k}\) is the actual coupling at stage \(k\). If
\(R_{N,p,d}\) is not bijective, the lift is zero. Charting ensures
bijectivity on the cutoff tails used below.

The vane is \emph{checked at level \(j\)} when:

\begin{enumerate}
\def\labelenumi{\arabic{enumi}.}
\tightlist
\item
  \(5<b\le\log\log\hat\beta_k\);
\item
  \(T_u\le k\) for every \(u\le D\);
\item
  for every \(N>k\) and \(p\), its tables obey the charting and sharper
  coordinate bounds in VII.6.d;
\item
  every generic rescaled stage-\(k\) lattice pair of norm at most
  \(100\) fits; for every \(w\le D\), every pair of free basis vectors
  in degree \(w\), every \(N>k\) and \(p\), and every fitted pair whose
  rounded configuration is in the radius-\(100\) window, all nonempty
  recentred parts have chessboard root at most \(1/(j+1)\);
\item
  its station is admissible for arity at most \(j\), window radius
  \(j+1\), and tolerance \(1/(j+1)\).
\end{enumerate}

\emph{Lean:}
\href{https://github.com/yang-mills-group/OpenAI-YM-20261028/blob/3456e5e75cf8d4cd558955a5c7638fba2d68f464/YM/TraceExtension/Prom/Zoom/Boats.lean\#L65-L69}{\texttt{Vane}},
\href{https://github.com/yang-mills-group/OpenAI-YM-20261028/blob/3456e5e75cf8d4cd558955a5c7638fba2d68f464/YM/TraceExtension/Prom/Zoom/Boats.lean\#L71}{\texttt{Vane.wheel}},
\href{https://github.com/yang-mills-group/OpenAI-YM-20261028/blob/3456e5e75cf8d4cd558955a5c7638fba2d68f464/YM/MergeProof/East/Pint.lean\#L25-L31}{\texttt{wok}},
\href{https://github.com/yang-mills-group/OpenAI-YM-20261028/blob/3456e5e75cf8d4cd558955a5c7638fba2d68f464/YM/MergeProof/East/Ladle.lean\#L17-L19}{\texttt{ladd}},
\href{https://github.com/yang-mills-group/OpenAI-YM-20261028/blob/3456e5e75cf8d4cd558955a5c7638fba2d68f464/YM/Further/CORRECTED_LIFT_COEFFICIENTS/Basis.lean\#L115-L117}{\texttt{ladder}},
\href{https://github.com/yang-mills-group/OpenAI-YM-20261028/blob/3456e5e75cf8d4cd558955a5c7638fba2d68f464/YM/TraceExtension/Prom/Zoom/Boats.lean\#L21-L29}{\texttt{SeqPair}},
\href{https://github.com/yang-mills-group/OpenAI-YM-20261028/blob/3456e5e75cf8d4cd558955a5c7638fba2d68f464/YM/TraceExtension/Prom/Zoom/Boats.lean\#L42-L47}{\texttt{TimberSky}},
\href{https://github.com/yang-mills-group/OpenAI-YM-20261028/blob/3456e5e75cf8d4cd558955a5c7638fba2d68f464/YM/TraceExtension/Prom/Zoom/Boats.lean\#L72}{\texttt{Vane.checked}}.
\emph{(definitions read)}

\textbf{Lemma VII.6.o (a checked vane below any radius).} Assume
simplicity, (S), \(z\ge c_i\) for \(i\le d\), \(\mathrm{Fire}_d(z)\),
and the limit-table hypotheses of VII.6.2. Let \(l\le d+2\), \(s_i>0\)
with \(s_i\to0\), \(Y\in\mathbb N\) and \(\theta>0\). There are an index
\(i\ge Y\), a stage \(k\ge Y\), and a vane of depth \(d\), stage \(k\),
checked at level \(l\), with \[\ell(b,k)=s_i<\theta.\]

\emph{Proof sketch.} VII.6.e gives increasing stages realizing a
subsequence of \(s\) at bands in \([b_0,b_0+1]\). The two-point
condition follows from VII.6.k applied to basis pairs, with arity budget
\(d+2\), together with uniform fitting and
\(\mathrm{bud}_{d+2}(k)\to0\). There are only finitely many degrees and
basis pairs. VII.6.m and the station theorem give admissibility through
arity \(d+2\), hence through arity \(l\), at window radius \(l+1\) and
tolerance \(1/(l+1)\). Choose a stage beyond the reference stages and
band ceiling, with both indices at least \(Y\) and radius below
\(\theta\).

Inputs:
\href{https://github.com/yang-mills-group/OpenAI-YM-20261028/blob/3456e5e75cf8d4cd558955a5c7638fba2d68f464/YM/TraceExtension/Prom/Zoom/Boats.lean\#L31-L39}{\texttt{await}}.

\emph{Lean:}
\href{https://github.com/yang-mills-group/OpenAI-YM-20261028/blob/3456e5e75cf8d4cd558955a5c7638fba2d68f464/YM/TraceExtension/Prom/Zoom/Boats.lean\#L49-L62}{\texttt{boatsSky}}.

\textbf{Lemma VII.6.p (= Proposition VII.6.2).}

\emph{Proof sketch.}

\begin{enumerate}
\def\labelenumi{\arabic{enumi}.}
\tightlist
\item
  For each \(d\), \(N\) and \(y>0\), VII.6.o with \(l=d\),
  \(Y=\max(d,N+1)\) and \(\theta=y\) gives \(k>N\) with \(s_k<y\), and a
  vane of depth \(d\), stage \(\ge d\), checked at level \(d\), with
  radius \(s_k\).
\item
  Choose recursively \(p(d+1)>p(d)\) with \(s_{p(d+1)}<s_{p(d)}\).
\item
  The stages are \(\ge d\), so they tend to infinity.
\end{enumerate}

\emph{Lean:}
\href{https://github.com/yang-mills-group/OpenAI-YM-20261028/blob/3456e5e75cf8d4cd558955a5c7638fba2d68f464/YM/TraceExtension/Prom/Zoom/Boats.lean\#L75-L87}{\texttt{wakesSky}}.

\paragraph{4.5 Limits along the
ultrafilter}\label{limits-along-the-ultrafilter}

\textbf{Lemma VII.6.q (limits of station lifts).} Under the
ordered-limit and hierarchy hypotheses of VII.6.1, let a station of
depth \(D\), with radius \(r>0\), be admissible through arity
\(n_{\max}\), on the radius-\(R\) window, with tolerance
\(\varepsilon\). Let \(\mathcal L\) be an ultrafilter finer than
\(\Phi_F\). Assume its degree-\(d\) lift sends
\(\mathbf 1_d^{\rm free}\) to the scalar label eventually along
\(\Phi_F\), for every \(d\le D\). Then the lift tuple converges along
\(\mathcal L\) to
\[Y\in\prod_{d\le D}\mathrm{Hom}_{\mathbb C}(\mathcal F_d^{\rm free},\mathrm{Label}_d).\]
The maps \(U_d=\pi_d\circ Y_d\) are real, signed-covariant and
unit-preserving for \(d\le D\). For \(n\le n_{\max}\), \(d_i\le D\),
\(v_i\in\mathcal F^{\rm free}_{d_i}\), \(M>0\), and every smooth test
\(t\) supported in the radius-\(R\) window with \(\|t\|_\infty\le M\),
\[\left|\int t(x)\left(\langle Uv_1(rx_1)\cdots Uv_n(rx_n)\rangle^S-\langle v_1(x_1)\cdots v_n(x_n)\rangle^{\rm free}\right)\,dx\right|\le\varepsilon M\prod_i\|v_i\|.\]

\emph{Proof sketch.}

\begin{enumerate}
\def\labelenumi{\arabic{enumi}.}
\tightlist
\item
  The lifts are eventually bounded, so a nontrivial filter
  \(\mathfrak b\) finer than \(\mathcal L\) along which they converge
  exists
  (\href{https://github.com/yang-mills-group/OpenAI-YM-20261028/blob/3456e5e75cf8d4cd558955a5c7638fba2d68f464/YM/Next/UV_STRESS/Vary.lean\#L74-L80}{\texttt{finite\_extract}}).
\item
  Since \(\mathcal L\) is an ultrafilter, \(\mathfrak b=\mathcal L\).
\item
  Symmetry and unit identities pass to the limit, as in VI.2's
  \href{https://github.com/yang-mills-group/OpenAI-YM-20261028/blob/3456e5e75cf8d4cd558955a5c7638fba2d68f464/YM/Next/UV_STRESS/Station.lean\#L132-L180}{\texttt{harvested}},
  and ordered smeared convergence gives the displayed estimate. The
  extra unit hypothesis covers every degree through \(D\).
\end{enumerate}

\emph{Lean:}
\href{https://github.com/yang-mills-group/OpenAI-YM-20261028/blob/3456e5e75cf8d4cd558955a5c7638fba2d68f464/YM/TraceExtension/Prom/Zoom/Pepper.lean\#L11-L60}{\texttt{harvestedSky}}.

\textbf{Lemma VII.6.r (vanes give approaching maps).} Under the
ordered-limit and hierarchy hypotheses of VII.6.1, let \(r_k\) be a
strictly decreasing null sequence. Let \(V_k\) be vanes of depth \(k\)
whose stages tend to infinity, checked at level \(k\), with
\(\ell(V_k)=r_k\), and let \(\mathcal L\) be an ultrafilter finer than
\(\Phi_F\). Then there are complex-linear maps
\(U_k(d):\mathcal F_d^{\rm free}\to\mathcal F_d^S\) such that:

\begin{itemize}
\tightlist
\item
  \((r_k,U_k)\) is approaching;
\item
  \(U_k(d)\) is real, signed-covariant and unit-preserving for
  \(d\le k\);
\item
  \(U_k(d)=\pi_d\circ Y_{k,d}\) for \(d\le k\), where \(Y_k\) is the
  \(\mathcal L\)-limit of the lifts of \(V_k\).
\end{itemize}

\emph{Proof sketch.}

\begin{enumerate}
\def\labelenumi{\arabic{enumi}.}
\tightlist
\item
  Conditions 2--3 of VII.6.n give unit preservation of the lifts in
  every degree through \(k\). VII.6.q then applies to each \(V_k\).
\item
  Fix \(n\), \(v\), a compact \(K\subset\mathrm{Conf}_n\), and
  \(M,\varepsilon\). For \(k\ge n\), \(k\ge\max d_i\), \(K\) contained
  in the radius-\((k+1)\) window, and
  \(\prod\|v_i\|M/(k+1)\le\varepsilon\), the estimate at tolerance
  \(1/(k+1)\) bounds \(|\int tK_k|\) by \(\varepsilon\).
\end{enumerate}

\emph{Lean:}
\href{https://github.com/yang-mills-group/OpenAI-YM-20261028/blob/3456e5e75cf8d4cd558955a5c7638fba2d68f464/YM/TraceExtension/Prom/Zoom/Mast.lean\#L13-L31}{\texttt{breadSky}}.

\textbf{Lemma VII.6.s (isomorphisms).} Assume in addition the dimension
cap (iv). Then there is a scaling with radii \(r_k\) and \(\xi_k=1\)
such that:

\begin{itemize}
\tightlist
\item
  \(U_k(d)\) coincides with the map of VII.6.r for all large \(k\) (for
  each \(d\));
\item
  (a)--(c) of VII.6.1 hold.
\end{itemize}

\emph{Proof sketch.}

\begin{enumerate}
\def\labelenumi{\arabic{enumi}.}
\tightlist
\item
  The maps are eventually injective in each degree by VI.2. Finitely
  many free test functionals separate \(\mathcal F^{\rm free}_d\), their
  rescaled versions converge by the approaching property, and
  injectivity is an open condition.
\item
  Injectivity and \(\dim\mathcal F^S_d\le\dim\mathcal F^{\rm free}_d\)
  give bijectivity.
\item
  Where the map is not yet bijective, one of the later isomorphisms in
  that degree is substituted. This does not affect any eventual
  property.
\end{enumerate}

\emph{Lean:}
\href{https://github.com/yang-mills-group/OpenAI-YM-20261028/blob/3456e5e75cf8d4cd558955a5c7638fba2d68f464/YM/TraceExtension/Prom/Zoom/Mast.lean\#L32-L58}{\texttt{dockedSky}}.

\textbf{Lemma VII.6.t (= Proposition VII.6.1).}

\emph{Proof sketch.} VII.6.p gives \(p\) and \(V\). The radii
\(s\circ p\) form a strictly decreasing null sequence. VII.6.s then
gives the scaling, with \(f=p\).

\emph{Lean:}
\href{https://github.com/yang-mills-group/OpenAI-YM-20261028/blob/3456e5e75cf8d4cd558955a5c7638fba2d68f464/YM/TraceExtension/Prom/Zoom/Running.lean\#L25-L31}{\texttt{pays\_sky}}.

\emph{Consumer.}
\href{https://github.com/yang-mills-group/OpenAI-YM-20261028/blob/3456e5e75cf8d4cd558955a5c7638fba2d68f464/YM/TraceExtension/Chart/Line.lean\#L49-L56}{\texttt{gang}}
(VII.2) applies VII.6.1 with:

\begin{itemize}
\tightlist
\item
  the seed's ordered limits, hierarchy and point limits;
\item
  the free hierarchy supplied by the free-field comparison theorem;
\item
  hypothesis (iv) from the rank-width bound.
\end{itemize}

It selects the lift limits and applies the trace-comparison theorem of
VII.2. The run
\href{https://github.com/yang-mills-group/OpenAI-YM-20261028/blob/3456e5e75cf8d4cd558955a5c7638fba2d68f464/YM/TraceExtension/Made/Bardo.lean\#L25-L76}{\texttt{merano}},
together with the other runs cited in §2, supplies the rest:

\begin{itemize}
\tightlist
\item
  \(z_d=\sum_{i\le d}c_i\);
\item
  \(\mathrm{Fire}_d(z_d)\) from the uniform-band construction;
\item
  \(\mathcal L\) from the existence of an ultrafilter finer than
  \(\Phi_F\).
\end{itemize}

\paragraph{4.6 Degrees below four}\label{degrees-below-four}

\textbf{Lemma VII.6.u (= Proposition VII.6.5).}

\emph{Proof sketch.}

\begin{itemize}
\tightlist
\item
  \begin{enumerate}
  \def\labelenumi{(\alph{enumi})}
  \tightlist
  \item
    For simple \(G\), every free polynomial label of degree \(d<4\) is
    \(c\cdot1\). The label included in the degree-\(4\) filtration is
    \(c\cdot1+Q\), where \(Q\) lies in the span of the quadratic
    invariants and is homogeneous of weight \(4\) (by the degree-four
    classification). A label of degree \(<4\) has no weight-\(4\) part,
    so \(Q=0\). Since labels map onto \(\mathcal F^{\rm free}_d\) and
    the unit is nonzero, \(\dim\mathcal F^{\rm free}_d=1\). For
    comparison, \(\dim\mathcal F^{\rm free}_4=22\)
    (\href{https://github.com/yang-mills-group/OpenAI-YM-20261028/blob/3456e5e75cf8d4cd558955a5c7638fba2d68f464/YM/UVNext/FREE_STRESS_CERTIFICATE/Size.lean\#L112-L114}{\texttt{four\_size}}).
  \end{enumerate}
\item
  (b) The comparator of the final record is an isomorphism
  \(\mathcal F^S_d\cong\mathcal F^{\rm free}_d\), so
  \(\dim\mathcal F^S_d=1\). Since \(\mathbf 1_d\ne0\),
  \(\mathcal F^S_d=\mathbb C\mathbf 1_d\).
\end{itemize}

\emph{Lean:}
\href{https://github.com/yang-mills-group/OpenAI-YM-20261028/blob/3456e5e75cf8d4cd558955a5c7638fba2d68f464/YM/UVNext/FREE_STRESS_CERTIFICATE/Size.lean\#L51-L54}{\texttt{dim\_before}},
\href{https://github.com/yang-mills-group/OpenAI-YM-20261028/blob/3456e5e75cf8d4cd558955a5c7638fba2d68f464/YM/Trace2/Transfer.lean\#L31-L37}{\texttt{first\_grade}}.

\textbf{Lemma VII.6.v (the unit is supplied).} Assume the construction
record for \(C\), let \(\varphi\) be finer than \(\Phi_F\), let \(G\) be
simple, and fix \(d\in\mathbb N\) and \(t\in\mathbb C\). Then
\(t\mathbf 1_d\) is supplied along \(\varphi\) with \(p=1\).

\emph{Proof sketch.} Let \(e_d^{\mathrm{ref}}\) be the reference unit
label. Take the constant family \(q\equiv y=t\,e_d^{\mathrm{ref}}\), for
which \(\pi_dy=t\mathbf 1_d\). For \(k\ge T_d\) and \(N\ge k\), the
quota is homogeneous and the unit has stage-\(k\) coordinate
\(a_k^d e_0\); this follows from the construction record. Hence the
quota is \(|t|\,a_k^d\). Along \(\varphi\) the cutoff eventually exceeds
\(k\), and for large \(k\) both \(a_k\le1\) and \(\hat\beta_k\ge|t|\)
hold.

\emph{Lean:}
\href{https://github.com/yang-mills-group/OpenAI-YM-20261028/blob/3456e5e75cf8d4cd558955a5c7638fba2d68f464/YM/SharpGram/Clauses.lean\#L54-L61}{\texttt{supplyUnit}},
\href{https://github.com/yang-mills-group/OpenAI-YM-20261028/blob/3456e5e75cf8d4cd558955a5c7638fba2d68f464/YM/SharpGram/Clauses.lean\#L33-L37}{\texttt{quota\_smul}},
\href{https://github.com/yang-mills-group/OpenAI-YM-20261028/blob/3456e5e75cf8d4cd558955a5c7638fba2d68f464/YM/SharpGram/Clauses.lean\#L38-L48}{\texttt{quota\_unit}}.

\textbf{Lemma VII.6.w (= Proposition VII.6.4).}

\emph{Proof sketch.} By VII.6.u(b), \(\pi_dx=t\mathbf 1_d\). Then
VII.6.v applies, with the construction record taken from the final
record.

\emph{Lean:}
\href{https://github.com/yang-mills-group/OpenAI-YM-20261028/blob/3456e5e75cf8d4cd558955a5c7638fba2d68f464/YM/SharpGram/Clauses.lean\#L126-L129}{\texttt{inletFirst}}.

\emph{Companion statements (VII.4).} The same file contains the dual
``release'' bounds:

\begin{itemize}
\tightlist
\item
  \href{https://github.com/yang-mills-group/OpenAI-YM-20261028/blob/3456e5e75cf8d4cd558955a5c7638fba2d68f464/YM/SharpGram/Clauses.lean\#L65-L85}{\texttt{coreRead}};
\item
  \href{https://github.com/yang-mills-group/OpenAI-YM-20261028/blob/3456e5e75cf8d4cd558955a5c7638fba2d68f464/YM/SharpGram/Clauses.lean\#L87-L117}{\texttt{physicalRelease}},
  which supplies an exponent \(p\) such that, for all large \(k\), every
  nontrivial finer filter along which \(q\to x\) and
  \(\mathrm{quota}_k(q)\le z\) eventually satisfies
  \(|L(\pi_dx)|\le\hat\beta_k^{\,p}a_k^{-(d-w)}z\). It assumes
  construction and work item data, simplicity of \(G\), \(w\le3\), and
  \(L(\mathbf 1_d)=0\) when \(w>0\);
\item
  \href{https://github.com/yang-mills-group/OpenAI-YM-20261028/blob/3456e5e75cf8d4cd558955a5c7638fba2d68f464/YM/SharpGram/Clauses.lean\#L120-L124}{\texttt{outletFirst}},
  where \(L\) vanishes on the images of \(\mathcal F^S_e\), \(e<w\),
  under the inclusions of the derivative plan.
\end{itemize}

These results belong to VII.4.

\subsubsection{5. Where the difficulty
lies}\label{where-the-difficulty-lies-30}

The analytic work extends the estimates of III.11 across a window with
two neighbouring price indices. The smaller block factor controls the
stopping depth; the larger and smaller neighbouring price parameters
control the one-step estimates. At the final root estimate,
affordability is charged at the width stage itself. These distinctions
replace uses of constant-factor plateaus.

The interval covering argument requires one stage threshold for every
band in a unit interval. Without that uniformity, the band selected by
the target radius could escape the available estimates. The subsequent
diagonal choice preserves exact equality with the prescribed radii.

The scaling maps are real, signed-covariant and unit-preserving
eventually in each degree. This result does not assert compatibility
with inclusions between degrees or stress-tensor normalization. The
dimension cap supplies surjectivity. The ultrafilter keeps all lift
limits on the same cutoff filter, and the checked tolerances give local
order-zero convergence without an explicit rate in the prescribed
radius. The fixed-field supply argument uses the comparator and the unit
coefficient identity; it requires none of the moving-band estimates.

\subsubsection{6. Notation}\label{notation-30}

{\def\LTcaptype{none} % do not increment counter
\begin{longtable}[]{@{}
  >{\raggedright\arraybackslash}p{(\linewidth - 4\tabcolsep) * \real{0.3333}}
  >{\raggedright\arraybackslash}p{(\linewidth - 4\tabcolsep) * \real{0.3333}}
  >{\raggedright\arraybackslash}p{(\linewidth - 4\tabcolsep) * \real{0.3333}}@{}}
\toprule\noalign{}
\begin{minipage}[b]{\linewidth}\raggedright
object
\end{minipage} & \begin{minipage}[b]{\linewidth}\raggedright
code name
\end{minipage} & \begin{minipage}[b]{\linewidth}\raggedright
proposed notation
\end{minipage} \\
\midrule\noalign{}
\endhead
\bottomrule\noalign{}
\endlastfoot
scaling (radii, constants, isomorphisms) &
\href{https://github.com/yang-mills-group/OpenAI-YM-20261028/blob/3456e5e75cf8d4cd558955a5c7638fba2d68f464/YM/Climb/Cells.lean\#L24-L31}{\texttt{Climb.Scaling}};
radii only:
\href{https://github.com/yang-mills-group/OpenAI-YM-20261028/blob/3456e5e75cf8d4cd558955a5c7638fba2d68f464/YM/Next/UV_STRESS/Rates.lean\#L15-L19}{\texttt{Zoom}}
& \((r_k,\xi_k,U_k)\) \\
approaching; weak vanishing on compacts &
\href{https://github.com/yang-mills-group/OpenAI-YM-20261028/blob/3456e5e75cf8d4cd558955a5c7638fba2d68f464/YM/Next/UV_STRESS/Rates.lean\#L52-L53}{\texttt{Approaches}},
\href{https://github.com/yang-mills-group/OpenAI-YM-20261028/blob/3456e5e75cf8d4cd558955a5c7638fba2d68f464/YM/Next/UV_STRESS/Weak.lean\#L24-L27}{\texttt{Fades}}
& \(K_k\to0\) locally with an order-zero bound \\
scaling sequence along \(\mathcal L\) (VII.6.1 (a)--(c)) &
\href{https://github.com/yang-mills-group/OpenAI-YM-20261028/blob/3456e5e75cf8d4cd558955a5c7638fba2d68f464/YM/TraceExtension/Prom/Zoom/Running.lean\#L19-L23}{\texttt{StarPath}}
& scaling sequence along \(\mathcal L\) \\
free Maxwell reference table &
\href{https://github.com/yang-mills-group/OpenAI-YM-20261028/blob/3456e5e75cf8d4cd558955a5c7638fba2d68f464/YM/Euclid/Comparator.lean\#L23-L24}{\texttt{freeTable}},
\href{https://github.com/yang-mills-group/OpenAI-YM-20261028/blob/3456e5e75cf8d4cd558955a5c7638fba2d68f464/YM/Climb/Cells.lean\#L20}{\texttt{Bel}},
\href{https://github.com/yang-mills-group/OpenAI-YM-20261028/blob/3456e5e75cf8d4cd558955a5c7638fba2d68f464/YM/Cells/Corrected.lean\#L22}{\texttt{Cells.standard}}
& \(\mathcal F^{\rm free}_d\), \(H_0\) \\
running coupling; growth scale &
\href{https://github.com/yang-mills-group/OpenAI-YM-20261028/blob/3456e5e75cf8d4cd558955a5c7638fba2d68f464/YM/RG/History.lean\#L79}{\texttt{hatBeta}},
\href{https://github.com/yang-mills-group/OpenAI-YM-20261028/blob/3456e5e75cf8d4cd558955a5c7638fba2d68f464/YM/RG/History.lean\#L81}{\texttt{monitor}},
\href{https://github.com/yang-mills-group/OpenAI-YM-20261028/blob/3456e5e75cf8d4cd558955a5c7638fba2d68f464/YM/Next/RANK/Reign.lean\#L25}{\texttt{Perch.bill}}
& \(\hat\beta_k\); \(\mu_k\); \(\mu_k^{a}\) \\
band; reading scale & \texttt{Panel.arc}, \texttt{Purse.band};
\href{https://github.com/yang-mills-group/OpenAI-YM-20261028/blob/3456e5e75cf8d4cd558955a5c7638fba2d68f464/YM/Next/RANK/Station.lean\#L49}{\texttt{Station.gap}}
& \(b\); \(\ell(b,k)=a_k\mu_k^b\) \\
middle and width stages; allowed; fit &
\href{https://github.com/yang-mills-group/OpenAI-YM-20261028/blob/3456e5e75cf8d4cd558955a5c7638fba2d68f464/YM/TraceExtension/Prom/Hill/Station.lean\#L31}{\texttt{middleSky}},
\href{https://github.com/yang-mills-group/OpenAI-YM-20261028/blob/3456e5e75cf8d4cd558955a5c7638fba2d68f464/YM/TraceExtension/Prom/Hill/Station.lean\#L32}{\texttt{widthSky}},
\href{https://github.com/yang-mills-group/OpenAI-YM-20261028/blob/3456e5e75cf8d4cd558955a5c7638fba2d68f464/YM/TraceExtension/Prom/Hill/Station.lean\#L42-L46}{\texttt{AllowedSky}},
\href{https://github.com/yang-mills-group/OpenAI-YM-20261028/blob/3456e5e75cf8d4cd558955a5c7638fba2d68f464/YM/TraceExtension/Prom/Hill/Tuning.lean\#L15-L17}{\texttt{FitSky}}
& \(m_1(b,k)\), \(m_0(b,k)\) \\
position weight; coupling dilation &
\href{https://github.com/yang-mills-group/OpenAI-YM-20261028/blob/3456e5e75cf8d4cd558955a5c7638fba2d68f464/YM/Next/RANK/Matrices.lean\#L89-L92}{\texttt{weightTag}},
\href{https://github.com/yang-mills-group/OpenAI-YM-20261028/blob/3456e5e75cf8d4cd558955a5c7638fba2d68f464/YM/Labels/Reference.lean\#L58-L61}{\texttt{Frame.scaled}},
\href{https://github.com/yang-mills-group/OpenAI-YM-20261028/blob/3456e5e75cf8d4cd558955a5c7638fba2d68f464/YM/Next/RANK/Matrices.lean\#L104}{\texttt{shrink}}
& \(W_{\rho,d}\); \(D_{k,N,p,d}\); \(s_{k,N,p}\) \\
normalized probe size &
\href{https://github.com/yang-mills-group/OpenAI-YM-20261028/blob/3456e5e75cf8d4cd558955a5c7638fba2d68f464/YM/Next/RANK/Insiders.lean\#L50}{\texttt{moneys}}
& \(M_{k,N,p,d}\) \\
stage package &
\href{https://github.com/yang-mills-group/OpenAI-YM-20261028/blob/3456e5e75cf8d4cd558955a5c7638fba2d68f464/YM/Next/RANK/Prepaid.lean\#L37-L47}{\texttt{Tapes.Panel}}
& \((b,\theta,e,\text{purses})\) \\
stage-\(k\) cost record; poised &
\href{https://github.com/yang-mills-group/OpenAI-YM-20261028/blob/3456e5e75cf8d4cd558955a5c7638fba2d68f464/YM/TraceExtension/Prom/Front/West/Lys.lean\#L38-L45}{\texttt{LegacyCostSky}},
\href{https://github.com/yang-mills-group/OpenAI-YM-20261028/blob/3456e5e75cf8d4cd558955a5c7638fba2d68f464/YM/TraceExtension/Prom/Orion/Plume.lean\#L23-L28}{\texttt{PoisedSky}}
& poised at stage \(k\) \\
uniform-band property &
\href{https://github.com/yang-mills-group/OpenAI-YM-20261028/blob/3456e5e75cf8d4cd558955a5c7638fba2d68f464/YM/TraceExtension/Prom/Order/State.lean\#L29-L32}{\texttt{FiresSky}}
& \(\mathrm{Fire}_d(z)\) \\
leap; packages with moving band &
\href{https://github.com/yang-mills-group/OpenAI-YM-20261028/blob/3456e5e75cf8d4cd558955a5c7638fba2d68f464/YM/TraceExtension/Prom/Order/State.lean\#L56-L60}{\texttt{LeapSky}},
\href{https://github.com/yang-mills-group/OpenAI-YM-20261028/blob/3456e5e75cf8d4cd558955a5c7638fba2d68f464/YM/TraceExtension/Prom/Orion/Reach.lean\#L21-L24}{\texttt{ArcSky}}
& --- \\
budget; chessboard root &
\href{https://github.com/yang-mills-group/OpenAI-YM-20261028/blob/3456e5e75cf8d4cd558955a5c7638fba2d68f464/YM/Next/RANK/Allowance.lean\#L28}{\texttt{Ship.budget}},
\href{https://github.com/yang-mills-group/OpenAI-YM-20261028/blob/3456e5e75cf8d4cd558955a5c7638fba2d68f464/YM/Cells/Tiling.lean\#L226-L227}{\texttt{rootTest}}
& \(\mathrm{bud}_n(k)=\mu_k^{-1/(16n)}\); \(\rho_{\mathcal T}\) \\
recentred probe product &
\href{https://github.com/yang-mills-group/OpenAI-YM-20261028/blob/3456e5e75cf8d4cd558955a5c7638fba2d68f464/YM/Next/RANK/Finite.lean\#L99}{\texttt{Read.cylinder}},
\href{https://github.com/yang-mills-group/OpenAI-YM-20261028/blob/3456e5e75cf8d4cd558955a5c7638fba2d68f464/YM/Next/RANK/Finite.lean\#L226}{\texttt{Read.cumulants}}
& \(\Phi_s\) \\
mark exponents &
\href{https://github.com/yang-mills-group/OpenAI-YM-20261028/blob/3456e5e75cf8d4cd558955a5c7638fba2d68f464/YM/MergeProof/SteamW.lean\#L18}{\texttt{generator}}
& \(c_d=2(w_d+1)\) \\
one-step cost package &
\href{https://github.com/yang-mills-group/OpenAI-YM-20261028/blob/3456e5e75cf8d4cd558955a5c7638fba2d68f464/YM/MergeProof/Biscuit.lean\#L29-L34}{\texttt{Scrim}}
& (S) \\
vane; its station; checked &
\href{https://github.com/yang-mills-group/OpenAI-YM-20261028/blob/3456e5e75cf8d4cd558955a5c7638fba2d68f464/YM/TraceExtension/Prom/Zoom/Boats.lean\#L65-L69}{\texttt{Vane}},
\href{https://github.com/yang-mills-group/OpenAI-YM-20261028/blob/3456e5e75cf8d4cd558955a5c7638fba2d68f464/YM/MergeProof/East/Pint.lean\#L25-L31}{\texttt{wok}},
\href{https://github.com/yang-mills-group/OpenAI-YM-20261028/blob/3456e5e75cf8d4cd558955a5c7638fba2d68f464/YM/TraceExtension/Prom/Zoom/Boats.lean\#L42-L47}{\texttt{TimberSky}},
\href{https://github.com/yang-mills-group/OpenAI-YM-20261028/blob/3456e5e75cf8d4cd558955a5c7638fba2d68f464/YM/TraceExtension/Prom/Zoom/Boats.lean\#L21-L29}{\texttt{SeqPair}}
& \(V=(k,b,c,R)\); checked at level \(j\) \\
dimension cap &
\href{https://github.com/yang-mills-group/OpenAI-YM-20261028/blob/3456e5e75cf8d4cd558955a5c7638fba2d68f464/YM/Next/UV_STRESS/Invert.lean\#L108}{\texttt{DimAbove}}
& \(\dim\mathcal F^S_d\le\dim\mathcal F^{\rm free}_d\) \\
gallery property (VII.2) &
\href{https://github.com/yang-mills-group/OpenAI-YM-20261028/blob/3456e5e75cf8d4cd558955a5c7638fba2d68f464/YM/TraceExtension/Chart/Line.lean\#L41-L44}{\texttt{Trace.Over.Gallery}}
& --- \\
stage-\(k\) coefficient size &
\href{https://github.com/yang-mills-group/OpenAI-YM-20261028/blob/3456e5e75cf8d4cd558955a5c7638fba2d68f464/YM/Next/UV_STRESS/Least.lean\#L20-L21}{\texttt{quota}}
& \(\mathrm{quota}_k^{N,p}(u)\) \\
supply; release; killing lower degrees &
\href{https://github.com/yang-mills-group/OpenAI-YM-20261028/blob/3456e5e75cf8d4cd558955a5c7638fba2d68f464/YM/SharpGram/Clauses.lean\#L18-L21}{\texttt{Supplies}},
\href{https://github.com/yang-mills-group/OpenAI-YM-20261028/blob/3456e5e75cf8d4cd558955a5c7638fba2d68f464/YM/SharpGram/Clauses.lean\#L23-L26}{\texttt{Releases}},
\href{https://github.com/yang-mills-group/OpenAI-YM-20261028/blob/3456e5e75cf8d4cd558955a5c7638fba2d68f464/YM/SharpGram/Clauses.lean\#L28-L29}{\texttt{Kills}}
& \(x\) supplied along \(\varphi\); \(L\) released at weight \(w\) \\
\end{longtable}
}

\subsubsection{7. What remains
unexplained}\label{what-remains-unexplained-30}

\begin{enumerate}
\def\labelenumi{\arabic{enumi}.}
\tightlist
\item
  \textbf{Internal analytic packages.} VII.6.g--j retain the room,
  background-field, source-expansion, small-field and Hessian-control
  hypotheses from III.9--III.11. Their direct implications and the final
  root estimate are verified, but the numerical definitions and proofs
  of all constituent one-step estimates have not been expanded here. In
  particular, the small-field input explicitly bounds the combined
  nonlinear errors by \(x^{-1/32}\); the eventual feasibility of that
  input remains part of III.11. The formula and guarded domain of
  VII.6.j are verified.
\item
  \textbf{Downstream trace argument.} The checked two-point bound is
  retained when the vanes and their lift limits are passed to VII.2. Its
  subsequent analytic use in deriving the trace anomaly is outside this
  unit. No such conclusion is inferred from the weak convergence alone.
\item
  \textbf{Construction inputs.} The passage from an enlarged paid
  waiting window and covariant coordinate maps to the uniform-band
  property is verified at the cited interfaces. The full construction
  that pays every numerical request belongs to III.2; it is not an
  additional conclusion of VII.6.1.
\end{enumerate}

\subsection{VIII.1. Assembly of the main
theorem}\label{viii.1.-assembly-of-the-main-theorem}

\textbf{Part:} VIII, Assembly of the main theorem. \textbf{Lean:} the
top-level files
\href{https://github.com/yang-mills-group/OpenAI-YM-20261028/blob/3456e5e75cf8d4cd558955a5c7638fba2d68f464/Solution.lean}{\texttt{Solution.lean}},
\href{https://github.com/yang-mills-group/OpenAI-YM-20261028/blob/3456e5e75cf8d4cd558955a5c7638fba2d68f464/Construction.lean}{\texttt{Construction.lean}}
and
\href{https://github.com/yang-mills-group/OpenAI-YM-20261028/blob/3456e5e75cf8d4cd558955a5c7638fba2d68f464/ConcreteSU2.lean}{\texttt{ConcreteSU2.lean}};
the bridges
\href{https://github.com/yang-mills-group/OpenAI-YM-20261028/blob/3456e5e75cf8d4cd558955a5c7638fba2d68f464/WilsonBridge.lean}{\texttt{WilsonBridge.lean}},
\href{https://github.com/yang-mills-group/OpenAI-YM-20261028/blob/3456e5e75cf8d4cd558955a5c7638fba2d68f464/HierarchyBridge.lean}{\texttt{HierarchyBridge.lean}},
\href{https://github.com/yang-mills-group/OpenAI-YM-20261028/blob/3456e5e75cf8d4cd558955a5c7638fba2d68f464/TraceBridge.lean}{\texttt{TraceBridge.lean}}
and
\href{https://github.com/yang-mills-group/OpenAI-YM-20261028/blob/3456e5e75cf8d4cd558955a5c7638fba2d68f464/ContactBridge.lean}{\texttt{ContactBridge.lean}};
the public definitions
\href{https://github.com/yang-mills-group/OpenAI-YM-20261028/blob/3456e5e75cf8d4cd558955a5c7638fba2d68f464/PublicTheory.lean}{\texttt{PublicTheory}},
\href{https://github.com/yang-mills-group/OpenAI-YM-20261028/blob/3456e5e75cf8d4cd558955a5c7638fba2d68f464/PublicWilson.lean}{\texttt{PublicWilson}},
\href{https://github.com/yang-mills-group/OpenAI-YM-20261028/blob/3456e5e75cf8d4cd558955a5c7638fba2d68f464/PublicView.lean}{\texttt{PublicView}},
\href{https://github.com/yang-mills-group/OpenAI-YM-20261028/blob/3456e5e75cf8d4cd558955a5c7638fba2d68f464/ChallengeSpec.lean}{\texttt{ChallengeSpec}}
and
\href{https://github.com/yang-mills-group/OpenAI-YM-20261028/blob/3456e5e75cf8d4cd558955a5c7638fba2d68f464/ContactSpec.lean}{\texttt{ContactSpec}};
the folders
\href{https://github.com/yang-mills-group/OpenAI-YM-20261028/tree/3456e5e75cf8d4cd558955a5c7638fba2d68f464/YM/Nonperturbative}{\texttt{YM/Nonperturbative}},
\href{https://github.com/yang-mills-group/OpenAI-YM-20261028/tree/3456e5e75cf8d4cd558955a5c7638fba2d68f464/YM/Independent}{\texttt{YM/Independent}},
\href{https://github.com/yang-mills-group/OpenAI-YM-20261028/tree/3456e5e75cf8d4cd558955a5c7638fba2d68f464/YM/Spec}{\texttt{YM/Spec}},
\href{https://github.com/yang-mills-group/OpenAI-YM-20261028/tree/3456e5e75cf8d4cd558955a5c7638fba2d68f464/YM/Org}{\texttt{YM/Org}}
and
\href{https://github.com/yang-mills-group/OpenAI-YM-20261028/tree/3456e5e75cf8d4cd558955a5c7638fba2d68f464/YM/Pert/Parity}{\texttt{YM/Pert/Parity}};
the files
\href{https://github.com/yang-mills-group/OpenAI-YM-20261028/blob/3456e5e75cf8d4cd558955a5c7638fba2d68f464/YM/R2I/Campaign/Accepted.lean}{\texttt{YM/R2I/Campaign/Accepted.lean}}
and
\href{https://github.com/yang-mills-group/OpenAI-YM-20261028/blob/3456e5e75cf8d4cd558955a5c7638fba2d68f464/YM/RG17/Joint.lean}{\texttt{YM/RG17/Joint.lean}};
and parts of YM/FullVolume, YM/EveryPhase, YM/TraceExtension, YM/OPE,
YM/Contacts and YM/Pert. \textbf{Size:} 12,060 lines in 1,674
declarations.

The construction chooses one Wilson family for each gauge group and
block-spin parameter \(b>0\). At each requested extraction order, it
assembles the continuum table, phases, fields, reconstruction and OPE on
one refinement. Fixed-volume ultraviolet convergence removes the mesh
subsequence. The bridges then give the public theorem in
\href{https://github.com/scottnarmstrong/yang-mills-informalization/blob/main/statement.md}{statement.md}.

\subsubsection{1. What this unit proves}\label{what-this-unit-proves-31}

\textbf{Conventions.} The notation follows statement.md:

\begin{itemize}
\tightlist
\item
  \((m,G,\mathfrak g)\) is a gauge group: \(m\ge1\), \(G\) is a closed
  subgroup of \(U(m)\), and
  \(\mathfrak g=\{X:\exp(tX)\in G\ \text{for all }t\in\mathbb R\}\).
  Simplicity includes nonabelianness.
\item
  \(\mathcal W=(L_p,D_N,\beta_{N,p},h,\mathcal T,I_{N,p})\) is a Wilson
  family, with \(a_N=1/D_N\), smeared moments \(C_{N,p}(v;f)\) and
  separated configurations \(\mathrm{Conf}_n\).
\item
  A hierarchy has a Wightman theory \(\mathfrak W\), finite-dimensional
  field spaces \(F_d\), label maps \(\iota_d\), token projection
  \(\pi\), phases \((X,\nu,\xi_*)\), correlations \(S^s\), time
  continuations \(\mathsf F_n\), OPE data
  \((\mathrm{Ch}_j,O_{j,\kappa},c_{j,\kappa})\), and \(\Theta\in F_4\).
  Write \(\mathbb F=\bigsqcup_dF_d\). The function \(G_\Theta(r)\) is
  the selected two-point correlation of \(\Theta\) at time-axis points
  \(-r/2\) and \(r/2\).
\item
  (Reg1)--(Reg4), (R1)--(R15), (O1)--(O2), (T1)--(T4), (P1)--(P4) and
  (D1)--(D16) are as in statement.md.
\item
  \(b_0=11C_{\mathfrak g}/(48\pi^2)\), where \(C_{\mathfrak g}\) is the
  adjoint Casimir in the construction normalization. All Lie-algebra
  dimensions are real.
\end{itemize}

Internal objects:

\begin{itemize}
\tightlist
\item
  An \emph{extraction order} is a pair
  \(\sigma=(\sigma^{\rm m},\sigma^{\rm v})\) of strictly increasing maps
  \(\mathbb N\to\mathbb N\): a mesh sequence and a volume sequence.
\item
  Its \emph{refinement} by \(\tau\) is
  \(\sigma\circ\tau:=(\sigma^{\rm m}\circ\tau^{\rm m},\sigma^{\rm v}\circ\tau^{\rm v})\).
  The full sequence is written \(\mathrm{id}\).
\item
  An \emph{internal family} \(\mathcal C\) is the record of a
  renormalization-group construction: trajectory, couplings
  \(\beta_{N,p}(m)\) at levels \(m\le N\), labels and local lattice
  sections. (D3) is its construction predicate, and
  \(\mathcal W(\mathcal C)\) is its Wilson family, with
  \(\beta_{N,p}:=\beta_{N,p}(N)\).
\item
  A \emph{table} \(S\) is a family of continuum correlation functions
  \(S(v;x)\) indexed by token lists.
\end{itemize}

\emph{Ordered limit.} \(\mathrm{Ord}_\sigma(\mathcal C,S)\) means the
following. For every \(n\), \(v\in\mathcal T^n\) and
\(t\in C^\infty_c(\mathrm{Conf}_n)\):

\begin{itemize}
\tightlist
\item
  \(t\) is \(\sigma^{\rm v}(k)\)-admissible for large \(k\);
\item
  there are numbers \(R_p\) with \(C_{\sigma^{\rm m}(j),p}(v;t)\to R_p\)
  for every admissible \(p\);
\item
  \(R_{\sigma^{\rm v}(k)}\to\int tS(v;\cdot)\).
\end{itemize}

\emph{Final object.} A final object \(\mathfrak F\) on
\((\mathcal C,\sigma,S)\) consists of five parts:

\begin{enumerate}
\def\labelenumi{\arabic{enumi}.}
\tightlist
\item
  Chosen renormalization data and step witnesses whose trajectory is
  exactly that of \(\mathcal C\) (D6).
\item
  A \emph{seed}: (D3), \(\mathrm{Ord}_\sigma(\mathcal C,S)\), point
  limits, the Euclidean hierarchy axioms for \(S\), cell hypotheses and
  corrected-field widths.
\item
  A field algebra: derivatives, rotations and a stress tensor.
\item
  A phase selection with its certificate.
\item
  Analytic continuations of the mixed and selected correlations, and a
  Wightman theory realizing the selected one.
\end{enumerate}

\textbf{Proposition VIII.1.1 (main theorem, public form).} Let
\((m,G,\mathfrak g)\) be a gauge group with \(\mathfrak g\) simple. Then
there is a Wilson family \(\mathcal W\) for \(G\) satisfying
(Reg1)--(Reg4), with the following property. For every strictly
increasing \(p:\mathbb N\to\mathbb N\) there are a strictly increasing
\(q\) and a hierarchy over \(\mathcal W\) satisfying (P1)--(P4) along
\(p\circ q\):

\begin{itemize}
\tightlist
\item
  \textbf{(P1)} For every \(n\), every \(v\in\mathcal T^n\) and every
  \(f\in C^\infty_c(\mathrm{Conf}_n)\), \(f\) is \(p(q(k))\)-admissible
  for all large \(k\). There are numbers \(R_{p'}\) such that
  \(C_{N,p'}(v;f)\to R_{p'}\) along all \(N\) for every \(p'\) for which
  \(f\) is admissible, and
  \(R_{p(q(k))}\to\int f\,S^{\rm mix}(\pi(v);\cdot)\).
\item
  \textbf{(P2)} The selected phase carries a Wightman theory satisfying
  (W0)--(W3), with faithful fields. The mixed correlations are the
  \(\nu\)-average of continuous phase correlations, with \(\nu\) a
  probability measure of full support. The continuation is holomorphic
  in the complex times at fixed real spatial positions, and its boundary
  values on product Schwartz tests are the selected Wightman functions.
  All clauses (R1)--(R15) hold.
\item
  \textbf{(P3)} The OPE bounds (O1)--(O2) hold with one system of
  coefficients, uniformly in the state \(s\in X\sqcup\{\mathrm{mix}\}\).
\item
  \textbf{(P4)} \(\Theta\) has zero vacuum expectation,
  \(A_{\mathfrak g}:=3\dim_{\mathbb R}\mathfrak g/\pi^4>0\), and
  \(r^8(\log 1/r)^2G_\Theta(r)\to A_{\mathfrak g}\) as \(r\downarrow0\).
  For every finite Wick profile and every \(\delta>0\), the profile
  differs from \(G_\Theta\) at some \(r\in(0,\delta)\).
\end{itemize}

\emph{Status:} theorem. \emph{Lean:}
\href{https://github.com/yang-mills-group/OpenAI-YM-20261028/blob/3456e5e75cf8d4cd558955a5c7638fba2d68f464/Solution.lean\#L4-L10}{\texttt{YangMills.jointExistence}},
statement
\href{https://github.com/yang-mills-group/OpenAI-YM-20261028/blob/3456e5e75cf8d4cd558955a5c7638fba2d68f464/ChallengeSpec.lean\#L186-L188}{\texttt{JointExistence}},
with
\href{https://github.com/yang-mills-group/OpenAI-YM-20261028/blob/3456e5e75cf8d4cd558955a5c7638fba2d68f464/ChallengeSpec.lean\#L12-L22}{\texttt{WilsonRegularity}},
\href{https://github.com/yang-mills-group/OpenAI-YM-20261028/blob/3456e5e75cf8d4cd558955a5c7638fba2d68f464/ChallengeSpec.lean\#L24-L48}{\texttt{Hierarchy}},
\href{https://github.com/yang-mills-group/OpenAI-YM-20261028/blob/3456e5e75cf8d4cd558955a5c7638fba2d68f464/ChallengeSpec.lean\#L177-L181}{\texttt{Properties}},
\href{https://github.com/yang-mills-group/OpenAI-YM-20261028/blob/3456e5e75cf8d4cd558955a5c7638fba2d68f464/PublicWilson.lean\#L85-L90}{\texttt{Converges}},
\href{https://github.com/yang-mills-group/OpenAI-YM-20261028/blob/3456e5e75cf8d4cd558955a5c7638fba2d68f464/ChallengeSpec.lean\#L117-L143}{\texttt{Reconstructed}},
\href{https://github.com/yang-mills-group/OpenAI-YM-20261028/blob/3456e5e75cf8d4cd558955a5c7638fba2d68f464/ChallengeSpec.lean\#L70-L80}{\texttt{OperatorProductExpansion}}
and
\href{https://github.com/yang-mills-group/OpenAI-YM-20261028/blob/3456e5e75cf8d4cd558955a5c7638fba2d68f464/ChallengeSpec.lean\#L168-L175}{\texttt{TraceTest}}.

\textbf{Proposition VIII.1.2 (detailed construction theorem).} Fix a
gauge group \((m,G,\mathfrak g)\) with \(\mathfrak g\) simple and a real
\(b>0\). Then there is an internal family \(\mathcal C\) such that:

\begin{itemize}
\tightlist
\item
  its block-spin kernel has penalty \(b\) (D16);
\item
  (D3) holds;
\item
  \(\mathcal W(\mathcal C)\) satisfies (Reg1)--(Reg4).
\end{itemize}

Moreover, for every extraction order \(\sigma\) there are:

\begin{itemize}
\tightlist
\item
  a refinement \(\tau\) and a table \(S\);
\item
  a final object \(\mathfrak F\) on \((\mathcal C,\sigma\circ\tau,S)\);
\item
  an every-phase reconstruction \(\mathfrak L\), an OPE reader \(Q\),
  contact devices \(\mathsf M\) and scaling data
  \(\Sigma=(r_k,\xi_k,U_k)\), with \(r_k\downarrow0\) and \(\xi_k>0\);
\end{itemize}

such that the following hold:

\begin{enumerate}
\def\labelenumi{\arabic{enumi}.}
\tightlist
\item
  The public hierarchy \(T(\mathfrak F,Q)\) over
  \(\mathcal W(\mathcal C)\) has (P1)--(P4) along
  \((\sigma\circ\tau)^{\rm v}\), and every correlation \(S^s(a;\cdot)\)
  extends to a distribution on \((\mathbb R^4)^n\) (D1, D2).
\item
  \(\mathrm{Ord}_{(\mathrm{id},(\sigma\circ\tau)^{\rm v})}(\mathcal C,S)\)
  holds: the ordered limit with the full mesh sequence (D4).
\item
  Positional ultraviolet convergence at every fixed volume holds along
  the full mesh sequence (D5). See VIII.1.6.
\item
  Chosen renormalization data and step witnesses generate exactly the
  trajectory of \(\mathcal C\) (D6).
\item
  The bundle (Core) holds: OPE services, reconstruction in every phase,
  the trace tests, the OPE foundation, and integrated OPE remainder
  bounds for integrable complex densities against positive measures,
  with bounded \(L^1\) mass.
\item
  The internal Wightman axioms hold for the selected phase and for every
  phase \(\xi\in X\), and the public (W0)--(W3) hold for every phase.
  The \(\nu\)-average of the phase Wightman distributions is the
  boundary distribution of the mixed continuation (D7).
\item
  The separated OPE and the nested OPE over finite trees hold (D8, D10).
\item
  The seven headline trace clauses hold, together with the unique
  high-energy spectral law contained in (Core) (D9; VIII.1.n and
  VIII.1.o).
\item
  The contact devices use the reader's prescription and give a full
  contact resolution, real bases of the OPE strata and contact
  thresholds (D11)--(D13).
\item
  The leading scaling comparison with free Maxwell fields along
  \(\Sigma\) and the normalization at degree four hold (D14, D15).
\end{enumerate}

The family \(\mathcal C\) depends only on \((G,b)\). Items 3 and 4 are
properties of \(\mathcal C\) alone, although they are stated inside the
quantifier over \(\sigma\). \emph{Status:} theorem. \emph{Lean:}
\href{https://github.com/yang-mills-group/OpenAI-YM-20261028/blob/3456e5e75cf8d4cd558955a5c7638fba2d68f464/Construction.lean\#L15-L40}{\texttt{YangMills.constructed}}.
Internal form:
\href{https://github.com/yang-mills-group/OpenAI-YM-20261028/blob/3456e5e75cf8d4cd558955a5c7638fba2d68f464/YM/Nonperturbative/Headlines.lean\#L19-L36}{\texttt{Nonperturbative.existence}},
statement
\href{https://github.com/yang-mills-group/OpenAI-YM-20261028/blob/3456e5e75cf8d4cd558955a5c7638fba2d68f464/YM/Nonperturbative/HeadlineSpec.lean\#L143-L155}{\texttt{Headline.WilsonExistence}}.

\textbf{Proposition VIII.1.3 (continuum theorem for the funded run).}
Fix a gauge group with \(\mathfrak g\) simple and \(b>0\). Let
\(\mathcal C_{G,b}\) be the family of the run \(\mathsf R_{G,b}\) chosen
in VIII.1.a. Its penalty is \(b\). For every extraction order \(\sigma\)
there are \(\tau,S,\mathfrak F,\mathfrak L,Q,\mathsf M\), with
\(\mathfrak F\) on \((\mathcal C_{G,b},\sigma\circ\tau,S)\), satisfying:

\begin{itemize}
\tightlist
\item
  items 2, 3, 5 and 9 of VIII.1.2;
\item
  the uniform nested OPE over admissible finite trees;
\item
  a scaling sequence \(\Sigma\) with (D14) and (D15).
\end{itemize}

\emph{Status:} theorem. \emph{Lean:}
\href{https://github.com/yang-mills-group/OpenAI-YM-20261028/blob/3456e5e75cf8d4cd558955a5c7638fba2d68f464/YM/Nonperturbative/Continuum.lean\#L28-L48}{\texttt{continuum\_with\_position\_tests}}.
The form without \(\Sigma\) is
\href{https://github.com/yang-mills-group/OpenAI-YM-20261028/blob/3456e5e75cf8d4cd558955a5c7638fba2d68f464/YM/Nonperturbative/Continuum.lean\#L52-L66}{\texttt{continuum}}.

\textbf{Proposition VIII.1.4 (a funded run supplies everything).} Let
\(\mathsf R\) be a run with campaign \(T\) and pillars \(W\), let
\(\Delta\) be a choice of local step estimates, and let \(\mathfrak g\)
be simple. Assume \(\mathsf R\) is \emph{funded} (Fd): see §2, item 1.
Then the assembly package (Goal) holds:

\begin{enumerate}
\def\labelenumi{\arabic{enumi}.}
\tightlist
\item
  (Fd).
\item
  The ultraviolet package (UV\(_\infty\)) of VIII.1.5.
\item
  The logarithmic OPE package.
\item
  The dense-coverage property of the sharp Gram family.
\item
  For every \(\sigma\), every table \(S\) and every final object
  \(\mathfrak F\) on \((\mathcal C_{\mathsf R},\sigma,S)\): the trace
  result of VIII.1.h.
\item
  For every \(\sigma\) there are \(\tau\), numbers
  \((B_m)_{m\in\mathbb N}\) and data \(S,\mathfrak F\) (on
  \(\sigma\circ\tau\)), \(\Sigma,\mathfrak L,Q\) such that:

  \begin{itemize}
  \tightlist
  \item
    for every level \(m\) the couplings \(\beta_{N,p}(m)\) have the
    ordered limit \(B_m\) along \(\sigma\circ\tau\);
  \item
    \(\mathrm{Ord}_{(\mathrm{id},(\sigma\circ\tau)^{\rm v})}(\mathcal C_{\mathsf R},S)\)
    and (D5) hold;
  \item
    (D14) holds;
  \item
    the reconstruction in every phase holds;
  \item
    the trace result holds;
  \item
    the OPE facilities, the OPE foundation for \(Q\) and the uniform
    nested OPE hold;
  \item
    the all-shapes OPE holds for every coherent prescription;
  \item
    for every ultrafilter on \(\mathbb N^2\) finer than the extraction
    filter of \(\sigma\circ\tau\) there is a compatible allocation of
    the sharp Gram data.
  \end{itemize}
\end{enumerate}

\emph{Status:} theorem. \emph{Lean:}
\href{https://github.com/yang-mills-group/OpenAI-YM-20261028/blob/3456e5e75cf8d4cd558955a5c7638fba2d68f464/YM/R2I/Campaign/Accepted.lean\#L150-L167}{\texttt{R2IGoals.fromFunded}},
statement
\href{https://github.com/yang-mills-group/OpenAI-YM-20261028/blob/3456e5e75cf8d4cd558955a5c7638fba2d68f464/YM/R2I/Campaign/Accepted.lean\#L130-L148}{\texttt{Broad}}.

\textbf{Proposition VIII.1.5 (ultraviolet endpoint).} Let
\(\mathfrak g\) be simple, let \(\mathsf R\) be a run with campaign
\(T\) and pillars \(W\), and assume the stationary-point and
amplification condition, the delay condition, and the weight-and-tableau
data \((\vartheta,v)\) of (Fd). Then the ultraviolet package
(UV\(_\infty\)) holds:

\begin{enumerate}
\def\labelenumi{\arabic{enumi}.}
\item
  (D3) holds for \(\mathcal C_{\mathsf R}\), and the weight data exist.
\item
  \emph{Sampler.} Every \(\sigma\) has a refinement \(\tau\) carrying a
  table \(S\) and a final object on
  \((\mathcal C_{\mathsf R},\sigma\circ\tau,S)\), and
  \(\mathrm{Ord}_{(\mathrm{id},(\sigma\circ\tau)^{\rm v})}\) holds.
\item
  The trajectory at every \((p,N)\) competes in, and is issued by, the
  fixed tuning recursion.
\item
  \emph{Tuning and convergence.} The parameter \(s\in[-1,1]\) specifies
  the target coupling \(\bar\beta+ws\) at level \(0\), where \(w>0\) is
  the fixed tuning width. For every \((p,N,s)\) there is a successful
  trajectory, and its bare coupling at level \(N\) is unique within the
  prescribed recursion. These bare couplings are values of a function
  analytic on the open unit disc in \(s\). At each fixed \((m,p)\), the
  couplings and activities of successful trajectories converge as
  \(N\to\infty\) whenever their level-zero targets converge. For complex
  targets in the disc, convergence of the analytic couplings and
  activities, with every parameter derivative, is uniform on compact
  subsets; activities are measured in the fixed weighted norm at
  \((m,p)\).
\item
  For every fixed \((m,p)\) there is \(c_{m,p}>0\) such that

  \[e^{-\mathrm{scalar}_{N,p,m}}Z_{N,p}\ge c_{m,p}\qquad(N\ge m),\]
  where \(Z_{N,p}\) is the bare Wilson partition function relative to
  probability Haar measure and \(\mathrm{scalar}_{N,p,m}\) is the
  field-independent scalar in the level-\(m\) action. Its scalar density
  \(v_{N,p,m}\) tends to \(-\infty\) as \(N\to\infty\).
\item
  The volume-uniform bounds hold, as stated in VIII.1.c, step 1.
\item
  (D5) holds along the full mesh sequence.
\item
  For every fixed \(p\) there is a family of successful trajectories
  with target \(\bar\beta+\lambda ws\), \(s\in[-1,1]\), for some
  \(\lambda>0\), whose central member is the chosen family. They obey
  the same step relations and insertion normalization. For each token
  list, each compact \(K\subset\mathrm{NoHit}_p\) and each finite
  spatial derivative order \(h\), there are an open neighbourhood
  \(\Omega\) of \(K\) in \(\mathrm{NoHit}_p\) and a smaller positive
  parameter width, independent of \(N\). In the resulting unit-disc
  coordinate, the limiting correlation is holomorphic in the parameter
  and \(C^h\) in position on \(\Omega\). All fixed coordinate
  differences of order at most \(h\), their interpolants and their
  admissible smearings converge to the corresponding derivatives,
  uniformly on compact parameter and position sets, with every parameter
  derivative.
\end{enumerate}

\emph{Status:} theorem. It is assembled here; its ingredients are in
III.12--III.15. \emph{Lean:}
\href{https://github.com/yang-mills-group/OpenAI-YM-20261028/blob/3456e5e75cf8d4cd558955a5c7638fba2d68f464/YM/RG17/Joint.lean\#L102-L111}{\texttt{FullVolume.Final.jewelOf17}}
and
\href{https://github.com/yang-mills-group/OpenAI-YM-20261028/blob/3456e5e75cf8d4cd558955a5c7638fba2d68f464/YM/FullVolume/Finale/Program.lean\#L133-L149}{\texttt{Jewel}}.

\textbf{Proposition VIII.1.6 (the ultraviolet limit along the full
spacing sequence).} For \(p\in\mathbb N\), let \(\mathrm{NoHit}_p\) be
the set of configurations no two of whose points are congruent modulo
\(L_p\mathbb Z^4\).

\begin{enumerate}
\def\labelenumi{(\alph{enumi})}
\tightlist
\item
  For the family of a run satisfying VIII.1.5's hypotheses, the
  following holds for every \(p\), \(n\) and \(v\in\mathcal T^n\). There
  is \(f_{p,v}\in C^\infty(\mathrm{NoHit}_p)\) such that:
\end{enumerate}

\begin{itemize}
\tightlist
\item
  \(C_{N,p}(v;t)\to\int t\,f_{p,v}\) as \(N\to\infty\) along all \(N\),
  for every \(p\)-admissible \(t\);
\item
  for every fixed finite list of normalized forward, backward or
  centered coordinate differences, with fixed integer stencil widths,
  the interpolated lattice differences converge locally uniformly on
  \(\mathrm{NoHit}_p\) to the corresponding derivatives of \(f_{p,v}\).
  The same limit holds whenever lattice configurations approach a point
  of \(\mathrm{NoHit}_p\), along any cutoff indices tending to infinity.
  The smeared differences converge to the integral of the derivative
  against the test, and that integrand is integrable.
\end{itemize}

The cutoff thresholds may depend on \(p,n,v\), the compact set, the
tests and the difference stencils.

\begin{enumerate}
\def\labelenumi{(\alph{enumi})}
\setcounter{enumi}{1}
\tightlist
\item
  Suppose a final object exists on \((\mathcal C,\sigma,S)\) and (a)
  holds for \(\mathcal C\). Then
  \(\mathrm{Ord}_{(\mathrm{id},\sigma^{\rm v})}(\mathcal C,S)\) holds.
\end{enumerate}

\emph{Status:} theorems. \emph{Lean:} (a)
\href{https://github.com/yang-mills-group/OpenAI-YM-20261028/blob/3456e5e75cf8d4cd558955a5c7638fba2d68f464/YM/FullVolume/Fields/Filled.lean\#L45}{\texttt{Campaign.realUV}},
with the definition
\href{https://github.com/yang-mills-group/OpenAI-YM-20261028/blob/3456e5e75cf8d4cd558955a5c7638fba2d68f464/YM/FullVolume/Profiles.lean\#L25-L36}{\texttt{PositionsUV}}.
(b)
\href{https://github.com/yang-mills-group/OpenAI-YM-20261028/blob/3456e5e75cf8d4cd558955a5c7638fba2d68f464/YM/FullVolume/Finale/Program.lean\#L117-L126}{\texttt{FullVolume.Final.from\_spec}}.

\textbf{Proposition VIII.1.7 (public bridges).} Let \(\mathfrak F\) be a
final object on \((\mathcal C,\sigma,S)\) and \(Q\) a reader, and let
\(T(\mathfrak F,Q)\) be the public hierarchy built from them (VIII.1.p).

\begin{enumerate}
\def\labelenumi{\arabic{enumi}.}
\tightlist
\item
  The public Wilson energy, Wilson measure, lattice moments and smeared
  moments of \(\mathcal W(\mathcal C)\) equal the internal ones.
  \(\mathrm{Ord}_{(\mathrm{id},s)}(\mathcal C,S)\) implies (P1) along
  \(s\).
\item
  If (D3) holds and \(\mathfrak g\) is simple, then
  \(\mathcal W(\mathcal C)\) satisfies (Reg1)--(Reg4), with
  \(b_*=2b_0\).
\item
  The reconstruction in every phase implies (R1)--(R15).
\item
  The internal separated OPE implies (O1)--(O2).
\item
  If \(\mathfrak g\) is simple, the seven internal headline trace
  clauses imply (T1)--(T4), with \(3\dim\mathfrak g/\pi^4\) identified
  with the internal amplitude.
\item
  The full contact resolution implies (D2).
\end{enumerate}

\emph{Status:} theorems. \emph{Lean:}
\href{https://github.com/yang-mills-group/OpenAI-YM-20261028/blob/3456e5e75cf8d4cd558955a5c7638fba2d68f464/WilsonBridge.lean\#L20-L30}{\texttt{wilsonFamily}}
and
\href{https://github.com/yang-mills-group/OpenAI-YM-20261028/blob/3456e5e75cf8d4cd558955a5c7638fba2d68f464/WilsonBridge.lean\#L62-L67}{\texttt{fromWilson}}
(1);
\href{https://github.com/yang-mills-group/OpenAI-YM-20261028/blob/3456e5e75cf8d4cd558955a5c7638fba2d68f464/HierarchyBridge.lean\#L16-L25}{\texttt{regular}}
(2);
\href{https://github.com/yang-mills-group/OpenAI-YM-20261028/blob/3456e5e75cf8d4cd558955a5c7638fba2d68f464/HierarchyBridge.lean\#L79-L133}{\texttt{publicReconstruction}}
(3);
\href{https://github.com/yang-mills-group/OpenAI-YM-20261028/blob/3456e5e75cf8d4cd558955a5c7638fba2d68f464/HierarchyBridge.lean\#L68-L77}{\texttt{publicOPE}}
(4);
\href{https://github.com/yang-mills-group/OpenAI-YM-20261028/blob/3456e5e75cf8d4cd558955a5c7638fba2d68f464/TraceBridge.lean\#L35-L44}{\texttt{publicTrace}}
(5);
\href{https://github.com/yang-mills-group/OpenAI-YM-20261028/blob/3456e5e75cf8d4cd558955a5c7638fba2d68f464/ContactBridge.lean\#L14-L49}{\texttt{contactDistribution}}
(6).

\textbf{Proposition VIII.1.8 (exclusion of Wick descriptions in momentum
space).} Let \(\mathfrak F\) be a final object. Assume
\(s^8(\log 1/s)^2G_\Theta(s)\to A>0\) as \(s\downarrow0\). Let \(W\) be
any finite list of complex coefficients \(c_\alpha\) and products of
mass-shell modes, in the class defining Wick profiles in statement.md.
There are Schwartz functions \(f,g\) on \(\mathbb R^4\) for which

\[v_{f,g}(p):=\overline{\mathcal F_+f(p)}\,\mathcal F_+g(p)\] has
compact support and

\[\mathcal W^{\xi_*}_2(\Theta^*,\Theta)(\bar f\otimes g)\ne\sum_\alpha c_\alpha\,\mathrm{read}_\alpha(v_{f,g}).\]
Here \(\mathcal F_+\) is the unnormalized positive-sign Fourier
transform used in the spectral representation. The four-dimensional test
\(v_{f,g}\) is the transform of the joint test restricted to relative
momenta \((-p,p)\). \emph{Status:} theorem. It enters the trace result
of VII.3. The public (T4) follows directly from the Euclidean exclusion
(VIII.1.r). \emph{Lean:}
\href{https://github.com/yang-mills-group/OpenAI-YM-20261028/blob/3456e5e75cf8d4cd558955a5c7638fba2d68f464/YM/TraceExtension/Wick/Refute.lean\#L86-L105}{\texttt{TraceExtension.Wick.polynomial\_excludes}}.

\textbf{Proposition VIII.1.9 (\(SU(2)\)).}
\(\mathfrak{su}(2)\subset M_2(\mathbb C)\) is simple. Hence VIII.1.1
holds for \((2,SU(2),\mathfrak{su}(2))\). \emph{Status:} theorem.
\emph{Lean:}
\href{https://github.com/yang-mills-group/OpenAI-YM-20261028/blob/3456e5e75cf8d4cd558955a5c7638fba2d68f464/ConcreteSU2.lean\#L12}{\texttt{su2Simple}}
(from
\href{https://github.com/yang-mills-group/OpenAI-YM-20261028/blob/3456e5e75cf8d4cd558955a5c7638fba2d68f464/YM/Examples/SU2.lean\#L79-L106}{\texttt{Examples.SU2.simple}})
and
\href{https://github.com/yang-mills-group/OpenAI-YM-20261028/blob/3456e5e75cf8d4cd558955a5c7638fba2d68f464/ConcreteSU2.lean\#L15-L18}{\texttt{su2Theory}}.

\textbf{Proposition VIII.1.10 (definitions filed here and used
elsewhere).} Several definitions used by III.10--III.15, V.1--V.3,
VI.1--VI.2 and VII.1--VII.4 are filed in this unit:

\begin{itemize}
\tightlist
\item
  the structure of a final object and its parts:
  \href{https://github.com/yang-mills-group/OpenAI-YM-20261028/blob/3456e5e75cf8d4cd558955a5c7638fba2d68f464/YM/Spec/Closed.lean\#L64-L73}{\texttt{Spec.Final}},
  \href{https://github.com/yang-mills-group/OpenAI-YM-20261028/blob/3456e5e75cf8d4cd558955a5c7638fba2d68f464/YM/Spec/Closed.lean\#L20-L26}{\texttt{Seed}},
  \href{https://github.com/yang-mills-group/OpenAI-YM-20261028/blob/3456e5e75cf8d4cd558955a5c7638fba2d68f464/YM/Spec/Closed.lean\#L34-L43}{\texttt{Locals}},
  \href{https://github.com/yang-mills-group/OpenAI-YM-20261028/blob/3456e5e75cf8d4cd558955a5c7638fba2d68f464/YM/Spec/Closed.lean\#L45-L47}{\texttt{Materializes}},
  \href{https://github.com/yang-mills-group/OpenAI-YM-20261028/blob/3456e5e75cf8d4cd558955a5c7638fba2d68f464/YM/Spec/Closed.lean\#L50-L54}{\texttt{Relativistic}}
  and
  \href{https://github.com/yang-mills-group/OpenAI-YM-20261028/blob/3456e5e75cf8d4cd558955a5c7638fba2d68f464/YM/Spec/Closed.lean\#L58-L61}{\texttt{Calibrated}};
\item
  the input bundles
  \href{https://github.com/yang-mills-group/OpenAI-YM-20261028/blob/3456e5e75cf8d4cd558955a5c7638fba2d68f464/YM/Org/Input.lean\#L47-L54}{\texttt{Org.Alloy}},
  \href{https://github.com/yang-mills-group/OpenAI-YM-20261028/blob/3456e5e75cf8d4cd558955a5c7638fba2d68f464/YM/Org/Input.lean\#L64-L66}{\texttt{Assembly}},
  \href{https://github.com/yang-mills-group/OpenAI-YM-20261028/blob/3456e5e75cf8d4cd558955a5c7638fba2d68f464/YM/Org/Input.lean\#L159-L170}{\texttt{MaxwellData}},
  \href{https://github.com/yang-mills-group/OpenAI-YM-20261028/blob/3456e5e75cf8d4cd558955a5c7638fba2d68f464/YM/Org/Input.lean\#L172-L175}{\texttt{IR}}
  and
  \href{https://github.com/yang-mills-group/OpenAI-YM-20261028/blob/3456e5e75cf8d4cd558955a5c7638fba2d68f464/YM/Org/Input.lean\#L177-L186}{\texttt{UV}};
\item
  the volume-uniform bounds of YM/FullVolume, for example
  \href{https://github.com/yang-mills-group/OpenAI-YM-20261028/blob/3456e5e75cf8d4cd558955a5c7638fba2d68f464/YM/FullVolume/Discrete.lean\#L50-L52}{\texttt{FixedBounds}}
  and
  \href{https://github.com/yang-mills-group/OpenAI-YM-20261028/blob/3456e5e75cf8d4cd558955a5c7638fba2d68f464/YM/FullVolume/Literal.lean\#L62}{\texttt{genuine}}.
\end{itemize}

\emph{Status:} definitions, and theorems described in the units that use
them.

\subsubsection{2. What it uses}\label{what-it-uses-31}

\textbf{From other units.}

\begin{enumerate}
\def\labelenumi{\arabic{enumi}.}
\tightlist
\item
  \textbf{III.14 and III.2 (selection of the run).} For simple
  \(\mathfrak g\) and \(b>0\) there are a run \(\mathsf R\), a campaign
  \(T\) and pillars \(W\) such that the penalty is \(b\) and
  \(\mathsf R\) is funded for the fixed local-estimate choice
  \(\Delta\). A run records levels, schedule of constants, step budgets,
  the trajectory and its tuning, the reference levels and the transfer
  data of the labels.

  \begin{itemize}
  \item
    (Fd) comprises thirteen conditions: local step budgets;
    phase-comparison and dense-coverage data; divergence of a schedule
    \(Q\); stationary-point and amplification data; a delay condition;
    frame compatibility; polynomial schedule compatibility; support
    exponent \(100\) at every level; the logarithmic OPE budget; the
    halo condition; a pacing condition; thermal OPE estimates; and
    weight-and-tableau data \((\vartheta,v)\). The detailed numerical
    conditions belong to III.15.
  \item
    The run is produced by a selection theorem for runs with a fixed
    treasury period, a preselected component recipe and the default
    dynamics (III.2).
  \item
    \emph{Lean:}
    \href{https://github.com/yang-mills-group/OpenAI-YM-20261028/blob/3456e5e75cf8d4cd558955a5c7638fba2d68f464/YM/Independent/ConstructionRun.lean\#L14-L28}{\texttt{Independent.constructionRun}},
    \href{https://github.com/yang-mills-group/OpenAI-YM-20261028/blob/3456e5e75cf8d4cd558955a5c7638fba2d68f464/YM/R2I/Campaign/Accepted.lean\#L60-L73}{\texttt{Funded}},
    and the selection theorem in
    \href{https://github.com/yang-mills-group/OpenAI-YM-20261028/blob/3456e5e75cf8d4cd558955a5c7638fba2d68f464/YM/Integration/NativePreselectedPostnumericPeriodRun.lean}{\texttt{NativePreselectedPostnumericPeriodRun.lean}}.
  \end{itemize}
\item
  \textbf{III.12--III.15 (ultraviolet bounds and convergence).} Used for
  VIII.1.5 and (D5):

  \begin{itemize}
  \tightlist
  \item
    the convergence and alignment of the activities along the tuning
    recursion;
  \item
    uniqueness and analyticity of the trial couplings;
  \item
    the volume-uniform bounds on correlation kernels and their lattice
    differences at separated points
    (\href{https://github.com/yang-mills-group/OpenAI-YM-20261028/blob/3456e5e75cf8d4cd558955a5c7638fba2d68f464/YM/FullVolume/Literal.lean\#L62}{\texttt{FullVolume.genuine}});
  \item
    the pointwise Cauchy property of lattice correlations along the full
    spacing sequence at fixed volume (as used in
    \href{https://github.com/yang-mills-group/OpenAI-YM-20261028/blob/3456e5e75cf8d4cd558955a5c7638fba2d68f464/YM/FullVolume/Fields/Filled.lean\#L30-L38}{\texttt{Campaign.sees}});
  \item
    the property that every extraction order has a refinement carrying a
    final object
    (\href{https://github.com/yang-mills-group/OpenAI-YM-20261028/blob/3456e5e75cf8d4cd558955a5c7638fba2d68f464/YM/FullVolume/Cross.lean\#L73-L75}{\texttt{Run.again}},
    through
    \href{https://github.com/yang-mills-group/OpenAI-YM-20261028/blob/3456e5e75cf8d4cd558955a5c7638fba2d68f464/YM/FullVolume/Forge.lean\#L60-L72}{\texttt{ahead}}).
  \end{itemize}
\item
  \textbf{I.1, II.1--II.2 (setting and construction record).} The
  internal family, (D3) including the running law: for every
  \(\varepsilon>0\),
  \(|\beta_{N,p}(m)-\hat\beta_m|\le\varepsilon\log(1/a_m)\) for all
  sufficiently large \(m\) and all \(N\ge m\) and all \(p\), where
  \(\hat\beta_m=\bar\beta+2b_0\log(1/a_m)\); the band condition
  \(\beta_{N,p}(m)\in[\hat\beta_m/2,2\hat\beta_m]\); tokens, local
  sections and their stencils; tables.
\item
  \textbf{IV.1 (infinite volume).} The seed: a subsequential table \(S\)
  with \(\mathrm{Ord}_{\sigma\circ\tau}\), point limits, the Euclidean
  hierarchy axioms, cell hypotheses and corrected-field widths. It is
  produced inside the run's extraction property and used in VIII.1.g.
\item
  \textbf{V.1--V.3 (phases and reconstruction).}

  \begin{itemize}
  \tightlist
  \item
    The phase selection with its certificate (work items B01--B18).
  \item
    The Osterwalder--Schrader--Wightman reconstruction of the selected
    phase
    (\href{https://github.com/yang-mills-group/OpenAI-YM-20261028/blob/3456e5e75cf8d4cd558955a5c7638fba2d68f464/YM/Next/BOUNDARY/Stage.lean\#L247-L255}{\texttt{after\_materializes}}).
  \item
    The reconstruction in every phase,
    \href{https://github.com/yang-mills-group/OpenAI-YM-20261028/blob/3456e5e75cf8d4cd558955a5c7638fba2d68f464/YM/EveryPhase/Extension.lean\#L75-L80}{\texttt{EveryPhase.all\_characters}}
    with
    \href{https://github.com/yang-mills-group/OpenAI-YM-20261028/blob/3456e5e75cf8d4cd558955a5c7638fba2d68f464/YM/EveryPhase/Extension.lean\#L51-L71}{\texttt{Completed}}
    (V.3).
  \end{itemize}
\item
  \textbf{VI.1--VI.2 (symmetry).} The field algebra (derivatives,
  rotations, conserved stress tensor) and its local laws: the local-law
  package built by
  \href{https://github.com/yang-mills-group/OpenAI-YM-20261028/blob/3456e5e75cf8d4cd558955a5c7638fba2d68f464/YM/Bid/Town.lean\#L159-L168}{\texttt{Bid.station}}
  from work items A01--A17, and
  \href{https://github.com/yang-mills-group/OpenAI-YM-20261028/blob/3456e5e75cf8d4cd558955a5c7638fba2d68f464/YM/Org/Twist.lean\#L19-L24}{\texttt{Org.Tab}}.
\item
  \textbf{VII.1 (free Maxwell comparison).} The free-Maxwell data
  \href{https://github.com/yang-mills-group/OpenAI-YM-20261028/blob/3456e5e75cf8d4cd558955a5c7638fba2d68f464/YM/Next/MAXWELL/Lit.lean\#L45-L48}{\texttt{maxwell\_supplies}}
  and
  \href{https://github.com/yang-mills-group/OpenAI-YM-20261028/blob/3456e5e75cf8d4cd558955a5c7638fba2d68f464/YM/UVNext/FREE_STRESS_CERTIFICATE/Choose.lean\#L137-L150}{\texttt{FreeStress.fire}};
  the leading short-distance comparison; the degree-four normalization
  (\href{https://github.com/yang-mills-group/OpenAI-YM-20261028/blob/3456e5e75cf8d4cd558955a5c7638fba2d68f464/YM/Pert/Position/Crest.lean\#L50-L52}{\texttt{credible}},
  \href{https://github.com/yang-mills-group/OpenAI-YM-20261028/blob/3456e5e75cf8d4cd558955a5c7638fba2d68f464/YM/Pert/Position/Four.lean\#L13-L15}{\texttt{square\_four}}).
\item
  \textbf{VII.2--VII.3 (trace).}

  \begin{itemize}
  \tightlist
  \item
    The gallery property
    \href{https://github.com/yang-mills-group/OpenAI-YM-20261028/blob/3456e5e75cf8d4cd558955a5c7638fba2d68f464/YM/TraceExtension/Chart/Line.lean\#L41-L44}{\texttt{Trace.Over.Gallery}}
    and the chain from it to the trace result.
  \item
    The Euclidean Wick exclusion
    \href{https://github.com/yang-mills-group/OpenAI-YM-20261028/blob/3456e5e75cf8d4cd558955a5c7638fba2d68f464/YM/TraceExtension/Wick/Refute.lean\#L75-L80}{\texttt{euclidean\_excludes}},
    with
    \href{https://github.com/yang-mills-group/OpenAI-YM-20261028/blob/3456e5e75cf8d4cd558955a5c7638fba2d68f464/YM/TraceExtension/Wick/Refute.lean\#L18-L22}{\texttt{real\_profile}}
    and
    \href{https://github.com/yang-mills-group/OpenAI-YM-20261028/blob/3456e5e75cf8d4cd558955a5c7638fba2d68f464/YM/TraceExtension/Wick/Log.lean\#L96-L108}{\texttt{no\_powerlog}}.
  \end{itemize}
\item
  \textbf{VII.4 (OPE and contacts).}

  \begin{itemize}
  \tightlist
  \item
    The existence of a coherent prescription
    (\href{https://github.com/yang-mills-group/OpenAI-YM-20261028/blob/3456e5e75cf8d4cd558955a5c7638fba2d68f464/YM/OPE/Splitting/Legacy.lean\#L20-L31}{\texttt{Prescription.chosen}}).
  \item
    The OPE services, which accompany the final object in VIII.1.g.
  \item
    The integrated estimates
    \href{https://github.com/yang-mills-group/OpenAI-YM-20261028/blob/3456e5e75cf8d4cd558955a5c7638fba2d68f464/YM/OPE/Orbits/Profile.lean\#L98-L108}{\texttt{estimates\_weighted}}.
  \item
    The contact resolution
    \href{https://github.com/yang-mills-group/OpenAI-YM-20261028/blob/3456e5e75cf8d4cd558955a5c7638fba2d68f464/YM/Contacts/Solstice/Result.lean\#L65-L69}{\texttt{Solstice.provided}},
    \href{https://github.com/yang-mills-group/OpenAI-YM-20261028/blob/3456e5e75cf8d4cd558955a5c7638fba2d68f464/YM/Contacts/Transport/Register.lean\#L52-L61}{\texttt{real\_basis}},
    \href{https://github.com/yang-mills-group/OpenAI-YM-20261028/blob/3456e5e75cf8d4cd558955a5c7638fba2d68f464/YM/Contacts/Solstice/Gates.lean\#L30-L36}{\texttt{installs}}
    and
    \href{https://github.com/yang-mills-group/OpenAI-YM-20261028/blob/3456e5e75cf8d4cd558955a5c7638fba2d68f464/YM/MultipleContact/Detector.lean\#L28-L35}{\texttt{detector\_exists}}.
  \end{itemize}
\item
  \textbf{VII.5 (asymptotic freedom).} The value of the slope,
  \(2b_0=11C_{\mathfrak g}/(24\pi^2)\), and \(C_{\mathfrak g}>0\)
  (\href{https://github.com/yang-mills-group/OpenAI-YM-20261028/blob/3456e5e75cf8d4cd558955a5c7638fba2d68f464/YM/Gauge/Casimir.lean\#L160-L172}{\texttt{casimir\_pos}}).
\end{enumerate}

\textbf{Standard mathematics used.}

\begin{itemize}
\tightlist
\item
  uniqueness of limits, and stability of limits under passage to
  subsequences;
\item
  sequential compactness of countable products of compact intervals
  (Tychonoff), with a diagonal extraction;
\item
  extension of convergence from a dense set by equicontinuity, and
  discrete Taylor estimates to recover derivatives, for (D5);
\item
  dominated convergence, and the Laplace transform of a positive measure
  (Källén--Lehmann representation of the trace two-point function:
  Källén 1952, Lehmann 1954), for VIII.1.8;
\item
  asymptotic scales \(r^a(\log r)^k\), for (T4).
\end{itemize}

The inputs rest on further standard mathematics:

\begin{itemize}
\tightlist
\item
  Bałaban's block-spin renormalization group for lattice gauge fields
  (Bałaban, CMP 1984--1989);
\item
  cluster expansions (Brydges, Les Houches 1984);
\item
  chessboard estimates (Fröhlich--Israel--Lieb--Simon, CMP 1978);
\item
  Osterwalder--Schrader reconstruction (CMP 1973, 1975) and Glaser's
  linear-growth argument;
\item
  decomposition of a translation-invariant state through the commutative
  algebra on the zero-energy subspace (Borchers; Ruelle; see
  Glimm--Jaffe, \emph{Quantum Physics}, §19.7);
\item
  one-loop asymptotic freedom (Gross--Wilczek, Politzer 1973) and the
  trace anomaly (Collins--Duncan--Joglekar 1977).
\end{itemize}

\textbf{Imported constants and assumptions.} The final existence theorem
has no additional analytic hypotheses. The run, local-law packages,
final object and (Core) are supplied by the preceding construction
results.

Parameters and auxiliary choices:

\begin{itemize}
\tightlist
\item
  The penalty \(b>0\) of the block-spin kernel. The coarse link \(V\) is
  drawn around the block prediction \(Q\) with weight
  \(\exp(-\beta b\,\varrho(VQ^{-1}))\)
  (\href{https://github.com/yang-mills-group/OpenAI-YM-20261028/blob/3456e5e75cf8d4cd558955a5c7638fba2d68f464/YM/RG/Block.lean}{\texttt{penaltyWeight}}).
  It plays the role of Bałaban's averaging parameter and is unrelated to
  \(b_*\).
\item
  The implication VIII.1.4 accepts a local-estimate choice \(\Delta\);
  VIII.1.a uses the fixed choice supplied by III.14.
\item
  The run-selection theorem is instantiated with the fixed dynamics and
  auxiliary estimates of III.2.
\end{itemize}

Constants:

\begin{itemize}
\tightlist
\item
  \(b_0=11C_{\mathfrak g}/(48\pi^2)\), and \(b_*=2b_0\) in (Reg1).
\item
  The trace amplitude
  \(A_{\mathfrak g}=12(d_2/b_0)^2\dim\mathfrak g/\pi^4\) with
  \(d_2=-b_0/2\), so \(A_{\mathfrak g}=3\dim\mathfrak g/\pi^4\).
\item
  \(\bar\beta>0\), the reference coupling at level \(0\).
\end{itemize}

\subsubsection{3. Strategy}\label{strategy-31}

Choose one run \(\mathsf R_{G,b}\) before any extraction order,
observable or test. Given \(\sigma\), first extract a refinement
\(\sigma\circ g\) on which the couplings at every fixed level have
ordered limits. Next extract \((\sigma\circ g)\circ w\), carrying a
continuum table and a final object. The later reconstruction, trace, OPE
and contact results attach to that object. Passing to further
subsequences preserves the earlier limits, so \(\tau=g\circ w\) serves
every conclusion.

At each fixed volume, the renormalization-step estimates make the
lattice correlations Cauchy on a dense union of grids. Equicontinuity
gives convergence on separated configurations; discrete Taylor estimates
recover all derivatives. Dominated convergence gives the smeared limits.
Uniqueness of limits identifies these full-sequence limits with the rows
of the extracted table, removing the mesh subsequence.

Specialize the detailed theorem to \(b=1\) and
\(\sigma=(\mathrm{id},p)\). The volume part of \(\tau\) is the required
\(q\). The public hierarchy is read from the final object and its OPE
reader. The bridges identify the Wilson measures and moments,
reconstruction axioms, OPE estimates and trace tests. The plaquette
identity, equality of Wick profiles and evaluation of the trace
amplitude are the required computations.

\subsubsection{4. Steps}\label{steps-31}

The internal statements use a fixed gauge group \((m,G,\mathfrak g)\).
When a run is specified, its campaign and pillars are those appearing in
(Fd) and VIII.1.5.

\paragraph{4.1 The run and the ultraviolet
endpoint}\label{the-run-and-the-ultraviolet-endpoint}

\textbf{Lemma VIII.1.a (the chosen run).} Fix a gauge group with
\(\mathfrak g\) simple and \(b>0\). There are a run \(\mathsf R_{G,b}\),
a campaign and pillars with penalty \(b\) satisfying (Fd). Write
\(\mathcal C_{G,b}\) for its family. \emph{Proof sketch.} Apply the
run-selection theorem of III.2 with the fixed dynamics, local estimates
and component prescription of III.14. The required schedule predicate is
identically true, so it is satisfied after any delay extension. Choose
one witness before specifying an extraction order.

\emph{Lean:}
\href{https://github.com/yang-mills-group/OpenAI-YM-20261028/blob/3456e5e75cf8d4cd558955a5c7638fba2d68f464/YM/Independent/ConstructionRun.lean\#L14-L28}{\texttt{constructionRun}}
(III.14) and
\href{https://github.com/yang-mills-group/OpenAI-YM-20261028/blob/3456e5e75cf8d4cd558955a5c7638fba2d68f464/YM/Independent/ConstructionRun.lean\#L30-L31}{\texttt{fundedRun}}.
\emph{(proof read; selection theorem statement only)}

\textbf{Lemma VIII.1.b (ultraviolet endpoint).} This is VIII.1.5.
\emph{Proof sketch.} The convergence theorem of III.13--III.15 supplies
the analytic families of couplings and activities. Its uniqueness clause
identifies successful trajectories with the same level-zero target, and
its analytic clause gives the bare-coupling interpolation. Removing the
field-independent scalar from the partition function gives the positive
lower bound; decay of the bare partition function then forces the scalar
density to \(-\infty\). The volume-uniform bounds and VIII.1.c give
(D5). The parameter-family construction gives item 8. Finally, refine to
a final object and apply VIII.1.d to obtain the sampler with the full
mesh sequence. The full funding hypothesis supplies the three
assumptions of VIII.1.5. \emph{Lean:}
\href{https://github.com/yang-mills-group/OpenAI-YM-20261028/blob/3456e5e75cf8d4cd558955a5c7638fba2d68f464/YM/FullVolume/Finale/Program.lean\#L101-L113}{\texttt{clients}};
\href{https://github.com/yang-mills-group/OpenAI-YM-20261028/blob/3456e5e75cf8d4cd558955a5c7638fba2d68f464/YM/RG17/Joint.lean\#L102-L111}{\texttt{jewelOf17}},
\href{https://github.com/yang-mills-group/OpenAI-YM-20261028/blob/3456e5e75cf8d4cd558955a5c7638fba2d68f464/YM/FullVolume/Finale/Program.lean\#L60-L68}{\texttt{wins}},
\href{https://github.com/yang-mills-group/OpenAI-YM-20261028/blob/3456e5e75cf8d4cd558955a5c7638fba2d68f464/YM/FullVolume/Finale/Program.lean\#L70-L77}{\texttt{bareTrial}}
and
\href{https://github.com/yang-mills-group/OpenAI-YM-20261028/blob/3456e5e75cf8d4cd558955a5c7638fba2d68f464/YM/R2I/Campaign/Accepted.lean\#L89-L91}{\texttt{Funded.jewel}}.
Further sources:
\href{https://github.com/yang-mills-group/OpenAI-YM-20261028/blob/3456e5e75cf8d4cd558955a5c7638fba2d68f464/YM/FullVolume/Glide/Vac.lean\#L31-L63}{\texttt{YM.PhysicalNext.Repeat.Legato.Campaign.Winter}},
\href{https://github.com/yang-mills-group/OpenAI-YM-20261028/blob/3456e5e75cf8d4cd558955a5c7638fba2d68f464/YM/FullVolume/Glide/Vac.lean\#L65-L89}{\texttt{YM.PhysicalNext.Repeat.Legato.Campaign.offers}},
\href{https://github.com/yang-mills-group/OpenAI-YM-20261028/blob/3456e5e75cf8d4cd558955a5c7638fba2d68f464/YM/FullVolume/Scalar.lean\#L87-L88}{\texttt{YM.FullVolume.stripped}},
\href{https://github.com/yang-mills-group/OpenAI-YM-20261028/blob/3456e5e75cf8d4cd558955a5c7638fba2d68f464/YM/FullVolume/Scalar.lean\#L97-L128}{\texttt{YM.FullVolume.saved}},
\href{https://github.com/yang-mills-group/OpenAI-YM-20261028/blob/3456e5e75cf8d4cd558955a5c7638fba2d68f464/YM/FullVolume/Scalar.lean\#L130-L143}{\texttt{YM.FullVolume.field\_independent}}.
\emph{(assembly, analytic convergence, tuning, costs and
parameter-family proofs read; their estimates are inputs from
III.12--III.15)}

\textbf{Lemma VIII.1.c (convergence along all spacings at fixed
volume).} This is VIII.1.6(a). \emph{Proof sketch.} Fix \(p,n,v\).

\begin{enumerate}
\def\labelenumi{\arabic{enumi}.}
\tightlist
\item
  The volume-uniform bounds control every fixed list of lattice
  differences on separated configurations. The bound and its cutoff
  threshold may depend on \(n,v\), the differences and the separation,
  but are uniform in \(p\).
\item
  The union of the configuration grids is dense. Grid nesting makes each
  point of this union a grid point at all sufficiently fine spacings. At
  these configurations, the renormalization-step comparison from III.12
  gives a Cauchy sequence along all \(N\). Equicontinuity extends
  convergence locally uniformly to \(\mathrm{NoHit}_p\).
\item
  A discrete Taylor estimate compares each first difference with a
  difference quotient over a fixed small physical displacement. The
  function values already converge; first let \(N\to\infty\), then let
  the displacement shrink. This proves that the first differences are
  Cauchy. Iterating gives limits of all differences, and the Taylor
  estimate identifies them as the derivatives of a smooth function
  \(f_{p,v}\). Multilinear interpolation has the same limits.
\item
  The lattice sums are integrals of piecewise constant functions on a
  fundamental cell. A uniform bound on the test support permits
  dominated convergence, giving the smeared limits.
\end{enumerate}

\emph{Lean:}
\href{https://github.com/yang-mills-group/OpenAI-YM-20261028/blob/3456e5e75cf8d4cd558955a5c7638fba2d68f464/YM/FullVolume/Fields/Filled.lean\#L45}{\texttt{realUV}},
\href{https://github.com/yang-mills-group/OpenAI-YM-20261028/blob/3456e5e75cf8d4cd558955a5c7638fba2d68f464/YM/FullVolume/Fields/Filled.lean\#L40-L44}{\texttt{purchased}},
\href{https://github.com/yang-mills-group/OpenAI-YM-20261028/blob/3456e5e75cf8d4cd558955a5c7638fba2d68f464/YM/FullVolume/Profiles.lean\#L42-L47}{\texttt{positioned}}
and
\href{https://github.com/yang-mills-group/OpenAI-YM-20261028/blob/3456e5e75cf8d4cd558955a5c7638fba2d68f464/YM/FullVolume/Discrete.lean\#L50-L52}{\texttt{FixedBounds}}.
Further sources:
\href{https://github.com/yang-mills-group/OpenAI-YM-20261028/blob/3456e5e75cf8d4cd558955a5c7638fba2d68f464/YM/FullVolume/Fields/Eves.lean\#L13-L14}{\texttt{YM.FullVolume.Eves.dense}},
\href{https://github.com/yang-mills-group/OpenAI-YM-20261028/blob/3456e5e75cf8d4cd558955a5c7638fba2d68f464/YM/FullVolume/Fields/Eves.lean\#L15-L24}{\texttt{YM.FullVolume.Eves.lateAnch}},
\href{https://github.com/yang-mills-group/OpenAI-YM-20261028/blob/3456e5e75cf8d4cd558955a5c7638fba2d68f464/YM/FullVolume/Tuned.lean\#L71-L125}{\texttt{YM.FullVolume.from\_previous}},
\href{https://github.com/yang-mills-group/OpenAI-YM-20261028/blob/3456e5e75cf8d4cd558955a5c7638fba2d68f464/YM/FullVolume/Wish.lean\#L22-L49}{\texttt{YM.FullVolume.honors}},
\href{https://github.com/yang-mills-group/OpenAI-YM-20261028/blob/3456e5e75cf8d4cd558955a5c7638fba2d68f464/YM/FullVolume/Wish.lean\#L61-L81}{\texttt{YM.FullVolume.augments}},
\href{https://github.com/yang-mills-group/OpenAI-YM-20261028/blob/3456e5e75cf8d4cd558955a5c7638fba2d68f464/YM/FullVolume/Smears.lean\#L49-L82}{\texttt{YM.FullVolume.pudding}}.
\emph{(proof read, including dense-grid convergence, Taylor recovery and
smearing; the renormalization-step estimate is an input)}

\textbf{Lemma VIII.1.d (removing the mesh subsequence).} This is
VIII.1.6(b). \emph{Proof sketch.} The seed of \(\mathfrak F\) gives, for
each admissible \(p\), \(C_{\sigma^{\rm m}(j),p}(v;t)\to R_p\), and also
\(R_{\sigma^{\rm v}(k)}\to\int tS(v;\cdot)\). By VIII.1.c,
\(C_{N,p}(v;t)\to\int t f_{p,v}\) along all \(N\). Since
\(\sigma^{\rm m}\) is cofinal, uniqueness of limits gives
\(R_p=\int t f_{p,v}\). Thus the same rows \(R_p\) are limits along all
\(N\), and the volume limit is unchanged. \emph{Lean:}
\href{https://github.com/yang-mills-group/OpenAI-YM-20261028/blob/3456e5e75cf8d4cd558955a5c7638fba2d68f464/YM/FullVolume/Finale/Program.lean\#L117-L126}{\texttt{from\_spec}}.

\textbf{Lemma VIII.1.e (OPE and logarithmic packages).} Let
\(\mathsf R\) be a run with campaign \(T\) and pillars \(W\). Assume
\(\mathfrak g\) simple and (Fd). Then the family carries the OPE package
and its logarithmic extension. \emph{Proof sketch.} The pillars,
polynomial schedule compatibility, halo condition and weight data supply
the ordinary package. The logarithmic extension adds the logarithmic
exit estimate, its budget, the fixed support exponent, and divergence of
the composed schedule \(m\mapsto Q(J(m))\). The last property follows
from \(J(m)\to\infty\) and \(Q(k)\to\infty\).

\emph{Lean:}
\href{https://github.com/yang-mills-group/OpenAI-YM-20261028/blob/3456e5e75cf8d4cd558955a5c7638fba2d68f464/YM/R2I/Campaign/Accepted.lean\#L100-L106}{\texttt{Funded.logPackage}}
and
\href{https://github.com/yang-mills-group/OpenAI-YM-20261028/blob/3456e5e75cf8d4cd558955a5c7638fba2d68f464/YM/RG17/Witness.lean\#L98-L106}{\texttt{opePackage17}}.
\emph{(assembly proofs read)}

\paragraph{4.2 The data at one extraction
order}\label{the-data-at-one-extraction-order}

\textbf{Lemma VIII.1.f (limits of the couplings at fixed level).} Assume
(D3) holds for \(\mathcal C\), and let \(\sigma\) be an extraction
order. Then there are \(g\) and \((B_m)_m\) such that, for every \(m\),
\(\beta_{N,p}(m)\) has the ordered limit \(B_m\) along
\(\sigma\circ g\). Moreover \(B_m/\log(1/a_m)\to 2b_0\). \emph{Proof
sketch.} Set \(c_N(p,m):=\beta_{N,p}(m)\) for \(m\le N\) and
\(c_N(p,m):=\hat\beta_m\) for \(m>N\). By the band condition of (D3),
\(c_N\in\prod_{(p,m)}[\hat\beta_m/2,2\hat\beta_m]\), a compact
metrizable space.

\begin{enumerate}
\def\labelenumi{\arabic{enumi}.}
\tightlist
\item
  Extract \(g^{\rm m}\) with \(c_{\sigma^{\rm m}(g^{\rm m}(j))}\to a\)
  in the product.
\item
  Extract \(g^{\rm v}\) with
  \(a(\sigma^{\rm v}(g^{\rm v}(k)),\cdot)\to B\) in
  \(\prod_m[\hat\beta_m/2,2\hat\beta_m]\).
\item
  For fixed \(m\) the truncation is eventually inactive.
\item
  The running law of (D3) gives
  \(|B_m-\hat\beta_m|\le\varepsilon\log(1/a_m)\) for large \(m\), hence
  the slope.
\end{enumerate}

The slope is not used by the main theorem. \emph{Lean:}
\href{https://github.com/yang-mills-group/OpenAI-YM-20261028/blob/3456e5e75cf8d4cd558955a5c7638fba2d68f464/YM/PertExtension/Running.lean\#L40-L56}{\texttt{Running.extraction}}
and
\href{https://github.com/yang-mills-group/OpenAI-YM-20261028/blob/3456e5e75cf8d4cd558955a5c7638fba2d68f464/YM/PertExtension/Running.lean\#L64-L84}{\texttt{Running.discrete}}.

\textbf{Lemma VIII.1.g (a final object on a refinement).} Let
\(\mathfrak g\) be simple, let \(\mathsf R\) be a run with the OPE
package, and let \(\sigma\) be an extraction order. Then there are a
refinement \(w\), a table \(S\), a final object \(\mathfrak F\) on
\((\mathcal C_{\mathsf R},\sigma\circ w,S)\) and scaling data
\(\Sigma\), with OPE services and (D14). \emph{Proof sketch.}

\begin{enumerate}
\def\labelenumi{\arabic{enumi}.}
\tightlist
\item
  The run's extraction property supplies a table and seed on a
  refinement, together with the trajectory witnesses (D6).
\item
  Add the free Maxwell data, with the translation normalization in which
  \(\kappa^{-1}J\) is the free Maxwell stress tensor.
\item
  The weak comparison estimates give the infrared data, scaling data
  \(\Sigma\) and the first comparison.
\item
  The symmetry results of VI.1--VI.2 produce the derivatives, rotations
  and stress tensor with their local laws.
\item
  The phase results B01--B18 produce the selection and its certificate.
  The reconstruction theorem then supplies the mixed and selected
  continuations and the selected Wightman theory. These data form
  \(\mathfrak F\) with the prescribed field algebra.
\item
  The leading mixed comparison, the one-dimensional space of conserved
  tensors, and the limits of derivatives and rotations give (D14).
  \emph{Lean:}
  \href{https://github.com/yang-mills-group/OpenAI-YM-20261028/blob/3456e5e75cf8d4cd558955a5c7638fba2d68f464/YM/Pert/Construction.lean\#L21-L31}{\texttt{Pert.withFinal}};
  \href{https://github.com/yang-mills-group/OpenAI-YM-20261028/blob/3456e5e75cf8d4cd558955a5c7638fba2d68f464/YM/RG17/Joint.lean\#L60-L90}{\texttt{Run.extractible17}}.
  \emph{(assembly and final-object proofs read; analytic and symmetry
  inputs as stated)}
\end{enumerate}

\textbf{Lemma VIII.1.h (trace result for every final object).} Assume
(Fd) and \(\mathfrak g\) simple. For every extraction order \(\sigma\),
every table \(S\) and every final object \(\mathfrak F\) on
\((\mathcal C_{\mathsf R},\sigma,S)\):

\begin{itemize}
\item
  \(d_2=-b_0/2\) and \(A_{\mathfrak g}>0\);
\item
  \(\Theta\) is fixed by the adjoint and Lorentz actions, and
  \(\langle\Omega,\Theta(f)\Omega\rangle=0\) for every Schwartz test
  \(f\);
\item
  on the common invariant domain,

  \[\Theta(f)=\Big(\sum_aT_{aa}\Big)(f)-\Big\langle\sum_aT_{aa}\Big\rangle_{\xi_*}\Big(\int f\Big)\mathrm{id},\]
  and \(\Theta\ne0\);
\item
  \(s^8(\log 1/s)^2G_\Theta(s)\to A_{\mathfrak g}\) as \(s\downarrow0\);
\item
  the trace has a unique spectral measure \(\mu\), with

  \[\rho^{-8}(\log\rho)^2\int\varphi(p/\rho)\,\mu(dp)\longrightarrow
  \frac{A_{\mathfrak g}}{1536\pi}\int_{\overline V_+}\mathsf m(p)^2\varphi(p)\,dp
  \qquad(\rho\to\infty),\] for every compactly supported continuous
  \(\varphi\) and every Schwartz \(\varphi\), where
  \(\mathsf m(p)=(p^4)^2-|\mathbf p|^2\) and
  \(\overline V_+=\{p:p^4\ge|\mathbf p|\}\);
\item
  there exist a scaling hierarchy, a Wilson coupling and an ultrafilter
  compatible with \(\sigma\) for which a normalized position-space
  scheme exists. Every scheme with that hierarchy and ultrafilter
  satisfying the stated alignment, leading-derivative and scalar
  normalization conditions admits the stress comparison of VII.3;
\item
  the momentum-space Wick countertest of VIII.1.8 and the Maxwell
  relative-locality countertests of VIII.1.n hold, including the version
  localized in any ball.
\end{itemize}

\emph{Proof sketch.} The phase-comparison condition in (Fd) supplies the
trace comparison data at \(\sigma\). The trace asymptotic, spectral law
and exclusions of VII.2--VII.3 then apply to this final object.
Centering and scalar covariance follow from the stress-tensor laws. The
stress-comparison theorem first constructs one admissible normalized
scheme, then proves the comparison for every scheme satisfying those
same conditions.

\emph{Lean:}
\href{https://github.com/yang-mills-group/OpenAI-YM-20261028/blob/3456e5e75cf8d4cd558955a5c7638fba2d68f464/YM/TraceExtension/Endpoint.lean\#L99-L120}{\texttt{Trace.Target.Result}},
\href{https://github.com/yang-mills-group/OpenAI-YM-20261028/blob/3456e5e75cf8d4cd558955a5c7638fba2d68f464/YM/TraceExtension/Endpoint.lean\#L138-L152}{\texttt{assembled}},
\href{https://github.com/yang-mills-group/OpenAI-YM-20261028/blob/3456e5e75cf8d4cd558955a5c7638fba2d68f464/YM/TraceExtension/Agora.lean\#L69-L74}{\texttt{arenaT}},
\href{https://github.com/yang-mills-group/OpenAI-YM-20261028/blob/3456e5e75cf8d4cd558955a5c7638fba2d68f464/YM/TraceExtension/Belfry.lean\#L31-L37}{\texttt{crest\_ready}}
and
\href{https://github.com/yang-mills-group/OpenAI-YM-20261028/blob/3456e5e75cf8d4cd558955a5c7638fba2d68f464/YM/R2I/Campaign/Accepted.lean\#L93-L95}{\texttt{Funded.gallery}}.
\emph{(assembly proofs read; the trace estimates are inputs from
VII.2--VII.3)}

\textbf{Lemma VIII.1.i (Wick exclusion in momentum space).} This is
VIII.1.8. \emph{Proof sketch.}

\begin{enumerate}
\def\labelenumi{\arabic{enumi}.}
\item
  Let \(\mu\) be the trace spectral measure. It is tempered and
  supported in the forward cone. For a joint Schwartz test \(\varphi\),
  the selected two-point function is the integral of its
  relative-momentum transform against \(\mu\), and

  \[G_\Theta(s)=\int e^{-sp^4}\,\mu(dp)\qquad(s>0).\]
\item
  Suppose the Wick identity held for every allowed pair \((f,g)\). Then
  \(\int v\,d\mu=\sum_\alpha c_\alpha\mathrm{read}_\alpha(v)\) for every
  compactly supported Schwartz \(v\). To see this, choose \(u\) with
  \(|\mathcal F_+u|=1\) on a ball containing \(\mathrm{supp}\,v\), and
  use Fourier surjectivity to choose \(w\) with
  \(\mathcal F_+w=(\mathcal F_+u)v\). Thus
  \(\overline{\mathcal F_+u}\mathcal F_+w=v\).
\item
  For each \(s>0\), approximate \(e^{-sp^4}\) by compactly supported
  smooth cutoffs bounded in absolute value by \(e^{-sp^4}\). Exponential
  integrability for \(\mu\) and for every mass-shell term permits
  dominated convergence on both sides. Hence
  \(G_\Theta(s)=\sum_\alpha c_\alpha K_\alpha(s)\) for every \(s>0\).
\item
  The resulting Wick profile contradicts VIII.1.r(c).
\end{enumerate}

\emph{Lean:}
\href{https://github.com/yang-mills-group/OpenAI-YM-20261028/blob/3456e5e75cf8d4cd558955a5c7638fba2d68f464/YM/TraceExtension/Wick/Refute.lean\#L86-L105}{\texttt{polynomial\_excludes}},
\href{https://github.com/yang-mills-group/OpenAI-YM-20261028/blob/3456e5e75cf8d4cd558955a5c7638fba2d68f464/YM/TraceExtension/Wick/Refute.lean\#L62-L73}{\texttt{identifies\_euclid}}
and
\href{https://github.com/yang-mills-group/OpenAI-YM-20261028/blob/3456e5e75cf8d4cd558955a5c7638fba2d68f464/YM/TraceExtension/WilsonTrace.lean\#L99-L119}{\texttt{traceSpectrum\_exists}}.
Further sources:
\href{https://github.com/yang-mills-group/OpenAI-YM-20261028/blob/3456e5e75cf8d4cd558955a5c7638fba2d68f464/YM/TraceExtension/Momentum.lean\#L22}{\texttt{YM.TraceExtension.productTest}},
\href{https://github.com/yang-mills-group/OpenAI-YM-20261028/blob/3456e5e75cf8d4cd558955a5c7638fba2d68f464/YM/TraceExtension/Momentum.lean\#L26-L35}{\texttt{YM.TraceExtension.waveTest\_product}}.

\textbf{Lemma VIII.1.j (reconstruction in every phase; input from V.3).}
For every final object \(\mathfrak F\) there are reconstruction data
\(\mathfrak L\) with the following properties. The phasewise assertions
hold for every \(\xi\in X\).

\begin{enumerate}
\def\labelenumi{\arabic{enumi}.}
\item
  For every \(\xi\in X\), a Wightman theory \(\mathfrak W_\xi\) realizes
  the continuation of that phase, with the same labels. At \(\xi_*\) it
  is the selected theory of \(\mathfrak F\).
\item
  The comparison map from \(\mathfrak W_\xi\) to its
  Osterwalder--Schrader space is onto.
\item
  The Euclidean ground state is unique, and the zero-energy subspace is
  \(\mathbb C\Omega_\xi\).
\item
  For every finite sum \(A\) of field words there is \(M_A\ge0\),
  independent of \(\xi\), such that

  \[\mu_{A\Omega_\xi}(\{0<p^4\le u\})\le M_Au^2\quad(0<u\le1),\]
  \[0\le\langle A\Omega_\xi,(e^{-sH_\xi}-P_{0,\xi})A\Omega_\xi\rangle\le M_As^{-2}\quad(s\ge1).\]
\item
  The bounded weak commutant, defined using the unbounded adjoints on
  the field-word domain, is \(\mathbb C\,\mathrm{id}\).
\item
  The stress tensor continues to every phase, and each Wightman
  distribution evaluated on a fixed joint Schwartz test is continuous in
  \(\xi\).
\item
  \(\int S^\xi\,d\nu=S\) on separated words, and
  \(\int\mathcal W^\xi\,d\nu=\mathcal W^{\rm mix}\) on every joint
  Schwartz test.
\end{enumerate}

\emph{Proof sketch (V.1, V.3).} The phases are characters of the
commutative \(C^*\)-algebra generated by time-slab word operators
compressed to the mixed Osterwalder--Schrader zero-energy subspace. Its
vacuum measure is \(\nu\). Translate a separated word into a slab to
obtain the Euclidean barycentre identity. Apply the reconstruction,
low-energy, heat-decay, irreducibility and stress-continuation results
of V.3 to every character, retaining the original selected theory at
\(\xi_*\). Their continuity and boundary barycentre identities give the
last two assertions.

\emph{Lean:}
\href{https://github.com/yang-mills-group/OpenAI-YM-20261028/blob/3456e5e75cf8d4cd558955a5c7638fba2d68f464/YM/EveryPhase/Extension.lean\#L75-L80}{\texttt{all\_characters}}
and
\href{https://github.com/yang-mills-group/OpenAI-YM-20261028/blob/3456e5e75cf8d4cd558955a5c7638fba2d68f464/YM/EveryPhase/Extension.lean\#L51-L71}{\texttt{Completed}}.
\emph{(statement and top-level proof read)}

\textbf{Lemma VIII.1.k (the OPE reader and its estimates).} Let
\(\mathfrak F\) be a final object on \((\mathcal C,\sigma,S)\) with OPE
services for that family, order and table. Choose a coherent
prescription \(P\) for the field derivatives, adapted to the symmetric
laws of the table and to the transport injection. Let \(Q\) be the
reader formed by \(P\), its coherence, the Euclidean kernel and the
transport data (injection, orientation, local laws). Then:

\begin{itemize}
\tightlist
\item
  the OPE foundation holds for \(Q\);
\item
  the uniform nested OPE holds;
\item
  the all-shapes OPE holds for every coherent prescription;
\item
  under the dense-coverage property of the sharp Gram family, a
  compatible sharp Gram allocation exists for every ultrafilter finer
  than the extraction filter of \(\sigma\).
\end{itemize}

\emph{Proof sketch.} Extend the finite-degree prescriptions compatibly
and take their common restriction in each degree. The OPE services give
the coefficient and remainder estimates, which also supply the nested
expansions and the all-shapes expansion. Dense coverage supplies the
test portfolios used to construct the sharp Gram allocation.

\emph{Lean:}
\href{https://github.com/yang-mills-group/OpenAI-YM-20261028/blob/3456e5e75cf8d4cd558955a5c7638fba2d68f464/YM/OPE/Splitting/Legacy.lean\#L20-L31}{\texttt{Prescription.chosen}}
(VII.4),
\href{https://github.com/yang-mills-group/OpenAI-YM-20261028/blob/3456e5e75cf8d4cd558955a5c7638fba2d68f464/YM/OPE/Orbits/Kernel.lean\#L67-L68}{\texttt{Orbits.produce}},
\href{https://github.com/yang-mills-group/OpenAI-YM-20261028/blob/3456e5e75cf8d4cd558955a5c7638fba2d68f464/YM/OPE/Nesting/Endpoint.lean\#L31-L36}{\texttt{Nesting.uniform}}
and
\href{https://github.com/yang-mills-group/OpenAI-YM-20261028/blob/3456e5e75cf8d4cd558955a5c7638fba2d68f464/YM/RG17/Canvas.lean\#L54-L58}{\texttt{attains\_mantle17}}.
\emph{(assembly proofs read; OPE estimates as stated)}

\textbf{Lemma VIII.1.l (assembly of the full package).} This is
VIII.1.4. \emph{Proof sketch.} The proof has three parts.

\begin{itemize}
\tightlist
\item
  \emph{Global items.} (UV\(_\infty\)) comes from VIII.1.b and the
  logarithmic package from VIII.1.e. The trace result for every
  \((\sigma,S,\mathfrak F)\) comes from VIII.1.h. Coverage comes from
  the phase-comparison and coverage clause of (Fd).
\item
  \emph{At an order \(\sigma\).} Apply VIII.1.f to obtain \((g,B)\).
  Apply VIII.1.g to \(\sigma\circ g\) to obtain
  \((w,S,\mathfrak F,\Sigma)\). Apply VIII.1.j to obtain
  \(\mathfrak L\), and VIII.1.k to obtain \(Q\). Then:

  \begin{itemize}
  \tightlist
  \item
    subsequence stability carries the coupling limits from
    \(\sigma\circ g\) to
    \((\sigma\circ g)\circ w=\sigma\circ(g\circ w)\);
  \item
    VIII.1.d, with the full-mesh clause of (UV\(_\infty\)), gives
    \(\mathrm{Ord}_{(\mathrm{id},(\sigma\circ g\circ w)^{\rm v})}\);
  \item
    the OPE facilities consist of the services, the run parameters and
    the logarithmic package.
  \end{itemize}
\item
  \emph{Output.} The refinement returned is \(\tau=g\circ w\).
\end{itemize}

\emph{Lean:}
\href{https://github.com/yang-mills-group/OpenAI-YM-20261028/blob/3456e5e75cf8d4cd558955a5c7638fba2d68f464/YM/R2I/Campaign/Accepted.lean\#L150-L167}{\texttt{fromFunded}}
and
\href{https://github.com/yang-mills-group/OpenAI-YM-20261028/blob/3456e5e75cf8d4cd558955a5c7638fba2d68f464/YM/Extraction/Orders.lean\#L66-L68}{\texttt{Order.retained}}.

\paragraph{4.3 The internal headline}\label{the-internal-headline}

\textbf{Lemma VIII.1.m (contact devices).} Let \(\mathfrak F\) carry a
reader \(Q\) and OPE services. Then there are devices \(\mathsf M\) with
the following properties:

\begin{itemize}
\tightlist
\item
  the prescription of \(\mathsf M\) is the reader's prescription;
\item
  the full contact resolution holds;
\item
  each OPE stratum has a complex basis of real fields, and the
  degree-zero basis contains the unit field;
\item
  there are positive thresholds \(t_n\to0\) with \(4t_{n+1}<t_n\) and
  \(2(n+1)t_n\le\mathrm{clearance}_{\mathsf M}(n)\).
\end{itemize}

\emph{Proof sketch.} Use the reader's coherent prescription, the
rotations, the Euclidean hierarchy and the phase transport laws,
together with a multiple-contact detector in each degree. The contact
theorem of VII.4 gives the resolution from the OPE services. Real
conjugation permits a real basis in each stratum, containing the unit in
degree zero. Successively choose each positive threshold below one
eighth of its predecessor, \(2^{-n}\) and
\(\mathrm{clearance}_{\mathsf M}(n)/(2(n+1))\). \emph{Lean:}
\href{https://github.com/yang-mills-group/OpenAI-YM-20261028/blob/3456e5e75cf8d4cd558955a5c7638fba2d68f464/YM/Nonperturbative/Continuum.lean\#L14-L24}{\texttt{contacts\_on}}
and
\href{https://github.com/yang-mills-group/OpenAI-YM-20261028/blob/3456e5e75cf8d4cd558955a5c7638fba2d68f464/YM/Contacts/Solstice/Result.lean\#L52-L63}{\texttt{MultipleLight}}.
Further sources:
\href{https://github.com/yang-mills-group/OpenAI-YM-20261028/blob/3456e5e75cf8d4cd558955a5c7638fba2d68f464/YM/Contacts/Solstice/Gates.lean\#L23-L28}{\texttt{YM.Contacts.Solstice.Thresholds}}.

\textbf{Lemma VIII.1.n (continuum theorem).} This is VIII.1.3.
\emph{Proof sketch.} Take \((T,W)\) and (Fd) from the specification of
\(\mathsf R_{G,b}\), and apply the refinement clause of VIII.1.4 at
\(\sigma\). The coupling limits, the all-shapes OPE and the sharp Gram
allocations are omitted. Contacts come from VIII.1.m. (Core) comprises:

\begin{itemize}
\tightlist
\item
  the OPE services;
\item
  the reconstruction in every phase;
\item
  the eight trace tests from VIII.1.h;
\item
  the OPE foundation;
\item
  the integrated remainder estimates of VII.4: fix cluster fields
  \(a_i\in F_{d_i}\) and spectator fields. For a measurable parameter
  space, an integrable complex density \(w\) against a positive measure
  \(\mu\), with \(\int|w|\,d\mu\le B\) and an almost everywhere
  measurable configuration map that is admissible at scale \(r\)
  wherever \(w\ne0\), the weighted raw correlation, finite OPE sum and
  their difference are integrable. The norm of the integrated difference
  is at most \(B\,r^{D+1-\sum_i d_i-\varepsilon}\), for every almost
  everywhere measurable choice of phases and for the mixed state. The
  radius threshold depends on the finite field lists, geometric
  constants, \(\varepsilon>0\), truncation degree \(D\), parameter space
  and mass bound \(B\), and precedes the choice of \(\mu,w\),
  configurations and phases.
\end{itemize}

The internal trace-test bundle has eight clauses:

\begin{itemize}
\item
  (T1);
\item
  the operator identity of VIII.1.h;
\item
  the asymptotics;
\item
  \(A_{\mathfrak g}>0\);
\item
  the unique spectral law;
\item
  the Wick countertest;
\item
  for every \(n\) and every complex-linear isometry \(U\) from
  \(\mathcal H\) into the \(n\)-colour free Maxwell Fock space,
  preserving the vacuum, intertwining translations and taking values in
  symmetric transverse waves, there are compactly supported Schwartz
  tests \(f,h_1,\ldots,h_n\), with each \(h_i\) real and spacelike
  separated from \(f\), such that

  \[\langle W(-h)\Omega,U\Theta(f)\Omega\rangle\ne
  \langle U\Theta(\bar f)\Omega,W(h)\Omega\rangle;\]
\item
  for every such \(U\), every probability measure on an index space and
  every family of maps preserving inner products on physical waves and
  fixing \(U\mathcal H\) pointwise, the analogous identity with the
  coherent vectors averaged under those maps also fails. For every
  \(r>0\), the test \(f\) can be supported in an open set whose closure
  lies in \(B(0,r)\), with the real tests \(h_i\) supported in a common
  open set whose closure is spacelike separated from the closure of the
  support region for \(f\).
\end{itemize}

Finally, the degree-four determining tests of VII.1 convert (D14) into
(D15). \emph{Lean:}
\href{https://github.com/yang-mills-group/OpenAI-YM-20261028/blob/3456e5e75cf8d4cd558955a5c7638fba2d68f464/YM/Nonperturbative/Continuum.lean\#L28-L48}{\texttt{continuum\_with\_position\_tests}},
\href{https://github.com/yang-mills-group/OpenAI-YM-20261028/blob/3456e5e75cf8d4cd558955a5c7638fba2d68f464/YM/Independent/Core.lean\#L18-L34}{\texttt{Core}},
\href{https://github.com/yang-mills-group/OpenAI-YM-20261028/blob/3456e5e75cf8d4cd558955a5c7638fba2d68f464/YM/Independent/Core.lean\#L38-L44}{\texttt{Core.of\_receipts}},
\href{https://github.com/yang-mills-group/OpenAI-YM-20261028/blob/3456e5e75cf8d4cd558955a5c7638fba2d68f464/YM/Integration/ManuscriptTraceTests.lean\#L77-L89}{\texttt{ManuscriptTraceTests.of\_result}}
and
\href{https://github.com/yang-mills-group/OpenAI-YM-20261028/blob/3456e5e75cf8d4cd558955a5c7638fba2d68f464/YM/Pert/Position/Four.lean\#L49-L53}{\texttt{Pert.Position.four\_on}}.
\emph{(assembly and degree-four normalization proofs read)}

\textbf{Lemma VIII.1.o (internal headline).} For a gauge group with
simple \(\mathfrak g\) and \(b>0\), the internal detailed existence
statement holds with \(\mathcal C=\mathcal C_{G,b}\). \emph{Proof
sketch.} Item by item:

\begin{itemize}
\tightlist
\item
  (D3) comes from the run, and (D6) from the final object.
\item
  \emph{Reconstruction.} It consists of:

  \begin{itemize}
  \tightlist
  \item
    the mixed and selected continuations;
  \item
    the internal Wightman axioms for the selected theory and for every
    \(\mathfrak W_\xi\): normalized cyclic vacuum, dense domain,
    temperedness, adjoints, the moment identity, positivity, locality,
    covariance, strong continuity, spectral representation, forward
    support, and uniqueness of the invariant vector;
  \item
    the \(\nu\)-average of the phase Wightman distributions.
  \end{itemize}
\item
  \emph{Separated OPE.} It consists of the coefficient bounds from the
  OPE foundation, the single-centre remainder bound and the integrated
  bounds of (Core).
\item
  \emph{Trace clauses.} The seven headline clauses are centering, the
  operator identity, short-distance asymptotics, positive amplitude, and
  the Wick, Maxwell and averaged Maxwell countertests. The unique
  spectral law is retained separately in (Core).
\end{itemize}

\emph{Lean:}
\href{https://github.com/yang-mills-group/OpenAI-YM-20261028/blob/3456e5e75cf8d4cd558955a5c7638fba2d68f464/YM/Nonperturbative/Headlines.lean\#L19-L36}{\texttt{Nonperturbative.existence}},
\href{https://github.com/yang-mills-group/OpenAI-YM-20261028/blob/3456e5e75cf8d4cd558955a5c7638fba2d68f464/YM/Nonperturbative/HeadlineSpec.lean\#L75-L81}{\texttt{Reconstruction}},
\href{https://github.com/yang-mills-group/OpenAI-YM-20261028/blob/3456e5e75cf8d4cd558955a5c7638fba2d68f464/YM/Nonperturbative/HeadlineSpec.lean\#L52-L71}{\texttt{WightmanAxioms}},
\href{https://github.com/yang-mills-group/OpenAI-YM-20261028/blob/3456e5e75cf8d4cd558955a5c7638fba2d68f464/YM/Nonperturbative/HeadlineSpec.lean\#L24-L48}{\texttt{SeparatedOPE}},
\href{https://github.com/yang-mills-group/OpenAI-YM-20261028/blob/3456e5e75cf8d4cd558955a5c7638fba2d68f464/YM/Nonperturbative/HeadlineSpec.lean\#L128-L138}{\texttt{NestedOPE}},
\href{https://github.com/yang-mills-group/OpenAI-YM-20261028/blob/3456e5e75cf8d4cd558955a5c7638fba2d68f464/YM/Nonperturbative/HeadlineSpec.lean\#L84-L120}{\texttt{TraceExclusions}}
and
\href{https://github.com/yang-mills-group/OpenAI-YM-20261028/blob/3456e5e75cf8d4cd558955a5c7638fba2d68f464/YM/Independent/Core.lean\#L49-L63}{\texttt{Core.point}}.

\paragraph{4.4 Translation to the public
statement}\label{translation-to-the-public-statement}

\textbf{Lemma VIII.1.p (the public Wilson family and hierarchy).} Let
\(\mathfrak F\) be a final object on \((\mathcal C,\sigma,S)\) and \(Q\)
a reader. The public Wilson family \(\mathcal W(\mathcal C)\) is defined
by:

\begin{itemize}
\tightlist
\item
  \(L_p\) and \(D_N\) are the period and denominator sequences of
  \(\mathcal C\);
\item
  \(\beta_{N,p}=\beta_{N,p}(N)\);
\item
  \(h\) is the normalized Haar measure;
\item
  \(\mathcal T\) is the set of internal tokens;
\item
  \(I_{N,p}\) is the internal insertion.
\end{itemize}

The public Wilson energy equals the internal one. The proof matches the
plaquette \(U(y,\mu)U(y+e_\mu,\nu)U(y+e_\nu,\mu)^{-1}U(y,\nu)^{-1}\)
with the internal face value. Hence the measures, the lattice moments
and the smeared moments agree.

The hierarchy \(T(\mathfrak F,Q)\) has the following components:

\begin{itemize}
\tightlist
\item
  \(\mathcal H\) and \(\mathfrak W\) are those of the selected
  reconstruction, with the same Hilbert space, domain, labels and
  operators;
\item
  \(\mathcal L\) is the filtered union of the field spaces, with
  \(F_d=\) the quotient field spaces of \(S\), \(\iota_d=\) the
  field-label inclusion and \(\pi=\) the class map of tokens;
\item
  \((X,\nu,\xi_*)\) are the phase space, its measure and the selected
  point;
\item
  \(S^s\) is the table read in state \(s\);
\item
  \(\mathsf F_n\) is the selected continuation;
\item
  \((\mathrm{Ch}_j,O_{j,\kappa},c_{j,\kappa})\) come from \(Q\);
\item
  \(\Theta=\sum_aT_{aa}-\langle\sum_aT_{aa}\rangle_{\xi_*}\mathbf 1_4\).
\end{itemize}

\emph{Lean:}
\href{https://github.com/yang-mills-group/OpenAI-YM-20261028/blob/3456e5e75cf8d4cd558955a5c7638fba2d68f464/WilsonBridge.lean\#L32-L41}{\texttt{energy\_eq}},
\href{https://github.com/yang-mills-group/OpenAI-YM-20261028/blob/3456e5e75cf8d4cd558955a5c7638fba2d68f464/WilsonBridge.lean\#L43-L46}{\texttt{law\_eq}},
\href{https://github.com/yang-mills-group/OpenAI-YM-20261028/blob/3456e5e75cf8d4cd558955a5c7638fba2d68f464/WilsonBridge.lean\#L49-L55}{\texttt{moment\_eq}},
\href{https://github.com/yang-mills-group/OpenAI-YM-20261028/blob/3456e5e75cf8d4cd558955a5c7638fba2d68f464/WilsonBridge.lean\#L57-L60}{\texttt{smear\_eq}}
and
\href{https://github.com/yang-mills-group/OpenAI-YM-20261028/blob/3456e5e75cf8d4cd558955a5c7638fba2d68f464/HierarchyBridge.lean\#L27-L42}{\texttt{hierarchy}}.

\textbf{Lemma VIII.1.q (regularity and the Wilson limit).} Let
\(\mathfrak F\) be a final object on \((\mathcal C,\sigma,S)\) and \(Q\)
a reader.

\begin{enumerate}
\def\labelenumi{(\alph{enumi})}
\tightlist
\item
  If (D3) holds and \(\mathfrak g\) is simple, then (Reg1)--(Reg4) hold:
\end{enumerate}

\begin{itemize}
\tightlist
\item
  (Reg1): the running law of (D3) at \(m=N\) gives
  \(|\beta_{N,p}/\log D_N-2b_0|<\varepsilon\) for every
  \(\varepsilon>0\), all sufficiently large \(N\) and all \(p\), and
  \(2b_0>0\) by \(C_{\mathfrak g}>0\);
\item
  (Reg2) and (Reg3): continuity and gauge invariance of the local
  sections;
\item
  (Reg4): the stencil of the local section.
\end{itemize}

\begin{enumerate}
\def\labelenumi{(\alph{enumi})}
\setcounter{enumi}{1}
\tightlist
\item
  If \(\mathrm{Ord}_{(\mathrm{id},s)}(\mathcal C,S)\) holds, then (P1)
  holds along \(s\). The smeared moments agree by VIII.1.p. The mixed
  correlation of \(\pi(v)\) equals \(S(v;\cdot)\) on
  \(\mathrm{Conf}_n\), and \(t=0\) off \(\mathrm{Conf}_n\). \emph{Lean:}
  \href{https://github.com/yang-mills-group/OpenAI-YM-20261028/blob/3456e5e75cf8d4cd558955a5c7638fba2d68f464/HierarchyBridge.lean\#L16-L25}{\texttt{regular}},
  \href{https://github.com/yang-mills-group/OpenAI-YM-20261028/blob/3456e5e75cf8d4cd558955a5c7638fba2d68f464/YM/Pert/Asymptotics/RawAxis.lean\#L15-L30}{\texttt{zeroUp\_bound}},
  \href{https://github.com/yang-mills-group/OpenAI-YM-20261028/blob/3456e5e75cf8d4cd558955a5c7638fba2d68f464/WilsonBridge.lean\#L62-L67}{\texttt{fromWilson}}
  and
  \href{https://github.com/yang-mills-group/OpenAI-YM-20261028/blob/3456e5e75cf8d4cd558955a5c7638fba2d68f464/HierarchyBridge.lean\#L44-L56}{\texttt{wilsonLimit}}.
\end{enumerate}

\textbf{Lemma VIII.1.r (reconstruction, OPE and trace clauses).} Let
\(\mathfrak F\) be a final object and \(Q\) a reader for it. The
conclusions concern \(T(\mathfrak F,Q)\).

\begin{enumerate}
\def\labelenumi{(\alph{enumi})}
\tightlist
\item
  The reconstruction in every phase implies (R1)--(R15). Clause by
  clause:
\end{enumerate}

\begin{itemize}
\tightlist
\item
  (R1): the reconstruction bridge assembles the internal domain,
  adjoint, spectral, covariance and locality properties into (W0)--(W3).
\item
  (R2): faithfulness of the selected theory.
\item
  (R3): injectivity of the field-label inclusion.
\item
  (R4): the unit field of degree \(0\) is nonzero, and faithfulness
  applies to it.
\item
  (R5): the labels are covered by the derivative plan.
\item
  (R6): token representatives.
\item
  (R7), (R8): the phase decomposition.
\item
  (R9), (R10): permutation symmetry and smoothness of each state.
\item
  (R11): continuity of the characters, after moving a separated word
  into a time slab.
\item
  (R12): the barycentre identity of VIII.1.j.
\item
  (R13), (R14): holomorphy and the values at Euclidean points of the
  selected continuation.
\item
  (R15): the boundary values, together with the moment identity, which
  identifies the public vacuum expectations with the internal Wightman
  distributions.
\end{itemize}

\begin{enumerate}
\def\labelenumi{(\alph{enumi})}
\setcounter{enumi}{1}
\item
  The internal separated OPE implies (O1)--(O2), after enumerating
  \(I\sqcup J\) once.
\item
  Assume \(\mathfrak g\) simple and the seven internal headline trace
  clauses. Then (T1)--(T4) hold. Clause (T1) is the centered clause. For
  (T2) and (T3), the internal amplitude is
  \(12(d_2/b_0)^2\dim\mathfrak g/\pi^4=3\dim\mathfrak g/\pi^4\), and
  \(G_\Theta\) is by definition the internal Euclidean two-point
  function. For (T4), public and internal Wick profiles coincide; then
  apply the Euclidean Wick exclusion. The real part of a finite Wick
  profile has, for every \(k\in\mathbb Z\), an expansion modulo
  \(O(r^k)\) in the span of \(r^a(\log r)^j\) with \(a\in\mathbb Z\) and
  \(j\in\mathbb N\). No function with such an expansion modulo \(O(1)\)
  satisfies \(r^8(\log 1/r)^2f(r)\to A\ne0\), since that would need the
  power \((\log r)^{-2}\). Thus (T3), together with
  \(A_{\mathfrak g}>0\) from (T2), implies (T4). \emph{Lean:}
  \href{https://github.com/yang-mills-group/OpenAI-YM-20261028/blob/3456e5e75cf8d4cd558955a5c7638fba2d68f464/HierarchyBridge.lean\#L79-L133}{\texttt{publicReconstruction}},
  \href{https://github.com/yang-mills-group/OpenAI-YM-20261028/blob/3456e5e75cf8d4cd558955a5c7638fba2d68f464/PublicView.lean\#L102-L142}{\texttt{view\_axioms}},
  \href{https://github.com/yang-mills-group/OpenAI-YM-20261028/blob/3456e5e75cf8d4cd558955a5c7638fba2d68f464/PublicView.lean\#L144-L172}{\texttt{publicMoment}},
  \href{https://github.com/yang-mills-group/OpenAI-YM-20261028/blob/3456e5e75cf8d4cd558955a5c7638fba2d68f464/HierarchyBridge.lean\#L68-L77}{\texttt{publicOPE}},
  \href{https://github.com/yang-mills-group/OpenAI-YM-20261028/blob/3456e5e75cf8d4cd558955a5c7638fba2d68f464/TraceBridge.lean\#L35-L44}{\texttt{publicTrace}}
  and
  \href{https://github.com/yang-mills-group/OpenAI-YM-20261028/blob/3456e5e75cf8d4cd558955a5c7638fba2d68f464/YM/TraceExtension/Wick/Refute.lean\#L75-L80}{\texttt{euclidean\_excludes}}.
  \emph{(proofs read, including the power-log exclusion)}
\end{enumerate}

\textbf{Lemma VIII.1.s (contact distributions).} Let \(\mathfrak F\) be
a final object with contact devices \(\mathsf M\) satisfying the full
contact resolution. For every state \(s\), every \(k\) and every
\(a\in\mathbb F^k\), there is a continuous linear functional
\(\mathcal D\) on \(C^\infty_c((\mathbb R^4)^k)\) with
\(\mathcal D(f)=\int fS^s(a;\cdot)\) whenever
\(\mathrm{supp}f\subset\mathrm{Conf}_k\), using the usual test-function
topology. No uniqueness is asserted. \emph{Proof sketch.} Apply the
distribution clause of the resolution with:

\begin{itemize}
\tightlist
\item
  the identity embedding of \((\mathbb R^4)^k\);
\item
  the full collision group, centred at \(0\);
\item
  \(B=\min(1,\mathrm{clearance}_{\mathsf M}(\sum_i d_i))\) and
  \(r=B/4\).
\end{itemize}

The contact theorem supplies \(\mathcal D\). On a separated test, the
subtraction formula reduces to the original field word. Reading that
word gives the correlation; off the support of the test both integrands
vanish. \emph{Lean:}
\href{https://github.com/yang-mills-group/OpenAI-YM-20261028/blob/3456e5e75cf8d4cd558955a5c7638fba2d68f464/ContactBridge.lean\#L14-L49}{\texttt{contactDistribution}}.

\paragraph{4.5 The headline theorems}\label{the-headline-theorems}

\textbf{Lemma VIII.1.t (detailed theorem).} This is VIII.1.2.
\emph{Proof sketch.} Apply VIII.1.o to obtain \(\mathcal C\), the
penalty and (D3). Regularity follows from VIII.1.q(a). For each
\(\sigma\), unpack the data and conclude as follows:

\begin{itemize}
\tightlist
\item
  (P1) by VIII.1.q(b);
\item
  (P2) by VIII.1.r(a), from (Core);
\item
  (P3) by VIII.1.r(b);
\item
  (P4) by VIII.1.r(c);
\item
  (D2) by VIII.1.s;
\item
  (W0)--(W3) for every phase by the reconstruction bridge.
\end{itemize}

\emph{Lean:}
\href{https://github.com/yang-mills-group/OpenAI-YM-20261028/blob/3456e5e75cf8d4cd558955a5c7638fba2d68f464/Construction.lean\#L15-L40}{\texttt{constructed}}.

\textbf{Lemma VIII.1.u (public theorem).} This is VIII.1.1. \emph{Proof
sketch.} Apply VIII.1.t with \(b=1\). Given \(p\), let
\(\sigma=(\mathrm{id},p)\), and let \(\tau\) be the refinement produced.
Take \(q:=\tau^{\rm v}\) and the hierarchy \(T(\mathfrak F,Q)\). Then
\((\sigma\circ\tau)^{\rm v}=p\circ q\), and the mesh part of \(\tau\)
plays no role because (P1) is proved along the full mesh sequence
(VIII.1.d). \emph{Lean:}
\href{https://github.com/yang-mills-group/OpenAI-YM-20261028/blob/3456e5e75cf8d4cd558955a5c7638fba2d68f464/Solution.lean\#L4-L10}{\texttt{jointExistence}}
and
\href{https://github.com/yang-mills-group/OpenAI-YM-20261028/blob/3456e5e75cf8d4cd558955a5c7638fba2d68f464/YM/Extraction/Orders.lean\#L40}{\texttt{Order.refine}}.

\textbf{Lemma VIII.1.v (\(SU(2)\)).} This is VIII.1.9. \emph{Proof
sketch.} The matrices

\[A=\begin{pmatrix}i&0\\0&-i\end{pmatrix},\qquad
B=\begin{pmatrix}0&1\\-1&0\end{pmatrix},\qquad
C=\begin{pmatrix}0&i\\i&0\end{pmatrix}\] span \(\mathfrak{su}(2)\) over
\(\mathbb R\) and satisfy \([A,B]=2C\), \([B,C]=2A\), \([C,A]=2B\). If
an ideal contains \(aA+bB+cC\ne0\), a suitable operator
\(1+\tfrac14(\mathrm{ad}\,A)^2\), or its cyclic counterpart, isolates a
nonzero multiple of one basis vector. The bracket identities then put
all three in the ideal. Apply VIII.1.u. \emph{Lean:}
\href{https://github.com/yang-mills-group/OpenAI-YM-20261028/blob/3456e5e75cf8d4cd558955a5c7638fba2d68f464/ConcreteSU2.lean\#L8-L10}{\texttt{su2}}
and
\href{https://github.com/yang-mills-group/OpenAI-YM-20261028/blob/3456e5e75cf8d4cd558955a5c7638fba2d68f464/ConcreteSU2.lean\#L15-L18}{\texttt{su2Theory}}.
\emph{(proof read, including simplicity)}

\paragraph{4.6 Endpoints filed here and not on the path of the main
theorem}\label{endpoints-filed-here-and-not-on-the-path-of-the-main-theorem}

These separate results and alternative existence theorems are not inputs
to the assembly above:

\begin{itemize}
\tightlist
\item
  the degree-four normalization
  \href{https://github.com/yang-mills-group/OpenAI-YM-20261028/blob/3456e5e75cf8d4cd558955a5c7638fba2d68f464/YM/Pert/Position/Four.lean\#L56-L63}{\texttt{normalized\_four}},
  obtained by applying the normalization result to a separately
  constructed family;
\item
  the parity results
  \href{https://github.com/yang-mills-group/OpenAI-YM-20261028/blob/3456e5e75cf8d4cd558955a5c7638fba2d68f464/YM/Pert/Parity/End.lean\#L21-L30}{\texttt{parityEndpoint}}
  and
  \href{https://github.com/yang-mills-group/OpenAI-YM-20261028/blob/3456e5e75cf8d4cd558955a5c7638fba2d68f464/YM/Pert/Parity/EvenLow.lean\#L44-L49}{\texttt{low\_read}}:
  parity is retained by the degree-four normalization, and every odd
  perturbative coefficient of the finite-volume ledger vanishes at each
  field-degree horizon \(d\le5\) for simple \(\mathfrak g\) (VII.6);
\item
  the contact endpoint
  \href{https://github.com/yang-mills-group/OpenAI-YM-20261028/blob/3456e5e75cf8d4cd558955a5c7638fba2d68f464/YM/Contacts/Solstice/Result.lean\#L78-L79}{\texttt{celestial}},
  built on
  \href{https://github.com/yang-mills-group/OpenAI-YM-20261028/blob/3456e5e75cf8d4cd558955a5c7638fba2d68f464/YM/Contacts/Layers/Conclusion.lean\#L20-L27}{\texttt{signed\_placements}}
  and
  \href{https://github.com/yang-mills-group/OpenAI-YM-20261028/blob/3456e5e75cf8d4cd558955a5c7638fba2d68f464/YM/OPE/Single.lean\#L57-L58}{\texttt{OPE.single}}
  (VII.4);
\item
  the alternative ultraviolet endpoint
  \href{https://github.com/yang-mills-group/OpenAI-YM-20261028/blob/3456e5e75cf8d4cd558955a5c7638fba2d68f464/YM/FullVolume/Finale/Program.lean\#L154-L164}{\texttt{july}}.
\end{itemize}

The volume-uniform ultraviolet bounds belong mathematically to
III.12--III.15; the supporting reconstruction and symmetry lemmas belong
to V.2--V.3 and VI.1.

\subsubsection{5. Where the difficulty
lies}\label{where-the-difficulty-lies-31}

The assembly preserves the order of choices. A family fixed before
\(\sigma\) must support the coupling limits, continuum table and all
later structures on one refinement. In particular, the fixed-volume
ultraviolet convergence must concern this same family. Uniqueness of row
limits then removes the mesh subsequence without replacing the final
object or its point-limit data.

The momentum-space Wick exclusion requires a further argument: every
compactly supported relative-momentum test must be realizable as
\(\overline{\mathcal F_+u}\mathcal F_+w\). Fourier surjectivity and an
auxiliary transform of modulus one on the test support give that
realization. Exponential integrability then permits passage to Euclidean
Laplace transforms.

The public theorem retains less than the detailed construction. Its
Wilson limit concerns the mixed state, while its trace test concerns the
selected phase. The Euclidean Wick exclusion follows from (T2)--(T3);
the detailed theorem also gives momentum-space and Maxwell
relative-locality countertests. The low-energy and heat bounds are
polynomial and imply no mass gap. Coupling limits, all-shapes OPE and
sharp Gram allocations occur in (Goal) but are omitted from VIII.1.3.
Neither independence of the volume subsequence nor agreement between
different penalties or run choices is asserted.

\subsubsection{6. Notation}\label{notation-31}

{\def\LTcaptype{none} % do not increment counter
\begin{longtable}[]{@{}
  >{\raggedright\arraybackslash}p{(\linewidth - 4\tabcolsep) * \real{0.3333}}
  >{\raggedright\arraybackslash}p{(\linewidth - 4\tabcolsep) * \real{0.3333}}
  >{\raggedright\arraybackslash}p{(\linewidth - 4\tabcolsep) * \real{0.3333}}@{}}
\toprule\noalign{}
\begin{minipage}[b]{\linewidth}\raggedright
object
\end{minipage} & \begin{minipage}[b]{\linewidth}\raggedright
code name
\end{minipage} & \begin{minipage}[b]{\linewidth}\raggedright
proposed notation
\end{minipage} \\
\midrule\noalign{}
\endhead
\bottomrule\noalign{}
\endlastfoot
gauge group &
\href{https://github.com/yang-mills-group/OpenAI-YM-20261028/blob/3456e5e75cf8d4cd558955a5c7638fba2d68f464/PublicTheory.lean\#L25-L33}{\texttt{GaugeGroup}},
\href{https://github.com/yang-mills-group/OpenAI-YM-20261028/blob/3456e5e75cf8d4cd558955a5c7638fba2d68f464/PublicTheory.lean\#L39-L45}{\texttt{IsSimple}}
& \((m,G,\mathfrak g)\), \(\mathfrak g\) simple \\
public Wilson family &
\href{https://github.com/yang-mills-group/OpenAI-YM-20261028/blob/3456e5e75cf8d4cd558955a5c7638fba2d68f464/WilsonBridge.lean\#L20-L30}{\texttt{wilsonFamily}}
&
\(\mathcal W(\mathcal C)=(L_p,D_N,\beta_{N,p},h,\mathcal T,I_{N,p})\) \\
regularity &
\href{https://github.com/yang-mills-group/OpenAI-YM-20261028/blob/3456e5e75cf8d4cd558955a5c7638fba2d68f464/ChallengeSpec.lean\#L12-L22}{\texttt{WilsonRegularity}}
& (Reg1)--(Reg4) \\
public hierarchy &
\href{https://github.com/yang-mills-group/OpenAI-YM-20261028/blob/3456e5e75cf8d4cd558955a5c7638fba2d68f464/HierarchyBridge.lean\#L27-L42}{\texttt{hierarchy}}
& \(T(\mathfrak F,Q)\) \\
extraction order, refinement &
\href{https://github.com/yang-mills-group/OpenAI-YM-20261028/blob/3456e5e75cf8d4cd558955a5c7638fba2d68f464/YM/Extraction/Orders.lean\#L34-L36}{\texttt{Extraction.Order}},
\href{https://github.com/yang-mills-group/OpenAI-YM-20261028/blob/3456e5e75cf8d4cd558955a5c7638fba2d68f464/YM/Extraction/Orders.lean\#L40}{\texttt{refine}}
& \(\sigma=(\sigma^{\rm m},\sigma^{\rm v})\), \(\sigma\circ\tau\) \\
ordered limit &
\href{https://github.com/yang-mills-group/OpenAI-YM-20261028/blob/3456e5e75cf8d4cd558955a5c7638fba2d68f464/YM/Extraction/Limits.lean\#L32-L35}{\texttt{Ordered}}
& \(\mathrm{Ord}_\sigma(\mathcal C,S)\) \\
internal family, construction predicate & \texttt{Family},
\texttt{Makes} & \(\mathcal C\), (D3) \\
run, funding conditions &
\href{https://github.com/yang-mills-group/OpenAI-YM-20261028/blob/3456e5e75cf8d4cd558955a5c7638fba2d68f464/YM/FullVolume/Cross.lean\#L30-L48}{\texttt{Run}},
\href{https://github.com/yang-mills-group/OpenAI-YM-20261028/blob/3456e5e75cf8d4cd558955a5c7638fba2d68f464/YM/Independent/ConstructionRun.lean\#L30-L31}{\texttt{fundedRun}},
\href{https://github.com/yang-mills-group/OpenAI-YM-20261028/blob/3456e5e75cf8d4cd558955a5c7638fba2d68f464/YM/R2I/Campaign/Accepted.lean\#L60-L73}{\texttt{Funded}}
& \(\mathsf R_{G,b}\), (Fd) \\
ultraviolet endpoint &
\href{https://github.com/yang-mills-group/OpenAI-YM-20261028/blob/3456e5e75cf8d4cd558955a5c7638fba2d68f464/YM/FullVolume/Finale/Program.lean\#L133-L149}{\texttt{Jewel}}
& (UV\(_\infty\)) \\
positional convergence &
\href{https://github.com/yang-mills-group/OpenAI-YM-20261028/blob/3456e5e75cf8d4cd558955a5c7638fba2d68f464/YM/FullVolume/Profiles.lean\#L25-L36}{\texttt{PositionsUV}}
& (D5), \(f_{p,v}\) \\
full assembly package &
\href{https://github.com/yang-mills-group/OpenAI-YM-20261028/blob/3456e5e75cf8d4cd558955a5c7638fba2d68f464/YM/R2I/Campaign/Accepted.lean\#L130-L148}{\texttt{Broad}}
& (Goal) \\
table & \texttt{Table} & \(S\), read as \(S^s\) \\
final object &
\href{https://github.com/yang-mills-group/OpenAI-YM-20261028/blob/3456e5e75cf8d4cd558955a5c7638fba2d68f464/YM/Spec/Closed.lean\#L64-L73}{\texttt{Spec.Final}}
& \(\mathfrak F\) on \((\mathcal C,\sigma,S)\) \\
reconstruction in every phase &
\href{https://github.com/yang-mills-group/OpenAI-YM-20261028/blob/3456e5e75cf8d4cd558955a5c7638fba2d68f464/YM/EveryPhase/Pack.lean\#L117-L132}{\texttt{Almanac}},
\href{https://github.com/yang-mills-group/OpenAI-YM-20261028/blob/3456e5e75cf8d4cd558955a5c7638fba2d68f464/YM/EveryPhase/Extension.lean\#L51-L71}{\texttt{Completed}}
& \(\mathfrak L=(\mathfrak W_\xi)_{\xi\in X}\) \\
OPE reader &
\href{https://github.com/yang-mills-group/OpenAI-YM-20261028/blob/3456e5e75cf8d4cd558955a5c7638fba2d68f464/YM/OPE/Orbits/Cuts.lean\#L17-L23}{\texttt{Orbits.Reader}}
& \(Q\); \((\mathrm{Ch}_j,O_{j,\kappa},c_{j,\kappa})\) \\
contact devices &
\href{https://github.com/yang-mills-group/OpenAI-YM-20261028/blob/3456e5e75cf8d4cd558955a5c7638fba2d68f464/YM/Contacts/Tomography.lean\#L17-L26}{\texttt{Contacts.Devices}}
& \(\mathsf M\) \\
scaling sequence &
\href{https://github.com/yang-mills-group/OpenAI-YM-20261028/blob/3456e5e75cf8d4cd558955a5c7638fba2d68f464/YM/Climb/Cells.lean\#L24-L31}{\texttt{Climb.Scaling}}
& \(\Sigma=(r_k,\xi_k,U_k)\) \\
combined continuum properties &
\href{https://github.com/yang-mills-group/OpenAI-YM-20261028/blob/3456e5e75cf8d4cd558955a5c7638fba2d68f464/YM/Independent/Core.lean\#L18-L34}{\texttt{Independent.Core}}
& (Core) \\
first comparison, degree-four normalization &
\href{https://github.com/yang-mills-group/OpenAI-YM-20261028/blob/3456e5e75cf8d4cd558955a5c7638fba2d68f464/YM/Pert/FirstTarget.lean\#L26-L39}{\texttt{FirstExtension}},
\href{https://github.com/yang-mills-group/OpenAI-YM-20261028/blob/3456e5e75cf8d4cd558955a5c7638fba2d68f464/YM/Pert/Position/Four.lean\#L45-L47}{\texttt{FourMeters}}
& (D14), (D15) \\
centered trace, its two-point function &
\href{https://github.com/yang-mills-group/OpenAI-YM-20261028/blob/3456e5e75cf8d4cd558955a5c7638fba2d68f464/YM/TraceExtension/WilsonTrace.lean\#L48}{\texttt{TraceExtension.center}},
\href{https://github.com/yang-mills-group/OpenAI-YM-20261028/blob/3456e5e75cf8d4cd558955a5c7638fba2d68f464/YM/TraceExtension/WilsonTrace.lean\#L53}{\texttt{euclidean}}
& \(\Theta\), \(G_\Theta\) \\
trace amplitude &
\href{https://github.com/yang-mills-group/OpenAI-YM-20261028/blob/3456e5e75cf8d4cd558955a5c7638fba2d68f464/YM/TraceExtension/Verse/Day.lean\#L41}{\texttt{Trace.Verse.award}}
& \(A_{\mathfrak g}=3\dim\mathfrak g/\pi^4\) \\
one-loop constants &
\href{https://github.com/yang-mills-group/OpenAI-YM-20261028/blob/3456e5e75cf8d4cd558955a5c7638fba2d68f464/YM/Pert/Asymptotics/RawAxis.lean\#L13}{\texttt{Pert.Asymptotics.b0}},
\texttt{Trace.Over.d2} & \(b_0=11C_{\mathfrak g}/(48\pi^2)\),
\(b_*=2b_0\), \(d_2=-b_0/2\) \\
reference coupling & \texttt{Calendar.hatBeta} &
\(\hat\beta_m=\bar\beta+2b_0\log(1/a_m)\) \\
penalty of the block-spin kernel & \texttt{KernelChoices.penalty} &
\(b\) \\
\end{longtable}
}

\subsubsection{7. What remains
unexplained}\label{what-remains-unexplained-31}

\begin{enumerate}
\def\labelenumi{\arabic{enumi}.}
\tightlist
\item
  \textbf{Numerical run conditions.} The thirteen funding conditions in
  §2 are specified by their roles. The exact inequalities for the frame
  and polynomial schedules, delays, pacing, thermal estimates and weight
  tableau must be supplied from III.15. The numerical proofs are inputs.
\item
  \textbf{Renormalization comparison.} VIII.1.c uses the fixed-volume
  Cauchy estimate from III.12. The underlying renormalization-step
  estimate is an input; the dense-grid, Taylor and smearing arguments
  are given in VIII.1.c.
\item
  \textbf{Normalized trace schemes.} VIII.1.h retains the existence and
  universal quantifiers of the stress-comparison theorem. The alignment,
  leading-derivative and scalar normalization conditions require the
  definitions and notation of VII.3. No comparison for an arbitrary
  renormalization prescription follows.
\item
  \textbf{Upstream reconstruction.} The phase selection and selected
  reconstruction are produced inside the final-object construction. The
  every-phase theorem is an input from V.3; its low-energy, decay,
  irreducibility and boundary arguments remain inputs from V.3.
\end{enumerate}

\end{document}
